% ═══════════════════════════════════════════════════════════════════════
\section{Trouver un SIF : le théorème du collage}
% ═══════════════════════════════════════════════════════════════════════

\begin{frame}{Théorème 3.4 --- Théorème du collage (Barnsley)}
  \small Le problème inverse : à partir d'une image (feuille, arbre,
  texture), trouver un SIF dont l'attracteur s'en approche.
  \begin{Enonce}{Théorème 3.4 --- Théorème du collage}
    Soit $L$ un compact non vide du plan, et un SIF de rapport $k$,
    d'attracteur $A$ et d'opérateur $W$. Alors
    \[
      \dH(L,A)\leq\frac{1}{1-k}\,\dH\big(L,W(L)\big).
    \]
  \end{Enonce}
  \begin{AuTableau}
    \small Point (iii) du théorème 3.3 avec $B=L$ et $n=0$.
  \end{AuTableau}
  \small Recouvrir $L$ par des photocopies réduites et déformées
  d'elle-même (un collage) : si le collage est fidèle, l'attracteur
  ressemble à $L$. L'estimation se dégrade si $k$ est proche de $1$.
\end{frame}

\begin{frame}{Trouver chaque $w_i$ : retour au bulldozer}
  \small
  \begin{itemize}
    \item Par le théorème des trois points (chapitre 1), une application
      affine est déterminée par les images de trois points non alignés.
    \item On repère trois points caractéristiques de $L$ (base de la
      tige, pointe, extrémité d'une feuille) et leurs positions sur la
      pièce $L_i$ voulue.
    \item On résout les deux systèmes $3\times3$ du chapitre 1 : on
      obtient $w_i$.
  \end{itemize}
\end{frame}

\begin{frame}{Exemple 3.7 --- Un arbre}
  \begin{columns}[c]
    \begin{column}{0.5\textwidth}
      \footnotesize
      \begin{itemize}
        \item Le tronc :
          $w_1\col{x}{y} = \mat{\num{0.03}}{0}{0}{\num{0.42}}\col{x}{y}$,
          l'arbre réduit à un mince segment.
        \item Deux branches : similitudes de rapport $\num{0.58}$,
          tournées de $\pm38^\circ$, posées au sommet du tronc :
          $w_{2,3}\col{x}{y} = \num{0.58}\,r_{\pm38^\circ}\col{x}{y}
          + \col{0}{\num{0.42}}$.
        \item La cime :
          $w_4\col{x}{y} = \num{0.42}\col{x}{y}+\col{0}{\num{0.6}}$.
      \end{itemize}
      Rapport de contraction : $\num{0.58}$. Comme $w_1$ n'est pas une
      similitude, l'attracteur n'est pas autosimilaire au sens de la
      définition 3.1.
    \end{column}
    \begin{column}{0.48\textwidth}
      \FigBox[0.6\textheight]{%
        \begin{tikzpicture}[scale=3.1]
          % Arbre fractal
% Fichier généré par generer.py : ne pas modifier à la main.
\fill[nuage]
  (-0.154,0.717)rectangle++(0.008,0.008) (-0.33,0.692)rectangle++(0.008,0.008) (-0.141,0.888)rectangle++(0.008,0.008) (-0.008,0.371)rectangle++(0.008,0.008) (-0.004,0.153)rectangle++(0.008,0.008)
  (-0.004,0.662)rectangle++(0.008,0.008) (0.234,0.72)rectangle++(0.008,0.008) (0.003,0.3)rectangle++(0.008,0.008) (-0.001,0.724)rectangle++(0.008,0.008) (-0.003,0.902)rectangle++(0.008,0.008)
  (-0.327,0.83)rectangle++(0.008,0.008) (-0.449,0.682)rectangle++(0.008,0.008) (-0.453,0.57)rectangle++(0.008,0.008) (-0.414,0.518)rectangle++(0.008,0.008) (-0.005,0.801)rectangle++(0.008,0.008)
  (-0.292,0.784)rectangle++(0.008,0.008) (0.146,0.879)rectangle++(0.008,0.008) (0,0.367)rectangle++(0.008,0.008) (-0.134,0.587)rectangle++(0.008,0.008) (-0.059,0.844)rectangle++(0.008,0.008)
  (-0.332,0.784)rectangle++(0.008,0.008) (-0.435,0.659)rectangle++(0.008,0.008) (0.036,0.873)rectangle++(0.008,0.008) (0.327,0.803)rectangle++(0.008,0.008) (-0.141,0.903)rectangle++(0.008,0.008)
  (-0.008,0.377)rectangle++(0.008,0.008) (0.13,0.592)rectangle++(0.008,0.008) (0.052,0.846)rectangle++(0.008,0.008) (0.02,0.953)rectangle++(0.008,0.008) (0.349,0.845)rectangle++(0.008,0.008)
  (-0.146,0.93)rectangle++(0.008,0.008) (-0.402,0.792)rectangle++(0.008,0.008) (0.098,0.922)rectangle++(0.008,0.008) (0.373,0.803)rectangle++(0.008,0.008) (-0.12,0.92)rectangle++(0.008,0.008)
  (-0.387,0.797)rectangle++(0.008,0.008) (-0.465,0.645)rectangle++(0.008,0.008) (0.017,0.877)rectangle++(0.008,0.008) (0.005,0.966)rectangle++(0.008,0.008) (0.347,0.856)rectangle++(0.008,0.008)
  (0.007,0.357)rectangle++(0.008,0.008) (0,0.748)rectangle++(0.008,0.008) (0.266,0.758)rectangle++(0.008,0.008) (0.392,0.668)rectangle++(0.008,0.008) (-0.063,0.864)rectangle++(0.008,0.008)
  (-0.341,0.792)rectangle++(0.008,0.008) (-0.442,0.659)rectangle++(0.008,0.008) (-0.188,0.875)rectangle++(0.008,0.008) (-0.402,0.752)rectangle++(0.008,0.008) (0.084,0.904)rectangle++(0.008,0.008)
  (-0.288,0.862)rectangle++(0.008,0.008) (0.176,0.913)rectangle++(0.008,0.008) (0.001,0.381)rectangle++(0.008,0.008) (-0.139,0.594)rectangle++(0.008,0.008) (-0.279,0.641)rectangle++(0.008,0.008)
  (0.101,0.809)rectangle++(0.008,0.008) (-0.001,0.338)rectangle++(0.008,0.008) (0.119,0.571)rectangle++(0.008,0.008) (-0.153,0.723)rectangle++(0.008,0.008) (-0.332,0.695)rectangle++(0.008,0.008)
  (0.096,0.852)rectangle++(0.008,0.008) (0.038,0.956)rectangle++(0.008,0.008) (0.358,0.84)rectangle++(0.008,0.008) (0.148,0.95)rectangle++(0.008,0.008) (-0.275,0.906)rectangle++(0.008,0.008)
  (-0.118,0.978)rectangle++(0.008,0.008) (-0.407,0.824)rectangle++(0.008,0.008) (-0.484,0.651)rectangle++(0.008,0.008) (0.01,0.887)rectangle++(0.008,0.008) (0.002,0.97)rectangle++(0.008,0.008)
  (-0.004,0.405)rectangle++(0.008,0.008) (-0.15,0.603)rectangle++(0.008,0.008) (-0.065,0.851)rectangle++(0.008,0.008) (-0.337,0.785)rectangle++(0.008,0.008) (-0.438,0.658)rectangle++(0.008,0.008)
  (-0.017,0.274)rectangle++(0.008,0.008) (-0.004,0.113)rectangle++(0.008,0.008) (0.037,0.469)rectangle++(0.008,0.008) (0.013,0.795)rectangle++(0.008,0.008) (0.003,0.932)rectangle++(0.008,0.008)
  (0.333,0.841)rectangle++(0.008,0.008) (0.006,0.351)rectangle++(0.008,0.008) (-0.126,0.582)rectangle++(0.008,0.008) (-0.055,0.842)rectangle++(0.008,0.008) (-0.026,0.951)rectangle++(0.008,0.008)
  (0.327,0.86)rectangle++(0.008,0.008) (-0.161,0.929)rectangle++(0.008,0.008) (0.257,0.899)rectangle++(0.008,0.008) (0.438,0.735)rectangle++(0.008,0.008) (0.462,0.596)rectangle++(0.008,0.008)
  (-0.005,0.857)rectangle++(0.008,0.008) (-0.005,0.957)rectangle++(0.008,0.008) (-0.004,1)rectangle++(0.008,0.008) (-0.363,0.875)rectangle++(0.008,0.008) (-0.482,0.69)rectangle++(0.008,0.008)
  (-0.205,0.887)rectangle++(0.008,0.008) (-0.414,0.752)rectangle++(0.008,0.008) (-0.176,0.913)rectangle++(0.008,0.008) (0.245,0.897)rectangle++(0.008,0.008) (0.431,0.739)rectangle++(0.008,0.008)
  (0.46,0.6)rectangle++(0.008,0.008) (-0.008,0.858)rectangle++(0.008,0.008) (-0.313,0.809)rectangle++(0.008,0.008) (-0.436,0.677)rectangle++(0.008,0.008) (0.042,0.881)rectangle++(0.008,0.008)
  (0.333,0.804)rectangle++(0.008,0.008) (-0.139,0.906)rectangle++(0.008,0.008) (-0.061,0.978)rectangle++(0.008,0.008) (-0.381,0.845)rectangle++(0.008,0.008) (-0.479,0.669)rectangle++(0.008,0.008)
  (0.019,0.893)rectangle++(0.008,0.008) (0.006,0.973)rectangle++(0.008,0.008) (0.349,0.859)rectangle++(0.008,0.008) (0.466,0.684)rectangle++(0.008,0.008) (0.193,0.885)rectangle++(0.008,0.008)
  (0.079,0.969)rectangle++(0.008,0.008) (0.031,1.005)rectangle++(0.008,0.008) (-0.348,0.89)rectangle++(0.008,0.008) (0.158,0.947)rectangle++(0.008,0.008) (0.064,0.996)rectangle++(0.008,0.008)
  (0.025,1.016)rectangle++(0.008,0.008) (0.373,0.872)rectangle++(0.008,0.008) (-0.144,0.951)rectangle++(0.008,0.008) (-0.409,0.802)rectangle++(0.008,0.008) (0.099,0.929)rectangle++(0.008,0.008)
  (0.376,0.806)rectangle++(0.008,0.008) (0.007,0.336)rectangle++(0.008,0.008) (0.123,0.567)rectangle++(0.008,0.008) (0.049,0.836)rectangle++(0.008,0.008) (0.018,0.949)rectangle++(0.008,0.008)
  (0.005,0.996)rectangle++(0.008,0.008) (0,1.016)rectangle++(0.008,0.008) (0.362,0.881)rectangle++(0.008,0.008) (0.15,0.968)rectangle++(0.008,0.008) (-0.281,0.915)rectangle++(0.008,0.008)
  (-0.012,0.382)rectangle++(0.008,0.008) (-0.146,0.589)rectangle++(0.008,0.008) (-0.281,0.637)rectangle++(0.008,0.008) (-0.12,0.865)rectangle++(0.008,0.008) (0.253,0.855)rectangle++(0.008,0.008)
  (0.42,0.717)rectangle++(0.008,0.008) (0.447,0.594)rectangle++(0.008,0.008) (-0.011,0.85)rectangle++(0.008,0.008) (0.298,0.809)rectangle++(0.008,0.008) (0.424,0.68)rectangle++(0.008,0.008)
  (-0.052,0.881)rectangle++(0.008,0.008) (-0.342,0.803)rectangle++(0.008,0.008) (-0.447,0.664)rectangle++(0.008,0.008) (-0.445,0.563)rectangle++(0.008,0.008) (-0.003,0.833)rectangle++(0.008,0.008)
  (-0.302,0.799)rectangle++(0.008,0.008) (-0.013,0.333)rectangle++(0.008,0.008) (-0.128,0.567)rectangle++(0.008,0.008) (0.143,0.721)rectangle++(0.008,0.008) (0.322,0.695)rectangle++(0.008,0.008)
  (0.395,0.619)rectangle++(0.008,0.008) (0.163,0.858)rectangle++(0.008,0.008) (-0.235,0.87)rectangle++(0.008,0.008) (-0.011,0.363)rectangle++(0.008,0.008) (0.124,0.586)rectangle++(0.008,0.008)
  (0.265,0.64)rectangle++(0.008,0.008) (-0.111,0.806)rectangle++(0.008,0.008) (0.237,0.825)rectangle++(0.008,0.008) (-0.19,0.881)rectangle++(0.008,0.008) (-0.405,0.754)rectangle++(0.008,0.008)
  (-0.458,0.619)rectangle++(0.008,0.008) (-0.195,0.858)rectangle++(0.008,0.008) (-0.01,0.358)rectangle++(0.008,0.008) (-0.136,0.579)rectangle++(0.008,0.008) (-0.273,0.636)rectangle++(0.008,0.008)
  (-0.355,0.612)rectangle++(0.008,0.008) (0.056,0.823)rectangle++(0.008,0.008) (-0.272,0.815)rectangle++(0.008,0.008) (-0.117,0.94)rectangle++(0.008,0.008) (0.282,0.888)rectangle++(0.008,0.008)
  (0.445,0.722)rectangle++(0.008,0.008) (-0.058,0.908)rectangle++(0.008,0.008) (0.297,0.852)rectangle++(0.008,0.008) (0.439,0.7)rectangle++(0.008,0.008) (0.182,0.892)rectangle++(0.008,0.008)
  (-0.239,0.892)rectangle++(0.008,0.008) (0.209,0.909)rectangle++(0.008,0.008) (0.085,0.98)rectangle++(0.008,0.008) (0.034,1.009)rectangle++(0.008,0.008) (-0.349,0.892)rectangle++(0.008,0.008)
  (0.159,0.949)rectangle++(0.008,0.008) (0.064,0.996)rectangle++(0.008,0.008) (-0.33,0.898)rectangle++(0.008,0.008) (0.169,0.944)rectangle++(0.008,0.008) (0.414,0.788)rectangle++(0.008,0.008)
  (0.171,0.929)rectangle++(0.008,0.008) (0.409,0.78)rectangle++(0.008,0.008) (0.465,0.627)rectangle++(0.008,0.008) (-0.015,0.872)rectangle++(0.008,0.008) (0.304,0.82)rectangle++(0.008,0.008)
  (-0.158,0.902)rectangle++(0.008,0.008) (0.249,0.885)rectangle++(0.008,0.008) (-0.206,0.913)rectangle++(0.008,0.008) (-0.424,0.763)rectangle++(0.008,0.008) (-0.47,0.617)rectangle++(0.008,0.008)
  (-0.438,0.533)rectangle++(0.008,0.008) (-0.011,0.817)rectangle++(0.008,0.008) (-0.3,0.789)rectangle++(0.008,0.008) (-0.013,0.329)rectangle++(0.008,0.008) (-0.008,0.736)rectangle++(0.008,0.008)
  (-0.27,0.753)rectangle++(0.008,0.008) (-0.116,0.914)rectangle++(0.008,0.008) (0.273,0.875)rectangle++(0.008,0.008) (-0.192,0.917)rectangle++(0.008,0.008) (-0.083,0.983)rectangle++(0.008,0.008)
  (-0.037,1.01)rectangle++(0.008,0.008) (-0.381,0.868)rectangle++(0.008,0.008) (-0.488,0.68)rectangle++(0.008,0.008) (0.019,0.901)rectangle++(0.008,0.008) (0.006,0.976)rectangle++(0.008,0.008)
  (0.35,0.861)rectangle++(0.008,0.008) (0.145,0.959)rectangle++(0.008,0.008) (-0.28,0.909)rectangle++(0.008,0.008) (0.196,0.932)rectangle++(0.008,0.008) (-0.247,0.915)rectangle++(0.008,0.008)
  (-0.443,0.749)rectangle++(0.008,0.008) (-0.474,0.604)rectangle++(0.008,0.008) (-0.436,0.526)rectangle++(0.008,0.008) (-0.391,0.504)rectangle++(0.008,0.008) (-0.166,0.809)rectangle++(0.008,0.008)
  (-0.369,0.73)rectangle++(0.008,0.008) (-0.433,0.621)rectangle++(0.008,0.008) (-0.423,0.549)rectangle++(0.008,0.008) (-0.18,0.828)rectangle++(0.008,0.008) (0.213,0.859)rectangle++(0.008,0.008)
  (0.003,0.359)rectangle++(0.008,0.008) (-0.001,0.748)rectangle++(0.008,0.008) (-0.003,0.912)rectangle++(0.008,0.008) (-0.004,0.981)rectangle++(0.008,0.008) (-0.355,0.866)rectangle++(0.008,0.008)
  (0.146,0.939)rectangle++(0.008,0.008) (-0.272,0.901)rectangle++(0.008,0.008) (-0.45,0.734)rectangle++(0.008,0.008) (-0.191,0.906)rectangle++(0.008,0.008) (-0.083,0.978)rectangle++(0.008,0.008)
  (-0.391,0.837)rectangle++(0.008,0.008) (-0.166,0.949)rectangle++(0.008,0.008) (-0.009,0.396)rectangle++(0.008,0.008) (-0.149,0.597)rectangle++(0.008,0.008) (-0.008,0.249)rectangle++(0.008,0.008)
  (-0.096,0.53)rectangle++(0.008,0.008) (-0.237,0.627)rectangle++(0.008,0.008) (0.115,0.788)rectangle++(0.008,0.008) (-0.232,0.82)rectangle++(0.008,0.008) (-0.403,0.711)rectangle++(0.008,0.008)
  (-0.442,0.601)rectangle++(0.008,0.008) (-0.42,0.536)rectangle++(0.008,0.008) (-0.179,0.823)rectangle++(0.008,0.008) (0.211,0.856)rectangle++(0.008,0.008) (-0.213,0.886)rectangle++(0.008,0.008)
  (0.218,0.897)rectangle++(0.008,0.008) (0.42,0.749)rectangle++(0.008,0.008) (0.458,0.609)rectangle++(0.008,0.008) (0.426,0.531)rectangle++(0.008,0.008) (0.384,0.507)rectangle++(0.008,0.008)
  (0.159,0.811)rectangle++(0.008,0.008) (0.361,0.73)rectangle++(0.008,0.008) (0.425,0.621)rectangle++(0.008,0.008) (0.415,0.548)rectangle++(0.008,0.008) (0.385,0.519)rectangle++(0.008,0.008)
  (-0.013,0.794)rectangle++(0.008,0.008) (-0.293,0.777)rectangle++(0.008,0.008) (0.143,0.876)rectangle++(0.008,0.008) (0.378,0.766)rectangle++(0.008,0.008) (-0.105,0.904)rectangle++(0.008,0.008)
  (-0.046,0.977)rectangle++(0.008,0.008) (0.327,0.88)rectangle++(0.008,0.008) (0.135,0.967)rectangle++(0.008,0.008) (0.406,0.81)rectangle++(0.008,0.008) (0.168,0.938)rectangle++(0.008,0.008)
  (0.411,0.785)rectangle++(0.008,0.008) (-0.096,0.925)rectangle++(0.008,0.008) (-0.378,0.808)rectangle++(0.008,0.008) (0.115,0.92)rectangle++(0.008,0.008) (0.38,0.796)rectangle++(0.008,0.008)
  (0.157,0.932)rectangle++(0.008,0.008) (-0.264,0.901)rectangle++(0.008,0.008) (-0.012,0.376)rectangle++(0.008,0.008) (-0.143,0.587)rectangle++(0.008,0.008) (-0.008,0.244)rectangle++(0.008,0.008)
  (-0.004,0.1)rectangle++(0.008,0.008) (-0.004,0.04)rectangle++(0.008,0.008) (-0.02,0.436)rectangle++(0.008,0.008) (-0.168,0.612)rectangle++(0.008,0.008) (0.141,0.756)rectangle++(0.008,0.008)
  (0.334,0.712)rectangle++(0.008,0.008) (0.138,0.897)rectangle++(0.008,0.008) (0.056,0.974)rectangle++(0.008,0.008) (0.373,0.842)rectangle++(0.008,0.008) (0.154,0.951)rectangle++(0.008,0.008)
  (0.001,0.397)rectangle++(0.008,0.008) (-0.145,0.601)rectangle++(0.008,0.008) (-0.063,0.85)rectangle++(0.008,0.008) (0.274,0.828)rectangle++(0.008,0.008) (-0.174,0.895)rectangle++(0.008,0.008)
  (0.239,0.888)rectangle++(0.008,0.008) (-0.211,0.91)rectangle++(0.008,0.008) (0.228,0.908)rectangle++(0.008,0.008) (-0.224,0.916)rectangle++(0.008,0.008) (-0.433,0.758)rectangle++(0.008,0.008)
  (-0.017,0.316)rectangle++(0.008,0.008) (0.104,0.567)rectangle++(0.008,0.008) (0.042,0.836)rectangle++(0.008,0.008) (0.015,0.949)rectangle++(0.008,0.008) (0.004,0.996)rectangle++(0.008,0.008)
  (-0.001,1.016)rectangle++(0.008,0.008) (-0.367,0.883)rectangle++(0.008,0.008) (-0.487,0.692)rectangle++(0.008,0.008) (-0.207,0.888)rectangle++(0.008,0.008) (0.222,0.896)rectangle++(0.008,0.008)
  (0.421,0.747)rectangle++(0.008,0.008) (0.009,0.311)rectangle++(0.008,0.008) (-0.111,0.565)rectangle++(0.008,0.008) (0.15,0.714)rectangle++(0.008,0.008) (0.323,0.689)rectangle++(0.008,0.008)
  (0.133,0.887)rectangle++(0.008,0.008) (0.377,0.774)rectangle++(0.008,0.008) (-0.108,0.908)rectangle++(0.008,0.008) (-0.377,0.796)rectangle++(0.008,0.008) (-0.161,0.932)rectangle++(0.008,0.008)
  (-0.07,0.989)rectangle++(0.008,0.008) (0.321,0.893)rectangle++(0.008,0.008) (-0.176,0.942)rectangle++(0.008,0.008) (0.255,0.91)rectangle++(0.008,0.008) (0.441,0.741)rectangle++(0.008,0.008)
  (-0.067,0.915)rectangle++(0.008,0.008) (-0.361,0.814)rectangle++(0.008,0.008) (-0.459,0.662)rectangle++(0.008,0.008) (0.026,0.883)rectangle++(0.008,0.008) (-0.307,0.832)rectangle++(0.008,0.008)
  (0.156,0.906)rectangle++(0.008,0.008) (0.063,0.978)rectangle++(0.008,0.008) (0.378,0.841)rectangle++(0.008,0.008) (-0.131,0.938)rectangle++(0.008,0.008) (-0.399,0.801)rectangle++(0.008,0.008)
  (-0.472,0.643)rectangle++(0.008,0.008) (0.013,0.879)rectangle++(0.008,0.008) (0.319,0.813)rectangle++(0.008,0.008) (-0.148,0.905)rectangle++(0.008,0.008) (0.255,0.883)rectangle++(0.008,0.008)
  (0.431,0.729)rectangle++(0.008,0.008) (0.457,0.596)rectangle++(0.008,0.008) (-0.008,0.855)rectangle++(0.008,0.008) (-0.312,0.807)rectangle++(0.008,0.008) (-0.435,0.677)rectangle++(0.008,0.008)
  (-0.185,0.882)rectangle++(0.008,0.008) (-0.403,0.756)rectangle++(0.008,0.008) (0.085,0.906)rectangle++(0.008,0.008) (0.362,0.8)rectangle++(0.008,0.008) (-0.124,0.914)rectangle++(0.008,0.008)
  (-0.387,0.793)rectangle++(0.008,0.008) (-0.463,0.644)rectangle++(0.008,0.008) (-0.445,0.548)rectangle++(0.008,0.008) (-0.403,0.511)rectangle++(0.008,0.008) (-0.37,0.509)rectangle++(0.008,0.008)
  (-0.158,0.811)rectangle++(0.008,0.008) (-0.365,0.734)rectangle++(0.008,0.008) (-0.433,0.624)rectangle++(0.008,0.008) (0.024,0.856)rectangle++(0.008,0.008) (-0.298,0.819)rectangle++(0.008,0.008)
  (0.156,0.897)rectangle++(0.008,0.008) (0.391,0.771)rectangle++(0.008,0.008) (0.008,0.321)rectangle++(0.008,0.008) (0.118,0.561)rectangle++(0.008,0.008) (-0.15,0.717)rectangle++(0.008,0.008)
  (-0.328,0.694)rectangle++(0.008,0.008) (-0.14,0.889)rectangle++(0.008,0.008) (-0.061,0.971)rectangle++(0.008,0.008) (0.318,0.882)rectangle++(0.008,0.008) (-0.173,0.936)rectangle++(0.008,0.008)
  (0.254,0.906)rectangle++(0.008,0.008) (-0.211,0.924)rectangle++(0.008,0.008) (-0.091,0.986)rectangle++(0.008,0.008) (-0.397,0.837)rectangle++(0.008,0.008) (-0.484,0.66)rectangle++(0.008,0.008)
  (0.014,0.891)rectangle++(0.008,0.008) (0.324,0.819)rectangle++(0.008,0.008) (0.44,0.675)rectangle++(0.008,0.008) (0.441,0.568)rectangle++(0.008,0.008) (0.183,0.836)rectangle++(0.008,0.008)
  (0.381,0.733)rectangle++(0.008,0.008) (0.158,0.906)rectangle++(0.008,0.008) (-0.255,0.89)rectangle++(0.008,0.008) (0.2,0.914)rectangle++(0.008,0.008) (-0.238,0.909)rectangle++(0.008,0.008)
  (0.215,0.917)rectangle++(0.008,0.008) (0.425,0.759)rectangle++(0.008,0.008) (-0.08,0.918)rectangle++(0.008,0.008) (0.29,0.865)rectangle++(0.008,0.008) (0.441,0.708)rectangle++(0.008,0.008)
  (0.453,0.583)rectangle++(0.008,0.008) (0.415,0.521)rectangle++(0.008,0.008) (0.375,0.506)rectangle++(0.008,0.008) (0.351,0.514)rectangle++(0.008,0.008) (-0.027,0.78)rectangle++(0.008,0.008)
  (-0.294,0.766)rectangle++(0.008,0.008) (-0.013,0.319)rectangle++(0.008,0.008) (-0.008,0.732)rectangle++(0.008,0.008) (-0.006,0.905)rectangle++(0.008,0.008) (-0.329,0.831)rectangle++(0.008,0.008)
  (-0.141,0.947)rectangle++(0.008,0.008) (-0.061,0.995)rectangle++(0.008,0.008) (-0.387,0.852)rectangle++(0.008,0.008) (-0.165,0.956)rectangle++(0.008,0.008) (-0.42,0.797)rectangle++(0.008,0.008)
  (-0.48,0.634)rectangle++(0.008,0.008) (-0.449,0.537)rectangle++(0.008,0.008) (-0.401,0.504)rectangle++(0.008,0.008) (-0.016,0.21)rectangle++(0.008,0.008) (-0.004,0.086)rectangle++(0.008,0.008)
  (-0.036,0.457)rectangle++(0.008,0.008) (-0.183,0.615)rectangle++(0.008,0.008) (-0.307,0.635)rectangle++(0.008,0.008) (-0.371,0.6)rectangle++(0.008,0.008) (-0.158,0.85)rectangle++(0.008,0.008)
  (-0.379,0.751)rectangle++(0.008,0.008) (0.094,0.895)rectangle++(0.008,0.008) (0.362,0.792)rectangle++(0.008,0.008) (0.15,0.93)rectangle++(0.008,0.008) (-0.267,0.898)rectangle++(0.008,0.008)
  (-0.446,0.734)rectangle++(0.008,0.008) (-0.19,0.906)rectangle++(0.008,0.008) (-0.082,0.978)rectangle++(0.008,0.008) (-0.037,1.009)rectangle++(0.008,0.008) (0.343,0.89)rectangle++(0.008,0.008)
  (-0.165,0.949)rectangle++(0.008,0.008) (-0.418,0.794)rectangle++(0.008,0.008) (-0.178,0.931)rectangle++(0.008,0.008) (-0.077,0.989)rectangle++(0.008,0.008) (0.317,0.896)rectangle++(0.008,0.008)
  (0.131,0.974)rectangle++(0.008,0.008) (-0.292,0.911)rectangle++(0.008,0.008) (0.191,0.937)rectangle++(0.008,0.008) (0.078,0.991)rectangle++(0.008,0.008) (-0.322,0.9)rectangle++(0.008,0.008)
  (0.174,0.943)rectangle++(0.008,0.008) (-0.261,0.912)rectangle++(0.008,0.008) (-0.112,0.981)rectangle++(0.008,0.008) (-0.049,1.01)rectangle++(0.008,0.008) (0.337,0.895)rectangle++(0.008,0.008)
  (0.473,0.705)rectangle++(0.008,0.008) (0.01,0.294)rectangle++(0.008,0.008) (0.109,0.547)rectangle++(0.008,0.008) (-0.149,0.708)rectangle++(0.008,0.008) (-0.325,0.69)rectangle++(0.008,0.008)
  (-0.014,0.287)rectangle++(0.008,0.008) (-0.008,0.718)rectangle++(0.008,0.008) (0.252,0.748)rectangle++(0.008,0.008) (0.381,0.668)rectangle++(0.008,0.008) (0.158,0.878)rectangle++(0.008,0.008)
  (0.064,0.967)rectangle++(0.008,0.008) (0.374,0.835)rectangle++(0.008,0.008) (0.468,0.665)rectangle++(0.008,0.008) (0.194,0.877)rectangle++(0.008,0.008) (-0.228,0.889)rectangle++(0.008,0.008)
  (-0.098,0.971)rectangle++(0.008,0.008) (-0.007,0.406)rectangle++(0.008,0.008) (-0.005,0.768)rectangle++(0.008,0.008) (0.271,0.769)rectangle++(0.008,0.008) (0.398,0.671)rectangle++(0.008,0.008)
  (0.165,0.88)rectangle++(0.008,0.008) (-0.242,0.88)rectangle++(0.008,0.008) (-0.429,0.735)rectangle++(0.008,0.008) (-0.182,0.906)rectangle++(0.008,0.008) (0.24,0.896)rectangle++(0.008,0.008)
  (0.429,0.74)rectangle++(0.008,0.008) (0.178,0.909)rectangle++(0.008,0.008) (0.405,0.768)rectangle++(0.008,0.008) (0.168,0.92)rectangle++(0.008,0.008) (0.405,0.777)rectangle++(0.008,0.008)
  (-0.096,0.919)rectangle++(0.008,0.008) (-0.376,0.805)rectangle++(0.008,0.008) (-0.015,0.336)rectangle++(0.008,0.008) (-0.13,0.567)rectangle++(0.008,0.008) (-0.008,0.236)rectangle++(0.008,0.008)
  (-0.006,0.697)rectangle++(0.008,0.008) (-0.255,0.736)rectangle++(0.008,0.008) (0.145,0.844)rectangle++(0.008,0.008) (-0.238,0.857)rectangle++(0.008,0.008) (-0.102,0.958)rectangle++(0.008,0.008)
  (0.294,0.891)rectangle++(0.008,0.008) (-0.187,0.931)rectangle++(0.008,0.008) (-0.081,0.989)rectangle++(0.008,0.008) (-0.036,1.013)rectangle++(0.008,0.008) (-0.018,1.023)rectangle++(0.008,0.008)
  (-0.377,0.881)rectangle++(0.008,0.008) (-0.49,0.687)rectangle++(0.008,0.008) (0.021,0.906)rectangle++(0.008,0.008) (-0.318,0.84)rectangle++(0.008,0.008) (0.154,0.914)rectangle++(0.008,0.008)
  (0.396,0.779)rectangle++(0.008,0.008) (0.458,0.631)rectangle++(0.008,0.008) (0.19,0.863)rectangle++(0.008,0.008) (0.394,0.743)rectangle++(0.008,0.008) (-0.089,0.9)rectangle++(0.008,0.008)
  (0.28,0.859)rectangle++(0.008,0.008) (0.434,0.709)rectangle++(0.008,0.008) (-0.058,0.898)rectangle++(0.008,0.008) (-0.006,0.375)rectangle++(0.008,0.008) (0.131,0.59)rectangle++(0.008,0.008)
  (-0.155,0.735)rectangle++(0.008,0.008) (-0.009,0.307)rectangle++(0.008,0.008) (0.105,0.56)rectangle++(0.008,0.008) (-0.155,0.712)rectangle++(0.008,0.008) (-0.329,0.689)rectangle++(0.008,0.008)
  (0.095,0.849)rectangle++(0.008,0.008) (0.346,0.77)rectangle++(0.008,0.008) (0.432,0.645)rectangle++(0.008,0.008) (0.179,0.869)rectangle++(0.008,0.008) (0.073,0.962)rectangle++(0.008,0.008)
  (-0.314,0.885)rectangle++(0.008,0.008) (-0.134,0.969)rectangle++(0.008,0.008) (-0.411,0.814)rectangle++(0.008,0.008) (-0.175,0.94)rectangle++(0.008,0.008) (-0.009,0.392)rectangle++(0.008,0.008)
  (-0.148,0.595)rectangle++(0.008,0.008) (-0.064,0.848)rectangle++(0.008,0.008) (0.273,0.827)rectangle++(0.008,0.008) (0.419,0.697)rectangle++(0.008,0.008) (-0.061,0.887)rectangle++(0.008,0.008)
  (-0.348,0.803)rectangle++(0.008,0.008) (-0.149,0.935)rectangle++(0.008,0.008) (-0.008,0.39)rectangle++(0.008,0.008) (-0.147,0.595)rectangle++(0.008,0.008) (0.145,0.741)rectangle++(0.008,0.008)
  (0.058,0.909)rectangle++(0.008,0.008) (0.35,0.811)rectangle++(0.008,0.008) (0.145,0.938)rectangle++(0.008,0.008) (0.059,0.992)rectangle++(0.008,0.008) (0.38,0.849)rectangle++(0.008,0.008)
  (0.157,0.954)rectangle++(0.008,0.008) (-0.272,0.912)rectangle++(0.008,0.008) (-0.012,0.381)rectangle++(0.008,0.008) (-0.007,0.758)rectangle++(0.008,0.008) (-0.005,0.916)rectangle++(0.008,0.008)
  (0.324,0.837)rectangle++(0.008,0.008) (0.446,0.683)rectangle++(0.008,0.008) (-0.044,0.891)rectangle++(0.008,0.008) (-0.342,0.811)rectangle++(0.008,0.008) (-0.014,0.338)rectangle++(0.008,0.008)
  (-0.004,0.14)rectangle++(0.008,0.008) (0.047,0.482)rectangle++(0.008,0.008) (0.017,0.8)rectangle++(0.008,0.008) (0.005,0.934)rectangle++(0.008,0.008) (0.335,0.841)rectangle++(0.008,0.008)
  (0.453,0.681)rectangle++(0.008,0.008) (0.449,0.566)rectangle++(0.008,0.008) (0.186,0.835)rectangle++(0.008,0.008) (0.383,0.732)rectangle++(0.008,0.008) (0.158,0.905)rectangle++(0.008,0.008)
  (0.064,0.978)rectangle++(0.008,0.008) (0.378,0.84)rectangle++(0.008,0.008) (0.472,0.666)rectangle++(0.008,0.008) (-0.026,0.892)rectangle++(0.008,0.008) (0.306,0.833)rectangle++(0.008,0.008)
  (0.437,0.688)rectangle++(0.008,0.008) (-0.05,0.89)rectangle++(0.008,0.008) (-0.344,0.808)rectangle++(0.008,0.008) (-0.449,0.666)rectangle++(0.008,0.008) (-0.447,0.563)rectangle++(0.008,0.008)
  (-0.004,0.833)rectangle++(0.008,0.008) (-0.303,0.799)rectangle++(0.008,0.008) (-0.13,0.933)rectangle++(0.008,0.008) (0.273,0.889)rectangle++(0.008,0.008) (0.442,0.725)rectangle++(0.008,0.008)
  (0.46,0.59)rectangle++(0.008,0.008) (-0.004,0.853)rectangle++(0.008,0.008) (-0.31,0.808)rectangle++(0.008,0.008) (-0.133,0.937)rectangle++(0.008,0.008) (-0.058,0.991)rectangle++(0.008,0.008)
  (0.327,0.89)rectangle++(0.008,0.008) (0.466,0.707)rectangle++(0.008,0.008) (-0.043,0.909)rectangle++(0.008,0.008) (0.304,0.847)rectangle++(0.008,0.008) (-0.167,0.915)rectangle++(0.008,0.008)
  (-0.407,0.778)rectangle++(0.008,0.008) (0.091,0.917)rectangle++(0.008,0.008) (-0.289,0.871)rectangle++(0.008,0.008) (-0.447,0.714)rectangle++(0.008,0.008) (0.05,0.902)rectangle++(0.008,0.008)
  (0.344,0.811)rectangle++(0.008,0.008) (0.142,0.938)rectangle++(0.008,0.008) (-0.274,0.899)rectangle++(0.008,0.008) (-0.012,0.375)rectangle++(0.008,0.008) (-0.007,0.755)rectangle++(0.008,0.008)
  (0.266,0.764)rectangle++(0.008,0.008) (0.394,0.671)rectangle++(0.008,0.008) (0.419,0.582)rectangle++(0.008,0.008) (-0.02,0.835)rectangle++(0.008,0.008) (-0.011,0.948)rectangle++(0.008,0.008)
  (0.333,0.854)rectangle++(0.008,0.008) (0.138,0.956)rectangle++(0.008,0.008) (-0.282,0.905)rectangle++(0.008,0.008) (-0.456,0.732)rectangle++(0.008,0.008) (0.052,0.914)rectangle++(0.008,0.008)
  (-0.306,0.856)rectangle++(0.008,0.008) (-0.131,0.957)rectangle++(0.008,0.008) (-0.405,0.81)rectangle++(0.008,0.008) (-0.016,0.338)rectangle++(0.008,0.008) (-0.132,0.568)rectangle++(0.008,0.008)
  (-0.267,0.632)rectangle++(0.008,0.008) (-0.114,0.863)rectangle++(0.008,0.008) (-0.05,0.96)rectangle++(0.008,0.008) (0.319,0.873)rectangle++(0.008,0.008) (-0.17,0.932)rectangle++(0.008,0.008)
  (-0.009,0.389)rectangle++(0.008,0.008) (-0.006,0.761)rectangle++(0.008,0.008) (-0.278,0.765)rectangle++(0.008,0.008) (-0.404,0.67)rectangle++(0.008,0.008) (0.054,0.867)rectangle++(0.008,0.008)
  (0.02,0.962)rectangle++(0.008,0.008) (-0.338,0.866)rectangle++(0.008,0.008) (-0.014,0.361)rectangle++(0.008,0.008) (-0.008,0.749)rectangle++(0.008,0.008) (-0.004,0.312)rectangle++(0.008,0.008)
  (0.109,0.561)rectangle++(0.008,0.008) (-0.154,0.714)rectangle++(0.008,0.008) (-0.329,0.691)rectangle++(0.008,0.008) (-0.401,0.617)rectangle++(0.008,0.008) (0.037,0.842)rectangle++(0.008,0.008)
  (0.013,0.951)rectangle++(0.008,0.008) (0.003,0.997)rectangle++(0.008,0.008) (0.357,0.871)rectangle++(0.008,0.008) (0.007,0.364)rectangle++(0.008,0.008) (-0.13,0.588)rectangle++(0.008,0.008)
  (-0.273,0.641)rectangle++(0.008,0.008) (-0.117,0.867)rectangle++(0.008,0.008) (-0.051,0.962)rectangle++(0.008,0.008) (-0.005,0.402)rectangle++(0.008,0.008) (-0.15,0.601)rectangle++(0.008,0.008)
  (-0.065,0.85)rectangle++(0.008,0.008) (-0.337,0.785)rectangle++(0.008,0.008) (-0.014,0.327)rectangle++(0.008,0.008) (-0.127,0.564)rectangle++(0.008,0.008) (-0.056,0.834)rectangle++(0.008,0.008)
  (-0.327,0.781)rectangle++(0.008,0.008) (0.129,0.89)rectangle++(0.008,0.008) (-0.263,0.872)rectangle++(0.008,0.008) (0.191,0.909)rectangle++(0.008,0.008) (-0.241,0.903)rectangle++(0.008,0.008)
  (0.211,0.915)rectangle++(0.008,0.008) (0.423,0.759)rectangle++(0.008,0.008) (0.175,0.916)rectangle++(0.008,0.008) (0.407,0.773)rectangle++(0.008,0.008) (0.168,0.922)rectangle++(0.008,0.008)
  (-0.256,0.901)rectangle++(0.008,0.008) (-0.11,0.976)rectangle++(0.008,0.008) (-0.048,1.008)rectangle++(0.008,0.008) (-0.023,1.021)rectangle++(0.008,0.008) (-0.378,0.878)rectangle++(0.008,0.008)
  (-0.49,0.685)rectangle++(0.008,0.008) (0.02,0.905)rectangle++(0.008,0.008) (0.006,0.978)rectangle++(0.008,0.008) (-0.35,0.868)rectangle++(0.008,0.008) (-0.149,0.962)rectangle++(0.008,0.008)
  (0.275,0.91)rectangle++(0.008,0.008) (0.004,0.38)rectangle++(0.008,0.008) (-0.004,0.157)rectangle++(0.008,0.008) (-0.004,0.064)rectangle++(0.008,0.008) (-0.028,0.447)rectangle++(0.008,0.008)
  (-0.014,0.785)rectangle++(0.008,0.008) (0.273,0.78)rectangle++(0.008,0.008) (0.112,0.925)rectangle++(0.008,0.008) (0.045,0.986)rectangle++(0.008,0.008) (-0.335,0.886)rectangle++(0.008,0.008)
  (0.162,0.941)rectangle++(0.008,0.008) (0.41,0.789)rectangle++(0.008,0.008) (0.468,0.631)rectangle++(0.008,0.008) (0.01,0.263)rectangle++(0.008,0.008) (0.002,0.708)rectangle++(0.008,0.008)
  (-0.256,0.744)rectangle++(0.008,0.008) (0.148,0.847)rectangle++(0.008,0.008) (0.37,0.751)rectangle++(0.008,0.008) (0.436,0.628)rectangle++(0.008,0.008) (0.181,0.861)rectangle++(0.008,0.008)
  (-0.228,0.878)rectangle++(0.008,0.008) (-0.421,0.739)rectangle++(0.008,0.008) (0.07,0.904)rectangle++(0.008,0.008) (-0.002,0.378)rectangle++(0.008,0.008) (-0.003,0.756)rectangle++(0.008,0.008)
  (-0.275,0.764)rectangle++(0.008,0.008) (-0.402,0.67)rectangle++(0.008,0.008) (-0.427,0.582)rectangle++(0.008,0.008) (0.012,0.835)rectangle++(0.008,0.008) (-0.296,0.805)rectangle++(0.008,0.008)
  (-0.426,0.681)rectangle++(0.008,0.008) (-0.017,0.284)rectangle++(0.008,0.008) (-0.113,0.543)rectangle++(0.008,0.008) (0.142,0.705)rectangle++(0.008,0.008) (0.057,0.894)rectangle++(0.008,0.008)
  (-0.002,0.373)rectangle++(0.008,0.008) (0.131,0.588)rectangle++(0.008,0.008) (-0.153,0.735)rectangle++(0.008,0.008) (-0.336,0.7)rectangle++(0.008,0.008) (-0.407,0.619)rectangle++(0.008,0.008)
  (-0.411,0.557)rectangle++(0.008,0.008) (-0.39,0.527)rectangle++(0.008,0.008) (-0.37,0.521)rectangle++(0.008,0.008) (-0.015,0.216)rectangle++(0.008,0.008) (-0.088,0.513)rectangle++(0.008,0.008)
  (-0.039,0.813)rectangle++(0.008,0.008) (0.272,0.802)rectangle++(0.008,0.008) (0.112,0.935)rectangle++(0.008,0.008) (0.045,0.99)rectangle++(0.008,0.008) (0.373,0.853)rectangle++(0.008,0.008)
  (-0.138,0.942)rectangle++(0.008,0.008) (-0.403,0.801)rectangle++(0.008,0.008) (-0.172,0.934)rectangle++(0.008,0.008) (0.254,0.905)rectangle++(0.008,0.008) (0.439,0.739)rectangle++(0.008,0.008)
  (-0.067,0.914)rectangle++(0.008,0.008) (-0.03,0.981)rectangle++(0.008,0.008) (0.336,0.876)rectangle++(0.008,0.008) (0.006,0.366)rectangle++(0.008,0.008) (-0.131,0.589)rectangle++(0.008,0.008)
  (-0.057,0.845)rectangle++(0.008,0.008) (0.275,0.823)rectangle++(0.008,0.008) (0.004,0.343)rectangle++(0.008,0.008) (-0.004,0.142)rectangle++(0.008,0.008) (0.048,0.483)rectangle++(0.008,0.008)
  (-0.154,0.657)rectangle++(0.008,0.008) (0.164,0.772)rectangle++(0.008,0.008) (-0.204,0.83)rectangle++(0.008,0.008) (-0.394,0.726)rectangle++(0.008,0.008) (0.079,0.889)rectangle++(0.008,0.008)
  (-0.285,0.853)rectangle++(0.008,0.008) (-0.439,0.708)rectangle++(0.008,0.008) (-0.457,0.586)rectangle++(0.008,0.008) (0,0.847)rectangle++(0.008,0.008) (-0.306,0.806)rectangle++(0.008,0.008)
  (-0.013,0.336)rectangle++(0.008,0.008) (0.113,0.575)rectangle++(0.008,0.008) (0,0.239)rectangle++(0.008,0.008) (0.084,0.526)rectangle++(0.008,0.008) (0.033,0.819)rectangle++(0.008,0.008)
  (0.307,0.779)rectangle++(0.008,0.008) (0.126,0.925)rectangle++(0.008,0.008) (-0.276,0.887)rectangle++(0.008,0.008) (-0.446,0.726)rectangle++(0.008,0.008) (-0.467,0.592)rectangle++(0.008,0.008)
  (-0.018,0.246)rectangle++(0.008,0.008) (-0.01,0.701)rectangle++(0.008,0.008) (-0.004,0.292)rectangle++(0.008,0.008) (-0.004,0.72)rectangle++(0.008,0.008) (0.255,0.747)rectangle++(0.008,0.008)
  (0.105,0.911)rectangle++(0.008,0.008) (-0.001,0.38)rectangle++(0.008,0.008) (-0.003,0.757)rectangle++(0.008,0.008) (-0.003,0.916)rectangle++(0.008,0.008) (-0.004,0.982)rectangle++(0.008,0.008)
  (0.348,0.867)rectangle++(0.008,0.008) (0.468,0.688)rectangle++(0.008,0.008) (0.459,0.564)rectangle++(0.008,0.008) (0.19,0.834)rectangle++(0.008,0.008) (0.384,0.73)rectangle++(0.008,0.008)
  (0.435,0.613)rectangle++(0.008,0.008) (-0.023,0.855)rectangle++(0.008,0.008) (0.294,0.815)rectangle++(0.008,0.008) (0.425,0.684)rectangle++(0.008,0.008) (0.176,0.885)rectangle++(0.008,0.008)
  (0.072,0.969)rectangle++(0.008,0.008) (-0.317,0.888)rectangle++(0.008,0.008) (0.171,0.935)rectangle++(0.008,0.008) (-0.259,0.908)rectangle++(0.008,0.008) (0.205,0.924)rectangle++(0.008,0.008)
  (0.423,0.766)rectangle++(0.008,0.008) (0.466,0.615)rectangle++(0.008,0.008) (0.01,0.256)rectangle++(0.008,0.008) (-0.09,0.54)rectangle++(0.008,0.008) (-0.238,0.634)rectangle++(0.008,0.008)
  (0.117,0.791)rectangle++(0.008,0.008) (-0.233,0.822)rectangle++(0.008,0.008) (0.187,0.875)rectangle++(0.008,0.008) (0.076,0.965)rectangle++(0.008,0.008) (-0.314,0.888)rectangle++(0.008,0.008)
  (-0.464,0.713)rectangle++(0.008,0.008) (0.042,0.908)rectangle++(0.008,0.008) (-0.309,0.849)rectangle++(0.008,0.008) (0.161,0.915)rectangle++(0.008,0.008) (-0.256,0.895)rectangle++(0.008,0.008)
  (0.202,0.917)rectangle++(0.008,0.008) (-0.239,0.91)rectangle++(0.008,0.008) (-0.103,0.98)rectangle++(0.008,0.008) (0.302,0.901)rectangle++(0.008,0.008) (0.459,0.72)rectangle++(0.008,0.008)
  (0.191,0.9)rectangle++(0.008,0.008) (0.078,0.976)rectangle++(0.008,0.008) (-0.002,0.407)rectangle++(0.008,0.008) (-0.003,0.769)rectangle++(0.008,0.008) (-0.279,0.77)rectangle++(0.008,0.008)
  (-0.12,0.921)rectangle++(0.008,0.008) (-0.053,0.984)rectangle++(0.008,0.008) (-0.379,0.85)rectangle++(0.008,0.008) (0.13,0.94)rectangle++(0.008,0.008) (0.394,0.8)rectangle++(0.008,0.008)
  (0.465,0.641)rectangle++(0.008,0.008) (0.441,0.543)rectangle++(0.008,0.008) (0.395,0.507)rectangle++(0.008,0.008) (0.361,0.507)rectangle++(0.008,0.008) (0.345,0.519)rectangle++(0.008,0.008)
  (0.006,0.216)rectangle++(0.008,0.008) (0.079,0.513)rectangle++(0.008,0.008) (0.219,0.622)rectangle++(0.008,0.008) (0.089,0.859)rectangle++(0.008,0.008) (0.347,0.777)rectangle++(0.008,0.008)
  (0.435,0.648)rectangle++(0.008,0.008) (-0.036,0.871)rectangle++(0.008,0.008) (-0.331,0.804)rectangle++(0.008,0.008) (-0.141,0.936)rectangle++(0.008,0.008) (-0.008,0.391)rectangle++(0.008,0.008)
  (-0.006,0.762)rectangle++(0.008,0.008) (-0.278,0.765)rectangle++(0.008,0.008) (-0.404,0.67)rectangle++(0.008,0.008) (0.054,0.867)rectangle++(0.008,0.008) (-0.002,0.362)rectangle++(0.008,0.008)
  (0.127,0.583)rectangle++(0.008,0.008) (0.051,0.842)rectangle++(0.008,0.008) (-0.281,0.823)rectangle++(0.008,0.008) (-0.426,0.695)rectangle++(0.008,0.008) (-0.017,0.29)rectangle++(0.008,0.008)
  (-0.009,0.719)rectangle++(0.008,0.008) (-0.265,0.745)rectangle++(0.008,0.008) (0.144,0.851)rectangle++(0.008,0.008) (0.369,0.754)rectangle++(0.008,0.008) (-0.104,0.896)rectangle++(0.008,0.008)
  (0.271,0.863)rectangle++(0.008,0.008) (0.431,0.714)rectangle++(0.008,0.008) (-0.061,0.9)rectangle++(0.008,0.008) (-0.353,0.809)rectangle++(0.008,0.008) (-0.151,0.937)rectangle++(0.008,0.008)
  (-0.407,0.794)rectangle++(0.008,0.008) (0.097,0.925)rectangle++(0.008,0.008) (0.374,0.804)rectangle++(0.008,0.008) (0.155,0.936)rectangle++(0.008,0.008) (0.063,0.991)rectangle++(0.008,0.008)
  (0.024,1.014)rectangle++(0.008,0.008) (0.372,0.871)rectangle++(0.008,0.008) (0.48,0.682)rectangle++(0.008,0.008) (0.011,0.284)rectangle++(0.008,0.008) (0.105,0.542)rectangle++(0.008,0.008)
  (0.241,0.627)rectangle++(0.008,0.008) (0.333,0.617)rectangle++(0.008,0.008) (-0.072,0.82)rectangle++(0.008,0.008) (-0.329,0.769)rectangle++(0.008,0.008) (-0.428,0.653)rectangle++(0.008,0.008)
  (-0.182,0.872)rectangle++(0.008,0.008) (0.227,0.88)rectangle++(0.008,0.008) (0.417,0.737)rectangle++(0.008,0.008) (0.009,0.307)rectangle++(0.008,0.008) (-0.109,0.563)rectangle++(0.008,0.008)
  (-0.255,0.637)rectangle++(0.008,0.008) (0.111,0.799)rectangle++(0.008,0.008) (-0.238,0.824)rectangle++(0.008,0.008) (-0.102,0.944)rectangle++(0.008,0.008) (-0.387,0.814)rectangle++(0.008,0.008)
  (-0.016,0.34)rectangle++(0.008,0.008) (-0.132,0.569)rectangle++(0.008,0.008) (0.142,0.724)rectangle++(0.008,0.008) (0.323,0.696)rectangle++(0.008,0.008) (0.133,0.89)rectangle++(0.008,0.008)
  (0,0.372)rectangle++(0.008,0.008) (-0.002,0.754)rectangle++(0.008,0.008) (-0.003,0.914)rectangle++(0.008,0.008) (0.324,0.835)rectangle++(0.008,0.008) (0.446,0.682)rectangle++(0.008,0.008)
  (-0.044,0.89)rectangle++(0.008,0.008) (-0.341,0.811)rectangle++(0.008,0.008) (-0.014,0.338)rectangle++(0.008,0.008) (-0.008,0.74)rectangle++(0.008,0.008) (-0.006,0.908)rectangle++(0.008,0.008)
  (0.321,0.834)rectangle++(0.008,0.008) (0.444,0.683)rectangle++(0.008,0.008) (0.446,0.57)rectangle++(0.008,0.008) (0.407,0.518)rectangle++(0.008,0.008) (0.008,0.215)rectangle++(0.008,0.008)
  (0.08,0.512)rectangle++(0.008,0.008) (-0.15,0.682)rectangle++(0.008,0.008) (0.174,0.781)rectangle++(0.008,0.008) (0.358,0.711)rectangle++(0.008,0.008) (0.148,0.896)rectangle++(0.008,0.008)
  (0.387,0.773)rectangle++(0.008,0.008) (-0.103,0.911)rectangle++(0.008,0.008) (-0.046,0.98)rectangle++(0.008,0.008) (-0.374,0.851)rectangle++(0.008,0.008) (0.132,0.939)rectangle++(0.008,0.008)
  (0.053,0.992)rectangle++(0.008,0.008) (0.02,1.014)rectangle++(0.008,0.008) (-0.357,0.89)rectangle++(0.008,0.008) (0.154,0.951)rectangle++(0.008,0.008) (-0.273,0.909)rectangle++(0.008,0.008)
  (-0.117,0.979)rectangle++(0.008,0.008) (-0.407,0.825)rectangle++(0.008,0.008) (-0.173,0.944)rectangle++(0.008,0.008) (-0.075,0.994)rectangle++(0.008,0.008) (-0.393,0.847)rectangle++(0.008,0.008)
  (-0.167,0.953)rectangle++(0.008,0.008) (0.263,0.912)rectangle++(0.008,0.008) (-0.209,0.93)rectangle++(0.008,0.008) (-0.431,0.77)rectangle++(0.008,0.008) (-0.476,0.617)rectangle++(0.008,0.008)
  (-0.202,0.857)rectangle++(0.008,0.008) (-0.402,0.739)rectangle++(0.008,0.008) (-0.451,0.613)rectangle++(0.008,0.008) (0.012,0.858)rectangle++(0.008,0.008) (0.003,0.958)rectangle++(0.008,0.008)
  (-0.344,0.858)rectangle++(0.008,0.008) (0.148,0.932)rectangle++(0.008,0.008) (-0.268,0.898)rectangle++(0.008,0.008) (-0.115,0.975)rectangle++(0.008,0.008) (0.295,0.903)rectangle++(0.008,0.008)
  (0.121,0.977)rectangle++(0.008,0.008) (0.404,0.82)rectangle++(0.008,0.008) (-0.112,0.938)rectangle++(0.008,0.008) (-0.049,0.992)rectangle++(0.008,0.008) (-0.38,0.855)rectangle++(0.008,0.008)
  (0.131,0.943)rectangle++(0.008,0.008) (0.053,0.994)rectangle++(0.008,0.008) (-0.334,0.892)rectangle++(0.008,0.008) (0.165,0.944)rectangle++(0.008,0.008) (0.412,0.789)rectangle++(0.008,0.008)
  (0.469,0.63)rectangle++(0.008,0.008) (-0.014,0.875)rectangle++(0.008,0.008) (-0.322,0.814)rectangle++(0.008,0.008) (-0.138,0.94)rectangle++(0.008,0.008) (0.272,0.895)rectangle++(0.008,0.008)
  (-0.199,0.925)rectangle++(0.008,0.008) (0.239,0.91)rectangle++(0.008,0.008) (0.098,0.98)rectangle++(0.008,0.008) (0.394,0.829)rectangle++(0.008,0.008) (0.163,0.946)rectangle++(0.008,0.008)
  (-0.267,0.91)rectangle++(0.008,0.008) (-0.45,0.74)rectangle++(0.008,0.008) (-0.192,0.908)rectangle++(0.008,0.008) (-0.416,0.766)rectangle++(0.008,0.008) (-0.467,0.621)rectangle++(0.008,0.008)
  (-0.439,0.536)rectangle++(0.008,0.008) (-0.187,0.823)rectangle++(0.008,0.008) (-0.081,0.943)rectangle++(0.008,0.008) (-0.377,0.822)rectangle++(0.008,0.008) (0.12,0.927)rectangle++(0.008,0.008)
  (-0.28,0.886)rectangle++(0.008,0.008) (0.188,0.921)rectangle++(0.008,0.008) (0.002,0.385)rectangle++(0.008,0.008) (0.137,0.592)rectangle++(0.008,0.008) (0.273,0.638)rectangle++(0.008,0.008)
  (0.352,0.61)rectangle++(0.008,0.008) (0.378,0.57)rectangle++(0.008,0.008) (-0.034,0.815)rectangle++(0.008,0.008) (0.274,0.801)rectangle++(0.008,0.008) (0.113,0.934)rectangle++(0.008,0.008)
  (-0.286,0.887)rectangle++(0.008,0.008) (0.185,0.924)rectangle++(0.008,0.008) (0.414,0.772)rectangle++(0.008,0.008) (0.171,0.922)rectangle++(0.008,0.008) (0.407,0.777)rectangle++(0.008,0.008)
  (-0.095,0.919)rectangle++(0.008,0.008) (-0.042,0.984)rectangle++(0.008,0.008) (-0.02,1.011)rectangle++(0.008,0.008) (-0.374,0.874)rectangle++(0.008,0.008) (0.141,0.949)rectangle++(0.008,0.008)
  (-0.278,0.903)rectangle++(0.008,0.008) (0.195,0.929)rectangle++(0.008,0.008) (-0.246,0.913)rectangle++(0.008,0.008) (0.213,0.922)rectangle++(0.008,0.008) (-0.235,0.917)rectangle++(0.008,0.008)
  (-0.101,0.983)rectangle++(0.008,0.008) (0.304,0.902)rectangle++(0.008,0.008) (0.46,0.72)rectangle++(0.008,0.008) (0.467,0.581)rectangle++(0.008,0.008) (0.42,0.515)rectangle++(0.008,0.008)
  (0.004,0.805)rectangle++(0.008,0.008) (0.289,0.783)rectangle++(0.008,0.008) (0.411,0.671)rectangle++(0.008,0.008) (-0.056,0.873)rectangle++(0.008,0.008) (0.285,0.835)rectangle++(0.008,0.008)
  (-0.171,0.903)rectangle++(0.008,0.008) (0.243,0.89)rectangle++(0.008,0.008) (0.1,0.972)rectangle++(0.008,0.008) (-0.305,0.899)rectangle++(0.008,0.008) (-0.464,0.721)rectangle++(0.008,0.008)
  (-0.197,0.901)rectangle++(0.008,0.008) (0.231,0.898)rectangle++(0.008,0.008) (-0.219,0.912)rectangle++(0.008,0.008) (-0.01,0.381)rectangle++(0.008,0.008) (0.13,0.594)rectangle++(0.008,0.008)
  (-0.156,0.737)rectangle++(0.008,0.008) (-0.068,0.907)rectangle++(0.008,0.008) (0.292,0.855)rectangle++(0.008,0.008) (-0.175,0.915)rectangle++(0.008,0.008) (-0.41,0.775)rectangle++(0.008,0.008)
  (-0.016,0.323)rectangle++(0.008,0.008) (-0.126,0.561)rectangle++(0.008,0.008) (0.142,0.718)rectangle++(0.008,0.008) (0.057,0.899)rectangle++(0.008,0.008) (-0.299,0.851)rectangle++(0.008,0.008)
  (-0.128,0.955)rectangle++(0.008,0.008) (-0.403,0.81)rectangle++(0.008,0.008) (-0.477,0.646)rectangle++(0.008,0.008) (-0.203,0.869)rectangle++(0.008,0.008) (-0.406,0.744)rectangle++(0.008,0.008)
  (-0.455,0.614)rectangle++(0.008,0.008) (0.011,0.86)rectangle++(0.008,0.008) (0.311,0.805)rectangle++(0.008,0.008) (-0.149,0.898)rectangle++(0.008,0.008) (-0.393,0.777)rectangle++(0.008,0.008)
  (0.097,0.912)rectangle++(0.008,0.008) (0.369,0.798)rectangle++(0.008,0.008) (0.453,0.649)rectangle++(0.008,0.008) (-0.028,0.878)rectangle++(0.008,0.008) (-0.014,0.966)rectangle++(0.008,0.008)
  (-0.355,0.856)rectangle++(0.008,0.008) (0.143,0.934)rectangle++(0.008,0.008) (0.058,0.99)rectangle++(0.008,0.008) (-0.331,0.892)rectangle++(0.008,0.008) (0.167,0.942)rectangle++(0.008,0.008)
  (-0.264,0.909)rectangle++(0.008,0.008) (0.203,0.926)rectangle++(0.008,0.008) (-0.241,0.915)rectangle++(0.008,0.008) (-0.441,0.751)rectangle++(0.008,0.008) (-0.473,0.605)rectangle++(0.008,0.008)
  (-0.001,0.862)rectangle++(0.008,0.008) (-0.312,0.813)rectangle++(0.008,0.008) (-0.133,0.939)rectangle++(0.008,0.008) (-0.058,0.992)rectangle++(0.008,0.008) (0.327,0.891)rectangle++(0.008,0.008)
  (-0.172,0.943)rectangle++(0.008,0.008) (-0.075,0.994)rectangle++(0.008,0.008) (0.32,0.897)rectangle++(0.008,0.008) (0.132,0.975)rectangle++(0.008,0.008) (0.053,1.007)rectangle++(0.008,0.008)
  (-0.339,0.898)rectangle++(0.008,0.008) (-0.145,0.975)rectangle++(0.008,0.008) (-0.063,1.007)rectangle++(0.008,0.008) (-0.392,0.857)rectangle++(0.008,0.008) (-0.489,0.671)rectangle++(0.008,0.008)
  (-0.019,0.279)rectangle++(0.008,0.008) (-0.01,0.715)rectangle++(0.008,0.008) (-0.264,0.742)rectangle++(0.008,0.008) (-0.012,0.31)rectangle++(0.008,0.008) (-0.004,0.128)rectangle++(0.008,0.008)
  (0.043,0.476)rectangle++(0.008,0.008) (0.189,0.619)rectangle++(0.008,0.008) (0.077,0.858)rectangle++(0.008,0.008) (0.341,0.781)rectangle++(0.008,0.008) (-0.127,0.898)rectangle++(0.008,0.008)
  (-0.056,0.975)rectangle++(0.008,0.008) (-0.026,1.007)rectangle++(0.008,0.008) (-0.013,1.021)rectangle++(0.008,0.008) (-0.008,1.026)rectangle++(0.008,0.008) (-0.006,1.029)rectangle++(0.008,0.008)
  (0.364,0.889)rectangle++(0.008,0.008) (0.151,0.971)rectangle++(0.008,0.008) (0.001,0.405)rectangle++(0.008,0.008) (-0.002,0.768)rectangle++(0.008,0.008) (-0.004,0.32)rectangle++(0.008,0.008)
  (-0.12,0.564)rectangle++(0.008,0.008) (0.146,0.717)rectangle++(0.008,0.008) (0.322,0.692)rectangle++(0.008,0.008) (-0.104,0.851)rectangle++(0.008,0.008) (0.256,0.842)rectangle++(0.008,0.008)
  (-0.187,0.895)rectangle++(0.008,0.008) (0.233,0.893)rectangle++(0.008,0.008) (0.003,0.373)rectangle++(0.008,0.008) (-0.135,0.591)rectangle++(0.008,0.008) (0.148,0.735)rectangle++(0.008,0.008)
  (0.329,0.699)rectangle++(0.008,0.008) (0.006,0.291)rectangle++(0.008,0.008) (0.106,0.547)rectangle++(0.008,0.008) (0.042,0.828)rectangle++(0.008,0.008) (0.314,0.78)rectangle++(0.008,0.008)
  (0.421,0.661)rectangle++(0.008,0.008) (-0.047,0.872)rectangle++(0.008,0.008) (-0.022,0.964)rectangle++(0.008,0.008) (0.333,0.865)rectangle++(0.008,0.008) (0.138,0.961)rectangle++(0.008,0.008)
  (0.405,0.806)rectangle++(0.008,0.008) (0.168,0.936)rectangle++(0.008,0.008) (-0.261,0.907)rectangle++(0.008,0.008) (-0.112,0.979)rectangle++(0.008,0.008) (-0.049,1.009)rectangle++(0.008,0.008)
  (-0.023,1.021)rectangle++(0.008,0.008) (-0.012,1.027)rectangle++(0.008,0.008) (-0.376,0.884)rectangle++(0.008,0.008) (-0.491,0.689)rectangle++(0.008,0.008) (-0.209,0.887)rectangle++(0.008,0.008)
  (-0.09,0.97)rectangle++(0.008,0.008) (0.305,0.892)rectangle++(0.008,0.008) (0.457,0.715)rectangle++(0.008,0.008) (0.01,0.298)rectangle++(0.008,0.008) (0.11,0.549)rectangle++(0.008,0.008)
  (-0.001,0.228)rectangle++(0.008,0.008) (0.081,0.521)rectangle++(0.008,0.008) (-0.153,0.686)rectangle++(0.008,0.008) (0.174,0.785)rectangle++(0.008,0.008) (-0.204,0.84)rectangle++(0.008,0.008)
  (-0.397,0.73)rectangle++(0.008,0.008) (-0.169,0.904)rectangle++(0.008,0.008) (0.245,0.89)rectangle++(0.008,0.008) (-0.209,0.914)rectangle++(0.008,0.008) (-0.09,0.981)rectangle++(0.008,0.008)
  (-0.007,0.41)rectangle++(0.008,0.008) (-0.153,0.604)rectangle++(0.008,0.008) (0.145,0.747)rectangle++(0.008,0.008) (0.332,0.706)rectangle++(0.008,0.008) (0.403,0.62)rectangle++(0.008,0.008)
  (0.167,0.858)rectangle++(0.008,0.008) (0.001,0.358)rectangle++(0.008,0.008) (0.128,0.58)rectangle++(0.008,0.008) (-0.152,0.73)rectangle++(0.008,0.008) (0.19,0.804)rectangle++(0.008,0.008)
  (0.373,0.716)rectangle++(0.008,0.008) (0.426,0.61)rectangle++(0.008,0.008) (0.176,0.854)rectangle++(0.008,0.008) (0.385,0.744)rectangle++(0.008,0.008) (-0.093,0.897)rectangle++(0.008,0.008)
  (0.277,0.859)rectangle++(0.008,0.008) (-0.184,0.911)rectangle++(0.008,0.008) (-0.08,0.98)rectangle++(0.008,0.008) (-0.39,0.839)rectangle++(0.008,0.008) (0.121,0.939)rectangle++(0.008,0.008)
  (0.048,0.992)rectangle++(0.008,0.008) (0.018,1.014)rectangle++(0.008,0.008) (0.005,1.024)rectangle++(0.008,0.008) (0.367,0.882)rectangle++(0.008,0.008) (0.482,0.689)rectangle++(0.008,0.008)
  (-0.029,0.906)rectangle++(0.008,0.008) (-0.015,0.978)rectangle++(0.008,0.008) (0.342,0.869)rectangle++(0.008,0.008) (0.006,0.363)rectangle++(0.008,0.008) (0.132,0.58)rectangle++(0.008,0.008)
  (0.266,0.634)rectangle++(0.008,0.008) (-0.108,0.804)rectangle++(0.008,0.008) (-0.048,0.936)rectangle++(0.008,0.008) (0.311,0.861)rectangle++(0.008,0.008) (0.128,0.959)rectangle++(0.008,0.008)
  (-0.287,0.904)rectangle++(0.008,0.008) (0.191,0.932)rectangle++(0.008,0.008) (0.419,0.774)rectangle++(0.008,0.008) (-0.089,0.923)rectangle++(0.008,0.008) (-0.374,0.809)rectangle++(0.008,0.008)
  (-0.015,0.338)rectangle++(0.008,0.008) (-0.009,0.739)rectangle++(0.008,0.008) (-0.272,0.754)rectangle++(0.008,0.008) (-0.397,0.667)rectangle++(0.008,0.008) (-0.423,0.582)rectangle++(0.008,0.008)
  (0.014,0.834)rectangle++(0.008,0.008) (0.303,0.793)rectangle++(0.008,0.008) (0.125,0.931)rectangle++(0.008,0.008) (0,0.389)rectangle++(0.008,0.008) (-0.142,0.597)rectangle++(0.008,0.008)
  (-0.062,0.848)rectangle++(0.008,0.008) (-0.028,0.954)rectangle++(0.008,0.008) (-0.005,0.398)rectangle++(0.008,0.008) (0.139,0.6)rectangle++(0.008,0.008) (0.056,0.85)rectangle++(0.008,0.008)
  (0.021,0.955)rectangle++(0.008,0.008) (-0.335,0.863)rectangle++(0.008,0.008) (0.154,0.93)rectangle++(0.008,0.008) (-0.265,0.9)rectangle++(0.008,0.008) (-0.012,0.376)rectangle++(0.008,0.008)
  (-0.007,0.755)rectangle++(0.008,0.008) (-0.277,0.762)rectangle++(0.008,0.008) (-0.119,0.918)rectangle++(0.008,0.008) (0.273,0.878)rectangle++(0.008,0.008) (0.112,0.967)rectangle++(0.008,0.008)
  (0.396,0.818)rectangle++(0.008,0.008) (0.164,0.941)rectangle++(0.008,0.008) (0.41,0.788)rectangle++(0.008,0.008) (-0.098,0.926)rectangle++(0.008,0.008) (-0.007,0.387)rectangle++(0.008,0.008)
  (-0.145,0.594)rectangle++(0.008,0.008) (0.145,0.739)rectangle++(0.008,0.008) (-0.201,0.809)rectangle++(0.008,0.008) (-0.01,0.337)rectangle++(0.008,0.008) (-0.129,0.57)rectangle++(0.008,0.008)
  (-0.008,0.237)rectangle++(0.008,0.008) (-0.092,0.525)rectangle++(0.008,0.008) (0.145,0.689)rectangle++(0.008,0.008) (-0.184,0.786)rectangle++(0.008,0.008) (-0.079,0.928)rectangle++(0.008,0.008)
  (-0.371,0.815)rectangle++(0.008,0.008) (-0.464,0.659)rectangle++(0.008,0.008) (-0.451,0.555)rectangle++(0.008,0.008) (-0.192,0.831)rectangle++(0.008,0.008) (-0.01,0.347)rectangle++(0.008,0.008)
  (-0.006,0.743)rectangle++(0.008,0.008) (-0.004,0.31)rectangle++(0.008,0.008) (-0.004,0.728)rectangle++(0.008,0.008) (0.257,0.75)rectangle++(0.008,0.008) (0.385,0.668)rectangle++(0.008,0.008)
  (0.414,0.584)rectangle++(0.008,0.008) (0.397,0.536)rectangle++(0.008,0.008) (0.372,0.52)rectangle++(0.008,0.008) (-0.019,0.79)rectangle++(0.008,0.008) (-0.294,0.773)rectangle++(0.008,0.008)
  (-0.126,0.922)rectangle++(0.008,0.008) (-0.391,0.796)rectangle++(0.008,0.008) (-0.466,0.644)rectangle++(0.008,0.008) (-0.018,0.268)rectangle++(0.008,0.008) (0.087,0.545)rectangle++(0.008,0.008)
  (0.034,0.827)rectangle++(0.008,0.008) (-0.003,0.345)rectangle++(0.008,0.008) (-0.004,0.743)rectangle++(0.008,0.008) (-0.27,0.757)rectangle++(0.008,0.008) (0.146,0.859)rectangle++(0.008,0.008)
  (0.373,0.757)rectangle++(0.008,0.008) (-0.103,0.898)rectangle++(0.008,0.008) (-0.372,0.793)rectangle++(0.008,0.008) (0.113,0.911)rectangle++(0.008,0.008) (0.045,0.981)rectangle++(0.008,0.008)
  (-0.333,0.883)rectangle++(0.008,0.008) (-0.471,0.704)rectangle++(0.008,0.008) (-0.2,0.893)rectangle++(0.008,0.008) (-0.414,0.756)rectangle++(0.008,0.008) (-0.463,0.617)rectangle++(0.008,0.008)
  (-0.197,0.857)rectangle++(0.008,0.008) (-0.085,0.958)rectangle++(0.008,0.008) (-0.384,0.827)rectangle++(0.008,0.008) (0.119,0.931)rectangle++(0.008,0.008) (0.048,0.989)rectangle++(0.008,0.008)
  (-0.335,0.888)rectangle++(0.008,0.008) (-0.474,0.706)rectangle++(0.008,0.008) (-0.018,0.294)rectangle++(0.008,0.008) (-0.117,0.547)rectangle++(0.008,0.008) (0.141,0.708)rectangle++(0.008,0.008)
  (0.317,0.69)rectangle++(0.008,0.008) (0.131,0.887)rectangle++(0.008,0.008) (0.376,0.775)rectangle++(0.008,0.008) (-0.109,0.908)rectangle++(0.008,0.008) (-0.377,0.795)rectangle++(0.008,0.008)
  (-0.46,0.648)rectangle++(0.008,0.008) (0.02,0.877)rectangle++(0.008,0.008) (0.322,0.81)rectangle++(0.008,0.008) (-0.146,0.904)rectangle++(0.008,0.008) (0.256,0.882)rectangle++(0.008,0.008)
  (0.105,0.968)rectangle++(0.008,0.008) (0.393,0.821)rectangle++(0.008,0.008) (-0.117,0.935)rectangle++(0.008,0.008) (0.279,0.886)rectangle++(0.008,0.008) (-0.192,0.924)rectangle++(0.008,0.008)
  (-0.421,0.773)rectangle++(0.008,0.008) (-0.179,0.922)rectangle++(0.008,0.008) (0.247,0.902)rectangle++(0.008,0.008) (0.434,0.741)rectangle++(0.008,0.008) (0.462,0.6)rectangle++(0.008,0.008)
  (-0.007,0.858)rectangle++(0.008,0.008) (-0.005,0.958)rectangle++(0.008,0.008) (0.339,0.856)rectangle++(0.008,0.008) (0.14,0.957)rectangle++(0.008,0.008) (0.405,0.804)rectangle++(0.008,0.008)
  (-0.106,0.931)rectangle++(0.008,0.008) (-0.047,0.989)rectangle++(0.008,0.008) (-0.378,0.855)rectangle++(0.008,0.008) (0.132,0.942)rectangle++(0.008,0.008) (0.053,0.993)rectangle++(0.008,0.008)
  (0.02,1.015)rectangle++(0.008,0.008) (0.371,0.873)rectangle++(0.008,0.008) (-0.146,0.951)rectangle++(0.008,0.008) (-0.41,0.802)rectangle++(0.008,0.008) (-0.174,0.934)rectangle++(0.008,0.008)
  (-0.417,0.784)rectangle++(0.008,0.008) (-0.474,0.629)rectangle++(0.008,0.008) (0.007,0.873)rectangle++(0.008,0.008) (-0.312,0.821)rectangle++(0.008,0.008) (0.15,0.903)rectangle++(0.008,0.008)
  (-0.258,0.885)rectangle++(0.008,0.008) (0.198,0.913)rectangle++(0.008,0.008) (0.081,0.981)rectangle++(0.008,0.008) (-0.317,0.897)rectangle++(0.008,0.008) (-0.469,0.716)rectangle++(0.008,0.008)
  (-0.473,0.579)rectangle++(0.008,0.008) (-0.427,0.515)rectangle++(0.008,0.008) (-0.382,0.502)rectangle++(0.008,0.008) (-0.358,0.512)rectangle++(0.008,0.008) (0.019,0.778)rectangle++(0.008,0.008)
  (-0.273,0.782)rectangle++(0.008,0.008) (0.154,0.871)rectangle++(0.008,0.008) (-0.244,0.872)rectangle++(0.008,0.008) (-0.427,0.731)rectangle++(0.008,0.008) (-0.182,0.905)rectangle++(0.008,0.008)
  (0.239,0.895)rectangle++(0.008,0.008) (0.428,0.74)rectangle++(0.008,0.008) (-0.072,0.91)rectangle++(0.008,0.008) (0.291,0.858)rectangle++(0.008,0.008) (0.439,0.705)rectangle++(0.008,0.008)
  (0.182,0.894)rectangle++(0.008,0.008) (0.402,0.76)rectangle++(0.008,0.008) (0.008,0.317)rectangle++(0.008,0.008) (-0.113,0.567)rectangle++(0.008,0.008) (-0.05,0.836)rectangle++(0.008,0.008)
  (0.275,0.816)rectangle++(0.008,0.008) (0.416,0.691)rectangle++(0.008,0.008) (-0.06,0.884)rectangle++(0.008,0.008) (-0.028,0.969)rectangle++(0.008,0.008) (-0.014,1.005)rectangle++(0.008,0.008)
  (-0.369,0.873)rectangle++(0.008,0.008) (-0.484,0.687)rectangle++(0.008,0.008) (-0.018,0.286)rectangle++(0.008,0.008) (-0.114,0.543)rectangle++(0.008,0.008) (0.141,0.706)rectangle++(0.008,0.008)
  (0.316,0.688)rectangle++(0.008,0.008) (0.006,0.287)rectangle++(0.008,0.008) (-0.004,0.118)rectangle++(0.008,0.008) (-0.047,0.472)rectangle++(0.008,0.008) (-0.022,0.796)rectangle++(0.008,0.008)
  (-0.012,0.932)rectangle++(0.008,0.008) (-0.007,0.989)rectangle++(0.008,0.008) (0.349,0.871)rectangle++(0.008,0.008) (-0.155,0.942)rectangle++(0.008,0.008) (-0.067,0.993)rectangle++(0.008,0.008)
  (0.323,0.894)rectangle++(0.008,0.008) (0.466,0.71)rectangle++(0.008,0.008) (0.466,0.574)rectangle++(0.008,0.008) (0.193,0.839)rectangle++(0.008,0.008) (0.387,0.731)rectangle++(0.008,0.008)
  (-0.088,0.892)rectangle++(0.008,0.008) (0.278,0.855)rectangle++(0.008,0.008) (0.114,0.957)rectangle++(0.008,0.008) (-0.293,0.897)rectangle++(0.008,0.008) (-0.458,0.725)rectangle++(0.008,0.008)
  (-0.195,0.902)rectangle++(0.008,0.008) (-0.084,0.977)rectangle++(0.008,0.008) (-0.006,0.408)rectangle++(0.008,0.008) (-0.005,0.769)rectangle++(0.008,0.008) (-0.28,0.769)rectangle++(0.008,0.008)
  (-0.406,0.671)rectangle++(0.008,0.008) (-0.016,0.279)rectangle++(0.008,0.008) (0.092,0.55)rectangle++(0.008,0.008) (0.237,0.635)rectangle++(0.008,0.008) (-0.122,0.794)rectangle++(0.008,0.008)
  (0.227,0.823)rectangle++(0.008,0.008) (-0.194,0.876)rectangle++(0.008,0.008) (-0.084,0.966)rectangle++(0.008,0.008) (-0.037,1.003)rectangle++(0.008,0.008) (-0.379,0.864)rectangle++(0.008,0.008)
  (-0.015,0.361)rectangle++(0.008,0.008) (-0.004,0.149)rectangle++(0.008,0.008) (0.051,0.486)rectangle++(0.008,0.008) (-0.154,0.659)rectangle++(0.008,0.008) (-0.067,0.875)rectangle++(0.008,0.008)
  (-0.347,0.795)rectangle++(0.008,0.008) (-0.446,0.659)rectangle++(0.008,0.008) (-0.443,0.561)rectangle++(0.008,0.008) (-0.188,0.833)rectangle++(0.008,0.008) (-0.081,0.948)rectangle++(0.008,0.008)
  (0.3,0.879)rectangle++(0.008,0.008) (-0.18,0.928)rectangle++(0.008,0.008) (0.248,0.905)rectangle++(0.008,0.008) (0.436,0.741)rectangle++(0.008,0.008) (0.463,0.6)rectangle++(0.008,0.008)
  (-0.006,0.859)rectangle++(0.008,0.008) (-0.313,0.81)rectangle++(0.008,0.008) (-0.134,0.938)rectangle++(0.008,0.008) (-0.008,0.392)rectangle++(0.008,0.008) (0.135,0.598)rectangle++(0.008,0.008)
  (0.275,0.641)rectangle++(0.008,0.008) (0.354,0.611)rectangle++(0.008,0.008) (0.379,0.569)rectangle++(0.008,0.008) (0.157,0.837)rectangle++(0.008,0.008) (-0.231,0.858)rectangle++(0.008,0.008)
  (0.2,0.891)rectangle++(0.008,0.008) (0.409,0.752)rectangle++(0.008,0.008) (0.455,0.614)rectangle++(0.008,0.008) (0.01,0.256)rectangle++(0.008,0.008) (0.095,0.53)rectangle++(0.008,0.008)
  (-0.149,0.695)rectangle++(0.008,0.008) (-0.32,0.684)rectangle++(0.008,0.008) (-0.137,0.885)rectangle++(0.008,0.008) (-0.382,0.775)rectangle++(0.008,0.008) (0.101,0.907)rectangle++(0.008,0.008)
  (-0.281,0.87)rectangle++(0.008,0.008) (-0.443,0.716)rectangle++(0.008,0.008) (-0.462,0.589)rectangle++(0.008,0.008) (-0.002,0.85)rectangle++(0.008,0.008) (0.302,0.806)rectangle++(0.008,0.008)
  (0.425,0.677)rectangle++(0.008,0.008) (-0.051,0.88)rectangle++(0.008,0.008) (-0.341,0.803)rectangle++(0.008,0.008) (-0.146,0.935)rectangle++(0.008,0.008) (0.267,0.896)rectangle++(0.008,0.008)
  (0.441,0.731)rectangle++(0.008,0.008) (0.462,0.593)rectangle++(0.008,0.008) (-0.004,0.855)rectangle++(0.008,0.008) (-0.004,0.357)rectangle++(0.008,0.008) (-0.004,0.748)rectangle++(0.008,0.008)
  (-0.004,0.912)rectangle++(0.008,0.008) (0.323,0.834)rectangle++(0.008,0.008) (-0.154,0.916)rectangle++(0.008,0.008) (0.256,0.89)rectangle++(0.008,0.008) (0.434,0.732)rectangle++(0.008,0.008)
  (-0.067,0.909)rectangle++(0.008,0.008) (-0.358,0.811)rectangle++(0.008,0.008) (-0.153,0.938)rectangle++(0.008,0.008) (-0.409,0.793)rectangle++(0.008,0.008) (0.096,0.925)rectangle++(0.008,0.008)
  (-0.001,0.386)rectangle++(0.008,0.008) (-0.004,0.16)rectangle++(0.008,0.008) (-0.004,0.065)rectangle++(0.008,0.008) (-0.029,0.447)rectangle++(0.008,0.008) (-0.014,0.786)rectangle++(0.008,0.008)
  (0.273,0.781)rectangle++(0.008,0.008) (0.403,0.676)rectangle++(0.008,0.008) (-0.061,0.872)rectangle++(0.008,0.008) (0.283,0.837)rectangle++(0.008,0.008) (-0.173,0.903)rectangle++(0.008,0.008)
  (-0.009,0.377)rectangle++(0.008,0.008) (-0.142,0.588)rectangle++(0.008,0.008) (-0.279,0.637)rectangle++(0.008,0.008) (-0.119,0.865)rectangle++(0.008,0.008) (0.254,0.855)rectangle++(0.008,0.008)
  (0.42,0.716)rectangle++(0.008,0.008) (-0.067,0.897)rectangle++(0.008,0.008) (-0.355,0.805)rectangle++(0.008,0.008) (-0.453,0.661)rectangle++(0.008,0.008) (-0.447,0.559)rectangle++(0.008,0.008)
  (-0.19,0.833)rectangle++(0.008,0.008) (0.21,0.865)rectangle++(0.008,0.008) (-0.217,0.889)rectangle++(0.008,0.008) (0.218,0.9)rectangle++(0.008,0.008) (0.089,0.976)rectangle++(0.008,0.008)
  (0.035,1.008)rectangle++(0.008,0.008) (0.375,0.864)rectangle++(0.008,0.008) (-0.141,0.948)rectangle++(0.008,0.008) (-0.008,0.396)rectangle++(0.008,0.008) (-0.004,0.164)rectangle++(0.008,0.008)
  (-0.004,0.667)rectangle++(0.008,0.008) (-0.004,0.878)rectangle++(0.008,0.008) (0.311,0.819)rectangle++(0.008,0.008) (0.434,0.68)rectangle++(0.008,0.008) (-0.048,0.885)rectangle++(0.008,0.008)
  (0.293,0.838)rectangle++(0.008,0.008) (0.005,0.35)rectangle++(0.008,0.008) (-0.004,0.145)rectangle++(0.008,0.008) (0.049,0.484)rectangle++(0.008,0.008) (-0.002,0.201)rectangle++(0.008,0.008)
  (-0.076,0.51)rectangle++(0.008,0.008) (-0.221,0.625)rectangle++(0.008,0.008) (-0.095,0.86)rectangle++(0.008,0.008) (-0.354,0.778)rectangle++(0.008,0.008) (-0.443,0.649)rectangle++(0.008,0.008)
  (0.028,0.871)rectangle++(0.008,0.008) (-0.302,0.827)rectangle++(0.008,0.008) (-0.129,0.945)rectangle++(0.008,0.008) (-0.057,0.995)rectangle++(0.008,0.008) (-0.026,1.015)rectangle++(0.008,0.008)
  (0.35,0.89)rectangle++(0.008,0.008) (-0.161,0.951)rectangle++(0.008,0.008) (-0.417,0.796)rectangle++(0.008,0.008) (0.093,0.929)rectangle++(0.008,0.008) (0.374,0.808)rectangle++(0.008,0.008)
  (0.458,0.652)rectangle++(0.008,0.008) (-0.027,0.881)rectangle++(0.008,0.008) (0.302,0.829)rectangle++(0.008,0.008) (0.124,0.946)rectangle++(0.008,0.008) (-0.284,0.896)rectangle++(0.008,0.008)
  (-0.122,0.974)rectangle++(0.008,0.008) (-0.407,0.821)rectangle++(0.008,0.008) (0.106,0.937)rectangle++(0.008,0.008) (-0.001,0.391)rectangle++(0.008,0.008) (-0.144,0.598)rectangle++(0.008,0.008)
  (-0.008,0.249)rectangle++(0.008,0.008) (-0.096,0.53)rectangle++(0.008,0.008) (0.145,0.693)rectangle++(0.008,0.008) (0.058,0.889)rectangle++(0.008,0.008) (0.022,0.971)rectangle++(0.008,0.008)
  (0.007,1.005)rectangle++(0.008,0.008) (-0.359,0.881)rectangle++(0.008,0.008) (-0.153,0.968)rectangle++(0.008,0.008) (-0.419,0.807)rectangle++(0.008,0.008) (0.096,0.935)rectangle++(0.008,0.008)
  (0.377,0.809)rectangle++(0.008,0.008) (-0.12,0.924)rectangle++(0.008,0.008) (-0.389,0.798)rectangle++(0.008,0.008) (0.107,0.92)rectangle++(0.008,0.008) (-0.283,0.878)rectangle++(0.008,0.008)
  (0.183,0.919)rectangle++(0.008,0.008) (0.075,0.984)rectangle++(0.008,0.008) (0.385,0.839)rectangle++(0.008,0.008) (0.008,0.35)rectangle++(0.008,0.008) (0.001,0.745)rectangle++(0.008,0.008)
  (-0.002,0.91)rectangle++(0.008,0.008) (-0.003,0.98)rectangle++(0.008,0.008) (-0.355,0.866)rectangle++(0.008,0.008) (-0.151,0.961)rectangle++(0.008,0.008) (-0.416,0.805)rectangle++(0.008,0.008)
  (0.096,0.933)rectangle++(0.008,0.008) (-0.293,0.88)rectangle++(0.008,0.008) (-0.452,0.717)rectangle++(0.008,0.008) (-0.466,0.586)rectangle++(0.008,0.008) (-0.005,0.85)rectangle++(0.008,0.008)
  (-0.309,0.806)rectangle++(0.008,0.008) (0.146,0.895)rectangle++(0.008,0.008) (-0.257,0.881)rectangle++(0.008,0.008) (-0.435,0.73)rectangle++(0.008,0.008) (0.061,0.906)rectangle++(0.008,0.008);
\fill[nuage]
  (-0.299,0.855)rectangle++(0.008,0.008) (-0.013,0.357)rectangle++(0.008,0.008) (-0.137,0.578)rectangle++(0.008,0.008) (-0.272,0.634)rectangle++(0.008,0.008) (0.101,0.804)rectangle++(0.008,0.008)
  (-0.001,0.335)rectangle++(0.008,0.008) (-0.003,0.738)rectangle++(0.008,0.008) (-0.004,0.308)rectangle++(0.008,0.008) (-0.004,0.727)rectangle++(0.008,0.008) (-0.004,0.303)rectangle++(0.008,0.008)
  (-0.114,0.556)rectangle++(0.008,0.008) (-0.254,0.633)rectangle++(0.008,0.008) (0.109,0.796)rectangle++(0.008,0.008) (-0.001,0.332)rectangle++(0.008,0.008) (-0.122,0.571)rectangle++(0.008,0.008)
  (-0.054,0.837)rectangle++(0.008,0.008) (-0.025,0.949)rectangle++(0.008,0.008) (-0.005,0.396)rectangle++(0.008,0.008) (-0.004,0.164)rectangle++(0.008,0.008) (-0.064,0.493)rectangle++(0.008,0.008)
  (-0.209,0.622)rectangle++(0.008,0.008) (0.126,0.775)rectangle++(0.008,0.008) (0.051,0.923)rectangle++(0.008,0.008) (0.352,0.82)rectangle++(0.008,0.008) (-0.136,0.92)rectangle++(0.008,0.008)
  (-0.394,0.791)rectangle++(0.008,0.008) (-0.016,0.33)rectangle++(0.008,0.008) (-0.004,0.136)rectangle++(0.008,0.008) (-0.004,0.655)rectangle++(0.008,0.008) (-0.004,0.873)rectangle++(0.008,0.008)
  (-0.317,0.817)rectangle++(0.008,0.008) (-0.44,0.679)rectangle++(0.008,0.008) (-0.187,0.883)rectangle++(0.008,0.008) (-0.404,0.756)rectangle++(0.008,0.008) (-0.458,0.62)rectangle++(0.008,0.008)
  (-0.195,0.858)rectangle++(0.008,0.008) (-0.01,0.358)rectangle++(0.008,0.008) (-0.136,0.579)rectangle++(0.008,0.008) (-0.273,0.636)rectangle++(0.008,0.008) (0.102,0.804)rectangle++(0.008,0.008)
  (0.04,0.935)rectangle++(0.008,0.008) (0.015,0.991)rectangle++(0.008,0.008) (0.36,0.864)rectangle++(0.008,0.008) (0.149,0.961)rectangle++(0.008,0.008) (-0.279,0.911)rectangle++(0.008,0.008)
  (0.197,0.932)rectangle++(0.008,0.008) (0.422,0.772)rectangle++(0.008,0.008) (0.175,0.922)rectangle++(0.008,0.008) (-0.253,0.903)rectangle++(0.008,0.008) (-0.108,0.977)rectangle++(0.008,0.008)
  (-0.402,0.827)rectangle++(0.008,0.008) (0.111,0.938)rectangle++(0.008,0.008) (0.044,0.992)rectangle++(0.008,0.008) (0.374,0.854)rectangle++(0.008,0.008) (0.475,0.673)rectangle++(0.008,0.008)
  (-0.027,0.897)rectangle++(0.008,0.008) (-0.336,0.819)rectangle++(0.008,0.008) (0.138,0.911)rectangle++(0.008,0.008) (0.056,0.98)rectangle++(0.008,0.008) (-0.328,0.887)rectangle++(0.008,0.008)
  (0.166,0.939)rectangle++(0.008,0.008) (0.41,0.786)rectangle++(0.008,0.008) (0.468,0.629)rectangle++(0.008,0.008) (-0.015,0.874)rectangle++(0.008,0.008) (-0.004,0.365)rectangle++(0.008,0.008)
  (-0.136,0.584)rectangle++(0.008,0.008) (-0.274,0.638)rectangle++(0.008,0.008) (0.102,0.806)rectangle++(0.008,0.008) (0.333,0.748)rectangle++(0.008,0.008) (0.006,0.312)rectangle++(0.008,0.008)
  (0,0.729)rectangle++(0.008,0.008) (-0.002,0.904)rectangle++(0.008,0.008) (-0.327,0.832)rectangle++(0.008,0.008) (-0.45,0.682)rectangle++(0.008,0.008) (-0.191,0.884)rectangle++(0.008,0.008)
  (-0.01,0.369)rectangle++(0.008,0.008) (-0.14,0.585)rectangle++(0.008,0.008) (-0.276,0.636)rectangle++(0.008,0.008) (-0.118,0.865)rectangle++(0.008,0.008) (-0.007,0.361)rectangle++(0.008,0.008)
  (0.125,0.584)rectangle++(0.008,0.008) (0.265,0.639)rectangle++(0.008,0.008) (0.348,0.614)rectangle++(0.008,0.008) (0.378,0.573)rectangle++(0.008,0.008) (0.376,0.543)rectangle++(0.008,0.008)
  (0.365,0.53)rectangle++(0.008,0.008) (0.356,0.528)rectangle++(0.008,0.008) (-0.03,0.788)rectangle++(0.008,0.008) (0.267,0.787)rectangle++(0.008,0.008) (0.11,0.928)rectangle++(0.008,0.008)
  (-0.001,0.388)rectangle++(0.008,0.008) (0.137,0.594)rectangle++(0.008,0.008) (-0.153,0.74)rectangle++(0.008,0.008) (-0.338,0.703)rectangle++(0.008,0.008) (-0.409,0.62)rectangle++(0.008,0.008)
  (-0.016,0.258)rectangle++(0.008,0.008) (-0.009,0.706)rectangle++(0.008,0.008) (0.247,0.742)rectangle++(0.008,0.008) (0.377,0.667)rectangle++(0.008,0.008) (0.156,0.878)rectangle++(0.008,0.008)
  (0.063,0.966)rectangle++(0.008,0.008) (-0.32,0.884)rectangle++(0.008,0.008) (-0.013,0.369)rectangle++(0.008,0.008) (0.125,0.59)rectangle++(0.008,0.008) (-0.157,0.733)rectangle++(0.008,0.008)
  (0.189,0.808)rectangle++(0.008,0.008) (0.002,0.337)rectangle++(0.008,0.008) (-0.002,0.739)rectangle++(0.008,0.008) (-0.003,0.908)rectangle++(0.008,0.008) (-0.329,0.833)rectangle++(0.008,0.008)
  (-0.141,0.948)rectangle++(0.008,0.008) (-0.008,0.396)rectangle++(0.008,0.008) (-0.149,0.597)rectangle++(0.008,0.008) (0.145,0.742)rectangle++(0.008,0.008) (-0.203,0.81)rectangle++(0.008,0.008)
  (-0.386,0.717)rectangle++(0.008,0.008) (-0.436,0.609)rectangle++(0.008,0.008) (-0.185,0.854)rectangle++(0.008,0.008) (-0.08,0.956)rectangle++(0.008,0.008) (-0.036,0.999)rectangle++(0.008,0.008)
  (-0.377,0.863)rectangle++(0.008,0.008) (-0.484,0.679)rectangle++(0.008,0.008) (-0.206,0.883)rectangle++(0.008,0.008) (-0.413,0.749)rectangle++(0.008,0.008) (-0.46,0.614)rectangle++(0.008,0.008)
  (0.008,0.861)rectangle++(0.008,0.008) (-0.307,0.816)rectangle++(0.008,0.008) (0.15,0.899)rectangle++(0.008,0.008) (-0.256,0.884)rectangle++(0.008,0.008) (-0.436,0.732)rectangle++(0.008,0.008)
  (0.061,0.907)rectangle++(0.008,0.008) (0.023,0.978)rectangle++(0.008,0.008) (-0.342,0.875)rectangle++(0.008,0.008) (-0.472,0.697)rectangle++(0.008,0.008) (0.032,0.904)rectangle++(0.008,0.008)
  (-0.312,0.844)rectangle++(0.008,0.008) (-0.133,0.952)rectangle++(0.008,0.008) (0.278,0.899)rectangle++(0.008,0.008) (0.004,0.375)rectangle++(0.008,0.008) (-0.136,0.592)rectangle++(0.008,0.008)
  (-0.277,0.642)rectangle++(0.008,0.008) (-0.359,0.614)rectangle++(0.008,0.008) (0.054,0.825)rectangle++(0.008,0.008) (0.02,0.944)rectangle++(0.008,0.008) (-0.331,0.858)rectangle++(0.008,0.008)
  (0.154,0.927)rectangle++(0.008,0.008) (0.062,0.987)rectangle++(0.008,0.008) (0.38,0.845)rectangle++(0.008,0.008) (-0.132,0.941)rectangle++(0.008,0.008) (-0.058,0.993)rectangle++(0.008,0.008)
  (0.328,0.891)rectangle++(0.008,0.008) (-0.172,0.943)rectangle++(0.008,0.008) (-0.075,0.994)rectangle++(0.008,0.008) (-0.393,0.847)rectangle++(0.008,0.008) (-0.485,0.666)rectangle++(0.008,0.008)
  (-0.463,0.55)rectangle++(0.008,0.008) (-0.412,0.505)rectangle++(0.008,0.008) (-0.175,0.81)rectangle++(0.008,0.008) (-0.373,0.727)rectangle++(0.008,0.008) (-0.015,0.303)rectangle++(0.008,0.008)
  (-0.119,0.552)rectangle++(0.008,0.008) (-0.052,0.83)rectangle++(0.008,0.008) (0.272,0.814)rectangle++(0.008,0.008) (0.112,0.94)rectangle++(0.008,0.008) (-0.288,0.889)rectangle++(0.008,0.008)
  (-0.453,0.723)rectangle++(0.008,0.008) (-0.468,0.588)rectangle++(0.008,0.008) (-0.005,0.852)rectangle++(0.008,0.008) (0.301,0.808)rectangle++(0.008,0.008) (0.124,0.937)rectangle++(0.008,0.008)
  (0.05,0.991)rectangle++(0.008,0.008) (0.376,0.852)rectangle++(0.008,0.008) (0.007,0.355)rectangle++(0.008,0.008) (0.001,0.747)rectangle++(0.008,0.008) (-0.002,0.911)rectangle++(0.008,0.008)
  (0.324,0.834)rectangle++(0.008,0.008) (0.445,0.682)rectangle++(0.008,0.008) (0.185,0.884)rectangle++(0.008,0.008) (-0.235,0.889)rectangle++(0.008,0.008) (-0.011,0.371)rectangle++(0.008,0.008)
  (-0.007,0.754)rectangle++(0.008,0.008) (0.265,0.763)rectangle++(0.008,0.008) (-0.155,0.863)rectangle++(0.008,0.008) (-0.067,0.96)rectangle++(0.008,0.008) (-0.031,1.001)rectangle++(0.008,0.008)
  (0.343,0.885)rectangle++(0.008,0.008) (0.472,0.698)rectangle++(0.008,0.008) (0.196,0.891)rectangle++(0.008,0.008) (0.407,0.754)rectangle++(0.008,0.008) (0.454,0.616)rectangle++(0.008,0.008)
  (0.189,0.856)rectangle++(0.008,0.008) (-0.223,0.878)rectangle++(0.008,0.008) (0.211,0.897)rectangle++(0.008,0.008) (-0.228,0.905)rectangle++(0.008,0.008) (-0.431,0.751)rectangle++(0.008,0.008)
  (-0.183,0.913)rectangle++(0.008,0.008) (-0.009,0.381)rectangle++(0.008,0.008) (-0.006,0.758)rectangle++(0.008,0.008) (0.267,0.765)rectangle++(0.008,0.008) (0.11,0.919)rectangle++(0.008,0.008)
  (0.044,0.984)rectangle++(0.008,0.008) (0.371,0.85)rectangle++(0.008,0.008) (-0.138,0.94)rectangle++(0.008,0.008) (-0.008,0.393)rectangle++(0.008,0.008) (-0.006,0.763)rectangle++(0.008,0.008)
  (-0.005,0.918)rectangle++(0.008,0.008) (-0.334,0.837)rectangle++(0.008,0.008) (-0.455,0.683)rectangle++(0.008,0.008) (-0.193,0.884)rectangle++(0.008,0.008) (-0.408,0.754)rectangle++(0.008,0.008)
  (-0.459,0.618)rectangle++(0.008,0.008) (-0.434,0.538)rectangle++(0.008,0.008) (-0.007,0.817)rectangle++(0.008,0.008) (-0.005,0.941)rectangle++(0.008,0.008) (-0.342,0.847)rectangle++(0.008,0.008)
  (0.146,0.926)rectangle++(0.008,0.008) (0.396,0.788)rectangle++(0.008,0.008) (0.462,0.635)rectangle++(0.008,0.008) (-0.019,0.874)rectangle++(0.008,0.008) (-0.01,0.965)rectangle++(0.008,0.008)
  (-0.353,0.857)rectangle++(0.008,0.008) (-0.471,0.685)rectangle++(0.008,0.008) (0.029,0.897)rectangle++(0.008,0.008) (0.01,0.975)rectangle++(0.008,0.008) (0.352,0.858)rectangle++(0.008,0.008)
  (0.145,0.958)rectangle++(0.008,0.008) (0.059,1)rectangle++(0.008,0.008) (0.383,0.853)rectangle++(0.008,0.008) (0.159,0.956)rectangle++(0.008,0.008) (0.064,0.999)rectangle++(0.008,0.008)
  (0.025,1.017)rectangle++(0.008,0.008) (-0.356,0.893)rectangle++(0.008,0.008) (-0.485,0.7)rectangle++(0.008,0.008) (0.028,0.91)rectangle++(0.008,0.008) (0.337,0.822)rectangle++(0.008,0.008)
  (0.447,0.672)rectangle++(0.008,0.008) (0.185,0.88)rectangle++(0.008,0.008) (0.398,0.752)rectangle++(0.008,0.008) (0.165,0.914)rectangle++(0.008,0.008) (0.067,0.981)rectangle++(0.008,0.008)
  (-0.002,0.41)rectangle++(0.008,0.008) (-0.151,0.606)rectangle++(0.008,0.008) (-0.066,0.852)rectangle++(0.008,0.008) (-0.338,0.785)rectangle++(0.008,0.008) (-0.438,0.658)rectangle++(0.008,0.008)
  (-0.439,0.563)rectangle++(0.008,0.008) (0,0.83)rectangle++(0.008,0.008) (0.296,0.796)rectangle++(0.008,0.008) (0.122,0.932)rectangle++(0.008,0.008) (0.049,0.989)rectangle++(0.008,0.008)
  (0.018,1.013)rectangle++(0.008,0.008) (-0.357,0.889)rectangle++(0.008,0.008) (0.153,0.95)rectangle++(0.008,0.008) (0.062,0.997)rectangle++(0.008,0.008) (-0.331,0.897)rectangle++(0.008,0.008)
  (0.168,0.945)rectangle++(0.008,0.008) (0.068,0.994)rectangle++(0.008,0.008) (0.386,0.846)rectangle++(0.008,0.008) (0.008,0.353)rectangle++(0.008,0.008) (0.129,0.575)rectangle++(0.008,0.008)
  (-0.15,0.728)rectangle++(0.008,0.008) (-0.008,0.303)rectangle++(0.008,0.008) (-0.006,0.725)rectangle++(0.008,0.008) (0.256,0.75)rectangle++(0.008,0.008) (-0.155,0.853)rectangle++(0.008,0.008)
  (-0.067,0.956)rectangle++(0.008,0.008) (-0.376,0.832)rectangle++(0.008,0.008) (-0.015,0.347)rectangle++(0.008,0.008) (-0.004,0.144)rectangle++(0.008,0.008) (0.049,0.484)rectangle++(0.008,0.008)
  (-0.154,0.658)rectangle++(0.008,0.008) (0.164,0.772)rectangle++(0.008,0.008) (0.066,0.922)rectangle++(0.008,0.008) (0.359,0.814)rectangle++(0.008,0.008) (0.148,0.94)rectangle++(0.008,0.008)
  (-0.271,0.902)rectangle++(0.008,0.008) (0.197,0.925)rectangle++(0.008,0.008) (-0.244,0.913)rectangle++(0.008,0.008) (-0.105,0.981)rectangle++(0.008,0.008) (0.302,0.902)rectangle++(0.008,0.008)
  (-0.188,0.939)rectangle++(0.008,0.008) (-0.425,0.781)rectangle++(0.008,0.008) (-0.477,0.625)rectangle++(0.008,0.008) (0.004,0.872)rectangle++(0.008,0.008) (0.313,0.813)rectangle++(0.008,0.008)
  (-0.151,0.903)rectangle++(0.008,0.008) (-0.395,0.778)rectangle++(0.008,0.008) (-0.462,0.634)rectangle++(0.008,0.008) (-0.441,0.544)rectangle++(0.008,0.008) (-0.399,0.51)rectangle++(0.008,0.008)
  (-0.001,0.792)rectangle++(0.008,0.008) (0.282,0.779)rectangle++(0.008,0.008) (0.005,0.325)rectangle++(0.008,0.008) (0.117,0.563)rectangle++(0.008,0.008) (-0.151,0.719)rectangle++(0.008,0.008)
  (0.187,0.799)rectangle++(0.008,0.008) (0.076,0.933)rectangle++(0.008,0.008) (-0.302,0.873)rectangle++(0.008,0.008) (-0.129,0.964)rectangle++(0.008,0.008) (-0.008,0.403)rectangle++(0.008,0.008)
  (-0.151,0.601)rectangle++(0.008,0.008) (-0.008,0.25)rectangle++(0.008,0.008) (-0.097,0.53)rectangle++(0.008,0.008) (0.144,0.693)rectangle++(0.008,0.008) (0.058,0.889)rectangle++(0.008,0.008)
  (0.022,0.971)rectangle++(0.008,0.008) (-0.34,0.871)rectangle++(0.008,0.008) (0.155,0.936)rectangle++(0.008,0.008) (0.404,0.789)rectangle++(0.008,0.008) (-0.101,0.924)rectangle++(0.008,0.008)
  (-0.007,0.386)rectangle++(0.008,0.008) (-0.005,0.76)rectangle++(0.008,0.008) (-0.005,0.917)rectangle++(0.008,0.008) (-0.333,0.837)rectangle++(0.008,0.008) (-0.455,0.683)rectangle++(0.008,0.008)
  (-0.193,0.884)rectangle++(0.008,0.008) (-0.408,0.754)rectangle++(0.008,0.008) (-0.016,0.315)rectangle++(0.008,0.008) (-0.009,0.73)rectangle++(0.008,0.008) (-0.268,0.75)rectangle++(0.008,0.008)
  (-0.394,0.666)rectangle++(0.008,0.008) (0.057,0.861)rectangle++(0.008,0.008) (0.333,0.79)rectangle++(0.008,0.008) (0.433,0.658)rectangle++(0.008,0.008) (-0.041,0.875)rectangle++(0.008,0.008)
  (-0.019,0.965)rectangle++(0.008,0.008) (-0.357,0.853)rectangle++(0.008,0.008) (-0.472,0.682)rectangle++(0.008,0.008) (-0.2,0.884)rectangle++(0.008,0.008) (-0.411,0.752)rectangle++(0.008,0.008)
  (-0.46,0.616)rectangle++(0.008,0.008) (-0.434,0.537)rectangle++(0.008,0.008) (-0.007,0.817)rectangle++(0.008,0.008) (0.287,0.792)rectangle++(0.008,0.008) (-0.155,0.884)rectangle++(0.008,0.008)
  (-0.067,0.969)rectangle++(0.008,0.008) (0.314,0.883)rectangle++(0.008,0.008) (0.13,0.969)rectangle++(0.008,0.008) (0.404,0.813)rectangle++(0.008,0.008) (-0.109,0.935)rectangle++(0.008,0.008)
  (0.283,0.883)rectangle++(0.008,0.008) (0.005,0.368)rectangle++(0.008,0.008) (-0.133,0.589)rectangle++(0.008,0.008) (-0.275,0.641)rectangle++(0.008,0.008) (-0.358,0.614)rectangle++(0.008,0.008)
  (-0.387,0.572)rectangle++(0.008,0.008) (-0.385,0.543)rectangle++(0.008,0.008) (-0.164,0.826)rectangle++(0.008,0.008) (-0.009,0.344)rectangle++(0.008,0.008) (-0.131,0.574)rectangle++(0.008,0.008)
  (-0.008,0.239)rectangle++(0.008,0.008) (0.081,0.528)rectangle++(0.008,0.008) (0.032,0.82)rectangle++(0.008,0.008) (0.306,0.78)rectangle++(0.008,0.008) (-0.142,0.885)rectangle++(0.008,0.008)
  (-0.062,0.969)rectangle++(0.008,0.008) (0.317,0.882)rectangle++(0.008,0.008) (0.459,0.706)rectangle++(0.008,0.008) (0.461,0.575)rectangle++(0.008,0.008) (0.002,0.847)rectangle++(0.008,0.008)
  (-0.004,0.353)rectangle++(0.008,0.008) (-0.132,0.579)rectangle++(0.008,0.008) (-0.271,0.637)rectangle++(0.008,0.008) (0.103,0.804)rectangle++(0.008,0.008) (0.041,0.935)rectangle++(0.008,0.008)
  (-0.003,0.391)rectangle++(0.008,0.008) (0.138,0.596)rectangle++(0.008,0.008) (0.275,0.64)rectangle++(0.008,0.008) (0.353,0.611)rectangle++(0.008,0.008) (-0.06,0.824)rectangle++(0.008,0.008)
  (-0.028,0.944)rectangle++(0.008,0.008) (-0.014,0.994)rectangle++(0.008,0.008) (-0.008,1.015)rectangle++(0.008,0.008) (-0.006,1.024)rectangle++(0.008,0.008) (-0.005,1.028)rectangle++(0.008,0.008)
  (-0.004,1.029)rectangle++(0.008,0.008) (-0.004,1.03)rectangle++(0.008,0.008) (0.365,0.889)rectangle++(0.008,0.008) (0.007,0.371)rectangle++(0.008,0.008) (-0.133,0.591)rectangle++(0.008,0.008)
  (-0.275,0.642)rectangle++(0.008,0.008) (-0.118,0.867)rectangle++(0.008,0.008) (-0.007,0.362)rectangle++(0.008,0.008) (0.125,0.584)rectangle++(0.008,0.008) (-0.155,0.731)rectangle++(0.008,0.008)
  (-0.336,0.698)rectangle++(0.008,0.008) (-0.406,0.618)rectangle++(0.008,0.008) (-0.41,0.557)rectangle++(0.008,0.008) (0.011,0.817)rectangle++(0.008,0.008) (-0.291,0.797)rectangle++(0.008,0.008)
  (-0.421,0.68)rectangle++(0.008,0.008) (-0.179,0.883)rectangle++(0.008,0.008) (-0.401,0.759)rectangle++(0.008,0.008) (-0.171,0.916)rectangle++(0.008,0.008) (-0.409,0.777)rectangle++(0.008,0.008)
  (-0.468,0.628)rectangle++(0.008,0.008) (-0.442,0.539)rectangle++(0.008,0.008) (-0.398,0.508)rectangle++(0.008,0.008) (-0.001,0.791)rectangle++(0.008,0.008) (-0.003,0.93)rectangle++(0.008,0.008)
  (0.33,0.842)rectangle++(0.008,0.008) (0.136,0.951)rectangle++(0.008,0.008) (-0.281,0.903)rectangle++(0.008,0.008) (-0.012,0.377)rectangle++(0.008,0.008) (-0.007,0.756)rectangle++(0.008,0.008)
  (-0.277,0.762)rectangle++(0.008,0.008) (-0.402,0.669)rectangle++(0.008,0.008) (0.054,0.866)rectangle++(0.008,0.008) (0.02,0.961)rectangle++(0.008,0.008) (0.352,0.848)rectangle++(0.008,0.008)
  (0.463,0.679)rectangle++(0.008,0.008) (-0.034,0.895)rectangle++(0.008,0.008) (-0.005,0.373)rectangle++(0.008,0.008) (-0.139,0.588)rectangle++(0.008,0.008) (-0.277,0.638)rectangle++(0.008,0.008)
  (0.1,0.807)rectangle++(0.008,0.008) (-0.246,0.824)rectangle++(0.008,0.008) (-0.41,0.708)rectangle++(0.008,0.008) (-0.444,0.596)rectangle++(0.008,0.008) (-0.017,0.248)rectangle++(0.008,0.008)
  (-0.01,0.702)rectangle++(0.008,0.008) (-0.259,0.737)rectangle++(0.008,0.008) (0.144,0.845)rectangle++(0.008,0.008) (0.367,0.751)rectangle++(0.008,0.008) (0.435,0.629)rectangle++(0.008,0.008)
  (0.423,0.548)rectangle++(0.008,0.008) (-0.006,0.821)rectangle++(0.008,0.008) (0.29,0.794)rectangle++(0.008,0.008) (-0.155,0.885)rectangle++(0.008,0.008) (-0.39,0.769)rectangle++(0.008,0.008)
  (-0.166,0.921)rectangle++(0.008,0.008) (-0.072,0.984)rectangle++(0.008,0.008) (-0.033,1.011)rectangle++(0.008,0.008) (-0.016,1.022)rectangle++(0.008,0.008) (0.357,0.889)rectangle++(0.008,0.008)
  (0.148,0.971)rectangle++(0.008,0.008) (0.414,0.808)rectangle++(0.008,0.008) (0.171,0.937)rectangle++(0.008,0.008) (-0.26,0.909)rectangle++(0.008,0.008) (-0.447,0.742)rectangle++(0.008,0.008)
  (-0.017,0.309)rectangle++(0.008,0.008) (-0.122,0.554)rectangle++(0.008,0.008) (-0.054,0.831)rectangle++(0.008,0.008) (0.271,0.815)rectangle++(0.008,0.008) (-0.171,0.889)rectangle++(0.008,0.008)
  (0.239,0.884)rectangle++(0.008,0.008) (-0.21,0.908)rectangle++(0.008,0.008) (0.228,0.907)rectangle++(0.008,0.008) (0.427,0.749)rectangle++(0.008,0.008) (0.462,0.606)rectangle++(0.008,0.008)
  (0.427,0.529)rectangle++(0.008,0.008) (0.003,0.813)rectangle++(0.008,0.008) (-0.293,0.792)rectangle++(0.008,0.008) (-0.013,0.33)rectangle++(0.008,0.008) (-0.127,0.566)rectangle++(0.008,0.008)
  (-0.264,0.632)rectangle++(0.008,0.008) (-0.113,0.863)rectangle++(0.008,0.008) (-0.05,0.96)rectangle++(0.008,0.008) (-0.369,0.84)rectangle++(0.008,0.008) (0.131,0.932)rectangle++(0.008,0.008)
  (-0.277,0.892)rectangle++(0.008,0.008) (-0.449,0.728)rectangle++(0.008,0.008) (-0.191,0.903)rectangle++(0.008,0.008) (0.235,0.897)rectangle++(0.008,0.008) (0.427,0.743)rectangle++(0.008,0.008)
  (0.177,0.91)rectangle++(0.008,0.008) (-0.248,0.898)rectangle++(0.008,0.008) (-0.437,0.741)rectangle++(0.008,0.008) (-0.468,0.602)rectangle++(0.008,0.008) (-0.433,0.527)rectangle++(0.008,0.008)
  (-0.01,0.812)rectangle++(0.008,0.008) (-0.004,0.339)rectangle++(0.008,0.008) (0.118,0.573)rectangle++(0.008,0.008) (-0.154,0.723)rectangle++(0.008,0.008) (0.187,0.802)rectangle++(0.008,0.008)
  (0.002,0.334)rectangle++(0.008,0.008) (0.119,0.569)rectangle++(0.008,0.008) (0.257,0.634)rectangle++(0.008,0.008) (0.106,0.864)rectangle++(0.008,0.008) (-0.264,0.852)rectangle++(0.008,0.008)
  (-0.428,0.714)rectangle++(0.008,0.008) (0.059,0.896)rectangle++(0.008,0.008) (-0.297,0.85)rectangle++(0.008,0.008) (-0.127,0.955)rectangle++(0.008,0.008) (0.282,0.898)rectangle++(0.008,0.008)
  (0.116,0.975)rectangle++(0.008,0.008) (0.4,0.82)rectangle++(0.008,0.008) (0.475,0.648)rectangle++(0.008,0.008) (-0.018,0.885)rectangle++(0.008,0.008) (0.307,0.827)rectangle++(0.008,0.008)
  (-0.159,0.907)rectangle++(0.008,0.008) (0.251,0.888)rectangle++(0.008,0.008) (0.431,0.733)rectangle++(0.008,0.008) (0.179,0.905)rectangle++(0.008,0.008) (0.404,0.766)rectangle++(0.008,0.008)
  (0.458,0.622)rectangle++(0.008,0.008) (-0.017,0.867)rectangle++(0.008,0.008) (-0.009,0.962)rectangle++(0.008,0.008) (-0.006,1.002)rectangle++(0.008,0.008) (-0.364,0.875)rectangle++(0.008,0.008)
  (0.145,0.946)rectangle++(0.008,0.008) (0.404,0.797)rectangle++(0.008,0.008) (-0.104,0.928)rectangle++(0.008,0.008) (-0.382,0.806)rectangle++(0.008,0.008) (-0.015,0.336)rectangle++(0.008,0.008)
  (-0.131,0.567)rectangle++(0.008,0.008) (-0.266,0.632)rectangle++(0.008,0.008) (-0.351,0.613)rectangle++(0.008,0.008) (-0.383,0.574)rectangle++(0.008,0.008) (-0.163,0.839)rectangle++(0.008,0.008)
  (-0.378,0.744)rectangle++(0.008,0.008) (-0.442,0.625)rectangle++(0.008,0.008) (-0.429,0.547)rectangle++(0.008,0.008) (-0.395,0.516)rectangle++(0.008,0.008) (-0.016,0.214)rectangle++(0.008,0.008)
  (-0.087,0.512)rectangle++(0.008,0.008) (-0.039,0.813)rectangle++(0.008,0.008) (0.272,0.802)rectangle++(0.008,0.008) (0.41,0.686)rectangle++(0.008,0.008) (0.431,0.584)rectangle++(0.008,0.008)
  (0.405,0.529)rectangle++(0.008,0.008) (-0.008,0.806)rectangle++(0.008,0.008) (-0.004,0.336)rectangle++(0.008,0.008) (0.117,0.571)rectangle++(0.008,0.008) (-0.154,0.722)rectangle++(0.008,0.008)
  (-0.009,0.301)rectangle++(0.008,0.008) (-0.006,0.724)rectangle++(0.008,0.008) (0.255,0.749)rectangle++(0.008,0.008) (0.383,0.668)rectangle++(0.008,0.008) (-0.067,0.861)rectangle++(0.008,0.008)
  (-0.342,0.789)rectangle++(0.008,0.008) (-0.146,0.929)rectangle++(0.008,0.008) (-0.402,0.792)rectangle++(0.008,0.008) (-0.47,0.638)rectangle++(0.008,0.008) (-0.2,0.865)rectangle++(0.008,0.008)
  (0.217,0.883)rectangle++(0.008,0.008) (0.089,0.969)rectangle++(0.008,0.008) (0.035,1.005)rectangle++(0.008,0.008) (-0.346,0.891)rectangle++(0.008,0.008) (-0.48,0.703)rectangle++(0.008,0.008)
  (-0.474,0.569)rectangle++(0.008,0.008) (-0.201,0.837)rectangle++(0.008,0.008) (-0.087,0.949)rectangle++(0.008,0.008) (-0.382,0.822)rectangle++(0.008,0.008) (0.118,0.929)rectangle++(0.008,0.008)
  (0,0.388)rectangle++(0.008,0.008) (0.138,0.594)rectangle++(0.008,0.008) (0.055,0.847)rectangle++(0.008,0.008) (0.327,0.784)rectangle++(0.008,0.008) (-0.134,0.894)rectangle++(0.008,0.008)
  (-0.384,0.78)rectangle++(0.008,0.008) (-0.164,0.925)rectangle++(0.008,0.008) (0.255,0.898)rectangle++(0.008,0.008) (0.004,0.375)rectangle++(0.008,0.008) (-0.136,0.592)rectangle++(0.008,0.008)
  (-0.008,0.246)rectangle++(0.008,0.008) (0.084,0.532)rectangle++(0.008,0.008) (-0.155,0.692)rectangle++(0.008,0.008) (-0.322,0.68)rectangle++(0.008,0.008) (-0.394,0.615)rectangle++(0.008,0.008)
  (-0.168,0.856)rectangle++(0.008,0.008) (0.228,0.868)rectangle++(0.008,0.008) (0.003,0.362)rectangle++(0.008,0.008) (0.13,0.581)rectangle++(0.008,0.008) (-0.152,0.731)rectangle++(0.008,0.008)
  (-0.066,0.905)rectangle++(0.008,0.008) (0.292,0.853)rectangle++(0.008,0.008) (0.438,0.702)rectangle++(0.008,0.008) (0.45,0.581)rectangle++(0.008,0.008) (0.412,0.521)rectangle++(0.008,0.008)
  (-0.001,0.805)rectangle++(0.008,0.008) (0.286,0.785)rectangle++(0.008,0.008) (-0.153,0.88)rectangle++(0.008,0.008) (-0.388,0.767)rectangle++(0.008,0.008) (0.096,0.905)rectangle++(0.008,0.008)
  (0.038,0.978)rectangle++(0.008,0.008) (0.366,0.85)rectangle++(0.008,0.008) (-0.14,0.938)rectangle++(0.008,0.008) (-0.403,0.798)rectangle++(0.008,0.008) (-0.473,0.64)rectangle++(0.008,0.008)
  (0.012,0.878)rectangle++(0.008,0.008) (-0.312,0.825)rectangle++(0.008,0.008) (-0.441,0.685)rectangle++(0.008,0.008) (0.042,0.887)rectangle++(0.008,0.008) (0.335,0.807)rectangle++(0.008,0.008)
  (-0.138,0.908)rectangle++(0.008,0.008) (-0.008,0.379)rectangle++(0.008,0.008) (0.131,0.592)rectangle++(0.008,0.008) (-0.155,0.737)rectangle++(0.008,0.008) (0.191,0.809)rectangle++(0.008,0.008)
  (-0.205,0.857)rectangle++(0.008,0.008) (-0.403,0.738)rectangle++(0.008,0.008) (-0.451,0.612)rectangle++(0.008,0.008) (-0.192,0.855)rectangle++(0.008,0.008) (-0.397,0.741)rectangle++(0.008,0.008)
  (-0.169,0.909)rectangle++(0.008,0.008) (-0.405,0.774)rectangle++(0.008,0.008) (-0.465,0.628)rectangle++(0.008,0.008) (0.011,0.87)rectangle++(0.008,0.008) (0.002,0.963)rectangle++(0.008,0.008)
  (0.344,0.856)rectangle++(0.008,0.008) (-0.152,0.933)rectangle++(0.008,0.008) (-0.066,0.99)rectangle++(0.008,0.008) (-0.387,0.848)rectangle++(0.008,0.008) (-0.165,0.954)rectangle++(0.008,0.008)
  (0.264,0.911)rectangle++(0.008,0.008) (0.109,0.98)rectangle++(0.008,0.008) (0.043,1.009)rectangle++(0.008,0.008) (0.38,0.862)rectangle++(0.008,0.008) (0.481,0.675)rectangle++(0.008,0.008)
  (0.2,0.881)rectangle++(0.008,0.008) (0.081,0.968)rectangle++(0.008,0.008) (-0.312,0.891)rectangle++(0.008,0.008) (-0.133,0.972)rectangle++(0.008,0.008) (0.285,0.908)rectangle++(0.008,0.008)
  (0.454,0.73)rectangle++(0.008,0.008) (-0.057,0.915)rectangle++(0.008,0.008) (0.3,0.855)rectangle++(0.008,0.008) (0.442,0.7)rectangle++(0.008,0.008) (-0.052,0.897)rectangle++(0.008,0.008)
  (0.296,0.845)rectangle++(0.008,0.008) (-0.17,0.911)rectangle++(0.008,0.008) (0.247,0.893)rectangle++(0.008,0.008) (0.101,0.973)rectangle++(0.008,0.008) (0.393,0.825)rectangle++(0.008,0.008)
  (0.008,0.344)rectangle++(0.008,0.008) (0.001,0.742)rectangle++(0.008,0.008) (-0.004,0.309)rectangle++(0.008,0.008) (-0.116,0.559)rectangle++(0.008,0.008) (-0.007,0.233)rectangle++(0.008,0.008)
  (-0.005,0.695)rectangle++(0.008,0.008) (-0.254,0.735)rectangle++(0.008,0.008) (-0.012,0.306)rectangle++(0.008,0.008) (-0.118,0.555)rectangle++(0.008,0.008) (0.143,0.712)rectangle++(0.008,0.008)
  (-0.192,0.796)rectangle++(0.008,0.008) (-0.01,0.332)rectangle++(0.008,0.008) (-0.127,0.568)rectangle++(0.008,0.008) (0.144,0.721)rectangle++(0.008,0.008) (0.058,0.901)rectangle++(0.008,0.008)
  (0.022,0.976)rectangle++(0.008,0.008) (0.007,1.008)rectangle++(0.008,0.008) (0.362,0.874)rectangle++(0.008,0.008) (-0.15,0.948)rectangle++(0.008,0.008) (-0.411,0.799)rectangle++(0.008,0.008)
  (-0.175,0.933)rectangle++(0.008,0.008) (0.253,0.905)rectangle++(0.008,0.008) (0.438,0.74)rectangle++(0.008,0.008) (0.182,0.908)rectangle++(0.008,0.008) (-0.245,0.899)rectangle++(0.008,0.008)
  (-0.437,0.743)rectangle++(0.008,0.008) (0.065,0.912)rectangle++(0.008,0.008) (-0.3,0.859)rectangle++(0.008,0.008) (-0.128,0.959)rectangle++(0.008,0.008) (-0.056,1)rectangle++(0.008,0.008)
  (-0.386,0.856)rectangle++(0.008,0.008) (0.128,0.946)rectangle++(0.008,0.008) (-0.283,0.897)rectangle++(0.008,0.008) (-0.453,0.728)rectangle++(0.008,0.008) (-0.017,0.304)rectangle++(0.008,0.008)
  (-0.12,0.552)rectangle++(0.008,0.008) (-0.255,0.629)rectangle++(0.008,0.008) (-0.11,0.862)rectangle++(0.008,0.008) (0.257,0.849)rectangle++(0.008,0.008) (0.106,0.954)rectangle++(0.008,0.008)
  (0.388,0.815)rectangle++(0.008,0.008) (-0.117,0.93)rectangle++(0.008,0.008) (0.278,0.883)rectangle++(0.008,0.008) (0.114,0.969)rectangle++(0.008,0.008) (-0.297,0.903)rectangle++(0.008,0.008)
  (-0.462,0.726)rectangle++(0.008,0.008) (-0.474,0.586)rectangle++(0.008,0.008) (-0.008,0.853)rectangle++(0.008,0.008) (-0.006,0.956)rectangle++(0.008,0.008) (-0.005,0.999)rectangle++(0.008,0.008)
  (-0.004,1.017)rectangle++(0.008,0.008) (-0.369,0.883)rectangle++(0.008,0.008) (-0.157,0.968)rectangle++(0.008,0.008) (-0.421,0.806)rectangle++(0.008,0.008) (-0.484,0.637)rectangle++(0.008,0.008)
  (-0.452,0.538)rectangle++(0.008,0.008) (-0.192,0.823)rectangle++(0.008,0.008) (-0.386,0.727)rectangle++(0.008,0.008) (-0.164,0.903)rectangle++(0.008,0.008) (-0.401,0.773)rectangle++(0.008,0.008)
  (0.092,0.913)rectangle++(0.008,0.008) (0.036,0.981)rectangle++(0.008,0.008) (0.366,0.852)rectangle++(0.008,0.008) (0.471,0.675)rectangle++(0.008,0.008) (0.195,0.881)rectangle++(0.008,0.008)
  (-0.229,0.892)rectangle++(0.008,0.008) (0.213,0.906)rectangle++(0.008,0.008) (-0.23,0.909)rectangle++(0.008,0.008) (-0.011,0.38)rectangle++(0.008,0.008) (-0.004,0.157)rectangle++(0.008,0.008)
  (0.053,0.49)rectangle++(0.008,0.008) (-0.154,0.662)rectangle++(0.008,0.008) (-0.31,0.667)rectangle++(0.008,0.008) (0.096,0.832)rectangle++(0.008,0.008) (-0.257,0.834)rectangle++(0.008,0.008)
  (-0.11,0.948)rectangle++(0.008,0.008) (-0.392,0.813)rectangle++(0.008,0.008) (-0.167,0.939)rectangle++(0.008,0.008) (0.258,0.905)rectangle++(0.008,0.008) (-0.209,0.925)rectangle++(0.008,0.008)
  (0.234,0.914)rectangle++(0.008,0.008) (0.003,0.382)rectangle++(0.008,0.008) (0.137,0.59)rectangle++(0.008,0.008) (-0.152,0.738)rectangle++(0.008,0.008) (-0.066,0.907)rectangle++(0.008,0.008)
  (-0.006,0.379)rectangle++(0.008,0.008) (0.132,0.592)rectangle++(0.008,0.008) (-0.155,0.737)rectangle++(0.008,0.008) (-0.337,0.701)rectangle++(0.008,0.008) (-0.144,0.892)rectangle++(0.008,0.008)
  (-0.388,0.776)rectangle++(0.008,0.008) (0.099,0.909)rectangle++(0.008,0.008) (0.039,0.98)rectangle++(0.008,0.008) (0.014,1.009)rectangle++(0.008,0.008) (0.366,0.873)rectangle++(0.008,0.008)
  (0.151,0.964)rectangle++(0.008,0.008) (0.061,1.003)rectangle++(0.008,0.008) (0.385,0.853)rectangle++(0.008,0.008) (-0.132,0.947)rectangle++(0.008,0.008) (-0.402,0.805)rectangle++(0.008,0.008)
  (-0.475,0.644)rectangle++(0.008,0.008) (0.012,0.88)rectangle++(0.008,0.008) (0.319,0.814)rectangle++(0.008,0.008) (0.436,0.675)rectangle++(0.008,0.008) (-0.045,0.883)rectangle++(0.008,0.008)
  (-0.34,0.807)rectangle++(0.008,0.008) (-0.447,0.667)rectangle++(0.008,0.008) (-0.017,0.278)rectangle++(0.008,0.008) (0.091,0.549)rectangle++(0.008,0.008) (-0.158,0.703)rectangle++(0.008,0.008)
  (-0.327,0.684)rectangle++(0.008,0.008) (0.094,0.846)rectangle++(0.008,0.008) (-0.263,0.839)rectangle++(0.008,0.008) (-0.423,0.709)rectangle++(0.008,0.008) (0.059,0.892)rectangle++(0.008,0.008)
  (0.022,0.972)rectangle++(0.008,0.008) (0.357,0.853)rectangle++(0.008,0.008) (0.147,0.956)rectangle++(0.008,0.008) (0.06,0.999)rectangle++(0.008,0.008) (-0.002,0.417)rectangle++(0.008,0.008)
  (0.147,0.608)rectangle++(0.008,0.008) (0.06,0.853)rectangle++(0.008,0.008) (-0.281,0.83)rectangle++(0.008,0.008) (0.167,0.896)rectangle++(0.008,0.008) (0.396,0.766)rectangle++(0.008,0.008)
  (0.164,0.92)rectangle++(0.008,0.008) (0.403,0.778)rectangle++(0.008,0.008) (0.167,0.924)rectangle++(0.008,0.008) (0.406,0.779)rectangle++(0.008,0.008) (0.008,0.325)rectangle++(0.008,0.008)
  (0.119,0.562)rectangle++(0.008,0.008) (-0.15,0.719)rectangle++(0.008,0.008) (-0.008,0.3)rectangle++(0.008,0.008) (-0.114,0.553)rectangle++(0.008,0.008) (-0.253,0.631)rectangle++(0.008,0.008)
  (-0.109,0.863)rectangle++(0.008,0.008) (-0.048,0.96)rectangle++(0.008,0.008) (0.32,0.872)rectangle++(0.008,0.008) (0.132,0.964)rectangle++(0.008,0.008) (-0.287,0.907)rectangle++(0.008,0.008)
  (0.192,0.934)rectangle++(0.008,0.008) (0.42,0.775)rectangle++(0.008,0.008) (0.174,0.923)rectangle++(0.008,0.008) (0.408,0.776)rectangle++(0.008,0.008) (-0.094,0.92)rectangle++(0.008,0.008)
  (0.285,0.87)rectangle++(0.008,0.008) (0.117,0.963)rectangle++(0.008,0.008) (0.047,1.002)rectangle++(0.008,0.008) (-0.34,0.894)rectangle++(0.008,0.008) (-0.145,0.973)rectangle++(0.008,0.008)
  (0.28,0.913)rectangle++(0.008,0.008) (0.115,0.981)rectangle++(0.008,0.008) (0.046,1.01)rectangle++(0.008,0.008) (0.017,1.022)rectangle++(0.008,0.008) (-0.361,0.892)rectangle++(0.008,0.008)
  (0.153,0.953)rectangle++(0.008,0.008) (-0.274,0.909)rectangle++(0.008,0.008) (-0.454,0.737)rectangle++(0.008,0.008) (-0.474,0.594)rectangle++(0.008,0.008) (-0.432,0.522)rectangle++(0.008,0.008)
  (-0.017,0.217)rectangle++(0.008,0.008) (0.069,0.521)rectangle++(0.008,0.008) (0.027,0.817)rectangle++(0.008,0.008) (0.303,0.78)rectangle++(0.008,0.008) (0.125,0.925)rectangle++(0.008,0.008)
  (0.387,0.795)rectangle++(0.008,0.008) (0.46,0.641)rectangle++(0.008,0.008) (0.438,0.545)rectangle++(0.008,0.008) (0.009,0.227)rectangle++(0.008,0.008) (0.084,0.517)rectangle++(0.008,0.008)
  (-0.15,0.686)rectangle++(0.008,0.008) (-0.065,0.886)rectangle++(0.008,0.008) (-0.35,0.801)rectangle++(0.008,0.008) (0.125,0.907)rectangle++(0.008,0.008) (0.381,0.786)rectangle++(0.008,0.008)
  (0.454,0.64)rectangle++(0.008,0.008) (0.435,0.547)rectangle++(0.008,0.008) (0.18,0.827)rectangle++(0.008,0.008) (0.073,0.945)rectangle++(0.008,0.008) (-0.308,0.877)rectangle++(0.008,0.008)
  (0.172,0.927)rectangle++(0.008,0.008) (0.409,0.779)rectangle++(0.008,0.008) (-0.095,0.921)rectangle++(0.008,0.008) (-0.376,0.806)rectangle++(0.008,0.008) (0.115,0.919)rectangle++(0.008,0.008)
  (-0.279,0.881)rectangle++(0.008,0.008) (0.186,0.919)rectangle++(0.008,0.008) (0.412,0.77)rectangle++(0.008,0.008) (0.008,0.321)rectangle++(0.008,0.008) (0.118,0.56)rectangle++(0.008,0.008)
  (0.047,0.833)rectangle++(0.008,0.008) (0.017,0.948)rectangle++(0.008,0.008) (0.346,0.843)rectangle++(0.008,0.008) (-0.147,0.928)rectangle++(0.008,0.008) (-0.402,0.791)rectangle++(0.008,0.008)
  (-0.47,0.637)rectangle++(0.008,0.008) (-0.2,0.865)rectangle++(0.008,0.008) (-0.086,0.961)rectangle++(0.008,0.008) (0.303,0.886)rectangle++(0.008,0.008) (-0.182,0.933)rectangle++(0.008,0.008)
  (0.249,0.908)rectangle++(0.008,0.008) (0.102,0.979)rectangle++(0.008,0.008) (0.041,1.009)rectangle++(0.008,0.008) (0.015,1.021)rectangle++(0.008,0.008) (-0.362,0.891)rectangle++(0.008,0.008)
  (-0.154,0.972)rectangle++(0.008,0.008) (0.276,0.916)rectangle++(0.008,0.008) (0.114,0.982)rectangle++(0.008,0.008) (-0.302,0.909)rectangle++(0.008,0.008) (-0.129,0.979)rectangle++(0.008,0.008)
  (-0.057,1.009)rectangle++(0.008,0.008) (0.334,0.898)rectangle++(0.008,0.008) (0.472,0.708)rectangle++(0.008,0.008) (0.468,0.571)rectangle++(0.008,0.008) (0.006,0.847)rectangle++(0.008,0.008)
  (0,0.954)rectangle++(0.008,0.008) (0.34,0.852)rectangle++(0.008,0.008) (0.459,0.684)rectangle++(0.008,0.008) (0.453,0.565)rectangle++(0.008,0.008) (0.188,0.835)rectangle++(0.008,0.008)
  (0.383,0.731)rectangle++(0.008,0.008) (0.159,0.905)rectangle++(0.008,0.008) (0.395,0.773)rectangle++(0.008,0.008) (0.164,0.922)rectangle++(0.008,0.008) (-0.258,0.899)rectangle++(0.008,0.008)
  (-0.111,0.975)rectangle++(0.008,0.008) (-0.403,0.825)rectangle++(0.008,0.008) (0.11,0.937)rectangle++(0.008,0.008) (0.384,0.806)rectangle++(0.008,0.008) (0.159,0.936)rectangle++(0.008,0.008)
  (0.406,0.787)rectangle++(0.008,0.008) (-0.099,0.924)rectangle++(0.008,0.008) (-0.379,0.806)rectangle++(0.008,0.008) (0.114,0.92)rectangle++(0.008,0.008) (-0.28,0.881)rectangle++(0.008,0.008)
  (-0.446,0.722)rectangle++(0.008,0.008) (-0.465,0.59)rectangle++(0.008,0.008) (-0.003,0.852)rectangle++(0.008,0.008) (0.302,0.807)rectangle++(0.008,0.008) (-0.154,0.896)rectangle++(0.008,0.008)
  (-0.394,0.774)rectangle++(0.008,0.008) (0.096,0.911)rectangle++(0.008,0.008) (0.368,0.798)rectangle++(0.008,0.008) (0.152,0.933)rectangle++(0.008,0.008) (-0.267,0.9)rectangle++(0.008,0.008)
  (-0.115,0.976)rectangle++(0.008,0.008) (0.295,0.903)rectangle++(0.008,0.008) (0.457,0.724)rectangle++(0.008,0.008) (0.19,0.902)rectangle++(0.008,0.008) (-0.239,0.899)rectangle++(0.008,0.008)
  (0.211,0.913)rectangle++(0.008,0.008) (0.086,0.981)rectangle++(0.008,0.008) (-0.314,0.898)rectangle++(0.008,0.008) (-0.468,0.718)rectangle++(0.008,0.008) (0.042,0.912)rectangle++(0.008,0.008)
  (-0.31,0.851)rectangle++(0.008,0.008) (0.161,0.916)rectangle++(0.008,0.008) (0.001,0.382)rectangle++(0.008,0.008) (-0.14,0.594)rectangle++(0.008,0.008) (0.148,0.738)rectangle++(0.008,0.008)
  (-0.2,0.809)rectangle++(0.008,0.008) (0.197,0.858)rectangle++(0.008,0.008) (0.396,0.738)rectangle++(0.008,0.008) (0.444,0.612)rectangle++(0.008,0.008) (0.009,0.255)rectangle++(0.008,0.008)
  (-0.09,0.539)rectangle++(0.008,0.008) (-0.237,0.633)rectangle++(0.008,0.008) (0.117,0.791)rectangle++(0.008,0.008) (-0.232,0.822)rectangle++(0.008,0.008) (-0.404,0.712)rectangle++(0.008,0.008)
  (-0.442,0.601)rectangle++(0.008,0.008) (-0.188,0.85)rectangle++(0.008,0.008) (-0.393,0.741)rectangle++(0.008,0.008) (-0.167,0.909)rectangle++(0.008,0.008) (-0.405,0.775)rectangle++(0.008,0.008)
  (0.091,0.915)rectangle++(0.008,0.008) (-0.289,0.87)rectangle++(0.008,0.008) (-0.446,0.714)rectangle++(0.008,0.008) (-0.462,0.586)rectangle++(0.008,0.008) (-0.003,0.849)rectangle++(0.008,0.008)
  (-0.308,0.806)rectangle++(0.008,0.008) (-0.432,0.678)rectangle++(0.008,0.008) (-0.443,0.575)rectangle++(0.008,0.008) (0.002,0.837)rectangle++(0.008,0.008) (-0.302,0.803)rectangle++(0.008,0.008)
  (-0.428,0.678)rectangle++(0.008,0.008) (-0.182,0.883)rectangle++(0.008,0.008) (0.231,0.885)rectangle++(0.008,0.008) (0.095,0.969)rectangle++(0.008,0.008) (0.389,0.826)rectangle++(0.008,0.008)
  (-0.121,0.935)rectangle++(0.008,0.008) (-0.053,0.991)rectangle++(0.008,0.008) (0.329,0.888)rectangle++(0.008,0.008) (0.006,0.371)rectangle++(0.008,0.008) (0.134,0.584)rectangle++(0.008,0.008)
  (0.269,0.635)rectangle++(0.008,0.008) (0.004,0.264)rectangle++(0.008,0.008) (0.096,0.536)rectangle++(0.008,0.008) (-0.151,0.698)rectangle++(0.008,0.008) (-0.322,0.684)rectangle++(0.008,0.008)
  (0.096,0.844)rectangle++(0.008,0.008) (0.345,0.768)rectangle++(0.008,0.008) (0.142,0.92)rectangle++(0.008,0.008) (0.058,0.984)rectangle++(0.008,0.008) (0.377,0.846)rectangle++(0.008,0.008)
  (-0.133,0.94)rectangle++(0.008,0.008) (0.274,0.894)rectangle++(0.008,0.008) (-0.197,0.926)rectangle++(0.008,0.008) (0.24,0.91)rectangle++(0.008,0.008) (0.434,0.747)rectangle++(0.008,0.008)
  (0.464,0.603)rectangle++(0.008,0.008) (0.427,0.526)rectangle++(0.008,0.008) (0.382,0.505)rectangle++(0.008,0.008) (0.158,0.81)rectangle++(0.008,0.008) (0.361,0.73)rectangle++(0.008,0.008)
  (-0.099,0.882)rectangle++(0.008,0.008) (0.269,0.855)rectangle++(0.008,0.008) (0.111,0.957)rectangle++(0.008,0.008) (-0.295,0.896)rectangle++(0.008,0.008) (-0.013,0.374)rectangle++(0.008,0.008)
  (-0.004,0.155)rectangle++(0.008,0.008) (-0.004,0.663)rectangle++(0.008,0.008) (-0.004,0.876)rectangle++(0.008,0.008) (0.31,0.818)rectangle++(0.008,0.008) (-0.154,0.904)rectangle++(0.008,0.008)
  (-0.397,0.777)rectangle++(0.008,0.008) (0.096,0.913)rectangle++(0.008,0.008) (0.369,0.8)rectangle++(0.008,0.008) (0.454,0.65)rectangle++(0.008,0.008) (0.188,0.871)rectangle++(0.008,0.008)
  (0.077,0.963)rectangle++(0.008,0.008) (-0.313,0.887)rectangle++(0.008,0.008) (-0.463,0.713)rectangle++(0.008,0.008) (-0.018,0.297)rectangle++(0.008,0.008) (-0.01,0.722)rectangle++(0.008,0.008)
  (-0.006,0.901)rectangle++(0.008,0.008) (-0.328,0.829)rectangle++(0.008,0.008) (-0.14,0.946)rectangle++(0.008,0.008) (-0.061,0.995)rectangle++(0.008,0.008) (-0.028,1.016)rectangle++(0.008,0.008)
  (0.349,0.891)rectangle++(0.008,0.008) (0.477,0.699)rectangle++(0.008,0.008) (0.467,0.566)rectangle++(0.008,0.008) (0.008,0.844)rectangle++(0.008,0.008) (-0.302,0.808)rectangle++(0.008,0.008)
  (-0.43,0.681)rectangle++(0.008,0.008) (0.046,0.881)rectangle++(0.008,0.008) (-0.297,0.838)rectangle++(0.008,0.008) (-0.127,0.95)rectangle++(0.008,0.008) (-0.401,0.808)rectangle++(0.008,0.008)
  (-0.475,0.645)rectangle++(0.008,0.008) (0.012,0.881)rectangle++(0.008,0.008) (0.003,0.968)rectangle++(0.008,0.008) (0.346,0.858)rectangle++(0.008,0.008) (-0.152,0.935)rectangle++(0.008,0.008)
  (-0.008,0.39)rectangle++(0.008,0.008) (-0.004,0.162)rectangle++(0.008,0.008) (-0.063,0.492)rectangle++(0.008,0.008) (-0.029,0.804)rectangle++(0.008,0.008) (0.273,0.794)rectangle++(0.008,0.008)
  (0.408,0.682)rectangle++(0.008,0.008) (0.008,0.284)rectangle++(0.008,0.008) (-0.101,0.552)rectangle++(0.008,0.008) (-0.247,0.635)rectangle++(0.008,0.008) (-0.106,0.865)rectangle++(0.008,0.008)
  (-0.361,0.777)rectangle++(0.008,0.008) (-0.446,0.645)rectangle++(0.008,0.008) (-0.19,0.869)rectangle++(0.008,0.008) (-0.01,0.363)rectangle++(0.008,0.008) (-0.004,0.15)rectangle++(0.008,0.008)
  (0.051,0.486)rectangle++(0.008,0.008) (0.019,0.802)rectangle++(0.008,0.008) (-0.281,0.793)rectangle++(0.008,0.008) (-0.415,0.681)rectangle++(0.008,0.008) (-0.437,0.582)rectangle++(0.008,0.008)
  (0.008,0.838)rectangle++(0.008,0.008) (0.302,0.797)rectangle++(0.008,0.008) (0.422,0.673)rectangle++(0.008,0.008) (0.009,0.28)rectangle++(0.008,0.008) (-0.1,0.55)rectangle++(0.008,0.008)
  (0.15,0.704)rectangle++(0.008,0.008) (-0.186,0.794)rectangle++(0.008,0.008) (0.198,0.846)rectangle++(0.008,0.008) (0.081,0.953)rectangle++(0.008,0.008) (0.376,0.823)rectangle++(0.008,0.008)
  (0.156,0.943)rectangle++(0.008,0.008) (0.001,0.394)rectangle++(0.008,0.008) (-0.144,0.6)rectangle++(0.008,0.008) (-0.008,0.25)rectangle++(0.008,0.008) (-0.006,0.702)rectangle++(0.008,0.008)
  (-0.257,0.738)rectangle++(0.008,0.008) (-0.385,0.665)rectangle++(0.008,0.008) (-0.417,0.586)rectangle++(0.008,0.008) (0.018,0.833)rectangle++(0.008,0.008) (-0.003,0.348)rectangle++(0.008,0.008)
  (0.122,0.576)rectangle++(0.008,0.008) (0,0.24)rectangle++(0.008,0.008) (-0.002,0.698)rectangle++(0.008,0.008) (0.248,0.736)rectangle++(0.008,0.008) (0.375,0.665)rectangle++(0.008,0.008)
  (0.007,0.277)rectangle++(0.008,0.008) (-0.004,0.114)rectangle++(0.008,0.008) (-0.046,0.47)rectangle++(0.008,0.008) (0.146,0.648)rectangle++(0.008,0.008) (0.297,0.66)rectangle++(0.008,0.008)
  (0.371,0.612)rectangle++(0.008,0.008) (-0.053,0.831)rectangle++(0.008,0.008) (-0.024,0.947)rectangle++(0.008,0.008) (0.326,0.858)rectangle++(0.008,0.008) (0.135,0.958)rectangle++(0.008,0.008)
  (0.054,1)rectangle++(0.008,0.008) (0.02,1.018)rectangle++(0.008,0.008) (-0.358,0.892)rectangle++(0.008,0.008) (0.154,0.952)rectangle++(0.008,0.008) (0.062,0.997)rectangle++(0.008,0.008)
  (0.384,0.85)rectangle++(0.008,0.008) (0.478,0.668)rectangle++(0.008,0.008) (0.456,0.551)rectangle++(0.008,0.008) (0.01,0.229)rectangle++(0.008,0.008) (0.086,0.518)rectangle++(0.008,0.008)
  (0.034,0.815)rectangle++(0.008,0.008) (0.012,0.94)rectangle++(0.008,0.008) (-0.334,0.853)rectangle++(0.008,0.008) (-0.143,0.956)rectangle++(0.008,0.008) (-0.41,0.805)rectangle++(0.008,0.008)
  (-0.175,0.936)rectangle++(0.008,0.008) (-0.076,0.991)rectangle++(0.008,0.008) (-0.392,0.845)rectangle++(0.008,0.008) (0.122,0.943)rectangle++(0.008,0.008) (0.392,0.804)rectangle++(0.008,0.008)
  (0.162,0.935)rectangle++(0.008,0.008) (0.001,0.39)rectangle++(0.008,0.008) (-0.143,0.598)rectangle++(0.008,0.008) (-0.062,0.849)rectangle++(0.008,0.008) (-0.028,0.954)rectangle++(0.008,0.008)
  (-0.014,0.998)rectangle++(0.008,0.008) (0.349,0.878)rectangle++(0.008,0.008) (0.144,0.966)rectangle++(0.008,0.008) (0.41,0.807)rectangle++(0.008,0.008) (-0.104,0.934)rectangle++(0.008,0.008)
  (0.285,0.881)rectangle++(0.008,0.008) (0.118,0.968)rectangle++(0.008,0.008) (0.047,1.004)rectangle++(0.008,0.008) (0.379,0.858)rectangle++(0.008,0.008) (0.157,0.958)rectangle++(0.008,0.008)
  (0.064,1)rectangle++(0.008,0.008) (-0.002,0.418)rectangle++(0.008,0.008) (-0.154,0.609)rectangle++(0.008,0.008) (-0.067,0.854)rectangle++(0.008,0.008) (-0.03,0.956)rectangle++(0.008,0.008)
  (-0.359,0.845)rectangle++(0.008,0.008) (-0.015,0.353)rectangle++(0.008,0.008) (-0.008,0.746)rectangle++(0.008,0.008) (0.262,0.76)rectangle++(0.008,0.008) (0.108,0.917)rectangle++(0.008,0.008)
  (0.043,0.983)rectangle++(0.008,0.008) (-0.335,0.884)rectangle++(0.008,0.008) (0.162,0.94)rectangle++(0.008,0.008) (0.409,0.788)rectangle++(0.008,0.008) (0.008,0.329)rectangle++(0.008,0.008)
  (0.001,0.736)rectangle++(0.008,0.008) (-0.266,0.756)rectangle++(0.008,0.008) (-0.395,0.67)rectangle++(0.008,0.008) (-0.423,0.584)rectangle++(0.008,0.008) (-0.18,0.843)rectangle++(0.008,0.008)
  (-0.078,0.952)rectangle++(0.008,0.008) (-0.379,0.826)rectangle++(0.008,0.008) (-0.162,0.945)rectangle++(0.008,0.008) (0.263,0.906)rectangle++(0.008,0.008) (0.443,0.737)rectangle++(0.008,0.008)
  (0.009,0.307)rectangle++(0.008,0.008) (-0.109,0.563)rectangle++(0.008,0.008) (-0.048,0.834)rectangle++(0.008,0.008) (-0.023,0.948)rectangle++(0.008,0.008) (0.327,0.858)rectangle++(0.008,0.008)
  (0.455,0.691)rectangle++(0.008,0.008) (-0.042,0.898)rectangle++(0.008,0.008) (-0.02,0.975)rectangle++(0.008,0.008) (-0.011,1.007)rectangle++(0.008,0.008) (-0.007,1.021)rectangle++(0.008,0.008)
  (0.361,0.885)rectangle++(0.008,0.008) (-0.155,0.953)rectangle++(0.008,0.008) (0.269,0.907)rectangle++(0.008,0.008) (0.111,0.979)rectangle++(0.008,0.008) (0.399,0.824)rectangle++(0.008,0.008)
  (0.476,0.651)rectangle++(0.008,0.008) (0.198,0.871)rectangle++(0.008,0.008) (0.401,0.744)rectangle++(0.008,0.008) (-0.086,0.902)rectangle++(0.008,0.008) (-0.006,0.377)rectangle++(0.008,0.008)
  (0.131,0.591)rectangle++(0.008,0.008) (0.053,0.846)rectangle++(0.008,0.008) (-0.282,0.825)rectangle++(0.008,0.008) (0.165,0.894)rectangle++(0.008,0.008) (-0.247,0.887)rectangle++(0.008,0.008)
  (0.203,0.91)rectangle++(0.008,0.008) (-0.236,0.908)rectangle++(0.008,0.008) (-0.435,0.75)rectangle++(0.008,0.008) (0.068,0.915)rectangle++(0.008,0.008) (-0.299,0.862)rectangle++(0.008,0.008)
  (-0.128,0.96)rectangle++(0.008,0.008) (-0.056,1.001)rectangle++(0.008,0.008) (-0.387,0.857)rectangle++(0.008,0.008) (-0.486,0.673)rectangle++(0.008,0.008) (0.017,0.897)rectangle++(0.008,0.008)
  (-0.316,0.836)rectangle++(0.008,0.008) (-0.135,0.949)rectangle++(0.008,0.008) (-0.008,0.396)rectangle++(0.008,0.008) (-0.006,0.764)rectangle++(0.008,0.008) (0.27,0.768)rectangle++(0.008,0.008)
  (0.111,0.92)rectangle++(0.008,0.008) (-0.001,0.384)rectangle++(0.008,0.008) (-0.004,0.159)rectangle++(0.008,0.008) (-0.004,0.664)rectangle++(0.008,0.008) (-0.004,0.877)rectangle++(0.008,0.008)
  (-0.318,0.819)rectangle++(0.008,0.008) (0.146,0.904)rectangle++(0.008,0.008) (0.059,0.977)rectangle++(0.008,0.008) (0.375,0.842)rectangle++(0.008,0.008) (0.471,0.667)rectangle++(0.008,0.008)
  (0.01,0.278)rectangle++(0.008,0.008) (-0.098,0.55)rectangle++(0.008,0.008) (-0.044,0.829)rectangle++(0.008,0.008) (0.275,0.811)rectangle++(0.008,0.008) (0.004,0.338)rectangle++(0.008,0.008)
  (-0.122,0.575)rectangle++(0.008,0.008) (-0.008,0.239)rectangle++(0.008,0.008) (0.081,0.528)rectangle++(0.008,0.008) (-0.001,0.22)rectangle++(0.008,0.008) (-0.083,0.519)rectangle++(0.008,0.008)
  (-0.227,0.627)rectangle++(0.008,0.008) (-0.098,0.861)rectangle++(0.008,0.008) (-0.356,0.778)rectangle++(0.008,0.008) (-0.444,0.648)rectangle++(0.008,0.008) (0.028,0.871)rectangle++(0.008,0.008)
  (0.009,0.964)rectangle++(0.008,0.008) (0.002,1.002)rectangle++(0.008,0.008) (0.358,0.874)rectangle++(0.008,0.008) (0.007,0.365)rectangle++(0.008,0.008) (0.001,0.751)rectangle++(0.008,0.008)
  (-0.004,0.313)rectangle++(0.008,0.008) (0.109,0.561)rectangle++(0.008,0.008) (-0.154,0.715)rectangle++(0.008,0.008) (-0.067,0.898)rectangle++(0.008,0.008) (-0.355,0.806)rectangle++(0.008,0.008)
  (0.125,0.911)rectangle++(0.008,0.008) (0,0.38)rectangle++(0.008,0.008) (-0.14,0.593)rectangle++(0.008,0.008) (-0.061,0.847)rectangle++(0.008,0.008) (-0.006,0.353)rectangle++(0.008,0.008)
  (-0.004,0.146)rectangle++(0.008,0.008) (0.05,0.485)rectangle++(0.008,0.008) (0.195,0.62)rectangle++(0.008,0.008) (-0.136,0.772)rectangle++(0.008,0.008) (-0.342,0.724)rectangle++(0.008,0.008)
  (0.102,0.869)rectangle++(0.008,0.008) (-0.268,0.853)rectangle++(0.008,0.008) (-0.43,0.713)rectangle++(0.008,0.008) (-0.455,0.592)rectangle++(0.008,0.008) (-0.193,0.846)rectangle++(0.008,0.008)
  (-0.084,0.953)rectangle++(0.008,0.008) (-0.037,0.998)rectangle++(0.008,0.008) (0.339,0.886)rectangle++(0.008,0.008) (-0.165,0.945)rectangle++(0.008,0.008) (0.261,0.907)rectangle++(0.008,0.008)
  (-0.208,0.927)rectangle++(0.008,0.008) (-0.09,0.987)rectangle++(0.008,0.008) (-0.397,0.838)rectangle++(0.008,0.008) (-0.169,0.95)rectangle++(0.008,0.008) (0.261,0.911)rectangle++(0.008,0.008)
  (-0.21,0.929)rectangle++(0.008,0.008) (-0.09,0.988)rectangle++(0.008,0.008) (0.311,0.9)rectangle++(0.008,0.008) (-0.183,0.942)rectangle++(0.008,0.008) (-0.423,0.784)rectangle++(0.008,0.008)
  (0.086,0.926)rectangle++(0.008,0.008) (0.369,0.809)rectangle++(0.008,0.008) (0.153,0.937)rectangle++(0.008,0.008) (0.062,0.991)rectangle++(0.008,0.008) (-0.329,0.894)rectangle++(0.008,0.008)
  (-0.474,0.71)rectangle++(0.008,0.008) (0.037,0.91)rectangle++(0.008,0.008) (0.341,0.819)rectangle++(0.008,0.008) (-0.14,0.915)rectangle++(0.008,0.008) (-0.061,0.982)rectangle++(0.008,0.008)
  (-0.028,1.01)rectangle++(0.008,0.008) (-0.005,0.422)rectangle++(0.008,0.008) (-0.156,0.61)rectangle++(0.008,0.008) (-0.009,0.254)rectangle++(0.008,0.008) (-0.098,0.532)rectangle++(0.008,0.008)
  (-0.007,0.221)rectangle++(0.008,0.008) (-0.005,0.691)rectangle++(0.008,0.008) (0.243,0.734)rectangle++(0.008,0.008) (-0.154,0.842)rectangle++(0.008,0.008) (-0.375,0.749)rectangle++(0.008,0.008)
  (0.095,0.892)rectangle++(0.008,0.008) (-0.279,0.861)rectangle++(0.008,0.008) (-0.012,0.359)rectangle++(0.008,0.008) (-0.138,0.579)rectangle++(0.008,0.008) (0.143,0.73)rectangle++(0.008,0.008)
  (-0.199,0.804)rectangle++(0.008,0.008) (-0.382,0.716)rectangle++(0.008,0.008) (0.08,0.88)rectangle++(0.008,0.008) (0.031,0.967)rectangle++(0.008,0.008) (-0.335,0.873)rectangle++(0.008,0.008)
  (-0.468,0.699)rectangle++(0.008,0.008) (0.035,0.903)rectangle++(0.008,0.008) (0.338,0.817)rectangle++(0.008,0.008) (0.445,0.669)rectangle++(0.008,0.008) (-0.039,0.884)rectangle++(0.008,0.008)
  (-0.019,0.969)rectangle++(0.008,0.008) (-0.004,0.405)rectangle++(0.008,0.008) (0.142,0.603)rectangle++(0.008,0.008) (-0.154,0.745)rectangle++(0.008,0.008) (-0.34,0.705)rectangle++(0.008,0.008)
  (-0.411,0.62)rectangle++(0.008,0.008) (-0.175,0.858)rectangle++(0.008,0.008) (-0.39,0.749)rectangle++(0.008,0.008) (-0.166,0.912)rectangle++(0.008,0.008) (-0.072,0.981)rectangle++(0.008,0.008)
  (-0.387,0.842)rectangle++(0.008,0.008) (-0.165,0.951)rectangle++(0.008,0.008) (-0.419,0.795)rectangle++(0.008,0.008) (0.092,0.929)rectangle++(0.008,0.008) (0.036,0.988)rectangle++(0.008,0.008)
  (0.369,0.855)rectangle++(0.008,0.008) (0.473,0.676)rectangle++(0.008,0.008) (0.457,0.556)rectangle++(0.008,0.008) (0.01,0.231)rectangle++(0.008,0.008) (0.002,0.695)rectangle++(0.008,0.008)
  (-0.251,0.737)rectangle++(0.008,0.008) (0.148,0.843)rectangle++(0.008,0.008) (0.368,0.749)rectangle++(0.008,0.008) (0.435,0.627)rectangle++(0.008,0.008) (0.422,0.548)rectangle++(0.008,0.008)
  (0.175,0.828)rectangle++(0.008,0.008) (-0.219,0.86)rectangle++(0.008,0.008) (-0.01,0.359)rectangle++(0.008,0.008) (0.123,0.584)rectangle++(0.008,0.008) (-0.156,0.73)rectangle++(0.008,0.008)
  (-0.068,0.904)rectangle++(0.008,0.008) (0.291,0.854)rectangle++(0.008,0.008) (-0.175,0.914)rectangle++(0.008,0.008) (0.245,0.897)rectangle++(0.008,0.008) (-0.212,0.917)rectangle++(0.008,0.008)
  (-0.091,0.983)rectangle++(0.008,0.008) (-0.396,0.836)rectangle++(0.008,0.008) (0.117,0.94)rectangle++(0.008,0.008) (-0.286,0.89)rectangle++(0.008,0.008) (-0.122,0.972)rectangle++(0.008,0.008)
  (-0.407,0.82)rectangle++(0.008,0.008) (-0.482,0.649)rectangle++(0.008,0.008) (-0.018,0.27)rectangle++(0.008,0.008) (-0.01,0.711)rectangle++(0.008,0.008) (-0.007,0.896)rectangle++(0.008,0.008)
  (-0.327,0.827)rectangle++(0.008,0.008) (-0.448,0.68)rectangle++(0.008,0.008) (-0.451,0.57)rectangle++(0.008,0.008) (-0.003,0.838)rectangle++(0.008,0.008) (-0.004,0.35)rectangle++(0.008,0.008)
  (-0.13,0.578)rectangle++(0.008,0.008) (-0.008,0.24)rectangle++(0.008,0.008) (0.082,0.529)rectangle++(0.008,0.008) (-0.155,0.69)rectangle++(0.008,0.008) (-0.321,0.679)rectangle++(0.008,0.008)
  (-0.137,0.883)rectangle++(0.008,0.008) (0.252,0.869)rectangle++(0.008,0.008) (-0.199,0.906)rectangle++(0.008,0.008) (0.232,0.902)rectangle++(0.008,0.008) (0.095,0.976)rectangle++(0.008,0.008)
  (0.038,1.008)rectangle++(0.008,0.008) (0.013,1.021)rectangle++(0.008,0.008) (-0.362,0.891)rectangle++(0.008,0.008) (0.152,0.953)rectangle++(0.008,0.008) (-0.274,0.909)rectangle++(0.008,0.008)
  (-0.118,0.979)rectangle++(0.008,0.008) (-0.407,0.825)rectangle++(0.008,0.008) (-0.173,0.944)rectangle++(0.008,0.008) (-0.075,0.994)rectangle++(0.008,0.008) (-0.393,0.847)rectangle++(0.008,0.008)
  (0.122,0.944)rectangle++(0.008,0.008) (0.392,0.804)rectangle++(0.008,0.008) (0.466,0.644)rectangle++(0.008,0.008) (0.442,0.544)rectangle++(0.008,0.008) (0.004,0.826)rectangle++(0.008,0.008)
  (-0.001,0.945)rectangle++(0.008,0.008) (0.336,0.848)rectangle++(0.008,0.008) (0.139,0.954)rectangle++(0.008,0.008) (0.056,0.998)rectangle++(0.008,0.008) (-0.334,0.896)rectangle++(0.008,0.008)
  (-0.476,0.709)rectangle++(0.008,0.008) (-0.474,0.573)rectangle++(0.008,0.008) (-0.425,0.512)rectangle++(0.008,0.008) (-0.381,0.501)rectangle++(0.008,0.008) (0.004,0.781)rectangle++(0.008,0.008)
  (-0.281,0.778)rectangle++(0.008,0.008) (0.149,0.872)rectangle++(0.008,0.008) (0.379,0.762)rectangle++(0.008,0.008) (0.157,0.918)rectangle++(0.008,0.008) (0.399,0.78)rectangle++(0.008,0.008)
  (0.165,0.925)rectangle++(0.008,0.008) (-0.259,0.901)rectangle++(0.008,0.008) (-0.111,0.976)rectangle++(0.008,0.008) (-0.049,1.008)rectangle++(0.008,0.008) (-0.023,1.021)rectangle++(0.008,0.008)
  (0.353,0.891)rectangle++(0.008,0.008) (0.146,0.972)rectangle++(0.008,0.008) (0.059,1.006)rectangle++(0.008,0.008) (-0.336,0.9)rectangle++(0.008,0.008) (0.167,0.948)rectangle++(0.008,0.008)
  (0.414,0.79)rectangle++(0.008,0.008) (-0.096,0.928)rectangle++(0.008,0.008) (-0.379,0.809)rectangle++(0.008,0.008) (0.115,0.922)rectangle++(0.008,0.008) (0.381,0.797)rectangle++(0.008,0.008)
  (0.458,0.644)rectangle++(0.008,0.008) (0.19,0.868)rectangle++(0.008,0.008) (0.396,0.745)rectangle++(0.008,0.008) (0.164,0.911)rectangle++(0.008,0.008) (0.399,0.774)rectangle++(0.008,0.008)
  (0.165,0.923)rectangle++(0.008,0.008) (-0.257,0.9)rectangle++(0.008,0.008) (-0.443,0.739)rectangle++(0.008,0.008) (-0.188,0.908)rectangle++(0.008,0.008) (0.237,0.899)rectangle++(0.008,0.008)
  (0.003,0.375)rectangle++(0.008,0.008) (-0.001,0.755)rectangle++(0.008,0.008) (0.268,0.762)rectangle++(0.008,0.008) (0.394,0.669)rectangle++(0.008,0.008) (-0.062,0.866)rectangle++(0.008,0.008)
  (-0.341,0.793)rectangle++(0.008,0.008) (-0.146,0.931)rectangle++(0.008,0.008) (-0.008,0.389)rectangle++(0.008,0.008) (-0.006,0.761)rectangle++(0.008,0.008) (-0.004,0.317)rectangle++(0.008,0.008)
  (-0.119,0.563)rectangle++(0.008,0.008) (-0.007,0.234)rectangle++(0.008,0.008) (0.079,0.526)rectangle++(0.008,0.008) (0.223,0.628)rectangle++(0.008,0.008) (0.326,0.624)rectangle++(0.008,0.008)
  (0.371,0.585)rectangle++(0.008,0.008) (0.007,0.243)rectangle++(0.008,0.008) (0.001,0.7)rectangle++(0.008,0.008) (-0.253,0.739)rectangle++(0.008,0.008) (0.148,0.845)rectangle++(0.008,0.008)
  (-0.238,0.858)rectangle++(0.008,0.008) (0.197,0.893)rectangle++(0.008,0.008) (0.08,0.973)rectangle++(0.008,0.008) (0.383,0.832)rectangle++(0.008,0.008) (-0.126,0.937)rectangle++(0.008,0.008)
  (-0.395,0.802)rectangle++(0.008,0.008) (-0.471,0.645)rectangle++(0.008,0.008) (0.014,0.879)rectangle++(0.008,0.008) (0.004,0.967)rectangle++(0.008,0.008) (-0.001,1.004)rectangle++(0.008,0.008)
  (0.357,0.875)rectangle++(0.008,0.008) (0.148,0.965)rectangle++(0.008,0.008) (0.06,1.003)rectangle++(0.008,0.008) (-0.334,0.899)rectangle++(0.008,0.008) (-0.478,0.711)rectangle++(0.008,0.008)
  (-0.203,0.896)rectangle++(0.008,0.008) (-0.088,0.974)rectangle++(0.008,0.008) (-0.039,1.007)rectangle++(0.008,0.008) (-0.381,0.865)rectangle++(0.008,0.008) (-0.487,0.679)rectangle++(0.008,0.008)
  (-0.207,0.883)rectangle++(0.008,0.008) (-0.413,0.749)rectangle++(0.008,0.008) (0.078,0.906)rectangle++(0.008,0.008) (-0.292,0.861)rectangle++(0.008,0.008) (-0.013,0.359)rectangle++(0.008,0.008)
  (-0.008,0.749)rectangle++(0.008,0.008) (0.263,0.761)rectangle++(0.008,0.008) (0.391,0.67)rectangle++(0.008,0.008) (0.418,0.583)rectangle++(0.008,0.008) (-0.021,0.835)rectangle++(0.008,0.008)
  (-0.311,0.793)rectangle++(0.008,0.008) (-0.013,0.331)rectangle++(0.008,0.008) (-0.128,0.566)rectangle++(0.008,0.008) (0.143,0.721)rectangle++(0.008,0.008) (0.322,0.695)rectangle++(0.008,0.008)
  (-0.105,0.852)rectangle++(0.008,0.008) (0.256,0.843)rectangle++(0.008,0.008) (0.004,0.352)rectangle++(0.008,0.008) (0.127,0.576)rectangle++(0.008,0.008) (0.051,0.84)rectangle++(0.008,0.008)
  (-0.28,0.821)rectangle++(0.008,0.008) (-0.12,0.943)rectangle++(0.008,0.008) (-0.395,0.807)rectangle++(0.008,0.008) (0.107,0.926)rectangle++(0.008,0.008) (-0.286,0.881)rectangle++(0.008,0.008)
  (0.183,0.921)rectangle++(0.008,0.008) (0.412,0.772)rectangle++(0.008,0.008) (-0.091,0.919)rectangle++(0.008,0.008) (0.286,0.869)rectangle++(0.008,0.008) (-0.183,0.918)rectangle++(0.008,0.008)
  (-0.415,0.774)rectangle++(0.008,0.008) (-0.177,0.923)rectangle++(0.008,0.008) (-0.077,0.985)rectangle++(0.008,0.008) (-0.034,1.011)rectangle++(0.008,0.008) (-0.017,1.022)rectangle++(0.008,0.008)
  (-0.004,0.427)rectangle++(0.008,0.008) (-0.158,0.613)rectangle++(0.008,0.008) (-0.295,0.643)rectangle++(0.008,0.008) (-0.368,0.608)rectangle++(0.008,0.008) (0.048,0.826)rectangle++(0.008,0.008)
  (0.018,0.944)rectangle++(0.008,0.008) (0.345,0.842)rectangle++(0.008,0.008) (0.142,0.951)rectangle++(0.008,0.008) (0.404,0.8)rectangle++(0.008,0.008) (0.008,0.334)rectangle++(0.008,0.008)
  (0.122,0.566)rectangle++(0.008,0.008) (0.049,0.835)rectangle++(0.008,0.008) (-0.28,0.819)rectangle++(0.008,0.008) (0.164,0.89)rectangle++(0.008,0.008) (0.066,0.972)rectangle++(0.008,0.008)
  (0.026,1.006)rectangle++(0.008,0.008) (0.37,0.867)rectangle++(0.008,0.008) (-0.144,0.948)rectangle++(0.008,0.008) (0.272,0.901)rectangle++(0.008,0.008) (-0.201,0.928)rectangle++(0.008,0.008)
  (-0.427,0.772)rectangle++(0.008,0.008) (-0.182,0.922)rectangle++(0.008,0.008) (-0.416,0.776)rectangle++(0.008,0.008) (-0.471,0.625)rectangle++(0.008,0.008) (-0.018,0.26)rectangle++(0.008,0.008)
  (0.084,0.542)rectangle++(0.008,0.008) (-0.159,0.697)rectangle++(0.008,0.008) (0.176,0.792)rectangle++(0.008,0.008) (0.071,0.93)rectangle++(0.008,0.008) (0.028,0.988)rectangle++(0.008,0.008)
  (0.009,1.013)rectangle++(0.008,0.008) (0.002,1.023)rectangle++(0.008,0.008) (-0.368,0.887)rectangle++(0.008,0.008) (-0.157,0.97)rectangle++(0.008,0.008) (-0.422,0.807)rectangle++(0.008,0.008)
  (-0.484,0.637)rectangle++(0.008,0.008) (-0.206,0.865)rectangle++(0.008,0.008) (-0.407,0.741)rectangle++(0.008,0.008) (-0.016,0.309)rectangle++(0.008,0.008) (0.102,0.563)rectangle++(0.008,0.008)
  (-0.001,0.234)rectangle++(0.008,0.008) (-0.088,0.526)rectangle++(0.008,0.008) (-0.231,0.628)rectangle++(0.008,0.008) (-0.011,0.262)rectangle++(0.008,0.008) (0.088,0.54)rectangle++(0.008,0.008)
  (0.035,0.824)rectangle++(0.008,0.008) (0.309,0.781)rectangle++(0.008,0.008) (0.42,0.663)rectangle++(0.008,0.008) (-0.049,0.872)rectangle++(0.008,0.008) (0.288,0.832)rectangle++(0.008,0.008)
  (-0.169,0.903)rectangle++(0.008,0.008) (0.244,0.889)rectangle++(0.008,0.008) (-0.209,0.913)rectangle++(0.008,0.008) (-0.09,0.981)rectangle++(0.008,0.008) (0.308,0.897)rectangle++(0.008,0.008)
  (0.461,0.716)rectangle++(0.008,0.008) (0.191,0.899)rectangle++(0.008,0.008) (-0.237,0.898)rectangle++(0.008,0.008) (-0.433,0.745)rectangle++(0.008,0.008) (-0.184,0.911)rectangle++(0.008,0.008)
  (-0.009,0.38)rectangle++(0.008,0.008) (-0.144,0.59)rectangle++(0.008,0.008) (-0.28,0.637)rectangle++(0.008,0.008) (0.099,0.808)rectangle++(0.008,0.008) (0.333,0.75)rectangle++(0.008,0.008)
  (-0.119,0.881)rectangle++(0.008,0.008) (-0.007,0.368)rectangle++(0.008,0.008) (-0.138,0.585)rectangle++(0.008,0.008) (0.145,0.733)rectangle++(0.008,0.008) (0,0.306)rectangle++(0.008,0.008)
  (-0.112,0.559)rectangle++(0.008,0.008) (0.147,0.712)rectangle++(0.008,0.008) (-0.19,0.797)rectangle++(0.008,0.008) (-0.01,0.333)rectangle++(0.008,0.008) (0.114,0.572)rectangle++(0.008,0.008);
\fill[nuage]
  (0.255,0.637)rectangle++(0.008,0.008) (0.105,0.865)rectangle++(0.008,0.008) (0.356,0.774)rectangle++(0.008,0.008) (0.147,0.923)rectangle++(0.008,0.008) (-0.266,0.894)rectangle++(0.008,0.008)
  (-0.444,0.733)rectangle++(0.008,0.008) (0.058,0.91)rectangle++(0.008,0.008) (0.351,0.812)rectangle++(0.008,0.008) (0.449,0.662)rectangle++(0.008,0.008) (0.186,0.876)rectangle++(0.008,0.008)
  (0.397,0.75)rectangle++(0.008,0.008) (0.449,0.617)rectangle++(0.008,0.008) (0.01,0.257)rectangle++(0.008,0.008) (-0.091,0.54)rectangle++(0.008,0.008) (-0.007,0.225)rectangle++(0.008,0.008)
  (-0.087,0.52)rectangle++(0.008,0.008) (-0.229,0.626)rectangle++(0.008,0.008) (-0.332,0.624)rectangle++(0.008,0.008) (0.07,0.82)rectangle++(0.008,0.008) (0.027,0.942)rectangle++(0.008,0.008)
  (0.348,0.837)rectangle++(0.008,0.008) (-0.143,0.926)rectangle++(0.008,0.008) (-0.063,0.987)rectangle++(0.008,0.008) (-0.029,1.012)rectangle++(0.008,0.008) (-0.014,1.023)rectangle++(0.008,0.008)
  (-0.008,1.027)rectangle++(0.008,0.008) (-0.374,0.886)rectangle++(0.008,0.008) (0.145,0.955)rectangle++(0.008,0.008) (0.406,0.801)rectangle++(0.008,0.008) (0.008,0.334)rectangle++(0.008,0.008)
  (-0.119,0.575)rectangle++(0.008,0.008) (0.15,0.722)rectangle++(0.008,0.008) (0.061,0.901)rectangle++(0.008,0.008) (-0.298,0.853)rectangle++(0.008,0.008) (-0.444,0.703)rectangle++(0.008,0.008)
  (-0.189,0.893)rectangle++(0.008,0.008) (0.232,0.892)rectangle++(0.008,0.008) (0.003,0.372)rectangle++(0.008,0.008) (-0.135,0.59)rectangle++(0.008,0.008) (0.148,0.735)rectangle++(0.008,0.008)
  (0.06,0.906)rectangle++(0.008,0.008) (0.35,0.809)rectangle++(0.008,0.008) (0.448,0.661)rectangle++(0.008,0.008) (-0.035,0.882)rectangle++(0.008,0.008) (-0.017,0.968)rectangle++(0.008,0.008)
  (-0.009,1.004)rectangle++(0.008,0.008) (-0.366,0.875)rectangle++(0.008,0.008) (-0.156,0.965)rectangle++(0.008,0.008) (-0.42,0.805)rectangle++(0.008,0.008) (-0.483,0.637)rectangle++(0.008,0.008)
  (-0.205,0.865)rectangle++(0.008,0.008) (-0.01,0.361)rectangle++(0.008,0.008) (0.124,0.585)rectangle++(0.008,0.008) (0.265,0.64)rectangle++(0.008,0.008) (0.004,0.266)rectangle++(0.008,0.008)
  (-0.097,0.542)rectangle++(0.008,0.008) (-0.242,0.633)rectangle++(0.008,0.008) (-0.34,0.622)rectangle++(0.008,0.008) (0.066,0.822)rectangle++(0.008,0.008) (0.025,0.943)rectangle++(0.008,0.008)
  (0.348,0.838)rectangle++(0.008,0.008) (0.457,0.675)rectangle++(0.008,0.008) (0.01,0.281)rectangle++(0.008,0.008) (-0.1,0.551)rectangle++(0.008,0.008) (0.151,0.704)rectangle++(0.008,0.008)
  (0.061,0.893)rectangle++(0.008,0.008) (-0.295,0.849)rectangle++(0.008,0.008) (-0.013,0.354)rectangle++(0.008,0.008) (0.12,0.583)rectangle++(0.008,0.008) (0,0.243)rectangle++(0.008,0.008)
  (-0.002,0.7)rectangle++(0.008,0.008) (0.248,0.737)rectangle++(0.008,0.008) (-0.153,0.845)rectangle++(0.008,0.008) (-0.067,0.952)rectangle++(0.008,0.008) (-0.374,0.831)rectangle++(0.008,0.008)
  (-0.471,0.665)rectangle++(0.008,0.008) (-0.2,0.877)rectangle++(0.008,0.008) (0.221,0.889)rectangle++(0.008,0.008) (0.09,0.971)rectangle++(0.008,0.008) (0.036,1.005)rectangle++(0.008,0.008)
  (-0.346,0.892)rectangle++(0.008,0.008) (0.159,0.948)rectangle++(0.008,0.008) (-0.269,0.909)rectangle++(0.008,0.008) (0.201,0.928)rectangle++(0.008,0.008) (-0.243,0.915)rectangle++(0.008,0.008)
  (-0.442,0.751)rectangle++(0.008,0.008) (0.066,0.917)rectangle++(0.008,0.008) (-0.301,0.862)rectangle++(0.008,0.008) (0.169,0.918)rectangle++(0.008,0.008) (0.404,0.775)rectangle++(0.008,0.008)
  (0.168,0.923)rectangle++(0.008,0.008) (-0.257,0.901)rectangle++(0.008,0.008) (-0.443,0.739)rectangle++(0.008,0.008) (-0.017,0.308)rectangle++(0.008,0.008) (-0.122,0.554)rectangle++(0.008,0.008)
  (-0.257,0.629)rectangle++(0.008,0.008) (-0.012,0.262)rectangle++(0.008,0.008) (-0.007,0.708)rectangle++(0.008,0.008) (0.249,0.742)rectangle++(0.008,0.008) (0.102,0.909)rectangle++(0.008,0.008)
  (-0.282,0.871)rectangle++(0.008,0.008) (-0.444,0.717)rectangle++(0.008,0.008) (-0.189,0.899)rectangle++(0.008,0.008) (-0.411,0.763)rectangle++(0.008,0.008) (-0.464,0.621)rectangle++(0.008,0.008)
  (0.009,0.866)rectangle++(0.008,0.008) (0.313,0.809)rectangle++(0.008,0.008) (-0.15,0.901)rectangle++(0.008,0.008) (0.252,0.881)rectangle++(0.008,0.008) (0.429,0.729)rectangle++(0.008,0.008)
  (-0.068,0.906)rectangle++(0.008,0.008) (-0.006,0.378)rectangle++(0.008,0.008) (-0.005,0.756)rectangle++(0.008,0.008) (-0.276,0.763)rectangle++(0.008,0.008) (0.146,0.864)rectangle++(0.008,0.008)
  (0.374,0.759)rectangle++(0.008,0.008) (0.155,0.917)rectangle++(0.008,0.008) (0.397,0.78)rectangle++(0.008,0.008) (0.459,0.631)rectangle++(0.008,0.008) (-0.019,0.872)rectangle++(0.008,0.008)
  (0.302,0.822)rectangle++(0.008,0.008) (0.124,0.943)rectangle++(0.008,0.008) (-0.283,0.895)rectangle++(0.008,0.008) (0.189,0.926)rectangle++(0.008,0.008) (0.417,0.772)rectangle++(0.008,0.008)
  (0.465,0.621)rectangle++(0.008,0.008) (0.434,0.534)rectangle++(0.008,0.008) (0.388,0.506)rectangle++(0.008,0.008) (0.357,0.509)rectangle++(0.008,0.008) (0.007,0.211)rectangle++(0.008,0.008)
  (-0.004,0.086)rectangle++(0.008,0.008) (-0.036,0.457)rectangle++(0.008,0.008) (-0.005,0.19)rectangle++(0.008,0.008) (0.065,0.505)rectangle++(0.008,0.008) (-0.154,0.673)rectangle++(0.008,0.008)
  (0.169,0.779)rectangle++(0.008,0.008) (0.355,0.712)rectangle++(0.008,0.008) (0.416,0.615)rectangle++(0.008,0.008) (0.409,0.549)rectangle++(0.008,0.008) (0.382,0.521)rectangle++(0.008,0.008)
  (0.36,0.518)rectangle++(0.008,0.008) (-0.024,0.785)rectangle++(0.008,0.008) (0.268,0.784)rectangle++(0.008,0.008) (0.402,0.679)rectangle++(0.008,0.008) (0.425,0.583)rectangle++(0.008,0.008)
  (-0.017,0.838)rectangle++(0.008,0.008) (-0.01,0.949)rectangle++(0.008,0.008) (-0.347,0.85)rectangle++(0.008,0.008) (-0.148,0.955)rectangle++(0.008,0.008) (-0.412,0.803)rectangle++(0.008,0.008)
  (-0.016,0.335)rectangle++(0.008,0.008) (-0.131,0.566)rectangle++(0.008,0.008) (-0.266,0.632)rectangle++(0.008,0.008) (0.103,0.8)rectangle++(0.008,0.008) (0.332,0.745)rectangle++(0.008,0.008)
  (-0.118,0.878)rectangle++(0.008,0.008) (-0.371,0.779)rectangle++(0.008,0.008) (-0.451,0.643)rectangle++(0.008,0.008) (0.022,0.871)rectangle++(0.008,0.008) (0.321,0.807)rectangle++(0.008,0.008)
  (-0.145,0.902)rectangle++(0.008,0.008) (-0.063,0.977)rectangle++(0.008,0.008) (-0.006,0.408)rectangle++(0.008,0.008) (0.142,0.605)rectangle++(0.008,0.008) (0.057,0.852)rectangle++(0.008,0.008)
  (0.022,0.955)rectangle++(0.008,0.008) (0.35,0.845)rectangle++(0.008,0.008) (-0.145,0.931)rectangle++(0.008,0.008) (0.265,0.894)rectangle++(0.008,0.008) (-0.202,0.922)rectangle++(0.008,0.008)
  (-0.01,0.385)rectangle++(0.008,0.008) (-0.146,0.592)rectangle++(0.008,0.008) (-0.063,0.846)rectangle++(0.008,0.008) (-0.335,0.783)rectangle++(0.008,0.008) (0.126,0.894)rectangle++(0.008,0.008)
  (-0.265,0.873)rectangle++(0.008,0.008) (0.19,0.91)rectangle++(0.008,0.008) (0.411,0.765)rectangle++(0.008,0.008) (0.17,0.919)rectangle++(0.008,0.008) (0.405,0.776)rectangle++(0.008,0.008)
  (0.461,0.626)rectangle++(0.008,0.008) (0.01,0.261)rectangle++(0.008,0.008) (-0.092,0.542)rectangle++(0.008,0.008) (-0.239,0.634)rectangle++(0.008,0.008) (-0.103,0.864)rectangle++(0.008,0.008)
  (-0.045,0.961)rectangle++(0.008,0.008) (0.321,0.872)rectangle++(0.008,0.008) (0.457,0.7)rectangle++(0.008,0.008) (0.19,0.892)rectangle++(0.008,0.008) (0.404,0.756)rectangle++(0.008,0.008)
  (0.454,0.618)rectangle++(0.008,0.008) (0.188,0.857)rectangle++(0.008,0.008) (0.391,0.741)rectangle++(0.008,0.008) (0.443,0.615)rectangle++(0.008,0.008) (0.009,0.256)rectangle++(0.008,0.008)
  (0.095,0.53)rectangle++(0.008,0.008) (0.232,0.625)rectangle++(0.008,0.008) (0.095,0.86)rectangle++(0.008,0.008) (0.038,0.959)rectangle++(0.008,0.008) (-0.329,0.871)rectangle++(0.008,0.008)
  (-0.465,0.7)rectangle++(0.008,0.008) (0.037,0.902)rectangle++(0.008,0.008) (-0.309,0.845)rectangle++(0.008,0.008) (-0.446,0.695)rectangle++(0.008,0.008) (0.043,0.893)rectangle++(0.008,0.008)
  (-0.303,0.843)rectangle++(0.008,0.008) (-0.443,0.696)rectangle++(0.008,0.008) (0.045,0.893)rectangle++(0.008,0.008) (-0.302,0.844)rectangle++(0.008,0.008) (0.163,0.91)rectangle++(0.008,0.008)
  (0.066,0.98)rectangle++(0.008,0.008) (-0.323,0.891)rectangle++(0.008,0.008) (0.17,0.939)rectangle++(0.008,0.008) (-0.261,0.909)rectangle++(0.008,0.008) (-0.112,0.979)rectangle++(0.008,0.008)
  (0.298,0.904)rectangle++(0.008,0.008) (0.458,0.723)rectangle++(0.008,0.008) (0.467,0.583)rectangle++(0.008,0.008) (0.194,0.843)rectangle++(0.008,0.008) (0.389,0.732)rectangle++(0.008,0.008)
  (-0.087,0.893)rectangle++(0.008,0.008) (0.278,0.856)rectangle++(0.008,0.008) (0.114,0.957)rectangle++(0.008,0.008) (0.393,0.813)rectangle++(0.008,0.008) (0.469,0.648)rectangle++(0.008,0.008)
  (-0.02,0.883)rectangle++(0.008,0.008) (-0.004,0.368)rectangle++(0.008,0.008) (-0.137,0.586)rectangle++(0.008,0.008) (-0.276,0.638)rectangle++(0.008,0.008) (-0.118,0.866)rectangle++(0.008,0.008)
  (-0.052,0.961)rectangle++(0.008,0.008) (-0.024,1.001)rectangle++(0.008,0.008) (-0.005,0.418)rectangle++(0.008,0.008) (-0.004,0.773)rectangle++(0.008,0.008) (0.273,0.771)rectangle++(0.008,0.008)
  (-0.154,0.869)rectangle++(0.008,0.008) (0.239,0.869)rectangle++(0.008,0.008) (0.003,0.363)rectangle++(0.008,0.008) (-0.001,0.75)rectangle++(0.008,0.008) (-0.003,0.913)rectangle++(0.008,0.008)
  (-0.004,0.381)rectangle++(0.008,0.008) (0.133,0.592)rectangle++(0.008,0.008) (0.272,0.639)rectangle++(0.008,0.008) (0.352,0.612)rectangle++(0.008,0.008) (-0.061,0.824)rectangle++(0.008,0.008)
  (-0.028,0.944)rectangle++(0.008,0.008) (-0.014,0.994)rectangle++(0.008,0.008) (-0.365,0.869)rectangle++(0.008,0.008) (0.143,0.944)rectangle++(0.008,0.008) (0.401,0.797)rectangle++(0.008,0.008)
  (-0.105,0.927)rectangle++(0.008,0.008) (-0.046,0.987)rectangle++(0.008,0.008) (0.331,0.884)rectangle++(0.008,0.008) (-0.168,0.941)rectangle++(0.008,0.008) (-0.417,0.789)rectangle++(0.008,0.008)
  (-0.476,0.631)rectangle++(0.008,0.008) (-0.447,0.538)rectangle++(0.008,0.008) (-0.4,0.506)rectangle++(0.008,0.008) (-0.17,0.81)rectangle++(0.008,0.008) (0.211,0.847)rectangle++(0.008,0.008)
  (0.002,0.354)rectangle++(0.008,0.008) (-0.129,0.582)rectangle++(0.008,0.008) (-0.056,0.842)rectangle++(0.008,0.008) (-0.006,0.351)rectangle++(0.008,0.008) (-0.005,0.745)rectangle++(0.008,0.008)
  (-0.272,0.758)rectangle++(0.008,0.008) (0.146,0.86)rectangle++(0.008,0.008) (0,0.359)rectangle++(0.008,0.008) (-0.002,0.748)rectangle++(0.008,0.008) (0.266,0.759)rectangle++(0.008,0.008)
  (0.109,0.917)rectangle++(0.008,0.008) (0.376,0.796)rectangle++(0.008,0.008) (0.156,0.932)rectangle++(0.008,0.008) (0.063,0.989)rectangle++(0.008,0.008) (0.024,1.013)rectangle++(0.008,0.008)
  (0.008,1.023)rectangle++(0.008,0.008) (0.368,0.881)rectangle++(0.008,0.008) (0.152,0.968)rectangle++(0.008,0.008) (0.062,1.004)rectangle++(0.008,0.008) (-0.334,0.9)rectangle++(0.008,0.008)
  (0.168,0.947)rectangle++(0.008,0.008) (0.001,0.395)rectangle++(0.008,0.008) (0.141,0.597)rectangle++(0.008,0.008) (-0.152,0.742)rectangle++(0.008,0.008) (0.195,0.81)rectangle++(0.008,0.008)
  (0.079,0.938)rectangle++(0.008,0.008) (-0.302,0.876)rectangle++(0.008,0.008) (0.174,0.925)rectangle++(0.008,0.008) (0.409,0.777)rectangle++(0.008,0.008) (0.169,0.924)rectangle++(0.008,0.008)
  (0.069,0.986)rectangle++(0.008,0.008) (-0.324,0.894)rectangle++(0.008,0.008) (0.17,0.941)rectangle++(0.008,0.008) (0.069,0.993)rectangle++(0.008,0.008) (0.385,0.845)rectangle++(0.008,0.008)
  (-0.129,0.943)rectangle++(0.008,0.008) (0.277,0.894)rectangle++(0.008,0.008) (0.445,0.726)rectangle++(0.008,0.008) (0.009,0.303)rectangle++(0.008,0.008) (0.112,0.551)rectangle++(0.008,0.008)
  (-0.149,0.711)rectangle++(0.008,0.008) (0.185,0.795)rectangle++(0.008,0.008) (0.002,0.331)rectangle++(0.008,0.008) (0.118,0.567)rectangle++(0.008,0.008) (0.047,0.836)rectangle++(0.008,0.008)
  (0.018,0.949)rectangle++(0.008,0.008) (-0.334,0.859)rectangle++(0.008,0.008) (-0.143,0.959)rectangle++(0.008,0.008) (-0.008,0.4)rectangle++(0.008,0.008) (-0.15,0.599)rectangle++(0.008,0.008)
  (-0.065,0.849)rectangle++(0.008,0.008) (-0.337,0.784)rectangle++(0.008,0.008) (-0.144,0.927)rectangle++(0.008,0.008) (0.265,0.891)rectangle++(0.008,0.008) (-0.201,0.921)rectangle++(0.008,0.008)
  (-0.424,0.768)rectangle++(0.008,0.008) (-0.017,0.32)rectangle++(0.008,0.008) (-0.009,0.732)rectangle++(0.008,0.008) (-0.269,0.751)rectangle++(0.008,0.008) (-0.395,0.666)rectangle++(0.008,0.008)
  (-0.422,0.583)rectangle++(0.008,0.008) (-0.405,0.535)rectangle++(0.008,0.008) (0.005,0.805)rectangle++(0.008,0.008) (0.289,0.783)rectangle++(0.008,0.008) (0.411,0.671)rectangle++(0.008,0.008)
  (-0.055,0.873)rectangle++(0.008,0.008) (-0.34,0.798)rectangle++(0.008,0.008) (-0.014,0.333)rectangle++(0.008,0.008) (0.112,0.574)rectangle++(0.008,0.008) (-0.157,0.721)rectangle++(0.008,0.008)
  (0.185,0.802)rectangle++(0.008,0.008) (-0.206,0.852)rectangle++(0.008,0.008) (-0.01,0.356)rectangle++(0.008,0.008) (-0.135,0.578)rectangle++(0.008,0.008) (0.144,0.729)rectangle++(0.008,0.008)
  (-0.198,0.804)rectangle++(0.008,0.008) (-0.086,0.935)rectangle++(0.008,0.008) (0.294,0.874)rectangle++(0.008,0.008) (-0.181,0.924)rectangle++(0.008,0.008) (0.246,0.903)rectangle++(0.008,0.008)
  (0.101,0.977)rectangle++(0.008,0.008) (0.394,0.827)rectangle++(0.008,0.008) (0.008,0.345)rectangle++(0.008,0.008) (0.126,0.571)rectangle++(0.008,0.008) (0.051,0.838)rectangle++(0.008,0.008)
  (0.019,0.949)rectangle++(0.008,0.008) (-0.334,0.86)rectangle++(0.008,0.008) (0.154,0.929)rectangle++(0.008,0.008) (0.401,0.786)rectangle++(0.008,0.008) (0.166,0.928)rectangle++(0.008,0.008)
  (0.406,0.781)rectangle++(0.008,0.008) (0.464,0.628)rectangle++(0.008,0.008) (0.193,0.862)rectangle++(0.008,0.008) (0.079,0.96)rectangle++(0.008,0.008) (0.031,1.001)rectangle++(0.008,0.008)
  (-0.003,0.418)rectangle++(0.008,0.008) (-0.154,0.609)rectangle++(0.008,0.008) (-0.009,0.254)rectangle++(0.008,0.008) (-0.006,0.704)rectangle++(0.008,0.008) (0.248,0.74)rectangle++(0.008,0.008)
  (0.004,0.309)rectangle++(0.008,0.008) (-0.112,0.562)rectangle++(0.008,0.008) (-0.255,0.636)rectangle++(0.008,0.008) (-0.347,0.619)rectangle++(0.008,0.008) (-0.383,0.578)rectangle++(0.008,0.008)
  (-0.163,0.84)rectangle++(0.008,0.008) (-0.071,0.951)rectangle++(0.008,0.008) (-0.006,0.397)rectangle++(0.008,0.008) (-0.148,0.599)rectangle++(0.008,0.008) (-0.285,0.64)rectangle++(0.008,0.008)
  (0.098,0.811)rectangle++(0.008,0.008) (0.039,0.938)rectangle++(0.008,0.008) (-0.321,0.862)rectangle++(0.008,0.008) (0.16,0.925)rectangle++(0.008,0.008) (0.403,0.782)rectangle++(0.008,0.008)
  (0.463,0.63)rectangle++(0.008,0.008) (0.436,0.539)rectangle++(0.008,0.008) (0.181,0.824)rectangle++(0.008,0.008) (0.002,0.344)rectangle++(0.008,0.008) (-0.126,0.577)rectangle++(0.008,0.008)
  (-0.008,0.24)rectangle++(0.008,0.008) (0.081,0.529)rectangle++(0.008,0.008) (0.225,0.629)rectangle++(0.008,0.008) (-0.125,0.787)rectangle++(0.008,0.008) (-0.342,0.734)rectangle++(0.008,0.008)
  (-0.146,0.906)rectangle++(0.008,0.008) (0.256,0.883)rectangle++(0.008,0.008) (0.431,0.728)rectangle++(0.008,0.008) (-0.066,0.906)rectangle++(0.008,0.008) (-0.03,0.978)rectangle++(0.008,0.008)
  (-0.015,1.009)rectangle++(0.008,0.008) (-0.371,0.875)rectangle++(0.008,0.008) (-0.158,0.965)rectangle++(0.008,0.008) (-0.009,0.403)rectangle++(0.008,0.008) (-0.151,0.6)rectangle++(0.008,0.008)
  (0.144,0.745)rectangle++(0.008,0.008) (0.058,0.911)rectangle++(0.008,0.008) (-0.302,0.856)rectangle++(0.008,0.008) (-0.013,0.357)rectangle++(0.008,0.008) (-0.137,0.578)rectangle++(0.008,0.008)
  (-0.008,0.24)rectangle++(0.008,0.008) (-0.006,0.699)rectangle++(0.008,0.008) (0.246,0.738)rectangle++(0.008,0.008) (-0.155,0.844)rectangle++(0.008,0.008) (0.23,0.857)rectangle++(0.008,0.008)
  (0.094,0.958)rectangle++(0.008,0.008) (-0.303,0.891)rectangle++(0.008,0.008) (0.179,0.932)rectangle++(0.008,0.008) (0.001,0.389)rectangle++(0.008,0.008) (-0.002,0.761)rectangle++(0.008,0.008)
  (-0.004,0.317)rectangle++(0.008,0.008) (0.111,0.563)rectangle++(0.008,0.008) (0.044,0.834)rectangle++(0.008,0.008) (-0.281,0.816)rectangle++(0.008,0.008) (-0.12,0.941)rectangle++(0.008,0.008)
  (-0.394,0.806)rectangle++(0.008,0.008) (-0.168,0.936)rectangle++(0.008,0.008) (-0.073,0.991)rectangle++(0.008,0.008) (0.32,0.895)rectangle++(0.008,0.008) (0.132,0.974)rectangle++(0.008,0.008)
  (0.407,0.814)rectangle++(0.008,0.008) (0.169,0.94)rectangle++(0.008,0.008) (0.069,0.992)rectangle++(0.008,0.008) (0.385,0.845)rectangle++(0.008,0.008) (0.477,0.665)rectangle++(0.008,0.008)
  (-0.023,0.894)rectangle++(0.008,0.008) (0.308,0.833)rectangle++(0.008,0.008) (0.437,0.687)rectangle++(0.008,0.008) (0.009,0.286)rectangle++(0.008,0.008) (-0.102,0.553)rectangle++(0.008,0.008)
  (-0.248,0.636)rectangle++(0.008,0.008) (-0.106,0.865)rectangle++(0.008,0.008) (0.259,0.85)rectangle++(0.008,0.008) (-0.188,0.9)rectangle++(0.008,0.008) (-0.411,0.763)rectangle++(0.008,0.008)
  (-0.016,0.318)rectangle++(0.008,0.008) (0.106,0.568)rectangle++(0.008,0.008) (0.042,0.836)rectangle++(0.008,0.008) (0.015,0.949)rectangle++(0.008,0.008) (0.004,0.996)rectangle++(0.008,0.008)
  (-0.001,1.016)rectangle++(0.008,0.008) (-0.004,0.424)rectangle++(0.008,0.008) (-0.157,0.612)rectangle++(0.008,0.008) (-0.294,0.643)rectangle++(0.008,0.008) (0.095,0.815)rectangle++(0.008,0.008)
  (0.037,0.94)rectangle++(0.008,0.008) (0.352,0.833)rectangle++(0.008,0.008) (0.457,0.671)rectangle++(0.008,0.008) (0.19,0.88)rectangle++(0.008,0.008) (0.4,0.751)rectangle++(0.008,0.008)
  (0.45,0.617)rectangle++(0.008,0.008) (0.425,0.537)rectangle++(0.008,0.008) (0.009,0.223)rectangle++(0.008,0.008) (-0.079,0.525)rectangle++(0.008,0.008) (-0.006,0.218)rectangle++(0.008,0.008)
  (-0.004,0.089)rectangle++(0.008,0.008) (0.029,0.459)rectangle++(0.008,0.008) (0.176,0.616)rectangle++(0.008,0.008) (0.072,0.856)rectangle++(0.008,0.008) (0.338,0.782)rectangle++(0.008,0.008)
  (0.433,0.653)rectangle++(0.008,0.008) (-0.039,0.872)rectangle++(0.008,0.008) (0.293,0.829)rectangle++(0.008,0.008) (0.121,0.946)rectangle++(0.008,0.008) (0.392,0.806)rectangle++(0.008,0.008)
  (0.162,0.936)rectangle++(0.008,0.008) (0.408,0.786)rectangle++(0.008,0.008) (0.466,0.63)rectangle++(0.008,0.008) (0.437,0.538)rectangle++(0.008,0.008) (0.004,0.821)rectangle++(0.008,0.008)
  (0.294,0.79)rectangle++(0.008,0.008) (0.416,0.672)rectangle++(0.008,0.008) (0.172,0.88)rectangle++(0.008,0.008) (-0.239,0.883)rectangle++(0.008,0.008) (-0.103,0.969)rectangle++(0.008,0.008)
  (-0.045,1.004)rectangle++(0.008,0.008) (-0.383,0.862)rectangle++(0.008,0.008) (-0.015,0.36)rectangle++(0.008,0.008) (-0.004,0.149)rectangle++(0.008,0.008) (-0.004,0.66)rectangle++(0.008,0.008)
  (-0.241,0.72)rectangle++(0.008,0.008) (-0.104,0.9)rectangle++(0.008,0.008) (0.273,0.865)rectangle++(0.008,0.008) (0.433,0.714)rectangle++(0.008,0.008) (0.452,0.588)rectangle++(0.008,0.008)
  (0.416,0.524)rectangle++(0.008,0.008) (0.009,0.218)rectangle++(0.008,0.008) (-0.077,0.522)rectangle++(0.008,0.008) (0.15,0.683)rectangle++(0.008,0.008) (-0.179,0.785)rectangle++(0.008,0.008)
  (0.198,0.839)rectangle++(0.008,0.008) (0.389,0.729)rectangle++(0.008,0.008) (0.008,0.304)rectangle++(0.008,0.008) (0.111,0.553)rectangle++(0.008,0.008) (-0.15,0.712)rectangle++(0.008,0.008)
  (0.185,0.795)rectangle++(0.008,0.008) (0.075,0.932)rectangle++(0.008,0.008) (0.029,0.989)rectangle++(0.008,0.008) (-0.343,0.882)rectangle++(0.008,0.008) (-0.147,0.968)rectangle++(0.008,0.008)
  (0.278,0.911)rectangle++(0.008,0.008) (-0.202,0.935)rectangle++(0.008,0.008) (0.241,0.916)rectangle++(0.008,0.008) (-0.221,0.924)rectangle++(0.008,0.008) (-0.434,0.763)rectangle++(0.008,0.008)
  (0.073,0.92)rectangle++(0.008,0.008) (0.028,0.984)rectangle++(0.008,0.008) (0.364,0.856)rectangle++(0.008,0.008) (-0.143,0.94)rectangle++(0.008,0.008) (-0.405,0.798)rectangle++(0.008,0.008)
  (0.099,0.926)rectangle++(0.008,0.008) (0.039,0.986)rectangle++(0.008,0.008) (0.369,0.853)rectangle++(0.008,0.008) (0.473,0.674)rectangle++(0.008,0.008) (0.196,0.881)rectangle++(0.008,0.008)
  (0.08,0.968)rectangle++(0.008,0.008) (-0.313,0.89)rectangle++(0.008,0.008) (-0.134,0.972)rectangle++(0.008,0.008) (-0.412,0.816)rectangle++(0.008,0.008) (-0.175,0.94)rectangle++(0.008,0.008)
  (-0.419,0.786)rectangle++(0.008,0.008) (-0.476,0.629)rectangle++(0.008,0.008) (-0.202,0.862)rectangle++(0.008,0.008) (-0.087,0.96)rectangle++(0.008,0.008) (0.302,0.886)rectangle++(0.008,0.008)
  (0.125,0.97)rectangle++(0.008,0.008) (0.403,0.815)rectangle++(0.008,0.008) (-0.111,0.936)rectangle++(0.008,0.008) (-0.388,0.807)rectangle++(0.008,0.008) (-0.165,0.937)rectangle++(0.008,0.008)
  (-0.414,0.788)rectangle++(0.008,0.008) (-0.176,0.929)rectangle++(0.008,0.008) (0.25,0.904)rectangle++(0.008,0.008) (-0.212,0.922)rectangle++(0.008,0.008) (-0.43,0.765)rectangle++(0.008,0.008)
  (-0.017,0.319)rectangle++(0.008,0.008) (-0.004,0.132)rectangle++(0.008,0.008) (-0.053,0.478)rectangle++(0.008,0.008) (-0.024,0.798)rectangle++(0.008,0.008) (-0.3,0.775)rectangle++(0.008,0.008)
  (0.139,0.878)rectangle++(0.008,0.008) (0.376,0.768)rectangle++(0.008,0.008) (0.007,0.32)rectangle++(0.008,0.008) (-0.115,0.568)rectangle++(0.008,0.008) (-0.05,0.836)rectangle++(0.008,0.008)
  (-0.024,0.949)rectangle++(0.008,0.008) (0.327,0.859)rectangle++(0.008,0.008) (0.006,0.358)rectangle++(0.008,0.008) (-0.129,0.585)rectangle++(0.008,0.008) (0.149,0.73)rectangle++(0.008,0.008)
  (-0.196,0.806)rectangle++(0.008,0.008) (-0.01,0.336)rectangle++(0.008,0.008) (0.115,0.574)rectangle++(0.008,0.008) (0.046,0.839)rectangle++(0.008,0.008) (0.32,0.783)rectangle++(0.008,0.008)
  (0.132,0.927)rectangle++(0.008,0.008) (0.39,0.793)rectangle++(0.008,0.008) (-0.108,0.921)rectangle++(0.008,0.008) (0.279,0.876)rectangle++(0.008,0.008) (-0.189,0.919)rectangle++(0.008,0.008)
  (-0.082,0.984)rectangle++(0.008,0.008) (-0.037,1.011)rectangle++(0.008,0.008) (-0.018,1.022)rectangle++(0.008,0.008) (0.356,0.89)rectangle++(0.008,0.008) (-0.159,0.953)rectangle++(0.008,0.008)
  (-0.416,0.798)rectangle++(0.008,0.008) (0.094,0.93)rectangle++(0.008,0.008) (-0.293,0.878)rectangle++(0.008,0.008) (0.179,0.922)rectangle++(0.008,0.008) (0.073,0.985)rectangle++(0.008,0.008)
  (0.028,1.011)rectangle++(0.008,0.008) (-0.352,0.892)rectangle++(0.008,0.008) (0.157,0.95)rectangle++(0.008,0.008) (0.41,0.794)rectangle++(0.008,0.008) (-0.1,0.929)rectangle++(0.008,0.008)
  (-0.007,0.388)rectangle++(0.008,0.008) (0.135,0.596)rectangle++(0.008,0.008) (0.054,0.848)rectangle++(0.008,0.008) (-0.282,0.826)rectangle++(0.008,0.008) (0.166,0.895)rectangle++(0.008,0.008)
  (-0.247,0.887)rectangle++(0.008,0.008) (-0.433,0.736)rectangle++(0.008,0.008) (0.064,0.908)rectangle++(0.008,0.008) (-0.298,0.857)rectangle++(0.008,0.008) (0.169,0.915)rectangle++(0.008,0.008)
  (-0.253,0.898)rectangle++(0.008,0.008) (-0.011,0.375)rectangle++(0.008,0.008) (0.128,0.592)rectangle++(0.008,0.008) (0.051,0.846)rectangle++(0.008,0.008) (-0.282,0.824)rectangle++(0.008,0.008)
  (0.165,0.894)rectangle++(0.008,0.008) (0.001,0.373)rectangle++(0.008,0.008) (-0.136,0.59)rectangle++(0.008,0.008) (0.148,0.735)rectangle++(0.008,0.008) (0.329,0.7)rectangle++(0.008,0.008)
  (0.399,0.619)rectangle++(0.008,0.008) (0.403,0.556)rectangle++(0.008,0.008) (0.167,0.831)rectangle++(0.008,0.008) (-0.224,0.859)rectangle++(0.008,0.008) (-0.096,0.958)rectangle++(0.008,0.008)
  (-0.043,1)rectangle++(0.008,0.008) (-0.02,1.018)rectangle++(0.008,0.008) (-0.011,1.025)rectangle++(0.008,0.008) (-0.004,0.428)rectangle++(0.008,0.008) (-0.004,0.778)rectangle++(0.008,0.008)
  (-0.283,0.773)rectangle++(0.008,0.008) (0.146,0.871)rectangle++(0.008,0.008) (-0.248,0.869)rectangle++(0.008,0.008) (-0.106,0.963)rectangle++(0.008,0.008) (0.294,0.894)rectangle++(0.008,0.008)
  (0.453,0.72)rectangle++(0.008,0.008) (0.01,0.3)rectangle++(0.008,0.008) (0.002,0.724)rectangle++(0.008,0.008) (-0.261,0.751)rectangle++(0.008,0.008) (-0.391,0.669)rectangle++(0.008,0.008)
  (-0.016,0.279)rectangle++(0.008,0.008) (-0.009,0.715)rectangle++(0.008,0.008) (-0.006,0.898)rectangle++(0.008,0.008) (-0.327,0.827)rectangle++(0.008,0.008) (0.145,0.911)rectangle++(0.008,0.008)
  (0,0.38)rectangle++(0.008,0.008) (-0.002,0.757)rectangle++(0.008,0.008) (0.269,0.763)rectangle++(0.008,0.008) (0.111,0.918)rectangle++(0.008,0.008) (-0.281,0.878)rectangle++(0.008,0.008)
  (-0.012,0.367)rectangle++(0.008,0.008) (0.125,0.588)rectangle++(0.008,0.008) (0.05,0.845)rectangle++(0.008,0.008) (0.019,0.952)rectangle++(0.008,0.008) (-0.003,0.398)rectangle++(0.008,0.008)
  (-0.147,0.6)rectangle++(0.008,0.008) (0.146,0.743)rectangle++(0.008,0.008) (0.331,0.704)rectangle++(0.008,0.008) (0.402,0.62)rectangle++(0.008,0.008) (0.167,0.858)rectangle++(0.008,0.008)
  (-0.234,0.871)rectangle++(0.008,0.008) (-0.101,0.963)rectangle++(0.008,0.008) (0.297,0.893)rectangle++(0.008,0.008) (0.123,0.973)rectangle++(0.008,0.008) (0.403,0.817)rectangle++(0.008,0.008)
  (0.475,0.646)rectangle++(0.008,0.008) (0.447,0.542)rectangle++(0.008,0.008) (0.397,0.505)rectangle++(0.008,0.008) (0.164,0.81)rectangle++(0.008,0.008) (-0.218,0.848)rectangle++(0.008,0.008)
  (-0.406,0.729)rectangle++(0.008,0.008) (-0.449,0.608)rectangle++(0.008,0.008) (-0.426,0.536)rectangle++(0.008,0.008) (-0.004,0.814)rectangle++(0.008,0.008) (0.288,0.79)rectangle++(0.008,0.008)
  (0.413,0.674)rectangle++(0.008,0.008) (0.429,0.577)rectangle++(0.008,0.008) (-0.014,0.836)rectangle++(0.008,0.008) (0.292,0.803)rectangle++(0.008,0.008) (0.12,0.935)rectangle++(0.008,0.008)
  (-0.283,0.89)rectangle++(0.008,0.008) (-0.45,0.725)rectangle++(0.008,0.008) (-0.468,0.59)rectangle++(0.008,0.008) (-0.428,0.522)rectangle++(0.008,0.008) (-0.182,0.817)rectangle++(0.008,0.008)
  (-0.009,0.341)rectangle++(0.008,0.008) (-0.13,0.572)rectangle++(0.008,0.008) (-0.057,0.838)rectangle++(0.008,0.008) (0.272,0.82)rectangle++(0.008,0.008) (0.004,0.342)rectangle++(0.008,0.008)
  (0.123,0.571)rectangle++(0.008,0.008) (0.26,0.633)rectangle++(0.008,0.008) (0.344,0.613)rectangle++(0.008,0.008) (0.142,0.855)rectangle++(0.008,0.008) (0.057,0.957)rectangle++(0.008,0.008)
  (-0.002,0.4)rectangle++(0.008,0.008) (-0.147,0.601)rectangle++(0.008,0.008) (-0.064,0.85)rectangle++(0.008,0.008) (-0.029,0.955)rectangle++(0.008,0.008) (0.327,0.863)rectangle++(0.008,0.008)
  (-0.162,0.93)rectangle++(0.008,0.008) (0.257,0.9)rectangle++(0.008,0.008) (0.004,0.376)rectangle++(0.008,0.008) (-0.136,0.592)rectangle++(0.008,0.008) (-0.277,0.641)rectangle++(0.008,0.008)
  (-0.012,0.267)rectangle++(0.008,0.008) (-0.004,0.11)rectangle++(0.008,0.008) (-0.045,0.468)rectangle++(0.008,0.008) (-0.191,0.617)rectangle++(0.008,0.008) (-0.311,0.633)rectangle++(0.008,0.008)
  (-0.133,0.864)rectangle++(0.008,0.008) (-0.373,0.766)rectangle++(0.008,0.008) (0.103,0.9)rectangle++(0.008,0.008) (0.041,0.976)rectangle++(0.008,0.008) (-0.333,0.88)rectangle++(0.008,0.008)
  (-0.47,0.702)rectangle++(0.008,0.008) (-0.469,0.572)rectangle++(0.008,0.008) (-0.018,0.238)rectangle++(0.008,0.008) (-0.097,0.522)rectangle++(0.008,0.008) (-0.234,0.623)rectangle++(0.008,0.008)
  (0.115,0.785)rectangle++(0.008,0.008) (0,0.327)rectangle++(0.008,0.008) (-0.003,0.735)rectangle++(0.008,0.008) (0.261,0.753)rectangle++(0.008,0.008) (0.387,0.668)rectangle++(0.008,0.008)
  (-0.065,0.863)rectangle++(0.008,0.008) (0.278,0.834)rectangle++(0.008,0.008) (0.424,0.698)rectangle++(0.008,0.008) (-0.059,0.89)rectangle++(0.008,0.008) (-0.348,0.805)rectangle++(0.008,0.008)
  (-0.45,0.663)rectangle++(0.008,0.008) (0.03,0.88)rectangle++(0.008,0.008) (0.327,0.808)rectangle++(0.008,0.008) (-0.142,0.905)rectangle++(0.008,0.008) (-0.062,0.978)rectangle++(0.008,0.008)
  (0.32,0.886)rectangle++(0.008,0.008) (0.132,0.97)rectangle++(0.008,0.008) (0.406,0.812)rectangle++(0.008,0.008) (0.168,0.939)rectangle++(0.008,0.008) (-0.262,0.908)rectangle++(0.008,0.008)
  (-0.448,0.741)rectangle++(0.008,0.008) (0.059,0.915)rectangle++(0.008,0.008) (0.353,0.813)rectangle++(0.008,0.008) (0.451,0.662)rectangle++(0.008,0.008) (0.187,0.876)rectangle++(0.008,0.008)
  (0.398,0.75)rectangle++(0.008,0.008) (0.165,0.913)rectangle++(0.008,0.008) (0.4,0.775)rectangle++(0.008,0.008) (0.008,0.323)rectangle++(0.008,0.008) (-0.115,0.57)rectangle++(0.008,0.008)
  (0.15,0.718)rectangle++(0.008,0.008) (0.061,0.899)rectangle++(0.008,0.008) (0.348,0.806)rectangle++(0.008,0.008) (0.144,0.936)rectangle++(0.008,0.008) (0.399,0.793)rectangle++(0.008,0.008)
  (0.465,0.636)rectangle++(0.008,0.008) (0.439,0.541)rectangle++(0.008,0.008) (0.393,0.507)rectangle++(0.008,0.008) (-0.005,0.791)rectangle++(0.008,0.008) (-0.004,0.93)rectangle++(0.008,0.008)
  (-0.338,0.843)rectangle++(0.008,0.008) (-0.014,0.352)rectangle++(0.008,0.008) (0.118,0.582)rectangle++(0.008,0.008) (0,0.242)rectangle++(0.008,0.008) (0.086,0.527)rectangle++(0.008,0.008)
  (-0.153,0.691)rectangle++(0.008,0.008) (-0.32,0.68)rectangle++(0.008,0.008) (-0.137,0.883)rectangle++(0.008,0.008) (-0.06,0.969)rectangle++(0.008,0.008) (0.318,0.88)rectangle++(0.008,0.008)
  (-0.173,0.935)rectangle++(0.008,0.008) (-0.009,0.39)rectangle++(0.008,0.008) (-0.147,0.594)rectangle++(0.008,0.008) (-0.064,0.847)rectangle++(0.008,0.008) (0.273,0.827)rectangle++(0.008,0.008)
  (-0.174,0.894)rectangle++(0.008,0.008) (-0.075,0.973)rectangle++(0.008,0.008) (0.312,0.888)rectangle++(0.008,0.008) (-0.178,0.937)rectangle++(0.008,0.008) (-0.009,0.391)rectangle++(0.008,0.008)
  (-0.147,0.595)rectangle++(0.008,0.008) (-0.283,0.638)rectangle++(0.008,0.008) (-0.361,0.61)rectangle++(0.008,0.008) (-0.386,0.569)rectangle++(0.008,0.008) (-0.165,0.837)rectangle++(0.008,0.008)
  (0.223,0.858)rectangle++(0.008,0.008) (0.407,0.729)rectangle++(0.008,0.008) (0.169,0.904)rectangle++(0.008,0.008) (0.069,0.977)rectangle++(0.008,0.008) (0.38,0.839)rectangle++(0.008,0.008)
  (0.472,0.664)rectangle++(0.008,0.008) (-0.025,0.891)rectangle++(0.008,0.008) (-0.333,0.818)rectangle++(0.008,0.008) (-0.448,0.674)rectangle++(0.008,0.008) (0.035,0.884)rectangle++(0.008,0.008)
  (0.331,0.808)rectangle++(0.008,0.008) (0.439,0.667)rectangle++(0.008,0.008) (0.438,0.565)rectangle++(0.008,0.008) (0.401,0.518)rectangle++(0.008,0.008) (0.368,0.51)rectangle++(0.008,0.008)
  (-0.018,0.784)rectangle++(0.008,0.008) (-0.291,0.771)rectangle++(0.008,0.008) (-0.125,0.922)rectangle++(0.008,0.008) (-0.055,0.985)rectangle++(0.008,0.008) (0.326,0.886)rectangle++(0.008,0.008)
  (0.006,0.37)rectangle++(0.008,0.008) (-0.133,0.59)rectangle++(0.008,0.008) (-0.058,0.846)rectangle++(0.008,0.008) (-0.332,0.785)rectangle++(0.008,0.008) (0.128,0.894)rectangle++(0.008,0.008)
  (-0.264,0.873)rectangle++(0.008,0.008) (0.19,0.91)rectangle++(0.008,0.008) (0.078,0.98)rectangle++(0.008,0.008) (-0.318,0.895)rectangle++(0.008,0.008) (-0.013,0.374)rectangle++(0.008,0.008)
  (0.126,0.592)rectangle++(0.008,0.008) (0.051,0.846)rectangle++(0.008,0.008) (0.325,0.785)rectangle++(0.008,0.008) (-0.136,0.894)rectangle++(0.008,0.008) (-0.059,0.973)rectangle++(0.008,0.008)
  (-0.378,0.843)rectangle++(0.008,0.008) (0.127,0.937)rectangle++(0.008,0.008) (0.051,0.991)rectangle++(0.008,0.008) (0.377,0.851)rectangle++(0.008,0.008) (0.156,0.955)rectangle++(0.008,0.008)
  (0.063,0.999)rectangle++(0.008,0.008) (0.385,0.85)rectangle++(0.008,0.008) (0.008,0.355)rectangle++(0.008,0.008) (0.001,0.747)rectangle++(0.008,0.008) (-0.27,0.761)rectangle++(0.008,0.008)
  (-0.399,0.671)rectangle++(0.008,0.008) (-0.17,0.879)rectangle++(0.008,0.008) (-0.395,0.761)rectangle++(0.008,0.008) (-0.168,0.917)rectangle++(0.008,0.008) (0.25,0.896)rectangle++(0.008,0.008)
  (0.103,0.974)rectangle++(0.008,0.008) (-0.304,0.901)rectangle++(0.008,0.008) (-0.464,0.722)rectangle++(0.008,0.008) (-0.197,0.901)rectangle++(0.008,0.008) (-0.085,0.976)rectangle++(0.008,0.008)
  (0.309,0.893)rectangle++(0.008,0.008) (-0.181,0.938)rectangle++(0.008,0.008) (0.251,0.91)rectangle++(0.008,0.008) (0.103,0.98)rectangle++(0.008,0.008) (-0.306,0.904)rectangle++(0.008,0.008)
  (-0.013,0.377)rectangle++(0.008,0.008) (0.128,0.594)rectangle++(0.008,0.008) (0.051,0.847)rectangle++(0.008,0.008) (0.325,0.785)rectangle++(0.008,0.008) (0.134,0.927)rectangle++(0.008,0.008)
  (0.054,0.987)rectangle++(0.008,0.008) (0.02,1.012)rectangle++(0.008,0.008) (0.37,0.872)rectangle++(0.008,0.008) (-0.146,0.95)rectangle++(0.008,0.008) (-0.008,0.397)rectangle++(0.008,0.008)
  (0.137,0.601)rectangle++(0.008,0.008) (0.055,0.85)rectangle++(0.008,0.008) (-0.282,0.827)rectangle++(0.008,0.008) (-0.428,0.697)rectangle++(0.008,0.008) (-0.017,0.29)rectangle++(0.008,0.008)
  (-0.004,0.12)rectangle++(0.008,0.008) (0.04,0.473)rectangle++(0.008,0.008) (-0.154,0.65)rectangle++(0.008,0.008) (-0.009,0.27)rectangle++(0.008,0.008) (-0.104,0.54)rectangle++(0.008,0.008)
  (-0.007,0.224)rectangle++(0.008,0.008) (-0.005,0.692)rectangle++(0.008,0.008) (0.244,0.735)rectangle++(0.008,0.008) (0.373,0.665)rectangle++(0.008,0.008) (-0.071,0.856)rectangle++(0.008,0.008)
  (0.273,0.833)rectangle++(0.008,0.008) (-0.176,0.897)rectangle++(0.008,0.008) (-0.405,0.766)rectangle++(0.008,0.008) (-0.462,0.625)rectangle++(0.008,0.008) (-0.196,0.86)rectangle++(0.008,0.008)
  (0.217,0.88)rectangle++(0.008,0.008) (0.412,0.741)rectangle++(0.008,0.008) (0.452,0.608)rectangle++(0.008,0.008) (-0.014,0.859)rectangle++(0.008,0.008) (-0.008,0.958)rectangle++(0.008,0.008)
  (0.338,0.857)rectangle++(0.008,0.008) (-0.155,0.932)rectangle++(0.008,0.008) (-0.407,0.79)rectangle++(0.008,0.008) (-0.472,0.635)rectangle++(0.008,0.008) (0.01,0.875)rectangle++(0.008,0.008)
  (-0.311,0.823)rectangle++(0.008,0.008) (0.151,0.904)rectangle++(0.008,0.008) (-0.257,0.886)rectangle++(0.008,0.008) (-0.438,0.732)rectangle++(0.008,0.008) (0.061,0.907)rectangle++(0.008,0.008)
  (0.023,0.979)rectangle++(0.008,0.008) (0.359,0.855)rectangle++(0.008,0.008) (-0.145,0.939)rectangle++(0.008,0.008) (0.268,0.897)rectangle++(0.008,0.008) (0.11,0.974)rectangle++(0.008,0.008)
  (0.398,0.822)rectangle++(0.008,0.008) (0.475,0.65)rectangle++(0.008,0.008) (0.197,0.871)rectangle++(0.008,0.008) (0.4,0.744)rectangle++(0.008,0.008) (0.448,0.614)rectangle++(0.008,0.008)
  (0.01,0.255)rectangle++(0.008,0.008) (0.095,0.53)rectangle++(0.008,0.008) (-0.001,0.22)rectangle++(0.008,0.008) (-0.003,0.69)rectangle++(0.008,0.008) (-0.251,0.734)rectangle++(0.008,0.008)
  (-0.011,0.306)rectangle++(0.008,0.008) (-0.118,0.555)rectangle++(0.008,0.008) (0.143,0.712)rectangle++(0.008,0.008) (-0.192,0.796)rectangle++(0.008,0.008) (0.196,0.849)rectangle++(0.008,0.008)
  (0.08,0.954)rectangle++(0.008,0.008) (0.376,0.824)rectangle++(0.008,0.008) (0.007,0.344)rectangle++(0.008,0.008) (-0.123,0.579)rectangle++(0.008,0.008) (-0.267,0.64)rectangle++(0.008,0.008)
  (0.106,0.804)rectangle++(0.008,0.008) (-0.242,0.825)rectangle++(0.008,0.008) (0.183,0.88)rectangle++(0.008,0.008) (-0.234,0.887)rectangle++(0.008,0.008) (-0.427,0.741)rectangle++(0.008,0.008)
  (-0.463,0.605)rectangle++(0.008,0.008) (0.004,0.859)rectangle++(0.008,0.008) (-0.004,0.358)rectangle++(0.008,0.008) (-0.004,0.748)rectangle++(0.008,0.008) (-0.273,0.76)rectangle++(0.008,0.008)
  (-0.399,0.669)rectangle++(0.008,0.008) (0.056,0.865)rectangle++(0.008,0.008) (-0.002,0.361)rectangle++(0.008,0.008) (-0.004,0.149)rectangle++(0.008,0.008) (-0.004,0.66)rectangle++(0.008,0.008)
  (-0.004,0.275)rectangle++(0.008,0.008) (-0.104,0.544)rectangle++(0.008,0.008) (-0.245,0.631)rectangle++(0.008,0.008) (-0.341,0.62)rectangle++(0.008,0.008) (-0.145,0.858)rectangle++(0.008,0.008)
  (0.239,0.861)rectangle++(0.008,0.008) (-0.202,0.898)rectangle++(0.008,0.008) (-0.01,0.375)rectangle++(0.008,0.008) (-0.006,0.755)rectangle++(0.008,0.008) (-0.004,0.315)rectangle++(0.008,0.008)
  (-0.004,0.73)rectangle++(0.008,0.008) (0.258,0.751)rectangle++(0.008,0.008) (-0.154,0.855)rectangle++(0.008,0.008) (-0.067,0.957)rectangle++(0.008,0.008) (-0.03,1)rectangle++(0.008,0.008)
  (0.342,0.884)rectangle++(0.008,0.008) (0.141,0.969)rectangle++(0.008,0.008) (-0.285,0.913)rectangle++(0.008,0.008) (-0.122,0.981)rectangle++(0.008,0.008) (0.294,0.908)rectangle++(0.008,0.008)
  (0.121,0.979)rectangle++(0.008,0.008) (-0.298,0.91)rectangle++(0.008,0.008) (-0.465,0.729)rectangle++(0.008,0.008) (-0.197,0.904)rectangle++(0.008,0.008) (0.232,0.9)rectangle++(0.008,0.008)
  (0.427,0.745)rectangle++(0.008,0.008) (0.46,0.605)rectangle++(0.008,0.008) (-0.009,0.86)rectangle++(0.008,0.008) (-0.315,0.809)rectangle++(0.008,0.008) (0.144,0.899)rectangle++(0.008,0.008)
  (-0.259,0.881)rectangle++(0.008,0.008) (0.196,0.912)rectangle++(0.008,0.008) (0.08,0.981)rectangle++(0.008,0.008) (0.386,0.836)rectangle++(0.008,0.008) (0.16,0.949)rectangle++(0.008,0.008)
  (0.411,0.793)rectangle++(0.008,0.008) (-0.099,0.928)rectangle++(0.008,0.008) (-0.38,0.808)rectangle++(0.008,0.008) (-0.466,0.653)rectangle++(0.008,0.008) (-0.45,0.551)rectangle++(0.008,0.008)
  (-0.406,0.511)rectangle++(0.008,0.008) (-0.173,0.812)rectangle++(0.008,0.008) (-0.009,0.339)rectangle++(0.008,0.008) (-0.004,0.14)rectangle++(0.008,0.008) (-0.004,0.656)rectangle++(0.008,0.008)
  (-0.24,0.718)rectangle++(0.008,0.008) (0.146,0.83)rectangle++(0.008,0.008) (-0.233,0.851)rectangle++(0.008,0.008) (0.196,0.889)rectangle++(0.008,0.008) (-0.231,0.896)rectangle++(0.008,0.008)
  (-0.011,0.374)rectangle++(0.008,0.008) (-0.007,0.755)rectangle++(0.008,0.008) (0.266,0.764)rectangle++(0.008,0.008) (-0.155,0.863)rectangle++(0.008,0.008) (-0.383,0.758)rectangle++(0.008,0.008)
  (-0.449,0.629)rectangle++(0.008,0.008) (-0.191,0.862)rectangle++(0.008,0.008) (-0.083,0.96)rectangle++(0.008,0.008) (-0.384,0.828)rectangle++(0.008,0.008) (0.12,0.932)rectangle++(0.008,0.008)
  (-0.282,0.888)rectangle++(0.008,0.008) (0.188,0.923)rectangle++(0.008,0.008) (-0.247,0.908)rectangle++(0.008,0.008) (-0.441,0.746)rectangle++(0.008,0.008) (-0.471,0.603)rectangle++(0.008,0.008)
  (-0.2,0.851)rectangle++(0.008,0.008) (0.212,0.877)rectangle++(0.008,0.008) (-0.22,0.896)rectangle++(0.008,0.008) (-0.424,0.75)rectangle++(0.008,0.008) (-0.18,0.913)rectangle++(0.008,0.008)
  (-0.009,0.381)rectangle++(0.008,0.008) (-0.006,0.758)rectangle++(0.008,0.008) (0.267,0.765)rectangle++(0.008,0.008) (-0.155,0.864)rectangle++(0.008,0.008) (-0.067,0.961)rectangle++(0.008,0.008)
  (0.312,0.879)rectangle++(0.008,0.008) (-0.175,0.932)rectangle++(0.008,0.008) (0.252,0.905)rectangle++(0.008,0.008) (-0.212,0.923)rectangle++(0.008,0.008) (-0.43,0.766)rectangle++(0.008,0.008)
  (-0.183,0.919)rectangle++(0.008,0.008) (0.244,0.902)rectangle++(0.008,0.008) (0.1,0.976)rectangle++(0.008,0.008) (0.394,0.827)rectangle++(0.008,0.008) (-0.119,0.938)rectangle++(0.008,0.008)
  (0.28,0.887)rectangle++(0.008,0.008) (0.115,0.97)rectangle++(0.008,0.008) (-0.297,0.904)rectangle++(0.008,0.008) (-0.462,0.726)rectangle++(0.008,0.008) (-0.197,0.903)rectangle++(0.008,0.008)
  (0.232,0.899)rectangle++(0.008,0.008) (0.426,0.745)rectangle++(0.008,0.008) (0.177,0.91)rectangle++(0.008,0.008) (-0.248,0.898)rectangle++(0.008,0.008) (-0.011,0.375)rectangle++(0.008,0.008)
  (-0.007,0.755)rectangle++(0.008,0.008) (-0.276,0.762)rectangle++(0.008,0.008) (-0.012,0.318)rectangle++(0.008,0.008) (-0.007,0.731)rectangle++(0.008,0.008) (-0.005,0.905)rectangle++(0.008,0.008)
  (-0.329,0.831)rectangle++(0.008,0.008) (0.145,0.914)rectangle++(0.008,0.008) (0,0.381)rectangle++(0.008,0.008) (-0.002,0.758)rectangle++(0.008,0.008) (-0.275,0.765)rectangle++(0.008,0.008)
  (-0.403,0.671)rectangle++(0.008,0.008) (0.055,0.867)rectangle++(0.008,0.008) (-0.288,0.835)rectangle++(0.008,0.008) (-0.123,0.948)rectangle++(0.008,0.008) (-0.399,0.809)rectangle++(0.008,0.008)
  (-0.475,0.647)rectangle++(0.008,0.008) (0.013,0.881)rectangle++(0.008,0.008) (0.003,0.968)rectangle++(0.008,0.008) (-0.348,0.863)rectangle++(0.008,0.008) (-0.471,0.689)rectangle++(0.008,0.008)
  (-0.018,0.287)rectangle++(0.008,0.008) (-0.114,0.544)rectangle++(0.008,0.008) (0.141,0.706)rectangle++(0.008,0.008) (0.057,0.894)rectangle++(0.008,0.008) (0.022,0.973)rectangle++(0.008,0.008)
  (0.007,1.006)rectangle++(0.008,0.008) (0.362,0.874)rectangle++(0.008,0.008) (0.15,0.965)rectangle++(0.008,0.008) (0.412,0.804)rectangle++(0.008,0.008) (-0.102,0.934)rectangle++(0.008,0.008)
  (-0.384,0.81)rectangle++(0.008,0.008) (-0.164,0.938)rectangle++(0.008,0.008) (-0.413,0.789)rectangle++(0.008,0.008) (-0.474,0.633)rectangle++(0.008,0.008) (-0.202,0.863)rectangle++(0.008,0.008)
  (-0.404,0.742)rectangle++(0.008,0.008) (-0.453,0.614)rectangle++(0.008,0.008) (-0.43,0.538)rectangle++(0.008,0.008) (-0.005,0.816)rectangle++(0.008,0.008) (-0.297,0.79)rectangle++(0.008,0.008)
  (-0.127,0.93)rectangle++(0.008,0.008) (-0.394,0.799)rectangle++(0.008,0.008) (-0.168,0.933)rectangle++(0.008,0.008) (0.256,0.903)rectangle++(0.008,0.008) (-0.209,0.923)rectangle++(0.008,0.008)
  (0.233,0.913)rectangle++(0.008,0.008) (0.003,0.381)rectangle++(0.008,0.008) (-0.001,0.758)rectangle++(0.008,0.008) (-0.003,0.916)rectangle++(0.008,0.008) (-0.332,0.837)rectangle++(0.008,0.008)
  (0.146,0.917)rectangle++(0.008,0.008) (0.001,0.383)rectangle++(0.008,0.008) (-0.002,0.759)rectangle++(0.008,0.008) (-0.275,0.765)rectangle++(0.008,0.008) (0.147,0.864)rectangle++(0.008,0.008)
  (0.001,0.361)rectangle++(0.008,0.008) (0.128,0.581)rectangle++(0.008,0.008) (0.052,0.842)rectangle++(0.008,0.008) (0.323,0.783)rectangle++(0.008,0.008) (0.427,0.659)rectangle++(0.008,0.008)
  (-0.044,0.873)rectangle++(0.008,0.008) (-0.335,0.802)rectangle++(0.008,0.008) (0.133,0.903)rectangle++(0.008,0.008) (0.053,0.977)rectangle++(0.008,0.008) (0.372,0.844)rectangle++(0.008,0.008)
  (-0.135,0.938)rectangle++(0.008,0.008) (-0.4,0.8)rectangle++(0.008,0.008) (-0.472,0.642)rectangle++(0.008,0.008) (-0.449,0.544)rectangle++(0.008,0.008) (-0.011,0.825)rectangle++(0.008,0.008)
  (-0.304,0.792)rectangle++(0.008,0.008) (0.143,0.887)rectangle++(0.008,0.008) (-0.255,0.876)rectangle++(0.008,0.008) (0.196,0.908)rectangle++(0.008,0.008) (0.08,0.979)rectangle++(0.008,0.008)
  (-0.001,0.409)rectangle++(0.008,0.008) (0.145,0.604)rectangle++(0.008,0.008) (0.058,0.851)rectangle++(0.008,0.008) (-0.281,0.829)rectangle++(0.008,0.008) (-0.428,0.698)rectangle++(0.008,0.008)
  (-0.182,0.891)rectangle++(0.008,0.008) (-0.079,0.972)rectangle++(0.008,0.008) (-0.387,0.835)rectangle++(0.008,0.008) (0.121,0.936)rectangle++(0.008,0.008) (0.389,0.801)rectangle++(0.008,0.008)
  (-0.112,0.924)rectangle++(0.008,0.008) (-0.385,0.802)rectangle++(0.008,0.008) (0.11,0.92)rectangle++(0.008,0.008) (0.378,0.798)rectangle++(0.008,0.008) (0.457,0.646)rectangle++(0.008,0.008)
  (0.439,0.549)rectangle++(0.008,0.008) (0.001,0.827)rectangle++(0.008,0.008) (-0.298,0.797)rectangle++(0.008,0.008) (-0.425,0.677)rectangle++(0.008,0.008) (0.047,0.878)rectangle++(0.008,0.008)
  (0.334,0.801)rectangle++(0.008,0.008) (0.006,0.334)rectangle++(0.008,0.008) (0,0.738)rectangle++(0.008,0.008) (-0.267,0.757)rectangle++(0.008,0.008) (-0.012,0.315)rectangle++(0.008,0.008)
  (0.106,0.565)rectangle++(0.008,0.008) (-0.157,0.715)rectangle++(0.008,0.008) (0.183,0.799)rectangle++(0.008,0.008) (-0.205,0.85)rectangle++(0.008,0.008) (0.209,0.878)rectangle++(0.008,0.008)
  (0.085,0.967)rectangle++(0.008,0.008) (0.383,0.828)rectangle++(0.008,0.008) (0.47,0.658)rectangle++(0.008,0.008) (-0.024,0.888)rectangle++(0.008,0.008) (-0.012,0.971)rectangle++(0.008,0.008)
  (-0.356,0.858)rectangle++(0.008,0.008) (-0.152,0.958)rectangle++(0.008,0.008) (0.272,0.909)rectangle++(0.008,0.008) (0.448,0.734)rectangle++(0.008,0.008) (-0.061,0.915)rectangle++(0.008,0.008)
  (-0.028,0.982)rectangle++(0.008,0.008) (-0.014,1.01)rectangle++(0.008,0.008) (-0.008,1.022)rectangle++(0.008,0.008) (0.36,0.886)rectangle++(0.008,0.008) (0.149,0.97)rectangle++(0.008,0.008)
  (-0.282,0.916)rectangle++(0.008,0.008) (-0.459,0.737)rectangle++(0.008,0.008) (-0.195,0.907)rectangle++(0.008,0.008) (-0.01,0.379)rectangle++(0.008,0.008) (0.13,0.593)rectangle++(0.008,0.008)
  (0.052,0.847)rectangle++(0.008,0.008) (-0.002,0.353)rectangle++(0.008,0.008) (-0.003,0.746)rectangle++(0.008,0.008) (-0.272,0.759)rectangle++(0.008,0.008) (0.146,0.86)rectangle++(0.008,0.008)
  (0.373,0.757)rectangle++(0.008,0.008) (0.44,0.629)rectangle++(0.008,0.008) (0.425,0.547)rectangle++(0.008,0.008) (-0.005,0.821)rectangle++(0.008,0.008) (0.29,0.793)rectangle++(0.008,0.008)
  (0.12,0.931)rectangle++(0.008,0.008) (0.048,0.989)rectangle++(0.008,0.008) (0.374,0.851)rectangle++(0.008,0.008) (-0.137,0.942)rectangle++(0.008,0.008) (0.273,0.896)rectangle++(0.008,0.008)
  (0.004,0.374)rectangle++(0.008,0.008) (0.135,0.586)rectangle++(0.008,0.008) (0.27,0.636)rectangle++(0.008,0.008) (0.111,0.865)rectangle++(0.008,0.008) (0.359,0.772)rectangle++(0.008,0.008)
  (0.148,0.922)rectangle++(0.008,0.008) (0.396,0.785)rectangle++(0.008,0.008) (0.164,0.927)rectangle++(0.008,0.008) (-0.26,0.902)rectangle++(0.008,0.008) (-0.111,0.976)rectangle++(0.008,0.008)
  (0.297,0.902)rectangle++(0.008,0.008) (-0.19,0.938)rectangle++(0.008,0.008) (0.247,0.913)rectangle++(0.008,0.008) (0.004,0.381)rectangle++(0.008,0.008) (-0.138,0.595)rectangle++(0.008,0.008)
  (-0.279,0.642)rectangle++(0.008,0.008) (-0.36,0.613)rectangle++(0.008,0.008) (-0.154,0.855)rectangle++(0.008,0.008) (-0.067,0.957)rectangle++(0.008,0.008) (0.31,0.878)rectangle++(0.008,0.008)
  (-0.175,0.931)rectangle++(0.008,0.008) (-0.076,0.989)rectangle++(0.008,0.008) (0.318,0.895)rectangle++(0.008,0.008) (-0.178,0.942)rectangle++(0.008,0.008) (0.254,0.911)rectangle++(0.008,0.008)
  (0.441,0.742)rectangle++(0.008,0.008) (-0.067,0.916)rectangle++(0.008,0.008) (-0.031,0.982)rectangle++(0.008,0.008) (0.336,0.876)rectangle++(0.008,0.008) (-0.163,0.94)rectangle++(0.008,0.008)
  (-0.071,0.992)rectangle++(0.008,0.008) (-0.39,0.848)rectangle++(0.008,0.008) (-0.166,0.954)rectangle++(0.008,0.008) (-0.42,0.796)rectangle++(0.008,0.008) (-0.016,0.332)rectangle++(0.008,0.008)
  (0.11,0.574)rectangle++(0.008,0.008) (0.044,0.839)rectangle++(0.008,0.008) (-0.283,0.818)rectangle++(0.008,0.008) (0.162,0.891)rectangle++(0.008,0.008) (0.392,0.766)rectangle++(0.008,0.008)
  (0.452,0.627)rectangle++(0.008,0.008) (0.429,0.541)rectangle++(0.008,0.008) (0.389,0.511)rectangle++(0.008,0.008) (0.008,0.212)rectangle++(0.008,0.008) (0.001,0.687)rectangle++(0.008,0.008)
  (-0.002,0.886)rectangle++(0.008,0.008) (-0.003,0.97)rectangle++(0.008,0.008) (-0.351,0.861)rectangle++(0.008,0.008) (-0.15,0.959)rectangle++(0.008,0.008) (-0.415,0.804)rectangle++(0.008,0.008)
  (-0.177,0.935)rectangle++(0.008,0.008) (-0.009,0.391)rectangle++(0.008,0.008) (0.135,0.598)rectangle++(0.008,0.008) (0,0.249)rectangle++(0.008,0.008) (-0.004,0.102)rectangle++(0.008,0.008)
  (0.034,0.465)rectangle++(0.008,0.008) (-0.154,0.644)rectangle++(0.008,0.008) (0.159,0.766)rectangle++(0.008,0.008) (0.345,0.71)rectangle++(0.008,0.008) (0.143,0.896)rectangle++(0.008,0.008)
  (0.058,0.974)rectangle++(0.008,0.008) (0.022,1.007)rectangle++(0.008,0.008) (0.007,1.02)rectangle++(0.008,0.008) (0.367,0.88)rectangle++(0.008,0.008) (0.481,0.688)rectangle++(0.008,0.008)
  (0.2,0.887)rectangle++(0.008,0.008) (-0.229,0.896)rectangle++(0.008,0.008) (-0.011,0.374)rectangle++(0.008,0.008) (-0.004,0.155)rectangle++(0.008,0.008) (-0.004,0.063)rectangle++(0.008,0.008)
  (-0.028,0.446)rectangle++(0.008,0.008) (-0.176,0.613)rectangle++(0.008,0.008) (-0.303,0.637)rectangle++(0.008,0.008) (-0.369,0.602)rectangle++(0.008,0.008) (-0.387,0.563)rectangle++(0.008,0.008)
  (-0.382,0.538)rectangle++(0.008,0.008) (0.017,0.799)rectangle++(0.008,0.008) (0.292,0.775)rectangle++(0.008,0.008) (0.12,0.923)rectangle++(0.008,0.008) (-0.278,0.884)rectangle++(0.008,0.008)
  (0.188,0.92)rectangle++(0.008,0.008) (-0.246,0.907)rectangle++(0.008,0.008) (-0.44,0.746)rectangle++(0.008,0.008) (-0.471,0.603)rectangle++(0.008,0.008) (-0.434,0.527)rectangle++(0.008,0.008)
  (-0.185,0.819)rectangle++(0.008,0.008) (0.207,0.857)rectangle++(0.008,0.008) (-0.215,0.885)rectangle++(0.008,0.008) (0.217,0.897)rectangle++(0.008,0.008) (-0.225,0.907)rectangle++(0.008,0.008)
  (0.22,0.911)rectangle++(0.008,0.008) (0.003,0.38)rectangle++(0.008,0.008) (-0.004,0.157)rectangle++(0.008,0.008) (-0.062,0.49)rectangle++(0.008,0.008) (-0.207,0.621)rectangle++(0.008,0.008)
  (-0.32,0.629)rectangle++(0.008,0.008) (-0.375,0.593)rectangle++(0.008,0.008) (-0.16,0.847)rectangle++(0.008,0.008) (0.229,0.86)rectangle++(0.008,0.008) (0.411,0.728)rectangle++(0.008,0.008)
  (0.17,0.903)rectangle++(0.008,0.008) (0.4,0.769)rectangle++(0.008,0.008) (0.008,0.32)rectangle++(0.008,0.008) (0.117,0.56)rectangle++(0.008,0.008) (0.253,0.63)rectangle++(0.008,0.008)
  (0.104,0.862)rectangle++(0.008,0.008) (-0.264,0.851)rectangle++(0.008,0.008) (-0.113,0.955)rectangle++(0.008,0.008) (0.288,0.893)rectangle++(0.008,0.008) (0.45,0.722)rectangle++(0.008,0.008)
  (0.187,0.901)rectangle++(0.008,0.008) (0.002,0.376)rectangle++(0.008,0.008) (-0.137,0.592)rectangle++(0.008,0.008) (-0.06,0.846)rectangle++(0.008,0.008) (0.274,0.825)rectangle++(0.008,0.008)
  (-0.173,0.894)rectangle++(0.008,0.008) (0.24,0.887)rectangle++(0.008,0.008) (-0.211,0.91)rectangle++(0.008,0.008) (0.228,0.908)rectangle++(0.008,0.008) (0.093,0.979)rectangle++(0.008,0.008)
  (0.391,0.83)rectangle++(0.008,0.008) (0.475,0.656)rectangle++(0.008,0.008) (-0.021,0.889)rectangle++(0.008,0.008) (0.307,0.83)rectangle++(0.008,0.008) (0.436,0.686)rectangle++(0.008,0.008)
  (0.444,0.574)rectangle++(0.008,0.008) (0.184,0.839)rectangle++(0.008,0.008) (-0.219,0.868)rectangle++(0.008,0.008) (-0.094,0.962)rectangle++(0.008,0.008) (0.3,0.89)rectangle++(0.008,0.008)
  (-0.184,0.933)rectangle++(0.008,0.008) (-0.08,0.99)rectangle++(0.008,0.008) (0.316,0.897)rectangle++(0.008,0.008) (0.464,0.714)rectangle++(0.008,0.008) (0.466,0.577)rectangle++(0.008,0.008)
  (0.193,0.84)rectangle++(0.008,0.008) (-0.215,0.872)rectangle++(0.008,0.008) (-0.413,0.741)rectangle++(0.008,0.008) (-0.016,0.309)rectangle++(0.008,0.008) (0.102,0.563)rectangle++(0.008,0.008)
  (0.041,0.834)rectangle++(0.008,0.008) (0.015,0.948)rectangle++(0.008,0.008) (0.345,0.844)rectangle++(0.008,0.008) (-0.148,0.928)rectangle++(0.008,0.008) (-0.064,0.988)rectangle++(0.008,0.008)
  (-0.006,0.412)rectangle++(0.008,0.008) (-0.154,0.606)rectangle++(0.008,0.008) (-0.067,0.852)rectangle++(0.008,0.008) (-0.03,0.956)rectangle++(0.008,0.008) (-0.359,0.845)rectangle++(0.008,0.008)
  (-0.153,0.953)rectangle++(0.008,0.008) (-0.414,0.8)rectangle++(0.008,0.008) (-0.478,0.637)rectangle++(0.008,0.008) (0.008,0.878)rectangle++(0.008,0.008) (-0.314,0.824)rectangle++(0.008,0.008)
  (0.15,0.905)rectangle++(0.008,0.008) (-0.258,0.886)rectangle++(0.008,0.008) (0.198,0.914)rectangle++(0.008,0.008) (0.081,0.981)rectangle++(0.008,0.008) (-0.317,0.897)rectangle++(0.008,0.008)
  (-0.136,0.974)rectangle++(0.008,0.008) (-0.413,0.816)rectangle++(0.008,0.008) (-0.176,0.94)rectangle++(0.008,0.008) (-0.42,0.786)rectangle++(0.008,0.008) (-0.476,0.629)rectangle++(0.008,0.008)
  (-0.446,0.537)rectangle++(0.008,0.008) (-0.013,0.821)rectangle++(0.008,0.008) (-0.303,0.79)rectangle++(0.008,0.008) (-0.424,0.672)rectangle++(0.008,0.008) (-0.437,0.575)rectangle++(0.008,0.008)
  (-0.409,0.526)rectangle++(0.008,0.008) (-0.174,0.819)rectangle++(0.008,0.008) (-0.375,0.731)rectangle++(0.008,0.008) (0.089,0.885)rectangle++(0.008,0.008) (0.356,0.789)rectangle++(0.008,0.008)
  (-0.123,0.907)rectangle++(0.008,0.008) (-0.054,0.979)rectangle++(0.008,0.008) (0.324,0.883)rectangle++(0.008,0.008) (-0.171,0.938)rectangle++(0.008,0.008) (-0.074,0.992)rectangle++(0.008,0.008)
  (0.32,0.896)rectangle++(0.008,0.008) (0.465,0.712)rectangle++(0.008,0.008) (-0.045,0.911)rectangle++(0.008,0.008) (-0.349,0.819)rectangle++(0.008,0.008) (0.132,0.916)rectangle++(0.008,0.008)
  (0.387,0.788)rectangle++(0.008,0.008) (0.008,0.329)rectangle++(0.008,0.008) (0.001,0.736)rectangle++(0.008,0.008) (-0.266,0.756)rectangle++(0.008,0.008) (0.148,0.857)rectangle++(0.008,0.008)
  (0.373,0.755)rectangle++(0.008,0.008) (0.439,0.629)rectangle++(0.008,0.008) (0.424,0.547)rectangle++(0.008,0.008) (0.176,0.827)rectangle++(0.008,0.008) (0.072,0.945)rectangle++(0.008,0.008)
  (0.028,0.995)rectangle++(0.008,0.008) (0.367,0.861)rectangle++(0.008,0.008) (0.007,0.359)rectangle++(0.008,0.008) (0.131,0.578)rectangle++(0.008,0.008) (0.053,0.84)rectangle++(0.008,0.008)
  (0.02,0.951)rectangle++(0.008,0.008) (-0.334,0.861)rectangle++(0.008,0.008) (0.154,0.929)rectangle++(0.008,0.008) (0.401,0.786)rectangle++(0.008,0.008) (0.463,0.632)rectangle++(0.008,0.008)
  (-0.018,0.874)rectangle++(0.008,0.008) (-0.01,0.965)rectangle++(0.008,0.008) (0.339,0.861)rectangle++(0.008,0.008) (0.14,0.959)rectangle++(0.008,0.008) (0.406,0.805)rectangle++(0.008,0.008)
  (0.168,0.936)rectangle++(0.008,0.008) (-0.261,0.907)rectangle++(0.008,0.008) (-0.112,0.979)rectangle++(0.008,0.008) (-0.049,1.009)rectangle++(0.008,0.008) (-0.386,0.863)rectangle++(0.008,0.008)
  (0.131,0.949)rectangle++(0.008,0.008) (0.398,0.803)rectangle++(0.008,0.008) (0.468,0.642)rectangle++(0.008,0.008) (0.194,0.867)rectangle++(0.008,0.008) (0.398,0.743)rectangle++(0.008,0.008)
  (0.446,0.614)rectangle++(0.008,0.008) (-0.019,0.859)rectangle++(0.008,0.008) (0.298,0.816)rectangle++(0.008,0.008) (-0.159,0.898)rectangle++(0.008,0.008) (-0.397,0.773)rectangle++(0.008,0.008)
  (0.094,0.912)rectangle++(0.008,0.008) (0.368,0.8)rectangle++(0.008,0.008) (-0.121,0.916)rectangle++(0.008,0.008) (0.271,0.878)rectangle++(0.008,0.008) (0.111,0.967)rectangle++(0.008,0.008)
  (0.045,1.004)rectangle++(0.008,0.008) (-0.342,0.894)rectangle++(0.008,0.008) (0.162,0.947)rectangle++(0.008,0.008) (0.066,0.995)rectangle++(0.008,0.008) (0.025,1.016)rectangle++(0.008,0.008)
  (0.374,0.872)rectangle++(0.008,0.008) (-0.144,0.951)rectangle++(0.008,0.008) (-0.063,0.997)rectangle++(0.008,0.008) (0.327,0.895)rectangle++(0.008,0.008) (0.468,0.709)rectangle++(0.008,0.008)
  (0.466,0.573)rectangle++(0.008,0.008) (0.005,0.848)rectangle++(0.008,0.008) (-0.304,0.808)rectangle++(0.008,0.008) (-0.013,0.337)rectangle++(0.008,0.008) (0.114,0.575)rectangle++(0.008,0.008)
  (0.257,0.639)rectangle++(0.008,0.008) (0.345,0.617)rectangle++(0.008,0.008) (-0.066,0.824)rectangle++(0.008,0.008) (-0.03,0.944)rectangle++(0.008,0.008) (-0.005,0.394)rectangle++(0.008,0.008)
  (-0.147,0.598)rectangle++(0.008,0.008) (0.146,0.742)rectangle++(0.008,0.008) (0.331,0.703)rectangle++(0.008,0.008) (0.402,0.62)rectangle++(0.008,0.008) (-0.041,0.846)rectangle++(0.008,0.008)
  (0.282,0.818)rectangle++(0.008,0.008) (0.116,0.941)rectangle++(0.008,0.008) (0.388,0.805)rectangle++(0.008,0.008) (0.161,0.936)rectangle++(0.008,0.008) (0.065,0.991)rectangle++(0.008,0.008)
  (0.383,0.846)rectangle++(0.008,0.008) (-0.131,0.943)rectangle++(0.008,0.008) (-0.4,0.803)rectangle++(0.008,0.008) (-0.17,0.935)rectangle++(0.008,0.008) (0.255,0.905)rectangle++(0.008,0.008)
  (-0.21,0.924)rectangle++(0.008,0.008) (0.233,0.914)rectangle++(0.008,0.008) (0.432,0.751)rectangle++(0.008,0.008) (0.465,0.605)rectangle++(0.008,0.008) (-0.007,0.862)rectangle++(0.008,0.008)
  (-0.004,0.36)rectangle++(0.008,0.008) (-0.004,0.149)rectangle++(0.008,0.008) (-0.059,0.486)rectangle++(0.008,0.008) (-0.006,0.202)rectangle++(0.008,0.008) (-0.078,0.509)rectangle++(0.008,0.008)
  (-0.221,0.624)rectangle++(0.008,0.008) (0.121,0.781)rectangle++(0.008,0.008) (-0.227,0.819)rectangle++(0.008,0.008) (0.188,0.872)rectangle++(0.008,0.008) (0.002,0.364)rectangle++(0.008,0.008)
  (-0.133,0.586)rectangle++(0.008,0.008) (-0.274,0.64)rectangle++(0.008,0.008) (-0.357,0.614)rectangle++(0.008,0.008) (0.055,0.825)rectangle++(0.008,0.008) (0.319,0.774)rectangle++(0.008,0.008)
  (-0.134,0.887)rectangle++(0.008,0.008) (-0.059,0.97)rectangle++(0.008,0.008) (-0.377,0.842)rectangle++(0.008,0.008) (0.128,0.936)rectangle++(0.008,0.008) (0.392,0.798)rectangle++(0.008,0.008)
  (0.463,0.641)rectangle++(0.008,0.008) (-0.021,0.878)rectangle++(0.008,0.008) (-0.327,0.813)rectangle++(0.008,0.008) (-0.014,0.339)rectangle++(0.008,0.008) (-0.131,0.569)rectangle++(0.008,0.008)
  (-0.267,0.633)rectangle++(0.008,0.008) (0.103,0.801)rectangle++(0.008,0.008) (0.041,0.934)rectangle++(0.008,0.008) (0.015,0.99)rectangle++(0.008,0.008) (-0.35,0.877)rectangle++(0.008,0.008)
  (-0.477,0.695)rectangle++(0.008,0.008) (0.029,0.904)rectangle++(0.008,0.008) (0.336,0.819)rectangle++(0.008,0.008) (0.445,0.671)rectangle++(0.008,0.008) (0.442,0.564)rectangle++(0.008,0.008)
  (0.183,0.835)rectangle++(0.008,0.008) (0.075,0.948)rectangle++(0.008,0.008) (0.372,0.823)rectangle++(0.008,0.008) (0.463,0.66)rectangle++(0.008,0.008) (0.447,0.553)rectangle++(0.008,0.008)
  (0.003,0.831)rectangle++(0.008,0.008) (-0.001,0.947)rectangle++(0.008,0.008) (-0.003,0.995)rectangle++(0.008,0.008) (-0.36,0.873)rectangle++(0.008,0.008) (-0.015,0.364)rectangle++(0.008,0.008)
  (-0.008,0.751)rectangle++(0.008,0.008) (0.263,0.763)rectangle++(0.008,0.008) (0.004,0.318)rectangle++(0.008,0.008) (-0.001,0.731)rectangle++(0.008,0.008) (-0.265,0.753)rectangle++(0.008,0.008)
  (0.147,0.855)rectangle++(0.008,0.008) (-0.242,0.863)rectangle++(0.008,0.008) (-0.422,0.727)rectangle++(0.008,0.008) (-0.18,0.903)rectangle++(0.008,0.008) (0.24,0.893)rectangle++(0.008,0.008)
  (0.098,0.973)rectangle++(0.008,0.008) (0.392,0.826)rectangle++(0.008,0.008) (0.473,0.654)rectangle++(0.008,0.008) (0.449,0.546)rectangle++(0.008,0.008) (0.4,0.506)rectangle++(0.008,0.008)
  (0.362,0.505)rectangle++(0.008,0.008) (-0.018,0.779)rectangle++(0.008,0.008) (-0.01,0.925)rectangle++(0.008,0.008) (-0.004,0.386)rectangle++(0.008,0.008) (-0.143,0.594)rectangle++(0.008,0.008)
  (0.146,0.739)rectangle++(0.008,0.008) (0,0.308)rectangle++(0.008,0.008) (0.11,0.557)rectangle++(0.008,0.008) (-0.152,0.713)rectangle++(0.008,0.008) (-0.008,0.297)rectangle++(0.008,0.008)
  (-0.006,0.722)rectangle++(0.008,0.008) (-0.005,0.901)rectangle++(0.008,0.008) (0.319,0.83)rectangle++(0.008,0.008) (-0.154,0.912)rectangle++(0.008,0.008) (0.255,0.889)rectangle++(0.008,0.008)
  (0.433,0.732)rectangle++(0.008,0.008) (-0.067,0.908)rectangle++(0.008,0.008) (-0.03,0.979)rectangle++(0.008,0.008) (-0.367,0.856)rectangle++(0.008,0.008) (-0.477,0.679)rectangle++(0.008,0.008)
  (0.024,0.897)rectangle++(0.008,0.008) (0.008,0.975)rectangle++(0.008,0.008) (0.351,0.859)rectangle++(0.008,0.008) (-0.15,0.937)rectangle++(0.008,0.008) (-0.065,0.991)rectangle++(0.008,0.008)
  (0.323,0.893)rectangle++(0.008,0.008) (0.466,0.709)rectangle++(0.008,0.008) (0.465,0.574)rectangle++(0.008,0.008) (0.004,0.848)rectangle++(0.008,0.008) (-0.004,0.354)rectangle++(0.008,0.008);
\fill[nuage]
  (-0.004,0.746)rectangle++(0.008,0.008) (-0.272,0.759)rectangle++(0.008,0.008) (0.146,0.86)rectangle++(0.008,0.008) (-0.244,0.865)rectangle++(0.008,0.008) (-0.105,0.961)rectangle++(0.008,0.008)
  (-0.395,0.821)rectangle++(0.008,0.008) (0.112,0.933)rectangle++(0.008,0.008) (-0.285,0.885)rectangle++(0.008,0.008) (-0.122,0.97)rectangle++(0.008,0.008) (-0.054,1.005)rectangle++(0.008,0.008)
  (-0.025,1.02)rectangle++(0.008,0.008) (-0.013,1.026)rectangle++(0.008,0.008) (-0.376,0.884)rectangle++(0.008,0.008) (-0.16,0.969)rectangle++(0.008,0.008) (0.272,0.916)rectangle++(0.008,0.008)
  (0.004,0.383)rectangle++(0.008,0.008) (-0.138,0.596)rectangle++(0.008,0.008) (-0.279,0.642)rectangle++(0.008,0.008) (-0.012,0.267)rectangle++(0.008,0.008) (0.089,0.543)rectangle++(0.008,0.008)
  (0.234,0.633)rectangle++(0.008,0.008) (-0.123,0.792)rectangle++(0.008,0.008) (-0.342,0.737)rectangle++(0.008,0.008) (0.106,0.876)rectangle++(0.008,0.008) (-0.268,0.857)rectangle++(0.008,0.008)
  (0.183,0.904)rectangle++(0.008,0.008) (0.406,0.764)rectangle++(0.008,0.008) (0.168,0.919)rectangle++(0.008,0.008) (0.068,0.984)rectangle++(0.008,0.008) (-0.324,0.893)rectangle++(0.008,0.008)
  (-0.47,0.712)rectangle++(0.008,0.008) (-0.473,0.577)rectangle++(0.008,0.008) (-0.011,0.849)rectangle++(0.008,0.008) (-0.312,0.803)rectangle++(0.008,0.008) (-0.433,0.675)rectangle++(0.008,0.008)
  (-0.017,0.281)rectangle++(0.008,0.008) (0.092,0.551)rectangle++(0.008,0.008) (0.238,0.635)rectangle++(0.008,0.008) (-0.122,0.795)rectangle++(0.008,0.008) (0.227,0.823)rectangle++(0.008,0.008)
  (0.397,0.711)rectangle++(0.008,0.008) (0.164,0.896)rectangle++(0.008,0.008) (0.395,0.767)rectangle++(0.008,0.008) (0.454,0.626)rectangle++(0.008,0.008) (0.01,0.261)rectangle++(0.008,0.008)
  (0.097,0.532)rectangle++(0.008,0.008) (0.038,0.821)rectangle++(0.008,0.008) (-0.279,0.808)rectangle++(0.008,0.008) (-0.12,0.937)rectangle++(0.008,0.008) (0.279,0.887)rectangle++(0.008,0.008)
  (-0.193,0.925)rectangle++(0.008,0.008) (0.241,0.908)rectangle++(0.008,0.008) (0.434,0.745)rectangle++(0.008,0.008) (0.464,0.602)rectangle++(0.008,0.008) (-0.007,0.86)rectangle++(0.008,0.008)
  (-0.004,0.359)rectangle++(0.008,0.008) (0.126,0.582)rectangle++(0.008,0.008) (0.264,0.638)rectangle++(0.008,0.008) (0.004,0.265)rectangle++(0.008,0.008) (0.096,0.536)rectangle++(0.008,0.008)
  (-0.151,0.699)rectangle++(0.008,0.008) (-0.066,0.891)rectangle++(0.008,0.008) (0.287,0.847)rectangle++(0.008,0.008) (0.118,0.953)rectangle++(0.008,0.008) (0.394,0.81)rectangle++(0.008,0.008)
  (-0.113,0.93)rectangle++(0.008,0.008) (0.28,0.882)rectangle++(0.008,0.008) (0.005,0.368)rectangle++(0.008,0.008) (0.133,0.583)rectangle++(0.008,0.008) (-0.151,0.733)rectangle++(0.008,0.008)
  (-0.334,0.7)rectangle++(0.008,0.008) (-0.407,0.62)rectangle++(0.008,0.008) (-0.016,0.258)rectangle++(0.008,0.008) (0.084,0.54)rectangle++(0.008,0.008) (-0.158,0.696)rectangle++(0.008,0.008)
  (-0.069,0.89)rectangle++(0.008,0.008) (-0.006,0.371)rectangle++(0.008,0.008) (0.129,0.588)rectangle++(0.008,0.008) (0.052,0.845)rectangle++(0.008,0.008) (-0.282,0.824)rectangle++(0.008,0.008)
  (0.165,0.893)rectangle++(0.008,0.008) (0.067,0.973)rectangle++(0.008,0.008) (-0.32,0.888)rectangle++(0.008,0.008) (0.17,0.937)rectangle++(0.008,0.008) (0.411,0.784)rectangle++(0.008,0.008)
  (-0.096,0.924)rectangle++(0.008,0.008) (-0.042,0.986)rectangle++(0.008,0.008) (0.332,0.882)rectangle++(0.008,0.008) (-0.167,0.941)rectangle++(0.008,0.008) (0.259,0.906)rectangle++(0.008,0.008)
  (0.441,0.738)rectangle++(0.008,0.008) (-0.066,0.914)rectangle++(0.008,0.008) (-0.36,0.814)rectangle++(0.008,0.008) (0.125,0.917)rectangle++(0.008,0.008) (0.05,0.983)rectangle++(0.008,0.008)
  (0.019,1.01)rectangle++(0.008,0.008) (0.369,0.871)rectangle++(0.008,0.008) (0.479,0.683)rectangle++(0.008,0.008) (0.01,0.285)rectangle++(0.008,0.008) (-0.1,0.553)rectangle++(0.008,0.008)
  (-0.044,0.83)rectangle++(0.008,0.008) (-0.32,0.783)rectangle++(0.008,0.008) (-0.137,0.926)rectangle++(0.008,0.008) (0.268,0.889)rectangle++(0.008,0.008) (0.439,0.727)rectangle++(0.008,0.008)
  (0.182,0.903)rectangle++(0.008,0.008) (0.074,0.977)rectangle++(0.008,0.008) (-0.319,0.892)rectangle++(0.008,0.008) (0.172,0.938)rectangle++(0.008,0.008) (0.413,0.784)rectangle++(0.008,0.008)
  (0.171,0.927)rectangle++(0.008,0.008) (-0.256,0.904)rectangle++(0.008,0.008) (-0.11,0.977)rectangle++(0.008,0.008) (-0.403,0.827)rectangle++(0.008,0.008) (-0.016,0.345)rectangle++(0.008,0.008)
  (-0.134,0.571)rectangle++(0.008,0.008) (-0.269,0.632)rectangle++(0.008,0.008) (-0.012,0.263)rectangle++(0.008,0.008) (0.088,0.541)rectangle++(0.008,0.008) (0.233,0.632)rectangle++(0.008,0.008)
  (0.095,0.863)rectangle++(0.008,0.008) (-0.268,0.848)rectangle++(0.008,0.008) (-0.429,0.711)rectangle++(0.008,0.008) (0.057,0.895)rectangle++(0.008,0.008) (0.345,0.805)rectangle++(0.008,0.008)
  (0.444,0.661)rectangle++(0.008,0.008) (-0.037,0.88)rectangle++(0.008,0.008) (-0.018,0.967)rectangle++(0.008,0.008) (-0.357,0.855)rectangle++(0.008,0.008) (-0.472,0.683)rectangle++(0.008,0.008)
  (-0.463,0.563)rectangle++(0.008,0.008) (-0.416,0.511)rectangle++(0.008,0.008) (-0.376,0.504)rectangle++(0.008,0.008) (-0.356,0.515)rectangle++(0.008,0.008) (-0.152,0.814)rectangle++(0.008,0.008)
  (0.221,0.843)rectangle++(0.008,0.008) (0.003,0.352)rectangle++(0.008,0.008) (-0.128,0.581)rectangle++(0.008,0.008) (-0.269,0.639)rectangle++(0.008,0.008) (0.104,0.805)rectangle++(0.008,0.008)
  (-0.243,0.824)rectangle++(0.008,0.008) (-0.105,0.944)rectangle++(0.008,0.008) (0.289,0.885)rectangle++(0.008,0.008) (-0.188,0.927)rectangle++(0.008,0.008) (-0.42,0.776)rectangle++(0.008,0.008)
  (-0.473,0.624)rectangle++(0.008,0.008) (-0.442,0.536)rectangle++(0.008,0.008) (-0.397,0.506)rectangle++(0.008,0.008) (-0.002,0.789)rectangle++(0.008,0.008) (0.28,0.778)rectangle++(0.008,0.008)
  (0.115,0.924)rectangle++(0.008,0.008) (0.046,0.986)rectangle++(0.008,0.008) (0.017,1.012)rectangle++(0.008,0.008) (-0.357,0.888)rectangle++(0.008,0.008) (-0.484,0.698)rectangle++(0.008,0.008)
  (-0.474,0.565)rectangle++(0.008,0.008) (-0.015,0.844)rectangle++(0.008,0.008) (0.294,0.808)rectangle++(0.008,0.008) (-0.158,0.893)rectangle++(0.008,0.008) (-0.069,0.973)rectangle++(0.008,0.008)
  (-0.382,0.839)rectangle++(0.008,0.008) (-0.163,0.95)rectangle++(0.008,0.008) (-0.417,0.795)rectangle++(0.008,0.008) (-0.016,0.332)rectangle++(0.008,0.008) (0.11,0.574)rectangle++(0.008,0.008)
  (-0.001,0.239)rectangle++(0.008,0.008) (0.084,0.526)rectangle++(0.008,0.008) (-0.001,0.218)rectangle++(0.008,0.008) (0.077,0.517)rectangle++(0.008,0.008) (-0.153,0.683)rectangle++(0.008,0.008)
  (-0.317,0.677)rectangle++(0.008,0.008) (-0.39,0.615)rectangle++(0.008,0.008) (-0.402,0.561)rectangle++(0.008,0.008) (-0.171,0.833)rectangle++(0.008,0.008) (-0.379,0.739)rectangle++(0.008,0.008)
  (0.09,0.89)rectangle++(0.008,0.008) (0.358,0.791)rectangle++(0.008,0.008) (0.148,0.93)rectangle++(0.008,0.008) (0.399,0.789)rectangle++(0.008,0.008) (0.463,0.634)rectangle++(0.008,0.008)
  (0.192,0.864)rectangle++(0.008,0.008) (0.396,0.743)rectangle++(0.008,0.008) (-0.088,0.9)rectangle++(0.008,0.008) (0.28,0.859)rectangle++(0.008,0.008) (0.434,0.709)rectangle++(0.008,0.008)
  (0.009,0.295)rectangle++(0.008,0.008) (0.109,0.548)rectangle++(0.008,0.008) (-0.15,0.709)rectangle++(0.008,0.008) (-0.065,0.895)rectangle++(0.008,0.008) (-0.006,0.374)rectangle++(0.008,0.008)
  (-0.14,0.588)rectangle++(0.008,0.008) (0.145,0.735)rectangle++(0.008,0.008) (0,0.306)rectangle++(0.008,0.008) (0.109,0.556)rectangle++(0.008,0.008) (0.043,0.831)rectangle++(0.008,0.008)
  (0.016,0.947)rectangle++(0.008,0.008) (0.004,0.995)rectangle++(0.008,0.008) (0.357,0.87)rectangle++(0.008,0.008) (-0.151,0.944)rectangle++(0.008,0.008) (-0.41,0.797)rectangle++(0.008,0.008)
  (-0.174,0.932)rectangle++(0.008,0.008) (0.252,0.905)rectangle++(0.008,0.008) (-0.211,0.923)rectangle++(0.008,0.008) (-0.43,0.766)rectangle++(0.008,0.008) (-0.183,0.919)rectangle++(0.008,0.008)
  (-0.009,0.384)rectangle++(0.008,0.008) (-0.006,0.759)rectangle++(0.008,0.008) (0.267,0.765)rectangle++(0.008,0.008) (-0.155,0.865)rectangle++(0.008,0.008) (0.237,0.867)rectangle++(0.008,0.008)
  (-0.205,0.9)rectangle++(0.008,0.008) (-0.088,0.976)rectangle++(0.008,0.008) (0.307,0.894)rectangle++(0.008,0.008) (0.459,0.715)rectangle++(0.008,0.008) (-0.049,0.91)rectangle++(0.008,0.008)
  (0.302,0.85)rectangle++(0.008,0.008) (-0.169,0.915)rectangle++(0.008,0.008) (-0.073,0.982)rectangle++(0.008,0.008) (0.316,0.891)rectangle++(0.008,0.008) (0.131,0.972)rectangle++(0.008,0.008)
  (-0.291,0.91)rectangle++(0.008,0.008) (0.191,0.936)rectangle++(0.008,0.008) (0.078,0.991)rectangle++(0.008,0.008) (-0.322,0.9)rectangle++(0.008,0.008) (-0.137,0.976)rectangle++(0.008,0.008)
  (0.285,0.911)rectangle++(0.008,0.008) (0.455,0.731)rectangle++(0.008,0.008) (-0.057,0.916)rectangle++(0.008,0.008) (-0.357,0.818)rectangle++(0.008,0.008) (0.128,0.917)rectangle++(0.008,0.008)
  (-0.273,0.884)rectangle++(0.008,0.008) (0.19,0.918)rectangle++(0.008,0.008) (0.414,0.768)rectangle++(0.008,0.008) (0.172,0.92)rectangle++(0.008,0.008) (0.406,0.776)rectangle++(0.008,0.008)
  (-0.095,0.919)rectangle++(0.008,0.008) (-0.042,0.984)rectangle++(0.008,0.008) (-0.005,0.411)rectangle++(0.008,0.008) (0.144,0.606)rectangle++(0.008,0.008) (-0.154,0.747)rectangle++(0.008,0.008)
  (0.196,0.813)rectangle++(0.008,0.008) (0.379,0.718)rectangle++(0.008,0.008) (0.429,0.609)rectangle++(0.008,0.008) (0.009,0.254)rectangle++(0.008,0.008) (0.094,0.529)rectangle++(0.008,0.008)
  (0.037,0.82)rectangle++(0.008,0.008) (0.013,0.942)rectangle++(0.008,0.008) (-0.334,0.855)rectangle++(0.008,0.008) (0.152,0.926)rectangle++(0.008,0.008) (0.001,0.387)rectangle++(0.008,0.008)
  (0.138,0.593)rectangle++(0.008,0.008) (0.274,0.638)rectangle++(0.008,0.008) (0.352,0.61)rectangle++(0.008,0.008) (0.146,0.854)rectangle++(0.008,0.008) (-0.242,0.862)rectangle++(0.008,0.008)
  (-0.104,0.96)rectangle++(0.008,0.008) (-0.007,0.401)rectangle++(0.008,0.008) (0.139,0.602)rectangle++(0.008,0.008) (0.278,0.642)rectangle++(0.008,0.008) (0.355,0.611)rectangle++(0.008,0.008)
  (0.147,0.854)rectangle++(0.008,0.008) (-0.241,0.862)rectangle++(0.008,0.008) (0.197,0.897)rectangle++(0.008,0.008) (0.409,0.756)rectangle++(0.008,0.008) (0.17,0.915)rectangle++(0.008,0.008)
  (0.069,0.982)rectangle++(0.008,0.008) (-0.323,0.893)rectangle++(0.008,0.008) (-0.47,0.712)rectangle++(0.008,0.008) (-0.2,0.897)rectangle++(0.008,0.008) (-0.086,0.974)rectangle++(0.008,0.008)
  (-0.391,0.834)rectangle++(0.008,0.008) (0.118,0.937)rectangle++(0.008,0.008) (-0.284,0.89)rectangle++(0.008,0.008) (0.187,0.925)rectangle++(0.008,0.008) (-0.248,0.909)rectangle++(0.008,0.008)
  (-0.107,0.979)rectangle++(0.008,0.008) (0.3,0.902)rectangle++(0.008,0.008) (0.459,0.721)rectangle++(0.008,0.008) (-0.052,0.913)rectangle++(0.008,0.008) (0.302,0.852)rectangle++(0.008,0.008)
  (0.005,0.356)rectangle++(0.008,0.008) (-0.128,0.584)rectangle++(0.008,0.008) (-0.056,0.843)rectangle++(0.008,0.008) (0.275,0.822)rectangle++(0.008,0.008) (0.113,0.943)rectangle++(0.008,0.008)
  (0.045,0.994)rectangle++(0.008,0.008) (-0.338,0.89)rectangle++(0.008,0.008) (0.163,0.944)rectangle++(0.008,0.008) (0.066,0.994)rectangle++(0.008,0.008) (0.025,1.015)rectangle++(0.008,0.008)
  (-0.354,0.892)rectangle++(0.008,0.008) (-0.484,0.7)rectangle++(0.008,0.008) (0.028,0.909)rectangle++(0.008,0.008) (0.009,0.98)rectangle++(0.008,0.008) (0.002,1.009)rectangle++(0.008,0.008)
  (-0.363,0.881)rectangle++(0.008,0.008) (0.148,0.949)rectangle++(0.008,0.008) (0.06,0.996)rectangle++(0.008,0.008) (0.382,0.85)rectangle++(0.008,0.008) (-0.133,0.944)rectangle++(0.008,0.008)
  (0.276,0.895)rectangle++(0.008,0.008) (0.445,0.727)rectangle++(0.008,0.008) (0.462,0.59)rectangle++(0.008,0.008) (-0.003,0.854)rectangle++(0.008,0.008) (0.303,0.808)rectangle++(0.008,0.008)
  (0.426,0.677)rectangle++(0.008,0.008) (0.436,0.574)rectangle++(0.008,0.008) (0.403,0.523)rectangle++(0.008,0.008) (-0.006,0.802)rectangle++(0.008,0.008) (0.283,0.785)rectangle++(0.008,0.008)
  (0.409,0.674)rectangle++(0.008,0.008) (0.169,0.881)rectangle++(0.008,0.008) (0.391,0.758)rectangle++(0.008,0.008) (0.162,0.916)rectangle++(0.008,0.008) (0.4,0.777)rectangle++(0.008,0.008)
  (0.46,0.629)rectangle++(0.008,0.008) (0.434,0.54)rectangle++(0.008,0.008) (0.002,0.821)rectangle++(0.008,0.008) (-0.004,0.342)rectangle++(0.008,0.008) (-0.004,0.741)rectangle++(0.008,0.008)
  (0.262,0.757)rectangle++(0.008,0.008) (0.389,0.669)rectangle++(0.008,0.008) (0.416,0.583)rectangle++(0.008,0.008) (0.398,0.534)rectangle++(0.008,0.008) (0.008,0.222)rectangle++(0.008,0.008)
  (-0.079,0.524)rectangle++(0.008,0.008) (-0.227,0.63)rectangle++(0.008,0.008) (-0.098,0.862)rectangle++(0.008,0.008) (-0.356,0.779)rectangle++(0.008,0.008) (-0.015,0.325)rectangle++(0.008,0.008)
  (0.109,0.57)rectangle++(0.008,0.008) (-0.001,0.237)rectangle++(0.008,0.008) (-0.089,0.527)rectangle++(0.008,0.008) (0.147,0.689)rectangle++(0.008,0.008) (-0.182,0.787)rectangle++(0.008,0.008)
  (-0.368,0.714)rectangle++(0.008,0.008) (-0.157,0.897)rectangle++(0.008,0.008) (0.248,0.883)rectangle++(0.008,0.008) (0.428,0.731)rectangle++(0.008,0.008) (-0.069,0.906)rectangle++(0.008,0.008)
  (-0.359,0.809)rectangle++(0.008,0.008) (0.124,0.914)rectangle++(0.008,0.008) (-0.273,0.881)rectangle++(0.008,0.008) (-0.117,0.968)rectangle++(0.008,0.008) (-0.403,0.82)rectangle++(0.008,0.008)
  (-0.48,0.65)rectangle++(0.008,0.008) (-0.455,0.545)rectangle++(0.008,0.008) (-0.406,0.506)rectangle++(0.008,0.008) (-0.006,0.793)rectangle++(0.008,0.008) (0.28,0.781)rectangle++(0.008,0.008)
  (0.005,0.326)rectangle++(0.008,0.008) (-0.118,0.57)rectangle++(0.008,0.008) (-0.261,0.638)rectangle++(0.008,0.008) (-0.112,0.865)rectangle++(0.008,0.008) (-0.049,0.961)rectangle++(0.008,0.008)
  (0.32,0.873)rectangle++(0.008,0.008) (-0.169,0.933)rectangle++(0.008,0.008) (0.255,0.903)rectangle++(0.008,0.008) (-0.21,0.923)rectangle++(0.008,0.008) (-0.429,0.766)rectangle++(0.008,0.008)
  (0.077,0.92)rectangle++(0.008,0.008) (0.363,0.809)rectangle++(0.008,0.008) (0.007,0.338)rectangle++(0.008,0.008) (0.123,0.568)rectangle++(0.008,0.008) (-0.15,0.723)rectangle++(0.008,0.008)
  (-0.065,0.901)rectangle++(0.008,0.008) (0.291,0.852)rectangle++(0.008,0.008) (0.12,0.955)rectangle++(0.008,0.008) (-0.29,0.899)rectangle++(0.008,0.008) (-0.457,0.727)rectangle++(0.008,0.008)
  (-0.018,0.303)rectangle++(0.008,0.008) (-0.01,0.725)rectangle++(0.008,0.008) (-0.006,0.902)rectangle++(0.008,0.008) (0.318,0.831)rectangle++(0.008,0.008) (0.131,0.947)rectangle++(0.008,0.008)
  (0.053,0.995)rectangle++(0.008,0.008) (-0.335,0.893)rectangle++(0.008,0.008) (-0.143,0.973)rectangle++(0.008,0.008) (0.281,0.912)rectangle++(0.008,0.008) (-0.201,0.937)rectangle++(0.008,0.008)
  (-0.01,0.391)rectangle++(0.008,0.008) (-0.006,0.762)rectangle++(0.008,0.008) (-0.279,0.765)rectangle++(0.008,0.008) (-0.119,0.919)rectangle++(0.008,0.008) (-0.386,0.797)rectangle++(0.008,0.008)
  (-0.165,0.932)rectangle++(0.008,0.008) (-0.071,0.989)rectangle++(0.008,0.008) (0.32,0.894)rectangle++(0.008,0.008) (-0.177,0.942)rectangle++(0.008,0.008) (-0.421,0.787)rectangle++(0.008,0.008)
  (-0.179,0.928)rectangle++(0.008,0.008) (-0.009,0.387)rectangle++(0.008,0.008) (-0.146,0.593)rectangle++(0.008,0.008) (-0.064,0.847)rectangle++(0.008,0.008) (0.272,0.826)rectangle++(0.008,0.008)
  (-0.174,0.894)rectangle++(0.008,0.008) (-0.402,0.766)rectangle++(0.008,0.008) (0.089,0.91)rectangle++(0.008,0.008) (0.035,0.98)rectangle++(0.008,0.008) (0.365,0.852)rectangle++(0.008,0.008)
  (-0.141,0.939)rectangle++(0.008,0.008) (0.27,0.896)rectangle++(0.008,0.008) (0.443,0.729)rectangle++(0.008,0.008) (0.462,0.592)rectangle++(0.008,0.008) (0.422,0.522)rectangle++(0.008,0.008)
  (0.175,0.817)rectangle++(0.008,0.008) (0.371,0.727)rectangle++(0.008,0.008) (0.428,0.616)rectangle++(0.008,0.008) (0.009,0.257)rectangle++(0.008,0.008) (-0.091,0.54)rectangle++(0.008,0.008)
  (-0.238,0.633)rectangle++(0.008,0.008) (0.117,0.791)rectangle++(0.008,0.008) (0.335,0.736)rectangle++(0.008,0.008) (-0.113,0.875)rectangle++(0.008,0.008) (-0.368,0.779)rectangle++(0.008,0.008)
  (-0.45,0.644)rectangle++(0.008,0.008) (-0.439,0.553)rectangle++(0.008,0.008) (-0.004,0.826)rectangle++(0.008,0.008) (-0.3,0.795)rectangle++(0.008,0.008) (0.146,0.887)rectangle++(0.008,0.008)
  (-0.254,0.877)rectangle++(0.008,0.008) (-0.109,0.966)rectangle++(0.008,0.008) (-0.048,1.003)rectangle++(0.008,0.008) (0.336,0.892)rectangle++(0.008,0.008) (0.006,0.372)rectangle++(0.008,0.008)
  (-0.134,0.592)rectangle++(0.008,0.008) (0.149,0.735)rectangle++(0.008,0.008) (0.33,0.699)rectangle++(0.008,0.008) (0.006,0.291)rectangle++(0.008,0.008) (0.106,0.547)rectangle++(0.008,0.008)
  (-0.151,0.707)rectangle++(0.008,0.008) (-0.066,0.895)rectangle++(0.008,0.008) (-0.03,0.973)rectangle++(0.008,0.008) (0.333,0.872)rectangle++(0.008,0.008) (0.138,0.964)rectangle++(0.008,0.008)
  (-0.285,0.909)rectangle++(0.008,0.008) (-0.458,0.733)rectangle++(0.008,0.008) (-0.195,0.906)rectangle++(0.008,0.008) (-0.416,0.764)rectangle++(0.008,0.008) (-0.177,0.918)rectangle++(0.008,0.008)
  (0.246,0.899)rectangle++(0.008,0.008) (0.433,0.74)rectangle++(0.008,0.008) (0.461,0.6)rectangle++(0.008,0.008) (0.01,0.25)rectangle++(0.008,0.008) (-0.088,0.537)rectangle++(0.008,0.008)
  (0.151,0.693)rectangle++(0.008,0.008) (0.316,0.679)rectangle++(0.008,0.008) (-0.102,0.843)rectangle++(0.008,0.008) (-0.351,0.768)rectangle++(0.008,0.008) (0.113,0.893)rectangle++(0.008,0.008)
  (0.37,0.784)rectangle++(0.008,0.008) (0.448,0.643)rectangle++(0.008,0.008) (0.434,0.55)rectangle++(0.008,0.008) (-0.002,0.826)rectangle++(0.008,0.008) (-0.004,0.344)rectangle++(0.008,0.008)
  (0.12,0.575)rectangle++(0.008,0.008) (-0.154,0.725)rectangle++(0.008,0.008) (-0.067,0.902)rectangle++(0.008,0.008) (-0.03,0.977)rectangle++(0.008,0.008) (-0.366,0.855)rectangle++(0.008,0.008)
  (0.137,0.938)rectangle++(0.008,0.008) (0.397,0.796)rectangle++(0.008,0.008) (-0.107,0.925)rectangle++(0.008,0.008) (-0.047,0.986)rectangle++(0.008,0.008) (-0.022,1.012)rectangle++(0.008,0.008)
  (0.35,0.887)rectangle++(0.008,0.008) (0.145,0.97)rectangle++(0.008,0.008) (-0.284,0.914)rectangle++(0.008,0.008) (0.196,0.936)rectangle++(0.008,0.008) (0.423,0.774)rectangle++(0.008,0.008)
  (-0.087,0.924)rectangle++(0.008,0.008) (-0.006,0.386)rectangle++(0.008,0.008) (-0.005,0.76)rectangle++(0.008,0.008) (-0.277,0.765)rectangle++(0.008,0.008) (0.146,0.865)rectangle++(0.008,0.008)
  (0.375,0.76)rectangle++(0.008,0.008) (0.155,0.917)rectangle++(0.008,0.008) (0.001,0.383)rectangle++(0.008,0.008) (-0.004,0.158)rectangle++(0.008,0.008) (-0.062,0.49)rectangle++(0.008,0.008)
  (0.146,0.663)rectangle++(0.008,0.008) (-0.173,0.774)rectangle++(0.008,0.008) (-0.075,0.923)rectangle++(0.008,0.008) (-0.034,0.985)rectangle++(0.008,0.008) (-0.005,0.411)rectangle++(0.008,0.008)
  (-0.004,0.771)rectangle++(0.008,0.008) (-0.281,0.77)rectangle++(0.008,0.008) (0.146,0.869)rectangle++(0.008,0.008) (0.376,0.761)rectangle++(0.008,0.008) (0.443,0.63)rectangle++(0.008,0.008)
  (-0.026,0.865)rectangle++(0.008,0.008) (-0.005,0.361)rectangle++(0.008,0.008) (-0.135,0.583)rectangle++(0.008,0.008) (-0.273,0.637)rectangle++(0.008,0.008) (0.102,0.805)rectangle++(0.008,0.008)
  (0.333,0.748)rectangle++(0.008,0.008) (0.138,0.912)rectangle++(0.008,0.008) (0.388,0.784)rectangle++(0.008,0.008) (-0.106,0.916)rectangle++(0.008,0.008) (-0.047,0.982)rectangle++(0.008,0.008)
  (-0.376,0.851)rectangle++(0.008,0.008) (-0.479,0.674)rectangle++(0.008,0.008) (0.021,0.896)rectangle++(0.008,0.008) (0.329,0.818)rectangle++(0.008,0.008) (-0.146,0.911)rectangle++(0.008,0.008)
  (-0.063,0.98)rectangle++(0.008,0.008) (0.32,0.887)rectangle++(0.008,0.008) (-0.174,0.939)rectangle++(0.008,0.008) (-0.418,0.786)rectangle++(0.008,0.008) (0.089,0.925)rectangle++(0.008,0.008)
  (0.035,0.986)rectangle++(0.008,0.008) (0.367,0.855)rectangle++(0.008,0.008) (0.472,0.676)rectangle++(0.008,0.008) (0.456,0.557)rectangle++(0.008,0.008) (0.407,0.508)rectangle++(0.008,0.008)
  (0.366,0.503)rectangle++(0.008,0.008) (0.346,0.516)rectangle++(0.008,0.008) (0.143,0.814)rectangle++(0.008,0.008) (0.355,0.737)rectangle++(0.008,0.008) (0.425,0.627)rectangle++(0.008,0.008)
  (0.176,0.861)rectangle++(0.008,0.008) (0.387,0.747)rectangle++(0.008,0.008) (0.443,0.62)rectangle++(0.008,0.008) (0.423,0.541)rectangle++(0.008,0.008) (0.009,0.225)rectangle++(0.008,0.008)
  (0.001,0.692)rectangle++(0.008,0.008) (0.247,0.732)rectangle++(0.008,0.008) (0.374,0.663)rectangle++(0.008,0.008) (0.007,0.276)rectangle++(0.008,0.008) (-0.004,0.114)rectangle++(0.008,0.008)
  (-0.046,0.47)rectangle++(0.008,0.008) (-0.022,0.795)rectangle++(0.008,0.008) (-0.011,0.932)rectangle++(0.008,0.008) (0.327,0.846)rectangle++(0.008,0.008) (0.451,0.687)rectangle++(0.008,0.008)
  (-0.043,0.894)rectangle++(0.008,0.008) (-0.005,0.373)rectangle++(0.008,0.008) (0.13,0.589)rectangle++(0.008,0.008) (-0.154,0.735)rectangle++(0.008,0.008) (-0.067,0.906)rectangle++(0.008,0.008)
  (0.292,0.855)rectangle++(0.008,0.008) (0.005,0.357)rectangle++(0.008,0.008) (-0.129,0.584)rectangle++(0.008,0.008) (-0.056,0.843)rectangle++(0.008,0.008) (0.274,0.822)rectangle++(0.008,0.008)
  (-0.172,0.893)rectangle++(0.008,0.008) (-0.401,0.766)rectangle++(0.008,0.008) (0.09,0.91)rectangle++(0.008,0.008) (-0.287,0.867)rectangle++(0.008,0.008) (0.177,0.915)rectangle++(0.008,0.008)
  (0.407,0.771)rectangle++(0.008,0.008) (0.008,0.322)rectangle++(0.008,0.008) (-0.115,0.569)rectangle++(0.008,0.008) (0.15,0.718)rectangle++(0.008,0.008) (0.001,0.299)rectangle++(0.008,0.008)
  (-0.002,0.723)rectangle++(0.008,0.008) (-0.003,0.901)rectangle++(0.008,0.008) (-0.004,0.976)rectangle++(0.008,0.008) (-0.354,0.864)rectangle++(0.008,0.008) (-0.474,0.688)rectangle++(0.008,0.008)
  (0.028,0.9)rectangle++(0.008,0.008) (-0.312,0.841)rectangle++(0.008,0.008) (0.157,0.912)rectangle++(0.008,0.008) (0.397,0.777)rectangle++(0.008,0.008) (0.458,0.63)rectangle++(0.008,0.008)
  (-0.019,0.871)rectangle++(0.008,0.008) (-0.323,0.81)rectangle++(0.008,0.008) (0.141,0.902)rectangle++(0.008,0.008) (0.386,0.778)rectangle++(0.008,0.008) (0.454,0.634)rectangle++(0.008,0.008)
  (-0.023,0.871)rectangle++(0.008,0.008) (-0.325,0.809)rectangle++(0.008,0.008) (-0.441,0.673)rectangle++(0.008,0.008) (0.038,0.882)rectangle++(0.008,0.008) (-0.301,0.836)rectangle++(0.008,0.008)
  (-0.44,0.694)rectangle++(0.008,0.008) (0.046,0.89)rectangle++(0.008,0.008) (0.338,0.807)rectangle++(0.008,0.008) (0.442,0.664)rectangle++(0.008,0.008) (0.009,0.277)rectangle++(0.008,0.008)
  (0.102,0.54)rectangle++(0.008,0.008) (0.239,0.626)rectangle++(0.008,0.008) (0.098,0.861)rectangle++(0.008,0.008) (0.351,0.775)rectangle++(0.008,0.008) (-0.12,0.899)rectangle++(0.008,0.008)
  (-0.053,0.975)rectangle++(0.008,0.008) (-0.376,0.846)rectangle++(0.008,0.008) (-0.478,0.672)rectangle++(0.008,0.008) (-0.203,0.88)rectangle++(0.008,0.008) (0.221,0.891)rectangle++(0.008,0.008)
  (0.003,0.372)rectangle++(0.008,0.008) (-0.001,0.754)rectangle++(0.008,0.008) (-0.273,0.763)rectangle++(0.008,0.008) (-0.117,0.918)rectangle++(0.008,0.008) (0.274,0.878)rectangle++(0.008,0.008)
  (-0.192,0.918)rectangle++(0.008,0.008) (-0.419,0.77)rectangle++(0.008,0.008) (0.083,0.918)rectangle++(0.008,0.008) (0.365,0.807)rectangle++(0.008,0.008) (-0.125,0.918)rectangle++(0.008,0.008)
  (-0.055,0.983)rectangle++(0.008,0.008) (-0.025,1.011)rectangle++(0.008,0.008) (-0.376,0.872)rectangle++(0.008,0.008) (0.139,0.949)rectangle++(0.008,0.008) (-0.279,0.903)rectangle++(0.008,0.008)
  (0.194,0.929)rectangle++(0.008,0.008) (-0.247,0.913)rectangle++(0.008,0.008) (-0.442,0.748)rectangle++(0.008,0.008) (0.064,0.916)rectangle++(0.008,0.008) (0.025,0.983)rectangle++(0.008,0.008)
  (0.008,1.01)rectangle++(0.008,0.008) (-0.361,0.884)rectangle++(0.008,0.008) (0.15,0.949)rectangle++(0.008,0.008) (-0.274,0.907)rectangle++(0.008,0.008) (-0.453,0.736)rectangle++(0.008,0.008)
  (-0.017,0.307)rectangle++(0.008,0.008) (0.101,0.563)rectangle++(0.008,0.008) (0.04,0.834)rectangle++(0.008,0.008) (0.315,0.783)rectangle++(0.008,0.008) (-0.139,0.89)rectangle++(0.008,0.008)
  (0.253,0.873)rectangle++(0.008,0.008) (0.427,0.725)rectangle++(0.008,0.008) (-0.067,0.903)rectangle++(0.008,0.008) (-0.357,0.808)rectangle++(0.008,0.008) (-0.152,0.937)rectangle++(0.008,0.008)
  (-0.008,0.391)rectangle++(0.008,0.008) (-0.147,0.595)rectangle++(0.008,0.008) (0.144,0.741)rectangle++(0.008,0.008) (-0.202,0.809)rectangle++(0.008,0.008) (-0.087,0.938)rectangle++(0.008,0.008)
  (-0.378,0.817)rectangle++(0.008,0.008) (-0.161,0.941)rectangle++(0.008,0.008) (0.261,0.904)rectangle++(0.008,0.008) (-0.207,0.926)rectangle++(0.008,0.008) (-0.089,0.986)rectangle++(0.008,0.008)
  (0.311,0.899)rectangle++(0.008,0.008) (0.462,0.716)rectangle++(0.008,0.008) (0.466,0.579)rectangle++(0.008,0.008) (0.419,0.514)rectangle++(0.008,0.008) (0.374,0.502)rectangle++(0.008,0.008)
  (-0.012,0.782)rectangle++(0.008,0.008) (0.273,0.778)rectangle++(0.008,0.008) (0.402,0.674)rectangle++(0.008,0.008) (0.167,0.881)rectangle++(0.008,0.008) (0.39,0.76)rectangle++(0.008,0.008)
  (0.449,0.624)rectangle++(0.008,0.008) (0.427,0.542)rectangle++(0.008,0.008) (-0.002,0.819)rectangle++(0.008,0.008) (-0.004,0.342)rectangle++(0.008,0.008) (0.12,0.574)rectangle++(0.008,0.008)
  (0.048,0.839)rectangle++(0.008,0.008) (-0.281,0.82)rectangle++(0.008,0.008) (-0.425,0.693)rectangle++(0.008,0.008) (-0.445,0.585)rectangle++(0.008,0.008) (-0.416,0.527)rectangle++(0.008,0.008)
  (-0.003,0.806)rectangle++(0.008,0.008) (-0.003,0.936)rectangle++(0.008,0.008) (0.332,0.845)rectangle++(0.008,0.008) (0.137,0.953)rectangle++(0.008,0.008) (0.055,0.998)rectangle++(0.008,0.008)
  (0.381,0.853)rectangle++(0.008,0.008) (0.008,0.356)rectangle++(0.008,0.008) (0.001,0.747)rectangle++(0.008,0.008) (-0.002,0.911)rectangle++(0.008,0.008) (0.324,0.834)rectangle++(0.008,0.008)
  (-0.153,0.916)rectangle++(0.008,0.008) (-0.067,0.982)rectangle++(0.008,0.008) (-0.385,0.844)rectangle++(0.008,0.008) (-0.164,0.952)rectangle++(0.008,0.008) (-0.419,0.796)rectangle++(0.008,0.008)
  (-0.479,0.634)rectangle++(0.008,0.008) (0.007,0.877)rectangle++(0.008,0.008) (0.315,0.815)rectangle++(0.008,0.008) (-0.15,0.904)rectangle++(0.008,0.008) (-0.066,0.978)rectangle++(0.008,0.008)
  (0.318,0.887)rectangle++(0.008,0.008) (0.131,0.97)rectangle++(0.008,0.008) (0.406,0.813)rectangle++(0.008,0.008) (0.168,0.939)rectangle++(0.008,0.008) (0.411,0.786)rectangle++(0.008,0.008)
  (-0.096,0.925)rectangle++(0.008,0.008) (0.286,0.874)rectangle++(0.008,0.008) (-0.185,0.921)rectangle++(0.008,0.008) (-0.417,0.774)rectangle++(0.008,0.008) (-0.177,0.923)rectangle++(0.008,0.008)
  (0.248,0.901)rectangle++(0.008,0.008) (-0.212,0.92)rectangle++(0.008,0.008) (0.231,0.913)rectangle++(0.008,0.008) (-0.224,0.919)rectangle++(0.008,0.008) (-0.011,0.384)rectangle++(0.008,0.008)
  (-0.145,0.591)rectangle++(0.008,0.008) (-0.063,0.846)rectangle++(0.008,0.008) (-0.006,0.353)rectangle++(0.008,0.008) (-0.005,0.746)rectangle++(0.008,0.008) (-0.272,0.758)rectangle++(0.008,0.008)
  (-0.399,0.669)rectangle++(0.008,0.008) (-0.425,0.582)rectangle++(0.008,0.008) (0.013,0.834)rectangle++(0.008,0.008) (-0.295,0.805)rectangle++(0.008,0.008) (-0.126,0.936)rectangle++(0.008,0.008)
  (-0.055,0.991)rectangle++(0.008,0.008) (-0.383,0.852)rectangle++(0.008,0.008) (-0.483,0.672)rectangle++(0.008,0.008) (0.019,0.896)rectangle++(0.008,0.008) (-0.003,0.374)rectangle++(0.008,0.008)
  (0.131,0.589)rectangle++(0.008,0.008) (-0.154,0.735)rectangle++(0.008,0.008) (-0.336,0.7)rectangle++(0.008,0.008) (-0.144,0.892)rectangle++(0.008,0.008) (-0.063,0.972)rectangle++(0.008,0.008)
  (-0.379,0.841)rectangle++(0.008,0.008) (0.126,0.936)rectangle++(0.008,0.008) (0.051,0.991)rectangle++(0.008,0.008) (-0.002,0.414)rectangle++(0.008,0.008) (-0.003,0.772)rectangle++(0.008,0.008)
  (-0.004,0.322)rectangle++(0.008,0.008) (-0.004,0.733)rectangle++(0.008,0.008) (-0.004,0.905)rectangle++(0.008,0.008) (-0.329,0.832)rectangle++(0.008,0.008) (-0.14,0.947)rectangle++(0.008,0.008)
  (-0.406,0.802)rectangle++(0.008,0.008) (0.1,0.928)rectangle++(0.008,0.008) (-0.289,0.879)rectangle++(0.008,0.008) (-0.124,0.967)rectangle++(0.008,0.008) (-0.405,0.817)rectangle++(0.008,0.008)
  (0.106,0.935)rectangle++(0.008,0.008) (-0.289,0.884)rectangle++(0.008,0.008) (-0.124,0.969)rectangle++(0.008,0.008) (-0.054,1.005)rectangle++(0.008,0.008) (-0.006,0.42)rectangle++(0.008,0.008)
  (-0.156,0.609)rectangle++(0.008,0.008) (0.145,0.75)rectangle++(0.008,0.008) (0,0.313)rectangle++(0.008,0.008) (-0.002,0.729)rectangle++(0.008,0.008) (-0.003,0.904)rectangle++(0.008,0.008)
  (-0.004,0.977)rectangle++(0.008,0.008) (-0.004,0.408)rectangle++(0.008,0.008) (0.143,0.604)rectangle++(0.008,0.008) (0.058,0.852)rectangle++(0.008,0.008) (0.022,0.955)rectangle++(0.008,0.008)
  (0.35,0.845)rectangle++(0.008,0.008) (-0.145,0.931)rectangle++(0.008,0.008) (-0.063,0.989)rectangle++(0.008,0.008) (-0.386,0.848)rectangle++(0.008,0.008) (-0.164,0.954)rectangle++(0.008,0.008)
  (-0.071,0.998)rectangle++(0.008,0.008) (0.323,0.898)rectangle++(0.008,0.008) (0.468,0.712)rectangle++(0.008,0.008) (0.194,0.897)rectangle++(0.008,0.008) (-0.235,0.898)rectangle++(0.008,0.008)
  (0.213,0.911)rectangle++(0.008,0.008) (-0.232,0.911)rectangle++(0.008,0.008) (-0.011,0.381)rectangle++(0.008,0.008) (0.13,0.594)rectangle++(0.008,0.008) (0.271,0.641)rectangle++(0.008,0.008)
  (-0.109,0.809)rectangle++(0.008,0.008) (-0.342,0.75)rectangle++(0.008,0.008) (-0.146,0.913)rectangle++(0.008,0.008) (0.258,0.886)rectangle++(0.008,0.008) (-0.202,0.916)rectangle++(0.008,0.008)
  (-0.423,0.766)rectangle++(0.008,0.008) (-0.47,0.618)rectangle++(0.008,0.008) (0.005,0.867)rectangle++(0.008,0.008) (-0.004,0.362)rectangle++(0.008,0.008) (0.127,0.583)rectangle++(0.008,0.008)
  (0.265,0.638)rectangle++(0.008,0.008) (-0.11,0.805)rectangle++(0.008,0.008) (-0.341,0.748)rectangle++(0.008,0.008) (-0.427,0.639)rectangle++(0.008,0.008) (0.032,0.861)rectangle++(0.008,0.008)
  (0.011,0.959)rectangle++(0.008,0.008) (-0.341,0.862)rectangle++(0.008,0.008) (-0.467,0.691)rectangle++(0.008,0.008) (-0.464,0.568)rectangle++(0.008,0.008) (-0.419,0.513)rectangle++(0.008,0.008)
  (-0.009,0.801)rectangle++(0.008,0.008) (-0.006,0.934)rectangle++(0.008,0.008) (-0.34,0.844)rectangle++(0.008,0.008) (0.145,0.923)rectangle++(0.008,0.008) (0.395,0.787)rectangle++(0.008,0.008)
  (0.461,0.635)rectangle++(0.008,0.008) (0.437,0.542)rectangle++(0.008,0.008) (0.392,0.508)rectangle++(0.008,0.008) (-0.006,0.792)rectangle++(0.008,0.008) (-0.005,0.93)rectangle++(0.008,0.008)
  (0.329,0.843)rectangle++(0.008,0.008) (0.451,0.684)rectangle++(0.008,0.008) (-0.042,0.893)rectangle++(0.008,0.008) (0.299,0.839)rectangle++(0.008,0.008) (-0.167,0.91)rectangle++(0.008,0.008)
  (-0.405,0.775)rectangle++(0.008,0.008) (-0.465,0.629)rectangle++(0.008,0.008) (-0.198,0.862)rectangle++(0.008,0.008) (0.217,0.881)rectangle++(0.008,0.008) (0.003,0.368)rectangle++(0.008,0.008)
  (-0.134,0.588)rectangle++(0.008,0.008) (0.148,0.733)rectangle++(0.008,0.008) (0.06,0.906)rectangle++(0.008,0.008) (-0.002,0.378)rectangle++(0.008,0.008) (-0.14,0.591)rectangle++(0.008,0.008)
  (-0.279,0.64)rectangle++(0.008,0.008) (0.1,0.808)rectangle++(0.008,0.008) (0.334,0.75)rectangle++(0.008,0.008) (0.138,0.913)rectangle++(0.008,0.008) (0,0.381)rectangle++(0.008,0.008)
  (0.135,0.59)rectangle++(0.008,0.008) (0.272,0.638)rectangle++(0.008,0.008) (0.112,0.866)rectangle++(0.008,0.008) (0.359,0.772)rectangle++(0.008,0.008) (0.439,0.641)rectangle++(0.008,0.008)
  (0.429,0.552)rectangle++(0.008,0.008) (-0.005,0.825)rectangle++(0.008,0.008) (-0.004,0.944)rectangle++(0.008,0.008) (0.334,0.849)rectangle++(0.008,0.008) (0.138,0.954)rectangle++(0.008,0.008)
  (0.403,0.803)rectangle++(0.008,0.008) (-0.106,0.93)rectangle++(0.008,0.008) (0.283,0.88)rectangle++(0.008,0.008) (-0.188,0.922)rectangle++(0.008,0.008) (-0.419,0.774)rectangle++(0.008,0.008)
  (0.084,0.92)rectangle++(0.008,0.008) (0.366,0.807)rectangle++(0.008,0.008) (0.007,0.336)rectangle++(0.008,0.008) (0.123,0.568)rectangle++(0.008,0.008) (0,0.236)rectangle++(0.008,0.008)
  (-0.002,0.697)rectangle++(0.008,0.008) (-0.254,0.737)rectangle++(0.008,0.008) (-0.383,0.666)rectangle++(0.008,0.008) (0.062,0.857)rectangle++(0.008,0.008) (0.024,0.958)rectangle++(0.008,0.008)
  (-0.335,0.865)rectangle++(0.008,0.008) (0.155,0.931)rectangle++(0.008,0.008) (0.403,0.787)rectangle++(0.008,0.008) (-0.1,0.923)rectangle++(0.008,0.008) (0.283,0.874)rectangle++(0.008,0.008)
  (0.116,0.965)rectangle++(0.008,0.008) (-0.295,0.902)rectangle++(0.008,0.008) (0.187,0.934)rectangle++(0.008,0.008) (0.418,0.777)rectangle++(0.008,0.008) (0.468,0.622)rectangle++(0.008,0.008)
  (-0.012,0.871)rectangle++(0.008,0.008) (0.305,0.819)rectangle++(0.008,0.008) (-0.157,0.902)rectangle++(0.008,0.008) (0.25,0.885)rectangle++(0.008,0.008) (0.429,0.732)rectangle++(0.008,0.008)
  (0.009,0.305)rectangle++(0.008,0.008) (0.001,0.726)rectangle++(0.008,0.008) (-0.262,0.751)rectangle++(0.008,0.008) (0.148,0.853)rectangle++(0.008,0.008) (0.372,0.754)rectangle++(0.008,0.008)
  (0.007,0.314)rectangle++(0.008,0.008) (-0.112,0.565)rectangle++(0.008,0.008) (-0.05,0.835)rectangle++(0.008,0.008) (-0.324,0.783)rectangle++(0.008,0.008) (-0.432,0.661)rectangle++(0.008,0.008)
  (-0.017,0.275)rectangle++(0.008,0.008) (-0.009,0.713)rectangle++(0.008,0.008) (-0.006,0.897)rectangle++(0.008,0.008) (-0.004,0.375)rectangle++(0.008,0.008) (-0.139,0.589)rectangle++(0.008,0.008)
  (-0.061,0.845)rectangle++(0.008,0.008) (-0.028,0.953)rectangle++(0.008,0.008) (-0.014,0.998)rectangle++(0.008,0.008) (-0.366,0.87)rectangle++(0.008,0.008) (-0.482,0.686)rectangle++(0.008,0.008)
  (-0.205,0.886)rectangle++(0.008,0.008) (0.222,0.894)rectangle++(0.008,0.008) (0.091,0.973)rectangle++(0.008,0.008) (0.388,0.829)rectangle++(0.008,0.008) (0.008,0.346)rectangle++(0.008,0.008)
  (0.001,0.743)rectangle++(0.008,0.008) (0.265,0.756)rectangle++(0.008,0.008) (-0.152,0.859)rectangle++(0.008,0.008) (0.236,0.863)rectangle++(0.008,0.008) (0.416,0.727)rectangle++(0.008,0.008)
  (0.449,0.6)rectangle++(0.008,0.008) (0.01,0.25)rectangle++(0.008,0.008) (0.093,0.527)rectangle++(0.008,0.008) (-0.149,0.693)rectangle++(0.008,0.008) (-0.319,0.683)rectangle++(0.008,0.008)
  (-0.393,0.617)rectangle++(0.008,0.008) (0.04,0.839)rectangle++(0.008,0.008) (0.317,0.786)rectangle++(0.008,0.008) (-0.139,0.892)rectangle++(0.008,0.008) (0.254,0.874)rectangle++(0.008,0.008)
  (0.104,0.965)rectangle++(0.008,0.008) (-0.3,0.897)rectangle++(0.008,0.008) (-0.128,0.975)rectangle++(0.008,0.008) (-0.41,0.819)rectangle++(0.008,0.008) (-0.484,0.647)rectangle++(0.008,0.008)
  (-0.456,0.542)rectangle++(0.008,0.008) (-0.405,0.504)rectangle++(0.008,0.008) (-0.006,0.792)rectangle++(0.008,0.008) (0.279,0.78)rectangle++(0.008,0.008) (-0.155,0.876)rectangle++(0.008,0.008)
  (-0.067,0.965)rectangle++(0.008,0.008) (0.313,0.882)rectangle++(0.008,0.008) (0.457,0.708)rectangle++(0.008,0.008) (-0.047,0.906)rectangle++(0.008,0.008) (0.301,0.847)rectangle++(0.008,0.008)
  (0.439,0.696)rectangle++(0.008,0.008) (-0.051,0.894)rectangle++(0.008,0.008) (0.295,0.843)rectangle++(0.008,0.008) (0.435,0.697)rectangle++(0.008,0.008) (0.181,0.89)rectangle++(0.008,0.008)
  (-0.239,0.891)rectangle++(0.008,0.008) (0.208,0.909)rectangle++(0.008,0.008) (0.085,0.979)rectangle++(0.008,0.008) (0.388,0.834)rectangle++(0.008,0.008) (0.161,0.948)rectangle++(0.008,0.008)
  (0.001,0.396)rectangle++(0.008,0.008) (0.141,0.597)rectangle++(0.008,0.008) (-0.152,0.742)rectangle++(0.008,0.008) (0.195,0.81)rectangle++(0.008,0.008) (-0.204,0.859)rectangle++(0.008,0.008)
  (0.213,0.882)rectangle++(0.008,0.008) (0.411,0.743)rectangle++(0.008,0.008) (-0.081,0.906)rectangle++(0.008,0.008) (0.286,0.859)rectangle++(0.008,0.008) (0.437,0.707)rectangle++(0.008,0.008)
  (-0.057,0.898)rectangle++(0.008,0.008) (0.294,0.847)rectangle++(0.008,0.008) (-0.172,0.912)rectangle++(0.008,0.008) (-0.408,0.775)rectangle++(0.008,0.008) (-0.466,0.628)rectangle++(0.008,0.008)
  (-0.441,0.54)rectangle++(0.008,0.008) (-0.398,0.508)rectangle++(0.008,0.008) (-0.367,0.51)rectangle++(0.008,0.008) (0.013,0.78)rectangle++(0.008,0.008) (0.284,0.768)rectangle++(0.008,0.008)
  (-0.148,0.872)rectangle++(0.008,0.008) (-0.065,0.964)rectangle++(0.008,0.008) (-0.029,1.002)rectangle++(0.008,0.008) (-0.005,0.419)rectangle++(0.008,0.008) (-0.004,0.774)rectangle++(0.008,0.008)
  (-0.004,0.923)rectangle++(0.008,0.008) (-0.335,0.839)rectangle++(0.008,0.008) (-0.456,0.683)rectangle++(0.008,0.008) (-0.194,0.885)rectangle++(0.008,0.008) (-0.084,0.969)rectangle++(0.008,0.008)
  (-0.038,1.005)rectangle++(0.008,0.008) (-0.018,1.02)rectangle++(0.008,0.008) (-0.01,1.026)rectangle++(0.008,0.008) (-0.374,0.885)rectangle++(0.008,0.008) (0.144,0.954)rectangle++(0.008,0.008)
  (0.058,0.999)rectangle++(0.008,0.008) (0.022,1.017)rectangle++(0.008,0.008) (0.007,1.025)rectangle++(0.008,0.008) (-0.366,0.89)rectangle++(0.008,0.008) (-0.489,0.695)rectangle++(0.008,0.008)
  (-0.475,0.562)rectangle++(0.008,0.008) (-0.018,0.234)rectangle++(0.008,0.008) (0.074,0.53)rectangle++(0.008,0.008) (-0.159,0.688)rectangle++(0.008,0.008) (-0.322,0.677)rectangle++(0.008,0.008)
  (-0.137,0.882)rectangle++(0.008,0.008) (0.251,0.869)rectangle++(0.008,0.008) (0.103,0.963)rectangle++(0.008,0.008) (0.041,1.002)rectangle++(0.008,0.008) (0.376,0.86)rectangle++(0.008,0.008)
  (0.478,0.675)rectangle++(0.008,0.008) (-0.026,0.898)rectangle++(0.008,0.008) (-0.005,0.375)rectangle++(0.008,0.008) (-0.004,0.755)rectangle++(0.008,0.008) (-0.004,0.915)rectangle++(0.008,0.008)
  (-0.004,0.982)rectangle++(0.008,0.008) (0.348,0.867)rectangle++(0.008,0.008) (0.144,0.962)rectangle++(0.008,0.008) (-0.281,0.91)rectangle++(0.008,0.008) (0.196,0.933)rectangle++(0.008,0.008)
  (-0.247,0.915)rectangle++(0.008,0.008) (0.213,0.923)rectangle++(0.008,0.008) (0.003,0.385)rectangle++(0.008,0.008) (-0.001,0.76)rectangle++(0.008,0.008) (0.27,0.764)rectangle++(0.008,0.008)
  (-0.153,0.865)rectangle++(0.008,0.008) (-0.067,0.961)rectangle++(0.008,0.008) (-0.377,0.835)rectangle++(0.008,0.008) (-0.161,0.948)rectangle++(0.008,0.008) (-0.416,0.795)rectangle++(0.008,0.008)
  (0.093,0.928)rectangle++(0.008,0.008) (0.373,0.807)rectangle++(0.008,0.008) (0.154,0.937)rectangle++(0.008,0.008) (0.063,0.991)rectangle++(0.008,0.008) (0.024,1.014)rectangle++(0.008,0.008)
  (0.372,0.871)rectangle++(0.008,0.008) (0.154,0.964)rectangle++(0.008,0.008) (0.414,0.802)rectangle++(0.008,0.008) (-0.101,0.933)rectangle++(0.008,0.008) (-0.383,0.81)rectangle++(0.008,0.008)
  (-0.468,0.653)rectangle++(0.008,0.008) (-0.45,0.55)rectangle++(0.008,0.008) (-0.01,0.829)rectangle++(0.008,0.008) (-0.007,0.946)rectangle++(0.008,0.008) (-0.005,0.995)rectangle++(0.008,0.008)
  (-0.361,0.872)rectangle++(0.008,0.008) (0.146,0.944)rectangle++(0.008,0.008) (0.403,0.796)rectangle++(0.008,0.008) (0.167,0.932)rectangle++(0.008,0.008) (-0.26,0.905)rectangle++(0.008,0.008)
  (-0.446,0.74)rectangle++(0.008,0.008) (-0.189,0.908)rectangle++(0.008,0.008) (0.237,0.899)rectangle++(0.008,0.008) (0.003,0.375)rectangle++(0.008,0.008) (0.135,0.587)rectangle++(0.008,0.008)
  (-0.152,0.736)rectangle++(0.008,0.008) (-0.066,0.907)rectangle++(0.008,0.008) (0.293,0.854)rectangle++(0.008,0.008) (0.438,0.702)rectangle++(0.008,0.008) (0.182,0.893)rectangle++(0.008,0.008)
  (0.074,0.973)rectangle++(0.008,0.008) (0.38,0.834)rectangle++(0.008,0.008) (0.157,0.948)rectangle++(0.008,0.008) (0.41,0.794)rectangle++(0.008,0.008) (-0.1,0.928)rectangle++(0.008,0.008)
  (0.285,0.876)rectangle++(0.008,0.008) (0.442,0.715)rectangle++(0.008,0.008) (0.184,0.898)rectangle++(0.008,0.008) (-0.24,0.895)rectangle++(0.008,0.008) (-0.433,0.743)rectangle++(0.008,0.008)
  (0.066,0.91)rectangle++(0.008,0.008) (-0.298,0.859)rectangle++(0.008,0.008) (0.17,0.916)rectangle++(0.008,0.008) (0.404,0.774)rectangle++(0.008,0.008) (0.46,0.626)rectangle++(0.008,0.008)
  (0.01,0.261)rectangle++(0.008,0.008) (0.097,0.532)rectangle++(0.008,0.008) (0.233,0.625)rectangle++(0.008,0.008) (0.329,0.619)rectangle++(0.008,0.008) (0.136,0.858)rectangle++(0.008,0.008)
  (-0.248,0.86)rectangle++(0.008,0.008) (-0.106,0.959)rectangle++(0.008,0.008) (-0.395,0.819)rectangle++(0.008,0.008) (-0.016,0.342)rectangle++(0.008,0.008) (-0.009,0.741)rectangle++(0.008,0.008)
  (0.26,0.758)rectangle++(0.008,0.008) (0.389,0.67)rectangle++(0.008,0.008) (0.416,0.584)rectangle++(0.008,0.008) (0.173,0.843)rectangle++(0.008,0.008) (0.07,0.952)rectangle++(0.008,0.008)
  (-0.311,0.879)rectangle++(0.008,0.008) (0.171,0.929)rectangle++(0.008,0.008) (-0.257,0.905)rectangle++(0.008,0.008) (-0.11,0.978)rectangle++(0.008,0.008) (-0.007,0.408)rectangle++(0.008,0.008)
  (0.142,0.606)rectangle++(0.008,0.008) (-0.155,0.747)rectangle++(0.008,0.008) (0.195,0.813)rectangle++(0.008,0.008) (0.002,0.339)rectangle++(0.008,0.008) (-0.124,0.575)rectangle++(0.008,0.008)
  (0.148,0.723)rectangle++(0.008,0.008) (-0.194,0.803)rectangle++(0.008,0.008) (-0.084,0.935)rectangle++(0.008,0.008) (0.295,0.874)rectangle++(0.008,0.008) (-0.181,0.924)rectangle++(0.008,0.008)
  (-0.009,0.386)rectangle++(0.008,0.008) (0.133,0.596)rectangle++(0.008,0.008) (-0.156,0.739)rectangle++(0.008,0.008) (-0.068,0.908)rectangle++(0.008,0.008) (-0.359,0.81)rectangle++(0.008,0.008)
  (0.125,0.915)rectangle++(0.008,0.008) (0.05,0.982)rectangle++(0.008,0.008) (-0.331,0.886)rectangle++(0.008,0.008) (0.164,0.94)rectangle++(0.008,0.008) (0.067,0.992)rectangle++(0.008,0.008)
  (-0.327,0.897)rectangle++(0.008,0.008) (-0.14,0.974)rectangle++(0.008,0.008) (-0.061,1.007)rectangle++(0.008,0.008) (0.331,0.898)rectangle++(0.008,0.008) (0.006,0.375)rectangle++(0.008,0.008)
  (-0.135,0.593)rectangle++(0.008,0.008) (0.149,0.735)rectangle++(0.008,0.008) (0.06,0.907)rectangle++(0.008,0.008) (-0.3,0.855)rectangle++(0.008,0.008) (-0.446,0.703)rectangle++(0.008,0.008)
  (-0.458,0.581)rectangle++(0.008,0.008) (-0.421,0.521)rectangle++(0.008,0.008) (-0.382,0.507)rectangle++(0.008,0.008) (0.006,0.785)rectangle++(0.008,0.008) (0,0.927)rectangle++(0.008,0.008)
  (-0.335,0.843)rectangle++(0.008,0.008) (-0.014,0.352)rectangle++(0.008,0.008) (-0.136,0.575)rectangle++(0.008,0.008) (0.143,0.728)rectangle++(0.008,0.008) (0.324,0.698)rectangle++(0.008,0.008)
  (0.134,0.891)rectangle++(0.008,0.008) (0.379,0.776)rectangle++(0.008,0.008) (-0.108,0.909)rectangle++(0.008,0.008) (-0.377,0.796)rectangle++(0.008,0.008) (-0.46,0.648)rectangle++(0.008,0.008)
  (0.02,0.877)rectangle++(0.008,0.008) (-0.003,0.366)rectangle++(0.008,0.008) (0.128,0.585)rectangle++(0.008,0.008) (0,0.243)rectangle++(0.008,0.008) (0.086,0.528)rectangle++(0.008,0.008)
  (-0.153,0.691)rectangle++(0.008,0.008) (0.176,0.787)rectangle++(0.008,0.008) (-0.204,0.842)rectangle++(0.008,0.008) (0.207,0.874)rectangle++(0.008,0.008) (0.084,0.965)rectangle++(0.008,0.008)
  (0.033,1.003)rectangle++(0.008,0.008) (0.012,1.019)rectangle++(0.008,0.008) (-0.362,0.889)rectangle++(0.008,0.008) (0.151,0.952)rectangle++(0.008,0.008) (0.408,0.798)rectangle++(0.008,0.008)
  (0.169,0.933)rectangle++(0.008,0.008) (-0.259,0.906)rectangle++(0.008,0.008) (-0.446,0.741)rectangle++(0.008,0.008) (-0.472,0.599)rectangle++(0.008,0.008) (-0.433,0.524)rectangle++(0.008,0.008)
  (-0.389,0.504)rectangle++(0.008,0.008) (-0.166,0.809)rectangle++(0.008,0.008) (0.213,0.846)rectangle++(0.008,0.008) (0.398,0.727)rectangle++(0.008,0.008) (0.165,0.903)rectangle++(0.008,0.008)
  (-0.251,0.891)rectangle++(0.008,0.008) (0.203,0.913)rectangle++(0.008,0.008) (0.083,0.981)rectangle++(0.008,0.008) (0.387,0.835)rectangle++(0.008,0.008) (-0.125,0.939)rectangle++(0.008,0.008)
  (-0.396,0.804)rectangle++(0.008,0.008) (0.105,0.925)rectangle++(0.008,0.008) (0.042,0.986)rectangle++(0.008,0.008) (0.371,0.852)rectangle++(0.008,0.008) (0.473,0.674)rectangle++(0.008,0.008)
  (0.196,0.881)rectangle++(0.008,0.008) (-0.228,0.892)rectangle++(0.008,0.008) (-0.426,0.745)rectangle++(0.008,0.008) (-0.181,0.911)rectangle++(0.008,0.008) (0.242,0.897)rectangle++(0.008,0.008)
  (0.003,0.375)rectangle++(0.008,0.008) (-0.001,0.755)rectangle++(0.008,0.008) (0.268,0.762)rectangle++(0.008,0.008) (-0.153,0.863)rectangle++(0.008,0.008) (0.238,0.866)rectangle++(0.008,0.008)
  (0.417,0.727)rectangle++(0.008,0.008) (0.449,0.6)rectangle++(0.008,0.008) (0.186,0.85)rectangle++(0.008,0.008) (0.002,0.355)rectangle++(0.008,0.008) (0.127,0.578)rectangle++(0.008,0.008)
  (-0.152,0.729)rectangle++(0.008,0.008) (-0.333,0.698)rectangle++(0.008,0.008) (-0.405,0.619)rectangle++(0.008,0.008) (-0.016,0.258)rectangle++(0.008,0.008) (-0.103,0.531)rectangle++(0.008,0.008)
  (-0.24,0.625)rectangle++(0.008,0.008) (-0.011,0.26)rectangle++(0.008,0.008) (-0.004,0.107)rectangle++(0.008,0.008) (-0.004,0.043)rectangle++(0.008,0.008) (0.013,0.437)rectangle++(0.008,0.008)
  (-0.154,0.624)rectangle++(0.008,0.008) (-0.297,0.649)rectangle++(0.008,0.008) (0.096,0.819)rectangle++(0.008,0.008) (0.335,0.757)rectangle++(0.008,0.008) (0.423,0.642)rectangle++(0.008,0.008)
  (0.422,0.559)rectangle++(0.008,0.008) (-0.01,0.825)rectangle++(0.008,0.008) (0.289,0.797)rectangle++(0.008,0.008) (0.119,0.933)rectangle++(0.008,0.008) (0.387,0.8)rectangle++(0.008,0.008)
  (0.008,0.334)rectangle++(0.008,0.008) (-0.119,0.575)rectangle++(0.008,0.008) (-0.007,0.239)rectangle++(0.008,0.008) (-0.092,0.526)rectangle++(0.008,0.008) (-0.041,0.819)rectangle++(0.008,0.008)
  (-0.02,0.941)rectangle++(0.008,0.008) (-0.011,0.993)rectangle++(0.008,0.008) (-0.007,1.015)rectangle++(0.008,0.008) (-0.369,0.881)rectangle++(0.008,0.008) (-0.157,0.968)rectangle++(0.008,0.008)
  (-0.421,0.805)rectangle++(0.008,0.008) (-0.179,0.936)rectangle++(0.008,0.008) (-0.009,0.391)rectangle++(0.008,0.008) (-0.006,0.762)rectangle++(0.008,0.008) (-0.278,0.765)rectangle++(0.008,0.008)
  (-0.119,0.919)rectangle++(0.008,0.008) (-0.052,0.984)rectangle++(0.008,0.008) (-0.024,1.011)rectangle++(0.008,0.008) (0.349,0.887)rectangle++(0.008,0.008) (0.144,0.97)rectangle++(0.008,0.008)
  (0.412,0.808)rectangle++(0.008,0.008) (-0.104,0.936)rectangle++(0.008,0.008) (-0.007,0.391)rectangle++(0.008,0.008) (-0.004,0.162)rectangle++(0.008,0.008) (-0.063,0.492)rectangle++(0.008,0.008)
  (-0.029,0.804)rectangle++(0.008,0.008) (0.273,0.794)rectangle++(0.008,0.008) (0.004,0.331)rectangle++(0.008,0.008) (0.12,0.566)rectangle++(0.008,0.008) (-0.151,0.721)rectangle++(0.008,0.008)
  (0.188,0.8)rectangle++(0.008,0.008) (0.371,0.715)rectangle++(0.008,0.008) (-0.09,0.878)rectangle++(0.008,0.008) (0.272,0.85)rectangle++(0.008,0.008) (0.112,0.955)rectangle++(0.008,0.008)
  (-0.293,0.896)rectangle++(0.008,0.008) (-0.126,0.974)rectangle++(0.008,0.008) (-0.409,0.82)rectangle++(0.008,0.008) (-0.483,0.648)rectangle++(0.008,0.008) (0.01,0.885)rectangle++(0.008,0.008)
  (0.32,0.817)rectangle++(0.008,0.008) (0.006,0.341)rectangle++(0.008,0.008) (0.124,0.57)rectangle++(0.008,0.008) (0.05,0.837)rectangle++(0.008,0.008) (0.019,0.949)rectangle++(0.008,0.008)
  (-0.003,0.396)rectangle++(0.008,0.008) (-0.004,0.764)rectangle++(0.008,0.008) (0.27,0.767)rectangle++(0.008,0.008) (-0.154,0.866)rectangle++(0.008,0.008) (0.238,0.867)rectangle++(0.008,0.008)
  (0.098,0.962)rectangle++(0.008,0.008) (-0.302,0.894)rectangle++(0.008,0.008) (-0.129,0.973)rectangle++(0.008,0.008) (-0.008,0.406)rectangle++(0.008,0.008) (-0.152,0.602)rectangle++(0.008,0.008)
  (-0.066,0.851)rectangle++(0.008,0.008) (0.273,0.829)rectangle++(0.008,0.008) (-0.175,0.895)rectangle++(0.008,0.008) (-0.403,0.766)rectangle++(0.008,0.008) (-0.461,0.625)rectangle++(0.008,0.008)
  (-0.196,0.86)rectangle++(0.008,0.008) (0.217,0.88)rectangle++(0.008,0.008) (-0.219,0.899)rectangle++(0.008,0.008) (-0.01,0.375)rectangle++(0.008,0.008) (-0.142,0.587)rectangle++(0.008,0.008)
  (-0.062,0.844)rectangle++(0.008,0.008) (-0.333,0.783)rectangle++(0.008,0.008) (-0.142,0.927)rectangle++(0.008,0.008) (-0.4,0.792)rectangle++(0.008,0.008) (0.099,0.921)rectangle++(0.008,0.008)
  (-0.001,0.384)rectangle++(0.008,0.008) (-0.004,0.159)rectangle++(0.008,0.008) (-0.062,0.491)rectangle++(0.008,0.008) (-0.028,0.804)rectangle++(0.008,0.008) (-0.014,0.935)rectangle++(0.008,0.008)
  (-0.004,0.39)rectangle++(0.008,0.008) (-0.145,0.596)rectangle++(0.008,0.008) (0.146,0.741)rectangle++(0.008,0.008) (0.059,0.909)rectangle++(0.008,0.008) (-0.301,0.856)rectangle++(0.008,0.008)
  (-0.013,0.357)rectangle++(0.008,0.008) (-0.008,0.748)rectangle++(0.008,0.008) (-0.006,0.912)rectangle++(0.008,0.008) (0.322,0.835)rectangle++(0.008,0.008) (-0.155,0.916)rectangle++(0.008,0.008)
  (-0.401,0.783)rectangle++(0.008,0.008) (-0.171,0.926)rectangle++(0.008,0.008) (-0.413,0.782)rectangle++(0.008,0.008) (-0.176,0.926)rectangle++(0.008,0.008) (-0.414,0.78)rectangle++(0.008,0.008)
  (0.088,0.921)rectangle++(0.008,0.008) (-0.001,0.384)rectangle++(0.008,0.008) (0.136,0.593)rectangle++(0.008,0.008) (0.273,0.639)rectangle++(0.008,0.008) (0.352,0.611)rectangle++(0.008,0.008)
  (0.378,0.57)rectangle++(0.008,0.008) (0.376,0.542)rectangle++(0.008,0.008) (0.007,0.225)rectangle++(0.008,0.008) (-0.081,0.525)rectangle++(0.008,0.008) (-0.228,0.63)rectangle++(0.008,0.008)
  (-0.098,0.862)rectangle++(0.008,0.008) (0.262,0.846)rectangle++(0.008,0.008) (0.004,0.353)rectangle++(0.008,0.008) (-0.128,0.582)rectangle++(0.008,0.008) (-0.27,0.64)rectangle++(0.008,0.008)
  (-0.355,0.615)rectangle++(0.008,0.008) (-0.386,0.574)rectangle++(0.008,0.008) (0.028,0.816)rectangle++(0.008,0.008) (-0.282,0.802)rectangle++(0.008,0.008) (-0.121,0.935)rectangle++(0.008,0.008)
  (-0.393,0.803)rectangle++(0.008,0.008) (0.107,0.924)rectangle++(0.008,0.008) (0.378,0.8)rectangle++(0.008,0.008) (0.007,0.334)rectangle++(0.008,0.008) (-0.119,0.575)rectangle++(0.008,0.008)
  (0.15,0.722)rectangle++(0.008,0.008) (0.325,0.693)rectangle++(0.008,0.008) (0.134,0.889)rectangle++(0.008,0.008) (0.054,0.971)rectangle++(0.008,0.008) (0.371,0.841)rectangle++(0.008,0.008)
  (0.153,0.951)rectangle++(0.008,0.008) (0.062,0.997)rectangle++(0.008,0.008) (0.384,0.85)rectangle++(0.008,0.008) (-0.132,0.945)rectangle++(0.008,0.008) (-0.058,0.994)rectangle++(0.008,0.008)
  (0.328,0.891)rectangle++(0.008,0.008) (0.468,0.707)rectangle++(0.008,0.008) (-0.042,0.909)rectangle++(0.008,0.008) (-0.02,0.98)rectangle++(0.008,0.008) (-0.011,1.009)rectangle++(0.008,0.008)
  (-0.007,1.021)rectangle++(0.008,0.008) (0.361,0.886)rectangle++(0.008,0.008) (-0.155,0.953)rectangle++(0.008,0.008) (-0.415,0.799)rectangle++(0.008,0.008) (-0.479,0.637)rectangle++(0.008,0.008)
  (-0.203,0.865)rectangle++(0.008,0.008) (-0.405,0.742)rectangle++(0.008,0.008) (-0.173,0.909)rectangle++(0.008,0.008) (-0.075,0.98)rectangle++(0.008,0.008) (0.315,0.891)rectangle++(0.008,0.008)
  (0.461,0.711)rectangle++(0.008,0.008) (0.464,0.577)rectangle++(0.008,0.008) (0.003,0.849)rectangle++(0.008,0.008) (-0.305,0.808)rectangle++(0.008,0.008) (-0.432,0.679)rectangle++(0.008,0.008)
  (-0.444,0.576)rectangle++(0.008,0.008) (-0.189,0.839)rectangle++(0.008,0.008) (-0.39,0.736)rectangle++(0.008,0.008) (0.084,0.892)rectangle++(0.008,0.008) (-0.284,0.857)rectangle++(0.008,0.008)
  (-0.121,0.958)rectangle++(0.008,0.008) (-0.053,1)rectangle++(0.008,0.008) (-0.385,0.857)rectangle++(0.008,0.008) (-0.164,0.958)rectangle++(0.008,0.008) (0.266,0.913)rectangle++(0.008,0.008)
  (0.447,0.738)rectangle++(0.008,0.008) (0.01,0.308)rectangle++(0.008,0.008) (0.002,0.727)rectangle++(0.008,0.008) (-0.262,0.752)rectangle++(0.008,0.008) (-0.012,0.314)rectangle++(0.008,0.008)
  (-0.121,0.558)rectangle++(0.008,0.008) (0.143,0.715)rectangle++(0.008,0.008) (0.058,0.898)rectangle++(0.008,0.008) (-0.298,0.85)rectangle++(0.008,0.008) (0.167,0.911)rectangle++(0.008,0.008)
  (-0.253,0.895)rectangle++(0.008,0.008) (0.203,0.916)rectangle++(0.008,0.008) (0.419,0.762)rectangle++(0.008,0.008) (0.463,0.615)rectangle++(0.008,0.008) (0.01,0.256)rectangle++(0.008,0.008)
  (-0.09,0.54)rectangle++(0.008,0.008) (-0.238,0.634)rectangle++(0.008,0.008) (0.117,0.791)rectangle++(0.008,0.008) (0,0.33)rectangle++(0.008,0.008) (-0.122,0.57)rectangle++(0.008,0.008)
  (-0.008,0.237)rectangle++(0.008,0.008) (-0.092,0.525)rectangle++(0.008,0.008) (-0.233,0.626)rectangle++(0.008,0.008) (0.116,0.786)rectangle++(0.008,0.008) (0.333,0.734)rectangle++(0.008,0.008)
  (0.414,0.633)rectangle++(0.008,0.008) (-0.041,0.856)rectangle++(0.008,0.008) (-0.019,0.957)rectangle++(0.008,0.008) (0.332,0.861)rectangle++(0.008,0.008) (0.458,0.691)rectangle++(0.008,0.008)
  (-0.041,0.899)rectangle++(0.008,0.008) (-0.343,0.815)rectangle++(0.008,0.008) (0.134,0.912)rectangle++(0.008,0.008) (0.054,0.981)rectangle++(0.008,0.008) (-0.329,0.887)rectangle++(0.008,0.008)
  (-0.141,0.97)rectangle++(0.008,0.008) (-0.008,0.405)rectangle++(0.008,0.008) (-0.006,0.768)rectangle++(0.008,0.008) (0.271,0.769)rectangle++(0.008,0.008) (0.111,0.921)rectangle++(0.008,0.008)
  (-0.281,0.88)rectangle++(0.008,0.008) (0.185,0.919)rectangle++(0.008,0.008) (0.075,0.984)rectangle++(0.008,0.008) (0.385,0.839)rectangle++(0.008,0.008) (0.159,0.95)rectangle++(0.008,0.008)
  (0.411,0.794)rectangle++(0.008,0.008) (-0.099,0.929)rectangle++(0.008,0.008) (-0.007,0.388)rectangle++(0.008,0.008) (0.135,0.596)rectangle++(0.008,0.008) (-0.155,0.74)rectangle++(0.008,0.008)
  (-0.339,0.702)rectangle++(0.008,0.008) (-0.409,0.619)rectangle++(0.008,0.008) (-0.174,0.858)rectangle++(0.008,0.008) (-0.389,0.749)rectangle++(0.008,0.008) (-0.449,0.623)rectangle++(0.008,0.008)
  (-0.431,0.543)rectangle++(0.008,0.008) (-0.183,0.826)rectangle++(0.008,0.008) (0.21,0.859)rectangle++(0.008,0.008) (-0.214,0.887)rectangle++(0.008,0.008) (0.218,0.898)rectangle++(0.008,0.008)
  (0.003,0.375)rectangle++(0.008,0.008) (-0.136,0.592)rectangle++(0.008,0.008) (-0.008,0.246)rectangle++(0.008,0.008) (-0.095,0.529)rectangle++(0.008,0.008) (-0.042,0.82)rectangle++(0.008,0.008)
  (-0.316,0.779)rectangle++(0.008,0.008) (0.133,0.885)rectangle++(0.008,0.008) (0.376,0.773)rectangle++(0.008,0.008) (0.156,0.923)rectangle++(0.008,0.008) (0.063,0.985)rectangle++(0.008,0.008)
  (-0.327,0.892)rectangle++(0.008,0.008) (-0.139,0.972)rectangle++(0.008,0.008) (0.283,0.911)rectangle++(0.008,0.008) (-0.2,0.936)rectangle++(0.008,0.008) (-0.429,0.776)rectangle++(0.008,0.008)
  (0.08,0.924)rectangle++(0.008,0.008) (-0.297,0.87)rectangle++(0.008,0.008) (-0.127,0.963)rectangle++(0.008,0.008) (0.285,0.902)rectangle++(0.008,0.008) (-0.195,0.933)rectangle++(0.008,0.008)
  (0.243,0.913)rectangle++(0.008,0.008) (0.436,0.747)rectangle++(0.008,0.008) (0.181,0.911)rectangle++(0.008,0.008) (0.407,0.768)rectangle++(0.008,0.008) (0.008,0.32)rectangle++(0.008,0.008)
  (-0.114,0.569)rectangle++(0.008,0.008) (0.15,0.717)rectangle++(0.008,0.008) (-0.191,0.801)rectangle++(0.008,0.008) (0.198,0.851)rectangle++(0.008,0.008) (-0.217,0.879)rectangle++(0.008,0.008)
  (-0.01,0.367)rectangle++(0.008,0.008) (0.125,0.588)rectangle++(0.008,0.008) (0.05,0.845)rectangle++(0.008,0.008) (0.019,0.952)rectangle++(0.008,0.008) (-0.335,0.861)rectangle++(0.008,0.008)
  (-0.143,0.959)rectangle++(0.008,0.008) (-0.412,0.807)rectangle++(0.008,0.008) (0.099,0.932)rectangle++(0.008,0.008) (0.039,0.989)rectangle++(0.008,0.008) (0.37,0.854)rectangle++(0.008,0.008)
  (0.474,0.675)rectangle++(0.008,0.008) (-0.028,0.897)rectangle++(0.008,0.008) (-0.337,0.819)rectangle++(0.008,0.008) (0.138,0.911)rectangle++(0.008,0.008) (0.056,0.98)rectangle++(0.008,0.008)
  (0.375,0.845)rectangle++(0.008,0.008) (-0.134,0.939)rectangle++(0.008,0.008) (0.273,0.893)rectangle++(0.008,0.008) (0.443,0.727)rectangle++(0.008,0.008) (0.461,0.59)rectangle++(0.008,0.008)
  (0.421,0.521)rectangle++(0.008,0.008) (0.378,0.504)rectangle++(0.008,0.008) (-0.011,0.785)rectangle++(0.008,0.008) (-0.007,0.927)rectangle++(0.008,0.008) (-0.338,0.841)rectangle++(0.008,0.008)
  (-0.014,0.351)rectangle++(0.008,0.008) (0.118,0.582)rectangle++(0.008,0.008) (0,0.242)rectangle++(0.008,0.008) (-0.002,0.699)rectangle++(0.008,0.008) (0.248,0.737)rectangle++(0.008,0.008)
  (0.102,0.907)rectangle++(0.008,0.008) (0.37,0.795)rectangle++(0.008,0.008) (-0.118,0.914)rectangle++(0.008,0.008) (-0.384,0.795)rectangle++(0.008,0.008) (-0.463,0.645)rectangle++(0.008,0.008)
  (-0.197,0.869)rectangle++(0.008,0.008) (-0.404,0.746)rectangle++(0.008,0.008) (-0.455,0.616)rectangle++(0.008,0.008) (0.011,0.86)rectangle++(0.008,0.008) (0.003,0.959)rectangle++(0.008,0.008)
  (-0.004,0.4)rectangle++(0.008,0.008) (0.141,0.601)rectangle++(0.008,0.008) (0.057,0.85)rectangle++(0.008,0.008) (0.329,0.785)rectangle++(0.008,0.008) (-0.134,0.895)rectangle++(0.008,0.008)
  (-0.008,0.374)rectangle++(0.008,0.008) (-0.006,0.755)rectangle++(0.008,0.008) (-0.005,0.915)rectangle++(0.008,0.008) (-0.332,0.836)rectangle++(0.008,0.008) (-0.142,0.949)rectangle++(0.008,0.008)
  (-0.407,0.802)rectangle++(0.008,0.008) (-0.476,0.64)rectangle++(0.008,0.008) (-0.018,0.267)rectangle++(0.008,0.008) (-0.107,0.535)rectangle++(0.008,0.008) (0.141,0.699)rectangle++(0.008,0.008)
  (0.313,0.685)rectangle++(0.008,0.008) (0.387,0.618)rectangle++(0.008,0.008) (-0.047,0.84)rectangle++(0.008,0.008) (-0.325,0.786)rectangle++(0.008,0.008) (0.131,0.892)rectangle++(0.008,0.008)
  (0.378,0.777)rectangle++(0.008,0.008) (0.449,0.637)rectangle++(0.008,0.008) (0.432,0.547)rectangle++(0.008,0.008) (-0.001,0.823)rectangle++(0.008,0.008) (0.293,0.793)rectangle++(0.008,0.008)
  (0.121,0.931)rectangle++(0.008,0.008) (0.048,0.989)rectangle++(0.008,0.008) (0.374,0.851)rectangle++(0.008,0.008) (-0.136,0.942)rectangle++(0.008,0.008) (0.273,0.896)rectangle++(0.008,0.008)
  (0.444,0.728)rectangle++(0.008,0.008) (0.184,0.904)rectangle++(0.008,0.008) (0.002,0.377)rectangle++(0.008,0.008) (-0.138,0.592)rectangle++(0.008,0.008) (-0.278,0.641)rectangle++(0.008,0.008)
  (-0.119,0.867)rectangle++(0.008,0.008) (-0.007,0.362)rectangle++(0.008,0.008) (0.125,0.584)rectangle++(0.008,0.008) (0.265,0.639)rectangle++(0.008,0.008) (0.109,0.866)rectangle++(0.008,0.008)
  (0.043,0.961)rectangle++(0.008,0.008) (0.362,0.84)rectangle++(0.008,0.008) (0.15,0.951)rectangle++(0.008,0.008) (-0.275,0.907)rectangle++(0.008,0.008) (-0.453,0.736)rectangle++(0.008,0.008)
  (0.055,0.915)rectangle++(0.008,0.008) (0.351,0.815)rectangle++(0.008,0.008) (0.451,0.663)rectangle++(0.008,0.008) (-0.035,0.883)rectangle++(0.008,0.008) (-0.335,0.811)rectangle++(0.008,0.008)
  (0.136,0.906)rectangle++(0.008,0.008) (-0.265,0.882)rectangle++(0.008,0.008) (-0.012,0.368)rectangle++(0.008,0.008) (-0.004,0.152)rectangle++(0.008,0.008) (0.052,0.488)rectangle++(0.008,0.008);
\fill[nuage]
  (0.019,0.802)rectangle++(0.008,0.008) (-0.281,0.793)rectangle++(0.008,0.008) (-0.012,0.331)rectangle++(0.008,0.008) (0.112,0.572)rectangle++(0.008,0.008) (-0.001,0.238)rectangle++(0.008,0.008)
  (-0.089,0.528)rectangle++(0.008,0.008) (0.147,0.689)rectangle++(0.008,0.008) (0.313,0.679)rectangle++(0.008,0.008) (-0.103,0.841)rectangle++(0.008,0.008) (0.252,0.838)rectangle++(0.008,0.008)
  (0.414,0.709)rectangle++(0.008,0.008) (-0.068,0.891)rectangle++(0.008,0.008) (0.287,0.848)rectangle++(0.008,0.008) (0.433,0.702)rectangle++(0.008,0.008) (0.18,0.892)rectangle++(0.008,0.008)
  (-0.24,0.891)rectangle++(0.008,0.008) (-0.432,0.741)rectangle++(0.008,0.008) (-0.184,0.909)rectangle++(0.008,0.008) (0.24,0.897)rectangle++(0.008,0.008) (0.098,0.975)rectangle++(0.008,0.008)
  (0.039,1.007)rectangle++(0.008,0.008) (-0.345,0.893)rectangle++(0.008,0.008) (-0.48,0.704)rectangle++(0.008,0.008) (-0.204,0.893)rectangle++(0.008,0.008) (-0.416,0.755)rectangle++(0.008,0.008)
  (-0.177,0.915)rectangle++(0.008,0.008) (-0.411,0.774)rectangle++(0.008,0.008) (-0.016,0.323)rectangle++(0.008,0.008) (0.107,0.57)rectangle++(0.008,0.008) (0.043,0.837)rectangle++(0.008,0.008)
  (-0.283,0.817)rectangle++(0.008,0.008) (0.162,0.891)rectangle++(0.008,0.008) (0.001,0.372)rectangle++(0.008,0.008) (-0.136,0.59)rectangle++(0.008,0.008) (-0.276,0.64)rectangle++(0.008,0.008)
  (-0.358,0.613)rectangle++(0.008,0.008) (0.054,0.825)rectangle++(0.008,0.008) (-0.273,0.816)rectangle++(0.008,0.008) (-0.42,0.694)rectangle++(0.008,0.008) (0.055,0.884)rectangle++(0.008,0.008)
  (-0.294,0.843)rectangle++(0.008,0.008) (-0.126,0.952)rectangle++(0.008,0.008) (-0.401,0.809)rectangle++(0.008,0.008) (0.105,0.929)rectangle++(0.008,0.008) (-0.287,0.882)rectangle++(0.008,0.008)
  (-0.123,0.968)rectangle++(0.008,0.008) (-0.405,0.818)rectangle++(0.008,0.008) (0.106,0.935)rectangle++(0.008,0.008) (0.382,0.806)rectangle++(0.008,0.008) (-0.117,0.924)rectangle++(0.008,0.008)
  (-0.387,0.8)rectangle++(0.008,0.008) (0.108,0.92)rectangle++(0.008,0.008) (-0.283,0.878)rectangle++(0.008,0.008) (0.184,0.919)rectangle++(0.008,0.008) (0.075,0.984)rectangle++(0.008,0.008)
  (-0.321,0.896)rectangle++(0.008,0.008) (-0.137,0.974)rectangle++(0.008,0.008) (0.284,0.91)rectangle++(0.008,0.008) (-0.199,0.937)rectangle++(0.008,0.008) (0.243,0.916)rectangle++(0.008,0.008)
  (0.003,0.382)rectangle++(0.008,0.008) (-0.004,0.158)rectangle++(0.008,0.008) (0.054,0.49)rectangle++(0.008,0.008) (-0.002,0.204)rectangle++(0.008,0.008) (-0.003,0.683)rectangle++(0.008,0.008)
  (-0.004,0.885)rectangle++(0.008,0.008) (-0.321,0.822)rectangle++(0.008,0.008) (0.146,0.907)rectangle++(0.008,0.008) (-0.261,0.886)rectangle++(0.008,0.008) (0.196,0.914)rectangle++(0.008,0.008)
  (0.08,0.982)rectangle++(0.008,0.008) (-0.318,0.897)rectangle++(0.008,0.008) (0.174,0.94)rectangle++(0.008,0.008) (-0.259,0.911)rectangle++(0.008,0.008) (-0.447,0.743)rectangle++(0.008,0.008)
  (-0.473,0.599)rectangle++(0.008,0.008) (-0.434,0.524)rectangle++(0.008,0.008) (-0.185,0.818)rectangle++(0.008,0.008) (-0.08,0.941)rectangle++(0.008,0.008) (-0.036,0.993)rectangle++(0.008,0.008)
  (0.337,0.883)rectangle++(0.008,0.008) (0.469,0.699)rectangle++(0.008,0.008) (0.463,0.569)rectangle++(0.008,0.008) (0.01,0.237)rectangle++(0.008,0.008) (-0.083,0.531)rectangle++(0.008,0.008)
  (-0.231,0.632)rectangle++(0.008,0.008) (-0.099,0.863)rectangle++(0.008,0.008) (-0.044,0.96)rectangle++(0.008,0.008) (0.322,0.871)rectangle++(0.008,0.008) (0.133,0.964)rectangle++(0.008,0.008)
  (0.404,0.809)rectangle++(0.008,0.008) (0.167,0.938)rectangle++(0.008,0.008) (0.068,0.991)rectangle++(0.008,0.008) (-0.327,0.897)rectangle++(0.008,0.008) (-0.473,0.712)rectangle++(0.008,0.008)
  (0.037,0.911)rectangle++(0.008,0.008) (-0.312,0.849)rectangle++(0.008,0.008) (-0.449,0.696)rectangle++(0.008,0.008) (0.042,0.895)rectangle++(0.008,0.008) (-0.003,0.374)rectangle++(0.008,0.008)
  (-0.138,0.589)rectangle++(0.008,0.008) (-0.06,0.845)rectangle++(0.008,0.008) (-0.333,0.784)rectangle++(0.008,0.008) (-0.436,0.659)rectangle++(0.008,0.008) (0.035,0.873)rectangle++(0.008,0.008)
  (-0.299,0.831)rectangle++(0.008,0.008) (-0.128,0.947)rectangle++(0.008,0.008) (-0.4,0.806)rectangle++(0.008,0.008) (0.104,0.928)rectangle++(0.008,0.008) (-0.001,0.387)rectangle++(0.008,0.008)
  (-0.142,0.596)rectangle++(0.008,0.008) (0.147,0.74)rectangle++(0.008,0.008) (0.331,0.702)rectangle++(0.008,0.008) (0.137,0.892)rectangle++(0.008,0.008) (0.38,0.776)rectangle++(0.008,0.008)
  (-0.107,0.91)rectangle++(0.008,0.008) (-0.007,0.38)rectangle++(0.008,0.008) (-0.005,0.757)rectangle++(0.008,0.008) (-0.276,0.763)rectangle++(0.008,0.008) (-0.118,0.918)rectangle++(0.008,0.008)
  (-0.386,0.797)rectangle++(0.008,0.008) (-0.464,0.646)rectangle++(0.008,0.008) (0.018,0.877)rectangle++(0.008,0.008) (0.005,0.966)rectangle++(0.008,0.008) (-0.346,0.863)rectangle++(0.008,0.008)
  (-0.014,0.36)rectangle++(0.008,0.008) (-0.139,0.579)rectangle++(0.008,0.008) (0.143,0.73)rectangle++(0.008,0.008) (0.325,0.699)rectangle++(0.008,0.008) (0.006,0.291)rectangle++(0.008,0.008)
  (0.106,0.547)rectangle++(0.008,0.008) (0.042,0.828)rectangle++(0.008,0.008) (-0.28,0.813)rectangle++(0.008,0.008) (-0.012,0.339)rectangle++(0.008,0.008) (-0.007,0.74)rectangle++(0.008,0.008)
  (-0.271,0.755)rectangle++(0.008,0.008) (-0.012,0.315)rectangle++(0.008,0.008) (0.106,0.565)rectangle++(0.008,0.008) (-0.157,0.715)rectangle++(0.008,0.008) (-0.068,0.898)rectangle++(0.008,0.008)
  (-0.006,0.375)rectangle++(0.008,0.008) (0.13,0.59)rectangle++(0.008,0.008) (0.052,0.845)rectangle++(0.008,0.008) (0.325,0.784)rectangle++(0.008,0.008) (0.428,0.659)rectangle++(0.008,0.008)
  (0.43,0.565)rectangle++(0.008,0.008) (0.397,0.521)rectangle++(0.008,0.008) (0.008,0.216)rectangle++(0.008,0.008) (0.08,0.512)rectangle++(0.008,0.008) (0.031,0.813)rectangle++(0.008,0.008)
  (-0.28,0.802)rectangle++(0.008,0.008) (0.158,0.883)rectangle++(0.008,0.008) (0.001,0.368)rectangle++(0.008,0.008) (0.131,0.584)rectangle++(0.008,0.008) (-0.152,0.733)rectangle++(0.008,0.008)
  (0.191,0.806)rectangle++(0.008,0.008) (0.078,0.936)rectangle++(0.008,0.008) (-0.302,0.875)rectangle++(0.008,0.008) (0.174,0.924)rectangle++(0.008,0.008) (0.071,0.986)rectangle++(0.008,0.008)
  (0.384,0.842)rectangle++(0.008,0.008) (0.159,0.951)rectangle++(0.008,0.008) (0.001,0.397)rectangle++(0.008,0.008) (-0.002,0.765)rectangle++(0.008,0.008) (-0.277,0.768)rectangle++(0.008,0.008)
  (-0.119,0.92)rectangle++(0.008,0.008) (-0.052,0.984)rectangle++(0.008,0.008) (-0.024,1.011)rectangle++(0.008,0.008) (-0.013,1.022)rectangle++(0.008,0.008) (-0.374,0.882)rectangle++(0.008,0.008)
  (0.143,0.953)rectangle++(0.008,0.008) (0.058,0.998)rectangle++(0.008,0.008) (0.382,0.852)rectangle++(0.008,0.008) (0.478,0.669)rectangle++(0.008,0.008) (-0.024,0.896)rectangle++(0.008,0.008)
  (-0.335,0.82)rectangle++(0.008,0.008) (-0.449,0.675)rectangle++(0.008,0.008) (0.035,0.885)rectangle++(0.008,0.008) (0.012,0.969)rectangle++(0.008,0.008) (-0.344,0.867)rectangle++(0.008,0.008)
  (-0.47,0.692)rectangle++(0.008,0.008) (-0.466,0.568)rectangle++(0.008,0.008) (-0.018,0.236)rectangle++(0.008,0.008) (-0.096,0.521)rectangle++(0.008,0.008) (-0.233,0.623)rectangle++(0.008,0.008)
  (0.115,0.785)rectangle++(0.008,0.008) (0.046,0.927)rectangle++(0.008,0.008) (0.017,0.987)rectangle++(0.008,0.008) (-0.348,0.876)rectangle++(0.008,0.008) (-0.149,0.966)rectangle++(0.008,0.008)
  (-0.416,0.808)rectangle++(0.008,0.008) (0.097,0.934)rectangle++(0.008,0.008) (0.039,0.99)rectangle++(0.008,0.008) (-0.003,0.413)rectangle++(0.008,0.008) (-0.003,0.771)rectangle++(0.008,0.008)
  (0.273,0.77)rectangle++(0.008,0.008) (0.112,0.921)rectangle++(0.008,0.008) (-0.281,0.88)rectangle++(0.008,0.008) (-0.446,0.721)rectangle++(0.008,0.008) (-0.465,0.589)rectangle++(0.008,0.008)
  (-0.003,0.852)rectangle++(0.008,0.008) (-0.309,0.808)rectangle++(0.008,0.008) (0.146,0.896)rectangle++(0.008,0.008) (0.386,0.774)rectangle++(0.008,0.008) (-0.103,0.911)rectangle++(0.008,0.008)
  (-0.007,0.38)rectangle++(0.008,0.008) (-0.004,0.157)rectangle++(0.008,0.008) (0.054,0.49)rectangle++(0.008,0.008) (-0.002,0.203)rectangle++(0.008,0.008) (-0.004,0.083)rectangle++(0.008,0.008)
  (-0.035,0.456)rectangle++(0.008,0.008) (-0.182,0.615)rectangle++(0.008,0.008) (0.136,0.763)rectangle++(0.008,0.008) (0.055,0.918)rectangle++(0.008,0.008) (0.352,0.816)rectangle++(0.008,0.008)
  (0.146,0.941)rectangle++(0.008,0.008) (0.059,0.993)rectangle++(0.008,0.008) (0.381,0.849)rectangle++(0.008,0.008) (0.476,0.669)rectangle++(0.008,0.008) (0.456,0.552)rectangle++(0.008,0.008)
  (0.008,0.834)rectangle++(0.008,0.008) (-0.298,0.803)rectangle++(0.008,0.008) (-0.427,0.68)rectangle++(0.008,0.008) (-0.182,0.883)rectangle++(0.008,0.008) (0.232,0.885)rectangle++(0.008,0.008)
  (0.095,0.969)rectangle++(0.008,0.008) (-0.306,0.896)rectangle++(0.008,0.008) (-0.464,0.719)rectangle++(0.008,0.008) (0.044,0.911)rectangle++(0.008,0.008) (0.016,0.98)rectangle++(0.008,0.008)
  (0.357,0.859)rectangle++(0.008,0.008) (0.469,0.681)rectangle++(0.008,0.008) (0.457,0.56)rectangle++(0.008,0.008) (0.408,0.509)rectangle++(0.008,0.008) (0.008,0.212)rectangle++(0.008,0.008)
  (0.001,0.687)rectangle++(0.008,0.008) (0.245,0.73)rectangle++(0.008,0.008) (0.372,0.662)rectangle++(0.008,0.008) (0.406,0.586)rectangle++(0.008,0.008) (0.394,0.54)rectangle++(0.008,0.008)
  (0.372,0.522)rectangle++(0.008,0.008) (-0.02,0.791)rectangle++(0.008,0.008) (0.272,0.785)rectangle++(0.008,0.008) (0.004,0.327)rectangle++(0.008,0.008) (-0.001,0.735)rectangle++(0.008,0.008)
  (-0.003,0.906)rectangle++(0.008,0.008) (-0.003,0.978)rectangle++(0.008,0.008) (-0.004,1.009)rectangle++(0.008,0.008) (0.358,0.879)rectangle++(0.008,0.008) (-0.154,0.949)rectangle++(0.008,0.008)
  (0.268,0.905)rectangle++(0.008,0.008) (0.11,0.978)rectangle++(0.008,0.008) (-0.302,0.905)rectangle++(0.008,0.008) (-0.465,0.725)rectangle++(0.008,0.008) (-0.475,0.585)rectangle++(0.008,0.008)
  (-0.202,0.843)rectangle++(0.008,0.008) (0.208,0.874)rectangle++(0.008,0.008) (0.406,0.741)rectangle++(0.008,0.008) (0.45,0.61)rectangle++(0.008,0.008) (-0.016,0.859)rectangle++(0.008,0.008)
  (-0.318,0.806)rectangle++(0.008,0.008) (-0.437,0.674)rectangle++(0.008,0.008) (-0.186,0.881)rectangle++(0.008,0.008) (0.229,0.885)rectangle++(0.008,0.008) (0.42,0.739)rectangle++(0.008,0.008)
  (0.455,0.604)rectangle++(0.008,0.008) (0.01,0.251)rectangle++(0.008,0.008) (0.094,0.528)rectangle++(0.008,0.008) (-0.149,0.694)rectangle++(0.008,0.008) (-0.065,0.889)rectangle++(0.008,0.008)
  (-0.351,0.802)rectangle++(0.008,0.008) (-0.45,0.661)rectangle++(0.008,0.008) (0.029,0.879)rectangle++(0.008,0.008) (0.01,0.967)rectangle++(0.008,0.008) (0.349,0.855)rectangle++(0.008,0.008)
  (0.007,0.357)rectangle++(0.008,0.008) (-0.128,0.585)rectangle++(0.008,0.008) (-0.271,0.641)rectangle++(0.008,0.008) (-0.116,0.867)rectangle++(0.008,0.008) (0.256,0.854)rectangle++(0.008,0.008)
  (0.421,0.715)rectangle++(0.008,0.008) (0.447,0.593)rectangle++(0.008,0.008) (-0.011,0.85)rectangle++(0.008,0.008) (-0.007,0.955)rectangle++(0.008,0.008) (-0.348,0.853)rectangle++(0.008,0.008)
  (-0.467,0.685)rectangle++(0.008,0.008) (-0.199,0.885)rectangle++(0.008,0.008) (-0.086,0.97)rectangle++(0.008,0.008) (-0.389,0.832)rectangle++(0.008,0.008) (-0.478,0.661)rectangle++(0.008,0.008)
  (0.016,0.889)rectangle++(0.008,0.008) (0.324,0.817)rectangle++(0.008,0.008) (-0.147,0.908)rectangle++(0.008,0.008) (0.256,0.884)rectangle++(0.008,0.008) (0.105,0.969)rectangle++(0.008,0.008)
  (-0.001,0.405)rectangle++(0.008,0.008) (-0.004,0.168)rectangle++(0.008,0.008) (-0.065,0.494)rectangle++(0.008,0.008) (-0.03,0.805)rectangle++(0.008,0.008) (-0.015,0.936)rectangle++(0.008,0.008)
  (0.327,0.849)rectangle++(0.008,0.008) (-0.158,0.924)rectangle++(0.008,0.008) (0.257,0.895)rectangle++(0.008,0.008) (0.106,0.974)rectangle++(0.008,0.008) (0.042,1.007)rectangle++(0.008,0.008)
  (0.015,1.02)rectangle++(0.008,0.008) (-0.361,0.891)rectangle++(0.008,0.008) (-0.487,0.698)rectangle++(0.008,0.008) (-0.018,0.291)rectangle++(0.008,0.008) (-0.116,0.546)rectangle++(0.008,0.008)
  (0.141,0.707)rectangle++(0.008,0.008) (-0.192,0.793)rectangle++(0.008,0.008) (-0.083,0.931)rectangle++(0.008,0.008) (-0.037,0.989)rectangle++(0.008,0.008) (-0.374,0.858)rectangle++(0.008,0.008)
  (-0.481,0.678)rectangle++(0.008,0.008) (0.022,0.898)rectangle++(0.008,0.008) (-0.314,0.837)rectangle++(0.008,0.008) (-0.134,0.949)rectangle++(0.008,0.008) (0.277,0.898)rectangle++(0.008,0.008)
  (0.114,0.975)rectangle++(0.008,0.008) (-0.3,0.906)rectangle++(0.008,0.008) (0.186,0.937)rectangle++(0.008,0.008) (0.076,0.991)rectangle++(0.008,0.008) (0.388,0.842)rectangle++(0.008,0.008)
  (-0.127,0.943)rectangle++(0.008,0.008) (-0.008,0.394)rectangle++(0.008,0.008) (0.136,0.599)rectangle++(0.008,0.008) (-0.155,0.742)rectangle++(0.008,0.008) (-0.068,0.909)rectangle++(0.008,0.008)
  (-0.359,0.811)rectangle++(0.008,0.008) (-0.153,0.938)rectangle++(0.008,0.008) (-0.067,0.992)rectangle++(0.008,0.008) (0.323,0.893)rectangle++(0.008,0.008) (0.133,0.973)rectangle++(0.008,0.008)
  (-0.29,0.912)rectangle++(0.008,0.008) (-0.124,0.981)rectangle++(0.008,0.008) (-0.41,0.823)rectangle++(0.008,0.008) (-0.175,0.943)rectangle++(0.008,0.008) (0.256,0.91)rectangle++(0.008,0.008)
  (0.004,0.38)rectangle++(0.008,0.008) (-0.138,0.594)rectangle++(0.008,0.008) (-0.279,0.642)rectangle++(0.008,0.008) (-0.36,0.613)rectangle++(0.008,0.008) (-0.015,0.255)rectangle++(0.008,0.008)
  (0.084,0.538)rectangle++(0.008,0.008) (-0.158,0.695)rectangle++(0.008,0.008) (-0.324,0.681)rectangle++(0.008,0.008) (0.094,0.843)rectangle++(0.008,0.008) (0.037,0.952)rectangle++(0.008,0.008)
  (0.013,0.997)rectangle++(0.008,0.008) (-0.354,0.88)rectangle++(0.008,0.008) (-0.014,0.367)rectangle++(0.008,0.008) (-0.008,0.752)rectangle++(0.008,0.008) (-0.276,0.76)rectangle++(0.008,0.008)
  (-0.401,0.668)rectangle++(0.008,0.008) (-0.171,0.878)rectangle++(0.008,0.008) (-0.009,0.367)rectangle++(0.008,0.008) (-0.004,0.152)rectangle++(0.008,0.008) (-0.004,0.661)rectangle++(0.008,0.008)
  (-0.004,0.275)rectangle++(0.008,0.008) (0.096,0.544)rectangle++(0.008,0.008) (-0.001,0.226)rectangle++(0.008,0.008) (-0.085,0.522)rectangle++(0.008,0.008) (-0.038,0.817)rectangle++(0.008,0.008)
  (-0.313,0.779)rectangle++(0.008,0.008) (0.135,0.884)rectangle++(0.008,0.008) (0.054,0.969)rectangle++(0.008,0.008) (0.02,1.005)rectangle++(0.008,0.008) (0.006,1.02)rectangle++(0.008,0.008)
  (-0.365,0.888)rectangle++(0.008,0.008) (0.149,0.952)rectangle++(0.008,0.008) (0.001,0.398)rectangle++(0.008,0.008) (-0.145,0.601)rectangle++(0.008,0.008) (-0.285,0.642)rectangle++(0.008,0.008)
  (-0.122,0.867)rectangle++(0.008,0.008) (-0.008,0.362)rectangle++(0.008,0.008) (-0.136,0.582)rectangle++(0.008,0.008) (-0.274,0.637)rectangle++(0.008,0.008) (-0.356,0.612)rectangle++(0.008,0.008)
  (-0.152,0.855)rectangle++(0.008,0.008) (-0.008,0.357)rectangle++(0.008,0.008) (-0.006,0.748)rectangle++(0.008,0.008) (-0.273,0.759)rectangle++(0.008,0.008) (-0.117,0.916)rectangle++(0.008,0.008)
  (-0.384,0.796)rectangle++(0.008,0.008) (-0.164,0.932)rectangle++(0.008,0.008) (-0.411,0.787)rectangle++(0.008,0.008) (0.092,0.923)rectangle++(0.008,0.008) (0.036,0.985)rectangle++(0.008,0.008)
  (0.368,0.854)rectangle++(0.008,0.008) (0.472,0.675)rectangle++(0.008,0.008) (-0.029,0.896)rectangle++(0.008,0.008) (-0.014,0.974)rectangle++(0.008,0.008) (-0.358,0.859)rectangle++(0.008,0.008)
  (-0.474,0.684)rectangle++(0.008,0.008) (-0.201,0.885)rectangle++(0.008,0.008) (0.223,0.893)rectangle++(0.008,0.008) (-0.22,0.907)rectangle++(0.008,0.008) (-0.428,0.755)rectangle++(0.008,0.008)
  (-0.469,0.611)rectangle++(0.008,0.008) (-0.436,0.531)rectangle++(0.008,0.008) (-0.01,0.815)rectangle++(0.008,0.008) (0.286,0.793)rectangle++(0.008,0.008) (0.118,0.931)rectangle++(0.008,0.008)
  (0.047,0.989)rectangle++(0.008,0.008) (0.017,1.013)rectangle++(0.008,0.008) (-0.357,0.888)rectangle++(0.008,0.008) (-0.152,0.971)rectangle++(0.008,0.008) (-0.066,1.005)rectangle++(0.008,0.008)
  (0.328,0.9)rectangle++(0.008,0.008) (-0.175,0.948)rectangle++(0.008,0.008) (-0.076,0.996)rectangle++(0.008,0.008) (-0.394,0.847)rectangle++(0.008,0.008) (0.122,0.944)rectangle++(0.008,0.008)
  (0.049,0.994)rectangle++(0.008,0.008) (0.018,1.015)rectangle++(0.008,0.008) (-0.358,0.89)rectangle++(0.008,0.008) (-0.485,0.698)rectangle++(0.008,0.008) (-0.475,0.565)rectangle++(0.008,0.008)
  (-0.422,0.508)rectangle++(0.008,0.008) (-0.012,0.799)rectangle++(0.008,0.008) (0.279,0.786)rectangle++(0.008,0.008) (0.408,0.676)rectangle++(0.008,0.008) (0.169,0.882)rectangle++(0.008,0.008)
  (0.001,0.368)rectangle++(0.008,0.008) (-0.134,0.588)rectangle++(0.008,0.008) (-0.275,0.64)rectangle++(0.008,0.008) (0.102,0.807)rectangle++(0.008,0.008) (-0.245,0.825)rectangle++(0.008,0.008)
  (-0.105,0.944)rectangle++(0.008,0.008) (-0.047,0.994)rectangle++(0.008,0.008) (0.333,0.887)rectangle++(0.008,0.008) (-0.168,0.944)rectangle++(0.008,0.008) (-0.073,0.994)rectangle++(0.008,0.008)
  (0.321,0.897)rectangle++(0.008,0.008) (-0.177,0.944)rectangle++(0.008,0.008) (-0.422,0.787)rectangle++(0.008,0.008) (0.088,0.927)rectangle++(0.008,0.008) (-0.294,0.874)rectangle++(0.008,0.008)
  (0.177,0.921)rectangle++(0.008,0.008) (0.409,0.774)rectangle++(0.008,0.008) (-0.093,0.919)rectangle++(0.008,0.008) (-0.041,0.984)rectangle++(0.008,0.008) (-0.374,0.854)rectangle++(0.008,0.008)
  (0.133,0.94)rectangle++(0.008,0.008) (0.396,0.799)rectangle++(0.008,0.008) (0.465,0.64)rectangle++(0.008,0.008) (0.193,0.866)rectangle++(0.008,0.008) (-0.225,0.884)rectangle++(0.008,0.008)
  (-0.422,0.743)rectangle++(0.008,0.008) (0.072,0.907)rectangle++(0.008,0.008) (-0.295,0.859)rectangle++(0.008,0.008) (0.171,0.914)rectangle++(0.008,0.008) (-0.252,0.898)rectangle++(0.008,0.008)
  (-0.439,0.74)rectangle++(0.008,0.008) (-0.469,0.601)rectangle++(0.008,0.008) (-0.001,0.858)rectangle++(0.008,0.008) (-0.31,0.811)rectangle++(0.008,0.008) (0.147,0.898)rectangle++(0.008,0.008)
  (0.001,0.375)rectangle++(0.008,0.008) (0.133,0.588)rectangle++(0.008,0.008) (-0.152,0.735)rectangle++(0.008,0.008) (-0.066,0.907)rectangle++(0.008,0.008) (0.293,0.854)rectangle++(0.008,0.008)
  (0.438,0.702)rectangle++(0.008,0.008) (0.182,0.893)rectangle++(0.008,0.008) (0.074,0.973)rectangle++(0.008,0.008) (-0.317,0.89)rectangle++(0.008,0.008) (0.172,0.936)rectangle++(0.008,0.008)
  (0.412,0.783)rectangle++(0.008,0.008) (0.467,0.627)rectangle++(0.008,0.008) (0.437,0.536)rectangle++(0.008,0.008) (0.39,0.505)rectangle++(0.008,0.008) (-0.006,0.79)rectangle++(0.008,0.008)
  (0.279,0.779)rectangle++(0.008,0.008) (-0.155,0.875)rectangle++(0.008,0.008) (-0.387,0.764)rectangle++(0.008,0.008) (-0.453,0.63)rectangle++(0.008,0.008) (-0.436,0.546)rectangle++(0.008,0.008)
  (-0.005,0.821)rectangle++(0.008,0.008) (-0.299,0.793)rectangle++(0.008,0.008) (-0.423,0.675)rectangle++(0.008,0.008) (-0.18,0.881)rectangle++(0.008,0.008) (0.232,0.883)rectangle++(0.008,0.008)
  (0.421,0.737)rectangle++(0.008,0.008) (0.009,0.307)rectangle++(0.008,0.008) (0.113,0.554)rectangle++(0.008,0.008) (0.249,0.629)rectangle++(0.008,0.008) (0.102,0.862)rectangle++(0.008,0.008)
  (0.041,0.96)rectangle++(0.008,0.008) (-0.328,0.872)rectangle++(0.008,0.008) (0.161,0.932)rectangle++(0.008,0.008) (0.406,0.785)rectangle++(0.008,0.008) (-0.098,0.923)rectangle++(0.008,0.008)
  (0.284,0.873)rectangle++(0.008,0.008) (0.441,0.714)rectangle++(0.008,0.008) (0.183,0.898)rectangle++(0.008,0.008) (-0.241,0.895)rectangle++(0.008,0.008) (-0.433,0.742)rectangle++(0.008,0.008)
  (0.066,0.91)rectangle++(0.008,0.008) (0.355,0.809)rectangle++(0.008,0.008) (-0.13,0.916)rectangle++(0.008,0.008) (-0.057,0.982)rectangle++(0.008,0.008) (-0.38,0.848)rectangle++(0.008,0.008)
  (-0.015,0.354)rectangle++(0.008,0.008) (-0.009,0.746)rectangle++(0.008,0.008) (-0.006,0.911)rectangle++(0.008,0.008) (0.322,0.835)rectangle++(0.008,0.008) (-0.155,0.916)rectangle++(0.008,0.008)
  (0.256,0.89)rectangle++(0.008,0.008) (0.434,0.732)rectangle++(0.008,0.008) (-0.067,0.909)rectangle++(0.008,0.008) (-0.359,0.811)rectangle++(0.008,0.008) (-0.153,0.938)rectangle++(0.008,0.008)
  (0.264,0.9)rectangle++(0.008,0.008) (0.109,0.976)rectangle++(0.008,0.008) (0.397,0.823)rectangle++(0.008,0.008) (-0.116,0.937)rectangle++(0.008,0.008) (-0.051,0.991)rectangle++(0.008,0.008)
  (-0.381,0.854)rectangle++(0.008,0.008) (0.13,0.943)rectangle++(0.008,0.008) (0.052,0.994)rectangle++(0.008,0.008) (-0.334,0.892)rectangle++(0.008,0.008) (0.165,0.944)rectangle++(0.008,0.008)
  (0.412,0.789)rectangle++(0.008,0.008) (0.469,0.63)rectangle++(0.008,0.008) (-0.014,0.875)rectangle++(0.008,0.008) (-0.008,0.965)rectangle++(0.008,0.008) (0.34,0.86)rectangle++(0.008,0.008)
  (0.462,0.688)rectangle++(0.008,0.008) (-0.038,0.899)rectangle++(0.008,0.008) (-0.342,0.816)rectangle++(0.008,0.008) (-0.014,0.341)rectangle++(0.008,0.008) (-0.008,0.741)rectangle++(0.008,0.008)
  (0.26,0.758)rectangle++(0.008,0.008) (0.004,0.316)rectangle++(0.008,0.008) (-0.004,0.13)rectangle++(0.008,0.008) (-0.004,0.652)rectangle++(0.008,0.008) (-0.238,0.716)rectangle++(0.008,0.008)
  (-0.368,0.661)rectangle++(0.008,0.008) (0.067,0.85)rectangle++(0.008,0.008) (-0.277,0.832)rectangle++(0.008,0.008) (-0.118,0.947)rectangle++(0.008,0.008) (0.283,0.892)rectangle++(0.008,0.008)
  (0.005,0.372)rectangle++(0.008,0.008) (0.134,0.585)rectangle++(0.008,0.008) (-0.151,0.734)rectangle++(0.008,0.008) (-0.335,0.701)rectangle++(0.008,0.008) (0.096,0.856)rectangle++(0.008,0.008)
  (-0.265,0.845)rectangle++(0.008,0.008) (0.18,0.897)rectangle++(0.008,0.008) (0.002,0.375)rectangle++(0.008,0.008) (-0.002,0.755)rectangle++(0.008,0.008) (-0.274,0.764)rectangle++(0.008,0.008)
  (-0.402,0.67)rectangle++(0.008,0.008) (0.055,0.866)rectangle++(0.008,0.008) (-0.288,0.835)rectangle++(0.008,0.008) (-0.123,0.948)rectangle++(0.008,0.008) (-0.054,0.996)rectangle++(0.008,0.008)
  (0.33,0.891)rectangle++(0.008,0.008) (0.468,0.706)rectangle++(0.008,0.008) (0.465,0.572)rectangle++(0.008,0.008) (0.005,0.847)rectangle++(0.008,0.008) (0,0.953)rectangle++(0.008,0.008)
  (-0.344,0.855)rectangle++(0.008,0.008) (0.147,0.93)rectangle++(0.008,0.008) (0.399,0.789)rectangle++(0.008,0.008) (0.165,0.929)rectangle++(0.008,0.008) (0.067,0.988)rectangle++(0.008,0.008)
  (0.383,0.844)rectangle++(0.008,0.008) (0.476,0.665)rectangle++(0.008,0.008) (0.197,0.877)rectangle++(0.008,0.008) (0.002,0.366)rectangle++(0.008,0.008) (-0.004,0.151)rectangle++(0.008,0.008)
  (-0.059,0.487)rectangle++(0.008,0.008) (-0.205,0.621)rectangle++(0.008,0.008) (0.127,0.773)rectangle++(0.008,0.008) (0.334,0.724)rectangle++(0.008,0.008) (0.41,0.628)rectangle++(0.008,0.008)
  (-0.04,0.853)rectangle++(0.008,0.008) (0.285,0.821)rectangle++(0.008,0.008) (-0.166,0.896)rectangle++(0.008,0.008) (-0.4,0.77)rectangle++(0.008,0.008) (0.091,0.911)rectangle++(0.008,0.008)
  (-0.287,0.868)rectangle++(0.008,0.008) (0.178,0.916)rectangle++(0.008,0.008) (0.408,0.771)rectangle++(0.008,0.008) (-0.093,0.917)rectangle++(0.008,0.008) (-0.041,0.983)rectangle++(0.008,0.008)
  (0.331,0.88)rectangle++(0.008,0.008) (0.465,0.7)rectangle++(0.008,0.008) (0.193,0.892)rectangle++(0.008,0.008) (0.406,0.755)rectangle++(0.008,0.008) (-0.088,0.909)rectangle++(0.008,0.008)
  (-0.368,0.804)rectangle++(0.008,0.008) (-0.015,0.335)rectangle++(0.008,0.008) (-0.009,0.738)rectangle++(0.008,0.008) (-0.271,0.754)rectangle++(0.008,0.008) (-0.397,0.667)rectangle++(0.008,0.008)
  (-0.016,0.278)rectangle++(0.008,0.008) (-0.009,0.714)rectangle++(0.008,0.008) (-0.263,0.743)rectangle++(0.008,0.008) (0.144,0.85)rectangle++(0.008,0.008) (0.058,0.955)rectangle++(0.008,0.008)
  (0.367,0.832)rectangle++(0.008,0.008) (0.007,0.347)rectangle++(0.008,0.008) (0.126,0.572)rectangle++(0.008,0.008) (0.261,0.633)rectangle++(0.008,0.008) (0.107,0.863)rectangle++(0.008,0.008)
  (0.043,0.96)rectangle++(0.008,0.008) (-0.327,0.873)rectangle++(0.008,0.008) (-0.465,0.702)rectangle++(0.008,0.008) (-0.467,0.574)rectangle++(0.008,0.008) (-0.422,0.515)rectangle++(0.008,0.008)
  (-0.38,0.504)rectangle++(0.008,0.008) (0.005,0.783)rectangle++(0.008,0.008) (0.281,0.772)rectangle++(0.008,0.008) (0.403,0.669)rectangle++(0.008,0.008) (-0.058,0.869)rectangle++(0.008,0.008)
  (0.283,0.834)rectangle++(0.008,0.008) (0.427,0.697)rectangle++(0.008,0.008) (-0.057,0.89)rectangle++(0.008,0.008) (0.291,0.844)rectangle++(0.008,0.008) (0.12,0.952)rectangle++(0.008,0.008)
  (0,0.398)rectangle++(0.008,0.008) (-0.004,0.165)rectangle++(0.008,0.008) (-0.004,0.667)rectangle++(0.008,0.008) (-0.244,0.723)rectangle++(0.008,0.008) (-0.105,0.901)rectangle++(0.008,0.008)
  (-0.373,0.794)rectangle++(0.008,0.008) (-0.458,0.649)rectangle++(0.008,0.008) (0.022,0.876)rectangle++(0.008,0.008) (0.322,0.809)rectangle++(0.008,0.008) (0.435,0.671)rectangle++(0.008,0.008)
  (-0.044,0.882)rectangle++(0.008,0.008) (-0.339,0.806)rectangle++(0.008,0.008) (-0.446,0.667)rectangle++(0.008,0.008) (-0.446,0.565)rectangle++(0.008,0.008) (-0.19,0.835)rectangle++(0.008,0.008)
  (-0.082,0.948)rectangle++(0.008,0.008) (-0.38,0.823)rectangle++(0.008,0.008) (-0.471,0.66)rectangle++(0.008,0.008) (-0.2,0.875)rectangle++(0.008,0.008) (-0.408,0.748)rectangle++(0.008,0.008)
  (-0.173,0.912)rectangle++(0.008,0.008) (-0.408,0.774)rectangle++(0.008,0.008) (-0.016,0.323)rectangle++(0.008,0.008) (0.107,0.57)rectangle++(0.008,0.008) (0.043,0.837)rectangle++(0.008,0.008)
  (-0.283,0.817)rectangle++(0.008,0.008) (-0.425,0.692)rectangle++(0.008,0.008) (-0.181,0.888)rectangle++(0.008,0.008) (-0.078,0.971)rectangle++(0.008,0.008) (0.31,0.888)rectangle++(0.008,0.008)
  (0.005,0.371)rectangle++(0.008,0.008) (0,0.753)rectangle++(0.008,0.008) (-0.002,0.914)rectangle++(0.008,0.008) (0.325,0.835)rectangle++(0.008,0.008) (-0.153,0.917)rectangle++(0.008,0.008)
  (-0.401,0.783)rectangle++(0.008,0.008) (0.096,0.918)rectangle++(0.008,0.008) (0.371,0.802)rectangle++(0.008,0.008) (-0.12,0.918)rectangle++(0.008,0.008) (-0.053,0.983)rectangle++(0.008,0.008)
  (-0.379,0.85)rectangle++(0.008,0.008) (0.13,0.94)rectangle++(0.008,0.008) (0.394,0.8)rectangle++(0.008,0.008) (0.465,0.641)rectangle++(0.008,0.008) (-0.02,0.878)rectangle++(0.008,0.008)
  (0.304,0.825)rectangle++(0.008,0.008) (-0.159,0.905)rectangle++(0.008,0.008) (0.25,0.887)rectangle++(0.008,0.008) (-0.206,0.914)rectangle++(0.008,0.008) (-0.424,0.763)rectangle++(0.008,0.008)
  (0.078,0.917)rectangle++(0.008,0.008) (-0.295,0.866)rectangle++(0.008,0.008) (-0.448,0.71)rectangle++(0.008,0.008) (-0.462,0.584)rectangle++(0.008,0.008) (-0.423,0.521)rectangle++(0.008,0.008)
  (-0.383,0.506)rectangle++(0.008,0.008) (-0.359,0.514)rectangle++(0.008,0.008) (0.018,0.78)rectangle++(0.008,0.008) (-0.274,0.782)rectangle++(0.008,0.008) (-0.117,0.926)rectangle++(0.008,0.008)
  (-0.388,0.801)rectangle++(0.008,0.008) (-0.016,0.334)rectangle++(0.008,0.008) (-0.13,0.566)rectangle++(0.008,0.008) (-0.008,0.236)rectangle++(0.008,0.008) (-0.091,0.524)rectangle++(0.008,0.008)
  (-0.232,0.626)rectangle++(0.008,0.008) (-0.1,0.861)rectangle++(0.008,0.008) (-0.044,0.959)rectangle++(0.008,0.008) (-0.366,0.842)rectangle++(0.008,0.008) (-0.472,0.673)rectangle++(0.008,0.008)
  (-0.46,0.559)rectangle++(0.008,0.008) (-0.195,0.832)rectangle++(0.008,0.008) (-0.39,0.73)rectangle++(0.008,0.008) (-0.166,0.904)rectangle++(0.008,0.008) (-0.402,0.773)rectangle++(0.008,0.008)
  (0.091,0.913)rectangle++(0.008,0.008) (0.036,0.981)rectangle++(0.008,0.008) (0.366,0.852)rectangle++(0.008,0.008) (-0.14,0.939)rectangle++(0.008,0.008) (0.271,0.896)rectangle++(0.008,0.008)
  (-0.2,0.925)rectangle++(0.008,0.008) (-0.086,0.986)rectangle++(0.008,0.008) (0.312,0.898)rectangle++(0.008,0.008) (0.129,0.975)rectangle++(0.008,0.008) (0.406,0.816)rectangle++(0.008,0.008)
  (0.168,0.94)rectangle++(0.008,0.008) (0.412,0.786)rectangle++(0.008,0.008) (-0.096,0.926)rectangle++(0.008,0.008) (-0.378,0.808)rectangle++(0.008,0.008) (0.115,0.921)rectangle++(0.008,0.008)
  (0.381,0.796)rectangle++(0.008,0.008) (-0.114,0.919)rectangle++(0.008,0.008) (-0.05,0.984)rectangle++(0.008,0.008) (-0.005,0.411)rectangle++(0.008,0.008) (-0.153,0.605)rectangle++(0.008,0.008)
  (0.146,0.748)rectangle++(0.008,0.008) (0.333,0.706)rectangle++(0.008,0.008) (0.137,0.894)rectangle++(0.008,0.008) (0.381,0.776)rectangle++(0.008,0.008) (0.158,0.924)rectangle++(0.008,0.008)
  (0.401,0.782)rectangle++(0.008,0.008) (0.462,0.631)rectangle++(0.008,0.008) (-0.018,0.872)rectangle++(0.008,0.008) (-0.323,0.812)rectangle++(0.008,0.008) (0.141,0.903)rectangle++(0.008,0.008)
  (-0.261,0.882)rectangle++(0.008,0.008) (-0.012,0.368)rectangle++(0.008,0.008) (0.125,0.589)rectangle++(0.008,0.008) (0.05,0.845)rectangle++(0.008,0.008) (0.324,0.785)rectangle++(0.008,0.008)
  (0.428,0.659)rectangle++(0.008,0.008) (0.43,0.565)rectangle++(0.008,0.008) (-0.009,0.831)rectangle++(0.008,0.008) (-0.006,0.947)rectangle++(0.008,0.008) (0.335,0.851)rectangle++(0.008,0.008)
  (0.456,0.686)rectangle++(0.008,0.008) (0.189,0.886)rectangle++(0.008,0.008) (0.402,0.754)rectangle++(0.008,0.008) (0.167,0.914)rectangle++(0.008,0.008) (0.402,0.775)rectangle++(0.008,0.008)
  (0.46,0.627)rectangle++(0.008,0.008) (0.191,0.861)rectangle++(0.008,0.008) (0.394,0.742)rectangle++(0.008,0.008) (0.444,0.615)rectangle++(0.008,0.008) (-0.02,0.859)rectangle++(0.008,0.008)
  (-0.319,0.805)rectangle++(0.008,0.008) (0.141,0.898)rectangle++(0.008,0.008) (0.057,0.975)rectangle++(0.008,0.008) (0.373,0.842)rectangle++(0.008,0.008) (0.47,0.668)rectangle++(0.008,0.008)
  (-0.027,0.892)rectangle++(0.008,0.008) (0.306,0.834)rectangle++(0.008,0.008) (0.126,0.948)rectangle++(0.008,0.008) (0.395,0.805)rectangle++(0.008,0.008) (0.164,0.936)rectangle++(0.008,0.008)
  (0.408,0.786)rectangle++(0.008,0.008) (0.466,0.63)rectangle++(0.008,0.008) (0.437,0.538)rectangle++(0.008,0.008) (0.391,0.506)rectangle++(0.008,0.008) (0.162,0.81)rectangle++(0.008,0.008)
  (0.066,0.938)rectangle++(0.008,0.008) (0.364,0.822)rectangle++(0.008,0.008) (0.459,0.662)rectangle++(0.008,0.008) (-0.03,0.886)rectangle++(0.008,0.008) (-0.005,0.37)rectangle++(0.008,0.008)
  (0.129,0.587)rectangle++(0.008,0.008) (0,0.244)rectangle++(0.008,0.008) (-0.004,0.1)rectangle++(0.008,0.008) (-0.041,0.464)rectangle++(0.008,0.008) (-0.005,0.192)rectangle++(0.008,0.008)
  (-0.075,0.505)rectangle++(0.008,0.008) (-0.218,0.624)rectangle++(0.008,0.008) (0.122,0.779)rectangle++(0.008,0.008) (-0.226,0.819)rectangle++(0.008,0.008) (-0.399,0.713)rectangle++(0.008,0.008)
  (-0.17,0.897)rectangle++(0.008,0.008) (0.242,0.887)rectangle++(0.008,0.008) (0.099,0.97)rectangle++(0.008,0.008) (-0.305,0.898)rectangle++(0.008,0.008) (-0.013,0.375)rectangle++(0.008,0.008)
  (-0.143,0.586)rectangle++(0.008,0.008) (-0.008,0.244)rectangle++(0.008,0.008) (0.083,0.531)rectangle++(0.008,0.008) (-0.001,0.221)rectangle++(0.008,0.008) (-0.003,0.69)rectangle++(0.008,0.008)
  (0.244,0.733)rectangle++(0.008,0.008) (0.1,0.906)rectangle++(0.008,0.008) (0.04,0.978)rectangle++(0.008,0.008) (-0.335,0.88)rectangle++(0.008,0.008) (-0.471,0.702)rectangle++(0.008,0.008)
  (-0.47,0.572)rectangle++(0.008,0.008) (-0.018,0.238)rectangle++(0.008,0.008) (-0.097,0.522)rectangle++(0.008,0.008) (-0.234,0.623)rectangle++(0.008,0.008) (-0.333,0.62)rectangle++(0.008,0.008)
  (-0.377,0.584)rectangle++(0.008,0.008) (-0.385,0.551)rectangle++(0.008,0.008) (-0.164,0.829)rectangle++(0.008,0.008) (-0.375,0.74)rectangle++(0.008,0.008) (-0.16,0.908)rectangle++(0.008,0.008)
  (0.251,0.889)rectangle++(0.008,0.008) (0.103,0.971)rectangle++(0.008,0.008) (0.041,1.005)rectangle++(0.008,0.008) (0.377,0.861)rectangle++(0.008,0.008) (0.479,0.675)rectangle++(0.008,0.008)
  (-0.026,0.899)rectangle++(0.008,0.008) (-0.013,0.975)rectangle++(0.008,0.008) (-0.004,0.407)rectangle++(0.008,0.008) (0.143,0.604)rectangle++(0.008,0.008) (0.28,0.642)rectangle++(0.008,0.008)
  (0.356,0.61)rectangle++(0.008,0.008) (0.38,0.568)rectangle++(0.008,0.008) (-0.033,0.814)rectangle++(0.008,0.008) (-0.005,0.34)rectangle++(0.008,0.008) (-0.127,0.573)rectangle++(0.008,0.008)
  (-0.056,0.838)rectangle++(0.008,0.008) (-0.026,0.95)rectangle++(0.008,0.008) (-0.354,0.844)rectangle++(0.008,0.008) (0.139,0.929)rectangle++(0.008,0.008) (0.056,0.988)rectangle++(0.008,0.008)
  (0.378,0.848)rectangle++(0.008,0.008) (-0.134,0.942)rectangle++(0.008,0.008) (-0.401,0.802)rectangle++(0.008,0.008) (0.102,0.926)rectangle++(0.008,0.008) (0.377,0.803)rectangle++(0.008,0.008)
  (-0.118,0.921)rectangle++(0.008,0.008) (-0.386,0.798)rectangle++(0.008,0.008) (0.108,0.919)rectangle++(0.008,0.008) (-0.283,0.878)rectangle++(0.008,0.008) (-0.121,0.966)rectangle++(0.008,0.008)
  (-0.404,0.818)rectangle++(0.008,0.008) (0.107,0.934)rectangle++(0.008,0.008) (0.382,0.805)rectangle++(0.008,0.008) (0.158,0.936)rectangle++(0.008,0.008) (-0.266,0.903)rectangle++(0.008,0.008)
  (-0.114,0.977)rectangle++(0.008,0.008) (-0.05,1.008)rectangle++(0.008,0.008) (0.336,0.895)rectangle++(0.008,0.008) (0.006,0.374)rectangle++(0.008,0.008) (0.135,0.585)rectangle++(0.008,0.008)
  (-0.151,0.735)rectangle++(0.008,0.008) (-0.066,0.906)rectangle++(0.008,0.008) (0.293,0.854)rectangle++(0.008,0.008) (-0.175,0.914)rectangle++(0.008,0.008) (0.246,0.897)rectangle++(0.008,0.008)
  (-0.211,0.917)rectangle++(0.008,0.008) (0.23,0.911)rectangle++(0.008,0.008) (-0.224,0.918)rectangle++(0.008,0.008) (-0.096,0.983)rectangle++(0.008,0.008) (0.306,0.9)rectangle++(0.008,0.008)
  (0.005,0.376)rectangle++(0.008,0.008) (0.136,0.586)rectangle++(0.008,0.008) (0,0.244)rectangle++(0.008,0.008) (-0.002,0.7)rectangle++(0.008,0.008) (0.248,0.737)rectangle++(0.008,0.008)
  (0.102,0.907)rectangle++(0.008,0.008) (0.37,0.795)rectangle++(0.008,0.008) (-0.118,0.915)rectangle++(0.008,0.008) (-0.384,0.795)rectangle++(0.008,0.008) (-0.463,0.645)rectangle++(0.008,0.008)
  (-0.197,0.869)rectangle++(0.008,0.008) (-0.01,0.363)rectangle++(0.008,0.008) (0.124,0.586)rectangle++(0.008,0.008) (-0.156,0.731)rectangle++(0.008,0.008) (-0.068,0.905)rectangle++(0.008,0.008)
  (0.291,0.854)rectangle++(0.008,0.008) (0.437,0.703)rectangle++(0.008,0.008) (0.181,0.893)rectangle++(0.008,0.008) (-0.24,0.892)rectangle++(0.008,0.008) (-0.432,0.741)rectangle++(0.008,0.008)
  (-0.466,0.604)rectangle++(0.008,0.008) (-0.018,0.251)rectangle++(0.008,0.008) (0.081,0.538)rectangle++(0.008,0.008) (-0.159,0.694)rectangle++(0.008,0.008) (-0.069,0.889)rectangle++(0.008,0.008)
  (0.285,0.847)rectangle++(0.008,0.008) (-0.176,0.908)rectangle++(0.008,0.008) (-0.076,0.979)rectangle++(0.008,0.008) (-0.034,1.009)rectangle++(0.008,0.008) (-0.005,0.421)rectangle++(0.008,0.008)
  (0.148,0.611)rectangle++(0.008,0.008) (0.06,0.854)rectangle++(0.008,0.008) (0.332,0.786)rectangle++(0.008,0.008) (0.137,0.928)rectangle++(0.008,0.008) (0,0.387)rectangle++(0.008,0.008)
  (-0.142,0.596)rectangle++(0.008,0.008) (-0.281,0.641)rectangle++(0.008,0.008) (-0.361,0.612)rectangle++(0.008,0.008) (-0.387,0.57)rectangle++(0.008,0.008) (0.026,0.815)rectangle++(0.008,0.008)
  (-0.283,0.801)rectangle++(0.008,0.008) (-0.121,0.934)rectangle++(0.008,0.008) (-0.393,0.803)rectangle++(0.008,0.008) (-0.167,0.935)rectangle++(0.008,0.008) (-0.414,0.787)rectangle++(0.008,0.008)
  (0.091,0.924)rectangle++(0.008,0.008) (0.036,0.986)rectangle++(0.008,0.008) (0.368,0.854)rectangle++(0.008,0.008) (-0.141,0.941)rectangle++(0.008,0.008) (0.271,0.897)rectangle++(0.008,0.008)
  (0.443,0.729)rectangle++(0.008,0.008) (0.009,0.304)rectangle++(0.008,0.008) (-0.004,0.125)rectangle++(0.008,0.008) (0.042,0.475)rectangle++(0.008,0.008) (-0.003,0.197)rectangle++(0.008,0.008)
  (-0.075,0.508)rectangle++(0.008,0.008) (-0.006,0.211)rectangle++(0.008,0.008) (-0.082,0.514)rectangle++(0.008,0.008) (0.145,0.68)rectangle++(0.008,0.008) (-0.18,0.782)rectangle++(0.008,0.008)
  (0.196,0.838)rectangle++(0.008,0.008) (0.08,0.95)rectangle++(0.008,0.008) (0.031,0.997)rectangle++(0.008,0.008) (0.369,0.861)rectangle++(0.008,0.008) (0.475,0.678)rectangle++(0.008,0.008)
  (0.197,0.882)rectangle++(0.008,0.008) (-0.228,0.893)rectangle++(0.008,0.008) (-0.098,0.973)rectangle++(0.008,0.008) (-0.044,1.006)rectangle++(0.008,0.008) (-0.383,0.864)rectangle++(0.008,0.008)
  (-0.015,0.36)rectangle++(0.008,0.008) (-0.009,0.749)rectangle++(0.008,0.008) (-0.275,0.758)rectangle++(0.008,0.008) (-0.4,0.668)rectangle++(0.008,0.008) (-0.17,0.878)rectangle++(0.008,0.008)
  (-0.395,0.76)rectangle++(0.008,0.008) (0.09,0.905)rectangle++(0.008,0.008) (0.035,0.978)rectangle++(0.008,0.008) (0.365,0.851)rectangle++(0.008,0.008) (0.151,0.955)rectangle++(0.008,0.008)
  (0.061,0.999)rectangle++(0.008,0.008) (0.023,1.017)rectangle++(0.008,0.008) (0.373,0.873)rectangle++(0.008,0.008) (-0.145,0.951)rectangle++(0.008,0.008) (-0.063,0.997)rectangle++(0.008,0.008)
  (-0.029,1.017)rectangle++(0.008,0.008) (-0.005,0.425)rectangle++(0.008,0.008) (-0.004,0.176)rectangle++(0.008,0.008) (-0.068,0.498)rectangle++(0.008,0.008) (-0.031,0.807)rectangle++(0.008,0.008)
  (0.273,0.796)rectangle++(0.008,0.008) (-0.163,0.881)rectangle++(0.008,0.008) (0.239,0.877)rectangle++(0.008,0.008) (0.003,0.366)rectangle++(0.008,0.008) (-0.004,0.151)rectangle++(0.008,0.008)
  (-0.004,0.061)rectangle++(0.008,0.008) (-0.027,0.446)rectangle++(0.008,0.008) (0.146,0.63)rectangle++(0.008,0.008) (-0.162,0.759)rectangle++(0.008,0.008) (-0.07,0.917)rectangle++(0.008,0.008)
  (0.294,0.86)rectangle++(0.008,0.008) (0.121,0.959)rectangle++(0.008,0.008) (0.049,1)rectangle++(0.008,0.008) (-0.339,0.894)rectangle++(0.008,0.008) (0.164,0.946)rectangle++(0.008,0.008)
  (0.066,0.995)rectangle++(0.008,0.008) (0.026,1.016)rectangle++(0.008,0.008) (-0.355,0.893)rectangle++(0.008,0.008) (0.156,0.951)rectangle++(0.008,0.008) (0.063,0.997)rectangle++(0.008,0.008)
  (0.384,0.85)rectangle++(0.008,0.008) (0.478,0.668)rectangle++(0.008,0.008) (0.01,0.278)rectangle++(0.008,0.008) (0.103,0.54)rectangle++(0.008,0.008) (0.239,0.626)rectangle++(0.008,0.008)
  (0.098,0.861)rectangle++(0.008,0.008) (-0.266,0.848)rectangle++(0.008,0.008) (0.18,0.899)rectangle++(0.008,0.008) (0.002,0.375)rectangle++(0.008,0.008) (0.134,0.587)rectangle++(0.008,0.008)
  (0.27,0.637)rectangle++(0.008,0.008) (0.111,0.865)rectangle++(0.008,0.008) (-0.001,0.361)rectangle++(0.008,0.008) (-0.003,0.749)rectangle++(0.008,0.008) (0.266,0.76)rectangle++(0.008,0.008)
  (0.109,0.917)rectangle++(0.008,0.008) (0.377,0.796)rectangle++(0.008,0.008) (-0.116,0.918)rectangle++(0.008,0.008) (0.274,0.877)rectangle++(0.008,0.008) (-0.192,0.918)rectangle++(0.008,0.008)
  (-0.01,0.383)rectangle++(0.008,0.008) (-0.004,0.159)rectangle++(0.008,0.008) (0.054,0.49)rectangle++(0.008,0.008) (0.02,0.804)rectangle++(0.008,0.008) (-0.281,0.794)rectangle++(0.008,0.008)
  (0.154,0.88)rectangle++(0.008,0.008) (0.384,0.763)rectangle++(0.008,0.008) (0.159,0.918)rectangle++(0.008,0.008) (-0.259,0.896)rectangle++(0.008,0.008) (-0.111,0.974)rectangle++(0.008,0.008)
  (-0.402,0.825)rectangle++(0.008,0.008) (-0.482,0.653)rectangle++(0.008,0.008) (0.012,0.887)rectangle++(0.008,0.008) (0.003,0.97)rectangle++(0.008,0.008) (-0.349,0.864)rectangle++(0.008,0.008)
  (-0.149,0.96)rectangle++(0.008,0.008) (0.274,0.908)rectangle++(0.008,0.008) (-0.203,0.932)rectangle++(0.008,0.008) (-0.429,0.773)rectangle++(0.008,0.008) (-0.183,0.922)rectangle++(0.008,0.008)
  (-0.079,0.985)rectangle++(0.008,0.008) (-0.391,0.841)rectangle++(0.008,0.008) (-0.167,0.951)rectangle++(0.008,0.008) (0.263,0.911)rectangle++(0.008,0.008) (0.444,0.739)rectangle++(0.008,0.008)
  (0.466,0.595)rectangle++(0.008,0.008) (0.01,0.248)rectangle++(0.008,0.008) (-0.087,0.536)rectangle++(0.008,0.008) (-0.039,0.823)rectangle++(0.008,0.008) (-0.315,0.781)rectangle++(0.008,0.008)
  (0.134,0.886)rectangle++(0.008,0.008) (0.377,0.773)rectangle++(0.008,0.008) (0.156,0.923)rectangle++(0.008,0.008) (0.4,0.782)rectangle++(0.008,0.008) (0.166,0.926)rectangle++(0.008,0.008)
  (0.067,0.987)rectangle++(0.008,0.008) (0.026,1.012)rectangle++(0.008,0.008) (-0.353,0.891)rectangle++(0.008,0.008) (-0.151,0.972)rectangle++(0.008,0.008) (0.277,0.914)rectangle++(0.008,0.008)
  (0.004,0.382)rectangle++(0.008,0.008) (0,0.758)rectangle++(0.008,0.008) (0.27,0.763)rectangle++(0.008,0.008) (0.395,0.669)rectangle++(0.008,0.008) (-0.062,0.866)rectangle++(0.008,0.008)
  (0.28,0.834)rectangle++(0.008,0.008) (-0.173,0.901)rectangle++(0.008,0.008) (-0.404,0.769)rectangle++(0.008,0.008) (-0.172,0.921)rectangle++(0.008,0.008) (-0.075,0.984)rectangle++(0.008,0.008)
  (-0.389,0.842)rectangle++(0.008,0.008) (-0.166,0.952)rectangle++(0.008,0.008) (0.263,0.91)rectangle++(0.008,0.008) (0.108,0.98)rectangle++(0.008,0.008) (-0.304,0.906)rectangle++(0.008,0.008)
  (-0.13,0.978)rectangle++(0.008,0.008) (-0.008,0.408)rectangle++(0.008,0.008) (-0.006,0.769)rectangle++(0.008,0.008) (-0.281,0.769)rectangle++(0.008,0.008) (-0.406,0.67)rectangle++(0.008,0.008)
  (0.053,0.868)rectangle++(0.008,0.008) (-0.289,0.835)rectangle++(0.008,0.008) (-0.434,0.697)rectangle++(0.008,0.008) (-0.451,0.583)rectangle++(0.008,0.008) (-0.192,0.843)rectangle++(0.008,0.008)
  (0.212,0.87)rectangle++(0.008,0.008) (-0.217,0.893)rectangle++(0.008,0.008) (-0.094,0.973)rectangle++(0.008,0.008) (-0.394,0.83)rectangle++(0.008,0.008) (-0.48,0.658)rectangle++(0.008,0.008)
  (0.015,0.889)rectangle++(0.008,0.008) (0.004,0.971)rectangle++(0.008,0.008) (-0.001,1.005)rectangle++(0.008,0.008) (-0.363,0.879)rectangle++(0.008,0.008) (0.147,0.948)rectangle++(0.008,0.008)
  (0.059,0.996)rectangle++(0.008,0.008) (-0.332,0.896)rectangle++(0.008,0.008) (0.167,0.944)rectangle++(0.008,0.008) (0.413,0.788)rectangle++(0.008,0.008) (0.469,0.629)rectangle++(0.008,0.008)
  (0.195,0.862)rectangle++(0.008,0.008) (0.396,0.741)rectangle++(0.008,0.008) (-0.087,0.899)rectangle++(0.008,0.008) (-0.365,0.799)rectangle++(0.008,0.008) (0.118,0.912)rectangle++(0.008,0.008)
  (0.379,0.791)rectangle++(0.008,0.008) (0.007,0.33)rectangle++(0.008,0.008) (-0.004,0.136)rectangle++(0.008,0.008) (-0.054,0.48)rectangle++(0.008,0.008) (-0.2,0.619)rectangle++(0.008,0.008)
  (-0.316,0.631)rectangle++(0.008,0.008) (0.08,0.818)rectangle++(0.008,0.008) (0.328,0.762)rectangle++(0.008,0.008) (0.135,0.918)rectangle++(0.008,0.008) (-0.269,0.887)rectangle++(0.008,0.008)
  (0.193,0.918)rectangle++(0.008,0.008) (0.002,0.383)rectangle++(0.008,0.008) (0.137,0.591)rectangle++(0.008,0.008) (0.273,0.638)rectangle++(0.008,0.008) (-0.107,0.808)rectangle++(0.008,0.008)
  (-0.341,0.751)rectangle++(0.008,0.008) (-0.145,0.913)rectangle++(0.008,0.008) (-0.396,0.785)rectangle++(0.008,0.008) (-0.465,0.636)rectangle++(0.008,0.008) (-0.443,0.544)rectangle++(0.008,0.008)
  (-0.401,0.51)rectangle++(0.008,0.008) (-0.002,0.792)rectangle++(0.008,0.008) (0.281,0.779)rectangle++(0.008,0.008) (-0.153,0.876)rectangle++(0.008,0.008) (-0.067,0.966)rectangle++(0.008,0.008)
  (-0.379,0.837)rectangle++(0.008,0.008) (0.125,0.934)rectangle++(0.008,0.008) (-0.28,0.891)rectangle++(0.008,0.008) (-0.45,0.726)rectangle++(0.008,0.008) (-0.017,0.303)rectangle++(0.008,0.008)
  (-0.004,0.125)rectangle++(0.008,0.008) (-0.05,0.475)rectangle++(0.008,0.008) (0.146,0.651)rectangle++(0.008,0.008) (0.059,0.871)rectangle++(0.008,0.008) (-0.288,0.838)rectangle++(0.008,0.008)
  (-0.123,0.95)rectangle++(0.008,0.008) (-0.399,0.809)rectangle++(0.008,0.008) (0.106,0.929)rectangle++(0.008,0.008) (-0.001,0.388)rectangle++(0.008,0.008) (-0.003,0.761)rectangle++(0.008,0.008)
  (-0.004,0.317)rectangle++(0.008,0.008) (0.111,0.563)rectangle++(0.008,0.008) (0.044,0.834)rectangle++(0.008,0.008) (0.016,0.948)rectangle++(0.008,0.008) (0.004,0.996)rectangle++(0.008,0.008)
  (0,1.016)rectangle++(0.008,0.008) (-0.367,0.883)rectangle++(0.008,0.008) (-0.487,0.692)rectangle++(0.008,0.008) (-0.473,0.562)rectangle++(0.008,0.008) (-0.201,0.834)rectangle++(0.008,0.008)
  (-0.393,0.728)rectangle++(0.008,0.008) (-0.016,0.304)rectangle++(0.008,0.008) (-0.119,0.552)rectangle++(0.008,0.008) (0.142,0.711)rectangle++(0.008,0.008) (-0.193,0.795)rectangle++(0.008,0.008)
  (-0.376,0.714)rectangle++(0.008,0.008) (-0.43,0.611)rectangle++(0.008,0.008) (0.021,0.849)rectangle++(0.008,0.008) (0.006,0.954)rectangle++(0.008,0.008) (-0.341,0.858)rectangle++(0.008,0.008)
  (-0.466,0.689)rectangle++(0.008,0.008) (-0.463,0.568)rectangle++(0.008,0.008) (-0.009,0.841)rectangle++(0.008,0.008) (-0.004,0.351)rectangle++(0.008,0.008) (-0.131,0.578)rectangle++(0.008,0.008)
  (0.146,0.727)rectangle++(0.008,0.008) (0.326,0.697)rectangle++(0.008,0.008) (-0.104,0.854)rectangle++(0.008,0.008) (-0.007,0.356)rectangle++(0.008,0.008) (-0.134,0.58)rectangle++(0.008,0.008)
  (-0.059,0.841)rectangle++(0.008,0.008) (-0.331,0.783)rectangle++(0.008,0.008) (-0.014,0.326)rectangle++(0.008,0.008) (0.11,0.571)rectangle++(0.008,0.008) (-0.157,0.719)rectangle++(0.008,0.008)
  (-0.068,0.9)rectangle++(0.008,0.008) (0.289,0.852)rectangle++(0.008,0.008) (-0.176,0.912)rectangle++(0.008,0.008) (-0.41,0.773)rectangle++(0.008,0.008) (-0.467,0.626)rectangle++(0.008,0.008)
  (0.01,0.869)rectangle++(0.008,0.008) (-0.31,0.82)rectangle++(0.008,0.008) (-0.132,0.942)rectangle++(0.008,0.008) (-0.008,0.393)rectangle++(0.008,0.008) (-0.006,0.763)rectangle++(0.008,0.008)
  (-0.279,0.766)rectangle++(0.008,0.008) (-0.119,0.919)rectangle++(0.008,0.008) (-0.007,0.384)rectangle++(0.008,0.008) (0.133,0.594)rectangle++(0.008,0.008) (0.272,0.641)rectangle++(0.008,0.008)
  (0.112,0.867)rectangle++(0.008,0.008) (0.36,0.773)rectangle++(0.008,0.008) (0.149,0.922)rectangle++(0.008,0.008) (-0.265,0.894)rectangle++(0.008,0.008) (-0.114,0.973)rectangle++(0.008,0.008)
  (-0.403,0.823)rectangle++(0.008,0.008) (0.109,0.937)rectangle++(0.008,0.008) (-0.001,0.391)rectangle++(0.008,0.008) (-0.004,0.162)rectangle++(0.008,0.008) (-0.004,0.666)rectangle++(0.008,0.008)
  (0.235,0.722)rectangle++(0.008,0.008) (0.003,0.301)rectangle++(0.008,0.008) (-0.11,0.558)rectangle++(0.008,0.008) (-0.253,0.635)rectangle++(0.008,0.008) (-0.346,0.619)rectangle++(0.008,0.008)
  (-0.148,0.858)rectangle++(0.008,0.008) (-0.377,0.759)rectangle++(0.008,0.008) (-0.447,0.631)rectangle++(0.008,0.008) (-0.433,0.548)rectangle++(0.008,0.008) (-0.397,0.515)rectangle++(0.008,0.008)
  (-0.169,0.814)rectangle++(0.008,0.008) (-0.073,0.94)rectangle++(0.008,0.008) (0.301,0.872)rectangle++(0.008,0.008) (0.005,0.364)rectangle++(0.008,0.008) (-0.131,0.587)rectangle++(0.008,0.008)
  (-0.273,0.641)rectangle++(0.008,0.008) (-0.357,0.615)rectangle++(0.008,0.008) (-0.386,0.573)rectangle++(0.008,0.008) (0.027,0.816)rectangle++(0.008,0.008) (0.303,0.78)rectangle++(0.008,0.008)
  (0.125,0.925)rectangle++(0.008,0.008) (-0.277,0.887)rectangle++(0.008,0.008) (-0.447,0.726)rectangle++(0.008,0.008) (-0.19,0.902)rectangle++(0.008,0.008) (-0.082,0.977)rectangle++(0.008,0.008)
  (0.31,0.892)rectangle++(0.008,0.008) (0.46,0.713)rectangle++(0.008,0.008) (0.464,0.578)rectangle++(0.008,0.008) (0.002,0.849)rectangle++(0.008,0.008) (-0.001,0.954)rectangle++(0.008,0.008)
  (-0.345,0.855)rectangle++(0.008,0.008) (0.147,0.93)rectangle++(0.008,0.008) (-0.269,0.897)rectangle++(0.008,0.008) (-0.447,0.733)rectangle++(0.008,0.008) (-0.19,0.906)rectangle++(0.008,0.008)
  (0.236,0.898)rectangle++(0.008,0.008) (0.003,0.375)rectangle++(0.008,0.008) (-0.136,0.592)rectangle++(0.008,0.008) (-0.277,0.641)rectangle++(0.008,0.008) (-0.359,0.613)rectangle++(0.008,0.008)
  (-0.153,0.855)rectangle++(0.008,0.008) (-0.067,0.957)rectangle++(0.008,0.008) (-0.03,1)rectangle++(0.008,0.008) (-0.374,0.865)rectangle++(0.008,0.008) (0.137,0.946)rectangle++(0.008,0.008)
  (0.4,0.8)rectangle++(0.008,0.008) (0.467,0.639)rectangle++(0.008,0.008) (0.01,0.266)rectangle++(0.008,0.008) (0.099,0.534)rectangle++(0.008,0.008) (0.235,0.625)rectangle++(0.008,0.008)
  (0.33,0.618)rectangle++(0.008,0.008) (0.006,0.257)rectangle++(0.008,0.008) (-0.093,0.539)rectangle++(0.008,0.008) (-0.007,0.224)rectangle++(0.008,0.008) (-0.005,0.692)rectangle++(0.008,0.008)
  (-0.253,0.734)rectangle++(0.008,0.008) (-0.381,0.664)rectangle++(0.008,0.008) (-0.415,0.587)rectangle++(0.008,0.008) (0.019,0.833)rectangle++(0.008,0.008) (-0.292,0.807)rectangle++(0.008,0.008)
  (-0.425,0.684)rectangle++(0.008,0.008) (-0.442,0.58)rectangle++(0.008,0.008) (-0.188,0.841)rectangle++(0.008,0.008) (0.214,0.868)rectangle++(0.008,0.008) (0.407,0.737)rectangle++(0.008,0.008)
  (0.169,0.907)rectangle++(0.008,0.008) (0.068,0.979)rectangle++(0.008,0.008) (0.026,1.009)rectangle++(0.008,0.008) (0.372,0.868)rectangle++(0.008,0.008) (0.479,0.68)rectangle++(0.008,0.008)
  (0.461,0.556)rectangle++(0.008,0.008) (0.409,0.506)rectangle++(0.008,0.008) (0.003,0.796)rectangle++(0.008,0.008) (-0.001,0.932)rectangle++(0.008,0.008) (-0.004,0.389)rectangle++(0.008,0.008)
  (-0.144,0.596)rectangle++(0.008,0.008) (0.146,0.74)rectangle++(0.008,0.008) (0.001,0.309)rectangle++(0.008,0.008) (-0.002,0.727)rectangle++(0.008,0.008) (0.258,0.75)rectangle++(0.008,0.008)
  (0.385,0.667)rectangle++(0.008,0.008) (0.413,0.584)rectangle++(0.008,0.008) (0.397,0.536)rectangle++(0.008,0.008) (0.372,0.52)rectangle++(0.008,0.008) (0.154,0.816)rectangle++(0.008,0.008)
  (0.361,0.734)rectangle++(0.008,0.008) (-0.101,0.884)rectangle++(0.008,0.008) (-0.365,0.787)rectangle++(0.008,0.008) (-0.156,0.928)rectangle++(0.008,0.008) (-0.406,0.788)rectangle++(0.008,0.008)
  (0.095,0.922)rectangle++(0.008,0.008) (-0.289,0.874)rectangle++(0.008,0.008) (-0.448,0.716)rectangle++(0.008,0.008) (0.05,0.903)rectangle++(0.008,0.008) (0.019,0.977)rectangle++(0.008,0.008)
  (0.357,0.856)rectangle++(0.008,0.008) (0.468,0.68)rectangle++(0.008,0.008) (-0.033,0.897)rectangle++(0.008,0.008) (-0.016,0.975)rectangle++(0.008,0.008) (-0.359,0.859)rectangle++(0.008,0.008)
  (0.142,0.937)rectangle++(0.008,0.008) (0.399,0.794)rectangle++(0.008,0.008) (0.008,0.331)rectangle++(0.008,0.008) (-0.118,0.574)rectangle++(0.008,0.008) (-0.262,0.639)rectangle++(0.008,0.008)
  (0.108,0.802)rectangle++(0.008,0.008) (0.335,0.745)rectangle++(0.008,0.008) (-0.116,0.879)rectangle++(0.008,0.008) (-0.371,0.78)rectangle++(0.008,0.008) (0.108,0.905)rectangle++(0.008,0.008)
  (0.043,0.978)rectangle++(0.008,0.008) (-0.003,0.408)rectangle++(0.008,0.008) (-0.003,0.769)rectangle++(0.008,0.008) (0.272,0.769)rectangle++(0.008,0.008) (0.112,0.921)rectangle++(0.008,0.008)
  (0.045,0.984)rectangle++(0.008,0.008) (0.371,0.85)rectangle++(0.008,0.008) (-0.138,0.94)rectangle++(0.008,0.008) (-0.06,0.993)rectangle++(0.008,0.008) (-0.028,1.015)rectangle++(0.008,0.008)
  (-0.014,1.024)rectangle++(0.008,0.008) (-0.004,0.428)rectangle++(0.008,0.008) (-0.004,0.177)rectangle++(0.008,0.008) (-0.069,0.499)rectangle++(0.008,0.008) (-0.006,0.207)rectangle++(0.008,0.008)
  (0.071,0.513)rectangle++(0.008,0.008) (0.215,0.626)rectangle++(0.008,0.008) (0.321,0.626)rectangle++(0.008,0.008) (0.369,0.588)rectangle++(0.008,0.008) (-0.045,0.82)rectangle++(0.008,0.008)
  (-0.005,0.342)rectangle++(0.008,0.008) (-0.005,0.741)rectangle++(0.008,0.008) (-0.004,0.909)rectangle++(0.008,0.008) (-0.33,0.833)rectangle++(0.008,0.008) (0.146,0.915)rectangle++(0.008,0.008)
  (0.393,0.783)rectangle++(0.008,0.008) (0.008,0.326)rectangle++(0.008,0.008) (-0.117,0.571)rectangle++(0.008,0.008) (-0.261,0.639)rectangle++(0.008,0.008) (0.108,0.801)rectangle++(0.008,0.008)
  (-0.24,0.824)rectangle++(0.008,0.008) (-0.103,0.944)rectangle++(0.008,0.008) (0.289,0.885)rectangle++(0.008,0.008) (0.119,0.969)rectangle++(0.008,0.008) (-0.295,0.905)rectangle++(0.008,0.008)
  (-0.462,0.727)rectangle++(0.008,0.008) (0.048,0.914)rectangle++(0.008,0.008) (-0.308,0.854)rectangle++(0.008,0.008) (-0.449,0.7)rectangle++(0.008,0.008) (-0.017,0.292)rectangle++(0.008,0.008)
  (-0.116,0.546)rectangle++(0.008,0.008) (0.141,0.707)rectangle++(0.008,0.008) (-0.192,0.793)rectangle++(0.008,0.008) (-0.083,0.931)rectangle++(0.008,0.008) (-0.374,0.815)rectangle++(0.008,0.008)
  (0.119,0.922)rectangle++(0.008,0.008) (-0.278,0.884)rectangle++(0.008,0.008) (-0.446,0.724)rectangle++(0.008,0.008) (-0.466,0.591)rectangle++(0.008,0.008) (-0.198,0.846)rectangle++(0.008,0.008)
  (-0.396,0.735)rectangle++(0.008,0.008) (0.081,0.894)rectangle++(0.008,0.008) (-0.286,0.857)rectangle++(0.008,0.008) (-0.012,0.357)rectangle++(0.008,0.008) (-0.137,0.578)rectangle++(0.008,0.008)
  (-0.273,0.635)rectangle++(0.008,0.008) (-0.117,0.864)rectangle++(0.008,0.008) (-0.051,0.961)rectangle++(0.008,0.008) (0.319,0.874)rectangle++(0.008,0.008) (-0.17,0.932)rectangle++(0.008,0.008)
  (-0.074,0.989)rectangle++(0.008,0.008) (-0.033,1.013)rectangle++(0.008,0.008) (0.346,0.891)rectangle++(0.008,0.008) (-0.164,0.95)rectangle++(0.008,0.008) (0.264,0.909)rectangle++(0.008,0.008)
  (-0.208,0.929)rectangle++(0.008,0.008) (0.236,0.915)rectangle++(0.008,0.008) (0.097,0.982)rectangle++(0.008,0.008) (0.394,0.831)rectangle++(0.008,0.008) (0.008,0.347)rectangle++(0.008,0.008)
  (-0.124,0.58)rectangle++(0.008,0.008) (-0.267,0.64)rectangle++(0.008,0.008) (-0.012,0.267)rectangle++(0.008,0.008) (-0.004,0.11)rectangle++(0.008,0.008) (-0.045,0.468)rectangle++(0.008,0.008)
  (-0.005,0.194)rectangle++(0.008,0.008) (0.066,0.507)rectangle++(0.008,0.008) (0.211,0.624)rectangle++(0.008,0.008) (-0.13,0.78)rectangle++(0.008,0.008) (-0.057,0.925)rectangle++(0.008,0.008)
  (-0.36,0.822)rectangle++(0.008,0.008) (-0.462,0.666)rectangle++(0.008,0.008) (-0.196,0.878)rectangle++(0.008,0.008) (0.223,0.888)rectangle++(0.008,0.008) (-0.219,0.904)rectangle++(0.008,0.008)
  (-0.094,0.978)rectangle++(0.008,0.008) (0.305,0.897)rectangle++(0.008,0.008) (0.126,0.974)rectangle++(0.008,0.008) (0.051,1.007)rectangle++(0.008,0.008) (0.019,1.021)rectangle++(0.008,0.008)
  (-0.359,0.892)rectangle++(0.008,0.008) (-0.487,0.699)rectangle++(0.008,0.008) (-0.018,0.291)rectangle++(0.008,0.008) (0.095,0.556)rectangle++(0.008,0.008) (0.241,0.637)rectangle++(0.008,0.008)
  (0.099,0.865)rectangle++(0.008,0.008) (0.353,0.776)rectangle++(0.008,0.008) (0.438,0.645)rectangle++(0.008,0.008) (-0.034,0.871)rectangle++(0.008,0.008) (-0.016,0.963)rectangle++(0.008,0.008)
  (-0.355,0.854)rectangle++(0.008,0.008) (-0.471,0.683)rectangle++(0.008,0.008) (-0.018,0.284)rectangle++(0.008,0.008) (-0.01,0.717)rectangle++(0.008,0.008) (0.251,0.748)rectangle++(0.008,0.008)
  (0.103,0.912)rectangle++(0.008,0.008) (-0.282,0.873)rectangle++(0.008,0.008) (0.182,0.916)rectangle++(0.008,0.008) (-0.248,0.903)rectangle++(0.008,0.008) (-0.106,0.977)rectangle++(0.008,0.008)
  (-0.047,1.008)rectangle++(0.008,0.008) (-0.385,0.863)rectangle++(0.008,0.008) (-0.164,0.96)rectangle++(0.008,0.008) (0.267,0.914)rectangle++(0.008,0.008) (-0.208,0.932)rectangle++(0.008,0.008)
  (-0.431,0.771)rectangle++(0.008,0.008) (-0.184,0.922)rectangle++(0.008,0.008) (-0.079,0.985)rectangle++(0.008,0.008) (-0.036,1.011)rectangle++(0.008,0.008) (0.344,0.891)rectangle++(0.008,0.008)
  (0.142,0.972)rectangle++(0.008,0.008) (-0.286,0.914)rectangle++(0.008,0.008) (-0.461,0.735)rectangle++(0.008,0.008) (-0.477,0.591)rectangle++(0.008,0.008) (-0.203,0.846)rectangle++(0.008,0.008)
  (-0.398,0.734)rectangle++(0.008,0.008) (0.079,0.894)rectangle++(0.008,0.008) (0.355,0.797)rectangle++(0.008,0.008) (0.147,0.932)rectangle++(0.008,0.008) (0.059,0.989)rectangle++(0.008,0.008)
  (0.023,1.013)rectangle++(0.008,0.008) (-0.355,0.89)rectangle++(0.008,0.008) (0.155,0.95)rectangle++(0.008,0.008) (0.409,0.795)rectangle++(0.008,0.008) (0.17,0.932)rectangle++(0.008,0.008)
  (-0.259,0.906)rectangle++(0.008,0.008) (-0.445,0.741)rectangle++(0.008,0.008) (-0.472,0.599)rectangle++(0.008,0.008) (-0.433,0.525)rectangle++(0.008,0.008) (-0.389,0.504)rectangle++(0.008,0.008)
  (0.002,0.786)rectangle++(0.008,0.008) (0.281,0.775)rectangle++(0.008,0.008) (-0.152,0.874)rectangle++(0.008,0.008) (0.242,0.87)rectangle++(0.008,0.008) (-0.204,0.903)rectangle++(0.008,0.008)
  (0.229,0.902)rectangle++(0.008,0.008) (0.003,0.377)rectangle++(0.008,0.008) (-0.137,0.592)rectangle++(0.008,0.008) (-0.278,0.641)rectangle++(0.008,0.008) (-0.119,0.867)rectangle++(0.008,0.008)
  (-0.368,0.773)rectangle++(0.008,0.008) (-0.015,0.322)rectangle++(0.008,0.008) (-0.009,0.733)rectangle++(0.008,0.008) (0.257,0.755)rectangle++(0.008,0.008) (0.386,0.669)rectangle++(0.008,0.008)
  (-0.066,0.863)rectangle++(0.008,0.008) (-0.342,0.79)rectangle++(0.008,0.008) (-0.442,0.658)rectangle++(0.008,0.008) (0.032,0.875)rectangle++(0.008,0.008) (0.011,0.965)rectangle++(0.008,0.008)
  (-0.004,0.403)rectangle++(0.008,0.008) (-0.004,0.767)rectangle++(0.008,0.008) (-0.004,0.92)rectangle++(0.008,0.008) (-0.004,0.984)rectangle++(0.008,0.008) (-0.357,0.868)rectangle++(0.008,0.008)
  (-0.015,0.362)rectangle++(0.008,0.008) (-0.008,0.75)rectangle++(0.008,0.008) (-0.275,0.759)rectangle++(0.008,0.008) (-0.118,0.916)rectangle++(0.008,0.008) (-0.007,0.383)rectangle++(0.008,0.008)
  (-0.144,0.591)rectangle++(0.008,0.008) (-0.008,0.246)rectangle++(0.008,0.008) (-0.006,0.701)rectangle++(0.008,0.008) (0.247,0.739)rectangle++(0.008,0.008) (0.376,0.666)rectangle++(0.008,0.008)
  (0.409,0.587)rectangle++(0.008,0.008) (0.396,0.538)rectangle++(0.008,0.008) (-0.015,0.807)rectangle++(0.008,0.008) (0.28,0.79)rectangle++(0.008,0.008) (-0.158,0.881)rectangle++(0.008,0.008)
  (-0.009,0.368)rectangle++(0.008,0.008) (-0.006,0.752)rectangle++(0.008,0.008) (-0.005,0.914)rectangle++(0.008,0.008) (-0.004,0.381)rectangle++(0.008,0.008) (-0.004,0.758)rectangle++(0.008,0.008)
  (-0.276,0.764)rectangle++(0.008,0.008) (-0.403,0.67)rectangle++(0.008,0.008) (-0.427,0.582)rectangle++(0.008,0.008) (-0.406,0.533)rectangle++(0.008,0.008) (0.004,0.805)rectangle++(0.008,0.008)
  (0.288,0.783)rectangle++(0.008,0.008) (0.411,0.671)rectangle++(0.008,0.008) (0.427,0.577)rectangle++(0.008,0.008) (0.4,0.528)rectangle++(0.008,0.008) (0.371,0.515)rectangle++(0.008,0.008)
  (0.007,0.214)rectangle++(0.008,0.008) (-0.004,0.087)rectangle++(0.008,0.008) (0.029,0.458)rectangle++(0.008,0.008) (-0.003,0.19)rectangle++(0.008,0.008) (-0.004,0.677)rectangle++(0.008,0.008)
  (-0.247,0.728)rectangle++(0.008,0.008) (0.146,0.837)rectangle++(0.008,0.008) (0.059,0.949)rectangle++(0.008,0.008) (0.022,0.996)rectangle++(0.008,0.008) (0.007,1.016)rectangle++(0.008,0.008)
  (0.365,0.878)rectangle++(0.008,0.008) (-0.15,0.951)rectangle++(0.008,0.008) (-0.008,0.397)rectangle++(0.008,0.008) (-0.149,0.598)rectangle++(0.008,0.008) (-0.065,0.849)rectangle++(0.008,0.008)
  (-0.03,0.954)rectangle++(0.008,0.008) (0.326,0.863)rectangle++(0.008,0.008) (-0.163,0.93)rectangle++(0.008,0.008) (0.257,0.9)rectangle++(0.008,0.008) (0.106,0.976)rectangle++(0.008,0.008)
  (-0.001,0.407)rectangle++(0.008,0.008) (-0.003,0.769)rectangle++(0.008,0.008) (-0.279,0.77)rectangle++(0.008,0.008) (-0.406,0.671)rectangle++(0.008,0.008) (-0.016,0.28)rectangle++(0.008,0.008);
\fill[nuage]
  (0.092,0.55)rectangle++(0.008,0.008) (0.238,0.635)rectangle++(0.008,0.008) (0.335,0.622)rectangle++(0.008,0.008) (0.374,0.581)rectangle++(0.008,0.008) (0.378,0.548)rectangle++(0.008,0.008)
  (0.156,0.828)rectangle++(0.008,0.008) (-0.228,0.854)rectangle++(0.008,0.008) (0.2,0.888)rectangle++(0.008,0.008) (-0.229,0.896)rectangle++(0.008,0.008) (-0.011,0.374)rectangle++(0.008,0.008)
  (0.128,0.591)rectangle++(0.008,0.008) (-0.156,0.735)rectangle++(0.008,0.008) (-0.338,0.699)rectangle++(0.008,0.008) (-0.144,0.891)rectangle++(0.008,0.008) (0.252,0.875)rectangle++(0.008,0.008)
  (0.103,0.965)rectangle++(0.008,0.008) (-0.001,0.403)rectangle++(0.008,0.008) (-0.004,0.167)rectangle++(0.008,0.008) (-0.065,0.494)rectangle++(0.008,0.008) (-0.21,0.622)rectangle++(0.008,0.008)
  (0.125,0.776)rectangle++(0.008,0.008) (-0.223,0.819)rectangle++(0.008,0.008) (-0.398,0.714)rectangle++(0.008,0.008) (0.072,0.885)rectangle++(0.008,0.008) (0.348,0.795)rectangle++(0.008,0.008)
  (-0.128,0.907)rectangle++(0.008,0.008) (-0.386,0.788)rectangle++(0.008,0.008) (-0.164,0.929)rectangle++(0.008,0.008) (-0.009,0.388)rectangle++(0.008,0.008) (-0.146,0.593)rectangle++(0.008,0.008)
  (-0.064,0.847)rectangle++(0.008,0.008) (-0.029,0.953)rectangle++(0.008,0.008) (-0.015,0.998)rectangle++(0.008,0.008) (-0.008,1.017)rectangle++(0.008,0.008) (-0.371,0.881)rectangle++(0.008,0.008)
  (0.144,0.951)rectangle++(0.008,0.008) (0.058,0.997)rectangle++(0.008,0.008) (-0.333,0.896)rectangle++(0.008,0.008) (0.167,0.945)rectangle++(0.008,0.008) (0.413,0.789)rectangle++(0.008,0.008)
  (0.47,0.629)rectangle++(0.008,0.008) (-0.014,0.875)rectangle++(0.008,0.008) (0.305,0.821)rectangle++(0.008,0.008) (0.126,0.943)rectangle++(0.008,0.008) (0,0.394)rectangle++(0.008,0.008)
  (-0.144,0.599)rectangle++(0.008,0.008) (0.147,0.742)rectangle++(0.008,0.008) (-0.201,0.811)rectangle++(0.008,0.008) (-0.087,0.938)rectangle++(0.008,0.008) (-0.378,0.817)rectangle++(0.008,0.008)
  (0.118,0.925)rectangle++(0.008,0.008) (0.384,0.797)rectangle++(0.008,0.008) (-0.113,0.92)rectangle++(0.008,0.008) (0.276,0.877)rectangle++(0.008,0.008) (-0.191,0.919)rectangle++(0.008,0.008)
  (0.24,0.904)rectangle++(0.008,0.008) (0.099,0.978)rectangle++(0.008,0.008) (0.393,0.828)rectangle++(0.008,0.008) (0.475,0.654)rectangle++(0.008,0.008) (-0.02,0.888)rectangle++(0.008,0.008)
  (-0.33,0.818)rectangle++(0.008,0.008) (0.14,0.908)rectangle++(0.008,0.008) (0.388,0.781)rectangle++(0.008,0.008) (0.455,0.635)rectangle++(0.008,0.008) (0.01,0.264)rectangle++(0.008,0.008)
  (-0.094,0.544)rectangle++(0.008,0.008) (-0.007,0.226)rectangle++(0.008,0.008) (0.077,0.522)rectangle++(0.008,0.008) (-0.002,0.217)rectangle++(0.008,0.008) (0.076,0.516)rectangle++(0.008,0.008)
  (0.218,0.625)rectangle++(0.008,0.008) (0.089,0.86)rectangle++(0.008,0.008) (-0.27,0.844)rectangle++(0.008,0.008) (-0.428,0.709)rectangle++(0.008,0.008) (0.057,0.893)rectangle++(0.008,0.008)
  (0.021,0.973)rectangle++(0.008,0.008) (0.007,1.006)rectangle++(0.008,0.008) (-0.36,0.882)rectangle++(0.008,0.008) (0.15,0.948)rectangle++(0.008,0.008) (0.06,0.996)rectangle++(0.008,0.008)
  (-0.332,0.896)rectangle++(0.008,0.008) (-0.142,0.974)rectangle++(0.008,0.008) (-0.008,0.407)rectangle++(0.008,0.008) (-0.006,0.769)rectangle++(0.008,0.008) (-0.281,0.768)rectangle++(0.008,0.008)
  (-0.406,0.67)rectangle++(0.008,0.008) (-0.173,0.879)rectangle++(0.008,0.008) (-0.397,0.759)rectangle++(0.008,0.008) (0.089,0.905)rectangle++(0.008,0.008) (-0.286,0.865)rectangle++(0.008,0.008)
  (-0.443,0.712)rectangle++(0.008,0.008) (-0.188,0.897)rectangle++(0.008,0.008) (-0.081,0.974)rectangle++(0.008,0.008) (0.31,0.891)rectangle++(0.008,0.008) (0.128,0.972)rectangle++(0.008,0.008)
  (0.405,0.815)rectangle++(0.008,0.008) (-0.11,0.936)rectangle++(0.008,0.008) (-0.048,0.991)rectangle++(0.008,0.008) (-0.005,0.414)rectangle++(0.008,0.008) (-0.154,0.607)rectangle++(0.008,0.008)
  (0.146,0.749)rectangle++(0.008,0.008) (-0.204,0.813)rectangle++(0.008,0.008) (-0.387,0.718)rectangle++(0.008,0.008) (-0.165,0.899)rectangle++(0.008,0.008) (-0.009,0.375)rectangle++(0.008,0.008)
  (0.129,0.591)rectangle++(0.008,0.008) (-0.156,0.736)rectangle++(0.008,0.008) (-0.009,0.307)rectangle++(0.008,0.008) (0.105,0.56)rectangle++(0.008,0.008) (0.042,0.833)rectangle++(0.008,0.008)
  (-0.282,0.815)rectangle++(0.008,0.008) (0.161,0.889)rectangle++(0.008,0.008) (0.001,0.371)rectangle++(0.008,0.008) (-0.136,0.589)rectangle++(0.008,0.008) (-0.059,0.845)rectangle++(0.008,0.008)
  (-0.333,0.784)rectangle++(0.008,0.008) (0.127,0.894)rectangle++(0.008,0.008) (-0.264,0.873)rectangle++(0.008,0.008) (0.19,0.91)rectangle++(0.008,0.008) (0.411,0.764)rectangle++(0.008,0.008)
  (-0.089,0.915)rectangle++(0.008,0.008) (-0.371,0.806)rectangle++(0.008,0.008) (-0.015,0.336)rectangle++(0.008,0.008) (-0.009,0.739)rectangle++(0.008,0.008) (-0.271,0.754)rectangle++(0.008,0.008)
  (-0.397,0.667)rectangle++(0.008,0.008) (0.056,0.863)rectangle++(0.008,0.008) (-0.286,0.834)rectangle++(0.008,0.008) (0.166,0.9)rectangle++(0.008,0.008) (0.396,0.768)rectangle++(0.008,0.008)
  (0.164,0.92)rectangle++(0.008,0.008) (0.067,0.984)rectangle++(0.008,0.008) (-0.325,0.893)rectangle++(0.008,0.008) (-0.014,0.373)rectangle++(0.008,0.008) (0.126,0.592)rectangle++(0.008,0.008)
  (0.051,0.846)rectangle++(0.008,0.008) (0.325,0.785)rectangle++(0.008,0.008) (0.428,0.659)rectangle++(0.008,0.008) (-0.043,0.873)rectangle++(0.008,0.008) (0.291,0.831)rectangle++(0.008,0.008)
  (0.12,0.947)rectangle++(0.008,0.008) (-0.287,0.895)rectangle++(0.008,0.008) (0.188,0.928)rectangle++(0.008,0.008) (0.002,0.387)rectangle++(0.008,0.008) (-0.141,0.597)rectangle++(0.008,0.008)
  (0.148,0.74)rectangle++(0.008,0.008) (0.001,0.308)rectangle++(0.008,0.008) (-0.002,0.727)rectangle++(0.008,0.008) (0.258,0.749)rectangle++(0.008,0.008) (0.385,0.667)rectangle++(0.008,0.008)
  (0.159,0.878)rectangle++(0.008,0.008) (-0.244,0.877)rectangle++(0.008,0.008) (0.201,0.905)rectangle++(0.008,0.008) (-0.235,0.904)rectangle++(0.008,0.008) (-0.434,0.749)rectangle++(0.008,0.008)
  (-0.469,0.607)rectangle++(0.008,0.008) (-0.435,0.529)rectangle++(0.008,0.008) (-0.017,0.22)rectangle++(0.008,0.008) (-0.004,0.09)rectangle++(0.008,0.008) (-0.004,0.035)rectangle++(0.008,0.008)
  (0.01,0.434)rectangle++(0.008,0.008) (0.159,0.611)rectangle++(0.008,0.008) (0.29,0.639)rectangle++(0.008,0.008) (-0.099,0.815)rectangle++(0.008,0.008) (-0.044,0.94)rectangle++(0.008,0.008)
  (0.315,0.862)rectangle++(0.008,0.008) (0.451,0.698)rectangle++(0.008,0.008) (-0.047,0.899)rectangle++(0.008,0.008) (0.299,0.844)rectangle++(0.008,0.008) (0.437,0.695)rectangle++(0.008,0.008)
  (0.447,0.578)rectangle++(0.008,0.008) (0.186,0.84)rectangle++(0.008,0.008) (-0.219,0.87)rectangle++(0.008,0.008) (0.21,0.892)rectangle++(0.008,0.008) (0.414,0.749)rectangle++(0.008,0.008)
  (0.456,0.611)rectangle++(0.008,0.008) (0.426,0.533)rectangle++(0.008,0.008) (0.384,0.508)rectangle++(0.008,0.008) (-0.009,0.789)rectangle++(0.008,0.008) (-0.006,0.929)rectangle++(0.008,0.008)
  (0.328,0.843)rectangle++(0.008,0.008) (0.45,0.685)rectangle++(0.008,0.008) (-0.042,0.893)rectangle++(0.008,0.008) (-0.02,0.973)rectangle++(0.008,0.008) (0.337,0.868)rectangle++(0.008,0.008)
  (0.139,0.962)rectangle++(0.008,0.008) (0.056,1.002)rectangle++(0.008,0.008) (0.383,0.854)rectangle++(0.008,0.008) (0.479,0.67)rectangle++(0.008,0.008) (0.01,0.279)rectangle++(0.008,0.008)
  (0.002,0.715)rectangle++(0.008,0.008) (0.255,0.742)rectangle++(0.008,0.008) (-0.152,0.85)rectangle++(0.008,0.008) (0.233,0.859)rectangle++(0.008,0.008) (0.096,0.958)rectangle++(0.008,0.008)
  (-0.302,0.891)rectangle++(0.008,0.008) (0.179,0.932)rectangle++(0.008,0.008) (0.073,0.989)rectangle++(0.008,0.008) (-0.323,0.897)rectangle++(0.008,0.008) (-0.138,0.975)rectangle++(0.008,0.008)
  (-0.06,1.007)rectangle++(0.008,0.008) (0.331,0.898)rectangle++(0.008,0.008) (-0.173,0.948)rectangle++(0.008,0.008) (-0.421,0.791)rectangle++(0.008,0.008) (-0.478,0.63)rectangle++(0.008,0.008)
  (-0.203,0.862)rectangle++(0.008,0.008) (-0.088,0.96)rectangle++(0.008,0.008) (-0.007,0.401)rectangle++(0.008,0.008) (0.139,0.602)rectangle++(0.008,0.008) (0.056,0.85)rectangle++(0.008,0.008)
  (-0.002,0.355)rectangle++(0.008,0.008) (0.125,0.579)rectangle++(0.008,0.008) (-0.153,0.729)rectangle++(0.008,0.008) (0.189,0.804)rectangle++(0.008,0.008) (-0.204,0.854)rectangle++(0.008,0.008)
  (-0.088,0.957)rectangle++(0.008,0.008) (-0.039,0.999)rectangle++(0.008,0.008) (-0.019,1.017)rectangle++(0.008,0.008) (-0.376,0.878)rectangle++(0.008,0.008) (0.141,0.952)rectangle++(0.008,0.008)
  (0.057,0.997)rectangle++(0.008,0.008) (0.022,1.017)rectangle++(0.008,0.008) (0.007,1.025)rectangle++(0.008,0.008) (-0.366,0.89)rectangle++(0.008,0.008) (-0.489,0.695)rectangle++(0.008,0.008)
  (-0.475,0.562)rectangle++(0.008,0.008) (-0.202,0.834)rectangle++(0.008,0.008) (-0.394,0.728)rectangle++(0.008,0.008) (-0.444,0.612)rectangle++(0.008,0.008) (-0.189,0.855)rectangle++(0.008,0.008)
  (-0.395,0.742)rectangle++(0.008,0.008) (-0.449,0.618)rectangle++(0.008,0.008) (-0.429,0.541)rectangle++(0.008,0.008) (-0.004,0.817)rectangle++(0.008,0.008) (-0.297,0.791)rectangle++(0.008,0.008)
  (0.146,0.884)rectangle++(0.008,0.008) (0.059,0.969)rectangle++(0.008,0.008) (0.372,0.838)rectangle++(0.008,0.008) (0.154,0.95)rectangle++(0.008,0.008) (-0.272,0.908)rectangle++(0.008,0.008)
  (-0.117,0.979)rectangle++(0.008,0.008) (-0.007,0.409)rectangle++(0.008,0.008) (0.142,0.606)rectangle++(0.008,0.008) (0.057,0.852)rectangle++(0.008,0.008) (-0.002,0.356)rectangle++(0.008,0.008)
  (-0.132,0.581)rectangle++(0.008,0.008) (0.147,0.729)rectangle++(0.008,0.008) (0.059,0.904)rectangle++(0.008,0.008) (-0.299,0.854)rectangle++(0.008,0.008) (0.167,0.913)rectangle++(0.008,0.008)
  (0.402,0.774)rectangle++(0.008,0.008) (0.459,0.627)rectangle++(0.008,0.008) (-0.017,0.87)rectangle++(0.008,0.008) (0.302,0.82)rectangle++(0.008,0.008) (-0.158,0.902)rectangle++(0.008,0.008)
  (-0.398,0.775)rectangle++(0.008,0.008) (-0.17,0.923)rectangle++(0.008,0.008) (0.251,0.899)rectangle++(0.008,0.008) (0.435,0.737)rectangle++(0.008,0.008) (-0.068,0.912)rectangle++(0.008,0.008)
  (-0.031,0.981)rectangle++(0.008,0.008) (0.335,0.876)rectangle++(0.008,0.008) (-0.163,0.939)rectangle++(0.008,0.008) (-0.009,0.392)rectangle++(0.008,0.008) (-0.006,0.762)rectangle++(0.008,0.008)
  (-0.279,0.766)rectangle++(0.008,0.008) (-0.119,0.919)rectangle++(0.008,0.008) (-0.386,0.797)rectangle++(0.008,0.008) (0.107,0.919)rectangle++(0.008,0.008) (0.043,0.983)rectangle++(0.008,0.008)
  (0.016,1.011)rectangle++(0.008,0.008) (0.004,1.022)rectangle++(0.008,0.008) (0.366,0.882)rectangle++(0.008,0.008) (0.482,0.689)rectangle++(0.008,0.008) (0.2,0.887)rectangle++(0.008,0.008)
  (0.082,0.97)rectangle++(0.008,0.008) (-0.313,0.892)rectangle++(0.008,0.008) (-0.134,0.972)rectangle++(0.008,0.008) (-0.058,1.006)rectangle++(0.008,0.008) (0.332,0.897)rectangle++(0.008,0.008)
  (0.006,0.374)rectangle++(0.008,0.008) (0,0.755)rectangle++(0.008,0.008) (-0.004,0.315)rectangle++(0.008,0.008) (0.11,0.562)rectangle++(0.008,0.008) (0.25,0.634)rectangle++(0.008,0.008)
  (0.34,0.617)rectangle++(0.008,0.008) (-0.069,0.823)rectangle++(0.008,0.008) (-0.031,0.943)rectangle++(0.008,0.008) (0.322,0.859)rectangle++(0.008,0.008) (0.133,0.958)rectangle++(0.008,0.008)
  (0.053,1)rectangle++(0.008,0.008) (-0.002,0.418)rectangle++(0.008,0.008) (-0.004,0.173)rectangle++(0.008,0.008) (-0.067,0.497)rectangle++(0.008,0.008) (0.146,0.668)rectangle++(0.008,0.008)
  (-0.175,0.776)rectangle++(0.008,0.008) (0.196,0.834)rectangle++(0.008,0.008) (0.387,0.727)rectangle++(0.008,0.008) (0.436,0.611)rectangle++(0.008,0.008) (0.417,0.54)rectangle++(0.008,0.008)
  (-0.006,0.815)rectangle++(0.008,0.008) (-0.297,0.789)rectangle++(0.008,0.008) (-0.421,0.674)rectangle++(0.008,0.008) (-0.437,0.577)rectangle++(0.008,0.008) (-0.409,0.527)rectangle++(0.008,0.008)
  (-0.379,0.514)rectangle++(0.008,0.008) (-0.36,0.519)rectangle++(0.008,0.008) (-0.154,0.816)rectangle++(0.008,0.008) (-0.067,0.94)rectangle++(0.008,0.008) (-0.03,0.993)rectangle++(0.008,0.008)
  (0.34,0.881)rectangle++(0.008,0.008) (-0.163,0.943)rectangle++(0.008,0.008) (-0.415,0.792)rectangle++(0.008,0.008) (-0.476,0.633)rectangle++(0.008,0.008) (-0.447,0.539)rectangle++(0.008,0.008)
  (-0.013,0.822)rectangle++(0.008,0.008) (-0.303,0.791)rectangle++(0.008,0.008) (-0.424,0.672)rectangle++(0.008,0.008) (0.045,0.875)rectangle++(0.008,0.008) (0.333,0.8)rectangle++(0.008,0.008)
  (0.437,0.663)rectangle++(0.008,0.008) (0.181,0.876)rectangle++(0.008,0.008) (0.002,0.366)rectangle++(0.008,0.008) (0.131,0.583)rectangle++(0.008,0.008) (-0.152,0.732)rectangle++(0.008,0.008)
  (-0.066,0.905)rectangle++(0.008,0.008) (0.292,0.854)rectangle++(0.008,0.008) (-0.175,0.914)rectangle++(0.008,0.008) (-0.41,0.774)rectangle++(0.008,0.008) (-0.467,0.627)rectangle++(0.008,0.008)
  (-0.441,0.539)rectangle++(0.008,0.008) (-0.398,0.508)rectangle++(0.008,0.008) (-0.169,0.811)rectangle++(0.008,0.008) (-0.371,0.729)rectangle++(0.008,0.008) (-0.433,0.62)rectangle++(0.008,0.008)
  (0.023,0.855)rectangle++(0.008,0.008) (0.007,0.957)rectangle++(0.008,0.008) (0.344,0.851)rectangle++(0.008,0.008) (0.006,0.355)rectangle++(0.008,0.008) (0,0.747)rectangle++(0.008,0.008)
  (-0.002,0.911)rectangle++(0.008,0.008) (-0.003,0.98)rectangle++(0.008,0.008) (-0.004,1.009)rectangle++(0.008,0.008) (-0.366,0.879)rectangle++(0.008,0.008) (-0.485,0.691)rectangle++(0.008,0.008)
  (-0.018,0.288)rectangle++(0.008,0.008) (-0.115,0.544)rectangle++(0.008,0.008) (-0.25,0.627)rectangle++(0.008,0.008) (0.109,0.792)rectangle++(0.008,0.008) (-0.237,0.82)rectangle++(0.008,0.008)
  (-0.405,0.71)rectangle++(0.008,0.008) (-0.442,0.599)rectangle++(0.008,0.008) (-0.017,0.249)rectangle++(0.008,0.008) (0.08,0.536)rectangle++(0.008,0.008) (0.031,0.823)rectangle++(0.008,0.008)
  (0.011,0.943)rectangle++(0.008,0.008) (0.002,0.994)rectangle++(0.008,0.008) (-0.001,1.015)rectangle++(0.008,0.008) (-0.003,1.024)rectangle++(0.008,0.008) (-0.004,1.028)rectangle++(0.008,0.008)
  (-0.372,0.888)rectangle++(0.008,0.008) (-0.491,0.692)rectangle++(0.008,0.008) (-0.475,0.56)rectangle++(0.008,0.008) (-0.018,0.842)rectangle++(0.008,0.008) (0.292,0.808)rectangle++(0.008,0.008)
  (-0.159,0.893)rectangle++(0.008,0.008) (-0.395,0.771)rectangle++(0.008,0.008) (0.094,0.91)rectangle++(0.008,0.008) (0.367,0.799)rectangle++(0.008,0.008) (-0.121,0.915)rectangle++(0.008,0.008)
  (-0.386,0.794)rectangle++(0.008,0.008) (0.107,0.917)rectangle++(0.008,0.008) (0.042,0.983)rectangle++(0.008,0.008) (0.016,1.011)rectangle++(0.008,0.008) (0.004,1.022)rectangle++(0.008,0.008)
  (0.366,0.882)rectangle++(0.008,0.008) (0.007,0.368)rectangle++(0.008,0.008) (0.134,0.582)rectangle++(0.008,0.008) (0.054,0.842)rectangle++(0.008,0.008) (0.325,0.782)rectangle++(0.008,0.008)
  (0.134,0.926)rectangle++(0.008,0.008) (-0.273,0.89)rectangle++(0.008,0.008) (-0.012,0.372)rectangle++(0.008,0.008) (-0.142,0.585)rectangle++(0.008,0.008) (0.143,0.734)rectangle++(0.008,0.008)
  (0.058,0.906)rectangle++(0.008,0.008) (-0.301,0.854)rectangle++(0.008,0.008) (0.167,0.914)rectangle++(0.008,0.008) (0.402,0.775)rectangle++(0.008,0.008) (0.46,0.627)rectangle++(0.008,0.008)
  (0.191,0.861)rectangle++(0.008,0.008) (-0.224,0.881)rectangle++(0.008,0.008) (-0.096,0.968)rectangle++(0.008,0.008) (-0.393,0.827)rectangle++(0.008,0.008) (0.115,0.935)rectangle++(0.008,0.008)
  (-0.285,0.888)rectangle++(0.008,0.008) (0.186,0.924)rectangle++(0.008,0.008) (-0.248,0.908)rectangle++(0.008,0.008) (0.21,0.92)rectangle++(0.008,0.008) (-0.236,0.915)rectangle++(0.008,0.008)
  (-0.102,0.982)rectangle++(0.008,0.008) (0.303,0.901)rectangle++(0.008,0.008) (0.125,0.976)rectangle++(0.008,0.008) (0.405,0.818)rectangle++(0.008,0.008) (0.008,0.341)rectangle++(0.008,0.008)
  (-0.122,0.578)rectangle++(0.008,0.008) (-0.266,0.64)rectangle++(0.008,0.008) (-0.354,0.617)rectangle++(0.008,0.008) (-0.385,0.575)rectangle++(0.008,0.008) (0.028,0.817)rectangle++(0.008,0.008)
  (0.304,0.78)rectangle++(0.008,0.008) (-0.143,0.884)rectangle++(0.008,0.008) (0.25,0.872)rectangle++(0.008,0.008) (0.425,0.726)rectangle++(0.008,0.008) (0.452,0.596)rectangle++(0.008,0.008)
  (0.419,0.527)rectangle++(0.008,0.008) (0.009,0.219)rectangle++(0.008,0.008) (0.082,0.513)rectangle++(0.008,0.008) (0.032,0.813)rectangle++(0.008,0.008) (-0.003,0.339)rectangle++(0.008,0.008)
  (-0.126,0.573)rectangle++(0.008,0.008) (0.146,0.723)rectangle++(0.008,0.008) (-0.195,0.802)rectangle++(0.008,0.008) (0.197,0.853)rectangle++(0.008,0.008) (-0.218,0.879)rectangle++(0.008,0.008)
  (0.213,0.896)rectangle++(0.008,0.008) (0.417,0.75)rectangle++(0.008,0.008) (-0.081,0.911)rectangle++(0.008,0.008) (-0.036,0.98)rectangle++(0.008,0.008) (0.333,0.877)rectangle++(0.008,0.008)
  (-0.165,0.939)rectangle++(0.008,0.008) (0.259,0.904)rectangle++(0.008,0.008) (0.441,0.737)rectangle++(0.008,0.008) (0.464,0.596)rectangle++(0.008,0.008) (0.424,0.523)rectangle++(0.008,0.008)
  (0.38,0.504)rectangle++(0.008,0.008) (-0.01,0.785)rectangle++(0.008,0.008) (0.275,0.779)rectangle++(0.008,0.008) (0.113,0.925)rectangle++(0.008,0.008) (0,0.386)rectangle++(0.008,0.008)
  (0.137,0.593)rectangle++(0.008,0.008) (0.274,0.639)rectangle++(0.008,0.008) (0.352,0.611)rectangle++(0.008,0.008) (-0.061,0.824)rectangle++(0.008,0.008) (-0.028,0.944)rectangle++(0.008,0.008)
  (-0.014,0.994)rectangle++(0.008,0.008) (0.348,0.876)rectangle++(0.008,0.008) (-0.157,0.944)rectangle++(0.008,0.008) (-0.068,0.994)rectangle++(0.008,0.008) (-0.39,0.849)rectangle++(0.008,0.008)
  (-0.485,0.668)rectangle++(0.008,0.008) (0.016,0.895)rectangle++(0.008,0.008) (0.326,0.82)rectangle++(0.008,0.008) (0.441,0.675)rectangle++(0.008,0.008) (0.442,0.567)rectangle++(0.008,0.008)
  (0.404,0.518)rectangle++(0.008,0.008) (0.167,0.815)rectangle++(0.008,0.008) (-0.218,0.852)rectangle++(0.008,0.008) (0.204,0.884)rectangle++(0.008,0.008) (-0.226,0.896)rectangle++(0.008,0.008)
  (-0.097,0.974)rectangle++(0.008,0.008) (-0.043,1.007)rectangle++(0.008,0.008) (0.339,0.892)rectangle++(0.008,0.008) (-0.167,0.948)rectangle++(0.008,0.008) (-0.418,0.793)rectangle++(0.008,0.008)
  (-0.178,0.931)rectangle++(0.008,0.008) (-0.417,0.781)rectangle++(0.008,0.008) (-0.473,0.627)rectangle++(0.008,0.008) (-0.444,0.537)rectangle++(0.008,0.008) (-0.189,0.823)rectangle++(0.008,0.008)
  (-0.384,0.728)rectangle++(0.008,0.008) (0.084,0.886)rectangle++(0.008,0.008) (0.033,0.97)rectangle++(0.008,0.008) (-0.003,0.405)rectangle++(0.008,0.008) (0.143,0.603)rectangle++(0.008,0.008)
  (0.058,0.851)rectangle++(0.008,0.008) (-0.281,0.829)rectangle++(0.008,0.008) (-0.428,0.698)rectangle++(0.008,0.008) (0.053,0.888)rectangle++(0.008,0.008) (-0.297,0.844)rectangle++(0.008,0.008)
  (-0.441,0.699)rectangle++(0.008,0.008) (-0.455,0.581)rectangle++(0.008,0.008) (-0.193,0.842)rectangle++(0.008,0.008) (-0.393,0.735)rectangle++(0.008,0.008) (0.082,0.893)rectangle++(0.008,0.008)
  (0.356,0.795)rectangle++(0.008,0.008) (-0.125,0.91)rectangle++(0.008,0.008) (-0.008,0.38)rectangle++(0.008,0.008) (-0.143,0.59)rectangle++(0.008,0.008) (-0.062,0.846)rectangle++(0.008,0.008)
  (0.273,0.825)rectangle++(0.008,0.008) (-0.174,0.894)rectangle++(0.008,0.008) (0.239,0.887)rectangle++(0.008,0.008) (0.425,0.736)rectangle++(0.008,0.008) (0.176,0.907)rectangle++(0.008,0.008)
  (-0.247,0.897)rectangle++(0.008,0.008) (-0.437,0.741)rectangle++(0.008,0.008) (0.064,0.911)rectangle++(0.008,0.008) (0.025,0.98)rectangle++(0.008,0.008) (-0.003,0.409)rectangle++(0.008,0.008)
  (-0.004,0.77)rectangle++(0.008,0.008) (-0.004,0.321)rectangle++(0.008,0.008) (-0.004,0.732)rectangle++(0.008,0.008) (-0.267,0.753)rectangle++(0.008,0.008) (-0.114,0.914)rectangle++(0.008,0.008)
  (-0.382,0.796)rectangle++(0.008,0.008) (-0.463,0.647)rectangle++(0.008,0.008) (-0.197,0.869)rectangle++(0.008,0.008) (-0.01,0.363)rectangle++(0.008,0.008) (-0.004,0.15)rectangle++(0.008,0.008)
  (-0.004,0.061)rectangle++(0.008,0.008) (0.019,0.446)rectangle++(0.008,0.008) (0.006,0.785)rectangle++(0.008,0.008) (0.282,0.773)rectangle++(0.008,0.008) (0.116,0.922)rectangle++(0.008,0.008)
  (0.382,0.796)rectangle++(0.008,0.008) (0.158,0.932)rectangle++(0.008,0.008) (-0.264,0.902)rectangle++(0.008,0.008) (-0.012,0.376)rectangle++(0.008,0.008) (-0.143,0.587)rectangle++(0.008,0.008)
  (0.143,0.736)rectangle++(0.008,0.008) (0.328,0.702)rectangle++(0.008,0.008) (0.399,0.62)rectangle++(0.008,0.008) (0.403,0.557)rectangle++(0.008,0.008) (0.383,0.527)rectangle++(0.008,0.008)
  (-0.017,0.797)rectangle++(0.008,0.008) (0.276,0.787)rectangle++(0.008,0.008) (0.004,0.328)rectangle++(0.008,0.008) (0.118,0.565)rectangle++(0.008,0.008) (-0.151,0.72)rectangle++(0.008,0.008)
  (0.187,0.799)rectangle++(0.008,0.008) (0.37,0.715)rectangle++(0.008,0.008) (0.153,0.898)rectangle++(0.008,0.008) (-0.254,0.884)rectangle++(0.008,0.008) (-0.436,0.733)rectangle++(0.008,0.008)
  (0.062,0.907)rectangle++(0.008,0.008) (0.024,0.979)rectangle++(0.008,0.008) (-0.342,0.875)rectangle++(0.008,0.008) (-0.472,0.697)rectangle++(0.008,0.008) (-0.018,0.29)rectangle++(0.008,0.008)
  (-0.116,0.546)rectangle++(0.008,0.008) (-0.251,0.627)rectangle++(0.008,0.008) (-0.108,0.861)rectangle++(0.008,0.008) (-0.048,0.959)rectangle++(0.008,0.008) (0.32,0.872)rectangle++(0.008,0.008)
  (0.457,0.701)rectangle++(0.008,0.008) (0.19,0.892)rectangle++(0.008,0.008) (0.404,0.756)rectangle++(0.008,0.008) (0.168,0.915)rectangle++(0.008,0.008) (-0.254,0.897)rectangle++(0.008,0.008)
  (-0.44,0.739)rectangle++(0.008,0.008) (-0.187,0.908)rectangle++(0.008,0.008) (-0.413,0.767)rectangle++(0.008,0.008) (-0.176,0.92)rectangle++(0.008,0.008) (-0.413,0.777)rectangle++(0.008,0.008)
  (-0.176,0.924)rectangle++(0.008,0.008) (-0.414,0.779)rectangle++(0.008,0.008) (-0.471,0.627)rectangle++(0.008,0.008) (-0.443,0.538)rectangle++(0.008,0.008) (-0.188,0.824)rectangle++(0.008,0.008)
  (0.207,0.86)rectangle++(0.008,0.008) (0.401,0.735)rectangle++(0.008,0.008) (0.166,0.907)rectangle++(0.008,0.008) (0.001,0.378)rectangle++(0.008,0.008) (0.135,0.589)rectangle++(0.008,0.008)
  (-0.152,0.737)rectangle++(0.008,0.008) (0.193,0.807)rectangle++(0.008,0.008) (0.079,0.937)rectangle++(0.008,0.008) (-0.302,0.875)rectangle++(0.008,0.008) (-0.129,0.965)rectangle++(0.008,0.008)
  (0.285,0.904)rectangle++(0.008,0.008) (-0.196,0.934)rectangle++(0.008,0.008) (-0.427,0.776)rectangle++(0.008,0.008) (-0.182,0.924)rectangle++(0.008,0.008) (-0.416,0.777)rectangle++(0.008,0.008)
  (-0.016,0.324)rectangle++(0.008,0.008) (-0.127,0.561)rectangle++(0.008,0.008) (0.142,0.718)rectangle++(0.008,0.008) (-0.195,0.798)rectangle++(0.008,0.008) (-0.084,0.933)rectangle++(0.008,0.008)
  (-0.375,0.816)rectangle++(0.008,0.008) (-0.16,0.94)rectangle++(0.008,0.008) (-0.412,0.792)rectangle++(0.008,0.008) (-0.176,0.93)rectangle++(0.008,0.008) (-0.416,0.782)rectangle++(0.008,0.008)
  (-0.473,0.628)rectangle++(0.008,0.008) (-0.444,0.537)rectangle++(0.008,0.008) (-0.398,0.506)rectangle++(0.008,0.008) (-0.17,0.81)rectangle++(0.008,0.008) (-0.009,0.338)rectangle++(0.008,0.008)
  (0.116,0.574)rectangle++(0.008,0.008) (0.046,0.839)rectangle++(0.008,0.008) (0.017,0.95)rectangle++(0.008,0.008) (-0.335,0.86)rectangle++(0.008,0.008) (0.153,0.929)rectangle++(0.008,0.008)
  (0.401,0.786)rectangle++(0.008,0.008) (-0.101,0.922)rectangle++(0.008,0.008) (-0.045,0.985)rectangle++(0.008,0.008) (-0.376,0.853)rectangle++(0.008,0.008) (-0.16,0.956)rectangle++(0.008,0.008)
  (0.267,0.911)rectangle++(0.008,0.008) (0.11,0.98)rectangle++(0.008,0.008) (-0.001,0.409)rectangle++(0.008,0.008) (0.145,0.604)rectangle++(0.008,0.008) (0.059,0.851)rectangle++(0.008,0.008)
  (-0.002,0.355)rectangle++(0.008,0.008) (0.125,0.579)rectangle++(0.008,0.008) (0.263,0.637)rectangle++(0.008,0.008) (0.347,0.613)rectangle++(0.008,0.008) (0.377,0.573)rectangle++(0.008,0.008)
  (0.376,0.544)rectangle++(0.008,0.008) (0.365,0.531)rectangle++(0.008,0.008) (0.356,0.528)rectangle++(0.008,0.008) (0.351,0.531)rectangle++(0.008,0.008) (-0.033,0.787)rectangle++(0.008,0.008)
  (-0.3,0.767)rectangle++(0.008,0.008) (0.136,0.874)rectangle++(0.008,0.008) (-0.253,0.867)rectangle++(0.008,0.008) (-0.429,0.725)rectangle++(0.008,0.008) (-0.459,0.597)rectangle++(0.008,0.008)
  (-0.427,0.529)rectangle++(0.008,0.008) (-0.007,0.81)rectangle++(0.008,0.008) (0.285,0.789)rectangle++(0.008,0.008) (0.412,0.675)rectangle++(0.008,0.008) (0.171,0.881)rectangle++(0.008,0.008)
  (-0.24,0.883)rectangle++(0.008,0.008) (-0.429,0.737)rectangle++(0.008,0.008) (0.066,0.906)rectangle++(0.008,0.008) (-0.297,0.857)rectangle++(0.008,0.008) (-0.127,0.958)rectangle++(0.008,0.008)
  (-0.056,1)rectangle++(0.008,0.008) (0.331,0.893)rectangle++(0.008,0.008) (0.006,0.373)rectangle++(0.008,0.008) (-0.134,0.592)rectangle++(0.008,0.008) (-0.059,0.846)rectangle++(0.008,0.008)
  (-0.027,0.953)rectangle++(0.008,0.008) (-0.005,0.398)rectangle++(0.008,0.008) (0.139,0.6)rectangle++(0.008,0.008) (0.056,0.85)rectangle++(0.008,0.008) (-0.002,0.355)rectangle++(0.008,0.008)
  (-0.131,0.581)rectangle++(0.008,0.008) (0.147,0.729)rectangle++(0.008,0.008) (0.326,0.697)rectangle++(0.008,0.008) (-0.103,0.854)rectangle++(0.008,0.008) (0.257,0.844)rectangle++(0.008,0.008)
  (0.418,0.71)rectangle++(0.008,0.008) (0.444,0.592)rectangle++(0.008,0.008) (0.413,0.528)rectangle++(0.008,0.008) (0.171,0.82)rectangle++(0.008,0.008) (0.37,0.73)rectangle++(0.008,0.008)
  (0.153,0.904)rectangle++(0.008,0.008) (-0.256,0.887)rectangle++(0.008,0.008) (-0.11,0.97)rectangle++(0.008,0.008) (-0.4,0.823)rectangle++(0.008,0.008) (-0.481,0.653)rectangle++(0.008,0.008)
  (-0.204,0.872)rectangle++(0.008,0.008) (-0.408,0.745)rectangle++(0.008,0.008) (-0.456,0.614)rectangle++(0.008,0.008) (-0.431,0.537)rectangle++(0.008,0.008) (-0.017,0.223)rectangle++(0.008,0.008)
  (-0.004,0.091)rectangle++(0.008,0.008) (-0.004,0.636)rectangle++(0.008,0.008) (0.224,0.709)rectangle++(0.008,0.008) (-0.154,0.823)rectangle++(0.008,0.008) (0.223,0.848)rectangle++(0.008,0.008)
  (-0.204,0.886)rectangle++(0.008,0.008) (0.222,0.894)rectangle++(0.008,0.008) (-0.221,0.907)rectangle++(0.008,0.008) (-0.095,0.979)rectangle++(0.008,0.008) (-0.042,1.009)rectangle++(0.008,0.008)
  (-0.005,0.421)rectangle++(0.008,0.008) (-0.156,0.61)rectangle++(0.008,0.008) (-0.293,0.642)rectangle++(0.008,0.008) (0.095,0.815)rectangle++(0.008,0.008) (0.037,0.94)rectangle++(0.008,0.008)
  (0.352,0.833)rectangle++(0.008,0.008) (0.007,0.347)rectangle++(0.008,0.008) (0.126,0.573)rectangle++(0.008,0.008) (0.051,0.838)rectangle++(0.008,0.008) (0.322,0.781)rectangle++(0.008,0.008)
  (-0.136,0.891)rectangle++(0.008,0.008) (0.256,0.872)rectangle++(0.008,0.008) (0.105,0.964)rectangle++(0.008,0.008) (-0.3,0.897)rectangle++(0.008,0.008) (-0.013,0.375)rectangle++(0.008,0.008)
  (-0.143,0.586)rectangle++(0.008,0.008) (0.143,0.735)rectangle++(0.008,0.008) (-0.201,0.806)rectangle++(0.008,0.008) (-0.383,0.716)rectangle++(0.008,0.008) (-0.435,0.61)rectangle++(0.008,0.008)
  (-0.42,0.543)rectangle++(0.008,0.008) (0.001,0.814)rectangle++(0.008,0.008) (0.291,0.788)rectangle++(0.008,0.008) (0.005,0.329)rectangle++(0.008,0.008) (-0.119,0.571)rectangle++(0.008,0.008)
  (0.149,0.72)rectangle++(0.008,0.008) (0.06,0.9)rectangle++(0.008,0.008) (-0.297,0.852)rectangle++(0.008,0.008) (0.168,0.912)rectangle++(0.008,0.008) (-0.253,0.896)rectangle++(0.008,0.008)
  (-0.108,0.974)rectangle++(0.008,0.008) (-0.401,0.826)rectangle++(0.008,0.008) (-0.482,0.653)rectangle++(0.008,0.008) (-0.018,0.272)rectangle++(0.008,0.008) (0.088,0.547)rectangle++(0.008,0.008)
  (0.035,0.828)rectangle++(0.008,0.008) (0.311,0.782)rectangle++(0.008,0.008) (0.421,0.663)rectangle++(0.008,0.008) (0.428,0.569)rectangle++(0.008,0.008) (0.178,0.837)rectangle++(0.008,0.008)
  (0.379,0.735)rectangle++(0.008,0.008) (0.157,0.907)rectangle++(0.008,0.008) (-0.256,0.89)rectangle++(0.008,0.008) (-0.438,0.735)rectangle++(0.008,0.008) (-0.466,0.599)rectangle++(0.008,0.008)
  (-0.43,0.526)rectangle++(0.008,0.008) (-0.388,0.506)rectangle++(0.008,0.008) (-0.362,0.512)rectangle++(0.008,0.008) (-0.154,0.813)rectangle++(0.008,0.008) (-0.364,0.736)rectangle++(0.008,0.008)
  (-0.433,0.625)rectangle++(0.008,0.008) (-0.184,0.86)rectangle++(0.008,0.008) (-0.009,0.359)rectangle++(0.008,0.008) (0.123,0.584)rectangle++(0.008,0.008) (-0.156,0.73)rectangle++(0.008,0.008)
  (-0.068,0.904)rectangle++(0.008,0.008) (0.291,0.854)rectangle++(0.008,0.008) (0.437,0.703)rectangle++(0.008,0.008) (0.45,0.581)rectangle++(0.008,0.008) (0.187,0.842)rectangle++(0.008,0.008)
  (0.385,0.735)rectangle++(0.008,0.008) (0.159,0.906)rectangle++(0.008,0.008) (0.065,0.978)rectangle++(0.008,0.008) (0.378,0.84)rectangle++(0.008,0.008) (0.007,0.351)rectangle++(0.008,0.008)
  (0.001,0.745)rectangle++(0.008,0.008) (-0.002,0.911)rectangle++(0.008,0.008) (0.323,0.833)rectangle++(0.008,0.008) (-0.153,0.916)rectangle++(0.008,0.008) (-0.067,0.982)rectangle++(0.008,0.008)
  (0.32,0.889)rectangle++(0.008,0.008) (0.132,0.971)rectangle++(0.008,0.008) (0.406,0.813)rectangle++(0.008,0.008) (-0.108,0.936)rectangle++(0.008,0.008) (-0.387,0.808)rectangle++(0.008,0.008)
  (-0.469,0.65)rectangle++(0.008,0.008) (-0.018,0.271)rectangle++(0.008,0.008) (-0.004,0.111)rectangle++(0.008,0.008) (0.037,0.469)rectangle++(0.008,0.008) (-0.154,0.647)rectangle++(0.008,0.008)
  (-0.009,0.269)rectangle++(0.008,0.008) (-0.004,0.111)rectangle++(0.008,0.008) (-0.004,0.044)rectangle++(0.008,0.008) (0.013,0.438)rectangle++(0.008,0.008) (-0.154,0.624)rectangle++(0.008,0.008)
  (-0.067,0.86)rectangle++(0.008,0.008) (0.276,0.833)rectangle++(0.008,0.008) (0.113,0.948)rectangle++(0.008,0.008) (-0.29,0.893)rectangle++(0.008,0.008) (-0.124,0.973)rectangle++(0.008,0.008)
  (-0.054,1.006)rectangle++(0.008,0.008) (-0.388,0.86)rectangle++(0.008,0.008) (0.129,0.948)rectangle++(0.008,0.008) (-0.283,0.899)rectangle++(0.008,0.008) (0.191,0.928)rectangle++(0.008,0.008)
  (-0.248,0.912)rectangle++(0.008,0.008) (-0.442,0.747)rectangle++(0.008,0.008) (-0.473,0.603)rectangle++(0.008,0.008) (-0.001,0.861)rectangle++(0.008,0.008) (-0.312,0.812)rectangle++(0.008,0.008)
  (-0.436,0.679)rectangle++(0.008,0.008) (0.042,0.882)rectangle++(0.008,0.008) (-0.299,0.838)rectangle++(0.008,0.008) (0.162,0.906)rectangle++(0.008,0.008) (0.066,0.978)rectangle++(0.008,0.008)
  (0.025,1.009)rectangle++(0.008,0.008) (0.371,0.868)rectangle++(0.008,0.008) (-0.144,0.949)rectangle++(0.008,0.008) (0.272,0.901)rectangle++(0.008,0.008) (0.445,0.731)rectangle++(0.008,0.008)
  (0.464,0.592)rectangle++(0.008,0.008) (0.423,0.521)rectangle++(0.008,0.008) (0.003,0.808)rectangle++(0.008,0.008) (0.289,0.785)rectangle++(0.008,0.008) (0.119,0.927)rectangle++(0.008,0.008)
  (0.385,0.798)rectangle++(0.008,0.008) (0.008,0.333)rectangle++(0.008,0.008) (-0.119,0.574)rectangle++(0.008,0.008) (-0.263,0.639)rectangle++(0.008,0.008) (-0.012,0.266)rectangle++(0.008,0.008)
  (0.089,0.542)rectangle++(0.008,0.008) (0.035,0.825)rectangle++(0.008,0.008) (0.31,0.781)rectangle++(0.008,0.008) (0.42,0.663)rectangle++(0.008,0.008) (0.174,0.876)rectangle++(0.008,0.008)
  (0.392,0.755)rectangle++(0.008,0.008) (0.448,0.621)rectangle++(0.008,0.008) (0.186,0.859)rectangle++(0.008,0.008) (-0.225,0.878)rectangle++(0.008,0.008) (-0.42,0.74)rectangle++(0.008,0.008)
  (-0.179,0.909)rectangle++(0.008,0.008) (-0.41,0.771)rectangle++(0.008,0.008) (0.087,0.915)rectangle++(0.008,0.008) (0.366,0.803)rectangle++(0.008,0.008) (0.151,0.935)rectangle++(0.008,0.008)
  (-0.268,0.901)rectangle++(0.008,0.008) (-0.448,0.735)rectangle++(0.008,0.008) (-0.19,0.906)rectangle++(0.008,0.008) (-0.414,0.766)rectangle++(0.008,0.008) (-0.466,0.621)rectangle++(0.008,0.008)
  (0.008,0.867)rectangle++(0.008,0.008) (-0.309,0.818)rectangle++(0.008,0.008) (-0.437,0.683)rectangle++(0.008,0.008) (0.043,0.885)rectangle++(0.008,0.008) (0.016,0.969)rectangle++(0.008,0.008)
  (-0.342,0.868)rectangle++(0.008,0.008) (-0.146,0.962)rectangle++(0.008,0.008) (0.276,0.908)rectangle++(0.008,0.008) (-0.202,0.933)rectangle++(0.008,0.008) (0.24,0.915)rectangle++(0.008,0.008)
  (-0.221,0.923)rectangle++(0.008,0.008) (-0.01,0.385)rectangle++(0.008,0.008) (0.132,0.596)rectangle++(0.008,0.008) (0.273,0.642)rectangle++(0.008,0.008) (0.112,0.867)rectangle++(0.008,0.008)
  (0.36,0.773)rectangle++(0.008,0.008) (0.44,0.641)rectangle++(0.008,0.008) (0.429,0.552)rectangle++(0.008,0.008) (-0.005,0.825)rectangle++(0.008,0.008) (-0.3,0.795)rectangle++(0.008,0.008)
  (-0.425,0.675)rectangle++(0.008,0.008) (-0.439,0.576)rectangle++(0.008,0.008) (0.004,0.836)rectangle++(0.008,0.008) (-0.004,0.349)rectangle++(0.008,0.008) (-0.004,0.744)rectangle++(0.008,0.008)
  (0.263,0.758)rectangle++(0.008,0.008) (0.108,0.916)rectangle++(0.008,0.008) (-0.001,0.382)rectangle++(0.008,0.008) (-0.003,0.758)rectangle++(0.008,0.008) (0.269,0.764)rectangle++(0.008,0.008)
  (-0.153,0.864)rectangle++(0.008,0.008) (-0.382,0.76)rectangle++(0.008,0.008) (-0.45,0.63)rectangle++(0.008,0.008) (-0.191,0.862)rectangle++(0.008,0.008) (-0.399,0.745)rectangle++(0.008,0.008)
  (0.083,0.899)rectangle++(0.008,0.008) (-0.287,0.86)rectangle++(0.008,0.008) (-0.123,0.959)rectangle++(0.008,0.008) (-0.054,1)rectangle++(0.008,0.008) (0.332,0.893)rectangle++(0.008,0.008)
  (0.006,0.373)rectangle++(0.008,0.008) (-0.134,0.592)rectangle++(0.008,0.008) (-0.008,0.246)rectangle++(0.008,0.008) (-0.095,0.529)rectangle++(0.008,0.008) (-0.236,0.627)rectangle++(0.008,0.008)
  (-0.101,0.861)rectangle++(0.008,0.008) (-0.357,0.777)rectangle++(0.008,0.008) (-0.152,0.924)rectangle++(0.008,0.008) (-0.403,0.787)rectangle++(0.008,0.008) (0.096,0.92)rectangle++(0.008,0.008)
  (-0.001,0.384)rectangle++(0.008,0.008) (0.136,0.592)rectangle++(0.008,0.008) (0,0.246)rectangle++(0.008,0.008) (-0.002,0.701)rectangle++(0.008,0.008) (-0.255,0.739)rectangle++(0.008,0.008)
  (-0.109,0.908)rectangle++(0.008,0.008) (0.273,0.87)rectangle++(0.008,0.008) (-0.189,0.915)rectangle++(0.008,0.008) (-0.417,0.77)rectangle++(0.008,0.008) (-0.016,0.321)rectangle++(0.008,0.008)
  (-0.126,0.56)rectangle++(0.008,0.008) (-0.261,0.63)rectangle++(0.008,0.008) (-0.348,0.614)rectangle++(0.008,0.008) (-0.148,0.856)rectangle++(0.008,0.008) (0.237,0.86)rectangle++(0.008,0.008)
  (-0.203,0.897)rectangle++(0.008,0.008) (-0.01,0.374)rectangle++(0.008,0.008) (0.128,0.591)rectangle++(0.008,0.008) (-0.156,0.735)rectangle++(0.008,0.008) (-0.337,0.7)rectangle++(0.008,0.008)
  (0.095,0.857)rectangle++(0.008,0.008) (0.038,0.957)rectangle++(0.008,0.008) (-0.328,0.87)rectangle++(0.008,0.008) (-0.464,0.7)rectangle++(0.008,0.008) (-0.466,0.573)rectangle++(0.008,0.008)
  (-0.421,0.515)rectangle++(0.008,0.008) (-0.38,0.504)rectangle++(0.008,0.008) (-0.162,0.809)rectangle++(0.008,0.008) (-0.07,0.938)rectangle++(0.008,0.008) (-0.371,0.823)rectangle++(0.008,0.008)
  (-0.467,0.663)rectangle++(0.008,0.008) (0.023,0.886)rectangle++(0.008,0.008) (0.007,0.97)rectangle++(0.008,0.008) (-0.347,0.865)rectangle++(0.008,0.008) (-0.148,0.961)rectangle++(0.008,0.008)
  (0.275,0.908)rectangle++(0.008,0.008) (-0.202,0.933)rectangle++(0.008,0.008) (-0.429,0.773)rectangle++(0.008,0.008) (-0.183,0.922)rectangle++(0.008,0.008) (-0.079,0.985)rectangle++(0.008,0.008)
  (-0.391,0.841)rectangle++(0.008,0.008) (-0.483,0.664)rectangle++(0.008,0.008) (-0.461,0.55)rectangle++(0.008,0.008) (-0.411,0.506)rectangle++(0.008,0.008) (-0.008,0.794)rectangle++(0.008,0.008)
  (0.279,0.782)rectangle++(0.008,0.008) (0.115,0.926)rectangle++(0.008,0.008) (0.383,0.799)rectangle++(0.008,0.008) (0.459,0.645)rectangle++(0.008,0.008) (-0.024,0.878)rectangle++(0.008,0.008)
  (0.302,0.826)rectangle++(0.008,0.008) (-0.161,0.905)rectangle++(0.008,0.008) (-0.07,0.978)rectangle++(0.008,0.008) (-0.385,0.841)rectangle++(0.008,0.008) (-0.48,0.666)rectangle++(0.008,0.008)
  (-0.461,0.553)rectangle++(0.008,0.008) (-0.412,0.507)rectangle++(0.008,0.008) (-0.373,0.504)rectangle++(0.008,0.008) (0.009,0.78)rectangle++(0.008,0.008) (0.001,0.925)rectangle++(0.008,0.008)
  (-0.002,0.986)rectangle++(0.008,0.008) (0.351,0.868)rectangle++(0.008,0.008) (-0.153,0.941)rectangle++(0.008,0.008) (-0.067,0.993)rectangle++(0.008,0.008) (-0.389,0.849)rectangle++(0.008,0.008)
  (0.125,0.943)rectangle++(0.008,0.008) (0.393,0.803)rectangle++(0.008,0.008) (-0.111,0.927)rectangle++(0.008,0.008) (-0.385,0.803)rectangle++(0.008,0.008) (-0.164,0.935)rectangle++(0.008,0.008)
  (-0.412,0.788)rectangle++(0.008,0.008) (0.092,0.924)rectangle++(0.008,0.008) (-0.291,0.874)rectangle++(0.008,0.008) (-0.449,0.715)rectangle++(0.008,0.008) (-0.017,0.298)rectangle++(0.008,0.008)
  (-0.004,0.123)rectangle++(0.008,0.008) (-0.004,0.049)rectangle++(0.008,0.008) (-0.004,0.018)rectangle++(0.008,0.008) (-0.012,0.426)rectangle++(0.008,0.008) (-0.161,0.61)rectangle++(0.008,0.008)
  (-0.295,0.64)rectangle++(0.008,0.008) (0.093,0.814)rectangle++(0.008,0.008) (0.333,0.755)rectangle++(0.008,0.008) (0.137,0.915)rectangle++(0.008,0.008) (0.055,0.982)rectangle++(0.008,0.008)
  (0.375,0.845)rectangle++(0.008,0.008) (-0.134,0.94)rectangle++(0.008,0.008) (0.274,0.894)rectangle++(0.008,0.008) (-0.198,0.925)rectangle++(0.008,0.008) (-0.424,0.772)rectangle++(0.008,0.008)
  (-0.473,0.62)rectangle++(0.008,0.008) (0.005,0.869)rectangle++(0.008,0.008) (0,0.963)rectangle++(0.008,0.008) (-0.348,0.859)rectangle++(0.008,0.008) (-0.148,0.958)rectangle++(0.008,0.008)
  (0.274,0.907)rectangle++(0.008,0.008) (-0.203,0.932)rectangle++(0.008,0.008) (-0.087,0.989)rectangle++(0.008,0.008) (-0.039,1.013)rectangle++(0.008,0.008) (-0.019,1.023)rectangle++(0.008,0.008)
  (-0.004,0.427)rectangle++(0.008,0.008) (-0.158,0.613)rectangle++(0.008,0.008) (-0.295,0.643)rectangle++(0.008,0.008) (-0.368,0.608)rectangle++(0.008,0.008) (-0.389,0.566)rectangle++(0.008,0.008)
  (0.024,0.814)rectangle++(0.008,0.008) (-0.283,0.8)rectangle++(0.008,0.008) (-0.121,0.934)rectangle++(0.008,0.008) (0.277,0.886)rectangle++(0.008,0.008) (0.004,0.37)rectangle++(0.008,0.008)
  (-0.134,0.59)rectangle++(0.008,0.008) (-0.058,0.845)rectangle++(0.008,0.008) (-0.027,0.953)rectangle++(0.008,0.008) (0.327,0.861)rectangle++(0.008,0.008) (0.135,0.959)rectangle++(0.008,0.008)
  (-0.284,0.906)rectangle++(0.008,0.008) (-0.457,0.732)rectangle++(0.008,0.008) (-0.194,0.905)rectangle++(0.008,0.008) (0.234,0.899)rectangle++(0.008,0.008) (-0.218,0.914)rectangle++(0.008,0.008)
  (-0.094,0.981)rectangle++(0.008,0.008) (-0.397,0.834)rectangle++(0.008,0.008) (-0.483,0.659)rectangle++(0.008,0.008) (-0.46,0.548)rectangle++(0.008,0.008) (-0.195,0.828)rectangle++(0.008,0.008)
  (0.206,0.864)rectangle++(0.008,0.008) (0.402,0.738)rectangle++(0.008,0.008) (0.447,0.61)rectangle++(0.008,0.008) (0.421,0.536)rectangle++(0.008,0.008) (0.383,0.511)rectangle++(0.008,0.008)
  (0.357,0.513)rectangle++(0.008,0.008) (0.148,0.813)rectangle++(0.008,0.008) (0.357,0.735)rectangle++(0.008,0.008) (0.425,0.625)rectangle++(0.008,0.008) (0.417,0.55)rectangle++(0.008,0.008)
  (0.386,0.519)rectangle++(0.008,0.008) (-0.012,0.794)rectangle++(0.008,0.008) (-0.008,0.931)rectangle++(0.008,0.008) (0.328,0.845)rectangle++(0.008,0.008) (0.006,0.352)rectangle++(0.008,0.008)
  (0.128,0.575)rectangle++(0.008,0.008) (0.263,0.634)rectangle++(0.008,0.008) (0.108,0.864)rectangle++(0.008,0.008) (-0.263,0.853)rectangle++(0.008,0.008) (-0.428,0.715)rectangle++(0.008,0.008)
  (-0.455,0.593)rectangle++(0.008,0.008) (0.003,0.85)rectangle++(0.008,0.008) (-0.306,0.809)rectangle++(0.008,0.008) (-0.131,0.937)rectangle++(0.008,0.008) (-0.057,0.991)rectangle++(0.008,0.008)
  (-0.026,1.014)rectangle++(0.008,0.008) (0.349,0.889)rectangle++(0.008,0.008) (-0.161,0.95)rectangle++(0.008,0.008) (-0.009,0.397)rectangle++(0.008,0.008) (-0.006,0.764)rectangle++(0.008,0.008)
  (-0.005,0.919)rectangle++(0.008,0.008) (-0.004,0.984)rectangle++(0.008,0.008) (-0.004,1.011)rectangle++(0.008,0.008) (-0.004,1.022)rectangle++(0.008,0.008) (-0.004,1.027)rectangle++(0.008,0.008)
  (-0.372,0.887)rectangle++(0.008,0.008) (-0.159,0.97)rectangle++(0.008,0.008) (0.273,0.917)rectangle++(0.008,0.008) (-0.206,0.936)rectangle++(0.008,0.008) (0.239,0.918)rectangle++(0.008,0.008)
  (-0.222,0.924)rectangle++(0.008,0.008) (0.228,0.918)rectangle++(0.008,0.008) (0.093,0.983)rectangle++(0.008,0.008) (0.393,0.832)rectangle++(0.008,0.008) (0.008,0.347)rectangle++(0.008,0.008)
  (-0.124,0.581)rectangle++(0.008,0.008) (-0.054,0.842)rectangle++(0.008,0.008) (0.275,0.82)rectangle++(0.008,0.008) (0.418,0.693)rectangle++(0.008,0.008) (-0.06,0.885)rectangle++(0.008,0.008)
  (0.288,0.843)rectangle++(0.008,0.008) (0.432,0.699)rectangle++(0.008,0.008) (-0.056,0.893)rectangle++(0.008,0.008) (0.293,0.844)rectangle++(0.008,0.008) (-0.171,0.91)rectangle++(0.008,0.008)
  (-0.407,0.774)rectangle++(0.008,0.008) (-0.173,0.923)rectangle++(0.008,0.008) (0.25,0.9)rectangle++(0.008,0.008) (0.103,0.976)rectangle++(0.008,0.008) (-0.305,0.902)rectangle++(0.008,0.008)
  (-0.13,0.976)rectangle++(0.008,0.008) (0.288,0.909)rectangle++(0.008,0.008) (0.456,0.729)rectangle++(0.008,0.008) (0.01,0.304)rectangle++(0.008,0.008) (-0.004,0.125)rectangle++(0.008,0.008)
  (0.042,0.475)rectangle++(0.008,0.008) (0.188,0.618)rectangle++(0.008,0.008) (0.077,0.857)rectangle++(0.008,0.008) (-0.275,0.839)rectangle++(0.008,0.008) (-0.118,0.95)rectangle++(0.008,0.008)
  (-0.397,0.811)rectangle++(0.008,0.008) (-0.169,0.938)rectangle++(0.008,0.008) (-0.416,0.788)rectangle++(0.008,0.008) (-0.177,0.929)rectangle++(0.008,0.008) (-0.416,0.78)rectangle++(0.008,0.008)
  (-0.177,0.925)rectangle++(0.008,0.008) (0.249,0.903)rectangle++(0.008,0.008) (0.102,0.977)rectangle++(0.008,0.008) (0.041,1.008)rectangle++(0.008,0.008) (-0.345,0.894)rectangle++(0.008,0.008)
  (-0.147,0.973)rectangle++(0.008,0.008) (0.28,0.914)rectangle++(0.008,0.008) (-0.202,0.937)rectangle++(0.008,0.008) (-0.087,0.991)rectangle++(0.008,0.008) (0.313,0.901)rectangle++(0.008,0.008)
  (0.464,0.716)rectangle++(0.008,0.008) (0.193,0.898)rectangle++(0.008,0.008) (0.408,0.758)rectangle++(0.008,0.008) (-0.088,0.912)rectangle++(0.008,0.008) (-0.039,0.981)rectangle++(0.008,0.008)
  (0.331,0.879)rectangle++(0.008,0.008) (0.137,0.967)rectangle++(0.008,0.008) (0,0.404)rectangle++(0.008,0.008) (-0.148,0.604)rectangle++(0.008,0.008) (-0.287,0.643)rectangle++(0.008,0.008)
  (-0.364,0.611)rectangle++(0.008,0.008) (-0.388,0.568)rectangle++(0.008,0.008) (0.025,0.815)rectangle++(0.008,0.008) (0.302,0.78)rectangle++(0.008,0.008) (-0.144,0.883)rectangle++(0.008,0.008)
  (0.249,0.872)rectangle++(0.008,0.008) (-0.201,0.906)rectangle++(0.008,0.008) (0.231,0.903)rectangle++(0.008,0.008) (0.427,0.746)rectangle++(0.008,0.008) (0.177,0.911)rectangle++(0.008,0.008)
  (-0.248,0.899)rectangle++(0.008,0.008) (-0.011,0.375)rectangle++(0.008,0.008) (-0.143,0.587)rectangle++(0.008,0.008) (-0.278,0.636)rectangle++(0.008,0.008) (-0.119,0.865)rectangle++(0.008,0.008)
  (-0.007,0.361)rectangle++(0.008,0.008) (-0.004,0.149)rectangle++(0.008,0.008) (-0.059,0.486)rectangle++(0.008,0.008) (-0.204,0.62)rectangle++(0.008,0.008) (0.128,0.773)rectangle++(0.008,0.008)
  (-0.221,0.818)rectangle++(0.008,0.008) (-0.397,0.714)rectangle++(0.008,0.008) (0.073,0.884)rectangle++(0.008,0.008) (-0.286,0.85)rectangle++(0.008,0.008) (-0.438,0.705)rectangle++(0.008,0.008)
  (-0.456,0.585)rectangle++(0.008,0.008) (-0.421,0.524)rectangle++(0.008,0.008) (-0.179,0.818)rectangle++(0.008,0.008) (-0.377,0.729)rectangle++(0.008,0.008) (0.087,0.884)rectangle++(0.008,0.008)
  (0.034,0.969)rectangle++(0.008,0.008) (-0.003,0.405)rectangle++(0.008,0.008) (0.142,0.602)rectangle++(0.008,0.008) (0.058,0.851)rectangle++(0.008,0.008) (0.329,0.785)rectangle++(0.008,0.008)
  (0.006,0.327)rectangle++(0.008,0.008) (-0.118,0.571)rectangle++(0.008,0.008) (-0.007,0.237)rectangle++(0.008,0.008) (-0.005,0.697)rectangle++(0.008,0.008) (0.246,0.737)rectangle++(0.008,0.008)
  (0.375,0.666)rectangle++(0.008,0.008) (0.155,0.877)rectangle++(0.008,0.008) (-0.246,0.876)rectangle++(0.008,0.008) (-0.106,0.965)rectangle++(0.008,0.008) (-0.047,1.003)rectangle++(0.008,0.008)
  (-0.005,0.419)rectangle++(0.008,0.008) (-0.005,0.774)rectangle++(0.008,0.008) (-0.282,0.771)rectangle++(0.008,0.008) (-0.121,0.922)rectangle++(0.008,0.008) (-0.388,0.797)rectangle++(0.008,0.008)
  (-0.165,0.933)rectangle++(0.008,0.008) (0.257,0.902)rectangle++(0.008,0.008) (0.106,0.976)rectangle++(0.008,0.008) (-0.304,0.903)rectangle++(0.008,0.008) (0.183,0.938)rectangle++(0.008,0.008)
  (0.418,0.78)rectangle++(0.008,0.008) (0.469,0.624)rectangle++(0.008,0.008) (-0.012,0.872)rectangle++(0.008,0.008) (-0.004,0.364)rectangle++(0.008,0.008) (0.127,0.584)rectangle++(0.008,0.008)
  (0,0.243)rectangle++(0.008,0.008) (-0.002,0.7)rectangle++(0.008,0.008) (-0.255,0.738)rectangle++(0.008,0.008) (-0.109,0.908)rectangle++(0.008,0.008) (-0.378,0.795)rectangle++(0.008,0.008)
  (-0.161,0.932)rectangle++(0.008,0.008) (-0.07,0.989)rectangle++(0.008,0.008) (-0.032,1.013)rectangle++(0.008,0.008) (-0.016,1.023)rectangle++(0.008,0.008) (-0.376,0.881)rectangle++(0.008,0.008)
  (0.142,0.953)rectangle++(0.008,0.008) (-0.279,0.906)rectangle++(0.008,0.008) (-0.12,0.978)rectangle++(0.008,0.008) (0.294,0.906)rectangle++(0.008,0.008) (0.457,0.726)rectangle++(0.008,0.008)
  (0.01,0.302)rectangle++(0.008,0.008) (0.112,0.551)rectangle++(0.008,0.008) (0.247,0.628)rectangle++(0.008,0.008) (-0.115,0.795)rectangle++(0.008,0.008) (-0.34,0.741)rectangle++(0.008,0.008)
  (-0.145,0.909)rectangle++(0.008,0.008) (-0.395,0.783)rectangle++(0.008,0.008) (-0.168,0.927)rectangle++(0.008,0.008) (0.253,0.9)rectangle++(0.008,0.008) (0.436,0.737)rectangle++(0.008,0.008)
  (-0.067,0.912)rectangle++(0.008,0.008) (-0.36,0.812)rectangle++(0.008,0.008) (-0.458,0.662)rectangle++(0.008,0.008) (-0.018,0.276)rectangle++(0.008,0.008) (-0.11,0.539)rectangle++(0.008,0.008)
  (-0.049,0.824)rectangle++(0.008,0.008) (-0.023,0.944)rectangle++(0.008,0.008) (-0.351,0.842)rectangle++(0.008,0.008) (-0.465,0.679)rectangle++(0.008,0.008) (0.029,0.893)rectangle++(0.008,0.008)
  (0.331,0.814)rectangle++(0.008,0.008) (0.137,0.94)rectangle++(0.008,0.008) (-0.277,0.898)rectangle++(0.008,0.008) (0.193,0.925)rectangle++(0.008,0.008) (0.079,0.986)rectangle++(0.008,0.008)
  (-0.32,0.898)rectangle++(0.008,0.008) (0.174,0.941)rectangle++(0.008,0.008) (0.001,0.393)rectangle++(0.008,0.008) (-0.004,0.163)rectangle++(0.008,0.008) (-0.063,0.492)rectangle++(0.008,0.008)
  (0.146,0.664)rectangle++(0.008,0.008) (-0.174,0.775)rectangle++(0.008,0.008) (0.196,0.833)rectangle++(0.008,0.008) (0.08,0.947)rectangle++(0.008,0.008) (-0.305,0.881)rectangle++(0.008,0.008)
  (-0.458,0.713)rectangle++(0.008,0.008) (0.045,0.906)rectangle++(0.008,0.008) (-0.307,0.849)rectangle++(0.008,0.008) (-0.013,0.354)rectangle++(0.008,0.008) (0.12,0.583)rectangle++(0.008,0.008)
  (-0.157,0.728)rectangle++(0.008,0.008) (-0.336,0.696)rectangle++(0.008,0.008) (0.094,0.854)rectangle++(0.008,0.008) (-0.265,0.843)rectangle++(0.008,0.008) (-0.426,0.71)rectangle++(0.008,0.008)
  (-0.452,0.592)rectangle++(0.008,0.008) (0.004,0.848)rectangle++(0.008,0.008) (-0.305,0.808)rectangle++(0.008,0.008) (-0.13,0.937)rectangle++(0.008,0.008) (-0.008,0.391)rectangle++(0.008,0.008)
  (-0.147,0.595)rectangle++(0.008,0.008) (-0.064,0.848)rectangle++(0.008,0.008) (0.273,0.827)rectangle++(0.008,0.008) (-0.174,0.894)rectangle++(0.008,0.008) (0.239,0.887)rectangle++(0.008,0.008)
  (-0.211,0.91)rectangle++(0.008,0.008) (0.228,0.908)rectangle++(0.008,0.008) (0.093,0.979)rectangle++(0.008,0.008) (0.391,0.831)rectangle++(0.008,0.008) (-0.121,0.939)rectangle++(0.008,0.008)
  (0.279,0.889)rectangle++(0.008,0.008) (-0.193,0.925)rectangle++(0.008,0.008) (-0.084,0.986)rectangle++(0.008,0.008) (0.313,0.897)rectangle++(0.008,0.008) (0.006,0.374)rectangle++(0.008,0.008)
  (-0.135,0.592)rectangle++(0.008,0.008) (0.149,0.735)rectangle++(0.008,0.008) (0.06,0.906)rectangle++(0.008,0.008) (0.351,0.809)rectangle++(0.008,0.008) (-0.132,0.914)rectangle++(0.008,0.008)
  (-0.008,0.382)rectangle++(0.008,0.008) (-0.143,0.591)rectangle++(0.008,0.008) (-0.28,0.638)rectangle++(0.008,0.008) (0.099,0.808)rectangle++(0.008,0.008) (0.333,0.75)rectangle++(0.008,0.008)
  (0.138,0.913)rectangle++(0.008,0.008) (0.055,0.981)rectangle++(0.008,0.008) (-0.329,0.887)rectangle++(0.008,0.008) (0.166,0.939)rectangle++(0.008,0.008) (0.067,0.992)rectangle++(0.008,0.008)
  (-0.327,0.897)rectangle++(0.008,0.008) (-0.14,0.974)rectangle++(0.008,0.008) (-0.415,0.815)rectangle++(0.008,0.008) (-0.177,0.94)rectangle++(0.008,0.008) (0.254,0.909)rectangle++(0.008,0.008)
  (0.104,0.979)rectangle++(0.008,0.008) (0.397,0.827)rectangle++(0.008,0.008) (-0.118,0.939)rectangle++(0.008,0.008) (-0.393,0.806)rectangle++(0.008,0.008) (0.108,0.925)rectangle++(0.008,0.008)
  (-0.001,0.386)rectangle++(0.008,0.008) (0.137,0.593)rectangle++(0.008,0.008) (0.274,0.639)rectangle++(0.008,0.008) (-0.107,0.809)rectangle++(0.008,0.008) (-0.007,0.337)rectangle++(0.008,0.008)
  (0.116,0.573)rectangle++(0.008,0.008) (-0.155,0.723)rectangle++(0.008,0.008) (-0.009,0.301)rectangle++(0.008,0.008) (-0.115,0.554)rectangle++(0.008,0.008) (-0.254,0.631)rectangle++(0.008,0.008)
  (0.109,0.796)rectangle++(0.008,0.008) (0.333,0.741)rectangle++(0.008,0.008) (0.416,0.636)rectangle++(0.008,0.008) (-0.041,0.859)rectangle++(0.008,0.008) (0.287,0.823)rectangle++(0.008,0.008)
  (-0.166,0.898)rectangle++(0.008,0.008) (0.244,0.886)rectangle++(0.008,0.008) (0.427,0.734)rectangle++(0.008,0.008) (-0.071,0.907)rectangle++(0.008,0.008) (0.291,0.856)rectangle++(0.008,0.008)
  (0.12,0.957)rectangle++(0.008,0.008) (0.396,0.811)rectangle++(0.008,0.008) (-0.112,0.931)rectangle++(0.008,0.008) (0.281,0.882)rectangle++(0.008,0.008) (0.005,0.368)rectangle++(0.008,0.008)
  (0.133,0.583)rectangle++(0.008,0.008) (0.268,0.635)rectangle++(0.008,0.008) (-0.108,0.805)rectangle++(0.008,0.008) (0.238,0.823)rectangle++(0.008,0.008) (0.097,0.943)rectangle++(0.008,0.008)
  (0.381,0.813)rectangle++(0.008,0.008) (-0.12,0.927)rectangle++(0.008,0.008) (-0.389,0.8)rectangle++(0.008,0.008) (-0.467,0.646)rectangle++(0.008,0.008) (-0.199,0.869)rectangle++(0.008,0.008)
  (0.219,0.884)rectangle++(0.008,0.008) (0.09,0.969)rectangle++(0.008,0.008) (-0.309,0.894)rectangle++(0.008,0.008) (-0.013,0.373)rectangle++(0.008,0.008) (-0.143,0.585)rectangle++(0.008,0.008)
  (-0.062,0.843)rectangle++(0.008,0.008) (0.272,0.824)rectangle++(0.008,0.008) (-0.174,0.893)rectangle++(0.008,0.008) (-0.009,0.373)rectangle++(0.008,0.008) (0.128,0.59)rectangle++(0.008,0.008)
  (0.269,0.64)rectangle++(0.008,0.008) (0.11,0.867)rectangle++(0.008,0.008) (0.359,0.773)rectangle++(0.008,0.008) (0.439,0.641)rectangle++(0.008,0.008) (-0.032,0.869)rectangle++(0.008,0.008)
  (-0.016,0.963)rectangle++(0.008,0.008) (0.336,0.862)rectangle++(0.008,0.008) (0.006,0.36)rectangle++(0.008,0.008) (0.131,0.579)rectangle++(0.008,0.008) (-0.151,0.73)rectangle++(0.008,0.008)
  (0.191,0.804)rectangle++(0.008,0.008) (0.078,0.935)rectangle++(0.008,0.008) (-0.302,0.875)rectangle++(0.008,0.008) (-0.129,0.965)rectangle++(0.008,0.008) (0.285,0.904)rectangle++(0.008,0.008)
  (0.117,0.977)rectangle++(0.008,0.008) (-0.299,0.908)rectangle++(0.008,0.008) (0.187,0.938)rectangle++(0.008,0.008) (-0.253,0.915)rectangle++(0.008,0.008) (-0.446,0.747)rectangle++(0.008,0.008)
  (0.062,0.917)rectangle++(0.008,0.008) (-0.303,0.861)rectangle++(0.008,0.008) (-0.449,0.705)rectangle++(0.008,0.008) (0.046,0.899)rectangle++(0.008,0.008) (-0.304,0.846)rectangle++(0.008,0.008)
  (-0.445,0.698)rectangle++(0.008,0.008) (0.045,0.894)rectangle++(0.008,0.008) (-0.302,0.844)rectangle++(0.008,0.008) (-0.443,0.697)rectangle++(0.008,0.008) (0.046,0.893)rectangle++(0.008,0.008)
  (0.339,0.808)rectangle++(0.008,0.008) (0.443,0.665)rectangle++(0.008,0.008) (0.439,0.562)rectangle++(0.008,0.008) (-0.004,0.833)rectangle++(0.008,0.008) (0.295,0.798)rectangle++(0.008,0.008)
  (0.419,0.676)rectangle++(0.008,0.008) (0.432,0.576)rectangle++(0.008,0.008) (0.402,0.525)rectangle++(0.008,0.008) (0.371,0.513)rectangle++(0.008,0.008) (0.153,0.813)rectangle++(0.008,0.008)
  (0.36,0.733)rectangle++(0.008,0.008) (0.425,0.623)rectangle++(0.008,0.008) (0.009,0.259)rectangle++(0.008,0.008) (0.001,0.707)rectangle++(0.008,0.008) (-0.255,0.743)rectangle++(0.008,0.008)
  (-0.012,0.31)rectangle++(0.008,0.008) (-0.119,0.557)rectangle++(0.008,0.008) (-0.257,0.631)rectangle++(0.008,0.008) (0.107,0.797)rectangle++(0.008,0.008) (0.333,0.742)rectangle++(0.008,0.008)
  (0.416,0.637)rectangle++(0.008,0.008) (0.417,0.559)rectangle++(0.008,0.008) (-0.013,0.824)rectangle++(0.008,0.008) (-0.303,0.791)rectangle++(0.008,0.008) (-0.425,0.673)rectangle++(0.008,0.008)
  (-0.438,0.575)rectangle++(0.008,0.008) (0.004,0.836)rectangle++(0.008,0.008) (0.3,0.797)rectangle++(0.008,0.008) (0.421,0.674)rectangle++(0.008,0.008) (0.174,0.881)rectangle++(0.008,0.008)
  (-0.238,0.884)rectangle++(0.008,0.008) (-0.428,0.738)rectangle++(0.008,0.008) (-0.463,0.604)rectangle++(0.008,0.008) (-0.431,0.53)rectangle++(0.008,0.008) (-0.39,0.508)rectangle++(0.008,0.008)
  (0.002,0.788)rectangle++(0.008,0.008) (-0.284,0.78)rectangle++(0.008,0.008) (-0.012,0.325)rectangle++(0.008,0.008) (-0.125,0.564)rectangle++(0.008,0.008) (-0.262,0.632)rectangle++(0.008,0.008)
  (0.105,0.799)rectangle++(0.008,0.008) (-0.241,0.822)rectangle++(0.008,0.008) (0.183,0.878)rectangle++(0.008,0.008) (0.396,0.752)rectangle++(0.008,0.008) (-0.091,0.905)rectangle++(0.008,0.008)
  (0.281,0.862)rectangle++(0.008,0.008) (-0.183,0.914)rectangle++(0.008,0.008) (-0.414,0.771)rectangle++(0.008,0.008) (0.086,0.917)rectangle++(0.008,0.008) (-0.292,0.869)rectangle++(0.008,0.008)
  (0.176,0.918)rectangle++(0.008,0.008) (-0.251,0.902)rectangle++(0.008,0.008) (0.207,0.918)rectangle++(0.008,0.008) (-0.237,0.913)rectangle++(0.008,0.008) (-0.438,0.752)rectangle++(0.008,0.008)
  (-0.472,0.607)rectangle++(0.008,0.008) (0,0.862)rectangle++(0.008,0.008) (-0.002,0.96)rectangle++(0.008,0.008) (-0.347,0.857)rectangle++(0.008,0.008) (-0.468,0.687)rectangle++(0.008,0.008)
  (-0.018,0.286)rectangle++(0.008,0.008) (-0.114,0.544)rectangle++(0.008,0.008) (0.141,0.706)rectangle++(0.008,0.008) (-0.191,0.792)rectangle++(0.008,0.008) (-0.083,0.93)rectangle++(0.008,0.008)
  (-0.037,0.988)rectangle++(0.008,0.008) (-0.373,0.858)rectangle++(0.008,0.008) (-0.481,0.678)rectangle++(0.008,0.008) (0.022,0.898)rectangle++(0.008,0.008) (-0.003,0.375)rectangle++(0.008,0.008)
  (-0.139,0.589)rectangle++(0.008,0.008) (-0.278,0.639)rectangle++(0.008,0.008) (0.101,0.808)rectangle++(0.008,0.008) (-0.246,0.824)rectangle++(0.008,0.008) (-0.106,0.944)rectangle++(0.008,0.008)
  (-0.389,0.813)rectangle++(0.008,0.008) (-0.166,0.939)rectangle++(0.008,0.008) (0.259,0.905)rectangle++(0.008,0.008) (0.441,0.737)rectangle++(0.008,0.008) (0.183,0.907)rectangle++(0.008,0.008)
  (0.407,0.766)rectangle++(0.008,0.008) (0.008,0.319)rectangle++(0.008,0.008) (-0.114,0.568)rectangle++(0.008,0.008) (0.15,0.717)rectangle++(0.008,0.008) (-0.191,0.8)rectangle++(0.008,0.008)
  (-0.01,0.334)rectangle++(0.008,0.008) (-0.127,0.568)rectangle++(0.008,0.008) (0.144,0.722)rectangle++(0.008,0.008) (0.323,0.695)rectangle++(0.008,0.008) (0.395,0.619)rectangle++(0.008,0.008)
  (-0.044,0.843)rectangle++(0.008,0.008) (-0.325,0.789)rectangle++(0.008,0.008) (-0.139,0.929)rectangle++(0.008,0.008) (-0.399,0.794)rectangle++(0.008,0.008) (-0.469,0.64)rectangle++(0.008,0.008)
  (-0.447,0.544)rectangle++(0.008,0.008) (-0.402,0.508)rectangle++(0.008,0.008) (-0.171,0.811)rectangle++(0.008,0.008) (-0.372,0.729)rectangle++(0.008,0.008) (-0.434,0.62)rectangle++(0.008,0.008)
  (-0.423,0.548)rectangle++(0.008,0.008) (0.001,0.818)rectangle++(0.008,0.008) (0.292,0.79)rectangle++(0.008,0.008) (0.12,0.929)rectangle++(0.008,0.008) (-0.28,0.887)rectangle++(0.008,0.008)
  (-0.12,0.97)rectangle++(0.008,0.008) (-0.007,0.405)rectangle++(0.008,0.008) (-0.005,0.768)rectangle++(0.008,0.008) (0.271,0.769)rectangle++(0.008,0.008) (0.004,0.321)rectangle++(0.008,0.008)
  (-0.001,0.732)rectangle++(0.008,0.008) (0.261,0.751)rectangle++(0.008,0.008) (0.107,0.913)rectangle++(0.008,0.008) (-0.281,0.875)rectangle++(0.008,0.008) (0.183,0.917)rectangle++(0.008,0.008)
  (0.41,0.77)rectangle++(0.008,0.008) (-0.091,0.918)rectangle++(0.008,0.008) (-0.373,0.806)rectangle++(0.008,0.008) (0.117,0.918)rectangle++(0.008,0.008) (-0.278,0.881)rectangle++(0.008,0.008)
  (0.187,0.918)rectangle++(0.008,0.008) (-0.246,0.906)rectangle++(0.008,0.008) (0.21,0.918)rectangle++(0.008,0.008) (-0.235,0.914)rectangle++(0.008,0.008) (0.218,0.918)rectangle++(0.008,0.008)
  (-0.232,0.917)rectangle++(0.008,0.008) (-0.437,0.755)rectangle++(0.008,0.008) (-0.473,0.609)rectangle++(0.008,0.008) (-0.437,0.528)rectangle++(0.008,0.008) (-0.186,0.82)rectangle++(0.008,0.008)
  (-0.08,0.942)rectangle++(0.008,0.008) (0.299,0.876)rectangle++(0.008,0.008) (-0.18,0.926)rectangle++(0.008,0.008) (0.248,0.904)rectangle++(0.008,0.008) (0.102,0.977)rectangle++(0.008,0.008)
  (-0.001,0.408)rectangle++(0.008,0.008) (0.145,0.603)rectangle++(0.008,0.008) (-0.153,0.747)rectangle++(0.008,0.008) (-0.34,0.706)rectangle++(0.008,0.008) (-0.145,0.894)rectangle++(0.008,0.008)
  (0.252,0.877)rectangle++(0.008,0.008) (0.428,0.727)rectangle++(0.008,0.008) (-0.068,0.904)rectangle++(0.008,0.008) (-0.031,0.977)rectangle++(0.008,0.008) (0.334,0.874)rectangle++(0.008,0.008)
  (-0.163,0.938)rectangle++(0.008,0.008) (-0.071,0.992)rectangle++(0.008,0.008) (-0.006,0.414)rectangle++(0.008,0.008) (0.144,0.608)rectangle++(0.008,0.008) (0.058,0.853)rectangle++(0.008,0.008)
  (0.331,0.785)rectangle++(0.008,0.008) (0.431,0.657)rectangle++(0.008,0.008) (0.179,0.874)rectangle++(0.008,0.008) (0.393,0.752)rectangle++(0.008,0.008) (0.447,0.62)rectangle++(0.008,0.008)
  (0.01,0.258)rectangle++(0.008,0.008) (0.096,0.531)rectangle++(0.008,0.008) (0.038,0.821)rectangle++(0.008,0.008) (0.31,0.778)rectangle++(0.008,0.008) (0.005,0.324)rectangle++(0.008,0.008)
  (0,0.734)rectangle++(0.008,0.008) (-0.004,0.306)rectangle++(0.008,0.008) (-0.115,0.558)rectangle++(0.008,0.008) (-0.05,0.832)rectangle++(0.008,0.008) (-0.324,0.781)rectangle++(0.008,0.008)
  (-0.138,0.926)rectangle++(0.008,0.008) (-0.397,0.793)rectangle++(0.008,0.008) (-0.468,0.64)rectangle++(0.008,0.008) (-0.018,0.266)rectangle++(0.008,0.008) (-0.107,0.535)rectangle++(0.008,0.008)
  (-0.243,0.625)rectangle++(0.008,0.008) (-0.338,0.618)rectangle++(0.008,0.008) (-0.379,0.581)rectangle++(0.008,0.008) (-0.384,0.55)rectangle++(0.008,0.008) (-0.164,0.828)rectangle++(0.008,0.008)
  (-0.071,0.946)rectangle++(0.008,0.008) (0.304,0.874)rectangle++(0.008,0.008) (0.45,0.707)rectangle++(0.008,0.008) (0.187,0.895)rectangle++(0.008,0.008) (0.076,0.973)rectangle++(0.008,0.008);
\fill[nuage]
  (0.03,1.007)rectangle++(0.008,0.008) (0.372,0.866)rectangle++(0.008,0.008) (-0.143,0.948)rectangle++(0.008,0.008) (-0.008,0.396)rectangle++(0.008,0.008) (0.137,0.6)rectangle++(0.008,0.008)
  (-0.155,0.742)rectangle++(0.008,0.008) (-0.068,0.91)rectangle++(0.008,0.008) (-0.006,0.38)rectangle++(0.008,0.008) (-0.142,0.591)rectangle++(0.008,0.008) (-0.062,0.846)rectangle++(0.008,0.008)
  (-0.334,0.784)rectangle++(0.008,0.008) (-0.436,0.658)rectangle++(0.008,0.008) (-0.185,0.874)rectangle++(0.008,0.008) (0.227,0.882)rectangle++(0.008,0.008) (-0.215,0.903)rectangle++(0.008,0.008)
  (0.224,0.906)rectangle++(0.008,0.008) (0.425,0.751)rectangle++(0.008,0.008) (0.462,0.608)rectangle++(0.008,0.008) (-0.01,0.862)rectangle++(0.008,0.008) (-0.006,0.96)rectangle++(0.008,0.008)
  (-0.004,0.401)rectangle++(0.008,0.008) (-0.149,0.601)rectangle++(0.008,0.008) (-0.286,0.641)rectangle++(0.008,0.008) (0.097,0.811)rectangle++(0.008,0.008) (0.334,0.753)rectangle++(0.008,0.008)
  (0.138,0.914)rectangle++(0.008,0.008) (0.388,0.785)rectangle++(0.008,0.008) (0.161,0.927)rectangle++(0.008,0.008) (0.404,0.783)rectangle++(0.008,0.008) (0.463,0.63)rectangle++(0.008,0.008)
  (0.01,0.262)rectangle++(0.008,0.008) (-0.093,0.543)rectangle++(0.008,0.008) (0.151,0.698)rectangle++(0.008,0.008) (-0.184,0.792)rectangle++(0.008,0.008) (0.198,0.844)rectangle++(0.008,0.008)
  (0.081,0.952)rectangle++(0.008,0.008) (-0.001,0.398)rectangle++(0.008,0.008) (-0.146,0.6)rectangle++(0.008,0.008) (-0.064,0.85)rectangle++(0.008,0.008) (0.274,0.828)rectangle++(0.008,0.008)
  (-0.174,0.895)rectangle++(0.008,0.008) (-0.403,0.766)rectangle++(0.008,0.008) (-0.461,0.626)rectangle++(0.008,0.008) (0.012,0.867)rectangle++(0.008,0.008) (-0.308,0.82)rectangle++(0.008,0.008)
  (0.151,0.901)rectangle++(0.008,0.008) (-0.256,0.885)rectangle++(0.008,0.008) (-0.11,0.969)rectangle++(0.008,0.008) (-0.4,0.823)rectangle++(0.008,0.008) (-0.17,0.943)rectangle++(0.008,0.008)
  (-0.418,0.79)rectangle++(0.008,0.008) (-0.178,0.929)rectangle++(0.008,0.008) (-0.417,0.78)rectangle++(0.008,0.008) (-0.177,0.925)rectangle++(0.008,0.008) (0.249,0.903)rectangle++(0.008,0.008)
  (0.435,0.74)rectangle++(0.008,0.008) (0.462,0.599)rectangle++(0.008,0.008) (0.425,0.525)rectangle++(0.008,0.008) (0.003,0.811)rectangle++(0.008,0.008) (-0.292,0.791)rectangle++(0.008,0.008)
  (-0.419,0.677)rectangle++(0.008,0.008) (-0.178,0.882)rectangle++(0.008,0.008) (-0.4,0.759)rectangle++(0.008,0.008) (0.087,0.906)rectangle++(0.008,0.008) (-0.287,0.864)rectangle++(0.008,0.008)
  (0.177,0.914)rectangle++(0.008,0.008) (0.072,0.982)rectangle++(0.008,0.008) (0.383,0.839)rectangle++(0.008,0.008) (-0.128,0.94)rectangle++(0.008,0.008) (-0.398,0.803)rectangle++(0.008,0.008)
  (-0.169,0.935)rectangle++(0.008,0.008) (-0.415,0.786)rectangle++(0.008,0.008) (0.09,0.924)rectangle++(0.008,0.008) (0.036,0.986)rectangle++(0.008,0.008) (-0.339,0.882)rectangle++(0.008,0.008)
  (0.159,0.941)rectangle++(0.008,0.008) (-0.267,0.906)rectangle++(0.008,0.008) (-0.449,0.738)rectangle++(0.008,0.008) (-0.191,0.908)rectangle++(0.008,0.008) (0.236,0.899)rectangle++(0.008,0.008)
  (0.097,0.975)rectangle++(0.008,0.008) (0.392,0.828)rectangle++(0.008,0.008) (-0.12,0.937)rectangle++(0.008,0.008) (-0.393,0.805)rectangle++(0.008,0.008) (-0.471,0.647)rectangle++(0.008,0.008)
  (-0.45,0.547)rectangle++(0.008,0.008) (-0.011,0.827)rectangle++(0.008,0.008) (-0.304,0.793)rectangle++(0.008,0.008) (0.144,0.887)rectangle++(0.008,0.008) (0.058,0.97)rectangle++(0.008,0.008)
  (0.372,0.839)rectangle++(0.008,0.008) (-0.133,0.936)rectangle++(0.008,0.008) (-0.399,0.799)rectangle++(0.008,0.008) (-0.17,0.933)rectangle++(0.008,0.008) (-0.074,0.99)rectangle++(0.008,0.008)
  (-0.391,0.845)rectangle++(0.008,0.008) (-0.016,0.353)rectangle++(0.008,0.008) (-0.009,0.746)rectangle++(0.008,0.008) (0.262,0.76)rectangle++(0.008,0.008) (-0.156,0.86)rectangle++(0.008,0.008)
  (-0.068,0.959)rectangle++(0.008,0.008) (-0.377,0.833)rectangle++(0.008,0.008) (-0.015,0.348)rectangle++(0.008,0.008) (-0.135,0.573)rectangle++(0.008,0.008) (-0.27,0.633)rectangle++(0.008,0.008)
  (-0.116,0.864)rectangle++(0.008,0.008) (-0.365,0.773)rectangle++(0.008,0.008) (0.108,0.9)rectangle++(0.008,0.008) (-0.275,0.869)rectangle++(0.008,0.008) (-0.012,0.363)rectangle++(0.008,0.008)
  (0.123,0.587)rectangle++(0.008,0.008) (0,0.244)rectangle++(0.008,0.008) (0.086,0.528)rectangle++(0.008,0.008) (0.034,0.819)rectangle++(0.008,0.008) (0.307,0.779)rectangle++(0.008,0.008)
  (0.418,0.663)rectangle++(0.008,0.008) (-0.049,0.871)rectangle++(0.008,0.008) (0.288,0.832)rectangle++(0.008,0.008) (0.428,0.694)rectangle++(0.008,0.008) (-0.056,0.889)rectangle++(0.008,0.008)
  (0.291,0.843)rectangle++(0.008,0.008) (0.433,0.698)rectangle++(0.008,0.008) (0.009,0.291)rectangle++(0.008,0.008) (0.107,0.546)rectangle++(0.008,0.008) (0.243,0.628)rectangle++(0.008,0.008)
  (0.335,0.616)rectangle++(0.008,0.008) (0.372,0.579)rectangle++(0.008,0.008) (0.376,0.548)rectangle++(0.008,0.008) (0.156,0.828)rectangle++(0.008,0.008) (-0.228,0.853)rectangle++(0.008,0.008)
  (-0.412,0.728)rectangle++(0.008,0.008) (0.071,0.896)rectangle++(0.008,0.008) (-0.291,0.854)rectangle++(0.008,0.008) (-0.442,0.706)rectangle++(0.008,0.008) (-0.188,0.894)rectangle++(0.008,0.008)
  (-0.409,0.761)rectangle++(0.008,0.008) (0.084,0.91)rectangle++(0.008,0.008) (0.033,0.98)rectangle++(0.008,0.008) (-0.338,0.879)rectangle++(0.008,0.008) (0.158,0.939)rectangle++(0.008,0.008)
  (0.064,0.992)rectangle++(0.008,0.008) (-0.002,0.414)rectangle++(0.008,0.008) (-0.152,0.608)rectangle++(0.008,0.008) (0.147,0.749)rectangle++(0.008,0.008) (0.001,0.312)rectangle++(0.008,0.008)
  (-0.115,0.562)rectangle++(0.008,0.008) (-0.257,0.635)rectangle++(0.008,0.008) (0.109,0.798)rectangle++(0.008,0.008) (0.043,0.933)rectangle++(0.008,0.008) (0.016,0.99)rectangle++(0.008,0.008)
  (-0.35,0.877)rectangle++(0.008,0.008) (-0.477,0.695)rectangle++(0.008,0.008) (0.03,0.904)rectangle++(0.008,0.008) (0.01,0.978)rectangle++(0.008,0.008) (0.353,0.86)rectangle++(0.008,0.008)
  (-0.149,0.938)rectangle++(0.008,0.008) (-0.065,0.992)rectangle++(0.008,0.008) (-0.03,1.014)rectangle++(0.008,0.008) (-0.005,0.424)rectangle++(0.008,0.008) (-0.157,0.611)rectangle++(0.008,0.008)
  (-0.068,0.854)rectangle++(0.008,0.008) (-0.34,0.785)rectangle++(0.008,0.008) (-0.439,0.657)rectangle++(0.008,0.008) (-0.017,0.274)rectangle++(0.008,0.008) (0.089,0.548)rectangle++(0.008,0.008)
  (0.236,0.635)rectangle++(0.008,0.008) (0.097,0.864)rectangle++(0.008,0.008) (-0.001,0.361)rectangle++(0.008,0.008) (-0.003,0.749)rectangle++(0.008,0.008) (-0.272,0.761)rectangle++(0.008,0.008)
  (-0.4,0.67)rectangle++(0.008,0.008) (0.056,0.865)rectangle++(0.008,0.008) (-0.287,0.835)rectangle++(0.008,0.008) (0.166,0.9)rectangle++(0.008,0.008) (0.397,0.769)rectangle++(0.008,0.008)
  (0.164,0.92)rectangle++(0.008,0.008) (0.403,0.778)rectangle++(0.008,0.008) (-0.097,0.919)rectangle++(0.008,0.008) (-0.376,0.804)rectangle++(0.008,0.008) (0.115,0.918)rectangle++(0.008,0.008)
  (0.046,0.983)rectangle++(0.008,0.008) (0.371,0.85)rectangle++(0.008,0.008) (0.154,0.954)rectangle++(0.008,0.008) (-0.274,0.91)rectangle++(0.008,0.008) (0.199,0.93)rectangle++(0.008,0.008)
  (0.002,0.388)rectangle++(0.008,0.008) (0.139,0.593)rectangle++(0.008,0.008) (0.275,0.638)rectangle++(0.008,0.008) (0.113,0.866)rectangle++(0.008,0.008) (-0.261,0.855)rectangle++(0.008,0.008)
  (-0.112,0.957)rectangle++(0.008,0.008) (0.29,0.894)rectangle++(0.008,0.008) (0.451,0.721)rectangle++(0.008,0.008) (0.187,0.901)rectangle++(0.008,0.008) (-0.24,0.898)rectangle++(0.008,0.008)
  (-0.103,0.975)rectangle++(0.008,0.008) (-0.399,0.828)rectangle++(0.008,0.008) (0.113,0.937)rectangle++(0.008,0.008) (-0.287,0.888)rectangle++(0.008,0.008) (-0.012,0.371)rectangle++(0.008,0.008)
  (-0.008,0.753)rectangle++(0.008,0.008) (-0.276,0.761)rectangle++(0.008,0.008) (-0.118,0.917)rectangle++(0.008,0.008) (-0.385,0.796)rectangle++(0.008,0.008) (-0.464,0.646)rectangle++(0.008,0.008)
  (-0.018,0.269)rectangle++(0.008,0.008) (0.087,0.546)rectangle++(0.008,0.008) (-0.001,0.227)rectangle++(0.008,0.008) (-0.085,0.522)rectangle++(0.008,0.008) (0.147,0.686)rectangle++(0.008,0.008)
  (-0.181,0.785)rectangle++(0.008,0.008) (-0.367,0.713)rectangle++(0.008,0.008) (-0.426,0.614)rectangle++(0.008,0.008) (-0.181,0.856)rectangle++(0.008,0.008) (-0.009,0.357)rectangle++(0.008,0.008)
  (-0.006,0.748)rectangle++(0.008,0.008) (-0.005,0.912)rectangle++(0.008,0.008) (0.323,0.835)rectangle++(0.008,0.008) (-0.154,0.916)rectangle++(0.008,0.008) (-0.067,0.982)rectangle++(0.008,0.008)
  (0.319,0.889)rectangle++(0.008,0.008) (-0.175,0.94)rectangle++(0.008,0.008) (-0.076,0.992)rectangle++(0.008,0.008) (-0.393,0.846)rectangle++(0.008,0.008) (-0.485,0.666)rectangle++(0.008,0.008)
  (0.015,0.894)rectangle++(0.008,0.008) (-0.316,0.833)rectangle++(0.008,0.008) (-0.135,0.948)rectangle++(0.008,0.008) (-0.404,0.804)rectangle++(0.008,0.008) (-0.475,0.643)rectangle++(0.008,0.008)
  (0.012,0.88)rectangle++(0.008,0.008) (-0.313,0.825)rectangle++(0.008,0.008) (0.151,0.905)rectangle++(0.008,0.008) (0.061,0.978)rectangle++(0.008,0.008) (-0.002,0.408)rectangle++(0.008,0.008)
  (-0.15,0.605)rectangle++(0.008,0.008) (-0.065,0.852)rectangle++(0.008,0.008) (-0.338,0.785)rectangle++(0.008,0.008) (0.125,0.896)rectangle++(0.008,0.008) (0,0.374)rectangle++(0.008,0.008)
  (-0.004,0.155)rectangle++(0.008,0.008) (-0.004,0.663)rectangle++(0.008,0.008) (-0.004,0.876)rectangle++(0.008,0.008) (-0.318,0.818)rectangle++(0.008,0.008) (-0.441,0.68)rectangle++(0.008,0.008)
  (0.04,0.885)rectangle++(0.008,0.008) (0.015,0.969)rectangle++(0.008,0.008) (-0.343,0.867)rectangle++(0.008,0.008) (-0.014,0.362)rectangle++(0.008,0.008) (0.122,0.587)rectangle++(0.008,0.008)
  (0.265,0.641)rectangle++(0.008,0.008) (0.349,0.615)rectangle++(0.008,0.008) (0.007,0.256)rectangle++(0.008,0.008) (-0.092,0.539)rectangle++(0.008,0.008) (-0.238,0.633)rectangle++(0.008,0.008)
  (-0.338,0.623)rectangle++(0.008,0.008) (-0.144,0.86)rectangle++(0.008,0.008) (-0.063,0.959)rectangle++(0.008,0.008) (0.313,0.877)rectangle++(0.008,0.008) (0.129,0.966)rectangle++(0.008,0.008)
  (0.052,1.003)rectangle++(0.008,0.008) (-0.338,0.896)rectangle++(0.008,0.008) (-0.144,0.974)rectangle++(0.008,0.008) (-0.063,1.007)rectangle++(0.008,0.008) (-0.029,1.021)rectangle++(0.008,0.008)
  (-0.381,0.875)rectangle++(0.008,0.008) (-0.162,0.965)rectangle++(0.008,0.008) (-0.423,0.802)rectangle++(0.008,0.008) (-0.483,0.635)rectangle++(0.008,0.008) (0.005,0.879)rectangle++(0.008,0.008)
  (-0.315,0.823)rectangle++(0.008,0.008) (-0.442,0.683)rectangle++(0.008,0.008) (0.041,0.886)rectangle++(0.008,0.008) (0.015,0.97)rectangle++(0.008,0.008) (-0.003,0.405)rectangle++(0.008,0.008)
  (-0.15,0.603)rectangle++(0.008,0.008) (0.146,0.746)rectangle++(0.008,0.008) (0.059,0.911)rectangle++(0.008,0.008) (0.022,0.98)rectangle++(0.008,0.008) (0.36,0.856)rectangle++(0.008,0.008)
  (0.469,0.679)rectangle++(0.008,0.008) (-0.032,0.897)rectangle++(0.008,0.008) (-0.016,0.975)rectangle++(0.008,0.008) (-0.009,1.007)rectangle++(0.008,0.008) (0.355,0.88)rectangle++(0.008,0.008)
  (0.476,0.692)rectangle++(0.008,0.008) (-0.033,0.905)rectangle++(0.008,0.008) (0.307,0.842)rectangle++(0.008,0.008) (-0.164,0.914)rectangle++(0.008,0.008) (0.251,0.893)rectangle++(0.008,0.008)
  (-0.208,0.917)rectangle++(0.008,0.008) (-0.01,0.383)rectangle++(0.008,0.008) (-0.007,0.758)rectangle++(0.008,0.008) (-0.277,0.764)rectangle++(0.008,0.008) (-0.012,0.318)rectangle++(0.008,0.008)
  (-0.004,0.131)rectangle++(0.008,0.008) (-0.004,0.653)rectangle++(0.008,0.008) (-0.239,0.716)rectangle++(0.008,0.008) (-0.103,0.898)rectangle++(0.008,0.008) (-0.371,0.793)rectangle++(0.008,0.008)
  (-0.457,0.649)rectangle++(0.008,0.008) (-0.444,0.553)rectangle++(0.008,0.008) (-0.017,0.23)rectangle++(0.008,0.008) (0.073,0.528)rectangle++(0.008,0.008) (0.221,0.631)rectangle++(0.008,0.008)
  (0.326,0.626)rectangle++(0.008,0.008) (0.135,0.861)rectangle++(0.008,0.008) (-0.249,0.861)rectangle++(0.008,0.008) (-0.425,0.724)rectangle++(0.008,0.008) (-0.181,0.902)rectangle++(0.008,0.008)
  (-0.009,0.376)rectangle++(0.008,0.008) (-0.006,0.756)rectangle++(0.008,0.008) (0.266,0.764)rectangle++(0.008,0.008) (0.394,0.671)rectangle++(0.008,0.008) (0.419,0.582)rectangle++(0.008,0.008)
  (-0.02,0.835)rectangle++(0.008,0.008) (-0.311,0.794)rectangle++(0.008,0.008) (0.141,0.89)rectangle++(0.008,0.008) (0.057,0.972)rectangle++(0.008,0.008) (0.021,1.006)rectangle++(0.008,0.008)
  (0.007,1.02)rectangle++(0.008,0.008) (0.367,0.88)rectangle++(0.008,0.008) (0.481,0.688)rectangle++(0.008,0.008) (0.2,0.887)rectangle++(0.008,0.008) (0.082,0.97)rectangle++(0.008,0.008)
  (0.032,1.005)rectangle++(0.008,0.008) (-0.003,0.42)rectangle++(0.008,0.008) (-0.004,0.774)rectangle++(0.008,0.008) (-0.004,0.323)rectangle++(0.008,0.008) (-0.004,0.733)rectangle++(0.008,0.008)
  (-0.004,0.906)rectangle++(0.008,0.008) (0.321,0.832)rectangle++(0.008,0.008) (0.132,0.947)rectangle++(0.008,0.008) (0.398,0.802)rectangle++(0.008,0.008) (0.467,0.641)rectangle++(0.008,0.008)
  (0.442,0.542)rectangle++(0.008,0.008) (0.005,0.825)rectangle++(0.008,0.008) (0.296,0.792)rectangle++(0.008,0.008) (0.417,0.673)rectangle++(0.008,0.008) (0.43,0.575)rectangle++(0.008,0.008)
  (-0.012,0.836)rectangle++(0.008,0.008) (0.292,0.803)rectangle++(0.008,0.008) (-0.157,0.89)rectangle++(0.008,0.008) (0.246,0.879)rectangle++(0.008,0.008) (0.425,0.731)rectangle++(0.008,0.008)
  (-0.07,0.905)rectangle++(0.008,0.008) (0.29,0.855)rectangle++(0.008,0.008) (0.12,0.957)rectangle++(0.008,0.008) (0.048,1)rectangle++(0.008,0.008) (-0.339,0.893)rectangle++(0.008,0.008)
  (-0.145,0.973)rectangle++(0.008,0.008) (-0.417,0.812)rectangle++(0.008,0.008) (-0.484,0.642)rectangle++(0.008,0.008) (-0.206,0.867)rectangle++(0.008,0.008) (-0.407,0.742)rectangle++(0.008,0.008)
  (-0.173,0.909)rectangle++(0.008,0.008) (0.245,0.894)rectangle++(0.008,0.008) (-0.211,0.915)rectangle++(0.008,0.008) (-0.01,0.382)rectangle++(0.008,0.008) (-0.007,0.758)rectangle++(0.008,0.008)
  (-0.277,0.763)rectangle++(0.008,0.008) (-0.403,0.669)rectangle++(0.008,0.008) (-0.172,0.879)rectangle++(0.008,0.008) (-0.074,0.967)rectangle++(0.008,0.008) (0.31,0.885)rectangle++(0.008,0.008)
  (0.457,0.71)rectangle++(0.008,0.008) (0.462,0.578)rectangle++(0.008,0.008) (0.417,0.516)rectangle++(0.008,0.008) (0.009,0.214)rectangle++(0.008,0.008) (-0.076,0.52)rectangle++(0.008,0.008)
  (0.15,0.681)rectangle++(0.008,0.008) (-0.178,0.784)rectangle++(0.008,0.008) (0.198,0.839)rectangle++(0.008,0.008) (-0.213,0.873)rectangle++(0.008,0.008) (0.214,0.891)rectangle++(0.008,0.008)
  (0.415,0.747)rectangle++(0.008,0.008) (0.456,0.61)rectangle++(0.008,0.008) (-0.013,0.861)rectangle++(0.008,0.008) (-0.008,0.959)rectangle++(0.008,0.008) (-0.006,1.001)rectangle++(0.008,0.008)
  (0.354,0.876)rectangle++(0.008,0.008) (0.474,0.69)rectangle++(0.008,0.008) (0.462,0.563)rectangle++(0.008,0.008) (0.192,0.834)rectangle++(0.008,0.008) (0.078,0.948)rectangle++(0.008,0.008)
  (0.031,0.996)rectangle++(0.008,0.008) (0.369,0.861)rectangle++(0.008,0.008) (0.475,0.678)rectangle++(0.008,0.008) (-0.029,0.899)rectangle++(0.008,0.008) (-0.014,0.975)rectangle++(0.008,0.008)
  (-0.358,0.86)rectangle++(0.008,0.008) (-0.474,0.684)rectangle++(0.008,0.008) (-0.202,0.885)rectangle++(0.008,0.008) (-0.087,0.969)rectangle++(0.008,0.008) (-0.39,0.831)rectangle++(0.008,0.008)
  (0.118,0.935)rectangle++(0.008,0.008) (0.387,0.802)rectangle++(0.008,0.008) (0.463,0.645)rectangle++(0.008,0.008) (-0.022,0.879)rectangle++(0.008,0.008) (-0.005,0.367)rectangle++(0.008,0.008)
  (0.128,0.586)rectangle++(0.008,0.008) (0.052,0.844)rectangle++(0.008,0.008) (-0.281,0.823)rectangle++(0.008,0.008) (-0.12,0.943)rectangle++(0.008,0.008) (-0.396,0.807)rectangle++(0.008,0.008)
  (-0.168,0.937)rectangle++(0.008,0.008) (-0.073,0.991)rectangle++(0.008,0.008) (-0.391,0.846)rectangle++(0.008,0.008) (-0.167,0.953)rectangle++(0.008,0.008) (0.263,0.911)rectangle++(0.008,0.008)
  (-0.209,0.93)rectangle++(0.008,0.008) (-0.01,0.388)rectangle++(0.008,0.008) (0.133,0.597)rectangle++(0.008,0.008) (-0.156,0.74)rectangle++(0.008,0.008) (-0.339,0.702)rectangle++(0.008,0.008)
  (-0.145,0.892)rectangle++(0.008,0.008) (-0.063,0.972)rectangle++(0.008,0.008) (-0.38,0.841)rectangle++(0.008,0.008) (-0.478,0.668)rectangle++(0.008,0.008) (-0.018,0.278)rectangle++(0.008,0.008)
  (-0.004,0.115)rectangle++(0.008,0.008) (0.038,0.47)rectangle++(0.008,0.008) (-0.154,0.648)rectangle++(0.008,0.008) (0.16,0.768)rectangle++(0.008,0.008) (0.065,0.92)rectangle++(0.008,0.008)
  (0.357,0.814)rectangle++(0.008,0.008) (0.453,0.661)rectangle++(0.008,0.008) (-0.032,0.883)rectangle++(0.008,0.008) (-0.005,0.369)rectangle++(0.008,0.008) (-0.004,0.752)rectangle++(0.008,0.008)
  (-0.274,0.762)rectangle++(0.008,0.008) (-0.012,0.318)rectangle++(0.008,0.008) (-0.123,0.56)rectangle++(0.008,0.008) (-0.26,0.631)rectangle++(0.008,0.008) (-0.348,0.615)rectangle++(0.008,0.008)
  (-0.148,0.856)rectangle++(0.008,0.008) (-0.065,0.957)rectangle++(0.008,0.008) (-0.029,1)rectangle++(0.008,0.008) (0.343,0.884)rectangle++(0.008,0.008) (-0.163,0.946)rectangle++(0.008,0.008)
  (0.263,0.907)rectangle++(0.008,0.008) (0.443,0.737)rectangle++(0.008,0.008) (0.465,0.595)rectangle++(0.008,0.008) (-0.004,0.857)rectangle++(0.008,0.008) (-0.311,0.81)rectangle++(0.008,0.008)
  (-0.435,0.678)rectangle++(0.008,0.008) (0.043,0.882)rectangle++(0.008,0.008) (0.016,0.968)rectangle++(0.008,0.008) (-0.003,0.404)rectangle++(0.008,0.008) (-0.004,0.767)rectangle++(0.008,0.008)
  (0.272,0.769)rectangle++(0.008,0.008) (0.398,0.671)rectangle++(0.008,0.008) (0.421,0.581)rectangle++(0.008,0.008) (-0.019,0.835)rectangle++(0.008,0.008) (-0.31,0.794)rectangle++(0.008,0.008)
  (-0.429,0.671)rectangle++(0.008,0.008) (-0.439,0.573)rectangle++(0.008,0.008) (-0.409,0.524)rectangle++(0.008,0.008) (-0.001,0.802)rectangle++(0.008,0.008) (-0.29,0.786)rectangle++(0.008,0.008)
  (-0.013,0.328)rectangle++(0.008,0.008) (-0.126,0.565)rectangle++(0.008,0.008) (-0.055,0.835)rectangle++(0.008,0.008) (-0.327,0.781)rectangle++(0.008,0.008) (-0.14,0.926)rectangle++(0.008,0.008)
  (-0.061,0.986)rectangle++(0.008,0.008) (-0.028,1.012)rectangle++(0.008,0.008) (0.348,0.889)rectangle++(0.008,0.008) (0.007,0.371)rectangle++(0.008,0.008) (0.135,0.584)rectangle++(0.008,0.008)
  (-0.15,0.734)rectangle++(0.008,0.008) (-0.065,0.906)rectangle++(0.008,0.008) (-0.006,0.378)rectangle++(0.008,0.008) (-0.141,0.59)rectangle++(0.008,0.008) (-0.062,0.845)rectangle++(0.008,0.008)
  (-0.334,0.784)rectangle++(0.008,0.008) (0.127,0.894)rectangle++(0.008,0.008) (0.376,0.78)rectangle++(0.008,0.008) (0.45,0.638)rectangle++(0.008,0.008) (-0.026,0.872)rectangle++(0.008,0.008)
  (-0.005,0.364)rectangle++(0.008,0.008) (0.127,0.584)rectangle++(0.008,0.008) (-0.154,0.732)rectangle++(0.008,0.008) (0.19,0.806)rectangle++(0.008,0.008) (-0.205,0.855)rectangle++(0.008,0.008)
  (-0.403,0.737)rectangle++(0.008,0.008) (-0.171,0.907)rectangle++(0.008,0.008) (0.245,0.892)rectangle++(0.008,0.008) (0.101,0.972)rectangle++(0.008,0.008) (0.392,0.825)rectangle++(0.008,0.008)
  (0.163,0.944)rectangle++(0.008,0.008) (0.001,0.394)rectangle++(0.008,0.008) (-0.002,0.763)rectangle++(0.008,0.008) (0.271,0.766)rectangle++(0.008,0.008) (-0.153,0.866)rectangle++(0.008,0.008)
  (-0.067,0.961)rectangle++(0.008,0.008) (-0.377,0.835)rectangle++(0.008,0.008) (0.125,0.933)rectangle++(0.008,0.008) (-0.28,0.89)rectangle++(0.008,0.008) (0.189,0.923)rectangle++(0.008,0.008)
  (0.002,0.385)rectangle++(0.008,0.008) (0.138,0.592)rectangle++(0.008,0.008) (0.056,0.846)rectangle++(0.008,0.008) (0.327,0.783)rectangle++(0.008,0.008) (0.006,0.327)rectangle++(0.008,0.008)
  (-0.118,0.571)rectangle++(0.008,0.008) (0.149,0.719)rectangle++(0.008,0.008) (0.324,0.692)rectangle++(0.008,0.008) (0.134,0.888)rectangle++(0.008,0.008) (0,0.371)rectangle++(0.008,0.008)
  (-0.002,0.753)rectangle++(0.008,0.008) (-0.274,0.763)rectangle++(0.008,0.008) (-0.117,0.918)rectangle++(0.008,0.008) (-0.052,0.983)rectangle++(0.008,0.008) (-0.024,1.011)rectangle++(0.008,0.008)
  (-0.005,0.422)rectangle++(0.008,0.008) (-0.156,0.611)rectangle++(0.008,0.008) (-0.068,0.854)rectangle++(0.008,0.008) (0.273,0.831)rectangle++(0.008,0.008) (0.112,0.947)rectangle++(0.008,0.008)
  (0.045,0.995)rectangle++(0.008,0.008) (0.375,0.855)rectangle++(0.008,0.008) (0.007,0.357)rectangle++(0.008,0.008) (0.001,0.748)rectangle++(0.008,0.008) (-0.004,0.312)rectangle++(0.008,0.008)
  (-0.117,0.56)rectangle++(0.008,0.008) (-0.257,0.634)rectangle++(0.008,0.008) (-0.11,0.864)rectangle++(0.008,0.008) (-0.049,0.96)rectangle++(0.008,0.008) (-0.023,1.001)rectangle++(0.008,0.008)
  (-0.371,0.869)rectangle++(0.008,0.008) (-0.158,0.963)rectangle++(0.008,0.008) (-0.42,0.803)rectangle++(0.008,0.008) (0.094,0.933)rectangle++(0.008,0.008) (0.037,0.99)rectangle++(0.008,0.008)
  (0.37,0.855)rectangle++(0.008,0.008) (0.474,0.675)rectangle++(0.008,0.008) (0.457,0.556)rectangle++(0.008,0.008) (0.407,0.507)rectangle++(0.008,0.008) (0.001,0.796)rectangle++(0.008,0.008)
  (-0.002,0.932)rectangle++(0.008,0.008) (0.331,0.843)rectangle++(0.008,0.008) (-0.153,0.923)rectangle++(0.008,0.008) (0.259,0.893)rectangle++(0.008,0.008) (-0.204,0.92)rectangle++(0.008,0.008)
  (-0.088,0.984)rectangle++(0.008,0.008) (-0.039,1.011)rectangle++(0.008,0.008) (-0.383,0.867)rectangle++(0.008,0.008) (-0.488,0.679)rectangle++(0.008,0.008) (0.019,0.901)rectangle++(0.008,0.008)
  (0.33,0.822)rectangle++(0.008,0.008) (-0.146,0.912)rectangle++(0.008,0.008) (0.258,0.886)rectangle++(0.008,0.008) (0.106,0.97)rectangle++(0.008,0.008) (0.042,1.005)rectangle++(0.008,0.008)
  (-0.343,0.894)rectangle++(0.008,0.008) (-0.014,0.373)rectangle++(0.008,0.008) (-0.008,0.754)rectangle++(0.008,0.008) (-0.006,0.915)rectangle++(0.008,0.008) (-0.004,0.382)rectangle++(0.008,0.008)
  (0.134,0.592)rectangle++(0.008,0.008) (0.272,0.639)rectangle++(0.008,0.008) (0.352,0.612)rectangle++(0.008,0.008) (-0.061,0.824)rectangle++(0.008,0.008) (0.266,0.815)rectangle++(0.008,0.008)
  (-0.173,0.887)rectangle++(0.008,0.008) (0.237,0.883)rectangle++(0.008,0.008) (-0.211,0.908)rectangle++(0.008,0.008) (-0.424,0.759)rectangle++(0.008,0.008) (0.076,0.915)rectangle++(0.008,0.008)
  (0.361,0.807)rectangle++(0.008,0.008) (0.452,0.656)rectangle++(0.008,0.008) (0.44,0.555)rectangle++(0.008,0.008) (0,0.83)rectangle++(0.008,0.008) (-0.003,0.946)rectangle++(0.008,0.008)
  (0.336,0.85)rectangle++(0.008,0.008) (-0.153,0.928)rectangle++(0.008,0.008) (0.26,0.895)rectangle++(0.008,0.008) (0.438,0.733)rectangle++(0.008,0.008) (0.461,0.595)rectangle++(0.008,0.008)
  (0.422,0.524)rectangle++(0.008,0.008) (0.379,0.505)rectangle++(0.008,0.008) (-0.011,0.785)rectangle++(0.008,0.008) (-0.289,0.774)rectangle++(0.008,0.008) (-0.412,0.67)rectangle++(0.008,0.008)
  (-0.175,0.879)rectangle++(0.008,0.008) (-0.009,0.367)rectangle++(0.008,0.008) (-0.006,0.752)rectangle++(0.008,0.008) (-0.275,0.761)rectangle++(0.008,0.008) (-0.401,0.669)rectangle++(0.008,0.008)
  (-0.016,0.279)rectangle++(0.008,0.008) (-0.009,0.715)rectangle++(0.008,0.008) (-0.006,0.898)rectangle++(0.008,0.008) (-0.327,0.827)rectangle++(0.008,0.008) (-0.449,0.681)rectangle++(0.008,0.008)
  (0.037,0.888)rectangle++(0.008,0.008) (0.333,0.809)rectangle++(0.008,0.008) (0.44,0.667)rectangle++(0.008,0.008) (0.439,0.564)rectangle++(0.008,0.008) (0.182,0.835)rectangle++(0.008,0.008)
  (0.38,0.733)rectangle++(0.008,0.008) (0.008,0.305)rectangle++(0.008,0.008) (-0.109,0.562)rectangle++(0.008,0.008) (-0.254,0.637)rectangle++(0.008,0.008) (-0.347,0.62)rectangle++(0.008,0.008)
  (-0.014,0.258)rectangle++(0.008,0.008) (-0.102,0.532)rectangle++(0.008,0.008) (-0.24,0.626)rectangle++(0.008,0.008) (0.113,0.788)rectangle++(0.008,0.008) (0.332,0.736)rectangle++(0.008,0.008)
  (0.006,0.307)rectangle++(0.008,0.008) (0,0.727)rectangle++(0.008,0.008) (0.259,0.748)rectangle++(0.008,0.008) (0.385,0.666)rectangle++(0.008,0.008) (0.159,0.877)rectangle++(0.008,0.008)
  (0.385,0.761)rectangle++(0.008,0.008) (0.008,0.317)rectangle++(0.008,0.008) (-0.113,0.567)rectangle++(0.008,0.008) (0.15,0.716)rectangle++(0.008,0.008) (-0.191,0.8)rectangle++(0.008,0.008)
  (-0.376,0.717)rectangle++(0.008,0.008) (-0.16,0.899)rectangle++(0.008,0.008) (-0.398,0.773)rectangle++(0.008,0.008) (0.093,0.912)rectangle++(0.008,0.008) (-0.286,0.869)rectangle++(0.008,0.008)
  (-0.123,0.963)rectangle++(0.008,0.008) (0.287,0.9)rectangle++(0.008,0.008) (-0.194,0.933)rectangle++(0.008,0.008) (-0.425,0.777)rectangle++(0.008,0.008) (0.082,0.923)rectangle++(0.008,0.008)
  (-0.296,0.871)rectangle++(0.008,0.008) (-0.127,0.963)rectangle++(0.008,0.008) (-0.405,0.814)rectangle++(0.008,0.008) (-0.48,0.647)rectangle++(0.008,0.008) (-0.204,0.869)rectangle++(0.008,0.008)
  (0.217,0.886)rectangle++(0.008,0.008) (-0.221,0.902)rectangle++(0.008,0.008) (-0.427,0.752)rectangle++(0.008,0.008) (-0.467,0.611)rectangle++(0.008,0.008) (-0.199,0.854)rectangle++(0.008,0.008)
  (-0.01,0.356)rectangle++(0.008,0.008) (-0.135,0.579)rectangle++(0.008,0.008) (-0.272,0.635)rectangle++(0.008,0.008) (-0.355,0.612)rectangle++(0.008,0.008) (-0.151,0.855)rectangle++(0.008,0.008)
  (-0.378,0.756)rectangle++(0.008,0.008) (-0.015,0.315)rectangle++(0.008,0.008) (-0.123,0.558)rectangle++(0.008,0.008) (-0.054,0.832)rectangle++(0.008,0.008) (0.272,0.816)rectangle++(0.008,0.008)
  (-0.171,0.889)rectangle++(0.008,0.008) (0.239,0.884)rectangle++(0.008,0.008) (0.424,0.735)rectangle++(0.008,0.008) (0.456,0.601)rectangle++(0.008,0.008) (0.189,0.85)rectangle++(0.008,0.008)
  (0.389,0.737)rectangle++(0.008,0.008) (0.161,0.907)rectangle++(0.008,0.008) (0.065,0.979)rectangle++(0.008,0.008) (0.379,0.84)rectangle++(0.008,0.008) (0.157,0.951)rectangle++(0.008,0.008)
  (0.063,0.997)rectangle++(0.008,0.008) (-0.002,0.416)rectangle++(0.008,0.008) (0.147,0.607)rectangle++(0.008,0.008) (0.283,0.642)rectangle++(0.008,0.008) (0.358,0.608)rectangle++(0.008,0.008)
  (0.38,0.567)rectangle++(0.008,0.008) (-0.032,0.814)rectangle++(0.008,0.008) (-0.309,0.78)rectangle++(0.008,0.008) (-0.423,0.665)rectangle++(0.008,0.008) (-0.435,0.572)rectangle++(0.008,0.008)
  (-0.407,0.526)rectangle++(0.008,0.008) (-0.173,0.818)rectangle++(0.008,0.008) (0.212,0.852)rectangle++(0.008,0.008) (0.002,0.356)rectangle++(0.008,0.008) (0.127,0.578)rectangle++(0.008,0.008)
  (0.264,0.635)rectangle++(0.008,0.008) (0.347,0.612)rectangle++(0.008,0.008) (0.376,0.573)rectangle++(0.008,0.008) (-0.036,0.815)rectangle++(0.008,0.008) (-0.017,0.94)rectangle++(0.008,0.008)
  (-0.01,0.993)rectangle++(0.008,0.008) (0.349,0.873)rectangle++(0.008,0.008) (0.007,0.365)rectangle++(0.008,0.008) (-0.131,0.588)rectangle++(0.008,0.008) (-0.273,0.641)rectangle++(0.008,0.008)
  (-0.117,0.867)rectangle++(0.008,0.008) (-0.007,0.362)rectangle++(0.008,0.008) (-0.004,0.15)rectangle++(0.008,0.008) (0.051,0.486)rectangle++(0.008,0.008) (-0.002,0.202)rectangle++(0.008,0.008)
  (0.07,0.51)rectangle++(0.008,0.008) (0.213,0.624)rectangle++(0.008,0.008) (0.087,0.86)rectangle++(0.008,0.008) (0.034,0.959)rectangle++(0.008,0.008) (-0.33,0.87)rectangle++(0.008,0.008)
  (0.159,0.932)rectangle++(0.008,0.008) (0.064,0.989)rectangle++(0.008,0.008) (0.382,0.845)rectangle++(0.008,0.008) (0.158,0.953)rectangle++(0.008,0.008) (0.064,0.998)rectangle++(0.008,0.008)
  (-0.331,0.898)rectangle++(0.008,0.008) (0.169,0.945)rectangle++(0.008,0.008) (0.414,0.788)rectangle++(0.008,0.008) (0.009,0.329)rectangle++(0.008,0.008) (-0.117,0.573)rectangle++(0.008,0.008)
  (-0.051,0.838)rectangle++(0.008,0.008) (0.275,0.818)rectangle++(0.008,0.008) (0.417,0.692)rectangle++(0.008,0.008) (-0.06,0.884)rectangle++(0.008,0.008) (-0.028,0.969)rectangle++(0.008,0.008)
  (-0.362,0.852)rectangle++(0.008,0.008) (-0.015,0.356)rectangle++(0.008,0.008) (-0.137,0.577)rectangle++(0.008,0.008) (-0.06,0.84)rectangle++(0.008,0.008) (-0.028,0.95)rectangle++(0.008,0.008)
  (0.326,0.861)rectangle++(0.008,0.008) (0.456,0.693)rectangle++(0.008,0.008) (0.01,0.289)rectangle++(0.008,0.008) (0.107,0.545)rectangle++(0.008,0.008) (0.043,0.827)rectangle++(0.008,0.008)
  (0.016,0.945)rectangle++(0.008,0.008) (0.344,0.843)rectangle++(0.008,0.008) (-0.147,0.927)rectangle++(0.008,0.008) (0.263,0.893)rectangle++(0.008,0.008) (0.004,0.373)rectangle++(0.008,0.008)
  (-0.001,0.754)rectangle++(0.008,0.008) (-0.003,0.914)rectangle++(0.008,0.008) (0.325,0.835)rectangle++(0.008,0.008) (-0.153,0.917)rectangle++(0.008,0.008) (-0.401,0.784)rectangle++(0.008,0.008)
  (-0.467,0.634)rectangle++(0.008,0.008) (0.012,0.873)rectangle++(0.008,0.008) (-0.31,0.823)rectangle++(0.008,0.008) (-0.132,0.943)rectangle++(0.008,0.008) (-0.058,0.994)rectangle++(0.008,0.008)
  (-0.027,1.015)rectangle++(0.008,0.008) (-0.014,1.024)rectangle++(0.008,0.008) (0.359,0.889)rectangle++(0.008,0.008) (0.148,0.971)rectangle++(0.008,0.008) (-0.283,0.916)rectangle++(0.008,0.008)
  (-0.46,0.737)rectangle++(0.008,0.008) (-0.195,0.907)rectangle++(0.008,0.008) (-0.417,0.764)rectangle++(0.008,0.008) (-0.467,0.62)rectangle++(0.008,0.008) (0.007,0.866)rectangle++(0.008,0.008)
  (0.001,0.962)rectangle++(0.008,0.008) (0.343,0.856)rectangle++(0.008,0.008) (0.462,0.685)rectangle++(0.008,0.008) (-0.037,0.897)rectangle++(0.008,0.008) (-0.018,0.974)rectangle++(0.008,0.008)
  (-0.004,0.407)rectangle++(0.008,0.008) (-0.004,0.769)rectangle++(0.008,0.008) (-0.28,0.769)rectangle++(0.008,0.008) (-0.12,0.921)rectangle++(0.008,0.008) (-0.387,0.797)rectangle++(0.008,0.008)
  (0.107,0.919)rectangle++(0.008,0.008) (0.376,0.798)rectangle++(0.008,0.008) (0.456,0.647)rectangle++(0.008,0.008) (-0.026,0.878)rectangle++(0.008,0.008) (0.301,0.827)rectangle++(0.008,0.008)
  (0.124,0.945)rectangle++(0.008,0.008) (0.05,0.995)rectangle++(0.008,0.008) (0.019,1.015)rectangle++(0.008,0.008) (0.37,0.874)rectangle++(0.008,0.008) (-0.146,0.951)rectangle++(0.008,0.008)
  (-0.064,0.997)rectangle++(0.008,0.008) (-0.029,1.016)rectangle++(0.008,0.008) (-0.38,0.873)rectangle++(0.008,0.008) (0.138,0.951)rectangle++(0.008,0.008) (0.402,0.802)rectangle++(0.008,0.008)
  (0.166,0.935)rectangle++(0.008,0.008) (0.001,0.39)rectangle++(0.008,0.008) (-0.002,0.762)rectangle++(0.008,0.008) (0.27,0.765)rectangle++(0.008,0.008) (0.396,0.67)rectangle++(0.008,0.008)
  (0.419,0.581)rectangle++(0.008,0.008) (-0.019,0.835)rectangle++(0.008,0.008) (-0.31,0.794)rectangle++(0.008,0.008) (-0.133,0.931)rectangle++(0.008,0.008) (0.271,0.889)rectangle++(0.008,0.008)
  (0.441,0.726)rectangle++(0.008,0.008) (0.46,0.591)rectangle++(0.008,0.008) (-0.004,0.854)rectangle++(0.008,0.008) (0.302,0.808)rectangle++(0.008,0.008) (-0.154,0.896)rectangle++(0.008,0.008)
  (0.249,0.881)rectangle++(0.008,0.008) (0.004,0.368)rectangle++(0.008,0.008) (-0.001,0.752)rectangle++(0.008,0.008) (0.267,0.76)rectangle++(0.008,0.008) (0.004,0.317)rectangle++(0.008,0.008)
  (-0.115,0.566)rectangle++(0.008,0.008) (-0.007,0.235)rectangle++(0.008,0.008) (-0.091,0.524)rectangle++(0.008,0.008) (-0.041,0.818)rectangle++(0.008,0.008) (0.273,0.805)rectangle++(0.008,0.008)
  (-0.166,0.884)rectangle++(0.008,0.008) (-0.395,0.764)rectangle++(0.008,0.008) (-0.016,0.319)rectangle++(0.008,0.008) (0.106,0.568)rectangle++(0.008,0.008) (0.25,0.638)rectangle++(0.008,0.008)
  (-0.117,0.8)rectangle++(0.008,0.008) (0.232,0.824)rectangle++(0.008,0.008) (-0.192,0.879)rectangle++(0.008,0.008) (-0.083,0.967)rectangle++(0.008,0.008) (-0.037,1.004)rectangle++(0.008,0.008)
  (-0.379,0.865)rectangle++(0.008,0.008) (0.135,0.947)rectangle++(0.008,0.008) (0.054,0.995)rectangle++(0.008,0.008) (-0.334,0.894)rectangle++(0.008,0.008) (-0.143,0.973)rectangle++(0.008,0.008)
  (-0.416,0.813)rectangle++(0.008,0.008) (-0.177,0.939)rectangle++(0.008,0.008) (-0.077,0.992)rectangle++(0.008,0.008) (0.318,0.897)rectangle++(0.008,0.008) (-0.178,0.943)rectangle++(0.008,0.008)
  (0.254,0.911)rectangle++(0.008,0.008) (0.004,0.38)rectangle++(0.008,0.008) (0.137,0.589)rectangle++(0.008,0.008) (0.055,0.845)rectangle++(0.008,0.008) (-0.28,0.825)rectangle++(0.008,0.008)
  (-0.426,0.696)rectangle++(0.008,0.008) (0.053,0.887)rectangle++(0.008,0.008) (-0.296,0.844)rectangle++(0.008,0.008) (-0.44,0.699)rectangle++(0.008,0.008) (-0.454,0.582)rectangle++(0.008,0.008)
  (-0.001,0.844)rectangle++(0.008,0.008) (-0.305,0.805)rectangle++(0.008,0.008) (-0.431,0.678)rectangle++(0.008,0.008) (-0.183,0.882)rectangle++(0.008,0.008) (-0.402,0.757)rectangle++(0.008,0.008)
  (-0.016,0.316)rectangle++(0.008,0.008) (-0.009,0.73)rectangle++(0.008,0.008) (0.256,0.753)rectangle++(0.008,0.008) (-0.156,0.855)rectangle++(0.008,0.008) (-0.009,0.357)rectangle++(0.008,0.008)
  (-0.006,0.748)rectangle++(0.008,0.008) (0.263,0.76)rectangle++(0.008,0.008) (-0.155,0.861)rectangle++(0.008,0.008) (-0.009,0.359)rectangle++(0.008,0.008) (-0.006,0.749)rectangle++(0.008,0.008)
  (-0.005,0.912)rectangle++(0.008,0.008) (-0.331,0.834)rectangle++(0.008,0.008) (0.146,0.916)rectangle++(0.008,0.008) (0.393,0.783)rectangle++(0.008,0.008) (-0.104,0.917)rectangle++(0.008,0.008)
  (0.28,0.873)rectangle++(0.008,0.008) (0.439,0.715)rectangle++(0.008,0.008) (0.182,0.898)rectangle++(0.008,0.008) (0.074,0.975)rectangle++(0.008,0.008) (0.381,0.836)rectangle++(0.008,0.008)
  (-0.128,0.937)rectangle++(0.008,0.008) (-0.008,0.391)rectangle++(0.008,0.008) (-0.004,0.162)rectangle++(0.008,0.008) (0.055,0.492)rectangle++(0.008,0.008) (0.021,0.804)rectangle++(0.008,0.008)
  (0.006,0.935)rectangle++(0.008,0.008) (-0.335,0.849)rectangle++(0.008,0.008) (-0.014,0.354)rectangle++(0.008,0.008) (0.119,0.583)rectangle++(0.008,0.008) (0.262,0.64)rectangle++(0.008,0.008)
  (0.108,0.867)rectangle++(0.008,0.008) (0.358,0.774)rectangle++(0.008,0.008) (0.439,0.642)rectangle++(0.008,0.008) (0.182,0.867)rectangle++(0.008,0.008) (-0.23,0.881)rectangle++(0.008,0.008)
  (-0.099,0.968)rectangle++(0.008,0.008) (-0.044,1.004)rectangle++(0.008,0.008) (-0.382,0.862)rectangle++(0.008,0.008) (-0.486,0.677)rectangle++(0.008,0.008) (-0.207,0.882)rectangle++(0.008,0.008)
  (0.22,0.893)rectangle++(0.008,0.008) (0.09,0.973)rectangle++(0.008,0.008) (-0.31,0.896)rectangle++(0.008,0.008) (0.178,0.937)rectangle++(0.008,0.008) (0.415,0.781)rectangle++(0.008,0.008)
  (0.468,0.625)rectangle++(0.008,0.008) (0.436,0.535)rectangle++(0.008,0.008) (0.005,0.82)rectangle++(0.008,0.008) (0,0.942)rectangle++(0.008,0.008) (-0.004,0.393)rectangle++(0.008,0.008)
  (-0.146,0.598)rectangle++(0.008,0.008) (-0.284,0.64)rectangle++(0.008,0.008) (-0.362,0.611)rectangle++(0.008,0.008) (0.052,0.825)rectangle++(0.008,0.008) (0.019,0.944)rectangle++(0.008,0.008)
  (-0.332,0.858)rectangle++(0.008,0.008) (-0.462,0.693)rectangle++(0.008,0.008) (-0.196,0.889)rectangle++(0.008,0.008) (-0.085,0.971)rectangle++(0.008,0.008) (-0.038,1.005)rectangle++(0.008,0.008)
  (-0.005,0.42)rectangle++(0.008,0.008) (-0.004,0.774)rectangle++(0.008,0.008) (-0.004,0.923)rectangle++(0.008,0.008) (-0.004,0.985)rectangle++(0.008,0.008) (-0.004,1.011)rectangle++(0.008,0.008)
  (-0.004,1.023)rectangle++(0.008,0.008) (0.363,0.885)rectangle++(0.008,0.008) (0.15,0.969)rectangle++(0.008,0.008) (0.061,1.005)rectangle++(0.008,0.008) (0.023,1.02)rectangle++(0.008,0.008)
  (-0.003,0.426)rectangle++(0.008,0.008) (-0.004,0.177)rectangle++(0.008,0.008) (0.06,0.499)rectangle++(0.008,0.008) (0.205,0.623)rectangle++(0.008,0.008) (-0.132,0.777)rectangle++(0.008,0.008)
  (-0.058,0.924)rectangle++(0.008,0.008) (-0.027,0.986)rectangle++(0.008,0.008) (-0.368,0.86)rectangle++(0.008,0.008) (0.138,0.941)rectangle++(0.008,0.008) (0.056,0.993)rectangle++(0.008,0.008)
  (-0.333,0.893)rectangle++(0.008,0.008) (-0.474,0.709)rectangle++(0.008,0.008) (-0.202,0.895)rectangle++(0.008,0.008) (-0.415,0.756)rectangle++(0.008,0.008) (-0.464,0.617)rectangle++(0.008,0.008)
  (-0.197,0.857)rectangle++(0.008,0.008) (-0.4,0.74)rectangle++(0.008,0.008) (0.081,0.897)rectangle++(0.008,0.008) (-0.287,0.858)rectangle++(0.008,0.008) (-0.123,0.958)rectangle++(0.008,0.008)
  (0.285,0.898)rectangle++(0.008,0.008) (0.117,0.975)rectangle++(0.008,0.008) (0.047,1.007)rectangle++(0.008,0.008) (0.38,0.86)rectangle++(0.008,0.008) (0.48,0.674)rectangle++(0.008,0.008)
  (0.199,0.881)rectangle++(0.008,0.008) (-0.227,0.893)rectangle++(0.008,0.008) (-0.098,0.973)rectangle++(0.008,0.008) (-0.396,0.829)rectangle++(0.008,0.008) (-0.168,0.946)rectangle++(0.008,0.008)
  (0.26,0.909)rectangle++(0.008,0.008) (-0.209,0.927)rectangle++(0.008,0.008) (0.235,0.915)rectangle++(0.008,0.008) (0.096,0.982)rectangle++(0.008,0.008) (0.394,0.831)rectangle++(0.008,0.008)
  (0.163,0.947)rectangle++(0.008,0.008) (0.412,0.791)rectangle++(0.008,0.008) (-0.098,0.928)rectangle++(0.008,0.008) (-0.043,0.987)rectangle++(0.008,0.008) (0.332,0.883)rectangle++(0.008,0.008)
  (-0.167,0.941)rectangle++(0.008,0.008) (-0.416,0.79)rectangle++(0.008,0.008) (0.091,0.926)rectangle++(0.008,0.008) (-0.293,0.875)rectangle++(0.008,0.008) (0.178,0.921)rectangle++(0.008,0.008)
  (0.409,0.774)rectangle++(0.008,0.008) (0.463,0.624)rectangle++(0.008,0.008) (0.192,0.86)rectangle++(0.008,0.008) (-0.223,0.881)rectangle++(0.008,0.008) (-0.011,0.368)rectangle++(0.008,0.008)
  (-0.007,0.752)rectangle++(0.008,0.008) (-0.005,0.914)rectangle++(0.008,0.008) (-0.004,0.981)rectangle++(0.008,0.008) (-0.356,0.866)rectangle++(0.008,0.008) (-0.476,0.688)rectangle++(0.008,0.008)
  (-0.018,0.287)rectangle++(0.008,0.008) (-0.01,0.718)rectangle++(0.008,0.008) (-0.265,0.744)rectangle++(0.008,0.008) (-0.012,0.31)rectangle++(0.008,0.008) (0.105,0.562)rectangle++(0.008,0.008)
  (0.042,0.834)rectangle++(0.008,0.008) (0.015,0.948)rectangle++(0.008,0.008) (0.345,0.844)rectangle++(0.008,0.008) (0.142,0.952)rectangle++(0.008,0.008) (0.404,0.801)rectangle++(0.008,0.008)
  (0.168,0.934)rectangle++(0.008,0.008) (-0.261,0.906)rectangle++(0.008,0.008) (0.204,0.924)rectangle++(0.008,0.008) (0.083,0.986)rectangle++(0.008,0.008) (-0.317,0.899)rectangle++(0.008,0.008)
  (0.175,0.941)rectangle++(0.008,0.008) (0.001,0.393)rectangle++(0.008,0.008) (0.14,0.595)rectangle++(0.008,0.008) (-0.152,0.741)rectangle++(0.008,0.008) (0.194,0.81)rectangle++(0.008,0.008)
  (0.377,0.717)rectangle++(0.008,0.008) (0.428,0.609)rectangle++(0.008,0.008) (0.412,0.542)rectangle++(0.008,0.008) (-0.009,0.814)rectangle++(0.008,0.008) (-0.004,0.34)rectangle++(0.008,0.008)
  (-0.127,0.573)rectangle++(0.008,0.008) (0.146,0.724)rectangle++(0.008,0.008) (0.324,0.695)rectangle++(0.008,0.008) (-0.104,0.853)rectangle++(0.008,0.008) (0.256,0.843)rectangle++(0.008,0.008)
  (0.418,0.71)rectangle++(0.008,0.008) (0.444,0.592)rectangle++(0.008,0.008) (-0.012,0.848)rectangle++(0.008,0.008) (-0.007,0.954)rectangle++(0.008,0.008) (-0.004,0.398)rectangle++(0.008,0.008)
  (-0.148,0.6)rectangle++(0.008,0.008) (-0.285,0.641)rectangle++(0.008,0.008) (0.098,0.811)rectangle++(0.008,0.008) (0.334,0.752)rectangle++(0.008,0.008) (0.42,0.641)rectangle++(0.008,0.008)
  (0.009,0.267)rectangle++(0.008,0.008) (0.099,0.535)rectangle++(0.008,0.008) (-0.15,0.699)rectangle++(0.008,0.008) (-0.322,0.685)rectangle++(0.008,0.008) (0.097,0.844)rectangle++(0.008,0.008);

          \begin{scope}[shift={(1.9,0)}]
            % Arbre fractal : les quatre pièces w_i(A)
% Fichier généré par generer.py : ne pas modifier à la main.
\fill[piece1]
  (-0.008,0.371)rectangle++(0.008,0.008) (-0.004,0.153)rectangle++(0.008,0.008) (0.003,0.3)rectangle++(0.008,0.008) (0,0.367)rectangle++(0.008,0.008) (-0.008,0.377)rectangle++(0.008,0.008)
  (0.007,0.357)rectangle++(0.008,0.008) (0.001,0.381)rectangle++(0.008,0.008) (-0.001,0.338)rectangle++(0.008,0.008) (-0.004,0.405)rectangle++(0.008,0.008) (-0.017,0.274)rectangle++(0.008,0.008)
  (-0.004,0.113)rectangle++(0.008,0.008) (0.006,0.351)rectangle++(0.008,0.008) (0.007,0.336)rectangle++(0.008,0.008) (-0.012,0.382)rectangle++(0.008,0.008) (-0.013,0.333)rectangle++(0.008,0.008)
  (-0.011,0.363)rectangle++(0.008,0.008) (-0.01,0.358)rectangle++(0.008,0.008) (-0.013,0.329)rectangle++(0.008,0.008) (0.003,0.359)rectangle++(0.008,0.008) (-0.009,0.396)rectangle++(0.008,0.008)
  (-0.008,0.249)rectangle++(0.008,0.008) (-0.012,0.376)rectangle++(0.008,0.008) (-0.008,0.244)rectangle++(0.008,0.008) (-0.004,0.1)rectangle++(0.008,0.008) (-0.004,0.04)rectangle++(0.008,0.008)
  (0.001,0.397)rectangle++(0.008,0.008) (-0.017,0.316)rectangle++(0.008,0.008) (0.009,0.311)rectangle++(0.008,0.008) (0.008,0.321)rectangle++(0.008,0.008) (-0.013,0.319)rectangle++(0.008,0.008)
  (-0.016,0.21)rectangle++(0.008,0.008) (-0.004,0.086)rectangle++(0.008,0.008) (0.01,0.294)rectangle++(0.008,0.008) (-0.014,0.287)rectangle++(0.008,0.008) (-0.007,0.406)rectangle++(0.008,0.008)
  (-0.015,0.336)rectangle++(0.008,0.008) (-0.008,0.236)rectangle++(0.008,0.008) (-0.006,0.375)rectangle++(0.008,0.008) (-0.009,0.307)rectangle++(0.008,0.008) (-0.009,0.392)rectangle++(0.008,0.008)
  (-0.008,0.39)rectangle++(0.008,0.008) (-0.012,0.381)rectangle++(0.008,0.008) (-0.014,0.338)rectangle++(0.008,0.008) (-0.004,0.14)rectangle++(0.008,0.008) (-0.012,0.375)rectangle++(0.008,0.008)
  (-0.016,0.338)rectangle++(0.008,0.008) (-0.009,0.389)rectangle++(0.008,0.008) (-0.014,0.361)rectangle++(0.008,0.008) (-0.004,0.312)rectangle++(0.008,0.008) (0.007,0.364)rectangle++(0.008,0.008)
  (-0.005,0.402)rectangle++(0.008,0.008) (-0.014,0.327)rectangle++(0.008,0.008) (0.004,0.38)rectangle++(0.008,0.008) (-0.004,0.157)rectangle++(0.008,0.008) (-0.004,0.064)rectangle++(0.008,0.008)
  (0.01,0.263)rectangle++(0.008,0.008) (-0.002,0.378)rectangle++(0.008,0.008) (-0.017,0.284)rectangle++(0.008,0.008) (-0.002,0.373)rectangle++(0.008,0.008) (-0.015,0.216)rectangle++(0.008,0.008)
  (0.006,0.366)rectangle++(0.008,0.008) (0.004,0.343)rectangle++(0.008,0.008) (-0.004,0.142)rectangle++(0.008,0.008) (-0.013,0.336)rectangle++(0.008,0.008) (0,0.239)rectangle++(0.008,0.008)
  (-0.018,0.246)rectangle++(0.008,0.008) (-0.004,0.292)rectangle++(0.008,0.008) (-0.001,0.38)rectangle++(0.008,0.008) (0.01,0.256)rectangle++(0.008,0.008) (-0.002,0.407)rectangle++(0.008,0.008)
  (0.006,0.216)rectangle++(0.008,0.008) (-0.008,0.391)rectangle++(0.008,0.008) (-0.002,0.362)rectangle++(0.008,0.008) (-0.017,0.29)rectangle++(0.008,0.008) (0.011,0.284)rectangle++(0.008,0.008)
  (0.009,0.307)rectangle++(0.008,0.008) (-0.016,0.34)rectangle++(0.008,0.008) (0,0.372)rectangle++(0.008,0.008) (-0.014,0.338)rectangle++(0.008,0.008) (0.008,0.215)rectangle++(0.008,0.008)
  (0.002,0.385)rectangle++(0.008,0.008) (-0.01,0.381)rectangle++(0.008,0.008) (-0.016,0.323)rectangle++(0.008,0.008) (-0.019,0.279)rectangle++(0.008,0.008) (-0.012,0.31)rectangle++(0.008,0.008)
  (-0.004,0.128)rectangle++(0.008,0.008) (0.001,0.405)rectangle++(0.008,0.008) (-0.004,0.32)rectangle++(0.008,0.008) (0.003,0.373)rectangle++(0.008,0.008) (0.006,0.291)rectangle++(0.008,0.008)
  (0.01,0.298)rectangle++(0.008,0.008) (-0.001,0.228)rectangle++(0.008,0.008) (-0.007,0.41)rectangle++(0.008,0.008) (0.001,0.358)rectangle++(0.008,0.008) (0.006,0.363)rectangle++(0.008,0.008)
  (-0.015,0.338)rectangle++(0.008,0.008) (0,0.389)rectangle++(0.008,0.008) (-0.005,0.398)rectangle++(0.008,0.008) (-0.012,0.376)rectangle++(0.008,0.008) (-0.007,0.387)rectangle++(0.008,0.008)
  (-0.01,0.337)rectangle++(0.008,0.008) (-0.008,0.237)rectangle++(0.008,0.008) (-0.01,0.347)rectangle++(0.008,0.008) (-0.004,0.31)rectangle++(0.008,0.008) (-0.018,0.268)rectangle++(0.008,0.008)
  (-0.003,0.345)rectangle++(0.008,0.008) (-0.018,0.294)rectangle++(0.008,0.008) (0.008,0.317)rectangle++(0.008,0.008) (-0.018,0.286)rectangle++(0.008,0.008) (0.006,0.287)rectangle++(0.008,0.008)
  (-0.004,0.118)rectangle++(0.008,0.008) (-0.006,0.408)rectangle++(0.008,0.008) (-0.016,0.279)rectangle++(0.008,0.008) (-0.015,0.361)rectangle++(0.008,0.008) (-0.004,0.149)rectangle++(0.008,0.008)
  (-0.008,0.392)rectangle++(0.008,0.008) (0.01,0.256)rectangle++(0.008,0.008) (-0.004,0.357)rectangle++(0.008,0.008) (-0.001,0.386)rectangle++(0.008,0.008) (-0.004,0.16)rectangle++(0.008,0.008)
  (-0.004,0.065)rectangle++(0.008,0.008) (-0.009,0.377)rectangle++(0.008,0.008) (-0.008,0.396)rectangle++(0.008,0.008) (-0.004,0.164)rectangle++(0.008,0.008) (0.005,0.35)rectangle++(0.008,0.008)
  (-0.004,0.145)rectangle++(0.008,0.008) (-0.002,0.201)rectangle++(0.008,0.008) (-0.001,0.391)rectangle++(0.008,0.008) (-0.008,0.249)rectangle++(0.008,0.008) (0.008,0.35)rectangle++(0.008,0.008)
  (-0.013,0.357)rectangle++(0.008,0.008) (-0.001,0.335)rectangle++(0.008,0.008) (-0.004,0.308)rectangle++(0.008,0.008) (-0.004,0.303)rectangle++(0.008,0.008) (-0.001,0.332)rectangle++(0.008,0.008)
  (-0.005,0.396)rectangle++(0.008,0.008) (-0.004,0.164)rectangle++(0.008,0.008) (-0.016,0.33)rectangle++(0.008,0.008) (-0.004,0.136)rectangle++(0.008,0.008) (-0.01,0.358)rectangle++(0.008,0.008)
  (-0.004,0.365)rectangle++(0.008,0.008) (0.006,0.312)rectangle++(0.008,0.008) (-0.01,0.369)rectangle++(0.008,0.008) (-0.007,0.361)rectangle++(0.008,0.008) (-0.001,0.388)rectangle++(0.008,0.008)
  (-0.016,0.258)rectangle++(0.008,0.008) (-0.013,0.369)rectangle++(0.008,0.008) (0.002,0.337)rectangle++(0.008,0.008) (-0.008,0.396)rectangle++(0.008,0.008) (0.004,0.375)rectangle++(0.008,0.008)
  (-0.015,0.303)rectangle++(0.008,0.008) (0.007,0.355)rectangle++(0.008,0.008) (-0.011,0.371)rectangle++(0.008,0.008) (-0.009,0.381)rectangle++(0.008,0.008) (-0.008,0.393)rectangle++(0.008,0.008)
  (-0.002,0.41)rectangle++(0.008,0.008) (0.008,0.353)rectangle++(0.008,0.008) (-0.008,0.303)rectangle++(0.008,0.008) (-0.015,0.347)rectangle++(0.008,0.008) (-0.004,0.144)rectangle++(0.008,0.008)
  (0.005,0.325)rectangle++(0.008,0.008) (-0.008,0.403)rectangle++(0.008,0.008) (-0.008,0.25)rectangle++(0.008,0.008) (-0.007,0.386)rectangle++(0.008,0.008) (-0.016,0.315)rectangle++(0.008,0.008)
  (0.005,0.368)rectangle++(0.008,0.008) (-0.009,0.344)rectangle++(0.008,0.008) (-0.008,0.239)rectangle++(0.008,0.008) (-0.004,0.353)rectangle++(0.008,0.008) (-0.003,0.391)rectangle++(0.008,0.008)
  (0.007,0.371)rectangle++(0.008,0.008) (-0.007,0.362)rectangle++(0.008,0.008) (-0.012,0.377)rectangle++(0.008,0.008) (-0.005,0.373)rectangle++(0.008,0.008) (-0.017,0.248)rectangle++(0.008,0.008)
  (-0.017,0.309)rectangle++(0.008,0.008) (-0.013,0.33)rectangle++(0.008,0.008) (-0.004,0.339)rectangle++(0.008,0.008) (0.002,0.334)rectangle++(0.008,0.008) (-0.015,0.336)rectangle++(0.008,0.008)
  (-0.016,0.214)rectangle++(0.008,0.008) (-0.004,0.336)rectangle++(0.008,0.008) (-0.009,0.301)rectangle++(0.008,0.008) (0,0.388)rectangle++(0.008,0.008) (0.004,0.375)rectangle++(0.008,0.008)
  (-0.008,0.246)rectangle++(0.008,0.008) (0.003,0.362)rectangle++(0.008,0.008) (-0.008,0.379)rectangle++(0.008,0.008) (0.008,0.344)rectangle++(0.008,0.008) (-0.004,0.309)rectangle++(0.008,0.008)
  (-0.007,0.233)rectangle++(0.008,0.008) (-0.012,0.306)rectangle++(0.008,0.008) (-0.01,0.332)rectangle++(0.008,0.008) (-0.017,0.304)rectangle++(0.008,0.008) (-0.011,0.38)rectangle++(0.008,0.008)
  (-0.004,0.157)rectangle++(0.008,0.008) (0.003,0.382)rectangle++(0.008,0.008) (-0.006,0.379)rectangle++(0.008,0.008) (-0.017,0.278)rectangle++(0.008,0.008) (-0.002,0.417)rectangle++(0.008,0.008)
  (0.008,0.325)rectangle++(0.008,0.008) (-0.008,0.3)rectangle++(0.008,0.008) (-0.017,0.217)rectangle++(0.008,0.008) (0.009,0.227)rectangle++(0.008,0.008) (0.008,0.321)rectangle++(0.008,0.008)
  (0.001,0.382)rectangle++(0.008,0.008) (0.009,0.255)rectangle++(0.008,0.008) (0.006,0.371)rectangle++(0.008,0.008) (0.004,0.264)rectangle++(0.008,0.008) (-0.013,0.374)rectangle++(0.008,0.008)
  (-0.004,0.155)rectangle++(0.008,0.008) (-0.018,0.297)rectangle++(0.008,0.008) (-0.008,0.39)rectangle++(0.008,0.008) (-0.004,0.162)rectangle++(0.008,0.008) (0.008,0.284)rectangle++(0.008,0.008)
  (-0.01,0.363)rectangle++(0.008,0.008) (-0.004,0.15)rectangle++(0.008,0.008) (0.009,0.28)rectangle++(0.008,0.008) (0.001,0.394)rectangle++(0.008,0.008) (-0.008,0.25)rectangle++(0.008,0.008)
  (-0.003,0.348)rectangle++(0.008,0.008) (0,0.24)rectangle++(0.008,0.008) (0.007,0.277)rectangle++(0.008,0.008) (-0.004,0.114)rectangle++(0.008,0.008) (0.01,0.229)rectangle++(0.008,0.008)
  (0.001,0.39)rectangle++(0.008,0.008) (-0.002,0.418)rectangle++(0.008,0.008) (-0.015,0.353)rectangle++(0.008,0.008) (0.008,0.329)rectangle++(0.008,0.008) (0.009,0.307)rectangle++(0.008,0.008)
  (-0.006,0.377)rectangle++(0.008,0.008) (-0.008,0.396)rectangle++(0.008,0.008) (-0.001,0.384)rectangle++(0.008,0.008) (-0.004,0.159)rectangle++(0.008,0.008) (0.01,0.278)rectangle++(0.008,0.008)
  (0.004,0.338)rectangle++(0.008,0.008) (-0.008,0.239)rectangle++(0.008,0.008) (-0.001,0.22)rectangle++(0.008,0.008) (0.007,0.365)rectangle++(0.008,0.008) (-0.004,0.313)rectangle++(0.008,0.008)
  (0,0.38)rectangle++(0.008,0.008) (-0.006,0.353)rectangle++(0.008,0.008) (-0.004,0.146)rectangle++(0.008,0.008) (-0.005,0.422)rectangle++(0.008,0.008) (-0.009,0.254)rectangle++(0.008,0.008)
  (-0.007,0.221)rectangle++(0.008,0.008) (-0.012,0.359)rectangle++(0.008,0.008) (-0.004,0.405)rectangle++(0.008,0.008) (0.01,0.231)rectangle++(0.008,0.008) (-0.01,0.359)rectangle++(0.008,0.008)
  (-0.018,0.27)rectangle++(0.008,0.008) (-0.004,0.35)rectangle++(0.008,0.008) (-0.008,0.24)rectangle++(0.008,0.008) (0.003,0.375)rectangle++(0.008,0.008) (-0.008,0.389)rectangle++(0.008,0.008)
  (-0.004,0.317)rectangle++(0.008,0.008) (-0.007,0.234)rectangle++(0.008,0.008) (0.007,0.243)rectangle++(0.008,0.008) (-0.013,0.359)rectangle++(0.008,0.008) (-0.013,0.331)rectangle++(0.008,0.008)
  (0.004,0.352)rectangle++(0.008,0.008) (-0.004,0.427)rectangle++(0.008,0.008) (0.008,0.334)rectangle++(0.008,0.008) (-0.018,0.26)rectangle++(0.008,0.008) (-0.016,0.309)rectangle++(0.008,0.008)
  (-0.001,0.234)rectangle++(0.008,0.008) (-0.011,0.262)rectangle++(0.008,0.008) (-0.009,0.38)rectangle++(0.008,0.008) (-0.007,0.368)rectangle++(0.008,0.008) (0,0.306)rectangle++(0.008,0.008)
  (-0.01,0.333)rectangle++(0.008,0.008) (0.01,0.257)rectangle++(0.008,0.008) (-0.007,0.225)rectangle++(0.008,0.008) (0.008,0.334)rectangle++(0.008,0.008) (0.003,0.372)rectangle++(0.008,0.008)
  (-0.01,0.361)rectangle++(0.008,0.008) (0.004,0.266)rectangle++(0.008,0.008) (0.01,0.281)rectangle++(0.008,0.008) (-0.013,0.354)rectangle++(0.008,0.008) (0,0.243)rectangle++(0.008,0.008)
  (-0.017,0.308)rectangle++(0.008,0.008) (-0.012,0.262)rectangle++(0.008,0.008) (-0.006,0.378)rectangle++(0.008,0.008) (0.007,0.211)rectangle++(0.008,0.008) (-0.004,0.086)rectangle++(0.008,0.008)
  (-0.005,0.19)rectangle++(0.008,0.008) (-0.016,0.335)rectangle++(0.008,0.008) (-0.006,0.408)rectangle++(0.008,0.008) (-0.01,0.385)rectangle++(0.008,0.008) (0.01,0.261)rectangle++(0.008,0.008)
  (0.009,0.256)rectangle++(0.008,0.008) (-0.004,0.368)rectangle++(0.008,0.008) (-0.005,0.418)rectangle++(0.008,0.008) (0.003,0.363)rectangle++(0.008,0.008) (-0.004,0.381)rectangle++(0.008,0.008)
  (0.002,0.354)rectangle++(0.008,0.008) (-0.006,0.351)rectangle++(0.008,0.008) (0,0.359)rectangle++(0.008,0.008) (0.001,0.395)rectangle++(0.008,0.008) (0.009,0.303)rectangle++(0.008,0.008)
  (0.002,0.331)rectangle++(0.008,0.008) (-0.008,0.4)rectangle++(0.008,0.008) (-0.017,0.32)rectangle++(0.008,0.008) (-0.014,0.333)rectangle++(0.008,0.008) (-0.01,0.356)rectangle++(0.008,0.008)
  (0.008,0.345)rectangle++(0.008,0.008) (-0.003,0.418)rectangle++(0.008,0.008) (-0.009,0.254)rectangle++(0.008,0.008) (0.004,0.309)rectangle++(0.008,0.008) (-0.006,0.397)rectangle++(0.008,0.008)
  (0.002,0.344)rectangle++(0.008,0.008) (-0.008,0.24)rectangle++(0.008,0.008) (-0.009,0.403)rectangle++(0.008,0.008) (-0.013,0.357)rectangle++(0.008,0.008) (-0.008,0.24)rectangle++(0.008,0.008)
  (0.001,0.389)rectangle++(0.008,0.008) (-0.004,0.317)rectangle++(0.008,0.008) (0.009,0.286)rectangle++(0.008,0.008) (-0.016,0.318)rectangle++(0.008,0.008) (-0.004,0.424)rectangle++(0.008,0.008)
  (0.009,0.223)rectangle++(0.008,0.008) (-0.006,0.218)rectangle++(0.008,0.008) (-0.004,0.089)rectangle++(0.008,0.008) (-0.015,0.36)rectangle++(0.008,0.008) (-0.004,0.149)rectangle++(0.008,0.008)
  (0.009,0.218)rectangle++(0.008,0.008) (0.008,0.304)rectangle++(0.008,0.008) (-0.017,0.319)rectangle++(0.008,0.008) (-0.004,0.132)rectangle++(0.008,0.008) (0.007,0.32)rectangle++(0.008,0.008)
  (0.006,0.358)rectangle++(0.008,0.008) (-0.01,0.336)rectangle++(0.008,0.008) (-0.007,0.388)rectangle++(0.008,0.008) (-0.011,0.375)rectangle++(0.008,0.008) (0.001,0.373)rectangle++(0.008,0.008)
  (-0.004,0.428)rectangle++(0.008,0.008) (0.01,0.3)rectangle++(0.008,0.008) (-0.016,0.279)rectangle++(0.008,0.008) (0,0.38)rectangle++(0.008,0.008) (-0.012,0.367)rectangle++(0.008,0.008)
  (-0.003,0.398)rectangle++(0.008,0.008) (-0.009,0.341)rectangle++(0.008,0.008) (0.004,0.342)rectangle++(0.008,0.008) (-0.002,0.4)rectangle++(0.008,0.008) (0.004,0.376)rectangle++(0.008,0.008)
  (-0.012,0.267)rectangle++(0.008,0.008) (-0.004,0.11)rectangle++(0.008,0.008) (-0.018,0.238)rectangle++(0.008,0.008) (0,0.327)rectangle++(0.008,0.008) (0.008,0.323)rectangle++(0.008,0.008)
  (-0.014,0.352)rectangle++(0.008,0.008) (0,0.242)rectangle++(0.008,0.008) (-0.009,0.39)rectangle++(0.008,0.008) (-0.009,0.391)rectangle++(0.008,0.008) (0.006,0.37)rectangle++(0.008,0.008)
  (-0.013,0.374)rectangle++(0.008,0.008) (0.008,0.355)rectangle++(0.008,0.008) (-0.013,0.377)rectangle++(0.008,0.008) (-0.008,0.397)rectangle++(0.008,0.008) (-0.017,0.29)rectangle++(0.008,0.008)
  (-0.004,0.12)rectangle++(0.008,0.008) (-0.009,0.27)rectangle++(0.008,0.008) (-0.007,0.224)rectangle++(0.008,0.008) (0.01,0.255)rectangle++(0.008,0.008) (-0.001,0.22)rectangle++(0.008,0.008)
  (-0.011,0.306)rectangle++(0.008,0.008) (0.007,0.344)rectangle++(0.008,0.008) (-0.004,0.358)rectangle++(0.008,0.008) (-0.002,0.361)rectangle++(0.008,0.008) (-0.004,0.149)rectangle++(0.008,0.008)
  (-0.004,0.275)rectangle++(0.008,0.008) (-0.01,0.375)rectangle++(0.008,0.008) (-0.004,0.315)rectangle++(0.008,0.008) (-0.009,0.339)rectangle++(0.008,0.008) (-0.004,0.14)rectangle++(0.008,0.008)
  (-0.011,0.374)rectangle++(0.008,0.008) (-0.009,0.381)rectangle++(0.008,0.008) (-0.011,0.375)rectangle++(0.008,0.008) (-0.012,0.318)rectangle++(0.008,0.008) (0,0.381)rectangle++(0.008,0.008)
  (-0.018,0.287)rectangle++(0.008,0.008) (0.003,0.381)rectangle++(0.008,0.008) (0.001,0.383)rectangle++(0.008,0.008) (0.001,0.361)rectangle++(0.008,0.008) (-0.001,0.409)rectangle++(0.008,0.008)
  (0.006,0.334)rectangle++(0.008,0.008) (-0.012,0.315)rectangle++(0.008,0.008) (-0.01,0.379)rectangle++(0.008,0.008) (-0.002,0.353)rectangle++(0.008,0.008) (0.004,0.374)rectangle++(0.008,0.008)
  (0.004,0.381)rectangle++(0.008,0.008) (-0.016,0.332)rectangle++(0.008,0.008) (0.008,0.212)rectangle++(0.008,0.008) (-0.009,0.391)rectangle++(0.008,0.008) (0,0.249)rectangle++(0.008,0.008)
  (-0.004,0.102)rectangle++(0.008,0.008) (-0.011,0.374)rectangle++(0.008,0.008) (-0.004,0.155)rectangle++(0.008,0.008) (-0.004,0.063)rectangle++(0.008,0.008) (0.003,0.38)rectangle++(0.008,0.008)
  (-0.004,0.157)rectangle++(0.008,0.008) (0.008,0.32)rectangle++(0.008,0.008) (0.002,0.376)rectangle++(0.008,0.008) (-0.016,0.309)rectangle++(0.008,0.008) (-0.006,0.412)rectangle++(0.008,0.008)
  (0.008,0.329)rectangle++(0.008,0.008) (0.007,0.359)rectangle++(0.008,0.008) (-0.013,0.337)rectangle++(0.008,0.008) (-0.005,0.394)rectangle++(0.008,0.008) (-0.004,0.36)rectangle++(0.008,0.008)
  (-0.004,0.149)rectangle++(0.008,0.008) (-0.006,0.202)rectangle++(0.008,0.008) (0.002,0.364)rectangle++(0.008,0.008) (-0.014,0.339)rectangle++(0.008,0.008) (-0.015,0.364)rectangle++(0.008,0.008)
  (0.004,0.318)rectangle++(0.008,0.008) (-0.004,0.386)rectangle++(0.008,0.008) (0,0.308)rectangle++(0.008,0.008) (-0.008,0.297)rectangle++(0.008,0.008) (-0.004,0.354)rectangle++(0.008,0.008)
  (0.004,0.383)rectangle++(0.008,0.008) (-0.012,0.267)rectangle++(0.008,0.008) (-0.017,0.281)rectangle++(0.008,0.008) (0.01,0.261)rectangle++(0.008,0.008) (-0.004,0.359)rectangle++(0.008,0.008)
  (0.004,0.265)rectangle++(0.008,0.008) (0.005,0.368)rectangle++(0.008,0.008) (-0.016,0.258)rectangle++(0.008,0.008) (-0.006,0.371)rectangle++(0.008,0.008) (0.01,0.285)rectangle++(0.008,0.008)
  (-0.016,0.345)rectangle++(0.008,0.008) (-0.012,0.263)rectangle++(0.008,0.008) (0.003,0.352)rectangle++(0.008,0.008) (-0.016,0.332)rectangle++(0.008,0.008) (-0.001,0.239)rectangle++(0.008,0.008)
  (-0.001,0.218)rectangle++(0.008,0.008) (0.009,0.295)rectangle++(0.008,0.008) (-0.006,0.374)rectangle++(0.008,0.008) (0,0.306)rectangle++(0.008,0.008) (-0.009,0.384)rectangle++(0.008,0.008)
  (-0.005,0.411)rectangle++(0.008,0.008) (0.009,0.254)rectangle++(0.008,0.008) (0.001,0.387)rectangle++(0.008,0.008) (-0.007,0.401)rectangle++(0.008,0.008) (0.005,0.356)rectangle++(0.008,0.008)
  (-0.004,0.342)rectangle++(0.008,0.008) (0.008,0.222)rectangle++(0.008,0.008) (-0.015,0.325)rectangle++(0.008,0.008) (-0.001,0.237)rectangle++(0.008,0.008) (0.005,0.326)rectangle++(0.008,0.008)
  (0.007,0.338)rectangle++(0.008,0.008) (-0.018,0.303)rectangle++(0.008,0.008) (-0.01,0.391)rectangle++(0.008,0.008) (-0.009,0.387)rectangle++(0.008,0.008) (0.009,0.257)rectangle++(0.008,0.008)
  (0.006,0.372)rectangle++(0.008,0.008) (0.006,0.291)rectangle++(0.008,0.008) (0.01,0.25)rectangle++(0.008,0.008) (-0.004,0.344)rectangle++(0.008,0.008) (-0.006,0.386)rectangle++(0.008,0.008)
  (0.001,0.383)rectangle++(0.008,0.008) (-0.004,0.158)rectangle++(0.008,0.008) (-0.005,0.411)rectangle++(0.008,0.008) (-0.005,0.361)rectangle++(0.008,0.008) (0.009,0.225)rectangle++(0.008,0.008)
  (0.007,0.276)rectangle++(0.008,0.008) (-0.004,0.114)rectangle++(0.008,0.008) (-0.005,0.373)rectangle++(0.008,0.008) (0.005,0.357)rectangle++(0.008,0.008) (0.008,0.322)rectangle++(0.008,0.008)
  (0.001,0.299)rectangle++(0.008,0.008) (0.009,0.277)rectangle++(0.008,0.008) (0.003,0.372)rectangle++(0.008,0.008) (-0.017,0.307)rectangle++(0.008,0.008) (-0.008,0.391)rectangle++(0.008,0.008)
  (-0.004,0.342)rectangle++(0.008,0.008) (0.008,0.356)rectangle++(0.008,0.008) (-0.011,0.384)rectangle++(0.008,0.008) (-0.006,0.353)rectangle++(0.008,0.008) (-0.003,0.374)rectangle++(0.008,0.008)
  (-0.002,0.414)rectangle++(0.008,0.008) (-0.004,0.322)rectangle++(0.008,0.008) (-0.006,0.42)rectangle++(0.008,0.008) (0,0.313)rectangle++(0.008,0.008) (-0.004,0.408)rectangle++(0.008,0.008)
  (-0.011,0.381)rectangle++(0.008,0.008) (-0.004,0.362)rectangle++(0.008,0.008) (0.003,0.368)rectangle++(0.008,0.008) (-0.002,0.378)rectangle++(0.008,0.008) (0,0.381)rectangle++(0.008,0.008)
  (0.007,0.336)rectangle++(0.008,0.008) (0,0.236)rectangle++(0.008,0.008) (0.009,0.305)rectangle++(0.008,0.008) (0.007,0.314)rectangle++(0.008,0.008) (-0.017,0.275)rectangle++(0.008,0.008)
  (-0.004,0.375)rectangle++(0.008,0.008) (0.008,0.346)rectangle++(0.008,0.008) (0.01,0.25)rectangle++(0.008,0.008) (0.001,0.396)rectangle++(0.008,0.008) (-0.005,0.419)rectangle++(0.008,0.008)
  (-0.018,0.234)rectangle++(0.008,0.008) (-0.005,0.375)rectangle++(0.008,0.008) (0.003,0.385)rectangle++(0.008,0.008) (0.003,0.375)rectangle++(0.008,0.008) (0.01,0.261)rectangle++(0.008,0.008)
  (-0.016,0.342)rectangle++(0.008,0.008) (-0.007,0.408)rectangle++(0.008,0.008) (0.002,0.339)rectangle++(0.008,0.008) (-0.009,0.386)rectangle++(0.008,0.008) (0.006,0.375)rectangle++(0.008,0.008)
  (-0.014,0.352)rectangle++(0.008,0.008) (-0.003,0.366)rectangle++(0.008,0.008) (0,0.243)rectangle++(0.008,0.008) (0.003,0.375)rectangle++(0.008,0.008) (0.002,0.355)rectangle++(0.008,0.008)
  (-0.016,0.258)rectangle++(0.008,0.008) (-0.011,0.26)rectangle++(0.008,0.008) (-0.004,0.107)rectangle++(0.008,0.008) (-0.004,0.043)rectangle++(0.008,0.008) (0.008,0.334)rectangle++(0.008,0.008)
  (-0.007,0.239)rectangle++(0.008,0.008) (-0.009,0.391)rectangle++(0.008,0.008) (-0.007,0.391)rectangle++(0.008,0.008) (-0.004,0.162)rectangle++(0.008,0.008) (0.004,0.331)rectangle++(0.008,0.008)
  (0.006,0.341)rectangle++(0.008,0.008) (-0.003,0.396)rectangle++(0.008,0.008) (-0.008,0.406)rectangle++(0.008,0.008) (-0.01,0.375)rectangle++(0.008,0.008) (-0.001,0.384)rectangle++(0.008,0.008)
  (-0.004,0.159)rectangle++(0.008,0.008) (-0.004,0.39)rectangle++(0.008,0.008) (-0.013,0.357)rectangle++(0.008,0.008) (-0.001,0.384)rectangle++(0.008,0.008) (0.007,0.225)rectangle++(0.008,0.008)
  (0.004,0.353)rectangle++(0.008,0.008) (0.007,0.334)rectangle++(0.008,0.008) (0.01,0.308)rectangle++(0.008,0.008) (-0.012,0.314)rectangle++(0.008,0.008) (0.01,0.256)rectangle++(0.008,0.008)
  (0,0.33)rectangle++(0.008,0.008) (-0.008,0.237)rectangle++(0.008,0.008) (-0.008,0.405)rectangle++(0.008,0.008) (-0.007,0.388)rectangle++(0.008,0.008) (0.003,0.375)rectangle++(0.008,0.008)
  (-0.008,0.246)rectangle++(0.008,0.008) (0.008,0.32)rectangle++(0.008,0.008) (-0.01,0.367)rectangle++(0.008,0.008) (-0.014,0.351)rectangle++(0.008,0.008) (0,0.242)rectangle++(0.008,0.008)
  (-0.004,0.4)rectangle++(0.008,0.008) (-0.008,0.374)rectangle++(0.008,0.008) (-0.018,0.267)rectangle++(0.008,0.008) (0.002,0.377)rectangle++(0.008,0.008) (-0.007,0.362)rectangle++(0.008,0.008)
  (-0.012,0.368)rectangle++(0.008,0.008) (-0.004,0.152)rectangle++(0.008,0.008) (-0.012,0.331)rectangle++(0.008,0.008) (-0.001,0.238)rectangle++(0.008,0.008) (-0.016,0.323)rectangle++(0.008,0.008)
  (0.001,0.372)rectangle++(0.008,0.008) (0.003,0.382)rectangle++(0.008,0.008) (-0.004,0.158)rectangle++(0.008,0.008) (-0.002,0.204)rectangle++(0.008,0.008) (0.01,0.237)rectangle++(0.008,0.008)
  (-0.003,0.374)rectangle++(0.008,0.008) (-0.001,0.387)rectangle++(0.008,0.008) (-0.007,0.38)rectangle++(0.008,0.008) (-0.014,0.36)rectangle++(0.008,0.008) (0.006,0.291)rectangle++(0.008,0.008)
  (-0.012,0.339)rectangle++(0.008,0.008) (-0.012,0.315)rectangle++(0.008,0.008) (-0.006,0.375)rectangle++(0.008,0.008) (0.008,0.216)rectangle++(0.008,0.008) (0.001,0.368)rectangle++(0.008,0.008)
  (0.001,0.397)rectangle++(0.008,0.008) (-0.018,0.236)rectangle++(0.008,0.008) (-0.003,0.413)rectangle++(0.008,0.008) (-0.007,0.38)rectangle++(0.008,0.008) (-0.004,0.157)rectangle++(0.008,0.008)
  (-0.002,0.203)rectangle++(0.008,0.008) (-0.004,0.083)rectangle++(0.008,0.008) (0.008,0.212)rectangle++(0.008,0.008) (0.004,0.327)rectangle++(0.008,0.008) (0.01,0.251)rectangle++(0.008,0.008)
  (0.007,0.357)rectangle++(0.008,0.008) (-0.001,0.405)rectangle++(0.008,0.008) (-0.004,0.168)rectangle++(0.008,0.008) (-0.018,0.291)rectangle++(0.008,0.008) (-0.008,0.394)rectangle++(0.008,0.008)
  (0.004,0.38)rectangle++(0.008,0.008) (-0.015,0.255)rectangle++(0.008,0.008) (-0.014,0.367)rectangle++(0.008,0.008) (-0.009,0.367)rectangle++(0.008,0.008) (-0.004,0.152)rectangle++(0.008,0.008)
  (-0.004,0.275)rectangle++(0.008,0.008) (-0.001,0.226)rectangle++(0.008,0.008) (0.001,0.398)rectangle++(0.008,0.008) (-0.008,0.362)rectangle++(0.008,0.008) (-0.008,0.357)rectangle++(0.008,0.008)
  (0.001,0.368)rectangle++(0.008,0.008) (0.001,0.375)rectangle++(0.008,0.008) (0.009,0.307)rectangle++(0.008,0.008) (-0.015,0.354)rectangle++(0.008,0.008) (-0.014,0.341)rectangle++(0.008,0.008)
  (0.004,0.316)rectangle++(0.008,0.008) (-0.004,0.13)rectangle++(0.008,0.008) (0.005,0.372)rectangle++(0.008,0.008) (0.002,0.375)rectangle++(0.008,0.008) (0.002,0.366)rectangle++(0.008,0.008)
  (-0.004,0.151)rectangle++(0.008,0.008) (-0.015,0.335)rectangle++(0.008,0.008) (-0.016,0.278)rectangle++(0.008,0.008) (0.007,0.347)rectangle++(0.008,0.008) (0,0.398)rectangle++(0.008,0.008)
  (-0.004,0.165)rectangle++(0.008,0.008) (-0.016,0.323)rectangle++(0.008,0.008) (0.005,0.371)rectangle++(0.008,0.008) (-0.016,0.334)rectangle++(0.008,0.008) (-0.008,0.236)rectangle++(0.008,0.008)
  (-0.005,0.411)rectangle++(0.008,0.008) (-0.012,0.368)rectangle++(0.008,0.008) (-0.005,0.37)rectangle++(0.008,0.008) (0,0.244)rectangle++(0.008,0.008) (-0.004,0.1)rectangle++(0.008,0.008)
  (-0.005,0.192)rectangle++(0.008,0.008) (-0.013,0.375)rectangle++(0.008,0.008) (-0.008,0.244)rectangle++(0.008,0.008) (-0.001,0.221)rectangle++(0.008,0.008) (-0.018,0.238)rectangle++(0.008,0.008)
  (-0.004,0.407)rectangle++(0.008,0.008) (-0.005,0.34)rectangle++(0.008,0.008) (0.006,0.374)rectangle++(0.008,0.008) (0.005,0.376)rectangle++(0.008,0.008) (0,0.244)rectangle++(0.008,0.008)
  (-0.01,0.363)rectangle++(0.008,0.008) (-0.018,0.251)rectangle++(0.008,0.008) (-0.005,0.421)rectangle++(0.008,0.008) (0,0.387)rectangle++(0.008,0.008) (0.009,0.304)rectangle++(0.008,0.008)
  (-0.004,0.125)rectangle++(0.008,0.008) (-0.003,0.197)rectangle++(0.008,0.008) (-0.006,0.211)rectangle++(0.008,0.008) (-0.015,0.36)rectangle++(0.008,0.008) (-0.005,0.425)rectangle++(0.008,0.008)
  (-0.004,0.176)rectangle++(0.008,0.008) (0.003,0.366)rectangle++(0.008,0.008) (-0.004,0.151)rectangle++(0.008,0.008) (-0.004,0.061)rectangle++(0.008,0.008) (0.01,0.278)rectangle++(0.008,0.008)
  (0.002,0.375)rectangle++(0.008,0.008) (-0.001,0.361)rectangle++(0.008,0.008) (-0.01,0.383)rectangle++(0.008,0.008) (-0.004,0.159)rectangle++(0.008,0.008) (0.01,0.248)rectangle++(0.008,0.008)
  (0.004,0.382)rectangle++(0.008,0.008) (-0.008,0.408)rectangle++(0.008,0.008) (0.007,0.33)rectangle++(0.008,0.008) (-0.004,0.136)rectangle++(0.008,0.008) (0.002,0.383)rectangle++(0.008,0.008)
  (-0.017,0.303)rectangle++(0.008,0.008) (-0.004,0.125)rectangle++(0.008,0.008) (-0.001,0.388)rectangle++(0.008,0.008) (-0.004,0.317)rectangle++(0.008,0.008) (-0.016,0.304)rectangle++(0.008,0.008)
  (-0.004,0.351)rectangle++(0.008,0.008) (-0.007,0.356)rectangle++(0.008,0.008) (-0.014,0.326)rectangle++(0.008,0.008) (-0.008,0.393)rectangle++(0.008,0.008) (-0.007,0.384)rectangle++(0.008,0.008)
  (-0.001,0.391)rectangle++(0.008,0.008) (-0.004,0.162)rectangle++(0.008,0.008) (0.003,0.301)rectangle++(0.008,0.008) (0.005,0.364)rectangle++(0.008,0.008) (0.003,0.375)rectangle++(0.008,0.008)
  (0.01,0.266)rectangle++(0.008,0.008) (0.006,0.257)rectangle++(0.008,0.008) (-0.007,0.224)rectangle++(0.008,0.008) (-0.004,0.389)rectangle++(0.008,0.008) (0.001,0.309)rectangle++(0.008,0.008)
  (0.008,0.331)rectangle++(0.008,0.008) (-0.003,0.408)rectangle++(0.008,0.008) (-0.004,0.428)rectangle++(0.008,0.008) (-0.004,0.177)rectangle++(0.008,0.008) (-0.006,0.207)rectangle++(0.008,0.008)
  (-0.005,0.342)rectangle++(0.008,0.008) (0.008,0.326)rectangle++(0.008,0.008) (-0.017,0.292)rectangle++(0.008,0.008) (-0.012,0.357)rectangle++(0.008,0.008) (0.008,0.347)rectangle++(0.008,0.008)
  (-0.012,0.267)rectangle++(0.008,0.008) (-0.004,0.11)rectangle++(0.008,0.008) (-0.005,0.194)rectangle++(0.008,0.008) (-0.018,0.291)rectangle++(0.008,0.008) (-0.018,0.284)rectangle++(0.008,0.008)
  (0.003,0.377)rectangle++(0.008,0.008) (-0.015,0.322)rectangle++(0.008,0.008) (-0.004,0.403)rectangle++(0.008,0.008) (-0.015,0.362)rectangle++(0.008,0.008) (-0.007,0.383)rectangle++(0.008,0.008)
  (-0.008,0.246)rectangle++(0.008,0.008) (-0.009,0.368)rectangle++(0.008,0.008) (-0.004,0.381)rectangle++(0.008,0.008) (0.007,0.214)rectangle++(0.008,0.008) (-0.004,0.087)rectangle++(0.008,0.008)
  (-0.003,0.19)rectangle++(0.008,0.008) (-0.008,0.397)rectangle++(0.008,0.008) (-0.001,0.407)rectangle++(0.008,0.008) (-0.016,0.28)rectangle++(0.008,0.008) (-0.011,0.374)rectangle++(0.008,0.008)
  (-0.001,0.403)rectangle++(0.008,0.008) (-0.004,0.167)rectangle++(0.008,0.008) (-0.009,0.388)rectangle++(0.008,0.008) (0,0.394)rectangle++(0.008,0.008) (0.01,0.264)rectangle++(0.008,0.008)
  (-0.007,0.226)rectangle++(0.008,0.008) (-0.002,0.217)rectangle++(0.008,0.008) (-0.008,0.407)rectangle++(0.008,0.008) (-0.005,0.414)rectangle++(0.008,0.008) (-0.009,0.375)rectangle++(0.008,0.008)
  (-0.009,0.307)rectangle++(0.008,0.008) (0.001,0.371)rectangle++(0.008,0.008) (-0.015,0.336)rectangle++(0.008,0.008) (-0.014,0.373)rectangle++(0.008,0.008) (0.002,0.387)rectangle++(0.008,0.008)
  (0.001,0.308)rectangle++(0.008,0.008) (-0.017,0.22)rectangle++(0.008,0.008) (-0.004,0.09)rectangle++(0.008,0.008) (-0.004,0.035)rectangle++(0.008,0.008) (0.01,0.279)rectangle++(0.008,0.008)
  (-0.007,0.401)rectangle++(0.008,0.008) (-0.002,0.355)rectangle++(0.008,0.008) (-0.007,0.409)rectangle++(0.008,0.008) (-0.002,0.356)rectangle++(0.008,0.008) (-0.009,0.392)rectangle++(0.008,0.008)
  (0.006,0.374)rectangle++(0.008,0.008) (-0.004,0.315)rectangle++(0.008,0.008) (-0.002,0.418)rectangle++(0.008,0.008) (-0.004,0.173)rectangle++(0.008,0.008) (0.002,0.366)rectangle++(0.008,0.008)
  (0.006,0.355)rectangle++(0.008,0.008) (-0.018,0.288)rectangle++(0.008,0.008) (-0.017,0.249)rectangle++(0.008,0.008) (0.007,0.368)rectangle++(0.008,0.008) (-0.012,0.372)rectangle++(0.008,0.008)
  (0.008,0.341)rectangle++(0.008,0.008) (0.009,0.219)rectangle++(0.008,0.008) (-0.003,0.339)rectangle++(0.008,0.008) (0,0.386)rectangle++(0.008,0.008) (-0.003,0.405)rectangle++(0.008,0.008)
  (-0.008,0.38)rectangle++(0.008,0.008) (-0.003,0.409)rectangle++(0.008,0.008) (-0.004,0.321)rectangle++(0.008,0.008) (-0.01,0.363)rectangle++(0.008,0.008) (-0.004,0.15)rectangle++(0.008,0.008)
  (-0.004,0.061)rectangle++(0.008,0.008) (-0.012,0.376)rectangle++(0.008,0.008) (0.004,0.328)rectangle++(0.008,0.008) (-0.018,0.29)rectangle++(0.008,0.008) (0.001,0.378)rectangle++(0.008,0.008)
  (-0.016,0.324)rectangle++(0.008,0.008) (-0.009,0.338)rectangle++(0.008,0.008) (-0.001,0.409)rectangle++(0.008,0.008) (-0.002,0.355)rectangle++(0.008,0.008) (0.006,0.373)rectangle++(0.008,0.008)
  (-0.005,0.398)rectangle++(0.008,0.008) (-0.002,0.355)rectangle++(0.008,0.008) (-0.017,0.223)rectangle++(0.008,0.008) (-0.004,0.091)rectangle++(0.008,0.008) (-0.005,0.421)rectangle++(0.008,0.008)
  (0.007,0.347)rectangle++(0.008,0.008) (-0.013,0.375)rectangle++(0.008,0.008) (0.005,0.329)rectangle++(0.008,0.008) (-0.018,0.272)rectangle++(0.008,0.008) (-0.009,0.359)rectangle++(0.008,0.008)
  (0.007,0.351)rectangle++(0.008,0.008) (-0.018,0.271)rectangle++(0.008,0.008) (-0.004,0.111)rectangle++(0.008,0.008) (-0.009,0.269)rectangle++(0.008,0.008) (-0.004,0.111)rectangle++(0.008,0.008)
  (-0.004,0.044)rectangle++(0.008,0.008) (0.008,0.333)rectangle++(0.008,0.008) (-0.012,0.266)rectangle++(0.008,0.008) (-0.01,0.385)rectangle++(0.008,0.008) (-0.004,0.349)rectangle++(0.008,0.008)
  (-0.001,0.382)rectangle++(0.008,0.008) (0.006,0.373)rectangle++(0.008,0.008) (-0.008,0.246)rectangle++(0.008,0.008) (-0.001,0.384)rectangle++(0.008,0.008) (0,0.246)rectangle++(0.008,0.008)
  (-0.016,0.321)rectangle++(0.008,0.008) (-0.01,0.374)rectangle++(0.008,0.008) (-0.017,0.298)rectangle++(0.008,0.008) (-0.004,0.123)rectangle++(0.008,0.008) (-0.004,0.049)rectangle++(0.008,0.008)
  (-0.004,0.018)rectangle++(0.008,0.008) (-0.004,0.427)rectangle++(0.008,0.008) (0.004,0.37)rectangle++(0.008,0.008) (0.006,0.352)rectangle++(0.008,0.008) (-0.009,0.397)rectangle++(0.008,0.008)
  (0.008,0.347)rectangle++(0.008,0.008) (0.01,0.304)rectangle++(0.008,0.008) (-0.004,0.125)rectangle++(0.008,0.008) (0,0.404)rectangle++(0.008,0.008) (-0.011,0.375)rectangle++(0.008,0.008)
  (-0.007,0.361)rectangle++(0.008,0.008) (-0.004,0.149)rectangle++(0.008,0.008) (-0.003,0.405)rectangle++(0.008,0.008) (0.006,0.327)rectangle++(0.008,0.008) (-0.007,0.237)rectangle++(0.008,0.008)
  (-0.005,0.419)rectangle++(0.008,0.008) (-0.004,0.364)rectangle++(0.008,0.008) (0,0.243)rectangle++(0.008,0.008) (0.01,0.302)rectangle++(0.008,0.008) (-0.018,0.276)rectangle++(0.008,0.008)
  (0.001,0.393)rectangle++(0.008,0.008) (-0.004,0.163)rectangle++(0.008,0.008) (-0.013,0.354)rectangle++(0.008,0.008) (-0.008,0.391)rectangle++(0.008,0.008) (0.006,0.374)rectangle++(0.008,0.008)
  (-0.008,0.382)rectangle++(0.008,0.008) (-0.001,0.386)rectangle++(0.008,0.008) (-0.007,0.337)rectangle++(0.008,0.008) (-0.009,0.301)rectangle++(0.008,0.008) (0.005,0.368)rectangle++(0.008,0.008)
  (-0.013,0.373)rectangle++(0.008,0.008) (-0.009,0.373)rectangle++(0.008,0.008) (0.006,0.36)rectangle++(0.008,0.008) (0.009,0.259)rectangle++(0.008,0.008) (-0.012,0.31)rectangle++(0.008,0.008)
  (-0.012,0.325)rectangle++(0.008,0.008) (-0.018,0.286)rectangle++(0.008,0.008) (-0.003,0.375)rectangle++(0.008,0.008) (0.008,0.319)rectangle++(0.008,0.008) (-0.01,0.334)rectangle++(0.008,0.008)
  (-0.007,0.405)rectangle++(0.008,0.008) (0.004,0.321)rectangle++(0.008,0.008) (-0.001,0.408)rectangle++(0.008,0.008) (-0.006,0.414)rectangle++(0.008,0.008) (0.01,0.258)rectangle++(0.008,0.008)
  (0.005,0.324)rectangle++(0.008,0.008) (-0.004,0.306)rectangle++(0.008,0.008) (-0.018,0.266)rectangle++(0.008,0.008) (-0.008,0.396)rectangle++(0.008,0.008) (-0.006,0.38)rectangle++(0.008,0.008)
  (-0.004,0.401)rectangle++(0.008,0.008) (0.01,0.262)rectangle++(0.008,0.008) (-0.001,0.398)rectangle++(0.008,0.008) (-0.016,0.353)rectangle++(0.008,0.008) (-0.015,0.348)rectangle++(0.008,0.008)
  (-0.012,0.363)rectangle++(0.008,0.008) (0,0.244)rectangle++(0.008,0.008) (0.009,0.291)rectangle++(0.008,0.008) (-0.002,0.414)rectangle++(0.008,0.008) (0.001,0.312)rectangle++(0.008,0.008)
  (-0.005,0.424)rectangle++(0.008,0.008) (-0.017,0.274)rectangle++(0.008,0.008) (-0.001,0.361)rectangle++(0.008,0.008) (0.002,0.388)rectangle++(0.008,0.008) (-0.012,0.371)rectangle++(0.008,0.008)
  (-0.018,0.269)rectangle++(0.008,0.008) (-0.001,0.227)rectangle++(0.008,0.008) (-0.009,0.357)rectangle++(0.008,0.008) (-0.002,0.408)rectangle++(0.008,0.008) (0,0.374)rectangle++(0.008,0.008)
  (-0.004,0.155)rectangle++(0.008,0.008) (-0.014,0.362)rectangle++(0.008,0.008) (0.007,0.256)rectangle++(0.008,0.008) (-0.003,0.405)rectangle++(0.008,0.008) (-0.01,0.383)rectangle++(0.008,0.008)
  (-0.012,0.318)rectangle++(0.008,0.008) (-0.004,0.131)rectangle++(0.008,0.008) (-0.017,0.23)rectangle++(0.008,0.008) (-0.009,0.376)rectangle++(0.008,0.008) (-0.003,0.42)rectangle++(0.008,0.008)
  (-0.004,0.323)rectangle++(0.008,0.008) (-0.01,0.382)rectangle++(0.008,0.008) (0.009,0.214)rectangle++(0.008,0.008) (-0.005,0.367)rectangle++(0.008,0.008) (-0.01,0.388)rectangle++(0.008,0.008)
  (-0.018,0.278)rectangle++(0.008,0.008) (-0.004,0.115)rectangle++(0.008,0.008) (-0.005,0.369)rectangle++(0.008,0.008) (-0.012,0.318)rectangle++(0.008,0.008) (-0.003,0.404)rectangle++(0.008,0.008)
  (-0.013,0.328)rectangle++(0.008,0.008) (0.007,0.371)rectangle++(0.008,0.008) (-0.006,0.378)rectangle++(0.008,0.008) (-0.005,0.364)rectangle++(0.008,0.008) (0.001,0.394)rectangle++(0.008,0.008)
  (0.002,0.385)rectangle++(0.008,0.008) (0.006,0.327)rectangle++(0.008,0.008) (0,0.371)rectangle++(0.008,0.008) (-0.005,0.422)rectangle++(0.008,0.008) (0.007,0.357)rectangle++(0.008,0.008)
  (-0.004,0.312)rectangle++(0.008,0.008) (-0.014,0.373)rectangle++(0.008,0.008) (-0.004,0.382)rectangle++(0.008,0.008) (-0.009,0.367)rectangle++(0.008,0.008) (-0.016,0.279)rectangle++(0.008,0.008)
  (0.008,0.305)rectangle++(0.008,0.008) (-0.014,0.258)rectangle++(0.008,0.008) (0.006,0.307)rectangle++(0.008,0.008) (0.008,0.317)rectangle++(0.008,0.008) (-0.01,0.356)rectangle++(0.008,0.008)
  (-0.015,0.315)rectangle++(0.008,0.008) (-0.002,0.416)rectangle++(0.008,0.008) (0.002,0.356)rectangle++(0.008,0.008) (0.007,0.365)rectangle++(0.008,0.008) (-0.007,0.362)rectangle++(0.008,0.008)
  (-0.004,0.15)rectangle++(0.008,0.008) (-0.002,0.202)rectangle++(0.008,0.008) (0.009,0.329)rectangle++(0.008,0.008) (-0.015,0.356)rectangle++(0.008,0.008) (0.01,0.289)rectangle++(0.008,0.008)
  (0.004,0.373)rectangle++(0.008,0.008) (-0.004,0.407)rectangle++(0.008,0.008) (0.001,0.39)rectangle++(0.008,0.008) (0.004,0.368)rectangle++(0.008,0.008) (0.004,0.317)rectangle++(0.008,0.008)
  (-0.007,0.235)rectangle++(0.008,0.008) (-0.016,0.319)rectangle++(0.008,0.008) (0.004,0.38)rectangle++(0.008,0.008) (-0.016,0.316)rectangle++(0.008,0.008) (-0.009,0.357)rectangle++(0.008,0.008)
  (-0.009,0.359)rectangle++(0.008,0.008) (-0.008,0.391)rectangle++(0.008,0.008) (-0.004,0.162)rectangle++(0.008,0.008) (-0.014,0.354)rectangle++(0.008,0.008) (-0.004,0.393)rectangle++(0.008,0.008)
  (-0.005,0.42)rectangle++(0.008,0.008) (-0.003,0.426)rectangle++(0.008,0.008) (-0.004,0.177)rectangle++(0.008,0.008) (-0.011,0.368)rectangle++(0.008,0.008) (-0.018,0.287)rectangle++(0.008,0.008)
  (-0.012,0.31)rectangle++(0.008,0.008) (0.001,0.393)rectangle++(0.008,0.008) (-0.004,0.34)rectangle++(0.008,0.008) (-0.004,0.398)rectangle++(0.008,0.008) (0.009,0.267)rectangle++(0.008,0.008);
\fill[piece2]
  (-0.154,0.717)rectangle++(0.008,0.008) (-0.33,0.692)rectangle++(0.008,0.008) (-0.327,0.83)rectangle++(0.008,0.008) (-0.449,0.682)rectangle++(0.008,0.008) (-0.453,0.57)rectangle++(0.008,0.008)
  (-0.414,0.518)rectangle++(0.008,0.008) (-0.292,0.784)rectangle++(0.008,0.008) (-0.134,0.587)rectangle++(0.008,0.008) (-0.332,0.784)rectangle++(0.008,0.008) (-0.435,0.659)rectangle++(0.008,0.008)
  (-0.141,0.903)rectangle++(0.008,0.008) (-0.146,0.93)rectangle++(0.008,0.008) (-0.402,0.792)rectangle++(0.008,0.008) (-0.12,0.92)rectangle++(0.008,0.008) (-0.387,0.797)rectangle++(0.008,0.008)
  (-0.465,0.645)rectangle++(0.008,0.008) (-0.063,0.864)rectangle++(0.008,0.008) (-0.341,0.792)rectangle++(0.008,0.008) (-0.442,0.659)rectangle++(0.008,0.008) (-0.402,0.752)rectangle++(0.008,0.008)
  (-0.288,0.862)rectangle++(0.008,0.008) (-0.139,0.594)rectangle++(0.008,0.008) (-0.279,0.641)rectangle++(0.008,0.008) (-0.153,0.723)rectangle++(0.008,0.008) (-0.332,0.695)rectangle++(0.008,0.008)
  (-0.275,0.906)rectangle++(0.008,0.008) (-0.407,0.824)rectangle++(0.008,0.008) (-0.484,0.651)rectangle++(0.008,0.008) (-0.15,0.603)rectangle++(0.008,0.008) (-0.337,0.785)rectangle++(0.008,0.008)
  (-0.438,0.658)rectangle++(0.008,0.008) (-0.126,0.582)rectangle++(0.008,0.008) (-0.161,0.929)rectangle++(0.008,0.008) (-0.005,0.857)rectangle++(0.008,0.008) (-0.363,0.875)rectangle++(0.008,0.008)
  (-0.482,0.69)rectangle++(0.008,0.008) (-0.414,0.752)rectangle++(0.008,0.008) (-0.008,0.858)rectangle++(0.008,0.008) (-0.313,0.809)rectangle++(0.008,0.008) (-0.436,0.677)rectangle++(0.008,0.008)
  (-0.139,0.906)rectangle++(0.008,0.008) (-0.381,0.845)rectangle++(0.008,0.008) (-0.479,0.669)rectangle++(0.008,0.008) (-0.348,0.89)rectangle++(0.008,0.008) (-0.144,0.951)rectangle++(0.008,0.008)
  (-0.409,0.802)rectangle++(0.008,0.008) (-0.281,0.915)rectangle++(0.008,0.008) (-0.146,0.589)rectangle++(0.008,0.008) (-0.281,0.637)rectangle++(0.008,0.008) (-0.011,0.85)rectangle++(0.008,0.008)
  (-0.052,0.881)rectangle++(0.008,0.008) (-0.342,0.803)rectangle++(0.008,0.008) (-0.447,0.664)rectangle++(0.008,0.008) (-0.445,0.563)rectangle++(0.008,0.008) (-0.302,0.799)rectangle++(0.008,0.008)
  (-0.128,0.567)rectangle++(0.008,0.008) (-0.235,0.87)rectangle++(0.008,0.008) (-0.111,0.806)rectangle++(0.008,0.008) (-0.19,0.881)rectangle++(0.008,0.008) (-0.405,0.754)rectangle++(0.008,0.008)
  (-0.458,0.619)rectangle++(0.008,0.008) (-0.136,0.579)rectangle++(0.008,0.008) (-0.273,0.636)rectangle++(0.008,0.008) (-0.355,0.612)rectangle++(0.008,0.008) (-0.272,0.815)rectangle++(0.008,0.008)
  (-0.058,0.908)rectangle++(0.008,0.008) (-0.239,0.892)rectangle++(0.008,0.008) (-0.349,0.892)rectangle++(0.008,0.008) (-0.33,0.898)rectangle++(0.008,0.008) (-0.015,0.872)rectangle++(0.008,0.008)
  (-0.158,0.902)rectangle++(0.008,0.008) (-0.206,0.913)rectangle++(0.008,0.008) (-0.424,0.763)rectangle++(0.008,0.008) (-0.47,0.617)rectangle++(0.008,0.008) (-0.438,0.533)rectangle++(0.008,0.008)
  (-0.3,0.789)rectangle++(0.008,0.008) (-0.27,0.753)rectangle++(0.008,0.008) (-0.192,0.917)rectangle++(0.008,0.008) (-0.381,0.868)rectangle++(0.008,0.008) (-0.488,0.68)rectangle++(0.008,0.008)
  (-0.28,0.909)rectangle++(0.008,0.008) (-0.247,0.915)rectangle++(0.008,0.008) (-0.443,0.749)rectangle++(0.008,0.008) (-0.474,0.604)rectangle++(0.008,0.008) (-0.436,0.526)rectangle++(0.008,0.008)
  (-0.391,0.504)rectangle++(0.008,0.008) (-0.369,0.73)rectangle++(0.008,0.008) (-0.433,0.621)rectangle++(0.008,0.008) (-0.423,0.549)rectangle++(0.008,0.008) (-0.355,0.866)rectangle++(0.008,0.008)
  (-0.272,0.901)rectangle++(0.008,0.008) (-0.45,0.734)rectangle++(0.008,0.008) (-0.391,0.837)rectangle++(0.008,0.008) (-0.149,0.597)rectangle++(0.008,0.008) (-0.096,0.53)rectangle++(0.008,0.008)
  (-0.237,0.627)rectangle++(0.008,0.008) (-0.232,0.82)rectangle++(0.008,0.008) (-0.403,0.711)rectangle++(0.008,0.008) (-0.442,0.601)rectangle++(0.008,0.008) (-0.42,0.536)rectangle++(0.008,0.008)
  (-0.213,0.886)rectangle++(0.008,0.008) (-0.013,0.794)rectangle++(0.008,0.008) (-0.293,0.777)rectangle++(0.008,0.008) (-0.105,0.904)rectangle++(0.008,0.008) (-0.096,0.925)rectangle++(0.008,0.008)
  (-0.378,0.808)rectangle++(0.008,0.008) (-0.264,0.901)rectangle++(0.008,0.008) (-0.143,0.587)rectangle++(0.008,0.008) (-0.02,0.436)rectangle++(0.008,0.008) (-0.168,0.612)rectangle++(0.008,0.008)
  (-0.145,0.601)rectangle++(0.008,0.008) (-0.174,0.895)rectangle++(0.008,0.008) (-0.211,0.91)rectangle++(0.008,0.008) (-0.224,0.916)rectangle++(0.008,0.008) (-0.433,0.758)rectangle++(0.008,0.008)
  (-0.367,0.883)rectangle++(0.008,0.008) (-0.487,0.692)rectangle++(0.008,0.008) (-0.111,0.565)rectangle++(0.008,0.008) (-0.108,0.908)rectangle++(0.008,0.008) (-0.377,0.796)rectangle++(0.008,0.008)
  (-0.176,0.942)rectangle++(0.008,0.008) (-0.067,0.915)rectangle++(0.008,0.008) (-0.361,0.814)rectangle++(0.008,0.008) (-0.459,0.662)rectangle++(0.008,0.008) (-0.307,0.832)rectangle++(0.008,0.008)
  (-0.131,0.938)rectangle++(0.008,0.008) (-0.399,0.801)rectangle++(0.008,0.008) (-0.472,0.643)rectangle++(0.008,0.008) (-0.148,0.905)rectangle++(0.008,0.008) (-0.008,0.855)rectangle++(0.008,0.008)
  (-0.312,0.807)rectangle++(0.008,0.008) (-0.435,0.677)rectangle++(0.008,0.008) (-0.403,0.756)rectangle++(0.008,0.008) (-0.124,0.914)rectangle++(0.008,0.008) (-0.387,0.793)rectangle++(0.008,0.008)
  (-0.463,0.644)rectangle++(0.008,0.008) (-0.445,0.548)rectangle++(0.008,0.008) (-0.403,0.511)rectangle++(0.008,0.008) (-0.37,0.509)rectangle++(0.008,0.008) (-0.365,0.734)rectangle++(0.008,0.008)
  (-0.433,0.624)rectangle++(0.008,0.008) (-0.298,0.819)rectangle++(0.008,0.008) (-0.15,0.717)rectangle++(0.008,0.008) (-0.328,0.694)rectangle++(0.008,0.008) (-0.173,0.936)rectangle++(0.008,0.008)
  (-0.211,0.924)rectangle++(0.008,0.008) (-0.397,0.837)rectangle++(0.008,0.008) (-0.484,0.66)rectangle++(0.008,0.008) (-0.255,0.89)rectangle++(0.008,0.008) (-0.238,0.909)rectangle++(0.008,0.008)
  (-0.08,0.918)rectangle++(0.008,0.008) (-0.027,0.78)rectangle++(0.008,0.008) (-0.294,0.766)rectangle++(0.008,0.008) (-0.329,0.831)rectangle++(0.008,0.008) (-0.387,0.852)rectangle++(0.008,0.008)
  (-0.42,0.797)rectangle++(0.008,0.008) (-0.48,0.634)rectangle++(0.008,0.008) (-0.449,0.537)rectangle++(0.008,0.008) (-0.401,0.504)rectangle++(0.008,0.008) (-0.036,0.457)rectangle++(0.008,0.008)
  (-0.183,0.615)rectangle++(0.008,0.008) (-0.307,0.635)rectangle++(0.008,0.008) (-0.371,0.6)rectangle++(0.008,0.008) (-0.379,0.751)rectangle++(0.008,0.008) (-0.267,0.898)rectangle++(0.008,0.008)
  (-0.446,0.734)rectangle++(0.008,0.008) (-0.165,0.949)rectangle++(0.008,0.008) (-0.418,0.794)rectangle++(0.008,0.008) (-0.292,0.911)rectangle++(0.008,0.008) (-0.322,0.9)rectangle++(0.008,0.008)
  (-0.261,0.912)rectangle++(0.008,0.008) (-0.149,0.708)rectangle++(0.008,0.008) (-0.325,0.69)rectangle++(0.008,0.008) (-0.228,0.889)rectangle++(0.008,0.008) (-0.242,0.88)rectangle++(0.008,0.008)
  (-0.429,0.735)rectangle++(0.008,0.008) (-0.096,0.919)rectangle++(0.008,0.008) (-0.376,0.805)rectangle++(0.008,0.008) (-0.13,0.567)rectangle++(0.008,0.008) (-0.255,0.736)rectangle++(0.008,0.008)
  (-0.238,0.857)rectangle++(0.008,0.008) (-0.187,0.931)rectangle++(0.008,0.008) (-0.377,0.881)rectangle++(0.008,0.008) (-0.49,0.687)rectangle++(0.008,0.008) (-0.318,0.84)rectangle++(0.008,0.008)
  (-0.089,0.9)rectangle++(0.008,0.008) (-0.058,0.898)rectangle++(0.008,0.008) (-0.155,0.735)rectangle++(0.008,0.008) (-0.155,0.712)rectangle++(0.008,0.008) (-0.329,0.689)rectangle++(0.008,0.008)
  (-0.314,0.885)rectangle++(0.008,0.008) (-0.411,0.814)rectangle++(0.008,0.008) (-0.148,0.595)rectangle++(0.008,0.008) (-0.061,0.887)rectangle++(0.008,0.008) (-0.348,0.803)rectangle++(0.008,0.008)
  (-0.147,0.595)rectangle++(0.008,0.008) (-0.272,0.912)rectangle++(0.008,0.008) (-0.044,0.891)rectangle++(0.008,0.008) (-0.342,0.811)rectangle++(0.008,0.008) (-0.026,0.892)rectangle++(0.008,0.008)
  (-0.05,0.89)rectangle++(0.008,0.008) (-0.344,0.808)rectangle++(0.008,0.008) (-0.449,0.666)rectangle++(0.008,0.008) (-0.447,0.563)rectangle++(0.008,0.008) (-0.303,0.799)rectangle++(0.008,0.008)
  (-0.004,0.853)rectangle++(0.008,0.008) (-0.31,0.808)rectangle++(0.008,0.008) (-0.043,0.909)rectangle++(0.008,0.008) (-0.167,0.915)rectangle++(0.008,0.008) (-0.407,0.778)rectangle++(0.008,0.008)
  (-0.289,0.871)rectangle++(0.008,0.008) (-0.447,0.714)rectangle++(0.008,0.008) (-0.274,0.899)rectangle++(0.008,0.008) (-0.02,0.835)rectangle++(0.008,0.008) (-0.282,0.905)rectangle++(0.008,0.008)
  (-0.456,0.732)rectangle++(0.008,0.008) (-0.306,0.856)rectangle++(0.008,0.008) (-0.405,0.81)rectangle++(0.008,0.008) (-0.132,0.568)rectangle++(0.008,0.008) (-0.267,0.632)rectangle++(0.008,0.008)
  (-0.17,0.932)rectangle++(0.008,0.008) (-0.278,0.765)rectangle++(0.008,0.008) (-0.404,0.67)rectangle++(0.008,0.008) (-0.338,0.866)rectangle++(0.008,0.008) (-0.154,0.714)rectangle++(0.008,0.008)
  (-0.329,0.691)rectangle++(0.008,0.008) (-0.401,0.617)rectangle++(0.008,0.008) (-0.13,0.588)rectangle++(0.008,0.008) (-0.273,0.641)rectangle++(0.008,0.008) (-0.15,0.601)rectangle++(0.008,0.008)
  (-0.337,0.785)rectangle++(0.008,0.008) (-0.127,0.564)rectangle++(0.008,0.008) (-0.327,0.781)rectangle++(0.008,0.008) (-0.263,0.872)rectangle++(0.008,0.008) (-0.241,0.903)rectangle++(0.008,0.008)
  (-0.256,0.901)rectangle++(0.008,0.008) (-0.378,0.878)rectangle++(0.008,0.008) (-0.49,0.685)rectangle++(0.008,0.008) (-0.35,0.868)rectangle++(0.008,0.008) (-0.028,0.447)rectangle++(0.008,0.008)
  (-0.335,0.886)rectangle++(0.008,0.008) (-0.256,0.744)rectangle++(0.008,0.008) (-0.228,0.878)rectangle++(0.008,0.008) (-0.421,0.739)rectangle++(0.008,0.008) (-0.275,0.764)rectangle++(0.008,0.008)
  (-0.402,0.67)rectangle++(0.008,0.008) (-0.427,0.582)rectangle++(0.008,0.008) (-0.296,0.805)rectangle++(0.008,0.008) (-0.426,0.681)rectangle++(0.008,0.008) (-0.113,0.543)rectangle++(0.008,0.008)
  (-0.153,0.735)rectangle++(0.008,0.008) (-0.336,0.7)rectangle++(0.008,0.008) (-0.407,0.619)rectangle++(0.008,0.008) (-0.411,0.557)rectangle++(0.008,0.008) (-0.39,0.527)rectangle++(0.008,0.008)
  (-0.37,0.521)rectangle++(0.008,0.008) (-0.088,0.513)rectangle++(0.008,0.008) (-0.138,0.942)rectangle++(0.008,0.008) (-0.403,0.801)rectangle++(0.008,0.008) (-0.067,0.914)rectangle++(0.008,0.008)
  (-0.131,0.589)rectangle++(0.008,0.008) (-0.154,0.657)rectangle++(0.008,0.008) (-0.204,0.83)rectangle++(0.008,0.008) (-0.394,0.726)rectangle++(0.008,0.008) (-0.285,0.853)rectangle++(0.008,0.008)
  (-0.439,0.708)rectangle++(0.008,0.008) (-0.457,0.586)rectangle++(0.008,0.008) (-0.306,0.806)rectangle++(0.008,0.008) (-0.276,0.887)rectangle++(0.008,0.008) (-0.446,0.726)rectangle++(0.008,0.008)
  (-0.467,0.592)rectangle++(0.008,0.008) (-0.023,0.855)rectangle++(0.008,0.008) (-0.317,0.888)rectangle++(0.008,0.008) (-0.259,0.908)rectangle++(0.008,0.008) (-0.09,0.54)rectangle++(0.008,0.008)
  (-0.238,0.634)rectangle++(0.008,0.008) (-0.233,0.822)rectangle++(0.008,0.008) (-0.314,0.888)rectangle++(0.008,0.008) (-0.464,0.713)rectangle++(0.008,0.008) (-0.309,0.849)rectangle++(0.008,0.008)
  (-0.256,0.895)rectangle++(0.008,0.008) (-0.239,0.91)rectangle++(0.008,0.008) (-0.279,0.77)rectangle++(0.008,0.008) (-0.379,0.85)rectangle++(0.008,0.008) (-0.036,0.871)rectangle++(0.008,0.008)
  (-0.331,0.804)rectangle++(0.008,0.008) (-0.278,0.765)rectangle++(0.008,0.008) (-0.404,0.67)rectangle++(0.008,0.008) (-0.281,0.823)rectangle++(0.008,0.008) (-0.426,0.695)rectangle++(0.008,0.008)
  (-0.265,0.745)rectangle++(0.008,0.008) (-0.104,0.896)rectangle++(0.008,0.008) (-0.061,0.9)rectangle++(0.008,0.008) (-0.353,0.809)rectangle++(0.008,0.008) (-0.407,0.794)rectangle++(0.008,0.008)
  (-0.072,0.82)rectangle++(0.008,0.008) (-0.329,0.769)rectangle++(0.008,0.008) (-0.428,0.653)rectangle++(0.008,0.008) (-0.109,0.563)rectangle++(0.008,0.008) (-0.255,0.637)rectangle++(0.008,0.008)
  (-0.238,0.824)rectangle++(0.008,0.008) (-0.387,0.814)rectangle++(0.008,0.008) (-0.132,0.569)rectangle++(0.008,0.008) (-0.044,0.89)rectangle++(0.008,0.008) (-0.341,0.811)rectangle++(0.008,0.008)
  (-0.15,0.682)rectangle++(0.008,0.008) (-0.103,0.911)rectangle++(0.008,0.008) (-0.374,0.851)rectangle++(0.008,0.008) (-0.357,0.89)rectangle++(0.008,0.008) (-0.273,0.909)rectangle++(0.008,0.008)
  (-0.407,0.825)rectangle++(0.008,0.008) (-0.393,0.847)rectangle++(0.008,0.008) (-0.209,0.93)rectangle++(0.008,0.008) (-0.431,0.77)rectangle++(0.008,0.008) (-0.476,0.617)rectangle++(0.008,0.008)
  (-0.402,0.739)rectangle++(0.008,0.008) (-0.451,0.613)rectangle++(0.008,0.008) (-0.344,0.858)rectangle++(0.008,0.008) (-0.268,0.898)rectangle++(0.008,0.008) (-0.112,0.938)rectangle++(0.008,0.008)
  (-0.38,0.855)rectangle++(0.008,0.008) (-0.334,0.892)rectangle++(0.008,0.008) (-0.014,0.875)rectangle++(0.008,0.008) (-0.322,0.814)rectangle++(0.008,0.008) (-0.199,0.925)rectangle++(0.008,0.008)
  (-0.267,0.91)rectangle++(0.008,0.008) (-0.45,0.74)rectangle++(0.008,0.008) (-0.416,0.766)rectangle++(0.008,0.008) (-0.467,0.621)rectangle++(0.008,0.008) (-0.439,0.536)rectangle++(0.008,0.008)
  (-0.377,0.822)rectangle++(0.008,0.008) (-0.28,0.886)rectangle++(0.008,0.008) (-0.034,0.815)rectangle++(0.008,0.008) (-0.286,0.887)rectangle++(0.008,0.008) (-0.095,0.919)rectangle++(0.008,0.008)
  (-0.374,0.874)rectangle++(0.008,0.008) (-0.278,0.903)rectangle++(0.008,0.008) (-0.246,0.913)rectangle++(0.008,0.008) (-0.235,0.917)rectangle++(0.008,0.008) (0.004,0.805)rectangle++(0.008,0.008)
  (-0.056,0.873)rectangle++(0.008,0.008) (-0.171,0.903)rectangle++(0.008,0.008) (-0.305,0.899)rectangle++(0.008,0.008) (-0.464,0.721)rectangle++(0.008,0.008) (-0.219,0.912)rectangle++(0.008,0.008)
  (-0.156,0.737)rectangle++(0.008,0.008) (-0.175,0.915)rectangle++(0.008,0.008) (-0.41,0.775)rectangle++(0.008,0.008) (-0.126,0.561)rectangle++(0.008,0.008) (-0.299,0.851)rectangle++(0.008,0.008)
  (-0.403,0.81)rectangle++(0.008,0.008) (-0.477,0.646)rectangle++(0.008,0.008) (-0.406,0.744)rectangle++(0.008,0.008) (-0.455,0.614)rectangle++(0.008,0.008) (-0.149,0.898)rectangle++(0.008,0.008)
  (-0.393,0.777)rectangle++(0.008,0.008) (-0.028,0.878)rectangle++(0.008,0.008) (-0.355,0.856)rectangle++(0.008,0.008) (-0.331,0.892)rectangle++(0.008,0.008) (-0.264,0.909)rectangle++(0.008,0.008)
  (-0.241,0.915)rectangle++(0.008,0.008) (-0.441,0.751)rectangle++(0.008,0.008) (-0.473,0.605)rectangle++(0.008,0.008) (-0.312,0.813)rectangle++(0.008,0.008) (-0.172,0.943)rectangle++(0.008,0.008)
  (-0.339,0.898)rectangle++(0.008,0.008) (-0.392,0.857)rectangle++(0.008,0.008) (-0.489,0.671)rectangle++(0.008,0.008) (-0.264,0.742)rectangle++(0.008,0.008) (-0.127,0.898)rectangle++(0.008,0.008)
  (-0.12,0.564)rectangle++(0.008,0.008) (-0.104,0.851)rectangle++(0.008,0.008) (-0.187,0.895)rectangle++(0.008,0.008) (-0.135,0.591)rectangle++(0.008,0.008) (-0.047,0.872)rectangle++(0.008,0.008)
  (-0.261,0.907)rectangle++(0.008,0.008) (-0.376,0.884)rectangle++(0.008,0.008) (-0.491,0.689)rectangle++(0.008,0.008) (-0.153,0.686)rectangle++(0.008,0.008) (-0.204,0.84)rectangle++(0.008,0.008)
  (-0.397,0.73)rectangle++(0.008,0.008) (-0.209,0.914)rectangle++(0.008,0.008) (-0.153,0.604)rectangle++(0.008,0.008) (-0.152,0.73)rectangle++(0.008,0.008) (-0.093,0.897)rectangle++(0.008,0.008)
  (-0.184,0.911)rectangle++(0.008,0.008) (-0.39,0.839)rectangle++(0.008,0.008) (-0.029,0.906)rectangle++(0.008,0.008) (-0.108,0.804)rectangle++(0.008,0.008) (-0.287,0.904)rectangle++(0.008,0.008)
  (-0.089,0.923)rectangle++(0.008,0.008) (-0.374,0.809)rectangle++(0.008,0.008) (-0.272,0.754)rectangle++(0.008,0.008) (-0.397,0.667)rectangle++(0.008,0.008) (-0.423,0.582)rectangle++(0.008,0.008)
  (-0.142,0.597)rectangle++(0.008,0.008) (-0.335,0.863)rectangle++(0.008,0.008) (-0.265,0.9)rectangle++(0.008,0.008) (-0.277,0.762)rectangle++(0.008,0.008) (-0.098,0.926)rectangle++(0.008,0.008)
  (-0.145,0.594)rectangle++(0.008,0.008) (-0.201,0.809)rectangle++(0.008,0.008) (-0.129,0.57)rectangle++(0.008,0.008) (-0.092,0.525)rectangle++(0.008,0.008) (-0.184,0.786)rectangle++(0.008,0.008)
  (-0.371,0.815)rectangle++(0.008,0.008) (-0.464,0.659)rectangle++(0.008,0.008) (-0.451,0.555)rectangle++(0.008,0.008) (-0.019,0.79)rectangle++(0.008,0.008) (-0.294,0.773)rectangle++(0.008,0.008)
  (-0.391,0.796)rectangle++(0.008,0.008) (-0.466,0.644)rectangle++(0.008,0.008) (-0.27,0.757)rectangle++(0.008,0.008) (-0.103,0.898)rectangle++(0.008,0.008) (-0.372,0.793)rectangle++(0.008,0.008)
  (-0.333,0.883)rectangle++(0.008,0.008) (-0.471,0.704)rectangle++(0.008,0.008) (-0.414,0.756)rectangle++(0.008,0.008) (-0.463,0.617)rectangle++(0.008,0.008) (-0.384,0.827)rectangle++(0.008,0.008)
  (-0.335,0.888)rectangle++(0.008,0.008) (-0.474,0.706)rectangle++(0.008,0.008) (-0.117,0.547)rectangle++(0.008,0.008) (-0.109,0.908)rectangle++(0.008,0.008) (-0.377,0.795)rectangle++(0.008,0.008)
  (-0.46,0.648)rectangle++(0.008,0.008) (-0.146,0.904)rectangle++(0.008,0.008) (-0.117,0.935)rectangle++(0.008,0.008) (-0.192,0.924)rectangle++(0.008,0.008) (-0.421,0.773)rectangle++(0.008,0.008)
  (-0.007,0.858)rectangle++(0.008,0.008) (-0.106,0.931)rectangle++(0.008,0.008) (-0.378,0.855)rectangle++(0.008,0.008) (-0.146,0.951)rectangle++(0.008,0.008) (-0.41,0.802)rectangle++(0.008,0.008)
  (-0.417,0.784)rectangle++(0.008,0.008) (-0.474,0.629)rectangle++(0.008,0.008) (-0.312,0.821)rectangle++(0.008,0.008) (-0.258,0.885)rectangle++(0.008,0.008) (-0.317,0.897)rectangle++(0.008,0.008)
  (-0.469,0.716)rectangle++(0.008,0.008) (-0.473,0.579)rectangle++(0.008,0.008) (-0.427,0.515)rectangle++(0.008,0.008) (-0.382,0.502)rectangle++(0.008,0.008) (-0.358,0.512)rectangle++(0.008,0.008)
  (-0.273,0.782)rectangle++(0.008,0.008) (-0.244,0.872)rectangle++(0.008,0.008) (-0.427,0.731)rectangle++(0.008,0.008) (-0.072,0.91)rectangle++(0.008,0.008) (-0.113,0.567)rectangle++(0.008,0.008)
  (-0.06,0.884)rectangle++(0.008,0.008) (-0.369,0.873)rectangle++(0.008,0.008) (-0.484,0.687)rectangle++(0.008,0.008) (-0.114,0.543)rectangle++(0.008,0.008) (-0.047,0.472)rectangle++(0.008,0.008)
  (-0.155,0.942)rectangle++(0.008,0.008) (-0.088,0.892)rectangle++(0.008,0.008) (-0.293,0.897)rectangle++(0.008,0.008) (-0.458,0.725)rectangle++(0.008,0.008) (-0.28,0.769)rectangle++(0.008,0.008)
  (-0.406,0.671)rectangle++(0.008,0.008) (-0.122,0.794)rectangle++(0.008,0.008) (-0.194,0.876)rectangle++(0.008,0.008) (-0.379,0.864)rectangle++(0.008,0.008) (-0.154,0.659)rectangle++(0.008,0.008)
  (-0.347,0.795)rectangle++(0.008,0.008) (-0.446,0.659)rectangle++(0.008,0.008) (-0.443,0.561)rectangle++(0.008,0.008) (-0.18,0.928)rectangle++(0.008,0.008) (-0.006,0.859)rectangle++(0.008,0.008)
  (-0.313,0.81)rectangle++(0.008,0.008) (-0.231,0.858)rectangle++(0.008,0.008) (-0.149,0.695)rectangle++(0.008,0.008) (-0.32,0.684)rectangle++(0.008,0.008) (-0.382,0.775)rectangle++(0.008,0.008)
  (-0.281,0.87)rectangle++(0.008,0.008) (-0.443,0.716)rectangle++(0.008,0.008) (-0.462,0.589)rectangle++(0.008,0.008) (-0.051,0.88)rectangle++(0.008,0.008) (-0.341,0.803)rectangle++(0.008,0.008)
  (-0.004,0.855)rectangle++(0.008,0.008) (-0.154,0.916)rectangle++(0.008,0.008) (-0.067,0.909)rectangle++(0.008,0.008) (-0.358,0.811)rectangle++(0.008,0.008) (-0.409,0.793)rectangle++(0.008,0.008)
  (-0.029,0.447)rectangle++(0.008,0.008) (-0.061,0.872)rectangle++(0.008,0.008) (-0.173,0.903)rectangle++(0.008,0.008) (-0.142,0.588)rectangle++(0.008,0.008) (-0.279,0.637)rectangle++(0.008,0.008)
  (-0.067,0.897)rectangle++(0.008,0.008) (-0.355,0.805)rectangle++(0.008,0.008) (-0.453,0.661)rectangle++(0.008,0.008) (-0.447,0.559)rectangle++(0.008,0.008) (-0.217,0.889)rectangle++(0.008,0.008)
  (-0.141,0.948)rectangle++(0.008,0.008) (-0.048,0.885)rectangle++(0.008,0.008) (-0.076,0.51)rectangle++(0.008,0.008) (-0.221,0.625)rectangle++(0.008,0.008) (-0.354,0.778)rectangle++(0.008,0.008)
  (-0.443,0.649)rectangle++(0.008,0.008) (-0.302,0.827)rectangle++(0.008,0.008) (-0.161,0.951)rectangle++(0.008,0.008) (-0.417,0.796)rectangle++(0.008,0.008) (-0.027,0.881)rectangle++(0.008,0.008)
  (-0.284,0.896)rectangle++(0.008,0.008) (-0.407,0.821)rectangle++(0.008,0.008) (-0.144,0.598)rectangle++(0.008,0.008) (-0.096,0.53)rectangle++(0.008,0.008) (-0.359,0.881)rectangle++(0.008,0.008)
  (-0.419,0.807)rectangle++(0.008,0.008) (-0.12,0.924)rectangle++(0.008,0.008) (-0.389,0.798)rectangle++(0.008,0.008) (-0.283,0.878)rectangle++(0.008,0.008) (-0.355,0.866)rectangle++(0.008,0.008)
  (-0.416,0.805)rectangle++(0.008,0.008) (-0.293,0.88)rectangle++(0.008,0.008) (-0.452,0.717)rectangle++(0.008,0.008) (-0.466,0.586)rectangle++(0.008,0.008) (-0.309,0.806)rectangle++(0.008,0.008)
  (-0.257,0.881)rectangle++(0.008,0.008) (-0.435,0.73)rectangle++(0.008,0.008) (-0.299,0.855)rectangle++(0.008,0.008) (-0.137,0.578)rectangle++(0.008,0.008) (-0.272,0.634)rectangle++(0.008,0.008)
  (-0.114,0.556)rectangle++(0.008,0.008) (-0.254,0.633)rectangle++(0.008,0.008) (-0.122,0.571)rectangle++(0.008,0.008) (-0.064,0.493)rectangle++(0.008,0.008) (-0.209,0.622)rectangle++(0.008,0.008)
  (-0.136,0.92)rectangle++(0.008,0.008) (-0.394,0.791)rectangle++(0.008,0.008) (-0.317,0.817)rectangle++(0.008,0.008) (-0.44,0.679)rectangle++(0.008,0.008) (-0.404,0.756)rectangle++(0.008,0.008)
  (-0.458,0.62)rectangle++(0.008,0.008) (-0.136,0.579)rectangle++(0.008,0.008) (-0.273,0.636)rectangle++(0.008,0.008) (-0.279,0.911)rectangle++(0.008,0.008) (-0.253,0.903)rectangle++(0.008,0.008)
  (-0.402,0.827)rectangle++(0.008,0.008) (-0.027,0.897)rectangle++(0.008,0.008) (-0.336,0.819)rectangle++(0.008,0.008) (-0.328,0.887)rectangle++(0.008,0.008) (-0.015,0.874)rectangle++(0.008,0.008)
  (-0.136,0.584)rectangle++(0.008,0.008) (-0.274,0.638)rectangle++(0.008,0.008) (-0.327,0.832)rectangle++(0.008,0.008) (-0.45,0.682)rectangle++(0.008,0.008) (-0.14,0.585)rectangle++(0.008,0.008)
  (-0.276,0.636)rectangle++(0.008,0.008) (-0.03,0.788)rectangle++(0.008,0.008) (-0.153,0.74)rectangle++(0.008,0.008) (-0.338,0.703)rectangle++(0.008,0.008) (-0.409,0.62)rectangle++(0.008,0.008)
  (-0.32,0.884)rectangle++(0.008,0.008) (-0.157,0.733)rectangle++(0.008,0.008) (-0.329,0.833)rectangle++(0.008,0.008) (-0.149,0.597)rectangle++(0.008,0.008) (-0.203,0.81)rectangle++(0.008,0.008)
  (-0.386,0.717)rectangle++(0.008,0.008) (-0.436,0.609)rectangle++(0.008,0.008) (-0.377,0.863)rectangle++(0.008,0.008) (-0.484,0.679)rectangle++(0.008,0.008) (-0.413,0.749)rectangle++(0.008,0.008)
  (-0.46,0.614)rectangle++(0.008,0.008) (-0.307,0.816)rectangle++(0.008,0.008) (-0.256,0.884)rectangle++(0.008,0.008) (-0.436,0.732)rectangle++(0.008,0.008) (-0.342,0.875)rectangle++(0.008,0.008)
  (-0.472,0.697)rectangle++(0.008,0.008) (-0.312,0.844)rectangle++(0.008,0.008) (-0.136,0.592)rectangle++(0.008,0.008) (-0.277,0.642)rectangle++(0.008,0.008) (-0.359,0.614)rectangle++(0.008,0.008)
  (-0.331,0.858)rectangle++(0.008,0.008) (-0.132,0.941)rectangle++(0.008,0.008) (-0.172,0.943)rectangle++(0.008,0.008) (-0.393,0.847)rectangle++(0.008,0.008) (-0.485,0.666)rectangle++(0.008,0.008)
  (-0.463,0.55)rectangle++(0.008,0.008) (-0.412,0.505)rectangle++(0.008,0.008) (-0.373,0.727)rectangle++(0.008,0.008) (-0.119,0.552)rectangle++(0.008,0.008) (-0.288,0.889)rectangle++(0.008,0.008)
  (-0.453,0.723)rectangle++(0.008,0.008) (-0.468,0.588)rectangle++(0.008,0.008) (-0.235,0.889)rectangle++(0.008,0.008) (-0.155,0.863)rectangle++(0.008,0.008) (-0.223,0.878)rectangle++(0.008,0.008)
  (-0.228,0.905)rectangle++(0.008,0.008) (-0.431,0.751)rectangle++(0.008,0.008) (-0.138,0.94)rectangle++(0.008,0.008) (-0.334,0.837)rectangle++(0.008,0.008) (-0.455,0.683)rectangle++(0.008,0.008)
  (-0.408,0.754)rectangle++(0.008,0.008) (-0.459,0.618)rectangle++(0.008,0.008) (-0.434,0.538)rectangle++(0.008,0.008) (-0.342,0.847)rectangle++(0.008,0.008) (-0.019,0.874)rectangle++(0.008,0.008)
  (-0.353,0.857)rectangle++(0.008,0.008) (-0.471,0.685)rectangle++(0.008,0.008) (-0.356,0.893)rectangle++(0.008,0.008) (-0.485,0.7)rectangle++(0.008,0.008) (-0.151,0.606)rectangle++(0.008,0.008)
  (-0.338,0.785)rectangle++(0.008,0.008) (-0.438,0.658)rectangle++(0.008,0.008) (-0.439,0.563)rectangle++(0.008,0.008) (-0.357,0.889)rectangle++(0.008,0.008) (-0.331,0.897)rectangle++(0.008,0.008)
  (-0.15,0.728)rectangle++(0.008,0.008) (-0.155,0.853)rectangle++(0.008,0.008) (-0.376,0.832)rectangle++(0.008,0.008) (-0.154,0.658)rectangle++(0.008,0.008) (-0.271,0.902)rectangle++(0.008,0.008)
  (-0.244,0.913)rectangle++(0.008,0.008) (-0.188,0.939)rectangle++(0.008,0.008) (-0.425,0.781)rectangle++(0.008,0.008) (-0.477,0.625)rectangle++(0.008,0.008) (-0.151,0.903)rectangle++(0.008,0.008)
  (-0.395,0.778)rectangle++(0.008,0.008) (-0.462,0.634)rectangle++(0.008,0.008) (-0.441,0.544)rectangle++(0.008,0.008) (-0.399,0.51)rectangle++(0.008,0.008) (-0.151,0.719)rectangle++(0.008,0.008)
  (-0.302,0.873)rectangle++(0.008,0.008) (-0.151,0.601)rectangle++(0.008,0.008) (-0.097,0.53)rectangle++(0.008,0.008) (-0.34,0.871)rectangle++(0.008,0.008) (-0.101,0.924)rectangle++(0.008,0.008)
  (-0.333,0.837)rectangle++(0.008,0.008) (-0.455,0.683)rectangle++(0.008,0.008) (-0.408,0.754)rectangle++(0.008,0.008) (-0.268,0.75)rectangle++(0.008,0.008) (-0.394,0.666)rectangle++(0.008,0.008)
  (-0.041,0.875)rectangle++(0.008,0.008) (-0.357,0.853)rectangle++(0.008,0.008) (-0.472,0.682)rectangle++(0.008,0.008) (-0.411,0.752)rectangle++(0.008,0.008) (-0.46,0.616)rectangle++(0.008,0.008)
  (-0.434,0.537)rectangle++(0.008,0.008) (-0.155,0.884)rectangle++(0.008,0.008) (-0.109,0.935)rectangle++(0.008,0.008) (-0.133,0.589)rectangle++(0.008,0.008) (-0.275,0.641)rectangle++(0.008,0.008)
  (-0.358,0.614)rectangle++(0.008,0.008) (-0.387,0.572)rectangle++(0.008,0.008) (-0.385,0.543)rectangle++(0.008,0.008) (-0.131,0.574)rectangle++(0.008,0.008) (-0.142,0.885)rectangle++(0.008,0.008)
  (0.002,0.847)rectangle++(0.008,0.008) (-0.132,0.579)rectangle++(0.008,0.008) (-0.271,0.637)rectangle++(0.008,0.008) (-0.06,0.824)rectangle++(0.008,0.008) (-0.133,0.591)rectangle++(0.008,0.008)
  (-0.275,0.642)rectangle++(0.008,0.008) (-0.155,0.731)rectangle++(0.008,0.008) (-0.336,0.698)rectangle++(0.008,0.008) (-0.406,0.618)rectangle++(0.008,0.008) (-0.41,0.557)rectangle++(0.008,0.008)
  (-0.291,0.797)rectangle++(0.008,0.008) (-0.421,0.68)rectangle++(0.008,0.008) (-0.401,0.759)rectangle++(0.008,0.008) (-0.409,0.777)rectangle++(0.008,0.008) (-0.468,0.628)rectangle++(0.008,0.008)
  (-0.442,0.539)rectangle++(0.008,0.008) (-0.398,0.508)rectangle++(0.008,0.008) (-0.281,0.903)rectangle++(0.008,0.008) (-0.277,0.762)rectangle++(0.008,0.008) (-0.402,0.669)rectangle++(0.008,0.008)
  (-0.034,0.895)rectangle++(0.008,0.008) (-0.139,0.588)rectangle++(0.008,0.008) (-0.277,0.638)rectangle++(0.008,0.008) (-0.246,0.824)rectangle++(0.008,0.008) (-0.41,0.708)rectangle++(0.008,0.008)
  (-0.444,0.596)rectangle++(0.008,0.008) (-0.259,0.737)rectangle++(0.008,0.008) (-0.006,0.821)rectangle++(0.008,0.008) (-0.155,0.885)rectangle++(0.008,0.008) (-0.39,0.769)rectangle++(0.008,0.008)
  (-0.26,0.909)rectangle++(0.008,0.008) (-0.447,0.742)rectangle++(0.008,0.008) (-0.122,0.554)rectangle++(0.008,0.008) (-0.171,0.889)rectangle++(0.008,0.008) (-0.21,0.908)rectangle++(0.008,0.008)
  (0.003,0.813)rectangle++(0.008,0.008) (-0.293,0.792)rectangle++(0.008,0.008) (-0.127,0.566)rectangle++(0.008,0.008) (-0.264,0.632)rectangle++(0.008,0.008) (-0.369,0.84)rectangle++(0.008,0.008)
  (-0.277,0.892)rectangle++(0.008,0.008) (-0.449,0.728)rectangle++(0.008,0.008) (-0.248,0.898)rectangle++(0.008,0.008) (-0.437,0.741)rectangle++(0.008,0.008) (-0.468,0.602)rectangle++(0.008,0.008)
  (-0.433,0.527)rectangle++(0.008,0.008) (-0.154,0.723)rectangle++(0.008,0.008) (-0.264,0.852)rectangle++(0.008,0.008) (-0.428,0.714)rectangle++(0.008,0.008) (-0.297,0.85)rectangle++(0.008,0.008)
  (-0.018,0.885)rectangle++(0.008,0.008) (-0.159,0.907)rectangle++(0.008,0.008) (-0.017,0.867)rectangle++(0.008,0.008) (-0.364,0.875)rectangle++(0.008,0.008) (-0.104,0.928)rectangle++(0.008,0.008)
  (-0.382,0.806)rectangle++(0.008,0.008) (-0.131,0.567)rectangle++(0.008,0.008) (-0.266,0.632)rectangle++(0.008,0.008) (-0.351,0.613)rectangle++(0.008,0.008) (-0.383,0.574)rectangle++(0.008,0.008)
  (-0.378,0.744)rectangle++(0.008,0.008) (-0.442,0.625)rectangle++(0.008,0.008) (-0.429,0.547)rectangle++(0.008,0.008) (-0.395,0.516)rectangle++(0.008,0.008) (-0.087,0.512)rectangle++(0.008,0.008)
  (-0.008,0.806)rectangle++(0.008,0.008) (-0.154,0.722)rectangle++(0.008,0.008) (-0.067,0.861)rectangle++(0.008,0.008) (-0.342,0.789)rectangle++(0.008,0.008) (-0.402,0.792)rectangle++(0.008,0.008)
  (-0.47,0.638)rectangle++(0.008,0.008) (-0.346,0.891)rectangle++(0.008,0.008) (-0.48,0.703)rectangle++(0.008,0.008) (-0.474,0.569)rectangle++(0.008,0.008) (-0.382,0.822)rectangle++(0.008,0.008)
  (-0.134,0.894)rectangle++(0.008,0.008) (-0.384,0.78)rectangle++(0.008,0.008) (-0.136,0.592)rectangle++(0.008,0.008) (-0.155,0.692)rectangle++(0.008,0.008) (-0.322,0.68)rectangle++(0.008,0.008)
  (-0.394,0.615)rectangle++(0.008,0.008) (-0.152,0.731)rectangle++(0.008,0.008) (-0.001,0.805)rectangle++(0.008,0.008) (-0.153,0.88)rectangle++(0.008,0.008) (-0.388,0.767)rectangle++(0.008,0.008)
  (-0.14,0.938)rectangle++(0.008,0.008) (-0.403,0.798)rectangle++(0.008,0.008) (-0.473,0.64)rectangle++(0.008,0.008) (-0.312,0.825)rectangle++(0.008,0.008) (-0.441,0.685)rectangle++(0.008,0.008)
  (-0.138,0.908)rectangle++(0.008,0.008) (-0.155,0.737)rectangle++(0.008,0.008) (-0.205,0.857)rectangle++(0.008,0.008) (-0.403,0.738)rectangle++(0.008,0.008) (-0.451,0.612)rectangle++(0.008,0.008)
  (-0.397,0.741)rectangle++(0.008,0.008) (-0.405,0.774)rectangle++(0.008,0.008) (-0.465,0.628)rectangle++(0.008,0.008) (-0.152,0.933)rectangle++(0.008,0.008) (-0.387,0.848)rectangle++(0.008,0.008)
  (-0.312,0.891)rectangle++(0.008,0.008) (-0.057,0.915)rectangle++(0.008,0.008) (-0.052,0.897)rectangle++(0.008,0.008) (-0.17,0.911)rectangle++(0.008,0.008) (-0.116,0.559)rectangle++(0.008,0.008)
  (-0.254,0.735)rectangle++(0.008,0.008) (-0.118,0.555)rectangle++(0.008,0.008) (-0.192,0.796)rectangle++(0.008,0.008) (-0.127,0.568)rectangle++(0.008,0.008) (-0.15,0.948)rectangle++(0.008,0.008)
  (-0.411,0.799)rectangle++(0.008,0.008) (-0.245,0.899)rectangle++(0.008,0.008) (-0.437,0.743)rectangle++(0.008,0.008) (-0.3,0.859)rectangle++(0.008,0.008) (-0.386,0.856)rectangle++(0.008,0.008)
  (-0.283,0.897)rectangle++(0.008,0.008) (-0.453,0.728)rectangle++(0.008,0.008) (-0.12,0.552)rectangle++(0.008,0.008) (-0.255,0.629)rectangle++(0.008,0.008) (-0.117,0.93)rectangle++(0.008,0.008)
  (-0.297,0.903)rectangle++(0.008,0.008) (-0.462,0.726)rectangle++(0.008,0.008) (-0.474,0.586)rectangle++(0.008,0.008) (-0.369,0.883)rectangle++(0.008,0.008) (-0.421,0.806)rectangle++(0.008,0.008)
  (-0.484,0.637)rectangle++(0.008,0.008) (-0.452,0.538)rectangle++(0.008,0.008) (-0.386,0.727)rectangle++(0.008,0.008) (-0.401,0.773)rectangle++(0.008,0.008) (-0.229,0.892)rectangle++(0.008,0.008)
  (-0.23,0.909)rectangle++(0.008,0.008) (-0.154,0.662)rectangle++(0.008,0.008) (-0.31,0.667)rectangle++(0.008,0.008) (-0.257,0.834)rectangle++(0.008,0.008) (-0.392,0.813)rectangle++(0.008,0.008)
  (-0.209,0.925)rectangle++(0.008,0.008) (-0.152,0.738)rectangle++(0.008,0.008) (-0.155,0.737)rectangle++(0.008,0.008) (-0.337,0.701)rectangle++(0.008,0.008) (-0.388,0.776)rectangle++(0.008,0.008)
  (-0.132,0.947)rectangle++(0.008,0.008) (-0.402,0.805)rectangle++(0.008,0.008) (-0.475,0.644)rectangle++(0.008,0.008) (-0.045,0.883)rectangle++(0.008,0.008) (-0.34,0.807)rectangle++(0.008,0.008)
  (-0.447,0.667)rectangle++(0.008,0.008) (-0.158,0.703)rectangle++(0.008,0.008) (-0.327,0.684)rectangle++(0.008,0.008) (-0.263,0.839)rectangle++(0.008,0.008) (-0.423,0.709)rectangle++(0.008,0.008)
  (-0.281,0.83)rectangle++(0.008,0.008) (-0.15,0.719)rectangle++(0.008,0.008) (-0.114,0.553)rectangle++(0.008,0.008) (-0.253,0.631)rectangle++(0.008,0.008) (-0.287,0.907)rectangle++(0.008,0.008)
  (-0.094,0.92)rectangle++(0.008,0.008) (-0.34,0.894)rectangle++(0.008,0.008) (-0.361,0.892)rectangle++(0.008,0.008) (-0.274,0.909)rectangle++(0.008,0.008) (-0.454,0.737)rectangle++(0.008,0.008)
  (-0.474,0.594)rectangle++(0.008,0.008) (-0.432,0.522)rectangle++(0.008,0.008) (-0.15,0.686)rectangle++(0.008,0.008) (-0.35,0.801)rectangle++(0.008,0.008) (-0.308,0.877)rectangle++(0.008,0.008)
  (-0.095,0.921)rectangle++(0.008,0.008) (-0.376,0.806)rectangle++(0.008,0.008) (-0.279,0.881)rectangle++(0.008,0.008) (-0.147,0.928)rectangle++(0.008,0.008) (-0.402,0.791)rectangle++(0.008,0.008)
  (-0.47,0.637)rectangle++(0.008,0.008) (-0.182,0.933)rectangle++(0.008,0.008) (-0.362,0.891)rectangle++(0.008,0.008) (-0.302,0.909)rectangle++(0.008,0.008) (0.006,0.847)rectangle++(0.008,0.008)
  (-0.258,0.899)rectangle++(0.008,0.008) (-0.403,0.825)rectangle++(0.008,0.008) (-0.099,0.924)rectangle++(0.008,0.008) (-0.379,0.806)rectangle++(0.008,0.008) (-0.28,0.881)rectangle++(0.008,0.008)
  (-0.446,0.722)rectangle++(0.008,0.008) (-0.465,0.59)rectangle++(0.008,0.008) (-0.154,0.896)rectangle++(0.008,0.008) (-0.394,0.774)rectangle++(0.008,0.008) (-0.267,0.9)rectangle++(0.008,0.008)
  (-0.239,0.899)rectangle++(0.008,0.008) (-0.314,0.898)rectangle++(0.008,0.008) (-0.468,0.718)rectangle++(0.008,0.008) (-0.31,0.851)rectangle++(0.008,0.008) (-0.14,0.594)rectangle++(0.008,0.008)
  (-0.2,0.809)rectangle++(0.008,0.008) (-0.09,0.539)rectangle++(0.008,0.008) (-0.237,0.633)rectangle++(0.008,0.008) (-0.232,0.822)rectangle++(0.008,0.008) (-0.404,0.712)rectangle++(0.008,0.008)
  (-0.442,0.601)rectangle++(0.008,0.008) (-0.393,0.741)rectangle++(0.008,0.008) (-0.405,0.775)rectangle++(0.008,0.008) (-0.289,0.87)rectangle++(0.008,0.008) (-0.446,0.714)rectangle++(0.008,0.008)
  (-0.462,0.586)rectangle++(0.008,0.008) (-0.308,0.806)rectangle++(0.008,0.008) (-0.432,0.678)rectangle++(0.008,0.008) (-0.443,0.575)rectangle++(0.008,0.008) (-0.302,0.803)rectangle++(0.008,0.008)
  (-0.428,0.678)rectangle++(0.008,0.008) (-0.121,0.935)rectangle++(0.008,0.008) (-0.151,0.698)rectangle++(0.008,0.008) (-0.322,0.684)rectangle++(0.008,0.008) (-0.133,0.94)rectangle++(0.008,0.008)
  (-0.197,0.926)rectangle++(0.008,0.008) (-0.099,0.882)rectangle++(0.008,0.008) (-0.295,0.896)rectangle++(0.008,0.008) (-0.154,0.904)rectangle++(0.008,0.008) (-0.397,0.777)rectangle++(0.008,0.008)
  (-0.313,0.887)rectangle++(0.008,0.008) (-0.463,0.713)rectangle++(0.008,0.008) (-0.328,0.829)rectangle++(0.008,0.008) (0.008,0.844)rectangle++(0.008,0.008) (-0.302,0.808)rectangle++(0.008,0.008)
  (-0.43,0.681)rectangle++(0.008,0.008) (-0.297,0.838)rectangle++(0.008,0.008) (-0.401,0.808)rectangle++(0.008,0.008) (-0.475,0.645)rectangle++(0.008,0.008) (-0.152,0.935)rectangle++(0.008,0.008)
  (-0.063,0.492)rectangle++(0.008,0.008) (-0.101,0.552)rectangle++(0.008,0.008) (-0.247,0.635)rectangle++(0.008,0.008) (-0.361,0.777)rectangle++(0.008,0.008) (-0.446,0.645)rectangle++(0.008,0.008)
  (-0.281,0.793)rectangle++(0.008,0.008) (-0.415,0.681)rectangle++(0.008,0.008) (-0.437,0.582)rectangle++(0.008,0.008) (-0.1,0.55)rectangle++(0.008,0.008) (-0.186,0.794)rectangle++(0.008,0.008)
  (-0.144,0.6)rectangle++(0.008,0.008) (-0.257,0.738)rectangle++(0.008,0.008) (-0.385,0.665)rectangle++(0.008,0.008) (-0.417,0.586)rectangle++(0.008,0.008) (-0.046,0.47)rectangle++(0.008,0.008)
  (-0.053,0.831)rectangle++(0.008,0.008) (-0.358,0.892)rectangle++(0.008,0.008) (-0.334,0.853)rectangle++(0.008,0.008) (-0.41,0.805)rectangle++(0.008,0.008) (-0.392,0.845)rectangle++(0.008,0.008)
  (-0.143,0.598)rectangle++(0.008,0.008) (-0.104,0.934)rectangle++(0.008,0.008) (-0.154,0.609)rectangle++(0.008,0.008) (-0.359,0.845)rectangle++(0.008,0.008) (-0.335,0.884)rectangle++(0.008,0.008)
  (-0.266,0.756)rectangle++(0.008,0.008) (-0.395,0.67)rectangle++(0.008,0.008) (-0.423,0.584)rectangle++(0.008,0.008) (-0.379,0.826)rectangle++(0.008,0.008) (-0.109,0.563)rectangle++(0.008,0.008)
  (-0.042,0.898)rectangle++(0.008,0.008) (-0.155,0.953)rectangle++(0.008,0.008) (-0.086,0.902)rectangle++(0.008,0.008) (-0.282,0.825)rectangle++(0.008,0.008) (-0.247,0.887)rectangle++(0.008,0.008)
  (-0.236,0.908)rectangle++(0.008,0.008) (-0.435,0.75)rectangle++(0.008,0.008) (-0.299,0.862)rectangle++(0.008,0.008) (-0.387,0.857)rectangle++(0.008,0.008) (-0.486,0.673)rectangle++(0.008,0.008)
  (-0.316,0.836)rectangle++(0.008,0.008) (-0.318,0.819)rectangle++(0.008,0.008) (-0.098,0.55)rectangle++(0.008,0.008) (-0.122,0.575)rectangle++(0.008,0.008) (-0.083,0.519)rectangle++(0.008,0.008)
  (-0.227,0.627)rectangle++(0.008,0.008) (-0.356,0.778)rectangle++(0.008,0.008) (-0.444,0.648)rectangle++(0.008,0.008) (-0.154,0.715)rectangle++(0.008,0.008) (-0.355,0.806)rectangle++(0.008,0.008)
  (-0.14,0.593)rectangle++(0.008,0.008) (-0.136,0.772)rectangle++(0.008,0.008) (-0.342,0.724)rectangle++(0.008,0.008) (-0.268,0.853)rectangle++(0.008,0.008) (-0.43,0.713)rectangle++(0.008,0.008)
  (-0.455,0.592)rectangle++(0.008,0.008) (-0.165,0.945)rectangle++(0.008,0.008) (-0.208,0.927)rectangle++(0.008,0.008) (-0.397,0.838)rectangle++(0.008,0.008) (-0.21,0.929)rectangle++(0.008,0.008)
  (-0.183,0.942)rectangle++(0.008,0.008) (-0.423,0.784)rectangle++(0.008,0.008) (-0.329,0.894)rectangle++(0.008,0.008) (-0.474,0.71)rectangle++(0.008,0.008) (-0.14,0.915)rectangle++(0.008,0.008)
  (-0.156,0.61)rectangle++(0.008,0.008) (-0.098,0.532)rectangle++(0.008,0.008) (-0.154,0.842)rectangle++(0.008,0.008) (-0.375,0.749)rectangle++(0.008,0.008) (-0.279,0.861)rectangle++(0.008,0.008)
  (-0.138,0.579)rectangle++(0.008,0.008) (-0.199,0.804)rectangle++(0.008,0.008) (-0.382,0.716)rectangle++(0.008,0.008) (-0.335,0.873)rectangle++(0.008,0.008) (-0.468,0.699)rectangle++(0.008,0.008)
  (-0.039,0.884)rectangle++(0.008,0.008) (-0.154,0.745)rectangle++(0.008,0.008) (-0.34,0.705)rectangle++(0.008,0.008) (-0.411,0.62)rectangle++(0.008,0.008) (-0.39,0.749)rectangle++(0.008,0.008)
  (-0.387,0.842)rectangle++(0.008,0.008) (-0.419,0.795)rectangle++(0.008,0.008) (-0.251,0.737)rectangle++(0.008,0.008) (-0.219,0.86)rectangle++(0.008,0.008) (-0.156,0.73)rectangle++(0.008,0.008)
  (-0.175,0.914)rectangle++(0.008,0.008) (-0.212,0.917)rectangle++(0.008,0.008) (-0.396,0.836)rectangle++(0.008,0.008) (-0.286,0.89)rectangle++(0.008,0.008) (-0.407,0.82)rectangle++(0.008,0.008)
  (-0.482,0.649)rectangle++(0.008,0.008) (-0.327,0.827)rectangle++(0.008,0.008) (-0.448,0.68)rectangle++(0.008,0.008) (-0.451,0.57)rectangle++(0.008,0.008) (-0.13,0.578)rectangle++(0.008,0.008)
  (-0.155,0.69)rectangle++(0.008,0.008) (-0.321,0.679)rectangle++(0.008,0.008) (-0.199,0.906)rectangle++(0.008,0.008) (-0.362,0.891)rectangle++(0.008,0.008) (-0.274,0.909)rectangle++(0.008,0.008)
  (-0.407,0.825)rectangle++(0.008,0.008) (-0.393,0.847)rectangle++(0.008,0.008) (0.004,0.826)rectangle++(0.008,0.008) (-0.334,0.896)rectangle++(0.008,0.008) (-0.476,0.709)rectangle++(0.008,0.008)
  (-0.474,0.573)rectangle++(0.008,0.008) (-0.425,0.512)rectangle++(0.008,0.008) (-0.381,0.501)rectangle++(0.008,0.008) (-0.281,0.778)rectangle++(0.008,0.008) (-0.259,0.901)rectangle++(0.008,0.008)
  (-0.336,0.9)rectangle++(0.008,0.008) (-0.096,0.928)rectangle++(0.008,0.008) (-0.379,0.809)rectangle++(0.008,0.008) (-0.257,0.9)rectangle++(0.008,0.008) (-0.443,0.739)rectangle++(0.008,0.008)
  (-0.062,0.866)rectangle++(0.008,0.008) (-0.341,0.793)rectangle++(0.008,0.008) (-0.119,0.563)rectangle++(0.008,0.008) (-0.253,0.739)rectangle++(0.008,0.008) (-0.238,0.858)rectangle++(0.008,0.008)
  (-0.126,0.937)rectangle++(0.008,0.008) (-0.395,0.802)rectangle++(0.008,0.008) (-0.471,0.645)rectangle++(0.008,0.008) (-0.334,0.899)rectangle++(0.008,0.008) (-0.478,0.711)rectangle++(0.008,0.008)
  (-0.381,0.865)rectangle++(0.008,0.008) (-0.487,0.679)rectangle++(0.008,0.008) (-0.413,0.749)rectangle++(0.008,0.008) (-0.292,0.861)rectangle++(0.008,0.008) (-0.021,0.835)rectangle++(0.008,0.008)
  (-0.311,0.793)rectangle++(0.008,0.008) (-0.128,0.566)rectangle++(0.008,0.008) (-0.105,0.852)rectangle++(0.008,0.008) (-0.28,0.821)rectangle++(0.008,0.008) (-0.395,0.807)rectangle++(0.008,0.008)
  (-0.286,0.881)rectangle++(0.008,0.008) (-0.091,0.919)rectangle++(0.008,0.008) (-0.183,0.918)rectangle++(0.008,0.008) (-0.415,0.774)rectangle++(0.008,0.008) (-0.158,0.613)rectangle++(0.008,0.008)
  (-0.295,0.643)rectangle++(0.008,0.008) (-0.368,0.608)rectangle++(0.008,0.008) (-0.28,0.819)rectangle++(0.008,0.008) (-0.144,0.948)rectangle++(0.008,0.008) (-0.201,0.928)rectangle++(0.008,0.008)
  (-0.427,0.772)rectangle++(0.008,0.008) (-0.416,0.776)rectangle++(0.008,0.008) (-0.471,0.625)rectangle++(0.008,0.008) (-0.159,0.697)rectangle++(0.008,0.008) (-0.368,0.887)rectangle++(0.008,0.008)
  (-0.422,0.807)rectangle++(0.008,0.008) (-0.484,0.637)rectangle++(0.008,0.008) (-0.407,0.741)rectangle++(0.008,0.008) (-0.088,0.526)rectangle++(0.008,0.008) (-0.231,0.628)rectangle++(0.008,0.008)
  (-0.049,0.872)rectangle++(0.008,0.008) (-0.169,0.903)rectangle++(0.008,0.008) (-0.209,0.913)rectangle++(0.008,0.008) (-0.237,0.898)rectangle++(0.008,0.008) (-0.433,0.745)rectangle++(0.008,0.008)
  (-0.144,0.59)rectangle++(0.008,0.008) (-0.28,0.637)rectangle++(0.008,0.008) (-0.119,0.881)rectangle++(0.008,0.008) (-0.138,0.585)rectangle++(0.008,0.008) (-0.112,0.559)rectangle++(0.008,0.008)
  (-0.19,0.797)rectangle++(0.008,0.008) (-0.266,0.894)rectangle++(0.008,0.008) (-0.444,0.733)rectangle++(0.008,0.008) (-0.091,0.54)rectangle++(0.008,0.008) (-0.087,0.52)rectangle++(0.008,0.008)
  (-0.229,0.626)rectangle++(0.008,0.008) (-0.332,0.624)rectangle++(0.008,0.008) (-0.143,0.926)rectangle++(0.008,0.008) (-0.374,0.886)rectangle++(0.008,0.008) (-0.119,0.575)rectangle++(0.008,0.008)
  (-0.298,0.853)rectangle++(0.008,0.008) (-0.444,0.703)rectangle++(0.008,0.008) (-0.135,0.59)rectangle++(0.008,0.008) (-0.035,0.882)rectangle++(0.008,0.008) (-0.366,0.875)rectangle++(0.008,0.008)
  (-0.42,0.805)rectangle++(0.008,0.008) (-0.483,0.637)rectangle++(0.008,0.008) (-0.097,0.542)rectangle++(0.008,0.008) (-0.242,0.633)rectangle++(0.008,0.008) (-0.34,0.622)rectangle++(0.008,0.008)
  (-0.1,0.551)rectangle++(0.008,0.008) (-0.295,0.849)rectangle++(0.008,0.008) (-0.153,0.845)rectangle++(0.008,0.008) (-0.374,0.831)rectangle++(0.008,0.008) (-0.471,0.665)rectangle++(0.008,0.008)
  (-0.346,0.892)rectangle++(0.008,0.008) (-0.269,0.909)rectangle++(0.008,0.008) (-0.243,0.915)rectangle++(0.008,0.008) (-0.442,0.751)rectangle++(0.008,0.008) (-0.301,0.862)rectangle++(0.008,0.008)
  (-0.257,0.901)rectangle++(0.008,0.008) (-0.443,0.739)rectangle++(0.008,0.008) (-0.122,0.554)rectangle++(0.008,0.008) (-0.257,0.629)rectangle++(0.008,0.008) (-0.282,0.871)rectangle++(0.008,0.008)
  (-0.444,0.717)rectangle++(0.008,0.008) (-0.411,0.763)rectangle++(0.008,0.008) (-0.464,0.621)rectangle++(0.008,0.008) (-0.15,0.901)rectangle++(0.008,0.008) (-0.068,0.906)rectangle++(0.008,0.008)
  (-0.276,0.763)rectangle++(0.008,0.008) (-0.019,0.872)rectangle++(0.008,0.008) (-0.283,0.895)rectangle++(0.008,0.008) (-0.036,0.457)rectangle++(0.008,0.008) (-0.154,0.673)rectangle++(0.008,0.008)
  (-0.024,0.785)rectangle++(0.008,0.008) (-0.017,0.838)rectangle++(0.008,0.008) (-0.347,0.85)rectangle++(0.008,0.008) (-0.412,0.803)rectangle++(0.008,0.008) (-0.131,0.566)rectangle++(0.008,0.008)
  (-0.266,0.632)rectangle++(0.008,0.008) (-0.118,0.878)rectangle++(0.008,0.008) (-0.371,0.779)rectangle++(0.008,0.008) (-0.451,0.643)rectangle++(0.008,0.008) (-0.145,0.902)rectangle++(0.008,0.008)
  (-0.145,0.931)rectangle++(0.008,0.008) (-0.202,0.922)rectangle++(0.008,0.008) (-0.146,0.592)rectangle++(0.008,0.008) (-0.335,0.783)rectangle++(0.008,0.008) (-0.265,0.873)rectangle++(0.008,0.008)
  (-0.092,0.542)rectangle++(0.008,0.008) (-0.239,0.634)rectangle++(0.008,0.008) (-0.329,0.871)rectangle++(0.008,0.008) (-0.465,0.7)rectangle++(0.008,0.008) (-0.309,0.845)rectangle++(0.008,0.008)
  (-0.446,0.695)rectangle++(0.008,0.008) (-0.303,0.843)rectangle++(0.008,0.008) (-0.443,0.696)rectangle++(0.008,0.008) (-0.302,0.844)rectangle++(0.008,0.008) (-0.323,0.891)rectangle++(0.008,0.008)
  (-0.261,0.909)rectangle++(0.008,0.008) (-0.087,0.893)rectangle++(0.008,0.008) (-0.02,0.883)rectangle++(0.008,0.008) (-0.137,0.586)rectangle++(0.008,0.008) (-0.276,0.638)rectangle++(0.008,0.008)
  (-0.154,0.869)rectangle++(0.008,0.008) (-0.061,0.824)rectangle++(0.008,0.008) (-0.365,0.869)rectangle++(0.008,0.008) (-0.105,0.927)rectangle++(0.008,0.008) (-0.168,0.941)rectangle++(0.008,0.008)
  (-0.417,0.789)rectangle++(0.008,0.008) (-0.476,0.631)rectangle++(0.008,0.008) (-0.447,0.538)rectangle++(0.008,0.008) (-0.4,0.506)rectangle++(0.008,0.008) (-0.129,0.582)rectangle++(0.008,0.008)
  (-0.272,0.758)rectangle++(0.008,0.008) (-0.334,0.9)rectangle++(0.008,0.008) (-0.152,0.742)rectangle++(0.008,0.008) (-0.302,0.876)rectangle++(0.008,0.008) (-0.324,0.894)rectangle++(0.008,0.008)
  (-0.129,0.943)rectangle++(0.008,0.008) (-0.149,0.711)rectangle++(0.008,0.008) (-0.334,0.859)rectangle++(0.008,0.008) (-0.15,0.599)rectangle++(0.008,0.008) (-0.337,0.784)rectangle++(0.008,0.008)
  (-0.201,0.921)rectangle++(0.008,0.008) (-0.424,0.768)rectangle++(0.008,0.008) (-0.269,0.751)rectangle++(0.008,0.008) (-0.395,0.666)rectangle++(0.008,0.008) (-0.422,0.583)rectangle++(0.008,0.008)
  (-0.405,0.535)rectangle++(0.008,0.008) (-0.055,0.873)rectangle++(0.008,0.008) (-0.34,0.798)rectangle++(0.008,0.008) (-0.157,0.721)rectangle++(0.008,0.008) (-0.206,0.852)rectangle++(0.008,0.008)
  (-0.135,0.578)rectangle++(0.008,0.008) (-0.198,0.804)rectangle++(0.008,0.008) (-0.181,0.924)rectangle++(0.008,0.008) (-0.334,0.86)rectangle++(0.008,0.008) (-0.154,0.609)rectangle++(0.008,0.008)
  (-0.112,0.562)rectangle++(0.008,0.008) (-0.255,0.636)rectangle++(0.008,0.008) (-0.347,0.619)rectangle++(0.008,0.008) (-0.383,0.578)rectangle++(0.008,0.008) (-0.148,0.599)rectangle++(0.008,0.008)
  (-0.285,0.64)rectangle++(0.008,0.008) (-0.321,0.862)rectangle++(0.008,0.008) (-0.126,0.577)rectangle++(0.008,0.008) (-0.125,0.787)rectangle++(0.008,0.008) (-0.342,0.734)rectangle++(0.008,0.008)
  (-0.066,0.906)rectangle++(0.008,0.008) (-0.371,0.875)rectangle++(0.008,0.008) (-0.151,0.6)rectangle++(0.008,0.008) (-0.302,0.856)rectangle++(0.008,0.008) (-0.137,0.578)rectangle++(0.008,0.008)
  (-0.155,0.844)rectangle++(0.008,0.008) (-0.303,0.891)rectangle++(0.008,0.008) (-0.281,0.816)rectangle++(0.008,0.008) (-0.394,0.806)rectangle++(0.008,0.008) (-0.023,0.894)rectangle++(0.008,0.008)
  (-0.102,0.553)rectangle++(0.008,0.008) (-0.248,0.636)rectangle++(0.008,0.008) (-0.188,0.9)rectangle++(0.008,0.008) (-0.411,0.763)rectangle++(0.008,0.008) (-0.157,0.612)rectangle++(0.008,0.008)
  (-0.294,0.643)rectangle++(0.008,0.008) (-0.079,0.525)rectangle++(0.008,0.008) (-0.039,0.872)rectangle++(0.008,0.008) (0.004,0.821)rectangle++(0.008,0.008) (-0.239,0.883)rectangle++(0.008,0.008)
  (-0.383,0.862)rectangle++(0.008,0.008) (-0.241,0.72)rectangle++(0.008,0.008) (-0.077,0.522)rectangle++(0.008,0.008) (-0.179,0.785)rectangle++(0.008,0.008) (-0.15,0.712)rectangle++(0.008,0.008)
  (-0.343,0.882)rectangle++(0.008,0.008) (-0.202,0.935)rectangle++(0.008,0.008) (-0.221,0.924)rectangle++(0.008,0.008) (-0.434,0.763)rectangle++(0.008,0.008) (-0.143,0.94)rectangle++(0.008,0.008)
  (-0.405,0.798)rectangle++(0.008,0.008) (-0.313,0.89)rectangle++(0.008,0.008) (-0.412,0.816)rectangle++(0.008,0.008) (-0.419,0.786)rectangle++(0.008,0.008) (-0.476,0.629)rectangle++(0.008,0.008)
  (-0.111,0.936)rectangle++(0.008,0.008) (-0.388,0.807)rectangle++(0.008,0.008) (-0.414,0.788)rectangle++(0.008,0.008) (-0.212,0.922)rectangle++(0.008,0.008) (-0.43,0.765)rectangle++(0.008,0.008)
  (-0.053,0.478)rectangle++(0.008,0.008) (-0.3,0.775)rectangle++(0.008,0.008) (-0.115,0.568)rectangle++(0.008,0.008) (-0.129,0.585)rectangle++(0.008,0.008) (-0.196,0.806)rectangle++(0.008,0.008)
  (-0.108,0.921)rectangle++(0.008,0.008) (-0.189,0.919)rectangle++(0.008,0.008) (-0.159,0.953)rectangle++(0.008,0.008) (-0.416,0.798)rectangle++(0.008,0.008) (-0.293,0.878)rectangle++(0.008,0.008)
  (-0.352,0.892)rectangle++(0.008,0.008) (-0.1,0.929)rectangle++(0.008,0.008) (-0.282,0.826)rectangle++(0.008,0.008) (-0.247,0.887)rectangle++(0.008,0.008) (-0.433,0.736)rectangle++(0.008,0.008)
  (-0.298,0.857)rectangle++(0.008,0.008) (-0.253,0.898)rectangle++(0.008,0.008) (-0.282,0.824)rectangle++(0.008,0.008) (-0.136,0.59)rectangle++(0.008,0.008) (-0.224,0.859)rectangle++(0.008,0.008)
  (-0.283,0.773)rectangle++(0.008,0.008) (-0.248,0.869)rectangle++(0.008,0.008) (-0.261,0.751)rectangle++(0.008,0.008) (-0.391,0.669)rectangle++(0.008,0.008) (-0.327,0.827)rectangle++(0.008,0.008)
  (-0.281,0.878)rectangle++(0.008,0.008) (-0.147,0.6)rectangle++(0.008,0.008) (-0.234,0.871)rectangle++(0.008,0.008) (-0.218,0.848)rectangle++(0.008,0.008) (-0.406,0.729)rectangle++(0.008,0.008)
  (-0.449,0.608)rectangle++(0.008,0.008) (-0.426,0.536)rectangle++(0.008,0.008) (-0.014,0.836)rectangle++(0.008,0.008) (-0.283,0.89)rectangle++(0.008,0.008) (-0.45,0.725)rectangle++(0.008,0.008)
  (-0.468,0.59)rectangle++(0.008,0.008) (-0.428,0.522)rectangle++(0.008,0.008) (-0.13,0.572)rectangle++(0.008,0.008) (-0.147,0.601)rectangle++(0.008,0.008) (-0.162,0.93)rectangle++(0.008,0.008)
  (-0.136,0.592)rectangle++(0.008,0.008) (-0.277,0.641)rectangle++(0.008,0.008) (-0.045,0.468)rectangle++(0.008,0.008) (-0.191,0.617)rectangle++(0.008,0.008) (-0.311,0.633)rectangle++(0.008,0.008)
  (-0.373,0.766)rectangle++(0.008,0.008) (-0.333,0.88)rectangle++(0.008,0.008) (-0.47,0.702)rectangle++(0.008,0.008) (-0.469,0.572)rectangle++(0.008,0.008) (-0.097,0.522)rectangle++(0.008,0.008)
  (-0.234,0.623)rectangle++(0.008,0.008) (-0.065,0.863)rectangle++(0.008,0.008) (-0.059,0.89)rectangle++(0.008,0.008) (-0.348,0.805)rectangle++(0.008,0.008) (-0.45,0.663)rectangle++(0.008,0.008)
  (-0.142,0.905)rectangle++(0.008,0.008) (-0.262,0.908)rectangle++(0.008,0.008) (-0.448,0.741)rectangle++(0.008,0.008) (-0.115,0.57)rectangle++(0.008,0.008) (-0.005,0.791)rectangle++(0.008,0.008)
  (-0.338,0.843)rectangle++(0.008,0.008) (-0.153,0.691)rectangle++(0.008,0.008) (-0.32,0.68)rectangle++(0.008,0.008) (-0.173,0.935)rectangle++(0.008,0.008) (-0.147,0.594)rectangle++(0.008,0.008)
  (-0.174,0.894)rectangle++(0.008,0.008) (-0.178,0.937)rectangle++(0.008,0.008) (-0.147,0.595)rectangle++(0.008,0.008) (-0.283,0.638)rectangle++(0.008,0.008) (-0.361,0.61)rectangle++(0.008,0.008)
  (-0.386,0.569)rectangle++(0.008,0.008) (-0.025,0.891)rectangle++(0.008,0.008) (-0.333,0.818)rectangle++(0.008,0.008) (-0.448,0.674)rectangle++(0.008,0.008) (-0.018,0.784)rectangle++(0.008,0.008)
  (-0.291,0.771)rectangle++(0.008,0.008) (-0.133,0.59)rectangle++(0.008,0.008) (-0.332,0.785)rectangle++(0.008,0.008) (-0.264,0.873)rectangle++(0.008,0.008) (-0.318,0.895)rectangle++(0.008,0.008)
  (-0.136,0.894)rectangle++(0.008,0.008) (-0.378,0.843)rectangle++(0.008,0.008) (-0.27,0.761)rectangle++(0.008,0.008) (-0.399,0.671)rectangle++(0.008,0.008) (-0.395,0.761)rectangle++(0.008,0.008)
  (-0.304,0.901)rectangle++(0.008,0.008) (-0.464,0.722)rectangle++(0.008,0.008) (-0.181,0.938)rectangle++(0.008,0.008) (-0.306,0.904)rectangle++(0.008,0.008) (-0.146,0.95)rectangle++(0.008,0.008)
  (-0.282,0.827)rectangle++(0.008,0.008) (-0.428,0.697)rectangle++(0.008,0.008) (-0.154,0.65)rectangle++(0.008,0.008) (-0.104,0.54)rectangle++(0.008,0.008) (-0.071,0.856)rectangle++(0.008,0.008)
  (-0.176,0.897)rectangle++(0.008,0.008) (-0.405,0.766)rectangle++(0.008,0.008) (-0.462,0.625)rectangle++(0.008,0.008) (-0.014,0.859)rectangle++(0.008,0.008) (-0.155,0.932)rectangle++(0.008,0.008)
  (-0.407,0.79)rectangle++(0.008,0.008) (-0.472,0.635)rectangle++(0.008,0.008) (-0.311,0.823)rectangle++(0.008,0.008) (-0.257,0.886)rectangle++(0.008,0.008) (-0.438,0.732)rectangle++(0.008,0.008)
  (-0.145,0.939)rectangle++(0.008,0.008) (-0.251,0.734)rectangle++(0.008,0.008) (-0.118,0.555)rectangle++(0.008,0.008) (-0.192,0.796)rectangle++(0.008,0.008) (-0.123,0.579)rectangle++(0.008,0.008)
  (-0.267,0.64)rectangle++(0.008,0.008) (-0.242,0.825)rectangle++(0.008,0.008) (-0.234,0.887)rectangle++(0.008,0.008) (-0.427,0.741)rectangle++(0.008,0.008) (-0.463,0.605)rectangle++(0.008,0.008)
  (-0.273,0.76)rectangle++(0.008,0.008) (-0.399,0.669)rectangle++(0.008,0.008) (-0.104,0.544)rectangle++(0.008,0.008) (-0.245,0.631)rectangle++(0.008,0.008) (-0.341,0.62)rectangle++(0.008,0.008)
  (-0.202,0.898)rectangle++(0.008,0.008) (-0.154,0.855)rectangle++(0.008,0.008) (-0.285,0.913)rectangle++(0.008,0.008) (-0.298,0.91)rectangle++(0.008,0.008) (-0.465,0.729)rectangle++(0.008,0.008)
  (-0.009,0.86)rectangle++(0.008,0.008) (-0.315,0.809)rectangle++(0.008,0.008) (-0.259,0.881)rectangle++(0.008,0.008) (-0.099,0.928)rectangle++(0.008,0.008) (-0.38,0.808)rectangle++(0.008,0.008)
  (-0.466,0.653)rectangle++(0.008,0.008) (-0.45,0.551)rectangle++(0.008,0.008) (-0.406,0.511)rectangle++(0.008,0.008) (-0.24,0.718)rectangle++(0.008,0.008) (-0.233,0.851)rectangle++(0.008,0.008)
  (-0.231,0.896)rectangle++(0.008,0.008) (-0.155,0.863)rectangle++(0.008,0.008) (-0.383,0.758)rectangle++(0.008,0.008) (-0.449,0.629)rectangle++(0.008,0.008) (-0.384,0.828)rectangle++(0.008,0.008)
  (-0.282,0.888)rectangle++(0.008,0.008) (-0.247,0.908)rectangle++(0.008,0.008) (-0.441,0.746)rectangle++(0.008,0.008) (-0.471,0.603)rectangle++(0.008,0.008) (-0.22,0.896)rectangle++(0.008,0.008)
  (-0.424,0.75)rectangle++(0.008,0.008) (-0.155,0.864)rectangle++(0.008,0.008) (-0.175,0.932)rectangle++(0.008,0.008) (-0.212,0.923)rectangle++(0.008,0.008) (-0.43,0.766)rectangle++(0.008,0.008)
  (-0.119,0.938)rectangle++(0.008,0.008) (-0.297,0.904)rectangle++(0.008,0.008) (-0.462,0.726)rectangle++(0.008,0.008) (-0.248,0.898)rectangle++(0.008,0.008) (-0.276,0.762)rectangle++(0.008,0.008)
  (-0.329,0.831)rectangle++(0.008,0.008) (-0.275,0.765)rectangle++(0.008,0.008) (-0.403,0.671)rectangle++(0.008,0.008) (-0.288,0.835)rectangle++(0.008,0.008) (-0.399,0.809)rectangle++(0.008,0.008)
  (-0.475,0.647)rectangle++(0.008,0.008) (-0.348,0.863)rectangle++(0.008,0.008) (-0.471,0.689)rectangle++(0.008,0.008) (-0.114,0.544)rectangle++(0.008,0.008) (-0.102,0.934)rectangle++(0.008,0.008)
  (-0.384,0.81)rectangle++(0.008,0.008) (-0.413,0.789)rectangle++(0.008,0.008) (-0.474,0.633)rectangle++(0.008,0.008) (-0.404,0.742)rectangle++(0.008,0.008) (-0.453,0.614)rectangle++(0.008,0.008)
  (-0.43,0.538)rectangle++(0.008,0.008) (-0.297,0.79)rectangle++(0.008,0.008) (-0.394,0.799)rectangle++(0.008,0.008) (-0.209,0.923)rectangle++(0.008,0.008) (-0.332,0.837)rectangle++(0.008,0.008)
  (-0.275,0.765)rectangle++(0.008,0.008) (-0.044,0.873)rectangle++(0.008,0.008) (-0.335,0.802)rectangle++(0.008,0.008) (-0.135,0.938)rectangle++(0.008,0.008) (-0.4,0.8)rectangle++(0.008,0.008)
  (-0.472,0.642)rectangle++(0.008,0.008) (-0.449,0.544)rectangle++(0.008,0.008) (-0.304,0.792)rectangle++(0.008,0.008) (-0.255,0.876)rectangle++(0.008,0.008) (-0.281,0.829)rectangle++(0.008,0.008)
  (-0.428,0.698)rectangle++(0.008,0.008) (-0.387,0.835)rectangle++(0.008,0.008) (-0.112,0.924)rectangle++(0.008,0.008) (-0.385,0.802)rectangle++(0.008,0.008) (0.001,0.827)rectangle++(0.008,0.008)
  (-0.298,0.797)rectangle++(0.008,0.008) (-0.425,0.677)rectangle++(0.008,0.008) (-0.267,0.757)rectangle++(0.008,0.008) (-0.157,0.715)rectangle++(0.008,0.008) (-0.205,0.85)rectangle++(0.008,0.008)
  (-0.024,0.888)rectangle++(0.008,0.008) (-0.356,0.858)rectangle++(0.008,0.008) (-0.061,0.915)rectangle++(0.008,0.008) (-0.282,0.916)rectangle++(0.008,0.008) (-0.459,0.737)rectangle++(0.008,0.008)
  (-0.272,0.759)rectangle++(0.008,0.008) (-0.005,0.821)rectangle++(0.008,0.008) (-0.137,0.942)rectangle++(0.008,0.008) (-0.26,0.902)rectangle++(0.008,0.008) (-0.19,0.938)rectangle++(0.008,0.008)
  (-0.138,0.595)rectangle++(0.008,0.008) (-0.279,0.642)rectangle++(0.008,0.008) (-0.36,0.613)rectangle++(0.008,0.008) (-0.175,0.931)rectangle++(0.008,0.008) (-0.178,0.942)rectangle++(0.008,0.008)
  (-0.067,0.916)rectangle++(0.008,0.008) (-0.163,0.94)rectangle++(0.008,0.008) (-0.39,0.848)rectangle++(0.008,0.008) (-0.42,0.796)rectangle++(0.008,0.008) (-0.283,0.818)rectangle++(0.008,0.008)
  (-0.351,0.861)rectangle++(0.008,0.008) (-0.415,0.804)rectangle++(0.008,0.008) (-0.154,0.644)rectangle++(0.008,0.008) (-0.229,0.896)rectangle++(0.008,0.008) (-0.028,0.446)rectangle++(0.008,0.008)
  (-0.176,0.613)rectangle++(0.008,0.008) (-0.303,0.637)rectangle++(0.008,0.008) (-0.369,0.602)rectangle++(0.008,0.008) (-0.387,0.563)rectangle++(0.008,0.008) (-0.382,0.538)rectangle++(0.008,0.008)
  (-0.278,0.884)rectangle++(0.008,0.008) (-0.246,0.907)rectangle++(0.008,0.008) (-0.44,0.746)rectangle++(0.008,0.008) (-0.471,0.603)rectangle++(0.008,0.008) (-0.434,0.527)rectangle++(0.008,0.008)
  (-0.215,0.885)rectangle++(0.008,0.008) (-0.225,0.907)rectangle++(0.008,0.008) (-0.062,0.49)rectangle++(0.008,0.008) (-0.207,0.621)rectangle++(0.008,0.008) (-0.32,0.629)rectangle++(0.008,0.008)
  (-0.375,0.593)rectangle++(0.008,0.008) (-0.264,0.851)rectangle++(0.008,0.008) (-0.137,0.592)rectangle++(0.008,0.008) (-0.173,0.894)rectangle++(0.008,0.008) (-0.211,0.91)rectangle++(0.008,0.008)
  (-0.021,0.889)rectangle++(0.008,0.008) (-0.219,0.868)rectangle++(0.008,0.008) (-0.184,0.933)rectangle++(0.008,0.008) (-0.215,0.872)rectangle++(0.008,0.008) (-0.413,0.741)rectangle++(0.008,0.008)
  (-0.148,0.928)rectangle++(0.008,0.008) (-0.154,0.606)rectangle++(0.008,0.008) (-0.359,0.845)rectangle++(0.008,0.008) (-0.414,0.8)rectangle++(0.008,0.008) (-0.478,0.637)rectangle++(0.008,0.008)
  (-0.314,0.824)rectangle++(0.008,0.008) (-0.258,0.886)rectangle++(0.008,0.008) (-0.317,0.897)rectangle++(0.008,0.008) (-0.413,0.816)rectangle++(0.008,0.008) (-0.42,0.786)rectangle++(0.008,0.008)
  (-0.476,0.629)rectangle++(0.008,0.008) (-0.446,0.537)rectangle++(0.008,0.008) (-0.303,0.79)rectangle++(0.008,0.008) (-0.424,0.672)rectangle++(0.008,0.008) (-0.437,0.575)rectangle++(0.008,0.008)
  (-0.409,0.526)rectangle++(0.008,0.008) (-0.375,0.731)rectangle++(0.008,0.008) (-0.123,0.907)rectangle++(0.008,0.008) (-0.171,0.938)rectangle++(0.008,0.008) (-0.045,0.911)rectangle++(0.008,0.008)
  (-0.349,0.819)rectangle++(0.008,0.008) (-0.266,0.756)rectangle++(0.008,0.008) (-0.334,0.861)rectangle++(0.008,0.008) (-0.018,0.874)rectangle++(0.008,0.008) (-0.261,0.907)rectangle++(0.008,0.008)
  (-0.386,0.863)rectangle++(0.008,0.008) (-0.019,0.859)rectangle++(0.008,0.008) (-0.159,0.898)rectangle++(0.008,0.008) (-0.397,0.773)rectangle++(0.008,0.008) (-0.121,0.916)rectangle++(0.008,0.008)
  (-0.342,0.894)rectangle++(0.008,0.008) (-0.144,0.951)rectangle++(0.008,0.008) (0.005,0.848)rectangle++(0.008,0.008) (-0.304,0.808)rectangle++(0.008,0.008) (-0.066,0.824)rectangle++(0.008,0.008)
  (-0.147,0.598)rectangle++(0.008,0.008) (-0.041,0.846)rectangle++(0.008,0.008) (-0.131,0.943)rectangle++(0.008,0.008) (-0.4,0.803)rectangle++(0.008,0.008) (-0.21,0.924)rectangle++(0.008,0.008)
  (-0.007,0.862)rectangle++(0.008,0.008) (-0.059,0.486)rectangle++(0.008,0.008) (-0.078,0.509)rectangle++(0.008,0.008) (-0.221,0.624)rectangle++(0.008,0.008) (-0.227,0.819)rectangle++(0.008,0.008)
  (-0.133,0.586)rectangle++(0.008,0.008) (-0.274,0.64)rectangle++(0.008,0.008) (-0.357,0.614)rectangle++(0.008,0.008) (-0.134,0.887)rectangle++(0.008,0.008) (-0.377,0.842)rectangle++(0.008,0.008)
  (-0.021,0.878)rectangle++(0.008,0.008) (-0.327,0.813)rectangle++(0.008,0.008) (-0.131,0.569)rectangle++(0.008,0.008) (-0.267,0.633)rectangle++(0.008,0.008) (-0.35,0.877)rectangle++(0.008,0.008);
\fill[piece2]
  (-0.477,0.695)rectangle++(0.008,0.008) (0.003,0.831)rectangle++(0.008,0.008) (-0.36,0.873)rectangle++(0.008,0.008) (-0.265,0.753)rectangle++(0.008,0.008) (-0.242,0.863)rectangle++(0.008,0.008)
  (-0.422,0.727)rectangle++(0.008,0.008) (-0.018,0.779)rectangle++(0.008,0.008) (-0.143,0.594)rectangle++(0.008,0.008) (-0.152,0.713)rectangle++(0.008,0.008) (-0.154,0.912)rectangle++(0.008,0.008)
  (-0.067,0.908)rectangle++(0.008,0.008) (-0.367,0.856)rectangle++(0.008,0.008) (-0.477,0.679)rectangle++(0.008,0.008) (-0.15,0.937)rectangle++(0.008,0.008) (0.004,0.848)rectangle++(0.008,0.008)
  (-0.272,0.759)rectangle++(0.008,0.008) (-0.244,0.865)rectangle++(0.008,0.008) (-0.395,0.821)rectangle++(0.008,0.008) (-0.285,0.885)rectangle++(0.008,0.008) (-0.376,0.884)rectangle++(0.008,0.008)
  (-0.138,0.596)rectangle++(0.008,0.008) (-0.279,0.642)rectangle++(0.008,0.008) (-0.123,0.792)rectangle++(0.008,0.008) (-0.342,0.737)rectangle++(0.008,0.008) (-0.268,0.857)rectangle++(0.008,0.008)
  (-0.324,0.893)rectangle++(0.008,0.008) (-0.47,0.712)rectangle++(0.008,0.008) (-0.473,0.577)rectangle++(0.008,0.008) (-0.312,0.803)rectangle++(0.008,0.008) (-0.433,0.675)rectangle++(0.008,0.008)
  (-0.122,0.795)rectangle++(0.008,0.008) (-0.279,0.808)rectangle++(0.008,0.008) (-0.193,0.925)rectangle++(0.008,0.008) (-0.007,0.86)rectangle++(0.008,0.008) (-0.151,0.699)rectangle++(0.008,0.008)
  (-0.113,0.93)rectangle++(0.008,0.008) (-0.151,0.733)rectangle++(0.008,0.008) (-0.334,0.7)rectangle++(0.008,0.008) (-0.407,0.62)rectangle++(0.008,0.008) (-0.158,0.696)rectangle++(0.008,0.008)
  (-0.282,0.824)rectangle++(0.008,0.008) (-0.32,0.888)rectangle++(0.008,0.008) (-0.096,0.924)rectangle++(0.008,0.008) (-0.167,0.941)rectangle++(0.008,0.008) (-0.066,0.914)rectangle++(0.008,0.008)
  (-0.36,0.814)rectangle++(0.008,0.008) (-0.1,0.553)rectangle++(0.008,0.008) (-0.32,0.783)rectangle++(0.008,0.008) (-0.319,0.892)rectangle++(0.008,0.008) (-0.256,0.904)rectangle++(0.008,0.008)
  (-0.403,0.827)rectangle++(0.008,0.008) (-0.134,0.571)rectangle++(0.008,0.008) (-0.269,0.632)rectangle++(0.008,0.008) (-0.268,0.848)rectangle++(0.008,0.008) (-0.429,0.711)rectangle++(0.008,0.008)
  (-0.037,0.88)rectangle++(0.008,0.008) (-0.357,0.855)rectangle++(0.008,0.008) (-0.472,0.683)rectangle++(0.008,0.008) (-0.463,0.563)rectangle++(0.008,0.008) (-0.416,0.511)rectangle++(0.008,0.008)
  (-0.376,0.504)rectangle++(0.008,0.008) (-0.356,0.515)rectangle++(0.008,0.008) (-0.128,0.581)rectangle++(0.008,0.008) (-0.269,0.639)rectangle++(0.008,0.008) (-0.243,0.824)rectangle++(0.008,0.008)
  (-0.188,0.927)rectangle++(0.008,0.008) (-0.42,0.776)rectangle++(0.008,0.008) (-0.473,0.624)rectangle++(0.008,0.008) (-0.442,0.536)rectangle++(0.008,0.008) (-0.397,0.506)rectangle++(0.008,0.008)
  (-0.357,0.888)rectangle++(0.008,0.008) (-0.484,0.698)rectangle++(0.008,0.008) (-0.474,0.565)rectangle++(0.008,0.008) (-0.158,0.893)rectangle++(0.008,0.008) (-0.382,0.839)rectangle++(0.008,0.008)
  (-0.417,0.795)rectangle++(0.008,0.008) (-0.153,0.683)rectangle++(0.008,0.008) (-0.317,0.677)rectangle++(0.008,0.008) (-0.39,0.615)rectangle++(0.008,0.008) (-0.402,0.561)rectangle++(0.008,0.008)
  (-0.379,0.739)rectangle++(0.008,0.008) (-0.088,0.9)rectangle++(0.008,0.008) (-0.15,0.709)rectangle++(0.008,0.008) (-0.14,0.588)rectangle++(0.008,0.008) (-0.151,0.944)rectangle++(0.008,0.008)
  (-0.41,0.797)rectangle++(0.008,0.008) (-0.211,0.923)rectangle++(0.008,0.008) (-0.43,0.766)rectangle++(0.008,0.008) (-0.155,0.865)rectangle++(0.008,0.008) (-0.205,0.9)rectangle++(0.008,0.008)
  (-0.049,0.91)rectangle++(0.008,0.008) (-0.169,0.915)rectangle++(0.008,0.008) (-0.291,0.91)rectangle++(0.008,0.008) (-0.322,0.9)rectangle++(0.008,0.008) (-0.057,0.916)rectangle++(0.008,0.008)
  (-0.357,0.818)rectangle++(0.008,0.008) (-0.273,0.884)rectangle++(0.008,0.008) (-0.095,0.919)rectangle++(0.008,0.008) (-0.154,0.747)rectangle++(0.008,0.008) (-0.334,0.855)rectangle++(0.008,0.008)
  (-0.242,0.862)rectangle++(0.008,0.008) (-0.241,0.862)rectangle++(0.008,0.008) (-0.323,0.893)rectangle++(0.008,0.008) (-0.47,0.712)rectangle++(0.008,0.008) (-0.391,0.834)rectangle++(0.008,0.008)
  (-0.284,0.89)rectangle++(0.008,0.008) (-0.248,0.909)rectangle++(0.008,0.008) (-0.052,0.913)rectangle++(0.008,0.008) (-0.128,0.584)rectangle++(0.008,0.008) (-0.338,0.89)rectangle++(0.008,0.008)
  (-0.354,0.892)rectangle++(0.008,0.008) (-0.484,0.7)rectangle++(0.008,0.008) (-0.363,0.881)rectangle++(0.008,0.008) (-0.133,0.944)rectangle++(0.008,0.008) (-0.003,0.854)rectangle++(0.008,0.008)
  (-0.006,0.802)rectangle++(0.008,0.008) (0.002,0.821)rectangle++(0.008,0.008) (-0.079,0.524)rectangle++(0.008,0.008) (-0.227,0.63)rectangle++(0.008,0.008) (-0.356,0.779)rectangle++(0.008,0.008)
  (-0.089,0.527)rectangle++(0.008,0.008) (-0.182,0.787)rectangle++(0.008,0.008) (-0.368,0.714)rectangle++(0.008,0.008) (-0.069,0.906)rectangle++(0.008,0.008) (-0.359,0.809)rectangle++(0.008,0.008)
  (-0.273,0.881)rectangle++(0.008,0.008) (-0.403,0.82)rectangle++(0.008,0.008) (-0.48,0.65)rectangle++(0.008,0.008) (-0.455,0.545)rectangle++(0.008,0.008) (-0.406,0.506)rectangle++(0.008,0.008)
  (-0.118,0.57)rectangle++(0.008,0.008) (-0.261,0.638)rectangle++(0.008,0.008) (-0.169,0.933)rectangle++(0.008,0.008) (-0.21,0.923)rectangle++(0.008,0.008) (-0.429,0.766)rectangle++(0.008,0.008)
  (-0.15,0.723)rectangle++(0.008,0.008) (-0.29,0.899)rectangle++(0.008,0.008) (-0.457,0.727)rectangle++(0.008,0.008) (-0.335,0.893)rectangle++(0.008,0.008) (-0.201,0.937)rectangle++(0.008,0.008)
  (-0.279,0.765)rectangle++(0.008,0.008) (-0.386,0.797)rectangle++(0.008,0.008) (-0.177,0.942)rectangle++(0.008,0.008) (-0.421,0.787)rectangle++(0.008,0.008) (-0.146,0.593)rectangle++(0.008,0.008)
  (-0.174,0.894)rectangle++(0.008,0.008) (-0.402,0.766)rectangle++(0.008,0.008) (-0.141,0.939)rectangle++(0.008,0.008) (-0.091,0.54)rectangle++(0.008,0.008) (-0.238,0.633)rectangle++(0.008,0.008)
  (-0.113,0.875)rectangle++(0.008,0.008) (-0.368,0.779)rectangle++(0.008,0.008) (-0.45,0.644)rectangle++(0.008,0.008) (-0.439,0.553)rectangle++(0.008,0.008) (-0.3,0.795)rectangle++(0.008,0.008)
  (-0.254,0.877)rectangle++(0.008,0.008) (-0.134,0.592)rectangle++(0.008,0.008) (-0.151,0.707)rectangle++(0.008,0.008) (-0.285,0.909)rectangle++(0.008,0.008) (-0.458,0.733)rectangle++(0.008,0.008)
  (-0.416,0.764)rectangle++(0.008,0.008) (-0.088,0.537)rectangle++(0.008,0.008) (-0.102,0.843)rectangle++(0.008,0.008) (-0.351,0.768)rectangle++(0.008,0.008) (-0.002,0.826)rectangle++(0.008,0.008)
  (-0.154,0.725)rectangle++(0.008,0.008) (-0.366,0.855)rectangle++(0.008,0.008) (-0.107,0.925)rectangle++(0.008,0.008) (-0.284,0.914)rectangle++(0.008,0.008) (-0.087,0.924)rectangle++(0.008,0.008)
  (-0.277,0.765)rectangle++(0.008,0.008) (-0.062,0.49)rectangle++(0.008,0.008) (-0.173,0.774)rectangle++(0.008,0.008) (-0.281,0.77)rectangle++(0.008,0.008) (-0.026,0.865)rectangle++(0.008,0.008)
  (-0.135,0.583)rectangle++(0.008,0.008) (-0.273,0.637)rectangle++(0.008,0.008) (-0.106,0.916)rectangle++(0.008,0.008) (-0.376,0.851)rectangle++(0.008,0.008) (-0.479,0.674)rectangle++(0.008,0.008)
  (-0.146,0.911)rectangle++(0.008,0.008) (-0.174,0.939)rectangle++(0.008,0.008) (-0.418,0.786)rectangle++(0.008,0.008) (-0.046,0.47)rectangle++(0.008,0.008) (-0.043,0.894)rectangle++(0.008,0.008)
  (-0.154,0.735)rectangle++(0.008,0.008) (-0.129,0.584)rectangle++(0.008,0.008) (-0.172,0.893)rectangle++(0.008,0.008) (-0.401,0.766)rectangle++(0.008,0.008) (-0.287,0.867)rectangle++(0.008,0.008)
  (-0.115,0.569)rectangle++(0.008,0.008) (-0.354,0.864)rectangle++(0.008,0.008) (-0.474,0.688)rectangle++(0.008,0.008) (-0.312,0.841)rectangle++(0.008,0.008) (-0.019,0.871)rectangle++(0.008,0.008)
  (-0.323,0.81)rectangle++(0.008,0.008) (-0.023,0.871)rectangle++(0.008,0.008) (-0.325,0.809)rectangle++(0.008,0.008) (-0.441,0.673)rectangle++(0.008,0.008) (-0.301,0.836)rectangle++(0.008,0.008)
  (-0.44,0.694)rectangle++(0.008,0.008) (-0.12,0.899)rectangle++(0.008,0.008) (-0.376,0.846)rectangle++(0.008,0.008) (-0.478,0.672)rectangle++(0.008,0.008) (-0.273,0.763)rectangle++(0.008,0.008)
  (-0.192,0.918)rectangle++(0.008,0.008) (-0.419,0.77)rectangle++(0.008,0.008) (-0.125,0.918)rectangle++(0.008,0.008) (-0.376,0.872)rectangle++(0.008,0.008) (-0.279,0.903)rectangle++(0.008,0.008)
  (-0.247,0.913)rectangle++(0.008,0.008) (-0.442,0.748)rectangle++(0.008,0.008) (-0.361,0.884)rectangle++(0.008,0.008) (-0.274,0.907)rectangle++(0.008,0.008) (-0.453,0.736)rectangle++(0.008,0.008)
  (-0.139,0.89)rectangle++(0.008,0.008) (-0.067,0.903)rectangle++(0.008,0.008) (-0.357,0.808)rectangle++(0.008,0.008) (-0.147,0.595)rectangle++(0.008,0.008) (-0.202,0.809)rectangle++(0.008,0.008)
  (-0.378,0.817)rectangle++(0.008,0.008) (-0.207,0.926)rectangle++(0.008,0.008) (-0.012,0.782)rectangle++(0.008,0.008) (-0.002,0.819)rectangle++(0.008,0.008) (-0.281,0.82)rectangle++(0.008,0.008)
  (-0.425,0.693)rectangle++(0.008,0.008) (-0.445,0.585)rectangle++(0.008,0.008) (-0.416,0.527)rectangle++(0.008,0.008) (-0.153,0.916)rectangle++(0.008,0.008) (-0.385,0.844)rectangle++(0.008,0.008)
  (-0.419,0.796)rectangle++(0.008,0.008) (-0.479,0.634)rectangle++(0.008,0.008) (-0.15,0.904)rectangle++(0.008,0.008) (-0.096,0.925)rectangle++(0.008,0.008) (-0.185,0.921)rectangle++(0.008,0.008)
  (-0.417,0.774)rectangle++(0.008,0.008) (-0.212,0.92)rectangle++(0.008,0.008) (-0.224,0.919)rectangle++(0.008,0.008) (-0.145,0.591)rectangle++(0.008,0.008) (-0.272,0.758)rectangle++(0.008,0.008)
  (-0.399,0.669)rectangle++(0.008,0.008) (-0.425,0.582)rectangle++(0.008,0.008) (-0.295,0.805)rectangle++(0.008,0.008) (-0.383,0.852)rectangle++(0.008,0.008) (-0.483,0.672)rectangle++(0.008,0.008)
  (-0.154,0.735)rectangle++(0.008,0.008) (-0.336,0.7)rectangle++(0.008,0.008) (-0.379,0.841)rectangle++(0.008,0.008) (-0.329,0.832)rectangle++(0.008,0.008) (-0.406,0.802)rectangle++(0.008,0.008)
  (-0.289,0.879)rectangle++(0.008,0.008) (-0.405,0.817)rectangle++(0.008,0.008) (-0.289,0.884)rectangle++(0.008,0.008) (-0.156,0.609)rectangle++(0.008,0.008) (-0.145,0.931)rectangle++(0.008,0.008)
  (-0.386,0.848)rectangle++(0.008,0.008) (-0.235,0.898)rectangle++(0.008,0.008) (-0.232,0.911)rectangle++(0.008,0.008) (-0.109,0.809)rectangle++(0.008,0.008) (-0.342,0.75)rectangle++(0.008,0.008)
  (-0.202,0.916)rectangle++(0.008,0.008) (-0.423,0.766)rectangle++(0.008,0.008) (-0.47,0.618)rectangle++(0.008,0.008) (-0.11,0.805)rectangle++(0.008,0.008) (-0.341,0.748)rectangle++(0.008,0.008)
  (-0.427,0.639)rectangle++(0.008,0.008) (-0.341,0.862)rectangle++(0.008,0.008) (-0.467,0.691)rectangle++(0.008,0.008) (-0.464,0.568)rectangle++(0.008,0.008) (-0.419,0.513)rectangle++(0.008,0.008)
  (-0.34,0.844)rectangle++(0.008,0.008) (-0.006,0.792)rectangle++(0.008,0.008) (-0.042,0.893)rectangle++(0.008,0.008) (-0.167,0.91)rectangle++(0.008,0.008) (-0.405,0.775)rectangle++(0.008,0.008)
  (-0.465,0.629)rectangle++(0.008,0.008) (-0.134,0.588)rectangle++(0.008,0.008) (-0.14,0.591)rectangle++(0.008,0.008) (-0.279,0.64)rectangle++(0.008,0.008) (-0.005,0.825)rectangle++(0.008,0.008)
  (-0.106,0.93)rectangle++(0.008,0.008) (-0.188,0.922)rectangle++(0.008,0.008) (-0.419,0.774)rectangle++(0.008,0.008) (-0.254,0.737)rectangle++(0.008,0.008) (-0.383,0.666)rectangle++(0.008,0.008)
  (-0.335,0.865)rectangle++(0.008,0.008) (-0.1,0.923)rectangle++(0.008,0.008) (-0.295,0.902)rectangle++(0.008,0.008) (-0.012,0.871)rectangle++(0.008,0.008) (-0.157,0.902)rectangle++(0.008,0.008)
  (-0.262,0.751)rectangle++(0.008,0.008) (-0.112,0.565)rectangle++(0.008,0.008) (-0.324,0.783)rectangle++(0.008,0.008) (-0.432,0.661)rectangle++(0.008,0.008) (-0.139,0.589)rectangle++(0.008,0.008)
  (-0.366,0.87)rectangle++(0.008,0.008) (-0.482,0.686)rectangle++(0.008,0.008) (-0.152,0.859)rectangle++(0.008,0.008) (-0.149,0.693)rectangle++(0.008,0.008) (-0.319,0.683)rectangle++(0.008,0.008)
  (-0.393,0.617)rectangle++(0.008,0.008) (-0.139,0.892)rectangle++(0.008,0.008) (-0.3,0.897)rectangle++(0.008,0.008) (-0.41,0.819)rectangle++(0.008,0.008) (-0.484,0.647)rectangle++(0.008,0.008)
  (-0.456,0.542)rectangle++(0.008,0.008) (-0.405,0.504)rectangle++(0.008,0.008) (-0.155,0.876)rectangle++(0.008,0.008) (-0.047,0.906)rectangle++(0.008,0.008) (-0.051,0.894)rectangle++(0.008,0.008)
  (-0.239,0.891)rectangle++(0.008,0.008) (-0.152,0.742)rectangle++(0.008,0.008) (-0.204,0.859)rectangle++(0.008,0.008) (-0.081,0.906)rectangle++(0.008,0.008) (-0.057,0.898)rectangle++(0.008,0.008)
  (-0.172,0.912)rectangle++(0.008,0.008) (-0.408,0.775)rectangle++(0.008,0.008) (-0.466,0.628)rectangle++(0.008,0.008) (-0.441,0.54)rectangle++(0.008,0.008) (-0.398,0.508)rectangle++(0.008,0.008)
  (-0.367,0.51)rectangle++(0.008,0.008) (-0.148,0.872)rectangle++(0.008,0.008) (-0.335,0.839)rectangle++(0.008,0.008) (-0.456,0.683)rectangle++(0.008,0.008) (-0.374,0.885)rectangle++(0.008,0.008)
  (-0.366,0.89)rectangle++(0.008,0.008) (-0.489,0.695)rectangle++(0.008,0.008) (-0.475,0.562)rectangle++(0.008,0.008) (-0.159,0.688)rectangle++(0.008,0.008) (-0.322,0.677)rectangle++(0.008,0.008)
  (-0.026,0.898)rectangle++(0.008,0.008) (-0.281,0.91)rectangle++(0.008,0.008) (-0.247,0.915)rectangle++(0.008,0.008) (-0.153,0.865)rectangle++(0.008,0.008) (-0.377,0.835)rectangle++(0.008,0.008)
  (-0.416,0.795)rectangle++(0.008,0.008) (-0.101,0.933)rectangle++(0.008,0.008) (-0.383,0.81)rectangle++(0.008,0.008) (-0.468,0.653)rectangle++(0.008,0.008) (-0.45,0.55)rectangle++(0.008,0.008)
  (-0.361,0.872)rectangle++(0.008,0.008) (-0.26,0.905)rectangle++(0.008,0.008) (-0.446,0.74)rectangle++(0.008,0.008) (-0.152,0.736)rectangle++(0.008,0.008) (-0.1,0.928)rectangle++(0.008,0.008)
  (-0.24,0.895)rectangle++(0.008,0.008) (-0.433,0.743)rectangle++(0.008,0.008) (-0.298,0.859)rectangle++(0.008,0.008) (-0.248,0.86)rectangle++(0.008,0.008) (-0.395,0.819)rectangle++(0.008,0.008)
  (-0.311,0.879)rectangle++(0.008,0.008) (-0.257,0.905)rectangle++(0.008,0.008) (-0.155,0.747)rectangle++(0.008,0.008) (-0.124,0.575)rectangle++(0.008,0.008) (-0.194,0.803)rectangle++(0.008,0.008)
  (-0.181,0.924)rectangle++(0.008,0.008) (-0.156,0.739)rectangle++(0.008,0.008) (-0.359,0.81)rectangle++(0.008,0.008) (-0.331,0.886)rectangle++(0.008,0.008) (-0.327,0.897)rectangle++(0.008,0.008)
  (-0.135,0.593)rectangle++(0.008,0.008) (-0.3,0.855)rectangle++(0.008,0.008) (-0.446,0.703)rectangle++(0.008,0.008) (-0.458,0.581)rectangle++(0.008,0.008) (-0.421,0.521)rectangle++(0.008,0.008)
  (-0.382,0.507)rectangle++(0.008,0.008) (-0.335,0.843)rectangle++(0.008,0.008) (-0.136,0.575)rectangle++(0.008,0.008) (-0.108,0.909)rectangle++(0.008,0.008) (-0.377,0.796)rectangle++(0.008,0.008)
  (-0.46,0.648)rectangle++(0.008,0.008) (-0.153,0.691)rectangle++(0.008,0.008) (-0.204,0.842)rectangle++(0.008,0.008) (-0.362,0.889)rectangle++(0.008,0.008) (-0.259,0.906)rectangle++(0.008,0.008)
  (-0.446,0.741)rectangle++(0.008,0.008) (-0.472,0.599)rectangle++(0.008,0.008) (-0.433,0.524)rectangle++(0.008,0.008) (-0.389,0.504)rectangle++(0.008,0.008) (-0.251,0.891)rectangle++(0.008,0.008)
  (-0.125,0.939)rectangle++(0.008,0.008) (-0.396,0.804)rectangle++(0.008,0.008) (-0.228,0.892)rectangle++(0.008,0.008) (-0.426,0.745)rectangle++(0.008,0.008) (-0.153,0.863)rectangle++(0.008,0.008)
  (-0.152,0.729)rectangle++(0.008,0.008) (-0.333,0.698)rectangle++(0.008,0.008) (-0.405,0.619)rectangle++(0.008,0.008) (-0.103,0.531)rectangle++(0.008,0.008) (-0.24,0.625)rectangle++(0.008,0.008)
  (-0.154,0.624)rectangle++(0.008,0.008) (-0.297,0.649)rectangle++(0.008,0.008) (-0.01,0.825)rectangle++(0.008,0.008) (-0.119,0.575)rectangle++(0.008,0.008) (-0.092,0.526)rectangle++(0.008,0.008)
  (-0.369,0.881)rectangle++(0.008,0.008) (-0.421,0.805)rectangle++(0.008,0.008) (-0.278,0.765)rectangle++(0.008,0.008) (-0.104,0.936)rectangle++(0.008,0.008) (-0.063,0.492)rectangle++(0.008,0.008)
  (-0.151,0.721)rectangle++(0.008,0.008) (-0.09,0.878)rectangle++(0.008,0.008) (-0.293,0.896)rectangle++(0.008,0.008) (-0.409,0.82)rectangle++(0.008,0.008) (-0.483,0.648)rectangle++(0.008,0.008)
  (-0.154,0.866)rectangle++(0.008,0.008) (-0.302,0.894)rectangle++(0.008,0.008) (-0.152,0.602)rectangle++(0.008,0.008) (-0.175,0.895)rectangle++(0.008,0.008) (-0.403,0.766)rectangle++(0.008,0.008)
  (-0.461,0.625)rectangle++(0.008,0.008) (-0.219,0.899)rectangle++(0.008,0.008) (-0.142,0.587)rectangle++(0.008,0.008) (-0.333,0.783)rectangle++(0.008,0.008) (-0.4,0.792)rectangle++(0.008,0.008)
  (-0.062,0.491)rectangle++(0.008,0.008) (-0.145,0.596)rectangle++(0.008,0.008) (-0.301,0.856)rectangle++(0.008,0.008) (-0.155,0.916)rectangle++(0.008,0.008) (-0.401,0.783)rectangle++(0.008,0.008)
  (-0.413,0.782)rectangle++(0.008,0.008) (-0.414,0.78)rectangle++(0.008,0.008) (-0.081,0.525)rectangle++(0.008,0.008) (-0.228,0.63)rectangle++(0.008,0.008) (-0.128,0.582)rectangle++(0.008,0.008)
  (-0.27,0.64)rectangle++(0.008,0.008) (-0.355,0.615)rectangle++(0.008,0.008) (-0.386,0.574)rectangle++(0.008,0.008) (-0.282,0.802)rectangle++(0.008,0.008) (-0.393,0.803)rectangle++(0.008,0.008)
  (-0.119,0.575)rectangle++(0.008,0.008) (-0.132,0.945)rectangle++(0.008,0.008) (-0.042,0.909)rectangle++(0.008,0.008) (-0.155,0.953)rectangle++(0.008,0.008) (-0.415,0.799)rectangle++(0.008,0.008)
  (-0.479,0.637)rectangle++(0.008,0.008) (-0.405,0.742)rectangle++(0.008,0.008) (0.003,0.849)rectangle++(0.008,0.008) (-0.305,0.808)rectangle++(0.008,0.008) (-0.432,0.679)rectangle++(0.008,0.008)
  (-0.444,0.576)rectangle++(0.008,0.008) (-0.39,0.736)rectangle++(0.008,0.008) (-0.284,0.857)rectangle++(0.008,0.008) (-0.385,0.857)rectangle++(0.008,0.008) (-0.262,0.752)rectangle++(0.008,0.008)
  (-0.121,0.558)rectangle++(0.008,0.008) (-0.298,0.85)rectangle++(0.008,0.008) (-0.253,0.895)rectangle++(0.008,0.008) (-0.09,0.54)rectangle++(0.008,0.008) (-0.238,0.634)rectangle++(0.008,0.008)
  (-0.122,0.57)rectangle++(0.008,0.008) (-0.092,0.525)rectangle++(0.008,0.008) (-0.233,0.626)rectangle++(0.008,0.008) (-0.041,0.856)rectangle++(0.008,0.008) (-0.041,0.899)rectangle++(0.008,0.008)
  (-0.343,0.815)rectangle++(0.008,0.008) (-0.329,0.887)rectangle++(0.008,0.008) (-0.281,0.88)rectangle++(0.008,0.008) (-0.099,0.929)rectangle++(0.008,0.008) (-0.155,0.74)rectangle++(0.008,0.008)
  (-0.339,0.702)rectangle++(0.008,0.008) (-0.409,0.619)rectangle++(0.008,0.008) (-0.389,0.749)rectangle++(0.008,0.008) (-0.449,0.623)rectangle++(0.008,0.008) (-0.431,0.543)rectangle++(0.008,0.008)
  (-0.214,0.887)rectangle++(0.008,0.008) (-0.136,0.592)rectangle++(0.008,0.008) (-0.095,0.529)rectangle++(0.008,0.008) (-0.316,0.779)rectangle++(0.008,0.008) (-0.327,0.892)rectangle++(0.008,0.008)
  (-0.2,0.936)rectangle++(0.008,0.008) (-0.429,0.776)rectangle++(0.008,0.008) (-0.297,0.87)rectangle++(0.008,0.008) (-0.195,0.933)rectangle++(0.008,0.008) (-0.114,0.569)rectangle++(0.008,0.008)
  (-0.191,0.801)rectangle++(0.008,0.008) (-0.217,0.879)rectangle++(0.008,0.008) (-0.335,0.861)rectangle++(0.008,0.008) (-0.412,0.807)rectangle++(0.008,0.008) (-0.028,0.897)rectangle++(0.008,0.008)
  (-0.337,0.819)rectangle++(0.008,0.008) (-0.134,0.939)rectangle++(0.008,0.008) (-0.011,0.785)rectangle++(0.008,0.008) (-0.338,0.841)rectangle++(0.008,0.008) (-0.118,0.914)rectangle++(0.008,0.008)
  (-0.384,0.795)rectangle++(0.008,0.008) (-0.463,0.645)rectangle++(0.008,0.008) (-0.404,0.746)rectangle++(0.008,0.008) (-0.455,0.616)rectangle++(0.008,0.008) (-0.134,0.895)rectangle++(0.008,0.008)
  (-0.332,0.836)rectangle++(0.008,0.008) (-0.407,0.802)rectangle++(0.008,0.008) (-0.476,0.64)rectangle++(0.008,0.008) (-0.107,0.535)rectangle++(0.008,0.008) (-0.047,0.84)rectangle++(0.008,0.008)
  (-0.325,0.786)rectangle++(0.008,0.008) (-0.001,0.823)rectangle++(0.008,0.008) (-0.136,0.942)rectangle++(0.008,0.008) (-0.138,0.592)rectangle++(0.008,0.008) (-0.278,0.641)rectangle++(0.008,0.008)
  (-0.275,0.907)rectangle++(0.008,0.008) (-0.453,0.736)rectangle++(0.008,0.008) (-0.035,0.883)rectangle++(0.008,0.008) (-0.335,0.811)rectangle++(0.008,0.008) (-0.265,0.882)rectangle++(0.008,0.008)
  (-0.281,0.793)rectangle++(0.008,0.008) (-0.089,0.528)rectangle++(0.008,0.008) (-0.103,0.841)rectangle++(0.008,0.008) (-0.068,0.891)rectangle++(0.008,0.008) (-0.24,0.891)rectangle++(0.008,0.008)
  (-0.432,0.741)rectangle++(0.008,0.008) (-0.345,0.893)rectangle++(0.008,0.008) (-0.48,0.704)rectangle++(0.008,0.008) (-0.416,0.755)rectangle++(0.008,0.008) (-0.411,0.774)rectangle++(0.008,0.008)
  (-0.283,0.817)rectangle++(0.008,0.008) (-0.136,0.59)rectangle++(0.008,0.008) (-0.276,0.64)rectangle++(0.008,0.008) (-0.358,0.613)rectangle++(0.008,0.008) (-0.273,0.816)rectangle++(0.008,0.008)
  (-0.42,0.694)rectangle++(0.008,0.008) (-0.294,0.843)rectangle++(0.008,0.008) (-0.401,0.809)rectangle++(0.008,0.008) (-0.287,0.882)rectangle++(0.008,0.008) (-0.405,0.818)rectangle++(0.008,0.008)
  (-0.117,0.924)rectangle++(0.008,0.008) (-0.387,0.8)rectangle++(0.008,0.008) (-0.283,0.878)rectangle++(0.008,0.008) (-0.321,0.896)rectangle++(0.008,0.008) (-0.199,0.937)rectangle++(0.008,0.008)
  (-0.321,0.822)rectangle++(0.008,0.008) (-0.261,0.886)rectangle++(0.008,0.008) (-0.318,0.897)rectangle++(0.008,0.008) (-0.259,0.911)rectangle++(0.008,0.008) (-0.447,0.743)rectangle++(0.008,0.008)
  (-0.473,0.599)rectangle++(0.008,0.008) (-0.434,0.524)rectangle++(0.008,0.008) (-0.083,0.531)rectangle++(0.008,0.008) (-0.231,0.632)rectangle++(0.008,0.008) (-0.327,0.897)rectangle++(0.008,0.008)
  (-0.473,0.712)rectangle++(0.008,0.008) (-0.312,0.849)rectangle++(0.008,0.008) (-0.449,0.696)rectangle++(0.008,0.008) (-0.138,0.589)rectangle++(0.008,0.008) (-0.333,0.784)rectangle++(0.008,0.008)
  (-0.436,0.659)rectangle++(0.008,0.008) (-0.299,0.831)rectangle++(0.008,0.008) (-0.4,0.806)rectangle++(0.008,0.008) (-0.142,0.596)rectangle++(0.008,0.008) (-0.107,0.91)rectangle++(0.008,0.008)
  (-0.276,0.763)rectangle++(0.008,0.008) (-0.386,0.797)rectangle++(0.008,0.008) (-0.464,0.646)rectangle++(0.008,0.008) (-0.346,0.863)rectangle++(0.008,0.008) (-0.139,0.579)rectangle++(0.008,0.008)
  (-0.28,0.813)rectangle++(0.008,0.008) (-0.271,0.755)rectangle++(0.008,0.008) (-0.157,0.715)rectangle++(0.008,0.008) (-0.28,0.802)rectangle++(0.008,0.008) (-0.152,0.733)rectangle++(0.008,0.008)
  (-0.302,0.875)rectangle++(0.008,0.008) (-0.277,0.768)rectangle++(0.008,0.008) (-0.374,0.882)rectangle++(0.008,0.008) (-0.024,0.896)rectangle++(0.008,0.008) (-0.335,0.82)rectangle++(0.008,0.008)
  (-0.449,0.675)rectangle++(0.008,0.008) (-0.344,0.867)rectangle++(0.008,0.008) (-0.47,0.692)rectangle++(0.008,0.008) (-0.466,0.568)rectangle++(0.008,0.008) (-0.096,0.521)rectangle++(0.008,0.008)
  (-0.233,0.623)rectangle++(0.008,0.008) (-0.348,0.876)rectangle++(0.008,0.008) (-0.416,0.808)rectangle++(0.008,0.008) (-0.281,0.88)rectangle++(0.008,0.008) (-0.446,0.721)rectangle++(0.008,0.008)
  (-0.465,0.589)rectangle++(0.008,0.008) (-0.309,0.808)rectangle++(0.008,0.008) (-0.103,0.911)rectangle++(0.008,0.008) (-0.035,0.456)rectangle++(0.008,0.008) (-0.182,0.615)rectangle++(0.008,0.008)
  (0.008,0.834)rectangle++(0.008,0.008) (-0.298,0.803)rectangle++(0.008,0.008) (-0.427,0.68)rectangle++(0.008,0.008) (-0.306,0.896)rectangle++(0.008,0.008) (-0.464,0.719)rectangle++(0.008,0.008)
  (-0.02,0.791)rectangle++(0.008,0.008) (-0.154,0.949)rectangle++(0.008,0.008) (-0.302,0.905)rectangle++(0.008,0.008) (-0.465,0.725)rectangle++(0.008,0.008) (-0.475,0.585)rectangle++(0.008,0.008)
  (-0.016,0.859)rectangle++(0.008,0.008) (-0.318,0.806)rectangle++(0.008,0.008) (-0.437,0.674)rectangle++(0.008,0.008) (-0.149,0.694)rectangle++(0.008,0.008) (-0.351,0.802)rectangle++(0.008,0.008)
  (-0.45,0.661)rectangle++(0.008,0.008) (-0.128,0.585)rectangle++(0.008,0.008) (-0.271,0.641)rectangle++(0.008,0.008) (-0.011,0.85)rectangle++(0.008,0.008) (-0.348,0.853)rectangle++(0.008,0.008)
  (-0.467,0.685)rectangle++(0.008,0.008) (-0.389,0.832)rectangle++(0.008,0.008) (-0.478,0.661)rectangle++(0.008,0.008) (-0.147,0.908)rectangle++(0.008,0.008) (-0.065,0.494)rectangle++(0.008,0.008)
  (-0.158,0.924)rectangle++(0.008,0.008) (-0.361,0.891)rectangle++(0.008,0.008) (-0.487,0.698)rectangle++(0.008,0.008) (-0.116,0.546)rectangle++(0.008,0.008) (-0.192,0.793)rectangle++(0.008,0.008)
  (-0.374,0.858)rectangle++(0.008,0.008) (-0.481,0.678)rectangle++(0.008,0.008) (-0.314,0.837)rectangle++(0.008,0.008) (-0.3,0.906)rectangle++(0.008,0.008) (-0.127,0.943)rectangle++(0.008,0.008)
  (-0.155,0.742)rectangle++(0.008,0.008) (-0.359,0.811)rectangle++(0.008,0.008) (-0.29,0.912)rectangle++(0.008,0.008) (-0.41,0.823)rectangle++(0.008,0.008) (-0.138,0.594)rectangle++(0.008,0.008)
  (-0.279,0.642)rectangle++(0.008,0.008) (-0.36,0.613)rectangle++(0.008,0.008) (-0.158,0.695)rectangle++(0.008,0.008) (-0.324,0.681)rectangle++(0.008,0.008) (-0.354,0.88)rectangle++(0.008,0.008)
  (-0.276,0.76)rectangle++(0.008,0.008) (-0.401,0.668)rectangle++(0.008,0.008) (-0.085,0.522)rectangle++(0.008,0.008) (-0.313,0.779)rectangle++(0.008,0.008) (-0.365,0.888)rectangle++(0.008,0.008)
  (-0.145,0.601)rectangle++(0.008,0.008) (-0.285,0.642)rectangle++(0.008,0.008) (-0.136,0.582)rectangle++(0.008,0.008) (-0.274,0.637)rectangle++(0.008,0.008) (-0.356,0.612)rectangle++(0.008,0.008)
  (-0.273,0.759)rectangle++(0.008,0.008) (-0.384,0.796)rectangle++(0.008,0.008) (-0.411,0.787)rectangle++(0.008,0.008) (-0.029,0.896)rectangle++(0.008,0.008) (-0.358,0.859)rectangle++(0.008,0.008)
  (-0.474,0.684)rectangle++(0.008,0.008) (-0.22,0.907)rectangle++(0.008,0.008) (-0.428,0.755)rectangle++(0.008,0.008) (-0.469,0.611)rectangle++(0.008,0.008) (-0.436,0.531)rectangle++(0.008,0.008)
  (-0.357,0.888)rectangle++(0.008,0.008) (-0.175,0.948)rectangle++(0.008,0.008) (-0.394,0.847)rectangle++(0.008,0.008) (-0.358,0.89)rectangle++(0.008,0.008) (-0.485,0.698)rectangle++(0.008,0.008)
  (-0.475,0.565)rectangle++(0.008,0.008) (-0.422,0.508)rectangle++(0.008,0.008) (-0.134,0.588)rectangle++(0.008,0.008) (-0.275,0.64)rectangle++(0.008,0.008) (-0.245,0.825)rectangle++(0.008,0.008)
  (-0.168,0.944)rectangle++(0.008,0.008) (-0.177,0.944)rectangle++(0.008,0.008) (-0.422,0.787)rectangle++(0.008,0.008) (-0.294,0.874)rectangle++(0.008,0.008) (-0.093,0.919)rectangle++(0.008,0.008)
  (-0.374,0.854)rectangle++(0.008,0.008) (-0.225,0.884)rectangle++(0.008,0.008) (-0.422,0.743)rectangle++(0.008,0.008) (-0.295,0.859)rectangle++(0.008,0.008) (-0.252,0.898)rectangle++(0.008,0.008)
  (-0.439,0.74)rectangle++(0.008,0.008) (-0.469,0.601)rectangle++(0.008,0.008) (-0.31,0.811)rectangle++(0.008,0.008) (-0.152,0.735)rectangle++(0.008,0.008) (-0.317,0.89)rectangle++(0.008,0.008)
  (-0.006,0.79)rectangle++(0.008,0.008) (-0.155,0.875)rectangle++(0.008,0.008) (-0.387,0.764)rectangle++(0.008,0.008) (-0.453,0.63)rectangle++(0.008,0.008) (-0.436,0.546)rectangle++(0.008,0.008)
  (-0.299,0.793)rectangle++(0.008,0.008) (-0.423,0.675)rectangle++(0.008,0.008) (-0.328,0.872)rectangle++(0.008,0.008) (-0.098,0.923)rectangle++(0.008,0.008) (-0.241,0.895)rectangle++(0.008,0.008)
  (-0.433,0.742)rectangle++(0.008,0.008) (-0.13,0.916)rectangle++(0.008,0.008) (-0.38,0.848)rectangle++(0.008,0.008) (-0.155,0.916)rectangle++(0.008,0.008) (-0.067,0.909)rectangle++(0.008,0.008)
  (-0.359,0.811)rectangle++(0.008,0.008) (-0.116,0.937)rectangle++(0.008,0.008) (-0.381,0.854)rectangle++(0.008,0.008) (-0.334,0.892)rectangle++(0.008,0.008) (-0.014,0.875)rectangle++(0.008,0.008)
  (-0.038,0.899)rectangle++(0.008,0.008) (-0.342,0.816)rectangle++(0.008,0.008) (-0.238,0.716)rectangle++(0.008,0.008) (-0.368,0.661)rectangle++(0.008,0.008) (-0.277,0.832)rectangle++(0.008,0.008)
  (-0.151,0.734)rectangle++(0.008,0.008) (-0.335,0.701)rectangle++(0.008,0.008) (-0.265,0.845)rectangle++(0.008,0.008) (-0.274,0.764)rectangle++(0.008,0.008) (-0.402,0.67)rectangle++(0.008,0.008)
  (-0.288,0.835)rectangle++(0.008,0.008) (0.005,0.847)rectangle++(0.008,0.008) (-0.344,0.855)rectangle++(0.008,0.008) (-0.059,0.487)rectangle++(0.008,0.008) (-0.205,0.621)rectangle++(0.008,0.008)
  (-0.04,0.853)rectangle++(0.008,0.008) (-0.166,0.896)rectangle++(0.008,0.008) (-0.4,0.77)rectangle++(0.008,0.008) (-0.287,0.868)rectangle++(0.008,0.008) (-0.093,0.917)rectangle++(0.008,0.008)
  (-0.088,0.909)rectangle++(0.008,0.008) (-0.368,0.804)rectangle++(0.008,0.008) (-0.271,0.754)rectangle++(0.008,0.008) (-0.397,0.667)rectangle++(0.008,0.008) (-0.263,0.743)rectangle++(0.008,0.008)
  (-0.327,0.873)rectangle++(0.008,0.008) (-0.465,0.702)rectangle++(0.008,0.008) (-0.467,0.574)rectangle++(0.008,0.008) (-0.422,0.515)rectangle++(0.008,0.008) (-0.38,0.504)rectangle++(0.008,0.008)
  (-0.058,0.869)rectangle++(0.008,0.008) (-0.057,0.89)rectangle++(0.008,0.008) (-0.244,0.723)rectangle++(0.008,0.008) (-0.373,0.794)rectangle++(0.008,0.008) (-0.458,0.649)rectangle++(0.008,0.008)
  (-0.044,0.882)rectangle++(0.008,0.008) (-0.339,0.806)rectangle++(0.008,0.008) (-0.446,0.667)rectangle++(0.008,0.008) (-0.446,0.565)rectangle++(0.008,0.008) (-0.38,0.823)rectangle++(0.008,0.008)
  (-0.471,0.66)rectangle++(0.008,0.008) (-0.408,0.748)rectangle++(0.008,0.008) (-0.408,0.774)rectangle++(0.008,0.008) (-0.283,0.817)rectangle++(0.008,0.008) (-0.425,0.692)rectangle++(0.008,0.008)
  (-0.153,0.917)rectangle++(0.008,0.008) (-0.401,0.783)rectangle++(0.008,0.008) (-0.12,0.918)rectangle++(0.008,0.008) (-0.379,0.85)rectangle++(0.008,0.008) (-0.02,0.878)rectangle++(0.008,0.008)
  (-0.159,0.905)rectangle++(0.008,0.008) (-0.206,0.914)rectangle++(0.008,0.008) (-0.424,0.763)rectangle++(0.008,0.008) (-0.295,0.866)rectangle++(0.008,0.008) (-0.448,0.71)rectangle++(0.008,0.008)
  (-0.462,0.584)rectangle++(0.008,0.008) (-0.423,0.521)rectangle++(0.008,0.008) (-0.383,0.506)rectangle++(0.008,0.008) (-0.359,0.514)rectangle++(0.008,0.008) (-0.274,0.782)rectangle++(0.008,0.008)
  (-0.388,0.801)rectangle++(0.008,0.008) (-0.13,0.566)rectangle++(0.008,0.008) (-0.091,0.524)rectangle++(0.008,0.008) (-0.232,0.626)rectangle++(0.008,0.008) (-0.366,0.842)rectangle++(0.008,0.008)
  (-0.472,0.673)rectangle++(0.008,0.008) (-0.46,0.559)rectangle++(0.008,0.008) (-0.39,0.73)rectangle++(0.008,0.008) (-0.402,0.773)rectangle++(0.008,0.008) (-0.14,0.939)rectangle++(0.008,0.008)
  (-0.2,0.925)rectangle++(0.008,0.008) (-0.096,0.926)rectangle++(0.008,0.008) (-0.378,0.808)rectangle++(0.008,0.008) (-0.114,0.919)rectangle++(0.008,0.008) (-0.153,0.605)rectangle++(0.008,0.008)
  (-0.018,0.872)rectangle++(0.008,0.008) (-0.323,0.812)rectangle++(0.008,0.008) (-0.261,0.882)rectangle++(0.008,0.008) (-0.009,0.831)rectangle++(0.008,0.008) (-0.02,0.859)rectangle++(0.008,0.008)
  (-0.319,0.805)rectangle++(0.008,0.008) (-0.027,0.892)rectangle++(0.008,0.008) (-0.03,0.886)rectangle++(0.008,0.008) (-0.041,0.464)rectangle++(0.008,0.008) (-0.075,0.505)rectangle++(0.008,0.008)
  (-0.218,0.624)rectangle++(0.008,0.008) (-0.226,0.819)rectangle++(0.008,0.008) (-0.399,0.713)rectangle++(0.008,0.008) (-0.305,0.898)rectangle++(0.008,0.008) (-0.143,0.586)rectangle++(0.008,0.008)
  (-0.335,0.88)rectangle++(0.008,0.008) (-0.471,0.702)rectangle++(0.008,0.008) (-0.47,0.572)rectangle++(0.008,0.008) (-0.097,0.522)rectangle++(0.008,0.008) (-0.234,0.623)rectangle++(0.008,0.008)
  (-0.333,0.62)rectangle++(0.008,0.008) (-0.377,0.584)rectangle++(0.008,0.008) (-0.385,0.551)rectangle++(0.008,0.008) (-0.375,0.74)rectangle++(0.008,0.008) (-0.026,0.899)rectangle++(0.008,0.008)
  (-0.033,0.814)rectangle++(0.008,0.008) (-0.127,0.573)rectangle++(0.008,0.008) (-0.354,0.844)rectangle++(0.008,0.008) (-0.134,0.942)rectangle++(0.008,0.008) (-0.401,0.802)rectangle++(0.008,0.008)
  (-0.118,0.921)rectangle++(0.008,0.008) (-0.386,0.798)rectangle++(0.008,0.008) (-0.283,0.878)rectangle++(0.008,0.008) (-0.404,0.818)rectangle++(0.008,0.008) (-0.266,0.903)rectangle++(0.008,0.008)
  (-0.151,0.735)rectangle++(0.008,0.008) (-0.175,0.914)rectangle++(0.008,0.008) (-0.211,0.917)rectangle++(0.008,0.008) (-0.224,0.918)rectangle++(0.008,0.008) (-0.118,0.915)rectangle++(0.008,0.008)
  (-0.384,0.795)rectangle++(0.008,0.008) (-0.463,0.645)rectangle++(0.008,0.008) (-0.156,0.731)rectangle++(0.008,0.008) (-0.24,0.892)rectangle++(0.008,0.008) (-0.432,0.741)rectangle++(0.008,0.008)
  (-0.466,0.604)rectangle++(0.008,0.008) (-0.159,0.694)rectangle++(0.008,0.008) (-0.176,0.908)rectangle++(0.008,0.008) (-0.142,0.596)rectangle++(0.008,0.008) (-0.281,0.641)rectangle++(0.008,0.008)
  (-0.361,0.612)rectangle++(0.008,0.008) (-0.387,0.57)rectangle++(0.008,0.008) (-0.283,0.801)rectangle++(0.008,0.008) (-0.393,0.803)rectangle++(0.008,0.008) (-0.414,0.787)rectangle++(0.008,0.008)
  (-0.141,0.941)rectangle++(0.008,0.008) (-0.075,0.508)rectangle++(0.008,0.008) (-0.082,0.514)rectangle++(0.008,0.008) (-0.18,0.782)rectangle++(0.008,0.008) (-0.228,0.893)rectangle++(0.008,0.008)
  (-0.383,0.864)rectangle++(0.008,0.008) (-0.275,0.758)rectangle++(0.008,0.008) (-0.4,0.668)rectangle++(0.008,0.008) (-0.395,0.76)rectangle++(0.008,0.008) (-0.145,0.951)rectangle++(0.008,0.008)
  (-0.068,0.498)rectangle++(0.008,0.008) (-0.163,0.881)rectangle++(0.008,0.008) (-0.027,0.446)rectangle++(0.008,0.008) (-0.162,0.759)rectangle++(0.008,0.008) (-0.339,0.894)rectangle++(0.008,0.008)
  (-0.355,0.893)rectangle++(0.008,0.008) (-0.266,0.848)rectangle++(0.008,0.008) (-0.116,0.918)rectangle++(0.008,0.008) (-0.192,0.918)rectangle++(0.008,0.008) (-0.281,0.794)rectangle++(0.008,0.008)
  (-0.259,0.896)rectangle++(0.008,0.008) (-0.402,0.825)rectangle++(0.008,0.008) (-0.482,0.653)rectangle++(0.008,0.008) (-0.349,0.864)rectangle++(0.008,0.008) (-0.203,0.932)rectangle++(0.008,0.008)
  (-0.429,0.773)rectangle++(0.008,0.008) (-0.391,0.841)rectangle++(0.008,0.008) (-0.087,0.536)rectangle++(0.008,0.008) (-0.315,0.781)rectangle++(0.008,0.008) (-0.353,0.891)rectangle++(0.008,0.008)
  (-0.062,0.866)rectangle++(0.008,0.008) (-0.173,0.901)rectangle++(0.008,0.008) (-0.404,0.769)rectangle++(0.008,0.008) (-0.389,0.842)rectangle++(0.008,0.008) (-0.304,0.906)rectangle++(0.008,0.008)
  (-0.281,0.769)rectangle++(0.008,0.008) (-0.406,0.67)rectangle++(0.008,0.008) (-0.289,0.835)rectangle++(0.008,0.008) (-0.434,0.697)rectangle++(0.008,0.008) (-0.451,0.583)rectangle++(0.008,0.008)
  (-0.217,0.893)rectangle++(0.008,0.008) (-0.394,0.83)rectangle++(0.008,0.008) (-0.48,0.658)rectangle++(0.008,0.008) (-0.363,0.879)rectangle++(0.008,0.008) (-0.332,0.896)rectangle++(0.008,0.008)
  (-0.087,0.899)rectangle++(0.008,0.008) (-0.365,0.799)rectangle++(0.008,0.008) (-0.054,0.48)rectangle++(0.008,0.008) (-0.2,0.619)rectangle++(0.008,0.008) (-0.316,0.631)rectangle++(0.008,0.008)
  (-0.269,0.887)rectangle++(0.008,0.008) (-0.107,0.808)rectangle++(0.008,0.008) (-0.341,0.751)rectangle++(0.008,0.008) (-0.396,0.785)rectangle++(0.008,0.008) (-0.465,0.636)rectangle++(0.008,0.008)
  (-0.443,0.544)rectangle++(0.008,0.008) (-0.401,0.51)rectangle++(0.008,0.008) (-0.153,0.876)rectangle++(0.008,0.008) (-0.379,0.837)rectangle++(0.008,0.008) (-0.28,0.891)rectangle++(0.008,0.008)
  (-0.45,0.726)rectangle++(0.008,0.008) (-0.05,0.475)rectangle++(0.008,0.008) (-0.288,0.838)rectangle++(0.008,0.008) (-0.399,0.809)rectangle++(0.008,0.008) (-0.367,0.883)rectangle++(0.008,0.008)
  (-0.487,0.692)rectangle++(0.008,0.008) (-0.473,0.562)rectangle++(0.008,0.008) (-0.393,0.728)rectangle++(0.008,0.008) (-0.119,0.552)rectangle++(0.008,0.008) (-0.193,0.795)rectangle++(0.008,0.008)
  (-0.376,0.714)rectangle++(0.008,0.008) (-0.43,0.611)rectangle++(0.008,0.008) (-0.341,0.858)rectangle++(0.008,0.008) (-0.466,0.689)rectangle++(0.008,0.008) (-0.463,0.568)rectangle++(0.008,0.008)
  (-0.131,0.578)rectangle++(0.008,0.008) (-0.104,0.854)rectangle++(0.008,0.008) (-0.134,0.58)rectangle++(0.008,0.008) (-0.331,0.783)rectangle++(0.008,0.008) (-0.157,0.719)rectangle++(0.008,0.008)
  (-0.176,0.912)rectangle++(0.008,0.008) (-0.41,0.773)rectangle++(0.008,0.008) (-0.467,0.626)rectangle++(0.008,0.008) (-0.31,0.82)rectangle++(0.008,0.008) (-0.279,0.766)rectangle++(0.008,0.008)
  (-0.265,0.894)rectangle++(0.008,0.008) (-0.403,0.823)rectangle++(0.008,0.008) (-0.11,0.558)rectangle++(0.008,0.008) (-0.253,0.635)rectangle++(0.008,0.008) (-0.346,0.619)rectangle++(0.008,0.008)
  (-0.377,0.759)rectangle++(0.008,0.008) (-0.447,0.631)rectangle++(0.008,0.008) (-0.433,0.548)rectangle++(0.008,0.008) (-0.397,0.515)rectangle++(0.008,0.008) (-0.131,0.587)rectangle++(0.008,0.008)
  (-0.273,0.641)rectangle++(0.008,0.008) (-0.357,0.615)rectangle++(0.008,0.008) (-0.386,0.573)rectangle++(0.008,0.008) (-0.277,0.887)rectangle++(0.008,0.008) (-0.447,0.726)rectangle++(0.008,0.008)
  (0.002,0.849)rectangle++(0.008,0.008) (-0.345,0.855)rectangle++(0.008,0.008) (-0.269,0.897)rectangle++(0.008,0.008) (-0.447,0.733)rectangle++(0.008,0.008) (-0.136,0.592)rectangle++(0.008,0.008)
  (-0.277,0.641)rectangle++(0.008,0.008) (-0.359,0.613)rectangle++(0.008,0.008) (-0.374,0.865)rectangle++(0.008,0.008) (-0.093,0.539)rectangle++(0.008,0.008) (-0.253,0.734)rectangle++(0.008,0.008)
  (-0.381,0.664)rectangle++(0.008,0.008) (-0.415,0.587)rectangle++(0.008,0.008) (-0.292,0.807)rectangle++(0.008,0.008) (-0.425,0.684)rectangle++(0.008,0.008) (-0.442,0.58)rectangle++(0.008,0.008)
  (0.003,0.796)rectangle++(0.008,0.008) (-0.144,0.596)rectangle++(0.008,0.008) (-0.101,0.884)rectangle++(0.008,0.008) (-0.365,0.787)rectangle++(0.008,0.008) (-0.406,0.788)rectangle++(0.008,0.008)
  (-0.289,0.874)rectangle++(0.008,0.008) (-0.448,0.716)rectangle++(0.008,0.008) (-0.033,0.897)rectangle++(0.008,0.008) (-0.359,0.859)rectangle++(0.008,0.008) (-0.118,0.574)rectangle++(0.008,0.008)
  (-0.262,0.639)rectangle++(0.008,0.008) (-0.116,0.879)rectangle++(0.008,0.008) (-0.371,0.78)rectangle++(0.008,0.008) (-0.138,0.94)rectangle++(0.008,0.008) (-0.069,0.499)rectangle++(0.008,0.008)
  (-0.045,0.82)rectangle++(0.008,0.008) (-0.33,0.833)rectangle++(0.008,0.008) (-0.117,0.571)rectangle++(0.008,0.008) (-0.261,0.639)rectangle++(0.008,0.008) (-0.24,0.824)rectangle++(0.008,0.008)
  (-0.295,0.905)rectangle++(0.008,0.008) (-0.462,0.727)rectangle++(0.008,0.008) (-0.308,0.854)rectangle++(0.008,0.008) (-0.449,0.7)rectangle++(0.008,0.008) (-0.116,0.546)rectangle++(0.008,0.008)
  (-0.192,0.793)rectangle++(0.008,0.008) (-0.374,0.815)rectangle++(0.008,0.008) (-0.278,0.884)rectangle++(0.008,0.008) (-0.446,0.724)rectangle++(0.008,0.008) (-0.466,0.591)rectangle++(0.008,0.008)
  (-0.396,0.735)rectangle++(0.008,0.008) (-0.286,0.857)rectangle++(0.008,0.008) (-0.137,0.578)rectangle++(0.008,0.008) (-0.273,0.635)rectangle++(0.008,0.008) (-0.17,0.932)rectangle++(0.008,0.008)
  (-0.164,0.95)rectangle++(0.008,0.008) (-0.208,0.929)rectangle++(0.008,0.008) (-0.124,0.58)rectangle++(0.008,0.008) (-0.267,0.64)rectangle++(0.008,0.008) (-0.045,0.468)rectangle++(0.008,0.008)
  (-0.13,0.78)rectangle++(0.008,0.008) (-0.36,0.822)rectangle++(0.008,0.008) (-0.462,0.666)rectangle++(0.008,0.008) (-0.219,0.904)rectangle++(0.008,0.008) (-0.359,0.892)rectangle++(0.008,0.008)
  (-0.487,0.699)rectangle++(0.008,0.008) (-0.034,0.871)rectangle++(0.008,0.008) (-0.355,0.854)rectangle++(0.008,0.008) (-0.471,0.683)rectangle++(0.008,0.008) (-0.282,0.873)rectangle++(0.008,0.008)
  (-0.248,0.903)rectangle++(0.008,0.008) (-0.385,0.863)rectangle++(0.008,0.008) (-0.208,0.932)rectangle++(0.008,0.008) (-0.431,0.771)rectangle++(0.008,0.008) (-0.286,0.914)rectangle++(0.008,0.008)
  (-0.461,0.735)rectangle++(0.008,0.008) (-0.477,0.591)rectangle++(0.008,0.008) (-0.398,0.734)rectangle++(0.008,0.008) (-0.355,0.89)rectangle++(0.008,0.008) (-0.259,0.906)rectangle++(0.008,0.008)
  (-0.445,0.741)rectangle++(0.008,0.008) (-0.472,0.599)rectangle++(0.008,0.008) (-0.433,0.525)rectangle++(0.008,0.008) (-0.389,0.504)rectangle++(0.008,0.008) (-0.152,0.874)rectangle++(0.008,0.008)
  (-0.204,0.903)rectangle++(0.008,0.008) (-0.137,0.592)rectangle++(0.008,0.008) (-0.278,0.641)rectangle++(0.008,0.008) (-0.368,0.773)rectangle++(0.008,0.008) (-0.066,0.863)rectangle++(0.008,0.008)
  (-0.342,0.79)rectangle++(0.008,0.008) (-0.442,0.658)rectangle++(0.008,0.008) (-0.357,0.868)rectangle++(0.008,0.008) (-0.275,0.759)rectangle++(0.008,0.008) (-0.144,0.591)rectangle++(0.008,0.008)
  (-0.015,0.807)rectangle++(0.008,0.008) (-0.158,0.881)rectangle++(0.008,0.008) (-0.276,0.764)rectangle++(0.008,0.008) (-0.403,0.67)rectangle++(0.008,0.008) (-0.427,0.582)rectangle++(0.008,0.008)
  (-0.406,0.533)rectangle++(0.008,0.008) (-0.247,0.728)rectangle++(0.008,0.008) (-0.15,0.951)rectangle++(0.008,0.008) (-0.149,0.598)rectangle++(0.008,0.008) (-0.163,0.93)rectangle++(0.008,0.008)
  (-0.279,0.77)rectangle++(0.008,0.008) (-0.406,0.671)rectangle++(0.008,0.008) (-0.228,0.854)rectangle++(0.008,0.008) (-0.229,0.896)rectangle++(0.008,0.008) (-0.156,0.735)rectangle++(0.008,0.008)
  (-0.338,0.699)rectangle++(0.008,0.008) (-0.065,0.494)rectangle++(0.008,0.008) (-0.21,0.622)rectangle++(0.008,0.008) (-0.223,0.819)rectangle++(0.008,0.008) (-0.398,0.714)rectangle++(0.008,0.008)
  (-0.128,0.907)rectangle++(0.008,0.008) (-0.386,0.788)rectangle++(0.008,0.008) (-0.146,0.593)rectangle++(0.008,0.008) (-0.371,0.881)rectangle++(0.008,0.008) (-0.333,0.896)rectangle++(0.008,0.008)
  (-0.014,0.875)rectangle++(0.008,0.008) (-0.144,0.599)rectangle++(0.008,0.008) (-0.201,0.811)rectangle++(0.008,0.008) (-0.378,0.817)rectangle++(0.008,0.008) (-0.113,0.92)rectangle++(0.008,0.008)
  (-0.191,0.919)rectangle++(0.008,0.008) (-0.02,0.888)rectangle++(0.008,0.008) (-0.33,0.818)rectangle++(0.008,0.008) (-0.094,0.544)rectangle++(0.008,0.008) (-0.27,0.844)rectangle++(0.008,0.008)
  (-0.428,0.709)rectangle++(0.008,0.008) (-0.36,0.882)rectangle++(0.008,0.008) (-0.332,0.896)rectangle++(0.008,0.008) (-0.281,0.768)rectangle++(0.008,0.008) (-0.406,0.67)rectangle++(0.008,0.008)
  (-0.397,0.759)rectangle++(0.008,0.008) (-0.286,0.865)rectangle++(0.008,0.008) (-0.443,0.712)rectangle++(0.008,0.008) (-0.11,0.936)rectangle++(0.008,0.008) (-0.154,0.607)rectangle++(0.008,0.008)
  (-0.204,0.813)rectangle++(0.008,0.008) (-0.387,0.718)rectangle++(0.008,0.008) (-0.156,0.736)rectangle++(0.008,0.008) (-0.282,0.815)rectangle++(0.008,0.008) (-0.136,0.589)rectangle++(0.008,0.008)
  (-0.333,0.784)rectangle++(0.008,0.008) (-0.264,0.873)rectangle++(0.008,0.008) (-0.089,0.915)rectangle++(0.008,0.008) (-0.371,0.806)rectangle++(0.008,0.008) (-0.271,0.754)rectangle++(0.008,0.008)
  (-0.397,0.667)rectangle++(0.008,0.008) (-0.286,0.834)rectangle++(0.008,0.008) (-0.325,0.893)rectangle++(0.008,0.008) (-0.043,0.873)rectangle++(0.008,0.008) (-0.287,0.895)rectangle++(0.008,0.008)
  (-0.141,0.597)rectangle++(0.008,0.008) (-0.244,0.877)rectangle++(0.008,0.008) (-0.235,0.904)rectangle++(0.008,0.008) (-0.434,0.749)rectangle++(0.008,0.008) (-0.469,0.607)rectangle++(0.008,0.008)
  (-0.435,0.529)rectangle++(0.008,0.008) (-0.099,0.815)rectangle++(0.008,0.008) (-0.047,0.899)rectangle++(0.008,0.008) (-0.219,0.87)rectangle++(0.008,0.008) (-0.009,0.789)rectangle++(0.008,0.008)
  (-0.042,0.893)rectangle++(0.008,0.008) (-0.152,0.85)rectangle++(0.008,0.008) (-0.302,0.891)rectangle++(0.008,0.008) (-0.323,0.897)rectangle++(0.008,0.008) (-0.173,0.948)rectangle++(0.008,0.008)
  (-0.421,0.791)rectangle++(0.008,0.008) (-0.478,0.63)rectangle++(0.008,0.008) (-0.153,0.729)rectangle++(0.008,0.008) (-0.204,0.854)rectangle++(0.008,0.008) (-0.376,0.878)rectangle++(0.008,0.008)
  (-0.366,0.89)rectangle++(0.008,0.008) (-0.489,0.695)rectangle++(0.008,0.008) (-0.475,0.562)rectangle++(0.008,0.008) (-0.394,0.728)rectangle++(0.008,0.008) (-0.444,0.612)rectangle++(0.008,0.008)
  (-0.395,0.742)rectangle++(0.008,0.008) (-0.449,0.618)rectangle++(0.008,0.008) (-0.429,0.541)rectangle++(0.008,0.008) (-0.297,0.791)rectangle++(0.008,0.008) (-0.272,0.908)rectangle++(0.008,0.008)
  (-0.132,0.581)rectangle++(0.008,0.008) (-0.299,0.854)rectangle++(0.008,0.008) (-0.017,0.87)rectangle++(0.008,0.008) (-0.158,0.902)rectangle++(0.008,0.008) (-0.398,0.775)rectangle++(0.008,0.008)
  (-0.068,0.912)rectangle++(0.008,0.008) (-0.163,0.939)rectangle++(0.008,0.008) (-0.279,0.766)rectangle++(0.008,0.008) (-0.386,0.797)rectangle++(0.008,0.008) (-0.313,0.892)rectangle++(0.008,0.008)
  (-0.069,0.823)rectangle++(0.008,0.008) (-0.067,0.497)rectangle++(0.008,0.008) (-0.175,0.776)rectangle++(0.008,0.008) (-0.006,0.815)rectangle++(0.008,0.008) (-0.297,0.789)rectangle++(0.008,0.008)
  (-0.421,0.674)rectangle++(0.008,0.008) (-0.437,0.577)rectangle++(0.008,0.008) (-0.409,0.527)rectangle++(0.008,0.008) (-0.379,0.514)rectangle++(0.008,0.008) (-0.36,0.519)rectangle++(0.008,0.008)
  (-0.163,0.943)rectangle++(0.008,0.008) (-0.415,0.792)rectangle++(0.008,0.008) (-0.476,0.633)rectangle++(0.008,0.008) (-0.447,0.539)rectangle++(0.008,0.008) (-0.303,0.791)rectangle++(0.008,0.008)
  (-0.424,0.672)rectangle++(0.008,0.008) (-0.152,0.732)rectangle++(0.008,0.008) (-0.175,0.914)rectangle++(0.008,0.008) (-0.41,0.774)rectangle++(0.008,0.008) (-0.467,0.627)rectangle++(0.008,0.008)
  (-0.441,0.539)rectangle++(0.008,0.008) (-0.398,0.508)rectangle++(0.008,0.008) (-0.371,0.729)rectangle++(0.008,0.008) (-0.433,0.62)rectangle++(0.008,0.008) (-0.366,0.879)rectangle++(0.008,0.008)
  (-0.485,0.691)rectangle++(0.008,0.008) (-0.115,0.544)rectangle++(0.008,0.008) (-0.25,0.627)rectangle++(0.008,0.008) (-0.237,0.82)rectangle++(0.008,0.008) (-0.405,0.71)rectangle++(0.008,0.008)
  (-0.442,0.599)rectangle++(0.008,0.008) (-0.372,0.888)rectangle++(0.008,0.008) (-0.491,0.692)rectangle++(0.008,0.008) (-0.475,0.56)rectangle++(0.008,0.008) (-0.159,0.893)rectangle++(0.008,0.008)
  (-0.395,0.771)rectangle++(0.008,0.008) (-0.121,0.915)rectangle++(0.008,0.008) (-0.386,0.794)rectangle++(0.008,0.008) (-0.273,0.89)rectangle++(0.008,0.008) (-0.142,0.585)rectangle++(0.008,0.008)
  (-0.301,0.854)rectangle++(0.008,0.008) (-0.224,0.881)rectangle++(0.008,0.008) (-0.393,0.827)rectangle++(0.008,0.008) (-0.285,0.888)rectangle++(0.008,0.008) (-0.248,0.908)rectangle++(0.008,0.008)
  (-0.236,0.915)rectangle++(0.008,0.008) (-0.122,0.578)rectangle++(0.008,0.008) (-0.266,0.64)rectangle++(0.008,0.008) (-0.354,0.617)rectangle++(0.008,0.008) (-0.385,0.575)rectangle++(0.008,0.008)
  (-0.143,0.884)rectangle++(0.008,0.008) (-0.126,0.573)rectangle++(0.008,0.008) (-0.195,0.802)rectangle++(0.008,0.008) (-0.218,0.879)rectangle++(0.008,0.008) (-0.081,0.911)rectangle++(0.008,0.008)
  (-0.165,0.939)rectangle++(0.008,0.008) (-0.01,0.785)rectangle++(0.008,0.008) (-0.061,0.824)rectangle++(0.008,0.008) (-0.157,0.944)rectangle++(0.008,0.008) (-0.39,0.849)rectangle++(0.008,0.008)
  (-0.485,0.668)rectangle++(0.008,0.008) (-0.218,0.852)rectangle++(0.008,0.008) (-0.226,0.896)rectangle++(0.008,0.008) (-0.167,0.948)rectangle++(0.008,0.008) (-0.418,0.793)rectangle++(0.008,0.008)
  (-0.417,0.781)rectangle++(0.008,0.008) (-0.473,0.627)rectangle++(0.008,0.008) (-0.444,0.537)rectangle++(0.008,0.008) (-0.384,0.728)rectangle++(0.008,0.008) (-0.281,0.829)rectangle++(0.008,0.008)
  (-0.428,0.698)rectangle++(0.008,0.008) (-0.297,0.844)rectangle++(0.008,0.008) (-0.441,0.699)rectangle++(0.008,0.008) (-0.455,0.581)rectangle++(0.008,0.008) (-0.393,0.735)rectangle++(0.008,0.008)
  (-0.125,0.91)rectangle++(0.008,0.008) (-0.143,0.59)rectangle++(0.008,0.008) (-0.174,0.894)rectangle++(0.008,0.008) (-0.247,0.897)rectangle++(0.008,0.008) (-0.437,0.741)rectangle++(0.008,0.008)
  (-0.267,0.753)rectangle++(0.008,0.008) (-0.382,0.796)rectangle++(0.008,0.008) (-0.463,0.647)rectangle++(0.008,0.008) (-0.264,0.902)rectangle++(0.008,0.008) (-0.143,0.587)rectangle++(0.008,0.008)
  (-0.017,0.797)rectangle++(0.008,0.008) (-0.151,0.72)rectangle++(0.008,0.008) (-0.254,0.884)rectangle++(0.008,0.008) (-0.436,0.733)rectangle++(0.008,0.008) (-0.342,0.875)rectangle++(0.008,0.008)
  (-0.472,0.697)rectangle++(0.008,0.008) (-0.116,0.546)rectangle++(0.008,0.008) (-0.251,0.627)rectangle++(0.008,0.008) (-0.254,0.897)rectangle++(0.008,0.008) (-0.44,0.739)rectangle++(0.008,0.008)
  (-0.413,0.767)rectangle++(0.008,0.008) (-0.413,0.777)rectangle++(0.008,0.008) (-0.414,0.779)rectangle++(0.008,0.008) (-0.471,0.627)rectangle++(0.008,0.008) (-0.443,0.538)rectangle++(0.008,0.008)
  (-0.152,0.737)rectangle++(0.008,0.008) (-0.302,0.875)rectangle++(0.008,0.008) (-0.196,0.934)rectangle++(0.008,0.008) (-0.427,0.776)rectangle++(0.008,0.008) (-0.416,0.777)rectangle++(0.008,0.008)
  (-0.127,0.561)rectangle++(0.008,0.008) (-0.195,0.798)rectangle++(0.008,0.008) (-0.375,0.816)rectangle++(0.008,0.008) (-0.412,0.792)rectangle++(0.008,0.008) (-0.416,0.782)rectangle++(0.008,0.008)
  (-0.473,0.628)rectangle++(0.008,0.008) (-0.444,0.537)rectangle++(0.008,0.008) (-0.398,0.506)rectangle++(0.008,0.008) (-0.335,0.86)rectangle++(0.008,0.008) (-0.101,0.922)rectangle++(0.008,0.008)
  (-0.376,0.853)rectangle++(0.008,0.008) (-0.033,0.787)rectangle++(0.008,0.008) (-0.3,0.767)rectangle++(0.008,0.008) (-0.253,0.867)rectangle++(0.008,0.008) (-0.429,0.725)rectangle++(0.008,0.008)
  (-0.459,0.597)rectangle++(0.008,0.008) (-0.427,0.529)rectangle++(0.008,0.008) (-0.24,0.883)rectangle++(0.008,0.008) (-0.429,0.737)rectangle++(0.008,0.008) (-0.297,0.857)rectangle++(0.008,0.008)
  (-0.134,0.592)rectangle++(0.008,0.008) (-0.131,0.581)rectangle++(0.008,0.008) (-0.103,0.854)rectangle++(0.008,0.008) (-0.256,0.887)rectangle++(0.008,0.008) (-0.4,0.823)rectangle++(0.008,0.008)
  (-0.481,0.653)rectangle++(0.008,0.008) (-0.408,0.745)rectangle++(0.008,0.008) (-0.456,0.614)rectangle++(0.008,0.008) (-0.431,0.537)rectangle++(0.008,0.008) (-0.154,0.823)rectangle++(0.008,0.008)
  (-0.204,0.886)rectangle++(0.008,0.008) (-0.221,0.907)rectangle++(0.008,0.008) (-0.156,0.61)rectangle++(0.008,0.008) (-0.293,0.642)rectangle++(0.008,0.008) (-0.136,0.891)rectangle++(0.008,0.008)
  (-0.3,0.897)rectangle++(0.008,0.008) (-0.143,0.586)rectangle++(0.008,0.008) (-0.201,0.806)rectangle++(0.008,0.008) (-0.383,0.716)rectangle++(0.008,0.008) (-0.435,0.61)rectangle++(0.008,0.008)
  (-0.42,0.543)rectangle++(0.008,0.008) (-0.119,0.571)rectangle++(0.008,0.008) (-0.297,0.852)rectangle++(0.008,0.008) (-0.253,0.896)rectangle++(0.008,0.008) (-0.401,0.826)rectangle++(0.008,0.008)
  (-0.482,0.653)rectangle++(0.008,0.008) (-0.256,0.89)rectangle++(0.008,0.008) (-0.438,0.735)rectangle++(0.008,0.008) (-0.466,0.599)rectangle++(0.008,0.008) (-0.43,0.526)rectangle++(0.008,0.008)
  (-0.388,0.506)rectangle++(0.008,0.008) (-0.362,0.512)rectangle++(0.008,0.008) (-0.364,0.736)rectangle++(0.008,0.008) (-0.433,0.625)rectangle++(0.008,0.008) (-0.156,0.73)rectangle++(0.008,0.008)
  (-0.153,0.916)rectangle++(0.008,0.008) (-0.108,0.936)rectangle++(0.008,0.008) (-0.387,0.808)rectangle++(0.008,0.008) (-0.469,0.65)rectangle++(0.008,0.008) (-0.154,0.647)rectangle++(0.008,0.008)
  (-0.154,0.624)rectangle++(0.008,0.008) (-0.29,0.893)rectangle++(0.008,0.008) (-0.388,0.86)rectangle++(0.008,0.008) (-0.283,0.899)rectangle++(0.008,0.008) (-0.248,0.912)rectangle++(0.008,0.008)
  (-0.442,0.747)rectangle++(0.008,0.008) (-0.473,0.603)rectangle++(0.008,0.008) (-0.312,0.812)rectangle++(0.008,0.008) (-0.436,0.679)rectangle++(0.008,0.008) (-0.299,0.838)rectangle++(0.008,0.008)
  (-0.144,0.949)rectangle++(0.008,0.008) (0.003,0.808)rectangle++(0.008,0.008) (-0.119,0.574)rectangle++(0.008,0.008) (-0.263,0.639)rectangle++(0.008,0.008) (-0.225,0.878)rectangle++(0.008,0.008)
  (-0.42,0.74)rectangle++(0.008,0.008) (-0.41,0.771)rectangle++(0.008,0.008) (-0.268,0.901)rectangle++(0.008,0.008) (-0.448,0.735)rectangle++(0.008,0.008) (-0.414,0.766)rectangle++(0.008,0.008)
  (-0.466,0.621)rectangle++(0.008,0.008) (-0.309,0.818)rectangle++(0.008,0.008) (-0.437,0.683)rectangle++(0.008,0.008) (-0.342,0.868)rectangle++(0.008,0.008) (-0.202,0.933)rectangle++(0.008,0.008)
  (-0.221,0.923)rectangle++(0.008,0.008) (-0.005,0.825)rectangle++(0.008,0.008) (-0.3,0.795)rectangle++(0.008,0.008) (-0.425,0.675)rectangle++(0.008,0.008) (-0.439,0.576)rectangle++(0.008,0.008)
  (-0.153,0.864)rectangle++(0.008,0.008) (-0.382,0.76)rectangle++(0.008,0.008) (-0.45,0.63)rectangle++(0.008,0.008) (-0.399,0.745)rectangle++(0.008,0.008) (-0.287,0.86)rectangle++(0.008,0.008)
  (-0.134,0.592)rectangle++(0.008,0.008) (-0.095,0.529)rectangle++(0.008,0.008) (-0.236,0.627)rectangle++(0.008,0.008) (-0.357,0.777)rectangle++(0.008,0.008) (-0.403,0.787)rectangle++(0.008,0.008)
  (-0.255,0.739)rectangle++(0.008,0.008) (-0.189,0.915)rectangle++(0.008,0.008) (-0.417,0.77)rectangle++(0.008,0.008) (-0.126,0.56)rectangle++(0.008,0.008) (-0.261,0.63)rectangle++(0.008,0.008)
  (-0.348,0.614)rectangle++(0.008,0.008) (-0.203,0.897)rectangle++(0.008,0.008) (-0.156,0.735)rectangle++(0.008,0.008) (-0.337,0.7)rectangle++(0.008,0.008) (-0.328,0.87)rectangle++(0.008,0.008)
  (-0.464,0.7)rectangle++(0.008,0.008) (-0.466,0.573)rectangle++(0.008,0.008) (-0.421,0.515)rectangle++(0.008,0.008) (-0.38,0.504)rectangle++(0.008,0.008) (-0.371,0.823)rectangle++(0.008,0.008)
  (-0.467,0.663)rectangle++(0.008,0.008) (-0.347,0.865)rectangle++(0.008,0.008) (-0.202,0.933)rectangle++(0.008,0.008) (-0.429,0.773)rectangle++(0.008,0.008) (-0.391,0.841)rectangle++(0.008,0.008)
  (-0.483,0.664)rectangle++(0.008,0.008) (-0.461,0.55)rectangle++(0.008,0.008) (-0.411,0.506)rectangle++(0.008,0.008) (-0.024,0.878)rectangle++(0.008,0.008) (-0.161,0.905)rectangle++(0.008,0.008)
  (-0.385,0.841)rectangle++(0.008,0.008) (-0.48,0.666)rectangle++(0.008,0.008) (-0.461,0.553)rectangle++(0.008,0.008) (-0.412,0.507)rectangle++(0.008,0.008) (-0.373,0.504)rectangle++(0.008,0.008)
  (-0.153,0.941)rectangle++(0.008,0.008) (-0.389,0.849)rectangle++(0.008,0.008) (-0.111,0.927)rectangle++(0.008,0.008) (-0.385,0.803)rectangle++(0.008,0.008) (-0.412,0.788)rectangle++(0.008,0.008)
  (-0.291,0.874)rectangle++(0.008,0.008) (-0.449,0.715)rectangle++(0.008,0.008) (-0.012,0.426)rectangle++(0.008,0.008) (-0.161,0.61)rectangle++(0.008,0.008) (-0.295,0.64)rectangle++(0.008,0.008)
  (-0.134,0.94)rectangle++(0.008,0.008) (-0.198,0.925)rectangle++(0.008,0.008) (-0.424,0.772)rectangle++(0.008,0.008) (-0.473,0.62)rectangle++(0.008,0.008) (-0.348,0.859)rectangle++(0.008,0.008)
  (-0.203,0.932)rectangle++(0.008,0.008) (-0.158,0.613)rectangle++(0.008,0.008) (-0.295,0.643)rectangle++(0.008,0.008) (-0.368,0.608)rectangle++(0.008,0.008) (-0.389,0.566)rectangle++(0.008,0.008)
  (-0.283,0.8)rectangle++(0.008,0.008) (-0.134,0.59)rectangle++(0.008,0.008) (-0.284,0.906)rectangle++(0.008,0.008) (-0.457,0.732)rectangle++(0.008,0.008) (-0.218,0.914)rectangle++(0.008,0.008)
  (-0.397,0.834)rectangle++(0.008,0.008) (-0.483,0.659)rectangle++(0.008,0.008) (-0.46,0.548)rectangle++(0.008,0.008) (-0.012,0.794)rectangle++(0.008,0.008) (-0.263,0.853)rectangle++(0.008,0.008)
  (-0.428,0.715)rectangle++(0.008,0.008) (-0.455,0.593)rectangle++(0.008,0.008) (-0.306,0.809)rectangle++(0.008,0.008) (-0.161,0.95)rectangle++(0.008,0.008) (-0.372,0.887)rectangle++(0.008,0.008)
  (-0.206,0.936)rectangle++(0.008,0.008) (-0.222,0.924)rectangle++(0.008,0.008) (-0.124,0.581)rectangle++(0.008,0.008) (-0.06,0.885)rectangle++(0.008,0.008) (-0.056,0.893)rectangle++(0.008,0.008)
  (-0.171,0.91)rectangle++(0.008,0.008) (-0.407,0.774)rectangle++(0.008,0.008) (-0.305,0.902)rectangle++(0.008,0.008) (-0.275,0.839)rectangle++(0.008,0.008) (-0.397,0.811)rectangle++(0.008,0.008)
  (-0.416,0.788)rectangle++(0.008,0.008) (-0.416,0.78)rectangle++(0.008,0.008) (-0.345,0.894)rectangle++(0.008,0.008) (-0.202,0.937)rectangle++(0.008,0.008) (-0.088,0.912)rectangle++(0.008,0.008)
  (-0.148,0.604)rectangle++(0.008,0.008) (-0.287,0.643)rectangle++(0.008,0.008) (-0.364,0.611)rectangle++(0.008,0.008) (-0.388,0.568)rectangle++(0.008,0.008) (-0.144,0.883)rectangle++(0.008,0.008)
  (-0.201,0.906)rectangle++(0.008,0.008) (-0.248,0.899)rectangle++(0.008,0.008) (-0.143,0.587)rectangle++(0.008,0.008) (-0.278,0.636)rectangle++(0.008,0.008) (-0.059,0.486)rectangle++(0.008,0.008)
  (-0.204,0.62)rectangle++(0.008,0.008) (-0.221,0.818)rectangle++(0.008,0.008) (-0.397,0.714)rectangle++(0.008,0.008) (-0.286,0.85)rectangle++(0.008,0.008) (-0.438,0.705)rectangle++(0.008,0.008)
  (-0.456,0.585)rectangle++(0.008,0.008) (-0.421,0.524)rectangle++(0.008,0.008) (-0.377,0.729)rectangle++(0.008,0.008) (-0.118,0.571)rectangle++(0.008,0.008) (-0.246,0.876)rectangle++(0.008,0.008)
  (-0.282,0.771)rectangle++(0.008,0.008) (-0.388,0.797)rectangle++(0.008,0.008) (-0.304,0.903)rectangle++(0.008,0.008) (-0.012,0.872)rectangle++(0.008,0.008) (-0.255,0.738)rectangle++(0.008,0.008)
  (-0.378,0.795)rectangle++(0.008,0.008) (-0.376,0.881)rectangle++(0.008,0.008) (-0.279,0.906)rectangle++(0.008,0.008) (-0.115,0.795)rectangle++(0.008,0.008) (-0.34,0.741)rectangle++(0.008,0.008)
  (-0.395,0.783)rectangle++(0.008,0.008) (-0.067,0.912)rectangle++(0.008,0.008) (-0.36,0.812)rectangle++(0.008,0.008) (-0.458,0.662)rectangle++(0.008,0.008) (-0.11,0.539)rectangle++(0.008,0.008)
  (-0.351,0.842)rectangle++(0.008,0.008) (-0.465,0.679)rectangle++(0.008,0.008) (-0.277,0.898)rectangle++(0.008,0.008) (-0.32,0.898)rectangle++(0.008,0.008) (-0.063,0.492)rectangle++(0.008,0.008)
  (-0.174,0.775)rectangle++(0.008,0.008) (-0.305,0.881)rectangle++(0.008,0.008) (-0.458,0.713)rectangle++(0.008,0.008) (-0.307,0.849)rectangle++(0.008,0.008) (-0.157,0.728)rectangle++(0.008,0.008)
  (-0.336,0.696)rectangle++(0.008,0.008) (-0.265,0.843)rectangle++(0.008,0.008) (-0.426,0.71)rectangle++(0.008,0.008) (-0.452,0.592)rectangle++(0.008,0.008) (-0.305,0.808)rectangle++(0.008,0.008)
  (-0.147,0.595)rectangle++(0.008,0.008) (-0.174,0.894)rectangle++(0.008,0.008) (-0.211,0.91)rectangle++(0.008,0.008) (-0.121,0.939)rectangle++(0.008,0.008) (-0.193,0.925)rectangle++(0.008,0.008)
  (-0.135,0.592)rectangle++(0.008,0.008) (-0.132,0.914)rectangle++(0.008,0.008) (-0.143,0.591)rectangle++(0.008,0.008) (-0.28,0.638)rectangle++(0.008,0.008) (-0.329,0.887)rectangle++(0.008,0.008)
  (-0.327,0.897)rectangle++(0.008,0.008) (-0.415,0.815)rectangle++(0.008,0.008) (-0.118,0.939)rectangle++(0.008,0.008) (-0.393,0.806)rectangle++(0.008,0.008) (-0.107,0.809)rectangle++(0.008,0.008)
  (-0.155,0.723)rectangle++(0.008,0.008) (-0.115,0.554)rectangle++(0.008,0.008) (-0.254,0.631)rectangle++(0.008,0.008) (-0.041,0.859)rectangle++(0.008,0.008) (-0.166,0.898)rectangle++(0.008,0.008)
  (-0.071,0.907)rectangle++(0.008,0.008) (-0.112,0.931)rectangle++(0.008,0.008) (-0.108,0.805)rectangle++(0.008,0.008) (-0.12,0.927)rectangle++(0.008,0.008) (-0.389,0.8)rectangle++(0.008,0.008)
  (-0.467,0.646)rectangle++(0.008,0.008) (-0.309,0.894)rectangle++(0.008,0.008) (-0.143,0.585)rectangle++(0.008,0.008) (-0.174,0.893)rectangle++(0.008,0.008) (-0.032,0.869)rectangle++(0.008,0.008)
  (-0.151,0.73)rectangle++(0.008,0.008) (-0.302,0.875)rectangle++(0.008,0.008) (-0.299,0.908)rectangle++(0.008,0.008) (-0.253,0.915)rectangle++(0.008,0.008) (-0.446,0.747)rectangle++(0.008,0.008)
  (-0.303,0.861)rectangle++(0.008,0.008) (-0.449,0.705)rectangle++(0.008,0.008) (-0.304,0.846)rectangle++(0.008,0.008) (-0.445,0.698)rectangle++(0.008,0.008) (-0.302,0.844)rectangle++(0.008,0.008)
  (-0.443,0.697)rectangle++(0.008,0.008) (-0.004,0.833)rectangle++(0.008,0.008) (-0.255,0.743)rectangle++(0.008,0.008) (-0.119,0.557)rectangle++(0.008,0.008) (-0.257,0.631)rectangle++(0.008,0.008)
  (-0.013,0.824)rectangle++(0.008,0.008) (-0.303,0.791)rectangle++(0.008,0.008) (-0.425,0.673)rectangle++(0.008,0.008) (-0.438,0.575)rectangle++(0.008,0.008) (-0.238,0.884)rectangle++(0.008,0.008)
  (-0.428,0.738)rectangle++(0.008,0.008) (-0.463,0.604)rectangle++(0.008,0.008) (-0.431,0.53)rectangle++(0.008,0.008) (-0.39,0.508)rectangle++(0.008,0.008) (-0.284,0.78)rectangle++(0.008,0.008);
\fill[piece2]
  (-0.125,0.564)rectangle++(0.008,0.008) (-0.262,0.632)rectangle++(0.008,0.008) (-0.241,0.822)rectangle++(0.008,0.008) (-0.091,0.905)rectangle++(0.008,0.008) (-0.183,0.914)rectangle++(0.008,0.008)
  (-0.414,0.771)rectangle++(0.008,0.008) (-0.292,0.869)rectangle++(0.008,0.008) (-0.251,0.902)rectangle++(0.008,0.008) (-0.237,0.913)rectangle++(0.008,0.008) (-0.438,0.752)rectangle++(0.008,0.008)
  (-0.472,0.607)rectangle++(0.008,0.008) (-0.347,0.857)rectangle++(0.008,0.008) (-0.468,0.687)rectangle++(0.008,0.008) (-0.114,0.544)rectangle++(0.008,0.008) (-0.191,0.792)rectangle++(0.008,0.008)
  (-0.373,0.858)rectangle++(0.008,0.008) (-0.481,0.678)rectangle++(0.008,0.008) (-0.139,0.589)rectangle++(0.008,0.008) (-0.278,0.639)rectangle++(0.008,0.008) (-0.246,0.824)rectangle++(0.008,0.008)
  (-0.389,0.813)rectangle++(0.008,0.008) (-0.114,0.568)rectangle++(0.008,0.008) (-0.191,0.8)rectangle++(0.008,0.008) (-0.127,0.568)rectangle++(0.008,0.008) (-0.044,0.843)rectangle++(0.008,0.008)
  (-0.325,0.789)rectangle++(0.008,0.008) (-0.399,0.794)rectangle++(0.008,0.008) (-0.469,0.64)rectangle++(0.008,0.008) (-0.447,0.544)rectangle++(0.008,0.008) (-0.402,0.508)rectangle++(0.008,0.008)
  (-0.372,0.729)rectangle++(0.008,0.008) (-0.434,0.62)rectangle++(0.008,0.008) (-0.423,0.548)rectangle++(0.008,0.008) (-0.28,0.887)rectangle++(0.008,0.008) (-0.281,0.875)rectangle++(0.008,0.008)
  (-0.091,0.918)rectangle++(0.008,0.008) (-0.373,0.806)rectangle++(0.008,0.008) (-0.278,0.881)rectangle++(0.008,0.008) (-0.246,0.906)rectangle++(0.008,0.008) (-0.235,0.914)rectangle++(0.008,0.008)
  (-0.232,0.917)rectangle++(0.008,0.008) (-0.437,0.755)rectangle++(0.008,0.008) (-0.473,0.609)rectangle++(0.008,0.008) (-0.437,0.528)rectangle++(0.008,0.008) (-0.18,0.926)rectangle++(0.008,0.008)
  (-0.153,0.747)rectangle++(0.008,0.008) (-0.34,0.706)rectangle++(0.008,0.008) (-0.068,0.904)rectangle++(0.008,0.008) (-0.163,0.938)rectangle++(0.008,0.008) (-0.115,0.558)rectangle++(0.008,0.008)
  (-0.324,0.781)rectangle++(0.008,0.008) (-0.397,0.793)rectangle++(0.008,0.008) (-0.468,0.64)rectangle++(0.008,0.008) (-0.107,0.535)rectangle++(0.008,0.008) (-0.243,0.625)rectangle++(0.008,0.008)
  (-0.338,0.618)rectangle++(0.008,0.008) (-0.379,0.581)rectangle++(0.008,0.008) (-0.384,0.55)rectangle++(0.008,0.008) (-0.143,0.948)rectangle++(0.008,0.008) (-0.155,0.742)rectangle++(0.008,0.008)
  (-0.142,0.591)rectangle++(0.008,0.008) (-0.334,0.784)rectangle++(0.008,0.008) (-0.436,0.658)rectangle++(0.008,0.008) (-0.215,0.903)rectangle++(0.008,0.008) (-0.01,0.862)rectangle++(0.008,0.008)
  (-0.149,0.601)rectangle++(0.008,0.008) (-0.286,0.641)rectangle++(0.008,0.008) (-0.093,0.543)rectangle++(0.008,0.008) (-0.184,0.792)rectangle++(0.008,0.008) (-0.146,0.6)rectangle++(0.008,0.008)
  (-0.174,0.895)rectangle++(0.008,0.008) (-0.403,0.766)rectangle++(0.008,0.008) (-0.461,0.626)rectangle++(0.008,0.008) (-0.308,0.82)rectangle++(0.008,0.008) (-0.256,0.885)rectangle++(0.008,0.008)
  (-0.4,0.823)rectangle++(0.008,0.008) (-0.418,0.79)rectangle++(0.008,0.008) (-0.417,0.78)rectangle++(0.008,0.008) (0.003,0.811)rectangle++(0.008,0.008) (-0.292,0.791)rectangle++(0.008,0.008)
  (-0.419,0.677)rectangle++(0.008,0.008) (-0.4,0.759)rectangle++(0.008,0.008) (-0.287,0.864)rectangle++(0.008,0.008) (-0.128,0.94)rectangle++(0.008,0.008) (-0.398,0.803)rectangle++(0.008,0.008)
  (-0.415,0.786)rectangle++(0.008,0.008) (-0.339,0.882)rectangle++(0.008,0.008) (-0.267,0.906)rectangle++(0.008,0.008) (-0.449,0.738)rectangle++(0.008,0.008) (-0.12,0.937)rectangle++(0.008,0.008)
  (-0.393,0.805)rectangle++(0.008,0.008) (-0.471,0.647)rectangle++(0.008,0.008) (-0.45,0.547)rectangle++(0.008,0.008) (-0.304,0.793)rectangle++(0.008,0.008) (-0.133,0.936)rectangle++(0.008,0.008)
  (-0.399,0.799)rectangle++(0.008,0.008) (-0.391,0.845)rectangle++(0.008,0.008) (-0.156,0.86)rectangle++(0.008,0.008) (-0.377,0.833)rectangle++(0.008,0.008) (-0.135,0.573)rectangle++(0.008,0.008)
  (-0.27,0.633)rectangle++(0.008,0.008) (-0.365,0.773)rectangle++(0.008,0.008) (-0.275,0.869)rectangle++(0.008,0.008) (-0.049,0.871)rectangle++(0.008,0.008) (-0.056,0.889)rectangle++(0.008,0.008)
  (-0.228,0.853)rectangle++(0.008,0.008) (-0.412,0.728)rectangle++(0.008,0.008) (-0.291,0.854)rectangle++(0.008,0.008) (-0.442,0.706)rectangle++(0.008,0.008) (-0.409,0.761)rectangle++(0.008,0.008)
  (-0.338,0.879)rectangle++(0.008,0.008) (-0.152,0.608)rectangle++(0.008,0.008) (-0.115,0.562)rectangle++(0.008,0.008) (-0.257,0.635)rectangle++(0.008,0.008) (-0.35,0.877)rectangle++(0.008,0.008)
  (-0.477,0.695)rectangle++(0.008,0.008) (-0.149,0.938)rectangle++(0.008,0.008) (-0.157,0.611)rectangle++(0.008,0.008) (-0.34,0.785)rectangle++(0.008,0.008) (-0.439,0.657)rectangle++(0.008,0.008)
  (-0.272,0.761)rectangle++(0.008,0.008) (-0.4,0.67)rectangle++(0.008,0.008) (-0.287,0.835)rectangle++(0.008,0.008) (-0.097,0.919)rectangle++(0.008,0.008) (-0.376,0.804)rectangle++(0.008,0.008)
  (-0.274,0.91)rectangle++(0.008,0.008) (-0.261,0.855)rectangle++(0.008,0.008) (-0.24,0.898)rectangle++(0.008,0.008) (-0.399,0.828)rectangle++(0.008,0.008) (-0.287,0.888)rectangle++(0.008,0.008)
  (-0.276,0.761)rectangle++(0.008,0.008) (-0.385,0.796)rectangle++(0.008,0.008) (-0.464,0.646)rectangle++(0.008,0.008) (-0.085,0.522)rectangle++(0.008,0.008) (-0.181,0.785)rectangle++(0.008,0.008)
  (-0.367,0.713)rectangle++(0.008,0.008) (-0.426,0.614)rectangle++(0.008,0.008) (-0.154,0.916)rectangle++(0.008,0.008) (-0.175,0.94)rectangle++(0.008,0.008) (-0.393,0.846)rectangle++(0.008,0.008)
  (-0.485,0.666)rectangle++(0.008,0.008) (-0.316,0.833)rectangle++(0.008,0.008) (-0.404,0.804)rectangle++(0.008,0.008) (-0.475,0.643)rectangle++(0.008,0.008) (-0.313,0.825)rectangle++(0.008,0.008)
  (-0.15,0.605)rectangle++(0.008,0.008) (-0.338,0.785)rectangle++(0.008,0.008) (-0.318,0.818)rectangle++(0.008,0.008) (-0.441,0.68)rectangle++(0.008,0.008) (-0.343,0.867)rectangle++(0.008,0.008)
  (-0.092,0.539)rectangle++(0.008,0.008) (-0.238,0.633)rectangle++(0.008,0.008) (-0.338,0.623)rectangle++(0.008,0.008) (-0.338,0.896)rectangle++(0.008,0.008) (-0.381,0.875)rectangle++(0.008,0.008)
  (-0.423,0.802)rectangle++(0.008,0.008) (-0.483,0.635)rectangle++(0.008,0.008) (-0.315,0.823)rectangle++(0.008,0.008) (-0.442,0.683)rectangle++(0.008,0.008) (-0.15,0.603)rectangle++(0.008,0.008)
  (-0.032,0.897)rectangle++(0.008,0.008) (-0.033,0.905)rectangle++(0.008,0.008) (-0.164,0.914)rectangle++(0.008,0.008) (-0.208,0.917)rectangle++(0.008,0.008) (-0.277,0.764)rectangle++(0.008,0.008)
  (-0.239,0.716)rectangle++(0.008,0.008) (-0.371,0.793)rectangle++(0.008,0.008) (-0.457,0.649)rectangle++(0.008,0.008) (-0.444,0.553)rectangle++(0.008,0.008) (-0.249,0.861)rectangle++(0.008,0.008)
  (-0.425,0.724)rectangle++(0.008,0.008) (-0.02,0.835)rectangle++(0.008,0.008) (-0.311,0.794)rectangle++(0.008,0.008) (0.005,0.825)rectangle++(0.008,0.008) (-0.012,0.836)rectangle++(0.008,0.008)
  (-0.157,0.89)rectangle++(0.008,0.008) (-0.07,0.905)rectangle++(0.008,0.008) (-0.339,0.893)rectangle++(0.008,0.008) (-0.417,0.812)rectangle++(0.008,0.008) (-0.484,0.642)rectangle++(0.008,0.008)
  (-0.407,0.742)rectangle++(0.008,0.008) (-0.211,0.915)rectangle++(0.008,0.008) (-0.277,0.763)rectangle++(0.008,0.008) (-0.403,0.669)rectangle++(0.008,0.008) (-0.076,0.52)rectangle++(0.008,0.008)
  (-0.178,0.784)rectangle++(0.008,0.008) (-0.213,0.873)rectangle++(0.008,0.008) (-0.013,0.861)rectangle++(0.008,0.008) (-0.029,0.899)rectangle++(0.008,0.008) (-0.358,0.86)rectangle++(0.008,0.008)
  (-0.474,0.684)rectangle++(0.008,0.008) (-0.39,0.831)rectangle++(0.008,0.008) (-0.022,0.879)rectangle++(0.008,0.008) (-0.281,0.823)rectangle++(0.008,0.008) (-0.396,0.807)rectangle++(0.008,0.008)
  (-0.391,0.846)rectangle++(0.008,0.008) (-0.209,0.93)rectangle++(0.008,0.008) (-0.156,0.74)rectangle++(0.008,0.008) (-0.339,0.702)rectangle++(0.008,0.008) (-0.38,0.841)rectangle++(0.008,0.008)
  (-0.478,0.668)rectangle++(0.008,0.008) (-0.154,0.648)rectangle++(0.008,0.008) (-0.032,0.883)rectangle++(0.008,0.008) (-0.274,0.762)rectangle++(0.008,0.008) (-0.123,0.56)rectangle++(0.008,0.008)
  (-0.26,0.631)rectangle++(0.008,0.008) (-0.348,0.615)rectangle++(0.008,0.008) (-0.163,0.946)rectangle++(0.008,0.008) (-0.004,0.857)rectangle++(0.008,0.008) (-0.311,0.81)rectangle++(0.008,0.008)
  (-0.435,0.678)rectangle++(0.008,0.008) (-0.019,0.835)rectangle++(0.008,0.008) (-0.31,0.794)rectangle++(0.008,0.008) (-0.429,0.671)rectangle++(0.008,0.008) (-0.439,0.573)rectangle++(0.008,0.008)
  (-0.409,0.524)rectangle++(0.008,0.008) (-0.29,0.786)rectangle++(0.008,0.008) (-0.126,0.565)rectangle++(0.008,0.008) (-0.327,0.781)rectangle++(0.008,0.008) (-0.15,0.734)rectangle++(0.008,0.008)
  (-0.141,0.59)rectangle++(0.008,0.008) (-0.334,0.784)rectangle++(0.008,0.008) (-0.026,0.872)rectangle++(0.008,0.008) (-0.154,0.732)rectangle++(0.008,0.008) (-0.205,0.855)rectangle++(0.008,0.008)
  (-0.403,0.737)rectangle++(0.008,0.008) (-0.153,0.866)rectangle++(0.008,0.008) (-0.377,0.835)rectangle++(0.008,0.008) (-0.28,0.89)rectangle++(0.008,0.008) (-0.118,0.571)rectangle++(0.008,0.008)
  (-0.274,0.763)rectangle++(0.008,0.008) (-0.156,0.611)rectangle++(0.008,0.008) (-0.117,0.56)rectangle++(0.008,0.008) (-0.257,0.634)rectangle++(0.008,0.008) (-0.371,0.869)rectangle++(0.008,0.008)
  (-0.42,0.803)rectangle++(0.008,0.008) (0.001,0.796)rectangle++(0.008,0.008) (-0.153,0.923)rectangle++(0.008,0.008) (-0.204,0.92)rectangle++(0.008,0.008) (-0.383,0.867)rectangle++(0.008,0.008)
  (-0.488,0.679)rectangle++(0.008,0.008) (-0.146,0.912)rectangle++(0.008,0.008) (-0.343,0.894)rectangle++(0.008,0.008) (-0.061,0.824)rectangle++(0.008,0.008) (-0.173,0.887)rectangle++(0.008,0.008)
  (-0.211,0.908)rectangle++(0.008,0.008) (-0.424,0.759)rectangle++(0.008,0.008) (0,0.83)rectangle++(0.008,0.008) (-0.153,0.928)rectangle++(0.008,0.008) (-0.011,0.785)rectangle++(0.008,0.008)
  (-0.289,0.774)rectangle++(0.008,0.008) (-0.412,0.67)rectangle++(0.008,0.008) (-0.275,0.761)rectangle++(0.008,0.008) (-0.401,0.669)rectangle++(0.008,0.008) (-0.327,0.827)rectangle++(0.008,0.008)
  (-0.449,0.681)rectangle++(0.008,0.008) (-0.109,0.562)rectangle++(0.008,0.008) (-0.254,0.637)rectangle++(0.008,0.008) (-0.347,0.62)rectangle++(0.008,0.008) (-0.102,0.532)rectangle++(0.008,0.008)
  (-0.24,0.626)rectangle++(0.008,0.008) (-0.113,0.567)rectangle++(0.008,0.008) (-0.191,0.8)rectangle++(0.008,0.008) (-0.376,0.717)rectangle++(0.008,0.008) (-0.398,0.773)rectangle++(0.008,0.008)
  (-0.286,0.869)rectangle++(0.008,0.008) (-0.194,0.933)rectangle++(0.008,0.008) (-0.425,0.777)rectangle++(0.008,0.008) (-0.296,0.871)rectangle++(0.008,0.008) (-0.405,0.814)rectangle++(0.008,0.008)
  (-0.48,0.647)rectangle++(0.008,0.008) (-0.221,0.902)rectangle++(0.008,0.008) (-0.427,0.752)rectangle++(0.008,0.008) (-0.467,0.611)rectangle++(0.008,0.008) (-0.135,0.579)rectangle++(0.008,0.008)
  (-0.272,0.635)rectangle++(0.008,0.008) (-0.355,0.612)rectangle++(0.008,0.008) (-0.378,0.756)rectangle++(0.008,0.008) (-0.123,0.558)rectangle++(0.008,0.008) (-0.171,0.889)rectangle++(0.008,0.008)
  (-0.032,0.814)rectangle++(0.008,0.008) (-0.309,0.78)rectangle++(0.008,0.008) (-0.423,0.665)rectangle++(0.008,0.008) (-0.435,0.572)rectangle++(0.008,0.008) (-0.407,0.526)rectangle++(0.008,0.008)
  (-0.036,0.815)rectangle++(0.008,0.008) (-0.131,0.588)rectangle++(0.008,0.008) (-0.273,0.641)rectangle++(0.008,0.008) (-0.33,0.87)rectangle++(0.008,0.008) (-0.331,0.898)rectangle++(0.008,0.008)
  (-0.117,0.573)rectangle++(0.008,0.008) (-0.06,0.884)rectangle++(0.008,0.008) (-0.362,0.852)rectangle++(0.008,0.008) (-0.137,0.577)rectangle++(0.008,0.008) (-0.147,0.927)rectangle++(0.008,0.008)
  (-0.153,0.917)rectangle++(0.008,0.008) (-0.401,0.784)rectangle++(0.008,0.008) (-0.467,0.634)rectangle++(0.008,0.008) (-0.31,0.823)rectangle++(0.008,0.008) (-0.283,0.916)rectangle++(0.008,0.008)
  (-0.46,0.737)rectangle++(0.008,0.008) (-0.417,0.764)rectangle++(0.008,0.008) (-0.467,0.62)rectangle++(0.008,0.008) (-0.037,0.897)rectangle++(0.008,0.008) (-0.28,0.769)rectangle++(0.008,0.008)
  (-0.387,0.797)rectangle++(0.008,0.008) (-0.026,0.878)rectangle++(0.008,0.008) (-0.146,0.951)rectangle++(0.008,0.008) (-0.38,0.873)rectangle++(0.008,0.008) (-0.019,0.835)rectangle++(0.008,0.008)
  (-0.31,0.794)rectangle++(0.008,0.008) (-0.004,0.854)rectangle++(0.008,0.008) (-0.154,0.896)rectangle++(0.008,0.008) (-0.115,0.566)rectangle++(0.008,0.008) (-0.091,0.524)rectangle++(0.008,0.008)
  (-0.166,0.884)rectangle++(0.008,0.008) (-0.395,0.764)rectangle++(0.008,0.008) (-0.117,0.8)rectangle++(0.008,0.008) (-0.192,0.879)rectangle++(0.008,0.008) (-0.379,0.865)rectangle++(0.008,0.008)
  (-0.334,0.894)rectangle++(0.008,0.008) (-0.416,0.813)rectangle++(0.008,0.008) (-0.178,0.943)rectangle++(0.008,0.008) (-0.28,0.825)rectangle++(0.008,0.008) (-0.426,0.696)rectangle++(0.008,0.008)
  (-0.296,0.844)rectangle++(0.008,0.008) (-0.44,0.699)rectangle++(0.008,0.008) (-0.454,0.582)rectangle++(0.008,0.008) (-0.305,0.805)rectangle++(0.008,0.008) (-0.431,0.678)rectangle++(0.008,0.008)
  (-0.402,0.757)rectangle++(0.008,0.008) (-0.156,0.855)rectangle++(0.008,0.008) (-0.155,0.861)rectangle++(0.008,0.008) (-0.331,0.834)rectangle++(0.008,0.008) (-0.104,0.917)rectangle++(0.008,0.008)
  (-0.128,0.937)rectangle++(0.008,0.008) (-0.335,0.849)rectangle++(0.008,0.008) (-0.23,0.881)rectangle++(0.008,0.008) (-0.382,0.862)rectangle++(0.008,0.008) (-0.486,0.677)rectangle++(0.008,0.008)
  (-0.31,0.896)rectangle++(0.008,0.008) (0.005,0.82)rectangle++(0.008,0.008) (-0.146,0.598)rectangle++(0.008,0.008) (-0.284,0.64)rectangle++(0.008,0.008) (-0.362,0.611)rectangle++(0.008,0.008)
  (-0.332,0.858)rectangle++(0.008,0.008) (-0.462,0.693)rectangle++(0.008,0.008) (-0.132,0.777)rectangle++(0.008,0.008) (-0.368,0.86)rectangle++(0.008,0.008) (-0.333,0.893)rectangle++(0.008,0.008)
  (-0.474,0.709)rectangle++(0.008,0.008) (-0.415,0.756)rectangle++(0.008,0.008) (-0.464,0.617)rectangle++(0.008,0.008) (-0.4,0.74)rectangle++(0.008,0.008) (-0.287,0.858)rectangle++(0.008,0.008)
  (-0.227,0.893)rectangle++(0.008,0.008) (-0.396,0.829)rectangle++(0.008,0.008) (-0.209,0.927)rectangle++(0.008,0.008) (-0.098,0.928)rectangle++(0.008,0.008) (-0.167,0.941)rectangle++(0.008,0.008)
  (-0.416,0.79)rectangle++(0.008,0.008) (-0.293,0.875)rectangle++(0.008,0.008) (-0.223,0.881)rectangle++(0.008,0.008) (-0.356,0.866)rectangle++(0.008,0.008) (-0.476,0.688)rectangle++(0.008,0.008)
  (-0.265,0.744)rectangle++(0.008,0.008) (-0.261,0.906)rectangle++(0.008,0.008) (-0.317,0.899)rectangle++(0.008,0.008) (-0.152,0.741)rectangle++(0.008,0.008) (-0.009,0.814)rectangle++(0.008,0.008)
  (-0.127,0.573)rectangle++(0.008,0.008) (-0.104,0.853)rectangle++(0.008,0.008) (-0.012,0.848)rectangle++(0.008,0.008) (-0.148,0.6)rectangle++(0.008,0.008) (-0.285,0.641)rectangle++(0.008,0.008)
  (-0.15,0.699)rectangle++(0.008,0.008) (-0.322,0.685)rectangle++(0.008,0.008);
\fill[piece3]
  (0.234,0.72)rectangle++(0.008,0.008) (-0.005,0.801)rectangle++(0.008,0.008) (0.146,0.879)rectangle++(0.008,0.008) (0.036,0.873)rectangle++(0.008,0.008) (0.327,0.803)rectangle++(0.008,0.008)
  (0.13,0.592)rectangle++(0.008,0.008) (0.349,0.845)rectangle++(0.008,0.008) (0.098,0.922)rectangle++(0.008,0.008) (0.373,0.803)rectangle++(0.008,0.008) (0.017,0.877)rectangle++(0.008,0.008)
  (0.347,0.856)rectangle++(0.008,0.008) (0.266,0.758)rectangle++(0.008,0.008) (0.392,0.668)rectangle++(0.008,0.008) (0.084,0.904)rectangle++(0.008,0.008) (0.176,0.913)rectangle++(0.008,0.008)
  (0.101,0.809)rectangle++(0.008,0.008) (0.119,0.571)rectangle++(0.008,0.008) (0.096,0.852)rectangle++(0.008,0.008) (0.358,0.84)rectangle++(0.008,0.008) (0.01,0.887)rectangle++(0.008,0.008)
  (0.037,0.469)rectangle++(0.008,0.008) (0.333,0.841)rectangle++(0.008,0.008) (0.327,0.86)rectangle++(0.008,0.008) (0.257,0.899)rectangle++(0.008,0.008) (0.438,0.735)rectangle++(0.008,0.008)
  (0.462,0.596)rectangle++(0.008,0.008) (0.245,0.897)rectangle++(0.008,0.008) (0.431,0.739)rectangle++(0.008,0.008) (0.46,0.6)rectangle++(0.008,0.008) (0.042,0.881)rectangle++(0.008,0.008)
  (0.333,0.804)rectangle++(0.008,0.008) (0.019,0.893)rectangle++(0.008,0.008) (0.349,0.859)rectangle++(0.008,0.008) (0.466,0.684)rectangle++(0.008,0.008) (0.158,0.947)rectangle++(0.008,0.008)
  (0.373,0.872)rectangle++(0.008,0.008) (0.099,0.929)rectangle++(0.008,0.008) (0.376,0.806)rectangle++(0.008,0.008) (0.123,0.567)rectangle++(0.008,0.008) (0.362,0.881)rectangle++(0.008,0.008)
  (0.253,0.855)rectangle++(0.008,0.008) (0.42,0.717)rectangle++(0.008,0.008) (0.447,0.594)rectangle++(0.008,0.008) (0.298,0.809)rectangle++(0.008,0.008) (0.424,0.68)rectangle++(0.008,0.008)
  (-0.003,0.833)rectangle++(0.008,0.008) (0.143,0.721)rectangle++(0.008,0.008) (0.322,0.695)rectangle++(0.008,0.008) (0.395,0.619)rectangle++(0.008,0.008) (0.124,0.586)rectangle++(0.008,0.008)
  (0.265,0.64)rectangle++(0.008,0.008) (0.237,0.825)rectangle++(0.008,0.008) (0.056,0.823)rectangle++(0.008,0.008) (0.282,0.888)rectangle++(0.008,0.008) (0.445,0.722)rectangle++(0.008,0.008)
  (0.297,0.852)rectangle++(0.008,0.008) (0.439,0.7)rectangle++(0.008,0.008) (0.209,0.909)rectangle++(0.008,0.008) (0.159,0.949)rectangle++(0.008,0.008) (0.169,0.944)rectangle++(0.008,0.008)
  (0.414,0.788)rectangle++(0.008,0.008) (0.409,0.78)rectangle++(0.008,0.008) (0.465,0.627)rectangle++(0.008,0.008) (0.304,0.82)rectangle++(0.008,0.008) (0.249,0.885)rectangle++(0.008,0.008)
  (-0.011,0.817)rectangle++(0.008,0.008) (0.273,0.875)rectangle++(0.008,0.008) (0.019,0.901)rectangle++(0.008,0.008) (0.35,0.861)rectangle++(0.008,0.008) (0.196,0.932)rectangle++(0.008,0.008)
  (0.213,0.859)rectangle++(0.008,0.008) (0.146,0.939)rectangle++(0.008,0.008) (0.115,0.788)rectangle++(0.008,0.008) (0.211,0.856)rectangle++(0.008,0.008) (0.218,0.897)rectangle++(0.008,0.008)
  (0.42,0.749)rectangle++(0.008,0.008) (0.458,0.609)rectangle++(0.008,0.008) (0.426,0.531)rectangle++(0.008,0.008) (0.384,0.507)rectangle++(0.008,0.008) (0.361,0.73)rectangle++(0.008,0.008)
  (0.425,0.621)rectangle++(0.008,0.008) (0.415,0.548)rectangle++(0.008,0.008) (0.385,0.519)rectangle++(0.008,0.008) (0.143,0.876)rectangle++(0.008,0.008) (0.378,0.766)rectangle++(0.008,0.008)
  (0.327,0.88)rectangle++(0.008,0.008) (0.406,0.81)rectangle++(0.008,0.008) (0.411,0.785)rectangle++(0.008,0.008) (0.115,0.92)rectangle++(0.008,0.008) (0.38,0.796)rectangle++(0.008,0.008)
  (0.141,0.756)rectangle++(0.008,0.008) (0.334,0.712)rectangle++(0.008,0.008) (0.373,0.842)rectangle++(0.008,0.008) (0.274,0.828)rectangle++(0.008,0.008) (0.239,0.888)rectangle++(0.008,0.008)
  (0.228,0.908)rectangle++(0.008,0.008) (0.104,0.567)rectangle++(0.008,0.008) (0.222,0.896)rectangle++(0.008,0.008) (0.421,0.747)rectangle++(0.008,0.008) (0.15,0.714)rectangle++(0.008,0.008)
  (0.323,0.689)rectangle++(0.008,0.008) (0.377,0.774)rectangle++(0.008,0.008) (0.321,0.893)rectangle++(0.008,0.008) (0.255,0.91)rectangle++(0.008,0.008) (0.441,0.741)rectangle++(0.008,0.008)
  (0.026,0.883)rectangle++(0.008,0.008) (0.156,0.906)rectangle++(0.008,0.008) (0.378,0.841)rectangle++(0.008,0.008) (0.013,0.879)rectangle++(0.008,0.008) (0.319,0.813)rectangle++(0.008,0.008)
  (0.255,0.883)rectangle++(0.008,0.008) (0.431,0.729)rectangle++(0.008,0.008) (0.457,0.596)rectangle++(0.008,0.008) (0.085,0.906)rectangle++(0.008,0.008) (0.362,0.8)rectangle++(0.008,0.008)
  (0.024,0.856)rectangle++(0.008,0.008) (0.156,0.897)rectangle++(0.008,0.008) (0.391,0.771)rectangle++(0.008,0.008) (0.118,0.561)rectangle++(0.008,0.008) (0.318,0.882)rectangle++(0.008,0.008)
  (0.254,0.906)rectangle++(0.008,0.008) (0.014,0.891)rectangle++(0.008,0.008) (0.324,0.819)rectangle++(0.008,0.008) (0.44,0.675)rectangle++(0.008,0.008) (0.441,0.568)rectangle++(0.008,0.008)
  (0.381,0.733)rectangle++(0.008,0.008) (0.2,0.914)rectangle++(0.008,0.008) (0.215,0.917)rectangle++(0.008,0.008) (0.425,0.759)rectangle++(0.008,0.008) (0.29,0.865)rectangle++(0.008,0.008)
  (0.441,0.708)rectangle++(0.008,0.008) (0.453,0.583)rectangle++(0.008,0.008) (0.415,0.521)rectangle++(0.008,0.008) (0.375,0.506)rectangle++(0.008,0.008) (0.351,0.514)rectangle++(0.008,0.008)
  (0.094,0.895)rectangle++(0.008,0.008) (0.362,0.792)rectangle++(0.008,0.008) (0.343,0.89)rectangle++(0.008,0.008) (0.317,0.896)rectangle++(0.008,0.008) (0.191,0.937)rectangle++(0.008,0.008)
  (0.174,0.943)rectangle++(0.008,0.008) (0.337,0.895)rectangle++(0.008,0.008) (0.473,0.705)rectangle++(0.008,0.008) (0.109,0.547)rectangle++(0.008,0.008) (0.252,0.748)rectangle++(0.008,0.008)
  (0.381,0.668)rectangle++(0.008,0.008) (0.374,0.835)rectangle++(0.008,0.008) (0.468,0.665)rectangle++(0.008,0.008) (0.271,0.769)rectangle++(0.008,0.008) (0.398,0.671)rectangle++(0.008,0.008)
  (0.24,0.896)rectangle++(0.008,0.008) (0.429,0.74)rectangle++(0.008,0.008) (0.405,0.768)rectangle++(0.008,0.008) (0.405,0.777)rectangle++(0.008,0.008) (0.145,0.844)rectangle++(0.008,0.008)
  (0.294,0.891)rectangle++(0.008,0.008) (0.021,0.906)rectangle++(0.008,0.008) (0.154,0.914)rectangle++(0.008,0.008) (0.396,0.779)rectangle++(0.008,0.008) (0.458,0.631)rectangle++(0.008,0.008)
  (0.394,0.743)rectangle++(0.008,0.008) (0.28,0.859)rectangle++(0.008,0.008) (0.434,0.709)rectangle++(0.008,0.008) (0.131,0.59)rectangle++(0.008,0.008) (0.105,0.56)rectangle++(0.008,0.008)
  (0.095,0.849)rectangle++(0.008,0.008) (0.346,0.77)rectangle++(0.008,0.008) (0.432,0.645)rectangle++(0.008,0.008) (0.273,0.827)rectangle++(0.008,0.008) (0.419,0.697)rectangle++(0.008,0.008)
  (0.145,0.741)rectangle++(0.008,0.008) (0.35,0.811)rectangle++(0.008,0.008) (0.38,0.849)rectangle++(0.008,0.008) (0.324,0.837)rectangle++(0.008,0.008) (0.446,0.683)rectangle++(0.008,0.008)
  (0.047,0.482)rectangle++(0.008,0.008) (0.335,0.841)rectangle++(0.008,0.008) (0.453,0.681)rectangle++(0.008,0.008) (0.449,0.566)rectangle++(0.008,0.008) (0.383,0.732)rectangle++(0.008,0.008)
  (0.378,0.84)rectangle++(0.008,0.008) (0.472,0.666)rectangle++(0.008,0.008) (0.306,0.833)rectangle++(0.008,0.008) (0.437,0.688)rectangle++(0.008,0.008) (-0.004,0.833)rectangle++(0.008,0.008)
  (0.273,0.889)rectangle++(0.008,0.008) (0.442,0.725)rectangle++(0.008,0.008) (0.46,0.59)rectangle++(0.008,0.008) (0.327,0.89)rectangle++(0.008,0.008) (0.466,0.707)rectangle++(0.008,0.008)
  (0.304,0.847)rectangle++(0.008,0.008) (0.091,0.917)rectangle++(0.008,0.008) (0.05,0.902)rectangle++(0.008,0.008) (0.344,0.811)rectangle++(0.008,0.008) (0.266,0.764)rectangle++(0.008,0.008)
  (0.394,0.671)rectangle++(0.008,0.008) (0.419,0.582)rectangle++(0.008,0.008) (0.333,0.854)rectangle++(0.008,0.008) (0.052,0.914)rectangle++(0.008,0.008) (0.319,0.873)rectangle++(0.008,0.008)
  (0.054,0.867)rectangle++(0.008,0.008) (0.109,0.561)rectangle++(0.008,0.008) (0.037,0.842)rectangle++(0.008,0.008) (0.357,0.871)rectangle++(0.008,0.008) (0.129,0.89)rectangle++(0.008,0.008)
  (0.191,0.909)rectangle++(0.008,0.008) (0.211,0.915)rectangle++(0.008,0.008) (0.423,0.759)rectangle++(0.008,0.008) (0.407,0.773)rectangle++(0.008,0.008) (0.02,0.905)rectangle++(0.008,0.008)
  (0.275,0.91)rectangle++(0.008,0.008) (0.273,0.78)rectangle++(0.008,0.008) (0.162,0.941)rectangle++(0.008,0.008) (0.41,0.789)rectangle++(0.008,0.008) (0.468,0.631)rectangle++(0.008,0.008)
  (0.148,0.847)rectangle++(0.008,0.008) (0.37,0.751)rectangle++(0.008,0.008) (0.436,0.628)rectangle++(0.008,0.008) (0.07,0.904)rectangle++(0.008,0.008) (0.012,0.835)rectangle++(0.008,0.008)
  (0.142,0.705)rectangle++(0.008,0.008) (0.131,0.588)rectangle++(0.008,0.008) (0.272,0.802)rectangle++(0.008,0.008) (0.373,0.853)rectangle++(0.008,0.008) (0.254,0.905)rectangle++(0.008,0.008)
  (0.439,0.739)rectangle++(0.008,0.008) (0.336,0.876)rectangle++(0.008,0.008) (0.275,0.823)rectangle++(0.008,0.008) (0.048,0.483)rectangle++(0.008,0.008) (0.164,0.772)rectangle++(0.008,0.008)
  (0.079,0.889)rectangle++(0.008,0.008) (0,0.847)rectangle++(0.008,0.008) (0.113,0.575)rectangle++(0.008,0.008) (0.084,0.526)rectangle++(0.008,0.008) (0.307,0.779)rectangle++(0.008,0.008)
  (0.255,0.747)rectangle++(0.008,0.008) (0.348,0.867)rectangle++(0.008,0.008) (0.468,0.688)rectangle++(0.008,0.008) (0.459,0.564)rectangle++(0.008,0.008) (0.384,0.73)rectangle++(0.008,0.008)
  (0.435,0.613)rectangle++(0.008,0.008) (0.294,0.815)rectangle++(0.008,0.008) (0.425,0.684)rectangle++(0.008,0.008) (0.171,0.935)rectangle++(0.008,0.008) (0.205,0.924)rectangle++(0.008,0.008)
  (0.423,0.766)rectangle++(0.008,0.008) (0.466,0.615)rectangle++(0.008,0.008) (0.117,0.791)rectangle++(0.008,0.008) (0.187,0.875)rectangle++(0.008,0.008) (0.042,0.908)rectangle++(0.008,0.008)
  (0.161,0.915)rectangle++(0.008,0.008) (0.202,0.917)rectangle++(0.008,0.008) (0.302,0.901)rectangle++(0.008,0.008) (0.459,0.72)rectangle++(0.008,0.008) (0.13,0.94)rectangle++(0.008,0.008)
  (0.394,0.8)rectangle++(0.008,0.008) (0.465,0.641)rectangle++(0.008,0.008) (0.441,0.543)rectangle++(0.008,0.008) (0.395,0.507)rectangle++(0.008,0.008) (0.361,0.507)rectangle++(0.008,0.008)
  (0.345,0.519)rectangle++(0.008,0.008) (0.079,0.513)rectangle++(0.008,0.008) (0.219,0.622)rectangle++(0.008,0.008) (0.347,0.777)rectangle++(0.008,0.008) (0.435,0.648)rectangle++(0.008,0.008)
  (0.054,0.867)rectangle++(0.008,0.008) (0.127,0.583)rectangle++(0.008,0.008) (0.144,0.851)rectangle++(0.008,0.008) (0.369,0.754)rectangle++(0.008,0.008) (0.271,0.863)rectangle++(0.008,0.008)
  (0.431,0.714)rectangle++(0.008,0.008) (0.097,0.925)rectangle++(0.008,0.008) (0.374,0.804)rectangle++(0.008,0.008) (0.372,0.871)rectangle++(0.008,0.008) (0.48,0.682)rectangle++(0.008,0.008)
  (0.105,0.542)rectangle++(0.008,0.008) (0.241,0.627)rectangle++(0.008,0.008) (0.333,0.617)rectangle++(0.008,0.008) (0.227,0.88)rectangle++(0.008,0.008) (0.417,0.737)rectangle++(0.008,0.008)
  (0.111,0.799)rectangle++(0.008,0.008) (0.142,0.724)rectangle++(0.008,0.008) (0.323,0.696)rectangle++(0.008,0.008) (0.324,0.835)rectangle++(0.008,0.008) (0.446,0.682)rectangle++(0.008,0.008)
  (0.321,0.834)rectangle++(0.008,0.008) (0.444,0.683)rectangle++(0.008,0.008) (0.446,0.57)rectangle++(0.008,0.008) (0.407,0.518)rectangle++(0.008,0.008) (0.08,0.512)rectangle++(0.008,0.008)
  (0.174,0.781)rectangle++(0.008,0.008) (0.358,0.711)rectangle++(0.008,0.008) (0.387,0.773)rectangle++(0.008,0.008) (0.132,0.939)rectangle++(0.008,0.008) (0.154,0.951)rectangle++(0.008,0.008)
  (0.263,0.912)rectangle++(0.008,0.008) (0.012,0.858)rectangle++(0.008,0.008) (0.148,0.932)rectangle++(0.008,0.008) (0.295,0.903)rectangle++(0.008,0.008) (0.404,0.82)rectangle++(0.008,0.008)
  (0.131,0.943)rectangle++(0.008,0.008) (0.165,0.944)rectangle++(0.008,0.008) (0.412,0.789)rectangle++(0.008,0.008) (0.469,0.63)rectangle++(0.008,0.008) (0.272,0.895)rectangle++(0.008,0.008)
  (0.239,0.91)rectangle++(0.008,0.008) (0.394,0.829)rectangle++(0.008,0.008) (0.12,0.927)rectangle++(0.008,0.008) (0.188,0.921)rectangle++(0.008,0.008) (0.137,0.592)rectangle++(0.008,0.008)
  (0.273,0.638)rectangle++(0.008,0.008) (0.352,0.61)rectangle++(0.008,0.008) (0.378,0.57)rectangle++(0.008,0.008) (0.274,0.801)rectangle++(0.008,0.008) (0.185,0.924)rectangle++(0.008,0.008)
  (0.414,0.772)rectangle++(0.008,0.008) (0.407,0.777)rectangle++(0.008,0.008) (0.141,0.949)rectangle++(0.008,0.008) (0.195,0.929)rectangle++(0.008,0.008) (0.213,0.922)rectangle++(0.008,0.008)
  (0.304,0.902)rectangle++(0.008,0.008) (0.46,0.72)rectangle++(0.008,0.008) (0.467,0.581)rectangle++(0.008,0.008) (0.42,0.515)rectangle++(0.008,0.008) (0.289,0.783)rectangle++(0.008,0.008)
  (0.411,0.671)rectangle++(0.008,0.008) (0.285,0.835)rectangle++(0.008,0.008) (0.243,0.89)rectangle++(0.008,0.008) (0.231,0.898)rectangle++(0.008,0.008) (0.13,0.594)rectangle++(0.008,0.008)
  (0.292,0.855)rectangle++(0.008,0.008) (0.142,0.718)rectangle++(0.008,0.008) (0.011,0.86)rectangle++(0.008,0.008) (0.311,0.805)rectangle++(0.008,0.008) (0.097,0.912)rectangle++(0.008,0.008)
  (0.369,0.798)rectangle++(0.008,0.008) (0.453,0.649)rectangle++(0.008,0.008) (0.143,0.934)rectangle++(0.008,0.008) (0.167,0.942)rectangle++(0.008,0.008) (0.203,0.926)rectangle++(0.008,0.008)
  (-0.001,0.862)rectangle++(0.008,0.008) (0.327,0.891)rectangle++(0.008,0.008) (0.32,0.897)rectangle++(0.008,0.008) (0.043,0.476)rectangle++(0.008,0.008) (0.189,0.619)rectangle++(0.008,0.008)
  (0.341,0.781)rectangle++(0.008,0.008) (0.364,0.889)rectangle++(0.008,0.008) (0.146,0.717)rectangle++(0.008,0.008) (0.322,0.692)rectangle++(0.008,0.008) (0.256,0.842)rectangle++(0.008,0.008)
  (0.233,0.893)rectangle++(0.008,0.008) (0.148,0.735)rectangle++(0.008,0.008) (0.329,0.699)rectangle++(0.008,0.008) (0.106,0.547)rectangle++(0.008,0.008) (0.314,0.78)rectangle++(0.008,0.008)
  (0.421,0.661)rectangle++(0.008,0.008) (0.333,0.865)rectangle++(0.008,0.008) (0.405,0.806)rectangle++(0.008,0.008) (0.305,0.892)rectangle++(0.008,0.008) (0.457,0.715)rectangle++(0.008,0.008)
  (0.11,0.549)rectangle++(0.008,0.008) (0.081,0.521)rectangle++(0.008,0.008) (0.174,0.785)rectangle++(0.008,0.008) (0.245,0.89)rectangle++(0.008,0.008) (0.145,0.747)rectangle++(0.008,0.008)
  (0.332,0.706)rectangle++(0.008,0.008) (0.403,0.62)rectangle++(0.008,0.008) (0.128,0.58)rectangle++(0.008,0.008) (0.19,0.804)rectangle++(0.008,0.008) (0.373,0.716)rectangle++(0.008,0.008)
  (0.426,0.61)rectangle++(0.008,0.008) (0.385,0.744)rectangle++(0.008,0.008) (0.277,0.859)rectangle++(0.008,0.008) (0.121,0.939)rectangle++(0.008,0.008) (0.367,0.882)rectangle++(0.008,0.008)
  (0.482,0.689)rectangle++(0.008,0.008) (0.342,0.869)rectangle++(0.008,0.008) (0.132,0.58)rectangle++(0.008,0.008) (0.266,0.634)rectangle++(0.008,0.008) (0.311,0.861)rectangle++(0.008,0.008)
  (0.191,0.932)rectangle++(0.008,0.008) (0.419,0.774)rectangle++(0.008,0.008) (0.014,0.834)rectangle++(0.008,0.008) (0.303,0.793)rectangle++(0.008,0.008) (0.139,0.6)rectangle++(0.008,0.008)
  (0.154,0.93)rectangle++(0.008,0.008) (0.273,0.878)rectangle++(0.008,0.008) (0.396,0.818)rectangle++(0.008,0.008) (0.41,0.788)rectangle++(0.008,0.008) (0.145,0.739)rectangle++(0.008,0.008)
  (0.145,0.689)rectangle++(0.008,0.008) (0.257,0.75)rectangle++(0.008,0.008) (0.385,0.668)rectangle++(0.008,0.008) (0.414,0.584)rectangle++(0.008,0.008) (0.397,0.536)rectangle++(0.008,0.008)
  (0.372,0.52)rectangle++(0.008,0.008) (0.087,0.545)rectangle++(0.008,0.008) (0.146,0.859)rectangle++(0.008,0.008) (0.373,0.757)rectangle++(0.008,0.008) (0.113,0.911)rectangle++(0.008,0.008)
  (0.119,0.931)rectangle++(0.008,0.008) (0.141,0.708)rectangle++(0.008,0.008) (0.317,0.69)rectangle++(0.008,0.008) (0.376,0.775)rectangle++(0.008,0.008) (0.02,0.877)rectangle++(0.008,0.008)
  (0.322,0.81)rectangle++(0.008,0.008) (0.256,0.882)rectangle++(0.008,0.008) (0.393,0.821)rectangle++(0.008,0.008) (0.279,0.886)rectangle++(0.008,0.008) (0.247,0.902)rectangle++(0.008,0.008)
  (0.434,0.741)rectangle++(0.008,0.008) (0.462,0.6)rectangle++(0.008,0.008) (0.339,0.856)rectangle++(0.008,0.008) (0.405,0.804)rectangle++(0.008,0.008) (0.132,0.942)rectangle++(0.008,0.008)
  (0.371,0.873)rectangle++(0.008,0.008) (0.007,0.873)rectangle++(0.008,0.008) (0.15,0.903)rectangle++(0.008,0.008) (0.198,0.913)rectangle++(0.008,0.008) (0.019,0.778)rectangle++(0.008,0.008)
  (0.154,0.871)rectangle++(0.008,0.008) (0.239,0.895)rectangle++(0.008,0.008) (0.428,0.74)rectangle++(0.008,0.008) (0.291,0.858)rectangle++(0.008,0.008) (0.439,0.705)rectangle++(0.008,0.008)
  (0.402,0.76)rectangle++(0.008,0.008) (0.275,0.816)rectangle++(0.008,0.008) (0.416,0.691)rectangle++(0.008,0.008) (0.141,0.706)rectangle++(0.008,0.008) (0.316,0.688)rectangle++(0.008,0.008)
  (0.349,0.871)rectangle++(0.008,0.008) (0.323,0.894)rectangle++(0.008,0.008) (0.466,0.71)rectangle++(0.008,0.008) (0.466,0.574)rectangle++(0.008,0.008) (0.387,0.731)rectangle++(0.008,0.008)
  (0.278,0.855)rectangle++(0.008,0.008) (0.092,0.55)rectangle++(0.008,0.008) (0.237,0.635)rectangle++(0.008,0.008) (0.227,0.823)rectangle++(0.008,0.008) (0.051,0.486)rectangle++(0.008,0.008)
  (0.3,0.879)rectangle++(0.008,0.008) (0.248,0.905)rectangle++(0.008,0.008) (0.436,0.741)rectangle++(0.008,0.008) (0.463,0.6)rectangle++(0.008,0.008) (0.135,0.598)rectangle++(0.008,0.008)
  (0.275,0.641)rectangle++(0.008,0.008) (0.354,0.611)rectangle++(0.008,0.008) (0.379,0.569)rectangle++(0.008,0.008) (0.2,0.891)rectangle++(0.008,0.008) (0.409,0.752)rectangle++(0.008,0.008)
  (0.455,0.614)rectangle++(0.008,0.008) (0.095,0.53)rectangle++(0.008,0.008) (0.101,0.907)rectangle++(0.008,0.008) (-0.002,0.85)rectangle++(0.008,0.008) (0.302,0.806)rectangle++(0.008,0.008)
  (0.425,0.677)rectangle++(0.008,0.008) (0.267,0.896)rectangle++(0.008,0.008) (0.441,0.731)rectangle++(0.008,0.008) (0.462,0.593)rectangle++(0.008,0.008) (0.323,0.834)rectangle++(0.008,0.008)
  (0.256,0.89)rectangle++(0.008,0.008) (0.434,0.732)rectangle++(0.008,0.008) (0.096,0.925)rectangle++(0.008,0.008) (0.273,0.781)rectangle++(0.008,0.008) (0.403,0.676)rectangle++(0.008,0.008)
  (0.283,0.837)rectangle++(0.008,0.008) (0.254,0.855)rectangle++(0.008,0.008) (0.42,0.716)rectangle++(0.008,0.008) (0.21,0.865)rectangle++(0.008,0.008) (0.218,0.9)rectangle++(0.008,0.008)
  (0.375,0.864)rectangle++(0.008,0.008) (0.311,0.819)rectangle++(0.008,0.008) (0.434,0.68)rectangle++(0.008,0.008) (0.293,0.838)rectangle++(0.008,0.008) (0.049,0.484)rectangle++(0.008,0.008)
  (0.028,0.871)rectangle++(0.008,0.008) (0.35,0.89)rectangle++(0.008,0.008) (0.093,0.929)rectangle++(0.008,0.008) (0.374,0.808)rectangle++(0.008,0.008) (0.458,0.652)rectangle++(0.008,0.008)
  (0.302,0.829)rectangle++(0.008,0.008) (0.106,0.937)rectangle++(0.008,0.008) (0.145,0.693)rectangle++(0.008,0.008) (0.096,0.935)rectangle++(0.008,0.008) (0.377,0.809)rectangle++(0.008,0.008)
  (0.107,0.92)rectangle++(0.008,0.008) (0.183,0.919)rectangle++(0.008,0.008) (0.385,0.839)rectangle++(0.008,0.008) (0.096,0.933)rectangle++(0.008,0.008) (-0.005,0.85)rectangle++(0.008,0.008)
  (0.146,0.895)rectangle++(0.008,0.008) (0.061,0.906)rectangle++(0.008,0.008) (0.101,0.804)rectangle++(0.008,0.008) (0.109,0.796)rectangle++(0.008,0.008) (0.126,0.775)rectangle++(0.008,0.008)
  (0.352,0.82)rectangle++(0.008,0.008) (0.102,0.804)rectangle++(0.008,0.008) (0.36,0.864)rectangle++(0.008,0.008) (0.197,0.932)rectangle++(0.008,0.008) (0.422,0.772)rectangle++(0.008,0.008)
  (0.111,0.938)rectangle++(0.008,0.008) (0.374,0.854)rectangle++(0.008,0.008) (0.475,0.673)rectangle++(0.008,0.008) (0.138,0.911)rectangle++(0.008,0.008) (0.166,0.939)rectangle++(0.008,0.008)
  (0.41,0.786)rectangle++(0.008,0.008) (0.468,0.629)rectangle++(0.008,0.008) (0.102,0.806)rectangle++(0.008,0.008) (0.333,0.748)rectangle++(0.008,0.008) (0.125,0.584)rectangle++(0.008,0.008)
  (0.265,0.639)rectangle++(0.008,0.008) (0.348,0.614)rectangle++(0.008,0.008) (0.378,0.573)rectangle++(0.008,0.008) (0.376,0.543)rectangle++(0.008,0.008) (0.365,0.53)rectangle++(0.008,0.008)
  (0.356,0.528)rectangle++(0.008,0.008) (0.267,0.787)rectangle++(0.008,0.008) (0.137,0.594)rectangle++(0.008,0.008) (0.247,0.742)rectangle++(0.008,0.008) (0.377,0.667)rectangle++(0.008,0.008)
  (0.125,0.59)rectangle++(0.008,0.008) (0.189,0.808)rectangle++(0.008,0.008) (0.145,0.742)rectangle++(0.008,0.008) (0.008,0.861)rectangle++(0.008,0.008) (0.15,0.899)rectangle++(0.008,0.008)
  (0.061,0.907)rectangle++(0.008,0.008) (0.032,0.904)rectangle++(0.008,0.008) (0.278,0.899)rectangle++(0.008,0.008) (0.054,0.825)rectangle++(0.008,0.008) (0.154,0.927)rectangle++(0.008,0.008)
  (0.38,0.845)rectangle++(0.008,0.008) (0.328,0.891)rectangle++(0.008,0.008) (0.272,0.814)rectangle++(0.008,0.008) (-0.005,0.852)rectangle++(0.008,0.008) (0.301,0.808)rectangle++(0.008,0.008)
  (0.376,0.852)rectangle++(0.008,0.008) (0.324,0.834)rectangle++(0.008,0.008) (0.445,0.682)rectangle++(0.008,0.008) (0.265,0.763)rectangle++(0.008,0.008) (0.343,0.885)rectangle++(0.008,0.008)
  (0.472,0.698)rectangle++(0.008,0.008) (0.407,0.754)rectangle++(0.008,0.008) (0.454,0.616)rectangle++(0.008,0.008) (0.211,0.897)rectangle++(0.008,0.008) (0.267,0.765)rectangle++(0.008,0.008)
  (0.371,0.85)rectangle++(0.008,0.008) (-0.007,0.817)rectangle++(0.008,0.008) (0.146,0.926)rectangle++(0.008,0.008) (0.396,0.788)rectangle++(0.008,0.008) (0.462,0.635)rectangle++(0.008,0.008)
  (0.029,0.897)rectangle++(0.008,0.008) (0.352,0.858)rectangle++(0.008,0.008) (0.383,0.853)rectangle++(0.008,0.008) (0.028,0.91)rectangle++(0.008,0.008) (0.337,0.822)rectangle++(0.008,0.008)
  (0.447,0.672)rectangle++(0.008,0.008) (0.398,0.752)rectangle++(0.008,0.008) (0,0.83)rectangle++(0.008,0.008) (0.296,0.796)rectangle++(0.008,0.008) (0.153,0.95)rectangle++(0.008,0.008)
  (0.168,0.945)rectangle++(0.008,0.008) (0.386,0.846)rectangle++(0.008,0.008) (0.129,0.575)rectangle++(0.008,0.008) (0.256,0.75)rectangle++(0.008,0.008) (0.049,0.484)rectangle++(0.008,0.008)
  (0.164,0.772)rectangle++(0.008,0.008) (0.359,0.814)rectangle++(0.008,0.008) (0.197,0.925)rectangle++(0.008,0.008) (0.302,0.902)rectangle++(0.008,0.008) (0.004,0.872)rectangle++(0.008,0.008)
  (0.313,0.813)rectangle++(0.008,0.008) (-0.001,0.792)rectangle++(0.008,0.008) (0.282,0.779)rectangle++(0.008,0.008) (0.117,0.563)rectangle++(0.008,0.008) (0.187,0.799)rectangle++(0.008,0.008)
  (0.144,0.693)rectangle++(0.008,0.008) (0.155,0.936)rectangle++(0.008,0.008) (0.404,0.789)rectangle++(0.008,0.008) (0.057,0.861)rectangle++(0.008,0.008) (0.333,0.79)rectangle++(0.008,0.008)
  (0.433,0.658)rectangle++(0.008,0.008) (-0.007,0.817)rectangle++(0.008,0.008) (0.287,0.792)rectangle++(0.008,0.008) (0.314,0.883)rectangle++(0.008,0.008) (0.404,0.813)rectangle++(0.008,0.008)
  (0.283,0.883)rectangle++(0.008,0.008) (0.081,0.528)rectangle++(0.008,0.008) (0.306,0.78)rectangle++(0.008,0.008) (0.317,0.882)rectangle++(0.008,0.008) (0.459,0.706)rectangle++(0.008,0.008)
  (0.461,0.575)rectangle++(0.008,0.008) (0.103,0.804)rectangle++(0.008,0.008) (0.138,0.596)rectangle++(0.008,0.008) (0.275,0.64)rectangle++(0.008,0.008) (0.353,0.611)rectangle++(0.008,0.008)
  (0.365,0.889)rectangle++(0.008,0.008) (0.125,0.584)rectangle++(0.008,0.008) (0.011,0.817)rectangle++(0.008,0.008) (-0.001,0.791)rectangle++(0.008,0.008) (0.33,0.842)rectangle++(0.008,0.008)
  (0.054,0.866)rectangle++(0.008,0.008) (0.352,0.848)rectangle++(0.008,0.008) (0.463,0.679)rectangle++(0.008,0.008) (0.1,0.807)rectangle++(0.008,0.008) (0.144,0.845)rectangle++(0.008,0.008)
  (0.367,0.751)rectangle++(0.008,0.008) (0.435,0.629)rectangle++(0.008,0.008) (0.423,0.548)rectangle++(0.008,0.008) (0.29,0.794)rectangle++(0.008,0.008) (0.357,0.889)rectangle++(0.008,0.008)
  (0.414,0.808)rectangle++(0.008,0.008) (0.271,0.815)rectangle++(0.008,0.008) (0.239,0.884)rectangle++(0.008,0.008) (0.228,0.907)rectangle++(0.008,0.008) (0.427,0.749)rectangle++(0.008,0.008)
  (0.462,0.606)rectangle++(0.008,0.008) (0.427,0.529)rectangle++(0.008,0.008) (0.131,0.932)rectangle++(0.008,0.008) (0.235,0.897)rectangle++(0.008,0.008) (0.427,0.743)rectangle++(0.008,0.008)
  (-0.01,0.812)rectangle++(0.008,0.008) (0.118,0.573)rectangle++(0.008,0.008) (0.187,0.802)rectangle++(0.008,0.008) (0.119,0.569)rectangle++(0.008,0.008) (0.257,0.634)rectangle++(0.008,0.008)
  (0.059,0.896)rectangle++(0.008,0.008) (0.282,0.898)rectangle++(0.008,0.008) (0.4,0.82)rectangle++(0.008,0.008) (0.475,0.648)rectangle++(0.008,0.008) (0.307,0.827)rectangle++(0.008,0.008)
  (0.251,0.888)rectangle++(0.008,0.008) (0.431,0.733)rectangle++(0.008,0.008) (0.404,0.766)rectangle++(0.008,0.008) (0.458,0.622)rectangle++(0.008,0.008) (0.145,0.946)rectangle++(0.008,0.008)
  (0.404,0.797)rectangle++(0.008,0.008) (0.272,0.802)rectangle++(0.008,0.008) (0.41,0.686)rectangle++(0.008,0.008) (0.431,0.584)rectangle++(0.008,0.008) (0.405,0.529)rectangle++(0.008,0.008)
  (0.117,0.571)rectangle++(0.008,0.008) (0.255,0.749)rectangle++(0.008,0.008) (0.383,0.668)rectangle++(0.008,0.008) (0.217,0.883)rectangle++(0.008,0.008) (0.118,0.929)rectangle++(0.008,0.008)
  (0.138,0.594)rectangle++(0.008,0.008) (0.327,0.784)rectangle++(0.008,0.008) (0.255,0.898)rectangle++(0.008,0.008) (0.084,0.532)rectangle++(0.008,0.008) (0.228,0.868)rectangle++(0.008,0.008)
  (0.13,0.581)rectangle++(0.008,0.008) (0.292,0.853)rectangle++(0.008,0.008) (0.438,0.702)rectangle++(0.008,0.008) (0.45,0.581)rectangle++(0.008,0.008) (0.412,0.521)rectangle++(0.008,0.008)
  (0.286,0.785)rectangle++(0.008,0.008) (0.096,0.905)rectangle++(0.008,0.008) (0.366,0.85)rectangle++(0.008,0.008) (0.012,0.878)rectangle++(0.008,0.008) (0.042,0.887)rectangle++(0.008,0.008)
  (0.335,0.807)rectangle++(0.008,0.008) (0.131,0.592)rectangle++(0.008,0.008) (0.191,0.809)rectangle++(0.008,0.008) (0.011,0.87)rectangle++(0.008,0.008) (0.344,0.856)rectangle++(0.008,0.008)
  (0.264,0.911)rectangle++(0.008,0.008) (0.38,0.862)rectangle++(0.008,0.008) (0.481,0.675)rectangle++(0.008,0.008) (0.285,0.908)rectangle++(0.008,0.008) (0.454,0.73)rectangle++(0.008,0.008)
  (0.3,0.855)rectangle++(0.008,0.008) (0.442,0.7)rectangle++(0.008,0.008) (0.296,0.845)rectangle++(0.008,0.008) (0.247,0.893)rectangle++(0.008,0.008) (0.393,0.825)rectangle++(0.008,0.008)
  (0.143,0.712)rectangle++(0.008,0.008) (0.144,0.721)rectangle++(0.008,0.008) (0.362,0.874)rectangle++(0.008,0.008) (0.253,0.905)rectangle++(0.008,0.008) (0.438,0.74)rectangle++(0.008,0.008)
  (0.065,0.912)rectangle++(0.008,0.008) (0.128,0.946)rectangle++(0.008,0.008) (0.257,0.849)rectangle++(0.008,0.008) (0.388,0.815)rectangle++(0.008,0.008) (0.278,0.883)rectangle++(0.008,0.008)
  (-0.008,0.853)rectangle++(0.008,0.008) (0.092,0.913)rectangle++(0.008,0.008) (0.366,0.852)rectangle++(0.008,0.008) (0.471,0.675)rectangle++(0.008,0.008) (0.213,0.906)rectangle++(0.008,0.008)
  (0.053,0.49)rectangle++(0.008,0.008) (0.096,0.832)rectangle++(0.008,0.008) (0.258,0.905)rectangle++(0.008,0.008) (0.234,0.914)rectangle++(0.008,0.008) (0.137,0.59)rectangle++(0.008,0.008)
  (0.132,0.592)rectangle++(0.008,0.008) (0.099,0.909)rectangle++(0.008,0.008) (0.366,0.873)rectangle++(0.008,0.008) (0.385,0.853)rectangle++(0.008,0.008) (0.012,0.88)rectangle++(0.008,0.008)
  (0.319,0.814)rectangle++(0.008,0.008) (0.436,0.675)rectangle++(0.008,0.008) (0.091,0.549)rectangle++(0.008,0.008) (0.094,0.846)rectangle++(0.008,0.008) (0.059,0.892)rectangle++(0.008,0.008)
  (0.357,0.853)rectangle++(0.008,0.008) (0.147,0.608)rectangle++(0.008,0.008) (0.167,0.896)rectangle++(0.008,0.008) (0.396,0.766)rectangle++(0.008,0.008) (0.403,0.778)rectangle++(0.008,0.008)
  (0.406,0.779)rectangle++(0.008,0.008) (0.119,0.562)rectangle++(0.008,0.008) (0.32,0.872)rectangle++(0.008,0.008) (0.192,0.934)rectangle++(0.008,0.008) (0.42,0.775)rectangle++(0.008,0.008)
  (0.408,0.776)rectangle++(0.008,0.008) (0.285,0.87)rectangle++(0.008,0.008) (0.28,0.913)rectangle++(0.008,0.008) (0.153,0.953)rectangle++(0.008,0.008) (0.069,0.521)rectangle++(0.008,0.008)
  (0.303,0.78)rectangle++(0.008,0.008) (0.387,0.795)rectangle++(0.008,0.008) (0.46,0.641)rectangle++(0.008,0.008) (0.438,0.545)rectangle++(0.008,0.008) (0.084,0.517)rectangle++(0.008,0.008)
  (0.125,0.907)rectangle++(0.008,0.008) (0.381,0.786)rectangle++(0.008,0.008) (0.454,0.64)rectangle++(0.008,0.008) (0.435,0.547)rectangle++(0.008,0.008) (0.172,0.927)rectangle++(0.008,0.008)
  (0.409,0.779)rectangle++(0.008,0.008) (0.115,0.919)rectangle++(0.008,0.008) (0.186,0.919)rectangle++(0.008,0.008) (0.412,0.77)rectangle++(0.008,0.008) (0.118,0.56)rectangle++(0.008,0.008)
  (0.346,0.843)rectangle++(0.008,0.008) (0.303,0.886)rectangle++(0.008,0.008) (0.249,0.908)rectangle++(0.008,0.008) (0.276,0.916)rectangle++(0.008,0.008) (0.334,0.898)rectangle++(0.008,0.008)
  (0.472,0.708)rectangle++(0.008,0.008) (0.468,0.571)rectangle++(0.008,0.008) (0.34,0.852)rectangle++(0.008,0.008) (0.459,0.684)rectangle++(0.008,0.008) (0.453,0.565)rectangle++(0.008,0.008)
  (0.383,0.731)rectangle++(0.008,0.008) (0.395,0.773)rectangle++(0.008,0.008) (0.11,0.937)rectangle++(0.008,0.008) (0.384,0.806)rectangle++(0.008,0.008) (0.406,0.787)rectangle++(0.008,0.008)
  (0.114,0.92)rectangle++(0.008,0.008) (-0.003,0.852)rectangle++(0.008,0.008) (0.302,0.807)rectangle++(0.008,0.008) (0.096,0.911)rectangle++(0.008,0.008) (0.368,0.798)rectangle++(0.008,0.008)
  (0.295,0.903)rectangle++(0.008,0.008) (0.457,0.724)rectangle++(0.008,0.008) (0.211,0.913)rectangle++(0.008,0.008) (0.042,0.912)rectangle++(0.008,0.008) (0.161,0.916)rectangle++(0.008,0.008)
  (0.148,0.738)rectangle++(0.008,0.008) (0.197,0.858)rectangle++(0.008,0.008) (0.396,0.738)rectangle++(0.008,0.008) (0.444,0.612)rectangle++(0.008,0.008) (0.117,0.791)rectangle++(0.008,0.008)
  (0.091,0.915)rectangle++(0.008,0.008) (-0.003,0.849)rectangle++(0.008,0.008) (0.002,0.837)rectangle++(0.008,0.008) (0.231,0.885)rectangle++(0.008,0.008) (0.389,0.826)rectangle++(0.008,0.008)
  (0.329,0.888)rectangle++(0.008,0.008) (0.134,0.584)rectangle++(0.008,0.008) (0.269,0.635)rectangle++(0.008,0.008) (0.096,0.536)rectangle++(0.008,0.008) (0.096,0.844)rectangle++(0.008,0.008)
  (0.345,0.768)rectangle++(0.008,0.008) (0.377,0.846)rectangle++(0.008,0.008) (0.274,0.894)rectangle++(0.008,0.008) (0.24,0.91)rectangle++(0.008,0.008) (0.434,0.747)rectangle++(0.008,0.008)
  (0.464,0.603)rectangle++(0.008,0.008) (0.427,0.526)rectangle++(0.008,0.008) (0.382,0.505)rectangle++(0.008,0.008) (0.361,0.73)rectangle++(0.008,0.008) (0.269,0.855)rectangle++(0.008,0.008)
  (0.31,0.818)rectangle++(0.008,0.008) (0.096,0.913)rectangle++(0.008,0.008) (0.369,0.8)rectangle++(0.008,0.008) (0.454,0.65)rectangle++(0.008,0.008) (0.349,0.891)rectangle++(0.008,0.008)
  (0.477,0.699)rectangle++(0.008,0.008) (0.467,0.566)rectangle++(0.008,0.008) (0.046,0.881)rectangle++(0.008,0.008) (0.012,0.881)rectangle++(0.008,0.008) (0.346,0.858)rectangle++(0.008,0.008)
  (0.273,0.794)rectangle++(0.008,0.008) (0.408,0.682)rectangle++(0.008,0.008) (0.051,0.486)rectangle++(0.008,0.008) (0.008,0.838)rectangle++(0.008,0.008) (0.302,0.797)rectangle++(0.008,0.008)
  (0.422,0.673)rectangle++(0.008,0.008) (0.15,0.704)rectangle++(0.008,0.008) (0.198,0.846)rectangle++(0.008,0.008) (0.376,0.823)rectangle++(0.008,0.008) (0.018,0.833)rectangle++(0.008,0.008)
  (0.122,0.576)rectangle++(0.008,0.008) (0.248,0.736)rectangle++(0.008,0.008) (0.375,0.665)rectangle++(0.008,0.008) (0.146,0.648)rectangle++(0.008,0.008) (0.297,0.66)rectangle++(0.008,0.008)
  (0.371,0.612)rectangle++(0.008,0.008) (0.326,0.858)rectangle++(0.008,0.008) (0.154,0.952)rectangle++(0.008,0.008) (0.384,0.85)rectangle++(0.008,0.008) (0.478,0.668)rectangle++(0.008,0.008)
  (0.456,0.551)rectangle++(0.008,0.008) (0.086,0.518)rectangle++(0.008,0.008) (0.122,0.943)rectangle++(0.008,0.008) (0.392,0.804)rectangle++(0.008,0.008) (0.349,0.878)rectangle++(0.008,0.008)
  (0.41,0.807)rectangle++(0.008,0.008) (0.285,0.881)rectangle++(0.008,0.008) (0.379,0.858)rectangle++(0.008,0.008) (0.262,0.76)rectangle++(0.008,0.008) (0.162,0.94)rectangle++(0.008,0.008)
  (0.409,0.788)rectangle++(0.008,0.008) (0.263,0.906)rectangle++(0.008,0.008) (0.443,0.737)rectangle++(0.008,0.008) (0.327,0.858)rectangle++(0.008,0.008) (0.455,0.691)rectangle++(0.008,0.008)
  (0.361,0.885)rectangle++(0.008,0.008) (0.269,0.907)rectangle++(0.008,0.008) (0.399,0.824)rectangle++(0.008,0.008) (0.476,0.651)rectangle++(0.008,0.008) (0.401,0.744)rectangle++(0.008,0.008)
  (0.131,0.591)rectangle++(0.008,0.008) (0.165,0.894)rectangle++(0.008,0.008) (0.203,0.91)rectangle++(0.008,0.008) (0.068,0.915)rectangle++(0.008,0.008) (0.017,0.897)rectangle++(0.008,0.008)
  (0.27,0.768)rectangle++(0.008,0.008) (0.146,0.904)rectangle++(0.008,0.008) (0.375,0.842)rectangle++(0.008,0.008) (0.471,0.667)rectangle++(0.008,0.008) (0.275,0.811)rectangle++(0.008,0.008)
  (0.081,0.528)rectangle++(0.008,0.008) (0.028,0.871)rectangle++(0.008,0.008) (0.358,0.874)rectangle++(0.008,0.008) (0.109,0.561)rectangle++(0.008,0.008) (0.125,0.911)rectangle++(0.008,0.008)
  (0.05,0.485)rectangle++(0.008,0.008) (0.195,0.62)rectangle++(0.008,0.008) (0.102,0.869)rectangle++(0.008,0.008) (0.339,0.886)rectangle++(0.008,0.008) (0.261,0.907)rectangle++(0.008,0.008)
  (0.261,0.911)rectangle++(0.008,0.008) (0.311,0.9)rectangle++(0.008,0.008) (0.086,0.926)rectangle++(0.008,0.008) (0.369,0.809)rectangle++(0.008,0.008) (0.037,0.91)rectangle++(0.008,0.008)
  (0.341,0.819)rectangle++(0.008,0.008) (0.243,0.734)rectangle++(0.008,0.008) (0.095,0.892)rectangle++(0.008,0.008) (0.143,0.73)rectangle++(0.008,0.008) (0.08,0.88)rectangle++(0.008,0.008)
  (0.035,0.903)rectangle++(0.008,0.008) (0.338,0.817)rectangle++(0.008,0.008) (0.445,0.669)rectangle++(0.008,0.008) (0.142,0.603)rectangle++(0.008,0.008) (0.092,0.929)rectangle++(0.008,0.008)
  (0.369,0.855)rectangle++(0.008,0.008) (0.473,0.676)rectangle++(0.008,0.008) (0.457,0.556)rectangle++(0.008,0.008) (0.148,0.843)rectangle++(0.008,0.008) (0.368,0.749)rectangle++(0.008,0.008)
  (0.435,0.627)rectangle++(0.008,0.008) (0.422,0.548)rectangle++(0.008,0.008) (0.123,0.584)rectangle++(0.008,0.008) (0.291,0.854)rectangle++(0.008,0.008) (0.245,0.897)rectangle++(0.008,0.008)
  (0.117,0.94)rectangle++(0.008,0.008) (-0.003,0.838)rectangle++(0.008,0.008) (0.082,0.529)rectangle++(0.008,0.008) (0.252,0.869)rectangle++(0.008,0.008) (0.232,0.902)rectangle++(0.008,0.008)
  (0.152,0.953)rectangle++(0.008,0.008) (0.122,0.944)rectangle++(0.008,0.008) (0.392,0.804)rectangle++(0.008,0.008) (0.466,0.644)rectangle++(0.008,0.008) (0.442,0.544)rectangle++(0.008,0.008)
  (0.336,0.848)rectangle++(0.008,0.008) (0.004,0.781)rectangle++(0.008,0.008) (0.149,0.872)rectangle++(0.008,0.008) (0.379,0.762)rectangle++(0.008,0.008) (0.399,0.78)rectangle++(0.008,0.008)
  (0.353,0.891)rectangle++(0.008,0.008) (0.167,0.948)rectangle++(0.008,0.008) (0.414,0.79)rectangle++(0.008,0.008) (0.115,0.922)rectangle++(0.008,0.008) (0.381,0.797)rectangle++(0.008,0.008)
  (0.458,0.644)rectangle++(0.008,0.008) (0.396,0.745)rectangle++(0.008,0.008) (0.399,0.774)rectangle++(0.008,0.008) (0.237,0.899)rectangle++(0.008,0.008) (0.268,0.762)rectangle++(0.008,0.008)
  (0.394,0.669)rectangle++(0.008,0.008) (0.079,0.526)rectangle++(0.008,0.008) (0.223,0.628)rectangle++(0.008,0.008) (0.326,0.624)rectangle++(0.008,0.008) (0.371,0.585)rectangle++(0.008,0.008)
  (0.148,0.845)rectangle++(0.008,0.008) (0.197,0.893)rectangle++(0.008,0.008) (0.383,0.832)rectangle++(0.008,0.008) (0.014,0.879)rectangle++(0.008,0.008) (0.357,0.875)rectangle++(0.008,0.008)
  (0.078,0.906)rectangle++(0.008,0.008) (0.263,0.761)rectangle++(0.008,0.008) (0.391,0.67)rectangle++(0.008,0.008) (0.418,0.583)rectangle++(0.008,0.008) (0.143,0.721)rectangle++(0.008,0.008)
  (0.322,0.695)rectangle++(0.008,0.008) (0.256,0.843)rectangle++(0.008,0.008) (0.127,0.576)rectangle++(0.008,0.008) (0.107,0.926)rectangle++(0.008,0.008) (0.183,0.921)rectangle++(0.008,0.008)
  (0.412,0.772)rectangle++(0.008,0.008) (0.286,0.869)rectangle++(0.008,0.008) (0.048,0.826)rectangle++(0.008,0.008) (0.345,0.842)rectangle++(0.008,0.008) (0.404,0.8)rectangle++(0.008,0.008)
  (0.122,0.566)rectangle++(0.008,0.008) (0.164,0.89)rectangle++(0.008,0.008) (0.37,0.867)rectangle++(0.008,0.008) (0.272,0.901)rectangle++(0.008,0.008) (0.084,0.542)rectangle++(0.008,0.008)
  (0.176,0.792)rectangle++(0.008,0.008) (0.102,0.563)rectangle++(0.008,0.008) (0.088,0.54)rectangle++(0.008,0.008) (0.309,0.781)rectangle++(0.008,0.008) (0.42,0.663)rectangle++(0.008,0.008)
  (0.288,0.832)rectangle++(0.008,0.008) (0.244,0.889)rectangle++(0.008,0.008) (0.308,0.897)rectangle++(0.008,0.008) (0.461,0.716)rectangle++(0.008,0.008) (0.099,0.808)rectangle++(0.008,0.008)
  (0.333,0.75)rectangle++(0.008,0.008) (0.145,0.733)rectangle++(0.008,0.008) (0.147,0.712)rectangle++(0.008,0.008) (0.114,0.572)rectangle++(0.008,0.008) (0.255,0.637)rectangle++(0.008,0.008)
  (0.356,0.774)rectangle++(0.008,0.008) (0.058,0.91)rectangle++(0.008,0.008) (0.351,0.812)rectangle++(0.008,0.008) (0.449,0.662)rectangle++(0.008,0.008) (0.397,0.75)rectangle++(0.008,0.008)
  (0.449,0.617)rectangle++(0.008,0.008) (0.07,0.82)rectangle++(0.008,0.008) (0.348,0.837)rectangle++(0.008,0.008) (0.145,0.955)rectangle++(0.008,0.008) (0.406,0.801)rectangle++(0.008,0.008)
  (0.15,0.722)rectangle++(0.008,0.008) (0.232,0.892)rectangle++(0.008,0.008) (0.148,0.735)rectangle++(0.008,0.008) (0.35,0.809)rectangle++(0.008,0.008) (0.448,0.661)rectangle++(0.008,0.008)
  (0.124,0.585)rectangle++(0.008,0.008) (0.265,0.64)rectangle++(0.008,0.008) (0.066,0.822)rectangle++(0.008,0.008) (0.348,0.838)rectangle++(0.008,0.008) (0.457,0.675)rectangle++(0.008,0.008)
  (0.151,0.704)rectangle++(0.008,0.008) (0.12,0.583)rectangle++(0.008,0.008) (0.248,0.737)rectangle++(0.008,0.008) (0.221,0.889)rectangle++(0.008,0.008) (0.159,0.948)rectangle++(0.008,0.008)
  (0.201,0.928)rectangle++(0.008,0.008) (0.066,0.917)rectangle++(0.008,0.008) (0.169,0.918)rectangle++(0.008,0.008) (0.404,0.775)rectangle++(0.008,0.008) (0.249,0.742)rectangle++(0.008,0.008)
  (0.009,0.866)rectangle++(0.008,0.008) (0.313,0.809)rectangle++(0.008,0.008) (0.252,0.881)rectangle++(0.008,0.008) (0.429,0.729)rectangle++(0.008,0.008) (0.146,0.864)rectangle++(0.008,0.008)
  (0.374,0.759)rectangle++(0.008,0.008) (0.397,0.78)rectangle++(0.008,0.008) (0.459,0.631)rectangle++(0.008,0.008) (0.302,0.822)rectangle++(0.008,0.008) (0.189,0.926)rectangle++(0.008,0.008)
  (0.417,0.772)rectangle++(0.008,0.008) (0.465,0.621)rectangle++(0.008,0.008) (0.434,0.534)rectangle++(0.008,0.008) (0.388,0.506)rectangle++(0.008,0.008) (0.357,0.509)rectangle++(0.008,0.008)
  (0.065,0.505)rectangle++(0.008,0.008) (0.169,0.779)rectangle++(0.008,0.008) (0.355,0.712)rectangle++(0.008,0.008) (0.416,0.615)rectangle++(0.008,0.008) (0.409,0.549)rectangle++(0.008,0.008)
  (0.382,0.521)rectangle++(0.008,0.008) (0.36,0.518)rectangle++(0.008,0.008) (0.268,0.784)rectangle++(0.008,0.008) (0.402,0.679)rectangle++(0.008,0.008) (0.425,0.583)rectangle++(0.008,0.008)
  (0.103,0.8)rectangle++(0.008,0.008) (0.332,0.745)rectangle++(0.008,0.008) (0.022,0.871)rectangle++(0.008,0.008) (0.321,0.807)rectangle++(0.008,0.008) (0.142,0.605)rectangle++(0.008,0.008)
  (0.35,0.845)rectangle++(0.008,0.008) (0.265,0.894)rectangle++(0.008,0.008) (0.126,0.894)rectangle++(0.008,0.008) (0.19,0.91)rectangle++(0.008,0.008) (0.411,0.765)rectangle++(0.008,0.008)
  (0.405,0.776)rectangle++(0.008,0.008) (0.461,0.626)rectangle++(0.008,0.008) (0.321,0.872)rectangle++(0.008,0.008) (0.457,0.7)rectangle++(0.008,0.008) (0.404,0.756)rectangle++(0.008,0.008)
  (0.454,0.618)rectangle++(0.008,0.008) (0.391,0.741)rectangle++(0.008,0.008) (0.443,0.615)rectangle++(0.008,0.008) (0.095,0.53)rectangle++(0.008,0.008) (0.232,0.625)rectangle++(0.008,0.008)
  (0.037,0.902)rectangle++(0.008,0.008) (0.043,0.893)rectangle++(0.008,0.008) (0.045,0.893)rectangle++(0.008,0.008) (0.163,0.91)rectangle++(0.008,0.008) (0.17,0.939)rectangle++(0.008,0.008)
  (0.298,0.904)rectangle++(0.008,0.008) (0.458,0.723)rectangle++(0.008,0.008) (0.467,0.583)rectangle++(0.008,0.008) (0.389,0.732)rectangle++(0.008,0.008) (0.278,0.856)rectangle++(0.008,0.008)
  (0.393,0.813)rectangle++(0.008,0.008) (0.469,0.648)rectangle++(0.008,0.008) (0.273,0.771)rectangle++(0.008,0.008) (0.239,0.869)rectangle++(0.008,0.008) (0.133,0.592)rectangle++(0.008,0.008)
  (0.272,0.639)rectangle++(0.008,0.008) (0.352,0.612)rectangle++(0.008,0.008) (0.143,0.944)rectangle++(0.008,0.008) (0.401,0.797)rectangle++(0.008,0.008) (0.331,0.884)rectangle++(0.008,0.008)
  (0.211,0.847)rectangle++(0.008,0.008) (0.146,0.86)rectangle++(0.008,0.008) (0.266,0.759)rectangle++(0.008,0.008) (0.376,0.796)rectangle++(0.008,0.008) (0.368,0.881)rectangle++(0.008,0.008)
  (0.168,0.947)rectangle++(0.008,0.008) (0.141,0.597)rectangle++(0.008,0.008) (0.195,0.81)rectangle++(0.008,0.008) (0.174,0.925)rectangle++(0.008,0.008) (0.409,0.777)rectangle++(0.008,0.008)
  (0.17,0.941)rectangle++(0.008,0.008) (0.385,0.845)rectangle++(0.008,0.008) (0.277,0.894)rectangle++(0.008,0.008) (0.445,0.726)rectangle++(0.008,0.008) (0.112,0.551)rectangle++(0.008,0.008)
  (0.185,0.795)rectangle++(0.008,0.008) (0.118,0.567)rectangle++(0.008,0.008) (0.265,0.891)rectangle++(0.008,0.008) (0.005,0.805)rectangle++(0.008,0.008) (0.289,0.783)rectangle++(0.008,0.008)
  (0.411,0.671)rectangle++(0.008,0.008) (0.112,0.574)rectangle++(0.008,0.008) (0.185,0.802)rectangle++(0.008,0.008) (0.144,0.729)rectangle++(0.008,0.008) (0.294,0.874)rectangle++(0.008,0.008)
  (0.246,0.903)rectangle++(0.008,0.008) (0.394,0.827)rectangle++(0.008,0.008) (0.126,0.571)rectangle++(0.008,0.008) (0.154,0.929)rectangle++(0.008,0.008) (0.401,0.786)rectangle++(0.008,0.008)
  (0.406,0.781)rectangle++(0.008,0.008) (0.464,0.628)rectangle++(0.008,0.008) (0.248,0.74)rectangle++(0.008,0.008) (0.098,0.811)rectangle++(0.008,0.008) (0.16,0.925)rectangle++(0.008,0.008)
  (0.403,0.782)rectangle++(0.008,0.008) (0.463,0.63)rectangle++(0.008,0.008) (0.436,0.539)rectangle++(0.008,0.008) (0.081,0.529)rectangle++(0.008,0.008) (0.225,0.629)rectangle++(0.008,0.008)
  (0.256,0.883)rectangle++(0.008,0.008) (0.431,0.728)rectangle++(0.008,0.008) (0.144,0.745)rectangle++(0.008,0.008) (0.246,0.738)rectangle++(0.008,0.008) (0.23,0.857)rectangle++(0.008,0.008)
  (0.179,0.932)rectangle++(0.008,0.008) (0.111,0.563)rectangle++(0.008,0.008) (0.32,0.895)rectangle++(0.008,0.008) (0.407,0.814)rectangle++(0.008,0.008) (0.385,0.845)rectangle++(0.008,0.008)
  (0.477,0.665)rectangle++(0.008,0.008) (0.308,0.833)rectangle++(0.008,0.008) (0.437,0.687)rectangle++(0.008,0.008) (0.259,0.85)rectangle++(0.008,0.008) (0.106,0.568)rectangle++(0.008,0.008)
  (0.095,0.815)rectangle++(0.008,0.008) (0.352,0.833)rectangle++(0.008,0.008) (0.457,0.671)rectangle++(0.008,0.008) (0.4,0.751)rectangle++(0.008,0.008) (0.45,0.617)rectangle++(0.008,0.008)
  (0.425,0.537)rectangle++(0.008,0.008) (0.029,0.459)rectangle++(0.008,0.008) (0.176,0.616)rectangle++(0.008,0.008) (0.338,0.782)rectangle++(0.008,0.008) (0.433,0.653)rectangle++(0.008,0.008)
  (0.293,0.829)rectangle++(0.008,0.008) (0.392,0.806)rectangle++(0.008,0.008) (0.408,0.786)rectangle++(0.008,0.008) (0.466,0.63)rectangle++(0.008,0.008) (0.437,0.538)rectangle++(0.008,0.008)
  (0.294,0.79)rectangle++(0.008,0.008) (0.416,0.672)rectangle++(0.008,0.008) (0.273,0.865)rectangle++(0.008,0.008) (0.433,0.714)rectangle++(0.008,0.008) (0.452,0.588)rectangle++(0.008,0.008)
  (0.416,0.524)rectangle++(0.008,0.008) (0.15,0.683)rectangle++(0.008,0.008) (0.198,0.839)rectangle++(0.008,0.008) (0.389,0.729)rectangle++(0.008,0.008) (0.111,0.553)rectangle++(0.008,0.008)
  (0.185,0.795)rectangle++(0.008,0.008) (0.278,0.911)rectangle++(0.008,0.008) (0.241,0.916)rectangle++(0.008,0.008) (0.073,0.92)rectangle++(0.008,0.008) (0.364,0.856)rectangle++(0.008,0.008)
  (0.099,0.926)rectangle++(0.008,0.008) (0.369,0.853)rectangle++(0.008,0.008) (0.473,0.674)rectangle++(0.008,0.008) (0.302,0.886)rectangle++(0.008,0.008) (0.403,0.815)rectangle++(0.008,0.008)
  (0.25,0.904)rectangle++(0.008,0.008) (0.139,0.878)rectangle++(0.008,0.008) (0.376,0.768)rectangle++(0.008,0.008) (0.327,0.859)rectangle++(0.008,0.008) (0.149,0.73)rectangle++(0.008,0.008)
  (0.115,0.574)rectangle++(0.008,0.008) (0.32,0.783)rectangle++(0.008,0.008) (0.39,0.793)rectangle++(0.008,0.008) (0.279,0.876)rectangle++(0.008,0.008) (0.356,0.89)rectangle++(0.008,0.008)
  (0.094,0.93)rectangle++(0.008,0.008) (0.179,0.922)rectangle++(0.008,0.008) (0.157,0.95)rectangle++(0.008,0.008) (0.41,0.794)rectangle++(0.008,0.008) (0.135,0.596)rectangle++(0.008,0.008)
  (0.166,0.895)rectangle++(0.008,0.008) (0.064,0.908)rectangle++(0.008,0.008) (0.169,0.915)rectangle++(0.008,0.008) (0.128,0.592)rectangle++(0.008,0.008) (0.165,0.894)rectangle++(0.008,0.008)
  (0.148,0.735)rectangle++(0.008,0.008) (0.329,0.7)rectangle++(0.008,0.008) (0.399,0.619)rectangle++(0.008,0.008) (0.403,0.556)rectangle++(0.008,0.008) (0.146,0.871)rectangle++(0.008,0.008)
  (0.294,0.894)rectangle++(0.008,0.008) (0.453,0.72)rectangle++(0.008,0.008) (0.145,0.911)rectangle++(0.008,0.008) (0.269,0.763)rectangle++(0.008,0.008) (0.125,0.588)rectangle++(0.008,0.008)
  (0.146,0.743)rectangle++(0.008,0.008) (0.331,0.704)rectangle++(0.008,0.008) (0.402,0.62)rectangle++(0.008,0.008) (0.297,0.893)rectangle++(0.008,0.008) (0.403,0.817)rectangle++(0.008,0.008)
  (0.475,0.646)rectangle++(0.008,0.008) (0.447,0.542)rectangle++(0.008,0.008) (0.397,0.505)rectangle++(0.008,0.008) (-0.004,0.814)rectangle++(0.008,0.008) (0.288,0.79)rectangle++(0.008,0.008)
  (0.413,0.674)rectangle++(0.008,0.008) (0.429,0.577)rectangle++(0.008,0.008) (0.292,0.803)rectangle++(0.008,0.008) (0.272,0.82)rectangle++(0.008,0.008) (0.123,0.571)rectangle++(0.008,0.008)
  (0.26,0.633)rectangle++(0.008,0.008) (0.344,0.613)rectangle++(0.008,0.008) (0.327,0.863)rectangle++(0.008,0.008) (0.257,0.9)rectangle++(0.008,0.008) (0.103,0.9)rectangle++(0.008,0.008)
  (0.115,0.785)rectangle++(0.008,0.008) (0.261,0.753)rectangle++(0.008,0.008) (0.387,0.668)rectangle++(0.008,0.008) (0.278,0.834)rectangle++(0.008,0.008) (0.424,0.698)rectangle++(0.008,0.008)
  (0.03,0.88)rectangle++(0.008,0.008) (0.327,0.808)rectangle++(0.008,0.008) (0.32,0.886)rectangle++(0.008,0.008) (0.406,0.812)rectangle++(0.008,0.008) (0.059,0.915)rectangle++(0.008,0.008)
  (0.353,0.813)rectangle++(0.008,0.008) (0.451,0.662)rectangle++(0.008,0.008) (0.398,0.75)rectangle++(0.008,0.008) (0.4,0.775)rectangle++(0.008,0.008) (0.15,0.718)rectangle++(0.008,0.008)
  (0.348,0.806)rectangle++(0.008,0.008) (0.399,0.793)rectangle++(0.008,0.008) (0.465,0.636)rectangle++(0.008,0.008) (0.439,0.541)rectangle++(0.008,0.008) (0.393,0.507)rectangle++(0.008,0.008)
  (0.118,0.582)rectangle++(0.008,0.008) (0.086,0.527)rectangle++(0.008,0.008) (0.318,0.88)rectangle++(0.008,0.008) (0.273,0.827)rectangle++(0.008,0.008) (0.312,0.888)rectangle++(0.008,0.008)
  (0.223,0.858)rectangle++(0.008,0.008) (0.407,0.729)rectangle++(0.008,0.008) (0.38,0.839)rectangle++(0.008,0.008) (0.472,0.664)rectangle++(0.008,0.008) (0.035,0.884)rectangle++(0.008,0.008)
  (0.331,0.808)rectangle++(0.008,0.008) (0.439,0.667)rectangle++(0.008,0.008) (0.438,0.565)rectangle++(0.008,0.008) (0.401,0.518)rectangle++(0.008,0.008) (0.368,0.51)rectangle++(0.008,0.008)
  (0.326,0.886)rectangle++(0.008,0.008) (0.128,0.894)rectangle++(0.008,0.008) (0.19,0.91)rectangle++(0.008,0.008) (0.126,0.592)rectangle++(0.008,0.008) (0.325,0.785)rectangle++(0.008,0.008)
  (0.127,0.937)rectangle++(0.008,0.008) (0.377,0.851)rectangle++(0.008,0.008) (0.385,0.85)rectangle++(0.008,0.008) (0.25,0.896)rectangle++(0.008,0.008) (0.309,0.893)rectangle++(0.008,0.008)
  (0.251,0.91)rectangle++(0.008,0.008) (0.128,0.594)rectangle++(0.008,0.008) (0.325,0.785)rectangle++(0.008,0.008) (0.37,0.872)rectangle++(0.008,0.008) (0.137,0.601)rectangle++(0.008,0.008)
  (0.04,0.473)rectangle++(0.008,0.008) (0.244,0.735)rectangle++(0.008,0.008) (0.373,0.665)rectangle++(0.008,0.008) (0.273,0.833)rectangle++(0.008,0.008) (0.217,0.88)rectangle++(0.008,0.008)
  (0.412,0.741)rectangle++(0.008,0.008) (0.452,0.608)rectangle++(0.008,0.008) (0.338,0.857)rectangle++(0.008,0.008) (0.01,0.875)rectangle++(0.008,0.008) (0.151,0.904)rectangle++(0.008,0.008)
  (0.061,0.907)rectangle++(0.008,0.008) (0.359,0.855)rectangle++(0.008,0.008) (0.268,0.897)rectangle++(0.008,0.008) (0.398,0.822)rectangle++(0.008,0.008) (0.475,0.65)rectangle++(0.008,0.008)
  (0.4,0.744)rectangle++(0.008,0.008) (0.448,0.614)rectangle++(0.008,0.008) (0.095,0.53)rectangle++(0.008,0.008) (0.143,0.712)rectangle++(0.008,0.008) (0.196,0.849)rectangle++(0.008,0.008)
  (0.376,0.824)rectangle++(0.008,0.008) (0.106,0.804)rectangle++(0.008,0.008) (0.183,0.88)rectangle++(0.008,0.008) (0.004,0.859)rectangle++(0.008,0.008) (0.056,0.865)rectangle++(0.008,0.008)
  (0.239,0.861)rectangle++(0.008,0.008) (0.258,0.751)rectangle++(0.008,0.008) (0.342,0.884)rectangle++(0.008,0.008) (0.294,0.908)rectangle++(0.008,0.008) (0.232,0.9)rectangle++(0.008,0.008)
  (0.427,0.745)rectangle++(0.008,0.008) (0.46,0.605)rectangle++(0.008,0.008) (0.144,0.899)rectangle++(0.008,0.008) (0.196,0.912)rectangle++(0.008,0.008) (0.386,0.836)rectangle++(0.008,0.008)
  (0.411,0.793)rectangle++(0.008,0.008) (0.146,0.83)rectangle++(0.008,0.008) (0.196,0.889)rectangle++(0.008,0.008) (0.266,0.764)rectangle++(0.008,0.008) (0.12,0.932)rectangle++(0.008,0.008)
  (0.188,0.923)rectangle++(0.008,0.008) (0.212,0.877)rectangle++(0.008,0.008) (0.267,0.765)rectangle++(0.008,0.008) (0.312,0.879)rectangle++(0.008,0.008) (0.252,0.905)rectangle++(0.008,0.008)
  (0.244,0.902)rectangle++(0.008,0.008) (0.394,0.827)rectangle++(0.008,0.008) (0.28,0.887)rectangle++(0.008,0.008) (0.232,0.899)rectangle++(0.008,0.008) (0.426,0.745)rectangle++(0.008,0.008)
  (0.145,0.914)rectangle++(0.008,0.008) (0.055,0.867)rectangle++(0.008,0.008) (0.013,0.881)rectangle++(0.008,0.008) (0.141,0.706)rectangle++(0.008,0.008) (0.362,0.874)rectangle++(0.008,0.008)
  (0.412,0.804)rectangle++(0.008,0.008) (-0.005,0.816)rectangle++(0.008,0.008) (0.256,0.903)rectangle++(0.008,0.008) (0.233,0.913)rectangle++(0.008,0.008) (0.146,0.917)rectangle++(0.008,0.008)
  (0.147,0.864)rectangle++(0.008,0.008) (0.128,0.581)rectangle++(0.008,0.008) (0.323,0.783)rectangle++(0.008,0.008) (0.427,0.659)rectangle++(0.008,0.008) (0.133,0.903)rectangle++(0.008,0.008)
  (0.372,0.844)rectangle++(0.008,0.008) (-0.011,0.825)rectangle++(0.008,0.008) (0.143,0.887)rectangle++(0.008,0.008) (0.196,0.908)rectangle++(0.008,0.008) (0.145,0.604)rectangle++(0.008,0.008)
  (0.121,0.936)rectangle++(0.008,0.008) (0.389,0.801)rectangle++(0.008,0.008) (0.11,0.92)rectangle++(0.008,0.008) (0.378,0.798)rectangle++(0.008,0.008) (0.457,0.646)rectangle++(0.008,0.008)
  (0.439,0.549)rectangle++(0.008,0.008) (0.047,0.878)rectangle++(0.008,0.008) (0.334,0.801)rectangle++(0.008,0.008) (0.106,0.565)rectangle++(0.008,0.008) (0.183,0.799)rectangle++(0.008,0.008)
  (0.209,0.878)rectangle++(0.008,0.008) (0.383,0.828)rectangle++(0.008,0.008) (0.47,0.658)rectangle++(0.008,0.008) (0.272,0.909)rectangle++(0.008,0.008) (0.448,0.734)rectangle++(0.008,0.008)
  (0.36,0.886)rectangle++(0.008,0.008) (0.13,0.593)rectangle++(0.008,0.008) (0.146,0.86)rectangle++(0.008,0.008) (0.373,0.757)rectangle++(0.008,0.008) (0.44,0.629)rectangle++(0.008,0.008)
  (0.425,0.547)rectangle++(0.008,0.008) (0.29,0.793)rectangle++(0.008,0.008) (0.374,0.851)rectangle++(0.008,0.008) (0.273,0.896)rectangle++(0.008,0.008) (0.135,0.586)rectangle++(0.008,0.008)
  (0.27,0.636)rectangle++(0.008,0.008) (0.359,0.772)rectangle++(0.008,0.008) (0.396,0.785)rectangle++(0.008,0.008) (0.297,0.902)rectangle++(0.008,0.008) (0.247,0.913)rectangle++(0.008,0.008)
  (0.31,0.878)rectangle++(0.008,0.008) (0.318,0.895)rectangle++(0.008,0.008) (0.254,0.911)rectangle++(0.008,0.008) (0.441,0.742)rectangle++(0.008,0.008) (0.336,0.876)rectangle++(0.008,0.008)
  (0.11,0.574)rectangle++(0.008,0.008) (0.162,0.891)rectangle++(0.008,0.008) (0.392,0.766)rectangle++(0.008,0.008) (0.452,0.627)rectangle++(0.008,0.008) (0.429,0.541)rectangle++(0.008,0.008)
  (0.389,0.511)rectangle++(0.008,0.008) (0.135,0.598)rectangle++(0.008,0.008) (0.034,0.465)rectangle++(0.008,0.008) (0.159,0.766)rectangle++(0.008,0.008) (0.345,0.71)rectangle++(0.008,0.008)
  (0.367,0.88)rectangle++(0.008,0.008) (0.481,0.688)rectangle++(0.008,0.008) (0.017,0.799)rectangle++(0.008,0.008) (0.292,0.775)rectangle++(0.008,0.008) (0.188,0.92)rectangle++(0.008,0.008)
  (0.207,0.857)rectangle++(0.008,0.008) (0.217,0.897)rectangle++(0.008,0.008) (0.22,0.911)rectangle++(0.008,0.008) (0.229,0.86)rectangle++(0.008,0.008) (0.411,0.728)rectangle++(0.008,0.008)
  (0.4,0.769)rectangle++(0.008,0.008) (0.117,0.56)rectangle++(0.008,0.008) (0.253,0.63)rectangle++(0.008,0.008) (0.288,0.893)rectangle++(0.008,0.008) (0.45,0.722)rectangle++(0.008,0.008)
  (0.274,0.825)rectangle++(0.008,0.008) (0.24,0.887)rectangle++(0.008,0.008) (0.228,0.908)rectangle++(0.008,0.008) (0.391,0.83)rectangle++(0.008,0.008) (0.475,0.656)rectangle++(0.008,0.008)
  (0.307,0.83)rectangle++(0.008,0.008) (0.436,0.686)rectangle++(0.008,0.008) (0.444,0.574)rectangle++(0.008,0.008) (0.3,0.89)rectangle++(0.008,0.008) (0.316,0.897)rectangle++(0.008,0.008)
  (0.464,0.714)rectangle++(0.008,0.008) (0.466,0.577)rectangle++(0.008,0.008) (0.102,0.563)rectangle++(0.008,0.008) (0.345,0.844)rectangle++(0.008,0.008) (0.008,0.878)rectangle++(0.008,0.008)
  (0.15,0.905)rectangle++(0.008,0.008) (0.198,0.914)rectangle++(0.008,0.008) (-0.013,0.821)rectangle++(0.008,0.008) (0.089,0.885)rectangle++(0.008,0.008) (0.356,0.789)rectangle++(0.008,0.008)
  (0.324,0.883)rectangle++(0.008,0.008) (0.32,0.896)rectangle++(0.008,0.008) (0.465,0.712)rectangle++(0.008,0.008) (0.132,0.916)rectangle++(0.008,0.008) (0.387,0.788)rectangle++(0.008,0.008)
  (0.148,0.857)rectangle++(0.008,0.008) (0.373,0.755)rectangle++(0.008,0.008) (0.439,0.629)rectangle++(0.008,0.008) (0.424,0.547)rectangle++(0.008,0.008) (0.367,0.861)rectangle++(0.008,0.008)
  (0.131,0.578)rectangle++(0.008,0.008) (0.154,0.929)rectangle++(0.008,0.008) (0.401,0.786)rectangle++(0.008,0.008) (0.463,0.632)rectangle++(0.008,0.008) (0.339,0.861)rectangle++(0.008,0.008)
  (0.406,0.805)rectangle++(0.008,0.008) (0.131,0.949)rectangle++(0.008,0.008) (0.398,0.803)rectangle++(0.008,0.008) (0.468,0.642)rectangle++(0.008,0.008) (0.398,0.743)rectangle++(0.008,0.008)
  (0.446,0.614)rectangle++(0.008,0.008) (0.298,0.816)rectangle++(0.008,0.008) (0.094,0.912)rectangle++(0.008,0.008) (0.368,0.8)rectangle++(0.008,0.008) (0.271,0.878)rectangle++(0.008,0.008)
  (0.162,0.947)rectangle++(0.008,0.008) (0.374,0.872)rectangle++(0.008,0.008) (0.327,0.895)rectangle++(0.008,0.008) (0.468,0.709)rectangle++(0.008,0.008) (0.466,0.573)rectangle++(0.008,0.008)
  (0.114,0.575)rectangle++(0.008,0.008) (0.257,0.639)rectangle++(0.008,0.008) (0.345,0.617)rectangle++(0.008,0.008) (0.146,0.742)rectangle++(0.008,0.008) (0.331,0.703)rectangle++(0.008,0.008)
  (0.402,0.62)rectangle++(0.008,0.008) (0.282,0.818)rectangle++(0.008,0.008) (0.388,0.805)rectangle++(0.008,0.008) (0.383,0.846)rectangle++(0.008,0.008) (0.255,0.905)rectangle++(0.008,0.008)
  (0.233,0.914)rectangle++(0.008,0.008) (0.432,0.751)rectangle++(0.008,0.008) (0.465,0.605)rectangle++(0.008,0.008) (0.121,0.781)rectangle++(0.008,0.008) (0.188,0.872)rectangle++(0.008,0.008)
  (0.055,0.825)rectangle++(0.008,0.008) (0.319,0.774)rectangle++(0.008,0.008) (0.128,0.936)rectangle++(0.008,0.008) (0.392,0.798)rectangle++(0.008,0.008) (0.463,0.641)rectangle++(0.008,0.008)
  (0.103,0.801)rectangle++(0.008,0.008) (0.029,0.904)rectangle++(0.008,0.008) (0.336,0.819)rectangle++(0.008,0.008) (0.445,0.671)rectangle++(0.008,0.008) (0.442,0.564)rectangle++(0.008,0.008)
  (0.372,0.823)rectangle++(0.008,0.008) (0.463,0.66)rectangle++(0.008,0.008) (0.447,0.553)rectangle++(0.008,0.008) (0.263,0.763)rectangle++(0.008,0.008) (0.147,0.855)rectangle++(0.008,0.008)
  (0.24,0.893)rectangle++(0.008,0.008) (0.392,0.826)rectangle++(0.008,0.008) (0.473,0.654)rectangle++(0.008,0.008) (0.449,0.546)rectangle++(0.008,0.008) (0.4,0.506)rectangle++(0.008,0.008)
  (0.362,0.505)rectangle++(0.008,0.008) (0.146,0.739)rectangle++(0.008,0.008) (0.11,0.557)rectangle++(0.008,0.008) (0.319,0.83)rectangle++(0.008,0.008) (0.255,0.889)rectangle++(0.008,0.008)
  (0.433,0.732)rectangle++(0.008,0.008) (0.024,0.897)rectangle++(0.008,0.008) (0.351,0.859)rectangle++(0.008,0.008) (0.323,0.893)rectangle++(0.008,0.008) (0.466,0.709)rectangle++(0.008,0.008)
  (0.465,0.574)rectangle++(0.008,0.008) (0.146,0.86)rectangle++(0.008,0.008) (0.112,0.933)rectangle++(0.008,0.008) (0.272,0.916)rectangle++(0.008,0.008) (0.089,0.543)rectangle++(0.008,0.008)
  (0.234,0.633)rectangle++(0.008,0.008) (0.106,0.876)rectangle++(0.008,0.008) (0.183,0.904)rectangle++(0.008,0.008) (0.406,0.764)rectangle++(0.008,0.008) (-0.011,0.849)rectangle++(0.008,0.008)
  (0.092,0.551)rectangle++(0.008,0.008) (0.238,0.635)rectangle++(0.008,0.008) (0.227,0.823)rectangle++(0.008,0.008) (0.397,0.711)rectangle++(0.008,0.008) (0.395,0.767)rectangle++(0.008,0.008)
  (0.454,0.626)rectangle++(0.008,0.008) (0.097,0.532)rectangle++(0.008,0.008) (0.279,0.887)rectangle++(0.008,0.008) (0.241,0.908)rectangle++(0.008,0.008) (0.434,0.745)rectangle++(0.008,0.008)
  (0.464,0.602)rectangle++(0.008,0.008) (0.126,0.582)rectangle++(0.008,0.008) (0.264,0.638)rectangle++(0.008,0.008) (0.096,0.536)rectangle++(0.008,0.008) (0.287,0.847)rectangle++(0.008,0.008)
  (0.394,0.81)rectangle++(0.008,0.008) (0.28,0.882)rectangle++(0.008,0.008) (0.133,0.583)rectangle++(0.008,0.008) (0.084,0.54)rectangle++(0.008,0.008) (0.129,0.588)rectangle++(0.008,0.008)
  (0.165,0.893)rectangle++(0.008,0.008) (0.17,0.937)rectangle++(0.008,0.008) (0.411,0.784)rectangle++(0.008,0.008) (0.332,0.882)rectangle++(0.008,0.008) (0.259,0.906)rectangle++(0.008,0.008)
  (0.441,0.738)rectangle++(0.008,0.008) (0.125,0.917)rectangle++(0.008,0.008) (0.369,0.871)rectangle++(0.008,0.008) (0.479,0.683)rectangle++(0.008,0.008) (0.268,0.889)rectangle++(0.008,0.008);
\fill[piece3]
  (0.439,0.727)rectangle++(0.008,0.008) (0.172,0.938)rectangle++(0.008,0.008) (0.413,0.784)rectangle++(0.008,0.008) (0.088,0.541)rectangle++(0.008,0.008) (0.233,0.632)rectangle++(0.008,0.008)
  (0.057,0.895)rectangle++(0.008,0.008) (0.345,0.805)rectangle++(0.008,0.008) (0.444,0.661)rectangle++(0.008,0.008) (0.221,0.843)rectangle++(0.008,0.008) (0.104,0.805)rectangle++(0.008,0.008)
  (0.289,0.885)rectangle++(0.008,0.008) (-0.002,0.789)rectangle++(0.008,0.008) (0.28,0.778)rectangle++(0.008,0.008) (-0.015,0.844)rectangle++(0.008,0.008) (0.294,0.808)rectangle++(0.008,0.008)
  (0.11,0.574)rectangle++(0.008,0.008) (0.084,0.526)rectangle++(0.008,0.008) (0.077,0.517)rectangle++(0.008,0.008) (0.09,0.89)rectangle++(0.008,0.008) (0.358,0.791)rectangle++(0.008,0.008)
  (0.399,0.789)rectangle++(0.008,0.008) (0.463,0.634)rectangle++(0.008,0.008) (0.396,0.743)rectangle++(0.008,0.008) (0.28,0.859)rectangle++(0.008,0.008) (0.434,0.709)rectangle++(0.008,0.008)
  (0.109,0.548)rectangle++(0.008,0.008) (0.145,0.735)rectangle++(0.008,0.008) (0.109,0.556)rectangle++(0.008,0.008) (0.357,0.87)rectangle++(0.008,0.008) (0.252,0.905)rectangle++(0.008,0.008)
  (0.267,0.765)rectangle++(0.008,0.008) (0.237,0.867)rectangle++(0.008,0.008) (0.307,0.894)rectangle++(0.008,0.008) (0.459,0.715)rectangle++(0.008,0.008) (0.302,0.85)rectangle++(0.008,0.008)
  (0.316,0.891)rectangle++(0.008,0.008) (0.191,0.936)rectangle++(0.008,0.008) (0.285,0.911)rectangle++(0.008,0.008) (0.455,0.731)rectangle++(0.008,0.008) (0.128,0.917)rectangle++(0.008,0.008)
  (0.19,0.918)rectangle++(0.008,0.008) (0.414,0.768)rectangle++(0.008,0.008) (0.406,0.776)rectangle++(0.008,0.008) (0.144,0.606)rectangle++(0.008,0.008) (0.196,0.813)rectangle++(0.008,0.008)
  (0.379,0.718)rectangle++(0.008,0.008) (0.429,0.609)rectangle++(0.008,0.008) (0.094,0.529)rectangle++(0.008,0.008) (0.152,0.926)rectangle++(0.008,0.008) (0.138,0.593)rectangle++(0.008,0.008)
  (0.274,0.638)rectangle++(0.008,0.008) (0.352,0.61)rectangle++(0.008,0.008) (0.139,0.602)rectangle++(0.008,0.008) (0.278,0.642)rectangle++(0.008,0.008) (0.355,0.611)rectangle++(0.008,0.008)
  (0.197,0.897)rectangle++(0.008,0.008) (0.409,0.756)rectangle++(0.008,0.008) (0.118,0.937)rectangle++(0.008,0.008) (0.187,0.925)rectangle++(0.008,0.008) (0.3,0.902)rectangle++(0.008,0.008)
  (0.459,0.721)rectangle++(0.008,0.008) (0.302,0.852)rectangle++(0.008,0.008) (0.275,0.822)rectangle++(0.008,0.008) (0.163,0.944)rectangle++(0.008,0.008) (0.028,0.909)rectangle++(0.008,0.008)
  (0.148,0.949)rectangle++(0.008,0.008) (0.382,0.85)rectangle++(0.008,0.008) (0.276,0.895)rectangle++(0.008,0.008) (0.445,0.727)rectangle++(0.008,0.008) (0.462,0.59)rectangle++(0.008,0.008)
  (0.303,0.808)rectangle++(0.008,0.008) (0.426,0.677)rectangle++(0.008,0.008) (0.436,0.574)rectangle++(0.008,0.008) (0.403,0.523)rectangle++(0.008,0.008) (0.283,0.785)rectangle++(0.008,0.008)
  (0.409,0.674)rectangle++(0.008,0.008) (0.391,0.758)rectangle++(0.008,0.008) (0.4,0.777)rectangle++(0.008,0.008) (0.46,0.629)rectangle++(0.008,0.008) (0.434,0.54)rectangle++(0.008,0.008)
  (0.262,0.757)rectangle++(0.008,0.008) (0.389,0.669)rectangle++(0.008,0.008) (0.416,0.583)rectangle++(0.008,0.008) (0.398,0.534)rectangle++(0.008,0.008) (0.109,0.57)rectangle++(0.008,0.008)
  (0.147,0.689)rectangle++(0.008,0.008) (0.248,0.883)rectangle++(0.008,0.008) (0.428,0.731)rectangle++(0.008,0.008) (0.124,0.914)rectangle++(0.008,0.008) (-0.006,0.793)rectangle++(0.008,0.008)
  (0.28,0.781)rectangle++(0.008,0.008) (0.32,0.873)rectangle++(0.008,0.008) (0.255,0.903)rectangle++(0.008,0.008) (0.077,0.92)rectangle++(0.008,0.008) (0.363,0.809)rectangle++(0.008,0.008)
  (0.123,0.568)rectangle++(0.008,0.008) (0.291,0.852)rectangle++(0.008,0.008) (0.318,0.831)rectangle++(0.008,0.008) (0.281,0.912)rectangle++(0.008,0.008) (0.32,0.894)rectangle++(0.008,0.008)
  (0.272,0.826)rectangle++(0.008,0.008) (0.089,0.91)rectangle++(0.008,0.008) (0.365,0.852)rectangle++(0.008,0.008) (0.27,0.896)rectangle++(0.008,0.008) (0.443,0.729)rectangle++(0.008,0.008)
  (0.462,0.592)rectangle++(0.008,0.008) (0.422,0.522)rectangle++(0.008,0.008) (0.371,0.727)rectangle++(0.008,0.008) (0.428,0.616)rectangle++(0.008,0.008) (0.117,0.791)rectangle++(0.008,0.008)
  (0.335,0.736)rectangle++(0.008,0.008) (-0.004,0.826)rectangle++(0.008,0.008) (0.146,0.887)rectangle++(0.008,0.008) (0.336,0.892)rectangle++(0.008,0.008) (0.149,0.735)rectangle++(0.008,0.008)
  (0.33,0.699)rectangle++(0.008,0.008) (0.106,0.547)rectangle++(0.008,0.008) (0.333,0.872)rectangle++(0.008,0.008) (0.246,0.899)rectangle++(0.008,0.008) (0.433,0.74)rectangle++(0.008,0.008)
  (0.461,0.6)rectangle++(0.008,0.008) (0.151,0.693)rectangle++(0.008,0.008) (0.316,0.679)rectangle++(0.008,0.008) (0.113,0.893)rectangle++(0.008,0.008) (0.37,0.784)rectangle++(0.008,0.008)
  (0.448,0.643)rectangle++(0.008,0.008) (0.434,0.55)rectangle++(0.008,0.008) (0.12,0.575)rectangle++(0.008,0.008) (0.137,0.938)rectangle++(0.008,0.008) (0.397,0.796)rectangle++(0.008,0.008)
  (0.35,0.887)rectangle++(0.008,0.008) (0.196,0.936)rectangle++(0.008,0.008) (0.423,0.774)rectangle++(0.008,0.008) (0.146,0.865)rectangle++(0.008,0.008) (0.375,0.76)rectangle++(0.008,0.008)
  (0.146,0.663)rectangle++(0.008,0.008) (0.146,0.869)rectangle++(0.008,0.008) (0.376,0.761)rectangle++(0.008,0.008) (0.443,0.63)rectangle++(0.008,0.008) (0.102,0.805)rectangle++(0.008,0.008)
  (0.333,0.748)rectangle++(0.008,0.008) (0.388,0.784)rectangle++(0.008,0.008) (0.021,0.896)rectangle++(0.008,0.008) (0.329,0.818)rectangle++(0.008,0.008) (0.32,0.887)rectangle++(0.008,0.008)
  (0.089,0.925)rectangle++(0.008,0.008) (0.367,0.855)rectangle++(0.008,0.008) (0.472,0.676)rectangle++(0.008,0.008) (0.456,0.557)rectangle++(0.008,0.008) (0.407,0.508)rectangle++(0.008,0.008)
  (0.366,0.503)rectangle++(0.008,0.008) (0.346,0.516)rectangle++(0.008,0.008) (0.355,0.737)rectangle++(0.008,0.008) (0.425,0.627)rectangle++(0.008,0.008) (0.387,0.747)rectangle++(0.008,0.008)
  (0.443,0.62)rectangle++(0.008,0.008) (0.423,0.541)rectangle++(0.008,0.008) (0.247,0.732)rectangle++(0.008,0.008) (0.374,0.663)rectangle++(0.008,0.008) (0.327,0.846)rectangle++(0.008,0.008)
  (0.451,0.687)rectangle++(0.008,0.008) (0.13,0.589)rectangle++(0.008,0.008) (0.292,0.855)rectangle++(0.008,0.008) (0.274,0.822)rectangle++(0.008,0.008) (0.09,0.91)rectangle++(0.008,0.008)
  (0.177,0.915)rectangle++(0.008,0.008) (0.407,0.771)rectangle++(0.008,0.008) (0.15,0.718)rectangle++(0.008,0.008) (0.028,0.9)rectangle++(0.008,0.008) (0.157,0.912)rectangle++(0.008,0.008)
  (0.397,0.777)rectangle++(0.008,0.008) (0.458,0.63)rectangle++(0.008,0.008) (0.141,0.902)rectangle++(0.008,0.008) (0.386,0.778)rectangle++(0.008,0.008) (0.454,0.634)rectangle++(0.008,0.008)
  (0.038,0.882)rectangle++(0.008,0.008) (0.046,0.89)rectangle++(0.008,0.008) (0.338,0.807)rectangle++(0.008,0.008) (0.442,0.664)rectangle++(0.008,0.008) (0.102,0.54)rectangle++(0.008,0.008)
  (0.239,0.626)rectangle++(0.008,0.008) (0.351,0.775)rectangle++(0.008,0.008) (0.221,0.891)rectangle++(0.008,0.008) (0.274,0.878)rectangle++(0.008,0.008) (0.083,0.918)rectangle++(0.008,0.008)
  (0.365,0.807)rectangle++(0.008,0.008) (0.139,0.949)rectangle++(0.008,0.008) (0.194,0.929)rectangle++(0.008,0.008) (0.064,0.916)rectangle++(0.008,0.008) (0.15,0.949)rectangle++(0.008,0.008)
  (0.101,0.563)rectangle++(0.008,0.008) (0.315,0.783)rectangle++(0.008,0.008) (0.253,0.873)rectangle++(0.008,0.008) (0.427,0.725)rectangle++(0.008,0.008) (0.144,0.741)rectangle++(0.008,0.008)
  (0.261,0.904)rectangle++(0.008,0.008) (0.311,0.899)rectangle++(0.008,0.008) (0.462,0.716)rectangle++(0.008,0.008) (0.466,0.579)rectangle++(0.008,0.008) (0.419,0.514)rectangle++(0.008,0.008)
  (0.374,0.502)rectangle++(0.008,0.008) (0.273,0.778)rectangle++(0.008,0.008) (0.402,0.674)rectangle++(0.008,0.008) (0.39,0.76)rectangle++(0.008,0.008) (0.449,0.624)rectangle++(0.008,0.008)
  (0.427,0.542)rectangle++(0.008,0.008) (0.12,0.574)rectangle++(0.008,0.008) (-0.003,0.806)rectangle++(0.008,0.008) (0.332,0.845)rectangle++(0.008,0.008) (0.381,0.853)rectangle++(0.008,0.008)
  (0.324,0.834)rectangle++(0.008,0.008) (0.007,0.877)rectangle++(0.008,0.008) (0.315,0.815)rectangle++(0.008,0.008) (0.318,0.887)rectangle++(0.008,0.008) (0.406,0.813)rectangle++(0.008,0.008)
  (0.411,0.786)rectangle++(0.008,0.008) (0.286,0.874)rectangle++(0.008,0.008) (0.248,0.901)rectangle++(0.008,0.008) (0.231,0.913)rectangle++(0.008,0.008) (0.013,0.834)rectangle++(0.008,0.008)
  (0.019,0.896)rectangle++(0.008,0.008) (0.131,0.589)rectangle++(0.008,0.008) (0.126,0.936)rectangle++(0.008,0.008) (0.1,0.928)rectangle++(0.008,0.008) (0.106,0.935)rectangle++(0.008,0.008)
  (0.145,0.75)rectangle++(0.008,0.008) (0.143,0.604)rectangle++(0.008,0.008) (0.35,0.845)rectangle++(0.008,0.008) (0.323,0.898)rectangle++(0.008,0.008) (0.468,0.712)rectangle++(0.008,0.008)
  (0.213,0.911)rectangle++(0.008,0.008) (0.13,0.594)rectangle++(0.008,0.008) (0.271,0.641)rectangle++(0.008,0.008) (0.258,0.886)rectangle++(0.008,0.008) (0.005,0.867)rectangle++(0.008,0.008)
  (0.127,0.583)rectangle++(0.008,0.008) (0.265,0.638)rectangle++(0.008,0.008) (0.032,0.861)rectangle++(0.008,0.008) (-0.009,0.801)rectangle++(0.008,0.008) (0.145,0.923)rectangle++(0.008,0.008)
  (0.395,0.787)rectangle++(0.008,0.008) (0.461,0.635)rectangle++(0.008,0.008) (0.437,0.542)rectangle++(0.008,0.008) (0.392,0.508)rectangle++(0.008,0.008) (0.329,0.843)rectangle++(0.008,0.008)
  (0.451,0.684)rectangle++(0.008,0.008) (0.299,0.839)rectangle++(0.008,0.008) (0.217,0.881)rectangle++(0.008,0.008) (0.148,0.733)rectangle++(0.008,0.008) (0.1,0.808)rectangle++(0.008,0.008)
  (0.334,0.75)rectangle++(0.008,0.008) (0.135,0.59)rectangle++(0.008,0.008) (0.272,0.638)rectangle++(0.008,0.008) (0.359,0.772)rectangle++(0.008,0.008) (0.439,0.641)rectangle++(0.008,0.008)
  (0.429,0.552)rectangle++(0.008,0.008) (0.334,0.849)rectangle++(0.008,0.008) (0.403,0.803)rectangle++(0.008,0.008) (0.283,0.88)rectangle++(0.008,0.008) (0.084,0.92)rectangle++(0.008,0.008)
  (0.366,0.807)rectangle++(0.008,0.008) (0.123,0.568)rectangle++(0.008,0.008) (0.062,0.857)rectangle++(0.008,0.008) (0.155,0.931)rectangle++(0.008,0.008) (0.403,0.787)rectangle++(0.008,0.008)
  (0.283,0.874)rectangle++(0.008,0.008) (0.187,0.934)rectangle++(0.008,0.008) (0.418,0.777)rectangle++(0.008,0.008) (0.468,0.622)rectangle++(0.008,0.008) (0.305,0.819)rectangle++(0.008,0.008)
  (0.25,0.885)rectangle++(0.008,0.008) (0.429,0.732)rectangle++(0.008,0.008) (0.148,0.853)rectangle++(0.008,0.008) (0.372,0.754)rectangle++(0.008,0.008) (0.222,0.894)rectangle++(0.008,0.008)
  (0.388,0.829)rectangle++(0.008,0.008) (0.265,0.756)rectangle++(0.008,0.008) (0.236,0.863)rectangle++(0.008,0.008) (0.416,0.727)rectangle++(0.008,0.008) (0.449,0.6)rectangle++(0.008,0.008)
  (0.093,0.527)rectangle++(0.008,0.008) (0.04,0.839)rectangle++(0.008,0.008) (0.317,0.786)rectangle++(0.008,0.008) (0.254,0.874)rectangle++(0.008,0.008) (-0.006,0.792)rectangle++(0.008,0.008)
  (0.279,0.78)rectangle++(0.008,0.008) (0.313,0.882)rectangle++(0.008,0.008) (0.457,0.708)rectangle++(0.008,0.008) (0.301,0.847)rectangle++(0.008,0.008) (0.439,0.696)rectangle++(0.008,0.008)
  (0.295,0.843)rectangle++(0.008,0.008) (0.435,0.697)rectangle++(0.008,0.008) (0.208,0.909)rectangle++(0.008,0.008) (0.388,0.834)rectangle++(0.008,0.008) (0.141,0.597)rectangle++(0.008,0.008)
  (0.195,0.81)rectangle++(0.008,0.008) (0.213,0.882)rectangle++(0.008,0.008) (0.411,0.743)rectangle++(0.008,0.008) (0.286,0.859)rectangle++(0.008,0.008) (0.437,0.707)rectangle++(0.008,0.008)
  (0.294,0.847)rectangle++(0.008,0.008) (0.013,0.78)rectangle++(0.008,0.008) (0.284,0.768)rectangle++(0.008,0.008) (0.144,0.954)rectangle++(0.008,0.008) (0.074,0.53)rectangle++(0.008,0.008)
  (0.251,0.869)rectangle++(0.008,0.008) (0.376,0.86)rectangle++(0.008,0.008) (0.478,0.675)rectangle++(0.008,0.008) (0.348,0.867)rectangle++(0.008,0.008) (0.196,0.933)rectangle++(0.008,0.008)
  (0.213,0.923)rectangle++(0.008,0.008) (0.27,0.764)rectangle++(0.008,0.008) (0.093,0.928)rectangle++(0.008,0.008) (0.373,0.807)rectangle++(0.008,0.008) (0.372,0.871)rectangle++(0.008,0.008)
  (0.414,0.802)rectangle++(0.008,0.008) (-0.01,0.829)rectangle++(0.008,0.008) (0.146,0.944)rectangle++(0.008,0.008) (0.403,0.796)rectangle++(0.008,0.008) (0.237,0.899)rectangle++(0.008,0.008)
  (0.135,0.587)rectangle++(0.008,0.008) (0.293,0.854)rectangle++(0.008,0.008) (0.438,0.702)rectangle++(0.008,0.008) (0.38,0.834)rectangle++(0.008,0.008) (0.41,0.794)rectangle++(0.008,0.008)
  (0.285,0.876)rectangle++(0.008,0.008) (0.442,0.715)rectangle++(0.008,0.008) (0.066,0.91)rectangle++(0.008,0.008) (0.17,0.916)rectangle++(0.008,0.008) (0.404,0.774)rectangle++(0.008,0.008)
  (0.46,0.626)rectangle++(0.008,0.008) (0.097,0.532)rectangle++(0.008,0.008) (0.233,0.625)rectangle++(0.008,0.008) (0.329,0.619)rectangle++(0.008,0.008) (0.26,0.758)rectangle++(0.008,0.008)
  (0.389,0.67)rectangle++(0.008,0.008) (0.416,0.584)rectangle++(0.008,0.008) (0.171,0.929)rectangle++(0.008,0.008) (0.142,0.606)rectangle++(0.008,0.008) (0.195,0.813)rectangle++(0.008,0.008)
  (0.148,0.723)rectangle++(0.008,0.008) (0.295,0.874)rectangle++(0.008,0.008) (0.133,0.596)rectangle++(0.008,0.008) (0.125,0.915)rectangle++(0.008,0.008) (0.164,0.94)rectangle++(0.008,0.008)
  (0.331,0.898)rectangle++(0.008,0.008) (0.149,0.735)rectangle++(0.008,0.008) (0.006,0.785)rectangle++(0.008,0.008) (0.143,0.728)rectangle++(0.008,0.008) (0.324,0.698)rectangle++(0.008,0.008)
  (0.379,0.776)rectangle++(0.008,0.008) (0.02,0.877)rectangle++(0.008,0.008) (0.128,0.585)rectangle++(0.008,0.008) (0.086,0.528)rectangle++(0.008,0.008) (0.176,0.787)rectangle++(0.008,0.008)
  (0.207,0.874)rectangle++(0.008,0.008) (0.151,0.952)rectangle++(0.008,0.008) (0.408,0.798)rectangle++(0.008,0.008) (0.213,0.846)rectangle++(0.008,0.008) (0.398,0.727)rectangle++(0.008,0.008)
  (0.203,0.913)rectangle++(0.008,0.008) (0.387,0.835)rectangle++(0.008,0.008) (0.105,0.925)rectangle++(0.008,0.008) (0.371,0.852)rectangle++(0.008,0.008) (0.473,0.674)rectangle++(0.008,0.008)
  (0.242,0.897)rectangle++(0.008,0.008) (0.268,0.762)rectangle++(0.008,0.008) (0.238,0.866)rectangle++(0.008,0.008) (0.417,0.727)rectangle++(0.008,0.008) (0.449,0.6)rectangle++(0.008,0.008)
  (0.127,0.578)rectangle++(0.008,0.008) (0.013,0.437)rectangle++(0.008,0.008) (0.096,0.819)rectangle++(0.008,0.008) (0.335,0.757)rectangle++(0.008,0.008) (0.423,0.642)rectangle++(0.008,0.008)
  (0.422,0.559)rectangle++(0.008,0.008) (0.289,0.797)rectangle++(0.008,0.008) (0.387,0.8)rectangle++(0.008,0.008) (0.349,0.887)rectangle++(0.008,0.008) (0.412,0.808)rectangle++(0.008,0.008)
  (0.273,0.794)rectangle++(0.008,0.008) (0.12,0.566)rectangle++(0.008,0.008) (0.188,0.8)rectangle++(0.008,0.008) (0.371,0.715)rectangle++(0.008,0.008) (0.272,0.85)rectangle++(0.008,0.008)
  (0.01,0.885)rectangle++(0.008,0.008) (0.32,0.817)rectangle++(0.008,0.008) (0.124,0.57)rectangle++(0.008,0.008) (0.27,0.767)rectangle++(0.008,0.008) (0.238,0.867)rectangle++(0.008,0.008)
  (0.273,0.829)rectangle++(0.008,0.008) (0.217,0.88)rectangle++(0.008,0.008) (0.099,0.921)rectangle++(0.008,0.008) (0.146,0.741)rectangle++(0.008,0.008) (0.322,0.835)rectangle++(0.008,0.008)
  (0.088,0.921)rectangle++(0.008,0.008) (0.136,0.593)rectangle++(0.008,0.008) (0.273,0.639)rectangle++(0.008,0.008) (0.352,0.611)rectangle++(0.008,0.008) (0.378,0.57)rectangle++(0.008,0.008)
  (0.376,0.542)rectangle++(0.008,0.008) (0.262,0.846)rectangle++(0.008,0.008) (0.028,0.816)rectangle++(0.008,0.008) (0.107,0.924)rectangle++(0.008,0.008) (0.378,0.8)rectangle++(0.008,0.008)
  (0.15,0.722)rectangle++(0.008,0.008) (0.325,0.693)rectangle++(0.008,0.008) (0.371,0.841)rectangle++(0.008,0.008) (0.384,0.85)rectangle++(0.008,0.008) (0.328,0.891)rectangle++(0.008,0.008)
  (0.468,0.707)rectangle++(0.008,0.008) (0.361,0.886)rectangle++(0.008,0.008) (0.315,0.891)rectangle++(0.008,0.008) (0.461,0.711)rectangle++(0.008,0.008) (0.464,0.577)rectangle++(0.008,0.008)
  (0.084,0.892)rectangle++(0.008,0.008) (0.266,0.913)rectangle++(0.008,0.008) (0.447,0.738)rectangle++(0.008,0.008) (0.143,0.715)rectangle++(0.008,0.008) (0.167,0.911)rectangle++(0.008,0.008)
  (0.203,0.916)rectangle++(0.008,0.008) (0.419,0.762)rectangle++(0.008,0.008) (0.463,0.615)rectangle++(0.008,0.008) (0.117,0.791)rectangle++(0.008,0.008) (0.116,0.786)rectangle++(0.008,0.008)
  (0.333,0.734)rectangle++(0.008,0.008) (0.414,0.633)rectangle++(0.008,0.008) (0.332,0.861)rectangle++(0.008,0.008) (0.458,0.691)rectangle++(0.008,0.008) (0.134,0.912)rectangle++(0.008,0.008)
  (0.271,0.769)rectangle++(0.008,0.008) (0.185,0.919)rectangle++(0.008,0.008) (0.385,0.839)rectangle++(0.008,0.008) (0.411,0.794)rectangle++(0.008,0.008) (0.135,0.596)rectangle++(0.008,0.008)
  (0.21,0.859)rectangle++(0.008,0.008) (0.218,0.898)rectangle++(0.008,0.008) (0.133,0.885)rectangle++(0.008,0.008) (0.376,0.773)rectangle++(0.008,0.008) (0.283,0.911)rectangle++(0.008,0.008)
  (0.08,0.924)rectangle++(0.008,0.008) (0.285,0.902)rectangle++(0.008,0.008) (0.243,0.913)rectangle++(0.008,0.008) (0.436,0.747)rectangle++(0.008,0.008) (0.407,0.768)rectangle++(0.008,0.008)
  (0.15,0.717)rectangle++(0.008,0.008) (0.198,0.851)rectangle++(0.008,0.008) (0.125,0.588)rectangle++(0.008,0.008) (0.099,0.932)rectangle++(0.008,0.008) (0.37,0.854)rectangle++(0.008,0.008)
  (0.474,0.675)rectangle++(0.008,0.008) (0.138,0.911)rectangle++(0.008,0.008) (0.375,0.845)rectangle++(0.008,0.008) (0.273,0.893)rectangle++(0.008,0.008) (0.443,0.727)rectangle++(0.008,0.008)
  (0.461,0.59)rectangle++(0.008,0.008) (0.421,0.521)rectangle++(0.008,0.008) (0.378,0.504)rectangle++(0.008,0.008) (0.118,0.582)rectangle++(0.008,0.008) (0.248,0.737)rectangle++(0.008,0.008)
  (0.37,0.795)rectangle++(0.008,0.008) (0.011,0.86)rectangle++(0.008,0.008) (0.141,0.601)rectangle++(0.008,0.008) (0.329,0.785)rectangle++(0.008,0.008) (0.141,0.699)rectangle++(0.008,0.008)
  (0.313,0.685)rectangle++(0.008,0.008) (0.387,0.618)rectangle++(0.008,0.008) (0.131,0.892)rectangle++(0.008,0.008) (0.378,0.777)rectangle++(0.008,0.008) (0.449,0.637)rectangle++(0.008,0.008)
  (0.432,0.547)rectangle++(0.008,0.008) (0.293,0.793)rectangle++(0.008,0.008) (0.374,0.851)rectangle++(0.008,0.008) (0.273,0.896)rectangle++(0.008,0.008) (0.444,0.728)rectangle++(0.008,0.008)
  (0.125,0.584)rectangle++(0.008,0.008) (0.265,0.639)rectangle++(0.008,0.008) (0.362,0.84)rectangle++(0.008,0.008) (0.055,0.915)rectangle++(0.008,0.008) (0.351,0.815)rectangle++(0.008,0.008)
  (0.451,0.663)rectangle++(0.008,0.008) (0.136,0.906)rectangle++(0.008,0.008) (0.052,0.488)rectangle++(0.008,0.008) (0.112,0.572)rectangle++(0.008,0.008) (0.147,0.689)rectangle++(0.008,0.008)
  (0.313,0.679)rectangle++(0.008,0.008) (0.252,0.838)rectangle++(0.008,0.008) (0.414,0.709)rectangle++(0.008,0.008) (0.287,0.848)rectangle++(0.008,0.008) (0.433,0.702)rectangle++(0.008,0.008)
  (0.24,0.897)rectangle++(0.008,0.008) (0.107,0.57)rectangle++(0.008,0.008) (0.162,0.891)rectangle++(0.008,0.008) (0.054,0.825)rectangle++(0.008,0.008) (0.055,0.884)rectangle++(0.008,0.008)
  (0.105,0.929)rectangle++(0.008,0.008) (0.106,0.935)rectangle++(0.008,0.008) (0.382,0.806)rectangle++(0.008,0.008) (0.108,0.92)rectangle++(0.008,0.008) (0.184,0.919)rectangle++(0.008,0.008)
  (0.284,0.91)rectangle++(0.008,0.008) (0.243,0.916)rectangle++(0.008,0.008) (0.054,0.49)rectangle++(0.008,0.008) (0.146,0.907)rectangle++(0.008,0.008) (0.196,0.914)rectangle++(0.008,0.008)
  (0.174,0.94)rectangle++(0.008,0.008) (0.337,0.883)rectangle++(0.008,0.008) (0.469,0.699)rectangle++(0.008,0.008) (0.463,0.569)rectangle++(0.008,0.008) (0.322,0.871)rectangle++(0.008,0.008)
  (0.404,0.809)rectangle++(0.008,0.008) (0.037,0.911)rectangle++(0.008,0.008) (0.042,0.895)rectangle++(0.008,0.008) (0.035,0.873)rectangle++(0.008,0.008) (0.104,0.928)rectangle++(0.008,0.008)
  (0.147,0.74)rectangle++(0.008,0.008) (0.331,0.702)rectangle++(0.008,0.008) (0.38,0.776)rectangle++(0.008,0.008) (0.018,0.877)rectangle++(0.008,0.008) (0.143,0.73)rectangle++(0.008,0.008)
  (0.325,0.699)rectangle++(0.008,0.008) (0.106,0.547)rectangle++(0.008,0.008) (0.106,0.565)rectangle++(0.008,0.008) (0.13,0.59)rectangle++(0.008,0.008) (0.325,0.784)rectangle++(0.008,0.008)
  (0.428,0.659)rectangle++(0.008,0.008) (0.43,0.565)rectangle++(0.008,0.008) (0.397,0.521)rectangle++(0.008,0.008) (0.08,0.512)rectangle++(0.008,0.008) (0.158,0.883)rectangle++(0.008,0.008)
  (0.131,0.584)rectangle++(0.008,0.008) (0.191,0.806)rectangle++(0.008,0.008) (0.174,0.924)rectangle++(0.008,0.008) (0.384,0.842)rectangle++(0.008,0.008) (0.143,0.953)rectangle++(0.008,0.008)
  (0.382,0.852)rectangle++(0.008,0.008) (0.478,0.669)rectangle++(0.008,0.008) (0.035,0.885)rectangle++(0.008,0.008) (0.115,0.785)rectangle++(0.008,0.008) (0.097,0.934)rectangle++(0.008,0.008)
  (0.273,0.77)rectangle++(0.008,0.008) (-0.003,0.852)rectangle++(0.008,0.008) (0.146,0.896)rectangle++(0.008,0.008) (0.386,0.774)rectangle++(0.008,0.008) (0.054,0.49)rectangle++(0.008,0.008)
  (0.136,0.763)rectangle++(0.008,0.008) (0.352,0.816)rectangle++(0.008,0.008) (0.381,0.849)rectangle++(0.008,0.008) (0.476,0.669)rectangle++(0.008,0.008) (0.456,0.552)rectangle++(0.008,0.008)
  (0.232,0.885)rectangle++(0.008,0.008) (0.044,0.911)rectangle++(0.008,0.008) (0.357,0.859)rectangle++(0.008,0.008) (0.469,0.681)rectangle++(0.008,0.008) (0.457,0.56)rectangle++(0.008,0.008)
  (0.408,0.509)rectangle++(0.008,0.008) (0.245,0.73)rectangle++(0.008,0.008) (0.372,0.662)rectangle++(0.008,0.008) (0.406,0.586)rectangle++(0.008,0.008) (0.394,0.54)rectangle++(0.008,0.008)
  (0.372,0.522)rectangle++(0.008,0.008) (0.272,0.785)rectangle++(0.008,0.008) (0.358,0.879)rectangle++(0.008,0.008) (0.268,0.905)rectangle++(0.008,0.008) (0.208,0.874)rectangle++(0.008,0.008)
  (0.406,0.741)rectangle++(0.008,0.008) (0.45,0.61)rectangle++(0.008,0.008) (0.229,0.885)rectangle++(0.008,0.008) (0.42,0.739)rectangle++(0.008,0.008) (0.455,0.604)rectangle++(0.008,0.008)
  (0.094,0.528)rectangle++(0.008,0.008) (0.029,0.879)rectangle++(0.008,0.008) (0.349,0.855)rectangle++(0.008,0.008) (0.256,0.854)rectangle++(0.008,0.008) (0.421,0.715)rectangle++(0.008,0.008)
  (0.447,0.593)rectangle++(0.008,0.008) (0.016,0.889)rectangle++(0.008,0.008) (0.324,0.817)rectangle++(0.008,0.008) (0.256,0.884)rectangle++(0.008,0.008) (0.327,0.849)rectangle++(0.008,0.008)
  (0.257,0.895)rectangle++(0.008,0.008) (0.141,0.707)rectangle++(0.008,0.008) (0.022,0.898)rectangle++(0.008,0.008) (0.277,0.898)rectangle++(0.008,0.008) (0.186,0.937)rectangle++(0.008,0.008)
  (0.388,0.842)rectangle++(0.008,0.008) (0.136,0.599)rectangle++(0.008,0.008) (0.323,0.893)rectangle++(0.008,0.008) (0.256,0.91)rectangle++(0.008,0.008) (0.084,0.538)rectangle++(0.008,0.008)
  (0.094,0.843)rectangle++(0.008,0.008) (0.096,0.544)rectangle++(0.008,0.008) (0.135,0.884)rectangle++(0.008,0.008) (0.149,0.952)rectangle++(0.008,0.008) (0.092,0.923)rectangle++(0.008,0.008)
  (0.368,0.854)rectangle++(0.008,0.008) (0.472,0.675)rectangle++(0.008,0.008) (0.223,0.893)rectangle++(0.008,0.008) (-0.01,0.815)rectangle++(0.008,0.008) (0.286,0.793)rectangle++(0.008,0.008)
  (0.328,0.9)rectangle++(0.008,0.008) (0.122,0.944)rectangle++(0.008,0.008) (-0.012,0.799)rectangle++(0.008,0.008) (0.279,0.786)rectangle++(0.008,0.008) (0.408,0.676)rectangle++(0.008,0.008)
  (0.102,0.807)rectangle++(0.008,0.008) (0.333,0.887)rectangle++(0.008,0.008) (0.321,0.897)rectangle++(0.008,0.008) (0.088,0.927)rectangle++(0.008,0.008) (0.177,0.921)rectangle++(0.008,0.008)
  (0.409,0.774)rectangle++(0.008,0.008) (0.133,0.94)rectangle++(0.008,0.008) (0.396,0.799)rectangle++(0.008,0.008) (0.465,0.64)rectangle++(0.008,0.008) (0.072,0.907)rectangle++(0.008,0.008)
  (0.171,0.914)rectangle++(0.008,0.008) (-0.001,0.858)rectangle++(0.008,0.008) (0.147,0.898)rectangle++(0.008,0.008) (0.133,0.588)rectangle++(0.008,0.008) (0.293,0.854)rectangle++(0.008,0.008)
  (0.438,0.702)rectangle++(0.008,0.008) (0.172,0.936)rectangle++(0.008,0.008) (0.412,0.783)rectangle++(0.008,0.008) (0.467,0.627)rectangle++(0.008,0.008) (0.437,0.536)rectangle++(0.008,0.008)
  (0.39,0.505)rectangle++(0.008,0.008) (0.279,0.779)rectangle++(0.008,0.008) (-0.005,0.821)rectangle++(0.008,0.008) (0.232,0.883)rectangle++(0.008,0.008) (0.421,0.737)rectangle++(0.008,0.008)
  (0.113,0.554)rectangle++(0.008,0.008) (0.249,0.629)rectangle++(0.008,0.008) (0.161,0.932)rectangle++(0.008,0.008) (0.406,0.785)rectangle++(0.008,0.008) (0.284,0.873)rectangle++(0.008,0.008)
  (0.441,0.714)rectangle++(0.008,0.008) (0.066,0.91)rectangle++(0.008,0.008) (0.355,0.809)rectangle++(0.008,0.008) (0.322,0.835)rectangle++(0.008,0.008) (0.256,0.89)rectangle++(0.008,0.008)
  (0.434,0.732)rectangle++(0.008,0.008) (0.264,0.9)rectangle++(0.008,0.008) (0.397,0.823)rectangle++(0.008,0.008) (0.13,0.943)rectangle++(0.008,0.008) (0.165,0.944)rectangle++(0.008,0.008)
  (0.412,0.789)rectangle++(0.008,0.008) (0.469,0.63)rectangle++(0.008,0.008) (0.34,0.86)rectangle++(0.008,0.008) (0.462,0.688)rectangle++(0.008,0.008) (0.26,0.758)rectangle++(0.008,0.008)
  (0.067,0.85)rectangle++(0.008,0.008) (0.283,0.892)rectangle++(0.008,0.008) (0.134,0.585)rectangle++(0.008,0.008) (0.096,0.856)rectangle++(0.008,0.008) (0.18,0.897)rectangle++(0.008,0.008)
  (0.055,0.866)rectangle++(0.008,0.008) (0.33,0.891)rectangle++(0.008,0.008) (0.468,0.706)rectangle++(0.008,0.008) (0.465,0.572)rectangle++(0.008,0.008) (0.147,0.93)rectangle++(0.008,0.008)
  (0.399,0.789)rectangle++(0.008,0.008) (0.383,0.844)rectangle++(0.008,0.008) (0.476,0.665)rectangle++(0.008,0.008) (0.127,0.773)rectangle++(0.008,0.008) (0.334,0.724)rectangle++(0.008,0.008)
  (0.41,0.628)rectangle++(0.008,0.008) (0.285,0.821)rectangle++(0.008,0.008) (0.091,0.911)rectangle++(0.008,0.008) (0.178,0.916)rectangle++(0.008,0.008) (0.408,0.771)rectangle++(0.008,0.008)
  (0.331,0.88)rectangle++(0.008,0.008) (0.465,0.7)rectangle++(0.008,0.008) (0.406,0.755)rectangle++(0.008,0.008) (0.144,0.85)rectangle++(0.008,0.008) (0.367,0.832)rectangle++(0.008,0.008)
  (0.126,0.572)rectangle++(0.008,0.008) (0.261,0.633)rectangle++(0.008,0.008) (0.005,0.783)rectangle++(0.008,0.008) (0.281,0.772)rectangle++(0.008,0.008) (0.403,0.669)rectangle++(0.008,0.008)
  (0.283,0.834)rectangle++(0.008,0.008) (0.427,0.697)rectangle++(0.008,0.008) (0.291,0.844)rectangle++(0.008,0.008) (0.022,0.876)rectangle++(0.008,0.008) (0.322,0.809)rectangle++(0.008,0.008)
  (0.435,0.671)rectangle++(0.008,0.008) (0.107,0.57)rectangle++(0.008,0.008) (0.31,0.888)rectangle++(0.008,0.008) (0.325,0.835)rectangle++(0.008,0.008) (0.096,0.918)rectangle++(0.008,0.008)
  (0.371,0.802)rectangle++(0.008,0.008) (0.13,0.94)rectangle++(0.008,0.008) (0.394,0.8)rectangle++(0.008,0.008) (0.465,0.641)rectangle++(0.008,0.008) (0.304,0.825)rectangle++(0.008,0.008)
  (0.25,0.887)rectangle++(0.008,0.008) (0.078,0.917)rectangle++(0.008,0.008) (0.018,0.78)rectangle++(0.008,0.008) (0.091,0.913)rectangle++(0.008,0.008) (0.366,0.852)rectangle++(0.008,0.008)
  (0.271,0.896)rectangle++(0.008,0.008) (0.312,0.898)rectangle++(0.008,0.008) (0.406,0.816)rectangle++(0.008,0.008) (0.412,0.786)rectangle++(0.008,0.008) (0.115,0.921)rectangle++(0.008,0.008)
  (0.381,0.796)rectangle++(0.008,0.008) (0.146,0.748)rectangle++(0.008,0.008) (0.333,0.706)rectangle++(0.008,0.008) (0.381,0.776)rectangle++(0.008,0.008) (0.401,0.782)rectangle++(0.008,0.008)
  (0.462,0.631)rectangle++(0.008,0.008) (0.141,0.903)rectangle++(0.008,0.008) (0.125,0.589)rectangle++(0.008,0.008) (0.324,0.785)rectangle++(0.008,0.008) (0.428,0.659)rectangle++(0.008,0.008)
  (0.43,0.565)rectangle++(0.008,0.008) (0.335,0.851)rectangle++(0.008,0.008) (0.456,0.686)rectangle++(0.008,0.008) (0.402,0.754)rectangle++(0.008,0.008) (0.402,0.775)rectangle++(0.008,0.008)
  (0.46,0.627)rectangle++(0.008,0.008) (0.394,0.742)rectangle++(0.008,0.008) (0.444,0.615)rectangle++(0.008,0.008) (0.141,0.898)rectangle++(0.008,0.008) (0.373,0.842)rectangle++(0.008,0.008)
  (0.47,0.668)rectangle++(0.008,0.008) (0.306,0.834)rectangle++(0.008,0.008) (0.395,0.805)rectangle++(0.008,0.008) (0.408,0.786)rectangle++(0.008,0.008) (0.466,0.63)rectangle++(0.008,0.008)
  (0.437,0.538)rectangle++(0.008,0.008) (0.391,0.506)rectangle++(0.008,0.008) (0.364,0.822)rectangle++(0.008,0.008) (0.459,0.662)rectangle++(0.008,0.008) (0.129,0.587)rectangle++(0.008,0.008)
  (0.122,0.779)rectangle++(0.008,0.008) (0.242,0.887)rectangle++(0.008,0.008) (0.083,0.531)rectangle++(0.008,0.008) (0.244,0.733)rectangle++(0.008,0.008) (0.251,0.889)rectangle++(0.008,0.008)
  (0.377,0.861)rectangle++(0.008,0.008) (0.479,0.675)rectangle++(0.008,0.008) (0.143,0.604)rectangle++(0.008,0.008) (0.28,0.642)rectangle++(0.008,0.008) (0.356,0.61)rectangle++(0.008,0.008)
  (0.38,0.568)rectangle++(0.008,0.008) (0.139,0.929)rectangle++(0.008,0.008) (0.378,0.848)rectangle++(0.008,0.008) (0.102,0.926)rectangle++(0.008,0.008) (0.377,0.803)rectangle++(0.008,0.008)
  (0.108,0.919)rectangle++(0.008,0.008) (0.107,0.934)rectangle++(0.008,0.008) (0.382,0.805)rectangle++(0.008,0.008) (0.336,0.895)rectangle++(0.008,0.008) (0.135,0.585)rectangle++(0.008,0.008)
  (0.293,0.854)rectangle++(0.008,0.008) (0.246,0.897)rectangle++(0.008,0.008) (0.23,0.911)rectangle++(0.008,0.008) (0.306,0.9)rectangle++(0.008,0.008) (0.136,0.586)rectangle++(0.008,0.008)
  (0.248,0.737)rectangle++(0.008,0.008) (0.37,0.795)rectangle++(0.008,0.008) (0.124,0.586)rectangle++(0.008,0.008) (0.291,0.854)rectangle++(0.008,0.008) (0.437,0.703)rectangle++(0.008,0.008)
  (0.081,0.538)rectangle++(0.008,0.008) (0.285,0.847)rectangle++(0.008,0.008) (0.148,0.611)rectangle++(0.008,0.008) (0.332,0.786)rectangle++(0.008,0.008) (0.026,0.815)rectangle++(0.008,0.008)
  (0.091,0.924)rectangle++(0.008,0.008) (0.368,0.854)rectangle++(0.008,0.008) (0.271,0.897)rectangle++(0.008,0.008) (0.443,0.729)rectangle++(0.008,0.008) (0.042,0.475)rectangle++(0.008,0.008)
  (0.145,0.68)rectangle++(0.008,0.008) (0.196,0.838)rectangle++(0.008,0.008) (0.369,0.861)rectangle++(0.008,0.008) (0.475,0.678)rectangle++(0.008,0.008) (0.09,0.905)rectangle++(0.008,0.008)
  (0.365,0.851)rectangle++(0.008,0.008) (0.373,0.873)rectangle++(0.008,0.008) (0.273,0.796)rectangle++(0.008,0.008) (0.239,0.877)rectangle++(0.008,0.008) (0.146,0.63)rectangle++(0.008,0.008)
  (0.294,0.86)rectangle++(0.008,0.008) (0.164,0.946)rectangle++(0.008,0.008) (0.156,0.951)rectangle++(0.008,0.008) (0.384,0.85)rectangle++(0.008,0.008) (0.478,0.668)rectangle++(0.008,0.008)
  (0.103,0.54)rectangle++(0.008,0.008) (0.239,0.626)rectangle++(0.008,0.008) (0.18,0.899)rectangle++(0.008,0.008) (0.134,0.587)rectangle++(0.008,0.008) (0.27,0.637)rectangle++(0.008,0.008)
  (0.266,0.76)rectangle++(0.008,0.008) (0.377,0.796)rectangle++(0.008,0.008) (0.274,0.877)rectangle++(0.008,0.008) (0.054,0.49)rectangle++(0.008,0.008) (0.154,0.88)rectangle++(0.008,0.008)
  (0.384,0.763)rectangle++(0.008,0.008) (0.012,0.887)rectangle++(0.008,0.008) (0.274,0.908)rectangle++(0.008,0.008) (0.263,0.911)rectangle++(0.008,0.008) (0.444,0.739)rectangle++(0.008,0.008)
  (0.466,0.595)rectangle++(0.008,0.008) (0.134,0.886)rectangle++(0.008,0.008) (0.377,0.773)rectangle++(0.008,0.008) (0.4,0.782)rectangle++(0.008,0.008) (0.277,0.914)rectangle++(0.008,0.008)
  (0.27,0.763)rectangle++(0.008,0.008) (0.395,0.669)rectangle++(0.008,0.008) (0.28,0.834)rectangle++(0.008,0.008) (0.263,0.91)rectangle++(0.008,0.008) (0.053,0.868)rectangle++(0.008,0.008)
  (0.212,0.87)rectangle++(0.008,0.008) (0.015,0.889)rectangle++(0.008,0.008) (0.147,0.948)rectangle++(0.008,0.008) (0.167,0.944)rectangle++(0.008,0.008) (0.413,0.788)rectangle++(0.008,0.008)
  (0.469,0.629)rectangle++(0.008,0.008) (0.396,0.741)rectangle++(0.008,0.008) (0.118,0.912)rectangle++(0.008,0.008) (0.379,0.791)rectangle++(0.008,0.008) (0.08,0.818)rectangle++(0.008,0.008)
  (0.328,0.762)rectangle++(0.008,0.008) (0.193,0.918)rectangle++(0.008,0.008) (0.137,0.591)rectangle++(0.008,0.008) (0.273,0.638)rectangle++(0.008,0.008) (-0.002,0.792)rectangle++(0.008,0.008)
  (0.281,0.779)rectangle++(0.008,0.008) (0.125,0.934)rectangle++(0.008,0.008) (0.146,0.651)rectangle++(0.008,0.008) (0.106,0.929)rectangle++(0.008,0.008) (0.111,0.563)rectangle++(0.008,0.008)
  (0.142,0.711)rectangle++(0.008,0.008) (0.021,0.849)rectangle++(0.008,0.008) (-0.009,0.841)rectangle++(0.008,0.008) (0.146,0.727)rectangle++(0.008,0.008) (0.326,0.697)rectangle++(0.008,0.008)
  (0.11,0.571)rectangle++(0.008,0.008) (0.289,0.852)rectangle++(0.008,0.008) (0.01,0.869)rectangle++(0.008,0.008) (0.133,0.594)rectangle++(0.008,0.008) (0.272,0.641)rectangle++(0.008,0.008)
  (0.36,0.773)rectangle++(0.008,0.008) (0.109,0.937)rectangle++(0.008,0.008) (0.235,0.722)rectangle++(0.008,0.008) (0.301,0.872)rectangle++(0.008,0.008) (0.027,0.816)rectangle++(0.008,0.008)
  (0.303,0.78)rectangle++(0.008,0.008) (0.31,0.892)rectangle++(0.008,0.008) (0.46,0.713)rectangle++(0.008,0.008) (0.464,0.578)rectangle++(0.008,0.008) (0.147,0.93)rectangle++(0.008,0.008)
  (0.236,0.898)rectangle++(0.008,0.008) (0.137,0.946)rectangle++(0.008,0.008) (0.4,0.8)rectangle++(0.008,0.008) (0.467,0.639)rectangle++(0.008,0.008) (0.099,0.534)rectangle++(0.008,0.008)
  (0.235,0.625)rectangle++(0.008,0.008) (0.33,0.618)rectangle++(0.008,0.008) (0.019,0.833)rectangle++(0.008,0.008) (0.214,0.868)rectangle++(0.008,0.008) (0.407,0.737)rectangle++(0.008,0.008)
  (0.372,0.868)rectangle++(0.008,0.008) (0.479,0.68)rectangle++(0.008,0.008) (0.461,0.556)rectangle++(0.008,0.008) (0.409,0.506)rectangle++(0.008,0.008) (0.146,0.74)rectangle++(0.008,0.008)
  (0.258,0.75)rectangle++(0.008,0.008) (0.385,0.667)rectangle++(0.008,0.008) (0.413,0.584)rectangle++(0.008,0.008) (0.397,0.536)rectangle++(0.008,0.008) (0.372,0.52)rectangle++(0.008,0.008)
  (0.361,0.734)rectangle++(0.008,0.008) (0.095,0.922)rectangle++(0.008,0.008) (0.05,0.903)rectangle++(0.008,0.008) (0.357,0.856)rectangle++(0.008,0.008) (0.468,0.68)rectangle++(0.008,0.008)
  (0.142,0.937)rectangle++(0.008,0.008) (0.399,0.794)rectangle++(0.008,0.008) (0.108,0.802)rectangle++(0.008,0.008) (0.335,0.745)rectangle++(0.008,0.008) (0.108,0.905)rectangle++(0.008,0.008)
  (0.272,0.769)rectangle++(0.008,0.008) (0.371,0.85)rectangle++(0.008,0.008) (0.071,0.513)rectangle++(0.008,0.008) (0.215,0.626)rectangle++(0.008,0.008) (0.321,0.626)rectangle++(0.008,0.008)
  (0.369,0.588)rectangle++(0.008,0.008) (0.146,0.915)rectangle++(0.008,0.008) (0.393,0.783)rectangle++(0.008,0.008) (0.108,0.801)rectangle++(0.008,0.008) (0.289,0.885)rectangle++(0.008,0.008)
  (0.048,0.914)rectangle++(0.008,0.008) (0.141,0.707)rectangle++(0.008,0.008) (0.119,0.922)rectangle++(0.008,0.008) (0.081,0.894)rectangle++(0.008,0.008) (0.319,0.874)rectangle++(0.008,0.008)
  (0.346,0.891)rectangle++(0.008,0.008) (0.264,0.909)rectangle++(0.008,0.008) (0.236,0.915)rectangle++(0.008,0.008) (0.394,0.831)rectangle++(0.008,0.008) (0.066,0.507)rectangle++(0.008,0.008)
  (0.211,0.624)rectangle++(0.008,0.008) (0.223,0.888)rectangle++(0.008,0.008) (0.305,0.897)rectangle++(0.008,0.008) (0.095,0.556)rectangle++(0.008,0.008) (0.241,0.637)rectangle++(0.008,0.008)
  (0.353,0.776)rectangle++(0.008,0.008) (0.438,0.645)rectangle++(0.008,0.008) (0.251,0.748)rectangle++(0.008,0.008) (0.182,0.916)rectangle++(0.008,0.008) (0.267,0.914)rectangle++(0.008,0.008)
  (0.344,0.891)rectangle++(0.008,0.008) (0.079,0.894)rectangle++(0.008,0.008) (0.355,0.797)rectangle++(0.008,0.008) (0.155,0.95)rectangle++(0.008,0.008) (0.409,0.795)rectangle++(0.008,0.008)
  (0.002,0.786)rectangle++(0.008,0.008) (0.281,0.775)rectangle++(0.008,0.008) (0.242,0.87)rectangle++(0.008,0.008) (0.229,0.902)rectangle++(0.008,0.008) (0.257,0.755)rectangle++(0.008,0.008)
  (0.386,0.669)rectangle++(0.008,0.008) (0.032,0.875)rectangle++(0.008,0.008) (0.247,0.739)rectangle++(0.008,0.008) (0.376,0.666)rectangle++(0.008,0.008) (0.409,0.587)rectangle++(0.008,0.008)
  (0.396,0.538)rectangle++(0.008,0.008) (0.28,0.79)rectangle++(0.008,0.008) (0.004,0.805)rectangle++(0.008,0.008) (0.288,0.783)rectangle++(0.008,0.008) (0.411,0.671)rectangle++(0.008,0.008)
  (0.427,0.577)rectangle++(0.008,0.008) (0.4,0.528)rectangle++(0.008,0.008) (0.371,0.515)rectangle++(0.008,0.008) (0.029,0.458)rectangle++(0.008,0.008) (0.146,0.837)rectangle++(0.008,0.008)
  (0.365,0.878)rectangle++(0.008,0.008) (0.326,0.863)rectangle++(0.008,0.008) (0.257,0.9)rectangle++(0.008,0.008) (0.092,0.55)rectangle++(0.008,0.008) (0.238,0.635)rectangle++(0.008,0.008)
  (0.335,0.622)rectangle++(0.008,0.008) (0.374,0.581)rectangle++(0.008,0.008) (0.378,0.548)rectangle++(0.008,0.008) (0.2,0.888)rectangle++(0.008,0.008) (0.128,0.591)rectangle++(0.008,0.008)
  (0.252,0.875)rectangle++(0.008,0.008) (0.125,0.776)rectangle++(0.008,0.008) (0.072,0.885)rectangle++(0.008,0.008) (0.348,0.795)rectangle++(0.008,0.008) (0.144,0.951)rectangle++(0.008,0.008)
  (0.167,0.945)rectangle++(0.008,0.008) (0.413,0.789)rectangle++(0.008,0.008) (0.47,0.629)rectangle++(0.008,0.008) (0.305,0.821)rectangle++(0.008,0.008) (0.147,0.742)rectangle++(0.008,0.008)
  (0.118,0.925)rectangle++(0.008,0.008) (0.384,0.797)rectangle++(0.008,0.008) (0.276,0.877)rectangle++(0.008,0.008) (0.24,0.904)rectangle++(0.008,0.008) (0.393,0.828)rectangle++(0.008,0.008)
  (0.475,0.654)rectangle++(0.008,0.008) (0.14,0.908)rectangle++(0.008,0.008) (0.388,0.781)rectangle++(0.008,0.008) (0.455,0.635)rectangle++(0.008,0.008) (0.077,0.522)rectangle++(0.008,0.008)
  (0.076,0.516)rectangle++(0.008,0.008) (0.218,0.625)rectangle++(0.008,0.008) (0.057,0.893)rectangle++(0.008,0.008) (0.15,0.948)rectangle++(0.008,0.008) (0.089,0.905)rectangle++(0.008,0.008)
  (0.31,0.891)rectangle++(0.008,0.008) (0.405,0.815)rectangle++(0.008,0.008) (0.146,0.749)rectangle++(0.008,0.008) (0.129,0.591)rectangle++(0.008,0.008) (0.105,0.56)rectangle++(0.008,0.008)
  (0.161,0.889)rectangle++(0.008,0.008) (0.127,0.894)rectangle++(0.008,0.008) (0.19,0.91)rectangle++(0.008,0.008) (0.411,0.764)rectangle++(0.008,0.008) (0.056,0.863)rectangle++(0.008,0.008)
  (0.166,0.9)rectangle++(0.008,0.008) (0.396,0.768)rectangle++(0.008,0.008) (0.126,0.592)rectangle++(0.008,0.008) (0.325,0.785)rectangle++(0.008,0.008) (0.428,0.659)rectangle++(0.008,0.008)
  (0.291,0.831)rectangle++(0.008,0.008) (0.188,0.928)rectangle++(0.008,0.008) (0.148,0.74)rectangle++(0.008,0.008) (0.258,0.749)rectangle++(0.008,0.008) (0.385,0.667)rectangle++(0.008,0.008)
  (0.201,0.905)rectangle++(0.008,0.008) (0.01,0.434)rectangle++(0.008,0.008) (0.159,0.611)rectangle++(0.008,0.008) (0.29,0.639)rectangle++(0.008,0.008) (0.315,0.862)rectangle++(0.008,0.008)
  (0.451,0.698)rectangle++(0.008,0.008) (0.299,0.844)rectangle++(0.008,0.008) (0.437,0.695)rectangle++(0.008,0.008) (0.447,0.578)rectangle++(0.008,0.008) (0.21,0.892)rectangle++(0.008,0.008)
  (0.414,0.749)rectangle++(0.008,0.008) (0.456,0.611)rectangle++(0.008,0.008) (0.426,0.533)rectangle++(0.008,0.008) (0.384,0.508)rectangle++(0.008,0.008) (0.328,0.843)rectangle++(0.008,0.008)
  (0.45,0.685)rectangle++(0.008,0.008) (0.337,0.868)rectangle++(0.008,0.008) (0.383,0.854)rectangle++(0.008,0.008) (0.479,0.67)rectangle++(0.008,0.008) (0.255,0.742)rectangle++(0.008,0.008)
  (0.233,0.859)rectangle++(0.008,0.008) (0.179,0.932)rectangle++(0.008,0.008) (0.331,0.898)rectangle++(0.008,0.008) (0.139,0.602)rectangle++(0.008,0.008) (0.125,0.579)rectangle++(0.008,0.008)
  (0.189,0.804)rectangle++(0.008,0.008) (0.141,0.952)rectangle++(0.008,0.008) (-0.004,0.817)rectangle++(0.008,0.008) (0.146,0.884)rectangle++(0.008,0.008) (0.372,0.838)rectangle++(0.008,0.008)
  (0.142,0.606)rectangle++(0.008,0.008) (0.147,0.729)rectangle++(0.008,0.008) (0.167,0.913)rectangle++(0.008,0.008) (0.402,0.774)rectangle++(0.008,0.008) (0.459,0.627)rectangle++(0.008,0.008)
  (0.302,0.82)rectangle++(0.008,0.008) (0.251,0.899)rectangle++(0.008,0.008) (0.435,0.737)rectangle++(0.008,0.008) (0.335,0.876)rectangle++(0.008,0.008) (0.107,0.919)rectangle++(0.008,0.008)
  (0.366,0.882)rectangle++(0.008,0.008) (0.482,0.689)rectangle++(0.008,0.008) (0.332,0.897)rectangle++(0.008,0.008) (0.11,0.562)rectangle++(0.008,0.008) (0.25,0.634)rectangle++(0.008,0.008)
  (0.34,0.617)rectangle++(0.008,0.008) (0.322,0.859)rectangle++(0.008,0.008) (0.146,0.668)rectangle++(0.008,0.008) (0.196,0.834)rectangle++(0.008,0.008) (0.387,0.727)rectangle++(0.008,0.008)
  (0.436,0.611)rectangle++(0.008,0.008) (0.417,0.54)rectangle++(0.008,0.008) (0.34,0.881)rectangle++(0.008,0.008) (-0.013,0.822)rectangle++(0.008,0.008) (0.045,0.875)rectangle++(0.008,0.008)
  (0.333,0.8)rectangle++(0.008,0.008) (0.437,0.663)rectangle++(0.008,0.008) (0.131,0.583)rectangle++(0.008,0.008) (0.292,0.854)rectangle++(0.008,0.008) (0.023,0.855)rectangle++(0.008,0.008)
  (0.344,0.851)rectangle++(0.008,0.008) (0.109,0.792)rectangle++(0.008,0.008) (0.08,0.536)rectangle++(0.008,0.008) (-0.018,0.842)rectangle++(0.008,0.008) (0.292,0.808)rectangle++(0.008,0.008)
  (0.094,0.91)rectangle++(0.008,0.008) (0.367,0.799)rectangle++(0.008,0.008) (0.107,0.917)rectangle++(0.008,0.008) (0.366,0.882)rectangle++(0.008,0.008) (0.134,0.582)rectangle++(0.008,0.008)
  (0.325,0.782)rectangle++(0.008,0.008) (0.143,0.734)rectangle++(0.008,0.008) (0.167,0.914)rectangle++(0.008,0.008) (0.402,0.775)rectangle++(0.008,0.008) (0.46,0.627)rectangle++(0.008,0.008)
  (0.115,0.935)rectangle++(0.008,0.008) (0.186,0.924)rectangle++(0.008,0.008) (0.21,0.92)rectangle++(0.008,0.008) (0.303,0.901)rectangle++(0.008,0.008) (0.405,0.818)rectangle++(0.008,0.008)
  (0.028,0.817)rectangle++(0.008,0.008) (0.304,0.78)rectangle++(0.008,0.008) (0.25,0.872)rectangle++(0.008,0.008) (0.425,0.726)rectangle++(0.008,0.008) (0.452,0.596)rectangle++(0.008,0.008)
  (0.419,0.527)rectangle++(0.008,0.008) (0.082,0.513)rectangle++(0.008,0.008) (0.146,0.723)rectangle++(0.008,0.008) (0.197,0.853)rectangle++(0.008,0.008) (0.213,0.896)rectangle++(0.008,0.008)
  (0.417,0.75)rectangle++(0.008,0.008) (0.333,0.877)rectangle++(0.008,0.008) (0.259,0.904)rectangle++(0.008,0.008) (0.441,0.737)rectangle++(0.008,0.008) (0.464,0.596)rectangle++(0.008,0.008)
  (0.424,0.523)rectangle++(0.008,0.008) (0.38,0.504)rectangle++(0.008,0.008) (0.275,0.779)rectangle++(0.008,0.008) (0.137,0.593)rectangle++(0.008,0.008) (0.274,0.639)rectangle++(0.008,0.008)
  (0.352,0.611)rectangle++(0.008,0.008) (0.348,0.876)rectangle++(0.008,0.008) (0.016,0.895)rectangle++(0.008,0.008) (0.326,0.82)rectangle++(0.008,0.008) (0.441,0.675)rectangle++(0.008,0.008)
  (0.442,0.567)rectangle++(0.008,0.008) (0.404,0.518)rectangle++(0.008,0.008) (0.204,0.884)rectangle++(0.008,0.008) (0.339,0.892)rectangle++(0.008,0.008) (0.084,0.886)rectangle++(0.008,0.008)
  (0.143,0.603)rectangle++(0.008,0.008) (0.053,0.888)rectangle++(0.008,0.008) (0.082,0.893)rectangle++(0.008,0.008) (0.356,0.795)rectangle++(0.008,0.008) (0.273,0.825)rectangle++(0.008,0.008)
  (0.239,0.887)rectangle++(0.008,0.008) (0.425,0.736)rectangle++(0.008,0.008) (0.064,0.911)rectangle++(0.008,0.008) (0.019,0.446)rectangle++(0.008,0.008) (0.282,0.773)rectangle++(0.008,0.008)
  (0.382,0.796)rectangle++(0.008,0.008) (0.143,0.736)rectangle++(0.008,0.008) (0.328,0.702)rectangle++(0.008,0.008) (0.399,0.62)rectangle++(0.008,0.008) (0.403,0.557)rectangle++(0.008,0.008)
  (0.383,0.527)rectangle++(0.008,0.008) (0.276,0.787)rectangle++(0.008,0.008) (0.118,0.565)rectangle++(0.008,0.008) (0.187,0.799)rectangle++(0.008,0.008) (0.37,0.715)rectangle++(0.008,0.008)
  (0.062,0.907)rectangle++(0.008,0.008) (0.32,0.872)rectangle++(0.008,0.008) (0.457,0.701)rectangle++(0.008,0.008) (0.404,0.756)rectangle++(0.008,0.008) (0.207,0.86)rectangle++(0.008,0.008)
  (0.401,0.735)rectangle++(0.008,0.008) (0.135,0.589)rectangle++(0.008,0.008) (0.193,0.807)rectangle++(0.008,0.008) (0.285,0.904)rectangle++(0.008,0.008) (0.142,0.718)rectangle++(0.008,0.008)
  (0.116,0.574)rectangle++(0.008,0.008) (0.153,0.929)rectangle++(0.008,0.008) (0.401,0.786)rectangle++(0.008,0.008) (0.267,0.911)rectangle++(0.008,0.008) (0.145,0.604)rectangle++(0.008,0.008)
  (0.125,0.579)rectangle++(0.008,0.008) (0.263,0.637)rectangle++(0.008,0.008) (0.347,0.613)rectangle++(0.008,0.008) (0.377,0.573)rectangle++(0.008,0.008) (0.376,0.544)rectangle++(0.008,0.008)
  (0.365,0.531)rectangle++(0.008,0.008) (0.356,0.528)rectangle++(0.008,0.008) (0.351,0.531)rectangle++(0.008,0.008) (0.136,0.874)rectangle++(0.008,0.008) (-0.007,0.81)rectangle++(0.008,0.008)
  (0.285,0.789)rectangle++(0.008,0.008) (0.412,0.675)rectangle++(0.008,0.008) (0.066,0.906)rectangle++(0.008,0.008) (0.331,0.893)rectangle++(0.008,0.008) (0.139,0.6)rectangle++(0.008,0.008)
  (0.147,0.729)rectangle++(0.008,0.008) (0.326,0.697)rectangle++(0.008,0.008) (0.257,0.844)rectangle++(0.008,0.008) (0.418,0.71)rectangle++(0.008,0.008) (0.444,0.592)rectangle++(0.008,0.008)
  (0.413,0.528)rectangle++(0.008,0.008) (0.37,0.73)rectangle++(0.008,0.008) (0.224,0.709)rectangle++(0.008,0.008) (0.223,0.848)rectangle++(0.008,0.008) (0.222,0.894)rectangle++(0.008,0.008)
  (0.095,0.815)rectangle++(0.008,0.008) (0.352,0.833)rectangle++(0.008,0.008) (0.126,0.573)rectangle++(0.008,0.008) (0.322,0.781)rectangle++(0.008,0.008) (0.256,0.872)rectangle++(0.008,0.008)
  (0.143,0.735)rectangle++(0.008,0.008) (0.001,0.814)rectangle++(0.008,0.008) (0.291,0.788)rectangle++(0.008,0.008) (0.149,0.72)rectangle++(0.008,0.008) (0.168,0.912)rectangle++(0.008,0.008)
  (0.088,0.547)rectangle++(0.008,0.008) (0.311,0.782)rectangle++(0.008,0.008) (0.421,0.663)rectangle++(0.008,0.008) (0.428,0.569)rectangle++(0.008,0.008) (0.379,0.735)rectangle++(0.008,0.008)
  (0.123,0.584)rectangle++(0.008,0.008) (0.291,0.854)rectangle++(0.008,0.008) (0.437,0.703)rectangle++(0.008,0.008) (0.45,0.581)rectangle++(0.008,0.008) (0.385,0.735)rectangle++(0.008,0.008)
  (0.378,0.84)rectangle++(0.008,0.008) (0.323,0.833)rectangle++(0.008,0.008) (0.32,0.889)rectangle++(0.008,0.008) (0.406,0.813)rectangle++(0.008,0.008) (0.037,0.469)rectangle++(0.008,0.008)
  (0.013,0.438)rectangle++(0.008,0.008) (0.276,0.833)rectangle++(0.008,0.008) (0.129,0.948)rectangle++(0.008,0.008) (0.191,0.928)rectangle++(0.008,0.008) (-0.001,0.861)rectangle++(0.008,0.008)
  (0.042,0.882)rectangle++(0.008,0.008) (0.162,0.906)rectangle++(0.008,0.008) (0.371,0.868)rectangle++(0.008,0.008) (0.272,0.901)rectangle++(0.008,0.008) (0.445,0.731)rectangle++(0.008,0.008)
  (0.464,0.592)rectangle++(0.008,0.008) (0.423,0.521)rectangle++(0.008,0.008) (0.289,0.785)rectangle++(0.008,0.008) (0.385,0.798)rectangle++(0.008,0.008) (0.089,0.542)rectangle++(0.008,0.008)
  (0.31,0.781)rectangle++(0.008,0.008) (0.42,0.663)rectangle++(0.008,0.008) (0.392,0.755)rectangle++(0.008,0.008) (0.448,0.621)rectangle++(0.008,0.008) (0.087,0.915)rectangle++(0.008,0.008)
  (0.366,0.803)rectangle++(0.008,0.008) (0.008,0.867)rectangle++(0.008,0.008) (0.043,0.885)rectangle++(0.008,0.008) (0.276,0.908)rectangle++(0.008,0.008) (0.24,0.915)rectangle++(0.008,0.008)
  (0.132,0.596)rectangle++(0.008,0.008) (0.273,0.642)rectangle++(0.008,0.008) (0.36,0.773)rectangle++(0.008,0.008) (0.44,0.641)rectangle++(0.008,0.008) (0.429,0.552)rectangle++(0.008,0.008)
  (0.004,0.836)rectangle++(0.008,0.008) (0.263,0.758)rectangle++(0.008,0.008) (0.269,0.764)rectangle++(0.008,0.008) (0.083,0.899)rectangle++(0.008,0.008) (0.332,0.893)rectangle++(0.008,0.008)
  (0.096,0.92)rectangle++(0.008,0.008) (0.136,0.592)rectangle++(0.008,0.008) (0.273,0.87)rectangle++(0.008,0.008) (0.237,0.86)rectangle++(0.008,0.008) (0.128,0.591)rectangle++(0.008,0.008)
  (0.095,0.857)rectangle++(0.008,0.008) (0.023,0.886)rectangle++(0.008,0.008) (0.275,0.908)rectangle++(0.008,0.008) (-0.008,0.794)rectangle++(0.008,0.008) (0.279,0.782)rectangle++(0.008,0.008)
  (0.383,0.799)rectangle++(0.008,0.008) (0.459,0.645)rectangle++(0.008,0.008) (0.302,0.826)rectangle++(0.008,0.008) (0.009,0.78)rectangle++(0.008,0.008) (0.351,0.868)rectangle++(0.008,0.008)
  (0.125,0.943)rectangle++(0.008,0.008) (0.393,0.803)rectangle++(0.008,0.008) (0.092,0.924)rectangle++(0.008,0.008) (0.093,0.814)rectangle++(0.008,0.008) (0.333,0.755)rectangle++(0.008,0.008)
  (0.375,0.845)rectangle++(0.008,0.008) (0.274,0.894)rectangle++(0.008,0.008) (0.005,0.869)rectangle++(0.008,0.008) (0.274,0.907)rectangle++(0.008,0.008) (0.024,0.814)rectangle++(0.008,0.008)
  (0.277,0.886)rectangle++(0.008,0.008) (0.327,0.861)rectangle++(0.008,0.008) (0.234,0.899)rectangle++(0.008,0.008) (0.206,0.864)rectangle++(0.008,0.008) (0.402,0.738)rectangle++(0.008,0.008)
  (0.447,0.61)rectangle++(0.008,0.008) (0.421,0.536)rectangle++(0.008,0.008) (0.383,0.511)rectangle++(0.008,0.008) (0.357,0.513)rectangle++(0.008,0.008) (0.357,0.735)rectangle++(0.008,0.008)
  (0.425,0.625)rectangle++(0.008,0.008) (0.417,0.55)rectangle++(0.008,0.008) (0.386,0.519)rectangle++(0.008,0.008) (0.328,0.845)rectangle++(0.008,0.008) (0.128,0.575)rectangle++(0.008,0.008)
  (0.263,0.634)rectangle++(0.008,0.008) (0.003,0.85)rectangle++(0.008,0.008) (0.349,0.889)rectangle++(0.008,0.008) (0.273,0.917)rectangle++(0.008,0.008) (0.239,0.918)rectangle++(0.008,0.008)
  (0.228,0.918)rectangle++(0.008,0.008) (0.393,0.832)rectangle++(0.008,0.008) (0.275,0.82)rectangle++(0.008,0.008) (0.418,0.693)rectangle++(0.008,0.008) (0.288,0.843)rectangle++(0.008,0.008)
  (0.432,0.699)rectangle++(0.008,0.008) (0.293,0.844)rectangle++(0.008,0.008) (0.25,0.9)rectangle++(0.008,0.008) (0.288,0.909)rectangle++(0.008,0.008) (0.456,0.729)rectangle++(0.008,0.008)
  (0.042,0.475)rectangle++(0.008,0.008) (0.188,0.618)rectangle++(0.008,0.008) (0.249,0.903)rectangle++(0.008,0.008) (0.28,0.914)rectangle++(0.008,0.008) (0.313,0.901)rectangle++(0.008,0.008)
  (0.464,0.716)rectangle++(0.008,0.008) (0.408,0.758)rectangle++(0.008,0.008) (0.331,0.879)rectangle++(0.008,0.008) (0.025,0.815)rectangle++(0.008,0.008) (0.302,0.78)rectangle++(0.008,0.008)
  (0.249,0.872)rectangle++(0.008,0.008) (0.231,0.903)rectangle++(0.008,0.008) (0.427,0.746)rectangle++(0.008,0.008) (0.128,0.773)rectangle++(0.008,0.008) (0.073,0.884)rectangle++(0.008,0.008)
  (0.087,0.884)rectangle++(0.008,0.008) (0.142,0.602)rectangle++(0.008,0.008) (0.329,0.785)rectangle++(0.008,0.008) (0.246,0.737)rectangle++(0.008,0.008) (0.375,0.666)rectangle++(0.008,0.008)
  (0.257,0.902)rectangle++(0.008,0.008) (0.183,0.938)rectangle++(0.008,0.008) (0.418,0.78)rectangle++(0.008,0.008) (0.469,0.624)rectangle++(0.008,0.008) (0.127,0.584)rectangle++(0.008,0.008)
  (0.142,0.953)rectangle++(0.008,0.008) (0.294,0.906)rectangle++(0.008,0.008) (0.457,0.726)rectangle++(0.008,0.008) (0.112,0.551)rectangle++(0.008,0.008) (0.247,0.628)rectangle++(0.008,0.008)
  (0.253,0.9)rectangle++(0.008,0.008) (0.436,0.737)rectangle++(0.008,0.008) (0.029,0.893)rectangle++(0.008,0.008) (0.331,0.814)rectangle++(0.008,0.008) (0.193,0.925)rectangle++(0.008,0.008)
  (0.174,0.941)rectangle++(0.008,0.008) (0.146,0.664)rectangle++(0.008,0.008) (0.196,0.833)rectangle++(0.008,0.008) (0.045,0.906)rectangle++(0.008,0.008) (0.12,0.583)rectangle++(0.008,0.008)
  (0.094,0.854)rectangle++(0.008,0.008) (0.004,0.848)rectangle++(0.008,0.008) (0.273,0.827)rectangle++(0.008,0.008) (0.239,0.887)rectangle++(0.008,0.008) (0.228,0.908)rectangle++(0.008,0.008)
  (0.391,0.831)rectangle++(0.008,0.008) (0.279,0.889)rectangle++(0.008,0.008) (0.313,0.897)rectangle++(0.008,0.008) (0.149,0.735)rectangle++(0.008,0.008) (0.351,0.809)rectangle++(0.008,0.008)
  (0.099,0.808)rectangle++(0.008,0.008) (0.333,0.75)rectangle++(0.008,0.008) (0.166,0.939)rectangle++(0.008,0.008) (0.254,0.909)rectangle++(0.008,0.008) (0.397,0.827)rectangle++(0.008,0.008)
  (0.108,0.925)rectangle++(0.008,0.008) (0.137,0.593)rectangle++(0.008,0.008) (0.274,0.639)rectangle++(0.008,0.008) (0.116,0.573)rectangle++(0.008,0.008) (0.109,0.796)rectangle++(0.008,0.008)
  (0.333,0.741)rectangle++(0.008,0.008) (0.416,0.636)rectangle++(0.008,0.008) (0.287,0.823)rectangle++(0.008,0.008) (0.244,0.886)rectangle++(0.008,0.008) (0.427,0.734)rectangle++(0.008,0.008)
  (0.291,0.856)rectangle++(0.008,0.008) (0.396,0.811)rectangle++(0.008,0.008) (0.281,0.882)rectangle++(0.008,0.008) (0.133,0.583)rectangle++(0.008,0.008) (0.268,0.635)rectangle++(0.008,0.008)
  (0.238,0.823)rectangle++(0.008,0.008) (0.381,0.813)rectangle++(0.008,0.008) (0.219,0.884)rectangle++(0.008,0.008) (0.272,0.824)rectangle++(0.008,0.008) (0.128,0.59)rectangle++(0.008,0.008)
  (0.269,0.64)rectangle++(0.008,0.008) (0.359,0.773)rectangle++(0.008,0.008) (0.439,0.641)rectangle++(0.008,0.008) (0.336,0.862)rectangle++(0.008,0.008) (0.131,0.579)rectangle++(0.008,0.008)
  (0.191,0.804)rectangle++(0.008,0.008) (0.285,0.904)rectangle++(0.008,0.008) (0.187,0.938)rectangle++(0.008,0.008) (0.062,0.917)rectangle++(0.008,0.008) (0.046,0.899)rectangle++(0.008,0.008)
  (0.045,0.894)rectangle++(0.008,0.008) (0.046,0.893)rectangle++(0.008,0.008) (0.339,0.808)rectangle++(0.008,0.008) (0.443,0.665)rectangle++(0.008,0.008) (0.439,0.562)rectangle++(0.008,0.008)
  (0.295,0.798)rectangle++(0.008,0.008) (0.419,0.676)rectangle++(0.008,0.008) (0.432,0.576)rectangle++(0.008,0.008) (0.402,0.525)rectangle++(0.008,0.008) (0.371,0.513)rectangle++(0.008,0.008)
  (0.36,0.733)rectangle++(0.008,0.008) (0.425,0.623)rectangle++(0.008,0.008) (0.107,0.797)rectangle++(0.008,0.008) (0.333,0.742)rectangle++(0.008,0.008) (0.416,0.637)rectangle++(0.008,0.008)
  (0.417,0.559)rectangle++(0.008,0.008) (0.004,0.836)rectangle++(0.008,0.008) (0.3,0.797)rectangle++(0.008,0.008) (0.421,0.674)rectangle++(0.008,0.008) (0.002,0.788)rectangle++(0.008,0.008)
  (0.105,0.799)rectangle++(0.008,0.008) (0.183,0.878)rectangle++(0.008,0.008) (0.396,0.752)rectangle++(0.008,0.008) (0.281,0.862)rectangle++(0.008,0.008) (0.086,0.917)rectangle++(0.008,0.008)
  (0.176,0.918)rectangle++(0.008,0.008) (0.207,0.918)rectangle++(0.008,0.008) (0,0.862)rectangle++(0.008,0.008) (0.141,0.706)rectangle++(0.008,0.008) (0.022,0.898)rectangle++(0.008,0.008)
  (0.101,0.808)rectangle++(0.008,0.008) (0.259,0.905)rectangle++(0.008,0.008) (0.441,0.737)rectangle++(0.008,0.008) (0.407,0.766)rectangle++(0.008,0.008) (0.15,0.717)rectangle++(0.008,0.008)
  (0.144,0.722)rectangle++(0.008,0.008) (0.323,0.695)rectangle++(0.008,0.008) (0.395,0.619)rectangle++(0.008,0.008) (0.001,0.818)rectangle++(0.008,0.008) (0.292,0.79)rectangle++(0.008,0.008)
  (0.271,0.769)rectangle++(0.008,0.008) (0.261,0.751)rectangle++(0.008,0.008) (0.183,0.917)rectangle++(0.008,0.008) (0.41,0.77)rectangle++(0.008,0.008) (0.117,0.918)rectangle++(0.008,0.008)
  (0.187,0.918)rectangle++(0.008,0.008) (0.21,0.918)rectangle++(0.008,0.008) (0.218,0.918)rectangle++(0.008,0.008) (0.299,0.876)rectangle++(0.008,0.008) (0.248,0.904)rectangle++(0.008,0.008)
  (0.145,0.603)rectangle++(0.008,0.008) (0.252,0.877)rectangle++(0.008,0.008) (0.428,0.727)rectangle++(0.008,0.008) (0.334,0.874)rectangle++(0.008,0.008) (0.144,0.608)rectangle++(0.008,0.008)
  (0.331,0.785)rectangle++(0.008,0.008) (0.431,0.657)rectangle++(0.008,0.008) (0.393,0.752)rectangle++(0.008,0.008) (0.447,0.62)rectangle++(0.008,0.008) (0.096,0.531)rectangle++(0.008,0.008)
  (0.31,0.778)rectangle++(0.008,0.008) (0.304,0.874)rectangle++(0.008,0.008) (0.45,0.707)rectangle++(0.008,0.008) (0.372,0.866)rectangle++(0.008,0.008) (0.137,0.6)rectangle++(0.008,0.008)
  (0.227,0.882)rectangle++(0.008,0.008) (0.224,0.906)rectangle++(0.008,0.008) (0.425,0.751)rectangle++(0.008,0.008) (0.462,0.608)rectangle++(0.008,0.008) (0.097,0.811)rectangle++(0.008,0.008)
  (0.334,0.753)rectangle++(0.008,0.008) (0.388,0.785)rectangle++(0.008,0.008) (0.404,0.783)rectangle++(0.008,0.008) (0.463,0.63)rectangle++(0.008,0.008) (0.151,0.698)rectangle++(0.008,0.008)
  (0.198,0.844)rectangle++(0.008,0.008) (0.274,0.828)rectangle++(0.008,0.008) (0.012,0.867)rectangle++(0.008,0.008) (0.151,0.901)rectangle++(0.008,0.008) (0.249,0.903)rectangle++(0.008,0.008)
  (0.435,0.74)rectangle++(0.008,0.008) (0.462,0.599)rectangle++(0.008,0.008) (0.425,0.525)rectangle++(0.008,0.008) (0.087,0.906)rectangle++(0.008,0.008) (0.177,0.914)rectangle++(0.008,0.008)
  (0.383,0.839)rectangle++(0.008,0.008) (0.09,0.924)rectangle++(0.008,0.008) (0.159,0.941)rectangle++(0.008,0.008) (0.236,0.899)rectangle++(0.008,0.008) (0.392,0.828)rectangle++(0.008,0.008)
  (-0.011,0.827)rectangle++(0.008,0.008) (0.144,0.887)rectangle++(0.008,0.008) (0.372,0.839)rectangle++(0.008,0.008) (0.262,0.76)rectangle++(0.008,0.008) (0.108,0.9)rectangle++(0.008,0.008)
  (0.123,0.587)rectangle++(0.008,0.008) (0.086,0.528)rectangle++(0.008,0.008) (0.307,0.779)rectangle++(0.008,0.008) (0.418,0.663)rectangle++(0.008,0.008) (0.288,0.832)rectangle++(0.008,0.008)
  (0.428,0.694)rectangle++(0.008,0.008) (0.291,0.843)rectangle++(0.008,0.008) (0.433,0.698)rectangle++(0.008,0.008) (0.107,0.546)rectangle++(0.008,0.008) (0.243,0.628)rectangle++(0.008,0.008)
  (0.335,0.616)rectangle++(0.008,0.008) (0.372,0.579)rectangle++(0.008,0.008) (0.376,0.548)rectangle++(0.008,0.008) (0.071,0.896)rectangle++(0.008,0.008) (0.084,0.91)rectangle++(0.008,0.008)
  (0.158,0.939)rectangle++(0.008,0.008) (0.147,0.749)rectangle++(0.008,0.008) (0.109,0.798)rectangle++(0.008,0.008) (0.03,0.904)rectangle++(0.008,0.008) (0.353,0.86)rectangle++(0.008,0.008)
  (0.089,0.548)rectangle++(0.008,0.008) (0.236,0.635)rectangle++(0.008,0.008) (0.056,0.865)rectangle++(0.008,0.008) (0.166,0.9)rectangle++(0.008,0.008) (0.397,0.769)rectangle++(0.008,0.008)
  (0.403,0.778)rectangle++(0.008,0.008) (0.115,0.918)rectangle++(0.008,0.008) (0.371,0.85)rectangle++(0.008,0.008) (0.199,0.93)rectangle++(0.008,0.008) (0.139,0.593)rectangle++(0.008,0.008)
  (0.275,0.638)rectangle++(0.008,0.008) (0.29,0.894)rectangle++(0.008,0.008) (0.451,0.721)rectangle++(0.008,0.008) (0.113,0.937)rectangle++(0.008,0.008) (0.087,0.546)rectangle++(0.008,0.008)
  (0.147,0.686)rectangle++(0.008,0.008) (0.323,0.835)rectangle++(0.008,0.008) (0.319,0.889)rectangle++(0.008,0.008) (0.015,0.894)rectangle++(0.008,0.008) (0.012,0.88)rectangle++(0.008,0.008);
\fill[piece3]
  (0.151,0.905)rectangle++(0.008,0.008) (0.125,0.896)rectangle++(0.008,0.008) (0.04,0.885)rectangle++(0.008,0.008) (0.122,0.587)rectangle++(0.008,0.008) (0.265,0.641)rectangle++(0.008,0.008)
  (0.349,0.615)rectangle++(0.008,0.008) (0.313,0.877)rectangle++(0.008,0.008) (0.005,0.879)rectangle++(0.008,0.008) (0.041,0.886)rectangle++(0.008,0.008) (0.146,0.746)rectangle++(0.008,0.008)
  (0.36,0.856)rectangle++(0.008,0.008) (0.469,0.679)rectangle++(0.008,0.008) (0.355,0.88)rectangle++(0.008,0.008) (0.476,0.692)rectangle++(0.008,0.008) (0.307,0.842)rectangle++(0.008,0.008)
  (0.251,0.893)rectangle++(0.008,0.008) (0.073,0.528)rectangle++(0.008,0.008) (0.221,0.631)rectangle++(0.008,0.008) (0.326,0.626)rectangle++(0.008,0.008) (0.266,0.764)rectangle++(0.008,0.008)
  (0.394,0.671)rectangle++(0.008,0.008) (0.419,0.582)rectangle++(0.008,0.008) (0.141,0.89)rectangle++(0.008,0.008) (0.367,0.88)rectangle++(0.008,0.008) (0.481,0.688)rectangle++(0.008,0.008)
  (0.321,0.832)rectangle++(0.008,0.008) (0.398,0.802)rectangle++(0.008,0.008) (0.467,0.641)rectangle++(0.008,0.008) (0.442,0.542)rectangle++(0.008,0.008) (0.296,0.792)rectangle++(0.008,0.008)
  (0.417,0.673)rectangle++(0.008,0.008) (0.43,0.575)rectangle++(0.008,0.008) (0.292,0.803)rectangle++(0.008,0.008) (0.246,0.879)rectangle++(0.008,0.008) (0.425,0.731)rectangle++(0.008,0.008)
  (0.29,0.855)rectangle++(0.008,0.008) (0.245,0.894)rectangle++(0.008,0.008) (0.31,0.885)rectangle++(0.008,0.008) (0.457,0.71)rectangle++(0.008,0.008) (0.462,0.578)rectangle++(0.008,0.008)
  (0.417,0.516)rectangle++(0.008,0.008) (0.15,0.681)rectangle++(0.008,0.008) (0.198,0.839)rectangle++(0.008,0.008) (0.214,0.891)rectangle++(0.008,0.008) (0.415,0.747)rectangle++(0.008,0.008)
  (0.456,0.61)rectangle++(0.008,0.008) (0.354,0.876)rectangle++(0.008,0.008) (0.474,0.69)rectangle++(0.008,0.008) (0.462,0.563)rectangle++(0.008,0.008) (0.369,0.861)rectangle++(0.008,0.008)
  (0.475,0.678)rectangle++(0.008,0.008) (0.118,0.935)rectangle++(0.008,0.008) (0.387,0.802)rectangle++(0.008,0.008) (0.463,0.645)rectangle++(0.008,0.008) (0.128,0.586)rectangle++(0.008,0.008)
  (0.263,0.911)rectangle++(0.008,0.008) (0.133,0.597)rectangle++(0.008,0.008) (0.038,0.47)rectangle++(0.008,0.008) (0.16,0.768)rectangle++(0.008,0.008) (0.357,0.814)rectangle++(0.008,0.008)
  (0.453,0.661)rectangle++(0.008,0.008) (0.343,0.884)rectangle++(0.008,0.008) (0.263,0.907)rectangle++(0.008,0.008) (0.443,0.737)rectangle++(0.008,0.008) (0.465,0.595)rectangle++(0.008,0.008)
  (0.043,0.882)rectangle++(0.008,0.008) (0.272,0.769)rectangle++(0.008,0.008) (0.398,0.671)rectangle++(0.008,0.008) (0.421,0.581)rectangle++(0.008,0.008) (-0.001,0.802)rectangle++(0.008,0.008)
  (0.348,0.889)rectangle++(0.008,0.008) (0.135,0.584)rectangle++(0.008,0.008) (0.127,0.894)rectangle++(0.008,0.008) (0.376,0.78)rectangle++(0.008,0.008) (0.45,0.638)rectangle++(0.008,0.008)
  (0.127,0.584)rectangle++(0.008,0.008) (0.19,0.806)rectangle++(0.008,0.008) (0.245,0.892)rectangle++(0.008,0.008) (0.392,0.825)rectangle++(0.008,0.008) (0.271,0.766)rectangle++(0.008,0.008)
  (0.125,0.933)rectangle++(0.008,0.008) (0.189,0.923)rectangle++(0.008,0.008) (0.138,0.592)rectangle++(0.008,0.008) (0.327,0.783)rectangle++(0.008,0.008) (0.149,0.719)rectangle++(0.008,0.008)
  (0.324,0.692)rectangle++(0.008,0.008) (0.273,0.831)rectangle++(0.008,0.008) (0.375,0.855)rectangle++(0.008,0.008) (0.094,0.933)rectangle++(0.008,0.008) (0.37,0.855)rectangle++(0.008,0.008)
  (0.474,0.675)rectangle++(0.008,0.008) (0.457,0.556)rectangle++(0.008,0.008) (0.407,0.507)rectangle++(0.008,0.008) (0.331,0.843)rectangle++(0.008,0.008) (0.259,0.893)rectangle++(0.008,0.008)
  (0.019,0.901)rectangle++(0.008,0.008) (0.33,0.822)rectangle++(0.008,0.008) (0.258,0.886)rectangle++(0.008,0.008) (0.134,0.592)rectangle++(0.008,0.008) (0.272,0.639)rectangle++(0.008,0.008)
  (0.352,0.612)rectangle++(0.008,0.008) (0.266,0.815)rectangle++(0.008,0.008) (0.237,0.883)rectangle++(0.008,0.008) (0.076,0.915)rectangle++(0.008,0.008) (0.361,0.807)rectangle++(0.008,0.008)
  (0.452,0.656)rectangle++(0.008,0.008) (0.44,0.555)rectangle++(0.008,0.008) (0.336,0.85)rectangle++(0.008,0.008) (0.26,0.895)rectangle++(0.008,0.008) (0.438,0.733)rectangle++(0.008,0.008)
  (0.461,0.595)rectangle++(0.008,0.008) (0.422,0.524)rectangle++(0.008,0.008) (0.379,0.505)rectangle++(0.008,0.008) (0.037,0.888)rectangle++(0.008,0.008) (0.333,0.809)rectangle++(0.008,0.008)
  (0.44,0.667)rectangle++(0.008,0.008) (0.439,0.564)rectangle++(0.008,0.008) (0.38,0.733)rectangle++(0.008,0.008) (0.113,0.788)rectangle++(0.008,0.008) (0.332,0.736)rectangle++(0.008,0.008)
  (0.259,0.748)rectangle++(0.008,0.008) (0.385,0.666)rectangle++(0.008,0.008) (0.385,0.761)rectangle++(0.008,0.008) (0.15,0.716)rectangle++(0.008,0.008) (0.093,0.912)rectangle++(0.008,0.008)
  (0.287,0.9)rectangle++(0.008,0.008) (0.082,0.923)rectangle++(0.008,0.008) (0.217,0.886)rectangle++(0.008,0.008) (0.272,0.816)rectangle++(0.008,0.008) (0.239,0.884)rectangle++(0.008,0.008)
  (0.424,0.735)rectangle++(0.008,0.008) (0.456,0.601)rectangle++(0.008,0.008) (0.389,0.737)rectangle++(0.008,0.008) (0.379,0.84)rectangle++(0.008,0.008) (0.147,0.607)rectangle++(0.008,0.008)
  (0.283,0.642)rectangle++(0.008,0.008) (0.358,0.608)rectangle++(0.008,0.008) (0.38,0.567)rectangle++(0.008,0.008) (0.212,0.852)rectangle++(0.008,0.008) (0.127,0.578)rectangle++(0.008,0.008)
  (0.264,0.635)rectangle++(0.008,0.008) (0.347,0.612)rectangle++(0.008,0.008) (0.376,0.573)rectangle++(0.008,0.008) (0.349,0.873)rectangle++(0.008,0.008) (0.051,0.486)rectangle++(0.008,0.008)
  (0.07,0.51)rectangle++(0.008,0.008) (0.213,0.624)rectangle++(0.008,0.008) (0.159,0.932)rectangle++(0.008,0.008) (0.382,0.845)rectangle++(0.008,0.008) (0.169,0.945)rectangle++(0.008,0.008)
  (0.414,0.788)rectangle++(0.008,0.008) (0.275,0.818)rectangle++(0.008,0.008) (0.417,0.692)rectangle++(0.008,0.008) (0.326,0.861)rectangle++(0.008,0.008) (0.456,0.693)rectangle++(0.008,0.008)
  (0.107,0.545)rectangle++(0.008,0.008) (0.344,0.843)rectangle++(0.008,0.008) (0.263,0.893)rectangle++(0.008,0.008) (0.325,0.835)rectangle++(0.008,0.008) (0.012,0.873)rectangle++(0.008,0.008)
  (0.359,0.889)rectangle++(0.008,0.008) (0.007,0.866)rectangle++(0.008,0.008) (0.343,0.856)rectangle++(0.008,0.008) (0.462,0.685)rectangle++(0.008,0.008) (0.107,0.919)rectangle++(0.008,0.008)
  (0.376,0.798)rectangle++(0.008,0.008) (0.456,0.647)rectangle++(0.008,0.008) (0.301,0.827)rectangle++(0.008,0.008) (0.37,0.874)rectangle++(0.008,0.008) (0.138,0.951)rectangle++(0.008,0.008)
  (0.402,0.802)rectangle++(0.008,0.008) (0.27,0.765)rectangle++(0.008,0.008) (0.396,0.67)rectangle++(0.008,0.008) (0.419,0.581)rectangle++(0.008,0.008) (0.271,0.889)rectangle++(0.008,0.008)
  (0.441,0.726)rectangle++(0.008,0.008) (0.46,0.591)rectangle++(0.008,0.008) (0.302,0.808)rectangle++(0.008,0.008) (0.249,0.881)rectangle++(0.008,0.008) (0.267,0.76)rectangle++(0.008,0.008)
  (0.273,0.805)rectangle++(0.008,0.008) (0.106,0.568)rectangle++(0.008,0.008) (0.25,0.638)rectangle++(0.008,0.008) (0.232,0.824)rectangle++(0.008,0.008) (0.135,0.947)rectangle++(0.008,0.008)
  (0.318,0.897)rectangle++(0.008,0.008) (0.254,0.911)rectangle++(0.008,0.008) (0.137,0.589)rectangle++(0.008,0.008) (0.053,0.887)rectangle++(0.008,0.008) (-0.001,0.844)rectangle++(0.008,0.008)
  (0.256,0.753)rectangle++(0.008,0.008) (0.263,0.76)rectangle++(0.008,0.008) (0.146,0.916)rectangle++(0.008,0.008) (0.393,0.783)rectangle++(0.008,0.008) (0.28,0.873)rectangle++(0.008,0.008)
  (0.439,0.715)rectangle++(0.008,0.008) (0.381,0.836)rectangle++(0.008,0.008) (0.055,0.492)rectangle++(0.008,0.008) (0.119,0.583)rectangle++(0.008,0.008) (0.262,0.64)rectangle++(0.008,0.008)
  (0.358,0.774)rectangle++(0.008,0.008) (0.439,0.642)rectangle++(0.008,0.008) (0.22,0.893)rectangle++(0.008,0.008) (0.178,0.937)rectangle++(0.008,0.008) (0.415,0.781)rectangle++(0.008,0.008)
  (0.468,0.625)rectangle++(0.008,0.008) (0.436,0.535)rectangle++(0.008,0.008) (0.052,0.825)rectangle++(0.008,0.008) (0.363,0.885)rectangle++(0.008,0.008) (0.06,0.499)rectangle++(0.008,0.008)
  (0.205,0.623)rectangle++(0.008,0.008) (0.138,0.941)rectangle++(0.008,0.008) (0.081,0.897)rectangle++(0.008,0.008) (0.285,0.898)rectangle++(0.008,0.008) (0.38,0.86)rectangle++(0.008,0.008)
  (0.48,0.674)rectangle++(0.008,0.008) (0.26,0.909)rectangle++(0.008,0.008) (0.235,0.915)rectangle++(0.008,0.008) (0.394,0.831)rectangle++(0.008,0.008) (0.412,0.791)rectangle++(0.008,0.008)
  (0.332,0.883)rectangle++(0.008,0.008) (0.091,0.926)rectangle++(0.008,0.008) (0.178,0.921)rectangle++(0.008,0.008) (0.409,0.774)rectangle++(0.008,0.008) (0.463,0.624)rectangle++(0.008,0.008)
  (0.105,0.562)rectangle++(0.008,0.008) (0.345,0.844)rectangle++(0.008,0.008) (0.404,0.801)rectangle++(0.008,0.008) (0.204,0.924)rectangle++(0.008,0.008) (0.175,0.941)rectangle++(0.008,0.008)
  (0.14,0.595)rectangle++(0.008,0.008) (0.194,0.81)rectangle++(0.008,0.008) (0.377,0.717)rectangle++(0.008,0.008) (0.428,0.609)rectangle++(0.008,0.008) (0.412,0.542)rectangle++(0.008,0.008)
  (0.146,0.724)rectangle++(0.008,0.008) (0.324,0.695)rectangle++(0.008,0.008) (0.256,0.843)rectangle++(0.008,0.008) (0.418,0.71)rectangle++(0.008,0.008) (0.444,0.592)rectangle++(0.008,0.008)
  (0.098,0.811)rectangle++(0.008,0.008) (0.334,0.752)rectangle++(0.008,0.008) (0.42,0.641)rectangle++(0.008,0.008) (0.099,0.535)rectangle++(0.008,0.008) (0.097,0.844)rectangle++(0.008,0.008);
\fill[piece4]
  (-0.141,0.888)rectangle++(0.008,0.008) (-0.004,0.662)rectangle++(0.008,0.008) (-0.001,0.724)rectangle++(0.008,0.008) (-0.003,0.902)rectangle++(0.008,0.008) (-0.059,0.844)rectangle++(0.008,0.008)
  (0.052,0.846)rectangle++(0.008,0.008) (0.02,0.953)rectangle++(0.008,0.008) (0.005,0.966)rectangle++(0.008,0.008) (0,0.748)rectangle++(0.008,0.008) (-0.188,0.875)rectangle++(0.008,0.008)
  (0.038,0.956)rectangle++(0.008,0.008) (0.148,0.95)rectangle++(0.008,0.008) (-0.118,0.978)rectangle++(0.008,0.008) (0.002,0.97)rectangle++(0.008,0.008) (-0.065,0.851)rectangle++(0.008,0.008)
  (0.013,0.795)rectangle++(0.008,0.008) (0.003,0.932)rectangle++(0.008,0.008) (-0.055,0.842)rectangle++(0.008,0.008) (-0.026,0.951)rectangle++(0.008,0.008) (-0.005,0.957)rectangle++(0.008,0.008)
  (-0.004,1)rectangle++(0.008,0.008) (-0.205,0.887)rectangle++(0.008,0.008) (-0.176,0.913)rectangle++(0.008,0.008) (-0.061,0.978)rectangle++(0.008,0.008) (0.006,0.973)rectangle++(0.008,0.008)
  (0.193,0.885)rectangle++(0.008,0.008) (0.079,0.969)rectangle++(0.008,0.008) (0.031,1.005)rectangle++(0.008,0.008) (0.064,0.996)rectangle++(0.008,0.008) (0.025,1.016)rectangle++(0.008,0.008)
  (0.049,0.836)rectangle++(0.008,0.008) (0.018,0.949)rectangle++(0.008,0.008) (0.005,0.996)rectangle++(0.008,0.008) (0,1.016)rectangle++(0.008,0.008) (0.15,0.968)rectangle++(0.008,0.008)
  (-0.12,0.865)rectangle++(0.008,0.008) (0.163,0.858)rectangle++(0.008,0.008) (-0.195,0.858)rectangle++(0.008,0.008) (-0.117,0.94)rectangle++(0.008,0.008) (0.182,0.892)rectangle++(0.008,0.008)
  (0.085,0.98)rectangle++(0.008,0.008) (0.034,1.009)rectangle++(0.008,0.008) (0.064,0.996)rectangle++(0.008,0.008) (0.171,0.929)rectangle++(0.008,0.008) (-0.008,0.736)rectangle++(0.008,0.008)
  (-0.116,0.914)rectangle++(0.008,0.008) (-0.083,0.983)rectangle++(0.008,0.008) (-0.037,1.01)rectangle++(0.008,0.008) (0.006,0.976)rectangle++(0.008,0.008) (0.145,0.959)rectangle++(0.008,0.008)
  (-0.166,0.809)rectangle++(0.008,0.008) (-0.18,0.828)rectangle++(0.008,0.008) (-0.001,0.748)rectangle++(0.008,0.008) (-0.003,0.912)rectangle++(0.008,0.008) (-0.004,0.981)rectangle++(0.008,0.008)
  (-0.191,0.906)rectangle++(0.008,0.008) (-0.083,0.978)rectangle++(0.008,0.008) (-0.166,0.949)rectangle++(0.008,0.008) (-0.179,0.823)rectangle++(0.008,0.008) (0.159,0.811)rectangle++(0.008,0.008)
  (-0.046,0.977)rectangle++(0.008,0.008) (0.135,0.967)rectangle++(0.008,0.008) (0.168,0.938)rectangle++(0.008,0.008) (0.157,0.932)rectangle++(0.008,0.008) (0.138,0.897)rectangle++(0.008,0.008)
  (0.056,0.974)rectangle++(0.008,0.008) (0.154,0.951)rectangle++(0.008,0.008) (-0.063,0.85)rectangle++(0.008,0.008) (0.042,0.836)rectangle++(0.008,0.008) (0.015,0.949)rectangle++(0.008,0.008)
  (0.004,0.996)rectangle++(0.008,0.008) (-0.001,1.016)rectangle++(0.008,0.008) (-0.207,0.888)rectangle++(0.008,0.008) (0.133,0.887)rectangle++(0.008,0.008) (-0.161,0.932)rectangle++(0.008,0.008)
  (-0.07,0.989)rectangle++(0.008,0.008) (0.063,0.978)rectangle++(0.008,0.008) (-0.185,0.882)rectangle++(0.008,0.008) (-0.158,0.811)rectangle++(0.008,0.008) (-0.14,0.889)rectangle++(0.008,0.008)
  (-0.061,0.971)rectangle++(0.008,0.008) (-0.091,0.986)rectangle++(0.008,0.008) (0.183,0.836)rectangle++(0.008,0.008) (0.158,0.906)rectangle++(0.008,0.008) (-0.008,0.732)rectangle++(0.008,0.008)
  (-0.006,0.905)rectangle++(0.008,0.008) (-0.141,0.947)rectangle++(0.008,0.008) (-0.061,0.995)rectangle++(0.008,0.008) (-0.165,0.956)rectangle++(0.008,0.008) (-0.158,0.85)rectangle++(0.008,0.008)
  (0.15,0.93)rectangle++(0.008,0.008) (-0.19,0.906)rectangle++(0.008,0.008) (-0.082,0.978)rectangle++(0.008,0.008) (-0.037,1.009)rectangle++(0.008,0.008) (-0.178,0.931)rectangle++(0.008,0.008)
  (-0.077,0.989)rectangle++(0.008,0.008) (0.131,0.974)rectangle++(0.008,0.008) (0.078,0.991)rectangle++(0.008,0.008) (-0.112,0.981)rectangle++(0.008,0.008) (-0.049,1.01)rectangle++(0.008,0.008)
  (-0.008,0.718)rectangle++(0.008,0.008) (0.158,0.878)rectangle++(0.008,0.008) (0.064,0.967)rectangle++(0.008,0.008) (0.194,0.877)rectangle++(0.008,0.008) (-0.098,0.971)rectangle++(0.008,0.008)
  (-0.005,0.768)rectangle++(0.008,0.008) (0.165,0.88)rectangle++(0.008,0.008) (-0.182,0.906)rectangle++(0.008,0.008) (0.178,0.909)rectangle++(0.008,0.008) (0.168,0.92)rectangle++(0.008,0.008)
  (-0.006,0.697)rectangle++(0.008,0.008) (-0.102,0.958)rectangle++(0.008,0.008) (-0.081,0.989)rectangle++(0.008,0.008) (-0.036,1.013)rectangle++(0.008,0.008) (-0.018,1.023)rectangle++(0.008,0.008)
  (0.19,0.863)rectangle++(0.008,0.008) (0.179,0.869)rectangle++(0.008,0.008) (0.073,0.962)rectangle++(0.008,0.008) (-0.134,0.969)rectangle++(0.008,0.008) (-0.175,0.94)rectangle++(0.008,0.008)
  (-0.064,0.848)rectangle++(0.008,0.008) (-0.149,0.935)rectangle++(0.008,0.008) (0.058,0.909)rectangle++(0.008,0.008) (0.145,0.938)rectangle++(0.008,0.008) (0.059,0.992)rectangle++(0.008,0.008)
  (0.157,0.954)rectangle++(0.008,0.008) (-0.007,0.758)rectangle++(0.008,0.008) (-0.005,0.916)rectangle++(0.008,0.008) (0.017,0.8)rectangle++(0.008,0.008) (0.005,0.934)rectangle++(0.008,0.008)
  (0.186,0.835)rectangle++(0.008,0.008) (0.158,0.905)rectangle++(0.008,0.008) (0.064,0.978)rectangle++(0.008,0.008) (-0.13,0.933)rectangle++(0.008,0.008) (-0.133,0.937)rectangle++(0.008,0.008)
  (-0.058,0.991)rectangle++(0.008,0.008) (0.142,0.938)rectangle++(0.008,0.008) (-0.007,0.755)rectangle++(0.008,0.008) (-0.011,0.948)rectangle++(0.008,0.008) (0.138,0.956)rectangle++(0.008,0.008)
  (-0.131,0.957)rectangle++(0.008,0.008) (-0.114,0.863)rectangle++(0.008,0.008) (-0.05,0.96)rectangle++(0.008,0.008) (-0.006,0.761)rectangle++(0.008,0.008) (0.02,0.962)rectangle++(0.008,0.008)
  (-0.008,0.749)rectangle++(0.008,0.008) (0.013,0.951)rectangle++(0.008,0.008) (0.003,0.997)rectangle++(0.008,0.008) (-0.117,0.867)rectangle++(0.008,0.008) (-0.051,0.962)rectangle++(0.008,0.008)
  (-0.065,0.85)rectangle++(0.008,0.008) (-0.056,0.834)rectangle++(0.008,0.008) (0.175,0.916)rectangle++(0.008,0.008) (0.168,0.922)rectangle++(0.008,0.008) (-0.11,0.976)rectangle++(0.008,0.008)
  (-0.048,1.008)rectangle++(0.008,0.008) (-0.023,1.021)rectangle++(0.008,0.008) (0.006,0.978)rectangle++(0.008,0.008) (-0.149,0.962)rectangle++(0.008,0.008) (-0.014,0.785)rectangle++(0.008,0.008)
  (0.112,0.925)rectangle++(0.008,0.008) (0.045,0.986)rectangle++(0.008,0.008) (0.002,0.708)rectangle++(0.008,0.008) (0.181,0.861)rectangle++(0.008,0.008) (-0.003,0.756)rectangle++(0.008,0.008)
  (0.057,0.894)rectangle++(0.008,0.008) (-0.039,0.813)rectangle++(0.008,0.008) (0.112,0.935)rectangle++(0.008,0.008) (0.045,0.99)rectangle++(0.008,0.008) (-0.172,0.934)rectangle++(0.008,0.008)
  (-0.03,0.981)rectangle++(0.008,0.008) (-0.057,0.845)rectangle++(0.008,0.008) (0.033,0.819)rectangle++(0.008,0.008) (0.126,0.925)rectangle++(0.008,0.008) (-0.01,0.701)rectangle++(0.008,0.008)
  (-0.004,0.72)rectangle++(0.008,0.008) (0.105,0.911)rectangle++(0.008,0.008) (-0.003,0.757)rectangle++(0.008,0.008) (-0.003,0.916)rectangle++(0.008,0.008) (-0.004,0.982)rectangle++(0.008,0.008)
  (0.19,0.834)rectangle++(0.008,0.008) (0.176,0.885)rectangle++(0.008,0.008) (0.072,0.969)rectangle++(0.008,0.008) (0.076,0.965)rectangle++(0.008,0.008) (-0.103,0.98)rectangle++(0.008,0.008)
  (0.191,0.9)rectangle++(0.008,0.008) (0.078,0.976)rectangle++(0.008,0.008) (-0.003,0.769)rectangle++(0.008,0.008) (-0.12,0.921)rectangle++(0.008,0.008) (-0.053,0.984)rectangle++(0.008,0.008)
  (0.089,0.859)rectangle++(0.008,0.008) (-0.141,0.936)rectangle++(0.008,0.008) (-0.006,0.762)rectangle++(0.008,0.008) (0.051,0.842)rectangle++(0.008,0.008) (-0.009,0.719)rectangle++(0.008,0.008)
  (-0.151,0.937)rectangle++(0.008,0.008) (0.155,0.936)rectangle++(0.008,0.008) (0.063,0.991)rectangle++(0.008,0.008) (0.024,1.014)rectangle++(0.008,0.008) (-0.182,0.872)rectangle++(0.008,0.008)
  (-0.102,0.944)rectangle++(0.008,0.008) (0.133,0.89)rectangle++(0.008,0.008) (-0.002,0.754)rectangle++(0.008,0.008) (-0.003,0.914)rectangle++(0.008,0.008) (-0.008,0.74)rectangle++(0.008,0.008)
  (-0.006,0.908)rectangle++(0.008,0.008) (0.148,0.896)rectangle++(0.008,0.008) (-0.046,0.98)rectangle++(0.008,0.008) (0.053,0.992)rectangle++(0.008,0.008) (0.02,1.014)rectangle++(0.008,0.008)
  (-0.117,0.979)rectangle++(0.008,0.008) (-0.173,0.944)rectangle++(0.008,0.008) (-0.075,0.994)rectangle++(0.008,0.008) (-0.167,0.953)rectangle++(0.008,0.008) (-0.202,0.857)rectangle++(0.008,0.008)
  (0.003,0.958)rectangle++(0.008,0.008) (-0.115,0.975)rectangle++(0.008,0.008) (0.121,0.977)rectangle++(0.008,0.008) (-0.049,0.992)rectangle++(0.008,0.008) (0.053,0.994)rectangle++(0.008,0.008)
  (-0.138,0.94)rectangle++(0.008,0.008) (0.098,0.98)rectangle++(0.008,0.008) (0.163,0.946)rectangle++(0.008,0.008) (-0.192,0.908)rectangle++(0.008,0.008) (-0.187,0.823)rectangle++(0.008,0.008)
  (-0.081,0.943)rectangle++(0.008,0.008) (0.113,0.934)rectangle++(0.008,0.008) (0.171,0.922)rectangle++(0.008,0.008) (-0.042,0.984)rectangle++(0.008,0.008) (-0.02,1.011)rectangle++(0.008,0.008)
  (-0.101,0.983)rectangle++(0.008,0.008) (0.1,0.972)rectangle++(0.008,0.008) (-0.197,0.901)rectangle++(0.008,0.008) (-0.068,0.907)rectangle++(0.008,0.008) (0.057,0.899)rectangle++(0.008,0.008)
  (-0.128,0.955)rectangle++(0.008,0.008) (-0.203,0.869)rectangle++(0.008,0.008) (-0.014,0.966)rectangle++(0.008,0.008) (0.058,0.99)rectangle++(0.008,0.008) (-0.133,0.939)rectangle++(0.008,0.008)
  (-0.058,0.992)rectangle++(0.008,0.008) (-0.075,0.994)rectangle++(0.008,0.008) (0.132,0.975)rectangle++(0.008,0.008) (0.053,1.007)rectangle++(0.008,0.008) (-0.145,0.975)rectangle++(0.008,0.008)
  (-0.063,1.007)rectangle++(0.008,0.008) (-0.01,0.715)rectangle++(0.008,0.008) (0.077,0.858)rectangle++(0.008,0.008) (-0.056,0.975)rectangle++(0.008,0.008) (-0.026,1.007)rectangle++(0.008,0.008)
  (-0.013,1.021)rectangle++(0.008,0.008) (-0.008,1.026)rectangle++(0.008,0.008) (-0.006,1.029)rectangle++(0.008,0.008) (0.151,0.971)rectangle++(0.008,0.008) (-0.002,0.768)rectangle++(0.008,0.008)
  (0.042,0.828)rectangle++(0.008,0.008) (-0.022,0.964)rectangle++(0.008,0.008) (0.138,0.961)rectangle++(0.008,0.008) (0.168,0.936)rectangle++(0.008,0.008) (-0.112,0.979)rectangle++(0.008,0.008)
  (-0.049,1.009)rectangle++(0.008,0.008) (-0.023,1.021)rectangle++(0.008,0.008) (-0.012,1.027)rectangle++(0.008,0.008) (-0.209,0.887)rectangle++(0.008,0.008) (-0.09,0.97)rectangle++(0.008,0.008)
  (-0.169,0.904)rectangle++(0.008,0.008) (-0.09,0.981)rectangle++(0.008,0.008) (0.167,0.858)rectangle++(0.008,0.008) (0.176,0.854)rectangle++(0.008,0.008) (-0.08,0.98)rectangle++(0.008,0.008)
  (0.048,0.992)rectangle++(0.008,0.008) (0.018,1.014)rectangle++(0.008,0.008) (0.005,1.024)rectangle++(0.008,0.008) (-0.015,0.978)rectangle++(0.008,0.008) (-0.048,0.936)rectangle++(0.008,0.008)
  (0.128,0.959)rectangle++(0.008,0.008) (-0.009,0.739)rectangle++(0.008,0.008) (0.125,0.931)rectangle++(0.008,0.008) (-0.062,0.848)rectangle++(0.008,0.008) (-0.028,0.954)rectangle++(0.008,0.008)
  (0.056,0.85)rectangle++(0.008,0.008) (0.021,0.955)rectangle++(0.008,0.008) (-0.007,0.755)rectangle++(0.008,0.008) (-0.119,0.918)rectangle++(0.008,0.008) (0.112,0.967)rectangle++(0.008,0.008)
  (0.164,0.941)rectangle++(0.008,0.008) (-0.079,0.928)rectangle++(0.008,0.008) (-0.192,0.831)rectangle++(0.008,0.008) (-0.006,0.743)rectangle++(0.008,0.008) (-0.004,0.728)rectangle++(0.008,0.008)
  (-0.126,0.922)rectangle++(0.008,0.008) (0.034,0.827)rectangle++(0.008,0.008) (-0.004,0.743)rectangle++(0.008,0.008) (0.045,0.981)rectangle++(0.008,0.008) (-0.2,0.893)rectangle++(0.008,0.008)
  (-0.197,0.857)rectangle++(0.008,0.008) (-0.085,0.958)rectangle++(0.008,0.008) (0.048,0.989)rectangle++(0.008,0.008) (0.131,0.887)rectangle++(0.008,0.008) (0.105,0.968)rectangle++(0.008,0.008)
  (-0.179,0.922)rectangle++(0.008,0.008) (-0.005,0.958)rectangle++(0.008,0.008) (0.14,0.957)rectangle++(0.008,0.008) (-0.047,0.989)rectangle++(0.008,0.008) (0.053,0.993)rectangle++(0.008,0.008)
  (0.02,1.015)rectangle++(0.008,0.008) (-0.174,0.934)rectangle++(0.008,0.008) (0.081,0.981)rectangle++(0.008,0.008) (-0.182,0.905)rectangle++(0.008,0.008) (0.182,0.894)rectangle++(0.008,0.008)
  (-0.05,0.836)rectangle++(0.008,0.008) (-0.028,0.969)rectangle++(0.008,0.008) (-0.014,1.005)rectangle++(0.008,0.008) (-0.022,0.796)rectangle++(0.008,0.008) (-0.012,0.932)rectangle++(0.008,0.008)
  (-0.007,0.989)rectangle++(0.008,0.008) (-0.067,0.993)rectangle++(0.008,0.008) (0.193,0.839)rectangle++(0.008,0.008) (0.114,0.957)rectangle++(0.008,0.008) (-0.195,0.902)rectangle++(0.008,0.008)
  (-0.084,0.977)rectangle++(0.008,0.008) (-0.005,0.769)rectangle++(0.008,0.008) (-0.084,0.966)rectangle++(0.008,0.008) (-0.037,1.003)rectangle++(0.008,0.008) (-0.067,0.875)rectangle++(0.008,0.008)
  (-0.188,0.833)rectangle++(0.008,0.008) (-0.081,0.948)rectangle++(0.008,0.008) (-0.134,0.938)rectangle++(0.008,0.008) (0.157,0.837)rectangle++(0.008,0.008) (-0.137,0.885)rectangle++(0.008,0.008)
  (-0.146,0.935)rectangle++(0.008,0.008) (-0.004,0.748)rectangle++(0.008,0.008) (-0.004,0.912)rectangle++(0.008,0.008) (-0.153,0.938)rectangle++(0.008,0.008) (-0.014,0.786)rectangle++(0.008,0.008)
  (-0.119,0.865)rectangle++(0.008,0.008) (-0.19,0.833)rectangle++(0.008,0.008) (0.089,0.976)rectangle++(0.008,0.008) (0.035,1.008)rectangle++(0.008,0.008) (-0.004,0.667)rectangle++(0.008,0.008)
  (-0.004,0.878)rectangle++(0.008,0.008) (-0.095,0.86)rectangle++(0.008,0.008) (-0.129,0.945)rectangle++(0.008,0.008) (-0.057,0.995)rectangle++(0.008,0.008) (-0.026,1.015)rectangle++(0.008,0.008)
  (0.124,0.946)rectangle++(0.008,0.008) (-0.122,0.974)rectangle++(0.008,0.008) (0.058,0.889)rectangle++(0.008,0.008) (0.022,0.971)rectangle++(0.008,0.008) (0.007,1.005)rectangle++(0.008,0.008)
  (-0.153,0.968)rectangle++(0.008,0.008) (0.075,0.984)rectangle++(0.008,0.008) (0.001,0.745)rectangle++(0.008,0.008) (-0.002,0.91)rectangle++(0.008,0.008) (-0.003,0.98)rectangle++(0.008,0.008)
  (-0.151,0.961)rectangle++(0.008,0.008) (-0.003,0.738)rectangle++(0.008,0.008) (-0.004,0.727)rectangle++(0.008,0.008) (-0.054,0.837)rectangle++(0.008,0.008) (-0.025,0.949)rectangle++(0.008,0.008)
  (0.051,0.923)rectangle++(0.008,0.008) (-0.004,0.655)rectangle++(0.008,0.008) (-0.004,0.873)rectangle++(0.008,0.008) (-0.187,0.883)rectangle++(0.008,0.008) (-0.195,0.858)rectangle++(0.008,0.008)
  (0.04,0.935)rectangle++(0.008,0.008) (0.015,0.991)rectangle++(0.008,0.008) (0.149,0.961)rectangle++(0.008,0.008) (0.175,0.922)rectangle++(0.008,0.008) (-0.108,0.977)rectangle++(0.008,0.008)
  (0.044,0.992)rectangle++(0.008,0.008) (0.056,0.98)rectangle++(0.008,0.008) (0,0.729)rectangle++(0.008,0.008) (-0.002,0.904)rectangle++(0.008,0.008) (-0.191,0.884)rectangle++(0.008,0.008)
  (-0.118,0.865)rectangle++(0.008,0.008) (0.11,0.928)rectangle++(0.008,0.008) (-0.009,0.706)rectangle++(0.008,0.008) (0.156,0.878)rectangle++(0.008,0.008) (0.063,0.966)rectangle++(0.008,0.008)
  (-0.002,0.739)rectangle++(0.008,0.008) (-0.003,0.908)rectangle++(0.008,0.008) (-0.141,0.948)rectangle++(0.008,0.008) (-0.185,0.854)rectangle++(0.008,0.008) (-0.08,0.956)rectangle++(0.008,0.008)
  (-0.036,0.999)rectangle++(0.008,0.008) (-0.206,0.883)rectangle++(0.008,0.008) (0.023,0.978)rectangle++(0.008,0.008) (-0.133,0.952)rectangle++(0.008,0.008) (0.02,0.944)rectangle++(0.008,0.008)
  (0.062,0.987)rectangle++(0.008,0.008) (-0.058,0.993)rectangle++(0.008,0.008) (-0.075,0.994)rectangle++(0.008,0.008) (-0.175,0.81)rectangle++(0.008,0.008) (-0.052,0.83)rectangle++(0.008,0.008)
  (0.112,0.94)rectangle++(0.008,0.008) (0.124,0.937)rectangle++(0.008,0.008) (0.05,0.991)rectangle++(0.008,0.008) (0.001,0.747)rectangle++(0.008,0.008) (-0.002,0.911)rectangle++(0.008,0.008)
  (0.185,0.884)rectangle++(0.008,0.008) (-0.007,0.754)rectangle++(0.008,0.008) (-0.067,0.96)rectangle++(0.008,0.008) (-0.031,1.001)rectangle++(0.008,0.008) (0.196,0.891)rectangle++(0.008,0.008)
  (0.189,0.856)rectangle++(0.008,0.008) (-0.183,0.913)rectangle++(0.008,0.008) (-0.006,0.758)rectangle++(0.008,0.008) (0.11,0.919)rectangle++(0.008,0.008) (0.044,0.984)rectangle++(0.008,0.008)
  (-0.006,0.763)rectangle++(0.008,0.008) (-0.005,0.918)rectangle++(0.008,0.008) (-0.193,0.884)rectangle++(0.008,0.008) (-0.005,0.941)rectangle++(0.008,0.008) (-0.01,0.965)rectangle++(0.008,0.008)
  (0.01,0.975)rectangle++(0.008,0.008) (0.145,0.958)rectangle++(0.008,0.008) (0.059,1)rectangle++(0.008,0.008) (0.159,0.956)rectangle++(0.008,0.008) (0.064,0.999)rectangle++(0.008,0.008)
  (0.025,1.017)rectangle++(0.008,0.008) (0.185,0.88)rectangle++(0.008,0.008) (0.165,0.914)rectangle++(0.008,0.008) (0.067,0.981)rectangle++(0.008,0.008) (-0.066,0.852)rectangle++(0.008,0.008)
  (0.122,0.932)rectangle++(0.008,0.008) (0.049,0.989)rectangle++(0.008,0.008) (0.018,1.013)rectangle++(0.008,0.008) (0.062,0.997)rectangle++(0.008,0.008) (0.068,0.994)rectangle++(0.008,0.008)
  (-0.006,0.725)rectangle++(0.008,0.008) (-0.067,0.956)rectangle++(0.008,0.008) (0.066,0.922)rectangle++(0.008,0.008) (0.148,0.94)rectangle++(0.008,0.008) (-0.105,0.981)rectangle++(0.008,0.008)
  (0.076,0.933)rectangle++(0.008,0.008) (-0.129,0.964)rectangle++(0.008,0.008) (0.058,0.889)rectangle++(0.008,0.008) (0.022,0.971)rectangle++(0.008,0.008) (-0.005,0.76)rectangle++(0.008,0.008)
  (-0.005,0.917)rectangle++(0.008,0.008) (-0.193,0.884)rectangle++(0.008,0.008) (-0.009,0.73)rectangle++(0.008,0.008) (-0.019,0.965)rectangle++(0.008,0.008) (-0.2,0.884)rectangle++(0.008,0.008)
  (-0.067,0.969)rectangle++(0.008,0.008) (0.13,0.969)rectangle++(0.008,0.008) (-0.164,0.826)rectangle++(0.008,0.008) (0.032,0.82)rectangle++(0.008,0.008) (-0.062,0.969)rectangle++(0.008,0.008)
  (0.041,0.935)rectangle++(0.008,0.008) (-0.028,0.944)rectangle++(0.008,0.008) (-0.014,0.994)rectangle++(0.008,0.008) (-0.008,1.015)rectangle++(0.008,0.008) (-0.006,1.024)rectangle++(0.008,0.008)
  (-0.005,1.028)rectangle++(0.008,0.008) (-0.004,1.029)rectangle++(0.008,0.008) (-0.004,1.03)rectangle++(0.008,0.008) (-0.118,0.867)rectangle++(0.008,0.008) (-0.179,0.883)rectangle++(0.008,0.008)
  (-0.171,0.916)rectangle++(0.008,0.008) (-0.003,0.93)rectangle++(0.008,0.008) (0.136,0.951)rectangle++(0.008,0.008) (-0.007,0.756)rectangle++(0.008,0.008) (0.02,0.961)rectangle++(0.008,0.008)
  (-0.01,0.702)rectangle++(0.008,0.008) (-0.166,0.921)rectangle++(0.008,0.008) (-0.072,0.984)rectangle++(0.008,0.008) (-0.033,1.011)rectangle++(0.008,0.008) (-0.016,1.022)rectangle++(0.008,0.008)
  (0.148,0.971)rectangle++(0.008,0.008) (0.171,0.937)rectangle++(0.008,0.008) (-0.054,0.831)rectangle++(0.008,0.008) (-0.113,0.863)rectangle++(0.008,0.008) (-0.05,0.96)rectangle++(0.008,0.008)
  (-0.191,0.903)rectangle++(0.008,0.008) (0.177,0.91)rectangle++(0.008,0.008) (0.106,0.864)rectangle++(0.008,0.008) (-0.127,0.955)rectangle++(0.008,0.008) (0.116,0.975)rectangle++(0.008,0.008)
  (0.179,0.905)rectangle++(0.008,0.008) (-0.009,0.962)rectangle++(0.008,0.008) (-0.006,1.002)rectangle++(0.008,0.008) (-0.163,0.839)rectangle++(0.008,0.008) (-0.039,0.813)rectangle++(0.008,0.008)
  (-0.006,0.724)rectangle++(0.008,0.008) (-0.146,0.929)rectangle++(0.008,0.008) (-0.2,0.865)rectangle++(0.008,0.008) (0.089,0.969)rectangle++(0.008,0.008) (0.035,1.005)rectangle++(0.008,0.008)
  (-0.201,0.837)rectangle++(0.008,0.008) (-0.087,0.949)rectangle++(0.008,0.008) (0.055,0.847)rectangle++(0.008,0.008) (-0.164,0.925)rectangle++(0.008,0.008) (-0.168,0.856)rectangle++(0.008,0.008)
  (-0.066,0.905)rectangle++(0.008,0.008) (0.038,0.978)rectangle++(0.008,0.008) (-0.192,0.855)rectangle++(0.008,0.008) (-0.169,0.909)rectangle++(0.008,0.008) (0.002,0.963)rectangle++(0.008,0.008)
  (-0.066,0.99)rectangle++(0.008,0.008) (-0.165,0.954)rectangle++(0.008,0.008) (0.109,0.98)rectangle++(0.008,0.008) (0.043,1.009)rectangle++(0.008,0.008) (0.2,0.881)rectangle++(0.008,0.008)
  (0.081,0.968)rectangle++(0.008,0.008) (-0.133,0.972)rectangle++(0.008,0.008) (0.101,0.973)rectangle++(0.008,0.008) (0.001,0.742)rectangle++(0.008,0.008) (-0.005,0.695)rectangle++(0.008,0.008)
  (0.058,0.901)rectangle++(0.008,0.008) (0.022,0.976)rectangle++(0.008,0.008) (0.007,1.008)rectangle++(0.008,0.008) (-0.175,0.933)rectangle++(0.008,0.008) (0.182,0.908)rectangle++(0.008,0.008)
  (-0.128,0.959)rectangle++(0.008,0.008) (-0.056,1)rectangle++(0.008,0.008) (-0.11,0.862)rectangle++(0.008,0.008) (0.106,0.954)rectangle++(0.008,0.008) (0.114,0.969)rectangle++(0.008,0.008)
  (-0.006,0.956)rectangle++(0.008,0.008) (-0.005,0.999)rectangle++(0.008,0.008) (-0.004,1.017)rectangle++(0.008,0.008) (-0.157,0.968)rectangle++(0.008,0.008) (-0.192,0.823)rectangle++(0.008,0.008)
  (-0.164,0.903)rectangle++(0.008,0.008) (0.036,0.981)rectangle++(0.008,0.008) (0.195,0.881)rectangle++(0.008,0.008) (-0.11,0.948)rectangle++(0.008,0.008) (-0.167,0.939)rectangle++(0.008,0.008)
  (-0.066,0.907)rectangle++(0.008,0.008) (-0.144,0.892)rectangle++(0.008,0.008) (0.039,0.98)rectangle++(0.008,0.008) (0.014,1.009)rectangle++(0.008,0.008) (0.151,0.964)rectangle++(0.008,0.008)
  (0.061,1.003)rectangle++(0.008,0.008) (0.022,0.972)rectangle++(0.008,0.008) (0.147,0.956)rectangle++(0.008,0.008) (0.06,0.999)rectangle++(0.008,0.008) (0.06,0.853)rectangle++(0.008,0.008)
  (0.164,0.92)rectangle++(0.008,0.008) (0.167,0.924)rectangle++(0.008,0.008) (-0.109,0.863)rectangle++(0.008,0.008) (-0.048,0.96)rectangle++(0.008,0.008) (0.132,0.964)rectangle++(0.008,0.008)
  (0.174,0.923)rectangle++(0.008,0.008) (0.117,0.963)rectangle++(0.008,0.008) (0.047,1.002)rectangle++(0.008,0.008) (-0.145,0.973)rectangle++(0.008,0.008) (0.115,0.981)rectangle++(0.008,0.008)
  (0.046,1.01)rectangle++(0.008,0.008) (0.017,1.022)rectangle++(0.008,0.008) (0.027,0.817)rectangle++(0.008,0.008) (0.125,0.925)rectangle++(0.008,0.008) (-0.065,0.886)rectangle++(0.008,0.008)
  (0.18,0.827)rectangle++(0.008,0.008) (0.073,0.945)rectangle++(0.008,0.008) (0.047,0.833)rectangle++(0.008,0.008) (0.017,0.948)rectangle++(0.008,0.008) (-0.2,0.865)rectangle++(0.008,0.008)
  (-0.086,0.961)rectangle++(0.008,0.008) (0.102,0.979)rectangle++(0.008,0.008) (0.041,1.009)rectangle++(0.008,0.008) (0.015,1.021)rectangle++(0.008,0.008) (-0.154,0.972)rectangle++(0.008,0.008)
  (0.114,0.982)rectangle++(0.008,0.008) (-0.129,0.979)rectangle++(0.008,0.008) (-0.057,1.009)rectangle++(0.008,0.008) (0,0.954)rectangle++(0.008,0.008) (0.188,0.835)rectangle++(0.008,0.008)
  (0.159,0.905)rectangle++(0.008,0.008) (0.164,0.922)rectangle++(0.008,0.008) (-0.111,0.975)rectangle++(0.008,0.008) (0.159,0.936)rectangle++(0.008,0.008) (0.152,0.933)rectangle++(0.008,0.008)
  (-0.115,0.976)rectangle++(0.008,0.008) (0.19,0.902)rectangle++(0.008,0.008) (0.086,0.981)rectangle++(0.008,0.008) (-0.188,0.85)rectangle++(0.008,0.008) (-0.167,0.909)rectangle++(0.008,0.008)
  (-0.182,0.883)rectangle++(0.008,0.008) (0.095,0.969)rectangle++(0.008,0.008) (-0.053,0.991)rectangle++(0.008,0.008) (0.142,0.92)rectangle++(0.008,0.008) (0.058,0.984)rectangle++(0.008,0.008)
  (0.158,0.81)rectangle++(0.008,0.008) (0.111,0.957)rectangle++(0.008,0.008) (-0.004,0.663)rectangle++(0.008,0.008) (-0.004,0.876)rectangle++(0.008,0.008) (0.188,0.871)rectangle++(0.008,0.008)
  (0.077,0.963)rectangle++(0.008,0.008) (-0.01,0.722)rectangle++(0.008,0.008) (-0.006,0.901)rectangle++(0.008,0.008) (-0.14,0.946)rectangle++(0.008,0.008) (-0.061,0.995)rectangle++(0.008,0.008)
  (-0.028,1.016)rectangle++(0.008,0.008) (-0.127,0.95)rectangle++(0.008,0.008) (0.003,0.968)rectangle++(0.008,0.008) (-0.029,0.804)rectangle++(0.008,0.008) (-0.106,0.865)rectangle++(0.008,0.008)
  (-0.19,0.869)rectangle++(0.008,0.008) (0.019,0.802)rectangle++(0.008,0.008) (0.081,0.953)rectangle++(0.008,0.008) (0.156,0.943)rectangle++(0.008,0.008) (-0.006,0.702)rectangle++(0.008,0.008)
  (-0.002,0.698)rectangle++(0.008,0.008) (-0.024,0.947)rectangle++(0.008,0.008) (0.135,0.958)rectangle++(0.008,0.008) (0.054,1)rectangle++(0.008,0.008) (0.02,1.018)rectangle++(0.008,0.008)
  (0.062,0.997)rectangle++(0.008,0.008) (0.034,0.815)rectangle++(0.008,0.008) (0.012,0.94)rectangle++(0.008,0.008) (-0.143,0.956)rectangle++(0.008,0.008) (-0.175,0.936)rectangle++(0.008,0.008)
  (-0.076,0.991)rectangle++(0.008,0.008) (0.162,0.935)rectangle++(0.008,0.008) (-0.062,0.849)rectangle++(0.008,0.008) (-0.028,0.954)rectangle++(0.008,0.008) (-0.014,0.998)rectangle++(0.008,0.008)
  (0.144,0.966)rectangle++(0.008,0.008) (0.118,0.968)rectangle++(0.008,0.008) (0.047,1.004)rectangle++(0.008,0.008) (0.157,0.958)rectangle++(0.008,0.008) (0.064,1)rectangle++(0.008,0.008)
  (-0.067,0.854)rectangle++(0.008,0.008) (-0.03,0.956)rectangle++(0.008,0.008) (-0.008,0.746)rectangle++(0.008,0.008) (0.108,0.917)rectangle++(0.008,0.008) (0.043,0.983)rectangle++(0.008,0.008)
  (0.001,0.736)rectangle++(0.008,0.008) (-0.18,0.843)rectangle++(0.008,0.008) (-0.078,0.952)rectangle++(0.008,0.008) (-0.162,0.945)rectangle++(0.008,0.008) (-0.048,0.834)rectangle++(0.008,0.008)
  (-0.023,0.948)rectangle++(0.008,0.008) (-0.02,0.975)rectangle++(0.008,0.008) (-0.011,1.007)rectangle++(0.008,0.008) (-0.007,1.021)rectangle++(0.008,0.008) (0.111,0.979)rectangle++(0.008,0.008)
  (0.198,0.871)rectangle++(0.008,0.008) (0.053,0.846)rectangle++(0.008,0.008) (-0.128,0.96)rectangle++(0.008,0.008) (-0.056,1.001)rectangle++(0.008,0.008) (-0.135,0.949)rectangle++(0.008,0.008)
  (-0.006,0.764)rectangle++(0.008,0.008) (0.111,0.92)rectangle++(0.008,0.008) (-0.004,0.664)rectangle++(0.008,0.008) (-0.004,0.877)rectangle++(0.008,0.008) (0.059,0.977)rectangle++(0.008,0.008)
  (-0.044,0.829)rectangle++(0.008,0.008) (-0.098,0.861)rectangle++(0.008,0.008) (0.009,0.964)rectangle++(0.008,0.008) (0.002,1.002)rectangle++(0.008,0.008) (0.001,0.751)rectangle++(0.008,0.008)
  (-0.067,0.898)rectangle++(0.008,0.008) (-0.061,0.847)rectangle++(0.008,0.008) (-0.193,0.846)rectangle++(0.008,0.008) (-0.084,0.953)rectangle++(0.008,0.008) (-0.037,0.998)rectangle++(0.008,0.008)
  (-0.09,0.987)rectangle++(0.008,0.008) (-0.169,0.95)rectangle++(0.008,0.008) (-0.09,0.988)rectangle++(0.008,0.008) (0.153,0.937)rectangle++(0.008,0.008) (0.062,0.991)rectangle++(0.008,0.008)
  (-0.061,0.982)rectangle++(0.008,0.008) (-0.028,1.01)rectangle++(0.008,0.008) (-0.005,0.691)rectangle++(0.008,0.008) (0.031,0.967)rectangle++(0.008,0.008) (-0.019,0.969)rectangle++(0.008,0.008)
  (-0.175,0.858)rectangle++(0.008,0.008) (-0.166,0.912)rectangle++(0.008,0.008) (-0.072,0.981)rectangle++(0.008,0.008) (-0.165,0.951)rectangle++(0.008,0.008) (0.036,0.988)rectangle++(0.008,0.008)
  (0.002,0.695)rectangle++(0.008,0.008) (0.175,0.828)rectangle++(0.008,0.008) (-0.068,0.904)rectangle++(0.008,0.008) (-0.091,0.983)rectangle++(0.008,0.008) (-0.122,0.972)rectangle++(0.008,0.008)
  (-0.01,0.711)rectangle++(0.008,0.008) (-0.007,0.896)rectangle++(0.008,0.008) (-0.137,0.883)rectangle++(0.008,0.008) (0.095,0.976)rectangle++(0.008,0.008) (0.038,1.008)rectangle++(0.008,0.008)
  (0.013,1.021)rectangle++(0.008,0.008) (-0.118,0.979)rectangle++(0.008,0.008) (-0.173,0.944)rectangle++(0.008,0.008) (-0.075,0.994)rectangle++(0.008,0.008) (-0.001,0.945)rectangle++(0.008,0.008)
  (0.139,0.954)rectangle++(0.008,0.008) (0.056,0.998)rectangle++(0.008,0.008) (0.157,0.918)rectangle++(0.008,0.008) (0.165,0.925)rectangle++(0.008,0.008) (-0.111,0.976)rectangle++(0.008,0.008)
  (-0.049,1.008)rectangle++(0.008,0.008) (-0.023,1.021)rectangle++(0.008,0.008) (0.146,0.972)rectangle++(0.008,0.008) (0.059,1.006)rectangle++(0.008,0.008) (0.19,0.868)rectangle++(0.008,0.008)
  (0.164,0.911)rectangle++(0.008,0.008) (0.165,0.923)rectangle++(0.008,0.008) (-0.188,0.908)rectangle++(0.008,0.008) (-0.001,0.755)rectangle++(0.008,0.008) (-0.146,0.931)rectangle++(0.008,0.008)
  (-0.006,0.761)rectangle++(0.008,0.008) (0.001,0.7)rectangle++(0.008,0.008) (0.08,0.973)rectangle++(0.008,0.008) (0.004,0.967)rectangle++(0.008,0.008) (-0.001,1.004)rectangle++(0.008,0.008)
  (0.148,0.965)rectangle++(0.008,0.008) (0.06,1.003)rectangle++(0.008,0.008) (-0.203,0.896)rectangle++(0.008,0.008) (-0.088,0.974)rectangle++(0.008,0.008) (-0.039,1.007)rectangle++(0.008,0.008)
  (-0.207,0.883)rectangle++(0.008,0.008) (-0.008,0.749)rectangle++(0.008,0.008) (0.051,0.84)rectangle++(0.008,0.008) (-0.12,0.943)rectangle++(0.008,0.008) (-0.177,0.923)rectangle++(0.008,0.008)
  (-0.077,0.985)rectangle++(0.008,0.008) (-0.034,1.011)rectangle++(0.008,0.008) (-0.017,1.022)rectangle++(0.008,0.008) (0.018,0.944)rectangle++(0.008,0.008) (0.142,0.951)rectangle++(0.008,0.008)
  (0.049,0.835)rectangle++(0.008,0.008) (0.066,0.972)rectangle++(0.008,0.008) (0.026,1.006)rectangle++(0.008,0.008) (-0.182,0.922)rectangle++(0.008,0.008) (0.071,0.93)rectangle++(0.008,0.008)
  (0.028,0.988)rectangle++(0.008,0.008) (0.009,1.013)rectangle++(0.008,0.008) (0.002,1.023)rectangle++(0.008,0.008) (-0.157,0.97)rectangle++(0.008,0.008) (-0.206,0.865)rectangle++(0.008,0.008)
  (0.035,0.824)rectangle++(0.008,0.008) (-0.09,0.981)rectangle++(0.008,0.008) (0.191,0.899)rectangle++(0.008,0.008) (-0.184,0.911)rectangle++(0.008,0.008) (0.105,0.865)rectangle++(0.008,0.008)
  (0.147,0.923)rectangle++(0.008,0.008) (0.186,0.876)rectangle++(0.008,0.008) (0.027,0.942)rectangle++(0.008,0.008) (-0.063,0.987)rectangle++(0.008,0.008) (-0.029,1.012)rectangle++(0.008,0.008)
  (-0.014,1.023)rectangle++(0.008,0.008) (-0.008,1.027)rectangle++(0.008,0.008) (0.061,0.901)rectangle++(0.008,0.008) (-0.189,0.893)rectangle++(0.008,0.008) (0.06,0.906)rectangle++(0.008,0.008)
  (-0.017,0.968)rectangle++(0.008,0.008) (-0.009,1.004)rectangle++(0.008,0.008) (-0.156,0.965)rectangle++(0.008,0.008) (-0.205,0.865)rectangle++(0.008,0.008) (0.025,0.943)rectangle++(0.008,0.008)
  (0.061,0.893)rectangle++(0.008,0.008) (-0.002,0.7)rectangle++(0.008,0.008) (-0.067,0.952)rectangle++(0.008,0.008) (-0.2,0.877)rectangle++(0.008,0.008) (0.09,0.971)rectangle++(0.008,0.008)
  (0.036,1.005)rectangle++(0.008,0.008) (0.168,0.923)rectangle++(0.008,0.008) (-0.007,0.708)rectangle++(0.008,0.008) (0.102,0.909)rectangle++(0.008,0.008) (-0.189,0.899)rectangle++(0.008,0.008)
  (-0.005,0.756)rectangle++(0.008,0.008) (0.155,0.917)rectangle++(0.008,0.008) (0.124,0.943)rectangle++(0.008,0.008) (-0.01,0.949)rectangle++(0.008,0.008) (-0.148,0.955)rectangle++(0.008,0.008)
  (-0.063,0.977)rectangle++(0.008,0.008) (0.057,0.852)rectangle++(0.008,0.008) (0.022,0.955)rectangle++(0.008,0.008) (-0.063,0.846)rectangle++(0.008,0.008) (0.17,0.919)rectangle++(0.008,0.008)
  (-0.103,0.864)rectangle++(0.008,0.008) (-0.045,0.961)rectangle++(0.008,0.008) (0.19,0.892)rectangle++(0.008,0.008) (0.188,0.857)rectangle++(0.008,0.008) (0.095,0.86)rectangle++(0.008,0.008)
  (0.038,0.959)rectangle++(0.008,0.008) (0.066,0.98)rectangle++(0.008,0.008) (-0.112,0.979)rectangle++(0.008,0.008) (0.194,0.843)rectangle++(0.008,0.008) (0.114,0.957)rectangle++(0.008,0.008)
  (-0.118,0.866)rectangle++(0.008,0.008) (-0.052,0.961)rectangle++(0.008,0.008) (-0.024,1.001)rectangle++(0.008,0.008) (-0.004,0.773)rectangle++(0.008,0.008) (-0.001,0.75)rectangle++(0.008,0.008)
  (-0.003,0.913)rectangle++(0.008,0.008) (-0.028,0.944)rectangle++(0.008,0.008) (-0.014,0.994)rectangle++(0.008,0.008) (-0.046,0.987)rectangle++(0.008,0.008) (-0.17,0.81)rectangle++(0.008,0.008)
  (-0.056,0.842)rectangle++(0.008,0.008) (-0.005,0.745)rectangle++(0.008,0.008) (-0.002,0.748)rectangle++(0.008,0.008) (0.109,0.917)rectangle++(0.008,0.008) (0.156,0.932)rectangle++(0.008,0.008)
  (0.063,0.989)rectangle++(0.008,0.008) (0.024,1.013)rectangle++(0.008,0.008) (0.008,1.023)rectangle++(0.008,0.008) (0.152,0.968)rectangle++(0.008,0.008) (0.062,1.004)rectangle++(0.008,0.008)
  (0.079,0.938)rectangle++(0.008,0.008) (0.169,0.924)rectangle++(0.008,0.008) (0.069,0.986)rectangle++(0.008,0.008) (0.069,0.993)rectangle++(0.008,0.008) (0.047,0.836)rectangle++(0.008,0.008)
  (0.018,0.949)rectangle++(0.008,0.008) (-0.143,0.959)rectangle++(0.008,0.008) (-0.065,0.849)rectangle++(0.008,0.008) (-0.144,0.927)rectangle++(0.008,0.008) (-0.009,0.732)rectangle++(0.008,0.008)
  (-0.086,0.935)rectangle++(0.008,0.008) (0.101,0.977)rectangle++(0.008,0.008) (0.051,0.838)rectangle++(0.008,0.008) (0.019,0.949)rectangle++(0.008,0.008) (0.166,0.928)rectangle++(0.008,0.008)
  (0.193,0.862)rectangle++(0.008,0.008) (0.079,0.96)rectangle++(0.008,0.008) (0.031,1.001)rectangle++(0.008,0.008) (-0.006,0.704)rectangle++(0.008,0.008) (-0.163,0.84)rectangle++(0.008,0.008)
  (-0.071,0.951)rectangle++(0.008,0.008) (0.039,0.938)rectangle++(0.008,0.008) (0.181,0.824)rectangle++(0.008,0.008) (-0.146,0.906)rectangle++(0.008,0.008) (-0.03,0.978)rectangle++(0.008,0.008)
  (-0.015,1.009)rectangle++(0.008,0.008) (-0.158,0.965)rectangle++(0.008,0.008) (0.058,0.911)rectangle++(0.008,0.008) (-0.006,0.699)rectangle++(0.008,0.008) (0.094,0.958)rectangle++(0.008,0.008)
  (-0.002,0.761)rectangle++(0.008,0.008) (0.044,0.834)rectangle++(0.008,0.008) (-0.12,0.941)rectangle++(0.008,0.008) (-0.168,0.936)rectangle++(0.008,0.008) (-0.073,0.991)rectangle++(0.008,0.008)
  (0.132,0.974)rectangle++(0.008,0.008) (0.169,0.94)rectangle++(0.008,0.008) (0.069,0.992)rectangle++(0.008,0.008) (-0.106,0.865)rectangle++(0.008,0.008) (0.042,0.836)rectangle++(0.008,0.008)
  (0.015,0.949)rectangle++(0.008,0.008) (0.004,0.996)rectangle++(0.008,0.008) (-0.001,1.016)rectangle++(0.008,0.008) (0.037,0.94)rectangle++(0.008,0.008) (0.19,0.88)rectangle++(0.008,0.008)
  (0.072,0.856)rectangle++(0.008,0.008) (0.121,0.946)rectangle++(0.008,0.008) (0.162,0.936)rectangle++(0.008,0.008) (0.172,0.88)rectangle++(0.008,0.008) (-0.103,0.969)rectangle++(0.008,0.008)
  (-0.045,1.004)rectangle++(0.008,0.008) (-0.004,0.66)rectangle++(0.008,0.008) (-0.104,0.9)rectangle++(0.008,0.008) (0.075,0.932)rectangle++(0.008,0.008) (0.029,0.989)rectangle++(0.008,0.008)
  (-0.147,0.968)rectangle++(0.008,0.008) (0.028,0.984)rectangle++(0.008,0.008) (0.039,0.986)rectangle++(0.008,0.008) (0.196,0.881)rectangle++(0.008,0.008) (0.08,0.968)rectangle++(0.008,0.008)
  (-0.134,0.972)rectangle++(0.008,0.008) (-0.175,0.94)rectangle++(0.008,0.008) (-0.202,0.862)rectangle++(0.008,0.008) (-0.087,0.96)rectangle++(0.008,0.008) (0.125,0.97)rectangle++(0.008,0.008)
  (-0.165,0.937)rectangle++(0.008,0.008) (-0.176,0.929)rectangle++(0.008,0.008) (-0.024,0.798)rectangle++(0.008,0.008) (-0.05,0.836)rectangle++(0.008,0.008) (-0.024,0.949)rectangle++(0.008,0.008)
  (0.046,0.839)rectangle++(0.008,0.008) (0.132,0.927)rectangle++(0.008,0.008) (-0.082,0.984)rectangle++(0.008,0.008) (-0.037,1.011)rectangle++(0.008,0.008) (-0.018,1.022)rectangle++(0.008,0.008)
  (0.073,0.985)rectangle++(0.008,0.008) (0.028,1.011)rectangle++(0.008,0.008) (0.054,0.848)rectangle++(0.008,0.008) (0.051,0.846)rectangle++(0.008,0.008) (0.167,0.831)rectangle++(0.008,0.008)
  (-0.096,0.958)rectangle++(0.008,0.008) (-0.043,1)rectangle++(0.008,0.008) (-0.02,1.018)rectangle++(0.008,0.008) (-0.011,1.025)rectangle++(0.008,0.008) (-0.004,0.778)rectangle++(0.008,0.008)
  (-0.106,0.963)rectangle++(0.008,0.008) (0.002,0.724)rectangle++(0.008,0.008) (-0.009,0.715)rectangle++(0.008,0.008) (-0.006,0.898)rectangle++(0.008,0.008) (-0.002,0.757)rectangle++(0.008,0.008)
  (0.111,0.918)rectangle++(0.008,0.008) (0.05,0.845)rectangle++(0.008,0.008) (0.019,0.952)rectangle++(0.008,0.008) (0.167,0.858)rectangle++(0.008,0.008) (-0.101,0.963)rectangle++(0.008,0.008)
  (0.123,0.973)rectangle++(0.008,0.008) (0.164,0.81)rectangle++(0.008,0.008) (0.12,0.935)rectangle++(0.008,0.008) (-0.182,0.817)rectangle++(0.008,0.008) (-0.057,0.838)rectangle++(0.008,0.008)
  (0.142,0.855)rectangle++(0.008,0.008) (0.057,0.957)rectangle++(0.008,0.008) (-0.064,0.85)rectangle++(0.008,0.008) (-0.029,0.955)rectangle++(0.008,0.008) (-0.133,0.864)rectangle++(0.008,0.008)
  (0.041,0.976)rectangle++(0.008,0.008) (-0.003,0.735)rectangle++(0.008,0.008) (-0.062,0.978)rectangle++(0.008,0.008) (0.132,0.97)rectangle++(0.008,0.008) (0.168,0.939)rectangle++(0.008,0.008)
  (0.187,0.876)rectangle++(0.008,0.008) (0.165,0.913)rectangle++(0.008,0.008) (0.061,0.899)rectangle++(0.008,0.008) (0.144,0.936)rectangle++(0.008,0.008) (-0.004,0.93)rectangle++(0.008,0.008)
  (-0.137,0.883)rectangle++(0.008,0.008) (-0.06,0.969)rectangle++(0.008,0.008) (-0.064,0.847)rectangle++(0.008,0.008) (-0.075,0.973)rectangle++(0.008,0.008) (-0.165,0.837)rectangle++(0.008,0.008)
  (0.169,0.904)rectangle++(0.008,0.008) (0.069,0.977)rectangle++(0.008,0.008) (-0.125,0.922)rectangle++(0.008,0.008) (-0.055,0.985)rectangle++(0.008,0.008) (-0.058,0.846)rectangle++(0.008,0.008)
  (0.078,0.98)rectangle++(0.008,0.008) (0.051,0.846)rectangle++(0.008,0.008) (-0.059,0.973)rectangle++(0.008,0.008) (0.051,0.991)rectangle++(0.008,0.008) (0.156,0.955)rectangle++(0.008,0.008)
  (0.063,0.999)rectangle++(0.008,0.008) (0.001,0.747)rectangle++(0.008,0.008) (-0.17,0.879)rectangle++(0.008,0.008) (-0.168,0.917)rectangle++(0.008,0.008) (0.103,0.974)rectangle++(0.008,0.008)
  (-0.197,0.901)rectangle++(0.008,0.008) (-0.085,0.976)rectangle++(0.008,0.008) (0.103,0.98)rectangle++(0.008,0.008) (0.051,0.847)rectangle++(0.008,0.008) (0.134,0.927)rectangle++(0.008,0.008)
  (0.054,0.987)rectangle++(0.008,0.008) (0.02,1.012)rectangle++(0.008,0.008) (0.055,0.85)rectangle++(0.008,0.008) (-0.005,0.692)rectangle++(0.008,0.008) (-0.196,0.86)rectangle++(0.008,0.008)
  (-0.008,0.958)rectangle++(0.008,0.008) (0.023,0.979)rectangle++(0.008,0.008) (0.11,0.974)rectangle++(0.008,0.008) (0.197,0.871)rectangle++(0.008,0.008) (-0.003,0.69)rectangle++(0.008,0.008)
  (0.08,0.954)rectangle++(0.008,0.008) (-0.004,0.748)rectangle++(0.008,0.008) (-0.004,0.66)rectangle++(0.008,0.008) (-0.145,0.858)rectangle++(0.008,0.008) (-0.006,0.755)rectangle++(0.008,0.008)
  (-0.004,0.73)rectangle++(0.008,0.008) (-0.067,0.957)rectangle++(0.008,0.008) (-0.03,1)rectangle++(0.008,0.008) (0.141,0.969)rectangle++(0.008,0.008) (-0.122,0.981)rectangle++(0.008,0.008)
  (0.121,0.979)rectangle++(0.008,0.008) (-0.197,0.904)rectangle++(0.008,0.008) (0.08,0.981)rectangle++(0.008,0.008) (0.16,0.949)rectangle++(0.008,0.008) (-0.173,0.812)rectangle++(0.008,0.008)
  (-0.004,0.656)rectangle++(0.008,0.008) (-0.007,0.755)rectangle++(0.008,0.008) (-0.191,0.862)rectangle++(0.008,0.008) (-0.083,0.96)rectangle++(0.008,0.008) (-0.2,0.851)rectangle++(0.008,0.008)
  (-0.18,0.913)rectangle++(0.008,0.008) (-0.006,0.758)rectangle++(0.008,0.008) (-0.067,0.961)rectangle++(0.008,0.008) (-0.183,0.919)rectangle++(0.008,0.008) (0.1,0.976)rectangle++(0.008,0.008)
  (0.115,0.97)rectangle++(0.008,0.008) (-0.197,0.903)rectangle++(0.008,0.008) (0.177,0.91)rectangle++(0.008,0.008) (-0.007,0.755)rectangle++(0.008,0.008) (-0.007,0.731)rectangle++(0.008,0.008)
  (-0.005,0.905)rectangle++(0.008,0.008) (-0.002,0.758)rectangle++(0.008,0.008) (-0.123,0.948)rectangle++(0.008,0.008) (0.003,0.968)rectangle++(0.008,0.008) (0.057,0.894)rectangle++(0.008,0.008)
  (0.022,0.973)rectangle++(0.008,0.008) (0.007,1.006)rectangle++(0.008,0.008) (0.15,0.965)rectangle++(0.008,0.008) (-0.164,0.938)rectangle++(0.008,0.008) (-0.202,0.863)rectangle++(0.008,0.008)
  (-0.127,0.93)rectangle++(0.008,0.008) (-0.168,0.933)rectangle++(0.008,0.008) (-0.001,0.758)rectangle++(0.008,0.008) (-0.003,0.916)rectangle++(0.008,0.008) (-0.002,0.759)rectangle++(0.008,0.008)
  (0.052,0.842)rectangle++(0.008,0.008) (0.053,0.977)rectangle++(0.008,0.008) (0.08,0.979)rectangle++(0.008,0.008) (0.058,0.851)rectangle++(0.008,0.008) (-0.182,0.891)rectangle++(0.008,0.008)
  (-0.079,0.972)rectangle++(0.008,0.008) (0,0.738)rectangle++(0.008,0.008) (0.085,0.967)rectangle++(0.008,0.008) (-0.012,0.971)rectangle++(0.008,0.008) (-0.152,0.958)rectangle++(0.008,0.008)
  (-0.028,0.982)rectangle++(0.008,0.008) (-0.014,1.01)rectangle++(0.008,0.008) (-0.008,1.022)rectangle++(0.008,0.008) (0.149,0.97)rectangle++(0.008,0.008) (-0.195,0.907)rectangle++(0.008,0.008)
  (0.052,0.847)rectangle++(0.008,0.008) (-0.003,0.746)rectangle++(0.008,0.008) (0.12,0.931)rectangle++(0.008,0.008) (0.048,0.989)rectangle++(0.008,0.008) (0.111,0.865)rectangle++(0.008,0.008)
  (0.148,0.922)rectangle++(0.008,0.008) (0.164,0.927)rectangle++(0.008,0.008) (-0.111,0.976)rectangle++(0.008,0.008) (-0.154,0.855)rectangle++(0.008,0.008) (-0.067,0.957)rectangle++(0.008,0.008)
  (-0.076,0.989)rectangle++(0.008,0.008) (-0.031,0.982)rectangle++(0.008,0.008) (-0.071,0.992)rectangle++(0.008,0.008) (-0.166,0.954)rectangle++(0.008,0.008) (0.044,0.839)rectangle++(0.008,0.008)
  (0.001,0.687)rectangle++(0.008,0.008) (-0.002,0.886)rectangle++(0.008,0.008) (-0.003,0.97)rectangle++(0.008,0.008) (-0.15,0.959)rectangle++(0.008,0.008) (-0.177,0.935)rectangle++(0.008,0.008)
  (0.143,0.896)rectangle++(0.008,0.008) (0.058,0.974)rectangle++(0.008,0.008) (0.022,1.007)rectangle++(0.008,0.008) (0.007,1.02)rectangle++(0.008,0.008) (0.2,0.887)rectangle++(0.008,0.008)
  (0.12,0.923)rectangle++(0.008,0.008) (-0.185,0.819)rectangle++(0.008,0.008) (-0.16,0.847)rectangle++(0.008,0.008) (0.17,0.903)rectangle++(0.008,0.008) (0.104,0.862)rectangle++(0.008,0.008)
  (-0.113,0.955)rectangle++(0.008,0.008) (0.187,0.901)rectangle++(0.008,0.008) (-0.06,0.846)rectangle++(0.008,0.008) (0.093,0.979)rectangle++(0.008,0.008) (0.184,0.839)rectangle++(0.008,0.008)
  (-0.094,0.962)rectangle++(0.008,0.008) (-0.08,0.99)rectangle++(0.008,0.008) (0.193,0.84)rectangle++(0.008,0.008) (0.041,0.834)rectangle++(0.008,0.008) (0.015,0.948)rectangle++(0.008,0.008)
  (-0.064,0.988)rectangle++(0.008,0.008) (-0.067,0.852)rectangle++(0.008,0.008) (-0.03,0.956)rectangle++(0.008,0.008) (-0.153,0.953)rectangle++(0.008,0.008) (0.081,0.981)rectangle++(0.008,0.008)
  (-0.136,0.974)rectangle++(0.008,0.008) (-0.176,0.94)rectangle++(0.008,0.008) (-0.174,0.819)rectangle++(0.008,0.008) (-0.054,0.979)rectangle++(0.008,0.008) (-0.074,0.992)rectangle++(0.008,0.008)
  (0.001,0.736)rectangle++(0.008,0.008) (0.176,0.827)rectangle++(0.008,0.008) (0.072,0.945)rectangle++(0.008,0.008) (0.028,0.995)rectangle++(0.008,0.008) (0.053,0.84)rectangle++(0.008,0.008)
  (0.02,0.951)rectangle++(0.008,0.008) (-0.01,0.965)rectangle++(0.008,0.008) (0.14,0.959)rectangle++(0.008,0.008) (0.168,0.936)rectangle++(0.008,0.008) (-0.112,0.979)rectangle++(0.008,0.008)
  (-0.049,1.009)rectangle++(0.008,0.008) (0.194,0.867)rectangle++(0.008,0.008) (0.111,0.967)rectangle++(0.008,0.008) (0.045,1.004)rectangle++(0.008,0.008) (0.066,0.995)rectangle++(0.008,0.008)
  (0.025,1.016)rectangle++(0.008,0.008) (-0.063,0.997)rectangle++(0.008,0.008) (-0.03,0.944)rectangle++(0.008,0.008) (0.116,0.941)rectangle++(0.008,0.008) (0.161,0.936)rectangle++(0.008,0.008)
  (0.065,0.991)rectangle++(0.008,0.008) (-0.17,0.935)rectangle++(0.008,0.008) (-0.059,0.97)rectangle++(0.008,0.008) (0.041,0.934)rectangle++(0.008,0.008) (0.015,0.99)rectangle++(0.008,0.008)
  (0.183,0.835)rectangle++(0.008,0.008) (0.075,0.948)rectangle++(0.008,0.008) (-0.001,0.947)rectangle++(0.008,0.008) (-0.003,0.995)rectangle++(0.008,0.008) (-0.008,0.751)rectangle++(0.008,0.008)
  (-0.001,0.731)rectangle++(0.008,0.008) (-0.18,0.903)rectangle++(0.008,0.008) (0.098,0.973)rectangle++(0.008,0.008) (-0.01,0.925)rectangle++(0.008,0.008) (-0.006,0.722)rectangle++(0.008,0.008)
  (-0.005,0.901)rectangle++(0.008,0.008) (-0.03,0.979)rectangle++(0.008,0.008) (0.008,0.975)rectangle++(0.008,0.008) (-0.065,0.991)rectangle++(0.008,0.008) (-0.004,0.746)rectangle++(0.008,0.008)
  (-0.105,0.961)rectangle++(0.008,0.008) (-0.122,0.97)rectangle++(0.008,0.008) (-0.054,1.005)rectangle++(0.008,0.008) (-0.025,1.02)rectangle++(0.008,0.008) (-0.013,1.026)rectangle++(0.008,0.008)
  (-0.16,0.969)rectangle++(0.008,0.008) (0.168,0.919)rectangle++(0.008,0.008) (0.068,0.984)rectangle++(0.008,0.008) (0.164,0.896)rectangle++(0.008,0.008) (0.038,0.821)rectangle++(0.008,0.008)
  (-0.12,0.937)rectangle++(0.008,0.008) (-0.066,0.891)rectangle++(0.008,0.008) (0.118,0.953)rectangle++(0.008,0.008) (-0.069,0.89)rectangle++(0.008,0.008) (0.052,0.845)rectangle++(0.008,0.008)
  (0.067,0.973)rectangle++(0.008,0.008) (-0.042,0.986)rectangle++(0.008,0.008) (0.05,0.983)rectangle++(0.008,0.008) (0.019,1.01)rectangle++(0.008,0.008) (-0.044,0.83)rectangle++(0.008,0.008)
  (-0.137,0.926)rectangle++(0.008,0.008) (0.182,0.903)rectangle++(0.008,0.008) (0.074,0.977)rectangle++(0.008,0.008) (0.171,0.927)rectangle++(0.008,0.008) (-0.11,0.977)rectangle++(0.008,0.008)
  (0.095,0.863)rectangle++(0.008,0.008) (-0.018,0.967)rectangle++(0.008,0.008) (-0.152,0.814)rectangle++(0.008,0.008) (-0.105,0.944)rectangle++(0.008,0.008) (0.115,0.924)rectangle++(0.008,0.008)
  (0.046,0.986)rectangle++(0.008,0.008) (0.017,1.012)rectangle++(0.008,0.008) (-0.069,0.973)rectangle++(0.008,0.008) (-0.163,0.95)rectangle++(0.008,0.008) (-0.171,0.833)rectangle++(0.008,0.008)
  (0.148,0.93)rectangle++(0.008,0.008) (0.192,0.864)rectangle++(0.008,0.008) (-0.065,0.895)rectangle++(0.008,0.008) (0.043,0.831)rectangle++(0.008,0.008) (0.016,0.947)rectangle++(0.008,0.008)
  (0.004,0.995)rectangle++(0.008,0.008) (-0.174,0.932)rectangle++(0.008,0.008) (-0.183,0.919)rectangle++(0.008,0.008) (-0.006,0.759)rectangle++(0.008,0.008) (-0.088,0.976)rectangle++(0.008,0.008)
  (-0.073,0.982)rectangle++(0.008,0.008) (0.131,0.972)rectangle++(0.008,0.008) (0.078,0.991)rectangle++(0.008,0.008) (-0.137,0.976)rectangle++(0.008,0.008) (0.172,0.92)rectangle++(0.008,0.008)
  (-0.042,0.984)rectangle++(0.008,0.008) (0.037,0.82)rectangle++(0.008,0.008) (0.013,0.942)rectangle++(0.008,0.008) (0.146,0.854)rectangle++(0.008,0.008) (-0.104,0.96)rectangle++(0.008,0.008)
  (0.147,0.854)rectangle++(0.008,0.008) (0.17,0.915)rectangle++(0.008,0.008) (0.069,0.982)rectangle++(0.008,0.008) (-0.2,0.897)rectangle++(0.008,0.008) (-0.086,0.974)rectangle++(0.008,0.008)
  (-0.107,0.979)rectangle++(0.008,0.008) (-0.056,0.843)rectangle++(0.008,0.008) (0.113,0.943)rectangle++(0.008,0.008) (0.045,0.994)rectangle++(0.008,0.008) (0.066,0.994)rectangle++(0.008,0.008)
  (0.025,1.015)rectangle++(0.008,0.008) (0.009,0.98)rectangle++(0.008,0.008) (0.002,1.009)rectangle++(0.008,0.008) (0.06,0.996)rectangle++(0.008,0.008) (0.169,0.881)rectangle++(0.008,0.008)
  (0.162,0.916)rectangle++(0.008,0.008) (-0.004,0.741)rectangle++(0.008,0.008) (-0.098,0.862)rectangle++(0.008,0.008) (-0.157,0.897)rectangle++(0.008,0.008) (-0.117,0.968)rectangle++(0.008,0.008)
  (-0.112,0.865)rectangle++(0.008,0.008) (-0.049,0.961)rectangle++(0.008,0.008) (-0.065,0.901)rectangle++(0.008,0.008) (0.12,0.955)rectangle++(0.008,0.008) (-0.01,0.725)rectangle++(0.008,0.008)
  (-0.006,0.902)rectangle++(0.008,0.008) (0.131,0.947)rectangle++(0.008,0.008) (0.053,0.995)rectangle++(0.008,0.008) (-0.143,0.973)rectangle++(0.008,0.008) (-0.006,0.762)rectangle++(0.008,0.008)
  (-0.119,0.919)rectangle++(0.008,0.008) (-0.165,0.932)rectangle++(0.008,0.008) (-0.071,0.989)rectangle++(0.008,0.008) (-0.179,0.928)rectangle++(0.008,0.008) (-0.064,0.847)rectangle++(0.008,0.008)
  (0.035,0.98)rectangle++(0.008,0.008) (0.175,0.817)rectangle++(0.008,0.008) (-0.109,0.966)rectangle++(0.008,0.008) (-0.048,1.003)rectangle++(0.008,0.008) (-0.066,0.895)rectangle++(0.008,0.008)
  (-0.03,0.973)rectangle++(0.008,0.008) (0.138,0.964)rectangle++(0.008,0.008) (-0.195,0.906)rectangle++(0.008,0.008) (-0.177,0.918)rectangle++(0.008,0.008) (-0.067,0.902)rectangle++(0.008,0.008)
  (-0.03,0.977)rectangle++(0.008,0.008) (-0.047,0.986)rectangle++(0.008,0.008) (-0.022,1.012)rectangle++(0.008,0.008) (0.145,0.97)rectangle++(0.008,0.008) (-0.005,0.76)rectangle++(0.008,0.008)
  (0.155,0.917)rectangle++(0.008,0.008) (-0.075,0.923)rectangle++(0.008,0.008) (-0.034,0.985)rectangle++(0.008,0.008) (-0.004,0.771)rectangle++(0.008,0.008) (0.138,0.912)rectangle++(0.008,0.008)
  (-0.047,0.982)rectangle++(0.008,0.008) (-0.063,0.98)rectangle++(0.008,0.008) (0.035,0.986)rectangle++(0.008,0.008) (0.143,0.814)rectangle++(0.008,0.008) (0.176,0.861)rectangle++(0.008,0.008)
  (0.001,0.692)rectangle++(0.008,0.008) (-0.022,0.795)rectangle++(0.008,0.008) (-0.011,0.932)rectangle++(0.008,0.008) (-0.067,0.906)rectangle++(0.008,0.008) (-0.056,0.843)rectangle++(0.008,0.008)
  (-0.002,0.723)rectangle++(0.008,0.008) (-0.003,0.901)rectangle++(0.008,0.008) (-0.004,0.976)rectangle++(0.008,0.008) (0.098,0.861)rectangle++(0.008,0.008) (-0.053,0.975)rectangle++(0.008,0.008)
  (-0.203,0.88)rectangle++(0.008,0.008) (-0.001,0.754)rectangle++(0.008,0.008) (-0.117,0.918)rectangle++(0.008,0.008) (-0.055,0.983)rectangle++(0.008,0.008) (-0.025,1.011)rectangle++(0.008,0.008)
  (0.025,0.983)rectangle++(0.008,0.008) (0.008,1.01)rectangle++(0.008,0.008) (0.04,0.834)rectangle++(0.008,0.008) (-0.152,0.937)rectangle++(0.008,0.008) (-0.087,0.938)rectangle++(0.008,0.008)
  (-0.161,0.941)rectangle++(0.008,0.008) (-0.089,0.986)rectangle++(0.008,0.008) (0.167,0.881)rectangle++(0.008,0.008) (0.048,0.839)rectangle++(0.008,0.008) (-0.003,0.936)rectangle++(0.008,0.008)
  (0.137,0.953)rectangle++(0.008,0.008) (0.055,0.998)rectangle++(0.008,0.008) (0.001,0.747)rectangle++(0.008,0.008) (-0.002,0.911)rectangle++(0.008,0.008) (-0.067,0.982)rectangle++(0.008,0.008)
  (-0.164,0.952)rectangle++(0.008,0.008) (-0.066,0.978)rectangle++(0.008,0.008) (0.131,0.97)rectangle++(0.008,0.008) (0.168,0.939)rectangle++(0.008,0.008) (-0.177,0.923)rectangle++(0.008,0.008)
  (-0.063,0.846)rectangle++(0.008,0.008) (-0.005,0.746)rectangle++(0.008,0.008) (-0.126,0.936)rectangle++(0.008,0.008) (-0.055,0.991)rectangle++(0.008,0.008) (-0.144,0.892)rectangle++(0.008,0.008)
  (-0.063,0.972)rectangle++(0.008,0.008) (0.051,0.991)rectangle++(0.008,0.008) (-0.003,0.772)rectangle++(0.008,0.008) (-0.004,0.733)rectangle++(0.008,0.008) (-0.004,0.905)rectangle++(0.008,0.008)
  (-0.14,0.947)rectangle++(0.008,0.008) (-0.124,0.967)rectangle++(0.008,0.008) (-0.124,0.969)rectangle++(0.008,0.008) (-0.054,1.005)rectangle++(0.008,0.008) (-0.002,0.729)rectangle++(0.008,0.008)
  (-0.003,0.904)rectangle++(0.008,0.008) (-0.004,0.977)rectangle++(0.008,0.008) (0.058,0.852)rectangle++(0.008,0.008) (0.022,0.955)rectangle++(0.008,0.008) (-0.063,0.989)rectangle++(0.008,0.008)
  (-0.164,0.954)rectangle++(0.008,0.008) (-0.071,0.998)rectangle++(0.008,0.008) (0.194,0.897)rectangle++(0.008,0.008) (-0.146,0.913)rectangle++(0.008,0.008) (0.011,0.959)rectangle++(0.008,0.008)
  (-0.006,0.934)rectangle++(0.008,0.008) (-0.005,0.93)rectangle++(0.008,0.008) (-0.198,0.862)rectangle++(0.008,0.008) (0.06,0.906)rectangle++(0.008,0.008) (0.138,0.913)rectangle++(0.008,0.008)
  (0.112,0.866)rectangle++(0.008,0.008) (-0.004,0.944)rectangle++(0.008,0.008) (0.138,0.954)rectangle++(0.008,0.008) (-0.002,0.697)rectangle++(0.008,0.008) (0.024,0.958)rectangle++(0.008,0.008)
  (0.116,0.965)rectangle++(0.008,0.008) (0.001,0.726)rectangle++(0.008,0.008) (-0.05,0.835)rectangle++(0.008,0.008) (-0.009,0.713)rectangle++(0.008,0.008) (-0.006,0.897)rectangle++(0.008,0.008)
  (-0.061,0.845)rectangle++(0.008,0.008) (-0.028,0.953)rectangle++(0.008,0.008) (-0.014,0.998)rectangle++(0.008,0.008) (-0.205,0.886)rectangle++(0.008,0.008) (0.091,0.973)rectangle++(0.008,0.008)
  (0.001,0.743)rectangle++(0.008,0.008) (0.104,0.965)rectangle++(0.008,0.008) (-0.128,0.975)rectangle++(0.008,0.008) (-0.067,0.965)rectangle++(0.008,0.008) (0.181,0.89)rectangle++(0.008,0.008)
  (0.085,0.979)rectangle++(0.008,0.008) (0.161,0.948)rectangle++(0.008,0.008) (-0.065,0.964)rectangle++(0.008,0.008) (-0.029,1.002)rectangle++(0.008,0.008) (-0.004,0.774)rectangle++(0.008,0.008)
  (-0.004,0.923)rectangle++(0.008,0.008) (-0.194,0.885)rectangle++(0.008,0.008) (-0.084,0.969)rectangle++(0.008,0.008) (-0.038,1.005)rectangle++(0.008,0.008) (-0.018,1.02)rectangle++(0.008,0.008)
  (-0.01,1.026)rectangle++(0.008,0.008) (0.058,0.999)rectangle++(0.008,0.008) (0.022,1.017)rectangle++(0.008,0.008) (0.007,1.025)rectangle++(0.008,0.008) (-0.137,0.882)rectangle++(0.008,0.008)
  (0.103,0.963)rectangle++(0.008,0.008) (0.041,1.002)rectangle++(0.008,0.008) (-0.004,0.755)rectangle++(0.008,0.008) (-0.004,0.915)rectangle++(0.008,0.008) (-0.004,0.982)rectangle++(0.008,0.008)
  (0.144,0.962)rectangle++(0.008,0.008) (-0.001,0.76)rectangle++(0.008,0.008) (-0.067,0.961)rectangle++(0.008,0.008) (-0.161,0.948)rectangle++(0.008,0.008) (0.154,0.937)rectangle++(0.008,0.008)
  (0.063,0.991)rectangle++(0.008,0.008) (0.024,1.014)rectangle++(0.008,0.008) (0.154,0.964)rectangle++(0.008,0.008) (-0.007,0.946)rectangle++(0.008,0.008) (-0.005,0.995)rectangle++(0.008,0.008)
  (0.167,0.932)rectangle++(0.008,0.008) (-0.189,0.908)rectangle++(0.008,0.008) (-0.066,0.907)rectangle++(0.008,0.008) (0.182,0.893)rectangle++(0.008,0.008) (0.074,0.973)rectangle++(0.008,0.008)
  (0.157,0.948)rectangle++(0.008,0.008) (0.184,0.898)rectangle++(0.008,0.008) (0.136,0.858)rectangle++(0.008,0.008) (-0.106,0.959)rectangle++(0.008,0.008) (-0.009,0.741)rectangle++(0.008,0.008)
  (0.173,0.843)rectangle++(0.008,0.008) (0.07,0.952)rectangle++(0.008,0.008) (-0.11,0.978)rectangle++(0.008,0.008) (-0.084,0.935)rectangle++(0.008,0.008) (-0.068,0.908)rectangle++(0.008,0.008)
  (0.05,0.982)rectangle++(0.008,0.008) (0.067,0.992)rectangle++(0.008,0.008) (-0.14,0.974)rectangle++(0.008,0.008) (-0.061,1.007)rectangle++(0.008,0.008) (0.06,0.907)rectangle++(0.008,0.008)
  (0,0.927)rectangle++(0.008,0.008) (0.134,0.891)rectangle++(0.008,0.008) (0.084,0.965)rectangle++(0.008,0.008) (0.033,1.003)rectangle++(0.008,0.008) (0.012,1.019)rectangle++(0.008,0.008)
  (0.169,0.933)rectangle++(0.008,0.008) (-0.166,0.809)rectangle++(0.008,0.008) (0.165,0.903)rectangle++(0.008,0.008) (0.083,0.981)rectangle++(0.008,0.008) (0.042,0.986)rectangle++(0.008,0.008)
  (0.196,0.881)rectangle++(0.008,0.008) (-0.181,0.911)rectangle++(0.008,0.008) (-0.001,0.755)rectangle++(0.008,0.008) (0.186,0.85)rectangle++(0.008,0.008) (0.119,0.933)rectangle++(0.008,0.008)
  (-0.041,0.819)rectangle++(0.008,0.008) (-0.02,0.941)rectangle++(0.008,0.008) (-0.011,0.993)rectangle++(0.008,0.008) (-0.007,1.015)rectangle++(0.008,0.008) (-0.157,0.968)rectangle++(0.008,0.008)
  (-0.179,0.936)rectangle++(0.008,0.008) (-0.006,0.762)rectangle++(0.008,0.008) (-0.119,0.919)rectangle++(0.008,0.008) (-0.052,0.984)rectangle++(0.008,0.008) (-0.024,1.011)rectangle++(0.008,0.008)
  (0.144,0.97)rectangle++(0.008,0.008) (-0.029,0.804)rectangle++(0.008,0.008) (0.112,0.955)rectangle++(0.008,0.008) (-0.126,0.974)rectangle++(0.008,0.008) (0.05,0.837)rectangle++(0.008,0.008)
  (0.019,0.949)rectangle++(0.008,0.008) (-0.004,0.764)rectangle++(0.008,0.008) (0.098,0.962)rectangle++(0.008,0.008) (-0.129,0.973)rectangle++(0.008,0.008) (-0.066,0.851)rectangle++(0.008,0.008)
  (-0.196,0.86)rectangle++(0.008,0.008) (-0.062,0.844)rectangle++(0.008,0.008) (-0.142,0.927)rectangle++(0.008,0.008) (-0.028,0.804)rectangle++(0.008,0.008) (-0.014,0.935)rectangle++(0.008,0.008)
  (0.059,0.909)rectangle++(0.008,0.008) (-0.008,0.748)rectangle++(0.008,0.008) (-0.006,0.912)rectangle++(0.008,0.008) (-0.171,0.926)rectangle++(0.008,0.008) (-0.176,0.926)rectangle++(0.008,0.008)
  (-0.098,0.862)rectangle++(0.008,0.008) (-0.121,0.935)rectangle++(0.008,0.008) (0.134,0.889)rectangle++(0.008,0.008) (0.054,0.971)rectangle++(0.008,0.008) (0.153,0.951)rectangle++(0.008,0.008)
  (0.062,0.997)rectangle++(0.008,0.008) (-0.058,0.994)rectangle++(0.008,0.008) (-0.02,0.98)rectangle++(0.008,0.008) (-0.011,1.009)rectangle++(0.008,0.008) (-0.007,1.021)rectangle++(0.008,0.008)
  (-0.203,0.865)rectangle++(0.008,0.008) (-0.173,0.909)rectangle++(0.008,0.008) (-0.075,0.98)rectangle++(0.008,0.008) (-0.189,0.839)rectangle++(0.008,0.008) (-0.121,0.958)rectangle++(0.008,0.008)
  (-0.053,1)rectangle++(0.008,0.008) (-0.164,0.958)rectangle++(0.008,0.008) (0.002,0.727)rectangle++(0.008,0.008) (0.058,0.898)rectangle++(0.008,0.008) (-0.019,0.957)rectangle++(0.008,0.008)
  (0.054,0.981)rectangle++(0.008,0.008) (-0.141,0.97)rectangle++(0.008,0.008) (-0.006,0.768)rectangle++(0.008,0.008) (0.111,0.921)rectangle++(0.008,0.008) (0.075,0.984)rectangle++(0.008,0.008)
  (0.159,0.95)rectangle++(0.008,0.008) (-0.174,0.858)rectangle++(0.008,0.008) (-0.183,0.826)rectangle++(0.008,0.008) (-0.042,0.82)rectangle++(0.008,0.008) (0.156,0.923)rectangle++(0.008,0.008)
  (0.063,0.985)rectangle++(0.008,0.008) (-0.139,0.972)rectangle++(0.008,0.008) (-0.127,0.963)rectangle++(0.008,0.008) (0.181,0.911)rectangle++(0.008,0.008) (0.05,0.845)rectangle++(0.008,0.008)
  (0.019,0.952)rectangle++(0.008,0.008) (-0.143,0.959)rectangle++(0.008,0.008) (0.039,0.989)rectangle++(0.008,0.008) (0.056,0.98)rectangle++(0.008,0.008) (-0.007,0.927)rectangle++(0.008,0.008)
  (-0.002,0.699)rectangle++(0.008,0.008) (0.102,0.907)rectangle++(0.008,0.008) (-0.197,0.869)rectangle++(0.008,0.008) (0.003,0.959)rectangle++(0.008,0.008) (0.057,0.85)rectangle++(0.008,0.008)
  (-0.006,0.755)rectangle++(0.008,0.008) (-0.005,0.915)rectangle++(0.008,0.008) (-0.142,0.949)rectangle++(0.008,0.008) (0.121,0.931)rectangle++(0.008,0.008) (0.048,0.989)rectangle++(0.008,0.008)
  (0.184,0.904)rectangle++(0.008,0.008) (-0.119,0.867)rectangle++(0.008,0.008) (0.109,0.866)rectangle++(0.008,0.008) (0.043,0.961)rectangle++(0.008,0.008) (0.15,0.951)rectangle++(0.008,0.008)
  (0.019,0.802)rectangle++(0.008,0.008) (0.18,0.892)rectangle++(0.008,0.008) (-0.184,0.909)rectangle++(0.008,0.008) (0.098,0.975)rectangle++(0.008,0.008) (0.039,1.007)rectangle++(0.008,0.008)
  (-0.204,0.893)rectangle++(0.008,0.008) (-0.177,0.915)rectangle++(0.008,0.008) (0.043,0.837)rectangle++(0.008,0.008) (-0.126,0.952)rectangle++(0.008,0.008) (-0.123,0.968)rectangle++(0.008,0.008)
  (0.075,0.984)rectangle++(0.008,0.008) (-0.137,0.974)rectangle++(0.008,0.008) (-0.003,0.683)rectangle++(0.008,0.008) (-0.004,0.885)rectangle++(0.008,0.008) (0.08,0.982)rectangle++(0.008,0.008)
  (-0.185,0.818)rectangle++(0.008,0.008) (-0.08,0.941)rectangle++(0.008,0.008) (-0.036,0.993)rectangle++(0.008,0.008) (-0.099,0.863)rectangle++(0.008,0.008) (-0.044,0.96)rectangle++(0.008,0.008)
  (0.133,0.964)rectangle++(0.008,0.008) (0.167,0.938)rectangle++(0.008,0.008) (0.068,0.991)rectangle++(0.008,0.008) (-0.06,0.845)rectangle++(0.008,0.008) (-0.128,0.947)rectangle++(0.008,0.008)
  (0.137,0.892)rectangle++(0.008,0.008) (-0.005,0.757)rectangle++(0.008,0.008) (-0.118,0.918)rectangle++(0.008,0.008) (0.005,0.966)rectangle++(0.008,0.008) (0.042,0.828)rectangle++(0.008,0.008)
  (-0.007,0.74)rectangle++(0.008,0.008) (-0.068,0.898)rectangle++(0.008,0.008) (0.052,0.845)rectangle++(0.008,0.008) (0.031,0.813)rectangle++(0.008,0.008) (0.078,0.936)rectangle++(0.008,0.008)
  (0.071,0.986)rectangle++(0.008,0.008) (0.159,0.951)rectangle++(0.008,0.008) (-0.002,0.765)rectangle++(0.008,0.008) (-0.119,0.92)rectangle++(0.008,0.008) (-0.052,0.984)rectangle++(0.008,0.008)
  (-0.024,1.011)rectangle++(0.008,0.008) (-0.013,1.022)rectangle++(0.008,0.008) (0.058,0.998)rectangle++(0.008,0.008) (0.012,0.969)rectangle++(0.008,0.008) (0.046,0.927)rectangle++(0.008,0.008);
\fill[piece4]
  (0.017,0.987)rectangle++(0.008,0.008) (-0.149,0.966)rectangle++(0.008,0.008) (0.039,0.99)rectangle++(0.008,0.008) (-0.003,0.771)rectangle++(0.008,0.008) (0.112,0.921)rectangle++(0.008,0.008)
  (0.055,0.918)rectangle++(0.008,0.008) (0.146,0.941)rectangle++(0.008,0.008) (0.059,0.993)rectangle++(0.008,0.008) (-0.182,0.883)rectangle++(0.008,0.008) (0.095,0.969)rectangle++(0.008,0.008)
  (0.016,0.98)rectangle++(0.008,0.008) (0.001,0.687)rectangle++(0.008,0.008) (-0.001,0.735)rectangle++(0.008,0.008) (-0.003,0.906)rectangle++(0.008,0.008) (-0.003,0.978)rectangle++(0.008,0.008)
  (-0.004,1.009)rectangle++(0.008,0.008) (0.11,0.978)rectangle++(0.008,0.008) (-0.202,0.843)rectangle++(0.008,0.008) (-0.186,0.881)rectangle++(0.008,0.008) (-0.065,0.889)rectangle++(0.008,0.008)
  (0.01,0.967)rectangle++(0.008,0.008) (-0.116,0.867)rectangle++(0.008,0.008) (-0.007,0.955)rectangle++(0.008,0.008) (-0.199,0.885)rectangle++(0.008,0.008) (-0.086,0.97)rectangle++(0.008,0.008)
  (0.105,0.969)rectangle++(0.008,0.008) (-0.03,0.805)rectangle++(0.008,0.008) (-0.015,0.936)rectangle++(0.008,0.008) (0.106,0.974)rectangle++(0.008,0.008) (0.042,1.007)rectangle++(0.008,0.008)
  (0.015,1.02)rectangle++(0.008,0.008) (-0.083,0.931)rectangle++(0.008,0.008) (-0.037,0.989)rectangle++(0.008,0.008) (-0.134,0.949)rectangle++(0.008,0.008) (0.114,0.975)rectangle++(0.008,0.008)
  (0.076,0.991)rectangle++(0.008,0.008) (-0.068,0.909)rectangle++(0.008,0.008) (-0.153,0.938)rectangle++(0.008,0.008) (-0.067,0.992)rectangle++(0.008,0.008) (0.133,0.973)rectangle++(0.008,0.008)
  (-0.124,0.981)rectangle++(0.008,0.008) (-0.175,0.943)rectangle++(0.008,0.008) (0.037,0.952)rectangle++(0.008,0.008) (0.013,0.997)rectangle++(0.008,0.008) (-0.008,0.752)rectangle++(0.008,0.008)
  (-0.171,0.878)rectangle++(0.008,0.008) (-0.004,0.661)rectangle++(0.008,0.008) (-0.038,0.817)rectangle++(0.008,0.008) (0.054,0.969)rectangle++(0.008,0.008) (0.02,1.005)rectangle++(0.008,0.008)
  (0.006,1.02)rectangle++(0.008,0.008) (-0.122,0.867)rectangle++(0.008,0.008) (-0.152,0.855)rectangle++(0.008,0.008) (-0.006,0.748)rectangle++(0.008,0.008) (-0.117,0.916)rectangle++(0.008,0.008)
  (-0.164,0.932)rectangle++(0.008,0.008) (0.036,0.985)rectangle++(0.008,0.008) (-0.014,0.974)rectangle++(0.008,0.008) (-0.201,0.885)rectangle++(0.008,0.008) (0.118,0.931)rectangle++(0.008,0.008)
  (0.047,0.989)rectangle++(0.008,0.008) (0.017,1.013)rectangle++(0.008,0.008) (-0.152,0.971)rectangle++(0.008,0.008) (-0.066,1.005)rectangle++(0.008,0.008) (-0.076,0.996)rectangle++(0.008,0.008)
  (0.049,0.994)rectangle++(0.008,0.008) (0.018,1.015)rectangle++(0.008,0.008) (0.169,0.882)rectangle++(0.008,0.008) (-0.105,0.944)rectangle++(0.008,0.008) (-0.047,0.994)rectangle++(0.008,0.008)
  (-0.073,0.994)rectangle++(0.008,0.008) (-0.041,0.984)rectangle++(0.008,0.008) (0.193,0.866)rectangle++(0.008,0.008) (-0.066,0.907)rectangle++(0.008,0.008) (0.182,0.893)rectangle++(0.008,0.008)
  (0.074,0.973)rectangle++(0.008,0.008) (-0.18,0.881)rectangle++(0.008,0.008) (0.102,0.862)rectangle++(0.008,0.008) (0.041,0.96)rectangle++(0.008,0.008) (0.183,0.898)rectangle++(0.008,0.008)
  (-0.057,0.982)rectangle++(0.008,0.008) (-0.009,0.746)rectangle++(0.008,0.008) (-0.006,0.911)rectangle++(0.008,0.008) (-0.153,0.938)rectangle++(0.008,0.008) (0.109,0.976)rectangle++(0.008,0.008)
  (-0.051,0.991)rectangle++(0.008,0.008) (0.052,0.994)rectangle++(0.008,0.008) (-0.008,0.965)rectangle++(0.008,0.008) (-0.008,0.741)rectangle++(0.008,0.008) (-0.004,0.652)rectangle++(0.008,0.008)
  (-0.118,0.947)rectangle++(0.008,0.008) (-0.002,0.755)rectangle++(0.008,0.008) (-0.123,0.948)rectangle++(0.008,0.008) (-0.054,0.996)rectangle++(0.008,0.008) (0,0.953)rectangle++(0.008,0.008)
  (0.165,0.929)rectangle++(0.008,0.008) (0.067,0.988)rectangle++(0.008,0.008) (0.197,0.877)rectangle++(0.008,0.008) (-0.041,0.983)rectangle++(0.008,0.008) (0.193,0.892)rectangle++(0.008,0.008)
  (-0.009,0.738)rectangle++(0.008,0.008) (-0.009,0.714)rectangle++(0.008,0.008) (0.058,0.955)rectangle++(0.008,0.008) (0.107,0.863)rectangle++(0.008,0.008) (0.043,0.96)rectangle++(0.008,0.008)
  (0.12,0.952)rectangle++(0.008,0.008) (-0.004,0.667)rectangle++(0.008,0.008) (-0.105,0.901)rectangle++(0.008,0.008) (-0.19,0.835)rectangle++(0.008,0.008) (-0.082,0.948)rectangle++(0.008,0.008)
  (-0.2,0.875)rectangle++(0.008,0.008) (-0.173,0.912)rectangle++(0.008,0.008) (0.043,0.837)rectangle++(0.008,0.008) (-0.181,0.888)rectangle++(0.008,0.008) (-0.078,0.971)rectangle++(0.008,0.008)
  (0,0.753)rectangle++(0.008,0.008) (-0.002,0.914)rectangle++(0.008,0.008) (-0.053,0.983)rectangle++(0.008,0.008) (-0.117,0.926)rectangle++(0.008,0.008) (-0.1,0.861)rectangle++(0.008,0.008)
  (-0.044,0.959)rectangle++(0.008,0.008) (-0.195,0.832)rectangle++(0.008,0.008) (-0.166,0.904)rectangle++(0.008,0.008) (0.036,0.981)rectangle++(0.008,0.008) (-0.086,0.986)rectangle++(0.008,0.008)
  (0.129,0.975)rectangle++(0.008,0.008) (0.168,0.94)rectangle++(0.008,0.008) (-0.05,0.984)rectangle++(0.008,0.008) (0.137,0.894)rectangle++(0.008,0.008) (0.158,0.924)rectangle++(0.008,0.008)
  (0.05,0.845)rectangle++(0.008,0.008) (-0.006,0.947)rectangle++(0.008,0.008) (0.189,0.886)rectangle++(0.008,0.008) (0.167,0.914)rectangle++(0.008,0.008) (0.191,0.861)rectangle++(0.008,0.008)
  (0.057,0.975)rectangle++(0.008,0.008) (0.126,0.948)rectangle++(0.008,0.008) (0.164,0.936)rectangle++(0.008,0.008) (0.162,0.81)rectangle++(0.008,0.008) (0.066,0.938)rectangle++(0.008,0.008)
  (-0.17,0.897)rectangle++(0.008,0.008) (0.099,0.97)rectangle++(0.008,0.008) (-0.003,0.69)rectangle++(0.008,0.008) (0.1,0.906)rectangle++(0.008,0.008) (0.04,0.978)rectangle++(0.008,0.008)
  (-0.164,0.829)rectangle++(0.008,0.008) (-0.16,0.908)rectangle++(0.008,0.008) (0.103,0.971)rectangle++(0.008,0.008) (0.041,1.005)rectangle++(0.008,0.008) (-0.013,0.975)rectangle++(0.008,0.008)
  (-0.056,0.838)rectangle++(0.008,0.008) (-0.026,0.95)rectangle++(0.008,0.008) (0.056,0.988)rectangle++(0.008,0.008) (-0.121,0.966)rectangle++(0.008,0.008) (0.158,0.936)rectangle++(0.008,0.008)
  (-0.114,0.977)rectangle++(0.008,0.008) (-0.05,1.008)rectangle++(0.008,0.008) (-0.066,0.906)rectangle++(0.008,0.008) (-0.096,0.983)rectangle++(0.008,0.008) (-0.002,0.7)rectangle++(0.008,0.008)
  (0.102,0.907)rectangle++(0.008,0.008) (-0.197,0.869)rectangle++(0.008,0.008) (-0.068,0.905)rectangle++(0.008,0.008) (0.181,0.893)rectangle++(0.008,0.008) (-0.069,0.889)rectangle++(0.008,0.008)
  (-0.076,0.979)rectangle++(0.008,0.008) (-0.034,1.009)rectangle++(0.008,0.008) (0.06,0.854)rectangle++(0.008,0.008) (0.137,0.928)rectangle++(0.008,0.008) (-0.121,0.934)rectangle++(0.008,0.008)
  (-0.167,0.935)rectangle++(0.008,0.008) (0.036,0.986)rectangle++(0.008,0.008) (0.08,0.95)rectangle++(0.008,0.008) (0.031,0.997)rectangle++(0.008,0.008) (0.197,0.882)rectangle++(0.008,0.008)
  (-0.098,0.973)rectangle++(0.008,0.008) (-0.044,1.006)rectangle++(0.008,0.008) (-0.009,0.749)rectangle++(0.008,0.008) (-0.17,0.878)rectangle++(0.008,0.008) (0.035,0.978)rectangle++(0.008,0.008)
  (0.151,0.955)rectangle++(0.008,0.008) (0.061,0.999)rectangle++(0.008,0.008) (0.023,1.017)rectangle++(0.008,0.008) (-0.063,0.997)rectangle++(0.008,0.008) (-0.029,1.017)rectangle++(0.008,0.008)
  (-0.031,0.807)rectangle++(0.008,0.008) (-0.07,0.917)rectangle++(0.008,0.008) (0.121,0.959)rectangle++(0.008,0.008) (0.049,1)rectangle++(0.008,0.008) (0.066,0.995)rectangle++(0.008,0.008)
  (0.026,1.016)rectangle++(0.008,0.008) (0.063,0.997)rectangle++(0.008,0.008) (0.098,0.861)rectangle++(0.008,0.008) (0.111,0.865)rectangle++(0.008,0.008) (-0.003,0.749)rectangle++(0.008,0.008)
  (0.109,0.917)rectangle++(0.008,0.008) (0.02,0.804)rectangle++(0.008,0.008) (0.159,0.918)rectangle++(0.008,0.008) (-0.111,0.974)rectangle++(0.008,0.008) (0.003,0.97)rectangle++(0.008,0.008)
  (-0.149,0.96)rectangle++(0.008,0.008) (-0.183,0.922)rectangle++(0.008,0.008) (-0.079,0.985)rectangle++(0.008,0.008) (-0.167,0.951)rectangle++(0.008,0.008) (-0.039,0.823)rectangle++(0.008,0.008)
  (0.156,0.923)rectangle++(0.008,0.008) (0.166,0.926)rectangle++(0.008,0.008) (0.067,0.987)rectangle++(0.008,0.008) (0.026,1.012)rectangle++(0.008,0.008) (-0.151,0.972)rectangle++(0.008,0.008)
  (0,0.758)rectangle++(0.008,0.008) (-0.172,0.921)rectangle++(0.008,0.008) (-0.075,0.984)rectangle++(0.008,0.008) (-0.166,0.952)rectangle++(0.008,0.008) (0.108,0.98)rectangle++(0.008,0.008)
  (-0.13,0.978)rectangle++(0.008,0.008) (-0.006,0.769)rectangle++(0.008,0.008) (-0.192,0.843)rectangle++(0.008,0.008) (-0.094,0.973)rectangle++(0.008,0.008) (0.004,0.971)rectangle++(0.008,0.008)
  (-0.001,1.005)rectangle++(0.008,0.008) (0.059,0.996)rectangle++(0.008,0.008) (0.195,0.862)rectangle++(0.008,0.008) (0.135,0.918)rectangle++(0.008,0.008) (-0.145,0.913)rectangle++(0.008,0.008)
  (-0.067,0.966)rectangle++(0.008,0.008) (0.059,0.871)rectangle++(0.008,0.008) (-0.123,0.95)rectangle++(0.008,0.008) (-0.003,0.761)rectangle++(0.008,0.008) (0.044,0.834)rectangle++(0.008,0.008)
  (0.016,0.948)rectangle++(0.008,0.008) (0.004,0.996)rectangle++(0.008,0.008) (0,1.016)rectangle++(0.008,0.008) (-0.201,0.834)rectangle++(0.008,0.008) (0.006,0.954)rectangle++(0.008,0.008)
  (-0.059,0.841)rectangle++(0.008,0.008) (-0.068,0.9)rectangle++(0.008,0.008) (-0.132,0.942)rectangle++(0.008,0.008) (-0.006,0.763)rectangle++(0.008,0.008) (-0.119,0.919)rectangle++(0.008,0.008)
  (0.112,0.867)rectangle++(0.008,0.008) (0.149,0.922)rectangle++(0.008,0.008) (-0.114,0.973)rectangle++(0.008,0.008) (-0.004,0.666)rectangle++(0.008,0.008) (-0.148,0.858)rectangle++(0.008,0.008)
  (-0.169,0.814)rectangle++(0.008,0.008) (-0.073,0.94)rectangle++(0.008,0.008) (0.125,0.925)rectangle++(0.008,0.008) (-0.19,0.902)rectangle++(0.008,0.008) (-0.082,0.977)rectangle++(0.008,0.008)
  (-0.001,0.954)rectangle++(0.008,0.008) (-0.19,0.906)rectangle++(0.008,0.008) (-0.153,0.855)rectangle++(0.008,0.008) (-0.067,0.957)rectangle++(0.008,0.008) (-0.03,1)rectangle++(0.008,0.008)
  (-0.005,0.692)rectangle++(0.008,0.008) (-0.188,0.841)rectangle++(0.008,0.008) (0.169,0.907)rectangle++(0.008,0.008) (0.068,0.979)rectangle++(0.008,0.008) (0.026,1.009)rectangle++(0.008,0.008)
  (-0.001,0.932)rectangle++(0.008,0.008) (-0.002,0.727)rectangle++(0.008,0.008) (0.154,0.816)rectangle++(0.008,0.008) (-0.156,0.928)rectangle++(0.008,0.008) (0.019,0.977)rectangle++(0.008,0.008)
  (-0.016,0.975)rectangle++(0.008,0.008) (0.043,0.978)rectangle++(0.008,0.008) (-0.003,0.769)rectangle++(0.008,0.008) (0.112,0.921)rectangle++(0.008,0.008) (0.045,0.984)rectangle++(0.008,0.008)
  (-0.06,0.993)rectangle++(0.008,0.008) (-0.028,1.015)rectangle++(0.008,0.008) (-0.014,1.024)rectangle++(0.008,0.008) (-0.005,0.741)rectangle++(0.008,0.008) (-0.004,0.909)rectangle++(0.008,0.008)
  (-0.103,0.944)rectangle++(0.008,0.008) (0.119,0.969)rectangle++(0.008,0.008) (-0.083,0.931)rectangle++(0.008,0.008) (-0.198,0.846)rectangle++(0.008,0.008) (-0.117,0.864)rectangle++(0.008,0.008)
  (-0.051,0.961)rectangle++(0.008,0.008) (-0.074,0.989)rectangle++(0.008,0.008) (-0.033,1.013)rectangle++(0.008,0.008) (0.097,0.982)rectangle++(0.008,0.008) (-0.057,0.925)rectangle++(0.008,0.008)
  (-0.196,0.878)rectangle++(0.008,0.008) (-0.094,0.978)rectangle++(0.008,0.008) (0.126,0.974)rectangle++(0.008,0.008) (0.051,1.007)rectangle++(0.008,0.008) (0.019,1.021)rectangle++(0.008,0.008)
  (0.099,0.865)rectangle++(0.008,0.008) (-0.016,0.963)rectangle++(0.008,0.008) (-0.01,0.717)rectangle++(0.008,0.008) (0.103,0.912)rectangle++(0.008,0.008) (-0.106,0.977)rectangle++(0.008,0.008)
  (-0.047,1.008)rectangle++(0.008,0.008) (-0.164,0.96)rectangle++(0.008,0.008) (-0.184,0.922)rectangle++(0.008,0.008) (-0.079,0.985)rectangle++(0.008,0.008) (-0.036,1.011)rectangle++(0.008,0.008)
  (0.142,0.972)rectangle++(0.008,0.008) (-0.203,0.846)rectangle++(0.008,0.008) (0.147,0.932)rectangle++(0.008,0.008) (0.059,0.989)rectangle++(0.008,0.008) (0.023,1.013)rectangle++(0.008,0.008)
  (0.17,0.932)rectangle++(0.008,0.008) (-0.119,0.867)rectangle++(0.008,0.008) (-0.009,0.733)rectangle++(0.008,0.008) (0.011,0.965)rectangle++(0.008,0.008) (-0.004,0.767)rectangle++(0.008,0.008)
  (-0.004,0.92)rectangle++(0.008,0.008) (-0.004,0.984)rectangle++(0.008,0.008) (-0.008,0.75)rectangle++(0.008,0.008) (-0.118,0.916)rectangle++(0.008,0.008) (-0.006,0.701)rectangle++(0.008,0.008)
  (-0.006,0.752)rectangle++(0.008,0.008) (-0.005,0.914)rectangle++(0.008,0.008) (-0.004,0.758)rectangle++(0.008,0.008) (-0.004,0.677)rectangle++(0.008,0.008) (0.059,0.949)rectangle++(0.008,0.008)
  (0.022,0.996)rectangle++(0.008,0.008) (0.007,1.016)rectangle++(0.008,0.008) (-0.065,0.849)rectangle++(0.008,0.008) (-0.03,0.954)rectangle++(0.008,0.008) (0.106,0.976)rectangle++(0.008,0.008)
  (-0.003,0.769)rectangle++(0.008,0.008) (0.156,0.828)rectangle++(0.008,0.008) (-0.144,0.891)rectangle++(0.008,0.008) (0.103,0.965)rectangle++(0.008,0.008) (-0.164,0.929)rectangle++(0.008,0.008)
  (-0.064,0.847)rectangle++(0.008,0.008) (-0.029,0.953)rectangle++(0.008,0.008) (-0.015,0.998)rectangle++(0.008,0.008) (-0.008,1.017)rectangle++(0.008,0.008) (0.058,0.997)rectangle++(0.008,0.008)
  (0.126,0.943)rectangle++(0.008,0.008) (-0.087,0.938)rectangle++(0.008,0.008) (0.099,0.978)rectangle++(0.008,0.008) (0.089,0.86)rectangle++(0.008,0.008) (0.021,0.973)rectangle++(0.008,0.008)
  (0.007,1.006)rectangle++(0.008,0.008) (0.06,0.996)rectangle++(0.008,0.008) (-0.142,0.974)rectangle++(0.008,0.008) (-0.006,0.769)rectangle++(0.008,0.008) (-0.173,0.879)rectangle++(0.008,0.008)
  (-0.188,0.897)rectangle++(0.008,0.008) (-0.081,0.974)rectangle++(0.008,0.008) (0.128,0.972)rectangle++(0.008,0.008) (-0.048,0.991)rectangle++(0.008,0.008) (-0.165,0.899)rectangle++(0.008,0.008)
  (0.042,0.833)rectangle++(0.008,0.008) (-0.059,0.845)rectangle++(0.008,0.008) (-0.009,0.739)rectangle++(0.008,0.008) (0.164,0.92)rectangle++(0.008,0.008) (0.067,0.984)rectangle++(0.008,0.008)
  (0.051,0.846)rectangle++(0.008,0.008) (0.12,0.947)rectangle++(0.008,0.008) (-0.002,0.727)rectangle++(0.008,0.008) (0.159,0.878)rectangle++(0.008,0.008) (-0.044,0.94)rectangle++(0.008,0.008)
  (0.186,0.84)rectangle++(0.008,0.008) (-0.006,0.929)rectangle++(0.008,0.008) (-0.02,0.973)rectangle++(0.008,0.008) (0.139,0.962)rectangle++(0.008,0.008) (0.056,1.002)rectangle++(0.008,0.008)
  (0.002,0.715)rectangle++(0.008,0.008) (0.096,0.958)rectangle++(0.008,0.008) (0.073,0.989)rectangle++(0.008,0.008) (-0.138,0.975)rectangle++(0.008,0.008) (-0.06,1.007)rectangle++(0.008,0.008)
  (-0.203,0.862)rectangle++(0.008,0.008) (-0.088,0.96)rectangle++(0.008,0.008) (0.056,0.85)rectangle++(0.008,0.008) (-0.088,0.957)rectangle++(0.008,0.008) (-0.039,0.999)rectangle++(0.008,0.008)
  (-0.019,1.017)rectangle++(0.008,0.008) (0.057,0.997)rectangle++(0.008,0.008) (0.022,1.017)rectangle++(0.008,0.008) (0.007,1.025)rectangle++(0.008,0.008) (-0.202,0.834)rectangle++(0.008,0.008)
  (-0.189,0.855)rectangle++(0.008,0.008) (0.059,0.969)rectangle++(0.008,0.008) (0.154,0.95)rectangle++(0.008,0.008) (-0.117,0.979)rectangle++(0.008,0.008) (0.057,0.852)rectangle++(0.008,0.008)
  (0.059,0.904)rectangle++(0.008,0.008) (-0.17,0.923)rectangle++(0.008,0.008) (-0.031,0.981)rectangle++(0.008,0.008) (-0.006,0.762)rectangle++(0.008,0.008) (-0.119,0.919)rectangle++(0.008,0.008)
  (0.043,0.983)rectangle++(0.008,0.008) (0.016,1.011)rectangle++(0.008,0.008) (0.004,1.022)rectangle++(0.008,0.008) (0.2,0.887)rectangle++(0.008,0.008) (0.082,0.97)rectangle++(0.008,0.008)
  (-0.134,0.972)rectangle++(0.008,0.008) (-0.058,1.006)rectangle++(0.008,0.008) (0,0.755)rectangle++(0.008,0.008) (-0.031,0.943)rectangle++(0.008,0.008) (0.133,0.958)rectangle++(0.008,0.008)
  (0.053,1)rectangle++(0.008,0.008) (-0.154,0.816)rectangle++(0.008,0.008) (-0.067,0.94)rectangle++(0.008,0.008) (-0.03,0.993)rectangle++(0.008,0.008) (0.181,0.876)rectangle++(0.008,0.008)
  (-0.066,0.905)rectangle++(0.008,0.008) (-0.169,0.811)rectangle++(0.008,0.008) (0.007,0.957)rectangle++(0.008,0.008) (0,0.747)rectangle++(0.008,0.008) (-0.002,0.911)rectangle++(0.008,0.008)
  (-0.003,0.98)rectangle++(0.008,0.008) (-0.004,1.009)rectangle++(0.008,0.008) (0.031,0.823)rectangle++(0.008,0.008) (0.011,0.943)rectangle++(0.008,0.008) (0.002,0.994)rectangle++(0.008,0.008)
  (-0.001,1.015)rectangle++(0.008,0.008) (-0.003,1.024)rectangle++(0.008,0.008) (-0.004,1.028)rectangle++(0.008,0.008) (0.042,0.983)rectangle++(0.008,0.008) (0.016,1.011)rectangle++(0.008,0.008)
  (0.004,1.022)rectangle++(0.008,0.008) (0.054,0.842)rectangle++(0.008,0.008) (0.134,0.926)rectangle++(0.008,0.008) (0.058,0.906)rectangle++(0.008,0.008) (0.191,0.861)rectangle++(0.008,0.008)
  (-0.096,0.968)rectangle++(0.008,0.008) (-0.102,0.982)rectangle++(0.008,0.008) (0.125,0.976)rectangle++(0.008,0.008) (0.032,0.813)rectangle++(0.008,0.008) (-0.036,0.98)rectangle++(0.008,0.008)
  (0.113,0.925)rectangle++(0.008,0.008) (-0.028,0.944)rectangle++(0.008,0.008) (-0.014,0.994)rectangle++(0.008,0.008) (-0.068,0.994)rectangle++(0.008,0.008) (0.167,0.815)rectangle++(0.008,0.008)
  (-0.097,0.974)rectangle++(0.008,0.008) (-0.043,1.007)rectangle++(0.008,0.008) (-0.178,0.931)rectangle++(0.008,0.008) (-0.189,0.823)rectangle++(0.008,0.008) (0.033,0.97)rectangle++(0.008,0.008)
  (0.058,0.851)rectangle++(0.008,0.008) (-0.193,0.842)rectangle++(0.008,0.008) (-0.062,0.846)rectangle++(0.008,0.008) (0.176,0.907)rectangle++(0.008,0.008) (0.025,0.98)rectangle++(0.008,0.008)
  (-0.004,0.77)rectangle++(0.008,0.008) (-0.004,0.732)rectangle++(0.008,0.008) (-0.114,0.914)rectangle++(0.008,0.008) (-0.197,0.869)rectangle++(0.008,0.008) (0.006,0.785)rectangle++(0.008,0.008)
  (0.116,0.922)rectangle++(0.008,0.008) (0.158,0.932)rectangle++(0.008,0.008) (0.153,0.898)rectangle++(0.008,0.008) (0.024,0.979)rectangle++(0.008,0.008) (-0.108,0.861)rectangle++(0.008,0.008)
  (-0.048,0.959)rectangle++(0.008,0.008) (0.19,0.892)rectangle++(0.008,0.008) (0.168,0.915)rectangle++(0.008,0.008) (-0.187,0.908)rectangle++(0.008,0.008) (-0.176,0.92)rectangle++(0.008,0.008)
  (-0.176,0.924)rectangle++(0.008,0.008) (-0.188,0.824)rectangle++(0.008,0.008) (0.166,0.907)rectangle++(0.008,0.008) (0.079,0.937)rectangle++(0.008,0.008) (-0.129,0.965)rectangle++(0.008,0.008)
  (-0.182,0.924)rectangle++(0.008,0.008) (-0.084,0.933)rectangle++(0.008,0.008) (-0.16,0.94)rectangle++(0.008,0.008) (-0.176,0.93)rectangle++(0.008,0.008) (-0.17,0.81)rectangle++(0.008,0.008)
  (0.046,0.839)rectangle++(0.008,0.008) (0.017,0.95)rectangle++(0.008,0.008) (-0.045,0.985)rectangle++(0.008,0.008) (-0.16,0.956)rectangle++(0.008,0.008) (0.11,0.98)rectangle++(0.008,0.008)
  (0.059,0.851)rectangle++(0.008,0.008) (0.171,0.881)rectangle++(0.008,0.008) (-0.127,0.958)rectangle++(0.008,0.008) (-0.056,1)rectangle++(0.008,0.008) (-0.059,0.846)rectangle++(0.008,0.008)
  (-0.027,0.953)rectangle++(0.008,0.008) (0.056,0.85)rectangle++(0.008,0.008) (0.171,0.82)rectangle++(0.008,0.008) (0.153,0.904)rectangle++(0.008,0.008) (-0.11,0.97)rectangle++(0.008,0.008)
  (-0.204,0.872)rectangle++(0.008,0.008) (-0.004,0.636)rectangle++(0.008,0.008) (-0.095,0.979)rectangle++(0.008,0.008) (-0.042,1.009)rectangle++(0.008,0.008) (0.037,0.94)rectangle++(0.008,0.008)
  (0.051,0.838)rectangle++(0.008,0.008) (0.105,0.964)rectangle++(0.008,0.008) (0.06,0.9)rectangle++(0.008,0.008) (-0.108,0.974)rectangle++(0.008,0.008) (0.035,0.828)rectangle++(0.008,0.008)
  (0.178,0.837)rectangle++(0.008,0.008) (0.157,0.907)rectangle++(0.008,0.008) (-0.154,0.813)rectangle++(0.008,0.008) (-0.184,0.86)rectangle++(0.008,0.008) (-0.068,0.904)rectangle++(0.008,0.008)
  (0.187,0.842)rectangle++(0.008,0.008) (0.159,0.906)rectangle++(0.008,0.008) (0.065,0.978)rectangle++(0.008,0.008) (0.001,0.745)rectangle++(0.008,0.008) (-0.002,0.911)rectangle++(0.008,0.008)
  (-0.067,0.982)rectangle++(0.008,0.008) (0.132,0.971)rectangle++(0.008,0.008) (-0.067,0.86)rectangle++(0.008,0.008) (0.113,0.948)rectangle++(0.008,0.008) (-0.124,0.973)rectangle++(0.008,0.008)
  (-0.054,1.006)rectangle++(0.008,0.008) (0.066,0.978)rectangle++(0.008,0.008) (0.025,1.009)rectangle++(0.008,0.008) (0.119,0.927)rectangle++(0.008,0.008) (0.035,0.825)rectangle++(0.008,0.008)
  (0.174,0.876)rectangle++(0.008,0.008) (0.186,0.859)rectangle++(0.008,0.008) (-0.179,0.909)rectangle++(0.008,0.008) (0.151,0.935)rectangle++(0.008,0.008) (-0.19,0.906)rectangle++(0.008,0.008)
  (0.016,0.969)rectangle++(0.008,0.008) (-0.146,0.962)rectangle++(0.008,0.008) (0.112,0.867)rectangle++(0.008,0.008) (-0.004,0.744)rectangle++(0.008,0.008) (0.108,0.916)rectangle++(0.008,0.008)
  (-0.003,0.758)rectangle++(0.008,0.008) (-0.191,0.862)rectangle++(0.008,0.008) (-0.123,0.959)rectangle++(0.008,0.008) (-0.054,1)rectangle++(0.008,0.008) (-0.101,0.861)rectangle++(0.008,0.008)
  (-0.152,0.924)rectangle++(0.008,0.008) (-0.002,0.701)rectangle++(0.008,0.008) (-0.109,0.908)rectangle++(0.008,0.008) (-0.148,0.856)rectangle++(0.008,0.008) (0.038,0.957)rectangle++(0.008,0.008)
  (-0.162,0.809)rectangle++(0.008,0.008) (-0.07,0.938)rectangle++(0.008,0.008) (0.007,0.97)rectangle++(0.008,0.008) (-0.148,0.961)rectangle++(0.008,0.008) (-0.183,0.922)rectangle++(0.008,0.008)
  (-0.079,0.985)rectangle++(0.008,0.008) (0.115,0.926)rectangle++(0.008,0.008) (-0.07,0.978)rectangle++(0.008,0.008) (0.001,0.925)rectangle++(0.008,0.008) (-0.002,0.986)rectangle++(0.008,0.008)
  (-0.067,0.993)rectangle++(0.008,0.008) (-0.164,0.935)rectangle++(0.008,0.008) (0.137,0.915)rectangle++(0.008,0.008) (0.055,0.982)rectangle++(0.008,0.008) (0,0.963)rectangle++(0.008,0.008)
  (-0.148,0.958)rectangle++(0.008,0.008) (-0.087,0.989)rectangle++(0.008,0.008) (-0.039,1.013)rectangle++(0.008,0.008) (-0.019,1.023)rectangle++(0.008,0.008) (-0.121,0.934)rectangle++(0.008,0.008)
  (-0.058,0.845)rectangle++(0.008,0.008) (-0.027,0.953)rectangle++(0.008,0.008) (0.135,0.959)rectangle++(0.008,0.008) (-0.194,0.905)rectangle++(0.008,0.008) (-0.094,0.981)rectangle++(0.008,0.008)
  (-0.195,0.828)rectangle++(0.008,0.008) (0.148,0.813)rectangle++(0.008,0.008) (-0.008,0.931)rectangle++(0.008,0.008) (0.108,0.864)rectangle++(0.008,0.008) (-0.131,0.937)rectangle++(0.008,0.008)
  (-0.057,0.991)rectangle++(0.008,0.008) (-0.026,1.014)rectangle++(0.008,0.008) (-0.006,0.764)rectangle++(0.008,0.008) (-0.005,0.919)rectangle++(0.008,0.008) (-0.004,0.984)rectangle++(0.008,0.008)
  (-0.004,1.011)rectangle++(0.008,0.008) (-0.004,1.022)rectangle++(0.008,0.008) (-0.004,1.027)rectangle++(0.008,0.008) (-0.159,0.97)rectangle++(0.008,0.008) (0.093,0.983)rectangle++(0.008,0.008)
  (-0.054,0.842)rectangle++(0.008,0.008) (-0.173,0.923)rectangle++(0.008,0.008) (0.103,0.976)rectangle++(0.008,0.008) (-0.13,0.976)rectangle++(0.008,0.008) (0.077,0.857)rectangle++(0.008,0.008)
  (-0.118,0.95)rectangle++(0.008,0.008) (-0.169,0.938)rectangle++(0.008,0.008) (-0.177,0.929)rectangle++(0.008,0.008) (-0.177,0.925)rectangle++(0.008,0.008) (0.102,0.977)rectangle++(0.008,0.008)
  (0.041,1.008)rectangle++(0.008,0.008) (-0.147,0.973)rectangle++(0.008,0.008) (-0.087,0.991)rectangle++(0.008,0.008) (0.193,0.898)rectangle++(0.008,0.008) (-0.039,0.981)rectangle++(0.008,0.008)
  (0.137,0.967)rectangle++(0.008,0.008) (0.177,0.911)rectangle++(0.008,0.008) (-0.119,0.865)rectangle++(0.008,0.008) (-0.179,0.818)rectangle++(0.008,0.008) (0.034,0.969)rectangle++(0.008,0.008)
  (0.058,0.851)rectangle++(0.008,0.008) (-0.005,0.697)rectangle++(0.008,0.008) (0.155,0.877)rectangle++(0.008,0.008) (-0.106,0.965)rectangle++(0.008,0.008) (-0.047,1.003)rectangle++(0.008,0.008)
  (-0.005,0.774)rectangle++(0.008,0.008) (-0.121,0.922)rectangle++(0.008,0.008) (-0.165,0.933)rectangle++(0.008,0.008) (0.106,0.976)rectangle++(0.008,0.008) (-0.002,0.7)rectangle++(0.008,0.008)
  (-0.109,0.908)rectangle++(0.008,0.008) (-0.161,0.932)rectangle++(0.008,0.008) (-0.07,0.989)rectangle++(0.008,0.008) (-0.032,1.013)rectangle++(0.008,0.008) (-0.016,1.023)rectangle++(0.008,0.008)
  (-0.12,0.978)rectangle++(0.008,0.008) (-0.145,0.909)rectangle++(0.008,0.008) (-0.168,0.927)rectangle++(0.008,0.008) (-0.049,0.824)rectangle++(0.008,0.008) (-0.023,0.944)rectangle++(0.008,0.008)
  (0.137,0.94)rectangle++(0.008,0.008) (0.079,0.986)rectangle++(0.008,0.008) (0.08,0.947)rectangle++(0.008,0.008) (-0.13,0.937)rectangle++(0.008,0.008) (-0.064,0.848)rectangle++(0.008,0.008)
  (0.093,0.979)rectangle++(0.008,0.008) (-0.084,0.986)rectangle++(0.008,0.008) (0.06,0.906)rectangle++(0.008,0.008) (0.138,0.913)rectangle++(0.008,0.008) (0.055,0.981)rectangle++(0.008,0.008)
  (0.067,0.992)rectangle++(0.008,0.008) (-0.14,0.974)rectangle++(0.008,0.008) (-0.177,0.94)rectangle++(0.008,0.008) (0.104,0.979)rectangle++(0.008,0.008) (0.12,0.957)rectangle++(0.008,0.008)
  (0.097,0.943)rectangle++(0.008,0.008) (-0.199,0.869)rectangle++(0.008,0.008) (0.09,0.969)rectangle++(0.008,0.008) (-0.062,0.843)rectangle++(0.008,0.008) (0.11,0.867)rectangle++(0.008,0.008)
  (-0.016,0.963)rectangle++(0.008,0.008) (0.078,0.935)rectangle++(0.008,0.008) (-0.129,0.965)rectangle++(0.008,0.008) (0.117,0.977)rectangle++(0.008,0.008) (0.153,0.813)rectangle++(0.008,0.008)
  (0.001,0.707)rectangle++(0.008,0.008) (0.174,0.881)rectangle++(0.008,0.008) (-0.002,0.96)rectangle++(0.008,0.008) (-0.083,0.93)rectangle++(0.008,0.008) (-0.037,0.988)rectangle++(0.008,0.008)
  (-0.106,0.944)rectangle++(0.008,0.008) (-0.166,0.939)rectangle++(0.008,0.008) (0.183,0.907)rectangle++(0.008,0.008) (-0.139,0.929)rectangle++(0.008,0.008) (-0.171,0.811)rectangle++(0.008,0.008)
  (0.12,0.929)rectangle++(0.008,0.008) (-0.12,0.97)rectangle++(0.008,0.008) (-0.005,0.768)rectangle++(0.008,0.008) (-0.001,0.732)rectangle++(0.008,0.008) (0.107,0.913)rectangle++(0.008,0.008)
  (-0.186,0.82)rectangle++(0.008,0.008) (-0.08,0.942)rectangle++(0.008,0.008) (0.102,0.977)rectangle++(0.008,0.008) (-0.145,0.894)rectangle++(0.008,0.008) (-0.031,0.977)rectangle++(0.008,0.008)
  (-0.071,0.992)rectangle++(0.008,0.008) (0.058,0.853)rectangle++(0.008,0.008) (0.179,0.874)rectangle++(0.008,0.008) (0.038,0.821)rectangle++(0.008,0.008) (0,0.734)rectangle++(0.008,0.008)
  (-0.05,0.832)rectangle++(0.008,0.008) (-0.138,0.926)rectangle++(0.008,0.008) (-0.164,0.828)rectangle++(0.008,0.008) (-0.071,0.946)rectangle++(0.008,0.008) (0.187,0.895)rectangle++(0.008,0.008)
  (0.076,0.973)rectangle++(0.008,0.008) (0.03,1.007)rectangle++(0.008,0.008) (-0.068,0.91)rectangle++(0.008,0.008) (-0.062,0.846)rectangle++(0.008,0.008) (-0.185,0.874)rectangle++(0.008,0.008)
  (-0.006,0.96)rectangle++(0.008,0.008) (0.138,0.914)rectangle++(0.008,0.008) (0.161,0.927)rectangle++(0.008,0.008) (0.081,0.952)rectangle++(0.008,0.008) (-0.064,0.85)rectangle++(0.008,0.008)
  (-0.11,0.969)rectangle++(0.008,0.008) (-0.17,0.943)rectangle++(0.008,0.008) (-0.178,0.929)rectangle++(0.008,0.008) (-0.177,0.925)rectangle++(0.008,0.008) (-0.178,0.882)rectangle++(0.008,0.008)
  (0.072,0.982)rectangle++(0.008,0.008) (-0.169,0.935)rectangle++(0.008,0.008) (0.036,0.986)rectangle++(0.008,0.008) (-0.191,0.908)rectangle++(0.008,0.008) (0.097,0.975)rectangle++(0.008,0.008)
  (0.058,0.97)rectangle++(0.008,0.008) (-0.17,0.933)rectangle++(0.008,0.008) (-0.074,0.99)rectangle++(0.008,0.008) (-0.009,0.746)rectangle++(0.008,0.008) (-0.068,0.959)rectangle++(0.008,0.008)
  (-0.116,0.864)rectangle++(0.008,0.008) (0.034,0.819)rectangle++(0.008,0.008) (0.156,0.828)rectangle++(0.008,0.008) (-0.188,0.894)rectangle++(0.008,0.008) (0.033,0.98)rectangle++(0.008,0.008)
  (0.064,0.992)rectangle++(0.008,0.008) (0.043,0.933)rectangle++(0.008,0.008) (0.016,0.99)rectangle++(0.008,0.008) (0.01,0.978)rectangle++(0.008,0.008) (-0.065,0.992)rectangle++(0.008,0.008)
  (-0.03,1.014)rectangle++(0.008,0.008) (-0.068,0.854)rectangle++(0.008,0.008) (0.097,0.864)rectangle++(0.008,0.008) (-0.003,0.749)rectangle++(0.008,0.008) (0.164,0.92)rectangle++(0.008,0.008)
  (0.046,0.983)rectangle++(0.008,0.008) (0.154,0.954)rectangle++(0.008,0.008) (0.113,0.866)rectangle++(0.008,0.008) (-0.112,0.957)rectangle++(0.008,0.008) (0.187,0.901)rectangle++(0.008,0.008)
  (-0.103,0.975)rectangle++(0.008,0.008) (-0.008,0.753)rectangle++(0.008,0.008) (-0.118,0.917)rectangle++(0.008,0.008) (-0.181,0.856)rectangle++(0.008,0.008) (-0.006,0.748)rectangle++(0.008,0.008)
  (-0.005,0.912)rectangle++(0.008,0.008) (-0.067,0.982)rectangle++(0.008,0.008) (-0.076,0.992)rectangle++(0.008,0.008) (-0.135,0.948)rectangle++(0.008,0.008) (0.061,0.978)rectangle++(0.008,0.008)
  (-0.065,0.852)rectangle++(0.008,0.008) (-0.004,0.663)rectangle++(0.008,0.008) (-0.004,0.876)rectangle++(0.008,0.008) (0.015,0.969)rectangle++(0.008,0.008) (-0.144,0.86)rectangle++(0.008,0.008)
  (-0.063,0.959)rectangle++(0.008,0.008) (0.129,0.966)rectangle++(0.008,0.008) (0.052,1.003)rectangle++(0.008,0.008) (-0.144,0.974)rectangle++(0.008,0.008) (-0.063,1.007)rectangle++(0.008,0.008)
  (-0.029,1.021)rectangle++(0.008,0.008) (-0.162,0.965)rectangle++(0.008,0.008) (0.015,0.97)rectangle++(0.008,0.008) (0.059,0.911)rectangle++(0.008,0.008) (0.022,0.98)rectangle++(0.008,0.008)
  (-0.016,0.975)rectangle++(0.008,0.008) (-0.009,1.007)rectangle++(0.008,0.008) (-0.007,0.758)rectangle++(0.008,0.008) (-0.004,0.653)rectangle++(0.008,0.008) (-0.103,0.898)rectangle++(0.008,0.008)
  (0.135,0.861)rectangle++(0.008,0.008) (-0.181,0.902)rectangle++(0.008,0.008) (-0.006,0.756)rectangle++(0.008,0.008) (0.057,0.972)rectangle++(0.008,0.008) (0.021,1.006)rectangle++(0.008,0.008)
  (0.007,1.02)rectangle++(0.008,0.008) (0.2,0.887)rectangle++(0.008,0.008) (0.082,0.97)rectangle++(0.008,0.008) (0.032,1.005)rectangle++(0.008,0.008) (-0.004,0.774)rectangle++(0.008,0.008)
  (-0.004,0.733)rectangle++(0.008,0.008) (-0.004,0.906)rectangle++(0.008,0.008) (0.132,0.947)rectangle++(0.008,0.008) (0.12,0.957)rectangle++(0.008,0.008) (0.048,1)rectangle++(0.008,0.008)
  (-0.145,0.973)rectangle++(0.008,0.008) (-0.206,0.867)rectangle++(0.008,0.008) (-0.173,0.909)rectangle++(0.008,0.008) (-0.007,0.758)rectangle++(0.008,0.008) (-0.172,0.879)rectangle++(0.008,0.008)
  (-0.074,0.967)rectangle++(0.008,0.008) (-0.008,0.959)rectangle++(0.008,0.008) (-0.006,1.001)rectangle++(0.008,0.008) (0.192,0.834)rectangle++(0.008,0.008) (0.078,0.948)rectangle++(0.008,0.008)
  (0.031,0.996)rectangle++(0.008,0.008) (-0.014,0.975)rectangle++(0.008,0.008) (-0.202,0.885)rectangle++(0.008,0.008) (-0.087,0.969)rectangle++(0.008,0.008) (0.052,0.844)rectangle++(0.008,0.008)
  (-0.12,0.943)rectangle++(0.008,0.008) (-0.168,0.937)rectangle++(0.008,0.008) (-0.073,0.991)rectangle++(0.008,0.008) (-0.167,0.953)rectangle++(0.008,0.008) (-0.145,0.892)rectangle++(0.008,0.008)
  (-0.063,0.972)rectangle++(0.008,0.008) (0.065,0.92)rectangle++(0.008,0.008) (-0.004,0.752)rectangle++(0.008,0.008) (-0.148,0.856)rectangle++(0.008,0.008) (-0.065,0.957)rectangle++(0.008,0.008)
  (-0.029,1)rectangle++(0.008,0.008) (0.016,0.968)rectangle++(0.008,0.008) (-0.004,0.767)rectangle++(0.008,0.008) (-0.055,0.835)rectangle++(0.008,0.008) (-0.14,0.926)rectangle++(0.008,0.008)
  (-0.061,0.986)rectangle++(0.008,0.008) (-0.028,1.012)rectangle++(0.008,0.008) (-0.065,0.906)rectangle++(0.008,0.008) (-0.062,0.845)rectangle++(0.008,0.008) (-0.171,0.907)rectangle++(0.008,0.008)
  (0.101,0.972)rectangle++(0.008,0.008) (0.163,0.944)rectangle++(0.008,0.008) (-0.002,0.763)rectangle++(0.008,0.008) (-0.067,0.961)rectangle++(0.008,0.008) (0.056,0.846)rectangle++(0.008,0.008)
  (0.134,0.888)rectangle++(0.008,0.008) (-0.002,0.753)rectangle++(0.008,0.008) (-0.117,0.918)rectangle++(0.008,0.008) (-0.052,0.983)rectangle++(0.008,0.008) (-0.024,1.011)rectangle++(0.008,0.008)
  (-0.068,0.854)rectangle++(0.008,0.008) (0.112,0.947)rectangle++(0.008,0.008) (0.045,0.995)rectangle++(0.008,0.008) (0.001,0.748)rectangle++(0.008,0.008) (-0.11,0.864)rectangle++(0.008,0.008)
  (-0.049,0.96)rectangle++(0.008,0.008) (-0.023,1.001)rectangle++(0.008,0.008) (-0.158,0.963)rectangle++(0.008,0.008) (0.037,0.99)rectangle++(0.008,0.008) (-0.002,0.932)rectangle++(0.008,0.008)
  (-0.088,0.984)rectangle++(0.008,0.008) (-0.039,1.011)rectangle++(0.008,0.008) (0.106,0.97)rectangle++(0.008,0.008) (0.042,1.005)rectangle++(0.008,0.008) (-0.008,0.754)rectangle++(0.008,0.008)
  (-0.006,0.915)rectangle++(0.008,0.008) (-0.003,0.946)rectangle++(0.008,0.008) (-0.175,0.879)rectangle++(0.008,0.008) (-0.006,0.752)rectangle++(0.008,0.008) (-0.009,0.715)rectangle++(0.008,0.008)
  (-0.006,0.898)rectangle++(0.008,0.008) (0.182,0.835)rectangle++(0.008,0.008) (0,0.727)rectangle++(0.008,0.008) (0.159,0.877)rectangle++(0.008,0.008) (-0.16,0.899)rectangle++(0.008,0.008)
  (-0.123,0.963)rectangle++(0.008,0.008) (-0.127,0.963)rectangle++(0.008,0.008) (-0.204,0.869)rectangle++(0.008,0.008) (-0.199,0.854)rectangle++(0.008,0.008) (-0.151,0.855)rectangle++(0.008,0.008)
  (-0.054,0.832)rectangle++(0.008,0.008) (0.189,0.85)rectangle++(0.008,0.008) (0.161,0.907)rectangle++(0.008,0.008) (0.065,0.979)rectangle++(0.008,0.008) (0.157,0.951)rectangle++(0.008,0.008)
  (0.063,0.997)rectangle++(0.008,0.008) (-0.173,0.818)rectangle++(0.008,0.008) (-0.017,0.94)rectangle++(0.008,0.008) (-0.01,0.993)rectangle++(0.008,0.008) (-0.117,0.867)rectangle++(0.008,0.008)
  (0.087,0.86)rectangle++(0.008,0.008) (0.034,0.959)rectangle++(0.008,0.008) (0.064,0.989)rectangle++(0.008,0.008) (0.158,0.953)rectangle++(0.008,0.008) (0.064,0.998)rectangle++(0.008,0.008)
  (-0.051,0.838)rectangle++(0.008,0.008) (-0.028,0.969)rectangle++(0.008,0.008) (-0.06,0.84)rectangle++(0.008,0.008) (-0.028,0.95)rectangle++(0.008,0.008) (0.043,0.827)rectangle++(0.008,0.008)
  (0.016,0.945)rectangle++(0.008,0.008) (-0.001,0.754)rectangle++(0.008,0.008) (-0.003,0.914)rectangle++(0.008,0.008) (-0.132,0.943)rectangle++(0.008,0.008) (-0.058,0.994)rectangle++(0.008,0.008)
  (-0.027,1.015)rectangle++(0.008,0.008) (-0.014,1.024)rectangle++(0.008,0.008) (0.148,0.971)rectangle++(0.008,0.008) (-0.195,0.907)rectangle++(0.008,0.008) (0.001,0.962)rectangle++(0.008,0.008)
  (-0.018,0.974)rectangle++(0.008,0.008) (-0.004,0.769)rectangle++(0.008,0.008) (-0.12,0.921)rectangle++(0.008,0.008) (0.124,0.945)rectangle++(0.008,0.008) (0.05,0.995)rectangle++(0.008,0.008)
  (0.019,1.015)rectangle++(0.008,0.008) (-0.064,0.997)rectangle++(0.008,0.008) (-0.029,1.016)rectangle++(0.008,0.008) (0.166,0.935)rectangle++(0.008,0.008) (-0.002,0.762)rectangle++(0.008,0.008)
  (-0.133,0.931)rectangle++(0.008,0.008) (-0.001,0.752)rectangle++(0.008,0.008) (-0.041,0.818)rectangle++(0.008,0.008) (-0.083,0.967)rectangle++(0.008,0.008) (-0.037,1.004)rectangle++(0.008,0.008)
  (0.054,0.995)rectangle++(0.008,0.008) (-0.143,0.973)rectangle++(0.008,0.008) (-0.177,0.939)rectangle++(0.008,0.008) (-0.077,0.992)rectangle++(0.008,0.008) (0.055,0.845)rectangle++(0.008,0.008)
  (-0.183,0.882)rectangle++(0.008,0.008) (-0.009,0.73)rectangle++(0.008,0.008) (-0.006,0.748)rectangle++(0.008,0.008) (-0.006,0.749)rectangle++(0.008,0.008) (-0.005,0.912)rectangle++(0.008,0.008)
  (0.182,0.898)rectangle++(0.008,0.008) (0.074,0.975)rectangle++(0.008,0.008) (0.021,0.804)rectangle++(0.008,0.008) (0.006,0.935)rectangle++(0.008,0.008) (0.108,0.867)rectangle++(0.008,0.008)
  (0.182,0.867)rectangle++(0.008,0.008) (-0.099,0.968)rectangle++(0.008,0.008) (-0.044,1.004)rectangle++(0.008,0.008) (-0.207,0.882)rectangle++(0.008,0.008) (0.09,0.973)rectangle++(0.008,0.008)
  (0,0.942)rectangle++(0.008,0.008) (0.019,0.944)rectangle++(0.008,0.008) (-0.196,0.889)rectangle++(0.008,0.008) (-0.085,0.971)rectangle++(0.008,0.008) (-0.038,1.005)rectangle++(0.008,0.008)
  (-0.004,0.774)rectangle++(0.008,0.008) (-0.004,0.923)rectangle++(0.008,0.008) (-0.004,0.985)rectangle++(0.008,0.008) (-0.004,1.011)rectangle++(0.008,0.008) (-0.004,1.023)rectangle++(0.008,0.008)
  (0.15,0.969)rectangle++(0.008,0.008) (0.061,1.005)rectangle++(0.008,0.008) (0.023,1.02)rectangle++(0.008,0.008) (-0.058,0.924)rectangle++(0.008,0.008) (-0.027,0.986)rectangle++(0.008,0.008)
  (0.056,0.993)rectangle++(0.008,0.008) (-0.202,0.895)rectangle++(0.008,0.008) (-0.197,0.857)rectangle++(0.008,0.008) (-0.123,0.958)rectangle++(0.008,0.008) (0.117,0.975)rectangle++(0.008,0.008)
  (0.047,1.007)rectangle++(0.008,0.008) (0.199,0.881)rectangle++(0.008,0.008) (-0.098,0.973)rectangle++(0.008,0.008) (-0.168,0.946)rectangle++(0.008,0.008) (0.096,0.982)rectangle++(0.008,0.008)
  (0.163,0.947)rectangle++(0.008,0.008) (-0.043,0.987)rectangle++(0.008,0.008) (0.192,0.86)rectangle++(0.008,0.008) (-0.007,0.752)rectangle++(0.008,0.008) (-0.005,0.914)rectangle++(0.008,0.008)
  (-0.004,0.981)rectangle++(0.008,0.008) (-0.01,0.718)rectangle++(0.008,0.008) (0.042,0.834)rectangle++(0.008,0.008) (0.015,0.948)rectangle++(0.008,0.008) (0.142,0.952)rectangle++(0.008,0.008)
  (0.168,0.934)rectangle++(0.008,0.008) (0.083,0.986)rectangle++(0.008,0.008) (-0.007,0.954)rectangle++(0.008,0.008);

            \node[etiq, firebrick3, right] at (0.03,0.2) {$w_1(A)$ : le tronc};
            \node[etiq, steelblue4, left] at (-0.72,0.9) {$w_2(A)$};
            \node[etiq, coveraccent!80!black, right] at (0.72,0.9) {$w_3(A)$};
            \node[etiq, vertpiece, above] at (0,1.34) {$w_4(A)$ : la cime};
          \end{scope}
        \end{tikzpicture}}
      \begin{center}\scriptsize Figure 3.10\end{center}
    \end{column}
  \end{columns}
\end{frame}

\begin{frame}{Exemple 3.8 --- La fougère de Barnsley (1988)}
  \small
  $w_i\col{x}{y} = \mat{a_i}{b_i}{c_i}{d_i}\col{x}{y}+\col{e_i}{f_i}$ :
  \begin{center}
  \small
  \begin{tabular}{crrrrrrr}
    \toprule
    $i$ & $a_i$ & $b_i$ & $c_i$ & $d_i$ & $e_i$ & $f_i$ & $p_i$ \\
    \midrule
    $1$ & $0$ & $0$ & $0$ & $\num{0.16}$ & $0$ & $0$ & $\num{0.01}$ \\
    $2$ & $\num{0.85}$ & $\num{0.04}$ & $\num{-0.04}$ & $\num{0.85}$ & $0$ & $\num{1.6}$ & $\num{0.85}$ \\
    $3$ & $\num{0.20}$ & $\num{-0.26}$ & $\num{0.23}$ & $\num{0.22}$ & $0$ & $\num{1.6}$ & $\num{0.07}$ \\
    $4$ & $\num{-0.15}$ & $\num{0.28}$ & $\num{0.26}$ & $\num{0.24}$ & $0$ & $\num{0.44}$ & $\num{0.07}$ \\
    \bottomrule
  \end{tabular}
  \end{center}
  \footnotesize
  \begin{itemize}
    \item $w_1$ écrase tout sur un petit segment vertical, la base de la
      tige : déterminant nul (application non inversible, ce que la
      définition 3.3 permet).
    \item $w_2$ : similitude de rapport
      $\sqrt{\num{0.85}^2+\num{0.04}^2}\approx\num{0.851}$, rotation
      d'environ $3^\circ$ en sens horaire ; elle envoie la fougère sur la
      fougère privée de sa première paire de folioles.
    \item $w_3$, $w_4$ : les deux folioles du bas ; $w_4$ a un déterminant
      négatif, elle retourne la fougère.
    \item La colonne $p_i$ servira au jeu du chaos (section 9).
  \end{itemize}
\end{frame}

\begin{frame}{Figure 3.11 --- La fougère et ses quatre pièces}
  \FigBox[0.66\textheight]{%
    \begin{tikzpicture}[scale=0.52]
      % Fougère de Barnsley, 20 000 points
% Fichier généré par generer.py : ne pas modifier à la main.
\fill[nuage]
  (1.402,9.252)rectangle++(0.04,0.04) (1.56,9.404)rectangle++(0.04,0.04) (1.7,9.527)rectangle++(0.04,0.04) (1.824,9.626)rectangle++(0.04,0.04) (1.933,9.706)rectangle++(0.04,0.04)
  (2.029,9.769)rectangle++(0.04,0.04) (2.113,9.818)rectangle++(0.04,0.04) (2.187,9.857)rectangle++(0.04,0.04) (2.251,9.887)rectangle++(0.04,0.04) (2.307,9.911)rectangle++(0.04,0.04)
  (2.412,3.408)rectangle++(0.04,0.04) (2.184,4.397)rectangle++(0.04,0.04) (2.03,5.246)rectangle++(0.04,0.04) (1.147,2.217)rectangle++(0.04,0.04) (1.061,3.435)rectangle++(0.04,0.04)
  (1.037,4.473)rectangle++(0.04,0.04) (1.059,5.357)rectangle++(0.04,0.04) (1.112,6.107)rectangle++(0.04,0.04) (1.187,6.743)rectangle++(0.04,0.04) (1.277,7.28)rectangle++(0.04,0.04)
  (1.374,7.733)rectangle++(0.04,0.04) (1.475,8.115)rectangle++(0.04,0.04) (1.576,8.435)rectangle++(0.04,0.04) (1.675,8.703)rectangle++(0.04,0.04) (1.77,8.926)rectangle++(0.04,0.04)
  (1.859,9.113)rectangle++(0.04,0.04) (1.942,9.268)rectangle++(0.04,0.04) (2.02,9.396)rectangle++(0.04,0.04) (-2.06,4.121)rectangle++(0.04,0.04) (-1.589,5.181)rectangle++(0.04,0.04)
  (-1.145,6.064)rectangle++(0.04,0.04) (-0.733,6.796)rectangle++(0.04,0.04) (-0.354,7.402)rectangle++(0.04,0.04) (-2.016,3.136)rectangle++(0.04,0.04) (-1.591,4.343)rectangle++(0.04,0.04)
  (-1.181,5.351)rectangle++(0.04,0.04) (-0.792,6.192)rectangle++(0.04,0.04) (-0.427,6.891)rectangle++(0.04,0.04) (-0.09,7.471)rectangle++(0.04,0.04) (0.22,7.95)rectangle++(0.04,0.04)
  (0.503,8.345)rectangle++(0.04,0.04) (2.244,2.563)rectangle++(0.04,0.04) (2.007,3.685)rectangle++(0.04,0.04) (1.852,4.649)rectangle++(0.04,0.04) (1.758,5.473)rectangle++(0.04,0.04)
  (1.711,6.178)rectangle++(0.04,0.04) (1.699,6.779)rectangle++(0.04,0.04) (-1.444,3.471)rectangle++(0.04,0.04) (-1.091,4.604)rectangle++(0.04,0.04) (1.435,1.251)rectangle++(0.04,0.04)
  (1.268,2.603)rectangle++(0.04,0.04) (-0.444,2.453)rectangle++(0.04,0.04) (-0.282,3.699)rectangle++(0.04,0.04) (-0.094,4.752)rectangle++(0.04,0.04) (0.108,5.639)rectangle++(0.04,0.04)
  (0.315,6.385)rectangle++(0.04,0.04) (0.521,7.011)rectangle++(0.04,0.04) (0.721,7.535)rectangle++(0.04,0.04) (0.912,7.972)rectangle++(0.04,0.04) (1.092,8.336)rectangle++(0.04,0.04)
  (1.26,8.638)rectangle++(0.04,0.04) (1.414,8.888)rectangle++(0.04,0.04) (1.555,9.094)rectangle++(0.04,0.04) (2.296,3.017)rectangle++(0.04,0.04) (2.07,4.069)rectangle++(0.04,0.04)
  (1.92,4.972)rectangle++(0.04,0.04) (1.087,2.122)rectangle++(0.04,0.04) (1.006,3.357)rectangle++(0.04,0.04) (0.988,4.409)rectangle++(0.04,0.04) (1.014,5.305)rectangle++(0.04,0.04)
  (1.072,6.065)rectangle++(0.04,0.04) (1.151,6.708)rectangle++(0.04,0.04) (1.245,7.252)rectangle++(0.04,0.04) (1.346,7.711)rectangle++(0.04,0.04) (1.45,8.096)rectangle++(0.04,0.04)
  (1.554,8.42)rectangle++(0.04,0.04) (1.656,8.691)rectangle++(0.04,0.04) (1.753,8.917)rectangle++(0.04,0.04) (1.844,9.106)rectangle++(0.04,0.04) (1.93,9.262)rectangle++(0.04,0.04)
  (2.009,9.392)rectangle++(0.04,0.04) (2.081,9.499)rectangle++(0.04,0.04) (2.146,9.587)rectangle++(0.04,0.04) (-2.085,4.192)rectangle++(0.04,0.04) (1.469,0.894)rectangle++(0.04,0.04)
  (1.282,2.297)rectangle++(0.04,0.04) (1.18,3.498)rectangle++(0.04,0.04) (1.14,4.522)rectangle++(0.04,0.04) (1.148,5.394)rectangle++(0.04,0.04) (1.189,6.135)rectangle++(0.04,0.04)
  (1.254,6.764)rectangle++(0.04,0.04) (1.334,7.295)rectangle++(0.04,0.04) (1.424,7.744)rectangle++(0.04,0.04) (1.518,8.121)rectangle++(0.04,0.04) (-1.829,3.725)rectangle++(0.04,0.04)
  (-1.408,4.835)rectangle++(0.04,0.04) (-1.006,5.763)rectangle++(0.04,0.04) (-0.626,6.535)rectangle++(0.04,0.04) (-0.273,7.176)rectangle++(0.04,0.04) (0.053,7.707)rectangle++(0.04,0.04)
  (-2.014,3.297)rectangle++(0.04,0.04) (-1.583,4.479)rectangle++(0.04,0.04) (-1.168,5.466)rectangle++(0.04,0.04) (-0.777,6.289)rectangle++(0.04,0.04) (-0.411,6.973)rectangle++(0.04,0.04)
  (-0.072,7.54)rectangle++(0.04,0.04) (0.238,8.008)rectangle++(0.04,0.04) (0.52,8.394)rectangle++(0.04,0.04) (0.776,8.71)rectangle++(0.04,0.04) (1.006,8.969)rectangle++(0.04,0.04)
  (1.211,9.179)rectangle++(0.04,0.04) (1.395,9.35)rectangle++(0.04,0.04) (1.557,9.488)rectangle++(0.04,0.04) (1.701,9.599)rectangle++(0.04,0.04) (1.828,9.687)rectangle++(0.04,0.04)
  (1.939,9.757)rectangle++(0.04,0.04) (2.036,9.812)rectangle++(0.04,0.04) (2.121,9.855)rectangle++(0.04,0.04) (2.195,9.888)rectangle++(0.04,0.04) (2.259,9.913)rectangle++(0.04,0.04)
  (2.314,9.932)rectangle++(0.04,0.04) (2.362,9.946)rectangle++(0.04,0.04) (2.404,9.956)rectangle++(0.04,0.04) (2.439,9.963)rectangle++(0.04,0.04) (2.406,3.455)rectangle++(0.04,0.04)
  (2.181,4.437)rectangle++(0.04,0.04) (-0.739,3.067)rectangle++(0.04,0.04) (-0.507,4.233)rectangle++(0.04,0.04) (-0.264,5.214)rectangle++(0.04,0.04) (-1.43,2.675)rectangle++(0.04,0.04)
  (-1.11,3.927)rectangle++(0.04,0.04) (-0.789,4.979)rectangle++(0.04,0.04) (-1.474,2.503)rectangle++(0.04,0.04) (-1.155,3.783)rectangle++(0.04,0.04) (-0.832,4.858)rectangle++(0.04,0.04)
  (-0.515,5.758)rectangle++(0.04,0.04) (-1.621,2.737)rectangle++(0.04,0.04) (-1.271,3.988)rectangle++(0.04,0.04) (-0.923,5.037)rectangle++(0.04,0.04) (-0.585,5.914)rectangle++(0.04,0.04)
  (-0.263,6.647)rectangle++(0.04,0.04) (0.04,7.256)rectangle++(0.04,0.04) (0.322,7.763)rectangle++(0.04,0.04) (-1.975,3.371)rectangle++(0.04,0.04) (-1.293,1.876)rectangle++(0.04,0.04)
  (-1.026,3.243)rectangle++(0.04,0.04) (-0.744,4.394)rectangle++(0.04,0.04) (-0.459,5.361)rectangle++(0.04,0.04) (1.552,1.597)rectangle++(0.04,0.04) (1.381,2.892)rectangle++(0.04,0.04)
  (1.288,3.999)rectangle++(0.04,0.04) (1.252,4.944)rectangle++(0.04,0.04) (-1.056,2.965)rectangle++(0.04,0.04) (-0.781,4.158)rectangle++(0.04,0.04) (-0.5,5.162)rectangle++(0.04,0.04)
  (-0.221,6.004)rectangle++(0.04,0.04) (0.05,6.708)rectangle++(0.04,0.04) (0.309,7.296)rectangle++(0.04,0.04) (0.552,7.786)rectangle++(0.04,0.04) (0.779,8.192)rectangle++(0.04,0.04)
  (0.987,8.528)rectangle++(0.04,0.04) (-2.041,3.692)rectangle++(0.04,0.04) (-1.589,4.816)rectangle++(0.04,0.04) (-1.161,5.754)rectangle++(0.04,0.04) (-0.759,6.533)rectangle++(0.04,0.04)
  (1.926,1.801)rectangle++(0.04,0.04) (0.198,1.363)rectangle++(0.04,0.04) (-0.336,1.934)rectangle++(0.04,0.04) (-0.21,3.254)rectangle++(0.04,0.04) (-0.051,4.37)rectangle++(0.04,0.04)
  (0.129,5.313)rectangle++(0.04,0.04) (0.32,6.107)rectangle++(0.04,0.04) (0.514,6.774)rectangle++(0.04,0.04) (-0.02,1.067)rectangle++(0.04,0.04) (0.023,2.504)rectangle++(0.04,0.04)
  (0.118,3.724)rectangle++(0.04,0.04) (0.247,4.757)rectangle++(0.04,0.04) (0.398,5.629)rectangle++(0.04,0.04) (0.561,6.365)rectangle++(0.04,0.04) (0.729,6.984)rectangle++(0.04,0.04)
  (0.897,7.504)rectangle++(0.04,0.04) (1.061,7.938)rectangle++(0.04,0.04) (2.046,2.611)rectangle++(0.04,0.04) (-0.291,2.634)rectangle++(0.04,0.04) (-0.144,3.847)rectangle++(0.04,0.04)
  (0.029,4.872)rectangle++(0.04,0.04) (0.218,5.736)rectangle++(0.04,0.04) (0.412,6.463)rectangle++(0.04,0.04) (0.607,7.073)rectangle++(0.04,0.04) (0.796,7.584)rectangle++(0.04,0.04)
  (0.978,8.011)rectangle++(0.04,0.04) (1.15,8.366)rectangle++(0.04,0.04) (1.31,8.662)rectangle++(0.04,0.04) (1.457,8.906)rectangle++(0.04,0.04) (1.593,9.108)rectangle++(0.04,0.04)
  (1.716,9.274)rectangle++(0.04,0.04) (1.827,9.411)rectangle++(0.04,0.04) (1.928,9.522)rectangle++(0.04,0.04) (2.017,9.613)rectangle++(0.04,0.04) (2.097,9.687)rectangle++(0.04,0.04)
  (2.168,9.746)rectangle++(0.04,0.04) (2.23,9.794)rectangle++(0.04,0.04) (-2.122,4.257)rectangle++(0.04,0.04) (-1.635,5.299)rectangle++(0.04,0.04) (-1.18,6.166)rectangle++(0.04,0.04)
  (-0.759,6.884)rectangle++(0.04,0.04) (-0.372,7.478)rectangle++(0.04,0.04) (-0.019,7.968)rectangle++(0.04,0.04) (-2.097,3.337)rectangle++(0.04,0.04) (-1.651,4.517)rectangle++(0.04,0.04)
  (-1.225,5.502)rectangle++(0.04,0.04) (-0.823,6.322)rectangle++(0.04,0.04) (-0.449,7.002)rectangle++(0.04,0.04) (-0.104,7.566)rectangle++(0.04,0.04) (0.212,8.032)rectangle++(0.04,0.04)
  (-2.067,3.405)rectangle++(0.04,0.04) (-1.623,4.573)rectangle++(0.04,0.04) (-1.199,5.548)rectangle++(0.04,0.04) (-0.799,6.36)rectangle++(0.04,0.04) (1.883,1.749)rectangle++(0.04,0.04)
  (1.669,3.007)rectangle++(0.04,0.04) (1.536,4.086)rectangle++(0.04,0.04) (1.467,5.007)rectangle++(0.04,0.04) (1.445,5.794)rectangle++(0.04,0.04) (1.458,6.463)rectangle++(0.04,0.04)
  (-1.41,3.346)rectangle++(0.04,0.04) (-1.067,4.497)rectangle++(0.04,0.04) (-0.729,5.461)rectangle++(0.04,0.04) (-0.404,6.267)rectangle++(0.04,0.04) (-0.095,6.94)rectangle++(0.04,0.04)
  (1.94,2.071)rectangle++(0.04,0.04) (1.73,3.279)rectangle++(0.04,0.04) (1.599,4.314)rectangle++(0.04,0.04) (0.951,1.881)rectangle++(0.04,0.04) (0.881,3.157)rectangle++(0.04,0.04)
  (0.873,4.245)rectangle++(0.04,0.04) (0.91,5.169)rectangle++(0.04,0.04) (0.978,5.954)rectangle++(0.04,0.04) (1.067,6.618)rectangle++(0.04,0.04) (1.17,7.178)rectangle++(0.04,0.04)
  (1.279,7.651)rectangle++(0.04,0.04) (1.391,8.049)rectangle++(0.04,0.04) (1.502,8.382)rectangle++(0.04,0.04) (1.61,8.661)rectangle++(0.04,0.04) (2.166,2.927)rectangle++(0.04,0.04)
  (-0.349,2.731)rectangle++(0.04,0.04) (-0.19,3.932)rectangle++(0.04,0.04) (-0.006,4.946)rectangle++(0.04,0.04) (0.19,5.8)rectangle++(0.04,0.04) (0.392,6.519)rectangle++(0.04,0.04)
  (0.591,7.122)rectangle++(0.04,0.04) (0.785,7.626)rectangle++(0.04,0.04) (0.97,8.047)rectangle++(0.04,0.04) (1.145,8.397)rectangle++(0.04,0.04) (1.307,8.688)rectangle++(0.04,0.04)
  (1.456,8.929)rectangle++(0.04,0.04) (-0.02,1.412)rectangle++(0.04,0.04) (0.037,2.797)rectangle++(0.04,0.04) (0.76,1.111)rectangle++(0.04,0.04) (0.688,2.51)rectangle++(0.04,0.04)
  (-0.02,0.385)rectangle++(0.04,0.04) (-0.125,1.669)rectangle++(0.04,0.04) (-0.042,3.02)rectangle++(0.04,0.04) (0.083,4.165)rectangle++(0.04,0.04) (0.235,5.133)rectangle++(0.04,0.04)
  (1.385,1.723)rectangle++(0.04,0.04) (0.257,1.204)rectangle++(0.04,0.04) (0.265,2.609)rectangle++(0.04,0.04) (-0.647,2.224)rectangle++(0.04,0.04) (-0.463,3.512)rectangle++(0.04,0.04)
  (-0.255,4.6)rectangle++(0.04,0.04) (-0.035,5.517)rectangle++(0.04,0.04) (0.189,6.287)rectangle++(0.04,0.04) (1.715,1.988)rectangle++(0.04,0.04) (1.535,3.217)rectangle++(0.04,0.04)
  (1.431,4.27)rectangle++(0.04,0.04) (1.385,5.168)rectangle++(0.04,0.04) (1.382,5.934)rectangle++(0.04,0.04) (1.41,6.585)rectangle++(0.04,0.04) (1.459,7.137)rectangle++(0.04,0.04)
  (1.762,2.522)rectangle++(0.04,0.04) (-0.325,2.549)rectangle++(0.04,0.04) (-0.176,3.776)rectangle++(0.04,0.04) (-0.001,4.813)rectangle++(0.04,0.04) (0.19,5.687)rectangle++(0.04,0.04)
  (0.386,6.423)rectangle++(0.04,0.04) (0.583,7.04)rectangle++(0.04,0.04) (0.775,7.557)rectangle++(0.04,0.04) (0.959,7.989)rectangle++(0.04,0.04) (1.132,8.348)rectangle++(0.04,0.04)
  (1.294,8.647)rectangle++(0.04,0.04) (2.21,2.842)rectangle++(0.04,0.04) (0.447,1.687)rectangle++(0.04,0.04) (0.445,3.012)rectangle++(0.04,0.04) (-0.715,2.354)rectangle++(0.04,0.04)
  (-0.516,3.626)rectangle++(0.04,0.04) (-0.296,4.699)rectangle++(0.04,0.04) (-0.066,5.602)rectangle++(0.04,0.04) (0.166,6.36)rectangle++(0.04,0.04) (1.739,2)rectangle++(0.04,0.04)
  (1.556,3.226)rectangle++(0.04,0.04) (1.449,4.276)rectangle++(0.04,0.04) (1.401,5.173)rectangle++(0.04,0.04) (1.395,5.937)rectangle++(0.04,0.04) (1.436,2.218)rectangle++(0.04,0.04)
  (1.307,3.424)rectangle++(0.04,0.04) (0.745,1.592)rectangle++(0.04,0.04) (0.695,2.919)rectangle++(0.04,0.04) (0.705,4.05)rectangle++(0.04,0.04) (0.759,5.01)rectangle++(0.04,0.04)
  (0.844,5.825)rectangle++(0.04,0.04) (0.948,6.513)rectangle++(0.04,0.04) (1.064,7.095)rectangle++(0.04,0.04) (1.186,7.584)rectangle++(0.04,0.04) (1.309,7.995)rectangle++(0.04,0.04)
  (1.43,8.34)rectangle++(0.04,0.04) (1.547,8.628)rectangle++(0.04,0.04) (2.166,2.903)rectangle++(0.04,0.04) (1.955,3.977)rectangle++(0.04,0.04) (1.819,4.898)rectangle++(0.04,0.04)
  (1.74,5.687)rectangle++(0.04,0.04) (1.704,6.361)rectangle++(0.04,0.04) (1.701,6.935)rectangle++(0.04,0.04) (1.721,7.423)rectangle++(0.04,0.04) (-1.607,3.618)rectangle++(0.04,0.04)
  (-1.223,4.736)rectangle++(0.04,0.04) (-0.853,5.67)rectangle++(0.04,0.04) (-0.5,6.45)rectangle++(0.04,0.04) (-0.169,7.099)rectangle++(0.04,0.04) (0.138,7.637)rectangle++(0.04,0.04)
  (0.42,8.082)rectangle++(0.04,0.04) (2.183,2.479)rectangle++(0.04,0.04) (-0.229,2.636)rectangle++(0.04,0.04) (-0.092,3.846)rectangle++(0.04,0.04) (0.074,4.869)rectangle++(0.04,0.04)
  (-1.272,2.677)rectangle++(0.04,0.04) (-0.977,3.923)rectangle++(0.04,0.04) (-0.675,4.97)rectangle++(0.04,0.04) (1.475,1.447)rectangle++(0.04,0.04) (1.31,2.767)rectangle++(0.04,0.04)
  (1.222,3.896)rectangle++(0.04,0.04) (1.192,4.859)rectangle++(0.04,0.04) (-1.046,2.932)rectangle++(0.04,0.04) (-0.774,4.13)rectangle++(0.04,0.04) (-0.495,5.138)rectangle++(0.04,0.04)
  (-0.217,5.983)rectangle++(0.04,0.04) (0.052,6.691)rectangle++(0.04,0.04) (0.31,7.281)rectangle++(0.04,0.04) (0.552,7.773)rectangle++(0.04,0.04) (-1.932,3.426)rectangle++(0.04,0.04)
  (-1.507,4.586)rectangle++(0.04,0.04) (-1.1,5.554)rectangle++(0.04,0.04) (-0.715,6.361)rectangle++(0.04,0.04) (-1.818,2.824)rectangle++(0.04,0.04) (-1.435,4.069)rectangle++(0.04,0.04)
  (-1.059,5.113)rectangle++(0.04,0.04) (-0.698,5.984)rectangle++(0.04,0.04) (-0.356,6.711)rectangle++(0.04,0.04) (-0.036,7.315)rectangle++(0.04,0.04) (-1.93,3.19)rectangle++(0.04,0.04)
  (-1.515,4.385)rectangle++(0.04,0.04) (-1.115,5.384)rectangle++(0.04,0.04) (-0.734,6.217)rectangle++(0.04,0.04) (-0.378,6.91)rectangle++(0.04,0.04) (-0.047,7.485)rectangle++(0.04,0.04)
  (0.257,7.96)rectangle++(0.04,0.04) (0.535,8.352)rectangle++(0.04,0.04) (0.787,8.674)rectangle++(0.04,0.04) (1.013,8.938)rectangle++(0.04,0.04) (1.217,9.153)rectangle++(0.04,0.04)
  (1.398,9.327)rectangle++(0.04,0.04) (1.559,9.469)rectangle++(0.04,0.04) (1.702,9.582)rectangle++(0.04,0.04) (1.828,9.673)rectangle++(0.04,0.04) (1.938,9.745)rectangle++(0.04,0.04)
  (2.035,9.802)rectangle++(0.04,0.04) (2.12,9.846)rectangle++(0.04,0.04) (2.193,9.881)rectangle++(0.04,0.04) (2.257,9.907)rectangle++(0.04,0.04) (2.313,9.927)rectangle++(0.04,0.04)
  (2.415,3.414)rectangle++(0.04,0.04) (2.187,4.401)rectangle++(0.04,0.04) (2.033,5.25)rectangle++(0.04,0.04) (1.936,5.977)rectangle++(0.04,0.04) (1.882,6.599)rectangle++(0.04,0.04)
  (1.862,7.13)rectangle++(0.04,0.04) (1.866,7.583)rectangle++(0.04,0.04) (1.887,7.967)rectangle++(0.04,0.04) (1.92,8.292)rectangle++(0.04,0.04) (1.962,8.568)rectangle++(0.04,0.04)
  (2.008,8.801)rectangle++(0.04,0.04) (2.057,8.996)rectangle++(0.04,0.04) (2.106,9.161)rectangle++(0.04,0.04) (2.232,3.176)rectangle++(0.04,0.04) (2.022,4.207)rectangle++(0.04,0.04)
  (-0.711,2.979)rectangle++(0.04,0.04) (-0.487,4.157)rectangle++(0.04,0.04) (-0.25,5.149)rectangle++(0.04,0.04) (-0.009,5.983)rectangle++(0.04,0.04) (0.23,6.682)rectangle++(0.04,0.04)
  (0.46,7.267)rectangle++(0.04,0.04) (0.68,7.755)rectangle++(0.04,0.04) (0.886,8.16)rectangle++(0.04,0.04) (1.077,8.497)rectangle++(0.04,0.04) (1.253,8.776)rectangle++(0.04,0.04)
  (1.414,9.005)rectangle++(0.04,0.04) (1.56,9.194)rectangle++(0.04,0.04) (1.692,9.349)rectangle++(0.04,0.04) (1.81,9.475)rectangle++(0.04,0.04) (-2.123,4.09)rectangle++(0.04,0.04)
  (-1.643,5.157)rectangle++(0.04,0.04) (-1.192,6.046)rectangle++(0.04,0.04) (-0.774,6.783)rectangle++(0.04,0.04) (-0.389,7.392)rectangle++(0.04,0.04) (-2.021,3.126)rectangle++(0.04,0.04)
  (-1.595,4.334)rectangle++(0.04,0.04) (-1.185,5.344)rectangle++(0.04,0.04) (-0.795,6.186)rectangle++(0.04,0.04) (-0.431,6.886)rectangle++(0.04,0.04) (-0.093,7.467)rectangle++(0.04,0.04)
  (0.217,7.947)rectangle++(0.04,0.04) (0.501,8.342)rectangle++(0.04,0.04) (2.243,2.562)rectangle++(0.04,0.04) (2.007,3.684)rectangle++(0.04,0.04) (0.713,1.836)rectangle++(0.04,0.04)
  (0.677,3.128)rectangle++(0.04,0.04) (-0.699,2.433)rectangle++(0.04,0.04) (-0.499,3.692)rectangle++(0.04,0.04) (-0.279,4.755)rectangle++(0.04,0.04) (-0.049,5.649)rectangle++(0.04,0.04)
  (0.182,6.4)rectangle++(0.04,0.04) (0.409,7.029)rectangle++(0.04,0.04) (0.626,7.554)rectangle++(0.04,0.04) (0.832,7.992)rectangle++(0.04,0.04) (1.025,8.356)rectangle++(0.04,0.04)
  (1.203,8.658)rectangle++(0.04,0.04) (1.367,8.907)rectangle++(0.04,0.04) (1.516,9.113)rectangle++(0.04,0.04) (1.651,9.281)rectangle++(0.04,0.04) (1.772,9.419)rectangle++(0.04,0.04)
  (1.881,9.532)rectangle++(0.04,0.04) (2.369,3.207)rectangle++(0.04,0.04) (2.14,4.227)rectangle++(0.04,0.04) (1.986,5.104)rectangle++(0.04,0.04) (1.89,5.855)rectangle++(0.04,0.04)
  (-1.165,3.312)rectangle++(0.04,0.04) (-0.86,4.458)rectangle++(0.04,0.04) (-1.352,2.372)rectangle++(0.04,0.04) (0.85,0.648)rectangle++(0.04,0.04) (0.746,2.113)rectangle++(0.04,0.04)
  (0.716,3.362)rectangle++(0.04,0.04) (-0.752,2.493)rectangle++(0.04,0.04) (0.794,0.833)rectangle++(0.04,0.04) (-0.079,1.955)rectangle++(0.04,0.04) (0.009,3.261)rectangle++(0.04,0.04)
  (0.136,4.368)rectangle++(0.04,0.04) (0.288,5.303)rectangle++(0.04,0.04) (0.455,6.092)rectangle++(0.04,0.04) (1.62,2.01)rectangle++(0.04,0.04) (1.455,3.24)rectangle++(0.04,0.04)
  (1.365,4.292)rectangle++(0.04,0.04) (1.329,5.19)rectangle++(0.04,0.04) (1.335,5.954)rectangle++(0.04,0.04) (1.371,6.604)rectangle++(0.04,0.04) (1.427,7.155)rectangle++(0.04,0.04)
  (1.772,2.518)rectangle++(0.04,0.04) (1.605,3.666)rectangle++(0.04,0.04) (0.768,1.727)rectangle++(0.04,0.04) (0.72,3.033)rectangle++(0.04,0.04) (-0.666,2.422)rectangle++(0.04,0.04)
  (0.761,0.838)rectangle++(0.04,0.04) (0.678,2.278)rectangle++(0.04,0.04) (0.665,3.506)rectangle++(0.04,0.04) (0.703,4.549)rectangle++(0.04,0.04) (0.778,5.435)rectangle++(0.04,0.04)
  (0.876,6.185)rectangle++(0.04,0.04) (0.99,6.818)rectangle++(0.04,0.04) (-0.02,1.074)rectangle++(0.04,0.04) (0.024,2.51)rectangle++(0.04,0.04) (0.118,3.729)rectangle++(0.04,0.04)
  (0.248,4.761)rectangle++(0.04,0.04) (0.399,5.633)rectangle++(0.04,0.04) (0.562,6.368)rectangle++(0.04,0.04) (0.73,6.987)rectangle++(0.04,0.04) (0.898,7.506)rectangle++(0.04,0.04)
  (1.061,7.94)rectangle++(0.04,0.04) (1.218,8.303)rectangle++(0.04,0.04) (2.125,2.739)rectangle++(0.04,0.04) (1.913,3.84)rectangle++(0.04,0.04) (-0.637,2.874)rectangle++(0.04,0.04)
  (-0.429,4.064)rectangle++(0.04,0.04) (-0.204,5.068)rectangle++(0.04,0.04) (0.027,5.912)rectangle++(0.04,0.04) (0.257,6.621)rectangle++(0.04,0.04) (0.481,7.213)rectangle++(0.04,0.04)
  (0.696,7.708)rectangle++(0.04,0.04) (0.897,8.12)rectangle++(0.04,0.04) (1.085,8.463)rectangle++(0.04,0.04) (-2.004,3.7)rectangle++(0.04,0.04) (-1.558,4.822)rectangle++(0.04,0.04)
  (-1.134,5.757)rectangle++(0.04,0.04) (-0.735,6.535)rectangle++(0.04,0.04) (-0.366,7.18)rectangle++(0.04,0.04) (-0.026,7.714)rectangle++(0.04,0.04) (0.284,8.154)rectangle++(0.04,0.04)
  (0.566,8.516)rectangle++(0.04,0.04) (0.819,8.812)rectangle++(0.04,0.04) (1.047,9.054)rectangle++(0.04,0.04) (1.25,9.25)rectangle++(0.04,0.04) (-2.176,3.911)rectangle++(0.04,0.04)
  (-1.696,5.008)rectangle++(0.04,0.04) (-1.243,5.921)rectangle++(0.04,0.04) (-1.809,2.606)rectangle++(0.04,0.04) (-0.02,0.4)rectangle++(0.04,0.04) (-0.003,1.937)rectangle++(0.04,0.04)
  (0.073,3.243)rectangle++(0.04,0.04) (0.88,1.227)rectangle++(0.04,0.04) (0.795,2.604)rectangle++(0.04,0.04) (0.777,3.778)rectangle++(0.04,0.04) (0.81,4.776)rectangle++(0.04,0.04)
  (0.877,5.624)rectangle++(0.04,0.04) (0.968,6.341)rectangle++(0.04,0.04) (1.613,2.204)rectangle++(0.04,0.04) (1.457,3.405)rectangle++(0.04,0.04) (1.372,4.432)rectangle++(0.04,0.04)
  (1.018,1.85)rectangle++(0.04,0.04) (0.937,3.128)rectangle++(0.04,0.04) (-0.647,2.493)rectangle++(0.04,0.04) (-0.799,1.989)rectangle++(0.04,0.04) (-0.602,3.318)rectangle++(0.04,0.04)
  (-0.381,4.441)rectangle++(0.04,0.04) (1.283,1.397)rectangle++(0.04,0.04) (-0.128,2.191)rectangle++(0.04,0.04) (-0.023,3.464)rectangle++(0.04,0.04) (0.117,4.542)rectangle++(0.04,0.04)
  (0.279,5.452)rectangle++(0.04,0.04) (0.453,6.219)rectangle++(0.04,0.04) (0.631,6.864)rectangle++(0.04,0.04) (0.809,7.406)rectangle++(0.04,0.04) (0.982,7.859)rectangle++(0.04,0.04)
  (1.147,8.237)rectangle++(0.04,0.04) (1.302,8.552)rectangle++(0.04,0.04) (2.182,2.821)rectangle++(0.04,0.04) (1.965,3.907)rectangle++(0.04,0.04) (1.824,4.838)rectangle++(0.04,0.04)
  (1.742,5.636)rectangle++(0.04,0.04) (1.704,6.317)rectangle++(0.04,0.04) (1.699,6.897)rectangle++(0.04,0.04) (1.718,7.391)rectangle++(0.04,0.04) (1.754,7.81)rectangle++(0.04,0.04)
  (1.906,2.76)rectangle++(0.04,0.04) (1.729,3.866)rectangle++(0.04,0.04) (0.806,1.807)rectangle++(0.04,0.04) (0.755,3.1)rectangle++(0.04,0.04) (0.764,4.201)rectangle++(0.04,0.04)
  (0.815,5.137)rectangle++(0.04,0.04) (0.896,5.93)rectangle++(0.04,0.04) (0.997,6.601)rectangle++(0.04,0.04) (1.109,7.167)rectangle++(0.04,0.04) (-1.663,3.421)rectangle++(0.04,0.04)
  (-1.279,4.57)rectangle++(0.04,0.04) (-0.906,5.532)rectangle++(0.04,0.04) (1.668,1.522)rectangle++(0.04,0.04) (1.476,2.823)rectangle++(0.04,0.04) (1.365,3.937)rectangle++(0.04,0.04)
  (1.316,4.888)rectangle++(0.04,0.04) (1.312,5.698)rectangle++(0.04,0.04) (1.341,6.387)rectangle++(0.04,0.04) (1.393,6.972)rectangle++(0.04,0.04) (1.461,7.467)rectangle++(0.04,0.04)
  (-1.67,3.568)rectangle++(0.04,0.04) (-1.279,4.695)rectangle++(0.04,0.04) (-0.902,5.639)rectangle++(0.04,0.04) (-0.543,6.425)rectangle++(0.04,0.04) (-0.207,7.079)rectangle++(0.04,0.04)
  (0.105,7.622)rectangle++(0.04,0.04) (0.392,8.071)rectangle++(0.04,0.04) (0.654,8.44)rectangle++(0.04,0.04) (0.891,8.744)rectangle++(0.04,0.04) (1.105,8.993)rectangle++(0.04,0.04)
  (1.297,9.196)rectangle++(0.04,0.04) (2.363,2.974)rectangle++(0.04,0.04) (2.125,4.03)rectangle++(0.04,0.04) (1.966,4.937)rectangle++(0.04,0.04) (1.866,5.714)rectangle++(0.04,0.04)
  (1.812,6.378)rectangle++(0.04,0.04) (1.793,6.945)rectangle++(0.04,0.04) (-1.468,3.529)rectangle++(0.04,0.04) (-1.109,4.655)rectangle++(0.04,0.04) (-0.759,5.597)rectangle++(0.04,0.04)
  (-0.423,6.384)rectangle++(0.04,0.04) (-0.107,7.04)rectangle++(0.04,0.04) (0.189,7.584)rectangle++(0.04,0.04) (0.462,8.035)rectangle++(0.04,0.04) (2.163,2.478)rectangle++(0.04,0.04)
  (1.936,3.616)rectangle++(0.04,0.04) (1.788,4.593)rectangle++(0.04,0.04) (1,1.997)rectangle++(0.04,0.04) (0.928,3.254)rectangle++(0.04,0.04) (0.917,4.325)rectangle++(0.04,0.04)
  (0.95,5.236)rectangle++(0.04,0.04) (1.015,6.008)rectangle++(0.04,0.04) (1.101,6.663)rectangle++(0.04,0.04) (1.2,7.216)rectangle++(0.04,0.04) (1.306,7.681)rectangle++(0.04,0.04)
  (1.415,8.073)rectangle++(0.04,0.04) (1.524,8.402)rectangle++(0.04,0.04) (1.629,8.677)rectangle++(0.04,0.04) (1.73,8.906)rectangle++(0.04,0.04) (-1.991,3.946)rectangle++(0.04,0.04)
  (-1.537,5.03)rectangle++(0.04,0.04) (-1.107,5.933)rectangle++(0.04,0.04) (-0.706,6.684)rectangle++(0.04,0.04) (-0.335,7.306)rectangle++(0.04,0.04) (0.005,7.819)rectangle++(0.04,0.04)
  (0.315,8.242)rectangle++(0.04,0.04) (0.595,8.59)rectangle++(0.04,0.04) (0.847,8.874)rectangle++(0.04,0.04) (2.34,2.78)rectangle++(0.04,0.04) (2.098,3.866)rectangle++(0.04,0.04)
  (1.936,4.798)rectangle++(0.04,0.04) (1.835,5.597)rectangle++(0.04,0.04) (1.782,6.28)rectangle++(0.04,0.04) (1.763,6.863)rectangle++(0.04,0.04) (-1.453,3.504)rectangle++(0.04,0.04)
  (1.182,0.893)rectangle++(0.04,0.04) (1.038,2.308)rectangle++(0.04,0.04) (0.972,3.517)rectangle++(0.04,0.04) (-0.741,2.586)rectangle++(0.04,0.04) (-0.529,3.824)rectangle++(0.04,0.04)
  (1.133,1.21)rectangle++(0.04,0.04) (1.009,2.58)rectangle++(0.04,0.04) (0.959,3.749)rectangle++(0.04,0.04) (0.963,4.744)rectangle++(0.04,0.04) (1.006,5.59)rectangle++(0.04,0.04)
  (1.076,6.308)rectangle++(0.04,0.04) (-1.446,3.224)rectangle++(0.04,0.04) (-1.102,4.395)rectangle++(0.04,0.04) (-0.763,5.376)rectangle++(0.04,0.04) (-0.436,6.196)rectangle++(0.04,0.04)
  (-0.125,6.88)rectangle++(0.04,0.04) (0.167,7.449)rectangle++(0.04,0.04) (0.438,7.922)rectangle++(0.04,0.04) (0.687,8.312)rectangle++(0.04,0.04) (-2.045,3.576)rectangle++(0.04,0.04)
  (-1.597,4.717)rectangle++(0.04,0.04) (-1.171,5.67)rectangle++(0.04,0.04) (-0.771,6.462)rectangle++(0.04,0.04) (-0.399,7.12)rectangle++(0.04,0.04) (-0.057,7.664)rectangle++(0.04,0.04)
  (0.256,8.113)rectangle++(0.04,0.04) (0.54,8.482)rectangle++(0.04,0.04) (0.796,8.784)rectangle++(0.04,0.04) (1.026,9.031)rectangle++(0.04,0.04) (1.231,9.232)rectangle++(0.04,0.04)
  (1.413,9.394)rectangle++(0.04,0.04) (1.575,9.524)rectangle++(0.04,0.04) (1.718,9.629)rectangle++(0.04,0.04) (1.843,9.712)rectangle++(0.04,0.04) (-2.178,4.15)rectangle++(0.04,0.04)
  (-1.687,5.21)rectangle++(0.04,0.04) (-1.228,6.093)rectangle++(0.04,0.04) (-0.802,6.824)rectangle++(0.04,0.04) (-0.411,7.429)rectangle++(0.04,0.04) (-0.055,7.927)rectangle++(0.04,0.04)
  (2.21,2.318)rectangle++(0.04,0.04) (1.969,3.478)rectangle++(0.04,0.04) (1.811,4.474)rectangle++(0.04,0.04) (1.716,5.327)rectangle++(0.04,0.04) (1.669,6.055)rectangle++(0.04,0.04)
  (1.659,6.676)rectangle++(0.04,0.04) (1.675,7.205)rectangle++(0.04,0.04) (1.71,7.653)rectangle++(0.04,0.04) (1.757,8.033)rectangle++(0.04,0.04) (1.813,8.354)rectangle++(0.04,0.04)
  (1.873,8.625)rectangle++(0.04,0.04) (1.935,8.852)rectangle++(0.04,0.04) (1.996,9.043)rectangle++(0.04,0.04) (2.056,9.203)rectangle++(0.04,0.04) (2.251,3.173)rectangle++(0.04,0.04)
  (2.038,4.204)rectangle++(0.04,0.04) (1.898,5.088)rectangle++(0.04,0.04) (1.122,2.145)rectangle++(0.04,0.04) (1.038,3.374)rectangle++(0.04,0.04) (1.015,4.423)rectangle++(0.04,0.04)
  (1.037,5.315)rectangle++(0.04,0.04) (-1.196,2.997)rectangle++(0.04,0.04) (-0.899,4.191)rectangle++(0.04,0.04) (-0.598,5.195)rectangle++(0.04,0.04) (-1.492,2.594)rectangle++(0.04,0.04)
  (-1.166,3.861)rectangle++(0.04,0.04) (-0.839,4.925)rectangle++(0.04,0.04) (-0.518,5.816)rectangle++(0.04,0.04) (-0.21,6.56)rectangle++(0.04,0.04) (0.082,7.181)rectangle++(0.04,0.04)
  (0.354,7.697)rectangle++(0.04,0.04) (-0.02,1.215)rectangle++(0.04,0.04) (0.029,2.629)rectangle++(0.04,0.04) (0.128,3.83)rectangle++(0.04,0.04) (0.26,4.847)rectangle++(0.04,0.04)
  (0.412,5.705)rectangle++(0.04,0.04) (0.577,6.429)rectangle++(0.04,0.04) (0.745,7.038)rectangle++(0.04,0.04) (0.913,7.549)rectangle++(0.04,0.04) (1.959,2.479)rectangle++(0.04,0.04)
  (1.762,3.625)rectangle++(0.04,0.04) (-0.611,2.792)rectangle++(0.04,0.04) (-0.41,3.994)rectangle++(0.04,0.04) (-0.191,5.007)rectangle++(0.04,0.04) (0.036,5.86)rectangle++(0.04,0.04)
  (-0.02,0.921)rectangle++(0.04,0.04) (0.018,2.38)rectangle++(0.04,0.04) (0.108,3.618)rectangle++(0.04,0.04) (0.234,4.667)rectangle++(0.04,0.04) (1.254,1.611)rectangle++(0.04,0.04)
  (1.128,2.915)rectangle++(0.04,0.04) (1.074,4.029)rectangle++(0.04,0.04) (1.072,4.978)rectangle++(0.04,0.04) (1.108,5.785)rectangle++(0.04,0.04) (-1.304,3.116)rectangle++(0.04,0.04)
  (-0.986,4.297)rectangle++(0.04,0.04) (-0.668,5.288)rectangle++(0.04,0.04) (-0.359,6.118)rectangle++(0.04,0.04) (-0.062,6.811)rectangle++(0.04,0.04) (0.217,7.388)rectangle++(0.04,0.04)
  (0.478,7.867)rectangle++(0.04,0.04) (0.719,8.264)rectangle++(0.04,0.04) (0.939,8.592)rectangle++(0.04,0.04) (1.14,8.862)rectangle++(0.04,0.04) (1.321,9.083)rectangle++(0.04,0.04)
  (1.484,9.264)rectangle++(0.04,0.04) (1.63,9.411)rectangle++(0.04,0.04) (1.76,9.531)rectangle++(0.04,0.04) (1.875,9.627)rectangle++(0.04,0.04) (2.397,3.228)rectangle++(0.04,0.04)
  (2.164,4.244)rectangle++(0.04,0.04) (2.007,5.117)rectangle++(0.04,0.04) (1.909,5.865)rectangle++(0.04,0.04) (1.855,6.505)rectangle++(0.04,0.04) (1.835,7.052)rectangle++(0.04,0.04)
  (1.839,7.517)rectangle++(0.04,0.04) (1.862,7.912)rectangle++(0.04,0.04) (1.897,8.247)rectangle++(0.04,0.04) (1.94,8.53)rectangle++(0.04,0.04) (1.988,8.769)rectangle++(0.04,0.04)
  (2.038,8.97)rectangle++(0.04,0.04) (2.089,9.14)rectangle++(0.04,0.04) (2.139,9.281)rectangle++(0.04,0.04) (2.187,9.4)rectangle++(0.04,0.04) (2.233,9.498)rectangle++(0.04,0.04)
  (2.276,9.581)rectangle++(0.04,0.04) (2.315,9.649)rectangle++(0.04,0.04) (2.352,9.705)rectangle++(0.04,0.04) (2.385,9.751)rectangle++(0.04,0.04) (2.415,9.789)rectangle++(0.04,0.04)
  (2.442,9.821)rectangle++(0.04,0.04) (2.467,9.846)rectangle++(0.04,0.04) (2.488,9.867)rectangle++(0.04,0.04) (2.507,9.883)rectangle++(0.04,0.04) (2.524,9.897)rectangle++(0.04,0.04)
  (2.539,9.907)rectangle++(0.04,0.04) (2.553,9.916)rectangle++(0.04,0.04) (2.564,9.923)rectangle++(0.04,0.04) (2.574,9.928)rectangle++(0.04,0.04) (2.583,9.932)rectangle++(0.04,0.04)
  (2.591,9.935)rectangle++(0.04,0.04) (2.597,9.937)rectangle++(0.04,0.04) (2.603,9.939)rectangle++(0.04,0.04) (2.608,9.94)rectangle++(0.04,0.04) (2.612,9.941)rectangle++(0.04,0.04)
  (2.616,9.942)rectangle++(0.04,0.04) (2.619,9.942)rectangle++(0.04,0.04) (2.621,9.942)rectangle++(0.04,0.04) (2.624,9.942)rectangle++(0.04,0.04) (2.626,9.942)rectangle++(0.04,0.04)
  (2.627,9.942)rectangle++(0.04,0.04) (2.629,9.942)rectangle++(0.04,0.04) (2.63,9.942)rectangle++(0.04,0.04) (-2.08,4.381)rectangle++(0.04,0.04) (-1.595,5.403)rectangle++(0.04,0.04)
  (1.735,1.312)rectangle++(0.04,0.04) (1.525,2.642)rectangle++(0.04,0.04) (1.4,3.781)rectangle++(0.04,0.04) (-0.724,2.743)rectangle++(0.04,0.04) (-0.508,3.956)rectangle++(0.04,0.04)
  (-0.276,4.98)rectangle++(0.04,0.04) (-0.038,5.84)rectangle++(0.04,0.04) (0.199,6.562)rectangle++(0.04,0.04) (-0.02,1.033)rectangle++(0.04,0.04) (0.022,2.475)rectangle++(0.04,0.04)
  (0.116,3.699)rectangle++(0.04,0.04) (0.244,4.736)rectangle++(0.04,0.04) (0.395,5.612)rectangle++(0.04,0.04) (0.558,6.351)rectangle++(0.04,0.04) (-1.561,3.114)rectangle++(0.04,0.04)
  (-1.204,4.306)rectangle++(0.04,0.04) (-0.854,5.304)rectangle++(0.04,0.04) (-0.516,6.139)rectangle++(0.04,0.04) (1.779,1.769)rectangle++(0.04,0.04) (0.211,1.317)rectangle++(0.04,0.04)
  (0.23,2.707)rectangle++(0.04,0.04) (0.302,3.888)rectangle++(0.04,0.04) (0.41,4.889)rectangle++(0.04,0.04) (0.542,5.736)rectangle++(0.04,0.04) (0.688,6.45)rectangle++(0.04,0.04)
  (0.84,7.051)rectangle++(0.04,0.04) (0.994,7.556)rectangle++(0.04,0.04) (1.145,7.979)rectangle++(0.04,0.04) (1.29,8.333)rectangle++(0.04,0.04) (1.428,8.627)rectangle++(0.04,0.04)
  (2.184,2.872)rectangle++(0.04,0.04) (1.969,3.95)rectangle++(0.04,0.04) (-0.654,2.911)rectangle++(0.04,0.04) (-0.442,4.097)rectangle++(0.04,0.04) (-0.214,5.096)rectangle++(0.04,0.04)
  (0.02,5.936)rectangle++(0.04,0.04) (0.252,6.641)rectangle++(0.04,0.04) (-1.698,3.108)rectangle++(0.04,0.04) (-1.321,4.306)rectangle++(0.04,0.04) (-0.953,5.309)rectangle++(0.04,0.04)
  (-0.6,6.147)rectangle++(0.04,0.04) (1.794,1.749)rectangle++(0.04,0.04) (1.592,3.011)rectangle++(0.04,0.04) (1.472,4.092)rectangle++(0.04,0.04) (-0.791,2.828)rectangle++(0.04,0.04)
  (-0.561,4.031)rectangle++(0.04,0.04) (-0.318,5.045)rectangle++(0.04,0.04) (-0.071,5.898)rectangle++(0.04,0.04) (0.174,6.612)rectangle++(0.04,0.04) (0.41,7.209)rectangle++(0.04,0.04)
  (0.635,7.708)rectangle++(0.04,0.04) (0.845,8.122)rectangle++(0.04,0.04) (1.041,8.466)rectangle++(0.04,0.04) (1.222,8.751)rectangle++(0.04,0.04) (1.386,8.986)rectangle++(0.04,0.04)
  (-2.08,3.885)rectangle++(0.04,0.04) (-1.615,4.981)rectangle++(0.04,0.04) (-1.639,2.313)rectangle++(0.04,0.04) (-1.303,3.628)rectangle++(0.04,0.04) (-0.965,4.732)rectangle++(0.04,0.04)
  (-0.633,5.657)rectangle++(0.04,0.04) (-0.314,6.43)rectangle++(0.04,0.04) (-0.012,7.074)rectangle++(0.04,0.04) (0.271,7.61)rectangle++(0.04,0.04) (0.532,8.054)rectangle++(0.04,0.04)
  (-2.009,3.483)rectangle++(0.04,0.04) (-1.57,4.637)rectangle++(0.04,0.04) (-1.151,5.601)rectangle++(0.04,0.04) (-1.708,2.556)rectangle++(0.04,0.04) (-1.351,3.837)rectangle++(0.04,0.04)
  (-0.997,4.912)rectangle++(0.04,0.04) (-0.654,5.811)rectangle++(0.04,0.04) (1.708,1.655)rectangle++(0.04,0.04) (1.516,2.934)rectangle++(0.04,0.04) (1.403,4.03)rectangle++(0.04,0.04)
  (1.352,4.965)rectangle++(0.04,0.04) (1.346,5.763)rectangle++(0.04,0.04) (1.372,6.441)rectangle++(0.04,0.04) (1.422,7.016)rectangle++(0.04,0.04) (-1.561,3.459)rectangle++(0.04,0.04)
  (-1.233,1.991)rectangle++(0.04,0.04) (-0.97,3.338)rectangle++(0.04,0.04) (-0.694,4.472)rectangle++(0.04,0.04) (-0.413,5.425)rectangle++(0.04,0.04) (-0.136,6.224)rectangle++(0.04,0.04)
  (0.131,6.892)rectangle++(0.04,0.04) (0.385,7.449)rectangle++(0.04,0.04) (0.623,7.913)rectangle++(0.04,0.04) (0.844,8.297)rectangle++(0.04,0.04) (1.047,8.615)rectangle++(0.04,0.04)
  (1.232,8.877)rectangle++(0.04,0.04) (1.4,9.092)rectangle++(0.04,0.04) (1.552,9.269)rectangle++(0.04,0.04) (1.688,9.413)rectangle++(0.04,0.04) (1.809,9.529)rectangle++(0.04,0.04)
  (-2.137,4.101)rectangle++(0.04,0.04) (-1.655,5.168)rectangle++(0.04,0.04) (-1.202,6.055)rectangle++(0.04,0.04) (-0.782,6.791)rectangle++(0.04,0.04) (-0.395,7.4)rectangle++(0.04,0.04)
  (-0.042,7.902)rectangle++(0.04,0.04) (0.278,8.315)rectangle++(0.04,0.04) (0.567,8.652)rectangle++(0.04,0.04) (0.826,8.928)rectangle++(0.04,0.04) (1.057,9.152)rectangle++(0.04,0.04)
  (1.262,9.333)rectangle++(0.04,0.04) (1.444,9.479)rectangle++(0.04,0.04) (1.604,9.596)rectangle++(0.04,0.04) (1.745,9.688)rectangle++(0.04,0.04) (2.434,3.209)rectangle++(0.04,0.04)
  (2.195,4.226)rectangle++(0.04,0.04) (2.032,5.101)rectangle++(0.04,0.04) (1.106,2.183)rectangle++(0.04,0.04) (-0.367,2.324)rectangle++(0.04,0.04) (-0.222,3.586)rectangle++(0.04,0.04)
  (1.02,1.233)rectangle++(0.04,0.04) (0.914,2.603)rectangle++(0.04,0.04) (0.879,3.773)rectangle++(0.04,0.04) (0.896,4.768)rectangle++(0.04,0.04) (0.95,5.613)rectangle++(0.04,0.04)
  (1.03,6.329)rectangle++(0.04,0.04) (1.126,6.935)rectangle++(0.04,0.04) (1.232,7.446)rectangle++(0.04,0.04) (-1.711,3.511)rectangle++(0.04,0.04) (1.222,0.828)rectangle++(0.04,0.04)
  (1.07,2.251)rectangle++(0.04,0.04) (0.997,3.467)rectangle++(0.04,0.04) (0.984,4.503)rectangle++(0.04,0.04) (1.014,5.384)rectangle++(0.04,0.04) (1.075,6.132)rectangle++(0.04,0.04)
  (1.538,2.181)rectangle++(0.04,0.04) (1.393,3.389)rectangle++(0.04,0.04) (1.317,4.421)rectangle++(0.04,0.04) (1.294,5.301)rectangle++(0.04,0.04) (1.31,6.051)rectangle++(0.04,0.04)
  (1.353,6.687)rectangle++(0.04,0.04) (1.416,7.226)rectangle++(0.04,0.04) (-1.617,3.504)rectangle++(0.04,0.04) (-1.236,4.639)rectangle++(0.04,0.04) (-0.867,5.589)rectangle++(0.04,0.04)
  (-0.516,6.382)rectangle++(0.04,0.04) (-0.186,7.041)rectangle++(0.04,0.04) (-1.889,3.095)rectangle++(0.04,0.04) (-1.484,4.303)rectangle++(0.04,0.04) (-1.092,5.313)rectangle++(0.04,0.04)
  (-0.717,6.156)rectangle++(0.04,0.04) (-0.366,6.857)rectangle++(0.04,0.04) (-1.877,3.013)rectangle++(0.04,0.04) (-1.18,1.82)rectangle++(0.04,0.04) (-0.933,3.191)rectangle++(0.04,0.04)
  (-0.667,4.345)rectangle++(0.04,0.04) (-0.396,5.317)rectangle++(0.04,0.04) (-0.126,6.131)rectangle++(0.04,0.04) (0.136,6.813)rectangle++(0.04,0.04) (0.386,7.382)rectangle++(0.04,0.04)
  (0.621,7.855)rectangle++(0.04,0.04) (0.84,8.248)rectangle++(0.04,0.04) (1.042,8.574)rectangle++(0.04,0.04) (-2.042,3.715)rectangle++(0.04,0.04) (-1.395,1.937)rectangle++(0.04,0.04)
  (-1.111,3.298)rectangle++(0.04,0.04) (-0.815,4.444)rectangle++(0.04,0.04) (-0.517,5.406)rectangle++(0.04,0.04) (-1.53,2.66)rectangle++(0.04,0.04) (-1.196,3.918)rectangle++(0.04,0.04)
  (-0.862,4.974)rectangle++(0.04,0.04) (-0.536,5.859)rectangle++(0.04,0.04) (-0.224,6.598)rectangle++(0.04,0.04) (0.072,7.213)rectangle++(0.04,0.04) (0.347,7.725)rectangle++(0.04,0.04)
  (0.602,8.148)rectangle++(0.04,0.04) (0.835,8.498)rectangle++(0.04,0.04) (-2.064,3.651)rectangle++(0.04,0.04) (-1.61,4.782)rectangle++(0.04,0.04) (1.563,1.159)rectangle++(0.04,0.04)
  (1.373,2.519)rectangle++(0.04,0.04) (1.265,3.682)rectangle++(0.04,0.04) (1.221,4.676)rectangle++(0.04,0.04) (1.222,5.522)rectangle++(0.04,0.04) (1.258,6.241)rectangle++(0.04,0.04)
  (1.316,6.85)rectangle++(0.04,0.04) (1.391,7.366)rectangle++(0.04,0.04) (1.475,7.802)rectangle++(0.04,0.04) (1.946,2.686)rectangle++(0.04,0.04) (1.759,3.801)rectangle++(0.04,0.04)
  (1.645,4.757)rectangle++(0.04,0.04) (1.068,1.999)rectangle++(0.04,0.04) (0.985,3.253)rectangle++(0.04,0.04) (-0.67,2.531)rectangle++(0.04,0.04) (-0.47,3.775)rectangle++(0.04,0.04)
  (-1.097,2.311)rectangle++(0.04,0.04) (-0.842,3.605)rectangle++(0.04,0.04) (-0.574,4.694)rectangle++(0.04,0.04) (-0.302,5.609)rectangle++(0.04,0.04) (-0.035,6.376)rectangle++(0.04,0.04)
  (0.223,7.017)rectangle++(0.04,0.04) (0.468,7.552)rectangle++(0.04,0.04) (0.698,7.996)rectangle++(0.04,0.04) (0.911,8.365)rectangle++(0.04,0.04) (-2.014,3.639)rectangle++(0.04,0.04)
  (-1.569,4.77)rectangle++(0.04,0.04) (-1.145,5.713)rectangle++(0.04,0.04) (-0.747,6.498)rectangle++(0.04,0.04) (-1.86,2.847)rectangle++(0.04,0.04) (-1.469,4.09)rectangle++(0.04,0.04)
  (-1.088,5.132)rectangle++(0.04,0.04) (1.583,1.379)rectangle++(0.04,0.04) (1.398,2.705)rectangle++(0.04,0.04) (1.294,3.839)rectangle++(0.04,0.04) (1.252,4.808)rectangle++(0.04,0.04)
  (1.254,5.633)rectangle++(0.04,0.04) (-1.235,3.117)rectangle++(0.04,0.04) (1.041,0.857)rectangle++(0.04,0.04) (0.917,2.283)rectangle++(0.04,0.04) (0.868,3.5)rectangle++(0.04,0.04)
  (0.876,4.537)rectangle++(0.04,0.04) (0.924,5.417)rectangle++(0.04,0.04) (1,6.164)rectangle++(0.04,0.04) (1.094,6.796)rectangle++(0.04,0.04) (1.2,7.329)rectangle++(0.04,0.04)
  (1.311,7.778)rectangle++(0.04,0.04) (1.423,8.155)rectangle++(0.04,0.04) (1.533,8.471)rectangle++(0.04,0.04) (1.64,8.735)rectangle++(0.04,0.04) (1.741,8.955)rectangle++(0.04,0.04)
  (1.836,9.139)rectangle++(0.04,0.04) (1.924,9.291)rectangle++(0.04,0.04) (2.005,9.416)rectangle++(0.04,0.04) (2.079,9.52)rectangle++(0.04,0.04) (2.145,9.605)rectangle++(0.04,0.04)
  (-2.089,4.196)rectangle++(0.04,0.04) (-1.61,5.246)rectangle++(0.04,0.04) (-1.161,6.12)rectangle++(0.04,0.04) (-0.744,6.844)rectangle++(0.04,0.04) (-0.361,7.444)rectangle++(0.04,0.04)
  (-0.011,7.938)rectangle++(0.04,0.04) (0.306,8.344)rectangle++(0.04,0.04) (0.591,8.676)rectangle++(0.04,0.04) (0.847,8.947)rectangle++(0.04,0.04) (1.076,9.168)rectangle++(0.04,0.04)
  (-2.19,3.853)rectangle++(0.04,0.04) (1.39,0.786)rectangle++(0.04,0.04) (1.211,2.208)rectangle++(0.04,0.04) (1.115,3.425)rectangle++(0.04,0.04) (1.083,4.463)rectangle++(0.04,0.04)
  (1.097,5.346)rectangle++(0.04,0.04) (1.144,6.097)rectangle++(0.04,0.04) (1.214,6.733)rectangle++(0.04,0.04) (1.686,2.361)rectangle++(0.04,0.04) (1.525,3.536)rectangle++(0.04,0.04)
  (1.436,4.541)rectangle++(0.04,0.04) (1.4,5.398)rectangle++(0.04,0.04) (1.403,6.129)rectangle++(0.04,0.04) (1.436,6.75)rectangle++(0.04,0.04) (1.488,7.276)rectangle++(0.04,0.04)
  (1.554,7.721)rectangle++(0.04,0.04) (1.911,2.687)rectangle++(0.04,0.04) (1.73,3.804)rectangle++(0.04,0.04) (1.62,4.76)rectangle++(0.04,0.04) (1.072,1.994)rectangle++(0.04,0.04)
  (0.989,3.248)rectangle++(0.04,0.04) (-0.668,2.531)rectangle++(0.04,0.04) (-0.469,3.774)rectangle++(0.04,0.04) (-0.25,4.823)rectangle++(0.04,0.04) (-0.021,5.706)rectangle++(0.04,0.04)
  (0.208,6.447)rectangle++(0.04,0.04) (0.432,7.068)rectangle++(0.04,0.04) (0.648,7.587)rectangle++(0.04,0.04) (-1.864,3.407)rectangle++(0.04,0.04) (-1.45,4.567)rectangle++(0.04,0.04)
  (-1.052,5.536)rectangle++(0.04,0.04) (-0.675,6.344)rectangle++(0.04,0.04) (-0.322,7.016)rectangle++(0.04,0.04) (0.004,7.572)rectangle++(0.04,0.04) (0.304,8.032)rectangle++(0.04,0.04)
  (0.578,8.412)rectangle++(0.04,0.04) (0.825,8.723)rectangle++(0.04,0.04) (-2.124,3.698)rectangle++(0.04,0.04) (-1.66,4.824)rectangle++(0.04,0.04) (-1.22,5.763)rectangle++(0.04,0.04)
  (-0.809,6.544)rectangle++(0.04,0.04) (-0.428,7.191)rectangle++(0.04,0.04) (-0.078,7.725)rectangle++(0.04,0.04) (0.24,8.166)rectangle++(0.04,0.04) (0.529,8.528)rectangle++(0.04,0.04)
  (0.788,8.824)rectangle++(0.04,0.04) (1.021,9.065)rectangle++(0.04,0.04) (1.228,9.26)rectangle++(0.04,0.04) (1.412,9.418)rectangle++(0.04,0.04) (1.575,9.545)rectangle++(0.04,0.04)
  (1.718,9.647)rectangle++(0.04,0.04) (1.844,9.727)rectangle++(0.04,0.04) (1.954,9.791)rectangle++(0.04,0.04) (2.051,9.84)rectangle++(0.04,0.04) (2.134,9.878)rectangle++(0.04,0.04)
  (2.428,3.356)rectangle++(0.04,0.04) (2.196,4.351)rectangle++(0.04,0.04) (2.039,5.207)rectangle++(0.04,0.04) (-0.02,0.816)rectangle++(0.04,0.04) (0.214,0.621)rectangle++(0.04,0.04)
  (0.205,2.115)rectangle++(0.04,0.04) (0.256,3.386)rectangle++(0.04,0.04) (0.351,4.464)rectangle++(0.04,0.04) (0.475,5.377)rectangle++(0.04,0.04) (1.417,1.844)rectangle++(0.04,0.04)
  (1.276,3.107)rectangle++(0.04,0.04) (1.207,4.186)rectangle++(0.04,0.04) (1.191,5.106)rectangle++(0.04,0.04) (1.214,5.889)rectangle++(0.04,0.04) (1.265,6.553)rectangle++(0.04,0.04)
  (1.336,7.116)rectangle++(0.04,0.04) (1.418,7.591)rectangle++(0.04,0.04) (1.506,7.992)rectangle++(0.04,0.04) (1.598,8.329)rectangle++(0.04,0.04) (1.689,8.612)rectangle++(0.04,0.04)
  (1.778,8.849)rectangle++(0.04,0.04) (1.863,9.047)rectangle++(0.04,0.04) (1.943,9.211)rectangle++(0.04,0.04) (2.018,9.348)rectangle++(0.04,0.04) (2.297,3.198)rectangle++(0.04,0.04)
  (2.078,4.223)rectangle++(0.04,0.04) (1.933,5.102)rectangle++(0.04,0.04) (1.845,5.856)rectangle++(0.04,0.04) (1.801,6.5)rectangle++(0.04,0.04) (1.788,7.049)rectangle++(0.04,0.04)
  (1.8,7.516)rectangle++(0.04,0.04) (1.817,2.702)rectangle++(0.04,0.04) (1.651,3.82)rectangle++(0.04,0.04) (1.554,4.777)rectangle++(0.04,0.04) (1.509,5.595)rectangle++(0.04,0.04)
  (1.505,6.291)rectangle++(0.04,0.04) (1.528,6.884)rectangle++(0.04,0.04) (1.572,7.386)rectangle++(0.04,0.04) (1.63,7.812)rectangle++(0.04,0.04) (1.695,8.171)rectangle++(0.04,0.04)
  (1.766,8.474)rectangle++(0.04,0.04) (1.838,8.728)rectangle++(0.04,0.04) (-1.923,3.932)rectangle++(0.04,0.04) (-1.479,5.015)rectangle++(0.04,0.04) (-1.059,5.918)rectangle++(0.04,0.04)
  (-0.666,6.669)rectangle++(0.04,0.04) (-0.301,7.292)rectangle++(0.04,0.04) (-1.977,3.124)rectangle++(0.04,0.04) (-1.558,4.331)rectangle++(0.04,0.04) (1.429,1.064)rectangle++(0.04,0.04)
  (1.255,2.444)rectangle++(0.04,0.04) (1.162,3.623)rectangle++(0.04,0.04) (1.131,4.629)rectangle++(0.04,0.04) (-0.999,2.867)rectangle++(0.04,0.04) (-0.736,4.074)rectangle++(0.04,0.04)
  (-0.465,5.088)rectangle++(0.04,0.04) (1.477,1.53)rectangle++(0.04,0.04) (1.315,2.838)rectangle++(0.04,0.04) (1.229,3.956)rectangle++(0.04,0.04) (1.2,4.909)rectangle++(0.04,0.04)
  (1.214,5.721)rectangle++(0.04,0.04) (1.402,2.119)rectangle++(0.04,0.04) (1.275,3.341)rectangle++(0.04,0.04) (1.215,4.385)rectangle++(0.04,0.04) (1.206,5.275)rectangle++(0.04,0.04)
  (1.234,6.032)rectangle++(0.04,0.04) (-1.343,3.2)rectangle++(0.04,0.04) (-1.016,4.37)rectangle++(0.04,0.04) (-0.691,5.351)rectangle++(0.04,0.04) (-0.375,6.172)rectangle++(0.04,0.04)
  (-0.074,6.858)rectangle++(0.04,0.04) (0.209,7.428)rectangle++(0.04,0.04) (0.473,7.902)rectangle++(0.04,0.04) (0.716,8.294)rectangle++(0.04,0.04) (-0.02,1.31)rectangle++(0.04,0.04)
  (0.033,2.711)rectangle++(0.04,0.04) (0.134,3.899)rectangle++(0.04,0.04) (0.268,4.905)rectangle++(0.04,0.04) (0.422,5.755)rectangle++(0.04,0.04) (-1.433,2.952)rectangle++(0.04,0.04)
  (1.024,0.766)rectangle++(0.04,0.04) (-0.016,1.993)rectangle++(0.04,0.04) (0.064,3.291)rectangle++(0.04,0.04) (0.184,4.391)rectangle++(0.04,0.04) (-1.126,2.597)rectangle++(0.04,0.04)
  (-0.855,3.849)rectangle++(0.04,0.04) (-0.575,4.902)rectangle++(0.04,0.04) (-0.295,5.786)rectangle++(0.04,0.04) (-0.022,6.526)rectangle++(0.04,0.04) (0.24,7.144)rectangle++(0.04,0.04)
  (0.488,7.659)rectangle++(0.04,0.04) (0.719,8.087)rectangle++(0.04,0.04) (0.932,8.441)rectangle++(0.04,0.04) (1.128,8.734)rectangle++(0.04,0.04) (2.259,2.819)rectangle++(0.04,0.04)
  (2.031,3.902)rectangle++(0.04,0.04) (1.88,4.832)rectangle++(0.04,0.04) (1.789,5.628)rectangle++(0.04,0.04) (1.744,6.309)rectangle++(0.04,0.04) (-1.313,3.378)rectangle++(0.04,0.04)
  (-0.983,4.52)rectangle++(0.04,0.04) (-0.657,5.477)rectangle++(0.04,0.04) (-0.341,6.278)rectangle++(0.04,0.04) (-0.041,6.946)rectangle++(0.04,0.04) (0.241,7.502)rectangle++(0.04,0.04)
  (0.502,7.964)rectangle++(0.04,0.04) (0.743,8.345)rectangle++(0.04,0.04) (0.963,8.66)rectangle++(0.04,0.04) (1.163,8.919)rectangle++(0.04,0.04) (-2.107,3.819)rectangle++(0.04,0.04)
  (1.368,0.799)rectangle++(0.04,0.04) (1.192,2.22)rectangle++(0.04,0.04) (1.1,3.436)rectangle++(0.04,0.04) (1.07,4.473)rectangle++(0.04,0.04) (1.087,5.355)rectangle++(0.04,0.04)
  (-1.196,3.017)rectangle++(0.04,0.04) (-0.898,4.209)rectangle++(0.04,0.04) (-0.597,5.209)rectangle++(0.04,0.04) (-0.302,6.048)rectangle++(0.04,0.04) (-0.017,6.749)rectangle++(0.04,0.04)
  (0.254,7.334)rectangle++(0.04,0.04) (0.507,7.82)rectangle++(0.04,0.04) (0.741,8.223)rectangle++(0.04,0.04) (0.957,8.556)rectangle++(0.04,0.04) (1.153,8.83)rectangle++(0.04,0.04)
  (1.331,9.056)rectangle++(0.04,0.04) (1.492,9.24)rectangle++(0.04,0.04) (1.635,9.391)rectangle++(0.04,0.04) (1.763,9.513)rectangle++(0.04,0.04) (1.877,9.612)rectangle++(0.04,0.04)
  (1.978,9.691)rectangle++(0.04,0.04) (-2.145,4.176)rectangle++(0.04,0.04) (-1.659,5.232)rectangle++(0.04,0.04) (-1.203,6.109)rectangle++(0.04,0.04) (-0.78,6.837)rectangle++(0.04,0.04)
  (-0.392,7.439)rectangle++(0.04,0.04) (-0.038,7.935)rectangle++(0.04,0.04) (2.21,2.325)rectangle++(0.04,0.04) (1.969,3.484)rectangle++(0.04,0.04) (1.811,4.479)rectangle++(0.04,0.04)
  (1.716,5.331)rectangle++(0.04,0.04) (1.67,6.059)rectangle++(0.04,0.04) (1.66,6.679)rectangle++(0.04,0.04) (1.676,7.207)rectangle++(0.04,0.04) (1.71,7.655)rectangle++(0.04,0.04)
  (1.758,8.035)rectangle++(0.04,0.04) (1.969,2.815)rectangle++(0.04,0.04) (1.784,3.911)rectangle++(0.04,0.04) (1.67,4.849)rectangle++(0.04,0.04) (1.612,5.651)rectangle++(0.04,0.04)
  (1.594,6.335)rectangle++(0.04,0.04) (1.606,6.917)rectangle++(0.04,0.04) (1.639,7.412)rectangle++(0.04,0.04) (1.688,7.83)rectangle++(0.04,0.04) (1.746,8.185)rectangle++(0.04,0.04)
  (1.809,8.483)rectangle++(0.04,0.04) (-0.02,1.341)rectangle++(0.04,0.04) (0.034,2.736)rectangle++(0.04,0.04) (0.137,3.921)rectangle++(0.04,0.04) (1.06,1.406)rectangle++(0.04,0.04)
  (0.955,2.749)rectangle++(0.04,0.04) (0.92,3.895)rectangle++(0.04,0.04) (0.935,4.87)rectangle++(0.04,0.04) (0.988,5.698)rectangle++(0.04,0.04) (-1.305,3.07)rectangle++(0.04,0.04)
  (1.038,0.827)rectangle++(0.04,0.04) (-0.029,2.01)rectangle++(0.04,0.04) (0.054,3.306)rectangle++(0.04,0.04) (0.176,4.404)rectangle++(0.04,0.04) (0.323,5.332)rectangle++(0.04,0.04)
  (0.486,6.116)rectangle++(0.04,0.04) (0.655,6.775)rectangle++(0.04,0.04) (0.826,7.329)rectangle++(0.04,0.04) (0.993,7.793)rectangle++(0.04,0.04) (1.154,8.18)rectangle++(0.04,0.04)
  (1.306,8.503)rectangle++(0.04,0.04) (1.448,8.772)rectangle++(0.04,0.04) (1.579,8.994)rectangle++(0.04,0.04) (1.7,9.178)rectangle++(0.04,0.04) (1.81,9.33)rectangle++(0.04,0.04)
  (1.909,9.454)rectangle++(0.04,0.04) (1.999,9.556)rectangle++(0.04,0.04) (2.079,9.639)rectangle++(0.04,0.04) (2.151,9.706)rectangle++(0.04,0.04) (2.214,9.76)rectangle++(0.04,0.04)
  (2.27,9.804)rectangle++(0.04,0.04) (2.32,9.839)rectangle++(0.04,0.04) (-2.115,4.287)rectangle++(0.04,0.04) (-1.629,5.325)rectangle++(0.04,0.04) (-1.174,6.187)rectangle++(0.04,0.04)
  (-0.752,6.902)rectangle++(0.04,0.04) (2.028,1.891)rectangle++(0.04,0.04) (1.797,3.122)rectangle++(0.04,0.04) (1.65,4.178)rectangle++(0.04,0.04) (1.568,5.082)rectangle++(0.04,0.04)
  (1.534,5.853)rectangle++(0.04,0.04) (1.536,6.51)rectangle++(0.04,0.04) (1.563,7.068)rectangle++(0.04,0.04) (1.609,7.542)rectangle++(0.04,0.04) (1.667,7.942)rectangle++(0.04,0.04)
  (-1.753,3.72)rectangle++(0.04,0.04) (-1.343,4.828)rectangle++(0.04,0.04) (-0.951,5.754)rectangle++(0.04,0.04) (-0.58,6.525)rectangle++(0.04,0.04) (-0.234,7.166)rectangle++(0.04,0.04)
  (0.085,7.696)rectangle++(0.04,0.04) (2.125,2.299)rectangle++(0.04,0.04) (1.896,3.466)rectangle++(0.04,0.04) (1.748,4.466)rectangle++(0.04,0.04) (1.662,5.322)rectangle++(0.04,0.04)
  (1.624,6.054)rectangle++(0.04,0.04) (-1.27,3.294)rectangle++(0.04,0.04) (-0.95,4.447)rectangle++(0.04,0.04) (-0.632,5.414)rectangle++(0.04,0.04) (-0.323,6.224)rectangle++(0.04,0.04)
  (-0.028,6.899)rectangle++(0.04,0.04) (0.25,7.462)rectangle++(0.04,0.04) (0.509,7.929)rectangle++(0.04,0.04) (0.748,8.315)rectangle++(0.04,0.04) (0.966,8.634)rectangle++(0.04,0.04)
  (1.164,8.897)rectangle++(0.04,0.04) (2.299,2.868)rectangle++(0.04,0.04) (2.067,3.942)rectangle++(0.04,0.04) (0.776,1.913)rectangle++(0.04,0.04) (0.734,3.192)rectangle++(0.04,0.04)
  (0.75,4.28)rectangle++(0.04,0.04) (0.806,5.204)rectangle++(0.04,0.04) (0.891,5.987)rectangle++(0.04,0.04) (0.995,6.65)rectangle++(0.04,0.04) (1.109,7.209)rectangle++(0.04,0.04)
  (1.229,7.679)rectangle++(0.04,0.04) (1.35,8.074)rectangle++(0.04,0.04) (1.468,8.405)rectangle++(0.04,0.04) (1.582,8.682)rectangle++(0.04,0.04) (2.176,2.925)rectangle++(0.04,0.04)
  (1.965,3.995)rectangle++(0.04,0.04) (1.828,4.914)rectangle++(0.04,0.04) (1.748,5.7)rectangle++(0.04,0.04) (1.711,6.371)rectangle++(0.04,0.04) (1.707,6.943)rectangle++(0.04,0.04)
  (1.727,7.43)rectangle++(0.04,0.04) (1.763,7.842)rectangle++(0.04,0.04) (1.81,8.192)rectangle++(0.04,0.04) (1.864,8.487)rectangle++(0.04,0.04) (1.922,8.735)rectangle++(0.04,0.04)
  (1.98,8.944)rectangle++(0.04,0.04) (2.039,9.12)rectangle++(0.04,0.04) (2.096,9.266)rectangle++(0.04,0.04) (2.263,3.199)rectangle++(0.04,0.04) (2.049,4.225)rectangle++(0.04,0.04)
  (1.909,5.105)rectangle++(0.04,0.04) (-0.967,3.151)rectangle++(0.04,0.04) (-0.698,4.313)rectangle++(0.04,0.04) (-0.423,5.29)rectangle++(0.04,0.04) (-0.15,6.11)rectangle++(0.04,0.04)
  (0.115,6.796)rectangle++(0.04,0.04) (0.367,7.368)rectangle++(0.04,0.04) (0.605,7.844)rectangle++(0.04,0.04) (-1.94,3.454)rectangle++(0.04,0.04) (-1.513,4.61)rectangle++(0.04,0.04)
  (-1.104,5.575)rectangle++(0.04,0.04) (1.709,1.481)rectangle++(0.04,0.04) (1.51,2.787)rectangle++(0.04,0.04) (1.393,3.904)rectangle++(0.04,0.04) (-0.758,2.768)rectangle++(0.04,0.04)
  (-0.536,3.98)rectangle++(0.04,0.04) (-0.298,5)rectangle++(0.04,0.04) (-0.056,5.858)rectangle++(0.04,0.04) (0.185,6.578)rectangle++(0.04,0.04) (0.418,7.18)rectangle++(0.04,0.04)
  (0.64,7.683)rectangle++(0.04,0.04) (0.849,8.101)rectangle++(0.04,0.04) (1.044,8.448)rectangle++(0.04,0.04) (1.223,8.735)rectangle++(0.04,0.04) (1.387,8.972)rectangle++(0.04,0.04)
  (1.535,9.167)rectangle++(0.04,0.04) (1.67,9.327)rectangle++(0.04,0.04) (2.344,3.103)rectangle++(0.04,0.04) (2.114,4.14)rectangle++(0.04,0.04) (1.96,5.03)rectangle++(0.04,0.04)
  (1.865,5.794)rectangle++(0.04,0.04) (-0.02,0.91)rectangle++(0.04,0.04) (0.017,2.371)rectangle++(0.04,0.04) (0.107,3.611)rectangle++(0.04,0.04) (0.233,4.661)rectangle++(0.04,0.04)
  (-1.186,2.668)rectangle++(0.04,0.04) (0.908,0.762)rectangle++(0.04,0.04) (-0.038,1.965)rectangle++(0.04,0.04) (0.044,3.268)rectangle++(0.04,0.04) (0.166,4.372)rectangle++(0.04,0.04)
  (0.314,5.306)rectangle++(0.04,0.04) (0.477,6.094)rectangle++(0.04,0.04) (0.647,6.757)rectangle++(0.04,0.04) (0.818,7.314)rectangle++(0.04,0.04) (0.986,7.78)rectangle++(0.04,0.04)
  (1.147,8.17)rectangle++(0.04,0.04) (1.299,8.495)rectangle++(0.04,0.04) (1.442,8.765)rectangle++(0.04,0.04) (1.574,8.989)rectangle++(0.04,0.04) (1.695,9.174)rectangle++(0.04,0.04)
  (1.806,9.326)rectangle++(0.04,0.04) (1.906,9.451)rectangle++(0.04,0.04) (1.996,9.553)rectangle++(0.04,0.04) (2.076,9.637)rectangle++(0.04,0.04) (2.148,9.704)rectangle++(0.04,0.04)
  (2.212,9.759)rectangle++(0.04,0.04) (2.268,9.803)rectangle++(0.04,0.04) (2.318,9.838)rectangle++(0.04,0.04) (-2.115,4.286)rectangle++(0.04,0.04) (1.5,0.909)rectangle++(0.04,0.04)
  (1.309,2.309)rectangle++(0.04,0.04) (-0.36,2.398)rectangle++(0.04,0.04) (-0.717,2.034)rectangle++(0.04,0.04) (-0.53,3.354)rectangle++(0.04,0.04) (-0.319,4.468)rectangle++(0.04,0.04)
  (-0.094,5.407)rectangle++(0.04,0.04) (0.134,6.196)rectangle++(0.04,0.04) (0.36,6.857)rectangle++(0.04,0.04) (0.578,7.41)rectangle++(0.04,0.04) (0.785,7.872)rectangle++(0.04,0.04)
  (-1.911,3.501)rectangle++(0.04,0.04) (-0.02,0.543)rectangle++(0.04,0.04) (0.138,0.555)rectangle++(0.04,0.04) (0.137,2.063)rectangle++(0.04,0.04) (0.197,3.344)rectangle++(0.04,0.04)
  (0.299,4.431)rectangle++(0.04,0.04) (0.429,5.35)rectangle++(0.04,0.04) (0.577,6.127)rectangle++(0.04,0.04) (0.733,6.781)rectangle++(0.04,0.04) (0.892,7.331)rectangle++(0.04,0.04)
  (1.049,7.792)rectangle++(0.04,0.04) (1.201,8.177)rectangle++(0.04,0.04) (1.346,8.499)rectangle++(0.04,0.04) (2.16,2.82)rectangle++(0.04,0.04) (-0.02,0.434)rectangle++(0.04,0.04)
  (0.107,0.529)rectangle++(0.04,0.04) (0.11,2.042)rectangle++(0.04,0.04) (0.173,3.327)rectangle++(0.04,0.04) (0.278,4.417)rectangle++(0.04,0.04) (1.178,1.562)rectangle++(0.04,0.04)
  (1.061,2.877)rectangle++(0.04,0.04) (1.015,3.999)rectangle++(0.04,0.04) (1.021,4.955)rectangle++(0.04,0.04) (1.063,5.767)rectangle++(0.04,0.04) (1.132,6.456)rectangle++(0.04,0.04)
  (1.62,2.274)rectangle++(0.04,0.04) (1.466,3.464)rectangle++(0.04,0.04) (1.382,4.482)rectangle++(0.04,0.04) (1.352,5.351)rectangle++(0.04,0.04) (1.361,6.09)rectangle++(0.04,0.04)
  (1.398,6.718)rectangle++(0.04,0.04) (1.455,7.251)rectangle++(0.04,0.04) (1.525,7.701)rectangle++(0.04,0.04) (1.602,8.081)rectangle++(0.04,0.04) (-1.802,3.735)rectangle++(0.04,0.04)
  (-1.384,4.843)rectangle++(0.04,0.04) (1.546,1.232)rectangle++(0.04,0.04) (1.362,2.582)rectangle++(0.04,0.04) (1.258,3.736)rectangle++(0.04,0.04) (1.217,4.722)rectangle++(0.04,0.04)
  (1.221,5.561)rectangle++(0.04,0.04) (1.258,6.274)rectangle++(0.04,0.04) (1.318,6.879)rectangle++(0.04,0.04) (1.393,7.391)rectangle++(0.04,0.04) (1.478,7.822)rectangle++(0.04,0.04)
  (1.567,8.186)rectangle++(0.04,0.04) (1.657,8.492)rectangle++(0.04,0.04) (1.746,8.748)rectangle++(0.04,0.04) (1.832,8.962)rectangle++(0.04,0.04) (1.913,9.141)rectangle++(0.04,0.04)
  (-0.02,1.446)rectangle++(0.04,0.04) (0.039,2.826)rectangle++(0.04,0.04) (0.144,3.997)rectangle++(0.04,0.04) (0.28,4.988)rectangle++(0.04,0.04) (0.435,5.824)rectangle++(0.04,0.04)
  (0.601,6.53)rectangle++(0.04,0.04) (0.77,7.122)rectangle++(0.04,0.04) (0.937,7.619)rectangle++(0.04,0.04) (1.099,8.035)rectangle++(0.04,0.04) (-1.891,3.609)rectangle++(0.04,0.04)
  (-1.465,4.74)rectangle++(0.04,0.04) (1.529,1.187)rectangle++(0.04,0.04) (1.345,2.544)rectangle++(0.04,0.04) (1.243,3.705)rectangle++(0.04,0.04) (0.833,1.642)rectangle++(0.04,0.04)
  (0.772,2.959)rectangle++(0.04,0.04) (0.772,4.08)rectangle++(0.04,0.04) (0.817,5.034)rectangle++(0.04,0.04) (-1.166,2.884)rectangle++(0.04,0.04) (-0.878,4.095)rectangle++(0.04,0.04)
  (-0.585,5.112)rectangle++(0.04,0.04) (-0.295,5.965)rectangle++(0.04,0.04) (-0.014,6.678)rectangle++(0.04,0.04) (0.253,7.273)rectangle++(0.04,0.04) (0.504,7.768)rectangle++(0.04,0.04)
  (-1.94,3.414)rectangle++(0.04,0.04) (-1.515,4.576)rectangle++(0.04,0.04) (-1.107,5.546)rectangle++(0.04,0.04) (-0.02,0.871)rectangle++(0.04,0.04) (0.229,0.634)rectangle++(0.04,0.04)
  (0.218,2.126)rectangle++(0.04,0.04) (0.268,3.394)rectangle++(0.04,0.04) (0.893,1.314)rectangle++(0.04,0.04) (0.809,2.678)rectangle++(0.04,0.04) (0.793,3.84)rectangle++(0.04,0.04)
  (0.825,4.828)rectangle++(0.04,0.04) (0.892,5.667)rectangle++(0.04,0.04) (0.983,6.378)rectangle++(0.04,0.04) (1.088,6.978)rectangle++(0.04,0.04) (1.202,7.484)rectangle++(0.04,0.04)
  (1.319,7.909)rectangle++(0.04,0.04) (1.435,8.266)rectangle++(0.04,0.04) (1.548,8.565)rectangle++(0.04,0.04) (1.657,8.815)rectangle++(0.04,0.04) (1.759,9.022)rectangle++(0.04,0.04)
  (1.853,9.195)rectangle++(0.04,0.04) (1.941,9.338)rectangle++(0.04,0.04) (2.021,9.456)rectangle++(0.04,0.04) (2.094,9.553)rectangle++(0.04,0.04) (2.343,3.267)rectangle++(0.04,0.04)
  (2.12,4.279)rectangle++(0.04,0.04) (1.971,5.149)rectangle++(0.04,0.04) (1.879,5.894)rectangle++(0.04,0.04) (1.831,6.531)rectangle++(0.04,0.04) (1.815,7.074)rectangle++(0.04,0.04);
\fill[nuage]
  (1.824,7.537)rectangle++(0.04,0.04) (1.849,7.929)rectangle++(0.04,0.04) (-1.713,3.759)rectangle++(0.04,0.04) (-1.308,4.86)rectangle++(0.04,0.04) (-0.919,5.779)rectangle++(0.04,0.04)
  (-0.553,6.545)rectangle++(0.04,0.04) (1.898,1.857)rectangle++(0.04,0.04) (1.686,3.099)rectangle++(0.04,0.04) (1.554,4.163)rectangle++(0.04,0.04) (1.486,5.072)rectangle++(0.04,0.04)
  (1.463,5.848)rectangle++(0.04,0.04) (1.401,2.214)rectangle++(0.04,0.04) (1.277,3.422)rectangle++(0.04,0.04) (-0.656,2.636)rectangle++(0.04,0.04) (-0.454,3.863)rectangle++(0.04,0.04)
  (-0.234,4.898)rectangle++(0.04,0.04) (-0.005,5.769)rectangle++(0.04,0.04) (0.224,6.5)rectangle++(0.04,0.04) (1.769,2.048)rectangle++(0.04,0.04) (1.583,3.266)rectangle++(0.04,0.04)
  (1.474,4.309)rectangle++(0.04,0.04) (1.423,5.2)rectangle++(0.04,0.04) (1.416,5.959)rectangle++(0.04,0.04) (1.439,6.605)rectangle++(0.04,0.04) (1.486,7.153)rectangle++(0.04,0.04)
  (1.547,7.617)rectangle++(0.04,0.04) (1.617,8.009)rectangle++(0.04,0.04) (1.693,8.339)rectangle++(0.04,0.04) (1.77,8.617)rectangle++(0.04,0.04) (1.847,8.849)rectangle++(0.04,0.04)
  (1.922,9.044)rectangle++(0.04,0.04) (1.993,9.207)rectangle++(0.04,0.04) (-2.016,4.073)rectangle++(0.04,0.04) (-1.553,5.139)rectangle++(0.04,0.04) (-1.117,6.026)rectangle++(0.04,0.04)
  (-0.71,6.763)rectangle++(0.04,0.04) (-0.336,7.373)rectangle++(0.04,0.04) (-0.02,1.163)rectangle++(0.04,0.04) (0.027,2.586)rectangle++(0.04,0.04) (-0.688,2.164)rectangle++(0.04,0.04)
  (-0.5,3.463)rectangle++(0.04,0.04) (-0.289,4.56)rectangle++(0.04,0.04) (-0.065,5.484)rectangle++(0.04,0.04) (0.161,6.26)rectangle++(0.04,0.04) (0.385,6.911)rectangle++(0.04,0.04)
  (0.602,7.455)rectangle++(0.04,0.04) (0.808,7.909)rectangle++(0.04,0.04) (1.001,8.286)rectangle++(0.04,0.04) (1.18,8.6)rectangle++(0.04,0.04) (1.345,8.859)rectangle++(0.04,0.04)
  (1.495,9.072)rectangle++(0.04,0.04) (1.631,9.248)rectangle++(0.04,0.04) (1.754,9.392)rectangle++(0.04,0.04) (1.865,9.509)rectangle++(0.04,0.04) (1.963,9.604)rectangle++(0.04,0.04)
  (2.051,9.681)rectangle++(0.04,0.04) (2.128,9.743)rectangle++(0.04,0.04) (-2.129,4.222)rectangle++(0.04,0.04) (-1.643,5.27)rectangle++(0.04,0.04) (-1.188,6.141)rectangle++(0.04,0.04)
  (1.88,1.595)rectangle++(0.04,0.04) (1.66,2.877)rectangle++(0.04,0.04) (1.524,3.975)rectangle++(0.04,0.04) (1.452,4.914)rectangle++(0.04,0.04) (1.429,5.715)rectangle++(0.04,0.04)
  (1.441,6.397)rectangle++(0.04,0.04) (1.478,6.976)rectangle++(0.04,0.04) (1.533,7.467)rectangle++(0.04,0.04) (1.6,7.881)rectangle++(0.04,0.04) (1.673,8.231)rectangle++(0.04,0.04)
  (1.749,8.526)rectangle++(0.04,0.04) (1.826,8.773)rectangle++(0.04,0.04) (1.9,8.981)rectangle++(0.04,0.04) (1.972,9.154)rectangle++(0.04,0.04) (2.04,9.298)rectangle++(0.04,0.04)
  (2.104,9.418)rectangle++(0.04,0.04) (2.163,9.517)rectangle++(0.04,0.04) (2.217,9.599)rectangle++(0.04,0.04) (2.266,9.667)rectangle++(0.04,0.04) (2.311,9.722)rectangle++(0.04,0.04)
  (2.351,9.768)rectangle++(0.04,0.04) (2.387,9.805)rectangle++(0.04,0.04) (2.419,9.835)rectangle++(0.04,0.04) (2.447,9.859)rectangle++(0.04,0.04) (2.472,9.879)rectangle++(0.04,0.04)
  (2.378,3.444)rectangle++(0.04,0.04) (-0.441,2.893)rectangle++(0.04,0.04) (-0.261,4.073)rectangle++(0.04,0.04) (-0.061,5.069)rectangle++(0.04,0.04) (0.148,5.907)rectangle++(0.04,0.04)
  (0.36,6.611)rectangle++(0.04,0.04) (0.568,7.202)rectangle++(0.04,0.04) (0.769,7.695)rectangle++(0.04,0.04) (0.959,8.106)rectangle++(0.04,0.04) (1.137,8.448)rectangle++(0.04,0.04)
  (1.303,8.731)rectangle++(0.04,0.04) (1.454,8.966)rectangle++(0.04,0.04) (1.593,9.159)rectangle++(0.04,0.04) (2.308,3.042)rectangle++(0.04,0.04) (2.081,4.09)rectangle++(0.04,0.04)
  (1.931,4.989)rectangle++(0.04,0.04) (1.838,5.76)rectangle++(0.04,0.04) (1.791,6.419)rectangle++(0.04,0.04) (1.777,6.98)rectangle++(0.04,0.04) (1.787,7.458)rectangle++(0.04,0.04)
  (1.815,7.864)rectangle++(0.04,0.04) (1.855,8.208)rectangle++(0.04,0.04) (-1.784,3.822)rectangle++(0.04,0.04) (-1.366,4.916)rectangle++(0.04,0.04) (-1.573,2.356)rectangle++(0.04,0.04)
  (0.878,0.587)rectangle++(0.04,0.04) (0.768,2.06)rectangle++(0.04,0.04) (0.733,3.316)rectangle++(0.04,0.04) (0.753,4.386)rectangle++(0.04,0.04) (-1.011,2.727)rectangle++(0.04,0.04)
  (-0.752,3.955)rectangle++(0.04,0.04) (-0.483,4.988)rectangle++(0.04,0.04) (-0.214,5.855)rectangle++(0.04,0.04) (0.05,6.582)rectangle++(0.04,0.04) (0.304,7.189)rectangle++(0.04,0.04)
  (1.95,2.234)rectangle++(0.04,0.04) (1.745,3.417)rectangle++(0.04,0.04) (0.678,1.704)rectangle++(0.04,0.04) (0.642,3.017)rectangle++(0.04,0.04) (0.664,4.135)rectangle++(0.04,0.04)
  (0.728,5.085)rectangle++(0.04,0.04) (1.297,1.84)rectangle++(0.04,0.04) (1.174,3.108)rectangle++(0.04,0.04) (-0.594,2.543)rectangle++(0.04,0.04) (0.784,0.886)rectangle++(0.04,0.04)
  (0.699,2.318)rectangle++(0.04,0.04) (0.527,1.168)rectangle++(0.04,0.04) (-0.22,1.967)rectangle++(0.04,0.04) (-0.11,3.277)rectangle++(0.04,0.04) (0.035,4.386)rectangle++(0.04,0.04)
  (0.203,5.323)rectangle++(0.04,0.04) (1.443,1.76)rectangle++(0.04,0.04) (1.294,3.035)rectangle++(0.04,0.04) (1.219,4.124)rectangle++(0.04,0.04) (1.199,5.053)rectangle++(0.04,0.04)
  (1.219,5.843)rectangle++(0.04,0.04) (1.268,6.514)rectangle++(0.04,0.04) (1.616,2.323)rectangle++(0.04,0.04) (1.465,3.506)rectangle++(0.04,0.04) (0.745,1.652)rectangle++(0.04,0.04)
  (0.697,2.971)rectangle++(0.04,0.04) (-0.654,2.403)rectangle++(0.04,0.04) (-0.462,3.665)rectangle++(0.04,0.04) (-0.248,4.73)rectangle++(0.04,0.04) (-0.024,5.626)rectangle++(0.04,0.04)
  (0.202,6.38)rectangle++(0.04,0.04) (0.425,7.011)rectangle++(0.04,0.04) (0.639,7.538)rectangle++(0.04,0.04) (0.843,7.978)rectangle++(0.04,0.04) (2.09,2.564)rectangle++(0.04,0.04)
  (1.877,3.692)rectangle++(0.04,0.04) (1.741,4.659)rectangle++(0.04,0.04) (1.664,5.487)rectangle++(0.04,0.04) (1.632,6.194)rectangle++(0.04,0.04) (1.632,6.795)rectangle++(0.04,0.04)
  (1.657,7.307)rectangle++(0.04,0.04) (1.78,2.615)rectangle++(0.04,0.04) (1.615,3.747)rectangle++(0.04,0.04) (0.79,1.749)rectangle++(0.04,0.04) (0.739,3.052)rectangle++(0.04,0.04)
  (0.748,4.16)rectangle++(0.04,0.04) (0.8,5.103)rectangle++(0.04,0.04) (1.291,1.863)rectangle++(0.04,0.04) (1.17,3.128)rectangle++(0.04,0.04) (1.117,4.208)rectangle++(0.04,0.04)
  (1.116,5.128)rectangle++(0.04,0.04) (1.151,5.911)rectangle++(0.04,0.04) (1.213,6.574)rectangle++(0.04,0.04) (1.292,7.136)rectangle++(0.04,0.04) (1.381,7.61)rectangle++(0.04,0.04)
  (1.476,8.009)rectangle++(0.04,0.04) (1.573,8.345)rectangle++(0.04,0.04) (1.669,8.627)rectangle++(0.04,0.04) (1.761,8.862)rectangle++(0.04,0.04) (1.849,9.059)rectangle++(0.04,0.04)
  (1.932,9.222)rectangle++(0.04,0.04) (2.009,9.358)rectangle++(0.04,0.04) (2.08,9.47)rectangle++(0.04,0.04) (2.144,9.562)rectangle++(0.04,0.04) (2.338,3.282)rectangle++(0.04,0.04)
  (0.551,1.826)rectangle++(0.04,0.04) (0.539,3.126)rectangle++(0.04,0.04) (0.581,4.232)rectangle++(0.04,0.04) (0.661,5.17)rectangle++(0.04,0.04) (1.331,1.843)rectangle++(0.04,0.04)
  (-0.234,2.301)rectangle++(0.04,0.04) (-0.109,3.561)rectangle++(0.04,0.04) (0.047,4.627)rectangle++(0.04,0.04) (0.223,5.528)rectangle++(0.04,0.04) (0.409,6.286)rectangle++(0.04,0.04)
  (0.597,6.923)rectangle++(0.04,0.04) (0.782,7.457)rectangle++(0.04,0.04) (-1.804,3.409)rectangle++(0.04,0.04) (-1.399,4.566)rectangle++(0.04,0.04) (-1.009,5.533)rectangle++(0.04,0.04)
  (-0.638,6.34)rectangle++(0.04,0.04) (-0.291,7.011)rectangle++(0.04,0.04) (0.031,7.567)rectangle++(0.04,0.04) (0.327,8.027)rectangle++(0.04,0.04) (2.181,2.441)rectangle++(0.04,0.04)
  (1.949,3.584)rectangle++(0.04,0.04) (1.798,4.565)rectangle++(0.04,0.04) (1.709,5.404)rectangle++(0.04,0.04) (1.666,6.122)rectangle++(0.04,0.04) (-1.279,3.319)rectangle++(0.04,0.04)
  (-0.957,4.469)rectangle++(0.04,0.04) (1.377,1.254)rectangle++(0.04,0.04) (1.219,2.607)rectangle++(0.04,0.04) (1.138,3.763)rectangle++(0.04,0.04) (1.116,4.749)rectangle++(0.04,0.04)
  (1.136,5.589)rectangle++(0.04,0.04) (1.187,6.301)rectangle++(0.04,0.04) (1.259,6.905)rectangle++(0.04,0.04) (1.344,7.415)rectangle++(0.04,0.04) (1.857,2.559)rectangle++(0.04,0.04)
  (1.679,3.697)rectangle++(0.04,0.04) (1.573,4.672)rectangle++(0.04,0.04) (1.521,5.504)rectangle++(0.04,0.04) (1.511,6.214)rectangle++(0.04,0.04) (1.531,6.818)rectangle++(0.04,0.04)
  (1.572,7.33)rectangle++(0.04,0.04) (1.627,7.764)rectangle++(0.04,0.04) (1.691,8.13)rectangle++(0.04,0.04) (1.761,8.439)rectangle++(0.04,0.04) (-1.863,3.851)rectangle++(0.04,0.04)
  (-1.432,4.944)rectangle++(0.04,0.04) (-1.022,5.856)rectangle++(0.04,0.04) (-0.636,6.614)rectangle++(0.04,0.04) (-0.279,7.244)rectangle++(0.04,0.04) (0.051,7.765)rectangle++(0.04,0.04)
  (0.352,8.194)rectangle++(0.04,0.04) (0.624,8.547)rectangle++(0.04,0.04) (0.87,8.836)rectangle++(0.04,0.04) (1.091,9.072)rectangle++(0.04,0.04) (1.288,9.264)rectangle++(0.04,0.04)
  (1.463,9.419)rectangle++(0.04,0.04) (-2.177,3.998)rectangle++(0.04,0.04) (-1.693,5.081)rectangle++(0.04,0.04) (-1.238,5.983)rectangle++(0.04,0.04) (-0.815,6.731)rectangle++(0.04,0.04)
  (-0.426,7.35)rectangle++(0.04,0.04) (-0.07,7.861)rectangle++(0.04,0.04) (0.253,8.281)rectangle++(0.04,0.04) (0.544,8.625)rectangle++(0.04,0.04) (0.805,8.906)rectangle++(0.04,0.04)
  (1.038,9.134)rectangle++(0.04,0.04) (1.246,9.318)rectangle++(0.04,0.04) (1.429,9.467)rectangle++(0.04,0.04) (1.591,9.586)rectangle++(0.04,0.04) (2.428,3.144)rectangle++(0.04,0.04)
  (2.187,4.172)rectangle++(0.04,0.04) (2.024,5.055)rectangle++(0.04,0.04) (1.92,5.812)rectangle++(0.04,0.04) (1.863,6.459)rectangle++(0.04,0.04) (1.839,7.012)rectangle++(0.04,0.04)
  (1.842,7.483)rectangle++(0.04,0.04) (1.863,7.883)rectangle++(0.04,0.04) (1.896,8.222)rectangle++(0.04,0.04) (1.939,8.509)rectangle++(0.04,0.04) (2.074,2.976)rectangle++(0.04,0.04)
  (1.88,4.043)rectangle++(0.04,0.04) (1.758,4.958)rectangle++(0.04,0.04) (1.69,5.74)rectangle++(0.04,0.04) (1.664,6.407)rectangle++(0.04,0.04) (1.668,6.976)rectangle++(0.04,0.04)
  (1.695,7.459)rectangle++(0.04,0.04) (1.737,7.869)rectangle++(0.04,0.04) (1.789,8.215)rectangle++(0.04,0.04) (1.847,8.507)rectangle++(0.04,0.04) (1.908,8.754)rectangle++(0.04,0.04)
  (1.97,8.96)rectangle++(0.04,0.04) (2.031,9.134)rectangle++(0.04,0.04) (2.089,9.279)rectangle++(0.04,0.04) (2.145,9.4)rectangle++(0.04,0.04) (2.293,3.244)rectangle++(0.04,0.04)
  (2.076,4.261)rectangle++(0.04,0.04) (1.933,5.135)rectangle++(0.04,0.04) (-0.97,3.163)rectangle++(0.04,0.04) (-0.7,4.324)rectangle++(0.04,0.04) (1.298,1.286)rectangle++(0.04,0.04)
  (1.153,2.637)rectangle++(0.04,0.04) (1.083,3.792)rectangle++(0.04,0.04) (1.07,4.776)rectangle++(0.04,0.04) (1.098,5.613)rectangle++(0.04,0.04) (1.156,6.323)rectangle++(0.04,0.04)
  (1.233,6.925)rectangle++(0.04,0.04) (1.323,7.433)rectangle++(0.04,0.04) (1.42,7.861)rectangle++(0.04,0.04) (1.519,8.221)rectangle++(0.04,0.04) (1.618,8.524)rectangle++(0.04,0.04)
  (1.714,8.777)rectangle++(0.04,0.04) (1.806,8.988)rectangle++(0.04,0.04) (1.892,9.164)rectangle++(0.04,0.04) (1.973,9.31)rectangle++(0.04,0.04) (2.047,9.43)rectangle++(0.04,0.04)
  (2.115,9.53)rectangle++(0.04,0.04) (2.177,9.612)rectangle++(0.04,0.04) (-2.085,4.204)rectangle++(0.04,0.04) (-1.606,5.253)rectangle++(0.04,0.04) (-1.157,6.126)rectangle++(0.04,0.04)
  (-0.741,6.849)rectangle++(0.04,0.04) (-1.95,2.925)rectangle++(0.04,0.04) (-1.543,4.161)rectangle++(0.04,0.04) (-1.147,5.195)rectangle++(0.04,0.04) (-0.77,6.058)rectangle++(0.04,0.04)
  (-0.414,6.776)rectangle++(0.04,0.04) (-0.083,7.372)rectangle++(0.04,0.04) (0.222,7.866)rectangle++(0.04,0.04) (-2.022,3.371)rectangle++(0.04,0.04) (-1.586,4.542)rectangle++(0.04,0.04)
  (-1.169,5.52)rectangle++(0.04,0.04) (-0.775,6.335)rectangle++(0.04,0.04) (-0.407,7.012)rectangle++(0.04,0.04) (-0.068,7.573)rectangle++(0.04,0.04) (2.113,2.23)rectangle++(0.04,0.04)
  (1.883,3.407)rectangle++(0.04,0.04) (1.735,4.417)rectangle++(0.04,0.04) (-0.823,2.96)rectangle++(0.04,0.04) (-0.583,4.145)rectangle++(0.04,0.04) (-0.332,5.143)rectangle++(0.04,0.04)
  (1.472,1.578)rectangle++(0.04,0.04) (-0.137,2.275)rectangle++(0.04,0.04) (-0.028,3.535)rectangle++(0.04,0.04) (0.116,4.602)rectangle++(0.04,0.04) (0.28,5.503)rectangle++(0.04,0.04)
  (1.482,1.824)rectangle++(0.04,0.04) (0.271,1.253)rectangle++(0.04,0.04) (0.278,2.65)rectangle++(0.04,0.04) (0.34,3.838)rectangle++(0.04,0.04) (0.441,4.845)rectangle++(0.04,0.04)
  (0.566,5.697)rectangle++(0.04,0.04) (0.707,6.416)rectangle++(0.04,0.04) (1.673,2.154)rectangle++(0.04,0.04) (1.506,3.36)rectangle++(0.04,0.04) (1.412,4.392)rectangle++(0.04,0.04)
  (1.374,5.273)rectangle++(0.04,0.04) (1.377,6.023)rectangle++(0.04,0.04) (1.409,6.661)rectangle++(0.04,0.04) (1.462,7.201)rectangle++(0.04,0.04) (1.528,7.659)rectangle++(0.04,0.04)
  (1.603,8.045)rectangle++(0.04,0.04) (1.682,8.371)rectangle++(0.04,0.04) (1.763,8.644)rectangle++(0.04,0.04) (1.842,8.873)rectangle++(0.04,0.04) (1.918,9.065)rectangle++(0.04,0.04)
  (1.991,9.224)rectangle++(0.04,0.04) (2.059,9.357)rectangle++(0.04,0.04) (2.122,9.467)rectangle++(0.04,0.04) (2.18,9.559)rectangle++(0.04,0.04) (2.233,9.634)rectangle++(0.04,0.04)
  (2.282,9.696)rectangle++(0.04,0.04) (2.325,9.746)rectangle++(0.04,0.04) (2.364,9.787)rectangle++(0.04,0.04) (2.399,9.821)rectangle++(0.04,0.04) (2.429,9.848)rectangle++(0.04,0.04)
  (2.457,9.87)rectangle++(0.04,0.04) (2.481,9.887)rectangle++(0.04,0.04) (2.502,9.901)rectangle++(0.04,0.04) (2.521,9.912)rectangle++(0.04,0.04) (2.537,9.921)rectangle++(0.04,0.04)
  (2.551,9.927)rectangle++(0.04,0.04) (2.38,3.476)rectangle++(0.04,0.04) (2.16,4.455)rectangle++(0.04,0.04) (0.906,2.061)rectangle++(0.04,0.04) (0.85,3.312)rectangle++(0.04,0.04)
  (0.782,1.446)rectangle++(0.04,0.04) (0.721,2.794)rectangle++(0.04,0.04) (0.722,3.942)rectangle++(0.04,0.04) (0.769,4.918)rectangle++(0.04,0.04) (0.848,5.746)rectangle++(0.04,0.04)
  (0.949,6.446)rectangle++(0.04,0.04) (1.062,7.038)rectangle++(0.04,0.04) (1.182,7.536)rectangle++(0.04,0.04) (1.915,2.546)rectangle++(0.04,0.04) (1.728,3.684)rectangle++(0.04,0.04)
  (1.614,4.658)rectangle++(0.04,0.04) (1.556,5.491)rectangle++(0.04,0.04) (1.54,6.201)rectangle++(0.04,0.04) (1.555,6.806)rectangle++(0.04,0.04) (1.592,7.319)rectangle++(0.04,0.04)
  (1.643,7.754)rectangle++(0.04,0.04) (-1.708,3.673)rectangle++(0.04,0.04) (-1.307,4.786)rectangle++(0.04,0.04) (-0.922,5.717)rectangle++(0.04,0.04) (-0.557,6.492)rectangle++(0.04,0.04)
  (-0.216,7.137)rectangle++(0.04,0.04) (0.099,7.671)rectangle++(0.04,0.04) (0.389,8.113)rectangle++(0.04,0.04) (0.653,8.477)rectangle++(0.04,0.04) (0.892,8.775)rectangle++(0.04,0.04)
  (-2.124,3.725)rectangle++(0.04,0.04) (-1.659,4.847)rectangle++(0.04,0.04) (1.589,1.162)rectangle++(0.04,0.04) (1.395,2.52)rectangle++(0.04,0.04) (1.284,3.683)rectangle++(0.04,0.04)
  (1.237,4.675)rectangle++(0.04,0.04) (1.236,5.521)rectangle++(0.04,0.04) (1.269,6.239)rectangle++(0.04,0.04) (1.326,6.849)rectangle++(0.04,0.04) (1.399,7.365)rectangle++(0.04,0.04)
  (1.481,7.8)rectangle++(0.04,0.04) (1.569,8.167)rectangle++(0.04,0.04) (1.658,8.475)rectangle++(0.04,0.04) (1.746,8.734)rectangle++(0.04,0.04) (1.832,8.95)rectangle++(0.04,0.04)
  (1.913,9.131)rectangle++(0.04,0.04) (1.989,9.281)rectangle++(0.04,0.04) (2.059,9.405)rectangle++(0.04,0.04) (2.125,9.508)rectangle++(0.04,0.04) (2.184,9.593)rectangle++(0.04,0.04)
  (2.238,9.663)rectangle++(0.04,0.04) (2.287,9.72)rectangle++(0.04,0.04) (2.33,9.767)rectangle++(0.04,0.04) (2.369,9.805)rectangle++(0.04,0.04) (2.404,9.836)rectangle++(0.04,0.04)
  (2.434,9.86)rectangle++(0.04,0.04) (2.461,9.88)rectangle++(0.04,0.04) (2.485,9.896)rectangle++(0.04,0.04) (2.506,9.908)rectangle++(0.04,0.04) (2.524,9.918)rectangle++(0.04,0.04)
  (2.54,9.926)rectangle++(0.04,0.04) (2.554,9.931)rectangle++(0.04,0.04) (2.38,3.478)rectangle++(0.04,0.04) (2.16,4.457)rectangle++(0.04,0.04) (0.907,2.061)rectangle++(0.04,0.04)
  (0.851,3.312)rectangle++(0.04,0.04) (-0.712,2.513)rectangle++(0.04,0.04) (-0.507,3.761)rectangle++(0.04,0.04) (-0.283,4.813)rectangle++(0.04,0.04) (-0.02,0.753)rectangle++(0.04,0.04)
  (-0.221,1.75)rectangle++(0.04,0.04) (-0.12,3.093)rectangle++(0.04,0.04) (0.019,4.23)rectangle++(0.04,0.04) (1.164,1.45)rectangle++(0.04,0.04) (1.045,2.782)rectangle++(0.04,0.04)
  (0.998,3.919)rectangle++(0.04,0.04) (1.002,4.888)rectangle++(0.04,0.04) (1.045,5.711)rectangle++(0.04,0.04) (1.115,6.408)rectangle++(0.04,0.04) (1.202,6.999)rectangle++(0.04,0.04)
  (1.299,7.497)rectangle++(0.04,0.04) (1.402,7.917)rectangle++(0.04,0.04) (-1.799,3.653)rectangle++(0.04,0.04) (-1.385,4.773)rectangle++(0.04,0.04) (-0.989,5.709)rectangle++(0.04,0.04)
  (-0.614,6.488)rectangle++(0.04,0.04) (-0.265,7.136)rectangle++(0.04,0.04) (0.058,7.672)rectangle++(0.04,0.04) (0.354,8.115)rectangle++(0.04,0.04) (0.623,8.48)rectangle++(0.04,0.04)
  (0.867,8.779)rectangle++(0.04,0.04) (1.086,9.024)rectangle++(0.04,0.04) (1.282,9.223)rectangle++(0.04,0.04) (1.456,9.385)rectangle++(0.04,0.04) (1.611,9.515)rectangle++(0.04,0.04)
  (1.748,9.619)rectangle++(0.04,0.04) (1.868,9.703)rectangle++(0.04,0.04) (1.974,9.769)rectangle++(0.04,0.04) (2.066,9.821)rectangle++(0.04,0.04) (2.147,9.861)rectangle++(0.04,0.04)
  (2.217,9.892)rectangle++(0.04,0.04) (2.278,9.916)rectangle++(0.04,0.04) (2.417,3.402)rectangle++(0.04,0.04) (2.189,4.391)rectangle++(0.04,0.04) (2.034,5.241)rectangle++(0.04,0.04)
  (1.145,2.217)rectangle++(0.04,0.04) (1.06,3.435)rectangle++(0.04,0.04) (1.036,4.473)rectangle++(0.04,0.04) (1.08,1.773)rectangle++(0.04,0.04) (0.986,3.06)rectangle++(0.04,0.04)
  (0.959,4.158)rectangle++(0.04,0.04) (0.979,5.092)rectangle++(0.04,0.04) (1.034,5.885)rectangle++(0.04,0.04) (1.112,6.557)rectangle++(0.04,0.04) (1.205,7.125)rectangle++(0.04,0.04)
  (1.307,7.605)rectangle++(0.04,0.04) (1.413,8.008)rectangle++(0.04,0.04) (1.519,8.346)rectangle++(0.04,0.04) (1.623,8.63)rectangle++(0.04,0.04) (1.723,8.867)rectangle++(0.04,0.04)
  (1.817,9.064)rectangle++(0.04,0.04) (1.904,9.228)rectangle++(0.04,0.04) (1.986,9.364)rectangle++(0.04,0.04) (2.06,9.476)rectangle++(0.04,0.04) (2.128,9.568)rectangle++(0.04,0.04)
  (2.189,9.644)rectangle++(0.04,0.04) (-2.091,4.214)rectangle++(0.04,0.04) (-1.611,5.262)rectangle++(0.04,0.04) (-1.161,6.133)rectangle++(0.04,0.04) (-0.744,6.856)rectangle++(0.04,0.04)
  (-0.36,7.453)rectangle++(0.04,0.04) (-0.01,7.946)rectangle++(0.04,0.04) (0.307,8.351)rectangle++(0.04,0.04) (0.593,8.682)rectangle++(0.04,0.04) (0.849,8.952)rectangle++(0.04,0.04)
  (1.078,9.172)rectangle++(0.04,0.04) (1.281,9.349)rectangle++(0.04,0.04) (1.46,9.492)rectangle++(0.04,0.04) (1.619,9.606)rectangle++(0.04,0.04) (2.429,3.156)rectangle++(0.04,0.04)
  (2.189,4.182)rectangle++(0.04,0.04) (2.026,5.063)rectangle++(0.04,0.04) (1.922,5.819)rectangle++(0.04,0.04) (1.864,6.465)rectangle++(0.04,0.04) (1.841,7.017)rectangle++(0.04,0.04)
  (1.843,7.487)rectangle++(0.04,0.04) (1.864,7.887)rectangle++(0.04,0.04) (-1.699,3.753)rectangle++(0.04,0.04) (-1.296,4.854)rectangle++(0.04,0.04) (-0.91,5.774)rectangle++(0.04,0.04)
  (-0.545,6.54)rectangle++(0.04,0.04) (-0.203,7.177)rectangle++(0.04,0.04) (2.023,2.1)rectangle++(0.04,0.04) (-0.02,0.319)rectangle++(0.04,0.04) (-0.006,1.868)rectangle++(0.04,0.04)
  (0.067,3.184)rectangle++(0.04,0.04) (0.182,4.3)rectangle++(0.04,0.04) (0.325,5.244)rectangle++(0.04,0.04) (0.484,6.041)rectangle++(0.04,0.04) (0.65,6.712)rectangle++(0.04,0.04)
  (0.819,7.275)rectangle++(0.04,0.04) (0.985,7.747)rectangle++(0.04,0.04) (1.145,8.142)rectangle++(0.04,0.04) (1.297,8.471)rectangle++(0.04,0.04) (1.439,8.745)rectangle++(0.04,0.04)
  (-2.007,3.844)rectangle++(0.04,0.04) (-1.554,4.944)rectangle++(0.04,0.04) (-1.126,5.861)rectangle++(0.04,0.04) (-0.725,6.623)rectangle++(0.04,0.04) (-0.353,7.254)rectangle++(0.04,0.04)
  (-0.012,7.777)rectangle++(0.04,0.04) (0.298,8.207)rectangle++(0.04,0.04) (0.58,8.56)rectangle++(0.04,0.04) (-2.131,3.606)rectangle++(0.04,0.04) (-1.669,4.746)rectangle++(0.04,0.04)
  (-1.231,5.697)rectangle++(0.04,0.04) (-0.821,6.488)rectangle++(0.04,0.04) (-0.44,7.144)rectangle++(0.04,0.04) (-0.091,7.686)rectangle++(0.04,0.04) (2.148,2.251)rectangle++(0.04,0.04)
  (1.914,3.424)rectangle++(0.04,0.04) (1.762,4.43)rectangle++(0.04,0.04) (1.672,5.291)rectangle++(0.04,0.04) (1.631,6.027)rectangle++(0.04,0.04) (-0.02,0.947)rectangle++(0.04,0.04)
  (0.019,2.402)rectangle++(0.04,0.04) (0.11,3.637)rectangle++(0.04,0.04) (0.237,4.684)rectangle++(0.04,0.04) (0.386,5.568)rectangle++(0.04,0.04) (0.549,6.313)rectangle++(0.04,0.04)
  (1.668,2.088)rectangle++(0.04,0.04) (1.499,3.304)rectangle++(0.04,0.04) (1.404,4.345)rectangle++(0.04,0.04) (1.365,5.233)rectangle++(0.04,0.04) (1.368,5.99)rectangle++(0.04,0.04)
  (1.4,6.633)rectangle++(0.04,0.04) (-1.466,3.37)rectangle++(0.04,0.04) (-1.113,4.519)rectangle++(0.04,0.04) (-0.768,5.482)rectangle++(0.04,0.04) (-0.435,6.287)rectangle++(0.04,0.04)
  (-1.743,2.872)rectangle++(0.04,0.04) (-1.369,4.107)rectangle++(0.04,0.04) (-1.001,5.142)rectangle++(0.04,0.04) (-0.648,6.007)rectangle++(0.04,0.04) (-0.312,6.728)rectangle++(0.04,0.04)
  (0.001,7.327)rectangle++(0.04,0.04) (0.292,7.825)rectangle++(0.04,0.04) (0.559,8.235)rectangle++(0.04,0.04) (0.802,8.574)rectangle++(0.04,0.04) (-2.09,3.66)rectangle++(0.04,0.04)
  (-1.632,4.791)rectangle++(0.04,0.04) (-1.593,2.268)rectangle++(0.04,0.04) (-1.266,3.587)rectangle++(0.04,0.04) (-0.935,4.696)rectangle++(0.04,0.04) (-0.609,5.625)rectangle++(0.04,0.04)
  (-0.295,6.402)rectangle++(0.04,0.04) (0.003,7.05)rectangle++(0.04,0.04) (0.283,7.588)rectangle++(0.04,0.04) (0.542,8.035)rectangle++(0.04,0.04) (0.78,8.404)rectangle++(0.04,0.04)
  (0.997,8.709)rectangle++(0.04,0.04) (1.193,8.959)rectangle++(0.04,0.04) (1.37,9.163)rectangle++(0.04,0.04) (1.529,9.33)rectangle++(0.04,0.04) (1.671,9.466)rectangle++(0.04,0.04)
  (1.797,9.575)rectangle++(0.04,0.04) (1.908,9.663)rectangle++(0.04,0.04) (2.402,3.245)rectangle++(0.04,0.04) (2.169,4.259)rectangle++(0.04,0.04) (-0.695,3.025)rectangle++(0.04,0.04)
  (-0.472,4.195)rectangle++(0.04,0.04) (-0.235,5.181)rectangle++(0.04,0.04) (0.005,6.009)rectangle++(0.04,0.04) (0.243,6.704)rectangle++(0.04,0.04) (0.472,7.285)rectangle++(0.04,0.04)
  (1.952,2.301)rectangle++(0.04,0.04) (1.749,3.474)rectangle++(0.04,0.04) (1.623,4.479)rectangle++(0.04,0.04) (1.557,5.339)rectangle++(0.04,0.04) (1.534,6.072)rectangle++(0.04,0.04)
  (1.545,6.696)rectangle++(0.04,0.04) (1.579,7.226)rectangle++(0.04,0.04) (1.629,7.675)rectangle++(0.04,0.04) (1.689,8.055)rectangle++(0.04,0.04) (1.756,8.375)rectangle++(0.04,0.04)
  (1.825,8.645)rectangle++(0.04,0.04) (-0.02,1.366)rectangle++(0.04,0.04) (0.035,2.758)rectangle++(0.04,0.04) (0.138,3.939)rectangle++(0.04,0.04) (0.273,4.939)rectangle++(0.04,0.04)
  (0.427,5.784)rectangle++(0.04,0.04) (0.592,6.495)rectangle++(0.04,0.04) (0.761,7.093)rectangle++(0.04,0.04) (0.928,7.595)rectangle++(0.04,0.04) (1.091,8.015)rectangle++(0.04,0.04)
  (1.246,8.365)rectangle++(0.04,0.04) (1.391,8.657)rectangle++(0.04,0.04) (1.527,8.899)rectangle++(0.04,0.04) (1.651,9.099)rectangle++(0.04,0.04) (1.765,9.264)rectangle++(0.04,0.04)
  (1.869,9.4)rectangle++(0.04,0.04) (1.962,9.512)rectangle++(0.04,0.04) (2.046,9.603)rectangle++(0.04,0.04) (2.121,9.677)rectangle++(0.04,0.04) (2.188,9.736)rectangle++(0.04,0.04)
  (2.247,9.785)rectangle++(0.04,0.04) (2.299,9.823)rectangle++(0.04,0.04) (2.388,3.385)rectangle++(0.04,0.04) (2.163,4.378)rectangle++(0.04,0.04) (2.012,5.231)rectangle++(0.04,0.04)
  (1.917,5.962)rectangle++(0.04,0.04) (1.866,6.587)rectangle++(0.04,0.04) (1.847,7.121)rectangle++(0.04,0.04) (1.853,7.575)rectangle++(0.04,0.04) (1.876,7.961)rectangle++(0.04,0.04)
  (1.911,8.288)rectangle++(0.04,0.04) (-1.794,3.852)rectangle++(0.04,0.04) (-1.373,4.942)rectangle++(0.04,0.04) (-0.972,5.852)rectangle++(0.04,0.04) (1.767,1.582)rectangle++(0.04,0.04)
  (1.563,2.87)rectangle++(0.04,0.04) (1.441,3.973)rectangle++(0.04,0.04) (1.382,4.916)rectangle++(0.04,0.04) (-0.02,0.77)rectangle++(0.04,0.04) (0.012,2.251)rectangle++(0.04,0.04)
  (0.098,3.509)rectangle++(0.04,0.04) (0.221,4.575)rectangle++(0.04,0.04) (0.369,5.476)rectangle++(0.04,0.04) (0.53,6.236)rectangle++(0.04,0.04) (0.698,6.876)rectangle++(0.04,0.04)
  (0.866,7.413)rectangle++(0.04,0.04) (1.031,7.862)rectangle++(0.04,0.04) (1.188,8.238)rectangle++(0.04,0.04) (1.337,8.551)rectangle++(0.04,0.04) (1.477,8.811)rectangle++(0.04,0.04)
  (1.605,9.027)rectangle++(0.04,0.04) (1.723,9.205)rectangle++(0.04,0.04) (1.831,9.351)rectangle++(0.04,0.04) (1.928,9.471)rectangle++(0.04,0.04) (2.016,9.57)rectangle++(0.04,0.04)
  (2.094,9.65)rectangle++(0.04,0.04) (2.164,9.715)rectangle++(0.04,0.04) (2.225,9.767)rectangle++(0.04,0.04) (2.28,9.809)rectangle++(0.04,0.04) (2.328,9.843)rectangle++(0.04,0.04)
  (2.371,9.87)rectangle++(0.04,0.04) (2.408,9.891)rectangle++(0.04,0.04) (2.44,9.907)rectangle++(0.04,0.04) (-2.109,4.33)rectangle++(0.04,0.04) (-1.569,2.056)rectangle++(0.04,0.04)
  (-1.253,3.407)rectangle++(0.04,0.04) (-0.931,4.542)rectangle++(0.04,0.04) (-0.612,5.494)rectangle++(0.04,0.04) (-0.303,6.291)rectangle++(0.04,0.04) (-0.008,6.956)rectangle++(0.04,0.04)
  (0.269,7.509)rectangle++(0.04,0.04) (0.527,7.968)rectangle++(0.04,0.04) (0.765,8.348)rectangle++(0.04,0.04) (0.982,8.661)rectangle++(0.04,0.04) (1.179,8.919)rectangle++(0.04,0.04)
  (1.356,9.13)rectangle++(0.04,0.04) (1.516,9.303)rectangle++(0.04,0.04) (1.658,9.443)rectangle++(0.04,0.04) (1.785,9.556)rectangle++(0.04,0.04) (1.897,9.648)rectangle++(0.04,0.04)
  (1.997,9.721)rectangle++(0.04,0.04) (2.084,9.779)rectangle++(0.04,0.04) (2.16,9.825)rectangle++(0.04,0.04) (2.227,9.861)rectangle++(0.04,0.04) (2.285,9.889)rectangle++(0.04,0.04)
  (2.336,9.91)rectangle++(0.04,0.04) (2.38,9.927)rectangle++(0.04,0.04) (2.417,9.939)rectangle++(0.04,0.04) (2.45,9.947)rectangle++(0.04,0.04) (2.478,9.953)rectangle++(0.04,0.04)
  (2.503,9.958)rectangle++(0.04,0.04) (2.523,9.96)rectangle++(0.04,0.04) (-2.106,4.361)rectangle++(0.04,0.04) (-1.618,5.387)rectangle++(0.04,0.04) (-1.162,6.24)rectangle++(0.04,0.04)
  (-0.74,6.947)rectangle++(0.04,0.04) (-0.354,7.53)rectangle++(0.04,0.04) (2.144,2.145)rectangle++(0.04,0.04) (1.906,3.334)rectangle++(0.04,0.04) (1.751,4.354)rectangle++(0.04,0.04)
  (1.661,5.227)rectangle++(0.04,0.04) (1.618,5.973)rectangle++(0.04,0.04) (1.612,6.608)rectangle++(0.04,0.04) (1.633,7.149)rectangle++(0.04,0.04) (-1.553,3.537)rectangle++(0.04,0.04)
  (-1.181,4.665)rectangle++(0.04,0.04) (-0.819,5.609)rectangle++(0.04,0.04) (-0.474,6.396)rectangle++(0.04,0.04) (-0.15,7.052)rectangle++(0.04,0.04) (0.153,7.596)rectangle++(0.04,0.04)
  (0.431,8.047)rectangle++(0.04,0.04) (0.686,8.419)rectangle++(0.04,0.04) (0.918,8.725)rectangle++(0.04,0.04) (1.127,8.976)rectangle++(0.04,0.04) (1.315,9.18)rectangle++(0.04,0.04)
  (1.483,9.347)rectangle++(0.04,0.04) (1.632,9.482)rectangle++(0.04,0.04) (2.393,3.13)rectangle++(0.04,0.04) (2.157,4.161)rectangle++(0.04,0.04) (0.824,1.989)rectangle++(0.04,0.04)
  (0.778,3.254)rectangle++(0.04,0.04) (0.789,4.331)rectangle++(0.04,0.04) (0.842,5.246)rectangle++(0.04,0.04) (0.923,6.022)rectangle++(0.04,0.04) (1.023,6.678)rectangle++(0.04,0.04)
  (1.135,7.231)rectangle++(0.04,0.04) (1.252,7.697)rectangle++(0.04,0.04) (1.37,8.089)rectangle++(0.04,0.04) (1.486,8.417)rectangle++(0.04,0.04) (-1.913,3.782)rectangle++(0.04,0.04)
  (-1.477,4.888)rectangle++(0.04,0.04) (-1.062,5.81)rectangle++(0.04,0.04) (-0.672,6.577)rectangle++(0.04,0.04) (-0.311,7.214)rectangle++(0.04,0.04) (0.022,7.74)rectangle++(0.04,0.04)
  (0.326,8.174)rectangle++(0.04,0.04) (0.602,8.531)rectangle++(0.04,0.04) (0.851,8.824)rectangle++(0.04,0.04) (1.074,9.062)rectangle++(0.04,0.04) (1.273,9.256)rectangle++(0.04,0.04)
  (1.45,9.413)rectangle++(0.04,0.04) (1.607,9.539)rectangle++(0.04,0.04) (1.745,9.64)rectangle++(0.04,0.04) (1.867,9.721)rectangle++(0.04,0.04) (1.974,9.784)rectangle++(0.04,0.04)
  (2.067,9.834)rectangle++(0.04,0.04) (-2.165,4.228)rectangle++(0.04,0.04) (-1.673,5.276)rectangle++(0.04,0.04) (-1.213,6.148)rectangle++(0.04,0.04) (-0.788,6.871)rectangle++(0.04,0.04)
  (-0.397,7.468)rectangle++(0.04,0.04) (-0.041,7.96)rectangle++(0.04,0.04) (0.282,8.363)rectangle++(0.04,0.04) (0.572,8.694)rectangle++(0.04,0.04) (2.331,2.665)rectangle++(0.04,0.04)
  (2.086,3.768)rectangle++(0.04,0.04) (1.922,4.716)rectangle++(0.04,0.04) (1.82,5.528)rectangle++(0.04,0.04) (1.766,6.222)rectangle++(0.04,0.04) (1.748,6.814)rectangle++(0.04,0.04)
  (1.628,2.52)rectangle++(0.04,0.04) (1.483,3.673)rectangle++(0.04,0.04) (1.405,4.659)rectangle++(0.04,0.04) (1.378,5.5)rectangle++(0.04,0.04) (1.39,6.216)rectangle++(0.04,0.04)
  (1.428,6.824)rectangle++(0.04,0.04) (1.484,7.34)rectangle++(0.04,0.04) (1.553,7.776)rectangle++(0.04,0.04) (1.629,8.143)rectangle++(0.04,0.04) (-1.813,3.755)rectangle++(0.04,0.04)
  (-1.393,4.861)rectangle++(0.04,0.04) (-0.992,5.783)rectangle++(0.04,0.04) (-0.614,6.552)rectangle++(0.04,0.04) (-0.262,7.19)rectangle++(0.04,0.04) (-1.943,3.111)rectangle++(0.04,0.04)
  (-1.529,4.318)rectangle++(0.04,0.04) (1.421,1.069)rectangle++(0.04,0.04) (1.248,2.448)rectangle++(0.04,0.04) (1.157,3.627)rectangle++(0.04,0.04) (1.126,4.633)rectangle++(0.04,0.04)
  (1.14,5.489)rectangle++(0.04,0.04) (1.187,6.216)rectangle++(0.04,0.04) (1.255,6.833)rectangle++(0.04,0.04) (1.338,7.354)rectangle++(0.04,0.04) (1.429,7.793)rectangle++(0.04,0.04)
  (1.524,8.163)rectangle++(0.04,0.04) (1.62,8.474)rectangle++(0.04,0.04) (1.714,8.734)rectangle++(0.04,0.04) (1.804,8.952)rectangle++(0.04,0.04) (1.889,9.133)rectangle++(0.04,0.04)
  (1.969,9.284)rectangle++(0.04,0.04) (2.043,9.409)rectangle++(0.04,0.04) (2.11,9.512)rectangle++(0.04,0.04) (2.172,9.597)rectangle++(0.04,0.04) (2.228,9.667)rectangle++(0.04,0.04)
  (2.355,3.329)rectangle++(0.04,0.04) (2.133,4.332)rectangle++(0.04,0.04) (1.984,5.193)rectangle++(0.04,0.04) (1.892,5.931)rectangle++(0.04,0.04) (1.843,6.562)rectangle++(0.04,0.04)
  (1.827,7.1)rectangle++(0.04,0.04) (1.835,7.558)rectangle++(0.04,0.04) (1.86,7.947)rectangle++(0.04,0.04) (1.896,8.277)rectangle++(0.04,0.04) (1.941,8.556)rectangle++(0.04,0.04)
  (1.99,8.791)rectangle++(0.04,0.04) (2.041,8.989)rectangle++(0.04,0.04) (2.092,9.155)rectangle++(0.04,0.04) (2.142,9.294)rectangle++(0.04,0.04) (2.19,9.411)rectangle++(0.04,0.04)
  (2.236,9.508)rectangle++(0.04,0.04) (2.279,9.588)rectangle++(0.04,0.04) (2.318,9.655)rectangle++(0.04,0.04) (2.355,9.71)rectangle++(0.04,0.04) (2.388,9.756)rectangle++(0.04,0.04)
  (2.417,9.793)rectangle++(0.04,0.04) (2.444,9.824)rectangle++(0.04,0.04) (2.468,9.849)rectangle++(0.04,0.04) (2.37,3.435)rectangle++(0.04,0.04) (2.15,4.422)rectangle++(0.04,0.04)
  (-0.741,3.056)rectangle++(0.04,0.04) (-0.51,4.224)rectangle++(0.04,0.04) (-1.221,2.401)rectangle++(0.04,0.04) (-0.944,3.686)rectangle++(0.04,0.04) (-0.657,4.767)rectangle++(0.04,0.04)
  (-0.37,5.674)rectangle++(0.04,0.04) (-0.09,6.434)rectangle++(0.04,0.04) (0.179,7.069)rectangle++(0.04,0.04) (0.432,7.598)rectangle++(0.04,0.04) (0.669,8.037)rectangle++(0.04,0.04)
  (0.888,8.401)rectangle++(0.04,0.04) (1.089,8.701)rectangle++(0.04,0.04) (-2.066,3.754)rectangle++(0.04,0.04) (-1.608,4.869)rectangle++(0.04,0.04) (-1.174,5.8)rectangle++(0.04,0.04)
  (-0.768,6.573)rectangle++(0.04,0.04) (-0.392,7.214)rectangle++(0.04,0.04) (-0.047,7.744)rectangle++(0.04,0.04) (0.267,8.18)rectangle++(0.04,0.04) (0.552,8.539)rectangle++(0.04,0.04)
  (0.809,8.832)rectangle++(0.04,0.04) (1.039,9.071)rectangle++(0.04,0.04) (1.243,9.265)rectangle++(0.04,0.04) (1.425,9.422)rectangle++(0.04,0.04) (1.586,9.548)rectangle++(0.04,0.04)
  (1.728,9.648)rectangle++(0.04,0.04) (1.853,9.728)rectangle++(0.04,0.04) (2.429,3.246)rectangle++(0.04,0.04) (2.192,4.258)rectangle++(0.04,0.04) (0.846,2.022)rectangle++(0.04,0.04)
  (0.798,3.281)rectangle++(0.04,0.04) (0.807,4.353)rectangle++(0.04,0.04) (0.858,5.264)rectangle++(0.04,0.04) (0.938,6.036)rectangle++(0.04,0.04) (1.036,6.69)rectangle++(0.04,0.04)
  (1.146,7.241)rectangle++(0.04,0.04) (1.262,7.705)rectangle++(0.04,0.04) (1.379,8.095)rectangle++(0.04,0.04) (1.493,8.422)rectangle++(0.04,0.04) (1.604,8.695)rectangle++(0.04,0.04)
  (2.177,2.934)rectangle++(0.04,0.04) (1.965,4.003)rectangle++(0.04,0.04) (1.828,4.92)rectangle++(0.04,0.04) (1.749,5.705)rectangle++(0.04,0.04) (1.712,6.376)rectangle++(0.04,0.04)
  (1.708,6.947)rectangle++(0.04,0.04) (1.728,7.433)rectangle++(0.04,0.04) (1.764,7.845)rectangle++(0.04,0.04) (1.811,8.194)rectangle++(0.04,0.04) (1.865,8.489)rectangle++(0.04,0.04)
  (1.922,8.737)rectangle++(0.04,0.04) (2.141,3.027)rectangle++(0.04,0.04) (1.938,4.083)rectangle++(0.04,0.04) (1.809,4.989)rectangle++(0.04,0.04) (1.735,5.765)rectangle++(0.04,0.04)
  (1.703,6.427)rectangle++(0.04,0.04) (1.702,6.991)rectangle++(0.04,0.04) (1.724,7.47)rectangle++(0.04,0.04) (1.762,7.877)rectangle++(0.04,0.04) (-1.717,3.727)rectangle++(0.04,0.04)
  (-1.312,4.833)rectangle++(0.04,0.04) (-0.924,5.757)rectangle++(0.04,0.04) (-0.558,6.526)rectangle++(0.04,0.04) (-0.215,7.166)rectangle++(0.04,0.04) (-1.927,3.116)rectangle++(0.04,0.04)
  (-1.516,4.322)rectangle++(0.04,0.04) (-1.118,5.33)rectangle++(0.04,0.04) (-0.739,6.172)rectangle++(0.04,0.04) (-0.384,6.872)rectangle++(0.04,0.04) (-0.053,7.453)rectangle++(0.04,0.04)
  (-1.97,3.216)rectangle++(0.04,0.04) (-1.548,4.409)rectangle++(0.04,0.04) (-1.141,5.406)rectangle++(0.04,0.04) (-0.756,6.237)rectangle++(0.04,0.04) (-0.395,6.928)rectangle++(0.04,0.04)
  (-1.901,3.022)rectangle++(0.04,0.04) (-1.498,4.241)rectangle++(0.04,0.04) (-1.105,5.261)rectangle++(0.04,0.04) (-1.61,2.492)rectangle++(0.04,0.04) (-1.271,3.779)rectangle++(0.04,0.04)
  (-0.932,4.859)rectangle++(0.04,0.04) (-0.6,5.764)rectangle++(0.04,0.04) (-0.281,6.519)rectangle++(0.04,0.04) (0.019,7.149)rectangle++(0.04,0.04) (0.3,7.672)rectangle++(0.04,0.04)
  (-0.02,1.211)rectangle++(0.04,0.04) (-0.34,1.851)rectangle++(0.04,0.04) (-0.217,3.183)rectangle++(0.04,0.04) (-0.059,4.31)rectangle++(0.04,0.04) (0.12,5.262)rectangle++(0.04,0.04)
  (0.31,6.064)rectangle++(0.04,0.04) (0.504,6.739)rectangle++(0.04,0.04) (0.696,7.304)rectangle++(0.04,0.04) (0.881,7.777)rectangle++(0.04,0.04) (1.058,8.171)rectangle++(0.04,0.04)
  (1.224,8.499)rectangle++(0.04,0.04) (1.378,8.772)rectangle++(0.04,0.04) (1.52,8.997)rectangle++(0.04,0.04) (1.65,9.183)rectangle++(0.04,0.04) (1.767,9.336)rectangle++(0.04,0.04)
  (1.873,9.461)rectangle++(0.04,0.04) (1.969,9.563)rectangle++(0.04,0.04) (2.054,9.646)rectangle++(0.04,0.04) (-2.118,4.183)rectangle++(0.04,0.04) (-1.636,5.237)rectangle++(0.04,0.04)
  (-1.183,6.113)rectangle++(0.04,0.04) (-0.763,6.84)rectangle++(0.04,0.04) (-1.952,2.918)rectangle++(0.04,0.04) (-1.545,4.155)rectangle++(0.04,0.04) (-1.41,2.148)rectangle++(0.04,0.04)
  (-1.115,3.478)rectangle++(0.04,0.04) (-0.811,4.597)rectangle++(0.04,0.04) (1.391,1.323)rectangle++(0.04,0.04) (1.233,2.665)rectangle++(0.04,0.04) (1.153,3.812)rectangle++(0.04,0.04)
  (1.13,4.79)rectangle++(0.04,0.04) (1.15,5.623)rectangle++(0.04,0.04) (1.2,6.329)rectangle++(0.04,0.04) (-1.427,3.258)rectangle++(0.04,0.04) (-1.085,4.422)rectangle++(0.04,0.04)
  (-0.747,5.398)rectangle++(0.04,0.04) (-0.421,6.215)rectangle++(0.04,0.04) (-0.112,6.896)rectangle++(0.04,0.04) (0.179,7.462)rectangle++(0.04,0.04) (0.448,7.932)rectangle++(0.04,0.04)
  (0.696,8.32)rectangle++(0.04,0.04) (0.922,8.641)rectangle++(0.04,0.04) (1.127,8.904)rectangle++(0.04,0.04) (1.312,9.119)rectangle++(0.04,0.04) (1.478,9.295)rectangle++(0.04,0.04)
  (1.626,9.438)rectangle++(0.04,0.04) (1.757,9.553)rectangle++(0.04,0.04) (1.874,9.646)rectangle++(0.04,0.04) (1.976,9.721)rectangle++(0.04,0.04) (2.066,9.78)rectangle++(0.04,0.04)
  (-2.151,4.216)rectangle++(0.04,0.04) (-1.662,5.266)rectangle++(0.04,0.04) (-1.204,6.138)rectangle++(0.04,0.04) (-0.78,6.862)rectangle++(0.04,0.04) (-0.391,7.46)rectangle++(0.04,0.04)
  (-0.036,7.953)rectangle++(0.04,0.04) (2.215,2.329)rectangle++(0.04,0.04) (1.974,3.488)rectangle++(0.04,0.04) (1.815,4.482)rectangle++(0.04,0.04) (1.72,5.333)rectangle++(0.04,0.04)
  (1.673,6.061)rectangle++(0.04,0.04) (1.662,6.681)rectangle++(0.04,0.04) (1.678,7.208)rectangle++(0.04,0.04) (1.712,7.656)rectangle++(0.04,0.04) (1.76,8.035)rectangle++(0.04,0.04)
  (1.815,8.356)rectangle++(0.04,0.04) (1.875,8.626)rectangle++(0.04,0.04) (1.936,8.853)rectangle++(0.04,0.04) (1.998,9.044)rectangle++(0.04,0.04) (2.058,9.204)rectangle++(0.04,0.04)
  (-2.003,4.087)rectangle++(0.04,0.04) (-1.541,5.15)rectangle++(0.04,0.04) (-1.106,6.036)rectangle++(0.04,0.04) (-0.701,6.771)rectangle++(0.04,0.04) (-0.327,7.379)rectangle++(0.04,0.04)
  (0.015,7.882)rectangle++(0.04,0.04) (0.326,8.295)rectangle++(0.04,0.04) (-2.113,3.489)rectangle++(0.04,0.04) (-1.659,4.646)rectangle++(0.04,0.04) (-1.226,5.612)rectangle++(0.04,0.04)
  (-0.82,6.415)rectangle++(0.04,0.04) (-0.442,7.082)rectangle++(0.04,0.04) (-0.095,7.634)rectangle++(0.04,0.04) (0.222,8.089)rectangle++(0.04,0.04) (0.51,8.463)rectangle++(0.04,0.04)
  (2.276,2.594)rectangle++(0.04,0.04) (0.368,1.644)rectangle++(0.04,0.04) (0.388,0.92)rectangle++(0.04,0.04) (0.364,2.363)rectangle++(0.04,0.04) (0.402,3.59)rectangle++(0.04,0.04)
  (0.483,4.632)rectangle++(0.04,0.04) (0.594,5.514)rectangle++(0.04,0.04) (0.723,6.259)rectangle++(0.04,0.04) (0.863,6.888)rectangle++(0.04,0.04) (1.007,7.416)rectangle++(0.04,0.04)
  (-1.748,3.452)rectangle++(0.04,0.04) (-1.35,4.6)rectangle++(0.04,0.04) (-0.966,5.561)rectangle++(0.04,0.04) (-1.66,2.59)rectangle++(0.04,0.04) (-1.31,3.864)rectangle++(0.04,0.04)
  (-0.961,4.933)rectangle++(0.04,0.04) (-0.622,5.828)rectangle++(0.04,0.04) (-0.297,6.575)rectangle++(0.04,0.04) (0.008,7.197)rectangle++(0.04,0.04) (0.292,7.713)rectangle++(0.04,0.04)
  (0.555,8.141)rectangle++(0.04,0.04) (0.795,8.493)rectangle++(0.04,0.04) (1.013,8.784)rectangle++(0.04,0.04) (1.211,9.022)rectangle++(0.04,0.04) (1.388,9.216)rectangle++(0.04,0.04)
  (1.546,9.375)rectangle++(0.04,0.04) (2.376,3.082)rectangle++(0.04,0.04) (2.14,4.121)rectangle++(0.04,0.04) (1.982,5.013)rectangle++(0.04,0.04) (1.883,5.778)rectangle++(0.04,0.04)
  (1.829,6.432)rectangle++(0.04,0.04) (1.81,6.99)rectangle++(0.04,0.04) (1.816,7.466)rectangle++(0.04,0.04) (1.84,7.869)rectangle++(0.04,0.04) (1.877,8.212)rectangle++(0.04,0.04)
  (2,2.889)rectangle++(0.04,0.04) (1.814,3.972)rectangle++(0.04,0.04) (1.698,4.9)rectangle++(0.04,0.04) (1.637,5.693)rectangle++(0.04,0.04) (1.617,6.37)rectangle++(0.04,0.04)
  (1.627,6.946)rectangle++(0.04,0.04) (-1.502,3.491)rectangle++(0.04,0.04) (-1.139,4.624)rectangle++(0.04,0.04) (-0.785,5.572)rectangle++(0.04,0.04) (-0.447,6.364)rectangle++(0.04,0.04)
  (-0.127,7.023)rectangle++(0.04,0.04) (-1.873,3.105)rectangle++(0.04,0.04) (1.133,0.688)rectangle++(0.04,0.04) (0.988,2.136)rectangle++(0.04,0.04) (0.923,3.372)rectangle++(0.04,0.04)
  (0.917,4.426)rectangle++(0.04,0.04) (0.955,5.321)rectangle++(0.04,0.04) (1.022,6.081)rectangle++(0.04,0.04) (1.11,6.724)rectangle++(0.04,0.04) (1.21,7.267)rectangle++(0.04,0.04)
  (1.317,7.725)rectangle++(0.04,0.04) (1.426,8.11)rectangle++(0.04,0.04) (1.535,8.433)rectangle++(0.04,0.04) (1.64,8.702)rectangle++(0.04,0.04) (1.739,8.928)rectangle++(0.04,0.04)
  (1.833,9.115)rectangle++(0.04,0.04) (1.921,9.271)rectangle++(0.04,0.04) (2.29,3.154)rectangle++(0.04,0.04) (2.071,4.186)rectangle++(0.04,0.04) (1.925,5.071)rectangle++(0.04,0.04)
  (1.837,5.83)rectangle++(0.04,0.04) (1.793,6.478)rectangle++(0.04,0.04) (1.781,7.031)rectangle++(0.04,0.04) (-1.493,3.545)rectangle++(0.04,0.04) (-1.13,4.669)rectangle++(0.04,0.04)
  (-0.776,5.61)rectangle++(0.04,0.04) (1.67,1.575)rectangle++(0.04,0.04) (1.48,2.868)rectangle++(0.04,0.04) (1.371,3.975)rectangle++(0.04,0.04) (1.322,4.92)rectangle++(0.04,0.04)
  (1.318,5.725)rectangle++(0.04,0.04) (1.347,6.41)rectangle++(0.04,0.04) (1.399,6.991)rectangle++(0.04,0.04) (1.467,7.482)rectangle++(0.04,0.04) (1.544,7.898)rectangle++(0.04,0.04)
  (1.626,8.247)rectangle++(0.04,0.04) (1.71,8.541)rectangle++(0.04,0.04) (1.793,8.788)rectangle++(0.04,0.04) (1.873,8.994)rectangle++(0.04,0.04) (-1.985,3.999)rectangle++(0.04,0.04)
  (-1.53,5.074)rectangle++(0.04,0.04) (-1.099,5.971)rectangle++(0.04,0.04) (1.819,1.577)rectangle++(0.04,0.04) (1.607,2.864)rectangle++(0.04,0.04) (1.479,3.966)rectangle++(0.04,0.04)
  (0.871,1.766)rectangle++(0.04,0.04) (0.809,3.063)rectangle++(0.04,0.04) (0.808,4.167)rectangle++(0.04,0.04) (0.851,5.106)rectangle++(0.04,0.04) (0.926,5.902)rectangle++(0.04,0.04)
  (1.021,6.576)rectangle++(0.04,0.04) (1.128,7.145)rectangle++(0.04,0.04) (1.243,7.624)rectangle++(0.04,0.04) (1.359,8.027)rectangle++(0.04,0.04) (1.474,8.365)rectangle++(0.04,0.04)
  (1.585,8.647)rectangle++(0.04,0.04) (1.691,8.883)rectangle++(0.04,0.04) (1.791,9.079)rectangle++(0.04,0.04) (1.883,9.242)rectangle++(0.04,0.04) (1.968,9.376)rectangle++(0.04,0.04)
  (2.046,9.487)rectangle++(0.04,0.04) (2.116,9.579)rectangle++(0.04,0.04) (2.347,3.279)rectangle++(0.04,0.04) (0.549,1.827)rectangle++(0.04,0.04) (-0.387,2.117)rectangle++(0.04,0.04)
  (-0.246,3.411)rectangle++(0.04,0.04) (0.975,1.185)rectangle++(0.04,0.04) (0.874,2.564)rectangle++(0.04,0.04) (0.843,3.741)rectangle++(0.04,0.04) (0.904,1.547)rectangle++(0.04,0.04)
  (0.828,2.875)rectangle++(0.04,0.04) (0.816,4.007)rectangle++(0.04,0.04) (0.852,4.969)rectangle++(0.04,0.04) (0.921,5.786)rectangle++(0.04,0.04) (1.012,6.478)rectangle++(0.04,0.04)
  (1.117,7.062)rectangle++(0.04,0.04) (1.23,7.554)rectangle++(0.04,0.04) (-1.739,3.534)rectangle++(0.04,0.04) (1.233,0.826)rectangle++(0.04,0.04) (1.079,2.249)rectangle++(0.04,0.04)
  (1.005,3.465)rectangle++(0.04,0.04) (0.99,4.501)rectangle++(0.04,0.04) (1.02,5.382)rectangle++(0.04,0.04) (1.08,6.13)rectangle++(0.04,0.04) (1.161,6.764)rectangle++(0.04,0.04)
  (1.255,7.299)rectangle++(0.04,0.04) (1.357,7.75)rectangle++(0.04,0.04) (1.461,8.13)rectangle++(0.04,0.04) (1.565,8.448)rectangle++(0.04,0.04) (1.666,8.714)rectangle++(0.04,0.04)
  (1.762,8.937)rectangle++(0.04,0.04) (1.853,9.122)rectangle++(0.04,0.04) (-2.022,4.022)rectangle++(0.04,0.04) (-1.56,5.096)rectangle++(0.04,0.04) (-1.125,5.99)rectangle++(0.04,0.04)
  (-0.718,6.733)rectangle++(0.04,0.04) (-0.344,7.348)rectangle++(0.04,0.04) (0,7.856)rectangle++(0.04,0.04) (0.312,8.273)rectangle++(0.04,0.04) (0.594,8.616)rectangle++(0.04,0.04)
  (-2.143,3.621)rectangle++(0.04,0.04) (-1.679,4.76)rectangle++(0.04,0.04) (-1.239,5.709)rectangle++(0.04,0.04) (-0.827,6.499)rectangle++(0.04,0.04) (-0.445,7.153)rectangle++(0.04,0.04)
  (-0.094,7.694)rectangle++(0.04,0.04) (0.225,8.14)rectangle++(0.04,0.04) (0.515,8.506)rectangle++(0.04,0.04) (-2.13,3.579)rectangle++(0.04,0.04) (-1.669,4.723)rectangle++(0.04,0.04)
  (-1.232,5.678)rectangle++(0.04,0.04) (-0.822,6.472)rectangle++(0.04,0.04) (-0.442,7.13)rectangle++(0.04,0.04) (-0.093,7.674)rectangle++(0.04,0.04) (0.226,8.123)rectangle++(0.04,0.04)
  (0.515,8.492)rectangle++(0.04,0.04) (0.775,8.794)rectangle++(0.04,0.04) (1.008,9.04)rectangle++(0.04,0.04) (1.216,9.24)rectangle++(0.04,0.04) (1.401,9.401)rectangle++(0.04,0.04)
  (1.565,9.531)rectangle++(0.04,0.04) (2.417,3.124)rectangle++(0.04,0.04) (2.177,4.155)rectangle++(0.04,0.04) (2.014,5.041)rectangle++(0.04,0.04) (1.912,5.801)rectangle++(0.04,0.04)
  (1.32,2.319)rectangle++(0.04,0.04) (1.213,3.515)rectangle++(0.04,0.04) (1.169,4.535)rectangle++(0.04,0.04) (1.173,5.404)rectangle++(0.04,0.04) (1.211,6.143)rectangle++(0.04,0.04)
  (1.273,6.769)rectangle++(0.04,0.04) (1.35,7.299)rectangle++(0.04,0.04) (1.438,7.746)rectangle++(0.04,0.04) (1.53,8.123)rectangle++(0.04,0.04) (1.623,8.44)rectangle++(0.04,0.04)
  (1.715,8.705)rectangle++(0.04,0.04) (1.804,8.927)rectangle++(0.04,0.04) (1.888,9.112)rectangle++(0.04,0.04) (1.967,9.266)rectangle++(0.04,0.04) (2.04,9.393)rectangle++(0.04,0.04)
  (2.108,9.499)rectangle++(0.04,0.04) (2.169,9.586)rectangle++(0.04,0.04) (2.225,9.658)rectangle++(0.04,0.04) (2.276,9.716)rectangle++(0.04,0.04) (2.321,9.764)rectangle++(0.04,0.04)
  (2.361,9.803)rectangle++(0.04,0.04) (2.397,9.834)rectangle++(0.04,0.04) (2.428,9.859)rectangle++(0.04,0.04) (2.456,9.879)rectangle++(0.04,0.04) (2.481,9.895)rectangle++(0.04,0.04)
  (2.502,9.908)rectangle++(0.04,0.04) (2.521,9.918)rectangle++(0.04,0.04) (-2.096,4.351)rectangle++(0.04,0.04) (-1.609,5.378)rectangle++(0.04,0.04) (-1.155,6.232)rectangle++(0.04,0.04)
  (-0.735,6.94)rectangle++(0.04,0.04) (-0.349,7.524)rectangle++(0.04,0.04) (0.002,8.006)rectangle++(0.04,0.04) (0.32,8.401)rectangle++(0.04,0.04) (0.606,8.724)rectangle++(0.04,0.04)
  (0.862,8.988)rectangle++(0.04,0.04) (1.09,9.201)rectangle++(0.04,0.04) (1.292,9.374)rectangle++(0.04,0.04) (1.471,9.512)rectangle++(0.04,0.04) (1.629,9.623)rectangle++(0.04,0.04)
  (1.767,9.71)rectangle++(0.04,0.04) (1.888,9.779)rectangle++(0.04,0.04) (1.994,9.833)rectangle++(0.04,0.04) (2.086,9.875)rectangle++(0.04,0.04) (-2.171,4.241)rectangle++(0.04,0.04)
  (-1.558,2.023)rectangle++(0.04,0.04) (-1.246,3.378)rectangle++(0.04,0.04) (-0.926,4.517)rectangle++(0.04,0.04) (-0.609,5.473)rectangle++(0.04,0.04) (1.606,1.585)rectangle++(0.04,0.04)
  (1.427,2.879)rectangle++(0.04,0.04) (1.326,3.987)rectangle++(0.04,0.04) (1.284,4.932)rectangle++(0.04,0.04) (-1.047,2.969)rectangle++(0.04,0.04) (-0.773,4.162)rectangle++(0.04,0.04)
  (-0.493,5.165)rectangle++(0.04,0.04) (-0.215,6.006)rectangle++(0.04,0.04) (0.056,6.71)rectangle++(0.04,0.04) (0.314,7.297)rectangle++(0.04,0.04) (0.556,7.786)rectangle++(0.04,0.04)
  (0.782,8.192)rectangle++(0.04,0.04) (0.99,8.528)rectangle++(0.04,0.04) (1.181,8.806)rectangle++(0.04,0.04) (2.271,2.85)rectangle++(0.04,0.04) (2.042,3.928)rectangle++(0.04,0.04)
  (1.891,4.853)rectangle++(0.04,0.04) (-0.905,3.092)rectangle++(0.04,0.04) (-0.648,4.26)rectangle++(0.04,0.04) (-0.382,5.243)rectangle++(0.04,0.04) (-0.117,6.068)rectangle++(0.04,0.04)
  (1.699,1.856)rectangle++(0.04,0.04) (1.516,3.106)rectangle++(0.04,0.04) (1.411,4.175)rectangle++(0.04,0.04) (1.364,5.089)rectangle++(0.04,0.04) (1.361,5.867)rectangle++(0.04,0.04)
  (-1.274,3.193)rectangle++(0.04,0.04) (-0.958,4.361)rectangle++(0.04,0.04) (-0.642,5.341)rectangle++(0.04,0.04) (1.574,1.545)rectangle++(0.04,0.04) (1.398,2.847)rectangle++(0.04,0.04)
  (1.3,3.96)rectangle++(0.04,0.04) (1.261,4.91)rectangle++(0.04,0.04) (1.266,5.719)rectangle++(0.04,0.04) (1.303,6.407)rectangle++(0.04,0.04) (1.361,6.99)rectangle++(0.04,0.04)
  (1.435,7.483)rectangle++(0.04,0.04) (1.517,7.9)rectangle++(0.04,0.04) (1.603,8.25)rectangle++(0.04,0.04) (1.69,8.545)rectangle++(0.04,0.04) (1.776,8.792)rectangle++(0.04,0.04)
  (1.859,8.998)rectangle++(0.04,0.04) (-1.989,3.996)rectangle++(0.04,0.04) (-1.533,5.073)rectangle++(0.04,0.04) (-1.102,5.969)rectangle++(0.04,0.04) (-0.7,6.714)rectangle++(0.04,0.04)
  (-1.907,2.905)rectangle++(0.04,0.04) (-1.507,4.142)rectangle++(0.04,0.04) (-1.117,5.177)rectangle++(0.04,0.04) (-0.745,6.041)rectangle++(0.04,0.04) (-0.394,6.761)rectangle++(0.04,0.04)
  (-0.066,7.359)rectangle++(0.04,0.04) (0.236,7.854)rectangle++(0.04,0.04) (0.512,8.263)rectangle++(0.04,0.04) (0.764,8.599)rectangle++(0.04,0.04) (0.991,8.875)rectangle++(0.04,0.04)
  (1.195,9.1)rectangle++(0.04,0.04) (1.378,9.283)rectangle++(0.04,0.04) (1.54,9.432)rectangle++(0.04,0.04) (1.684,9.552)rectangle++(0.04,0.04) (1.811,9.648)rectangle++(0.04,0.04)
  (-2.167,4.128)rectangle++(0.04,0.04) (-1.679,5.192)rectangle++(0.04,0.04) (1.688,1.239)rectangle++(0.04,0.04) (1.482,2.582)rectangle++(0.04,0.04) (1.361,3.732)rectangle++(0.04,0.04)
  (1.304,4.714)rectangle++(0.04,0.04) (1.295,5.551)rectangle++(0.04,0.04) (-1.205,3.108)rectangle++(0.04,0.04) (-0.903,4.286)rectangle++(0.04,0.04) (-0.598,5.276)rectangle++(0.04,0.04)
  (-0.299,6.104)rectangle++(0.04,0.04) (-0.012,6.797)rectangle++(0.04,0.04) (0.259,7.374)rectangle++(0.04,0.04) (0.513,7.854)rectangle++(0.04,0.04) (0.748,8.251)rectangle++(0.04,0.04)
  (0.964,8.58)rectangle++(0.04,0.04) (-2.059,3.698)rectangle++(0.04,0.04) (-1.395,1.929)rectangle++(0.04,0.04) (-1.11,3.292)rectangle++(0.04,0.04) (-0.814,4.438)rectangle++(0.04,0.04)
  (1.348,1.283)rectangle++(0.04,0.04) (1.195,2.633)rectangle++(0.04,0.04) (1.118,3.787)rectangle++(0.04,0.04) (0.875,1.63)rectangle++(0.04,0.04) (0.807,2.946)rectangle++(0.04,0.04)
  (-0.626,2.423)rectangle++(0.04,0.04) (-0.437,3.681)rectangle++(0.04,0.04) (-0.227,4.742)rectangle++(0.04,0.04) (-0.005,5.636)rectangle++(0.04,0.04) (1.561,1.781)rectangle++(0.04,0.04)
  (1.396,3.048)rectangle++(0.04,0.04) (1.307,4.131)rectangle++(0.04,0.04) (1.274,5.055)rectangle++(0.04,0.04) (1.283,5.842)rectangle++(0.04,0.04) (1.322,6.511)rectangle++(0.04,0.04)
  (-1.45,3.325)rectangle++(0.04,0.04) (-1.176,1.987)rectangle++(0.04,0.04) (-0.922,3.332)rectangle++(0.04,0.04) (-0.653,4.466)rectangle++(0.04,0.04) (-0.378,5.418)rectangle++(0.04,0.04)
  (1.556,1.632)rectangle++(0.04,0.04) (1.386,2.921)rectangle++(0.04,0.04) (1.293,4.024)rectangle++(0.04,0.04) (1.258,4.965)rectangle++(0.04,0.04) (-1.06,2.97)rectangle++(0.04,0.04)
  (-0.785,4.164)rectangle++(0.04,0.04) (-0.503,5.167)rectangle++(0.04,0.04) (-0.223,6.008)rectangle++(0.04,0.04) (0.049,6.712)rectangle++(0.04,0.04) (0.308,7.299)rectangle++(0.04,0.04)
  (1.98,2.262)rectangle++(0.04,0.04) (1.771,3.44)rectangle++(0.04,0.04) (1.641,4.449)rectangle++(0.04,0.04) (-0.02,0.695)rectangle++(0.04,0.04) (0.18,0.592)rectangle++(0.04,0.04)
  (0.175,2.092)rectangle++(0.04,0.04) (0.23,3.367)rectangle++(0.04,0.04) (0.328,4.449)rectangle++(0.04,0.04) (0.455,5.365)rectangle++(0.04,0.04) (1.417,1.836)rectangle++(0.04,0.04)
  (1.275,3.1)rectangle++(0.04,0.04) (1.206,4.18)rectangle++(0.04,0.04) (1.19,5.101)rectangle++(0.04,0.04) (1.213,5.885)rectangle++(0.04,0.04) (1.264,6.55)rectangle++(0.04,0.04)
  (1.335,7.113)rectangle++(0.04,0.04) (1.417,7.589)rectangle++(0.04,0.04) (1.506,7.99)rectangle++(0.04,0.04) (1.994,2.739)rectangle++(0.04,0.04) (-0.335,2.65)rectangle++(0.04,0.04)
  (-0.777,2.095)rectangle++(0.04,0.04) (-0.579,3.408)rectangle++(0.04,0.04) (-1.023,2.206)rectangle++(0.04,0.04) (-0.784,3.512)rectangle++(0.04,0.04) (-0.528,4.613)rectangle++(0.04,0.04)
  (-0.266,5.538)rectangle++(0.04,0.04) (-1.514,2.746)rectangle++(0.04,0.04) (-1.18,3.991)rectangle++(0.04,0.04) (-0.845,5.036)rectangle++(0.04,0.04) (-1.5,2.502)rectangle++(0.04,0.04)
  (-1.177,3.783)rectangle++(0.04,0.04) (-0.851,4.859)rectangle++(0.04,0.04) (-0.531,5.76)rectangle++(0.04,0.04) (-0.223,6.514)rectangle++(0.04,0.04) (0.069,7.142)rectangle++(0.04,0.04)
  (0.342,7.664)rectangle++(0.04,0.04) (0.595,8.097)rectangle++(0.04,0.04) (0.827,8.455)rectangle++(0.04,0.04) (1.039,8.75)rectangle++(0.04,0.04) (1.231,8.992)rectangle++(0.04,0.04)
  (-2.113,3.85)rectangle++(0.04,0.04) (-1.644,4.954)rectangle++(0.04,0.04) (-1.202,5.872)rectangle++(0.04,0.04) (-0.789,6.636)rectangle++(0.04,0.04) (-0.407,7.268)rectangle++(0.04,0.04)
  (-0.058,7.79)rectangle++(0.04,0.04) (0.261,8.22)rectangle++(0.04,0.04) (0.548,8.573)rectangle++(0.04,0.04) (2.301,2.63)rectangle++(0.04,0.04) (-0.245,2.697)rectangle++(0.04,0.04)
  (-0.102,3.898)rectangle++(0.04,0.04) (-0.02,0.607)rectangle++(0.04,0.04) (0.005,2.113)rectangle++(0.04,0.04) (0.087,3.392)rectangle++(0.04,0.04) (0.207,4.476)rectangle++(0.04,0.04)
  (0.353,5.392)rectangle++(0.04,0.04) (1.44,1.816)rectangle++(0.04,0.04) (1.294,3.082)rectangle++(0.04,0.04) (1.221,4.164)rectangle++(0.04,0.04) (1.202,5.087)rectangle++(0.04,0.04)
  (-1.103,2.985)rectangle++(0.04,0.04) (-0.821,4.177)rectangle++(0.04,0.04) (-0.533,5.18)rectangle++(0.04,0.04) (-0.248,6.02)rectangle++(0.04,0.04) (1.705,1.81)rectangle++(0.04,0.04)
  (1.52,3.067)rectangle++(0.04,0.04) (1.412,4.142)rectangle++(0.04,0.04) (1.364,5.061)rectangle++(0.04,0.04) (1.36,5.843)rectangle++(0.04,0.04) (1.387,6.508)rectangle++(0.04,0.04)
  (1.437,7.073)rectangle++(0.04,0.04) (1.502,7.551)rectangle++(0.04,0.04) (1.577,7.954)rectangle++(0.04,0.04) (1.656,8.294)rectangle++(0.04,0.04) (1.737,8.58)rectangle++(0.04,0.04)
  (1.818,8.82)rectangle++(0.04,0.04) (1.896,9.02)rectangle++(0.04,0.04) (1.97,9.188)rectangle++(0.04,0.04) (2.04,9.327)rectangle++(0.04,0.04) (2.105,9.442)rectangle++(0.04,0.04)
  (2.164,9.538)rectangle++(0.04,0.04) (2.219,9.617)rectangle++(0.04,0.04) (2.269,9.682)rectangle++(0.04,0.04) (2.314,9.735)rectangle++(0.04,0.04) (2.354,9.778)rectangle++(0.04,0.04)
  (2.39,9.814)rectangle++(0.04,0.04) (-2.095,4.298)rectangle++(0.04,0.04) (-1.611,5.333)rectangle++(0.04,0.04) (-1.158,6.194)rectangle++(0.04,0.04) (-0.739,6.907)rectangle++(0.04,0.04)
  (-1.965,2.939)rectangle++(0.04,0.04) (-1.555,4.173)rectangle++(0.04,0.04) (-1.157,5.205)rectangle++(0.04,0.04) (-0.777,6.067)rectangle++(0.04,0.04) (-0.42,6.784)rectangle++(0.04,0.04)
  (-0.088,7.379)rectangle++(0.04,0.04) (0.218,7.872)rectangle++(0.04,0.04) (2.154,2.376)rectangle++(0.04,0.04) (-0.208,2.607)rectangle++(0.04,0.04) (-0.075,3.821)rectangle++(0.04,0.04)
  (-1.03,2.412)rectangle++(0.04,0.04) (-0.781,3.688)rectangle++(0.04,0.04) (-0.518,4.762)rectangle++(0.04,0.04) (-0.252,5.665)rectangle++(0.04,0.04) (0.01,6.421)rectangle++(0.04,0.04)
  (0.263,7.054)rectangle++(0.04,0.04) (0.504,7.582)rectangle++(0.04,0.04) (0.729,8.02)rectangle++(0.04,0.04) (-1.961,3.521)rectangle++(0.04,0.04) (-1.528,4.668)rectangle++(0.04,0.04)
  (1.519,1.153)rectangle++(0.04,0.04) (1.335,2.515)rectangle++(0.04,0.04) (1.233,3.681)rectangle++(0.04,0.04) (1.193,4.676)rectangle++(0.04,0.04) (1.199,5.523)rectangle++(0.04,0.04)
  (1.238,6.243)rectangle++(0.04,0.04) (1.3,6.853)rectangle++(0.04,0.04) (1.706,2.413)rectangle++(0.04,0.04) (1.545,3.579)rectangle++(0.04,0.04) (1.454,4.576)rectangle++(0.04,0.04)
  (-0.02,0.715)rectangle++(0.04,0.04) (0.009,2.205)rectangle++(0.04,0.04) (0.094,3.47)rectangle++(0.04,0.04) (0.217,4.542)rectangle++(0.04,0.04) (0.364,5.448)rectangle++(0.04,0.04)
  (0.525,6.213)rectangle++(0.04,0.04) (0.692,6.856)rectangle++(0.04,0.04) (0.861,7.396)rectangle++(0.04,0.04) (1.025,7.848)rectangle++(0.04,0.04) (1.183,8.226)rectangle++(0.04,0.04)
  (1.332,8.541)rectangle++(0.04,0.04) (1.472,8.803)rectangle++(0.04,0.04) (1.601,9.02)rectangle++(0.04,0.04) (1.72,9.199)rectangle++(0.04,0.04) (1.827,9.347)rectangle++(0.04,0.04)
  (1.925,9.468)rectangle++(0.04,0.04) (2.345,3.203)rectangle++(0.04,0.04) (2.119,4.225)rectangle++(0.04,0.04) (1.968,5.102)rectangle++(0.04,0.04) (1.116,2.166)rectangle++(0.04,0.04)
  (1.033,3.393)rectangle++(0.04,0.04) (1.012,4.439)rectangle++(0.04,0.04) (1.035,5.329)rectangle++(0.04,0.04) (1.091,6.084)rectangle++(0.04,0.04) (1.168,6.724)rectangle++(0.04,0.04)
  (1.26,7.265)rectangle++(0.04,0.04) (1.359,7.721)rectangle++(0.04,0.04) (1.462,8.105)rectangle++(0.04,0.04) (1.565,8.427)rectangle++(0.04,0.04) (1.665,8.696)rectangle++(0.04,0.04)
  (1.761,8.921)rectangle++(0.04,0.04) (1.851,9.109)rectangle++(0.04,0.04) (1.936,9.265)rectangle++(0.04,0.04) (2.014,9.394)rectangle++(0.04,0.04) (2.085,9.5)rectangle++(0.04,0.04)
  (2.15,9.588)rectangle++(0.04,0.04) (2.209,9.66)rectangle++(0.04,0.04) (2.262,9.719)rectangle++(0.04,0.04) (2.309,9.767)rectangle++(0.04,0.04) (2.351,9.806)rectangle++(0.04,0.04)
  (2.389,9.837)rectangle++(0.04,0.04) (2.422,9.862)rectangle++(0.04,0.04) (2.451,9.882)rectangle++(0.04,0.04) (2.476,9.898)rectangle++(0.04,0.04) (2.498,9.91)rectangle++(0.04,0.04)
  (2.518,9.92)rectangle++(0.04,0.04) (2.535,9.928)rectangle++(0.04,0.04) (2.549,9.933)rectangle++(0.04,0.04) (2.562,9.937)rectangle++(0.04,0.04) (2.573,9.941)rectangle++(0.04,0.04);
\fill[nuage]
  (2.583,9.943)rectangle++(0.04,0.04) (2.591,9.944)rectangle++(0.04,0.04) (2.598,9.945)rectangle++(0.04,0.04) (2.604,9.946)rectangle++(0.04,0.04) (-2.086,4.376)rectangle++(0.04,0.04)
  (1.521,0.938)rectangle++(0.04,0.04) (1.328,2.332)rectangle++(0.04,0.04) (0.436,1.335)rectangle++(0.04,0.04) (-0.02,0.197)rectangle++(0.04,0.04) (-0.02,0.015)rectangle++(0.04,0.04)
  (-0.019,1.609)rectangle++(0.04,0.04) (0.046,2.965)rectangle++(0.04,0.04) (0.156,4.115)rectangle++(0.04,0.04) (0.295,5.087)rectangle++(0.04,0.04) (1.363,1.728)rectangle++(0.04,0.04)
  (1.225,3.01)rectangle++(0.04,0.04) (1.16,4.106)rectangle++(0.04,0.04) (1.148,5.04)rectangle++(0.04,0.04) (1.175,5.834)rectangle++(0.04,0.04) (1.23,6.508)rectangle++(0.04,0.04)
  (1.304,7.079)rectangle++(0.04,0.04) (1.389,7.561)rectangle++(0.04,0.04) (1.481,7.968)rectangle++(0.04,0.04) (1.991,2.727)rectangle++(0.04,0.04) (1.8,3.835)rectangle++(0.04,0.04)
  (-0.658,2.847)rectangle++(0.04,0.04) (-0.448,4.042)rectangle++(0.04,0.04) (-0.221,5.05)rectangle++(0.04,0.04) (0.012,5.897)rectangle++(0.04,0.04) (0.244,6.609)rectangle++(0.04,0.04)
  (0.469,7.204)rectangle++(0.04,0.04) (0.685,7.701)rectangle++(0.04,0.04) (0.888,8.114)rectangle++(0.04,0.04) (1.077,8.458)rectangle++(0.04,0.04) (1.252,8.742)rectangle++(0.04,0.04)
  (-2.044,3.8)rectangle++(0.04,0.04) (-1.587,4.908)rectangle++(0.04,0.04) (1.595,1.195)rectangle++(0.04,0.04) (1.401,2.548)rectangle++(0.04,0.04) (1.291,3.706)rectangle++(0.04,0.04)
  (1.243,4.695)rectangle++(0.04,0.04) (1.242,5.537)rectangle++(0.04,0.04) (1.275,6.253)rectangle++(0.04,0.04) (1.332,6.86)rectangle++(0.04,0.04) (1.404,7.374)rectangle++(0.04,0.04)
  (1.486,7.808)rectangle++(0.04,0.04) (1.574,8.174)rectangle++(0.04,0.04) (2.035,2.801)rectangle++(0.04,0.04) (1.84,3.895)rectangle++(0.04,0.04) (1.717,4.834)rectangle++(0.04,0.04)
  (1.651,5.636)rectangle++(0.04,0.04) (1.627,6.321)rectangle++(0.04,0.04) (1.633,6.904)rectangle++(0.04,0.04) (1.662,7.399)rectangle++(0.04,0.04) (1.707,7.819)rectangle++(0.04,0.04)
  (1.761,8.174)rectangle++(0.04,0.04) (1.822,8.474)rectangle++(0.04,0.04) (1.885,8.726)rectangle++(0.04,0.04) (1.949,8.938)rectangle++(0.04,0.04) (2.012,9.115)rectangle++(0.04,0.04)
  (2.073,9.264)rectangle++(0.04,0.04) (2.13,9.388)rectangle++(0.04,0.04) (2.184,9.49)rectangle++(0.04,0.04) (2.234,9.576)rectangle++(0.04,0.04) (2.28,9.646)rectangle++(0.04,0.04)
  (2.321,9.704)rectangle++(0.04,0.04) (2.359,9.752)rectangle++(0.04,0.04) (2.393,9.791)rectangle++(0.04,0.04) (2.365,3.402)rectangle++(0.04,0.04) (2.144,4.393)rectangle++(0.04,0.04)
  (1.996,5.245)rectangle++(0.04,0.04) (1.904,5.974)rectangle++(0.04,0.04) (1.855,6.598)rectangle++(0.04,0.04) (1.839,7.131)rectangle++(0.04,0.04) (1.846,7.584)rectangle++(0.04,0.04)
  (1.87,7.968)rectangle++(0.04,0.04) (1.906,8.295)rectangle++(0.04,0.04) (1.95,8.57)rectangle++(0.04,0.04) (1.998,8.803)rectangle++(0.04,0.04) (-1.91,3.985)rectangle++(0.04,0.04)
  (-1.467,5.06)rectangle++(0.04,0.04) (-1.046,5.956)rectangle++(0.04,0.04) (-0.653,6.701)rectangle++(0.04,0.04) (-0.29,7.318)rectangle++(0.04,0.04) (0.044,7.828)rectangle++(0.04,0.04)
  (0.349,8.248)rectangle++(0.04,0.04) (0.624,8.593)rectangle++(0.04,0.04) (0.872,8.875)rectangle++(0.04,0.04) (2.337,2.787)rectangle++(0.04,0.04) (2.096,3.872)rectangle++(0.04,0.04)
  (1.934,4.803)rectangle++(0.04,0.04) (1.834,5.602)rectangle++(0.04,0.04) (1.276,2.251)rectangle++(0.04,0.04) (0.422,1.302)rectangle++(0.04,0.04) (0.408,2.686)rectangle++(0.04,0.04)
  (0.452,3.863)rectangle++(0.04,0.04) (0.537,4.862)rectangle++(0.04,0.04) (0.648,5.707)rectangle++(0.04,0.04) (0.777,6.421)rectangle++(0.04,0.04) (0.915,7.023)rectangle++(0.04,0.04)
  (1.057,7.529)rectangle++(0.04,0.04) (1.197,7.954)rectangle++(0.04,0.04) (1.334,8.309)rectangle++(0.04,0.04) (1.464,8.606)rectangle++(0.04,0.04) (1.586,8.852)rectangle++(0.04,0.04)
  (1.7,9.057)rectangle++(0.04,0.04) (1.805,9.227)rectangle++(0.04,0.04) (-0.02,1.46)rectangle++(0.04,0.04) (-0.405,1.905)rectangle++(0.04,0.04) (-0.27,3.232)rectangle++(0.04,0.04)
  (-0.102,4.354)rectangle++(0.04,0.04) (-1.174,2.523)rectangle++(0.04,0.04) (-0.899,3.788)rectangle++(0.04,0.04) (-0.615,4.852)rectangle++(0.04,0.04) (-0.331,5.745)rectangle++(0.04,0.04)
  (-1.581,2.777)rectangle++(0.04,0.04) (-1.235,4.02)rectangle++(0.04,0.04) (-0.891,5.062)rectangle++(0.04,0.04) (-0.557,5.935)rectangle++(0.04,0.04) (-0.238,6.663)rectangle++(0.04,0.04)
  (0.062,7.269)rectangle++(0.04,0.04) (0.341,7.773)rectangle++(0.04,0.04) (0.599,8.189)rectangle++(0.04,0.04) (0.834,8.533)rectangle++(0.04,0.04) (1.048,8.816)rectangle++(0.04,0.04)
  (1.241,9.048)rectangle++(0.04,0.04) (1.415,9.237)rectangle++(0.04,0.04) (1.57,9.391)rectangle++(0.04,0.04) (1.708,9.516)rectangle++(0.04,0.04) (1.83,9.616)rectangle++(0.04,0.04)
  (1.938,9.697)rectangle++(0.04,0.04) (-2.155,4.168)rectangle++(0.04,0.04) (-1.667,5.225)rectangle++(0.04,0.04) (-1.21,6.104)rectangle++(0.04,0.04) (-0.787,6.833)rectangle++(0.04,0.04)
  (-0.398,7.436)rectangle++(0.04,0.04) (-0.043,7.933)rectangle++(0.04,0.04) (0.279,8.341)rectangle++(0.04,0.04) (0.568,8.675)rectangle++(0.04,0.04) (0.828,8.947)rectangle++(0.04,0.04)
  (2.364,2.793)rectangle++(0.04,0.04) (2.119,3.875)rectangle++(0.04,0.04) (1.954,4.805)rectangle++(0.04,0.04) (1.851,5.603)rectangle++(0.04,0.04) (1.795,6.284)rectangle++(0.04,0.04)
  (1.775,6.866)rectangle++(0.04,0.04) (1.781,7.361)rectangle++(0.04,0.04) (1.806,7.782)rectangle++(0.04,0.04) (1.844,8.139)rectangle++(0.04,0.04) (1.891,8.44)rectangle++(0.04,0.04)
  (1.943,8.695)rectangle++(0.04,0.04) (1.997,8.909)rectangle++(0.04,0.04) (2.052,9.089)rectangle++(0.04,0.04) (2.105,9.24)rectangle++(0.04,0.04) (2.157,9.366)rectangle++(0.04,0.04)
  (2.206,9.471)rectangle++(0.04,0.04) (2.252,9.558)rectangle++(0.04,0.04) (2.294,9.631)rectangle++(0.04,0.04) (2.333,9.691)rectangle++(0.04,0.04) (2.368,9.74)rectangle++(0.04,0.04)
  (2.4,9.78)rectangle++(0.04,0.04) (2.429,9.813)rectangle++(0.04,0.04) (2.455,9.84)rectangle++(0.04,0.04) (2.37,3.43)rectangle++(0.04,0.04) (2.149,4.417)rectangle++(0.04,0.04)
  (2.001,5.265)rectangle++(0.04,0.04) (1.909,5.991)rectangle++(0.04,0.04) (1.861,6.612)rectangle++(0.04,0.04) (1.555,2.501)rectangle++(0.04,0.04) (1.42,3.66)rectangle++(0.04,0.04)
  (1.351,4.65)rectangle++(0.04,0.04) (1.332,5.495)rectangle++(0.04,0.04) (-1.183,3.104)rectangle++(0.04,0.04) (-1.065,2)rectangle++(0.04,0.04) (0.702,0.633)rectangle++(0.04,0.04)
  (0.62,2.106)rectangle++(0.04,0.04) (0.609,3.362)rectangle++(0.04,0.04) (0.833,1.395)rectangle++(0.04,0.04) (0.761,2.749)rectangle++(0.04,0.04) (0.755,3.902)rectangle++(0.04,0.04)
  (0.795,4.883)rectangle++(0.04,0.04) (0.869,5.715)rectangle++(0.04,0.04) (0.965,6.419)rectangle++(0.04,0.04) (1.075,7.014)rectangle++(0.04,0.04) (1.192,7.515)rectangle++(0.04,0.04)
  (1.312,7.936)rectangle++(0.04,0.04) (1.43,8.289)rectangle++(0.04,0.04) (1.545,8.585)rectangle++(0.04,0.04) (1.654,8.832)rectangle++(0.04,0.04) (1.757,9.037)rectangle++(0.04,0.04)
  (1.853,9.207)rectangle++(0.04,0.04) (1.941,9.348)rectangle++(0.04,0.04) (-2.064,4.092)rectangle++(0.04,0.04) (-1.593,5.157)rectangle++(0.04,0.04) (-1.15,6.043)rectangle++(0.04,0.04)
  (-0.738,6.779)rectangle++(0.04,0.04) (-0.358,7.388)rectangle++(0.04,0.04) (-0.011,7.89)rectangle++(0.04,0.04) (0.304,8.303)rectangle++(0.04,0.04) (2.262,2.502)rectangle++(0.04,0.04)
  (2.021,3.632)rectangle++(0.04,0.04) (1.861,4.603)rectangle++(0.04,0.04) (1.763,5.434)rectangle++(0.04,0.04) (1.714,6.145)rectangle++(0.04,0.04) (1.701,6.751)rectangle++(0.04,0.04)
  (1.713,7.266)rectangle++(0.04,0.04) (1.745,7.704)rectangle++(0.04,0.04) (1.789,8.075)rectangle++(0.04,0.04) (1.841,8.388)rectangle++(0.04,0.04) (1.899,8.653)rectangle++(0.04,0.04)
  (1.958,8.875)rectangle++(0.04,0.04) (2.017,9.062)rectangle++(0.04,0.04) (2.075,9.218)rectangle++(0.04,0.04) (2.252,3.182)rectangle++(0.04,0.04) (2.04,4.211)rectangle++(0.04,0.04)
  (1.9,5.094)rectangle++(0.04,0.04) (1.124,2.146)rectangle++(0.04,0.04) (1.039,3.376)rectangle++(0.04,0.04) (1.016,4.424)rectangle++(0.04,0.04) (1.038,5.316)rectangle++(0.04,0.04)
  (1.093,6.073)rectangle++(0.04,0.04) (1.17,6.715)rectangle++(0.04,0.04) (1.261,7.257)rectangle++(0.04,0.04) (1.825,2.499)rectangle++(0.04,0.04) (1.649,3.648)rectangle++(0.04,0.04)
  (1.546,4.631)rectangle++(0.04,0.04) (1.497,5.471)rectangle++(0.04,0.04) (1.489,6.186)rectangle++(0.04,0.04) (1.511,6.795)rectangle++(0.04,0.04) (1.554,7.311)rectangle++(0.04,0.04)
  (1.611,7.749)rectangle++(0.04,0.04) (1.677,8.118)rectangle++(0.04,0.04) (1.748,8.43)rectangle++(0.04,0.04) (1.821,8.691)rectangle++(0.04,0.04) (-1.917,3.92)rectangle++(0.04,0.04)
  (-1.475,5.005)rectangle++(0.04,0.04) (-1.055,5.909)rectangle++(0.04,0.04) (1.796,1.574)rectangle++(0.04,0.04) (1.587,2.862)rectangle++(0.04,0.04) (1.461,3.966)rectangle++(0.04,0.04)
  (1.398,4.908)rectangle++(0.04,0.04) (1.383,5.712)rectangle++(0.04,0.04) (1.402,6.396)rectangle++(0.04,0.04) (1.445,6.977)rectangle++(0.04,0.04) (1.505,7.469)rectangle++(0.04,0.04)
  (1.576,7.885)rectangle++(0.04,0.04) (1.653,8.235)rectangle++(0.04,0.04) (2.041,2.836)rectangle++(0.04,0.04) (1.846,3.925)rectangle++(0.04,0.04) (1.724,4.859)rectangle++(0.04,0.04)
  (-0.94,3.054)rectangle++(0.04,0.04) (0.979,0.919)rectangle++(0.04,0.04) (0.867,2.338)rectangle++(0.04,0.04) (0.828,3.549)rectangle++(0.04,0.04) (0.843,4.58)rectangle++(0.04,0.04)
  (0.898,5.455)rectangle++(0.04,0.04) (1.375,1.973)rectangle++(0.04,0.04) (1.246,3.218)rectangle++(0.04,0.04) (1.185,4.282)rectangle++(0.04,0.04) (1.177,5.188)rectangle++(0.04,0.04)
  (1.205,5.959)rectangle++(0.04,0.04) (1.261,6.613)rectangle++(0.04,0.04) (-1.488,3.334)rectangle++(0.04,0.04) (-1.134,4.49)rectangle++(0.04,0.04) (-0.787,5.458)rectangle++(0.04,0.04)
  (-0.452,6.267)rectangle++(0.04,0.04) (-0.136,6.941)rectangle++(0.04,0.04) (-1.853,3.085)rectangle++(0.04,0.04) (-1.454,4.292)rectangle++(0.04,0.04) (1.403,1.082)rectangle++(0.04,0.04)
  (1.233,2.46)rectangle++(0.04,0.04) (1.144,3.638)rectangle++(0.04,0.04) (1.116,4.643)rectangle++(0.04,0.04) (1.132,5.498)rectangle++(0.04,0.04) (1.18,6.224)rectangle++(0.04,0.04)
  (1.25,6.839)rectangle++(0.04,0.04) (1.334,7.36)rectangle++(0.04,0.04) (1.426,7.799)rectangle++(0.04,0.04) (1.522,8.168)rectangle++(0.04,0.04) (1.618,8.478)rectangle++(0.04,0.04)
  (1.712,8.738)rectangle++(0.04,0.04) (1.803,8.955)rectangle++(0.04,0.04) (1.888,9.136)rectangle++(0.04,0.04) (1.968,9.286)rectangle++(0.04,0.04) (2.042,9.411)rectangle++(0.04,0.04)
  (2.11,9.514)rectangle++(0.04,0.04) (2.172,9.598)rectangle++(0.04,0.04) (2.228,9.668)rectangle++(0.04,0.04) (2.278,9.725)rectangle++(0.04,0.04) (2.323,9.771)rectangle++(0.04,0.04)
  (2.37,3.379)rectangle++(0.04,0.04) (0.573,1.857)rectangle++(0.04,0.04) (-0.389,2.129)rectangle++(0.04,0.04) (-0.248,3.422)rectangle++(0.04,0.04) (-0.076,4.515)rectangle++(0.04,0.04)
  (0.114,5.437)rectangle++(0.04,0.04) (1.488,1.764)rectangle++(0.04,0.04) (1.333,3.036)rectangle++(0.04,0.04) (1.252,4.124)rectangle++(0.04,0.04) (1.227,5.051)rectangle++(0.04,0.04)
  (1.243,5.841)rectangle++(0.04,0.04) (1.288,6.511)rectangle++(0.04,0.04) (1.353,7.079)rectangle++(0.04,0.04) (1.431,7.559)rectangle++(0.04,0.04) (1.885,2.616)rectangle++(0.04,0.04)
  (-0.325,2.598)rectangle++(0.04,0.04) (-0.174,3.818)rectangle++(0.04,0.04) (0.002,4.848)rectangle++(0.04,0.04) (0.194,5.717)rectangle++(0.04,0.04) (0.391,6.448)rectangle++(0.04,0.04)
  (0.588,7.061)rectangle++(0.04,0.04) (0.78,7.575)rectangle++(0.04,0.04) (0.964,8.003)rectangle++(0.04,0.04) (2.079,2.601)rectangle++(0.04,0.04) (1.869,3.724)rectangle++(0.04,0.04)
  (1.735,4.687)rectangle++(0.04,0.04) (-0.893,3.019)rectangle++(0.04,0.04) (0.962,0.923)rectangle++(0.04,0.04) (0.852,2.342)rectangle++(0.04,0.04) (0.816,3.553)rectangle++(0.04,0.04)
  (0.833,4.583)rectangle++(0.04,0.04) (0.89,5.459)rectangle++(0.04,0.04) (-1.263,2.995)rectangle++(0.04,0.04) (-0.956,4.192)rectangle++(0.04,0.04) (-0.647,5.198)rectangle++(0.04,0.04)
  (-0.344,6.04)rectangle++(0.04,0.04) (-0.053,6.744)rectangle++(0.04,0.04) (0.222,7.331)rectangle++(0.04,0.04) (-1.883,3.253)rectangle++(0.04,0.04) (-1.472,4.437)rectangle++(0.04,0.04)
  (-1.076,5.426)rectangle++(0.04,0.04) (-0.02,0.851)rectangle++(0.04,0.04) (0.015,2.321)rectangle++(0.04,0.04) (0.103,3.568)rectangle++(0.04,0.04) (0.228,4.625)rectangle++(0.04,0.04)
  (0.377,5.518)rectangle++(0.04,0.04) (-1.381,2.89)rectangle++(0.04,0.04) (-1.06,4.108)rectangle++(0.04,0.04) (1.292,1.14)rectangle++(0.04,0.04) (1.141,2.514)rectangle++(0.04,0.04)
  (0.515,1.33)rectangle++(0.04,0.04) (-0.264,2)rectangle++(0.04,0.04) (-0.147,3.307)rectangle++(0.04,0.04) (0.005,4.413)rectangle++(0.04,0.04) (0.179,5.347)rectangle++(0.04,0.04)
  (0.364,6.134)rectangle++(0.04,0.04) (0.552,6.795)rectangle++(0.04,0.04) (0.739,7.35)rectangle++(0.04,0.04) (1.93,2.386)rectangle++(0.04,0.04) (1.734,3.547)rectangle++(0.04,0.04)
  (-0.597,2.768)rectangle++(0.04,0.04) (-0.399,3.973)rectangle++(0.04,0.04) (-0.182,4.989)rectangle++(0.04,0.04) (0.042,5.844)rectangle++(0.04,0.04) (0.268,6.562)rectangle++(0.04,0.04)
  (0.488,7.163)rectangle++(0.04,0.04) (0.699,7.666)rectangle++(0.04,0.04) (0.899,8.084)rectangle++(0.04,0.04) (1.085,8.432)rectangle++(0.04,0.04) (1.257,8.72)rectangle++(0.04,0.04)
  (1.415,8.958)rectangle++(0.04,0.04) (1.559,9.154)rectangle++(0.04,0.04) (1.689,9.314)rectangle++(0.04,0.04) (1.806,9.446)rectangle++(0.04,0.04) (1.911,9.553)rectangle++(0.04,0.04)
  (2.371,3.22)rectangle++(0.04,0.04) (2.142,4.238)rectangle++(0.04,0.04) (1.988,5.113)rectangle++(0.04,0.04) (1.116,2.174)rectangle++(0.04,0.04) (1.033,3.399)rectangle++(0.04,0.04)
  (1.012,4.444)rectangle++(0.04,0.04) (1.036,5.333)rectangle++(0.04,0.04) (1.092,6.088)rectangle++(0.04,0.04) (1.169,6.727)rectangle++(0.04,0.04) (1.261,7.268)rectangle++(0.04,0.04)
  (1.36,7.723)rectangle++(0.04,0.04) (1.463,8.107)rectangle++(0.04,0.04) (-1.836,3.709)rectangle++(0.04,0.04) (-1.415,4.822)rectangle++(0.04,0.04) (-1.012,5.752)rectangle++(0.04,0.04)
  (-0.632,6.526)rectangle++(0.04,0.04) (-0.279,7.168)rectangle++(0.04,0.04) (-1.941,3.102)rectangle++(0.04,0.04) (-1.528,4.31)rectangle++(0.04,0.04) (-1.128,5.321)rectangle++(0.04,0.04)
  (-0.748,6.164)rectangle++(0.04,0.04) (-0.392,6.866)rectangle++(0.04,0.04) (-0.061,7.448)rectangle++(0.04,0.04) (2.077,2.202)rectangle++(0.04,0.04) (-0.178,2.551)rectangle++(0.04,0.04)
  (-0.052,3.772)rectangle++(0.04,0.04) (0.105,4.804)rectangle++(0.04,0.04) (0.279,5.676)rectangle++(0.04,0.04) (1.53,1.865)rectangle++(0.04,0.04) (1.373,3.12)rectangle++(0.04,0.04)
  (0.65,1.536)rectangle++(0.04,0.04) (0.612,2.876)rectangle++(0.04,0.04) (-0.646,2.362)rectangle++(0.04,0.04) (-0.457,3.63)rectangle++(0.04,0.04) (-0.246,4.7)rectangle++(0.04,0.04)
  (-0.023,5.601)rectangle++(0.04,0.04) (0.202,6.358)rectangle++(0.04,0.04) (0.424,6.992)rectangle++(0.04,0.04) (0.638,7.523)rectangle++(0.04,0.04) (-1.85,3.391)rectangle++(0.04,0.04)
  (-1.439,4.552)rectangle++(0.04,0.04) (-1.043,5.523)rectangle++(0.04,0.04) (-0.668,6.333)rectangle++(0.04,0.04) (-0.317,7.006)rectangle++(0.04,0.04) (0.009,7.564)rectangle++(0.04,0.04)
  (0.308,8.025)rectangle++(0.04,0.04) (-2.046,3.425)rectangle++(0.04,0.04) (-1.604,4.59)rectangle++(0.04,0.04) (-1.182,5.562)rectangle++(0.04,0.04) (-0.785,6.371)rectangle++(0.04,0.04)
  (-0.414,7.043)rectangle++(0.04,0.04) (-0.073,7.599)rectangle++(0.04,0.04) (0.24,8.058)rectangle++(0.04,0.04) (0.524,8.436)rectangle++(0.04,0.04) (0.781,8.746)rectangle++(0.04,0.04)
  (2.314,2.732)rectangle++(0.04,0.04) (2.074,3.826)rectangle++(0.04,0.04) (1.914,4.765)rectangle++(0.04,0.04) (1.815,5.57)rectangle++(0.04,0.04) (-1.106,3.232)rectangle++(0.04,0.04)
  (-0.813,4.388)rectangle++(0.04,0.04) (-0.518,5.358)rectangle++(0.04,0.04) (-0.228,6.171)rectangle++(0.04,0.04) (0.051,6.851)rectangle++(0.04,0.04) (0.315,7.418)rectangle++(0.04,0.04)
  (0.562,7.888)rectangle++(0.04,0.04) (0.791,8.279)rectangle++(0.04,0.04) (1.001,8.602)rectangle++(0.04,0.04) (1.193,8.868)rectangle++(0.04,0.04) (1.367,9.086)rectangle++(0.04,0.04)
  (2.322,2.966)rectangle++(0.04,0.04) (2.09,4.024)rectangle++(0.04,0.04) (1.935,4.933)rectangle++(0.04,0.04) (1.84,5.712)rectangle++(0.04,0.04) (1.79,6.378)rectangle++(0.04,0.04)
  (1.775,6.946)rectangle++(0.04,0.04) (1.784,7.429)rectangle++(0.04,0.04) (1.811,7.84)rectangle++(0.04,0.04) (1.851,8.187)rectangle++(0.04,0.04) (1.899,8.481)rectangle++(0.04,0.04)
  (1.951,8.729)rectangle++(0.04,0.04) (2.005,8.938)rectangle++(0.04,0.04) (2.06,9.113)rectangle++(0.04,0.04) (2.113,9.26)rectangle++(0.04,0.04) (2.164,9.383)rectangle++(0.04,0.04)
  (2.285,3.245)rectangle++(0.04,0.04) (2.07,4.263)rectangle++(0.04,0.04) (1.928,5.137)rectangle++(0.04,0.04) (1.842,5.885)rectangle++(0.04,0.04) (1.799,6.525)rectangle++(0.04,0.04)
  (1.788,7.071)rectangle++(0.04,0.04) (1.8,7.535)rectangle++(0.04,0.04) (1.829,7.929)rectangle++(0.04,0.04) (1.87,8.262)rectangle++(0.04,0.04) (1.918,8.544)rectangle++(0.04,0.04)
  (2.087,2.979)rectangle++(0.04,0.04) (1.891,4.045)rectangle++(0.04,0.04) (1.767,4.959)rectangle++(0.04,0.04) (1.698,5.741)rectangle++(0.04,0.04) (1.671,6.408)rectangle++(0.04,0.04)
  (1.674,6.976)rectangle++(0.04,0.04) (1.7,7.459)rectangle++(0.04,0.04) (1.741,7.868)rectangle++(0.04,0.04) (1.793,8.214)rectangle++(0.04,0.04) (1.85,8.507)rectangle++(0.04,0.04)
  (1.911,8.753)rectangle++(0.04,0.04) (-1.915,3.954)rectangle++(0.04,0.04) (-1.432,2.018)rectangle++(0.04,0.04) (-1.139,3.369)rectangle++(0.04,0.04) (-0.835,4.506)rectangle++(0.04,0.04)
  (-0.532,5.459)rectangle++(0.04,0.04) (-0.236,6.258)rectangle++(0.04,0.04) (1.77,1.871)rectangle++(0.04,0.04) (1.577,3.115)rectangle++(0.04,0.04) (1.463,4.181)rectangle++(0.04,0.04)
  (1.409,5.092)rectangle++(0.04,0.04) (1.399,5.868)rectangle++(0.04,0.04) (-1.267,3.202)rectangle++(0.04,0.04) (-0.951,4.368)rectangle++(0.04,0.04) (1.348,1.231)rectangle++(0.04,0.04)
  (1.193,2.589)rectangle++(0.04,0.04) (1.116,3.749)rectangle++(0.04,0.04) (1.096,4.738)rectangle++(0.04,0.04) (1.119,5.58)rectangle++(0.04,0.04) (1.172,6.294)rectangle++(0.04,0.04)
  (1.246,6.899)rectangle++(0.04,0.04) (1.333,7.411)rectangle++(0.04,0.04) (1.427,7.842)rectangle++(0.04,0.04) (1.964,2.683)rectangle++(0.04,0.04) (1.775,3.798)rectangle++(0.04,0.04)
  (0.78,1.803)rectangle++(0.04,0.04) (0.733,3.098)rectangle++(0.04,0.04) (0.745,4.2)rectangle++(0.04,0.04) (0.799,5.136)rectangle++(0.04,0.04) (0.882,5.93)rectangle++(0.04,0.04)
  (0.985,6.601)rectangle++(0.04,0.04) (1.099,7.168)rectangle++(0.04,0.04) (1.219,7.645)rectangle++(0.04,0.04) (1.339,8.046)rectangle++(0.04,0.04) (1.458,8.382)rectangle++(0.04,0.04)
  (1.573,8.662)rectangle++(0.04,0.04) (1.681,8.896)rectangle++(0.04,0.04) (1.782,9.091)rectangle++(0.04,0.04) (-2.028,3.999)rectangle++(0.04,0.04) (-1.566,5.076)rectangle++(0.04,0.04)
  (-1.13,5.974)rectangle++(0.04,0.04) (-0.724,6.719)rectangle++(0.04,0.04) (-1.913,2.901)rectangle++(0.04,0.04) (-1.512,4.138)rectangle++(0.04,0.04) (-1.122,5.174)rectangle++(0.04,0.04)
  (1.6,1.38)rectangle++(0.04,0.04) (1.413,2.705)rectangle++(0.04,0.04) (1.307,3.839)rectangle++(0.04,0.04) (-0.02,0.597)rectangle++(0.04,0.04) (0.005,2.105)rectangle++(0.04,0.04)
  (0.086,3.385)rectangle++(0.04,0.04) (0.918,1.265)rectangle++(0.04,0.04) (0.828,2.635)rectangle++(0.04,0.04) (0.807,3.802)rectangle++(0.04,0.04) (0.926,1.552)rectangle++(0.04,0.04)
  (0.847,2.879)rectangle++(0.04,0.04) (0.833,4.009)rectangle++(0.04,0.04) (0.866,4.971)rectangle++(0.04,0.04) (-1.14,2.882)rectangle++(0.04,0.04) (-0.999,1.961)rectangle++(0.04,0.04)
  (-0.773,3.303)rectangle++(0.04,0.04) (-0.527,4.434)rectangle++(0.04,0.04) (-0.273,5.387)rectangle++(0.04,0.04) (-0.018,6.186)rectangle++(0.04,0.04) (0.23,6.855)rectangle++(0.04,0.04)
  (0.467,7.414)rectangle++(0.04,0.04) (-1.855,3.327)rectangle++(0.04,0.04) (-1.446,4.499)rectangle++(0.04,0.04) (-1.051,5.478)rectangle++(0.04,0.04) (-0.677,6.295)rectangle++(0.04,0.04)
  (-0.326,6.974)rectangle++(0.04,0.04) (0,7.537)rectangle++(0.04,0.04) (0.299,8.002)rectangle++(0.04,0.04) (0.572,8.386)rectangle++(0.04,0.04) (2.245,2.591)rectangle++(0.04,0.04)
  (2.01,3.709)rectangle++(0.04,0.04) (1.854,4.669)rectangle++(0.04,0.04) (1.761,5.49)rectangle++(0.04,0.04) (1.714,6.193)rectangle++(0.04,0.04) (1.459,2.362)rectangle++(0.04,0.04)
  (1.333,3.545)rectangle++(0.04,0.04) (1.272,4.556)rectangle++(0.04,0.04) (-0.951,2.884)rectangle++(0.04,0.04) (-0.696,4.086)rectangle++(0.04,0.04) (-0.43,5.097)rectangle++(0.04,0.04)
  (-0.164,5.946)rectangle++(0.04,0.04) (1.672,1.814)rectangle++(0.04,0.04) (1.492,3.072)rectangle++(0.04,0.04) (1.388,4.147)rectangle++(0.04,0.04) (1.344,5.066)rectangle++(0.04,0.04)
  (1.199,1.995)rectangle++(0.04,0.04) (1.097,3.244)rectangle++(0.04,0.04) (1.06,4.31)rectangle++(0.04,0.04) (1.071,5.217)rectangle++(0.04,0.04) (1.117,5.988)rectangle++(0.04,0.04)
  (1.492,2.158)rectangle++(0.04,0.04) (-0.284,2.407)rectangle++(0.04,0.04) (-0.147,3.653)rectangle++(0.04,0.04) (0.019,4.707)rectangle++(0.04,0.04) (0.202,5.597)rectangle++(0.04,0.04)
  (0.393,6.345)rectangle++(0.04,0.04) (0.586,6.974)rectangle++(0.04,0.04) (1.847,2.256)rectangle++(0.04,0.04) (1.658,3.44)rectangle++(0.04,0.04) (1.545,4.454)rectangle++(0.04,0.04)
  (1.489,5.32)rectangle++(0.04,0.04) (1.476,6.059)rectangle++(0.04,0.04) (1.495,6.687)rectangle++(0.04,0.04) (1.536,7.22)rectangle++(0.04,0.04) (1.592,7.672)rectangle++(0.04,0.04)
  (1.658,8.054)rectangle++(0.04,0.04) (1.729,8.376)rectangle++(0.04,0.04) (1.803,8.646)rectangle++(0.04,0.04) (1.876,8.873)rectangle++(0.04,0.04) (1.947,9.064)rectangle++(0.04,0.04)
  (2.016,9.222)rectangle++(0.04,0.04) (2.08,9.355)rectangle++(0.04,0.04) (2.14,9.464)rectangle++(0.04,0.04) (2.312,3.258)rectangle++(0.04,0.04) (2.093,4.273)rectangle++(0.04,0.04)
  (1.948,5.144)rectangle++(0.04,0.04) (1.859,5.891)rectangle++(0.04,0.04) (1.353,2.327)rectangle++(0.04,0.04) (1.241,3.52)rectangle++(0.04,0.04) (1.194,4.539)rectangle++(0.04,0.04)
  (1.194,5.406)rectangle++(0.04,0.04) (1.317,2.038)rectangle++(0.04,0.04) (1.199,3.276)rectangle++(0.04,0.04) (1.148,4.333)rectangle++(0.04,0.04) (1.147,5.233)rectangle++(0.04,0.04)
  (1.182,5.998)rectangle++(0.04,0.04) (1.242,6.648)rectangle++(0.04,0.04) (1.658,2.348)rectangle++(0.04,0.04) (1.501,3.526)rectangle++(0.04,0.04) (1.414,4.533)rectangle++(0.04,0.04)
  (1.381,5.393)rectangle++(0.04,0.04) (1.388,6.125)rectangle++(0.04,0.04) (1.422,6.747)rectangle++(0.04,0.04) (1.477,7.274)rectangle++(0.04,0.04) (1.544,7.72)rectangle++(0.04,0.04)
  (1.619,8.097)rectangle++(0.04,0.04) (-1.803,3.743)rectangle++(0.04,0.04) (-1.385,4.85)rectangle++(0.04,0.04) (-0.985,5.774)rectangle++(0.04,0.04) (-0.609,6.543)rectangle++(0.04,0.04)
  (-0.258,7.182)rectangle++(0.04,0.04) (2.032,2.087)rectangle++(0.04,0.04) (1.809,3.289)rectangle++(0.04,0.04) (1.667,4.319)rectangle++(0.04,0.04) (1.587,5.201)rectangle++(0.04,0.04)
  (-1.056,3.098)rectangle++(0.04,0.04) (-0.776,4.272)rectangle++(0.04,0.04) (-0.491,5.258)rectangle++(0.04,0.04) (-0.209,6.085)rectangle++(0.04,0.04) (0.064,6.777)rectangle++(0.04,0.04)
  (0.323,7.354)rectangle++(0.04,0.04) (0.566,7.834)rectangle++(0.04,0.04) (0.793,8.233)rectangle++(0.04,0.04) (1.001,8.562)rectangle++(0.04,0.04) (1.191,8.834)rectangle++(0.04,0.04)
  (-2.08,3.806)rectangle++(0.04,0.04) (-1.618,4.915)rectangle++(0.04,0.04) (1.601,1.189)rectangle++(0.04,0.04) (1.407,2.543)rectangle++(0.04,0.04) (1.295,3.701)rectangle++(0.04,0.04)
  (-0.725,2.701)rectangle++(0.04,0.04) (-0.51,3.921)rectangle++(0.04,0.04) (-1.143,2.334)rectangle++(0.04,0.04) (-0.88,3.626)rectangle++(0.04,0.04) (-0.605,4.714)rectangle++(0.04,0.04)
  (-0.328,5.627)rectangle++(0.04,0.04) (-0.056,6.392)rectangle++(0.04,0.04) (0.206,7.032)rectangle++(0.04,0.04) (0.454,7.565)rectangle++(0.04,0.04) (2.033,2.364)rectangle++(0.04,0.04)
  (1.82,3.524)rectangle++(0.04,0.04) (1.686,4.519)rectangle++(0.04,0.04) (1.612,5.37)rectangle++(0.04,0.04) (1.582,6.096)rectangle++(0.04,0.04) (1.452,2.304)rectangle++(0.04,0.04)
  (1.324,3.497)rectangle++(0.04,0.04) (1.263,4.516)rectangle++(0.04,0.04) (-0.943,2.873)rectangle++(0.04,0.04) (-0.688,4.076)rectangle++(0.04,0.04) (-0.424,5.088)rectangle++(0.04,0.04)
  (-0.159,5.938)rectangle++(0.04,0.04) (0.1,6.65)rectangle++(0.04,0.04) (0.349,7.245)rectangle++(0.04,0.04) (0.584,7.74)rectangle++(0.04,0.04) (0.804,8.152)rectangle++(0.04,0.04)
  (1.007,8.493)rectangle++(0.04,0.04) (1.194,8.775)rectangle++(0.04,0.04) (-2.064,3.794)rectangle++(0.04,0.04) (-1.605,4.904)rectangle++(0.04,0.04) (-1.17,5.829)rectangle++(0.04,0.04)
  (-0.764,6.597)rectangle++(0.04,0.04) (-0.387,7.234)rectangle++(0.04,0.04) (-0.042,7.761)rectangle++(0.04,0.04) (0.272,8.195)rectangle++(0.04,0.04) (0.557,8.551)rectangle++(0.04,0.04)
  (0.813,8.842)rectangle++(0.04,0.04) (1.043,9.079)rectangle++(0.04,0.04) (1.247,9.272)rectangle++(0.04,0.04) (1.429,9.428)rectangle++(0.04,0.04) (-2.187,3.992)rectangle++(0.04,0.04)
  (-1.701,5.077)rectangle++(0.04,0.04) (-1.245,5.979)rectangle++(0.04,0.04) (-0.821,6.728)rectangle++(0.04,0.04) (-0.431,7.348)rectangle++(0.04,0.04) (-0.075,7.859)rectangle++(0.04,0.04)
  (0.249,8.28)rectangle++(0.04,0.04) (0.54,8.624)rectangle++(0.04,0.04) (0.802,8.905)rectangle++(0.04,0.04) (1.036,9.133)rectangle++(0.04,0.04) (1.243,9.318)rectangle++(0.04,0.04)
  (-2.195,3.925)rectangle++(0.04,0.04) (-1.711,5.02)rectangle++(0.04,0.04) (-1.256,5.932)rectangle++(0.04,0.04) (-0.832,6.689)rectangle++(0.04,0.04) (-0.442,7.315)rectangle++(0.04,0.04)
  (-0.085,7.831)rectangle++(0.04,0.04) (0.238,8.256)rectangle++(0.04,0.04) (0.531,8.605)rectangle++(0.04,0.04) (0.793,8.889)rectangle++(0.04,0.04) (1.027,9.12)rectangle++(0.04,0.04)
  (-2.187,3.832)rectangle++(0.04,0.04) (-1.708,4.941)rectangle++(0.04,0.04) (-1.256,5.864)rectangle++(0.04,0.04) (1.813,1.511)rectangle++(0.04,0.04) (1.599,2.808)rectangle++(0.04,0.04)
  (1.469,3.919)rectangle++(0.04,0.04) (1.404,4.868)rectangle++(0.04,0.04) (1.386,5.678)rectangle++(0.04,0.04) (1.403,6.367)rectangle++(0.04,0.04) (1.445,6.952)rectangle++(0.04,0.04)
  (1.504,7.448)rectangle++(0.04,0.04) (1.574,7.867)rectangle++(0.04,0.04) (1.65,8.22)rectangle++(0.04,0.04) (1.729,8.517)rectangle++(0.04,0.04) (1.809,8.767)rectangle++(0.04,0.04)
  (2.166,3.004)rectangle++(0.04,0.04) (1.959,4.063)rectangle++(0.04,0.04) (1.826,4.972)rectangle++(0.04,0.04) (1.748,5.749)rectangle++(0.04,0.04) (-1.166,3.256)rectangle++(0.04,0.04)
  (-0.02,0.504)rectangle++(0.04,0.04) (0.001,2.026)rectangle++(0.04,0.04) (0.08,3.318)rectangle++(0.04,0.04) (0.198,4.413)rectangle++(0.04,0.04) (0.343,5.339)rectangle++(0.04,0.04)
  (0.503,6.121)rectangle++(0.04,0.04) (0.67,6.779)rectangle++(0.04,0.04) (0.838,7.332)rectangle++(0.04,0.04) (1.004,7.794)rectangle++(0.04,0.04) (1.163,8.181)rectangle++(0.04,0.04)
  (1.313,8.504)rectangle++(0.04,0.04) (1.454,8.772)rectangle++(0.04,0.04) (1.585,8.994)rectangle++(0.04,0.04) (1.705,9.178)rectangle++(0.04,0.04) (1.814,9.329)rectangle++(0.04,0.04)
  (2.323,3.141)rectangle++(0.04,0.04) (2.098,4.173)rectangle++(0.04,0.04) (1.948,5.059)rectangle++(0.04,0.04) (1.856,5.819)rectangle++(0.04,0.04) (1.808,6.468)rectangle++(0.04,0.04)
  (1.793,7.021)rectangle++(0.04,0.04) (1.803,7.493)rectangle++(0.04,0.04) (1.83,7.893)rectangle++(0.04,0.04) (1.869,8.232)rectangle++(0.04,0.04) (-0.02,1.3)rectangle++(0.04,0.04)
  (0.033,2.702)rectangle++(0.04,0.04) (0.134,3.892)rectangle++(0.04,0.04) (-1.006,2.476)rectangle++(0.04,0.04) (-0.759,3.741)rectangle++(0.04,0.04) (-0.497,4.806)rectangle++(0.04,0.04)
  (-0.233,5.702)rectangle++(0.04,0.04) (0.028,6.452)rectangle++(0.04,0.04) (0.28,7.079)rectangle++(0.04,0.04) (0.519,7.602)rectangle++(0.04,0.04) (0.743,8.037)rectangle++(0.04,0.04)
  (0.951,8.398)rectangle++(0.04,0.04) (1.142,8.697)rectangle++(0.04,0.04) (1.316,8.943)rectangle++(0.04,0.04) (-2.083,3.859)rectangle++(0.04,0.04) (-1.441,1.959)rectangle++(0.04,0.04)
  (-1.149,3.319)rectangle++(0.04,0.04) (-0.846,4.463)rectangle++(0.04,0.04) (-1.351,2.376)rectangle++(0.04,0.04) (-1.055,3.67)rectangle++(0.04,0.04) (-0.752,4.758)rectangle++(0.04,0.04)
  (-0.451,5.671)rectangle++(0.04,0.04) (-0.159,6.434)rectangle++(0.04,0.04) (0.12,7.072)rectangle++(0.04,0.04) (0.383,7.602)rectangle++(0.04,0.04) (0.627,8.043)rectangle++(0.04,0.04)
  (-1.987,3.503)rectangle++(0.04,0.04) (-1.551,4.653)rectangle++(0.04,0.04) (-1.134,5.613)rectangle++(0.04,0.04) (-0.742,6.413)rectangle++(0.04,0.04) (-0.376,7.077)rectangle++(0.04,0.04)
  (-0.02,1.115)rectangle++(0.04,0.04) (0.025,2.545)rectangle++(0.04,0.04) (0.121,3.759)rectangle++(0.04,0.04) (0.251,4.786)rectangle++(0.04,0.04) (0.403,5.654)rectangle++(0.04,0.04)
  (0.566,6.386)rectangle++(0.04,0.04) (0.735,7.002)rectangle++(0.04,0.04) (0.902,7.518)rectangle++(0.04,0.04) (1.952,2.469)rectangle++(0.04,0.04) (1.756,3.617)rectangle++(0.04,0.04)
  (-0.61,2.789)rectangle++(0.04,0.04) (-0.409,3.991)rectangle++(0.04,0.04) (-0.191,5.005)rectangle++(0.04,0.04) (1.413,1.582)rectangle++(0.04,0.04) (1.262,2.884)rectangle++(0.04,0.04)
  (1.186,3.997)rectangle++(0.04,0.04) (-0.823,2.741)rectangle++(0.04,0.04) (-0.592,3.959)rectangle++(0.04,0.04) (-0.347,4.985)rectangle++(0.04,0.04) (1.431,1.536)rectangle++(0.04,0.04)
  (1.275,2.845)rectangle++(0.04,0.04) (1.196,3.963)rectangle++(0.04,0.04) (1.173,4.917)rectangle++(0.04,0.04) (-1.065,2.94)rectangle++(0.04,0.04) (-0.79,4.138)rectangle++(0.04,0.04)
  (-0.508,5.145)rectangle++(0.04,0.04) (-0.228,5.99)rectangle++(0.04,0.04) (0.043,6.697)rectangle++(0.04,0.04) (0.303,7.287)rectangle++(0.04,0.04) (0.546,7.778)rectangle++(0.04,0.04)
  (0.773,8.186)rectangle++(0.04,0.04) (0.983,8.523)rectangle++(0.04,0.04) (1.174,8.801)rectangle++(0.04,0.04) (2.271,2.848)rectangle++(0.04,0.04) (2.042,3.926)rectangle++(0.04,0.04)
  (1.891,4.851)rectangle++(0.04,0.04) (1.799,5.644)rectangle++(0.04,0.04) (1.293,2.252)rectangle++(0.04,0.04) (1.187,3.459)rectangle++(0.04,0.04) (1.145,4.489)rectangle++(0.04,0.04)
  (1.151,5.366)rectangle++(0.04,0.04) (1.191,6.111)rectangle++(0.04,0.04) (1.254,6.743)rectangle++(0.04,0.04) (1.683,2.374)rectangle++(0.04,0.04) (1.523,3.547)rectangle++(0.04,0.04)
  (1.434,4.55)rectangle++(0.04,0.04) (1.399,5.407)rectangle++(0.04,0.04) (1.403,6.136)rectangle++(0.04,0.04) (1.436,6.756)rectangle++(0.04,0.04) (1.489,7.281)rectangle++(0.04,0.04)
  (1.554,7.726)rectangle++(0.04,0.04) (1.628,8.101)rectangle++(0.04,0.04) (1.706,8.417)rectangle++(0.04,0.04) (1.784,8.682)rectangle++(0.04,0.04) (2.146,2.978)rectangle++(0.04,0.04)
  (1.941,4.041)rectangle++(0.04,0.04) (-0.684,2.925)rectangle++(0.04,0.04) (-0.02,0.451)rectangle++(0.04,0.04) (-0.001,1.98)rectangle++(0.04,0.04) (0.076,3.28)rectangle++(0.04,0.04)
  (0.194,4.381)rectangle++(0.04,0.04) (0.338,5.312)rectangle++(0.04,0.04) (0.497,6.098)rectangle++(0.04,0.04) (0.664,6.76)rectangle++(0.04,0.04) (0.833,7.315)rectangle++(0.04,0.04)
  (0.998,7.781)rectangle++(0.04,0.04) (1.158,8.17)rectangle++(0.04,0.04) (1.309,8.494)rectangle++(0.04,0.04) (1.45,8.764)rectangle++(0.04,0.04) (1.581,8.988)rectangle++(0.04,0.04)
  (2.262,2.998)rectangle++(0.04,0.04) (2.04,4.054)rectangle++(0.04,0.04) (1.894,4.961)rectangle++(0.04,0.04) (1.806,5.737)rectangle++(0.04,0.04) (1.763,6.4)rectangle++(0.04,0.04)
  (1.752,6.966)rectangle++(0.04,0.04) (1.766,7.447)rectangle++(0.04,0.04) (1.797,7.856)rectangle++(0.04,0.04) (1.839,8.202)rectangle++(0.04,0.04) (1.889,8.494)rectangle++(0.04,0.04)
  (-1.852,3.892)rectangle++(0.04,0.04) (-1.421,4.979)rectangle++(0.04,0.04) (-1.011,5.885)rectangle++(0.04,0.04) (-0.626,6.639)rectangle++(0.04,0.04) (-0.269,7.264)rectangle++(0.04,0.04)
  (0.06,7.781)rectangle++(0.04,0.04) (0.36,8.208)rectangle++(0.04,0.04) (0.632,8.559)rectangle++(0.04,0.04) (0.878,8.846)rectangle++(0.04,0.04) (1.098,9.08)rectangle++(0.04,0.04)
  (1.294,9.27)rectangle++(0.04,0.04) (1.468,9.424)rectangle++(0.04,0.04) (1.623,9.548)rectangle++(0.04,0.04) (1.759,9.647)rectangle++(0.04,0.04) (1.879,9.726)rectangle++(0.04,0.04)
  (1.984,9.788)rectangle++(0.04,0.04) (2.076,9.837)rectangle++(0.04,0.04) (2.156,9.874)rectangle++(0.04,0.04) (2.225,9.903)rectangle++(0.04,0.04) (2.285,9.925)rectangle++(0.04,0.04)
  (2.337,9.941)rectangle++(0.04,0.04) (2.415,3.424)rectangle++(0.04,0.04) (2.188,4.41)rectangle++(0.04,0.04) (-0.02,0.689)rectangle++(0.04,0.04) (0.008,2.182)rectangle++(0.04,0.04)
  (0.092,3.451)rectangle++(0.04,0.04) (0.214,4.526)rectangle++(0.04,0.04) (0.361,5.435)rectangle++(0.04,0.04) (-1.362,2.868)rectangle++(0.04,0.04) (-1.045,4.088)rectangle++(0.04,0.04)
  (-0.727,5.113)rectangle++(0.04,0.04) (-0.416,5.971)rectangle++(0.04,0.04) (-0.117,6.688)rectangle++(0.04,0.04) (0.166,7.286)rectangle++(0.04,0.04) (1.998,2.222)rectangle++(0.04,0.04)
  (0.305,1.483)rectangle++(0.04,0.04) (0.316,2.844)rectangle++(0.04,0.04) (0.381,4.001)rectangle++(0.04,0.04) (0.481,4.982)rectangle++(0.04,0.04) (0.606,5.812)rectangle++(0.04,0.04)
  (0.746,6.512)rectangle++(0.04,0.04) (0.892,7.101)rectangle++(0.04,0.04) (1.04,7.597)rectangle++(0.04,0.04) (1.186,8.012)rectangle++(0.04,0.04) (1.326,8.359)rectangle++(0.04,0.04)
  (1.459,8.648)rectangle++(0.04,0.04) (1.584,8.889)rectangle++(0.04,0.04) (1.7,9.088)rectangle++(0.04,0.04) (1.806,9.253)rectangle++(0.04,0.04) (1.903,9.389)rectangle++(0.04,0.04)
  (1.991,9.501)rectangle++(0.04,0.04) (2.07,9.592)rectangle++(0.04,0.04) (2.141,9.667)rectangle++(0.04,0.04) (2.205,9.727)rectangle++(0.04,0.04) (2.261,9.776)rectangle++(0.04,0.04)
  (-2.111,4.26)rectangle++(0.04,0.04) (-1.626,5.301)rectangle++(0.04,0.04) (1.711,1.28)rectangle++(0.04,0.04) (1.503,2.615)rectangle++(0.04,0.04) (1.38,3.759)rectangle++(0.04,0.04)
  (1.321,4.736)rectangle++(0.04,0.04) (1.31,5.569)rectangle++(0.04,0.04) (1.334,6.278)rectangle++(0.04,0.04) (1.383,6.879)rectangle++(0.04,0.04) (1.449,7.388)rectangle++(0.04,0.04)
  (1.525,7.818)rectangle++(0.04,0.04) (-1.749,3.66)rectangle++(0.04,0.04) (-1.342,4.777)rectangle++(0.04,0.04) (-0.952,5.71)rectangle++(0.04,0.04) (-0.583,6.488)rectangle++(0.04,0.04)
  (-0.238,7.134)rectangle++(0.04,0.04) (0.081,7.67)rectangle++(0.04,0.04) (2.118,2.292)rectangle++(0.04,0.04) (1.89,3.459)rectangle++(0.04,0.04) (1.743,4.461)rectangle++(0.04,0.04)
  (1.657,5.318)rectangle++(0.04,0.04) (1.619,6.051)rectangle++(0.04,0.04) (1.616,6.674)rectangle++(0.04,0.04) (1.639,7.205)rectangle++(0.04,0.04) (-1.567,3.551)rectangle++(0.04,0.04)
  (-1.192,4.677)rectangle++(0.04,0.04) (-0.828,5.619)rectangle++(0.04,0.04) (-0.481,6.406)rectangle++(0.04,0.04) (-0.155,7.06)rectangle++(0.04,0.04) (0.148,7.604)rectangle++(0.04,0.04)
  (0.428,8.053)rectangle++(0.04,0.04) (2.173,2.474)rectangle++(0.04,0.04) (1.944,3.612)rectangle++(0.04,0.04) (1.795,4.589)rectangle++(0.04,0.04) (1.707,5.425)rectangle++(0.04,0.04)
  (1.666,6.139)rectangle++(0.04,0.04) (1.659,6.748)rectangle++(0.04,0.04) (1.678,7.265)rectangle++(0.04,0.04) (1.715,7.705)rectangle++(0.04,0.04) (1.764,8.077)rectangle++(0.04,0.04)
  (1.82,8.391)rectangle++(0.04,0.04) (1.88,8.656)rectangle++(0.04,0.04) (1.942,8.878)rectangle++(0.04,0.04) (2.004,9.065)rectangle++(0.04,0.04) (2.064,9.221)rectangle++(0.04,0.04)
  (-2.006,4.092)rectangle++(0.04,0.04) (-1.544,5.155)rectangle++(0.04,0.04) (-1.108,6.04)rectangle++(0.04,0.04) (-0.702,6.774)rectangle++(0.04,0.04) (-0.328,7.382)rectangle++(0.04,0.04)
  (0.014,7.884)rectangle++(0.04,0.04) (0.325,8.297)rectangle++(0.04,0.04) (0.606,8.636)rectangle++(0.04,0.04) (0.858,8.913)rectangle++(0.04,0.04) (1.084,9.137)rectangle++(0.04,0.04)
  (1.285,9.32)rectangle++(0.04,0.04) (1.463,9.467)rectangle++(0.04,0.04) (1.62,9.584)rectangle++(0.04,0.04) (1.758,9.678)rectangle++(0.04,0.04) (1.879,9.752)rectangle++(0.04,0.04)
  (-2.181,4.167)rectangle++(0.04,0.04) (-1.689,5.225)rectangle++(0.04,0.04) (-1.229,6.105)rectangle++(0.04,0.04) (-0.803,6.835)rectangle++(0.04,0.04) (-0.411,7.438)rectangle++(0.04,0.04)
  (-0.054,7.935)rectangle++(0.04,0.04) (0.269,8.343)rectangle++(0.04,0.04) (0.56,8.677)rectangle++(0.04,0.04) (0.821,8.949)rectangle++(0.04,0.04) (1.054,9.17)rectangle++(0.04,0.04)
  (1.26,9.349)rectangle++(0.04,0.04) (1.443,9.492)rectangle++(0.04,0.04) (1.604,9.607)rectangle++(0.04,0.04) (1.745,9.698)rectangle++(0.04,0.04) (1.869,9.77)rectangle++(0.04,0.04)
  (1.978,9.826)rectangle++(0.04,0.04) (2.072,9.869)rectangle++(0.04,0.04) (2.154,9.902)rectangle++(0.04,0.04) (2.224,9.927)rectangle++(0.04,0.04) (2.286,9.945)rectangle++(0.04,0.04)
  (-2.15,4.303)rectangle++(0.04,0.04) (-1.657,5.339)rectangle++(0.04,0.04) (-1.197,6.201)rectangle++(0.04,0.04) (-0.772,6.915)rectangle++(0.04,0.04) (-0.382,7.505)rectangle++(0.04,0.04)
  (-0.027,7.991)rectangle++(0.04,0.04) (0.295,8.389)rectangle++(0.04,0.04) (0.584,8.715)rectangle++(0.04,0.04) (0.843,8.981)rectangle++(0.04,0.04) (1.073,9.196)rectangle++(0.04,0.04)
  (1.278,9.37)rectangle++(0.04,0.04) (1.459,9.51)rectangle++(0.04,0.04) (-0.02,1.505)rectangle++(0.04,0.04) (0.041,2.876)rectangle++(0.04,0.04) (0.148,4.039)rectangle++(0.04,0.04)
  (0.285,5.024)rectangle++(0.04,0.04) (0.441,5.855)rectangle++(0.04,0.04) (0.607,6.555)rectangle++(0.04,0.04) (0.776,7.144)rectangle++(0.04,0.04) (0.943,7.637)rectangle++(0.04,0.04)
  (1.105,8.05)rectangle++(0.04,0.04) (1.259,8.395)rectangle++(0.04,0.04) (1.404,8.681)rectangle++(0.04,0.04) (1.538,8.919)rectangle++(0.04,0.04) (1.662,9.116)rectangle++(0.04,0.04)
  (1.775,9.278)rectangle++(0.04,0.04) (1.878,9.412)rectangle++(0.04,0.04) (-2.093,4.091)rectangle++(0.04,0.04) (-1.617,5.158)rectangle++(0.04,0.04) (-1.171,6.045)rectangle++(0.04,0.04)
  (-0.755,6.781)rectangle++(0.04,0.04) (-0.373,7.39)rectangle++(0.04,0.04) (-0.024,7.893)rectangle++(0.04,0.04) (0.293,8.306)rectangle++(0.04,0.04) (0.579,8.645)rectangle++(0.04,0.04)
  (0.836,8.921)rectangle++(0.04,0.04) (1.065,9.146)rectangle++(0.04,0.04) (1.269,9.327)rectangle++(0.04,0.04) (1.45,9.474)rectangle++(0.04,0.04) (1.609,9.591)rectangle++(0.04,0.04)
  (1.749,9.684)rectangle++(0.04,0.04) (1.872,9.758)rectangle++(0.04,0.04) (1.979,9.815)rectangle++(0.04,0.04) (2.073,9.86)rectangle++(0.04,0.04) (2.433,3.335)rectangle++(0.04,0.04)
  (2.199,4.334)rectangle++(0.04,0.04) (2.04,5.192)rectangle++(0.04,0.04) (1.94,5.928)rectangle++(0.04,0.04) (1.884,6.557)rectangle++(0.04,0.04) (1.861,7.095)rectangle++(0.04,0.04)
  (1.864,7.552)rectangle++(0.04,0.04) (1.884,7.941)rectangle++(0.04,0.04) (1.917,8.271)rectangle++(0.04,0.04) (-1.788,3.849)rectangle++(0.04,0.04) (-1.368,4.94)rectangle++(0.04,0.04)
  (-0.968,5.85)rectangle++(0.04,0.04) (-0.591,6.607)rectangle++(0.04,0.04) (-0.24,7.236)rectangle++(0.04,0.04) (0.083,7.756)rectangle++(0.04,0.04) (0.379,8.186)rectangle++(0.04,0.04)
  (0.647,8.539)rectangle++(0.04,0.04) (0.889,8.828)rectangle++(0.04,0.04) (1.107,9.065)rectangle++(0.04,0.04) (1.301,9.257)rectangle++(0.04,0.04) (1.474,9.413)rectangle++(0.04,0.04)
  (1.627,9.538)rectangle++(0.04,0.04) (1.763,9.638)rectangle++(0.04,0.04) (1.882,9.718)rectangle++(0.04,0.04) (1.986,9.781)rectangle++(0.04,0.04) (2.077,9.831)rectangle++(0.04,0.04)
  (-0.02,1.556)rectangle++(0.04,0.04) (-0.43,1.927)rectangle++(0.04,0.04) (-0.29,3.251)rectangle++(0.04,0.04) (0.936,1.135)rectangle++(0.04,0.04) (0.839,2.523)rectangle++(0.04,0.04)
  (-0.509,2.337)rectangle++(0.04,0.04) (-0.342,3.603)rectangle++(0.04,0.04) (-0.149,4.673)rectangle++(0.04,0.04) (0.058,5.574)rectangle++(0.04,0.04) (0.27,6.332)rectangle++(0.04,0.04)
  (1.715,2.02)rectangle++(0.04,0.04) (1.536,3.245)rectangle++(0.04,0.04) (1.433,4.293)rectangle++(0.04,0.04) (1.388,5.188)rectangle++(0.04,0.04) (1.385,5.95)rectangle++(0.04,0.04)
  (1.413,6.598)rectangle++(0.04,0.04) (1.618,2.381)rectangle++(0.04,0.04) (1.469,3.555)rectangle++(0.04,0.04) (1.388,4.559)rectangle++(0.04,0.04) (1.36,5.416)rectangle++(0.04,0.04)
  (1.371,6.146)rectangle++(0.04,0.04) (-1.345,3.256)rectangle++(0.04,0.04) (-1.015,4.418)rectangle++(0.04,0.04) (-0.688,5.392)rectangle++(0.04,0.04) (-0.372,6.207)rectangle++(0.04,0.04)
  (-0.07,6.887)rectangle++(0.04,0.04) (0.214,7.453)rectangle++(0.04,0.04) (0.478,7.923)rectangle++(0.04,0.04) (0.721,8.311)rectangle++(0.04,0.04) (0.943,8.632)rectangle++(0.04,0.04)
  (1.145,8.896)rectangle++(0.04,0.04) (1.327,9.112)rectangle++(0.04,0.04) (1.49,9.288)rectangle++(0.04,0.04) (1.636,9.431)rectangle++(0.04,0.04) (1.765,9.548)rectangle++(0.04,0.04)
  (1.88,9.641)rectangle++(0.04,0.04) (1.982,9.716)rectangle++(0.04,0.04) (2.071,9.775)rectangle++(0.04,0.04) (2.149,9.822)rectangle++(0.04,0.04) (2.217,9.859)rectangle++(0.04,0.04)
  (2.277,9.888)rectangle++(0.04,0.04) (2.329,9.91)rectangle++(0.04,0.04) (2.374,9.926)rectangle++(0.04,0.04) (2.412,9.939)rectangle++(0.04,0.04) (2.446,9.948)rectangle++(0.04,0.04)
  (-2.118,4.34)rectangle++(0.04,0.04) (-1.629,5.37)rectangle++(0.04,0.04) (-1.172,6.226)rectangle++(0.04,0.04) (1.902,1.619)rectangle++(0.04,0.04) (1.679,2.897)rectangle++(0.04,0.04)
  (1.541,3.991)rectangle++(0.04,0.04) (-0.751,2.821)rectangle++(0.04,0.04) (-0.527,4.024)rectangle++(0.04,0.04) (-0.29,5.038)rectangle++(0.04,0.04) (-0.047,5.89)rectangle++(0.04,0.04)
  (0.194,6.605)rectangle++(0.04,0.04) (0.427,7.202)rectangle++(0.04,0.04) (0.648,7.701)rectangle++(0.04,0.04) (0.857,8.116)rectangle++(0.04,0.04) (1.051,8.461)rectangle++(0.04,0.04)
  (1.23,8.746)rectangle++(0.04,0.04) (1.393,8.981)rectangle++(0.04,0.04) (1.541,9.174)rectangle++(0.04,0.04) (1.675,9.333)rectangle++(0.04,0.04) (1.794,9.462)rectangle++(0.04,0.04)
  (1.902,9.567)rectangle++(0.04,0.04) (1.997,9.652)rectangle++(0.04,0.04) (2.081,9.721)rectangle++(0.04,0.04) (2.156,9.776)rectangle++(0.04,0.04) (2.221,9.819)rectangle++(0.04,0.04)
  (2.279,9.854)rectangle++(0.04,0.04) (2.329,9.881)rectangle++(0.04,0.04) (2.4,3.407)rectangle++(0.04,0.04) (2.174,4.396)rectangle++(0.04,0.04) (2.022,5.246)rectangle++(0.04,0.04)
  (1.926,5.974)rectangle++(0.04,0.04) (1.874,6.597)rectangle++(0.04,0.04) (1.854,7.129)rectangle++(0.04,0.04) (1.859,7.582)rectangle++(0.04,0.04) (1.881,7.966)rectangle++(0.04,0.04)
  (1.916,8.292)rectangle++(0.04,0.04) (1.958,8.568)rectangle++(0.04,0.04) (2.005,8.801)rectangle++(0.04,0.04) (2.054,8.997)rectangle++(0.04,0.04) (-1.95,4.041)rectangle++(0.04,0.04)
  (-1.498,5.109)rectangle++(0.04,0.04) (-1.071,5.999)rectangle++(0.04,0.04) (-0.673,6.738)rectangle++(0.04,0.04) (-0.304,7.35)rectangle++(0.04,0.04) (0.033,7.856)rectangle++(0.04,0.04)
  (0.34,8.273)rectangle++(0.04,0.04) (0.618,8.614)rectangle++(0.04,0.04) (0.868,8.894)rectangle++(0.04,0.04) (1.091,9.121)rectangle++(0.04,0.04) (-2.174,3.847)rectangle++(0.04,0.04)
  (-1.697,4.953)rectangle++(0.04,0.04) (-1.246,5.874)rectangle++(0.04,0.04) (-0.827,6.639)rectangle++(0.04,0.04) (-0.439,7.272)rectangle++(0.04,0.04) (-0.085,7.795)rectangle++(0.04,0.04)
  (0.238,8.226)rectangle++(0.04,0.04) (0.529,8.578)rectangle++(0.04,0.04) (0.79,8.867)rectangle++(0.04,0.04) (1.024,9.101)rectangle++(0.04,0.04) (1.233,9.291)rectangle++(0.04,0.04)
  (1.417,9.444)rectangle++(0.04,0.04) (1.58,9.567)rectangle++(0.04,0.04) (1.724,9.665)rectangle++(0.04,0.04) (1.849,9.743)rectangle++(0.04,0.04) (2.433,3.249)rectangle++(0.04,0.04)
  (-0.379,2.863)rectangle++(0.04,0.04) (-0.21,4.045)rectangle++(0.04,0.04) (-0.019,5.043)rectangle++(0.04,0.04) (-0.02,0.79)rectangle++(0.04,0.04) (0.012,2.269)rectangle++(0.04,0.04)
  (-0.609,2.091)rectangle++(0.04,0.04) (-0.436,3.398)rectangle++(0.04,0.04) (-0.237,4.502)rectangle++(0.04,0.04) (-0.023,5.432)rectangle++(0.04,0.04) (0.195,6.215)rectangle++(0.04,0.04)
  (0.412,6.871)rectangle++(0.04,0.04) (0.623,7.42)rectangle++(0.04,0.04) (0.824,7.878)rectangle++(0.04,0.04) (1.014,8.26)rectangle++(0.04,0.04) (1.19,8.576)rectangle++(0.04,0.04)
  (1.352,8.839)rectangle++(0.04,0.04) (1.501,9.055)rectangle++(0.04,0.04) (1.636,9.233)rectangle++(0.04,0.04) (1.757,9.379)rectangle++(0.04,0.04) (1.867,9.498)rectangle++(0.04,0.04)
  (1.964,9.595)rectangle++(0.04,0.04) (2.051,9.673)rectangle++(0.04,0.04) (2.128,9.736)rectangle++(0.04,0.04) (2.196,9.787)rectangle++(0.04,0.04) (2.256,9.827)rectangle++(0.04,0.04)
  (2.309,9.859)rectangle++(0.04,0.04) (2.354,9.884)rectangle++(0.04,0.04) (2.397,3.414)rectangle++(0.04,0.04) (2.172,4.403)rectangle++(0.04,0.04) (2.02,5.251)rectangle++(0.04,0.04)
  (-0.983,3.209)rectangle++(0.04,0.04) (-0.709,4.363)rectangle++(0.04,0.04) (-0.43,5.333)rectangle++(0.04,0.04) (-0.155,6.147)rectangle++(0.04,0.04) (0.112,6.827)rectangle++(0.04,0.04)
  (0.366,7.395)rectangle++(0.04,0.04) (0.605,7.867)rectangle++(0.04,0.04) (0.827,8.259)rectangle++(0.04,0.04) (1.031,8.583)rectangle++(0.04,0.04) (-2.047,3.714)rectangle++(0.04,0.04)
  (-1.593,4.835)rectangle++(0.04,0.04) (-1.163,5.77)rectangle++(0.04,0.04) (-0.76,6.547)rectangle++(0.04,0.04) (-0.386,7.192)rectangle++(0.04,0.04) (-0.043,7.725)rectangle++(0.04,0.04)
  (0.27,8.164)rectangle++(0.04,0.04) (0.554,8.525)rectangle++(0.04,0.04) (0.81,8.82)rectangle++(0.04,0.04) (1.039,9.061)rectangle++(0.04,0.04) (1.243,9.256)rectangle++(0.04,0.04)
  (1.425,9.414)rectangle++(0.04,0.04) (1.586,9.541)rectangle++(0.04,0.04) (1.727,9.643)rectangle++(0.04,0.04) (1.852,9.724)rectangle++(0.04,0.04) (1.961,9.787)rectangle++(0.04,0.04)
  (2.056,9.837)rectangle++(0.04,0.04) (2.139,9.875)rectangle++(0.04,0.04) (2.211,9.905)rectangle++(0.04,0.04) (2.273,9.927)rectangle++(0.04,0.04) (2.327,9.943)rectangle++(0.04,0.04)
  (2.373,9.955)rectangle++(0.04,0.04) (-2.135,4.325)rectangle++(0.04,0.04) (-1.644,5.358)rectangle++(0.04,0.04) (-1.185,6.216)rectangle++(0.04,0.04) (-0.761,6.927)rectangle++(0.04,0.04)
  (-0.372,7.515)rectangle++(0.04,0.04) (-0.018,7.999)rectangle++(0.04,0.04) (0.303,8.396)rectangle++(0.04,0.04) (0.591,8.72)rectangle++(0.04,0.04) (0.849,8.985)rectangle++(0.04,0.04)
  (1.079,9.199)rectangle++(0.04,0.04) (1.283,9.373)rectangle++(0.04,0.04) (1.463,9.512)rectangle++(0.04,0.04) (1.622,9.623)rectangle++(0.04,0.04) (1.761,9.71)rectangle++(0.04,0.04)
  (1.883,9.78)rectangle++(0.04,0.04) (1.99,9.834)rectangle++(0.04,0.04) (2.082,9.875)rectangle++(0.04,0.04) (2.163,9.907)rectangle++(0.04,0.04) (2.233,9.93)rectangle++(0.04,0.04)
  (2.293,9.948)rectangle++(0.04,0.04) (2.344,9.96)rectangle++(0.04,0.04) (2.389,9.969)rectangle++(0.04,0.04) (2.427,9.974)rectangle++(0.04,0.04) (2.46,9.977)rectangle++(0.04,0.04)
  (2.488,9.978)rectangle++(0.04,0.04) (-2.118,4.356)rectangle++(0.04,0.04) (-1.628,5.384)rectangle++(0.04,0.04) (-1.171,6.238)rectangle++(0.04,0.04) (-0.748,6.945)rectangle++(0.04,0.04)
  (-0.36,7.529)rectangle++(0.04,0.04) (-0.007,8.011)rectangle++(0.04,0.04) (0.312,8.405)rectangle++(0.04,0.04) (0.599,8.728)rectangle++(0.04,0.04) (0.856,8.991)rectangle++(0.04,0.04)
  (1.085,9.205)rectangle++(0.04,0.04) (2.397,2.921)rectangle++(0.04,0.04) (2.152,3.983)rectangle++(0.04,0.04) (-0.626,2.96)rectangle++(0.04,0.04) (-0.416,4.138)rectangle++(0.04,0.04)
  (-0.191,5.13)rectangle++(0.04,0.04) (0.041,5.964)rectangle++(0.04,0.04) (0.271,6.664)rectangle++(0.04,0.04) (1.808,2.1)rectangle++(0.04,0.04) (1.618,3.309)rectangle++(0.04,0.04)
  (1.506,4.344)rectangle++(0.04,0.04) (1.452,5.228)rectangle++(0.04,0.04) (1.441,5.982)rectangle++(0.04,0.04) (1.462,6.623)rectangle++(0.04,0.04) (1.505,7.168)rectangle++(0.04,0.04)
  (1.564,7.629)rectangle++(0.04,0.04) (1.632,8.018)rectangle++(0.04,0.04) (1.983,2.779)rectangle++(0.04,0.04) (1.794,3.879)rectangle++(0.04,0.04) (1.678,4.821)rectangle++(0.04,0.04)
  (1.617,5.627)rectangle++(0.04,0.04) (1.597,6.315)rectangle++(0.04,0.04) (1.608,6.9)rectangle++(0.04,0.04) (1.641,7.397)rectangle++(0.04,0.04) (1.688,7.818)rectangle++(0.04,0.04)
  (1.746,8.174)rectangle++(0.04,0.04) (-1.797,3.789)rectangle++(0.04,0.04) (-1.378,4.888)rectangle++(0.04,0.04) (-0.978,5.807)rectangle++(0.04,0.04) (-0.601,6.571)rectangle++(0.04,0.04)
  (-0.251,7.206)rectangle++(0.04,0.04) (-1.945,3.117)rectangle++(0.04,0.04) (-1.531,4.323)rectangle++(0.04,0.04) (-1.13,5.332)rectangle++(0.04,0.04) (-0.75,6.174)rectangle++(0.04,0.04)
  (-0.392,6.874)rectangle++(0.04,0.04) (-0.061,7.455)rectangle++(0.04,0.04) (-1.972,3.215)rectangle++(0.04,0.04) (-1.549,4.408)rectangle++(0.04,0.04) (-1.143,5.405)rectangle++(0.04,0.04)
  (-0.757,6.236)rectangle++(0.04,0.04) (-0.397,6.927)rectangle++(0.04,0.04) (-1.902,3.022)rectangle++(0.04,0.04) (-1.498,4.241)rectangle++(0.04,0.04) (-1.106,5.261)rectangle++(0.04,0.04)
  (-0.732,6.112)rectangle++(0.04,0.04) (-0.38,6.821)rectangle++(0.04,0.04) (-0.052,7.409)rectangle++(0.04,0.04) (2.065,2.195)rectangle++(0.04,0.04) (1.841,3.379)rectangle++(0.04,0.04)
  (1.698,4.395)rectangle++(0.04,0.04) (1.617,5.264)rectangle++(0.04,0.04) (1.582,6.006)rectangle++(0.04,0.04) (1.583,6.638)rectangle++(0.04,0.04) (1.609,7.175)rectangle++(0.04,0.04)
  (1.652,7.631)rectangle++(0.04,0.04) (1.708,8.016)rectangle++(0.04,0.04) (1.971,2.798)rectangle++(0.04,0.04) (0.47,1.614)rectangle++(0.04,0.04) (0.462,2.949)rectangle++(0.04,0.04)
  (0.509,4.085)rectangle++(0.04,0.04) (0.593,5.048)rectangle++(0.04,0.04) (0.704,5.863)rectangle++(0.04,0.04) (0.831,6.552)rectangle++(0.04,0.04) (0.966,7.132)rectangle++(0.04,0.04)
  (-1.682,3.38)rectangle++(0.04,0.04) (-1.297,4.537)rectangle++(0.04,0.04) (-0.923,5.504)rectangle++(0.04,0.04) (-0.567,6.312)rectangle++(0.04,0.04) (1.835,1.797)rectangle++(0.04,0.04)
  (1.629,3.051)rectangle++(0.04,0.04) (1.505,4.124)rectangle++(0.04,0.04) (1.442,5.041)rectangle++(0.04,0.04) (1.425,5.824)rectangle++(0.04,0.04) (1.4,2.198)rectangle++(0.04,0.04)
  (1.275,3.409)rectangle++(0.04,0.04) (1.218,4.443)rectangle++(0.04,0.04) (1.211,5.324)rectangle++(0.04,0.04) (1.24,6.073)rectangle++(0.04,0.04) (-1.352,3.21)rectangle++(0.04,0.04)
  (-1.023,4.379)rectangle++(0.04,0.04) (-0.697,5.359)rectangle++(0.04,0.04) (-0.38,6.179)rectangle++(0.04,0.04) (1.77,1.814)rectangle++(0.04,0.04) (1.575,3.068)rectangle++(0.04,0.04)
  (1.459,4.141)rectangle++(0.04,0.04) (1.404,5.057)rectangle++(0.04,0.04) (1.393,5.839)rectangle++(0.04,0.04) (1.416,6.503)rectangle++(0.04,0.04) (1.461,7.068)rectangle++(0.04,0.04)
  (1.522,7.545)rectangle++(0.04,0.04) (1.594,7.949)rectangle++(0.04,0.04) (1.67,8.289)rectangle++(0.04,0.04) (1.749,8.575)rectangle++(0.04,0.04) (-1.901,3.878)rectangle++(0.04,0.04)
  (-1.463,4.968)rectangle++(0.04,0.04) (-1.047,5.878)rectangle++(0.04,0.04) (-0.657,6.634)rectangle++(0.04,0.04) (-0.295,7.262)rectangle++(0.04,0.04) (0.037,7.78)rectangle++(0.04,0.04)
  (0.341,8.208)rectangle++(0.04,0.04) (0.616,8.559)rectangle++(0.04,0.04) (-2.123,3.614)rectangle++(0.04,0.04) (-1.663,4.753)rectangle++(0.04,0.04) (-1.225,5.703)rectangle++(0.04,0.04)
  (-0.816,6.492)rectangle++(0.04,0.04) (1.923,1.776)rectangle++(0.04,0.04) (1.703,3.029)rectangle++(0.04,0.04) (1.567,4.103)rectangle++(0.04,0.04) (1.494,5.021)rectangle++(0.04,0.04)
  (1.468,5.804)rectangle++(0.04,0.04) (1.388,2.205)rectangle++(0.04,0.04) (-0.02,0.336)rectangle++(0.04,0.04) (-0.006,1.883)rectangle++(0.04,0.04) (0.068,3.197)rectangle++(0.04,0.04)
  (0.184,4.311)rectangle++(0.04,0.04) (0.326,5.253)rectangle++(0.04,0.04) (0.485,6.048)rectangle++(0.04,0.04) (0.652,6.718)rectangle++(0.04,0.04) (1.766,2.212)rectangle++(0.04,0.04)
  (1.587,3.406)rectangle++(0.04,0.04) (1.483,4.427)rectangle++(0.04,0.04) (-0.876,2.904)rectangle++(0.04,0.04) (-0.63,4.1)rectangle++(0.04,0.04) (-0.374,5.106)rectangle++(0.04,0.04)
  (-0.116,5.951)rectangle++(0.04,0.04) (0.137,6.66)rectangle++(0.04,0.04) (0.381,7.251)rectangle++(0.04,0.04) (0.612,7.745)rectangle++(0.04,0.04) (0.827,8.155)rectangle++(0.04,0.04)
  (1.027,8.495)rectangle++(0.04,0.04) (1.211,8.775)rectangle++(0.04,0.04) (1.378,9.007)rectangle++(0.04,0.04) (1.529,9.197)rectangle++(0.04,0.04) (1.666,9.352)rectangle++(0.04,0.04)
  (1.788,9.479)rectangle++(0.04,0.04) (1.897,9.582)rectangle++(0.04,0.04) (1.993,9.665)rectangle++(0.04,0.04) (2.079,9.732)rectangle++(0.04,0.04) (2.154,9.785)rectangle++(0.04,0.04)
  (2.22,9.827)rectangle++(0.04,0.04) (2.278,9.861)rectangle++(0.04,0.04) (2.328,9.887)rectangle++(0.04,0.04) (2.372,9.907)rectangle++(0.04,0.04) (2.411,9.922)rectangle++(0.04,0.04)
  (2.444,9.933)rectangle++(0.04,0.04) (2.472,9.942)rectangle++(0.04,0.04) (2.497,9.948)rectangle++(0.04,0.04) (2.518,9.952)rectangle++(0.04,0.04) (2.536,9.955)rectangle++(0.04,0.04)
  (2.552,9.956)rectangle++(0.04,0.04) (2.388,3.483)rectangle++(0.04,0.04) (2.167,4.461)rectangle++(0.04,0.04) (0.907,2.064)rectangle++(0.04,0.04) (0.425,1.161)rectangle++(0.04,0.04)
  (0.405,2.566)rectangle++(0.04,0.04) (-0.607,2.247)rectangle++(0.04,0.04) (-0.429,3.53)rectangle++(0.04,0.04) (1.035,1.166)rectangle++(0.04,0.04) (0.924,2.546)rectangle++(0.04,0.04)
  (0.885,3.723)rectangle++(0.04,0.04) (0.899,4.725)rectangle++(0.04,0.04) (0.951,5.577)rectangle++(0.04,0.04) (1.029,6.298)rectangle++(0.04,0.04) (1.592,2.209)rectangle++(0.04,0.04)
  (1.439,3.41)rectangle++(0.04,0.04) (1.358,4.437)rectangle++(0.04,0.04) (1.329,5.314)rectangle++(0.04,0.04) (1.34,6.06)rectangle++(0.04,0.04) (1.379,6.693)rectangle++(0.04,0.04)
  (1.438,7.23)rectangle++(0.04,0.04) (1.509,7.685)rectangle++(0.04,0.04) (-1.717,3.627)rectangle++(0.04,0.04) (-1.317,4.748)rectangle++(0.04,0.04) (-0.932,5.684)rectangle++(0.04,0.04)
  (-0.567,6.465)rectangle++(0.04,0.04) (1.878,1.834)rectangle++(0.04,0.04) (0.215,1.358)rectangle++(0.04,0.04) (0.234,2.742)rectangle++(0.04,0.04) (0.307,3.918)rectangle++(0.04,0.04)
  (-0.978,2.521)rectangle++(0.04,0.04) (-0.872,1.919)rectangle++(0.04,0.04) (0.651,0.664)rectangle++(0.04,0.04) (0.577,2.134)rectangle++(0.04,0.04) (0.574,3.387)rectangle++(0.04,0.04)
  (0.621,4.452)rectangle++(0.04,0.04) (0.704,5.356)rectangle++(0.04,0.04) (0.81,6.121)rectangle++(0.04,0.04) (0.931,6.766)rectangle++(0.04,0.04) (1.06,7.31)rectangle++(0.04,0.04)
  (1.191,7.768)rectangle++(0.04,0.04) (1.321,8.151)rectangle++(0.04,0.04) (-1.876,3.686)rectangle++(0.04,0.04) (-1.45,4.804)rectangle++(0.04,0.04) (-1.042,5.738)rectangle++(0.04,0.04)
  (-0.658,6.515)rectangle++(0.04,0.04) (-0.301,7.16)rectangle++(0.04,0.04) (0.028,7.695)rectangle++(0.04,0.04) (0.329,8.135)rectangle++(0.04,0.04) (0.603,8.498)rectangle++(0.04,0.04)
  (0.851,8.796)rectangle++(0.04,0.04) (1.073,9.038)rectangle++(0.04,0.04) (1.271,9.236)rectangle++(0.04,0.04) (1.448,9.396)rectangle++(0.04,0.04) (1.604,9.525)rectangle++(0.04,0.04)
  (1.742,9.628)rectangle++(0.04,0.04) (1.864,9.71)rectangle++(0.04,0.04) (1.97,9.775)rectangle++(0.04,0.04) (2.064,9.827)rectangle++(0.04,0.04) (2.145,9.866)rectangle++(0.04,0.04)
  (2.216,9.897)rectangle++(0.04,0.04) (-2.151,4.276)rectangle++(0.04,0.04) (-1.66,5.317)rectangle++(0.04,0.04) (-1.735,2.377)rectangle++(0.04,0.04) (0.908,0.549)rectangle++(0.04,0.04)
  (0.018,1.919)rectangle++(0.04,0.04) (0.09,3.226)rectangle++(0.04,0.04) (0.203,4.335)rectangle++(0.04,0.04) (0.344,5.273)rectangle++(0.04,0.04) (0.501,6.064)rectangle++(0.04,0.04)
  (0.666,6.731)rectangle++(0.04,0.04) (0.833,7.291)rectangle++(0.04,0.04) (0.998,7.76)rectangle++(0.04,0.04) (1.156,8.152)rectangle++(0.04,0.04) (2.092,2.687)rectangle++(0.04,0.04)
  (0.421,1.619)rectangle++(0.04,0.04) (0.421,2.955)rectangle++(0.04,0.04) (-0.02,0.456)rectangle++(0.04,0.04) (-0.144,1.685)rectangle++(0.04,0.04) (0.476,0.797)rectangle++(0.04,0.04)
  (0.434,2.255)rectangle++(0.04,0.04) (0.457,3.495)rectangle++(0.04,0.04) (0.893,1.388)rectangle++(0.04,0.04) (0.812,2.74)rectangle++(0.04,0.04) (0.798,3.893)rectangle++(0.04,0.04)
  (0.832,4.873)rectangle++(0.04,0.04) (-1.122,2.852)rectangle++(0.04,0.04) (-0.842,4.066)rectangle++(0.04,0.04) (-0.555,5.086)rectangle++(0.04,0.04) (-0.271,5.941)rectangle++(0.04,0.04);
\fill[nuage]
  (0.005,6.657)rectangle++(0.04,0.04) (0.269,7.254)rectangle++(0.04,0.04) (0.516,7.752)rectangle++(0.04,0.04) (0.747,8.165)rectangle++(0.04,0.04) (-1.995,3.557)rectangle++(0.04,0.04)
  (-1.555,4.699)rectangle++(0.04,0.04) (1.532,1.153)rectangle++(0.04,0.04) (0.076,1.105)rectangle++(0.04,0.04) (0.106,2.532)rectangle++(0.04,0.04) (0.19,3.745)rectangle++(0.04,0.04)
  (0.309,4.771)rectangle++(0.04,0.04) (0.451,5.64)rectangle++(0.04,0.04) (-1.397,2.933)rectangle++(0.04,0.04) (-1.073,4.146)rectangle++(0.04,0.04) (-0.748,5.163)rectangle++(0.04,0.04)
  (-0.432,6.015)rectangle++(0.04,0.04) (-0.128,6.726)rectangle++(0.04,0.04) (0.158,7.318)rectangle++(0.04,0.04) (0.425,7.81)rectangle++(0.04,0.04) (0.671,8.218)rectangle++(0.04,0.04)
  (2.183,2.577)rectangle++(0.04,0.04) (1.956,3.699)rectangle++(0.04,0.04) (1.809,4.662)rectangle++(0.04,0.04) (1.722,5.487)rectangle++(0.04,0.04) (1.681,6.191)rectangle++(0.04,0.04)
  (-1.295,3.338)rectangle++(0.04,0.04) (-0.969,4.485)rectangle++(0.04,0.04) (-0.647,5.447)rectangle++(0.04,0.04) (-0.334,6.252)rectangle++(0.04,0.04) (-0.036,6.924)rectangle++(0.04,0.04)
  (0.244,7.483)rectangle++(0.04,0.04) (2.041,2.289)rectangle++(0.04,0.04) (-0.02,0.35)rectangle++(0.04,0.04) (-0.005,1.894)rectangle++(0.04,0.04) (0.069,3.206)rectangle++(0.04,0.04)
  (0.185,4.319)rectangle++(0.04,0.04) (1.164,1.515)rectangle++(0.04,0.04) (1.048,2.837)rectangle++(0.04,0.04) (1.002,3.966)rectangle++(0.04,0.04) (1.008,4.927)rectangle++(0.04,0.04)
  (1.052,5.744)rectangle++(0.04,0.04) (1.122,6.436)rectangle++(0.04,0.04) (1.209,7.022)rectangle++(0.04,0.04) (1.306,7.517)rectangle++(0.04,0.04) (1.409,7.933)rectangle++(0.04,0.04)
  (1.512,8.283)rectangle++(0.04,0.04) (1.615,8.576)rectangle++(0.04,0.04) (1.713,8.821)rectangle++(0.04,0.04) (1.807,9.026)rectangle++(0.04,0.04) (1.895,9.196)rectangle++(0.04,0.04)
  (1.976,9.337)rectangle++(0.04,0.04) (2.051,9.454)rectangle++(0.04,0.04) (2.322,3.232)rectangle++(0.04,0.04) (-0.397,2.834)rectangle++(0.04,0.04) (-0.226,4.021)rectangle++(0.04,0.04)
  (1.142,1.336)rectangle++(0.04,0.04) (1.022,2.686)rectangle++(0.04,0.04) (0.974,3.839)rectangle++(0.04,0.04) (0.979,4.82)rectangle++(0.04,0.04) (1.023,5.654)rectangle++(0.04,0.04)
  (1.094,6.361)rectangle++(0.04,0.04) (1.182,6.96)rectangle++(0.04,0.04) (1.281,7.465)rectangle++(0.04,0.04) (-1.706,3.526)rectangle++(0.04,0.04) (-1.311,4.661)rectangle++(0.04,0.04)
  (-0.93,5.611)rectangle++(0.04,0.04) (-0.568,6.403)rectangle++(0.04,0.04) (-0.229,7.061)rectangle++(0.04,0.04) (0.085,7.607)rectangle++(0.04,0.04) (0.375,8.059)rectangle++(0.04,0.04)
  (0.639,8.431)rectangle++(0.04,0.04) (0.878,8.737)rectangle++(0.04,0.04) (1.093,8.988)rectangle++(0.04,0.04) (1.287,9.192)rectangle++(0.04,0.04) (1.459,9.358)rectangle++(0.04,0.04)
  (-2.162,3.983)rectangle++(0.04,0.04) (-1.681,5.069)rectangle++(0.04,0.04) (-1.228,5.972)rectangle++(0.04,0.04) (-1.819,2.62)rectangle++(0.04,0.04) (-1.444,3.896)rectangle++(0.04,0.04)
  (-1.074,4.966)rectangle++(0.04,0.04) (-0.716,5.86)rectangle++(0.04,0.04) (-0.377,6.606)rectangle++(0.04,0.04) (-0.058,7.226)rectangle++(0.04,0.04) (0.237,7.741)rectangle++(0.04,0.04)
  (0.509,8.166)rectangle++(0.04,0.04) (0.757,8.517)rectangle++(0.04,0.04) (2.254,2.671)rectangle++(0.04,0.04) (2.02,3.776)rectangle++(0.04,0.04) (1.866,4.725)rectangle++(0.04,0.04)
  (1.773,5.538)rectangle++(0.04,0.04) (1.726,6.233)rectangle++(0.04,0.04) (1.715,6.825)rectangle++(0.04,0.04) (1.728,7.329)rectangle++(0.04,0.04) (1.76,7.757)rectangle++(0.04,0.04)
  (1.89,2.749)rectangle++(0.04,0.04) (1.715,3.857)rectangle++(0.04,0.04) (0.805,1.802)rectangle++(0.04,0.04) (0.366,1.072)rectangle++(0.04,0.04) (0.352,2.493)rectangle++(0.04,0.04)
  (0.397,3.701)rectangle++(0.04,0.04) (-0.02,0.575)rectangle++(0.04,0.04) (0.004,2.086)rectangle++(0.04,0.04) (0.084,3.369)rectangle++(0.04,0.04) (-0.88,2.35)rectangle++(0.04,0.04)
  (-0.656,3.629)rectangle++(0.04,0.04) (-0.415,4.707)rectangle++(0.04,0.04) (-0.167,5.614)rectangle++(0.04,0.04) (0.081,6.374)rectangle++(0.04,0.04) (0.321,7.011)rectangle++(0.04,0.04)
  (0.551,7.543)rectangle++(0.04,0.04) (0.768,7.986)rectangle++(0.04,0.04) (0.97,8.353)rectangle++(0.04,0.04) (1.157,8.658)rectangle++(0.04,0.04) (1.327,8.909)rectangle++(0.04,0.04)
  (1.482,9.116)rectangle++(0.04,0.04) (1.622,9.285)rectangle++(0.04,0.04) (1.748,9.424)rectangle++(0.04,0.04) (1.861,9.536)rectangle++(0.04,0.04) (1.961,9.628)rectangle++(0.04,0.04)
  (2.05,9.701)rectangle++(0.04,0.04) (2.128,9.76)rectangle++(0.04,0.04) (2.197,9.807)rectangle++(0.04,0.04) (2.258,9.845)rectangle++(0.04,0.04) (2.311,9.874)rectangle++(0.04,0.04)
  (2.357,9.897)rectangle++(0.04,0.04) (2.397,9.914)rectangle++(0.04,0.04) (2.432,9.927)rectangle++(0.04,0.04) (2.462,9.937)rectangle++(0.04,0.04) (2.488,9.944)rectangle++(0.04,0.04)
  (2.51,9.949)rectangle++(0.04,0.04) (2.529,9.953)rectangle++(0.04,0.04) (2.546,9.955)rectangle++(0.04,0.04) (2.56,9.956)rectangle++(0.04,0.04) (2.386,3.485)rectangle++(0.04,0.04)
  (2.166,4.463)rectangle++(0.04,0.04) (2.017,5.303)rectangle++(0.04,0.04) (1.924,6.023)rectangle++(0.04,0.04) (1.874,6.639)rectangle++(0.04,0.04) (1.857,7.164)rectangle++(0.04,0.04)
  (-1.513,3.592)rectangle++(0.04,0.04) (-1.144,4.71)rectangle++(0.04,0.04) (-0.786,5.646)rectangle++(0.04,0.04) (-0.445,6.426)rectangle++(0.04,0.04) (1.849,1.857)rectangle++(0.04,0.04)
  (1.643,3.1)rectangle++(0.04,0.04) (1.519,4.166)rectangle++(0.04,0.04) (1.455,5.076)rectangle++(0.04,0.04) (1.438,5.853)rectangle++(0.04,0.04) (1.454,6.514)rectangle++(0.04,0.04)
  (1.494,7.075)rectangle++(0.04,0.04) (1.551,7.55)rectangle++(0.04,0.04) (-1.674,3.607)rectangle++(0.04,0.04) (-0.02,0.56)rectangle++(0.04,0.04) (0.003,2.073)rectangle++(0.04,0.04)
  (0.083,3.358)rectangle++(0.04,0.04) (0.203,4.447)rectangle++(0.04,0.04) (0.348,5.368)rectangle++(0.04,0.04) (0.509,6.145)rectangle++(0.04,0.04) (0.676,6.799)rectangle++(0.04,0.04)
  (0.844,7.349)rectangle++(0.04,0.04) (1.009,7.809)rectangle++(0.04,0.04) (1.168,8.193)rectangle++(0.04,0.04) (1.318,8.514)rectangle++(0.04,0.04) (1.459,8.78)rectangle++(0.04,0.04)
  (1.589,9.001)rectangle++(0.04,0.04) (1.709,9.183)rectangle++(0.04,0.04) (1.817,9.334)rectangle++(0.04,0.04) (1.916,9.457)rectangle++(0.04,0.04) (2.005,9.558)rectangle++(0.04,0.04)
  (2.084,9.64)rectangle++(0.04,0.04) (2.155,9.707)rectangle++(0.04,0.04) (2.218,9.761)rectangle++(0.04,0.04) (2.273,9.804)rectangle++(0.04,0.04) (2.322,9.839)rectangle++(0.04,0.04)
  (2.365,9.867)rectangle++(0.04,0.04) (2.403,9.888)rectangle++(0.04,0.04) (2.436,9.905)rectangle++(0.04,0.04) (2.464,9.918)rectangle++(0.04,0.04) (2.489,9.928)rectangle++(0.04,0.04)
  (-2.105,4.346)rectangle++(0.04,0.04) (-1.617,5.374)rectangle++(0.04,0.04) (-1.162,6.229)rectangle++(0.04,0.04) (-0.741,6.937)rectangle++(0.04,0.04) (-0.354,7.523)rectangle++(0.04,0.04)
  (-0.002,8.005)rectangle++(0.04,0.04) (0.316,8.4)rectangle++(0.04,0.04) (0.602,8.724)rectangle++(0.04,0.04) (2.335,2.68)rectangle++(0.04,0.04) (0.383,1.68)rectangle++(0.04,0.04)
  (-0.382,2.047)rectangle++(0.04,0.04) (-0.245,3.351)rectangle++(0.04,0.04) (-0.076,4.454)rectangle++(0.04,0.04) (-0.02,0.696)rectangle++(0.04,0.04) (0.009,2.189)rectangle++(0.04,0.04)
  (0.093,3.456)rectangle++(0.04,0.04) (0.215,4.53)rectangle++(0.04,0.04) (0.362,5.438)rectangle++(0.04,0.04) (0.523,6.204)rectangle++(0.04,0.04) (0.69,6.849)rectangle++(0.04,0.04)
  (0.858,7.39)rectangle++(0.04,0.04) (1.023,7.844)rectangle++(0.04,0.04) (1.181,8.222)rectangle++(0.04,0.04) (1.331,8.538)rectangle++(0.04,0.04) (2.174,2.825)rectangle++(0.04,0.04)
  (1.958,3.911)rectangle++(0.04,0.04) (1.819,4.842)rectangle++(0.04,0.04) (1.737,5.639)rectangle++(0.04,0.04) (1.7,6.32)rectangle++(0.04,0.04) (1.696,6.9)rectangle++(0.04,0.04)
  (1.715,7.393)rectangle++(0.04,0.04) (1.751,7.812)rectangle++(0.04,0.04) (1.799,8.166)rectangle++(0.04,0.04) (1.854,8.466)rectangle++(0.04,0.04) (1.912,8.718)rectangle++(0.04,0.04)
  (2.137,3.019)rectangle++(0.04,0.04) (1.935,4.077)rectangle++(0.04,0.04) (1.806,4.984)rectangle++(0.04,0.04) (1.732,5.761)rectangle++(0.04,0.04) (1.7,6.424)rectangle++(0.04,0.04)
  (1.7,6.988)rectangle++(0.04,0.04) (1.722,7.468)rectangle++(0.04,0.04) (-1.618,3.628)rectangle++(0.04,0.04) (-1.233,4.745)rectangle++(0.04,0.04) (-0.86,5.679)rectangle++(0.04,0.04)
  (-0.506,6.457)rectangle++(0.04,0.04) (-0.174,7.105)rectangle++(0.04,0.04) (0.134,7.643)rectangle++(0.04,0.04) (0.417,8.087)rectangle++(0.04,0.04) (0.676,8.454)rectangle++(0.04,0.04)
  (0.911,8.755)rectangle++(0.04,0.04) (1.122,9.001)rectangle++(0.04,0.04) (1.312,9.202)rectangle++(0.04,0.04) (1.481,9.366)rectangle++(0.04,0.04) (1.631,9.498)rectangle++(0.04,0.04)
  (1.764,9.604)rectangle++(0.04,0.04) (1.881,9.689)rectangle++(0.04,0.04) (1.985,9.757)rectangle++(0.04,0.04) (-2.161,4.192)rectangle++(0.04,0.04) (-1.671,5.246)rectangle++(0.04,0.04)
  (-1.213,6.122)rectangle++(0.04,0.04) (-0.788,6.848)rectangle++(0.04,0.04) (-0.398,7.449)rectangle++(0.04,0.04) (-0.043,7.944)rectangle++(0.04,0.04) (0.279,8.35)rectangle++(0.04,0.04)
  (0.569,8.683)rectangle++(0.04,0.04) (0.829,8.954)rectangle++(0.04,0.04) (1.06,9.174)rectangle++(0.04,0.04) (1.266,9.351)rectangle++(0.04,0.04) (1.448,9.494)rectangle++(0.04,0.04)
  (1.608,9.608)rectangle++(0.04,0.04) (1.749,9.699)rectangle++(0.04,0.04) (1.873,9.77)rectangle++(0.04,0.04) (1.98,9.826)rectangle++(0.04,0.04) (2.074,9.869)rectangle++(0.04,0.04)
  (2.156,9.902)rectangle++(0.04,0.04) (2.226,9.927)rectangle++(0.04,0.04) (2.287,9.945)rectangle++(0.04,0.04) (2.34,9.958)rectangle++(0.04,0.04) (2.42,3.428)rectangle++(0.04,0.04)
  (2.192,4.413)rectangle++(0.04,0.04) (2.037,5.26)rectangle++(0.04,0.04) (1.94,5.986)rectangle++(0.04,0.04) (1.368,2.371)rectangle++(0.04,0.04) (1.255,3.557)rectangle++(0.04,0.04)
  (1.207,4.569)rectangle++(0.04,0.04) (1.206,5.432)rectangle++(0.04,0.04) (1.241,6.165)rectangle++(0.04,0.04) (1.299,6.787)rectangle++(0.04,0.04) (1.373,7.313)rectangle++(0.04,0.04)
  (1.458,7.757)rectangle++(0.04,0.04) (1.547,8.132)rectangle++(0.04,0.04) (1.638,8.446)rectangle++(0.04,0.04) (1.728,8.71)rectangle++(0.04,0.04) (1.815,8.931)rectangle++(0.04,0.04)
  (1.898,9.115)rectangle++(0.04,0.04) (1.975,9.268)rectangle++(0.04,0.04) (2.048,9.395)rectangle++(0.04,0.04) (2.114,9.5)rectangle++(0.04,0.04) (2.175,9.586)rectangle++(0.04,0.04)
  (2.23,9.658)rectangle++(0.04,0.04) (2.279,9.716)rectangle++(0.04,0.04) (2.324,9.764)rectangle++(0.04,0.04) (2.364,9.802)rectangle++(0.04,0.04) (2.399,9.834)rectangle++(0.04,0.04)
  (2.43,9.859)rectangle++(0.04,0.04) (-2.098,4.317)rectangle++(0.04,0.04) (-0.02,0.674)rectangle++(0.04,0.04) (0.008,2.17)rectangle++(0.04,0.04) (0.091,3.44)rectangle++(0.04,0.04)
  (0.213,4.517)rectangle++(0.04,0.04) (0.359,5.427)rectangle++(0.04,0.04) (0.52,6.195)rectangle++(0.04,0.04) (0.688,6.841)rectangle++(0.04,0.04) (-1.662,3.252)rectangle++(0.04,0.04)
  (-1.285,4.427)rectangle++(0.04,0.04) (-0.917,5.411)rectangle++(0.04,0.04) (-0.566,6.232)rectangle++(0.04,0.04) (-0.234,6.916)rectangle++(0.04,0.04) (0.076,7.484)rectangle++(0.04,0.04)
  (0.362,7.955)rectangle++(0.04,0.04) (0.623,8.343)rectangle++(0.04,0.04) (0.861,8.663)rectangle++(0.04,0.04) (1.077,8.925)rectangle++(0.04,0.04) (1.27,9.14)rectangle++(0.04,0.04)
  (1.443,9.314)rectangle++(0.04,0.04) (1.597,9.455)rectangle++(0.04,0.04) (1.733,9.569)rectangle++(0.04,0.04) (-2.163,4.093)rectangle++(0.04,0.04) (-1.677,5.162)rectangle++(0.04,0.04)
  (-1.221,6.051)rectangle++(0.04,0.04) (-0.798,6.788)rectangle++(0.04,0.04) (-0.409,7.398)rectangle++(0.04,0.04) (-0.054,7.901)rectangle++(0.04,0.04) (0.268,8.314)rectangle++(0.04,0.04)
  (0.558,8.652)rectangle++(0.04,0.04) (0.818,8.928)rectangle++(0.04,0.04) (1.051,9.153)rectangle++(0.04,0.04) (1.257,9.334)rectangle++(0.04,0.04) (2.408,2.997)rectangle++(0.04,0.04)
  (2.164,4.047)rectangle++(0.04,0.04) (1.999,4.95)rectangle++(0.04,0.04) (1.895,5.724)rectangle++(0.04,0.04) (1.838,6.385)rectangle++(0.04,0.04) (1.815,6.95)rectangle++(0.04,0.04)
  (1.819,7.431)rectangle++(0.04,0.04) (1.841,7.84)rectangle++(0.04,0.04) (1.876,8.187)rectangle++(0.04,0.04) (1.92,8.48)rectangle++(0.04,0.04) (1.969,8.727)rectangle++(0.04,0.04)
  (2.021,8.936)rectangle++(0.04,0.04) (2.073,9.111)rectangle++(0.04,0.04) (2.124,9.257)rectangle++(0.04,0.04) (2.174,9.38)rectangle++(0.04,0.04) (2.22,9.482)rectangle++(0.04,0.04)
  (2.265,9.567)rectangle++(0.04,0.04) (2.305,9.638)rectangle++(0.04,0.04) (2.343,9.696)rectangle++(0.04,0.04) (2.377,9.744)rectangle++(0.04,0.04) (2.408,9.784)rectangle++(0.04,0.04)
  (2.436,9.816)rectangle++(0.04,0.04) (-2.086,4.309)rectangle++(0.04,0.04) (-1.559,2.057)rectangle++(0.04,0.04) (-1.245,3.407)rectangle++(0.04,0.04) (-0.924,4.542)rectangle++(0.04,0.04)
  (1.393,1.28)rectangle++(0.04,0.04) (1.233,2.628)rectangle++(0.04,0.04) (1.151,3.781)rectangle++(0.04,0.04) (1.127,4.764)rectangle++(0.04,0.04) (1.147,5.6)rectangle++(0.04,0.04)
  (1.196,6.311)rectangle++(0.04,0.04) (-0.02,0.993)rectangle++(0.04,0.04) (0.021,2.441)rectangle++(0.04,0.04) (0.113,3.67)rectangle++(0.04,0.04) (0.241,4.711)rectangle++(0.04,0.04)
  (0.391,5.591)rectangle++(0.04,0.04) (0.554,6.333)rectangle++(0.04,0.04) (0.722,6.957)rectangle++(0.04,0.04) (0.889,7.481)rectangle++(0.04,0.04) (-1.788,3.439)rectangle++(0.04,0.04)
  (-1.273,1.934)rectangle++(0.04,0.04) (-1.007,3.291)rectangle++(0.04,0.04) (-0.726,4.434)rectangle++(0.04,0.04) (-0.442,5.394)rectangle++(0.04,0.04) (-0.162,6.199)rectangle++(0.04,0.04)
  (0.108,6.872)rectangle++(0.04,0.04) (0.364,7.433)rectangle++(0.04,0.04) (0.605,7.9)rectangle++(0.04,0.04) (0.828,8.287)rectangle++(0.04,0.04) (1.033,8.607)rectangle++(0.04,0.04)
  (1.22,8.871)rectangle++(0.04,0.04) (1.39,9.087)rectangle++(0.04,0.04) (1.543,9.265)rectangle++(0.04,0.04) (1.68,9.41)rectangle++(0.04,0.04) (1.802,9.527)rectangle++(0.04,0.04)
  (1.91,9.622)rectangle++(0.04,0.04) (2.007,9.699)rectangle++(0.04,0.04) (2.091,9.76)rectangle++(0.04,0.04) (2.166,9.808)rectangle++(0.04,0.04) (-2.138,4.245)rectangle++(0.04,0.04)
  (-1.65,5.29)rectangle++(0.04,0.04) (-1.193,6.159)rectangle++(0.04,0.04) (-0.77,6.879)rectangle++(0.04,0.04) (-0.381,7.474)rectangle++(0.04,0.04) (-0.027,7.964)rectangle++(0.04,0.04)
  (0.293,8.367)rectangle++(0.04,0.04) (0.582,8.696)rectangle++(0.04,0.04) (0.84,8.965)rectangle++(0.04,0.04) (1.07,9.183)rectangle++(0.04,0.04) (1.275,9.359)rectangle++(0.04,0.04)
  (1.456,9.5)rectangle++(0.04,0.04) (1.615,9.613)rectangle++(0.04,0.04) (1.755,9.703)rectangle++(0.04,0.04) (1.878,9.773)rectangle++(0.04,0.04) (1.985,9.828)rectangle++(0.04,0.04)
  (2.078,9.871)rectangle++(0.04,0.04) (2.435,3.339)rectangle++(0.04,0.04) (2.201,4.337)rectangle++(0.04,0.04) (2.042,5.195)rectangle++(0.04,0.04) (1.941,5.93)rectangle++(0.04,0.04)
  (1.885,6.559)rectangle++(0.04,0.04) (1.863,7.096)rectangle++(0.04,0.04) (1.69,2.617)rectangle++(0.04,0.04) (1.539,3.753)rectangle++(0.04,0.04) (1.456,4.725)rectangle++(0.04,0.04)
  (1.425,5.554)rectangle++(0.04,0.04) (1.431,6.26)rectangle++(0.04,0.04) (1.464,6.86)rectangle++(0.04,0.04) (1.517,7.369)rectangle++(0.04,0.04) (1.582,7.799)rectangle++(0.04,0.04)
  (1.654,8.162)rectangle++(0.04,0.04) (1.731,8.468)rectangle++(0.04,0.04) (1.807,8.725)rectangle++(0.04,0.04) (1.883,8.94)rectangle++(0.04,0.04) (1.956,9.12)rectangle++(0.04,0.04)
  (-2.001,4.045)rectangle++(0.04,0.04) (-1.541,5.115)rectangle++(0.04,0.04) (-1.108,6.005)rectangle++(0.04,0.04) (-0.704,6.745)rectangle++(0.04,0.04) (-0.33,7.358)rectangle++(0.04,0.04)
  (0.011,7.863)rectangle++(0.04,0.04) (0.322,8.28)rectangle++(0.04,0.04) (-2.11,3.485)rectangle++(0.04,0.04) (-1.656,4.642)rectangle++(0.04,0.04) (-1.224,5.609)rectangle++(0.04,0.04)
  (-0.818,6.412)rectangle++(0.04,0.04) (-1.852,2.812)rectangle++(0.04,0.04) (-1.464,4.06)rectangle++(0.04,0.04) (-1.084,5.106)rectangle++(0.04,0.04) (-0.72,5.98)rectangle++(0.04,0.04)
  (-0.375,6.708)rectangle++(0.04,0.04) (-0.052,7.313)rectangle++(0.04,0.04) (0.246,7.814)rectangle++(0.04,0.04) (0.519,8.228)rectangle++(0.04,0.04) (0.768,8.569)rectangle++(0.04,0.04)
  (0.994,8.85)rectangle++(0.04,0.04) (1.196,9.079)rectangle++(0.04,0.04) (1.378,9.265)rectangle++(0.04,0.04) (1.54,9.416)rectangle++(0.04,0.04) (1.683,9.539)rectangle++(0.04,0.04)
  (1.81,9.637)rectangle++(0.04,0.04) (1.922,9.715)rectangle++(0.04,0.04) (2.02,9.777)rectangle++(0.04,0.04) (2.106,9.826)rectangle++(0.04,0.04) (2.181,9.864)rectangle++(0.04,0.04)
  (2.246,9.893)rectangle++(0.04,0.04) (2.303,9.916)rectangle++(0.04,0.04) (2.414,3.408)rectangle++(0.04,0.04) (2.186,4.397)rectangle++(0.04,0.04) (2.032,5.246)rectangle++(0.04,0.04)
  (1.934,5.974)rectangle++(0.04,0.04) (1.881,6.597)rectangle++(0.04,0.04) (1.861,7.128)rectangle++(0.04,0.04) (1.864,7.581)rectangle++(0.04,0.04) (1.886,7.965)rectangle++(0.04,0.04)
  (1.919,8.291)rectangle++(0.04,0.04) (1.961,8.567)rectangle++(0.04,0.04) (2.007,8.8)rectangle++(0.04,0.04) (-1.908,3.987)rectangle++(0.04,0.04) (-1.464,5.061)rectangle++(0.04,0.04)
  (-1.044,5.957)rectangle++(0.04,0.04) (1.807,1.588)rectangle++(0.04,0.04) (0.156,1.281)rectangle++(0.04,0.04) (0.182,2.679)rectangle++(0.04,0.04) (0.259,3.866)rectangle++(0.04,0.04)
  (0.373,4.872)rectangle++(0.04,0.04) (-1.213,2.747)rectangle++(0.04,0.04) (-0.924,3.979)rectangle++(0.04,0.04) (-0.628,5.016)rectangle++(0.04,0.04) (-0.335,5.885)rectangle++(0.04,0.04)
  (-0.052,6.612)rectangle++(0.04,0.04) (0.218,7.218)rectangle++(0.04,0.04) (0.472,7.723)rectangle++(0.04,0.04) (0.708,8.142)rectangle++(0.04,0.04) (-1.996,3.543)rectangle++(0.04,0.04)
  (-1.557,4.688)rectangle++(0.04,0.04) (-1.139,5.643)rectangle++(0.04,0.04) (-0.744,6.438)rectangle++(0.04,0.04) (-1.844,2.834)rectangle++(0.04,0.04) (-1.456,4.079)rectangle++(0.04,0.04)
  (-1.077,5.122)rectangle++(0.04,0.04) (-0.713,5.993)rectangle++(0.04,0.04) (-0.368,6.718)rectangle++(0.04,0.04) (-0.046,7.322)rectangle++(0.04,0.04) (0.251,7.821)rectangle++(0.04,0.04)
  (0.524,8.234)rectangle++(0.04,0.04) (0.773,8.574)rectangle++(0.04,0.04) (0.998,8.854)rectangle++(0.04,0.04) (1.2,9.082)rectangle++(0.04,0.04) (1.381,9.268)rectangle++(0.04,0.04)
  (1.542,9.419)rectangle++(0.04,0.04) (1.686,9.54)rectangle++(0.04,0.04) (-2.165,4.076)rectangle++(0.04,0.04) (-1.679,5.147)rectangle++(0.04,0.04) (-1.224,6.038)rectangle++(0.04,0.04)
  (-0.801,6.778)rectangle++(0.04,0.04) (-0.412,7.389)rectangle++(0.04,0.04) (-0.057,7.894)rectangle++(0.04,0.04) (-2.085,3.313)rectangle++(0.04,0.04) (-1.642,4.495)rectangle++(0.04,0.04)
  (-1.218,5.483)rectangle++(0.04,0.04) (-0.818,6.305)rectangle++(0.04,0.04) (-0.445,6.988)rectangle++(0.04,0.04) (-0.101,7.554)rectangle++(0.04,0.04) (0.214,8.021)rectangle++(0.04,0.04)
  (0.5,8.406)rectangle++(0.04,0.04) (-2.107,3.553)rectangle++(0.04,0.04) (-1.651,4.701)rectangle++(0.04,0.04) (1.546,1.129)rectangle++(0.04,0.04) (1.357,2.494)rectangle++(0.04,0.04)
  (1.251,3.662)rectangle++(0.04,0.04) (-0.723,2.682)rectangle++(0.04,0.04) (-0.509,3.905)rectangle++(0.04,0.04) (-0.279,4.936)rectangle++(0.04,0.04) (-0.042,5.803)rectangle++(0.04,0.04)
  (0.194,6.53)rectangle++(0.04,0.04) (0.424,7.139)rectangle++(0.04,0.04) (1.918,2.254)rectangle++(0.04,0.04) (1.718,3.435)rectangle++(0.04,0.04) (1.596,4.447)rectangle++(0.04,0.04)
  (1.532,5.313)rectangle++(0.04,0.04) (1.513,6.051)rectangle++(0.04,0.04) (1.525,6.679)rectangle++(0.04,0.04) (1.562,7.212)rectangle++(0.04,0.04) (1.614,7.664)rectangle++(0.04,0.04)
  (1.676,8.046)rectangle++(0.04,0.04) (-1.778,3.745)rectangle++(0.04,0.04) (-1.364,4.85)rectangle++(0.04,0.04) (-0.967,5.773)rectangle++(0.04,0.04) (1.744,1.564)rectangle++(0.04,0.04)
  (1.543,2.856)rectangle++(0.04,0.04) (1.424,3.962)rectangle++(0.04,0.04) (1.366,4.907)rectangle++(0.04,0.04) (1.355,5.712)rectangle++(0.04,0.04) (1.378,6.398)rectangle++(0.04,0.04)
  (1.425,6.979)rectangle++(0.04,0.04) (1.489,7.471)rectangle++(0.04,0.04) (1.562,7.887)rectangle++(0.04,0.04) (1.641,8.238)rectangle++(0.04,0.04) (1.722,8.533)rectangle++(0.04,0.04)
  (1.803,8.78)rectangle++(0.04,0.04) (1.881,8.987)rectangle++(0.04,0.04) (-1.982,3.999)rectangle++(0.04,0.04) (-1.527,5.075)rectangle++(0.04,0.04) (-1.097,5.971)rectangle++(0.04,0.04)
  (-0.696,6.715)rectangle++(0.04,0.04) (-0.325,7.332)rectangle++(0.04,0.04) (0.015,7.841)rectangle++(0.04,0.04) (0.324,8.261)rectangle++(0.04,0.04) (0.604,8.605)rectangle++(0.04,0.04)
  (0.855,8.886)rectangle++(0.04,0.04) (1.08,9.115)rectangle++(0.04,0.04) (1.281,9.301)rectangle++(0.04,0.04) (1.458,9.451)rectangle++(0.04,0.04) (1.615,9.571)rectangle++(0.04,0.04)
  (2.42,3.147)rectangle++(0.04,0.04) (-0.355,2.838)rectangle++(0.04,0.04) (-0.191,4.023)rectangle++(0.04,0.04) (-0.003,5.023)rectangle++(0.04,0.04) (0.196,5.866)rectangle++(0.04,0.04)
  (0.399,6.574)rectangle++(0.04,0.04) (0.6,7.169)rectangle++(0.04,0.04) (1.9,2.306)rectangle++(0.04,0.04) (1.705,3.481)rectangle++(0.04,0.04) (1.586,4.487)rectangle++(0.04,0.04)
  (1.526,5.346)rectangle++(0.04,0.04) (1.508,6.08)rectangle++(0.04,0.04) (1.459,2.281)rectangle++(0.04,0.04) (1.329,3.477)rectangle++(0.04,0.04) (1.266,4.498)rectangle++(0.04,0.04)
  (1.254,5.369)rectangle++(0.04,0.04) (1.279,6.11)rectangle++(0.04,0.04) (1.329,6.738)rectangle++(0.04,0.04) (-0.02,1.061)rectangle++(0.04,0.04) (0.283,0.68)rectangle++(0.04,0.04)
  (-0.141,1.804)rectangle++(0.04,0.04) (-0.05,3.135)rectangle++(0.04,0.04) (0.081,4.263)rectangle++(0.04,0.04) (1.164,1.474)rectangle++(0.04,0.04) (0.221,1.086)rectangle++(0.04,0.04)
  (0.229,2.511)rectangle++(0.04,0.04) (0.293,3.721)rectangle++(0.04,0.04) (-0.93,2.475)rectangle++(0.04,0.04) (-0.694,3.737)rectangle++(0.04,0.04) (-0.442,4.801)rectangle++(0.04,0.04)
  (-0.186,5.694)rectangle++(0.04,0.04) (0.067,6.444)rectangle++(0.04,0.04) (1.777,1.994)rectangle++(0.04,0.04) (1.588,3.22)rectangle++(0.04,0.04) (1.476,4.27)rectangle++(0.04,0.04)
  (1.423,5.166)rectangle++(0.04,0.04) (1.414,5.931)rectangle++(0.04,0.04) (1.437,6.581)rectangle++(0.04,0.04) (1.483,7.132)rectangle++(0.04,0.04) (1.543,7.599)rectangle++(0.04,0.04)
  (1.614,7.994)rectangle++(0.04,0.04) (1.689,8.326)rectangle++(0.04,0.04) (1.767,8.606)rectangle++(0.04,0.04) (1.844,8.841)rectangle++(0.04,0.04) (1.919,9.037)rectangle++(0.04,0.04)
  (1.99,9.201)rectangle++(0.04,0.04) (2.057,9.337)rectangle++(0.04,0.04) (2.12,9.451)rectangle++(0.04,0.04) (2.178,9.545)rectangle++(0.04,0.04) (2.231,9.622)rectangle++(0.04,0.04)
  (2.279,9.686)rectangle++(0.04,0.04) (2.322,9.738)rectangle++(0.04,0.04) (2.361,9.78)rectangle++(0.04,0.04) (2.396,9.815)rectangle++(0.04,0.04) (2.427,9.843)rectangle++(0.04,0.04)
  (2.455,9.866)rectangle++(0.04,0.04) (2.479,9.884)rectangle++(0.04,0.04) (-2.095,4.334)rectangle++(0.04,0.04) (-1.61,5.364)rectangle++(0.04,0.04) (-1.156,6.22)rectangle++(0.04,0.04)
  (-0.736,6.929)rectangle++(0.04,0.04) (-0.351,7.515)rectangle++(0.04,0.04) (0,7.998)rectangle++(0.04,0.04) (-0.02,1.263)rectangle++(0.04,0.04) (0.031,2.67)rectangle++(0.04,0.04)
  (0.131,3.865)rectangle++(0.04,0.04) (0.264,4.876)rectangle++(0.04,0.04) (0.417,5.73)rectangle++(0.04,0.04) (0.582,6.45)rectangle++(0.04,0.04) (0.75,7.056)rectangle++(0.04,0.04)
  (0.918,7.564)rectangle++(0.04,0.04) (1.08,7.988)rectangle++(0.04,0.04) (1.236,8.343)rectangle++(0.04,0.04) (1.382,8.638)rectangle++(0.04,0.04) (1.518,8.884)rectangle++(0.04,0.04)
  (1.643,9.087)rectangle++(0.04,0.04) (1.758,9.254)rectangle++(0.04,0.04) (2.31,3.108)rectangle++(0.04,0.04) (2.086,4.146)rectangle++(0.04,0.04) (1.936,5.037)rectangle++(0.04,0.04)
  (1.845,5.8)rectangle++(0.04,0.04) (1.798,6.452)rectangle++(0.04,0.04) (1.519,2.446)rectangle++(0.04,0.04) (1.387,3.615)rectangle++(0.04,0.04) (1.322,4.613)rectangle++(0.04,0.04)
  (1.306,5.464)rectangle++(0.04,0.04) (1.326,6.189)rectangle++(0.04,0.04) (1.373,6.804)rectangle++(0.04,0.04) (1.437,7.324)rectangle++(0.04,0.04) (-1.638,3.531)rectangle++(0.04,0.04)
  (-1.253,4.663)rectangle++(0.04,0.04) (-0.881,5.61)rectangle++(0.04,0.04) (-0.527,6.4)rectangle++(0.04,0.04) (-1.79,2.876)rectangle++(0.04,0.04) (-1.409,4.112)rectangle++(0.04,0.04)
  (-1.035,5.148)rectangle++(0.04,0.04) (-1.567,2.483)rectangle++(0.04,0.04) (-1.235,3.77)rectangle++(0.04,0.04) (1.223,1.014)rectangle++(0.04,0.04) (1.078,2.409)rectangle++(0.04,0.04)
  (1.011,3.601)rectangle++(0.04,0.04) (0.839,1.557)rectangle++(0.04,0.04) (0.773,2.886)rectangle++(0.04,0.04) (0.771,4.018)rectangle++(0.04,0.04) (0.814,4.981)rectangle++(0.04,0.04)
  (0.889,5.798)rectangle++(0.04,0.04) (0.985,6.489)rectangle++(0.04,0.04) (1.095,7.072)rectangle++(0.04,0.04) (1.211,7.564)rectangle++(0.04,0.04) (1.33,7.977)rectangle++(0.04,0.04)
  (1.447,8.323)rectangle++(0.04,0.04) (1.561,8.613)rectangle++(0.04,0.04) (1.669,8.855)rectangle++(0.04,0.04) (1.771,9.056)rectangle++(0.04,0.04) (1.865,9.223)rectangle++(0.04,0.04)
  (1.952,9.361)rectangle++(0.04,0.04) (2.032,9.475)rectangle++(0.04,0.04) (-2.078,4.141)rectangle++(0.04,0.04) (1.454,0.883)rectangle++(0.04,0.04) (1.269,2.289)rectangle++(0.04,0.04)
  (1.168,3.491)rectangle++(0.04,0.04) (1.13,4.517)rectangle++(0.04,0.04) (1.139,5.39)rectangle++(0.04,0.04) (1.182,6.132)rectangle++(0.04,0.04) (-1.379,3.21)rectangle++(0.04,0.04)
  (-1.046,4.38)rectangle++(0.04,0.04) (-0.716,5.361)rectangle++(0.04,0.04) (-1.558,2.604)rectangle++(0.04,0.04) (-1.223,3.872)rectangle++(0.04,0.04) (-0.887,4.936)rectangle++(0.04,0.04)
  (1.498,1.384)rectangle++(0.04,0.04) (0.146,1.152)rectangle++(0.04,0.04) (0.168,2.569)rectangle++(0.04,0.04) (0.243,3.773)rectangle++(0.04,0.04) (0.355,4.794)rectangle++(0.04,0.04)
  (0.492,5.657)rectangle++(0.04,0.04) (1.493,1.915)rectangle++(0.04,0.04) (1.343,3.165)rectangle++(0.04,0.04) (1.266,4.232)rectangle++(0.04,0.04) (1.243,5.143)rectangle++(0.04,0.04)
  (1.26,5.918)rectangle++(0.04,0.04) (1.306,6.576)rectangle++(0.04,0.04) (1.371,7.134)rectangle++(0.04,0.04) (1.448,7.605)rectangle++(0.04,0.04) (1.533,8.003)rectangle++(0.04,0.04)
  (1.621,8.337)rectangle++(0.04,0.04) (1.709,8.618)rectangle++(0.04,0.04) (1.795,8.853)rectangle++(0.04,0.04) (1.878,9.049)rectangle++(0.04,0.04) (1.956,9.213)rectangle++(0.04,0.04)
  (2.029,9.349)rectangle++(0.04,0.04) (2.096,9.462)rectangle++(0.04,0.04) (2.158,9.555)rectangle++(0.04,0.04) (-2.074,4.187)rectangle++(0.04,0.04) (-1.597,5.238)rectangle++(0.04,0.04)
  (-1.15,6.113)rectangle++(0.04,0.04) (-0.736,6.838)rectangle++(0.04,0.04) (-0.354,7.438)rectangle++(0.04,0.04) (-0.006,7.933)rectangle++(0.04,0.04) (0.31,8.339)rectangle++(0.04,0.04)
  (0.595,8.672)rectangle++(0.04,0.04) (-2.157,3.634)rectangle++(0.04,0.04) (-1.69,4.771)rectangle++(0.04,0.04) (-1.248,5.719)rectangle++(0.04,0.04) (-0.834,6.508)rectangle++(0.04,0.04)
  (-0.451,7.161)rectangle++(0.04,0.04) (-0.099,7.701)rectangle++(0.04,0.04) (0.222,8.146)rectangle++(0.04,0.04) (0.512,8.511)rectangle++(0.04,0.04) (0.773,8.81)rectangle++(0.04,0.04)
  (1.008,9.054)rectangle++(0.04,0.04) (1.216,9.252)rectangle++(0.04,0.04) (2.391,2.967)rectangle++(0.04,0.04) (2.149,4.022)rectangle++(0.04,0.04) (1.985,4.929)rectangle++(0.04,0.04)
  (1.882,5.707)rectangle++(0.04,0.04) (1.826,6.372)rectangle++(0.04,0.04) (1.805,6.939)rectangle++(0.04,0.04) (1.809,7.422)rectangle++(0.04,0.04) (1.833,7.833)rectangle++(0.04,0.04)
  (1.869,8.181)rectangle++(0.04,0.04) (1.914,8.475)rectangle++(0.04,0.04) (1.963,8.723)rectangle++(0.04,0.04) (2.016,8.933)rectangle++(0.04,0.04) (2.181,3.098)rectangle++(0.04,0.04)
  (1.976,4.142)rectangle++(0.04,0.04) (1.843,5.038)rectangle++(0.04,0.04) (1.766,5.805)rectangle++(0.04,0.04) (1.731,6.46)rectangle++(0.04,0.04) (1.728,7.018)rectangle++(0.04,0.04)
  (1.747,7.492)rectangle++(0.04,0.04) (1.782,7.895)rectangle++(0.04,0.04) (1.829,8.235)rectangle++(0.04,0.04) (1.881,8.523)rectangle++(0.04,0.04) (1.938,8.766)rectangle++(0.04,0.04)
  (1.996,8.969)rectangle++(0.04,0.04) (2.053,9.14)rectangle++(0.04,0.04) (2.108,9.283)rectangle++(0.04,0.04) (2.161,9.403)rectangle++(0.04,0.04) (-0.02,1.488)rectangle++(0.04,0.04)
  (0.04,2.861)rectangle++(0.04,0.04) (0.147,4.027)rectangle++(0.04,0.04) (0.283,5.013)rectangle++(0.04,0.04) (-1.268,2.757)rectangle++(0.04,0.04) (-0.02,0.424)rectangle++(0.04,0.04)
  (-0.002,1.958)rectangle++(0.04,0.04) (0.074,3.26)rectangle++(0.04,0.04) (0.191,4.365)rectangle++(0.04,0.04) (0.335,5.298)rectangle++(0.04,0.04) (1.416,1.789)rectangle++(0.04,0.04)
  (-0.203,2.308)rectangle++(0.04,0.04) (-0.082,3.566)rectangle++(0.04,0.04) (0.07,4.631)rectangle++(0.04,0.04) (0.243,5.53)rectangle++(0.04,0.04) (0.425,6.287)rectangle++(0.04,0.04)
  (0.611,6.923)rectangle++(0.04,0.04) (0.794,7.456)rectangle++(0.04,0.04) (0.971,7.902)rectangle++(0.04,0.04) (1.139,8.274)rectangle++(0.04,0.04) (2.129,2.712)rectangle++(0.04,0.04)
  (1.916,3.816)rectangle++(0.04,0.04) (0.764,1.844)rectangle++(0.04,0.04) (0.721,3.133)rectangle++(0.04,0.04) (0.752,1.369)rectangle++(0.04,0.04) (0.692,2.73)rectangle++(0.04,0.04)
  (0.695,3.889)rectangle++(0.04,0.04) (0.744,4.874)rectangle++(0.04,0.04) (0.825,5.709)rectangle++(0.04,0.04) (0.928,6.416)rectangle++(0.04,0.04) (1.64,2.211)rectangle++(0.04,0.04)
  (1.48,3.41)rectangle++(0.04,0.04) (1.392,4.435)rectangle++(0.04,0.04) (-0.896,2.885)rectangle++(0.04,0.04) (-0.648,4.084)rectangle++(0.04,0.04) (-0.39,5.094)rectangle++(0.04,0.04)
  (-0.13,5.942)rectangle++(0.04,0.04) (0.125,6.652)rectangle++(0.04,0.04) (0.37,7.245)rectangle++(0.04,0.04) (0.602,7.74)rectangle++(0.04,0.04) (0.819,8.151)rectangle++(0.04,0.04)
  (1.02,8.492)rectangle++(0.04,0.04) (1.205,8.773)rectangle++(0.04,0.04) (1.373,9.005)rectangle++(0.04,0.04) (1.525,9.196)rectangle++(0.04,0.04) (1.662,9.352)rectangle++(0.04,0.04)
  (1.784,9.479)rectangle++(0.04,0.04) (1.894,9.582)rectangle++(0.04,0.04) (-2.134,4.133)rectangle++(0.04,0.04) (-1.651,5.194)rectangle++(0.04,0.04) (-1.197,6.077)rectangle++(0.04,0.04)
  (-0.777,6.81)rectangle++(0.04,0.04) (-0.39,7.416)rectangle++(0.04,0.04) (-0.037,7.915)rectangle++(0.04,0.04) (0.283,8.325)rectangle++(0.04,0.04) (0.571,8.662)rectangle++(0.04,0.04)
  (0.83,8.936)rectangle++(0.04,0.04) (1.061,9.158)rectangle++(0.04,0.04) (1.266,9.338)rectangle++(0.04,0.04) (2.408,3)rectangle++(0.04,0.04) (2.164,4.05)rectangle++(0.04,0.04)
  (1.999,4.952)rectangle++(0.04,0.04) (1.895,5.726)rectangle++(0.04,0.04) (1.838,6.387)rectangle++(0.04,0.04) (1.815,6.952)rectangle++(0.04,0.04) (1.819,7.433)rectangle++(0.04,0.04)
  (1.791,2.687)rectangle++(0.04,0.04) (1.628,3.808)rectangle++(0.04,0.04) (1.534,4.768)rectangle++(0.04,0.04) (1.492,5.588)rectangle++(0.04,0.04) (-1.176,3.161)rectangle++(0.04,0.04)
  (-0.875,4.33)rectangle++(0.04,0.04) (-0.573,5.312)rectangle++(0.04,0.04) (-0.277,6.134)rectangle++(0.04,0.04) (0.008,6.822)rectangle++(0.04,0.04) (0.278,7.394)rectangle++(0.04,0.04)
  (0.529,7.87)rectangle++(0.04,0.04) (2.107,2.456)rectangle++(0.04,0.04) (1.887,3.6)rectangle++(0.04,0.04) (1.746,4.581)rectangle++(0.04,0.04) (1.003,1.983)rectangle++(0.04,0.04)
  (0.93,3.242)rectangle++(0.04,0.04) (0.751,1.45)rectangle++(0.04,0.04) (0.694,2.799)rectangle++(0.04,0.04) (0.7,3.947)rectangle++(0.04,0.04) (0.75,4.923)rectangle++(0.04,0.04)
  (0.833,5.751)rectangle++(0.04,0.04) (0.936,6.451)rectangle++(0.04,0.04) (1.051,7.042)rectangle++(0.04,0.04) (1.173,7.54)rectangle++(0.04,0.04) (-1.747,3.518)rectangle++(0.04,0.04)
  (-1.347,4.656)rectangle++(0.04,0.04) (-0.96,5.608)rectangle++(0.04,0.04) (-0.594,6.401)rectangle++(0.04,0.04) (-0.251,7.061)rectangle++(0.04,0.04) (0.067,7.608)rectangle++(0.04,0.04)
  (0.359,8.06)rectangle++(0.04,0.04) (-2.045,3.445)rectangle++(0.04,0.04) (-1.326,1.876)rectangle++(0.04,0.04) (-1.054,3.244)rectangle++(0.04,0.04) (-1.076,2.06)rectangle++(0.04,0.04)
  (0.721,0.645)rectangle++(0.04,0.04) (0.636,2.115)rectangle++(0.04,0.04) (-0.444,2.201)rectangle++(0.04,0.04) (-0.292,3.485)rectangle++(0.04,0.04) (-0.111,4.57)rectangle++(0.04,0.04)
  (0.087,5.485)rectangle++(0.04,0.04) (0.291,6.255)rectangle++(0.04,0.04) (0.495,6.901)rectangle++(0.04,0.04) (0.695,7.442)rectangle++(0.04,0.04) (0.886,7.895)rectangle++(0.04,0.04)
  (1.067,8.271)rectangle++(0.04,0.04) (1.235,8.584)rectangle++(0.04,0.04) (1.391,8.843)rectangle++(0.04,0.04) (1.534,9.057)rectangle++(0.04,0.04) (1.664,9.233)rectangle++(0.04,0.04)
  (1.782,9.378)rectangle++(0.04,0.04) (2.341,3.144)rectangle++(0.04,0.04) (2.114,4.175)rectangle++(0.04,0.04) (1.961,5.06)rectangle++(0.04,0.04) (1.867,5.819)rectangle++(0.04,0.04)
  (1.332,2.312)rectangle++(0.04,0.04) (1.222,3.508)rectangle++(0.04,0.04) (1.177,4.529)rectangle++(0.04,0.04) (1.074,1.823)rectangle++(0.04,0.04) (0.984,3.103)rectangle++(0.04,0.04)
  (0.958,4.194)rectangle++(0.04,0.04) (0.98,5.123)rectangle++(0.04,0.04) (1.036,5.912)rectangle++(0.04,0.04) (1.115,6.58)rectangle++(0.04,0.04) (1.208,7.144)rectangle++(0.04,0.04)
  (1.311,7.62)rectangle++(0.04,0.04) (1.417,8.021)rectangle++(0.04,0.04) (1.523,8.358)rectangle++(0.04,0.04) (1.627,8.639)rectangle++(0.04,0.04) (1.726,8.874)rectangle++(0.04,0.04)
  (1.82,9.07)rectangle++(0.04,0.04) (1.907,9.233)rectangle++(0.04,0.04) (1.988,9.368)rectangle++(0.04,0.04) (2.063,9.48)rectangle++(0.04,0.04) (2.13,9.571)rectangle++(0.04,0.04)
  (2.191,9.647)rectangle++(0.04,0.04) (2.246,9.708)rectangle++(0.04,0.04) (2.296,9.758)rectangle++(0.04,0.04) (2.339,9.799)rectangle++(0.04,0.04) (-0.02,1.551)rectangle++(0.04,0.04)
  (0.42,0.797)rectangle++(0.04,0.04) (0.387,2.257)rectangle++(0.04,0.04) (0.417,3.499)rectangle++(0.04,0.04) (0.492,4.554)rectangle++(0.04,0.04) (0.598,5.447)rectangle++(0.04,0.04)
  (0.724,6.202)rectangle++(0.04,0.04) (0.861,6.839)rectangle++(0.04,0.04) (1.004,7.375)rectangle++(0.04,0.04) (1.146,7.825)rectangle++(0.04,0.04) (1.285,8.202)rectangle++(0.04,0.04)
  (1.418,8.516)rectangle++(0.04,0.04) (1.544,8.778)rectangle++(0.04,0.04) (1.661,8.996)rectangle++(0.04,0.04) (1.77,9.176)rectangle++(0.04,0.04) (1.869,9.325)rectangle++(0.04,0.04)
  (1.959,9.448)rectangle++(0.04,0.04) (2.041,9.549)rectangle++(0.04,0.04) (2.115,9.631)rectangle++(0.04,0.04) (2.181,9.698)rectangle++(0.04,0.04) (2.371,3.324)rectangle++(0.04,0.04)
  (2.146,4.327)rectangle++(0.04,0.04) (1.995,5.188)rectangle++(0.04,0.04) (1.901,5.927)rectangle++(0.04,0.04) (1.851,6.558)rectangle++(0.04,0.04) (1.833,7.096)rectangle++(0.04,0.04)
  (1.84,7.555)rectangle++(0.04,0.04) (1.864,7.944)rectangle++(0.04,0.04) (1.9,8.274)rectangle++(0.04,0.04) (1.944,8.553)rectangle++(0.04,0.04) (1.992,8.789)rectangle++(0.04,0.04)
  (2.043,8.987)rectangle++(0.04,0.04) (2.093,9.153)rectangle++(0.04,0.04) (-1.982,4.084)rectangle++(0.04,0.04) (-1.524,5.147)rectangle++(0.04,0.04) (-1.092,6.032)rectangle++(0.04,0.04)
  (-0.689,6.767)rectangle++(0.04,0.04) (-0.317,7.376)rectangle++(0.04,0.04) (0.023,7.878)rectangle++(0.04,0.04) (0.333,8.292)rectangle++(0.04,0.04) (0.612,8.631)rectangle++(0.04,0.04)
  (-0.02,1.364)rectangle++(0.04,0.04) (-0.38,1.885)rectangle++(0.04,0.04) (-0.25,3.213)rectangle++(0.04,0.04) (-0.086,4.337)rectangle++(0.04,0.04) (0.098,5.286)rectangle++(0.04,0.04)
  (0.293,6.086)rectangle++(0.04,0.04) (0.49,6.757)rectangle++(0.04,0.04) (0.685,7.32)rectangle++(0.04,0.04) (0.873,7.791)rectangle++(0.04,0.04) (2.033,2.527)rectangle++(0.04,0.04)
  (1.827,3.663)rectangle++(0.04,0.04) (1.697,4.636)rectangle++(0.04,0.04) (1.626,5.469)rectangle++(0.04,0.04) (1.27,2.165)rectangle++(0.04,0.04) (1.164,3.386)rectangle++(0.04,0.04)
  (1.123,4.428)rectangle++(0.04,0.04) (1.129,5.315)rectangle++(0.04,0.04) (1.17,6.069)rectangle++(0.04,0.04) (1.235,6.708)rectangle++(0.04,0.04) (1.316,7.248)rectangle++(0.04,0.04)
  (1.406,7.705)rectangle++(0.04,0.04) (1.501,8.089)rectangle++(0.04,0.04) (1.598,8.412)rectangle++(0.04,0.04) (-1.889,3.807)rectangle++(0.04,0.04) (-1.455,4.908)rectangle++(0.04,0.04)
  (-1.043,5.826)rectangle++(0.04,0.04) (-0.656,6.59)rectangle++(0.04,0.04) (-0.296,7.224)rectangle++(0.04,0.04) (0.035,7.748)rectangle++(0.04,0.04) (0.338,8.181)rectangle++(0.04,0.04)
  (0.612,8.536)rectangle++(0.04,0.04) (0.86,8.828)rectangle++(0.04,0.04) (1.081,9.065)rectangle++(0.04,0.04) (-2.162,3.832)rectangle++(0.04,0.04) (-1.687,4.94)rectangle++(0.04,0.04)
  (1.619,1.177)rectangle++(0.04,0.04) (1.421,2.532)rectangle++(0.04,0.04) (1.307,3.692)rectangle++(0.04,0.04) (1.256,4.682)rectangle++(0.04,0.04) (1.253,5.525)rectangle++(0.04,0.04)
  (-1.207,3.093)rectangle++(0.04,0.04) (-0.905,4.273)rectangle++(0.04,0.04) (-0.6,5.265)rectangle++(0.04,0.04) (-0.302,6.095)rectangle++(0.04,0.04) (-0.015,6.789)rectangle++(0.04,0.04)
  (0.257,7.368)rectangle++(0.04,0.04) (0.511,7.848)rectangle++(0.04,0.04) (0.746,8.247)rectangle++(0.04,0.04) (2.18,2.603)rectangle++(0.04,0.04) (1.955,3.722)rectangle++(0.04,0.04)
  (1.808,4.681)rectangle++(0.04,0.04) (1.722,5.503)rectangle++(0.04,0.04) (1.682,6.205)rectangle++(0.04,0.04) (1.675,6.803)rectangle++(0.04,0.04) (1.694,7.312)rectangle++(0.04,0.04)
  (1.73,7.744)rectangle++(0.04,0.04) (1.778,8.109)rectangle++(0.04,0.04) (1.834,8.418)rectangle++(0.04,0.04) (1.893,8.678)rectangle++(0.04,0.04) (1.954,8.897)rectangle++(0.04,0.04)
  (2.015,9.08)rectangle++(0.04,0.04) (-1.979,4.05)rectangle++(0.04,0.04) (-1.522,5.118)rectangle++(0.04,0.04) (-1.092,6.007)rectangle++(0.04,0.04) (-0.69,6.746)rectangle++(0.04,0.04)
  (-0.319,7.358)rectangle++(0.04,0.04) (0.021,7.863)rectangle++(0.04,0.04) (0.33,8.279)rectangle++(0.04,0.04) (0.61,8.62)rectangle++(0.04,0.04) (0.861,8.899)rectangle++(0.04,0.04)
  (2.345,2.79)rectangle++(0.04,0.04) (2.103,3.874)rectangle++(0.04,0.04) (1.94,4.805)rectangle++(0.04,0.04) (1.839,5.603)rectangle++(0.04,0.04) (1.785,6.285)rectangle++(0.04,0.04)
  (1.767,6.867)rectangle++(0.04,0.04) (1.774,7.362)rectangle++(0.04,0.04) (1.8,7.783)rectangle++(0.04,0.04) (1.839,8.14)rectangle++(0.04,0.04) (-1.77,3.803)rectangle++(0.04,0.04)
  (-1.364,2.019)rectangle++(0.04,0.04) (-1.081,3.367)rectangle++(0.04,0.04) (-0.786,4.501)rectangle++(0.04,0.04) (-0.49,5.453)rectangle++(0.04,0.04) (-1.537,2.676)rectangle++(0.04,0.04)
  (-1.202,3.932)rectangle++(0.04,0.04) (-0.866,4.987)rectangle++(0.04,0.04) (-0.539,5.87)rectangle++(0.04,0.04) (-0.226,6.607)rectangle++(0.04,0.04) (0.07,7.221)rectangle++(0.04,0.04)
  (0.346,7.731)rectangle++(0.04,0.04) (0.601,8.154)rectangle++(0.04,0.04) (-2.021,3.521)rectangle++(0.04,0.04) (-1.579,4.67)rectangle++(0.04,0.04) (-1.158,5.629)rectangle++(0.04,0.04)
  (-0.761,6.427)rectangle++(0.04,0.04) (-0.392,7.09)rectangle++(0.04,0.04) (-1.943,3.059)rectangle++(0.04,0.04) (-1.531,4.274)rectangle++(0.04,0.04) (-1.133,5.29)rectangle++(0.04,0.04)
  (-0.754,6.138)rectangle++(0.04,0.04) (-0.397,6.844)rectangle++(0.04,0.04) (-0.066,7.429)rectangle++(0.04,0.04) (0.239,7.914)rectangle++(0.04,0.04) (0.517,8.313)rectangle++(0.04,0.04)
  (0.77,8.642)rectangle++(0.04,0.04) (0.998,8.911)rectangle++(0.04,0.04) (1.203,9.131)rectangle++(0.04,0.04) (1.385,9.309)rectangle++(0.04,0.04) (-2.165,3.956)rectangle++(0.04,0.04)
  (-1.684,5.045)rectangle++(0.04,0.04) (-1.232,5.952)rectangle++(0.04,0.04) (1.834,1.538)rectangle++(0.04,0.04) (1.618,2.83)rectangle++(0.04,0.04) (1.486,3.937)rectangle++(0.04,0.04)
  (-0.748,2.797)rectangle++(0.04,0.04) (-0.526,4.004)rectangle++(0.04,0.04) (-0.289,5.02)rectangle++(0.04,0.04) (-1.384,2.627)rectangle++(0.04,0.04) (-1.074,3.885)rectangle++(0.04,0.04)
  (-1.246,2.197)rectangle++(0.04,0.04) (-0.842,1.786)rectangle++(0.04,0.04) (0.609,0.64)rectangle++(0.04,0.04) (0.07,0.742)rectangle++(0.04,0.04) (0.087,2.224)rectangle++(0.04,0.04)
  (0.161,3.483)rectangle++(0.04,0.04) (0.274,4.55)rectangle++(0.04,0.04) (0.413,5.453)rectangle++(0.04,0.04) (0.567,6.215)rectangle++(0.04,0.04) (0.728,6.856)rectangle++(0.04,0.04)
  (0.891,7.395)rectangle++(0.04,0.04) (1.051,7.846)rectangle++(0.04,0.04) (1.205,8.223)rectangle++(0.04,0.04) (1.351,8.538)rectangle++(0.04,0.04) (1.488,8.799)rectangle++(0.04,0.04)
  (1.614,9.016)rectangle++(0.04,0.04) (1.731,9.195)rectangle++(0.04,0.04) (1.837,9.343)rectangle++(0.04,0.04) (1.933,9.464)rectangle++(0.04,0.04) (2.019,9.564)rectangle++(0.04,0.04)
  (2.097,9.644)rectangle++(0.04,0.04) (2.166,9.71)rectangle++(0.04,0.04) (2.227,9.763)rectangle++(0.04,0.04) (2.281,9.806)rectangle++(0.04,0.04) (2.329,9.84)rectangle++(0.04,0.04)
  (2.371,9.867)rectangle++(0.04,0.04) (2.408,9.888)rectangle++(0.04,0.04) (-2.111,4.318)rectangle++(0.04,0.04) (-1.623,5.351)rectangle++(0.04,0.04) (-1.168,6.21)rectangle++(0.04,0.04)
  (-0.747,6.921)rectangle++(0.04,0.04) (-0.36,7.509)rectangle++(0.04,0.04) (-0.008,7.993)rectangle++(0.04,0.04) (0.311,8.391)rectangle++(0.04,0.04) (0.598,8.716)rectangle++(0.04,0.04)
  (0.854,8.981)rectangle++(0.04,0.04) (1.083,9.196)rectangle++(0.04,0.04) (1.286,9.369)rectangle++(0.04,0.04) (1.466,9.509)rectangle++(0.04,0.04) (-2.2,4.018)rectangle++(0.04,0.04)
  (-1.712,5.1)rectangle++(0.04,0.04) (-1.253,5.999)rectangle++(0.04,0.04) (-0.827,6.746)rectangle++(0.04,0.04) (1.996,1.834)rectangle++(0.04,0.04) (1.767,3.075)rectangle++(0.04,0.04)
  (1.623,4.139)rectangle++(0.04,0.04) (1.543,5.05)rectangle++(0.04,0.04) (1.511,5.827)rectangle++(0.04,0.04) (1.387,2.221)rectangle++(0.04,0.04) (-0.321,2.397)rectangle++(0.04,0.04)
  (-0.179,3.646)rectangle++(0.04,0.04) (-0.009,4.703)rectangle++(0.04,0.04) (0.178,5.594)rectangle++(0.04,0.04) (0.373,6.344)rectangle++(0.04,0.04) (0.569,6.974)rectangle++(0.04,0.04)
  (0.76,7.501)rectangle++(0.04,0.04) (0.944,7.942)rectangle++(0.04,0.04) (1.118,8.309)rectangle++(0.04,0.04) (1.28,8.614)rectangle++(0.04,0.04) (1.431,8.867)rectangle++(0.04,0.04)
  (1.569,9.076)rectangle++(0.04,0.04) (2.289,3.016)rectangle++(0.04,0.04) (-0.348,2.779)rectangle++(0.04,0.04) (-0.187,3.972)rectangle++(0.04,0.04) (-0.002,4.98)rectangle++(0.04,0.04)
  (0.195,5.829)rectangle++(0.04,0.04) (1.585,1.88)rectangle++(0.04,0.04) (1.421,3.131)rectangle++(0.04,0.04) (1.331,4.2)rectangle++(0.04,0.04) (1.297,5.113)rectangle++(0.04,0.04)
  (1.305,5.891)rectangle++(0.04,0.04) (1.342,6.551)rectangle++(0.04,0.04) (-1.456,3.339)rectangle++(0.04,0.04) (-1.106,4.493)rectangle++(0.04,0.04) (-0.763,5.459)rectangle++(0.04,0.04)
  (-0.432,6.267)rectangle++(0.04,0.04) (-0.119,6.94)rectangle++(0.04,0.04) (0.174,7.5)rectangle++(0.04,0.04) (0.446,7.964)rectangle++(0.04,0.04) (0.695,8.348)rectangle++(0.04,0.04)
  (-0.02,1.319)rectangle++(0.04,0.04) (0.034,2.718)rectangle++(0.04,0.04) (0.135,3.905)rectangle++(0.04,0.04) (0.269,4.91)rectangle++(0.04,0.04) (-1.244,2.731)rectangle++(0.04,0.04)
  (-0.95,3.967)rectangle++(0.04,0.04) (-0.651,5.006)rectangle++(0.04,0.04) (-0.356,5.878)rectangle++(0.04,0.04) (-0.069,6.607)rectangle++(0.04,0.04) (0.203,7.215)rectangle++(0.04,0.04)
  (0.459,7.72)rectangle++(0.04,0.04) (0.697,8.14)rectangle++(0.04,0.04) (0.916,8.487)rectangle++(0.04,0.04) (1.116,8.774)rectangle++(0.04,0.04) (1.297,9.009)rectangle++(0.04,0.04)
  (-2.104,3.869)rectangle++(0.04,0.04) (-1.636,4.969)rectangle++(0.04,0.04) (-1.194,5.886)rectangle++(0.04,0.04) (-0.782,6.647)rectangle++(0.04,0.04) (-0.401,7.277)rectangle++(0.04,0.04)
  (-0.052,7.798)rectangle++(0.04,0.04) (0.266,8.226)rectangle++(0.04,0.04) (0.553,8.578)rectangle++(0.04,0.04) (2.302,2.632)rectangle++(0.04,0.04) (2.059,3.742)rectangle++(0.04,0.04)
  (1.898,4.694)rectangle++(0.04,0.04) (1.799,5.51)rectangle++(0.04,0.04) (1.747,6.208)rectangle++(0.04,0.04) (1.731,6.803)rectangle++(0.04,0.04) (1.742,7.31)rectangle++(0.04,0.04)
  (1.77,7.74)rectangle++(0.04,0.04) (1.812,8.104)rectangle++(0.04,0.04) (1.862,8.412)rectangle++(0.04,0.04) (1.917,8.672)rectangle++(0.04,0.04) (1.974,8.891)rectangle++(0.04,0.04)
  (2.032,9.074)rectangle++(0.04,0.04) (2.088,9.228)rectangle++(0.04,0.04) (2.141,9.357)rectangle++(0.04,0.04) (2.192,9.464)rectangle++(0.04,0.04) (2.24,9.553)rectangle++(0.04,0.04)
  (2.284,9.626)rectangle++(0.04,0.04) (2.324,9.687)rectangle++(0.04,0.04) (2.361,9.737)rectangle++(0.04,0.04) (2.394,9.779)rectangle++(0.04,0.04) (2.424,9.812)rectangle++(0.04,0.04)
  (2.45,9.84)rectangle++(0.04,0.04) (2.474,9.862)rectangle++(0.04,0.04) (2.495,9.88)rectangle++(0.04,0.04) (2.514,9.894)rectangle++(0.04,0.04) (2.531,9.906)rectangle++(0.04,0.04)
  (2.545,9.915)rectangle++(0.04,0.04) (2.558,9.922)rectangle++(0.04,0.04) (2.569,9.928)rectangle++(0.04,0.04) (2.578,9.932)rectangle++(0.04,0.04) (2.587,9.935)rectangle++(0.04,0.04)
  (2.594,9.938)rectangle++(0.04,0.04) (2.6,9.939)rectangle++(0.04,0.04) (2.605,9.941)rectangle++(0.04,0.04) (2.61,9.942)rectangle++(0.04,0.04) (2.614,9.942)rectangle++(0.04,0.04)
  (2.617,9.942)rectangle++(0.04,0.04) (2.62,9.943)rectangle++(0.04,0.04) (2.373,3.498)rectangle++(0.04,0.04) (2.155,4.474)rectangle++(0.04,0.04) (2.009,5.313)rectangle++(0.04,0.04)
  (1.918,6.032)rectangle++(0.04,0.04) (1.869,6.647)rectangle++(0.04,0.04) (1.852,7.171)rectangle++(0.04,0.04) (1.859,7.618)rectangle++(0.04,0.04) (1.883,7.997)rectangle++(0.04,0.04)
  (1.918,8.318)rectangle++(0.04,0.04) (1.961,8.59)rectangle++(0.04,0.04) (2.008,8.819)rectangle++(0.04,0.04) (-1.913,3.991)rectangle++(0.04,0.04) (-1.468,5.065)rectangle++(0.04,0.04)
  (-1.048,5.96)rectangle++(0.04,0.04) (-0.654,6.704)rectangle++(0.04,0.04) (-0.29,7.321)rectangle++(0.04,0.04) (-1.983,3.133)rectangle++(0.04,0.04) (-1.562,4.338)rectangle++(0.04,0.04)
  (-1.157,5.346)rectangle++(0.04,0.04) (-0.771,6.187)rectangle++(0.04,0.04) (-0.41,6.886)rectangle++(0.04,0.04) (-1.894,3.01)rectangle++(0.04,0.04) (-1.491,4.23)rectangle++(0.04,0.04)
  (-1.101,5.251)rectangle++(0.04,0.04) (-0.728,6.104)rectangle++(0.04,0.04) (-0.377,6.814)rectangle++(0.04,0.04) (-0.05,7.403)rectangle++(0.04,0.04) (0.252,7.891)rectangle++(0.04,0.04)
  (0.527,8.293)rectangle++(0.04,0.04) (0.778,8.624)rectangle++(0.04,0.04) (1.004,8.896)rectangle++(0.04,0.04) (1.207,9.117)rectangle++(0.04,0.04) (1.388,9.298)rectangle++(0.04,0.04)
  (2.378,3.022)rectangle++(0.04,0.04) (-0.331,2.801)rectangle++(0.04,0.04) (-0.172,3.99)rectangle++(0.04,0.04) (0.011,4.995)rectangle++(0.04,0.04) (0.207,5.841)rectangle++(0.04,0.04)
  (0.408,6.553)rectangle++(0.04,0.04) (0.606,7.15)rectangle++(0.04,0.04) (0.799,7.649)rectangle++(0.04,0.04) (0.983,8.066)rectangle++(0.04,0.04) (1.156,8.413)rectangle++(0.04,0.04)
  (1.317,8.701)rectangle++(0.04,0.04) (1.465,8.939)rectangle++(0.04,0.04) (1.601,9.136)rectangle++(0.04,0.04) (2.301,3.039)rectangle++(0.04,0.04) (2.075,4.087)rectangle++(0.04,0.04)
  (0.816,1.95)rectangle++(0.04,0.04) (0.769,3.221)rectangle++(0.04,0.04) (0.769,1.403)rectangle++(0.04,0.04) (0.708,2.758)rectangle++(0.04,0.04) (-0.597,2.359)rectangle++(0.04,0.04)
  (-0.415,3.625)rectangle++(0.04,0.04) (-0.21,4.694)rectangle++(0.04,0.04) (0.007,5.594)rectangle++(0.04,0.04) (0.228,6.351)rectangle++(0.04,0.04) (0.445,6.986)rectangle++(0.04,0.04)
  (0.656,7.516)rectangle++(0.04,0.04) (0.856,7.959)rectangle++(0.04,0.04) (1.044,8.327)rectangle++(0.04,0.04) (1.218,8.632)rectangle++(0.04,0.04) (1.378,8.885)rectangle++(0.04,0.04)
  (1.525,9.093)rectangle++(0.04,0.04) (1.658,9.264)rectangle++(0.04,0.04) (-2.098,4.008)rectangle++(0.04,0.04) (-1.626,5.087)rectangle++(0.04,0.04) (-1.18,5.985)rectangle++(0.04,0.04)
  (1.836,1.56)rectangle++(0.04,0.04) (1.62,2.848)rectangle++(0.04,0.04) (1.489,3.953)rectangle++(0.04,0.04) (1.422,4.896)rectangle++(0.04,0.04) (1.402,5.701)rectangle++(0.04,0.04)
  (1.418,6.386)rectangle++(0.04,0.04) (1.458,6.968)rectangle++(0.04,0.04) (1.516,7.46)rectangle++(0.04,0.04) (-1.658,3.579)rectangle++(0.04,0.04) (-1.268,4.705)rectangle++(0.04,0.04)
  (-0.892,5.646)rectangle++(0.04,0.04) (-0.534,6.431)rectangle++(0.04,0.04) (-0.199,7.084)rectangle++(0.04,0.04) (0.112,7.625)rectangle++(0.04,0.04) (0.398,8.073)rectangle++(0.04,0.04)
  (0.659,8.443)rectangle++(0.04,0.04) (0.896,8.746)rectangle++(0.04,0.04) (1.109,8.995)rectangle++(0.04,0.04) (1.3,9.197)rectangle++(0.04,0.04) (1.471,9.362)rectangle++(0.04,0.04)
  (1.622,9.495)rectangle++(0.04,0.04) (1.757,9.602)rectangle++(0.04,0.04) (1.875,9.688)rectangle++(0.04,0.04) (1.979,9.756)rectangle++(0.04,0.04) (2.07,9.809)rectangle++(0.04,0.04)
  (2.419,3.323)rectangle++(0.04,0.04) (0.55,1.856)rectangle++(0.04,0.04) (-0.394,2.124)rectangle++(0.04,0.04) (-0.652,1.966)rectangle++(0.04,0.04) (-0.478,3.293)rectangle++(0.04,0.04)
  (-0.277,4.414)rectangle++(0.04,0.04) (1.26,1.418)rectangle++(0.04,0.04) (1.126,2.751)rectangle++(0.04,0.04) (1.065,3.889)rectangle++(0.04,0.04) (0.912,1.64)rectangle++(0.04,0.04)
  (-0.265,2.16)rectangle++(0.04,0.04) (-0.141,3.442)rectangle++(0.04,0.04) (0.015,4.528)rectangle++(0.04,0.04) (0.192,5.444)rectangle++(0.04,0.04) (0.379,6.216)rectangle++(0.04,0.04)
  (0.568,6.865)rectangle++(0.04,0.04) (0.756,7.409)rectangle++(0.04,0.04) (-1.796,3.393)rectangle++(0.04,0.04) (1.202,0.777)rectangle++(0.04,0.04) (1.051,2.209)rectangle++(0.04,0.04)
  (0.979,3.432)rectangle++(0.04,0.04) (0.967,4.474)rectangle++(0.04,0.04) (0.999,5.36)rectangle++(0.04,0.04) (1.061,6.113)rectangle++(0.04,0.04) (1.144,6.749)rectangle++(0.04,0.04)
  (1.241,7.287)rectangle++(0.04,0.04) (1.344,7.741)rectangle++(0.04,0.04) (1.45,8.122)rectangle++(0.04,0.04) (-1.843,3.709)rectangle++(0.04,0.04) (1.298,0.841)rectangle++(0.04,0.04)
  (1.134,2.259)rectangle++(0.04,0.04) (1.052,3.471)rectangle++(0.04,0.04) (1.031,4.505)rectangle++(0.04,0.04) (1.055,5.384)rectangle++(0.04,0.04) (1.11,6.13)rectangle++(0.04,0.04)
  (1.186,6.763)rectangle++(0.04,0.04) (1.276,7.297)rectangle++(0.04,0.04) (1.834,2.513)rectangle++(0.04,0.04) (1.657,3.659)rectangle++(0.04,0.04) (-0.641,2.775)rectangle++(0.04,0.04)
  (-0.436,3.981)rectangle++(0.04,0.04) (-0.214,4.997)rectangle++(0.04,0.04) (0.016,5.852)rectangle++(0.04,0.04) (0.246,6.57)rectangle++(0.04,0.04) (0.469,7.171)rectangle++(0.04,0.04)
  (-1.792,3.275)rectangle++(0.04,0.04) (-1.394,4.451)rectangle++(0.04,0.04) (-1.009,5.436)rectangle++(0.04,0.04) (-0.643,6.257)rectangle++(0.04,0.04) (-0.298,6.94)rectangle++(0.04,0.04)
  (-0.02,1.094)rectangle++(0.04,0.04) (0.025,2.527)rectangle++(0.04,0.04) (0.12,3.743)rectangle++(0.04,0.04) (0.249,4.773)rectangle++(0.04,0.04) (0.401,5.643)rectangle++(0.04,0.04)
  (0.564,6.377)rectangle++(0.04,0.04) (0.732,6.994)rectangle++(0.04,0.04) (0.9,7.512)rectangle++(0.04,0.04) (1.063,7.945)rectangle++(0.04,0.04) (1.219,8.307)rectangle++(0.04,0.04)
  (1.367,8.608)rectangle++(0.04,0.04) (1.504,8.859)rectangle++(0.04,0.04) (1.63,9.066)rectangle++(0.04,0.04) (1.746,9.237)rectangle++(0.04,0.04) (1.852,9.378)rectangle++(0.04,0.04)
  (1.947,9.493)rectangle++(0.04,0.04) (2.032,9.588)rectangle++(0.04,0.04) (2.109,9.664)rectangle++(0.04,0.04) (2.177,9.727)rectangle++(0.04,0.04) (2.237,9.777)rectangle++(0.04,0.04)
  (2.29,9.817)rectangle++(0.04,0.04) (2.337,9.849)rectangle++(0.04,0.04) (2.379,9.874)rectangle++(0.04,0.04) (2.415,9.894)rectangle++(0.04,0.04) (2.446,9.91)rectangle++(0.04,0.04)
  (2.473,9.922)rectangle++(0.04,0.04) (2.497,9.931)rectangle++(0.04,0.04) (2.517,9.937)rectangle++(0.04,0.04) (2.535,9.942)rectangle++(0.04,0.04) (2.55,9.946)rectangle++(0.04,0.04)
  (2.385,3.48)rectangle++(0.04,0.04) (2.164,4.459)rectangle++(0.04,0.04) (2.016,5.3)rectangle++(0.04,0.04) (1.923,6.02)rectangle++(0.04,0.04) (1.873,6.637)rectangle++(0.04,0.04)
  (1.856,7.162)rectangle++(0.04,0.04) (1.861,7.61)rectangle++(0.04,0.04) (1.884,7.99)rectangle++(0.04,0.04) (1.919,8.312)rectangle++(0.04,0.04) (1.962,8.585)rectangle++(0.04,0.04)
  (2.009,8.815)rectangle++(0.04,0.04) (2.058,9.009)rectangle++(0.04,0.04) (2.107,9.171)rectangle++(0.04,0.04) (2.156,9.307)rectangle++(0.04,0.04) (2.202,9.421)rectangle++(0.04,0.04)
  (-2.03,4.168)rectangle++(0.04,0.04) (-1.561,5.22)rectangle++(0.04,0.04) (-1.12,6.096)rectangle++(0.04,0.04) (-0.711,6.823)rectangle++(0.04,0.04) (2,1.883)rectangle++(0.04,0.04)
  (1.773,3.116)rectangle++(0.04,0.04) (1.629,4.174)rectangle++(0.04,0.04) (1.55,5.079)rectangle++(0.04,0.04) (1.518,5.852)rectangle++(0.04,0.04) (1.522,6.509)rectangle++(0.04,0.04)
  (1.552,7.068)rectangle++(0.04,0.04) (1.6,7.542)rectangle++(0.04,0.04) (1.659,7.943)rectangle++(0.04,0.04) (1.726,8.281)rectangle++(0.04,0.04) (1.796,8.566)rectangle++(0.04,0.04)
  (1.867,8.806)rectangle++(0.04,0.04) (1.937,9.006)rectangle++(0.04,0.04) (2.005,9.174)rectangle++(0.04,0.04) (2.069,9.314)rectangle++(0.04,0.04) (2.129,9.43)rectangle++(0.04,0.04)
  (2.184,9.527)rectangle++(0.04,0.04) (2.236,9.607)rectangle++(0.04,0.04) (2.282,9.672)rectangle++(0.04,0.04) (2.325,9.726)rectangle++(0.04,0.04) (2.363,9.771)rectangle++(0.04,0.04)
  (2.397,9.807)rectangle++(0.04,0.04) (2.428,9.836)rectangle++(0.04,0.04) (2.455,9.86)rectangle++(0.04,0.04) (2.479,9.879)rectangle++(0.04,0.04) (2.5,9.894)rectangle++(0.04,0.04)
  (2.518,9.906)rectangle++(0.04,0.04) (2.535,9.916)rectangle++(0.04,0.04) (2.379,3.469)rectangle++(0.04,0.04) (2.159,4.45)rectangle++(0.04,0.04) (2.011,5.292)rectangle++(0.04,0.04)
  (1.918,6.014)rectangle++(0.04,0.04) (1.869,6.631)rectangle++(0.04,0.04) (1.852,7.158)rectangle++(0.04,0.04) (1.858,7.606)rectangle++(0.04,0.04) (1.881,7.987)rectangle++(0.04,0.04)
  (1.917,8.31)rectangle++(0.04,0.04) (1.959,8.583)rectangle++(0.04,0.04) (2.006,8.814)rectangle++(0.04,0.04) (2.056,9.007)rectangle++(0.04,0.04) (2.106,9.17)rectangle++(0.04,0.04)
  (2.234,3.178)rectangle++(0.04,0.04) (2.024,4.208)rectangle++(0.04,0.04) (1.887,5.092)rectangle++(0.04,0.04) (1.805,5.849)rectangle++(0.04,0.04) (1.766,6.496)rectangle++(0.04,0.04)
  (1.537,2.448)rectangle++(0.04,0.04) (1.402,3.616)rectangle++(0.04,0.04) (-0.681,2.707)rectangle++(0.04,0.04) (-0.473,3.924)rectangle++(0.04,0.04) (1.152,1.249)rectangle++(0.04,0.04)
  (1.027,2.612)rectangle++(0.04,0.04) (0.975,3.775)rectangle++(0.04,0.04) (0.893,1.59)rectangle++(0.04,0.04) (0.821,2.912)rectangle++(0.04,0.04) (0.812,4.038)rectangle++(0.04,0.04)
  (0.849,4.996)rectangle++(0.04,0.04) (0.92,5.809)rectangle++(0.04,0.04) (1.012,6.497)rectangle++(0.04,0.04) (1.118,7.078)rectangle++(0.04,0.04) (1.231,7.568)rectangle++(0.04,0.04)
  (-1.743,3.537)rectangle++(0.04,0.04) (-1.342,4.672)rectangle++(0.04,0.04) (-0.956,5.621)rectangle++(0.04,0.04) (-0.59,6.413)rectangle++(0.04,0.04) (-0.247,7.071)rectangle++(0.04,0.04)
  (0.071,7.616)rectangle++(0.04,0.04) (0.362,8.067)rectangle++(0.04,0.04) (0.629,8.439)rectangle++(0.04,0.04) (0.87,8.744)rectangle++(0.04,0.04) (1.087,8.994)rectangle++(0.04,0.04)
  (1.281,9.197)rectangle++(0.04,0.04) (2.366,2.971)rectangle++(0.04,0.04) (2.127,4.027)rectangle++(0.04,0.04) (-0.643,2.964)rectangle++(0.04,0.04) (-0.43,4.141)rectangle++(0.04,0.04)
  (-0.202,5.134)rectangle++(0.04,0.04) (0.032,5.968)rectangle++(0.04,0.04) (0.263,6.668)rectangle++(0.04,0.04) (0.488,7.253)rectangle++(0.04,0.04) (0.703,7.742)rectangle++(0.04,0.04)
  (0.905,8.149)rectangle++(0.04,0.04) (1.093,8.486)rectangle++(0.04,0.04) (1.266,8.766)rectangle++(0.04,0.04) (1.425,8.997)rectangle++(0.04,0.04) (1.569,9.186)rectangle++(0.04,0.04)
  (-2.096,3.971)rectangle++(0.04,0.04) (1.409,0.838)rectangle++(0.04,0.04) (1.229,2.252)rectangle++(0.04,0.04) (1.132,3.461)rectangle++(0.04,0.04) (1.099,4.493)rectangle++(0.04,0.04)
  (1.111,5.371)rectangle++(0.04,0.04) (1.32,2.008)rectangle++(0.04,0.04) (1.2,3.25)rectangle++(0.04,0.04) (1.148,4.311)rectangle++(0.04,0.04) (1.146,5.215)rectangle++(0.04,0.04)
  (-1.148,3)rectangle++(0.04,0.04) (-0.858,4.192)rectangle++(0.04,0.04) (-0.564,5.194)rectangle++(0.04,0.04) (-0.274,6.033)rectangle++(0.04,0.04) (0.007,6.735)rectangle++(0.04,0.04)
  (0.273,7.321)rectangle++(0.04,0.04) (0.523,7.808)rectangle++(0.04,0.04) (0.754,8.212)rectangle++(0.04,0.04) (0.967,8.546)rectangle++(0.04,0.04) (1.162,8.822)rectangle++(0.04,0.04);
\fill[nuage]
  (1.338,9.048)rectangle++(0.04,0.04) (1.497,9.234)rectangle++(0.04,0.04) (1.64,9.385)rectangle++(0.04,0.04) (1.767,9.508)rectangle++(0.04,0.04) (1.88,9.607)rectangle++(0.04,0.04)
  (1.98,9.687)rectangle++(0.04,0.04) (2.068,9.751)rectangle++(0.04,0.04) (2.146,9.802)rectangle++(0.04,0.04) (2.405,3.34)rectangle++(0.04,0.04) (2.176,4.339)rectangle++(0.04,0.04)
  (2.021,5.198)rectangle++(0.04,0.04) (1.923,5.933)rectangle++(0.04,0.04) (1.87,6.563)rectangle++(0.04,0.04) (1.85,7.1)rectangle++(0.04,0.04) (1.854,7.557)rectangle++(0.04,0.04)
  (1.876,7.945)rectangle++(0.04,0.04) (1.926,2.825)rectangle++(0.04,0.04) (1.748,3.92)rectangle++(0.04,0.04) (1.64,4.858)rectangle++(0.04,0.04) (1.586,5.66)rectangle++(0.04,0.04)
  (1.573,6.344)rectangle++(0.04,0.04) (1.588,6.926)rectangle++(0.04,0.04) (1.625,7.419)rectangle++(0.04,0.04) (1.676,7.838)rectangle++(0.04,0.04) (1.926,2.747)rectangle++(0.04,0.04)
  (1.745,3.854)rectangle++(0.04,0.04) (1.635,4.802)rectangle++(0.04,0.04) (1.082,2.008)rectangle++(0.04,0.04) (0.998,3.259)rectangle++(0.04,0.04) (0.976,4.327)rectangle++(0.04,0.04)
  (1.001,5.235)rectangle++(0.04,0.04) (1.058,6.006)rectangle++(0.04,0.04) (1.137,6.659)rectangle++(0.04,0.04) (1.231,7.211)rectangle++(0.04,0.04) (1.332,7.676)rectangle++(0.04,0.04)
  (1.437,8.068)rectangle++(0.04,0.04) (1.542,8.396)rectangle++(0.04,0.04) (1.645,8.671)rectangle++(0.04,0.04) (2.164,2.939)rectangle++(0.04,0.04) (1.955,4.008)rectangle++(0.04,0.04)
  (1.82,4.924)rectangle++(0.04,0.04) (1.741,5.709)rectangle++(0.04,0.04) (1.706,6.379)rectangle++(0.04,0.04) (1.703,6.95)rectangle++(0.04,0.04) (1.724,7.436)rectangle++(0.04,0.04)
  (1.76,7.848)rectangle++(0.04,0.04) (1.808,8.196)rectangle++(0.04,0.04) (1.862,8.491)rectangle++(0.04,0.04) (1.921,8.739)rectangle++(0.04,0.04) (1.98,8.947)rectangle++(0.04,0.04)
  (2.191,3.092)rectangle++(0.04,0.04) (1.984,4.137)rectangle++(0.04,0.04) (-0.02,0.645)rectangle++(0.04,0.04) (0.007,2.145)rectangle++(0.04,0.04) (0.582,0.947)rectangle++(0.04,0.04)
  (0.531,2.378)rectangle++(0.04,0.04) (0.544,3.596)rectangle++(0.04,0.04) (0.908,1.434)rectangle++(0.04,0.04) (0.827,2.779)rectangle++(0.04,0.04) (0.812,3.925)rectangle++(0.04,0.04)
  (0.845,4.9)rectangle++(0.04,0.04) (-1.126,2.861)rectangle++(0.04,0.04) (-0.845,4.073)rectangle++(0.04,0.04) (-0.558,5.092)rectangle++(0.04,0.04) (-0.272,5.947)rectangle++(0.04,0.04)
  (0.004,6.662)rectangle++(0.04,0.04) (0.268,7.259)rectangle++(0.04,0.04) (0.516,7.755)rectangle++(0.04,0.04) (0.746,8.168)rectangle++(0.04,0.04) (0.959,8.509)rectangle++(0.04,0.04)
  (1.153,8.79)rectangle++(0.04,0.04) (1.33,9.022)rectangle++(0.04,0.04) (1.489,9.212)rectangle++(0.04,0.04) (1.632,9.367)rectangle++(0.04,0.04) (1.76,9.492)rectangle++(0.04,0.04)
  (1.873,9.594)rectangle++(0.04,0.04) (1.974,9.677)rectangle++(0.04,0.04) (2.063,9.742)rectangle++(0.04,0.04) (2.401,3.304)rectangle++(0.04,0.04) (0.548,1.847)rectangle++(0.04,0.04)
  (0.537,3.145)rectangle++(0.04,0.04) (-0.731,2.404)rectangle++(0.04,0.04) (-0.528,3.669)rectangle++(0.04,0.04) (-0.304,4.736)rectangle++(0.04,0.04) (-0.02,0.741)rectangle++(0.04,0.04)
  (0.01,2.227)rectangle++(0.04,0.04) (0.096,3.489)rectangle++(0.04,0.04) (0.219,4.558)rectangle++(0.04,0.04) (0.366,5.461)rectangle++(0.04,0.04) (0.527,6.224)rectangle++(0.04,0.04)
  (0.695,6.865)rectangle++(0.04,0.04) (1.801,2.258)rectangle++(0.04,0.04) (0.345,1.44)rectangle++(0.04,0.04) (0.349,2.807)rectangle++(0.04,0.04) (0.406,3.968)rectangle++(0.04,0.04)
  (0.502,4.953)rectangle++(0.04,0.04) (0.623,5.786)rectangle++(0.04,0.04) (-0.02,0.909)rectangle++(0.04,0.04) (0.017,2.37)rectangle++(0.04,0.04) (0.107,3.61)rectangle++(0.04,0.04)
  (0.233,4.66)rectangle++(0.04,0.04) (0.382,5.548)rectangle++(0.04,0.04) (0.545,6.297)rectangle++(0.04,0.04) (0.713,6.927)rectangle++(0.04,0.04) (0.881,7.455)rectangle++(0.04,0.04)
  (1.045,7.898)rectangle++(0.04,0.04) (1.202,8.268)rectangle++(0.04,0.04) (1.35,8.576)rectangle++(0.04,0.04) (1.488,8.832)rectangle++(0.04,0.04) (1.616,9.043)rectangle++(0.04,0.04)
  (1.733,9.218)rectangle++(0.04,0.04) (1.84,9.363)rectangle++(0.04,0.04) (1.936,9.481)rectangle++(0.04,0.04) (2.023,9.577)rectangle++(0.04,0.04) (2.1,9.656)rectangle++(0.04,0.04)
  (2.169,9.72)rectangle++(0.04,0.04) (2.23,9.771)rectangle++(0.04,0.04) (2.285,9.813)rectangle++(0.04,0.04) (2.332,9.846)rectangle++(0.04,0.04) (2.374,9.872)rectangle++(0.04,0.04)
  (2.411,9.892)rectangle++(0.04,0.04) (2.442,9.908)rectangle++(0.04,0.04) (2.391,3.443)rectangle++(0.04,0.04) (2.167,4.427)rectangle++(0.04,0.04) (2.017,5.273)rectangle++(0.04,0.04)
  (1.923,5.997)rectangle++(0.04,0.04) (1.873,6.617)rectangle++(0.04,0.04) (1.854,7.146)rectangle++(0.04,0.04) (-1.508,3.587)rectangle++(0.04,0.04) (-1.141,4.706)rectangle++(0.04,0.04)
  (-1.473,2.362)rectangle++(0.04,0.04) (-1.16,3.663)rectangle++(0.04,0.04) (-0.841,4.756)rectangle++(0.04,0.04) (-0.527,5.672)rectangle++(0.04,0.04) (-0.223,6.439)rectangle++(0.04,0.04)
  (0.065,7.078)rectangle++(0.04,0.04) (1.955,2.146)rectangle++(0.04,0.04) (-0.02,0.327)rectangle++(0.04,0.04) (-0.006,1.875)rectangle++(0.04,0.04) (0.068,3.19)rectangle++(0.04,0.04)
  (0.183,4.305)rectangle++(0.04,0.04) (0.325,5.248)rectangle++(0.04,0.04) (-1.321,2.818)rectangle++(0.04,0.04) (-1.012,4.045)rectangle++(0.04,0.04) (-1.275,2.246)rectangle++(0.04,0.04)
  (-0.996,3.556)rectangle++(0.04,0.04) (-0.707,4.659)rectangle++(0.04,0.04) (1.393,1.364)rectangle++(0.04,0.04) (1.237,2.7)rectangle++(0.04,0.04) (1.157,3.842)rectangle++(0.04,0.04)
  (1.135,4.816)rectangle++(0.04,0.04) (1.155,5.644)rectangle++(0.04,0.04) (1.39,2.085)rectangle++(0.04,0.04) (1.262,3.313)rectangle++(0.04,0.04) (1.203,4.362)rectangle++(0.04,0.04)
  (1.195,5.255)rectangle++(0.04,0.04) (1.224,6.015)rectangle++(0.04,0.04) (1.279,6.66)rectangle++(0.04,0.04) (1.351,7.206)rectangle++(0.04,0.04) (1.435,7.668)rectangle++(0.04,0.04)
  (1.524,8.056)rectangle++(0.04,0.04) (1.615,8.383)rectangle++(0.04,0.04) (1.706,8.657)rectangle++(0.04,0.04) (1.794,8.887)rectangle++(0.04,0.04) (1.878,9.078)rectangle++(0.04,0.04)
  (1.958,9.237)rectangle++(0.04,0.04) (2.031,9.37)rectangle++(0.04,0.04) (2.099,9.479)rectangle++(0.04,0.04) (2.161,9.57)rectangle++(0.04,0.04) (2.218,9.644)rectangle++(0.04,0.04)
  (2.269,9.705)rectangle++(0.04,0.04) (2.314,9.755)rectangle++(0.04,0.04) (2.355,9.795)rectangle++(0.04,0.04) (2.391,9.828)rectangle++(0.04,0.04) (2.424,9.854)rectangle++(0.04,0.04)
  (2.452,9.875)rectangle++(0.04,0.04) (2.477,9.892)rectangle++(0.04,0.04) (2.499,9.905)rectangle++(0.04,0.04) (2.518,9.916)rectangle++(0.04,0.04) (2.535,9.924)rectangle++(0.04,0.04)
  (2.549,9.93)rectangle++(0.04,0.04) (2.562,9.935)rectangle++(0.04,0.04) (2.573,9.938)rectangle++(0.04,0.04) (2.582,9.941)rectangle++(0.04,0.04) (2.59,9.943)rectangle++(0.04,0.04)
  (2.597,9.944)rectangle++(0.04,0.04) (2.603,9.945)rectangle++(0.04,0.04) (2.608,9.945)rectangle++(0.04,0.04) (2.613,9.945)rectangle++(0.04,0.04) (2.616,9.945)rectangle++(0.04,0.04)
  (2.62,9.945)rectangle++(0.04,0.04) (2.374,3.498)rectangle++(0.04,0.04) (2.156,4.474)rectangle++(0.04,0.04) (2.009,5.313)rectangle++(0.04,0.04) (-1.001,3.22)rectangle++(0.04,0.04)
  (1.034,0.943)rectangle++(0.04,0.04) (0.915,2.356)rectangle++(0.04,0.04) (0.505,1.233)rectangle++(0.04,0.04) (0.252,0.857)rectangle++(0.04,0.04) (0.246,2.315)rectangle++(0.04,0.04)
  (0.3,3.554)rectangle++(0.04,0.04) (0.395,4.605)rectangle++(0.04,0.04) (0.518,5.495)rectangle++(0.04,0.04) (0.658,6.246)rectangle++(0.04,0.04) (-1.514,3.114)rectangle++(0.04,0.04)
  (-1.164,4.304)rectangle++(0.04,0.04) (-0.82,5.301)rectangle++(0.04,0.04) (-0.487,6.135)rectangle++(0.04,0.04) (-0.171,6.83)rectangle++(0.04,0.04) (0.126,7.409)rectangle++(0.04,0.04)
  (0.401,7.889)rectangle++(0.04,0.04) (0.654,8.286)rectangle++(0.04,0.04) (0.885,8.613)rectangle++(0.04,0.04) (1.095,8.882)rectangle++(0.04,0.04) (1.284,9.102)rectangle++(0.04,0.04)
  (1.453,9.281)rectangle++(0.04,0.04) (1.604,9.427)rectangle++(0.04,0.04) (1.738,9.545)rectangle++(0.04,0.04) (1.857,9.64)rectangle++(0.04,0.04) (1.962,9.716)rectangle++(0.04,0.04)
  (2.054,9.776)rectangle++(0.04,0.04) (2.135,9.824)rectangle++(0.04,0.04) (2.205,9.861)rectangle++(0.04,0.04) (2.267,9.89)rectangle++(0.04,0.04) (2.32,9.912)rectangle++(0.04,0.04)
  (2.366,9.929)rectangle++(0.04,0.04) (2.406,9.941)rectangle++(0.04,0.04) (2.441,9.95)rectangle++(0.04,0.04) (2.471,9.956)rectangle++(0.04,0.04) (2.496,9.96)rectangle++(0.04,0.04)
  (2.518,9.962)rectangle++(0.04,0.04) (2.536,9.963)rectangle++(0.04,0.04) (2.552,9.964)rectangle++(0.04,0.04) (2.566,9.963)rectangle++(0.04,0.04) (2.577,9.962)rectangle++(0.04,0.04)
  (2.587,9.961)rectangle++(0.04,0.04) (2.595,9.96)rectangle++(0.04,0.04) (2.602,9.958)rectangle++(0.04,0.04) (2.608,9.956)rectangle++(0.04,0.04) (2.613,9.955)rectangle++(0.04,0.04)
  (-2.087,4.38)rectangle++(0.04,0.04) (-1.601,5.403)rectangle++(0.04,0.04) (-1.147,6.253)rectangle++(0.04,0.04) (-0.727,6.957)rectangle++(0.04,0.04) (-0.342,7.538)rectangle++(0.04,0.04)
  (0.009,8.018)rectangle++(0.04,0.04) (0.326,8.411)rectangle++(0.04,0.04) (0.611,8.732)rectangle++(0.04,0.04) (0.867,8.994)rectangle++(0.04,0.04) (1.094,9.207)rectangle++(0.04,0.04)
  (1.296,9.378)rectangle++(0.04,0.04) (-2.2,3.95)rectangle++(0.04,0.04) (1.419,0.806)rectangle++(0.04,0.04) (1.236,2.225)rectangle++(0.04,0.04) (1.137,3.438)rectangle++(0.04,0.04)
  (0.775,1.551)rectangle++(0.04,0.04) (0.718,2.883)rectangle++(0.04,0.04) (0.682,1.309)rectangle++(0.04,0.04) (0.63,2.681)rectangle++(0.04,0.04) (0.639,1.237)rectangle++(0.04,0.04)
  (0.59,2.622)rectangle++(0.04,0.04) (0.604,3.802)rectangle++(0.04,0.04) (0.664,4.803)rectangle++(0.04,0.04) (1.228,1.755)rectangle++(0.04,0.04) (1.112,3.039)rectangle++(0.04,0.04)
  (-0.02,0.469)rectangle++(0.04,0.04) (0,1.996)rectangle++(0.04,0.04) (0.077,3.293)rectangle++(0.04,0.04) (0.195,4.392)rectangle++(0.04,0.04) (0.339,5.322)rectangle++(0.04,0.04)
  (0.499,6.106)rectangle++(0.04,0.04) (0.666,6.766)rectangle++(0.04,0.04) (0.835,7.321)rectangle++(0.04,0.04) (1,7.786)rectangle++(0.04,0.04) (1.159,8.174)rectangle++(0.04,0.04)
  (1.31,8.498)rectangle++(0.04,0.04) (1.451,8.767)rectangle++(0.04,0.04) (1.582,8.99)rectangle++(0.04,0.04) (2.262,2.999)rectangle++(0.04,0.04) (2.041,4.055)rectangle++(0.04,0.04)
  (0.812,1.934)rectangle++(0.04,0.04) (-0.362,2.201)rectangle++(0.04,0.04) (-0.222,3.482)rectangle++(0.04,0.04) (-0.051,4.564)rectangle++(0.04,0.04) (1.268,1.512)rectangle++(0.04,0.04)
  (1.136,2.831)rectangle++(0.04,0.04) (1.077,3.957)rectangle++(0.04,0.04) (-0.02,0.616)rectangle++(0.04,0.04) (0.005,2.121)rectangle++(0.04,0.04) (-0.572,2.057)rectangle++(0.04,0.04)
  (-0.406,3.367)rectangle++(0.04,0.04) (-0.212,4.475)rectangle++(0.04,0.04) (-0.004,5.408)rectangle++(0.04,0.04) (0.211,6.193)rectangle++(0.04,0.04) (0.425,6.852)rectangle++(0.04,0.04)
  (0.633,7.403)rectangle++(0.04,0.04) (0.832,7.864)rectangle++(0.04,0.04) (1.02,8.247)rectangle++(0.04,0.04) (1.194,8.566)rectangle++(0.04,0.04) (-0.02,1.354)rectangle++(0.04,0.04)
  (0.035,2.748)rectangle++(0.04,0.04) (0.137,3.93)rectangle++(0.04,0.04) (0.272,4.931)rectangle++(0.04,0.04) (0.426,5.777)rectangle++(0.04,0.04) (0.591,6.49)rectangle++(0.04,0.04)
  (0.76,7.089)rectangle++(0.04,0.04) (0.927,7.591)rectangle++(0.04,0.04) (1.09,8.012)rectangle++(0.04,0.04) (1.244,8.363)rectangle++(0.04,0.04) (1.39,8.655)rectangle++(0.04,0.04)
  (1.526,8.897)rectangle++(0.04,0.04) (-2.029,3.897)rectangle++(0.04,0.04) (-1.571,4.99)rectangle++(0.04,0.04) (-1.138,5.901)rectangle++(0.04,0.04) (-0.734,6.657)rectangle++(0.04,0.04)
  (-1.899,2.885)rectangle++(0.04,0.04) (-1.501,4.124)rectangle++(0.04,0.04) (-1.113,5.162)rectangle++(0.04,0.04) (-0.742,6.028)rectangle++(0.04,0.04) (-0.391,6.75)rectangle++(0.04,0.04)
  (-0.065,7.349)rectangle++(0.04,0.04) (0.237,7.846)rectangle++(0.04,0.04) (0.513,8.256)rectangle++(0.04,0.04) (0.764,8.593)rectangle++(0.04,0.04) (0.991,8.87)rectangle++(0.04,0.04)
  (1.195,9.096)rectangle++(0.04,0.04) (1.377,9.28)rectangle++(0.04,0.04) (1.54,9.429)rectangle++(0.04,0.04) (1.684,9.549)rectangle++(0.04,0.04) (1.811,9.646)rectangle++(0.04,0.04)
  (1.923,9.723)rectangle++(0.04,0.04) (2.021,9.784)rectangle++(0.04,0.04) (2.107,9.831)rectangle++(0.04,0.04) (2.182,9.869)rectangle++(0.04,0.04) (2.247,9.897)rectangle++(0.04,0.04)
  (2.304,9.919)rectangle++(0.04,0.04) (2.353,9.935)rectangle++(0.04,0.04) (2.395,9.947)rectangle++(0.04,0.04) (-2.128,4.328)rectangle++(0.04,0.04) (1.514,0.915)rectangle++(0.04,0.04)
  (1.321,2.314)rectangle++(0.04,0.04) (1.213,3.51)rectangle++(0.04,0.04) (1.169,4.531)rectangle++(0.04,0.04) (1.173,5.401)rectangle++(0.04,0.04) (1.211,6.14)rectangle++(0.04,0.04)
  (1.273,6.767)rectangle++(0.04,0.04) (1.35,7.297)rectangle++(0.04,0.04) (1.437,7.745)rectangle++(0.04,0.04) (1.529,8.122)rectangle++(0.04,0.04) (1.623,8.438)rectangle++(0.04,0.04)
  (1.715,8.704)rectangle++(0.04,0.04) (1.803,8.926)rectangle++(0.04,0.04) (1.888,9.111)rectangle++(0.04,0.04) (2.251,3.107)rectangle++(0.04,0.04) (2.035,4.148)rectangle++(0.04,0.04)
  (1.894,5.04)rectangle++(0.04,0.04) (1.11,2.132)rectangle++(0.04,0.04) (1.026,3.364)rectangle++(0.04,0.04) (1.005,4.415)rectangle++(0.04,0.04) (1.028,5.308)rectangle++(0.04,0.04)
  (1.084,6.067)rectangle++(0.04,0.04) (1.162,6.71)rectangle++(0.04,0.04) (1.254,7.253)rectangle++(0.04,0.04) (1.354,7.711)rectangle++(0.04,0.04) (-0.02,1.217)rectangle++(0.04,0.04)
  (0.029,2.631)rectangle++(0.04,0.04) (0.128,3.832)rectangle++(0.04,0.04) (0.26,4.848)rectangle++(0.04,0.04) (0.413,5.707)rectangle++(0.04,0.04) (-1.422,2.939)rectangle++(0.04,0.04)
  (-1.094,4.152)rectangle++(0.04,0.04) (-0.766,5.169)rectangle++(0.04,0.04) (-0.446,6.02)rectangle++(0.04,0.04) (-0.141,6.731)rectangle++(0.04,0.04) (0.147,7.323)rectangle++(0.04,0.04)
  (0.416,7.815)rectangle++(0.04,0.04) (0.664,8.223)rectangle++(0.04,0.04) (0.891,8.559)rectangle++(0.04,0.04) (1.098,8.836)rectangle++(0.04,0.04) (1.284,9.062)rectangle++(0.04,0.04)
  (1.452,9.248)rectangle++(0.04,0.04) (1.602,9.399)rectangle++(0.04,0.04) (1.735,9.521)rectangle++(0.04,0.04) (1.854,9.62)rectangle++(0.04,0.04) (1.958,9.699)rectangle++(0.04,0.04)
  (2.05,9.762)rectangle++(0.04,0.04) (2.131,9.812)rectangle++(0.04,0.04) (2.202,9.851)rectangle++(0.04,0.04) (2.263,9.881)rectangle++(0.04,0.04) (2.317,9.905)rectangle++(0.04,0.04)
  (2.363,9.923)rectangle++(0.04,0.04) (2.403,9.936)rectangle++(0.04,0.04) (2.438,9.946)rectangle++(0.04,0.04) (2.468,9.952)rectangle++(0.04,0.04) (2.494,9.957)rectangle++(0.04,0.04)
  (2.516,9.96)rectangle++(0.04,0.04) (2.535,9.962)rectangle++(0.04,0.04) (2.551,9.962)rectangle++(0.04,0.04) (2.564,9.962)rectangle++(0.04,0.04) (2.387,3.488)rectangle++(0.04,0.04)
  (0.601,1.888)rectangle++(0.04,0.04) (0.584,3.177)rectangle++(0.04,0.04) (0.621,4.273)rectangle++(0.04,0.04) (0.697,5.203)rectangle++(0.04,0.04) (0.798,5.991)rectangle++(0.04,0.04)
  (0.916,6.657)rectangle++(0.04,0.04) (1.043,7.218)rectangle++(0.04,0.04) (1.173,7.69)rectangle++(0.04,0.04) (1.302,8.086)rectangle++(0.04,0.04) (1.428,8.417)rectangle++(0.04,0.04)
  (1.548,8.693)rectangle++(0.04,0.04) (1.662,8.924)rectangle++(0.04,0.04) (1.767,9.115)rectangle++(0.04,0.04) (1.864,9.273)rectangle++(0.04,0.04) (2.299,3.14)rectangle++(0.04,0.04)
  (2.078,4.173)rectangle++(0.04,0.04) (1.931,5.061)rectangle++(0.04,0.04) (1.842,5.82)rectangle++(0.04,0.04) (1.796,6.47)rectangle++(0.04,0.04) (1.783,7.024)rectangle++(0.04,0.04)
  (1.794,7.495)rectangle++(0.04,0.04) (1.823,7.895)rectangle++(0.04,0.04) (1.863,8.234)rectangle++(0.04,0.04) (1.911,8.521)rectangle++(0.04,0.04) (1.963,8.762)rectangle++(0.04,0.04)
  (2.017,8.966)rectangle++(0.04,0.04) (-0.02,1.418)rectangle++(0.04,0.04) (0.038,2.802)rectangle++(0.04,0.04) (0.142,3.976)rectangle++(0.04,0.04) (0.277,4.971)rectangle++(0.04,0.04)
  (0.432,5.81)rectangle++(0.04,0.04) (0.598,6.517)rectangle++(0.04,0.04) (0.767,7.112)rectangle++(0.04,0.04) (1.859,2.336)rectangle++(0.04,0.04) (1.671,3.508)rectangle++(0.04,0.04)
  (1.559,4.511)rectangle++(0.04,0.04) (1.503,5.368)rectangle++(0.04,0.04) (1.49,6.099)rectangle++(0.04,0.04) (-0.02,0.959)rectangle++(0.04,0.04) (0.019,2.412)rectangle++(0.04,0.04)
  (0.111,3.646)rectangle++(0.04,0.04) (0.238,4.691)rectangle++(0.04,0.04) (-1.193,2.676)rectangle++(0.04,0.04) (-0.909,3.918)rectangle++(0.04,0.04) (-0.618,4.963)rectangle++(0.04,0.04)
  (-0.329,5.84)rectangle++(0.04,0.04) (-0.049,6.573)rectangle++(0.04,0.04) (0.219,7.185)rectangle++(0.04,0.04) (-1.845,3.22)rectangle++(0.04,0.04) (-1.442,4.407)rectangle++(0.04,0.04)
  (-1.052,5.4)rectangle++(0.04,0.04) (-0.68,6.228)rectangle++(0.04,0.04) (-0.331,6.917)rectangle++(0.04,0.04) (-0.007,7.489)rectangle++(0.04,0.04) (2.081,2.226)rectangle++(0.04,0.04)
  (1.855,3.405)rectangle++(0.04,0.04) (1.711,4.416)rectangle++(0.04,0.04) (1.629,5.281)rectangle++(0.04,0.04) (1.594,6.02)rectangle++(0.04,0.04) (1.429,2.289)rectangle++(0.04,0.04)
  (1.304,3.485)rectangle++(0.04,0.04) (1.246,4.506)rectangle++(0.04,0.04) (-0.02,0.704)rectangle++(0.04,0.04) (0.009,2.196)rectangle++(0.04,0.04) (0.093,3.462)rectangle++(0.04,0.04)
  (0.216,4.535)rectangle++(0.04,0.04) (1.22,1.574)rectangle++(0.04,0.04) (1.098,2.886)rectangle++(0.04,0.04) (-0.02,0.445)rectangle++(0.04,0.04) (-0.141,1.682)rectangle++(0.04,0.04)
  (-0.055,3.032)rectangle++(0.04,0.04) (0.073,4.175)rectangle++(0.04,0.04) (0.227,5.142)rectangle++(0.04,0.04) (0.396,5.958)rectangle++(0.04,0.04) (0.573,6.645)rectangle++(0.04,0.04)
  (0.75,7.221)rectangle++(0.04,0.04) (0.925,7.704)rectangle++(0.04,0.04) (-1.839,3.497)rectangle++(0.04,0.04) (-1.426,4.642)rectangle++(0.04,0.04) (-1.513,2.282)rectangle++(0.04,0.04)
  (-1.197,3.597)rectangle++(0.04,0.04) (-0.876,4.701)rectangle++(0.04,0.04) (-0.559,5.627)rectangle++(0.04,0.04) (-0.252,6.402)rectangle++(0.04,0.04) (0.04,7.048)rectangle++(0.04,0.04)
  (0.313,7.585)rectangle++(0.04,0.04) (0.568,8.031)rectangle++(0.04,0.04) (0.802,8.4)rectangle++(0.04,0.04) (1.015,8.704)rectangle++(0.04,0.04) (1.209,8.954)rectangle++(0.04,0.04)
  (1.383,9.159)rectangle++(0.04,0.04) (1.54,9.326)rectangle++(0.04,0.04) (-2.138,3.995)rectangle++(0.04,0.04) (-1.66,5.077)rectangle++(0.04,0.04) (-1.21,5.978)rectangle++(0.04,0.04)
  (1.838,1.55)rectangle++(0.04,0.04) (1.622,2.84)rectangle++(0.04,0.04) (1.49,3.946)rectangle++(0.04,0.04) (1.422,4.89)rectangle++(0.04,0.04) (1.402,5.696)rectangle++(0.04,0.04)
  (1.418,6.382)rectangle++(0.04,0.04) (1.557,2.33)rectangle++(0.04,0.04) (1.414,3.515)rectangle++(0.04,0.04) (1.341,4.527)rectangle++(0.04,0.04) (1.049,1.865)rectangle++(0.04,0.04)
  (0.964,3.14)rectangle++(0.04,0.04) (0.943,4.226)rectangle++(0.04,0.04) (0.968,5.151)rectangle++(0.04,0.04) (1.027,5.936)rectangle++(0.04,0.04) (1.108,6.6)rectangle++(0.04,0.04)
  (1.204,7.162)rectangle++(0.04,0.04) (1.307,7.636)rectangle++(0.04,0.04) (1.415,8.034)rectangle++(0.04,0.04) (2.02,2.726)rectangle++(0.04,0.04) (1.824,3.833)rectangle++(0.04,0.04)
  (1.701,4.781)rectangle++(0.04,0.04) (1.635,5.592)rectangle++(0.04,0.04) (1.611,6.284)rectangle++(0.04,0.04) (-1.333,3.342)rectangle++(0.04,0.04) (-1.001,4.49)rectangle++(0.04,0.04)
  (-0.674,5.453)rectangle++(0.04,0.04) (1.61,1.564)rectangle++(0.04,0.04) (0.179,1.224)rectangle++(0.04,0.04) (0.298,0.77)rectangle++(0.04,0.04) (0.282,2.239)rectangle++(0.04,0.04)
  (0.327,3.488)rectangle++(0.04,0.04) (-0.863,2.432)rectangle++(0.04,0.04) (-0.638,3.698)rectangle++(0.04,0.04) (-0.397,4.765)rectangle++(0.04,0.04) (-0.149,5.662)rectangle++(0.04,0.04)
  (-1.523,2.8)rectangle++(0.04,0.04) (-1.185,4.037)rectangle++(0.04,0.04) (-0.848,5.075)rectangle++(0.04,0.04) (-0.52,5.944)rectangle++(0.04,0.04) (-0.206,6.67)rectangle++(0.04,0.04)
  (0.089,7.274)rectangle++(0.04,0.04) (0.365,7.775)rectangle++(0.04,0.04) (0.619,8.191)rectangle++(0.04,0.04) (-2.027,3.533)rectangle++(0.04,0.04) (-1.584,4.681)rectangle++(0.04,0.04)
  (1.531,1.142)rectangle++(0.04,0.04) (1.345,2.505)rectangle++(0.04,0.04) (1.241,3.672)rectangle++(0.04,0.04) (1.199,4.668)rectangle++(0.04,0.04) (1.204,5.516)rectangle++(0.04,0.04)
  (1.242,6.236)rectangle++(0.04,0.04) (1.303,6.847)rectangle++(0.04,0.04) (1.379,7.364)rectangle++(0.04,0.04) (1.465,7.801)rectangle++(0.04,0.04) (-1.756,3.642)rectangle++(0.04,0.04)
  (-1.35,4.762)rectangle++(0.04,0.04) (-0.959,5.698)rectangle++(0.04,0.04) (-1.694,2.622)rectangle++(0.04,0.04) (-1.338,3.893)rectangle++(0.04,0.04) (-0.983,4.959)rectangle++(0.04,0.04)
  (-0.64,5.85)rectangle++(0.04,0.04) (1.717,1.668)rectangle++(0.04,0.04) (1.524,2.945)rectangle++(0.04,0.04) (1.411,4.039)rectangle++(0.04,0.04) (1.358,4.973)rectangle++(0.04,0.04)
  (1.351,5.769)rectangle++(0.04,0.04) (1.377,6.445)rectangle++(0.04,0.04) (1.426,7.02)rectangle++(0.04,0.04) (1.734,2.486)rectangle++(0.04,0.04) (1.571,3.64)rectangle++(0.04,0.04)
  (0.766,1.712)rectangle++(0.04,0.04) (0.717,3.021)rectangle++(0.04,0.04) (0.728,4.135)rectangle++(0.04,0.04) (0.782,5.082)rectangle++(0.04,0.04) (0.866,5.885)rectangle++(0.04,0.04)
  (0.969,6.563)rectangle++(0.04,0.04) (1.084,7.136)rectangle++(0.04,0.04) (-1.66,3.408)rectangle++(0.04,0.04) (-1.277,4.56)rectangle++(0.04,0.04) (-0.905,5.523)rectangle++(0.04,0.04)
  (-0.551,6.327)rectangle++(0.04,0.04) (-0.217,6.996)rectangle++(0.04,0.04) (0.093,7.552)rectangle++(0.04,0.04) (0.379,8.011)rectangle++(0.04,0.04) (0.64,8.391)rectangle++(0.04,0.04)
  (0.878,8.703)rectangle++(0.04,0.04) (1.092,8.958)rectangle++(0.04,0.04) (1.284,9.167)rectangle++(0.04,0.04) (1.456,9.337)rectangle++(0.04,0.04) (1.609,9.474)rectangle++(0.04,0.04)
  (1.744,9.585)rectangle++(0.04,0.04) (1.864,9.674)rectangle++(0.04,0.04) (1.969,9.744)rectangle++(0.04,0.04) (2.061,9.8)rectangle++(0.04,0.04) (2.142,9.844)rectangle++(0.04,0.04)
  (2.212,9.878)rectangle++(0.04,0.04) (2.273,9.904)rectangle++(0.04,0.04) (2.326,9.923)rectangle++(0.04,0.04) (2.372,9.938)rectangle++(0.04,0.04) (2.409,3.432)rectangle++(0.04,0.04)
  (2.183,4.417)rectangle++(0.04,0.04) (2.03,5.263)rectangle++(0.04,0.04) (1.934,5.989)rectangle++(0.04,0.04) (1.881,6.609)rectangle++(0.04,0.04) (1.861,7.139)rectangle++(0.04,0.04)
  (1.865,7.59)rectangle++(0.04,0.04) (1.887,7.973)rectangle++(0.04,0.04) (1.921,8.298)rectangle++(0.04,0.04) (1.962,8.572)rectangle++(0.04,0.04) (2.009,8.804)rectangle++(0.04,0.04)
  (2.057,8.999)rectangle++(0.04,0.04) (2.106,9.163)rectangle++(0.04,0.04) (-1.982,4.089)rectangle++(0.04,0.04) (1.425,0.896)rectangle++(0.04,0.04) (1.245,2.301)rectangle++(0.04,0.04)
  (1.148,3.502)rectangle++(0.04,0.04) (1.114,4.527)rectangle++(0.04,0.04) (1.083,1.806)rectangle++(0.04,0.04) (0.991,3.088)rectangle++(0.04,0.04) (0.963,4.181)rectangle++(0.04,0.04)
  (0.984,5.112)rectangle++(0.04,0.04) (1.039,5.902)rectangle++(0.04,0.04) (1.117,6.571)rectangle++(0.04,0.04) (1.21,7.137)rectangle++(0.04,0.04) (1.312,7.614)rectangle++(0.04,0.04)
  (1.417,8.016)rectangle++(0.04,0.04) (1.523,8.353)rectangle++(0.04,0.04) (1.627,8.635)rectangle++(0.04,0.04) (1.726,8.871)rectangle++(0.04,0.04) (1.82,9.068)rectangle++(0.04,0.04)
  (1.907,9.231)rectangle++(0.04,0.04) (1.988,9.366)rectangle++(0.04,0.04) (2.062,9.478)rectangle++(0.04,0.04) (2.13,9.57)rectangle++(0.04,0.04) (2.191,9.645)rectangle++(0.04,0.04)
  (2.246,9.707)rectangle++(0.04,0.04) (2.295,9.757)rectangle++(0.04,0.04) (2.339,9.798)rectangle++(0.04,0.04) (2.375,3.39)rectangle++(0.04,0.04) (2.152,4.382)rectangle++(0.04,0.04)
  (2.003,5.235)rectangle++(0.04,0.04) (1.909,5.966)rectangle++(0.04,0.04) (1.859,6.591)rectangle++(0.04,0.04) (1.842,7.124)rectangle++(0.04,0.04) (1.848,7.578)rectangle++(0.04,0.04)
  (1.872,7.964)rectangle++(0.04,0.04) (1.908,8.29)rectangle++(0.04,0.04) (2.018,2.916)rectangle++(0.04,0.04) (1.83,3.994)rectangle++(0.04,0.04) (1.713,4.918)rectangle++(0.04,0.04)
  (1.65,5.708)rectangle++(0.04,0.04) (1.629,6.382)rectangle++(0.04,0.04) (1.638,6.956)rectangle++(0.04,0.04) (1.668,7.443)rectangle++(0.04,0.04) (1.713,7.856)rectangle++(0.04,0.04)
  (1.768,8.205)rectangle++(0.04,0.04) (1.829,8.5)rectangle++(0.04,0.04) (1.893,8.748)rectangle++(0.04,0.04) (1.956,8.956)rectangle++(0.04,0.04) (2.019,9.131)rectangle++(0.04,0.04)
  (2.079,9.277)rectangle++(0.04,0.04) (2.136,9.398)rectangle++(0.04,0.04) (2.189,9.499)rectangle++(0.04,0.04) (2.239,9.583)rectangle++(0.04,0.04) (2.284,9.652)rectangle++(0.04,0.04)
  (2.325,9.709)rectangle++(0.04,0.04) (-2.081,4.26)rectangle++(0.04,0.04) (-1.6,5.3)rectangle++(0.04,0.04) (-1.15,6.165)rectangle++(0.04,0.04) (-0.733,6.883)rectangle++(0.04,0.04)
  (-0.35,7.476)rectangle++(0.04,0.04) (-0.001,7.965)rectangle++(0.04,0.04) (0.316,8.366)rectangle++(0.04,0.04) (0.601,8.695)rectangle++(0.04,0.04) (0.856,8.963)rectangle++(0.04,0.04)
  (1.084,9.18)rectangle++(0.04,0.04) (1.287,9.356)rectangle++(0.04,0.04) (1.466,9.497)rectangle++(0.04,0.04) (1.623,9.61)rectangle++(0.04,0.04) (1.762,9.7)rectangle++(0.04,0.04)
  (1.884,9.771)rectangle++(0.04,0.04) (1.99,9.826)rectangle++(0.04,0.04) (-2.178,4.208)rectangle++(0.04,0.04) (-1.685,5.26)rectangle++(0.04,0.04) (-1.224,6.135)rectangle++(0.04,0.04)
  (-0.797,6.86)rectangle++(0.04,0.04) (-0.406,7.459)rectangle++(0.04,0.04) (-0.049,7.953)rectangle++(0.04,0.04) (0.275,8.358)rectangle++(0.04,0.04) (0.566,8.689)rectangle++(0.04,0.04)
  (0.826,8.96)rectangle++(0.04,0.04) (1.058,9.179)rectangle++(0.04,0.04) (1.265,9.356)rectangle++(0.04,0.04) (1.447,9.498)rectangle++(0.04,0.04) (1.608,9.612)rectangle++(0.04,0.04)
  (1.749,9.702)rectangle++(0.04,0.04) (-2.194,4.126)rectangle++(0.04,0.04) (-1.702,5.191)rectangle++(0.04,0.04) (1.691,1.233)rectangle++(0.04,0.04) (1.485,2.577)rectangle++(0.04,0.04)
  (1.363,3.727)rectangle++(0.04,0.04) (1.305,4.71)rectangle++(0.04,0.04) (1.296,5.547)rectangle++(0.04,0.04) (1.321,6.26)rectangle++(0.04,0.04) (1.371,6.864)rectangle++(0.04,0.04)
  (1.438,7.376)rectangle++(0.04,0.04) (1.832,2.574)rectangle++(0.04,0.04) (1.658,3.711)rectangle++(0.04,0.04) (1.556,4.684)rectangle++(0.04,0.04) (1.507,5.515)rectangle++(0.04,0.04)
  (1.5,6.224)rectangle++(0.04,0.04) (1.522,6.827)rectangle++(0.04,0.04) (1.564,7.338)rectangle++(0.04,0.04) (1.621,7.771)rectangle++(0.04,0.04) (1.686,8.137)rectangle++(0.04,0.04)
  (1.757,8.445)rectangle++(0.04,0.04) (1.829,8.704)rectangle++(0.04,0.04) (1.9,8.922)rectangle++(0.04,0.04) (1.97,9.103)rectangle++(0.04,0.04) (2.036,9.255)rectangle++(0.04,0.04)
  (2.099,9.382)rectangle++(0.04,0.04) (2.295,3.227)rectangle++(0.04,0.04) (2.077,4.248)rectangle++(0.04,0.04) (1.933,5.124)rectangle++(0.04,0.04) (1.846,5.874)rectangle++(0.04,0.04)
  (1.802,6.515)rectangle++(0.04,0.04) (1.537,2.462)rectangle++(0.04,0.04) (1.402,3.628)rectangle++(0.04,0.04) (1.335,4.624)rectangle++(0.04,0.04) (1.317,5.473)rectangle++(0.04,0.04)
  (1.317,2.086)rectangle++(0.04,0.04) (1.201,3.317)rectangle++(0.04,0.04) (1.151,4.367)rectangle++(0.04,0.04) (1.033,1.777)rectangle++(0.04,0.04) (0.947,3.066)rectangle++(0.04,0.04)
  (0.925,4.164)rectangle++(0.04,0.04) (0.951,5.099)rectangle++(0.04,0.04) (1.01,5.892)rectangle++(0.04,0.04) (1.092,6.564)rectangle++(0.04,0.04) (1.188,7.132)rectangle++(0.04,0.04)
  (1.293,7.611)rectangle++(0.04,0.04) (1.402,8.014)rectangle++(0.04,0.04) (1.51,8.352)rectangle++(0.04,0.04) (1.615,8.635)rectangle++(0.04,0.04) (1.716,8.871)rectangle++(0.04,0.04)
  (1.811,9.068)rectangle++(0.04,0.04) (1.9,9.232)rectangle++(0.04,0.04) (-2.041,4.057)rectangle++(0.04,0.04) (-1.575,5.126)rectangle++(0.04,0.04) (-1.136,6.017)rectangle++(0.04,0.04)
  (-0.727,6.756)rectangle++(0.04,0.04) (-0.35,7.368)rectangle++(0.04,0.04) (-0.005,7.873)rectangle++(0.04,0.04) (0.308,8.288)rectangle++(0.04,0.04) (0.591,8.629)rectangle++(0.04,0.04)
  (0.846,8.907)rectangle++(0.04,0.04) (1.073,9.133)rectangle++(0.04,0.04) (1.275,9.317)rectangle++(0.04,0.04) (1.454,9.464)rectangle++(0.04,0.04) (1.613,9.583)rectangle++(0.04,0.04)
  (1.752,9.677)rectangle++(0.04,0.04) (-2.187,4.121)rectangle++(0.04,0.04) (-1.696,5.186)rectangle++(0.04,0.04) (-1.237,6.072)rectangle++(0.04,0.04) (1.868,1.566)rectangle++(0.04,0.04)
  (-0.055,2.363)rectangle++(0.04,0.04) (0.046,3.607)rectangle++(0.04,0.04) (0.181,4.66)rectangle++(0.04,0.04) (1.26,1.596)rectangle++(0.04,0.04) (1.133,2.902)rectangle++(0.04,0.04)
  (1.077,4.018)rectangle++(0.04,0.04) (-0.85,2.721)rectangle++(0.04,0.04) (-0.616,3.943)rectangle++(0.04,0.04) (-0.368,4.972)rectangle++(0.04,0.04) (-0.116,5.837)rectangle++(0.04,0.04)
  (0.132,6.563)rectangle++(0.04,0.04) (0.373,7.169)rectangle++(0.04,0.04) (0.601,7.675)rectangle++(0.04,0.04) (0.816,8.096)rectangle++(0.04,0.04) (1.015,8.445)rectangle++(0.04,0.04)
  (-2.014,3.68)rectangle++(0.04,0.04) (-1.567,4.805)rectangle++(0.04,0.04) (-1.142,5.743)rectangle++(0.04,0.04) (-0.743,6.524)rectangle++(0.04,0.04) (-0.373,7.171)rectangle++(0.04,0.04)
  (-0.032,7.706)rectangle++(0.04,0.04) (0.279,8.148)rectangle++(0.04,0.04) (0.561,8.511)rectangle++(0.04,0.04) (0.815,8.808)rectangle++(0.04,0.04) (1.043,9.05)rectangle++(0.04,0.04)
  (1.246,9.247)rectangle++(0.04,0.04) (2.385,2.973)rectangle++(0.04,0.04) (2.144,4.028)rectangle++(0.04,0.04) (1.981,4.934)rectangle++(0.04,0.04) (1.879,5.711)rectangle++(0.04,0.04)
  (1.824,6.376)rectangle++(0.04,0.04) (1.803,6.942)rectangle++(0.04,0.04) (-1.466,3.531)rectangle++(0.04,0.04) (-1.107,4.656)rectangle++(0.04,0.04) (-1.453,2.359)rectangle++(0.04,0.04)
  (-1.143,3.659)rectangle++(0.04,0.04) (-0.827,4.752)rectangle++(0.04,0.04) (-0.515,5.669)rectangle++(0.04,0.04) (-0.214,6.435)rectangle++(0.04,0.04) (1.817,1.919)rectangle++(0.04,0.04)
  (1.619,3.155)rectangle++(0.04,0.04) (1.5,4.213)rectangle++(0.04,0.04) (1.441,5.117)rectangle++(0.04,0.04) (1.199,2.033)rectangle++(0.04,0.04) (1.098,3.276)rectangle++(0.04,0.04)
  (1.063,4.337)rectangle++(0.04,0.04) (1.074,5.24)rectangle++(0.04,0.04) (1.121,6.007)rectangle++(0.04,0.04) (1.191,6.658)rectangle++(0.04,0.04) (1.276,7.208)rectangle++(0.04,0.04)
  (1.371,7.672)rectangle++(0.04,0.04) (1.47,8.062)rectangle++(0.04,0.04) (1.57,8.39)rectangle++(0.04,0.04) (1.668,8.665)rectangle++(0.04,0.04) (1.762,8.895)rectangle++(0.04,0.04)
  (2.209,3.023)rectangle++(0.04,0.04) (1.996,4.077)rectangle++(0.04,0.04) (1.858,4.982)rectangle++(0.04,0.04) (-0.945,3.112)rectangle++(0.04,0.04) (-0.681,4.279)rectangle++(0.04,0.04)
  (-0.41,5.261)rectangle++(0.04,0.04) (-0.14,6.084)rectangle++(0.04,0.04) (0.122,6.774)rectangle++(0.04,0.04) (0.373,7.349)rectangle++(0.04,0.04) (0.608,7.828)rectangle++(0.04,0.04)
  (0.828,8.226)rectangle++(0.04,0.04) (1.031,8.555)rectangle++(0.04,0.04) (1.216,8.827)rectangle++(0.04,0.04) (1.385,9.05)rectangle++(0.04,0.04) (1.537,9.233)rectangle++(0.04,0.04)
  (1.673,9.383)rectangle++(0.04,0.04) (1.795,9.505)rectangle++(0.04,0.04) (2.375,3.178)rectangle++(0.04,0.04) (2.143,4.203)rectangle++(0.04,0.04) (1.988,5.083)rectangle++(0.04,0.04)
  (1.891,5.837)rectangle++(0.04,0.04) (-1.161,3.308)rectangle++(0.04,0.04) (-0.856,4.454)rectangle++(0.04,0.04) (-0.552,5.417)rectangle++(0.04,0.04) (-0.255,6.222)rectangle++(0.04,0.04)
  (0.03,6.895)rectangle++(0.04,0.04) (1.909,2.093)rectangle++(0.04,0.04) (1.704,3.299)rectangle++(0.04,0.04) (1.578,4.332)rectangle++(0.04,0.04) (1.513,5.215)rectangle++(0.04,0.04)
  (1.492,5.969)rectangle++(0.04,0.04) (1.505,6.61)rectangle++(0.04,0.04) (1.541,7.154)rectangle++(0.04,0.04) (1.594,7.616)rectangle++(0.04,0.04) (1.657,8.006)rectangle++(0.04,0.04)
  (1.976,2.782)rectangle++(0.04,0.04) (1.788,3.882)rectangle++(0.04,0.04) (0.801,1.827)rectangle++(0.04,0.04) (0.752,3.117)rectangle++(0.04,0.04) (0.743,1.374)rectangle++(0.04,0.04)
  (0.684,2.734)rectangle++(0.04,0.04) (0.688,3.893)rectangle++(0.04,0.04) (-0.896,2.604)rectangle++(0.04,0.04) (-0.659,3.845)rectangle++(0.04,0.04) (-0.409,4.891)rectangle++(0.04,0.04)
  (-1.375,2.571)rectangle++(0.04,0.04) (-1.068,3.837)rectangle++(0.04,0.04) (-0.756,4.9)rectangle++(0.04,0.04) (-1.446,2.493)rectangle++(0.04,0.04) (-1.132,3.773)rectangle++(0.04,0.04)
  (-0.813,4.849)rectangle++(0.04,0.04) (-1.445,2.469)rectangle++(0.04,0.04) (-1.131,3.752)rectangle++(0.04,0.04) (-0.814,4.831)rectangle++(0.04,0.04) (-0.501,5.735)rectangle++(0.04,0.04)
  (-0.198,6.491)rectangle++(0.04,0.04) (0.089,7.121)rectangle++(0.04,0.04) (0.358,7.646)rectangle++(0.04,0.04) (0.608,8.081)rectangle++(0.04,0.04) (0.838,8.441)rectangle++(0.04,0.04)
  (1.048,8.737)rectangle++(0.04,0.04) (2.272,2.799)rectangle++(0.04,0.04) (2.041,3.885)rectangle++(0.04,0.04) (1.888,4.817)rectangle++(0.04,0.04) (1.795,5.615)rectangle++(0.04,0.04)
  (-1.122,3.237)rectangle++(0.04,0.04) (-0.826,4.393)rectangle++(0.04,0.04) (-0.529,5.363)rectangle++(0.04,0.04) (-0.237,6.176)rectangle++(0.04,0.04) (0.043,6.855)rectangle++(0.04,0.04)
  (1.896,2.086)rectangle++(0.04,0.04) (1.693,3.294)rectangle++(0.04,0.04) (1.568,4.328)rectangle++(0.04,0.04) (1.504,5.213)rectangle++(0.04,0.04) (-1.076,3.082)rectangle++(0.04,0.04)
  (-0.793,4.259)rectangle++(0.04,0.04) (-0.506,5.248)rectangle++(0.04,0.04) (-0.222,6.077)rectangle++(0.04,0.04) (0.052,6.771)rectangle++(0.04,0.04) (0.313,7.349)rectangle++(0.04,0.04)
  (0.557,7.83)rectangle++(0.04,0.04) (0.785,8.23)rectangle++(0.04,0.04) (0.994,8.56)rectangle++(0.04,0.04) (2.23,2.743)rectangle++(0.04,0.04) (-0.288,2.705)rectangle++(0.04,0.04)
  (-0.139,3.907)rectangle++(0.04,0.04) (1.098,1.332)rectangle++(0.04,0.04) (0.984,2.684)rectangle++(0.04,0.04) (0.942,3.838)rectangle++(0.04,0.04) (-0.02,0.597)rectangle++(0.04,0.04)
  (0.005,2.105)rectangle++(0.04,0.04) (0.086,3.385)rectangle++(0.04,0.04) (0.206,4.47)rectangle++(0.04,0.04) (0.352,5.387)rectangle++(0.04,0.04) (0.512,6.161)rectangle++(0.04,0.04)
  (0.68,6.813)rectangle++(0.04,0.04) (0.848,7.36)rectangle++(0.04,0.04) (1.013,7.818)rectangle++(0.04,0.04) (1.172,8.201)rectangle++(0.04,0.04) (2.103,2.703)rectangle++(0.04,0.04)
  (-0.303,2.667)rectangle++(0.04,0.04) (0.775,0.991)rectangle++(0.04,0.04) (0.696,2.408)rectangle++(0.04,0.04) (0.552,1.189)rectangle++(0.04,0.04) (0.515,2.585)rectangle++(0.04,0.04)
  (0.539,3.773)rectangle++(0.04,0.04) (0.958,1.476)rectangle++(0.04,0.04) (0.252,1.033)rectangle++(0.04,0.04) (0.253,2.464)rectangle++(0.04,0.04) (0.635,1.087)rectangle++(0.04,0.04)
  (0.581,2.495)rectangle++(0.04,0.04) (0.591,3.694)rectangle++(0.04,0.04) (0.648,4.712)rectangle++(0.04,0.04) (0.737,5.576)rectangle++(0.04,0.04) (0.847,6.306)rectangle++(0.04,0.04)
  (1.621,2.164)rectangle++(0.04,0.04) (1.462,3.371)rectangle++(0.04,0.04) (1.376,4.403)rectangle++(0.04,0.04) (1.343,5.283)rectangle++(0.04,0.04) (1.351,6.033)rectangle++(0.04,0.04)
  (1.469,2.229)rectangle++(0.04,0.04) (1.336,3.432)rectangle++(0.04,0.04) (1.271,4.46)rectangle++(0.04,0.04) (1.256,5.337)rectangle++(0.04,0.04) (1.279,6.082)rectangle++(0.04,0.04)
  (1.328,6.715)rectangle++(0.04,0.04) (1.395,7.251)rectangle++(0.04,0.04) (1.474,7.703)rectangle++(0.04,0.04) (1.559,8.085)rectangle++(0.04,0.04) (1.646,8.406)rectangle++(0.04,0.04)
  (1.733,8.676)rectangle++(0.04,0.04) (1.818,8.901)rectangle++(0.04,0.04) (1.899,9.089)rectangle++(0.04,0.04) (1.976,9.246)rectangle++(0.04,0.04) (2.047,9.376)rectangle++(0.04,0.04)
  (2.113,9.484)rectangle++(0.04,0.04) (2.173,9.573)rectangle++(0.04,0.04) (-0.02,1.515)rectangle++(0.04,0.04) (0.041,2.885)rectangle++(0.04,0.04) (0.148,4.047)rectangle++(0.04,0.04)
  (0.286,5.03)rectangle++(0.04,0.04) (0.442,5.86)rectangle++(0.04,0.04) (0.608,6.56)rectangle++(0.04,0.04) (0.777,7.148)rectangle++(0.04,0.04) (0.944,7.641)rectangle++(0.04,0.04)
  (1.106,8.053)rectangle++(0.04,0.04) (1.26,8.397)rectangle++(0.04,0.04) (-1.952,3.726)rectangle++(0.04,0.04) (-1.513,4.841)rectangle++(0.04,0.04) (-1.094,5.772)rectangle++(0.04,0.04)
  (-0.702,6.546)rectangle++(0.04,0.04) (-0.337,7.188)rectangle++(0.04,0.04) (-0.001,7.72)rectangle++(0.04,0.04) (0.306,8.158)rectangle++(0.04,0.04) (0.584,8.518)rectangle++(0.04,0.04)
  (0.835,8.813)rectangle++(0.04,0.04) (1.06,9.054)rectangle++(0.04,0.04) (1.261,9.25)rectangle++(0.04,0.04) (1.44,9.408)rectangle++(0.04,0.04) (1.598,9.536)rectangle++(0.04,0.04)
  (1.737,9.638)rectangle++(0.04,0.04) (1.86,9.719)rectangle++(0.04,0.04) (1.968,9.783)rectangle++(0.04,0.04) (2.062,9.833)rectangle++(0.04,0.04) (-2.165,4.226)rectangle++(0.04,0.04)
  (-1.674,5.275)rectangle++(0.04,0.04) (-1.214,6.147)rectangle++(0.04,0.04) (-0.788,6.87)rectangle++(0.04,0.04) (-0.397,7.467)rectangle++(0.04,0.04) (-0.041,7.959)rectangle++(0.04,0.04)
  (0.281,8.363)rectangle++(0.04,0.04) (0.571,8.694)rectangle++(0.04,0.04) (-2.167,3.633)rectangle++(0.04,0.04) (-1.699,4.771)rectangle++(0.04,0.04) (-1.256,5.719)rectangle++(0.04,0.04)
  (-0.841,6.508)rectangle++(0.04,0.04) (-0.456,7.162)rectangle++(0.04,0.04) (-0.104,7.702)rectangle++(0.04,0.04) (0.218,8.147)rectangle++(0.04,0.04) (0.509,8.512)rectangle++(0.04,0.04)
  (0.771,8.811)rectangle++(0.04,0.04) (1.005,9.055)rectangle++(0.04,0.04) (1.215,9.253)rectangle++(0.04,0.04) (-2.184,3.904)rectangle++(0.04,0.04) (-1.702,5.002)rectangle++(0.04,0.04)
  (-1.249,5.916)rectangle++(0.04,0.04) (-0.827,6.675)rectangle++(0.04,0.04) (-0.438,7.303)rectangle++(0.04,0.04) (-0.083,7.821)rectangle++(0.04,0.04) (0.24,8.247)rectangle++(0.04,0.04)
  (0.532,8.597)rectangle++(0.04,0.04) (0.794,8.882)rectangle++(0.04,0.04) (1.028,9.114)rectangle++(0.04,0.04) (1.236,9.302)rectangle++(0.04,0.04) (1.421,9.454)rectangle++(0.04,0.04)
  (1.583,9.575)rectangle++(0.04,0.04) (2.426,3.14)rectangle++(0.04,0.04) (2.186,4.168)rectangle++(0.04,0.04) (2.022,5.051)rectangle++(0.04,0.04) (-0.93,3.165)rectangle++(0.04,0.04)
  (-0.666,4.324)rectangle++(0.04,0.04) (-0.396,5.298)rectangle++(0.04,0.04) (1.525,1.599)rectangle++(0.04,0.04) (1.358,2.894)rectangle++(0.04,0.04) (1.268,4.002)rectangle++(0.04,0.04)
  (1.236,4.947)rectangle++(0.04,0.04) (1.246,5.752)rectangle++(0.04,0.04) (1.287,6.435)rectangle++(0.04,0.04) (-1.437,3.301)rectangle++(0.04,0.04) (-1.092,4.459)rectangle++(0.04,0.04)
  (-0.752,5.43)rectangle++(0.04,0.04) (-0.424,6.242)rectangle++(0.04,0.04) (-0.113,6.919)rectangle++(0.04,0.04) (0.179,7.482)rectangle++(0.04,0.04) (0.449,7.949)rectangle++(0.04,0.04)
  (0.697,8.335)rectangle++(0.04,0.04) (0.924,8.653)rectangle++(0.04,0.04) (2.267,2.747)rectangle++(0.04,0.04) (2.034,3.84)rectangle++(0.04,0.04) (-0.613,2.902)rectangle++(0.04,0.04)
  (-0.407,4.087)rectangle++(0.04,0.04) (-1.165,2.395)rectangle++(0.04,0.04) (-0.897,3.678)rectangle++(0.04,0.04) (-0.617,4.759)rectangle++(0.04,0.04) (-0.337,5.666)rectangle++(0.04,0.04)
  (-0.062,6.425)rectangle++(0.04,0.04) (0.202,7.06)rectangle++(0.04,0.04) (0.452,7.589)rectangle++(0.04,0.04) (0.686,8.029)rectangle++(0.04,0.04) (-1.972,3.513)rectangle++(0.04,0.04)
  (-0.02,0.545)rectangle++(0.04,0.04) (0.003,2.061)rectangle++(0.04,0.04) (0.082,3.348)rectangle++(0.04,0.04) (0.202,4.438)rectangle++(0.04,0.04) (0.347,5.361)rectangle++(0.04,0.04)
  (0.507,6.139)rectangle++(0.04,0.04) (0.674,6.794)rectangle++(0.04,0.04) (0.843,7.344)rectangle++(0.04,0.04) (-1.762,3.399)rectangle++(0.04,0.04) (-1.364,4.555)rectangle++(0.04,0.04)
  (-0.979,5.523)rectangle++(0.04,0.04) (-0.614,6.33)rectangle++(0.04,0.04) (-1.79,2.84)rectangle++(0.04,0.04) (-1.41,4.082)rectangle++(0.04,0.04) (-1.037,5.122)rectangle++(0.04,0.04)
  (-0.679,5.992)rectangle++(0.04,0.04) (-0.34,6.716)rectangle++(0.04,0.04) (-0.022,7.319)rectangle++(0.04,0.04) (0.272,7.818)rectangle++(0.04,0.04) (2.131,2.377)rectangle++(0.04,0.04)
  (1.904,3.531)rectangle++(0.04,0.04) (1.758,4.522)rectangle++(0.04,0.04) (1.673,5.369)rectangle++(0.04,0.04) (1.634,6.093)rectangle++(0.04,0.04) (-1.279,3.305)rectangle++(0.04,0.04)
  (-0.957,4.457)rectangle++(0.04,0.04) (-0.637,5.423)rectangle++(0.04,0.04) (1.597,1.566)rectangle++(0.04,0.04) (1.418,2.863)rectangle++(0.04,0.04) (1.317,3.973)rectangle++(0.04,0.04)
  (0.898,1.726)rectangle++(0.04,0.04) (0.83,3.027)rectangle++(0.04,0.04) (0.824,4.136)rectangle++(0.04,0.04) (0.864,5.079)rectangle++(0.04,0.04) (0.935,5.879)rectangle++(0.04,0.04)
  (1.028,6.556)rectangle++(0.04,0.04) (1.134,7.128)rectangle++(0.04,0.04) (1.247,7.609)rectangle++(0.04,0.04) (1.362,8.014)rectangle++(0.04,0.04) (1.476,8.354)rectangle++(0.04,0.04)
  (1.586,8.638)rectangle++(0.04,0.04) (-1.95,3.854)rectangle++(0.04,0.04) (1.354,0.848)rectangle++(0.04,0.04) (1.183,2.263)rectangle++(0.04,0.04) (1.094,3.472)rectangle++(0.04,0.04)
  (1.066,4.504)rectangle++(0.04,0.04) (1.084,5.382)rectangle++(0.04,0.04) (1.135,6.127)rectangle++(0.04,0.04) (1.207,6.759)rectangle++(0.04,0.04) (1.694,2.366)rectangle++(0.04,0.04)
  (1.532,3.54)rectangle++(0.04,0.04) (1.442,4.544)rectangle++(0.04,0.04) (1.039,1.895)rectangle++(0.04,0.04) (0.956,3.166)rectangle++(0.04,0.04) (0.937,4.249)rectangle++(0.04,0.04)
  (0.964,5.17)rectangle++(0.04,0.04) (1.024,5.952)rectangle++(0.04,0.04) (-1.364,3.134)rectangle++(0.04,0.04) (-1.036,4.315)rectangle++(0.04,0.04) (-0.71,5.305)rectangle++(0.04,0.04)
  (-0.394,6.134)rectangle++(0.04,0.04) (-1.695,2.848)rectangle++(0.04,0.04) (-1.329,4.085)rectangle++(0.04,0.04) (-0.968,5.121)rectangle++(0.04,0.04) (-0.62,5.988)rectangle++(0.04,0.04)
  (-0.29,6.711)rectangle++(0.04,0.04) (0.02,7.312)rectangle++(0.04,0.04) (0.307,7.811)rectangle++(0.04,0.04) (0.571,8.223)rectangle++(0.04,0.04) (2.199,2.552)rectangle++(0.04,0.04)
  (1.969,3.677)rectangle++(0.04,0.04) (1.819,4.643)rectangle++(0.04,0.04) (1.73,5.47)rectangle++(0.04,0.04) (1.687,6.177)rectangle++(0.04,0.04) (1.679,6.779)rectangle++(0.04,0.04)
  (1.696,7.291)rectangle++(0.04,0.04) (1.731,7.726)rectangle++(0.04,0.04) (1.778,8.094)rectangle++(0.04,0.04) (1.833,8.405)rectangle++(0.04,0.04) (1.892,8.667)rectangle++(0.04,0.04)
  (1.953,8.888)rectangle++(0.04,0.04) (2.013,9.072)rectangle++(0.04,0.04) (2.072,9.227)rectangle++(0.04,0.04) (2.128,9.357)rectangle++(0.04,0.04) (2.181,9.464)rectangle++(0.04,0.04)
  (2.305,3.268)rectangle++(0.04,0.04) (2.088,4.282)rectangle++(0.04,0.04) (1.944,5.152)rectangle++(0.04,0.04) (1.856,5.898)rectangle++(0.04,0.04) (1.812,6.535)rectangle++(0.04,0.04)
  (1.799,7.079)rectangle++(0.04,0.04) (1.81,7.541)rectangle++(0.04,0.04) (1.838,7.934)rectangle++(0.04,0.04) (1.878,8.266)rectangle++(0.04,0.04) (1.924,8.548)rectangle++(0.04,0.04)
  (1.975,8.785)rectangle++(0.04,0.04) (2.028,8.984)rectangle++(0.04,0.04) (2.081,9.152)rectangle++(0.04,0.04) (-1.984,4.081)rectangle++(0.04,0.04) (-1.526,5.144)rectangle++(0.04,0.04)
  (-1.664,2.37)rectangle++(0.04,0.04) (-1.322,3.677)rectangle++(0.04,0.04) (1.21,0.969)rectangle++(0.04,0.04) (1.065,2.371)rectangle++(0.04,0.04) (-0.425,2.356)rectangle++(0.04,0.04)
  (-0.269,3.616)rectangle++(0.04,0.04) (-0.086,4.68)rectangle++(0.04,0.04) (1.306,1.531)rectangle++(0.04,0.04) (1.169,2.845)rectangle++(0.04,0.04) (1.105,3.968)rectangle++(0.04,0.04)
  (0.928,1.67)rectangle++(0.04,0.04) (0.853,2.978)rectangle++(0.04,0.04) (0.842,4.094)rectangle++(0.04,0.04) (0.877,5.042)rectangle++(0.04,0.04) (0.945,5.847)rectangle++(0.04,0.04)
  (1.035,6.528)rectangle++(0.04,0.04) (1.139,7.104)rectangle++(0.04,0.04) (1.25,7.589)rectangle++(0.04,0.04) (-1.744,3.546)rectangle++(0.04,0.04) (-1.343,4.68)rectangle++(0.04,0.04)
  (-0.957,5.628)rectangle++(0.04,0.04) (-0.59,6.418)rectangle++(0.04,0.04) (-0.247,7.075)rectangle++(0.04,0.04) (0.071,7.62)rectangle++(0.04,0.04) (0.363,8.07)rectangle++(0.04,0.04)
  (2.188,2.461)rectangle++(0.04,0.04) (1.956,3.601)rectangle++(0.04,0.04) (1.804,4.579)rectangle++(0.04,0.04) (1.715,5.416)rectangle++(0.04,0.04) (1.672,6.131)rectangle++(0.04,0.04)
  (1.664,6.741)rectangle++(0.04,0.04) (-1.441,3.455)rectangle++(0.04,0.04) (-1.089,4.59)rectangle++(0.04,0.04) (-0.744,5.542)rectangle++(0.04,0.04) (-0.413,6.336)rectangle++(0.04,0.04)
  (-0.1,6.999)rectangle++(0.04,0.04) (0.193,7.549)rectangle++(0.04,0.04) (0.464,8.005)rectangle++(0.04,0.04) (0.712,8.382)rectangle++(0.04,0.04) (0.938,8.692)rectangle++(0.04,0.04)
  (1.143,8.947)rectangle++(0.04,0.04) (1.327,9.156)rectangle++(0.04,0.04) (1.492,9.325)rectangle++(0.04,0.04) (1.639,9.463)rectangle++(0.04,0.04) (1.77,9.574)rectangle++(0.04,0.04)
  (1.885,9.664)rectangle++(0.04,0.04) (-0.02,1.529)rectangle++(0.04,0.04) (0.042,2.897)rectangle++(0.04,0.04) (0.149,4.057)rectangle++(0.04,0.04) (-1.046,2.516)rectangle++(0.04,0.04)
  (-0.791,3.777)rectangle++(0.04,0.04) (-0.523,4.838)rectangle++(0.04,0.04) (-0.253,5.729)rectangle++(0.04,0.04) (0.012,6.476)rectangle++(0.04,0.04) (0.267,7.101)rectangle++(0.04,0.04)
  (0.508,7.621)rectangle++(0.04,0.04) (0.735,8.054)rectangle++(0.04,0.04) (2.127,2.554)rectangle++(0.04,0.04) (1.908,3.682)rectangle++(0.04,0.04) (1.767,4.65)rectangle++(0.04,0.04)
  (1.686,5.478)rectangle++(0.04,0.04) (1.65,6.185)rectangle++(0.04,0.04) (1.648,6.787)rectangle++(0.04,0.04) (-1.456,3.461)rectangle++(0.04,0.04) (-1.212,2.015)rectangle++(0.04,0.04)
  (-0.952,3.358)rectangle++(0.04,0.04) (-0.677,4.488)rectangle++(0.04,0.04) (-0.02,0.701)rectangle++(0.04,0.04) (0.009,2.193)rectangle++(0.04,0.04) (0.093,3.46)rectangle++(0.04,0.04)
  (-0.902,2.372)rectangle++(0.04,0.04) (-0.674,3.648)rectangle++(0.04,0.04) (-0.429,4.724)rectangle++(0.04,0.04) (-0.178,5.629)rectangle++(0.04,0.04) (0.072,6.388)rectangle++(0.04,0.04)
  (0.314,7.023)rectangle++(0.04,0.04) (0.546,7.553)rectangle++(0.04,0.04) (-1.876,3.376)rectangle++(0.04,0.04) (-1.462,4.541)rectangle++(0.04,0.04) (-1.063,5.515)rectangle++(0.04,0.04)
  (-0.685,6.326)rectangle++(0.04,0.04) (-0.332,7.001)rectangle++(0.04,0.04) (-1.908,3.053)rectangle++(0.04,0.04) (-1.502,4.267)rectangle++(0.04,0.04) (-1.108,5.284)rectangle++(0.04,0.04)
  (-0.733,6.132)rectangle++(0.04,0.04) (-0.38,6.837)rectangle++(0.04,0.04) (-0.051,7.423)rectangle++(0.04,0.04) (-1.961,3.21)rectangle++(0.04,0.04) (-1.541,4.403)rectangle++(0.04,0.04)
  (-1.136,5.401)rectangle++(0.04,0.04) (-0.752,6.232)rectangle++(0.04,0.04) (-0.392,6.924)rectangle++(0.04,0.04) (-0.058,7.497)rectangle++(0.04,0.04) (0.248,7.971)rectangle++(0.04,0.04)
  (0.528,8.362)rectangle++(0.04,0.04) (0.781,8.682)rectangle++(0.04,0.04) (1.009,8.945)rectangle++(0.04,0.04) (1.213,9.159)rectangle++(0.04,0.04) (-2.16,3.883)rectangle++(0.04,0.04)
  (-1.683,4.983)rectangle++(0.04,0.04) (-1.233,5.899)rectangle++(0.04,0.04) (-0.815,6.66)rectangle++(0.04,0.04) (-0.428,7.29)rectangle++(0.04,0.04) (-0.075,7.81)rectangle++(0.04,0.04)
  (0.247,8.237)rectangle++(0.04,0.04) (0.537,8.588)rectangle++(0.04,0.04) (0.798,8.875)rectangle++(0.04,0.04) (1.031,9.108)rectangle++(0.04,0.04) (1.238,9.296)rectangle++(0.04,0.04)
  (1.422,9.449)rectangle++(0.04,0.04) (2.415,3.067)rectangle++(0.04,0.04) (2.173,4.107)rectangle++(0.04,0.04) (2.009,5)rectangle++(0.04,0.04) (1.906,5.766)rectangle++(0.04,0.04)
  (1.848,6.421)rectangle++(0.04,0.04) (1.826,6.98)rectangle++(0.04,0.04) (1.829,7.456)rectangle++(0.04,0.04) (1.851,7.861)rectangle++(0.04,0.04) (-1.695,3.744)rectangle++(0.04,0.04)
  (-1.293,4.846)rectangle++(0.04,0.04) (-0.908,5.767)rectangle++(0.04,0.04) (-0.543,6.535)rectangle++(0.04,0.04) (-0.202,7.172)rectangle++(0.04,0.04) (-1.926,3.12)rectangle++(0.04,0.04)
  (-1.515,4.326)rectangle++(0.04,0.04) (-1.117,5.334)rectangle++(0.04,0.04) (-0.738,6.174)rectangle++(0.04,0.04) (1.822,1.72)rectangle++(0.04,0.04) (1.615,2.985)rectangle++(0.04,0.04)
  (-0.02,0.461)rectangle++(0.04,0.04) (-0.001,1.989)rectangle++(0.04,0.04) (0.077,3.287)rectangle++(0.04,0.04) (0.194,4.387)rectangle++(0.04,0.04) (0.339,5.317)rectangle++(0.04,0.04)
  (0.498,6.102)rectangle++(0.04,0.04) (0.665,6.763)rectangle++(0.04,0.04) (0.834,7.318)rectangle++(0.04,0.04) (-1.757,3.391)rectangle++(0.04,0.04) (-1.36,4.549)rectangle++(0.04,0.04)
  (-0.976,5.517)rectangle++(0.04,0.04) (-0.611,6.325)rectangle++(0.04,0.04) (1.845,1.789)rectangle++(0.04,0.04) (1.638,3.043)rectangle++(0.04,0.04) (1.512,4.117)rectangle++(0.04,0.04)
  (1.447,5.035)rectangle++(0.04,0.04) (1.43,5.818)rectangle++(0.04,0.04) (-0.02,0.914)rectangle++(0.04,0.04) (0.017,2.374)rectangle++(0.04,0.04) (0.645,1.004)rectangle++(0.04,0.04)
  (0.586,2.424)rectangle++(0.04,0.04) (0.593,3.633)rectangle++(0.04,0.04) (0.647,4.661)rectangle++(0.04,0.04) (0.734,5.532)rectangle++(0.04,0.04) (0.843,6.269)rectangle++(0.04,0.04)
  (0.965,6.891)rectangle++(0.04,0.04) (1.094,7.415)rectangle++(0.04,0.04) (1.224,7.855)rectangle++(0.04,0.04) (1.353,8.224)rectangle++(0.04,0.04) (-1.889,3.709)rectangle++(0.04,0.04)
  (-1.459,4.825)rectangle++(0.04,0.04) (-1.05,5.756)rectangle++(0.04,0.04) (-0.664,6.53)rectangle++(0.04,0.04) (-0.306,7.174)rectangle++(0.04,0.04) (0.025,7.706)rectangle++(0.04,0.04)
  (0.327,8.145)rectangle++(0.04,0.04) (0.602,8.507)rectangle++(0.04,0.04) (0.85,8.803)rectangle++(0.04,0.04) (1.072,9.045)rectangle++(0.04,0.04) (1.271,9.241)rectangle++(0.04,0.04)
  (1.448,9.4)rectangle++(0.04,0.04) (1.604,9.529)rectangle++(0.04,0.04) (1.743,9.631)rectangle++(0.04,0.04) (1.864,9.713)rectangle++(0.04,0.04) (1.971,9.778)rectangle++(0.04,0.04)
  (2.064,9.828)rectangle++(0.04,0.04) (2.146,9.868)rectangle++(0.04,0.04) (2.216,9.898)rectangle++(0.04,0.04) (2.278,9.921)rectangle++(0.04,0.04) (2.331,9.938)rectangle++(0.04,0.04)
  (2.376,9.95)rectangle++(0.04,0.04) (2.416,9.959)rectangle++(0.04,0.04) (2.449,9.965)rectangle++(0.04,0.04) (2.405,3.458)rectangle++(0.04,0.04) (2.181,4.44)rectangle++(0.04,0.04)
  (-0.739,3.067)rectangle++(0.04,0.04) (-0.508,4.233)rectangle++(0.04,0.04) (-0.265,5.215)rectangle++(0.04,0.04) (-0.019,6.039)rectangle++(0.04,0.04) (0.224,6.73)rectangle++(0.04,0.04)
  (0.457,7.308)rectangle++(0.04,0.04) (0.679,7.79)rectangle++(0.04,0.04) (0.886,8.19)rectangle++(0.04,0.04) (1.079,8.522)rectangle++(0.04,0.04) (1.256,8.797)rectangle++(0.04,0.04)
  (1.417,9.024)rectangle++(0.04,0.04) (1.563,9.21)rectangle++(0.04,0.04) (1.695,9.362)rectangle++(0.04,0.04) (1.813,9.486)rectangle++(0.04,0.04) (1.918,9.587)rectangle++(0.04,0.04)
  (2.012,9.668)rectangle++(0.04,0.04) (2.094,9.734)rectangle++(0.04,0.04) (2.167,9.786)rectangle++(0.04,0.04) (2.232,9.828)rectangle++(0.04,0.04) (2.288,9.86)rectangle++(0.04,0.04)
  (2.4,3.391)rectangle++(0.04,0.04) (2.174,4.383)rectangle++(0.04,0.04) (2.021,5.235)rectangle++(0.04,0.04) (1.925,5.965)rectangle++(0.04,0.04) (1.364,2.362)rectangle++(0.04,0.04)
  (1.252,3.549)rectangle++(0.04,0.04) (0.789,1.607)rectangle++(0.04,0.04) (0.732,2.931)rectangle++(0.04,0.04) (0.738,4.058)rectangle++(0.04,0.04) (0.787,5.016)rectangle++(0.04,0.04)
  (0.867,5.828)rectangle++(0.04,0.04) (0.968,6.516)rectangle++(0.04,0.04) (1.081,7.096)rectangle++(0.04,0.04) (1.201,7.584)rectangle++(0.04,0.04) (1.322,7.995)rectangle++(0.04,0.04)
  (1.441,8.339)rectangle++(0.04,0.04) (1.556,8.627)rectangle++(0.04,0.04) (1.666,8.867)rectangle++(0.04,0.04) (1.768,9.066)rectangle++(0.04,0.04) (1.864,9.232)rectangle++(0.04,0.04)
  (2.288,3.13)rectangle++(0.04,0.04) (-0.02,0.484)rectangle++(0.04,0.04) (0,2.008)rectangle++(0.04,0.04) (0.078,3.303)rectangle++(0.04,0.04) (0.196,4.401)rectangle++(0.04,0.04)
  (0.341,5.329)rectangle++(0.04,0.04) (0.501,6.112)rectangle++(0.04,0.04) (0.668,6.772)rectangle++(0.04,0.04) (0.836,7.325)rectangle++(0.04,0.04) (1.002,7.789)rectangle++(0.04,0.04)
  (1.161,8.177)rectangle++(0.04,0.04) (1.312,8.5)rectangle++(0.04,0.04) (1.453,8.769)rectangle++(0.04,0.04) (1.583,8.992)rectangle++(0.04,0.04) (1.703,9.176)rectangle++(0.04,0.04)
  (1.813,9.328)rectangle++(0.04,0.04) (-2.084,4.058)rectangle++(0.04,0.04) (-1.611,5.129)rectangle++(0.04,0.04) (-1.167,6.02)rectangle++(0.04,0.04) (-0.753,6.76)rectangle++(0.04,0.04)
  (-0.372,7.372)rectangle++(0.04,0.04) (-0.023,7.878)rectangle++(0.04,0.04) (0.293,8.293)rectangle++(0.04,0.04) (0.579,8.634)rectangle++(0.04,0.04) (0.835,8.912)rectangle++(0.04,0.04)
  (2.353,2.786)rectangle++(0.04,0.04) (-0.275,2.743)rectangle++(0.04,0.04) (0.792,1.017)rectangle++(0.04,0.04) (0.712,2.429)rectangle++(0.04,0.04) (0.7,3.632)rectangle++(0.04,0.04)
  (-0.826,2.549)rectangle++(0.04,0.04) (-0.602,3.796)rectangle++(0.04,0.04) (-0.362,4.847)rectangle++(0.04,0.04) (-0.116,5.73)rectangle++(0.04,0.04) (0.128,6.472)rectangle++(0.04,0.04)
  (-1.678,3.042)rectangle++(0.04,0.04) (-1.307,4.249)rectangle++(0.04,0.04) (-0.943,5.26)rectangle++(0.04,0.04) (1.597,1.447)rectangle++(0.04,0.04) (-0.078,2.275)rectangle++(0.04,0.04)
  (0.022,3.533)rectangle++(0.04,0.04) (-0.935,2.371)rectangle++(0.04,0.04) (-0.702,3.649)rectangle++(0.04,0.04) (-0.453,4.726)rectangle++(0.04,0.04) (-0.198,5.632)rectangle++(0.04,0.04)
  (0.054,6.391)rectangle++(0.04,0.04) (0.3,7.026)rectangle++(0.04,0.04) (0.534,7.557)rectangle++(0.04,0.04) (0.754,7.998)rectangle++(0.04,0.04) (0.958,8.364)rectangle++(0.04,0.04)
  (1.147,8.668)rectangle++(0.04,0.04) (2.237,2.808)rectangle++(0.04,0.04) (-0.304,2.721)rectangle++(0.04,0.04) (-0.152,3.922)rectangle++(0.04,0.04) (0.026,4.936)rectangle++(0.04,0.04)
  (0.217,5.79)rectangle++(0.04,0.04) (0.414,6.509)rectangle++(0.04,0.04) (0.61,7.113)rectangle++(0.04,0.04) (0.801,7.618)rectangle++(0.04,0.04) (0.983,8.039)rectangle++(0.04,0.04)
  (1.155,8.39)rectangle++(0.04,0.04) (1.315,8.682)rectangle++(0.04,0.04) (2.216,2.856)rectangle++(0.04,0.04) (1.996,3.935)rectangle++(0.04,0.04) (1.852,4.861)rectangle++(0.04,0.04)
  (1.766,5.654)rectangle++(0.04,0.04) (1.725,6.331)rectangle++(0.04,0.04) (1.717,6.909)rectangle++(0.04,0.04) (1.734,7.4)rectangle++(0.04,0.04) (1.768,7.817)rectangle++(0.04,0.04)
  (1.813,8.17)rectangle++(0.04,0.04) (1.866,8.468)rectangle++(0.04,0.04) (1.922,8.719)rectangle++(0.04,0.04) (1.981,8.931)rectangle++(0.04,0.04) (2.038,9.108)rectangle++(0.04,0.04)
  (2.095,9.257)rectangle++(0.04,0.04) (2.26,3.196)rectangle++(0.04,0.04) (2.047,4.223)rectangle++(0.04,0.04) (1.907,5.104)rectangle++(0.04,0.04) (1.126,2.151)rectangle++(0.04,0.04)
  (1.041,3.379)rectangle++(0.04,0.04) (1.017,4.427)rectangle++(0.04,0.04) (1.04,5.318)rectangle++(0.04,0.04) (1.094,6.075)rectangle++(0.04,0.04) (1.52,2.173)rectangle++(0.04,0.04)
  (0.363,1.346)rectangle++(0.04,0.04) (0.36,2.726)rectangle++(0.04,0.04) (0.413,3.899)rectangle++(0.04,0.04) (0.505,4.894)rectangle++(0.04,0.04) (0.623,5.736)rectangle++(0.04,0.04)
  (0.756,6.447)rectangle++(0.04,0.04) (1.674,2.174)rectangle++(0.04,0.04) (1.508,3.377)rectangle++(0.04,0.04) (1.415,4.406)rectangle++(0.04,0.04) (1.376,5.285)rectangle++(0.04,0.04)
  (1.379,6.033)rectangle++(0.04,0.04) (1.465,2.237)rectangle++(0.04,0.04) (1.333,3.439)rectangle++(0.04,0.04) (0.746,1.602)rectangle++(0.04,0.04) (0.696,2.928)rectangle++(0.04,0.04)
  (0.706,4.057)rectangle++(0.04,0.04) (0.76,5.016)rectangle++(0.04,0.04) (1.273,1.832)rectangle++(0.04,0.04) (1.153,3.102)rectangle++(0.04,0.04) (1.102,4.187)rectangle++(0.04,0.04)
  (1.102,5.111)rectangle++(0.04,0.04) (1.139,5.896)rectangle++(0.04,0.04) (1.202,6.563)rectangle++(0.04,0.04) (1.282,7.126)rectangle++(0.04,0.04) (1.372,7.602)rectangle++(0.04,0.04)
  (1.468,8.003)rectangle++(0.04,0.04) (2.003,2.733)rectangle++(0.04,0.04) (1.81,3.839)rectangle++(0.04,0.04) (1.69,4.787)rectangle++(0.04,0.04) (1.626,5.597)rectangle++(0.04,0.04)
  (1.603,6.289)rectangle++(0.04,0.04) (1.612,6.878)rectangle++(0.04,0.04) (1.643,7.378)rectangle++(0.04,0.04) (1.69,7.802)rectangle++(0.04,0.04) (1.746,8.16)rectangle++(0.04,0.04)
  (2.005,2.842)rectangle++(0.04,0.04) (0.478,1.634)rectangle++(0.04,0.04) (0.469,2.966)rectangle++(0.04,0.04) (0.515,4.098)rectangle++(0.04,0.04) (1.053,1.548)rectangle++(0.04,0.04)
  (0.955,2.869)rectangle++(0.04,0.04) (0.924,3.997)rectangle++(0.04,0.04) (0.943,4.957)rectangle++(0.04,0.04) (0.998,5.772)rectangle++(0.04,0.04) (1.077,6.462)rectangle++(0.04,0.04);
\fill[nuage]
  (1.171,7.046)rectangle++(0.04,0.04) (1.275,7.538)rectangle++(0.04,0.04) (1.383,7.953)rectangle++(0.04,0.04) (1.492,8.301)rectangle++(0.04,0.04) (1.598,8.592)rectangle++(0.04,0.04)
  (1.7,8.836)rectangle++(0.04,0.04) (2.202,2.992)rectangle++(0.04,0.04) (1.989,4.052)rectangle++(0.04,0.04) (1.85,4.961)rectangle++(0.04,0.04) (1.769,5.739)rectangle++(0.04,0.04)
  (1.731,6.403)rectangle++(0.04,0.04) (1.725,6.97)rectangle++(0.04,0.04) (-1.488,3.519)rectangle++(0.04,0.04) (-1.126,4.647)rectangle++(0.04,0.04) (-0.774,5.591)rectangle++(0.04,0.04)
  (-0.436,6.38)rectangle++(0.04,0.04) (-0.118,7.036)rectangle++(0.04,0.04) (0.179,7.582)rectangle++(0.04,0.04) (0.453,8.034)rectangle++(0.04,0.04) (0.704,8.407)rectangle++(0.04,0.04)
  (0.933,8.714)rectangle++(0.04,0.04) (1.139,8.966)rectangle++(0.04,0.04) (1.325,9.171)rectangle++(0.04,0.04) (1.491,9.339)rectangle++(0.04,0.04) (1.638,9.475)rectangle++(0.04,0.04)
  (1.769,9.584)rectangle++(0.04,0.04) (1.885,9.672)rectangle++(0.04,0.04) (1.987,9.742)rectangle++(0.04,0.04) (2.077,9.797)rectangle++(0.04,0.04) (2.155,9.841)rectangle++(0.04,0.04)
  (2.223,9.875)rectangle++(0.04,0.04) (2.282,9.901)rectangle++(0.04,0.04) (2.334,9.921)rectangle++(0.04,0.04) (2.378,9.935)rectangle++(0.04,0.04) (2.417,9.946)rectangle++(0.04,0.04)
  (2.45,9.954)rectangle++(0.04,0.04) (2.478,9.959)rectangle++(0.04,0.04) (2.503,9.962)rectangle++(0.04,0.04) (2.524,9.964)rectangle++(0.04,0.04) (2.541,9.965)rectangle++(0.04,0.04)
  (2.557,9.964)rectangle++(0.04,0.04) (2.57,9.964)rectangle++(0.04,0.04) (2.58,9.963)rectangle++(0.04,0.04) (2.59,9.961)rectangle++(0.04,0.04) (2.597,9.96)rectangle++(0.04,0.04)
  (2.604,9.958)rectangle++(0.04,0.04) (-2.089,4.379)rectangle++(0.04,0.04) (-1.603,5.402)rectangle++(0.04,0.04) (-1.149,6.252)rectangle++(0.04,0.04) (-0.729,6.956)rectangle++(0.04,0.04)
  (-0.343,7.538)rectangle++(0.04,0.04) (0.008,8.017)rectangle++(0.04,0.04) (0.325,8.411)rectangle++(0.04,0.04) (2.289,2.533)rectangle++(0.04,0.04) (-0.222,2.673)rectangle++(0.04,0.04)
  (-0.084,3.877)rectangle++(0.04,0.04) (0.081,4.895)rectangle++(0.04,0.04) (0.263,5.754)rectangle++(0.04,0.04) (0.451,6.476)rectangle++(0.04,0.04) (0.641,7.083)rectangle++(0.04,0.04)
  (0.826,7.591)rectangle++(0.04,0.04) (-0.02,1.198)rectangle++(0.04,0.04) (0.029,2.615)rectangle++(0.04,0.04) (0.127,3.818)rectangle++(0.04,0.04) (0.258,4.836)rectangle++(0.04,0.04)
  (0.411,5.697)rectangle++(0.04,0.04) (-1.42,2.937)rectangle++(0.04,0.04) (-1.092,4.149)rectangle++(0.04,0.04) (-0.764,5.167)rectangle++(0.04,0.04) (-0.445,6.019)rectangle++(0.04,0.04)
  (-0.14,6.73)rectangle++(0.04,0.04) (0.148,7.322)rectangle++(0.04,0.04) (0.417,7.814)rectangle++(0.04,0.04) (0.664,8.221)rectangle++(0.04,0.04) (-2.026,3.551)rectangle++(0.04,0.04)
  (-1.582,4.695)rectangle++(0.04,0.04) (-1.159,5.65)rectangle++(0.04,0.04) (-0.762,6.445)rectangle++(0.04,0.04) (1.902,1.779)rectangle++(0.04,0.04) (1.685,3.032)rectangle++(0.04,0.04)
  (1.552,4.106)rectangle++(0.04,0.04) (1.481,5.024)rectangle++(0.04,0.04) (1.458,5.808)rectangle++(0.04,0.04) (1.469,6.474)rectangle++(0.04,0.04) (1.505,7.041)rectangle++(0.04,0.04)
  (1.559,7.521)rectangle++(0.04,0.04) (1.624,7.926)rectangle++(0.04,0.04) (-1.757,3.706)rectangle++(0.04,0.04) (-1.336,2)rectangle++(0.04,0.04) (-0.809,1.722)rectangle++(0.04,0.04)
  (-0.621,3.092)rectangle++(0.04,0.04) (-0.406,4.249)rectangle++(0.04,0.04) (-0.177,5.224)rectangle++(0.04,0.04) (0.056,6.044)rectangle++(0.04,0.04) (0.287,6.731)rectangle++(0.04,0.04)
  (0.511,7.306)rectangle++(0.04,0.04) (-1.819,3.314)rectangle++(0.04,0.04) (-1.415,4.486)rectangle++(0.04,0.04) (-1.026,5.466)rectangle++(0.04,0.04) (-0.656,6.283)rectangle++(0.04,0.04)
  (-1.786,2.821)rectangle++(0.04,0.04) (-1.407,4.065)rectangle++(0.04,0.04) (-1.036,5.108)rectangle++(0.04,0.04) (-0.678,5.979)rectangle++(0.04,0.04) (-0.34,6.706)rectangle++(0.04,0.04)
  (-0.023,7.31)rectangle++(0.04,0.04) (0.271,7.81)rectangle++(0.04,0.04) (0.54,8.224)rectangle++(0.04,0.04) (0.786,8.565)rectangle++(0.04,0.04) (2.263,2.69)rectangle++(0.04,0.04)
  (2.029,3.792)rectangle++(0.04,0.04) (1.874,4.738)rectangle++(0.04,0.04) (1.78,5.549)rectangle++(0.04,0.04) (1.733,6.242)rectangle++(0.04,0.04) (1.72,6.832)rectangle++(0.04,0.04)
  (1.734,7.335)rectangle++(0.04,0.04) (-1.582,3.601)rectangle++(0.04,0.04) (-1.202,4.721)rectangle++(0.04,0.04) (-0.835,5.657)rectangle++(0.04,0.04) (1.692,1.57)rectangle++(0.04,0.04)
  (0.169,1.247)rectangle++(0.04,0.04) (0.191,2.649)rectangle++(0.04,0.04) (0.266,3.84)rectangle++(0.04,0.04) (0.378,4.85)rectangle++(0.04,0.04) (0.513,5.704)rectangle++(0.04,0.04)
  (0.662,6.424)rectangle++(0.04,0.04) (0.817,7.03)rectangle++(0.04,0.04) (0.974,7.539)rectangle++(0.04,0.04) (1.127,7.965)rectangle++(0.04,0.04) (1.274,8.322)rectangle++(0.04,0.04)
  (1.414,8.619)rectangle++(0.04,0.04) (1.544,8.865)rectangle++(0.04,0.04) (1.665,9.07)rectangle++(0.04,0.04) (1.776,9.239)rectangle++(0.04,0.04) (1.877,9.378)rectangle++(0.04,0.04)
  (2.327,3.169)rectangle++(0.04,0.04) (2.103,4.197)rectangle++(0.04,0.04) (1.953,5.079)rectangle++(0.04,0.04) (1.861,5.835)rectangle++(0.04,0.04) (1.813,6.482)rectangle++(0.04,0.04)
  (-1.344,3.432)rectangle++(0.04,0.04) (-1.007,4.567)rectangle++(0.04,0.04) (-0.676,5.519)rectangle++(0.04,0.04) (-0.356,6.314)rectangle++(0.04,0.04) (-0.052,6.977)rectangle++(0.04,0.04)
  (-1.846,3.112)rectangle++(0.04,0.04) (-1.447,4.315)rectangle++(0.04,0.04) (-1.059,5.322)rectangle++(0.04,0.04) (-0.69,6.162)rectangle++(0.04,0.04) (-0.342,6.862)rectangle++(0.04,0.04)
  (-0.018,7.442)rectangle++(0.04,0.04) (0.28,7.923)rectangle++(0.04,0.04) (-2.025,3.396)rectangle++(0.04,0.04) (-1.309,1.87)rectangle++(0.04,0.04) (-1.04,3.238)rectangle++(0.04,0.04)
  (-0.757,4.39)rectangle++(0.04,0.04) (-0.47,5.358)rectangle++(0.04,0.04) (1.553,1.594)rectangle++(0.04,0.04) (1.382,2.889)rectangle++(0.04,0.04) (1.288,3.996)rectangle++(0.04,0.04)
  (0.908,1.724)rectangle++(0.04,0.04) (0.839,3.025)rectangle++(0.04,0.04) (0.832,4.134)rectangle++(0.04,0.04) (1.015,1.638)rectangle++(0.04,0.04) (0.926,2.948)rectangle++(0.04,0.04)
  (0.903,4.065)rectangle++(0.04,0.04) (0.928,5.015)rectangle++(0.04,0.04) (0.987,5.822)rectangle++(0.04,0.04) (1.07,6.506)rectangle++(0.04,0.04) (-1.499,3.266)rectangle++(0.04,0.04)
  (-1.145,4.433)rectangle++(0.04,0.04) (-0.799,5.41)rectangle++(0.04,0.04) (-0.465,6.226)rectangle++(0.04,0.04) (-0.148,6.907)rectangle++(0.04,0.04) (0.148,7.473)rectangle++(0.04,0.04)
  (0.423,7.943)rectangle++(0.04,0.04) (0.675,8.33)rectangle++(0.04,0.04) (0.905,8.65)rectangle++(0.04,0.04) (1.113,8.913)rectangle++(0.04,0.04) (1.3,9.127)rectangle++(0.04,0.04)
  (1.468,9.302)rectangle++(0.04,0.04) (2.367,3.044)rectangle++(0.04,0.04) (2.132,4.089)rectangle++(0.04,0.04) (1.973,4.987)rectangle++(0.04,0.04) (1.875,5.756)rectangle++(0.04,0.04)
  (-1.143,3.286)rectangle++(0.04,0.04) (-0.842,4.435)rectangle++(0.04,0.04) (-0.541,5.4)rectangle++(0.04,0.04) (-0.246,6.208)rectangle++(0.04,0.04) (0.037,6.883)rectangle++(0.04,0.04)
  (0.305,7.445)rectangle++(0.04,0.04) (0.555,7.912)rectangle++(0.04,0.04) (-1.967,3.457)rectangle++(0.04,0.04) (-1.536,4.614)rectangle++(0.04,0.04) (-1.123,5.579)rectangle++(0.04,0.04)
  (-0.734,6.383)rectangle++(0.04,0.04) (-0.371,7.051)rectangle++(0.04,0.04) (-0.035,7.605)rectangle++(0.04,0.04) (2.117,2.246)rectangle++(0.04,0.04) (1.887,3.421)rectangle++(0.04,0.04)
  (1.739,4.428)rectangle++(0.04,0.04) (1.653,5.291)rectangle++(0.04,0.04) (-1.066,3.133)rectangle++(0.04,0.04) (-1.049,2.033)rectangle++(0.04,0.04) (-0.813,3.366)rectangle++(0.04,0.04)
  (-0.558,4.49)rectangle++(0.04,0.04) (-0.297,5.435)rectangle++(0.04,0.04) (-1.494,2.716)rectangle++(0.04,0.04) (-1.163,3.965)rectangle++(0.04,0.04) (-0.832,5.013)rectangle++(0.04,0.04)
  (-0.509,5.89)rectangle++(0.04,0.04) (-0.199,6.623)rectangle++(0.04,0.04) (0.093,7.234)rectangle++(0.04,0.04) (0.366,7.741)rectangle++(0.04,0.04) (0.619,8.162)rectangle++(0.04,0.04)
  (0.85,8.509)rectangle++(0.04,0.04) (-2.063,3.657)rectangle++(0.04,0.04) (-1.61,4.787)rectangle++(0.04,0.04) (-1.179,5.729)rectangle++(0.04,0.04) (-0.775,6.513)rectangle++(0.04,0.04)
  (-0.401,7.164)rectangle++(0.04,0.04) (-0.056,7.701)rectangle++(0.04,0.04) (-2.035,3.27)rectangle++(0.04,0.04) (-1.601,4.457)rectangle++(0.04,0.04) (-1.185,5.449)rectangle++(0.04,0.04)
  (-0.791,6.275)rectangle++(0.04,0.04) (-0.424,6.962)rectangle++(0.04,0.04) (-0.084,7.531)rectangle++(0.04,0.04) (0.228,8.001)rectangle++(0.04,0.04) (0.511,8.388)rectangle++(0.04,0.04)
  (0.768,8.705)rectangle++(0.04,0.04) (0.999,8.965)rectangle++(0.04,0.04) (1.205,9.176)rectangle++(0.04,0.04) (2.371,2.946)rectangle++(0.04,0.04) (2.131,4.005)rectangle++(0.04,0.04)
  (1.969,4.915)rectangle++(0.04,0.04) (1.868,5.696)rectangle++(0.04,0.04) (1.814,6.363)rectangle++(0.04,0.04) (1.794,6.932)rectangle++(0.04,0.04) (1.8,7.417)rectangle++(0.04,0.04)
  (1.824,7.828)rectangle++(0.04,0.04) (1.862,8.177)rectangle++(0.04,0.04) (1.907,8.472)rectangle++(0.04,0.04) (1.958,8.721)rectangle++(0.04,0.04) (2.011,8.931)rectangle++(0.04,0.04)
  (-1.941,4.016)rectangle++(0.04,0.04) (-1.491,5.088)rectangle++(0.04,0.04) (-1.066,5.98)rectangle++(0.04,0.04) (-0.669,6.722)rectangle++(0.04,0.04) (-0.302,7.337)rectangle++(0.04,0.04)
  (-1.989,3.134)rectangle++(0.04,0.04) (-1.568,4.339)rectangle++(0.04,0.04) (-1.161,5.347)rectangle++(0.04,0.04) (-0.775,6.188)rectangle++(0.04,0.04) (-0.414,6.887)rectangle++(0.04,0.04)
  (-0.078,7.467)rectangle++(0.04,0.04) (0.23,7.946)rectangle++(0.04,0.04) (0.511,8.341)rectangle++(0.04,0.04) (0.766,8.666)rectangle++(0.04,0.04) (0.995,8.931)rectangle++(0.04,0.04)
  (1.201,9.148)rectangle++(0.04,0.04) (1.385,9.324)rectangle++(0.04,0.04) (1.548,9.466)rectangle++(0.04,0.04) (1.692,9.581)rectangle++(0.04,0.04) (-2.174,4.086)rectangle++(0.04,0.04)
  (-1.686,5.156)rectangle++(0.04,0.04) (-1.229,6.046)rectangle++(0.04,0.04) (-0.805,6.785)rectangle++(0.04,0.04) (-0.415,7.396)rectangle++(0.04,0.04) (-0.059,7.899)rectangle++(0.04,0.04)
  (0.263,8.313)rectangle++(0.04,0.04) (0.554,8.651)rectangle++(0.04,0.04) (0.815,8.928)rectangle++(0.04,0.04) (1.047,9.152)rectangle++(0.04,0.04) (1.254,9.334)rectangle++(0.04,0.04)
  (1.437,9.48)rectangle++(0.04,0.04) (1.599,9.596)rectangle++(0.04,0.04) (1.741,9.689)rectangle++(0.04,0.04) (1.865,9.762)rectangle++(0.04,0.04) (1.973,9.82)rectangle++(0.04,0.04)
  (2.068,9.864)rectangle++(0.04,0.04) (2.15,9.898)rectangle++(0.04,0.04) (2.221,9.923)rectangle++(0.04,0.04) (2.283,9.942)rectangle++(0.04,0.04) (2.336,9.956)rectangle++(0.04,0.04)
  (2.382,9.965)rectangle++(0.04,0.04) (2.421,9.971)rectangle++(0.04,0.04) (2.454,9.975)rectangle++(0.04,0.04) (2.407,3.462)rectangle++(0.04,0.04) (2.183,4.443)rectangle++(0.04,0.04)
  (0.899,2.064)rectangle++(0.04,0.04) (0.845,3.314)rectangle++(0.04,0.04) (0.848,4.38)rectangle++(0.04,0.04) (0.894,5.285)rectangle++(0.04,0.04) (1.328,1.931)rectangle++(0.04,0.04)
  (1.204,3.184)rectangle++(0.04,0.04) (1.149,4.255)rectangle++(0.04,0.04) (1.144,5.167)rectangle++(0.04,0.04) (1.177,5.942)rectangle++(0.04,0.04) (-1.331,3.167)rectangle++(0.04,0.04)
  (-1.007,4.341)rectangle++(0.04,0.04) (-0.684,5.327)rectangle++(0.04,0.04) (-0.371,6.151)rectangle++(0.04,0.04) (-0.071,6.84)rectangle++(0.04,0.04) (0.211,7.413)rectangle++(0.04,0.04)
  (-1.906,3.268)rectangle++(0.04,0.04) (1.184,0.719)rectangle++(0.04,0.04) (1.033,2.16)rectangle++(0.04,0.04) (0.962,3.391)rectangle++(0.04,0.04) (0.951,4.44)rectangle++(0.04,0.04)
  (0.984,5.332)rectangle++(0.04,0.04) (1.047,6.089)rectangle++(0.04,0.04) (1.53,2.164)rectangle++(0.04,0.04) (1.385,3.374)rectangle++(0.04,0.04) (1.31,4.409)rectangle++(0.04,0.04)
  (1.288,5.291)rectangle++(0.04,0.04) (1.304,6.042)rectangle++(0.04,0.04) (1.348,6.68)rectangle++(0.04,0.04) (1.411,7.22)rectangle++(0.04,0.04) (1.793,2.53)rectangle++(0.04,0.04)
  (1.623,3.675)rectangle++(0.04,0.04) (-0.02,0.571)rectangle++(0.04,0.04) (0.004,2.082)rectangle++(0.04,0.04) (0.084,3.366)rectangle++(0.04,0.04) (0.204,4.454)rectangle++(0.04,0.04)
  (0.349,5.374)rectangle++(0.04,0.04) (0.51,6.15)rectangle++(0.04,0.04) (0.677,6.803)rectangle++(0.04,0.04) (0.845,7.352)rectangle++(0.04,0.04) (1.011,7.812)rectangle++(0.04,0.04)
  (1.169,8.196)rectangle++(0.04,0.04) (1.319,8.516)rectangle++(0.04,0.04) (1.46,8.782)rectangle++(0.04,0.04) (1.59,9.002)rectangle++(0.04,0.04) (1.709,9.185)rectangle++(0.04,0.04)
  (1.818,9.335)rectangle++(0.04,0.04) (1.917,9.458)rectangle++(0.04,0.04) (2.005,9.559)rectangle++(0.04,0.04) (2.085,9.641)rectangle++(0.04,0.04) (2.155,9.708)rectangle++(0.04,0.04)
  (-2.114,4.22)rectangle++(0.04,0.04) (-1.63,5.268)rectangle++(0.04,0.04) (-1.177,6.139)rectangle++(0.04,0.04) (-0.757,6.862)rectangle++(0.04,0.04) (-0.371,7.459)rectangle++(0.04,0.04)
  (-0.02,7.951)rectangle++(0.04,0.04) (2.212,2.333)rectangle++(0.04,0.04) (1.971,3.491)rectangle++(0.04,0.04) (1.813,4.485)rectangle++(0.04,0.04) (1.718,5.336)rectangle++(0.04,0.04)
  (1.672,6.063)rectangle++(0.04,0.04) (1.661,6.683)rectangle++(0.04,0.04) (1.677,7.21)rectangle++(0.04,0.04) (1.712,7.658)rectangle++(0.04,0.04) (1.759,8.037)rectangle++(0.04,0.04)
  (1.815,8.357)rectangle++(0.04,0.04) (1.874,8.627)rectangle++(0.04,0.04) (1.936,8.854)rectangle++(0.04,0.04) (-1.936,3.982)rectangle++(0.04,0.04) (-1.489,5.059)rectangle++(0.04,0.04)
  (-1.065,5.956)rectangle++(0.04,0.04) (-0.669,6.701)rectangle++(0.04,0.04) (-1.897,2.909)rectangle++(0.04,0.04) (-1.157,1.793)rectangle++(0.04,0.04) (-0.914,3.166)rectangle++(0.04,0.04)
  (-0.652,4.324)rectangle++(0.04,0.04) (-0.384,5.298)rectangle++(0.04,0.04) (-0.117,6.115)rectangle++(0.04,0.04) (0.143,6.798)rectangle++(0.04,0.04) (0.392,7.369)rectangle++(0.04,0.04)
  (0.625,7.844)rectangle++(0.04,0.04) (0.843,8.239)rectangle++(0.04,0.04) (1.044,8.565)rectangle++(0.04,0.04) (1.228,8.835)rectangle++(0.04,0.04) (1.395,9.057)rectangle++(0.04,0.04)
  (1.546,9.239)rectangle++(0.04,0.04) (2.338,3.049)rectangle++(0.04,0.04) (2.107,4.095)rectangle++(0.04,0.04) (1.952,4.992)rectangle++(0.04,0.04) (1.857,5.762)rectangle++(0.04,0.04)
  (1.807,6.419)rectangle++(0.04,0.04) (1.79,6.98)rectangle++(0.04,0.04) (1.799,7.458)rectangle++(0.04,0.04) (1.825,7.863)rectangle++(0.04,0.04) (1.864,8.207)rectangle++(0.04,0.04)
  (1.91,8.498)rectangle++(0.04,0.04) (1.961,8.743)rectangle++(0.04,0.04) (2.015,8.949)rectangle++(0.04,0.04) (2.186,3.102)rectangle++(0.04,0.04) (1.98,4.145)rectangle++(0.04,0.04)
  (1.847,5.04)rectangle++(0.04,0.04) (1.769,5.807)rectangle++(0.04,0.04) (1.734,6.461)rectangle++(0.04,0.04) (1.73,7.019)rectangle++(0.04,0.04) (1.749,7.493)rectangle++(0.04,0.04)
  (1.784,7.895)rectangle++(0.04,0.04) (1.83,8.236)rectangle++(0.04,0.04) (2.014,2.882)rectangle++(0.04,0.04) (0.488,1.645)rectangle++(0.04,0.04) (0.478,2.975)rectangle++(0.04,0.04)
  (0.523,4.106)rectangle++(0.04,0.04) (0.607,5.065)rectangle++(0.04,0.04) (1.31,1.803)rectangle++(0.04,0.04) (1.183,3.077)rectangle++(0.04,0.04) (1.127,4.164)rectangle++(0.04,0.04)
  (0.98,1.722)rectangle++(0.04,0.04) (0.899,3.021)rectangle++(0.04,0.04) (0.883,4.128)rectangle++(0.04,0.04) (0.914,5.07)rectangle++(0.04,0.04) (0.977,5.869)rectangle++(0.04,0.04)
  (1.479,2.093)rectangle++(0.04,0.04) (1.339,3.316)rectangle++(0.04,0.04) (1.269,4.361)rectangle++(0.04,0.04) (1.25,5.252)rectangle++(0.04,0.04) (1.271,6.011)rectangle++(0.04,0.04)
  (1.318,6.654)rectangle++(0.04,0.04) (1.385,7.2)rectangle++(0.04,0.04) (1.463,7.661)rectangle++(0.04,0.04) (1.548,8.049)rectangle++(0.04,0.04) (1.635,8.376)rectangle++(0.04,0.04)
  (1.723,8.65)rectangle++(0.04,0.04) (1.808,8.88)rectangle++(0.04,0.04) (1.89,9.072)rectangle++(0.04,0.04) (1.967,9.232)rectangle++(0.04,0.04) (2.039,9.365)rectangle++(0.04,0.04)
  (2.106,9.475)rectangle++(0.04,0.04) (2.167,9.565)rectangle++(0.04,0.04) (2.222,9.64)rectangle++(0.04,0.04) (2.349,3.321)rectangle++(0.04,0.04) (2.127,4.325)rectangle++(0.04,0.04)
  (1.979,5.188)rectangle++(0.04,0.04) (1.887,5.927)rectangle++(0.04,0.04) (1.839,6.558)rectangle++(0.04,0.04) (1.823,7.097)rectangle++(0.04,0.04) (-1.502,3.57)rectangle++(0.04,0.04)
  (-1.136,4.691)rectangle++(0.04,0.04) (-1.468,2.36)rectangle++(0.04,0.04) (-1.156,3.661)rectangle++(0.04,0.04) (-0.838,4.754)rectangle++(0.04,0.04) (-0.524,5.671)rectangle++(0.04,0.04)
  (-0.221,6.437)rectangle++(0.04,0.04) (0.067,7.077)rectangle++(0.04,0.04) (0.338,7.609)rectangle++(0.04,0.04) (0.59,8.05)rectangle++(0.04,0.04) (0.821,8.415)rectangle++(0.04,0.04)
  (2.216,2.663)rectangle++(0.04,0.04) (1.988,3.771)rectangle++(0.04,0.04) (1.838,4.722)rectangle++(0.04,0.04) (1.749,5.537)rectangle++(0.04,0.04) (1.706,6.232)rectangle++(0.04,0.04)
  (1.697,6.825)rectangle++(0.04,0.04) (1.713,7.33)rectangle++(0.04,0.04) (1.747,7.758)rectangle++(0.04,0.04) (1.793,8.121)rectangle++(0.04,0.04) (1.847,8.427)rectangle++(0.04,0.04)
  (1.905,8.685)rectangle++(0.04,0.04) (1.964,8.903)rectangle++(0.04,0.04) (-1.943,3.999)rectangle++(0.04,0.04) (-1.494,5.073)rectangle++(0.04,0.04) (-1.069,5.968)rectangle++(0.04,0.04)
  (-0.672,6.712)rectangle++(0.04,0.04) (-0.305,7.328)rectangle++(0.04,0.04) (0.032,7.837)rectangle++(0.04,0.04) (0.338,8.257)rectangle++(0.04,0.04) (0.616,8.601)rectangle++(0.04,0.04)
  (0.865,8.882)rectangle++(0.04,0.04) (1.088,9.112)rectangle++(0.04,0.04) (1.287,9.298)rectangle++(0.04,0.04) (1.464,9.448)rectangle++(0.04,0.04) (1.62,9.568)rectangle++(0.04,0.04)
  (1.758,9.664)rectangle++(0.04,0.04) (1.878,9.741)rectangle++(0.04,0.04) (1.984,9.801)rectangle++(0.04,0.04) (2.076,9.847)rectangle++(0.04,0.04) (2.156,9.883)rectangle++(0.04,0.04)
  (2.226,9.911)rectangle++(0.04,0.04) (2.286,9.931)rectangle++(0.04,0.04) (2.339,9.946)rectangle++(0.04,0.04) (2.383,9.957)rectangle++(0.04,0.04) (2.422,9.964)rectangle++(0.04,0.04)
  (2.455,9.969)rectangle++(0.04,0.04) (2.483,9.972)rectangle++(0.04,0.04) (2.508,9.973)rectangle++(0.04,0.04) (2.528,9.973)rectangle++(0.04,0.04) (-2.108,4.364)rectangle++(0.04,0.04)
  (-1.62,5.39)rectangle++(0.04,0.04) (-1.163,6.243)rectangle++(0.04,0.04) (-0.741,6.949)rectangle++(0.04,0.04) (-0.354,7.533)rectangle++(0.04,0.04) (-0.002,8.013)rectangle++(0.04,0.04)
  (0.316,8.407)rectangle++(0.04,0.04) (0.603,8.73)rectangle++(0.04,0.04) (-2.17,3.648)rectangle++(0.04,0.04) (-1.701,4.784)rectangle++(0.04,0.04) (-1.257,5.731)rectangle++(0.04,0.04)
  (-1.763,2.561)rectangle++(0.04,0.04) (-1.398,3.843)rectangle++(0.04,0.04) (-1.037,4.919)rectangle++(0.04,0.04) (-0.687,5.819)rectangle++(0.04,0.04) (-0.353,6.57)rectangle++(0.04,0.04)
  (-0.02,1.034)rectangle++(0.04,0.04) (0.022,2.476)rectangle++(0.04,0.04) (0.116,3.7)rectangle++(0.04,0.04) (0.244,4.737)rectangle++(0.04,0.04) (0.395,5.613)rectangle++(0.04,0.04)
  (0.558,6.351)rectangle++(0.04,0.04) (0.726,6.972)rectangle++(0.04,0.04) (0.894,7.494)rectangle++(0.04,0.04) (1.057,7.93)rectangle++(0.04,0.04) (1.214,8.294)rectangle++(0.04,0.04)
  (1.361,8.598)rectangle++(0.04,0.04) (1.499,8.85)rectangle++(0.04,0.04) (-2.022,3.881)rectangle++(0.04,0.04) (-1.566,4.976)rectangle++(0.04,0.04) (-1.134,5.888)rectangle++(0.04,0.04)
  (-0.731,6.647)rectangle++(0.04,0.04) (-0.358,7.275)rectangle++(0.04,0.04) (-0.015,7.794)rectangle++(0.04,0.04) (0.297,8.222)rectangle++(0.04,0.04) (0.579,8.573)rectangle++(0.04,0.04)
  (0.833,8.86)rectangle++(0.04,0.04) (1.06,9.094)rectangle++(0.04,0.04) (-2.174,3.833)rectangle++(0.04,0.04) (-1.696,4.942)rectangle++(0.04,0.04) (-1.247,5.864)rectangle++(0.04,0.04)
  (-0.827,6.631)rectangle++(0.04,0.04) (-0.44,7.265)rectangle++(0.04,0.04) (-0.086,7.789)rectangle++(0.04,0.04) (0.237,8.221)rectangle++(0.04,0.04) (0.528,8.574)rectangle++(0.04,0.04)
  (0.789,8.863)rectangle++(0.04,0.04) (-2.168,3.72)rectangle++(0.04,0.04) (-1.696,4.845)rectangle++(0.04,0.04) (-1.25,5.783)rectangle++(0.04,0.04) (-0.833,6.561)rectangle++(0.04,0.04)
  (-0.448,7.207)rectangle++(0.04,0.04) (-0.095,7.74)rectangle++(0.04,0.04) (0.227,8.179)rectangle++(0.04,0.04) (0.518,8.539)rectangle++(0.04,0.04) (0.779,8.834)rectangle++(0.04,0.04)
  (1.014,9.074)rectangle++(0.04,0.04) (1.222,9.268)rectangle++(0.04,0.04) (1.408,9.425)rectangle++(0.04,0.04) (1.571,9.551)rectangle++(0.04,0.04) (1.715,9.652)rectangle++(0.04,0.04)
  (1.842,9.732)rectangle++(0.04,0.04) (1.953,9.795)rectangle++(0.04,0.04) (2.049,9.844)rectangle++(0.04,0.04) (2.134,9.881)rectangle++(0.04,0.04) (2.207,9.91)rectangle++(0.04,0.04)
  (2.27,9.931)rectangle++(0.04,0.04) (2.324,9.947)rectangle++(0.04,0.04) (-0.02,1.575)rectangle++(0.04,0.04) (-0.435,1.931)rectangle++(0.04,0.04) (-0.294,3.255)rectangle++(0.04,0.04)
  (-0.926,2.237)rectangle++(0.04,0.04) (-0.7,3.535)rectangle++(0.04,0.04) (-1.08,2.206)rectangle++(0.04,0.04) (-0.832,3.514)rectangle++(0.04,0.04) (-0.569,4.617)rectangle++(0.04,0.04)
  (-0.301,5.543)rectangle++(0.04,0.04) (-0.037,6.32)rectangle++(0.04,0.04) (-1.672,2.971)rectangle++(0.04,0.04) (-1.304,4.188)rectangle++(0.04,0.04) (1.351,1.096)rectangle++(0.04,0.04)
  (1.19,2.474)rectangle++(0.04,0.04) (1.108,3.651)rectangle++(0.04,0.04) (1.086,4.656)rectangle++(0.04,0.04) (1.107,5.51)rectangle++(0.04,0.04) (1.159,6.235)rectangle++(0.04,0.04)
  (-1.411,3.227)rectangle++(0.04,0.04) (-1.072,4.396)rectangle++(0.04,0.04) (-0.738,5.376)rectangle++(0.04,0.04) (-0.414,6.195)rectangle++(0.04,0.04) (1.779,1.809)rectangle++(0.04,0.04)
  (1.583,3.063)rectangle++(0.04,0.04) (1.466,4.136)rectangle++(0.04,0.04) (0.921,1.804)rectangle++(0.04,0.04) (0.853,3.093)rectangle++(0.04,0.04) (0.846,4.191)rectangle++(0.04,0.04)
  (0.885,5.124)rectangle++(0.04,0.04) (0.955,5.917)rectangle++(0.04,0.04) (1.046,6.587)rectangle++(0.04,0.04) (1.15,7.153)rectangle++(0.04,0.04) (-1.651,3.427)rectangle++(0.04,0.04)
  (-1.268,4.575)rectangle++(0.04,0.04) (-0.897,5.536)rectangle++(0.04,0.04) (-0.544,6.338)rectangle++(0.04,0.04) (-0.211,7.005)rectangle++(0.04,0.04) (0.099,7.559)rectangle++(0.04,0.04)
  (0.384,8.017)rectangle++(0.04,0.04) (0.645,8.396)rectangle++(0.04,0.04) (-2.075,3.584)rectangle++(0.04,0.04) (-1.623,4.726)rectangle++(0.04,0.04) (-1.192,5.678)rectangle++(0.04,0.04)
  (-0.789,6.47)rectangle++(0.04,0.04) (-0.414,7.128)rectangle++(0.04,0.04) (-0.069,7.671)rectangle++(0.04,0.04) (0.246,8.119)rectangle++(0.04,0.04) (0.532,8.488)rectangle++(0.04,0.04)
  (0.789,8.79)rectangle++(0.04,0.04) (1.02,9.036)rectangle++(0.04,0.04) (1.227,9.236)rectangle++(0.04,0.04) (1.41,9.398)rectangle++(0.04,0.04) (1.572,9.528)rectangle++(0.04,0.04)
  (1.715,9.632)rectangle++(0.04,0.04) (1.841,9.715)rectangle++(0.04,0.04) (1.951,9.78)rectangle++(0.04,0.04) (2.048,9.831)rectangle++(0.04,0.04) (2.131,9.871)rectangle++(0.04,0.04)
  (2.204,9.901)rectangle++(0.04,0.04) (2.268,9.924)rectangle++(0.04,0.04) (2.322,9.941)rectangle++(0.04,0.04) (2.369,9.953)rectangle++(0.04,0.04) (2.41,9.962)rectangle++(0.04,0.04)
  (2.445,9.967)rectangle++(0.04,0.04) (2.474,9.97)rectangle++(0.04,0.04) (2.5,9.972)rectangle++(0.04,0.04) (2.522,9.973)rectangle++(0.04,0.04) (2.54,9.972)rectangle++(0.04,0.04)
  (2.556,9.971)rectangle++(0.04,0.04) (-2.102,4.37)rectangle++(0.04,0.04) (-1.614,5.395)rectangle++(0.04,0.04) (-1.747,2.405)rectangle++(0.04,0.04) (-1.391,3.71)rectangle++(0.04,0.04)
  (-1.036,4.805)rectangle++(0.04,0.04) (-0.691,5.722)rectangle++(0.04,0.04) (-0.36,6.488)rectangle++(0.04,0.04) (-0.049,7.125)rectangle++(0.04,0.04) (-1.884,3.145)rectangle++(0.04,0.04)
  (-1.477,4.345)rectangle++(0.04,0.04) (-1.084,5.349)rectangle++(0.04,0.04) (-0.71,6.186)rectangle++(0.04,0.04) (-0.358,6.883)rectangle++(0.04,0.04) (-0.031,7.461)rectangle++(0.04,0.04)
  (0.27,7.939)rectangle++(0.04,0.04) (0.545,8.334)rectangle++(0.04,0.04) (0.794,8.658)rectangle++(0.04,0.04) (1.019,8.924)rectangle++(0.04,0.04) (1.221,9.141)rectangle++(0.04,0.04)
  (1.401,9.317)rectangle++(0.04,0.04) (1.562,9.459)rectangle++(0.04,0.04) (1.703,9.574)rectangle++(0.04,0.04) (1.829,9.666)rectangle++(0.04,0.04) (-2.169,4.136)rectangle++(0.04,0.04)
  (-1.68,5.199)rectangle++(0.04,0.04) (-1.222,6.082)rectangle++(0.04,0.04) (-0.798,6.815)rectangle++(0.04,0.04) (-0.408,7.421)rectangle++(0.04,0.04) (-2.032,3.128)rectangle++(0.04,0.04)
  (-1.604,4.336)rectangle++(0.04,0.04) (-1.193,5.346)rectangle++(0.04,0.04) (-0.802,6.188)rectangle++(0.04,0.04) (-0.436,6.888)rectangle++(0.04,0.04) (-0.098,7.469)rectangle++(0.04,0.04)
  (0.214,7.948)rectangle++(0.04,0.04) (0.497,8.344)rectangle++(0.04,0.04) (0.754,8.669)rectangle++(0.04,0.04) (0.986,8.934)rectangle++(0.04,0.04) (-0.02,1.413)rectangle++(0.04,0.04)
  (0.037,2.798)rectangle++(0.04,0.04) (0.141,3.973)rectangle++(0.04,0.04) (0.277,4.967)rectangle++(0.04,0.04) (0.432,5.807)rectangle++(0.04,0.04) (0.597,6.515)rectangle++(0.04,0.04)
  (0.766,7.11)rectangle++(0.04,0.04) (0.933,7.609)rectangle++(0.04,0.04) (1.096,8.027)rectangle++(0.04,0.04) (-1.889,3.607)rectangle++(0.04,0.04) (-1.464,4.738)rectangle++(0.04,0.04)
  (1.529,1.186)rectangle++(0.04,0.04) (1.345,2.544)rectangle++(0.04,0.04) (-0.414,2.458)rectangle++(0.04,0.04) (-0.255,3.702)rectangle++(0.04,0.04) (-0.071,4.753)rectangle++(0.04,0.04)
  (0.127,5.639)rectangle++(0.04,0.04) (1.542,1.817)rectangle++(0.04,0.04) (1.382,3.079)rectangle++(0.04,0.04) (1.295,4.158)rectangle++(0.04,0.04) (1.265,5.078)rectangle++(0.04,0.04)
  (1.276,5.862)rectangle++(0.04,0.04) (1.317,6.528)rectangle++(0.04,0.04) (1.378,7.092)rectangle++(0.04,0.04) (1.762,2.491)rectangle++(0.04,0.04) (1.595,3.643)rectangle++(0.04,0.04)
  (1.499,4.629)rectangle++(0.04,0.04) (-0.925,2.952)rectangle++(0.04,0.04) (-0.67,4.142)rectangle++(0.04,0.04) (-0.406,5.144)rectangle++(0.04,0.04) (-0.142,5.985)rectangle++(0.04,0.04)
  (0.117,6.689)rectangle++(0.04,0.04) (-1.737,3.087)rectangle++(0.04,0.04) (-1.355,4.29)rectangle++(0.04,0.04) (-0.982,5.297)rectangle++(0.04,0.04) (1.613,1.446)rectangle++(0.04,0.04)
  (1.427,2.761)rectangle++(0.04,0.04) (1.321,3.886)rectangle++(0.04,0.04) (1.276,4.846)rectangle++(0.04,0.04) (1.276,5.664)rectangle++(0.04,0.04) (1.309,6.36)rectangle++(0.04,0.04)
  (1.365,6.95)rectangle++(0.04,0.04) (1.436,7.449)rectangle++(0.04,0.04) (1.516,7.87)rectangle++(0.04,0.04) (1.602,8.225)rectangle++(0.04,0.04) (1.688,8.524)rectangle++(0.04,0.04)
  (1.774,8.774)rectangle++(0.04,0.04) (2.173,2.997)rectangle++(0.04,0.04) (1.965,4.057)rectangle++(0.04,0.04) (1.83,4.966)rectangle++(0.04,0.04) (1.752,5.744)rectangle++(0.04,0.04)
  (1.717,6.408)rectangle++(0.04,0.04) (1.519,2.414)rectangle++(0.04,0.04) (1.386,3.588)rectangle++(0.04,0.04) (-0.677,2.697)rectangle++(0.04,0.04) (-0.47,3.916)rectangle++(0.04,0.04)
  (-0.245,4.943)rectangle++(0.04,0.04) (-0.012,5.808)rectangle++(0.04,0.04) (0.219,6.533)rectangle++(0.04,0.04) (0.446,7.141)rectangle++(0.04,0.04) (0.662,7.648)rectangle++(0.04,0.04)
  (0.867,8.071)rectangle++(0.04,0.04) (1.057,8.421)rectangle++(0.04,0.04) (1.233,8.712)rectangle++(0.04,0.04) (1.395,8.952)rectangle++(0.04,0.04) (1.541,9.15)rectangle++(0.04,0.04)
  (1.674,9.312)rectangle++(0.04,0.04) (1.793,9.444)rectangle++(0.04,0.04) (1.9,9.552)rectangle++(0.04,0.04) (1.995,9.64)rectangle++(0.04,0.04) (2.079,9.71)rectangle++(0.04,0.04)
  (2.153,9.767)rectangle++(0.04,0.04) (2.219,9.812)rectangle++(0.04,0.04) (2.276,9.847)rectangle++(0.04,0.04) (2.398,3.385)rectangle++(0.04,0.04) (-0.422,2.885)rectangle++(0.04,0.04)
  (-0.02,0.445)rectangle++(0.04,0.04) (-0.02,0.054)rectangle++(0.04,0.04) (-0.017,1.643)rectangle++(0.04,0.04) (0.049,2.994)rectangle++(0.04,0.04) (0.159,4.139)rectangle++(0.04,0.04)
  (0.299,5.108)rectangle++(0.04,0.04) (0.456,5.926)rectangle++(0.04,0.04) (0.622,6.615)rectangle++(0.04,0.04) (0.791,7.194)rectangle++(0.04,0.04) (0.958,7.679)rectangle++(0.04,0.04)
  (1.12,8.085)rectangle++(0.04,0.04) (1.273,8.424)rectangle++(0.04,0.04) (1.417,8.706)rectangle++(0.04,0.04) (1.55,8.939)rectangle++(0.04,0.04) (1.673,9.133)rectangle++(0.04,0.04)
  (1.785,9.292)rectangle++(0.04,0.04) (1.887,9.423)rectangle++(0.04,0.04) (1.979,9.53)rectangle++(0.04,0.04) (2.061,9.618)rectangle++(0.04,0.04) (2.134,9.689)rectangle++(0.04,0.04)
  (-2.113,4.211)rectangle++(0.04,0.04) (-1.63,5.26)rectangle++(0.04,0.04) (-1.177,6.133)rectangle++(0.04,0.04) (-0.758,6.856)rectangle++(0.04,0.04) (-0.372,7.454)rectangle++(0.04,0.04)
  (-0.02,7.947)rectangle++(0.04,0.04) (0.298,8.352)rectangle++(0.04,0.04) (0.586,8.684)rectangle++(0.04,0.04) (0.843,8.954)rectangle++(0.04,0.04) (1.072,9.173)rectangle++(0.04,0.04)
  (-2.192,3.854)rectangle++(0.04,0.04) (-1.711,4.96)rectangle++(0.04,0.04) (-1.653,2.287)rectangle++(0.04,0.04) (-1.316,3.606)rectangle++(0.04,0.04) (-0.976,4.714)rectangle++(0.04,0.04)
  (-0.644,5.642)rectangle++(0.04,0.04) (-0.324,6.418)rectangle++(0.04,0.04) (-0.02,7.064)rectangle++(0.04,0.04) (0.263,7.602)rectangle++(0.04,0.04) (0.525,8.047)rectangle++(0.04,0.04)
  (0.766,8.415)rectangle++(0.04,0.04) (2.224,2.649)rectangle++(0.04,0.04) (0.391,1.644)rectangle++(0.04,0.04) (0.396,2.978)rectangle++(0.04,0.04) (0.453,4.112)rectangle++(0.04,0.04)
  (0.547,5.073)rectangle++(0.04,0.04) (0.666,5.886)rectangle++(0.04,0.04) (0.799,6.573)rectangle++(0.04,0.04) (0.94,7.151)rectangle++(0.04,0.04) (1.083,7.637)rectangle++(0.04,0.04)
  (1.224,8.044)rectangle++(0.04,0.04) (1.36,8.385)rectangle++(0.04,0.04) (1.489,8.669)rectangle++(0.04,0.04) (2.187,2.898)rectangle++(0.04,0.04) (1.972,3.972)rectangle++(0.04,0.04)
  (-0.659,2.916)rectangle++(0.04,0.04) (0.898,0.958)rectangle++(0.04,0.04) (-0.091,2.006)rectangle++(0.04,0.04) (0.001,3.305)rectangle++(0.04,0.04) (0.131,4.406)rectangle++(0.04,0.04)
  (0.285,5.336)rectangle++(0.04,0.04) (0.454,6.12)rectangle++(0.04,0.04) (1.628,2.017)rectangle++(0.04,0.04) (1.462,3.245)rectangle++(0.04,0.04) (1.371,4.296)rectangle++(0.04,0.04)
  (1.335,5.193)rectangle++(0.04,0.04) (1.34,5.957)rectangle++(0.04,0.04) (1.375,6.606)rectangle++(0.04,0.04) (1.431,7.156)rectangle++(0.04,0.04) (-1.596,3.493)rectangle++(0.04,0.04)
  (-1.219,4.629)rectangle++(0.04,0.04) (-1.468,2.327)rectangle++(0.04,0.04) (-1.157,3.633)rectangle++(0.04,0.04) (-0.841,4.73)rectangle++(0.04,0.04) (-0.527,5.651)rectangle++(0.04,0.04)
  (-0.225,6.42)rectangle++(0.04,0.04) (0.064,7.063)rectangle++(0.04,0.04) (0.335,7.597)rectangle++(0.04,0.04) (0.586,8.04)rectangle++(0.04,0.04) (0.818,8.407)rectangle++(0.04,0.04)
  (1.029,8.709)rectangle++(0.04,0.04) (1.221,8.958)rectangle++(0.04,0.04) (2.308,2.897)rectangle++(0.04,0.04) (2.075,3.967)rectangle++(0.04,0.04) (1.92,4.885)rectangle++(0.04,0.04)
  (1.826,5.671)rectangle++(0.04,0.04) (1.776,6.344)rectangle++(0.04,0.04) (1.761,6.917)rectangle++(0.04,0.04) (-1.467,3.516)rectangle++(0.04,0.04) (-1.109,4.643)rectangle++(0.04,0.04)
  (-0.759,5.588)rectangle++(0.04,0.04) (-0.424,6.376)rectangle++(0.04,0.04) (-0.107,7.033)rectangle++(0.04,0.04) (0.188,7.578)rectangle++(0.04,0.04) (0.461,8.03)rectangle++(0.04,0.04)
  (0.71,8.403)rectangle++(0.04,0.04) (2.229,2.632)rectangle++(0.04,0.04) (1.998,3.744)rectangle++(0.04,0.04) (1.846,4.699)rectangle++(0.04,0.04) (1.755,5.516)rectangle++(0.04,0.04)
  (1.71,6.215)rectangle++(0.04,0.04) (1.7,6.81)rectangle++(0.04,0.04) (1.715,7.317)rectangle++(0.04,0.04) (1.748,7.747)rectangle++(0.04,0.04) (1.794,8.111)rectangle++(0.04,0.04)
  (1.847,8.419)rectangle++(0.04,0.04) (1.904,8.679)rectangle++(0.04,0.04) (1.964,8.897)rectangle++(0.04,0.04) (2.023,9.08)rectangle++(0.04,0.04) (2.08,9.233)rectangle++(0.04,0.04)
  (2.135,9.361)rectangle++(0.04,0.04) (2.187,9.468)rectangle++(0.04,0.04) (2.236,9.556)rectangle++(0.04,0.04) (-2.059,4.206)rectangle++(0.04,0.04) (-1.584,5.253)rectangle++(0.04,0.04)
  (-1.138,6.125)rectangle++(0.04,0.04) (-0.725,6.848)rectangle++(0.04,0.04) (-0.344,7.446)rectangle++(0.04,0.04) (2.119,2.127)rectangle++(0.04,0.04) (1.884,3.32)rectangle++(0.04,0.04)
  (1.732,4.343)rectangle++(0.04,0.04) (1.644,5.218)rectangle++(0.04,0.04) (1.604,5.966)rectangle++(0.04,0.04) (1.6,6.603)rectangle++(0.04,0.04) (1.622,7.145)rectangle++(0.04,0.04)
  (1.662,7.604)rectangle++(0.04,0.04) (1.715,7.993)rectangle++(0.04,0.04) (1.775,8.322)rectangle++(0.04,0.04) (1.839,8.599)rectangle++(0.04,0.04) (1.905,8.832)rectangle++(0.04,0.04)
  (1.971,9.027)rectangle++(0.04,0.04) (2.034,9.19)rectangle++(0.04,0.04) (2.094,9.327)rectangle++(0.04,0.04) (2.151,9.44)rectangle++(0.04,0.04) (2.204,9.534)rectangle++(0.04,0.04)
  (2.252,9.612)rectangle++(0.04,0.04) (2.297,9.676)rectangle++(0.04,0.04) (2.337,9.729)rectangle++(0.04,0.04) (2.373,9.773)rectangle++(0.04,0.04) (2.406,9.808)rectangle++(0.04,0.04)
  (2.435,9.837)rectangle++(0.04,0.04) (2.461,9.86)rectangle++(0.04,0.04) (2.484,9.879)rectangle++(0.04,0.04) (2.505,9.894)rectangle++(0.04,0.04) (-2.093,4.342)rectangle++(0.04,0.04)
  (1.512,0.928)rectangle++(0.04,0.04) (1.32,2.324)rectangle++(0.04,0.04) (1.213,3.519)rectangle++(0.04,0.04) (1.17,4.539)rectangle++(0.04,0.04) (1.174,5.408)rectangle++(0.04,0.04)
  (1.212,6.146)rectangle++(0.04,0.04) (1.273,6.772)rectangle++(0.04,0.04) (1.351,7.301)rectangle++(0.04,0.04) (1.438,7.748)rectangle++(0.04,0.04) (1.53,8.125)rectangle++(0.04,0.04)
  (1.624,8.441)rectangle++(0.04,0.04) (1.715,8.706)rectangle++(0.04,0.04) (1.804,8.928)rectangle++(0.04,0.04) (1.888,9.113)rectangle++(0.04,0.04) (1.967,9.266)rectangle++(0.04,0.04)
  (2.041,9.394)rectangle++(0.04,0.04) (2.108,9.499)rectangle++(0.04,0.04) (2.17,9.586)rectangle++(0.04,0.04) (2.341,3.295)rectangle++(0.04,0.04) (2.12,4.303)rectangle++(0.04,0.04)
  (1.972,5.169)rectangle++(0.04,0.04) (1.88,5.911)rectangle++(0.04,0.04) (1.833,6.545)rectangle++(0.04,0.04) (1.817,7.086)rectangle++(0.04,0.04) (1.694,2.603)rectangle++(0.04,0.04)
  (0.457,1.495)rectangle++(0.04,0.04) (0.446,2.849)rectangle++(0.04,0.04) (0.491,4)rectangle++(0.04,0.04) (0.575,4.976)rectangle++(0.04,0.04) (1.29,1.774)rectangle++(0.04,0.04)
  (1.165,3.052)rectangle++(0.04,0.04) (1.11,4.144)rectangle++(0.04,0.04) (1.107,5.074)rectangle++(0.04,0.04) (1.142,5.865)rectangle++(0.04,0.04) (-1.318,3.142)rectangle++(0.04,0.04)
  (-0.997,4.32)rectangle++(0.04,0.04) (-0.677,5.308)rectangle++(0.04,0.04) (-0.365,6.135)rectangle++(0.04,0.04) (-0.067,6.825)rectangle++(0.04,0.04) (0.214,7.4)rectangle++(0.04,0.04)
  (-1.903,3.266)rectangle++(0.04,0.04) (-1.489,4.449)rectangle++(0.04,0.04) (-1.09,5.437)rectangle++(0.04,0.04) (-1.653,2.535)rectangle++(0.04,0.04) (-1.011,1.766)rectangle++(0.04,0.04)
  (-0.791,3.138)rectangle++(0.04,0.04) (-0.549,4.295)rectangle++(0.04,0.04) (1.268,1.318)rectangle++(0.04,0.04) (1.128,2.666)rectangle++(0.04,0.04) (1.063,3.817)rectangle++(0.04,0.04)
  (-0.801,2.673)rectangle++(0.04,0.04) (-0.576,3.901)rectangle++(0.04,0.04) (-0.336,4.935)rectangle++(0.04,0.04) (-0.09,5.804)rectangle++(0.04,0.04) (0.153,6.533)rectangle++(0.04,0.04)
  (0.389,7.143)rectangle++(0.04,0.04) (0.614,7.653)rectangle++(0.04,0.04) (0.826,8.076)rectangle++(0.04,0.04) (1.023,8.428)rectangle++(0.04,0.04) (1.205,8.719)rectangle++(0.04,0.04)
  (1.37,8.959)rectangle++(0.04,0.04) (1.521,9.157)rectangle++(0.04,0.04) (2.318,3.023)rectangle++(0.04,0.04) (2.089,4.073)rectangle++(0.04,0.04) (1.937,4.975)rectangle++(0.04,0.04)
  (1.843,5.747)rectangle++(0.04,0.04) (-1.147,3.277)rectangle++(0.04,0.04) (-0.846,4.428)rectangle++(0.04,0.04) (-0.544,5.394)rectangle++(0.04,0.04) (-0.249,6.203)rectangle++(0.04,0.04)
  (0.034,6.878)rectangle++(0.04,0.04) (0.302,7.441)rectangle++(0.04,0.04) (0.552,7.909)rectangle++(0.04,0.04) (0.784,8.297)rectangle++(0.04,0.04) (-2.022,3.595)rectangle++(0.04,0.04)
  (-1.577,4.732)rectangle++(0.04,0.04) (-1.567,2.267)rectangle++(0.04,0.04) (-0.924,1.727)rectangle++(0.04,0.04) (-0.719,3.101)rectangle++(0.04,0.04) (-0.489,4.261)rectangle++(0.04,0.04)
  (-0.247,5.238)rectangle++(0.04,0.04) (-0.003,6.058)rectangle++(0.04,0.04) (0.238,6.746)rectangle++(0.04,0.04) (0.47,7.321)rectangle++(0.04,0.04) (0.69,7.8)rectangle++(0.04,0.04)
  (0.896,8.199)rectangle++(0.04,0.04) (1.087,8.529)rectangle++(0.04,0.04) (1.263,8.802)rectangle++(0.04,0.04) (1.424,9.028)rectangle++(0.04,0.04) (1.569,9.213)rectangle++(0.04,0.04)
  (2.327,3.049)rectangle++(0.04,0.04) (2.098,4.095)rectangle++(0.04,0.04) (1.945,4.993)rectangle++(0.04,0.04) (1.85,5.762)rectangle++(0.04,0.04) (1.801,6.42)rectangle++(0.04,0.04)
  (1.786,6.981)rectangle++(0.04,0.04) (1.795,7.459)rectangle++(0.04,0.04) (-1.602,3.643)rectangle++(0.04,0.04) (-1.218,4.757)rectangle++(0.04,0.04) (-0.847,5.688)rectangle++(0.04,0.04)
  (-0.495,6.465)rectangle++(0.04,0.04) (-0.164,7.111)rectangle++(0.04,0.04) (0.143,7.647)rectangle++(0.04,0.04) (0.425,8.091)rectangle++(0.04,0.04) (0.683,8.456)rectangle++(0.04,0.04)
  (0.916,8.757)rectangle++(0.04,0.04) (1.127,9.003)rectangle++(0.04,0.04) (2.334,2.884)rectangle++(0.04,0.04) (2.097,3.954)rectangle++(0.04,0.04) (1.939,4.873)rectangle++(0.04,0.04)
  (1.841,5.661)rectangle++(0.04,0.04) (1.292,2.267)rectangle++(0.04,0.04) (1.186,3.472)rectangle++(0.04,0.04) (1.145,4.5)rectangle++(0.04,0.04) (-0.962,2.842)rectangle++(0.04,0.04)
  (-0.953,1.993)rectangle++(0.04,0.04) (-0.732,3.328)rectangle++(0.04,0.04) (-0.491,4.455)rectangle++(0.04,0.04) (-0.242,5.402)rectangle++(0.04,0.04) (0.008,6.198)rectangle++(0.04,0.04)
  (0.253,6.864)rectangle++(0.04,0.04) (0.487,7.42)rectangle++(0.04,0.04) (0.709,7.884)rectangle++(0.04,0.04) (0.916,8.269)rectangle++(0.04,0.04) (1.107,8.589)rectangle++(0.04,0.04)
  (1.282,8.852)rectangle++(0.04,0.04) (1.442,9.069)rectangle++(0.04,0.04) (1.586,9.247)rectangle++(0.04,0.04) (1.716,9.393)rectangle++(0.04,0.04) (1.832,9.512)rectangle++(0.04,0.04)
  (1.935,9.608)rectangle++(0.04,0.04) (2.027,9.685)rectangle++(0.04,0.04) (2.108,9.748)rectangle++(0.04,0.04) (2.18,9.797)rectangle++(0.04,0.04) (2.243,9.837)rectangle++(0.04,0.04)
  (2.297,9.868)rectangle++(0.04,0.04) (2.345,9.892)rectangle++(0.04,0.04) (2.387,9.911)rectangle++(0.04,0.04) (2.423,9.925)rectangle++(0.04,0.04) (2.454,9.935)rectangle++(0.04,0.04)
  (2.482,9.943)rectangle++(0.04,0.04) (2.505,9.948)rectangle++(0.04,0.04) (2.525,9.952)rectangle++(0.04,0.04) (2.542,9.955)rectangle++(0.04,0.04) (2.557,9.956)rectangle++(0.04,0.04)
  (2.569,9.956)rectangle++(0.04,0.04) (2.58,9.956)rectangle++(0.04,0.04) (2.589,9.956)rectangle++(0.04,0.04) (2.597,9.955)rectangle++(0.04,0.04) (2.603,9.954)rectangle++(0.04,0.04)
  (2.609,9.953)rectangle++(0.04,0.04) (2.613,9.952)rectangle++(0.04,0.04) (2.617,9.951)rectangle++(0.04,0.04) (2.376,3.499)rectangle++(0.04,0.04) (2.158,4.475)rectangle++(0.04,0.04)
  (2.011,5.314)rectangle++(0.04,0.04) (1.919,6.032)rectangle++(0.04,0.04) (1.871,6.647)rectangle++(0.04,0.04) (1.854,7.171)rectangle++(0.04,0.04) (1.86,7.618)rectangle++(0.04,0.04)
  (1.884,7.997)rectangle++(0.04,0.04) (1.919,8.318)rectangle++(0.04,0.04) (1.962,8.59)rectangle++(0.04,0.04) (2.009,8.819)rectangle++(0.04,0.04) (2.058,9.012)rectangle++(0.04,0.04)
  (2.108,9.174)rectangle++(0.04,0.04) (2.156,9.31)rectangle++(0.04,0.04) (2.203,9.423)rectangle++(0.04,0.04) (2.247,9.518)rectangle++(0.04,0.04) (2.289,9.597)rectangle++(0.04,0.04)
  (2.326,3.328)rectangle++(0.04,0.04) (2.108,4.332)rectangle++(0.04,0.04) (1.963,5.194)rectangle++(0.04,0.04) (1.874,5.933)rectangle++(0.04,0.04) (1.828,6.564)rectangle++(0.04,0.04)
  (1.814,7.103)rectangle++(0.04,0.04) (1.824,7.561)rectangle++(0.04,0.04) (1.851,7.95)rectangle++(0.04,0.04) (1.889,8.28)rectangle++(0.04,0.04) (1.935,8.558)rectangle++(0.04,0.04)
  (1.984,8.793)rectangle++(0.04,0.04) (2.036,8.991)rectangle++(0.04,0.04) (2.088,9.157)rectangle++(0.04,0.04) (2.139,9.296)rectangle++(0.04,0.04) (2.188,9.413)rectangle++(0.04,0.04)
  (2.234,9.509)rectangle++(0.04,0.04) (2.277,9.59)rectangle++(0.04,0.04) (2.317,9.656)rectangle++(0.04,0.04) (2.353,9.711)rectangle++(0.04,0.04) (2.387,9.757)rectangle++(0.04,0.04)
  (2.417,9.794)rectangle++(0.04,0.04) (2.444,9.824)rectangle++(0.04,0.04) (2.468,9.849)rectangle++(0.04,0.04) (2.49,9.869)rectangle++(0.04,0.04) (2.509,9.886)rectangle++(0.04,0.04)
  (-2.09,4.341)rectangle++(0.04,0.04) (-1.605,5.369)rectangle++(0.04,0.04) (-1.152,6.224)rectangle++(0.04,0.04) (-0.732,6.933)rectangle++(0.04,0.04) (-0.347,7.519)rectangle++(0.04,0.04)
  (2.14,2.144)rectangle++(0.04,0.04) (1.902,3.333)rectangle++(0.04,0.04) (1.748,4.353)rectangle++(0.04,0.04) (-0.803,2.949)rectangle++(0.04,0.04) (0.929,0.929)rectangle++(0.04,0.04)
  (0.824,2.349)rectangle++(0.04,0.04) (0.792,3.559)rectangle++(0.04,0.04) (-0.788,2.554)rectangle++(0.04,0.04) (-0.57,3.799)rectangle++(0.04,0.04) (-0.335,4.848)rectangle++(0.04,0.04)
  (-0.093,5.73)rectangle++(0.04,0.04) (0.148,6.471)rectangle++(0.04,0.04) (0.383,7.09)rectangle++(0.04,0.04) (0.607,7.608)rectangle++(0.04,0.04) (0.818,8.039)rectangle++(0.04,0.04)
  (1.014,8.396)rectangle++(0.04,0.04) (1.196,8.692)rectangle++(0.04,0.04) (1.362,8.937)rectangle++(0.04,0.04) (1.513,9.138)rectangle++(0.04,0.04) (1.649,9.303)rectangle++(0.04,0.04)
  (1.772,9.438)rectangle++(0.04,0.04) (1.881,9.547)rectangle++(0.04,0.04) (1.979,9.636)rectangle++(0.04,0.04) (2.065,9.708)rectangle++(0.04,0.04) (-2.132,4.2)rectangle++(0.04,0.04)
  (-1.647,5.251)rectangle++(0.04,0.04) (1.7,1.262)rectangle++(0.04,0.04) (1.493,2.601)rectangle++(0.04,0.04) (1.371,3.747)rectangle++(0.04,0.04) (1.313,4.727)rectangle++(0.04,0.04)
  (1.303,5.561)rectangle++(0.04,0.04) (1.328,6.271)rectangle++(0.04,0.04) (1.377,6.874)rectangle++(0.04,0.04) (1.443,7.384)rectangle++(0.04,0.04) (1.52,7.815)rectangle++(0.04,0.04)
  (1.602,8.178)rectangle++(0.04,0.04) (1.687,8.483)rectangle++(0.04,0.04) (1.771,8.739)rectangle++(0.04,0.04) (2.164,2.988)rectangle++(0.04,0.04) (1.957,4.049)rectangle++(0.04,0.04)
  (1.823,4.96)rectangle++(0.04,0.04) (1.746,5.739)rectangle++(0.04,0.04) (1.711,6.405)rectangle++(0.04,0.04) (1.709,6.972)rectangle++(0.04,0.04) (1.729,7.454)rectangle++(0.04,0.04)
  (1.766,7.863)rectangle++(0.04,0.04) (1.813,8.209)rectangle++(0.04,0.04) (1.867,8.501)rectangle++(0.04,0.04) (1.925,8.748)rectangle++(0.04,0.04) (1.984,8.955)rectangle++(0.04,0.04)
  (2.042,9.128)rectangle++(0.04,0.04) (2.099,9.274)rectangle++(0.04,0.04) (2.153,9.395)rectangle++(0.04,0.04) (2.203,9.496)rectangle++(0.04,0.04) (2.251,9.579)rectangle++(0.04,0.04)
  (2.294,9.649)rectangle++(0.04,0.04) (2.334,9.706)rectangle++(0.04,0.04) (2.37,9.753)rectangle++(0.04,0.04) (2.402,9.791)rectangle++(0.04,0.04) (2.431,9.823)rectangle++(0.04,0.04)
  (2.457,9.848)rectangle++(0.04,0.04) (2.48,9.869)rectangle++(0.04,0.04) (2.501,9.886)rectangle++(0.04,0.04) (2.519,9.899)rectangle++(0.04,0.04) (2.376,3.461)rectangle++(0.04,0.04)
  (0.595,1.878)rectangle++(0.04,0.04) (0.579,3.169)rectangle++(0.04,0.04) (0.617,4.267)rectangle++(0.04,0.04) (0.693,5.198)rectangle++(0.04,0.04) (0.794,5.987)rectangle++(0.04,0.04)
  (0.913,6.653)rectangle++(0.04,0.04) (1.04,7.215)rectangle++(0.04,0.04) (1.17,7.687)rectangle++(0.04,0.04) (1.3,8.084)rectangle++(0.04,0.04) (1.426,8.415)rectangle++(0.04,0.04)
  (1.547,8.692)rectangle++(0.04,0.04) (1.66,8.923)rectangle++(0.04,0.04) (1.766,9.114)rectangle++(0.04,0.04) (2.27,3.076)rectangle++(0.04,0.04) (-0.367,2.788)rectangle++(0.04,0.04)
  (-0.819,2.118)rectangle++(0.04,0.04) (-0.614,3.429)rectangle++(0.04,0.04) (-0.387,4.536)rectangle++(0.04,0.04) (-0.15,5.467)rectangle++(0.04,0.04) (0.089,6.249)rectangle++(0.04,0.04)
  (0.324,6.904)rectangle++(0.04,0.04) (0.549,7.452)rectangle++(0.04,0.04) (0.763,7.908)rectangle++(0.04,0.04) (0.962,8.288)rectangle++(0.04,0.04) (1.147,8.602)rectangle++(0.04,0.04)
  (1.317,8.862)rectangle++(0.04,0.04) (1.472,9.076)rectangle++(0.04,0.04) (1.612,9.252)rectangle++(0.04,0.04) (1.738,9.396)rectangle++(0.04,0.04) (-0.02,1.487)rectangle++(0.04,0.04)
  (0.04,2.861)rectangle++(0.04,0.04) (0.146,4.026)rectangle++(0.04,0.04) (0.283,5.013)rectangle++(0.04,0.04) (0.439,5.846)rectangle++(0.04,0.04) (0.605,6.547)rectangle++(0.04,0.04)
  (1.725,2.159)rectangle++(0.04,0.04) (0.328,1.397)rectangle++(0.04,0.04) (0.333,2.77)rectangle++(0.04,0.04) (0.391,3.938)rectangle++(0.04,0.04) (0.488,4.927)rectangle++(0.04,0.04)
  (0.61,5.765)rectangle++(0.04,0.04) (0.747,6.472)rectangle++(0.04,0.04) (-1.555,3.185)rectangle++(0.04,0.04) (-1.196,4.365)rectangle++(0.04,0.04) (-0.844,5.355)rectangle++(0.04,0.04)
  (-0.506,6.181)rectangle++(0.04,0.04) (-0.185,6.871)rectangle++(0.04,0.04) (0.116,7.444)rectangle++(0.04,0.04) (0.394,7.919)rectangle++(0.04,0.04) (0.649,8.311)rectangle++(0.04,0.04)
  (0.882,8.635)rectangle++(0.04,0.04) (1.093,8.901)rectangle++(0.04,0.04) (1.283,9.118)rectangle++(0.04,0.04) (1.453,9.295)rectangle++(0.04,0.04) (2.367,3.039)rectangle++(0.04,0.04)
  (2.132,4.084)rectangle++(0.04,0.04) (1.973,4.983)rectangle++(0.04,0.04) (1.874,5.752)rectangle++(0.04,0.04) (1.821,6.411)rectangle++(0.04,0.04) (1.802,6.973)rectangle++(0.04,0.04)
  (1.808,7.451)rectangle++(0.04,0.04) (1.833,7.857)rectangle++(0.04,0.04) (1.87,8.201)rectangle++(0.04,0.04) (1.915,8.493)rectangle++(0.04,0.04) (1.966,8.738)rectangle++(0.04,0.04)
  (2.018,8.945)rectangle++(0.04,0.04) (2.071,9.119)rectangle++(0.04,0.04) (2.225,3.157)rectangle++(0.04,0.04) (2.016,4.191)rectangle++(0.04,0.04) (-0.708,2.975)rectangle++(0.04,0.04)
  (-0.485,4.153)rectangle++(0.04,0.04) (-0.248,5.145)rectangle++(0.04,0.04) (-0.007,5.98)rectangle++(0.04,0.04) (0.231,6.679)rectangle++(0.04,0.04) (0.461,7.264)rectangle++(0.04,0.04)
  (0.68,7.752)rectangle++(0.04,0.04) (2.051,2.467)rectangle++(0.04,0.04) (1.84,3.612)rectangle++(0.04,0.04) (-0.592,2.807)rectangle++(0.04,0.04) (-0.393,4.006)rectangle++(0.04,0.04)
  (-0.176,5.017)rectangle++(0.04,0.04) (0.049,5.867)rectangle++(0.04,0.04) (0.274,6.582)rectangle++(0.04,0.04) (0.494,7.18)rectangle++(0.04,0.04) (0.705,7.679)rectangle++(0.04,0.04)
  (0.904,8.095)rectangle++(0.04,0.04) (2.114,2.608)rectangle++(0.04,0.04) (1.899,3.728)rectangle++(0.04,0.04) (1.761,4.689)rectangle++(0.04,0.04) (1.682,5.512)rectangle++(0.04,0.04)
  (1.648,6.214)rectangle++(0.04,0.04) (1.647,6.812)rectangle++(0.04,0.04) (-1.463,3.467)rectangle++(0.04,0.04) (1.173,0.882)rectangle++(0.04,0.04) (1.03,2.299)rectangle++(0.04,0.04)
  (0.965,3.509)rectangle++(0.04,0.04) (0.958,4.54)rectangle++(0.04,0.04) (0.994,5.417)rectangle++(0.04,0.04) (1.059,6.161)rectangle++(0.04,0.04) (1.145,6.791)rectangle++(0.04,0.04)
  (1.242,7.322)rectangle++(0.04,0.04) (1.347,7.771)rectangle++(0.04,0.04) (1.453,8.147)rectangle++(0.04,0.04) (2.046,2.763)rectangle++(0.04,0.04) (1.847,3.863)rectangle++(0.04,0.04)
  (1.723,4.806)rectangle++(0.04,0.04) (1.654,5.612)rectangle++(0.04,0.04) (1.628,6.301)rectangle++(0.04,0.04) (1.634,6.887)rectangle++(0.04,0.04) (-1.485,3.48)rectangle++(0.04,0.04)
  (-1.125,4.613)rectangle++(0.04,0.04) (-0.774,5.563)rectangle++(0.04,0.04) (-0.438,6.355)rectangle++(0.04,0.04) (-0.12,7.016)rectangle++(0.04,0.04) (0.176,7.564)rectangle++(0.04,0.04)
  (0.45,8.019)rectangle++(0.04,0.04) (0.701,8.394)rectangle++(0.04,0.04) (0.93,8.703)rectangle++(0.04,0.04) (1.136,8.957)rectangle++(0.04,0.04) (1.322,9.164)rectangle++(0.04,0.04)
  (1.488,9.333)rectangle++(0.04,0.04) (1.636,9.47)rectangle++(0.04,0.04) (1.767,9.58)rectangle++(0.04,0.04) (1.883,9.668)rectangle++(0.04,0.04) (-2.158,4.149)rectangle++(0.04,0.04)
  (-1.671,5.209)rectangle++(0.04,0.04) (-1.214,6.091)rectangle++(0.04,0.04) (-0.791,6.822)rectangle++(0.04,0.04) (-0.401,7.427)rectangle++(0.04,0.04) (-0.046,7.925)rectangle++(0.04,0.04)
  (0.276,8.334)rectangle++(0.04,0.04) (0.565,8.669)rectangle++(0.04,0.04) (-2.162,3.626)rectangle++(0.04,0.04) (-1.695,4.765)rectangle++(0.04,0.04) (1.571,1.133)rectangle++(0.04,0.04)
  (1.379,2.496)rectangle++(0.04,0.04) (1.269,3.663)rectangle++(0.04,0.04) (-0.02,0.569)rectangle++(0.04,0.04) (0.004,2.081)rectangle++(0.04,0.04) (0.084,3.365)rectangle++(0.04,0.04)
  (0.204,4.453)rectangle++(0.04,0.04) (0.349,5.373)rectangle++(0.04,0.04) (-1.348,2.851)rectangle++(0.04,0.04) (-1.034,4.074)rectangle++(0.04,0.04) (-0.718,5.1)rectangle++(0.04,0.04)
  (-0.409,5.96)rectangle++(0.04,0.04) (-0.111,6.679)rectangle++(0.04,0.04) (0.17,7.278)rectangle++(0.04,0.04) (0.434,7.775)rectangle++(0.04,0.04) (2.095,2.409)rectangle++(0.04,0.04)
  (1.875,3.56)rectangle++(0.04,0.04) (1.734,4.547)rectangle++(0.04,0.04) (1.653,5.392)rectangle++(0.04,0.04) (1.619,6.113)rectangle++(0.04,0.04) (1.618,6.728)rectangle++(0.04,0.04)
  (1.642,7.25)rectangle++(0.04,0.04) (1.684,7.693)rectangle++(0.04,0.04) (1.737,8.068)rectangle++(0.04,0.04) (1.797,8.384)rectangle++(0.04,0.04) (1.86,8.651)rectangle++(0.04,0.04)
  (1.925,8.875)rectangle++(0.04,0.04) (-1.944,3.984)rectangle++(0.04,0.04) (-1.446,2.019)rectangle++(0.04,0.04) (-0.835,1.701)rectangle++(0.04,0.04) (-0.644,3.075)rectangle++(0.04,0.04)
  (-0.427,4.236)rectangle++(0.04,0.04) (-0.195,5.214)rectangle++(0.04,0.04) (0.04,6.036)rectangle++(0.04,0.04) (0.273,6.725)rectangle++(0.04,0.04) (0.499,7.301)rectangle++(0.04,0.04)
  (-0.02,1.151)rectangle++(0.04,0.04) (0.027,2.576)rectangle++(0.04,0.04) (0.124,3.784)rectangle++(0.04,0.04) (0.254,4.808)rectangle++(0.04,0.04) (1.291,1.65)rectangle++(0.04,0.04)
  (1.161,2.947)rectangle++(0.04,0.04) (1.102,4.055)rectangle++(0.04,0.04) (1.097,4.999)rectangle++(0.04,0.04) (1.13,5.801)rectangle++(0.04,0.04) (1.191,6.482)rectangle++(0.04,0.04)
  (1.269,7.058)rectangle++(0.04,0.04) (1.359,7.545)rectangle++(0.04,0.04) (1.455,7.955)rectangle++(0.04,0.04) (-1.799,3.674)rectangle++(0.04,0.04) (1.281,0.844)rectangle++(0.04,0.04)
  (1.12,2.262)rectangle++(0.04,0.04) (1.041,3.474)rectangle++(0.04,0.04) (1.021,4.508)rectangle++(0.04,0.04) (1.046,5.387)rectangle++(0.04,0.04) (1.103,6.133)rectangle++(0.04,0.04)
  (1.18,6.765)rectangle++(0.04,0.04) (1.272,7.3)rectangle++(0.04,0.04) (-1.665,3.487)rectangle++(0.04,0.04) (-1.278,4.627)rectangle++(0.04,0.04) (-0.903,5.58)rectangle++(0.04,0.04)
  (-0.547,6.376)rectangle++(0.04,0.04) (-0.212,7.037)rectangle++(0.04,0.04) (-1.893,3.088)rectangle++(0.04,0.04) (-1.488,4.297)rectangle++(0.04,0.04) (-1.095,5.308)rectangle++(0.04,0.04)
  (-0.721,6.152)rectangle++(0.04,0.04) (1.813,1.719)rectangle++(0.04,0.04) (1.608,2.985)rectangle++(0.04,0.04) (1.484,4.069)rectangle++(0.04,0.04) (1.422,4.996)rectangle++(0.04,0.04)
  (1.168,1.999)rectangle++(0.04,0.04) (1.071,3.248)rectangle++(0.04,0.04) (1.038,4.314)rectangle++(0.04,0.04) (1.052,5.222)rectangle++(0.04,0.04) (1.101,5.993)rectangle++(0.04,0.04)
  (1.174,6.646)rectangle++(0.04,0.04) (1.261,7.198)rectangle++(0.04,0.04) (-1.641,3.463)rectangle++(0.04,0.04) (1.198,0.835)rectangle++(0.04,0.04) (0.001,2.048)rectangle++(0.04,0.04)
  (0.081,3.337)rectangle++(0.04,0.04) (-0.873,2.342)rectangle++(0.04,0.04) (-0.65,3.622)rectangle++(0.04,0.04) (-0.41,4.701)rectangle++(0.04,0.04) (-0.163,5.608)rectangle++(0.04,0.04)
  (0.084,6.37)rectangle++(0.04,0.04) (0.324,7.007)rectangle++(0.04,0.04) (0.553,7.539)rectangle++(0.04,0.04) (0.77,7.982)rectangle++(0.04,0.04) (-1.943,3.522)rectangle++(0.04,0.04)
  (-1.513,4.668)rectangle++(0.04,0.04) (-1.101,5.624)rectangle++(0.04,0.04) (-0.713,6.421)rectangle++(0.04,0.04) (-0.352,7.082)rectangle++(0.04,0.04) (-1.933,3.066)rectangle++(0.04,0.04)
  (-1.523,4.28)rectangle++(0.04,0.04) (-1.438,2.18)rectangle++(0.04,0.04) (-0.02,0.332)rectangle++(0.04,0.04) (-0.006,1.879)rectangle++(0.04,0.04) (0.068,3.194)rectangle++(0.04,0.04)
  (0.183,4.308)rectangle++(0.04,0.04) (-0.02,0.673)rectangle++(0.04,0.04) (0.174,0.586)rectangle++(0.04,0.04) (0.169,2.088)rectangle++(0.04,0.04) (0.225,3.364)rectangle++(0.04,0.04)
  (0.324,4.446)rectangle++(0.04,0.04) (0.451,5.363)rectangle++(0.04,0.04) (0.595,6.137)rectangle++(0.04,0.04) (1.612,2.058)rectangle++(0.04,0.04) (1.45,3.281)rectangle++(0.04,0.04)
  (1.361,4.327)rectangle++(0.04,0.04) (1.328,5.22)rectangle++(0.04,0.04) (1.335,5.98)rectangle++(0.04,0.04) (1.372,6.625)rectangle++(0.04,0.04) (-1.469,3.362)rectangle++(0.04,0.04)
  (-1.117,4.513)rectangle++(0.04,0.04) (-1.418,2.325)rectangle++(0.04,0.04) (-1.114,3.629)rectangle++(0.04,0.04) (-0.804,4.726)rectangle++(0.04,0.04) (-0.497,5.645)rectangle++(0.04,0.04)
  (-0.199,6.414)rectangle++(0.04,0.04) (1.808,1.918)rectangle++(0.04,0.04) (-0.158,2.427)rectangle++(0.04,0.04) (-0.04,3.665)rectangle++(0.04,0.04) (0.111,4.713)rectangle++(0.04,0.04)
  (0.281,5.598)rectangle++(0.04,0.04) (0.46,6.343)rectangle++(0.04,0.04) (0.643,6.97)rectangle++(0.04,0.04) (-1.705,3.27)rectangle++(0.04,0.04) (-1.32,4.444)rectangle++(0.04,0.04)
  (-0.947,5.426)rectangle++(0.04,0.04) (-0.59,6.247)rectangle++(0.04,0.04) (-0.254,6.929)rectangle++(0.04,0.04) (0.059,7.496)rectangle++(0.04,0.04) (0.348,7.966)rectangle++(0.04,0.04)
  (0.612,8.353)rectangle++(0.04,0.04) (0.852,8.672)rectangle++(0.04,0.04) (1.069,8.933)rectangle++(0.04,0.04) (1.264,9.147)rectangle++(0.04,0.04) (-2.147,3.892)rectangle++(0.04,0.04)
  (-1.671,4.99)rectangle++(0.04,0.04) (-1.223,5.905)rectangle++(0.04,0.04) (-0.806,6.664)rectangle++(0.04,0.04) (-0.42,7.293)rectangle++(0.04,0.04) (-0.068,7.812)rectangle++(0.04,0.04)
  (0.253,8.239)rectangle++(0.04,0.04) (0.542,8.589)rectangle++(0.04,0.04) (0.802,8.875)rectangle++(0.04,0.04) (1.035,9.108)rectangle++(0.04,0.04) (1.242,9.297)rectangle++(0.04,0.04)
  (1.425,9.449)rectangle++(0.04,0.04) (1.587,9.571)rectangle++(0.04,0.04) (1.73,9.668)rectangle++(0.04,0.04) (1.855,9.745)rectangle++(0.04,0.04) (1.964,9.805)rectangle++(0.04,0.04)
  (2.059,9.852)rectangle++(0.04,0.04) (2.142,9.888)rectangle++(0.04,0.04) (2.214,9.915)rectangle++(0.04,0.04) (2.277,9.936)rectangle++(0.04,0.04) (2.33,9.95)rectangle++(0.04,0.04)
  (2.377,9.961)rectangle++(0.04,0.04) (2.416,9.968)rectangle++(0.04,0.04) (2.45,9.972)rectangle++(0.04,0.04) (2.48,9.975)rectangle++(0.04,0.04) (2.504,9.975)rectangle++(0.04,0.04)
  (2.526,9.975)rectangle++(0.04,0.04) (2.544,9.974)rectangle++(0.04,0.04) (2.559,9.972)rectangle++(0.04,0.04) (2.572,9.97)rectangle++(0.04,0.04) (2.583,9.968)rectangle++(0.04,0.04)
  (2.592,9.966)rectangle++(0.04,0.04) (2.599,9.963)rectangle++(0.04,0.04) (2.606,9.961)rectangle++(0.04,0.04) (2.611,9.959)rectangle++(0.04,0.04) (2.616,9.957)rectangle++(0.04,0.04)
  (2.619,9.955)rectangle++(0.04,0.04) (2.622,9.953)rectangle++(0.04,0.04) (-2.085,4.382)rectangle++(0.04,0.04) (-1.599,5.404)rectangle++(0.04,0.04) (-1.145,6.254)rectangle++(0.04,0.04)
  (-0.725,6.958)rectangle++(0.04,0.04) (-0.34,7.539)rectangle++(0.04,0.04) (0.01,8.018)rectangle++(0.04,0.04) (0.327,8.411)rectangle++(0.04,0.04) (0.612,8.733)rectangle++(0.04,0.04)
  (0.868,8.994)rectangle++(0.04,0.04) (1.095,9.207)rectangle++(0.04,0.04) (1.297,9.378)rectangle++(0.04,0.04) (1.475,9.516)rectangle++(0.04,0.04) (-0.02,1.506)rectangle++(0.04,0.04)
  (-0.417,1.916)rectangle++(0.04,0.04) (-0.28,3.241)rectangle++(0.04,0.04) (-0.11,4.362)rectangle++(0.04,0.04) (0.079,5.309)rectangle++(0.04,0.04) (0.277,6.105)rectangle++(0.04,0.04)
  (0.477,6.775)rectangle++(0.04,0.04) (0.675,7.336)rectangle++(0.04,0.04) (0.865,7.804)rectangle++(0.04,0.04) (1.045,8.195)rectangle++(0.04,0.04) (1.214,8.521)rectangle++(0.04,0.04)
  (1.37,8.79)rectangle++(0.04,0.04) (1.514,9.013)rectangle++(0.04,0.04) (1.645,9.197)rectangle++(0.04,0.04) (1.764,9.348)rectangle++(0.04,0.04) (1.871,9.471)rectangle++(0.04,0.04)
  (-2.109,4.103)rectangle++(0.04,0.04) (-1.631,5.168)rectangle++(0.04,0.04) (-1.182,6.054)rectangle++(0.04,0.04) (-0.765,6.79)rectangle++(0.04,0.04) (-0.381,7.398)rectangle++(0.04,0.04)
  (-0.03,7.9)rectangle++(0.04,0.04) (0.289,8.312)rectangle++(0.04,0.04) (0.576,8.65)rectangle++(0.04,0.04) (0.833,8.926)rectangle++(0.04,0.04) (1.063,9.15)rectangle++(0.04,0.04)
  (1.267,9.331)rectangle++(0.04,0.04) (1.448,9.477)rectangle++(0.04,0.04) (1.608,9.594)rectangle++(0.04,0.04) (2.428,3.15)rectangle++(0.04,0.04) (2.187,4.177)rectangle++(0.04,0.04);
\fill[nuage]
  (2.024,5.059)rectangle++(0.04,0.04) (-0.932,3.168)rectangle++(0.04,0.04) (-0.667,4.326)rectangle++(0.04,0.04) (-0.397,5.3)rectangle++(0.04,0.04) (-0.127,6.117)rectangle++(0.04,0.04)
  (0.134,6.801)rectangle++(0.04,0.04) (0.384,7.371)rectangle++(0.04,0.04) (-1.861,3.299)rectangle++(0.04,0.04) (-1.452,4.475)rectangle++(0.04,0.04) (-1.057,5.458)rectangle++(0.04,0.04)
  (-1.652,2.547)rectangle++(0.04,0.04) (-1.014,1.769)rectangle++(0.04,0.04) (-0.793,3.141)rectangle++(0.04,0.04) (-0.551,4.297)rectangle++(0.04,0.04) (-0.298,5.271)rectangle++(0.04,0.04)
  (-0.045,6.089)rectangle++(0.04,0.04) (0.203,6.773)rectangle++(0.04,0.04) (1.849,2.108)rectangle++(0.04,0.04) (1.653,3.314)rectangle++(0.04,0.04) (0.663,1.655)rectangle++(0.04,0.04)
  (0.627,2.977)rectangle++(0.04,0.04) (0.65,4.101)rectangle++(0.04,0.04) (0.714,5.056)rectangle++(0.04,0.04) (0.807,5.866)rectangle++(0.04,0.04) (0.919,6.55)rectangle++(0.04,0.04)
  (1.041,7.127)rectangle++(0.04,0.04) (-1.666,3.396)rectangle++(0.04,0.04) (-1.282,4.55)rectangle++(0.04,0.04) (-0.91,5.515)rectangle++(0.04,0.04) (-0.555,6.32)rectangle++(0.04,0.04)
  (-0.221,6.99)rectangle++(0.04,0.04) (0.089,7.547)rectangle++(0.04,0.04) (0.375,8.008)rectangle++(0.04,0.04) (0.637,8.388)rectangle++(0.04,0.04) (0.875,8.7)rectangle++(0.04,0.04)
  (1.09,8.956)rectangle++(0.04,0.04) (1.282,9.166)rectangle++(0.04,0.04) (-2.148,3.9)rectangle++(0.04,0.04) (-1.672,4.997)rectangle++(0.04,0.04) (-1.223,5.911)rectangle++(0.04,0.04)
  (-0.806,6.669)rectangle++(0.04,0.04) (-0.42,7.297)rectangle++(0.04,0.04) (2.089,2.072)rectangle++(0.04,0.04) (1.856,3.274)rectangle++(0.04,0.04) (1.707,4.305)rectangle++(0.04,0.04)
  (1.621,5.187)rectangle++(0.04,0.04) (1.583,5.94)rectangle++(0.04,0.04) (1.581,6.582)rectangle++(0.04,0.04) (1.605,7.128)rectangle++(0.04,0.04) (-1.553,3.526)rectangle++(0.04,0.04)
  (-1.182,4.656)rectangle++(0.04,0.04) (-0.82,5.601)rectangle++(0.04,0.04) (-0.475,6.39)rectangle++(0.04,0.04) (-0.151,7.046)rectangle++(0.04,0.04) (0.152,7.592)rectangle++(0.04,0.04)
  (0.43,8.043)rectangle++(0.04,0.04) (0.685,8.416)rectangle++(0.04,0.04) (0.917,8.722)rectangle++(0.04,0.04) (1.126,8.973)rectangle++(0.04,0.04) (1.314,9.178)rectangle++(0.04,0.04)
  (1.482,9.345)rectangle++(0.04,0.04) (1.631,9.48)rectangle++(0.04,0.04) (1.763,9.589)rectangle++(0.04,0.04) (1.88,9.677)rectangle++(0.04,0.04) (1.983,9.746)rectangle++(0.04,0.04)
  (-2.159,4.189)rectangle++(0.04,0.04) (-1.669,5.243)rectangle++(0.04,0.04) (-1.211,6.12)rectangle++(0.04,0.04) (-0.787,6.847)rectangle++(0.04,0.04) (2.018,1.869)rectangle++(0.04,0.04)
  (1.788,3.104)rectangle++(0.04,0.04) (1.641,4.163)rectangle++(0.04,0.04) (1.56,5.069)rectangle++(0.04,0.04) (1.526,5.842)rectangle++(0.04,0.04) (-1.235,3.225)rectangle++(0.04,0.04)
  (-0.923,4.387)rectangle++(0.04,0.04) (-0.611,5.362)rectangle++(0.04,0.04) (-0.307,6.179)rectangle++(0.04,0.04) (-0.016,6.86)rectangle++(0.04,0.04) (0.258,7.428)rectangle++(0.04,0.04)
  (0.515,7.9)rectangle++(0.04,0.04) (-0.02,1.247)rectangle++(0.04,0.04) (0.031,2.657)rectangle++(0.04,0.04) (0.13,3.853)rectangle++(0.04,0.04) (0.263,4.866)rectangle++(0.04,0.04)
  (0.416,5.722)rectangle++(0.04,0.04) (1.522,1.911)rectangle++(0.04,0.04) (1.368,3.16)rectangle++(0.04,0.04) (1.287,4.227)rectangle++(0.04,0.04) (1.261,5.138)rectangle++(0.04,0.04)
  (1.275,5.913)rectangle++(0.04,0.04) (1.318,6.571)rectangle++(0.04,0.04) (1.381,7.129)rectangle++(0.04,0.04) (1.457,7.601)rectangle++(0.04,0.04) (-1.706,3.596)rectangle++(0.04,0.04)
  (-1.308,4.721)rectangle++(0.04,0.04) (-0.926,5.662)rectangle++(0.04,0.04) (-0.562,6.446)rectangle++(0.04,0.04) (-0.222,7.097)rectangle++(0.04,0.04) (0.093,7.638)rectangle++(0.04,0.04)
  (0.382,8.085)rectangle++(0.04,0.04) (0.646,8.453)rectangle++(0.04,0.04) (0.885,8.755)rectangle++(0.04,0.04) (1.1,9.003)rectangle++(0.04,0.04) (1.293,9.205)rectangle++(0.04,0.04)
  (1.465,9.368)rectangle++(0.04,0.04) (-2.164,3.987)rectangle++(0.04,0.04) (-1.682,5.072)rectangle++(0.04,0.04) (-1.229,5.974)rectangle++(0.04,0.04) (-0.808,6.724)rectangle++(0.04,0.04)
  (-0.42,7.344)rectangle++(0.04,0.04) (2.102,2.083)rectangle++(0.04,0.04) (1.868,3.283)rectangle++(0.04,0.04) (1.717,4.312)rectangle++(0.04,0.04) (1.629,5.193)rectangle++(0.04,0.04)
  (1.591,5.945)rectangle++(0.04,0.04) (1.409,2.27)rectangle++(0.04,0.04) (1.286,3.47)rectangle++(0.04,0.04) (1.23,4.494)rectangle++(0.04,0.04) (1.223,5.367)rectangle++(0.04,0.04)
  (1.252,6.109)rectangle++(0.04,0.04) (1.306,6.739)rectangle++(0.04,0.04) (1.378,7.272)rectangle++(0.04,0.04) (1.46,7.722)rectangle++(0.04,0.04) (1.547,8.102)rectangle++(0.04,0.04)
  (1.637,8.421)rectangle++(0.04,0.04) (1.726,8.688)rectangle++(0.04,0.04) (-1.935,3.897)rectangle++(0.04,0.04) (-1.491,4.986)rectangle++(0.04,0.04) (-1.07,5.894)rectangle++(0.04,0.04)
  (-0.676,6.649)rectangle++(0.04,0.04) (-0.311,7.275)rectangle++(0.04,0.04) (0.025,7.792)rectangle++(0.04,0.04) (0.33,8.219)rectangle++(0.04,0.04) (0.607,8.569)rectangle++(0.04,0.04)
  (0.857,8.856)rectangle++(0.04,0.04) (-2.152,3.734)rectangle++(0.04,0.04) (-1.682,4.856)rectangle++(0.04,0.04) (-1.238,5.791)rectangle++(0.04,0.04) (-0.823,6.568)rectangle++(0.04,0.04)
  (-0.439,7.212)rectangle++(0.04,0.04) (-0.087,7.744)rectangle++(0.04,0.04) (0.234,8.182)rectangle++(0.04,0.04) (0.524,8.542)rectangle++(0.04,0.04) (0.785,8.836)rectangle++(0.04,0.04)
  (1.018,9.075)rectangle++(0.04,0.04) (2.371,2.873)rectangle++(0.04,0.04) (2.128,3.943)rectangle++(0.04,0.04) (1.964,4.863)rectangle++(0.04,0.04) (-0.02,0.761)rectangle++(0.04,0.04)
  (0.011,2.244)rectangle++(0.04,0.04) (0.097,3.503)rectangle++(0.04,0.04) (0.22,4.57)rectangle++(0.04,0.04) (0.368,5.472)rectangle++(0.04,0.04) (0.529,6.233)rectangle++(0.04,0.04)
  (0.697,6.873)rectangle++(0.04,0.04) (0.865,7.41)rectangle++(0.04,0.04) (1.03,7.86)rectangle++(0.04,0.04) (1.187,8.236)rectangle++(0.04,0.04) (1.337,8.549)rectangle++(0.04,0.04)
  (1.476,8.81)rectangle++(0.04,0.04) (2.228,2.928)rectangle++(0.04,0.04) (2.009,3.996)rectangle++(0.04,0.04) (1.865,4.912)rectangle++(0.04,0.04) (1.78,5.697)rectangle++(0.04,0.04)
  (1.738,6.368)rectangle++(0.04,0.04) (1.73,6.939)rectangle++(0.04,0.04) (1.746,7.425)rectangle++(0.04,0.04) (1.779,7.838)rectangle++(0.04,0.04) (1.823,8.187)rectangle++(0.04,0.04)
  (1.875,8.482)rectangle++(0.04,0.04) (1.931,8.731)rectangle++(0.04,0.04) (2.138,3.028)rectangle++(0.04,0.04) (1.936,4.084)rectangle++(0.04,0.04) (1.807,4.99)rectangle++(0.04,0.04)
  (1.733,5.766)rectangle++(0.04,0.04) (-0.02,0.906)rectangle++(0.04,0.04) (0.017,2.367)rectangle++(0.04,0.04) (0.107,3.607)rectangle++(0.04,0.04) (0.233,4.658)rectangle++(0.04,0.04)
  (1.252,1.609)rectangle++(0.04,0.04) (1.126,2.913)rectangle++(0.04,0.04) (1.072,4.028)rectangle++(0.04,0.04) (1.07,4.977)rectangle++(0.04,0.04) (1.106,5.784)rectangle++(0.04,0.04)
  (1.169,6.468)rectangle++(0.04,0.04) (1.251,7.047)rectangle++(0.04,0.04) (1.343,7.536)rectangle++(0.04,0.04) (1.891,2.588)rectangle++(0.04,0.04) (1.709,3.72)rectangle++(0.04,0.04)
  (1.599,4.69)rectangle++(0.04,0.04) (1.545,5.519)rectangle++(0.04,0.04) (-1.147,3.158)rectangle++(0.04,0.04) (-0.851,4.327)rectangle++(0.04,0.04) (-0.552,5.308)rectangle++(0.04,0.04)
  (-0.259,6.13)rectangle++(0.04,0.04) (0.022,6.817)rectangle++(0.04,0.04) (0.29,7.39)rectangle++(0.04,0.04) (0.54,7.866)rectangle++(0.04,0.04) (-0.02,1.242)rectangle++(0.04,0.04)
  (0.03,2.652)rectangle++(0.04,0.04) (0.13,3.85)rectangle++(0.04,0.04) (0.262,4.863)rectangle++(0.04,0.04) (1.305,1.665)rectangle++(0.04,0.04) (0.253,1.169)rectangle++(0.04,0.04)
  (0.26,2.58)rectangle++(0.04,0.04) (0.322,3.779)rectangle++(0.04,0.04) (0.992,1.421)rectangle++(0.04,0.04) (0.898,2.764)rectangle++(0.04,0.04) (0.622,1.327)rectangle++(0.04,0.04)
  (0.579,2.699)rectangle++(0.04,0.04) (0.598,3.867)rectangle++(0.04,0.04) (0.661,4.859)rectangle++(0.04,0.04) (0.754,5.7)rectangle++(0.04,0.04) (0.867,6.411)rectangle++(0.04,0.04)
  (0.991,7.011)rectangle++(0.04,0.04) (1.121,7.516)rectangle++(0.04,0.04) (1.251,7.94)rectangle++(0.04,0.04) (1.379,8.295)rectangle++(0.04,0.04) (1.501,8.592)rectangle++(0.04,0.04)
  (1.618,8.839)rectangle++(0.04,0.04) (1.726,9.045)rectangle++(0.04,0.04) (1.827,9.215)rectangle++(0.04,0.04) (1.919,9.356)rectangle++(0.04,0.04) (2.004,9.472)rectangle++(0.04,0.04)
  (2.08,9.567)rectangle++(0.04,0.04) (2.148,9.645)rectangle++(0.04,0.04) (2.21,9.709)rectangle++(0.04,0.04) (2.264,9.76)rectangle++(0.04,0.04) (2.313,9.802)rectangle++(0.04,0.04)
  (2.356,9.835)rectangle++(0.04,0.04) (2.394,9.862)rectangle++(0.04,0.04) (2.427,9.883)rectangle++(0.04,0.04) (2.456,9.9)rectangle++(0.04,0.04) (2.481,9.913)rectangle++(0.04,0.04)
  (2.503,9.923)rectangle++(0.04,0.04) (2.523,9.93)rectangle++(0.04,0.04) (2.539,9.936)rectangle++(0.04,0.04) (2.554,9.94)rectangle++(0.04,0.04) (2.566,9.943)rectangle++(0.04,0.04)
  (2.577,9.945)rectangle++(0.04,0.04) (2.586,9.947)rectangle++(0.04,0.04) (2.594,9.947)rectangle++(0.04,0.04) (2.6,9.948)rectangle++(0.04,0.04) (2.606,9.948)rectangle++(0.04,0.04)
  (2.611,9.948)rectangle++(0.04,0.04) (2.615,9.947)rectangle++(0.04,0.04) (2.618,9.947)rectangle++(0.04,0.04) (2.621,9.946)rectangle++(0.04,0.04) (2.624,9.946)rectangle++(0.04,0.04)
  (2.626,9.945)rectangle++(0.04,0.04) (2.627,9.944)rectangle++(0.04,0.04) (2.629,9.944)rectangle++(0.04,0.04) (2.63,9.943)rectangle++(0.04,0.04) (2.631,9.943)rectangle++(0.04,0.04)
  (2.632,9.942)rectangle++(0.04,0.04) (2.633,9.942)rectangle++(0.04,0.04) (2.633,9.942)rectangle++(0.04,0.04) (2.634,9.941)rectangle++(0.04,0.04) (2.634,9.941)rectangle++(0.04,0.04)
  (2.634,9.941)rectangle++(0.04,0.04) (2.371,3.501)rectangle++(0.04,0.04) (2.153,4.477)rectangle++(0.04,0.04) (2.007,5.315)rectangle++(0.04,0.04) (1.916,6.034)rectangle++(0.04,0.04)
  (1.868,6.649)rectangle++(0.04,0.04) (-0.02,1.047)rectangle++(0.04,0.04) (0.023,2.487)rectangle++(0.04,0.04) (0.117,3.709)rectangle++(0.04,0.04) (0.245,4.744)rectangle++(0.04,0.04)
  (-1.206,2.689)rectangle++(0.04,0.04) (-0.919,3.93)rectangle++(0.04,0.04) (-1.227,2.242)rectangle++(0.04,0.04) (-0.955,3.551)rectangle++(0.04,0.04) (-0.672,4.653)rectangle++(0.04,0.04)
  (-0.388,5.578)rectangle++(0.04,0.04) (-0.108,6.353)rectangle++(0.04,0.04) (0.16,7.001)rectangle++(0.04,0.04) (0.414,7.54)rectangle++(0.04,0.04) (0.651,7.989)rectangle++(0.04,0.04)
  (0.871,8.361)rectangle++(0.04,0.04) (1.072,8.668)rectangle++(0.04,0.04) (1.256,8.921)rectangle++(0.04,0.04) (1.422,9.129)rectangle++(0.04,0.04) (1.572,9.299)rectangle++(0.04,0.04)
  (1.706,9.437)rectangle++(0.04,0.04) (1.825,9.55)rectangle++(0.04,0.04) (1.931,9.64)rectangle++(0.04,0.04) (2.025,9.713)rectangle++(0.04,0.04) (2.108,9.772)rectangle++(0.04,0.04)
  (2.18,9.818)rectangle++(0.04,0.04) (2.244,9.854)rectangle++(0.04,0.04) (2.405,3.378)rectangle++(0.04,0.04) (2.177,4.372)rectangle++(0.04,0.04) (2.023,5.225)rectangle++(0.04,0.04)
  (1.927,5.956)rectangle++(0.04,0.04) (1.874,6.582)rectangle++(0.04,0.04) (1.854,7.116)rectangle++(0.04,0.04) (1.858,7.571)rectangle++(0.04,0.04) (1.88,7.957)rectangle++(0.04,0.04)
  (1.914,8.284)rectangle++(0.04,0.04) (1.956,8.561)rectangle++(0.04,0.04) (2.003,8.795)rectangle++(0.04,0.04) (2.052,8.992)rectangle++(0.04,0.04) (2.193,3.122)rectangle++(0.04,0.04)
  (1.986,4.162)rectangle++(0.04,0.04) (1.853,5.054)rectangle++(0.04,0.04) (1.775,5.818)rectangle++(0.04,0.04) (1.739,6.471)rectangle++(0.04,0.04) (1.534,2.435)rectangle++(0.04,0.04)
  (-0.348,2.477)rectangle++(0.04,0.04) (-0.199,3.716)rectangle++(0.04,0.04) (-0.022,4.763)rectangle++(0.04,0.04) (0.169,5.645)rectangle++(0.04,0.04) (0.368,6.388)rectangle++(0.04,0.04)
  (0.566,7.011)rectangle++(0.04,0.04) (0.759,7.533)rectangle++(0.04,0.04) (0.944,7.969)rectangle++(0.04,0.04) (1.119,8.332)rectangle++(0.04,0.04) (2.148,2.721)rectangle++(0.04,0.04)
  (1.932,3.823)rectangle++(0.04,0.04) (1.793,4.768)rectangle++(0.04,0.04) (1.713,5.578)rectangle++(0.04,0.04) (1.677,6.269)rectangle++(0.04,0.04) (1.674,6.857)rectangle++(0.04,0.04)
  (1.695,7.358)rectangle++(0.04,0.04) (1.733,7.783)rectangle++(0.04,0.04) (1.782,8.142)rectangle++(0.04,0.04) (1.838,8.446)rectangle++(0.04,0.04) (1.898,8.702)rectangle++(0.04,0.04)
  (1.959,8.917)rectangle++(0.04,0.04) (2.02,9.097)rectangle++(0.04,0.04) (2.078,9.248)rectangle++(0.04,0.04) (2.134,9.374)rectangle++(0.04,0.04) (2.187,9.479)rectangle++(0.04,0.04)
  (-2.048,4.177)rectangle++(0.04,0.04) (-1.576,5.229)rectangle++(0.04,0.04) (-1.133,6.104)rectangle++(0.04,0.04) (-0.721,6.83)rectangle++(0.04,0.04) (-0.342,7.43)rectangle++(0.04,0.04)
  (0.005,7.926)rectangle++(0.04,0.04) (0.319,8.333)rectangle++(0.04,0.04) (0.602,8.666)rectangle++(0.04,0.04) (2.319,2.666)rectangle++(0.04,0.04) (-0.251,2.709)rectangle++(0.04,0.04)
  (-0.107,3.909)rectangle++(0.04,0.04) (0.063,4.923)rectangle++(0.04,0.04) (0.248,5.778)rectangle++(0.04,0.04) (0.44,6.498)rectangle++(0.04,0.04) (0.632,7.102)rectangle++(0.04,0.04)
  (-1.741,3.297)rectangle++(0.04,0.04) (-1.35,4.468)rectangle++(0.04,0.04) (-0.971,5.448)rectangle++(0.04,0.04) (-0.61,6.266)rectangle++(0.04,0.04) (-0.27,6.947)rectangle++(0.04,0.04)
  (0.046,7.512)rectangle++(0.04,0.04) (0.338,7.979)rectangle++(0.04,0.04) (0.604,8.365)rectangle++(0.04,0.04) (0.846,8.682)rectangle++(0.04,0.04) (2.287,2.734)rectangle++(0.04,0.04)
  (2.051,3.828)rectangle++(0.04,0.04) (1.894,4.768)rectangle++(0.04,0.04) (1.799,5.573)rectangle++(0.04,0.04) (1.75,6.262)rectangle++(0.04,0.04) (1.473,2.388)rectangle++(0.04,0.04)
  (1.346,3.567)rectangle++(0.04,0.04) (1.284,4.574)rectangle++(0.04,0.04) (1.272,5.433)rectangle++(0.04,0.04) (1.297,6.163)rectangle++(0.04,0.04) (1.347,6.783)rectangle++(0.04,0.04)
  (1.414,7.308)rectangle++(0.04,0.04) (1.492,7.751)rectangle++(0.04,0.04) (1.576,8.125)rectangle++(0.04,0.04) (1.662,8.44)rectangle++(0.04,0.04) (-1.883,3.828)rectangle++(0.04,0.04)
  (-1.45,4.925)rectangle++(0.04,0.04) (-1.037,5.841)rectangle++(0.04,0.04) (-0.65,6.602)rectangle++(0.04,0.04) (-0.291,7.234)rectangle++(0.04,0.04) (0.04,7.757)rectangle++(0.04,0.04)
  (0.342,8.188)rectangle++(0.04,0.04) (0.616,8.542)rectangle++(0.04,0.04) (0.863,8.833)rectangle++(0.04,0.04) (1.085,9.069)rectangle++(0.04,0.04) (1.283,9.262)rectangle++(0.04,0.04)
  (1.458,9.417)rectangle++(0.04,0.04) (1.614,9.543)rectangle++(0.04,0.04) (1.752,9.643)rectangle++(0.04,0.04) (-2.178,4.113)rectangle++(0.04,0.04) (-1.689,5.18)rectangle++(0.04,0.04)
  (-1.706,2.34)rectangle++(0.04,0.04) (-1.358,3.653)rectangle++(0.04,0.04) (-1.011,4.756)rectangle++(0.04,0.04) (-0.671,5.679)rectangle++(0.04,0.04) (1.673,1.619)rectangle++(0.04,0.04)
  (1.485,2.905)rectangle++(0.04,0.04) (1.376,4.006)rectangle++(0.04,0.04) (1.328,4.946)rectangle++(0.04,0.04) (1.324,5.747)rectangle++(0.04,0.04) (1.353,6.429)rectangle++(0.04,0.04)
  (1.405,7.006)rectangle++(0.04,0.04) (1.473,7.495)rectangle++(0.04,0.04) (1.549,7.908)rectangle++(0.04,0.04) (1.631,8.256)rectangle++(0.04,0.04) (-1.842,3.781)rectangle++(0.04,0.04)
  (-0.02,0.588)rectangle++(0.04,0.04) (0.004,2.097)rectangle++(0.04,0.04) (0.085,3.378)rectangle++(0.04,0.04) (0.205,4.464)rectangle++(0.04,0.04) (0.351,5.383)rectangle++(0.04,0.04)
  (0.511,6.157)rectangle++(0.04,0.04) (0.679,6.81)rectangle++(0.04,0.04) (0.847,7.357)rectangle++(0.04,0.04) (1.916,2.416)rectangle++(0.04,0.04) (1.723,3.573)rectangle++(0.04,0.04)
  (0.725,1.735)rectangle++(0.04,0.04) (0.683,3.042)rectangle++(0.04,0.04) (0.7,4.155)rectangle++(0.04,0.04) (0.759,5.1)rectangle++(0.04,0.04) (1.297,1.851)rectangle++(0.04,0.04)
  (1.174,3.118)rectangle++(0.04,0.04) (1.12,4.2)rectangle++(0.04,0.04) (1.118,5.121)rectangle++(0.04,0.04) (1.153,5.904)rectangle++(0.04,0.04) (1.214,6.569)rectangle++(0.04,0.04)
  (1.293,7.131)rectangle++(0.04,0.04) (1.382,7.606)rectangle++(0.04,0.04) (1.476,8.006)rectangle++(0.04,0.04) (1.573,8.342)rectangle++(0.04,0.04) (1.669,8.624)rectangle++(0.04,0.04)
  (1.761,8.86)rectangle++(0.04,0.04) (1.849,9.057)rectangle++(0.04,0.04) (1.932,9.22)rectangle++(0.04,0.04) (-2.032,4.062)rectangle++(0.04,0.04) (-1.567,5.13)rectangle++(0.04,0.04)
  (-1.129,6.019)rectangle++(0.04,0.04) (1.837,1.581)rectangle++(0.04,0.04) (1.623,2.867)rectangle++(0.04,0.04) (1.492,3.968)rectangle++(0.04,0.04) (1.425,4.909)rectangle++(0.04,0.04)
  (1.405,5.712)rectangle++(0.04,0.04) (1.421,6.395)rectangle++(0.04,0.04) (1.461,6.975)rectangle++(0.04,0.04) (1.519,7.467)rectangle++(0.04,0.04) (1.587,7.882)rectangle++(0.04,0.04)
  (-1.753,3.688)rectangle++(0.04,0.04) (-1.345,4.801)rectangle++(0.04,0.04) (-0.953,5.731)rectangle++(0.04,0.04) (-0.583,6.506)rectangle++(0.04,0.04) (-1.829,2.886)rectangle++(0.04,0.04)
  (-1.442,4.123)rectangle++(0.04,0.04) (-1.063,5.158)rectangle++(0.04,0.04) (-0.699,6.023)rectangle++(0.04,0.04) (-0.356,6.744)rectangle++(0.04,0.04) (-0.035,7.343)rectangle++(0.04,0.04)
  (0.262,7.839)rectangle++(0.04,0.04) (2.138,2.379)rectangle++(0.04,0.04) (1.91,3.533)rectangle++(0.04,0.04) (1.763,4.523)rectangle++(0.04,0.04) (1.677,5.37)rectangle++(0.04,0.04)
  (1.638,6.094)rectangle++(0.04,0.04) (1.634,6.71)rectangle++(0.04,0.04) (1.655,7.235)rectangle++(0.04,0.04) (-1.571,3.561)rectangle++(0.04,0.04) (-1.261,2.011)rectangle++(0.04,0.04)
  (-0.994,3.356)rectangle++(0.04,0.04) (1.071,0.977)rectangle++(0.04,0.04) (0.948,2.384)rectangle++(0.04,0.04) (0.899,3.585)rectangle++(0.04,0.04) (0.905,4.607)rectangle++(0.04,0.04)
  (1.137,1.771)rectangle++(0.04,0.04) (1.035,3.056)rectangle++(0.04,0.04) (1,4.152)rectangle++(0.04,0.04) (1.014,5.086)rectangle++(0.04,0.04) (1.063,5.879)rectangle++(0.04,0.04)
  (1.136,6.55)rectangle++(0.04,0.04) (1.646,2.298)rectangle++(0.04,0.04) (1.489,3.483)rectangle++(0.04,0.04) (1.403,4.497)rectangle++(0.04,0.04) (1.37,5.363)rectangle++(0.04,0.04)
  (-1.142,3.084)rectangle++(0.04,0.04) (-0.849,4.263)rectangle++(0.04,0.04) (-0.553,5.254)rectangle++(0.04,0.04) (-0.262,6.084)rectangle++(0.04,0.04) (0.018,6.778)rectangle++(0.04,0.04)
  (0.284,7.357)rectangle++(0.04,0.04) (0.534,7.838)rectangle++(0.04,0.04) (0.765,8.237)rectangle++(0.04,0.04) (0.978,8.567)rectangle++(0.04,0.04) (1.171,8.839)rectangle++(0.04,0.04)
  (1.347,9.063)rectangle++(0.04,0.04) (1.505,9.246)rectangle++(0.04,0.04) (1.647,9.395)rectangle++(0.04,0.04) (1.774,9.516)rectangle++(0.04,0.04) (1.886,9.614)rectangle++(0.04,0.04)
  (1.985,9.692)rectangle++(0.04,0.04) (2.073,9.755)rectangle++(0.04,0.04) (2.15,9.805)rectangle++(0.04,0.04) (-2.141,4.241)rectangle++(0.04,0.04) (-1.652,5.286)rectangle++(0.04,0.04)
  (-1.726,2.372)rectangle++(0.04,0.04) (-0.983,1.714)rectangle++(0.04,0.04) (-0.663,1.74)rectangle++(0.04,0.04) (-0.497,3.102)rectangle++(0.04,0.04) (-0.3,4.252)rectangle++(0.04,0.04)
  (-0.087,5.223)rectangle++(0.04,0.04) (0.133,6.039)rectangle++(0.04,0.04) (0.352,6.724)rectangle++(0.04,0.04) (0.566,7.298)rectangle++(0.04,0.04) (0.771,7.777)rectangle++(0.04,0.04)
  (0.964,8.175)rectangle++(0.04,0.04) (1.144,8.507)rectangle++(0.04,0.04) (1.311,8.781)rectangle++(0.04,0.04) (1.463,9.008)rectangle++(0.04,0.04) (1.602,9.194)rectangle++(0.04,0.04)
  (2.317,3.053)rectangle++(0.04,0.04) (2.089,4.099)rectangle++(0.04,0.04) (1.938,4.996)rectangle++(0.04,0.04) (1.845,5.766)rectangle++(0.04,0.04) (1.796,6.423)rectangle++(0.04,0.04)
  (1.782,6.984)rectangle++(0.04,0.04) (1.792,7.461)rectangle++(0.04,0.04) (1.819,7.867)rectangle++(0.04,0.04) (1.859,8.21)rectangle++(0.04,0.04) (1.906,8.501)rectangle++(0.04,0.04)
  (1.958,8.745)rectangle++(0.04,0.04) (2.012,8.951)rectangle++(0.04,0.04) (2.066,9.124)rectangle++(0.04,0.04) (2.228,3.157)rectangle++(0.04,0.04) (2.018,4.191)rectangle++(0.04,0.04)
  (1.88,5.077)rectangle++(0.04,0.04) (1.799,5.837)rectangle++(0.04,0.04) (1.347,2.299)rectangle++(0.04,0.04) (-0.349,2.405)rectangle++(0.04,0.04) (-0.203,3.654)rectangle++(0.04,0.04)
  (-0.029,4.71)rectangle++(0.04,0.04) (0.162,5.601)rectangle++(0.04,0.04) (0.359,6.351)rectangle++(0.04,0.04) (0.557,6.98)rectangle++(0.04,0.04) (0.751,7.507)rectangle++(0.04,0.04)
  (0.936,7.947)rectangle++(0.04,0.04) (1.111,8.314)rectangle++(0.04,0.04) (-1.96,3.674)rectangle++(0.04,0.04) (-1.522,4.797)rectangle++(0.04,0.04) (-1.104,5.735)rectangle++(0.04,0.04)
  (-0.711,6.515)rectangle++(0.04,0.04) (-0.346,7.162)rectangle++(0.04,0.04) (-0.01,7.698)rectangle++(0.04,0.04) (-2.025,3.28)rectangle++(0.04,0.04) (-1.592,4.465)rectangle++(0.04,0.04)
  (-1.177,5.456)rectangle++(0.04,0.04) (-0.784,6.28)rectangle++(0.04,0.04) (-0.418,6.966)rectangle++(0.04,0.04) (-0.078,7.534)rectangle++(0.04,0.04) (0.232,8.003)rectangle++(0.04,0.04)
  (0.516,8.39)rectangle++(0.04,0.04) (0.772,8.707)rectangle++(0.04,0.04) (2.305,2.72)rectangle++(0.04,0.04) (2.066,3.816)rectangle++(0.04,0.04) (1.906,4.757)rectangle++(0.04,0.04)
  (1.808,5.564)rectangle++(0.04,0.04) (1.269,2.235)rectangle++(0.04,0.04) (1.166,3.446)rectangle++(0.04,0.04) (1.127,4.478)rectangle++(0.04,0.04) (1.068,1.798)rectangle++(0.04,0.04)
  (0.977,3.082)rectangle++(0.04,0.04) (0.952,4.176)rectangle++(0.04,0.04) (1.009,1.68)rectangle++(0.04,0.04) (0.923,2.984)rectangle++(0.04,0.04) (0.902,4.095)rectangle++(0.04,0.04)
  (0.928,5.041)rectangle++(0.04,0.04) (0.988,5.844)rectangle++(0.04,0.04) (1.072,6.524)rectangle++(0.04,0.04) (1.17,7.099)rectangle++(0.04,0.04) (1.276,7.583)rectangle++(0.04,0.04)
  (1.386,7.991)rectangle++(0.04,0.04) (1.495,8.333)rectangle++(0.04,0.04) (1.602,8.62)rectangle++(0.04,0.04) (1.704,8.859)rectangle++(0.04,0.04) (-1.984,3.93)rectangle++(0.04,0.04)
  (-1.531,5.016)rectangle++(0.04,0.04) (-1.103,5.921)rectangle++(0.04,0.04) (-0.703,6.673)rectangle++(0.04,0.04) (-0.333,7.297)rectangle++(0.04,0.04) (-1.985,3.118)rectangle++(0.04,0.04)
  (-1.565,4.326)rectangle++(0.04,0.04) (-1.159,5.336)rectangle++(0.04,0.04) (-0.774,6.178)rectangle++(0.04,0.04) (-0.413,6.878)rectangle++(0.04,0.04) (-0.078,7.459)rectangle++(0.04,0.04)
  (0.23,7.94)rectangle++(0.04,0.04) (0.511,8.336)rectangle++(0.04,0.04) (0.765,8.661)rectangle++(0.04,0.04) (0.995,8.928)rectangle++(0.04,0.04) (-2.143,3.782)rectangle++(0.04,0.04)
  (-1.673,4.897)rectangle++(0.04,0.04) (-1.228,5.825)rectangle++(0.04,0.04) (-0.813,6.597)rectangle++(0.04,0.04) (-0.43,7.236)rectangle++(0.04,0.04) (2.073,2.055)rectangle++(0.04,0.04)
  (1.842,3.26)rectangle++(0.04,0.04) (0.619,1.691)rectangle++(0.04,0.04) (0.592,3.009)rectangle++(0.04,0.04) (0.621,4.13)rectangle++(0.04,0.04) (0.691,5.082)rectangle++(0.04,0.04)
  (-1.204,2.866)rectangle++(0.04,0.04) (-0.911,4.08)rectangle++(0.04,0.04) (-0.614,5.101)rectangle++(0.04,0.04) (-0.32,5.957)rectangle++(0.04,0.04) (-0.036,6.672)rectangle++(0.04,0.04)
  (0.234,7.269)rectangle++(0.04,0.04) (0.488,7.765)rectangle++(0.04,0.04) (0.723,8.177)rectangle++(0.04,0.04) (0.939,8.518)rectangle++(0.04,0.04) (1.137,8.799)rectangle++(0.04,0.04)
  (1.316,9.03)rectangle++(0.04,0.04) (1.478,9.219)rectangle++(0.04,0.04) (1.623,9.373)rectangle++(0.04,0.04) (1.752,9.498)rectangle++(0.04,0.04) (-2.14,4.082)rectangle++(0.04,0.04)
  (-1.658,5.151)rectangle++(0.04,0.04) (-1.206,6.041)rectangle++(0.04,0.04) (-0.785,6.779)rectangle++(0.04,0.04) (-0.399,7.39)rectangle++(0.04,0.04) (-0.045,7.894)rectangle++(0.04,0.04)
  (0.275,8.308)rectangle++(0.04,0.04) (2.267,2.495)rectangle++(0.04,0.04) (2.025,3.627)rectangle++(0.04,0.04) (1.864,4.598)rectangle++(0.04,0.04) (1.766,5.43)rectangle++(0.04,0.04)
  (1.716,6.141)rectangle++(0.04,0.04) (-1.275,3.335)rectangle++(0.04,0.04) (-0.952,4.482)rectangle++(0.04,0.04) (-0.632,5.444)rectangle++(0.04,0.04) (-0.322,6.249)rectangle++(0.04,0.04)
  (-1.71,2.89)rectangle++(0.04,0.04) (1.048,0.679)rectangle++(0.04,0.04) (0.916,2.131)rectangle++(0.04,0.04) (-0.392,2.269)rectangle++(0.04,0.04) (0.677,0.872)rectangle++(0.04,0.04)
  (0.608,2.311)rectangle++(0.04,0.04) (0.607,3.536)rectangle++(0.04,0.04) (0.655,4.578)rectangle++(0.04,0.04) (0.738,5.461)rectangle++(0.04,0.04) (0.843,6.208)rectangle++(0.04,0.04)
  (0.963,6.84)rectangle++(0.04,0.04) (1.09,7.371)rectangle++(0.04,0.04) (1.219,7.818)rectangle++(0.04,0.04) (1.347,8.193)rectangle++(0.04,0.04) (-0.02,1.294)rectangle++(0.04,0.04)
  (0.033,2.697)rectangle++(0.04,0.04) (0.133,3.887)rectangle++(0.04,0.04) (0.267,4.895)rectangle++(0.04,0.04) (0.42,5.746)rectangle++(0.04,0.04) (0.585,6.464)rectangle++(0.04,0.04)
  (-1.585,3.146)rectangle++(0.04,0.04) (-1.156,1.917)rectangle++(0.04,0.04) (-0.908,3.271)rectangle++(0.04,0.04) (-0.643,4.413)rectangle++(0.04,0.04) (-0.372,5.373)rectangle++(0.04,0.04)
  (-0.104,6.178)rectangle++(0.04,0.04) (0.157,6.852)rectangle++(0.04,0.04) (0.405,7.414)rectangle++(0.04,0.04) (0.639,7.882)rectangle++(0.04,0.04) (0.856,8.27)rectangle++(0.04,0.04)
  (1.056,8.592)rectangle++(0.04,0.04) (1.239,8.857)rectangle++(0.04,0.04) (1.405,9.075)rectangle++(0.04,0.04) (1.555,9.254)rectangle++(0.04,0.04) (-2.116,3.983)rectangle++(0.04,0.04)
  (-1.48,1.978)rectangle++(0.04,0.04) (-1.181,3.337)rectangle++(0.04,0.04) (-0.873,4.48)rectangle++(0.04,0.04) (-0.565,5.439)rectangle++(0.04,0.04) (-0.265,6.242)rectangle++(0.04,0.04)
  (0.023,6.912)rectangle++(0.04,0.04) (1.915,2.095)rectangle++(0.04,0.04) (1.709,3.3)rectangle++(0.04,0.04) (1.583,4.333)rectangle++(0.04,0.04) (-0.831,2.906)rectangle++(0.04,0.04)
  (-0.593,4.1)rectangle++(0.04,0.04) (1.219,1.26)rectangle++(0.04,0.04) (1.085,2.618)rectangle++(0.04,0.04) (1.025,3.778)rectangle++(0.04,0.04) (1.02,4.767)rectangle++(0.04,0.04)
  (1.055,5.607)rectangle++(0.04,0.04) (1.119,6.32)rectangle++(0.04,0.04) (1.202,6.924)rectangle++(0.04,0.04) (1.296,7.433)rectangle++(0.04,0.04) (1.397,7.863)rectangle++(0.04,0.04)
  (1.5,8.223)rectangle++(0.04,0.04) (2.06,2.794)rectangle++(0.04,0.04) (1.861,3.888)rectangle++(0.04,0.04) (1.735,4.827)rectangle++(0.04,0.04) (1.666,5.63)rectangle++(0.04,0.04)
  (1.639,6.315)rectangle++(0.04,0.04) (1.643,6.898)rectangle++(0.04,0.04) (1.671,7.394)rectangle++(0.04,0.04) (1.714,7.814)rectangle++(0.04,0.04) (1.767,8.17)rectangle++(0.04,0.04)
  (1.826,8.47)rectangle++(0.04,0.04) (1.889,8.722)rectangle++(0.04,0.04) (1.952,8.935)rectangle++(0.04,0.04) (2.015,9.113)rectangle++(0.04,0.04) (2.075,9.261)rectangle++(0.04,0.04)
  (2.132,9.385)rectangle++(0.04,0.04) (2.185,9.488)rectangle++(0.04,0.04) (2.235,9.574)rectangle++(0.04,0.04) (2.28,9.645)rectangle++(0.04,0.04) (-2.073,4.235)rectangle++(0.04,0.04)
  (-1.595,5.279)rectangle++(0.04,0.04) (-1.146,6.147)rectangle++(0.04,0.04) (-0.731,6.867)rectangle++(0.04,0.04) (-0.349,7.463)rectangle++(0.04,0.04) (0,7.953)rectangle++(0.04,0.04)
  (0.316,8.357)rectangle++(0.04,0.04) (0.601,8.687)rectangle++(0.04,0.04) (0.856,8.956)rectangle++(0.04,0.04) (1.083,9.174)rectangle++(0.04,0.04) (1.286,9.351)rectangle++(0.04,0.04)
  (1.465,9.493)rectangle++(0.04,0.04) (1.622,9.607)rectangle++(0.04,0.04) (1.761,9.697)rectangle++(0.04,0.04) (1.883,9.768)rectangle++(0.04,0.04) (1.989,9.824)rectangle++(0.04,0.04)
  (2.081,9.867)rectangle++(0.04,0.04) (2.162,9.9)rectangle++(0.04,0.04) (2.231,9.925)rectangle++(0.04,0.04) (2.291,9.943)rectangle++(0.04,0.04) (2.343,9.956)rectangle++(0.04,0.04)
  (2.388,9.965)rectangle++(0.04,0.04) (2.426,9.971)rectangle++(0.04,0.04) (2.459,9.975)rectangle++(0.04,0.04) (2.487,9.976)rectangle++(0.04,0.04) (2.51,9.977)rectangle++(0.04,0.04)
  (2.531,9.976)rectangle++(0.04,0.04) (2.548,9.974)rectangle++(0.04,0.04) (2.563,9.973)rectangle++(0.04,0.04) (2.575,9.97)rectangle++(0.04,0.04) (2.585,9.968)rectangle++(0.04,0.04)
  (2.594,9.966)rectangle++(0.04,0.04) (2.601,9.963)rectangle++(0.04,0.04) (2.607,9.961)rectangle++(0.04,0.04) (2.381,3.499)rectangle++(0.04,0.04) (2.161,4.475)rectangle++(0.04,0.04)
  (-0.752,3.071)rectangle++(0.04,0.04) (-0.519,4.236)rectangle++(0.04,0.04) (-0.274,5.218)rectangle++(0.04,0.04) (-0.026,6.042)rectangle++(0.04,0.04) (1.678,1.873)rectangle++(0.04,0.04)
  (0.255,1.316)rectangle++(0.04,0.04) (0.268,2.705)rectangle++(0.04,0.04) (0.333,3.884)rectangle++(0.04,0.04) (0.437,4.885)rectangle++(0.04,0.04) (0.564,5.731)rectangle++(0.04,0.04)
  (-1.398,2.98)rectangle++(0.04,0.04) (-1.072,4.185)rectangle++(0.04,0.04) (-0.746,5.196)rectangle++(0.04,0.04) (-0.428,6.043)rectangle++(0.04,0.04) (-0.124,6.75)rectangle++(0.04,0.04)
  (-1.801,3.045)rectangle++(0.04,0.04) (-1.411,4.257)rectangle++(0.04,0.04) (-1.031,5.271)rectangle++(0.04,0.04) (-0.668,6.118)rectangle++(0.04,0.04) (-0.325,6.823)rectangle++(0.04,0.04)
  (-0.006,7.409)rectangle++(0.04,0.04) (0.289,7.894)rectangle++(0.04,0.04) (0.559,8.294)rectangle++(0.04,0.04) (0.805,8.624)rectangle++(0.04,0.04) (1.027,8.894)rectangle++(0.04,0.04)
  (1.227,9.115)rectangle++(0.04,0.04) (1.405,9.295)rectangle++(0.04,0.04) (1.564,9.441)rectangle++(0.04,0.04) (1.705,9.558)rectangle++(0.04,0.04) (1.829,9.653)rectangle++(0.04,0.04)
  (1.939,9.728)rectangle++(0.04,0.04) (2.035,9.787)rectangle++(0.04,0.04) (2.119,9.834)rectangle++(0.04,0.04) (2.192,9.87)rectangle++(0.04,0.04) (2.256,9.898)rectangle++(0.04,0.04)
  (2.311,9.92)rectangle++(0.04,0.04) (2.359,9.935)rectangle++(0.04,0.04) (2.401,9.947)rectangle++(0.04,0.04) (2.436,9.955)rectangle++(0.04,0.04) (2.467,9.961)rectangle++(0.04,0.04)
  (2.493,9.964)rectangle++(0.04,0.04) (2.515,9.966)rectangle++(0.04,0.04) (2.396,3.476)rectangle++(0.04,0.04) (0.596,1.887)rectangle++(0.04,0.04) (0.58,3.176)rectangle++(0.04,0.04)
  (0.618,4.273)rectangle++(0.04,0.04) (-1.009,2.671)rectangle++(0.04,0.04) (-0.917,1.945)rectangle++(0.04,0.04) (-0.704,3.286)rectangle++(0.04,0.04) (-0.469,4.417)rectangle++(0.04,0.04)
  (-0.224,5.37)rectangle++(0.04,0.04) (1.52,1.66)rectangle++(0.04,0.04) (-0.149,2.304)rectangle++(0.04,0.04) (-0.037,3.56)rectangle++(0.04,0.04) (0.109,4.624)rectangle++(0.04,0.04)
  (0.275,5.522)rectangle++(0.04,0.04) (0.453,6.279)rectangle++(0.04,0.04) (1.673,2.055)rectangle++(0.04,0.04) (1.502,3.276)rectangle++(0.04,0.04) (1.405,4.321)rectangle++(0.04,0.04)
  (-0.863,2.863)rectangle++(0.04,0.04) (-0.622,4.064)rectangle++(0.04,0.04) (-0.368,5.076)rectangle++(0.04,0.04) (-0.112,5.925)rectangle++(0.04,0.04) (0.14,6.637)rectangle++(0.04,0.04)
  (-1.719,3.081)rectangle++(0.04,0.04) (-1.34,4.284)rectangle++(0.04,0.04) (-0.97,5.291)rectangle++(0.04,0.04) (-0.615,6.133)rectangle++(0.04,0.04) (-0.28,6.833)rectangle++(0.04,0.04)
  (0.033,7.416)rectangle++(0.04,0.04) (2.054,2.219)rectangle++(0.04,0.04) (1.832,3.4)rectangle++(0.04,0.04) (1.691,4.413)rectangle++(0.04,0.04) (1.612,5.279)rectangle++(0.04,0.04)
  (1.579,6.019)rectangle++(0.04,0.04) (1.581,6.649)rectangle++(0.04,0.04) (1.607,7.185)rectangle++(0.04,0.04) (1.652,7.639)rectangle++(0.04,0.04) (1.707,8.023)rectangle++(0.04,0.04)
  (1.77,8.348)rectangle++(0.04,0.04) (1.836,8.621)rectangle++(0.04,0.04) (1.903,8.851)rectangle++(0.04,0.04) (1.97,9.043)rectangle++(0.04,0.04) (2.034,9.204)rectangle++(0.04,0.04)
  (2.095,9.338)rectangle++(0.04,0.04) (2.152,9.45)rectangle++(0.04,0.04) (2.306,3.257)rectangle++(0.04,0.04) (2.088,4.273)rectangle++(0.04,0.04) (-0.715,3.009)rectangle++(0.04,0.04)
  (-0.489,4.183)rectangle++(0.04,0.04) (-0.251,5.171)rectangle++(0.04,0.04) (-0.008,6.002)rectangle++(0.04,0.04) (0.231,6.698)rectangle++(0.04,0.04) (0.462,7.28)rectangle++(0.04,0.04)
  (1.952,2.297)rectangle++(0.04,0.04) (1.749,3.471)rectangle++(0.04,0.04) (-0.574,2.755)rectangle++(0.04,0.04) (-0.38,3.961)rectangle++(0.04,0.04) (-0.167,4.978)rectangle++(0.04,0.04)
  (0.055,5.834)rectangle++(0.04,0.04) (0.278,6.553)rectangle++(0.04,0.04) (0.496,7.155)rectangle++(0.04,0.04) (0.706,7.658)rectangle++(0.04,0.04) (0.904,8.077)rectangle++(0.04,0.04)
  (1.089,8.426)rectangle++(0.04,0.04) (1.261,8.715)rectangle++(0.04,0.04) (1.418,8.953)rectangle++(0.04,0.04) (1.561,9.15)rectangle++(0.04,0.04) (1.691,9.311)rectangle++(0.04,0.04)
  (2.336,3.104)rectangle++(0.04,0.04) (2.108,4.141)rectangle++(0.04,0.04) (1.955,5.032)rectangle++(0.04,0.04) (1.861,5.795)rectangle++(0.04,0.04) (1.811,6.448)rectangle++(0.04,0.04)
  (1.795,7.004)rectangle++(0.04,0.04) (1.804,7.478)rectangle++(0.04,0.04) (1.83,7.88)rectangle++(0.04,0.04) (1.869,8.221)rectangle++(0.04,0.04) (1.915,8.51)rectangle++(0.04,0.04)
  (1.966,8.753)rectangle++(0.04,0.04) (2.019,8.957)rectangle++(0.04,0.04) (-1.946,4.024)rectangle++(0.04,0.04) (-1.496,5.094)rectangle++(0.04,0.04) (-1.07,5.986)rectangle++(0.04,0.04)
  (-0.672,6.727)rectangle++(0.04,0.04) (-1.905,2.914)rectangle++(0.04,0.04) (-1.505,4.15)rectangle++(0.04,0.04) (-1.115,5.184)rectangle++(0.04,0.04) (-1.592,2.473)rectangle++(0.04,0.04)
  (-1.256,3.762)rectangle++(0.04,0.04) (-0.92,4.844)rectangle++(0.04,0.04) (-0.59,5.75)rectangle++(0.04,0.04) (-0.274,6.508)rectangle++(0.04,0.04) (0.025,7.139)rectangle++(0.04,0.04)
  (0.305,7.663)rectangle++(0.04,0.04) (0.563,8.098)rectangle++(0.04,0.04) (0.801,8.457)rectangle++(0.04,0.04) (1.017,8.752)rectangle++(0.04,0.04) (1.212,8.995)rectangle++(0.04,0.04)
  (1.388,9.193)rectangle++(0.04,0.04) (1.545,9.355)rectangle++(0.04,0.04) (1.685,9.486)rectangle++(0.04,0.04) (1.81,9.592)rectangle++(0.04,0.04) (2.397,3.203)rectangle++(0.04,0.04)
  (2.163,4.223)rectangle++(0.04,0.04) (2.005,5.099)rectangle++(0.04,0.04) (1.906,5.85)rectangle++(0.04,0.04) (-0.02,0.919)rectangle++(0.04,0.04) (0.018,2.378)rectangle++(0.04,0.04)
  (0.108,3.617)rectangle++(0.04,0.04) (-0.94,2.41)rectangle++(0.04,0.04) (-0.705,3.682)rectangle++(0.04,0.04) (-0.454,4.754)rectangle++(0.04,0.04) (-0.198,5.655)rectangle++(0.04,0.04)
  (-0.02,0.888)rectangle++(0.04,0.04) (0.016,2.352)rectangle++(0.04,0.04) (-0.629,2.11)rectangle++(0.04,0.04) (-0.453,3.415)rectangle++(0.04,0.04) (-0.25,4.517)rectangle++(0.04,0.04)
  (-0.034,5.446)rectangle++(0.04,0.04) (0.186,6.226)rectangle++(0.04,0.04) (0.405,6.881)rectangle++(0.04,0.04) (0.618,7.429)rectangle++(0.04,0.04) (0.82,7.886)rectangle++(0.04,0.04)
  (2.068,2.536)rectangle++(0.04,0.04) (1.857,3.669)rectangle++(0.04,0.04) (1.723,4.641)rectangle++(0.04,0.04) (1.648,5.472)rectangle++(0.04,0.04) (1.617,6.181)rectangle++(0.04,0.04)
  (1.471,2.334)rectangle++(0.04,0.04) (0.416,1.373)rectangle++(0.04,0.04) (0.305,0.867)rectangle++(0.04,0.04) (0.291,2.321)rectangle++(0.04,0.04) (0.338,3.558)rectangle++(0.04,0.04)
  (0.428,4.607)rectangle++(0.04,0.04) (0.546,5.495)rectangle++(0.04,0.04) (0.681,6.245)rectangle++(0.04,0.04) (0.827,6.877)rectangle++(0.04,0.04) (0.976,7.409)rectangle++(0.04,0.04)
  (1.123,7.855)rectangle++(0.04,0.04) (1.267,8.228)rectangle++(0.04,0.04) (1.404,8.539)rectangle++(0.04,0.04) (1.533,8.798)rectangle++(0.04,0.04) (1.652,9.013)rectangle++(0.04,0.04)
  (1.763,9.191)rectangle++(0.04,0.04) (1.864,9.338)rectangle++(0.04,0.04) (1.956,9.459)rectangle++(0.04,0.04) (2.038,9.558)rectangle++(0.04,0.04) (2.113,9.639)rectangle++(0.04,0.04)
  (2.179,9.705)rectangle++(0.04,0.04) (2.238,9.758)rectangle++(0.04,0.04) (2.291,9.801)rectangle++(0.04,0.04) (2.337,9.836)rectangle++(0.04,0.04) (2.378,9.863)rectangle++(0.04,0.04)
  (2.413,9.885)rectangle++(0.04,0.04) (2.445,9.902)rectangle++(0.04,0.04) (2.472,9.915)rectangle++(0.04,0.04) (2.495,9.925)rectangle++(0.04,0.04) (2.516,9.933)rectangle++(0.04,0.04)
  (2.534,9.938)rectangle++(0.04,0.04) (2.549,9.942)rectangle++(0.04,0.04) (2.562,9.945)rectangle++(0.04,0.04) (2.573,9.947)rectangle++(0.04,0.04) (2.583,9.948)rectangle++(0.04,0.04)
  (2.591,9.949)rectangle++(0.04,0.04) (2.598,9.949)rectangle++(0.04,0.04) (2.604,9.949)rectangle++(0.04,0.04) (2.609,9.949)rectangle++(0.04,0.04) (2.614,9.948)rectangle++(0.04,0.04)
  (2.617,9.948)rectangle++(0.04,0.04) (2.621,9.947)rectangle++(0.04,0.04) (2.623,9.946)rectangle++(0.04,0.04) (2.625,9.946)rectangle++(0.04,0.04) (2.627,9.945)rectangle++(0.04,0.04)
  (2.629,9.944)rectangle++(0.04,0.04) (2.63,9.944)rectangle++(0.04,0.04) (2.631,9.943)rectangle++(0.04,0.04) (2.632,9.943)rectangle++(0.04,0.04) (2.633,9.942)rectangle++(0.04,0.04)
  (2.633,9.942)rectangle++(0.04,0.04) (2.634,9.941)rectangle++(0.04,0.04) (2.371,3.501)rectangle++(0.04,0.04) (2.153,4.477)rectangle++(0.04,0.04) (2.007,5.315)rectangle++(0.04,0.04)
  (1.917,6.034)rectangle++(0.04,0.04) (1.868,6.648)rectangle++(0.04,0.04) (1.852,7.173)rectangle++(0.04,0.04) (1.713,2.633)rectangle++(0.04,0.04) (1.559,3.766)rectangle++(0.04,0.04)
  (1.474,4.735)rectangle++(0.04,0.04) (1.44,5.562)rectangle++(0.04,0.04) (1.444,6.266)rectangle++(0.04,0.04) (1.476,6.865)rectangle++(0.04,0.04) (1.527,7.372)rectangle++(0.04,0.04)
  (1.591,7.801)rectangle++(0.04,0.04) (1.928,2.716)rectangle++(0.04,0.04) (1.746,3.828)rectangle++(0.04,0.04) (1.635,4.78)rectangle++(0.04,0.04) (1.578,5.594)rectangle++(0.04,0.04)
  (1.563,6.288)rectangle++(0.04,0.04) (1.578,6.878)rectangle++(0.04,0.04) (1.614,7.38)rectangle++(0.04,0.04) (-1.617,3.584)rectangle++(0.04,0.04) (-1.233,4.707)rectangle++(0.04,0.04)
  (-0.862,5.647)rectangle++(0.04,0.04) (-0.509,6.43)rectangle++(0.04,0.04) (-0.178,7.082)rectangle++(0.04,0.04) (0.13,7.623)rectangle++(0.04,0.04) (0.413,8.071)rectangle++(0.04,0.04)
  (0.672,8.44)rectangle++(0.04,0.04) (-0.02,1.334)rectangle++(0.04,0.04) (0.034,2.731)rectangle++(0.04,0.04) (-0.724,2.198)rectangle++(0.04,0.04) (-0.737,1.906)rectangle++(0.04,0.04)
  (-0.553,3.246)rectangle++(0.04,0.04) (-0.342,4.377)rectangle++(0.04,0.04) (-0.118,5.33)rectangle++(0.04,0.04) (1.493,1.679)rectangle++(0.04,0.04) (1.334,2.963)rectangle++(0.04,0.04)
  (1.25,4.062)rectangle++(0.04,0.04) (1.223,4.999)rectangle++(0.04,0.04) (1.237,5.796)rectangle++(0.04,0.04) (1.281,6.473)rectangle++(0.04,0.04) (1.346,7.047)rectangle++(0.04,0.04)
  (1.424,7.533)rectangle++(0.04,0.04) (1.509,7.942)rectangle++(0.04,0.04) (1.598,8.287)rectangle++(0.04,0.04) (1.688,8.576)rectangle++(0.04,0.04) (1.775,8.818)rectangle++(0.04,0.04)
  (1.86,9.021)rectangle++(0.04,0.04) (1.939,9.189)rectangle++(0.04,0.04) (2.014,9.33)rectangle++(0.04,0.04) (-2.044,4.105)rectangle++(0.04,0.04) (-1.576,5.167)rectangle++(0.04,0.04)
  (-1.135,6.051)rectangle++(0.04,0.04) (1.847,1.587)rectangle++(0.04,0.04) (1.631,2.871)rectangle++(0.04,0.04) (0.542,1.543)rectangle++(0.04,0.04) (0.333,0.941)rectangle++(0.04,0.04)
  (0.196,0.743)rectangle++(0.04,0.04) (0.194,2.22)rectangle++(0.04,0.04) (0.252,3.475)rectangle++(0.04,0.04) (0.351,4.54)rectangle++(0.04,0.04) (0.478,5.441)rectangle++(0.04,0.04)
  (1.434,1.86)rectangle++(0.04,0.04) (1.292,3.12)rectangle++(0.04,0.04) (1.22,4.196)rectangle++(0.04,0.04) (1.203,5.114)rectangle++(0.04,0.04) (1.225,5.895)rectangle++(0.04,0.04)
  (1.275,6.558)rectangle++(0.04,0.04) (1.344,7.12)rectangle++(0.04,0.04) (1.425,7.594)rectangle++(0.04,0.04) (1.513,7.994)rectangle++(0.04,0.04) (1.603,8.331)rectangle++(0.04,0.04)
  (1.694,8.613)rectangle++(0.04,0.04) (1.782,8.85)rectangle++(0.04,0.04) (1.867,9.047)rectangle++(0.04,0.04) (1.946,9.212)rectangle++(0.04,0.04) (2.021,9.348)rectangle++(0.04,0.04)
  (2.089,9.461)rectangle++(0.04,0.04) (2.152,9.555)rectangle++(0.04,0.04) (2.209,9.632)rectangle++(0.04,0.04) (2.261,9.695)rectangle++(0.04,0.04) (2.307,9.746)rectangle++(0.04,0.04)
  (2.349,9.788)rectangle++(0.04,0.04) (2.386,9.822)rectangle++(0.04,0.04) (2.419,9.85)rectangle++(0.04,0.04) (2.448,9.872)rectangle++(0.04,0.04) (2.473,9.889)rectangle++(0.04,0.04)
  (2.381,3.446)rectangle++(0.04,0.04) (2.159,4.43)rectangle++(0.04,0.04) (2.01,5.276)rectangle++(0.04,0.04) (-0.991,3.212)rectangle++(0.04,0.04) (-0.716,4.366)rectangle++(0.04,0.04)
  (-0.436,5.336)rectangle++(0.04,0.04) (-0.159,6.149)rectangle++(0.04,0.04) (0.108,6.829)rectangle++(0.04,0.04) (1.879,2.097)rectangle++(0.04,0.04) (0.288,1.422)rectangle++(0.04,0.04)
  (0.299,2.793)rectangle++(0.04,0.04) (0.364,3.958)rectangle++(0.04,0.04) (-0.02,0.617)rectangle++(0.04,0.04) (0.005,2.121)rectangle++(0.04,0.04) (0.087,3.399)rectangle++(0.04,0.04)
  (0.208,4.482)rectangle++(0.04,0.04) (0.354,5.397)rectangle++(0.04,0.04) (0.514,6.17)rectangle++(0.04,0.04) (0.682,6.82)rectangle++(0.04,0.04) (0.85,7.366)rectangle++(0.04,0.04)
  (1.015,7.823)rectangle++(0.04,0.04) (1.174,8.205)rectangle++(0.04,0.04) (1.324,8.524)rectangle++(0.04,0.04) (2.171,2.82)rectangle++(0.04,0.04) (1.956,3.906)rectangle++(0.04,0.04)
  (1.816,4.838)rectangle++(0.04,0.04) (1.735,5.636)rectangle++(0.04,0.04) (1.698,6.317)rectangle++(0.04,0.04) (1.694,6.898)rectangle++(0.04,0.04) (1.714,7.392)rectangle++(0.04,0.04)
  (1.75,7.811)rectangle++(0.04,0.04) (1.798,8.165)rectangle++(0.04,0.04) (1.853,8.465)rectangle++(0.04,0.04) (1.911,8.717)rectangle++(0.04,0.04) (1.971,8.929)rectangle++(0.04,0.04)
  (2.03,9.107)rectangle++(0.04,0.04) (2.088,9.256)rectangle++(0.04,0.04) (2.143,9.38)rectangle++(0.04,0.04) (2.194,9.484)rectangle++(0.04,0.04) (2.242,9.57)rectangle++(0.04,0.04)
  (2.326,3.31)rectangle++(0.04,0.04) (2.107,4.316)rectangle++(0.04,0.04) (1.961,5.181)rectangle++(0.04,0.04) (1.872,5.921)rectangle++(0.04,0.04) (-1.186,3.322)rectangle++(0.04,0.04)
  (-0.878,4.468)rectangle++(0.04,0.04) (-0.57,5.429)rectangle++(0.04,0.04) (-0.269,6.233)rectangle++(0.04,0.04) (0.018,6.905)rectangle++(0.04,0.04) (0.29,7.465)rectangle++(0.04,0.04)
  (0.543,7.93)rectangle++(0.04,0.04) (0.776,8.315)rectangle++(0.04,0.04) (0.99,8.633)rectangle++(0.04,0.04) (1.185,8.895)rectangle++(0.04,0.04) (1.361,9.109)rectangle++(0.04,0.04)
  (1.519,9.285)rectangle++(0.04,0.04) (1.66,9.427)rectangle++(0.04,0.04) (1.786,9.543)rectangle++(0.04,0.04) (1.898,9.636)rectangle++(0.04,0.04) (1.996,9.711)rectangle++(0.04,0.04)
  (2.083,9.771)rectangle++(0.04,0.04) (2.159,9.818)rectangle++(0.04,0.04) (2.226,9.855)rectangle++(0.04,0.04) (2.284,9.884)rectangle++(0.04,0.04) (2.335,9.906)rectangle++(0.04,0.04)
  (2.378,9.923)rectangle++(0.04,0.04) (2.416,9.936)rectangle++(0.04,0.04) (2.449,9.945)rectangle++(0.04,0.04) (-2.117,4.34)rectangle++(0.04,0.04) (-1.628,5.37)rectangle++(0.04,0.04)
  (-1.171,6.226)rectangle++(0.04,0.04) (-0.749,6.935)rectangle++(0.04,0.04) (-0.361,7.521)rectangle++(0.04,0.04) (-0.008,8.003)rectangle++(0.04,0.04) (0.311,8.399)rectangle++(0.04,0.04)
  (0.598,8.723)rectangle++(0.04,0.04) (0.855,8.987)rectangle++(0.04,0.04) (2.371,2.809)rectangle++(0.04,0.04) (2.125,3.889)rectangle++(0.04,0.04) (1.96,4.817)rectangle++(0.04,0.04)
  (1.856,5.612)rectangle++(0.04,0.04) (1.8,6.292)rectangle++(0.04,0.04) (1.78,6.873)rectangle++(0.04,0.04) (1.64,2.542)rectangle++(0.04,0.04) (1.493,3.691)rectangle++(0.04,0.04)
  (1.415,4.674)rectangle++(0.04,0.04) (1.387,5.513)rectangle++(0.04,0.04) (1.398,6.226)rectangle++(0.04,0.04) (1.435,6.833)rectangle++(0.04,0.04) (1.491,7.347)rectangle++(0.04,0.04)
  (1.559,7.781)rectangle++(0.04,0.04) (1.634,8.148)rectangle++(0.04,0.04) (1.713,8.457)rectangle++(0.04,0.04) (1.792,8.716)rectangle++(0.04,0.04) (1.869,8.933)rectangle++(0.04,0.04)
  (-1.97,3.984)rectangle++(0.04,0.04) (-1.517,5.062)rectangle++(0.04,0.04) (-1.089,5.959)rectangle++(0.04,0.04) (-0.69,6.705)rectangle++(0.04,0.04) (-0.32,7.323)rectangle++(0.04,0.04)
  (2.081,2.104)rectangle++(0.04,0.04) (0.26,1.476)rectangle++(0.04,0.04) (0.278,2.841)rectangle++(0.04,0.04) (0.347,4)rectangle++(0.04,0.04) (0.453,4.982)rectangle++(0.04,0.04)
  (0.582,5.813)rectangle++(0.04,0.04) (0.725,6.514)rectangle++(0.04,0.04) (0.875,7.104)rectangle++(0.04,0.04) (1.025,7.599)rectangle++(0.04,0.04) (1.173,8.015)rectangle++(0.04,0.04)
  (1.316,8.362)rectangle++(0.04,0.04) (1.451,8.651)rectangle++(0.04,0.04) (1.577,8.892)rectangle++(0.04,0.04) (1.694,9.091)rectangle++(0.04,0.04) (1.801,9.256)rectangle++(0.04,0.04)
  (1.899,9.392)rectangle++(0.04,0.04) (1.988,9.503)rectangle++(0.04,0.04) (2.067,9.594)rectangle++(0.04,0.04) (2.139,9.669)rectangle++(0.04,0.04) (2.369,3.307)rectangle++(0.04,0.04)
  (2.144,4.312)rectangle++(0.04,0.04) (1.992,5.176)rectangle++(0.04,0.04) (1.898,5.916)rectangle++(0.04,0.04) (1.848,6.549)rectangle++(0.04,0.04) (1.831,7.089)rectangle++(0.04,0.04)
  (1.837,7.548)rectangle++(0.04,0.04) (-1.616,3.672)rectangle++(0.04,0.04) (-1.229,4.782)rectangle++(0.04,0.04) (-0.856,5.71)rectangle++(0.04,0.04) (-0.501,6.484)rectangle++(0.04,0.04)
  (1.873,1.856)rectangle++(0.04,0.04) (1.664,3.099)rectangle++(0.04,0.04) (1.536,4.164)rectangle++(0.04,0.04) (1.47,5.074)rectangle++(0.04,0.04) (1.451,5.85)rectangle++(0.04,0.04)
  (1.403,2.211)rectangle++(0.04,0.04) (1.279,3.42)rectangle++(0.04,0.04) (1.222,4.452)rectangle++(0.04,0.04) (1.214,5.331)rectangle++(0.04,0.04) (1.243,6.079)rectangle++(0.04,0.04)
  (1.298,6.714)rectangle++(0.04,0.04) (-1.507,3.364)rectangle++(0.04,0.04) (-1.149,4.516)rectangle++(0.04,0.04) (-0.798,5.481)rectangle++(0.04,0.04) (-0.461,6.287)rectangle++(0.04,0.04)
  (-0.143,6.959)rectangle++(0.04,0.04) (0.155,7.517)rectangle++(0.04,0.04) (0.43,7.979)rectangle++(0.04,0.04) (0.682,8.361)rectangle++(0.04,0.04) (0.912,8.676)rectangle++(0.04,0.04)
  (1.12,8.934)rectangle++(0.04,0.04) (1.307,9.146)rectangle++(0.04,0.04) (1.475,9.318)rectangle++(0.04,0.04) (1.624,9.457)rectangle++(0.04,0.04) (1.757,9.57)rectangle++(0.04,0.04)
  (1.874,9.66)rectangle++(0.04,0.04) (1.977,9.733)rectangle++(0.04,0.04) (2.067,9.79)rectangle++(0.04,0.04) (2.147,9.835)rectangle++(0.04,0.04) (2.216,9.87)rectangle++(0.04,0.04)
  (-2.144,4.27)rectangle++(0.04,0.04) (-1.654,5.312)rectangle++(0.04,0.04) (-1.196,6.177)rectangle++(0.04,0.04) (-0.771,6.895)rectangle++(0.04,0.04) (2.029,1.884)rectangle++(0.04,0.04)
  (1.798,3.117)rectangle++(0.04,0.04) (0.586,1.645)rectangle++(0.04,0.04) (0.561,2.971)rectangle++(0.04,0.04) (0.594,4.099)rectangle++(0.04,0.04) (0.667,5.057)rectangle++(0.04,0.04)
  (0.767,5.868)rectangle++(0.04,0.04) (0.884,6.553)rectangle++(0.04,0.04) (1.011,7.131)rectangle++(0.04,0.04) (1.143,7.617)rectangle++(0.04,0.04) (1.274,8.025)rectangle++(0.04,0.04)
  (1.402,8.367)rectangle++(0.04,0.04) (1.524,8.652)rectangle++(0.04,0.04) (1.639,8.889)rectangle++(0.04,0.04) (1.747,9.086)rectangle++(0.04,0.04) (-2.034,3.99)rectangle++(0.04,0.04)
  (-1.572,5.069)rectangle++(0.04,0.04) (-1.135,5.968)rectangle++(0.04,0.04) (-0.729,6.714)rectangle++(0.04,0.04) (-0.353,7.332)rectangle++(0.04,0.04) (-0.009,7.843)rectangle++(0.04,0.04)
  (0.304,8.263)rectangle++(0.04,0.04) (0.587,8.608)rectangle++(0.04,0.04) (0.841,8.889)rectangle++(0.04,0.04) (1.068,9.118)rectangle++(0.04,0.04) (1.27,9.304)rectangle++(0.04,0.04)
  (1.45,9.454)rectangle++(0.04,0.04) (1.608,9.574)rectangle++(0.04,0.04) (-2.189,4.065)rectangle++(0.04,0.04) (-1.7,5.139)rectangle++(0.04,0.04) (-1.242,6.032)rectangle++(0.04,0.04)
  (-0.816,6.773)rectangle++(0.04,0.04) (-0.425,7.386)rectangle++(0.04,0.04) (-0.068,7.892)rectangle++(0.04,0.04) (0.256,8.307)rectangle++(0.04,0.04) (0.547,8.647)rectangle++(0.04,0.04)
  (0.809,8.924)rectangle++(0.04,0.04) (1.042,9.149)rectangle++(0.04,0.04) (1.25,9.331)rectangle++(0.04,0.04) (1.433,9.478)rectangle++(0.04,0.04) (1.595,9.595)rectangle++(0.04,0.04)
  (1.738,9.688)rectangle++(0.04,0.04) (1.862,9.762)rectangle++(0.04,0.04) (1.971,9.819)rectangle++(0.04,0.04) (2.066,9.864)rectangle++(0.04,0.04) (2.149,9.898)rectangle++(0.04,0.04)
  (2.22,9.923)rectangle++(0.04,0.04) (2.282,9.942)rectangle++(0.04,0.04) (2.335,9.956)rectangle++(0.04,0.04) (2.381,9.965)rectangle++(0.04,0.04) (2.416,3.441)rectangle++(0.04,0.04)
  (2.189,4.424)rectangle++(0.04,0.04) (2.035,5.269)rectangle++(0.04,0.04) (1.939,5.994)rectangle++(0.04,0.04) (1.885,6.613)rectangle++(0.04,0.04) (1.865,7.142)rectangle++(0.04,0.04)
  (1.869,7.592)rectangle++(0.04,0.04) (1.89,7.975)rectangle++(0.04,0.04) (1.923,8.299)rectangle++(0.04,0.04) (1.964,8.574)rectangle++(0.04,0.04) (2.011,8.805)rectangle++(0.04,0.04)
  (2.059,9)rectangle++(0.04,0.04) (2.108,9.164)rectangle++(0.04,0.04) (2.156,9.301)rectangle++(0.04,0.04) (2.203,9.416)rectangle++(0.04,0.04) (2.247,9.512)rectangle++(0.04,0.04)
  (2.288,9.591)rectangle++(0.04,0.04) (2.326,9.657)rectangle++(0.04,0.04) (2.361,9.712)rectangle++(0.04,0.04) (2.393,9.757)rectangle++(0.04,0.04) (2.422,9.794)rectangle++(0.04,0.04)
  (2.361,3.41)rectangle++(0.04,0.04) (2.141,4.401)rectangle++(0.04,0.04) (1.994,5.251)rectangle++(0.04,0.04) (1.903,5.98)rectangle++(0.04,0.04) (1.854,6.603)rectangle++(0.04,0.04)
  (1.838,7.134)rectangle++(0.04,0.04) (-1.509,3.581)rectangle++(0.04,0.04) (1.212,0.897)rectangle++(0.04,0.04) (1.064,2.31)rectangle++(0.04,0.04) (0.994,3.518)rectangle++(0.04,0.04)
  (0.984,4.546)rectangle++(0.04,0.04) (1.016,5.421)rectangle++(0.04,0.04) (1.078,6.164)rectangle++(0.04,0.04) (1.161,6.792)rectangle++(0.04,0.04) (1.256,7.323)rectangle++(0.04,0.04)
  (1.358,7.771)rectangle++(0.04,0.04) (1.463,8.147)rectangle++(0.04,0.04) (1.567,8.463)rectangle++(0.04,0.04) (1.669,8.727)rectangle++(0.04,0.04) (1.765,8.947)rectangle++(0.04,0.04)
  (1.856,9.131)rectangle++(0.04,0.04) (1.941,9.283)rectangle++(0.04,0.04) (2.019,9.409)rectangle++(0.04,0.04) (2.09,9.513)rectangle++(0.04,0.04) (2.155,9.599)rectangle++(0.04,0.04)
  (2.213,9.669)rectangle++(0.04,0.04) (2.266,9.726)rectangle++(0.04,0.04) (2.313,9.773)rectangle++(0.04,0.04) (2.355,9.811)rectangle++(0.04,0.04) (2.392,9.841)rectangle++(0.04,0.04)
  (2.424,9.865)rectangle++(0.04,0.04) (2.453,9.885)rectangle++(0.04,0.04) (-2.101,4.328)rectangle++(0.04,0.04) (-1.615,5.359)rectangle++(0.04,0.04) (-1.16,6.216)rectangle++(0.04,0.04)
  (-0.74,6.926)rectangle++(0.04,0.04) (-0.354,7.513)rectangle++(0.04,0.04) (-0.003,7.996)rectangle++(0.04,0.04) (0.315,8.393)rectangle++(0.04,0.04) (0.602,8.718)rectangle++(0.04,0.04)
  (0.858,8.982)rectangle++(0.04,0.04) (1.086,9.197)rectangle++(0.04,0.04) (1.289,9.37)rectangle++(0.04,0.04) (1.468,9.509)rectangle++(0.04,0.04) (1.626,9.62)rectangle++(0.04,0.04)
  (1.765,9.708)rectangle++(0.04,0.04) (1.886,9.778)rectangle++(0.04,0.04) (2.437,3.267)rectangle++(0.04,0.04) (2.2,4.276)rectangle++(0.04,0.04) (0.85,2.028)rectangle++(0.04,0.04)
  (0.801,3.286)rectangle++(0.04,0.04) (0.81,4.357)rectangle++(0.04,0.04) (1.081,1.686)rectangle++(0.04,0.04) (0.984,2.986)rectangle++(0.04,0.04) (-0.601,2.472)rectangle++(0.04,0.04)
  (-0.414,3.722)rectangle++(0.04,0.04) (-0.205,4.776)rectangle++(0.04,0.04) (0.014,5.664)rectangle++(0.04,0.04) (0.237,6.41)rectangle++(0.04,0.04) (0.455,7.035)rectangle++(0.04,0.04)
  (0.666,7.558)rectangle++(0.04,0.04) (0.866,7.994)rectangle++(0.04,0.04) (1.054,8.356)rectangle++(0.04,0.04) (1.228,8.657)rectangle++(0.04,0.04) (1.388,8.906)rectangle++(0.04,0.04)
  (1.534,9.11)rectangle++(0.04,0.04) (1.666,9.279)rectangle++(0.04,0.04) (1.785,9.416)rectangle++(0.04,0.04) (1.892,9.529)rectangle++(0.04,0.04) (1.987,9.62)rectangle++(0.04,0.04)
  (2.071,9.694)rectangle++(0.04,0.04) (2.146,9.753)rectangle++(0.04,0.04) (2.212,9.8)rectangle++(0.04,0.04) (2.27,9.838)rectangle++(0.04,0.04) (2.321,9.868)rectangle++(0.04,0.04)
  (2.365,9.891)rectangle++(0.04,0.04) (2.404,9.909)rectangle++(0.04,0.04) (2.438,9.923)rectangle++(0.04,0.04) (2.467,9.933)rectangle++(0.04,0.04) (2.394,3.455)rectangle++(0.04,0.04)
  (0.591,1.882)rectangle++(0.04,0.04) (0.575,3.172)rectangle++(0.04,0.04) (0.614,4.269)rectangle++(0.04,0.04) (0.69,5.201)rectangle++(0.04,0.04) (0.793,5.989)rectangle++(0.04,0.04)
  (0.911,6.655)rectangle++(0.04,0.04) (1.038,7.217)rectangle++(0.04,0.04) (1.169,7.689)rectangle++(0.04,0.04) (1.299,8.085)rectangle++(0.04,0.04) (-1.863,3.666)rectangle++(0.04,0.04)
  (-1.439,4.787)rectangle++(0.04,0.04) (-1.034,5.723)rectangle++(0.04,0.04) (-0.652,6.502)rectangle++(0.04,0.04) (-0.297,7.149)rectangle++(0.04,0.04) (0.032,7.685)rectangle++(0.04,0.04)
  (0.332,8.127)rectangle++(0.04,0.04) (2.208,2.467)rectangle++(0.04,0.04) (1.974,3.605)rectangle++(0.04,0.04) (1.82,4.581)rectangle++(0.04,0.04) (1.728,5.417)rectangle++(0.04,0.04)
  (1.683,6.132)rectangle++(0.04,0.04) (-0.02,0.964)rectangle++(0.04,0.04) (0.019,2.417)rectangle++(0.04,0.04) (0.111,3.65)rectangle++(0.04,0.04) (0.238,4.694)rectangle++(0.04,0.04)
  (0.388,5.577)rectangle++(0.04,0.04) (0.551,6.321)rectangle++(0.04,0.04) (0.719,6.947)rectangle++(0.04,0.04) (0.887,7.472)rectangle++(0.04,0.04) (1.05,7.912)rectangle++(0.04,0.04)
  (1.207,8.279)rectangle++(0.04,0.04) (1.355,8.585)rectangle++(0.04,0.04) (1.493,8.84)rectangle++(0.04,0.04) (1.62,9.05)rectangle++(0.04,0.04) (1.737,9.224)rectangle++(0.04,0.04)
  (1.843,9.367)rectangle++(0.04,0.04) (1.939,9.485)rectangle++(0.04,0.04) (2.026,9.581)rectangle++(0.04,0.04) (2.103,9.659)rectangle++(0.04,0.04) (2.171,9.722)rectangle++(0.04,0.04)
  (2.232,9.773)rectangle++(0.04,0.04) (2.384,3.356)rectangle++(0.04,0.04) (2.159,4.353)rectangle++(0.04,0.04) (2.007,5.21)rectangle++(0.04,0.04) (1.912,5.945)rectangle++(0.04,0.04)
  (1.861,6.573)rectangle++(0.04,0.04) (1.842,7.109)rectangle++(0.04,0.04) (1.848,7.565)rectangle++(0.04,0.04) (1.871,7.952)rectangle++(0.04,0.04) (1.906,8.281)rectangle++(0.04,0.04)
  (2.015,2.913)rectangle++(0.04,0.04) (0.496,1.653)rectangle++(0.04,0.04) (0.486,2.982)rectangle++(0.04,0.04) (0.53,4.111)rectangle++(0.04,0.04) (0.613,5.069)rectangle++(0.04,0.04)
  (0.721,5.881)rectangle++(0.04,0.04) (0.846,6.566)rectangle++(0.04,0.04) (0.98,7.143)rectangle++(0.04,0.04) (1.116,7.629)rectangle++(0.04,0.04) (1.252,8.036)rectangle++(0.04,0.04)
  (-1.86,3.645)rectangle++(0.04,0.04) (-1.438,4.769)rectangle++(0.04,0.04) (-1.549,2.307)rectangle++(0.04,0.04) (-1.226,3.619)rectangle++(0.04,0.04) (-1.208,2.103)rectangle++(0.04,0.04)
  (-0.944,3.432)rectangle++(0.04,0.04) (-0.668,4.551)rectangle++(0.04,0.04) (-0.388,5.492)rectangle++(0.04,0.04) (-0.112,6.28)rectangle++(0.04,0.04) (0.154,6.938)rectangle++(0.04,0.04)
  (0.406,7.488)rectangle++(0.04,0.04) (0.642,7.944)rectangle++(0.04,0.04) (-1.958,3.485)rectangle++(0.04,0.04) (-1.527,4.636)rectangle++(0.04,0.04) (-1.115,5.598)rectangle++(0.04,0.04)
  (-0.726,6.399)rectangle++(0.04,0.04) (-0.363,7.065)rectangle++(0.04,0.04) (2.015,2.031)rectangle++(0.04,0.04) (1.792,3.242)rectangle++(0.04,0.04) (1.651,4.28)rectangle++(0.04,0.04)
  (1.572,5.168)rectangle++(0.04,0.04) (1.541,5.926)rectangle++(0.04,0.04) (-0.02,0.931)rectangle++(0.04,0.04) (0.018,2.389)rectangle++(0.04,0.04) (0.109,3.626)rectangle++(0.04,0.04)
  (0.235,4.674)rectangle++(0.04,0.04) (0.385,5.56)rectangle++(0.04,0.04) (0.547,6.306)rectangle++(0.04,0.04) (0.715,6.935)rectangle++(0.04,0.04) (0.883,7.462)rectangle++(0.04,0.04)
  (1.047,7.904)rectangle++(0.04,0.04) (1.204,8.272)rectangle++(0.04,0.04) (1.352,8.58)rectangle++(0.04,0.04) (2.182,2.841)rectangle++(0.04,0.04) (1.966,3.923)rectangle++(0.04,0.04)
  (0.786,1.883)rectangle++(0.04,0.04) (0.741,3.165)rectangle++(0.04,0.04) (0.755,4.257)rectangle++(0.04,0.04) (0.809,5.184)rectangle++(0.04,0.04) (0.893,5.971)rectangle++(0.04,0.04)
  (0.996,6.635)rectangle++(0.04,0.04) (1.11,7.197)rectangle++(0.04,0.04) (1.229,7.669)rectangle++(0.04,0.04) (1.349,8.066)rectangle++(0.04,0.04) (1.467,8.398)rectangle++(0.04,0.04)
  (1.581,8.676)rectangle++(0.04,0.04) (1.689,8.907)rectangle++(0.04,0.04) (1.789,9.1)rectangle++(0.04,0.04) (1.883,9.26)rectangle++(0.04,0.04) (-2.052,4.059)rectangle++(0.04,0.04)
  (-1.487,2.01)rectangle++(0.04,0.04) (0.768,0.526)rectangle++(0.04,0.04) (0.672,2.012)rectangle++(0.04,0.04) (0.65,3.28)rectangle++(0.04,0.04) (0.681,4.358)rectangle++(0.04,0.04)
  (1.101,1.653)rectangle++(0.04,0.04) (1,2.957)rectangle++(0.04,0.04) (0.966,4.07)rectangle++(0.04,0.04) (0.981,5.017)rectangle++(0.04,0.04) (1.033,5.821)rectangle++(0.04,0.04)
  (1.108,6.503)rectangle++(0.04,0.04) (1.2,7.079)rectangle++(0.04,0.04) (1.301,7.566)rectangle++(0.04,0.04) (1.406,7.975)rectangle++(0.04,0.04) (1.512,8.319)rectangle++(0.04,0.04)
  (-1.882,3.767)rectangle++(0.04,0.04) (-1.451,4.873)rectangle++(0.04,0.04) (-1.041,5.797)rectangle++(0.04,0.04) (-0.655,6.565)rectangle++(0.04,0.04) (-0.296,7.203)rectangle++(0.04,0.04);
\fill[nuage]
  (0.034,7.73)rectangle++(0.04,0.04) (0.336,8.166)rectangle++(0.04,0.04) (0.61,8.523)rectangle++(0.04,0.04) (0.857,8.817)rectangle++(0.04,0.04) (1.079,9.056)rectangle++(0.04,0.04)
  (1.277,9.251)rectangle++(0.04,0.04) (1.454,9.408)rectangle++(0.04,0.04) (1.61,9.535)rectangle++(0.04,0.04) (1.747,9.637)rectangle++(0.04,0.04) (2.419,3.197)rectangle++(0.04,0.04)
  (2.182,4.217)rectangle++(0.04,0.04) (2.021,5.093)rectangle++(0.04,0.04) (1.919,5.845)rectangle++(0.04,0.04) (1.863,6.487)rectangle++(0.04,0.04) (1.841,7.036)rectangle++(0.04,0.04)
  (1.844,7.503)rectangle++(0.04,0.04) (1.865,7.9)rectangle++(0.04,0.04) (1.899,8.237)rectangle++(0.04,0.04) (1.942,8.521)rectangle++(0.04,0.04) (1.989,8.762)rectangle++(0.04,0.04)
  (2.039,8.964)rectangle++(0.04,0.04) (2.089,9.134)rectangle++(0.04,0.04) (2.139,9.277)rectangle++(0.04,0.04) (2.187,9.396)rectangle++(0.04,0.04) (2.233,9.495)rectangle++(0.04,0.04)
  (-2.043,4.191)rectangle++(0.04,0.04) (1.463,0.905)rectangle++(0.04,0.04) (1.277,2.307)rectangle++(0.04,0.04) (1.176,3.506)rectangle++(0.04,0.04) (-0.698,2.631)rectangle++(0.04,0.04)
  (-0.49,3.86)rectangle++(0.04,0.04) (-0.264,4.897)rectangle++(0.04,0.04) (-0.031,5.769)rectangle++(0.04,0.04) (0.202,6.501)rectangle++(0.04,0.04) (0.43,7.114)rectangle++(0.04,0.04)
  (0.648,7.626)rectangle++(0.04,0.04) (2.021,2.429)rectangle++(0.04,0.04) (1.813,3.58)rectangle++(0.04,0.04) (0.713,1.76)rectangle++(0.04,0.04) (0.674,3.064)rectangle++(0.04,0.04)
  (0.694,4.174)rectangle++(0.04,0.04) (0.754,5.116)rectangle++(0.04,0.04) (-1.201,2.888)rectangle++(0.04,0.04) (-0.907,4.099)rectangle++(0.04,0.04) (-1.268,2.282)rectangle++(0.04,0.04)
  (-0.989,3.587)rectangle++(0.04,0.04) (-0.699,4.685)rectangle++(0.04,0.04) (-0.409,5.606)rectangle++(0.04,0.04) (-1.561,2.728)rectangle++(0.04,0.04) (-1.22,3.978)rectangle++(0.04,0.04)
  (-0.88,5.026)rectangle++(0.04,0.04) (-0.549,5.903)rectangle++(0.04,0.04) (-0.233,6.636)rectangle++(0.04,0.04) (0.065,7.246)rectangle++(0.04,0.04) (0.343,7.753)rectangle++(0.04,0.04)
  (0.6,8.172)rectangle++(0.04,0.04) (0.834,8.519)rectangle++(0.04,0.04) (1.048,8.804)rectangle++(0.04,0.04) (1.241,9.037)rectangle++(0.04,0.04) (1.414,9.228)rectangle++(0.04,0.04)
  (2.354,3.012)rectangle++(0.04,0.04) (2.12,4.063)rectangle++(0.04,0.04) (1.962,4.965)rectangle++(0.04,0.04) (1.864,5.738)rectangle++(0.04,0.04) (1.812,6.399)rectangle++(0.04,0.04)
  (1.794,6.963)rectangle++(0.04,0.04) (1.801,7.443)rectangle++(0.04,0.04) (1.826,7.85)rectangle++(0.04,0.04) (1.864,8.196)rectangle++(0.04,0.04) (1.91,8.488)rectangle++(0.04,0.04)
  (1.961,8.735)rectangle++(0.04,0.04) (2.014,8.942)rectangle++(0.04,0.04) (2.067,9.117)rectangle++(0.04,0.04) (2.12,9.263)rectangle++(0.04,0.04) (2.17,9.385)rectangle++(0.04,0.04)
  (-2.027,4.153)rectangle++(0.04,0.04) (-1.559,5.207)rectangle++(0.04,0.04) (-1.119,6.085)rectangle++(0.04,0.04) (-0.71,6.813)rectangle++(0.04,0.04) (-0.333,7.416)rectangle++(0.04,0.04)
  (0.011,7.913)rectangle++(0.04,0.04) (-2.076,3.332)rectangle++(0.04,0.04) (-1.303,1.845)rectangle++(0.04,0.04) (-1.036,3.216)rectangle++(0.04,0.04) (-0.754,4.371)rectangle++(0.04,0.04)
  (-0.468,5.342)rectangle++(0.04,0.04) (-0.187,6.156)rectangle++(0.04,0.04) (0.085,6.836)rectangle++(0.04,0.04) (0.344,7.403)rectangle++(0.04,0.04) (0.586,7.875)rectangle++(0.04,0.04)
  (0.811,8.267)rectangle++(0.04,0.04) (1.018,8.591)rectangle++(0.04,0.04) (2.235,2.756)rectangle++(0.04,0.04) (2.008,3.85)rectangle++(0.04,0.04) (1.859,4.788)rectangle++(0.04,0.04)
  (1.769,5.592)rectangle++(0.04,0.04) (1.725,6.278)rectangle++(0.04,0.04) (1.715,6.864)rectangle++(0.04,0.04) (1.73,7.362)rectangle++(0.04,0.04) (1.763,7.785)rectangle++(0.04,0.04)
  (1.898,2.757)rectangle++(0.04,0.04) (-0.358,2.632)rectangle++(0.04,0.04) (-0.777,2.086)rectangle++(0.04,0.04) (-0.579,3.4)rectangle++(0.04,0.04) (-0.359,4.509)rectangle++(0.04,0.04)
  (-0.127,5.444)rectangle++(0.04,0.04) (1.526,1.704)rectangle++(0.04,0.04) (-0.159,2.315)rectangle++(0.04,0.04) (-0.045,3.57)rectangle++(0.04,0.04) (0.103,4.633)rectangle++(0.04,0.04)
  (0.27,5.53)rectangle++(0.04,0.04) (0.449,6.286)rectangle++(0.04,0.04) (0.631,6.921)rectangle++(0.04,0.04) (0.811,7.454)rectangle++(0.04,0.04) (0.985,7.9)rectangle++(0.04,0.04)
  (1.151,8.271)rectangle++(0.04,0.04) (1.307,8.581)rectangle++(0.04,0.04) (1.452,8.838)rectangle++(0.04,0.04) (1.586,9.05)rectangle++(0.04,0.04) (1.708,9.225)rectangle++(0.04,0.04)
  (1.818,9.369)rectangle++(0.04,0.04) (1.918,9.488)rectangle++(0.04,0.04) (2.008,9.584)rectangle++(0.04,0.04) (2.088,9.662)rectangle++(0.04,0.04) (-2.116,4.195)rectangle++(0.04,0.04)
  (-1.633,5.246)rectangle++(0.04,0.04) (1.697,1.265)rectangle++(0.04,0.04) (1.49,2.603)rectangle++(0.04,0.04) (1.369,3.749)rectangle++(0.04,0.04) (1.311,4.728)rectangle++(0.04,0.04)
  (1.302,5.563)rectangle++(0.04,0.04) (1.327,6.273)rectangle++(0.04,0.04) (1.376,6.875)rectangle++(0.04,0.04) (1.443,7.385)rectangle++(0.04,0.04) (1.519,7.816)rectangle++(0.04,0.04)
  (1.602,8.179)rectangle++(0.04,0.04) (1.687,8.484)rectangle++(0.04,0.04) (1.771,8.74)rectangle++(0.04,0.04) (1.853,8.954)rectangle++(0.04,0.04) (1.931,9.133)rectangle++(0.04,0.04)
  (2.004,9.282)rectangle++(0.04,0.04) (2.073,9.406)rectangle++(0.04,0.04) (2.136,9.508)rectangle++(0.04,0.04) (2.325,3.267)rectangle++(0.04,0.04) (-0.406,2.842)rectangle++(0.04,0.04)
  (-0.841,2.121)rectangle++(0.04,0.04) (-0.633,3.433)rectangle++(0.04,0.04) (-0.403,4.539)rectangle++(0.04,0.04) (-0.163,5.471)rectangle++(0.04,0.04) (0.078,6.253)rectangle++(0.04,0.04)
  (0.314,6.908)rectangle++(0.04,0.04) (0.541,7.455)rectangle++(0.04,0.04) (1.989,2.36)rectangle++(0.04,0.04) (1.783,3.523)rectangle++(0.04,0.04) (1.654,4.519)rectangle++(0.04,0.04)
  (1.585,5.371)rectangle++(0.04,0.04) (1.559,6.098)rectangle++(0.04,0.04) (1.567,6.718)rectangle++(0.04,0.04) (1.599,7.243)rectangle++(0.04,0.04) (1.646,7.689)rectangle++(0.04,0.04)
  (1.705,8.066)rectangle++(0.04,0.04) (1.77,8.384)rectangle++(0.04,0.04) (1.837,8.652)rectangle++(0.04,0.04) (2.13,2.984)rectangle++(0.04,0.04) (0.499,1.7)rectangle++(0.04,0.04)
  (-0.363,2.078)rectangle++(0.04,0.04) (-0.228,3.377)rectangle++(0.04,0.04) (-0.061,4.476)rectangle++(0.04,0.04) (0.125,5.403)rectangle++(0.04,0.04) (0.32,6.184)rectangle++(0.04,0.04)
  (0.517,6.839)rectangle++(0.04,0.04) (0.711,7.389)rectangle++(0.04,0.04) (0.898,7.848)rectangle++(0.04,0.04) (1.075,8.231)rectangle++(0.04,0.04) (2.126,2.685)rectangle++(0.04,0.04)
  (1.912,3.793)rectangle++(0.04,0.04) (1.775,4.744)rectangle++(0.04,0.04) (1.696,5.558)rectangle++(0.04,0.04) (1.662,6.252)rectangle++(0.04,0.04) (1.661,6.844)rectangle++(0.04,0.04)
  (1.683,7.347)rectangle++(0.04,0.04) (1.722,7.774)rectangle++(0.04,0.04) (1.773,8.135)rectangle++(0.04,0.04) (1.83,8.44)rectangle++(0.04,0.04) (1.891,8.697)rectangle++(0.04,0.04)
  (1.953,8.913)rectangle++(0.04,0.04) (2.014,9.094)rectangle++(0.04,0.04) (2.074,9.246)rectangle++(0.04,0.04) (-2.01,4.1)rectangle++(0.04,0.04) (-1.547,5.162)rectangle++(0.04,0.04)
  (-1.111,6.045)rectangle++(0.04,0.04) (1.842,1.592)rectangle++(0.04,0.04) (1.627,2.876)rectangle++(0.04,0.04) (1.496,3.976)rectangle++(0.04,0.04) (-0.756,2.808)rectangle++(0.04,0.04)
  (-0.532,4.013)rectangle++(0.04,0.04) (-0.294,5.029)rectangle++(0.04,0.04) (-1.387,2.628)rectangle++(0.04,0.04) (-0.982,1.848)rectangle++(0.04,0.04) (0.647,0.618)rectangle++(0.04,0.04)
  (0.573,2.096)rectangle++(0.04,0.04) (0.484,1.082)rectangle++(0.04,0.04) (0.452,2.496)rectangle++(0.04,0.04) (0.482,3.7)rectangle++(0.04,0.04) (0.555,4.722)rectangle++(0.04,0.04)
  (0.659,5.588)rectangle++(0.04,0.04) (0.781,6.319)rectangle++(0.04,0.04) (0.915,6.936)rectangle++(0.04,0.04) (1.053,7.456)rectangle++(0.04,0.04) (1.191,7.891)rectangle++(0.04,0.04)
  (1.326,8.256)rectangle++(0.04,0.04) (1.455,8.561)rectangle++(0.04,0.04) (2.161,2.863)rectangle++(0.04,0.04) (1.95,3.943)rectangle++(0.04,0.04) (1.813,4.87)rectangle++(0.04,0.04)
  (1.733,5.663)rectangle++(0.04,0.04) (1.698,6.341)rectangle++(0.04,0.04) (1.694,6.918)rectangle++(0.04,0.04) (1.715,7.409)rectangle++(0.04,0.04) (1.752,7.825)rectangle++(0.04,0.04)
  (1.8,8.177)rectangle++(0.04,0.04) (1.855,8.475)rectangle++(0.04,0.04) (1.913,8.726)rectangle++(0.04,0.04) (1.973,8.936)rectangle++(0.04,0.04) (2.032,9.113)rectangle++(0.04,0.04)
  (2.09,9.261)rectangle++(0.04,0.04) (2.145,9.385)rectangle++(0.04,0.04) (2.196,9.487)rectangle++(0.04,0.04) (2.31,3.278)rectangle++(0.04,0.04) (2.092,4.29)rectangle++(0.04,0.04)
  (1.948,5.159)rectangle++(0.04,0.04) (1.86,5.904)rectangle++(0.04,0.04) (-1.184,3.316)rectangle++(0.04,0.04) (1.089,0.918)rectangle++(0.04,0.04) (0.96,2.333)rectangle++(0.04,0.04)
  (0.907,3.541)rectangle++(0.04,0.04) (0.91,4.57)rectangle++(0.04,0.04) (1.126,1.763)rectangle++(0.04,0.04) (1.025,3.05)rectangle++(0.04,0.04) (0.683,1.429)rectangle++(0.04,0.04)
  (0.635,2.783)rectangle++(0.04,0.04) (-0.618,2.347)rectangle++(0.04,0.04) (-0.433,3.616)rectangle++(0.04,0.04) (-0.226,4.687)rectangle++(0.04,0.04) (-0.007,5.589)rectangle++(0.04,0.04)
  (0.216,6.348)rectangle++(0.04,0.04) (-1.628,3.035)rectangle++(0.04,0.04) (-1.265,4.241)rectangle++(0.04,0.04) (-0.908,5.252)rectangle++(0.04,0.04) (-0.564,6.097)rectangle++(0.04,0.04)
  (-0.238,6.801)rectangle++(0.04,0.04) (0.068,7.386)rectangle++(0.04,0.04) (0.351,7.872)rectangle++(0.04,0.04) (-1.998,3.402)rectangle++(0.04,0.04) (-1.564,4.567)rectangle++(0.04,0.04)
  (-1.149,5.541)rectangle++(0.04,0.04) (-0.757,6.352)rectangle++(0.04,0.04) (-0.392,7.026)rectangle++(0.04,0.04) (-0.054,7.584)rectangle++(0.04,0.04) (0.255,8.045)rectangle++(0.04,0.04)
  (0.536,8.424)rectangle++(0.04,0.04) (0.791,8.735)rectangle++(0.04,0.04) (1.019,8.989)rectangle++(0.04,0.04) (1.224,9.196)rectangle++(0.04,0.04) (1.406,9.364)rectangle++(0.04,0.04)
  (1.567,9.5)rectangle++(0.04,0.04) (1.71,9.608)rectangle++(0.04,0.04) (1.836,9.695)rectangle++(0.04,0.04) (1.946,9.763)rectangle++(0.04,0.04) (2.042,9.817)rectangle++(0.04,0.04)
  (2.126,9.859)rectangle++(0.04,0.04) (2.2,9.891)rectangle++(0.04,0.04) (2.263,9.916)rectangle++(0.04,0.04) (2.318,9.934)rectangle++(0.04,0.04) (2.366,9.948)rectangle++(0.04,0.04)
  (2.406,9.957)rectangle++(0.04,0.04) (2.442,9.963)rectangle++(0.04,0.04) (2.472,9.967)rectangle++(0.04,0.04) (2.497,9.97)rectangle++(0.04,0.04) (2.519,9.97)rectangle++(0.04,0.04)
  (2.538,9.97)rectangle++(0.04,0.04) (2.394,3.483)rectangle++(0.04,0.04) (2.172,4.461)rectangle++(0.04,0.04) (2.022,5.301)rectangle++(0.04,0.04) (1.929,6.021)rectangle++(0.04,0.04)
  (-1.201,3.357)rectangle++(0.04,0.04) (-0.889,4.498)rectangle++(0.04,0.04) (-0.578,5.455)rectangle++(0.04,0.04) (-0.275,6.256)rectangle++(0.04,0.04) (-1.703,2.902)rectangle++(0.04,0.04)
  (-1.333,4.131)rectangle++(0.04,0.04) (-0.97,5.161)rectangle++(0.04,0.04) (-0.621,6.022)rectangle++(0.04,0.04) (-0.289,6.74)rectangle++(0.04,0.04) (-1.831,3.005)rectangle++(0.04,0.04)
  (-1.439,4.224)rectangle++(0.04,0.04) (-1.056,5.244)rectangle++(0.04,0.04) (-0.69,6.096)rectangle++(0.04,0.04) (-0.345,6.805)rectangle++(0.04,0.04) (-0.023,7.395)rectangle++(0.04,0.04)
  (0.274,7.882)rectangle++(0.04,0.04) (0.546,8.285)rectangle++(0.04,0.04) (0.793,8.617)rectangle++(0.04,0.04) (1.017,8.889)rectangle++(0.04,0.04) (2.319,2.828)rectangle++(0.04,0.04)
  (2.082,3.907)rectangle++(0.04,0.04) (1.924,4.834)rectangle++(0.04,0.04) (1.826,5.628)rectangle++(0.04,0.04) (1.775,6.307)rectangle++(0.04,0.04) (1.759,6.886)rectangle++(0.04,0.04)
  (1.768,7.379)rectangle++(0.04,0.04) (1.796,7.798)rectangle++(0.04,0.04) (1.836,8.152)rectangle++(0.04,0.04) (1.885,8.452)rectangle++(0.04,0.04) (1.938,8.705)rectangle++(0.04,0.04)
  (1.993,8.918)rectangle++(0.04,0.04) (2.049,9.097)rectangle++(0.04,0.04) (2.103,9.247)rectangle++(0.04,0.04) (2.155,9.372)rectangle++(0.04,0.04) (2.205,9.476)rectangle++(0.04,0.04)
  (2.251,9.563)rectangle++(0.04,0.04) (2.294,9.634)rectangle++(0.04,0.04) (2.333,9.694)rectangle++(0.04,0.04) (-2.075,4.258)rectangle++(0.04,0.04) (-1.596,5.299)rectangle++(0.04,0.04)
  (-1.718,2.388)rectangle++(0.04,0.04) (-1.367,3.694)rectangle++(0.04,0.04) (-1.016,4.791)rectangle++(0.04,0.04) (-0.674,5.709)rectangle++(0.04,0.04) (-0.347,6.476)rectangle++(0.04,0.04)
  (-0.038,7.115)rectangle++(0.04,0.04) (0.25,7.645)rectangle++(0.04,0.04) (0.516,8.085)rectangle++(0.04,0.04) (0.76,8.448)rectangle++(0.04,0.04) (0.982,8.746)rectangle++(0.04,0.04)
  (1.182,8.991)rectangle++(0.04,0.04) (1.362,9.191)rectangle++(0.04,0.04) (1.523,9.354)rectangle++(0.04,0.04) (1.667,9.487)rectangle++(0.04,0.04) (-2.154,4.059)rectangle++(0.04,0.04)
  (1.442,0.844)rectangle++(0.04,0.04) (-0.02,0.118)rectangle++(0.04,0.04) (-0.014,1.698)rectangle++(0.04,0.04) (0.053,3.04)rectangle++(0.04,0.04) (0.165,4.178)rectangle++(0.04,0.04)
  (0.305,5.141)rectangle++(0.04,0.04) (0.463,5.954)rectangle++(0.04,0.04) (0.629,6.638)rectangle++(0.04,0.04) (0.798,7.214)rectangle++(0.04,0.04) (0.965,7.696)rectangle++(0.04,0.04)
  (1.126,8.099)rectangle++(0.04,0.04) (-1.902,3.63)rectangle++(0.04,0.04) (-1.345,1.95)rectangle++(0.04,0.04) (-1.068,3.308)rectangle++(0.04,0.04) (1.069,0.946)rectangle++(0.04,0.04)
  (0.944,2.358)rectangle++(0.04,0.04) (0.895,3.563)rectangle++(0.04,0.04) (0.901,4.589)rectangle++(0.04,0.04) (0.947,5.46)rectangle++(0.04,0.04) (1.021,6.2)rectangle++(0.04,0.04)
  (-1.429,3.188)rectangle++(0.04,0.04) (-1.089,4.363)rectangle++(0.04,0.04) (-0.754,5.348)rectangle++(0.04,0.04) (-1.562,2.592)rectangle++(0.04,0.04) (-1.227,3.862)rectangle++(0.04,0.04)
  (-0.89,4.928)rectangle++(0.04,0.04) (-0.562,5.821)rectangle++(0.04,0.04) (-0.247,6.566)rectangle++(0.04,0.04) (0.051,7.187)rectangle++(0.04,0.04) (0.328,7.703)rectangle++(0.04,0.04)
  (-1.958,3.359)rectangle++(0.04,0.04) (-1.533,4.53)rectangle++(0.04,0.04) (-1.124,5.508)rectangle++(0.04,0.04) (-0.737,6.323)rectangle++(0.04,0.04) (1.864,1.756)rectangle++(0.04,0.04)
  (-0.105,2.404)rectangle++(0.04,0.04) (0.005,3.644)rectangle++(0.04,0.04) (-0.968,2.392)rectangle++(0.04,0.04) (0.797,0.752)rectangle++(0.04,0.04) (0.706,2.204)rectangle++(0.04,0.04)
  (0.686,3.441)rectangle++(0.04,0.04) (0.718,4.494)rectangle++(0.04,0.04) (0.788,5.387)rectangle++(0.04,0.04) (0.883,6.144)rectangle++(0.04,0.04) (0.994,6.783)rectangle++(0.04,0.04)
  (1.114,7.322)rectangle++(0.04,0.04) (1.238,7.775)rectangle++(0.04,0.04) (1.361,8.156)rectangle++(0.04,0.04) (1.481,8.474)rectangle++(0.04,0.04) (1.595,8.74)rectangle++(0.04,0.04)
  (1.704,8.961)rectangle++(0.04,0.04) (-2.01,3.952)rectangle++(0.04,0.04) (-1.553,5.036)rectangle++(0.04,0.04) (-1.121,5.939)rectangle++(0.04,0.04) (-1.789,2.638)rectangle++(0.04,0.04)
  (-1.418,3.91)rectangle++(0.04,0.04) (-1.051,4.976)rectangle++(0.04,0.04) (-0.696,5.868)rectangle++(0.04,0.04) (-0.359,6.612)rectangle++(0.04,0.04) (-0.043,7.231)rectangle++(0.04,0.04)
  (0.25,7.744)rectangle++(0.04,0.04) (0.52,8.169)rectangle++(0.04,0.04) (0.767,8.519)rectangle++(0.04,0.04) (0.99,8.806)rectangle++(0.04,0.04) (-0.02,1.392)rectangle++(0.04,0.04)
  (0.036,2.78)rectangle++(0.04,0.04) (-0.737,2.209)rectangle++(0.04,0.04) (-0.54,3.503)rectangle++(0.04,0.04) (-0.321,4.596)rectangle++(0.04,0.04) (-0.091,5.515)rectangle++(0.04,0.04)
  (-1.473,2.781)rectangle++(0.04,0.04) (-1.143,4.019)rectangle++(0.04,0.04) (-0.813,5.058)rectangle++(0.04,0.04) (-0.491,5.928)rectangle++(0.04,0.04) (-0.183,6.655)rectangle++(0.04,0.04)
  (-1.788,3.011)rectangle++(0.04,0.04) (-1.402,4.227)rectangle++(0.04,0.04) (-1.024,5.245)rectangle++(0.04,0.04) (-0.663,6.096)rectangle++(0.04,0.04) (-0.322,6.804)rectangle++(0.04,0.04)
  (-0.004,7.393)rectangle++(0.04,0.04) (0.29,7.88)rectangle++(0.04,0.04) (0.56,8.283)rectangle++(0.04,0.04) (0.805,8.614)rectangle++(0.04,0.04) (2.274,2.707)rectangle++(0.04,0.04)
  (2.039,3.806)rectangle++(0.04,0.04) (1.883,4.75)rectangle++(0.04,0.04) (1.788,5.558)rectangle++(0.04,0.04) (1.74,6.249)rectangle++(0.04,0.04) (1.727,6.838)rectangle++(0.04,0.04)
  (1.739,7.34)rectangle++(0.04,0.04) (1.77,7.765)rectangle++(0.04,0.04) (1.813,8.126)rectangle++(0.04,0.04) (1.864,8.431)rectangle++(0.04,0.04) (2.064,2.938)rectangle++(0.04,0.04)
  (1.869,4.011)rectangle++(0.04,0.04) (0.825,1.879)rectangle++(0.04,0.04) (0.774,3.16)rectangle++(0.04,0.04) (0.782,4.251)rectangle++(0.04,0.04) (0.833,5.178)rectangle++(0.04,0.04)
  (0.913,5.965)rectangle++(0.04,0.04) (1.012,6.63)rectangle++(0.04,0.04) (1.123,7.191)rectangle++(0.04,0.04) (1.24,7.663)rectangle++(0.04,0.04) (1.359,8.061)rectangle++(0.04,0.04)
  (1.475,8.393)rectangle++(0.04,0.04) (-1.908,3.775)rectangle++(0.04,0.04) (-1.384,1.981)rectangle++(0.04,0.04) (-1.1,3.335)rectangle++(0.04,0.04) (-0.804,4.475)rectangle++(0.04,0.04)
  (-0.506,5.432)rectangle++(0.04,0.04) (-0.215,6.234)rectangle++(0.04,0.04) (0.064,6.903)rectangle++(0.04,0.04) (0.329,7.462)rectangle++(0.04,0.04) (0.576,7.925)rectangle++(0.04,0.04)
  (0.804,8.31)rectangle++(0.04,0.04) (1.014,8.627)rectangle++(0.04,0.04) (1.204,8.889)rectangle++(0.04,0.04) (1.377,9.104)rectangle++(0.04,0.04) (2.325,2.973)rectangle++(0.04,0.04)
  (2.093,4.03)rectangle++(0.04,0.04) (1.938,4.938)rectangle++(0.04,0.04) (1.843,5.716)rectangle++(0.04,0.04) (1.793,6.381)rectangle++(0.04,0.04) (1.777,6.948)rectangle++(0.04,0.04)
  (1.786,7.431)rectangle++(0.04,0.04) (1.813,7.841)rectangle++(0.04,0.04) (1.853,8.189)rectangle++(0.04,0.04) (1.9,8.483)rectangle++(0.04,0.04) (1.952,8.73)rectangle++(0.04,0.04)
  (2.006,8.939)rectangle++(0.04,0.04) (2.061,9.114)rectangle++(0.04,0.04) (2.114,9.261)rectangle++(0.04,0.04) (2.165,9.383)rectangle++(0.04,0.04) (2.214,9.485)rectangle++(0.04,0.04)
  (2.306,3.282)rectangle++(0.04,0.04) (0.556,1.817)rectangle++(0.04,0.04) (0.543,3.119)rectangle++(0.04,0.04) (0.584,4.225)rectangle++(0.04,0.04) (1.078,1.596)rectangle++(0.04,0.04)
  (0.978,2.91)rectangle++(0.04,0.04) (0.946,4.03)rectangle++(0.04,0.04) (0.963,4.984)rectangle++(0.04,0.04) (1.234,1.876)rectangle++(0.04,0.04) (1.122,3.142)rectangle++(0.04,0.04)
  (1.077,4.222)rectangle++(0.04,0.04) (1.082,5.142)rectangle++(0.04,0.04) (1.123,5.923)rectangle++(0.04,0.04) (1.189,6.586)rectangle++(0.04,0.04) (1.272,7.147)rectangle++(0.04,0.04)
  (1.365,7.62)rectangle++(0.04,0.04) (-0.02,1.202)rectangle++(0.04,0.04) (0.029,2.619)rectangle++(0.04,0.04) (0.127,3.821)rectangle++(0.04,0.04) (0.259,4.839)rectangle++(0.04,0.04)
  (0.411,5.699)rectangle++(0.04,0.04) (0.575,6.424)rectangle++(0.04,0.04) (0.744,7.034)rectangle++(0.04,0.04) (0.911,7.545)rectangle++(0.04,0.04) (1.074,7.973)rectangle++(0.04,0.04)
  (1.23,8.33)rectangle++(0.04,0.04) (2.131,2.749)rectangle++(0.04,0.04) (1.919,3.848)rectangle++(0.04,0.04) (1.783,4.79)rectangle++(0.04,0.04) (1.705,5.596)rectangle++(0.04,0.04)
  (1.671,6.285)rectangle++(0.04,0.04) (1.669,6.872)rectangle++(0.04,0.04) (1.691,7.37)rectangle++(0.04,0.04) (1.73,7.793)rectangle++(0.04,0.04) (1.78,8.151)rectangle++(0.04,0.04)
  (-1.784,3.792)rectangle++(0.04,0.04) (-1.367,4.891)rectangle++(0.04,0.04) (1.557,1.248)rectangle++(0.04,0.04) (1.371,2.595)rectangle++(0.04,0.04) (1.267,3.747)rectangle++(0.04,0.04)
  (1.225,4.73)rectangle++(0.04,0.04) (-1.006,2.911)rectangle++(0.04,0.04) (-0.741,4.111)rectangle++(0.04,0.04) (-1.238,2.323)rectangle++(0.04,0.04) (-0.962,3.62)rectangle++(0.04,0.04)
  (-0.675,4.712)rectangle++(0.04,0.04) (-0.387,5.628)rectangle++(0.04,0.04) (-0.106,6.396)rectangle++(0.04,0.04) (-0.02,1.007)rectangle++(0.04,0.04) (0.021,2.453)rectangle++(0.04,0.04)
  (0.666,1.024)rectangle++(0.04,0.04) (0.605,2.44)rectangle++(0.04,0.04) (0.61,3.646)rectangle++(0.04,0.04) (0.662,4.671)rectangle++(0.04,0.04) (-1.103,2.769)rectangle++(0.04,0.04)
  (-0.962,1.944)rectangle++(0.04,0.04) (-0.742,3.287)rectangle++(0.04,0.04) (-0.501,4.42)rectangle++(0.04,0.04) (-0.252,5.373)rectangle++(0.04,0.04) (-0.001,6.174)rectangle++(0.04,0.04)
  (0.244,6.844)rectangle++(0.04,0.04) (0.479,7.404)rectangle++(0.04,0.04) (0.701,7.87)rectangle++(0.04,0.04) (2.081,2.501)rectangle++(0.04,0.04) (-0.255,2.618)rectangle++(0.04,0.04)
  (-0.114,3.832)rectangle++(0.04,0.04) (0.054,4.858)rectangle++(0.04,0.04) (0.238,5.723)rectangle++(0.04,0.04) (0.429,6.451)rectangle++(0.04,0.04) (0.62,7.063)rectangle++(0.04,0.04)
  (0.808,7.575)rectangle++(0.04,0.04) (1.982,2.458)rectangle++(0.04,0.04) (1.781,3.606)rectangle++(0.04,0.04) (0.725,1.759)rectangle++(0.04,0.04) (0.685,3.062)rectangle++(0.04,0.04)
  (0.702,4.171)rectangle++(0.04,0.04) (1.045,1.614)rectangle++(0.04,0.04) (0.951,2.926)rectangle++(0.04,0.04) (-0.592,2.451)rectangle++(0.04,0.04) (-0.407,3.704)rectangle++(0.04,0.04)
  (-0.2,4.761)rectangle++(0.04,0.04) (0.018,5.651)rectangle++(0.04,0.04) (0.239,6.399)rectangle++(0.04,0.04) (0.457,7.025)rectangle++(0.04,0.04) (0.667,7.549)rectangle++(0.04,0.04)
  (0.867,7.987)rectangle++(0.04,0.04) (1.054,8.35)rectangle++(0.04,0.04) (1.228,8.652)rectangle++(0.04,0.04) (1.388,8.901)rectangle++(0.04,0.04) (1.533,9.107)rectangle++(0.04,0.04)
  (1.665,9.275)rectangle++(0.04,0.04) (1.784,9.414)rectangle++(0.04,0.04) (1.891,9.526)rectangle++(0.04,0.04) (1.986,9.618)rectangle++(0.04,0.04) (2.071,9.692)rectangle++(0.04,0.04)
  (2.146,9.752)rectangle++(0.04,0.04) (2.212,9.799)rectangle++(0.04,0.04) (2.27,9.837)rectangle++(0.04,0.04) (2.321,9.867)rectangle++(0.04,0.04) (-2.122,4.293)rectangle++(0.04,0.04)
  (-1.635,5.331)rectangle++(0.04,0.04) (-1.178,6.193)rectangle++(0.04,0.04) (-0.756,6.907)rectangle++(0.04,0.04) (-0.369,7.497)rectangle++(0.04,0.04) (-0.016,7.984)rectangle++(0.04,0.04)
  (0.304,8.383)rectangle++(0.04,0.04) (0.591,8.71)rectangle++(0.04,0.04) (0.849,8.976)rectangle++(0.04,0.04) (-2.185,3.759)rectangle++(0.04,0.04) (-1.709,4.879)rectangle++(0.04,0.04)
  (-1.26,5.811)rectangle++(0.04,0.04) (-0.841,6.586)rectangle++(0.04,0.04) (1.953,1.792)rectangle++(0.04,0.04) (1.729,3.041)rectangle++(0.04,0.04) (1.589,4.112)rectangle++(0.04,0.04)
  (1.513,5.028)rectangle++(0.04,0.04) (1.485,5.809)rectangle++(0.04,0.04) (1.493,6.475)rectangle++(0.04,0.04) (-0.02,1.019)rectangle++(0.04,0.04) (0.022,2.463)rectangle++(0.04,0.04)
  (0.115,3.689)rectangle++(0.04,0.04) (0.243,4.727)rectangle++(0.04,0.04) (0.393,5.605)rectangle++(0.04,0.04) (0.556,6.345)rectangle++(0.04,0.04) (0.724,6.967)rectangle++(0.04,0.04)
  (0.892,7.489)rectangle++(0.04,0.04) (1.056,7.926)rectangle++(0.04,0.04) (-1.871,3.576)rectangle++(0.04,0.04) (-1.449,4.71)rectangle++(0.04,0.04) (-1.046,5.658)rectangle++(0.04,0.04)
  (-0.665,6.447)rectangle++(0.04,0.04) (-0.309,7.103)rectangle++(0.04,0.04) (2.018,2.054)rectangle++(0.04,0.04) (0.255,1.448)rectangle++(0.04,0.04) (0.273,2.817)rectangle++(0.04,0.04)
  (0.342,3.979)rectangle++(0.04,0.04) (0.448,4.965)rectangle++(0.04,0.04) (0.577,5.798)rectangle++(0.04,0.04) (0.72,6.502)rectangle++(0.04,0.04) (0.87,7.094)rectangle++(0.04,0.04)
  (-1.692,3.35)rectangle++(0.04,0.04) (-1.306,4.511)rectangle++(0.04,0.04) (-0.932,5.483)rectangle++(0.04,0.04) (-0.575,6.294)rectangle++(0.04,0.04) (-0.239,6.969)rectangle++(0.04,0.04)
  (1.97,2.04)rectangle++(0.04,0.04) (1.754,3.252)rectangle++(0.04,0.04) (1.619,4.29)rectangle++(0.04,0.04) (1.545,5.178)rectangle++(0.04,0.04) (-1.058,3.084)rectangle++(0.04,0.04)
  (-0.779,4.26)rectangle++(0.04,0.04) (-0.494,5.248)rectangle++(0.04,0.04) (-0.212,6.077)rectangle++(0.04,0.04) (-1.644,2.877)rectangle++(0.04,0.04) (-1.284,4.108)rectangle++(0.04,0.04)
  (-0.929,5.139)rectangle++(0.04,0.04) (1.561,1.422)rectangle++(0.04,0.04) (1.381,2.742)rectangle++(0.04,0.04) (1.282,3.872)rectangle++(0.04,0.04) (1.242,4.836)rectangle++(0.04,0.04)
  (1.15,1.914)rectangle++(0.04,0.04) (0.346,1.188)rectangle++(0.04,0.04) (0.339,2.592)rectangle++(0.04,0.04) (0.39,3.786)rectangle++(0.04,0.04) (0.481,4.799)rectangle++(0.04,0.04)
  (0.598,5.656)rectangle++(0.04,0.04) (0.733,6.38)rectangle++(0.04,0.04) (0.876,6.99)rectangle++(0.04,0.04) (1.022,7.503)rectangle++(0.04,0.04) (1.166,7.932)rectangle++(0.04,0.04)
  (1.307,8.292)rectangle++(0.04,0.04) (1.44,8.592)rectangle++(0.04,0.04) (1.566,8.842)rectangle++(0.04,0.04) (-2.007,3.894)rectangle++(0.04,0.04) (-1.552,4.987)rectangle++(0.04,0.04)
  (-1.122,5.897)rectangle++(0.04,0.04) (-0.72,6.653)rectangle++(0.04,0.04) (-0.348,7.28)rectangle++(0.04,0.04) (-0.007,7.799)rectangle++(0.04,0.04) (0.304,8.225)rectangle++(0.04,0.04)
  (0.585,8.576)rectangle++(0.04,0.04) (0.838,8.862)rectangle++(0.04,0.04) (1.065,9.095)rectangle++(0.04,0.04) (2.37,2.89)rectangle++(0.04,0.04) (2.128,3.958)rectangle++(0.04,0.04)
  (1.965,4.875)rectangle++(0.04,0.04) (1.863,5.661)rectangle++(0.04,0.04) (1.808,6.334)rectangle++(0.04,0.04) (1.788,6.908)rectangle++(0.04,0.04) (1.794,7.396)rectangle++(0.04,0.04)
  (1.818,7.811)rectangle++(0.04,0.04) (1.856,8.163)rectangle++(0.04,0.04) (1.902,8.461)rectangle++(0.04,0.04) (1.953,8.712)rectangle++(0.04,0.04) (2.006,8.923)rectangle++(0.04,0.04)
  (2.06,9.1)rectangle++(0.04,0.04) (2.113,9.249)rectangle++(0.04,0.04) (2.164,9.374)rectangle++(0.04,0.04) (2.212,9.477)rectangle++(0.04,0.04) (2.257,9.563)rectangle++(0.04,0.04)
  (2.299,9.635)rectangle++(0.04,0.04) (2.337,9.694)rectangle++(0.04,0.04) (2.372,9.742)rectangle++(0.04,0.04) (2.404,9.782)rectangle++(0.04,0.04) (-2.084,4.294)rectangle++(0.04,0.04)
  (-1.602,5.329)rectangle++(0.04,0.04) (-1.15,6.19)rectangle++(0.04,0.04) (-0.733,6.904)rectangle++(0.04,0.04) (-0.349,7.494)rectangle++(0.04,0.04) (0.001,7.98)rectangle++(0.04,0.04)
  (0.318,8.379)rectangle++(0.04,0.04) (0.603,8.706)rectangle++(0.04,0.04) (-2.164,3.643)rectangle++(0.04,0.04) (-0.02,0.566)rectangle++(0.04,0.04) (0.003,2.078)rectangle++(0.04,0.04)
  (0.084,3.363)rectangle++(0.04,0.04) (0.204,4.451)rectangle++(0.04,0.04) (0.349,5.371)rectangle++(0.04,0.04) (1.434,1.81)rectangle++(0.04,0.04) (1.289,3.077)rectangle++(0.04,0.04)
  (1.217,4.16)rectangle++(0.04,0.04) (1.198,5.084)rectangle++(0.04,0.04) (1.22,5.869)rectangle++(0.04,0.04) (1.443,2.156)rectangle++(0.04,0.04) (1.311,3.371)rectangle++(0.04,0.04)
  (1.247,4.409)rectangle++(0.04,0.04) (1.234,5.294)rectangle++(0.04,0.04) (1.258,6.047)rectangle++(0.04,0.04) (-1.342,3.209)rectangle++(0.04,0.04) (-1.014,4.377)rectangle++(0.04,0.04)
  (-0.689,5.357)rectangle++(0.04,0.04) (-0.374,6.178)rectangle++(0.04,0.04) (-0.073,6.862)rectangle++(0.04,0.04) (0.21,7.432)rectangle++(0.04,0.04) (0.474,7.905)rectangle++(0.04,0.04)
  (0.717,8.296)rectangle++(0.04,0.04) (-2.035,3.579)rectangle++(0.04,0.04) (-1.589,4.72)rectangle++(0.04,0.04) (-1.164,5.672)rectangle++(0.04,0.04) (-0.765,6.464)rectangle++(0.04,0.04)
  (-0.394,7.121)rectangle++(0.04,0.04) (-0.052,7.665)rectangle++(0.04,0.04) (0.26,8.113)rectangle++(0.04,0.04) (0.544,8.482)rectangle++(0.04,0.04) (2.276,2.607)rectangle++(0.04,0.04)
  (-0.244,2.686)rectangle++(0.04,0.04) (-0.102,3.889)rectangle++(0.04,0.04) (0.067,4.906)rectangle++(0.04,0.04) (0.251,5.764)rectangle++(0.04,0.04) (0.441,6.485)rectangle++(0.04,0.04)
  (0.632,7.091)rectangle++(0.04,0.04) (0.819,7.598)rectangle++(0.04,0.04) (1.987,2.467)rectangle++(0.04,0.04) (1.786,3.613)rectangle++(0.04,0.04) (1.66,4.596)rectangle++(0.04,0.04)
  (1.593,5.436)rectangle++(0.04,0.04) (1.569,6.153)rectangle++(0.04,0.04) (1.578,6.764)rectangle++(0.04,0.04) (1.609,7.282)rectangle++(0.04,0.04) (1.657,7.722)rectangle++(0.04,0.04)
  (1.715,8.093)rectangle++(0.04,0.04) (1.779,8.407)rectangle++(0.04,0.04) (1.847,8.671)rectangle++(0.04,0.04) (1.914,8.893)rectangle++(0.04,0.04) (1.981,9.078)rectangle++(0.04,0.04)
  (-1.985,4.042)rectangle++(0.04,0.04) (-1.528,5.111)rectangle++(0.04,0.04) (-1.097,6.002)rectangle++(0.04,0.04) (-0.694,6.742)rectangle++(0.04,0.04) (-1.913,2.912)rectangle++(0.04,0.04)
  (-1.512,4.148)rectangle++(0.04,0.04) (-1.121,5.183)rectangle++(0.04,0.04) (-0.748,6.046)rectangle++(0.04,0.04) (-0.396,6.766)rectangle++(0.04,0.04) (-0.068,7.363)rectangle++(0.04,0.04)
  (0.234,7.857)rectangle++(0.04,0.04) (0.511,8.265)rectangle++(0.04,0.04) (0.763,8.601)rectangle++(0.04,0.04) (0.99,8.877)rectangle++(0.04,0.04) (1.195,9.102)rectangle++(0.04,0.04)
  (1.377,9.285)rectangle++(0.04,0.04) (1.54,9.433)rectangle++(0.04,0.04) (1.684,9.553)rectangle++(0.04,0.04) (1.811,9.649)rectangle++(0.04,0.04) (1.923,9.725)rectangle++(0.04,0.04)
  (2.022,9.786)rectangle++(0.04,0.04) (2.419,3.304)rectangle++(0.04,0.04) (2.186,4.308)rectangle++(0.04,0.04) (2.029,5.171)rectangle++(0.04,0.04) (1.929,5.91)rectangle++(0.04,0.04)
  (1.874,6.543)rectangle++(0.04,0.04) (1.852,7.082)rectangle++(0.04,0.04) (1.855,7.542)rectangle++(0.04,0.04) (1.877,7.933)rectangle++(0.04,0.04) (1.91,8.264)rectangle++(0.04,0.04)
  (1.952,8.544)rectangle++(0.04,0.04) (1.999,8.781)rectangle++(0.04,0.04) (2.048,8.98)rectangle++(0.04,0.04) (2.098,9.147)rectangle++(0.04,0.04) (2.147,9.287)rectangle++(0.04,0.04)
  (2.194,9.405)rectangle++(0.04,0.04) (2.239,9.502)rectangle++(0.04,0.04) (2.281,9.584)rectangle++(0.04,0.04) (2.32,9.651)rectangle++(0.04,0.04) (2.356,9.707)rectangle++(0.04,0.04)
  (2.389,9.753)rectangle++(0.04,0.04) (2.418,9.79)rectangle++(0.04,0.04) (2.445,9.821)rectangle++(0.04,0.04) (2.469,9.847)rectangle++(0.04,0.04) (2.49,9.867)rectangle++(0.04,0.04)
  (2.509,9.884)rectangle++(0.04,0.04) (2.526,9.897)rectangle++(0.04,0.04) (2.541,9.908)rectangle++(0.04,0.04) (-2.089,4.353)rectangle++(0.04,0.04) (-1.604,5.38)rectangle++(0.04,0.04)
  (-1.15,6.233)rectangle++(0.04,0.04) (-0.731,6.94)rectangle++(0.04,0.04) (-0.346,7.525)rectangle++(0.04,0.04) (0.005,8.006)rectangle++(0.04,0.04) (2.224,2.353)rectangle++(0.04,0.04)
  (1.982,3.507)rectangle++(0.04,0.04) (1.823,4.498)rectangle++(0.04,0.04) (1.727,5.347)rectangle++(0.04,0.04) (1.68,6.072)rectangle++(0.04,0.04) (1.668,6.69)rectangle++(0.04,0.04)
  (1.684,7.216)rectangle++(0.04,0.04) (1.717,7.662)rectangle++(0.04,0.04) (1.764,8.041)rectangle++(0.04,0.04) (1.819,8.36)rectangle++(0.04,0.04) (1.878,8.63)rectangle++(0.04,0.04)
  (1.94,8.856)rectangle++(0.04,0.04) (2.001,9.046)rectangle++(0.04,0.04) (2.06,9.206)rectangle++(0.04,0.04) (2.117,9.339)rectangle++(0.04,0.04) (2.171,9.449)rectangle++(0.04,0.04)
  (2.221,9.541)rectangle++(0.04,0.04) (2.267,9.617)rectangle++(0.04,0.04) (2.31,9.68)rectangle++(0.04,0.04) (2.348,9.732)rectangle++(0.04,0.04) (-2.082,4.27)rectangle++(0.04,0.04)
  (-1.601,5.309)rectangle++(0.04,0.04) (-1.151,6.173)rectangle++(0.04,0.04) (-0.733,6.889)rectangle++(0.04,0.04) (-0.35,7.481)rectangle++(0.04,0.04) (2.13,2.135)rectangle++(0.04,0.04)
  (1.894,3.325)rectangle++(0.04,0.04) (1.74,4.347)rectangle++(0.04,0.04) (1.651,5.222)rectangle++(0.04,0.04) (1.61,5.968)rectangle++(0.04,0.04) (1.605,6.605)rectangle++(0.04,0.04)
  (1.626,7.146)rectangle++(0.04,0.04) (1.666,7.605)rectangle++(0.04,0.04) (1.718,7.994)rectangle++(0.04,0.04) (1.963,2.795)rectangle++(0.04,0.04) (1.778,3.894)rectangle++(0.04,0.04)
  (1.665,4.835)rectangle++(0.04,0.04) (1.607,5.639)rectangle++(0.04,0.04) (-1.166,3.199)rectangle++(0.04,0.04) (-0.865,4.362)rectangle++(0.04,0.04) (-1.328,2.35)rectangle++(0.04,0.04)
  (-1.037,3.647)rectangle++(0.04,0.04) (1.159,1.035)rectangle++(0.04,0.04) (1.025,2.43)rectangle++(0.04,0.04) (0.966,3.621)rectangle++(0.04,0.04) (0.964,4.635)rectangle++(0.04,0.04)
  (-1.034,2.83)rectangle++(0.04,0.04) (-0.768,4.043)rectangle++(0.04,0.04) (-0.493,5.064)rectangle++(0.04,0.04) (-0.219,5.92)rectangle++(0.04,0.04) (0.049,6.637)rectangle++(0.04,0.04)
  (0.305,7.236)rectangle++(0.04,0.04) (0.546,7.734)rectangle++(0.04,0.04) (0.772,8.149)rectangle++(0.04,0.04) (0.98,8.492)rectangle++(0.04,0.04) (1.17,8.775)rectangle++(0.04,0.04)
  (1.343,9.008)rectangle++(0.04,0.04) (1.5,9.199)rectangle++(0.04,0.04) (2.333,3.028)rectangle++(0.04,0.04) (2.102,4.077)rectangle++(0.04,0.04) (1.948,4.977)rectangle++(0.04,0.04)
  (1.853,5.749)rectangle++(0.04,0.04) (1.802,6.409)rectangle++(0.04,0.04) (1.507,2.437)rectangle++(0.04,0.04) (1.376,3.607)rectangle++(0.04,0.04) (1.312,4.607)rectangle++(0.04,0.04)
  (1.297,5.46)rectangle++(0.04,0.04) (1.319,6.185)rectangle++(0.04,0.04) (1.366,6.801)rectangle++(0.04,0.04) (1.431,7.322)rectangle++(0.04,0.04) (1.818,2.559)rectangle++(0.04,0.04)
  (0.427,1.517)rectangle++(0.04,0.04) (0.421,2.869)rectangle++(0.04,0.04) (0.47,4.018)rectangle++(0.04,0.04) (0.558,4.992)rectangle++(0.04,0.04) (0.672,5.817)rectangle++(0.04,0.04)
  (0.802,6.514)rectangle++(0.04,0.04) (0.94,7.101)rectangle++(0.04,0.04) (1.081,7.595)rectangle++(0.04,0.04) (1.22,8.008)rectangle++(0.04,0.04) (1.355,8.354)rectangle++(0.04,0.04)
  (1.484,8.643)rectangle++(0.04,0.04) (1.605,8.884)rectangle++(0.04,0.04) (-2.01,3.913)rectangle++(0.04,0.04) (-1.554,5.002)rectangle++(0.04,0.04) (-1.123,5.91)rectangle++(0.04,0.04)
  (-0.72,6.665)rectangle++(0.04,0.04) (-0.348,7.29)rectangle++(0.04,0.04) (-0.006,7.807)rectangle++(0.04,0.04) (0.305,8.232)rectangle++(0.04,0.04) (0.586,8.581)rectangle++(0.04,0.04)
  (0.839,8.867)rectangle++(0.04,0.04) (-2.159,3.733)rectangle++(0.04,0.04) (-1.688,4.855)rectangle++(0.04,0.04) (-1.243,5.791)rectangle++(0.04,0.04) (-0.827,6.568)rectangle++(0.04,0.04)
  (-0.02,1.034)rectangle++(0.04,0.04) (0.022,2.476)rectangle++(0.04,0.04) (0.116,3.7)rectangle++(0.04,0.04) (0.244,4.736)rectangle++(0.04,0.04) (0.395,5.612)rectangle++(0.04,0.04)
  (-0.02,0.881)rectangle++(0.04,0.04) (0.016,2.346)rectangle++(0.04,0.04) (0.637,0.997)rectangle++(0.04,0.04) (0.579,2.418)rectangle++(0.04,0.04) (0.587,3.629)rectangle++(0.04,0.04)
  (-0.847,2.522)rectangle++(0.04,0.04) (-0.621,3.774)rectangle++(0.04,0.04) (-0.38,4.829)rectangle++(0.04,0.04) (-0.132,5.716)rectangle++(0.04,0.04) (0.115,6.46)rectangle++(0.04,0.04)
  (0.354,7.083)rectangle++(0.04,0.04) (0.582,7.602)rectangle++(0.04,0.04) (0.796,8.035)rectangle++(0.04,0.04) (0.996,8.394)rectangle++(0.04,0.04) (1.18,8.691)rectangle++(0.04,0.04)
  (1.349,8.937)rectangle++(0.04,0.04) (1.502,9.138)rectangle++(0.04,0.04) (1.64,9.304)rectangle++(0.04,0.04) (1.764,9.439)rectangle++(0.04,0.04) (1.874,9.549)rectangle++(0.04,0.04)
  (1.973,9.638)rectangle++(0.04,0.04) (2.06,9.709)rectangle++(0.04,0.04) (2.138,9.767)rectangle++(0.04,0.04) (2.205,9.812)rectangle++(0.04,0.04) (2.265,9.848)rectangle++(0.04,0.04)
  (2.317,9.877)rectangle++(0.04,0.04) (2.362,9.899)rectangle++(0.04,0.04) (2.402,9.916)rectangle++(0.04,0.04) (2.436,9.928)rectangle++(0.04,0.04) (2.465,9.938)rectangle++(0.04,0.04)
  (2.395,3.456)rectangle++(0.04,0.04) (-0.441,2.9)rectangle++(0.04,0.04) (0.861,1.011)rectangle++(0.04,0.04) (0.77,2.422)rectangle++(0.04,0.04) (0.749,3.624)rectangle++(0.04,0.04)
  (0.779,4.646)rectangle++(0.04,0.04) (1.167,1.748)rectangle++(0.04,0.04) (1.059,3.035)rectangle++(0.04,0.04) (1.02,4.134)rectangle++(0.04,0.04) (-0.892,2.733)rectangle++(0.04,0.04)
  (-0.91,1.985)rectangle++(0.04,0.04) (-0.696,3.32)rectangle++(0.04,0.04) (-0.461,4.446)rectangle++(0.04,0.04) (-0.217,5.394)rectangle++(0.04,0.04) (-1.467,2.726)rectangle++(0.04,0.04)
  (-1.14,3.972)rectangle++(0.04,0.04) (-0.812,5.018)rectangle++(0.04,0.04) (-0.492,5.894)rectangle++(0.04,0.04) (-0.185,6.626)rectangle++(0.04,0.04) (0.106,7.235)rectangle++(0.04,0.04)
  (0.377,7.742)rectangle++(0.04,0.04) (0.628,8.162)rectangle++(0.04,0.04) (0.858,8.509)rectangle++(0.04,0.04) (1.068,8.794)rectangle++(0.04,0.04) (1.257,9.029)rectangle++(0.04,0.04)
  (-2.117,3.864)rectangle++(0.04,0.04) (-1.647,4.966)rectangle++(0.04,0.04) (-1.204,5.883)rectangle++(0.04,0.04) (-0.79,6.645)rectangle++(0.04,0.04) (-0.408,7.276)rectangle++(0.04,0.04)
  (-0.058,7.797)rectangle++(0.04,0.04) (-2.06,3.291)rectangle++(0.04,0.04) (-1.622,4.476)rectangle++(0.04,0.04) (-1.201,5.466)rectangle++(0.04,0.04) (-0.805,6.29)rectangle++(0.04,0.04)
  (-1.818,2.788)rectangle++(0.04,0.04) (-1.11,1.784)rectangle++(0.04,0.04) (-0.874,3.157)rectangle++(0.04,0.04) (-0.619,4.315)rectangle++(0.04,0.04) (-0.356,5.289)rectangle++(0.04,0.04)
  (-0.093,6.106)rectangle++(0.04,0.04) (0.163,6.79)rectangle++(0.04,0.04) (0.408,7.361)rectangle++(0.04,0.04) (0.639,7.837)rectangle++(0.04,0.04) (0.854,8.232)rectangle++(0.04,0.04)
  (1.053,8.559)rectangle++(0.04,0.04) (1.236,8.829)rectangle++(0.04,0.04) (1.401,9.052)rectangle++(0.04,0.04) (2.307,2.967)rectangle++(0.04,0.04) (-0.331,2.772)rectangle++(0.04,0.04)
  (-0.173,3.966)rectangle++(0.04,0.04) (1.119,1.337)rectangle++(0.04,0.04) (1.002,2.688)rectangle++(0.04,0.04) (-0.52,2.411)rectangle++(0.04,0.04) (-0.347,3.666)rectangle++(0.04,0.04)
  (-0.151,4.726)rectangle++(0.04,0.04) (0.059,5.62)rectangle++(0.04,0.04) (-0.02,0.882)rectangle++(0.04,0.04) (0.016,2.347)rectangle++(0.04,0.04) (0.105,3.591)rectangle++(0.04,0.04)
  (0.231,4.644)rectangle++(0.04,0.04) (0.38,5.534)rectangle++(0.04,0.04) (0.542,6.285)rectangle++(0.04,0.04) (0.71,6.917)rectangle++(0.04,0.04) (0.878,7.447)rectangle++(0.04,0.04)
  (1.042,7.891)rectangle++(0.04,0.04) (1.199,8.262)rectangle++(0.04,0.04) (1.348,8.571)rectangle++(0.04,0.04) (1.486,8.828)rectangle++(0.04,0.04) (1.614,9.04)rectangle++(0.04,0.04)
  (1.731,9.216)rectangle++(0.04,0.04) (1.838,9.36)rectangle++(0.04,0.04) (2.328,3.154)rectangle++(0.04,0.04) (0.517,1.792)rectangle++(0.04,0.04) (0.407,0.994)rectangle++(0.04,0.04)
  (0.383,2.425)rectangle++(0.04,0.04) (0.421,3.642)rectangle++(0.04,0.04) (0.501,4.675)rectangle++(0.04,0.04) (0.611,5.55)rectangle++(0.04,0.04) (-1.342,2.951)rectangle++(0.04,0.04)
  (-1.025,4.158)rectangle++(0.04,0.04) (-0.707,5.171)rectangle++(0.04,0.04) (-0.396,6.02)rectangle++(0.04,0.04) (-0.098,6.729)rectangle++(0.04,0.04) (0.183,7.32)rectangle++(0.04,0.04)
  (0.446,7.811)rectangle++(0.04,0.04) (2.103,2.421)rectangle++(0.04,0.04) (1.882,3.57)rectangle++(0.04,0.04) (1.74,4.555)rectangle++(0.04,0.04) (1.659,5.398)rectangle++(0.04,0.04)
  (1.624,6.119)rectangle++(0.04,0.04) (1.623,6.732)rectangle++(0.04,0.04) (1.647,7.253)rectangle++(0.04,0.04) (-1.578,3.563)rectangle++(0.04,0.04) (-1.201,4.688)rectangle++(0.04,0.04)
  (-0.835,5.629)rectangle++(0.04,0.04) (-0.487,6.414)rectangle++(0.04,0.04) (-0.16,7.068)rectangle++(0.04,0.04) (1.986,2.085)rectangle++(0.04,0.04) (1.769,3.289)rectangle++(0.04,0.04)
  (-0.523,2.719)rectangle++(0.04,0.04) (-0.338,3.929)rectangle++(0.04,0.04) (-0.132,4.949)rectangle++(0.04,0.04) (0.084,5.808)rectangle++(0.04,0.04) (-1.515,2.886)rectangle++(0.04,0.04)
  (-1.174,4.11)rectangle++(0.04,0.04) (-0.836,5.137)rectangle++(0.04,0.04) (-0.507,5.996)rectangle++(0.04,0.04) (-0.193,6.713)rectangle++(0.04,0.04) (0.102,7.31)rectangle++(0.04,0.04)
  (2.014,2.211)rectangle++(0.04,0.04) (1.798,3.395)rectangle++(0.04,0.04) (1.662,4.41)rectangle++(0.04,0.04) (0.968,1.921)rectangle++(0.04,0.04) (0.897,3.19)rectangle++(0.04,0.04)
  (0.888,4.272)rectangle++(0.04,0.04) (0.924,5.192)rectangle++(0.04,0.04) (0.991,5.972)rectangle++(0.04,0.04) (1.079,6.633)rectangle++(0.04,0.04) (1.18,7.191)rectangle++(0.04,0.04)
  (1.819,2.463)rectangle++(0.04,0.04) (1.643,3.617)rectangle++(0.04,0.04) (1.539,4.605)rectangle++(0.04,0.04) (1.49,5.449)rectangle++(0.04,0.04) (1.482,6.168)rectangle++(0.04,0.04)
  (1.504,6.78)rectangle++(0.04,0.04) (1.548,7.299)rectangle++(0.04,0.04) (1.605,7.738)rectangle++(0.04,0.04) (1.672,8.109)rectangle++(0.04,0.04) (2.002,2.811)rectangle++(0.04,0.04)
  (1.812,3.905)rectangle++(0.04,0.04) (1.695,4.843)rectangle++(0.04,0.04) (1.632,5.645)rectangle++(0.04,0.04) (1.611,6.329)rectangle++(0.04,0.04) (1.62,6.912)rectangle++(0.04,0.04)
  (1.651,7.406)rectangle++(0.04,0.04) (1.809,2.637)rectangle++(0.04,0.04) (1.641,3.765)rectangle++(0.04,0.04) (1.543,4.731)rectangle++(0.04,0.04) (1.499,5.556)rectangle++(0.04,0.04)
  (1.494,6.259)rectangle++(0.04,0.04) (1.518,6.856)rectangle++(0.04,0.04) (1.562,7.363)rectangle++(0.04,0.04) (1.62,7.793)rectangle++(0.04,0.04) (1.687,8.155)rectangle++(0.04,0.04)
  (1.758,8.461)rectangle++(0.04,0.04) (1.83,8.717)rectangle++(0.04,0.04) (-1.922,3.928)rectangle++(0.04,0.04) (-1.478,5.012)rectangle++(0.04,0.04) (-1.058,5.915)rectangle++(0.04,0.04)
  (-0.665,6.667)rectangle++(0.04,0.04) (-0.301,7.289)rectangle++(0.04,0.04) (0.034,7.804)rectangle++(0.04,0.04) (0.338,8.228)rectangle++(0.04,0.04) (0.615,8.577)rectangle++(0.04,0.04)
  (0.863,8.862)rectangle++(0.04,0.04) (1.086,9.094)rectangle++(0.04,0.04) (1.285,9.283)rectangle++(0.04,0.04) (1.461,9.435)rectangle++(0.04,0.04) (1.617,9.558)rectangle++(0.04,0.04)
  (2.416,3.144)rectangle++(0.04,0.04) (2.177,4.172)rectangle++(0.04,0.04) (-0.671,3.008)rectangle++(0.04,0.04) (-0.452,4.18)rectangle++(0.04,0.04) (-0.219,5.167)rectangle++(0.04,0.04)
  (0.018,5.997)rectangle++(0.04,0.04) (0.253,6.693)rectangle++(0.04,0.04) (1.819,2.102)rectangle++(0.04,0.04) (1.628,3.31)rectangle++(0.04,0.04) (1.514,4.345)rectangle++(0.04,0.04)
  (1.458,5.229)rectangle++(0.04,0.04) (1.446,5.982)rectangle++(0.04,0.04) (1.467,6.623)rectangle++(0.04,0.04) (1.509,7.167)rectangle++(0.04,0.04) (1.567,7.628)rectangle++(0.04,0.04)
  (1.883,2.668)rectangle++(0.04,0.04) (1.705,3.789)rectangle++(0.04,0.04) (1.599,4.749)rectangle++(0.04,0.04) (1.547,5.569)rectangle++(0.04,0.04) (1.535,6.268)rectangle++(0.04,0.04)
  (1.554,6.862)rectangle++(0.04,0.04) (1.593,7.367)rectangle++(0.04,0.04) (1.646,7.794)rectangle++(0.04,0.04) (1.709,8.156)rectangle++(0.04,0.04) (1.777,8.46)rectangle++(0.04,0.04)
  (1.846,8.716)rectangle++(0.04,0.04) (1.916,8.931)rectangle++(0.04,0.04) (1.984,9.111)rectangle++(0.04,0.04) (2.048,9.261)rectangle++(0.04,0.04) (2.109,9.386)rectangle++(0.04,0.04)
  (-0.02,1.485)rectangle++(0.04,0.04) (0.04,2.859)rectangle++(0.04,0.04) (0.146,4.025)rectangle++(0.04,0.04) (0.283,5.012)rectangle++(0.04,0.04) (0.439,5.845)rectangle++(0.04,0.04)
  (0.605,6.547)rectangle++(0.04,0.04) (1.725,2.158)rectangle++(0.04,0.04) (1.55,3.362)rectangle++(0.04,0.04) (1.45,4.392)rectangle++(0.04,0.04) (1.406,5.271)rectangle++(0.04,0.04)
  (1.404,6.02)rectangle++(0.04,0.04) (1.432,6.657)rectangle++(0.04,0.04) (1.481,7.198)rectangle++(0.04,0.04) (1.545,7.655)rectangle++(0.04,0.04) (1.617,8.041)rectangle++(0.04,0.04)
  (1.694,8.367)rectangle++(0.04,0.04) (1.772,8.64)rectangle++(0.04,0.04) (1.85,8.869)rectangle++(0.04,0.04) (1.925,9.061)rectangle++(0.04,0.04) (1.996,9.221)rectangle++(0.04,0.04)
  (2.064,9.354)rectangle++(0.04,0.04) (2.126,9.465)rectangle++(0.04,0.04) (2.184,9.556)rectangle++(0.04,0.04) (2.331,3.291)rectangle++(0.04,0.04) (2.111,4.301)rectangle++(0.04,0.04)
  (1.964,5.167)rectangle++(0.04,0.04) (1.874,5.91)rectangle++(0.04,0.04) (1.827,6.545)rectangle++(0.04,0.04) (1.812,7.086)rectangle++(0.04,0.04) (1.822,7.547)rectangle++(0.04,0.04)
  (1.848,7.938)rectangle++(0.04,0.04) (1.886,8.27)rectangle++(0.04,0.04) (1.932,8.55)rectangle++(0.04,0.04) (1.982,8.786)rectangle++(0.04,0.04) (2.034,8.985)rectangle++(0.04,0.04)
  (2.086,9.152)rectangle++(0.04,0.04) (2.137,9.292)rectangle++(0.04,0.04) (-2.01,4.125)rectangle++(0.04,0.04) (-1.546,5.183)rectangle++(0.04,0.04) (-1.109,6.063)rectangle++(0.04,0.04)
  (-0.702,6.794)rectangle++(0.04,0.04) (-0.327,7.399)rectangle++(0.04,0.04) (0.016,7.899)rectangle++(0.04,0.04) (0.327,8.31)rectangle++(0.04,0.04) (0.608,8.646)rectangle++(0.04,0.04)
  (0.861,8.921)rectangle++(0.04,0.04) (1.086,9.145)rectangle++(0.04,0.04) (1.287,9.326)rectangle++(0.04,0.04) (1.465,9.472)rectangle++(0.04,0.04) (1.622,9.589)rectangle++(0.04,0.04)
  (1.76,9.682)rectangle++(0.04,0.04) (2.429,3.211)rectangle++(0.04,0.04) (2.191,4.228)rectangle++(0.04,0.04) (2.03,5.103)rectangle++(0.04,0.04) (1.107,2.182)rectangle++(0.04,0.04)
  (0.428,1.242)rectangle++(0.04,0.04) (0.411,2.634)rectangle++(0.04,0.04) (0.452,3.819)rectangle++(0.04,0.04) (0.535,4.824)rectangle++(0.04,0.04) (0.646,5.675)rectangle++(0.04,0.04)
  (0.774,6.394)rectangle++(0.04,0.04) (0.911,7.001)rectangle++(0.04,0.04) (1.052,7.51)rectangle++(0.04,0.04) (1.193,7.938)rectangle++(0.04,0.04) (1.329,8.296)rectangle++(0.04,0.04)
  (1.459,8.594)rectangle++(0.04,0.04) (1.582,8.843)rectangle++(0.04,0.04) (1.696,9.049)rectangle++(0.04,0.04) (1.802,9.22)rectangle++(0.04,0.04) (-2.058,4.032)rectangle++(0.04,0.04)
  (-1.59,5.106)rectangle++(0.04,0.04) (-1.667,2.346)rectangle++(0.04,0.04) (-1.325,3.657)rectangle++(0.04,0.04) (-0.982,4.758)rectangle++(0.04,0.04) (-0.647,5.68)rectangle++(0.04,0.04)
  (-0.325,6.45)rectangle++(0.04,0.04) (-0.02,7.092)rectangle++(0.04,0.04) (-1.869,3.144)rectangle++(0.04,0.04) (-1.465,4.344)rectangle++(0.04,0.04) (-1.074,5.347)rectangle++(0.04,0.04)
  (-0.701,6.184)rectangle++(0.04,0.04) (-0.351,6.881)rectangle++(0.04,0.04) (-0.025,7.459)rectangle++(0.04,0.04) (0.275,7.937)rectangle++(0.04,0.04) (0.549,8.332)rectangle++(0.04,0.04)
  (0.798,8.656)rectangle++(0.04,0.04) (1.022,8.922)rectangle++(0.04,0.04) (1.223,9.139)rectangle++(0.04,0.04) (1.403,9.316)rectangle++(0.04,0.04) (1.563,9.458)rectangle++(0.04,0.04)
  (1.705,9.573)rectangle++(0.04,0.04) (1.83,9.665)rectangle++(0.04,0.04) (2.414,3.225)rectangle++(0.04,0.04) (2.179,4.241)rectangle++(0.04,0.04) (2.02,5.114)rectangle++(0.04,0.04)
  (-0.947,3.179)rectangle++(0.04,0.04) (-0.68,4.336)rectangle++(0.04,0.04) (-0.407,5.309)rectangle++(0.04,0.04) (-0.136,6.125)rectangle++(0.04,0.04) (1.718,1.865)rectangle++(0.04,0.04)
  (1.533,3.113)rectangle++(0.04,0.04) (1.425,4.181)rectangle++(0.04,0.04) (1.376,5.093)rectangle++(0.04,0.04) (1.371,5.87)rectangle++(0.04,0.04) (1.398,6.531)rectangle++(0.04,0.04)
  (1.448,7.091)rectangle++(0.04,0.04) (1.512,7.566)rectangle++(0.04,0.04) (1.586,7.967)rectangle++(0.04,0.04) (1.664,8.305)rectangle++(0.04,0.04) (1.745,8.589)rectangle++(0.04,0.04)
  (1.824,8.827)rectangle++(0.04,0.04) (1.901,9.026)rectangle++(0.04,0.04) (1.975,9.192)rectangle++(0.04,0.04) (2.044,9.331)rectangle++(0.04,0.04) (2.109,9.445)rectangle++(0.04,0.04)
  (2.168,9.54)rectangle++(0.04,0.04) (2.222,9.619)rectangle++(0.04,0.04) (2.271,9.683)rectangle++(0.04,0.04) (2.316,9.736)rectangle++(0.04,0.04) (2.356,9.779)rectangle++(0.04,0.04)
  (2.391,9.814)rectangle++(0.04,0.04) (2.423,9.843)rectangle++(0.04,0.04) (2.451,9.866)rectangle++(0.04,0.04) (2.476,9.884)rectangle++(0.04,0.04) (2.498,9.899)rectangle++(0.04,0.04)
  (2.517,9.91)rectangle++(0.04,0.04) (2.533,9.919)rectangle++(0.04,0.04) (2.548,9.926)rectangle++(0.04,0.04) (2.561,9.931)rectangle++(0.04,0.04) (2.572,9.936)rectangle++(0.04,0.04)
  (2.581,9.939)rectangle++(0.04,0.04) (2.589,9.941)rectangle++(0.04,0.04) (2.596,9.942)rectangle++(0.04,0.04) (2.602,9.943)rectangle++(0.04,0.04) (2.608,9.944)rectangle++(0.04,0.04)
  (2.612,9.944)rectangle++(0.04,0.04) (2.616,9.944)rectangle++(0.04,0.04) (-2.084,4.378)rectangle++(0.04,0.04) (-1.598,5.401)rectangle++(0.04,0.04) (-1.145,6.251)rectangle++(0.04,0.04)
  (-0.725,6.955)rectangle++(0.04,0.04) (-0.34,7.537)rectangle++(0.04,0.04) (0.01,8.017)rectangle++(0.04,0.04) (0.327,8.41)rectangle++(0.04,0.04) (0.612,8.731)rectangle++(0.04,0.04)
  (-2.169,3.651)rectangle++(0.04,0.04) (-1.7,4.786)rectangle++(0.04,0.04) (-1.256,5.732)rectangle++(0.04,0.04) (-0.84,6.519)rectangle++(0.04,0.04) (-0.456,7.171)rectangle++(0.04,0.04)
  (-0.103,7.71)rectangle++(0.04,0.04) (2.157,2.254)rectangle++(0.04,0.04) (1.921,3.426)rectangle++(0.04,0.04) (1.768,4.431)rectangle++(0.04,0.04) (1.678,5.292)rectangle++(0.04,0.04)
  (1.635,6.027)rectangle++(0.04,0.04) (1.629,6.654)rectangle++(0.04,0.04) (1.649,7.187)rectangle++(0.04,0.04) (1.687,7.639)rectangle++(0.04,0.04) (1.737,8.022)rectangle++(0.04,0.04)
  (1.795,8.345)rectangle++(0.04,0.04) (2.05,2.9)rectangle++(0.04,0.04) (1.856,3.979)rectangle++(0.04,0.04) (0.818,1.868)rectangle++(0.04,0.04) (-0.02,0.282)rectangle++(0.04,0.04)
  (-0.008,1.837)rectangle++(0.04,0.04) (0.065,3.158)rectangle++(0.04,0.04) (0.179,4.278)rectangle++(0.04,0.04) (0.321,5.225)rectangle++(0.04,0.04) (0.48,6.025)rectangle++(0.04,0.04)
  (0.647,6.698)rectangle++(0.04,0.04) (0.815,7.264)rectangle++(0.04,0.04) (1.894,2.385)rectangle++(0.04,0.04) (1.703,3.548)rectangle++(0.04,0.04) (1.587,4.544)rectangle++(0.04,0.04)
  (1.529,5.395)rectangle++(0.04,0.04) (1.513,6.121)rectangle++(0.04,0.04) (1.529,6.738)rectangle++(0.04,0.04) (1.567,7.263)rectangle++(0.04,0.04) (1.62,7.707)rectangle++(0.04,0.04)
  (1.683,8.082)rectangle++(0.04,0.04) (1.752,8.399)rectangle++(0.04,0.04) (1.823,8.665)rectangle++(0.04,0.04) (1.894,8.889)rectangle++(0.04,0.04) (1.963,9.076)rectangle++(0.04,0.04)
  (2.229,3.119)rectangle++(0.04,0.04) (2.017,4.158)rectangle++(0.04,0.04) (1.879,5.05)rectangle++(0.04,0.04) (-0.958,3.132)rectangle++(0.04,0.04) (-0.691,4.297)rectangle++(0.04,0.04)
  (-0.418,5.276)rectangle++(0.04,0.04) (-0.147,6.098)rectangle++(0.04,0.04) (-1.636,2.897)rectangle++(0.04,0.04) (-0.02,0.447)rectangle++(0.04,0.04) (-0.141,1.683)rectangle++(0.04,0.04)
  (-0.055,3.032)rectangle++(0.04,0.04) (0.072,4.176)rectangle++(0.04,0.04) (0.226,5.143)rectangle++(0.04,0.04) (0.396,5.958)rectangle++(0.04,0.04) (0.573,6.645)rectangle++(0.04,0.04)
  (1.757,2.174)rectangle++(0.04,0.04) (1.578,3.374)rectangle++(0.04,0.04) (0.69,1.65)rectangle++(0.04,0.04) (0.651,2.971)rectangle++(0.04,0.04) (0.67,4.096)rectangle++(0.04,0.04)
  (0.731,5.051)rectangle++(0.04,0.04) (0.821,5.86)rectangle++(0.04,0.04) (0.93,6.544)rectangle++(0.04,0.04) (1.05,7.122)rectangle++(0.04,0.04) (-1.663,3.397)rectangle++(0.04,0.04)
  (-1.28,4.55)rectangle++(0.04,0.04) (-0.908,5.515)rectangle++(0.04,0.04) (1.663,1.518)rectangle++(0.04,0.04) (1.472,2.82)rectangle++(0.04,0.04) (1.362,3.934)rectangle++(0.04,0.04)
  (1.313,4.886)rectangle++(0.04,0.04) (1.309,5.696)rectangle++(0.04,0.04) (1.338,6.386)rectangle++(0.04,0.04) (1.391,6.971)rectangle++(0.04,0.04) (1.726,2.465)rectangle++(0.04,0.04)
  (1.563,3.622)rectangle++(0.04,0.04) (-0.65,2.745)rectangle++(0.04,0.04) (-0.445,3.956)rectangle++(0.04,0.04) (-0.222,4.976)rectangle++(0.04,0.04) (0.008,5.835)rectangle++(0.04,0.04)
  (0.238,6.556)rectangle++(0.04,0.04) (1.783,2.065)rectangle++(0.04,0.04) (1.596,3.28)rectangle++(0.04,0.04) (1.485,4.321)rectangle++(0.04,0.04) (1.433,5.209)rectangle++(0.04,0.04)
  (1.424,5.967)rectangle++(0.04,0.04) (1.447,6.611)rectangle++(0.04,0.04) (1.492,7.158)rectangle++(0.04,0.04) (1.553,7.621)rectangle++(0.04,0.04) (1.622,8.012)rectangle++(0.04,0.04)
  (1.697,8.341)rectangle++(0.04,0.04) (1.774,8.618)rectangle++(0.04,0.04) (-0.02,1.362)rectangle++(0.04,0.04) (0.035,2.755)rectangle++(0.04,0.04) (0.138,3.936)rectangle++(0.04,0.04)
  (0.273,4.937)rectangle++(0.04,0.04) (0.427,5.781)rectangle++(0.04,0.04) (0.592,6.493)rectangle++(0.04,0.04) (0.761,7.092)rectangle++(0.04,0.04) (0.928,7.594)rectangle++(0.04,0.04)
  (1.09,8.014)rectangle++(0.04,0.04) (2.063,2.637)rectangle++(0.04,0.04) (-0.02,0.405)rectangle++(0.04,0.04) (-0.003,1.941)rectangle++(0.04,0.04) (0.073,3.246)rectangle++(0.04,0.04)
  (0.19,4.353)rectangle++(0.04,0.04) (1.173,1.524)rectangle++(0.04,0.04) (0.233,1.101)rectangle++(0.04,0.04) (0.24,2.522)rectangle++(0.04,0.04) (0.303,3.731)rectangle++(0.04,0.04)
  (-0.931,2.479)rectangle++(0.04,0.04) (-0.694,3.741)rectangle++(0.04,0.04) (-0.442,4.804)rectangle++(0.04,0.04) (-0.186,5.697)rectangle++(0.04,0.04) (0.067,6.446)rectangle++(0.04,0.04)
  (0.313,7.073)rectangle++(0.04,0.04) (0.547,7.596)rectangle++(0.04,0.04) (0.766,8.031)rectangle++(0.04,0.04) (0.97,8.391)rectangle++(0.04,0.04) (1.158,8.69)rectangle++(0.04,0.04)
  (1.33,8.936)rectangle++(0.04,0.04) (1.486,9.139)rectangle++(0.04,0.04) (1.626,9.305)rectangle++(0.04,0.04) (1.752,9.44)rectangle++(0.04,0.04) (1.865,9.55)rectangle++(0.04,0.04)
  (1.965,9.639)rectangle++(0.04,0.04) (2.054,9.711)rectangle++(0.04,0.04) (2.132,9.769)rectangle++(0.04,0.04) (2.201,9.814)rectangle++(0.04,0.04) (2.261,9.85)rectangle++(0.04,0.04)
  (2.314,9.878)rectangle++(0.04,0.04) (2.359,9.9)rectangle++(0.04,0.04) (2.399,9.917)rectangle++(0.04,0.04) (2.434,9.93)rectangle++(0.04,0.04) (-2.116,4.333)rectangle++(0.04,0.04)
  (-1.628,5.364)rectangle++(0.04,0.04) (-1.171,6.221)rectangle++(0.04,0.04) (-0.749,6.931)rectangle++(0.04,0.04) (-0.361,7.517)rectangle++(0.04,0.04) (2.142,2.14)rectangle++(0.04,0.04)
  (1.904,3.33)rectangle++(0.04,0.04) (1.749,4.35)rectangle++(0.04,0.04) (1.659,5.224)rectangle++(0.04,0.04) (1.617,5.97)rectangle++(0.04,0.04) (-1.25,3.274)rectangle++(0.04,0.04)
  (-0.934,4.429)rectangle++(0.04,0.04) (-0.619,5.398)rectangle++(0.04,0.04) (-0.312,6.21)rectangle++(0.04,0.04) (-0.019,6.887)rectangle++(0.04,0.04) (1.914,2.078)rectangle++(0.04,0.04)
  (1.708,3.286)rectangle++(0.04,0.04) (1.581,4.321)rectangle++(0.04,0.04) (-0.828,2.903)rectangle++(0.04,0.04) (-0.942,2.037)rectangle++(0.04,0.04) (-0.721,3.365)rectangle++(0.04,0.04)
  (-0.481,4.486)rectangle++(0.04,0.04) (-0.231,5.428)rectangle++(0.04,0.04) (0.018,6.219)rectangle++(0.04,0.04) (0.262,6.882)rectangle++(0.04,0.04) (0.496,7.435)rectangle++(0.04,0.04)
  (0.717,7.896)rectangle++(0.04,0.04) (-1.931,3.491)rectangle++(0.04,0.04) (-1.504,4.641)rectangle++(0.04,0.04) (-1.095,5.601)rectangle++(0.04,0.04) (-1.696,2.569)rectangle++(0.04,0.04)
  (-1.029,1.764)rectangle++(0.04,0.04) (-0.806,3.137)rectangle++(0.04,0.04) (-0.562,4.295)rectangle++(0.04,0.04) (-0.308,5.269)rectangle++(0.04,0.04) (-0.053,6.087)rectangle++(0.04,0.04)
  (0.196,6.773)rectangle++(0.04,0.04) (0.435,7.345)rectangle++(0.04,0.04) (0.662,7.822)rectangle++(0.04,0.04) (0.873,8.218)rectangle++(0.04,0.04) (1.069,8.547)rectangle++(0.04,0.04)
  (1.248,8.818)rectangle++(0.04,0.04) (1.411,9.042)rectangle++(0.04,0.04) (1.559,9.225)rectangle++(0.04,0.04) (1.692,9.375)rectangle++(0.04,0.04) (1.811,9.498)rectangle++(0.04,0.04)
  (1.917,9.597)rectangle++(0.04,0.04) (2.011,9.677)rectangle++(0.04,0.04) (2.39,3.275)rectangle++(0.04,0.04) (2.161,4.285)rectangle++(0.04,0.04) (2.006,5.152)rectangle++(0.04,0.04)
  (1.909,5.895)rectangle++(0.04,0.04) (1.856,6.531)rectangle++(0.04,0.04) (1.837,7.073)rectangle++(0.04,0.04) (-1.493,3.567)rectangle++(0.04,0.04) (-1.128,4.688)rectangle++(0.04,0.04)
  (1.465,1.262)rectangle++(0.04,0.04) (1.293,2.61)rectangle++(0.04,0.04) (0.519,1.393)rectangle++(0.04,0.04) (0.495,2.759)rectangle++(0.04,0.04) (0.529,3.922)rectangle++(0.04,0.04)
  (0.604,4.908)rectangle++(0.04,0.04) (0.708,5.744)rectangle++(0.04,0.04) (0.829,6.45)rectangle++(0.04,0.04) (0.961,7.046)rectangle++(0.04,0.04) (1.096,7.547)rectangle++(0.04,0.04)
  (-1.764,3.501)rectangle++(0.04,0.04) (-1.362,4.643)rectangle++(0.04,0.04) (-0.974,5.597)rectangle++(0.04,0.04) (-0.606,6.393)rectangle++(0.04,0.04) (-0.262,7.054)rectangle++(0.04,0.04)
  (0.058,7.603)rectangle++(0.04,0.04) (0.351,8.056)rectangle++(0.04,0.04) (0.618,8.43)rectangle++(0.04,0.04) (0.861,8.737)rectangle++(0.04,0.04) (1.079,8.988)rectangle++(0.04,0.04)
  (1.274,9.193)rectangle++(0.04,0.04) (1.449,9.359)rectangle++(0.04,0.04) (1.603,9.494)rectangle++(0.04,0.04) (1.741,9.602)rectangle++(0.04,0.04) (1.861,9.688)rectangle++(0.04,0.04)
  (1.967,9.757)rectangle++(0.04,0.04) (-2.164,4.188)rectangle++(0.04,0.04) (-1.674,5.243)rectangle++(0.04,0.04) (-1.216,6.119)rectangle++(0.04,0.04) (-0.791,6.846)rectangle++(0.04,0.04)
  (-0.401,7.447)rectangle++(0.04,0.04) (-0.045,7.942)rectangle++(0.04,0.04) (0.277,8.349)rectangle++(0.04,0.04) (0.568,8.682)rectangle++(0.04,0.04) (0.828,8.953)rectangle++(0.04,0.04)
  (1.059,9.173)rectangle++(0.04,0.04) (1.265,9.351)rectangle++(0.04,0.04) (1.447,9.494)rectangle++(0.04,0.04) (2.424,3.085)rectangle++(0.04,0.04) (2.181,4.121)rectangle++(0.04,0.04)
  (2.017,5.012)rectangle++(0.04,0.04) (1.913,5.776)rectangle++(0.04,0.04) (1.855,6.429)rectangle++(0.04,0.04) (1.831,6.987)rectangle++(0.04,0.04) (1.834,7.462)rectangle++(0.04,0.04)
  (1.855,7.865)rectangle++(0.04,0.04) (1.889,8.207)rectangle++(0.04,0.04) (1.932,8.497)rectangle++(0.04,0.04) (1.98,8.741)rectangle++(0.04,0.04) (2.03,8.947)rectangle++(0.04,0.04);
\fill[nuage]
  (-1.941,4.024)rectangle++(0.04,0.04) (-1.491,5.095)rectangle++(0.04,0.04) (-1.066,5.986)rectangle++(0.04,0.04) (-0.669,6.727)rectangle++(0.04,0.04) (-0.302,7.341)rectangle++(0.04,0.04)
  (0.035,7.848)rectangle++(0.04,0.04) (0.341,8.266)rectangle++(0.04,0.04) (0.619,8.608)rectangle++(0.04,0.04) (0.868,8.889)rectangle++(0.04,0.04) (1.091,9.117)rectangle++(0.04,0.04)
  (1.29,9.302)rectangle++(0.04,0.04) (1.466,9.451)rectangle++(0.04,0.04) (1.622,9.571)rectangle++(0.04,0.04) (1.76,9.667)rectangle++(0.04,0.04) (-2.183,4.12)rectangle++(0.04,0.04)
  (-1.693,5.186)rectangle++(0.04,0.04) (-1.233,6.072)rectangle++(0.04,0.04) (-0.808,6.807)rectangle++(0.04,0.04) (-0.417,7.414)rectangle++(0.04,0.04) (-0.06,7.915)rectangle++(0.04,0.04)
  (0.264,8.326)rectangle++(0.04,0.04) (0.555,8.663)rectangle++(0.04,0.04) (0.816,8.938)rectangle++(0.04,0.04) (1.049,9.16)rectangle++(0.04,0.04) (-0.02,1.449)rectangle++(0.04,0.04)
  (0.391,0.773)rectangle++(0.04,0.04) (0.361,2.237)rectangle++(0.04,0.04) (0.394,3.483)rectangle++(0.04,0.04) (-0.02,0.541)rectangle++(0.04,0.04) (0.002,2.056)rectangle++(0.04,0.04)
  (0.082,3.344)rectangle++(0.04,0.04) (-0.874,2.344)rectangle++(0.04,0.04) (-0.652,3.623)rectangle++(0.04,0.04) (-0.411,4.702)rectangle++(0.04,0.04) (-0.164,5.609)rectangle++(0.04,0.04)
  (0.083,6.371)rectangle++(0.04,0.04) (0.323,7.008)rectangle++(0.04,0.04) (0.553,7.54)rectangle++(0.04,0.04) (0.769,7.983)rectangle++(0.04,0.04) (0.971,8.351)rectangle++(0.04,0.04)
  (1.157,8.656)rectangle++(0.04,0.04) (2.233,2.808)rectangle++(0.04,0.04) (2.008,3.894)rectangle++(0.04,0.04) (1.86,4.826)rectangle++(0.04,0.04) (1.772,5.624)rectangle++(0.04,0.04)
  (1.729,6.305)rectangle++(0.04,0.04) (1.72,6.887)rectangle++(0.04,0.04) (1.735,7.381)rectangle++(0.04,0.04) (1.768,7.801)rectangle++(0.04,0.04) (1.812,8.156)rectangle++(0.04,0.04)
  (1.865,8.456)rectangle++(0.04,0.04) (1.921,8.71)rectangle++(0.04,0.04) (1.979,8.922)rectangle++(0.04,0.04) (2.037,9.101)rectangle++(0.04,0.04) (-1.98,4.06)rectangle++(0.04,0.04)
  (-1.523,5.126)rectangle++(0.04,0.04) (-1.092,6.014)rectangle++(0.04,0.04) (-0.69,6.752)rectangle++(0.04,0.04) (-0.318,7.363)rectangle++(0.04,0.04) (0.022,7.868)rectangle++(0.04,0.04)
  (0.331,8.283)rectangle++(0.04,0.04) (0.611,8.623)rectangle++(0.04,0.04) (0.862,8.902)rectangle++(0.04,0.04) (1.086,9.128)rectangle++(0.04,0.04) (1.286,9.312)rectangle++(0.04,0.04)
  (1.464,9.46)rectangle++(0.04,0.04) (1.62,9.578)rectangle++(0.04,0.04) (1.758,9.673)rectangle++(0.04,0.04) (1.879,9.748)rectangle++(0.04,0.04) (1.985,9.807)rectangle++(0.04,0.04)
  (2.431,3.3)rectangle++(0.04,0.04) (2.196,4.304)rectangle++(0.04,0.04) (2.036,5.167)rectangle++(0.04,0.04) (1.935,5.906)rectangle++(0.04,0.04) (1.346,2.351)rectangle++(0.04,0.04)
  (1.236,3.54)rectangle++(0.04,0.04) (1.19,4.556)rectangle++(0.04,0.04) (1.192,5.421)rectangle++(0.04,0.04) (1.227,6.157)rectangle++(0.04,0.04) (1.287,6.78)rectangle++(0.04,0.04)
  (1.363,7.308)rectangle++(0.04,0.04) (1.449,7.753)rectangle++(0.04,0.04) (1.54,8.129)rectangle++(0.04,0.04) (1.632,8.444)rectangle++(0.04,0.04) (1.722,8.708)rectangle++(0.04,0.04)
  (1.81,8.929)rectangle++(0.04,0.04) (1.894,9.114)rectangle++(0.04,0.04) (1.972,9.267)rectangle++(0.04,0.04) (-0.02,1.466)rectangle++(0.04,0.04) (0.039,2.843)rectangle++(0.04,0.04)
  (0.145,4.011)rectangle++(0.04,0.04) (0.282,5)rectangle++(0.04,0.04) (0.437,5.835)rectangle++(0.04,0.04) (0.603,6.538)rectangle++(0.04,0.04) (0.772,7.13)rectangle++(0.04,0.04)
  (0.939,7.626)rectangle++(0.04,0.04) (1.101,8.04)rectangle++(0.04,0.04) (1.255,8.386)rectangle++(0.04,0.04) (1.4,8.675)rectangle++(0.04,0.04) (1.535,8.914)rectangle++(0.04,0.04)
  (1.659,9.111)rectangle++(0.04,0.04) (1.772,9.274)rectangle++(0.04,0.04) (1.875,9.409)rectangle++(0.04,0.04) (1.968,9.518)rectangle++(0.04,0.04) (2.051,9.608)rectangle++(0.04,0.04)
  (2.126,9.681)rectangle++(0.04,0.04) (2.192,9.74)rectangle++(0.04,0.04) (2.251,9.788)rectangle++(0.04,0.04) (2.302,9.826)rectangle++(0.04,0.04) (2.388,3.387)rectangle++(0.04,0.04)
  (2.163,4.379)rectangle++(0.04,0.04) (2.012,5.232)rectangle++(0.04,0.04) (1.917,5.963)rectangle++(0.04,0.04) (1.866,6.588)rectangle++(0.04,0.04) (1.847,7.121)rectangle++(0.04,0.04)
  (-1.503,3.581)rectangle++(0.04,0.04) (-1.137,4.7)rectangle++(0.04,0.04) (-0.78,5.637)rectangle++(0.04,0.04) (1.678,1.58)rectangle++(0.04,0.04) (1.487,2.872)rectangle++(0.04,0.04)
  (1.377,3.978)rectangle++(0.04,0.04) (1.327,4.922)rectangle++(0.04,0.04) (1.323,5.727)rectangle++(0.04,0.04) (1.351,6.411)rectangle++(0.04,0.04) (1.403,6.992)rectangle++(0.04,0.04)
  (1.47,7.483)rectangle++(0.04,0.04) (1.547,7.898)rectangle++(0.04,0.04) (1.628,8.248)rectangle++(0.04,0.04) (1.712,8.542)rectangle++(0.04,0.04) (1.794,8.788)rectangle++(0.04,0.04)
  (-1.947,3.935)rectangle++(0.04,0.04) (-1.5,5.019)rectangle++(0.04,0.04) (-1.076,5.922)rectangle++(0.04,0.04) (-0.68,6.673)rectangle++(0.04,0.04) (-0.313,7.296)rectangle++(0.04,0.04)
  (0.023,7.81)rectangle++(0.04,0.04) (0.33,8.234)rectangle++(0.04,0.04) (-2.096,3.476)rectangle++(0.04,0.04) (-1.645,4.635)rectangle++(0.04,0.04) (-1.215,5.602)rectangle++(0.04,0.04)
  (-0.811,6.406)rectangle++(0.04,0.04) (-0.435,7.074)rectangle++(0.04,0.04) (-1.947,3.045)rectangle++(0.04,0.04) (-1.536,4.263)rectangle++(0.04,0.04) (-1.137,5.281)rectangle++(0.04,0.04)
  (-0.757,6.13)rectangle++(0.04,0.04) (-0.401,6.837)rectangle++(0.04,0.04) (-0.069,7.424)rectangle++(0.04,0.04) (0.236,7.909)rectangle++(0.04,0.04) (0.515,8.31)rectangle++(0.04,0.04)
  (0.768,8.639)rectangle++(0.04,0.04) (0.996,8.909)rectangle++(0.04,0.04) (-2.138,3.778)rectangle++(0.04,0.04) (-1.431,1.928)rectangle++(0.04,0.04) (-1.142,3.293)rectangle++(0.04,0.04)
  (-0.841,4.441)rectangle++(0.04,0.04) (1.352,1.277)rectangle++(0.04,0.04) (1.198,2.628)rectangle++(0.04,0.04) (1.121,3.782)rectangle++(0.04,0.04) (1.102,4.766)rectangle++(0.04,0.04)
  (1.125,5.603)rectangle++(0.04,0.04) (1.178,6.314)rectangle++(0.04,0.04) (1.252,6.916)rectangle++(0.04,0.04) (1.731,2.415)rectangle++(0.04,0.04) (1.566,3.58)rectangle++(0.04,0.04)
  (1.472,4.577)rectangle++(0.04,0.04) (1.432,5.427)rectangle++(0.04,0.04) (1.432,6.152)rectangle++(0.04,0.04) (1.461,6.768)rectangle++(0.04,0.04) (1.511,7.291)rectangle++(0.04,0.04)
  (1.573,7.733)rectangle++(0.04,0.04) (1.645,8.106)rectangle++(0.04,0.04) (1.72,8.421)rectangle++(0.04,0.04) (2.082,2.898)rectangle++(0.04,0.04) (-0.358,2.706)rectangle++(0.04,0.04)
  (-0.198,3.91)rectangle++(0.04,0.04) (-0.014,4.928)rectangle++(0.04,0.04) (0.183,5.785)rectangle++(0.04,0.04) (0.384,6.507)rectangle++(0.04,0.04) (-1.636,3.109)rectangle++(0.04,0.04)
  (-1.268,4.304)rectangle++(0.04,0.04) (-0.908,5.305)rectangle++(0.04,0.04) (-0.562,6.142)rectangle++(0.04,0.04) (-0.234,6.84)rectangle++(0.04,0.04) (-1.846,3.04)rectangle++(0.04,0.04)
  (-1.45,4.254)rectangle++(0.04,0.04) (-1.065,5.27)rectangle++(0.04,0.04) (-0.696,6.118)rectangle++(0.04,0.04) (-0.349,6.825)rectangle++(0.04,0.04) (-0.026,7.411)rectangle++(0.04,0.04)
  (-0.02,1.169)rectangle++(0.04,0.04) (0.028,2.591)rectangle++(0.04,0.04) (0.125,3.797)rectangle++(0.04,0.04) (-0.983,2.453)rectangle++(0.04,0.04) (-0.74,3.721)rectangle++(0.04,0.04)
  (1.135,1.131)rectangle++(0.04,0.04) (1.008,2.512)rectangle++(0.04,0.04) (0.955,3.691)rectangle++(0.04,0.04) (0.957,4.695)rectangle++(0.04,0.04) (0.999,5.549)rectangle++(0.04,0.04)
  (1.069,6.273)rectangle++(0.04,0.04) (1.158,6.885)rectangle++(0.04,0.04) (1.257,7.402)rectangle++(0.04,0.04) (1.362,7.838)rectangle++(0.04,0.04) (1.469,8.204)rectangle++(0.04,0.04)
  (1.575,8.511)rectangle++(0.04,0.04) (1.677,8.767)rectangle++(0.04,0.04) (1.774,8.981)rectangle++(0.04,0.04) (1.865,9.159)rectangle++(0.04,0.04) (1.949,9.307)rectangle++(0.04,0.04)
  (-2.051,4.085)rectangle++(0.04,0.04) (-1.582,5.15)rectangle++(0.04,0.04) (-1.141,6.037)rectangle++(0.04,0.04) (-1.819,2.655)rectangle++(0.04,0.04) (-1.442,3.926)rectangle++(0.04,0.04)
  (-1.071,4.991)rectangle++(0.04,0.04) (-0.713,5.881)rectangle++(0.04,0.04) (-0.373,6.624)rectangle++(0.04,0.04) (1.893,1.923)rectangle++(0.04,0.04) (1.684,3.155)rectangle++(0.04,0.04)
  (1.555,4.21)rectangle++(0.04,0.04) (0.928,1.845)rectangle++(0.04,0.04) (0.861,3.127)rectangle++(0.04,0.04) (0.854,4.22)rectangle++(0.04,0.04) (1.036,1.665)rectangle++(0.04,0.04)
  (0.945,2.97)rectangle++(0.04,0.04) (0.92,4.083)rectangle++(0.04,0.04) (0.988,1.649)rectangle++(0.04,0.04) (0.903,2.958)rectangle++(0.04,0.04) (0.884,4.075)rectangle++(0.04,0.04)
  (-0.904,2.689)rectangle++(0.04,0.04) (-0.663,3.918)rectangle++(0.04,0.04) (-0.409,4.953)rectangle++(0.04,0.04) (-0.152,5.822)rectangle++(0.04,0.04) (0.102,6.551)rectangle++(0.04,0.04)
  (0.346,7.161)rectangle++(0.04,0.04) (0.579,7.669)rectangle++(0.04,0.04) (0.796,8.092)rectangle++(0.04,0.04) (0.998,8.442)rectangle++(0.04,0.04) (1.184,8.732)rectangle++(0.04,0.04)
  (1.354,8.971)rectangle++(0.04,0.04) (1.507,9.168)rectangle++(0.04,0.04) (1.646,9.328)rectangle++(0.04,0.04) (1.77,9.459)rectangle++(0.04,0.04) (1.88,9.566)rectangle++(0.04,0.04)
  (1.979,9.652)rectangle++(0.04,0.04) (2.066,9.721)rectangle++(0.04,0.04) (2.143,9.777)rectangle++(0.04,0.04) (2.21,9.821)rectangle++(0.04,0.04) (2.269,9.855)rectangle++(0.04,0.04)
  (2.321,9.882)rectangle++(0.04,0.04) (-2.126,4.297)rectangle++(0.04,0.04) (-1.638,5.334)rectangle++(0.04,0.04) (-1.181,6.195)rectangle++(0.04,0.04) (1.894,1.61)rectangle++(0.04,0.04)
  (1.672,2.889)rectangle++(0.04,0.04) (1.535,3.985)rectangle++(0.04,0.04) (1.462,4.922)rectangle++(0.04,0.04) (1.437,5.721)rectangle++(0.04,0.04) (1.448,6.402)rectangle++(0.04,0.04)
  (1.558,2.343)rectangle++(0.04,0.04) (-0.319,2.463)rectangle++(0.04,0.04) (-0.175,3.702)rectangle++(0.04,0.04) (-0.003,4.75)rectangle++(0.04,0.04) (0.186,5.634)rectangle++(0.04,0.04)
  (0.381,6.378)rectangle++(0.04,0.04) (0.577,7.002)rectangle++(0.04,0.04) (0.768,7.525)rectangle++(0.04,0.04) (0.952,7.962)rectangle++(0.04,0.04) (1.125,8.325)rectangle++(0.04,0.04)
  (1.287,8.628)rectangle++(0.04,0.04) (1.437,8.878)rectangle++(0.04,0.04) (1.574,9.085)rectangle++(0.04,0.04) (1.699,9.256)rectangle++(0.04,0.04) (1.813,9.396)rectangle++(0.04,0.04)
  (1.914,9.51)rectangle++(0.04,0.04) (2.005,9.603)rectangle++(0.04,0.04) (2.086,9.679)rectangle++(0.04,0.04) (2.158,9.74)rectangle++(0.04,0.04) (2.222,9.788)rectangle++(0.04,0.04)
  (2.278,9.828)rectangle++(0.04,0.04) (2.327,9.858)rectangle++(0.04,0.04) (2.37,9.883)rectangle++(0.04,0.04) (-2.117,4.308)rectangle++(0.04,0.04) (-1.629,5.343)rectangle++(0.04,0.04)
  (-1.173,6.203)rectangle++(0.04,0.04) (-0.751,6.916)rectangle++(0.04,0.04) (-0.364,7.505)rectangle++(0.04,0.04) (-2.045,3.156)rectangle++(0.04,0.04) (-1.614,4.361)rectangle++(0.04,0.04)
  (-1.2,5.367)rectangle++(0.04,0.04) (-0.807,6.207)rectangle++(0.04,0.04) (-0.44,6.904)rectangle++(0.04,0.04) (1.982,1.972)rectangle++(0.04,0.04) (1.761,3.194)rectangle++(0.04,0.04)
  (1.623,4.24)rectangle++(0.04,0.04) (1.547,5.136)rectangle++(0.04,0.04) (1.518,5.9)rectangle++(0.04,0.04) (1.524,6.55)rectangle++(0.04,0.04) (1.555,7.103)rectangle++(0.04,0.04)
  (1.604,7.571)rectangle++(0.04,0.04) (1.664,7.968)rectangle++(0.04,0.04) (1.731,8.302)rectangle++(0.04,0.04) (2.048,2.873)rectangle++(0.04,0.04) (1.853,3.956)rectangle++(0.04,0.04)
  (1.731,4.885)rectangle++(0.04,0.04) (1.091,2.052)rectangle++(0.04,0.04) (1.007,3.297)rectangle++(0.04,0.04) (0.986,4.358)rectangle++(0.04,0.04) (1.01,5.261)rectangle++(0.04,0.04)
  (1.067,6.028)rectangle++(0.04,0.04) (1.146,6.677)rectangle++(0.04,0.04) (1.239,7.226)rectangle++(0.04,0.04) (1.82,2.486)rectangle++(0.04,0.04) (1.644,3.637)rectangle++(0.04,0.04)
  (1.541,4.622)rectangle++(0.04,0.04) (1.493,5.463)rectangle++(0.04,0.04) (1.485,6.18)rectangle++(0.04,0.04) (1.507,6.79)rectangle++(0.04,0.04) (1.551,7.307)rectangle++(0.04,0.04)
  (1.608,7.745)rectangle++(0.04,0.04) (1.674,8.115)rectangle++(0.04,0.04) (1.746,8.427)rectangle++(0.04,0.04) (1.819,8.69)rectangle++(0.04,0.04) (1.891,8.91)rectangle++(0.04,0.04)
  (1.962,9.094)rectangle++(0.04,0.04) (2.029,9.247)rectangle++(0.04,0.04) (2.092,9.375)rectangle++(0.04,0.04) (2.151,9.482)rectangle++(0.04,0.04) (2.206,9.569)rectangle++(0.04,0.04)
  (2.255,9.642)rectangle++(0.04,0.04) (2.301,9.702)rectangle++(0.04,0.04) (2.341,9.751)rectangle++(0.04,0.04) (2.378,9.791)rectangle++(0.04,0.04) (2.367,3.398)rectangle++(0.04,0.04)
  (2.146,4.39)rectangle++(0.04,0.04) (-0.733,3.048)rectangle++(0.04,0.04) (-0.504,4.217)rectangle++(0.04,0.04) (-0.262,5.2)rectangle++(0.04,0.04) (-0.017,6.027)rectangle++(0.04,0.04)
  (0.225,6.72)rectangle++(0.04,0.04) (0.458,7.299)rectangle++(0.04,0.04) (0.679,7.782)rectangle++(0.04,0.04) (-1.909,3.457)rectangle++(0.04,0.04) (-1.486,4.611)rectangle++(0.04,0.04)
  (-1.081,5.575)rectangle++(0.04,0.04) (-1.687,2.567)rectangle++(0.04,0.04) (-1.333,3.846)rectangle++(0.04,0.04) (-0.982,4.918)rectangle++(0.04,0.04) (-1.496,2.445)rectangle++(0.04,0.04)
  (-1.176,3.734)rectangle++(0.04,0.04) (-0.853,4.818)rectangle++(0.04,0.04) (-0.534,5.725)rectangle++(0.04,0.04) (-0.227,6.484)rectangle++(0.04,0.04) (-1.752,2.963)rectangle++(0.04,0.04)
  (-1.373,4.185)rectangle++(0.04,0.04) (-1.002,5.208)rectangle++(0.04,0.04) (-1.576,2.504)rectangle++(0.04,0.04) (-1.241,3.788)rectangle++(0.04,0.04) (-0.906,4.866)rectangle++(0.04,0.04)
  (-0.578,5.768)rectangle++(0.04,0.04) (-1.636,2.725)rectangle++(0.04,0.04) (-1.284,3.978)rectangle++(0.04,0.04) (-1.312,2.169)rectangle++(0.04,0.04) (-1.031,3.492)rectangle++(0.04,0.04)
  (-0.739,4.606)rectangle++(0.04,0.04) (-0.446,5.541)rectangle++(0.04,0.04) (1.601,1.644)rectangle++(0.04,0.04) (1.424,2.929)rectangle++(0.04,0.04) (1.326,4.029)rectangle++(0.04,0.04)
  (-0.02,0.628)rectangle++(0.04,0.04) (0.006,2.131)rectangle++(0.04,0.04) (0.088,3.407)rectangle++(0.04,0.04) (0.209,4.489)rectangle++(0.04,0.04) (0.355,5.403)rectangle++(0.04,0.04)
  (0.516,6.175)rectangle++(0.04,0.04) (1.634,2.046)rectangle++(0.04,0.04) (1.469,3.27)rectangle++(0.04,0.04) (0.678,1.597)rectangle++(0.04,0.04) (0.638,2.926)rectangle++(0.04,0.04)
  (0.657,4.058)rectangle++(0.04,0.04) (0.719,5.019)rectangle++(0.04,0.04) (0.809,5.834)rectangle++(0.04,0.04) (1.495,2.041)rectangle++(0.04,0.04) (1.35,3.271)rectangle++(0.04,0.04)
  (1.276,4.322)rectangle++(0.04,0.04) (1.255,5.219)rectangle++(0.04,0.04) (1.274,5.982)rectangle++(0.04,0.04) (1.32,6.63)rectangle++(0.04,0.04) (-1.481,3.351)rectangle++(0.04,0.04)
  (-1.127,4.504)rectangle++(0.04,0.04) (-0.78,5.47)rectangle++(0.04,0.04) (-0.446,6.277)rectangle++(0.04,0.04) (-0.131,6.949)rectangle++(0.04,0.04) (0.165,7.508)rectangle++(0.04,0.04)
  (0.438,7.972)rectangle++(0.04,0.04) (0.689,8.355)rectangle++(0.04,0.04) (0.918,8.67)rectangle++(0.04,0.04) (1.125,8.929)rectangle++(0.04,0.04) (1.311,9.141)rectangle++(0.04,0.04)
  (1.478,9.313)rectangle++(0.04,0.04) (2.369,3.049)rectangle++(0.04,0.04) (2.133,4.093)rectangle++(0.04,0.04) (1.975,4.99)rectangle++(0.04,0.04) (-0.924,3.141)rectangle++(0.04,0.04)
  (-0.662,4.303)rectangle++(0.04,0.04) (-0.393,5.28)rectangle++(0.04,0.04) (-0.125,6.1)rectangle++(0.04,0.04) (0.136,6.786)rectangle++(0.04,0.04) (0.385,7.359)rectangle++(0.04,0.04)
  (-1.858,3.296)rectangle++(0.04,0.04) (-1.449,4.473)rectangle++(0.04,0.04) (-1.055,5.456)rectangle++(0.04,0.04) (-0.681,6.276)rectangle++(0.04,0.04) (-0.33,6.958)rectangle++(0.04,0.04)
  (-0.004,7.524)rectangle++(0.04,0.04) (0.295,7.991)rectangle++(0.04,0.04) (0.568,8.377)rectangle++(0.04,0.04) (0.816,8.694)rectangle++(0.04,0.04) (1.039,8.953)rectangle++(0.04,0.04)
  (1.239,9.165)rectangle++(0.04,0.04) (1.418,9.337)rectangle++(0.04,0.04) (1.576,9.476)rectangle++(0.04,0.04) (1.717,9.588)rectangle++(0.04,0.04) (2.41,3.177)rectangle++(0.04,0.04)
  (2.173,4.201)rectangle++(0.04,0.04) (2.013,5.08)rectangle++(0.04,0.04) (1.103,2.173)rectangle++(0.04,0.04) (1.022,3.399)rectangle++(0.04,0.04) (1.003,4.444)rectangle++(0.04,0.04)
  (-0.976,2.797)rectangle++(0.04,0.04) (-0.72,4.013)rectangle++(0.04,0.04) (-0.454,5.036)rectangle++(0.04,0.04) (-0.186,5.895)rectangle++(0.04,0.04) (0.075,6.614)rectangle++(0.04,0.04)
  (0.326,7.215)rectangle++(0.04,0.04) (0.564,7.716)rectangle++(0.04,0.04) (0.786,8.132)rectangle++(0.04,0.04) (0.991,8.477)rectangle++(0.04,0.04) (1.179,8.762)rectangle++(0.04,0.04)
  (-0.02,1.385)rectangle++(0.04,0.04) (0.036,2.774)rectangle++(0.04,0.04) (0.14,3.953)rectangle++(0.04,0.04) (0.275,4.951)rectangle++(0.04,0.04) (0.429,5.793)rectangle++(0.04,0.04)
  (0.594,6.503)rectangle++(0.04,0.04) (1.714,2.145)rectangle++(0.04,0.04) (1.541,3.351)rectangle++(0.04,0.04) (1.442,4.383)rectangle++(0.04,0.04) (1.398,5.264)rectangle++(0.04,0.04)
  (1.397,6.015)rectangle++(0.04,0.04) (1.426,6.653)rectangle++(0.04,0.04) (1.476,7.194)rectangle++(0.04,0.04) (1.54,7.652)rectangle++(0.04,0.04) (1.613,8.039)rectangle++(0.04,0.04)
  (1.69,8.365)rectangle++(0.04,0.04) (1.769,8.639)rectangle++(0.04,0.04) (1.847,8.868)rectangle++(0.04,0.04) (1.923,9.06)rectangle++(0.04,0.04) (1.994,9.221)rectangle++(0.04,0.04)
  (2.062,9.354)rectangle++(0.04,0.04) (2.125,9.465)rectangle++(0.04,0.04) (2.182,9.556)rectangle++(0.04,0.04) (2.331,3.291)rectangle++(0.04,0.04) (2.111,4.3)rectangle++(0.04,0.04)
  (1.964,5.167)rectangle++(0.04,0.04) (1.874,5.91)rectangle++(0.04,0.04) (1.827,6.544)rectangle++(0.04,0.04) (1.812,7.086)rectangle++(0.04,0.04) (1.822,7.547)rectangle++(0.04,0.04)
  (1.848,7.938)rectangle++(0.04,0.04) (1.886,8.27)rectangle++(0.04,0.04) (-1.794,3.842)rectangle++(0.04,0.04) (-1.373,4.934)rectangle++(0.04,0.04) (-0.972,5.845)rectangle++(0.04,0.04)
  (-0.595,6.603)rectangle++(0.04,0.04) (-0.244,7.233)rectangle++(0.04,0.04) (0.08,7.754)rectangle++(0.04,0.04) (0.376,8.184)rectangle++(0.04,0.04) (0.645,8.537)rectangle++(0.04,0.04)
  (2.276,2.647)rectangle++(0.04,0.04) (2.039,3.755)rectangle++(0.04,0.04) (1.881,4.706)rectangle++(0.04,0.04) (1.785,5.521)rectangle++(0.04,0.04) (1.736,6.218)rectangle++(0.04,0.04)
  (1.722,6.812)rectangle++(0.04,0.04) (1.734,7.317)rectangle++(0.04,0.04) (1.764,7.747)rectangle++(0.04,0.04) (1.807,8.11)rectangle++(0.04,0.04) (1.858,8.418)rectangle++(0.04,0.04)
  (1.914,8.677)rectangle++(0.04,0.04) (-1.894,3.938)rectangle++(0.04,0.04) (-1.424,2.02)rectangle++(0.04,0.04) (-1.132,3.37)rectangle++(0.04,0.04) (-0.829,4.506)rectangle++(0.04,0.04)
  (-0.527,5.459)rectangle++(0.04,0.04) (-0.232,6.258)rectangle++(0.04,0.04) (0.051,6.925)rectangle++(0.04,0.04) (-1.811,3.124)rectangle++(0.04,0.04) (-1.417,4.324)rectangle++(0.04,0.04)
  (-1.034,5.328)rectangle++(0.04,0.04) (-1.613,2.524)rectangle++(0.04,0.04) (-1.273,3.806)rectangle++(0.04,0.04) (-1.265,2.134)rectangle++(0.04,0.04) (-0.992,3.46)rectangle++(0.04,0.04)
  (1.1,1.002)rectangle++(0.04,0.04) (-0.062,2.063)rectangle++(0.04,0.04) (0.028,3.352)rectangle++(0.04,0.04) (0.155,4.444)rectangle++(0.04,0.04) (0.308,5.368)rectangle++(0.04,0.04)
  (0.474,6.146)rectangle++(0.04,0.04) (0.647,6.802)rectangle++(0.04,0.04) (0.819,7.352)rectangle++(0.04,0.04) (0.988,7.812)rectangle++(0.04,0.04) (1.15,8.197)rectangle++(0.04,0.04)
  (1.304,8.518)rectangle++(0.04,0.04) (1.447,8.784)rectangle++(0.04,0.04) (1.579,9.005)rectangle++(0.04,0.04) (1.7,9.187)rectangle++(0.04,0.04) (1.81,9.337)rectangle++(0.04,0.04)
  (-2.087,4.06)rectangle++(0.04,0.04) (1.432,0.862)rectangle++(0.04,0.04) (1.25,2.271)rectangle++(0.04,0.04) (1.151,3.477)rectangle++(0.04,0.04) (1.115,4.506)rectangle++(0.04,0.04)
  (1.126,5.381)rectangle++(0.04,0.04) (1.17,6.125)rectangle++(0.04,0.04) (1.237,6.756)rectangle++(0.04,0.04) (1.32,7.289)rectangle++(0.04,0.04) (1.411,7.739)rectangle++(0.04,0.04)
  (1.507,8.118)rectangle++(0.04,0.04) (1.603,8.436)rectangle++(0.04,0.04) (1.698,8.703)rectangle++(0.04,0.04) (1.789,8.926)rectangle++(0.04,0.04) (1.876,9.112)rectangle++(0.04,0.04)
  (1.957,9.266)rectangle++(0.04,0.04) (2.032,9.394)rectangle++(0.04,0.04) (2.1,9.5)rectangle++(0.04,0.04) (2.163,9.587)rectangle++(0.04,0.04) (2.22,9.659)rectangle++(0.04,0.04)
  (2.271,9.717)rectangle++(0.04,0.04) (2.317,9.765)rectangle++(0.04,0.04) (2.358,9.804)rectangle++(0.04,0.04) (2.374,3.396)rectangle++(0.04,0.04) (2.152,4.388)rectangle++(0.04,0.04)
  (2.002,5.24)rectangle++(0.04,0.04) (1.909,5.97)rectangle++(0.04,0.04) (1.368,2.359)rectangle++(0.04,0.04) (1.255,3.547)rectangle++(0.04,0.04) (0.787,1.607)rectangle++(0.04,0.04)
  (0.731,2.931)rectangle++(0.04,0.04) (-0.637,2.402)rectangle++(0.04,0.04) (-0.448,3.663)rectangle++(0.04,0.04) (-1.063,2.292)rectangle++(0.04,0.04) (-0.814,3.587)rectangle++(0.04,0.04)
  (-0.551,4.678)rectangle++(0.04,0.04) (-0.283,5.594)rectangle++(0.04,0.04) (-0.019,6.363)rectangle++(0.04,0.04) (0.236,7.005)rectangle++(0.04,0.04) (0.479,7.541)rectangle++(0.04,0.04)
  (0.706,7.987)rectangle++(0.04,0.04) (0.918,8.357)rectangle++(0.04,0.04) (1.112,8.663)rectangle++(0.04,0.04) (1.29,8.915)rectangle++(0.04,0.04) (1.451,9.123)rectangle++(0.04,0.04)
  (1.596,9.292)rectangle++(0.04,0.04) (-0.02,1.47)rectangle++(0.04,0.04) (0.04,2.846)rectangle++(0.04,0.04) (0.145,4.014)rectangle++(0.04,0.04) (0.282,5.002)rectangle++(0.04,0.04)
  (0.437,5.837)rectangle++(0.04,0.04) (0.603,6.54)rectangle++(0.04,0.04) (1.723,2.156)rectangle++(0.04,0.04) (1.549,3.36)rectangle++(0.04,0.04) (1.449,4.39)rectangle++(0.04,0.04)
  (1.405,5.27)rectangle++(0.04,0.04) (1.403,6.02)rectangle++(0.04,0.04) (1.431,6.657)rectangle++(0.04,0.04) (1.632,2.4)rectangle++(0.04,0.04) (1.481,3.571)rectangle++(0.04,0.04)
  (1.399,4.572)rectangle++(0.04,0.04) (1.37,5.426)rectangle++(0.04,0.04) (1.379,6.154)rectangle++(0.04,0.04) (1.417,6.772)rectangle++(0.04,0.04) (1.473,7.296)rectangle++(0.04,0.04)
  (1.541,7.739)rectangle++(0.04,0.04) (1.618,8.112)rectangle++(0.04,0.04) (1.697,8.427)rectangle++(0.04,0.04) (1.778,8.691)rectangle++(0.04,0.04) (1.856,8.913)rectangle++(0.04,0.04)
  (1.932,9.098)rectangle++(0.04,0.04) (2.004,9.252)rectangle++(0.04,0.04) (2.273,3.172)rectangle++(0.04,0.04) (2.056,4.201)rectangle++(0.04,0.04) (1.914,5.085)rectangle++(0.04,0.04)
  (1.119,2.148)rectangle++(0.04,0.04) (-0.356,2.319)rectangle++(0.04,0.04) (-0.212,3.582)rectangle++(0.04,0.04) (-0.039,4.649)rectangle++(0.04,0.04) (0.151,5.549)rectangle++(0.04,0.04)
  (0.348,6.307)rectangle++(0.04,0.04) (0.546,6.943)rectangle++(0.04,0.04) (0.739,7.476)rectangle++(0.04,0.04) (0.925,7.921)rectangle++(0.04,0.04) (1.101,8.292)rectangle++(0.04,0.04)
  (1.266,8.601)rectangle++(0.04,0.04) (1.418,8.856)rectangle++(0.04,0.04) (2.25,2.924)rectangle++(0.04,0.04) (2.027,3.992)rectangle++(0.04,0.04) (1.88,4.908)rectangle++(0.04,0.04)
  (1.792,5.693)rectangle++(0.04,0.04) (1.308,2.262)rectangle++(0.04,0.04) (1.2,3.467)rectangle++(0.04,0.04) (1.156,4.495)rectangle++(0.04,0.04) (1.16,5.371)rectangle++(0.04,0.04)
  (1.199,6.115)rectangle++(0.04,0.04) (-1.371,3.21)rectangle++(0.04,0.04) (-1.039,4.38)rectangle++(0.04,0.04) (-1.368,2.313)rectangle++(0.04,0.04) (-1.072,3.617)rectangle++(0.04,0.04)
  (-0.769,4.714)rectangle++(0.04,0.04) (-0.467,5.634)rectangle++(0.04,0.04) (-0.174,6.404)rectangle++(0.04,0.04) (1.802,1.922)rectangle++(0.04,0.04) (1.606,3.158)rectangle++(0.04,0.04)
  (1.489,4.216)rectangle++(0.04,0.04) (1.432,5.12)rectangle++(0.04,0.04) (1.42,5.891)rectangle++(0.04,0.04) (1.441,6.547)rectangle++(0.04,0.04) (1.484,7.103)rectangle++(0.04,0.04)
  (1.543,7.575)rectangle++(0.04,0.04) (1.872,2.649)rectangle++(0.04,0.04) (1.695,3.773)rectangle++(0.04,0.04) (1.589,4.736)rectangle++(0.04,0.04) (1.07,1.98)rectangle++(0.04,0.04)
  (0.987,3.236)rectangle++(0.04,0.04) (0.741,1.463)rectangle++(0.04,0.04) (0.686,2.81)rectangle++(0.04,0.04) (0.693,3.958)rectangle++(0.04,0.04) (0.745,4.932)rectangle++(0.04,0.04)
  (0.829,5.759)rectangle++(0.04,0.04) (0.933,6.458)rectangle++(0.04,0.04) (1.049,7.048)rectangle++(0.04,0.04) (1.171,7.545)rectangle++(0.04,0.04) (1.295,7.963)rectangle++(0.04,0.04)
  (1.417,8.313)rectangle++(0.04,0.04) (1.535,8.605)rectangle++(0.04,0.04) (1.647,8.849)rectangle++(0.04,0.04) (1.751,9.052)rectangle++(0.04,0.04) (1.849,9.221)rectangle++(0.04,0.04)
  (1.938,9.36)rectangle++(0.04,0.04) (2.019,9.474)rectangle++(0.04,0.04) (2.093,9.569)rectangle++(0.04,0.04) (2.348,3.271)rectangle++(0.04,0.04) (2.124,4.282)rectangle++(0.04,0.04)
  (1.975,5.151)rectangle++(0.04,0.04) (1.882,5.896)rectangle++(0.04,0.04) (1.351,2.334)rectangle++(0.04,0.04) (1.24,3.526)rectangle++(0.04,0.04) (1.193,4.544)rectangle++(0.04,0.04)
  (1.193,5.411)rectangle++(0.04,0.04) (-0.02,0.849)rectangle++(0.04,0.04) (0.015,2.319)rectangle++(0.04,0.04) (0.103,3.566)rectangle++(0.04,0.04) (0.228,4.624)rectangle++(0.04,0.04)
  (0.377,5.517)rectangle++(0.04,0.04) (0.539,6.271)rectangle++(0.04,0.04) (0.706,6.905)rectangle++(0.04,0.04) (0.874,7.437)rectangle++(0.04,0.04) (1.039,7.883)rectangle++(0.04,0.04)
  (1.196,8.255)rectangle++(0.04,0.04) (1.345,8.565)rectangle++(0.04,0.04) (2.179,2.835)rectangle++(0.04,0.04) (0.45,1.677)rectangle++(0.04,0.04) (0.447,3.004)rectangle++(0.04,0.04)
  (0.498,4.131)rectangle++(0.04,0.04) (0.586,5.088)rectangle++(0.04,0.04) (0.7,5.898)rectangle++(0.04,0.04) (0.828,6.581)rectangle++(0.04,0.04) (0.965,7.157)rectangle++(0.04,0.04)
  (1.105,7.641)rectangle++(0.04,0.04) (1.956,2.551)rectangle++(0.04,0.04) (1.763,3.686)rectangle++(0.04,0.04) (1.644,4.659)rectangle++(0.04,0.04) (1.581,5.491)rectangle++(0.04,0.04)
  (1.561,6.2)rectangle++(0.04,0.04) (1.573,6.804)rectangle++(0.04,0.04) (1.607,7.316)rectangle++(0.04,0.04) (1.656,7.751)rectangle++(0.04,0.04) (1.716,8.118)rectangle++(0.04,0.04)
  (1.781,8.428)rectangle++(0.04,0.04) (1.849,8.689)rectangle++(0.04,0.04) (1.917,8.908)rectangle++(0.04,0.04) (1.983,9.091)rectangle++(0.04,0.04) (2.047,9.244)rectangle++(0.04,0.04)
  (2.108,9.372)rectangle++(0.04,0.04) (-2.036,4.136)rectangle++(0.04,0.04) (-1.568,5.193)rectangle++(0.04,0.04) (-1.127,6.073)rectangle++(0.04,0.04) (1.852,1.594)rectangle++(0.04,0.04)
  (1.636,2.877)rectangle++(0.04,0.04) (0.543,1.546)rectangle++(0.04,0.04) (0.521,2.888)rectangle++(0.04,0.04) (0.556,4.031)rectangle++(0.04,0.04) (1.028,1.542)rectangle++(0.04,0.04)
  (0.933,2.866)rectangle++(0.04,0.04) (0.906,3.995)rectangle++(0.04,0.04) (0.927,4.956)rectangle++(0.04,0.04) (-1.124,2.892)rectangle++(0.04,0.04) (-0.842,4.1)rectangle++(0.04,0.04)
  (-0.554,5.115)rectangle++(0.04,0.04) (1.498,1.513)rectangle++(0.04,0.04) (1.331,2.823)rectangle++(0.04,0.04) (1.242,3.942)rectangle++(0.04,0.04) (1.212,4.897)rectangle++(0.04,0.04)
  (1.224,5.711)rectangle++(0.04,0.04) (1.266,6.401)rectangle++(0.04,0.04) (1.33,6.987)rectangle++(0.04,0.04) (1.408,7.482)rectangle++(0.04,0.04) (1.866,2.592)rectangle++(0.04,0.04)
  (0.428,1.537)rectangle++(0.04,0.04) (0.423,2.886)rectangle++(0.04,0.04) (0.473,4.032)rectangle++(0.04,0.04) (0.561,5.005)rectangle++(0.04,0.04) (0.675,5.828)rectangle++(0.04,0.04)
  (1.513,2.004)rectangle++(0.04,0.04) (1.364,3.239)rectangle++(0.04,0.04) (-0.02,0.501)rectangle++(0.04,0.04) (0.001,2.023)rectangle++(0.04,0.04) (0.079,3.316)rectangle++(0.04,0.04)
  (0.198,4.412)rectangle++(0.04,0.04) (0.343,5.338)rectangle++(0.04,0.04) (0.502,6.12)rectangle++(0.04,0.04) (0.67,6.778)rectangle++(0.04,0.04) (0.838,7.331)rectangle++(0.04,0.04)
  (1.003,7.794)rectangle++(0.04,0.04) (1.163,8.181)rectangle++(0.04,0.04) (1.313,8.503)rectangle++(0.04,0.04) (1.454,8.772)rectangle++(0.04,0.04) (1.585,8.994)rectangle++(0.04,0.04)
  (1.705,9.178)rectangle++(0.04,0.04) (1.814,9.329)rectangle++(0.04,0.04) (1.913,9.453)rectangle++(0.04,0.04) (-2.097,4.109)rectangle++(0.04,0.04) (-1.62,5.172)rectangle++(0.04,0.04)
  (-1.69,2.354)rectangle++(0.04,0.04) (-1.345,3.665)rectangle++(0.04,0.04) (-0.998,4.765)rectangle++(0.04,0.04) (-0.66,5.687)rectangle++(0.04,0.04) (-0.336,6.456)rectangle++(0.04,0.04)
  (-0.03,7.097)rectangle++(0.04,0.04) (0.257,7.63)rectangle++(0.04,0.04) (0.521,8.072)rectangle++(0.04,0.04) (0.764,8.436)rectangle++(0.04,0.04) (0.984,8.736)rectangle++(0.04,0.04)
  (1.184,8.983)rectangle++(0.04,0.04) (1.363,9.184)rectangle++(0.04,0.04) (1.524,9.348)rectangle++(0.04,0.04) (1.667,9.481)rectangle++(0.04,0.04) (1.794,9.589)rectangle++(0.04,0.04)
  (1.906,9.675)rectangle++(0.04,0.04) (2.005,9.743)rectangle++(0.04,0.04) (2.092,9.798)rectangle++(0.04,0.04) (2.168,9.841)rectangle++(0.04,0.04) (2.234,9.874)rectangle++(0.04,0.04)
  (2.292,9.9)rectangle++(0.04,0.04) (2.342,9.919)rectangle++(0.04,0.04) (2.385,9.934)rectangle++(0.04,0.04) (2.423,9.945)rectangle++(0.04,0.04) (2.455,9.952)rectangle++(0.04,0.04)
  (2.482,9.957)rectangle++(0.04,0.04) (2.506,9.961)rectangle++(0.04,0.04) (2.526,9.963)rectangle++(0.04,0.04) (2.544,9.963)rectangle++(0.04,0.04) (2.559,9.963)rectangle++(0.04,0.04)
  (2.571,9.963)rectangle++(0.04,0.04) (2.582,9.962)rectangle++(0.04,0.04) (2.591,9.96)rectangle++(0.04,0.04) (2.598,9.959)rectangle++(0.04,0.04) (2.605,9.957)rectangle++(0.04,0.04)
  (2.61,9.956)rectangle++(0.04,0.04) (2.615,9.954)rectangle++(0.04,0.04) (2.618,9.953)rectangle++(0.04,0.04) (2.622,9.951)rectangle++(0.04,0.04) (2.624,9.95)rectangle++(0.04,0.04)
  (2.626,9.949)rectangle++(0.04,0.04) (2.628,9.947)rectangle++(0.04,0.04) (2.63,9.946)rectangle++(0.04,0.04) (-2.081,4.382)rectangle++(0.04,0.04) (-1.596,5.404)rectangle++(0.04,0.04)
  (-1.143,6.254)rectangle++(0.04,0.04) (-0.723,6.957)rectangle++(0.04,0.04) (-0.339,7.539)rectangle++(0.04,0.04) (0.011,8.018)rectangle++(0.04,0.04) (0.328,8.411)rectangle++(0.04,0.04)
  (-2.142,3.515)rectangle++(0.04,0.04) (-1.683,4.67)rectangle++(0.04,0.04) (-1.572,2.229)rectangle++(0.04,0.04) (-1.249,3.554)rectangle++(0.04,0.04) (-0.922,4.667)rectangle++(0.04,0.04)
  (-0.599,5.6)rectangle++(0.04,0.04) (-0.287,6.38)rectangle++(0.04,0.04) (0.009,7.031)rectangle++(0.04,0.04) (0.286,7.572)rectangle++(0.04,0.04) (0.544,8.021)rectangle++(0.04,0.04)
  (0.781,8.392)rectangle++(0.04,0.04) (0.998,8.698)rectangle++(0.04,0.04) (1.194,8.95)rectangle++(0.04,0.04) (1.37,9.156)rectangle++(0.04,0.04) (1.529,9.324)rectangle++(0.04,0.04)
  (1.67,9.46)rectangle++(0.04,0.04) (2.381,3.135)rectangle++(0.04,0.04) (0.503,1.801)rectangle++(0.04,0.04) (0.498,3.107)rectangle++(0.04,0.04) (0.545,4.217)rectangle++(0.04,0.04)
  (0.63,5.159)rectangle++(0.04,0.04) (1.333,1.832)rectangle++(0.04,0.04) (1.204,3.1)rectangle++(0.04,0.04) (1.145,4.183)rectangle++(0.04,0.04) (1.138,5.106)rectangle++(0.04,0.04)
  (1.17,5.891)rectangle++(0.04,0.04) (1.228,6.557)rectangle++(0.04,0.04) (1.304,7.12)rectangle++(0.04,0.04) (1.391,7.596)rectangle++(0.04,0.04) (1.484,7.997)rectangle++(0.04,0.04)
  (-0.02,1.263)rectangle++(0.04,0.04) (0.031,2.67)rectangle++(0.04,0.04) (-0.709,2.184)rectangle++(0.04,0.04) (-0.518,3.481)rectangle++(0.04,0.04) (-0.303,4.576)rectangle++(0.04,0.04)
  (-0.077,5.497)rectangle++(0.04,0.04) (0.152,6.272)rectangle++(0.04,0.04) (0.378,6.921)rectangle++(0.04,0.04) (0.596,7.464)rectangle++(0.04,0.04) (0.803,7.917)rectangle++(0.04,0.04)
  (0.997,8.294)rectangle++(0.04,0.04) (1.177,8.606)rectangle++(0.04,0.04) (1.343,8.864)rectangle++(0.04,0.04) (-2.057,3.848)rectangle++(0.04,0.04) (-1.597,4.949)rectangle++(0.04,0.04)
  (-1.162,5.867)rectangle++(0.04,0.04) (1.8,1.536)rectangle++(0.04,0.04) (0.143,1.267)rectangle++(0.04,0.04) (0.17,2.667)rectangle++(0.04,0.04) (-0.681,2.215)rectangle++(0.04,0.04)
  (-0.492,3.506)rectangle++(0.04,0.04) (-0.28,4.596)rectangle++(0.04,0.04) (-0.057,5.514)rectangle++(0.04,0.04) (0.17,6.285)rectangle++(0.04,0.04) (0.394,6.932)rectangle++(0.04,0.04)
  (0.61,7.473)rectangle++(0.04,0.04) (0.815,7.924)rectangle++(0.04,0.04) (2.079,2.544)rectangle++(0.04,0.04) (1.867,3.675)rectangle++(0.04,0.04) (1.731,4.645)rectangle++(0.04,0.04)
  (1.655,5.475)rectangle++(0.04,0.04) (1.624,6.184)rectangle++(0.04,0.04) (1.625,6.788)rectangle++(0.04,0.04) (1.651,7.301)rectangle++(0.04,0.04) (1.693,7.736)rectangle++(0.04,0.04)
  (1.746,8.104)rectangle++(0.04,0.04) (1.806,8.415)rectangle++(0.04,0.04) (1.87,8.676)rectangle++(0.04,0.04) (-1.903,3.928)rectangle++(0.04,0.04) (-1.463,5.011)rectangle++(0.04,0.04)
  (-1.045,5.914)rectangle++(0.04,0.04) (-0.654,6.665)rectangle++(0.04,0.04) (-0.291,7.288)rectangle++(0.04,0.04) (0.042,7.802)rectangle++(0.04,0.04) (0.345,8.226)rectangle++(0.04,0.04)
  (0.62,8.575)rectangle++(0.04,0.04) (0.868,8.86)rectangle++(0.04,0.04) (-2.151,3.738)rectangle++(0.04,0.04) (-1.681,4.859)rectangle++(0.04,0.04) (-1.237,5.794)rectangle++(0.04,0.04)
  (-0.822,6.571)rectangle++(0.04,0.04) (-0.438,7.214)rectangle++(0.04,0.04) (-0.086,7.746)rectangle++(0.04,0.04) (0.235,8.183)rectangle++(0.04,0.04) (0.525,8.543)rectangle++(0.04,0.04)
  (0.785,8.837)rectangle++(0.04,0.04) (1.019,9.076)rectangle++(0.04,0.04) (-2.177,3.82)rectangle++(0.04,0.04) (-1.7,4.93)rectangle++(0.04,0.04) (-1.25,5.855)rectangle++(0.04,0.04)
  (-0.83,6.623)rectangle++(0.04,0.04) (-0.443,7.259)rectangle++(0.04,0.04) (-0.089,7.784)rectangle++(0.04,0.04) (0.234,8.216)rectangle++(0.04,0.04) (0.525,8.571)rectangle++(0.04,0.04)
  (0.787,8.86)rectangle++(0.04,0.04) (1.021,9.096)rectangle++(0.04,0.04) (1.23,9.287)rectangle++(0.04,0.04) (1.415,9.441)rectangle++(0.04,0.04) (1.578,9.564)rectangle++(0.04,0.04)
  (1.721,9.663)rectangle++(0.04,0.04) (1.848,9.741)rectangle++(0.04,0.04) (1.958,9.802)rectangle++(0.04,0.04) (2.433,3.291)rectangle++(0.04,0.04) (2.198,4.297)rectangle++(0.04,0.04)
  (2.038,5.16)rectangle++(0.04,0.04) (1.936,5.901)rectangle++(0.04,0.04) (1.344,2.35)rectangle++(0.04,0.04) (1.235,3.54)rectangle++(0.04,0.04) (1.189,4.556)rectangle++(0.04,0.04)
  (1.19,5.421)rectangle++(0.04,0.04) (1.227,6.156)rectangle++(0.04,0.04) (1.287,6.78)rectangle++(0.04,0.04) (1.688,2.392)rectangle++(0.04,0.04) (1.528,3.562)rectangle++(0.04,0.04)
  (1.439,4.562)rectangle++(0.04,0.04) (1.404,5.417)rectangle++(0.04,0.04) (1.408,6.144)rectangle++(0.04,0.04) (1.44,6.763)rectangle++(0.04,0.04) (1.66,2.427)rectangle++(0.04,0.04)
  (0.413,1.444)rectangle++(0.04,0.04) (0.407,2.807)rectangle++(0.04,0.04) (0.456,3.966)rectangle++(0.04,0.04) (1.025,1.5)rectangle++(0.04,0.04) (0.929,2.831)rectangle++(0.04,0.04)
  (0.901,3.965)rectangle++(0.04,0.04) (0.958,1.616)rectangle++(0.04,0.04) (0.876,2.931)rectangle++(0.04,0.04) (0.86,4.053)rectangle++(0.04,0.04) (0.891,5.007)rectangle++(0.04,0.04)
  (0.955,5.816)rectangle++(0.04,0.04) (1.043,6.502)rectangle++(0.04,0.04) (1.144,7.081)rectangle++(0.04,0.04) (1.253,7.569)rectangle++(0.04,0.04) (1.366,7.98)rectangle++(0.04,0.04)
  (1.478,8.325)rectangle++(0.04,0.04) (1.587,8.613)rectangle++(0.04,0.04) (1.691,8.854)rectangle++(0.04,0.04) (1.79,9.054)rectangle++(0.04,0.04) (1.881,9.221)rectangle++(0.04,0.04)
  (-2.042,4.05)rectangle++(0.04,0.04) (1.423,0.871)rectangle++(0.04,0.04) (1.242,2.28)rectangle++(0.04,0.04) (1.145,3.484)rectangle++(0.04,0.04) (1.11,4.512)rectangle++(0.04,0.04)
  (1.122,5.387)rectangle++(0.04,0.04) (1.167,6.13)rectangle++(0.04,0.04) (1.524,2.205)rectangle++(0.04,0.04) (1.381,3.409)rectangle++(0.04,0.04) (1.308,4.439)rectangle++(0.04,0.04)
  (1.287,5.317)rectangle++(0.04,0.04) (1.305,6.064)rectangle++(0.04,0.04) (1.349,6.698)rectangle++(0.04,0.04) (1.413,7.236)rectangle++(0.04,0.04) (-1.62,3.506)rectangle++(0.04,0.04)
  (-1.239,4.641)rectangle++(0.04,0.04) (1.468,1.222)rectangle++(0.04,0.04) (1.294,2.576)rectangle++(0.04,0.04) (1.201,3.734)rectangle++(0.04,0.04) (1.168,4.722)rectangle++(0.04,0.04)
  (1.18,5.563)rectangle++(0.04,0.04) (1.223,6.278)rectangle++(0.04,0.04) (1.288,6.883)rectangle++(0.04,0.04) (1.368,7.396)rectangle++(0.04,0.04) (1.457,7.828)rectangle++(0.04,0.04)
  (1.956,2.687)rectangle++(0.04,0.04) (1.768,3.802)rectangle++(0.04,0.04) (1.652,4.757)rectangle++(0.04,0.04) (1.593,5.574)rectangle++(0.04,0.04) (1.575,6.27)rectangle++(0.04,0.04)
  (-1.337,3.331)rectangle++(0.04,0.04) (-1.005,4.481)rectangle++(0.04,0.04) (-0.677,5.445)rectangle++(0.04,0.04) (-0.36,6.252)rectangle++(0.04,0.04) (-0.058,6.924)rectangle++(0.04,0.04)
  (0.225,7.484)rectangle++(0.04,0.04) (0.489,7.949)rectangle++(0.04,0.04) (0.731,8.333)rectangle++(0.04,0.04) (0.953,8.65)rectangle++(0.04,0.04) (1.154,8.911)rectangle++(0.04,0.04)
  (1.335,9.124)rectangle++(0.04,0.04) (1.497,9.298)rectangle++(0.04,0.04) (1.642,9.44)rectangle++(0.04,0.04) (1.771,9.554)rectangle++(0.04,0.04) (1.886,9.647)rectangle++(0.04,0.04)
  (1.987,9.72)rectangle++(0.04,0.04) (2.075,9.779)rectangle++(0.04,0.04) (2.153,9.825)rectangle++(0.04,0.04) (2.221,9.862)rectangle++(0.04,0.04) (2.28,9.89)rectangle++(0.04,0.04)
  (2.331,9.911)rectangle++(0.04,0.04) (2.376,9.928)rectangle++(0.04,0.04) (2.414,9.94)rectangle++(0.04,0.04) (2.448,9.948)rectangle++(0.04,0.04) (2.476,9.954)rectangle++(0.04,0.04)
  (2.501,9.958)rectangle++(0.04,0.04) (2.522,9.961)rectangle++(0.04,0.04) (-2.107,4.36)rectangle++(0.04,0.04) (-1.618,5.387)rectangle++(0.04,0.04) (-1.162,6.24)rectangle++(0.04,0.04)
  (-0.741,6.946)rectangle++(0.04,0.04) (-0.354,7.53)rectangle++(0.04,0.04) (-0.002,8.011)rectangle++(0.04,0.04) (0.317,8.406)rectangle++(0.04,0.04) (0.603,8.728)rectangle++(0.04,0.04)
  (0.86,8.991)rectangle++(0.04,0.04) (-2.187,3.765)rectangle++(0.04,0.04) (-1.711,4.884)rectangle++(0.04,0.04) (1.607,1.157)rectangle++(0.04,0.04) (-0.001,2.213)rectangle++(0.04,0.04)
  (0.086,3.477)rectangle++(0.04,0.04) (0.21,4.549)rectangle++(0.04,0.04) (0.358,5.454)rectangle++(0.04,0.04) (0.52,6.218)rectangle++(0.04,0.04) (-1.534,3.077)rectangle++(0.04,0.04)
  (-1.183,4.273)rectangle++(0.04,0.04) (-0.837,5.275)rectangle++(0.04,0.04) (-0.502,6.114)rectangle++(0.04,0.04) (-0.185,6.813)rectangle++(0.04,0.04) (-0.02,1.073)rectangle++(0.04,0.04)
  (0.024,2.509)rectangle++(0.04,0.04) (0.118,3.728)rectangle++(0.04,0.04) (0.248,4.76)rectangle++(0.04,0.04) (0.399,5.633)rectangle++(0.04,0.04) (1.5,1.885)rectangle++(0.04,0.04)
  (1.348,3.139)rectangle++(0.04,0.04) (1.269,4.21)rectangle++(0.04,0.04) (-0.862,2.807)rectangle++(0.04,0.04) (-0.623,4.017)rectangle++(0.04,0.04) (-0.371,5.035)rectangle++(0.04,0.04)
  (-0.116,5.891)rectangle++(0.04,0.04) (0.135,6.608)rectangle++(0.04,0.04) (1.813,2.051)rectangle++(0.04,0.04) (1.621,3.267)rectangle++(0.04,0.04) (1.506,4.308)rectangle++(0.04,0.04)
  (1.45,5.198)rectangle++(0.04,0.04) (1.438,5.957)rectangle++(0.04,0.04) (1.459,6.602)rectangle++(0.04,0.04) (1.502,7.149)rectangle++(0.04,0.04) (1.56,7.613)rectangle++(0.04,0.04)
  (1.629,8.005)rectangle++(0.04,0.04) (1.702,8.335)rectangle++(0.04,0.04) (2.061,2.873)rectangle++(0.04,0.04) (1.865,3.956)rectangle++(0.04,0.04) (1.741,4.884)rectangle++(0.04,0.04)
  (1.673,5.678)rectangle++(0.04,0.04) (1.647,6.356)rectangle++(0.04,0.04) (1.652,6.933)rectangle++(0.04,0.04) (1.679,7.423)rectangle++(0.04,0.04) (1.722,7.838)rectangle++(0.04,0.04)
  (1.775,8.19)rectangle++(0.04,0.04) (1.834,8.487)rectangle++(0.04,0.04) (2.084,2.944)rectangle++(0.04,0.04) (1.887,4.015)rectangle++(0.04,0.04) (1.762,4.933)rectangle++(0.04,0.04)
  (1.693,5.719)rectangle++(0.04,0.04) (1.666,6.39)rectangle++(0.04,0.04) (1.669,6.961)rectangle++(0.04,0.04) (1.695,7.446)rectangle++(0.04,0.04) (1.736,7.858)rectangle++(0.04,0.04)
  (1.788,8.206)rectangle++(0.04,0.04) (1.846,8.5)rectangle++(0.04,0.04) (1.907,8.747)rectangle++(0.04,0.04) (2.146,3.025)rectangle++(0.04,0.04) (1.943,4.082)rectangle++(0.04,0.04)
  (1.812,4.988)rectangle++(0.04,0.04) (1.738,5.763)rectangle++(0.04,0.04) (1.705,6.426)rectangle++(0.04,0.04) (1.704,6.99)rectangle++(0.04,0.04) (1.726,7.469)rectangle++(0.04,0.04)
  (1.764,7.876)rectangle++(0.04,0.04) (1.812,8.22)rectangle++(0.04,0.04) (2.012,2.874)rectangle++(0.04,0.04) (-0.366,2.684)rectangle++(0.04,0.04) (-0.206,3.892)rectangle++(0.04,0.04)
  (-0.022,4.913)rectangle++(0.04,0.04) (0.176,5.773)rectangle++(0.04,0.04) (0.378,6.496)rectangle++(0.04,0.04) (1.745,2.087)rectangle++(0.04,0.04) (-0.215,2.45)rectangle++(0.04,0.04)
  (-0.087,3.687)rectangle++(0.04,0.04) (0.071,4.734)rectangle++(0.04,0.04) (0.248,5.617)rectangle++(0.04,0.04) (0.433,6.361)rectangle++(0.04,0.04) (0.62,6.985)rectangle++(0.04,0.04)
  (0.805,7.509)rectangle++(0.04,0.04) (0.982,7.947)rectangle++(0.04,0.04) (1.15,8.312)rectangle++(0.04,0.04) (1.308,8.615)rectangle++(0.04,0.04) (-1.999,3.785)rectangle++(0.04,0.04)
  (-1.55,4.894)rectangle++(0.04,0.04) (-1.604,2.309)rectangle++(0.04,0.04) (-1.273,3.623)rectangle++(0.04,0.04) (1.188,0.969)rectangle++(0.04,0.04) (1.046,2.372)rectangle++(0.04,0.04)
  (0.982,3.571)rectangle++(0.04,0.04) (0.975,4.592)rectangle++(0.04,0.04) (-1.02,2.824)rectangle++(0.04,0.04) (-0.756,4.037)rectangle++(0.04,0.04) (-0.484,5.058)rectangle++(0.04,0.04)
  (-1.433,2.591)rectangle++(0.04,0.04) (-1.117,3.855)rectangle++(0.04,0.04) (-0.797,4.918)rectangle++(0.04,0.04) (-0.483,5.808)rectangle++(0.04,0.04) (-0.18,6.553)rectangle++(0.04,0.04)
  (0.107,7.173)rectangle++(0.04,0.04) (0.375,7.689)rectangle++(0.04,0.04) (0.624,8.117)rectangle++(0.04,0.04) (-2.007,3.518)rectangle++(0.04,0.04) (-1.567,4.667)rectangle++(0.04,0.04)
  (-1.148,5.626)rectangle++(0.04,0.04) (-0.753,6.424)rectangle++(0.04,0.04) (-0.385,7.087)rectangle++(0.04,0.04) (2.025,2.031)rectangle++(0.04,0.04) (1.8,3.241)rectangle++(0.04,0.04)
  (1.657,4.279)rectangle++(0.04,0.04) (1.578,5.167)rectangle++(0.04,0.04) (1.546,5.925)rectangle++(0.04,0.04) (1.549,6.571)rectangle++(0.04,0.04) (1.577,7.12)rectangle++(0.04,0.04)
  (1.623,7.585)rectangle++(0.04,0.04) (1.681,7.978)rectangle++(0.04,0.04) (1.746,8.311)rectangle++(0.04,0.04) (1.814,8.59)rectangle++(0.04,0.04) (2.116,2.963)rectangle++(0.04,0.04)
  (1.915,4.03)rectangle++(0.04,0.04) (-0.686,2.916)rectangle++(0.04,0.04) (-0.469,4.102)rectangle++(0.04,0.04) (-0.237,5.102)rectangle++(0.04,0.04) (1.447,1.593)rectangle++(0.04,0.04)
  (1.291,2.892)rectangle++(0.04,0.04) (1.211,4.003)rectangle++(0.04,0.04) (1.187,4.95)rectangle++(0.04,0.04) (1.205,5.757)rectangle++(0.04,0.04) (1.252,6.441)rectangle++(0.04,0.04)
  (1.32,7.021)rectangle++(0.04,0.04) (1.401,7.511)rectangle++(0.04,0.04) (1.489,7.925)rectangle++(0.04,0.04) (1.58,8.273)rectangle++(0.04,0.04) (1.672,8.565)rectangle++(0.04,0.04)
  (1.761,8.809)rectangle++(0.04,0.04) (1.847,9.014)rectangle++(0.04,0.04) (1.929,9.184)rectangle++(0.04,0.04) (2.005,9.325)rectangle++(0.04,0.04) (2.075,9.443)rectangle++(0.04,0.04)
  (2.139,9.539)rectangle++(0.04,0.04) (2.198,9.619)rectangle++(0.04,0.04) (2.346,3.31)rectangle++(0.04,0.04) (2.125,4.316)rectangle++(0.04,0.04) (1.976,5.18)rectangle++(0.04,0.04)
  (1.885,5.92)rectangle++(0.04,0.04) (1.837,6.553)rectangle++(0.04,0.04) (1.821,7.093)rectangle++(0.04,0.04) (1.829,7.552)rectangle++(0.04,0.04) (1.823,2.718)rectangle++(0.04,0.04)
  (1.656,3.834)rectangle++(0.04,0.04) (1.559,4.789)rectangle++(0.04,0.04) (1.514,5.604)rectangle++(0.04,0.04) (1.509,6.299)rectangle++(0.04,0.04) (1.532,6.89)rectangle++(0.04,0.04)
  (1.576,7.392)rectangle++(0.04,0.04) (1.633,7.816)rectangle++(0.04,0.04) (1.699,8.174)rectangle++(0.04,0.04) (1.769,8.477)rectangle++(0.04,0.04) (-1.871,3.861)rectangle++(0.04,0.04)
  (-1.438,4.953)rectangle++(0.04,0.04) (-1.027,5.863)rectangle++(0.04,0.04) (-0.64,6.621)rectangle++(0.04,0.04) (-0.282,7.25)rectangle++(0.04,0.04) (0.048,7.77)rectangle++(0.04,0.04)
  (0.35,8.199)rectangle++(0.04,0.04) (0.623,8.551)rectangle++(0.04,0.04) (0.869,8.84)rectangle++(0.04,0.04) (1.09,9.075)rectangle++(0.04,0.04) (1.288,9.266)rectangle++(0.04,0.04)
  (1.463,9.421)rectangle++(0.04,0.04) (1.618,9.546)rectangle++(0.04,0.04) (1.755,9.645)rectangle++(0.04,0.04) (1.875,9.725)rectangle++(0.04,0.04) (1.981,9.787)rectangle++(0.04,0.04)
  (2.073,9.836)rectangle++(0.04,0.04) (2.153,9.874)rectangle++(0.04,0.04) (2.223,9.903)rectangle++(0.04,0.04) (2.284,9.925)rectangle++(0.04,0.04) (2.336,9.941)rectangle++(0.04,0.04)
  (2.381,9.952)rectangle++(0.04,0.04) (2.42,9.961)rectangle++(0.04,0.04) (2.453,9.966)rectangle++(0.04,0.04) (-2.122,4.346)rectangle++(0.04,0.04) (-1.632,5.375)rectangle++(0.04,0.04)
  (-1.174,6.23)rectangle++(0.04,0.04) (-0.751,6.939)rectangle++(0.04,0.04) (2.038,1.9)rectangle++(0.04,0.04) (1.806,3.13)rectangle++(0.04,0.04) (1.658,4.184)rectangle++(0.04,0.04)
  (1.575,5.086)rectangle++(0.04,0.04) (1.54,5.857)rectangle++(0.04,0.04) (1.541,6.513)rectangle++(0.04,0.04) (1.568,7.07)rectangle++(0.04,0.04) (1.613,7.543)rectangle++(0.04,0.04)
  (1.671,7.944)rectangle++(0.04,0.04) (1.736,8.281)rectangle++(0.04,0.04) (1.805,8.566)rectangle++(0.04,0.04) (1.874,8.805)rectangle++(0.04,0.04) (1.943,9.006)rectangle++(0.04,0.04)
  (2.01,9.173)rectangle++(0.04,0.04) (2.073,9.313)rectangle++(0.04,0.04) (-2.028,4.115)rectangle++(0.04,0.04) (-1.561,5.175)rectangle++(0.04,0.04) (-1.122,6.057)rectangle++(0.04,0.04)
  (-0.714,6.79)rectangle++(0.04,0.04) (-0.337,7.396)rectangle++(0.04,0.04) (0.007,7.896)rectangle++(0.04,0.04) (0.319,8.308)rectangle++(0.04,0.04) (0.602,8.645)rectangle++(0.04,0.04)
  (0.855,8.92)rectangle++(0.04,0.04) (1.081,9.144)rectangle++(0.04,0.04) (1.283,9.326)rectangle++(0.04,0.04) (1.461,9.472)rectangle++(0.04,0.04) (1.619,9.589)rectangle++(0.04,0.04)
  (1.757,9.682)rectangle++(0.04,0.04) (1.879,9.755)rectangle++(0.04,0.04) (1.985,9.813)rectangle++(0.04,0.04) (2.077,9.858)rectangle++(0.04,0.04) (2.431,3.336)rectangle++(0.04,0.04)
  (2.198,4.335)rectangle++(0.04,0.04) (0.867,2.042)rectangle++(0.04,0.04) (0.816,3.297)rectangle++(0.04,0.04) (-0.715,2.502)rectangle++(0.04,0.04) (-0.51,3.752)rectangle++(0.04,0.04)
  (-0.286,4.805)rectangle++(0.04,0.04) (-1.328,2.58)rectangle++(0.04,0.04) (-1.028,3.843)rectangle++(0.04,0.04) (-0.722,4.904)rectangle++(0.04,0.04) (-0.42,5.793)rectangle++(0.04,0.04)
  (-0.127,6.537)rectangle++(0.04,0.04) (0.151,7.158)rectangle++(0.04,0.04) (0.413,7.674)rectangle++(0.04,0.04) (0.655,8.103)rectangle++(0.04,0.04) (0.879,8.457)rectangle++(0.04,0.04)
  (1.083,8.75)rectangle++(0.04,0.04) (1.269,8.99)rectangle++(0.04,0.04) (1.436,9.187)rectangle++(0.04,0.04) (-0.02,1.453)rectangle++(0.04,0.04) (0.039,2.832)rectangle++(0.04,0.04)
  (0.144,4.002)rectangle++(0.04,0.04) (0.28,4.992)rectangle++(0.04,0.04) (0.436,5.828)rectangle++(0.04,0.04) (0.601,6.533)rectangle++(0.04,0.04) (0.77,7.125)rectangle++(0.04,0.04)
  (0.938,7.622)rectangle++(0.04,0.04) (1.1,8.037)rectangle++(0.04,0.04) (-1.891,3.61)rectangle++(0.04,0.04) (-1.465,4.74)rectangle++(0.04,0.04) (-1.058,5.684)rectangle++(0.04,0.04)
  (-0.674,6.47)rectangle++(0.04,0.04) (-0.316,7.123)rectangle++(0.04,0.04) (0.014,7.663)rectangle++(0.04,0.04) (0.316,8.109)rectangle++(0.04,0.04) (0.591,8.476)rectangle++(0.04,0.04)
  (0.839,8.778)rectangle++(0.04,0.04) (1.062,9.024)rectangle++(0.04,0.04) (1.262,9.224)rectangle++(0.04,0.04) (1.439,9.386)rectangle++(0.04,0.04) (1.596,9.517)rectangle++(0.04,0.04)
  (1.735,9.622)rectangle++(0.04,0.04) (1.858,9.705)rectangle++(0.04,0.04) (2.421,3.242)rectangle++(0.04,0.04) (2.186,4.255)rectangle++(0.04,0.04) (2.026,5.126)rectangle++(0.04,0.04)
  (1.925,5.872)rectangle++(0.04,0.04) (1.869,6.51)rectangle++(0.04,0.04) (1.847,7.055)rectangle++(0.04,0.04) (1.85,7.519)rectangle++(0.04,0.04) (1.871,7.914)rectangle++(0.04,0.04)
  (1.905,8.248)rectangle++(0.04,0.04) (1.947,8.531)rectangle++(0.04,0.04) (-1.85,3.913)rectangle++(0.04,0.04) (1.356,0.888)rectangle++(0.04,0.04) (1.186,2.297)rectangle++(0.04,0.04)
  (1.098,3.501)rectangle++(0.04,0.04) (1.071,4.528)rectangle++(0.04,0.04) (1.089,5.402)rectangle++(0.04,0.04) (1.14,6.145)rectangle++(0.04,0.04) (1.212,6.774)rectangle++(0.04,0.04)
  (1.299,7.305)rectangle++(0.04,0.04) (1.394,7.754)rectangle++(0.04,0.04) (1.493,8.131)rectangle++(0.04,0.04) (1.592,8.448)rectangle++(0.04,0.04) (1.689,8.713)rectangle++(0.04,0.04)
  (1.782,8.935)rectangle++(0.04,0.04) (2.217,3.038)rectangle++(0.04,0.04) (2.004,4.09)rectangle++(0.04,0.04) (1.865,4.992)rectangle++(0.04,0.04) (1.782,5.765)rectangle++(0.04,0.04)
  (1.743,6.425)rectangle++(0.04,0.04) (1.737,6.988)rectangle++(0.04,0.04) (1.754,7.466)rectangle++(0.04,0.04) (1.787,7.872)rectangle++(0.04,0.04) (-1.711,3.732)rectangle++(0.04,0.04)
  (-1.307,4.837)rectangle++(0.04,0.04) (-0.92,5.76)rectangle++(0.04,0.04) (-0.554,6.529)rectangle++(0.04,0.04) (-0.212,7.168)rectangle++(0.04,0.04) (-1.927,3.117)rectangle++(0.04,0.04)
  (-1.516,4.323)rectangle++(0.04,0.04) (-1.118,5.331)rectangle++(0.04,0.04) (-0.739,6.173)rectangle++(0.04,0.04) (-0.383,6.872)rectangle++(0.04,0.04) (-0.053,7.453)rectangle++(0.04,0.04)
  (0.251,7.933)rectangle++(0.04,0.04) (0.528,8.33)rectangle++(0.04,0.04) (0.78,8.655)rectangle++(0.04,0.04) (1.007,8.922)rectangle++(0.04,0.04) (1.211,9.14)rectangle++(0.04,0.04)
  (2.36,2.938)rectangle++(0.04,0.04) (2.121,3.999)rectangle++(0.04,0.04) (-0.637,2.957)rectangle++(0.04,0.04) (-0.425,4.135)rectangle++(0.04,0.04) (-0.198,5.128)rectangle++(0.04,0.04)
  (0.034,5.963)rectangle++(0.04,0.04) (0.266,6.663)rectangle++(0.04,0.04) (0.49,7.249)rectangle++(0.04,0.04) (0.704,7.739)rectangle++(0.04,0.04) (0.906,8.146)rectangle++(0.04,0.04)
  (1.094,8.484)rectangle++(0.04,0.04) (1.267,8.764)rectangle++(0.04,0.04) (1.425,8.995)rectangle++(0.04,0.04) (1.569,9.185)rectangle++(0.04,0.04) (1.699,9.34)rectangle++(0.04,0.04)
  (1.815,9.468)rectangle++(0.04,0.04) (-2.12,4.089)rectangle++(0.04,0.04) (-1.64,5.157)rectangle++(0.04,0.04) (-1.69,2.346)rectangle++(0.04,0.04) (-1.345,3.658)rectangle++(0.04,0.04)
  (-0.999,4.759)rectangle++(0.04,0.04) (-0.661,5.682)rectangle++(0.04,0.04) (-0.337,6.452)rectangle++(0.04,0.04) (-0.03,7.094)rectangle++(0.04,0.04) (0.256,7.627)rectangle++(0.04,0.04)
  (0.52,8.069)rectangle++(0.04,0.04) (0.763,8.434)rectangle++(0.04,0.04) (0.984,8.735)rectangle++(0.04,0.04) (1.183,8.981)rectangle++(0.04,0.04) (1.363,9.183)rectangle++(0.04,0.04)
  (1.523,9.347)rectangle++(0.04,0.04) (1.667,9.48)rectangle++(0.04,0.04) (1.794,9.588)rectangle++(0.04,0.04) (2.398,3.197)rectangle++(0.04,0.04) (2.164,4.218)rectangle++(0.04,0.04)
  (2.006,5.095)rectangle++(0.04,0.04) (1.907,5.847)rectangle++(0.04,0.04) (1.852,6.49)rectangle++(0.04,0.04) (1.832,7.038)rectangle++(0.04,0.04) (1.836,7.505)rectangle++(0.04,0.04)
  (1.859,7.902)rectangle++(0.04,0.04) (1.916,2.81)rectangle++(0.04,0.04) (1.739,3.908)rectangle++(0.04,0.04) (1.632,4.848)rectangle++(0.04,0.04) (1.579,5.652)rectangle++(0.04,0.04)
  (1.566,6.337)rectangle++(0.04,0.04) (1.583,6.92)rectangle++(0.04,0.04) (1.62,7.415)rectangle++(0.04,0.04) (1.671,7.834)rectangle++(0.04,0.04) (1.732,8.188)rectangle++(0.04,0.04)
  (2.016,2.845)rectangle++(0.04,0.04) (1.825,3.934)rectangle++(0.04,0.04) (1.706,4.867)rectangle++(0.04,0.04) (1.643,5.665)rectangle++(0.04,0.04) (1.621,6.346)rectangle++(0.04,0.04)
  (1.629,6.925)rectangle++(0.04,0.04) (1.66,7.418)rectangle++(0.04,0.04) (1.811,2.642)rectangle++(0.04,0.04) (1.642,3.769)rectangle++(0.04,0.04) (1.545,4.734)rectangle++(0.04,0.04)
  (1.5,5.559)rectangle++(0.04,0.04) (1.495,6.261)rectangle++(0.04,0.04) (1.519,6.858)rectangle++(0.04,0.04) (1.563,7.365)rectangle++(0.04,0.04) (1.621,7.794)rectangle++(0.04,0.04)
  (1.688,8.156)rectangle++(0.04,0.04) (1.759,8.461)rectangle++(0.04,0.04) (1.831,8.718)rectangle++(0.04,0.04) (1.903,8.933)rectangle++(0.04,0.04) (1.973,9.113)rectangle++(0.04,0.04)
  (2.039,9.264)rectangle++(0.04,0.04) (2.102,9.389)rectangle++(0.04,0.04) (2.16,9.493)rectangle++(0.04,0.04) (2.213,9.579)rectangle++(0.04,0.04) (-2.069,4.205)rectangle++(0.04,0.04)
  (-1.593,5.253)rectangle++(0.04,0.04) (-1.146,6.125)rectangle++(0.04,0.04) (-1.843,2.673)rectangle++(0.04,0.04) (-1.462,3.942)rectangle++(0.04,0.04) (-1.087,5.005)rectangle++(0.04,0.04)
  (-0.726,5.894)rectangle++(0.04,0.04) (-0.384,6.635)rectangle++(0.04,0.04) (-0.063,7.252)rectangle++(0.04,0.04) (-1.919,3.17)rectangle++(0.04,0.04) (-1.507,4.367)rectangle++(0.04,0.04)
  (-1.108,5.369)rectangle++(0.04,0.04) (-0.729,6.204)rectangle++(0.04,0.04) (-1.78,2.786)rectangle++(0.04,0.04) (-1.404,4.036)rectangle++(0.04,0.04) (-1.034,5.083)rectangle++(0.04,0.04)
  (-1.549,2.469)rectangle++(0.04,0.04) (0.906,0.62)rectangle++(0.04,0.04) (0.793,2.087)rectangle++(0.04,0.04) (0.755,3.338)rectangle++(0.04,0.04) (0.773,4.403)rectangle++(0.04,0.04)
  (0.831,5.308)rectangle++(0.04,0.04) (0.917,6.075)rectangle++(0.04,0.04) (1.02,6.723)rectangle++(0.04,0.04) (1.134,7.27)rectangle++(0.04,0.04) (1.252,7.73)rectangle++(0.04,0.04)
  (1.371,8.117)rectangle++(0.04,0.04) (1.488,8.441)rectangle++(0.04,0.04) (-1.918,3.788)rectangle++(0.04,0.04) (-1.481,4.893)rectangle++(0.04,0.04) (-1.065,5.814)rectangle++(0.04,0.04)
  (-0.675,6.581)rectangle++(0.04,0.04) (-0.313,7.217)rectangle++(0.04,0.04) (0.021,7.743)rectangle++(0.04,0.04) (0.325,8.177)rectangle++(0.04,0.04) (0.601,8.534)rectangle++(0.04,0.04)
  (-2.12,3.605)rectangle++(0.04,0.04) (-1.66,4.745)rectangle++(0.04,0.04) (-1.223,5.696)rectangle++(0.04,0.04) (-0.814,6.487)rectangle++(0.04,0.04) (-0.435,7.142)rectangle++(0.04,0.04)
  (-0.086,7.685)rectangle++(0.04,0.04) (0.232,8.132)rectangle++(0.04,0.04) (0.52,8.499)rectangle++(0.04,0.04) (0.78,8.799)rectangle++(0.04,0.04) (2.329,2.745)rectangle++(0.04,0.04)
  (2.088,3.836)rectangle++(0.04,0.04) (1.926,4.773)rectangle++(0.04,0.04) (1.826,5.576)rectangle++(0.04,0.04) (1.27,2.243)rectangle++(0.04,0.04) (1.167,3.452)rectangle++(0.04,0.04)
  (1.128,4.484)rectangle++(0.04,0.04) (1.136,5.362)rectangle++(0.04,0.04) (1.178,6.109)rectangle++(0.04,0.04) (1.243,6.741)rectangle++(0.04,0.04) (-1.525,3.358)rectangle++(0.04,0.04)
  (-1.164,4.512)rectangle++(0.04,0.04) (-0.811,5.478)rectangle++(0.04,0.04) (-0.473,6.285)rectangle++(0.04,0.04) (-0.153,6.957)rectangle++(0.04,0.04) (0.146,7.516)rectangle++(0.04,0.04)
  (0.423,7.979)rectangle++(0.04,0.04) (0.676,8.361)rectangle++(0.04,0.04) (0.907,8.676)rectangle++(0.04,0.04) (1.116,8.935)rectangle++(0.04,0.04) (1.304,9.146)rectangle++(0.04,0.04)
  (1.472,9.318)rectangle++(0.04,0.04) (1.622,9.458)rectangle++(0.04,0.04) (1.754,9.57)rectangle++(0.04,0.04) (1.872,9.661)rectangle++(0.04,0.04) (2.407,3.235)rectangle++(0.04,0.04)
  (2.173,4.25)rectangle++(0.04,0.04) (0.847,2.015)rectangle++(0.04,0.04) (0.798,3.275)rectangle++(0.04,0.04) (0.807,4.348)rectangle++(0.04,0.04) (0.858,5.26)rectangle++(0.04,0.04)
  (0.937,6.033)rectangle++(0.04,0.04) (1.036,6.687)rectangle++(0.04,0.04) (1.146,7.238)rectangle++(0.04,0.04) (1.261,7.703)rectangle++(0.04,0.04) (-1.772,3.574)rectangle++(0.04,0.04)
  (-1.365,4.705)rectangle++(0.04,0.04) (-0.974,5.65)rectangle++(0.04,0.04) (-0.604,6.438)rectangle++(0.04,0.04) (-1.816,2.866)rectangle++(0.04,0.04) (-1.431,4.105)rectangle++(0.04,0.04)
  (-1.054,5.143)rectangle++(0.04,0.04) (-0.693,6.01)rectangle++(0.04,0.04) (-0.351,6.732)rectangle++(0.04,0.04) (-0.031,7.333)rectangle++(0.04,0.04) (0.265,7.83)rectangle++(0.04,0.04)
  (0.536,8.241)rectangle++(0.04,0.04) (0.783,8.58)rectangle++(0.04,0.04) (1.007,8.858)rectangle++(0.04,0.04) (1.208,9.085)rectangle++(0.04,0.04) (1.388,9.27)rectangle++(0.04,0.04)
  (1.548,9.42)rectangle++(0.04,0.04) (1.691,9.542)rectangle++(0.04,0.04) (1.816,9.639)rectangle++(0.04,0.04) (1.927,9.717)rectangle++(0.04,0.04) (2.025,9.778)rectangle++(0.04,0.04)
  (2.11,9.827)rectangle++(0.04,0.04) (2.184,9.864)rectangle++(0.04,0.04) (2.249,9.894)rectangle++(0.04,0.04) (2.305,9.916)rectangle++(0.04,0.04) (2.354,9.932)rectangle++(0.04,0.04)
  (2.396,9.945)rectangle++(0.04,0.04) (2.432,9.953)rectangle++(0.04,0.04) (2.463,9.959)rectangle++(0.04,0.04) (2.49,9.963)rectangle++(0.04,0.04) (2.399,3.468)rectangle++(0.04,0.04)
  (2.175,4.448)rectangle++(0.04,0.04) (2.025,5.29)rectangle++(0.04,0.04) (1.931,6.012)rectangle++(0.04,0.04) (1.879,6.629)rectangle++(0.04,0.04) (1.86,7.156)rectangle++(0.04,0.04);
\fill[nuage]
  (1.865,7.604)rectangle++(0.04,0.04) (1.888,7.985)rectangle++(0.04,0.04) (1.922,8.308)rectangle++(0.04,0.04) (1.963,8.581)rectangle++(0.04,0.04) (2.01,8.812)rectangle++(0.04,0.04)
  (2.059,9.006)rectangle++(0.04,0.04) (2.195,3.127)rectangle++(0.04,0.04) (0.529,1.751)rectangle++(0.04,0.04) (0.517,3.064)rectangle++(0.04,0.04) (0.56,4.18)rectangle++(0.04,0.04)
  (0.641,5.126)rectangle++(0.04,0.04) (0.748,5.928)rectangle++(0.04,0.04) (0.87,6.605)rectangle++(0.04,0.04) (-1.564,3.242)rectangle++(0.04,0.04) (-1.202,4.415)rectangle++(0.04,0.04)
  (-0.848,5.397)rectangle++(0.04,0.04) (-1.594,2.581)rectangle++(0.04,0.04) (-1.254,3.854)rectangle++(0.04,0.04) (-0.914,4.922)rectangle++(0.04,0.04) (-0.582,5.817)rectangle++(0.04,0.04)
  (-0.264,6.564)rectangle++(0.04,0.04) (0.036,7.186)rectangle++(0.04,0.04) (0.316,7.703)rectangle++(0.04,0.04) (0.574,8.131)rectangle++(0.04,0.04) (0.811,8.485)rectangle++(0.04,0.04)
  (1.027,8.776)rectangle++(0.04,0.04) (2.286,2.803)rectangle++(0.04,0.04) (2.053,3.887)rectangle++(0.04,0.04) (1.898,4.818)rectangle++(0.04,0.04) (-0.894,3.086)rectangle++(0.04,0.04)
  (-0.639,4.255)rectangle++(0.04,0.04) (-0.375,5.238)rectangle++(0.04,0.04) (-0.112,6.064)rectangle++(0.04,0.04) (0.146,6.755)rectangle++(0.04,0.04) (0.392,7.332)rectangle++(0.04,0.04)
  (-1.849,3.292)rectangle++(0.04,0.04) (-1.247,1.888)rectangle++(0.04,0.04) (-0.987,3.251)rectangle++(0.04,0.04) (-0.711,4.399)rectangle++(0.04,0.04) (-0.43,5.364)rectangle++(0.04,0.04)
  (-0.154,6.173)rectangle++(0.04,0.04) (1.734,1.871)rectangle++(0.04,0.04) (1.547,3.118)rectangle++(0.04,0.04) (1.437,4.184)rectangle++(0.04,0.04) (1.387,5.095)rectangle++(0.04,0.04)
  (1.38,5.872)rectangle++(0.04,0.04) (1.406,6.532)rectangle++(0.04,0.04) (1.454,7.092)rectangle++(0.04,0.04) (1.517,7.566)rectangle++(0.04,0.04) (1.59,7.967)rectangle++(0.04,0.04)
  (1.668,8.304)rectangle++(0.04,0.04) (2.058,2.857)rectangle++(0.04,0.04) (1.861,3.942)rectangle++(0.04,0.04) (1.737,4.873)rectangle++(0.04,0.04) (1.086,2.051)rectangle++(0.04,0.04)
  (1.003,3.296)rectangle++(0.04,0.04) (0.982,4.358)rectangle++(0.04,0.04) (1.007,5.261)rectangle++(0.04,0.04) (1.064,6.028)rectangle++(0.04,0.04) (1.144,6.677)rectangle++(0.04,0.04)
  (1.237,7.226)rectangle++(0.04,0.04) (1.338,7.689)rectangle++(0.04,0.04) (1.443,8.078)rectangle++(0.04,0.04) (1.547,8.405)rectangle++(0.04,0.04) (1.649,8.679)rectangle++(0.04,0.04)
  (1.747,8.907)rectangle++(0.04,0.04) (1.839,9.097)rectangle++(0.04,0.04) (1.925,9.255)rectangle++(0.04,0.04) (2.004,9.386)rectangle++(0.04,0.04) (2.077,9.494)rectangle++(0.04,0.04)
  (2.143,9.583)rectangle++(0.04,0.04) (2.202,9.656)rectangle++(0.04,0.04) (2.256,9.716)rectangle++(0.04,0.04) (-2.096,4.245)rectangle++(0.04,0.04) (-1.614,5.289)rectangle++(0.04,0.04)
  (-1.163,6.156)rectangle++(0.04,0.04) (-0.744,6.875)rectangle++(0.04,0.04) (-0.36,7.47)rectangle++(0.04,0.04) (-2.035,3.15)rectangle++(0.04,0.04) (1.17,0.657)rectangle++(0.04,0.04)
  (1.018,2.108)rectangle++(0.04,0.04) (-0.366,2.287)rectangle++(0.04,0.04) (-0.221,3.555)rectangle++(0.04,0.04) (-0.048,4.627)rectangle++(0.04,0.04) (0.142,5.531)rectangle++(0.04,0.04)
  (0.34,6.292)rectangle++(0.04,0.04) (0.538,6.93)rectangle++(0.04,0.04) (0.732,7.466)rectangle++(0.04,0.04) (0.919,7.913)rectangle++(0.04,0.04) (1.095,8.285)rectangle++(0.04,0.04)
  (1.26,8.595)rectangle++(0.04,0.04) (1.413,8.851)rectangle++(0.04,0.04) (2.249,2.922)rectangle++(0.04,0.04) (2.026,3.99)rectangle++(0.04,0.04) (1.88,4.906)rectangle++(0.04,0.04)
  (1.792,5.691)rectangle++(0.04,0.04) (1.749,6.362)rectangle++(0.04,0.04) (1.739,6.934)rectangle++(0.04,0.04) (1.753,7.421)rectangle++(0.04,0.04) (1.785,7.834)rectangle++(0.04,0.04)
  (1.828,8.183)rectangle++(0.04,0.04) (1.879,8.479)rectangle++(0.04,0.04) (1.934,8.728)rectangle++(0.04,0.04) (1.991,8.938)rectangle++(0.04,0.04) (2.048,9.114)rectangle++(0.04,0.04)
  (2.103,9.261)rectangle++(0.04,0.04) (2.156,9.384)rectangle++(0.04,0.04) (2.205,9.486)rectangle++(0.04,0.04) (2.252,9.571)rectangle++(0.04,0.04) (2.295,9.642)rectangle++(0.04,0.04)
  (2.334,9.7)rectangle++(0.04,0.04) (2.37,9.748)rectangle++(0.04,0.04) (2.402,9.787)rectangle++(0.04,0.04) (2.431,9.819)rectangle++(0.04,0.04) (2.367,3.419)rectangle++(0.04,0.04)
  (2.147,4.407)rectangle++(0.04,0.04) (1.999,5.257)rectangle++(0.04,0.04) (1.907,5.984)rectangle++(0.04,0.04) (1.858,6.607)rectangle++(0.04,0.04) (1.842,7.137)rectangle++(0.04,0.04)
  (1.705,2.622)rectangle++(0.04,0.04) (1.552,3.757)rectangle++(0.04,0.04) (1.467,4.727)rectangle++(0.04,0.04) (1.434,5.556)rectangle++(0.04,0.04) (1.439,6.261)rectangle++(0.04,0.04)
  (1.52,2.307)rectangle++(0.04,0.04) (0.401,1.379)rectangle++(0.04,0.04) (0.393,2.752)rectangle++(0.04,0.04) (0.442,3.92)rectangle++(0.04,0.04) (0.531,4.91)rectangle++(0.04,0.04)
  (0.645,5.749)rectangle++(0.04,0.04) (0.776,6.457)rectangle++(0.04,0.04) (-1.545,3.188)rectangle++(0.04,0.04) (1.107,0.793)rectangle++(0.04,0.04) (0.97,2.226)rectangle++(0.04,0.04)
  (0.912,3.45)rectangle++(0.04,0.04) (0.911,4.492)rectangle++(0.04,0.04) (0.952,5.378)rectangle++(0.04,0.04) (1.022,6.129)rectangle++(0.04,0.04) (1.112,6.765)rectangle++(0.04,0.04)
  (1.71,2.343)rectangle++(0.04,0.04) (1.545,3.519)rectangle++(0.04,0.04) (1.452,4.526)rectangle++(0.04,0.04) (1.413,5.385)rectangle++(0.04,0.04) (1.414,6.117)rectangle++(0.04,0.04)
  (1.445,6.739)rectangle++(0.04,0.04) (1.653,2.423)rectangle++(0.04,0.04) (1.5,3.59)rectangle++(0.04,0.04) (1.416,4.587)rectangle++(0.04,0.04) (1.385,5.439)rectangle++(0.04,0.04)
  (1.393,6.164)rectangle++(0.04,0.04) (1.428,6.78)rectangle++(0.04,0.04) (1.483,7.302)rectangle++(0.04,0.04) (1.55,7.743)rectangle++(0.04,0.04) (1.625,8.116)rectangle++(0.04,0.04)
  (1.704,8.43)rectangle++(0.04,0.04) (2.087,2.896)rectangle++(0.04,0.04) (0.48,1.668)rectangle++(0.04,0.04) (0.473,2.995)rectangle++(0.04,0.04) (0.52,4.123)rectangle++(0.04,0.04)
  (0.604,5.08)rectangle++(0.04,0.04) (0.715,5.89)rectangle++(0.04,0.04) (-1.41,3.049)rectangle++(0.04,0.04) (-1.078,4.244)rectangle++(0.04,0.04) (-1.34,2.275)rectangle++(0.04,0.04)
  (-1.051,3.583)rectangle++(0.04,0.04) (-0.752,4.684)rectangle++(0.04,0.04) (-1.389,2.447)rectangle++(0.04,0.04) (-1.085,3.731)rectangle++(0.04,0.04) (-0.775,4.811)rectangle++(0.04,0.04)
  (-0.469,5.717)rectangle++(0.04,0.04) (-0.172,6.474)rectangle++(0.04,0.04) (0.11,7.106)rectangle++(0.04,0.04) (0.376,7.632)rectangle++(0.04,0.04) (0.623,8.068)rectangle++(0.04,0.04)
  (0.85,8.429)rectangle++(0.04,0.04) (1.057,8.727)rectangle++(0.04,0.04) (-2.079,3.752)rectangle++(0.04,0.04) (-1.619,4.869)rectangle++(0.04,0.04) (-1.184,5.799)rectangle++(0.04,0.04)
  (-0.776,6.573)rectangle++(0.04,0.04) (-1.885,2.856)rectangle++(0.04,0.04) (-1.491,4.1)rectangle++(0.04,0.04) (-1.105,5.141)rectangle++(0.04,0.04) (-0.736,6.01)rectangle++(0.04,0.04)
  (-0.387,6.734)rectangle++(0.04,0.04) (-0.062,7.336)rectangle++(0.04,0.04) (0.238,7.834)rectangle++(0.04,0.04) (0.514,8.246)rectangle++(0.04,0.04) (0.764,8.584)rectangle++(0.04,0.04)
  (0.991,8.862)rectangle++(0.04,0.04) (1.195,9.09)rectangle++(0.04,0.04) (1.377,9.275)rectangle++(0.04,0.04) (1.539,9.424)rectangle++(0.04,0.04) (-2.164,4.016)rectangle++(0.04,0.04)
  (1.432,0.831)rectangle++(0.04,0.04) (1.248,2.246)rectangle++(0.04,0.04) (1.148,3.455)rectangle++(0.04,0.04) (0.778,1.558)rectangle++(0.04,0.04) (-0.271,2.111)rectangle++(0.04,0.04)
  (-0.148,3.401)rectangle++(0.04,0.04) (0.008,4.493)rectangle++(0.04,0.04) (0.184,5.415)rectangle++(0.04,0.04) (0.371,6.192)rectangle++(0.04,0.04) (0.561,6.844)rectangle++(0.04,0.04)
  (0.748,7.391)rectangle++(0.04,0.04) (0.93,7.849)rectangle++(0.04,0.04) (1.102,8.231)rectangle++(0.04,0.04) (1.264,8.548)rectangle++(0.04,0.04) (1.414,8.812)rectangle++(0.04,0.04)
  (1.552,9.029)rectangle++(0.04,0.04) (1.678,9.209)rectangle++(0.04,0.04) (1.793,9.357)rectangle++(0.04,0.04) (1.896,9.478)rectangle++(0.04,0.04) (1.988,9.577)rectangle++(0.04,0.04)
  (2.071,9.657)rectangle++(0.04,0.04) (2.144,9.722)rectangle++(0.04,0.04) (-2.12,4.221)rectangle++(0.04,0.04) (-1.635,5.269)rectangle++(0.04,0.04) (-1.181,6.14)rectangle++(0.04,0.04)
  (-0.761,6.863)rectangle++(0.04,0.04) (-0.374,7.46)rectangle++(0.04,0.04) (-0.022,7.952)rectangle++(0.04,0.04) (0.297,8.356)rectangle++(0.04,0.04) (0.585,8.687)rectangle++(0.04,0.04)
  (0.842,8.957)rectangle++(0.04,0.04) (1.072,9.176)rectangle++(0.04,0.04) (1.276,9.353)rectangle++(0.04,0.04) (1.457,9.495)rectangle++(0.04,0.04) (1.616,9.609)rectangle++(0.04,0.04)
  (1.755,9.699)rectangle++(0.04,0.04) (1.878,9.77)rectangle++(0.04,0.04) (1.985,9.826)rectangle++(0.04,0.04) (2.078,9.869)rectangle++(0.04,0.04) (2.159,9.901)rectangle++(0.04,0.04)
  (2.229,9.926)rectangle++(0.04,0.04) (2.289,9.944)rectangle++(0.04,0.04) (2.342,9.957)rectangle++(0.04,0.04) (2.419,3.429)rectangle++(0.04,0.04) (2.191,4.414)rectangle++(0.04,0.04)
  (2.037,5.26)rectangle++(0.04,0.04) (1.94,5.986)rectangle++(0.04,0.04) (-1.19,3.352)rectangle++(0.04,0.04) (-0.879,4.493)rectangle++(0.04,0.04) (-0.57,5.45)rectangle++(0.04,0.04)
  (-0.269,6.252)rectangle++(0.04,0.04) (0.02,6.921)rectangle++(0.04,0.04) (0.291,7.478)rectangle++(0.04,0.04) (0.545,7.941)rectangle++(0.04,0.04) (-1.977,3.461)rectangle++(0.04,0.04)
  (-1.544,4.617)rectangle++(0.04,0.04) (-1.13,5.583)rectangle++(0.04,0.04) (-0.739,6.387)rectangle++(0.04,0.04) (-0.375,7.054)rectangle++(0.04,0.04) (-0.039,7.608)rectangle++(0.04,0.04)
  (0.269,8.064)rectangle++(0.04,0.04) (0.549,8.44)rectangle++(0.04,0.04) (-2.106,3.572)rectangle++(0.04,0.04) (-1.371,1.891)rectangle++(0.04,0.04) (-1.092,3.258)rectangle++(0.04,0.04)
  (-0.8,4.409)rectangle++(0.04,0.04) (1.337,1.28)rectangle++(0.04,0.04) (1.186,2.631)rectangle++(0.04,0.04) (1.111,3.785)rectangle++(0.04,0.04) (1.093,4.769)rectangle++(0.04,0.04)
  (1.154,1.859)rectangle++(0.04,0.04) (1.053,3.13)rectangle++(0.04,0.04) (1.018,4.215)rectangle++(0.04,0.04) (1.032,5.138)rectangle++(0.04,0.04) (1.08,5.922)rectangle++(0.04,0.04)
  (1.153,6.587)rectangle++(0.04,0.04) (1.654,2.311)rectangle++(0.04,0.04) (-0.291,2.478)rectangle++(0.04,0.04) (-0.151,3.714)rectangle++(0.04,0.04) (0.018,4.759)rectangle++(0.04,0.04)
  (0.204,5.641)rectangle++(0.04,0.04) (0.397,6.383)rectangle++(0.04,0.04) (0.59,7.006)rectangle++(0.04,0.04) (0.78,7.527)rectangle++(0.04,0.04) (0.962,7.963)rectangle++(0.04,0.04)
  (1.134,8.326)rectangle++(0.04,0.04) (2.144,2.723)rectangle++(0.04,0.04) (1.929,3.825)rectangle++(0.04,0.04) (1.791,4.77)rectangle++(0.04,0.04) (1.711,5.579)rectangle++(0.04,0.04)
  (1.675,6.27)rectangle++(0.04,0.04) (1.672,6.859)rectangle++(0.04,0.04) (1.694,7.359)rectangle++(0.04,0.04) (1.732,7.784)rectangle++(0.04,0.04) (1.781,8.143)rectangle++(0.04,0.04)
  (1.996,2.847)rectangle++(0.04,0.04) (1.808,3.937)rectangle++(0.04,0.04) (1.692,4.87)rectangle++(0.04,0.04) (1.631,5.668)rectangle++(0.04,0.04) (1.611,6.349)rectangle++(0.04,0.04)
  (1.621,6.928)rectangle++(0.04,0.04) (1.653,7.42)rectangle++(0.04,0.04) (1.699,7.837)rectangle++(0.04,0.04) (1.756,8.19)rectangle++(0.04,0.04) (1.818,8.488)rectangle++(0.04,0.04)
  (1.882,8.738)rectangle++(0.04,0.04) (1.947,8.948)rectangle++(0.04,0.04) (2.011,9.124)rectangle++(0.04,0.04) (2.072,9.271)rectangle++(0.04,0.04) (2.13,9.394)rectangle++(0.04,0.04)
  (2.184,9.496)rectangle++(0.04,0.04) (2.234,9.58)rectangle++(0.04,0.04) (2.28,9.65)rectangle++(0.04,0.04) (2.322,9.708)rectangle++(0.04,0.04) (2.36,9.755)rectangle++(0.04,0.04)
  (2.394,9.793)rectangle++(0.04,0.04) (2.424,9.825)rectangle++(0.04,0.04) (2.451,9.85)rectangle++(0.04,0.04) (2.475,9.871)rectangle++(0.04,0.04) (2.375,3.443)rectangle++(0.04,0.04)
  (2.154,4.427)rectangle++(0.04,0.04) (2.006,5.273)rectangle++(0.04,0.04) (1.914,5.998)rectangle++(0.04,0.04) (1.865,6.618)rectangle++(0.04,0.04) (1.847,7.147)rectangle++(0.04,0.04)
  (1.854,7.597)rectangle++(0.04,0.04) (1.878,7.98)rectangle++(0.04,0.04) (1.913,8.304)rectangle++(0.04,0.04) (1.956,8.578)rectangle++(0.04,0.04) (2.003,8.809)rectangle++(0.04,0.04)
  (2.053,9.004)rectangle++(0.04,0.04) (2.103,9.167)rectangle++(0.04,0.04) (2.152,9.304)rectangle++(0.04,0.04) (2.199,9.419)rectangle++(0.04,0.04) (2.244,9.514)rectangle++(0.04,0.04)
  (2.286,9.594)rectangle++(0.04,0.04) (2.324,9.659)rectangle++(0.04,0.04) (-0.02,1.529)rectangle++(0.04,0.04) (0.042,2.896)rectangle++(0.04,0.04) (0.149,4.056)rectangle++(0.04,0.04)
  (0.287,5.038)rectangle++(0.04,0.04) (0.443,5.867)rectangle++(0.04,0.04) (0.609,6.566)rectangle++(0.04,0.04) (1.73,2.164)rectangle++(0.04,0.04) (1.555,3.367)rectangle++(0.04,0.04)
  (1.454,4.396)rectangle++(0.04,0.04) (1.409,5.274)rectangle++(0.04,0.04) (1.248,2.062)rectangle++(0.04,0.04) (1.141,3.299)rectangle++(0.04,0.04) (1.1,4.355)rectangle++(0.04,0.04)
  (1.107,5.254)rectangle++(0.04,0.04) (1.149,6.018)rectangle++(0.04,0.04) (1.215,6.665)rectangle++(0.04,0.04) (1.297,7.213)rectangle++(0.04,0.04) (1.389,7.675)rectangle++(0.04,0.04)
  (1.485,8.065)rectangle++(0.04,0.04) (1.583,8.392)rectangle++(0.04,0.04) (1.679,8.666)rectangle++(0.04,0.04) (1.772,8.895)rectangle++(0.04,0.04) (-1.98,3.953)rectangle++(0.04,0.04)
  (-1.527,5.036)rectangle++(0.04,0.04) (-1.098,5.938)rectangle++(0.04,0.04) (-0.698,6.687)rectangle++(0.04,0.04) (-0.328,7.308)rectangle++(0.04,0.04) (0.011,7.821)rectangle++(0.04,0.04)
  (0.32,8.244)rectangle++(0.04,0.04) (0.6,8.591)rectangle++(0.04,0.04) (0.851,8.874)rectangle++(0.04,0.04) (1.076,9.105)rectangle++(0.04,0.04) (1.277,9.293)rectangle++(0.04,0.04)
  (1.455,9.444)rectangle++(0.04,0.04) (1.612,9.565)rectangle++(0.04,0.04) (1.751,9.662)rectangle++(0.04,0.04) (1.872,9.739)rectangle++(0.04,0.04) (1.979,9.8)rectangle++(0.04,0.04)
  (2.072,9.847)rectangle++(0.04,0.04) (-2.167,4.232)rectangle++(0.04,0.04) (-1.675,5.28)rectangle++(0.04,0.04) (-1.215,6.151)rectangle++(0.04,0.04) (-0.789,6.873)rectangle++(0.04,0.04)
  (-0.398,7.47)rectangle++(0.04,0.04) (-0.041,7.962)rectangle++(0.04,0.04) (0.281,8.365)rectangle++(0.04,0.04) (-2.14,3.494)rectangle++(0.04,0.04) (-1.681,4.652)rectangle++(0.04,0.04)
  (-1.245,5.617)rectangle++(0.04,0.04) (1.742,1.454)rectangle++(0.04,0.04) (1.537,2.763)rectangle++(0.04,0.04) (1.415,3.883)rectangle++(0.04,0.04) (1.356,4.84)rectangle++(0.04,0.04)
  (1.135,1.944)rectangle++(0.04,0.04) (1.04,3.203)rectangle++(0.04,0.04) (1.01,4.277)rectangle++(0.04,0.04) (1.027,5.192)rectangle++(0.04,0.04) (1.079,5.968)rectangle++(0.04,0.04)
  (1.153,6.626)rectangle++(0.04,0.04) (1.665,2.32)rectangle++(0.04,0.04) (1.506,3.502)rectangle++(0.04,0.04) (1.418,4.512)rectangle++(0.04,0.04) (1.383,5.375)rectangle++(0.04,0.04)
  (1.389,6.11)rectangle++(0.04,0.04) (1.423,6.734)rectangle++(0.04,0.04) (1.476,7.263)rectangle++(0.04,0.04) (1.543,7.711)rectangle++(0.04,0.04) (1.618,8.089)rectangle++(0.04,0.04)
  (1.697,8.407)rectangle++(0.04,0.04) (1.776,8.674)rectangle++(0.04,0.04) (-1.921,3.906)rectangle++(0.04,0.04) (-1.479,4.993)rectangle++(0.04,0.04) (-1.06,5.899)rectangle++(0.04,0.04)
  (-0.667,6.653)rectangle++(0.04,0.04) (-0.303,7.278)rectangle++(0.04,0.04) (0.031,7.795)rectangle++(0.04,0.04) (0.336,8.22)rectangle++(0.04,0.04) (2.234,2.49)rectangle++(0.04,0.04)
  (0.345,1.608)rectangle++(0.04,0.04) (0.355,2.95)rectangle++(0.04,0.04) (0.418,4.089)rectangle++(0.04,0.04) (0.516,5.055)rectangle++(0.04,0.04) (0.639,5.873)rectangle++(0.04,0.04)
  (-1.42,3.028)rectangle++(0.04,0.04) (-1.088,4.227)rectangle++(0.04,0.04) (-0.758,5.232)rectangle++(0.04,0.04) (-0.437,6.074)rectangle++(0.04,0.04) (-0.131,6.777)rectangle++(0.04,0.04)
  (0.158,7.362)rectangle++(0.04,0.04) (0.426,7.847)rectangle++(0.04,0.04) (0.674,8.249)rectangle++(0.04,0.04) (2.191,2.585)rectangle++(0.04,0.04) (1.964,3.706)rectangle++(0.04,0.04)
  (-0.592,2.856)rectangle++(0.04,0.04) (-0.391,4.047)rectangle++(0.04,0.04) (-0.173,5.052)rectangle++(0.04,0.04) (0.053,5.897)rectangle++(0.04,0.04) (0.279,6.607)rectangle++(0.04,0.04)
  (0.499,7.201)rectangle++(0.04,0.04) (0.71,7.697)rectangle++(0.04,0.04) (0.909,8.11)rectangle++(0.04,0.04) (1.095,8.454)rectangle++(0.04,0.04) (1.267,8.738)rectangle++(0.04,0.04)
  (1.424,8.973)rectangle++(0.04,0.04) (1.567,9.166)rectangle++(0.04,0.04) (1.697,9.325)rectangle++(0.04,0.04) (2.339,3.109)rectangle++(0.04,0.04) (2.11,4.145)rectangle++(0.04,0.04)
  (1.957,5.035)rectangle++(0.04,0.04) (1.863,5.798)rectangle++(0.04,0.04) (1.813,6.45)rectangle++(0.04,0.04) (1.797,7.006)rectangle++(0.04,0.04) (1.675,2.579)rectangle++(0.04,0.04)
  (-0.357,2.542)rectangle++(0.04,0.04) (-0.204,3.771)rectangle++(0.04,0.04) (1.069,1.282)rectangle++(0.04,0.04) (0.958,2.643)rectangle++(0.04,0.04) (0.918,3.805)rectangle++(0.04,0.04)
  (0.93,4.793)rectangle++(0.04,0.04) (0.98,5.633)rectangle++(0.04,0.04) (1.056,6.345)rectangle++(0.04,0.04) (1.601,2.227)rectangle++(0.04,0.04) (0.366,1.381)rectangle++(0.04,0.04)
  (0.364,2.755)rectangle++(0.04,0.04) (0.418,3.924)rectangle++(0.04,0.04) (0.51,4.915)rectangle++(0.04,0.04) (0.628,5.753)rectangle++(0.04,0.04) (0.761,6.461)rectangle++(0.04,0.04)
  (-1.549,3.186)rectangle++(0.04,0.04) (-1.191,4.366)rectangle++(0.04,0.04) (-0.84,5.355)rectangle++(0.04,0.04) (-1.582,2.574)rectangle++(0.04,0.04) (-1.244,3.847)rectangle++(0.04,0.04)
  (-0.905,4.916)rectangle++(0.04,0.04) (-0.575,5.811)rectangle++(0.04,0.04) (-0.259,6.559)rectangle++(0.04,0.04) (0.04,7.181)rectangle++(0.04,0.04) (0.319,7.699)rectangle++(0.04,0.04)
  (2.09,2.361)rectangle++(0.04,0.04) (1.869,3.519)rectangle++(0.04,0.04) (1.727,4.513)rectangle++(0.04,0.04) (1.646,5.363)rectangle++(0.04,0.04) (1.612,6.089)rectangle++(0.04,0.04)
  (1.611,6.707)rectangle++(0.04,0.04) (1.636,7.233)rectangle++(0.04,0.04) (1.678,7.679)rectangle++(0.04,0.04) (1.731,8.056)rectangle++(0.04,0.04) (1.791,8.375)rectangle++(0.04,0.04)
  (2.059,2.906)rectangle++(0.04,0.04) (1.864,3.984)rectangle++(0.04,0.04) (1.742,4.908)rectangle++(0.04,0.04) (1.096,2.061)rectangle++(0.04,0.04) (1.011,3.304)rectangle++(0.04,0.04)
  (0.99,4.364)rectangle++(0.04,0.04) (1.014,5.266)rectangle++(0.04,0.04) (1.07,6.032)rectangle++(0.04,0.04) (1.149,6.68)rectangle++(0.04,0.04) (-1.528,3.323)rectangle++(0.04,0.04)
  (-1.168,4.482)rectangle++(0.04,0.04) (-0.816,5.452)rectangle++(0.04,0.04) (-0.478,6.263)rectangle++(0.04,0.04) (-0.158,6.939)rectangle++(0.04,0.04) (0.141,7.501)rectangle++(0.04,0.04)
  (0.418,7.966)rectangle++(0.04,0.04) (0.672,8.351)rectangle++(0.04,0.04) (0.903,8.668)rectangle++(0.04,0.04) (1.112,8.927)rectangle++(0.04,0.04) (1.3,9.14)rectangle++(0.04,0.04)
  (1.468,9.313)rectangle++(0.04,0.04) (1.618,9.454)rectangle++(0.04,0.04) (1.752,9.567)rectangle++(0.04,0.04) (1.869,9.658)rectangle++(0.04,0.04) (1.973,9.731)rectangle++(0.04,0.04)
  (2.064,9.789)rectangle++(0.04,0.04) (2.144,9.834)rectangle++(0.04,0.04) (2.213,9.869)rectangle++(0.04,0.04) (2.274,9.897)rectangle++(0.04,0.04) (2.327,9.917)rectangle++(0.04,0.04)
  (2.372,9.933)rectangle++(0.04,0.04) (2.411,9.944)rectangle++(0.04,0.04) (2.445,9.952)rectangle++(0.04,0.04) (2.474,9.958)rectangle++(0.04,0.04) (2.499,9.961)rectangle++(0.04,0.04)
  (2.521,9.963)rectangle++(0.04,0.04) (2.539,9.964)rectangle++(0.04,0.04) (2.554,9.964)rectangle++(0.04,0.04) (2.568,9.964)rectangle++(0.04,0.04) (2.579,9.963)rectangle++(0.04,0.04)
  (2.588,9.961)rectangle++(0.04,0.04) (2.596,9.96)rectangle++(0.04,0.04) (2.603,9.958)rectangle++(0.04,0.04) (2.609,9.956)rectangle++(0.04,0.04) (2.613,9.955)rectangle++(0.04,0.04)
  (2.617,9.953)rectangle++(0.04,0.04) (2.621,9.952)rectangle++(0.04,0.04) (-2.085,4.381)rectangle++(0.04,0.04) (-1.599,5.404)rectangle++(0.04,0.04) (-1.145,6.253)rectangle++(0.04,0.04)
  (-0.725,6.957)rectangle++(0.04,0.04) (-0.34,7.539)rectangle++(0.04,0.04) (0.01,8.018)rectangle++(0.04,0.04) (-2.104,3.355)rectangle++(0.04,0.04) (-1.656,4.532)rectangle++(0.04,0.04)
  (-1.229,5.515)rectangle++(0.04,0.04) (1.711,1.434)rectangle++(0.04,0.04) (1.51,2.747)rectangle++(0.04,0.04) (-0.433,2.54)rectangle++(0.04,0.04) (-0.269,3.773)rectangle++(0.04,0.04)
  (-0.08,4.814)rectangle++(0.04,0.04) (0.122,5.691)rectangle++(0.04,0.04) (0.33,6.429)rectangle++(0.04,0.04) (0.535,7.048)rectangle++(0.04,0.04) (0.734,7.565)rectangle++(0.04,0.04)
  (0.925,7.997)rectangle++(0.04,0.04) (1.104,8.357)rectangle++(0.04,0.04) (1.27,8.655)rectangle++(0.04,0.04) (-0.02,1.368)rectangle++(0.04,0.04) (0.036,2.76)rectangle++(0.04,0.04)
  (0.138,3.941)rectangle++(0.04,0.04) (0.273,4.94)rectangle++(0.04,0.04) (1.325,1.687)rectangle++(0.04,0.04) (1.191,2.977)rectangle++(0.04,0.04) (1.13,4.079)rectangle++(0.04,0.04)
  (1.121,5.018)rectangle++(0.04,0.04) (1.151,5.817)rectangle++(0.04,0.04) (1.209,6.494)rectangle++(0.04,0.04) (1.285,7.068)rectangle++(0.04,0.04) (1.373,7.553)rectangle++(0.04,0.04)
  (1.467,7.961)rectangle++(0.04,0.04) (1.563,8.304)rectangle++(0.04,0.04) (1.659,8.592)rectangle++(0.04,0.04) (1.751,8.833)rectangle++(0.04,0.04) (1.84,9.035)rectangle++(0.04,0.04)
  (-2.002,4)rectangle++(0.04,0.04) (-1.462,2.008)rectangle++(0.04,0.04) (-1.164,3.362)rectangle++(0.04,0.04) (-0.857,4.5)rectangle++(0.04,0.04) (-0.551,5.456)rectangle++(0.04,0.04)
  (-0.252,6.256)rectangle++(0.04,0.04) (0.034,6.924)rectangle++(0.04,0.04) (0.303,7.48)rectangle++(0.04,0.04) (0.555,7.942)rectangle++(0.04,0.04) (-1.975,3.464)rectangle++(0.04,0.04)
  (-1.543,4.619)rectangle++(0.04,0.04) (-1.129,5.584)rectangle++(0.04,0.04) (-0.738,6.388)rectangle++(0.04,0.04) (-0.374,7.056)rectangle++(0.04,0.04) (-1.93,3.055)rectangle++(0.04,0.04)
  (-1.521,4.27)rectangle++(0.04,0.04) (-1.124,5.287)rectangle++(0.04,0.04) (-0.746,6.135)rectangle++(0.04,0.04) (-0.391,6.841)rectangle++(0.04,0.04) (-0.061,7.427)rectangle++(0.04,0.04)
  (0.243,7.911)rectangle++(0.04,0.04) (0.521,8.311)rectangle++(0.04,0.04) (2.232,2.56)rectangle++(0.04,0.04) (1.997,3.683)rectangle++(0.04,0.04) (1.843,4.647)rectangle++(0.04,0.04)
  (1.75,5.472)rectangle++(0.04,0.04) (1.704,6.178)rectangle++(0.04,0.04) (-1.287,3.34)rectangle++(0.04,0.04) (-1.147,2.028)rectangle++(0.04,0.04) (0.722,0.618)rectangle++(0.04,0.04)
  (0.637,2.093)rectangle++(0.04,0.04) (0.623,3.35)rectangle++(0.04,0.04) (-0.768,2.469)rectangle++(0.04,0.04) (-0.556,3.726)rectangle++(0.04,0.04) (-0.326,4.785)rectangle++(0.04,0.04)
  (-0.088,5.677)rectangle++(0.04,0.04) (0.15,6.425)rectangle++(0.04,0.04) (1.759,2.011)rectangle++(0.04,0.04) (1.573,3.235)rectangle++(0.04,0.04) (1.465,4.283)rectangle++(0.04,0.04)
  (1.414,5.178)rectangle++(0.04,0.04) (1.407,5.941)rectangle++(0.04,0.04) (1.431,6.59)rectangle++(0.04,0.04) (1.478,7.14)rectangle++(0.04,0.04) (1.54,7.606)rectangle++(0.04,0.04)
  (1.611,8)rectangle++(0.04,0.04) (1.687,8.332)rectangle++(0.04,0.04) (1.765,8.611)rectangle++(0.04,0.04) (2.129,2.955)rectangle++(0.04,0.04) (1.926,4.023)rectangle++(0.04,0.04)
  (1.795,4.939)rectangle++(0.04,0.04) (1.721,5.722)rectangle++(0.04,0.04) (1.69,6.391)rectangle++(0.04,0.04) (1.69,6.961)rectangle++(0.04,0.04) (1.713,7.446)rectangle++(0.04,0.04)
  (1.81,2.662)rectangle++(0.04,0.04) (1.643,3.787)rectangle++(0.04,0.04) (1.546,4.749)rectangle++(0.04,0.04) (1.502,5.571)rectangle++(0.04,0.04) (1.497,6.272)rectangle++(0.04,0.04)
  (1.521,6.867)rectangle++(0.04,0.04) (1.566,7.372)rectangle++(0.04,0.04) (-1.625,3.571)rectangle++(0.04,0.04) (-1.241,4.697)rectangle++(0.04,0.04) (-0.869,5.638)rectangle++(0.04,0.04)
  (-0.515,6.423)rectangle++(0.04,0.04) (-0.183,7.077)rectangle++(0.04,0.04) (0.125,7.619)rectangle++(0.04,0.04) (0.409,8.067)rectangle++(0.04,0.04) (0.668,8.437)rectangle++(0.04,0.04)
  (0.903,8.741)rectangle++(0.04,0.04) (2.295,2.763)rectangle++(0.04,0.04) (0.412,1.69)rectangle++(0.04,0.04) (0.394,0.943)rectangle++(0.04,0.04) (0.37,2.382)rectangle++(0.04,0.04)
  (0.408,3.606)rectangle++(0.04,0.04) (-0.877,2.476)rectangle++(0.04,0.04) (-0.649,3.736)rectangle++(0.04,0.04) (-0.404,4.798)rectangle++(0.04,0.04) (-1.349,2.552)rectangle++(0.04,0.04)
  (-1.047,3.819)rectangle++(0.04,0.04) (-1.224,2.188)rectangle++(0.04,0.04) (0.779,0.637)rectangle++(0.04,0.04) (0.685,2.107)rectangle++(0.04,0.04) (0.665,3.359)rectangle++(0.04,0.04)
  (0.697,4.425)rectangle++(0.04,0.04) (0.767,5.33)rectangle++(0.04,0.04) (-1.253,2.938)rectangle++(0.04,0.04) (-1.036,1.947)rectangle++(0.04,0.04) (0.683,0.628)rectangle++(0.04,0.04)
  (0.604,2.103)rectangle++(0.04,0.04) (0.595,3.359)rectangle++(0.04,0.04) (0.638,4.428)rectangle++(0.04,0.04) (0.717,5.334)rectangle++(0.04,0.04) (0.821,6.102)rectangle++(0.04,0.04)
  (0.939,6.75)rectangle++(0.04,0.04) (1.066,7.296)rectangle++(0.04,0.04) (1.196,7.755)rectangle++(0.04,0.04) (1.325,8.14)rectangle++(0.04,0.04) (1.449,8.462)rectangle++(0.04,0.04)
  (1.568,8.731)rectangle++(0.04,0.04) (1.68,8.955)rectangle++(0.04,0.04) (1.784,9.141)rectangle++(0.04,0.04) (2.274,3.088)rectangle++(0.04,0.04) (2.055,4.13)rectangle++(0.04,0.04)
  (1.909,5.024)rectangle++(0.04,0.04) (1.822,5.79)rectangle++(0.04,0.04) (1.778,6.445)rectangle++(0.04,0.04) (1.767,7.004)rectangle++(0.04,0.04) (1.78,7.479)rectangle++(0.04,0.04)
  (1.81,7.882)rectangle++(0.04,0.04) (1.851,8.223)rectangle++(0.04,0.04) (1.9,8.512)rectangle++(0.04,0.04) (1.954,8.755)rectangle++(0.04,0.04) (2.009,8.96)rectangle++(0.04,0.04)
  (2.063,9.132)rectangle++(0.04,0.04) (2.117,9.276)rectangle++(0.04,0.04) (2.168,9.396)rectangle++(0.04,0.04) (2.217,9.496)rectangle++(0.04,0.04) (2.262,9.579)rectangle++(0.04,0.04)
  (2.304,9.648)rectangle++(0.04,0.04) (2.342,9.705)rectangle++(0.04,0.04) (2.376,9.752)rectangle++(0.04,0.04) (2.408,9.79)rectangle++(0.04,0.04) (2.436,9.821)rectangle++(0.04,0.04)
  (2.461,9.847)rectangle++(0.04,0.04) (2.484,9.868)rectangle++(0.04,0.04) (2.504,9.884)rectangle++(0.04,0.04) (2.521,9.898)rectangle++(0.04,0.04) (2.537,9.908)rectangle++(0.04,0.04)
  (2.55,9.917)rectangle++(0.04,0.04) (2.562,9.924)rectangle++(0.04,0.04) (2.573,9.929)rectangle++(0.04,0.04) (2.582,9.933)rectangle++(0.04,0.04) (2.59,9.936)rectangle++(0.04,0.04)
  (2.596,9.938)rectangle++(0.04,0.04) (2.602,9.94)rectangle++(0.04,0.04) (2.607,9.941)rectangle++(0.04,0.04) (2.612,9.942)rectangle++(0.04,0.04) (2.615,9.942)rectangle++(0.04,0.04)
  (2.619,9.942)rectangle++(0.04,0.04) (2.621,9.942)rectangle++(0.04,0.04) (2.624,9.942)rectangle++(0.04,0.04) (2.626,9.942)rectangle++(0.04,0.04) (2.627,9.942)rectangle++(0.04,0.04)
  (2.629,9.942)rectangle++(0.04,0.04) (-2.08,4.381)rectangle++(0.04,0.04) (-1.595,5.403)rectangle++(0.04,0.04) (1.735,1.312)rectangle++(0.04,0.04) (1.525,2.642)rectangle++(0.04,0.04)
  (1.4,3.781)rectangle++(0.04,0.04) (0.831,1.701)rectangle++(0.04,0.04) (0.772,3.009)rectangle++(0.04,0.04) (0.775,4.123)rectangle++(0.04,0.04) (0.821,5.07)rectangle++(0.04,0.04)
  (0.899,5.873)rectangle++(0.04,0.04) (0.997,6.552)rectangle++(0.04,0.04) (1.107,7.126)rectangle++(0.04,0.04) (1.812,2.428)rectangle++(0.04,0.04) (1.635,3.587)rectangle++(0.04,0.04)
  (1.531,4.58)rectangle++(0.04,0.04) (1.482,5.428)rectangle++(0.04,0.04) (1.475,6.151)rectangle++(0.04,0.04) (1.484,2.29)rectangle++(0.04,0.04) (1.35,3.483)rectangle++(0.04,0.04)
  (1.285,4.503)rectangle++(0.04,0.04) (1.27,5.372)rectangle++(0.04,0.04) (1.296,2.05)rectangle++(0.04,0.04) (1.182,3.286)rectangle++(0.04,0.04) (1.134,4.342)rectangle++(0.04,0.04)
  (1.028,1.767)rectangle++(0.04,0.04) (0.943,3.057)rectangle++(0.04,0.04) (0.921,4.157)rectangle++(0.04,0.04) (0.947,5.093)rectangle++(0.04,0.04) (1.007,5.887)rectangle++(0.04,0.04)
  (1.089,6.56)rectangle++(0.04,0.04) (1.186,7.129)rectangle++(0.04,0.04) (1.291,7.608)rectangle++(0.04,0.04) (1.399,8.011)rectangle++(0.04,0.04) (1.508,8.35)rectangle++(0.04,0.04)
  (-1.891,3.773)rectangle++(0.04,0.04) (-1.38,1.984)rectangle++(0.04,0.04) (-1.096,3.338)rectangle++(0.04,0.04) (-0.8,4.477)rectangle++(0.04,0.04) (-0.503,5.434)rectangle++(0.04,0.04)
  (-0.213,6.235)rectangle++(0.04,0.04) (0.066,6.905)rectangle++(0.04,0.04) (0.33,7.462)rectangle++(0.04,0.04) (0.577,7.926)rectangle++(0.04,0.04) (0.805,8.31)rectangle++(0.04,0.04)
  (1.015,8.628)rectangle++(0.04,0.04) (1.206,8.889)rectangle++(0.04,0.04) (1.378,9.104)rectangle++(0.04,0.04) (-2.113,3.909)rectangle++(0.04,0.04) (-1.642,5.003)rectangle++(0.04,0.04)
  (-1.197,5.915)rectangle++(0.04,0.04) (-0.783,6.671)rectangle++(0.04,0.04) (-0.401,7.298)rectangle++(0.04,0.04) (-0.051,7.816)rectangle++(0.04,0.04) (0.267,8.242)rectangle++(0.04,0.04)
  (0.554,8.591)rectangle++(0.04,0.04) (0.813,8.876)rectangle++(0.04,0.04) (1.044,9.109)rectangle++(0.04,0.04) (1.249,9.297)rectangle++(0.04,0.04) (1.431,9.448)rectangle++(0.04,0.04)
  (1.592,9.57)rectangle++(0.04,0.04) (1.734,9.667)rectangle++(0.04,0.04) (1.859,9.744)rectangle++(0.04,0.04) (2.432,3.252)rectangle++(0.04,0.04) (2.195,4.263)rectangle++(0.04,0.04)
  (2.034,5.132)rectangle++(0.04,0.04) (1.932,5.877)rectangle++(0.04,0.04) (1.875,6.514)rectangle++(0.04,0.04) (1.852,7.058)rectangle++(0.04,0.04) (-0.02,1.113)rectangle++(0.04,0.04)
  (0.025,2.543)rectangle++(0.04,0.04) (0.121,3.756)rectangle++(0.04,0.04) (0.251,4.784)rectangle++(0.04,0.04) (0.402,5.653)rectangle++(0.04,0.04) (0.566,6.385)rectangle++(0.04,0.04)
  (0.734,7.001)rectangle++(0.04,0.04) (0.902,7.518)rectangle++(0.04,0.04) (1.065,7.95)rectangle++(0.04,0.04) (1.221,8.311)rectangle++(0.04,0.04) (2.127,2.742)rectangle++(0.04,0.04)
  (1.915,3.842)rectangle++(0.04,0.04) (1.779,4.785)rectangle++(0.04,0.04) (-0.91,3.051)rectangle++(0.04,0.04) (-0.653,4.226)rectangle++(0.04,0.04) (-0.388,5.214)rectangle++(0.04,0.04)
  (-0.124,6.044)rectangle++(0.04,0.04) (1.693,1.848)rectangle++(0.04,0.04) (1.511,3.1)rectangle++(0.04,0.04) (1.406,4.17)rectangle++(0.04,0.04) (1.36,5.085)rectangle++(0.04,0.04)
  (1.357,5.864)rectangle++(0.04,0.04) (1.386,6.526)rectangle++(0.04,0.04) (1.437,7.088)rectangle++(0.04,0.04) (1.503,7.564)rectangle++(0.04,0.04) (1.578,7.965)rectangle++(0.04,0.04)
  (1.657,8.303)rectangle++(0.04,0.04) (-1.849,3.797)rectangle++(0.04,0.04) (-1.422,4.898)rectangle++(0.04,0.04) (-1.015,5.816)rectangle++(0.04,0.04) (-0.632,6.58)rectangle++(0.04,0.04)
  (1.92,1.845)rectangle++(0.04,0.04) (1.704,3.088)rectangle++(0.04,0.04) (1.569,4.153)rectangle++(0.04,0.04) (1.498,5.063)rectangle++(0.04,0.04) (1.473,5.84)rectangle++(0.04,0.04)
  (1.484,6.501)rectangle++(0.04,0.04) (1.519,7.063)rectangle++(0.04,0.04) (1.572,7.539)rectangle++(0.04,0.04) (1.635,7.941)rectangle++(0.04,0.04) (-1.759,3.712)rectangle++(0.04,0.04)
  (-1.349,4.822)rectangle++(0.04,0.04) (-0.956,5.749)rectangle++(0.04,0.04) (-0.585,6.521)rectangle++(0.04,0.04) (-1.834,2.889)rectangle++(0.04,0.04) (1.067,0.647)rectangle++(0.04,0.04)
  (0.93,2.103)rectangle++(0.04,0.04) (0.873,3.347)rectangle++(0.04,0.04) (0.873,4.406)rectangle++(0.04,0.04) (0.916,5.306)rectangle++(0.04,0.04) (0.989,6.07)rectangle++(0.04,0.04)
  (1.081,6.716)rectangle++(0.04,0.04) (1.186,7.262)rectangle++(0.04,0.04) (1.296,7.721)rectangle++(0.04,0.04) (1.95,2.62)rectangle++(0.04,0.04) (0.424,1.566)rectangle++(0.04,0.04)
  (0.421,2.91)rectangle++(0.04,0.04) (-0.694,2.326)rectangle++(0.04,0.04) (-0.499,3.601)rectangle++(0.04,0.04) (-0.282,4.677)rectangle++(0.04,0.04) (-0.055,5.583)rectangle++(0.04,0.04)
  (0.174,6.344)rectangle++(0.04,0.04) (0.4,6.982)rectangle++(0.04,0.04) (0.617,7.515)rectangle++(0.04,0.04) (-1.852,3.384)rectangle++(0.04,0.04) (-1.441,4.547)rectangle++(0.04,0.04)
  (-1.045,5.519)rectangle++(0.04,0.04) (-0.67,6.329)rectangle++(0.04,0.04) (-1.801,2.827)rectangle++(0.04,0.04) (-1.42,4.071)rectangle++(0.04,0.04) (-1.046,5.114)rectangle++(0.04,0.04)
  (-0.687,5.985)rectangle++(0.04,0.04) (-0.347,6.711)rectangle++(0.04,0.04) (-0.028,7.314)rectangle++(0.04,0.04) (-1.929,3.192)rectangle++(0.04,0.04) (1.166,0.695)rectangle++(0.04,0.04)
  (1.016,2.14)rectangle++(0.04,0.04) (0.947,3.375)rectangle++(0.04,0.04) (0.785,1.486)rectangle++(0.04,0.04) (0.725,2.828)rectangle++(0.04,0.04) (-0.612,2.378)rectangle++(0.04,0.04)
  (-0.427,3.642)rectangle++(0.04,0.04) (-0.219,4.709)rectangle++(0.04,0.04) (1.334,1.503)rectangle++(0.04,0.04) (1.192,2.82)rectangle++(0.04,0.04) (1.124,3.946)rectangle++(0.04,0.04)
  (-0.822,2.716)rectangle++(0.04,0.04) (-0.593,3.937)rectangle++(0.04,0.04) (-0.348,4.967)rectangle++(0.04,0.04) (1.426,1.531)rectangle++(0.04,0.04) (1.271,2.841)rectangle++(0.04,0.04)
  (1.192,3.96)rectangle++(0.04,0.04) (-0.02,0.617)rectangle++(0.04,0.04) (0.005,2.121)rectangle++(0.04,0.04) (0.087,3.399)rectangle++(0.04,0.04) (0.208,4.482)rectangle++(0.04,0.04)
  (0.354,5.398)rectangle++(0.04,0.04) (0.514,6.17)rectangle++(0.04,0.04) (0.682,6.82)rectangle++(0.04,0.04) (0.85,7.366)rectangle++(0.04,0.04) (1.015,7.823)rectangle++(0.04,0.04)
  (1.174,8.205)rectangle++(0.04,0.04) (1.324,8.524)rectangle++(0.04,0.04) (1.464,8.789)rectangle++(0.04,0.04) (1.594,9.008)rectangle++(0.04,0.04) (1.713,9.189)rectangle++(0.04,0.04)
  (1.821,9.338)rectangle++(0.04,0.04) (1.919,9.461)rectangle++(0.04,0.04) (2.008,9.561)rectangle++(0.04,0.04) (2.359,3.247)rectangle++(0.04,0.04) (2.133,4.262)rectangle++(0.04,0.04)
  (1.981,5.133)rectangle++(0.04,0.04) (1.123,2.177)rectangle++(0.04,0.04) (1.039,3.402)rectangle++(0.04,0.04) (-0.698,2.576)rectangle++(0.04,0.04) (-0.492,3.814)rectangle++(0.04,0.04)
  (-0.268,4.858)rectangle++(0.04,0.04) (-0.036,5.736)rectangle++(0.04,0.04) (-1.52,2.843)rectangle++(0.04,0.04) (-1.18,4.073)rectangle++(0.04,0.04) (-0.842,5.106)rectangle++(0.04,0.04)
  (-0.514,5.97)rectangle++(0.04,0.04) (-0.2,6.691)rectangle++(0.04,0.04) (-1.801,3.015)rectangle++(0.04,0.04) (-1.412,4.231)rectangle++(0.04,0.04) (-1.034,5.249)rectangle++(0.04,0.04)
  (-0.671,6.099)rectangle++(0.04,0.04) (-1.741,2.777)rectangle++(0.04,0.04) (-1.371,4.026)rectangle++(0.04,0.04) (-1.007,5.073)rectangle++(0.04,0.04) (-1.542,2.474)rectangle++(0.04,0.04)
  (-1.214,3.76)rectangle++(0.04,0.04) (-0.883,4.841)rectangle++(0.04,0.04) (-0.559,5.746)rectangle++(0.04,0.04) (-0.248,6.503)rectangle++(0.04,0.04) (0.047,7.134)rectangle++(0.04,0.04)
  (0.323,7.658)rectangle++(0.04,0.04) (0.579,8.093)rectangle++(0.04,0.04) (0.814,8.452)rectangle++(0.04,0.04) (1.027,8.748)rectangle++(0.04,0.04) (1.221,8.991)rectangle++(0.04,0.04)
  (-2.115,3.848)rectangle++(0.04,0.04) (-1.646,4.951)rectangle++(0.04,0.04) (-1.203,5.871)rectangle++(0.04,0.04) (-0.79,6.634)rectangle++(0.04,0.04) (-0.408,7.267)rectangle++(0.04,0.04)
  (-0.059,7.79)rectangle++(0.04,0.04) (0.26,8.22)rectangle++(0.04,0.04) (0.547,8.572)rectangle++(0.04,0.04) (0.806,8.861)rectangle++(0.04,0.04) (1.037,9.096)rectangle++(0.04,0.04)
  (1.243,9.286)rectangle++(0.04,0.04) (1.426,9.44)rectangle++(0.04,0.04) (-2.19,3.994)rectangle++(0.04,0.04) (-1.704,5.078)rectangle++(0.04,0.04) (-1.248,5.981)rectangle++(0.04,0.04)
  (-0.823,6.73)rectangle++(0.04,0.04) (-0.433,7.35)rectangle++(0.04,0.04) (-0.076,7.861)rectangle++(0.04,0.04) (0.247,8.281)rectangle++(0.04,0.04) (0.539,8.625)rectangle++(0.04,0.04)
  (0.801,8.906)rectangle++(0.04,0.04) (1.035,9.134)rectangle++(0.04,0.04) (1.243,9.319)rectangle++(0.04,0.04) (1.427,9.467)rectangle++(0.04,0.04) (1.59,9.586)rectangle++(0.04,0.04)
  (1.732,9.681)rectangle++(0.04,0.04) (1.858,9.756)rectangle++(0.04,0.04) (1.967,9.814)rectangle++(0.04,0.04) (2.062,9.86)rectangle++(0.04,0.04) (2.145,9.894)rectangle++(0.04,0.04)
  (2.217,9.921)rectangle++(0.04,0.04) (2.279,9.94)rectangle++(0.04,0.04) (2.333,9.954)rectangle++(0.04,0.04) (2.379,9.964)rectangle++(0.04,0.04) (2.418,9.97)rectangle++(0.04,0.04)
  (2.452,9.974)rectangle++(0.04,0.04) (2.481,9.976)rectangle++(0.04,0.04) (2.506,9.977)rectangle++(0.04,0.04) (2.527,9.976)rectangle++(0.04,0.04) (2.545,9.975)rectangle++(0.04,0.04)
  (2.56,9.973)rectangle++(0.04,0.04) (2.572,9.971)rectangle++(0.04,0.04) (2.583,9.969)rectangle++(0.04,0.04) (2.592,9.966)rectangle++(0.04,0.04) (2.6,9.964)rectangle++(0.04,0.04)
  (2.606,9.961)rectangle++(0.04,0.04) (2.612,9.959)rectangle++(0.04,0.04) (2.616,9.957)rectangle++(0.04,0.04) (2.62,9.955)rectangle++(0.04,0.04) (2.623,9.953)rectangle++(0.04,0.04)
  (2.625,9.951)rectangle++(0.04,0.04) (2.627,9.95)rectangle++(0.04,0.04) (2.629,9.949)rectangle++(0.04,0.04) (2.63,9.947)rectangle++(0.04,0.04) (2.632,9.946)rectangle++(0.04,0.04)
  (2.632,9.945)rectangle++(0.04,0.04) (2.633,9.944)rectangle++(0.04,0.04) (2.634,9.944)rectangle++(0.04,0.04) (2.634,9.943)rectangle++(0.04,0.04) (2.371,3.501)rectangle++(0.04,0.04)
  (-0.457,2.905)rectangle++(0.04,0.04) (-0.275,4.083)rectangle++(0.04,0.04) (-0.072,5.078)rectangle++(0.04,0.04) (-1.356,2.69)rectangle++(0.04,0.04) (-1.047,3.937)rectangle++(0.04,0.04)
  (-0.735,4.984)rectangle++(0.04,0.04) (-0.427,5.862)rectangle++(0.04,0.04) (-1.631,2.78)rectangle++(0.04,0.04) (-1.277,4.025)rectangle++(0.04,0.04) (-0.927,5.068)rectangle++(0.04,0.04)
  (-0.587,5.941)rectangle++(0.04,0.04) (-0.264,6.67)rectangle++(0.04,0.04) (0.04,7.276)rectangle++(0.04,0.04) (0.323,7.779)rectangle++(0.04,0.04) (0.584,8.196)rectangle++(0.04,0.04)
  (0.822,8.539)rectangle++(0.04,0.04) (1.038,8.822)rectangle++(0.04,0.04) (1.233,9.053)rectangle++(0.04,0.04) (1.408,9.242)rectangle++(0.04,0.04) (1.564,9.396)rectangle++(0.04,0.04)
  (1.703,9.52)rectangle++(0.04,0.04) (1.826,9.62)rectangle++(0.04,0.04) (1.935,9.7)rectangle++(0.04,0.04) (2.03,9.764)rectangle++(0.04,0.04) (2.114,9.814)rectangle++(0.04,0.04)
  (2.188,9.854)rectangle++(0.04,0.04) (2.251,9.884)rectangle++(0.04,0.04) (2.307,9.908)rectangle++(0.04,0.04) (2.355,9.926)rectangle++(0.04,0.04) (2.396,9.939)rectangle++(0.04,0.04)
  (2.432,9.948)rectangle++(0.04,0.04) (2.463,9.955)rectangle++(0.04,0.04) (2.49,9.959)rectangle++(0.04,0.04) (-2.113,4.353)rectangle++(0.04,0.04) (-1.624,5.381)rectangle++(0.04,0.04)
  (-1.167,6.235)rectangle++(0.04,0.04) (1.903,1.623)rectangle++(0.04,0.04) (1.681,2.899)rectangle++(0.04,0.04) (1.542,3.994)rectangle++(0.04,0.04) (1.468,4.929)rectangle++(0.04,0.04)
  (1.142,1.995)rectangle++(0.04,0.04) (0.37,1.206)rectangle++(0.04,0.04) (0.36,2.606)rectangle++(0.04,0.04) (0.408,3.797)rectangle++(0.04,0.04) (0.497,4.807)rectangle++(0.04,0.04)
  (0.612,5.663)rectangle++(0.04,0.04) (0.745,6.385)rectangle++(0.04,0.04) (0.886,6.994)rectangle++(0.04,0.04) (1.031,7.505)rectangle++(0.04,0.04) (-0.02,1.184)rectangle++(0.04,0.04)
  (0.028,2.603)rectangle++(0.04,0.04) (0.126,3.808)rectangle++(0.04,0.04) (0.257,4.828)rectangle++(0.04,0.04) (0.409,5.69)rectangle++(0.04,0.04) (0.573,6.416)rectangle++(0.04,0.04)
  (0.742,7.027)rectangle++(0.04,0.04) (0.909,7.539)rectangle++(0.04,0.04) (1.072,7.968)rectangle++(0.04,0.04) (2.053,2.621)rectangle++(0.04,0.04) (1.848,3.742)rectangle++(0.04,0.04)
  (1.718,4.703)rectangle++(0.04,0.04) (1.646,5.525)rectangle++(0.04,0.04) (1.618,6.227)rectangle++(0.04,0.04) (1.622,6.824)rectangle++(0.04,0.04) (1.65,7.332)rectangle++(0.04,0.04)
  (1.693,7.762)rectangle++(0.04,0.04) (1.748,8.126)rectangle++(0.04,0.04) (1.808,8.434)rectangle++(0.04,0.04) (1.872,8.693)rectangle++(0.04,0.04) (1.937,8.91)rectangle++(0.04,0.04)
  (2.001,9.092)rectangle++(0.04,0.04) (2.062,9.245)rectangle++(0.04,0.04) (2.12,9.372)rectangle++(0.04,0.04) (-2.034,4.138)rectangle++(0.04,0.04) (-1.565,5.195)rectangle++(0.04,0.04)
  (-1.125,6.075)rectangle++(0.04,0.04) (-0.715,6.805)rectangle++(0.04,0.04) (-0.338,7.409)rectangle++(0.04,0.04) (0.007,7.907)rectangle++(0.04,0.04) (0.32,8.317)rectangle++(0.04,0.04)
  (0.602,8.653)rectangle++(0.04,0.04) (-2.15,3.631)rectangle++(0.04,0.04) (-1.685,4.769)rectangle++(0.04,0.04) (-1.244,5.717)rectangle++(0.04,0.04) (-0.831,6.505)rectangle++(0.04,0.04)
  (-0.448,7.159)rectangle++(0.04,0.04) (-0.097,7.699)rectangle++(0.04,0.04) (0.224,8.144)rectangle++(0.04,0.04) (0.514,8.51)rectangle++(0.04,0.04) (0.775,8.809)rectangle++(0.04,0.04)
  (1.009,9.053)rectangle++(0.04,0.04) (1.217,9.251)rectangle++(0.04,0.04) (1.403,9.411)rectangle++(0.04,0.04) (1.566,9.539)rectangle++(0.04,0.04) (1.711,9.642)rectangle++(0.04,0.04)
  (2.426,3.189)rectangle++(0.04,0.04) (2.187,4.21)rectangle++(0.04,0.04) (2.025,5.087)rectangle++(0.04,0.04) (1.923,5.839)rectangle++(0.04,0.04) (1.866,6.483)rectangle++(0.04,0.04)
  (1.843,7.032)rectangle++(0.04,0.04) (1.846,7.499)rectangle++(0.04,0.04) (1.867,7.897)rectangle++(0.04,0.04) (1.9,8.234)rectangle++(0.04,0.04) (1.942,8.519)rectangle++(0.04,0.04)
  (1.99,8.76)rectangle++(0.04,0.04) (2.039,8.962)rectangle++(0.04,0.04) (2.186,3.111)rectangle++(0.04,0.04) (1.98,4.153)rectangle++(0.04,0.04) (1.847,5.047)rectangle++(0.04,0.04)
  (1.77,5.812)rectangle++(0.04,0.04) (1.345,2.285)rectangle++(0.04,0.04) (1.232,3.485)rectangle++(0.04,0.04) (1.184,4.509)rectangle++(0.04,0.04) (1.185,5.381)rectangle++(0.04,0.04)
  (1.22,6.123)rectangle++(0.04,0.04) (1.28,6.752)rectangle++(0.04,0.04) (1.356,7.284)rectangle++(0.04,0.04) (1.442,7.734)rectangle++(0.04,0.04) (1.932,2.661)rectangle++(0.04,0.04)
  (1.746,3.781)rectangle++(0.04,0.04) (1.633,4.74)rectangle++(0.04,0.04) (-0.927,3.007)rectangle++(0.04,0.04) (-0.67,4.19)rectangle++(0.04,0.04) (-0.404,5.184)rectangle++(0.04,0.04)
  (-0.138,6.019)rectangle++(0.04,0.04) (0.121,6.718)rectangle++(0.04,0.04) (0.369,7.301)rectangle++(0.04,0.04) (0.604,7.788)rectangle++(0.04,0.04) (0.823,8.192)rectangle++(0.04,0.04)
  (1.025,8.526)rectangle++(0.04,0.04) (1.21,8.802)rectangle++(0.04,0.04) (1.378,9.03)rectangle++(0.04,0.04) (1.531,9.216)rectangle++(0.04,0.04) (1.667,9.369)rectangle++(0.04,0.04)
  (1.79,9.493)rectangle++(0.04,0.04) (1.899,9.594)rectangle++(0.04,0.04) (1.996,9.675)rectangle++(0.04,0.04) (2.081,9.74)rectangle++(0.04,0.04) (2.156,9.792)rectangle++(0.04,0.04)
  (2.222,9.833)rectangle++(0.04,0.04) (2.28,9.865)rectangle++(0.04,0.04) (2.331,9.891)rectangle++(0.04,0.04) (2.374,9.91)rectangle++(0.04,0.04) (2.401,3.426)rectangle++(0.04,0.04)
  (2.176,4.412)rectangle++(0.04,0.04) (2.024,5.259)rectangle++(0.04,0.04) (1.928,5.986)rectangle++(0.04,0.04) (1.876,6.607)rectangle++(0.04,0.04) (1.857,7.137)rectangle++(0.04,0.04)
  (1.862,7.588)rectangle++(0.04,0.04) (1.884,7.972)rectangle++(0.04,0.04) (1.918,8.297)rectangle++(0.04,0.04) (1.96,8.572)rectangle++(0.04,0.04) (-1.858,3.926)rectangle++(0.04,0.04)
  (-1.424,5.007)rectangle++(0.04,0.04) (-1.013,5.909)rectangle++(0.04,0.04) (-0.627,6.66)rectangle++(0.04,0.04) (-1.878,2.91)rectangle++(0.04,0.04) (-1.482,4.145)rectangle++(0.04,0.04)
  (-1.096,5.179)rectangle++(0.04,0.04) (1.597,1.388)rectangle++(0.04,0.04) (1.411,2.712)rectangle++(0.04,0.04) (1.305,3.845)rectangle++(0.04,0.04) (1.261,4.812)rectangle++(0.04,0.04)
  (1.262,5.636)rectangle++(0.04,0.04) (1.296,6.336)rectangle++(0.04,0.04) (1.353,6.93)rectangle++(0.04,0.04) (1.425,7.433)rectangle++(0.04,0.04) (1.506,7.857)rectangle++(0.04,0.04)
  (1.593,8.214)rectangle++(0.04,0.04) (1.68,8.515)rectangle++(0.04,0.04) (1.766,8.767)rectangle++(0.04,0.04) (1.85,8.977)rectangle++(0.04,0.04) (1.929,9.153)rectangle++(0.04,0.04)
  (-2.015,4.046)rectangle++(0.04,0.04) (-1.553,5.116)rectangle++(0.04,0.04) (-1.118,6.007)rectangle++(0.04,0.04) (-0.712,6.747)rectangle++(0.04,0.04) (-0.338,7.36)rectangle++(0.04,0.04)
  (0.005,7.865)rectangle++(0.04,0.04) (0.317,8.282)rectangle++(0.04,0.04) (0.598,8.623)rectangle++(0.04,0.04) (0.851,8.902)rectangle++(0.04,0.04) (1.078,9.129)rectangle++(0.04,0.04)
  (1.279,9.312)rectangle++(0.04,0.04) (1.457,9.461)rectangle++(0.04,0.04) (1.615,9.579)rectangle++(0.04,0.04) (1.754,9.674)rectangle++(0.04,0.04) (1.875,9.749)rectangle++(0.04,0.04)
  (1.982,9.808)rectangle++(0.04,0.04) (2.432,3.299)rectangle++(0.04,0.04) (2.197,4.303)rectangle++(0.04,0.04) (2.037,5.166)rectangle++(0.04,0.04) (1.936,5.906)rectangle++(0.04,0.04)
  (1.346,2.351)rectangle++(0.04,0.04) (1.236,3.541)rectangle++(0.04,0.04) (1.19,4.556)rectangle++(0.04,0.04) (1.191,5.421)rectangle++(0.04,0.04) (-1.192,3.056)rectangle++(0.04,0.04)
  (-0.894,4.241)rectangle++(0.04,0.04) (-0.592,5.237)rectangle++(0.04,0.04) (-0.296,6.071)rectangle++(0.04,0.04) (-0.011,6.769)rectangle++(0.04,0.04) (0.259,7.35)rectangle++(0.04,0.04)
  (0.512,7.833)rectangle++(0.04,0.04) (0.746,8.234)rectangle++(0.04,0.04) (0.962,8.565)rectangle++(0.04,0.04) (1.158,8.838)rectangle++(0.04,0.04) (1.335,9.062)rectangle++(0.04,0.04)
  (1.495,9.246)rectangle++(0.04,0.04) (1.639,9.395)rectangle++(0.04,0.04) (1.767,9.517)rectangle++(0.04,0.04) (1.88,9.615)rectangle++(0.04,0.04) (2.393,3.226)rectangle++(0.04,0.04)
  (2.161,4.243)rectangle++(0.04,0.04) (-0.692,3.019)rectangle++(0.04,0.04) (-0.47,4.19)rectangle++(0.04,0.04) (-0.234,5.177)rectangle++(0.04,0.04) (0.006,6.006)rectangle++(0.04,0.04)
  (0.243,6.701)rectangle++(0.04,0.04) (0.473,7.282)rectangle++(0.04,0.04) (0.691,7.767)rectangle++(0.04,0.04) (0.896,8.171)rectangle++(0.04,0.04) (1.086,8.505)rectangle++(0.04,0.04)
  (1.261,8.782)rectangle++(0.04,0.04) (1.421,9.011)rectangle++(0.04,0.04) (-2.08,3.898)rectangle++(0.04,0.04) (-1.614,4.993)rectangle++(0.04,0.04) (-1.174,5.905)rectangle++(0.04,0.04)
  (-0.764,6.662)rectangle++(0.04,0.04) (1.963,1.83)rectangle++(0.04,0.04) (1.739,3.073)rectangle++(0.04,0.04) (0.582,1.62)rectangle++(0.04,0.04) (0.558,2.95)rectangle++(0.04,0.04)
  (0.59,4.081)rectangle++(0.04,0.04) (0.662,5.042)rectangle++(0.04,0.04) (0.762,5.855)rectangle++(0.04,0.04) (0.88,6.543)rectangle++(0.04,0.04) (1.008,7.122)rectangle++(0.04,0.04)
  (1.139,7.61)rectangle++(0.04,0.04) (1.27,8.019)rectangle++(0.04,0.04) (2.037,2.685)rectangle++(0.04,0.04) (-0.02,0.413)rectangle++(0.04,0.04) (-0.003,1.948)rectangle++(0.04,0.04)
  (0.073,3.252)rectangle++(0.04,0.04) (0.19,4.357)rectangle++(0.04,0.04) (0.334,5.292)rectangle++(0.04,0.04) (0.493,6.081)rectangle++(0.04,0.04) (-1.504,3.04)rectangle++(0.04,0.04)
  (-1.159,4.241)rectangle++(0.04,0.04) (-0.818,5.247)rectangle++(0.04,0.04) (-0.487,6.089)rectangle++(0.04,0.04) (-0.173,6.791)rectangle++(0.04,0.04) (0.123,7.376)rectangle++(0.04,0.04)
  (0.397,7.861)rectangle++(0.04,0.04) (0.65,8.262)rectangle++(0.04,0.04) (-2.039,3.556)rectangle++(0.04,0.04) (1.284,0.753)rectangle++(0.04,0.04) (1.119,2.185)rectangle++(0.04,0.04)
  (1.037,3.409)rectangle++(0.04,0.04) (1.015,4.452)rectangle++(0.04,0.04) (1.039,5.34)rectangle++(0.04,0.04) (1.095,6.094)rectangle++(0.04,0.04) (1.172,6.732)rectangle++(0.04,0.04)
  (1.263,7.271)rectangle++(0.04,0.04) (1.362,7.726)rectangle++(0.04,0.04) (-0.02,1.219)rectangle++(0.04,0.04) (-0.342,1.853)rectangle++(0.04,0.04) (-0.219,3.185)rectangle++(0.04,0.04)
  (-0.061,4.312)rectangle++(0.04,0.04) (0.118,5.264)rectangle++(0.04,0.04) (0.309,6.066)rectangle++(0.04,0.04) (0.503,6.74)rectangle++(0.04,0.04) (0.695,7.305)rectangle++(0.04,0.04)
  (0.881,7.777)rectangle++(0.04,0.04) (1.058,8.172)rectangle++(0.04,0.04) (1.224,8.5)rectangle++(0.04,0.04) (2.179,2.788)rectangle++(0.04,0.04) (1.962,3.879)rectangle++(0.04,0.04)
  (-0.637,2.894)rectangle++(0.04,0.04) (-0.428,4.081)rectangle++(0.04,0.04) (-1.168,2.388)rectangle++(0.04,0.04) (-0.899,3.673)rectangle++(0.04,0.04) (1.146,1.078)rectangle++(0.04,0.04)
  (1.015,2.466)rectangle++(0.04,0.04) (0.959,3.652)rectangle++(0.04,0.04) (0.959,4.662)rectangle++(0.04,0.04) (-1.041,2.835)rectangle++(0.04,0.04) (-0.774,4.048)rectangle++(0.04,0.04)
  (-0.498,5.068)rectangle++(0.04,0.04) (-0.223,5.924)rectangle++(0.04,0.04) (0.045,6.64)rectangle++(0.04,0.04) (0.302,7.239)rectangle++(0.04,0.04) (0.544,7.737)rectangle++(0.04,0.04)
  (0.77,8.151)rectangle++(0.04,0.04) (2.149,2.586)rectangle++(0.04,0.04) (1.928,3.709)rectangle++(0.04,0.04) (1.785,4.671)rectangle++(0.04,0.04) (1.702,5.495)rectangle++(0.04,0.04)
  (1.664,6.199)rectangle++(0.04,0.04) (1.66,6.799)rectangle++(0.04,0.04) (1.681,7.309)rectangle++(0.04,0.04) (1.719,7.742)rectangle++(0.04,0.04) (1.769,8.108)rectangle++(0.04,0.04)
  (1.826,8.417)rectangle++(0.04,0.04) (1.886,8.678)rectangle++(0.04,0.04) (1.948,8.897)rectangle++(0.04,0.04) (2.01,9.081)rectangle++(0.04,0.04) (2.069,9.234)rectangle++(0.04,0.04)
  (-2.008,4.096)rectangle++(0.04,0.04) (-1.545,5.159)rectangle++(0.04,0.04) (-1.109,6.043)rectangle++(0.04,0.04) (-0.704,6.777)rectangle++(0.04,0.04) (-0.329,7.385)rectangle++(0.04,0.04)
  (0.013,7.886)rectangle++(0.04,0.04) (-2.069,3.327)rectangle++(0.04,0.04) (-1.628,4.507)rectangle++(0.04,0.04) (1.489,1.088)rectangle++(0.04,0.04) (1.307,2.462)rectangle++(0.04,0.04)
  (1.207,3.637)rectangle++(0.04,0.04) (1.169,4.639)rectangle++(0.04,0.04) (1.177,5.493)rectangle++(0.04,0.04) (1.344,2.054)rectangle++(0.04,0.04) (1.222,3.289)rectangle++(0.04,0.04)
  (1.168,4.343)rectangle++(0.04,0.04) (-0.917,2.813)rectangle++(0.04,0.04) (-0.669,4.024)rectangle++(0.04,0.04) (1.21,1.222)rectangle++(0.04,0.04) (1.075,2.586)rectangle++(0.04,0.04)
  (1.015,3.752)rectangle++(0.04,0.04) (1.011,4.745)rectangle++(0.04,0.04) (1.047,5.589)rectangle++(0.04,0.04) (1.111,6.305)rectangle++(0.04,0.04) (1.194,6.911)rectangle++(0.04,0.04)
  (1.289,7.423)rectangle++(0.04,0.04) (1.391,7.854)rectangle++(0.04,0.04) (1.494,8.216)rectangle++(0.04,0.04) (1.596,8.52)rectangle++(0.04,0.04) (1.696,8.775)rectangle++(0.04,0.04)
  (-1.963,3.909)rectangle++(0.04,0.04) (-1.515,4.998)rectangle++(0.04,0.04) (-1.09,5.905)rectangle++(0.04,0.04) (-0.692,6.659)rectangle++(0.04,0.04) (-0.324,7.284)rectangle++(0.04,0.04)
  (0.013,7.801)rectangle++(0.04,0.04) (0.321,8.226)rectangle++(0.04,0.04) (2.238,2.488)rectangle++(0.04,0.04) (1.999,3.621)rectangle++(0.04,0.04) (1.842,4.594)rectangle++(0.04,0.04)
  (1.747,5.428)rectangle++(0.04,0.04) (1.7,6.14)rectangle++(0.04,0.04) (1.447,2.346)rectangle++(0.04,0.04) (0.422,1.369)rectangle++(0.04,0.04) (0.412,2.743)rectangle++(0.04,0.04)
  (0.457,3.911)rectangle++(0.04,0.04) (-0.947,2.555)rectangle++(0.04,0.04) (-0.705,3.806)rectangle++(0.04,0.04) (-0.449,4.859)rectangle++(0.04,0.04) (-0.189,5.744)rectangle++(0.04,0.04)
  (0.067,6.487)rectangle++(0.04,0.04) (0.314,7.107)rectangle++(0.04,0.04) (0.549,7.625)rectangle++(0.04,0.04) (0.769,8.055)rectangle++(0.04,0.04) (2.123,2.563)rectangle++(0.04,0.04)
  (1.905,3.69)rectangle++(0.04,0.04) (1.764,4.657)rectangle++(0.04,0.04) (1.684,5.484)rectangle++(0.04,0.04) (1.648,6.19)rectangle++(0.04,0.04) (1.469,2.344)rectangle++(0.04,0.04)
  (1.34,3.53)rectangle++(0.04,0.04) (-0.671,2.674)rectangle++(0.04,0.04) (-0.466,3.896)rectangle++(0.04,0.04) (-1.127,2.339)rectangle++(0.04,0.04) (-0.867,3.629)rectangle++(0.04,0.04)
  (-0.594,4.716)rectangle++(0.04,0.04) (-1.366,2.49)rectangle++(0.04,0.04) (-1.064,3.767)rectangle++(0.04,0.04) (-1.213,2.173)rectangle++(0.04,0.04) (-0.947,3.492)rectangle++(0.04,0.04)
  (-0.667,4.602)rectangle++(0.04,0.04) (-0.385,5.535)rectangle++(0.04,0.04) (-1.537,2.718)rectangle++(0.04,0.04) (-1.2,3.968)rectangle++(0.04,0.04) (-0.864,5.017)rectangle++(0.04,0.04)
  (-0.536,5.895)rectangle++(0.04,0.04) (-0.222,6.629)rectangle++(0.04,0.04) (0.075,7.239)rectangle++(0.04,0.04) (1.998,2.187)rectangle++(0.04,0.04) (0.295,1.474)rectangle++(0.04,0.04)
  (0.308,2.838)rectangle++(0.04,0.04) (-0.697,2.284)rectangle++(0.04,0.04) (-0.504,3.566)rectangle++(0.04,0.04) (-0.288,4.647)rectangle++(0.04,0.04) (-0.061,5.558)rectangle++(0.04,0.04)
  (0.168,6.323)rectangle++(0.04,0.04) (0.394,6.964)rectangle++(0.04,0.04) (-1.753,3.212)rectangle++(0.04,0.04) (-1.364,4.396)rectangle++(0.04,0.04) (-0.986,5.388)rectangle++(0.04,0.04)
  (-0.624,6.215)rectangle++(0.04,0.04) (-0.284,6.904)rectangle++(0.04,0.04) (0.032,7.476)rectangle++(0.04,0.04) (0.324,7.949)rectangle++(0.04,0.04) (0.591,8.34)rectangle++(0.04,0.04)
  (2.229,2.585)rectangle++(0.04,0.04) (1.996,3.705)rectangle++(0.04,0.04) (0.72,1.838)rectangle++(0.04,0.04) (0.684,3.13)rectangle++(0.04,0.04) (0.704,4.229)rectangle++(0.04,0.04)
  (0.766,5.163)rectangle++(0.04,0.04) (0.855,5.954)rectangle++(0.04,0.04) (0.963,6.623)rectangle++(0.04,0.04) (1.081,7.187)rectangle++(0.04,0.04) (1.204,7.662)rectangle++(0.04,0.04)
  (1.328,8.061)rectangle++(0.04,0.04) (1.449,8.395)rectangle++(0.04,0.04) (1.565,8.674)rectangle++(0.04,0.04) (2.176,2.919)rectangle++(0.04,0.04) (1.965,3.99)rectangle++(0.04,0.04)
  (1.827,4.909)rectangle++(0.04,0.04) (1.747,5.696)rectangle++(0.04,0.04) (1.711,6.368)rectangle++(0.04,0.04) (1.509,2.403)rectangle++(0.04,0.04) (1.377,3.578)rectangle++(0.04,0.04)
  (1.311,4.583)rectangle++(0.04,0.04) (1.295,5.439)rectangle++(0.04,0.04) (-0.02,0.853)rectangle++(0.04,0.04) (0.015,2.322)rectangle++(0.04,0.04) (0.103,3.57)rectangle++(0.04,0.04)
  (-0.929,2.398)rectangle++(0.04,0.04) (-0.696,3.672)rectangle++(0.04,0.04) (-0.447,4.745)rectangle++(0.04,0.04) (-0.192,5.647)rectangle++(0.04,0.04) (0.06,6.404)rectangle++(0.04,0.04)
  (0.305,7.037)rectangle++(0.04,0.04) (0.539,7.566)rectangle++(0.04,0.04) (0.758,8.005)rectangle++(0.04,0.04) (-0.02,1.264)rectangle++(0.04,0.04) (0.031,2.671)rectangle++(0.04,0.04)
  (0.131,3.866)rectangle++(0.04,0.04) (0.264,4.877)rectangle++(0.04,0.04) (0.417,5.731)rectangle++(0.04,0.04) (0.582,6.451)rectangle++(0.04,0.04) (0.75,7.056)rectangle++(0.04,0.04)
  (0.918,7.564)rectangle++(0.04,0.04) (1.08,7.989)rectangle++(0.04,0.04) (1.236,8.343)rectangle++(0.04,0.04) (1.382,8.639)rectangle++(0.04,0.04) (2.194,2.863)rectangle++(0.04,0.04)
  (1.977,3.942)rectangle++(0.04,0.04) (1.836,4.867)rectangle++(0.04,0.04) (1.753,5.66)rectangle++(0.04,0.04) (1.714,6.337)rectangle++(0.04,0.04) (1.709,6.914)rectangle++(0.04,0.04)
  (1.727,7.405)rectangle++(0.04,0.04) (1.762,7.821)rectangle++(0.04,0.04) (1.808,8.174)rectangle++(0.04,0.04) (1.862,8.472)rectangle++(0.04,0.04) (1.919,8.723)rectangle++(0.04,0.04)
  (1.978,8.934)rectangle++(0.04,0.04) (2.036,9.111)rectangle++(0.04,0.04) (-1.983,4.062)rectangle++(0.04,0.04) (-1.474,2.027)rectangle++(0.04,0.04) (-1.174,3.378)rectangle++(0.04,0.04)
  (-0.865,4.514)rectangle++(0.04,0.04) (1.376,1.289)rectangle++(0.04,0.04) (1.219,2.636)rectangle++(0.04,0.04) (1.14,3.788)rectangle++(0.04,0.04) (1.118,4.771)rectangle++(0.04,0.04)
  (1.139,5.607)rectangle++(0.04,0.04) (1.19,6.316)rectangle++(0.04,0.04) (1.573,2.255)rectangle++(0.04,0.04) (1.425,3.45)rectangle++(0.04,0.04) (-0.633,2.676)rectangle++(0.04,0.04)
  (-0.844,2.032)rectangle++(0.04,0.04) (-0.638,3.357)rectangle++(0.04,0.04) (-0.41,4.475)rectangle++(0.04,0.04) (-0.172,5.417)rectangle++(0.04,0.04) (1.525,1.685)rectangle++(0.04,0.04)
  (-0.154,2.311)rectangle++(0.04,0.04) (-0.041,3.566)rectangle++(0.04,0.04) (0.106,4.629)rectangle++(0.04,0.04) (-1.204,2.632)rectangle++(0.04,0.04) (-0.92,3.881)rectangle++(0.04,0.04)
  (-0.629,4.932)rectangle++(0.04,0.04) (-0.34,5.814)rectangle++(0.04,0.04) (1.661,1.737)rectangle++(0.04,0.04) (1.479,3.006)rectangle++(0.04,0.04) (1.376,4.092)rectangle++(0.04,0.04)
  (-0.02,0.638)rectangle++(0.04,0.04) (0.006,2.139)rectangle++(0.04,0.04) (0.089,3.414)rectangle++(0.04,0.04) (-0.891,2.361)rectangle++(0.04,0.04) (-0.665,3.638)rectangle++(0.04,0.04)
  (-0.422,4.715)rectangle++(0.04,0.04) (1.366,1.452)rectangle++(0.04,0.04) (0.184,1.134)rectangle++(0.04,0.04) (0.2,2.552)rectangle++(0.04,0.04) (0.27,3.758)rectangle++(0.04,0.04)
  (0.377,4.78)rectangle++(0.04,0.04) (0.51,5.644)rectangle++(0.04,0.04) (-1.387,2.948)rectangle++(0.04,0.04) (-1.063,4.157)rectangle++(0.04,0.04) (-0.739,5.172)rectangle++(0.04,0.04)
  (-0.424,6.022)rectangle++(0.04,0.04) (-0.122,6.732)rectangle++(0.04,0.04) (0.164,7.323)rectangle++(0.04,0.04) (0.43,7.815)rectangle++(0.04,0.04) (0.676,8.221)rectangle++(0.04,0.04)
  (0.901,8.557)rectangle++(0.04,0.04) (1.106,8.834)rectangle++(0.04,0.04) (1.291,9.061)rectangle++(0.04,0.04) (1.458,9.246)rectangle++(0.04,0.04) (1.607,9.397)rectangle++(0.04,0.04)
  (1.739,9.52)rectangle++(0.04,0.04) (1.857,9.618)rectangle++(0.04,0.04) (1.961,9.697)rectangle++(0.04,0.04) (2.053,9.761)rectangle++(0.04,0.04) (2.133,9.811)rectangle++(0.04,0.04);
\fill[nuage]
  (2.203,9.85)rectangle++(0.04,0.04) (2.265,9.88)rectangle++(0.04,0.04) (2.318,9.904)rectangle++(0.04,0.04) (2.408,3.41)rectangle++(0.04,0.04) (2.181,4.398)rectangle++(0.04,0.04)
  (2.028,5.247)rectangle++(0.04,0.04) (1.931,5.975)rectangle++(0.04,0.04) (1.878,6.598)rectangle++(0.04,0.04) (1.858,7.129)rectangle++(0.04,0.04) (1.862,7.582)rectangle++(0.04,0.04)
  (-1.62,3.685)rectangle++(0.04,0.04) (-1.232,4.794)rectangle++(0.04,0.04) (-0.857,5.72)rectangle++(0.04,0.04) (-1.68,2.65)rectangle++(0.04,0.04) (-1.324,3.916)rectangle++(0.04,0.04)
  (-0.971,4.978)rectangle++(0.04,0.04) (-0.628,5.866)rectangle++(0.04,0.04) (-0.302,6.608)rectangle++(0.04,0.04) (-0.02,1.04)rectangle++(0.04,0.04) (0.022,2.481)rectangle++(0.04,0.04)
  (0.116,3.704)rectangle++(0.04,0.04) (0.245,4.74)rectangle++(0.04,0.04) (0.395,5.616)rectangle++(0.04,0.04) (0.559,6.354)rectangle++(0.04,0.04) (0.727,6.975)rectangle++(0.04,0.04)
  (0.894,7.495)rectangle++(0.04,0.04) (1.947,2.461)rectangle++(0.04,0.04) (-0.272,2.578)rectangle++(0.04,0.04) (-0.746,2.094)rectangle++(0.04,0.04) (-0.552,3.406)rectangle++(0.04,0.04)
  (1.019,1.104)rectangle++(0.04,0.04) (0.908,2.494)rectangle++(0.04,0.04) (0.869,3.679)rectangle++(0.04,0.04) (0.884,4.689)rectangle++(0.04,0.04) (0.937,5.546)rectangle++(0.04,0.04)
  (1.016,6.273)rectangle++(0.04,0.04) (1.587,2.2)rectangle++(0.04,0.04) (1.434,3.402)rectangle++(0.04,0.04) (1.353,4.431)rectangle++(0.04,0.04) (1.325,5.308)rectangle++(0.04,0.04)
  (1.337,6.055)rectangle++(0.04,0.04) (1.376,6.69)rectangle++(0.04,0.04) (1.435,7.227)rectangle++(0.04,0.04) (1.507,7.682)rectangle++(0.04,0.04) (-1.717,3.626)rectangle++(0.04,0.04)
  (-1.317,4.747)rectangle++(0.04,0.04) (-0.932,5.684)rectangle++(0.04,0.04) (-0.567,6.464)rectangle++(0.04,0.04) (-0.225,7.114)rectangle++(0.04,0.04) (0.091,7.652)rectangle++(0.04,0.04)
  (0.381,8.097)rectangle++(0.04,0.04) (0.646,8.463)rectangle++(0.04,0.04) (2.255,2.629)rectangle++(0.04,0.04) (2.02,3.741)rectangle++(0.04,0.04) (-0.59,2.877)rectangle++(0.04,0.04)
  (-0.388,4.065)rectangle++(0.04,0.04) (-0.17,5.067)rectangle++(0.04,0.04) (-1.373,2.665)rectangle++(0.04,0.04) (-0.989,1.86)rectangle++(0.04,0.04) (-0.768,3.216)rectangle++(0.04,0.04)
  (-0.526,4.361)rectangle++(0.04,0.04) (-0.275,5.324)rectangle++(0.04,0.04) (-0.023,6.133)rectangle++(0.04,0.04) (0.223,6.81)rectangle++(0.04,0.04) (0.46,7.376)rectangle++(0.04,0.04)
  (0.684,7.847)rectangle++(0.04,0.04) (0.893,8.239)rectangle++(0.04,0.04) (1.086,8.563)rectangle++(0.04,0.04) (1.264,8.832)rectangle++(0.04,0.04) (1.425,9.053)rectangle++(0.04,0.04)
  (1.571,9.234)rectangle++(0.04,0.04) (1.703,9.382)rectangle++(0.04,0.04) (1.821,9.503)rectangle++(0.04,0.04) (1.925,9.601)rectangle++(0.04,0.04) (-2.132,4.144)rectangle++(0.04,0.04)
  (-1.649,5.204)rectangle++(0.04,0.04) (1.687,1.25)rectangle++(0.04,0.04) (0.08,1.169)rectangle++(0.04,0.04) (0.112,2.586)rectangle++(0.04,0.04) (0.197,3.79)rectangle++(0.04,0.04)
  (0.317,4.81)rectangle++(0.04,0.04) (0.459,5.672)rectangle++(0.04,0.04) (0.615,6.399)rectangle++(0.04,0.04) (0.777,7.011)rectangle++(0.04,0.04) (0.938,7.524)rectangle++(0.04,0.04)
  (1.096,7.954)rectangle++(0.04,0.04) (2.045,2.624)rectangle++(0.04,0.04) (1.841,3.745)rectangle++(0.04,0.04) (1.713,4.706)rectangle++(0.04,0.04) (1.642,5.528)rectangle++(0.04,0.04)
  (1.614,6.229)rectangle++(0.04,0.04) (1.619,6.826)rectangle++(0.04,0.04) (1.651,2.489)rectangle++(0.04,0.04) (1.501,3.646)rectangle++(0.04,0.04) (1.419,4.635)rectangle++(0.04,0.04)
  (1.068,1.911)rectangle++(0.04,0.04) (0.982,3.178)rectangle++(0.04,0.04) (0.959,4.258)rectangle++(0.04,0.04) (0.984,5.178)rectangle++(0.04,0.04) (1.041,5.958)rectangle++(0.04,0.04)
  (1.121,6.619)rectangle++(0.04,0.04) (-1.518,3.303)rectangle++(0.04,0.04) (-1.16,4.464)rectangle++(0.04,0.04) (-1.414,2.304)rectangle++(0.04,0.04) (-1.112,3.611)rectangle++(0.04,0.04)
  (-0.803,4.71)rectangle++(0.04,0.04) (1.422,1.352)rectangle++(0.04,0.04) (1.261,2.688)rectangle++(0.04,0.04) (-0.468,2.47)rectangle++(0.04,0.04) (-0.301,3.715)rectangle++(0.04,0.04)
  (-0.11,4.766)rectangle++(0.04,0.04) (0.095,5.651)rectangle++(0.04,0.04) (0.305,6.396)rectangle++(0.04,0.04) (0.513,7.021)rectangle++(0.04,0.04) (0.714,7.543)rectangle++(0.04,0.04)
  (0.907,7.979)rectangle++(0.04,0.04) (1.088,8.342)rectangle++(0.04,0.04) (1.256,8.644)rectangle++(0.04,0.04) (1.411,8.893)rectangle++(0.04,0.04) (1.553,9.099)rectangle++(0.04,0.04)
  (1.682,9.268)rectangle++(0.04,0.04) (1.798,9.407)rectangle++(0.04,0.04) (1.902,9.52)rectangle++(0.04,0.04) (1.996,9.612)rectangle++(0.04,0.04) (2.079,9.687)rectangle++(0.04,0.04)
  (2.152,9.747)rectangle++(0.04,0.04) (2.217,9.795)rectangle++(0.04,0.04) (2.274,9.833)rectangle++(0.04,0.04) (2.324,9.863)rectangle++(0.04,0.04) (2.368,9.887)rectangle++(0.04,0.04)
  (2.406,9.906)rectangle++(0.04,0.04) (2.439,9.92)rectangle++(0.04,0.04) (2.468,9.93)rectangle++(0.04,0.04) (2.493,9.938)rectangle++(0.04,0.04) (2.514,9.944)rectangle++(0.04,0.04)
  (2.39,3.47)rectangle++(0.04,0.04) (2.168,4.45)rectangle++(0.04,0.04) (2.019,5.292)rectangle++(0.04,0.04) (1.925,6.014)rectangle++(0.04,0.04) (1.875,6.631)rectangle++(0.04,0.04)
  (1.857,7.158)rectangle++(0.04,0.04) (1.862,7.606)rectangle++(0.04,0.04) (-1.626,3.691)rectangle++(0.04,0.04) (-1.237,4.798)rectangle++(0.04,0.04) (-0.862,5.724)rectangle++(0.04,0.04)
  (-0.506,6.496)rectangle++(0.04,0.04) (-0.172,7.138)rectangle++(0.04,0.04) (-0.02,1.125)rectangle++(0.04,0.04) (0.026,2.554)rectangle++(0.04,0.04) (0.122,3.766)rectangle++(0.04,0.04)
  (0.252,4.792)rectangle++(0.04,0.04) (0.404,5.659)rectangle++(0.04,0.04) (-1.412,2.927)rectangle++(0.04,0.04) (1.014,0.765)rectangle++(0.04,0.04) (0.89,2.206)rectangle++(0.04,0.04)
  (0.843,3.436)rectangle++(0.04,0.04) (0.852,4.483)rectangle++(0.04,0.04) (0.901,5.373)rectangle++(0.04,0.04) (0.979,6.127)rectangle++(0.04,0.04) (1.075,6.765)rectangle++(0.04,0.04)
  (1.716,2.333)rectangle++(0.04,0.04) (1.549,3.511)rectangle++(0.04,0.04) (1.455,4.518)rectangle++(0.04,0.04) (1.415,5.379)rectangle++(0.04,0.04) (1.416,6.111)rectangle++(0.04,0.04)
  (1.446,6.734)rectangle++(0.04,0.04) (1.496,7.262)rectangle++(0.04,0.04) (1.792,2.562)rectangle++(0.04,0.04) (1.623,3.702)rectangle++(0.04,0.04) (1.526,4.678)rectangle++(0.04,0.04)
  (1.482,5.512)rectangle++(0.04,0.04) (1.478,6.222)rectangle++(0.04,0.04) (1.503,6.826)rectangle++(0.04,0.04) (1.548,7.338)rectangle++(0.04,0.04) (1.607,7.771)rectangle++(0.04,0.04)
  (1.675,8.138)rectangle++(0.04,0.04) (1.747,8.446)rectangle++(0.04,0.04) (1.821,8.706)rectangle++(0.04,0.04) (1.893,8.923)rectangle++(0.04,0.04) (1.964,9.105)rectangle++(0.04,0.04)
  (2.237,3.126)rectangle++(0.04,0.04) (2.025,4.164)rectangle++(0.04,0.04) (1.885,5.054)rectangle++(0.04,0.04) (1.115,2.133)rectangle++(0.04,0.04) (1.031,3.365)rectangle++(0.04,0.04)
  (1.009,4.415)rectangle++(0.04,0.04) (1.032,5.309)rectangle++(0.04,0.04) (1.087,6.067)rectangle++(0.04,0.04) (1.165,6.71)rectangle++(0.04,0.04) (1.256,7.253)rectangle++(0.04,0.04)
  (1.825,2.497)rectangle++(0.04,0.04) (1.649,3.646)rectangle++(0.04,0.04) (1.545,4.629)rectangle++(0.04,0.04) (1.496,5.469)rectangle++(0.04,0.04) (1.489,6.185)rectangle++(0.04,0.04)
  (1.51,6.794)rectangle++(0.04,0.04) (1.553,7.311)rectangle++(0.04,0.04) (-1.611,3.555)rectangle++(0.04,0.04) (-1.23,4.682)rectangle++(0.04,0.04) (-0.86,5.625)rectangle++(0.04,0.04)
  (-0.508,6.412)rectangle++(0.04,0.04) (-0.178,7.067)rectangle++(0.04,0.04) (0.129,7.61)rectangle++(0.04,0.04) (0.412,8.06)rectangle++(0.04,0.04) (0.671,8.43)rectangle++(0.04,0.04)
  (2.243,2.628)rectangle++(0.04,0.04) (2.009,3.74)rectangle++(0.04,0.04) (1.855,4.695)rectangle++(0.04,0.04) (1.762,5.513)rectangle++(0.04,0.04) (1.262,2.211)rectangle++(0.04,0.04)
  (1.159,3.425)rectangle++(0.04,0.04) (-0.68,2.609)rectangle++(0.04,0.04) (-0.836,2.007)rectangle++(0.04,0.04) (-0.632,3.335)rectangle++(0.04,0.04) (-0.406,4.456)rectangle++(0.04,0.04)
  (-0.169,5.4)rectangle++(0.04,0.04) (1.52,1.682)rectangle++(0.04,0.04) (1.357,2.965)rectangle++(0.04,0.04) (1.27,4.062)rectangle++(0.04,0.04) (1.24,4.998)rectangle++(0.04,0.04)
  (1.252,5.795)rectangle++(0.04,0.04) (1.293,6.472)rectangle++(0.04,0.04) (1.356,7.046)rectangle++(0.04,0.04) (1.432,7.531)rectangle++(0.04,0.04) (1.517,7.94)rectangle++(0.04,0.04)
  (1.604,8.285)rectangle++(0.04,0.04) (1.693,8.574)rectangle++(0.04,0.04) (1.78,8.816)rectangle++(0.04,0.04) (1.863,9.019)rectangle++(0.04,0.04) (1.942,9.188)rectangle++(0.04,0.04)
  (2.264,3.14)rectangle++(0.04,0.04) (2.048,4.175)rectangle++(0.04,0.04) (-0.697,2.978)rectangle++(0.04,0.04) (-0.476,4.156)rectangle++(0.04,0.04) (-0.24,5.148)rectangle++(0.04,0.04)
  (0,5.981)rectangle++(0.04,0.04) (0.237,6.68)rectangle++(0.04,0.04) (-1.711,3.113)rectangle++(0.04,0.04) (-1.332,4.311)rectangle++(0.04,0.04) (-0.962,5.314)rectangle++(0.04,0.04)
  (-0.607,6.151)rectangle++(0.04,0.04) (-0.272,6.849)rectangle++(0.04,0.04) (1.941,2.003)rectangle++(0.04,0.04) (1.728,3.221)rectangle++(0.04,0.04) (1.595,4.265)rectangle++(0.04,0.04)
  (1.524,5.158)rectangle++(0.04,0.04) (1.5,5.919)rectangle++(0.04,0.04) (1.509,6.568)rectangle++(0.04,0.04) (1.595,2.399)rectangle++(0.04,0.04) (1.45,3.571)rectangle++(0.04,0.04)
  (1.373,4.574)rectangle++(0.04,0.04) (1.348,5.429)rectangle++(0.04,0.04) (1.36,6.157)rectangle++(0.04,0.04) (-1.35,3.256)rectangle++(0.04,0.04) (-1.019,4.418)rectangle++(0.04,0.04)
  (-0.692,5.392)rectangle++(0.04,0.04) (-0.375,6.207)rectangle++(0.04,0.04) (-0.072,6.888)rectangle++(0.04,0.04) (0.212,7.453)rectangle++(0.04,0.04) (0.476,7.923)rectangle++(0.04,0.04)
  (0.719,8.312)rectangle++(0.04,0.04) (0.942,8.633)rectangle++(0.04,0.04) (1.144,8.896)rectangle++(0.04,0.04) (1.326,9.112)rectangle++(0.04,0.04) (1.489,9.289)rectangle++(0.04,0.04)
  (1.635,9.432)rectangle++(0.04,0.04) (1.765,9.548)rectangle++(0.04,0.04) (1.88,9.641)rectangle++(0.04,0.04) (1.981,9.716)rectangle++(0.04,0.04) (2.071,9.776)rectangle++(0.04,0.04)
  (2.149,9.823)rectangle++(0.04,0.04) (2.217,9.86)rectangle++(0.04,0.04) (2.277,9.888)rectangle++(0.04,0.04) (2.329,9.91)rectangle++(0.04,0.04) (2.374,9.927)rectangle++(0.04,0.04)
  (2.412,9.939)rectangle++(0.04,0.04) (2.446,9.948)rectangle++(0.04,0.04) (2.475,9.954)rectangle++(0.04,0.04) (2.499,9.958)rectangle++(0.04,0.04) (2.521,9.961)rectangle++(0.04,0.04)
  (2.539,9.962)rectangle++(0.04,0.04) (2.554,9.962)rectangle++(0.04,0.04) (2.567,9.962)rectangle++(0.04,0.04) (2.579,9.961)rectangle++(0.04,0.04) (2.588,9.96)rectangle++(0.04,0.04)
  (2.596,9.959)rectangle++(0.04,0.04) (2.603,9.957)rectangle++(0.04,0.04) (2.608,9.956)rectangle++(0.04,0.04) (2.613,9.954)rectangle++(0.04,0.04) (2.617,9.953)rectangle++(0.04,0.04)
  (2.621,9.951)rectangle++(0.04,0.04) (-0.02,1.575)rectangle++(0.04,0.04) (0.044,2.936)rectangle++(0.04,0.04) (0.152,4.09)rectangle++(0.04,0.04) (0.291,5.067)rectangle++(0.04,0.04)
  (0.448,5.891)rectangle++(0.04,0.04) (0.614,6.586)rectangle++(0.04,0.04) (0.783,7.17)rectangle++(0.04,0.04) (-1.729,3.346)rectangle++(0.04,0.04) (-1.338,4.51)rectangle++(0.04,0.04)
  (-0.959,5.483)rectangle++(0.04,0.04) (-0.598,6.295)rectangle++(0.04,0.04) (-1.778,2.836)rectangle++(0.04,0.04) (-1.4,4.078)rectangle++(0.04,0.04) (1.334,1.045)rectangle++(0.04,0.04)
  (1.174,2.431)rectangle++(0.04,0.04) (1.093,3.616)rectangle++(0.04,0.04) (1.071,4.626)rectangle++(0.04,0.04) (1.093,5.485)rectangle++(0.04,0.04) (1.147,6.215)rectangle++(0.04,0.04)
  (-1.408,3.22)rectangle++(0.04,0.04) (-1.07,4.39)rectangle++(0.04,0.04) (-0.736,5.37)rectangle++(0.04,0.04) (-0.413,6.19)rectangle++(0.04,0.04) (-0.106,6.874)rectangle++(0.04,0.04)
  (1.923,2.052)rectangle++(0.04,0.04) (-0.17,2.483)rectangle++(0.04,0.04) (-0.048,3.713)rectangle++(0.04,0.04) (1.029,1.309)rectangle++(0.04,0.04) (0.925,2.668)rectangle++(0.04,0.04)
  (0.891,3.827)rectangle++(0.04,0.04) (0.908,4.813)rectangle++(0.04,0.04) (0.962,5.651)rectangle++(0.04,0.04) (-0.02,0.887)rectangle++(0.04,0.04) (0.016,2.351)rectangle++(0.04,0.04)
  (0.106,3.594)rectangle++(0.04,0.04) (0.231,4.647)rectangle++(0.04,0.04) (0.38,5.537)rectangle++(0.04,0.04) (0.543,6.287)rectangle++(0.04,0.04) (0.711,6.919)rectangle++(0.04,0.04)
  (0.878,7.449)rectangle++(0.04,0.04) (1.042,7.892)rectangle++(0.04,0.04) (1.2,8.263)rectangle++(0.04,0.04) (1.348,8.572)rectangle++(0.04,0.04) (1.486,8.828)rectangle++(0.04,0.04)
  (1.614,9.041)rectangle++(0.04,0.04) (1.732,9.216)rectangle++(0.04,0.04) (1.838,9.361)rectangle++(0.04,0.04) (1.935,9.479)rectangle++(0.04,0.04) (2.022,9.576)rectangle++(0.04,0.04)
  (2.099,9.655)rectangle++(0.04,0.04) (2.371,3.293)rectangle++(0.04,0.04) (2.145,4.3)rectangle++(0.04,0.04) (1.993,5.166)rectangle++(0.04,0.04) (1.899,5.907)rectangle++(0.04,0.04)
  (1.848,6.542)rectangle++(0.04,0.04) (1.83,7.083)rectangle++(0.04,0.04) (1.837,7.543)rectangle++(0.04,0.04) (1.861,7.934)rectangle++(0.04,0.04) (1.897,8.266)rectangle++(0.04,0.04)
  (1.941,8.546)rectangle++(0.04,0.04) (1.989,8.783)rectangle++(0.04,0.04) (2.04,8.982)rectangle++(0.04,0.04) (2.091,9.15)rectangle++(0.04,0.04) (2.141,9.29)rectangle++(0.04,0.04)
  (2.189,9.407)rectangle++(0.04,0.04) (2.288,3.257)rectangle++(0.04,0.04) (0.551,1.807)rectangle++(0.04,0.04) (0.539,3.11)rectangle++(0.04,0.04) (0.58,4.218)rectangle++(0.04,0.04)
  (0.66,5.158)rectangle++(0.04,0.04) (0.765,5.954)rectangle++(0.04,0.04) (1.535,2.058)rectangle++(0.04,0.04) (1.385,3.284)rectangle++(0.04,0.04) (1.306,4.332)rectangle++(0.04,0.04)
  (1.281,5.226)rectangle++(0.04,0.04) (1.254,2.018)rectangle++(0.04,0.04) (1.144,3.261)rectangle++(0.04,0.04) (1.101,4.322)rectangle++(0.04,0.04) (1.106,5.226)rectangle++(0.04,0.04)
  (1.147,5.994)rectangle++(0.04,0.04) (1.213,6.645)rectangle++(0.04,0.04) (1.294,7.196)rectangle++(0.04,0.04) (1.386,7.661)rectangle++(0.04,0.04) (1.482,8.053)rectangle++(0.04,0.04)
  (-1.818,3.702)rectangle++(0.04,0.04) (-1.4,4.815)rectangle++(0.04,0.04) (-0.999,5.745)rectangle++(0.04,0.04) (-0.622,6.52)rectangle++(0.04,0.04) (-0.27,7.163)rectangle++(0.04,0.04)
  (0.055,7.695)rectangle++(0.04,0.04) (0.352,8.135)rectangle++(0.04,0.04) (0.623,8.497)rectangle++(0.04,0.04) (0.867,8.794)rectangle++(0.04,0.04) (1.086,9.036)rectangle++(0.04,0.04)
  (1.283,9.233)rectangle++(0.04,0.04) (2.376,2.98)rectangle++(0.04,0.04) (2.136,4.034)rectangle++(0.04,0.04) (1.975,4.939)rectangle++(0.04,0.04) (1.874,5.716)rectangle++(0.04,0.04)
  (1.302,2.289)rectangle++(0.04,0.04) (0.428,1.318)rectangle++(0.04,0.04) (0.415,2.699)rectangle++(0.04,0.04) (0.458,3.874)rectangle++(0.04,0.04) (-0.02,0.603)rectangle++(0.04,0.04)
  (-0.182,1.717)rectangle++(0.04,0.04) (-0.088,3.063)rectangle++(0.04,0.04) (0.045,4.203)rectangle++(0.04,0.04) (0.204,5.167)rectangle++(0.04,0.04) (0.378,5.98)rectangle++(0.04,0.04)
  (0.559,6.664)rectangle++(0.04,0.04) (0.739,7.238)rectangle++(0.04,0.04) (0.916,7.719)rectangle++(0.04,0.04) (1.085,8.121)rectangle++(0.04,0.04) (1.245,8.456)rectangle++(0.04,0.04)
  (1.394,8.734)rectangle++(0.04,0.04) (1.532,8.964)rectangle++(0.04,0.04) (-2.045,3.913)rectangle++(0.04,0.04) (-1.584,5.004)rectangle++(0.04,0.04) (-1.149,5.913)rectangle++(0.04,0.04)
  (-0.742,6.669)rectangle++(0.04,0.04) (-0.366,7.294)rectangle++(0.04,0.04) (2.08,2.085)rectangle++(0.04,0.04) (1.849,3.286)rectangle++(0.04,0.04) (1.701,4.315)rectangle++(0.04,0.04)
  (-0.803,2.93)rectangle++(0.04,0.04) (-0.567,4.118)rectangle++(0.04,0.04) (-0.32,5.12)rectangle++(0.04,0.04) (-0.069,5.961)rectangle++(0.04,0.04) (0.177,6.665)rectangle++(0.04,0.04)
  (0.415,7.255)rectangle++(0.04,0.04) (1.952,2.279)rectangle++(0.04,0.04) (-0.223,2.539)rectangle++(0.04,0.04) (-0.091,3.764)rectangle++(0.04,0.04) (0.071,4.799)rectangle++(0.04,0.04)
  (0.25,5.672)rectangle++(0.04,0.04) (1.533,1.856)rectangle++(0.04,0.04) (1.375,3.113)rectangle++(0.04,0.04) (1.291,4.187)rectangle++(0.04,0.04) (1.263,5.104)rectangle++(0.04,0.04)
  (1.275,5.884)rectangle++(0.04,0.04) (1.317,6.546)rectangle++(0.04,0.04) (1.379,7.108)rectangle++(0.04,0.04) (1.455,7.583)rectangle++(0.04,0.04) (1.537,7.983)rectangle++(0.04,0.04)
  (1.987,2.746)rectangle++(0.04,0.04) (1.797,3.851)rectangle++(0.04,0.04) (1.679,4.797)rectangle++(0.04,0.04) (1.617,5.607)rectangle++(0.04,0.04) (1.596,6.297)rectangle++(0.04,0.04)
  (1.607,6.885)rectangle++(0.04,0.04) (1.639,7.384)rectangle++(0.04,0.04) (1.686,7.807)rectangle++(0.04,0.04) (1.743,8.165)rectangle++(0.04,0.04) (1.806,8.467)rectangle++(0.04,0.04)
  (1.872,8.721)rectangle++(0.04,0.04) (1.938,8.934)rectangle++(0.04,0.04) (2.193,3.078)rectangle++(0.04,0.04) (1.985,4.125)rectangle++(0.04,0.04) (0.84,1.936)rectangle++(0.04,0.04)
  (0.789,3.208)rectangle++(0.04,0.04) (0.797,4.292)rectangle++(0.04,0.04) (0.847,5.212)rectangle++(0.04,0.04) (0.926,5.993)rectangle++(0.04,0.04) (1.025,6.653)rectangle++(0.04,0.04)
  (1.135,7.21)rectangle++(0.04,0.04) (-1.669,3.436)rectangle++(0.04,0.04) (-1.283,4.584)rectangle++(0.04,0.04) (-0.91,5.544)rectangle++(0.04,0.04) (1.671,1.524)rectangle++(0.04,0.04)
  (1.479,2.825)rectangle++(0.04,0.04) (-0.46,2.551)rectangle++(0.04,0.04) (-0.291,3.783)rectangle++(0.04,0.04) (-1.063,2.354)rectangle++(0.04,0.04) (-0.811,3.64)rectangle++(0.04,0.04)
  (-0.546,4.723)rectangle++(0.04,0.04) (-0.278,5.632)rectangle++(0.04,0.04) (-0.013,6.395)rectangle++(0.04,0.04) (0.243,7.032)rectangle++(0.04,0.04) (0.485,7.564)rectangle++(0.04,0.04)
  (0.713,8.006)rectangle++(0.04,0.04) (0.924,8.373)rectangle++(0.04,0.04) (1.118,8.676)rectangle++(0.04,0.04) (-2.053,3.755)rectangle++(0.04,0.04) (1.342,0.797)rectangle++(0.04,0.04)
  (0.005,0.97)rectangle++(0.04,0.04) (0.04,2.421)rectangle++(0.04,0.04) (0.129,3.652)rectangle++(0.04,0.04) (0.254,4.695)rectangle++(0.04,0.04) (1.259,1.623)rectangle++(0.04,0.04)
  (1.133,2.925)rectangle++(0.04,0.04) (1.078,4.037)rectangle++(0.04,0.04) (1.076,4.985)rectangle++(0.04,0.04) (1.111,5.79)rectangle++(0.04,0.04) (-1.304,3.118)rectangle++(0.04,0.04)
  (-0.986,4.299)rectangle++(0.04,0.04) (1.334,1.205)rectangle++(0.04,0.04) (1.18,2.567)rectangle++(0.04,0.04) (1.104,3.731)rectangle++(0.04,0.04) (1.085,4.724)rectangle++(0.04,0.04)
  (1.109,5.568)rectangle++(0.04,0.04) (1.163,6.285)rectangle++(0.04,0.04) (1.238,6.892)rectangle++(0.04,0.04) (1.326,7.404)rectangle++(0.04,0.04) (1.421,7.837)rectangle++(0.04,0.04)
  (1.519,8.201)rectangle++(0.04,0.04) (1.617,8.506)rectangle++(0.04,0.04) (1.712,8.762)rectangle++(0.04,0.04) (1.804,8.975)rectangle++(0.04,0.04) (1.89,9.153)rectangle++(0.04,0.04)
  (1.971,9.301)rectangle++(0.04,0.04) (2.045,9.423)rectangle++(0.04,0.04) (2.113,9.524)rectangle++(0.04,0.04) (2.175,9.607)rectangle++(0.04,0.04) (2.23,9.675)rectangle++(0.04,0.04)
  (2.281,9.731)rectangle++(0.04,0.04) (-2.095,4.254)rectangle++(0.04,0.04) (-1.613,5.296)rectangle++(0.04,0.04) (-1.161,6.162)rectangle++(0.04,0.04) (-0.743,6.881)rectangle++(0.04,0.04)
  (-0.358,7.475)rectangle++(0.04,0.04) (-0.008,7.964)rectangle++(0.04,0.04) (2.214,2.339)rectangle++(0.04,0.04) (-0.02,0.357)rectangle++(0.04,0.04) (-0.118,1.663)rectangle++(0.04,0.04)
  (-0.036,3.015)rectangle++(0.04,0.04) (0.088,4.16)rectangle++(0.04,0.04) (0.239,5.129)rectangle++(0.04,0.04) (0.406,5.946)rectangle++(0.04,0.04) (0.581,6.634)rectangle++(0.04,0.04)
  (1.753,2.173)rectangle++(0.04,0.04) (0.328,1.407)rectangle++(0.04,0.04) (0.333,2.779)rectangle++(0.04,0.04) (0.392,3.945)rectangle++(0.04,0.04) (0.489,4.934)rectangle++(0.04,0.04)
  (0.611,5.771)rectangle++(0.04,0.04) (0.748,6.477)rectangle++(0.04,0.04) (0.892,7.072)rectangle++(0.04,0.04) (1.039,7.571)rectangle++(0.04,0.04) (1.184,7.99)rectangle++(0.04,0.04)
  (1.324,8.341)rectangle++(0.04,0.04) (1.457,8.633)rectangle++(0.04,0.04) (1.581,8.876)rectangle++(0.04,0.04) (1.697,9.077)rectangle++(0.04,0.04) (1.803,9.244)rectangle++(0.04,0.04)
  (1.9,9.382)rectangle++(0.04,0.04) (-2.08,4.09)rectangle++(0.04,0.04) (-1.607,5.156)rectangle++(0.04,0.04) (-1.162,6.043)rectangle++(0.04,0.04) (-0.748,6.779)rectangle++(0.04,0.04)
  (-0.367,7.388)rectangle++(0.04,0.04) (-0.018,7.891)rectangle++(0.04,0.04) (0.298,8.304)rectangle++(0.04,0.04) (0.583,8.643)rectangle++(0.04,0.04) (0.839,8.919)rectangle++(0.04,0.04)
  (1.068,9.144)rectangle++(0.04,0.04) (1.271,9.326)rectangle++(0.04,0.04) (1.451,9.472)rectangle++(0.04,0.04) (1.61,9.59)rectangle++(0.04,0.04) (2.426,3.15)rectangle++(0.04,0.04)
  (2.186,4.177)rectangle++(0.04,0.04) (2.023,5.059)rectangle++(0.04,0.04) (1.92,5.815)rectangle++(0.04,0.04) (1.862,6.463)rectangle++(0.04,0.04) (1.513,2.465)rectangle++(0.04,0.04)
  (1.382,3.631)rectangle++(0.04,0.04) (1.318,4.627)rectangle++(0.04,0.04) (1.303,5.477)rectangle++(0.04,0.04) (1.325,6.199)rectangle++(0.04,0.04) (1.372,6.813)rectangle++(0.04,0.04)
  (1.436,7.332)rectangle++(0.04,0.04) (1.512,7.771)rectangle++(0.04,0.04) (1.594,8.141)rectangle++(0.04,0.04) (2.023,2.798)rectangle++(0.04,0.04) (1.829,3.894)rectangle++(0.04,0.04)
  (1.708,4.833)rectangle++(0.04,0.04) (1.643,5.636)rectangle++(0.04,0.04) (1.62,6.321)rectangle++(0.04,0.04) (1.628,6.904)rectangle++(0.04,0.04) (1.657,7.4)rectangle++(0.04,0.04)
  (1.703,7.82)rectangle++(0.04,0.04) (1.758,8.175)rectangle++(0.04,0.04) (1.819,8.474)rectangle++(0.04,0.04) (-1.861,3.872)rectangle++(0.04,0.04) (-1.429,4.962)rectangle++(0.04,0.04)
  (-1.018,5.871)rectangle++(0.04,0.04) (-0.633,6.627)rectangle++(0.04,0.04) (-0.275,7.255)rectangle++(0.04,0.04) (0.054,7.774)rectangle++(0.04,0.04) (0.355,8.202)rectangle++(0.04,0.04)
  (0.627,8.553)rectangle++(0.04,0.04) (2.283,2.646)rectangle++(0.04,0.04) (2.045,3.754)rectangle++(0.04,0.04) (1.886,4.705)rectangle++(0.04,0.04) (1.789,5.52)rectangle++(0.04,0.04)
  (1.739,6.217)rectangle++(0.04,0.04) (1.725,6.811)rectangle++(0.04,0.04) (1.631,2.513)rectangle++(0.04,0.04) (1.485,3.667)rectangle++(0.04,0.04) (1.406,4.654)rectangle++(0.04,0.04)
  (1.379,5.496)rectangle++(0.04,0.04) (1.39,6.212)rectangle++(0.04,0.04) (1.514,2.282)rectangle++(0.04,0.04) (0.395,1.371)rectangle++(0.04,0.04) (0.388,2.746)rectangle++(0.04,0.04)
  (0.438,3.915)rectangle++(0.04,0.04) (-0.952,2.551)rectangle++(0.04,0.04) (-0.709,3.803)rectangle++(0.04,0.04) (-1.152,2.262)rectangle++(0.04,0.04) (-0.891,3.565)rectangle++(0.04,0.04)
  (-1.126,2.169)rectangle++(0.04,0.04) (-0.873,3.485)rectangle++(0.04,0.04) (-0.605,4.593)rectangle++(0.04,0.04) (-0.332,5.524)rectangle++(0.04,0.04) (-0.064,6.305)rectangle++(0.04,0.04)
  (0.196,6.958)rectangle++(0.04,0.04) (0.443,7.503)rectangle++(0.04,0.04) (0.674,7.956)rectangle++(0.04,0.04) (0.889,8.332)rectangle++(0.04,0.04) (1.087,8.643)rectangle++(0.04,0.04)
  (1.267,8.899)rectangle++(0.04,0.04) (1.431,9.11)rectangle++(0.04,0.04) (1.578,9.282)rectangle++(0.04,0.04) (1.711,9.423)rectangle++(0.04,0.04) (1.829,9.537)rectangle++(0.04,0.04)
  (1.934,9.63)rectangle++(0.04,0.04) (2.027,9.704)rectangle++(0.04,0.04) (2.109,9.764)rectangle++(0.04,0.04) (2.181,9.811)rectangle++(0.04,0.04) (2.244,9.848)rectangle++(0.04,0.04)
  (2.299,9.877)rectangle++(0.04,0.04) (2.403,3.398)rectangle++(0.04,0.04) (2.177,4.389)rectangle++(0.04,0.04) (0.885,2.049)rectangle++(0.04,0.04) (0.832,3.303)rectangle++(0.04,0.04)
  (0.837,4.37)rectangle++(0.04,0.04) (0.884,5.277)rectangle++(0.04,0.04) (1.328,1.926)rectangle++(0.04,0.04) (1.203,3.181)rectangle++(0.04,0.04) (-0.607,2.565)rectangle++(0.04,0.04)
  (-0.416,3.801)rectangle++(0.04,0.04) (1.109,1.234)rectangle++(0.04,0.04) (0.99,2.601)rectangle++(0.04,0.04) (0.943,3.767)rectangle++(0.04,0.04) (-0.02,0.586)rectangle++(0.04,0.04)
  (0.004,2.095)rectangle++(0.04,0.04) (0.085,3.377)rectangle++(0.04,0.04) (0.205,4.463)rectangle++(0.04,0.04) (-0.02,0.697)rectangle++(0.04,0.04) (0.009,2.19)rectangle++(0.04,0.04)
  (0.093,3.457)rectangle++(0.04,0.04) (0.215,4.531)rectangle++(0.04,0.04) (0.362,5.439)rectangle++(0.04,0.04) (0.523,6.205)rectangle++(0.04,0.04) (0.69,6.849)rectangle++(0.04,0.04)
  (0.859,7.391)rectangle++(0.04,0.04) (-1.771,3.412)rectangle++(0.04,0.04) (-1.371,4.568)rectangle++(0.04,0.04) (-0.985,5.533)rectangle++(0.04,0.04) (-0.618,6.339)rectangle++(0.04,0.04)
  (-0.02,0.997)rectangle++(0.04,0.04) (0.021,2.445)rectangle++(0.04,0.04) (0.113,3.673)rectangle++(0.04,0.04) (0.241,4.714)rectangle++(0.04,0.04) (0.391,5.594)rectangle++(0.04,0.04)
  (0.554,6.335)rectangle++(0.04,0.04) (0.722,6.959)rectangle++(0.04,0.04) (0.89,7.482)rectangle++(0.04,0.04) (1.054,7.921)rectangle++(0.04,0.04) (1.21,8.287)rectangle++(0.04,0.04)
  (1.358,8.591)rectangle++(0.04,0.04) (1.496,8.845)rectangle++(0.04,0.04) (1.623,9.054)rectangle++(0.04,0.04) (2.274,3.025)rectangle++(0.04,0.04) (0.488,1.747)rectangle++(0.04,0.04)
  (0.483,3.062)rectangle++(0.04,0.04) (0.531,4.179)rectangle++(0.04,0.04) (0.616,5.128)rectangle++(0.04,0.04) (0.727,5.93)rectangle++(0.04,0.04) (0.853,6.608)rectangle++(0.04,0.04)
  (0.987,7.179)rectangle++(0.04,0.04) (1.124,7.659)rectangle++(0.04,0.04) (1.259,8.061)rectangle++(0.04,0.04) (2.051,2.692)rectangle++(0.04,0.04) (1.849,3.802)rectangle++(0.04,0.04)
  (1.721,4.754)rectangle++(0.04,0.04) (1.651,5.569)rectangle++(0.04,0.04) (1.624,6.263)rectangle++(0.04,0.04) (1.629,6.855)rectangle++(0.04,0.04) (1.656,7.358)rectangle++(0.04,0.04)
  (-1.603,3.589)rectangle++(0.04,0.04) (-1.221,4.711)rectangle++(0.04,0.04) (-0.852,5.649)rectangle++(0.04,0.04) (-0.5,6.432)rectangle++(0.04,0.04) (-0.17,7.083)rectangle++(0.04,0.04)
  (0.137,7.624)rectangle++(0.04,0.04) (0.419,8.071)rectangle++(0.04,0.04) (0.677,8.44)rectangle++(0.04,0.04) (0.911,8.743)rectangle++(0.04,0.04) (1.121,8.991)rectangle++(0.04,0.04)
  (1.311,9.194)rectangle++(0.04,0.04) (1.48,9.359)rectangle++(0.04,0.04) (1.63,9.492)rectangle++(0.04,0.04) (1.763,9.599)rectangle++(0.04,0.04) (1.88,9.685)rectangle++(0.04,0.04)
  (-2.163,4.152)rectangle++(0.04,0.04) (-1.675,5.212)rectangle++(0.04,0.04) (-1.217,6.093)rectangle++(0.04,0.04) (-0.793,6.824)rectangle++(0.04,0.04) (-0.403,7.429)rectangle++(0.04,0.04)
  (-0.048,7.927)rectangle++(0.04,0.04) (-2.092,3.322)rectangle++(0.04,0.04) (-1.647,4.503)rectangle++(0.04,0.04) (-1.222,5.49)rectangle++(0.04,0.04) (-0.822,6.312)rectangle++(0.04,0.04)
  (-0.448,6.994)rectangle++(0.04,0.04) (2.008,1.992)rectangle++(0.04,0.04) (1.784,3.209)rectangle++(0.04,0.04) (1.643,4.253)rectangle++(0.04,0.04) (-0.798,2.902)rectangle++(0.04,0.04)
  (-0.565,4.095)rectangle++(0.04,0.04) (-1.199,2.36)rectangle++(0.04,0.04) (-0.927,3.65)rectangle++(0.04,0.04) (-0.644,4.736)rectangle++(0.04,0.04) (-0.36,5.648)rectangle++(0.04,0.04)
  (-0.082,6.411)rectangle++(0.04,0.04) (0.184,7.049)rectangle++(0.04,0.04) (-1.817,3.182)rectangle++(0.04,0.04) (-1.419,4.374)rectangle++(0.04,0.04) (-1.034,5.371)rectangle++(0.04,0.04)
  (-0.666,6.203)rectangle++(0.04,0.04) (-0.32,6.895)rectangle++(0.04,0.04) (-1.878,3.032)rectangle++(0.04,0.04) (-1.185,1.824)rectangle++(0.04,0.04) (-0.937,3.194)rectangle++(0.04,0.04)
  (-0.671,4.349)rectangle++(0.04,0.04) (-0.398,5.319)rectangle++(0.04,0.04) (-1.484,2.668)rectangle++(0.04,0.04) (-1.157,3.923)rectangle++(0.04,0.04) (-0.829,4.977)rectangle++(0.04,0.04)
  (-0.507,5.86)rectangle++(0.04,0.04) (-0.199,6.597)rectangle++(0.04,0.04) (0.092,7.212)rectangle++(0.04,0.04) (0.365,7.723)rectangle++(0.04,0.04) (-1.956,3.372)rectangle++(0.04,0.04)
  (-1.53,4.541)rectangle++(0.04,0.04) (-1.121,5.517)rectangle++(0.04,0.04) (-0.734,6.33)rectangle++(0.04,0.04) (-0.373,7.006)rectangle++(0.04,0.04) (-0.039,7.567)rectangle++(0.04,0.04)
  (0.267,8.029)rectangle++(0.04,0.04) (0.546,8.41)rectangle++(0.04,0.04) (0.798,8.723)rectangle++(0.04,0.04) (1.025,8.979)rectangle++(0.04,0.04) (1.228,9.187)rectangle++(0.04,0.04)
  (1.41,9.356)rectangle++(0.04,0.04) (1.57,9.493)rectangle++(0.04,0.04) (1.712,9.602)rectangle++(0.04,0.04) (1.837,9.69)rectangle++(0.04,0.04) (1.947,9.759)rectangle++(0.04,0.04)
  (2.043,9.813)rectangle++(0.04,0.04) (2.127,9.856)rectangle++(0.04,0.04) (2.2,9.889)rectangle++(0.04,0.04) (2.263,9.913)rectangle++(0.04,0.04) (2.318,9.932)rectangle++(0.04,0.04)
  (2.366,9.946)rectangle++(0.04,0.04) (-0.02,1.575)rectangle++(0.04,0.04) (0.044,2.935)rectangle++(0.04,0.04) (0.152,4.089)rectangle++(0.04,0.04) (0.291,5.066)rectangle++(0.04,0.04)
  (0.448,5.891)rectangle++(0.04,0.04) (0.614,6.585)rectangle++(0.04,0.04) (0.783,7.169)rectangle++(0.04,0.04) (0.95,7.659)rectangle++(0.04,0.04) (1.112,8.068)rectangle++(0.04,0.04)
  (1.266,8.41)rectangle++(0.04,0.04) (-1.955,3.73)rectangle++(0.04,0.04) (1.32,0.817)rectangle++(0.04,0.04) (1.153,2.238)rectangle++(0.04,0.04) (1.067,3.452)rectangle++(0.04,0.04)
  (1.043,4.488)rectangle++(0.04,0.04) (1.064,5.369)rectangle++(0.04,0.04) (1.117,6.118)rectangle++(0.04,0.04) (1.192,6.751)rectangle++(0.04,0.04) (1.281,7.287)rectangle++(0.04,0.04)
  (1.378,7.739)rectangle++(0.04,0.04) (1.479,8.119)rectangle++(0.04,0.04) (1.579,8.438)rectangle++(0.04,0.04) (1.678,8.706)rectangle++(0.04,0.04) (1.772,8.929)rectangle++(0.04,0.04)
  (1.861,9.115)rectangle++(0.04,0.04) (1.945,9.269)rectangle++(0.04,0.04) (2.021,9.397)rectangle++(0.04,0.04) (2.092,9.503)rectangle++(0.04,0.04) (2.156,9.59)rectangle++(0.04,0.04)
  (2.214,9.662)rectangle++(0.04,0.04) (2.266,9.72)rectangle++(0.04,0.04) (2.313,9.768)rectangle++(0.04,0.04) (2.354,9.806)rectangle++(0.04,0.04) (2.391,9.837)rectangle++(0.04,0.04)
  (2.424,9.862)rectangle++(0.04,0.04) (2.453,9.882)rectangle++(0.04,0.04) (2.478,9.898)rectangle++(0.04,0.04) (2.5,9.91)rectangle++(0.04,0.04) (2.383,3.458)rectangle++(0.04,0.04)
  (2.161,4.441)rectangle++(0.04,0.04) (2.012,5.284)rectangle++(0.04,0.04) (1.92,6.007)rectangle++(0.04,0.04) (1.87,6.626)rectangle++(0.04,0.04) (1.852,7.153)rectangle++(0.04,0.04)
  (1.858,7.602)rectangle++(0.04,0.04) (1.881,7.984)rectangle++(0.04,0.04) (1.916,8.307)rectangle++(0.04,0.04) (1.959,8.581)rectangle++(0.04,0.04) (2.006,8.811)rectangle++(0.04,0.04)
  (2.056,9.006)rectangle++(0.04,0.04) (2.105,9.169)rectangle++(0.04,0.04) (-1.984,4.09)rectangle++(0.04,0.04) (-1.525,5.152)rectangle++(0.04,0.04) (-1.092,6.037)rectangle++(0.04,0.04)
  (-0.689,6.771)rectangle++(0.04,0.04) (-0.317,7.379)rectangle++(0.04,0.04) (2.096,2.119)rectangle++(0.04,0.04) (0.261,1.483)rectangle++(0.04,0.04) (0.279,2.847)rectangle++(0.04,0.04)
  (0.349,4.005)rectangle++(0.04,0.04) (0.455,4.986)rectangle++(0.04,0.04) (0.584,5.816)rectangle++(0.04,0.04) (0.727,6.517)rectangle++(0.04,0.04) (1.698,2.183)rectangle++(0.04,0.04)
  (1.529,3.384)rectangle++(0.04,0.04) (1.433,4.411)rectangle++(0.04,0.04) (1.392,5.288)rectangle++(0.04,0.04) (1.392,6.036)rectangle++(0.04,0.04) (1.423,6.671)rectangle++(0.04,0.04)
  (1.474,7.21)rectangle++(0.04,0.04) (1.78,2.544)rectangle++(0.04,0.04) (1.613,3.687)rectangle++(0.04,0.04) (-0.657,2.771)rectangle++(0.04,0.04) (-0.45,3.978)rectangle++(0.04,0.04)
  (-0.226,4.995)rectangle++(0.04,0.04) (0.006,5.851)rectangle++(0.04,0.04) (0.237,6.57)rectangle++(0.04,0.04) (1.787,2.068)rectangle++(0.04,0.04) (1.599,3.283)rectangle++(0.04,0.04)
  (1.488,4.323)rectangle++(0.04,0.04) (1.436,5.211)rectangle++(0.04,0.04) (1.427,5.968)rectangle++(0.04,0.04) (-1.288,3.23)rectangle++(0.04,0.04) (-1.119,2.003)rectangle++(0.04,0.04)
  (0.711,0.62)rectangle++(0.04,0.04) (0.627,2.095)rectangle++(0.04,0.04) (0.615,3.352)rectangle++(0.04,0.04) (0.654,4.421)rectangle++(0.04,0.04) (0.731,5.327)rectangle++(0.04,0.04)
  (0.832,6.095)rectangle++(0.04,0.04) (0.949,6.744)rectangle++(0.04,0.04) (1.074,7.291)rectangle++(0.04,0.04) (1.202,7.75)rectangle++(0.04,0.04) (1.33,8.136)rectangle++(0.04,0.04)
  (1.454,8.458)rectangle++(0.04,0.04) (1.572,8.728)rectangle++(0.04,0.04) (1.683,8.952)rectangle++(0.04,0.04) (1.786,9.138)rectangle++(0.04,0.04) (1.882,9.292)rectangle++(0.04,0.04)
  (-0.02,1.47)rectangle++(0.04,0.04) (0.04,2.846)rectangle++(0.04,0.04) (0.145,4.014)rectangle++(0.04,0.04) (0.282,5.002)rectangle++(0.04,0.04) (1.341,1.704)rectangle++(0.04,0.04)
  (-0.196,2.272)rectangle++(0.04,0.04) (-0.078,3.535)rectangle++(0.04,0.04) (0.073,4.604)rectangle++(0.04,0.04) (0.244,5.507)rectangle++(0.04,0.04) (0.426,6.267)rectangle++(0.04,0.04)
  (0.61,6.907)rectangle++(0.04,0.04) (0.793,7.442)rectangle++(0.04,0.04) (0.969,7.89)rectangle++(0.04,0.04) (2.047,2.576)rectangle++(0.04,0.04) (1.84,3.704)rectangle++(0.04,0.04)
  (1.71,4.671)rectangle++(0.04,0.04) (1.638,5.498)rectangle++(0.04,0.04) (1.61,6.204)rectangle++(0.04,0.04) (1.478,2.338)rectangle++(0.04,0.04) (1.348,3.524)rectangle++(0.04,0.04)
  (1.284,4.538)rectangle++(0.04,0.04) (1.271,5.402)rectangle++(0.04,0.04) (1.294,6.137)rectangle++(0.04,0.04) (1.343,6.761)rectangle++(0.04,0.04) (1.41,7.289)rectangle++(0.04,0.04)
  (1.488,7.736)rectangle++(0.04,0.04) (1.572,8.112)rectangle++(0.04,0.04) (1.658,8.428)rectangle++(0.04,0.04) (1.745,8.694)rectangle++(0.04,0.04) (1.828,8.916)rectangle++(0.04,0.04)
  (1.909,9.102)rectangle++(0.04,0.04) (1.984,9.257)rectangle++(0.04,0.04) (2.055,9.385)rectangle++(0.04,0.04) (2.12,9.491)rectangle++(0.04,0.04) (2.322,3.259)rectangle++(0.04,0.04)
  (2.102,4.273)rectangle++(0.04,0.04) (1.955,5.145)rectangle++(0.04,0.04) (-0.02,0.806)rectangle++(0.04,0.04) (0.013,2.282)rectangle++(0.04,0.04) (0.1,3.536)rectangle++(0.04,0.04)
  (-0.92,2.39)rectangle++(0.04,0.04) (-0.689,3.664)rectangle++(0.04,0.04) (-0.441,4.739)rectangle++(0.04,0.04) (-0.188,5.642)rectangle++(0.04,0.04) (0.064,6.399)rectangle++(0.04,0.04)
  (0.308,7.033)rectangle++(0.04,0.04) (0.541,7.562)rectangle++(0.04,0.04) (0.76,8.002)rectangle++(0.04,0.04) (-1.95,3.524)rectangle++(0.04,0.04) (-1.518,4.67)rectangle++(0.04,0.04)
  (-1.106,5.626)rectangle++(0.04,0.04) (-0.717,6.423)rectangle++(0.04,0.04) (-0.355,7.084)rectangle++(0.04,0.04) (2.019,2.038)rectangle++(0.04,0.04) (1.796,3.248)rectangle++(0.04,0.04)
  (1.654,4.285)rectangle++(0.04,0.04) (1.575,5.172)rectangle++(0.04,0.04) (1.544,5.93)rectangle++(0.04,0.04) (1.547,6.575)rectangle++(0.04,0.04) (1.576,7.123)rectangle++(0.04,0.04)
  (1.622,7.587)rectangle++(0.04,0.04) (1.68,7.981)rectangle++(0.04,0.04) (1.745,8.313)rectangle++(0.04,0.04) (1.814,8.592)rectangle++(0.04,0.04) (1.883,8.827)rectangle++(0.04,0.04)
  (2.172,3.038)rectangle++(0.04,0.04) (1.965,4.092)rectangle++(0.04,0.04) (1.832,4.996)rectangle++(0.04,0.04) (1.107,2.105)rectangle++(0.04,0.04) (1.023,3.341)rectangle++(0.04,0.04)
  (1.001,4.395)rectangle++(0.04,0.04) (1.063,1.745)rectangle++(0.04,0.04) (0.971,3.037)rectangle++(0.04,0.04) (0.945,4.139)rectangle++(0.04,0.04) (1,1.669)rectangle++(0.04,0.04)
  (0.914,2.975)rectangle++(0.04,0.04) (0.894,4.088)rectangle++(0.04,0.04) (-0.905,2.694)rectangle++(0.04,0.04) (0.873,0.841)rectangle++(0.04,0.04) (0.773,2.276)rectangle++(0.04,0.04)
  (0.746,3.5)rectangle++(0.04,0.04) (0.851,1.464)rectangle++(0.04,0.04) (0.779,2.807)rectangle++(0.04,0.04) (0.773,3.951)rectangle++(0.04,0.04) (0.813,4.923)rectangle++(0.04,0.04)
  (0.885,5.749)rectangle++(0.04,0.04) (0.98,6.447)rectangle++(0.04,0.04) (1.089,7.037)rectangle++(0.04,0.04) (-1.633,3.388)rectangle++(0.04,0.04) (-1.255,4.541)rectangle++(0.04,0.04)
  (-1.453,2.299)rectangle++(0.04,0.04) (-0.91,1.761)rectangle++(0.04,0.04) (-0.705,3.129)rectangle++(0.04,0.04) (-0.476,4.284)rectangle++(0.04,0.04) (-0.236,5.257)rectangle++(0.04,0.04)
  (0.008,6.074)rectangle++(0.04,0.04) (0.247,6.759)rectangle++(0.04,0.04) (0.478,7.331)rectangle++(0.04,0.04) (0.698,7.809)rectangle++(0.04,0.04) (0.903,8.206)rectangle++(0.04,0.04)
  (1.094,8.535)rectangle++(0.04,0.04) (1.269,8.807)rectangle++(0.04,0.04) (1.429,9.031)rectangle++(0.04,0.04) (1.573,9.216)rectangle++(0.04,0.04) (1.704,9.367)rectangle++(0.04,0.04)
  (1.821,9.49)rectangle++(0.04,0.04) (1.925,9.59)rectangle++(0.04,0.04) (2.018,9.67)rectangle++(0.04,0.04) (-2.132,4.181)rectangle++(0.04,0.04) (-1.647,5.235)rectangle++(0.04,0.04)
  (-1.193,6.112)rectangle++(0.04,0.04) (-0.772,6.839)rectangle++(0.04,0.04) (-0.385,7.44)rectangle++(0.04,0.04) (-0.031,7.936)rectangle++(0.04,0.04) (0.288,8.343)rectangle++(0.04,0.04)
  (0.577,8.676)rectangle++(0.04,0.04) (0.835,8.948)rectangle++(0.04,0.04) (1.066,9.168)rectangle++(0.04,0.04) (1.27,9.347)rectangle++(0.04,0.04) (1.451,9.49)rectangle++(0.04,0.04)
  (1.611,9.605)rectangle++(0.04,0.04) (2.43,3.154)rectangle++(0.04,0.04) (2.19,4.18)rectangle++(0.04,0.04) (-0.67,3.012)rectangle++(0.04,0.04) (-0.451,4.183)rectangle++(0.04,0.04)
  (-0.218,5.17)rectangle++(0.04,0.04) (0.019,6)rectangle++(0.04,0.04) (0.254,6.695)rectangle++(0.04,0.04) (0.481,7.277)rectangle++(0.04,0.04) (0.698,7.762)rectangle++(0.04,0.04)
  (0.902,8.166)rectangle++(0.04,0.04) (1.091,8.501)rectangle++(0.04,0.04) (1.265,8.779)rectangle++(0.04,0.04) (1.424,9.008)rectangle++(0.04,0.04) (1.569,9.196)rectangle++(0.04,0.04)
  (1.699,9.35)rectangle++(0.04,0.04) (1.816,9.476)rectangle++(0.04,0.04) (1.92,9.578)rectangle++(0.04,0.04) (2.013,9.66)rectangle++(0.04,0.04) (2.095,9.727)rectangle++(0.04,0.04)
  (2.392,3.309)rectangle++(0.04,0.04) (2.163,4.313)rectangle++(0.04,0.04) (0.866,2.028)rectangle++(0.04,0.04) (-0.375,2.234)rectangle++(0.04,0.04) (-0.677,1.994)rectangle++(0.04,0.04)
  (-0.675,1.872)rectangle++(0.04,0.04) (-0.501,3.214)rectangle++(0.04,0.04) (-0.957,2.181)rectangle++(0.04,0.04) (-0.729,3.488)rectangle++(0.04,0.04) (-0.482,4.59)rectangle++(0.04,0.04)
  (-0.228,5.517)rectangle++(0.04,0.04) (1.562,1.695)rectangle++(0.04,0.04) (0.223,1.243)rectangle++(0.04,0.04) (0.237,2.644)rectangle++(0.04,0.04) (0.305,3.834)rectangle++(0.04,0.04)
  (-0.957,2.503)rectangle++(0.04,0.04) (-0.716,3.762)rectangle++(0.04,0.04) (-0.46,4.822)rectangle++(0.04,0.04) (-0.2,5.714)rectangle++(0.04,0.04) (0.056,6.461)rectangle++(0.04,0.04)
  (0.304,7.086)rectangle++(0.04,0.04) (0.54,7.607)rectangle++(0.04,0.04) (0.761,8.04)rectangle++(0.04,0.04) (-1.96,3.533)rectangle++(0.04,0.04) (-1.527,4.677)rectangle++(0.04,0.04)
  (-1.113,5.633)rectangle++(0.04,0.04) (-0.723,6.429)rectangle++(0.04,0.04) (1.891,1.785)rectangle++(0.04,0.04) (1.677,3.038)rectangle++(0.04,0.04) (1.544,4.111)rectangle++(0.04,0.04)
  (1.475,5.029)rectangle++(0.04,0.04) (1.453,5.812)rectangle++(0.04,0.04) (1.465,6.478)rectangle++(0.04,0.04) (1.502,7.044)rectangle++(0.04,0.04) (1.73,2.511)rectangle++(0.04,0.04)
  (1.568,3.662)rectangle++(0.04,0.04) (1.477,4.646)rectangle++(0.04,0.04) (1.439,5.486)rectangle++(0.04,0.04) (1.441,6.202)rectangle++(0.04,0.04) (1.471,6.81)rectangle++(0.04,0.04)
  (1.52,7.326)rectangle++(0.04,0.04) (1.583,7.762)rectangle++(0.04,0.04) (1.654,8.131)rectangle++(0.04,0.04) (1.729,8.441)rectangle++(0.04,0.04) (1.805,8.702)rectangle++(0.04,0.04)
  (1.88,8.921)rectangle++(0.04,0.04) (1.953,9.104)rectangle++(0.04,0.04) (-1.998,4.041)rectangle++(0.04,0.04) (-1.471,2.019)rectangle++(0.04,0.04) (-1.172,3.371)rectangle++(0.04,0.04)
  (-0.864,4.508)rectangle++(0.04,0.04) (-0.556,5.463)rectangle++(0.04,0.04) (-0.256,6.262)rectangle++(0.04,0.04) (0.03,6.929)rectangle++(0.04,0.04) (0.301,7.485)rectangle++(0.04,0.04)
  (0.553,7.946)rectangle++(0.04,0.04) (0.786,8.328)rectangle++(0.04,0.04) (0.999,8.644)rectangle++(0.04,0.04) (1.192,8.903)rectangle++(0.04,0.04) (1.368,9.116)rectangle++(0.04,0.04)
  (1.525,9.29)rectangle++(0.04,0.04) (1.666,9.432)rectangle++(0.04,0.04) (1.791,9.547)rectangle++(0.04,0.04) (-2.145,4.101)rectangle++(0.04,0.04) (-1.662,5.168)rectangle++(0.04,0.04)
  (-1.208,6.055)rectangle++(0.04,0.04) (-0.787,6.792)rectangle++(0.04,0.04) (-0.399,7.401)rectangle++(0.04,0.04) (-0.045,7.903)rectangle++(0.04,0.04) (0.275,8.315)rectangle++(0.04,0.04)
  (0.564,8.653)rectangle++(0.04,0.04) (0.824,8.929)rectangle++(0.04,0.04) (1.055,9.153)rectangle++(0.04,0.04) (1.261,9.334)rectangle++(0.04,0.04) (-2.196,3.932)rectangle++(0.04,0.04)
  (1.413,0.803)rectangle++(0.04,0.04) (1.231,2.222)rectangle++(0.04,0.04) (1.133,3.436)rectangle++(0.04,0.04) (1.098,4.471)rectangle++(0.04,0.04) (1.11,5.353)rectangle++(0.04,0.04)
  (1.156,6.102)rectangle++(0.04,0.04) (1.224,6.736)rectangle++(0.04,0.04) (1.308,7.273)rectangle++(0.04,0.04) (1.4,7.726)rectangle++(0.04,0.04) (1.497,8.107)rectangle++(0.04,0.04)
  (1.595,8.428)rectangle++(0.04,0.04) (1.69,8.696)rectangle++(0.04,0.04) (1.782,8.92)rectangle++(0.04,0.04) (1.87,9.107)rectangle++(0.04,0.04) (-2.015,4.023)rectangle++(0.04,0.04)
  (-1.554,5.096)rectangle++(0.04,0.04) (-1.119,5.99)rectangle++(0.04,0.04) (-0.714,6.732)rectangle++(0.04,0.04) (-0.34,7.347)rectangle++(0.04,0.04) (-1.999,3.127)rectangle++(0.04,0.04)
  (-1.577,4.334)rectangle++(0.04,0.04) (-1.169,5.343)rectangle++(0.04,0.04) (-1.644,2.496)rectangle++(0.04,0.04) (-1.3,3.783)rectangle++(0.04,0.04) (1.237,1)rectangle++(0.04,0.04)
  (1.089,2.397)rectangle++(0.04,0.04) (1.019,3.59)rectangle++(0.04,0.04) (1.008,4.607)rectangle++(0.04,0.04) (1.039,5.472)rectangle++(0.04,0.04) (1.1,6.206)rectangle++(0.04,0.04)
  (1.181,6.827)rectangle++(0.04,0.04) (1.275,7.352)rectangle++(0.04,0.04) (1.375,7.794)rectangle++(0.04,0.04) (1.478,8.166)rectangle++(0.04,0.04) (1.581,8.478)rectangle++(0.04,0.04)
  (1.681,8.74)rectangle++(0.04,0.04) (1.776,8.958)rectangle++(0.04,0.04) (1.866,9.139)rectangle++(0.04,0.04) (1.949,9.29)rectangle++(0.04,0.04) (2.026,9.415)rectangle++(0.04,0.04)
  (2.097,9.518)rectangle++(0.04,0.04) (2.161,9.602)rectangle++(0.04,0.04) (2.219,9.672)rectangle++(0.04,0.04) (2.27,9.728)rectangle++(0.04,0.04) (2.317,9.775)rectangle++(0.04,0.04)
  (2.358,9.812)rectangle++(0.04,0.04) (2.395,9.842)rectangle++(0.04,0.04) (2.427,9.866)rectangle++(0.04,0.04) (2.455,9.885)rectangle++(0.04,0.04) (2.48,9.9)rectangle++(0.04,0.04)
  (2.383,3.451)rectangle++(0.04,0.04) (-0.02,0.535)rectangle++(0.04,0.04) (0.002,2.052)rectangle++(0.04,0.04) (0.082,3.34)rectangle++(0.04,0.04) (0.201,4.432)rectangle++(0.04,0.04)
  (0.346,5.356)rectangle++(0.04,0.04) (0.506,6.135)rectangle++(0.04,0.04) (0.673,6.79)rectangle++(0.04,0.04) (1.783,2.235)rectangle++(0.04,0.04) (1.603,3.424)rectangle++(0.04,0.04)
  (1.497,4.443)rectangle++(0.04,0.04) (1.448,5.313)rectangle++(0.04,0.04) (1.441,6.054)rectangle++(0.04,0.04) (1.465,6.685)rectangle++(0.04,0.04) (1.51,7.219)rectangle++(0.04,0.04)
  (1.57,7.672)rectangle++(0.04,0.04) (1.64,8.055)rectangle++(0.04,0.04) (1.714,8.377)rectangle++(0.04,0.04) (1.789,8.648)rectangle++(0.04,0.04) (-1.912,3.903)rectangle++(0.04,0.04)
  (-1.471,4.99)rectangle++(0.04,0.04) (1.601,1.245)rectangle++(0.04,0.04) (1.408,2.591)rectangle++(0.04,0.04) (1.298,3.742)rectangle++(0.04,0.04) (1.251,4.725)rectangle++(0.04,0.04)
  (1.25,5.562)rectangle++(0.04,0.04) (1.283,6.274)rectangle++(0.04,0.04) (1.547,2.269)rectangle++(0.04,0.04) (1.403,3.463)rectangle++(0.04,0.04) (1.329,4.484)rectangle++(0.04,0.04)
  (1.307,5.354)rectangle++(0.04,0.04) (1.323,6.095)rectangle++(0.04,0.04) (1.366,6.724)rectangle++(0.04,0.04) (1.428,7.257)rectangle++(0.04,0.04) (1.502,7.708)rectangle++(0.04,0.04)
  (1.915,2.67)rectangle++(0.04,0.04) (1.733,3.789)rectangle++(0.04,0.04) (1.622,4.748)rectangle++(0.04,0.04) (-0.931,3.007)rectangle++(0.04,0.04) (-0.673,4.189)rectangle++(0.04,0.04)
  (-0.407,5.184)rectangle++(0.04,0.04) (-0.141,6.019)rectangle++(0.04,0.04) (0.119,6.718)rectangle++(0.04,0.04) (0.368,7.302)rectangle++(0.04,0.04) (0.602,7.788)rectangle++(0.04,0.04)
  (0.821,8.192)rectangle++(0.04,0.04) (1.024,8.526)rectangle++(0.04,0.04) (-2.033,3.7)rectangle++(0.04,0.04) (-1.39,1.935)rectangle++(0.04,0.04) (-1.106,3.297)rectangle++(0.04,0.04)
  (-0.811,4.443)rectangle++(0.04,0.04) (-0.514,5.405)rectangle++(0.04,0.04) (-0.222,6.211)rectangle++(0.04,0.04) (0.057,6.884)rectangle++(0.04,0.04) (0.322,7.446)rectangle++(0.04,0.04)
  (-0.02,1.175)rectangle++(0.04,0.04) (0.028,2.595)rectangle++(0.04,0.04) (0.125,3.801)rectangle++(0.04,0.04) (0.256,4.822)rectangle++(0.04,0.04) (-0.02,0.755)rectangle++(0.04,0.04)
  (0.197,0.606)rectangle++(0.04,0.04) (-0.139,1.768)rectangle++(0.04,0.04) (-0.05,3.104)rectangle++(0.04,0.04) (-0.838,2.26)rectangle++(0.04,0.04) (-0.624,3.551)rectangle++(0.04,0.04)
  (-0.391,4.64)rectangle++(0.04,0.04) (-0.149,5.556)rectangle++(0.04,0.04) (0.094,6.324)rectangle++(0.04,0.04) (0.33,6.968)rectangle++(0.04,0.04) (0.557,7.506)rectangle++(0.04,0.04)
  (0.772,7.954)rectangle++(0.04,0.04) (0.972,8.326)rectangle++(0.04,0.04) (1.157,8.635)rectangle++(0.04,0.04) (1.327,8.889)rectangle++(0.04,0.04) (1.481,9.099)rectangle++(0.04,0.04)
  (1.621,9.271)rectangle++(0.04,0.04) (1.746,9.412)rectangle++(0.04,0.04) (1.859,9.526)rectangle++(0.04,0.04) (1.959,9.619)rectangle++(0.04,0.04) (2.047,9.694)rectangle++(0.04,0.04)
  (2.126,9.754)rectangle++(0.04,0.04) (2.195,9.802)rectangle++(0.04,0.04) (2.256,9.84)rectangle++(0.04,0.04) (2.4,3.378)rectangle++(0.04,0.04) (2.173,4.372)rectangle++(0.04,0.04)
  (2.019,5.225)rectangle++(0.04,0.04) (1.923,5.957)rectangle++(0.04,0.04) (1.871,6.583)rectangle++(0.04,0.04) (1.851,7.117)rectangle++(0.04,0.04) (1.856,7.571)rectangle++(0.04,0.04)
  (1.878,7.958)rectangle++(0.04,0.04) (1.913,8.285)rectangle++(0.04,0.04) (1.955,8.562)rectangle++(0.04,0.04) (2.087,2.993)rectangle++(0.04,0.04) (1.891,4.057)rectangle++(0.04,0.04)
  (1.768,4.969)rectangle++(0.04,0.04) (1.699,5.749)rectangle++(0.04,0.04) (1.672,6.415)rectangle++(0.04,0.04) (1.676,6.982)rectangle++(0.04,0.04) (-1.501,3.51)rectangle++(0.04,0.04)
  (-1.138,4.64)rectangle++(0.04,0.04) (-0.784,5.586)rectangle++(0.04,0.04) (-0.445,6.376)rectangle++(0.04,0.04) (-0.125,7.033)rectangle++(0.04,0.04) (0.172,7.579)rectangle++(0.04,0.04)
  (0.448,8.032)rectangle++(0.04,0.04) (0.7,8.405)rectangle++(0.04,0.04) (0.929,8.713)rectangle++(0.04,0.04) (1.136,8.965)rectangle++(0.04,0.04) (1.322,9.171)rectangle++(0.04,0.04)
  (1.488,9.339)rectangle++(0.04,0.04) (-0.02,1.477)rectangle++(0.04,0.04) (0.04,2.853)rectangle++(0.04,0.04) (0.146,4.019)rectangle++(0.04,0.04) (0.283,5.007)rectangle++(0.04,0.04)
  (0.438,5.841)rectangle++(0.04,0.04) (0.604,6.543)rectangle++(0.04,0.04) (0.773,7.134)rectangle++(0.04,0.04) (0.94,7.629)rectangle++(0.04,0.04) (1.978,2.505)rectangle++(0.04,0.04)
  (1.779,3.647)rectangle++(0.04,0.04) (1.656,4.625)rectangle++(0.04,0.04) (1.59,5.461)rectangle++(0.04,0.04) (1.568,6.174)rectangle++(0.04,0.04) (1.578,6.782)rectangle++(0.04,0.04)
  (-1.469,3.444)rectangle++(0.04,0.04) (-1.113,4.582)rectangle++(0.04,0.04) (-0.765,5.536)rectangle++(0.04,0.04) (-0.431,6.332)rectangle++(0.04,0.04) (-0.02,0.996)rectangle++(0.04,0.04)
  (0.021,2.444)rectangle++(0.04,0.04) (0.113,3.673)rectangle++(0.04,0.04) (0.994,1.341)rectangle++(0.04,0.04) (0.896,2.696)rectangle++(0.04,0.04) (0.868,3.852)rectangle++(0.04,0.04)
  (0.889,4.836)rectangle++(0.04,0.04) (0.947,5.671)rectangle++(0.04,0.04) (1.03,6.379)rectangle++(0.04,0.04) (1.128,6.977)rectangle++(0.04,0.04) (1.236,7.481)rectangle++(0.04,0.04)
  (1.348,7.906)rectangle++(0.04,0.04) (1.459,8.262)rectangle++(0.04,0.04) (1.569,8.561)rectangle++(0.04,0.04) (1.674,8.81)rectangle++(0.04,0.04) (1.773,9.018)rectangle++(0.04,0.04)
  (2.242,3.055)rectangle++(0.04,0.04) (-0.367,2.777)rectangle++(0.04,0.04) (-0.817,2.115)rectangle++(0.04,0.04) (-0.612,3.427)rectangle++(0.04,0.04) (-0.385,4.534)rectangle++(0.04,0.04)
  (-0.148,5.465)rectangle++(0.04,0.04) (0.09,6.248)rectangle++(0.04,0.04) (0.325,6.903)rectangle++(0.04,0.04) (0.55,7.451)rectangle++(0.04,0.04) (0.763,7.907)rectangle++(0.04,0.04)
  (0.963,8.287)rectangle++(0.04,0.04) (1.148,8.602)rectangle++(0.04,0.04) (1.317,8.862)rectangle++(0.04,0.04) (1.472,9.076)rectangle++(0.04,0.04) (1.612,9.252)rectangle++(0.04,0.04)
  (1.738,9.396)rectangle++(0.04,0.04) (1.851,9.513)rectangle++(0.04,0.04) (1.952,9.608)rectangle++(0.04,0.04) (2.041,9.685)rectangle++(0.04,0.04) (2.388,3.285)rectangle++(0.04,0.04)
  (2.159,4.293)rectangle++(0.04,0.04) (2.005,5.159)rectangle++(0.04,0.04) (1.908,5.901)rectangle++(0.04,0.04) (1.856,6.536)rectangle++(0.04,0.04) (1.837,7.077)rectangle++(0.04,0.04)
  (1.842,7.538)rectangle++(0.04,0.04) (1.865,7.93)rectangle++(0.04,0.04) (1.9,8.262)rectangle++(0.04,0.04) (1.944,8.543)rectangle++(0.04,0.04) (1.992,8.78)rectangle++(0.04,0.04)
  (-1.906,3.979)rectangle++(0.04,0.04) (-1.463,5.054)rectangle++(0.04,0.04) (-1.043,5.951)rectangle++(0.04,0.04) (-0.651,6.696)rectangle++(0.04,0.04) (-0.288,7.314)rectangle++(0.04,0.04)
  (0.046,7.825)rectangle++(0.04,0.04) (0.35,8.245)rectangle++(0.04,0.04) (0.625,8.591)rectangle++(0.04,0.04) (0.873,8.873)rectangle++(0.04,0.04) (-2.154,3.742)rectangle++(0.04,0.04)
  (-1.683,4.863)rectangle++(0.04,0.04) (-1.238,5.797)rectangle++(0.04,0.04) (-1.776,2.579)rectangle++(0.04,0.04) (-1.409,3.86)rectangle++(0.04,0.04) (-1.045,4.933)rectangle++(0.04,0.04)
  (-1.513,2.434)rectangle++(0.04,0.04) (0.891,0.621)rectangle++(0.04,0.04) (0.78,2.088)rectangle++(0.04,0.04) (-0.408,2.228)rectangle++(0.04,0.04) (-0.26,3.506)rectangle++(0.04,0.04)
  (-0.083,4.587)rectangle++(0.04,0.04) (-1.23,2.579)rectangle++(0.04,0.04) (-0.945,3.838)rectangle++(0.04,0.04) (-0.652,4.896)rectangle++(0.04,0.04) (-0.36,5.784)rectangle++(0.04,0.04)
  (-0.077,6.527)rectangle++(0.04,0.04) (-0.02,1.027)rectangle++(0.04,0.04) (0.022,2.47)rectangle++(0.04,0.04) (0.115,3.695)rectangle++(0.04,0.04) (0.244,4.732)rectangle++(0.04,0.04)
  (0.394,5.609)rectangle++(0.04,0.04) (0.557,6.348)rectangle++(0.04,0.04) (0.725,6.97)rectangle++(0.04,0.04) (0.893,7.492)rectangle++(0.04,0.04) (1.057,7.928)rectangle++(0.04,0.04)
  (1.213,8.293)rectangle++(0.04,0.04) (1.361,8.597)rectangle++(0.04,0.04) (1.498,8.849)rectangle++(0.04,0.04) (2.236,2.943)rectangle++(0.04,0.04) (2.016,4.009)rectangle++(0.04,0.04)
  (1.872,4.923)rectangle++(0.04,0.04) (1.786,5.706)rectangle++(0.04,0.04) (1.312,2.264)rectangle++(0.04,0.04) (1.204,3.468)rectangle++(0.04,0.04) (1.16,4.496)rectangle++(0.04,0.04)
  (1.163,5.371)rectangle++(0.04,0.04) (1.202,6.115)rectangle++(0.04,0.04) (1.264,6.746)rectangle++(0.04,0.04) (1.342,7.28)rectangle++(0.04,0.04) (1.43,7.73)rectangle++(0.04,0.04)
  (1.522,8.11)rectangle++(0.04,0.04) (1.616,8.429)rectangle++(0.04,0.04) (1.709,8.696)rectangle++(0.04,0.04) (1.798,8.919)rectangle++(0.04,0.04) (1.883,9.106)rectangle++(0.04,0.04)
  (1.962,9.261)rectangle++(0.04,0.04) (2.036,9.389)rectangle++(0.04,0.04) (-2.055,4.123)rectangle++(0.04,0.04) (-1.584,5.183)rectangle++(0.04,0.04) (-1.141,6.065)rectangle++(0.04,0.04)
  (-0.73,6.797)rectangle++(0.04,0.04) (1.995,1.872)rectangle++(0.04,0.04) (-0.109,2.46)rectangle++(0.04,0.04) (0.004,3.691)rectangle++(0.04,0.04) (0.149,4.734)rectangle++(0.04,0.04)
  (0.313,5.614)rectangle++(0.04,0.04) (0.489,6.355)rectangle++(0.04,0.04) (0.668,6.979)rectangle++(0.04,0.04) (0.844,7.501)rectangle++(0.04,0.04) (1.016,7.939)rectangle++(0.04,0.04)
  (1.179,8.303)rectangle++(0.04,0.04) (1.332,8.607)rectangle++(0.04,0.04) (1.474,8.859)rectangle++(0.04,0.04) (1.605,9.067)rectangle++(0.04,0.04) (-0.02,1.434)rectangle++(0.04,0.04)
  (0.038,2.816)rectangle++(0.04,0.04) (0.143,3.988)rectangle++(0.04,0.04) (0.279,4.98)rectangle++(0.04,0.04) (0.434,5.818)rectangle++(0.04,0.04) (0.599,6.524)rectangle++(0.04,0.04)
  (0.768,7.118)rectangle++(0.04,0.04) (0.936,7.616)rectangle++(0.04,0.04) (1.098,8.032)rectangle++(0.04,0.04) (1.252,8.38)rectangle++(0.04,0.04) (1.397,8.669)rectangle++(0.04,0.04)
  (1.532,8.909)rectangle++(0.04,0.04) (1.657,9.107)rectangle++(0.04,0.04) (1.77,9.271)rectangle++(0.04,0.04) (1.873,9.406)rectangle++(0.04,0.04) (1.966,9.516)rectangle++(0.04,0.04)
  (2.05,9.606)rectangle++(0.04,0.04) (2.124,9.68)rectangle++(0.04,0.04) (2.191,9.739)rectangle++(0.04,0.04) (2.249,9.787)rectangle++(0.04,0.04) (2.301,9.825)rectangle++(0.04,0.04)
  (2.347,9.855)rectangle++(0.04,0.04) (2.39,3.405)rectangle++(0.04,0.04) (2.166,4.395)rectangle++(0.04,0.04) (2.014,5.246)rectangle++(0.04,0.04) (1.92,5.974)rectangle++(0.04,0.04)
  (1.869,6.598)rectangle++(0.04,0.04) (1.85,7.129)rectangle++(0.04,0.04) (1.855,7.582)rectangle++(0.04,0.04) (1.878,7.967)rectangle++(0.04,0.04) (1.913,8.293)rectangle++(0.04,0.04)
  (1.956,8.569)rectangle++(0.04,0.04) (2.003,8.801)rectangle++(0.04,0.04) (2.052,8.997)rectangle++(0.04,0.04) (2.102,9.162)rectangle++(0.04,0.04) (2.151,9.3)rectangle++(0.04,0.04)
  (2.198,9.415)rectangle++(0.04,0.04) (2.243,9.511)rectangle++(0.04,0.04) (2.285,9.591)rectangle++(0.04,0.04) (2.323,9.657)rectangle++(0.04,0.04) (2.359,9.712)rectangle++(0.04,0.04)
  (2.391,9.757)rectangle++(0.04,0.04) (2.421,9.794)rectangle++(0.04,0.04) (2.447,9.824)rectangle++(0.04,0.04) (2.471,9.849)rectangle++(0.04,0.04) (2.492,9.869)rectangle++(0.04,0.04)
  (2.511,9.885)rectangle++(0.04,0.04) (2.527,9.898)rectangle++(0.04,0.04) (2.542,9.908)rectangle++(0.04,0.04) (2.555,9.917)rectangle++(0.04,0.04) (2.566,9.923)rectangle++(0.04,0.04)
  (-2.088,4.362)rectangle++(0.04,0.04) (-1.603,5.388)rectangle++(0.04,0.04) (-1.149,6.24)rectangle++(0.04,0.04) (-0.729,6.946)rectangle++(0.04,0.04) (-0.344,7.529)rectangle++(0.04,0.04)
  (0.006,8.01)rectangle++(0.04,0.04) (-0.02,1.265)rectangle++(0.04,0.04) (0.031,2.672)rectangle++(0.04,0.04) (0.131,3.866)rectangle++(0.04,0.04) (-1,2.47)rectangle++(0.04,0.04)
  (-0.754,3.736)rectangle++(0.04,0.04) (-0.493,4.802)rectangle++(0.04,0.04) (-0.229,5.697)rectangle++(0.04,0.04) (0.031,6.448)rectangle++(0.04,0.04) (0.282,7.076)rectangle++(0.04,0.04)
  (0.52,7.599)rectangle++(0.04,0.04) (0.744,8.035)rectangle++(0.04,0.04) (0.952,8.396)rectangle++(0.04,0.04) (1.143,8.695)rectangle++(0.04,0.04) (1.317,8.941)rectangle++(0.04,0.04)
  (1.475,9.143)rectangle++(0.04,0.04) (1.617,9.309)rectangle++(0.04,0.04) (-0.02,1.473)rectangle++(0.04,0.04) (0.398,0.778)rectangle++(0.04,0.04) (0.367,2.242)rectangle++(0.04,0.04)
  (0.4,3.487)rectangle++(0.04,0.04) (0.477,4.544)rectangle++(0.04,0.04) (0.585,5.44)rectangle++(0.04,0.04) (0.713,6.197)rectangle++(0.04,0.04) (0.851,6.835)rectangle++(0.04,0.04)
  (0.995,7.372)rectangle++(0.04,0.04) (1.138,7.822)rectangle++(0.04,0.04) (1.278,8.2)rectangle++(0.04,0.04) (1.412,8.515)rectangle++(0.04,0.04) (1.539,8.777)rectangle++(0.04,0.04)
  (1.657,8.995)rectangle++(0.04,0.04) (1.766,9.176)rectangle++(0.04,0.04) (1.866,9.325)rectangle++(0.04,0.04) (1.957,9.448)rectangle++(0.04,0.04) (2.039,9.549)rectangle++(0.04,0.04)
  (2.113,9.631)rectangle++(0.04,0.04) (2.179,9.698)rectangle++(0.04,0.04) (2.238,9.752)rectangle++(0.04,0.04) (2.29,9.796)rectangle++(0.04,0.04) (2.336,9.831)rectangle++(0.04,0.04)
  (-2.11,4.289)rectangle++(0.04,0.04) (-1.624,5.326)rectangle++(0.04,0.04) (-1.17,6.189)rectangle++(0.04,0.04) (-0.749,6.903)rectangle++(0.04,0.04) (-0.363,7.494)rectangle++(0.04,0.04)
  (-0.011,7.981)rectangle++(0.04,0.04) (0.308,8.38)rectangle++(0.04,0.04) (0.595,8.707)rectangle++(0.04,0.04) (2.331,2.674)rectangle++(0.04,0.04) (2.086,3.776)rectangle++(0.04,0.04)
  (1.922,4.722)rectangle++(0.04,0.04) (1.821,5.533)rectangle++(0.04,0.04) (-1.096,3.225)rectangle++(0.04,0.04) (-0.805,4.381)rectangle++(0.04,0.04) (-0.511,5.353)rectangle++(0.04,0.04)
  (-0.222,6.166)rectangle++(0.04,0.04) (0.055,6.846)rectangle++(0.04,0.04) (0.319,7.413)rectangle++(0.04,0.04) (2.011,2.292)rectangle++(0.04,0.04) (1.798,3.464)rectangle++(0.04,0.04)
  (1.665,4.469)rectangle++(0.04,0.04) (1.592,5.328)rectangle++(0.04,0.04) (1.564,6.061)rectangle++(0.04,0.04) (1.57,6.686)rectangle++(0.04,0.04) (1.599,7.216)rectangle++(0.04,0.04)
  (-1.578,3.544)rectangle++(0.04,0.04) (-1.201,4.672)rectangle++(0.04,0.04) (-0.836,5.616)rectangle++(0.04,0.04) (-0.489,6.403)rectangle++(0.04,0.04) (1.849,1.84)rectangle++(0.04,0.04)
  (1.643,3.086)rectangle++(0.04,0.04) (1.518,4.154)rectangle++(0.04,0.04) (-0.798,2.852)rectangle++(0.04,0.04) (-0.566,4.052)rectangle++(0.04,0.04) (-0.321,5.063)rectangle++(0.04,0.04);
\fill[nuage]
  (-0.073,5.913)rectangle++(0.04,0.04) (0.172,6.625)rectangle++(0.04,0.04) (0.409,7.221)rectangle++(0.04,0.04) (0.635,7.717)rectangle++(0.04,0.04) (0.846,8.13)rectangle++(0.04,0.04)
  (1.042,8.473)rectangle++(0.04,0.04) (1.222,8.757)rectangle++(0.04,0.04) (1.387,8.991)rectangle++(0.04,0.04) (1.537,9.183)rectangle++(0.04,0.04) (1.671,9.34)rectangle++(0.04,0.04)
  (2.347,3.106)rectangle++(0.04,0.04) (2.117,4.143)rectangle++(0.04,0.04) (0.825,1.975)rectangle++(0.04,0.04) (0.778,3.242)rectangle++(0.04,0.04) (0.789,4.32)rectangle++(0.04,0.04)
  (0.841,5.237)rectangle++(0.04,0.04) (0.922,6.014)rectangle++(0.04,0.04) (1.022,6.671)rectangle++(0.04,0.04) (1.134,7.226)rectangle++(0.04,0.04) (1.25,7.693)rectangle++(0.04,0.04)
  (1.949,2.601)rectangle++(0.04,0.04) (1.759,3.729)rectangle++(0.04,0.04) (-0.639,2.814)rectangle++(0.04,0.04) (-0.433,4.014)rectangle++(0.04,0.04) (-0.21,5.025)rectangle++(0.04,0.04)
  (0.021,5.876)rectangle++(0.04,0.04) (0.25,6.59)rectangle++(0.04,0.04) (0.474,7.188)rectangle++(0.04,0.04) (0.688,7.687)rectangle++(0.04,0.04) (0.89,8.102)rectangle++(0.04,0.04)
  (1.079,8.448)rectangle++(0.04,0.04) (1.253,8.733)rectangle++(0.04,0.04) (1.412,8.97)rectangle++(0.04,0.04) (1.557,9.164)rectangle++(0.04,0.04) (-2.092,3.963)rectangle++(0.04,0.04)
  (-1.47,1.98)rectangle++(0.04,0.04) (-1.173,3.338)rectangle++(0.04,0.04) (-0.865,4.48)rectangle++(0.04,0.04) (-0.559,5.439)rectangle++(0.04,0.04) (-1.547,2.657)rectangle++(0.04,0.04)
  (0.959,0.665)rectangle++(0.04,0.04) (0.839,2.124)rectangle++(0.04,0.04) (0.796,3.368)rectangle++(0.04,0.04) (0.809,4.427)rectangle++(0.04,0.04) (0.863,5.327)rectangle++(0.04,0.04)
  (0.944,6.089)rectangle++(0.04,0.04) (1.044,6.734)rectangle++(0.04,0.04) (1.154,7.279)rectangle++(0.04,0.04) (1.27,7.737)rectangle++(0.04,0.04) (1.387,8.122)rectangle++(0.04,0.04)
  (2.049,2.74)rectangle++(0.04,0.04) (1.849,3.843)rectangle++(0.04,0.04) (1.723,4.789)rectangle++(0.04,0.04) (1.654,5.598)rectangle++(0.04,0.04) (1.627,6.288)rectangle++(0.04,0.04)
  (1.633,6.876)rectangle++(0.04,0.04) (1.661,7.376)rectangle++(0.04,0.04) (1.704,7.799)rectangle++(0.04,0.04) (1.758,8.157)rectangle++(0.04,0.04) (1.819,8.459)rectangle++(0.04,0.04)
  (-1.857,3.868)rectangle++(0.04,0.04) (1.344,0.876)rectangle++(0.04,0.04) (1.175,2.287)rectangle++(0.04,0.04) (-0.381,2.362)rectangle++(0.04,0.04) (-0.231,3.619)rectangle++(0.04,0.04)
  (-0.054,4.682)rectangle++(0.04,0.04) (1.302,1.54)rectangle++(0.04,0.04) (1.166,2.853)rectangle++(0.04,0.04) (1.103,3.974)rectangle++(0.04,0.04) (1.094,4.93)rectangle++(0.04,0.04)
  (1.125,5.743)rectangle++(0.04,0.04) (1.184,6.433)rectangle++(0.04,0.04) (1.261,7.017)rectangle++(0.04,0.04) (1.351,7.51)rectangle++(0.04,0.04) (1.446,7.926)rectangle++(0.04,0.04)
  (1.544,8.275)rectangle++(0.04,0.04) (1.641,8.568)rectangle++(0.04,0.04) (1.736,8.814)rectangle++(0.04,0.04) (1.826,9.018)rectangle++(0.04,0.04) (2.234,3.069)rectangle++(0.04,0.04)
  (2.019,4.116)rectangle++(0.04,0.04) (1.879,5.014)rectangle++(0.04,0.04) (1.795,5.783)rectangle++(0.04,0.04) (1.755,6.44)rectangle++(0.04,0.04) (1.747,7)rectangle++(0.04,0.04)
  (1.763,7.476)rectangle++(0.04,0.04) (-1.612,3.639)rectangle++(0.04,0.04) (-1.227,4.754)rectangle++(0.04,0.04) (1.498,1.252)rectangle++(0.04,0.04) (1.321,2.6)rectangle++(0.04,0.04)
  (1.225,3.754)rectangle++(0.04,0.04) (1.189,4.738)rectangle++(0.04,0.04) (1.198,5.576)rectangle++(0.04,0.04) (1.239,6.288)rectangle++(0.04,0.04) (1.302,6.891)rectangle++(0.04,0.04)
  (1.381,7.402)rectangle++(0.04,0.04) (1.467,7.832)rectangle++(0.04,0.04) (1.558,8.195)rectangle++(0.04,0.04) (1.65,8.5)rectangle++(0.04,0.04) (1.74,8.755)rectangle++(0.04,0.04)
  (1.827,8.968)rectangle++(0.04,0.04) (1.91,9.146)rectangle++(0.04,0.04) (1.987,9.294)rectangle++(0.04,0.04) (2.058,9.417)rectangle++(0.04,0.04) (2.124,9.518)rectangle++(0.04,0.04)
  (2.184,9.602)rectangle++(0.04,0.04) (2.238,9.67)rectangle++(0.04,0.04) (2.287,9.726)rectangle++(0.04,0.04) (2.331,9.772)rectangle++(0.04,0.04) (-0.02,1.547)rectangle++(0.04,0.04)
  (0.419,0.796)rectangle++(0.04,0.04) (0.386,2.256)rectangle++(0.04,0.04) (0.556,1.072)rectangle++(0.04,0.04) (0.514,2.485)rectangle++(0.04,0.04) (0.534,3.688)rectangle++(0.04,0.04)
  (0.599,4.709)rectangle++(0.04,0.04) (0.695,5.575)rectangle++(0.04,0.04) (-1.332,2.976)rectangle++(0.04,0.04) (-1.015,4.179)rectangle++(0.04,0.04) (-0.698,5.189)rectangle++(0.04,0.04)
  (-0.388,6.034)rectangle++(0.04,0.04) (-0.091,6.741)rectangle++(0.04,0.04) (-0.02,1.062)rectangle++(0.04,0.04) (0.023,2.499)rectangle++(0.04,0.04) (0.118,3.72)rectangle++(0.04,0.04)
  (0.247,4.753)rectangle++(0.04,0.04) (0.397,5.627)rectangle++(0.04,0.04) (0.561,6.363)rectangle++(0.04,0.04) (0.729,6.982)rectangle++(0.04,0.04) (0.897,7.502)rectangle++(0.04,0.04)
  (1.06,7.937)rectangle++(0.04,0.04) (1.216,8.3)rectangle++(0.04,0.04) (1.364,8.603)rectangle++(0.04,0.04) (1.501,8.854)rectangle++(0.04,0.04) (1.628,9.062)rectangle++(0.04,0.04)
  (1.744,9.234)rectangle++(0.04,0.04) (-2.073,4.022)rectangle++(0.04,0.04) (-1.481,1.997)rectangle++(0.04,0.04) (-1.182,3.353)rectangle++(0.04,0.04) (-0.872,4.493)rectangle++(0.04,0.04)
  (-0.564,5.45)rectangle++(0.04,0.04) (-0.264,6.252)rectangle++(0.04,0.04) (0.024,6.921)rectangle++(0.04,0.04) (0.295,7.478)rectangle++(0.04,0.04) (0.548,7.941)rectangle++(0.04,0.04)
  (2.124,2.478)rectangle++(0.04,0.04) (0.358,1.577)rectangle++(0.04,0.04) (0.365,2.922)rectangle++(0.04,0.04) (0.425,4.066)rectangle++(0.04,0.04) (0.522,5.035)rectangle++(0.04,0.04)
  (0.643,5.855)rectangle++(0.04,0.04) (0.778,6.547)rectangle++(0.04,0.04) (0.921,7.13)rectangle++(0.04,0.04) (1.066,7.62)rectangle++(0.04,0.04) (1.209,8.031)rectangle++(0.04,0.04)
  (1.346,8.374)rectangle++(0.04,0.04) (1.477,8.66)rectangle++(0.04,0.04) (-1.977,3.834)rectangle++(0.04,0.04) (-1.53,4.934)rectangle++(0.04,0.04) (-1.105,5.851)rectangle++(0.04,0.04)
  (-0.707,6.614)rectangle++(0.04,0.04) (-0.339,7.247)rectangle++(0.04,0.04) (0,7.769)rectangle++(0.04,0.04) (0.308,8.2)rectangle++(0.04,0.04) (0.588,8.554)rectangle++(0.04,0.04)
  (-2.128,3.606)rectangle++(0.04,0.04) (-1.666,4.746)rectangle++(0.04,0.04) (1.562,1.136)rectangle++(0.04,0.04) (-0.004,2.198)rectangle++(0.04,0.04) (0.082,3.465)rectangle++(0.04,0.04)
  (0.206,4.538)rectangle++(0.04,0.04) (0.355,5.445)rectangle++(0.04,0.04) (-1.366,2.868)rectangle++(0.04,0.04) (-1.049,4.089)rectangle++(0.04,0.04) (-0.73,5.114)rectangle++(0.04,0.04)
  (-0.418,5.972)rectangle++(0.04,0.04) (-0.119,6.689)rectangle++(0.04,0.04) (0.164,7.287)rectangle++(0.04,0.04) (1.998,2.222)rectangle++(0.04,0.04) (1.785,3.405)rectangle++(0.04,0.04)
  (1.651,4.419)rectangle++(0.04,0.04) (1.578,5.286)rectangle++(0.04,0.04) (1.551,6.026)rectangle++(0.04,0.04) (-1.278,3.271)rectangle++(0.04,0.04) (-0.958,4.428)rectangle++(0.04,0.04)
  (-0.639,5.398)rectangle++(0.04,0.04) (-0.329,6.21)rectangle++(0.04,0.04) (-0.034,6.888)rectangle++(0.04,0.04) (0.245,7.452)rectangle++(0.04,0.04) (0.504,7.921)rectangle++(0.04,0.04)
  (0.743,8.309)rectangle++(0.04,0.04) (0.962,8.629)rectangle++(0.04,0.04) (2.254,2.751)rectangle++(0.04,0.04) (2.024,3.844)rectangle++(0.04,0.04) (1.872,4.783)rectangle++(0.04,0.04)
  (1.78,5.587)rectangle++(0.04,0.04) (1.735,6.274)rectangle++(0.04,0.04) (1.723,6.86)rectangle++(0.04,0.04) (1.737,7.358)rectangle++(0.04,0.04) (1.768,7.781)rectangle++(0.04,0.04)
  (1.812,8.139)rectangle++(0.04,0.04) (1.864,8.442)rectangle++(0.04,0.04) (1.92,8.697)rectangle++(0.04,0.04) (1.977,8.912)rectangle++(0.04,0.04) (2.035,9.092)rectangle++(0.04,0.04)
  (2.223,3.141)rectangle++(0.04,0.04) (-0.393,2.791)rectangle++(0.04,0.04) (-0.225,3.985)rectangle++(0.04,0.04) (-0.034,4.992)rectangle++(0.04,0.04) (-1.326,2.679)rectangle++(0.04,0.04)
  (-1.022,3.927)rectangle++(0.04,0.04) (-0.714,4.975)rectangle++(0.04,0.04) (-0.41,5.853)rectangle++(0.04,0.04) (-0.117,6.588)rectangle++(0.04,0.04) (0.162,7.201)rectangle++(0.04,0.04)
  (0.424,7.71)rectangle++(0.04,0.04) (2.078,2.391)rectangle++(0.04,0.04) (0.34,1.544)rectangle++(0.04,0.04) (0.349,2.895)rectangle++(0.04,0.04) (0.41,4.043)rectangle++(0.04,0.04)
  (0.508,5.016)rectangle++(0.04,0.04) (0.63,5.84)rectangle++(0.04,0.04) (0.767,6.535)rectangle++(0.04,0.04) (0.911,7.12)rectangle++(0.04,0.04) (1.057,7.612)rectangle++(0.04,0.04)
  (1.201,8.024)rectangle++(0.04,0.04) (1.34,8.369)rectangle++(0.04,0.04) (1.471,8.656)rectangle++(0.04,0.04) (1.594,8.895)rectangle++(0.04,0.04) (1.709,9.093)rectangle++(0.04,0.04)
  (1.814,9.257)rectangle++(0.04,0.04) (1.91,9.392)rectangle++(0.04,0.04) (2.326,3.181)rectangle++(0.04,0.04) (2.102,4.207)rectangle++(0.04,0.04) (1.953,5.088)rectangle++(0.04,0.04)
  (1.861,5.843)rectangle++(0.04,0.04) (1.814,6.488)rectangle++(0.04,0.04) (-1.345,3.433)rectangle++(0.04,0.04) (-1.183,2.035)rectangle++(0.04,0.04) (-0.926,3.373)rectangle++(0.04,0.04)
  (-0.655,4.5)rectangle++(0.04,0.04) (-0.379,5.448)rectangle++(0.04,0.04) (-1.513,2.7)rectangle++(0.04,0.04) (-1.181,3.952)rectangle++(0.04,0.04) (-0.848,5.003)rectangle++(0.04,0.04)
  (-0.523,5.882)rectangle++(0.04,0.04) (-0.211,6.617)rectangle++(0.04,0.04) (0.083,7.229)rectangle++(0.04,0.04) (0.358,7.738)rectangle++(0.04,0.04) (0.611,8.159)rectangle++(0.04,0.04)
  (0.844,8.507)rectangle++(0.04,0.04) (1.055,8.793)rectangle++(0.04,0.04) (1.246,9.028)rectangle++(0.04,0.04) (1.418,9.22)rectangle++(0.04,0.04) (1.572,9.377)rectangle++(0.04,0.04)
  (-2.145,4.014)rectangle++(0.04,0.04) (-1.665,5.093)rectangle++(0.04,0.04) (-1.213,5.992)rectangle++(0.04,0.04) (-0.794,6.738)rectangle++(0.04,0.04) (-0.408,7.355)rectangle++(0.04,0.04)
  (-0.054,7.865)rectangle++(0.04,0.04) (0.266,8.283)rectangle++(0.04,0.04) (-2.122,3.473)rectangle++(0.04,0.04) (-1.667,4.633)rectangle++(0.04,0.04) (-1.234,5.601)rectangle++(0.04,0.04)
  (-0.827,6.406)rectangle++(0.04,0.04) (-1.852,2.808)rectangle++(0.04,0.04) (-1.464,4.057)rectangle++(0.04,0.04) (-1.084,5.103)rectangle++(0.04,0.04) (-0.72,5.978)rectangle++(0.04,0.04)
  (-0.375,6.706)rectangle++(0.04,0.04) (-0.053,7.311)rectangle++(0.04,0.04) (0.245,7.813)rectangle++(0.04,0.04) (0.519,8.227)rectangle++(0.04,0.04) (0.768,8.569)rectangle++(0.04,0.04)
  (0.993,8.849)rectangle++(0.04,0.04) (1.196,9.078)rectangle++(0.04,0.04) (1.378,9.265)rectangle++(0.04,0.04) (1.539,9.416)rectangle++(0.04,0.04) (1.683,9.538)rectangle++(0.04,0.04)
  (1.81,9.636)rectangle++(0.04,0.04) (1.922,9.715)rectangle++(0.04,0.04) (2.02,9.777)rectangle++(0.04,0.04) (2.106,9.826)rectangle++(0.04,0.04) (2.181,9.864)rectangle++(0.04,0.04)
  (2.246,9.893)rectangle++(0.04,0.04) (-2.144,4.282)rectangle++(0.04,0.04) (-1.563,2.038)rectangle++(0.04,0.04) (-1.25,3.391)rectangle++(0.04,0.04) (-0.929,4.528)rectangle++(0.04,0.04)
  (-0.61,5.483)rectangle++(0.04,0.04) (-1.569,2.655)rectangle++(0.04,0.04) (-1.229,3.915)rectangle++(0.04,0.04) (-0.891,4.974)rectangle++(0.04,0.04) (-0.56,5.859)rectangle++(0.04,0.04)
  (-0.244,6.599)rectangle++(0.04,0.04) (0.054,7.215)rectangle++(0.04,0.04) (0.333,7.727)rectangle++(0.04,0.04) (-1.964,3.365)rectangle++(0.04,0.04) (-1.289,1.878)rectangle++(0.04,0.04)
  (-1.023,3.244)rectangle++(0.04,0.04) (-0.742,4.394)rectangle++(0.04,0.04) (-0.457,5.361)rectangle++(0.04,0.04) (1.552,1.598)rectangle++(0.04,0.04) (1.381,2.892)rectangle++(0.04,0.04)
  (1.287,3.999)rectangle++(0.04,0.04) (0.909,1.725)rectangle++(0.04,0.04) (0.84,3.026)rectangle++(0.04,0.04) (0.833,4.134)rectangle++(0.04,0.04) (0.871,5.077)rectangle++(0.04,0.04)
  (0.941,5.877)rectangle++(0.04,0.04) (1.033,6.554)rectangle++(0.04,0.04) (1.138,7.126)rectangle++(0.04,0.04) (1.25,7.608)rectangle++(0.04,0.04) (1.365,8.013)rectangle++(0.04,0.04)
  (1.478,8.352)rectangle++(0.04,0.04) (1.588,8.637)rectangle++(0.04,0.04) (1.693,8.874)rectangle++(0.04,0.04) (1.792,9.071)rectangle++(0.04,0.04) (1.884,9.235)rectangle++(0.04,0.04)
  (1.969,9.371)rectangle++(0.04,0.04) (2.046,9.482)rectangle++(0.04,0.04) (2.116,9.574)rectangle++(0.04,0.04) (2.179,9.65)rectangle++(0.04,0.04) (2.236,9.711)rectangle++(0.04,0.04)
  (2.287,9.761)rectangle++(0.04,0.04) (2.332,9.802)rectangle++(0.04,0.04) (2.372,9.835)rectangle++(0.04,0.04) (2.408,9.861)rectangle++(0.04,0.04) (-2.103,4.312)rectangle++(0.04,0.04)
  (-1.618,5.346)rectangle++(0.04,0.04) (-1.163,6.205)rectangle++(0.04,0.04) (-0.743,6.917)rectangle++(0.04,0.04) (-0.357,7.505)rectangle++(0.04,0.04) (2.138,2.138)rectangle++(0.04,0.04)
  (1.9,3.328)rectangle++(0.04,0.04) (1.746,4.349)rectangle++(0.04,0.04) (1.656,5.223)rectangle++(0.04,0.04) (1.614,5.97)rectangle++(0.04,0.04) (1.609,6.606)rectangle++(0.04,0.04)
  (1.63,7.147)rectangle++(0.04,0.04) (1.669,7.606)rectangle++(0.04,0.04) (1.72,7.994)rectangle++(0.04,0.04) (1.78,8.323)rectangle++(0.04,0.04) (-1.829,3.829)rectangle++(0.04,0.04)
  (-1.404,4.924)rectangle++(0.04,0.04) (-0.998,5.838)rectangle++(0.04,0.04) (-1.739,2.644)rectangle++(0.04,0.04) (-1.374,3.913)rectangle++(0.04,0.04) (-1.014,4.977)rectangle++(0.04,0.04)
  (-0.665,5.867)rectangle++(0.04,0.04) (-0.333,6.61)rectangle++(0.04,0.04) (-0.021,7.228)rectangle++(0.04,0.04) (0.269,7.741)rectangle++(0.04,0.04) (0.536,8.165)rectangle++(0.04,0.04)
  (0.78,8.515)rectangle++(0.04,0.04) (1.002,8.803)rectangle++(0.04,0.04) (1.201,9.039)rectangle++(0.04,0.04) (1.381,9.231)rectangle++(0.04,0.04) (-2.145,3.937)rectangle++(0.04,0.04)
  (-1.668,5.029)rectangle++(0.04,0.04) (-1.219,5.937)rectangle++(0.04,0.04) (-0.801,6.692)rectangle++(0.04,0.04) (-0.415,7.316)rectangle++(0.04,0.04) (-0.062,7.832)rectangle++(0.04,0.04)
  (0.258,8.256)rectangle++(0.04,0.04) (0.547,8.603)rectangle++(0.04,0.04) (0.807,8.887)rectangle++(0.04,0.04) (1.039,9.118)rectangle++(0.04,0.04) (1.246,9.305)rectangle++(0.04,0.04)
  (-2.191,3.923)rectangle++(0.04,0.04) (-1.708,5.018)rectangle++(0.04,0.04) (1.644,1.19)rectangle++(0.04,0.04) (1.443,2.542)rectangle++(0.04,0.04) (1.326,3.699)rectangle++(0.04,0.04)
  (1.273,4.688)rectangle++(0.04,0.04) (1.267,5.53)rectangle++(0.04,0.04) (1.296,6.246)rectangle++(0.04,0.04) (-1.386,3.261)rectangle++(0.04,0.04) (-1.05,4.424)rectangle++(0.04,0.04)
  (-0.718,5.398)rectangle++(0.04,0.04) (-0.396,6.213)rectangle++(0.04,0.04) (1.782,1.818)rectangle++(0.04,0.04) (-0.138,2.399)rectangle++(0.04,0.04) (-0.023,3.641)rectangle++(0.04,0.04)
  (0.124,4.692)rectangle++(0.04,0.04) (0.291,5.579)rectangle++(0.04,0.04) (-1.414,2.883)rectangle++(0.04,0.04) (-1.088,4.104)rectangle++(0.04,0.04) (-1.306,2.241)rectangle++(0.04,0.04)
  (-1.022,3.554)rectangle++(0.04,0.04) (1.131,1.017)rectangle++(0.04,0.04) (1,2.415)rectangle++(0.04,0.04) (0.944,3.609)rectangle++(0.04,0.04) (0.945,4.626)rectangle++(0.04,0.04)
  (0.986,5.491)rectangle++(0.04,0.04) (1.055,6.224)rectangle++(0.04,0.04) (1.144,6.844)rectangle++(0.04,0.04) (1.244,7.368)rectangle++(0.04,0.04) (1.35,7.809)rectangle++(0.04,0.04)
  (1.458,8.18)rectangle++(0.04,0.04) (1.564,8.491)rectangle++(0.04,0.04) (-1.916,3.817)rectangle++(0.04,0.04) (-1.478,4.917)rectangle++(0.04,0.04) (-1.062,5.835)rectangle++(0.04,0.04)
  (-0.671,6.598)rectangle++(0.04,0.04) (-0.309,7.232)rectangle++(0.04,0.04) (0.024,7.755)rectangle++(0.04,0.04) (0.329,8.187)rectangle++(0.04,0.04) (2.226,2.48)rectangle++(0.04,0.04)
  (1.989,3.616)rectangle++(0.04,0.04) (-0.563,2.842)rectangle++(0.04,0.04) (-0.367,4.034)rectangle++(0.04,0.04) (-0.153,5.04)rectangle++(0.04,0.04) (0.069,5.886)rectangle++(0.04,0.04)
  (0.292,6.597)rectangle++(0.04,0.04) (0.51,7.192)rectangle++(0.04,0.04) (0.719,7.689)rectangle++(0.04,0.04) (0.916,8.103)rectangle++(0.04,0.04) (1.101,8.447)rectangle++(0.04,0.04)
  (1.271,8.732)rectangle++(0.04,0.04) (1.428,8.968)rectangle++(0.04,0.04) (1.57,9.162)rectangle++(0.04,0.04) (2.312,3.037)rectangle++(0.04,0.04) (2.085,4.085)rectangle++(0.04,0.04)
  (1.933,4.985)rectangle++(0.04,0.04) (1.088,2.129)rectangle++(0.04,0.04) (1.008,3.362)rectangle++(0.04,0.04) (0.989,4.414)rectangle++(0.04,0.04) (1.015,5.308)rectangle++(0.04,0.04)
  (1.073,6.068)rectangle++(0.04,0.04) (1.153,6.711)rectangle++(0.04,0.04) (1.246,7.254)rectangle++(0.04,0.04) (1.347,7.713)rectangle++(0.04,0.04) (1.451,8.098)rectangle++(0.04,0.04)
  (1.555,8.421)rectangle++(0.04,0.04) (2.107,2.856)rectangle++(0.04,0.04) (0.466,1.663)rectangle++(0.04,0.04) (0.378,0.95)rectangle++(0.04,0.04) (0.357,2.389)rectangle++(0.04,0.04)
  (0.397,3.612)rectangle++(0.04,0.04) (0.48,4.651)rectangle++(0.04,0.04) (0.592,5.53)rectangle++(0.04,0.04) (0.722,6.273)rectangle++(0.04,0.04) (0.862,6.9)rectangle++(0.04,0.04)
  (-1.643,3.305)rectangle++(0.04,0.04) (-1.266,4.471)rectangle++(0.04,0.04) (-0.9,5.448)rectangle++(0.04,0.04) (-1.617,2.581)rectangle++(0.04,0.04) (0.948,0.629)rectangle++(0.04,0.04)
  (0.005,1.945)rectangle++(0.04,0.04) (0.08,3.25)rectangle++(0.04,0.04) (0.196,4.355)rectangle++(0.04,0.04) (0.338,5.29)rectangle++(0.04,0.04) (1.413,1.788)rectangle++(0.04,0.04)
  (1.27,3.059)rectangle++(0.04,0.04) (1.2,4.146)rectangle++(0.04,0.04) (1.184,5.072)rectangle++(0.04,0.04) (1.207,5.86)rectangle++(0.04,0.04) (1.258,6.529)rectangle++(0.04,0.04)
  (1.328,7.096)rectangle++(0.04,0.04) (1.411,7.574)rectangle++(0.04,0.04) (1.5,7.978)rectangle++(0.04,0.04) (1.592,8.317)rectangle++(0.04,0.04) (1.683,8.602)rectangle++(0.04,0.04)
  (1.773,8.841)rectangle++(0.04,0.04) (1.858,9.04)rectangle++(0.04,0.04) (1.939,9.206)rectangle++(0.04,0.04) (2.014,9.344)rectangle++(0.04,0.04) (2.084,9.458)rectangle++(0.04,0.04)
  (2.147,9.552)rectangle++(0.04,0.04) (2.205,9.629)rectangle++(0.04,0.04) (2.257,9.693)rectangle++(0.04,0.04) (2.304,9.745)rectangle++(0.04,0.04) (2.346,9.787)rectangle++(0.04,0.04)
  (2.383,9.822)rectangle++(0.04,0.04) (2.417,9.849)rectangle++(0.04,0.04) (2.378,3.422)rectangle++(0.04,0.04) (0.584,1.87)rectangle++(0.04,0.04) (0.569,3.162)rectangle++(0.04,0.04)
  (0.608,4.261)rectangle++(0.04,0.04) (0.685,5.194)rectangle++(0.04,0.04) (1.334,1.855)rectangle++(0.04,0.04) (1.206,3.119)rectangle++(0.04,0.04) (1.148,4.199)rectangle++(0.04,0.04)
  (1.141,5.12)rectangle++(0.04,0.04) (1.173,5.902)rectangle++(0.04,0.04) (-1.321,3.157)rectangle++(0.04,0.04) (-0.999,4.333)rectangle++(0.04,0.04) (-0.678,5.319)rectangle++(0.04,0.04)
  (-0.366,6.144)rectangle++(0.04,0.04) (-0.067,6.834)rectangle++(0.04,0.04) (-0.02,1.077)rectangle++(0.04,0.04) (0.024,2.512)rectangle++(0.04,0.04) (0.119,3.731)rectangle++(0.04,0.04)
  (0.248,4.762)rectangle++(0.04,0.04) (0.399,5.634)rectangle++(0.04,0.04) (0.562,6.369)rectangle++(0.04,0.04) (0.73,6.988)rectangle++(0.04,0.04) (0.898,7.507)rectangle++(0.04,0.04)
  (1.062,7.941)rectangle++(0.04,0.04) (1.218,8.303)rectangle++(0.04,0.04) (1.365,8.605)rectangle++(0.04,0.04) (-1.986,3.796)rectangle++(0.04,0.04) (-1.538,4.902)rectangle++(0.04,0.04)
  (-1.114,5.825)rectangle++(0.04,0.04) (-0.716,6.592)rectangle++(0.04,0.04) (-0.347,7.228)rectangle++(0.04,0.04) (-0.008,7.754)rectangle++(0.04,0.04) (0.301,8.187)rectangle++(0.04,0.04)
  (0.581,8.543)rectangle++(0.04,0.04) (0.834,8.835)rectangle++(0.04,0.04) (1.06,9.072)rectangle++(0.04,0.04) (1.262,9.265)rectangle++(0.04,0.04) (1.441,9.421)rectangle++(0.04,0.04)
  (1.599,9.547)rectangle++(0.04,0.04) (1.739,9.647)rectangle++(0.04,0.04) (1.862,9.726)rectangle++(0.04,0.04) (1.969,9.789)rectangle++(0.04,0.04) (2.428,3.291)rectangle++(0.04,0.04)
  (0.54,1.851)rectangle++(0.04,0.04) (0.531,3.148)rectangle++(0.04,0.04) (0.784,1.324)rectangle++(0.04,0.04) (0.718,2.69)rectangle++(0.04,0.04) (0.715,3.854)rectangle++(0.04,0.04)
  (0.76,4.843)rectangle++(0.04,0.04) (0.838,5.683)rectangle++(0.04,0.04) (-1.331,3.032)rectangle++(0.04,0.04) (-1.076,1.95)rectangle++(0.04,0.04) (-0.839,3.297)rectangle++(0.04,0.04)
  (-0.583,4.432)rectangle++(0.04,0.04) (-0.321,5.387)rectangle++(0.04,0.04) (1.539,1.639)rectangle++(0.04,0.04) (1.371,2.928)rectangle++(0.04,0.04) (1.281,4.03)rectangle++(0.04,0.04)
  (1.248,4.971)rectangle++(0.04,0.04) (1.257,5.771)rectangle++(0.04,0.04) (-1.27,3.148)rectangle++(0.04,0.04) (-0.956,4.323)rectangle++(0.04,0.04) (-0.642,5.309)rectangle++(0.04,0.04)
  (-0.336,6.134)rectangle++(0.04,0.04) (-1.683,2.861)rectangle++(0.04,0.04) (-1.318,4.096)rectangle++(0.04,0.04) (-0.959,5.13)rectangle++(0.04,0.04) (-0.612,5.995)rectangle++(0.04,0.04)
  (-0.283,6.717)rectangle++(0.04,0.04) (0.026,7.317)rectangle++(0.04,0.04) (-1.918,3.205)rectangle++(0.04,0.04) (-1.505,4.397)rectangle++(0.04,0.04) (-1.105,5.394)rectangle++(0.04,0.04)
  (-0.726,6.225)rectangle++(0.04,0.04) (-0.37,6.917)rectangle++(0.04,0.04) (-0.04,7.49)rectangle++(0.04,0.04) (0.263,7.964)rectangle++(0.04,0.04) (0.54,8.355)rectangle++(0.04,0.04)
  (0.791,8.677)rectangle++(0.04,0.04) (1.017,8.94)rectangle++(0.04,0.04) (1.22,9.154)rectangle++(0.04,0.04) (1.401,9.329)rectangle++(0.04,0.04) (-2.166,3.964)rectangle++(0.04,0.04)
  (-1.685,5.052)rectangle++(0.04,0.04) (-1.232,5.958)rectangle++(0.04,0.04) (-0.811,6.71)rectangle++(0.04,0.04) (-0.424,7.332)rectangle++(0.04,0.04) (-0.069,7.845)rectangle++(0.04,0.04)
  (0.253,8.267)rectangle++(0.04,0.04) (0.544,8.613)rectangle++(0.04,0.04) (0.804,8.896)rectangle++(0.04,0.04) (-0.02,1.407)rectangle++(0.04,0.04) (0.037,2.793)rectangle++(0.04,0.04)
  (0.141,3.968)rectangle++(0.04,0.04) (0.276,4.964)rectangle++(0.04,0.04) (0.431,5.804)rectangle++(0.04,0.04) (0.597,6.513)rectangle++(0.04,0.04) (0.765,7.108)rectangle++(0.04,0.04)
  (0.933,7.607)rectangle++(0.04,0.04) (1.973,2.498)rectangle++(0.04,0.04) (1.775,3.641)rectangle++(0.04,0.04) (1.652,4.62)rectangle++(0.04,0.04) (1.587,5.457)rectangle++(0.04,0.04)
  (1.565,6.171)rectangle++(0.04,0.04) (1.575,6.779)rectangle++(0.04,0.04) (1.607,7.295)rectangle++(0.04,0.04) (1.656,7.733)rectangle++(0.04,0.04) (-1.701,3.671)rectangle++(0.04,0.04)
  (-1.301,4.785)rectangle++(0.04,0.04) (-0.917,5.715)rectangle++(0.04,0.04) (1.72,1.563)rectangle++(0.04,0.04) (1.523,2.856)rectangle++(0.04,0.04) (1.406,3.963)rectangle++(0.04,0.04)
  (1.352,4.909)rectangle++(0.04,0.04) (1.343,5.714)rectangle++(0.04,0.04) (1.368,6.4)rectangle++(0.04,0.04) (1.417,6.981)rectangle++(0.04,0.04) (1.481,7.474)rectangle++(0.04,0.04)
  (-0.02,1.179)rectangle++(0.04,0.04) (0.028,2.599)rectangle++(0.04,0.04) (-0.691,2.167)rectangle++(0.04,0.04) (-0.503,3.466)rectangle++(0.04,0.04) (-1.023,2.236)rectangle++(0.04,0.04)
  (-0.782,3.538)rectangle++(0.04,0.04) (-0.526,4.634)rectangle++(0.04,0.04) (-1.331,2.488)rectangle++(0.04,0.04) (-1.034,3.764)rectangle++(0.04,0.04) (-0.731,4.837)rectangle++(0.04,0.04)
  (-0.43,5.737)rectangle++(0.04,0.04) (-0.138,6.49)rectangle++(0.04,0.04) (1.82,1.952)rectangle++(0.04,0.04) (-0.165,2.437)rectangle++(0.04,0.04) (-0.045,3.674)rectangle++(0.04,0.04)
  (0.107,4.721)rectangle++(0.04,0.04) (0.277,5.605)rectangle++(0.04,0.04) (0.458,6.349)rectangle++(0.04,0.04) (0.641,6.975)rectangle++(0.04,0.04) (0.822,7.499)rectangle++(0.04,0.04)
  (0.996,7.938)rectangle++(0.04,0.04) (1.162,8.303)rectangle++(0.04,0.04) (1.318,8.608)rectangle++(0.04,0.04) (2.195,2.838)rectangle++(0.04,0.04) (1.977,3.921)rectangle++(0.04,0.04)
  (1.835,4.85)rectangle++(0.04,0.04) (1.752,5.645)rectangle++(0.04,0.04) (1.713,6.325)rectangle++(0.04,0.04) (1.706,6.904)rectangle++(0.04,0.04) (1.724,7.396)rectangle++(0.04,0.04)
  (1.759,7.814)rectangle++(0.04,0.04) (-0.02,1.233)rectangle++(0.04,0.04) (0.03,2.645)rectangle++(0.04,0.04) (0.129,3.844)rectangle++(0.04,0.04) (0.261,4.858)rectangle++(0.04,0.04)
  (0.414,5.715)rectangle++(0.04,0.04) (0.579,6.437)rectangle++(0.04,0.04) (0.747,7.045)rectangle++(0.04,0.04) (0.915,7.554)rectangle++(0.04,0.04) (1.077,7.981)rectangle++(0.04,0.04)
  (1.233,8.337)rectangle++(0.04,0.04) (1.379,8.633)rectangle++(0.04,0.04) (1.515,8.879)rectangle++(0.04,0.04) (1.641,9.083)rectangle++(0.04,0.04) (-2.055,3.965)rectangle++(0.04,0.04)
  (-1.59,5.048)rectangle++(0.04,0.04) (-1.152,5.951)rectangle++(0.04,0.04) (-0.743,6.701)rectangle++(0.04,0.04) (-0.366,7.321)rectangle++(0.04,0.04) (-0.02,7.834)rectangle++(0.04,0.04)
  (0.294,8.256)rectangle++(0.04,0.04) (0.578,8.602)rectangle++(0.04,0.04) (0.833,8.885)rectangle++(0.04,0.04) (1.061,9.115)rectangle++(0.04,0.04) (1.264,9.301)rectangle++(0.04,0.04)
  (1.445,9.452)rectangle++(0.04,0.04) (1.604,9.572)rectangle++(0.04,0.04) (1.744,9.669)rectangle++(0.04,0.04) (1.867,9.745)rectangle++(0.04,0.04) (1.974,9.805)rectangle++(0.04,0.04)
  (2.432,3.296)rectangle++(0.04,0.04) (2.197,4.301)rectangle++(0.04,0.04) (2.037,5.164)rectangle++(0.04,0.04) (1.936,5.904)rectangle++(0.04,0.04) (1.879,6.537)rectangle++(0.04,0.04)
  (1.857,7.078)rectangle++(0.04,0.04) (1.859,7.538)rectangle++(0.04,0.04) (1.88,7.929)rectangle++(0.04,0.04) (1.913,8.261)rectangle++(0.04,0.04) (1.954,8.541)rectangle++(0.04,0.04)
  (2,8.778)rectangle++(0.04,0.04) (2.049,8.978)rectangle++(0.04,0.04) (2.099,9.145)rectangle++(0.04,0.04) (2.148,9.286)rectangle++(0.04,0.04) (2.195,9.403)rectangle++(0.04,0.04)
  (2.239,9.501)rectangle++(0.04,0.04) (2.281,9.583)rectangle++(0.04,0.04) (2.32,9.65)rectangle++(0.04,0.04) (2.356,9.706)rectangle++(0.04,0.04) (2.389,9.752)rectangle++(0.04,0.04)
  (2.418,9.79)rectangle++(0.04,0.04) (2.361,3.408)rectangle++(0.04,0.04) (2.141,4.399)rectangle++(0.04,0.04) (1.994,5.25)rectangle++(0.04,0.04) (1.902,5.979)rectangle++(0.04,0.04)
  (1.854,6.602)rectangle++(0.04,0.04) (1.838,7.134)rectangle++(0.04,0.04) (1.845,7.586)rectangle++(0.04,0.04) (1.87,7.971)rectangle++(0.04,0.04) (1.934,2.829)rectangle++(0.04,0.04)
  (1.755,3.924)rectangle++(0.04,0.04) (-0.69,2.856)rectangle++(0.04,0.04) (-0.475,4.051)rectangle++(0.04,0.04) (-0.244,5.059)rectangle++(0.04,0.04) (-0.007,5.906)rectangle++(0.04,0.04)
  (0.228,6.616)rectangle++(0.04,0.04) (0.456,7.211)rectangle++(0.04,0.04) (0.674,7.707)rectangle++(0.04,0.04) (0.879,8.12)rectangle++(0.04,0.04) (-1.957,3.578)rectangle++(0.04,0.04)
  (-1.522,4.716)rectangle++(0.04,0.04) (-1.108,5.665)rectangle++(0.04,0.04) (-0.717,6.456)rectangle++(0.04,0.04) (-0.353,7.112)rectangle++(0.04,0.04) (-0.018,7.656)rectangle++(0.04,0.04)
  (0.289,8.104)rectangle++(0.04,0.04) (0.567,8.473)rectangle++(0.04,0.04) (0.819,8.776)rectangle++(0.04,0.04) (1.045,9.023)rectangle++(0.04,0.04) (1.247,9.224)rectangle++(0.04,0.04)
  (2.378,2.968)rectangle++(0.04,0.04) (2.138,4.024)rectangle++(0.04,0.04) (1.976,4.931)rectangle++(0.04,0.04) (1.067,2.127)rectangle++(0.04,0.04) (0.99,3.362)rectangle++(0.04,0.04)
  (0.974,4.414)rectangle++(0.04,0.04) (1.002,5.309)rectangle++(0.04,0.04) (1.062,6.069)rectangle++(0.04,0.04) (1.143,6.712)rectangle++(0.04,0.04) (1.238,7.256)rectangle++(0.04,0.04)
  (1.34,7.714)rectangle++(0.04,0.04) (-1.759,3.594)rectangle++(0.04,0.04) (-1.353,4.722)rectangle++(0.04,0.04) (-0.964,5.664)rectangle++(0.04,0.04) (-0.595,6.449)rectangle++(0.04,0.04)
  (1.878,1.823)rectangle++(0.04,0.04) (1.667,3.071)rectangle++(0.04,0.04) (1.537,4.14)rectangle++(0.04,0.04) (1.47,5.053)rectangle++(0.04,0.04) (1.449,5.833)rectangle++(0.04,0.04)
  (1.398,2.207)rectangle++(0.04,0.04) (1.275,3.416)rectangle++(0.04,0.04) (1.218,4.449)rectangle++(0.04,0.04) (1.211,5.329)rectangle++(0.04,0.04) (1.24,6.077)rectangle++(0.04,0.04)
  (1.295,6.712)rectangle++(0.04,0.04) (1.367,7.25)rectangle++(0.04,0.04) (1.45,7.704)rectangle++(0.04,0.04) (1.538,8.087)rectangle++(0.04,0.04) (1.629,8.408)rectangle++(0.04,0.04)
  (-1.882,3.813)rectangle++(0.04,0.04) (-1.449,4.913)rectangle++(0.04,0.04) (-1.037,5.83)rectangle++(0.04,0.04) (-0.02,0.916)rectangle++(0.04,0.04) (0.017,2.376)rectangle++(0.04,0.04)
  (0.108,3.615)rectangle++(0.04,0.04) (0.234,4.664)rectangle++(0.04,0.04) (0.383,5.552)rectangle++(0.04,0.04) (1.48,1.862)rectangle++(0.04,0.04) (1.33,3.12)rectangle++(0.04,0.04)
  (1.253,4.195)rectangle++(0.04,0.04) (1.231,5.112)rectangle++(0.04,0.04) (1.248,5.892)rectangle++(0.04,0.04) (1.295,6.554)rectangle++(0.04,0.04) (1.36,7.116)rectangle++(0.04,0.04)
  (1.439,7.59)rectangle++(0.04,0.04) (1.524,7.99)rectangle++(0.04,0.04) (1.613,8.327)rectangle++(0.04,0.04) (1.702,8.61)rectangle++(0.04,0.04) (1.789,8.846)rectangle++(0.04,0.04)
  (1.872,9.044)rectangle++(0.04,0.04) (1.951,9.209)rectangle++(0.04,0.04) (2.024,9.346)rectangle++(0.04,0.04) (2.092,9.459)rectangle++(0.04,0.04) (2.155,9.553)rectangle++(0.04,0.04)
  (2.211,9.63)rectangle++(0.04,0.04) (2.263,9.693)rectangle++(0.04,0.04) (2.309,9.745)rectangle++(0.04,0.04) (2.35,9.787)rectangle++(0.04,0.04) (2.387,9.821)rectangle++(0.04,0.04)
  (2.374,3.408)rectangle++(0.04,0.04) (2.152,4.398)rectangle++(0.04,0.04) (0.891,2.045)rectangle++(0.04,0.04) (0.837,3.299)rectangle++(0.04,0.04) (0.841,4.367)rectangle++(0.04,0.04)
  (-0.988,2.743)rectangle++(0.04,0.04) (-0.733,3.967)rectangle++(0.04,0.04) (-0.466,4.998)rectangle++(0.04,0.04) (-0.199,5.863)rectangle++(0.04,0.04) (1.654,1.786)rectangle++(0.04,0.04)
  (1.475,3.048)rectangle++(0.04,0.04) (1.374,4.128)rectangle++(0.04,0.04) (1.33,5.05)rectangle++(0.04,0.04) (1.331,5.835)rectangle++(0.04,0.04) (-1.272,3.179)rectangle++(0.04,0.04)
  (-1.102,1.996)rectangle++(0.04,0.04) (-0.859,3.337)rectangle++(0.04,0.04) (-0.599,4.467)rectangle++(0.04,0.04) (-0.333,5.417)rectangle++(0.04,0.04) (-0.068,6.214)rectangle++(0.04,0.04)
  (0.188,6.881)rectangle++(0.04,0.04) (0.433,7.437)rectangle++(0.04,0.04) (2,2.328)rectangle++(0.04,0.04) (1.791,3.495)rectangle++(0.04,0.04) (1.66,4.495)rectangle++(0.04,0.04)
  (1.589,5.351)rectangle++(0.04,0.04) (1.562,6.081)rectangle++(0.04,0.04) (1.569,6.702)rectangle++(0.04,0.04) (1.599,7.23)rectangle++(0.04,0.04) (1.646,7.678)rectangle++(0.04,0.04)
  (1.704,8.057)rectangle++(0.04,0.04) (1.769,8.376)rectangle++(0.04,0.04) (1.836,8.645)rectangle++(0.04,0.04) (-1.902,3.913)rectangle++(0.04,0.04) (-1.462,4.999)rectangle++(0.04,0.04)
  (-1.045,5.903)rectangle++(0.04,0.04) (-0.654,6.656)rectangle++(0.04,0.04) (-0.292,7.28)rectangle++(0.04,0.04) (2.065,2.101)rectangle++(0.04,0.04) (-0.155,2.526)rectangle++(0.04,0.04)
  (-0.02,0.387)rectangle++(0.04,0.04) (-0.004,1.926)rectangle++(0.04,0.04) (0.072,3.234)rectangle++(0.04,0.04) (-0.848,2.317)rectangle++(0.04,0.04) (-0.63,3.599)rectangle++(0.04,0.04)
  (-0.394,4.681)rectangle++(0.04,0.04) (-0.15,5.591)rectangle++(0.04,0.04) (0.094,6.354)rectangle++(0.04,0.04) (0.332,6.994)rectangle++(0.04,0.04) (0.56,7.527)rectangle++(0.04,0.04)
  (0.775,7.972)rectangle++(0.04,0.04) (0.975,8.342)rectangle++(0.04,0.04) (1.16,8.648)rectangle++(0.04,0.04) (1.33,8.9)rectangle++(0.04,0.04) (1.484,9.108)rectangle++(0.04,0.04)
  (1.624,9.279)rectangle++(0.04,0.04) (1.749,9.418)rectangle++(0.04,0.04) (1.861,9.532)rectangle++(0.04,0.04) (1.961,9.624)rectangle++(0.04,0.04) (2.05,9.698)rectangle++(0.04,0.04)
  (2.128,9.757)rectangle++(0.04,0.04) (2.197,9.805)rectangle++(0.04,0.04) (-2.131,4.251)rectangle++(0.04,0.04) (-1.644,5.295)rectangle++(0.04,0.04) (-1.187,6.163)rectangle++(0.04,0.04)
  (-0.765,6.882)rectangle++(0.04,0.04) (-0.377,7.477)rectangle++(0.04,0.04) (-0.024,7.966)rectangle++(0.04,0.04) (-2.097,3.336)rectangle++(0.04,0.04) (-1.308,1.841)rectangle++(0.04,0.04)
  (-1.04,3.213)rectangle++(0.04,0.04) (-0.758,4.369)rectangle++(0.04,0.04) (1.32,1.281)rectangle++(0.04,0.04) (1.171,2.633)rectangle++(0.04,0.04) (1.098,3.787)rectangle++(0.04,0.04)
  (1.083,4.771)rectangle++(0.04,0.04) (1.109,5.609)rectangle++(0.04,0.04) (1.165,6.319)rectangle++(0.04,0.04) (1.241,6.921)rectangle++(0.04,0.04) (1.329,7.429)rectangle++(0.04,0.04)
  (1.425,7.858)rectangle++(0.04,0.04) (1.523,8.218)rectangle++(0.04,0.04) (1.621,8.521)rectangle++(0.04,0.04) (1.717,8.774)rectangle++(0.04,0.04) (1.808,8.986)rectangle++(0.04,0.04)
  (1.894,9.162)rectangle++(0.04,0.04) (1.974,9.308)rectangle++(0.04,0.04) (2.293,3.177)rectangle++(0.04,0.04) (-0.02,0.492)rectangle++(0.04,0.04) (0,2.015)rectangle++(0.04,0.04)
  (0.079,3.309)rectangle++(0.04,0.04) (0.197,4.406)rectangle++(0.04,0.04) (0.342,5.333)rectangle++(0.04,0.04) (0.501,6.116)rectangle++(0.04,0.04) (0.669,6.774)rectangle++(0.04,0.04)
  (0.837,7.328)rectangle++(0.04,0.04) (1.909,2.406)rectangle++(0.04,0.04) (1.717,3.565)rectangle++(0.04,0.04) (1.599,4.558)rectangle++(0.04,0.04) (1.54,5.406)rectangle++(0.04,0.04)
  (1.523,6.13)rectangle++(0.04,0.04) (1.537,6.746)rectangle++(0.04,0.04) (1.574,7.269)rectangle++(0.04,0.04) (-1.596,3.55)rectangle++(0.04,0.04) (1.216,0.867)rectangle++(0.04,0.04)
  (1.066,2.285)rectangle++(0.04,0.04) (0.995,3.495)rectangle++(0.04,0.04) (0.984,4.527)rectangle++(0.04,0.04) (1.015,5.405)rectangle++(0.04,0.04) (1.077,6.15)rectangle++(0.04,0.04)
  (1.159,6.781)rectangle++(0.04,0.04) (1.254,7.313)rectangle++(0.04,0.04) (1.356,7.762)rectangle++(0.04,0.04) (1.461,8.14)rectangle++(0.04,0.04) (1.565,8.457)rectangle++(0.04,0.04)
  (1.667,8.722)rectangle++(0.04,0.04) (1.763,8.943)rectangle++(0.04,0.04) (1.854,9.127)rectangle++(0.04,0.04) (1.939,9.28)rectangle++(0.04,0.04) (2.29,3.161)rectangle++(0.04,0.04)
  (0.524,1.784)rectangle++(0.04,0.04) (0.404,0.995)rectangle++(0.04,0.04) (0.201,0.774)rectangle++(0.04,0.04) (0.199,2.246)rectangle++(0.04,0.04) (0.257,3.497)rectangle++(0.04,0.04)
  (0.356,4.558)rectangle++(0.04,0.04) (0.483,5.457)rectangle++(0.04,0.04) (0.626,6.215)rectangle++(0.04,0.04) (0.779,6.854)rectangle++(0.04,0.04) (0.934,7.391)rectangle++(0.04,0.04)
  (1.087,7.841)rectangle++(0.04,0.04) (1.236,8.218)rectangle++(0.04,0.04) (1.377,8.532)rectangle++(0.04,0.04) (1.509,8.793)rectangle++(0.04,0.04) (1.633,9.01)rectangle++(0.04,0.04)
  (1.746,9.189)rectangle++(0.04,0.04) (1.849,9.337)rectangle++(0.04,0.04) (1.943,9.459)rectangle++(0.04,0.04) (2.028,9.559)rectangle++(0.04,0.04) (2.104,9.64)rectangle++(0.04,0.04)
  (2.172,9.706)rectangle++(0.04,0.04) (2.232,9.759)rectangle++(0.04,0.04) (2.285,9.802)rectangle++(0.04,0.04) (2.332,9.837)rectangle++(0.04,0.04) (2.374,9.864)rectangle++(0.04,0.04)
  (2.41,9.886)rectangle++(0.04,0.04) (-0.02,1.565)rectangle++(0.04,0.04) (0.043,2.927)rectangle++(0.04,0.04) (0.152,4.083)rectangle++(0.04,0.04) (0.29,5.06)rectangle++(0.04,0.04)
  (0.447,5.886)rectangle++(0.04,0.04) (1.564,1.959)rectangle++(0.04,0.04) (1.405,3.199)rectangle++(0.04,0.04) (1.32,4.259)rectangle++(0.04,0.04) (1.29,5.163)rectangle++(0.04,0.04)
  (-1.106,3.022)rectangle++(0.04,0.04) (-0.821,4.209)rectangle++(0.04,0.04) (-0.532,5.207)rectangle++(0.04,0.04) (-0.246,6.043)rectangle++(0.04,0.04) (0.03,6.743)rectangle++(0.04,0.04)
  (0.293,7.326)rectangle++(0.04,0.04) (0.54,7.812)rectangle++(0.04,0.04) (0.769,8.215)rectangle++(0.04,0.04) (0.98,8.548)rectangle++(0.04,0.04) (1.173,8.823)rectangle++(0.04,0.04)
  (2.277,2.852)rectangle++(0.04,0.04) (2.047,3.93)rectangle++(0.04,0.04) (1.895,4.855)rectangle++(0.04,0.04) (1.803,5.647)rectangle++(0.04,0.04) (1.756,6.324)rectangle++(0.04,0.04)
  (1.743,6.901)rectangle++(0.04,0.04) (1.756,7.392)rectangle++(0.04,0.04) (1.786,7.81)rectangle++(0.04,0.04) (1.828,8.163)rectangle++(0.04,0.04) (1.878,8.462)rectangle++(0.04,0.04)
  (1.933,8.713)rectangle++(0.04,0.04) (-1.9,3.95)rectangle++(0.04,0.04) (-1.459,5.03)rectangle++(0.04,0.04) (-1.041,5.93)rectangle++(0.04,0.04) (-0.65,6.679)rectangle++(0.04,0.04)
  (-0.288,7.299)rectangle++(0.04,0.04) (0.045,7.812)rectangle++(0.04,0.04) (0.349,8.234)rectangle++(0.04,0.04) (2.236,2.497)rectangle++(0.04,0.04) (1.998,3.629)rectangle++(0.04,0.04)
  (1.841,4.601)rectangle++(0.04,0.04) (1.747,5.433)rectangle++(0.04,0.04) (1.7,6.145)rectangle++(0.04,0.04) (1.689,6.751)rectangle++(0.04,0.04) (1.703,7.267)rectangle++(0.04,0.04)
  (1.736,7.705)rectangle++(0.04,0.04) (1.782,8.076)rectangle++(0.04,0.04) (1.835,8.39)rectangle++(0.04,0.04) (1.893,8.654)rectangle++(0.04,0.04) (1.953,8.876)rectangle++(0.04,0.04)
  (2.013,9.063)rectangle++(0.04,0.04) (2.072,9.219)rectangle++(0.04,0.04) (2.127,9.35)rectangle++(0.04,0.04) (2.18,9.458)rectangle++(0.04,0.04) (2.229,9.549)rectangle++(0.04,0.04)
  (2.275,9.623)rectangle++(0.04,0.04) (2.316,9.685)rectangle++(0.04,0.04) (2.354,9.736)rectangle++(0.04,0.04) (2.388,9.777)rectangle++(0.04,0.04) (2.419,9.812)rectangle++(0.04,0.04)
  (2.446,9.839)rectangle++(0.04,0.04) (2.471,9.862)rectangle++(0.04,0.04) (2.492,9.88)rectangle++(0.04,0.04) (2.375,3.449)rectangle++(0.04,0.04) (2.155,4.433)rectangle++(0.04,0.04)
  (2.007,5.278)rectangle++(0.04,0.04) (1.159,2.218)rectangle++(0.04,0.04) (1.072,3.435)rectangle++(0.04,0.04) (1.046,4.473)rectangle++(0.04,0.04) (1.066,5.357)rectangle++(0.04,0.04)
  (1.118,6.107)rectangle++(0.04,0.04) (1.193,6.742)rectangle++(0.04,0.04) (1.281,7.279)rectangle++(0.04,0.04) (1.378,7.732)rectangle++(0.04,0.04) (1.478,8.114)rectangle++(0.04,0.04)
  (1.579,8.434)rectangle++(0.04,0.04) (1.677,8.702)rectangle++(0.04,0.04) (1.772,8.926)rectangle++(0.04,0.04) (1.861,9.112)rectangle++(0.04,0.04) (1.944,9.267)rectangle++(0.04,0.04)
  (2.021,9.395)rectangle++(0.04,0.04) (2.31,3.21)rectangle++(0.04,0.04) (2.09,4.233)rectangle++(0.04,0.04) (1.943,5.11)rectangle++(0.04,0.04) (1.122,2.162)rectangle++(0.04,0.04)
  (1.038,3.389)rectangle++(0.04,0.04) (1.016,4.435)rectangle++(0.04,0.04) (1.038,5.325)rectangle++(0.04,0.04) (1.094,6.081)rectangle++(0.04,0.04) (1.171,6.722)rectangle++(0.04,0.04)
  (1.262,7.263)rectangle++(0.04,0.04) (1.361,7.719)rectangle++(0.04,0.04) (1.463,8.103)rectangle++(0.04,0.04) (-0.02,1.28)rectangle++(0.04,0.04) (0.032,2.685)rectangle++(0.04,0.04)
  (0.132,3.877)rectangle++(0.04,0.04) (0.265,4.886)rectangle++(0.04,0.04) (0.419,5.739)rectangle++(0.04,0.04) (1.527,1.916)rectangle++(0.04,0.04) (1.372,3.164)rectangle++(0.04,0.04)
  (1.291,4.231)rectangle++(0.04,0.04) (0.974,1.781)rectangle++(0.04,0.04) (0.335,1.111)rectangle++(0.04,0.04) (-0.243,1.91)rectangle++(0.04,0.04) (-0.566,1.953)rectangle++(0.04,0.04)
  (0.615,0.752)rectangle++(0.04,0.04) (0.55,2.21)rectangle++(0.04,0.04) (0.554,3.453)rectangle++(0.04,0.04) (0.866,1.403)rectangle++(0.04,0.04) (0.79,2.754)rectangle++(0.04,0.04)
  (0.78,3.905)rectangle++(0.04,0.04) (0.817,4.885)rectangle++(0.04,0.04) (0.887,5.715)rectangle++(0.04,0.04) (0.981,6.419)rectangle++(0.04,0.04) (1.088,7.013)rectangle++(0.04,0.04)
  (1.203,7.514)rectangle++(0.04,0.04) (1.321,7.935)rectangle++(0.04,0.04) (1.438,8.288)rectangle++(0.04,0.04) (-1.888,3.743)rectangle++(0.04,0.04) (-1.458,4.853)rectangle++(0.04,0.04)
  (-1.047,5.78)rectangle++(0.04,0.04) (-0.661,6.551)rectangle++(0.04,0.04) (-0.302,7.191)rectangle++(0.04,0.04) (0.029,7.721)rectangle++(0.04,0.04) (0.331,8.158)rectangle++(0.04,0.04)
  (0.605,8.517)rectangle++(0.04,0.04) (0.853,8.811)rectangle++(0.04,0.04) (1.075,9.052)rectangle++(0.04,0.04) (1.274,9.247)rectangle++(0.04,0.04) (1.451,9.405)rectangle++(0.04,0.04)
  (1.607,9.533)rectangle++(0.04,0.04) (1.745,9.635)rectangle++(0.04,0.04) (1.866,9.716)rectangle++(0.04,0.04) (1.973,9.78)rectangle++(0.04,0.04) (2.066,9.83)rectangle++(0.04,0.04)
  (-2.164,4.227)rectangle++(0.04,0.04) (-1.672,5.276)rectangle++(0.04,0.04) (-1.213,6.147)rectangle++(0.04,0.04) (-0.787,6.87)rectangle++(0.04,0.04) (2.024,1.874)rectangle++(0.04,0.04)
  (1.793,3.108)rectangle++(0.04,0.04) (0.584,1.642)rectangle++(0.04,0.04) (0.355,0.976)rectangle++(0.04,0.04) (0.203,0.756)rectangle++(0.04,0.04) (0.2,2.231)rectangle++(0.04,0.04)
  (0.257,3.485)rectangle++(0.04,0.04) (0.356,4.548)rectangle++(0.04,0.04) (0.482,5.448)rectangle++(0.04,0.04) (0.626,6.207)rectangle++(0.04,0.04) (0.778,6.847)rectangle++(0.04,0.04)
  (0.933,7.385)rectangle++(0.04,0.04) (1.086,7.837)rectangle++(0.04,0.04) (1.234,8.214)rectangle++(0.04,0.04) (1.376,8.529)rectangle++(0.04,0.04) (1.508,8.79)rectangle++(0.04,0.04)
  (1.631,9.008)rectangle++(0.04,0.04) (1.745,9.188)rectangle++(0.04,0.04) (1.848,9.336)rectangle++(0.04,0.04) (-2.079,4.068)rectangle++(0.04,0.04) (-1.495,2.006)rectangle++(0.04,0.04)
  (-1.192,3.361)rectangle++(0.04,0.04) (-0.881,4.501)rectangle++(0.04,0.04) (-0.571,5.457)rectangle++(0.04,0.04) (-0.27,6.258)rectangle++(0.04,0.04) (0.019,6.926)rectangle++(0.04,0.04)
  (0.291,7.482)rectangle++(0.04,0.04) (0.544,7.945)rectangle++(0.04,0.04) (0.778,8.327)rectangle++(0.04,0.04) (0.993,8.643)rectangle++(0.04,0.04) (1.187,8.903)rectangle++(0.04,0.04)
  (1.363,9.117)rectangle++(0.04,0.04) (1.521,9.291)rectangle++(0.04,0.04) (1.662,9.432)rectangle++(0.04,0.04) (1.788,9.547)rectangle++(0.04,0.04) (1.9,9.64)rectangle++(0.04,0.04)
  (-2.148,4.147)rectangle++(0.04,0.04) (-1.662,5.207)rectangle++(0.04,0.04) (-1.207,6.088)rectangle++(0.04,0.04) (-0.784,6.82)rectangle++(0.04,0.04) (-0.396,7.424)rectangle++(0.04,0.04)
  (-0.042,7.923)rectangle++(0.04,0.04) (-2.089,3.322)rectangle++(0.04,0.04) (-1.645,4.504)rectangle++(0.04,0.04) (-1.221,5.49)rectangle++(0.04,0.04) (-0.82,6.312)rectangle++(0.04,0.04)
  (-0.447,6.994)rectangle++(0.04,0.04) (-0.102,7.559)rectangle++(0.04,0.04) (0.213,8.025)rectangle++(0.04,0.04) (2.198,2.412)rectangle++(0.04,0.04) (1.962,3.558)rectangle++(0.04,0.04)
  (1.808,4.542)rectangle++(0.04,0.04) (1.716,5.385)rectangle++(0.04,0.04) (1.672,6.105)rectangle++(0.04,0.04) (1.663,6.718)rectangle++(0.04,0.04) (1.68,7.24)rectangle++(0.04,0.04)
  (1.758,2.605)rectangle++(0.04,0.04) (-0.347,2.566)rectangle++(0.04,0.04) (-0.194,3.791)rectangle++(0.04,0.04) (-0.016,4.827)rectangle++(0.04,0.04) (0.177,5.7)rectangle++(0.04,0.04)
  (0.377,6.434)rectangle++(0.04,0.04) (0.575,7.05)rectangle++(0.04,0.04) (1.87,2.272)rectangle++(0.04,0.04) (1.678,3.452)rectangle++(0.04,0.04) (1.562,4.463)rectangle++(0.04,0.04)
  (1.504,5.328)rectangle++(0.04,0.04) (1.49,6.064)rectangle++(0.04,0.04) (1.507,6.691)rectangle++(0.04,0.04) (1.546,7.224)rectangle++(0.04,0.04) (1.601,7.674)rectangle++(0.04,0.04)
  (1.666,8.055)rectangle++(0.04,0.04) (1.988,2.796)rectangle++(0.04,0.04) (1.8,3.894)rectangle++(0.04,0.04) (1.683,4.834)rectangle++(0.04,0.04) (1.622,5.638)rectangle++(0.04,0.04)
  (1.602,6.323)rectangle++(0.04,0.04) (1.612,6.907)rectangle++(0.04,0.04) (1.645,7.403)rectangle++(0.04,0.04) (1.692,7.823)rectangle++(0.04,0.04) (1.749,8.178)rectangle++(0.04,0.04)
  (1.811,8.477)rectangle++(0.04,0.04) (1.877,8.729)rectangle++(0.04,0.04) (1.942,8.941)rectangle++(0.04,0.04) (2.195,3.081)rectangle++(0.04,0.04) (1.987,4.127)rectangle++(0.04,0.04)
  (1.852,5.025)rectangle++(0.04,0.04) (1.112,2.117)rectangle++(0.04,0.04) (1.028,3.351)rectangle++(0.04,0.04) (1.005,4.404)rectangle++(0.04,0.04) (1.028,5.299)rectangle++(0.04,0.04)
  (1.084,6.059)rectangle++(0.04,0.04) (1.162,6.703)rectangle++(0.04,0.04) (1.253,7.248)rectangle++(0.04,0.04) (1.353,7.707)rectangle++(0.04,0.04) (1.456,8.093)rectangle++(0.04,0.04)
  (2.03,2.751)rectangle++(0.04,0.04) (1.833,3.853)rectangle++(0.04,0.04) (1.71,4.798)rectangle++(0.04,0.04) (1.644,5.606)rectangle++(0.04,0.04) (1.619,6.296)rectangle++(0.04,0.04)
  (1.626,6.883)rectangle++(0.04,0.04) (1.666,2.505)rectangle++(0.04,0.04) (1.514,3.658)rectangle++(0.04,0.04) (0.78,1.702)rectangle++(0.04,0.04) (0.729,3.011)rectangle++(0.04,0.04)
  (0.738,4.127)rectangle++(0.04,0.04) (0.79,5.074)rectangle++(0.04,0.04) (0.872,5.878)rectangle++(0.04,0.04) (-1.375,3.083)rectangle++(0.04,0.04) (-1.048,4.272)rectangle++(0.04,0.04)
  (-0.722,5.269)rectangle++(0.04,0.04) (-0.405,6.104)rectangle++(0.04,0.04) (-0.102,6.8)rectangle++(0.04,0.04) (0.183,7.381)rectangle++(0.04,0.04) (0.448,7.862)rectangle++(0.04,0.04)
  (0.693,8.261)rectangle++(0.04,0.04) (-2.03,3.566)rectangle++(0.04,0.04) (-1.585,4.709)rectangle++(0.04,0.04) (1.539,1.148)rectangle++(0.04,0.04) (1.352,2.51)rectangle++(0.04,0.04)
  (1.247,3.676)rectangle++(0.04,0.04) (1.205,4.671)rectangle++(0.04,0.04) (1.209,5.518)rectangle++(0.04,0.04) (1.346,2.069)rectangle++(0.04,0.04) (1.225,3.301)rectangle++(0.04,0.04)
  (1.171,4.353)rectangle++(0.04,0.04) (1.167,5.249)rectangle++(0.04,0.04) (1.2,6.011)rectangle++(0.04,0.04) (1.258,6.658)rectangle++(0.04,0.04) (1.334,7.205)rectangle++(0.04,0.04)
  (1.42,7.667)rectangle++(0.04,0.04) (1.511,8.056)rectangle++(0.04,0.04) (-1.814,3.709)rectangle++(0.04,0.04) (-1.395,4.821)rectangle++(0.04,0.04) (-0.995,5.75)rectangle++(0.04,0.04)
  (-0.618,6.524)rectangle++(0.04,0.04) (-0.267,7.166)rectangle++(0.04,0.04) (0.058,7.698)rectangle++(0.04,0.04) (-2.011,3.296)rectangle++(0.04,0.04) (-1.58,4.478)rectangle++(0.04,0.04)
  (-1.166,5.466)rectangle++(0.04,0.04) (-0.775,6.289)rectangle++(0.04,0.04) (-0.409,6.973)rectangle++(0.04,0.04) (-0.071,7.539)rectangle++(0.04,0.04) (0.239,8.007)rectangle++(0.04,0.04)
  (2.189,2.414)rectangle++(0.04,0.04) (1.955,3.56)rectangle++(0.04,0.04) (1.802,4.544)rectangle++(0.04,0.04) (1.711,5.387)rectangle++(0.04,0.04) (1.668,6.107)rectangle++(0.04,0.04)
  (1.66,6.72)rectangle++(0.04,0.04) (1.677,7.242)rectangle++(0.04,0.04) (-1.569,3.568)rectangle++(0.04,0.04) (-1.193,4.692)rectangle++(0.04,0.04) (-0.828,5.632)rectangle++(0.04,0.04)
  (-0.481,6.416)rectangle++(0.04,0.04) (-0.154,7.069)rectangle++(0.04,0.04) (0.149,7.611)rectangle++(0.04,0.04) (0.429,8.06)rectangle++(0.04,0.04) (0.685,8.43)rectangle++(0.04,0.04)
  (0.917,8.734)rectangle++(0.04,0.04) (1.127,8.984)rectangle++(0.04,0.04) (1.315,9.187)rectangle++(0.04,0.04) (2.358,2.977)rectangle++(0.04,0.04) (2.121,4.032)rectangle++(0.04,0.04)
  (1.962,4.939)rectangle++(0.04,0.04) (1.863,5.716)rectangle++(0.04,0.04) (1.81,6.38)rectangle++(0.04,0.04) (1.791,6.947)rectangle++(0.04,0.04) (1.798,7.429)rectangle++(0.04,0.04)
  (1.824,7.839)rectangle++(0.04,0.04) (1.904,2.786)rectangle++(0.04,0.04) (1.728,3.888)rectangle++(0.04,0.04) (1.622,4.832)rectangle++(0.04,0.04) (1.57,5.638)rectangle++(0.04,0.04)
  (1.557,6.326)rectangle++(0.04,0.04) (1.575,6.911)rectangle++(0.04,0.04) (1.613,7.408)rectangle++(0.04,0.04) (1.665,7.828)rectangle++(0.04,0.04) (1.726,8.183)rectangle++(0.04,0.04)
  (1.792,8.483)rectangle++(0.04,0.04) (1.861,8.735)rectangle++(0.04,0.04) (-0.02,1.381)rectangle++(0.04,0.04) (0.036,2.771)rectangle++(0.04,0.04) (0.139,3.95)rectangle++(0.04,0.04)
  (0.274,4.948)rectangle++(0.04,0.04) (0.429,5.791)rectangle++(0.04,0.04) (0.594,6.501)rectangle++(0.04,0.04) (0.763,7.099)rectangle++(0.04,0.04) (0.93,7.6)rectangle++(0.04,0.04)
  (1.092,8.019)rectangle++(0.04,0.04) (1.247,8.368)rectangle++(0.04,0.04) (1.392,8.659)rectangle++(0.04,0.04) (1.528,8.901)rectangle++(0.04,0.04) (1.652,9.101)rectangle++(0.04,0.04)
  (1.766,9.266)rectangle++(0.04,0.04) (1.87,9.402)rectangle++(0.04,0.04) (1.963,9.513)rectangle++(0.04,0.04) (-0.02,1.505)rectangle++(0.04,0.04) (0.041,2.876)rectangle++(0.04,0.04)
  (0.148,4.04)rectangle++(0.04,0.04) (0.285,5.024)rectangle++(0.04,0.04) (1.347,1.71)rectangle++(0.04,0.04) (1.211,2.996)rectangle++(0.04,0.04) (1.147,4.094)rectangle++(0.04,0.04)
  (1.136,5.03)rectangle++(0.04,0.04) (1.165,5.827)rectangle++(0.04,0.04) (1.221,6.502)rectangle++(0.04,0.04) (1.296,7.074)rectangle++(0.04,0.04) (1.382,7.557)rectangle++(0.04,0.04)
  (1.475,7.965)rectangle++(0.04,0.04) (1.57,8.307)rectangle++(0.04,0.04) (1.665,8.595)rectangle++(0.04,0.04) (1.757,8.835)rectangle++(0.04,0.04) (1.844,9.036)rectangle++(0.04,0.04)
  (-0.02,1.429)rectangle++(0.04,0.04) (0.038,2.812)rectangle++(0.04,0.04) (0.143,3.985)rectangle++(0.04,0.04) (1.077,1.423)rectangle++(0.04,0.04) (0.22,1.052)rectangle++(0.04,0.04)
  (0.227,2.481)rectangle++(0.04,0.04) (-0.621,2.187)rectangle++(0.04,0.04) (-0.443,3.48)rectangle++(0.04,0.04) (-0.239,4.572)rectangle++(0.04,0.04) (-0.023,5.492)rectangle++(0.04,0.04)
  (0.198,6.265)rectangle++(0.04,0.04) (-1.611,3.013)rectangle++(0.04,0.04) (-1.251,4.222)rectangle++(0.04,0.04) (-0.896,5.235)rectangle++(0.04,0.04) (-0.555,6.081)rectangle++(0.04,0.04)
  (-1.713,2.799)rectangle++(0.04,0.04) (-1.347,4.044)rectangle++(0.04,0.04) (-0.985,5.088)rectangle++(0.04,0.04) (-0.636,5.96)rectangle++(0.04,0.04) (-0.304,6.688)rectangle++(0.04,0.04)
  (0.007,7.293)rectangle++(0.04,0.04) (0.295,7.795)rectangle++(0.04,0.04) (0.56,8.21)rectangle++(0.04,0.04) (0.803,8.552)rectangle++(0.04,0.04) (1.022,8.834)rectangle++(0.04,0.04)
  (1.22,9.064)rectangle++(0.04,0.04) (2.337,2.923)rectangle++(0.04,0.04) (2.102,3.987)rectangle++(0.04,0.04) (1.944,4.901)rectangle++(0.04,0.04) (1.846,5.684)rectangle++(0.04,0.04)
  (1.794,6.354)rectangle++(0.04,0.04) (1.777,6.925)rectangle++(0.04,0.04) (1.785,7.412)rectangle++(0.04,0.04) (1.79,2.673)rectangle++(0.04,0.04) (1.626,3.797)rectangle++(0.04,0.04)
  (1.532,4.758)rectangle++(0.04,0.04) (-0.952,2.988)rectangle++(0.04,0.04) (-0.692,4.174)rectangle++(0.04,0.04) (-0.423,5.172)rectangle++(0.04,0.04) (-0.155,6.009)rectangle++(0.04,0.04)
  (0.106,6.71)rectangle++(0.04,0.04) (0.357,7.296)rectangle++(0.04,0.04) (0.593,7.783)rectangle++(0.04,0.04) (0.813,8.188)rectangle++(0.04,0.04) (2.153,2.607)rectangle++(0.04,0.04)
  (0.389,1.615)rectangle++(0.04,0.04) (0.393,2.954)rectangle++(0.04,0.04) (0.45,4.091)rectangle++(0.04,0.04) (1.061,1.529)rectangle++(0.04,0.04) (0.96,2.853)rectangle++(0.04,0.04)
  (0.928,3.983)rectangle++(0.04,0.04) (0.946,4.945)rectangle++(0.04,0.04) (1,5.761)rectangle++(0.04,0.04) (1.078,6.453)rectangle++(0.04,0.04) (1.172,7.038)rectangle++(0.04,0.04)
  (1.276,7.532)rectangle++(0.04,0.04) (1.384,7.947)rectangle++(0.04,0.04) (-0.02,1.255)rectangle++(0.04,0.04) (0.031,2.664)rectangle++(0.04,0.04) (0.131,3.859)rectangle++(0.04,0.04)
  (-0.998,2.468)rectangle++(0.04,0.04) (-0.752,3.734)rectangle++(0.04,0.04) (-0.492,4.8)rectangle++(0.04,0.04) (-0.229,5.696)rectangle++(0.04,0.04) (0.031,6.447)rectangle++(0.04,0.04)
  (0.282,7.075)rectangle++(0.04,0.04) (0.521,7.599)rectangle++(0.04,0.04) (0.744,8.034)rectangle++(0.04,0.04) (0.952,8.395)rectangle++(0.04,0.04) (1.143,8.694)rectangle++(0.04,0.04)
  (1.317,8.941)rectangle++(0.04,0.04) (1.475,9.143)rectangle++(0.04,0.04) (1.617,9.309)rectangle++(0.04,0.04) (1.745,9.444)rectangle++(0.04,0.04) (1.859,9.554)rectangle++(0.04,0.04)
  (1.96,9.643)rectangle++(0.04,0.04) (2.049,9.714)rectangle++(0.04,0.04) (2.128,9.771)rectangle++(0.04,0.04) (2.198,9.817)rectangle++(0.04,0.04) (2.258,9.852)rectangle++(0.04,0.04)
  (2.312,9.88)rectangle++(0.04,0.04) (2.358,9.902)rectangle++(0.04,0.04) (2.398,9.919)rectangle++(0.04,0.04) (2.433,9.931)rectangle++(0.04,0.04) (2.463,9.94)rectangle++(0.04,0.04)
  (2.489,9.947)rectangle++(0.04,0.04) (2.511,9.952)rectangle++(0.04,0.04) (2.53,9.955)rectangle++(0.04,0.04) (2.547,9.956)rectangle++(0.04,0.04) (2.561,9.957)rectangle++(0.04,0.04)
  (2.386,3.486)rectangle++(0.04,0.04) (2.166,4.463)rectangle++(0.04,0.04) (2.017,5.304)rectangle++(0.04,0.04) (1.165,2.227)rectangle++(0.04,0.04) (1.077,3.443)rectangle++(0.04,0.04)
  (1.051,4.48)rectangle++(0.04,0.04) (1.07,5.362)rectangle++(0.04,0.04) (1.122,6.111)rectangle++(0.04,0.04) (1.196,6.746)rectangle++(0.04,0.04) (1.284,7.282)rectangle++(0.04,0.04)
  (1.381,7.735)rectangle++(0.04,0.04) (1.481,8.115)rectangle++(0.04,0.04) (2.033,2.763)rectangle++(0.04,0.04) (1.836,3.863)rectangle++(0.04,0.04) (1.713,4.806)rectangle++(0.04,0.04)
  (1.646,5.613)rectangle++(0.04,0.04) (1.622,6.302)rectangle++(0.04,0.04) (1.628,6.888)rectangle++(0.04,0.04) (1.657,7.386)rectangle++(0.04,0.04) (1.702,7.808)rectangle++(0.04,0.04)
  (1.757,8.165)rectangle++(0.04,0.04) (1.818,8.466)rectangle++(0.04,0.04) (1.881,8.719)rectangle++(0.04,0.04) (1.946,8.932)rectangle++(0.04,0.04) (2.192,3.08)rectangle++(0.04,0.04)
  (1.984,4.126)rectangle++(0.04,0.04) (-0.697,2.953)rectangle++(0.04,0.04) (-0.477,4.134)rectangle++(0.04,0.04) (-0.242,5.129)rectangle++(0.04,0.04) (-0.003,5.966)rectangle++(0.04,0.04)
  (0.234,6.667)rectangle++(0.04,0.04) (0.463,7.254)rectangle++(0.04,0.04) (0.682,7.744)rectangle++(0.04,0.04) (0.887,8.151)rectangle++(0.04,0.04) (1.078,8.489)rectangle++(0.04,0.04)
  (2.198,2.748)rectangle++(0.04,0.04) (1.976,3.844)rectangle++(0.04,0.04) (1.831,4.784)rectangle++(0.04,0.04) (1.746,5.59)rectangle++(0.04,0.04) (1.705,6.278)rectangle++(0.04,0.04)
  (1.698,6.864)rectangle++(0.04,0.04) (1.716,7.363)rectangle++(0.04,0.04) (1.751,7.786)rectangle++(0.04,0.04) (1.797,8.144)rectangle++(0.04,0.04) (1.851,8.447)rectangle++(0.04,0.04)
  (1.909,8.702)rectangle++(0.04,0.04) (-1.902,3.943)rectangle++(0.04,0.04) (-1.427,2.019)rectangle++(0.04,0.04) (-1.134,3.369)rectangle++(0.04,0.04) (-0.831,4.506)rectangle++(0.04,0.04)
  (-0.529,5.459)rectangle++(0.04,0.04) (-0.233,6.258)rectangle++(0.04,0.04) (-1.695,2.912)rectangle++(0.04,0.04) (-1.326,4.139)rectangle++(0.04,0.04) (-0.964,5.168)rectangle++(0.04,0.04)
  (-1.558,2.504)rectangle++(0.04,0.04) (-1.226,3.787)rectangle++(0.04,0.04) (-0.893,4.864)rectangle++(0.04,0.04) (-0.567,5.766)rectangle++(0.04,0.04) (-0.253,6.52)rectangle++(0.04,0.04)
  (1.846,1.929)rectangle++(0.04,0.04) (1.644,3.162)rectangle++(0.04,0.04) (1.522,4.218)rectangle++(0.04,0.04) (-0.814,2.867)rectangle++(0.04,0.04) (-0.579,4.066)rectangle++(0.04,0.04)
  (-0.332,5.075)rectangle++(0.04,0.04) (1.453,1.562)rectangle++(0.04,0.04) (1.296,2.866)rectangle++(0.04,0.04) (0.591,1.455)rectangle++(0.04,0.04) (0.558,2.809)rectangle++(0.04,0.04)
  (0.584,3.962)rectangle++(0.04,0.04) (0.653,4.94)rectangle++(0.04,0.04) (0.751,5.769)rectangle++(0.04,0.04) (0.866,6.47)rectangle++(0.04,0.04) (0.993,7.061)rectangle++(0.04,0.04)
  (1.124,7.558)rectangle++(0.04,0.04) (1.256,7.976)rectangle++(0.04,0.04) (1.384,8.325)rectangle++(0.04,0.04) (1.507,8.617)rectangle++(0.04,0.04) (1.624,8.861)rectangle++(0.04,0.04)
  (1.733,9.063)rectangle++(0.04,0.04) (1.833,9.23)rectangle++(0.04,0.04) (2.292,3.122)rectangle++(0.04,0.04) (-0.374,2.803)rectangle++(0.04,0.04) (-0.208,3.994)rectangle++(0.04,0.04)
  (-0.02,0.622)rectangle++(0.04,0.04) (0.006,2.126)rectangle++(0.04,0.04) (0.088,3.403)rectangle++(0.04,0.04) (0.208,4.485)rectangle++(0.04,0.04) (0.354,5.4)rectangle++(0.04,0.04)
  (0.515,6.172)rectangle++(0.04,0.04) (0.682,6.822)rectangle++(0.04,0.04) (0.851,7.368)rectangle++(0.04,0.04) (-1.767,3.406)rectangle++(0.04,0.04) (-1.368,4.562)rectangle++(0.04,0.04)
  (-0.982,5.528)rectangle++(0.04,0.04) (-0.616,6.335)rectangle++(0.04,0.04) (-0.272,7.005)rectangle++(0.04,0.04) (0.046,7.561)rectangle++(0.04,0.04) (0.34,8.022)rectangle++(0.04,0.04)
  (0.607,8.401)rectangle++(0.04,0.04) (0.85,8.713)rectangle++(0.04,0.04) (1.069,8.968)rectangle++(0.04,0.04) (1.265,9.176)rectangle++(0.04,0.04) (1.44,9.345)rectangle++(0.04,0.04)
  (1.596,9.482)rectangle++(0.04,0.04) (1.734,9.592)rectangle++(0.04,0.04) (1.855,9.68)rectangle++(0.04,0.04) (-2.167,4.145)rectangle++(0.04,0.04) (-1.678,5.206)rectangle++(0.04,0.04)
  (-1.711,2.348)rectangle++(0.04,0.04) (-1.362,3.661)rectangle++(0.04,0.04) (-1.245,2.081)rectangle++(0.04,0.04) (-0.978,3.415)rectangle++(0.04,0.04) (-0.697,4.538)rectangle++(0.04,0.04)
  (-1.34,2.427)rectangle++(0.04,0.04) (0.863,0.664)rectangle++(0.04,0.04) (0.758,2.126)rectangle++(0.04,0.04) (0.727,3.373)rectangle++(0.04,0.04) (0.751,4.434)rectangle++(0.04,0.04)
  (0.813,5.335)rectangle++(0.04,0.04) (0.903,6.099)rectangle++(0.04,0.04) (1.009,6.744)rectangle++(0.04,0.04) (1.125,7.288)rectangle++(0.04,0.04) (1.246,7.746)rectangle++(0.04,0.04)
  (-1.786,3.58)rectangle++(0.04,0.04) (-1.377,4.71)rectangle++(0.04,0.04) (-0.984,5.655)rectangle++(0.04,0.04) (-0.613,6.442)rectangle++(0.04,0.04) (-0.265,7.097)rectangle++(0.04,0.04)
  (0.056,7.639)rectangle++(0.04,0.04) (0.351,8.087)rectangle++(0.04,0.04) (0.62,8.456)rectangle++(0.04,0.04) (-2.096,3.592)rectangle++(0.04,0.04) (-1.64,4.733)rectangle++(0.04,0.04)
  (-1.207,5.685)rectangle++(0.04,0.04) (-0.801,6.477)rectangle++(0.04,0.04) (-0.424,7.133)rectangle++(0.04,0.04) (-0.077,7.677)rectangle++(0.04,0.04) (0.239,8.124)rectangle++(0.04,0.04)
  (0.526,8.492)rectangle++(0.04,0.04) (0.785,8.794)rectangle++(0.04,0.04) (1.017,9.039)rectangle++(0.04,0.04) (1.224,9.239)rectangle++(0.04,0.04) (1.407,9.4)rectangle++(0.04,0.04);
\fill[nuage]
  (1.57,9.53)rectangle++(0.04,0.04) (1.714,9.634)rectangle++(0.04,0.04) (-0.02,1.525)rectangle++(0.04,0.04) (0.042,2.893)rectangle++(0.04,0.04) (0.149,4.054)rectangle++(0.04,0.04)
  (0.287,5.036)rectangle++(0.04,0.04) (0.443,5.865)rectangle++(0.04,0.04) (0.609,6.564)rectangle++(0.04,0.04) (0.778,7.151)rectangle++(0.04,0.04) (0.945,7.644)rectangle++(0.04,0.04)
  (-1.82,3.488)rectangle++(0.04,0.04) (-1.292,1.938)rectangle++(0.04,0.04) (-1.023,3.295)rectangle++(0.04,0.04) (-0.74,4.438)rectangle++(0.04,0.04) (-0.454,5.398)rectangle++(0.04,0.04)
  (-0.172,6.203)rectangle++(0.04,0.04) (0.1,6.875)rectangle++(0.04,0.04) (1.893,2.106)rectangle++(0.04,0.04) (0.288,1.428)rectangle++(0.04,0.04) (0.3,2.798)rectangle++(0.04,0.04)
  (-0.689,2.274)rectangle++(0.04,0.04) (-0.497,3.556)rectangle++(0.04,0.04) (-0.282,4.639)rectangle++(0.04,0.04) (-0.056,5.551)rectangle++(0.04,0.04) (0.172,6.316)rectangle++(0.04,0.04)
  (0.397,6.958)rectangle++(0.04,0.04) (0.613,7.495)rectangle++(0.04,0.04) (0.819,7.942)rectangle++(0.04,0.04) (1.011,8.314)rectangle++(0.04,0.04) (1.19,8.623)rectangle++(0.04,0.04)
  (1.354,8.878)rectangle++(0.04,0.04) (1.504,9.088)rectangle++(0.04,0.04) (1.64,9.261)rectangle++(0.04,0.04) (1.762,9.403)rectangle++(0.04,0.04) (1.872,9.518)rectangle++(0.04,0.04)
  (1.969,9.612)rectangle++(0.04,0.04) (2.056,9.687)rectangle++(0.04,0.04) (2.133,9.748)rectangle++(0.04,0.04) (2.392,3.324)rectangle++(0.04,0.04) (2.164,4.326)rectangle++(0.04,0.04)
  (2.01,5.187)rectangle++(0.04,0.04) (1.914,5.925)rectangle++(0.04,0.04) (1.862,6.556)rectangle++(0.04,0.04) (1.842,7.094)rectangle++(0.04,0.04) (1.848,7.552)rectangle++(0.04,0.04)
  (1.87,7.942)rectangle++(0.04,0.04) (1.905,8.272)rectangle++(0.04,0.04) (1.948,8.551)rectangle++(0.04,0.04) (1.996,8.787)rectangle++(0.04,0.04) (2.046,8.985)rectangle++(0.04,0.04)
  (2.096,9.152)rectangle++(0.04,0.04) (2.146,9.291)rectangle++(0.04,0.04) (2.193,9.408)rectangle++(0.04,0.04) (2.238,9.505)rectangle++(0.04,0.04) (-0.02,1.504)rectangle++(0.04,0.04)
  (0.041,2.875)rectangle++(0.04,0.04) (0.148,4.039)rectangle++(0.04,0.04) (0.285,5.023)rectangle++(0.04,0.04) (0.441,5.855)rectangle++(0.04,0.04) (0.607,6.555)rectangle++(0.04,0.04)
  (0.776,7.144)rectangle++(0.04,0.04) (0.943,7.637)rectangle++(0.04,0.04) (1.105,8.05)rectangle++(0.04,0.04) (1.259,8.395)rectangle++(0.04,0.04) (1.404,8.681)rectangle++(0.04,0.04)
  (1.538,8.919)rectangle++(0.04,0.04) (1.662,9.116)rectangle++(0.04,0.04) (1.775,9.278)rectangle++(0.04,0.04) (1.878,9.412)rectangle++(0.04,0.04) (1.97,9.521)rectangle++(0.04,0.04)
  (2.053,9.61)rectangle++(0.04,0.04) (2.128,9.683)rectangle++(0.04,0.04) (2.194,9.741)rectangle++(0.04,0.04) (2.252,9.789)rectangle++(0.04,0.04) (2.304,9.827)rectangle++(0.04,0.04)
  (2.349,9.857)rectangle++(0.04,0.04) (2.389,9.88)rectangle++(0.04,0.04) (2.391,3.422)rectangle++(0.04,0.04) (2.167,4.41)rectangle++(0.04,0.04) (0.892,2.052)rectangle++(0.04,0.04)
  (0.838,3.304)rectangle++(0.04,0.04) (-0.713,2.509)rectangle++(0.04,0.04) (-0.508,3.757)rectangle++(0.04,0.04) (-0.283,4.81)rectangle++(0.04,0.04) (-0.051,5.696)rectangle++(0.04,0.04)
  (1.585,1.784)rectangle++(0.04,0.04) (1.417,3.049)rectangle++(0.04,0.04) (1.324,4.131)rectangle++(0.04,0.04) (1.288,5.055)rectangle++(0.04,0.04) (1.295,5.841)rectangle++(0.04,0.04)
  (1.424,2.169)rectangle++(0.04,0.04) (1.295,3.383)rectangle++(0.04,0.04) (1.234,4.42)rectangle++(0.04,0.04) (1.223,5.303)rectangle++(0.04,0.04) (1.25,6.055)rectangle++(0.04,0.04)
  (1.302,6.693)rectangle++(0.04,0.04) (1.372,7.233)rectangle++(0.04,0.04) (1.802,2.523)rectangle++(0.04,0.04) (1.63,3.669)rectangle++(0.04,0.04) (1.53,4.649)rectangle++(0.04,0.04)
  (1.485,5.487)rectangle++(0.04,0.04) (1.479,6.201)rectangle++(0.04,0.04) (1.503,6.808)rectangle++(0.04,0.04) (1.663,2.455)rectangle++(0.04,0.04) (1.51,3.616)rectangle++(0.04,0.04)
  (1.426,4.61)rectangle++(0.04,0.04) (1.394,5.457)rectangle++(0.04,0.04) (1.401,6.179)rectangle++(0.04,0.04) (1.436,6.792)rectangle++(0.04,0.04) (1.669,2.433)rectangle++(0.04,0.04)
  (1.514,3.598)rectangle++(0.04,0.04) (0.763,1.687)rectangle++(0.04,0.04) (0.714,3)rectangle++(0.04,0.04) (0.715,1.336)rectangle++(0.04,0.04) (0.249,0.937)rectangle++(0.04,0.04)
  (0.247,2.382)rectangle++(0.04,0.04) (0.303,3.611)rectangle++(0.04,0.04) (0.4,4.654)rectangle++(0.04,0.04) (0.524,5.536)rectangle++(0.04,0.04) (0.665,6.281)rectangle++(0.04,0.04)
  (0.814,6.908)rectangle++(0.04,0.04) (0.966,7.436)rectangle++(0.04,0.04) (1.116,7.878)rectangle++(0.04,0.04) (2.021,2.611)rectangle++(0.04,0.04) (1.82,3.735)rectangle++(0.04,0.04)
  (1.694,4.698)rectangle++(0.04,0.04) (1.626,5.522)rectangle++(0.04,0.04) (-1.132,3.178)rectangle++(0.04,0.04) (-1.074,2.028)rectangle++(0.04,0.04) (-0.834,3.363)rectangle++(0.04,0.04)
  (-0.576,4.488)rectangle++(0.04,0.04) (-1.303,2.444)rectangle++(0.04,0.04) (0.862,0.678)rectangle++(0.04,0.04) (0.758,2.138)rectangle++(0.04,0.04) (0.728,3.383)rectangle++(0.04,0.04)
  (0.752,4.443)rectangle++(0.04,0.04) (0.814,5.342)rectangle++(0.04,0.04) (0.904,6.105)rectangle++(0.04,0.04) (1.01,6.749)rectangle++(0.04,0.04) (1.126,7.292)rectangle++(0.04,0.04)
  (1.247,7.75)rectangle++(0.04,0.04) (1.368,8.134)rectangle++(0.04,0.04) (1.486,8.455)rectangle++(0.04,0.04) (1.599,8.724)rectangle++(0.04,0.04) (1.706,8.947)rectangle++(0.04,0.04)
  (-2.006,3.95)rectangle++(0.04,0.04) (1.389,0.856)rectangle++(0.04,0.04) (1.213,2.268)rectangle++(0.04,0.04) (1.12,3.476)rectangle++(0.04,0.04) (1.089,4.506)rectangle++(0.04,0.04)
  (1.103,5.383)rectangle++(0.04,0.04) (1.151,6.127)rectangle++(0.04,0.04) (1.221,6.758)rectangle++(0.04,0.04) (1.306,7.292)rectangle++(0.04,0.04) (-1.656,3.494)rectangle++(0.04,0.04)
  (-1.27,4.632)rectangle++(0.04,0.04) (-0.896,5.584)rectangle++(0.04,0.04) (-0.541,6.379)rectangle++(0.04,0.04) (-0.207,7.04)rectangle++(0.04,0.04) (0.104,7.588)rectangle++(0.04,0.04)
  (0.389,8.042)rectangle++(0.04,0.04) (0.651,8.416)rectangle++(0.04,0.04) (0.887,8.724)rectangle++(0.04,0.04) (1.101,8.976)rectangle++(0.04,0.04) (1.293,9.182)rectangle++(0.04,0.04)
  (1.464,9.349)rectangle++(0.04,0.04) (1.616,9.484)rectangle++(0.04,0.04) (1.751,9.593)rectangle++(0.04,0.04) (1.87,9.68)rectangle++(0.04,0.04) (1.974,9.75)rectangle++(0.04,0.04)
  (2.066,9.805)rectangle++(0.04,0.04) (2.146,9.847)rectangle++(0.04,0.04) (2.216,9.881)rectangle++(0.04,0.04) (2.276,9.906)rectangle++(0.04,0.04) (2.329,9.925)rectangle++(0.04,0.04)
  (2.375,9.94)rectangle++(0.04,0.04) (2.414,9.95)rectangle++(0.04,0.04) (2.407,3.446)rectangle++(0.04,0.04) (2.181,4.429)rectangle++(0.04,0.04) (2.029,5.273)rectangle++(0.04,0.04)
  (1.933,5.997)rectangle++(0.04,0.04) (1.881,6.617)rectangle++(0.04,0.04) (1.861,7.145)rectangle++(0.04,0.04) (1.704,2.629)rectangle++(0.04,0.04) (1.551,3.762)rectangle++(0.04,0.04)
  (1.467,4.732)rectangle++(0.04,0.04) (1.434,5.56)rectangle++(0.04,0.04) (1.439,6.265)rectangle++(0.04,0.04) (1.472,6.864)rectangle++(0.04,0.04) (1.523,7.371)rectangle++(0.04,0.04)
  (1.587,7.801)rectangle++(0.04,0.04) (1.659,8.164)rectangle++(0.04,0.04) (1.735,8.469)rectangle++(0.04,0.04) (1.811,8.725)rectangle++(0.04,0.04) (1.886,8.94)rectangle++(0.04,0.04)
  (1.959,9.12)rectangle++(0.04,0.04) (2.027,9.27)rectangle++(0.04,0.04) (2.092,9.394)rectangle++(0.04,0.04) (2.152,9.498)rectangle++(0.04,0.04) (2.207,9.583)rectangle++(0.04,0.04)
  (2.257,9.654)rectangle++(0.04,0.04) (2.302,9.712)rectangle++(0.04,0.04) (2.343,9.759)rectangle++(0.04,0.04) (2.38,9.798)rectangle++(0.04,0.04) (2.369,3.4)rectangle++(0.04,0.04)
  (2.147,4.392)rectangle++(0.04,0.04) (1.999,5.243)rectangle++(0.04,0.04) (1.906,5.973)rectangle++(0.04,0.04) (1.857,6.597)rectangle++(0.04,0.04) (1.84,7.129)rectangle++(0.04,0.04)
  (1.847,7.582)rectangle++(0.04,0.04) (1.871,7.967)rectangle++(0.04,0.04) (1.907,8.294)rectangle++(0.04,0.04) (1.951,8.57)rectangle++(0.04,0.04) (1.999,8.802)rectangle++(0.04,0.04)
  (2.049,8.998)rectangle++(0.04,0.04) (-1.951,4.04)rectangle++(0.04,0.04) (-1.499,5.108)rectangle++(0.04,0.04) (-1.072,5.998)rectangle++(0.04,0.04) (1.823,1.591)rectangle++(0.04,0.04)
  (1.611,2.875)rectangle++(0.04,0.04) (1.482,3.976)rectangle++(0.04,0.04) (1.417,4.916)rectangle++(0.04,0.04) (1.399,5.719)rectangle++(0.04,0.04) (1.415,6.401)rectangle++(0.04,0.04)
  (1.457,6.98)rectangle++(0.04,0.04) (1.515,7.471)rectangle++(0.04,0.04) (1.585,7.886)rectangle++(0.04,0.04) (-1.755,3.688)rectangle++(0.04,0.04) (-1.346,4.802)rectangle++(0.04,0.04)
  (-0.954,5.731)rectangle++(0.04,0.04) (-0.584,6.506)rectangle++(0.04,0.04) (-0.238,7.15)rectangle++(0.04,0.04) (0.081,7.683)rectangle++(0.04,0.04) (0.374,8.123)rectangle++(0.04,0.04)
  (0.641,8.486)rectangle++(0.04,0.04) (0.882,8.784)rectangle++(0.04,0.04) (1.099,9.027)rectangle++(0.04,0.04) (1.293,9.225)rectangle++(0.04,0.04) (2.372,2.98)rectangle++(0.04,0.04)
  (0.461,1.762)rectangle++(0.04,0.04) (0.46,3.075)rectangle++(0.04,0.04) (0.512,4.192)rectangle++(0.04,0.04) (0.601,5.139)rectangle++(0.04,0.04) (0.714,5.94)rectangle++(0.04,0.04)
  (0.842,6.617)rectangle++(0.04,0.04) (-1.573,3.238)rectangle++(0.04,0.04) (-1.21,4.412)rectangle++(0.04,0.04) (-0.854,5.395)rectangle++(0.04,0.04) (-0.512,6.216)rectangle++(0.04,0.04)
  (-0.189,6.9)rectangle++(0.04,0.04) (0.113,7.469)rectangle++(0.04,0.04) (0.393,7.94)rectangle++(0.04,0.04) (-2.007,3.426)rectangle++(0.04,0.04) (-1.571,4.589)rectangle++(0.04,0.04)
  (-1.154,5.559)rectangle++(0.04,0.04) (-0.761,6.368)rectangle++(0.04,0.04) (-0.394,7.039)rectangle++(0.04,0.04) (-0.056,7.595)rectangle++(0.04,0.04) (0.254,8.055)rectangle++(0.04,0.04)
  (0.536,8.432)rectangle++(0.04,0.04) (0.791,8.742)rectangle++(0.04,0.04) (1.02,8.996)rectangle++(0.04,0.04) (1.224,9.202)rectangle++(0.04,0.04) (1.407,9.369)rectangle++(0.04,0.04)
  (-2.176,3.974)rectangle++(0.04,0.04) (-1.693,5.061)rectangle++(0.04,0.04) (-1.238,5.966)rectangle++(0.04,0.04) (1.839,1.54)rectangle++(0.04,0.04) (1.622,2.831)rectangle++(0.04,0.04)
  (1.49,3.938)rectangle++(0.04,0.04) (1.422,4.884)rectangle++(0.04,0.04) (1.402,5.691)rectangle++(0.04,0.04) (1.417,6.377)rectangle++(0.04,0.04) (1.457,6.96)rectangle++(0.04,0.04)
  (1.515,7.454)rectangle++(0.04,0.04) (1.584,7.872)rectangle++(0.04,0.04) (1.659,8.224)rectangle++(0.04,0.04) (1.737,8.52)rectangle++(0.04,0.04) (1.815,8.769)rectangle++(0.04,0.04)
  (1.891,8.977)rectangle++(0.04,0.04) (1.964,9.151)rectangle++(0.04,0.04) (2.033,9.296)rectangle++(0.04,0.04) (-2.031,4.102)rectangle++(0.04,0.04) (-1.565,5.164)rectangle++(0.04,0.04)
  (-1.126,6.048)rectangle++(0.04,0.04) (-0.717,6.782)rectangle++(0.04,0.04) (-0.341,7.39)rectangle++(0.04,0.04) (0.004,7.891)rectangle++(0.04,0.04) (0.317,8.303)rectangle++(0.04,0.04)
  (0.599,8.641)rectangle++(0.04,0.04) (0.853,8.918)rectangle++(0.04,0.04) (-0.02,1.41)rectangle++(0.04,0.04) (0.037,2.796)rectangle++(0.04,0.04) (0.76,1.111)rectangle++(0.04,0.04)
  (0.688,2.51)rectangle++(0.04,0.04) (0.582,1.211)rectangle++(0.04,0.04) (-0.02,0.177)rectangle++(0.04,0.04) (-0.012,1.747)rectangle++(0.04,0.04) (0.057,3.082)rectangle++(0.04,0.04)
  (0.17,4.214)rectangle++(0.04,0.04) (0.311,5.171)rectangle++(0.04,0.04) (0.469,5.979)rectangle++(0.04,0.04) (0.635,6.66)rectangle++(0.04,0.04) (0.804,7.232)rectangle++(0.04,0.04)
  (1.887,2.375)rectangle++(0.04,0.04) (-0.261,2.545)rectangle++(0.04,0.04) (-0.122,3.77)rectangle++(0.04,0.04) (0.045,4.806)rectangle++(0.04,0.04) (-1.262,2.657)rectangle++(0.04,0.04)
  (0.916,0.739)rectangle++(0.04,0.04) (0.806,2.188)rectangle++(0.04,0.04) (0.77,3.424)rectangle++(0.04,0.04) (0.789,4.476)rectangle++(0.04,0.04) (0.848,5.369)rectangle++(0.04,0.04)
  (0.933,6.126)rectangle++(0.04,0.04) (1.036,6.766)rectangle++(0.04,0.04) (1.149,7.306)rectangle++(0.04,0.04) (1.267,7.76)rectangle++(0.04,0.04) (1.385,8.142)rectangle++(0.04,0.04)
  (-1.861,3.699)rectangle++(0.04,0.04) (-1.436,4.815)rectangle++(0.04,0.04) (1.546,1.212)rectangle++(0.04,0.04) (1.36,2.565)rectangle++(0.04,0.04) (1.257,3.722)rectangle++(0.04,0.04)
  (1.215,4.709)rectangle++(0.04,0.04) (1.219,5.551)rectangle++(0.04,0.04) (1.256,6.265)rectangle++(0.04,0.04) (1.549,2.26)rectangle++(0.04,0.04) (1.404,3.455)rectangle++(0.04,0.04)
  (1.33,4.477)rectangle++(0.04,0.04) (1.307,5.349)rectangle++(0.04,0.04) (1.323,6.09)rectangle++(0.04,0.04) (1.366,6.72)rectangle++(0.04,0.04) (1.428,7.254)rectangle++(0.04,0.04)
  (1.501,7.705)rectangle++(0.04,0.04) (1.582,8.085)rectangle++(0.04,0.04) (2.009,2.782)rectangle++(0.04,0.04) (0.46,1.62)rectangle++(0.04,0.04) (0.454,2.955)rectangle++(0.04,0.04)
  (0.502,4.09)rectangle++(0.04,0.04) (0.588,5.052)rectangle++(0.04,0.04) (0.7,5.867)rectangle++(0.04,0.04) (0.827,6.555)rectangle++(0.04,0.04) (0.963,7.135)rectangle++(0.04,0.04)
  (1.102,7.623)rectangle++(0.04,0.04) (1.239,8.031)rectangle++(0.04,0.04) (1.372,8.373)rectangle++(0.04,0.04) (1.499,8.659)rectangle++(0.04,0.04) (1.619,8.896)rectangle++(0.04,0.04)
  (1.729,9.093)rectangle++(0.04,0.04) (1.831,9.256)rectangle++(0.04,0.04) (1.925,9.391)rectangle++(0.04,0.04) (2.01,9.501)rectangle++(0.04,0.04) (2.086,9.592)rectangle++(0.04,0.04)
  (2.155,9.666)rectangle++(0.04,0.04) (2.216,9.726)rectangle++(0.04,0.04) (2.27,9.775)rectangle++(0.04,0.04) (2.318,9.814)rectangle++(0.04,0.04) (2.361,9.845)rectangle++(0.04,0.04)
  (2.399,9.87)rectangle++(0.04,0.04) (-2.108,4.312)rectangle++(0.04,0.04) (-1.564,2.053)rectangle++(0.04,0.04) (-1.249,3.404)rectangle++(0.04,0.04) (-0.928,4.539)rectangle++(0.04,0.04)
  (-0.609,5.492)rectangle++(0.04,0.04) (-0.301,6.289)rectangle++(0.04,0.04) (-0.006,6.954)rectangle++(0.04,0.04) (0.271,7.507)rectangle++(0.04,0.04) (0.528,7.966)rectangle++(0.04,0.04)
  (0.765,8.346)rectangle++(0.04,0.04) (0.982,8.66)rectangle++(0.04,0.04) (1.179,8.918)rectangle++(0.04,0.04) (1.357,9.129)rectangle++(0.04,0.04) (2.335,2.974)rectangle++(0.04,0.04)
  (2.102,4.031)rectangle++(0.04,0.04) (1.945,4.938)rectangle++(0.04,0.04) (1.849,5.716)rectangle++(0.04,0.04) (1.798,6.381)rectangle++(0.04,0.04) (1.781,6.948)rectangle++(0.04,0.04)
  (1.79,7.431)rectangle++(0.04,0.04) (1.816,7.841)rectangle++(0.04,0.04) (1.855,8.188)rectangle++(0.04,0.04) (1.902,8.482)rectangle++(0.04,0.04) (1.954,8.73)rectangle++(0.04,0.04)
  (2.008,8.938)rectangle++(0.04,0.04) (2.062,9.113)rectangle++(0.04,0.04) (-1.978,4.068)rectangle++(0.04,0.04) (-1.521,5.133)rectangle++(0.04,0.04) (-1.09,6.02)rectangle++(0.04,0.04)
  (-0.688,6.757)rectangle++(0.04,0.04) (-0.316,7.367)rectangle++(0.04,0.04) (0.024,7.871)rectangle++(0.04,0.04) (0.333,8.286)rectangle++(0.04,0.04) (0.612,8.626)rectangle++(0.04,0.04)
  (-2.141,3.627)rectangle++(0.04,0.04) (-1.677,4.765)rectangle++(0.04,0.04) (-1.237,5.714)rectangle++(0.04,0.04) (1.768,1.48)rectangle++(0.04,0.04) (1.56,2.783)rectangle++(0.04,0.04)
  (1.435,3.899)rectangle++(0.04,0.04) (1.373,4.853)rectangle++(0.04,0.04) (1.359,5.667)rectangle++(0.04,0.04) (-0.02,0.89)rectangle++(0.04,0.04) (0.016,2.353)rectangle++(0.04,0.04)
  (0.106,3.596)rectangle++(0.04,0.04) (0.232,4.648)rectangle++(0.04,0.04) (1.249,1.606)rectangle++(0.04,0.04) (1.124,2.911)rectangle++(0.04,0.04) (1.07,4.026)rectangle++(0.04,0.04)
  (0.949,1.674)rectangle++(0.04,0.04) (0.872,2.981)rectangle++(0.04,0.04) (0.858,4.096)rectangle++(0.04,0.04) (0.891,5.043)rectangle++(0.04,0.04) (0.957,5.847)rectangle++(0.04,0.04)
  (1.045,6.528)rectangle++(0.04,0.04) (1.147,7.103)rectangle++(0.04,0.04) (1.257,7.588)rectangle++(0.04,0.04) (1.37,7.996)rectangle++(0.04,0.04) (1.482,8.338)rectangle++(0.04,0.04)
  (2.095,2.816)rectangle++(0.04,0.04) (1.891,3.906)rectangle++(0.04,0.04) (1.762,4.841)rectangle++(0.04,0.04) (1.689,5.641)rectangle++(0.04,0.04) (1.659,6.323)rectangle++(0.04,0.04)
  (1.661,6.904)rectangle++(0.04,0.04) (1.667,2.519)rectangle++(0.04,0.04) (1.515,3.671)rectangle++(0.04,0.04) (1.433,4.656)rectangle++(0.04,0.04) (1.402,5.496)rectangle++(0.04,0.04)
  (1.409,6.212)rectangle++(0.04,0.04) (1.444,6.82)rectangle++(0.04,0.04) (1.498,7.335)rectangle++(0.04,0.04) (1.565,7.771)rectangle++(0.04,0.04) (1.639,8.139)rectangle++(0.04,0.04)
  (2.016,2.809)rectangle++(0.04,0.04) (1.824,3.904)rectangle++(0.04,0.04) (1.704,4.841)rectangle++(0.04,0.04) (1.64,5.643)rectangle++(0.04,0.04) (1.617,6.327)rectangle++(0.04,0.04)
  (1.626,6.91)rectangle++(0.04,0.04) (1.656,7.404)rectangle++(0.04,0.04) (1.702,7.824)rectangle++(0.04,0.04) (1.757,8.178)rectangle++(0.04,0.04) (1.818,8.477)rectangle++(0.04,0.04)
  (1.883,8.729)rectangle++(0.04,0.04) (1.947,8.941)rectangle++(0.04,0.04) (2.011,9.118)rectangle++(0.04,0.04) (2.071,9.266)rectangle++(0.04,0.04) (2.129,9.39)rectangle++(0.04,0.04)
  (2.183,9.492)rectangle++(0.04,0.04) (2.233,9.577)rectangle++(0.04,0.04) (-2.065,4.21)rectangle++(0.04,0.04) (-1.589,5.257)rectangle++(0.04,0.04) (-1.142,6.128)rectangle++(0.04,0.04)
  (-0.728,6.851)rectangle++(0.04,0.04) (-0.347,7.449)rectangle++(0.04,0.04) (0.001,7.941)rectangle++(0.04,0.04) (0.316,8.346)rectangle++(0.04,0.04) (0.6,8.678)rectangle++(0.04,0.04)
  (-2.157,3.636)rectangle++(0.04,0.04) (-1.691,4.773)rectangle++(0.04,0.04) (-1.248,5.721)rectangle++(0.04,0.04) (-0.834,6.509)rectangle++(0.04,0.04) (-1.88,2.829)rectangle++(0.04,0.04)
  (-1.487,4.076)rectangle++(0.04,0.04) (-1.103,5.12)rectangle++(0.04,0.04) (-0.735,5.993)rectangle++(0.04,0.04) (-0.388,6.719)rectangle++(0.04,0.04) (-0.063,7.323)rectangle++(0.04,0.04)
  (-1.938,3.186)rectangle++(0.04,0.04) (-1.522,4.382)rectangle++(0.04,0.04) (-1.121,5.381)rectangle++(0.04,0.04) (-0.739,6.215)rectangle++(0.04,0.04) (-0.382,6.909)rectangle++(0.04,0.04)
  (-1.894,3.021)rectangle++(0.04,0.04) (-1.491,4.24)rectangle++(0.04,0.04) (-1.1,5.26)rectangle++(0.04,0.04) (-0.727,6.111)rectangle++(0.04,0.04) (-0.376,6.82)rectangle++(0.04,0.04)
  (-0.049,7.408)rectangle++(0.04,0.04) (0.253,7.895)rectangle++(0.04,0.04) (0.528,8.297)rectangle++(0.04,0.04) (0.779,8.627)rectangle++(0.04,0.04) (1.005,8.898)rectangle++(0.04,0.04)
  (1.208,9.119)rectangle++(0.04,0.04) (1.389,9.299)rectangle++(0.04,0.04) (1.551,9.445)rectangle++(0.04,0.04) (2.395,3.1)rectangle++(0.04,0.04) (0.491,1.797)rectangle++(0.04,0.04)
  (0.487,3.104)rectangle++(0.04,0.04) (0.536,4.215)rectangle++(0.04,0.04) (-1.01,2.64)rectangle++(0.04,0.04) (0.873,0.801)rectangle++(0.04,0.04) (0.076,0.849)rectangle++(0.04,0.04)
  (0.096,2.315)rectangle++(0.04,0.04) (0.616,1.011)rectangle++(0.04,0.04) (0.562,2.431)rectangle++(0.04,0.04) (0.573,3.64)rectangle++(0.04,0.04) (0.916,1.452)rectangle++(0.04,0.04)
  (0.834,2.794)rectangle++(0.04,0.04) (0.819,3.938)rectangle++(0.04,0.04) (0.851,4.911)rectangle++(0.04,0.04) (-1.128,2.865)rectangle++(0.04,0.04) (-0.846,4.077)rectangle++(0.04,0.04)
  (-0.558,5.095)rectangle++(0.04,0.04) (-0.273,5.949)rectangle++(0.04,0.04) (1.689,1.787)rectangle++(0.04,0.04) (1.505,3.047)rectangle++(0.04,0.04) (0.61,1.553)rectangle++(0.04,0.04)
  (0.326,0.961)rectangle++(0.04,0.04) (0.313,2.4)rectangle++(0.04,0.04) (0.36,3.624)rectangle++(0.04,0.04) (0.449,4.662)rectangle++(0.04,0.04) (0.566,5.541)rectangle++(0.04,0.04)
  (0.7,6.283)rectangle++(0.04,0.04) (0.844,6.909)rectangle++(0.04,0.04) (0.992,7.435)rectangle++(0.04,0.04) (1.138,7.876)rectangle++(0.04,0.04) (1.28,8.246)rectangle++(0.04,0.04)
  (1.416,8.554)rectangle++(0.04,0.04) (-1.962,3.797)rectangle++(0.04,0.04) (-1.518,4.902)rectangle++(0.04,0.04) (-1.096,5.823)rectangle++(0.04,0.04) (-0.701,6.59)rectangle++(0.04,0.04)
  (-0.335,7.226)rectangle++(0.04,0.04) (0.002,7.751)rectangle++(0.04,0.04) (2.153,2.291)rectangle++(0.04,0.04) (-0.186,2.588)rectangle++(0.04,0.04) (-0.057,3.804)rectangle++(0.04,0.04)
  (0.101,4.832)rectangle++(0.04,0.04) (0.277,5.699)rectangle++(0.04,0.04) (1.537,1.87)rectangle++(0.04,0.04) (0.276,1.278)rectangle++(0.04,0.04) (0.283,2.672)rectangle++(0.04,0.04)
  (0.345,3.856)rectangle++(0.04,0.04) (-0.955,2.517)rectangle++(0.04,0.04) (-0.866,1.923)rectangle++(0.04,0.04) (-0.662,3.266)rectangle++(0.04,0.04) (-0.434,4.398)rectangle++(0.04,0.04)
  (-0.195,5.352)rectangle++(0.04,0.04) (0.046,6.153)rectangle++(0.04,0.04) (0.283,6.825)rectangle++(0.04,0.04) (0.511,7.386)rectangle++(0.04,0.04) (1.974,2.336)rectangle++(0.04,0.04)
  (1.769,3.502)rectangle++(0.04,0.04) (1.642,4.503)rectangle++(0.04,0.04) (1.573,5.358)rectangle++(0.04,0.04) (1.549,6.087)rectangle++(0.04,0.04) (1.558,6.708)rectangle++(0.04,0.04)
  (1.591,7.236)rectangle++(0.04,0.04) (-1.584,3.547)rectangle++(0.04,0.04) (-1.207,4.674)rectangle++(0.04,0.04) (-0.841,5.618)rectangle++(0.04,0.04) (-0.493,6.405)rectangle++(0.04,0.04)
  (-0.165,7.06)rectangle++(0.04,0.04) (0.14,7.604)rectangle++(0.04,0.04) (0.421,8.054)rectangle++(0.04,0.04) (0.678,8.425)rectangle++(0.04,0.04) (0.911,8.73)rectangle++(0.04,0.04)
  (1.121,8.981)rectangle++(0.04,0.04) (1.31,9.185)rectangle++(0.04,0.04) (1.479,9.351)rectangle++(0.04,0.04) (2.379,3.059)rectangle++(0.04,0.04) (2.142,4.101)rectangle++(0.04,0.04)
  (1.983,4.996)rectangle++(0.04,0.04) (1.883,5.764)rectangle++(0.04,0.04) (1.829,6.42)rectangle++(0.04,0.04) (1.809,6.98)rectangle++(0.04,0.04) (1.815,7.457)rectangle++(0.04,0.04)
  (1.839,7.862)rectangle++(0.04,0.04) (1.875,8.205)rectangle++(0.04,0.04) (1.92,8.496)rectangle++(0.04,0.04) (1.97,8.741)rectangle++(0.04,0.04) (2.022,8.947)rectangle++(0.04,0.04)
  (2.074,9.12)rectangle++(0.04,0.04) (2.126,9.266)rectangle++(0.04,0.04) (2.175,9.387)rectangle++(0.04,0.04) (2.222,9.488)rectangle++(0.04,0.04) (2.266,9.572)rectangle++(0.04,0.04)
  (2.307,9.642)rectangle++(0.04,0.04) (2.344,9.7)rectangle++(0.04,0.04) (2.378,9.747)rectangle++(0.04,0.04) (2.409,9.786)rectangle++(0.04,0.04) (2.437,9.818)rectangle++(0.04,0.04)
  (-2.086,4.31)rectangle++(0.04,0.04) (1.502,0.922)rectangle++(0.04,0.04) (1.312,2.32)rectangle++(0.04,0.04) (1.205,3.515)rectangle++(0.04,0.04) (1.163,4.536)rectangle++(0.04,0.04)
  (1.168,5.405)rectangle++(0.04,0.04) (1.207,6.144)rectangle++(0.04,0.04) (1.269,6.77)rectangle++(0.04,0.04) (1.347,7.3)rectangle++(0.04,0.04) (1.435,7.748)rectangle++(0.04,0.04)
  (-1.749,3.624)rectangle++(0.04,0.04) (-1.344,4.746)rectangle++(0.04,0.04) (-0.954,5.684)rectangle++(0.04,0.04) (-0.586,6.466)rectangle++(0.04,0.04) (-0.242,7.116)rectangle++(0.04,0.04)
  (2.011,2.075)rectangle++(0.04,0.04) (1.79,3.279)rectangle++(0.04,0.04) (1.651,4.312)rectangle++(0.04,0.04) (0.942,1.894)rectangle++(0.04,0.04) (0.875,3.168)rectangle++(0.04,0.04)
  (0.868,4.254)rectangle++(0.04,0.04) (0.906,5.178)rectangle++(0.04,0.04) (0.975,5.961)rectangle++(0.04,0.04) (1.065,6.624)rectangle++(0.04,0.04) (1.168,7.184)rectangle++(0.04,0.04)
  (-1.655,3.438)rectangle++(0.04,0.04) (-1.246,1.965)rectangle++(0.04,0.04) (-0.983,3.316)rectangle++(0.04,0.04) (-0.705,4.454)rectangle++(0.04,0.04) (-0.423,5.41)rectangle++(0.04,0.04)
  (-0.146,6.212)rectangle++(0.04,0.04) (-1.665,2.922)rectangle++(0.04,0.04) (-1.301,4.147)rectangle++(0.04,0.04) (-0.942,5.173)rectangle++(0.04,0.04) (-0.596,6.031)rectangle++(0.04,0.04)
  (-0.268,6.746)rectangle++(0.04,0.04) (0.04,7.341)rectangle++(0.04,0.04) (0.326,7.835)rectangle++(0.04,0.04) (0.588,8.243)rectangle++(0.04,0.04) (0.827,8.579)rectangle++(0.04,0.04)
  (1.044,8.855)rectangle++(0.04,0.04) (1.24,9.081)rectangle++(0.04,0.04) (1.415,9.266)rectangle++(0.04,0.04) (1.571,9.416)rectangle++(0.04,0.04) (1.71,9.537)rectangle++(0.04,0.04)
  (1.832,9.634)rectangle++(0.04,0.04) (1.941,9.712)rectangle++(0.04,0.04) (2.036,9.774)rectangle++(0.04,0.04) (2.119,9.822)rectangle++(0.04,0.04) (-2.151,4.237)rectangle++(0.04,0.04)
  (-1.661,5.284)rectangle++(0.04,0.04) (-1.203,6.154)rectangle++(0.04,0.04) (1.886,1.594)rectangle++(0.04,0.04) (-0.058,2.374)rectangle++(0.04,0.04) (0.656,0.984)rectangle++(0.04,0.04)
  (0.595,2.407)rectangle++(0.04,0.04) (-0.528,2.255)rectangle++(0.04,0.04) (-0.361,3.534)rectangle++(0.04,0.04) (-0.168,4.615)rectangle++(0.04,0.04) (1.3,1.494)rectangle++(0.04,0.04)
  (1.162,2.814)rectangle++(0.04,0.04) (1.098,3.942)rectangle++(0.04,0.04) (1.089,4.903)rectangle++(0.04,0.04) (1.12,5.72)rectangle++(0.04,0.04) (1.178,6.413)rectangle++(0.04,0.04)
  (1.256,7)rectangle++(0.04,0.04) (1.345,7.496)rectangle++(0.04,0.04) (-1.701,3.548)rectangle++(0.04,0.04) (1.231,0.839)rectangle++(0.04,0.04) (1.078,2.26)rectangle++(0.04,0.04)
  (1.004,3.474)rectangle++(0.04,0.04) (0.99,4.509)rectangle++(0.04,0.04) (1.02,5.389)rectangle++(0.04,0.04) (1.08,6.136)rectangle++(0.04,0.04) (1.162,6.769)rectangle++(0.04,0.04)
  (1.256,7.303)rectangle++(0.04,0.04) (1.357,7.754)rectangle++(0.04,0.04) (1.462,8.133)rectangle++(0.04,0.04) (1.566,8.45)rectangle++(0.04,0.04) (1.667,8.716)rectangle++(0.04,0.04)
  (1.763,8.939)rectangle++(0.04,0.04) (1.854,9.123)rectangle++(0.04,0.04) (1.939,9.277)rectangle++(0.04,0.04) (2.017,9.404)rectangle++(0.04,0.04) (2.088,9.509)rectangle++(0.04,0.04)
  (2.153,9.595)rectangle++(0.04,0.04) (2.212,9.666)rectangle++(0.04,0.04) (2.264,9.724)rectangle++(0.04,0.04) (2.312,9.771)rectangle++(0.04,0.04) (2.353,9.809)rectangle++(0.04,0.04)
  (2.391,9.84)rectangle++(0.04,0.04) (-2.101,4.304)rectangle++(0.04,0.04) (-0.02,0.672)rectangle++(0.04,0.04) (0.008,2.168)rectangle++(0.04,0.04) (0.091,3.439)rectangle++(0.04,0.04)
  (0.213,4.515)rectangle++(0.04,0.04) (0.359,5.426)rectangle++(0.04,0.04) (0.52,6.194)rectangle++(0.04,0.04) (0.688,6.84)rectangle++(0.04,0.04) (1.795,2.25)rectangle++(0.04,0.04)
  (1.613,3.437)rectangle++(0.04,0.04) (1.507,4.453)rectangle++(0.04,0.04) (1.457,5.321)rectangle++(0.04,0.04) (1.449,6.061)rectangle++(0.04,0.04) (1.472,6.69)rectangle++(0.04,0.04)
  (1.516,7.224)rectangle++(0.04,0.04) (1.576,7.676)rectangle++(0.04,0.04) (1.644,8.058)rectangle++(0.04,0.04) (1.718,8.379)rectangle++(0.04,0.04) (1.793,8.65)rectangle++(0.04,0.04)
  (1.868,8.877)rectangle++(0.04,0.04) (-1.956,3.972)rectangle++(0.04,0.04) (-1.506,5.05)rectangle++(0.04,0.04) (-1.08,5.949)rectangle++(0.04,0.04) (1.81,1.577)rectangle++(0.04,0.04)
  (1.6,2.864)rectangle++(0.04,0.04) (0.545,1.533)rectangle++(0.04,0.04) (0.522,2.878)rectangle++(0.04,0.04) (0.557,4.021)rectangle++(0.04,0.04) (0.632,4.992)rectangle++(0.04,0.04)
  (0.735,5.814)rectangle++(0.04,0.04) (0.855,6.509)rectangle++(0.04,0.04) (0.985,7.095)rectangle++(0.04,0.04) (1.119,7.587)rectangle++(0.04,0.04) (1.252,8.001)rectangle++(0.04,0.04)
  (1.382,8.347)rectangle++(0.04,0.04) (1.506,8.636)rectangle++(0.04,0.04) (1.624,8.876)rectangle++(0.04,0.04) (1.733,9.076)rectangle++(0.04,0.04) (1.834,9.241)rectangle++(0.04,0.04)
  (1.926,9.378)rectangle++(0.04,0.04) (-0.02,1.484)rectangle++(0.04,0.04) (0.04,2.858)rectangle++(0.04,0.04) (-0.756,2.227)rectangle++(0.04,0.04) (-0.556,3.519)rectangle++(0.04,0.04)
  (-0.334,4.61)rectangle++(0.04,0.04) (-0.102,5.528)rectangle++(0.04,0.04) (0.132,6.299)rectangle++(0.04,0.04) (0.362,6.945)rectangle++(0.04,0.04) (0.584,7.485)rectangle++(0.04,0.04)
  (0.793,7.935)rectangle++(0.04,0.04) (0.989,8.309)rectangle++(0.04,0.04) (1.171,8.62)rectangle++(0.04,0.04) (-2.028,3.755)rectangle++(0.04,0.04) (-1.576,4.869)rectangle++(0.04,0.04)
  (-1.147,5.798)rectangle++(0.04,0.04) (-0.745,6.57)rectangle++(0.04,0.04) (-0.373,7.211)rectangle++(0.04,0.04) (-0.031,7.74)rectangle++(0.04,0.04) (0.281,8.177)rectangle++(0.04,0.04)
  (2.23,2.466)rectangle++(0.04,0.04) (1.992,3.603)rectangle++(0.04,0.04) (1.835,4.579)rectangle++(0.04,0.04) (-0.845,3.018)rectangle++(0.04,0.04) (-0.599,4.196)rectangle++(0.04,0.04)
  (-0.344,5.186)rectangle++(0.04,0.04) (-0.087,6.018)rectangle++(0.04,0.04) (0.165,6.715)rectangle++(0.04,0.04) (0.406,7.298)rectangle++(0.04,0.04) (0.635,7.783)rectangle++(0.04,0.04)
  (0.849,8.186)rectangle++(0.04,0.04) (1.047,8.521)rectangle++(0.04,0.04) (1.228,8.797)rectangle++(0.04,0.04) (1.394,9.024)rectangle++(0.04,0.04) (2.3,2.958)rectangle++(0.04,0.04)
  (2.071,4.019)rectangle++(0.04,0.04) (1.919,4.929)rectangle++(0.04,0.04) (-0.919,3.115)rectangle++(0.04,0.04) (-0.659,4.281)rectangle++(0.04,0.04) (-0.391,5.261)rectangle++(0.04,0.04)
  (-1.467,2.657)rectangle++(0.04,0.04) (-1.143,3.913)rectangle++(0.04,0.04) (-0.817,4.968)rectangle++(0.04,0.04) (-0.498,5.852)rectangle++(0.04,0.04) (-0.192,6.59)rectangle++(0.04,0.04)
  (0.099,7.205)rectangle++(0.04,0.04) (0.37,7.717)rectangle++(0.04,0.04) (0.621,8.141)rectangle++(0.04,0.04) (0.851,8.491)rectangle++(0.04,0.04) (1.061,8.779)rectangle++(0.04,0.04)
  (1.251,9.016)rectangle++(0.04,0.04) (1.422,9.21)rectangle++(0.04,0.04) (1.575,9.368)rectangle++(0.04,0.04) (1.711,9.496)rectangle++(0.04,0.04) (1.832,9.599)rectangle++(0.04,0.04)
  (1.939,9.682)rectangle++(0.04,0.04) (-2.151,4.165)rectangle++(0.04,0.04) (1.471,0.87)rectangle++(0.04,0.04) (1.283,2.277)rectangle++(0.04,0.04) (1.18,3.48)rectangle++(0.04,0.04)
  (1.14,4.507)rectangle++(0.04,0.04) (1.147,5.382)rectangle++(0.04,0.04) (1.317,2.02)rectangle++(0.04,0.04) (1.198,3.26)rectangle++(0.04,0.04) (1.147,4.32)rectangle++(0.04,0.04)
  (1.145,5.222)rectangle++(0.04,0.04) (1.18,5.989)rectangle++(0.04,0.04) (1.241,6.64)rectangle++(0.04,0.04) (1.318,7.19)rectangle++(0.04,0.04) (1.406,7.655)rectangle++(0.04,0.04)
  (1.499,8.047)rectangle++(0.04,0.04) (1.594,8.376)rectangle++(0.04,0.04) (1.687,8.652)rectangle++(0.04,0.04) (1.778,8.883)rectangle++(0.04,0.04) (1.865,9.076)rectangle++(0.04,0.04)
  (1.946,9.236)rectangle++(0.04,0.04) (2.021,9.369)rectangle++(0.04,0.04) (2.091,9.479)rectangle++(0.04,0.04) (2.154,9.57)rectangle++(0.04,0.04) (2.211,9.644)rectangle++(0.04,0.04)
  (-2.086,4.219)rectangle++(0.04,0.04) (-1.607,5.266)rectangle++(0.04,0.04) (-1.157,6.137)rectangle++(0.04,0.04) (-0.741,6.859)rectangle++(0.04,0.04) (-0.357,7.456)rectangle++(0.04,0.04)
  (2.124,2.126)rectangle++(0.04,0.04) (-0.149,2.545)rectangle++(0.04,0.04) (-0.027,3.766)rectangle++(0.04,0.04) (0.125,4.798)rectangle++(0.04,0.04) (0.296,5.67)rectangle++(0.04,0.04)
  (0.476,6.404)rectangle++(0.04,0.04) (0.659,7.02)rectangle++(0.04,0.04) (0.839,7.537)rectangle++(0.04,0.04) (1.012,7.969)rectangle++(0.04,0.04) (1.177,8.329)rectangle++(0.04,0.04)
  (1.331,8.629)rectangle++(0.04,0.04) (1.475,8.878)rectangle++(0.04,0.04) (2.247,2.944)rectangle++(0.04,0.04) (0.47,1.721)rectangle++(0.04,0.04) (0.466,3.04)rectangle++(0.04,0.04)
  (0.516,4.162)rectangle++(0.04,0.04) (0.602,5.113)rectangle++(0.04,0.04) (0.714,5.918)rectangle++(0.04,0.04) (0.842,6.598)rectangle++(0.04,0.04) (0.977,7.171)rectangle++(0.04,0.04)
  (1.115,7.652)rectangle++(0.04,0.04) (1.252,8.056)rectangle++(0.04,0.04) (1.384,8.394)rectangle++(0.04,0.04) (1.51,8.676)rectangle++(0.04,0.04) (1.628,8.91)rectangle++(0.04,0.04)
  (1.738,9.105)rectangle++(0.04,0.04) (1.84,9.266)rectangle++(0.04,0.04) (1.932,9.398)rectangle++(0.04,0.04) (2.016,9.508)rectangle++(0.04,0.04) (2.092,9.597)rectangle++(0.04,0.04)
  (2.16,9.67)rectangle++(0.04,0.04) (2.22,9.729)rectangle++(0.04,0.04) (2.374,3.342)rectangle++(0.04,0.04) (0.562,1.849)rectangle++(0.04,0.04) (0.55,3.146)rectangle++(0.04,0.04)
  (0.591,4.248)rectangle++(0.04,0.04) (0.67,5.183)rectangle++(0.04,0.04) (1.333,1.848)rectangle++(0.04,0.04) (-0.02,0.279)rectangle++(0.04,0.04) (-0.008,1.834)rectangle++(0.04,0.04)
  (0.064,3.155)rectangle++(0.04,0.04) (0.179,4.276)rectangle++(0.04,0.04) (0.321,5.223)rectangle++(0.04,0.04) (0.479,6.023)rectangle++(0.04,0.04) (0.646,6.697)rectangle++(0.04,0.04)
  (0.815,7.263)rectangle++(0.04,0.04) (1.894,2.385)rectangle++(0.04,0.04) (1.703,3.548)rectangle++(0.04,0.04) (1.587,4.544)rectangle++(0.04,0.04) (1.529,5.395)rectangle++(0.04,0.04)
  (-1.118,3.127)rectangle++(0.04,0.04) (-0.827,4.299)rectangle++(0.04,0.04) (-0.534,5.284)rectangle++(0.04,0.04) (-0.244,6.109)rectangle++(0.04,0.04) (0.034,6.798)rectangle++(0.04,0.04)
  (0.299,7.373)rectangle++(0.04,0.04) (0.547,7.852)rectangle++(0.04,0.04) (0.777,8.248)rectangle++(0.04,0.04) (0.988,8.576)rectangle++(0.04,0.04) (1.181,8.846)rectangle++(0.04,0.04)
  (1.355,9.068)rectangle++(0.04,0.04) (1.512,9.25)rectangle++(0.04,0.04) (-2.124,3.972)rectangle++(0.04,0.04) (-1.649,5.057)rectangle++(0.04,0.04) (-1.201,5.961)rectangle++(0.04,0.04)
  (-0.785,6.711)rectangle++(0.04,0.04) (-0.401,7.332)rectangle++(0.04,0.04) (-0.05,7.844)rectangle++(0.04,0.04) (0.269,8.266)rectangle++(0.04,0.04) (0.557,8.611)rectangle++(0.04,0.04)
  (0.816,8.894)rectangle++(0.04,0.04) (1.047,9.123)rectangle++(0.04,0.04) (1.253,9.309)rectangle++(0.04,0.04) (-2.191,3.925)rectangle++(0.04,0.04) (-1.708,5.02)rectangle++(0.04,0.04)
  (-1.253,5.932)rectangle++(0.04,0.04) (-0.83,6.688)rectangle++(0.04,0.04) (-0.44,7.314)rectangle++(0.04,0.04) (-0.084,7.831)rectangle++(0.04,0.04) (0.24,8.256)rectangle++(0.04,0.04)
  (0.532,8.604)rectangle++(0.04,0.04) (0.794,8.888)rectangle++(0.04,0.04) (1.028,9.12)rectangle++(0.04,0.04) (1.237,9.307)rectangle++(0.04,0.04) (1.421,9.457)rectangle++(0.04,0.04)
  (1.584,9.578)rectangle++(0.04,0.04) (1.727,9.674)rectangle++(0.04,0.04) (1.853,9.75)rectangle++(0.04,0.04) (1.963,9.81)rectangle++(0.04,0.04) (2.059,9.856)rectangle++(0.04,0.04)
  (2.142,9.891)rectangle++(0.04,0.04) (2.214,9.918)rectangle++(0.04,0.04) (-2.157,4.28)rectangle++(0.04,0.04) (-1.665,5.321)rectangle++(0.04,0.04) (-1.204,6.185)rectangle++(0.04,0.04)
  (-0.778,6.902)rectangle++(0.04,0.04) (-0.388,7.494)rectangle++(0.04,0.04) (2.139,2.128)rectangle++(0.04,0.04) (1.901,3.319)rectangle++(0.04,0.04) (1.747,4.341)rectangle++(0.04,0.04)
  (1.656,5.217)rectangle++(0.04,0.04) (1.614,5.964)rectangle++(0.04,0.04) (1.608,6.601)rectangle++(0.04,0.04) (1.629,7.143)rectangle++(0.04,0.04) (1.668,7.602)rectangle++(0.04,0.04)
  (1.861,2.688)rectangle++(0.04,0.04) (1.687,3.807)rectangle++(0.04,0.04) (1.584,4.764)rectangle++(0.04,0.04) (1.535,5.583)rectangle++(0.04,0.04) (1.526,6.28)rectangle++(0.04,0.04)
  (1.546,6.873)rectangle++(0.04,0.04) (1.587,7.377)rectangle++(0.04,0.04) (1.642,7.803)rectangle++(0.04,0.04) (1.705,8.163)rectangle++(0.04,0.04) (1.774,8.466)rectangle++(0.04,0.04)
  (1.844,8.722)rectangle++(0.04,0.04) (1.914,8.936)rectangle++(0.04,0.04) (1.982,9.115)rectangle++(0.04,0.04) (2.047,9.265)rectangle++(0.04,0.04) (2.109,9.389)rectangle++(0.04,0.04)
  (2.166,9.493)rectangle++(0.04,0.04) (2.218,9.578)rectangle++(0.04,0.04) (2.267,9.649)rectangle++(0.04,0.04) (2.31,9.707)rectangle++(0.04,0.04) (2.35,9.755)rectangle++(0.04,0.04)
  (2.385,9.794)rectangle++(0.04,0.04) (-2.091,4.292)rectangle++(0.04,0.04) (-1.607,5.328)rectangle++(0.04,0.04) (-1.155,6.19)rectangle++(0.04,0.04) (1.889,1.615)rectangle++(0.04,0.04)
  (1.668,2.893)rectangle++(0.04,0.04) (1.531,3.989)rectangle++(0.04,0.04) (1.459,4.926)rectangle++(0.04,0.04) (1.435,5.725)rectangle++(0.04,0.04) (1.447,6.405)rectangle++(0.04,0.04)
  (1.484,6.982)rectangle++(0.04,0.04) (1.538,7.472)rectangle++(0.04,0.04) (1.604,7.886)rectangle++(0.04,0.04) (1.677,8.235)rectangle++(0.04,0.04) (1.752,8.529)rectangle++(0.04,0.04)
  (1.828,8.776)rectangle++(0.04,0.04) (1.903,8.982)rectangle++(0.04,0.04) (1.975,9.155)rectangle++(0.04,0.04) (2.042,9.299)rectangle++(0.04,0.04) (2.106,9.419)rectangle++(0.04,0.04)
  (2.165,9.518)rectangle++(0.04,0.04) (2.218,9.6)rectangle++(0.04,0.04) (2.267,9.667)rectangle++(0.04,0.04) (2.312,9.723)rectangle++(0.04,0.04) (-2.087,4.26)rectangle++(0.04,0.04)
  (-1.606,5.3)rectangle++(0.04,0.04) (-1.155,6.166)rectangle++(0.04,0.04) (-0.737,6.883)rectangle++(0.04,0.04) (-0.354,7.476)rectangle++(0.04,0.04) (-0.004,7.965)rectangle++(0.04,0.04)
  (0.313,8.367)rectangle++(0.04,0.04) (0.599,8.696)rectangle++(0.04,0.04) (0.855,8.963)rectangle++(0.04,0.04) (1.083,9.181)rectangle++(0.04,0.04) (1.285,9.357)rectangle++(0.04,0.04)
  (1.465,9.498)rectangle++(0.04,0.04) (1.623,9.611)rectangle++(0.04,0.04) (1.762,9.701)rectangle++(0.04,0.04) (1.883,9.771)rectangle++(0.04,0.04) (1.989,9.826)rectangle++(0.04,0.04)
  (2.082,9.869)rectangle++(0.04,0.04) (-0.02,1.562)rectangle++(0.04,0.04) (0.043,2.925)rectangle++(0.04,0.04) (0.152,4.081)rectangle++(0.04,0.04) (0.29,5.059)rectangle++(0.04,0.04)
  (0.447,5.884)rectangle++(0.04,0.04) (-1.462,2.986)rectangle++(0.04,0.04) (-1.09,1.91)rectangle++(0.04,0.04) (-0.852,3.263)rectangle++(0.04,0.04) (-0.596,4.404)rectangle++(0.04,0.04)
  (-0.333,5.363)rectangle++(0.04,0.04) (-0.071,6.168)rectangle++(0.04,0.04) (0.185,6.842)rectangle++(0.04,0.04) (0.428,7.405)rectangle++(0.04,0.04) (-0.02,1.168)rectangle++(0.04,0.04)
  (0.028,2.59)rectangle++(0.04,0.04) (0.125,3.796)rectangle++(0.04,0.04) (0.256,4.818)rectangle++(0.04,0.04) (0.408,5.681)rectangle++(0.04,0.04) (-0.02,0.892)rectangle++(0.04,0.04)
  (0.016,2.355)rectangle++(0.04,0.04) (0.106,3.598)rectangle++(0.04,0.04) (0.232,4.65)rectangle++(0.04,0.04) (1.25,1.606)rectangle++(0.04,0.04) (1.124,2.912)rectangle++(0.04,0.04)
  (1.07,4.026)rectangle++(0.04,0.04) (1.068,4.976)rectangle++(0.04,0.04) (1.105,5.783)rectangle++(0.04,0.04) (1.168,6.467)rectangle++(0.04,0.04) (1.25,7.047)rectangle++(0.04,0.04)
  (-1.603,3.427)rectangle++(0.04,0.04) (-1.228,4.573)rectangle++(0.04,0.04) (-0.863,5.532)rectangle++(0.04,0.04) (-0.515,6.333)rectangle++(0.04,0.04) (-0.186,7)rectangle++(0.04,0.04)
  (0.119,7.554)rectangle++(0.04,0.04) (0.402,8.012)rectangle++(0.04,0.04) (0.66,8.39)rectangle++(0.04,0.04) (0.894,8.702)rectangle++(0.04,0.04) (1.106,8.957)rectangle++(0.04,0.04)
  (1.296,9.165)rectangle++(0.04,0.04) (1.466,9.335)rectangle++(0.04,0.04) (1.617,9.472)rectangle++(0.04,0.04) (2.392,3.124)rectangle++(0.04,0.04) (2.156,4.156)rectangle++(0.04,0.04)
  (1.997,5.042)rectangle++(0.04,0.04) (1.897,5.802)rectangle++(0.04,0.04) (1.842,6.452)rectangle++(0.04,0.04) (1.822,7.007)rectangle++(0.04,0.04) (1.827,7.479)rectangle++(0.04,0.04)
  (1.803,2.7)rectangle++(0.04,0.04) (1.638,3.819)rectangle++(0.04,0.04) (1.543,4.777)rectangle++(0.04,0.04) (1.5,5.595)rectangle++(0.04,0.04) (1.497,6.292)rectangle++(0.04,0.04)
  (1.522,6.884)rectangle++(0.04,0.04) (1.567,7.387)rectangle++(0.04,0.04) (1.625,7.812)rectangle++(0.04,0.04) (1.692,8.172)rectangle++(0.04,0.04) (1.763,8.475)rectangle++(0.04,0.04)
  (1.835,8.729)rectangle++(0.04,0.04) (1.907,8.943)rectangle++(0.04,0.04) (1.976,9.121)rectangle++(0.04,0.04) (2.042,9.27)rectangle++(0.04,0.04) (2.105,9.394)rectangle++(0.04,0.04)
  (2.162,9.497)rectangle++(0.04,0.04) (2.216,9.582)rectangle++(0.04,0.04) (2.265,9.652)rectangle++(0.04,0.04) (2.309,9.71)rectangle++(0.04,0.04) (2.349,9.757)rectangle++(0.04,0.04)
  (2.384,9.796)rectangle++(0.04,0.04) (2.416,9.828)rectangle++(0.04,0.04) (2.445,9.853)rectangle++(0.04,0.04) (2.47,9.873)rectangle++(0.04,0.04) (2.492,9.89)rectangle++(0.04,0.04)
  (2.512,9.903)rectangle++(0.04,0.04) (2.379,3.46)rectangle++(0.04,0.04) (2.158,4.442)rectangle++(0.04,0.04) (2.01,5.285)rectangle++(0.04,0.04) (1.918,6.008)rectangle++(0.04,0.04)
  (1.868,6.627)rectangle++(0.04,0.04) (1.851,7.154)rectangle++(0.04,0.04) (1.857,7.603)rectangle++(0.04,0.04) (1.88,7.985)rectangle++(0.04,0.04) (1.916,8.308)rectangle++(0.04,0.04)
  (1.958,8.581)rectangle++(0.04,0.04) (2.006,8.812)rectangle++(0.04,0.04) (2.055,9.006)rectangle++(0.04,0.04) (2.196,3.126)rectangle++(0.04,0.04) (0.528,1.751)rectangle++(0.04,0.04)
  (0.394,0.988)rectangle++(0.04,0.04) (-0.199,1.897)rectangle++(0.04,0.04) (-0.554,1.96)rectangle++(0.04,0.04) (-0.395,3.285)rectangle++(0.04,0.04) (-0.206,4.404)rectangle++(0.04,0.04)
  (-0.002,5.348)rectangle++(0.04,0.04) (0.21,6.142)rectangle++(0.04,0.04) (0.422,6.808)rectangle++(0.04,0.04) (0.629,7.367)rectangle++(0.04,0.04) (0.827,7.833)rectangle++(0.04,0.04)
  (1.014,8.221)rectangle++(0.04,0.04) (1.189,8.543)rectangle++(0.04,0.04) (1.35,8.81)rectangle++(0.04,0.04) (1.498,9.031)rectangle++(0.04,0.04) (1.632,9.213)rectangle++(0.04,0.04)
  (1.754,9.362)rectangle++(0.04,0.04) (1.863,9.484)rectangle++(0.04,0.04) (2.359,3.19)rectangle++(0.04,0.04) (2.13,4.214)rectangle++(0.04,0.04) (1.977,5.093)rectangle++(0.04,0.04)
  (-0.95,3.164)rectangle++(0.04,0.04) (-0.683,4.324)rectangle++(0.04,0.04) (-0.41,5.299)rectangle++(0.04,0.04) (-0.139,6.116)rectangle++(0.04,0.04) (0.125,6.801)rectangle++(0.04,0.04)
  (0.376,7.372)rectangle++(0.04,0.04) (0.612,7.847)rectangle++(0.04,0.04) (0.832,8.242)rectangle++(0.04,0.04) (1.035,8.568)rectangle++(0.04,0.04) (1.22,8.838)rectangle++(0.04,0.04)
  (1.388,9.06)rectangle++(0.04,0.04) (1.54,9.241)rectangle++(0.04,0.04) (1.677,9.39)rectangle++(0.04,0.04) (1.799,9.51)rectangle++(0.04,0.04) (1.907,9.608)rectangle++(0.04,0.04)
  (2.003,9.687)rectangle++(0.04,0.04) (2.088,9.75)rectangle++(0.04,0.04) (2.162,9.8)rectangle++(0.04,0.04) (2.228,9.84)rectangle++(0.04,0.04) (2.285,9.871)rectangle++(0.04,0.04)
  (2.335,9.895)rectangle++(0.04,0.04) (2.378,9.914)rectangle++(0.04,0.04) (2.416,9.928)rectangle++(0.04,0.04) (2.448,9.938)rectangle++(0.04,0.04) (2.477,9.946)rectangle++(0.04,0.04)
  (2.501,9.951)rectangle++(0.04,0.04) (2.521,9.954)rectangle++(0.04,0.04) (2.539,9.957)rectangle++(0.04,0.04) (2.554,9.958)rectangle++(0.04,0.04) (2.388,3.484)rectangle++(0.04,0.04)
  (2.167,4.462)rectangle++(0.04,0.04) (2.018,5.302)rectangle++(0.04,0.04) (-0.996,3.22)rectangle++(0.04,0.04) (-0.02,0.498)rectangle++(0.04,0.04) (0.001,2.021)rectangle++(0.04,0.04)
  (0.079,3.314)rectangle++(0.04,0.04) (0.198,4.41)rectangle++(0.04,0.04) (0.342,5.336)rectangle++(0.04,0.04) (0.502,6.119)rectangle++(0.04,0.04) (0.669,6.777)rectangle++(0.04,0.04)
  (-1.649,3.234)rectangle++(0.04,0.04) (-1.275,4.411)rectangle++(0.04,0.04) (-0.909,5.397)rectangle++(0.04,0.04) (-0.559,6.22)rectangle++(0.04,0.04) (-0.229,6.905)rectangle++(0.04,0.04)
  (0.08,7.475)rectangle++(0.04,0.04) (0.364,7.947)rectangle++(0.04,0.04) (0.625,8.336)rectangle++(0.04,0.04) (0.863,8.657)rectangle++(0.04,0.04) (1.077,8.92)rectangle++(0.04,0.04)
  (1.27,9.135)rectangle++(0.04,0.04) (1.443,9.31)rectangle++(0.04,0.04) (1.597,9.452)rectangle++(0.04,0.04) (2.39,3.114)rectangle++(0.04,0.04) (2.154,4.147)rectangle++(0.04,0.04)
  (1.994,5.035)rectangle++(0.04,0.04) (1.894,5.796)rectangle++(0.04,0.04) (1.84,6.447)rectangle++(0.04,0.04) (1.82,7.003)rectangle++(0.04,0.04) (1.825,7.476)rectangle++(0.04,0.04)
  (1.848,7.878)rectangle++(0.04,0.04) (1.883,8.218)rectangle++(0.04,0.04) (-1.781,3.83)rectangle++(0.04,0.04) (-1.363,4.923)rectangle++(0.04,0.04) (-0.964,5.835)rectangle++(0.04,0.04)
  (-1.731,2.651)rectangle++(0.04,0.04) (-1.368,3.919)rectangle++(0.04,0.04) (-1.008,4.982)rectangle++(0.04,0.04) (-0.66,5.871)rectangle++(0.04,0.04) (-0.328,6.613)rectangle++(0.04,0.04)
  (-0.017,7.23)rectangle++(0.04,0.04) (0.273,7.743)rectangle++(0.04,0.04) (0.54,8.167)rectangle++(0.04,0.04) (2.188,2.53)rectangle++(0.04,0.04) (1.959,3.659)rectangle++(0.04,0.04)
  (1.809,4.628)rectangle++(0.04,0.04) (1.721,5.458)rectangle++(0.04,0.04) (1.679,6.167)rectangle++(0.04,0.04) (1.672,6.771)rectangle++(0.04,0.04) (1.628,2.49)rectangle++(0.04,0.04)
  (1.481,3.647)rectangle++(0.04,0.04) (0.782,1.69)rectangle++(0.04,0.04) (0.339,1.039)rectangle++(0.04,0.04) (0.327,2.466)rectangle++(0.04,0.04) (0.375,3.679)rectangle++(0.04,0.04)
  (0.957,1.41)rectangle++(0.04,0.04) (0.867,2.757)rectangle++(0.04,0.04) (0.624,1.317)rectangle++(0.04,0.04) (-0.239,2.022)rectangle++(0.04,0.04) (-0.124,3.325)rectangle++(0.04,0.04)
  (0.932,1.196)rectangle++(0.04,0.04) (0.838,2.575)rectangle++(0.04,0.04) (0.813,3.752)rectangle++(0.04,0.04) (0.839,4.753)rectangle++(0.04,0.04) (0.901,5.602)rectangle++(0.04,0.04)
  (0.988,6.322)rectangle++(0.04,0.04) (1.09,6.93)rectangle++(0.04,0.04) (1.202,7.444)rectangle++(0.04,0.04) (1.317,7.875)rectangle++(0.04,0.04) (1.432,8.237)rectangle++(0.04,0.04)
  (1.545,8.541)rectangle++(0.04,0.04) (-1.933,3.823)rectangle++(0.04,0.04) (-1.492,4.923)rectangle++(0.04,0.04) (1.585,1.224)rectangle++(0.04,0.04) (1.394,2.573)rectangle++(0.04,0.04)
  (1.286,3.727)rectangle++(0.04,0.04) (1.24,4.713)rectangle++(0.04,0.04) (1.116,1.883)rectangle++(0.04,0.04) (1.022,3.152)rectangle++(0.04,0.04) (0.993,4.235)rectangle++(0.04,0.04)
  (1.011,5.156)rectangle++(0.04,0.04) (1.063,5.939)rectangle++(0.04,0.04) (1.139,6.601)rectangle++(0.04,0.04) (1.23,7.162)rectangle++(0.04,0.04) (-1.637,3.448)rectangle++(0.04,0.04)
  (-1.256,4.592)rectangle++(0.04,0.04) (-0.886,5.55)rectangle++(0.04,0.04) (-0.533,6.349)rectangle++(0.04,0.04) (1.84,1.815)rectangle++(0.04,0.04) (1.635,3.065)rectangle++(0.04,0.04)
  (1.51,4.136)rectangle++(0.04,0.04) (1.447,5.052)rectangle++(0.04,0.04) (1.43,5.832)rectangle++(0.04,0.04) (1.446,6.496)rectangle++(0.04,0.04) (1.487,7.06)rectangle++(0.04,0.04)
  (1.544,7.538)rectangle++(0.04,0.04) (1.612,7.942)rectangle++(0.04,0.04) (1.686,8.282)rectangle++(0.04,0.04) (1.762,8.569)rectangle++(0.04,0.04) (1.838,8.809)rectangle++(0.04,0.04)
  (1.913,9.01)rectangle++(0.04,0.04) (1.984,9.179)rectangle++(0.04,0.04) (2.051,9.319)rectangle++(0.04,0.04) (2.114,9.435)rectangle++(0.04,0.04) (2.172,9.531)rectangle++(0.04,0.04)
  (-2.065,4.185)rectangle++(0.04,0.04) (-1.59,5.236)rectangle++(0.04,0.04) (-1.144,6.111)rectangle++(0.04,0.04) (-0.73,6.836)rectangle++(0.04,0.04) (-0.35,7.436)rectangle++(0.04,0.04)
  (-0.002,7.931)rectangle++(0.04,0.04) (0.313,8.338)rectangle++(0.04,0.04) (0.598,8.671)rectangle++(0.04,0.04) (0.853,8.942)rectangle++(0.04,0.04) (1.08,9.163)rectangle++(0.04,0.04)
  (1.283,9.342)rectangle++(0.04,0.04) (1.462,9.485)rectangle++(0.04,0.04) (1.62,9.6)rectangle++(0.04,0.04) (2.428,3.155)rectangle++(0.04,0.04) (2.188,4.181)rectangle++(0.04,0.04)
  (-0.671,3.012)rectangle++(0.04,0.04) (-0.452,4.183)rectangle++(0.04,0.04) (-0.219,5.17)rectangle++(0.04,0.04) (0.018,5.999)rectangle++(0.04,0.04) (1.66,1.875)rectangle++(0.04,0.04)
  (1.484,3.123)rectangle++(0.04,0.04) (1.384,4.192)rectangle++(0.04,0.04) (1.342,5.104)rectangle++(0.04,0.04) (-1.08,3.02)rectangle++(0.04,0.04) (-1.022,2.005)rectangle++(0.04,0.04)
  (-0.791,3.341)rectangle++(0.04,0.04) (-0.541,4.468)rectangle++(0.04,0.04) (-1.291,2.448)rectangle++(0.04,0.04) (-1.002,3.728)rectangle++(0.04,0.04) (-0.705,4.805)rectangle++(0.04,0.04)
  (-0.409,5.709)rectangle++(0.04,0.04) (-0.121,6.465)rectangle++(0.04,0.04) (0.153,7.096)rectangle++(0.04,0.04) (0.412,7.622)rectangle++(0.04,0.04) (0.653,8.058)rectangle++(0.04,0.04)
  (0.875,8.42)rectangle++(0.04,0.04) (2.209,2.678)rectangle++(0.04,0.04) (1.982,3.784)rectangle++(0.04,0.04) (1.834,4.734)rectangle++(0.04,0.04) (1.746,5.546)rectangle++(0.04,0.04)
  (1.704,6.241)rectangle++(0.04,0.04) (1.696,6.833)rectangle++(0.04,0.04) (1.713,7.336)rectangle++(0.04,0.04) (1.747,7.763)rectangle++(0.04,0.04) (1.793,8.125)rectangle++(0.04,0.04)
  (1.847,8.431)rectangle++(0.04,0.04) (1.905,8.689)rectangle++(0.04,0.04) (-1.899,3.939)rectangle++(0.04,0.04) (-1.459,5.02)rectangle++(0.04,0.04) (1.607,1.255)rectangle++(0.04,0.04)
  (-0.026,2.235)rectangle++(0.04,0.04) (-0.607,2.075)rectangle++(0.04,0.04) (-0.436,3.384)rectangle++(0.04,0.04) (-0.237,4.49)rectangle++(0.04,0.04) (-0.024,5.422)rectangle++(0.04,0.04)
  (0.194,6.206)rectangle++(0.04,0.04) (0.411,6.864)rectangle++(0.04,0.04) (0.622,7.414)rectangle++(0.04,0.04) (0.823,7.873)rectangle++(0.04,0.04) (1.012,8.255)rectangle++(0.04,0.04)
  (1.188,8.573)rectangle++(0.04,0.04) (1.351,8.836)rectangle++(0.04,0.04) (1.499,9.052)rectangle++(0.04,0.04) (1.634,9.231)rectangle++(0.04,0.04) (1.756,9.377)rectangle++(0.04,0.04)
  (1.866,9.496)rectangle++(0.04,0.04) (1.963,9.593)rectangle++(0.04,0.04) (2.051,9.672)rectangle++(0.04,0.04) (2.128,9.735)rectangle++(0.04,0.04) (2.196,9.786)rectangle++(0.04,0.04)
  (2.256,9.827)rectangle++(0.04,0.04) (2.308,9.859)rectangle++(0.04,0.04) (2.354,9.884)rectangle++(0.04,0.04) (2.394,9.903)rectangle++(0.04,0.04) (2.429,9.918)rectangle++(0.04,0.04)
  (2.459,9.929)rectangle++(0.04,0.04) (2.485,9.938)rectangle++(0.04,0.04) (2.508,9.944)rectangle++(0.04,0.04) (2.527,9.948)rectangle++(0.04,0.04) (2.544,9.951)rectangle++(0.04,0.04)
  (2.387,3.48)rectangle++(0.04,0.04) (2.166,4.458)rectangle++(0.04,0.04) (-0.747,3.068)rectangle++(0.04,0.04) (-0.515,4.234)rectangle++(0.04,0.04) (-0.27,5.216)rectangle++(0.04,0.04)
  (-0.023,6.04)rectangle++(0.04,0.04) (0.22,6.731)rectangle++(0.04,0.04) (0.454,7.309)rectangle++(0.04,0.04) (0.676,7.791)rectangle++(0.04,0.04) (0.884,8.191)rectangle++(0.04,0.04)
  (1.077,8.523)rectangle++(0.04,0.04) (1.254,8.798)rectangle++(0.04,0.04) (1.416,9.024)rectangle++(0.04,0.04) (1.562,9.21)rectangle++(0.04,0.04) (1.694,9.362)rectangle++(0.04,0.04)
  (1.812,9.487)rectangle++(0.04,0.04) (1.918,9.587)rectangle++(0.04,0.04) (2.011,9.669)rectangle++(0.04,0.04) (2.388,3.273)rectangle++(0.04,0.04) (2.159,4.283)rectangle++(0.04,0.04)
  (2.004,5.15)rectangle++(0.04,0.04) (1.907,5.894)rectangle++(0.04,0.04) (1.855,6.53)rectangle++(0.04,0.04) (1.835,7.072)rectangle++(0.04,0.04) (1.841,7.534)rectangle++(0.04,0.04)
  (1.864,7.927)rectangle++(0.04,0.04) (1.899,8.259)rectangle++(0.04,0.04) (1.942,8.541)rectangle++(0.04,0.04) (1.991,8.778)rectangle++(0.04,0.04) (2.041,8.978)rectangle++(0.04,0.04)
  (2.092,9.146)rectangle++(0.04,0.04) (2.142,9.286)rectangle++(0.04,0.04) (2.19,9.404)rectangle++(0.04,0.04) (2.235,9.502)rectangle++(0.04,0.04) (2.278,9.584)rectangle++(0.04,0.04)
  (2.317,9.651)rectangle++(0.04,0.04) (2.353,9.707)rectangle++(0.04,0.04) (2.348,3.372)rectangle++(0.04,0.04) (-0.428,2.871)rectangle++(0.04,0.04) (-0.251,4.053)rectangle++(0.04,0.04)
  (-0.054,5.052)rectangle++(0.04,0.04) (0.154,5.892)rectangle++(0.04,0.04) (0.365,6.598)rectangle++(0.04,0.04) (0.572,7.19)rectangle++(0.04,0.04) (0.771,7.685)rectangle++(0.04,0.04)
  (2.019,2.475)rectangle++(0.04,0.04) (1.813,3.619)rectangle++(0.04,0.04) (1.683,4.6)rectangle++(0.04,0.04) (1.613,5.439)rectangle++(0.04,0.04) (1.586,6.155)rectangle++(0.04,0.04)
  (1.592,6.764)rectangle++(0.04,0.04) (1.622,7.282)rectangle++(0.04,0.04) (1.668,7.721)rectangle++(0.04,0.04) (1.724,8.092)rectangle++(0.04,0.04) (1.99,2.82)rectangle++(0.04,0.04)
  (1.802,3.914)rectangle++(0.04,0.04) (1.686,4.851)rectangle++(0.04,0.04) (1.625,5.652)rectangle++(0.04,0.04) (-1.166,3.206)rectangle++(0.04,0.04) (-0.865,4.368)rectangle++(0.04,0.04)
  (-0.563,5.344)rectangle++(0.04,0.04) (-0.02,0.838)rectangle++(0.04,0.04) (0.014,2.309)rectangle++(0.04,0.04) (0.102,3.559)rectangle++(0.04,0.04) (0.227,4.617)rectangle++(0.04,0.04)
  (0.376,5.512)rectangle++(0.04,0.04) (0.537,6.266)rectangle++(0.04,0.04) (0.705,6.901)rectangle++(0.04,0.04) (-1.674,3.269)rectangle++(0.04,0.04) (-1.295,4.442)rectangle++(0.04,0.04)
  (1.421,1.16)rectangle++(0.04,0.04) (-0.039,2.171)rectangle++(0.04,0.04) (0.052,3.443)rectangle++(0.04,0.04) (0.18,4.521)rectangle++(0.04,0.04) (0.331,5.432)rectangle++(0.04,0.04)
  (0.497,6.2)rectangle++(0.04,0.04) (0.668,6.846)rectangle++(0.04,0.04) (0.839,7.389)rectangle++(0.04,0.04) (1.007,7.843)rectangle++(0.04,0.04) (1.167,8.222)rectangle++(0.04,0.04)
  (1.319,8.539)rectangle++(0.04,0.04) (1.46,8.801)rectangle++(0.04,0.04) (2.228,2.922)rectangle++(0.04,0.04) (2.008,3.991)rectangle++(0.04,0.04) (1.865,4.908)rectangle++(0.04,0.04)
  (1.779,5.693)rectangle++(0.04,0.04) (1.31,2.259)rectangle++(0.04,0.04) (1.202,3.464)rectangle++(0.04,0.04) (1.158,4.492)rectangle++(0.04,0.04) (-0.958,2.844)rectangle++(0.04,0.04)
  (-0.703,4.052)rectangle++(0.04,0.04) (-0.437,5.068)rectangle++(0.04,0.04) (-0.171,5.922)rectangle++(0.04,0.04) (0.089,6.636)rectangle++(0.04,0.04) (0.339,7.234)rectangle++(0.04,0.04)
  (0.575,7.731)rectangle++(0.04,0.04) (0.796,8.145)rectangle++(0.04,0.04) (1,8.487)rectangle++(0.04,0.04) (1.188,8.77)rectangle++(0.04,0.04) (1.358,9.004)rectangle++(0.04,0.04)
  (1.512,9.195)rectangle++(0.04,0.04) (1.651,9.351)rectangle++(0.04,0.04) (1.775,9.479)rectangle++(0.04,0.04) (1.886,9.582)rectangle++(0.04,0.04) (1.984,9.666)rectangle++(0.04,0.04)
  (2.071,9.733)rectangle++(0.04,0.04) (2.397,3.304)rectangle++(0.04,0.04) (2.168,4.309)rectangle++(0.04,0.04) (2.013,5.172)rectangle++(0.04,0.04) (1.915,5.912)rectangle++(0.04,0.04)
  (1.862,6.545)rectangle++(0.04,0.04) (1.843,7.085)rectangle++(0.04,0.04) (1.847,7.545)rectangle++(0.04,0.04) (1.87,7.935)rectangle++(0.04,0.04) (1.905,8.266)rectangle++(0.04,0.04)
  (1.947,8.546)rectangle++(0.04,0.04) (1.995,8.783)rectangle++(0.04,0.04) (2.045,8.982)rectangle++(0.04,0.04) (2.095,9.149)rectangle++(0.04,0.04) (2.145,9.289)rectangle++(0.04,0.04)
  (2.192,9.406)rectangle++(0.04,0.04) (2.237,9.504)rectangle++(0.04,0.04) (2.28,9.585)rectangle++(0.04,0.04) (2.319,9.652)rectangle++(0.04,0.04) (2.355,9.708)rectangle++(0.04,0.04)
  (2.388,9.754)rectangle++(0.04,0.04) (2.418,9.791)rectangle++(0.04,0.04) (2.444,9.822)rectangle++(0.04,0.04) (2.366,3.423)rectangle++(0.04,0.04) (2.146,4.411)rectangle++(0.04,0.04)
  (-0.739,3.053)rectangle++(0.04,0.04) (-0.508,4.221)rectangle++(0.04,0.04) (-0.265,5.204)rectangle++(0.04,0.04) (-0.02,6.03)rectangle++(0.04,0.04) (0.222,6.723)rectangle++(0.04,0.04)
  (0.456,7.302)rectangle++(0.04,0.04) (-1.828,3.3)rectangle++(0.04,0.04) (-1.424,4.475)rectangle++(0.04,0.04) (-1.034,5.457)rectangle++(0.04,0.04) (-0.663,6.276)rectangle++(0.04,0.04)
  (-0.315,6.957)rectangle++(0.04,0.04) (0.009,7.522)rectangle++(0.04,0.04) (0.306,7.99)rectangle++(0.04,0.04) (2.174,2.427)rectangle++(0.04,0.04) (1.943,3.572)rectangle++(0.04,0.04)
  (1.792,4.555)rectangle++(0.04,0.04) (1.703,5.396)rectangle++(0.04,0.04) (1.661,6.115)rectangle++(0.04,0.04) (1.655,6.727)rectangle++(0.04,0.04) (1.673,7.248)rectangle++(0.04,0.04)
  (-1.571,3.568)rectangle++(0.04,0.04) (-1.195,4.692)rectangle++(0.04,0.04) (-0.83,5.632)rectangle++(0.04,0.04) (-0.483,6.417)rectangle++(0.04,0.04) (-1.786,2.89)rectangle++(0.04,0.04)
  (-1.405,4.124)rectangle++(0.04,0.04) (-1.031,5.158)rectangle++(0.04,0.04) (-0.673,6.022)rectangle++(0.04,0.04) (-0.333,6.741)rectangle++(0.04,0.04) (-0.016,7.34)rectangle++(0.04,0.04)
  (0.278,7.836)rectangle++(0.04,0.04) (0.548,8.245)rectangle++(0.04,0.04) (0.793,8.583)rectangle++(0.04,0.04) (1.015,8.86)rectangle++(0.04,0.04) (1.215,9.086)rectangle++(0.04,0.04)
  (1.394,9.271)rectangle++(0.04,0.04) (1.554,9.421)rectangle++(0.04,0.04) (1.695,9.542)rectangle++(0.04,0.04) (1.82,9.639)rectangle++(0.04,0.04) (1.931,9.716)rectangle++(0.04,0.04)
  (2.414,3.264)rectangle++(0.04,0.04) (2.18,4.274)rectangle++(0.04,0.04) (2.022,5.142)rectangle++(0.04,0.04) (1.922,5.886)rectangle++(0.04,0.04) (1.867,6.522)rectangle++(0.04,0.04)
  (1.846,7.066)rectangle++(0.04,0.04) (1.849,7.528)rectangle++(0.04,0.04) (1.871,7.921)rectangle++(0.04,0.04) (1.905,8.254)rectangle++(0.04,0.04) (1.947,8.536)rectangle++(0.04,0.04)
  (1.994,8.774)rectangle++(0.04,0.04) (2.14,3.054)rectangle++(0.04,0.04) (-0.387,2.753)rectangle++(0.04,0.04) (-0.221,3.952)rectangle++(0.04,0.04) (1.122,1.321)rectangle++(0.04,0.04)
  (1.005,2.674)rectangle++(0.04,0.04) (0.959,3.829)rectangle++(0.04,0.04) (0.911,1.598)rectangle++(0.04,0.04) (0.836,2.918)rectangle++(0.04,0.04) (0.825,4.043)rectangle++(0.04,0.04)
  (0.861,5)rectangle++(0.04,0.04) (0.93,5.812)rectangle++(0.04,0.04) (1.02,6.499)rectangle++(0.04,0.04) (1.125,7.08)rectangle++(0.04,0.04) (1.796,2.422)rectangle++(0.04,0.04)
  (1.621,3.583)rectangle++(0.04,0.04) (1.519,4.577)rectangle++(0.04,0.04) (1.472,5.426)rectangle++(0.04,0.04) (1.281,2.115)rectangle++(0.04,0.04) (1.171,3.343)rectangle++(0.04,0.04)
  (1.127,4.391)rectangle++(0.04,0.04) (1.131,5.283)rectangle++(0.04,0.04) (1.171,6.042)rectangle++(0.04,0.04) (1.235,6.685)rectangle++(0.04,0.04) (1.315,7.229)rectangle++(0.04,0.04)
  (1.404,7.688)rectangle++(0.04,0.04) (1.499,8.075)rectangle++(0.04,0.04) (1.595,8.4)rectangle++(0.04,0.04) (1.69,8.672)rectangle++(0.04,0.04) (1.781,8.9)rectangle++(0.04,0.04);
\fill[nuage]
  (1.867,9.09)rectangle++(0.04,0.04) (-2.011,4.018)rectangle++(0.04,0.04) (-1.551,5.092)rectangle++(0.04,0.04) (-1.117,5.987)rectangle++(0.04,0.04) (-0.712,6.73)rectangle++(0.04,0.04)
  (-1.913,2.906)rectangle++(0.04,0.04) (-1.512,4.143)rectangle++(0.04,0.04) (-1.122,5.178)rectangle++(0.04,0.04) (-0.749,6.042)rectangle++(0.04,0.04) (-0.397,6.762)rectangle++(0.04,0.04)
  (-0.069,7.36)rectangle++(0.04,0.04) (0.233,7.855)rectangle++(0.04,0.04) (0.51,8.263)rectangle++(0.04,0.04) (2.22,2.546)rectangle++(0.04,0.04) (0.362,1.618)rectangle++(0.04,0.04)
  (0.371,2.957)rectangle++(0.04,0.04) (0.431,4.095)rectangle++(0.04,0.04) (-1,2.589)rectangle++(0.04,0.04) (-0.748,3.837)rectangle++(0.04,0.04) (-0.485,4.887)rectangle++(0.04,0.04)
  (-0.219,5.77)rectangle++(0.04,0.04) (0.043,6.509)rectangle++(0.04,0.04) (0.294,7.127)rectangle++(0.04,0.04) (0.533,7.643)rectangle++(0.04,0.04) (-0.02,1.206)rectangle++(0.04,0.04)
  (0.029,2.622)rectangle++(0.04,0.04) (0.127,3.824)rectangle++(0.04,0.04) (0.259,4.841)rectangle++(0.04,0.04) (0.412,5.701)rectangle++(0.04,0.04) (0.576,6.426)rectangle++(0.04,0.04)
  (-0.02,1.011)rectangle++(0.04,0.04) (0.021,2.457)rectangle++(0.04,0.04) (0.114,3.683)rectangle++(0.04,0.04) (0.242,4.723)rectangle++(0.04,0.04) (0.393,5.601)rectangle++(0.04,0.04)
  (0.555,6.341)rectangle++(0.04,0.04) (0.724,6.964)rectangle++(0.04,0.04) (0.891,7.487)rectangle++(0.04,0.04) (1.055,7.924)rectangle++(0.04,0.04) (1.211,8.29)rectangle++(0.04,0.04)
  (1.359,8.594)rectangle++(0.04,0.04) (1.497,8.847)rectangle++(0.04,0.04) (1.624,9.056)rectangle++(0.04,0.04) (-2.051,3.955)rectangle++(0.04,0.04) (-1.587,5.04)rectangle++(0.04,0.04)
  (1.632,1.227)rectangle++(0.04,0.04) (0.081,1.149)rectangle++(0.04,0.04) (0.113,2.569)rectangle++(0.04,0.04) (0.197,3.776)rectangle++(0.04,0.04) (0.316,4.798)rectangle++(0.04,0.04)
  (0.458,5.662)rectangle++(0.04,0.04) (0.614,6.39)rectangle++(0.04,0.04) (0.775,7.003)rectangle++(0.04,0.04) (-1.687,3.308)rectangle++(0.04,0.04) (-1.304,4.475)rectangle++(0.04,0.04)
  (-0.931,5.453)rectangle++(0.04,0.04) (-0.576,6.268)rectangle++(0.04,0.04) (-0.241,6.947)rectangle++(0.04,0.04) (0.071,7.511)rectangle++(0.04,0.04) (0.358,7.978)rectangle++(0.04,0.04)
  (0.622,8.363)rectangle++(0.04,0.04) (0.861,8.68)rectangle++(0.04,0.04) (2.284,2.737)rectangle++(0.04,0.04) (-0.276,2.716)rectangle++(0.04,0.04) (-0.128,3.916)rectangle++(0.04,0.04)
  (0.046,4.93)rectangle++(0.04,0.04) (0.234,5.785)rectangle++(0.04,0.04) (0.428,6.504)rectangle++(0.04,0.04) (0.622,7.108)rectangle++(0.04,0.04) (0.81,7.613)rectangle++(0.04,0.04)
  (0.991,8.035)rectangle++(0.04,0.04) (2.084,2.616)rectangle++(0.04,0.04) (1.874,3.736)rectangle++(0.04,0.04) (1.74,4.697)rectangle++(0.04,0.04) (1.664,5.519)rectangle++(0.04,0.04)
  (1.633,6.221)rectangle++(0.04,0.04) (1.635,6.819)rectangle++(0.04,0.04) (1.647,2.492)rectangle++(0.04,0.04) (-0.34,2.516)rectangle++(0.04,0.04) (-0.19,3.748)rectangle++(0.04,0.04)
  (-0.014,4.79)rectangle++(0.04,0.04) (0.177,5.668)rectangle++(0.04,0.04) (-1.459,2.877)rectangle++(0.04,0.04) (-1.128,4.1)rectangle++(0.04,0.04) (-0.797,5.126)rectangle++(0.04,0.04)
  (-0.474,5.985)rectangle++(0.04,0.04) (-0.166,6.703)rectangle++(0.04,0.04) (0.125,7.3)rectangle++(0.04,0.04) (0.396,7.796)rectangle++(0.04,0.04) (0.646,8.207)rectangle++(0.04,0.04)
  (0.875,8.547)rectangle++(0.04,0.04) (-2.068,3.671)rectangle++(0.04,0.04) (-1.613,4.799)rectangle++(0.04,0.04) (-1.182,5.74)rectangle++(0.04,0.04) (-0.777,6.522)rectangle++(0.04,0.04)
  (-0.402,7.171)rectangle++(0.04,0.04) (-0.057,7.708)rectangle++(0.04,0.04) (2.149,2.265)rectangle++(0.04,0.04) (1.915,3.436)rectangle++(0.04,0.04) (1.763,4.44)rectangle++(0.04,0.04)
  (-0.823,2.971)rectangle++(0.04,0.04) (-0.583,4.155)rectangle++(0.04,0.04) (-0.331,5.151)rectangle++(0.04,0.04) (-0.078,5.988)rectangle++(0.04,0.04) (0.171,6.689)rectangle++(0.04,0.04)
  (0.411,7.275)rectangle++(0.04,0.04) (0.638,7.763)rectangle++(0.04,0.04) (0.851,8.17)rectangle++(0.04,0.04) (1.048,8.506)rectangle++(0.04,0.04) (1.229,8.785)rectangle++(0.04,0.04)
  (1.393,9.014)rectangle++(0.04,0.04) (1.543,9.202)rectangle++(0.04,0.04) (1.677,9.357)rectangle++(0.04,0.04) (1.798,9.482)rectangle++(0.04,0.04) (1.905,9.584)rectangle++(0.04,0.04)
  (2.001,9.667)rectangle++(0.04,0.04) (2.085,9.733)rectangle++(0.04,0.04) (2.159,9.786)rectangle++(0.04,0.04) (2.225,9.828)rectangle++(0.04,0.04) (2.282,9.861)rectangle++(0.04,0.04)
  (2.332,9.886)rectangle++(0.04,0.04) (2.375,9.906)rectangle++(0.04,0.04) (2.413,9.922)rectangle++(0.04,0.04) (2.446,9.933)rectangle++(0.04,0.04) (2.474,9.942)rectangle++(0.04,0.04)
  (-2.111,4.345)rectangle++(0.04,0.04) (-1.623,5.374)rectangle++(0.04,0.04) (-1.167,6.229)rectangle++(0.04,0.04) (-1.874,2.691)rectangle++(0.04,0.04) (-1.488,3.959)rectangle++(0.04,0.04)
  (-1.108,5.02)rectangle++(0.04,0.04) (1.555,1.347)rectangle++(0.04,0.04) (1.373,2.679)rectangle++(0.04,0.04) (1.272,3.818)rectangle++(0.04,0.04) (-0.76,2.722)rectangle++(0.04,0.04)
  (-0.539,3.94)rectangle++(0.04,0.04) (-0.303,4.967)rectangle++(0.04,0.04) (-0.061,5.83)rectangle++(0.04,0.04) (0.179,6.554)rectangle++(0.04,0.04) (0.412,7.16)rectangle++(0.04,0.04)
  (0.635,7.666)rectangle++(0.04,0.04) (0.844,8.087)rectangle++(0.04,0.04) (1.039,8.436)rectangle++(0.04,0.04) (1.218,8.725)rectangle++(0.04,0.04) (1.382,8.964)rectangle++(0.04,0.04)
  (1.531,9.16)rectangle++(0.04,0.04) (1.666,9.321)rectangle++(0.04,0.04) (1.787,9.453)rectangle++(0.04,0.04) (1.894,9.559)rectangle++(0.04,0.04) (1.99,9.646)rectangle++(0.04,0.04)
  (2.076,9.716)rectangle++(0.04,0.04) (2.151,9.771)rectangle++(0.04,0.04) (2.217,9.816)rectangle++(0.04,0.04) (2.275,9.851)rectangle++(0.04,0.04) (2.325,9.879)rectangle++(0.04,0.04)
  (2.369,9.9)rectangle++(0.04,0.04) (2.408,9.916)rectangle++(0.04,0.04) (2.441,9.929)rectangle++(0.04,0.04) (-0.02,1.572)rectangle++(0.04,0.04) (0.044,2.933)rectangle++(0.04,0.04)
  (0.152,4.088)rectangle++(0.04,0.04) (0.291,5.065)rectangle++(0.04,0.04) (0.447,5.889)rectangle++(0.04,0.04) (-1.463,2.988)rectangle++(0.04,0.04) (-1.126,4.194)rectangle++(0.04,0.04)
  (-0.792,5.206)rectangle++(0.04,0.04) (-0.467,6.053)rectangle++(0.04,0.04) (-0.157,6.76)rectangle++(0.04,0.04) (0.135,7.349)rectangle++(0.04,0.04) (2.02,2.229)rectangle++(0.04,0.04)
  (1.804,3.41)rectangle++(0.04,0.04) (1.668,4.422)rectangle++(0.04,0.04) (-0.838,2.945)rectangle++(0.04,0.04) (-0.596,4.133)rectangle++(0.04,0.04) (-0.344,5.133)rectangle++(0.04,0.04)
  (-0.089,5.973)rectangle++(0.04,0.04) (0.161,6.677)rectangle++(0.04,0.04) (0.402,7.265)rectangle++(0.04,0.04) (0.63,7.756)rectangle++(0.04,0.04) (0.843,8.163)rectangle++(0.04,0.04)
  (1.041,8.501)rectangle++(0.04,0.04) (1.223,8.781)rectangle++(0.04,0.04) (1.389,9.011)rectangle++(0.04,0.04) (1.538,9.2)rectangle++(0.04,0.04) (1.674,9.355)rectangle++(0.04,0.04)
  (1.794,9.481)rectangle++(0.04,0.04) (1.902,9.583)rectangle++(0.04,0.04) (1.998,9.666)rectangle++(0.04,0.04) (-2.135,4.175)rectangle++(0.04,0.04) (-1.65,5.23)rectangle++(0.04,0.04)
  (-1.195,6.108)rectangle++(0.04,0.04) (1.872,1.585)rectangle++(0.04,0.04) (1.652,2.869)rectangle++(0.04,0.04) (1.517,3.968)rectangle++(0.04,0.04) (-0.02,0.618)rectangle++(0.04,0.04)
  (0.006,2.122)rectangle++(0.04,0.04) (0.576,0.941)rectangle++(0.04,0.04) (0.525,2.373)rectangle++(0.04,0.04) (0.539,3.592)rectangle++(0.04,0.04) (-0.847,2.503)rectangle++(0.04,0.04)
  (-0.622,3.758)rectangle++(0.04,0.04) (-0.381,4.815)rectangle++(0.04,0.04) (-0.133,5.704)rectangle++(0.04,0.04) (-1.531,2.813)rectangle++(0.04,0.04) (-1.191,4.049)rectangle++(0.04,0.04)
  (-0.853,5.085)rectangle++(0.04,0.04) (-0.524,5.953)rectangle++(0.04,0.04) (-0.209,6.677)rectangle++(0.04,0.04) (0.087,7.28)rectangle++(0.04,0.04) (0.363,7.781)rectangle++(0.04,0.04)
  (0.618,8.195)rectangle++(0.04,0.04) (0.851,8.537)rectangle++(0.04,0.04) (2.246,2.7)rectangle++(0.04,0.04) (2.014,3.802)rectangle++(0.04,0.04) (0.745,1.866)rectangle++(0.04,0.04)
  (0.706,3.153)rectangle++(0.04,0.04) (0.724,4.248)rectangle++(0.04,0.04) (0.783,5.178)rectangle++(0.04,0.04) (0.87,5.966)rectangle++(0.04,0.04) (0.976,6.633)rectangle++(0.04,0.04)
  (-1.55,3.273)rectangle++(0.04,0.04) (-1.189,4.44)rectangle++(0.04,0.04) (1.404,1.186)rectangle++(0.04,0.04) (1.239,2.548)rectangle++(0.04,0.04) (1.153,3.713)rectangle++(0.04,0.04)
  (1.126,4.706)rectangle++(0.04,0.04) (1.143,5.551)rectangle++(0.04,0.04) (1.192,6.269)rectangle++(0.04,0.04) (1.261,6.877)rectangle++(0.04,0.04) (1.345,7.391)rectangle++(0.04,0.04)
  (1.437,7.825)rectangle++(0.04,0.04) (1.532,8.19)rectangle++(0.04,0.04) (1.628,8.496)rectangle++(0.04,0.04) (1.721,8.753)rectangle++(0.04,0.04) (1.811,8.967)rectangle++(0.04,0.04)
  (1.896,9.146)rectangle++(0.04,0.04) (2.259,3.118)rectangle++(0.04,0.04) (2.043,4.156)rectangle++(0.04,0.04) (1.9,5.047)rectangle++(0.04,0.04) (1.815,5.81)rectangle++(0.04,0.04)
  (1.773,6.462)rectangle++(0.04,0.04) (1.526,2.442)rectangle++(0.04,0.04) (1.393,3.611)rectangle++(0.04,0.04) (1.326,4.61)rectangle++(0.04,0.04) (1.309,5.461)rectangle++(0.04,0.04)
  (1.329,6.186)rectangle++(0.04,0.04) (1.375,6.801)rectangle++(0.04,0.04) (1.439,7.322)rectangle++(0.04,0.04) (-1.637,3.531)rectangle++(0.04,0.04) (-1.253,4.663)rectangle++(0.04,0.04)
  (-0.88,5.61)rectangle++(0.04,0.04) (-0.526,6.4)rectangle++(0.04,0.04) (-0.193,7.057)rectangle++(0.04,0.04) (0.116,7.602)rectangle++(0.04,0.04) (-0.02,1.2)rectangle++(0.04,0.04)
  (0.029,2.617)rectangle++(0.04,0.04) (0.127,3.819)rectangle++(0.04,0.04) (-0.989,2.458)rectangle++(0.04,0.04) (-0.744,3.725)rectangle++(0.04,0.04) (-0.486,4.793)rectangle++(0.04,0.04)
  (-0.223,5.689)rectangle++(0.04,0.04) (-1.545,2.789)rectangle++(0.04,0.04) (-1.204,4.029)rectangle++(0.04,0.04) (-0.864,5.069)rectangle++(0.04,0.04) (-0.534,5.939)rectangle++(0.04,0.04)
  (1.726,1.717)rectangle++(0.04,0.04) (-0.122,2.364)rectangle++(0.04,0.04) (-0.012,3.61)rectangle++(0.04,0.04) (0.132,4.665)rectangle++(0.04,0.04) (0.297,5.556)rectangle++(0.04,0.04)
  (0.472,6.307)rectangle++(0.04,0.04) (0.652,6.938)rectangle++(0.04,0.04) (1.828,2.265)rectangle++(0.04,0.04) (-0.244,2.508)rectangle++(0.04,0.04) (-0.11,3.737)rectangle++(0.04,0.04)
  (0.054,4.777)rectangle++(0.04,0.04) (0.235,5.655)rectangle++(0.04,0.04) (1.531,1.848)rectangle++(0.04,0.04) (1.373,3.106)rectangle++(0.04,0.04) (1.289,4.181)rectangle++(0.04,0.04)
  (1.261,5.099)rectangle++(0.04,0.04) (1.273,5.88)rectangle++(0.04,0.04) (-1.295,3.175)rectangle++(0.04,0.04) (-0.976,4.347)rectangle++(0.04,0.04) (-0.658,5.33)rectangle++(0.04,0.04)
  (-0.348,6.153)rectangle++(0.04,0.04) (-0.052,6.84)rectangle++(0.04,0.04) (0.227,7.413)rectangle++(0.04,0.04) (0.487,7.888)rectangle++(0.04,0.04) (0.728,8.281)rectangle++(0.04,0.04)
  (0.947,8.606)rectangle++(0.04,0.04) (1.147,8.874)rectangle++(0.04,0.04) (2.295,2.858)rectangle++(0.04,0.04) (2.063,3.934)rectangle++(0.04,0.04) (1.909,4.857)rectangle++(0.04,0.04)
  (1.814,5.649)rectangle++(0.04,0.04) (1.766,6.325)rectangle++(0.04,0.04) (1.752,6.902)rectangle++(0.04,0.04) (1.763,7.393)rectangle++(0.04,0.04) (1.792,7.809)rectangle++(0.04,0.04)
  (-1.693,3.719)rectangle++(0.04,0.04) (-1.327,2.018)rectangle++(0.04,0.04) (-0.811,1.728)rectangle++(0.04,0.04) (-0.623,3.097)rectangle++(0.04,0.04) (-0.951,2.127)rectangle++(0.04,0.04)
  (-0.725,3.442)rectangle++(0.04,0.04) (-1.061,2.179)rectangle++(0.04,0.04) (-0.817,3.491)rectangle++(0.04,0.04) (-0.557,4.596)rectangle++(0.04,0.04) (-0.292,5.525)rectangle++(0.04,0.04)
  (1.574,1.68)rectangle++(0.04,0.04) (0.217,1.242)rectangle++(0.04,0.04) (0.232,2.644)rectangle++(0.04,0.04) (-0.662,2.224)rectangle++(0.04,0.04) (-0.732,1.926)rectangle++(0.04,0.04)
  (-0.02,0.291)rectangle++(0.04,0.04) (-0.008,1.845)rectangle++(0.04,0.04) (0.065,3.164)rectangle++(0.04,0.04) (0.18,4.283)rectangle++(0.04,0.04) (0.322,5.23)rectangle++(0.04,0.04)
  (0.481,6.029)rectangle++(0.04,0.04) (1.599,2.002)rectangle++(0.04,0.04) (-0.222,2.397)rectangle++(0.04,0.04) (-0.095,3.643)rectangle++(0.04,0.04) (0.063,4.696)rectangle++(0.04,0.04)
  (0.239,5.585)rectangle++(0.04,0.04) (0.424,6.334)rectangle++(0.04,0.04) (-1.583,3.08)rectangle++(0.04,0.04) (-1.225,4.278)rectangle++(0.04,0.04) (-0.872,5.281)rectangle++(0.04,0.04)
  (-0.532,6.12)rectangle++(0.04,0.04) (-0.21,6.82)rectangle++(0.04,0.04) (0.092,7.401)rectangle++(0.04,0.04) (-1.927,3.238)rectangle++(0.04,0.04) (-1.511,4.426)rectangle++(0.04,0.04)
  (-1.109,5.419)rectangle++(0.04,0.04) (-0.728,6.246)rectangle++(0.04,0.04) (-0.371,6.935)rectangle++(0.04,0.04) (-0.04,7.506)rectangle++(0.04,0.04) (0.264,7.978)rectangle++(0.04,0.04)
  (0.541,8.367)rectangle++(0.04,0.04) (0.792,8.686)rectangle++(0.04,0.04) (1.019,8.948)rectangle++(0.04,0.04) (1.222,9.161)rectangle++(0.04,0.04) (1.403,9.334)rectangle++(0.04,0.04)
  (1.563,9.474)rectangle++(0.04,0.04) (1.706,9.587)rectangle++(0.04,0.04) (1.831,9.677)rectangle++(0.04,0.04) (1.941,9.748)rectangle++(0.04,0.04) (2.038,9.804)rectangle++(0.04,0.04)
  (2.122,9.848)rectangle++(0.04,0.04) (2.196,9.883)rectangle++(0.04,0.04) (2.259,9.909)rectangle++(0.04,0.04) (-2.146,4.289)rectangle++(0.04,0.04) (-1.654,5.327)rectangle++(0.04,0.04)
  (-1.737,2.38)rectangle++(0.04,0.04) (-1.384,3.689)rectangle++(0.04,0.04) (-1.031,4.787)rectangle++(0.04,0.04) (-1.472,2.405)rectangle++(0.04,0.04) (-1.157,3.699)rectangle++(0.04,0.04)
  (1.192,1.017)rectangle++(0.04,0.04) (-0.047,2.087)rectangle++(0.04,0.04) (0.041,3.372)rectangle++(0.04,0.04) (0.168,4.461)rectangle++(0.04,0.04) (0.319,5.381)rectangle++(0.04,0.04)
  (0.484,6.157)rectangle++(0.04,0.04) (0.655,6.811)rectangle++(0.04,0.04) (0.827,7.359)rectangle++(0.04,0.04) (1.919,2.411)rectangle++(0.04,0.04) (1.725,3.569)rectangle++(0.04,0.04)
  (1.607,4.561)rectangle++(0.04,0.04) (1.546,5.409)rectangle++(0.04,0.04) (1.529,6.132)rectangle++(0.04,0.04) (1.47,2.299)rectangle++(0.04,0.04) (1.339,3.492)rectangle++(0.04,0.04)
  (1.276,4.51)rectangle++(0.04,0.04) (1.263,5.379)rectangle++(0.04,0.04) (1.286,6.118)rectangle++(0.04,0.04) (1.336,6.745)rectangle++(0.04,0.04) (1.403,7.276)rectangle++(0.04,0.04)
  (1.481,7.725)rectangle++(0.04,0.04) (1.566,8.103)rectangle++(0.04,0.04) (-1.815,3.732)rectangle++(0.04,0.04) (-1.395,4.841)rectangle++(0.04,0.04) (-0.995,5.767)rectangle++(0.04,0.04)
  (-0.02,0.906)rectangle++(0.04,0.04) (-0.261,1.784)rectangle++(0.04,0.04) (-0.152,3.123)rectangle++(0.04,0.04) (-0.864,2.241)rectangle++(0.04,0.04) (-0.647,3.536)rectangle++(0.04,0.04)
  (-0.41,4.627)rectangle++(0.04,0.04) (-0.166,5.546)rectangle++(0.04,0.04) (0.079,6.317)rectangle++(0.04,0.04) (0.317,6.962)rectangle++(0.04,0.04) (0.546,7.501)rectangle++(0.04,0.04)
  (0.762,7.951)rectangle++(0.04,0.04) (0.963,8.324)rectangle++(0.04,0.04) (1.15,8.633)rectangle++(0.04,0.04) (1.32,8.888)rectangle++(0.04,0.04) (1.476,9.098)rectangle++(0.04,0.04)
  (1.616,9.271)rectangle++(0.04,0.04) (1.742,9.412)rectangle++(0.04,0.04) (1.855,9.526)rectangle++(0.04,0.04) (1.956,9.619)rectangle++(0.04,0.04) (2.045,9.695)rectangle++(0.04,0.04)
  (2.124,9.755)rectangle++(0.04,0.04) (2.193,9.803)rectangle++(0.04,0.04) (2.254,9.841)rectangle++(0.04,0.04) (2.307,9.871)rectangle++(0.04,0.04) (2.354,9.894)rectangle++(0.04,0.04)
  (2.394,9.912)rectangle++(0.04,0.04) (2.43,9.926)rectangle++(0.04,0.04) (2.46,9.936)rectangle++(0.04,0.04) (2.486,9.943)rectangle++(0.04,0.04) (-2.109,4.348)rectangle++(0.04,0.04)
  (-1.621,5.377)rectangle++(0.04,0.04) (-1.165,6.231)rectangle++(0.04,0.04) (-0.743,6.939)rectangle++(0.04,0.04) (-0.356,7.524)rectangle++(0.04,0.04) (-2.049,3.162)rectangle++(0.04,0.04)
  (-1.617,4.366)rectangle++(0.04,0.04) (-1.202,5.372)rectangle++(0.04,0.04) (-1.658,2.494)rectangle++(0.04,0.04) (-1.312,3.783)rectangle++(0.04,0.04) (-0.966,4.864)rectangle++(0.04,0.04)
  (1.489,1.346)rectangle++(0.04,0.04) (1.318,2.681)rectangle++(0.04,0.04) (1.225,3.822)rectangle++(0.04,0.04) (1.192,4.796)rectangle++(0.04,0.04) (1.203,5.625)rectangle++(0.04,0.04)
  (1.245,6.33)rectangle++(0.04,0.04) (1.309,6.926)rectangle++(0.04,0.04) (1.388,7.431)rectangle++(0.04,0.04) (1.475,7.857)rectangle++(0.04,0.04) (1.566,8.216)rectangle++(0.04,0.04)
  (2.048,2.809)rectangle++(0.04,0.04) (1.851,3.902)rectangle++(0.04,0.04) (1.727,4.839)rectangle++(0.04,0.04) (1.66,5.64)rectangle++(0.04,0.04) (1.313,2.215)rectangle++(0.04,0.04)
  (1.202,3.427)rectangle++(0.04,0.04) (1.157,4.461)rectangle++(0.04,0.04) (-0.95,2.836)rectangle++(0.04,0.04) (-0.696,4.045)rectangle++(0.04,0.04) (-0.432,5.062)rectangle++(0.04,0.04)
  (-0.167,5.917)rectangle++(0.04,0.04) (0.093,6.632)rectangle++(0.04,0.04) (0.342,7.23)rectangle++(0.04,0.04) (0.578,7.728)rectangle++(0.04,0.04) (0.798,8.142)rectangle++(0.04,0.04)
  (1.002,8.485)rectangle++(0.04,0.04) (1.189,8.768)rectangle++(0.04,0.04) (1.359,9.002)rectangle++(0.04,0.04) (1.513,9.193)rectangle++(0.04,0.04) (1.651,9.35)rectangle++(0.04,0.04)
  (1.776,9.478)rectangle++(0.04,0.04) (1.886,9.581)rectangle++(0.04,0.04) (1.984,9.665)rectangle++(0.04,0.04) (2.071,9.732)rectangle++(0.04,0.04) (2.147,9.785)rectangle++(0.04,0.04)
  (2.215,9.828)rectangle++(0.04,0.04) (2.273,9.861)rectangle++(0.04,0.04) (2.325,9.887)rectangle++(0.04,0.04) (2.369,9.908)rectangle++(0.04,0.04) (2.408,9.923)rectangle++(0.04,0.04)
  (2.441,9.934)rectangle++(0.04,0.04) (2.47,9.943)rectangle++(0.04,0.04) (2.495,9.949)rectangle++(0.04,0.04) (2.394,3.466)rectangle++(0.04,0.04) (2.171,4.447)rectangle++(0.04,0.04)
  (2.021,5.289)rectangle++(0.04,0.04) (1.927,6.011)rectangle++(0.04,0.04) (1.877,6.629)rectangle++(0.04,0.04) (1.858,7.155)rectangle++(0.04,0.04) (1.863,7.604)rectangle++(0.04,0.04)
  (1.832,2.739)rectangle++(0.04,0.04) (1.665,3.851)rectangle++(0.04,0.04) (0.811,1.787)rectangle++(0.04,0.04) (0.759,3.083)rectangle++(0.04,0.04) (0.766,4.186)rectangle++(0.04,0.04)
  (0.816,5.124)rectangle++(0.04,0.04) (0.897,5.919)rectangle++(0.04,0.04) (0.997,6.591)rectangle++(0.04,0.04) (1.679,2.271)rectangle++(0.04,0.04) (1.516,3.459)rectangle++(0.04,0.04)
  (1.424,4.476)rectangle++(0.04,0.04) (1.388,5.344)rectangle++(0.04,0.04) (1.391,6.083)rectangle++(0.04,0.04) (1.423,6.711)rectangle++(0.04,0.04) (1.476,7.244)rectangle++(0.04,0.04)
  (1.542,7.694)rectangle++(0.04,0.04) (-1.713,3.636)rectangle++(0.04,0.04) (-1.313,4.756)rectangle++(0.04,0.04) (-0.928,5.691)rectangle++(0.04,0.04) (-0.563,6.471)rectangle++(0.04,0.04)
  (-0.222,7.119)rectangle++(0.04,0.04) (0.094,7.656)rectangle++(0.04,0.04) (0.384,8.1)rectangle++(0.04,0.04) (0.648,8.466)rectangle++(0.04,0.04) (0.887,8.766)rectangle++(0.04,0.04);

      \begin{scope}[shift={(7.4,0)}]
        % Fougère : les quatre pièces w_i(A)
% Fichier généré par generer.py : ne pas modifier à la main.
\fill[piece1]
  (-0.02,1.067)rectangle++(0.04,0.04) (-0.02,1.412)rectangle++(0.04,0.04) (-0.02,0.385)rectangle++(0.04,0.04) (-0.02,1.074)rectangle++(0.04,0.04) (-0.02,0.4)rectangle++(0.04,0.04)
  (-0.02,1.215)rectangle++(0.04,0.04) (-0.02,0.921)rectangle++(0.04,0.04) (-0.02,1.033)rectangle++(0.04,0.04) (-0.02,0.816)rectangle++(0.04,0.04) (-0.02,1.31)rectangle++(0.04,0.04)
  (-0.02,1.341)rectangle++(0.04,0.04) (-0.02,0.91)rectangle++(0.04,0.04) (-0.02,0.543)rectangle++(0.04,0.04) (-0.02,0.434)rectangle++(0.04,0.04) (-0.02,1.446)rectangle++(0.04,0.04)
  (-0.02,0.871)rectangle++(0.04,0.04) (-0.02,1.163)rectangle++(0.04,0.04) (-0.02,0.753)rectangle++(0.04,0.04) (-0.02,0.319)rectangle++(0.04,0.04) (-0.02,0.947)rectangle++(0.04,0.04)
  (-0.02,1.366)rectangle++(0.04,0.04) (-0.02,0.77)rectangle++(0.04,0.04) (-0.02,1.211)rectangle++(0.04,0.04) (-0.02,0.695)rectangle++(0.04,0.04) (-0.02,0.607)rectangle++(0.04,0.04)
  (-0.02,0.715)rectangle++(0.04,0.04) (-0.02,0.197)rectangle++(0.04,0.04) (-0.02,0.015)rectangle++(0.04,0.04) (-0.02,1.46)rectangle++(0.04,0.04) (-0.02,0.851)rectangle++(0.04,0.04)
  (-0.02,0.597)rectangle++(0.04,0.04) (-0.02,0.504)rectangle++(0.04,0.04) (-0.02,1.3)rectangle++(0.04,0.04) (-0.02,1.115)rectangle++(0.04,0.04) (-0.02,0.451)rectangle++(0.04,0.04)
  (-0.02,0.689)rectangle++(0.04,0.04) (-0.02,1.505)rectangle++(0.04,0.04) (-0.02,1.556)rectangle++(0.04,0.04) (-0.02,0.79)rectangle++(0.04,0.04) (-0.02,0.336)rectangle++(0.04,0.04)
  (-0.02,0.456)rectangle++(0.04,0.04) (-0.02,0.35)rectangle++(0.04,0.04) (-0.02,0.575)rectangle++(0.04,0.04) (-0.02,0.56)rectangle++(0.04,0.04) (-0.02,0.696)rectangle++(0.04,0.04)
  (-0.02,0.674)rectangle++(0.04,0.04) (-0.02,0.993)rectangle++(0.04,0.04) (-0.02,1.061)rectangle++(0.04,0.04) (-0.02,1.263)rectangle++(0.04,0.04) (-0.02,1.488)rectangle++(0.04,0.04)
  (-0.02,0.424)rectangle++(0.04,0.04) (-0.02,1.551)rectangle++(0.04,0.04) (-0.02,1.364)rectangle++(0.04,0.04) (-0.02,1.319)rectangle++(0.04,0.04) (-0.02,1.094)rectangle++(0.04,0.04)
  (-0.02,0.645)rectangle++(0.04,0.04) (-0.02,0.741)rectangle++(0.04,0.04) (-0.02,0.909)rectangle++(0.04,0.04) (-0.02,0.327)rectangle++(0.04,0.04) (-0.02,0.469)rectangle++(0.04,0.04)
  (-0.02,0.616)rectangle++(0.04,0.04) (-0.02,1.354)rectangle++(0.04,0.04) (-0.02,1.217)rectangle++(0.04,0.04) (-0.02,1.418)rectangle++(0.04,0.04) (-0.02,0.959)rectangle++(0.04,0.04)
  (-0.02,0.704)rectangle++(0.04,0.04) (-0.02,0.445)rectangle++(0.04,0.04) (-0.02,0.597)rectangle++(0.04,0.04) (-0.02,1.515)rectangle++(0.04,0.04) (-0.02,0.545)rectangle++(0.04,0.04)
  (-0.02,1.529)rectangle++(0.04,0.04) (-0.02,0.701)rectangle++(0.04,0.04) (-0.02,0.461)rectangle++(0.04,0.04) (-0.02,0.914)rectangle++(0.04,0.04) (-0.02,0.484)rectangle++(0.04,0.04)
  (-0.02,1.198)rectangle++(0.04,0.04) (-0.02,0.571)rectangle++(0.04,0.04) (-0.02,1.034)rectangle++(0.04,0.04) (-0.02,1.575)rectangle++(0.04,0.04) (-0.02,1.413)rectangle++(0.04,0.04)
  (-0.02,0.445)rectangle++(0.04,0.04) (-0.02,0.054)rectangle++(0.04,0.04) (-0.02,1.487)rectangle++(0.04,0.04) (-0.02,0.569)rectangle++(0.04,0.04) (-0.02,1.151)rectangle++(0.04,0.04)
  (-0.02,0.332)rectangle++(0.04,0.04) (-0.02,0.673)rectangle++(0.04,0.04) (-0.02,1.506)rectangle++(0.04,0.04) (-0.02,1.247)rectangle++(0.04,0.04) (-0.02,0.761)rectangle++(0.04,0.04)
  (-0.02,0.906)rectangle++(0.04,0.04) (-0.02,1.242)rectangle++(0.04,0.04) (-0.02,1.047)rectangle++(0.04,0.04) (-0.02,0.588)rectangle++(0.04,0.04) (-0.02,1.294)rectangle++(0.04,0.04)
  (-0.02,0.919)rectangle++(0.04,0.04) (-0.02,0.888)rectangle++(0.04,0.04) (-0.02,1.334)rectangle++(0.04,0.04) (-0.02,0.617)rectangle++(0.04,0.04) (-0.02,0.964)rectangle++(0.04,0.04)
  (-0.02,0.931)rectangle++(0.04,0.04) (-0.02,0.118)rectangle++(0.04,0.04) (-0.02,1.392)rectangle++(0.04,0.04) (-0.02,1.202)rectangle++(0.04,0.04) (-0.02,1.007)rectangle++(0.04,0.04)
  (-0.02,1.019)rectangle++(0.04,0.04) (-0.02,0.566)rectangle++(0.04,0.04) (-0.02,1.034)rectangle++(0.04,0.04) (-0.02,0.881)rectangle++(0.04,0.04) (-0.02,0.882)rectangle++(0.04,0.04)
  (-0.02,1.485)rectangle++(0.04,0.04) (-0.02,0.282)rectangle++(0.04,0.04) (-0.02,0.447)rectangle++(0.04,0.04) (-0.02,1.362)rectangle++(0.04,0.04) (-0.02,0.405)rectangle++(0.04,0.04)
  (-0.02,1.449)rectangle++(0.04,0.04) (-0.02,0.541)rectangle++(0.04,0.04) (-0.02,1.466)rectangle++(0.04,0.04) (-0.02,1.169)rectangle++(0.04,0.04) (-0.02,0.628)rectangle++(0.04,0.04)
  (-0.02,1.385)rectangle++(0.04,0.04) (-0.02,1.47)rectangle++(0.04,0.04) (-0.02,0.849)rectangle++(0.04,0.04) (-0.02,0.501)rectangle++(0.04,0.04) (-0.02,1.263)rectangle++(0.04,0.04)
  (-0.02,1.073)rectangle++(0.04,0.04) (-0.02,1.453)rectangle++(0.04,0.04) (-0.02,1.529)rectangle++(0.04,0.04) (-0.02,1.368)rectangle++(0.04,0.04) (-0.02,1.113)rectangle++(0.04,0.04)
  (-0.02,0.617)rectangle++(0.04,0.04) (-0.02,1.184)rectangle++(0.04,0.04) (-0.02,0.413)rectangle++(0.04,0.04) (-0.02,1.219)rectangle++(0.04,0.04) (-0.02,0.853)rectangle++(0.04,0.04)
  (-0.02,1.264)rectangle++(0.04,0.04) (-0.02,0.638)rectangle++(0.04,0.04) (-0.02,1.04)rectangle++(0.04,0.04) (-0.02,1.125)rectangle++(0.04,0.04) (-0.02,1.575)rectangle++(0.04,0.04)
  (-0.02,0.887)rectangle++(0.04,0.04) (-0.02,0.603)rectangle++(0.04,0.04) (-0.02,0.357)rectangle++(0.04,0.04) (-0.02,0.586)rectangle++(0.04,0.04) (-0.02,0.697)rectangle++(0.04,0.04)
  (-0.02,0.997)rectangle++(0.04,0.04) (-0.02,1.575)rectangle++(0.04,0.04) (-0.02,1.47)rectangle++(0.04,0.04) (-0.02,0.806)rectangle++(0.04,0.04) (-0.02,0.535)rectangle++(0.04,0.04)
  (-0.02,1.175)rectangle++(0.04,0.04) (-0.02,0.755)rectangle++(0.04,0.04) (-0.02,1.477)rectangle++(0.04,0.04) (-0.02,0.996)rectangle++(0.04,0.04) (-0.02,1.027)rectangle++(0.04,0.04)
  (-0.02,1.434)rectangle++(0.04,0.04) (-0.02,1.265)rectangle++(0.04,0.04) (-0.02,1.473)rectangle++(0.04,0.04) (-0.02,1.547)rectangle++(0.04,0.04) (-0.02,1.062)rectangle++(0.04,0.04)
  (-0.02,1.077)rectangle++(0.04,0.04) (-0.02,1.407)rectangle++(0.04,0.04) (-0.02,1.179)rectangle++(0.04,0.04) (-0.02,1.233)rectangle++(0.04,0.04) (-0.02,0.916)rectangle++(0.04,0.04)
  (-0.02,0.387)rectangle++(0.04,0.04) (-0.02,0.492)rectangle++(0.04,0.04) (-0.02,1.565)rectangle++(0.04,0.04) (-0.02,1.28)rectangle++(0.04,0.04) (-0.02,1.381)rectangle++(0.04,0.04)
  (-0.02,1.505)rectangle++(0.04,0.04) (-0.02,1.429)rectangle++(0.04,0.04) (-0.02,1.255)rectangle++(0.04,0.04) (-0.02,0.622)rectangle++(0.04,0.04) (-0.02,1.525)rectangle++(0.04,0.04)
  (-0.02,1.504)rectangle++(0.04,0.04) (-0.02,1.41)rectangle++(0.04,0.04) (-0.02,0.177)rectangle++(0.04,0.04) (-0.02,0.89)rectangle++(0.04,0.04) (-0.02,0.672)rectangle++(0.04,0.04)
  (-0.02,1.484)rectangle++(0.04,0.04) (-0.02,0.279)rectangle++(0.04,0.04) (-0.02,1.562)rectangle++(0.04,0.04) (-0.02,1.168)rectangle++(0.04,0.04) (-0.02,0.892)rectangle++(0.04,0.04)
  (-0.02,0.498)rectangle++(0.04,0.04) (-0.02,0.838)rectangle++(0.04,0.04) (-0.02,1.206)rectangle++(0.04,0.04) (-0.02,1.011)rectangle++(0.04,0.04) (-0.02,1.572)rectangle++(0.04,0.04)
  (-0.02,0.618)rectangle++(0.04,0.04) (-0.02,1.2)rectangle++(0.04,0.04) (-0.02,0.291)rectangle++(0.04,0.04) (-0.02,0.906)rectangle++(0.04,0.04);
\fill[piece2]
  (1.402,9.252)rectangle++(0.04,0.04) (1.56,9.404)rectangle++(0.04,0.04) (1.7,9.527)rectangle++(0.04,0.04) (1.824,9.626)rectangle++(0.04,0.04) (1.933,9.706)rectangle++(0.04,0.04)
  (2.029,9.769)rectangle++(0.04,0.04) (2.113,9.818)rectangle++(0.04,0.04) (2.187,9.857)rectangle++(0.04,0.04) (2.251,9.887)rectangle++(0.04,0.04) (2.307,9.911)rectangle++(0.04,0.04)
  (2.184,4.397)rectangle++(0.04,0.04) (2.03,5.246)rectangle++(0.04,0.04) (1.061,3.435)rectangle++(0.04,0.04) (1.037,4.473)rectangle++(0.04,0.04) (1.059,5.357)rectangle++(0.04,0.04)
  (1.112,6.107)rectangle++(0.04,0.04) (1.187,6.743)rectangle++(0.04,0.04) (1.277,7.28)rectangle++(0.04,0.04) (1.374,7.733)rectangle++(0.04,0.04) (1.475,8.115)rectangle++(0.04,0.04)
  (1.576,8.435)rectangle++(0.04,0.04) (1.675,8.703)rectangle++(0.04,0.04) (1.77,8.926)rectangle++(0.04,0.04) (1.859,9.113)rectangle++(0.04,0.04) (1.942,9.268)rectangle++(0.04,0.04)
  (2.02,9.396)rectangle++(0.04,0.04) (-1.589,5.181)rectangle++(0.04,0.04) (-1.145,6.064)rectangle++(0.04,0.04) (-0.733,6.796)rectangle++(0.04,0.04) (-0.354,7.402)rectangle++(0.04,0.04)
  (-1.591,4.343)rectangle++(0.04,0.04) (-1.181,5.351)rectangle++(0.04,0.04) (-0.792,6.192)rectangle++(0.04,0.04) (-0.427,6.891)rectangle++(0.04,0.04) (-0.09,7.471)rectangle++(0.04,0.04)
  (0.22,7.95)rectangle++(0.04,0.04) (0.503,8.345)rectangle++(0.04,0.04) (2.007,3.685)rectangle++(0.04,0.04) (1.852,4.649)rectangle++(0.04,0.04) (1.758,5.473)rectangle++(0.04,0.04)
  (1.711,6.178)rectangle++(0.04,0.04) (1.699,6.779)rectangle++(0.04,0.04) (-1.091,4.604)rectangle++(0.04,0.04) (1.268,2.603)rectangle++(0.04,0.04) (-0.282,3.699)rectangle++(0.04,0.04)
  (-0.094,4.752)rectangle++(0.04,0.04) (0.108,5.639)rectangle++(0.04,0.04) (0.315,6.385)rectangle++(0.04,0.04) (0.521,7.011)rectangle++(0.04,0.04) (0.721,7.535)rectangle++(0.04,0.04)
  (0.912,7.972)rectangle++(0.04,0.04) (1.092,8.336)rectangle++(0.04,0.04) (1.26,8.638)rectangle++(0.04,0.04) (1.414,8.888)rectangle++(0.04,0.04) (1.555,9.094)rectangle++(0.04,0.04)
  (2.07,4.069)rectangle++(0.04,0.04) (1.92,4.972)rectangle++(0.04,0.04) (1.006,3.357)rectangle++(0.04,0.04) (0.988,4.409)rectangle++(0.04,0.04) (1.014,5.305)rectangle++(0.04,0.04)
  (1.072,6.065)rectangle++(0.04,0.04) (1.151,6.708)rectangle++(0.04,0.04) (1.245,7.252)rectangle++(0.04,0.04) (1.346,7.711)rectangle++(0.04,0.04) (1.45,8.096)rectangle++(0.04,0.04)
  (1.554,8.42)rectangle++(0.04,0.04) (1.656,8.691)rectangle++(0.04,0.04) (1.753,8.917)rectangle++(0.04,0.04) (1.844,9.106)rectangle++(0.04,0.04) (1.93,9.262)rectangle++(0.04,0.04)
  (2.009,9.392)rectangle++(0.04,0.04) (2.081,9.499)rectangle++(0.04,0.04) (2.146,9.587)rectangle++(0.04,0.04) (1.282,2.297)rectangle++(0.04,0.04) (1.18,3.498)rectangle++(0.04,0.04)
  (1.14,4.522)rectangle++(0.04,0.04) (1.148,5.394)rectangle++(0.04,0.04) (1.189,6.135)rectangle++(0.04,0.04) (1.254,6.764)rectangle++(0.04,0.04) (1.334,7.295)rectangle++(0.04,0.04)
  (1.424,7.744)rectangle++(0.04,0.04) (1.518,8.121)rectangle++(0.04,0.04) (-1.408,4.835)rectangle++(0.04,0.04) (-1.006,5.763)rectangle++(0.04,0.04) (-0.626,6.535)rectangle++(0.04,0.04)
  (-0.273,7.176)rectangle++(0.04,0.04) (0.053,7.707)rectangle++(0.04,0.04) (-1.583,4.479)rectangle++(0.04,0.04) (-1.168,5.466)rectangle++(0.04,0.04) (-0.777,6.289)rectangle++(0.04,0.04)
  (-0.411,6.973)rectangle++(0.04,0.04) (-0.072,7.54)rectangle++(0.04,0.04) (0.238,8.008)rectangle++(0.04,0.04) (0.52,8.394)rectangle++(0.04,0.04) (0.776,8.71)rectangle++(0.04,0.04)
  (1.006,8.969)rectangle++(0.04,0.04) (1.211,9.179)rectangle++(0.04,0.04) (1.395,9.35)rectangle++(0.04,0.04) (1.557,9.488)rectangle++(0.04,0.04) (1.701,9.599)rectangle++(0.04,0.04)
  (1.828,9.687)rectangle++(0.04,0.04) (1.939,9.757)rectangle++(0.04,0.04) (2.036,9.812)rectangle++(0.04,0.04) (2.121,9.855)rectangle++(0.04,0.04) (2.195,9.888)rectangle++(0.04,0.04)
  (2.259,9.913)rectangle++(0.04,0.04) (2.314,9.932)rectangle++(0.04,0.04) (2.362,9.946)rectangle++(0.04,0.04) (2.404,9.956)rectangle++(0.04,0.04) (2.439,9.963)rectangle++(0.04,0.04)
  (2.181,4.437)rectangle++(0.04,0.04) (-0.507,4.233)rectangle++(0.04,0.04) (-0.264,5.214)rectangle++(0.04,0.04) (-1.11,3.927)rectangle++(0.04,0.04) (-0.789,4.979)rectangle++(0.04,0.04)
  (-1.155,3.783)rectangle++(0.04,0.04) (-0.832,4.858)rectangle++(0.04,0.04) (-0.515,5.758)rectangle++(0.04,0.04) (-1.271,3.988)rectangle++(0.04,0.04) (-0.923,5.037)rectangle++(0.04,0.04)
  (-0.585,5.914)rectangle++(0.04,0.04) (-0.263,6.647)rectangle++(0.04,0.04) (0.04,7.256)rectangle++(0.04,0.04) (0.322,7.763)rectangle++(0.04,0.04) (-1.026,3.243)rectangle++(0.04,0.04)
  (-0.744,4.394)rectangle++(0.04,0.04) (-0.459,5.361)rectangle++(0.04,0.04) (1.381,2.892)rectangle++(0.04,0.04) (1.288,3.999)rectangle++(0.04,0.04) (1.252,4.944)rectangle++(0.04,0.04)
  (-0.781,4.158)rectangle++(0.04,0.04) (-0.5,5.162)rectangle++(0.04,0.04) (-0.221,6.004)rectangle++(0.04,0.04) (0.05,6.708)rectangle++(0.04,0.04) (0.309,7.296)rectangle++(0.04,0.04)
  (0.552,7.786)rectangle++(0.04,0.04) (0.779,8.192)rectangle++(0.04,0.04) (0.987,8.528)rectangle++(0.04,0.04) (-1.589,4.816)rectangle++(0.04,0.04) (-1.161,5.754)rectangle++(0.04,0.04)
  (-0.759,6.533)rectangle++(0.04,0.04) (-0.21,3.254)rectangle++(0.04,0.04) (-0.051,4.37)rectangle++(0.04,0.04) (0.129,5.313)rectangle++(0.04,0.04) (0.32,6.107)rectangle++(0.04,0.04)
  (0.514,6.774)rectangle++(0.04,0.04) (0.023,2.504)rectangle++(0.04,0.04) (0.118,3.724)rectangle++(0.04,0.04) (0.247,4.757)rectangle++(0.04,0.04) (0.398,5.629)rectangle++(0.04,0.04)
  (0.561,6.365)rectangle++(0.04,0.04) (0.729,6.984)rectangle++(0.04,0.04) (0.897,7.504)rectangle++(0.04,0.04) (1.061,7.938)rectangle++(0.04,0.04) (-0.144,3.847)rectangle++(0.04,0.04)
  (0.029,4.872)rectangle++(0.04,0.04) (0.218,5.736)rectangle++(0.04,0.04) (0.412,6.463)rectangle++(0.04,0.04) (0.607,7.073)rectangle++(0.04,0.04) (0.796,7.584)rectangle++(0.04,0.04)
  (0.978,8.011)rectangle++(0.04,0.04) (1.15,8.366)rectangle++(0.04,0.04) (1.31,8.662)rectangle++(0.04,0.04) (1.457,8.906)rectangle++(0.04,0.04) (1.593,9.108)rectangle++(0.04,0.04)
  (1.716,9.274)rectangle++(0.04,0.04) (1.827,9.411)rectangle++(0.04,0.04) (1.928,9.522)rectangle++(0.04,0.04) (2.017,9.613)rectangle++(0.04,0.04) (2.097,9.687)rectangle++(0.04,0.04)
  (2.168,9.746)rectangle++(0.04,0.04) (2.23,9.794)rectangle++(0.04,0.04) (-1.635,5.299)rectangle++(0.04,0.04) (-1.18,6.166)rectangle++(0.04,0.04) (-0.759,6.884)rectangle++(0.04,0.04)
  (-0.372,7.478)rectangle++(0.04,0.04) (-0.019,7.968)rectangle++(0.04,0.04) (-1.651,4.517)rectangle++(0.04,0.04) (-1.225,5.502)rectangle++(0.04,0.04) (-0.823,6.322)rectangle++(0.04,0.04)
  (-0.449,7.002)rectangle++(0.04,0.04) (-0.104,7.566)rectangle++(0.04,0.04) (0.212,8.032)rectangle++(0.04,0.04) (-1.623,4.573)rectangle++(0.04,0.04) (-1.199,5.548)rectangle++(0.04,0.04)
  (-0.799,6.36)rectangle++(0.04,0.04) (1.669,3.007)rectangle++(0.04,0.04) (1.536,4.086)rectangle++(0.04,0.04) (1.467,5.007)rectangle++(0.04,0.04) (1.445,5.794)rectangle++(0.04,0.04)
  (1.458,6.463)rectangle++(0.04,0.04) (-1.067,4.497)rectangle++(0.04,0.04) (-0.729,5.461)rectangle++(0.04,0.04) (-0.404,6.267)rectangle++(0.04,0.04) (-0.095,6.94)rectangle++(0.04,0.04)
  (1.73,3.279)rectangle++(0.04,0.04) (1.599,4.314)rectangle++(0.04,0.04) (0.881,3.157)rectangle++(0.04,0.04) (0.873,4.245)rectangle++(0.04,0.04) (0.91,5.169)rectangle++(0.04,0.04)
  (0.978,5.954)rectangle++(0.04,0.04) (1.067,6.618)rectangle++(0.04,0.04) (1.17,7.178)rectangle++(0.04,0.04) (1.279,7.651)rectangle++(0.04,0.04) (1.391,8.049)rectangle++(0.04,0.04)
  (1.502,8.382)rectangle++(0.04,0.04) (1.61,8.661)rectangle++(0.04,0.04) (-0.19,3.932)rectangle++(0.04,0.04) (-0.006,4.946)rectangle++(0.04,0.04) (0.19,5.8)rectangle++(0.04,0.04)
  (0.392,6.519)rectangle++(0.04,0.04) (0.591,7.122)rectangle++(0.04,0.04) (0.785,7.626)rectangle++(0.04,0.04) (0.97,8.047)rectangle++(0.04,0.04) (1.145,8.397)rectangle++(0.04,0.04)
  (1.307,8.688)rectangle++(0.04,0.04) (1.456,8.929)rectangle++(0.04,0.04) (0.037,2.797)rectangle++(0.04,0.04) (0.688,2.51)rectangle++(0.04,0.04) (-0.042,3.02)rectangle++(0.04,0.04)
  (0.083,4.165)rectangle++(0.04,0.04) (0.235,5.133)rectangle++(0.04,0.04) (0.265,2.609)rectangle++(0.04,0.04) (-0.463,3.512)rectangle++(0.04,0.04) (-0.255,4.6)rectangle++(0.04,0.04)
  (-0.035,5.517)rectangle++(0.04,0.04) (0.189,6.287)rectangle++(0.04,0.04) (1.535,3.217)rectangle++(0.04,0.04) (1.431,4.27)rectangle++(0.04,0.04) (1.385,5.168)rectangle++(0.04,0.04)
  (1.382,5.934)rectangle++(0.04,0.04) (1.41,6.585)rectangle++(0.04,0.04) (1.459,7.137)rectangle++(0.04,0.04) (-0.176,3.776)rectangle++(0.04,0.04) (-0.001,4.813)rectangle++(0.04,0.04)
  (0.19,5.687)rectangle++(0.04,0.04) (0.386,6.423)rectangle++(0.04,0.04) (0.583,7.04)rectangle++(0.04,0.04) (0.775,7.557)rectangle++(0.04,0.04) (0.959,7.989)rectangle++(0.04,0.04)
  (1.132,8.348)rectangle++(0.04,0.04) (1.294,8.647)rectangle++(0.04,0.04) (0.445,3.012)rectangle++(0.04,0.04) (-0.516,3.626)rectangle++(0.04,0.04) (-0.296,4.699)rectangle++(0.04,0.04)
  (-0.066,5.602)rectangle++(0.04,0.04) (0.166,6.36)rectangle++(0.04,0.04) (1.556,3.226)rectangle++(0.04,0.04) (1.449,4.276)rectangle++(0.04,0.04) (1.401,5.173)rectangle++(0.04,0.04)
  (1.395,5.937)rectangle++(0.04,0.04) (1.307,3.424)rectangle++(0.04,0.04) (0.695,2.919)rectangle++(0.04,0.04) (0.705,4.05)rectangle++(0.04,0.04) (0.759,5.01)rectangle++(0.04,0.04)
  (0.844,5.825)rectangle++(0.04,0.04) (0.948,6.513)rectangle++(0.04,0.04) (1.064,7.095)rectangle++(0.04,0.04) (1.186,7.584)rectangle++(0.04,0.04) (1.309,7.995)rectangle++(0.04,0.04)
  (1.43,8.34)rectangle++(0.04,0.04) (1.547,8.628)rectangle++(0.04,0.04) (1.955,3.977)rectangle++(0.04,0.04) (1.819,4.898)rectangle++(0.04,0.04) (1.74,5.687)rectangle++(0.04,0.04)
  (1.704,6.361)rectangle++(0.04,0.04) (1.701,6.935)rectangle++(0.04,0.04) (1.721,7.423)rectangle++(0.04,0.04) (-1.223,4.736)rectangle++(0.04,0.04) (-0.853,5.67)rectangle++(0.04,0.04)
  (-0.5,6.45)rectangle++(0.04,0.04) (-0.169,7.099)rectangle++(0.04,0.04) (0.138,7.637)rectangle++(0.04,0.04) (0.42,8.082)rectangle++(0.04,0.04) (-0.092,3.846)rectangle++(0.04,0.04)
  (0.074,4.869)rectangle++(0.04,0.04) (-0.977,3.923)rectangle++(0.04,0.04) (-0.675,4.97)rectangle++(0.04,0.04) (1.31,2.767)rectangle++(0.04,0.04) (1.222,3.896)rectangle++(0.04,0.04)
  (1.192,4.859)rectangle++(0.04,0.04) (-0.774,4.13)rectangle++(0.04,0.04) (-0.495,5.138)rectangle++(0.04,0.04) (-0.217,5.983)rectangle++(0.04,0.04) (0.052,6.691)rectangle++(0.04,0.04)
  (0.31,7.281)rectangle++(0.04,0.04) (0.552,7.773)rectangle++(0.04,0.04) (-1.507,4.586)rectangle++(0.04,0.04) (-1.1,5.554)rectangle++(0.04,0.04) (-0.715,6.361)rectangle++(0.04,0.04)
  (-1.435,4.069)rectangle++(0.04,0.04) (-1.059,5.113)rectangle++(0.04,0.04) (-0.698,5.984)rectangle++(0.04,0.04) (-0.356,6.711)rectangle++(0.04,0.04) (-0.036,7.315)rectangle++(0.04,0.04)
  (-1.515,4.385)rectangle++(0.04,0.04) (-1.115,5.384)rectangle++(0.04,0.04) (-0.734,6.217)rectangle++(0.04,0.04) (-0.378,6.91)rectangle++(0.04,0.04) (-0.047,7.485)rectangle++(0.04,0.04)
  (0.257,7.96)rectangle++(0.04,0.04) (0.535,8.352)rectangle++(0.04,0.04) (0.787,8.674)rectangle++(0.04,0.04) (1.013,8.938)rectangle++(0.04,0.04) (1.217,9.153)rectangle++(0.04,0.04)
  (1.398,9.327)rectangle++(0.04,0.04) (1.559,9.469)rectangle++(0.04,0.04) (1.702,9.582)rectangle++(0.04,0.04) (1.828,9.673)rectangle++(0.04,0.04) (1.938,9.745)rectangle++(0.04,0.04)
  (2.035,9.802)rectangle++(0.04,0.04) (2.12,9.846)rectangle++(0.04,0.04) (2.193,9.881)rectangle++(0.04,0.04) (2.257,9.907)rectangle++(0.04,0.04) (2.313,9.927)rectangle++(0.04,0.04)
  (2.187,4.401)rectangle++(0.04,0.04) (2.033,5.25)rectangle++(0.04,0.04) (1.936,5.977)rectangle++(0.04,0.04) (1.882,6.599)rectangle++(0.04,0.04) (1.862,7.13)rectangle++(0.04,0.04)
  (1.866,7.583)rectangle++(0.04,0.04) (1.887,7.967)rectangle++(0.04,0.04) (1.92,8.292)rectangle++(0.04,0.04) (1.962,8.568)rectangle++(0.04,0.04) (2.008,8.801)rectangle++(0.04,0.04)
  (2.057,8.996)rectangle++(0.04,0.04) (2.106,9.161)rectangle++(0.04,0.04) (2.022,4.207)rectangle++(0.04,0.04) (-0.487,4.157)rectangle++(0.04,0.04) (-0.25,5.149)rectangle++(0.04,0.04)
  (-0.009,5.983)rectangle++(0.04,0.04) (0.23,6.682)rectangle++(0.04,0.04) (0.46,7.267)rectangle++(0.04,0.04) (0.68,7.755)rectangle++(0.04,0.04) (0.886,8.16)rectangle++(0.04,0.04)
  (1.077,8.497)rectangle++(0.04,0.04) (1.253,8.776)rectangle++(0.04,0.04) (1.414,9.005)rectangle++(0.04,0.04) (1.56,9.194)rectangle++(0.04,0.04) (1.692,9.349)rectangle++(0.04,0.04)
  (1.81,9.475)rectangle++(0.04,0.04) (-1.643,5.157)rectangle++(0.04,0.04) (-1.192,6.046)rectangle++(0.04,0.04) (-0.774,6.783)rectangle++(0.04,0.04) (-0.389,7.392)rectangle++(0.04,0.04)
  (-1.595,4.334)rectangle++(0.04,0.04) (-1.185,5.344)rectangle++(0.04,0.04) (-0.795,6.186)rectangle++(0.04,0.04) (-0.431,6.886)rectangle++(0.04,0.04) (-0.093,7.467)rectangle++(0.04,0.04)
  (0.217,7.947)rectangle++(0.04,0.04) (0.501,8.342)rectangle++(0.04,0.04) (2.007,3.684)rectangle++(0.04,0.04) (0.677,3.128)rectangle++(0.04,0.04) (-0.499,3.692)rectangle++(0.04,0.04)
  (-0.279,4.755)rectangle++(0.04,0.04) (-0.049,5.649)rectangle++(0.04,0.04) (0.182,6.4)rectangle++(0.04,0.04) (0.409,7.029)rectangle++(0.04,0.04) (0.626,7.554)rectangle++(0.04,0.04)
  (0.832,7.992)rectangle++(0.04,0.04) (1.025,8.356)rectangle++(0.04,0.04) (1.203,8.658)rectangle++(0.04,0.04) (1.367,8.907)rectangle++(0.04,0.04) (1.516,9.113)rectangle++(0.04,0.04)
  (1.651,9.281)rectangle++(0.04,0.04) (1.772,9.419)rectangle++(0.04,0.04) (1.881,9.532)rectangle++(0.04,0.04) (2.14,4.227)rectangle++(0.04,0.04) (1.986,5.104)rectangle++(0.04,0.04)
  (1.89,5.855)rectangle++(0.04,0.04) (-0.86,4.458)rectangle++(0.04,0.04) (0.746,2.113)rectangle++(0.04,0.04) (0.716,3.362)rectangle++(0.04,0.04) (0.009,3.261)rectangle++(0.04,0.04)
  (0.136,4.368)rectangle++(0.04,0.04) (0.288,5.303)rectangle++(0.04,0.04) (0.455,6.092)rectangle++(0.04,0.04) (1.455,3.24)rectangle++(0.04,0.04) (1.365,4.292)rectangle++(0.04,0.04)
  (1.329,5.19)rectangle++(0.04,0.04) (1.335,5.954)rectangle++(0.04,0.04) (1.371,6.604)rectangle++(0.04,0.04) (1.427,7.155)rectangle++(0.04,0.04) (1.605,3.666)rectangle++(0.04,0.04)
  (0.72,3.033)rectangle++(0.04,0.04) (0.678,2.278)rectangle++(0.04,0.04) (0.665,3.506)rectangle++(0.04,0.04) (0.703,4.549)rectangle++(0.04,0.04) (0.778,5.435)rectangle++(0.04,0.04)
  (0.876,6.185)rectangle++(0.04,0.04) (0.99,6.818)rectangle++(0.04,0.04) (0.024,2.51)rectangle++(0.04,0.04) (0.118,3.729)rectangle++(0.04,0.04) (0.248,4.761)rectangle++(0.04,0.04)
  (0.399,5.633)rectangle++(0.04,0.04) (0.562,6.368)rectangle++(0.04,0.04) (0.73,6.987)rectangle++(0.04,0.04) (0.898,7.506)rectangle++(0.04,0.04) (1.061,7.94)rectangle++(0.04,0.04)
  (1.218,8.303)rectangle++(0.04,0.04) (1.913,3.84)rectangle++(0.04,0.04) (-0.429,4.064)rectangle++(0.04,0.04) (-0.204,5.068)rectangle++(0.04,0.04) (0.027,5.912)rectangle++(0.04,0.04)
  (0.257,6.621)rectangle++(0.04,0.04) (0.481,7.213)rectangle++(0.04,0.04) (0.696,7.708)rectangle++(0.04,0.04) (0.897,8.12)rectangle++(0.04,0.04) (1.085,8.463)rectangle++(0.04,0.04)
  (-1.558,4.822)rectangle++(0.04,0.04) (-1.134,5.757)rectangle++(0.04,0.04) (-0.735,6.535)rectangle++(0.04,0.04) (-0.366,7.18)rectangle++(0.04,0.04) (-0.026,7.714)rectangle++(0.04,0.04)
  (0.284,8.154)rectangle++(0.04,0.04) (0.566,8.516)rectangle++(0.04,0.04) (0.819,8.812)rectangle++(0.04,0.04) (1.047,9.054)rectangle++(0.04,0.04) (1.25,9.25)rectangle++(0.04,0.04)
  (-1.696,5.008)rectangle++(0.04,0.04) (-1.243,5.921)rectangle++(0.04,0.04) (-0.003,1.937)rectangle++(0.04,0.04) (0.073,3.243)rectangle++(0.04,0.04) (0.795,2.604)rectangle++(0.04,0.04)
  (0.777,3.778)rectangle++(0.04,0.04) (0.81,4.776)rectangle++(0.04,0.04) (0.877,5.624)rectangle++(0.04,0.04) (0.968,6.341)rectangle++(0.04,0.04) (1.457,3.405)rectangle++(0.04,0.04)
  (1.372,4.432)rectangle++(0.04,0.04) (0.937,3.128)rectangle++(0.04,0.04) (-0.602,3.318)rectangle++(0.04,0.04) (-0.381,4.441)rectangle++(0.04,0.04) (-0.023,3.464)rectangle++(0.04,0.04)
  (0.117,4.542)rectangle++(0.04,0.04) (0.279,5.452)rectangle++(0.04,0.04) (0.453,6.219)rectangle++(0.04,0.04) (0.631,6.864)rectangle++(0.04,0.04) (0.809,7.406)rectangle++(0.04,0.04)
  (0.982,7.859)rectangle++(0.04,0.04) (1.147,8.237)rectangle++(0.04,0.04) (1.302,8.552)rectangle++(0.04,0.04) (1.965,3.907)rectangle++(0.04,0.04) (1.824,4.838)rectangle++(0.04,0.04)
  (1.742,5.636)rectangle++(0.04,0.04) (1.704,6.317)rectangle++(0.04,0.04) (1.699,6.897)rectangle++(0.04,0.04) (1.718,7.391)rectangle++(0.04,0.04) (1.754,7.81)rectangle++(0.04,0.04)
  (1.729,3.866)rectangle++(0.04,0.04) (0.755,3.1)rectangle++(0.04,0.04) (0.764,4.201)rectangle++(0.04,0.04) (0.815,5.137)rectangle++(0.04,0.04) (0.896,5.93)rectangle++(0.04,0.04)
  (0.997,6.601)rectangle++(0.04,0.04) (1.109,7.167)rectangle++(0.04,0.04) (-1.279,4.57)rectangle++(0.04,0.04) (-0.906,5.532)rectangle++(0.04,0.04) (1.476,2.823)rectangle++(0.04,0.04)
  (1.365,3.937)rectangle++(0.04,0.04) (1.316,4.888)rectangle++(0.04,0.04) (1.312,5.698)rectangle++(0.04,0.04) (1.341,6.387)rectangle++(0.04,0.04) (1.393,6.972)rectangle++(0.04,0.04)
  (1.461,7.467)rectangle++(0.04,0.04) (-1.279,4.695)rectangle++(0.04,0.04) (-0.902,5.639)rectangle++(0.04,0.04) (-0.543,6.425)rectangle++(0.04,0.04) (-0.207,7.079)rectangle++(0.04,0.04)
  (0.105,7.622)rectangle++(0.04,0.04) (0.392,8.071)rectangle++(0.04,0.04) (0.654,8.44)rectangle++(0.04,0.04) (0.891,8.744)rectangle++(0.04,0.04) (1.105,8.993)rectangle++(0.04,0.04)
  (1.297,9.196)rectangle++(0.04,0.04) (2.125,4.03)rectangle++(0.04,0.04) (1.966,4.937)rectangle++(0.04,0.04) (1.866,5.714)rectangle++(0.04,0.04) (1.812,6.378)rectangle++(0.04,0.04)
  (1.793,6.945)rectangle++(0.04,0.04) (-1.109,4.655)rectangle++(0.04,0.04) (-0.759,5.597)rectangle++(0.04,0.04) (-0.423,6.384)rectangle++(0.04,0.04) (-0.107,7.04)rectangle++(0.04,0.04)
  (0.189,7.584)rectangle++(0.04,0.04) (0.462,8.035)rectangle++(0.04,0.04) (1.936,3.616)rectangle++(0.04,0.04) (1.788,4.593)rectangle++(0.04,0.04) (0.928,3.254)rectangle++(0.04,0.04)
  (0.917,4.325)rectangle++(0.04,0.04) (0.95,5.236)rectangle++(0.04,0.04) (1.015,6.008)rectangle++(0.04,0.04) (1.101,6.663)rectangle++(0.04,0.04) (1.2,7.216)rectangle++(0.04,0.04)
  (1.306,7.681)rectangle++(0.04,0.04) (1.415,8.073)rectangle++(0.04,0.04) (1.524,8.402)rectangle++(0.04,0.04) (1.629,8.677)rectangle++(0.04,0.04) (1.73,8.906)rectangle++(0.04,0.04)
  (-1.537,5.03)rectangle++(0.04,0.04) (-1.107,5.933)rectangle++(0.04,0.04) (-0.706,6.684)rectangle++(0.04,0.04) (-0.335,7.306)rectangle++(0.04,0.04) (0.005,7.819)rectangle++(0.04,0.04)
  (0.315,8.242)rectangle++(0.04,0.04) (0.595,8.59)rectangle++(0.04,0.04) (0.847,8.874)rectangle++(0.04,0.04) (2.098,3.866)rectangle++(0.04,0.04) (1.936,4.798)rectangle++(0.04,0.04)
  (1.835,5.597)rectangle++(0.04,0.04) (1.782,6.28)rectangle++(0.04,0.04) (1.763,6.863)rectangle++(0.04,0.04) (1.038,2.308)rectangle++(0.04,0.04) (0.972,3.517)rectangle++(0.04,0.04)
  (-0.529,3.824)rectangle++(0.04,0.04) (1.009,2.58)rectangle++(0.04,0.04) (0.959,3.749)rectangle++(0.04,0.04) (0.963,4.744)rectangle++(0.04,0.04) (1.006,5.59)rectangle++(0.04,0.04)
  (1.076,6.308)rectangle++(0.04,0.04) (-1.102,4.395)rectangle++(0.04,0.04) (-0.763,5.376)rectangle++(0.04,0.04) (-0.436,6.196)rectangle++(0.04,0.04) (-0.125,6.88)rectangle++(0.04,0.04)
  (0.167,7.449)rectangle++(0.04,0.04) (0.438,7.922)rectangle++(0.04,0.04) (0.687,8.312)rectangle++(0.04,0.04) (-1.597,4.717)rectangle++(0.04,0.04) (-1.171,5.67)rectangle++(0.04,0.04)
  (-0.771,6.462)rectangle++(0.04,0.04) (-0.399,7.12)rectangle++(0.04,0.04) (-0.057,7.664)rectangle++(0.04,0.04) (0.256,8.113)rectangle++(0.04,0.04) (0.54,8.482)rectangle++(0.04,0.04)
  (0.796,8.784)rectangle++(0.04,0.04) (1.026,9.031)rectangle++(0.04,0.04) (1.231,9.232)rectangle++(0.04,0.04) (1.413,9.394)rectangle++(0.04,0.04) (1.575,9.524)rectangle++(0.04,0.04)
  (1.718,9.629)rectangle++(0.04,0.04) (1.843,9.712)rectangle++(0.04,0.04) (-1.687,5.21)rectangle++(0.04,0.04) (-1.228,6.093)rectangle++(0.04,0.04) (-0.802,6.824)rectangle++(0.04,0.04)
  (-0.411,7.429)rectangle++(0.04,0.04) (-0.055,7.927)rectangle++(0.04,0.04) (1.969,3.478)rectangle++(0.04,0.04) (1.811,4.474)rectangle++(0.04,0.04) (1.716,5.327)rectangle++(0.04,0.04)
  (1.669,6.055)rectangle++(0.04,0.04) (1.659,6.676)rectangle++(0.04,0.04) (1.675,7.205)rectangle++(0.04,0.04) (1.71,7.653)rectangle++(0.04,0.04) (1.757,8.033)rectangle++(0.04,0.04)
  (1.813,8.354)rectangle++(0.04,0.04) (1.873,8.625)rectangle++(0.04,0.04) (1.935,8.852)rectangle++(0.04,0.04) (1.996,9.043)rectangle++(0.04,0.04) (2.056,9.203)rectangle++(0.04,0.04)
  (2.038,4.204)rectangle++(0.04,0.04) (1.898,5.088)rectangle++(0.04,0.04) (1.038,3.374)rectangle++(0.04,0.04) (1.015,4.423)rectangle++(0.04,0.04) (1.037,5.315)rectangle++(0.04,0.04)
  (-0.899,4.191)rectangle++(0.04,0.04) (-0.598,5.195)rectangle++(0.04,0.04) (-1.166,3.861)rectangle++(0.04,0.04) (-0.839,4.925)rectangle++(0.04,0.04) (-0.518,5.816)rectangle++(0.04,0.04)
  (-0.21,6.56)rectangle++(0.04,0.04) (0.082,7.181)rectangle++(0.04,0.04) (0.354,7.697)rectangle++(0.04,0.04) (0.029,2.629)rectangle++(0.04,0.04) (0.128,3.83)rectangle++(0.04,0.04)
  (0.26,4.847)rectangle++(0.04,0.04) (0.412,5.705)rectangle++(0.04,0.04) (0.577,6.429)rectangle++(0.04,0.04) (0.745,7.038)rectangle++(0.04,0.04) (0.913,7.549)rectangle++(0.04,0.04)
  (1.762,3.625)rectangle++(0.04,0.04) (-0.41,3.994)rectangle++(0.04,0.04) (-0.191,5.007)rectangle++(0.04,0.04) (0.036,5.86)rectangle++(0.04,0.04) (0.018,2.38)rectangle++(0.04,0.04)
  (0.108,3.618)rectangle++(0.04,0.04) (0.234,4.667)rectangle++(0.04,0.04) (1.128,2.915)rectangle++(0.04,0.04) (1.074,4.029)rectangle++(0.04,0.04) (1.072,4.978)rectangle++(0.04,0.04)
  (1.108,5.785)rectangle++(0.04,0.04) (-0.986,4.297)rectangle++(0.04,0.04) (-0.668,5.288)rectangle++(0.04,0.04) (-0.359,6.118)rectangle++(0.04,0.04) (-0.062,6.811)rectangle++(0.04,0.04)
  (0.217,7.388)rectangle++(0.04,0.04) (0.478,7.867)rectangle++(0.04,0.04) (0.719,8.264)rectangle++(0.04,0.04) (0.939,8.592)rectangle++(0.04,0.04) (1.14,8.862)rectangle++(0.04,0.04)
  (1.321,9.083)rectangle++(0.04,0.04) (1.484,9.264)rectangle++(0.04,0.04) (1.63,9.411)rectangle++(0.04,0.04) (1.76,9.531)rectangle++(0.04,0.04) (1.875,9.627)rectangle++(0.04,0.04)
  (2.164,4.244)rectangle++(0.04,0.04) (2.007,5.117)rectangle++(0.04,0.04) (1.909,5.865)rectangle++(0.04,0.04) (1.855,6.505)rectangle++(0.04,0.04) (1.835,7.052)rectangle++(0.04,0.04)
  (1.839,7.517)rectangle++(0.04,0.04) (1.862,7.912)rectangle++(0.04,0.04) (1.897,8.247)rectangle++(0.04,0.04) (1.94,8.53)rectangle++(0.04,0.04) (1.988,8.769)rectangle++(0.04,0.04)
  (2.038,8.97)rectangle++(0.04,0.04) (2.089,9.14)rectangle++(0.04,0.04) (2.139,9.281)rectangle++(0.04,0.04) (2.187,9.4)rectangle++(0.04,0.04) (2.233,9.498)rectangle++(0.04,0.04)
  (2.276,9.581)rectangle++(0.04,0.04) (2.315,9.649)rectangle++(0.04,0.04) (2.352,9.705)rectangle++(0.04,0.04) (2.385,9.751)rectangle++(0.04,0.04) (2.415,9.789)rectangle++(0.04,0.04)
  (2.442,9.821)rectangle++(0.04,0.04) (2.467,9.846)rectangle++(0.04,0.04) (2.488,9.867)rectangle++(0.04,0.04) (2.507,9.883)rectangle++(0.04,0.04) (2.524,9.897)rectangle++(0.04,0.04)
  (2.539,9.907)rectangle++(0.04,0.04) (2.553,9.916)rectangle++(0.04,0.04) (2.564,9.923)rectangle++(0.04,0.04) (2.574,9.928)rectangle++(0.04,0.04) (2.583,9.932)rectangle++(0.04,0.04)
  (2.591,9.935)rectangle++(0.04,0.04) (2.597,9.937)rectangle++(0.04,0.04) (2.603,9.939)rectangle++(0.04,0.04) (2.608,9.94)rectangle++(0.04,0.04) (2.612,9.941)rectangle++(0.04,0.04)
  (2.616,9.942)rectangle++(0.04,0.04) (2.619,9.942)rectangle++(0.04,0.04) (2.621,9.942)rectangle++(0.04,0.04) (2.624,9.942)rectangle++(0.04,0.04) (2.626,9.942)rectangle++(0.04,0.04)
  (2.627,9.942)rectangle++(0.04,0.04) (2.629,9.942)rectangle++(0.04,0.04) (2.63,9.942)rectangle++(0.04,0.04) (-1.595,5.403)rectangle++(0.04,0.04) (1.525,2.642)rectangle++(0.04,0.04)
  (1.4,3.781)rectangle++(0.04,0.04) (-0.508,3.956)rectangle++(0.04,0.04) (-0.276,4.98)rectangle++(0.04,0.04) (-0.038,5.84)rectangle++(0.04,0.04) (0.199,6.562)rectangle++(0.04,0.04)
  (0.022,2.475)rectangle++(0.04,0.04) (0.116,3.699)rectangle++(0.04,0.04) (0.244,4.736)rectangle++(0.04,0.04) (0.395,5.612)rectangle++(0.04,0.04) (0.558,6.351)rectangle++(0.04,0.04)
  (-1.204,4.306)rectangle++(0.04,0.04) (-0.854,5.304)rectangle++(0.04,0.04) (-0.516,6.139)rectangle++(0.04,0.04) (0.23,2.707)rectangle++(0.04,0.04) (0.302,3.888)rectangle++(0.04,0.04)
  (0.41,4.889)rectangle++(0.04,0.04) (0.542,5.736)rectangle++(0.04,0.04) (0.688,6.45)rectangle++(0.04,0.04) (0.84,7.051)rectangle++(0.04,0.04) (0.994,7.556)rectangle++(0.04,0.04)
  (1.145,7.979)rectangle++(0.04,0.04) (1.29,8.333)rectangle++(0.04,0.04) (1.428,8.627)rectangle++(0.04,0.04) (1.969,3.95)rectangle++(0.04,0.04) (-0.442,4.097)rectangle++(0.04,0.04)
  (-0.214,5.096)rectangle++(0.04,0.04) (0.02,5.936)rectangle++(0.04,0.04) (0.252,6.641)rectangle++(0.04,0.04) (-1.321,4.306)rectangle++(0.04,0.04) (-0.953,5.309)rectangle++(0.04,0.04)
  (-0.6,6.147)rectangle++(0.04,0.04) (1.592,3.011)rectangle++(0.04,0.04) (1.472,4.092)rectangle++(0.04,0.04) (-0.561,4.031)rectangle++(0.04,0.04) (-0.318,5.045)rectangle++(0.04,0.04)
  (-0.071,5.898)rectangle++(0.04,0.04) (0.174,6.612)rectangle++(0.04,0.04) (0.41,7.209)rectangle++(0.04,0.04) (0.635,7.708)rectangle++(0.04,0.04) (0.845,8.122)rectangle++(0.04,0.04)
  (1.041,8.466)rectangle++(0.04,0.04) (1.222,8.751)rectangle++(0.04,0.04) (1.386,8.986)rectangle++(0.04,0.04) (-1.615,4.981)rectangle++(0.04,0.04) (-1.303,3.628)rectangle++(0.04,0.04)
  (-0.965,4.732)rectangle++(0.04,0.04) (-0.633,5.657)rectangle++(0.04,0.04) (-0.314,6.43)rectangle++(0.04,0.04) (-0.012,7.074)rectangle++(0.04,0.04) (0.271,7.61)rectangle++(0.04,0.04)
  (0.532,8.054)rectangle++(0.04,0.04) (-1.57,4.637)rectangle++(0.04,0.04) (-1.151,5.601)rectangle++(0.04,0.04) (-1.351,3.837)rectangle++(0.04,0.04) (-0.997,4.912)rectangle++(0.04,0.04)
  (-0.654,5.811)rectangle++(0.04,0.04) (1.516,2.934)rectangle++(0.04,0.04) (1.403,4.03)rectangle++(0.04,0.04) (1.352,4.965)rectangle++(0.04,0.04) (1.346,5.763)rectangle++(0.04,0.04)
  (1.372,6.441)rectangle++(0.04,0.04) (1.422,7.016)rectangle++(0.04,0.04) (-0.97,3.338)rectangle++(0.04,0.04) (-0.694,4.472)rectangle++(0.04,0.04) (-0.413,5.425)rectangle++(0.04,0.04)
  (-0.136,6.224)rectangle++(0.04,0.04) (0.131,6.892)rectangle++(0.04,0.04) (0.385,7.449)rectangle++(0.04,0.04) (0.623,7.913)rectangle++(0.04,0.04) (0.844,8.297)rectangle++(0.04,0.04)
  (1.047,8.615)rectangle++(0.04,0.04) (1.232,8.877)rectangle++(0.04,0.04) (1.4,9.092)rectangle++(0.04,0.04) (1.552,9.269)rectangle++(0.04,0.04) (1.688,9.413)rectangle++(0.04,0.04)
  (1.809,9.529)rectangle++(0.04,0.04) (-1.655,5.168)rectangle++(0.04,0.04) (-1.202,6.055)rectangle++(0.04,0.04) (-0.782,6.791)rectangle++(0.04,0.04) (-0.395,7.4)rectangle++(0.04,0.04)
  (-0.042,7.902)rectangle++(0.04,0.04) (0.278,8.315)rectangle++(0.04,0.04) (0.567,8.652)rectangle++(0.04,0.04) (0.826,8.928)rectangle++(0.04,0.04) (1.057,9.152)rectangle++(0.04,0.04)
  (1.262,9.333)rectangle++(0.04,0.04) (1.444,9.479)rectangle++(0.04,0.04) (1.604,9.596)rectangle++(0.04,0.04) (1.745,9.688)rectangle++(0.04,0.04) (2.195,4.226)rectangle++(0.04,0.04)
  (2.032,5.101)rectangle++(0.04,0.04) (-0.222,3.586)rectangle++(0.04,0.04) (0.914,2.603)rectangle++(0.04,0.04) (0.879,3.773)rectangle++(0.04,0.04) (0.896,4.768)rectangle++(0.04,0.04)
  (0.95,5.613)rectangle++(0.04,0.04) (1.03,6.329)rectangle++(0.04,0.04) (1.126,6.935)rectangle++(0.04,0.04) (1.232,7.446)rectangle++(0.04,0.04) (1.07,2.251)rectangle++(0.04,0.04)
  (0.997,3.467)rectangle++(0.04,0.04) (0.984,4.503)rectangle++(0.04,0.04) (1.014,5.384)rectangle++(0.04,0.04) (1.075,6.132)rectangle++(0.04,0.04) (1.393,3.389)rectangle++(0.04,0.04)
  (1.317,4.421)rectangle++(0.04,0.04) (1.294,5.301)rectangle++(0.04,0.04) (1.31,6.051)rectangle++(0.04,0.04) (1.353,6.687)rectangle++(0.04,0.04) (1.416,7.226)rectangle++(0.04,0.04)
  (-1.236,4.639)rectangle++(0.04,0.04) (-0.867,5.589)rectangle++(0.04,0.04) (-0.516,6.382)rectangle++(0.04,0.04) (-0.186,7.041)rectangle++(0.04,0.04) (-1.484,4.303)rectangle++(0.04,0.04)
  (-1.092,5.313)rectangle++(0.04,0.04) (-0.717,6.156)rectangle++(0.04,0.04) (-0.366,6.857)rectangle++(0.04,0.04) (-0.933,3.191)rectangle++(0.04,0.04) (-0.667,4.345)rectangle++(0.04,0.04)
  (-0.396,5.317)rectangle++(0.04,0.04) (-0.126,6.131)rectangle++(0.04,0.04) (0.136,6.813)rectangle++(0.04,0.04) (0.386,7.382)rectangle++(0.04,0.04) (0.621,7.855)rectangle++(0.04,0.04)
  (0.84,8.248)rectangle++(0.04,0.04) (1.042,8.574)rectangle++(0.04,0.04) (-1.111,3.298)rectangle++(0.04,0.04) (-0.815,4.444)rectangle++(0.04,0.04) (-0.517,5.406)rectangle++(0.04,0.04)
  (-1.196,3.918)rectangle++(0.04,0.04) (-0.862,4.974)rectangle++(0.04,0.04) (-0.536,5.859)rectangle++(0.04,0.04) (-0.224,6.598)rectangle++(0.04,0.04) (0.072,7.213)rectangle++(0.04,0.04)
  (0.347,7.725)rectangle++(0.04,0.04) (0.602,8.148)rectangle++(0.04,0.04) (0.835,8.498)rectangle++(0.04,0.04) (-1.61,4.782)rectangle++(0.04,0.04) (1.373,2.519)rectangle++(0.04,0.04)
  (1.265,3.682)rectangle++(0.04,0.04) (1.221,4.676)rectangle++(0.04,0.04) (1.222,5.522)rectangle++(0.04,0.04) (1.258,6.241)rectangle++(0.04,0.04) (1.316,6.85)rectangle++(0.04,0.04)
  (1.391,7.366)rectangle++(0.04,0.04) (1.475,7.802)rectangle++(0.04,0.04) (1.759,3.801)rectangle++(0.04,0.04) (1.645,4.757)rectangle++(0.04,0.04) (0.985,3.253)rectangle++(0.04,0.04)
  (-0.47,3.775)rectangle++(0.04,0.04) (-0.842,3.605)rectangle++(0.04,0.04) (-0.574,4.694)rectangle++(0.04,0.04) (-0.302,5.609)rectangle++(0.04,0.04) (-0.035,6.376)rectangle++(0.04,0.04)
  (0.223,7.017)rectangle++(0.04,0.04) (0.468,7.552)rectangle++(0.04,0.04) (0.698,7.996)rectangle++(0.04,0.04) (0.911,8.365)rectangle++(0.04,0.04) (-1.569,4.77)rectangle++(0.04,0.04)
  (-1.145,5.713)rectangle++(0.04,0.04) (-0.747,6.498)rectangle++(0.04,0.04) (-1.469,4.09)rectangle++(0.04,0.04) (-1.088,5.132)rectangle++(0.04,0.04) (1.398,2.705)rectangle++(0.04,0.04)
  (1.294,3.839)rectangle++(0.04,0.04) (1.252,4.808)rectangle++(0.04,0.04) (1.254,5.633)rectangle++(0.04,0.04) (0.917,2.283)rectangle++(0.04,0.04) (0.868,3.5)rectangle++(0.04,0.04)
  (0.876,4.537)rectangle++(0.04,0.04) (0.924,5.417)rectangle++(0.04,0.04) (1,6.164)rectangle++(0.04,0.04) (1.094,6.796)rectangle++(0.04,0.04) (1.2,7.329)rectangle++(0.04,0.04)
  (1.311,7.778)rectangle++(0.04,0.04) (1.423,8.155)rectangle++(0.04,0.04) (1.533,8.471)rectangle++(0.04,0.04) (1.64,8.735)rectangle++(0.04,0.04) (1.741,8.955)rectangle++(0.04,0.04)
  (1.836,9.139)rectangle++(0.04,0.04) (1.924,9.291)rectangle++(0.04,0.04) (2.005,9.416)rectangle++(0.04,0.04) (2.079,9.52)rectangle++(0.04,0.04) (2.145,9.605)rectangle++(0.04,0.04)
  (-1.61,5.246)rectangle++(0.04,0.04) (-1.161,6.12)rectangle++(0.04,0.04) (-0.744,6.844)rectangle++(0.04,0.04) (-0.361,7.444)rectangle++(0.04,0.04) (-0.011,7.938)rectangle++(0.04,0.04)
  (0.306,8.344)rectangle++(0.04,0.04) (0.591,8.676)rectangle++(0.04,0.04) (0.847,8.947)rectangle++(0.04,0.04) (1.076,9.168)rectangle++(0.04,0.04) (1.211,2.208)rectangle++(0.04,0.04)
  (1.115,3.425)rectangle++(0.04,0.04) (1.083,4.463)rectangle++(0.04,0.04) (1.097,5.346)rectangle++(0.04,0.04) (1.144,6.097)rectangle++(0.04,0.04) (1.214,6.733)rectangle++(0.04,0.04)
  (1.525,3.536)rectangle++(0.04,0.04) (1.436,4.541)rectangle++(0.04,0.04) (1.4,5.398)rectangle++(0.04,0.04) (1.403,6.129)rectangle++(0.04,0.04) (1.436,6.75)rectangle++(0.04,0.04)
  (1.488,7.276)rectangle++(0.04,0.04) (1.554,7.721)rectangle++(0.04,0.04) (1.73,3.804)rectangle++(0.04,0.04) (1.62,4.76)rectangle++(0.04,0.04) (0.989,3.248)rectangle++(0.04,0.04)
  (-0.469,3.774)rectangle++(0.04,0.04) (-0.25,4.823)rectangle++(0.04,0.04) (-0.021,5.706)rectangle++(0.04,0.04) (0.208,6.447)rectangle++(0.04,0.04) (0.432,7.068)rectangle++(0.04,0.04)
  (0.648,7.587)rectangle++(0.04,0.04) (-1.45,4.567)rectangle++(0.04,0.04) (-1.052,5.536)rectangle++(0.04,0.04) (-0.675,6.344)rectangle++(0.04,0.04) (-0.322,7.016)rectangle++(0.04,0.04)
  (0.004,7.572)rectangle++(0.04,0.04) (0.304,8.032)rectangle++(0.04,0.04) (0.578,8.412)rectangle++(0.04,0.04) (0.825,8.723)rectangle++(0.04,0.04) (-1.66,4.824)rectangle++(0.04,0.04)
  (-1.22,5.763)rectangle++(0.04,0.04) (-0.809,6.544)rectangle++(0.04,0.04) (-0.428,7.191)rectangle++(0.04,0.04) (-0.078,7.725)rectangle++(0.04,0.04) (0.24,8.166)rectangle++(0.04,0.04)
  (0.529,8.528)rectangle++(0.04,0.04) (0.788,8.824)rectangle++(0.04,0.04) (1.021,9.065)rectangle++(0.04,0.04) (1.228,9.26)rectangle++(0.04,0.04) (1.412,9.418)rectangle++(0.04,0.04)
  (1.575,9.545)rectangle++(0.04,0.04) (1.718,9.647)rectangle++(0.04,0.04) (1.844,9.727)rectangle++(0.04,0.04) (1.954,9.791)rectangle++(0.04,0.04) (2.051,9.84)rectangle++(0.04,0.04)
  (2.134,9.878)rectangle++(0.04,0.04) (2.196,4.351)rectangle++(0.04,0.04) (2.039,5.207)rectangle++(0.04,0.04) (0.205,2.115)rectangle++(0.04,0.04) (0.256,3.386)rectangle++(0.04,0.04)
  (0.351,4.464)rectangle++(0.04,0.04) (0.475,5.377)rectangle++(0.04,0.04) (1.276,3.107)rectangle++(0.04,0.04) (1.207,4.186)rectangle++(0.04,0.04) (1.191,5.106)rectangle++(0.04,0.04)
  (1.214,5.889)rectangle++(0.04,0.04) (1.265,6.553)rectangle++(0.04,0.04) (1.336,7.116)rectangle++(0.04,0.04) (1.418,7.591)rectangle++(0.04,0.04) (1.506,7.992)rectangle++(0.04,0.04)
  (1.598,8.329)rectangle++(0.04,0.04) (1.689,8.612)rectangle++(0.04,0.04) (1.778,8.849)rectangle++(0.04,0.04) (1.863,9.047)rectangle++(0.04,0.04) (1.943,9.211)rectangle++(0.04,0.04)
  (2.018,9.348)rectangle++(0.04,0.04) (2.078,4.223)rectangle++(0.04,0.04) (1.933,5.102)rectangle++(0.04,0.04) (1.845,5.856)rectangle++(0.04,0.04) (1.801,6.5)rectangle++(0.04,0.04)
  (1.788,7.049)rectangle++(0.04,0.04) (1.8,7.516)rectangle++(0.04,0.04) (1.651,3.82)rectangle++(0.04,0.04) (1.554,4.777)rectangle++(0.04,0.04) (1.509,5.595)rectangle++(0.04,0.04)
  (1.505,6.291)rectangle++(0.04,0.04) (1.528,6.884)rectangle++(0.04,0.04) (1.572,7.386)rectangle++(0.04,0.04) (1.63,7.812)rectangle++(0.04,0.04) (1.695,8.171)rectangle++(0.04,0.04)
  (1.766,8.474)rectangle++(0.04,0.04) (1.838,8.728)rectangle++(0.04,0.04) (-1.479,5.015)rectangle++(0.04,0.04) (-1.059,5.918)rectangle++(0.04,0.04) (-0.666,6.669)rectangle++(0.04,0.04)
  (-0.301,7.292)rectangle++(0.04,0.04) (-1.558,4.331)rectangle++(0.04,0.04) (1.255,2.444)rectangle++(0.04,0.04) (1.162,3.623)rectangle++(0.04,0.04) (1.131,4.629)rectangle++(0.04,0.04)
  (-0.736,4.074)rectangle++(0.04,0.04) (-0.465,5.088)rectangle++(0.04,0.04) (1.315,2.838)rectangle++(0.04,0.04) (1.229,3.956)rectangle++(0.04,0.04) (1.2,4.909)rectangle++(0.04,0.04)
  (1.214,5.721)rectangle++(0.04,0.04) (1.275,3.341)rectangle++(0.04,0.04) (1.215,4.385)rectangle++(0.04,0.04) (1.206,5.275)rectangle++(0.04,0.04) (1.234,6.032)rectangle++(0.04,0.04)
  (-1.016,4.37)rectangle++(0.04,0.04) (-0.691,5.351)rectangle++(0.04,0.04) (-0.375,6.172)rectangle++(0.04,0.04) (-0.074,6.858)rectangle++(0.04,0.04) (0.209,7.428)rectangle++(0.04,0.04)
  (0.473,7.902)rectangle++(0.04,0.04) (0.716,8.294)rectangle++(0.04,0.04) (0.033,2.711)rectangle++(0.04,0.04) (0.134,3.899)rectangle++(0.04,0.04) (0.268,4.905)rectangle++(0.04,0.04)
  (0.422,5.755)rectangle++(0.04,0.04) (0.064,3.291)rectangle++(0.04,0.04) (0.184,4.391)rectangle++(0.04,0.04) (-0.855,3.849)rectangle++(0.04,0.04) (-0.575,4.902)rectangle++(0.04,0.04)
  (-0.295,5.786)rectangle++(0.04,0.04) (-0.022,6.526)rectangle++(0.04,0.04) (0.24,7.144)rectangle++(0.04,0.04) (0.488,7.659)rectangle++(0.04,0.04) (0.719,8.087)rectangle++(0.04,0.04)
  (0.932,8.441)rectangle++(0.04,0.04) (1.128,8.734)rectangle++(0.04,0.04) (2.031,3.902)rectangle++(0.04,0.04) (1.88,4.832)rectangle++(0.04,0.04) (1.789,5.628)rectangle++(0.04,0.04)
  (1.744,6.309)rectangle++(0.04,0.04) (-0.983,4.52)rectangle++(0.04,0.04) (-0.657,5.477)rectangle++(0.04,0.04) (-0.341,6.278)rectangle++(0.04,0.04) (-0.041,6.946)rectangle++(0.04,0.04)
  (0.241,7.502)rectangle++(0.04,0.04) (0.502,7.964)rectangle++(0.04,0.04) (0.743,8.345)rectangle++(0.04,0.04) (0.963,8.66)rectangle++(0.04,0.04) (1.163,8.919)rectangle++(0.04,0.04)
  (1.192,2.22)rectangle++(0.04,0.04) (1.1,3.436)rectangle++(0.04,0.04) (1.07,4.473)rectangle++(0.04,0.04) (1.087,5.355)rectangle++(0.04,0.04) (-0.898,4.209)rectangle++(0.04,0.04)
  (-0.597,5.209)rectangle++(0.04,0.04) (-0.302,6.048)rectangle++(0.04,0.04) (-0.017,6.749)rectangle++(0.04,0.04) (0.254,7.334)rectangle++(0.04,0.04) (0.507,7.82)rectangle++(0.04,0.04)
  (0.741,8.223)rectangle++(0.04,0.04) (0.957,8.556)rectangle++(0.04,0.04) (1.153,8.83)rectangle++(0.04,0.04) (1.331,9.056)rectangle++(0.04,0.04) (1.492,9.24)rectangle++(0.04,0.04)
  (1.635,9.391)rectangle++(0.04,0.04) (1.763,9.513)rectangle++(0.04,0.04) (1.877,9.612)rectangle++(0.04,0.04) (1.978,9.691)rectangle++(0.04,0.04) (-1.659,5.232)rectangle++(0.04,0.04)
  (-1.203,6.109)rectangle++(0.04,0.04) (-0.78,6.837)rectangle++(0.04,0.04) (-0.392,7.439)rectangle++(0.04,0.04) (-0.038,7.935)rectangle++(0.04,0.04) (1.969,3.484)rectangle++(0.04,0.04)
  (1.811,4.479)rectangle++(0.04,0.04) (1.716,5.331)rectangle++(0.04,0.04) (1.67,6.059)rectangle++(0.04,0.04) (1.66,6.679)rectangle++(0.04,0.04) (1.676,7.207)rectangle++(0.04,0.04)
  (1.71,7.655)rectangle++(0.04,0.04) (1.758,8.035)rectangle++(0.04,0.04) (1.784,3.911)rectangle++(0.04,0.04) (1.67,4.849)rectangle++(0.04,0.04) (1.612,5.651)rectangle++(0.04,0.04)
  (1.594,6.335)rectangle++(0.04,0.04) (1.606,6.917)rectangle++(0.04,0.04) (1.639,7.412)rectangle++(0.04,0.04) (1.688,7.83)rectangle++(0.04,0.04) (1.746,8.185)rectangle++(0.04,0.04)
  (1.809,8.483)rectangle++(0.04,0.04) (0.034,2.736)rectangle++(0.04,0.04) (0.137,3.921)rectangle++(0.04,0.04) (0.955,2.749)rectangle++(0.04,0.04) (0.92,3.895)rectangle++(0.04,0.04)
  (0.935,4.87)rectangle++(0.04,0.04) (0.988,5.698)rectangle++(0.04,0.04) (0.054,3.306)rectangle++(0.04,0.04) (0.176,4.404)rectangle++(0.04,0.04) (0.323,5.332)rectangle++(0.04,0.04)
  (0.486,6.116)rectangle++(0.04,0.04) (0.655,6.775)rectangle++(0.04,0.04) (0.826,7.329)rectangle++(0.04,0.04) (0.993,7.793)rectangle++(0.04,0.04) (1.154,8.18)rectangle++(0.04,0.04)
  (1.306,8.503)rectangle++(0.04,0.04) (1.448,8.772)rectangle++(0.04,0.04) (1.579,8.994)rectangle++(0.04,0.04) (1.7,9.178)rectangle++(0.04,0.04) (1.81,9.33)rectangle++(0.04,0.04)
  (1.909,9.454)rectangle++(0.04,0.04) (1.999,9.556)rectangle++(0.04,0.04) (2.079,9.639)rectangle++(0.04,0.04) (2.151,9.706)rectangle++(0.04,0.04) (2.214,9.76)rectangle++(0.04,0.04)
  (2.27,9.804)rectangle++(0.04,0.04) (2.32,9.839)rectangle++(0.04,0.04) (-1.629,5.325)rectangle++(0.04,0.04) (-1.174,6.187)rectangle++(0.04,0.04) (-0.752,6.902)rectangle++(0.04,0.04)
  (1.797,3.122)rectangle++(0.04,0.04) (1.65,4.178)rectangle++(0.04,0.04) (1.568,5.082)rectangle++(0.04,0.04) (1.534,5.853)rectangle++(0.04,0.04) (1.536,6.51)rectangle++(0.04,0.04)
  (1.563,7.068)rectangle++(0.04,0.04) (1.609,7.542)rectangle++(0.04,0.04) (1.667,7.942)rectangle++(0.04,0.04) (-1.343,4.828)rectangle++(0.04,0.04) (-0.951,5.754)rectangle++(0.04,0.04)
  (-0.58,6.525)rectangle++(0.04,0.04) (-0.234,7.166)rectangle++(0.04,0.04) (0.085,7.696)rectangle++(0.04,0.04) (1.896,3.466)rectangle++(0.04,0.04) (1.748,4.466)rectangle++(0.04,0.04)
  (1.662,5.322)rectangle++(0.04,0.04) (1.624,6.054)rectangle++(0.04,0.04) (-0.95,4.447)rectangle++(0.04,0.04) (-0.632,5.414)rectangle++(0.04,0.04) (-0.323,6.224)rectangle++(0.04,0.04)
  (-0.028,6.899)rectangle++(0.04,0.04) (0.25,7.462)rectangle++(0.04,0.04) (0.509,7.929)rectangle++(0.04,0.04) (0.748,8.315)rectangle++(0.04,0.04) (0.966,8.634)rectangle++(0.04,0.04)
  (1.164,8.897)rectangle++(0.04,0.04) (2.067,3.942)rectangle++(0.04,0.04) (0.734,3.192)rectangle++(0.04,0.04) (0.75,4.28)rectangle++(0.04,0.04) (0.806,5.204)rectangle++(0.04,0.04)
  (0.891,5.987)rectangle++(0.04,0.04) (0.995,6.65)rectangle++(0.04,0.04) (1.109,7.209)rectangle++(0.04,0.04) (1.229,7.679)rectangle++(0.04,0.04) (1.35,8.074)rectangle++(0.04,0.04)
  (1.468,8.405)rectangle++(0.04,0.04) (1.582,8.682)rectangle++(0.04,0.04) (1.965,3.995)rectangle++(0.04,0.04) (1.828,4.914)rectangle++(0.04,0.04) (1.748,5.7)rectangle++(0.04,0.04)
  (1.711,6.371)rectangle++(0.04,0.04) (1.707,6.943)rectangle++(0.04,0.04) (1.727,7.43)rectangle++(0.04,0.04) (1.763,7.842)rectangle++(0.04,0.04) (1.81,8.192)rectangle++(0.04,0.04)
  (1.864,8.487)rectangle++(0.04,0.04) (1.922,8.735)rectangle++(0.04,0.04) (1.98,8.944)rectangle++(0.04,0.04) (2.039,9.12)rectangle++(0.04,0.04) (2.096,9.266)rectangle++(0.04,0.04)
  (2.049,4.225)rectangle++(0.04,0.04) (1.909,5.105)rectangle++(0.04,0.04) (-0.698,4.313)rectangle++(0.04,0.04) (-0.423,5.29)rectangle++(0.04,0.04) (-0.15,6.11)rectangle++(0.04,0.04)
  (0.115,6.796)rectangle++(0.04,0.04) (0.367,7.368)rectangle++(0.04,0.04) (0.605,7.844)rectangle++(0.04,0.04) (-1.513,4.61)rectangle++(0.04,0.04) (-1.104,5.575)rectangle++(0.04,0.04)
  (1.51,2.787)rectangle++(0.04,0.04) (1.393,3.904)rectangle++(0.04,0.04) (-0.536,3.98)rectangle++(0.04,0.04) (-0.298,5)rectangle++(0.04,0.04) (-0.056,5.858)rectangle++(0.04,0.04)
  (0.185,6.578)rectangle++(0.04,0.04) (0.418,7.18)rectangle++(0.04,0.04) (0.64,7.683)rectangle++(0.04,0.04) (0.849,8.101)rectangle++(0.04,0.04) (1.044,8.448)rectangle++(0.04,0.04)
  (1.223,8.735)rectangle++(0.04,0.04) (1.387,8.972)rectangle++(0.04,0.04) (1.535,9.167)rectangle++(0.04,0.04) (1.67,9.327)rectangle++(0.04,0.04) (2.114,4.14)rectangle++(0.04,0.04)
  (1.96,5.03)rectangle++(0.04,0.04) (1.865,5.794)rectangle++(0.04,0.04) (0.017,2.371)rectangle++(0.04,0.04) (0.107,3.611)rectangle++(0.04,0.04) (0.233,4.661)rectangle++(0.04,0.04)
  (0.044,3.268)rectangle++(0.04,0.04) (0.166,4.372)rectangle++(0.04,0.04) (0.314,5.306)rectangle++(0.04,0.04) (0.477,6.094)rectangle++(0.04,0.04) (0.647,6.757)rectangle++(0.04,0.04)
  (0.818,7.314)rectangle++(0.04,0.04) (0.986,7.78)rectangle++(0.04,0.04) (1.147,8.17)rectangle++(0.04,0.04) (1.299,8.495)rectangle++(0.04,0.04) (1.442,8.765)rectangle++(0.04,0.04)
  (1.574,8.989)rectangle++(0.04,0.04) (1.695,9.174)rectangle++(0.04,0.04) (1.806,9.326)rectangle++(0.04,0.04) (1.906,9.451)rectangle++(0.04,0.04) (1.996,9.553)rectangle++(0.04,0.04)
  (2.076,9.637)rectangle++(0.04,0.04) (2.148,9.704)rectangle++(0.04,0.04) (2.212,9.759)rectangle++(0.04,0.04) (2.268,9.803)rectangle++(0.04,0.04) (2.318,9.838)rectangle++(0.04,0.04)
  (1.309,2.309)rectangle++(0.04,0.04) (-0.53,3.354)rectangle++(0.04,0.04) (-0.319,4.468)rectangle++(0.04,0.04) (-0.094,5.407)rectangle++(0.04,0.04) (0.134,6.196)rectangle++(0.04,0.04)
  (0.36,6.857)rectangle++(0.04,0.04) (0.578,7.41)rectangle++(0.04,0.04) (0.785,7.872)rectangle++(0.04,0.04) (0.137,2.063)rectangle++(0.04,0.04) (0.197,3.344)rectangle++(0.04,0.04)
  (0.299,4.431)rectangle++(0.04,0.04) (0.429,5.35)rectangle++(0.04,0.04) (0.577,6.127)rectangle++(0.04,0.04) (0.733,6.781)rectangle++(0.04,0.04) (0.892,7.331)rectangle++(0.04,0.04)
  (1.049,7.792)rectangle++(0.04,0.04) (1.201,8.177)rectangle++(0.04,0.04) (1.346,8.499)rectangle++(0.04,0.04) (0.11,2.042)rectangle++(0.04,0.04) (0.173,3.327)rectangle++(0.04,0.04)
  (0.278,4.417)rectangle++(0.04,0.04) (1.061,2.877)rectangle++(0.04,0.04) (1.015,3.999)rectangle++(0.04,0.04) (1.021,4.955)rectangle++(0.04,0.04) (1.063,5.767)rectangle++(0.04,0.04)
  (1.132,6.456)rectangle++(0.04,0.04) (1.466,3.464)rectangle++(0.04,0.04) (1.382,4.482)rectangle++(0.04,0.04) (1.352,5.351)rectangle++(0.04,0.04) (1.361,6.09)rectangle++(0.04,0.04)
  (1.398,6.718)rectangle++(0.04,0.04) (1.455,7.251)rectangle++(0.04,0.04) (1.525,7.701)rectangle++(0.04,0.04) (1.602,8.081)rectangle++(0.04,0.04) (-1.384,4.843)rectangle++(0.04,0.04)
  (1.362,2.582)rectangle++(0.04,0.04) (1.258,3.736)rectangle++(0.04,0.04) (1.217,4.722)rectangle++(0.04,0.04) (1.221,5.561)rectangle++(0.04,0.04) (1.258,6.274)rectangle++(0.04,0.04)
  (1.318,6.879)rectangle++(0.04,0.04) (1.393,7.391)rectangle++(0.04,0.04) (1.478,7.822)rectangle++(0.04,0.04) (1.567,8.186)rectangle++(0.04,0.04) (1.657,8.492)rectangle++(0.04,0.04)
  (1.746,8.748)rectangle++(0.04,0.04) (1.832,8.962)rectangle++(0.04,0.04) (1.913,9.141)rectangle++(0.04,0.04) (0.039,2.826)rectangle++(0.04,0.04) (0.144,3.997)rectangle++(0.04,0.04)
  (0.28,4.988)rectangle++(0.04,0.04) (0.435,5.824)rectangle++(0.04,0.04) (0.601,6.53)rectangle++(0.04,0.04) (0.77,7.122)rectangle++(0.04,0.04) (0.937,7.619)rectangle++(0.04,0.04)
  (1.099,8.035)rectangle++(0.04,0.04) (-1.465,4.74)rectangle++(0.04,0.04) (1.345,2.544)rectangle++(0.04,0.04) (1.243,3.705)rectangle++(0.04,0.04) (0.772,2.959)rectangle++(0.04,0.04)
  (0.772,4.08)rectangle++(0.04,0.04) (0.817,5.034)rectangle++(0.04,0.04) (-0.878,4.095)rectangle++(0.04,0.04) (-0.585,5.112)rectangle++(0.04,0.04) (-0.295,5.965)rectangle++(0.04,0.04)
  (-0.014,6.678)rectangle++(0.04,0.04) (0.253,7.273)rectangle++(0.04,0.04) (0.504,7.768)rectangle++(0.04,0.04) (-1.515,4.576)rectangle++(0.04,0.04) (-1.107,5.546)rectangle++(0.04,0.04)
  (0.218,2.126)rectangle++(0.04,0.04) (0.268,3.394)rectangle++(0.04,0.04) (0.809,2.678)rectangle++(0.04,0.04) (0.793,3.84)rectangle++(0.04,0.04) (0.825,4.828)rectangle++(0.04,0.04)
  (0.892,5.667)rectangle++(0.04,0.04) (0.983,6.378)rectangle++(0.04,0.04) (1.088,6.978)rectangle++(0.04,0.04) (1.202,7.484)rectangle++(0.04,0.04) (1.319,7.909)rectangle++(0.04,0.04)
  (1.435,8.266)rectangle++(0.04,0.04) (1.548,8.565)rectangle++(0.04,0.04) (1.657,8.815)rectangle++(0.04,0.04) (1.759,9.022)rectangle++(0.04,0.04) (1.853,9.195)rectangle++(0.04,0.04)
  (1.941,9.338)rectangle++(0.04,0.04) (2.021,9.456)rectangle++(0.04,0.04) (2.094,9.553)rectangle++(0.04,0.04) (2.12,4.279)rectangle++(0.04,0.04) (1.971,5.149)rectangle++(0.04,0.04)
  (1.879,5.894)rectangle++(0.04,0.04) (1.831,6.531)rectangle++(0.04,0.04) (1.815,7.074)rectangle++(0.04,0.04) (1.824,7.537)rectangle++(0.04,0.04) (1.849,7.929)rectangle++(0.04,0.04)
  (-1.308,4.86)rectangle++(0.04,0.04) (-0.919,5.779)rectangle++(0.04,0.04) (-0.553,6.545)rectangle++(0.04,0.04) (1.686,3.099)rectangle++(0.04,0.04) (1.554,4.163)rectangle++(0.04,0.04)
  (1.486,5.072)rectangle++(0.04,0.04) (1.463,5.848)rectangle++(0.04,0.04) (1.277,3.422)rectangle++(0.04,0.04) (-0.454,3.863)rectangle++(0.04,0.04) (-0.234,4.898)rectangle++(0.04,0.04)
  (-0.005,5.769)rectangle++(0.04,0.04) (0.224,6.5)rectangle++(0.04,0.04) (1.583,3.266)rectangle++(0.04,0.04) (1.474,4.309)rectangle++(0.04,0.04) (1.423,5.2)rectangle++(0.04,0.04)
  (1.416,5.959)rectangle++(0.04,0.04) (1.439,6.605)rectangle++(0.04,0.04) (1.486,7.153)rectangle++(0.04,0.04) (1.547,7.617)rectangle++(0.04,0.04) (1.617,8.009)rectangle++(0.04,0.04)
  (1.693,8.339)rectangle++(0.04,0.04) (1.77,8.617)rectangle++(0.04,0.04) (1.847,8.849)rectangle++(0.04,0.04) (1.922,9.044)rectangle++(0.04,0.04) (1.993,9.207)rectangle++(0.04,0.04)
  (-1.553,5.139)rectangle++(0.04,0.04) (-1.117,6.026)rectangle++(0.04,0.04) (-0.71,6.763)rectangle++(0.04,0.04) (-0.336,7.373)rectangle++(0.04,0.04) (0.027,2.586)rectangle++(0.04,0.04)
  (-0.5,3.463)rectangle++(0.04,0.04) (-0.289,4.56)rectangle++(0.04,0.04) (-0.065,5.484)rectangle++(0.04,0.04) (0.161,6.26)rectangle++(0.04,0.04) (0.385,6.911)rectangle++(0.04,0.04)
  (0.602,7.455)rectangle++(0.04,0.04) (0.808,7.909)rectangle++(0.04,0.04) (1.001,8.286)rectangle++(0.04,0.04) (1.18,8.6)rectangle++(0.04,0.04) (1.345,8.859)rectangle++(0.04,0.04)
  (1.495,9.072)rectangle++(0.04,0.04) (1.631,9.248)rectangle++(0.04,0.04) (1.754,9.392)rectangle++(0.04,0.04) (1.865,9.509)rectangle++(0.04,0.04) (1.963,9.604)rectangle++(0.04,0.04)
  (2.051,9.681)rectangle++(0.04,0.04) (2.128,9.743)rectangle++(0.04,0.04) (-1.643,5.27)rectangle++(0.04,0.04) (-1.188,6.141)rectangle++(0.04,0.04) (1.66,2.877)rectangle++(0.04,0.04)
  (1.524,3.975)rectangle++(0.04,0.04) (1.452,4.914)rectangle++(0.04,0.04) (1.429,5.715)rectangle++(0.04,0.04) (1.441,6.397)rectangle++(0.04,0.04) (1.478,6.976)rectangle++(0.04,0.04)
  (1.533,7.467)rectangle++(0.04,0.04) (1.6,7.881)rectangle++(0.04,0.04) (1.673,8.231)rectangle++(0.04,0.04) (1.749,8.526)rectangle++(0.04,0.04) (1.826,8.773)rectangle++(0.04,0.04)
  (1.9,8.981)rectangle++(0.04,0.04) (1.972,9.154)rectangle++(0.04,0.04) (2.04,9.298)rectangle++(0.04,0.04) (2.104,9.418)rectangle++(0.04,0.04) (2.163,9.517)rectangle++(0.04,0.04)
  (2.217,9.599)rectangle++(0.04,0.04) (2.266,9.667)rectangle++(0.04,0.04) (2.311,9.722)rectangle++(0.04,0.04) (2.351,9.768)rectangle++(0.04,0.04) (2.387,9.805)rectangle++(0.04,0.04)
  (2.419,9.835)rectangle++(0.04,0.04) (2.447,9.859)rectangle++(0.04,0.04) (2.472,9.879)rectangle++(0.04,0.04) (-0.261,4.073)rectangle++(0.04,0.04) (-0.061,5.069)rectangle++(0.04,0.04)
  (0.148,5.907)rectangle++(0.04,0.04) (0.36,6.611)rectangle++(0.04,0.04) (0.568,7.202)rectangle++(0.04,0.04) (0.769,7.695)rectangle++(0.04,0.04) (0.959,8.106)rectangle++(0.04,0.04)
  (1.137,8.448)rectangle++(0.04,0.04) (1.303,8.731)rectangle++(0.04,0.04) (1.454,8.966)rectangle++(0.04,0.04) (1.593,9.159)rectangle++(0.04,0.04) (2.081,4.09)rectangle++(0.04,0.04)
  (1.931,4.989)rectangle++(0.04,0.04) (1.838,5.76)rectangle++(0.04,0.04) (1.791,6.419)rectangle++(0.04,0.04) (1.777,6.98)rectangle++(0.04,0.04) (1.787,7.458)rectangle++(0.04,0.04)
  (1.815,7.864)rectangle++(0.04,0.04) (1.855,8.208)rectangle++(0.04,0.04) (-1.366,4.916)rectangle++(0.04,0.04) (0.768,2.06)rectangle++(0.04,0.04) (0.733,3.316)rectangle++(0.04,0.04)
  (0.753,4.386)rectangle++(0.04,0.04) (-0.752,3.955)rectangle++(0.04,0.04) (-0.483,4.988)rectangle++(0.04,0.04) (-0.214,5.855)rectangle++(0.04,0.04) (0.05,6.582)rectangle++(0.04,0.04)
  (0.304,7.189)rectangle++(0.04,0.04) (1.745,3.417)rectangle++(0.04,0.04) (0.642,3.017)rectangle++(0.04,0.04) (0.664,4.135)rectangle++(0.04,0.04) (0.728,5.085)rectangle++(0.04,0.04)
  (1.174,3.108)rectangle++(0.04,0.04) (0.699,2.318)rectangle++(0.04,0.04) (-0.11,3.277)rectangle++(0.04,0.04) (0.035,4.386)rectangle++(0.04,0.04) (0.203,5.323)rectangle++(0.04,0.04)
  (1.294,3.035)rectangle++(0.04,0.04) (1.219,4.124)rectangle++(0.04,0.04) (1.199,5.053)rectangle++(0.04,0.04) (1.219,5.843)rectangle++(0.04,0.04) (1.268,6.514)rectangle++(0.04,0.04)
  (1.465,3.506)rectangle++(0.04,0.04) (0.697,2.971)rectangle++(0.04,0.04) (-0.462,3.665)rectangle++(0.04,0.04) (-0.248,4.73)rectangle++(0.04,0.04) (-0.024,5.626)rectangle++(0.04,0.04)
  (0.202,6.38)rectangle++(0.04,0.04) (0.425,7.011)rectangle++(0.04,0.04) (0.639,7.538)rectangle++(0.04,0.04) (0.843,7.978)rectangle++(0.04,0.04) (1.877,3.692)rectangle++(0.04,0.04)
  (1.741,4.659)rectangle++(0.04,0.04) (1.664,5.487)rectangle++(0.04,0.04) (1.632,6.194)rectangle++(0.04,0.04) (1.632,6.795)rectangle++(0.04,0.04) (1.657,7.307)rectangle++(0.04,0.04)
  (1.615,3.747)rectangle++(0.04,0.04) (0.739,3.052)rectangle++(0.04,0.04) (0.748,4.16)rectangle++(0.04,0.04) (0.8,5.103)rectangle++(0.04,0.04) (1.17,3.128)rectangle++(0.04,0.04)
  (1.117,4.208)rectangle++(0.04,0.04) (1.116,5.128)rectangle++(0.04,0.04) (1.151,5.911)rectangle++(0.04,0.04) (1.213,6.574)rectangle++(0.04,0.04) (1.292,7.136)rectangle++(0.04,0.04)
  (1.381,7.61)rectangle++(0.04,0.04) (1.476,8.009)rectangle++(0.04,0.04) (1.573,8.345)rectangle++(0.04,0.04) (1.669,8.627)rectangle++(0.04,0.04) (1.761,8.862)rectangle++(0.04,0.04)
  (1.849,9.059)rectangle++(0.04,0.04) (1.932,9.222)rectangle++(0.04,0.04) (2.009,9.358)rectangle++(0.04,0.04) (2.08,9.47)rectangle++(0.04,0.04) (2.144,9.562)rectangle++(0.04,0.04)
  (0.539,3.126)rectangle++(0.04,0.04) (0.581,4.232)rectangle++(0.04,0.04) (0.661,5.17)rectangle++(0.04,0.04) (-0.109,3.561)rectangle++(0.04,0.04) (0.047,4.627)rectangle++(0.04,0.04)
  (0.223,5.528)rectangle++(0.04,0.04) (0.409,6.286)rectangle++(0.04,0.04) (0.597,6.923)rectangle++(0.04,0.04) (0.782,7.457)rectangle++(0.04,0.04) (-1.399,4.566)rectangle++(0.04,0.04)
  (-1.009,5.533)rectangle++(0.04,0.04) (-0.638,6.34)rectangle++(0.04,0.04) (-0.291,7.011)rectangle++(0.04,0.04) (0.031,7.567)rectangle++(0.04,0.04) (0.327,8.027)rectangle++(0.04,0.04)
  (1.949,3.584)rectangle++(0.04,0.04) (1.798,4.565)rectangle++(0.04,0.04) (1.709,5.404)rectangle++(0.04,0.04) (1.666,6.122)rectangle++(0.04,0.04) (-0.957,4.469)rectangle++(0.04,0.04)
  (1.219,2.607)rectangle++(0.04,0.04) (1.138,3.763)rectangle++(0.04,0.04) (1.116,4.749)rectangle++(0.04,0.04) (1.136,5.589)rectangle++(0.04,0.04) (1.187,6.301)rectangle++(0.04,0.04)
  (1.259,6.905)rectangle++(0.04,0.04) (1.344,7.415)rectangle++(0.04,0.04) (1.679,3.697)rectangle++(0.04,0.04) (1.573,4.672)rectangle++(0.04,0.04) (1.521,5.504)rectangle++(0.04,0.04)
  (1.511,6.214)rectangle++(0.04,0.04) (1.531,6.818)rectangle++(0.04,0.04) (1.572,7.33)rectangle++(0.04,0.04) (1.627,7.764)rectangle++(0.04,0.04) (1.691,8.13)rectangle++(0.04,0.04)
  (1.761,8.439)rectangle++(0.04,0.04) (-1.432,4.944)rectangle++(0.04,0.04) (-1.022,5.856)rectangle++(0.04,0.04) (-0.636,6.614)rectangle++(0.04,0.04) (-0.279,7.244)rectangle++(0.04,0.04)
  (0.051,7.765)rectangle++(0.04,0.04) (0.352,8.194)rectangle++(0.04,0.04) (0.624,8.547)rectangle++(0.04,0.04) (0.87,8.836)rectangle++(0.04,0.04) (1.091,9.072)rectangle++(0.04,0.04)
  (1.288,9.264)rectangle++(0.04,0.04) (1.463,9.419)rectangle++(0.04,0.04) (-1.693,5.081)rectangle++(0.04,0.04) (-1.238,5.983)rectangle++(0.04,0.04) (-0.815,6.731)rectangle++(0.04,0.04)
  (-0.426,7.35)rectangle++(0.04,0.04) (-0.07,7.861)rectangle++(0.04,0.04) (0.253,8.281)rectangle++(0.04,0.04) (0.544,8.625)rectangle++(0.04,0.04) (0.805,8.906)rectangle++(0.04,0.04)
  (1.038,9.134)rectangle++(0.04,0.04) (1.246,9.318)rectangle++(0.04,0.04) (1.429,9.467)rectangle++(0.04,0.04) (1.591,9.586)rectangle++(0.04,0.04) (2.187,4.172)rectangle++(0.04,0.04)
  (2.024,5.055)rectangle++(0.04,0.04) (1.92,5.812)rectangle++(0.04,0.04) (1.863,6.459)rectangle++(0.04,0.04) (1.839,7.012)rectangle++(0.04,0.04) (1.842,7.483)rectangle++(0.04,0.04)
  (1.863,7.883)rectangle++(0.04,0.04) (1.896,8.222)rectangle++(0.04,0.04) (1.939,8.509)rectangle++(0.04,0.04) (1.88,4.043)rectangle++(0.04,0.04) (1.758,4.958)rectangle++(0.04,0.04)
  (1.69,5.74)rectangle++(0.04,0.04) (1.664,6.407)rectangle++(0.04,0.04) (1.668,6.976)rectangle++(0.04,0.04) (1.695,7.459)rectangle++(0.04,0.04) (1.737,7.869)rectangle++(0.04,0.04)
  (1.789,8.215)rectangle++(0.04,0.04) (1.847,8.507)rectangle++(0.04,0.04) (1.908,8.754)rectangle++(0.04,0.04) (1.97,8.96)rectangle++(0.04,0.04) (2.031,9.134)rectangle++(0.04,0.04)
  (2.089,9.279)rectangle++(0.04,0.04) (2.145,9.4)rectangle++(0.04,0.04) (2.076,4.261)rectangle++(0.04,0.04) (1.933,5.135)rectangle++(0.04,0.04) (-0.7,4.324)rectangle++(0.04,0.04)
  (1.153,2.637)rectangle++(0.04,0.04) (1.083,3.792)rectangle++(0.04,0.04) (1.07,4.776)rectangle++(0.04,0.04) (1.098,5.613)rectangle++(0.04,0.04) (1.156,6.323)rectangle++(0.04,0.04)
  (1.233,6.925)rectangle++(0.04,0.04) (1.323,7.433)rectangle++(0.04,0.04) (1.42,7.861)rectangle++(0.04,0.04) (1.519,8.221)rectangle++(0.04,0.04) (1.618,8.524)rectangle++(0.04,0.04)
  (1.714,8.777)rectangle++(0.04,0.04) (1.806,8.988)rectangle++(0.04,0.04) (1.892,9.164)rectangle++(0.04,0.04) (1.973,9.31)rectangle++(0.04,0.04) (2.047,9.43)rectangle++(0.04,0.04)
  (2.115,9.53)rectangle++(0.04,0.04) (2.177,9.612)rectangle++(0.04,0.04) (-1.606,5.253)rectangle++(0.04,0.04) (-1.157,6.126)rectangle++(0.04,0.04) (-0.741,6.849)rectangle++(0.04,0.04);
\fill[piece2]
  (-1.543,4.161)rectangle++(0.04,0.04) (-1.147,5.195)rectangle++(0.04,0.04) (-0.77,6.058)rectangle++(0.04,0.04) (-0.414,6.776)rectangle++(0.04,0.04) (-0.083,7.372)rectangle++(0.04,0.04)
  (0.222,7.866)rectangle++(0.04,0.04) (-1.586,4.542)rectangle++(0.04,0.04) (-1.169,5.52)rectangle++(0.04,0.04) (-0.775,6.335)rectangle++(0.04,0.04) (-0.407,7.012)rectangle++(0.04,0.04)
  (-0.068,7.573)rectangle++(0.04,0.04) (1.883,3.407)rectangle++(0.04,0.04) (1.735,4.417)rectangle++(0.04,0.04) (-0.583,4.145)rectangle++(0.04,0.04) (-0.332,5.143)rectangle++(0.04,0.04)
  (-0.028,3.535)rectangle++(0.04,0.04) (0.116,4.602)rectangle++(0.04,0.04) (0.28,5.503)rectangle++(0.04,0.04) (0.278,2.65)rectangle++(0.04,0.04) (0.34,3.838)rectangle++(0.04,0.04)
  (0.441,4.845)rectangle++(0.04,0.04) (0.566,5.697)rectangle++(0.04,0.04) (0.707,6.416)rectangle++(0.04,0.04) (1.506,3.36)rectangle++(0.04,0.04) (1.412,4.392)rectangle++(0.04,0.04)
  (1.374,5.273)rectangle++(0.04,0.04) (1.377,6.023)rectangle++(0.04,0.04) (1.409,6.661)rectangle++(0.04,0.04) (1.462,7.201)rectangle++(0.04,0.04) (1.528,7.659)rectangle++(0.04,0.04)
  (1.603,8.045)rectangle++(0.04,0.04) (1.682,8.371)rectangle++(0.04,0.04) (1.763,8.644)rectangle++(0.04,0.04) (1.842,8.873)rectangle++(0.04,0.04) (1.918,9.065)rectangle++(0.04,0.04)
  (1.991,9.224)rectangle++(0.04,0.04) (2.059,9.357)rectangle++(0.04,0.04) (2.122,9.467)rectangle++(0.04,0.04) (2.18,9.559)rectangle++(0.04,0.04) (2.233,9.634)rectangle++(0.04,0.04)
  (2.282,9.696)rectangle++(0.04,0.04) (2.325,9.746)rectangle++(0.04,0.04) (2.364,9.787)rectangle++(0.04,0.04) (2.399,9.821)rectangle++(0.04,0.04) (2.429,9.848)rectangle++(0.04,0.04)
  (2.457,9.87)rectangle++(0.04,0.04) (2.481,9.887)rectangle++(0.04,0.04) (2.502,9.901)rectangle++(0.04,0.04) (2.521,9.912)rectangle++(0.04,0.04) (2.537,9.921)rectangle++(0.04,0.04)
  (2.551,9.927)rectangle++(0.04,0.04) (2.16,4.455)rectangle++(0.04,0.04) (0.85,3.312)rectangle++(0.04,0.04) (0.721,2.794)rectangle++(0.04,0.04) (0.722,3.942)rectangle++(0.04,0.04)
  (0.769,4.918)rectangle++(0.04,0.04) (0.848,5.746)rectangle++(0.04,0.04) (0.949,6.446)rectangle++(0.04,0.04) (1.062,7.038)rectangle++(0.04,0.04) (1.182,7.536)rectangle++(0.04,0.04)
  (1.728,3.684)rectangle++(0.04,0.04) (1.614,4.658)rectangle++(0.04,0.04) (1.556,5.491)rectangle++(0.04,0.04) (1.54,6.201)rectangle++(0.04,0.04) (1.555,6.806)rectangle++(0.04,0.04)
  (1.592,7.319)rectangle++(0.04,0.04) (1.643,7.754)rectangle++(0.04,0.04) (-1.307,4.786)rectangle++(0.04,0.04) (-0.922,5.717)rectangle++(0.04,0.04) (-0.557,6.492)rectangle++(0.04,0.04)
  (-0.216,7.137)rectangle++(0.04,0.04) (0.099,7.671)rectangle++(0.04,0.04) (0.389,8.113)rectangle++(0.04,0.04) (0.653,8.477)rectangle++(0.04,0.04) (0.892,8.775)rectangle++(0.04,0.04)
  (-1.659,4.847)rectangle++(0.04,0.04) (1.395,2.52)rectangle++(0.04,0.04) (1.284,3.683)rectangle++(0.04,0.04) (1.237,4.675)rectangle++(0.04,0.04) (1.236,5.521)rectangle++(0.04,0.04)
  (1.269,6.239)rectangle++(0.04,0.04) (1.326,6.849)rectangle++(0.04,0.04) (1.399,7.365)rectangle++(0.04,0.04) (1.481,7.8)rectangle++(0.04,0.04) (1.569,8.167)rectangle++(0.04,0.04)
  (1.658,8.475)rectangle++(0.04,0.04) (1.746,8.734)rectangle++(0.04,0.04) (1.832,8.95)rectangle++(0.04,0.04) (1.913,9.131)rectangle++(0.04,0.04) (1.989,9.281)rectangle++(0.04,0.04)
  (2.059,9.405)rectangle++(0.04,0.04) (2.125,9.508)rectangle++(0.04,0.04) (2.184,9.593)rectangle++(0.04,0.04) (2.238,9.663)rectangle++(0.04,0.04) (2.287,9.72)rectangle++(0.04,0.04)
  (2.33,9.767)rectangle++(0.04,0.04) (2.369,9.805)rectangle++(0.04,0.04) (2.404,9.836)rectangle++(0.04,0.04) (2.434,9.86)rectangle++(0.04,0.04) (2.461,9.88)rectangle++(0.04,0.04)
  (2.485,9.896)rectangle++(0.04,0.04) (2.506,9.908)rectangle++(0.04,0.04) (2.524,9.918)rectangle++(0.04,0.04) (2.54,9.926)rectangle++(0.04,0.04) (2.554,9.931)rectangle++(0.04,0.04)
  (2.16,4.457)rectangle++(0.04,0.04) (0.851,3.312)rectangle++(0.04,0.04) (-0.507,3.761)rectangle++(0.04,0.04) (-0.283,4.813)rectangle++(0.04,0.04) (-0.12,3.093)rectangle++(0.04,0.04)
  (0.019,4.23)rectangle++(0.04,0.04) (1.045,2.782)rectangle++(0.04,0.04) (0.998,3.919)rectangle++(0.04,0.04) (1.002,4.888)rectangle++(0.04,0.04) (1.045,5.711)rectangle++(0.04,0.04)
  (1.115,6.408)rectangle++(0.04,0.04) (1.202,6.999)rectangle++(0.04,0.04) (1.299,7.497)rectangle++(0.04,0.04) (1.402,7.917)rectangle++(0.04,0.04) (-1.385,4.773)rectangle++(0.04,0.04)
  (-0.989,5.709)rectangle++(0.04,0.04) (-0.614,6.488)rectangle++(0.04,0.04) (-0.265,7.136)rectangle++(0.04,0.04) (0.058,7.672)rectangle++(0.04,0.04) (0.354,8.115)rectangle++(0.04,0.04)
  (0.623,8.48)rectangle++(0.04,0.04) (0.867,8.779)rectangle++(0.04,0.04) (1.086,9.024)rectangle++(0.04,0.04) (1.282,9.223)rectangle++(0.04,0.04) (1.456,9.385)rectangle++(0.04,0.04)
  (1.611,9.515)rectangle++(0.04,0.04) (1.748,9.619)rectangle++(0.04,0.04) (1.868,9.703)rectangle++(0.04,0.04) (1.974,9.769)rectangle++(0.04,0.04) (2.066,9.821)rectangle++(0.04,0.04)
  (2.147,9.861)rectangle++(0.04,0.04) (2.217,9.892)rectangle++(0.04,0.04) (2.278,9.916)rectangle++(0.04,0.04) (2.189,4.391)rectangle++(0.04,0.04) (2.034,5.241)rectangle++(0.04,0.04)
  (1.06,3.435)rectangle++(0.04,0.04) (1.036,4.473)rectangle++(0.04,0.04) (0.986,3.06)rectangle++(0.04,0.04) (0.959,4.158)rectangle++(0.04,0.04) (0.979,5.092)rectangle++(0.04,0.04)
  (1.034,5.885)rectangle++(0.04,0.04) (1.112,6.557)rectangle++(0.04,0.04) (1.205,7.125)rectangle++(0.04,0.04) (1.307,7.605)rectangle++(0.04,0.04) (1.413,8.008)rectangle++(0.04,0.04)
  (1.519,8.346)rectangle++(0.04,0.04) (1.623,8.63)rectangle++(0.04,0.04) (1.723,8.867)rectangle++(0.04,0.04) (1.817,9.064)rectangle++(0.04,0.04) (1.904,9.228)rectangle++(0.04,0.04)
  (1.986,9.364)rectangle++(0.04,0.04) (2.06,9.476)rectangle++(0.04,0.04) (2.128,9.568)rectangle++(0.04,0.04) (2.189,9.644)rectangle++(0.04,0.04) (-1.611,5.262)rectangle++(0.04,0.04)
  (-1.161,6.133)rectangle++(0.04,0.04) (-0.744,6.856)rectangle++(0.04,0.04) (-0.36,7.453)rectangle++(0.04,0.04) (-0.01,7.946)rectangle++(0.04,0.04) (0.307,8.351)rectangle++(0.04,0.04)
  (0.593,8.682)rectangle++(0.04,0.04) (0.849,8.952)rectangle++(0.04,0.04) (1.078,9.172)rectangle++(0.04,0.04) (1.281,9.349)rectangle++(0.04,0.04) (1.46,9.492)rectangle++(0.04,0.04)
  (1.619,9.606)rectangle++(0.04,0.04) (2.189,4.182)rectangle++(0.04,0.04) (2.026,5.063)rectangle++(0.04,0.04) (1.922,5.819)rectangle++(0.04,0.04) (1.864,6.465)rectangle++(0.04,0.04)
  (1.841,7.017)rectangle++(0.04,0.04) (1.843,7.487)rectangle++(0.04,0.04) (1.864,7.887)rectangle++(0.04,0.04) (-1.296,4.854)rectangle++(0.04,0.04) (-0.91,5.774)rectangle++(0.04,0.04)
  (-0.545,6.54)rectangle++(0.04,0.04) (-0.203,7.177)rectangle++(0.04,0.04) (-0.006,1.868)rectangle++(0.04,0.04) (0.067,3.184)rectangle++(0.04,0.04) (0.182,4.3)rectangle++(0.04,0.04)
  (0.325,5.244)rectangle++(0.04,0.04) (0.484,6.041)rectangle++(0.04,0.04) (0.65,6.712)rectangle++(0.04,0.04) (0.819,7.275)rectangle++(0.04,0.04) (0.985,7.747)rectangle++(0.04,0.04)
  (1.145,8.142)rectangle++(0.04,0.04) (1.297,8.471)rectangle++(0.04,0.04) (1.439,8.745)rectangle++(0.04,0.04) (-1.554,4.944)rectangle++(0.04,0.04) (-1.126,5.861)rectangle++(0.04,0.04)
  (-0.725,6.623)rectangle++(0.04,0.04) (-0.353,7.254)rectangle++(0.04,0.04) (-0.012,7.777)rectangle++(0.04,0.04) (0.298,8.207)rectangle++(0.04,0.04) (0.58,8.56)rectangle++(0.04,0.04)
  (-1.669,4.746)rectangle++(0.04,0.04) (-1.231,5.697)rectangle++(0.04,0.04) (-0.821,6.488)rectangle++(0.04,0.04) (-0.44,7.144)rectangle++(0.04,0.04) (-0.091,7.686)rectangle++(0.04,0.04)
  (1.914,3.424)rectangle++(0.04,0.04) (1.762,4.43)rectangle++(0.04,0.04) (1.672,5.291)rectangle++(0.04,0.04) (1.631,6.027)rectangle++(0.04,0.04) (0.019,2.402)rectangle++(0.04,0.04)
  (0.11,3.637)rectangle++(0.04,0.04) (0.237,4.684)rectangle++(0.04,0.04) (0.386,5.568)rectangle++(0.04,0.04) (0.549,6.313)rectangle++(0.04,0.04) (1.499,3.304)rectangle++(0.04,0.04)
  (1.404,4.345)rectangle++(0.04,0.04) (1.365,5.233)rectangle++(0.04,0.04) (1.368,5.99)rectangle++(0.04,0.04) (1.4,6.633)rectangle++(0.04,0.04) (-1.113,4.519)rectangle++(0.04,0.04)
  (-0.768,5.482)rectangle++(0.04,0.04) (-0.435,6.287)rectangle++(0.04,0.04) (-1.369,4.107)rectangle++(0.04,0.04) (-1.001,5.142)rectangle++(0.04,0.04) (-0.648,6.007)rectangle++(0.04,0.04)
  (-0.312,6.728)rectangle++(0.04,0.04) (0.001,7.327)rectangle++(0.04,0.04) (0.292,7.825)rectangle++(0.04,0.04) (0.559,8.235)rectangle++(0.04,0.04) (0.802,8.574)rectangle++(0.04,0.04)
  (-1.632,4.791)rectangle++(0.04,0.04) (-1.266,3.587)rectangle++(0.04,0.04) (-0.935,4.696)rectangle++(0.04,0.04) (-0.609,5.625)rectangle++(0.04,0.04) (-0.295,6.402)rectangle++(0.04,0.04)
  (0.003,7.05)rectangle++(0.04,0.04) (0.283,7.588)rectangle++(0.04,0.04) (0.542,8.035)rectangle++(0.04,0.04) (0.78,8.404)rectangle++(0.04,0.04) (0.997,8.709)rectangle++(0.04,0.04)
  (1.193,8.959)rectangle++(0.04,0.04) (1.37,9.163)rectangle++(0.04,0.04) (1.529,9.33)rectangle++(0.04,0.04) (1.671,9.466)rectangle++(0.04,0.04) (1.797,9.575)rectangle++(0.04,0.04)
  (1.908,9.663)rectangle++(0.04,0.04) (2.169,4.259)rectangle++(0.04,0.04) (-0.472,4.195)rectangle++(0.04,0.04) (-0.235,5.181)rectangle++(0.04,0.04) (0.005,6.009)rectangle++(0.04,0.04)
  (0.243,6.704)rectangle++(0.04,0.04) (0.472,7.285)rectangle++(0.04,0.04) (1.749,3.474)rectangle++(0.04,0.04) (1.623,4.479)rectangle++(0.04,0.04) (1.557,5.339)rectangle++(0.04,0.04)
  (1.534,6.072)rectangle++(0.04,0.04) (1.545,6.696)rectangle++(0.04,0.04) (1.579,7.226)rectangle++(0.04,0.04) (1.629,7.675)rectangle++(0.04,0.04) (1.689,8.055)rectangle++(0.04,0.04)
  (1.756,8.375)rectangle++(0.04,0.04) (1.825,8.645)rectangle++(0.04,0.04) (0.035,2.758)rectangle++(0.04,0.04) (0.138,3.939)rectangle++(0.04,0.04) (0.273,4.939)rectangle++(0.04,0.04)
  (0.427,5.784)rectangle++(0.04,0.04) (0.592,6.495)rectangle++(0.04,0.04) (0.761,7.093)rectangle++(0.04,0.04) (0.928,7.595)rectangle++(0.04,0.04) (1.091,8.015)rectangle++(0.04,0.04)
  (1.246,8.365)rectangle++(0.04,0.04) (1.391,8.657)rectangle++(0.04,0.04) (1.527,8.899)rectangle++(0.04,0.04) (1.651,9.099)rectangle++(0.04,0.04) (1.765,9.264)rectangle++(0.04,0.04)
  (1.869,9.4)rectangle++(0.04,0.04) (1.962,9.512)rectangle++(0.04,0.04) (2.046,9.603)rectangle++(0.04,0.04) (2.121,9.677)rectangle++(0.04,0.04) (2.188,9.736)rectangle++(0.04,0.04)
  (2.247,9.785)rectangle++(0.04,0.04) (2.299,9.823)rectangle++(0.04,0.04) (2.163,4.378)rectangle++(0.04,0.04) (2.012,5.231)rectangle++(0.04,0.04) (1.917,5.962)rectangle++(0.04,0.04)
  (1.866,6.587)rectangle++(0.04,0.04) (1.847,7.121)rectangle++(0.04,0.04) (1.853,7.575)rectangle++(0.04,0.04) (1.876,7.961)rectangle++(0.04,0.04) (1.911,8.288)rectangle++(0.04,0.04)
  (-1.373,4.942)rectangle++(0.04,0.04) (-0.972,5.852)rectangle++(0.04,0.04) (1.563,2.87)rectangle++(0.04,0.04) (1.441,3.973)rectangle++(0.04,0.04) (1.382,4.916)rectangle++(0.04,0.04)
  (0.012,2.251)rectangle++(0.04,0.04) (0.098,3.509)rectangle++(0.04,0.04) (0.221,4.575)rectangle++(0.04,0.04) (0.369,5.476)rectangle++(0.04,0.04) (0.53,6.236)rectangle++(0.04,0.04)
  (0.698,6.876)rectangle++(0.04,0.04) (0.866,7.413)rectangle++(0.04,0.04) (1.031,7.862)rectangle++(0.04,0.04) (1.188,8.238)rectangle++(0.04,0.04) (1.337,8.551)rectangle++(0.04,0.04)
  (1.477,8.811)rectangle++(0.04,0.04) (1.605,9.027)rectangle++(0.04,0.04) (1.723,9.205)rectangle++(0.04,0.04) (1.831,9.351)rectangle++(0.04,0.04) (1.928,9.471)rectangle++(0.04,0.04)
  (2.016,9.57)rectangle++(0.04,0.04) (2.094,9.65)rectangle++(0.04,0.04) (2.164,9.715)rectangle++(0.04,0.04) (2.225,9.767)rectangle++(0.04,0.04) (2.28,9.809)rectangle++(0.04,0.04)
  (2.328,9.843)rectangle++(0.04,0.04) (2.371,9.87)rectangle++(0.04,0.04) (2.408,9.891)rectangle++(0.04,0.04) (2.44,9.907)rectangle++(0.04,0.04) (-1.253,3.407)rectangle++(0.04,0.04)
  (-0.931,4.542)rectangle++(0.04,0.04) (-0.612,5.494)rectangle++(0.04,0.04) (-0.303,6.291)rectangle++(0.04,0.04) (-0.008,6.956)rectangle++(0.04,0.04) (0.269,7.509)rectangle++(0.04,0.04)
  (0.527,7.968)rectangle++(0.04,0.04) (0.765,8.348)rectangle++(0.04,0.04) (0.982,8.661)rectangle++(0.04,0.04) (1.179,8.919)rectangle++(0.04,0.04) (1.356,9.13)rectangle++(0.04,0.04)
  (1.516,9.303)rectangle++(0.04,0.04) (1.658,9.443)rectangle++(0.04,0.04) (1.785,9.556)rectangle++(0.04,0.04) (1.897,9.648)rectangle++(0.04,0.04) (1.997,9.721)rectangle++(0.04,0.04)
  (2.084,9.779)rectangle++(0.04,0.04) (2.16,9.825)rectangle++(0.04,0.04) (2.227,9.861)rectangle++(0.04,0.04) (2.285,9.889)rectangle++(0.04,0.04) (2.336,9.91)rectangle++(0.04,0.04)
  (2.38,9.927)rectangle++(0.04,0.04) (2.417,9.939)rectangle++(0.04,0.04) (2.45,9.947)rectangle++(0.04,0.04) (2.478,9.953)rectangle++(0.04,0.04) (2.503,9.958)rectangle++(0.04,0.04)
  (2.523,9.96)rectangle++(0.04,0.04) (-1.618,5.387)rectangle++(0.04,0.04) (-1.162,6.24)rectangle++(0.04,0.04) (-0.74,6.947)rectangle++(0.04,0.04) (-0.354,7.53)rectangle++(0.04,0.04)
  (1.906,3.334)rectangle++(0.04,0.04) (1.751,4.354)rectangle++(0.04,0.04) (1.661,5.227)rectangle++(0.04,0.04) (1.618,5.973)rectangle++(0.04,0.04) (1.612,6.608)rectangle++(0.04,0.04)
  (1.633,7.149)rectangle++(0.04,0.04) (-1.181,4.665)rectangle++(0.04,0.04) (-0.819,5.609)rectangle++(0.04,0.04) (-0.474,6.396)rectangle++(0.04,0.04) (-0.15,7.052)rectangle++(0.04,0.04)
  (0.153,7.596)rectangle++(0.04,0.04) (0.431,8.047)rectangle++(0.04,0.04) (0.686,8.419)rectangle++(0.04,0.04) (0.918,8.725)rectangle++(0.04,0.04) (1.127,8.976)rectangle++(0.04,0.04)
  (1.315,9.18)rectangle++(0.04,0.04) (1.483,9.347)rectangle++(0.04,0.04) (1.632,9.482)rectangle++(0.04,0.04) (2.157,4.161)rectangle++(0.04,0.04) (0.778,3.254)rectangle++(0.04,0.04)
  (0.789,4.331)rectangle++(0.04,0.04) (0.842,5.246)rectangle++(0.04,0.04) (0.923,6.022)rectangle++(0.04,0.04) (1.023,6.678)rectangle++(0.04,0.04) (1.135,7.231)rectangle++(0.04,0.04)
  (1.252,7.697)rectangle++(0.04,0.04) (1.37,8.089)rectangle++(0.04,0.04) (1.486,8.417)rectangle++(0.04,0.04) (-1.477,4.888)rectangle++(0.04,0.04) (-1.062,5.81)rectangle++(0.04,0.04)
  (-0.672,6.577)rectangle++(0.04,0.04) (-0.311,7.214)rectangle++(0.04,0.04) (0.022,7.74)rectangle++(0.04,0.04) (0.326,8.174)rectangle++(0.04,0.04) (0.602,8.531)rectangle++(0.04,0.04)
  (0.851,8.824)rectangle++(0.04,0.04) (1.074,9.062)rectangle++(0.04,0.04) (1.273,9.256)rectangle++(0.04,0.04) (1.45,9.413)rectangle++(0.04,0.04) (1.607,9.539)rectangle++(0.04,0.04)
  (1.745,9.64)rectangle++(0.04,0.04) (1.867,9.721)rectangle++(0.04,0.04) (1.974,9.784)rectangle++(0.04,0.04) (2.067,9.834)rectangle++(0.04,0.04) (-1.673,5.276)rectangle++(0.04,0.04)
  (-1.213,6.148)rectangle++(0.04,0.04) (-0.788,6.871)rectangle++(0.04,0.04) (-0.397,7.468)rectangle++(0.04,0.04) (-0.041,7.96)rectangle++(0.04,0.04) (0.282,8.363)rectangle++(0.04,0.04)
  (0.572,8.694)rectangle++(0.04,0.04) (2.086,3.768)rectangle++(0.04,0.04) (1.922,4.716)rectangle++(0.04,0.04) (1.82,5.528)rectangle++(0.04,0.04) (1.766,6.222)rectangle++(0.04,0.04)
  (1.748,6.814)rectangle++(0.04,0.04) (1.483,3.673)rectangle++(0.04,0.04) (1.405,4.659)rectangle++(0.04,0.04) (1.378,5.5)rectangle++(0.04,0.04) (1.39,6.216)rectangle++(0.04,0.04)
  (1.428,6.824)rectangle++(0.04,0.04) (1.484,7.34)rectangle++(0.04,0.04) (1.553,7.776)rectangle++(0.04,0.04) (1.629,8.143)rectangle++(0.04,0.04) (-1.393,4.861)rectangle++(0.04,0.04)
  (-0.992,5.783)rectangle++(0.04,0.04) (-0.614,6.552)rectangle++(0.04,0.04) (-0.262,7.19)rectangle++(0.04,0.04) (-1.529,4.318)rectangle++(0.04,0.04) (1.248,2.448)rectangle++(0.04,0.04)
  (1.157,3.627)rectangle++(0.04,0.04) (1.126,4.633)rectangle++(0.04,0.04) (1.14,5.489)rectangle++(0.04,0.04) (1.187,6.216)rectangle++(0.04,0.04) (1.255,6.833)rectangle++(0.04,0.04)
  (1.338,7.354)rectangle++(0.04,0.04) (1.429,7.793)rectangle++(0.04,0.04) (1.524,8.163)rectangle++(0.04,0.04) (1.62,8.474)rectangle++(0.04,0.04) (1.714,8.734)rectangle++(0.04,0.04)
  (1.804,8.952)rectangle++(0.04,0.04) (1.889,9.133)rectangle++(0.04,0.04) (1.969,9.284)rectangle++(0.04,0.04) (2.043,9.409)rectangle++(0.04,0.04) (2.11,9.512)rectangle++(0.04,0.04)
  (2.172,9.597)rectangle++(0.04,0.04) (2.228,9.667)rectangle++(0.04,0.04) (2.133,4.332)rectangle++(0.04,0.04) (1.984,5.193)rectangle++(0.04,0.04) (1.892,5.931)rectangle++(0.04,0.04)
  (1.843,6.562)rectangle++(0.04,0.04) (1.827,7.1)rectangle++(0.04,0.04) (1.835,7.558)rectangle++(0.04,0.04) (1.86,7.947)rectangle++(0.04,0.04) (1.896,8.277)rectangle++(0.04,0.04)
  (1.941,8.556)rectangle++(0.04,0.04) (1.99,8.791)rectangle++(0.04,0.04) (2.041,8.989)rectangle++(0.04,0.04) (2.092,9.155)rectangle++(0.04,0.04) (2.142,9.294)rectangle++(0.04,0.04)
  (2.19,9.411)rectangle++(0.04,0.04) (2.236,9.508)rectangle++(0.04,0.04) (2.279,9.588)rectangle++(0.04,0.04) (2.318,9.655)rectangle++(0.04,0.04) (2.355,9.71)rectangle++(0.04,0.04)
  (2.388,9.756)rectangle++(0.04,0.04) (2.417,9.793)rectangle++(0.04,0.04) (2.444,9.824)rectangle++(0.04,0.04) (2.468,9.849)rectangle++(0.04,0.04) (2.15,4.422)rectangle++(0.04,0.04)
  (-0.51,4.224)rectangle++(0.04,0.04) (-0.944,3.686)rectangle++(0.04,0.04) (-0.657,4.767)rectangle++(0.04,0.04) (-0.37,5.674)rectangle++(0.04,0.04) (-0.09,6.434)rectangle++(0.04,0.04)
  (0.179,7.069)rectangle++(0.04,0.04) (0.432,7.598)rectangle++(0.04,0.04) (0.669,8.037)rectangle++(0.04,0.04) (0.888,8.401)rectangle++(0.04,0.04) (1.089,8.701)rectangle++(0.04,0.04)
  (-1.608,4.869)rectangle++(0.04,0.04) (-1.174,5.8)rectangle++(0.04,0.04) (-0.768,6.573)rectangle++(0.04,0.04) (-0.392,7.214)rectangle++(0.04,0.04) (-0.047,7.744)rectangle++(0.04,0.04)
  (0.267,8.18)rectangle++(0.04,0.04) (0.552,8.539)rectangle++(0.04,0.04) (0.809,8.832)rectangle++(0.04,0.04) (1.039,9.071)rectangle++(0.04,0.04) (1.243,9.265)rectangle++(0.04,0.04)
  (1.425,9.422)rectangle++(0.04,0.04) (1.586,9.548)rectangle++(0.04,0.04) (1.728,9.648)rectangle++(0.04,0.04) (1.853,9.728)rectangle++(0.04,0.04) (2.192,4.258)rectangle++(0.04,0.04)
  (0.798,3.281)rectangle++(0.04,0.04) (0.807,4.353)rectangle++(0.04,0.04) (0.858,5.264)rectangle++(0.04,0.04) (0.938,6.036)rectangle++(0.04,0.04) (1.036,6.69)rectangle++(0.04,0.04)
  (1.146,7.241)rectangle++(0.04,0.04) (1.262,7.705)rectangle++(0.04,0.04) (1.379,8.095)rectangle++(0.04,0.04) (1.493,8.422)rectangle++(0.04,0.04) (1.604,8.695)rectangle++(0.04,0.04)
  (1.965,4.003)rectangle++(0.04,0.04) (1.828,4.92)rectangle++(0.04,0.04) (1.749,5.705)rectangle++(0.04,0.04) (1.712,6.376)rectangle++(0.04,0.04) (1.708,6.947)rectangle++(0.04,0.04)
  (1.728,7.433)rectangle++(0.04,0.04) (1.764,7.845)rectangle++(0.04,0.04) (1.811,8.194)rectangle++(0.04,0.04) (1.865,8.489)rectangle++(0.04,0.04) (1.922,8.737)rectangle++(0.04,0.04)
  (1.938,4.083)rectangle++(0.04,0.04) (1.809,4.989)rectangle++(0.04,0.04) (1.735,5.765)rectangle++(0.04,0.04) (1.703,6.427)rectangle++(0.04,0.04) (1.702,6.991)rectangle++(0.04,0.04)
  (1.724,7.47)rectangle++(0.04,0.04) (1.762,7.877)rectangle++(0.04,0.04) (-1.312,4.833)rectangle++(0.04,0.04) (-0.924,5.757)rectangle++(0.04,0.04) (-0.558,6.526)rectangle++(0.04,0.04)
  (-0.215,7.166)rectangle++(0.04,0.04) (-1.516,4.322)rectangle++(0.04,0.04) (-1.118,5.33)rectangle++(0.04,0.04) (-0.739,6.172)rectangle++(0.04,0.04) (-0.384,6.872)rectangle++(0.04,0.04)
  (-0.053,7.453)rectangle++(0.04,0.04) (-1.548,4.409)rectangle++(0.04,0.04) (-1.141,5.406)rectangle++(0.04,0.04) (-0.756,6.237)rectangle++(0.04,0.04) (-0.395,6.928)rectangle++(0.04,0.04)
  (-1.498,4.241)rectangle++(0.04,0.04) (-1.105,5.261)rectangle++(0.04,0.04) (-1.271,3.779)rectangle++(0.04,0.04) (-0.932,4.859)rectangle++(0.04,0.04) (-0.6,5.764)rectangle++(0.04,0.04)
  (-0.281,6.519)rectangle++(0.04,0.04) (0.019,7.149)rectangle++(0.04,0.04) (0.3,7.672)rectangle++(0.04,0.04) (-0.217,3.183)rectangle++(0.04,0.04) (-0.059,4.31)rectangle++(0.04,0.04)
  (0.12,5.262)rectangle++(0.04,0.04) (0.31,6.064)rectangle++(0.04,0.04) (0.504,6.739)rectangle++(0.04,0.04) (0.696,7.304)rectangle++(0.04,0.04) (0.881,7.777)rectangle++(0.04,0.04)
  (1.058,8.171)rectangle++(0.04,0.04) (1.224,8.499)rectangle++(0.04,0.04) (1.378,8.772)rectangle++(0.04,0.04) (1.52,8.997)rectangle++(0.04,0.04) (1.65,9.183)rectangle++(0.04,0.04)
  (1.767,9.336)rectangle++(0.04,0.04) (1.873,9.461)rectangle++(0.04,0.04) (1.969,9.563)rectangle++(0.04,0.04) (2.054,9.646)rectangle++(0.04,0.04) (-1.636,5.237)rectangle++(0.04,0.04)
  (-1.183,6.113)rectangle++(0.04,0.04) (-0.763,6.84)rectangle++(0.04,0.04) (-1.545,4.155)rectangle++(0.04,0.04) (-1.115,3.478)rectangle++(0.04,0.04) (-0.811,4.597)rectangle++(0.04,0.04)
  (1.233,2.665)rectangle++(0.04,0.04) (1.153,3.812)rectangle++(0.04,0.04) (1.13,4.79)rectangle++(0.04,0.04) (1.15,5.623)rectangle++(0.04,0.04) (1.2,6.329)rectangle++(0.04,0.04)
  (-1.085,4.422)rectangle++(0.04,0.04) (-0.747,5.398)rectangle++(0.04,0.04) (-0.421,6.215)rectangle++(0.04,0.04) (-0.112,6.896)rectangle++(0.04,0.04) (0.179,7.462)rectangle++(0.04,0.04)
  (0.448,7.932)rectangle++(0.04,0.04) (0.696,8.32)rectangle++(0.04,0.04) (0.922,8.641)rectangle++(0.04,0.04) (1.127,8.904)rectangle++(0.04,0.04) (1.312,9.119)rectangle++(0.04,0.04)
  (1.478,9.295)rectangle++(0.04,0.04) (1.626,9.438)rectangle++(0.04,0.04) (1.757,9.553)rectangle++(0.04,0.04) (1.874,9.646)rectangle++(0.04,0.04) (1.976,9.721)rectangle++(0.04,0.04)
  (2.066,9.78)rectangle++(0.04,0.04) (-1.662,5.266)rectangle++(0.04,0.04) (-1.204,6.138)rectangle++(0.04,0.04) (-0.78,6.862)rectangle++(0.04,0.04) (-0.391,7.46)rectangle++(0.04,0.04)
  (-0.036,7.953)rectangle++(0.04,0.04) (1.974,3.488)rectangle++(0.04,0.04) (1.815,4.482)rectangle++(0.04,0.04) (1.72,5.333)rectangle++(0.04,0.04) (1.673,6.061)rectangle++(0.04,0.04)
  (1.662,6.681)rectangle++(0.04,0.04) (1.678,7.208)rectangle++(0.04,0.04) (1.712,7.656)rectangle++(0.04,0.04) (1.76,8.035)rectangle++(0.04,0.04) (1.815,8.356)rectangle++(0.04,0.04)
  (1.875,8.626)rectangle++(0.04,0.04) (1.936,8.853)rectangle++(0.04,0.04) (1.998,9.044)rectangle++(0.04,0.04) (2.058,9.204)rectangle++(0.04,0.04) (-1.541,5.15)rectangle++(0.04,0.04)
  (-1.106,6.036)rectangle++(0.04,0.04) (-0.701,6.771)rectangle++(0.04,0.04) (-0.327,7.379)rectangle++(0.04,0.04) (0.015,7.882)rectangle++(0.04,0.04) (0.326,8.295)rectangle++(0.04,0.04)
  (-1.659,4.646)rectangle++(0.04,0.04) (-1.226,5.612)rectangle++(0.04,0.04) (-0.82,6.415)rectangle++(0.04,0.04) (-0.442,7.082)rectangle++(0.04,0.04) (-0.095,7.634)rectangle++(0.04,0.04)
  (0.222,8.089)rectangle++(0.04,0.04) (0.51,8.463)rectangle++(0.04,0.04) (0.364,2.363)rectangle++(0.04,0.04) (0.402,3.59)rectangle++(0.04,0.04) (0.483,4.632)rectangle++(0.04,0.04)
  (0.594,5.514)rectangle++(0.04,0.04) (0.723,6.259)rectangle++(0.04,0.04) (0.863,6.888)rectangle++(0.04,0.04) (1.007,7.416)rectangle++(0.04,0.04) (-1.35,4.6)rectangle++(0.04,0.04)
  (-0.966,5.561)rectangle++(0.04,0.04) (-1.31,3.864)rectangle++(0.04,0.04) (-0.961,4.933)rectangle++(0.04,0.04) (-0.622,5.828)rectangle++(0.04,0.04) (-0.297,6.575)rectangle++(0.04,0.04)
  (0.008,7.197)rectangle++(0.04,0.04) (0.292,7.713)rectangle++(0.04,0.04) (0.555,8.141)rectangle++(0.04,0.04) (0.795,8.493)rectangle++(0.04,0.04) (1.013,8.784)rectangle++(0.04,0.04)
  (1.211,9.022)rectangle++(0.04,0.04) (1.388,9.216)rectangle++(0.04,0.04) (1.546,9.375)rectangle++(0.04,0.04) (2.14,4.121)rectangle++(0.04,0.04) (1.982,5.013)rectangle++(0.04,0.04)
  (1.883,5.778)rectangle++(0.04,0.04) (1.829,6.432)rectangle++(0.04,0.04) (1.81,6.99)rectangle++(0.04,0.04) (1.816,7.466)rectangle++(0.04,0.04) (1.84,7.869)rectangle++(0.04,0.04)
  (1.877,8.212)rectangle++(0.04,0.04) (1.814,3.972)rectangle++(0.04,0.04) (1.698,4.9)rectangle++(0.04,0.04) (1.637,5.693)rectangle++(0.04,0.04) (1.617,6.37)rectangle++(0.04,0.04)
  (1.627,6.946)rectangle++(0.04,0.04) (-1.139,4.624)rectangle++(0.04,0.04) (-0.785,5.572)rectangle++(0.04,0.04) (-0.447,6.364)rectangle++(0.04,0.04) (-0.127,7.023)rectangle++(0.04,0.04)
  (0.988,2.136)rectangle++(0.04,0.04) (0.923,3.372)rectangle++(0.04,0.04) (0.917,4.426)rectangle++(0.04,0.04) (0.955,5.321)rectangle++(0.04,0.04) (1.022,6.081)rectangle++(0.04,0.04)
  (1.11,6.724)rectangle++(0.04,0.04) (1.21,7.267)rectangle++(0.04,0.04) (1.317,7.725)rectangle++(0.04,0.04) (1.426,8.11)rectangle++(0.04,0.04) (1.535,8.433)rectangle++(0.04,0.04)
  (1.64,8.702)rectangle++(0.04,0.04) (1.739,8.928)rectangle++(0.04,0.04) (1.833,9.115)rectangle++(0.04,0.04) (1.921,9.271)rectangle++(0.04,0.04) (2.071,4.186)rectangle++(0.04,0.04)
  (1.925,5.071)rectangle++(0.04,0.04) (1.837,5.83)rectangle++(0.04,0.04) (1.793,6.478)rectangle++(0.04,0.04) (1.781,7.031)rectangle++(0.04,0.04) (-1.13,4.669)rectangle++(0.04,0.04)
  (-0.776,5.61)rectangle++(0.04,0.04) (1.48,2.868)rectangle++(0.04,0.04) (1.371,3.975)rectangle++(0.04,0.04) (1.322,4.92)rectangle++(0.04,0.04) (1.318,5.725)rectangle++(0.04,0.04)
  (1.347,6.41)rectangle++(0.04,0.04) (1.399,6.991)rectangle++(0.04,0.04) (1.467,7.482)rectangle++(0.04,0.04) (1.544,7.898)rectangle++(0.04,0.04) (1.626,8.247)rectangle++(0.04,0.04)
  (1.71,8.541)rectangle++(0.04,0.04) (1.793,8.788)rectangle++(0.04,0.04) (1.873,8.994)rectangle++(0.04,0.04) (-1.53,5.074)rectangle++(0.04,0.04) (-1.099,5.971)rectangle++(0.04,0.04)
  (1.607,2.864)rectangle++(0.04,0.04) (1.479,3.966)rectangle++(0.04,0.04) (0.809,3.063)rectangle++(0.04,0.04) (0.808,4.167)rectangle++(0.04,0.04) (0.851,5.106)rectangle++(0.04,0.04)
  (0.926,5.902)rectangle++(0.04,0.04) (1.021,6.576)rectangle++(0.04,0.04) (1.128,7.145)rectangle++(0.04,0.04) (1.243,7.624)rectangle++(0.04,0.04) (1.359,8.027)rectangle++(0.04,0.04)
  (1.474,8.365)rectangle++(0.04,0.04) (1.585,8.647)rectangle++(0.04,0.04) (1.691,8.883)rectangle++(0.04,0.04) (1.791,9.079)rectangle++(0.04,0.04) (1.883,9.242)rectangle++(0.04,0.04)
  (1.968,9.376)rectangle++(0.04,0.04) (2.046,9.487)rectangle++(0.04,0.04) (2.116,9.579)rectangle++(0.04,0.04) (-0.246,3.411)rectangle++(0.04,0.04) (0.874,2.564)rectangle++(0.04,0.04)
  (0.843,3.741)rectangle++(0.04,0.04) (0.828,2.875)rectangle++(0.04,0.04) (0.816,4.007)rectangle++(0.04,0.04) (0.852,4.969)rectangle++(0.04,0.04) (0.921,5.786)rectangle++(0.04,0.04)
  (1.012,6.478)rectangle++(0.04,0.04) (1.117,7.062)rectangle++(0.04,0.04) (1.23,7.554)rectangle++(0.04,0.04) (1.079,2.249)rectangle++(0.04,0.04) (1.005,3.465)rectangle++(0.04,0.04)
  (0.99,4.501)rectangle++(0.04,0.04) (1.02,5.382)rectangle++(0.04,0.04) (1.08,6.13)rectangle++(0.04,0.04) (1.161,6.764)rectangle++(0.04,0.04) (1.255,7.299)rectangle++(0.04,0.04)
  (1.357,7.75)rectangle++(0.04,0.04) (1.461,8.13)rectangle++(0.04,0.04) (1.565,8.448)rectangle++(0.04,0.04) (1.666,8.714)rectangle++(0.04,0.04) (1.762,8.937)rectangle++(0.04,0.04)
  (1.853,9.122)rectangle++(0.04,0.04) (-1.56,5.096)rectangle++(0.04,0.04) (-1.125,5.99)rectangle++(0.04,0.04) (-0.718,6.733)rectangle++(0.04,0.04) (-0.344,7.348)rectangle++(0.04,0.04)
  (0,7.856)rectangle++(0.04,0.04) (0.312,8.273)rectangle++(0.04,0.04) (0.594,8.616)rectangle++(0.04,0.04) (-1.679,4.76)rectangle++(0.04,0.04) (-1.239,5.709)rectangle++(0.04,0.04)
  (-0.827,6.499)rectangle++(0.04,0.04) (-0.445,7.153)rectangle++(0.04,0.04) (-0.094,7.694)rectangle++(0.04,0.04) (0.225,8.14)rectangle++(0.04,0.04) (0.515,8.506)rectangle++(0.04,0.04)
  (-1.669,4.723)rectangle++(0.04,0.04) (-1.232,5.678)rectangle++(0.04,0.04) (-0.822,6.472)rectangle++(0.04,0.04) (-0.442,7.13)rectangle++(0.04,0.04) (-0.093,7.674)rectangle++(0.04,0.04)
  (0.226,8.123)rectangle++(0.04,0.04) (0.515,8.492)rectangle++(0.04,0.04) (0.775,8.794)rectangle++(0.04,0.04) (1.008,9.04)rectangle++(0.04,0.04) (1.216,9.24)rectangle++(0.04,0.04)
  (1.401,9.401)rectangle++(0.04,0.04) (1.565,9.531)rectangle++(0.04,0.04) (2.177,4.155)rectangle++(0.04,0.04) (2.014,5.041)rectangle++(0.04,0.04) (1.912,5.801)rectangle++(0.04,0.04)
  (1.213,3.515)rectangle++(0.04,0.04) (1.169,4.535)rectangle++(0.04,0.04) (1.173,5.404)rectangle++(0.04,0.04) (1.211,6.143)rectangle++(0.04,0.04) (1.273,6.769)rectangle++(0.04,0.04)
  (1.35,7.299)rectangle++(0.04,0.04) (1.438,7.746)rectangle++(0.04,0.04) (1.53,8.123)rectangle++(0.04,0.04) (1.623,8.44)rectangle++(0.04,0.04) (1.715,8.705)rectangle++(0.04,0.04)
  (1.804,8.927)rectangle++(0.04,0.04) (1.888,9.112)rectangle++(0.04,0.04) (1.967,9.266)rectangle++(0.04,0.04) (2.04,9.393)rectangle++(0.04,0.04) (2.108,9.499)rectangle++(0.04,0.04)
  (2.169,9.586)rectangle++(0.04,0.04) (2.225,9.658)rectangle++(0.04,0.04) (2.276,9.716)rectangle++(0.04,0.04) (2.321,9.764)rectangle++(0.04,0.04) (2.361,9.803)rectangle++(0.04,0.04)
  (2.397,9.834)rectangle++(0.04,0.04) (2.428,9.859)rectangle++(0.04,0.04) (2.456,9.879)rectangle++(0.04,0.04) (2.481,9.895)rectangle++(0.04,0.04) (2.502,9.908)rectangle++(0.04,0.04)
  (2.521,9.918)rectangle++(0.04,0.04) (-1.609,5.378)rectangle++(0.04,0.04) (-1.155,6.232)rectangle++(0.04,0.04) (-0.735,6.94)rectangle++(0.04,0.04) (-0.349,7.524)rectangle++(0.04,0.04)
  (0.002,8.006)rectangle++(0.04,0.04) (0.32,8.401)rectangle++(0.04,0.04) (0.606,8.724)rectangle++(0.04,0.04) (0.862,8.988)rectangle++(0.04,0.04) (1.09,9.201)rectangle++(0.04,0.04)
  (1.292,9.374)rectangle++(0.04,0.04) (1.471,9.512)rectangle++(0.04,0.04) (1.629,9.623)rectangle++(0.04,0.04) (1.767,9.71)rectangle++(0.04,0.04) (1.888,9.779)rectangle++(0.04,0.04)
  (1.994,9.833)rectangle++(0.04,0.04) (2.086,9.875)rectangle++(0.04,0.04) (-1.246,3.378)rectangle++(0.04,0.04) (-0.926,4.517)rectangle++(0.04,0.04) (-0.609,5.473)rectangle++(0.04,0.04)
  (1.427,2.879)rectangle++(0.04,0.04) (1.326,3.987)rectangle++(0.04,0.04) (1.284,4.932)rectangle++(0.04,0.04) (-0.773,4.162)rectangle++(0.04,0.04) (-0.493,5.165)rectangle++(0.04,0.04)
  (-0.215,6.006)rectangle++(0.04,0.04) (0.056,6.71)rectangle++(0.04,0.04) (0.314,7.297)rectangle++(0.04,0.04) (0.556,7.786)rectangle++(0.04,0.04) (0.782,8.192)rectangle++(0.04,0.04)
  (0.99,8.528)rectangle++(0.04,0.04) (1.181,8.806)rectangle++(0.04,0.04) (2.042,3.928)rectangle++(0.04,0.04) (1.891,4.853)rectangle++(0.04,0.04) (-0.648,4.26)rectangle++(0.04,0.04)
  (-0.382,5.243)rectangle++(0.04,0.04) (-0.117,6.068)rectangle++(0.04,0.04) (1.516,3.106)rectangle++(0.04,0.04) (1.411,4.175)rectangle++(0.04,0.04) (1.364,5.089)rectangle++(0.04,0.04)
  (1.361,5.867)rectangle++(0.04,0.04) (-0.958,4.361)rectangle++(0.04,0.04) (-0.642,5.341)rectangle++(0.04,0.04) (1.398,2.847)rectangle++(0.04,0.04) (1.3,3.96)rectangle++(0.04,0.04)
  (1.261,4.91)rectangle++(0.04,0.04) (1.266,5.719)rectangle++(0.04,0.04) (1.303,6.407)rectangle++(0.04,0.04) (1.361,6.99)rectangle++(0.04,0.04) (1.435,7.483)rectangle++(0.04,0.04)
  (1.517,7.9)rectangle++(0.04,0.04) (1.603,8.25)rectangle++(0.04,0.04) (1.69,8.545)rectangle++(0.04,0.04) (1.776,8.792)rectangle++(0.04,0.04) (1.859,8.998)rectangle++(0.04,0.04)
  (-1.533,5.073)rectangle++(0.04,0.04) (-1.102,5.969)rectangle++(0.04,0.04) (-0.7,6.714)rectangle++(0.04,0.04) (-1.507,4.142)rectangle++(0.04,0.04) (-1.117,5.177)rectangle++(0.04,0.04)
  (-0.745,6.041)rectangle++(0.04,0.04) (-0.394,6.761)rectangle++(0.04,0.04) (-0.066,7.359)rectangle++(0.04,0.04) (0.236,7.854)rectangle++(0.04,0.04) (0.512,8.263)rectangle++(0.04,0.04)
  (0.764,8.599)rectangle++(0.04,0.04) (0.991,8.875)rectangle++(0.04,0.04) (1.195,9.1)rectangle++(0.04,0.04) (1.378,9.283)rectangle++(0.04,0.04) (1.54,9.432)rectangle++(0.04,0.04)
  (1.684,9.552)rectangle++(0.04,0.04) (1.811,9.648)rectangle++(0.04,0.04) (-1.679,5.192)rectangle++(0.04,0.04) (1.482,2.582)rectangle++(0.04,0.04) (1.361,3.732)rectangle++(0.04,0.04)
  (1.304,4.714)rectangle++(0.04,0.04) (1.295,5.551)rectangle++(0.04,0.04) (-0.903,4.286)rectangle++(0.04,0.04) (-0.598,5.276)rectangle++(0.04,0.04) (-0.299,6.104)rectangle++(0.04,0.04)
  (-0.012,6.797)rectangle++(0.04,0.04) (0.259,7.374)rectangle++(0.04,0.04) (0.513,7.854)rectangle++(0.04,0.04) (0.748,8.251)rectangle++(0.04,0.04) (0.964,8.58)rectangle++(0.04,0.04)
  (-1.11,3.292)rectangle++(0.04,0.04) (-0.814,4.438)rectangle++(0.04,0.04) (1.195,2.633)rectangle++(0.04,0.04) (1.118,3.787)rectangle++(0.04,0.04) (0.807,2.946)rectangle++(0.04,0.04)
  (-0.437,3.681)rectangle++(0.04,0.04) (-0.227,4.742)rectangle++(0.04,0.04) (-0.005,5.636)rectangle++(0.04,0.04) (1.396,3.048)rectangle++(0.04,0.04) (1.307,4.131)rectangle++(0.04,0.04)
  (1.274,5.055)rectangle++(0.04,0.04) (1.283,5.842)rectangle++(0.04,0.04) (1.322,6.511)rectangle++(0.04,0.04) (-0.922,3.332)rectangle++(0.04,0.04) (-0.653,4.466)rectangle++(0.04,0.04)
  (-0.378,5.418)rectangle++(0.04,0.04) (1.386,2.921)rectangle++(0.04,0.04) (1.293,4.024)rectangle++(0.04,0.04) (1.258,4.965)rectangle++(0.04,0.04) (-0.785,4.164)rectangle++(0.04,0.04)
  (-0.503,5.167)rectangle++(0.04,0.04) (-0.223,6.008)rectangle++(0.04,0.04) (0.049,6.712)rectangle++(0.04,0.04) (0.308,7.299)rectangle++(0.04,0.04) (1.771,3.44)rectangle++(0.04,0.04)
  (1.641,4.449)rectangle++(0.04,0.04) (0.175,2.092)rectangle++(0.04,0.04) (0.23,3.367)rectangle++(0.04,0.04) (0.328,4.449)rectangle++(0.04,0.04) (0.455,5.365)rectangle++(0.04,0.04)
  (1.275,3.1)rectangle++(0.04,0.04) (1.206,4.18)rectangle++(0.04,0.04) (1.19,5.101)rectangle++(0.04,0.04) (1.213,5.885)rectangle++(0.04,0.04) (1.264,6.55)rectangle++(0.04,0.04)
  (1.335,7.113)rectangle++(0.04,0.04) (1.417,7.589)rectangle++(0.04,0.04) (1.506,7.99)rectangle++(0.04,0.04) (-0.579,3.408)rectangle++(0.04,0.04) (-0.784,3.512)rectangle++(0.04,0.04)
  (-0.528,4.613)rectangle++(0.04,0.04) (-0.266,5.538)rectangle++(0.04,0.04) (-1.18,3.991)rectangle++(0.04,0.04) (-0.845,5.036)rectangle++(0.04,0.04) (-1.177,3.783)rectangle++(0.04,0.04)
  (-0.851,4.859)rectangle++(0.04,0.04) (-0.531,5.76)rectangle++(0.04,0.04) (-0.223,6.514)rectangle++(0.04,0.04) (0.069,7.142)rectangle++(0.04,0.04) (0.342,7.664)rectangle++(0.04,0.04)
  (0.595,8.097)rectangle++(0.04,0.04) (0.827,8.455)rectangle++(0.04,0.04) (1.039,8.75)rectangle++(0.04,0.04) (1.231,8.992)rectangle++(0.04,0.04) (-1.644,4.954)rectangle++(0.04,0.04)
  (-1.202,5.872)rectangle++(0.04,0.04) (-0.789,6.636)rectangle++(0.04,0.04) (-0.407,7.268)rectangle++(0.04,0.04) (-0.058,7.79)rectangle++(0.04,0.04) (0.261,8.22)rectangle++(0.04,0.04)
  (0.548,8.573)rectangle++(0.04,0.04) (-0.102,3.898)rectangle++(0.04,0.04) (0.005,2.113)rectangle++(0.04,0.04) (0.087,3.392)rectangle++(0.04,0.04) (0.207,4.476)rectangle++(0.04,0.04)
  (0.353,5.392)rectangle++(0.04,0.04) (1.294,3.082)rectangle++(0.04,0.04) (1.221,4.164)rectangle++(0.04,0.04) (1.202,5.087)rectangle++(0.04,0.04) (-0.821,4.177)rectangle++(0.04,0.04)
  (-0.533,5.18)rectangle++(0.04,0.04) (-0.248,6.02)rectangle++(0.04,0.04) (1.52,3.067)rectangle++(0.04,0.04) (1.412,4.142)rectangle++(0.04,0.04) (1.364,5.061)rectangle++(0.04,0.04)
  (1.36,5.843)rectangle++(0.04,0.04) (1.387,6.508)rectangle++(0.04,0.04) (1.437,7.073)rectangle++(0.04,0.04) (1.502,7.551)rectangle++(0.04,0.04) (1.577,7.954)rectangle++(0.04,0.04)
  (1.656,8.294)rectangle++(0.04,0.04) (1.737,8.58)rectangle++(0.04,0.04) (1.818,8.82)rectangle++(0.04,0.04) (1.896,9.02)rectangle++(0.04,0.04) (1.97,9.188)rectangle++(0.04,0.04)
  (2.04,9.327)rectangle++(0.04,0.04) (2.105,9.442)rectangle++(0.04,0.04) (2.164,9.538)rectangle++(0.04,0.04) (2.219,9.617)rectangle++(0.04,0.04) (2.269,9.682)rectangle++(0.04,0.04)
  (2.314,9.735)rectangle++(0.04,0.04) (2.354,9.778)rectangle++(0.04,0.04) (2.39,9.814)rectangle++(0.04,0.04) (-1.611,5.333)rectangle++(0.04,0.04) (-1.158,6.194)rectangle++(0.04,0.04)
  (-0.739,6.907)rectangle++(0.04,0.04) (-1.555,4.173)rectangle++(0.04,0.04) (-1.157,5.205)rectangle++(0.04,0.04) (-0.777,6.067)rectangle++(0.04,0.04) (-0.42,6.784)rectangle++(0.04,0.04)
  (-0.088,7.379)rectangle++(0.04,0.04) (0.218,7.872)rectangle++(0.04,0.04) (-0.075,3.821)rectangle++(0.04,0.04) (-0.781,3.688)rectangle++(0.04,0.04) (-0.518,4.762)rectangle++(0.04,0.04)
  (-0.252,5.665)rectangle++(0.04,0.04) (0.01,6.421)rectangle++(0.04,0.04) (0.263,7.054)rectangle++(0.04,0.04) (0.504,7.582)rectangle++(0.04,0.04) (0.729,8.02)rectangle++(0.04,0.04)
  (-1.528,4.668)rectangle++(0.04,0.04) (1.335,2.515)rectangle++(0.04,0.04) (1.233,3.681)rectangle++(0.04,0.04) (1.193,4.676)rectangle++(0.04,0.04) (1.199,5.523)rectangle++(0.04,0.04)
  (1.238,6.243)rectangle++(0.04,0.04) (1.3,6.853)rectangle++(0.04,0.04) (1.545,3.579)rectangle++(0.04,0.04) (1.454,4.576)rectangle++(0.04,0.04) (0.009,2.205)rectangle++(0.04,0.04)
  (0.094,3.47)rectangle++(0.04,0.04) (0.217,4.542)rectangle++(0.04,0.04) (0.364,5.448)rectangle++(0.04,0.04) (0.525,6.213)rectangle++(0.04,0.04) (0.692,6.856)rectangle++(0.04,0.04)
  (0.861,7.396)rectangle++(0.04,0.04) (1.025,7.848)rectangle++(0.04,0.04) (1.183,8.226)rectangle++(0.04,0.04) (1.332,8.541)rectangle++(0.04,0.04) (1.472,8.803)rectangle++(0.04,0.04)
  (1.601,9.02)rectangle++(0.04,0.04) (1.72,9.199)rectangle++(0.04,0.04) (1.827,9.347)rectangle++(0.04,0.04) (1.925,9.468)rectangle++(0.04,0.04) (2.119,4.225)rectangle++(0.04,0.04)
  (1.968,5.102)rectangle++(0.04,0.04) (1.033,3.393)rectangle++(0.04,0.04) (1.012,4.439)rectangle++(0.04,0.04) (1.035,5.329)rectangle++(0.04,0.04) (1.091,6.084)rectangle++(0.04,0.04)
  (1.168,6.724)rectangle++(0.04,0.04) (1.26,7.265)rectangle++(0.04,0.04) (1.359,7.721)rectangle++(0.04,0.04) (1.462,8.105)rectangle++(0.04,0.04) (1.565,8.427)rectangle++(0.04,0.04)
  (1.665,8.696)rectangle++(0.04,0.04) (1.761,8.921)rectangle++(0.04,0.04) (1.851,9.109)rectangle++(0.04,0.04) (1.936,9.265)rectangle++(0.04,0.04) (2.014,9.394)rectangle++(0.04,0.04)
  (2.085,9.5)rectangle++(0.04,0.04) (2.15,9.588)rectangle++(0.04,0.04) (2.209,9.66)rectangle++(0.04,0.04) (2.262,9.719)rectangle++(0.04,0.04) (2.309,9.767)rectangle++(0.04,0.04)
  (2.351,9.806)rectangle++(0.04,0.04) (2.389,9.837)rectangle++(0.04,0.04) (2.422,9.862)rectangle++(0.04,0.04) (2.451,9.882)rectangle++(0.04,0.04) (2.476,9.898)rectangle++(0.04,0.04)
  (2.498,9.91)rectangle++(0.04,0.04) (2.518,9.92)rectangle++(0.04,0.04) (2.535,9.928)rectangle++(0.04,0.04) (2.549,9.933)rectangle++(0.04,0.04) (2.562,9.937)rectangle++(0.04,0.04)
  (2.573,9.941)rectangle++(0.04,0.04) (2.583,9.943)rectangle++(0.04,0.04) (2.591,9.944)rectangle++(0.04,0.04) (2.598,9.945)rectangle++(0.04,0.04) (2.604,9.946)rectangle++(0.04,0.04)
  (1.328,2.332)rectangle++(0.04,0.04) (-0.019,1.609)rectangle++(0.04,0.04) (0.046,2.965)rectangle++(0.04,0.04) (0.156,4.115)rectangle++(0.04,0.04) (0.295,5.087)rectangle++(0.04,0.04)
  (1.225,3.01)rectangle++(0.04,0.04) (1.16,4.106)rectangle++(0.04,0.04) (1.148,5.04)rectangle++(0.04,0.04) (1.175,5.834)rectangle++(0.04,0.04) (1.23,6.508)rectangle++(0.04,0.04)
  (1.304,7.079)rectangle++(0.04,0.04) (1.389,7.561)rectangle++(0.04,0.04) (1.481,7.968)rectangle++(0.04,0.04) (1.8,3.835)rectangle++(0.04,0.04) (-0.448,4.042)rectangle++(0.04,0.04)
  (-0.221,5.05)rectangle++(0.04,0.04) (0.012,5.897)rectangle++(0.04,0.04) (0.244,6.609)rectangle++(0.04,0.04) (0.469,7.204)rectangle++(0.04,0.04) (0.685,7.701)rectangle++(0.04,0.04)
  (0.888,8.114)rectangle++(0.04,0.04) (1.077,8.458)rectangle++(0.04,0.04) (1.252,8.742)rectangle++(0.04,0.04) (-1.587,4.908)rectangle++(0.04,0.04) (1.401,2.548)rectangle++(0.04,0.04)
  (1.291,3.706)rectangle++(0.04,0.04) (1.243,4.695)rectangle++(0.04,0.04) (1.242,5.537)rectangle++(0.04,0.04) (1.275,6.253)rectangle++(0.04,0.04) (1.332,6.86)rectangle++(0.04,0.04)
  (1.404,7.374)rectangle++(0.04,0.04) (1.486,7.808)rectangle++(0.04,0.04) (1.574,8.174)rectangle++(0.04,0.04) (1.84,3.895)rectangle++(0.04,0.04) (1.717,4.834)rectangle++(0.04,0.04)
  (1.651,5.636)rectangle++(0.04,0.04) (1.627,6.321)rectangle++(0.04,0.04) (1.633,6.904)rectangle++(0.04,0.04) (1.662,7.399)rectangle++(0.04,0.04) (1.707,7.819)rectangle++(0.04,0.04)
  (1.761,8.174)rectangle++(0.04,0.04) (1.822,8.474)rectangle++(0.04,0.04) (1.885,8.726)rectangle++(0.04,0.04) (1.949,8.938)rectangle++(0.04,0.04) (2.012,9.115)rectangle++(0.04,0.04)
  (2.073,9.264)rectangle++(0.04,0.04) (2.13,9.388)rectangle++(0.04,0.04) (2.184,9.49)rectangle++(0.04,0.04) (2.234,9.576)rectangle++(0.04,0.04) (2.28,9.646)rectangle++(0.04,0.04)
  (2.321,9.704)rectangle++(0.04,0.04) (2.359,9.752)rectangle++(0.04,0.04) (2.393,9.791)rectangle++(0.04,0.04) (2.144,4.393)rectangle++(0.04,0.04) (1.996,5.245)rectangle++(0.04,0.04)
  (1.904,5.974)rectangle++(0.04,0.04) (1.855,6.598)rectangle++(0.04,0.04) (1.839,7.131)rectangle++(0.04,0.04) (1.846,7.584)rectangle++(0.04,0.04) (1.87,7.968)rectangle++(0.04,0.04)
  (1.906,8.295)rectangle++(0.04,0.04) (1.95,8.57)rectangle++(0.04,0.04) (1.998,8.803)rectangle++(0.04,0.04) (-1.467,5.06)rectangle++(0.04,0.04) (-1.046,5.956)rectangle++(0.04,0.04)
  (-0.653,6.701)rectangle++(0.04,0.04) (-0.29,7.318)rectangle++(0.04,0.04) (0.044,7.828)rectangle++(0.04,0.04) (0.349,8.248)rectangle++(0.04,0.04) (0.624,8.593)rectangle++(0.04,0.04)
  (0.872,8.875)rectangle++(0.04,0.04) (2.096,3.872)rectangle++(0.04,0.04) (1.934,4.803)rectangle++(0.04,0.04) (1.834,5.602)rectangle++(0.04,0.04) (0.408,2.686)rectangle++(0.04,0.04)
  (0.452,3.863)rectangle++(0.04,0.04) (0.537,4.862)rectangle++(0.04,0.04) (0.648,5.707)rectangle++(0.04,0.04) (0.777,6.421)rectangle++(0.04,0.04) (0.915,7.023)rectangle++(0.04,0.04)
  (1.057,7.529)rectangle++(0.04,0.04) (1.197,7.954)rectangle++(0.04,0.04) (1.334,8.309)rectangle++(0.04,0.04) (1.464,8.606)rectangle++(0.04,0.04) (1.586,8.852)rectangle++(0.04,0.04)
  (1.7,9.057)rectangle++(0.04,0.04) (1.805,9.227)rectangle++(0.04,0.04) (-0.27,3.232)rectangle++(0.04,0.04) (-0.102,4.354)rectangle++(0.04,0.04) (-0.899,3.788)rectangle++(0.04,0.04)
  (-0.615,4.852)rectangle++(0.04,0.04) (-0.331,5.745)rectangle++(0.04,0.04) (-1.235,4.02)rectangle++(0.04,0.04) (-0.891,5.062)rectangle++(0.04,0.04) (-0.557,5.935)rectangle++(0.04,0.04)
  (-0.238,6.663)rectangle++(0.04,0.04) (0.062,7.269)rectangle++(0.04,0.04) (0.341,7.773)rectangle++(0.04,0.04) (0.599,8.189)rectangle++(0.04,0.04) (0.834,8.533)rectangle++(0.04,0.04)
  (1.048,8.816)rectangle++(0.04,0.04) (1.241,9.048)rectangle++(0.04,0.04) (1.415,9.237)rectangle++(0.04,0.04) (1.57,9.391)rectangle++(0.04,0.04) (1.708,9.516)rectangle++(0.04,0.04)
  (1.83,9.616)rectangle++(0.04,0.04) (1.938,9.697)rectangle++(0.04,0.04) (-1.667,5.225)rectangle++(0.04,0.04) (-1.21,6.104)rectangle++(0.04,0.04) (-0.787,6.833)rectangle++(0.04,0.04)
  (-0.398,7.436)rectangle++(0.04,0.04) (-0.043,7.933)rectangle++(0.04,0.04) (0.279,8.341)rectangle++(0.04,0.04) (0.568,8.675)rectangle++(0.04,0.04) (0.828,8.947)rectangle++(0.04,0.04)
  (2.119,3.875)rectangle++(0.04,0.04) (1.954,4.805)rectangle++(0.04,0.04) (1.851,5.603)rectangle++(0.04,0.04) (1.795,6.284)rectangle++(0.04,0.04) (1.775,6.866)rectangle++(0.04,0.04)
  (1.781,7.361)rectangle++(0.04,0.04) (1.806,7.782)rectangle++(0.04,0.04) (1.844,8.139)rectangle++(0.04,0.04) (1.891,8.44)rectangle++(0.04,0.04) (1.943,8.695)rectangle++(0.04,0.04)
  (1.997,8.909)rectangle++(0.04,0.04) (2.052,9.089)rectangle++(0.04,0.04) (2.105,9.24)rectangle++(0.04,0.04) (2.157,9.366)rectangle++(0.04,0.04) (2.206,9.471)rectangle++(0.04,0.04)
  (2.252,9.558)rectangle++(0.04,0.04) (2.294,9.631)rectangle++(0.04,0.04) (2.333,9.691)rectangle++(0.04,0.04) (2.368,9.74)rectangle++(0.04,0.04) (2.4,9.78)rectangle++(0.04,0.04)
  (2.429,9.813)rectangle++(0.04,0.04) (2.455,9.84)rectangle++(0.04,0.04) (2.149,4.417)rectangle++(0.04,0.04) (2.001,5.265)rectangle++(0.04,0.04) (1.909,5.991)rectangle++(0.04,0.04)
  (1.861,6.612)rectangle++(0.04,0.04) (1.42,3.66)rectangle++(0.04,0.04) (1.351,4.65)rectangle++(0.04,0.04) (1.332,5.495)rectangle++(0.04,0.04) (0.62,2.106)rectangle++(0.04,0.04)
  (0.609,3.362)rectangle++(0.04,0.04) (0.761,2.749)rectangle++(0.04,0.04) (0.755,3.902)rectangle++(0.04,0.04) (0.795,4.883)rectangle++(0.04,0.04) (0.869,5.715)rectangle++(0.04,0.04)
  (0.965,6.419)rectangle++(0.04,0.04) (1.075,7.014)rectangle++(0.04,0.04) (1.192,7.515)rectangle++(0.04,0.04) (1.312,7.936)rectangle++(0.04,0.04) (1.43,8.289)rectangle++(0.04,0.04)
  (1.545,8.585)rectangle++(0.04,0.04) (1.654,8.832)rectangle++(0.04,0.04) (1.757,9.037)rectangle++(0.04,0.04) (1.853,9.207)rectangle++(0.04,0.04) (1.941,9.348)rectangle++(0.04,0.04)
  (-1.593,5.157)rectangle++(0.04,0.04) (-1.15,6.043)rectangle++(0.04,0.04) (-0.738,6.779)rectangle++(0.04,0.04) (-0.358,7.388)rectangle++(0.04,0.04) (-0.011,7.89)rectangle++(0.04,0.04)
  (0.304,8.303)rectangle++(0.04,0.04) (2.021,3.632)rectangle++(0.04,0.04) (1.861,4.603)rectangle++(0.04,0.04) (1.763,5.434)rectangle++(0.04,0.04) (1.714,6.145)rectangle++(0.04,0.04)
  (1.701,6.751)rectangle++(0.04,0.04) (1.713,7.266)rectangle++(0.04,0.04) (1.745,7.704)rectangle++(0.04,0.04) (1.789,8.075)rectangle++(0.04,0.04) (1.841,8.388)rectangle++(0.04,0.04)
  (1.899,8.653)rectangle++(0.04,0.04) (1.958,8.875)rectangle++(0.04,0.04) (2.017,9.062)rectangle++(0.04,0.04) (2.075,9.218)rectangle++(0.04,0.04) (2.04,4.211)rectangle++(0.04,0.04)
  (1.9,5.094)rectangle++(0.04,0.04) (1.039,3.376)rectangle++(0.04,0.04) (1.016,4.424)rectangle++(0.04,0.04) (1.038,5.316)rectangle++(0.04,0.04) (1.093,6.073)rectangle++(0.04,0.04)
  (1.17,6.715)rectangle++(0.04,0.04) (1.261,7.257)rectangle++(0.04,0.04) (1.649,3.648)rectangle++(0.04,0.04) (1.546,4.631)rectangle++(0.04,0.04) (1.497,5.471)rectangle++(0.04,0.04)
  (1.489,6.186)rectangle++(0.04,0.04) (1.511,6.795)rectangle++(0.04,0.04) (1.554,7.311)rectangle++(0.04,0.04) (1.611,7.749)rectangle++(0.04,0.04) (1.677,8.118)rectangle++(0.04,0.04)
  (1.748,8.43)rectangle++(0.04,0.04) (1.821,8.691)rectangle++(0.04,0.04) (-1.475,5.005)rectangle++(0.04,0.04) (-1.055,5.909)rectangle++(0.04,0.04) (1.587,2.862)rectangle++(0.04,0.04)
  (1.461,3.966)rectangle++(0.04,0.04) (1.398,4.908)rectangle++(0.04,0.04) (1.383,5.712)rectangle++(0.04,0.04) (1.402,6.396)rectangle++(0.04,0.04) (1.445,6.977)rectangle++(0.04,0.04)
  (1.505,7.469)rectangle++(0.04,0.04) (1.576,7.885)rectangle++(0.04,0.04) (1.653,8.235)rectangle++(0.04,0.04) (1.846,3.925)rectangle++(0.04,0.04) (1.724,4.859)rectangle++(0.04,0.04)
  (0.867,2.338)rectangle++(0.04,0.04) (0.828,3.549)rectangle++(0.04,0.04) (0.843,4.58)rectangle++(0.04,0.04) (0.898,5.455)rectangle++(0.04,0.04) (1.246,3.218)rectangle++(0.04,0.04)
  (1.185,4.282)rectangle++(0.04,0.04) (1.177,5.188)rectangle++(0.04,0.04) (1.205,5.959)rectangle++(0.04,0.04) (1.261,6.613)rectangle++(0.04,0.04) (-1.134,4.49)rectangle++(0.04,0.04)
  (-0.787,5.458)rectangle++(0.04,0.04) (-0.452,6.267)rectangle++(0.04,0.04) (-0.136,6.941)rectangle++(0.04,0.04) (-1.454,4.292)rectangle++(0.04,0.04) (1.233,2.46)rectangle++(0.04,0.04)
  (1.144,3.638)rectangle++(0.04,0.04) (1.116,4.643)rectangle++(0.04,0.04) (1.132,5.498)rectangle++(0.04,0.04) (1.18,6.224)rectangle++(0.04,0.04) (1.25,6.839)rectangle++(0.04,0.04)
  (1.334,7.36)rectangle++(0.04,0.04) (1.426,7.799)rectangle++(0.04,0.04) (1.522,8.168)rectangle++(0.04,0.04) (1.618,8.478)rectangle++(0.04,0.04) (1.712,8.738)rectangle++(0.04,0.04)
  (1.803,8.955)rectangle++(0.04,0.04) (1.888,9.136)rectangle++(0.04,0.04) (1.968,9.286)rectangle++(0.04,0.04) (2.042,9.411)rectangle++(0.04,0.04) (2.11,9.514)rectangle++(0.04,0.04)
  (2.172,9.598)rectangle++(0.04,0.04) (2.228,9.668)rectangle++(0.04,0.04) (2.278,9.725)rectangle++(0.04,0.04) (2.323,9.771)rectangle++(0.04,0.04) (-0.248,3.422)rectangle++(0.04,0.04)
  (-0.076,4.515)rectangle++(0.04,0.04) (0.114,5.437)rectangle++(0.04,0.04) (1.333,3.036)rectangle++(0.04,0.04) (1.252,4.124)rectangle++(0.04,0.04) (1.227,5.051)rectangle++(0.04,0.04)
  (1.243,5.841)rectangle++(0.04,0.04) (1.288,6.511)rectangle++(0.04,0.04) (1.353,7.079)rectangle++(0.04,0.04) (1.431,7.559)rectangle++(0.04,0.04) (-0.174,3.818)rectangle++(0.04,0.04)
  (0.002,4.848)rectangle++(0.04,0.04) (0.194,5.717)rectangle++(0.04,0.04) (0.391,6.448)rectangle++(0.04,0.04) (0.588,7.061)rectangle++(0.04,0.04) (0.78,7.575)rectangle++(0.04,0.04)
  (0.964,8.003)rectangle++(0.04,0.04) (1.869,3.724)rectangle++(0.04,0.04) (1.735,4.687)rectangle++(0.04,0.04) (0.852,2.342)rectangle++(0.04,0.04) (0.816,3.553)rectangle++(0.04,0.04)
  (0.833,4.583)rectangle++(0.04,0.04) (0.89,5.459)rectangle++(0.04,0.04) (-0.956,4.192)rectangle++(0.04,0.04) (-0.647,5.198)rectangle++(0.04,0.04) (-0.344,6.04)rectangle++(0.04,0.04)
  (-0.053,6.744)rectangle++(0.04,0.04) (0.222,7.331)rectangle++(0.04,0.04) (-1.472,4.437)rectangle++(0.04,0.04) (-1.076,5.426)rectangle++(0.04,0.04) (0.015,2.321)rectangle++(0.04,0.04)
  (0.103,3.568)rectangle++(0.04,0.04) (0.228,4.625)rectangle++(0.04,0.04) (0.377,5.518)rectangle++(0.04,0.04) (-1.06,4.108)rectangle++(0.04,0.04) (1.141,2.514)rectangle++(0.04,0.04)
  (-0.147,3.307)rectangle++(0.04,0.04) (0.005,4.413)rectangle++(0.04,0.04) (0.179,5.347)rectangle++(0.04,0.04) (0.364,6.134)rectangle++(0.04,0.04) (0.552,6.795)rectangle++(0.04,0.04)
  (0.739,7.35)rectangle++(0.04,0.04) (1.734,3.547)rectangle++(0.04,0.04) (-0.399,3.973)rectangle++(0.04,0.04) (-0.182,4.989)rectangle++(0.04,0.04) (0.042,5.844)rectangle++(0.04,0.04)
  (0.268,6.562)rectangle++(0.04,0.04) (0.488,7.163)rectangle++(0.04,0.04) (0.699,7.666)rectangle++(0.04,0.04) (0.899,8.084)rectangle++(0.04,0.04) (1.085,8.432)rectangle++(0.04,0.04)
  (1.257,8.72)rectangle++(0.04,0.04) (1.415,8.958)rectangle++(0.04,0.04) (1.559,9.154)rectangle++(0.04,0.04) (1.689,9.314)rectangle++(0.04,0.04) (1.806,9.446)rectangle++(0.04,0.04)
  (1.911,9.553)rectangle++(0.04,0.04) (2.142,4.238)rectangle++(0.04,0.04) (1.988,5.113)rectangle++(0.04,0.04) (1.033,3.399)rectangle++(0.04,0.04) (1.012,4.444)rectangle++(0.04,0.04)
  (1.036,5.333)rectangle++(0.04,0.04) (1.092,6.088)rectangle++(0.04,0.04) (1.169,6.727)rectangle++(0.04,0.04) (1.261,7.268)rectangle++(0.04,0.04) (1.36,7.723)rectangle++(0.04,0.04)
  (1.463,8.107)rectangle++(0.04,0.04) (-1.415,4.822)rectangle++(0.04,0.04) (-1.012,5.752)rectangle++(0.04,0.04) (-0.632,6.526)rectangle++(0.04,0.04) (-0.279,7.168)rectangle++(0.04,0.04)
  (-1.528,4.31)rectangle++(0.04,0.04) (-1.128,5.321)rectangle++(0.04,0.04) (-0.748,6.164)rectangle++(0.04,0.04) (-0.392,6.866)rectangle++(0.04,0.04) (-0.061,7.448)rectangle++(0.04,0.04)
  (-0.052,3.772)rectangle++(0.04,0.04) (0.105,4.804)rectangle++(0.04,0.04) (0.279,5.676)rectangle++(0.04,0.04) (1.373,3.12)rectangle++(0.04,0.04) (0.612,2.876)rectangle++(0.04,0.04)
  (-0.457,3.63)rectangle++(0.04,0.04) (-0.246,4.7)rectangle++(0.04,0.04) (-0.023,5.601)rectangle++(0.04,0.04) (0.202,6.358)rectangle++(0.04,0.04) (0.424,6.992)rectangle++(0.04,0.04)
  (0.638,7.523)rectangle++(0.04,0.04) (-1.439,4.552)rectangle++(0.04,0.04) (-1.043,5.523)rectangle++(0.04,0.04) (-0.668,6.333)rectangle++(0.04,0.04) (-0.317,7.006)rectangle++(0.04,0.04)
  (0.009,7.564)rectangle++(0.04,0.04) (0.308,8.025)rectangle++(0.04,0.04) (-1.604,4.59)rectangle++(0.04,0.04) (-1.182,5.562)rectangle++(0.04,0.04) (-0.785,6.371)rectangle++(0.04,0.04)
  (-0.414,7.043)rectangle++(0.04,0.04) (-0.073,7.599)rectangle++(0.04,0.04) (0.24,8.058)rectangle++(0.04,0.04) (0.524,8.436)rectangle++(0.04,0.04) (0.781,8.746)rectangle++(0.04,0.04)
  (2.074,3.826)rectangle++(0.04,0.04) (1.914,4.765)rectangle++(0.04,0.04) (1.815,5.57)rectangle++(0.04,0.04) (-0.813,4.388)rectangle++(0.04,0.04) (-0.518,5.358)rectangle++(0.04,0.04)
  (-0.228,6.171)rectangle++(0.04,0.04) (0.051,6.851)rectangle++(0.04,0.04) (0.315,7.418)rectangle++(0.04,0.04) (0.562,7.888)rectangle++(0.04,0.04) (0.791,8.279)rectangle++(0.04,0.04)
  (1.001,8.602)rectangle++(0.04,0.04) (1.193,8.868)rectangle++(0.04,0.04) (1.367,9.086)rectangle++(0.04,0.04) (2.09,4.024)rectangle++(0.04,0.04) (1.935,4.933)rectangle++(0.04,0.04)
  (1.84,5.712)rectangle++(0.04,0.04) (1.79,6.378)rectangle++(0.04,0.04) (1.775,6.946)rectangle++(0.04,0.04) (1.784,7.429)rectangle++(0.04,0.04) (1.811,7.84)rectangle++(0.04,0.04)
  (1.851,8.187)rectangle++(0.04,0.04) (1.899,8.481)rectangle++(0.04,0.04) (1.951,8.729)rectangle++(0.04,0.04) (2.005,8.938)rectangle++(0.04,0.04) (2.06,9.113)rectangle++(0.04,0.04)
  (2.113,9.26)rectangle++(0.04,0.04) (2.164,9.383)rectangle++(0.04,0.04) (2.07,4.263)rectangle++(0.04,0.04) (1.928,5.137)rectangle++(0.04,0.04) (1.842,5.885)rectangle++(0.04,0.04)
  (1.799,6.525)rectangle++(0.04,0.04) (1.788,7.071)rectangle++(0.04,0.04) (1.8,7.535)rectangle++(0.04,0.04) (1.829,7.929)rectangle++(0.04,0.04) (1.87,8.262)rectangle++(0.04,0.04)
  (1.918,8.544)rectangle++(0.04,0.04) (1.891,4.045)rectangle++(0.04,0.04) (1.767,4.959)rectangle++(0.04,0.04) (1.698,5.741)rectangle++(0.04,0.04) (1.671,6.408)rectangle++(0.04,0.04)
  (1.674,6.976)rectangle++(0.04,0.04) (1.7,7.459)rectangle++(0.04,0.04) (1.741,7.868)rectangle++(0.04,0.04) (1.793,8.214)rectangle++(0.04,0.04) (1.85,8.507)rectangle++(0.04,0.04)
  (1.911,8.753)rectangle++(0.04,0.04) (-1.139,3.369)rectangle++(0.04,0.04) (-0.835,4.506)rectangle++(0.04,0.04) (-0.532,5.459)rectangle++(0.04,0.04) (-0.236,6.258)rectangle++(0.04,0.04)
  (1.577,3.115)rectangle++(0.04,0.04) (1.463,4.181)rectangle++(0.04,0.04) (1.409,5.092)rectangle++(0.04,0.04) (1.399,5.868)rectangle++(0.04,0.04) (-0.951,4.368)rectangle++(0.04,0.04)
  (1.193,2.589)rectangle++(0.04,0.04) (1.116,3.749)rectangle++(0.04,0.04) (1.096,4.738)rectangle++(0.04,0.04) (1.119,5.58)rectangle++(0.04,0.04) (1.172,6.294)rectangle++(0.04,0.04)
  (1.246,6.899)rectangle++(0.04,0.04) (1.333,7.411)rectangle++(0.04,0.04) (1.427,7.842)rectangle++(0.04,0.04) (1.775,3.798)rectangle++(0.04,0.04) (0.733,3.098)rectangle++(0.04,0.04)
  (0.745,4.2)rectangle++(0.04,0.04) (0.799,5.136)rectangle++(0.04,0.04) (0.882,5.93)rectangle++(0.04,0.04) (0.985,6.601)rectangle++(0.04,0.04) (1.099,7.168)rectangle++(0.04,0.04)
  (1.219,7.645)rectangle++(0.04,0.04) (1.339,8.046)rectangle++(0.04,0.04) (1.458,8.382)rectangle++(0.04,0.04) (1.573,8.662)rectangle++(0.04,0.04) (1.681,8.896)rectangle++(0.04,0.04)
  (1.782,9.091)rectangle++(0.04,0.04) (-1.566,5.076)rectangle++(0.04,0.04) (-1.13,5.974)rectangle++(0.04,0.04) (-0.724,6.719)rectangle++(0.04,0.04) (-1.512,4.138)rectangle++(0.04,0.04)
  (-1.122,5.174)rectangle++(0.04,0.04) (1.413,2.705)rectangle++(0.04,0.04) (1.307,3.839)rectangle++(0.04,0.04) (0.005,2.105)rectangle++(0.04,0.04) (0.086,3.385)rectangle++(0.04,0.04)
  (0.828,2.635)rectangle++(0.04,0.04) (0.807,3.802)rectangle++(0.04,0.04) (0.847,2.879)rectangle++(0.04,0.04) (0.833,4.009)rectangle++(0.04,0.04) (0.866,4.971)rectangle++(0.04,0.04)
  (-0.773,3.303)rectangle++(0.04,0.04) (-0.527,4.434)rectangle++(0.04,0.04) (-0.273,5.387)rectangle++(0.04,0.04) (-0.018,6.186)rectangle++(0.04,0.04) (0.23,6.855)rectangle++(0.04,0.04)
  (0.467,7.414)rectangle++(0.04,0.04) (-1.446,4.499)rectangle++(0.04,0.04) (-1.051,5.478)rectangle++(0.04,0.04) (-0.677,6.295)rectangle++(0.04,0.04) (-0.326,6.974)rectangle++(0.04,0.04)
  (0,7.537)rectangle++(0.04,0.04) (0.299,8.002)rectangle++(0.04,0.04) (0.572,8.386)rectangle++(0.04,0.04) (2.01,3.709)rectangle++(0.04,0.04) (1.854,4.669)rectangle++(0.04,0.04)
  (1.761,5.49)rectangle++(0.04,0.04) (1.714,6.193)rectangle++(0.04,0.04) (1.333,3.545)rectangle++(0.04,0.04) (1.272,4.556)rectangle++(0.04,0.04) (-0.696,4.086)rectangle++(0.04,0.04)
  (-0.43,5.097)rectangle++(0.04,0.04) (-0.164,5.946)rectangle++(0.04,0.04) (1.492,3.072)rectangle++(0.04,0.04) (1.388,4.147)rectangle++(0.04,0.04) (1.344,5.066)rectangle++(0.04,0.04)
  (1.097,3.244)rectangle++(0.04,0.04) (1.06,4.31)rectangle++(0.04,0.04) (1.071,5.217)rectangle++(0.04,0.04) (1.117,5.988)rectangle++(0.04,0.04) (-0.147,3.653)rectangle++(0.04,0.04)
  (0.019,4.707)rectangle++(0.04,0.04) (0.202,5.597)rectangle++(0.04,0.04) (0.393,6.345)rectangle++(0.04,0.04) (0.586,6.974)rectangle++(0.04,0.04) (1.658,3.44)rectangle++(0.04,0.04);
\fill[piece2]
  (1.545,4.454)rectangle++(0.04,0.04) (1.489,5.32)rectangle++(0.04,0.04) (1.476,6.059)rectangle++(0.04,0.04) (1.495,6.687)rectangle++(0.04,0.04) (1.536,7.22)rectangle++(0.04,0.04)
  (1.592,7.672)rectangle++(0.04,0.04) (1.658,8.054)rectangle++(0.04,0.04) (1.729,8.376)rectangle++(0.04,0.04) (1.803,8.646)rectangle++(0.04,0.04) (1.876,8.873)rectangle++(0.04,0.04)
  (1.947,9.064)rectangle++(0.04,0.04) (2.016,9.222)rectangle++(0.04,0.04) (2.08,9.355)rectangle++(0.04,0.04) (2.14,9.464)rectangle++(0.04,0.04) (2.093,4.273)rectangle++(0.04,0.04)
  (1.948,5.144)rectangle++(0.04,0.04) (1.859,5.891)rectangle++(0.04,0.04) (1.241,3.52)rectangle++(0.04,0.04) (1.194,4.539)rectangle++(0.04,0.04) (1.194,5.406)rectangle++(0.04,0.04)
  (1.199,3.276)rectangle++(0.04,0.04) (1.148,4.333)rectangle++(0.04,0.04) (1.147,5.233)rectangle++(0.04,0.04) (1.182,5.998)rectangle++(0.04,0.04) (1.242,6.648)rectangle++(0.04,0.04)
  (1.501,3.526)rectangle++(0.04,0.04) (1.414,4.533)rectangle++(0.04,0.04) (1.381,5.393)rectangle++(0.04,0.04) (1.388,6.125)rectangle++(0.04,0.04) (1.422,6.747)rectangle++(0.04,0.04)
  (1.477,7.274)rectangle++(0.04,0.04) (1.544,7.72)rectangle++(0.04,0.04) (1.619,8.097)rectangle++(0.04,0.04) (-1.385,4.85)rectangle++(0.04,0.04) (-0.985,5.774)rectangle++(0.04,0.04)
  (-0.609,6.543)rectangle++(0.04,0.04) (-0.258,7.182)rectangle++(0.04,0.04) (1.809,3.289)rectangle++(0.04,0.04) (1.667,4.319)rectangle++(0.04,0.04) (1.587,5.201)rectangle++(0.04,0.04)
  (-0.776,4.272)rectangle++(0.04,0.04) (-0.491,5.258)rectangle++(0.04,0.04) (-0.209,6.085)rectangle++(0.04,0.04) (0.064,6.777)rectangle++(0.04,0.04) (0.323,7.354)rectangle++(0.04,0.04)
  (0.566,7.834)rectangle++(0.04,0.04) (0.793,8.233)rectangle++(0.04,0.04) (1.001,8.562)rectangle++(0.04,0.04) (1.191,8.834)rectangle++(0.04,0.04) (-1.618,4.915)rectangle++(0.04,0.04)
  (1.407,2.543)rectangle++(0.04,0.04) (1.295,3.701)rectangle++(0.04,0.04) (-0.51,3.921)rectangle++(0.04,0.04) (-0.88,3.626)rectangle++(0.04,0.04) (-0.605,4.714)rectangle++(0.04,0.04)
  (-0.328,5.627)rectangle++(0.04,0.04) (-0.056,6.392)rectangle++(0.04,0.04) (0.206,7.032)rectangle++(0.04,0.04) (0.454,7.565)rectangle++(0.04,0.04) (1.82,3.524)rectangle++(0.04,0.04)
  (1.686,4.519)rectangle++(0.04,0.04) (1.612,5.37)rectangle++(0.04,0.04) (1.582,6.096)rectangle++(0.04,0.04) (1.324,3.497)rectangle++(0.04,0.04) (1.263,4.516)rectangle++(0.04,0.04)
  (-0.688,4.076)rectangle++(0.04,0.04) (-0.424,5.088)rectangle++(0.04,0.04) (-0.159,5.938)rectangle++(0.04,0.04) (0.1,6.65)rectangle++(0.04,0.04) (0.349,7.245)rectangle++(0.04,0.04)
  (0.584,7.74)rectangle++(0.04,0.04) (0.804,8.152)rectangle++(0.04,0.04) (1.007,8.493)rectangle++(0.04,0.04) (1.194,8.775)rectangle++(0.04,0.04) (-1.605,4.904)rectangle++(0.04,0.04)
  (-1.17,5.829)rectangle++(0.04,0.04) (-0.764,6.597)rectangle++(0.04,0.04) (-0.387,7.234)rectangle++(0.04,0.04) (-0.042,7.761)rectangle++(0.04,0.04) (0.272,8.195)rectangle++(0.04,0.04)
  (0.557,8.551)rectangle++(0.04,0.04) (0.813,8.842)rectangle++(0.04,0.04) (1.043,9.079)rectangle++(0.04,0.04) (1.247,9.272)rectangle++(0.04,0.04) (1.429,9.428)rectangle++(0.04,0.04)
  (-1.701,5.077)rectangle++(0.04,0.04) (-1.245,5.979)rectangle++(0.04,0.04) (-0.821,6.728)rectangle++(0.04,0.04) (-0.431,7.348)rectangle++(0.04,0.04) (-0.075,7.859)rectangle++(0.04,0.04)
  (0.249,8.28)rectangle++(0.04,0.04) (0.54,8.624)rectangle++(0.04,0.04) (0.802,8.905)rectangle++(0.04,0.04) (1.036,9.133)rectangle++(0.04,0.04) (1.243,9.318)rectangle++(0.04,0.04)
  (-1.711,5.02)rectangle++(0.04,0.04) (-1.256,5.932)rectangle++(0.04,0.04) (-0.832,6.689)rectangle++(0.04,0.04) (-0.442,7.315)rectangle++(0.04,0.04) (-0.085,7.831)rectangle++(0.04,0.04)
  (0.238,8.256)rectangle++(0.04,0.04) (0.531,8.605)rectangle++(0.04,0.04) (0.793,8.889)rectangle++(0.04,0.04) (1.027,9.12)rectangle++(0.04,0.04) (-1.708,4.941)rectangle++(0.04,0.04)
  (-1.256,5.864)rectangle++(0.04,0.04) (1.599,2.808)rectangle++(0.04,0.04) (1.469,3.919)rectangle++(0.04,0.04) (1.404,4.868)rectangle++(0.04,0.04) (1.386,5.678)rectangle++(0.04,0.04)
  (1.403,6.367)rectangle++(0.04,0.04) (1.445,6.952)rectangle++(0.04,0.04) (1.504,7.448)rectangle++(0.04,0.04) (1.574,7.867)rectangle++(0.04,0.04) (1.65,8.22)rectangle++(0.04,0.04)
  (1.729,8.517)rectangle++(0.04,0.04) (1.809,8.767)rectangle++(0.04,0.04) (1.959,4.063)rectangle++(0.04,0.04) (1.826,4.972)rectangle++(0.04,0.04) (1.748,5.749)rectangle++(0.04,0.04)
  (0.001,2.026)rectangle++(0.04,0.04) (0.08,3.318)rectangle++(0.04,0.04) (0.198,4.413)rectangle++(0.04,0.04) (0.343,5.339)rectangle++(0.04,0.04) (0.503,6.121)rectangle++(0.04,0.04)
  (0.67,6.779)rectangle++(0.04,0.04) (0.838,7.332)rectangle++(0.04,0.04) (1.004,7.794)rectangle++(0.04,0.04) (1.163,8.181)rectangle++(0.04,0.04) (1.313,8.504)rectangle++(0.04,0.04)
  (1.454,8.772)rectangle++(0.04,0.04) (1.585,8.994)rectangle++(0.04,0.04) (1.705,9.178)rectangle++(0.04,0.04) (1.814,9.329)rectangle++(0.04,0.04) (2.098,4.173)rectangle++(0.04,0.04)
  (1.948,5.059)rectangle++(0.04,0.04) (1.856,5.819)rectangle++(0.04,0.04) (1.808,6.468)rectangle++(0.04,0.04) (1.793,7.021)rectangle++(0.04,0.04) (1.803,7.493)rectangle++(0.04,0.04)
  (1.83,7.893)rectangle++(0.04,0.04) (1.869,8.232)rectangle++(0.04,0.04) (0.033,2.702)rectangle++(0.04,0.04) (0.134,3.892)rectangle++(0.04,0.04) (-0.759,3.741)rectangle++(0.04,0.04)
  (-0.497,4.806)rectangle++(0.04,0.04) (-0.233,5.702)rectangle++(0.04,0.04) (0.028,6.452)rectangle++(0.04,0.04) (0.28,7.079)rectangle++(0.04,0.04) (0.519,7.602)rectangle++(0.04,0.04)
  (0.743,8.037)rectangle++(0.04,0.04) (0.951,8.398)rectangle++(0.04,0.04) (1.142,8.697)rectangle++(0.04,0.04) (1.316,8.943)rectangle++(0.04,0.04) (-1.149,3.319)rectangle++(0.04,0.04)
  (-0.846,4.463)rectangle++(0.04,0.04) (-1.055,3.67)rectangle++(0.04,0.04) (-0.752,4.758)rectangle++(0.04,0.04) (-0.451,5.671)rectangle++(0.04,0.04) (-0.159,6.434)rectangle++(0.04,0.04)
  (0.12,7.072)rectangle++(0.04,0.04) (0.383,7.602)rectangle++(0.04,0.04) (0.627,8.043)rectangle++(0.04,0.04) (-1.551,4.653)rectangle++(0.04,0.04) (-1.134,5.613)rectangle++(0.04,0.04)
  (-0.742,6.413)rectangle++(0.04,0.04) (-0.376,7.077)rectangle++(0.04,0.04) (0.025,2.545)rectangle++(0.04,0.04) (0.121,3.759)rectangle++(0.04,0.04) (0.251,4.786)rectangle++(0.04,0.04)
  (0.403,5.654)rectangle++(0.04,0.04) (0.566,6.386)rectangle++(0.04,0.04) (0.735,7.002)rectangle++(0.04,0.04) (0.902,7.518)rectangle++(0.04,0.04) (1.756,3.617)rectangle++(0.04,0.04)
  (-0.409,3.991)rectangle++(0.04,0.04) (-0.191,5.005)rectangle++(0.04,0.04) (1.262,2.884)rectangle++(0.04,0.04) (1.186,3.997)rectangle++(0.04,0.04) (-0.592,3.959)rectangle++(0.04,0.04)
  (-0.347,4.985)rectangle++(0.04,0.04) (1.275,2.845)rectangle++(0.04,0.04) (1.196,3.963)rectangle++(0.04,0.04) (1.173,4.917)rectangle++(0.04,0.04) (-0.79,4.138)rectangle++(0.04,0.04)
  (-0.508,5.145)rectangle++(0.04,0.04) (-0.228,5.99)rectangle++(0.04,0.04) (0.043,6.697)rectangle++(0.04,0.04) (0.303,7.287)rectangle++(0.04,0.04) (0.546,7.778)rectangle++(0.04,0.04)
  (0.773,8.186)rectangle++(0.04,0.04) (0.983,8.523)rectangle++(0.04,0.04) (1.174,8.801)rectangle++(0.04,0.04) (2.042,3.926)rectangle++(0.04,0.04) (1.891,4.851)rectangle++(0.04,0.04)
  (1.799,5.644)rectangle++(0.04,0.04) (1.187,3.459)rectangle++(0.04,0.04) (1.145,4.489)rectangle++(0.04,0.04) (1.151,5.366)rectangle++(0.04,0.04) (1.191,6.111)rectangle++(0.04,0.04)
  (1.254,6.743)rectangle++(0.04,0.04) (1.523,3.547)rectangle++(0.04,0.04) (1.434,4.55)rectangle++(0.04,0.04) (1.399,5.407)rectangle++(0.04,0.04) (1.403,6.136)rectangle++(0.04,0.04)
  (1.436,6.756)rectangle++(0.04,0.04) (1.489,7.281)rectangle++(0.04,0.04) (1.554,7.726)rectangle++(0.04,0.04) (1.628,8.101)rectangle++(0.04,0.04) (1.706,8.417)rectangle++(0.04,0.04)
  (1.784,8.682)rectangle++(0.04,0.04) (1.941,4.041)rectangle++(0.04,0.04) (-0.001,1.98)rectangle++(0.04,0.04) (0.076,3.28)rectangle++(0.04,0.04) (0.194,4.381)rectangle++(0.04,0.04)
  (0.338,5.312)rectangle++(0.04,0.04) (0.497,6.098)rectangle++(0.04,0.04) (0.664,6.76)rectangle++(0.04,0.04) (0.833,7.315)rectangle++(0.04,0.04) (0.998,7.781)rectangle++(0.04,0.04)
  (1.158,8.17)rectangle++(0.04,0.04) (1.309,8.494)rectangle++(0.04,0.04) (1.45,8.764)rectangle++(0.04,0.04) (1.581,8.988)rectangle++(0.04,0.04) (2.04,4.054)rectangle++(0.04,0.04)
  (1.894,4.961)rectangle++(0.04,0.04) (1.806,5.737)rectangle++(0.04,0.04) (1.763,6.4)rectangle++(0.04,0.04) (1.752,6.966)rectangle++(0.04,0.04) (1.766,7.447)rectangle++(0.04,0.04)
  (1.797,7.856)rectangle++(0.04,0.04) (1.839,8.202)rectangle++(0.04,0.04) (1.889,8.494)rectangle++(0.04,0.04) (-1.421,4.979)rectangle++(0.04,0.04) (-1.011,5.885)rectangle++(0.04,0.04)
  (-0.626,6.639)rectangle++(0.04,0.04) (-0.269,7.264)rectangle++(0.04,0.04) (0.06,7.781)rectangle++(0.04,0.04) (0.36,8.208)rectangle++(0.04,0.04) (0.632,8.559)rectangle++(0.04,0.04)
  (0.878,8.846)rectangle++(0.04,0.04) (1.098,9.08)rectangle++(0.04,0.04) (1.294,9.27)rectangle++(0.04,0.04) (1.468,9.424)rectangle++(0.04,0.04) (1.623,9.548)rectangle++(0.04,0.04)
  (1.759,9.647)rectangle++(0.04,0.04) (1.879,9.726)rectangle++(0.04,0.04) (1.984,9.788)rectangle++(0.04,0.04) (2.076,9.837)rectangle++(0.04,0.04) (2.156,9.874)rectangle++(0.04,0.04)
  (2.225,9.903)rectangle++(0.04,0.04) (2.285,9.925)rectangle++(0.04,0.04) (2.337,9.941)rectangle++(0.04,0.04) (2.188,4.41)rectangle++(0.04,0.04) (0.008,2.182)rectangle++(0.04,0.04)
  (0.092,3.451)rectangle++(0.04,0.04) (0.214,4.526)rectangle++(0.04,0.04) (0.361,5.435)rectangle++(0.04,0.04) (-1.045,4.088)rectangle++(0.04,0.04) (-0.727,5.113)rectangle++(0.04,0.04)
  (-0.416,5.971)rectangle++(0.04,0.04) (-0.117,6.688)rectangle++(0.04,0.04) (0.166,7.286)rectangle++(0.04,0.04) (0.316,2.844)rectangle++(0.04,0.04) (0.381,4.001)rectangle++(0.04,0.04)
  (0.481,4.982)rectangle++(0.04,0.04) (0.606,5.812)rectangle++(0.04,0.04) (0.746,6.512)rectangle++(0.04,0.04) (0.892,7.101)rectangle++(0.04,0.04) (1.04,7.597)rectangle++(0.04,0.04)
  (1.186,8.012)rectangle++(0.04,0.04) (1.326,8.359)rectangle++(0.04,0.04) (1.459,8.648)rectangle++(0.04,0.04) (1.584,8.889)rectangle++(0.04,0.04) (1.7,9.088)rectangle++(0.04,0.04)
  (1.806,9.253)rectangle++(0.04,0.04) (1.903,9.389)rectangle++(0.04,0.04) (1.991,9.501)rectangle++(0.04,0.04) (2.07,9.592)rectangle++(0.04,0.04) (2.141,9.667)rectangle++(0.04,0.04)
  (2.205,9.727)rectangle++(0.04,0.04) (2.261,9.776)rectangle++(0.04,0.04) (-1.626,5.301)rectangle++(0.04,0.04) (1.503,2.615)rectangle++(0.04,0.04) (1.38,3.759)rectangle++(0.04,0.04)
  (1.321,4.736)rectangle++(0.04,0.04) (1.31,5.569)rectangle++(0.04,0.04) (1.334,6.278)rectangle++(0.04,0.04) (1.383,6.879)rectangle++(0.04,0.04) (1.449,7.388)rectangle++(0.04,0.04)
  (1.525,7.818)rectangle++(0.04,0.04) (-1.342,4.777)rectangle++(0.04,0.04) (-0.952,5.71)rectangle++(0.04,0.04) (-0.583,6.488)rectangle++(0.04,0.04) (-0.238,7.134)rectangle++(0.04,0.04)
  (0.081,7.67)rectangle++(0.04,0.04) (1.89,3.459)rectangle++(0.04,0.04) (1.743,4.461)rectangle++(0.04,0.04) (1.657,5.318)rectangle++(0.04,0.04) (1.619,6.051)rectangle++(0.04,0.04)
  (1.616,6.674)rectangle++(0.04,0.04) (1.639,7.205)rectangle++(0.04,0.04) (-1.192,4.677)rectangle++(0.04,0.04) (-0.828,5.619)rectangle++(0.04,0.04) (-0.481,6.406)rectangle++(0.04,0.04)
  (-0.155,7.06)rectangle++(0.04,0.04) (0.148,7.604)rectangle++(0.04,0.04) (0.428,8.053)rectangle++(0.04,0.04) (1.944,3.612)rectangle++(0.04,0.04) (1.795,4.589)rectangle++(0.04,0.04)
  (1.707,5.425)rectangle++(0.04,0.04) (1.666,6.139)rectangle++(0.04,0.04) (1.659,6.748)rectangle++(0.04,0.04) (1.678,7.265)rectangle++(0.04,0.04) (1.715,7.705)rectangle++(0.04,0.04)
  (1.764,8.077)rectangle++(0.04,0.04) (1.82,8.391)rectangle++(0.04,0.04) (1.88,8.656)rectangle++(0.04,0.04) (1.942,8.878)rectangle++(0.04,0.04) (2.004,9.065)rectangle++(0.04,0.04)
  (2.064,9.221)rectangle++(0.04,0.04) (-1.544,5.155)rectangle++(0.04,0.04) (-1.108,6.04)rectangle++(0.04,0.04) (-0.702,6.774)rectangle++(0.04,0.04) (-0.328,7.382)rectangle++(0.04,0.04)
  (0.014,7.884)rectangle++(0.04,0.04) (0.325,8.297)rectangle++(0.04,0.04) (0.606,8.636)rectangle++(0.04,0.04) (0.858,8.913)rectangle++(0.04,0.04) (1.084,9.137)rectangle++(0.04,0.04)
  (1.285,9.32)rectangle++(0.04,0.04) (1.463,9.467)rectangle++(0.04,0.04) (1.62,9.584)rectangle++(0.04,0.04) (1.758,9.678)rectangle++(0.04,0.04) (1.879,9.752)rectangle++(0.04,0.04)
  (-1.689,5.225)rectangle++(0.04,0.04) (-1.229,6.105)rectangle++(0.04,0.04) (-0.803,6.835)rectangle++(0.04,0.04) (-0.411,7.438)rectangle++(0.04,0.04) (-0.054,7.935)rectangle++(0.04,0.04)
  (0.269,8.343)rectangle++(0.04,0.04) (0.56,8.677)rectangle++(0.04,0.04) (0.821,8.949)rectangle++(0.04,0.04) (1.054,9.17)rectangle++(0.04,0.04) (1.26,9.349)rectangle++(0.04,0.04)
  (1.443,9.492)rectangle++(0.04,0.04) (1.604,9.607)rectangle++(0.04,0.04) (1.745,9.698)rectangle++(0.04,0.04) (1.869,9.77)rectangle++(0.04,0.04) (1.978,9.826)rectangle++(0.04,0.04)
  (2.072,9.869)rectangle++(0.04,0.04) (2.154,9.902)rectangle++(0.04,0.04) (2.224,9.927)rectangle++(0.04,0.04) (2.286,9.945)rectangle++(0.04,0.04) (-1.657,5.339)rectangle++(0.04,0.04)
  (-1.197,6.201)rectangle++(0.04,0.04) (-0.772,6.915)rectangle++(0.04,0.04) (-0.382,7.505)rectangle++(0.04,0.04) (-0.027,7.991)rectangle++(0.04,0.04) (0.295,8.389)rectangle++(0.04,0.04)
  (0.584,8.715)rectangle++(0.04,0.04) (0.843,8.981)rectangle++(0.04,0.04) (1.073,9.196)rectangle++(0.04,0.04) (1.278,9.37)rectangle++(0.04,0.04) (1.459,9.51)rectangle++(0.04,0.04)
  (0.041,2.876)rectangle++(0.04,0.04) (0.148,4.039)rectangle++(0.04,0.04) (0.285,5.024)rectangle++(0.04,0.04) (0.441,5.855)rectangle++(0.04,0.04) (0.607,6.555)rectangle++(0.04,0.04)
  (0.776,7.144)rectangle++(0.04,0.04) (0.943,7.637)rectangle++(0.04,0.04) (1.105,8.05)rectangle++(0.04,0.04) (1.259,8.395)rectangle++(0.04,0.04) (1.404,8.681)rectangle++(0.04,0.04)
  (1.538,8.919)rectangle++(0.04,0.04) (1.662,9.116)rectangle++(0.04,0.04) (1.775,9.278)rectangle++(0.04,0.04) (1.878,9.412)rectangle++(0.04,0.04) (-1.617,5.158)rectangle++(0.04,0.04)
  (-1.171,6.045)rectangle++(0.04,0.04) (-0.755,6.781)rectangle++(0.04,0.04) (-0.373,7.39)rectangle++(0.04,0.04) (-0.024,7.893)rectangle++(0.04,0.04) (0.293,8.306)rectangle++(0.04,0.04)
  (0.579,8.645)rectangle++(0.04,0.04) (0.836,8.921)rectangle++(0.04,0.04) (1.065,9.146)rectangle++(0.04,0.04) (1.269,9.327)rectangle++(0.04,0.04) (1.45,9.474)rectangle++(0.04,0.04)
  (1.609,9.591)rectangle++(0.04,0.04) (1.749,9.684)rectangle++(0.04,0.04) (1.872,9.758)rectangle++(0.04,0.04) (1.979,9.815)rectangle++(0.04,0.04) (2.073,9.86)rectangle++(0.04,0.04)
  (2.199,4.334)rectangle++(0.04,0.04) (2.04,5.192)rectangle++(0.04,0.04) (1.94,5.928)rectangle++(0.04,0.04) (1.884,6.557)rectangle++(0.04,0.04) (1.861,7.095)rectangle++(0.04,0.04)
  (1.864,7.552)rectangle++(0.04,0.04) (1.884,7.941)rectangle++(0.04,0.04) (1.917,8.271)rectangle++(0.04,0.04) (-1.368,4.94)rectangle++(0.04,0.04) (-0.968,5.85)rectangle++(0.04,0.04)
  (-0.591,6.607)rectangle++(0.04,0.04) (-0.24,7.236)rectangle++(0.04,0.04) (0.083,7.756)rectangle++(0.04,0.04) (0.379,8.186)rectangle++(0.04,0.04) (0.647,8.539)rectangle++(0.04,0.04)
  (0.889,8.828)rectangle++(0.04,0.04) (1.107,9.065)rectangle++(0.04,0.04) (1.301,9.257)rectangle++(0.04,0.04) (1.474,9.413)rectangle++(0.04,0.04) (1.627,9.538)rectangle++(0.04,0.04)
  (1.763,9.638)rectangle++(0.04,0.04) (1.882,9.718)rectangle++(0.04,0.04) (1.986,9.781)rectangle++(0.04,0.04) (2.077,9.831)rectangle++(0.04,0.04) (-0.29,3.251)rectangle++(0.04,0.04)
  (0.839,2.523)rectangle++(0.04,0.04) (-0.342,3.603)rectangle++(0.04,0.04) (-0.149,4.673)rectangle++(0.04,0.04) (0.058,5.574)rectangle++(0.04,0.04) (0.27,6.332)rectangle++(0.04,0.04)
  (1.536,3.245)rectangle++(0.04,0.04) (1.433,4.293)rectangle++(0.04,0.04) (1.388,5.188)rectangle++(0.04,0.04) (1.385,5.95)rectangle++(0.04,0.04) (1.413,6.598)rectangle++(0.04,0.04)
  (1.469,3.555)rectangle++(0.04,0.04) (1.388,4.559)rectangle++(0.04,0.04) (1.36,5.416)rectangle++(0.04,0.04) (1.371,6.146)rectangle++(0.04,0.04) (-1.015,4.418)rectangle++(0.04,0.04)
  (-0.688,5.392)rectangle++(0.04,0.04) (-0.372,6.207)rectangle++(0.04,0.04) (-0.07,6.887)rectangle++(0.04,0.04) (0.214,7.453)rectangle++(0.04,0.04) (0.478,7.923)rectangle++(0.04,0.04)
  (0.721,8.311)rectangle++(0.04,0.04) (0.943,8.632)rectangle++(0.04,0.04) (1.145,8.896)rectangle++(0.04,0.04) (1.327,9.112)rectangle++(0.04,0.04) (1.49,9.288)rectangle++(0.04,0.04)
  (1.636,9.431)rectangle++(0.04,0.04) (1.765,9.548)rectangle++(0.04,0.04) (1.88,9.641)rectangle++(0.04,0.04) (1.982,9.716)rectangle++(0.04,0.04) (2.071,9.775)rectangle++(0.04,0.04)
  (2.149,9.822)rectangle++(0.04,0.04) (2.217,9.859)rectangle++(0.04,0.04) (2.277,9.888)rectangle++(0.04,0.04) (2.329,9.91)rectangle++(0.04,0.04) (2.374,9.926)rectangle++(0.04,0.04)
  (2.412,9.939)rectangle++(0.04,0.04) (2.446,9.948)rectangle++(0.04,0.04) (-1.629,5.37)rectangle++(0.04,0.04) (-1.172,6.226)rectangle++(0.04,0.04) (1.679,2.897)rectangle++(0.04,0.04)
  (1.541,3.991)rectangle++(0.04,0.04) (-0.527,4.024)rectangle++(0.04,0.04) (-0.29,5.038)rectangle++(0.04,0.04) (-0.047,5.89)rectangle++(0.04,0.04) (0.194,6.605)rectangle++(0.04,0.04)
  (0.427,7.202)rectangle++(0.04,0.04) (0.648,7.701)rectangle++(0.04,0.04) (0.857,8.116)rectangle++(0.04,0.04) (1.051,8.461)rectangle++(0.04,0.04) (1.23,8.746)rectangle++(0.04,0.04)
  (1.393,8.981)rectangle++(0.04,0.04) (1.541,9.174)rectangle++(0.04,0.04) (1.675,9.333)rectangle++(0.04,0.04) (1.794,9.462)rectangle++(0.04,0.04) (1.902,9.567)rectangle++(0.04,0.04)
  (1.997,9.652)rectangle++(0.04,0.04) (2.081,9.721)rectangle++(0.04,0.04) (2.156,9.776)rectangle++(0.04,0.04) (2.221,9.819)rectangle++(0.04,0.04) (2.279,9.854)rectangle++(0.04,0.04)
  (2.329,9.881)rectangle++(0.04,0.04) (2.174,4.396)rectangle++(0.04,0.04) (2.022,5.246)rectangle++(0.04,0.04) (1.926,5.974)rectangle++(0.04,0.04) (1.874,6.597)rectangle++(0.04,0.04)
  (1.854,7.129)rectangle++(0.04,0.04) (1.859,7.582)rectangle++(0.04,0.04) (1.881,7.966)rectangle++(0.04,0.04) (1.916,8.292)rectangle++(0.04,0.04) (1.958,8.568)rectangle++(0.04,0.04)
  (2.005,8.801)rectangle++(0.04,0.04) (2.054,8.997)rectangle++(0.04,0.04) (-1.498,5.109)rectangle++(0.04,0.04) (-1.071,5.999)rectangle++(0.04,0.04) (-0.673,6.738)rectangle++(0.04,0.04)
  (-0.304,7.35)rectangle++(0.04,0.04) (0.033,7.856)rectangle++(0.04,0.04) (0.34,8.273)rectangle++(0.04,0.04) (0.618,8.614)rectangle++(0.04,0.04) (0.868,8.894)rectangle++(0.04,0.04)
  (1.091,9.121)rectangle++(0.04,0.04) (-1.697,4.953)rectangle++(0.04,0.04) (-1.246,5.874)rectangle++(0.04,0.04) (-0.827,6.639)rectangle++(0.04,0.04) (-0.439,7.272)rectangle++(0.04,0.04)
  (-0.085,7.795)rectangle++(0.04,0.04) (0.238,8.226)rectangle++(0.04,0.04) (0.529,8.578)rectangle++(0.04,0.04) (0.79,8.867)rectangle++(0.04,0.04) (1.024,9.101)rectangle++(0.04,0.04)
  (1.233,9.291)rectangle++(0.04,0.04) (1.417,9.444)rectangle++(0.04,0.04) (1.58,9.567)rectangle++(0.04,0.04) (1.724,9.665)rectangle++(0.04,0.04) (1.849,9.743)rectangle++(0.04,0.04)
  (-0.21,4.045)rectangle++(0.04,0.04) (-0.019,5.043)rectangle++(0.04,0.04) (0.012,2.269)rectangle++(0.04,0.04) (-0.436,3.398)rectangle++(0.04,0.04) (-0.237,4.502)rectangle++(0.04,0.04)
  (-0.023,5.432)rectangle++(0.04,0.04) (0.195,6.215)rectangle++(0.04,0.04) (0.412,6.871)rectangle++(0.04,0.04) (0.623,7.42)rectangle++(0.04,0.04) (0.824,7.878)rectangle++(0.04,0.04)
  (1.014,8.26)rectangle++(0.04,0.04) (1.19,8.576)rectangle++(0.04,0.04) (1.352,8.839)rectangle++(0.04,0.04) (1.501,9.055)rectangle++(0.04,0.04) (1.636,9.233)rectangle++(0.04,0.04)
  (1.757,9.379)rectangle++(0.04,0.04) (1.867,9.498)rectangle++(0.04,0.04) (1.964,9.595)rectangle++(0.04,0.04) (2.051,9.673)rectangle++(0.04,0.04) (2.128,9.736)rectangle++(0.04,0.04)
  (2.196,9.787)rectangle++(0.04,0.04) (2.256,9.827)rectangle++(0.04,0.04) (2.309,9.859)rectangle++(0.04,0.04) (2.354,9.884)rectangle++(0.04,0.04) (2.172,4.403)rectangle++(0.04,0.04)
  (2.02,5.251)rectangle++(0.04,0.04) (-0.709,4.363)rectangle++(0.04,0.04) (-0.43,5.333)rectangle++(0.04,0.04) (-0.155,6.147)rectangle++(0.04,0.04) (0.112,6.827)rectangle++(0.04,0.04)
  (0.366,7.395)rectangle++(0.04,0.04) (0.605,7.867)rectangle++(0.04,0.04) (0.827,8.259)rectangle++(0.04,0.04) (1.031,8.583)rectangle++(0.04,0.04) (-1.593,4.835)rectangle++(0.04,0.04)
  (-1.163,5.77)rectangle++(0.04,0.04) (-0.76,6.547)rectangle++(0.04,0.04) (-0.386,7.192)rectangle++(0.04,0.04) (-0.043,7.725)rectangle++(0.04,0.04) (0.27,8.164)rectangle++(0.04,0.04)
  (0.554,8.525)rectangle++(0.04,0.04) (0.81,8.82)rectangle++(0.04,0.04) (1.039,9.061)rectangle++(0.04,0.04) (1.243,9.256)rectangle++(0.04,0.04) (1.425,9.414)rectangle++(0.04,0.04)
  (1.586,9.541)rectangle++(0.04,0.04) (1.727,9.643)rectangle++(0.04,0.04) (1.852,9.724)rectangle++(0.04,0.04) (1.961,9.787)rectangle++(0.04,0.04) (2.056,9.837)rectangle++(0.04,0.04)
  (2.139,9.875)rectangle++(0.04,0.04) (2.211,9.905)rectangle++(0.04,0.04) (2.273,9.927)rectangle++(0.04,0.04) (2.327,9.943)rectangle++(0.04,0.04) (2.373,9.955)rectangle++(0.04,0.04)
  (-1.644,5.358)rectangle++(0.04,0.04) (-1.185,6.216)rectangle++(0.04,0.04) (-0.761,6.927)rectangle++(0.04,0.04) (-0.372,7.515)rectangle++(0.04,0.04) (-0.018,7.999)rectangle++(0.04,0.04)
  (0.303,8.396)rectangle++(0.04,0.04) (0.591,8.72)rectangle++(0.04,0.04) (0.849,8.985)rectangle++(0.04,0.04) (1.079,9.199)rectangle++(0.04,0.04) (1.283,9.373)rectangle++(0.04,0.04)
  (1.463,9.512)rectangle++(0.04,0.04) (1.622,9.623)rectangle++(0.04,0.04) (1.761,9.71)rectangle++(0.04,0.04) (1.883,9.78)rectangle++(0.04,0.04) (1.99,9.834)rectangle++(0.04,0.04)
  (2.082,9.875)rectangle++(0.04,0.04) (2.163,9.907)rectangle++(0.04,0.04) (2.233,9.93)rectangle++(0.04,0.04) (2.293,9.948)rectangle++(0.04,0.04) (2.344,9.96)rectangle++(0.04,0.04)
  (2.389,9.969)rectangle++(0.04,0.04) (2.427,9.974)rectangle++(0.04,0.04) (2.46,9.977)rectangle++(0.04,0.04) (2.488,9.978)rectangle++(0.04,0.04) (-1.628,5.384)rectangle++(0.04,0.04)
  (-1.171,6.238)rectangle++(0.04,0.04) (-0.748,6.945)rectangle++(0.04,0.04) (-0.36,7.529)rectangle++(0.04,0.04) (-0.007,8.011)rectangle++(0.04,0.04) (0.312,8.405)rectangle++(0.04,0.04)
  (0.599,8.728)rectangle++(0.04,0.04) (0.856,8.991)rectangle++(0.04,0.04) (1.085,9.205)rectangle++(0.04,0.04) (2.152,3.983)rectangle++(0.04,0.04) (-0.416,4.138)rectangle++(0.04,0.04)
  (-0.191,5.13)rectangle++(0.04,0.04) (0.041,5.964)rectangle++(0.04,0.04) (0.271,6.664)rectangle++(0.04,0.04) (1.618,3.309)rectangle++(0.04,0.04) (1.506,4.344)rectangle++(0.04,0.04)
  (1.452,5.228)rectangle++(0.04,0.04) (1.441,5.982)rectangle++(0.04,0.04) (1.462,6.623)rectangle++(0.04,0.04) (1.505,7.168)rectangle++(0.04,0.04) (1.564,7.629)rectangle++(0.04,0.04)
  (1.632,8.018)rectangle++(0.04,0.04) (1.794,3.879)rectangle++(0.04,0.04) (1.678,4.821)rectangle++(0.04,0.04) (1.617,5.627)rectangle++(0.04,0.04) (1.597,6.315)rectangle++(0.04,0.04)
  (1.608,6.9)rectangle++(0.04,0.04) (1.641,7.397)rectangle++(0.04,0.04) (1.688,7.818)rectangle++(0.04,0.04) (1.746,8.174)rectangle++(0.04,0.04) (-1.378,4.888)rectangle++(0.04,0.04)
  (-0.978,5.807)rectangle++(0.04,0.04) (-0.601,6.571)rectangle++(0.04,0.04) (-0.251,7.206)rectangle++(0.04,0.04) (-1.531,4.323)rectangle++(0.04,0.04) (-1.13,5.332)rectangle++(0.04,0.04)
  (-0.75,6.174)rectangle++(0.04,0.04) (-0.392,6.874)rectangle++(0.04,0.04) (-0.061,7.455)rectangle++(0.04,0.04) (-1.549,4.408)rectangle++(0.04,0.04) (-1.143,5.405)rectangle++(0.04,0.04)
  (-0.757,6.236)rectangle++(0.04,0.04) (-0.397,6.927)rectangle++(0.04,0.04) (-1.498,4.241)rectangle++(0.04,0.04) (-1.106,5.261)rectangle++(0.04,0.04) (-0.732,6.112)rectangle++(0.04,0.04)
  (-0.38,6.821)rectangle++(0.04,0.04) (-0.052,7.409)rectangle++(0.04,0.04) (1.841,3.379)rectangle++(0.04,0.04) (1.698,4.395)rectangle++(0.04,0.04) (1.617,5.264)rectangle++(0.04,0.04)
  (1.582,6.006)rectangle++(0.04,0.04) (1.583,6.638)rectangle++(0.04,0.04) (1.609,7.175)rectangle++(0.04,0.04) (1.652,7.631)rectangle++(0.04,0.04) (1.708,8.016)rectangle++(0.04,0.04)
  (0.462,2.949)rectangle++(0.04,0.04) (0.509,4.085)rectangle++(0.04,0.04) (0.593,5.048)rectangle++(0.04,0.04) (0.704,5.863)rectangle++(0.04,0.04) (0.831,6.552)rectangle++(0.04,0.04)
  (0.966,7.132)rectangle++(0.04,0.04) (-1.297,4.537)rectangle++(0.04,0.04) (-0.923,5.504)rectangle++(0.04,0.04) (-0.567,6.312)rectangle++(0.04,0.04) (1.629,3.051)rectangle++(0.04,0.04)
  (1.505,4.124)rectangle++(0.04,0.04) (1.442,5.041)rectangle++(0.04,0.04) (1.425,5.824)rectangle++(0.04,0.04) (1.275,3.409)rectangle++(0.04,0.04) (1.218,4.443)rectangle++(0.04,0.04)
  (1.211,5.324)rectangle++(0.04,0.04) (1.24,6.073)rectangle++(0.04,0.04) (-1.023,4.379)rectangle++(0.04,0.04) (-0.697,5.359)rectangle++(0.04,0.04) (-0.38,6.179)rectangle++(0.04,0.04)
  (1.575,3.068)rectangle++(0.04,0.04) (1.459,4.141)rectangle++(0.04,0.04) (1.404,5.057)rectangle++(0.04,0.04) (1.393,5.839)rectangle++(0.04,0.04) (1.416,6.503)rectangle++(0.04,0.04)
  (1.461,7.068)rectangle++(0.04,0.04) (1.522,7.545)rectangle++(0.04,0.04) (1.594,7.949)rectangle++(0.04,0.04) (1.67,8.289)rectangle++(0.04,0.04) (1.749,8.575)rectangle++(0.04,0.04)
  (-1.463,4.968)rectangle++(0.04,0.04) (-1.047,5.878)rectangle++(0.04,0.04) (-0.657,6.634)rectangle++(0.04,0.04) (-0.295,7.262)rectangle++(0.04,0.04) (0.037,7.78)rectangle++(0.04,0.04)
  (0.341,8.208)rectangle++(0.04,0.04) (0.616,8.559)rectangle++(0.04,0.04) (-1.663,4.753)rectangle++(0.04,0.04) (-1.225,5.703)rectangle++(0.04,0.04) (-0.816,6.492)rectangle++(0.04,0.04)
  (1.703,3.029)rectangle++(0.04,0.04) (1.567,4.103)rectangle++(0.04,0.04) (1.494,5.021)rectangle++(0.04,0.04) (1.468,5.804)rectangle++(0.04,0.04) (-0.006,1.883)rectangle++(0.04,0.04)
  (0.068,3.197)rectangle++(0.04,0.04) (0.184,4.311)rectangle++(0.04,0.04) (0.326,5.253)rectangle++(0.04,0.04) (0.485,6.048)rectangle++(0.04,0.04) (0.652,6.718)rectangle++(0.04,0.04)
  (1.587,3.406)rectangle++(0.04,0.04) (1.483,4.427)rectangle++(0.04,0.04) (-0.63,4.1)rectangle++(0.04,0.04) (-0.374,5.106)rectangle++(0.04,0.04) (-0.116,5.951)rectangle++(0.04,0.04)
  (0.137,6.66)rectangle++(0.04,0.04) (0.381,7.251)rectangle++(0.04,0.04) (0.612,7.745)rectangle++(0.04,0.04) (0.827,8.155)rectangle++(0.04,0.04) (1.027,8.495)rectangle++(0.04,0.04)
  (1.211,8.775)rectangle++(0.04,0.04) (1.378,9.007)rectangle++(0.04,0.04) (1.529,9.197)rectangle++(0.04,0.04) (1.666,9.352)rectangle++(0.04,0.04) (1.788,9.479)rectangle++(0.04,0.04)
  (1.897,9.582)rectangle++(0.04,0.04) (1.993,9.665)rectangle++(0.04,0.04) (2.079,9.732)rectangle++(0.04,0.04) (2.154,9.785)rectangle++(0.04,0.04) (2.22,9.827)rectangle++(0.04,0.04)
  (2.278,9.861)rectangle++(0.04,0.04) (2.328,9.887)rectangle++(0.04,0.04) (2.372,9.907)rectangle++(0.04,0.04) (2.411,9.922)rectangle++(0.04,0.04) (2.444,9.933)rectangle++(0.04,0.04)
  (2.472,9.942)rectangle++(0.04,0.04) (2.497,9.948)rectangle++(0.04,0.04) (2.518,9.952)rectangle++(0.04,0.04) (2.536,9.955)rectangle++(0.04,0.04) (2.552,9.956)rectangle++(0.04,0.04)
  (2.167,4.461)rectangle++(0.04,0.04) (0.405,2.566)rectangle++(0.04,0.04) (-0.429,3.53)rectangle++(0.04,0.04) (0.924,2.546)rectangle++(0.04,0.04) (0.885,3.723)rectangle++(0.04,0.04)
  (0.899,4.725)rectangle++(0.04,0.04) (0.951,5.577)rectangle++(0.04,0.04) (1.029,6.298)rectangle++(0.04,0.04) (1.439,3.41)rectangle++(0.04,0.04) (1.358,4.437)rectangle++(0.04,0.04)
  (1.329,5.314)rectangle++(0.04,0.04) (1.34,6.06)rectangle++(0.04,0.04) (1.379,6.693)rectangle++(0.04,0.04) (1.438,7.23)rectangle++(0.04,0.04) (1.509,7.685)rectangle++(0.04,0.04)
  (-1.317,4.748)rectangle++(0.04,0.04) (-0.932,5.684)rectangle++(0.04,0.04) (-0.567,6.465)rectangle++(0.04,0.04) (0.234,2.742)rectangle++(0.04,0.04) (0.307,3.918)rectangle++(0.04,0.04)
  (0.577,2.134)rectangle++(0.04,0.04) (0.574,3.387)rectangle++(0.04,0.04) (0.621,4.452)rectangle++(0.04,0.04) (0.704,5.356)rectangle++(0.04,0.04) (0.81,6.121)rectangle++(0.04,0.04)
  (0.931,6.766)rectangle++(0.04,0.04) (1.06,7.31)rectangle++(0.04,0.04) (1.191,7.768)rectangle++(0.04,0.04) (1.321,8.151)rectangle++(0.04,0.04) (-1.45,4.804)rectangle++(0.04,0.04)
  (-1.042,5.738)rectangle++(0.04,0.04) (-0.658,6.515)rectangle++(0.04,0.04) (-0.301,7.16)rectangle++(0.04,0.04) (0.028,7.695)rectangle++(0.04,0.04) (0.329,8.135)rectangle++(0.04,0.04)
  (0.603,8.498)rectangle++(0.04,0.04) (0.851,8.796)rectangle++(0.04,0.04) (1.073,9.038)rectangle++(0.04,0.04) (1.271,9.236)rectangle++(0.04,0.04) (1.448,9.396)rectangle++(0.04,0.04)
  (1.604,9.525)rectangle++(0.04,0.04) (1.742,9.628)rectangle++(0.04,0.04) (1.864,9.71)rectangle++(0.04,0.04) (1.97,9.775)rectangle++(0.04,0.04) (2.064,9.827)rectangle++(0.04,0.04)
  (2.145,9.866)rectangle++(0.04,0.04) (2.216,9.897)rectangle++(0.04,0.04) (-1.66,5.317)rectangle++(0.04,0.04) (0.09,3.226)rectangle++(0.04,0.04) (0.203,4.335)rectangle++(0.04,0.04)
  (0.344,5.273)rectangle++(0.04,0.04) (0.501,6.064)rectangle++(0.04,0.04) (0.666,6.731)rectangle++(0.04,0.04) (0.833,7.291)rectangle++(0.04,0.04) (0.998,7.76)rectangle++(0.04,0.04)
  (1.156,8.152)rectangle++(0.04,0.04) (0.421,2.955)rectangle++(0.04,0.04) (0.434,2.255)rectangle++(0.04,0.04) (0.457,3.495)rectangle++(0.04,0.04) (0.812,2.74)rectangle++(0.04,0.04)
  (0.798,3.893)rectangle++(0.04,0.04) (0.832,4.873)rectangle++(0.04,0.04) (-0.842,4.066)rectangle++(0.04,0.04) (-0.555,5.086)rectangle++(0.04,0.04) (-0.271,5.941)rectangle++(0.04,0.04)
  (0.005,6.657)rectangle++(0.04,0.04) (0.269,7.254)rectangle++(0.04,0.04) (0.516,7.752)rectangle++(0.04,0.04) (0.747,8.165)rectangle++(0.04,0.04) (-1.555,4.699)rectangle++(0.04,0.04)
  (0.106,2.532)rectangle++(0.04,0.04) (0.19,3.745)rectangle++(0.04,0.04) (0.309,4.771)rectangle++(0.04,0.04) (0.451,5.64)rectangle++(0.04,0.04) (-1.073,4.146)rectangle++(0.04,0.04)
  (-0.748,5.163)rectangle++(0.04,0.04) (-0.432,6.015)rectangle++(0.04,0.04) (-0.128,6.726)rectangle++(0.04,0.04) (0.158,7.318)rectangle++(0.04,0.04) (0.425,7.81)rectangle++(0.04,0.04)
  (0.671,8.218)rectangle++(0.04,0.04) (1.956,3.699)rectangle++(0.04,0.04) (1.809,4.662)rectangle++(0.04,0.04) (1.722,5.487)rectangle++(0.04,0.04) (1.681,6.191)rectangle++(0.04,0.04)
  (-0.969,4.485)rectangle++(0.04,0.04) (-0.647,5.447)rectangle++(0.04,0.04) (-0.334,6.252)rectangle++(0.04,0.04) (-0.036,6.924)rectangle++(0.04,0.04) (0.244,7.483)rectangle++(0.04,0.04)
  (-0.005,1.894)rectangle++(0.04,0.04) (0.069,3.206)rectangle++(0.04,0.04) (0.185,4.319)rectangle++(0.04,0.04) (1.048,2.837)rectangle++(0.04,0.04) (1.002,3.966)rectangle++(0.04,0.04)
  (1.008,4.927)rectangle++(0.04,0.04) (1.052,5.744)rectangle++(0.04,0.04) (1.122,6.436)rectangle++(0.04,0.04) (1.209,7.022)rectangle++(0.04,0.04) (1.306,7.517)rectangle++(0.04,0.04)
  (1.409,7.933)rectangle++(0.04,0.04) (1.512,8.283)rectangle++(0.04,0.04) (1.615,8.576)rectangle++(0.04,0.04) (1.713,8.821)rectangle++(0.04,0.04) (1.807,9.026)rectangle++(0.04,0.04)
  (1.895,9.196)rectangle++(0.04,0.04) (1.976,9.337)rectangle++(0.04,0.04) (2.051,9.454)rectangle++(0.04,0.04) (-0.226,4.021)rectangle++(0.04,0.04) (1.022,2.686)rectangle++(0.04,0.04)
  (0.974,3.839)rectangle++(0.04,0.04) (0.979,4.82)rectangle++(0.04,0.04) (1.023,5.654)rectangle++(0.04,0.04) (1.094,6.361)rectangle++(0.04,0.04) (1.182,6.96)rectangle++(0.04,0.04)
  (1.281,7.465)rectangle++(0.04,0.04) (-1.311,4.661)rectangle++(0.04,0.04) (-0.93,5.611)rectangle++(0.04,0.04) (-0.568,6.403)rectangle++(0.04,0.04) (-0.229,7.061)rectangle++(0.04,0.04)
  (0.085,7.607)rectangle++(0.04,0.04) (0.375,8.059)rectangle++(0.04,0.04) (0.639,8.431)rectangle++(0.04,0.04) (0.878,8.737)rectangle++(0.04,0.04) (1.093,8.988)rectangle++(0.04,0.04)
  (1.287,9.192)rectangle++(0.04,0.04) (1.459,9.358)rectangle++(0.04,0.04) (-1.681,5.069)rectangle++(0.04,0.04) (-1.228,5.972)rectangle++(0.04,0.04) (-1.444,3.896)rectangle++(0.04,0.04)
  (-1.074,4.966)rectangle++(0.04,0.04) (-0.716,5.86)rectangle++(0.04,0.04) (-0.377,6.606)rectangle++(0.04,0.04) (-0.058,7.226)rectangle++(0.04,0.04) (0.237,7.741)rectangle++(0.04,0.04)
  (0.509,8.166)rectangle++(0.04,0.04) (0.757,8.517)rectangle++(0.04,0.04) (2.02,3.776)rectangle++(0.04,0.04) (1.866,4.725)rectangle++(0.04,0.04) (1.773,5.538)rectangle++(0.04,0.04)
  (1.726,6.233)rectangle++(0.04,0.04) (1.715,6.825)rectangle++(0.04,0.04) (1.728,7.329)rectangle++(0.04,0.04) (1.76,7.757)rectangle++(0.04,0.04) (1.715,3.857)rectangle++(0.04,0.04)
  (0.352,2.493)rectangle++(0.04,0.04) (0.397,3.701)rectangle++(0.04,0.04) (0.004,2.086)rectangle++(0.04,0.04) (0.084,3.369)rectangle++(0.04,0.04) (-0.656,3.629)rectangle++(0.04,0.04)
  (-0.415,4.707)rectangle++(0.04,0.04) (-0.167,5.614)rectangle++(0.04,0.04) (0.081,6.374)rectangle++(0.04,0.04) (0.321,7.011)rectangle++(0.04,0.04) (0.551,7.543)rectangle++(0.04,0.04)
  (0.768,7.986)rectangle++(0.04,0.04) (0.97,8.353)rectangle++(0.04,0.04) (1.157,8.658)rectangle++(0.04,0.04) (1.327,8.909)rectangle++(0.04,0.04) (1.482,9.116)rectangle++(0.04,0.04)
  (1.622,9.285)rectangle++(0.04,0.04) (1.748,9.424)rectangle++(0.04,0.04) (1.861,9.536)rectangle++(0.04,0.04) (1.961,9.628)rectangle++(0.04,0.04) (2.05,9.701)rectangle++(0.04,0.04)
  (2.128,9.76)rectangle++(0.04,0.04) (2.197,9.807)rectangle++(0.04,0.04) (2.258,9.845)rectangle++(0.04,0.04) (2.311,9.874)rectangle++(0.04,0.04) (2.357,9.897)rectangle++(0.04,0.04)
  (2.397,9.914)rectangle++(0.04,0.04) (2.432,9.927)rectangle++(0.04,0.04) (2.462,9.937)rectangle++(0.04,0.04) (2.488,9.944)rectangle++(0.04,0.04) (2.51,9.949)rectangle++(0.04,0.04)
  (2.529,9.953)rectangle++(0.04,0.04) (2.546,9.955)rectangle++(0.04,0.04) (2.56,9.956)rectangle++(0.04,0.04) (2.166,4.463)rectangle++(0.04,0.04) (2.017,5.303)rectangle++(0.04,0.04)
  (1.924,6.023)rectangle++(0.04,0.04) (1.874,6.639)rectangle++(0.04,0.04) (1.857,7.164)rectangle++(0.04,0.04) (-1.144,4.71)rectangle++(0.04,0.04) (-0.786,5.646)rectangle++(0.04,0.04)
  (-0.445,6.426)rectangle++(0.04,0.04) (1.643,3.1)rectangle++(0.04,0.04) (1.519,4.166)rectangle++(0.04,0.04) (1.455,5.076)rectangle++(0.04,0.04) (1.438,5.853)rectangle++(0.04,0.04)
  (1.454,6.514)rectangle++(0.04,0.04) (1.494,7.075)rectangle++(0.04,0.04) (1.551,7.55)rectangle++(0.04,0.04) (0.003,2.073)rectangle++(0.04,0.04) (0.083,3.358)rectangle++(0.04,0.04)
  (0.203,4.447)rectangle++(0.04,0.04) (0.348,5.368)rectangle++(0.04,0.04) (0.509,6.145)rectangle++(0.04,0.04) (0.676,6.799)rectangle++(0.04,0.04) (0.844,7.349)rectangle++(0.04,0.04)
  (1.009,7.809)rectangle++(0.04,0.04) (1.168,8.193)rectangle++(0.04,0.04) (1.318,8.514)rectangle++(0.04,0.04) (1.459,8.78)rectangle++(0.04,0.04) (1.589,9.001)rectangle++(0.04,0.04)
  (1.709,9.183)rectangle++(0.04,0.04) (1.817,9.334)rectangle++(0.04,0.04) (1.916,9.457)rectangle++(0.04,0.04) (2.005,9.558)rectangle++(0.04,0.04) (2.084,9.64)rectangle++(0.04,0.04)
  (2.155,9.707)rectangle++(0.04,0.04) (2.218,9.761)rectangle++(0.04,0.04) (2.273,9.804)rectangle++(0.04,0.04) (2.322,9.839)rectangle++(0.04,0.04) (2.365,9.867)rectangle++(0.04,0.04)
  (2.403,9.888)rectangle++(0.04,0.04) (2.436,9.905)rectangle++(0.04,0.04) (2.464,9.918)rectangle++(0.04,0.04) (2.489,9.928)rectangle++(0.04,0.04) (-1.617,5.374)rectangle++(0.04,0.04)
  (-1.162,6.229)rectangle++(0.04,0.04) (-0.741,6.937)rectangle++(0.04,0.04) (-0.354,7.523)rectangle++(0.04,0.04) (-0.002,8.005)rectangle++(0.04,0.04) (0.316,8.4)rectangle++(0.04,0.04)
  (0.602,8.724)rectangle++(0.04,0.04) (-0.245,3.351)rectangle++(0.04,0.04) (-0.076,4.454)rectangle++(0.04,0.04) (0.009,2.189)rectangle++(0.04,0.04) (0.093,3.456)rectangle++(0.04,0.04)
  (0.215,4.53)rectangle++(0.04,0.04) (0.362,5.438)rectangle++(0.04,0.04) (0.523,6.204)rectangle++(0.04,0.04) (0.69,6.849)rectangle++(0.04,0.04) (0.858,7.39)rectangle++(0.04,0.04)
  (1.023,7.844)rectangle++(0.04,0.04) (1.181,8.222)rectangle++(0.04,0.04) (1.331,8.538)rectangle++(0.04,0.04) (1.958,3.911)rectangle++(0.04,0.04) (1.819,4.842)rectangle++(0.04,0.04)
  (1.737,5.639)rectangle++(0.04,0.04) (1.7,6.32)rectangle++(0.04,0.04) (1.696,6.9)rectangle++(0.04,0.04) (1.715,7.393)rectangle++(0.04,0.04) (1.751,7.812)rectangle++(0.04,0.04)
  (1.799,8.166)rectangle++(0.04,0.04) (1.854,8.466)rectangle++(0.04,0.04) (1.912,8.718)rectangle++(0.04,0.04) (1.935,4.077)rectangle++(0.04,0.04) (1.806,4.984)rectangle++(0.04,0.04)
  (1.732,5.761)rectangle++(0.04,0.04) (1.7,6.424)rectangle++(0.04,0.04) (1.7,6.988)rectangle++(0.04,0.04) (1.722,7.468)rectangle++(0.04,0.04) (-1.233,4.745)rectangle++(0.04,0.04)
  (-0.86,5.679)rectangle++(0.04,0.04) (-0.506,6.457)rectangle++(0.04,0.04) (-0.174,7.105)rectangle++(0.04,0.04) (0.134,7.643)rectangle++(0.04,0.04) (0.417,8.087)rectangle++(0.04,0.04)
  (0.676,8.454)rectangle++(0.04,0.04) (0.911,8.755)rectangle++(0.04,0.04) (1.122,9.001)rectangle++(0.04,0.04) (1.312,9.202)rectangle++(0.04,0.04) (1.481,9.366)rectangle++(0.04,0.04)
  (1.631,9.498)rectangle++(0.04,0.04) (1.764,9.604)rectangle++(0.04,0.04) (1.881,9.689)rectangle++(0.04,0.04) (1.985,9.757)rectangle++(0.04,0.04) (-1.671,5.246)rectangle++(0.04,0.04)
  (-1.213,6.122)rectangle++(0.04,0.04) (-0.788,6.848)rectangle++(0.04,0.04) (-0.398,7.449)rectangle++(0.04,0.04) (-0.043,7.944)rectangle++(0.04,0.04) (0.279,8.35)rectangle++(0.04,0.04)
  (0.569,8.683)rectangle++(0.04,0.04) (0.829,8.954)rectangle++(0.04,0.04) (1.06,9.174)rectangle++(0.04,0.04) (1.266,9.351)rectangle++(0.04,0.04) (1.448,9.494)rectangle++(0.04,0.04)
  (1.608,9.608)rectangle++(0.04,0.04) (1.749,9.699)rectangle++(0.04,0.04) (1.873,9.77)rectangle++(0.04,0.04) (1.98,9.826)rectangle++(0.04,0.04) (2.074,9.869)rectangle++(0.04,0.04)
  (2.156,9.902)rectangle++(0.04,0.04) (2.226,9.927)rectangle++(0.04,0.04) (2.287,9.945)rectangle++(0.04,0.04) (2.34,9.958)rectangle++(0.04,0.04) (2.192,4.413)rectangle++(0.04,0.04)
  (2.037,5.26)rectangle++(0.04,0.04) (1.94,5.986)rectangle++(0.04,0.04) (1.255,3.557)rectangle++(0.04,0.04) (1.207,4.569)rectangle++(0.04,0.04) (1.206,5.432)rectangle++(0.04,0.04)
  (1.241,6.165)rectangle++(0.04,0.04) (1.299,6.787)rectangle++(0.04,0.04) (1.373,7.313)rectangle++(0.04,0.04) (1.458,7.757)rectangle++(0.04,0.04) (1.547,8.132)rectangle++(0.04,0.04)
  (1.638,8.446)rectangle++(0.04,0.04) (1.728,8.71)rectangle++(0.04,0.04) (1.815,8.931)rectangle++(0.04,0.04) (1.898,9.115)rectangle++(0.04,0.04) (1.975,9.268)rectangle++(0.04,0.04)
  (2.048,9.395)rectangle++(0.04,0.04) (2.114,9.5)rectangle++(0.04,0.04) (2.175,9.586)rectangle++(0.04,0.04) (2.23,9.658)rectangle++(0.04,0.04) (2.279,9.716)rectangle++(0.04,0.04)
  (2.324,9.764)rectangle++(0.04,0.04) (2.364,9.802)rectangle++(0.04,0.04) (2.399,9.834)rectangle++(0.04,0.04) (2.43,9.859)rectangle++(0.04,0.04) (0.008,2.17)rectangle++(0.04,0.04)
  (0.091,3.44)rectangle++(0.04,0.04) (0.213,4.517)rectangle++(0.04,0.04) (0.359,5.427)rectangle++(0.04,0.04) (0.52,6.195)rectangle++(0.04,0.04) (0.688,6.841)rectangle++(0.04,0.04)
  (-1.285,4.427)rectangle++(0.04,0.04) (-0.917,5.411)rectangle++(0.04,0.04) (-0.566,6.232)rectangle++(0.04,0.04) (-0.234,6.916)rectangle++(0.04,0.04) (0.076,7.484)rectangle++(0.04,0.04)
  (0.362,7.955)rectangle++(0.04,0.04) (0.623,8.343)rectangle++(0.04,0.04) (0.861,8.663)rectangle++(0.04,0.04) (1.077,8.925)rectangle++(0.04,0.04) (1.27,9.14)rectangle++(0.04,0.04)
  (1.443,9.314)rectangle++(0.04,0.04) (1.597,9.455)rectangle++(0.04,0.04) (1.733,9.569)rectangle++(0.04,0.04) (-1.677,5.162)rectangle++(0.04,0.04) (-1.221,6.051)rectangle++(0.04,0.04)
  (-0.798,6.788)rectangle++(0.04,0.04) (-0.409,7.398)rectangle++(0.04,0.04) (-0.054,7.901)rectangle++(0.04,0.04) (0.268,8.314)rectangle++(0.04,0.04) (0.558,8.652)rectangle++(0.04,0.04)
  (0.818,8.928)rectangle++(0.04,0.04) (1.051,9.153)rectangle++(0.04,0.04) (1.257,9.334)rectangle++(0.04,0.04) (2.164,4.047)rectangle++(0.04,0.04) (1.999,4.95)rectangle++(0.04,0.04)
  (1.895,5.724)rectangle++(0.04,0.04) (1.838,6.385)rectangle++(0.04,0.04) (1.815,6.95)rectangle++(0.04,0.04) (1.819,7.431)rectangle++(0.04,0.04) (1.841,7.84)rectangle++(0.04,0.04)
  (1.876,8.187)rectangle++(0.04,0.04) (1.92,8.48)rectangle++(0.04,0.04) (1.969,8.727)rectangle++(0.04,0.04) (2.021,8.936)rectangle++(0.04,0.04) (2.073,9.111)rectangle++(0.04,0.04)
  (2.124,9.257)rectangle++(0.04,0.04) (2.174,9.38)rectangle++(0.04,0.04) (2.22,9.482)rectangle++(0.04,0.04) (2.265,9.567)rectangle++(0.04,0.04) (2.305,9.638)rectangle++(0.04,0.04)
  (2.343,9.696)rectangle++(0.04,0.04) (2.377,9.744)rectangle++(0.04,0.04) (2.408,9.784)rectangle++(0.04,0.04) (2.436,9.816)rectangle++(0.04,0.04) (-1.245,3.407)rectangle++(0.04,0.04)
  (-0.924,4.542)rectangle++(0.04,0.04) (1.233,2.628)rectangle++(0.04,0.04) (1.151,3.781)rectangle++(0.04,0.04) (1.127,4.764)rectangle++(0.04,0.04) (1.147,5.6)rectangle++(0.04,0.04)
  (1.196,6.311)rectangle++(0.04,0.04) (0.021,2.441)rectangle++(0.04,0.04) (0.113,3.67)rectangle++(0.04,0.04) (0.241,4.711)rectangle++(0.04,0.04) (0.391,5.591)rectangle++(0.04,0.04)
  (0.554,6.333)rectangle++(0.04,0.04) (0.722,6.957)rectangle++(0.04,0.04) (0.889,7.481)rectangle++(0.04,0.04) (-1.007,3.291)rectangle++(0.04,0.04) (-0.726,4.434)rectangle++(0.04,0.04)
  (-0.442,5.394)rectangle++(0.04,0.04) (-0.162,6.199)rectangle++(0.04,0.04) (0.108,6.872)rectangle++(0.04,0.04) (0.364,7.433)rectangle++(0.04,0.04) (0.605,7.9)rectangle++(0.04,0.04)
  (0.828,8.287)rectangle++(0.04,0.04) (1.033,8.607)rectangle++(0.04,0.04) (1.22,8.871)rectangle++(0.04,0.04) (1.39,9.087)rectangle++(0.04,0.04) (1.543,9.265)rectangle++(0.04,0.04)
  (1.68,9.41)rectangle++(0.04,0.04) (1.802,9.527)rectangle++(0.04,0.04) (1.91,9.622)rectangle++(0.04,0.04) (2.007,9.699)rectangle++(0.04,0.04) (2.091,9.76)rectangle++(0.04,0.04)
  (2.166,9.808)rectangle++(0.04,0.04) (-1.65,5.29)rectangle++(0.04,0.04) (-1.193,6.159)rectangle++(0.04,0.04) (-0.77,6.879)rectangle++(0.04,0.04) (-0.381,7.474)rectangle++(0.04,0.04)
  (-0.027,7.964)rectangle++(0.04,0.04) (0.293,8.367)rectangle++(0.04,0.04) (0.582,8.696)rectangle++(0.04,0.04) (0.84,8.965)rectangle++(0.04,0.04) (1.07,9.183)rectangle++(0.04,0.04)
  (1.275,9.359)rectangle++(0.04,0.04) (1.456,9.5)rectangle++(0.04,0.04) (1.615,9.613)rectangle++(0.04,0.04) (1.755,9.703)rectangle++(0.04,0.04) (1.878,9.773)rectangle++(0.04,0.04)
  (1.985,9.828)rectangle++(0.04,0.04) (2.078,9.871)rectangle++(0.04,0.04) (2.201,4.337)rectangle++(0.04,0.04) (2.042,5.195)rectangle++(0.04,0.04) (1.941,5.93)rectangle++(0.04,0.04)
  (1.885,6.559)rectangle++(0.04,0.04) (1.863,7.096)rectangle++(0.04,0.04) (1.539,3.753)rectangle++(0.04,0.04) (1.456,4.725)rectangle++(0.04,0.04) (1.425,5.554)rectangle++(0.04,0.04)
  (1.431,6.26)rectangle++(0.04,0.04) (1.464,6.86)rectangle++(0.04,0.04) (1.517,7.369)rectangle++(0.04,0.04) (1.582,7.799)rectangle++(0.04,0.04) (1.654,8.162)rectangle++(0.04,0.04)
  (1.731,8.468)rectangle++(0.04,0.04) (1.807,8.725)rectangle++(0.04,0.04) (1.883,8.94)rectangle++(0.04,0.04) (1.956,9.12)rectangle++(0.04,0.04) (-1.541,5.115)rectangle++(0.04,0.04)
  (-1.108,6.005)rectangle++(0.04,0.04) (-0.704,6.745)rectangle++(0.04,0.04) (-0.33,7.358)rectangle++(0.04,0.04) (0.011,7.863)rectangle++(0.04,0.04) (0.322,8.28)rectangle++(0.04,0.04)
  (-1.656,4.642)rectangle++(0.04,0.04) (-1.224,5.609)rectangle++(0.04,0.04) (-0.818,6.412)rectangle++(0.04,0.04) (-1.464,4.06)rectangle++(0.04,0.04) (-1.084,5.106)rectangle++(0.04,0.04)
  (-0.72,5.98)rectangle++(0.04,0.04) (-0.375,6.708)rectangle++(0.04,0.04) (-0.052,7.313)rectangle++(0.04,0.04) (0.246,7.814)rectangle++(0.04,0.04) (0.519,8.228)rectangle++(0.04,0.04)
  (0.768,8.569)rectangle++(0.04,0.04) (0.994,8.85)rectangle++(0.04,0.04) (1.196,9.079)rectangle++(0.04,0.04) (1.378,9.265)rectangle++(0.04,0.04) (1.54,9.416)rectangle++(0.04,0.04)
  (1.683,9.539)rectangle++(0.04,0.04) (1.81,9.637)rectangle++(0.04,0.04) (1.922,9.715)rectangle++(0.04,0.04) (2.02,9.777)rectangle++(0.04,0.04) (2.106,9.826)rectangle++(0.04,0.04)
  (2.181,9.864)rectangle++(0.04,0.04) (2.246,9.893)rectangle++(0.04,0.04) (2.303,9.916)rectangle++(0.04,0.04) (2.186,4.397)rectangle++(0.04,0.04) (2.032,5.246)rectangle++(0.04,0.04)
  (1.934,5.974)rectangle++(0.04,0.04) (1.881,6.597)rectangle++(0.04,0.04) (1.861,7.128)rectangle++(0.04,0.04) (1.864,7.581)rectangle++(0.04,0.04) (1.886,7.965)rectangle++(0.04,0.04)
  (1.919,8.291)rectangle++(0.04,0.04) (1.961,8.567)rectangle++(0.04,0.04) (2.007,8.8)rectangle++(0.04,0.04) (-1.464,5.061)rectangle++(0.04,0.04) (-1.044,5.957)rectangle++(0.04,0.04)
  (0.182,2.679)rectangle++(0.04,0.04) (0.259,3.866)rectangle++(0.04,0.04) (0.373,4.872)rectangle++(0.04,0.04) (-0.924,3.979)rectangle++(0.04,0.04) (-0.628,5.016)rectangle++(0.04,0.04)
  (-0.335,5.885)rectangle++(0.04,0.04) (-0.052,6.612)rectangle++(0.04,0.04) (0.218,7.218)rectangle++(0.04,0.04) (0.472,7.723)rectangle++(0.04,0.04) (0.708,8.142)rectangle++(0.04,0.04)
  (-1.557,4.688)rectangle++(0.04,0.04) (-1.139,5.643)rectangle++(0.04,0.04) (-0.744,6.438)rectangle++(0.04,0.04) (-1.456,4.079)rectangle++(0.04,0.04) (-1.077,5.122)rectangle++(0.04,0.04)
  (-0.713,5.993)rectangle++(0.04,0.04) (-0.368,6.718)rectangle++(0.04,0.04) (-0.046,7.322)rectangle++(0.04,0.04) (0.251,7.821)rectangle++(0.04,0.04) (0.524,8.234)rectangle++(0.04,0.04)
  (0.773,8.574)rectangle++(0.04,0.04) (0.998,8.854)rectangle++(0.04,0.04) (1.2,9.082)rectangle++(0.04,0.04) (1.381,9.268)rectangle++(0.04,0.04) (1.542,9.419)rectangle++(0.04,0.04)
  (1.686,9.54)rectangle++(0.04,0.04) (-1.679,5.147)rectangle++(0.04,0.04) (-1.224,6.038)rectangle++(0.04,0.04) (-0.801,6.778)rectangle++(0.04,0.04) (-0.412,7.389)rectangle++(0.04,0.04)
  (-0.057,7.894)rectangle++(0.04,0.04) (-1.642,4.495)rectangle++(0.04,0.04) (-1.218,5.483)rectangle++(0.04,0.04) (-0.818,6.305)rectangle++(0.04,0.04) (-0.445,6.988)rectangle++(0.04,0.04)
  (-0.101,7.554)rectangle++(0.04,0.04) (0.214,8.021)rectangle++(0.04,0.04) (0.5,8.406)rectangle++(0.04,0.04) (-1.651,4.701)rectangle++(0.04,0.04) (1.357,2.494)rectangle++(0.04,0.04)
  (1.251,3.662)rectangle++(0.04,0.04) (-0.509,3.905)rectangle++(0.04,0.04) (-0.279,4.936)rectangle++(0.04,0.04) (-0.042,5.803)rectangle++(0.04,0.04) (0.194,6.53)rectangle++(0.04,0.04)
  (0.424,7.139)rectangle++(0.04,0.04) (1.718,3.435)rectangle++(0.04,0.04) (1.596,4.447)rectangle++(0.04,0.04) (1.532,5.313)rectangle++(0.04,0.04) (1.513,6.051)rectangle++(0.04,0.04)
  (1.525,6.679)rectangle++(0.04,0.04) (1.562,7.212)rectangle++(0.04,0.04) (1.614,7.664)rectangle++(0.04,0.04) (1.676,8.046)rectangle++(0.04,0.04) (-1.364,4.85)rectangle++(0.04,0.04)
  (-0.967,5.773)rectangle++(0.04,0.04) (1.543,2.856)rectangle++(0.04,0.04) (1.424,3.962)rectangle++(0.04,0.04) (1.366,4.907)rectangle++(0.04,0.04) (1.355,5.712)rectangle++(0.04,0.04)
  (1.378,6.398)rectangle++(0.04,0.04) (1.425,6.979)rectangle++(0.04,0.04) (1.489,7.471)rectangle++(0.04,0.04) (1.562,7.887)rectangle++(0.04,0.04) (1.641,8.238)rectangle++(0.04,0.04)
  (1.722,8.533)rectangle++(0.04,0.04) (1.803,8.78)rectangle++(0.04,0.04) (1.881,8.987)rectangle++(0.04,0.04) (-1.527,5.075)rectangle++(0.04,0.04) (-1.097,5.971)rectangle++(0.04,0.04)
  (-0.696,6.715)rectangle++(0.04,0.04) (-0.325,7.332)rectangle++(0.04,0.04) (0.015,7.841)rectangle++(0.04,0.04) (0.324,8.261)rectangle++(0.04,0.04) (0.604,8.605)rectangle++(0.04,0.04)
  (0.855,8.886)rectangle++(0.04,0.04) (1.08,9.115)rectangle++(0.04,0.04) (1.281,9.301)rectangle++(0.04,0.04) (1.458,9.451)rectangle++(0.04,0.04) (1.615,9.571)rectangle++(0.04,0.04)
  (-0.191,4.023)rectangle++(0.04,0.04) (-0.003,5.023)rectangle++(0.04,0.04) (0.196,5.866)rectangle++(0.04,0.04) (0.399,6.574)rectangle++(0.04,0.04) (0.6,7.169)rectangle++(0.04,0.04)
  (1.705,3.481)rectangle++(0.04,0.04) (1.586,4.487)rectangle++(0.04,0.04) (1.526,5.346)rectangle++(0.04,0.04) (1.508,6.08)rectangle++(0.04,0.04) (1.329,3.477)rectangle++(0.04,0.04)
  (1.266,4.498)rectangle++(0.04,0.04) (1.254,5.369)rectangle++(0.04,0.04) (1.279,6.11)rectangle++(0.04,0.04) (1.329,6.738)rectangle++(0.04,0.04) (-0.05,3.135)rectangle++(0.04,0.04)
  (0.081,4.263)rectangle++(0.04,0.04) (0.229,2.511)rectangle++(0.04,0.04) (0.293,3.721)rectangle++(0.04,0.04) (-0.694,3.737)rectangle++(0.04,0.04) (-0.442,4.801)rectangle++(0.04,0.04)
  (-0.186,5.694)rectangle++(0.04,0.04) (0.067,6.444)rectangle++(0.04,0.04) (1.588,3.22)rectangle++(0.04,0.04) (1.476,4.27)rectangle++(0.04,0.04) (1.423,5.166)rectangle++(0.04,0.04)
  (1.414,5.931)rectangle++(0.04,0.04) (1.437,6.581)rectangle++(0.04,0.04) (1.483,7.132)rectangle++(0.04,0.04) (1.543,7.599)rectangle++(0.04,0.04) (1.614,7.994)rectangle++(0.04,0.04)
  (1.689,8.326)rectangle++(0.04,0.04) (1.767,8.606)rectangle++(0.04,0.04) (1.844,8.841)rectangle++(0.04,0.04) (1.919,9.037)rectangle++(0.04,0.04) (1.99,9.201)rectangle++(0.04,0.04)
  (2.057,9.337)rectangle++(0.04,0.04) (2.12,9.451)rectangle++(0.04,0.04) (2.178,9.545)rectangle++(0.04,0.04) (2.231,9.622)rectangle++(0.04,0.04) (2.279,9.686)rectangle++(0.04,0.04)
  (2.322,9.738)rectangle++(0.04,0.04) (2.361,9.78)rectangle++(0.04,0.04) (2.396,9.815)rectangle++(0.04,0.04) (2.427,9.843)rectangle++(0.04,0.04) (2.455,9.866)rectangle++(0.04,0.04)
  (2.479,9.884)rectangle++(0.04,0.04) (-1.61,5.364)rectangle++(0.04,0.04) (-1.156,6.22)rectangle++(0.04,0.04) (-0.736,6.929)rectangle++(0.04,0.04) (-0.351,7.515)rectangle++(0.04,0.04)
  (0,7.998)rectangle++(0.04,0.04) (0.031,2.67)rectangle++(0.04,0.04) (0.131,3.865)rectangle++(0.04,0.04) (0.264,4.876)rectangle++(0.04,0.04) (0.417,5.73)rectangle++(0.04,0.04)
  (0.582,6.45)rectangle++(0.04,0.04) (0.75,7.056)rectangle++(0.04,0.04) (0.918,7.564)rectangle++(0.04,0.04) (1.08,7.988)rectangle++(0.04,0.04) (1.236,8.343)rectangle++(0.04,0.04)
  (1.382,8.638)rectangle++(0.04,0.04) (1.518,8.884)rectangle++(0.04,0.04) (1.643,9.087)rectangle++(0.04,0.04) (1.758,9.254)rectangle++(0.04,0.04) (2.086,4.146)rectangle++(0.04,0.04)
  (1.936,5.037)rectangle++(0.04,0.04) (1.845,5.8)rectangle++(0.04,0.04) (1.798,6.452)rectangle++(0.04,0.04) (1.387,3.615)rectangle++(0.04,0.04) (1.322,4.613)rectangle++(0.04,0.04)
  (1.306,5.464)rectangle++(0.04,0.04) (1.326,6.189)rectangle++(0.04,0.04) (1.373,6.804)rectangle++(0.04,0.04) (1.437,7.324)rectangle++(0.04,0.04) (-1.253,4.663)rectangle++(0.04,0.04)
  (-0.881,5.61)rectangle++(0.04,0.04) (-0.527,6.4)rectangle++(0.04,0.04) (-1.409,4.112)rectangle++(0.04,0.04) (-1.035,5.148)rectangle++(0.04,0.04) (-1.235,3.77)rectangle++(0.04,0.04)
  (1.078,2.409)rectangle++(0.04,0.04) (1.011,3.601)rectangle++(0.04,0.04) (0.773,2.886)rectangle++(0.04,0.04) (0.771,4.018)rectangle++(0.04,0.04) (0.814,4.981)rectangle++(0.04,0.04)
  (0.889,5.798)rectangle++(0.04,0.04) (0.985,6.489)rectangle++(0.04,0.04) (1.095,7.072)rectangle++(0.04,0.04) (1.211,7.564)rectangle++(0.04,0.04) (1.33,7.977)rectangle++(0.04,0.04)
  (1.447,8.323)rectangle++(0.04,0.04) (1.561,8.613)rectangle++(0.04,0.04) (1.669,8.855)rectangle++(0.04,0.04) (1.771,9.056)rectangle++(0.04,0.04) (1.865,9.223)rectangle++(0.04,0.04)
  (1.952,9.361)rectangle++(0.04,0.04) (2.032,9.475)rectangle++(0.04,0.04) (1.269,2.289)rectangle++(0.04,0.04) (1.168,3.491)rectangle++(0.04,0.04) (1.13,4.517)rectangle++(0.04,0.04)
  (1.139,5.39)rectangle++(0.04,0.04) (1.182,6.132)rectangle++(0.04,0.04) (-1.046,4.38)rectangle++(0.04,0.04) (-0.716,5.361)rectangle++(0.04,0.04) (-1.223,3.872)rectangle++(0.04,0.04)
  (-0.887,4.936)rectangle++(0.04,0.04) (0.168,2.569)rectangle++(0.04,0.04) (0.243,3.773)rectangle++(0.04,0.04) (0.355,4.794)rectangle++(0.04,0.04) (0.492,5.657)rectangle++(0.04,0.04)
  (1.343,3.165)rectangle++(0.04,0.04) (1.266,4.232)rectangle++(0.04,0.04) (1.243,5.143)rectangle++(0.04,0.04) (1.26,5.918)rectangle++(0.04,0.04) (1.306,6.576)rectangle++(0.04,0.04)
  (1.371,7.134)rectangle++(0.04,0.04) (1.448,7.605)rectangle++(0.04,0.04) (1.533,8.003)rectangle++(0.04,0.04) (1.621,8.337)rectangle++(0.04,0.04) (1.709,8.618)rectangle++(0.04,0.04)
  (1.795,8.853)rectangle++(0.04,0.04) (1.878,9.049)rectangle++(0.04,0.04) (1.956,9.213)rectangle++(0.04,0.04) (2.029,9.349)rectangle++(0.04,0.04) (2.096,9.462)rectangle++(0.04,0.04)
  (2.158,9.555)rectangle++(0.04,0.04) (-1.597,5.238)rectangle++(0.04,0.04) (-1.15,6.113)rectangle++(0.04,0.04) (-0.736,6.838)rectangle++(0.04,0.04) (-0.354,7.438)rectangle++(0.04,0.04)
  (-0.006,7.933)rectangle++(0.04,0.04) (0.31,8.339)rectangle++(0.04,0.04) (0.595,8.672)rectangle++(0.04,0.04) (-1.69,4.771)rectangle++(0.04,0.04) (-1.248,5.719)rectangle++(0.04,0.04)
  (-0.834,6.508)rectangle++(0.04,0.04) (-0.451,7.161)rectangle++(0.04,0.04) (-0.099,7.701)rectangle++(0.04,0.04) (0.222,8.146)rectangle++(0.04,0.04) (0.512,8.511)rectangle++(0.04,0.04)
  (0.773,8.81)rectangle++(0.04,0.04) (1.008,9.054)rectangle++(0.04,0.04) (1.216,9.252)rectangle++(0.04,0.04) (2.149,4.022)rectangle++(0.04,0.04) (1.985,4.929)rectangle++(0.04,0.04)
  (1.882,5.707)rectangle++(0.04,0.04) (1.826,6.372)rectangle++(0.04,0.04) (1.805,6.939)rectangle++(0.04,0.04) (1.809,7.422)rectangle++(0.04,0.04) (1.833,7.833)rectangle++(0.04,0.04)
  (1.869,8.181)rectangle++(0.04,0.04) (1.914,8.475)rectangle++(0.04,0.04) (1.963,8.723)rectangle++(0.04,0.04) (2.016,8.933)rectangle++(0.04,0.04) (1.976,4.142)rectangle++(0.04,0.04)
  (1.843,5.038)rectangle++(0.04,0.04) (1.766,5.805)rectangle++(0.04,0.04) (1.731,6.46)rectangle++(0.04,0.04) (1.728,7.018)rectangle++(0.04,0.04) (1.747,7.492)rectangle++(0.04,0.04)
  (1.782,7.895)rectangle++(0.04,0.04) (1.829,8.235)rectangle++(0.04,0.04) (1.881,8.523)rectangle++(0.04,0.04) (1.938,8.766)rectangle++(0.04,0.04) (1.996,8.969)rectangle++(0.04,0.04)
  (2.053,9.14)rectangle++(0.04,0.04) (2.108,9.283)rectangle++(0.04,0.04) (2.161,9.403)rectangle++(0.04,0.04) (0.04,2.861)rectangle++(0.04,0.04) (0.147,4.027)rectangle++(0.04,0.04)
  (0.283,5.013)rectangle++(0.04,0.04) (-0.002,1.958)rectangle++(0.04,0.04) (0.074,3.26)rectangle++(0.04,0.04) (0.191,4.365)rectangle++(0.04,0.04) (0.335,5.298)rectangle++(0.04,0.04)
  (-0.082,3.566)rectangle++(0.04,0.04) (0.07,4.631)rectangle++(0.04,0.04) (0.243,5.53)rectangle++(0.04,0.04) (0.425,6.287)rectangle++(0.04,0.04) (0.611,6.923)rectangle++(0.04,0.04)
  (0.794,7.456)rectangle++(0.04,0.04) (0.971,7.902)rectangle++(0.04,0.04) (1.139,8.274)rectangle++(0.04,0.04) (1.916,3.816)rectangle++(0.04,0.04) (0.721,3.133)rectangle++(0.04,0.04)
  (0.692,2.73)rectangle++(0.04,0.04) (0.695,3.889)rectangle++(0.04,0.04) (0.744,4.874)rectangle++(0.04,0.04) (0.825,5.709)rectangle++(0.04,0.04) (0.928,6.416)rectangle++(0.04,0.04)
  (1.48,3.41)rectangle++(0.04,0.04) (1.392,4.435)rectangle++(0.04,0.04) (-0.648,4.084)rectangle++(0.04,0.04) (-0.39,5.094)rectangle++(0.04,0.04) (-0.13,5.942)rectangle++(0.04,0.04)
  (0.125,6.652)rectangle++(0.04,0.04) (0.37,7.245)rectangle++(0.04,0.04) (0.602,7.74)rectangle++(0.04,0.04) (0.819,8.151)rectangle++(0.04,0.04) (1.02,8.492)rectangle++(0.04,0.04)
  (1.205,8.773)rectangle++(0.04,0.04) (1.373,9.005)rectangle++(0.04,0.04) (1.525,9.196)rectangle++(0.04,0.04) (1.662,9.352)rectangle++(0.04,0.04) (1.784,9.479)rectangle++(0.04,0.04)
  (1.894,9.582)rectangle++(0.04,0.04) (-1.651,5.194)rectangle++(0.04,0.04) (-1.197,6.077)rectangle++(0.04,0.04) (-0.777,6.81)rectangle++(0.04,0.04) (-0.39,7.416)rectangle++(0.04,0.04)
  (-0.037,7.915)rectangle++(0.04,0.04) (0.283,8.325)rectangle++(0.04,0.04) (0.571,8.662)rectangle++(0.04,0.04) (0.83,8.936)rectangle++(0.04,0.04) (1.061,9.158)rectangle++(0.04,0.04)
  (1.266,9.338)rectangle++(0.04,0.04) (2.164,4.05)rectangle++(0.04,0.04) (1.999,4.952)rectangle++(0.04,0.04) (1.895,5.726)rectangle++(0.04,0.04) (1.838,6.387)rectangle++(0.04,0.04);
\fill[piece2]
  (1.815,6.952)rectangle++(0.04,0.04) (1.819,7.433)rectangle++(0.04,0.04) (1.628,3.808)rectangle++(0.04,0.04) (1.534,4.768)rectangle++(0.04,0.04) (1.492,5.588)rectangle++(0.04,0.04)
  (-0.875,4.33)rectangle++(0.04,0.04) (-0.573,5.312)rectangle++(0.04,0.04) (-0.277,6.134)rectangle++(0.04,0.04) (0.008,6.822)rectangle++(0.04,0.04) (0.278,7.394)rectangle++(0.04,0.04)
  (0.529,7.87)rectangle++(0.04,0.04) (1.887,3.6)rectangle++(0.04,0.04) (1.746,4.581)rectangle++(0.04,0.04) (0.93,3.242)rectangle++(0.04,0.04) (0.694,2.799)rectangle++(0.04,0.04)
  (0.7,3.947)rectangle++(0.04,0.04) (0.75,4.923)rectangle++(0.04,0.04) (0.833,5.751)rectangle++(0.04,0.04) (0.936,6.451)rectangle++(0.04,0.04) (1.051,7.042)rectangle++(0.04,0.04)
  (1.173,7.54)rectangle++(0.04,0.04) (-1.347,4.656)rectangle++(0.04,0.04) (-0.96,5.608)rectangle++(0.04,0.04) (-0.594,6.401)rectangle++(0.04,0.04) (-0.251,7.061)rectangle++(0.04,0.04)
  (0.067,7.608)rectangle++(0.04,0.04) (0.359,8.06)rectangle++(0.04,0.04) (-1.054,3.244)rectangle++(0.04,0.04) (0.636,2.115)rectangle++(0.04,0.04) (-0.292,3.485)rectangle++(0.04,0.04)
  (-0.111,4.57)rectangle++(0.04,0.04) (0.087,5.485)rectangle++(0.04,0.04) (0.291,6.255)rectangle++(0.04,0.04) (0.495,6.901)rectangle++(0.04,0.04) (0.695,7.442)rectangle++(0.04,0.04)
  (0.886,7.895)rectangle++(0.04,0.04) (1.067,8.271)rectangle++(0.04,0.04) (1.235,8.584)rectangle++(0.04,0.04) (1.391,8.843)rectangle++(0.04,0.04) (1.534,9.057)rectangle++(0.04,0.04)
  (1.664,9.233)rectangle++(0.04,0.04) (1.782,9.378)rectangle++(0.04,0.04) (2.114,4.175)rectangle++(0.04,0.04) (1.961,5.06)rectangle++(0.04,0.04) (1.867,5.819)rectangle++(0.04,0.04)
  (1.222,3.508)rectangle++(0.04,0.04) (1.177,4.529)rectangle++(0.04,0.04) (0.984,3.103)rectangle++(0.04,0.04) (0.958,4.194)rectangle++(0.04,0.04) (0.98,5.123)rectangle++(0.04,0.04)
  (1.036,5.912)rectangle++(0.04,0.04) (1.115,6.58)rectangle++(0.04,0.04) (1.208,7.144)rectangle++(0.04,0.04) (1.311,7.62)rectangle++(0.04,0.04) (1.417,8.021)rectangle++(0.04,0.04)
  (1.523,8.358)rectangle++(0.04,0.04) (1.627,8.639)rectangle++(0.04,0.04) (1.726,8.874)rectangle++(0.04,0.04) (1.82,9.07)rectangle++(0.04,0.04) (1.907,9.233)rectangle++(0.04,0.04)
  (1.988,9.368)rectangle++(0.04,0.04) (2.063,9.48)rectangle++(0.04,0.04) (2.13,9.571)rectangle++(0.04,0.04) (2.191,9.647)rectangle++(0.04,0.04) (2.246,9.708)rectangle++(0.04,0.04)
  (2.296,9.758)rectangle++(0.04,0.04) (2.339,9.799)rectangle++(0.04,0.04) (0.387,2.257)rectangle++(0.04,0.04) (0.417,3.499)rectangle++(0.04,0.04) (0.492,4.554)rectangle++(0.04,0.04)
  (0.598,5.447)rectangle++(0.04,0.04) (0.724,6.202)rectangle++(0.04,0.04) (0.861,6.839)rectangle++(0.04,0.04) (1.004,7.375)rectangle++(0.04,0.04) (1.146,7.825)rectangle++(0.04,0.04)
  (1.285,8.202)rectangle++(0.04,0.04) (1.418,8.516)rectangle++(0.04,0.04) (1.544,8.778)rectangle++(0.04,0.04) (1.661,8.996)rectangle++(0.04,0.04) (1.77,9.176)rectangle++(0.04,0.04)
  (1.869,9.325)rectangle++(0.04,0.04) (1.959,9.448)rectangle++(0.04,0.04) (2.041,9.549)rectangle++(0.04,0.04) (2.115,9.631)rectangle++(0.04,0.04) (2.181,9.698)rectangle++(0.04,0.04)
  (2.146,4.327)rectangle++(0.04,0.04) (1.995,5.188)rectangle++(0.04,0.04) (1.901,5.927)rectangle++(0.04,0.04) (1.851,6.558)rectangle++(0.04,0.04) (1.833,7.096)rectangle++(0.04,0.04)
  (1.84,7.555)rectangle++(0.04,0.04) (1.864,7.944)rectangle++(0.04,0.04) (1.9,8.274)rectangle++(0.04,0.04) (1.944,8.553)rectangle++(0.04,0.04) (1.992,8.789)rectangle++(0.04,0.04)
  (2.043,8.987)rectangle++(0.04,0.04) (2.093,9.153)rectangle++(0.04,0.04) (-1.524,5.147)rectangle++(0.04,0.04) (-1.092,6.032)rectangle++(0.04,0.04) (-0.689,6.767)rectangle++(0.04,0.04)
  (-0.317,7.376)rectangle++(0.04,0.04) (0.023,7.878)rectangle++(0.04,0.04) (0.333,8.292)rectangle++(0.04,0.04) (0.612,8.631)rectangle++(0.04,0.04) (-0.25,3.213)rectangle++(0.04,0.04)
  (-0.086,4.337)rectangle++(0.04,0.04) (0.098,5.286)rectangle++(0.04,0.04) (0.293,6.086)rectangle++(0.04,0.04) (0.49,6.757)rectangle++(0.04,0.04) (0.685,7.32)rectangle++(0.04,0.04)
  (0.873,7.791)rectangle++(0.04,0.04) (1.827,3.663)rectangle++(0.04,0.04) (1.697,4.636)rectangle++(0.04,0.04) (1.626,5.469)rectangle++(0.04,0.04) (1.164,3.386)rectangle++(0.04,0.04)
  (1.123,4.428)rectangle++(0.04,0.04) (1.129,5.315)rectangle++(0.04,0.04) (1.17,6.069)rectangle++(0.04,0.04) (1.235,6.708)rectangle++(0.04,0.04) (1.316,7.248)rectangle++(0.04,0.04)
  (1.406,7.705)rectangle++(0.04,0.04) (1.501,8.089)rectangle++(0.04,0.04) (1.598,8.412)rectangle++(0.04,0.04) (-1.455,4.908)rectangle++(0.04,0.04) (-1.043,5.826)rectangle++(0.04,0.04)
  (-0.656,6.59)rectangle++(0.04,0.04) (-0.296,7.224)rectangle++(0.04,0.04) (0.035,7.748)rectangle++(0.04,0.04) (0.338,8.181)rectangle++(0.04,0.04) (0.612,8.536)rectangle++(0.04,0.04)
  (0.86,8.828)rectangle++(0.04,0.04) (1.081,9.065)rectangle++(0.04,0.04) (-1.687,4.94)rectangle++(0.04,0.04) (1.421,2.532)rectangle++(0.04,0.04) (1.307,3.692)rectangle++(0.04,0.04)
  (1.256,4.682)rectangle++(0.04,0.04) (1.253,5.525)rectangle++(0.04,0.04) (-0.905,4.273)rectangle++(0.04,0.04) (-0.6,5.265)rectangle++(0.04,0.04) (-0.302,6.095)rectangle++(0.04,0.04)
  (-0.015,6.789)rectangle++(0.04,0.04) (0.257,7.368)rectangle++(0.04,0.04) (0.511,7.848)rectangle++(0.04,0.04) (0.746,8.247)rectangle++(0.04,0.04) (1.955,3.722)rectangle++(0.04,0.04)
  (1.808,4.681)rectangle++(0.04,0.04) (1.722,5.503)rectangle++(0.04,0.04) (1.682,6.205)rectangle++(0.04,0.04) (1.675,6.803)rectangle++(0.04,0.04) (1.694,7.312)rectangle++(0.04,0.04)
  (1.73,7.744)rectangle++(0.04,0.04) (1.778,8.109)rectangle++(0.04,0.04) (1.834,8.418)rectangle++(0.04,0.04) (1.893,8.678)rectangle++(0.04,0.04) (1.954,8.897)rectangle++(0.04,0.04)
  (2.015,9.08)rectangle++(0.04,0.04) (-1.522,5.118)rectangle++(0.04,0.04) (-1.092,6.007)rectangle++(0.04,0.04) (-0.69,6.746)rectangle++(0.04,0.04) (-0.319,7.358)rectangle++(0.04,0.04)
  (0.021,7.863)rectangle++(0.04,0.04) (0.33,8.279)rectangle++(0.04,0.04) (0.61,8.62)rectangle++(0.04,0.04) (0.861,8.899)rectangle++(0.04,0.04) (2.103,3.874)rectangle++(0.04,0.04)
  (1.94,4.805)rectangle++(0.04,0.04) (1.839,5.603)rectangle++(0.04,0.04) (1.785,6.285)rectangle++(0.04,0.04) (1.767,6.867)rectangle++(0.04,0.04) (1.774,7.362)rectangle++(0.04,0.04)
  (1.8,7.783)rectangle++(0.04,0.04) (1.839,8.14)rectangle++(0.04,0.04) (-1.081,3.367)rectangle++(0.04,0.04) (-0.786,4.501)rectangle++(0.04,0.04) (-0.49,5.453)rectangle++(0.04,0.04)
  (-1.202,3.932)rectangle++(0.04,0.04) (-0.866,4.987)rectangle++(0.04,0.04) (-0.539,5.87)rectangle++(0.04,0.04) (-0.226,6.607)rectangle++(0.04,0.04) (0.07,7.221)rectangle++(0.04,0.04)
  (0.346,7.731)rectangle++(0.04,0.04) (0.601,8.154)rectangle++(0.04,0.04) (-1.579,4.67)rectangle++(0.04,0.04) (-1.158,5.629)rectangle++(0.04,0.04) (-0.761,6.427)rectangle++(0.04,0.04)
  (-0.392,7.09)rectangle++(0.04,0.04) (-1.531,4.274)rectangle++(0.04,0.04) (-1.133,5.29)rectangle++(0.04,0.04) (-0.754,6.138)rectangle++(0.04,0.04) (-0.397,6.844)rectangle++(0.04,0.04)
  (-0.066,7.429)rectangle++(0.04,0.04) (0.239,7.914)rectangle++(0.04,0.04) (0.517,8.313)rectangle++(0.04,0.04) (0.77,8.642)rectangle++(0.04,0.04) (0.998,8.911)rectangle++(0.04,0.04)
  (1.203,9.131)rectangle++(0.04,0.04) (1.385,9.309)rectangle++(0.04,0.04) (-1.684,5.045)rectangle++(0.04,0.04) (-1.232,5.952)rectangle++(0.04,0.04) (1.618,2.83)rectangle++(0.04,0.04)
  (1.486,3.937)rectangle++(0.04,0.04) (-0.526,4.004)rectangle++(0.04,0.04) (-0.289,5.02)rectangle++(0.04,0.04) (-1.074,3.885)rectangle++(0.04,0.04) (0.087,2.224)rectangle++(0.04,0.04)
  (0.161,3.483)rectangle++(0.04,0.04) (0.274,4.55)rectangle++(0.04,0.04) (0.413,5.453)rectangle++(0.04,0.04) (0.567,6.215)rectangle++(0.04,0.04) (0.728,6.856)rectangle++(0.04,0.04)
  (0.891,7.395)rectangle++(0.04,0.04) (1.051,7.846)rectangle++(0.04,0.04) (1.205,8.223)rectangle++(0.04,0.04) (1.351,8.538)rectangle++(0.04,0.04) (1.488,8.799)rectangle++(0.04,0.04)
  (1.614,9.016)rectangle++(0.04,0.04) (1.731,9.195)rectangle++(0.04,0.04) (1.837,9.343)rectangle++(0.04,0.04) (1.933,9.464)rectangle++(0.04,0.04) (2.019,9.564)rectangle++(0.04,0.04)
  (2.097,9.644)rectangle++(0.04,0.04) (2.166,9.71)rectangle++(0.04,0.04) (2.227,9.763)rectangle++(0.04,0.04) (2.281,9.806)rectangle++(0.04,0.04) (2.329,9.84)rectangle++(0.04,0.04)
  (2.371,9.867)rectangle++(0.04,0.04) (2.408,9.888)rectangle++(0.04,0.04) (-1.623,5.351)rectangle++(0.04,0.04) (-1.168,6.21)rectangle++(0.04,0.04) (-0.747,6.921)rectangle++(0.04,0.04)
  (-0.36,7.509)rectangle++(0.04,0.04) (-0.008,7.993)rectangle++(0.04,0.04) (0.311,8.391)rectangle++(0.04,0.04) (0.598,8.716)rectangle++(0.04,0.04) (0.854,8.981)rectangle++(0.04,0.04)
  (1.083,9.196)rectangle++(0.04,0.04) (1.286,9.369)rectangle++(0.04,0.04) (1.466,9.509)rectangle++(0.04,0.04) (-1.712,5.1)rectangle++(0.04,0.04) (-1.253,5.999)rectangle++(0.04,0.04)
  (-0.827,6.746)rectangle++(0.04,0.04) (1.767,3.075)rectangle++(0.04,0.04) (1.623,4.139)rectangle++(0.04,0.04) (1.543,5.05)rectangle++(0.04,0.04) (1.511,5.827)rectangle++(0.04,0.04)
  (-0.179,3.646)rectangle++(0.04,0.04) (-0.009,4.703)rectangle++(0.04,0.04) (0.178,5.594)rectangle++(0.04,0.04) (0.373,6.344)rectangle++(0.04,0.04) (0.569,6.974)rectangle++(0.04,0.04)
  (0.76,7.501)rectangle++(0.04,0.04) (0.944,7.942)rectangle++(0.04,0.04) (1.118,8.309)rectangle++(0.04,0.04) (1.28,8.614)rectangle++(0.04,0.04) (1.431,8.867)rectangle++(0.04,0.04)
  (1.569,9.076)rectangle++(0.04,0.04) (-0.187,3.972)rectangle++(0.04,0.04) (-0.002,4.98)rectangle++(0.04,0.04) (0.195,5.829)rectangle++(0.04,0.04) (1.421,3.131)rectangle++(0.04,0.04)
  (1.331,4.2)rectangle++(0.04,0.04) (1.297,5.113)rectangle++(0.04,0.04) (1.305,5.891)rectangle++(0.04,0.04) (1.342,6.551)rectangle++(0.04,0.04) (-1.106,4.493)rectangle++(0.04,0.04)
  (-0.763,5.459)rectangle++(0.04,0.04) (-0.432,6.267)rectangle++(0.04,0.04) (-0.119,6.94)rectangle++(0.04,0.04) (0.174,7.5)rectangle++(0.04,0.04) (0.446,7.964)rectangle++(0.04,0.04)
  (0.695,8.348)rectangle++(0.04,0.04) (0.034,2.718)rectangle++(0.04,0.04) (0.135,3.905)rectangle++(0.04,0.04) (0.269,4.91)rectangle++(0.04,0.04) (-0.95,3.967)rectangle++(0.04,0.04)
  (-0.651,5.006)rectangle++(0.04,0.04) (-0.356,5.878)rectangle++(0.04,0.04) (-0.069,6.607)rectangle++(0.04,0.04) (0.203,7.215)rectangle++(0.04,0.04) (0.459,7.72)rectangle++(0.04,0.04)
  (0.697,8.14)rectangle++(0.04,0.04) (0.916,8.487)rectangle++(0.04,0.04) (1.116,8.774)rectangle++(0.04,0.04) (1.297,9.009)rectangle++(0.04,0.04) (-1.636,4.969)rectangle++(0.04,0.04)
  (-1.194,5.886)rectangle++(0.04,0.04) (-0.782,6.647)rectangle++(0.04,0.04) (-0.401,7.277)rectangle++(0.04,0.04) (-0.052,7.798)rectangle++(0.04,0.04) (0.266,8.226)rectangle++(0.04,0.04)
  (0.553,8.578)rectangle++(0.04,0.04) (2.059,3.742)rectangle++(0.04,0.04) (1.898,4.694)rectangle++(0.04,0.04) (1.799,5.51)rectangle++(0.04,0.04) (1.747,6.208)rectangle++(0.04,0.04)
  (1.731,6.803)rectangle++(0.04,0.04) (1.742,7.31)rectangle++(0.04,0.04) (1.77,7.74)rectangle++(0.04,0.04) (1.812,8.104)rectangle++(0.04,0.04) (1.862,8.412)rectangle++(0.04,0.04)
  (1.917,8.672)rectangle++(0.04,0.04) (1.974,8.891)rectangle++(0.04,0.04) (2.032,9.074)rectangle++(0.04,0.04) (2.088,9.228)rectangle++(0.04,0.04) (2.141,9.357)rectangle++(0.04,0.04)
  (2.192,9.464)rectangle++(0.04,0.04) (2.24,9.553)rectangle++(0.04,0.04) (2.284,9.626)rectangle++(0.04,0.04) (2.324,9.687)rectangle++(0.04,0.04) (2.361,9.737)rectangle++(0.04,0.04)
  (2.394,9.779)rectangle++(0.04,0.04) (2.424,9.812)rectangle++(0.04,0.04) (2.45,9.84)rectangle++(0.04,0.04) (2.474,9.862)rectangle++(0.04,0.04) (2.495,9.88)rectangle++(0.04,0.04)
  (2.514,9.894)rectangle++(0.04,0.04) (2.531,9.906)rectangle++(0.04,0.04) (2.545,9.915)rectangle++(0.04,0.04) (2.558,9.922)rectangle++(0.04,0.04) (2.569,9.928)rectangle++(0.04,0.04)
  (2.578,9.932)rectangle++(0.04,0.04) (2.587,9.935)rectangle++(0.04,0.04) (2.594,9.938)rectangle++(0.04,0.04) (2.6,9.939)rectangle++(0.04,0.04) (2.605,9.941)rectangle++(0.04,0.04)
  (2.61,9.942)rectangle++(0.04,0.04) (2.614,9.942)rectangle++(0.04,0.04) (2.617,9.942)rectangle++(0.04,0.04) (2.62,9.943)rectangle++(0.04,0.04) (2.155,4.474)rectangle++(0.04,0.04)
  (2.009,5.313)rectangle++(0.04,0.04) (1.918,6.032)rectangle++(0.04,0.04) (1.869,6.647)rectangle++(0.04,0.04) (1.852,7.171)rectangle++(0.04,0.04) (1.859,7.618)rectangle++(0.04,0.04)
  (1.883,7.997)rectangle++(0.04,0.04) (1.918,8.318)rectangle++(0.04,0.04) (1.961,8.59)rectangle++(0.04,0.04) (2.008,8.819)rectangle++(0.04,0.04) (-1.468,5.065)rectangle++(0.04,0.04)
  (-1.048,5.96)rectangle++(0.04,0.04) (-0.654,6.704)rectangle++(0.04,0.04) (-0.29,7.321)rectangle++(0.04,0.04) (-1.562,4.338)rectangle++(0.04,0.04) (-1.157,5.346)rectangle++(0.04,0.04)
  (-0.771,6.187)rectangle++(0.04,0.04) (-0.41,6.886)rectangle++(0.04,0.04) (-1.491,4.23)rectangle++(0.04,0.04) (-1.101,5.251)rectangle++(0.04,0.04) (-0.728,6.104)rectangle++(0.04,0.04)
  (-0.377,6.814)rectangle++(0.04,0.04) (-0.05,7.403)rectangle++(0.04,0.04) (0.252,7.891)rectangle++(0.04,0.04) (0.527,8.293)rectangle++(0.04,0.04) (0.778,8.624)rectangle++(0.04,0.04)
  (1.004,8.896)rectangle++(0.04,0.04) (1.207,9.117)rectangle++(0.04,0.04) (1.388,9.298)rectangle++(0.04,0.04) (-0.172,3.99)rectangle++(0.04,0.04) (0.011,4.995)rectangle++(0.04,0.04)
  (0.207,5.841)rectangle++(0.04,0.04) (0.408,6.553)rectangle++(0.04,0.04) (0.606,7.15)rectangle++(0.04,0.04) (0.799,7.649)rectangle++(0.04,0.04) (0.983,8.066)rectangle++(0.04,0.04)
  (1.156,8.413)rectangle++(0.04,0.04) (1.317,8.701)rectangle++(0.04,0.04) (1.465,8.939)rectangle++(0.04,0.04) (1.601,9.136)rectangle++(0.04,0.04) (2.075,4.087)rectangle++(0.04,0.04)
  (0.769,3.221)rectangle++(0.04,0.04) (0.708,2.758)rectangle++(0.04,0.04) (-0.415,3.625)rectangle++(0.04,0.04) (-0.21,4.694)rectangle++(0.04,0.04) (0.007,5.594)rectangle++(0.04,0.04)
  (0.228,6.351)rectangle++(0.04,0.04) (0.445,6.986)rectangle++(0.04,0.04) (0.656,7.516)rectangle++(0.04,0.04) (0.856,7.959)rectangle++(0.04,0.04) (1.044,8.327)rectangle++(0.04,0.04)
  (1.218,8.632)rectangle++(0.04,0.04) (1.378,8.885)rectangle++(0.04,0.04) (1.525,9.093)rectangle++(0.04,0.04) (1.658,9.264)rectangle++(0.04,0.04) (-1.626,5.087)rectangle++(0.04,0.04)
  (-1.18,5.985)rectangle++(0.04,0.04) (1.62,2.848)rectangle++(0.04,0.04) (1.489,3.953)rectangle++(0.04,0.04) (1.422,4.896)rectangle++(0.04,0.04) (1.402,5.701)rectangle++(0.04,0.04)
  (1.418,6.386)rectangle++(0.04,0.04) (1.458,6.968)rectangle++(0.04,0.04) (1.516,7.46)rectangle++(0.04,0.04) (-1.268,4.705)rectangle++(0.04,0.04) (-0.892,5.646)rectangle++(0.04,0.04)
  (-0.534,6.431)rectangle++(0.04,0.04) (-0.199,7.084)rectangle++(0.04,0.04) (0.112,7.625)rectangle++(0.04,0.04) (0.398,8.073)rectangle++(0.04,0.04) (0.659,8.443)rectangle++(0.04,0.04)
  (0.896,8.746)rectangle++(0.04,0.04) (1.109,8.995)rectangle++(0.04,0.04) (1.3,9.197)rectangle++(0.04,0.04) (1.471,9.362)rectangle++(0.04,0.04) (1.622,9.495)rectangle++(0.04,0.04)
  (1.757,9.602)rectangle++(0.04,0.04) (1.875,9.688)rectangle++(0.04,0.04) (1.979,9.756)rectangle++(0.04,0.04) (2.07,9.809)rectangle++(0.04,0.04) (-0.478,3.293)rectangle++(0.04,0.04)
  (-0.277,4.414)rectangle++(0.04,0.04) (1.126,2.751)rectangle++(0.04,0.04) (1.065,3.889)rectangle++(0.04,0.04) (-0.141,3.442)rectangle++(0.04,0.04) (0.015,4.528)rectangle++(0.04,0.04)
  (0.192,5.444)rectangle++(0.04,0.04) (0.379,6.216)rectangle++(0.04,0.04) (0.568,6.865)rectangle++(0.04,0.04) (0.756,7.409)rectangle++(0.04,0.04) (1.051,2.209)rectangle++(0.04,0.04)
  (0.979,3.432)rectangle++(0.04,0.04) (0.967,4.474)rectangle++(0.04,0.04) (0.999,5.36)rectangle++(0.04,0.04) (1.061,6.113)rectangle++(0.04,0.04) (1.144,6.749)rectangle++(0.04,0.04)
  (1.241,7.287)rectangle++(0.04,0.04) (1.344,7.741)rectangle++(0.04,0.04) (1.45,8.122)rectangle++(0.04,0.04) (1.134,2.259)rectangle++(0.04,0.04) (1.052,3.471)rectangle++(0.04,0.04)
  (1.031,4.505)rectangle++(0.04,0.04) (1.055,5.384)rectangle++(0.04,0.04) (1.11,6.13)rectangle++(0.04,0.04) (1.186,6.763)rectangle++(0.04,0.04) (1.276,7.297)rectangle++(0.04,0.04)
  (1.657,3.659)rectangle++(0.04,0.04) (-0.436,3.981)rectangle++(0.04,0.04) (-0.214,4.997)rectangle++(0.04,0.04) (0.016,5.852)rectangle++(0.04,0.04) (0.246,6.57)rectangle++(0.04,0.04)
  (0.469,7.171)rectangle++(0.04,0.04) (-1.394,4.451)rectangle++(0.04,0.04) (-1.009,5.436)rectangle++(0.04,0.04) (-0.643,6.257)rectangle++(0.04,0.04) (-0.298,6.94)rectangle++(0.04,0.04)
  (0.025,2.527)rectangle++(0.04,0.04) (0.12,3.743)rectangle++(0.04,0.04) (0.249,4.773)rectangle++(0.04,0.04) (0.401,5.643)rectangle++(0.04,0.04) (0.564,6.377)rectangle++(0.04,0.04)
  (0.732,6.994)rectangle++(0.04,0.04) (0.9,7.512)rectangle++(0.04,0.04) (1.063,7.945)rectangle++(0.04,0.04) (1.219,8.307)rectangle++(0.04,0.04) (1.367,8.608)rectangle++(0.04,0.04)
  (1.504,8.859)rectangle++(0.04,0.04) (1.63,9.066)rectangle++(0.04,0.04) (1.746,9.237)rectangle++(0.04,0.04) (1.852,9.378)rectangle++(0.04,0.04) (1.947,9.493)rectangle++(0.04,0.04)
  (2.032,9.588)rectangle++(0.04,0.04) (2.109,9.664)rectangle++(0.04,0.04) (2.177,9.727)rectangle++(0.04,0.04) (2.237,9.777)rectangle++(0.04,0.04) (2.29,9.817)rectangle++(0.04,0.04)
  (2.337,9.849)rectangle++(0.04,0.04) (2.379,9.874)rectangle++(0.04,0.04) (2.415,9.894)rectangle++(0.04,0.04) (2.446,9.91)rectangle++(0.04,0.04) (2.473,9.922)rectangle++(0.04,0.04)
  (2.497,9.931)rectangle++(0.04,0.04) (2.517,9.937)rectangle++(0.04,0.04) (2.535,9.942)rectangle++(0.04,0.04) (2.55,9.946)rectangle++(0.04,0.04) (2.164,4.459)rectangle++(0.04,0.04)
  (2.016,5.3)rectangle++(0.04,0.04) (1.923,6.02)rectangle++(0.04,0.04) (1.873,6.637)rectangle++(0.04,0.04) (1.856,7.162)rectangle++(0.04,0.04) (1.861,7.61)rectangle++(0.04,0.04)
  (1.884,7.99)rectangle++(0.04,0.04) (1.919,8.312)rectangle++(0.04,0.04) (1.962,8.585)rectangle++(0.04,0.04) (2.009,8.815)rectangle++(0.04,0.04) (2.058,9.009)rectangle++(0.04,0.04)
  (2.107,9.171)rectangle++(0.04,0.04) (2.156,9.307)rectangle++(0.04,0.04) (2.202,9.421)rectangle++(0.04,0.04) (-1.561,5.22)rectangle++(0.04,0.04) (-1.12,6.096)rectangle++(0.04,0.04)
  (-0.711,6.823)rectangle++(0.04,0.04) (1.773,3.116)rectangle++(0.04,0.04) (1.629,4.174)rectangle++(0.04,0.04) (1.55,5.079)rectangle++(0.04,0.04) (1.518,5.852)rectangle++(0.04,0.04)
  (1.522,6.509)rectangle++(0.04,0.04) (1.552,7.068)rectangle++(0.04,0.04) (1.6,7.542)rectangle++(0.04,0.04) (1.659,7.943)rectangle++(0.04,0.04) (1.726,8.281)rectangle++(0.04,0.04)
  (1.796,8.566)rectangle++(0.04,0.04) (1.867,8.806)rectangle++(0.04,0.04) (1.937,9.006)rectangle++(0.04,0.04) (2.005,9.174)rectangle++(0.04,0.04) (2.069,9.314)rectangle++(0.04,0.04)
  (2.129,9.43)rectangle++(0.04,0.04) (2.184,9.527)rectangle++(0.04,0.04) (2.236,9.607)rectangle++(0.04,0.04) (2.282,9.672)rectangle++(0.04,0.04) (2.325,9.726)rectangle++(0.04,0.04)
  (2.363,9.771)rectangle++(0.04,0.04) (2.397,9.807)rectangle++(0.04,0.04) (2.428,9.836)rectangle++(0.04,0.04) (2.455,9.86)rectangle++(0.04,0.04) (2.479,9.879)rectangle++(0.04,0.04)
  (2.5,9.894)rectangle++(0.04,0.04) (2.518,9.906)rectangle++(0.04,0.04) (2.535,9.916)rectangle++(0.04,0.04) (2.159,4.45)rectangle++(0.04,0.04) (2.011,5.292)rectangle++(0.04,0.04)
  (1.918,6.014)rectangle++(0.04,0.04) (1.869,6.631)rectangle++(0.04,0.04) (1.852,7.158)rectangle++(0.04,0.04) (1.858,7.606)rectangle++(0.04,0.04) (1.881,7.987)rectangle++(0.04,0.04)
  (1.917,8.31)rectangle++(0.04,0.04) (1.959,8.583)rectangle++(0.04,0.04) (2.006,8.814)rectangle++(0.04,0.04) (2.056,9.007)rectangle++(0.04,0.04) (2.106,9.17)rectangle++(0.04,0.04)
  (2.024,4.208)rectangle++(0.04,0.04) (1.887,5.092)rectangle++(0.04,0.04) (1.805,5.849)rectangle++(0.04,0.04) (1.766,6.496)rectangle++(0.04,0.04) (1.402,3.616)rectangle++(0.04,0.04)
  (-0.473,3.924)rectangle++(0.04,0.04) (1.027,2.612)rectangle++(0.04,0.04) (0.975,3.775)rectangle++(0.04,0.04) (0.821,2.912)rectangle++(0.04,0.04) (0.812,4.038)rectangle++(0.04,0.04)
  (0.849,4.996)rectangle++(0.04,0.04) (0.92,5.809)rectangle++(0.04,0.04) (1.012,6.497)rectangle++(0.04,0.04) (1.118,7.078)rectangle++(0.04,0.04) (1.231,7.568)rectangle++(0.04,0.04)
  (-1.342,4.672)rectangle++(0.04,0.04) (-0.956,5.621)rectangle++(0.04,0.04) (-0.59,6.413)rectangle++(0.04,0.04) (-0.247,7.071)rectangle++(0.04,0.04) (0.071,7.616)rectangle++(0.04,0.04)
  (0.362,8.067)rectangle++(0.04,0.04) (0.629,8.439)rectangle++(0.04,0.04) (0.87,8.744)rectangle++(0.04,0.04) (1.087,8.994)rectangle++(0.04,0.04) (1.281,9.197)rectangle++(0.04,0.04)
  (2.127,4.027)rectangle++(0.04,0.04) (-0.43,4.141)rectangle++(0.04,0.04) (-0.202,5.134)rectangle++(0.04,0.04) (0.032,5.968)rectangle++(0.04,0.04) (0.263,6.668)rectangle++(0.04,0.04)
  (0.488,7.253)rectangle++(0.04,0.04) (0.703,7.742)rectangle++(0.04,0.04) (0.905,8.149)rectangle++(0.04,0.04) (1.093,8.486)rectangle++(0.04,0.04) (1.266,8.766)rectangle++(0.04,0.04)
  (1.425,8.997)rectangle++(0.04,0.04) (1.569,9.186)rectangle++(0.04,0.04) (1.229,2.252)rectangle++(0.04,0.04) (1.132,3.461)rectangle++(0.04,0.04) (1.099,4.493)rectangle++(0.04,0.04)
  (1.111,5.371)rectangle++(0.04,0.04) (1.2,3.25)rectangle++(0.04,0.04) (1.148,4.311)rectangle++(0.04,0.04) (1.146,5.215)rectangle++(0.04,0.04) (-0.858,4.192)rectangle++(0.04,0.04)
  (-0.564,5.194)rectangle++(0.04,0.04) (-0.274,6.033)rectangle++(0.04,0.04) (0.007,6.735)rectangle++(0.04,0.04) (0.273,7.321)rectangle++(0.04,0.04) (0.523,7.808)rectangle++(0.04,0.04)
  (0.754,8.212)rectangle++(0.04,0.04) (0.967,8.546)rectangle++(0.04,0.04) (1.162,8.822)rectangle++(0.04,0.04) (1.338,9.048)rectangle++(0.04,0.04) (1.497,9.234)rectangle++(0.04,0.04)
  (1.64,9.385)rectangle++(0.04,0.04) (1.767,9.508)rectangle++(0.04,0.04) (1.88,9.607)rectangle++(0.04,0.04) (1.98,9.687)rectangle++(0.04,0.04) (2.068,9.751)rectangle++(0.04,0.04)
  (2.146,9.802)rectangle++(0.04,0.04) (2.176,4.339)rectangle++(0.04,0.04) (2.021,5.198)rectangle++(0.04,0.04) (1.923,5.933)rectangle++(0.04,0.04) (1.87,6.563)rectangle++(0.04,0.04)
  (1.85,7.1)rectangle++(0.04,0.04) (1.854,7.557)rectangle++(0.04,0.04) (1.876,7.945)rectangle++(0.04,0.04) (1.748,3.92)rectangle++(0.04,0.04) (1.64,4.858)rectangle++(0.04,0.04)
  (1.586,5.66)rectangle++(0.04,0.04) (1.573,6.344)rectangle++(0.04,0.04) (1.588,6.926)rectangle++(0.04,0.04) (1.625,7.419)rectangle++(0.04,0.04) (1.676,7.838)rectangle++(0.04,0.04)
  (1.745,3.854)rectangle++(0.04,0.04) (1.635,4.802)rectangle++(0.04,0.04) (0.998,3.259)rectangle++(0.04,0.04) (0.976,4.327)rectangle++(0.04,0.04) (1.001,5.235)rectangle++(0.04,0.04)
  (1.058,6.006)rectangle++(0.04,0.04) (1.137,6.659)rectangle++(0.04,0.04) (1.231,7.211)rectangle++(0.04,0.04) (1.332,7.676)rectangle++(0.04,0.04) (1.437,8.068)rectangle++(0.04,0.04)
  (1.542,8.396)rectangle++(0.04,0.04) (1.645,8.671)rectangle++(0.04,0.04) (1.955,4.008)rectangle++(0.04,0.04) (1.82,4.924)rectangle++(0.04,0.04) (1.741,5.709)rectangle++(0.04,0.04)
  (1.706,6.379)rectangle++(0.04,0.04) (1.703,6.95)rectangle++(0.04,0.04) (1.724,7.436)rectangle++(0.04,0.04) (1.76,7.848)rectangle++(0.04,0.04) (1.808,8.196)rectangle++(0.04,0.04)
  (1.862,8.491)rectangle++(0.04,0.04) (1.921,8.739)rectangle++(0.04,0.04) (1.98,8.947)rectangle++(0.04,0.04) (1.984,4.137)rectangle++(0.04,0.04) (0.007,2.145)rectangle++(0.04,0.04)
  (0.531,2.378)rectangle++(0.04,0.04) (0.544,3.596)rectangle++(0.04,0.04) (0.827,2.779)rectangle++(0.04,0.04) (0.812,3.925)rectangle++(0.04,0.04) (0.845,4.9)rectangle++(0.04,0.04)
  (-0.845,4.073)rectangle++(0.04,0.04) (-0.558,5.092)rectangle++(0.04,0.04) (-0.272,5.947)rectangle++(0.04,0.04) (0.004,6.662)rectangle++(0.04,0.04) (0.268,7.259)rectangle++(0.04,0.04)
  (0.516,7.755)rectangle++(0.04,0.04) (0.746,8.168)rectangle++(0.04,0.04) (0.959,8.509)rectangle++(0.04,0.04) (1.153,8.79)rectangle++(0.04,0.04) (1.33,9.022)rectangle++(0.04,0.04)
  (1.489,9.212)rectangle++(0.04,0.04) (1.632,9.367)rectangle++(0.04,0.04) (1.76,9.492)rectangle++(0.04,0.04) (1.873,9.594)rectangle++(0.04,0.04) (1.974,9.677)rectangle++(0.04,0.04)
  (2.063,9.742)rectangle++(0.04,0.04) (0.537,3.145)rectangle++(0.04,0.04) (-0.528,3.669)rectangle++(0.04,0.04) (-0.304,4.736)rectangle++(0.04,0.04) (0.01,2.227)rectangle++(0.04,0.04)
  (0.096,3.489)rectangle++(0.04,0.04) (0.219,4.558)rectangle++(0.04,0.04) (0.366,5.461)rectangle++(0.04,0.04) (0.527,6.224)rectangle++(0.04,0.04) (0.695,6.865)rectangle++(0.04,0.04)
  (0.349,2.807)rectangle++(0.04,0.04) (0.406,3.968)rectangle++(0.04,0.04) (0.502,4.953)rectangle++(0.04,0.04) (0.623,5.786)rectangle++(0.04,0.04) (0.017,2.37)rectangle++(0.04,0.04)
  (0.107,3.61)rectangle++(0.04,0.04) (0.233,4.66)rectangle++(0.04,0.04) (0.382,5.548)rectangle++(0.04,0.04) (0.545,6.297)rectangle++(0.04,0.04) (0.713,6.927)rectangle++(0.04,0.04)
  (0.881,7.455)rectangle++(0.04,0.04) (1.045,7.898)rectangle++(0.04,0.04) (1.202,8.268)rectangle++(0.04,0.04) (1.35,8.576)rectangle++(0.04,0.04) (1.488,8.832)rectangle++(0.04,0.04)
  (1.616,9.043)rectangle++(0.04,0.04) (1.733,9.218)rectangle++(0.04,0.04) (1.84,9.363)rectangle++(0.04,0.04) (1.936,9.481)rectangle++(0.04,0.04) (2.023,9.577)rectangle++(0.04,0.04)
  (2.1,9.656)rectangle++(0.04,0.04) (2.169,9.72)rectangle++(0.04,0.04) (2.23,9.771)rectangle++(0.04,0.04) (2.285,9.813)rectangle++(0.04,0.04) (2.332,9.846)rectangle++(0.04,0.04)
  (2.374,9.872)rectangle++(0.04,0.04) (2.411,9.892)rectangle++(0.04,0.04) (2.442,9.908)rectangle++(0.04,0.04) (2.167,4.427)rectangle++(0.04,0.04) (2.017,5.273)rectangle++(0.04,0.04)
  (1.923,5.997)rectangle++(0.04,0.04) (1.873,6.617)rectangle++(0.04,0.04) (1.854,7.146)rectangle++(0.04,0.04) (-1.141,4.706)rectangle++(0.04,0.04) (-1.16,3.663)rectangle++(0.04,0.04)
  (-0.841,4.756)rectangle++(0.04,0.04) (-0.527,5.672)rectangle++(0.04,0.04) (-0.223,6.439)rectangle++(0.04,0.04) (0.065,7.078)rectangle++(0.04,0.04) (-0.006,1.875)rectangle++(0.04,0.04)
  (0.068,3.19)rectangle++(0.04,0.04) (0.183,4.305)rectangle++(0.04,0.04) (0.325,5.248)rectangle++(0.04,0.04) (-1.012,4.045)rectangle++(0.04,0.04) (-0.996,3.556)rectangle++(0.04,0.04)
  (-0.707,4.659)rectangle++(0.04,0.04) (1.237,2.7)rectangle++(0.04,0.04) (1.157,3.842)rectangle++(0.04,0.04) (1.135,4.816)rectangle++(0.04,0.04) (1.155,5.644)rectangle++(0.04,0.04)
  (1.262,3.313)rectangle++(0.04,0.04) (1.203,4.362)rectangle++(0.04,0.04) (1.195,5.255)rectangle++(0.04,0.04) (1.224,6.015)rectangle++(0.04,0.04) (1.279,6.66)rectangle++(0.04,0.04)
  (1.351,7.206)rectangle++(0.04,0.04) (1.435,7.668)rectangle++(0.04,0.04) (1.524,8.056)rectangle++(0.04,0.04) (1.615,8.383)rectangle++(0.04,0.04) (1.706,8.657)rectangle++(0.04,0.04)
  (1.794,8.887)rectangle++(0.04,0.04) (1.878,9.078)rectangle++(0.04,0.04) (1.958,9.237)rectangle++(0.04,0.04) (2.031,9.37)rectangle++(0.04,0.04) (2.099,9.479)rectangle++(0.04,0.04)
  (2.161,9.57)rectangle++(0.04,0.04) (2.218,9.644)rectangle++(0.04,0.04) (2.269,9.705)rectangle++(0.04,0.04) (2.314,9.755)rectangle++(0.04,0.04) (2.355,9.795)rectangle++(0.04,0.04)
  (2.391,9.828)rectangle++(0.04,0.04) (2.424,9.854)rectangle++(0.04,0.04) (2.452,9.875)rectangle++(0.04,0.04) (2.477,9.892)rectangle++(0.04,0.04) (2.499,9.905)rectangle++(0.04,0.04)
  (2.518,9.916)rectangle++(0.04,0.04) (2.535,9.924)rectangle++(0.04,0.04) (2.549,9.93)rectangle++(0.04,0.04) (2.562,9.935)rectangle++(0.04,0.04) (2.573,9.938)rectangle++(0.04,0.04)
  (2.582,9.941)rectangle++(0.04,0.04) (2.59,9.943)rectangle++(0.04,0.04) (2.597,9.944)rectangle++(0.04,0.04) (2.603,9.945)rectangle++(0.04,0.04) (2.608,9.945)rectangle++(0.04,0.04)
  (2.613,9.945)rectangle++(0.04,0.04) (2.616,9.945)rectangle++(0.04,0.04) (2.62,9.945)rectangle++(0.04,0.04) (2.156,4.474)rectangle++(0.04,0.04) (2.009,5.313)rectangle++(0.04,0.04)
  (0.915,2.356)rectangle++(0.04,0.04) (0.246,2.315)rectangle++(0.04,0.04) (0.3,3.554)rectangle++(0.04,0.04) (0.395,4.605)rectangle++(0.04,0.04) (0.518,5.495)rectangle++(0.04,0.04)
  (0.658,6.246)rectangle++(0.04,0.04) (-1.164,4.304)rectangle++(0.04,0.04) (-0.82,5.301)rectangle++(0.04,0.04) (-0.487,6.135)rectangle++(0.04,0.04) (-0.171,6.83)rectangle++(0.04,0.04)
  (0.126,7.409)rectangle++(0.04,0.04) (0.401,7.889)rectangle++(0.04,0.04) (0.654,8.286)rectangle++(0.04,0.04) (0.885,8.613)rectangle++(0.04,0.04) (1.095,8.882)rectangle++(0.04,0.04)
  (1.284,9.102)rectangle++(0.04,0.04) (1.453,9.281)rectangle++(0.04,0.04) (1.604,9.427)rectangle++(0.04,0.04) (1.738,9.545)rectangle++(0.04,0.04) (1.857,9.64)rectangle++(0.04,0.04)
  (1.962,9.716)rectangle++(0.04,0.04) (2.054,9.776)rectangle++(0.04,0.04) (2.135,9.824)rectangle++(0.04,0.04) (2.205,9.861)rectangle++(0.04,0.04) (2.267,9.89)rectangle++(0.04,0.04)
  (2.32,9.912)rectangle++(0.04,0.04) (2.366,9.929)rectangle++(0.04,0.04) (2.406,9.941)rectangle++(0.04,0.04) (2.441,9.95)rectangle++(0.04,0.04) (2.471,9.956)rectangle++(0.04,0.04)
  (2.496,9.96)rectangle++(0.04,0.04) (2.518,9.962)rectangle++(0.04,0.04) (2.536,9.963)rectangle++(0.04,0.04) (2.552,9.964)rectangle++(0.04,0.04) (2.566,9.963)rectangle++(0.04,0.04)
  (2.577,9.962)rectangle++(0.04,0.04) (2.587,9.961)rectangle++(0.04,0.04) (2.595,9.96)rectangle++(0.04,0.04) (2.602,9.958)rectangle++(0.04,0.04) (2.608,9.956)rectangle++(0.04,0.04)
  (2.613,9.955)rectangle++(0.04,0.04) (-1.601,5.403)rectangle++(0.04,0.04) (-1.147,6.253)rectangle++(0.04,0.04) (-0.727,6.957)rectangle++(0.04,0.04) (-0.342,7.538)rectangle++(0.04,0.04)
  (0.009,8.018)rectangle++(0.04,0.04) (0.326,8.411)rectangle++(0.04,0.04) (0.611,8.732)rectangle++(0.04,0.04) (0.867,8.994)rectangle++(0.04,0.04) (1.094,9.207)rectangle++(0.04,0.04)
  (1.296,9.378)rectangle++(0.04,0.04) (1.236,2.225)rectangle++(0.04,0.04) (1.137,3.438)rectangle++(0.04,0.04) (0.718,2.883)rectangle++(0.04,0.04) (0.63,2.681)rectangle++(0.04,0.04)
  (0.59,2.622)rectangle++(0.04,0.04) (0.604,3.802)rectangle++(0.04,0.04) (0.664,4.803)rectangle++(0.04,0.04) (1.112,3.039)rectangle++(0.04,0.04) (0,1.996)rectangle++(0.04,0.04)
  (0.077,3.293)rectangle++(0.04,0.04) (0.195,4.392)rectangle++(0.04,0.04) (0.339,5.322)rectangle++(0.04,0.04) (0.499,6.106)rectangle++(0.04,0.04) (0.666,6.766)rectangle++(0.04,0.04)
  (0.835,7.321)rectangle++(0.04,0.04) (1,7.786)rectangle++(0.04,0.04) (1.159,8.174)rectangle++(0.04,0.04) (1.31,8.498)rectangle++(0.04,0.04) (1.451,8.767)rectangle++(0.04,0.04)
  (1.582,8.99)rectangle++(0.04,0.04) (2.041,4.055)rectangle++(0.04,0.04) (-0.222,3.482)rectangle++(0.04,0.04) (-0.051,4.564)rectangle++(0.04,0.04) (1.136,2.831)rectangle++(0.04,0.04)
  (1.077,3.957)rectangle++(0.04,0.04) (0.005,2.121)rectangle++(0.04,0.04) (-0.406,3.367)rectangle++(0.04,0.04) (-0.212,4.475)rectangle++(0.04,0.04) (-0.004,5.408)rectangle++(0.04,0.04)
  (0.211,6.193)rectangle++(0.04,0.04) (0.425,6.852)rectangle++(0.04,0.04) (0.633,7.403)rectangle++(0.04,0.04) (0.832,7.864)rectangle++(0.04,0.04) (1.02,8.247)rectangle++(0.04,0.04)
  (1.194,8.566)rectangle++(0.04,0.04) (0.035,2.748)rectangle++(0.04,0.04) (0.137,3.93)rectangle++(0.04,0.04) (0.272,4.931)rectangle++(0.04,0.04) (0.426,5.777)rectangle++(0.04,0.04)
  (0.591,6.49)rectangle++(0.04,0.04) (0.76,7.089)rectangle++(0.04,0.04) (0.927,7.591)rectangle++(0.04,0.04) (1.09,8.012)rectangle++(0.04,0.04) (1.244,8.363)rectangle++(0.04,0.04)
  (1.39,8.655)rectangle++(0.04,0.04) (1.526,8.897)rectangle++(0.04,0.04) (-1.571,4.99)rectangle++(0.04,0.04) (-1.138,5.901)rectangle++(0.04,0.04) (-0.734,6.657)rectangle++(0.04,0.04)
  (-1.501,4.124)rectangle++(0.04,0.04) (-1.113,5.162)rectangle++(0.04,0.04) (-0.742,6.028)rectangle++(0.04,0.04) (-0.391,6.75)rectangle++(0.04,0.04) (-0.065,7.349)rectangle++(0.04,0.04)
  (0.237,7.846)rectangle++(0.04,0.04) (0.513,8.256)rectangle++(0.04,0.04) (0.764,8.593)rectangle++(0.04,0.04) (0.991,8.87)rectangle++(0.04,0.04) (1.195,9.096)rectangle++(0.04,0.04)
  (1.377,9.28)rectangle++(0.04,0.04) (1.54,9.429)rectangle++(0.04,0.04) (1.684,9.549)rectangle++(0.04,0.04) (1.811,9.646)rectangle++(0.04,0.04) (1.923,9.723)rectangle++(0.04,0.04)
  (2.021,9.784)rectangle++(0.04,0.04) (2.107,9.831)rectangle++(0.04,0.04) (2.182,9.869)rectangle++(0.04,0.04) (2.247,9.897)rectangle++(0.04,0.04) (2.304,9.919)rectangle++(0.04,0.04)
  (2.353,9.935)rectangle++(0.04,0.04) (2.395,9.947)rectangle++(0.04,0.04) (1.321,2.314)rectangle++(0.04,0.04) (1.213,3.51)rectangle++(0.04,0.04) (1.169,4.531)rectangle++(0.04,0.04)
  (1.173,5.401)rectangle++(0.04,0.04) (1.211,6.14)rectangle++(0.04,0.04) (1.273,6.767)rectangle++(0.04,0.04) (1.35,7.297)rectangle++(0.04,0.04) (1.437,7.745)rectangle++(0.04,0.04)
  (1.529,8.122)rectangle++(0.04,0.04) (1.623,8.438)rectangle++(0.04,0.04) (1.715,8.704)rectangle++(0.04,0.04) (1.803,8.926)rectangle++(0.04,0.04) (1.888,9.111)rectangle++(0.04,0.04)
  (2.035,4.148)rectangle++(0.04,0.04) (1.894,5.04)rectangle++(0.04,0.04) (1.026,3.364)rectangle++(0.04,0.04) (1.005,4.415)rectangle++(0.04,0.04) (1.028,5.308)rectangle++(0.04,0.04)
  (1.084,6.067)rectangle++(0.04,0.04) (1.162,6.71)rectangle++(0.04,0.04) (1.254,7.253)rectangle++(0.04,0.04) (1.354,7.711)rectangle++(0.04,0.04) (0.029,2.631)rectangle++(0.04,0.04)
  (0.128,3.832)rectangle++(0.04,0.04) (0.26,4.848)rectangle++(0.04,0.04) (0.413,5.707)rectangle++(0.04,0.04) (-1.094,4.152)rectangle++(0.04,0.04) (-0.766,5.169)rectangle++(0.04,0.04)
  (-0.446,6.02)rectangle++(0.04,0.04) (-0.141,6.731)rectangle++(0.04,0.04) (0.147,7.323)rectangle++(0.04,0.04) (0.416,7.815)rectangle++(0.04,0.04) (0.664,8.223)rectangle++(0.04,0.04)
  (0.891,8.559)rectangle++(0.04,0.04) (1.098,8.836)rectangle++(0.04,0.04) (1.284,9.062)rectangle++(0.04,0.04) (1.452,9.248)rectangle++(0.04,0.04) (1.602,9.399)rectangle++(0.04,0.04)
  (1.735,9.521)rectangle++(0.04,0.04) (1.854,9.62)rectangle++(0.04,0.04) (1.958,9.699)rectangle++(0.04,0.04) (2.05,9.762)rectangle++(0.04,0.04) (2.131,9.812)rectangle++(0.04,0.04)
  (2.202,9.851)rectangle++(0.04,0.04) (2.263,9.881)rectangle++(0.04,0.04) (2.317,9.905)rectangle++(0.04,0.04) (2.363,9.923)rectangle++(0.04,0.04) (2.403,9.936)rectangle++(0.04,0.04)
  (2.438,9.946)rectangle++(0.04,0.04) (2.468,9.952)rectangle++(0.04,0.04) (2.494,9.957)rectangle++(0.04,0.04) (2.516,9.96)rectangle++(0.04,0.04) (2.535,9.962)rectangle++(0.04,0.04)
  (2.551,9.962)rectangle++(0.04,0.04) (2.564,9.962)rectangle++(0.04,0.04) (0.584,3.177)rectangle++(0.04,0.04) (0.621,4.273)rectangle++(0.04,0.04) (0.697,5.203)rectangle++(0.04,0.04)
  (0.798,5.991)rectangle++(0.04,0.04) (0.916,6.657)rectangle++(0.04,0.04) (1.043,7.218)rectangle++(0.04,0.04) (1.173,7.69)rectangle++(0.04,0.04) (1.302,8.086)rectangle++(0.04,0.04)
  (1.428,8.417)rectangle++(0.04,0.04) (1.548,8.693)rectangle++(0.04,0.04) (1.662,8.924)rectangle++(0.04,0.04) (1.767,9.115)rectangle++(0.04,0.04) (1.864,9.273)rectangle++(0.04,0.04)
  (2.078,4.173)rectangle++(0.04,0.04) (1.931,5.061)rectangle++(0.04,0.04) (1.842,5.82)rectangle++(0.04,0.04) (1.796,6.47)rectangle++(0.04,0.04) (1.783,7.024)rectangle++(0.04,0.04)
  (1.794,7.495)rectangle++(0.04,0.04) (1.823,7.895)rectangle++(0.04,0.04) (1.863,8.234)rectangle++(0.04,0.04) (1.911,8.521)rectangle++(0.04,0.04) (1.963,8.762)rectangle++(0.04,0.04)
  (2.017,8.966)rectangle++(0.04,0.04) (0.038,2.802)rectangle++(0.04,0.04) (0.142,3.976)rectangle++(0.04,0.04) (0.277,4.971)rectangle++(0.04,0.04) (0.432,5.81)rectangle++(0.04,0.04)
  (0.598,6.517)rectangle++(0.04,0.04) (0.767,7.112)rectangle++(0.04,0.04) (1.671,3.508)rectangle++(0.04,0.04) (1.559,4.511)rectangle++(0.04,0.04) (1.503,5.368)rectangle++(0.04,0.04)
  (1.49,6.099)rectangle++(0.04,0.04) (0.019,2.412)rectangle++(0.04,0.04) (0.111,3.646)rectangle++(0.04,0.04) (0.238,4.691)rectangle++(0.04,0.04) (-0.909,3.918)rectangle++(0.04,0.04)
  (-0.618,4.963)rectangle++(0.04,0.04) (-0.329,5.84)rectangle++(0.04,0.04) (-0.049,6.573)rectangle++(0.04,0.04) (0.219,7.185)rectangle++(0.04,0.04) (-1.442,4.407)rectangle++(0.04,0.04)
  (-1.052,5.4)rectangle++(0.04,0.04) (-0.68,6.228)rectangle++(0.04,0.04) (-0.331,6.917)rectangle++(0.04,0.04) (-0.007,7.489)rectangle++(0.04,0.04) (1.855,3.405)rectangle++(0.04,0.04)
  (1.711,4.416)rectangle++(0.04,0.04) (1.629,5.281)rectangle++(0.04,0.04) (1.594,6.02)rectangle++(0.04,0.04) (1.304,3.485)rectangle++(0.04,0.04) (1.246,4.506)rectangle++(0.04,0.04)
  (0.009,2.196)rectangle++(0.04,0.04) (0.093,3.462)rectangle++(0.04,0.04) (0.216,4.535)rectangle++(0.04,0.04) (1.098,2.886)rectangle++(0.04,0.04) (-0.055,3.032)rectangle++(0.04,0.04)
  (0.073,4.175)rectangle++(0.04,0.04) (0.227,5.142)rectangle++(0.04,0.04) (0.396,5.958)rectangle++(0.04,0.04) (0.573,6.645)rectangle++(0.04,0.04) (0.75,7.221)rectangle++(0.04,0.04)
  (0.925,7.704)rectangle++(0.04,0.04) (-1.426,4.642)rectangle++(0.04,0.04) (-1.197,3.597)rectangle++(0.04,0.04) (-0.876,4.701)rectangle++(0.04,0.04) (-0.559,5.627)rectangle++(0.04,0.04)
  (-0.252,6.402)rectangle++(0.04,0.04) (0.04,7.048)rectangle++(0.04,0.04) (0.313,7.585)rectangle++(0.04,0.04) (0.568,8.031)rectangle++(0.04,0.04) (0.802,8.4)rectangle++(0.04,0.04)
  (1.015,8.704)rectangle++(0.04,0.04) (1.209,8.954)rectangle++(0.04,0.04) (1.383,9.159)rectangle++(0.04,0.04) (1.54,9.326)rectangle++(0.04,0.04) (-1.66,5.077)rectangle++(0.04,0.04)
  (-1.21,5.978)rectangle++(0.04,0.04) (1.622,2.84)rectangle++(0.04,0.04) (1.49,3.946)rectangle++(0.04,0.04) (1.422,4.89)rectangle++(0.04,0.04) (1.402,5.696)rectangle++(0.04,0.04)
  (1.418,6.382)rectangle++(0.04,0.04) (1.414,3.515)rectangle++(0.04,0.04) (1.341,4.527)rectangle++(0.04,0.04) (0.964,3.14)rectangle++(0.04,0.04) (0.943,4.226)rectangle++(0.04,0.04)
  (0.968,5.151)rectangle++(0.04,0.04) (1.027,5.936)rectangle++(0.04,0.04) (1.108,6.6)rectangle++(0.04,0.04) (1.204,7.162)rectangle++(0.04,0.04) (1.307,7.636)rectangle++(0.04,0.04)
  (1.415,8.034)rectangle++(0.04,0.04) (1.824,3.833)rectangle++(0.04,0.04) (1.701,4.781)rectangle++(0.04,0.04) (1.635,5.592)rectangle++(0.04,0.04) (1.611,6.284)rectangle++(0.04,0.04)
  (-1.001,4.49)rectangle++(0.04,0.04) (-0.674,5.453)rectangle++(0.04,0.04) (0.282,2.239)rectangle++(0.04,0.04) (0.327,3.488)rectangle++(0.04,0.04) (-0.638,3.698)rectangle++(0.04,0.04)
  (-0.397,4.765)rectangle++(0.04,0.04) (-0.149,5.662)rectangle++(0.04,0.04) (-1.185,4.037)rectangle++(0.04,0.04) (-0.848,5.075)rectangle++(0.04,0.04) (-0.52,5.944)rectangle++(0.04,0.04)
  (-0.206,6.67)rectangle++(0.04,0.04) (0.089,7.274)rectangle++(0.04,0.04) (0.365,7.775)rectangle++(0.04,0.04) (0.619,8.191)rectangle++(0.04,0.04) (-1.584,4.681)rectangle++(0.04,0.04)
  (1.345,2.505)rectangle++(0.04,0.04) (1.241,3.672)rectangle++(0.04,0.04) (1.199,4.668)rectangle++(0.04,0.04) (1.204,5.516)rectangle++(0.04,0.04) (1.242,6.236)rectangle++(0.04,0.04)
  (1.303,6.847)rectangle++(0.04,0.04) (1.379,7.364)rectangle++(0.04,0.04) (1.465,7.801)rectangle++(0.04,0.04) (-1.35,4.762)rectangle++(0.04,0.04) (-0.959,5.698)rectangle++(0.04,0.04)
  (-1.338,3.893)rectangle++(0.04,0.04) (-0.983,4.959)rectangle++(0.04,0.04) (-0.64,5.85)rectangle++(0.04,0.04) (1.524,2.945)rectangle++(0.04,0.04) (1.411,4.039)rectangle++(0.04,0.04)
  (1.358,4.973)rectangle++(0.04,0.04) (1.351,5.769)rectangle++(0.04,0.04) (1.377,6.445)rectangle++(0.04,0.04) (1.426,7.02)rectangle++(0.04,0.04) (1.571,3.64)rectangle++(0.04,0.04)
  (0.717,3.021)rectangle++(0.04,0.04) (0.728,4.135)rectangle++(0.04,0.04) (0.782,5.082)rectangle++(0.04,0.04) (0.866,5.885)rectangle++(0.04,0.04) (0.969,6.563)rectangle++(0.04,0.04)
  (1.084,7.136)rectangle++(0.04,0.04) (-1.277,4.56)rectangle++(0.04,0.04) (-0.905,5.523)rectangle++(0.04,0.04) (-0.551,6.327)rectangle++(0.04,0.04) (-0.217,6.996)rectangle++(0.04,0.04)
  (0.093,7.552)rectangle++(0.04,0.04) (0.379,8.011)rectangle++(0.04,0.04) (0.64,8.391)rectangle++(0.04,0.04) (0.878,8.703)rectangle++(0.04,0.04) (1.092,8.958)rectangle++(0.04,0.04)
  (1.284,9.167)rectangle++(0.04,0.04) (1.456,9.337)rectangle++(0.04,0.04) (1.609,9.474)rectangle++(0.04,0.04) (1.744,9.585)rectangle++(0.04,0.04) (1.864,9.674)rectangle++(0.04,0.04)
  (1.969,9.744)rectangle++(0.04,0.04) (2.061,9.8)rectangle++(0.04,0.04) (2.142,9.844)rectangle++(0.04,0.04) (2.212,9.878)rectangle++(0.04,0.04) (2.273,9.904)rectangle++(0.04,0.04)
  (2.326,9.923)rectangle++(0.04,0.04) (2.372,9.938)rectangle++(0.04,0.04) (2.183,4.417)rectangle++(0.04,0.04) (2.03,5.263)rectangle++(0.04,0.04) (1.934,5.989)rectangle++(0.04,0.04)
  (1.881,6.609)rectangle++(0.04,0.04) (1.861,7.139)rectangle++(0.04,0.04) (1.865,7.59)rectangle++(0.04,0.04) (1.887,7.973)rectangle++(0.04,0.04) (1.921,8.298)rectangle++(0.04,0.04)
  (1.962,8.572)rectangle++(0.04,0.04) (2.009,8.804)rectangle++(0.04,0.04) (2.057,8.999)rectangle++(0.04,0.04) (2.106,9.163)rectangle++(0.04,0.04) (1.245,2.301)rectangle++(0.04,0.04)
  (1.148,3.502)rectangle++(0.04,0.04) (1.114,4.527)rectangle++(0.04,0.04) (0.991,3.088)rectangle++(0.04,0.04) (0.963,4.181)rectangle++(0.04,0.04) (0.984,5.112)rectangle++(0.04,0.04)
  (1.039,5.902)rectangle++(0.04,0.04) (1.117,6.571)rectangle++(0.04,0.04) (1.21,7.137)rectangle++(0.04,0.04) (1.312,7.614)rectangle++(0.04,0.04) (1.417,8.016)rectangle++(0.04,0.04)
  (1.523,8.353)rectangle++(0.04,0.04) (1.627,8.635)rectangle++(0.04,0.04) (1.726,8.871)rectangle++(0.04,0.04) (1.82,9.068)rectangle++(0.04,0.04) (1.907,9.231)rectangle++(0.04,0.04)
  (1.988,9.366)rectangle++(0.04,0.04) (2.062,9.478)rectangle++(0.04,0.04) (2.13,9.57)rectangle++(0.04,0.04) (2.191,9.645)rectangle++(0.04,0.04) (2.246,9.707)rectangle++(0.04,0.04)
  (2.295,9.757)rectangle++(0.04,0.04) (2.339,9.798)rectangle++(0.04,0.04) (2.152,4.382)rectangle++(0.04,0.04) (2.003,5.235)rectangle++(0.04,0.04) (1.909,5.966)rectangle++(0.04,0.04)
  (1.859,6.591)rectangle++(0.04,0.04) (1.842,7.124)rectangle++(0.04,0.04) (1.848,7.578)rectangle++(0.04,0.04) (1.872,7.964)rectangle++(0.04,0.04) (1.908,8.29)rectangle++(0.04,0.04)
  (1.83,3.994)rectangle++(0.04,0.04) (1.713,4.918)rectangle++(0.04,0.04) (1.65,5.708)rectangle++(0.04,0.04) (1.629,6.382)rectangle++(0.04,0.04) (1.638,6.956)rectangle++(0.04,0.04)
  (1.668,7.443)rectangle++(0.04,0.04) (1.713,7.856)rectangle++(0.04,0.04) (1.768,8.205)rectangle++(0.04,0.04) (1.829,8.5)rectangle++(0.04,0.04) (1.893,8.748)rectangle++(0.04,0.04)
  (1.956,8.956)rectangle++(0.04,0.04) (2.019,9.131)rectangle++(0.04,0.04) (2.079,9.277)rectangle++(0.04,0.04) (2.136,9.398)rectangle++(0.04,0.04) (2.189,9.499)rectangle++(0.04,0.04)
  (2.239,9.583)rectangle++(0.04,0.04) (2.284,9.652)rectangle++(0.04,0.04) (2.325,9.709)rectangle++(0.04,0.04) (-1.6,5.3)rectangle++(0.04,0.04) (-1.15,6.165)rectangle++(0.04,0.04)
  (-0.733,6.883)rectangle++(0.04,0.04) (-0.35,7.476)rectangle++(0.04,0.04) (-0.001,7.965)rectangle++(0.04,0.04) (0.316,8.366)rectangle++(0.04,0.04) (0.601,8.695)rectangle++(0.04,0.04)
  (0.856,8.963)rectangle++(0.04,0.04) (1.084,9.18)rectangle++(0.04,0.04) (1.287,9.356)rectangle++(0.04,0.04) (1.466,9.497)rectangle++(0.04,0.04) (1.623,9.61)rectangle++(0.04,0.04)
  (1.762,9.7)rectangle++(0.04,0.04) (1.884,9.771)rectangle++(0.04,0.04) (1.99,9.826)rectangle++(0.04,0.04) (-1.685,5.26)rectangle++(0.04,0.04) (-1.224,6.135)rectangle++(0.04,0.04)
  (-0.797,6.86)rectangle++(0.04,0.04) (-0.406,7.459)rectangle++(0.04,0.04) (-0.049,7.953)rectangle++(0.04,0.04) (0.275,8.358)rectangle++(0.04,0.04) (0.566,8.689)rectangle++(0.04,0.04)
  (0.826,8.96)rectangle++(0.04,0.04) (1.058,9.179)rectangle++(0.04,0.04) (1.265,9.356)rectangle++(0.04,0.04) (1.447,9.498)rectangle++(0.04,0.04) (1.608,9.612)rectangle++(0.04,0.04)
  (1.749,9.702)rectangle++(0.04,0.04) (-1.702,5.191)rectangle++(0.04,0.04) (1.485,2.577)rectangle++(0.04,0.04) (1.363,3.727)rectangle++(0.04,0.04) (1.305,4.71)rectangle++(0.04,0.04)
  (1.296,5.547)rectangle++(0.04,0.04) (1.321,6.26)rectangle++(0.04,0.04) (1.371,6.864)rectangle++(0.04,0.04) (1.438,7.376)rectangle++(0.04,0.04) (1.658,3.711)rectangle++(0.04,0.04)
  (1.556,4.684)rectangle++(0.04,0.04) (1.507,5.515)rectangle++(0.04,0.04) (1.5,6.224)rectangle++(0.04,0.04) (1.522,6.827)rectangle++(0.04,0.04) (1.564,7.338)rectangle++(0.04,0.04)
  (1.621,7.771)rectangle++(0.04,0.04) (1.686,8.137)rectangle++(0.04,0.04) (1.757,8.445)rectangle++(0.04,0.04) (1.829,8.704)rectangle++(0.04,0.04) (1.9,8.922)rectangle++(0.04,0.04)
  (1.97,9.103)rectangle++(0.04,0.04) (2.036,9.255)rectangle++(0.04,0.04) (2.099,9.382)rectangle++(0.04,0.04) (2.077,4.248)rectangle++(0.04,0.04) (1.933,5.124)rectangle++(0.04,0.04)
  (1.846,5.874)rectangle++(0.04,0.04) (1.802,6.515)rectangle++(0.04,0.04) (1.402,3.628)rectangle++(0.04,0.04) (1.335,4.624)rectangle++(0.04,0.04) (1.317,5.473)rectangle++(0.04,0.04)
  (1.201,3.317)rectangle++(0.04,0.04) (1.151,4.367)rectangle++(0.04,0.04) (0.947,3.066)rectangle++(0.04,0.04) (0.925,4.164)rectangle++(0.04,0.04) (0.951,5.099)rectangle++(0.04,0.04)
  (1.01,5.892)rectangle++(0.04,0.04) (1.092,6.564)rectangle++(0.04,0.04) (1.188,7.132)rectangle++(0.04,0.04) (1.293,7.611)rectangle++(0.04,0.04) (1.402,8.014)rectangle++(0.04,0.04)
  (1.51,8.352)rectangle++(0.04,0.04) (1.615,8.635)rectangle++(0.04,0.04) (1.716,8.871)rectangle++(0.04,0.04) (1.811,9.068)rectangle++(0.04,0.04) (1.9,9.232)rectangle++(0.04,0.04)
  (-1.575,5.126)rectangle++(0.04,0.04) (-1.136,6.017)rectangle++(0.04,0.04) (-0.727,6.756)rectangle++(0.04,0.04) (-0.35,7.368)rectangle++(0.04,0.04) (-0.005,7.873)rectangle++(0.04,0.04)
  (0.308,8.288)rectangle++(0.04,0.04) (0.591,8.629)rectangle++(0.04,0.04) (0.846,8.907)rectangle++(0.04,0.04) (1.073,9.133)rectangle++(0.04,0.04) (1.275,9.317)rectangle++(0.04,0.04)
  (1.454,9.464)rectangle++(0.04,0.04) (1.613,9.583)rectangle++(0.04,0.04) (1.752,9.677)rectangle++(0.04,0.04) (-1.696,5.186)rectangle++(0.04,0.04) (-1.237,6.072)rectangle++(0.04,0.04)
  (0.046,3.607)rectangle++(0.04,0.04) (0.181,4.66)rectangle++(0.04,0.04) (1.133,2.902)rectangle++(0.04,0.04) (1.077,4.018)rectangle++(0.04,0.04) (-0.616,3.943)rectangle++(0.04,0.04)
  (-0.368,4.972)rectangle++(0.04,0.04) (-0.116,5.837)rectangle++(0.04,0.04) (0.132,6.563)rectangle++(0.04,0.04) (0.373,7.169)rectangle++(0.04,0.04) (0.601,7.675)rectangle++(0.04,0.04)
  (0.816,8.096)rectangle++(0.04,0.04) (1.015,8.445)rectangle++(0.04,0.04) (-1.567,4.805)rectangle++(0.04,0.04) (-1.142,5.743)rectangle++(0.04,0.04) (-0.743,6.524)rectangle++(0.04,0.04)
  (-0.373,7.171)rectangle++(0.04,0.04) (-0.032,7.706)rectangle++(0.04,0.04) (0.279,8.148)rectangle++(0.04,0.04) (0.561,8.511)rectangle++(0.04,0.04) (0.815,8.808)rectangle++(0.04,0.04)
  (1.043,9.05)rectangle++(0.04,0.04) (1.246,9.247)rectangle++(0.04,0.04) (2.144,4.028)rectangle++(0.04,0.04) (1.981,4.934)rectangle++(0.04,0.04) (1.879,5.711)rectangle++(0.04,0.04)
  (1.824,6.376)rectangle++(0.04,0.04) (1.803,6.942)rectangle++(0.04,0.04) (-1.107,4.656)rectangle++(0.04,0.04) (-1.143,3.659)rectangle++(0.04,0.04) (-0.827,4.752)rectangle++(0.04,0.04)
  (-0.515,5.669)rectangle++(0.04,0.04) (-0.214,6.435)rectangle++(0.04,0.04) (1.619,3.155)rectangle++(0.04,0.04) (1.5,4.213)rectangle++(0.04,0.04) (1.441,5.117)rectangle++(0.04,0.04)
  (1.098,3.276)rectangle++(0.04,0.04) (1.063,4.337)rectangle++(0.04,0.04) (1.074,5.24)rectangle++(0.04,0.04) (1.121,6.007)rectangle++(0.04,0.04) (1.191,6.658)rectangle++(0.04,0.04)
  (1.276,7.208)rectangle++(0.04,0.04) (1.371,7.672)rectangle++(0.04,0.04) (1.47,8.062)rectangle++(0.04,0.04) (1.57,8.39)rectangle++(0.04,0.04) (1.668,8.665)rectangle++(0.04,0.04)
  (1.762,8.895)rectangle++(0.04,0.04) (1.996,4.077)rectangle++(0.04,0.04) (1.858,4.982)rectangle++(0.04,0.04) (-0.681,4.279)rectangle++(0.04,0.04) (-0.41,5.261)rectangle++(0.04,0.04)
  (-0.14,6.084)rectangle++(0.04,0.04) (0.122,6.774)rectangle++(0.04,0.04) (0.373,7.349)rectangle++(0.04,0.04) (0.608,7.828)rectangle++(0.04,0.04) (0.828,8.226)rectangle++(0.04,0.04)
  (1.031,8.555)rectangle++(0.04,0.04) (1.216,8.827)rectangle++(0.04,0.04) (1.385,9.05)rectangle++(0.04,0.04) (1.537,9.233)rectangle++(0.04,0.04) (1.673,9.383)rectangle++(0.04,0.04)
  (1.795,9.505)rectangle++(0.04,0.04) (2.143,4.203)rectangle++(0.04,0.04) (1.988,5.083)rectangle++(0.04,0.04) (1.891,5.837)rectangle++(0.04,0.04) (-0.856,4.454)rectangle++(0.04,0.04)
  (-0.552,5.417)rectangle++(0.04,0.04) (-0.255,6.222)rectangle++(0.04,0.04) (0.03,6.895)rectangle++(0.04,0.04) (1.704,3.299)rectangle++(0.04,0.04) (1.578,4.332)rectangle++(0.04,0.04)
  (1.513,5.215)rectangle++(0.04,0.04) (1.492,5.969)rectangle++(0.04,0.04) (1.505,6.61)rectangle++(0.04,0.04) (1.541,7.154)rectangle++(0.04,0.04) (1.594,7.616)rectangle++(0.04,0.04)
  (1.657,8.006)rectangle++(0.04,0.04) (1.788,3.882)rectangle++(0.04,0.04) (0.752,3.117)rectangle++(0.04,0.04) (0.684,2.734)rectangle++(0.04,0.04) (0.688,3.893)rectangle++(0.04,0.04)
  (-0.659,3.845)rectangle++(0.04,0.04) (-0.409,4.891)rectangle++(0.04,0.04) (-1.068,3.837)rectangle++(0.04,0.04) (-0.756,4.9)rectangle++(0.04,0.04) (-1.132,3.773)rectangle++(0.04,0.04)
  (-0.813,4.849)rectangle++(0.04,0.04) (-1.131,3.752)rectangle++(0.04,0.04) (-0.814,4.831)rectangle++(0.04,0.04) (-0.501,5.735)rectangle++(0.04,0.04) (-0.198,6.491)rectangle++(0.04,0.04)
  (0.089,7.121)rectangle++(0.04,0.04) (0.358,7.646)rectangle++(0.04,0.04) (0.608,8.081)rectangle++(0.04,0.04) (0.838,8.441)rectangle++(0.04,0.04) (1.048,8.737)rectangle++(0.04,0.04)
  (2.041,3.885)rectangle++(0.04,0.04) (1.888,4.817)rectangle++(0.04,0.04) (1.795,5.615)rectangle++(0.04,0.04) (-0.826,4.393)rectangle++(0.04,0.04) (-0.529,5.363)rectangle++(0.04,0.04)
  (-0.237,6.176)rectangle++(0.04,0.04) (0.043,6.855)rectangle++(0.04,0.04) (1.693,3.294)rectangle++(0.04,0.04) (1.568,4.328)rectangle++(0.04,0.04) (1.504,5.213)rectangle++(0.04,0.04)
  (-0.793,4.259)rectangle++(0.04,0.04) (-0.506,5.248)rectangle++(0.04,0.04) (-0.222,6.077)rectangle++(0.04,0.04) (0.052,6.771)rectangle++(0.04,0.04) (0.313,7.349)rectangle++(0.04,0.04)
  (0.557,7.83)rectangle++(0.04,0.04) (0.785,8.23)rectangle++(0.04,0.04) (0.994,8.56)rectangle++(0.04,0.04) (-0.139,3.907)rectangle++(0.04,0.04) (0.984,2.684)rectangle++(0.04,0.04)
  (0.942,3.838)rectangle++(0.04,0.04) (0.005,2.105)rectangle++(0.04,0.04) (0.086,3.385)rectangle++(0.04,0.04) (0.206,4.47)rectangle++(0.04,0.04) (0.352,5.387)rectangle++(0.04,0.04)
  (0.512,6.161)rectangle++(0.04,0.04) (0.68,6.813)rectangle++(0.04,0.04) (0.848,7.36)rectangle++(0.04,0.04) (1.013,7.818)rectangle++(0.04,0.04) (1.172,8.201)rectangle++(0.04,0.04)
  (0.696,2.408)rectangle++(0.04,0.04) (0.515,2.585)rectangle++(0.04,0.04) (0.539,3.773)rectangle++(0.04,0.04) (0.253,2.464)rectangle++(0.04,0.04) (0.581,2.495)rectangle++(0.04,0.04)
  (0.591,3.694)rectangle++(0.04,0.04) (0.648,4.712)rectangle++(0.04,0.04) (0.737,5.576)rectangle++(0.04,0.04) (0.847,6.306)rectangle++(0.04,0.04) (1.462,3.371)rectangle++(0.04,0.04)
  (1.376,4.403)rectangle++(0.04,0.04) (1.343,5.283)rectangle++(0.04,0.04) (1.351,6.033)rectangle++(0.04,0.04) (1.336,3.432)rectangle++(0.04,0.04) (1.271,4.46)rectangle++(0.04,0.04)
  (1.256,5.337)rectangle++(0.04,0.04) (1.279,6.082)rectangle++(0.04,0.04) (1.328,6.715)rectangle++(0.04,0.04) (1.395,7.251)rectangle++(0.04,0.04) (1.474,7.703)rectangle++(0.04,0.04)
  (1.559,8.085)rectangle++(0.04,0.04) (1.646,8.406)rectangle++(0.04,0.04) (1.733,8.676)rectangle++(0.04,0.04) (1.818,8.901)rectangle++(0.04,0.04) (1.899,9.089)rectangle++(0.04,0.04)
  (1.976,9.246)rectangle++(0.04,0.04) (2.047,9.376)rectangle++(0.04,0.04) (2.113,9.484)rectangle++(0.04,0.04) (2.173,9.573)rectangle++(0.04,0.04) (0.041,2.885)rectangle++(0.04,0.04)
  (0.148,4.047)rectangle++(0.04,0.04) (0.286,5.03)rectangle++(0.04,0.04) (0.442,5.86)rectangle++(0.04,0.04) (0.608,6.56)rectangle++(0.04,0.04) (0.777,7.148)rectangle++(0.04,0.04)
  (0.944,7.641)rectangle++(0.04,0.04) (1.106,8.053)rectangle++(0.04,0.04) (1.26,8.397)rectangle++(0.04,0.04) (-1.513,4.841)rectangle++(0.04,0.04) (-1.094,5.772)rectangle++(0.04,0.04)
  (-0.702,6.546)rectangle++(0.04,0.04) (-0.337,7.188)rectangle++(0.04,0.04) (-0.001,7.72)rectangle++(0.04,0.04) (0.306,8.158)rectangle++(0.04,0.04) (0.584,8.518)rectangle++(0.04,0.04)
  (0.835,8.813)rectangle++(0.04,0.04) (1.06,9.054)rectangle++(0.04,0.04) (1.261,9.25)rectangle++(0.04,0.04) (1.44,9.408)rectangle++(0.04,0.04) (1.598,9.536)rectangle++(0.04,0.04)
  (1.737,9.638)rectangle++(0.04,0.04) (1.86,9.719)rectangle++(0.04,0.04) (1.968,9.783)rectangle++(0.04,0.04) (2.062,9.833)rectangle++(0.04,0.04) (-1.674,5.275)rectangle++(0.04,0.04)
  (-1.214,6.147)rectangle++(0.04,0.04) (-0.788,6.87)rectangle++(0.04,0.04) (-0.397,7.467)rectangle++(0.04,0.04) (-0.041,7.959)rectangle++(0.04,0.04) (0.281,8.363)rectangle++(0.04,0.04)
  (0.571,8.694)rectangle++(0.04,0.04) (-1.699,4.771)rectangle++(0.04,0.04) (-1.256,5.719)rectangle++(0.04,0.04) (-0.841,6.508)rectangle++(0.04,0.04) (-0.456,7.162)rectangle++(0.04,0.04)
  (-0.104,7.702)rectangle++(0.04,0.04) (0.218,8.147)rectangle++(0.04,0.04) (0.509,8.512)rectangle++(0.04,0.04) (0.771,8.811)rectangle++(0.04,0.04) (1.005,9.055)rectangle++(0.04,0.04)
  (1.215,9.253)rectangle++(0.04,0.04) (-1.702,5.002)rectangle++(0.04,0.04) (-1.249,5.916)rectangle++(0.04,0.04) (-0.827,6.675)rectangle++(0.04,0.04) (-0.438,7.303)rectangle++(0.04,0.04)
  (-0.083,7.821)rectangle++(0.04,0.04) (0.24,8.247)rectangle++(0.04,0.04) (0.532,8.597)rectangle++(0.04,0.04) (0.794,8.882)rectangle++(0.04,0.04) (1.028,9.114)rectangle++(0.04,0.04)
  (1.236,9.302)rectangle++(0.04,0.04) (1.421,9.454)rectangle++(0.04,0.04) (1.583,9.575)rectangle++(0.04,0.04) (2.186,4.168)rectangle++(0.04,0.04) (2.022,5.051)rectangle++(0.04,0.04)
  (-0.666,4.324)rectangle++(0.04,0.04) (-0.396,5.298)rectangle++(0.04,0.04) (1.358,2.894)rectangle++(0.04,0.04) (1.268,4.002)rectangle++(0.04,0.04) (1.236,4.947)rectangle++(0.04,0.04)
  (1.246,5.752)rectangle++(0.04,0.04) (1.287,6.435)rectangle++(0.04,0.04) (-1.092,4.459)rectangle++(0.04,0.04) (-0.752,5.43)rectangle++(0.04,0.04) (-0.424,6.242)rectangle++(0.04,0.04)
  (-0.113,6.919)rectangle++(0.04,0.04) (0.179,7.482)rectangle++(0.04,0.04) (0.449,7.949)rectangle++(0.04,0.04) (0.697,8.335)rectangle++(0.04,0.04) (0.924,8.653)rectangle++(0.04,0.04)
  (2.034,3.84)rectangle++(0.04,0.04) (-0.407,4.087)rectangle++(0.04,0.04) (-0.897,3.678)rectangle++(0.04,0.04) (-0.617,4.759)rectangle++(0.04,0.04) (-0.337,5.666)rectangle++(0.04,0.04)
  (-0.062,6.425)rectangle++(0.04,0.04) (0.202,7.06)rectangle++(0.04,0.04) (0.452,7.589)rectangle++(0.04,0.04) (0.686,8.029)rectangle++(0.04,0.04) (0.003,2.061)rectangle++(0.04,0.04)
  (0.082,3.348)rectangle++(0.04,0.04) (0.202,4.438)rectangle++(0.04,0.04) (0.347,5.361)rectangle++(0.04,0.04) (0.507,6.139)rectangle++(0.04,0.04) (0.674,6.794)rectangle++(0.04,0.04)
  (0.843,7.344)rectangle++(0.04,0.04) (-1.364,4.555)rectangle++(0.04,0.04) (-0.979,5.523)rectangle++(0.04,0.04) (-0.614,6.33)rectangle++(0.04,0.04) (-1.41,4.082)rectangle++(0.04,0.04)
  (-1.037,5.122)rectangle++(0.04,0.04) (-0.679,5.992)rectangle++(0.04,0.04) (-0.34,6.716)rectangle++(0.04,0.04) (-0.022,7.319)rectangle++(0.04,0.04) (0.272,7.818)rectangle++(0.04,0.04)
  (1.904,3.531)rectangle++(0.04,0.04) (1.758,4.522)rectangle++(0.04,0.04) (1.673,5.369)rectangle++(0.04,0.04) (1.634,6.093)rectangle++(0.04,0.04) (-0.957,4.457)rectangle++(0.04,0.04)
  (-0.637,5.423)rectangle++(0.04,0.04) (1.418,2.863)rectangle++(0.04,0.04) (1.317,3.973)rectangle++(0.04,0.04) (0.83,3.027)rectangle++(0.04,0.04) (0.824,4.136)rectangle++(0.04,0.04)
  (0.864,5.079)rectangle++(0.04,0.04) (0.935,5.879)rectangle++(0.04,0.04) (1.028,6.556)rectangle++(0.04,0.04) (1.134,7.128)rectangle++(0.04,0.04) (1.247,7.609)rectangle++(0.04,0.04);
\fill[piece2]
  (1.362,8.014)rectangle++(0.04,0.04) (1.476,8.354)rectangle++(0.04,0.04) (1.586,8.638)rectangle++(0.04,0.04) (1.183,2.263)rectangle++(0.04,0.04) (1.094,3.472)rectangle++(0.04,0.04)
  (1.066,4.504)rectangle++(0.04,0.04) (1.084,5.382)rectangle++(0.04,0.04) (1.135,6.127)rectangle++(0.04,0.04) (1.207,6.759)rectangle++(0.04,0.04) (1.532,3.54)rectangle++(0.04,0.04)
  (1.442,4.544)rectangle++(0.04,0.04) (0.956,3.166)rectangle++(0.04,0.04) (0.937,4.249)rectangle++(0.04,0.04) (0.964,5.17)rectangle++(0.04,0.04) (1.024,5.952)rectangle++(0.04,0.04)
  (-1.036,4.315)rectangle++(0.04,0.04) (-0.71,5.305)rectangle++(0.04,0.04) (-0.394,6.134)rectangle++(0.04,0.04) (-1.329,4.085)rectangle++(0.04,0.04) (-0.968,5.121)rectangle++(0.04,0.04)
  (-0.62,5.988)rectangle++(0.04,0.04) (-0.29,6.711)rectangle++(0.04,0.04) (0.02,7.312)rectangle++(0.04,0.04) (0.307,7.811)rectangle++(0.04,0.04) (0.571,8.223)rectangle++(0.04,0.04)
  (1.969,3.677)rectangle++(0.04,0.04) (1.819,4.643)rectangle++(0.04,0.04) (1.73,5.47)rectangle++(0.04,0.04) (1.687,6.177)rectangle++(0.04,0.04) (1.679,6.779)rectangle++(0.04,0.04)
  (1.696,7.291)rectangle++(0.04,0.04) (1.731,7.726)rectangle++(0.04,0.04) (1.778,8.094)rectangle++(0.04,0.04) (1.833,8.405)rectangle++(0.04,0.04) (1.892,8.667)rectangle++(0.04,0.04)
  (1.953,8.888)rectangle++(0.04,0.04) (2.013,9.072)rectangle++(0.04,0.04) (2.072,9.227)rectangle++(0.04,0.04) (2.128,9.357)rectangle++(0.04,0.04) (2.181,9.464)rectangle++(0.04,0.04)
  (2.088,4.282)rectangle++(0.04,0.04) (1.944,5.152)rectangle++(0.04,0.04) (1.856,5.898)rectangle++(0.04,0.04) (1.812,6.535)rectangle++(0.04,0.04) (1.799,7.079)rectangle++(0.04,0.04)
  (1.81,7.541)rectangle++(0.04,0.04) (1.838,7.934)rectangle++(0.04,0.04) (1.878,8.266)rectangle++(0.04,0.04) (1.924,8.548)rectangle++(0.04,0.04) (1.975,8.785)rectangle++(0.04,0.04)
  (2.028,8.984)rectangle++(0.04,0.04) (2.081,9.152)rectangle++(0.04,0.04) (-1.526,5.144)rectangle++(0.04,0.04) (-1.322,3.677)rectangle++(0.04,0.04) (1.065,2.371)rectangle++(0.04,0.04)
  (-0.269,3.616)rectangle++(0.04,0.04) (-0.086,4.68)rectangle++(0.04,0.04) (1.169,2.845)rectangle++(0.04,0.04) (1.105,3.968)rectangle++(0.04,0.04) (0.853,2.978)rectangle++(0.04,0.04)
  (0.842,4.094)rectangle++(0.04,0.04) (0.877,5.042)rectangle++(0.04,0.04) (0.945,5.847)rectangle++(0.04,0.04) (1.035,6.528)rectangle++(0.04,0.04) (1.139,7.104)rectangle++(0.04,0.04)
  (1.25,7.589)rectangle++(0.04,0.04) (-1.343,4.68)rectangle++(0.04,0.04) (-0.957,5.628)rectangle++(0.04,0.04) (-0.59,6.418)rectangle++(0.04,0.04) (-0.247,7.075)rectangle++(0.04,0.04)
  (0.071,7.62)rectangle++(0.04,0.04) (0.363,8.07)rectangle++(0.04,0.04) (1.956,3.601)rectangle++(0.04,0.04) (1.804,4.579)rectangle++(0.04,0.04) (1.715,5.416)rectangle++(0.04,0.04)
  (1.672,6.131)rectangle++(0.04,0.04) (1.664,6.741)rectangle++(0.04,0.04) (-1.089,4.59)rectangle++(0.04,0.04) (-0.744,5.542)rectangle++(0.04,0.04) (-0.413,6.336)rectangle++(0.04,0.04)
  (-0.1,6.999)rectangle++(0.04,0.04) (0.193,7.549)rectangle++(0.04,0.04) (0.464,8.005)rectangle++(0.04,0.04) (0.712,8.382)rectangle++(0.04,0.04) (0.938,8.692)rectangle++(0.04,0.04)
  (1.143,8.947)rectangle++(0.04,0.04) (1.327,9.156)rectangle++(0.04,0.04) (1.492,9.325)rectangle++(0.04,0.04) (1.639,9.463)rectangle++(0.04,0.04) (1.77,9.574)rectangle++(0.04,0.04)
  (1.885,9.664)rectangle++(0.04,0.04) (0.042,2.897)rectangle++(0.04,0.04) (0.149,4.057)rectangle++(0.04,0.04) (-0.791,3.777)rectangle++(0.04,0.04) (-0.523,4.838)rectangle++(0.04,0.04)
  (-0.253,5.729)rectangle++(0.04,0.04) (0.012,6.476)rectangle++(0.04,0.04) (0.267,7.101)rectangle++(0.04,0.04) (0.508,7.621)rectangle++(0.04,0.04) (0.735,8.054)rectangle++(0.04,0.04)
  (1.908,3.682)rectangle++(0.04,0.04) (1.767,4.65)rectangle++(0.04,0.04) (1.686,5.478)rectangle++(0.04,0.04) (1.65,6.185)rectangle++(0.04,0.04) (1.648,6.787)rectangle++(0.04,0.04)
  (-0.952,3.358)rectangle++(0.04,0.04) (-0.677,4.488)rectangle++(0.04,0.04) (0.009,2.193)rectangle++(0.04,0.04) (0.093,3.46)rectangle++(0.04,0.04) (-0.674,3.648)rectangle++(0.04,0.04)
  (-0.429,4.724)rectangle++(0.04,0.04) (-0.178,5.629)rectangle++(0.04,0.04) (0.072,6.388)rectangle++(0.04,0.04) (0.314,7.023)rectangle++(0.04,0.04) (0.546,7.553)rectangle++(0.04,0.04)
  (-1.462,4.541)rectangle++(0.04,0.04) (-1.063,5.515)rectangle++(0.04,0.04) (-0.685,6.326)rectangle++(0.04,0.04) (-0.332,7.001)rectangle++(0.04,0.04) (-1.502,4.267)rectangle++(0.04,0.04)
  (-1.108,5.284)rectangle++(0.04,0.04) (-0.733,6.132)rectangle++(0.04,0.04) (-0.38,6.837)rectangle++(0.04,0.04) (-0.051,7.423)rectangle++(0.04,0.04) (-1.541,4.403)rectangle++(0.04,0.04)
  (-1.136,5.401)rectangle++(0.04,0.04) (-0.752,6.232)rectangle++(0.04,0.04) (-0.392,6.924)rectangle++(0.04,0.04) (-0.058,7.497)rectangle++(0.04,0.04) (0.248,7.971)rectangle++(0.04,0.04)
  (0.528,8.362)rectangle++(0.04,0.04) (0.781,8.682)rectangle++(0.04,0.04) (1.009,8.945)rectangle++(0.04,0.04) (1.213,9.159)rectangle++(0.04,0.04) (-1.683,4.983)rectangle++(0.04,0.04)
  (-1.233,5.899)rectangle++(0.04,0.04) (-0.815,6.66)rectangle++(0.04,0.04) (-0.428,7.29)rectangle++(0.04,0.04) (-0.075,7.81)rectangle++(0.04,0.04) (0.247,8.237)rectangle++(0.04,0.04)
  (0.537,8.588)rectangle++(0.04,0.04) (0.798,8.875)rectangle++(0.04,0.04) (1.031,9.108)rectangle++(0.04,0.04) (1.238,9.296)rectangle++(0.04,0.04) (1.422,9.449)rectangle++(0.04,0.04)
  (2.173,4.107)rectangle++(0.04,0.04) (2.009,5)rectangle++(0.04,0.04) (1.906,5.766)rectangle++(0.04,0.04) (1.848,6.421)rectangle++(0.04,0.04) (1.826,6.98)rectangle++(0.04,0.04)
  (1.829,7.456)rectangle++(0.04,0.04) (1.851,7.861)rectangle++(0.04,0.04) (-1.293,4.846)rectangle++(0.04,0.04) (-0.908,5.767)rectangle++(0.04,0.04) (-0.543,6.535)rectangle++(0.04,0.04)
  (-0.202,7.172)rectangle++(0.04,0.04) (-1.515,4.326)rectangle++(0.04,0.04) (-1.117,5.334)rectangle++(0.04,0.04) (-0.738,6.174)rectangle++(0.04,0.04) (1.615,2.985)rectangle++(0.04,0.04)
  (-0.001,1.989)rectangle++(0.04,0.04) (0.077,3.287)rectangle++(0.04,0.04) (0.194,4.387)rectangle++(0.04,0.04) (0.339,5.317)rectangle++(0.04,0.04) (0.498,6.102)rectangle++(0.04,0.04)
  (0.665,6.763)rectangle++(0.04,0.04) (0.834,7.318)rectangle++(0.04,0.04) (-1.36,4.549)rectangle++(0.04,0.04) (-0.976,5.517)rectangle++(0.04,0.04) (-0.611,6.325)rectangle++(0.04,0.04)
  (1.638,3.043)rectangle++(0.04,0.04) (1.512,4.117)rectangle++(0.04,0.04) (1.447,5.035)rectangle++(0.04,0.04) (1.43,5.818)rectangle++(0.04,0.04) (0.017,2.374)rectangle++(0.04,0.04)
  (0.586,2.424)rectangle++(0.04,0.04) (0.593,3.633)rectangle++(0.04,0.04) (0.647,4.661)rectangle++(0.04,0.04) (0.734,5.532)rectangle++(0.04,0.04) (0.843,6.269)rectangle++(0.04,0.04)
  (0.965,6.891)rectangle++(0.04,0.04) (1.094,7.415)rectangle++(0.04,0.04) (1.224,7.855)rectangle++(0.04,0.04) (1.353,8.224)rectangle++(0.04,0.04) (-1.459,4.825)rectangle++(0.04,0.04)
  (-1.05,5.756)rectangle++(0.04,0.04) (-0.664,6.53)rectangle++(0.04,0.04) (-0.306,7.174)rectangle++(0.04,0.04) (0.025,7.706)rectangle++(0.04,0.04) (0.327,8.145)rectangle++(0.04,0.04)
  (0.602,8.507)rectangle++(0.04,0.04) (0.85,8.803)rectangle++(0.04,0.04) (1.072,9.045)rectangle++(0.04,0.04) (1.271,9.241)rectangle++(0.04,0.04) (1.448,9.4)rectangle++(0.04,0.04)
  (1.604,9.529)rectangle++(0.04,0.04) (1.743,9.631)rectangle++(0.04,0.04) (1.864,9.713)rectangle++(0.04,0.04) (1.971,9.778)rectangle++(0.04,0.04) (2.064,9.828)rectangle++(0.04,0.04)
  (2.146,9.868)rectangle++(0.04,0.04) (2.216,9.898)rectangle++(0.04,0.04) (2.278,9.921)rectangle++(0.04,0.04) (2.331,9.938)rectangle++(0.04,0.04) (2.376,9.95)rectangle++(0.04,0.04)
  (2.416,9.959)rectangle++(0.04,0.04) (2.449,9.965)rectangle++(0.04,0.04) (2.181,4.44)rectangle++(0.04,0.04) (-0.508,4.233)rectangle++(0.04,0.04) (-0.265,5.215)rectangle++(0.04,0.04)
  (-0.019,6.039)rectangle++(0.04,0.04) (0.224,6.73)rectangle++(0.04,0.04) (0.457,7.308)rectangle++(0.04,0.04) (0.679,7.79)rectangle++(0.04,0.04) (0.886,8.19)rectangle++(0.04,0.04)
  (1.079,8.522)rectangle++(0.04,0.04) (1.256,8.797)rectangle++(0.04,0.04) (1.417,9.024)rectangle++(0.04,0.04) (1.563,9.21)rectangle++(0.04,0.04) (1.695,9.362)rectangle++(0.04,0.04)
  (1.813,9.486)rectangle++(0.04,0.04) (1.918,9.587)rectangle++(0.04,0.04) (2.012,9.668)rectangle++(0.04,0.04) (2.094,9.734)rectangle++(0.04,0.04) (2.167,9.786)rectangle++(0.04,0.04)
  (2.232,9.828)rectangle++(0.04,0.04) (2.288,9.86)rectangle++(0.04,0.04) (2.174,4.383)rectangle++(0.04,0.04) (2.021,5.235)rectangle++(0.04,0.04) (1.925,5.965)rectangle++(0.04,0.04)
  (1.252,3.549)rectangle++(0.04,0.04) (0.732,2.931)rectangle++(0.04,0.04) (0.738,4.058)rectangle++(0.04,0.04) (0.787,5.016)rectangle++(0.04,0.04) (0.867,5.828)rectangle++(0.04,0.04)
  (0.968,6.516)rectangle++(0.04,0.04) (1.081,7.096)rectangle++(0.04,0.04) (1.201,7.584)rectangle++(0.04,0.04) (1.322,7.995)rectangle++(0.04,0.04) (1.441,8.339)rectangle++(0.04,0.04)
  (1.556,8.627)rectangle++(0.04,0.04) (1.666,8.867)rectangle++(0.04,0.04) (1.768,9.066)rectangle++(0.04,0.04) (1.864,9.232)rectangle++(0.04,0.04) (0,2.008)rectangle++(0.04,0.04)
  (0.078,3.303)rectangle++(0.04,0.04) (0.196,4.401)rectangle++(0.04,0.04) (0.341,5.329)rectangle++(0.04,0.04) (0.501,6.112)rectangle++(0.04,0.04) (0.668,6.772)rectangle++(0.04,0.04)
  (0.836,7.325)rectangle++(0.04,0.04) (1.002,7.789)rectangle++(0.04,0.04) (1.161,8.177)rectangle++(0.04,0.04) (1.312,8.5)rectangle++(0.04,0.04) (1.453,8.769)rectangle++(0.04,0.04)
  (1.583,8.992)rectangle++(0.04,0.04) (1.703,9.176)rectangle++(0.04,0.04) (1.813,9.328)rectangle++(0.04,0.04) (-1.611,5.129)rectangle++(0.04,0.04) (-1.167,6.02)rectangle++(0.04,0.04)
  (-0.753,6.76)rectangle++(0.04,0.04) (-0.372,7.372)rectangle++(0.04,0.04) (-0.023,7.878)rectangle++(0.04,0.04) (0.293,8.293)rectangle++(0.04,0.04) (0.579,8.634)rectangle++(0.04,0.04)
  (0.835,8.912)rectangle++(0.04,0.04) (0.712,2.429)rectangle++(0.04,0.04) (0.7,3.632)rectangle++(0.04,0.04) (-0.602,3.796)rectangle++(0.04,0.04) (-0.362,4.847)rectangle++(0.04,0.04)
  (-0.116,5.73)rectangle++(0.04,0.04) (0.128,6.472)rectangle++(0.04,0.04) (-1.307,4.249)rectangle++(0.04,0.04) (-0.943,5.26)rectangle++(0.04,0.04) (0.022,3.533)rectangle++(0.04,0.04)
  (-0.702,3.649)rectangle++(0.04,0.04) (-0.453,4.726)rectangle++(0.04,0.04) (-0.198,5.632)rectangle++(0.04,0.04) (0.054,6.391)rectangle++(0.04,0.04) (0.3,7.026)rectangle++(0.04,0.04)
  (0.534,7.557)rectangle++(0.04,0.04) (0.754,7.998)rectangle++(0.04,0.04) (0.958,8.364)rectangle++(0.04,0.04) (1.147,8.668)rectangle++(0.04,0.04) (-0.152,3.922)rectangle++(0.04,0.04)
  (0.026,4.936)rectangle++(0.04,0.04) (0.217,5.79)rectangle++(0.04,0.04) (0.414,6.509)rectangle++(0.04,0.04) (0.61,7.113)rectangle++(0.04,0.04) (0.801,7.618)rectangle++(0.04,0.04)
  (0.983,8.039)rectangle++(0.04,0.04) (1.155,8.39)rectangle++(0.04,0.04) (1.315,8.682)rectangle++(0.04,0.04) (1.996,3.935)rectangle++(0.04,0.04) (1.852,4.861)rectangle++(0.04,0.04)
  (1.766,5.654)rectangle++(0.04,0.04) (1.725,6.331)rectangle++(0.04,0.04) (1.717,6.909)rectangle++(0.04,0.04) (1.734,7.4)rectangle++(0.04,0.04) (1.768,7.817)rectangle++(0.04,0.04)
  (1.813,8.17)rectangle++(0.04,0.04) (1.866,8.468)rectangle++(0.04,0.04) (1.922,8.719)rectangle++(0.04,0.04) (1.981,8.931)rectangle++(0.04,0.04) (2.038,9.108)rectangle++(0.04,0.04)
  (2.095,9.257)rectangle++(0.04,0.04) (2.047,4.223)rectangle++(0.04,0.04) (1.907,5.104)rectangle++(0.04,0.04) (1.041,3.379)rectangle++(0.04,0.04) (1.017,4.427)rectangle++(0.04,0.04)
  (1.04,5.318)rectangle++(0.04,0.04) (1.094,6.075)rectangle++(0.04,0.04) (0.36,2.726)rectangle++(0.04,0.04) (0.413,3.899)rectangle++(0.04,0.04) (0.505,4.894)rectangle++(0.04,0.04)
  (0.623,5.736)rectangle++(0.04,0.04) (0.756,6.447)rectangle++(0.04,0.04) (1.508,3.377)rectangle++(0.04,0.04) (1.415,4.406)rectangle++(0.04,0.04) (1.376,5.285)rectangle++(0.04,0.04)
  (1.379,6.033)rectangle++(0.04,0.04) (1.333,3.439)rectangle++(0.04,0.04) (0.696,2.928)rectangle++(0.04,0.04) (0.706,4.057)rectangle++(0.04,0.04) (0.76,5.016)rectangle++(0.04,0.04)
  (1.153,3.102)rectangle++(0.04,0.04) (1.102,4.187)rectangle++(0.04,0.04) (1.102,5.111)rectangle++(0.04,0.04) (1.139,5.896)rectangle++(0.04,0.04) (1.202,6.563)rectangle++(0.04,0.04)
  (1.282,7.126)rectangle++(0.04,0.04) (1.372,7.602)rectangle++(0.04,0.04) (1.468,8.003)rectangle++(0.04,0.04) (1.81,3.839)rectangle++(0.04,0.04) (1.69,4.787)rectangle++(0.04,0.04)
  (1.626,5.597)rectangle++(0.04,0.04) (1.603,6.289)rectangle++(0.04,0.04) (1.612,6.878)rectangle++(0.04,0.04) (1.643,7.378)rectangle++(0.04,0.04) (1.69,7.802)rectangle++(0.04,0.04)
  (1.746,8.16)rectangle++(0.04,0.04) (0.469,2.966)rectangle++(0.04,0.04) (0.515,4.098)rectangle++(0.04,0.04) (0.955,2.869)rectangle++(0.04,0.04) (0.924,3.997)rectangle++(0.04,0.04)
  (0.943,4.957)rectangle++(0.04,0.04) (0.998,5.772)rectangle++(0.04,0.04) (1.077,6.462)rectangle++(0.04,0.04) (1.171,7.046)rectangle++(0.04,0.04) (1.275,7.538)rectangle++(0.04,0.04)
  (1.383,7.953)rectangle++(0.04,0.04) (1.492,8.301)rectangle++(0.04,0.04) (1.598,8.592)rectangle++(0.04,0.04) (1.7,8.836)rectangle++(0.04,0.04) (1.989,4.052)rectangle++(0.04,0.04)
  (1.85,4.961)rectangle++(0.04,0.04) (1.769,5.739)rectangle++(0.04,0.04) (1.731,6.403)rectangle++(0.04,0.04) (1.725,6.97)rectangle++(0.04,0.04) (-1.126,4.647)rectangle++(0.04,0.04)
  (-0.774,5.591)rectangle++(0.04,0.04) (-0.436,6.38)rectangle++(0.04,0.04) (-0.118,7.036)rectangle++(0.04,0.04) (0.179,7.582)rectangle++(0.04,0.04) (0.453,8.034)rectangle++(0.04,0.04)
  (0.704,8.407)rectangle++(0.04,0.04) (0.933,8.714)rectangle++(0.04,0.04) (1.139,8.966)rectangle++(0.04,0.04) (1.325,9.171)rectangle++(0.04,0.04) (1.491,9.339)rectangle++(0.04,0.04)
  (1.638,9.475)rectangle++(0.04,0.04) (1.769,9.584)rectangle++(0.04,0.04) (1.885,9.672)rectangle++(0.04,0.04) (1.987,9.742)rectangle++(0.04,0.04) (2.077,9.797)rectangle++(0.04,0.04)
  (2.155,9.841)rectangle++(0.04,0.04) (2.223,9.875)rectangle++(0.04,0.04) (2.282,9.901)rectangle++(0.04,0.04) (2.334,9.921)rectangle++(0.04,0.04) (2.378,9.935)rectangle++(0.04,0.04)
  (2.417,9.946)rectangle++(0.04,0.04) (2.45,9.954)rectangle++(0.04,0.04) (2.478,9.959)rectangle++(0.04,0.04) (2.503,9.962)rectangle++(0.04,0.04) (2.524,9.964)rectangle++(0.04,0.04)
  (2.541,9.965)rectangle++(0.04,0.04) (2.557,9.964)rectangle++(0.04,0.04) (2.57,9.964)rectangle++(0.04,0.04) (2.58,9.963)rectangle++(0.04,0.04) (2.59,9.961)rectangle++(0.04,0.04)
  (2.597,9.96)rectangle++(0.04,0.04) (2.604,9.958)rectangle++(0.04,0.04) (-1.603,5.402)rectangle++(0.04,0.04) (-1.149,6.252)rectangle++(0.04,0.04) (-0.729,6.956)rectangle++(0.04,0.04)
  (-0.343,7.538)rectangle++(0.04,0.04) (0.008,8.017)rectangle++(0.04,0.04) (0.325,8.411)rectangle++(0.04,0.04) (-0.084,3.877)rectangle++(0.04,0.04) (0.081,4.895)rectangle++(0.04,0.04)
  (0.263,5.754)rectangle++(0.04,0.04) (0.451,6.476)rectangle++(0.04,0.04) (0.641,7.083)rectangle++(0.04,0.04) (0.826,7.591)rectangle++(0.04,0.04) (0.029,2.615)rectangle++(0.04,0.04)
  (0.127,3.818)rectangle++(0.04,0.04) (0.258,4.836)rectangle++(0.04,0.04) (0.411,5.697)rectangle++(0.04,0.04) (-1.092,4.149)rectangle++(0.04,0.04) (-0.764,5.167)rectangle++(0.04,0.04)
  (-0.445,6.019)rectangle++(0.04,0.04) (-0.14,6.73)rectangle++(0.04,0.04) (0.148,7.322)rectangle++(0.04,0.04) (0.417,7.814)rectangle++(0.04,0.04) (0.664,8.221)rectangle++(0.04,0.04)
  (-1.582,4.695)rectangle++(0.04,0.04) (-1.159,5.65)rectangle++(0.04,0.04) (-0.762,6.445)rectangle++(0.04,0.04) (1.685,3.032)rectangle++(0.04,0.04) (1.552,4.106)rectangle++(0.04,0.04)
  (1.481,5.024)rectangle++(0.04,0.04) (1.458,5.808)rectangle++(0.04,0.04) (1.469,6.474)rectangle++(0.04,0.04) (1.505,7.041)rectangle++(0.04,0.04) (1.559,7.521)rectangle++(0.04,0.04)
  (1.624,7.926)rectangle++(0.04,0.04) (-0.621,3.092)rectangle++(0.04,0.04) (-0.406,4.249)rectangle++(0.04,0.04) (-0.177,5.224)rectangle++(0.04,0.04) (0.056,6.044)rectangle++(0.04,0.04)
  (0.287,6.731)rectangle++(0.04,0.04) (0.511,7.306)rectangle++(0.04,0.04) (-1.415,4.486)rectangle++(0.04,0.04) (-1.026,5.466)rectangle++(0.04,0.04) (-0.656,6.283)rectangle++(0.04,0.04)
  (-1.407,4.065)rectangle++(0.04,0.04) (-1.036,5.108)rectangle++(0.04,0.04) (-0.678,5.979)rectangle++(0.04,0.04) (-0.34,6.706)rectangle++(0.04,0.04) (-0.023,7.31)rectangle++(0.04,0.04)
  (0.271,7.81)rectangle++(0.04,0.04) (0.54,8.224)rectangle++(0.04,0.04) (0.786,8.565)rectangle++(0.04,0.04) (2.029,3.792)rectangle++(0.04,0.04) (1.874,4.738)rectangle++(0.04,0.04)
  (1.78,5.549)rectangle++(0.04,0.04) (1.733,6.242)rectangle++(0.04,0.04) (1.72,6.832)rectangle++(0.04,0.04) (1.734,7.335)rectangle++(0.04,0.04) (-1.202,4.721)rectangle++(0.04,0.04)
  (-0.835,5.657)rectangle++(0.04,0.04) (0.191,2.649)rectangle++(0.04,0.04) (0.266,3.84)rectangle++(0.04,0.04) (0.378,4.85)rectangle++(0.04,0.04) (0.513,5.704)rectangle++(0.04,0.04)
  (0.662,6.424)rectangle++(0.04,0.04) (0.817,7.03)rectangle++(0.04,0.04) (0.974,7.539)rectangle++(0.04,0.04) (1.127,7.965)rectangle++(0.04,0.04) (1.274,8.322)rectangle++(0.04,0.04)
  (1.414,8.619)rectangle++(0.04,0.04) (1.544,8.865)rectangle++(0.04,0.04) (1.665,9.07)rectangle++(0.04,0.04) (1.776,9.239)rectangle++(0.04,0.04) (1.877,9.378)rectangle++(0.04,0.04)
  (2.103,4.197)rectangle++(0.04,0.04) (1.953,5.079)rectangle++(0.04,0.04) (1.861,5.835)rectangle++(0.04,0.04) (1.813,6.482)rectangle++(0.04,0.04) (-1.007,4.567)rectangle++(0.04,0.04)
  (-0.676,5.519)rectangle++(0.04,0.04) (-0.356,6.314)rectangle++(0.04,0.04) (-0.052,6.977)rectangle++(0.04,0.04) (-1.447,4.315)rectangle++(0.04,0.04) (-1.059,5.322)rectangle++(0.04,0.04)
  (-0.69,6.162)rectangle++(0.04,0.04) (-0.342,6.862)rectangle++(0.04,0.04) (-0.018,7.442)rectangle++(0.04,0.04) (0.28,7.923)rectangle++(0.04,0.04) (-1.04,3.238)rectangle++(0.04,0.04)
  (-0.757,4.39)rectangle++(0.04,0.04) (-0.47,5.358)rectangle++(0.04,0.04) (1.382,2.889)rectangle++(0.04,0.04) (1.288,3.996)rectangle++(0.04,0.04) (0.839,3.025)rectangle++(0.04,0.04)
  (0.832,4.134)rectangle++(0.04,0.04) (0.926,2.948)rectangle++(0.04,0.04) (0.903,4.065)rectangle++(0.04,0.04) (0.928,5.015)rectangle++(0.04,0.04) (0.987,5.822)rectangle++(0.04,0.04)
  (1.07,6.506)rectangle++(0.04,0.04) (-1.145,4.433)rectangle++(0.04,0.04) (-0.799,5.41)rectangle++(0.04,0.04) (-0.465,6.226)rectangle++(0.04,0.04) (-0.148,6.907)rectangle++(0.04,0.04)
  (0.148,7.473)rectangle++(0.04,0.04) (0.423,7.943)rectangle++(0.04,0.04) (0.675,8.33)rectangle++(0.04,0.04) (0.905,8.65)rectangle++(0.04,0.04) (1.113,8.913)rectangle++(0.04,0.04)
  (1.3,9.127)rectangle++(0.04,0.04) (1.468,9.302)rectangle++(0.04,0.04) (2.132,4.089)rectangle++(0.04,0.04) (1.973,4.987)rectangle++(0.04,0.04) (1.875,5.756)rectangle++(0.04,0.04)
  (-0.842,4.435)rectangle++(0.04,0.04) (-0.541,5.4)rectangle++(0.04,0.04) (-0.246,6.208)rectangle++(0.04,0.04) (0.037,6.883)rectangle++(0.04,0.04) (0.305,7.445)rectangle++(0.04,0.04)
  (0.555,7.912)rectangle++(0.04,0.04) (-1.536,4.614)rectangle++(0.04,0.04) (-1.123,5.579)rectangle++(0.04,0.04) (-0.734,6.383)rectangle++(0.04,0.04) (-0.371,7.051)rectangle++(0.04,0.04)
  (-0.035,7.605)rectangle++(0.04,0.04) (1.887,3.421)rectangle++(0.04,0.04) (1.739,4.428)rectangle++(0.04,0.04) (1.653,5.291)rectangle++(0.04,0.04) (-0.813,3.366)rectangle++(0.04,0.04)
  (-0.558,4.49)rectangle++(0.04,0.04) (-0.297,5.435)rectangle++(0.04,0.04) (-1.163,3.965)rectangle++(0.04,0.04) (-0.832,5.013)rectangle++(0.04,0.04) (-0.509,5.89)rectangle++(0.04,0.04)
  (-0.199,6.623)rectangle++(0.04,0.04) (0.093,7.234)rectangle++(0.04,0.04) (0.366,7.741)rectangle++(0.04,0.04) (0.619,8.162)rectangle++(0.04,0.04) (0.85,8.509)rectangle++(0.04,0.04)
  (-1.61,4.787)rectangle++(0.04,0.04) (-1.179,5.729)rectangle++(0.04,0.04) (-0.775,6.513)rectangle++(0.04,0.04) (-0.401,7.164)rectangle++(0.04,0.04) (-0.056,7.701)rectangle++(0.04,0.04)
  (-1.601,4.457)rectangle++(0.04,0.04) (-1.185,5.449)rectangle++(0.04,0.04) (-0.791,6.275)rectangle++(0.04,0.04) (-0.424,6.962)rectangle++(0.04,0.04) (-0.084,7.531)rectangle++(0.04,0.04)
  (0.228,8.001)rectangle++(0.04,0.04) (0.511,8.388)rectangle++(0.04,0.04) (0.768,8.705)rectangle++(0.04,0.04) (0.999,8.965)rectangle++(0.04,0.04) (1.205,9.176)rectangle++(0.04,0.04)
  (2.131,4.005)rectangle++(0.04,0.04) (1.969,4.915)rectangle++(0.04,0.04) (1.868,5.696)rectangle++(0.04,0.04) (1.814,6.363)rectangle++(0.04,0.04) (1.794,6.932)rectangle++(0.04,0.04)
  (1.8,7.417)rectangle++(0.04,0.04) (1.824,7.828)rectangle++(0.04,0.04) (1.862,8.177)rectangle++(0.04,0.04) (1.907,8.472)rectangle++(0.04,0.04) (1.958,8.721)rectangle++(0.04,0.04)
  (2.011,8.931)rectangle++(0.04,0.04) (-1.491,5.088)rectangle++(0.04,0.04) (-1.066,5.98)rectangle++(0.04,0.04) (-0.669,6.722)rectangle++(0.04,0.04) (-0.302,7.337)rectangle++(0.04,0.04)
  (-1.568,4.339)rectangle++(0.04,0.04) (-1.161,5.347)rectangle++(0.04,0.04) (-0.775,6.188)rectangle++(0.04,0.04) (-0.414,6.887)rectangle++(0.04,0.04) (-0.078,7.467)rectangle++(0.04,0.04)
  (0.23,7.946)rectangle++(0.04,0.04) (0.511,8.341)rectangle++(0.04,0.04) (0.766,8.666)rectangle++(0.04,0.04) (0.995,8.931)rectangle++(0.04,0.04) (1.201,9.148)rectangle++(0.04,0.04)
  (1.385,9.324)rectangle++(0.04,0.04) (1.548,9.466)rectangle++(0.04,0.04) (1.692,9.581)rectangle++(0.04,0.04) (-1.686,5.156)rectangle++(0.04,0.04) (-1.229,6.046)rectangle++(0.04,0.04)
  (-0.805,6.785)rectangle++(0.04,0.04) (-0.415,7.396)rectangle++(0.04,0.04) (-0.059,7.899)rectangle++(0.04,0.04) (0.263,8.313)rectangle++(0.04,0.04) (0.554,8.651)rectangle++(0.04,0.04)
  (0.815,8.928)rectangle++(0.04,0.04) (1.047,9.152)rectangle++(0.04,0.04) (1.254,9.334)rectangle++(0.04,0.04) (1.437,9.48)rectangle++(0.04,0.04) (1.599,9.596)rectangle++(0.04,0.04)
  (1.741,9.689)rectangle++(0.04,0.04) (1.865,9.762)rectangle++(0.04,0.04) (1.973,9.82)rectangle++(0.04,0.04) (2.068,9.864)rectangle++(0.04,0.04) (2.15,9.898)rectangle++(0.04,0.04)
  (2.221,9.923)rectangle++(0.04,0.04) (2.283,9.942)rectangle++(0.04,0.04) (2.336,9.956)rectangle++(0.04,0.04) (2.382,9.965)rectangle++(0.04,0.04) (2.421,9.971)rectangle++(0.04,0.04)
  (2.454,9.975)rectangle++(0.04,0.04) (2.183,4.443)rectangle++(0.04,0.04) (0.845,3.314)rectangle++(0.04,0.04) (0.848,4.38)rectangle++(0.04,0.04) (0.894,5.285)rectangle++(0.04,0.04)
  (1.204,3.184)rectangle++(0.04,0.04) (1.149,4.255)rectangle++(0.04,0.04) (1.144,5.167)rectangle++(0.04,0.04) (1.177,5.942)rectangle++(0.04,0.04) (-1.007,4.341)rectangle++(0.04,0.04)
  (-0.684,5.327)rectangle++(0.04,0.04) (-0.371,6.151)rectangle++(0.04,0.04) (-0.071,6.84)rectangle++(0.04,0.04) (0.211,7.413)rectangle++(0.04,0.04) (1.033,2.16)rectangle++(0.04,0.04)
  (0.962,3.391)rectangle++(0.04,0.04) (0.951,4.44)rectangle++(0.04,0.04) (0.984,5.332)rectangle++(0.04,0.04) (1.047,6.089)rectangle++(0.04,0.04) (1.385,3.374)rectangle++(0.04,0.04)
  (1.31,4.409)rectangle++(0.04,0.04) (1.288,5.291)rectangle++(0.04,0.04) (1.304,6.042)rectangle++(0.04,0.04) (1.348,6.68)rectangle++(0.04,0.04) (1.411,7.22)rectangle++(0.04,0.04)
  (1.623,3.675)rectangle++(0.04,0.04) (0.004,2.082)rectangle++(0.04,0.04) (0.084,3.366)rectangle++(0.04,0.04) (0.204,4.454)rectangle++(0.04,0.04) (0.349,5.374)rectangle++(0.04,0.04)
  (0.51,6.15)rectangle++(0.04,0.04) (0.677,6.803)rectangle++(0.04,0.04) (0.845,7.352)rectangle++(0.04,0.04) (1.011,7.812)rectangle++(0.04,0.04) (1.169,8.196)rectangle++(0.04,0.04)
  (1.319,8.516)rectangle++(0.04,0.04) (1.46,8.782)rectangle++(0.04,0.04) (1.59,9.002)rectangle++(0.04,0.04) (1.709,9.185)rectangle++(0.04,0.04) (1.818,9.335)rectangle++(0.04,0.04)
  (1.917,9.458)rectangle++(0.04,0.04) (2.005,9.559)rectangle++(0.04,0.04) (2.085,9.641)rectangle++(0.04,0.04) (2.155,9.708)rectangle++(0.04,0.04) (-1.63,5.268)rectangle++(0.04,0.04)
  (-1.177,6.139)rectangle++(0.04,0.04) (-0.757,6.862)rectangle++(0.04,0.04) (-0.371,7.459)rectangle++(0.04,0.04) (-0.02,7.951)rectangle++(0.04,0.04) (1.971,3.491)rectangle++(0.04,0.04)
  (1.813,4.485)rectangle++(0.04,0.04) (1.718,5.336)rectangle++(0.04,0.04) (1.672,6.063)rectangle++(0.04,0.04) (1.661,6.683)rectangle++(0.04,0.04) (1.677,7.21)rectangle++(0.04,0.04)
  (1.712,7.658)rectangle++(0.04,0.04) (1.759,8.037)rectangle++(0.04,0.04) (1.815,8.357)rectangle++(0.04,0.04) (1.874,8.627)rectangle++(0.04,0.04) (1.936,8.854)rectangle++(0.04,0.04)
  (-1.489,5.059)rectangle++(0.04,0.04) (-1.065,5.956)rectangle++(0.04,0.04) (-0.669,6.701)rectangle++(0.04,0.04) (-0.914,3.166)rectangle++(0.04,0.04) (-0.652,4.324)rectangle++(0.04,0.04)
  (-0.384,5.298)rectangle++(0.04,0.04) (-0.117,6.115)rectangle++(0.04,0.04) (0.143,6.798)rectangle++(0.04,0.04) (0.392,7.369)rectangle++(0.04,0.04) (0.625,7.844)rectangle++(0.04,0.04)
  (0.843,8.239)rectangle++(0.04,0.04) (1.044,8.565)rectangle++(0.04,0.04) (1.228,8.835)rectangle++(0.04,0.04) (1.395,9.057)rectangle++(0.04,0.04) (1.546,9.239)rectangle++(0.04,0.04)
  (2.107,4.095)rectangle++(0.04,0.04) (1.952,4.992)rectangle++(0.04,0.04) (1.857,5.762)rectangle++(0.04,0.04) (1.807,6.419)rectangle++(0.04,0.04) (1.79,6.98)rectangle++(0.04,0.04)
  (1.799,7.458)rectangle++(0.04,0.04) (1.825,7.863)rectangle++(0.04,0.04) (1.864,8.207)rectangle++(0.04,0.04) (1.91,8.498)rectangle++(0.04,0.04) (1.961,8.743)rectangle++(0.04,0.04)
  (2.015,8.949)rectangle++(0.04,0.04) (1.98,4.145)rectangle++(0.04,0.04) (1.847,5.04)rectangle++(0.04,0.04) (1.769,5.807)rectangle++(0.04,0.04) (1.734,6.461)rectangle++(0.04,0.04)
  (1.73,7.019)rectangle++(0.04,0.04) (1.749,7.493)rectangle++(0.04,0.04) (1.784,7.895)rectangle++(0.04,0.04) (1.83,8.236)rectangle++(0.04,0.04) (0.478,2.975)rectangle++(0.04,0.04)
  (0.523,4.106)rectangle++(0.04,0.04) (0.607,5.065)rectangle++(0.04,0.04) (1.183,3.077)rectangle++(0.04,0.04) (1.127,4.164)rectangle++(0.04,0.04) (0.899,3.021)rectangle++(0.04,0.04)
  (0.883,4.128)rectangle++(0.04,0.04) (0.914,5.07)rectangle++(0.04,0.04) (0.977,5.869)rectangle++(0.04,0.04) (1.339,3.316)rectangle++(0.04,0.04) (1.269,4.361)rectangle++(0.04,0.04)
  (1.25,5.252)rectangle++(0.04,0.04) (1.271,6.011)rectangle++(0.04,0.04) (1.318,6.654)rectangle++(0.04,0.04) (1.385,7.2)rectangle++(0.04,0.04) (1.463,7.661)rectangle++(0.04,0.04)
  (1.548,8.049)rectangle++(0.04,0.04) (1.635,8.376)rectangle++(0.04,0.04) (1.723,8.65)rectangle++(0.04,0.04) (1.808,8.88)rectangle++(0.04,0.04) (1.89,9.072)rectangle++(0.04,0.04)
  (1.967,9.232)rectangle++(0.04,0.04) (2.039,9.365)rectangle++(0.04,0.04) (2.106,9.475)rectangle++(0.04,0.04) (2.167,9.565)rectangle++(0.04,0.04) (2.222,9.64)rectangle++(0.04,0.04)
  (2.127,4.325)rectangle++(0.04,0.04) (1.979,5.188)rectangle++(0.04,0.04) (1.887,5.927)rectangle++(0.04,0.04) (1.839,6.558)rectangle++(0.04,0.04) (1.823,7.097)rectangle++(0.04,0.04)
  (-1.136,4.691)rectangle++(0.04,0.04) (-1.156,3.661)rectangle++(0.04,0.04) (-0.838,4.754)rectangle++(0.04,0.04) (-0.524,5.671)rectangle++(0.04,0.04) (-0.221,6.437)rectangle++(0.04,0.04)
  (0.067,7.077)rectangle++(0.04,0.04) (0.338,7.609)rectangle++(0.04,0.04) (0.59,8.05)rectangle++(0.04,0.04) (0.821,8.415)rectangle++(0.04,0.04) (1.988,3.771)rectangle++(0.04,0.04)
  (1.838,4.722)rectangle++(0.04,0.04) (1.749,5.537)rectangle++(0.04,0.04) (1.706,6.232)rectangle++(0.04,0.04) (1.697,6.825)rectangle++(0.04,0.04) (1.713,7.33)rectangle++(0.04,0.04)
  (1.747,7.758)rectangle++(0.04,0.04) (1.793,8.121)rectangle++(0.04,0.04) (1.847,8.427)rectangle++(0.04,0.04) (1.905,8.685)rectangle++(0.04,0.04) (1.964,8.903)rectangle++(0.04,0.04)
  (-1.494,5.073)rectangle++(0.04,0.04) (-1.069,5.968)rectangle++(0.04,0.04) (-0.672,6.712)rectangle++(0.04,0.04) (-0.305,7.328)rectangle++(0.04,0.04) (0.032,7.837)rectangle++(0.04,0.04)
  (0.338,8.257)rectangle++(0.04,0.04) (0.616,8.601)rectangle++(0.04,0.04) (0.865,8.882)rectangle++(0.04,0.04) (1.088,9.112)rectangle++(0.04,0.04) (1.287,9.298)rectangle++(0.04,0.04)
  (1.464,9.448)rectangle++(0.04,0.04) (1.62,9.568)rectangle++(0.04,0.04) (1.758,9.664)rectangle++(0.04,0.04) (1.878,9.741)rectangle++(0.04,0.04) (1.984,9.801)rectangle++(0.04,0.04)
  (2.076,9.847)rectangle++(0.04,0.04) (2.156,9.883)rectangle++(0.04,0.04) (2.226,9.911)rectangle++(0.04,0.04) (2.286,9.931)rectangle++(0.04,0.04) (2.339,9.946)rectangle++(0.04,0.04)
  (2.383,9.957)rectangle++(0.04,0.04) (2.422,9.964)rectangle++(0.04,0.04) (2.455,9.969)rectangle++(0.04,0.04) (2.483,9.972)rectangle++(0.04,0.04) (2.508,9.973)rectangle++(0.04,0.04)
  (2.528,9.973)rectangle++(0.04,0.04) (-1.62,5.39)rectangle++(0.04,0.04) (-1.163,6.243)rectangle++(0.04,0.04) (-0.741,6.949)rectangle++(0.04,0.04) (-0.354,7.533)rectangle++(0.04,0.04)
  (-0.002,8.013)rectangle++(0.04,0.04) (0.316,8.407)rectangle++(0.04,0.04) (0.603,8.73)rectangle++(0.04,0.04) (-1.701,4.784)rectangle++(0.04,0.04) (-1.257,5.731)rectangle++(0.04,0.04)
  (-1.398,3.843)rectangle++(0.04,0.04) (-1.037,4.919)rectangle++(0.04,0.04) (-0.687,5.819)rectangle++(0.04,0.04) (-0.353,6.57)rectangle++(0.04,0.04) (0.022,2.476)rectangle++(0.04,0.04)
  (0.116,3.7)rectangle++(0.04,0.04) (0.244,4.737)rectangle++(0.04,0.04) (0.395,5.613)rectangle++(0.04,0.04) (0.558,6.351)rectangle++(0.04,0.04) (0.726,6.972)rectangle++(0.04,0.04)
  (0.894,7.494)rectangle++(0.04,0.04) (1.057,7.93)rectangle++(0.04,0.04) (1.214,8.294)rectangle++(0.04,0.04) (1.361,8.598)rectangle++(0.04,0.04) (1.499,8.85)rectangle++(0.04,0.04)
  (-1.566,4.976)rectangle++(0.04,0.04) (-1.134,5.888)rectangle++(0.04,0.04) (-0.731,6.647)rectangle++(0.04,0.04) (-0.358,7.275)rectangle++(0.04,0.04) (-0.015,7.794)rectangle++(0.04,0.04)
  (0.297,8.222)rectangle++(0.04,0.04) (0.579,8.573)rectangle++(0.04,0.04) (0.833,8.86)rectangle++(0.04,0.04) (1.06,9.094)rectangle++(0.04,0.04) (-1.696,4.942)rectangle++(0.04,0.04)
  (-1.247,5.864)rectangle++(0.04,0.04) (-0.827,6.631)rectangle++(0.04,0.04) (-0.44,7.265)rectangle++(0.04,0.04) (-0.086,7.789)rectangle++(0.04,0.04) (0.237,8.221)rectangle++(0.04,0.04)
  (0.528,8.574)rectangle++(0.04,0.04) (0.789,8.863)rectangle++(0.04,0.04) (-1.696,4.845)rectangle++(0.04,0.04) (-1.25,5.783)rectangle++(0.04,0.04) (-0.833,6.561)rectangle++(0.04,0.04)
  (-0.448,7.207)rectangle++(0.04,0.04) (-0.095,7.74)rectangle++(0.04,0.04) (0.227,8.179)rectangle++(0.04,0.04) (0.518,8.539)rectangle++(0.04,0.04) (0.779,8.834)rectangle++(0.04,0.04)
  (1.014,9.074)rectangle++(0.04,0.04) (1.222,9.268)rectangle++(0.04,0.04) (1.408,9.425)rectangle++(0.04,0.04) (1.571,9.551)rectangle++(0.04,0.04) (1.715,9.652)rectangle++(0.04,0.04)
  (1.842,9.732)rectangle++(0.04,0.04) (1.953,9.795)rectangle++(0.04,0.04) (2.049,9.844)rectangle++(0.04,0.04) (2.134,9.881)rectangle++(0.04,0.04) (2.207,9.91)rectangle++(0.04,0.04)
  (2.27,9.931)rectangle++(0.04,0.04) (2.324,9.947)rectangle++(0.04,0.04) (-0.294,3.255)rectangle++(0.04,0.04) (-0.7,3.535)rectangle++(0.04,0.04) (-0.832,3.514)rectangle++(0.04,0.04)
  (-0.569,4.617)rectangle++(0.04,0.04) (-0.301,5.543)rectangle++(0.04,0.04) (-0.037,6.32)rectangle++(0.04,0.04) (-1.304,4.188)rectangle++(0.04,0.04) (1.19,2.474)rectangle++(0.04,0.04)
  (1.108,3.651)rectangle++(0.04,0.04) (1.086,4.656)rectangle++(0.04,0.04) (1.107,5.51)rectangle++(0.04,0.04) (1.159,6.235)rectangle++(0.04,0.04) (-1.072,4.396)rectangle++(0.04,0.04)
  (-0.738,5.376)rectangle++(0.04,0.04) (-0.414,6.195)rectangle++(0.04,0.04) (1.583,3.063)rectangle++(0.04,0.04) (1.466,4.136)rectangle++(0.04,0.04) (0.853,3.093)rectangle++(0.04,0.04)
  (0.846,4.191)rectangle++(0.04,0.04) (0.885,5.124)rectangle++(0.04,0.04) (0.955,5.917)rectangle++(0.04,0.04) (1.046,6.587)rectangle++(0.04,0.04) (1.15,7.153)rectangle++(0.04,0.04)
  (-1.268,4.575)rectangle++(0.04,0.04) (-0.897,5.536)rectangle++(0.04,0.04) (-0.544,6.338)rectangle++(0.04,0.04) (-0.211,7.005)rectangle++(0.04,0.04) (0.099,7.559)rectangle++(0.04,0.04)
  (0.384,8.017)rectangle++(0.04,0.04) (0.645,8.396)rectangle++(0.04,0.04) (-1.623,4.726)rectangle++(0.04,0.04) (-1.192,5.678)rectangle++(0.04,0.04) (-0.789,6.47)rectangle++(0.04,0.04)
  (-0.414,7.128)rectangle++(0.04,0.04) (-0.069,7.671)rectangle++(0.04,0.04) (0.246,8.119)rectangle++(0.04,0.04) (0.532,8.488)rectangle++(0.04,0.04) (0.789,8.79)rectangle++(0.04,0.04)
  (1.02,9.036)rectangle++(0.04,0.04) (1.227,9.236)rectangle++(0.04,0.04) (1.41,9.398)rectangle++(0.04,0.04) (1.572,9.528)rectangle++(0.04,0.04) (1.715,9.632)rectangle++(0.04,0.04)
  (1.841,9.715)rectangle++(0.04,0.04) (1.951,9.78)rectangle++(0.04,0.04) (2.048,9.831)rectangle++(0.04,0.04) (2.131,9.871)rectangle++(0.04,0.04) (2.204,9.901)rectangle++(0.04,0.04)
  (2.268,9.924)rectangle++(0.04,0.04) (2.322,9.941)rectangle++(0.04,0.04) (2.369,9.953)rectangle++(0.04,0.04) (2.41,9.962)rectangle++(0.04,0.04) (2.445,9.967)rectangle++(0.04,0.04)
  (2.474,9.97)rectangle++(0.04,0.04) (2.5,9.972)rectangle++(0.04,0.04) (2.522,9.973)rectangle++(0.04,0.04) (2.54,9.972)rectangle++(0.04,0.04) (2.556,9.971)rectangle++(0.04,0.04)
  (-1.614,5.395)rectangle++(0.04,0.04) (-1.391,3.71)rectangle++(0.04,0.04) (-1.036,4.805)rectangle++(0.04,0.04) (-0.691,5.722)rectangle++(0.04,0.04) (-0.36,6.488)rectangle++(0.04,0.04)
  (-0.049,7.125)rectangle++(0.04,0.04) (-1.477,4.345)rectangle++(0.04,0.04) (-1.084,5.349)rectangle++(0.04,0.04) (-0.71,6.186)rectangle++(0.04,0.04) (-0.358,6.883)rectangle++(0.04,0.04)
  (-0.031,7.461)rectangle++(0.04,0.04) (0.27,7.939)rectangle++(0.04,0.04) (0.545,8.334)rectangle++(0.04,0.04) (0.794,8.658)rectangle++(0.04,0.04) (1.019,8.924)rectangle++(0.04,0.04)
  (1.221,9.141)rectangle++(0.04,0.04) (1.401,9.317)rectangle++(0.04,0.04) (1.562,9.459)rectangle++(0.04,0.04) (1.703,9.574)rectangle++(0.04,0.04) (1.829,9.666)rectangle++(0.04,0.04)
  (-1.68,5.199)rectangle++(0.04,0.04) (-1.222,6.082)rectangle++(0.04,0.04) (-0.798,6.815)rectangle++(0.04,0.04) (-0.408,7.421)rectangle++(0.04,0.04) (-1.604,4.336)rectangle++(0.04,0.04)
  (-1.193,5.346)rectangle++(0.04,0.04) (-0.802,6.188)rectangle++(0.04,0.04) (-0.436,6.888)rectangle++(0.04,0.04) (-0.098,7.469)rectangle++(0.04,0.04) (0.214,7.948)rectangle++(0.04,0.04)
  (0.497,8.344)rectangle++(0.04,0.04) (0.754,8.669)rectangle++(0.04,0.04) (0.986,8.934)rectangle++(0.04,0.04) (0.037,2.798)rectangle++(0.04,0.04) (0.141,3.973)rectangle++(0.04,0.04)
  (0.277,4.967)rectangle++(0.04,0.04) (0.432,5.807)rectangle++(0.04,0.04) (0.597,6.515)rectangle++(0.04,0.04) (0.766,7.11)rectangle++(0.04,0.04) (0.933,7.609)rectangle++(0.04,0.04)
  (1.096,8.027)rectangle++(0.04,0.04) (-1.464,4.738)rectangle++(0.04,0.04) (1.345,2.544)rectangle++(0.04,0.04) (-0.255,3.702)rectangle++(0.04,0.04) (-0.071,4.753)rectangle++(0.04,0.04)
  (0.127,5.639)rectangle++(0.04,0.04) (1.382,3.079)rectangle++(0.04,0.04) (1.295,4.158)rectangle++(0.04,0.04) (1.265,5.078)rectangle++(0.04,0.04) (1.276,5.862)rectangle++(0.04,0.04)
  (1.317,6.528)rectangle++(0.04,0.04) (1.378,7.092)rectangle++(0.04,0.04) (1.595,3.643)rectangle++(0.04,0.04) (1.499,4.629)rectangle++(0.04,0.04) (-0.67,4.142)rectangle++(0.04,0.04)
  (-0.406,5.144)rectangle++(0.04,0.04) (-0.142,5.985)rectangle++(0.04,0.04) (0.117,6.689)rectangle++(0.04,0.04) (-1.355,4.29)rectangle++(0.04,0.04) (-0.982,5.297)rectangle++(0.04,0.04)
  (1.427,2.761)rectangle++(0.04,0.04) (1.321,3.886)rectangle++(0.04,0.04) (1.276,4.846)rectangle++(0.04,0.04) (1.276,5.664)rectangle++(0.04,0.04) (1.309,6.36)rectangle++(0.04,0.04)
  (1.365,6.95)rectangle++(0.04,0.04) (1.436,7.449)rectangle++(0.04,0.04) (1.516,7.87)rectangle++(0.04,0.04) (1.602,8.225)rectangle++(0.04,0.04) (1.688,8.524)rectangle++(0.04,0.04)
  (1.774,8.774)rectangle++(0.04,0.04) (1.965,4.057)rectangle++(0.04,0.04) (1.83,4.966)rectangle++(0.04,0.04) (1.752,5.744)rectangle++(0.04,0.04) (1.717,6.408)rectangle++(0.04,0.04)
  (1.386,3.588)rectangle++(0.04,0.04) (-0.47,3.916)rectangle++(0.04,0.04) (-0.245,4.943)rectangle++(0.04,0.04) (-0.012,5.808)rectangle++(0.04,0.04) (0.219,6.533)rectangle++(0.04,0.04)
  (0.446,7.141)rectangle++(0.04,0.04) (0.662,7.648)rectangle++(0.04,0.04) (0.867,8.071)rectangle++(0.04,0.04) (1.057,8.421)rectangle++(0.04,0.04) (1.233,8.712)rectangle++(0.04,0.04)
  (1.395,8.952)rectangle++(0.04,0.04) (1.541,9.15)rectangle++(0.04,0.04) (1.674,9.312)rectangle++(0.04,0.04) (1.793,9.444)rectangle++(0.04,0.04) (1.9,9.552)rectangle++(0.04,0.04)
  (1.995,9.64)rectangle++(0.04,0.04) (2.079,9.71)rectangle++(0.04,0.04) (2.153,9.767)rectangle++(0.04,0.04) (2.219,9.812)rectangle++(0.04,0.04) (2.276,9.847)rectangle++(0.04,0.04)
  (-0.017,1.643)rectangle++(0.04,0.04) (0.049,2.994)rectangle++(0.04,0.04) (0.159,4.139)rectangle++(0.04,0.04) (0.299,5.108)rectangle++(0.04,0.04) (0.456,5.926)rectangle++(0.04,0.04)
  (0.622,6.615)rectangle++(0.04,0.04) (0.791,7.194)rectangle++(0.04,0.04) (0.958,7.679)rectangle++(0.04,0.04) (1.12,8.085)rectangle++(0.04,0.04) (1.273,8.424)rectangle++(0.04,0.04)
  (1.417,8.706)rectangle++(0.04,0.04) (1.55,8.939)rectangle++(0.04,0.04) (1.673,9.133)rectangle++(0.04,0.04) (1.785,9.292)rectangle++(0.04,0.04) (1.887,9.423)rectangle++(0.04,0.04)
  (1.979,9.53)rectangle++(0.04,0.04) (2.061,9.618)rectangle++(0.04,0.04) (2.134,9.689)rectangle++(0.04,0.04) (-1.63,5.26)rectangle++(0.04,0.04) (-1.177,6.133)rectangle++(0.04,0.04)
  (-0.758,6.856)rectangle++(0.04,0.04) (-0.372,7.454)rectangle++(0.04,0.04) (-0.02,7.947)rectangle++(0.04,0.04) (0.298,8.352)rectangle++(0.04,0.04) (0.586,8.684)rectangle++(0.04,0.04)
  (0.843,8.954)rectangle++(0.04,0.04) (1.072,9.173)rectangle++(0.04,0.04) (-1.711,4.96)rectangle++(0.04,0.04) (-1.316,3.606)rectangle++(0.04,0.04) (-0.976,4.714)rectangle++(0.04,0.04)
  (-0.644,5.642)rectangle++(0.04,0.04) (-0.324,6.418)rectangle++(0.04,0.04) (-0.02,7.064)rectangle++(0.04,0.04) (0.263,7.602)rectangle++(0.04,0.04) (0.525,8.047)rectangle++(0.04,0.04)
  (0.766,8.415)rectangle++(0.04,0.04) (0.396,2.978)rectangle++(0.04,0.04) (0.453,4.112)rectangle++(0.04,0.04) (0.547,5.073)rectangle++(0.04,0.04) (0.666,5.886)rectangle++(0.04,0.04)
  (0.799,6.573)rectangle++(0.04,0.04) (0.94,7.151)rectangle++(0.04,0.04) (1.083,7.637)rectangle++(0.04,0.04) (1.224,8.044)rectangle++(0.04,0.04) (1.36,8.385)rectangle++(0.04,0.04)
  (1.489,8.669)rectangle++(0.04,0.04) (1.972,3.972)rectangle++(0.04,0.04) (0.001,3.305)rectangle++(0.04,0.04) (0.131,4.406)rectangle++(0.04,0.04) (0.285,5.336)rectangle++(0.04,0.04)
  (0.454,6.12)rectangle++(0.04,0.04) (1.462,3.245)rectangle++(0.04,0.04) (1.371,4.296)rectangle++(0.04,0.04) (1.335,5.193)rectangle++(0.04,0.04) (1.34,5.957)rectangle++(0.04,0.04)
  (1.375,6.606)rectangle++(0.04,0.04) (1.431,7.156)rectangle++(0.04,0.04) (-1.219,4.629)rectangle++(0.04,0.04) (-1.157,3.633)rectangle++(0.04,0.04) (-0.841,4.73)rectangle++(0.04,0.04)
  (-0.527,5.651)rectangle++(0.04,0.04) (-0.225,6.42)rectangle++(0.04,0.04) (0.064,7.063)rectangle++(0.04,0.04) (0.335,7.597)rectangle++(0.04,0.04) (0.586,8.04)rectangle++(0.04,0.04)
  (0.818,8.407)rectangle++(0.04,0.04) (1.029,8.709)rectangle++(0.04,0.04) (1.221,8.958)rectangle++(0.04,0.04) (2.075,3.967)rectangle++(0.04,0.04) (1.92,4.885)rectangle++(0.04,0.04)
  (1.826,5.671)rectangle++(0.04,0.04) (1.776,6.344)rectangle++(0.04,0.04) (1.761,6.917)rectangle++(0.04,0.04) (-1.109,4.643)rectangle++(0.04,0.04) (-0.759,5.588)rectangle++(0.04,0.04)
  (-0.424,6.376)rectangle++(0.04,0.04) (-0.107,7.033)rectangle++(0.04,0.04) (0.188,7.578)rectangle++(0.04,0.04) (0.461,8.03)rectangle++(0.04,0.04) (0.71,8.403)rectangle++(0.04,0.04)
  (1.998,3.744)rectangle++(0.04,0.04) (1.846,4.699)rectangle++(0.04,0.04) (1.755,5.516)rectangle++(0.04,0.04) (1.71,6.215)rectangle++(0.04,0.04) (1.7,6.81)rectangle++(0.04,0.04)
  (1.715,7.317)rectangle++(0.04,0.04) (1.748,7.747)rectangle++(0.04,0.04) (1.794,8.111)rectangle++(0.04,0.04) (1.847,8.419)rectangle++(0.04,0.04) (1.904,8.679)rectangle++(0.04,0.04)
  (1.964,8.897)rectangle++(0.04,0.04) (2.023,9.08)rectangle++(0.04,0.04) (2.08,9.233)rectangle++(0.04,0.04) (2.135,9.361)rectangle++(0.04,0.04) (2.187,9.468)rectangle++(0.04,0.04)
  (2.236,9.556)rectangle++(0.04,0.04) (-1.584,5.253)rectangle++(0.04,0.04) (-1.138,6.125)rectangle++(0.04,0.04) (-0.725,6.848)rectangle++(0.04,0.04) (-0.344,7.446)rectangle++(0.04,0.04)
  (1.884,3.32)rectangle++(0.04,0.04) (1.732,4.343)rectangle++(0.04,0.04) (1.644,5.218)rectangle++(0.04,0.04) (1.604,5.966)rectangle++(0.04,0.04) (1.6,6.603)rectangle++(0.04,0.04)
  (1.622,7.145)rectangle++(0.04,0.04) (1.662,7.604)rectangle++(0.04,0.04) (1.715,7.993)rectangle++(0.04,0.04) (1.775,8.322)rectangle++(0.04,0.04) (1.839,8.599)rectangle++(0.04,0.04)
  (1.905,8.832)rectangle++(0.04,0.04) (1.971,9.027)rectangle++(0.04,0.04) (2.034,9.19)rectangle++(0.04,0.04) (2.094,9.327)rectangle++(0.04,0.04) (2.151,9.44)rectangle++(0.04,0.04)
  (2.204,9.534)rectangle++(0.04,0.04) (2.252,9.612)rectangle++(0.04,0.04) (2.297,9.676)rectangle++(0.04,0.04) (2.337,9.729)rectangle++(0.04,0.04) (2.373,9.773)rectangle++(0.04,0.04)
  (2.406,9.808)rectangle++(0.04,0.04) (2.435,9.837)rectangle++(0.04,0.04) (2.461,9.86)rectangle++(0.04,0.04) (2.484,9.879)rectangle++(0.04,0.04) (2.505,9.894)rectangle++(0.04,0.04)
  (1.32,2.324)rectangle++(0.04,0.04) (1.213,3.519)rectangle++(0.04,0.04) (1.17,4.539)rectangle++(0.04,0.04) (1.174,5.408)rectangle++(0.04,0.04) (1.212,6.146)rectangle++(0.04,0.04)
  (1.273,6.772)rectangle++(0.04,0.04) (1.351,7.301)rectangle++(0.04,0.04) (1.438,7.748)rectangle++(0.04,0.04) (1.53,8.125)rectangle++(0.04,0.04) (1.624,8.441)rectangle++(0.04,0.04)
  (1.715,8.706)rectangle++(0.04,0.04) (1.804,8.928)rectangle++(0.04,0.04) (1.888,9.113)rectangle++(0.04,0.04) (1.967,9.266)rectangle++(0.04,0.04) (2.041,9.394)rectangle++(0.04,0.04)
  (2.108,9.499)rectangle++(0.04,0.04) (2.17,9.586)rectangle++(0.04,0.04) (2.12,4.303)rectangle++(0.04,0.04) (1.972,5.169)rectangle++(0.04,0.04) (1.88,5.911)rectangle++(0.04,0.04)
  (1.833,6.545)rectangle++(0.04,0.04) (1.817,7.086)rectangle++(0.04,0.04) (0.446,2.849)rectangle++(0.04,0.04) (0.491,4)rectangle++(0.04,0.04) (0.575,4.976)rectangle++(0.04,0.04)
  (1.165,3.052)rectangle++(0.04,0.04) (1.11,4.144)rectangle++(0.04,0.04) (1.107,5.074)rectangle++(0.04,0.04) (1.142,5.865)rectangle++(0.04,0.04) (-0.997,4.32)rectangle++(0.04,0.04)
  (-0.677,5.308)rectangle++(0.04,0.04) (-0.365,6.135)rectangle++(0.04,0.04) (-0.067,6.825)rectangle++(0.04,0.04) (0.214,7.4)rectangle++(0.04,0.04) (-1.489,4.449)rectangle++(0.04,0.04)
  (-1.09,5.437)rectangle++(0.04,0.04) (-0.791,3.138)rectangle++(0.04,0.04) (-0.549,4.295)rectangle++(0.04,0.04) (1.128,2.666)rectangle++(0.04,0.04) (1.063,3.817)rectangle++(0.04,0.04)
  (-0.576,3.901)rectangle++(0.04,0.04) (-0.336,4.935)rectangle++(0.04,0.04) (-0.09,5.804)rectangle++(0.04,0.04) (0.153,6.533)rectangle++(0.04,0.04) (0.389,7.143)rectangle++(0.04,0.04)
  (0.614,7.653)rectangle++(0.04,0.04) (0.826,8.076)rectangle++(0.04,0.04) (1.023,8.428)rectangle++(0.04,0.04) (1.205,8.719)rectangle++(0.04,0.04) (1.37,8.959)rectangle++(0.04,0.04)
  (1.521,9.157)rectangle++(0.04,0.04) (2.089,4.073)rectangle++(0.04,0.04) (1.937,4.975)rectangle++(0.04,0.04) (1.843,5.747)rectangle++(0.04,0.04) (-0.846,4.428)rectangle++(0.04,0.04)
  (-0.544,5.394)rectangle++(0.04,0.04) (-0.249,6.203)rectangle++(0.04,0.04) (0.034,6.878)rectangle++(0.04,0.04) (0.302,7.441)rectangle++(0.04,0.04) (0.552,7.909)rectangle++(0.04,0.04)
  (0.784,8.297)rectangle++(0.04,0.04) (-1.577,4.732)rectangle++(0.04,0.04) (-0.719,3.101)rectangle++(0.04,0.04) (-0.489,4.261)rectangle++(0.04,0.04) (-0.247,5.238)rectangle++(0.04,0.04)
  (-0.003,6.058)rectangle++(0.04,0.04) (0.238,6.746)rectangle++(0.04,0.04) (0.47,7.321)rectangle++(0.04,0.04) (0.69,7.8)rectangle++(0.04,0.04) (0.896,8.199)rectangle++(0.04,0.04)
  (1.087,8.529)rectangle++(0.04,0.04) (1.263,8.802)rectangle++(0.04,0.04) (1.424,9.028)rectangle++(0.04,0.04) (1.569,9.213)rectangle++(0.04,0.04) (2.098,4.095)rectangle++(0.04,0.04)
  (1.945,4.993)rectangle++(0.04,0.04) (1.85,5.762)rectangle++(0.04,0.04) (1.801,6.42)rectangle++(0.04,0.04) (1.786,6.981)rectangle++(0.04,0.04) (1.795,7.459)rectangle++(0.04,0.04)
  (-1.218,4.757)rectangle++(0.04,0.04) (-0.847,5.688)rectangle++(0.04,0.04) (-0.495,6.465)rectangle++(0.04,0.04) (-0.164,7.111)rectangle++(0.04,0.04) (0.143,7.647)rectangle++(0.04,0.04)
  (0.425,8.091)rectangle++(0.04,0.04) (0.683,8.456)rectangle++(0.04,0.04) (0.916,8.757)rectangle++(0.04,0.04) (1.127,9.003)rectangle++(0.04,0.04) (2.097,3.954)rectangle++(0.04,0.04)
  (1.939,4.873)rectangle++(0.04,0.04) (1.841,5.661)rectangle++(0.04,0.04) (1.186,3.472)rectangle++(0.04,0.04) (1.145,4.5)rectangle++(0.04,0.04) (-0.732,3.328)rectangle++(0.04,0.04)
  (-0.491,4.455)rectangle++(0.04,0.04) (-0.242,5.402)rectangle++(0.04,0.04) (0.008,6.198)rectangle++(0.04,0.04) (0.253,6.864)rectangle++(0.04,0.04) (0.487,7.42)rectangle++(0.04,0.04)
  (0.709,7.884)rectangle++(0.04,0.04) (0.916,8.269)rectangle++(0.04,0.04) (1.107,8.589)rectangle++(0.04,0.04) (1.282,8.852)rectangle++(0.04,0.04) (1.442,9.069)rectangle++(0.04,0.04)
  (1.586,9.247)rectangle++(0.04,0.04) (1.716,9.393)rectangle++(0.04,0.04) (1.832,9.512)rectangle++(0.04,0.04) (1.935,9.608)rectangle++(0.04,0.04) (2.027,9.685)rectangle++(0.04,0.04)
  (2.108,9.748)rectangle++(0.04,0.04) (2.18,9.797)rectangle++(0.04,0.04) (2.243,9.837)rectangle++(0.04,0.04) (2.297,9.868)rectangle++(0.04,0.04) (2.345,9.892)rectangle++(0.04,0.04)
  (2.387,9.911)rectangle++(0.04,0.04) (2.423,9.925)rectangle++(0.04,0.04) (2.454,9.935)rectangle++(0.04,0.04) (2.482,9.943)rectangle++(0.04,0.04) (2.505,9.948)rectangle++(0.04,0.04)
  (2.525,9.952)rectangle++(0.04,0.04) (2.542,9.955)rectangle++(0.04,0.04) (2.557,9.956)rectangle++(0.04,0.04) (2.569,9.956)rectangle++(0.04,0.04) (2.58,9.956)rectangle++(0.04,0.04)
  (2.589,9.956)rectangle++(0.04,0.04) (2.597,9.955)rectangle++(0.04,0.04) (2.603,9.954)rectangle++(0.04,0.04) (2.609,9.953)rectangle++(0.04,0.04) (2.613,9.952)rectangle++(0.04,0.04)
  (2.617,9.951)rectangle++(0.04,0.04) (2.158,4.475)rectangle++(0.04,0.04) (2.011,5.314)rectangle++(0.04,0.04) (1.919,6.032)rectangle++(0.04,0.04) (1.871,6.647)rectangle++(0.04,0.04)
  (1.854,7.171)rectangle++(0.04,0.04) (1.86,7.618)rectangle++(0.04,0.04) (1.884,7.997)rectangle++(0.04,0.04) (1.919,8.318)rectangle++(0.04,0.04) (1.962,8.59)rectangle++(0.04,0.04)
  (2.009,8.819)rectangle++(0.04,0.04) (2.058,9.012)rectangle++(0.04,0.04) (2.108,9.174)rectangle++(0.04,0.04) (2.156,9.31)rectangle++(0.04,0.04) (2.203,9.423)rectangle++(0.04,0.04)
  (2.247,9.518)rectangle++(0.04,0.04) (2.289,9.597)rectangle++(0.04,0.04) (2.108,4.332)rectangle++(0.04,0.04) (1.963,5.194)rectangle++(0.04,0.04) (1.874,5.933)rectangle++(0.04,0.04)
  (1.828,6.564)rectangle++(0.04,0.04) (1.814,7.103)rectangle++(0.04,0.04) (1.824,7.561)rectangle++(0.04,0.04) (1.851,7.95)rectangle++(0.04,0.04) (1.889,8.28)rectangle++(0.04,0.04)
  (1.935,8.558)rectangle++(0.04,0.04) (1.984,8.793)rectangle++(0.04,0.04) (2.036,8.991)rectangle++(0.04,0.04) (2.088,9.157)rectangle++(0.04,0.04) (2.139,9.296)rectangle++(0.04,0.04)
  (2.188,9.413)rectangle++(0.04,0.04) (2.234,9.509)rectangle++(0.04,0.04) (2.277,9.59)rectangle++(0.04,0.04) (2.317,9.656)rectangle++(0.04,0.04) (2.353,9.711)rectangle++(0.04,0.04)
  (2.387,9.757)rectangle++(0.04,0.04) (2.417,9.794)rectangle++(0.04,0.04) (2.444,9.824)rectangle++(0.04,0.04) (2.468,9.849)rectangle++(0.04,0.04) (2.49,9.869)rectangle++(0.04,0.04)
  (2.509,9.886)rectangle++(0.04,0.04) (-1.605,5.369)rectangle++(0.04,0.04) (-1.152,6.224)rectangle++(0.04,0.04) (-0.732,6.933)rectangle++(0.04,0.04) (-0.347,7.519)rectangle++(0.04,0.04)
  (1.902,3.333)rectangle++(0.04,0.04) (1.748,4.353)rectangle++(0.04,0.04) (0.824,2.349)rectangle++(0.04,0.04) (0.792,3.559)rectangle++(0.04,0.04) (-0.57,3.799)rectangle++(0.04,0.04)
  (-0.335,4.848)rectangle++(0.04,0.04) (-0.093,5.73)rectangle++(0.04,0.04) (0.148,6.471)rectangle++(0.04,0.04) (0.383,7.09)rectangle++(0.04,0.04) (0.607,7.608)rectangle++(0.04,0.04)
  (0.818,8.039)rectangle++(0.04,0.04) (1.014,8.396)rectangle++(0.04,0.04) (1.196,8.692)rectangle++(0.04,0.04) (1.362,8.937)rectangle++(0.04,0.04) (1.513,9.138)rectangle++(0.04,0.04)
  (1.649,9.303)rectangle++(0.04,0.04) (1.772,9.438)rectangle++(0.04,0.04) (1.881,9.547)rectangle++(0.04,0.04) (1.979,9.636)rectangle++(0.04,0.04) (2.065,9.708)rectangle++(0.04,0.04)
  (-1.647,5.251)rectangle++(0.04,0.04) (1.493,2.601)rectangle++(0.04,0.04) (1.371,3.747)rectangle++(0.04,0.04) (1.313,4.727)rectangle++(0.04,0.04) (1.303,5.561)rectangle++(0.04,0.04)
  (1.328,6.271)rectangle++(0.04,0.04) (1.377,6.874)rectangle++(0.04,0.04) (1.443,7.384)rectangle++(0.04,0.04) (1.52,7.815)rectangle++(0.04,0.04) (1.602,8.178)rectangle++(0.04,0.04)
  (1.687,8.483)rectangle++(0.04,0.04) (1.771,8.739)rectangle++(0.04,0.04) (1.957,4.049)rectangle++(0.04,0.04) (1.823,4.96)rectangle++(0.04,0.04) (1.746,5.739)rectangle++(0.04,0.04)
  (1.711,6.405)rectangle++(0.04,0.04) (1.709,6.972)rectangle++(0.04,0.04) (1.729,7.454)rectangle++(0.04,0.04) (1.766,7.863)rectangle++(0.04,0.04) (1.813,8.209)rectangle++(0.04,0.04)
  (1.867,8.501)rectangle++(0.04,0.04) (1.925,8.748)rectangle++(0.04,0.04) (1.984,8.955)rectangle++(0.04,0.04) (2.042,9.128)rectangle++(0.04,0.04) (2.099,9.274)rectangle++(0.04,0.04)
  (2.153,9.395)rectangle++(0.04,0.04) (2.203,9.496)rectangle++(0.04,0.04) (2.251,9.579)rectangle++(0.04,0.04) (2.294,9.649)rectangle++(0.04,0.04) (2.334,9.706)rectangle++(0.04,0.04)
  (2.37,9.753)rectangle++(0.04,0.04) (2.402,9.791)rectangle++(0.04,0.04) (2.431,9.823)rectangle++(0.04,0.04) (2.457,9.848)rectangle++(0.04,0.04) (2.48,9.869)rectangle++(0.04,0.04)
  (2.501,9.886)rectangle++(0.04,0.04) (2.519,9.899)rectangle++(0.04,0.04) (0.579,3.169)rectangle++(0.04,0.04) (0.617,4.267)rectangle++(0.04,0.04) (0.693,5.198)rectangle++(0.04,0.04)
  (0.794,5.987)rectangle++(0.04,0.04) (0.913,6.653)rectangle++(0.04,0.04) (1.04,7.215)rectangle++(0.04,0.04) (1.17,7.687)rectangle++(0.04,0.04) (1.3,8.084)rectangle++(0.04,0.04)
  (1.426,8.415)rectangle++(0.04,0.04) (1.547,8.692)rectangle++(0.04,0.04) (1.66,8.923)rectangle++(0.04,0.04) (1.766,9.114)rectangle++(0.04,0.04) (-0.614,3.429)rectangle++(0.04,0.04)
  (-0.387,4.536)rectangle++(0.04,0.04) (-0.15,5.467)rectangle++(0.04,0.04) (0.089,6.249)rectangle++(0.04,0.04) (0.324,6.904)rectangle++(0.04,0.04) (0.549,7.452)rectangle++(0.04,0.04)
  (0.763,7.908)rectangle++(0.04,0.04) (0.962,8.288)rectangle++(0.04,0.04) (1.147,8.602)rectangle++(0.04,0.04) (1.317,8.862)rectangle++(0.04,0.04) (1.472,9.076)rectangle++(0.04,0.04)
  (1.612,9.252)rectangle++(0.04,0.04) (1.738,9.396)rectangle++(0.04,0.04) (0.04,2.861)rectangle++(0.04,0.04) (0.146,4.026)rectangle++(0.04,0.04) (0.283,5.013)rectangle++(0.04,0.04)
  (0.439,5.846)rectangle++(0.04,0.04) (0.605,6.547)rectangle++(0.04,0.04) (0.333,2.77)rectangle++(0.04,0.04) (0.391,3.938)rectangle++(0.04,0.04) (0.488,4.927)rectangle++(0.04,0.04)
  (0.61,5.765)rectangle++(0.04,0.04) (0.747,6.472)rectangle++(0.04,0.04) (-1.196,4.365)rectangle++(0.04,0.04) (-0.844,5.355)rectangle++(0.04,0.04) (-0.506,6.181)rectangle++(0.04,0.04)
  (-0.185,6.871)rectangle++(0.04,0.04) (0.116,7.444)rectangle++(0.04,0.04) (0.394,7.919)rectangle++(0.04,0.04) (0.649,8.311)rectangle++(0.04,0.04) (0.882,8.635)rectangle++(0.04,0.04)
  (1.093,8.901)rectangle++(0.04,0.04) (1.283,9.118)rectangle++(0.04,0.04) (1.453,9.295)rectangle++(0.04,0.04) (2.132,4.084)rectangle++(0.04,0.04) (1.973,4.983)rectangle++(0.04,0.04)
  (1.874,5.752)rectangle++(0.04,0.04) (1.821,6.411)rectangle++(0.04,0.04) (1.802,6.973)rectangle++(0.04,0.04) (1.808,7.451)rectangle++(0.04,0.04) (1.833,7.857)rectangle++(0.04,0.04)
  (1.87,8.201)rectangle++(0.04,0.04) (1.915,8.493)rectangle++(0.04,0.04) (1.966,8.738)rectangle++(0.04,0.04) (2.018,8.945)rectangle++(0.04,0.04) (2.071,9.119)rectangle++(0.04,0.04)
  (2.016,4.191)rectangle++(0.04,0.04) (-0.485,4.153)rectangle++(0.04,0.04) (-0.248,5.145)rectangle++(0.04,0.04) (-0.007,5.98)rectangle++(0.04,0.04) (0.231,6.679)rectangle++(0.04,0.04)
  (0.461,7.264)rectangle++(0.04,0.04) (0.68,7.752)rectangle++(0.04,0.04) (1.84,3.612)rectangle++(0.04,0.04) (-0.393,4.006)rectangle++(0.04,0.04) (-0.176,5.017)rectangle++(0.04,0.04)
  (0.049,5.867)rectangle++(0.04,0.04) (0.274,6.582)rectangle++(0.04,0.04) (0.494,7.18)rectangle++(0.04,0.04) (0.705,7.679)rectangle++(0.04,0.04) (0.904,8.095)rectangle++(0.04,0.04)
  (1.899,3.728)rectangle++(0.04,0.04) (1.761,4.689)rectangle++(0.04,0.04) (1.682,5.512)rectangle++(0.04,0.04) (1.648,6.214)rectangle++(0.04,0.04) (1.647,6.812)rectangle++(0.04,0.04)
  (1.03,2.299)rectangle++(0.04,0.04) (0.965,3.509)rectangle++(0.04,0.04) (0.958,4.54)rectangle++(0.04,0.04) (0.994,5.417)rectangle++(0.04,0.04) (1.059,6.161)rectangle++(0.04,0.04)
  (1.145,6.791)rectangle++(0.04,0.04) (1.242,7.322)rectangle++(0.04,0.04) (1.347,7.771)rectangle++(0.04,0.04) (1.453,8.147)rectangle++(0.04,0.04) (1.847,3.863)rectangle++(0.04,0.04)
  (1.723,4.806)rectangle++(0.04,0.04) (1.654,5.612)rectangle++(0.04,0.04) (1.628,6.301)rectangle++(0.04,0.04) (1.634,6.887)rectangle++(0.04,0.04) (-1.125,4.613)rectangle++(0.04,0.04)
  (-0.774,5.563)rectangle++(0.04,0.04) (-0.438,6.355)rectangle++(0.04,0.04) (-0.12,7.016)rectangle++(0.04,0.04) (0.176,7.564)rectangle++(0.04,0.04) (0.45,8.019)rectangle++(0.04,0.04)
  (0.701,8.394)rectangle++(0.04,0.04) (0.93,8.703)rectangle++(0.04,0.04) (1.136,8.957)rectangle++(0.04,0.04) (1.322,9.164)rectangle++(0.04,0.04) (1.488,9.333)rectangle++(0.04,0.04)
  (1.636,9.47)rectangle++(0.04,0.04) (1.767,9.58)rectangle++(0.04,0.04) (1.883,9.668)rectangle++(0.04,0.04) (-1.671,5.209)rectangle++(0.04,0.04) (-1.214,6.091)rectangle++(0.04,0.04)
  (-0.791,6.822)rectangle++(0.04,0.04) (-0.401,7.427)rectangle++(0.04,0.04) (-0.046,7.925)rectangle++(0.04,0.04) (0.276,8.334)rectangle++(0.04,0.04) (0.565,8.669)rectangle++(0.04,0.04)
  (-1.695,4.765)rectangle++(0.04,0.04) (1.379,2.496)rectangle++(0.04,0.04) (1.269,3.663)rectangle++(0.04,0.04) (0.004,2.081)rectangle++(0.04,0.04) (0.084,3.365)rectangle++(0.04,0.04)
  (0.204,4.453)rectangle++(0.04,0.04) (0.349,5.373)rectangle++(0.04,0.04) (-1.034,4.074)rectangle++(0.04,0.04) (-0.718,5.1)rectangle++(0.04,0.04) (-0.409,5.96)rectangle++(0.04,0.04)
  (-0.111,6.679)rectangle++(0.04,0.04) (0.17,7.278)rectangle++(0.04,0.04) (0.434,7.775)rectangle++(0.04,0.04) (1.875,3.56)rectangle++(0.04,0.04) (1.734,4.547)rectangle++(0.04,0.04)
  (1.653,5.392)rectangle++(0.04,0.04) (1.619,6.113)rectangle++(0.04,0.04) (1.618,6.728)rectangle++(0.04,0.04) (1.642,7.25)rectangle++(0.04,0.04) (1.684,7.693)rectangle++(0.04,0.04)
  (1.737,8.068)rectangle++(0.04,0.04) (1.797,8.384)rectangle++(0.04,0.04) (1.86,8.651)rectangle++(0.04,0.04) (1.925,8.875)rectangle++(0.04,0.04) (-0.644,3.075)rectangle++(0.04,0.04)
  (-0.427,4.236)rectangle++(0.04,0.04) (-0.195,5.214)rectangle++(0.04,0.04) (0.04,6.036)rectangle++(0.04,0.04) (0.273,6.725)rectangle++(0.04,0.04) (0.499,7.301)rectangle++(0.04,0.04)
  (0.027,2.576)rectangle++(0.04,0.04) (0.124,3.784)rectangle++(0.04,0.04) (0.254,4.808)rectangle++(0.04,0.04) (1.161,2.947)rectangle++(0.04,0.04) (1.102,4.055)rectangle++(0.04,0.04)
  (1.097,4.999)rectangle++(0.04,0.04) (1.13,5.801)rectangle++(0.04,0.04) (1.191,6.482)rectangle++(0.04,0.04) (1.269,7.058)rectangle++(0.04,0.04) (1.359,7.545)rectangle++(0.04,0.04)
  (1.455,7.955)rectangle++(0.04,0.04) (1.12,2.262)rectangle++(0.04,0.04) (1.041,3.474)rectangle++(0.04,0.04) (1.021,4.508)rectangle++(0.04,0.04) (1.046,5.387)rectangle++(0.04,0.04)
  (1.103,6.133)rectangle++(0.04,0.04) (1.18,6.765)rectangle++(0.04,0.04) (1.272,7.3)rectangle++(0.04,0.04) (-1.278,4.627)rectangle++(0.04,0.04) (-0.903,5.58)rectangle++(0.04,0.04);
\fill[piece2]
  (-0.547,6.376)rectangle++(0.04,0.04) (-0.212,7.037)rectangle++(0.04,0.04) (-1.488,4.297)rectangle++(0.04,0.04) (-1.095,5.308)rectangle++(0.04,0.04) (-0.721,6.152)rectangle++(0.04,0.04)
  (1.608,2.985)rectangle++(0.04,0.04) (1.484,4.069)rectangle++(0.04,0.04) (1.422,4.996)rectangle++(0.04,0.04) (1.071,3.248)rectangle++(0.04,0.04) (1.038,4.314)rectangle++(0.04,0.04)
  (1.052,5.222)rectangle++(0.04,0.04) (1.101,5.993)rectangle++(0.04,0.04) (1.174,6.646)rectangle++(0.04,0.04) (1.261,7.198)rectangle++(0.04,0.04) (0.081,3.337)rectangle++(0.04,0.04)
  (-0.65,3.622)rectangle++(0.04,0.04) (-0.41,4.701)rectangle++(0.04,0.04) (-0.163,5.608)rectangle++(0.04,0.04) (0.084,6.37)rectangle++(0.04,0.04) (0.324,7.007)rectangle++(0.04,0.04)
  (0.553,7.539)rectangle++(0.04,0.04) (0.77,7.982)rectangle++(0.04,0.04) (-1.513,4.668)rectangle++(0.04,0.04) (-1.101,5.624)rectangle++(0.04,0.04) (-0.713,6.421)rectangle++(0.04,0.04)
  (-0.352,7.082)rectangle++(0.04,0.04) (-1.523,4.28)rectangle++(0.04,0.04) (-0.006,1.879)rectangle++(0.04,0.04) (0.068,3.194)rectangle++(0.04,0.04) (0.183,4.308)rectangle++(0.04,0.04)
  (0.169,2.088)rectangle++(0.04,0.04) (0.225,3.364)rectangle++(0.04,0.04) (0.324,4.446)rectangle++(0.04,0.04) (0.451,5.363)rectangle++(0.04,0.04) (0.595,6.137)rectangle++(0.04,0.04)
  (1.45,3.281)rectangle++(0.04,0.04) (1.361,4.327)rectangle++(0.04,0.04) (1.328,5.22)rectangle++(0.04,0.04) (1.335,5.98)rectangle++(0.04,0.04) (1.372,6.625)rectangle++(0.04,0.04)
  (-1.117,4.513)rectangle++(0.04,0.04) (-1.114,3.629)rectangle++(0.04,0.04) (-0.804,4.726)rectangle++(0.04,0.04) (-0.497,5.645)rectangle++(0.04,0.04) (-0.199,6.414)rectangle++(0.04,0.04)
  (-0.04,3.665)rectangle++(0.04,0.04) (0.111,4.713)rectangle++(0.04,0.04) (0.281,5.598)rectangle++(0.04,0.04) (0.46,6.343)rectangle++(0.04,0.04) (0.643,6.97)rectangle++(0.04,0.04)
  (-1.32,4.444)rectangle++(0.04,0.04) (-0.947,5.426)rectangle++(0.04,0.04) (-0.59,6.247)rectangle++(0.04,0.04) (-0.254,6.929)rectangle++(0.04,0.04) (0.059,7.496)rectangle++(0.04,0.04)
  (0.348,7.966)rectangle++(0.04,0.04) (0.612,8.353)rectangle++(0.04,0.04) (0.852,8.672)rectangle++(0.04,0.04) (1.069,8.933)rectangle++(0.04,0.04) (1.264,9.147)rectangle++(0.04,0.04)
  (-1.671,4.99)rectangle++(0.04,0.04) (-1.223,5.905)rectangle++(0.04,0.04) (-0.806,6.664)rectangle++(0.04,0.04) (-0.42,7.293)rectangle++(0.04,0.04) (-0.068,7.812)rectangle++(0.04,0.04)
  (0.253,8.239)rectangle++(0.04,0.04) (0.542,8.589)rectangle++(0.04,0.04) (0.802,8.875)rectangle++(0.04,0.04) (1.035,9.108)rectangle++(0.04,0.04) (1.242,9.297)rectangle++(0.04,0.04)
  (1.425,9.449)rectangle++(0.04,0.04) (1.587,9.571)rectangle++(0.04,0.04) (1.73,9.668)rectangle++(0.04,0.04) (1.855,9.745)rectangle++(0.04,0.04) (1.964,9.805)rectangle++(0.04,0.04)
  (2.059,9.852)rectangle++(0.04,0.04) (2.142,9.888)rectangle++(0.04,0.04) (2.214,9.915)rectangle++(0.04,0.04) (2.277,9.936)rectangle++(0.04,0.04) (2.33,9.95)rectangle++(0.04,0.04)
  (2.377,9.961)rectangle++(0.04,0.04) (2.416,9.968)rectangle++(0.04,0.04) (2.45,9.972)rectangle++(0.04,0.04) (2.48,9.975)rectangle++(0.04,0.04) (2.504,9.975)rectangle++(0.04,0.04)
  (2.526,9.975)rectangle++(0.04,0.04) (2.544,9.974)rectangle++(0.04,0.04) (2.559,9.972)rectangle++(0.04,0.04) (2.572,9.97)rectangle++(0.04,0.04) (2.583,9.968)rectangle++(0.04,0.04)
  (2.592,9.966)rectangle++(0.04,0.04) (2.599,9.963)rectangle++(0.04,0.04) (2.606,9.961)rectangle++(0.04,0.04) (2.611,9.959)rectangle++(0.04,0.04) (2.616,9.957)rectangle++(0.04,0.04)
  (2.619,9.955)rectangle++(0.04,0.04) (2.622,9.953)rectangle++(0.04,0.04) (-1.599,5.404)rectangle++(0.04,0.04) (-1.145,6.254)rectangle++(0.04,0.04) (-0.725,6.958)rectangle++(0.04,0.04)
  (-0.34,7.539)rectangle++(0.04,0.04) (0.01,8.018)rectangle++(0.04,0.04) (0.327,8.411)rectangle++(0.04,0.04) (0.612,8.733)rectangle++(0.04,0.04) (0.868,8.994)rectangle++(0.04,0.04)
  (1.095,9.207)rectangle++(0.04,0.04) (1.297,9.378)rectangle++(0.04,0.04) (1.475,9.516)rectangle++(0.04,0.04) (-0.28,3.241)rectangle++(0.04,0.04) (-0.11,4.362)rectangle++(0.04,0.04)
  (0.079,5.309)rectangle++(0.04,0.04) (0.277,6.105)rectangle++(0.04,0.04) (0.477,6.775)rectangle++(0.04,0.04) (0.675,7.336)rectangle++(0.04,0.04) (0.865,7.804)rectangle++(0.04,0.04)
  (1.045,8.195)rectangle++(0.04,0.04) (1.214,8.521)rectangle++(0.04,0.04) (1.37,8.79)rectangle++(0.04,0.04) (1.514,9.013)rectangle++(0.04,0.04) (1.645,9.197)rectangle++(0.04,0.04)
  (1.764,9.348)rectangle++(0.04,0.04) (1.871,9.471)rectangle++(0.04,0.04) (-1.631,5.168)rectangle++(0.04,0.04) (-1.182,6.054)rectangle++(0.04,0.04) (-0.765,6.79)rectangle++(0.04,0.04)
  (-0.381,7.398)rectangle++(0.04,0.04) (-0.03,7.9)rectangle++(0.04,0.04) (0.289,8.312)rectangle++(0.04,0.04) (0.576,8.65)rectangle++(0.04,0.04) (0.833,8.926)rectangle++(0.04,0.04)
  (1.063,9.15)rectangle++(0.04,0.04) (1.267,9.331)rectangle++(0.04,0.04) (1.448,9.477)rectangle++(0.04,0.04) (1.608,9.594)rectangle++(0.04,0.04) (2.187,4.177)rectangle++(0.04,0.04)
  (2.024,5.059)rectangle++(0.04,0.04) (-0.667,4.326)rectangle++(0.04,0.04) (-0.397,5.3)rectangle++(0.04,0.04) (-0.127,6.117)rectangle++(0.04,0.04) (0.134,6.801)rectangle++(0.04,0.04)
  (0.384,7.371)rectangle++(0.04,0.04) (-1.452,4.475)rectangle++(0.04,0.04) (-1.057,5.458)rectangle++(0.04,0.04) (-0.793,3.141)rectangle++(0.04,0.04) (-0.551,4.297)rectangle++(0.04,0.04)
  (-0.298,5.271)rectangle++(0.04,0.04) (-0.045,6.089)rectangle++(0.04,0.04) (0.203,6.773)rectangle++(0.04,0.04) (1.653,3.314)rectangle++(0.04,0.04) (0.627,2.977)rectangle++(0.04,0.04)
  (0.65,4.101)rectangle++(0.04,0.04) (0.714,5.056)rectangle++(0.04,0.04) (0.807,5.866)rectangle++(0.04,0.04) (0.919,6.55)rectangle++(0.04,0.04) (1.041,7.127)rectangle++(0.04,0.04)
  (-1.282,4.55)rectangle++(0.04,0.04) (-0.91,5.515)rectangle++(0.04,0.04) (-0.555,6.32)rectangle++(0.04,0.04) (-0.221,6.99)rectangle++(0.04,0.04) (0.089,7.547)rectangle++(0.04,0.04)
  (0.375,8.008)rectangle++(0.04,0.04) (0.637,8.388)rectangle++(0.04,0.04) (0.875,8.7)rectangle++(0.04,0.04) (1.09,8.956)rectangle++(0.04,0.04) (1.282,9.166)rectangle++(0.04,0.04)
  (-1.672,4.997)rectangle++(0.04,0.04) (-1.223,5.911)rectangle++(0.04,0.04) (-0.806,6.669)rectangle++(0.04,0.04) (-0.42,7.297)rectangle++(0.04,0.04) (1.856,3.274)rectangle++(0.04,0.04)
  (1.707,4.305)rectangle++(0.04,0.04) (1.621,5.187)rectangle++(0.04,0.04) (1.583,5.94)rectangle++(0.04,0.04) (1.581,6.582)rectangle++(0.04,0.04) (1.605,7.128)rectangle++(0.04,0.04)
  (-1.182,4.656)rectangle++(0.04,0.04) (-0.82,5.601)rectangle++(0.04,0.04) (-0.475,6.39)rectangle++(0.04,0.04) (-0.151,7.046)rectangle++(0.04,0.04) (0.152,7.592)rectangle++(0.04,0.04)
  (0.43,8.043)rectangle++(0.04,0.04) (0.685,8.416)rectangle++(0.04,0.04) (0.917,8.722)rectangle++(0.04,0.04) (1.126,8.973)rectangle++(0.04,0.04) (1.314,9.178)rectangle++(0.04,0.04)
  (1.482,9.345)rectangle++(0.04,0.04) (1.631,9.48)rectangle++(0.04,0.04) (1.763,9.589)rectangle++(0.04,0.04) (1.88,9.677)rectangle++(0.04,0.04) (1.983,9.746)rectangle++(0.04,0.04)
  (-1.669,5.243)rectangle++(0.04,0.04) (-1.211,6.12)rectangle++(0.04,0.04) (-0.787,6.847)rectangle++(0.04,0.04) (1.788,3.104)rectangle++(0.04,0.04) (1.641,4.163)rectangle++(0.04,0.04)
  (1.56,5.069)rectangle++(0.04,0.04) (1.526,5.842)rectangle++(0.04,0.04) (-0.923,4.387)rectangle++(0.04,0.04) (-0.611,5.362)rectangle++(0.04,0.04) (-0.307,6.179)rectangle++(0.04,0.04)
  (-0.016,6.86)rectangle++(0.04,0.04) (0.258,7.428)rectangle++(0.04,0.04) (0.515,7.9)rectangle++(0.04,0.04) (0.031,2.657)rectangle++(0.04,0.04) (0.13,3.853)rectangle++(0.04,0.04)
  (0.263,4.866)rectangle++(0.04,0.04) (0.416,5.722)rectangle++(0.04,0.04) (1.368,3.16)rectangle++(0.04,0.04) (1.287,4.227)rectangle++(0.04,0.04) (1.261,5.138)rectangle++(0.04,0.04)
  (1.275,5.913)rectangle++(0.04,0.04) (1.318,6.571)rectangle++(0.04,0.04) (1.381,7.129)rectangle++(0.04,0.04) (1.457,7.601)rectangle++(0.04,0.04) (-1.308,4.721)rectangle++(0.04,0.04)
  (-0.926,5.662)rectangle++(0.04,0.04) (-0.562,6.446)rectangle++(0.04,0.04) (-0.222,7.097)rectangle++(0.04,0.04) (0.093,7.638)rectangle++(0.04,0.04) (0.382,8.085)rectangle++(0.04,0.04)
  (0.646,8.453)rectangle++(0.04,0.04) (0.885,8.755)rectangle++(0.04,0.04) (1.1,9.003)rectangle++(0.04,0.04) (1.293,9.205)rectangle++(0.04,0.04) (1.465,9.368)rectangle++(0.04,0.04)
  (-1.682,5.072)rectangle++(0.04,0.04) (-1.229,5.974)rectangle++(0.04,0.04) (-0.808,6.724)rectangle++(0.04,0.04) (-0.42,7.344)rectangle++(0.04,0.04) (1.868,3.283)rectangle++(0.04,0.04)
  (1.717,4.312)rectangle++(0.04,0.04) (1.629,5.193)rectangle++(0.04,0.04) (1.591,5.945)rectangle++(0.04,0.04) (1.286,3.47)rectangle++(0.04,0.04) (1.23,4.494)rectangle++(0.04,0.04)
  (1.223,5.367)rectangle++(0.04,0.04) (1.252,6.109)rectangle++(0.04,0.04) (1.306,6.739)rectangle++(0.04,0.04) (1.378,7.272)rectangle++(0.04,0.04) (1.46,7.722)rectangle++(0.04,0.04)
  (1.547,8.102)rectangle++(0.04,0.04) (1.637,8.421)rectangle++(0.04,0.04) (1.726,8.688)rectangle++(0.04,0.04) (-1.491,4.986)rectangle++(0.04,0.04) (-1.07,5.894)rectangle++(0.04,0.04)
  (-0.676,6.649)rectangle++(0.04,0.04) (-0.311,7.275)rectangle++(0.04,0.04) (0.025,7.792)rectangle++(0.04,0.04) (0.33,8.219)rectangle++(0.04,0.04) (0.607,8.569)rectangle++(0.04,0.04)
  (0.857,8.856)rectangle++(0.04,0.04) (-1.682,4.856)rectangle++(0.04,0.04) (-1.238,5.791)rectangle++(0.04,0.04) (-0.823,6.568)rectangle++(0.04,0.04) (-0.439,7.212)rectangle++(0.04,0.04)
  (-0.087,7.744)rectangle++(0.04,0.04) (0.234,8.182)rectangle++(0.04,0.04) (0.524,8.542)rectangle++(0.04,0.04) (0.785,8.836)rectangle++(0.04,0.04) (1.018,9.075)rectangle++(0.04,0.04)
  (2.128,3.943)rectangle++(0.04,0.04) (1.964,4.863)rectangle++(0.04,0.04) (0.011,2.244)rectangle++(0.04,0.04) (0.097,3.503)rectangle++(0.04,0.04) (0.22,4.57)rectangle++(0.04,0.04)
  (0.368,5.472)rectangle++(0.04,0.04) (0.529,6.233)rectangle++(0.04,0.04) (0.697,6.873)rectangle++(0.04,0.04) (0.865,7.41)rectangle++(0.04,0.04) (1.03,7.86)rectangle++(0.04,0.04)
  (1.187,8.236)rectangle++(0.04,0.04) (1.337,8.549)rectangle++(0.04,0.04) (1.476,8.81)rectangle++(0.04,0.04) (2.009,3.996)rectangle++(0.04,0.04) (1.865,4.912)rectangle++(0.04,0.04)
  (1.78,5.697)rectangle++(0.04,0.04) (1.738,6.368)rectangle++(0.04,0.04) (1.73,6.939)rectangle++(0.04,0.04) (1.746,7.425)rectangle++(0.04,0.04) (1.779,7.838)rectangle++(0.04,0.04)
  (1.823,8.187)rectangle++(0.04,0.04) (1.875,8.482)rectangle++(0.04,0.04) (1.931,8.731)rectangle++(0.04,0.04) (1.936,4.084)rectangle++(0.04,0.04) (1.807,4.99)rectangle++(0.04,0.04)
  (1.733,5.766)rectangle++(0.04,0.04) (0.017,2.367)rectangle++(0.04,0.04) (0.107,3.607)rectangle++(0.04,0.04) (0.233,4.658)rectangle++(0.04,0.04) (1.126,2.913)rectangle++(0.04,0.04)
  (1.072,4.028)rectangle++(0.04,0.04) (1.07,4.977)rectangle++(0.04,0.04) (1.106,5.784)rectangle++(0.04,0.04) (1.169,6.468)rectangle++(0.04,0.04) (1.251,7.047)rectangle++(0.04,0.04)
  (1.343,7.536)rectangle++(0.04,0.04) (1.709,3.72)rectangle++(0.04,0.04) (1.599,4.69)rectangle++(0.04,0.04) (1.545,5.519)rectangle++(0.04,0.04) (-0.851,4.327)rectangle++(0.04,0.04)
  (-0.552,5.308)rectangle++(0.04,0.04) (-0.259,6.13)rectangle++(0.04,0.04) (0.022,6.817)rectangle++(0.04,0.04) (0.29,7.39)rectangle++(0.04,0.04) (0.54,7.866)rectangle++(0.04,0.04)
  (0.03,2.652)rectangle++(0.04,0.04) (0.13,3.85)rectangle++(0.04,0.04) (0.262,4.863)rectangle++(0.04,0.04) (0.26,2.58)rectangle++(0.04,0.04) (0.322,3.779)rectangle++(0.04,0.04)
  (0.898,2.764)rectangle++(0.04,0.04) (0.579,2.699)rectangle++(0.04,0.04) (0.598,3.867)rectangle++(0.04,0.04) (0.661,4.859)rectangle++(0.04,0.04) (0.754,5.7)rectangle++(0.04,0.04)
  (0.867,6.411)rectangle++(0.04,0.04) (0.991,7.011)rectangle++(0.04,0.04) (1.121,7.516)rectangle++(0.04,0.04) (1.251,7.94)rectangle++(0.04,0.04) (1.379,8.295)rectangle++(0.04,0.04)
  (1.501,8.592)rectangle++(0.04,0.04) (1.618,8.839)rectangle++(0.04,0.04) (1.726,9.045)rectangle++(0.04,0.04) (1.827,9.215)rectangle++(0.04,0.04) (1.919,9.356)rectangle++(0.04,0.04)
  (2.004,9.472)rectangle++(0.04,0.04) (2.08,9.567)rectangle++(0.04,0.04) (2.148,9.645)rectangle++(0.04,0.04) (2.21,9.709)rectangle++(0.04,0.04) (2.264,9.76)rectangle++(0.04,0.04)
  (2.313,9.802)rectangle++(0.04,0.04) (2.356,9.835)rectangle++(0.04,0.04) (2.394,9.862)rectangle++(0.04,0.04) (2.427,9.883)rectangle++(0.04,0.04) (2.456,9.9)rectangle++(0.04,0.04)
  (2.481,9.913)rectangle++(0.04,0.04) (2.503,9.923)rectangle++(0.04,0.04) (2.523,9.93)rectangle++(0.04,0.04) (2.539,9.936)rectangle++(0.04,0.04) (2.554,9.94)rectangle++(0.04,0.04)
  (2.566,9.943)rectangle++(0.04,0.04) (2.577,9.945)rectangle++(0.04,0.04) (2.586,9.947)rectangle++(0.04,0.04) (2.594,9.947)rectangle++(0.04,0.04) (2.6,9.948)rectangle++(0.04,0.04)
  (2.606,9.948)rectangle++(0.04,0.04) (2.611,9.948)rectangle++(0.04,0.04) (2.615,9.947)rectangle++(0.04,0.04) (2.618,9.947)rectangle++(0.04,0.04) (2.621,9.946)rectangle++(0.04,0.04)
  (2.624,9.946)rectangle++(0.04,0.04) (2.626,9.945)rectangle++(0.04,0.04) (2.627,9.944)rectangle++(0.04,0.04) (2.629,9.944)rectangle++(0.04,0.04) (2.63,9.943)rectangle++(0.04,0.04)
  (2.631,9.943)rectangle++(0.04,0.04) (2.632,9.942)rectangle++(0.04,0.04) (2.633,9.942)rectangle++(0.04,0.04) (2.633,9.942)rectangle++(0.04,0.04) (2.634,9.941)rectangle++(0.04,0.04)
  (2.634,9.941)rectangle++(0.04,0.04) (2.634,9.941)rectangle++(0.04,0.04) (2.153,4.477)rectangle++(0.04,0.04) (2.007,5.315)rectangle++(0.04,0.04) (1.916,6.034)rectangle++(0.04,0.04)
  (1.868,6.649)rectangle++(0.04,0.04) (0.023,2.487)rectangle++(0.04,0.04) (0.117,3.709)rectangle++(0.04,0.04) (0.245,4.744)rectangle++(0.04,0.04) (-0.919,3.93)rectangle++(0.04,0.04)
  (-0.955,3.551)rectangle++(0.04,0.04) (-0.672,4.653)rectangle++(0.04,0.04) (-0.388,5.578)rectangle++(0.04,0.04) (-0.108,6.353)rectangle++(0.04,0.04) (0.16,7.001)rectangle++(0.04,0.04)
  (0.414,7.54)rectangle++(0.04,0.04) (0.651,7.989)rectangle++(0.04,0.04) (0.871,8.361)rectangle++(0.04,0.04) (1.072,8.668)rectangle++(0.04,0.04) (1.256,8.921)rectangle++(0.04,0.04)
  (1.422,9.129)rectangle++(0.04,0.04) (1.572,9.299)rectangle++(0.04,0.04) (1.706,9.437)rectangle++(0.04,0.04) (1.825,9.55)rectangle++(0.04,0.04) (1.931,9.64)rectangle++(0.04,0.04)
  (2.025,9.713)rectangle++(0.04,0.04) (2.108,9.772)rectangle++(0.04,0.04) (2.18,9.818)rectangle++(0.04,0.04) (2.244,9.854)rectangle++(0.04,0.04) (2.177,4.372)rectangle++(0.04,0.04)
  (2.023,5.225)rectangle++(0.04,0.04) (1.927,5.956)rectangle++(0.04,0.04) (1.874,6.582)rectangle++(0.04,0.04) (1.854,7.116)rectangle++(0.04,0.04) (1.858,7.571)rectangle++(0.04,0.04)
  (1.88,7.957)rectangle++(0.04,0.04) (1.914,8.284)rectangle++(0.04,0.04) (1.956,8.561)rectangle++(0.04,0.04) (2.003,8.795)rectangle++(0.04,0.04) (2.052,8.992)rectangle++(0.04,0.04)
  (1.986,4.162)rectangle++(0.04,0.04) (1.853,5.054)rectangle++(0.04,0.04) (1.775,5.818)rectangle++(0.04,0.04) (1.739,6.471)rectangle++(0.04,0.04) (-0.199,3.716)rectangle++(0.04,0.04)
  (-0.022,4.763)rectangle++(0.04,0.04) (0.169,5.645)rectangle++(0.04,0.04) (0.368,6.388)rectangle++(0.04,0.04) (0.566,7.011)rectangle++(0.04,0.04) (0.759,7.533)rectangle++(0.04,0.04)
  (0.944,7.969)rectangle++(0.04,0.04) (1.119,8.332)rectangle++(0.04,0.04) (1.932,3.823)rectangle++(0.04,0.04) (1.793,4.768)rectangle++(0.04,0.04) (1.713,5.578)rectangle++(0.04,0.04)
  (1.677,6.269)rectangle++(0.04,0.04) (1.674,6.857)rectangle++(0.04,0.04) (1.695,7.358)rectangle++(0.04,0.04) (1.733,7.783)rectangle++(0.04,0.04) (1.782,8.142)rectangle++(0.04,0.04)
  (1.838,8.446)rectangle++(0.04,0.04) (1.898,8.702)rectangle++(0.04,0.04) (1.959,8.917)rectangle++(0.04,0.04) (2.02,9.097)rectangle++(0.04,0.04) (2.078,9.248)rectangle++(0.04,0.04)
  (2.134,9.374)rectangle++(0.04,0.04) (2.187,9.479)rectangle++(0.04,0.04) (-1.576,5.229)rectangle++(0.04,0.04) (-1.133,6.104)rectangle++(0.04,0.04) (-0.721,6.83)rectangle++(0.04,0.04)
  (-0.342,7.43)rectangle++(0.04,0.04) (0.005,7.926)rectangle++(0.04,0.04) (0.319,8.333)rectangle++(0.04,0.04) (0.602,8.666)rectangle++(0.04,0.04) (-0.107,3.909)rectangle++(0.04,0.04)
  (0.063,4.923)rectangle++(0.04,0.04) (0.248,5.778)rectangle++(0.04,0.04) (0.44,6.498)rectangle++(0.04,0.04) (0.632,7.102)rectangle++(0.04,0.04) (-1.35,4.468)rectangle++(0.04,0.04)
  (-0.971,5.448)rectangle++(0.04,0.04) (-0.61,6.266)rectangle++(0.04,0.04) (-0.27,6.947)rectangle++(0.04,0.04) (0.046,7.512)rectangle++(0.04,0.04) (0.338,7.979)rectangle++(0.04,0.04)
  (0.604,8.365)rectangle++(0.04,0.04) (0.846,8.682)rectangle++(0.04,0.04) (2.051,3.828)rectangle++(0.04,0.04) (1.894,4.768)rectangle++(0.04,0.04) (1.799,5.573)rectangle++(0.04,0.04)
  (1.75,6.262)rectangle++(0.04,0.04) (1.346,3.567)rectangle++(0.04,0.04) (1.284,4.574)rectangle++(0.04,0.04) (1.272,5.433)rectangle++(0.04,0.04) (1.297,6.163)rectangle++(0.04,0.04)
  (1.347,6.783)rectangle++(0.04,0.04) (1.414,7.308)rectangle++(0.04,0.04) (1.492,7.751)rectangle++(0.04,0.04) (1.576,8.125)rectangle++(0.04,0.04) (1.662,8.44)rectangle++(0.04,0.04)
  (-1.45,4.925)rectangle++(0.04,0.04) (-1.037,5.841)rectangle++(0.04,0.04) (-0.65,6.602)rectangle++(0.04,0.04) (-0.291,7.234)rectangle++(0.04,0.04) (0.04,7.757)rectangle++(0.04,0.04)
  (0.342,8.188)rectangle++(0.04,0.04) (0.616,8.542)rectangle++(0.04,0.04) (0.863,8.833)rectangle++(0.04,0.04) (1.085,9.069)rectangle++(0.04,0.04) (1.283,9.262)rectangle++(0.04,0.04)
  (1.458,9.417)rectangle++(0.04,0.04) (1.614,9.543)rectangle++(0.04,0.04) (1.752,9.643)rectangle++(0.04,0.04) (-1.689,5.18)rectangle++(0.04,0.04) (-1.358,3.653)rectangle++(0.04,0.04)
  (-1.011,4.756)rectangle++(0.04,0.04) (-0.671,5.679)rectangle++(0.04,0.04) (1.485,2.905)rectangle++(0.04,0.04) (1.376,4.006)rectangle++(0.04,0.04) (1.328,4.946)rectangle++(0.04,0.04)
  (1.324,5.747)rectangle++(0.04,0.04) (1.353,6.429)rectangle++(0.04,0.04) (1.405,7.006)rectangle++(0.04,0.04) (1.473,7.495)rectangle++(0.04,0.04) (1.549,7.908)rectangle++(0.04,0.04)
  (1.631,8.256)rectangle++(0.04,0.04) (0.004,2.097)rectangle++(0.04,0.04) (0.085,3.378)rectangle++(0.04,0.04) (0.205,4.464)rectangle++(0.04,0.04) (0.351,5.383)rectangle++(0.04,0.04)
  (0.511,6.157)rectangle++(0.04,0.04) (0.679,6.81)rectangle++(0.04,0.04) (0.847,7.357)rectangle++(0.04,0.04) (1.723,3.573)rectangle++(0.04,0.04) (0.683,3.042)rectangle++(0.04,0.04)
  (0.7,4.155)rectangle++(0.04,0.04) (0.759,5.1)rectangle++(0.04,0.04) (1.174,3.118)rectangle++(0.04,0.04) (1.12,4.2)rectangle++(0.04,0.04) (1.118,5.121)rectangle++(0.04,0.04)
  (1.153,5.904)rectangle++(0.04,0.04) (1.214,6.569)rectangle++(0.04,0.04) (1.293,7.131)rectangle++(0.04,0.04) (1.382,7.606)rectangle++(0.04,0.04) (1.476,8.006)rectangle++(0.04,0.04)
  (1.573,8.342)rectangle++(0.04,0.04) (1.669,8.624)rectangle++(0.04,0.04) (1.761,8.86)rectangle++(0.04,0.04) (1.849,9.057)rectangle++(0.04,0.04) (1.932,9.22)rectangle++(0.04,0.04)
  (-1.567,5.13)rectangle++(0.04,0.04) (-1.129,6.019)rectangle++(0.04,0.04) (1.623,2.867)rectangle++(0.04,0.04) (1.492,3.968)rectangle++(0.04,0.04) (1.425,4.909)rectangle++(0.04,0.04)
  (1.405,5.712)rectangle++(0.04,0.04) (1.421,6.395)rectangle++(0.04,0.04) (1.461,6.975)rectangle++(0.04,0.04) (1.519,7.467)rectangle++(0.04,0.04) (1.587,7.882)rectangle++(0.04,0.04)
  (-1.345,4.801)rectangle++(0.04,0.04) (-0.953,5.731)rectangle++(0.04,0.04) (-0.583,6.506)rectangle++(0.04,0.04) (-1.442,4.123)rectangle++(0.04,0.04) (-1.063,5.158)rectangle++(0.04,0.04)
  (-0.699,6.023)rectangle++(0.04,0.04) (-0.356,6.744)rectangle++(0.04,0.04) (-0.035,7.343)rectangle++(0.04,0.04) (0.262,7.839)rectangle++(0.04,0.04) (1.91,3.533)rectangle++(0.04,0.04)
  (1.763,4.523)rectangle++(0.04,0.04) (1.677,5.37)rectangle++(0.04,0.04) (1.638,6.094)rectangle++(0.04,0.04) (1.634,6.71)rectangle++(0.04,0.04) (1.655,7.235)rectangle++(0.04,0.04)
  (-0.994,3.356)rectangle++(0.04,0.04) (0.948,2.384)rectangle++(0.04,0.04) (0.899,3.585)rectangle++(0.04,0.04) (0.905,4.607)rectangle++(0.04,0.04) (1.035,3.056)rectangle++(0.04,0.04)
  (1,4.152)rectangle++(0.04,0.04) (1.014,5.086)rectangle++(0.04,0.04) (1.063,5.879)rectangle++(0.04,0.04) (1.136,6.55)rectangle++(0.04,0.04) (1.489,3.483)rectangle++(0.04,0.04)
  (1.403,4.497)rectangle++(0.04,0.04) (1.37,5.363)rectangle++(0.04,0.04) (-0.849,4.263)rectangle++(0.04,0.04) (-0.553,5.254)rectangle++(0.04,0.04) (-0.262,6.084)rectangle++(0.04,0.04)
  (0.018,6.778)rectangle++(0.04,0.04) (0.284,7.357)rectangle++(0.04,0.04) (0.534,7.838)rectangle++(0.04,0.04) (0.765,8.237)rectangle++(0.04,0.04) (0.978,8.567)rectangle++(0.04,0.04)
  (1.171,8.839)rectangle++(0.04,0.04) (1.347,9.063)rectangle++(0.04,0.04) (1.505,9.246)rectangle++(0.04,0.04) (1.647,9.395)rectangle++(0.04,0.04) (1.774,9.516)rectangle++(0.04,0.04)
  (1.886,9.614)rectangle++(0.04,0.04) (1.985,9.692)rectangle++(0.04,0.04) (2.073,9.755)rectangle++(0.04,0.04) (2.15,9.805)rectangle++(0.04,0.04) (-1.652,5.286)rectangle++(0.04,0.04)
  (-0.497,3.102)rectangle++(0.04,0.04) (-0.3,4.252)rectangle++(0.04,0.04) (-0.087,5.223)rectangle++(0.04,0.04) (0.133,6.039)rectangle++(0.04,0.04) (0.352,6.724)rectangle++(0.04,0.04)
  (0.566,7.298)rectangle++(0.04,0.04) (0.771,7.777)rectangle++(0.04,0.04) (0.964,8.175)rectangle++(0.04,0.04) (1.144,8.507)rectangle++(0.04,0.04) (1.311,8.781)rectangle++(0.04,0.04)
  (1.463,9.008)rectangle++(0.04,0.04) (1.602,9.194)rectangle++(0.04,0.04) (2.089,4.099)rectangle++(0.04,0.04) (1.938,4.996)rectangle++(0.04,0.04) (1.845,5.766)rectangle++(0.04,0.04)
  (1.796,6.423)rectangle++(0.04,0.04) (1.782,6.984)rectangle++(0.04,0.04) (1.792,7.461)rectangle++(0.04,0.04) (1.819,7.867)rectangle++(0.04,0.04) (1.859,8.21)rectangle++(0.04,0.04)
  (1.906,8.501)rectangle++(0.04,0.04) (1.958,8.745)rectangle++(0.04,0.04) (2.012,8.951)rectangle++(0.04,0.04) (2.066,9.124)rectangle++(0.04,0.04) (2.018,4.191)rectangle++(0.04,0.04)
  (1.88,5.077)rectangle++(0.04,0.04) (1.799,5.837)rectangle++(0.04,0.04) (-0.203,3.654)rectangle++(0.04,0.04) (-0.029,4.71)rectangle++(0.04,0.04) (0.162,5.601)rectangle++(0.04,0.04)
  (0.359,6.351)rectangle++(0.04,0.04) (0.557,6.98)rectangle++(0.04,0.04) (0.751,7.507)rectangle++(0.04,0.04) (0.936,7.947)rectangle++(0.04,0.04) (1.111,8.314)rectangle++(0.04,0.04)
  (-1.522,4.797)rectangle++(0.04,0.04) (-1.104,5.735)rectangle++(0.04,0.04) (-0.711,6.515)rectangle++(0.04,0.04) (-0.346,7.162)rectangle++(0.04,0.04) (-0.01,7.698)rectangle++(0.04,0.04)
  (-1.592,4.465)rectangle++(0.04,0.04) (-1.177,5.456)rectangle++(0.04,0.04) (-0.784,6.28)rectangle++(0.04,0.04) (-0.418,6.966)rectangle++(0.04,0.04) (-0.078,7.534)rectangle++(0.04,0.04)
  (0.232,8.003)rectangle++(0.04,0.04) (0.516,8.39)rectangle++(0.04,0.04) (0.772,8.707)rectangle++(0.04,0.04) (2.066,3.816)rectangle++(0.04,0.04) (1.906,4.757)rectangle++(0.04,0.04)
  (1.808,5.564)rectangle++(0.04,0.04) (1.166,3.446)rectangle++(0.04,0.04) (1.127,4.478)rectangle++(0.04,0.04) (0.977,3.082)rectangle++(0.04,0.04) (0.952,4.176)rectangle++(0.04,0.04)
  (0.923,2.984)rectangle++(0.04,0.04) (0.902,4.095)rectangle++(0.04,0.04) (0.928,5.041)rectangle++(0.04,0.04) (0.988,5.844)rectangle++(0.04,0.04) (1.072,6.524)rectangle++(0.04,0.04)
  (1.17,7.099)rectangle++(0.04,0.04) (1.276,7.583)rectangle++(0.04,0.04) (1.386,7.991)rectangle++(0.04,0.04) (1.495,8.333)rectangle++(0.04,0.04) (1.602,8.62)rectangle++(0.04,0.04)
  (1.704,8.859)rectangle++(0.04,0.04) (-1.531,5.016)rectangle++(0.04,0.04) (-1.103,5.921)rectangle++(0.04,0.04) (-0.703,6.673)rectangle++(0.04,0.04) (-0.333,7.297)rectangle++(0.04,0.04)
  (-1.565,4.326)rectangle++(0.04,0.04) (-1.159,5.336)rectangle++(0.04,0.04) (-0.774,6.178)rectangle++(0.04,0.04) (-0.413,6.878)rectangle++(0.04,0.04) (-0.078,7.459)rectangle++(0.04,0.04)
  (0.23,7.94)rectangle++(0.04,0.04) (0.511,8.336)rectangle++(0.04,0.04) (0.765,8.661)rectangle++(0.04,0.04) (0.995,8.928)rectangle++(0.04,0.04) (-1.673,4.897)rectangle++(0.04,0.04)
  (-1.228,5.825)rectangle++(0.04,0.04) (-0.813,6.597)rectangle++(0.04,0.04) (-0.43,7.236)rectangle++(0.04,0.04) (1.842,3.26)rectangle++(0.04,0.04) (0.592,3.009)rectangle++(0.04,0.04)
  (0.621,4.13)rectangle++(0.04,0.04) (0.691,5.082)rectangle++(0.04,0.04) (-0.911,4.08)rectangle++(0.04,0.04) (-0.614,5.101)rectangle++(0.04,0.04) (-0.32,5.957)rectangle++(0.04,0.04)
  (-0.036,6.672)rectangle++(0.04,0.04) (0.234,7.269)rectangle++(0.04,0.04) (0.488,7.765)rectangle++(0.04,0.04) (0.723,8.177)rectangle++(0.04,0.04) (0.939,8.518)rectangle++(0.04,0.04)
  (1.137,8.799)rectangle++(0.04,0.04) (1.316,9.03)rectangle++(0.04,0.04) (1.478,9.219)rectangle++(0.04,0.04) (1.623,9.373)rectangle++(0.04,0.04) (1.752,9.498)rectangle++(0.04,0.04)
  (-1.658,5.151)rectangle++(0.04,0.04) (-1.206,6.041)rectangle++(0.04,0.04) (-0.785,6.779)rectangle++(0.04,0.04) (-0.399,7.39)rectangle++(0.04,0.04) (-0.045,7.894)rectangle++(0.04,0.04)
  (0.275,8.308)rectangle++(0.04,0.04) (2.025,3.627)rectangle++(0.04,0.04) (1.864,4.598)rectangle++(0.04,0.04) (1.766,5.43)rectangle++(0.04,0.04) (1.716,6.141)rectangle++(0.04,0.04)
  (-0.952,4.482)rectangle++(0.04,0.04) (-0.632,5.444)rectangle++(0.04,0.04) (-0.322,6.249)rectangle++(0.04,0.04) (0.916,2.131)rectangle++(0.04,0.04) (0.608,2.311)rectangle++(0.04,0.04)
  (0.607,3.536)rectangle++(0.04,0.04) (0.655,4.578)rectangle++(0.04,0.04) (0.738,5.461)rectangle++(0.04,0.04) (0.843,6.208)rectangle++(0.04,0.04) (0.963,6.84)rectangle++(0.04,0.04)
  (1.09,7.371)rectangle++(0.04,0.04) (1.219,7.818)rectangle++(0.04,0.04) (1.347,8.193)rectangle++(0.04,0.04) (0.033,2.697)rectangle++(0.04,0.04) (0.133,3.887)rectangle++(0.04,0.04)
  (0.267,4.895)rectangle++(0.04,0.04) (0.42,5.746)rectangle++(0.04,0.04) (0.585,6.464)rectangle++(0.04,0.04) (-0.908,3.271)rectangle++(0.04,0.04) (-0.643,4.413)rectangle++(0.04,0.04)
  (-0.372,5.373)rectangle++(0.04,0.04) (-0.104,6.178)rectangle++(0.04,0.04) (0.157,6.852)rectangle++(0.04,0.04) (0.405,7.414)rectangle++(0.04,0.04) (0.639,7.882)rectangle++(0.04,0.04)
  (0.856,8.27)rectangle++(0.04,0.04) (1.056,8.592)rectangle++(0.04,0.04) (1.239,8.857)rectangle++(0.04,0.04) (1.405,9.075)rectangle++(0.04,0.04) (1.555,9.254)rectangle++(0.04,0.04)
  (-1.181,3.337)rectangle++(0.04,0.04) (-0.873,4.48)rectangle++(0.04,0.04) (-0.565,5.439)rectangle++(0.04,0.04) (-0.265,6.242)rectangle++(0.04,0.04) (0.023,6.912)rectangle++(0.04,0.04)
  (1.709,3.3)rectangle++(0.04,0.04) (1.583,4.333)rectangle++(0.04,0.04) (-0.593,4.1)rectangle++(0.04,0.04) (1.085,2.618)rectangle++(0.04,0.04) (1.025,3.778)rectangle++(0.04,0.04)
  (1.02,4.767)rectangle++(0.04,0.04) (1.055,5.607)rectangle++(0.04,0.04) (1.119,6.32)rectangle++(0.04,0.04) (1.202,6.924)rectangle++(0.04,0.04) (1.296,7.433)rectangle++(0.04,0.04)
  (1.397,7.863)rectangle++(0.04,0.04) (1.5,8.223)rectangle++(0.04,0.04) (1.861,3.888)rectangle++(0.04,0.04) (1.735,4.827)rectangle++(0.04,0.04) (1.666,5.63)rectangle++(0.04,0.04)
  (1.639,6.315)rectangle++(0.04,0.04) (1.643,6.898)rectangle++(0.04,0.04) (1.671,7.394)rectangle++(0.04,0.04) (1.714,7.814)rectangle++(0.04,0.04) (1.767,8.17)rectangle++(0.04,0.04)
  (1.826,8.47)rectangle++(0.04,0.04) (1.889,8.722)rectangle++(0.04,0.04) (1.952,8.935)rectangle++(0.04,0.04) (2.015,9.113)rectangle++(0.04,0.04) (2.075,9.261)rectangle++(0.04,0.04)
  (2.132,9.385)rectangle++(0.04,0.04) (2.185,9.488)rectangle++(0.04,0.04) (2.235,9.574)rectangle++(0.04,0.04) (2.28,9.645)rectangle++(0.04,0.04) (-1.595,5.279)rectangle++(0.04,0.04)
  (-1.146,6.147)rectangle++(0.04,0.04) (-0.731,6.867)rectangle++(0.04,0.04) (-0.349,7.463)rectangle++(0.04,0.04) (0,7.953)rectangle++(0.04,0.04) (0.316,8.357)rectangle++(0.04,0.04)
  (0.601,8.687)rectangle++(0.04,0.04) (0.856,8.956)rectangle++(0.04,0.04) (1.083,9.174)rectangle++(0.04,0.04) (1.286,9.351)rectangle++(0.04,0.04) (1.465,9.493)rectangle++(0.04,0.04)
  (1.622,9.607)rectangle++(0.04,0.04) (1.761,9.697)rectangle++(0.04,0.04) (1.883,9.768)rectangle++(0.04,0.04) (1.989,9.824)rectangle++(0.04,0.04) (2.081,9.867)rectangle++(0.04,0.04)
  (2.162,9.9)rectangle++(0.04,0.04) (2.231,9.925)rectangle++(0.04,0.04) (2.291,9.943)rectangle++(0.04,0.04) (2.343,9.956)rectangle++(0.04,0.04) (2.388,9.965)rectangle++(0.04,0.04)
  (2.426,9.971)rectangle++(0.04,0.04) (2.459,9.975)rectangle++(0.04,0.04) (2.487,9.976)rectangle++(0.04,0.04) (2.51,9.977)rectangle++(0.04,0.04) (2.531,9.976)rectangle++(0.04,0.04)
  (2.548,9.974)rectangle++(0.04,0.04) (2.563,9.973)rectangle++(0.04,0.04) (2.575,9.97)rectangle++(0.04,0.04) (2.585,9.968)rectangle++(0.04,0.04) (2.594,9.966)rectangle++(0.04,0.04)
  (2.601,9.963)rectangle++(0.04,0.04) (2.607,9.961)rectangle++(0.04,0.04) (2.161,4.475)rectangle++(0.04,0.04) (-0.519,4.236)rectangle++(0.04,0.04) (-0.274,5.218)rectangle++(0.04,0.04)
  (-0.026,6.042)rectangle++(0.04,0.04) (0.268,2.705)rectangle++(0.04,0.04) (0.333,3.884)rectangle++(0.04,0.04) (0.437,4.885)rectangle++(0.04,0.04) (0.564,5.731)rectangle++(0.04,0.04)
  (-1.072,4.185)rectangle++(0.04,0.04) (-0.746,5.196)rectangle++(0.04,0.04) (-0.428,6.043)rectangle++(0.04,0.04) (-0.124,6.75)rectangle++(0.04,0.04) (-1.411,4.257)rectangle++(0.04,0.04)
  (-1.031,5.271)rectangle++(0.04,0.04) (-0.668,6.118)rectangle++(0.04,0.04) (-0.325,6.823)rectangle++(0.04,0.04) (-0.006,7.409)rectangle++(0.04,0.04) (0.289,7.894)rectangle++(0.04,0.04)
  (0.559,8.294)rectangle++(0.04,0.04) (0.805,8.624)rectangle++(0.04,0.04) (1.027,8.894)rectangle++(0.04,0.04) (1.227,9.115)rectangle++(0.04,0.04) (1.405,9.295)rectangle++(0.04,0.04)
  (1.564,9.441)rectangle++(0.04,0.04) (1.705,9.558)rectangle++(0.04,0.04) (1.829,9.653)rectangle++(0.04,0.04) (1.939,9.728)rectangle++(0.04,0.04) (2.035,9.787)rectangle++(0.04,0.04)
  (2.119,9.834)rectangle++(0.04,0.04) (2.192,9.87)rectangle++(0.04,0.04) (2.256,9.898)rectangle++(0.04,0.04) (2.311,9.92)rectangle++(0.04,0.04) (2.359,9.935)rectangle++(0.04,0.04)
  (2.401,9.947)rectangle++(0.04,0.04) (2.436,9.955)rectangle++(0.04,0.04) (2.467,9.961)rectangle++(0.04,0.04) (2.493,9.964)rectangle++(0.04,0.04) (2.515,9.966)rectangle++(0.04,0.04)
  (0.58,3.176)rectangle++(0.04,0.04) (0.618,4.273)rectangle++(0.04,0.04) (-0.704,3.286)rectangle++(0.04,0.04) (-0.469,4.417)rectangle++(0.04,0.04) (-0.224,5.37)rectangle++(0.04,0.04)
  (-0.037,3.56)rectangle++(0.04,0.04) (0.109,4.624)rectangle++(0.04,0.04) (0.275,5.522)rectangle++(0.04,0.04) (0.453,6.279)rectangle++(0.04,0.04) (1.502,3.276)rectangle++(0.04,0.04)
  (1.405,4.321)rectangle++(0.04,0.04) (-0.622,4.064)rectangle++(0.04,0.04) (-0.368,5.076)rectangle++(0.04,0.04) (-0.112,5.925)rectangle++(0.04,0.04) (0.14,6.637)rectangle++(0.04,0.04)
  (-1.34,4.284)rectangle++(0.04,0.04) (-0.97,5.291)rectangle++(0.04,0.04) (-0.615,6.133)rectangle++(0.04,0.04) (-0.28,6.833)rectangle++(0.04,0.04) (0.033,7.416)rectangle++(0.04,0.04)
  (1.832,3.4)rectangle++(0.04,0.04) (1.691,4.413)rectangle++(0.04,0.04) (1.612,5.279)rectangle++(0.04,0.04) (1.579,6.019)rectangle++(0.04,0.04) (1.581,6.649)rectangle++(0.04,0.04)
  (1.607,7.185)rectangle++(0.04,0.04) (1.652,7.639)rectangle++(0.04,0.04) (1.707,8.023)rectangle++(0.04,0.04) (1.77,8.348)rectangle++(0.04,0.04) (1.836,8.621)rectangle++(0.04,0.04)
  (1.903,8.851)rectangle++(0.04,0.04) (1.97,9.043)rectangle++(0.04,0.04) (2.034,9.204)rectangle++(0.04,0.04) (2.095,9.338)rectangle++(0.04,0.04) (2.152,9.45)rectangle++(0.04,0.04)
  (2.088,4.273)rectangle++(0.04,0.04) (-0.489,4.183)rectangle++(0.04,0.04) (-0.251,5.171)rectangle++(0.04,0.04) (-0.008,6.002)rectangle++(0.04,0.04) (0.231,6.698)rectangle++(0.04,0.04)
  (0.462,7.28)rectangle++(0.04,0.04) (1.749,3.471)rectangle++(0.04,0.04) (-0.38,3.961)rectangle++(0.04,0.04) (-0.167,4.978)rectangle++(0.04,0.04) (0.055,5.834)rectangle++(0.04,0.04)
  (0.278,6.553)rectangle++(0.04,0.04) (0.496,7.155)rectangle++(0.04,0.04) (0.706,7.658)rectangle++(0.04,0.04) (0.904,8.077)rectangle++(0.04,0.04) (1.089,8.426)rectangle++(0.04,0.04)
  (1.261,8.715)rectangle++(0.04,0.04) (1.418,8.953)rectangle++(0.04,0.04) (1.561,9.15)rectangle++(0.04,0.04) (1.691,9.311)rectangle++(0.04,0.04) (2.108,4.141)rectangle++(0.04,0.04)
  (1.955,5.032)rectangle++(0.04,0.04) (1.861,5.795)rectangle++(0.04,0.04) (1.811,6.448)rectangle++(0.04,0.04) (1.795,7.004)rectangle++(0.04,0.04) (1.804,7.478)rectangle++(0.04,0.04)
  (1.83,7.88)rectangle++(0.04,0.04) (1.869,8.221)rectangle++(0.04,0.04) (1.915,8.51)rectangle++(0.04,0.04) (1.966,8.753)rectangle++(0.04,0.04) (2.019,8.957)rectangle++(0.04,0.04)
  (-1.496,5.094)rectangle++(0.04,0.04) (-1.07,5.986)rectangle++(0.04,0.04) (-0.672,6.727)rectangle++(0.04,0.04) (-1.505,4.15)rectangle++(0.04,0.04) (-1.115,5.184)rectangle++(0.04,0.04)
  (-1.256,3.762)rectangle++(0.04,0.04) (-0.92,4.844)rectangle++(0.04,0.04) (-0.59,5.75)rectangle++(0.04,0.04) (-0.274,6.508)rectangle++(0.04,0.04) (0.025,7.139)rectangle++(0.04,0.04)
  (0.305,7.663)rectangle++(0.04,0.04) (0.563,8.098)rectangle++(0.04,0.04) (0.801,8.457)rectangle++(0.04,0.04) (1.017,8.752)rectangle++(0.04,0.04) (1.212,8.995)rectangle++(0.04,0.04)
  (1.388,9.193)rectangle++(0.04,0.04) (1.545,9.355)rectangle++(0.04,0.04) (1.685,9.486)rectangle++(0.04,0.04) (1.81,9.592)rectangle++(0.04,0.04) (2.163,4.223)rectangle++(0.04,0.04)
  (2.005,5.099)rectangle++(0.04,0.04) (1.906,5.85)rectangle++(0.04,0.04) (0.018,2.378)rectangle++(0.04,0.04) (0.108,3.617)rectangle++(0.04,0.04) (-0.705,3.682)rectangle++(0.04,0.04)
  (-0.454,4.754)rectangle++(0.04,0.04) (-0.198,5.655)rectangle++(0.04,0.04) (0.016,2.352)rectangle++(0.04,0.04) (-0.453,3.415)rectangle++(0.04,0.04) (-0.25,4.517)rectangle++(0.04,0.04)
  (-0.034,5.446)rectangle++(0.04,0.04) (0.186,6.226)rectangle++(0.04,0.04) (0.405,6.881)rectangle++(0.04,0.04) (0.618,7.429)rectangle++(0.04,0.04) (0.82,7.886)rectangle++(0.04,0.04)
  (1.857,3.669)rectangle++(0.04,0.04) (1.723,4.641)rectangle++(0.04,0.04) (1.648,5.472)rectangle++(0.04,0.04) (1.617,6.181)rectangle++(0.04,0.04) (0.291,2.321)rectangle++(0.04,0.04)
  (0.338,3.558)rectangle++(0.04,0.04) (0.428,4.607)rectangle++(0.04,0.04) (0.546,5.495)rectangle++(0.04,0.04) (0.681,6.245)rectangle++(0.04,0.04) (0.827,6.877)rectangle++(0.04,0.04)
  (0.976,7.409)rectangle++(0.04,0.04) (1.123,7.855)rectangle++(0.04,0.04) (1.267,8.228)rectangle++(0.04,0.04) (1.404,8.539)rectangle++(0.04,0.04) (1.533,8.798)rectangle++(0.04,0.04)
  (1.652,9.013)rectangle++(0.04,0.04) (1.763,9.191)rectangle++(0.04,0.04) (1.864,9.338)rectangle++(0.04,0.04) (1.956,9.459)rectangle++(0.04,0.04) (2.038,9.558)rectangle++(0.04,0.04)
  (2.113,9.639)rectangle++(0.04,0.04) (2.179,9.705)rectangle++(0.04,0.04) (2.238,9.758)rectangle++(0.04,0.04) (2.291,9.801)rectangle++(0.04,0.04) (2.337,9.836)rectangle++(0.04,0.04)
  (2.378,9.863)rectangle++(0.04,0.04) (2.413,9.885)rectangle++(0.04,0.04) (2.445,9.902)rectangle++(0.04,0.04) (2.472,9.915)rectangle++(0.04,0.04) (2.495,9.925)rectangle++(0.04,0.04)
  (2.516,9.933)rectangle++(0.04,0.04) (2.534,9.938)rectangle++(0.04,0.04) (2.549,9.942)rectangle++(0.04,0.04) (2.562,9.945)rectangle++(0.04,0.04) (2.573,9.947)rectangle++(0.04,0.04)
  (2.583,9.948)rectangle++(0.04,0.04) (2.591,9.949)rectangle++(0.04,0.04) (2.598,9.949)rectangle++(0.04,0.04) (2.604,9.949)rectangle++(0.04,0.04) (2.609,9.949)rectangle++(0.04,0.04)
  (2.614,9.948)rectangle++(0.04,0.04) (2.617,9.948)rectangle++(0.04,0.04) (2.621,9.947)rectangle++(0.04,0.04) (2.623,9.946)rectangle++(0.04,0.04) (2.625,9.946)rectangle++(0.04,0.04)
  (2.627,9.945)rectangle++(0.04,0.04) (2.629,9.944)rectangle++(0.04,0.04) (2.63,9.944)rectangle++(0.04,0.04) (2.631,9.943)rectangle++(0.04,0.04) (2.632,9.943)rectangle++(0.04,0.04)
  (2.633,9.942)rectangle++(0.04,0.04) (2.633,9.942)rectangle++(0.04,0.04) (2.634,9.941)rectangle++(0.04,0.04) (2.153,4.477)rectangle++(0.04,0.04) (2.007,5.315)rectangle++(0.04,0.04)
  (1.917,6.034)rectangle++(0.04,0.04) (1.868,6.648)rectangle++(0.04,0.04) (1.852,7.173)rectangle++(0.04,0.04) (1.559,3.766)rectangle++(0.04,0.04) (1.474,4.735)rectangle++(0.04,0.04)
  (1.44,5.562)rectangle++(0.04,0.04) (1.444,6.266)rectangle++(0.04,0.04) (1.476,6.865)rectangle++(0.04,0.04) (1.527,7.372)rectangle++(0.04,0.04) (1.591,7.801)rectangle++(0.04,0.04)
  (1.746,3.828)rectangle++(0.04,0.04) (1.635,4.78)rectangle++(0.04,0.04) (1.578,5.594)rectangle++(0.04,0.04) (1.563,6.288)rectangle++(0.04,0.04) (1.578,6.878)rectangle++(0.04,0.04)
  (1.614,7.38)rectangle++(0.04,0.04) (-1.233,4.707)rectangle++(0.04,0.04) (-0.862,5.647)rectangle++(0.04,0.04) (-0.509,6.43)rectangle++(0.04,0.04) (-0.178,7.082)rectangle++(0.04,0.04)
  (0.13,7.623)rectangle++(0.04,0.04) (0.413,8.071)rectangle++(0.04,0.04) (0.672,8.44)rectangle++(0.04,0.04) (0.034,2.731)rectangle++(0.04,0.04) (-0.553,3.246)rectangle++(0.04,0.04)
  (-0.342,4.377)rectangle++(0.04,0.04) (-0.118,5.33)rectangle++(0.04,0.04) (1.334,2.963)rectangle++(0.04,0.04) (1.25,4.062)rectangle++(0.04,0.04) (1.223,4.999)rectangle++(0.04,0.04)
  (1.237,5.796)rectangle++(0.04,0.04) (1.281,6.473)rectangle++(0.04,0.04) (1.346,7.047)rectangle++(0.04,0.04) (1.424,7.533)rectangle++(0.04,0.04) (1.509,7.942)rectangle++(0.04,0.04)
  (1.598,8.287)rectangle++(0.04,0.04) (1.688,8.576)rectangle++(0.04,0.04) (1.775,8.818)rectangle++(0.04,0.04) (1.86,9.021)rectangle++(0.04,0.04) (1.939,9.189)rectangle++(0.04,0.04)
  (2.014,9.33)rectangle++(0.04,0.04) (-1.576,5.167)rectangle++(0.04,0.04) (-1.135,6.051)rectangle++(0.04,0.04) (1.631,2.871)rectangle++(0.04,0.04) (0.194,2.22)rectangle++(0.04,0.04)
  (0.252,3.475)rectangle++(0.04,0.04) (0.351,4.54)rectangle++(0.04,0.04) (0.478,5.441)rectangle++(0.04,0.04) (1.292,3.12)rectangle++(0.04,0.04) (1.22,4.196)rectangle++(0.04,0.04)
  (1.203,5.114)rectangle++(0.04,0.04) (1.225,5.895)rectangle++(0.04,0.04) (1.275,6.558)rectangle++(0.04,0.04) (1.344,7.12)rectangle++(0.04,0.04) (1.425,7.594)rectangle++(0.04,0.04)
  (1.513,7.994)rectangle++(0.04,0.04) (1.603,8.331)rectangle++(0.04,0.04) (1.694,8.613)rectangle++(0.04,0.04) (1.782,8.85)rectangle++(0.04,0.04) (1.867,9.047)rectangle++(0.04,0.04)
  (1.946,9.212)rectangle++(0.04,0.04) (2.021,9.348)rectangle++(0.04,0.04) (2.089,9.461)rectangle++(0.04,0.04) (2.152,9.555)rectangle++(0.04,0.04) (2.209,9.632)rectangle++(0.04,0.04)
  (2.261,9.695)rectangle++(0.04,0.04) (2.307,9.746)rectangle++(0.04,0.04) (2.349,9.788)rectangle++(0.04,0.04) (2.386,9.822)rectangle++(0.04,0.04) (2.419,9.85)rectangle++(0.04,0.04)
  (2.448,9.872)rectangle++(0.04,0.04) (2.473,9.889)rectangle++(0.04,0.04) (2.159,4.43)rectangle++(0.04,0.04) (2.01,5.276)rectangle++(0.04,0.04) (-0.716,4.366)rectangle++(0.04,0.04)
  (-0.436,5.336)rectangle++(0.04,0.04) (-0.159,6.149)rectangle++(0.04,0.04) (0.108,6.829)rectangle++(0.04,0.04) (0.299,2.793)rectangle++(0.04,0.04) (0.364,3.958)rectangle++(0.04,0.04)
  (0.005,2.121)rectangle++(0.04,0.04) (0.087,3.399)rectangle++(0.04,0.04) (0.208,4.482)rectangle++(0.04,0.04) (0.354,5.397)rectangle++(0.04,0.04) (0.514,6.17)rectangle++(0.04,0.04)
  (0.682,6.82)rectangle++(0.04,0.04) (0.85,7.366)rectangle++(0.04,0.04) (1.015,7.823)rectangle++(0.04,0.04) (1.174,8.205)rectangle++(0.04,0.04) (1.324,8.524)rectangle++(0.04,0.04)
  (1.956,3.906)rectangle++(0.04,0.04) (1.816,4.838)rectangle++(0.04,0.04) (1.735,5.636)rectangle++(0.04,0.04) (1.698,6.317)rectangle++(0.04,0.04) (1.694,6.898)rectangle++(0.04,0.04)
  (1.714,7.392)rectangle++(0.04,0.04) (1.75,7.811)rectangle++(0.04,0.04) (1.798,8.165)rectangle++(0.04,0.04) (1.853,8.465)rectangle++(0.04,0.04) (1.911,8.717)rectangle++(0.04,0.04)
  (1.971,8.929)rectangle++(0.04,0.04) (2.03,9.107)rectangle++(0.04,0.04) (2.088,9.256)rectangle++(0.04,0.04) (2.143,9.38)rectangle++(0.04,0.04) (2.194,9.484)rectangle++(0.04,0.04)
  (2.242,9.57)rectangle++(0.04,0.04) (2.107,4.316)rectangle++(0.04,0.04) (1.961,5.181)rectangle++(0.04,0.04) (1.872,5.921)rectangle++(0.04,0.04) (-0.878,4.468)rectangle++(0.04,0.04)
  (-0.57,5.429)rectangle++(0.04,0.04) (-0.269,6.233)rectangle++(0.04,0.04) (0.018,6.905)rectangle++(0.04,0.04) (0.29,7.465)rectangle++(0.04,0.04) (0.543,7.93)rectangle++(0.04,0.04)
  (0.776,8.315)rectangle++(0.04,0.04) (0.99,8.633)rectangle++(0.04,0.04) (1.185,8.895)rectangle++(0.04,0.04) (1.361,9.109)rectangle++(0.04,0.04) (1.519,9.285)rectangle++(0.04,0.04)
  (1.66,9.427)rectangle++(0.04,0.04) (1.786,9.543)rectangle++(0.04,0.04) (1.898,9.636)rectangle++(0.04,0.04) (1.996,9.711)rectangle++(0.04,0.04) (2.083,9.771)rectangle++(0.04,0.04)
  (2.159,9.818)rectangle++(0.04,0.04) (2.226,9.855)rectangle++(0.04,0.04) (2.284,9.884)rectangle++(0.04,0.04) (2.335,9.906)rectangle++(0.04,0.04) (2.378,9.923)rectangle++(0.04,0.04)
  (2.416,9.936)rectangle++(0.04,0.04) (2.449,9.945)rectangle++(0.04,0.04) (-1.628,5.37)rectangle++(0.04,0.04) (-1.171,6.226)rectangle++(0.04,0.04) (-0.749,6.935)rectangle++(0.04,0.04)
  (-0.361,7.521)rectangle++(0.04,0.04) (-0.008,8.003)rectangle++(0.04,0.04) (0.311,8.399)rectangle++(0.04,0.04) (0.598,8.723)rectangle++(0.04,0.04) (0.855,8.987)rectangle++(0.04,0.04)
  (2.125,3.889)rectangle++(0.04,0.04) (1.96,4.817)rectangle++(0.04,0.04) (1.856,5.612)rectangle++(0.04,0.04) (1.8,6.292)rectangle++(0.04,0.04) (1.78,6.873)rectangle++(0.04,0.04)
  (1.493,3.691)rectangle++(0.04,0.04) (1.415,4.674)rectangle++(0.04,0.04) (1.387,5.513)rectangle++(0.04,0.04) (1.398,6.226)rectangle++(0.04,0.04) (1.435,6.833)rectangle++(0.04,0.04)
  (1.491,7.347)rectangle++(0.04,0.04) (1.559,7.781)rectangle++(0.04,0.04) (1.634,8.148)rectangle++(0.04,0.04) (1.713,8.457)rectangle++(0.04,0.04) (1.792,8.716)rectangle++(0.04,0.04)
  (1.869,8.933)rectangle++(0.04,0.04) (-1.517,5.062)rectangle++(0.04,0.04) (-1.089,5.959)rectangle++(0.04,0.04) (-0.69,6.705)rectangle++(0.04,0.04) (-0.32,7.323)rectangle++(0.04,0.04)
  (0.278,2.841)rectangle++(0.04,0.04) (0.347,4)rectangle++(0.04,0.04) (0.453,4.982)rectangle++(0.04,0.04) (0.582,5.813)rectangle++(0.04,0.04) (0.725,6.514)rectangle++(0.04,0.04)
  (0.875,7.104)rectangle++(0.04,0.04) (1.025,7.599)rectangle++(0.04,0.04) (1.173,8.015)rectangle++(0.04,0.04) (1.316,8.362)rectangle++(0.04,0.04) (1.451,8.651)rectangle++(0.04,0.04)
  (1.577,8.892)rectangle++(0.04,0.04) (1.694,9.091)rectangle++(0.04,0.04) (1.801,9.256)rectangle++(0.04,0.04) (1.899,9.392)rectangle++(0.04,0.04) (1.988,9.503)rectangle++(0.04,0.04)
  (2.067,9.594)rectangle++(0.04,0.04) (2.139,9.669)rectangle++(0.04,0.04) (2.144,4.312)rectangle++(0.04,0.04) (1.992,5.176)rectangle++(0.04,0.04) (1.898,5.916)rectangle++(0.04,0.04)
  (1.848,6.549)rectangle++(0.04,0.04) (1.831,7.089)rectangle++(0.04,0.04) (1.837,7.548)rectangle++(0.04,0.04) (-1.229,4.782)rectangle++(0.04,0.04) (-0.856,5.71)rectangle++(0.04,0.04)
  (-0.501,6.484)rectangle++(0.04,0.04) (1.664,3.099)rectangle++(0.04,0.04) (1.536,4.164)rectangle++(0.04,0.04) (1.47,5.074)rectangle++(0.04,0.04) (1.451,5.85)rectangle++(0.04,0.04)
  (1.279,3.42)rectangle++(0.04,0.04) (1.222,4.452)rectangle++(0.04,0.04) (1.214,5.331)rectangle++(0.04,0.04) (1.243,6.079)rectangle++(0.04,0.04) (1.298,6.714)rectangle++(0.04,0.04)
  (-1.149,4.516)rectangle++(0.04,0.04) (-0.798,5.481)rectangle++(0.04,0.04) (-0.461,6.287)rectangle++(0.04,0.04) (-0.143,6.959)rectangle++(0.04,0.04) (0.155,7.517)rectangle++(0.04,0.04)
  (0.43,7.979)rectangle++(0.04,0.04) (0.682,8.361)rectangle++(0.04,0.04) (0.912,8.676)rectangle++(0.04,0.04) (1.12,8.934)rectangle++(0.04,0.04) (1.307,9.146)rectangle++(0.04,0.04)
  (1.475,9.318)rectangle++(0.04,0.04) (1.624,9.457)rectangle++(0.04,0.04) (1.757,9.57)rectangle++(0.04,0.04) (1.874,9.66)rectangle++(0.04,0.04) (1.977,9.733)rectangle++(0.04,0.04)
  (2.067,9.79)rectangle++(0.04,0.04) (2.147,9.835)rectangle++(0.04,0.04) (2.216,9.87)rectangle++(0.04,0.04) (-1.654,5.312)rectangle++(0.04,0.04) (-1.196,6.177)rectangle++(0.04,0.04)
  (-0.771,6.895)rectangle++(0.04,0.04) (1.798,3.117)rectangle++(0.04,0.04) (0.561,2.971)rectangle++(0.04,0.04) (0.594,4.099)rectangle++(0.04,0.04) (0.667,5.057)rectangle++(0.04,0.04)
  (0.767,5.868)rectangle++(0.04,0.04) (0.884,6.553)rectangle++(0.04,0.04) (1.011,7.131)rectangle++(0.04,0.04) (1.143,7.617)rectangle++(0.04,0.04) (1.274,8.025)rectangle++(0.04,0.04)
  (1.402,8.367)rectangle++(0.04,0.04) (1.524,8.652)rectangle++(0.04,0.04) (1.639,8.889)rectangle++(0.04,0.04) (1.747,9.086)rectangle++(0.04,0.04) (-1.572,5.069)rectangle++(0.04,0.04)
  (-1.135,5.968)rectangle++(0.04,0.04) (-0.729,6.714)rectangle++(0.04,0.04) (-0.353,7.332)rectangle++(0.04,0.04) (-0.009,7.843)rectangle++(0.04,0.04) (0.304,8.263)rectangle++(0.04,0.04)
  (0.587,8.608)rectangle++(0.04,0.04) (0.841,8.889)rectangle++(0.04,0.04) (1.068,9.118)rectangle++(0.04,0.04) (1.27,9.304)rectangle++(0.04,0.04) (1.45,9.454)rectangle++(0.04,0.04)
  (1.608,9.574)rectangle++(0.04,0.04) (-1.7,5.139)rectangle++(0.04,0.04) (-1.242,6.032)rectangle++(0.04,0.04) (-0.816,6.773)rectangle++(0.04,0.04) (-0.425,7.386)rectangle++(0.04,0.04)
  (-0.068,7.892)rectangle++(0.04,0.04) (0.256,8.307)rectangle++(0.04,0.04) (0.547,8.647)rectangle++(0.04,0.04) (0.809,8.924)rectangle++(0.04,0.04) (1.042,9.149)rectangle++(0.04,0.04)
  (1.25,9.331)rectangle++(0.04,0.04) (1.433,9.478)rectangle++(0.04,0.04) (1.595,9.595)rectangle++(0.04,0.04) (1.738,9.688)rectangle++(0.04,0.04) (1.862,9.762)rectangle++(0.04,0.04)
  (1.971,9.819)rectangle++(0.04,0.04) (2.066,9.864)rectangle++(0.04,0.04) (2.149,9.898)rectangle++(0.04,0.04) (2.22,9.923)rectangle++(0.04,0.04) (2.282,9.942)rectangle++(0.04,0.04)
  (2.335,9.956)rectangle++(0.04,0.04) (2.381,9.965)rectangle++(0.04,0.04) (2.189,4.424)rectangle++(0.04,0.04) (2.035,5.269)rectangle++(0.04,0.04) (1.939,5.994)rectangle++(0.04,0.04)
  (1.885,6.613)rectangle++(0.04,0.04) (1.865,7.142)rectangle++(0.04,0.04) (1.869,7.592)rectangle++(0.04,0.04) (1.89,7.975)rectangle++(0.04,0.04) (1.923,8.299)rectangle++(0.04,0.04)
  (1.964,8.574)rectangle++(0.04,0.04) (2.011,8.805)rectangle++(0.04,0.04) (2.059,9)rectangle++(0.04,0.04) (2.108,9.164)rectangle++(0.04,0.04) (2.156,9.301)rectangle++(0.04,0.04)
  (2.203,9.416)rectangle++(0.04,0.04) (2.247,9.512)rectangle++(0.04,0.04) (2.288,9.591)rectangle++(0.04,0.04) (2.326,9.657)rectangle++(0.04,0.04) (2.361,9.712)rectangle++(0.04,0.04)
  (2.393,9.757)rectangle++(0.04,0.04) (2.422,9.794)rectangle++(0.04,0.04) (2.141,4.401)rectangle++(0.04,0.04) (1.994,5.251)rectangle++(0.04,0.04) (1.903,5.98)rectangle++(0.04,0.04)
  (1.854,6.603)rectangle++(0.04,0.04) (1.838,7.134)rectangle++(0.04,0.04) (1.064,2.31)rectangle++(0.04,0.04) (0.994,3.518)rectangle++(0.04,0.04) (0.984,4.546)rectangle++(0.04,0.04)
  (1.016,5.421)rectangle++(0.04,0.04) (1.078,6.164)rectangle++(0.04,0.04) (1.161,6.792)rectangle++(0.04,0.04) (1.256,7.323)rectangle++(0.04,0.04) (1.358,7.771)rectangle++(0.04,0.04)
  (1.463,8.147)rectangle++(0.04,0.04) (1.567,8.463)rectangle++(0.04,0.04) (1.669,8.727)rectangle++(0.04,0.04) (1.765,8.947)rectangle++(0.04,0.04) (1.856,9.131)rectangle++(0.04,0.04)
  (1.941,9.283)rectangle++(0.04,0.04) (2.019,9.409)rectangle++(0.04,0.04) (2.09,9.513)rectangle++(0.04,0.04) (2.155,9.599)rectangle++(0.04,0.04) (2.213,9.669)rectangle++(0.04,0.04)
  (2.266,9.726)rectangle++(0.04,0.04) (2.313,9.773)rectangle++(0.04,0.04) (2.355,9.811)rectangle++(0.04,0.04) (2.392,9.841)rectangle++(0.04,0.04) (2.424,9.865)rectangle++(0.04,0.04)
  (2.453,9.885)rectangle++(0.04,0.04) (-1.615,5.359)rectangle++(0.04,0.04) (-1.16,6.216)rectangle++(0.04,0.04) (-0.74,6.926)rectangle++(0.04,0.04) (-0.354,7.513)rectangle++(0.04,0.04)
  (-0.003,7.996)rectangle++(0.04,0.04) (0.315,8.393)rectangle++(0.04,0.04) (0.602,8.718)rectangle++(0.04,0.04) (0.858,8.982)rectangle++(0.04,0.04) (1.086,9.197)rectangle++(0.04,0.04)
  (1.289,9.37)rectangle++(0.04,0.04) (1.468,9.509)rectangle++(0.04,0.04) (1.626,9.62)rectangle++(0.04,0.04) (1.765,9.708)rectangle++(0.04,0.04) (1.886,9.778)rectangle++(0.04,0.04)
  (2.2,4.276)rectangle++(0.04,0.04) (0.801,3.286)rectangle++(0.04,0.04) (0.81,4.357)rectangle++(0.04,0.04) (0.984,2.986)rectangle++(0.04,0.04) (-0.414,3.722)rectangle++(0.04,0.04)
  (-0.205,4.776)rectangle++(0.04,0.04) (0.014,5.664)rectangle++(0.04,0.04) (0.237,6.41)rectangle++(0.04,0.04) (0.455,7.035)rectangle++(0.04,0.04) (0.666,7.558)rectangle++(0.04,0.04)
  (0.866,7.994)rectangle++(0.04,0.04) (1.054,8.356)rectangle++(0.04,0.04) (1.228,8.657)rectangle++(0.04,0.04) (1.388,8.906)rectangle++(0.04,0.04) (1.534,9.11)rectangle++(0.04,0.04)
  (1.666,9.279)rectangle++(0.04,0.04) (1.785,9.416)rectangle++(0.04,0.04) (1.892,9.529)rectangle++(0.04,0.04) (1.987,9.62)rectangle++(0.04,0.04) (2.071,9.694)rectangle++(0.04,0.04)
  (2.146,9.753)rectangle++(0.04,0.04) (2.212,9.8)rectangle++(0.04,0.04) (2.27,9.838)rectangle++(0.04,0.04) (2.321,9.868)rectangle++(0.04,0.04) (2.365,9.891)rectangle++(0.04,0.04)
  (2.404,9.909)rectangle++(0.04,0.04) (2.438,9.923)rectangle++(0.04,0.04) (2.467,9.933)rectangle++(0.04,0.04) (0.575,3.172)rectangle++(0.04,0.04) (0.614,4.269)rectangle++(0.04,0.04)
  (0.69,5.201)rectangle++(0.04,0.04) (0.793,5.989)rectangle++(0.04,0.04) (0.911,6.655)rectangle++(0.04,0.04) (1.038,7.217)rectangle++(0.04,0.04) (1.169,7.689)rectangle++(0.04,0.04)
  (1.299,8.085)rectangle++(0.04,0.04) (-1.439,4.787)rectangle++(0.04,0.04) (-1.034,5.723)rectangle++(0.04,0.04) (-0.652,6.502)rectangle++(0.04,0.04) (-0.297,7.149)rectangle++(0.04,0.04)
  (0.032,7.685)rectangle++(0.04,0.04) (0.332,8.127)rectangle++(0.04,0.04) (1.974,3.605)rectangle++(0.04,0.04) (1.82,4.581)rectangle++(0.04,0.04) (1.728,5.417)rectangle++(0.04,0.04)
  (1.683,6.132)rectangle++(0.04,0.04) (0.019,2.417)rectangle++(0.04,0.04) (0.111,3.65)rectangle++(0.04,0.04) (0.238,4.694)rectangle++(0.04,0.04) (0.388,5.577)rectangle++(0.04,0.04)
  (0.551,6.321)rectangle++(0.04,0.04) (0.719,6.947)rectangle++(0.04,0.04) (0.887,7.472)rectangle++(0.04,0.04) (1.05,7.912)rectangle++(0.04,0.04) (1.207,8.279)rectangle++(0.04,0.04)
  (1.355,8.585)rectangle++(0.04,0.04) (1.493,8.84)rectangle++(0.04,0.04) (1.62,9.05)rectangle++(0.04,0.04) (1.737,9.224)rectangle++(0.04,0.04) (1.843,9.367)rectangle++(0.04,0.04)
  (1.939,9.485)rectangle++(0.04,0.04) (2.026,9.581)rectangle++(0.04,0.04) (2.103,9.659)rectangle++(0.04,0.04) (2.171,9.722)rectangle++(0.04,0.04) (2.232,9.773)rectangle++(0.04,0.04)
  (2.159,4.353)rectangle++(0.04,0.04) (2.007,5.21)rectangle++(0.04,0.04) (1.912,5.945)rectangle++(0.04,0.04) (1.861,6.573)rectangle++(0.04,0.04) (1.842,7.109)rectangle++(0.04,0.04)
  (1.848,7.565)rectangle++(0.04,0.04) (1.871,7.952)rectangle++(0.04,0.04) (1.906,8.281)rectangle++(0.04,0.04) (0.486,2.982)rectangle++(0.04,0.04) (0.53,4.111)rectangle++(0.04,0.04)
  (0.613,5.069)rectangle++(0.04,0.04) (0.721,5.881)rectangle++(0.04,0.04) (0.846,6.566)rectangle++(0.04,0.04) (0.98,7.143)rectangle++(0.04,0.04) (1.116,7.629)rectangle++(0.04,0.04)
  (1.252,8.036)rectangle++(0.04,0.04) (-1.438,4.769)rectangle++(0.04,0.04) (-1.226,3.619)rectangle++(0.04,0.04) (-0.944,3.432)rectangle++(0.04,0.04) (-0.668,4.551)rectangle++(0.04,0.04)
  (-0.388,5.492)rectangle++(0.04,0.04) (-0.112,6.28)rectangle++(0.04,0.04) (0.154,6.938)rectangle++(0.04,0.04) (0.406,7.488)rectangle++(0.04,0.04) (0.642,7.944)rectangle++(0.04,0.04)
  (-1.527,4.636)rectangle++(0.04,0.04) (-1.115,5.598)rectangle++(0.04,0.04) (-0.726,6.399)rectangle++(0.04,0.04) (-0.363,7.065)rectangle++(0.04,0.04) (1.792,3.242)rectangle++(0.04,0.04)
  (1.651,4.28)rectangle++(0.04,0.04) (1.572,5.168)rectangle++(0.04,0.04) (1.541,5.926)rectangle++(0.04,0.04) (0.018,2.389)rectangle++(0.04,0.04) (0.109,3.626)rectangle++(0.04,0.04)
  (0.235,4.674)rectangle++(0.04,0.04) (0.385,5.56)rectangle++(0.04,0.04) (0.547,6.306)rectangle++(0.04,0.04) (0.715,6.935)rectangle++(0.04,0.04) (0.883,7.462)rectangle++(0.04,0.04)
  (1.047,7.904)rectangle++(0.04,0.04) (1.204,8.272)rectangle++(0.04,0.04) (1.352,8.58)rectangle++(0.04,0.04) (1.966,3.923)rectangle++(0.04,0.04) (0.741,3.165)rectangle++(0.04,0.04)
  (0.755,4.257)rectangle++(0.04,0.04) (0.809,5.184)rectangle++(0.04,0.04) (0.893,5.971)rectangle++(0.04,0.04) (0.996,6.635)rectangle++(0.04,0.04) (1.11,7.197)rectangle++(0.04,0.04)
  (1.229,7.669)rectangle++(0.04,0.04) (1.349,8.066)rectangle++(0.04,0.04) (1.467,8.398)rectangle++(0.04,0.04) (1.581,8.676)rectangle++(0.04,0.04) (1.689,8.907)rectangle++(0.04,0.04)
  (1.789,9.1)rectangle++(0.04,0.04) (1.883,9.26)rectangle++(0.04,0.04) (0.672,2.012)rectangle++(0.04,0.04) (0.65,3.28)rectangle++(0.04,0.04) (0.681,4.358)rectangle++(0.04,0.04)
  (1,2.957)rectangle++(0.04,0.04) (0.966,4.07)rectangle++(0.04,0.04) (0.981,5.017)rectangle++(0.04,0.04) (1.033,5.821)rectangle++(0.04,0.04) (1.108,6.503)rectangle++(0.04,0.04)
  (1.2,7.079)rectangle++(0.04,0.04) (1.301,7.566)rectangle++(0.04,0.04) (1.406,7.975)rectangle++(0.04,0.04) (1.512,8.319)rectangle++(0.04,0.04) (-1.451,4.873)rectangle++(0.04,0.04)
  (-1.041,5.797)rectangle++(0.04,0.04) (-0.655,6.565)rectangle++(0.04,0.04) (-0.296,7.203)rectangle++(0.04,0.04) (0.034,7.73)rectangle++(0.04,0.04) (0.336,8.166)rectangle++(0.04,0.04)
  (0.61,8.523)rectangle++(0.04,0.04) (0.857,8.817)rectangle++(0.04,0.04) (1.079,9.056)rectangle++(0.04,0.04) (1.277,9.251)rectangle++(0.04,0.04) (1.454,9.408)rectangle++(0.04,0.04)
  (1.61,9.535)rectangle++(0.04,0.04) (1.747,9.637)rectangle++(0.04,0.04) (2.182,4.217)rectangle++(0.04,0.04) (2.021,5.093)rectangle++(0.04,0.04) (1.919,5.845)rectangle++(0.04,0.04)
  (1.863,6.487)rectangle++(0.04,0.04) (1.841,7.036)rectangle++(0.04,0.04) (1.844,7.503)rectangle++(0.04,0.04) (1.865,7.9)rectangle++(0.04,0.04) (1.899,8.237)rectangle++(0.04,0.04)
  (1.942,8.521)rectangle++(0.04,0.04) (1.989,8.762)rectangle++(0.04,0.04) (2.039,8.964)rectangle++(0.04,0.04) (2.089,9.134)rectangle++(0.04,0.04) (2.139,9.277)rectangle++(0.04,0.04)
  (2.187,9.396)rectangle++(0.04,0.04) (2.233,9.495)rectangle++(0.04,0.04) (1.277,2.307)rectangle++(0.04,0.04) (1.176,3.506)rectangle++(0.04,0.04) (-0.49,3.86)rectangle++(0.04,0.04)
  (-0.264,4.897)rectangle++(0.04,0.04) (-0.031,5.769)rectangle++(0.04,0.04) (0.202,6.501)rectangle++(0.04,0.04) (0.43,7.114)rectangle++(0.04,0.04) (0.648,7.626)rectangle++(0.04,0.04)
  (1.813,3.58)rectangle++(0.04,0.04) (0.674,3.064)rectangle++(0.04,0.04) (0.694,4.174)rectangle++(0.04,0.04) (0.754,5.116)rectangle++(0.04,0.04) (-0.907,4.099)rectangle++(0.04,0.04)
  (-0.989,3.587)rectangle++(0.04,0.04) (-0.699,4.685)rectangle++(0.04,0.04) (-0.409,5.606)rectangle++(0.04,0.04) (-1.22,3.978)rectangle++(0.04,0.04) (-0.88,5.026)rectangle++(0.04,0.04)
  (-0.549,5.903)rectangle++(0.04,0.04) (-0.233,6.636)rectangle++(0.04,0.04) (0.065,7.246)rectangle++(0.04,0.04) (0.343,7.753)rectangle++(0.04,0.04) (0.6,8.172)rectangle++(0.04,0.04)
  (0.834,8.519)rectangle++(0.04,0.04) (1.048,8.804)rectangle++(0.04,0.04) (1.241,9.037)rectangle++(0.04,0.04) (1.414,9.228)rectangle++(0.04,0.04) (2.12,4.063)rectangle++(0.04,0.04)
  (1.962,4.965)rectangle++(0.04,0.04) (1.864,5.738)rectangle++(0.04,0.04) (1.812,6.399)rectangle++(0.04,0.04) (1.794,6.963)rectangle++(0.04,0.04) (1.801,7.443)rectangle++(0.04,0.04)
  (1.826,7.85)rectangle++(0.04,0.04) (1.864,8.196)rectangle++(0.04,0.04) (1.91,8.488)rectangle++(0.04,0.04) (1.961,8.735)rectangle++(0.04,0.04) (2.014,8.942)rectangle++(0.04,0.04)
  (2.067,9.117)rectangle++(0.04,0.04) (2.12,9.263)rectangle++(0.04,0.04) (2.17,9.385)rectangle++(0.04,0.04) (-1.559,5.207)rectangle++(0.04,0.04) (-1.119,6.085)rectangle++(0.04,0.04)
  (-0.71,6.813)rectangle++(0.04,0.04) (-0.333,7.416)rectangle++(0.04,0.04) (0.011,7.913)rectangle++(0.04,0.04) (-1.036,3.216)rectangle++(0.04,0.04) (-0.754,4.371)rectangle++(0.04,0.04);
\fill[piece2]
  (-0.468,5.342)rectangle++(0.04,0.04) (-0.187,6.156)rectangle++(0.04,0.04) (0.085,6.836)rectangle++(0.04,0.04) (0.344,7.403)rectangle++(0.04,0.04) (0.586,7.875)rectangle++(0.04,0.04)
  (0.811,8.267)rectangle++(0.04,0.04) (1.018,8.591)rectangle++(0.04,0.04) (2.008,3.85)rectangle++(0.04,0.04) (1.859,4.788)rectangle++(0.04,0.04) (1.769,5.592)rectangle++(0.04,0.04)
  (1.725,6.278)rectangle++(0.04,0.04) (1.715,6.864)rectangle++(0.04,0.04) (1.73,7.362)rectangle++(0.04,0.04) (1.763,7.785)rectangle++(0.04,0.04) (-0.579,3.4)rectangle++(0.04,0.04)
  (-0.359,4.509)rectangle++(0.04,0.04) (-0.127,5.444)rectangle++(0.04,0.04) (-0.045,3.57)rectangle++(0.04,0.04) (0.103,4.633)rectangle++(0.04,0.04) (0.27,5.53)rectangle++(0.04,0.04)
  (0.449,6.286)rectangle++(0.04,0.04) (0.631,6.921)rectangle++(0.04,0.04) (0.811,7.454)rectangle++(0.04,0.04) (0.985,7.9)rectangle++(0.04,0.04) (1.151,8.271)rectangle++(0.04,0.04)
  (1.307,8.581)rectangle++(0.04,0.04) (1.452,8.838)rectangle++(0.04,0.04) (1.586,9.05)rectangle++(0.04,0.04) (1.708,9.225)rectangle++(0.04,0.04) (1.818,9.369)rectangle++(0.04,0.04)
  (1.918,9.488)rectangle++(0.04,0.04) (2.008,9.584)rectangle++(0.04,0.04) (2.088,9.662)rectangle++(0.04,0.04) (-1.633,5.246)rectangle++(0.04,0.04) (1.49,2.603)rectangle++(0.04,0.04)
  (1.369,3.749)rectangle++(0.04,0.04) (1.311,4.728)rectangle++(0.04,0.04) (1.302,5.563)rectangle++(0.04,0.04) (1.327,6.273)rectangle++(0.04,0.04) (1.376,6.875)rectangle++(0.04,0.04)
  (1.443,7.385)rectangle++(0.04,0.04) (1.519,7.816)rectangle++(0.04,0.04) (1.602,8.179)rectangle++(0.04,0.04) (1.687,8.484)rectangle++(0.04,0.04) (1.771,8.74)rectangle++(0.04,0.04)
  (1.853,8.954)rectangle++(0.04,0.04) (1.931,9.133)rectangle++(0.04,0.04) (2.004,9.282)rectangle++(0.04,0.04) (2.073,9.406)rectangle++(0.04,0.04) (2.136,9.508)rectangle++(0.04,0.04)
  (-0.633,3.433)rectangle++(0.04,0.04) (-0.403,4.539)rectangle++(0.04,0.04) (-0.163,5.471)rectangle++(0.04,0.04) (0.078,6.253)rectangle++(0.04,0.04) (0.314,6.908)rectangle++(0.04,0.04)
  (0.541,7.455)rectangle++(0.04,0.04) (1.783,3.523)rectangle++(0.04,0.04) (1.654,4.519)rectangle++(0.04,0.04) (1.585,5.371)rectangle++(0.04,0.04) (1.559,6.098)rectangle++(0.04,0.04)
  (1.567,6.718)rectangle++(0.04,0.04) (1.599,7.243)rectangle++(0.04,0.04) (1.646,7.689)rectangle++(0.04,0.04) (1.705,8.066)rectangle++(0.04,0.04) (1.77,8.384)rectangle++(0.04,0.04)
  (1.837,8.652)rectangle++(0.04,0.04) (-0.228,3.377)rectangle++(0.04,0.04) (-0.061,4.476)rectangle++(0.04,0.04) (0.125,5.403)rectangle++(0.04,0.04) (0.32,6.184)rectangle++(0.04,0.04)
  (0.517,6.839)rectangle++(0.04,0.04) (0.711,7.389)rectangle++(0.04,0.04) (0.898,7.848)rectangle++(0.04,0.04) (1.075,8.231)rectangle++(0.04,0.04) (1.912,3.793)rectangle++(0.04,0.04)
  (1.775,4.744)rectangle++(0.04,0.04) (1.696,5.558)rectangle++(0.04,0.04) (1.662,6.252)rectangle++(0.04,0.04) (1.661,6.844)rectangle++(0.04,0.04) (1.683,7.347)rectangle++(0.04,0.04)
  (1.722,7.774)rectangle++(0.04,0.04) (1.773,8.135)rectangle++(0.04,0.04) (1.83,8.44)rectangle++(0.04,0.04) (1.891,8.697)rectangle++(0.04,0.04) (1.953,8.913)rectangle++(0.04,0.04)
  (2.014,9.094)rectangle++(0.04,0.04) (2.074,9.246)rectangle++(0.04,0.04) (-1.547,5.162)rectangle++(0.04,0.04) (-1.111,6.045)rectangle++(0.04,0.04) (1.627,2.876)rectangle++(0.04,0.04)
  (1.496,3.976)rectangle++(0.04,0.04) (-0.532,4.013)rectangle++(0.04,0.04) (-0.294,5.029)rectangle++(0.04,0.04) (0.573,2.096)rectangle++(0.04,0.04) (0.452,2.496)rectangle++(0.04,0.04)
  (0.482,3.7)rectangle++(0.04,0.04) (0.555,4.722)rectangle++(0.04,0.04) (0.659,5.588)rectangle++(0.04,0.04) (0.781,6.319)rectangle++(0.04,0.04) (0.915,6.936)rectangle++(0.04,0.04)
  (1.053,7.456)rectangle++(0.04,0.04) (1.191,7.891)rectangle++(0.04,0.04) (1.326,8.256)rectangle++(0.04,0.04) (1.455,8.561)rectangle++(0.04,0.04) (1.95,3.943)rectangle++(0.04,0.04)
  (1.813,4.87)rectangle++(0.04,0.04) (1.733,5.663)rectangle++(0.04,0.04) (1.698,6.341)rectangle++(0.04,0.04) (1.694,6.918)rectangle++(0.04,0.04) (1.715,7.409)rectangle++(0.04,0.04)
  (1.752,7.825)rectangle++(0.04,0.04) (1.8,8.177)rectangle++(0.04,0.04) (1.855,8.475)rectangle++(0.04,0.04) (1.913,8.726)rectangle++(0.04,0.04) (1.973,8.936)rectangle++(0.04,0.04)
  (2.032,9.113)rectangle++(0.04,0.04) (2.09,9.261)rectangle++(0.04,0.04) (2.145,9.385)rectangle++(0.04,0.04) (2.196,9.487)rectangle++(0.04,0.04) (2.092,4.29)rectangle++(0.04,0.04)
  (1.948,5.159)rectangle++(0.04,0.04) (1.86,5.904)rectangle++(0.04,0.04) (0.96,2.333)rectangle++(0.04,0.04) (0.907,3.541)rectangle++(0.04,0.04) (0.91,4.57)rectangle++(0.04,0.04)
  (1.025,3.05)rectangle++(0.04,0.04) (0.635,2.783)rectangle++(0.04,0.04) (-0.433,3.616)rectangle++(0.04,0.04) (-0.226,4.687)rectangle++(0.04,0.04) (-0.007,5.589)rectangle++(0.04,0.04)
  (0.216,6.348)rectangle++(0.04,0.04) (-1.265,4.241)rectangle++(0.04,0.04) (-0.908,5.252)rectangle++(0.04,0.04) (-0.564,6.097)rectangle++(0.04,0.04) (-0.238,6.801)rectangle++(0.04,0.04)
  (0.068,7.386)rectangle++(0.04,0.04) (0.351,7.872)rectangle++(0.04,0.04) (-1.564,4.567)rectangle++(0.04,0.04) (-1.149,5.541)rectangle++(0.04,0.04) (-0.757,6.352)rectangle++(0.04,0.04)
  (-0.392,7.026)rectangle++(0.04,0.04) (-0.054,7.584)rectangle++(0.04,0.04) (0.255,8.045)rectangle++(0.04,0.04) (0.536,8.424)rectangle++(0.04,0.04) (0.791,8.735)rectangle++(0.04,0.04)
  (1.019,8.989)rectangle++(0.04,0.04) (1.224,9.196)rectangle++(0.04,0.04) (1.406,9.364)rectangle++(0.04,0.04) (1.567,9.5)rectangle++(0.04,0.04) (1.71,9.608)rectangle++(0.04,0.04)
  (1.836,9.695)rectangle++(0.04,0.04) (1.946,9.763)rectangle++(0.04,0.04) (2.042,9.817)rectangle++(0.04,0.04) (2.126,9.859)rectangle++(0.04,0.04) (2.2,9.891)rectangle++(0.04,0.04)
  (2.263,9.916)rectangle++(0.04,0.04) (2.318,9.934)rectangle++(0.04,0.04) (2.366,9.948)rectangle++(0.04,0.04) (2.406,9.957)rectangle++(0.04,0.04) (2.442,9.963)rectangle++(0.04,0.04)
  (2.472,9.967)rectangle++(0.04,0.04) (2.497,9.97)rectangle++(0.04,0.04) (2.519,9.97)rectangle++(0.04,0.04) (2.538,9.97)rectangle++(0.04,0.04) (2.172,4.461)rectangle++(0.04,0.04)
  (2.022,5.301)rectangle++(0.04,0.04) (1.929,6.021)rectangle++(0.04,0.04) (-0.889,4.498)rectangle++(0.04,0.04) (-0.578,5.455)rectangle++(0.04,0.04) (-0.275,6.256)rectangle++(0.04,0.04)
  (-1.333,4.131)rectangle++(0.04,0.04) (-0.97,5.161)rectangle++(0.04,0.04) (-0.621,6.022)rectangle++(0.04,0.04) (-0.289,6.74)rectangle++(0.04,0.04) (-1.439,4.224)rectangle++(0.04,0.04)
  (-1.056,5.244)rectangle++(0.04,0.04) (-0.69,6.096)rectangle++(0.04,0.04) (-0.345,6.805)rectangle++(0.04,0.04) (-0.023,7.395)rectangle++(0.04,0.04) (0.274,7.882)rectangle++(0.04,0.04)
  (0.546,8.285)rectangle++(0.04,0.04) (0.793,8.617)rectangle++(0.04,0.04) (1.017,8.889)rectangle++(0.04,0.04) (2.082,3.907)rectangle++(0.04,0.04) (1.924,4.834)rectangle++(0.04,0.04)
  (1.826,5.628)rectangle++(0.04,0.04) (1.775,6.307)rectangle++(0.04,0.04) (1.759,6.886)rectangle++(0.04,0.04) (1.768,7.379)rectangle++(0.04,0.04) (1.796,7.798)rectangle++(0.04,0.04)
  (1.836,8.152)rectangle++(0.04,0.04) (1.885,8.452)rectangle++(0.04,0.04) (1.938,8.705)rectangle++(0.04,0.04) (1.993,8.918)rectangle++(0.04,0.04) (2.049,9.097)rectangle++(0.04,0.04)
  (2.103,9.247)rectangle++(0.04,0.04) (2.155,9.372)rectangle++(0.04,0.04) (2.205,9.476)rectangle++(0.04,0.04) (2.251,9.563)rectangle++(0.04,0.04) (2.294,9.634)rectangle++(0.04,0.04)
  (2.333,9.694)rectangle++(0.04,0.04) (-1.596,5.299)rectangle++(0.04,0.04) (-1.367,3.694)rectangle++(0.04,0.04) (-1.016,4.791)rectangle++(0.04,0.04) (-0.674,5.709)rectangle++(0.04,0.04)
  (-0.347,6.476)rectangle++(0.04,0.04) (-0.038,7.115)rectangle++(0.04,0.04) (0.25,7.645)rectangle++(0.04,0.04) (0.516,8.085)rectangle++(0.04,0.04) (0.76,8.448)rectangle++(0.04,0.04)
  (0.982,8.746)rectangle++(0.04,0.04) (1.182,8.991)rectangle++(0.04,0.04) (1.362,9.191)rectangle++(0.04,0.04) (1.523,9.354)rectangle++(0.04,0.04) (1.667,9.487)rectangle++(0.04,0.04)
  (-0.014,1.698)rectangle++(0.04,0.04) (0.053,3.04)rectangle++(0.04,0.04) (0.165,4.178)rectangle++(0.04,0.04) (0.305,5.141)rectangle++(0.04,0.04) (0.463,5.954)rectangle++(0.04,0.04)
  (0.629,6.638)rectangle++(0.04,0.04) (0.798,7.214)rectangle++(0.04,0.04) (0.965,7.696)rectangle++(0.04,0.04) (1.126,8.099)rectangle++(0.04,0.04) (-1.068,3.308)rectangle++(0.04,0.04)
  (0.944,2.358)rectangle++(0.04,0.04) (0.895,3.563)rectangle++(0.04,0.04) (0.901,4.589)rectangle++(0.04,0.04) (0.947,5.46)rectangle++(0.04,0.04) (1.021,6.2)rectangle++(0.04,0.04)
  (-1.089,4.363)rectangle++(0.04,0.04) (-0.754,5.348)rectangle++(0.04,0.04) (-1.227,3.862)rectangle++(0.04,0.04) (-0.89,4.928)rectangle++(0.04,0.04) (-0.562,5.821)rectangle++(0.04,0.04)
  (-0.247,6.566)rectangle++(0.04,0.04) (0.051,7.187)rectangle++(0.04,0.04) (0.328,7.703)rectangle++(0.04,0.04) (-1.533,4.53)rectangle++(0.04,0.04) (-1.124,5.508)rectangle++(0.04,0.04)
  (-0.737,6.323)rectangle++(0.04,0.04) (0.005,3.644)rectangle++(0.04,0.04) (0.706,2.204)rectangle++(0.04,0.04) (0.686,3.441)rectangle++(0.04,0.04) (0.718,4.494)rectangle++(0.04,0.04)
  (0.788,5.387)rectangle++(0.04,0.04) (0.883,6.144)rectangle++(0.04,0.04) (0.994,6.783)rectangle++(0.04,0.04) (1.114,7.322)rectangle++(0.04,0.04) (1.238,7.775)rectangle++(0.04,0.04)
  (1.361,8.156)rectangle++(0.04,0.04) (1.481,8.474)rectangle++(0.04,0.04) (1.595,8.74)rectangle++(0.04,0.04) (1.704,8.961)rectangle++(0.04,0.04) (-1.553,5.036)rectangle++(0.04,0.04)
  (-1.121,5.939)rectangle++(0.04,0.04) (-1.418,3.91)rectangle++(0.04,0.04) (-1.051,4.976)rectangle++(0.04,0.04) (-0.696,5.868)rectangle++(0.04,0.04) (-0.359,6.612)rectangle++(0.04,0.04)
  (-0.043,7.231)rectangle++(0.04,0.04) (0.25,7.744)rectangle++(0.04,0.04) (0.52,8.169)rectangle++(0.04,0.04) (0.767,8.519)rectangle++(0.04,0.04) (0.99,8.806)rectangle++(0.04,0.04)
  (0.036,2.78)rectangle++(0.04,0.04) (-0.54,3.503)rectangle++(0.04,0.04) (-0.321,4.596)rectangle++(0.04,0.04) (-0.091,5.515)rectangle++(0.04,0.04) (-1.143,4.019)rectangle++(0.04,0.04)
  (-0.813,5.058)rectangle++(0.04,0.04) (-0.491,5.928)rectangle++(0.04,0.04) (-0.183,6.655)rectangle++(0.04,0.04) (-1.402,4.227)rectangle++(0.04,0.04) (-1.024,5.245)rectangle++(0.04,0.04)
  (-0.663,6.096)rectangle++(0.04,0.04) (-0.322,6.804)rectangle++(0.04,0.04) (-0.004,7.393)rectangle++(0.04,0.04) (0.29,7.88)rectangle++(0.04,0.04) (0.56,8.283)rectangle++(0.04,0.04)
  (0.805,8.614)rectangle++(0.04,0.04) (2.039,3.806)rectangle++(0.04,0.04) (1.883,4.75)rectangle++(0.04,0.04) (1.788,5.558)rectangle++(0.04,0.04) (1.74,6.249)rectangle++(0.04,0.04)
  (1.727,6.838)rectangle++(0.04,0.04) (1.739,7.34)rectangle++(0.04,0.04) (1.77,7.765)rectangle++(0.04,0.04) (1.813,8.126)rectangle++(0.04,0.04) (1.864,8.431)rectangle++(0.04,0.04)
  (1.869,4.011)rectangle++(0.04,0.04) (0.774,3.16)rectangle++(0.04,0.04) (0.782,4.251)rectangle++(0.04,0.04) (0.833,5.178)rectangle++(0.04,0.04) (0.913,5.965)rectangle++(0.04,0.04)
  (1.012,6.63)rectangle++(0.04,0.04) (1.123,7.191)rectangle++(0.04,0.04) (1.24,7.663)rectangle++(0.04,0.04) (1.359,8.061)rectangle++(0.04,0.04) (1.475,8.393)rectangle++(0.04,0.04)
  (-1.1,3.335)rectangle++(0.04,0.04) (-0.804,4.475)rectangle++(0.04,0.04) (-0.506,5.432)rectangle++(0.04,0.04) (-0.215,6.234)rectangle++(0.04,0.04) (0.064,6.903)rectangle++(0.04,0.04)
  (0.329,7.462)rectangle++(0.04,0.04) (0.576,7.925)rectangle++(0.04,0.04) (0.804,8.31)rectangle++(0.04,0.04) (1.014,8.627)rectangle++(0.04,0.04) (1.204,8.889)rectangle++(0.04,0.04)
  (1.377,9.104)rectangle++(0.04,0.04) (2.093,4.03)rectangle++(0.04,0.04) (1.938,4.938)rectangle++(0.04,0.04) (1.843,5.716)rectangle++(0.04,0.04) (1.793,6.381)rectangle++(0.04,0.04)
  (1.777,6.948)rectangle++(0.04,0.04) (1.786,7.431)rectangle++(0.04,0.04) (1.813,7.841)rectangle++(0.04,0.04) (1.853,8.189)rectangle++(0.04,0.04) (1.9,8.483)rectangle++(0.04,0.04)
  (1.952,8.73)rectangle++(0.04,0.04) (2.006,8.939)rectangle++(0.04,0.04) (2.061,9.114)rectangle++(0.04,0.04) (2.114,9.261)rectangle++(0.04,0.04) (2.165,9.383)rectangle++(0.04,0.04)
  (2.214,9.485)rectangle++(0.04,0.04) (0.543,3.119)rectangle++(0.04,0.04) (0.584,4.225)rectangle++(0.04,0.04) (0.978,2.91)rectangle++(0.04,0.04) (0.946,4.03)rectangle++(0.04,0.04)
  (0.963,4.984)rectangle++(0.04,0.04) (1.122,3.142)rectangle++(0.04,0.04) (1.077,4.222)rectangle++(0.04,0.04) (1.082,5.142)rectangle++(0.04,0.04) (1.123,5.923)rectangle++(0.04,0.04)
  (1.189,6.586)rectangle++(0.04,0.04) (1.272,7.147)rectangle++(0.04,0.04) (1.365,7.62)rectangle++(0.04,0.04) (0.029,2.619)rectangle++(0.04,0.04) (0.127,3.821)rectangle++(0.04,0.04)
  (0.259,4.839)rectangle++(0.04,0.04) (0.411,5.699)rectangle++(0.04,0.04) (0.575,6.424)rectangle++(0.04,0.04) (0.744,7.034)rectangle++(0.04,0.04) (0.911,7.545)rectangle++(0.04,0.04)
  (1.074,7.973)rectangle++(0.04,0.04) (1.23,8.33)rectangle++(0.04,0.04) (1.919,3.848)rectangle++(0.04,0.04) (1.783,4.79)rectangle++(0.04,0.04) (1.705,5.596)rectangle++(0.04,0.04)
  (1.671,6.285)rectangle++(0.04,0.04) (1.669,6.872)rectangle++(0.04,0.04) (1.691,7.37)rectangle++(0.04,0.04) (1.73,7.793)rectangle++(0.04,0.04) (1.78,8.151)rectangle++(0.04,0.04)
  (-1.367,4.891)rectangle++(0.04,0.04) (1.371,2.595)rectangle++(0.04,0.04) (1.267,3.747)rectangle++(0.04,0.04) (1.225,4.73)rectangle++(0.04,0.04) (-0.741,4.111)rectangle++(0.04,0.04)
  (-0.962,3.62)rectangle++(0.04,0.04) (-0.675,4.712)rectangle++(0.04,0.04) (-0.387,5.628)rectangle++(0.04,0.04) (-0.106,6.396)rectangle++(0.04,0.04) (0.021,2.453)rectangle++(0.04,0.04)
  (0.605,2.44)rectangle++(0.04,0.04) (0.61,3.646)rectangle++(0.04,0.04) (0.662,4.671)rectangle++(0.04,0.04) (-0.742,3.287)rectangle++(0.04,0.04) (-0.501,4.42)rectangle++(0.04,0.04)
  (-0.252,5.373)rectangle++(0.04,0.04) (-0.001,6.174)rectangle++(0.04,0.04) (0.244,6.844)rectangle++(0.04,0.04) (0.479,7.404)rectangle++(0.04,0.04) (0.701,7.87)rectangle++(0.04,0.04)
  (-0.114,3.832)rectangle++(0.04,0.04) (0.054,4.858)rectangle++(0.04,0.04) (0.238,5.723)rectangle++(0.04,0.04) (0.429,6.451)rectangle++(0.04,0.04) (0.62,7.063)rectangle++(0.04,0.04)
  (0.808,7.575)rectangle++(0.04,0.04) (1.781,3.606)rectangle++(0.04,0.04) (0.685,3.062)rectangle++(0.04,0.04) (0.702,4.171)rectangle++(0.04,0.04) (0.951,2.926)rectangle++(0.04,0.04)
  (-0.407,3.704)rectangle++(0.04,0.04) (-0.2,4.761)rectangle++(0.04,0.04) (0.018,5.651)rectangle++(0.04,0.04) (0.239,6.399)rectangle++(0.04,0.04) (0.457,7.025)rectangle++(0.04,0.04)
  (0.667,7.549)rectangle++(0.04,0.04) (0.867,7.987)rectangle++(0.04,0.04) (1.054,8.35)rectangle++(0.04,0.04) (1.228,8.652)rectangle++(0.04,0.04) (1.388,8.901)rectangle++(0.04,0.04)
  (1.533,9.107)rectangle++(0.04,0.04) (1.665,9.275)rectangle++(0.04,0.04) (1.784,9.414)rectangle++(0.04,0.04) (1.891,9.526)rectangle++(0.04,0.04) (1.986,9.618)rectangle++(0.04,0.04)
  (2.071,9.692)rectangle++(0.04,0.04) (2.146,9.752)rectangle++(0.04,0.04) (2.212,9.799)rectangle++(0.04,0.04) (2.27,9.837)rectangle++(0.04,0.04) (2.321,9.867)rectangle++(0.04,0.04)
  (-1.635,5.331)rectangle++(0.04,0.04) (-1.178,6.193)rectangle++(0.04,0.04) (-0.756,6.907)rectangle++(0.04,0.04) (-0.369,7.497)rectangle++(0.04,0.04) (-0.016,7.984)rectangle++(0.04,0.04)
  (0.304,8.383)rectangle++(0.04,0.04) (0.591,8.71)rectangle++(0.04,0.04) (0.849,8.976)rectangle++(0.04,0.04) (-1.709,4.879)rectangle++(0.04,0.04) (-1.26,5.811)rectangle++(0.04,0.04)
  (-0.841,6.586)rectangle++(0.04,0.04) (1.729,3.041)rectangle++(0.04,0.04) (1.589,4.112)rectangle++(0.04,0.04) (1.513,5.028)rectangle++(0.04,0.04) (1.485,5.809)rectangle++(0.04,0.04)
  (1.493,6.475)rectangle++(0.04,0.04) (0.022,2.463)rectangle++(0.04,0.04) (0.115,3.689)rectangle++(0.04,0.04) (0.243,4.727)rectangle++(0.04,0.04) (0.393,5.605)rectangle++(0.04,0.04)
  (0.556,6.345)rectangle++(0.04,0.04) (0.724,6.967)rectangle++(0.04,0.04) (0.892,7.489)rectangle++(0.04,0.04) (1.056,7.926)rectangle++(0.04,0.04) (-1.449,4.71)rectangle++(0.04,0.04)
  (-1.046,5.658)rectangle++(0.04,0.04) (-0.665,6.447)rectangle++(0.04,0.04) (-0.309,7.103)rectangle++(0.04,0.04) (0.273,2.817)rectangle++(0.04,0.04) (0.342,3.979)rectangle++(0.04,0.04)
  (0.448,4.965)rectangle++(0.04,0.04) (0.577,5.798)rectangle++(0.04,0.04) (0.72,6.502)rectangle++(0.04,0.04) (0.87,7.094)rectangle++(0.04,0.04) (-1.306,4.511)rectangle++(0.04,0.04)
  (-0.932,5.483)rectangle++(0.04,0.04) (-0.575,6.294)rectangle++(0.04,0.04) (-0.239,6.969)rectangle++(0.04,0.04) (1.754,3.252)rectangle++(0.04,0.04) (1.619,4.29)rectangle++(0.04,0.04)
  (1.545,5.178)rectangle++(0.04,0.04) (-0.779,4.26)rectangle++(0.04,0.04) (-0.494,5.248)rectangle++(0.04,0.04) (-0.212,6.077)rectangle++(0.04,0.04) (-1.284,4.108)rectangle++(0.04,0.04)
  (-0.929,5.139)rectangle++(0.04,0.04) (1.381,2.742)rectangle++(0.04,0.04) (1.282,3.872)rectangle++(0.04,0.04) (1.242,4.836)rectangle++(0.04,0.04) (0.339,2.592)rectangle++(0.04,0.04)
  (0.39,3.786)rectangle++(0.04,0.04) (0.481,4.799)rectangle++(0.04,0.04) (0.598,5.656)rectangle++(0.04,0.04) (0.733,6.38)rectangle++(0.04,0.04) (0.876,6.99)rectangle++(0.04,0.04)
  (1.022,7.503)rectangle++(0.04,0.04) (1.166,7.932)rectangle++(0.04,0.04) (1.307,8.292)rectangle++(0.04,0.04) (1.44,8.592)rectangle++(0.04,0.04) (1.566,8.842)rectangle++(0.04,0.04)
  (-1.552,4.987)rectangle++(0.04,0.04) (-1.122,5.897)rectangle++(0.04,0.04) (-0.72,6.653)rectangle++(0.04,0.04) (-0.348,7.28)rectangle++(0.04,0.04) (-0.007,7.799)rectangle++(0.04,0.04)
  (0.304,8.225)rectangle++(0.04,0.04) (0.585,8.576)rectangle++(0.04,0.04) (0.838,8.862)rectangle++(0.04,0.04) (1.065,9.095)rectangle++(0.04,0.04) (2.128,3.958)rectangle++(0.04,0.04)
  (1.965,4.875)rectangle++(0.04,0.04) (1.863,5.661)rectangle++(0.04,0.04) (1.808,6.334)rectangle++(0.04,0.04) (1.788,6.908)rectangle++(0.04,0.04) (1.794,7.396)rectangle++(0.04,0.04)
  (1.818,7.811)rectangle++(0.04,0.04) (1.856,8.163)rectangle++(0.04,0.04) (1.902,8.461)rectangle++(0.04,0.04) (1.953,8.712)rectangle++(0.04,0.04) (2.006,8.923)rectangle++(0.04,0.04)
  (2.06,9.1)rectangle++(0.04,0.04) (2.113,9.249)rectangle++(0.04,0.04) (2.164,9.374)rectangle++(0.04,0.04) (2.212,9.477)rectangle++(0.04,0.04) (2.257,9.563)rectangle++(0.04,0.04)
  (2.299,9.635)rectangle++(0.04,0.04) (2.337,9.694)rectangle++(0.04,0.04) (2.372,9.742)rectangle++(0.04,0.04) (2.404,9.782)rectangle++(0.04,0.04) (-1.602,5.329)rectangle++(0.04,0.04)
  (-1.15,6.19)rectangle++(0.04,0.04) (-0.733,6.904)rectangle++(0.04,0.04) (-0.349,7.494)rectangle++(0.04,0.04) (0.001,7.98)rectangle++(0.04,0.04) (0.318,8.379)rectangle++(0.04,0.04)
  (0.603,8.706)rectangle++(0.04,0.04) (0.003,2.078)rectangle++(0.04,0.04) (0.084,3.363)rectangle++(0.04,0.04) (0.204,4.451)rectangle++(0.04,0.04) (0.349,5.371)rectangle++(0.04,0.04)
  (1.289,3.077)rectangle++(0.04,0.04) (1.217,4.16)rectangle++(0.04,0.04) (1.198,5.084)rectangle++(0.04,0.04) (1.22,5.869)rectangle++(0.04,0.04) (1.311,3.371)rectangle++(0.04,0.04)
  (1.247,4.409)rectangle++(0.04,0.04) (1.234,5.294)rectangle++(0.04,0.04) (1.258,6.047)rectangle++(0.04,0.04) (-1.014,4.377)rectangle++(0.04,0.04) (-0.689,5.357)rectangle++(0.04,0.04)
  (-0.374,6.178)rectangle++(0.04,0.04) (-0.073,6.862)rectangle++(0.04,0.04) (0.21,7.432)rectangle++(0.04,0.04) (0.474,7.905)rectangle++(0.04,0.04) (0.717,8.296)rectangle++(0.04,0.04)
  (-1.589,4.72)rectangle++(0.04,0.04) (-1.164,5.672)rectangle++(0.04,0.04) (-0.765,6.464)rectangle++(0.04,0.04) (-0.394,7.121)rectangle++(0.04,0.04) (-0.052,7.665)rectangle++(0.04,0.04)
  (0.26,8.113)rectangle++(0.04,0.04) (0.544,8.482)rectangle++(0.04,0.04) (-0.102,3.889)rectangle++(0.04,0.04) (0.067,4.906)rectangle++(0.04,0.04) (0.251,5.764)rectangle++(0.04,0.04)
  (0.441,6.485)rectangle++(0.04,0.04) (0.632,7.091)rectangle++(0.04,0.04) (0.819,7.598)rectangle++(0.04,0.04) (1.786,3.613)rectangle++(0.04,0.04) (1.66,4.596)rectangle++(0.04,0.04)
  (1.593,5.436)rectangle++(0.04,0.04) (1.569,6.153)rectangle++(0.04,0.04) (1.578,6.764)rectangle++(0.04,0.04) (1.609,7.282)rectangle++(0.04,0.04) (1.657,7.722)rectangle++(0.04,0.04)
  (1.715,8.093)rectangle++(0.04,0.04) (1.779,8.407)rectangle++(0.04,0.04) (1.847,8.671)rectangle++(0.04,0.04) (1.914,8.893)rectangle++(0.04,0.04) (1.981,9.078)rectangle++(0.04,0.04)
  (-1.528,5.111)rectangle++(0.04,0.04) (-1.097,6.002)rectangle++(0.04,0.04) (-0.694,6.742)rectangle++(0.04,0.04) (-1.512,4.148)rectangle++(0.04,0.04) (-1.121,5.183)rectangle++(0.04,0.04)
  (-0.748,6.046)rectangle++(0.04,0.04) (-0.396,6.766)rectangle++(0.04,0.04) (-0.068,7.363)rectangle++(0.04,0.04) (0.234,7.857)rectangle++(0.04,0.04) (0.511,8.265)rectangle++(0.04,0.04)
  (0.763,8.601)rectangle++(0.04,0.04) (0.99,8.877)rectangle++(0.04,0.04) (1.195,9.102)rectangle++(0.04,0.04) (1.377,9.285)rectangle++(0.04,0.04) (1.54,9.433)rectangle++(0.04,0.04)
  (1.684,9.553)rectangle++(0.04,0.04) (1.811,9.649)rectangle++(0.04,0.04) (1.923,9.725)rectangle++(0.04,0.04) (2.022,9.786)rectangle++(0.04,0.04) (2.186,4.308)rectangle++(0.04,0.04)
  (2.029,5.171)rectangle++(0.04,0.04) (1.929,5.91)rectangle++(0.04,0.04) (1.874,6.543)rectangle++(0.04,0.04) (1.852,7.082)rectangle++(0.04,0.04) (1.855,7.542)rectangle++(0.04,0.04)
  (1.877,7.933)rectangle++(0.04,0.04) (1.91,8.264)rectangle++(0.04,0.04) (1.952,8.544)rectangle++(0.04,0.04) (1.999,8.781)rectangle++(0.04,0.04) (2.048,8.98)rectangle++(0.04,0.04)
  (2.098,9.147)rectangle++(0.04,0.04) (2.147,9.287)rectangle++(0.04,0.04) (2.194,9.405)rectangle++(0.04,0.04) (2.239,9.502)rectangle++(0.04,0.04) (2.281,9.584)rectangle++(0.04,0.04)
  (2.32,9.651)rectangle++(0.04,0.04) (2.356,9.707)rectangle++(0.04,0.04) (2.389,9.753)rectangle++(0.04,0.04) (2.418,9.79)rectangle++(0.04,0.04) (2.445,9.821)rectangle++(0.04,0.04)
  (2.469,9.847)rectangle++(0.04,0.04) (2.49,9.867)rectangle++(0.04,0.04) (2.509,9.884)rectangle++(0.04,0.04) (2.526,9.897)rectangle++(0.04,0.04) (2.541,9.908)rectangle++(0.04,0.04)
  (-1.604,5.38)rectangle++(0.04,0.04) (-1.15,6.233)rectangle++(0.04,0.04) (-0.731,6.94)rectangle++(0.04,0.04) (-0.346,7.525)rectangle++(0.04,0.04) (0.005,8.006)rectangle++(0.04,0.04)
  (1.982,3.507)rectangle++(0.04,0.04) (1.823,4.498)rectangle++(0.04,0.04) (1.727,5.347)rectangle++(0.04,0.04) (1.68,6.072)rectangle++(0.04,0.04) (1.668,6.69)rectangle++(0.04,0.04)
  (1.684,7.216)rectangle++(0.04,0.04) (1.717,7.662)rectangle++(0.04,0.04) (1.764,8.041)rectangle++(0.04,0.04) (1.819,8.36)rectangle++(0.04,0.04) (1.878,8.63)rectangle++(0.04,0.04)
  (1.94,8.856)rectangle++(0.04,0.04) (2.001,9.046)rectangle++(0.04,0.04) (2.06,9.206)rectangle++(0.04,0.04) (2.117,9.339)rectangle++(0.04,0.04) (2.171,9.449)rectangle++(0.04,0.04)
  (2.221,9.541)rectangle++(0.04,0.04) (2.267,9.617)rectangle++(0.04,0.04) (2.31,9.68)rectangle++(0.04,0.04) (2.348,9.732)rectangle++(0.04,0.04) (-1.601,5.309)rectangle++(0.04,0.04)
  (-1.151,6.173)rectangle++(0.04,0.04) (-0.733,6.889)rectangle++(0.04,0.04) (-0.35,7.481)rectangle++(0.04,0.04) (1.894,3.325)rectangle++(0.04,0.04) (1.74,4.347)rectangle++(0.04,0.04)
  (1.651,5.222)rectangle++(0.04,0.04) (1.61,5.968)rectangle++(0.04,0.04) (1.605,6.605)rectangle++(0.04,0.04) (1.626,7.146)rectangle++(0.04,0.04) (1.666,7.605)rectangle++(0.04,0.04)
  (1.718,7.994)rectangle++(0.04,0.04) (1.778,3.894)rectangle++(0.04,0.04) (1.665,4.835)rectangle++(0.04,0.04) (1.607,5.639)rectangle++(0.04,0.04) (-0.865,4.362)rectangle++(0.04,0.04)
  (-1.037,3.647)rectangle++(0.04,0.04) (1.025,2.43)rectangle++(0.04,0.04) (0.966,3.621)rectangle++(0.04,0.04) (0.964,4.635)rectangle++(0.04,0.04) (-0.768,4.043)rectangle++(0.04,0.04)
  (-0.493,5.064)rectangle++(0.04,0.04) (-0.219,5.92)rectangle++(0.04,0.04) (0.049,6.637)rectangle++(0.04,0.04) (0.305,7.236)rectangle++(0.04,0.04) (0.546,7.734)rectangle++(0.04,0.04)
  (0.772,8.149)rectangle++(0.04,0.04) (0.98,8.492)rectangle++(0.04,0.04) (1.17,8.775)rectangle++(0.04,0.04) (1.343,9.008)rectangle++(0.04,0.04) (1.5,9.199)rectangle++(0.04,0.04)
  (2.102,4.077)rectangle++(0.04,0.04) (1.948,4.977)rectangle++(0.04,0.04) (1.853,5.749)rectangle++(0.04,0.04) (1.802,6.409)rectangle++(0.04,0.04) (1.376,3.607)rectangle++(0.04,0.04)
  (1.312,4.607)rectangle++(0.04,0.04) (1.297,5.46)rectangle++(0.04,0.04) (1.319,6.185)rectangle++(0.04,0.04) (1.366,6.801)rectangle++(0.04,0.04) (1.431,7.322)rectangle++(0.04,0.04)
  (0.421,2.869)rectangle++(0.04,0.04) (0.47,4.018)rectangle++(0.04,0.04) (0.558,4.992)rectangle++(0.04,0.04) (0.672,5.817)rectangle++(0.04,0.04) (0.802,6.514)rectangle++(0.04,0.04)
  (0.94,7.101)rectangle++(0.04,0.04) (1.081,7.595)rectangle++(0.04,0.04) (1.22,8.008)rectangle++(0.04,0.04) (1.355,8.354)rectangle++(0.04,0.04) (1.484,8.643)rectangle++(0.04,0.04)
  (1.605,8.884)rectangle++(0.04,0.04) (-1.554,5.002)rectangle++(0.04,0.04) (-1.123,5.91)rectangle++(0.04,0.04) (-0.72,6.665)rectangle++(0.04,0.04) (-0.348,7.29)rectangle++(0.04,0.04)
  (-0.006,7.807)rectangle++(0.04,0.04) (0.305,8.232)rectangle++(0.04,0.04) (0.586,8.581)rectangle++(0.04,0.04) (0.839,8.867)rectangle++(0.04,0.04) (-1.688,4.855)rectangle++(0.04,0.04)
  (-1.243,5.791)rectangle++(0.04,0.04) (-0.827,6.568)rectangle++(0.04,0.04) (0.022,2.476)rectangle++(0.04,0.04) (0.116,3.7)rectangle++(0.04,0.04) (0.244,4.736)rectangle++(0.04,0.04)
  (0.395,5.612)rectangle++(0.04,0.04) (0.016,2.346)rectangle++(0.04,0.04) (0.579,2.418)rectangle++(0.04,0.04) (0.587,3.629)rectangle++(0.04,0.04) (-0.621,3.774)rectangle++(0.04,0.04)
  (-0.38,4.829)rectangle++(0.04,0.04) (-0.132,5.716)rectangle++(0.04,0.04) (0.115,6.46)rectangle++(0.04,0.04) (0.354,7.083)rectangle++(0.04,0.04) (0.582,7.602)rectangle++(0.04,0.04)
  (0.796,8.035)rectangle++(0.04,0.04) (0.996,8.394)rectangle++(0.04,0.04) (1.18,8.691)rectangle++(0.04,0.04) (1.349,8.937)rectangle++(0.04,0.04) (1.502,9.138)rectangle++(0.04,0.04)
  (1.64,9.304)rectangle++(0.04,0.04) (1.764,9.439)rectangle++(0.04,0.04) (1.874,9.549)rectangle++(0.04,0.04) (1.973,9.638)rectangle++(0.04,0.04) (2.06,9.709)rectangle++(0.04,0.04)
  (2.138,9.767)rectangle++(0.04,0.04) (2.205,9.812)rectangle++(0.04,0.04) (2.265,9.848)rectangle++(0.04,0.04) (2.317,9.877)rectangle++(0.04,0.04) (2.362,9.899)rectangle++(0.04,0.04)
  (2.402,9.916)rectangle++(0.04,0.04) (2.436,9.928)rectangle++(0.04,0.04) (2.465,9.938)rectangle++(0.04,0.04) (0.77,2.422)rectangle++(0.04,0.04) (0.749,3.624)rectangle++(0.04,0.04)
  (0.779,4.646)rectangle++(0.04,0.04) (1.059,3.035)rectangle++(0.04,0.04) (1.02,4.134)rectangle++(0.04,0.04) (-0.696,3.32)rectangle++(0.04,0.04) (-0.461,4.446)rectangle++(0.04,0.04)
  (-0.217,5.394)rectangle++(0.04,0.04) (-1.14,3.972)rectangle++(0.04,0.04) (-0.812,5.018)rectangle++(0.04,0.04) (-0.492,5.894)rectangle++(0.04,0.04) (-0.185,6.626)rectangle++(0.04,0.04)
  (0.106,7.235)rectangle++(0.04,0.04) (0.377,7.742)rectangle++(0.04,0.04) (0.628,8.162)rectangle++(0.04,0.04) (0.858,8.509)rectangle++(0.04,0.04) (1.068,8.794)rectangle++(0.04,0.04)
  (1.257,9.029)rectangle++(0.04,0.04) (-1.647,4.966)rectangle++(0.04,0.04) (-1.204,5.883)rectangle++(0.04,0.04) (-0.79,6.645)rectangle++(0.04,0.04) (-0.408,7.276)rectangle++(0.04,0.04)
  (-0.058,7.797)rectangle++(0.04,0.04) (-1.622,4.476)rectangle++(0.04,0.04) (-1.201,5.466)rectangle++(0.04,0.04) (-0.805,6.29)rectangle++(0.04,0.04) (-0.874,3.157)rectangle++(0.04,0.04)
  (-0.619,4.315)rectangle++(0.04,0.04) (-0.356,5.289)rectangle++(0.04,0.04) (-0.093,6.106)rectangle++(0.04,0.04) (0.163,6.79)rectangle++(0.04,0.04) (0.408,7.361)rectangle++(0.04,0.04)
  (0.639,7.837)rectangle++(0.04,0.04) (0.854,8.232)rectangle++(0.04,0.04) (1.053,8.559)rectangle++(0.04,0.04) (1.236,8.829)rectangle++(0.04,0.04) (1.401,9.052)rectangle++(0.04,0.04)
  (-0.173,3.966)rectangle++(0.04,0.04) (1.002,2.688)rectangle++(0.04,0.04) (-0.347,3.666)rectangle++(0.04,0.04) (-0.151,4.726)rectangle++(0.04,0.04) (0.059,5.62)rectangle++(0.04,0.04)
  (0.016,2.347)rectangle++(0.04,0.04) (0.105,3.591)rectangle++(0.04,0.04) (0.231,4.644)rectangle++(0.04,0.04) (0.38,5.534)rectangle++(0.04,0.04) (0.542,6.285)rectangle++(0.04,0.04)
  (0.71,6.917)rectangle++(0.04,0.04) (0.878,7.447)rectangle++(0.04,0.04) (1.042,7.891)rectangle++(0.04,0.04) (1.199,8.262)rectangle++(0.04,0.04) (1.348,8.571)rectangle++(0.04,0.04)
  (1.486,8.828)rectangle++(0.04,0.04) (1.614,9.04)rectangle++(0.04,0.04) (1.731,9.216)rectangle++(0.04,0.04) (1.838,9.36)rectangle++(0.04,0.04) (0.383,2.425)rectangle++(0.04,0.04)
  (0.421,3.642)rectangle++(0.04,0.04) (0.501,4.675)rectangle++(0.04,0.04) (0.611,5.55)rectangle++(0.04,0.04) (-1.025,4.158)rectangle++(0.04,0.04) (-0.707,5.171)rectangle++(0.04,0.04)
  (-0.396,6.02)rectangle++(0.04,0.04) (-0.098,6.729)rectangle++(0.04,0.04) (0.183,7.32)rectangle++(0.04,0.04) (0.446,7.811)rectangle++(0.04,0.04) (1.882,3.57)rectangle++(0.04,0.04)
  (1.74,4.555)rectangle++(0.04,0.04) (1.659,5.398)rectangle++(0.04,0.04) (1.624,6.119)rectangle++(0.04,0.04) (1.623,6.732)rectangle++(0.04,0.04) (1.647,7.253)rectangle++(0.04,0.04)
  (-1.201,4.688)rectangle++(0.04,0.04) (-0.835,5.629)rectangle++(0.04,0.04) (-0.487,6.414)rectangle++(0.04,0.04) (-0.16,7.068)rectangle++(0.04,0.04) (1.769,3.289)rectangle++(0.04,0.04)
  (-0.338,3.929)rectangle++(0.04,0.04) (-0.132,4.949)rectangle++(0.04,0.04) (0.084,5.808)rectangle++(0.04,0.04) (-1.174,4.11)rectangle++(0.04,0.04) (-0.836,5.137)rectangle++(0.04,0.04)
  (-0.507,5.996)rectangle++(0.04,0.04) (-0.193,6.713)rectangle++(0.04,0.04) (0.102,7.31)rectangle++(0.04,0.04) (1.798,3.395)rectangle++(0.04,0.04) (1.662,4.41)rectangle++(0.04,0.04)
  (0.897,3.19)rectangle++(0.04,0.04) (0.888,4.272)rectangle++(0.04,0.04) (0.924,5.192)rectangle++(0.04,0.04) (0.991,5.972)rectangle++(0.04,0.04) (1.079,6.633)rectangle++(0.04,0.04)
  (1.18,7.191)rectangle++(0.04,0.04) (1.643,3.617)rectangle++(0.04,0.04) (1.539,4.605)rectangle++(0.04,0.04) (1.49,5.449)rectangle++(0.04,0.04) (1.482,6.168)rectangle++(0.04,0.04)
  (1.504,6.78)rectangle++(0.04,0.04) (1.548,7.299)rectangle++(0.04,0.04) (1.605,7.738)rectangle++(0.04,0.04) (1.672,8.109)rectangle++(0.04,0.04) (1.812,3.905)rectangle++(0.04,0.04)
  (1.695,4.843)rectangle++(0.04,0.04) (1.632,5.645)rectangle++(0.04,0.04) (1.611,6.329)rectangle++(0.04,0.04) (1.62,6.912)rectangle++(0.04,0.04) (1.651,7.406)rectangle++(0.04,0.04)
  (1.641,3.765)rectangle++(0.04,0.04) (1.543,4.731)rectangle++(0.04,0.04) (1.499,5.556)rectangle++(0.04,0.04) (1.494,6.259)rectangle++(0.04,0.04) (1.518,6.856)rectangle++(0.04,0.04)
  (1.562,7.363)rectangle++(0.04,0.04) (1.62,7.793)rectangle++(0.04,0.04) (1.687,8.155)rectangle++(0.04,0.04) (1.758,8.461)rectangle++(0.04,0.04) (1.83,8.717)rectangle++(0.04,0.04)
  (-1.478,5.012)rectangle++(0.04,0.04) (-1.058,5.915)rectangle++(0.04,0.04) (-0.665,6.667)rectangle++(0.04,0.04) (-0.301,7.289)rectangle++(0.04,0.04) (0.034,7.804)rectangle++(0.04,0.04)
  (0.338,8.228)rectangle++(0.04,0.04) (0.615,8.577)rectangle++(0.04,0.04) (0.863,8.862)rectangle++(0.04,0.04) (1.086,9.094)rectangle++(0.04,0.04) (1.285,9.283)rectangle++(0.04,0.04)
  (1.461,9.435)rectangle++(0.04,0.04) (1.617,9.558)rectangle++(0.04,0.04) (2.177,4.172)rectangle++(0.04,0.04) (-0.452,4.18)rectangle++(0.04,0.04) (-0.219,5.167)rectangle++(0.04,0.04)
  (0.018,5.997)rectangle++(0.04,0.04) (0.253,6.693)rectangle++(0.04,0.04) (1.628,3.31)rectangle++(0.04,0.04) (1.514,4.345)rectangle++(0.04,0.04) (1.458,5.229)rectangle++(0.04,0.04)
  (1.446,5.982)rectangle++(0.04,0.04) (1.467,6.623)rectangle++(0.04,0.04) (1.509,7.167)rectangle++(0.04,0.04) (1.567,7.628)rectangle++(0.04,0.04) (1.705,3.789)rectangle++(0.04,0.04)
  (1.599,4.749)rectangle++(0.04,0.04) (1.547,5.569)rectangle++(0.04,0.04) (1.535,6.268)rectangle++(0.04,0.04) (1.554,6.862)rectangle++(0.04,0.04) (1.593,7.367)rectangle++(0.04,0.04)
  (1.646,7.794)rectangle++(0.04,0.04) (1.709,8.156)rectangle++(0.04,0.04) (1.777,8.46)rectangle++(0.04,0.04) (1.846,8.716)rectangle++(0.04,0.04) (1.916,8.931)rectangle++(0.04,0.04)
  (1.984,9.111)rectangle++(0.04,0.04) (2.048,9.261)rectangle++(0.04,0.04) (2.109,9.386)rectangle++(0.04,0.04) (0.04,2.859)rectangle++(0.04,0.04) (0.146,4.025)rectangle++(0.04,0.04)
  (0.283,5.012)rectangle++(0.04,0.04) (0.439,5.845)rectangle++(0.04,0.04) (0.605,6.547)rectangle++(0.04,0.04) (1.55,3.362)rectangle++(0.04,0.04) (1.45,4.392)rectangle++(0.04,0.04)
  (1.406,5.271)rectangle++(0.04,0.04) (1.404,6.02)rectangle++(0.04,0.04) (1.432,6.657)rectangle++(0.04,0.04) (1.481,7.198)rectangle++(0.04,0.04) (1.545,7.655)rectangle++(0.04,0.04)
  (1.617,8.041)rectangle++(0.04,0.04) (1.694,8.367)rectangle++(0.04,0.04) (1.772,8.64)rectangle++(0.04,0.04) (1.85,8.869)rectangle++(0.04,0.04) (1.925,9.061)rectangle++(0.04,0.04)
  (1.996,9.221)rectangle++(0.04,0.04) (2.064,9.354)rectangle++(0.04,0.04) (2.126,9.465)rectangle++(0.04,0.04) (2.184,9.556)rectangle++(0.04,0.04) (2.111,4.301)rectangle++(0.04,0.04)
  (1.964,5.167)rectangle++(0.04,0.04) (1.874,5.91)rectangle++(0.04,0.04) (1.827,6.545)rectangle++(0.04,0.04) (1.812,7.086)rectangle++(0.04,0.04) (1.822,7.547)rectangle++(0.04,0.04)
  (1.848,7.938)rectangle++(0.04,0.04) (1.886,8.27)rectangle++(0.04,0.04) (1.932,8.55)rectangle++(0.04,0.04) (1.982,8.786)rectangle++(0.04,0.04) (2.034,8.985)rectangle++(0.04,0.04)
  (2.086,9.152)rectangle++(0.04,0.04) (2.137,9.292)rectangle++(0.04,0.04) (-1.546,5.183)rectangle++(0.04,0.04) (-1.109,6.063)rectangle++(0.04,0.04) (-0.702,6.794)rectangle++(0.04,0.04)
  (-0.327,7.399)rectangle++(0.04,0.04) (0.016,7.899)rectangle++(0.04,0.04) (0.327,8.31)rectangle++(0.04,0.04) (0.608,8.646)rectangle++(0.04,0.04) (0.861,8.921)rectangle++(0.04,0.04)
  (1.086,9.145)rectangle++(0.04,0.04) (1.287,9.326)rectangle++(0.04,0.04) (1.465,9.472)rectangle++(0.04,0.04) (1.622,9.589)rectangle++(0.04,0.04) (1.76,9.682)rectangle++(0.04,0.04)
  (2.191,4.228)rectangle++(0.04,0.04) (2.03,5.103)rectangle++(0.04,0.04) (0.411,2.634)rectangle++(0.04,0.04) (0.452,3.819)rectangle++(0.04,0.04) (0.535,4.824)rectangle++(0.04,0.04)
  (0.646,5.675)rectangle++(0.04,0.04) (0.774,6.394)rectangle++(0.04,0.04) (0.911,7.001)rectangle++(0.04,0.04) (1.052,7.51)rectangle++(0.04,0.04) (1.193,7.938)rectangle++(0.04,0.04)
  (1.329,8.296)rectangle++(0.04,0.04) (1.459,8.594)rectangle++(0.04,0.04) (1.582,8.843)rectangle++(0.04,0.04) (1.696,9.049)rectangle++(0.04,0.04) (1.802,9.22)rectangle++(0.04,0.04)
  (-1.59,5.106)rectangle++(0.04,0.04) (-1.325,3.657)rectangle++(0.04,0.04) (-0.982,4.758)rectangle++(0.04,0.04) (-0.647,5.68)rectangle++(0.04,0.04) (-0.325,6.45)rectangle++(0.04,0.04)
  (-0.02,7.092)rectangle++(0.04,0.04) (-1.465,4.344)rectangle++(0.04,0.04) (-1.074,5.347)rectangle++(0.04,0.04) (-0.701,6.184)rectangle++(0.04,0.04) (-0.351,6.881)rectangle++(0.04,0.04)
  (-0.025,7.459)rectangle++(0.04,0.04) (0.275,7.937)rectangle++(0.04,0.04) (0.549,8.332)rectangle++(0.04,0.04) (0.798,8.656)rectangle++(0.04,0.04) (1.022,8.922)rectangle++(0.04,0.04)
  (1.223,9.139)rectangle++(0.04,0.04) (1.403,9.316)rectangle++(0.04,0.04) (1.563,9.458)rectangle++(0.04,0.04) (1.705,9.573)rectangle++(0.04,0.04) (1.83,9.665)rectangle++(0.04,0.04)
  (2.179,4.241)rectangle++(0.04,0.04) (2.02,5.114)rectangle++(0.04,0.04) (-0.68,4.336)rectangle++(0.04,0.04) (-0.407,5.309)rectangle++(0.04,0.04) (-0.136,6.125)rectangle++(0.04,0.04)
  (1.533,3.113)rectangle++(0.04,0.04) (1.425,4.181)rectangle++(0.04,0.04) (1.376,5.093)rectangle++(0.04,0.04) (1.371,5.87)rectangle++(0.04,0.04) (1.398,6.531)rectangle++(0.04,0.04)
  (1.448,7.091)rectangle++(0.04,0.04) (1.512,7.566)rectangle++(0.04,0.04) (1.586,7.967)rectangle++(0.04,0.04) (1.664,8.305)rectangle++(0.04,0.04) (1.745,8.589)rectangle++(0.04,0.04)
  (1.824,8.827)rectangle++(0.04,0.04) (1.901,9.026)rectangle++(0.04,0.04) (1.975,9.192)rectangle++(0.04,0.04) (2.044,9.331)rectangle++(0.04,0.04) (2.109,9.445)rectangle++(0.04,0.04)
  (2.168,9.54)rectangle++(0.04,0.04) (2.222,9.619)rectangle++(0.04,0.04) (2.271,9.683)rectangle++(0.04,0.04) (2.316,9.736)rectangle++(0.04,0.04) (2.356,9.779)rectangle++(0.04,0.04)
  (2.391,9.814)rectangle++(0.04,0.04) (2.423,9.843)rectangle++(0.04,0.04) (2.451,9.866)rectangle++(0.04,0.04) (2.476,9.884)rectangle++(0.04,0.04) (2.498,9.899)rectangle++(0.04,0.04)
  (2.517,9.91)rectangle++(0.04,0.04) (2.533,9.919)rectangle++(0.04,0.04) (2.548,9.926)rectangle++(0.04,0.04) (2.561,9.931)rectangle++(0.04,0.04) (2.572,9.936)rectangle++(0.04,0.04)
  (2.581,9.939)rectangle++(0.04,0.04) (2.589,9.941)rectangle++(0.04,0.04) (2.596,9.942)rectangle++(0.04,0.04) (2.602,9.943)rectangle++(0.04,0.04) (2.608,9.944)rectangle++(0.04,0.04)
  (2.612,9.944)rectangle++(0.04,0.04) (2.616,9.944)rectangle++(0.04,0.04) (-1.598,5.401)rectangle++(0.04,0.04) (-1.145,6.251)rectangle++(0.04,0.04) (-0.725,6.955)rectangle++(0.04,0.04)
  (-0.34,7.537)rectangle++(0.04,0.04) (0.01,8.017)rectangle++(0.04,0.04) (0.327,8.41)rectangle++(0.04,0.04) (0.612,8.731)rectangle++(0.04,0.04) (-1.7,4.786)rectangle++(0.04,0.04)
  (-1.256,5.732)rectangle++(0.04,0.04) (-0.84,6.519)rectangle++(0.04,0.04) (-0.456,7.171)rectangle++(0.04,0.04) (-0.103,7.71)rectangle++(0.04,0.04) (1.921,3.426)rectangle++(0.04,0.04)
  (1.768,4.431)rectangle++(0.04,0.04) (1.678,5.292)rectangle++(0.04,0.04) (1.635,6.027)rectangle++(0.04,0.04) (1.629,6.654)rectangle++(0.04,0.04) (1.649,7.187)rectangle++(0.04,0.04)
  (1.687,7.639)rectangle++(0.04,0.04) (1.737,8.022)rectangle++(0.04,0.04) (1.795,8.345)rectangle++(0.04,0.04) (1.856,3.979)rectangle++(0.04,0.04) (-0.008,1.837)rectangle++(0.04,0.04)
  (0.065,3.158)rectangle++(0.04,0.04) (0.179,4.278)rectangle++(0.04,0.04) (0.321,5.225)rectangle++(0.04,0.04) (0.48,6.025)rectangle++(0.04,0.04) (0.647,6.698)rectangle++(0.04,0.04)
  (0.815,7.264)rectangle++(0.04,0.04) (1.703,3.548)rectangle++(0.04,0.04) (1.587,4.544)rectangle++(0.04,0.04) (1.529,5.395)rectangle++(0.04,0.04) (1.513,6.121)rectangle++(0.04,0.04)
  (1.529,6.738)rectangle++(0.04,0.04) (1.567,7.263)rectangle++(0.04,0.04) (1.62,7.707)rectangle++(0.04,0.04) (1.683,8.082)rectangle++(0.04,0.04) (1.752,8.399)rectangle++(0.04,0.04)
  (1.823,8.665)rectangle++(0.04,0.04) (1.894,8.889)rectangle++(0.04,0.04) (1.963,9.076)rectangle++(0.04,0.04) (2.017,4.158)rectangle++(0.04,0.04) (1.879,5.05)rectangle++(0.04,0.04)
  (-0.691,4.297)rectangle++(0.04,0.04) (-0.418,5.276)rectangle++(0.04,0.04) (-0.147,6.098)rectangle++(0.04,0.04) (-0.055,3.032)rectangle++(0.04,0.04) (0.072,4.176)rectangle++(0.04,0.04)
  (0.226,5.143)rectangle++(0.04,0.04) (0.396,5.958)rectangle++(0.04,0.04) (0.573,6.645)rectangle++(0.04,0.04) (1.578,3.374)rectangle++(0.04,0.04) (0.651,2.971)rectangle++(0.04,0.04)
  (0.67,4.096)rectangle++(0.04,0.04) (0.731,5.051)rectangle++(0.04,0.04) (0.821,5.86)rectangle++(0.04,0.04) (0.93,6.544)rectangle++(0.04,0.04) (1.05,7.122)rectangle++(0.04,0.04)
  (-1.28,4.55)rectangle++(0.04,0.04) (-0.908,5.515)rectangle++(0.04,0.04) (1.472,2.82)rectangle++(0.04,0.04) (1.362,3.934)rectangle++(0.04,0.04) (1.313,4.886)rectangle++(0.04,0.04)
  (1.309,5.696)rectangle++(0.04,0.04) (1.338,6.386)rectangle++(0.04,0.04) (1.391,6.971)rectangle++(0.04,0.04) (1.563,3.622)rectangle++(0.04,0.04) (-0.445,3.956)rectangle++(0.04,0.04)
  (-0.222,4.976)rectangle++(0.04,0.04) (0.008,5.835)rectangle++(0.04,0.04) (0.238,6.556)rectangle++(0.04,0.04) (1.596,3.28)rectangle++(0.04,0.04) (1.485,4.321)rectangle++(0.04,0.04)
  (1.433,5.209)rectangle++(0.04,0.04) (1.424,5.967)rectangle++(0.04,0.04) (1.447,6.611)rectangle++(0.04,0.04) (1.492,7.158)rectangle++(0.04,0.04) (1.553,7.621)rectangle++(0.04,0.04)
  (1.622,8.012)rectangle++(0.04,0.04) (1.697,8.341)rectangle++(0.04,0.04) (1.774,8.618)rectangle++(0.04,0.04) (0.035,2.755)rectangle++(0.04,0.04) (0.138,3.936)rectangle++(0.04,0.04)
  (0.273,4.937)rectangle++(0.04,0.04) (0.427,5.781)rectangle++(0.04,0.04) (0.592,6.493)rectangle++(0.04,0.04) (0.761,7.092)rectangle++(0.04,0.04) (0.928,7.594)rectangle++(0.04,0.04)
  (1.09,8.014)rectangle++(0.04,0.04) (-0.003,1.941)rectangle++(0.04,0.04) (0.073,3.246)rectangle++(0.04,0.04) (0.19,4.353)rectangle++(0.04,0.04) (0.24,2.522)rectangle++(0.04,0.04)
  (0.303,3.731)rectangle++(0.04,0.04) (-0.694,3.741)rectangle++(0.04,0.04) (-0.442,4.804)rectangle++(0.04,0.04) (-0.186,5.697)rectangle++(0.04,0.04) (0.067,6.446)rectangle++(0.04,0.04)
  (0.313,7.073)rectangle++(0.04,0.04) (0.547,7.596)rectangle++(0.04,0.04) (0.766,8.031)rectangle++(0.04,0.04) (0.97,8.391)rectangle++(0.04,0.04) (1.158,8.69)rectangle++(0.04,0.04)
  (1.33,8.936)rectangle++(0.04,0.04) (1.486,9.139)rectangle++(0.04,0.04) (1.626,9.305)rectangle++(0.04,0.04) (1.752,9.44)rectangle++(0.04,0.04) (1.865,9.55)rectangle++(0.04,0.04)
  (1.965,9.639)rectangle++(0.04,0.04) (2.054,9.711)rectangle++(0.04,0.04) (2.132,9.769)rectangle++(0.04,0.04) (2.201,9.814)rectangle++(0.04,0.04) (2.261,9.85)rectangle++(0.04,0.04)
  (2.314,9.878)rectangle++(0.04,0.04) (2.359,9.9)rectangle++(0.04,0.04) (2.399,9.917)rectangle++(0.04,0.04) (2.434,9.93)rectangle++(0.04,0.04) (-1.628,5.364)rectangle++(0.04,0.04)
  (-1.171,6.221)rectangle++(0.04,0.04) (-0.749,6.931)rectangle++(0.04,0.04) (-0.361,7.517)rectangle++(0.04,0.04) (1.904,3.33)rectangle++(0.04,0.04) (1.749,4.35)rectangle++(0.04,0.04)
  (1.659,5.224)rectangle++(0.04,0.04) (1.617,5.97)rectangle++(0.04,0.04) (-0.934,4.429)rectangle++(0.04,0.04) (-0.619,5.398)rectangle++(0.04,0.04) (-0.312,6.21)rectangle++(0.04,0.04)
  (-0.019,6.887)rectangle++(0.04,0.04) (1.708,3.286)rectangle++(0.04,0.04) (1.581,4.321)rectangle++(0.04,0.04) (-0.721,3.365)rectangle++(0.04,0.04) (-0.481,4.486)rectangle++(0.04,0.04)
  (-0.231,5.428)rectangle++(0.04,0.04) (0.018,6.219)rectangle++(0.04,0.04) (0.262,6.882)rectangle++(0.04,0.04) (0.496,7.435)rectangle++(0.04,0.04) (0.717,7.896)rectangle++(0.04,0.04)
  (-1.504,4.641)rectangle++(0.04,0.04) (-1.095,5.601)rectangle++(0.04,0.04) (-0.806,3.137)rectangle++(0.04,0.04) (-0.562,4.295)rectangle++(0.04,0.04) (-0.308,5.269)rectangle++(0.04,0.04)
  (-0.053,6.087)rectangle++(0.04,0.04) (0.196,6.773)rectangle++(0.04,0.04) (0.435,7.345)rectangle++(0.04,0.04) (0.662,7.822)rectangle++(0.04,0.04) (0.873,8.218)rectangle++(0.04,0.04)
  (1.069,8.547)rectangle++(0.04,0.04) (1.248,8.818)rectangle++(0.04,0.04) (1.411,9.042)rectangle++(0.04,0.04) (1.559,9.225)rectangle++(0.04,0.04) (1.692,9.375)rectangle++(0.04,0.04)
  (1.811,9.498)rectangle++(0.04,0.04) (1.917,9.597)rectangle++(0.04,0.04) (2.011,9.677)rectangle++(0.04,0.04) (2.161,4.285)rectangle++(0.04,0.04) (2.006,5.152)rectangle++(0.04,0.04)
  (1.909,5.895)rectangle++(0.04,0.04) (1.856,6.531)rectangle++(0.04,0.04) (1.837,7.073)rectangle++(0.04,0.04) (-1.128,4.688)rectangle++(0.04,0.04) (1.293,2.61)rectangle++(0.04,0.04)
  (0.495,2.759)rectangle++(0.04,0.04) (0.529,3.922)rectangle++(0.04,0.04) (0.604,4.908)rectangle++(0.04,0.04) (0.708,5.744)rectangle++(0.04,0.04) (0.829,6.45)rectangle++(0.04,0.04)
  (0.961,7.046)rectangle++(0.04,0.04) (1.096,7.547)rectangle++(0.04,0.04) (-1.362,4.643)rectangle++(0.04,0.04) (-0.974,5.597)rectangle++(0.04,0.04) (-0.606,6.393)rectangle++(0.04,0.04)
  (-0.262,7.054)rectangle++(0.04,0.04) (0.058,7.603)rectangle++(0.04,0.04) (0.351,8.056)rectangle++(0.04,0.04) (0.618,8.43)rectangle++(0.04,0.04) (0.861,8.737)rectangle++(0.04,0.04)
  (1.079,8.988)rectangle++(0.04,0.04) (1.274,9.193)rectangle++(0.04,0.04) (1.449,9.359)rectangle++(0.04,0.04) (1.603,9.494)rectangle++(0.04,0.04) (1.741,9.602)rectangle++(0.04,0.04)
  (1.861,9.688)rectangle++(0.04,0.04) (1.967,9.757)rectangle++(0.04,0.04) (-1.674,5.243)rectangle++(0.04,0.04) (-1.216,6.119)rectangle++(0.04,0.04) (-0.791,6.846)rectangle++(0.04,0.04)
  (-0.401,7.447)rectangle++(0.04,0.04) (-0.045,7.942)rectangle++(0.04,0.04) (0.277,8.349)rectangle++(0.04,0.04) (0.568,8.682)rectangle++(0.04,0.04) (0.828,8.953)rectangle++(0.04,0.04)
  (1.059,9.173)rectangle++(0.04,0.04) (1.265,9.351)rectangle++(0.04,0.04) (1.447,9.494)rectangle++(0.04,0.04) (2.181,4.121)rectangle++(0.04,0.04) (2.017,5.012)rectangle++(0.04,0.04)
  (1.913,5.776)rectangle++(0.04,0.04) (1.855,6.429)rectangle++(0.04,0.04) (1.831,6.987)rectangle++(0.04,0.04) (1.834,7.462)rectangle++(0.04,0.04) (1.855,7.865)rectangle++(0.04,0.04)
  (1.889,8.207)rectangle++(0.04,0.04) (1.932,8.497)rectangle++(0.04,0.04) (1.98,8.741)rectangle++(0.04,0.04) (2.03,8.947)rectangle++(0.04,0.04) (-1.491,5.095)rectangle++(0.04,0.04)
  (-1.066,5.986)rectangle++(0.04,0.04) (-0.669,6.727)rectangle++(0.04,0.04) (-0.302,7.341)rectangle++(0.04,0.04) (0.035,7.848)rectangle++(0.04,0.04) (0.341,8.266)rectangle++(0.04,0.04)
  (0.619,8.608)rectangle++(0.04,0.04) (0.868,8.889)rectangle++(0.04,0.04) (1.091,9.117)rectangle++(0.04,0.04) (1.29,9.302)rectangle++(0.04,0.04) (1.466,9.451)rectangle++(0.04,0.04)
  (1.622,9.571)rectangle++(0.04,0.04) (1.76,9.667)rectangle++(0.04,0.04) (-1.693,5.186)rectangle++(0.04,0.04) (-1.233,6.072)rectangle++(0.04,0.04) (-0.808,6.807)rectangle++(0.04,0.04)
  (-0.417,7.414)rectangle++(0.04,0.04) (-0.06,7.915)rectangle++(0.04,0.04) (0.264,8.326)rectangle++(0.04,0.04) (0.555,8.663)rectangle++(0.04,0.04) (0.816,8.938)rectangle++(0.04,0.04)
  (1.049,9.16)rectangle++(0.04,0.04) (0.361,2.237)rectangle++(0.04,0.04) (0.394,3.483)rectangle++(0.04,0.04) (0.002,2.056)rectangle++(0.04,0.04) (0.082,3.344)rectangle++(0.04,0.04)
  (-0.652,3.623)rectangle++(0.04,0.04) (-0.411,4.702)rectangle++(0.04,0.04) (-0.164,5.609)rectangle++(0.04,0.04) (0.083,6.371)rectangle++(0.04,0.04) (0.323,7.008)rectangle++(0.04,0.04)
  (0.553,7.54)rectangle++(0.04,0.04) (0.769,7.983)rectangle++(0.04,0.04) (0.971,8.351)rectangle++(0.04,0.04) (1.157,8.656)rectangle++(0.04,0.04) (2.008,3.894)rectangle++(0.04,0.04)
  (1.86,4.826)rectangle++(0.04,0.04) (1.772,5.624)rectangle++(0.04,0.04) (1.729,6.305)rectangle++(0.04,0.04) (1.72,6.887)rectangle++(0.04,0.04) (1.735,7.381)rectangle++(0.04,0.04)
  (1.768,7.801)rectangle++(0.04,0.04) (1.812,8.156)rectangle++(0.04,0.04) (1.865,8.456)rectangle++(0.04,0.04) (1.921,8.71)rectangle++(0.04,0.04) (1.979,8.922)rectangle++(0.04,0.04)
  (2.037,9.101)rectangle++(0.04,0.04) (-1.523,5.126)rectangle++(0.04,0.04) (-1.092,6.014)rectangle++(0.04,0.04) (-0.69,6.752)rectangle++(0.04,0.04) (-0.318,7.363)rectangle++(0.04,0.04)
  (0.022,7.868)rectangle++(0.04,0.04) (0.331,8.283)rectangle++(0.04,0.04) (0.611,8.623)rectangle++(0.04,0.04) (0.862,8.902)rectangle++(0.04,0.04) (1.086,9.128)rectangle++(0.04,0.04)
  (1.286,9.312)rectangle++(0.04,0.04) (1.464,9.46)rectangle++(0.04,0.04) (1.62,9.578)rectangle++(0.04,0.04) (1.758,9.673)rectangle++(0.04,0.04) (1.879,9.748)rectangle++(0.04,0.04)
  (1.985,9.807)rectangle++(0.04,0.04) (2.196,4.304)rectangle++(0.04,0.04) (2.036,5.167)rectangle++(0.04,0.04) (1.935,5.906)rectangle++(0.04,0.04) (1.236,3.54)rectangle++(0.04,0.04)
  (1.19,4.556)rectangle++(0.04,0.04) (1.192,5.421)rectangle++(0.04,0.04) (1.227,6.157)rectangle++(0.04,0.04) (1.287,6.78)rectangle++(0.04,0.04) (1.363,7.308)rectangle++(0.04,0.04)
  (1.449,7.753)rectangle++(0.04,0.04) (1.54,8.129)rectangle++(0.04,0.04) (1.632,8.444)rectangle++(0.04,0.04) (1.722,8.708)rectangle++(0.04,0.04) (1.81,8.929)rectangle++(0.04,0.04)
  (1.894,9.114)rectangle++(0.04,0.04) (1.972,9.267)rectangle++(0.04,0.04) (0.039,2.843)rectangle++(0.04,0.04) (0.145,4.011)rectangle++(0.04,0.04) (0.282,5)rectangle++(0.04,0.04)
  (0.437,5.835)rectangle++(0.04,0.04) (0.603,6.538)rectangle++(0.04,0.04) (0.772,7.13)rectangle++(0.04,0.04) (0.939,7.626)rectangle++(0.04,0.04) (1.101,8.04)rectangle++(0.04,0.04)
  (1.255,8.386)rectangle++(0.04,0.04) (1.4,8.675)rectangle++(0.04,0.04) (1.535,8.914)rectangle++(0.04,0.04) (1.659,9.111)rectangle++(0.04,0.04) (1.772,9.274)rectangle++(0.04,0.04)
  (1.875,9.409)rectangle++(0.04,0.04) (1.968,9.518)rectangle++(0.04,0.04) (2.051,9.608)rectangle++(0.04,0.04) (2.126,9.681)rectangle++(0.04,0.04) (2.192,9.74)rectangle++(0.04,0.04)
  (2.251,9.788)rectangle++(0.04,0.04) (2.302,9.826)rectangle++(0.04,0.04) (2.163,4.379)rectangle++(0.04,0.04) (2.012,5.232)rectangle++(0.04,0.04) (1.917,5.963)rectangle++(0.04,0.04)
  (1.866,6.588)rectangle++(0.04,0.04) (1.847,7.121)rectangle++(0.04,0.04) (-1.137,4.7)rectangle++(0.04,0.04) (-0.78,5.637)rectangle++(0.04,0.04) (1.487,2.872)rectangle++(0.04,0.04)
  (1.377,3.978)rectangle++(0.04,0.04) (1.327,4.922)rectangle++(0.04,0.04) (1.323,5.727)rectangle++(0.04,0.04) (1.351,6.411)rectangle++(0.04,0.04) (1.403,6.992)rectangle++(0.04,0.04)
  (1.47,7.483)rectangle++(0.04,0.04) (1.547,7.898)rectangle++(0.04,0.04) (1.628,8.248)rectangle++(0.04,0.04) (1.712,8.542)rectangle++(0.04,0.04) (1.794,8.788)rectangle++(0.04,0.04)
  (-1.5,5.019)rectangle++(0.04,0.04) (-1.076,5.922)rectangle++(0.04,0.04) (-0.68,6.673)rectangle++(0.04,0.04) (-0.313,7.296)rectangle++(0.04,0.04) (0.023,7.81)rectangle++(0.04,0.04)
  (0.33,8.234)rectangle++(0.04,0.04) (-1.645,4.635)rectangle++(0.04,0.04) (-1.215,5.602)rectangle++(0.04,0.04) (-0.811,6.406)rectangle++(0.04,0.04) (-0.435,7.074)rectangle++(0.04,0.04)
  (-1.536,4.263)rectangle++(0.04,0.04) (-1.137,5.281)rectangle++(0.04,0.04) (-0.757,6.13)rectangle++(0.04,0.04) (-0.401,6.837)rectangle++(0.04,0.04) (-0.069,7.424)rectangle++(0.04,0.04)
  (0.236,7.909)rectangle++(0.04,0.04) (0.515,8.31)rectangle++(0.04,0.04) (0.768,8.639)rectangle++(0.04,0.04) (0.996,8.909)rectangle++(0.04,0.04) (-1.142,3.293)rectangle++(0.04,0.04)
  (-0.841,4.441)rectangle++(0.04,0.04) (1.198,2.628)rectangle++(0.04,0.04) (1.121,3.782)rectangle++(0.04,0.04) (1.102,4.766)rectangle++(0.04,0.04) (1.125,5.603)rectangle++(0.04,0.04)
  (1.178,6.314)rectangle++(0.04,0.04) (1.252,6.916)rectangle++(0.04,0.04) (1.566,3.58)rectangle++(0.04,0.04) (1.472,4.577)rectangle++(0.04,0.04) (1.432,5.427)rectangle++(0.04,0.04)
  (1.432,6.152)rectangle++(0.04,0.04) (1.461,6.768)rectangle++(0.04,0.04) (1.511,7.291)rectangle++(0.04,0.04) (1.573,7.733)rectangle++(0.04,0.04) (1.645,8.106)rectangle++(0.04,0.04)
  (1.72,8.421)rectangle++(0.04,0.04) (-0.198,3.91)rectangle++(0.04,0.04) (-0.014,4.928)rectangle++(0.04,0.04) (0.183,5.785)rectangle++(0.04,0.04) (0.384,6.507)rectangle++(0.04,0.04)
  (-1.268,4.304)rectangle++(0.04,0.04) (-0.908,5.305)rectangle++(0.04,0.04) (-0.562,6.142)rectangle++(0.04,0.04) (-0.234,6.84)rectangle++(0.04,0.04) (-1.45,4.254)rectangle++(0.04,0.04)
  (-1.065,5.27)rectangle++(0.04,0.04) (-0.696,6.118)rectangle++(0.04,0.04) (-0.349,6.825)rectangle++(0.04,0.04) (-0.026,7.411)rectangle++(0.04,0.04) (0.028,2.591)rectangle++(0.04,0.04)
  (0.125,3.797)rectangle++(0.04,0.04) (-0.74,3.721)rectangle++(0.04,0.04) (1.008,2.512)rectangle++(0.04,0.04) (0.955,3.691)rectangle++(0.04,0.04) (0.957,4.695)rectangle++(0.04,0.04)
  (0.999,5.549)rectangle++(0.04,0.04) (1.069,6.273)rectangle++(0.04,0.04) (1.158,6.885)rectangle++(0.04,0.04) (1.257,7.402)rectangle++(0.04,0.04) (1.362,7.838)rectangle++(0.04,0.04)
  (1.469,8.204)rectangle++(0.04,0.04) (1.575,8.511)rectangle++(0.04,0.04) (1.677,8.767)rectangle++(0.04,0.04) (1.774,8.981)rectangle++(0.04,0.04) (1.865,9.159)rectangle++(0.04,0.04)
  (1.949,9.307)rectangle++(0.04,0.04) (-1.582,5.15)rectangle++(0.04,0.04) (-1.141,6.037)rectangle++(0.04,0.04) (-1.442,3.926)rectangle++(0.04,0.04) (-1.071,4.991)rectangle++(0.04,0.04)
  (-0.713,5.881)rectangle++(0.04,0.04) (-0.373,6.624)rectangle++(0.04,0.04) (1.684,3.155)rectangle++(0.04,0.04) (1.555,4.21)rectangle++(0.04,0.04) (0.861,3.127)rectangle++(0.04,0.04)
  (0.854,4.22)rectangle++(0.04,0.04) (0.945,2.97)rectangle++(0.04,0.04) (0.92,4.083)rectangle++(0.04,0.04) (0.903,2.958)rectangle++(0.04,0.04) (0.884,4.075)rectangle++(0.04,0.04)
  (-0.663,3.918)rectangle++(0.04,0.04) (-0.409,4.953)rectangle++(0.04,0.04) (-0.152,5.822)rectangle++(0.04,0.04) (0.102,6.551)rectangle++(0.04,0.04) (0.346,7.161)rectangle++(0.04,0.04)
  (0.579,7.669)rectangle++(0.04,0.04) (0.796,8.092)rectangle++(0.04,0.04) (0.998,8.442)rectangle++(0.04,0.04) (1.184,8.732)rectangle++(0.04,0.04) (1.354,8.971)rectangle++(0.04,0.04)
  (1.507,9.168)rectangle++(0.04,0.04) (1.646,9.328)rectangle++(0.04,0.04) (1.77,9.459)rectangle++(0.04,0.04) (1.88,9.566)rectangle++(0.04,0.04) (1.979,9.652)rectangle++(0.04,0.04)
  (2.066,9.721)rectangle++(0.04,0.04) (2.143,9.777)rectangle++(0.04,0.04) (2.21,9.821)rectangle++(0.04,0.04) (2.269,9.855)rectangle++(0.04,0.04) (2.321,9.882)rectangle++(0.04,0.04)
  (-1.638,5.334)rectangle++(0.04,0.04) (-1.181,6.195)rectangle++(0.04,0.04) (1.672,2.889)rectangle++(0.04,0.04) (1.535,3.985)rectangle++(0.04,0.04) (1.462,4.922)rectangle++(0.04,0.04)
  (1.437,5.721)rectangle++(0.04,0.04) (1.448,6.402)rectangle++(0.04,0.04) (-0.175,3.702)rectangle++(0.04,0.04) (-0.003,4.75)rectangle++(0.04,0.04) (0.186,5.634)rectangle++(0.04,0.04)
  (0.381,6.378)rectangle++(0.04,0.04) (0.577,7.002)rectangle++(0.04,0.04) (0.768,7.525)rectangle++(0.04,0.04) (0.952,7.962)rectangle++(0.04,0.04) (1.125,8.325)rectangle++(0.04,0.04)
  (1.287,8.628)rectangle++(0.04,0.04) (1.437,8.878)rectangle++(0.04,0.04) (1.574,9.085)rectangle++(0.04,0.04) (1.699,9.256)rectangle++(0.04,0.04) (1.813,9.396)rectangle++(0.04,0.04)
  (1.914,9.51)rectangle++(0.04,0.04) (2.005,9.603)rectangle++(0.04,0.04) (2.086,9.679)rectangle++(0.04,0.04) (2.158,9.74)rectangle++(0.04,0.04) (2.222,9.788)rectangle++(0.04,0.04)
  (2.278,9.828)rectangle++(0.04,0.04) (2.327,9.858)rectangle++(0.04,0.04) (2.37,9.883)rectangle++(0.04,0.04) (-1.629,5.343)rectangle++(0.04,0.04) (-1.173,6.203)rectangle++(0.04,0.04)
  (-0.751,6.916)rectangle++(0.04,0.04) (-0.364,7.505)rectangle++(0.04,0.04) (-1.614,4.361)rectangle++(0.04,0.04) (-1.2,5.367)rectangle++(0.04,0.04) (-0.807,6.207)rectangle++(0.04,0.04)
  (-0.44,6.904)rectangle++(0.04,0.04) (1.761,3.194)rectangle++(0.04,0.04) (1.623,4.24)rectangle++(0.04,0.04) (1.547,5.136)rectangle++(0.04,0.04) (1.518,5.9)rectangle++(0.04,0.04)
  (1.524,6.55)rectangle++(0.04,0.04) (1.555,7.103)rectangle++(0.04,0.04) (1.604,7.571)rectangle++(0.04,0.04) (1.664,7.968)rectangle++(0.04,0.04) (1.731,8.302)rectangle++(0.04,0.04)
  (1.853,3.956)rectangle++(0.04,0.04) (1.731,4.885)rectangle++(0.04,0.04) (1.007,3.297)rectangle++(0.04,0.04) (0.986,4.358)rectangle++(0.04,0.04) (1.01,5.261)rectangle++(0.04,0.04)
  (1.067,6.028)rectangle++(0.04,0.04) (1.146,6.677)rectangle++(0.04,0.04) (1.239,7.226)rectangle++(0.04,0.04) (1.644,3.637)rectangle++(0.04,0.04) (1.541,4.622)rectangle++(0.04,0.04)
  (1.493,5.463)rectangle++(0.04,0.04) (1.485,6.18)rectangle++(0.04,0.04) (1.507,6.79)rectangle++(0.04,0.04) (1.551,7.307)rectangle++(0.04,0.04) (1.608,7.745)rectangle++(0.04,0.04)
  (1.674,8.115)rectangle++(0.04,0.04) (1.746,8.427)rectangle++(0.04,0.04) (1.819,8.69)rectangle++(0.04,0.04) (1.891,8.91)rectangle++(0.04,0.04) (1.962,9.094)rectangle++(0.04,0.04)
  (2.029,9.247)rectangle++(0.04,0.04) (2.092,9.375)rectangle++(0.04,0.04) (2.151,9.482)rectangle++(0.04,0.04) (2.206,9.569)rectangle++(0.04,0.04) (2.255,9.642)rectangle++(0.04,0.04)
  (2.301,9.702)rectangle++(0.04,0.04) (2.341,9.751)rectangle++(0.04,0.04) (2.378,9.791)rectangle++(0.04,0.04) (2.146,4.39)rectangle++(0.04,0.04) (-0.504,4.217)rectangle++(0.04,0.04)
  (-0.262,5.2)rectangle++(0.04,0.04) (-0.017,6.027)rectangle++(0.04,0.04) (0.225,6.72)rectangle++(0.04,0.04) (0.458,7.299)rectangle++(0.04,0.04) (0.679,7.782)rectangle++(0.04,0.04)
  (-1.486,4.611)rectangle++(0.04,0.04) (-1.081,5.575)rectangle++(0.04,0.04) (-1.333,3.846)rectangle++(0.04,0.04) (-0.982,4.918)rectangle++(0.04,0.04) (-1.176,3.734)rectangle++(0.04,0.04)
  (-0.853,4.818)rectangle++(0.04,0.04) (-0.534,5.725)rectangle++(0.04,0.04) (-0.227,6.484)rectangle++(0.04,0.04) (-1.373,4.185)rectangle++(0.04,0.04) (-1.002,5.208)rectangle++(0.04,0.04);
\fill[piece2]
  (-1.241,3.788)rectangle++(0.04,0.04) (-0.906,4.866)rectangle++(0.04,0.04) (-0.578,5.768)rectangle++(0.04,0.04) (-1.284,3.978)rectangle++(0.04,0.04) (-1.031,3.492)rectangle++(0.04,0.04)
  (-0.739,4.606)rectangle++(0.04,0.04) (-0.446,5.541)rectangle++(0.04,0.04) (1.424,2.929)rectangle++(0.04,0.04) (1.326,4.029)rectangle++(0.04,0.04) (0.006,2.131)rectangle++(0.04,0.04)
  (0.088,3.407)rectangle++(0.04,0.04) (0.209,4.489)rectangle++(0.04,0.04) (0.355,5.403)rectangle++(0.04,0.04) (0.516,6.175)rectangle++(0.04,0.04) (1.469,3.27)rectangle++(0.04,0.04)
  (0.638,2.926)rectangle++(0.04,0.04) (0.657,4.058)rectangle++(0.04,0.04) (0.719,5.019)rectangle++(0.04,0.04) (0.809,5.834)rectangle++(0.04,0.04) (1.35,3.271)rectangle++(0.04,0.04)
  (1.276,4.322)rectangle++(0.04,0.04) (1.255,5.219)rectangle++(0.04,0.04) (1.274,5.982)rectangle++(0.04,0.04) (1.32,6.63)rectangle++(0.04,0.04) (-1.127,4.504)rectangle++(0.04,0.04)
  (-0.78,5.47)rectangle++(0.04,0.04) (-0.446,6.277)rectangle++(0.04,0.04) (-0.131,6.949)rectangle++(0.04,0.04) (0.165,7.508)rectangle++(0.04,0.04) (0.438,7.972)rectangle++(0.04,0.04)
  (0.689,8.355)rectangle++(0.04,0.04) (0.918,8.67)rectangle++(0.04,0.04) (1.125,8.929)rectangle++(0.04,0.04) (1.311,9.141)rectangle++(0.04,0.04) (1.478,9.313)rectangle++(0.04,0.04)
  (2.133,4.093)rectangle++(0.04,0.04) (1.975,4.99)rectangle++(0.04,0.04) (-0.662,4.303)rectangle++(0.04,0.04) (-0.393,5.28)rectangle++(0.04,0.04) (-0.125,6.1)rectangle++(0.04,0.04)
  (0.136,6.786)rectangle++(0.04,0.04) (0.385,7.359)rectangle++(0.04,0.04) (-1.449,4.473)rectangle++(0.04,0.04) (-1.055,5.456)rectangle++(0.04,0.04) (-0.681,6.276)rectangle++(0.04,0.04)
  (-0.33,6.958)rectangle++(0.04,0.04) (-0.004,7.524)rectangle++(0.04,0.04) (0.295,7.991)rectangle++(0.04,0.04) (0.568,8.377)rectangle++(0.04,0.04) (0.816,8.694)rectangle++(0.04,0.04)
  (1.039,8.953)rectangle++(0.04,0.04) (1.239,9.165)rectangle++(0.04,0.04) (1.418,9.337)rectangle++(0.04,0.04) (1.576,9.476)rectangle++(0.04,0.04) (1.717,9.588)rectangle++(0.04,0.04)
  (2.173,4.201)rectangle++(0.04,0.04) (2.013,5.08)rectangle++(0.04,0.04) (1.022,3.399)rectangle++(0.04,0.04) (1.003,4.444)rectangle++(0.04,0.04) (-0.72,4.013)rectangle++(0.04,0.04)
  (-0.454,5.036)rectangle++(0.04,0.04) (-0.186,5.895)rectangle++(0.04,0.04) (0.075,6.614)rectangle++(0.04,0.04) (0.326,7.215)rectangle++(0.04,0.04) (0.564,7.716)rectangle++(0.04,0.04)
  (0.786,8.132)rectangle++(0.04,0.04) (0.991,8.477)rectangle++(0.04,0.04) (1.179,8.762)rectangle++(0.04,0.04) (0.036,2.774)rectangle++(0.04,0.04) (0.14,3.953)rectangle++(0.04,0.04)
  (0.275,4.951)rectangle++(0.04,0.04) (0.429,5.793)rectangle++(0.04,0.04) (0.594,6.503)rectangle++(0.04,0.04) (1.541,3.351)rectangle++(0.04,0.04) (1.442,4.383)rectangle++(0.04,0.04)
  (1.398,5.264)rectangle++(0.04,0.04) (1.397,6.015)rectangle++(0.04,0.04) (1.426,6.653)rectangle++(0.04,0.04) (1.476,7.194)rectangle++(0.04,0.04) (1.54,7.652)rectangle++(0.04,0.04)
  (1.613,8.039)rectangle++(0.04,0.04) (1.69,8.365)rectangle++(0.04,0.04) (1.769,8.639)rectangle++(0.04,0.04) (1.847,8.868)rectangle++(0.04,0.04) (1.923,9.06)rectangle++(0.04,0.04)
  (1.994,9.221)rectangle++(0.04,0.04) (2.062,9.354)rectangle++(0.04,0.04) (2.125,9.465)rectangle++(0.04,0.04) (2.182,9.556)rectangle++(0.04,0.04) (2.111,4.3)rectangle++(0.04,0.04)
  (1.964,5.167)rectangle++(0.04,0.04) (1.874,5.91)rectangle++(0.04,0.04) (1.827,6.544)rectangle++(0.04,0.04) (1.812,7.086)rectangle++(0.04,0.04) (1.822,7.547)rectangle++(0.04,0.04)
  (1.848,7.938)rectangle++(0.04,0.04) (1.886,8.27)rectangle++(0.04,0.04) (-1.373,4.934)rectangle++(0.04,0.04) (-0.972,5.845)rectangle++(0.04,0.04) (-0.595,6.603)rectangle++(0.04,0.04)
  (-0.244,7.233)rectangle++(0.04,0.04) (0.08,7.754)rectangle++(0.04,0.04) (0.376,8.184)rectangle++(0.04,0.04) (0.645,8.537)rectangle++(0.04,0.04) (2.039,3.755)rectangle++(0.04,0.04)
  (1.881,4.706)rectangle++(0.04,0.04) (1.785,5.521)rectangle++(0.04,0.04) (1.736,6.218)rectangle++(0.04,0.04) (1.722,6.812)rectangle++(0.04,0.04) (1.734,7.317)rectangle++(0.04,0.04)
  (1.764,7.747)rectangle++(0.04,0.04) (1.807,8.11)rectangle++(0.04,0.04) (1.858,8.418)rectangle++(0.04,0.04) (1.914,8.677)rectangle++(0.04,0.04) (-1.132,3.37)rectangle++(0.04,0.04)
  (-0.829,4.506)rectangle++(0.04,0.04) (-0.527,5.459)rectangle++(0.04,0.04) (-0.232,6.258)rectangle++(0.04,0.04) (0.051,6.925)rectangle++(0.04,0.04) (-1.417,4.324)rectangle++(0.04,0.04)
  (-1.034,5.328)rectangle++(0.04,0.04) (-1.273,3.806)rectangle++(0.04,0.04) (-0.992,3.46)rectangle++(0.04,0.04) (0.028,3.352)rectangle++(0.04,0.04) (0.155,4.444)rectangle++(0.04,0.04)
  (0.308,5.368)rectangle++(0.04,0.04) (0.474,6.146)rectangle++(0.04,0.04) (0.647,6.802)rectangle++(0.04,0.04) (0.819,7.352)rectangle++(0.04,0.04) (0.988,7.812)rectangle++(0.04,0.04)
  (1.15,8.197)rectangle++(0.04,0.04) (1.304,8.518)rectangle++(0.04,0.04) (1.447,8.784)rectangle++(0.04,0.04) (1.579,9.005)rectangle++(0.04,0.04) (1.7,9.187)rectangle++(0.04,0.04)
  (1.81,9.337)rectangle++(0.04,0.04) (1.25,2.271)rectangle++(0.04,0.04) (1.151,3.477)rectangle++(0.04,0.04) (1.115,4.506)rectangle++(0.04,0.04) (1.126,5.381)rectangle++(0.04,0.04)
  (1.17,6.125)rectangle++(0.04,0.04) (1.237,6.756)rectangle++(0.04,0.04) (1.32,7.289)rectangle++(0.04,0.04) (1.411,7.739)rectangle++(0.04,0.04) (1.507,8.118)rectangle++(0.04,0.04)
  (1.603,8.436)rectangle++(0.04,0.04) (1.698,8.703)rectangle++(0.04,0.04) (1.789,8.926)rectangle++(0.04,0.04) (1.876,9.112)rectangle++(0.04,0.04) (1.957,9.266)rectangle++(0.04,0.04)
  (2.032,9.394)rectangle++(0.04,0.04) (2.1,9.5)rectangle++(0.04,0.04) (2.163,9.587)rectangle++(0.04,0.04) (2.22,9.659)rectangle++(0.04,0.04) (2.271,9.717)rectangle++(0.04,0.04)
  (2.317,9.765)rectangle++(0.04,0.04) (2.358,9.804)rectangle++(0.04,0.04) (2.152,4.388)rectangle++(0.04,0.04) (2.002,5.24)rectangle++(0.04,0.04) (1.909,5.97)rectangle++(0.04,0.04)
  (1.255,3.547)rectangle++(0.04,0.04) (0.731,2.931)rectangle++(0.04,0.04) (-0.448,3.663)rectangle++(0.04,0.04) (-0.814,3.587)rectangle++(0.04,0.04) (-0.551,4.678)rectangle++(0.04,0.04)
  (-0.283,5.594)rectangle++(0.04,0.04) (-0.019,6.363)rectangle++(0.04,0.04) (0.236,7.005)rectangle++(0.04,0.04) (0.479,7.541)rectangle++(0.04,0.04) (0.706,7.987)rectangle++(0.04,0.04)
  (0.918,8.357)rectangle++(0.04,0.04) (1.112,8.663)rectangle++(0.04,0.04) (1.29,8.915)rectangle++(0.04,0.04) (1.451,9.123)rectangle++(0.04,0.04) (1.596,9.292)rectangle++(0.04,0.04)
  (0.04,2.846)rectangle++(0.04,0.04) (0.145,4.014)rectangle++(0.04,0.04) (0.282,5.002)rectangle++(0.04,0.04) (0.437,5.837)rectangle++(0.04,0.04) (0.603,6.54)rectangle++(0.04,0.04)
  (1.549,3.36)rectangle++(0.04,0.04) (1.449,4.39)rectangle++(0.04,0.04) (1.405,5.27)rectangle++(0.04,0.04) (1.403,6.02)rectangle++(0.04,0.04) (1.431,6.657)rectangle++(0.04,0.04)
  (1.481,3.571)rectangle++(0.04,0.04) (1.399,4.572)rectangle++(0.04,0.04) (1.37,5.426)rectangle++(0.04,0.04) (1.379,6.154)rectangle++(0.04,0.04) (1.417,6.772)rectangle++(0.04,0.04)
  (1.473,7.296)rectangle++(0.04,0.04) (1.541,7.739)rectangle++(0.04,0.04) (1.618,8.112)rectangle++(0.04,0.04) (1.697,8.427)rectangle++(0.04,0.04) (1.778,8.691)rectangle++(0.04,0.04)
  (1.856,8.913)rectangle++(0.04,0.04) (1.932,9.098)rectangle++(0.04,0.04) (2.004,9.252)rectangle++(0.04,0.04) (2.056,4.201)rectangle++(0.04,0.04) (1.914,5.085)rectangle++(0.04,0.04)
  (-0.212,3.582)rectangle++(0.04,0.04) (-0.039,4.649)rectangle++(0.04,0.04) (0.151,5.549)rectangle++(0.04,0.04) (0.348,6.307)rectangle++(0.04,0.04) (0.546,6.943)rectangle++(0.04,0.04)
  (0.739,7.476)rectangle++(0.04,0.04) (0.925,7.921)rectangle++(0.04,0.04) (1.101,8.292)rectangle++(0.04,0.04) (1.266,8.601)rectangle++(0.04,0.04) (1.418,8.856)rectangle++(0.04,0.04)
  (2.027,3.992)rectangle++(0.04,0.04) (1.88,4.908)rectangle++(0.04,0.04) (1.792,5.693)rectangle++(0.04,0.04) (1.2,3.467)rectangle++(0.04,0.04) (1.156,4.495)rectangle++(0.04,0.04)
  (1.16,5.371)rectangle++(0.04,0.04) (1.199,6.115)rectangle++(0.04,0.04) (-1.039,4.38)rectangle++(0.04,0.04) (-1.072,3.617)rectangle++(0.04,0.04) (-0.769,4.714)rectangle++(0.04,0.04)
  (-0.467,5.634)rectangle++(0.04,0.04) (-0.174,6.404)rectangle++(0.04,0.04) (1.606,3.158)rectangle++(0.04,0.04) (1.489,4.216)rectangle++(0.04,0.04) (1.432,5.12)rectangle++(0.04,0.04)
  (1.42,5.891)rectangle++(0.04,0.04) (1.441,6.547)rectangle++(0.04,0.04) (1.484,7.103)rectangle++(0.04,0.04) (1.543,7.575)rectangle++(0.04,0.04) (1.695,3.773)rectangle++(0.04,0.04)
  (1.589,4.736)rectangle++(0.04,0.04) (0.987,3.236)rectangle++(0.04,0.04) (0.686,2.81)rectangle++(0.04,0.04) (0.693,3.958)rectangle++(0.04,0.04) (0.745,4.932)rectangle++(0.04,0.04)
  (0.829,5.759)rectangle++(0.04,0.04) (0.933,6.458)rectangle++(0.04,0.04) (1.049,7.048)rectangle++(0.04,0.04) (1.171,7.545)rectangle++(0.04,0.04) (1.295,7.963)rectangle++(0.04,0.04)
  (1.417,8.313)rectangle++(0.04,0.04) (1.535,8.605)rectangle++(0.04,0.04) (1.647,8.849)rectangle++(0.04,0.04) (1.751,9.052)rectangle++(0.04,0.04) (1.849,9.221)rectangle++(0.04,0.04)
  (1.938,9.36)rectangle++(0.04,0.04) (2.019,9.474)rectangle++(0.04,0.04) (2.093,9.569)rectangle++(0.04,0.04) (2.124,4.282)rectangle++(0.04,0.04) (1.975,5.151)rectangle++(0.04,0.04)
  (1.882,5.896)rectangle++(0.04,0.04) (1.24,3.526)rectangle++(0.04,0.04) (1.193,4.544)rectangle++(0.04,0.04) (1.193,5.411)rectangle++(0.04,0.04) (0.015,2.319)rectangle++(0.04,0.04)
  (0.103,3.566)rectangle++(0.04,0.04) (0.228,4.624)rectangle++(0.04,0.04) (0.377,5.517)rectangle++(0.04,0.04) (0.539,6.271)rectangle++(0.04,0.04) (0.706,6.905)rectangle++(0.04,0.04)
  (0.874,7.437)rectangle++(0.04,0.04) (1.039,7.883)rectangle++(0.04,0.04) (1.196,8.255)rectangle++(0.04,0.04) (1.345,8.565)rectangle++(0.04,0.04) (0.447,3.004)rectangle++(0.04,0.04)
  (0.498,4.131)rectangle++(0.04,0.04) (0.586,5.088)rectangle++(0.04,0.04) (0.7,5.898)rectangle++(0.04,0.04) (0.828,6.581)rectangle++(0.04,0.04) (0.965,7.157)rectangle++(0.04,0.04)
  (1.105,7.641)rectangle++(0.04,0.04) (1.763,3.686)rectangle++(0.04,0.04) (1.644,4.659)rectangle++(0.04,0.04) (1.581,5.491)rectangle++(0.04,0.04) (1.561,6.2)rectangle++(0.04,0.04)
  (1.573,6.804)rectangle++(0.04,0.04) (1.607,7.316)rectangle++(0.04,0.04) (1.656,7.751)rectangle++(0.04,0.04) (1.716,8.118)rectangle++(0.04,0.04) (1.781,8.428)rectangle++(0.04,0.04)
  (1.849,8.689)rectangle++(0.04,0.04) (1.917,8.908)rectangle++(0.04,0.04) (1.983,9.091)rectangle++(0.04,0.04) (2.047,9.244)rectangle++(0.04,0.04) (2.108,9.372)rectangle++(0.04,0.04)
  (-1.568,5.193)rectangle++(0.04,0.04) (-1.127,6.073)rectangle++(0.04,0.04) (1.636,2.877)rectangle++(0.04,0.04) (0.521,2.888)rectangle++(0.04,0.04) (0.556,4.031)rectangle++(0.04,0.04)
  (0.933,2.866)rectangle++(0.04,0.04) (0.906,3.995)rectangle++(0.04,0.04) (0.927,4.956)rectangle++(0.04,0.04) (-0.842,4.1)rectangle++(0.04,0.04) (-0.554,5.115)rectangle++(0.04,0.04)
  (1.331,2.823)rectangle++(0.04,0.04) (1.242,3.942)rectangle++(0.04,0.04) (1.212,4.897)rectangle++(0.04,0.04) (1.224,5.711)rectangle++(0.04,0.04) (1.266,6.401)rectangle++(0.04,0.04)
  (1.33,6.987)rectangle++(0.04,0.04) (1.408,7.482)rectangle++(0.04,0.04) (0.423,2.886)rectangle++(0.04,0.04) (0.473,4.032)rectangle++(0.04,0.04) (0.561,5.005)rectangle++(0.04,0.04)
  (0.675,5.828)rectangle++(0.04,0.04) (1.364,3.239)rectangle++(0.04,0.04) (0.001,2.023)rectangle++(0.04,0.04) (0.079,3.316)rectangle++(0.04,0.04) (0.198,4.412)rectangle++(0.04,0.04)
  (0.343,5.338)rectangle++(0.04,0.04) (0.502,6.12)rectangle++(0.04,0.04) (0.67,6.778)rectangle++(0.04,0.04) (0.838,7.331)rectangle++(0.04,0.04) (1.003,7.794)rectangle++(0.04,0.04)
  (1.163,8.181)rectangle++(0.04,0.04) (1.313,8.503)rectangle++(0.04,0.04) (1.454,8.772)rectangle++(0.04,0.04) (1.585,8.994)rectangle++(0.04,0.04) (1.705,9.178)rectangle++(0.04,0.04)
  (1.814,9.329)rectangle++(0.04,0.04) (1.913,9.453)rectangle++(0.04,0.04) (-1.62,5.172)rectangle++(0.04,0.04) (-1.345,3.665)rectangle++(0.04,0.04) (-0.998,4.765)rectangle++(0.04,0.04)
  (-0.66,5.687)rectangle++(0.04,0.04) (-0.336,6.456)rectangle++(0.04,0.04) (-0.03,7.097)rectangle++(0.04,0.04) (0.257,7.63)rectangle++(0.04,0.04) (0.521,8.072)rectangle++(0.04,0.04)
  (0.764,8.436)rectangle++(0.04,0.04) (0.984,8.736)rectangle++(0.04,0.04) (1.184,8.983)rectangle++(0.04,0.04) (1.363,9.184)rectangle++(0.04,0.04) (1.524,9.348)rectangle++(0.04,0.04)
  (1.667,9.481)rectangle++(0.04,0.04) (1.794,9.589)rectangle++(0.04,0.04) (1.906,9.675)rectangle++(0.04,0.04) (2.005,9.743)rectangle++(0.04,0.04) (2.092,9.798)rectangle++(0.04,0.04)
  (2.168,9.841)rectangle++(0.04,0.04) (2.234,9.874)rectangle++(0.04,0.04) (2.292,9.9)rectangle++(0.04,0.04) (2.342,9.919)rectangle++(0.04,0.04) (2.385,9.934)rectangle++(0.04,0.04)
  (2.423,9.945)rectangle++(0.04,0.04) (2.455,9.952)rectangle++(0.04,0.04) (2.482,9.957)rectangle++(0.04,0.04) (2.506,9.961)rectangle++(0.04,0.04) (2.526,9.963)rectangle++(0.04,0.04)
  (2.544,9.963)rectangle++(0.04,0.04) (2.559,9.963)rectangle++(0.04,0.04) (2.571,9.963)rectangle++(0.04,0.04) (2.582,9.962)rectangle++(0.04,0.04) (2.591,9.96)rectangle++(0.04,0.04)
  (2.598,9.959)rectangle++(0.04,0.04) (2.605,9.957)rectangle++(0.04,0.04) (2.61,9.956)rectangle++(0.04,0.04) (2.615,9.954)rectangle++(0.04,0.04) (2.618,9.953)rectangle++(0.04,0.04)
  (2.622,9.951)rectangle++(0.04,0.04) (2.624,9.95)rectangle++(0.04,0.04) (2.626,9.949)rectangle++(0.04,0.04) (2.628,9.947)rectangle++(0.04,0.04) (2.63,9.946)rectangle++(0.04,0.04)
  (-1.596,5.404)rectangle++(0.04,0.04) (-1.143,6.254)rectangle++(0.04,0.04) (-0.723,6.957)rectangle++(0.04,0.04) (-0.339,7.539)rectangle++(0.04,0.04) (0.011,8.018)rectangle++(0.04,0.04)
  (0.328,8.411)rectangle++(0.04,0.04) (-1.683,4.67)rectangle++(0.04,0.04) (-1.249,3.554)rectangle++(0.04,0.04) (-0.922,4.667)rectangle++(0.04,0.04) (-0.599,5.6)rectangle++(0.04,0.04)
  (-0.287,6.38)rectangle++(0.04,0.04) (0.009,7.031)rectangle++(0.04,0.04) (0.286,7.572)rectangle++(0.04,0.04) (0.544,8.021)rectangle++(0.04,0.04) (0.781,8.392)rectangle++(0.04,0.04)
  (0.998,8.698)rectangle++(0.04,0.04) (1.194,8.95)rectangle++(0.04,0.04) (1.37,9.156)rectangle++(0.04,0.04) (1.529,9.324)rectangle++(0.04,0.04) (1.67,9.46)rectangle++(0.04,0.04)
  (0.498,3.107)rectangle++(0.04,0.04) (0.545,4.217)rectangle++(0.04,0.04) (0.63,5.159)rectangle++(0.04,0.04) (1.204,3.1)rectangle++(0.04,0.04) (1.145,4.183)rectangle++(0.04,0.04)
  (1.138,5.106)rectangle++(0.04,0.04) (1.17,5.891)rectangle++(0.04,0.04) (1.228,6.557)rectangle++(0.04,0.04) (1.304,7.12)rectangle++(0.04,0.04) (1.391,7.596)rectangle++(0.04,0.04)
  (1.484,7.997)rectangle++(0.04,0.04) (0.031,2.67)rectangle++(0.04,0.04) (-0.518,3.481)rectangle++(0.04,0.04) (-0.303,4.576)rectangle++(0.04,0.04) (-0.077,5.497)rectangle++(0.04,0.04)
  (0.152,6.272)rectangle++(0.04,0.04) (0.378,6.921)rectangle++(0.04,0.04) (0.596,7.464)rectangle++(0.04,0.04) (0.803,7.917)rectangle++(0.04,0.04) (0.997,8.294)rectangle++(0.04,0.04)
  (1.177,8.606)rectangle++(0.04,0.04) (1.343,8.864)rectangle++(0.04,0.04) (-1.597,4.949)rectangle++(0.04,0.04) (-1.162,5.867)rectangle++(0.04,0.04) (0.17,2.667)rectangle++(0.04,0.04)
  (-0.492,3.506)rectangle++(0.04,0.04) (-0.28,4.596)rectangle++(0.04,0.04) (-0.057,5.514)rectangle++(0.04,0.04) (0.17,6.285)rectangle++(0.04,0.04) (0.394,6.932)rectangle++(0.04,0.04)
  (0.61,7.473)rectangle++(0.04,0.04) (0.815,7.924)rectangle++(0.04,0.04) (1.867,3.675)rectangle++(0.04,0.04) (1.731,4.645)rectangle++(0.04,0.04) (1.655,5.475)rectangle++(0.04,0.04)
  (1.624,6.184)rectangle++(0.04,0.04) (1.625,6.788)rectangle++(0.04,0.04) (1.651,7.301)rectangle++(0.04,0.04) (1.693,7.736)rectangle++(0.04,0.04) (1.746,8.104)rectangle++(0.04,0.04)
  (1.806,8.415)rectangle++(0.04,0.04) (1.87,8.676)rectangle++(0.04,0.04) (-1.463,5.011)rectangle++(0.04,0.04) (-1.045,5.914)rectangle++(0.04,0.04) (-0.654,6.665)rectangle++(0.04,0.04)
  (-0.291,7.288)rectangle++(0.04,0.04) (0.042,7.802)rectangle++(0.04,0.04) (0.345,8.226)rectangle++(0.04,0.04) (0.62,8.575)rectangle++(0.04,0.04) (0.868,8.86)rectangle++(0.04,0.04)
  (-1.681,4.859)rectangle++(0.04,0.04) (-1.237,5.794)rectangle++(0.04,0.04) (-0.822,6.571)rectangle++(0.04,0.04) (-0.438,7.214)rectangle++(0.04,0.04) (-0.086,7.746)rectangle++(0.04,0.04)
  (0.235,8.183)rectangle++(0.04,0.04) (0.525,8.543)rectangle++(0.04,0.04) (0.785,8.837)rectangle++(0.04,0.04) (1.019,9.076)rectangle++(0.04,0.04) (-1.7,4.93)rectangle++(0.04,0.04)
  (-1.25,5.855)rectangle++(0.04,0.04) (-0.83,6.623)rectangle++(0.04,0.04) (-0.443,7.259)rectangle++(0.04,0.04) (-0.089,7.784)rectangle++(0.04,0.04) (0.234,8.216)rectangle++(0.04,0.04)
  (0.525,8.571)rectangle++(0.04,0.04) (0.787,8.86)rectangle++(0.04,0.04) (1.021,9.096)rectangle++(0.04,0.04) (1.23,9.287)rectangle++(0.04,0.04) (1.415,9.441)rectangle++(0.04,0.04)
  (1.578,9.564)rectangle++(0.04,0.04) (1.721,9.663)rectangle++(0.04,0.04) (1.848,9.741)rectangle++(0.04,0.04) (1.958,9.802)rectangle++(0.04,0.04) (2.198,4.297)rectangle++(0.04,0.04)
  (2.038,5.16)rectangle++(0.04,0.04) (1.936,5.901)rectangle++(0.04,0.04) (1.235,3.54)rectangle++(0.04,0.04) (1.189,4.556)rectangle++(0.04,0.04) (1.19,5.421)rectangle++(0.04,0.04)
  (1.227,6.156)rectangle++(0.04,0.04) (1.287,6.78)rectangle++(0.04,0.04) (1.528,3.562)rectangle++(0.04,0.04) (1.439,4.562)rectangle++(0.04,0.04) (1.404,5.417)rectangle++(0.04,0.04)
  (1.408,6.144)rectangle++(0.04,0.04) (1.44,6.763)rectangle++(0.04,0.04) (0.407,2.807)rectangle++(0.04,0.04) (0.456,3.966)rectangle++(0.04,0.04) (0.929,2.831)rectangle++(0.04,0.04)
  (0.901,3.965)rectangle++(0.04,0.04) (0.876,2.931)rectangle++(0.04,0.04) (0.86,4.053)rectangle++(0.04,0.04) (0.891,5.007)rectangle++(0.04,0.04) (0.955,5.816)rectangle++(0.04,0.04)
  (1.043,6.502)rectangle++(0.04,0.04) (1.144,7.081)rectangle++(0.04,0.04) (1.253,7.569)rectangle++(0.04,0.04) (1.366,7.98)rectangle++(0.04,0.04) (1.478,8.325)rectangle++(0.04,0.04)
  (1.587,8.613)rectangle++(0.04,0.04) (1.691,8.854)rectangle++(0.04,0.04) (1.79,9.054)rectangle++(0.04,0.04) (1.881,9.221)rectangle++(0.04,0.04) (1.242,2.28)rectangle++(0.04,0.04)
  (1.145,3.484)rectangle++(0.04,0.04) (1.11,4.512)rectangle++(0.04,0.04) (1.122,5.387)rectangle++(0.04,0.04) (1.167,6.13)rectangle++(0.04,0.04) (1.381,3.409)rectangle++(0.04,0.04)
  (1.308,4.439)rectangle++(0.04,0.04) (1.287,5.317)rectangle++(0.04,0.04) (1.305,6.064)rectangle++(0.04,0.04) (1.349,6.698)rectangle++(0.04,0.04) (1.413,7.236)rectangle++(0.04,0.04)
  (-1.239,4.641)rectangle++(0.04,0.04) (1.294,2.576)rectangle++(0.04,0.04) (1.201,3.734)rectangle++(0.04,0.04) (1.168,4.722)rectangle++(0.04,0.04) (1.18,5.563)rectangle++(0.04,0.04)
  (1.223,6.278)rectangle++(0.04,0.04) (1.288,6.883)rectangle++(0.04,0.04) (1.368,7.396)rectangle++(0.04,0.04) (1.457,7.828)rectangle++(0.04,0.04) (1.768,3.802)rectangle++(0.04,0.04)
  (1.652,4.757)rectangle++(0.04,0.04) (1.593,5.574)rectangle++(0.04,0.04) (1.575,6.27)rectangle++(0.04,0.04) (-1.005,4.481)rectangle++(0.04,0.04) (-0.677,5.445)rectangle++(0.04,0.04)
  (-0.36,6.252)rectangle++(0.04,0.04) (-0.058,6.924)rectangle++(0.04,0.04) (0.225,7.484)rectangle++(0.04,0.04) (0.489,7.949)rectangle++(0.04,0.04) (0.731,8.333)rectangle++(0.04,0.04)
  (0.953,8.65)rectangle++(0.04,0.04) (1.154,8.911)rectangle++(0.04,0.04) (1.335,9.124)rectangle++(0.04,0.04) (1.497,9.298)rectangle++(0.04,0.04) (1.642,9.44)rectangle++(0.04,0.04)
  (1.771,9.554)rectangle++(0.04,0.04) (1.886,9.647)rectangle++(0.04,0.04) (1.987,9.72)rectangle++(0.04,0.04) (2.075,9.779)rectangle++(0.04,0.04) (2.153,9.825)rectangle++(0.04,0.04)
  (2.221,9.862)rectangle++(0.04,0.04) (2.28,9.89)rectangle++(0.04,0.04) (2.331,9.911)rectangle++(0.04,0.04) (2.376,9.928)rectangle++(0.04,0.04) (2.414,9.94)rectangle++(0.04,0.04)
  (2.448,9.948)rectangle++(0.04,0.04) (2.476,9.954)rectangle++(0.04,0.04) (2.501,9.958)rectangle++(0.04,0.04) (2.522,9.961)rectangle++(0.04,0.04) (-1.618,5.387)rectangle++(0.04,0.04)
  (-1.162,6.24)rectangle++(0.04,0.04) (-0.741,6.946)rectangle++(0.04,0.04) (-0.354,7.53)rectangle++(0.04,0.04) (-0.002,8.011)rectangle++(0.04,0.04) (0.317,8.406)rectangle++(0.04,0.04)
  (0.603,8.728)rectangle++(0.04,0.04) (0.86,8.991)rectangle++(0.04,0.04) (-1.711,4.884)rectangle++(0.04,0.04) (0.086,3.477)rectangle++(0.04,0.04) (0.21,4.549)rectangle++(0.04,0.04)
  (0.358,5.454)rectangle++(0.04,0.04) (0.52,6.218)rectangle++(0.04,0.04) (-1.183,4.273)rectangle++(0.04,0.04) (-0.837,5.275)rectangle++(0.04,0.04) (-0.502,6.114)rectangle++(0.04,0.04)
  (-0.185,6.813)rectangle++(0.04,0.04) (0.024,2.509)rectangle++(0.04,0.04) (0.118,3.728)rectangle++(0.04,0.04) (0.248,4.76)rectangle++(0.04,0.04) (0.399,5.633)rectangle++(0.04,0.04)
  (1.348,3.139)rectangle++(0.04,0.04) (1.269,4.21)rectangle++(0.04,0.04) (-0.623,4.017)rectangle++(0.04,0.04) (-0.371,5.035)rectangle++(0.04,0.04) (-0.116,5.891)rectangle++(0.04,0.04)
  (0.135,6.608)rectangle++(0.04,0.04) (1.621,3.267)rectangle++(0.04,0.04) (1.506,4.308)rectangle++(0.04,0.04) (1.45,5.198)rectangle++(0.04,0.04) (1.438,5.957)rectangle++(0.04,0.04)
  (1.459,6.602)rectangle++(0.04,0.04) (1.502,7.149)rectangle++(0.04,0.04) (1.56,7.613)rectangle++(0.04,0.04) (1.629,8.005)rectangle++(0.04,0.04) (1.702,8.335)rectangle++(0.04,0.04)
  (1.865,3.956)rectangle++(0.04,0.04) (1.741,4.884)rectangle++(0.04,0.04) (1.673,5.678)rectangle++(0.04,0.04) (1.647,6.356)rectangle++(0.04,0.04) (1.652,6.933)rectangle++(0.04,0.04)
  (1.679,7.423)rectangle++(0.04,0.04) (1.722,7.838)rectangle++(0.04,0.04) (1.775,8.19)rectangle++(0.04,0.04) (1.834,8.487)rectangle++(0.04,0.04) (1.887,4.015)rectangle++(0.04,0.04)
  (1.762,4.933)rectangle++(0.04,0.04) (1.693,5.719)rectangle++(0.04,0.04) (1.666,6.39)rectangle++(0.04,0.04) (1.669,6.961)rectangle++(0.04,0.04) (1.695,7.446)rectangle++(0.04,0.04)
  (1.736,7.858)rectangle++(0.04,0.04) (1.788,8.206)rectangle++(0.04,0.04) (1.846,8.5)rectangle++(0.04,0.04) (1.907,8.747)rectangle++(0.04,0.04) (1.943,4.082)rectangle++(0.04,0.04)
  (1.812,4.988)rectangle++(0.04,0.04) (1.738,5.763)rectangle++(0.04,0.04) (1.705,6.426)rectangle++(0.04,0.04) (1.704,6.99)rectangle++(0.04,0.04) (1.726,7.469)rectangle++(0.04,0.04)
  (1.764,7.876)rectangle++(0.04,0.04) (1.812,8.22)rectangle++(0.04,0.04) (-0.206,3.892)rectangle++(0.04,0.04) (-0.022,4.913)rectangle++(0.04,0.04) (0.176,5.773)rectangle++(0.04,0.04)
  (0.378,6.496)rectangle++(0.04,0.04) (-0.087,3.687)rectangle++(0.04,0.04) (0.071,4.734)rectangle++(0.04,0.04) (0.248,5.617)rectangle++(0.04,0.04) (0.433,6.361)rectangle++(0.04,0.04)
  (0.62,6.985)rectangle++(0.04,0.04) (0.805,7.509)rectangle++(0.04,0.04) (0.982,7.947)rectangle++(0.04,0.04) (1.15,8.312)rectangle++(0.04,0.04) (1.308,8.615)rectangle++(0.04,0.04)
  (-1.55,4.894)rectangle++(0.04,0.04) (-1.273,3.623)rectangle++(0.04,0.04) (1.046,2.372)rectangle++(0.04,0.04) (0.982,3.571)rectangle++(0.04,0.04) (0.975,4.592)rectangle++(0.04,0.04)
  (-0.756,4.037)rectangle++(0.04,0.04) (-0.484,5.058)rectangle++(0.04,0.04) (-1.117,3.855)rectangle++(0.04,0.04) (-0.797,4.918)rectangle++(0.04,0.04) (-0.483,5.808)rectangle++(0.04,0.04)
  (-0.18,6.553)rectangle++(0.04,0.04) (0.107,7.173)rectangle++(0.04,0.04) (0.375,7.689)rectangle++(0.04,0.04) (0.624,8.117)rectangle++(0.04,0.04) (-1.567,4.667)rectangle++(0.04,0.04)
  (-1.148,5.626)rectangle++(0.04,0.04) (-0.753,6.424)rectangle++(0.04,0.04) (-0.385,7.087)rectangle++(0.04,0.04) (1.8,3.241)rectangle++(0.04,0.04) (1.657,4.279)rectangle++(0.04,0.04)
  (1.578,5.167)rectangle++(0.04,0.04) (1.546,5.925)rectangle++(0.04,0.04) (1.549,6.571)rectangle++(0.04,0.04) (1.577,7.12)rectangle++(0.04,0.04) (1.623,7.585)rectangle++(0.04,0.04)
  (1.681,7.978)rectangle++(0.04,0.04) (1.746,8.311)rectangle++(0.04,0.04) (1.814,8.59)rectangle++(0.04,0.04) (1.915,4.03)rectangle++(0.04,0.04) (-0.469,4.102)rectangle++(0.04,0.04)
  (-0.237,5.102)rectangle++(0.04,0.04) (1.291,2.892)rectangle++(0.04,0.04) (1.211,4.003)rectangle++(0.04,0.04) (1.187,4.95)rectangle++(0.04,0.04) (1.205,5.757)rectangle++(0.04,0.04)
  (1.252,6.441)rectangle++(0.04,0.04) (1.32,7.021)rectangle++(0.04,0.04) (1.401,7.511)rectangle++(0.04,0.04) (1.489,7.925)rectangle++(0.04,0.04) (1.58,8.273)rectangle++(0.04,0.04)
  (1.672,8.565)rectangle++(0.04,0.04) (1.761,8.809)rectangle++(0.04,0.04) (1.847,9.014)rectangle++(0.04,0.04) (1.929,9.184)rectangle++(0.04,0.04) (2.005,9.325)rectangle++(0.04,0.04)
  (2.075,9.443)rectangle++(0.04,0.04) (2.139,9.539)rectangle++(0.04,0.04) (2.198,9.619)rectangle++(0.04,0.04) (2.125,4.316)rectangle++(0.04,0.04) (1.976,5.18)rectangle++(0.04,0.04)
  (1.885,5.92)rectangle++(0.04,0.04) (1.837,6.553)rectangle++(0.04,0.04) (1.821,7.093)rectangle++(0.04,0.04) (1.829,7.552)rectangle++(0.04,0.04) (1.656,3.834)rectangle++(0.04,0.04)
  (1.559,4.789)rectangle++(0.04,0.04) (1.514,5.604)rectangle++(0.04,0.04) (1.509,6.299)rectangle++(0.04,0.04) (1.532,6.89)rectangle++(0.04,0.04) (1.576,7.392)rectangle++(0.04,0.04)
  (1.633,7.816)rectangle++(0.04,0.04) (1.699,8.174)rectangle++(0.04,0.04) (1.769,8.477)rectangle++(0.04,0.04) (-1.438,4.953)rectangle++(0.04,0.04) (-1.027,5.863)rectangle++(0.04,0.04)
  (-0.64,6.621)rectangle++(0.04,0.04) (-0.282,7.25)rectangle++(0.04,0.04) (0.048,7.77)rectangle++(0.04,0.04) (0.35,8.199)rectangle++(0.04,0.04) (0.623,8.551)rectangle++(0.04,0.04)
  (0.869,8.84)rectangle++(0.04,0.04) (1.09,9.075)rectangle++(0.04,0.04) (1.288,9.266)rectangle++(0.04,0.04) (1.463,9.421)rectangle++(0.04,0.04) (1.618,9.546)rectangle++(0.04,0.04)
  (1.755,9.645)rectangle++(0.04,0.04) (1.875,9.725)rectangle++(0.04,0.04) (1.981,9.787)rectangle++(0.04,0.04) (2.073,9.836)rectangle++(0.04,0.04) (2.153,9.874)rectangle++(0.04,0.04)
  (2.223,9.903)rectangle++(0.04,0.04) (2.284,9.925)rectangle++(0.04,0.04) (2.336,9.941)rectangle++(0.04,0.04) (2.381,9.952)rectangle++(0.04,0.04) (2.42,9.961)rectangle++(0.04,0.04)
  (2.453,9.966)rectangle++(0.04,0.04) (-1.632,5.375)rectangle++(0.04,0.04) (-1.174,6.23)rectangle++(0.04,0.04) (-0.751,6.939)rectangle++(0.04,0.04) (1.806,3.13)rectangle++(0.04,0.04)
  (1.658,4.184)rectangle++(0.04,0.04) (1.575,5.086)rectangle++(0.04,0.04) (1.54,5.857)rectangle++(0.04,0.04) (1.541,6.513)rectangle++(0.04,0.04) (1.568,7.07)rectangle++(0.04,0.04)
  (1.613,7.543)rectangle++(0.04,0.04) (1.671,7.944)rectangle++(0.04,0.04) (1.736,8.281)rectangle++(0.04,0.04) (1.805,8.566)rectangle++(0.04,0.04) (1.874,8.805)rectangle++(0.04,0.04)
  (1.943,9.006)rectangle++(0.04,0.04) (2.01,9.173)rectangle++(0.04,0.04) (2.073,9.313)rectangle++(0.04,0.04) (-1.561,5.175)rectangle++(0.04,0.04) (-1.122,6.057)rectangle++(0.04,0.04)
  (-0.714,6.79)rectangle++(0.04,0.04) (-0.337,7.396)rectangle++(0.04,0.04) (0.007,7.896)rectangle++(0.04,0.04) (0.319,8.308)rectangle++(0.04,0.04) (0.602,8.645)rectangle++(0.04,0.04)
  (0.855,8.92)rectangle++(0.04,0.04) (1.081,9.144)rectangle++(0.04,0.04) (1.283,9.326)rectangle++(0.04,0.04) (1.461,9.472)rectangle++(0.04,0.04) (1.619,9.589)rectangle++(0.04,0.04)
  (1.757,9.682)rectangle++(0.04,0.04) (1.879,9.755)rectangle++(0.04,0.04) (1.985,9.813)rectangle++(0.04,0.04) (2.077,9.858)rectangle++(0.04,0.04) (2.198,4.335)rectangle++(0.04,0.04)
  (0.816,3.297)rectangle++(0.04,0.04) (-0.51,3.752)rectangle++(0.04,0.04) (-0.286,4.805)rectangle++(0.04,0.04) (-1.028,3.843)rectangle++(0.04,0.04) (-0.722,4.904)rectangle++(0.04,0.04)
  (-0.42,5.793)rectangle++(0.04,0.04) (-0.127,6.537)rectangle++(0.04,0.04) (0.151,7.158)rectangle++(0.04,0.04) (0.413,7.674)rectangle++(0.04,0.04) (0.655,8.103)rectangle++(0.04,0.04)
  (0.879,8.457)rectangle++(0.04,0.04) (1.083,8.75)rectangle++(0.04,0.04) (1.269,8.99)rectangle++(0.04,0.04) (1.436,9.187)rectangle++(0.04,0.04) (0.039,2.832)rectangle++(0.04,0.04)
  (0.144,4.002)rectangle++(0.04,0.04) (0.28,4.992)rectangle++(0.04,0.04) (0.436,5.828)rectangle++(0.04,0.04) (0.601,6.533)rectangle++(0.04,0.04) (0.77,7.125)rectangle++(0.04,0.04)
  (0.938,7.622)rectangle++(0.04,0.04) (1.1,8.037)rectangle++(0.04,0.04) (-1.465,4.74)rectangle++(0.04,0.04) (-1.058,5.684)rectangle++(0.04,0.04) (-0.674,6.47)rectangle++(0.04,0.04)
  (-0.316,7.123)rectangle++(0.04,0.04) (0.014,7.663)rectangle++(0.04,0.04) (0.316,8.109)rectangle++(0.04,0.04) (0.591,8.476)rectangle++(0.04,0.04) (0.839,8.778)rectangle++(0.04,0.04)
  (1.062,9.024)rectangle++(0.04,0.04) (1.262,9.224)rectangle++(0.04,0.04) (1.439,9.386)rectangle++(0.04,0.04) (1.596,9.517)rectangle++(0.04,0.04) (1.735,9.622)rectangle++(0.04,0.04)
  (1.858,9.705)rectangle++(0.04,0.04) (2.186,4.255)rectangle++(0.04,0.04) (2.026,5.126)rectangle++(0.04,0.04) (1.925,5.872)rectangle++(0.04,0.04) (1.869,6.51)rectangle++(0.04,0.04)
  (1.847,7.055)rectangle++(0.04,0.04) (1.85,7.519)rectangle++(0.04,0.04) (1.871,7.914)rectangle++(0.04,0.04) (1.905,8.248)rectangle++(0.04,0.04) (1.947,8.531)rectangle++(0.04,0.04)
  (1.186,2.297)rectangle++(0.04,0.04) (1.098,3.501)rectangle++(0.04,0.04) (1.071,4.528)rectangle++(0.04,0.04) (1.089,5.402)rectangle++(0.04,0.04) (1.14,6.145)rectangle++(0.04,0.04)
  (1.212,6.774)rectangle++(0.04,0.04) (1.299,7.305)rectangle++(0.04,0.04) (1.394,7.754)rectangle++(0.04,0.04) (1.493,8.131)rectangle++(0.04,0.04) (1.592,8.448)rectangle++(0.04,0.04)
  (1.689,8.713)rectangle++(0.04,0.04) (1.782,8.935)rectangle++(0.04,0.04) (2.004,4.09)rectangle++(0.04,0.04) (1.865,4.992)rectangle++(0.04,0.04) (1.782,5.765)rectangle++(0.04,0.04)
  (1.743,6.425)rectangle++(0.04,0.04) (1.737,6.988)rectangle++(0.04,0.04) (1.754,7.466)rectangle++(0.04,0.04) (1.787,7.872)rectangle++(0.04,0.04) (-1.307,4.837)rectangle++(0.04,0.04)
  (-0.92,5.76)rectangle++(0.04,0.04) (-0.554,6.529)rectangle++(0.04,0.04) (-0.212,7.168)rectangle++(0.04,0.04) (-1.516,4.323)rectangle++(0.04,0.04) (-1.118,5.331)rectangle++(0.04,0.04)
  (-0.739,6.173)rectangle++(0.04,0.04) (-0.383,6.872)rectangle++(0.04,0.04) (-0.053,7.453)rectangle++(0.04,0.04) (0.251,7.933)rectangle++(0.04,0.04) (0.528,8.33)rectangle++(0.04,0.04)
  (0.78,8.655)rectangle++(0.04,0.04) (1.007,8.922)rectangle++(0.04,0.04) (1.211,9.14)rectangle++(0.04,0.04) (2.121,3.999)rectangle++(0.04,0.04) (-0.425,4.135)rectangle++(0.04,0.04)
  (-0.198,5.128)rectangle++(0.04,0.04) (0.034,5.963)rectangle++(0.04,0.04) (0.266,6.663)rectangle++(0.04,0.04) (0.49,7.249)rectangle++(0.04,0.04) (0.704,7.739)rectangle++(0.04,0.04)
  (0.906,8.146)rectangle++(0.04,0.04) (1.094,8.484)rectangle++(0.04,0.04) (1.267,8.764)rectangle++(0.04,0.04) (1.425,8.995)rectangle++(0.04,0.04) (1.569,9.185)rectangle++(0.04,0.04)
  (1.699,9.34)rectangle++(0.04,0.04) (1.815,9.468)rectangle++(0.04,0.04) (-1.64,5.157)rectangle++(0.04,0.04) (-1.345,3.658)rectangle++(0.04,0.04) (-0.999,4.759)rectangle++(0.04,0.04)
  (-0.661,5.682)rectangle++(0.04,0.04) (-0.337,6.452)rectangle++(0.04,0.04) (-0.03,7.094)rectangle++(0.04,0.04) (0.256,7.627)rectangle++(0.04,0.04) (0.52,8.069)rectangle++(0.04,0.04)
  (0.763,8.434)rectangle++(0.04,0.04) (0.984,8.735)rectangle++(0.04,0.04) (1.183,8.981)rectangle++(0.04,0.04) (1.363,9.183)rectangle++(0.04,0.04) (1.523,9.347)rectangle++(0.04,0.04)
  (1.667,9.48)rectangle++(0.04,0.04) (1.794,9.588)rectangle++(0.04,0.04) (2.164,4.218)rectangle++(0.04,0.04) (2.006,5.095)rectangle++(0.04,0.04) (1.907,5.847)rectangle++(0.04,0.04)
  (1.852,6.49)rectangle++(0.04,0.04) (1.832,7.038)rectangle++(0.04,0.04) (1.836,7.505)rectangle++(0.04,0.04) (1.859,7.902)rectangle++(0.04,0.04) (1.739,3.908)rectangle++(0.04,0.04)
  (1.632,4.848)rectangle++(0.04,0.04) (1.579,5.652)rectangle++(0.04,0.04) (1.566,6.337)rectangle++(0.04,0.04) (1.583,6.92)rectangle++(0.04,0.04) (1.62,7.415)rectangle++(0.04,0.04)
  (1.671,7.834)rectangle++(0.04,0.04) (1.732,8.188)rectangle++(0.04,0.04) (1.825,3.934)rectangle++(0.04,0.04) (1.706,4.867)rectangle++(0.04,0.04) (1.643,5.665)rectangle++(0.04,0.04)
  (1.621,6.346)rectangle++(0.04,0.04) (1.629,6.925)rectangle++(0.04,0.04) (1.66,7.418)rectangle++(0.04,0.04) (1.642,3.769)rectangle++(0.04,0.04) (1.545,4.734)rectangle++(0.04,0.04)
  (1.5,5.559)rectangle++(0.04,0.04) (1.495,6.261)rectangle++(0.04,0.04) (1.519,6.858)rectangle++(0.04,0.04) (1.563,7.365)rectangle++(0.04,0.04) (1.621,7.794)rectangle++(0.04,0.04)
  (1.688,8.156)rectangle++(0.04,0.04) (1.759,8.461)rectangle++(0.04,0.04) (1.831,8.718)rectangle++(0.04,0.04) (1.903,8.933)rectangle++(0.04,0.04) (1.973,9.113)rectangle++(0.04,0.04)
  (2.039,9.264)rectangle++(0.04,0.04) (2.102,9.389)rectangle++(0.04,0.04) (2.16,9.493)rectangle++(0.04,0.04) (2.213,9.579)rectangle++(0.04,0.04) (-1.593,5.253)rectangle++(0.04,0.04)
  (-1.146,6.125)rectangle++(0.04,0.04) (-1.462,3.942)rectangle++(0.04,0.04) (-1.087,5.005)rectangle++(0.04,0.04) (-0.726,5.894)rectangle++(0.04,0.04) (-0.384,6.635)rectangle++(0.04,0.04)
  (-0.063,7.252)rectangle++(0.04,0.04) (-1.507,4.367)rectangle++(0.04,0.04) (-1.108,5.369)rectangle++(0.04,0.04) (-0.729,6.204)rectangle++(0.04,0.04) (-1.404,4.036)rectangle++(0.04,0.04)
  (-1.034,5.083)rectangle++(0.04,0.04) (0.793,2.087)rectangle++(0.04,0.04) (0.755,3.338)rectangle++(0.04,0.04) (0.773,4.403)rectangle++(0.04,0.04) (0.831,5.308)rectangle++(0.04,0.04)
  (0.917,6.075)rectangle++(0.04,0.04) (1.02,6.723)rectangle++(0.04,0.04) (1.134,7.27)rectangle++(0.04,0.04) (1.252,7.73)rectangle++(0.04,0.04) (1.371,8.117)rectangle++(0.04,0.04)
  (1.488,8.441)rectangle++(0.04,0.04) (-1.481,4.893)rectangle++(0.04,0.04) (-1.065,5.814)rectangle++(0.04,0.04) (-0.675,6.581)rectangle++(0.04,0.04) (-0.313,7.217)rectangle++(0.04,0.04)
  (0.021,7.743)rectangle++(0.04,0.04) (0.325,8.177)rectangle++(0.04,0.04) (0.601,8.534)rectangle++(0.04,0.04) (-1.66,4.745)rectangle++(0.04,0.04) (-1.223,5.696)rectangle++(0.04,0.04)
  (-0.814,6.487)rectangle++(0.04,0.04) (-0.435,7.142)rectangle++(0.04,0.04) (-0.086,7.685)rectangle++(0.04,0.04) (0.232,8.132)rectangle++(0.04,0.04) (0.52,8.499)rectangle++(0.04,0.04)
  (0.78,8.799)rectangle++(0.04,0.04) (2.088,3.836)rectangle++(0.04,0.04) (1.926,4.773)rectangle++(0.04,0.04) (1.826,5.576)rectangle++(0.04,0.04) (1.167,3.452)rectangle++(0.04,0.04)
  (1.128,4.484)rectangle++(0.04,0.04) (1.136,5.362)rectangle++(0.04,0.04) (1.178,6.109)rectangle++(0.04,0.04) (1.243,6.741)rectangle++(0.04,0.04) (-1.164,4.512)rectangle++(0.04,0.04)
  (-0.811,5.478)rectangle++(0.04,0.04) (-0.473,6.285)rectangle++(0.04,0.04) (-0.153,6.957)rectangle++(0.04,0.04) (0.146,7.516)rectangle++(0.04,0.04) (0.423,7.979)rectangle++(0.04,0.04)
  (0.676,8.361)rectangle++(0.04,0.04) (0.907,8.676)rectangle++(0.04,0.04) (1.116,8.935)rectangle++(0.04,0.04) (1.304,9.146)rectangle++(0.04,0.04) (1.472,9.318)rectangle++(0.04,0.04)
  (1.622,9.458)rectangle++(0.04,0.04) (1.754,9.57)rectangle++(0.04,0.04) (1.872,9.661)rectangle++(0.04,0.04) (2.173,4.25)rectangle++(0.04,0.04) (0.798,3.275)rectangle++(0.04,0.04)
  (0.807,4.348)rectangle++(0.04,0.04) (0.858,5.26)rectangle++(0.04,0.04) (0.937,6.033)rectangle++(0.04,0.04) (1.036,6.687)rectangle++(0.04,0.04) (1.146,7.238)rectangle++(0.04,0.04)
  (1.261,7.703)rectangle++(0.04,0.04) (-1.365,4.705)rectangle++(0.04,0.04) (-0.974,5.65)rectangle++(0.04,0.04) (-0.604,6.438)rectangle++(0.04,0.04) (-1.431,4.105)rectangle++(0.04,0.04)
  (-1.054,5.143)rectangle++(0.04,0.04) (-0.693,6.01)rectangle++(0.04,0.04) (-0.351,6.732)rectangle++(0.04,0.04) (-0.031,7.333)rectangle++(0.04,0.04) (0.265,7.83)rectangle++(0.04,0.04)
  (0.536,8.241)rectangle++(0.04,0.04) (0.783,8.58)rectangle++(0.04,0.04) (1.007,8.858)rectangle++(0.04,0.04) (1.208,9.085)rectangle++(0.04,0.04) (1.388,9.27)rectangle++(0.04,0.04)
  (1.548,9.42)rectangle++(0.04,0.04) (1.691,9.542)rectangle++(0.04,0.04) (1.816,9.639)rectangle++(0.04,0.04) (1.927,9.717)rectangle++(0.04,0.04) (2.025,9.778)rectangle++(0.04,0.04)
  (2.11,9.827)rectangle++(0.04,0.04) (2.184,9.864)rectangle++(0.04,0.04) (2.249,9.894)rectangle++(0.04,0.04) (2.305,9.916)rectangle++(0.04,0.04) (2.354,9.932)rectangle++(0.04,0.04)
  (2.396,9.945)rectangle++(0.04,0.04) (2.432,9.953)rectangle++(0.04,0.04) (2.463,9.959)rectangle++(0.04,0.04) (2.49,9.963)rectangle++(0.04,0.04) (2.175,4.448)rectangle++(0.04,0.04)
  (2.025,5.29)rectangle++(0.04,0.04) (1.931,6.012)rectangle++(0.04,0.04) (1.879,6.629)rectangle++(0.04,0.04) (1.86,7.156)rectangle++(0.04,0.04) (1.865,7.604)rectangle++(0.04,0.04)
  (1.888,7.985)rectangle++(0.04,0.04) (1.922,8.308)rectangle++(0.04,0.04) (1.963,8.581)rectangle++(0.04,0.04) (2.01,8.812)rectangle++(0.04,0.04) (2.059,9.006)rectangle++(0.04,0.04)
  (0.517,3.064)rectangle++(0.04,0.04) (0.56,4.18)rectangle++(0.04,0.04) (0.641,5.126)rectangle++(0.04,0.04) (0.748,5.928)rectangle++(0.04,0.04) (0.87,6.605)rectangle++(0.04,0.04)
  (-1.202,4.415)rectangle++(0.04,0.04) (-0.848,5.397)rectangle++(0.04,0.04) (-1.254,3.854)rectangle++(0.04,0.04) (-0.914,4.922)rectangle++(0.04,0.04) (-0.582,5.817)rectangle++(0.04,0.04)
  (-0.264,6.564)rectangle++(0.04,0.04) (0.036,7.186)rectangle++(0.04,0.04) (0.316,7.703)rectangle++(0.04,0.04) (0.574,8.131)rectangle++(0.04,0.04) (0.811,8.485)rectangle++(0.04,0.04)
  (1.027,8.776)rectangle++(0.04,0.04) (2.053,3.887)rectangle++(0.04,0.04) (1.898,4.818)rectangle++(0.04,0.04) (-0.639,4.255)rectangle++(0.04,0.04) (-0.375,5.238)rectangle++(0.04,0.04)
  (-0.112,6.064)rectangle++(0.04,0.04) (0.146,6.755)rectangle++(0.04,0.04) (0.392,7.332)rectangle++(0.04,0.04) (-0.987,3.251)rectangle++(0.04,0.04) (-0.711,4.399)rectangle++(0.04,0.04)
  (-0.43,5.364)rectangle++(0.04,0.04) (-0.154,6.173)rectangle++(0.04,0.04) (1.547,3.118)rectangle++(0.04,0.04) (1.437,4.184)rectangle++(0.04,0.04) (1.387,5.095)rectangle++(0.04,0.04)
  (1.38,5.872)rectangle++(0.04,0.04) (1.406,6.532)rectangle++(0.04,0.04) (1.454,7.092)rectangle++(0.04,0.04) (1.517,7.566)rectangle++(0.04,0.04) (1.59,7.967)rectangle++(0.04,0.04)
  (1.668,8.304)rectangle++(0.04,0.04) (1.861,3.942)rectangle++(0.04,0.04) (1.737,4.873)rectangle++(0.04,0.04) (1.003,3.296)rectangle++(0.04,0.04) (0.982,4.358)rectangle++(0.04,0.04)
  (1.007,5.261)rectangle++(0.04,0.04) (1.064,6.028)rectangle++(0.04,0.04) (1.144,6.677)rectangle++(0.04,0.04) (1.237,7.226)rectangle++(0.04,0.04) (1.338,7.689)rectangle++(0.04,0.04)
  (1.443,8.078)rectangle++(0.04,0.04) (1.547,8.405)rectangle++(0.04,0.04) (1.649,8.679)rectangle++(0.04,0.04) (1.747,8.907)rectangle++(0.04,0.04) (1.839,9.097)rectangle++(0.04,0.04)
  (1.925,9.255)rectangle++(0.04,0.04) (2.004,9.386)rectangle++(0.04,0.04) (2.077,9.494)rectangle++(0.04,0.04) (2.143,9.583)rectangle++(0.04,0.04) (2.202,9.656)rectangle++(0.04,0.04)
  (2.256,9.716)rectangle++(0.04,0.04) (-1.614,5.289)rectangle++(0.04,0.04) (-1.163,6.156)rectangle++(0.04,0.04) (-0.744,6.875)rectangle++(0.04,0.04) (-0.36,7.47)rectangle++(0.04,0.04)
  (1.018,2.108)rectangle++(0.04,0.04) (-0.221,3.555)rectangle++(0.04,0.04) (-0.048,4.627)rectangle++(0.04,0.04) (0.142,5.531)rectangle++(0.04,0.04) (0.34,6.292)rectangle++(0.04,0.04)
  (0.538,6.93)rectangle++(0.04,0.04) (0.732,7.466)rectangle++(0.04,0.04) (0.919,7.913)rectangle++(0.04,0.04) (1.095,8.285)rectangle++(0.04,0.04) (1.26,8.595)rectangle++(0.04,0.04)
  (1.413,8.851)rectangle++(0.04,0.04) (2.026,3.99)rectangle++(0.04,0.04) (1.88,4.906)rectangle++(0.04,0.04) (1.792,5.691)rectangle++(0.04,0.04) (1.749,6.362)rectangle++(0.04,0.04)
  (1.739,6.934)rectangle++(0.04,0.04) (1.753,7.421)rectangle++(0.04,0.04) (1.785,7.834)rectangle++(0.04,0.04) (1.828,8.183)rectangle++(0.04,0.04) (1.879,8.479)rectangle++(0.04,0.04)
  (1.934,8.728)rectangle++(0.04,0.04) (1.991,8.938)rectangle++(0.04,0.04) (2.048,9.114)rectangle++(0.04,0.04) (2.103,9.261)rectangle++(0.04,0.04) (2.156,9.384)rectangle++(0.04,0.04)
  (2.205,9.486)rectangle++(0.04,0.04) (2.252,9.571)rectangle++(0.04,0.04) (2.295,9.642)rectangle++(0.04,0.04) (2.334,9.7)rectangle++(0.04,0.04) (2.37,9.748)rectangle++(0.04,0.04)
  (2.402,9.787)rectangle++(0.04,0.04) (2.431,9.819)rectangle++(0.04,0.04) (2.147,4.407)rectangle++(0.04,0.04) (1.999,5.257)rectangle++(0.04,0.04) (1.907,5.984)rectangle++(0.04,0.04)
  (1.858,6.607)rectangle++(0.04,0.04) (1.842,7.137)rectangle++(0.04,0.04) (1.552,3.757)rectangle++(0.04,0.04) (1.467,4.727)rectangle++(0.04,0.04) (1.434,5.556)rectangle++(0.04,0.04)
  (1.439,6.261)rectangle++(0.04,0.04) (0.393,2.752)rectangle++(0.04,0.04) (0.442,3.92)rectangle++(0.04,0.04) (0.531,4.91)rectangle++(0.04,0.04) (0.645,5.749)rectangle++(0.04,0.04)
  (0.776,6.457)rectangle++(0.04,0.04) (0.97,2.226)rectangle++(0.04,0.04) (0.912,3.45)rectangle++(0.04,0.04) (0.911,4.492)rectangle++(0.04,0.04) (0.952,5.378)rectangle++(0.04,0.04)
  (1.022,6.129)rectangle++(0.04,0.04) (1.112,6.765)rectangle++(0.04,0.04) (1.545,3.519)rectangle++(0.04,0.04) (1.452,4.526)rectangle++(0.04,0.04) (1.413,5.385)rectangle++(0.04,0.04)
  (1.414,6.117)rectangle++(0.04,0.04) (1.445,6.739)rectangle++(0.04,0.04) (1.5,3.59)rectangle++(0.04,0.04) (1.416,4.587)rectangle++(0.04,0.04) (1.385,5.439)rectangle++(0.04,0.04)
  (1.393,6.164)rectangle++(0.04,0.04) (1.428,6.78)rectangle++(0.04,0.04) (1.483,7.302)rectangle++(0.04,0.04) (1.55,7.743)rectangle++(0.04,0.04) (1.625,8.116)rectangle++(0.04,0.04)
  (1.704,8.43)rectangle++(0.04,0.04) (0.473,2.995)rectangle++(0.04,0.04) (0.52,4.123)rectangle++(0.04,0.04) (0.604,5.08)rectangle++(0.04,0.04) (0.715,5.89)rectangle++(0.04,0.04)
  (-1.078,4.244)rectangle++(0.04,0.04) (-1.051,3.583)rectangle++(0.04,0.04) (-0.752,4.684)rectangle++(0.04,0.04) (-1.085,3.731)rectangle++(0.04,0.04) (-0.775,4.811)rectangle++(0.04,0.04)
  (-0.469,5.717)rectangle++(0.04,0.04) (-0.172,6.474)rectangle++(0.04,0.04) (0.11,7.106)rectangle++(0.04,0.04) (0.376,7.632)rectangle++(0.04,0.04) (0.623,8.068)rectangle++(0.04,0.04)
  (0.85,8.429)rectangle++(0.04,0.04) (1.057,8.727)rectangle++(0.04,0.04) (-1.619,4.869)rectangle++(0.04,0.04) (-1.184,5.799)rectangle++(0.04,0.04) (-0.776,6.573)rectangle++(0.04,0.04)
  (-1.491,4.1)rectangle++(0.04,0.04) (-1.105,5.141)rectangle++(0.04,0.04) (-0.736,6.01)rectangle++(0.04,0.04) (-0.387,6.734)rectangle++(0.04,0.04) (-0.062,7.336)rectangle++(0.04,0.04)
  (0.238,7.834)rectangle++(0.04,0.04) (0.514,8.246)rectangle++(0.04,0.04) (0.764,8.584)rectangle++(0.04,0.04) (0.991,8.862)rectangle++(0.04,0.04) (1.195,9.09)rectangle++(0.04,0.04)
  (1.377,9.275)rectangle++(0.04,0.04) (1.539,9.424)rectangle++(0.04,0.04) (1.248,2.246)rectangle++(0.04,0.04) (1.148,3.455)rectangle++(0.04,0.04) (-0.148,3.401)rectangle++(0.04,0.04)
  (0.008,4.493)rectangle++(0.04,0.04) (0.184,5.415)rectangle++(0.04,0.04) (0.371,6.192)rectangle++(0.04,0.04) (0.561,6.844)rectangle++(0.04,0.04) (0.748,7.391)rectangle++(0.04,0.04)
  (0.93,7.849)rectangle++(0.04,0.04) (1.102,8.231)rectangle++(0.04,0.04) (1.264,8.548)rectangle++(0.04,0.04) (1.414,8.812)rectangle++(0.04,0.04) (1.552,9.029)rectangle++(0.04,0.04)
  (1.678,9.209)rectangle++(0.04,0.04) (1.793,9.357)rectangle++(0.04,0.04) (1.896,9.478)rectangle++(0.04,0.04) (1.988,9.577)rectangle++(0.04,0.04) (2.071,9.657)rectangle++(0.04,0.04)
  (2.144,9.722)rectangle++(0.04,0.04) (-1.635,5.269)rectangle++(0.04,0.04) (-1.181,6.14)rectangle++(0.04,0.04) (-0.761,6.863)rectangle++(0.04,0.04) (-0.374,7.46)rectangle++(0.04,0.04)
  (-0.022,7.952)rectangle++(0.04,0.04) (0.297,8.356)rectangle++(0.04,0.04) (0.585,8.687)rectangle++(0.04,0.04) (0.842,8.957)rectangle++(0.04,0.04) (1.072,9.176)rectangle++(0.04,0.04)
  (1.276,9.353)rectangle++(0.04,0.04) (1.457,9.495)rectangle++(0.04,0.04) (1.616,9.609)rectangle++(0.04,0.04) (1.755,9.699)rectangle++(0.04,0.04) (1.878,9.77)rectangle++(0.04,0.04)
  (1.985,9.826)rectangle++(0.04,0.04) (2.078,9.869)rectangle++(0.04,0.04) (2.159,9.901)rectangle++(0.04,0.04) (2.229,9.926)rectangle++(0.04,0.04) (2.289,9.944)rectangle++(0.04,0.04)
  (2.342,9.957)rectangle++(0.04,0.04) (2.191,4.414)rectangle++(0.04,0.04) (2.037,5.26)rectangle++(0.04,0.04) (1.94,5.986)rectangle++(0.04,0.04) (-0.879,4.493)rectangle++(0.04,0.04)
  (-0.57,5.45)rectangle++(0.04,0.04) (-0.269,6.252)rectangle++(0.04,0.04) (0.02,6.921)rectangle++(0.04,0.04) (0.291,7.478)rectangle++(0.04,0.04) (0.545,7.941)rectangle++(0.04,0.04)
  (-1.544,4.617)rectangle++(0.04,0.04) (-1.13,5.583)rectangle++(0.04,0.04) (-0.739,6.387)rectangle++(0.04,0.04) (-0.375,7.054)rectangle++(0.04,0.04) (-0.039,7.608)rectangle++(0.04,0.04)
  (0.269,8.064)rectangle++(0.04,0.04) (0.549,8.44)rectangle++(0.04,0.04) (-1.092,3.258)rectangle++(0.04,0.04) (-0.8,4.409)rectangle++(0.04,0.04) (1.186,2.631)rectangle++(0.04,0.04)
  (1.111,3.785)rectangle++(0.04,0.04) (1.093,4.769)rectangle++(0.04,0.04) (1.053,3.13)rectangle++(0.04,0.04) (1.018,4.215)rectangle++(0.04,0.04) (1.032,5.138)rectangle++(0.04,0.04)
  (1.08,5.922)rectangle++(0.04,0.04) (1.153,6.587)rectangle++(0.04,0.04) (-0.151,3.714)rectangle++(0.04,0.04) (0.018,4.759)rectangle++(0.04,0.04) (0.204,5.641)rectangle++(0.04,0.04)
  (0.397,6.383)rectangle++(0.04,0.04) (0.59,7.006)rectangle++(0.04,0.04) (0.78,7.527)rectangle++(0.04,0.04) (0.962,7.963)rectangle++(0.04,0.04) (1.134,8.326)rectangle++(0.04,0.04)
  (1.929,3.825)rectangle++(0.04,0.04) (1.791,4.77)rectangle++(0.04,0.04) (1.711,5.579)rectangle++(0.04,0.04) (1.675,6.27)rectangle++(0.04,0.04) (1.672,6.859)rectangle++(0.04,0.04)
  (1.694,7.359)rectangle++(0.04,0.04) (1.732,7.784)rectangle++(0.04,0.04) (1.781,8.143)rectangle++(0.04,0.04) (1.808,3.937)rectangle++(0.04,0.04) (1.692,4.87)rectangle++(0.04,0.04)
  (1.631,5.668)rectangle++(0.04,0.04) (1.611,6.349)rectangle++(0.04,0.04) (1.621,6.928)rectangle++(0.04,0.04) (1.653,7.42)rectangle++(0.04,0.04) (1.699,7.837)rectangle++(0.04,0.04)
  (1.756,8.19)rectangle++(0.04,0.04) (1.818,8.488)rectangle++(0.04,0.04) (1.882,8.738)rectangle++(0.04,0.04) (1.947,8.948)rectangle++(0.04,0.04) (2.011,9.124)rectangle++(0.04,0.04)
  (2.072,9.271)rectangle++(0.04,0.04) (2.13,9.394)rectangle++(0.04,0.04) (2.184,9.496)rectangle++(0.04,0.04) (2.234,9.58)rectangle++(0.04,0.04) (2.28,9.65)rectangle++(0.04,0.04)
  (2.322,9.708)rectangle++(0.04,0.04) (2.36,9.755)rectangle++(0.04,0.04) (2.394,9.793)rectangle++(0.04,0.04) (2.424,9.825)rectangle++(0.04,0.04) (2.451,9.85)rectangle++(0.04,0.04)
  (2.475,9.871)rectangle++(0.04,0.04) (2.154,4.427)rectangle++(0.04,0.04) (2.006,5.273)rectangle++(0.04,0.04) (1.914,5.998)rectangle++(0.04,0.04) (1.865,6.618)rectangle++(0.04,0.04)
  (1.847,7.147)rectangle++(0.04,0.04) (1.854,7.597)rectangle++(0.04,0.04) (1.878,7.98)rectangle++(0.04,0.04) (1.913,8.304)rectangle++(0.04,0.04) (1.956,8.578)rectangle++(0.04,0.04)
  (2.003,8.809)rectangle++(0.04,0.04) (2.053,9.004)rectangle++(0.04,0.04) (2.103,9.167)rectangle++(0.04,0.04) (2.152,9.304)rectangle++(0.04,0.04) (2.199,9.419)rectangle++(0.04,0.04)
  (2.244,9.514)rectangle++(0.04,0.04) (2.286,9.594)rectangle++(0.04,0.04) (2.324,9.659)rectangle++(0.04,0.04) (0.042,2.896)rectangle++(0.04,0.04) (0.149,4.056)rectangle++(0.04,0.04)
  (0.287,5.038)rectangle++(0.04,0.04) (0.443,5.867)rectangle++(0.04,0.04) (0.609,6.566)rectangle++(0.04,0.04) (1.555,3.367)rectangle++(0.04,0.04) (1.454,4.396)rectangle++(0.04,0.04)
  (1.409,5.274)rectangle++(0.04,0.04) (1.141,3.299)rectangle++(0.04,0.04) (1.1,4.355)rectangle++(0.04,0.04) (1.107,5.254)rectangle++(0.04,0.04) (1.149,6.018)rectangle++(0.04,0.04)
  (1.215,6.665)rectangle++(0.04,0.04) (1.297,7.213)rectangle++(0.04,0.04) (1.389,7.675)rectangle++(0.04,0.04) (1.485,8.065)rectangle++(0.04,0.04) (1.583,8.392)rectangle++(0.04,0.04)
  (1.679,8.666)rectangle++(0.04,0.04) (1.772,8.895)rectangle++(0.04,0.04) (-1.527,5.036)rectangle++(0.04,0.04) (-1.098,5.938)rectangle++(0.04,0.04) (-0.698,6.687)rectangle++(0.04,0.04)
  (-0.328,7.308)rectangle++(0.04,0.04) (0.011,7.821)rectangle++(0.04,0.04) (0.32,8.244)rectangle++(0.04,0.04) (0.6,8.591)rectangle++(0.04,0.04) (0.851,8.874)rectangle++(0.04,0.04)
  (1.076,9.105)rectangle++(0.04,0.04) (1.277,9.293)rectangle++(0.04,0.04) (1.455,9.444)rectangle++(0.04,0.04) (1.612,9.565)rectangle++(0.04,0.04) (1.751,9.662)rectangle++(0.04,0.04)
  (1.872,9.739)rectangle++(0.04,0.04) (1.979,9.8)rectangle++(0.04,0.04) (2.072,9.847)rectangle++(0.04,0.04) (-1.675,5.28)rectangle++(0.04,0.04) (-1.215,6.151)rectangle++(0.04,0.04)
  (-0.789,6.873)rectangle++(0.04,0.04) (-0.398,7.47)rectangle++(0.04,0.04) (-0.041,7.962)rectangle++(0.04,0.04) (0.281,8.365)rectangle++(0.04,0.04) (-1.681,4.652)rectangle++(0.04,0.04)
  (-1.245,5.617)rectangle++(0.04,0.04) (1.537,2.763)rectangle++(0.04,0.04) (1.415,3.883)rectangle++(0.04,0.04) (1.356,4.84)rectangle++(0.04,0.04) (1.04,3.203)rectangle++(0.04,0.04)
  (1.01,4.277)rectangle++(0.04,0.04) (1.027,5.192)rectangle++(0.04,0.04) (1.079,5.968)rectangle++(0.04,0.04) (1.153,6.626)rectangle++(0.04,0.04) (1.506,3.502)rectangle++(0.04,0.04)
  (1.418,4.512)rectangle++(0.04,0.04) (1.383,5.375)rectangle++(0.04,0.04) (1.389,6.11)rectangle++(0.04,0.04) (1.423,6.734)rectangle++(0.04,0.04) (1.476,7.263)rectangle++(0.04,0.04)
  (1.543,7.711)rectangle++(0.04,0.04) (1.618,8.089)rectangle++(0.04,0.04) (1.697,8.407)rectangle++(0.04,0.04) (1.776,8.674)rectangle++(0.04,0.04) (-1.479,4.993)rectangle++(0.04,0.04)
  (-1.06,5.899)rectangle++(0.04,0.04) (-0.667,6.653)rectangle++(0.04,0.04) (-0.303,7.278)rectangle++(0.04,0.04) (0.031,7.795)rectangle++(0.04,0.04) (0.336,8.22)rectangle++(0.04,0.04)
  (0.355,2.95)rectangle++(0.04,0.04) (0.418,4.089)rectangle++(0.04,0.04) (0.516,5.055)rectangle++(0.04,0.04) (0.639,5.873)rectangle++(0.04,0.04) (-1.088,4.227)rectangle++(0.04,0.04)
  (-0.758,5.232)rectangle++(0.04,0.04) (-0.437,6.074)rectangle++(0.04,0.04) (-0.131,6.777)rectangle++(0.04,0.04) (0.158,7.362)rectangle++(0.04,0.04) (0.426,7.847)rectangle++(0.04,0.04)
  (0.674,8.249)rectangle++(0.04,0.04) (1.964,3.706)rectangle++(0.04,0.04) (-0.391,4.047)rectangle++(0.04,0.04) (-0.173,5.052)rectangle++(0.04,0.04) (0.053,5.897)rectangle++(0.04,0.04)
  (0.279,6.607)rectangle++(0.04,0.04) (0.499,7.201)rectangle++(0.04,0.04) (0.71,7.697)rectangle++(0.04,0.04) (0.909,8.11)rectangle++(0.04,0.04) (1.095,8.454)rectangle++(0.04,0.04)
  (1.267,8.738)rectangle++(0.04,0.04) (1.424,8.973)rectangle++(0.04,0.04) (1.567,9.166)rectangle++(0.04,0.04) (1.697,9.325)rectangle++(0.04,0.04) (2.11,4.145)rectangle++(0.04,0.04)
  (1.957,5.035)rectangle++(0.04,0.04) (1.863,5.798)rectangle++(0.04,0.04) (1.813,6.45)rectangle++(0.04,0.04) (1.797,7.006)rectangle++(0.04,0.04) (-0.204,3.771)rectangle++(0.04,0.04)
  (0.958,2.643)rectangle++(0.04,0.04) (0.918,3.805)rectangle++(0.04,0.04) (0.93,4.793)rectangle++(0.04,0.04) (0.98,5.633)rectangle++(0.04,0.04) (1.056,6.345)rectangle++(0.04,0.04)
  (0.364,2.755)rectangle++(0.04,0.04) (0.418,3.924)rectangle++(0.04,0.04) (0.51,4.915)rectangle++(0.04,0.04) (0.628,5.753)rectangle++(0.04,0.04) (0.761,6.461)rectangle++(0.04,0.04)
  (-1.191,4.366)rectangle++(0.04,0.04) (-0.84,5.355)rectangle++(0.04,0.04) (-1.244,3.847)rectangle++(0.04,0.04) (-0.905,4.916)rectangle++(0.04,0.04) (-0.575,5.811)rectangle++(0.04,0.04)
  (-0.259,6.559)rectangle++(0.04,0.04) (0.04,7.181)rectangle++(0.04,0.04) (0.319,7.699)rectangle++(0.04,0.04) (1.869,3.519)rectangle++(0.04,0.04) (1.727,4.513)rectangle++(0.04,0.04)
  (1.646,5.363)rectangle++(0.04,0.04) (1.612,6.089)rectangle++(0.04,0.04) (1.611,6.707)rectangle++(0.04,0.04) (1.636,7.233)rectangle++(0.04,0.04) (1.678,7.679)rectangle++(0.04,0.04)
  (1.731,8.056)rectangle++(0.04,0.04) (1.791,8.375)rectangle++(0.04,0.04) (1.864,3.984)rectangle++(0.04,0.04) (1.742,4.908)rectangle++(0.04,0.04) (1.011,3.304)rectangle++(0.04,0.04)
  (0.99,4.364)rectangle++(0.04,0.04) (1.014,5.266)rectangle++(0.04,0.04) (1.07,6.032)rectangle++(0.04,0.04) (1.149,6.68)rectangle++(0.04,0.04) (-1.168,4.482)rectangle++(0.04,0.04)
  (-0.816,5.452)rectangle++(0.04,0.04) (-0.478,6.263)rectangle++(0.04,0.04) (-0.158,6.939)rectangle++(0.04,0.04) (0.141,7.501)rectangle++(0.04,0.04) (0.418,7.966)rectangle++(0.04,0.04)
  (0.672,8.351)rectangle++(0.04,0.04) (0.903,8.668)rectangle++(0.04,0.04) (1.112,8.927)rectangle++(0.04,0.04) (1.3,9.14)rectangle++(0.04,0.04) (1.468,9.313)rectangle++(0.04,0.04)
  (1.618,9.454)rectangle++(0.04,0.04) (1.752,9.567)rectangle++(0.04,0.04) (1.869,9.658)rectangle++(0.04,0.04) (1.973,9.731)rectangle++(0.04,0.04) (2.064,9.789)rectangle++(0.04,0.04)
  (2.144,9.834)rectangle++(0.04,0.04) (2.213,9.869)rectangle++(0.04,0.04) (2.274,9.897)rectangle++(0.04,0.04) (2.327,9.917)rectangle++(0.04,0.04) (2.372,9.933)rectangle++(0.04,0.04)
  (2.411,9.944)rectangle++(0.04,0.04) (2.445,9.952)rectangle++(0.04,0.04) (2.474,9.958)rectangle++(0.04,0.04) (2.499,9.961)rectangle++(0.04,0.04) (2.521,9.963)rectangle++(0.04,0.04)
  (2.539,9.964)rectangle++(0.04,0.04) (2.554,9.964)rectangle++(0.04,0.04) (2.568,9.964)rectangle++(0.04,0.04) (2.579,9.963)rectangle++(0.04,0.04) (2.588,9.961)rectangle++(0.04,0.04)
  (2.596,9.96)rectangle++(0.04,0.04) (2.603,9.958)rectangle++(0.04,0.04) (2.609,9.956)rectangle++(0.04,0.04) (2.613,9.955)rectangle++(0.04,0.04) (2.617,9.953)rectangle++(0.04,0.04)
  (2.621,9.952)rectangle++(0.04,0.04) (-1.599,5.404)rectangle++(0.04,0.04) (-1.145,6.253)rectangle++(0.04,0.04) (-0.725,6.957)rectangle++(0.04,0.04) (-0.34,7.539)rectangle++(0.04,0.04)
  (0.01,8.018)rectangle++(0.04,0.04) (-1.656,4.532)rectangle++(0.04,0.04) (-1.229,5.515)rectangle++(0.04,0.04) (1.51,2.747)rectangle++(0.04,0.04) (-0.269,3.773)rectangle++(0.04,0.04)
  (-0.08,4.814)rectangle++(0.04,0.04) (0.122,5.691)rectangle++(0.04,0.04) (0.33,6.429)rectangle++(0.04,0.04) (0.535,7.048)rectangle++(0.04,0.04) (0.734,7.565)rectangle++(0.04,0.04)
  (0.925,7.997)rectangle++(0.04,0.04) (1.104,8.357)rectangle++(0.04,0.04) (1.27,8.655)rectangle++(0.04,0.04) (0.036,2.76)rectangle++(0.04,0.04) (0.138,3.941)rectangle++(0.04,0.04)
  (0.273,4.94)rectangle++(0.04,0.04) (1.191,2.977)rectangle++(0.04,0.04) (1.13,4.079)rectangle++(0.04,0.04) (1.121,5.018)rectangle++(0.04,0.04) (1.151,5.817)rectangle++(0.04,0.04)
  (1.209,6.494)rectangle++(0.04,0.04) (1.285,7.068)rectangle++(0.04,0.04) (1.373,7.553)rectangle++(0.04,0.04) (1.467,7.961)rectangle++(0.04,0.04) (1.563,8.304)rectangle++(0.04,0.04)
  (1.659,8.592)rectangle++(0.04,0.04) (1.751,8.833)rectangle++(0.04,0.04) (1.84,9.035)rectangle++(0.04,0.04) (-1.164,3.362)rectangle++(0.04,0.04) (-0.857,4.5)rectangle++(0.04,0.04)
  (-0.551,5.456)rectangle++(0.04,0.04) (-0.252,6.256)rectangle++(0.04,0.04) (0.034,6.924)rectangle++(0.04,0.04) (0.303,7.48)rectangle++(0.04,0.04) (0.555,7.942)rectangle++(0.04,0.04)
  (-1.543,4.619)rectangle++(0.04,0.04) (-1.129,5.584)rectangle++(0.04,0.04) (-0.738,6.388)rectangle++(0.04,0.04) (-0.374,7.056)rectangle++(0.04,0.04) (-1.521,4.27)rectangle++(0.04,0.04)
  (-1.124,5.287)rectangle++(0.04,0.04) (-0.746,6.135)rectangle++(0.04,0.04) (-0.391,6.841)rectangle++(0.04,0.04) (-0.061,7.427)rectangle++(0.04,0.04) (0.243,7.911)rectangle++(0.04,0.04)
  (0.521,8.311)rectangle++(0.04,0.04) (1.997,3.683)rectangle++(0.04,0.04) (1.843,4.647)rectangle++(0.04,0.04) (1.75,5.472)rectangle++(0.04,0.04) (1.704,6.178)rectangle++(0.04,0.04)
  (0.637,2.093)rectangle++(0.04,0.04) (0.623,3.35)rectangle++(0.04,0.04) (-0.556,3.726)rectangle++(0.04,0.04) (-0.326,4.785)rectangle++(0.04,0.04) (-0.088,5.677)rectangle++(0.04,0.04)
  (0.15,6.425)rectangle++(0.04,0.04) (1.573,3.235)rectangle++(0.04,0.04) (1.465,4.283)rectangle++(0.04,0.04) (1.414,5.178)rectangle++(0.04,0.04) (1.407,5.941)rectangle++(0.04,0.04)
  (1.431,6.59)rectangle++(0.04,0.04) (1.478,7.14)rectangle++(0.04,0.04) (1.54,7.606)rectangle++(0.04,0.04) (1.611,8)rectangle++(0.04,0.04) (1.687,8.332)rectangle++(0.04,0.04);
\fill[piece2]
  (1.765,8.611)rectangle++(0.04,0.04) (1.926,4.023)rectangle++(0.04,0.04) (1.795,4.939)rectangle++(0.04,0.04) (1.721,5.722)rectangle++(0.04,0.04) (1.69,6.391)rectangle++(0.04,0.04)
  (1.69,6.961)rectangle++(0.04,0.04) (1.713,7.446)rectangle++(0.04,0.04) (1.643,3.787)rectangle++(0.04,0.04) (1.546,4.749)rectangle++(0.04,0.04) (1.502,5.571)rectangle++(0.04,0.04)
  (1.497,6.272)rectangle++(0.04,0.04) (1.521,6.867)rectangle++(0.04,0.04) (1.566,7.372)rectangle++(0.04,0.04) (-1.241,4.697)rectangle++(0.04,0.04) (-0.869,5.638)rectangle++(0.04,0.04)
  (-0.515,6.423)rectangle++(0.04,0.04) (-0.183,7.077)rectangle++(0.04,0.04) (0.125,7.619)rectangle++(0.04,0.04) (0.409,8.067)rectangle++(0.04,0.04) (0.668,8.437)rectangle++(0.04,0.04)
  (0.903,8.741)rectangle++(0.04,0.04) (0.37,2.382)rectangle++(0.04,0.04) (0.408,3.606)rectangle++(0.04,0.04) (-0.649,3.736)rectangle++(0.04,0.04) (-0.404,4.798)rectangle++(0.04,0.04)
  (-1.047,3.819)rectangle++(0.04,0.04) (0.685,2.107)rectangle++(0.04,0.04) (0.665,3.359)rectangle++(0.04,0.04) (0.697,4.425)rectangle++(0.04,0.04) (0.767,5.33)rectangle++(0.04,0.04)
  (0.604,2.103)rectangle++(0.04,0.04) (0.595,3.359)rectangle++(0.04,0.04) (0.638,4.428)rectangle++(0.04,0.04) (0.717,5.334)rectangle++(0.04,0.04) (0.821,6.102)rectangle++(0.04,0.04)
  (0.939,6.75)rectangle++(0.04,0.04) (1.066,7.296)rectangle++(0.04,0.04) (1.196,7.755)rectangle++(0.04,0.04) (1.325,8.14)rectangle++(0.04,0.04) (1.449,8.462)rectangle++(0.04,0.04)
  (1.568,8.731)rectangle++(0.04,0.04) (1.68,8.955)rectangle++(0.04,0.04) (1.784,9.141)rectangle++(0.04,0.04) (2.055,4.13)rectangle++(0.04,0.04) (1.909,5.024)rectangle++(0.04,0.04)
  (1.822,5.79)rectangle++(0.04,0.04) (1.778,6.445)rectangle++(0.04,0.04) (1.767,7.004)rectangle++(0.04,0.04) (1.78,7.479)rectangle++(0.04,0.04) (1.81,7.882)rectangle++(0.04,0.04)
  (1.851,8.223)rectangle++(0.04,0.04) (1.9,8.512)rectangle++(0.04,0.04) (1.954,8.755)rectangle++(0.04,0.04) (2.009,8.96)rectangle++(0.04,0.04) (2.063,9.132)rectangle++(0.04,0.04)
  (2.117,9.276)rectangle++(0.04,0.04) (2.168,9.396)rectangle++(0.04,0.04) (2.217,9.496)rectangle++(0.04,0.04) (2.262,9.579)rectangle++(0.04,0.04) (2.304,9.648)rectangle++(0.04,0.04)
  (2.342,9.705)rectangle++(0.04,0.04) (2.376,9.752)rectangle++(0.04,0.04) (2.408,9.79)rectangle++(0.04,0.04) (2.436,9.821)rectangle++(0.04,0.04) (2.461,9.847)rectangle++(0.04,0.04)
  (2.484,9.868)rectangle++(0.04,0.04) (2.504,9.884)rectangle++(0.04,0.04) (2.521,9.898)rectangle++(0.04,0.04) (2.537,9.908)rectangle++(0.04,0.04) (2.55,9.917)rectangle++(0.04,0.04)
  (2.562,9.924)rectangle++(0.04,0.04) (2.573,9.929)rectangle++(0.04,0.04) (2.582,9.933)rectangle++(0.04,0.04) (2.59,9.936)rectangle++(0.04,0.04) (2.596,9.938)rectangle++(0.04,0.04)
  (2.602,9.94)rectangle++(0.04,0.04) (2.607,9.941)rectangle++(0.04,0.04) (2.612,9.942)rectangle++(0.04,0.04) (2.615,9.942)rectangle++(0.04,0.04) (2.619,9.942)rectangle++(0.04,0.04)
  (2.621,9.942)rectangle++(0.04,0.04) (2.624,9.942)rectangle++(0.04,0.04) (2.626,9.942)rectangle++(0.04,0.04) (2.627,9.942)rectangle++(0.04,0.04) (2.629,9.942)rectangle++(0.04,0.04)
  (-1.595,5.403)rectangle++(0.04,0.04) (1.525,2.642)rectangle++(0.04,0.04) (1.4,3.781)rectangle++(0.04,0.04) (0.772,3.009)rectangle++(0.04,0.04) (0.775,4.123)rectangle++(0.04,0.04)
  (0.821,5.07)rectangle++(0.04,0.04) (0.899,5.873)rectangle++(0.04,0.04) (0.997,6.552)rectangle++(0.04,0.04) (1.107,7.126)rectangle++(0.04,0.04) (1.635,3.587)rectangle++(0.04,0.04)
  (1.531,4.58)rectangle++(0.04,0.04) (1.482,5.428)rectangle++(0.04,0.04) (1.475,6.151)rectangle++(0.04,0.04) (1.35,3.483)rectangle++(0.04,0.04) (1.285,4.503)rectangle++(0.04,0.04)
  (1.27,5.372)rectangle++(0.04,0.04) (1.182,3.286)rectangle++(0.04,0.04) (1.134,4.342)rectangle++(0.04,0.04) (0.943,3.057)rectangle++(0.04,0.04) (0.921,4.157)rectangle++(0.04,0.04)
  (0.947,5.093)rectangle++(0.04,0.04) (1.007,5.887)rectangle++(0.04,0.04) (1.089,6.56)rectangle++(0.04,0.04) (1.186,7.129)rectangle++(0.04,0.04) (1.291,7.608)rectangle++(0.04,0.04)
  (1.399,8.011)rectangle++(0.04,0.04) (1.508,8.35)rectangle++(0.04,0.04) (-1.096,3.338)rectangle++(0.04,0.04) (-0.8,4.477)rectangle++(0.04,0.04) (-0.503,5.434)rectangle++(0.04,0.04)
  (-0.213,6.235)rectangle++(0.04,0.04) (0.066,6.905)rectangle++(0.04,0.04) (0.33,7.462)rectangle++(0.04,0.04) (0.577,7.926)rectangle++(0.04,0.04) (0.805,8.31)rectangle++(0.04,0.04)
  (1.015,8.628)rectangle++(0.04,0.04) (1.206,8.889)rectangle++(0.04,0.04) (1.378,9.104)rectangle++(0.04,0.04) (-1.642,5.003)rectangle++(0.04,0.04) (-1.197,5.915)rectangle++(0.04,0.04)
  (-0.783,6.671)rectangle++(0.04,0.04) (-0.401,7.298)rectangle++(0.04,0.04) (-0.051,7.816)rectangle++(0.04,0.04) (0.267,8.242)rectangle++(0.04,0.04) (0.554,8.591)rectangle++(0.04,0.04)
  (0.813,8.876)rectangle++(0.04,0.04) (1.044,9.109)rectangle++(0.04,0.04) (1.249,9.297)rectangle++(0.04,0.04) (1.431,9.448)rectangle++(0.04,0.04) (1.592,9.57)rectangle++(0.04,0.04)
  (1.734,9.667)rectangle++(0.04,0.04) (1.859,9.744)rectangle++(0.04,0.04) (2.195,4.263)rectangle++(0.04,0.04) (2.034,5.132)rectangle++(0.04,0.04) (1.932,5.877)rectangle++(0.04,0.04)
  (1.875,6.514)rectangle++(0.04,0.04) (1.852,7.058)rectangle++(0.04,0.04) (0.025,2.543)rectangle++(0.04,0.04) (0.121,3.756)rectangle++(0.04,0.04) (0.251,4.784)rectangle++(0.04,0.04)
  (0.402,5.653)rectangle++(0.04,0.04) (0.566,6.385)rectangle++(0.04,0.04) (0.734,7.001)rectangle++(0.04,0.04) (0.902,7.518)rectangle++(0.04,0.04) (1.065,7.95)rectangle++(0.04,0.04)
  (1.221,8.311)rectangle++(0.04,0.04) (1.915,3.842)rectangle++(0.04,0.04) (1.779,4.785)rectangle++(0.04,0.04) (-0.653,4.226)rectangle++(0.04,0.04) (-0.388,5.214)rectangle++(0.04,0.04)
  (-0.124,6.044)rectangle++(0.04,0.04) (1.511,3.1)rectangle++(0.04,0.04) (1.406,4.17)rectangle++(0.04,0.04) (1.36,5.085)rectangle++(0.04,0.04) (1.357,5.864)rectangle++(0.04,0.04)
  (1.386,6.526)rectangle++(0.04,0.04) (1.437,7.088)rectangle++(0.04,0.04) (1.503,7.564)rectangle++(0.04,0.04) (1.578,7.965)rectangle++(0.04,0.04) (1.657,8.303)rectangle++(0.04,0.04)
  (-1.422,4.898)rectangle++(0.04,0.04) (-1.015,5.816)rectangle++(0.04,0.04) (-0.632,6.58)rectangle++(0.04,0.04) (1.704,3.088)rectangle++(0.04,0.04) (1.569,4.153)rectangle++(0.04,0.04)
  (1.498,5.063)rectangle++(0.04,0.04) (1.473,5.84)rectangle++(0.04,0.04) (1.484,6.501)rectangle++(0.04,0.04) (1.519,7.063)rectangle++(0.04,0.04) (1.572,7.539)rectangle++(0.04,0.04)
  (1.635,7.941)rectangle++(0.04,0.04) (-1.349,4.822)rectangle++(0.04,0.04) (-0.956,5.749)rectangle++(0.04,0.04) (-0.585,6.521)rectangle++(0.04,0.04) (0.93,2.103)rectangle++(0.04,0.04)
  (0.873,3.347)rectangle++(0.04,0.04) (0.873,4.406)rectangle++(0.04,0.04) (0.916,5.306)rectangle++(0.04,0.04) (0.989,6.07)rectangle++(0.04,0.04) (1.081,6.716)rectangle++(0.04,0.04)
  (1.186,7.262)rectangle++(0.04,0.04) (1.296,7.721)rectangle++(0.04,0.04) (0.421,2.91)rectangle++(0.04,0.04) (-0.499,3.601)rectangle++(0.04,0.04) (-0.282,4.677)rectangle++(0.04,0.04)
  (-0.055,5.583)rectangle++(0.04,0.04) (0.174,6.344)rectangle++(0.04,0.04) (0.4,6.982)rectangle++(0.04,0.04) (0.617,7.515)rectangle++(0.04,0.04) (-1.441,4.547)rectangle++(0.04,0.04)
  (-1.045,5.519)rectangle++(0.04,0.04) (-0.67,6.329)rectangle++(0.04,0.04) (-1.42,4.071)rectangle++(0.04,0.04) (-1.046,5.114)rectangle++(0.04,0.04) (-0.687,5.985)rectangle++(0.04,0.04)
  (-0.347,6.711)rectangle++(0.04,0.04) (-0.028,7.314)rectangle++(0.04,0.04) (1.016,2.14)rectangle++(0.04,0.04) (0.947,3.375)rectangle++(0.04,0.04) (0.725,2.828)rectangle++(0.04,0.04)
  (-0.427,3.642)rectangle++(0.04,0.04) (-0.219,4.709)rectangle++(0.04,0.04) (1.192,2.82)rectangle++(0.04,0.04) (1.124,3.946)rectangle++(0.04,0.04) (-0.593,3.937)rectangle++(0.04,0.04)
  (-0.348,4.967)rectangle++(0.04,0.04) (1.271,2.841)rectangle++(0.04,0.04) (1.192,3.96)rectangle++(0.04,0.04) (0.005,2.121)rectangle++(0.04,0.04) (0.087,3.399)rectangle++(0.04,0.04)
  (0.208,4.482)rectangle++(0.04,0.04) (0.354,5.398)rectangle++(0.04,0.04) (0.514,6.17)rectangle++(0.04,0.04) (0.682,6.82)rectangle++(0.04,0.04) (0.85,7.366)rectangle++(0.04,0.04)
  (1.015,7.823)rectangle++(0.04,0.04) (1.174,8.205)rectangle++(0.04,0.04) (1.324,8.524)rectangle++(0.04,0.04) (1.464,8.789)rectangle++(0.04,0.04) (1.594,9.008)rectangle++(0.04,0.04)
  (1.713,9.189)rectangle++(0.04,0.04) (1.821,9.338)rectangle++(0.04,0.04) (1.919,9.461)rectangle++(0.04,0.04) (2.008,9.561)rectangle++(0.04,0.04) (2.133,4.262)rectangle++(0.04,0.04)
  (1.981,5.133)rectangle++(0.04,0.04) (1.039,3.402)rectangle++(0.04,0.04) (-0.492,3.814)rectangle++(0.04,0.04) (-0.268,4.858)rectangle++(0.04,0.04) (-0.036,5.736)rectangle++(0.04,0.04)
  (-1.18,4.073)rectangle++(0.04,0.04) (-0.842,5.106)rectangle++(0.04,0.04) (-0.514,5.97)rectangle++(0.04,0.04) (-0.2,6.691)rectangle++(0.04,0.04) (-1.412,4.231)rectangle++(0.04,0.04)
  (-1.034,5.249)rectangle++(0.04,0.04) (-0.671,6.099)rectangle++(0.04,0.04) (-1.371,4.026)rectangle++(0.04,0.04) (-1.007,5.073)rectangle++(0.04,0.04) (-1.214,3.76)rectangle++(0.04,0.04)
  (-0.883,4.841)rectangle++(0.04,0.04) (-0.559,5.746)rectangle++(0.04,0.04) (-0.248,6.503)rectangle++(0.04,0.04) (0.047,7.134)rectangle++(0.04,0.04) (0.323,7.658)rectangle++(0.04,0.04)
  (0.579,8.093)rectangle++(0.04,0.04) (0.814,8.452)rectangle++(0.04,0.04) (1.027,8.748)rectangle++(0.04,0.04) (1.221,8.991)rectangle++(0.04,0.04) (-1.646,4.951)rectangle++(0.04,0.04)
  (-1.203,5.871)rectangle++(0.04,0.04) (-0.79,6.634)rectangle++(0.04,0.04) (-0.408,7.267)rectangle++(0.04,0.04) (-0.059,7.79)rectangle++(0.04,0.04) (0.26,8.22)rectangle++(0.04,0.04)
  (0.547,8.572)rectangle++(0.04,0.04) (0.806,8.861)rectangle++(0.04,0.04) (1.037,9.096)rectangle++(0.04,0.04) (1.243,9.286)rectangle++(0.04,0.04) (1.426,9.44)rectangle++(0.04,0.04)
  (-1.704,5.078)rectangle++(0.04,0.04) (-1.248,5.981)rectangle++(0.04,0.04) (-0.823,6.73)rectangle++(0.04,0.04) (-0.433,7.35)rectangle++(0.04,0.04) (-0.076,7.861)rectangle++(0.04,0.04)
  (0.247,8.281)rectangle++(0.04,0.04) (0.539,8.625)rectangle++(0.04,0.04) (0.801,8.906)rectangle++(0.04,0.04) (1.035,9.134)rectangle++(0.04,0.04) (1.243,9.319)rectangle++(0.04,0.04)
  (1.427,9.467)rectangle++(0.04,0.04) (1.59,9.586)rectangle++(0.04,0.04) (1.732,9.681)rectangle++(0.04,0.04) (1.858,9.756)rectangle++(0.04,0.04) (1.967,9.814)rectangle++(0.04,0.04)
  (2.062,9.86)rectangle++(0.04,0.04) (2.145,9.894)rectangle++(0.04,0.04) (2.217,9.921)rectangle++(0.04,0.04) (2.279,9.94)rectangle++(0.04,0.04) (2.333,9.954)rectangle++(0.04,0.04)
  (2.379,9.964)rectangle++(0.04,0.04) (2.418,9.97)rectangle++(0.04,0.04) (2.452,9.974)rectangle++(0.04,0.04) (2.481,9.976)rectangle++(0.04,0.04) (2.506,9.977)rectangle++(0.04,0.04)
  (2.527,9.976)rectangle++(0.04,0.04) (2.545,9.975)rectangle++(0.04,0.04) (2.56,9.973)rectangle++(0.04,0.04) (2.572,9.971)rectangle++(0.04,0.04) (2.583,9.969)rectangle++(0.04,0.04)
  (2.592,9.966)rectangle++(0.04,0.04) (2.6,9.964)rectangle++(0.04,0.04) (2.606,9.961)rectangle++(0.04,0.04) (2.612,9.959)rectangle++(0.04,0.04) (2.616,9.957)rectangle++(0.04,0.04)
  (2.62,9.955)rectangle++(0.04,0.04) (2.623,9.953)rectangle++(0.04,0.04) (2.625,9.951)rectangle++(0.04,0.04) (2.627,9.95)rectangle++(0.04,0.04) (2.629,9.949)rectangle++(0.04,0.04)
  (2.63,9.947)rectangle++(0.04,0.04) (2.632,9.946)rectangle++(0.04,0.04) (2.632,9.945)rectangle++(0.04,0.04) (2.633,9.944)rectangle++(0.04,0.04) (2.634,9.944)rectangle++(0.04,0.04)
  (2.634,9.943)rectangle++(0.04,0.04) (-0.275,4.083)rectangle++(0.04,0.04) (-0.072,5.078)rectangle++(0.04,0.04) (-1.047,3.937)rectangle++(0.04,0.04) (-0.735,4.984)rectangle++(0.04,0.04)
  (-0.427,5.862)rectangle++(0.04,0.04) (-1.277,4.025)rectangle++(0.04,0.04) (-0.927,5.068)rectangle++(0.04,0.04) (-0.587,5.941)rectangle++(0.04,0.04) (-0.264,6.67)rectangle++(0.04,0.04)
  (0.04,7.276)rectangle++(0.04,0.04) (0.323,7.779)rectangle++(0.04,0.04) (0.584,8.196)rectangle++(0.04,0.04) (0.822,8.539)rectangle++(0.04,0.04) (1.038,8.822)rectangle++(0.04,0.04)
  (1.233,9.053)rectangle++(0.04,0.04) (1.408,9.242)rectangle++(0.04,0.04) (1.564,9.396)rectangle++(0.04,0.04) (1.703,9.52)rectangle++(0.04,0.04) (1.826,9.62)rectangle++(0.04,0.04)
  (1.935,9.7)rectangle++(0.04,0.04) (2.03,9.764)rectangle++(0.04,0.04) (2.114,9.814)rectangle++(0.04,0.04) (2.188,9.854)rectangle++(0.04,0.04) (2.251,9.884)rectangle++(0.04,0.04)
  (2.307,9.908)rectangle++(0.04,0.04) (2.355,9.926)rectangle++(0.04,0.04) (2.396,9.939)rectangle++(0.04,0.04) (2.432,9.948)rectangle++(0.04,0.04) (2.463,9.955)rectangle++(0.04,0.04)
  (2.49,9.959)rectangle++(0.04,0.04) (-1.624,5.381)rectangle++(0.04,0.04) (-1.167,6.235)rectangle++(0.04,0.04) (1.681,2.899)rectangle++(0.04,0.04) (1.542,3.994)rectangle++(0.04,0.04)
  (1.468,4.929)rectangle++(0.04,0.04) (0.36,2.606)rectangle++(0.04,0.04) (0.408,3.797)rectangle++(0.04,0.04) (0.497,4.807)rectangle++(0.04,0.04) (0.612,5.663)rectangle++(0.04,0.04)
  (0.745,6.385)rectangle++(0.04,0.04) (0.886,6.994)rectangle++(0.04,0.04) (1.031,7.505)rectangle++(0.04,0.04) (0.028,2.603)rectangle++(0.04,0.04) (0.126,3.808)rectangle++(0.04,0.04)
  (0.257,4.828)rectangle++(0.04,0.04) (0.409,5.69)rectangle++(0.04,0.04) (0.573,6.416)rectangle++(0.04,0.04) (0.742,7.027)rectangle++(0.04,0.04) (0.909,7.539)rectangle++(0.04,0.04)
  (1.072,7.968)rectangle++(0.04,0.04) (1.848,3.742)rectangle++(0.04,0.04) (1.718,4.703)rectangle++(0.04,0.04) (1.646,5.525)rectangle++(0.04,0.04) (1.618,6.227)rectangle++(0.04,0.04)
  (1.622,6.824)rectangle++(0.04,0.04) (1.65,7.332)rectangle++(0.04,0.04) (1.693,7.762)rectangle++(0.04,0.04) (1.748,8.126)rectangle++(0.04,0.04) (1.808,8.434)rectangle++(0.04,0.04)
  (1.872,8.693)rectangle++(0.04,0.04) (1.937,8.91)rectangle++(0.04,0.04) (2.001,9.092)rectangle++(0.04,0.04) (2.062,9.245)rectangle++(0.04,0.04) (2.12,9.372)rectangle++(0.04,0.04)
  (-1.565,5.195)rectangle++(0.04,0.04) (-1.125,6.075)rectangle++(0.04,0.04) (-0.715,6.805)rectangle++(0.04,0.04) (-0.338,7.409)rectangle++(0.04,0.04) (0.007,7.907)rectangle++(0.04,0.04)
  (0.32,8.317)rectangle++(0.04,0.04) (0.602,8.653)rectangle++(0.04,0.04) (-1.685,4.769)rectangle++(0.04,0.04) (-1.244,5.717)rectangle++(0.04,0.04) (-0.831,6.505)rectangle++(0.04,0.04)
  (-0.448,7.159)rectangle++(0.04,0.04) (-0.097,7.699)rectangle++(0.04,0.04) (0.224,8.144)rectangle++(0.04,0.04) (0.514,8.51)rectangle++(0.04,0.04) (0.775,8.809)rectangle++(0.04,0.04)
  (1.009,9.053)rectangle++(0.04,0.04) (1.217,9.251)rectangle++(0.04,0.04) (1.403,9.411)rectangle++(0.04,0.04) (1.566,9.539)rectangle++(0.04,0.04) (1.711,9.642)rectangle++(0.04,0.04)
  (2.187,4.21)rectangle++(0.04,0.04) (2.025,5.087)rectangle++(0.04,0.04) (1.923,5.839)rectangle++(0.04,0.04) (1.866,6.483)rectangle++(0.04,0.04) (1.843,7.032)rectangle++(0.04,0.04)
  (1.846,7.499)rectangle++(0.04,0.04) (1.867,7.897)rectangle++(0.04,0.04) (1.9,8.234)rectangle++(0.04,0.04) (1.942,8.519)rectangle++(0.04,0.04) (1.99,8.76)rectangle++(0.04,0.04)
  (2.039,8.962)rectangle++(0.04,0.04) (1.98,4.153)rectangle++(0.04,0.04) (1.847,5.047)rectangle++(0.04,0.04) (1.77,5.812)rectangle++(0.04,0.04) (1.232,3.485)rectangle++(0.04,0.04)
  (1.184,4.509)rectangle++(0.04,0.04) (1.185,5.381)rectangle++(0.04,0.04) (1.22,6.123)rectangle++(0.04,0.04) (1.28,6.752)rectangle++(0.04,0.04) (1.356,7.284)rectangle++(0.04,0.04)
  (1.442,7.734)rectangle++(0.04,0.04) (1.746,3.781)rectangle++(0.04,0.04) (1.633,4.74)rectangle++(0.04,0.04) (-0.67,4.19)rectangle++(0.04,0.04) (-0.404,5.184)rectangle++(0.04,0.04)
  (-0.138,6.019)rectangle++(0.04,0.04) (0.121,6.718)rectangle++(0.04,0.04) (0.369,7.301)rectangle++(0.04,0.04) (0.604,7.788)rectangle++(0.04,0.04) (0.823,8.192)rectangle++(0.04,0.04)
  (1.025,8.526)rectangle++(0.04,0.04) (1.21,8.802)rectangle++(0.04,0.04) (1.378,9.03)rectangle++(0.04,0.04) (1.531,9.216)rectangle++(0.04,0.04) (1.667,9.369)rectangle++(0.04,0.04)
  (1.79,9.493)rectangle++(0.04,0.04) (1.899,9.594)rectangle++(0.04,0.04) (1.996,9.675)rectangle++(0.04,0.04) (2.081,9.74)rectangle++(0.04,0.04) (2.156,9.792)rectangle++(0.04,0.04)
  (2.222,9.833)rectangle++(0.04,0.04) (2.28,9.865)rectangle++(0.04,0.04) (2.331,9.891)rectangle++(0.04,0.04) (2.374,9.91)rectangle++(0.04,0.04) (2.176,4.412)rectangle++(0.04,0.04)
  (2.024,5.259)rectangle++(0.04,0.04) (1.928,5.986)rectangle++(0.04,0.04) (1.876,6.607)rectangle++(0.04,0.04) (1.857,7.137)rectangle++(0.04,0.04) (1.862,7.588)rectangle++(0.04,0.04)
  (1.884,7.972)rectangle++(0.04,0.04) (1.918,8.297)rectangle++(0.04,0.04) (1.96,8.572)rectangle++(0.04,0.04) (-1.424,5.007)rectangle++(0.04,0.04) (-1.013,5.909)rectangle++(0.04,0.04)
  (-0.627,6.66)rectangle++(0.04,0.04) (-1.482,4.145)rectangle++(0.04,0.04) (-1.096,5.179)rectangle++(0.04,0.04) (1.411,2.712)rectangle++(0.04,0.04) (1.305,3.845)rectangle++(0.04,0.04)
  (1.261,4.812)rectangle++(0.04,0.04) (1.262,5.636)rectangle++(0.04,0.04) (1.296,6.336)rectangle++(0.04,0.04) (1.353,6.93)rectangle++(0.04,0.04) (1.425,7.433)rectangle++(0.04,0.04)
  (1.506,7.857)rectangle++(0.04,0.04) (1.593,8.214)rectangle++(0.04,0.04) (1.68,8.515)rectangle++(0.04,0.04) (1.766,8.767)rectangle++(0.04,0.04) (1.85,8.977)rectangle++(0.04,0.04)
  (1.929,9.153)rectangle++(0.04,0.04) (-1.553,5.116)rectangle++(0.04,0.04) (-1.118,6.007)rectangle++(0.04,0.04) (-0.712,6.747)rectangle++(0.04,0.04) (-0.338,7.36)rectangle++(0.04,0.04)
  (0.005,7.865)rectangle++(0.04,0.04) (0.317,8.282)rectangle++(0.04,0.04) (0.598,8.623)rectangle++(0.04,0.04) (0.851,8.902)rectangle++(0.04,0.04) (1.078,9.129)rectangle++(0.04,0.04)
  (1.279,9.312)rectangle++(0.04,0.04) (1.457,9.461)rectangle++(0.04,0.04) (1.615,9.579)rectangle++(0.04,0.04) (1.754,9.674)rectangle++(0.04,0.04) (1.875,9.749)rectangle++(0.04,0.04)
  (1.982,9.808)rectangle++(0.04,0.04) (2.197,4.303)rectangle++(0.04,0.04) (2.037,5.166)rectangle++(0.04,0.04) (1.936,5.906)rectangle++(0.04,0.04) (1.236,3.541)rectangle++(0.04,0.04)
  (1.19,4.556)rectangle++(0.04,0.04) (1.191,5.421)rectangle++(0.04,0.04) (-0.894,4.241)rectangle++(0.04,0.04) (-0.592,5.237)rectangle++(0.04,0.04) (-0.296,6.071)rectangle++(0.04,0.04)
  (-0.011,6.769)rectangle++(0.04,0.04) (0.259,7.35)rectangle++(0.04,0.04) (0.512,7.833)rectangle++(0.04,0.04) (0.746,8.234)rectangle++(0.04,0.04) (0.962,8.565)rectangle++(0.04,0.04)
  (1.158,8.838)rectangle++(0.04,0.04) (1.335,9.062)rectangle++(0.04,0.04) (1.495,9.246)rectangle++(0.04,0.04) (1.639,9.395)rectangle++(0.04,0.04) (1.767,9.517)rectangle++(0.04,0.04)
  (1.88,9.615)rectangle++(0.04,0.04) (2.161,4.243)rectangle++(0.04,0.04) (-0.47,4.19)rectangle++(0.04,0.04) (-0.234,5.177)rectangle++(0.04,0.04) (0.006,6.006)rectangle++(0.04,0.04)
  (0.243,6.701)rectangle++(0.04,0.04) (0.473,7.282)rectangle++(0.04,0.04) (0.691,7.767)rectangle++(0.04,0.04) (0.896,8.171)rectangle++(0.04,0.04) (1.086,8.505)rectangle++(0.04,0.04)
  (1.261,8.782)rectangle++(0.04,0.04) (1.421,9.011)rectangle++(0.04,0.04) (-1.614,4.993)rectangle++(0.04,0.04) (-1.174,5.905)rectangle++(0.04,0.04) (-0.764,6.662)rectangle++(0.04,0.04)
  (1.739,3.073)rectangle++(0.04,0.04) (0.558,2.95)rectangle++(0.04,0.04) (0.59,4.081)rectangle++(0.04,0.04) (0.662,5.042)rectangle++(0.04,0.04) (0.762,5.855)rectangle++(0.04,0.04)
  (0.88,6.543)rectangle++(0.04,0.04) (1.008,7.122)rectangle++(0.04,0.04) (1.139,7.61)rectangle++(0.04,0.04) (1.27,8.019)rectangle++(0.04,0.04) (-0.003,1.948)rectangle++(0.04,0.04)
  (0.073,3.252)rectangle++(0.04,0.04) (0.19,4.357)rectangle++(0.04,0.04) (0.334,5.292)rectangle++(0.04,0.04) (0.493,6.081)rectangle++(0.04,0.04) (-1.159,4.241)rectangle++(0.04,0.04)
  (-0.818,5.247)rectangle++(0.04,0.04) (-0.487,6.089)rectangle++(0.04,0.04) (-0.173,6.791)rectangle++(0.04,0.04) (0.123,7.376)rectangle++(0.04,0.04) (0.397,7.861)rectangle++(0.04,0.04)
  (0.65,8.262)rectangle++(0.04,0.04) (1.119,2.185)rectangle++(0.04,0.04) (1.037,3.409)rectangle++(0.04,0.04) (1.015,4.452)rectangle++(0.04,0.04) (1.039,5.34)rectangle++(0.04,0.04)
  (1.095,6.094)rectangle++(0.04,0.04) (1.172,6.732)rectangle++(0.04,0.04) (1.263,7.271)rectangle++(0.04,0.04) (1.362,7.726)rectangle++(0.04,0.04) (-0.219,3.185)rectangle++(0.04,0.04)
  (-0.061,4.312)rectangle++(0.04,0.04) (0.118,5.264)rectangle++(0.04,0.04) (0.309,6.066)rectangle++(0.04,0.04) (0.503,6.74)rectangle++(0.04,0.04) (0.695,7.305)rectangle++(0.04,0.04)
  (0.881,7.777)rectangle++(0.04,0.04) (1.058,8.172)rectangle++(0.04,0.04) (1.224,8.5)rectangle++(0.04,0.04) (1.962,3.879)rectangle++(0.04,0.04) (-0.428,4.081)rectangle++(0.04,0.04)
  (-0.899,3.673)rectangle++(0.04,0.04) (1.015,2.466)rectangle++(0.04,0.04) (0.959,3.652)rectangle++(0.04,0.04) (0.959,4.662)rectangle++(0.04,0.04) (-0.774,4.048)rectangle++(0.04,0.04)
  (-0.498,5.068)rectangle++(0.04,0.04) (-0.223,5.924)rectangle++(0.04,0.04) (0.045,6.64)rectangle++(0.04,0.04) (0.302,7.239)rectangle++(0.04,0.04) (0.544,7.737)rectangle++(0.04,0.04)
  (0.77,8.151)rectangle++(0.04,0.04) (1.928,3.709)rectangle++(0.04,0.04) (1.785,4.671)rectangle++(0.04,0.04) (1.702,5.495)rectangle++(0.04,0.04) (1.664,6.199)rectangle++(0.04,0.04)
  (1.66,6.799)rectangle++(0.04,0.04) (1.681,7.309)rectangle++(0.04,0.04) (1.719,7.742)rectangle++(0.04,0.04) (1.769,8.108)rectangle++(0.04,0.04) (1.826,8.417)rectangle++(0.04,0.04)
  (1.886,8.678)rectangle++(0.04,0.04) (1.948,8.897)rectangle++(0.04,0.04) (2.01,9.081)rectangle++(0.04,0.04) (2.069,9.234)rectangle++(0.04,0.04) (-1.545,5.159)rectangle++(0.04,0.04)
  (-1.109,6.043)rectangle++(0.04,0.04) (-0.704,6.777)rectangle++(0.04,0.04) (-0.329,7.385)rectangle++(0.04,0.04) (0.013,7.886)rectangle++(0.04,0.04) (-1.628,4.507)rectangle++(0.04,0.04)
  (1.307,2.462)rectangle++(0.04,0.04) (1.207,3.637)rectangle++(0.04,0.04) (1.169,4.639)rectangle++(0.04,0.04) (1.177,5.493)rectangle++(0.04,0.04) (1.222,3.289)rectangle++(0.04,0.04)
  (1.168,4.343)rectangle++(0.04,0.04) (-0.669,4.024)rectangle++(0.04,0.04) (1.075,2.586)rectangle++(0.04,0.04) (1.015,3.752)rectangle++(0.04,0.04) (1.011,4.745)rectangle++(0.04,0.04)
  (1.047,5.589)rectangle++(0.04,0.04) (1.111,6.305)rectangle++(0.04,0.04) (1.194,6.911)rectangle++(0.04,0.04) (1.289,7.423)rectangle++(0.04,0.04) (1.391,7.854)rectangle++(0.04,0.04)
  (1.494,8.216)rectangle++(0.04,0.04) (1.596,8.52)rectangle++(0.04,0.04) (1.696,8.775)rectangle++(0.04,0.04) (-1.515,4.998)rectangle++(0.04,0.04) (-1.09,5.905)rectangle++(0.04,0.04)
  (-0.692,6.659)rectangle++(0.04,0.04) (-0.324,7.284)rectangle++(0.04,0.04) (0.013,7.801)rectangle++(0.04,0.04) (0.321,8.226)rectangle++(0.04,0.04) (1.999,3.621)rectangle++(0.04,0.04)
  (1.842,4.594)rectangle++(0.04,0.04) (1.747,5.428)rectangle++(0.04,0.04) (1.7,6.14)rectangle++(0.04,0.04) (0.412,2.743)rectangle++(0.04,0.04) (0.457,3.911)rectangle++(0.04,0.04)
  (-0.705,3.806)rectangle++(0.04,0.04) (-0.449,4.859)rectangle++(0.04,0.04) (-0.189,5.744)rectangle++(0.04,0.04) (0.067,6.487)rectangle++(0.04,0.04) (0.314,7.107)rectangle++(0.04,0.04)
  (0.549,7.625)rectangle++(0.04,0.04) (0.769,8.055)rectangle++(0.04,0.04) (1.905,3.69)rectangle++(0.04,0.04) (1.764,4.657)rectangle++(0.04,0.04) (1.684,5.484)rectangle++(0.04,0.04)
  (1.648,6.19)rectangle++(0.04,0.04) (1.34,3.53)rectangle++(0.04,0.04) (-0.466,3.896)rectangle++(0.04,0.04) (-0.867,3.629)rectangle++(0.04,0.04) (-0.594,4.716)rectangle++(0.04,0.04)
  (-1.064,3.767)rectangle++(0.04,0.04) (-0.947,3.492)rectangle++(0.04,0.04) (-0.667,4.602)rectangle++(0.04,0.04) (-0.385,5.535)rectangle++(0.04,0.04) (-1.2,3.968)rectangle++(0.04,0.04)
  (-0.864,5.017)rectangle++(0.04,0.04) (-0.536,5.895)rectangle++(0.04,0.04) (-0.222,6.629)rectangle++(0.04,0.04) (0.075,7.239)rectangle++(0.04,0.04) (0.308,2.838)rectangle++(0.04,0.04)
  (-0.504,3.566)rectangle++(0.04,0.04) (-0.288,4.647)rectangle++(0.04,0.04) (-0.061,5.558)rectangle++(0.04,0.04) (0.168,6.323)rectangle++(0.04,0.04) (0.394,6.964)rectangle++(0.04,0.04)
  (-1.364,4.396)rectangle++(0.04,0.04) (-0.986,5.388)rectangle++(0.04,0.04) (-0.624,6.215)rectangle++(0.04,0.04) (-0.284,6.904)rectangle++(0.04,0.04) (0.032,7.476)rectangle++(0.04,0.04)
  (0.324,7.949)rectangle++(0.04,0.04) (0.591,8.34)rectangle++(0.04,0.04) (1.996,3.705)rectangle++(0.04,0.04) (0.684,3.13)rectangle++(0.04,0.04) (0.704,4.229)rectangle++(0.04,0.04)
  (0.766,5.163)rectangle++(0.04,0.04) (0.855,5.954)rectangle++(0.04,0.04) (0.963,6.623)rectangle++(0.04,0.04) (1.081,7.187)rectangle++(0.04,0.04) (1.204,7.662)rectangle++(0.04,0.04)
  (1.328,8.061)rectangle++(0.04,0.04) (1.449,8.395)rectangle++(0.04,0.04) (1.565,8.674)rectangle++(0.04,0.04) (1.965,3.99)rectangle++(0.04,0.04) (1.827,4.909)rectangle++(0.04,0.04)
  (1.747,5.696)rectangle++(0.04,0.04) (1.711,6.368)rectangle++(0.04,0.04) (1.377,3.578)rectangle++(0.04,0.04) (1.311,4.583)rectangle++(0.04,0.04) (1.295,5.439)rectangle++(0.04,0.04)
  (0.015,2.322)rectangle++(0.04,0.04) (0.103,3.57)rectangle++(0.04,0.04) (-0.696,3.672)rectangle++(0.04,0.04) (-0.447,4.745)rectangle++(0.04,0.04) (-0.192,5.647)rectangle++(0.04,0.04)
  (0.06,6.404)rectangle++(0.04,0.04) (0.305,7.037)rectangle++(0.04,0.04) (0.539,7.566)rectangle++(0.04,0.04) (0.758,8.005)rectangle++(0.04,0.04) (0.031,2.671)rectangle++(0.04,0.04)
  (0.131,3.866)rectangle++(0.04,0.04) (0.264,4.877)rectangle++(0.04,0.04) (0.417,5.731)rectangle++(0.04,0.04) (0.582,6.451)rectangle++(0.04,0.04) (0.75,7.056)rectangle++(0.04,0.04)
  (0.918,7.564)rectangle++(0.04,0.04) (1.08,7.989)rectangle++(0.04,0.04) (1.236,8.343)rectangle++(0.04,0.04) (1.382,8.639)rectangle++(0.04,0.04) (1.977,3.942)rectangle++(0.04,0.04)
  (1.836,4.867)rectangle++(0.04,0.04) (1.753,5.66)rectangle++(0.04,0.04) (1.714,6.337)rectangle++(0.04,0.04) (1.709,6.914)rectangle++(0.04,0.04) (1.727,7.405)rectangle++(0.04,0.04)
  (1.762,7.821)rectangle++(0.04,0.04) (1.808,8.174)rectangle++(0.04,0.04) (1.862,8.472)rectangle++(0.04,0.04) (1.919,8.723)rectangle++(0.04,0.04) (1.978,8.934)rectangle++(0.04,0.04)
  (2.036,9.111)rectangle++(0.04,0.04) (-1.174,3.378)rectangle++(0.04,0.04) (-0.865,4.514)rectangle++(0.04,0.04) (1.219,2.636)rectangle++(0.04,0.04) (1.14,3.788)rectangle++(0.04,0.04)
  (1.118,4.771)rectangle++(0.04,0.04) (1.139,5.607)rectangle++(0.04,0.04) (1.19,6.316)rectangle++(0.04,0.04) (1.425,3.45)rectangle++(0.04,0.04) (-0.638,3.357)rectangle++(0.04,0.04)
  (-0.41,4.475)rectangle++(0.04,0.04) (-0.172,5.417)rectangle++(0.04,0.04) (-0.041,3.566)rectangle++(0.04,0.04) (0.106,4.629)rectangle++(0.04,0.04) (-0.92,3.881)rectangle++(0.04,0.04)
  (-0.629,4.932)rectangle++(0.04,0.04) (-0.34,5.814)rectangle++(0.04,0.04) (1.479,3.006)rectangle++(0.04,0.04) (1.376,4.092)rectangle++(0.04,0.04) (0.006,2.139)rectangle++(0.04,0.04)
  (0.089,3.414)rectangle++(0.04,0.04) (-0.665,3.638)rectangle++(0.04,0.04) (-0.422,4.715)rectangle++(0.04,0.04) (0.2,2.552)rectangle++(0.04,0.04) (0.27,3.758)rectangle++(0.04,0.04)
  (0.377,4.78)rectangle++(0.04,0.04) (0.51,5.644)rectangle++(0.04,0.04) (-1.063,4.157)rectangle++(0.04,0.04) (-0.739,5.172)rectangle++(0.04,0.04) (-0.424,6.022)rectangle++(0.04,0.04)
  (-0.122,6.732)rectangle++(0.04,0.04) (0.164,7.323)rectangle++(0.04,0.04) (0.43,7.815)rectangle++(0.04,0.04) (0.676,8.221)rectangle++(0.04,0.04) (0.901,8.557)rectangle++(0.04,0.04)
  (1.106,8.834)rectangle++(0.04,0.04) (1.291,9.061)rectangle++(0.04,0.04) (1.458,9.246)rectangle++(0.04,0.04) (1.607,9.397)rectangle++(0.04,0.04) (1.739,9.52)rectangle++(0.04,0.04)
  (1.857,9.618)rectangle++(0.04,0.04) (1.961,9.697)rectangle++(0.04,0.04) (2.053,9.761)rectangle++(0.04,0.04) (2.133,9.811)rectangle++(0.04,0.04) (2.203,9.85)rectangle++(0.04,0.04)
  (2.265,9.88)rectangle++(0.04,0.04) (2.318,9.904)rectangle++(0.04,0.04) (2.181,4.398)rectangle++(0.04,0.04) (2.028,5.247)rectangle++(0.04,0.04) (1.931,5.975)rectangle++(0.04,0.04)
  (1.878,6.598)rectangle++(0.04,0.04) (1.858,7.129)rectangle++(0.04,0.04) (1.862,7.582)rectangle++(0.04,0.04) (-1.232,4.794)rectangle++(0.04,0.04) (-0.857,5.72)rectangle++(0.04,0.04)
  (-1.324,3.916)rectangle++(0.04,0.04) (-0.971,4.978)rectangle++(0.04,0.04) (-0.628,5.866)rectangle++(0.04,0.04) (-0.302,6.608)rectangle++(0.04,0.04) (0.022,2.481)rectangle++(0.04,0.04)
  (0.116,3.704)rectangle++(0.04,0.04) (0.245,4.74)rectangle++(0.04,0.04) (0.395,5.616)rectangle++(0.04,0.04) (0.559,6.354)rectangle++(0.04,0.04) (0.727,6.975)rectangle++(0.04,0.04)
  (0.894,7.495)rectangle++(0.04,0.04) (-0.552,3.406)rectangle++(0.04,0.04) (0.908,2.494)rectangle++(0.04,0.04) (0.869,3.679)rectangle++(0.04,0.04) (0.884,4.689)rectangle++(0.04,0.04)
  (0.937,5.546)rectangle++(0.04,0.04) (1.016,6.273)rectangle++(0.04,0.04) (1.434,3.402)rectangle++(0.04,0.04) (1.353,4.431)rectangle++(0.04,0.04) (1.325,5.308)rectangle++(0.04,0.04)
  (1.337,6.055)rectangle++(0.04,0.04) (1.376,6.69)rectangle++(0.04,0.04) (1.435,7.227)rectangle++(0.04,0.04) (1.507,7.682)rectangle++(0.04,0.04) (-1.317,4.747)rectangle++(0.04,0.04)
  (-0.932,5.684)rectangle++(0.04,0.04) (-0.567,6.464)rectangle++(0.04,0.04) (-0.225,7.114)rectangle++(0.04,0.04) (0.091,7.652)rectangle++(0.04,0.04) (0.381,8.097)rectangle++(0.04,0.04)
  (0.646,8.463)rectangle++(0.04,0.04) (2.02,3.741)rectangle++(0.04,0.04) (-0.388,4.065)rectangle++(0.04,0.04) (-0.17,5.067)rectangle++(0.04,0.04) (-0.768,3.216)rectangle++(0.04,0.04)
  (-0.526,4.361)rectangle++(0.04,0.04) (-0.275,5.324)rectangle++(0.04,0.04) (-0.023,6.133)rectangle++(0.04,0.04) (0.223,6.81)rectangle++(0.04,0.04) (0.46,7.376)rectangle++(0.04,0.04)
  (0.684,7.847)rectangle++(0.04,0.04) (0.893,8.239)rectangle++(0.04,0.04) (1.086,8.563)rectangle++(0.04,0.04) (1.264,8.832)rectangle++(0.04,0.04) (1.425,9.053)rectangle++(0.04,0.04)
  (1.571,9.234)rectangle++(0.04,0.04) (1.703,9.382)rectangle++(0.04,0.04) (1.821,9.503)rectangle++(0.04,0.04) (1.925,9.601)rectangle++(0.04,0.04) (-1.649,5.204)rectangle++(0.04,0.04)
  (0.112,2.586)rectangle++(0.04,0.04) (0.197,3.79)rectangle++(0.04,0.04) (0.317,4.81)rectangle++(0.04,0.04) (0.459,5.672)rectangle++(0.04,0.04) (0.615,6.399)rectangle++(0.04,0.04)
  (0.777,7.011)rectangle++(0.04,0.04) (0.938,7.524)rectangle++(0.04,0.04) (1.096,7.954)rectangle++(0.04,0.04) (1.841,3.745)rectangle++(0.04,0.04) (1.713,4.706)rectangle++(0.04,0.04)
  (1.642,5.528)rectangle++(0.04,0.04) (1.614,6.229)rectangle++(0.04,0.04) (1.619,6.826)rectangle++(0.04,0.04) (1.501,3.646)rectangle++(0.04,0.04) (1.419,4.635)rectangle++(0.04,0.04)
  (0.982,3.178)rectangle++(0.04,0.04) (0.959,4.258)rectangle++(0.04,0.04) (0.984,5.178)rectangle++(0.04,0.04) (1.041,5.958)rectangle++(0.04,0.04) (1.121,6.619)rectangle++(0.04,0.04)
  (-1.16,4.464)rectangle++(0.04,0.04) (-1.112,3.611)rectangle++(0.04,0.04) (-0.803,4.71)rectangle++(0.04,0.04) (1.261,2.688)rectangle++(0.04,0.04) (-0.301,3.715)rectangle++(0.04,0.04)
  (-0.11,4.766)rectangle++(0.04,0.04) (0.095,5.651)rectangle++(0.04,0.04) (0.305,6.396)rectangle++(0.04,0.04) (0.513,7.021)rectangle++(0.04,0.04) (0.714,7.543)rectangle++(0.04,0.04)
  (0.907,7.979)rectangle++(0.04,0.04) (1.088,8.342)rectangle++(0.04,0.04) (1.256,8.644)rectangle++(0.04,0.04) (1.411,8.893)rectangle++(0.04,0.04) (1.553,9.099)rectangle++(0.04,0.04)
  (1.682,9.268)rectangle++(0.04,0.04) (1.798,9.407)rectangle++(0.04,0.04) (1.902,9.52)rectangle++(0.04,0.04) (1.996,9.612)rectangle++(0.04,0.04) (2.079,9.687)rectangle++(0.04,0.04)
  (2.152,9.747)rectangle++(0.04,0.04) (2.217,9.795)rectangle++(0.04,0.04) (2.274,9.833)rectangle++(0.04,0.04) (2.324,9.863)rectangle++(0.04,0.04) (2.368,9.887)rectangle++(0.04,0.04)
  (2.406,9.906)rectangle++(0.04,0.04) (2.439,9.92)rectangle++(0.04,0.04) (2.468,9.93)rectangle++(0.04,0.04) (2.493,9.938)rectangle++(0.04,0.04) (2.514,9.944)rectangle++(0.04,0.04)
  (2.168,4.45)rectangle++(0.04,0.04) (2.019,5.292)rectangle++(0.04,0.04) (1.925,6.014)rectangle++(0.04,0.04) (1.875,6.631)rectangle++(0.04,0.04) (1.857,7.158)rectangle++(0.04,0.04)
  (1.862,7.606)rectangle++(0.04,0.04) (-1.237,4.798)rectangle++(0.04,0.04) (-0.862,5.724)rectangle++(0.04,0.04) (-0.506,6.496)rectangle++(0.04,0.04) (-0.172,7.138)rectangle++(0.04,0.04)
  (0.026,2.554)rectangle++(0.04,0.04) (0.122,3.766)rectangle++(0.04,0.04) (0.252,4.792)rectangle++(0.04,0.04) (0.404,5.659)rectangle++(0.04,0.04) (0.89,2.206)rectangle++(0.04,0.04)
  (0.843,3.436)rectangle++(0.04,0.04) (0.852,4.483)rectangle++(0.04,0.04) (0.901,5.373)rectangle++(0.04,0.04) (0.979,6.127)rectangle++(0.04,0.04) (1.075,6.765)rectangle++(0.04,0.04)
  (1.549,3.511)rectangle++(0.04,0.04) (1.455,4.518)rectangle++(0.04,0.04) (1.415,5.379)rectangle++(0.04,0.04) (1.416,6.111)rectangle++(0.04,0.04) (1.446,6.734)rectangle++(0.04,0.04)
  (1.496,7.262)rectangle++(0.04,0.04) (1.623,3.702)rectangle++(0.04,0.04) (1.526,4.678)rectangle++(0.04,0.04) (1.482,5.512)rectangle++(0.04,0.04) (1.478,6.222)rectangle++(0.04,0.04)
  (1.503,6.826)rectangle++(0.04,0.04) (1.548,7.338)rectangle++(0.04,0.04) (1.607,7.771)rectangle++(0.04,0.04) (1.675,8.138)rectangle++(0.04,0.04) (1.747,8.446)rectangle++(0.04,0.04)
  (1.821,8.706)rectangle++(0.04,0.04) (1.893,8.923)rectangle++(0.04,0.04) (1.964,9.105)rectangle++(0.04,0.04) (2.025,4.164)rectangle++(0.04,0.04) (1.885,5.054)rectangle++(0.04,0.04)
  (1.031,3.365)rectangle++(0.04,0.04) (1.009,4.415)rectangle++(0.04,0.04) (1.032,5.309)rectangle++(0.04,0.04) (1.087,6.067)rectangle++(0.04,0.04) (1.165,6.71)rectangle++(0.04,0.04)
  (1.256,7.253)rectangle++(0.04,0.04) (1.649,3.646)rectangle++(0.04,0.04) (1.545,4.629)rectangle++(0.04,0.04) (1.496,5.469)rectangle++(0.04,0.04) (1.489,6.185)rectangle++(0.04,0.04)
  (1.51,6.794)rectangle++(0.04,0.04) (1.553,7.311)rectangle++(0.04,0.04) (-1.23,4.682)rectangle++(0.04,0.04) (-0.86,5.625)rectangle++(0.04,0.04) (-0.508,6.412)rectangle++(0.04,0.04)
  (-0.178,7.067)rectangle++(0.04,0.04) (0.129,7.61)rectangle++(0.04,0.04) (0.412,8.06)rectangle++(0.04,0.04) (0.671,8.43)rectangle++(0.04,0.04) (2.009,3.74)rectangle++(0.04,0.04)
  (1.855,4.695)rectangle++(0.04,0.04) (1.762,5.513)rectangle++(0.04,0.04) (1.159,3.425)rectangle++(0.04,0.04) (-0.632,3.335)rectangle++(0.04,0.04) (-0.406,4.456)rectangle++(0.04,0.04)
  (-0.169,5.4)rectangle++(0.04,0.04) (1.357,2.965)rectangle++(0.04,0.04) (1.27,4.062)rectangle++(0.04,0.04) (1.24,4.998)rectangle++(0.04,0.04) (1.252,5.795)rectangle++(0.04,0.04)
  (1.293,6.472)rectangle++(0.04,0.04) (1.356,7.046)rectangle++(0.04,0.04) (1.432,7.531)rectangle++(0.04,0.04) (1.517,7.94)rectangle++(0.04,0.04) (1.604,8.285)rectangle++(0.04,0.04)
  (1.693,8.574)rectangle++(0.04,0.04) (1.78,8.816)rectangle++(0.04,0.04) (1.863,9.019)rectangle++(0.04,0.04) (1.942,9.188)rectangle++(0.04,0.04) (2.048,4.175)rectangle++(0.04,0.04)
  (-0.476,4.156)rectangle++(0.04,0.04) (-0.24,5.148)rectangle++(0.04,0.04) (0,5.981)rectangle++(0.04,0.04) (0.237,6.68)rectangle++(0.04,0.04) (-1.332,4.311)rectangle++(0.04,0.04)
  (-0.962,5.314)rectangle++(0.04,0.04) (-0.607,6.151)rectangle++(0.04,0.04) (-0.272,6.849)rectangle++(0.04,0.04) (1.728,3.221)rectangle++(0.04,0.04) (1.595,4.265)rectangle++(0.04,0.04)
  (1.524,5.158)rectangle++(0.04,0.04) (1.5,5.919)rectangle++(0.04,0.04) (1.509,6.568)rectangle++(0.04,0.04) (1.45,3.571)rectangle++(0.04,0.04) (1.373,4.574)rectangle++(0.04,0.04)
  (1.348,5.429)rectangle++(0.04,0.04) (1.36,6.157)rectangle++(0.04,0.04) (-1.019,4.418)rectangle++(0.04,0.04) (-0.692,5.392)rectangle++(0.04,0.04) (-0.375,6.207)rectangle++(0.04,0.04)
  (-0.072,6.888)rectangle++(0.04,0.04) (0.212,7.453)rectangle++(0.04,0.04) (0.476,7.923)rectangle++(0.04,0.04) (0.719,8.312)rectangle++(0.04,0.04) (0.942,8.633)rectangle++(0.04,0.04)
  (1.144,8.896)rectangle++(0.04,0.04) (1.326,9.112)rectangle++(0.04,0.04) (1.489,9.289)rectangle++(0.04,0.04) (1.635,9.432)rectangle++(0.04,0.04) (1.765,9.548)rectangle++(0.04,0.04)
  (1.88,9.641)rectangle++(0.04,0.04) (1.981,9.716)rectangle++(0.04,0.04) (2.071,9.776)rectangle++(0.04,0.04) (2.149,9.823)rectangle++(0.04,0.04) (2.217,9.86)rectangle++(0.04,0.04)
  (2.277,9.888)rectangle++(0.04,0.04) (2.329,9.91)rectangle++(0.04,0.04) (2.374,9.927)rectangle++(0.04,0.04) (2.412,9.939)rectangle++(0.04,0.04) (2.446,9.948)rectangle++(0.04,0.04)
  (2.475,9.954)rectangle++(0.04,0.04) (2.499,9.958)rectangle++(0.04,0.04) (2.521,9.961)rectangle++(0.04,0.04) (2.539,9.962)rectangle++(0.04,0.04) (2.554,9.962)rectangle++(0.04,0.04)
  (2.567,9.962)rectangle++(0.04,0.04) (2.579,9.961)rectangle++(0.04,0.04) (2.588,9.96)rectangle++(0.04,0.04) (2.596,9.959)rectangle++(0.04,0.04) (2.603,9.957)rectangle++(0.04,0.04)
  (2.608,9.956)rectangle++(0.04,0.04) (2.613,9.954)rectangle++(0.04,0.04) (2.617,9.953)rectangle++(0.04,0.04) (2.621,9.951)rectangle++(0.04,0.04) (0.044,2.936)rectangle++(0.04,0.04)
  (0.152,4.09)rectangle++(0.04,0.04) (0.291,5.067)rectangle++(0.04,0.04) (0.448,5.891)rectangle++(0.04,0.04) (0.614,6.586)rectangle++(0.04,0.04) (0.783,7.17)rectangle++(0.04,0.04)
  (-1.338,4.51)rectangle++(0.04,0.04) (-0.959,5.483)rectangle++(0.04,0.04) (-0.598,6.295)rectangle++(0.04,0.04) (-1.4,4.078)rectangle++(0.04,0.04) (1.174,2.431)rectangle++(0.04,0.04)
  (1.093,3.616)rectangle++(0.04,0.04) (1.071,4.626)rectangle++(0.04,0.04) (1.093,5.485)rectangle++(0.04,0.04) (1.147,6.215)rectangle++(0.04,0.04) (-1.07,4.39)rectangle++(0.04,0.04)
  (-0.736,5.37)rectangle++(0.04,0.04) (-0.413,6.19)rectangle++(0.04,0.04) (-0.106,6.874)rectangle++(0.04,0.04) (-0.048,3.713)rectangle++(0.04,0.04) (0.925,2.668)rectangle++(0.04,0.04)
  (0.891,3.827)rectangle++(0.04,0.04) (0.908,4.813)rectangle++(0.04,0.04) (0.962,5.651)rectangle++(0.04,0.04) (0.016,2.351)rectangle++(0.04,0.04) (0.106,3.594)rectangle++(0.04,0.04)
  (0.231,4.647)rectangle++(0.04,0.04) (0.38,5.537)rectangle++(0.04,0.04) (0.543,6.287)rectangle++(0.04,0.04) (0.711,6.919)rectangle++(0.04,0.04) (0.878,7.449)rectangle++(0.04,0.04)
  (1.042,7.892)rectangle++(0.04,0.04) (1.2,8.263)rectangle++(0.04,0.04) (1.348,8.572)rectangle++(0.04,0.04) (1.486,8.828)rectangle++(0.04,0.04) (1.614,9.041)rectangle++(0.04,0.04)
  (1.732,9.216)rectangle++(0.04,0.04) (1.838,9.361)rectangle++(0.04,0.04) (1.935,9.479)rectangle++(0.04,0.04) (2.022,9.576)rectangle++(0.04,0.04) (2.099,9.655)rectangle++(0.04,0.04)
  (2.145,4.3)rectangle++(0.04,0.04) (1.993,5.166)rectangle++(0.04,0.04) (1.899,5.907)rectangle++(0.04,0.04) (1.848,6.542)rectangle++(0.04,0.04) (1.83,7.083)rectangle++(0.04,0.04)
  (1.837,7.543)rectangle++(0.04,0.04) (1.861,7.934)rectangle++(0.04,0.04) (1.897,8.266)rectangle++(0.04,0.04) (1.941,8.546)rectangle++(0.04,0.04) (1.989,8.783)rectangle++(0.04,0.04)
  (2.04,8.982)rectangle++(0.04,0.04) (2.091,9.15)rectangle++(0.04,0.04) (2.141,9.29)rectangle++(0.04,0.04) (2.189,9.407)rectangle++(0.04,0.04) (0.539,3.11)rectangle++(0.04,0.04)
  (0.58,4.218)rectangle++(0.04,0.04) (0.66,5.158)rectangle++(0.04,0.04) (0.765,5.954)rectangle++(0.04,0.04) (1.385,3.284)rectangle++(0.04,0.04) (1.306,4.332)rectangle++(0.04,0.04)
  (1.281,5.226)rectangle++(0.04,0.04) (1.144,3.261)rectangle++(0.04,0.04) (1.101,4.322)rectangle++(0.04,0.04) (1.106,5.226)rectangle++(0.04,0.04) (1.147,5.994)rectangle++(0.04,0.04)
  (1.213,6.645)rectangle++(0.04,0.04) (1.294,7.196)rectangle++(0.04,0.04) (1.386,7.661)rectangle++(0.04,0.04) (1.482,8.053)rectangle++(0.04,0.04) (-1.4,4.815)rectangle++(0.04,0.04)
  (-0.999,5.745)rectangle++(0.04,0.04) (-0.622,6.52)rectangle++(0.04,0.04) (-0.27,7.163)rectangle++(0.04,0.04) (0.055,7.695)rectangle++(0.04,0.04) (0.352,8.135)rectangle++(0.04,0.04)
  (0.623,8.497)rectangle++(0.04,0.04) (0.867,8.794)rectangle++(0.04,0.04) (1.086,9.036)rectangle++(0.04,0.04) (1.283,9.233)rectangle++(0.04,0.04) (2.136,4.034)rectangle++(0.04,0.04)
  (1.975,4.939)rectangle++(0.04,0.04) (1.874,5.716)rectangle++(0.04,0.04) (0.415,2.699)rectangle++(0.04,0.04) (0.458,3.874)rectangle++(0.04,0.04) (-0.088,3.063)rectangle++(0.04,0.04)
  (0.045,4.203)rectangle++(0.04,0.04) (0.204,5.167)rectangle++(0.04,0.04) (0.378,5.98)rectangle++(0.04,0.04) (0.559,6.664)rectangle++(0.04,0.04) (0.739,7.238)rectangle++(0.04,0.04)
  (0.916,7.719)rectangle++(0.04,0.04) (1.085,8.121)rectangle++(0.04,0.04) (1.245,8.456)rectangle++(0.04,0.04) (1.394,8.734)rectangle++(0.04,0.04) (1.532,8.964)rectangle++(0.04,0.04)
  (-1.584,5.004)rectangle++(0.04,0.04) (-1.149,5.913)rectangle++(0.04,0.04) (-0.742,6.669)rectangle++(0.04,0.04) (-0.366,7.294)rectangle++(0.04,0.04) (1.849,3.286)rectangle++(0.04,0.04)
  (1.701,4.315)rectangle++(0.04,0.04) (-0.567,4.118)rectangle++(0.04,0.04) (-0.32,5.12)rectangle++(0.04,0.04) (-0.069,5.961)rectangle++(0.04,0.04) (0.177,6.665)rectangle++(0.04,0.04)
  (0.415,7.255)rectangle++(0.04,0.04) (-0.091,3.764)rectangle++(0.04,0.04) (0.071,4.799)rectangle++(0.04,0.04) (0.25,5.672)rectangle++(0.04,0.04) (1.375,3.113)rectangle++(0.04,0.04)
  (1.291,4.187)rectangle++(0.04,0.04) (1.263,5.104)rectangle++(0.04,0.04) (1.275,5.884)rectangle++(0.04,0.04) (1.317,6.546)rectangle++(0.04,0.04) (1.379,7.108)rectangle++(0.04,0.04)
  (1.455,7.583)rectangle++(0.04,0.04) (1.537,7.983)rectangle++(0.04,0.04) (1.797,3.851)rectangle++(0.04,0.04) (1.679,4.797)rectangle++(0.04,0.04) (1.617,5.607)rectangle++(0.04,0.04)
  (1.596,6.297)rectangle++(0.04,0.04) (1.607,6.885)rectangle++(0.04,0.04) (1.639,7.384)rectangle++(0.04,0.04) (1.686,7.807)rectangle++(0.04,0.04) (1.743,8.165)rectangle++(0.04,0.04)
  (1.806,8.467)rectangle++(0.04,0.04) (1.872,8.721)rectangle++(0.04,0.04) (1.938,8.934)rectangle++(0.04,0.04) (1.985,4.125)rectangle++(0.04,0.04) (0.789,3.208)rectangle++(0.04,0.04)
  (0.797,4.292)rectangle++(0.04,0.04) (0.847,5.212)rectangle++(0.04,0.04) (0.926,5.993)rectangle++(0.04,0.04) (1.025,6.653)rectangle++(0.04,0.04) (1.135,7.21)rectangle++(0.04,0.04)
  (-1.283,4.584)rectangle++(0.04,0.04) (-0.91,5.544)rectangle++(0.04,0.04) (1.479,2.825)rectangle++(0.04,0.04) (-0.291,3.783)rectangle++(0.04,0.04) (-0.811,3.64)rectangle++(0.04,0.04)
  (-0.546,4.723)rectangle++(0.04,0.04) (-0.278,5.632)rectangle++(0.04,0.04) (-0.013,6.395)rectangle++(0.04,0.04) (0.243,7.032)rectangle++(0.04,0.04) (0.485,7.564)rectangle++(0.04,0.04)
  (0.713,8.006)rectangle++(0.04,0.04) (0.924,8.373)rectangle++(0.04,0.04) (1.118,8.676)rectangle++(0.04,0.04) (0.04,2.421)rectangle++(0.04,0.04) (0.129,3.652)rectangle++(0.04,0.04)
  (0.254,4.695)rectangle++(0.04,0.04) (1.133,2.925)rectangle++(0.04,0.04) (1.078,4.037)rectangle++(0.04,0.04) (1.076,4.985)rectangle++(0.04,0.04) (1.111,5.79)rectangle++(0.04,0.04)
  (-0.986,4.299)rectangle++(0.04,0.04) (1.18,2.567)rectangle++(0.04,0.04) (1.104,3.731)rectangle++(0.04,0.04) (1.085,4.724)rectangle++(0.04,0.04) (1.109,5.568)rectangle++(0.04,0.04)
  (1.163,6.285)rectangle++(0.04,0.04) (1.238,6.892)rectangle++(0.04,0.04) (1.326,7.404)rectangle++(0.04,0.04) (1.421,7.837)rectangle++(0.04,0.04) (1.519,8.201)rectangle++(0.04,0.04)
  (1.617,8.506)rectangle++(0.04,0.04) (1.712,8.762)rectangle++(0.04,0.04) (1.804,8.975)rectangle++(0.04,0.04) (1.89,9.153)rectangle++(0.04,0.04) (1.971,9.301)rectangle++(0.04,0.04)
  (2.045,9.423)rectangle++(0.04,0.04) (2.113,9.524)rectangle++(0.04,0.04) (2.175,9.607)rectangle++(0.04,0.04) (2.23,9.675)rectangle++(0.04,0.04) (2.281,9.731)rectangle++(0.04,0.04)
  (-1.613,5.296)rectangle++(0.04,0.04) (-1.161,6.162)rectangle++(0.04,0.04) (-0.743,6.881)rectangle++(0.04,0.04) (-0.358,7.475)rectangle++(0.04,0.04) (-0.008,7.964)rectangle++(0.04,0.04)
  (-0.036,3.015)rectangle++(0.04,0.04) (0.088,4.16)rectangle++(0.04,0.04) (0.239,5.129)rectangle++(0.04,0.04) (0.406,5.946)rectangle++(0.04,0.04) (0.581,6.634)rectangle++(0.04,0.04)
  (0.333,2.779)rectangle++(0.04,0.04) (0.392,3.945)rectangle++(0.04,0.04) (0.489,4.934)rectangle++(0.04,0.04) (0.611,5.771)rectangle++(0.04,0.04) (0.748,6.477)rectangle++(0.04,0.04)
  (0.892,7.072)rectangle++(0.04,0.04) (1.039,7.571)rectangle++(0.04,0.04) (1.184,7.99)rectangle++(0.04,0.04) (1.324,8.341)rectangle++(0.04,0.04) (1.457,8.633)rectangle++(0.04,0.04)
  (1.581,8.876)rectangle++(0.04,0.04) (1.697,9.077)rectangle++(0.04,0.04) (1.803,9.244)rectangle++(0.04,0.04) (1.9,9.382)rectangle++(0.04,0.04) (-1.607,5.156)rectangle++(0.04,0.04)
  (-1.162,6.043)rectangle++(0.04,0.04) (-0.748,6.779)rectangle++(0.04,0.04) (-0.367,7.388)rectangle++(0.04,0.04) (-0.018,7.891)rectangle++(0.04,0.04) (0.298,8.304)rectangle++(0.04,0.04)
  (0.583,8.643)rectangle++(0.04,0.04) (0.839,8.919)rectangle++(0.04,0.04) (1.068,9.144)rectangle++(0.04,0.04) (1.271,9.326)rectangle++(0.04,0.04) (1.451,9.472)rectangle++(0.04,0.04)
  (1.61,9.59)rectangle++(0.04,0.04) (2.186,4.177)rectangle++(0.04,0.04) (2.023,5.059)rectangle++(0.04,0.04) (1.92,5.815)rectangle++(0.04,0.04) (1.862,6.463)rectangle++(0.04,0.04)
  (1.382,3.631)rectangle++(0.04,0.04) (1.318,4.627)rectangle++(0.04,0.04) (1.303,5.477)rectangle++(0.04,0.04) (1.325,6.199)rectangle++(0.04,0.04) (1.372,6.813)rectangle++(0.04,0.04)
  (1.436,7.332)rectangle++(0.04,0.04) (1.512,7.771)rectangle++(0.04,0.04) (1.594,8.141)rectangle++(0.04,0.04) (1.829,3.894)rectangle++(0.04,0.04) (1.708,4.833)rectangle++(0.04,0.04)
  (1.643,5.636)rectangle++(0.04,0.04) (1.62,6.321)rectangle++(0.04,0.04) (1.628,6.904)rectangle++(0.04,0.04) (1.657,7.4)rectangle++(0.04,0.04) (1.703,7.82)rectangle++(0.04,0.04)
  (1.758,8.175)rectangle++(0.04,0.04) (1.819,8.474)rectangle++(0.04,0.04) (-1.429,4.962)rectangle++(0.04,0.04) (-1.018,5.871)rectangle++(0.04,0.04) (-0.633,6.627)rectangle++(0.04,0.04)
  (-0.275,7.255)rectangle++(0.04,0.04) (0.054,7.774)rectangle++(0.04,0.04) (0.355,8.202)rectangle++(0.04,0.04) (0.627,8.553)rectangle++(0.04,0.04) (2.045,3.754)rectangle++(0.04,0.04)
  (1.886,4.705)rectangle++(0.04,0.04) (1.789,5.52)rectangle++(0.04,0.04) (1.739,6.217)rectangle++(0.04,0.04) (1.725,6.811)rectangle++(0.04,0.04) (1.485,3.667)rectangle++(0.04,0.04)
  (1.406,4.654)rectangle++(0.04,0.04) (1.379,5.496)rectangle++(0.04,0.04) (1.39,6.212)rectangle++(0.04,0.04) (0.388,2.746)rectangle++(0.04,0.04) (0.438,3.915)rectangle++(0.04,0.04)
  (-0.709,3.803)rectangle++(0.04,0.04) (-0.891,3.565)rectangle++(0.04,0.04) (-0.873,3.485)rectangle++(0.04,0.04) (-0.605,4.593)rectangle++(0.04,0.04) (-0.332,5.524)rectangle++(0.04,0.04)
  (-0.064,6.305)rectangle++(0.04,0.04) (0.196,6.958)rectangle++(0.04,0.04) (0.443,7.503)rectangle++(0.04,0.04) (0.674,7.956)rectangle++(0.04,0.04) (0.889,8.332)rectangle++(0.04,0.04)
  (1.087,8.643)rectangle++(0.04,0.04) (1.267,8.899)rectangle++(0.04,0.04) (1.431,9.11)rectangle++(0.04,0.04) (1.578,9.282)rectangle++(0.04,0.04) (1.711,9.423)rectangle++(0.04,0.04)
  (1.829,9.537)rectangle++(0.04,0.04) (1.934,9.63)rectangle++(0.04,0.04) (2.027,9.704)rectangle++(0.04,0.04) (2.109,9.764)rectangle++(0.04,0.04) (2.181,9.811)rectangle++(0.04,0.04)
  (2.244,9.848)rectangle++(0.04,0.04) (2.299,9.877)rectangle++(0.04,0.04) (2.177,4.389)rectangle++(0.04,0.04) (0.832,3.303)rectangle++(0.04,0.04) (0.837,4.37)rectangle++(0.04,0.04)
  (0.884,5.277)rectangle++(0.04,0.04) (1.203,3.181)rectangle++(0.04,0.04) (-0.416,3.801)rectangle++(0.04,0.04) (0.99,2.601)rectangle++(0.04,0.04) (0.943,3.767)rectangle++(0.04,0.04)
  (0.004,2.095)rectangle++(0.04,0.04) (0.085,3.377)rectangle++(0.04,0.04) (0.205,4.463)rectangle++(0.04,0.04) (0.009,2.19)rectangle++(0.04,0.04) (0.093,3.457)rectangle++(0.04,0.04)
  (0.215,4.531)rectangle++(0.04,0.04) (0.362,5.439)rectangle++(0.04,0.04) (0.523,6.205)rectangle++(0.04,0.04) (0.69,6.849)rectangle++(0.04,0.04) (0.859,7.391)rectangle++(0.04,0.04)
  (-1.371,4.568)rectangle++(0.04,0.04) (-0.985,5.533)rectangle++(0.04,0.04) (-0.618,6.339)rectangle++(0.04,0.04) (0.021,2.445)rectangle++(0.04,0.04) (0.113,3.673)rectangle++(0.04,0.04)
  (0.241,4.714)rectangle++(0.04,0.04) (0.391,5.594)rectangle++(0.04,0.04) (0.554,6.335)rectangle++(0.04,0.04) (0.722,6.959)rectangle++(0.04,0.04) (0.89,7.482)rectangle++(0.04,0.04)
  (1.054,7.921)rectangle++(0.04,0.04) (1.21,8.287)rectangle++(0.04,0.04) (1.358,8.591)rectangle++(0.04,0.04) (1.496,8.845)rectangle++(0.04,0.04) (1.623,9.054)rectangle++(0.04,0.04)
  (0.483,3.062)rectangle++(0.04,0.04) (0.531,4.179)rectangle++(0.04,0.04) (0.616,5.128)rectangle++(0.04,0.04) (0.727,5.93)rectangle++(0.04,0.04) (0.853,6.608)rectangle++(0.04,0.04)
  (0.987,7.179)rectangle++(0.04,0.04) (1.124,7.659)rectangle++(0.04,0.04) (1.259,8.061)rectangle++(0.04,0.04) (1.849,3.802)rectangle++(0.04,0.04) (1.721,4.754)rectangle++(0.04,0.04)
  (1.651,5.569)rectangle++(0.04,0.04) (1.624,6.263)rectangle++(0.04,0.04) (1.629,6.855)rectangle++(0.04,0.04) (1.656,7.358)rectangle++(0.04,0.04) (-1.221,4.711)rectangle++(0.04,0.04)
  (-0.852,5.649)rectangle++(0.04,0.04) (-0.5,6.432)rectangle++(0.04,0.04) (-0.17,7.083)rectangle++(0.04,0.04) (0.137,7.624)rectangle++(0.04,0.04) (0.419,8.071)rectangle++(0.04,0.04)
  (0.677,8.44)rectangle++(0.04,0.04) (0.911,8.743)rectangle++(0.04,0.04) (1.121,8.991)rectangle++(0.04,0.04) (1.311,9.194)rectangle++(0.04,0.04) (1.48,9.359)rectangle++(0.04,0.04)
  (1.63,9.492)rectangle++(0.04,0.04) (1.763,9.599)rectangle++(0.04,0.04) (1.88,9.685)rectangle++(0.04,0.04) (-1.675,5.212)rectangle++(0.04,0.04) (-1.217,6.093)rectangle++(0.04,0.04)
  (-0.793,6.824)rectangle++(0.04,0.04) (-0.403,7.429)rectangle++(0.04,0.04) (-0.048,7.927)rectangle++(0.04,0.04) (-1.647,4.503)rectangle++(0.04,0.04) (-1.222,5.49)rectangle++(0.04,0.04)
  (-0.822,6.312)rectangle++(0.04,0.04) (-0.448,6.994)rectangle++(0.04,0.04) (1.784,3.209)rectangle++(0.04,0.04) (1.643,4.253)rectangle++(0.04,0.04) (-0.565,4.095)rectangle++(0.04,0.04)
  (-0.927,3.65)rectangle++(0.04,0.04) (-0.644,4.736)rectangle++(0.04,0.04) (-0.36,5.648)rectangle++(0.04,0.04) (-0.082,6.411)rectangle++(0.04,0.04) (0.184,7.049)rectangle++(0.04,0.04)
  (-1.419,4.374)rectangle++(0.04,0.04) (-1.034,5.371)rectangle++(0.04,0.04) (-0.666,6.203)rectangle++(0.04,0.04) (-0.32,6.895)rectangle++(0.04,0.04) (-0.937,3.194)rectangle++(0.04,0.04)
  (-0.671,4.349)rectangle++(0.04,0.04) (-0.398,5.319)rectangle++(0.04,0.04) (-1.157,3.923)rectangle++(0.04,0.04) (-0.829,4.977)rectangle++(0.04,0.04) (-0.507,5.86)rectangle++(0.04,0.04)
  (-0.199,6.597)rectangle++(0.04,0.04) (0.092,7.212)rectangle++(0.04,0.04) (0.365,7.723)rectangle++(0.04,0.04) (-1.53,4.541)rectangle++(0.04,0.04) (-1.121,5.517)rectangle++(0.04,0.04)
  (-0.734,6.33)rectangle++(0.04,0.04) (-0.373,7.006)rectangle++(0.04,0.04) (-0.039,7.567)rectangle++(0.04,0.04) (0.267,8.029)rectangle++(0.04,0.04) (0.546,8.41)rectangle++(0.04,0.04)
  (0.798,8.723)rectangle++(0.04,0.04) (1.025,8.979)rectangle++(0.04,0.04) (1.228,9.187)rectangle++(0.04,0.04) (1.41,9.356)rectangle++(0.04,0.04) (1.57,9.493)rectangle++(0.04,0.04)
  (1.712,9.602)rectangle++(0.04,0.04) (1.837,9.69)rectangle++(0.04,0.04) (1.947,9.759)rectangle++(0.04,0.04) (2.043,9.813)rectangle++(0.04,0.04) (2.127,9.856)rectangle++(0.04,0.04)
  (2.2,9.889)rectangle++(0.04,0.04) (2.263,9.913)rectangle++(0.04,0.04) (2.318,9.932)rectangle++(0.04,0.04) (2.366,9.946)rectangle++(0.04,0.04) (0.044,2.935)rectangle++(0.04,0.04)
  (0.152,4.089)rectangle++(0.04,0.04) (0.291,5.066)rectangle++(0.04,0.04) (0.448,5.891)rectangle++(0.04,0.04) (0.614,6.585)rectangle++(0.04,0.04) (0.783,7.169)rectangle++(0.04,0.04)
  (0.95,7.659)rectangle++(0.04,0.04) (1.112,8.068)rectangle++(0.04,0.04) (1.266,8.41)rectangle++(0.04,0.04) (1.153,2.238)rectangle++(0.04,0.04) (1.067,3.452)rectangle++(0.04,0.04)
  (1.043,4.488)rectangle++(0.04,0.04) (1.064,5.369)rectangle++(0.04,0.04) (1.117,6.118)rectangle++(0.04,0.04) (1.192,6.751)rectangle++(0.04,0.04) (1.281,7.287)rectangle++(0.04,0.04)
  (1.378,7.739)rectangle++(0.04,0.04) (1.479,8.119)rectangle++(0.04,0.04) (1.579,8.438)rectangle++(0.04,0.04) (1.678,8.706)rectangle++(0.04,0.04) (1.772,8.929)rectangle++(0.04,0.04)
  (1.861,9.115)rectangle++(0.04,0.04) (1.945,9.269)rectangle++(0.04,0.04) (2.021,9.397)rectangle++(0.04,0.04) (2.092,9.503)rectangle++(0.04,0.04) (2.156,9.59)rectangle++(0.04,0.04)
  (2.214,9.662)rectangle++(0.04,0.04) (2.266,9.72)rectangle++(0.04,0.04) (2.313,9.768)rectangle++(0.04,0.04) (2.354,9.806)rectangle++(0.04,0.04) (2.391,9.837)rectangle++(0.04,0.04)
  (2.424,9.862)rectangle++(0.04,0.04) (2.453,9.882)rectangle++(0.04,0.04) (2.478,9.898)rectangle++(0.04,0.04) (2.5,9.91)rectangle++(0.04,0.04) (2.161,4.441)rectangle++(0.04,0.04)
  (2.012,5.284)rectangle++(0.04,0.04) (1.92,6.007)rectangle++(0.04,0.04) (1.87,6.626)rectangle++(0.04,0.04) (1.852,7.153)rectangle++(0.04,0.04) (1.858,7.602)rectangle++(0.04,0.04)
  (1.881,7.984)rectangle++(0.04,0.04) (1.916,8.307)rectangle++(0.04,0.04) (1.959,8.581)rectangle++(0.04,0.04) (2.006,8.811)rectangle++(0.04,0.04) (2.056,9.006)rectangle++(0.04,0.04)
  (2.105,9.169)rectangle++(0.04,0.04) (-1.525,5.152)rectangle++(0.04,0.04) (-1.092,6.037)rectangle++(0.04,0.04) (-0.689,6.771)rectangle++(0.04,0.04) (-0.317,7.379)rectangle++(0.04,0.04)
  (0.279,2.847)rectangle++(0.04,0.04) (0.349,4.005)rectangle++(0.04,0.04) (0.455,4.986)rectangle++(0.04,0.04) (0.584,5.816)rectangle++(0.04,0.04) (0.727,6.517)rectangle++(0.04,0.04)
  (1.529,3.384)rectangle++(0.04,0.04) (1.433,4.411)rectangle++(0.04,0.04) (1.392,5.288)rectangle++(0.04,0.04) (1.392,6.036)rectangle++(0.04,0.04) (1.423,6.671)rectangle++(0.04,0.04)
  (1.474,7.21)rectangle++(0.04,0.04) (1.613,3.687)rectangle++(0.04,0.04) (-0.45,3.978)rectangle++(0.04,0.04) (-0.226,4.995)rectangle++(0.04,0.04) (0.006,5.851)rectangle++(0.04,0.04)
  (0.237,6.57)rectangle++(0.04,0.04) (1.599,3.283)rectangle++(0.04,0.04) (1.488,4.323)rectangle++(0.04,0.04) (1.436,5.211)rectangle++(0.04,0.04) (1.427,5.968)rectangle++(0.04,0.04)
  (0.627,2.095)rectangle++(0.04,0.04) (0.615,3.352)rectangle++(0.04,0.04) (0.654,4.421)rectangle++(0.04,0.04) (0.731,5.327)rectangle++(0.04,0.04) (0.832,6.095)rectangle++(0.04,0.04)
  (0.949,6.744)rectangle++(0.04,0.04) (1.074,7.291)rectangle++(0.04,0.04) (1.202,7.75)rectangle++(0.04,0.04) (1.33,8.136)rectangle++(0.04,0.04) (1.454,8.458)rectangle++(0.04,0.04)
  (1.572,8.728)rectangle++(0.04,0.04) (1.683,8.952)rectangle++(0.04,0.04) (1.786,9.138)rectangle++(0.04,0.04) (1.882,9.292)rectangle++(0.04,0.04) (0.04,2.846)rectangle++(0.04,0.04)
  (0.145,4.014)rectangle++(0.04,0.04) (0.282,5.002)rectangle++(0.04,0.04) (-0.078,3.535)rectangle++(0.04,0.04) (0.073,4.604)rectangle++(0.04,0.04) (0.244,5.507)rectangle++(0.04,0.04)
  (0.426,6.267)rectangle++(0.04,0.04) (0.61,6.907)rectangle++(0.04,0.04) (0.793,7.442)rectangle++(0.04,0.04) (0.969,7.89)rectangle++(0.04,0.04) (1.84,3.704)rectangle++(0.04,0.04)
  (1.71,4.671)rectangle++(0.04,0.04) (1.638,5.498)rectangle++(0.04,0.04) (1.61,6.204)rectangle++(0.04,0.04) (1.348,3.524)rectangle++(0.04,0.04) (1.284,4.538)rectangle++(0.04,0.04)
  (1.271,5.402)rectangle++(0.04,0.04) (1.294,6.137)rectangle++(0.04,0.04) (1.343,6.761)rectangle++(0.04,0.04) (1.41,7.289)rectangle++(0.04,0.04) (1.488,7.736)rectangle++(0.04,0.04)
  (1.572,8.112)rectangle++(0.04,0.04) (1.658,8.428)rectangle++(0.04,0.04) (1.745,8.694)rectangle++(0.04,0.04) (1.828,8.916)rectangle++(0.04,0.04) (1.909,9.102)rectangle++(0.04,0.04)
  (1.984,9.257)rectangle++(0.04,0.04) (2.055,9.385)rectangle++(0.04,0.04) (2.12,9.491)rectangle++(0.04,0.04) (2.102,4.273)rectangle++(0.04,0.04) (1.955,5.145)rectangle++(0.04,0.04)
  (0.013,2.282)rectangle++(0.04,0.04) (0.1,3.536)rectangle++(0.04,0.04) (-0.689,3.664)rectangle++(0.04,0.04) (-0.441,4.739)rectangle++(0.04,0.04) (-0.188,5.642)rectangle++(0.04,0.04)
  (0.064,6.399)rectangle++(0.04,0.04) (0.308,7.033)rectangle++(0.04,0.04) (0.541,7.562)rectangle++(0.04,0.04) (0.76,8.002)rectangle++(0.04,0.04) (-1.518,4.67)rectangle++(0.04,0.04);
\fill[piece2]
  (-1.106,5.626)rectangle++(0.04,0.04) (-0.717,6.423)rectangle++(0.04,0.04) (-0.355,7.084)rectangle++(0.04,0.04) (1.796,3.248)rectangle++(0.04,0.04) (1.654,4.285)rectangle++(0.04,0.04)
  (1.575,5.172)rectangle++(0.04,0.04) (1.544,5.93)rectangle++(0.04,0.04) (1.547,6.575)rectangle++(0.04,0.04) (1.576,7.123)rectangle++(0.04,0.04) (1.622,7.587)rectangle++(0.04,0.04)
  (1.68,7.981)rectangle++(0.04,0.04) (1.745,8.313)rectangle++(0.04,0.04) (1.814,8.592)rectangle++(0.04,0.04) (1.883,8.827)rectangle++(0.04,0.04) (1.965,4.092)rectangle++(0.04,0.04)
  (1.832,4.996)rectangle++(0.04,0.04) (1.023,3.341)rectangle++(0.04,0.04) (1.001,4.395)rectangle++(0.04,0.04) (0.971,3.037)rectangle++(0.04,0.04) (0.945,4.139)rectangle++(0.04,0.04)
  (0.914,2.975)rectangle++(0.04,0.04) (0.894,4.088)rectangle++(0.04,0.04) (0.773,2.276)rectangle++(0.04,0.04) (0.746,3.5)rectangle++(0.04,0.04) (0.779,2.807)rectangle++(0.04,0.04)
  (0.773,3.951)rectangle++(0.04,0.04) (0.813,4.923)rectangle++(0.04,0.04) (0.885,5.749)rectangle++(0.04,0.04) (0.98,6.447)rectangle++(0.04,0.04) (1.089,7.037)rectangle++(0.04,0.04)
  (-1.255,4.541)rectangle++(0.04,0.04) (-0.705,3.129)rectangle++(0.04,0.04) (-0.476,4.284)rectangle++(0.04,0.04) (-0.236,5.257)rectangle++(0.04,0.04) (0.008,6.074)rectangle++(0.04,0.04)
  (0.247,6.759)rectangle++(0.04,0.04) (0.478,7.331)rectangle++(0.04,0.04) (0.698,7.809)rectangle++(0.04,0.04) (0.903,8.206)rectangle++(0.04,0.04) (1.094,8.535)rectangle++(0.04,0.04)
  (1.269,8.807)rectangle++(0.04,0.04) (1.429,9.031)rectangle++(0.04,0.04) (1.573,9.216)rectangle++(0.04,0.04) (1.704,9.367)rectangle++(0.04,0.04) (1.821,9.49)rectangle++(0.04,0.04)
  (1.925,9.59)rectangle++(0.04,0.04) (2.018,9.67)rectangle++(0.04,0.04) (-1.647,5.235)rectangle++(0.04,0.04) (-1.193,6.112)rectangle++(0.04,0.04) (-0.772,6.839)rectangle++(0.04,0.04)
  (-0.385,7.44)rectangle++(0.04,0.04) (-0.031,7.936)rectangle++(0.04,0.04) (0.288,8.343)rectangle++(0.04,0.04) (0.577,8.676)rectangle++(0.04,0.04) (0.835,8.948)rectangle++(0.04,0.04)
  (1.066,9.168)rectangle++(0.04,0.04) (1.27,9.347)rectangle++(0.04,0.04) (1.451,9.49)rectangle++(0.04,0.04) (1.611,9.605)rectangle++(0.04,0.04) (2.19,4.18)rectangle++(0.04,0.04)
  (-0.451,4.183)rectangle++(0.04,0.04) (-0.218,5.17)rectangle++(0.04,0.04) (0.019,6)rectangle++(0.04,0.04) (0.254,6.695)rectangle++(0.04,0.04) (0.481,7.277)rectangle++(0.04,0.04)
  (0.698,7.762)rectangle++(0.04,0.04) (0.902,8.166)rectangle++(0.04,0.04) (1.091,8.501)rectangle++(0.04,0.04) (1.265,8.779)rectangle++(0.04,0.04) (1.424,9.008)rectangle++(0.04,0.04)
  (1.569,9.196)rectangle++(0.04,0.04) (1.699,9.35)rectangle++(0.04,0.04) (1.816,9.476)rectangle++(0.04,0.04) (1.92,9.578)rectangle++(0.04,0.04) (2.013,9.66)rectangle++(0.04,0.04)
  (2.095,9.727)rectangle++(0.04,0.04) (2.163,4.313)rectangle++(0.04,0.04) (-0.501,3.214)rectangle++(0.04,0.04) (-0.729,3.488)rectangle++(0.04,0.04) (-0.482,4.59)rectangle++(0.04,0.04)
  (-0.228,5.517)rectangle++(0.04,0.04) (0.237,2.644)rectangle++(0.04,0.04) (0.305,3.834)rectangle++(0.04,0.04) (-0.716,3.762)rectangle++(0.04,0.04) (-0.46,4.822)rectangle++(0.04,0.04)
  (-0.2,5.714)rectangle++(0.04,0.04) (0.056,6.461)rectangle++(0.04,0.04) (0.304,7.086)rectangle++(0.04,0.04) (0.54,7.607)rectangle++(0.04,0.04) (0.761,8.04)rectangle++(0.04,0.04)
  (-1.527,4.677)rectangle++(0.04,0.04) (-1.113,5.633)rectangle++(0.04,0.04) (-0.723,6.429)rectangle++(0.04,0.04) (1.677,3.038)rectangle++(0.04,0.04) (1.544,4.111)rectangle++(0.04,0.04)
  (1.475,5.029)rectangle++(0.04,0.04) (1.453,5.812)rectangle++(0.04,0.04) (1.465,6.478)rectangle++(0.04,0.04) (1.502,7.044)rectangle++(0.04,0.04) (1.568,3.662)rectangle++(0.04,0.04)
  (1.477,4.646)rectangle++(0.04,0.04) (1.439,5.486)rectangle++(0.04,0.04) (1.441,6.202)rectangle++(0.04,0.04) (1.471,6.81)rectangle++(0.04,0.04) (1.52,7.326)rectangle++(0.04,0.04)
  (1.583,7.762)rectangle++(0.04,0.04) (1.654,8.131)rectangle++(0.04,0.04) (1.729,8.441)rectangle++(0.04,0.04) (1.805,8.702)rectangle++(0.04,0.04) (1.88,8.921)rectangle++(0.04,0.04)
  (1.953,9.104)rectangle++(0.04,0.04) (-1.172,3.371)rectangle++(0.04,0.04) (-0.864,4.508)rectangle++(0.04,0.04) (-0.556,5.463)rectangle++(0.04,0.04) (-0.256,6.262)rectangle++(0.04,0.04)
  (0.03,6.929)rectangle++(0.04,0.04) (0.301,7.485)rectangle++(0.04,0.04) (0.553,7.946)rectangle++(0.04,0.04) (0.786,8.328)rectangle++(0.04,0.04) (0.999,8.644)rectangle++(0.04,0.04)
  (1.192,8.903)rectangle++(0.04,0.04) (1.368,9.116)rectangle++(0.04,0.04) (1.525,9.29)rectangle++(0.04,0.04) (1.666,9.432)rectangle++(0.04,0.04) (1.791,9.547)rectangle++(0.04,0.04)
  (-1.662,5.168)rectangle++(0.04,0.04) (-1.208,6.055)rectangle++(0.04,0.04) (-0.787,6.792)rectangle++(0.04,0.04) (-0.399,7.401)rectangle++(0.04,0.04) (-0.045,7.903)rectangle++(0.04,0.04)
  (0.275,8.315)rectangle++(0.04,0.04) (0.564,8.653)rectangle++(0.04,0.04) (0.824,8.929)rectangle++(0.04,0.04) (1.055,9.153)rectangle++(0.04,0.04) (1.261,9.334)rectangle++(0.04,0.04)
  (1.231,2.222)rectangle++(0.04,0.04) (1.133,3.436)rectangle++(0.04,0.04) (1.098,4.471)rectangle++(0.04,0.04) (1.11,5.353)rectangle++(0.04,0.04) (1.156,6.102)rectangle++(0.04,0.04)
  (1.224,6.736)rectangle++(0.04,0.04) (1.308,7.273)rectangle++(0.04,0.04) (1.4,7.726)rectangle++(0.04,0.04) (1.497,8.107)rectangle++(0.04,0.04) (1.595,8.428)rectangle++(0.04,0.04)
  (1.69,8.696)rectangle++(0.04,0.04) (1.782,8.92)rectangle++(0.04,0.04) (1.87,9.107)rectangle++(0.04,0.04) (-1.554,5.096)rectangle++(0.04,0.04) (-1.119,5.99)rectangle++(0.04,0.04)
  (-0.714,6.732)rectangle++(0.04,0.04) (-0.34,7.347)rectangle++(0.04,0.04) (-1.577,4.334)rectangle++(0.04,0.04) (-1.169,5.343)rectangle++(0.04,0.04) (-1.3,3.783)rectangle++(0.04,0.04)
  (1.089,2.397)rectangle++(0.04,0.04) (1.019,3.59)rectangle++(0.04,0.04) (1.008,4.607)rectangle++(0.04,0.04) (1.039,5.472)rectangle++(0.04,0.04) (1.1,6.206)rectangle++(0.04,0.04)
  (1.181,6.827)rectangle++(0.04,0.04) (1.275,7.352)rectangle++(0.04,0.04) (1.375,7.794)rectangle++(0.04,0.04) (1.478,8.166)rectangle++(0.04,0.04) (1.581,8.478)rectangle++(0.04,0.04)
  (1.681,8.74)rectangle++(0.04,0.04) (1.776,8.958)rectangle++(0.04,0.04) (1.866,9.139)rectangle++(0.04,0.04) (1.949,9.29)rectangle++(0.04,0.04) (2.026,9.415)rectangle++(0.04,0.04)
  (2.097,9.518)rectangle++(0.04,0.04) (2.161,9.602)rectangle++(0.04,0.04) (2.219,9.672)rectangle++(0.04,0.04) (2.27,9.728)rectangle++(0.04,0.04) (2.317,9.775)rectangle++(0.04,0.04)
  (2.358,9.812)rectangle++(0.04,0.04) (2.395,9.842)rectangle++(0.04,0.04) (2.427,9.866)rectangle++(0.04,0.04) (2.455,9.885)rectangle++(0.04,0.04) (2.48,9.9)rectangle++(0.04,0.04)
  (0.002,2.052)rectangle++(0.04,0.04) (0.082,3.34)rectangle++(0.04,0.04) (0.201,4.432)rectangle++(0.04,0.04) (0.346,5.356)rectangle++(0.04,0.04) (0.506,6.135)rectangle++(0.04,0.04)
  (0.673,6.79)rectangle++(0.04,0.04) (1.603,3.424)rectangle++(0.04,0.04) (1.497,4.443)rectangle++(0.04,0.04) (1.448,5.313)rectangle++(0.04,0.04) (1.441,6.054)rectangle++(0.04,0.04)
  (1.465,6.685)rectangle++(0.04,0.04) (1.51,7.219)rectangle++(0.04,0.04) (1.57,7.672)rectangle++(0.04,0.04) (1.64,8.055)rectangle++(0.04,0.04) (1.714,8.377)rectangle++(0.04,0.04)
  (1.789,8.648)rectangle++(0.04,0.04) (-1.471,4.99)rectangle++(0.04,0.04) (1.408,2.591)rectangle++(0.04,0.04) (1.298,3.742)rectangle++(0.04,0.04) (1.251,4.725)rectangle++(0.04,0.04)
  (1.25,5.562)rectangle++(0.04,0.04) (1.283,6.274)rectangle++(0.04,0.04) (1.403,3.463)rectangle++(0.04,0.04) (1.329,4.484)rectangle++(0.04,0.04) (1.307,5.354)rectangle++(0.04,0.04)
  (1.323,6.095)rectangle++(0.04,0.04) (1.366,6.724)rectangle++(0.04,0.04) (1.428,7.257)rectangle++(0.04,0.04) (1.502,7.708)rectangle++(0.04,0.04) (1.733,3.789)rectangle++(0.04,0.04)
  (1.622,4.748)rectangle++(0.04,0.04) (-0.673,4.189)rectangle++(0.04,0.04) (-0.407,5.184)rectangle++(0.04,0.04) (-0.141,6.019)rectangle++(0.04,0.04) (0.119,6.718)rectangle++(0.04,0.04)
  (0.368,7.302)rectangle++(0.04,0.04) (0.602,7.788)rectangle++(0.04,0.04) (0.821,8.192)rectangle++(0.04,0.04) (1.024,8.526)rectangle++(0.04,0.04) (-1.106,3.297)rectangle++(0.04,0.04)
  (-0.811,4.443)rectangle++(0.04,0.04) (-0.514,5.405)rectangle++(0.04,0.04) (-0.222,6.211)rectangle++(0.04,0.04) (0.057,6.884)rectangle++(0.04,0.04) (0.322,7.446)rectangle++(0.04,0.04)
  (0.028,2.595)rectangle++(0.04,0.04) (0.125,3.801)rectangle++(0.04,0.04) (0.256,4.822)rectangle++(0.04,0.04) (-0.05,3.104)rectangle++(0.04,0.04) (-0.624,3.551)rectangle++(0.04,0.04)
  (-0.391,4.64)rectangle++(0.04,0.04) (-0.149,5.556)rectangle++(0.04,0.04) (0.094,6.324)rectangle++(0.04,0.04) (0.33,6.968)rectangle++(0.04,0.04) (0.557,7.506)rectangle++(0.04,0.04)
  (0.772,7.954)rectangle++(0.04,0.04) (0.972,8.326)rectangle++(0.04,0.04) (1.157,8.635)rectangle++(0.04,0.04) (1.327,8.889)rectangle++(0.04,0.04) (1.481,9.099)rectangle++(0.04,0.04)
  (1.621,9.271)rectangle++(0.04,0.04) (1.746,9.412)rectangle++(0.04,0.04) (1.859,9.526)rectangle++(0.04,0.04) (1.959,9.619)rectangle++(0.04,0.04) (2.047,9.694)rectangle++(0.04,0.04)
  (2.126,9.754)rectangle++(0.04,0.04) (2.195,9.802)rectangle++(0.04,0.04) (2.256,9.84)rectangle++(0.04,0.04) (2.173,4.372)rectangle++(0.04,0.04) (2.019,5.225)rectangle++(0.04,0.04)
  (1.923,5.957)rectangle++(0.04,0.04) (1.871,6.583)rectangle++(0.04,0.04) (1.851,7.117)rectangle++(0.04,0.04) (1.856,7.571)rectangle++(0.04,0.04) (1.878,7.958)rectangle++(0.04,0.04)
  (1.913,8.285)rectangle++(0.04,0.04) (1.955,8.562)rectangle++(0.04,0.04) (1.891,4.057)rectangle++(0.04,0.04) (1.768,4.969)rectangle++(0.04,0.04) (1.699,5.749)rectangle++(0.04,0.04)
  (1.672,6.415)rectangle++(0.04,0.04) (1.676,6.982)rectangle++(0.04,0.04) (-1.138,4.64)rectangle++(0.04,0.04) (-0.784,5.586)rectangle++(0.04,0.04) (-0.445,6.376)rectangle++(0.04,0.04)
  (-0.125,7.033)rectangle++(0.04,0.04) (0.172,7.579)rectangle++(0.04,0.04) (0.448,8.032)rectangle++(0.04,0.04) (0.7,8.405)rectangle++(0.04,0.04) (0.929,8.713)rectangle++(0.04,0.04)
  (1.136,8.965)rectangle++(0.04,0.04) (1.322,9.171)rectangle++(0.04,0.04) (1.488,9.339)rectangle++(0.04,0.04) (0.04,2.853)rectangle++(0.04,0.04) (0.146,4.019)rectangle++(0.04,0.04)
  (0.283,5.007)rectangle++(0.04,0.04) (0.438,5.841)rectangle++(0.04,0.04) (0.604,6.543)rectangle++(0.04,0.04) (0.773,7.134)rectangle++(0.04,0.04) (0.94,7.629)rectangle++(0.04,0.04)
  (1.779,3.647)rectangle++(0.04,0.04) (1.656,4.625)rectangle++(0.04,0.04) (1.59,5.461)rectangle++(0.04,0.04) (1.568,6.174)rectangle++(0.04,0.04) (1.578,6.782)rectangle++(0.04,0.04)
  (-1.113,4.582)rectangle++(0.04,0.04) (-0.765,5.536)rectangle++(0.04,0.04) (-0.431,6.332)rectangle++(0.04,0.04) (0.021,2.444)rectangle++(0.04,0.04) (0.113,3.673)rectangle++(0.04,0.04)
  (0.896,2.696)rectangle++(0.04,0.04) (0.868,3.852)rectangle++(0.04,0.04) (0.889,4.836)rectangle++(0.04,0.04) (0.947,5.671)rectangle++(0.04,0.04) (1.03,6.379)rectangle++(0.04,0.04)
  (1.128,6.977)rectangle++(0.04,0.04) (1.236,7.481)rectangle++(0.04,0.04) (1.348,7.906)rectangle++(0.04,0.04) (1.459,8.262)rectangle++(0.04,0.04) (1.569,8.561)rectangle++(0.04,0.04)
  (1.674,8.81)rectangle++(0.04,0.04) (1.773,9.018)rectangle++(0.04,0.04) (-0.612,3.427)rectangle++(0.04,0.04) (-0.385,4.534)rectangle++(0.04,0.04) (-0.148,5.465)rectangle++(0.04,0.04)
  (0.09,6.248)rectangle++(0.04,0.04) (0.325,6.903)rectangle++(0.04,0.04) (0.55,7.451)rectangle++(0.04,0.04) (0.763,7.907)rectangle++(0.04,0.04) (0.963,8.287)rectangle++(0.04,0.04)
  (1.148,8.602)rectangle++(0.04,0.04) (1.317,8.862)rectangle++(0.04,0.04) (1.472,9.076)rectangle++(0.04,0.04) (1.612,9.252)rectangle++(0.04,0.04) (1.738,9.396)rectangle++(0.04,0.04)
  (1.851,9.513)rectangle++(0.04,0.04) (1.952,9.608)rectangle++(0.04,0.04) (2.041,9.685)rectangle++(0.04,0.04) (2.159,4.293)rectangle++(0.04,0.04) (2.005,5.159)rectangle++(0.04,0.04)
  (1.908,5.901)rectangle++(0.04,0.04) (1.856,6.536)rectangle++(0.04,0.04) (1.837,7.077)rectangle++(0.04,0.04) (1.842,7.538)rectangle++(0.04,0.04) (1.865,7.93)rectangle++(0.04,0.04)
  (1.9,8.262)rectangle++(0.04,0.04) (1.944,8.543)rectangle++(0.04,0.04) (1.992,8.78)rectangle++(0.04,0.04) (-1.463,5.054)rectangle++(0.04,0.04) (-1.043,5.951)rectangle++(0.04,0.04)
  (-0.651,6.696)rectangle++(0.04,0.04) (-0.288,7.314)rectangle++(0.04,0.04) (0.046,7.825)rectangle++(0.04,0.04) (0.35,8.245)rectangle++(0.04,0.04) (0.625,8.591)rectangle++(0.04,0.04)
  (0.873,8.873)rectangle++(0.04,0.04) (-1.683,4.863)rectangle++(0.04,0.04) (-1.238,5.797)rectangle++(0.04,0.04) (-1.409,3.86)rectangle++(0.04,0.04) (-1.045,4.933)rectangle++(0.04,0.04)
  (0.78,2.088)rectangle++(0.04,0.04) (-0.26,3.506)rectangle++(0.04,0.04) (-0.083,4.587)rectangle++(0.04,0.04) (-0.945,3.838)rectangle++(0.04,0.04) (-0.652,4.896)rectangle++(0.04,0.04)
  (-0.36,5.784)rectangle++(0.04,0.04) (-0.077,6.527)rectangle++(0.04,0.04) (0.022,2.47)rectangle++(0.04,0.04) (0.115,3.695)rectangle++(0.04,0.04) (0.244,4.732)rectangle++(0.04,0.04)
  (0.394,5.609)rectangle++(0.04,0.04) (0.557,6.348)rectangle++(0.04,0.04) (0.725,6.97)rectangle++(0.04,0.04) (0.893,7.492)rectangle++(0.04,0.04) (1.057,7.928)rectangle++(0.04,0.04)
  (1.213,8.293)rectangle++(0.04,0.04) (1.361,8.597)rectangle++(0.04,0.04) (1.498,8.849)rectangle++(0.04,0.04) (2.016,4.009)rectangle++(0.04,0.04) (1.872,4.923)rectangle++(0.04,0.04)
  (1.786,5.706)rectangle++(0.04,0.04) (1.204,3.468)rectangle++(0.04,0.04) (1.16,4.496)rectangle++(0.04,0.04) (1.163,5.371)rectangle++(0.04,0.04) (1.202,6.115)rectangle++(0.04,0.04)
  (1.264,6.746)rectangle++(0.04,0.04) (1.342,7.28)rectangle++(0.04,0.04) (1.43,7.73)rectangle++(0.04,0.04) (1.522,8.11)rectangle++(0.04,0.04) (1.616,8.429)rectangle++(0.04,0.04)
  (1.709,8.696)rectangle++(0.04,0.04) (1.798,8.919)rectangle++(0.04,0.04) (1.883,9.106)rectangle++(0.04,0.04) (1.962,9.261)rectangle++(0.04,0.04) (2.036,9.389)rectangle++(0.04,0.04)
  (-1.584,5.183)rectangle++(0.04,0.04) (-1.141,6.065)rectangle++(0.04,0.04) (-0.73,6.797)rectangle++(0.04,0.04) (0.004,3.691)rectangle++(0.04,0.04) (0.149,4.734)rectangle++(0.04,0.04)
  (0.313,5.614)rectangle++(0.04,0.04) (0.489,6.355)rectangle++(0.04,0.04) (0.668,6.979)rectangle++(0.04,0.04) (0.844,7.501)rectangle++(0.04,0.04) (1.016,7.939)rectangle++(0.04,0.04)
  (1.179,8.303)rectangle++(0.04,0.04) (1.332,8.607)rectangle++(0.04,0.04) (1.474,8.859)rectangle++(0.04,0.04) (1.605,9.067)rectangle++(0.04,0.04) (0.038,2.816)rectangle++(0.04,0.04)
  (0.143,3.988)rectangle++(0.04,0.04) (0.279,4.98)rectangle++(0.04,0.04) (0.434,5.818)rectangle++(0.04,0.04) (0.599,6.524)rectangle++(0.04,0.04) (0.768,7.118)rectangle++(0.04,0.04)
  (0.936,7.616)rectangle++(0.04,0.04) (1.098,8.032)rectangle++(0.04,0.04) (1.252,8.38)rectangle++(0.04,0.04) (1.397,8.669)rectangle++(0.04,0.04) (1.532,8.909)rectangle++(0.04,0.04)
  (1.657,9.107)rectangle++(0.04,0.04) (1.77,9.271)rectangle++(0.04,0.04) (1.873,9.406)rectangle++(0.04,0.04) (1.966,9.516)rectangle++(0.04,0.04) (2.05,9.606)rectangle++(0.04,0.04)
  (2.124,9.68)rectangle++(0.04,0.04) (2.191,9.739)rectangle++(0.04,0.04) (2.249,9.787)rectangle++(0.04,0.04) (2.301,9.825)rectangle++(0.04,0.04) (2.347,9.855)rectangle++(0.04,0.04)
  (2.166,4.395)rectangle++(0.04,0.04) (2.014,5.246)rectangle++(0.04,0.04) (1.92,5.974)rectangle++(0.04,0.04) (1.869,6.598)rectangle++(0.04,0.04) (1.85,7.129)rectangle++(0.04,0.04)
  (1.855,7.582)rectangle++(0.04,0.04) (1.878,7.967)rectangle++(0.04,0.04) (1.913,8.293)rectangle++(0.04,0.04) (1.956,8.569)rectangle++(0.04,0.04) (2.003,8.801)rectangle++(0.04,0.04)
  (2.052,8.997)rectangle++(0.04,0.04) (2.102,9.162)rectangle++(0.04,0.04) (2.151,9.3)rectangle++(0.04,0.04) (2.198,9.415)rectangle++(0.04,0.04) (2.243,9.511)rectangle++(0.04,0.04)
  (2.285,9.591)rectangle++(0.04,0.04) (2.323,9.657)rectangle++(0.04,0.04) (2.359,9.712)rectangle++(0.04,0.04) (2.391,9.757)rectangle++(0.04,0.04) (2.421,9.794)rectangle++(0.04,0.04)
  (2.447,9.824)rectangle++(0.04,0.04) (2.471,9.849)rectangle++(0.04,0.04) (2.492,9.869)rectangle++(0.04,0.04) (2.511,9.885)rectangle++(0.04,0.04) (2.527,9.898)rectangle++(0.04,0.04)
  (2.542,9.908)rectangle++(0.04,0.04) (2.555,9.917)rectangle++(0.04,0.04) (2.566,9.923)rectangle++(0.04,0.04) (-1.603,5.388)rectangle++(0.04,0.04) (-1.149,6.24)rectangle++(0.04,0.04)
  (-0.729,6.946)rectangle++(0.04,0.04) (-0.344,7.529)rectangle++(0.04,0.04) (0.006,8.01)rectangle++(0.04,0.04) (0.031,2.672)rectangle++(0.04,0.04) (0.131,3.866)rectangle++(0.04,0.04)
  (-0.754,3.736)rectangle++(0.04,0.04) (-0.493,4.802)rectangle++(0.04,0.04) (-0.229,5.697)rectangle++(0.04,0.04) (0.031,6.448)rectangle++(0.04,0.04) (0.282,7.076)rectangle++(0.04,0.04)
  (0.52,7.599)rectangle++(0.04,0.04) (0.744,8.035)rectangle++(0.04,0.04) (0.952,8.396)rectangle++(0.04,0.04) (1.143,8.695)rectangle++(0.04,0.04) (1.317,8.941)rectangle++(0.04,0.04)
  (1.475,9.143)rectangle++(0.04,0.04) (1.617,9.309)rectangle++(0.04,0.04) (0.367,2.242)rectangle++(0.04,0.04) (0.4,3.487)rectangle++(0.04,0.04) (0.477,4.544)rectangle++(0.04,0.04)
  (0.585,5.44)rectangle++(0.04,0.04) (0.713,6.197)rectangle++(0.04,0.04) (0.851,6.835)rectangle++(0.04,0.04) (0.995,7.372)rectangle++(0.04,0.04) (1.138,7.822)rectangle++(0.04,0.04)
  (1.278,8.2)rectangle++(0.04,0.04) (1.412,8.515)rectangle++(0.04,0.04) (1.539,8.777)rectangle++(0.04,0.04) (1.657,8.995)rectangle++(0.04,0.04) (1.766,9.176)rectangle++(0.04,0.04)
  (1.866,9.325)rectangle++(0.04,0.04) (1.957,9.448)rectangle++(0.04,0.04) (2.039,9.549)rectangle++(0.04,0.04) (2.113,9.631)rectangle++(0.04,0.04) (2.179,9.698)rectangle++(0.04,0.04)
  (2.238,9.752)rectangle++(0.04,0.04) (2.29,9.796)rectangle++(0.04,0.04) (2.336,9.831)rectangle++(0.04,0.04) (-1.624,5.326)rectangle++(0.04,0.04) (-1.17,6.189)rectangle++(0.04,0.04)
  (-0.749,6.903)rectangle++(0.04,0.04) (-0.363,7.494)rectangle++(0.04,0.04) (-0.011,7.981)rectangle++(0.04,0.04) (0.308,8.38)rectangle++(0.04,0.04) (0.595,8.707)rectangle++(0.04,0.04)
  (2.086,3.776)rectangle++(0.04,0.04) (1.922,4.722)rectangle++(0.04,0.04) (1.821,5.533)rectangle++(0.04,0.04) (-0.805,4.381)rectangle++(0.04,0.04) (-0.511,5.353)rectangle++(0.04,0.04)
  (-0.222,6.166)rectangle++(0.04,0.04) (0.055,6.846)rectangle++(0.04,0.04) (0.319,7.413)rectangle++(0.04,0.04) (1.798,3.464)rectangle++(0.04,0.04) (1.665,4.469)rectangle++(0.04,0.04)
  (1.592,5.328)rectangle++(0.04,0.04) (1.564,6.061)rectangle++(0.04,0.04) (1.57,6.686)rectangle++(0.04,0.04) (1.599,7.216)rectangle++(0.04,0.04) (-1.201,4.672)rectangle++(0.04,0.04)
  (-0.836,5.616)rectangle++(0.04,0.04) (-0.489,6.403)rectangle++(0.04,0.04) (1.643,3.086)rectangle++(0.04,0.04) (1.518,4.154)rectangle++(0.04,0.04) (-0.566,4.052)rectangle++(0.04,0.04)
  (-0.321,5.063)rectangle++(0.04,0.04) (-0.073,5.913)rectangle++(0.04,0.04) (0.172,6.625)rectangle++(0.04,0.04) (0.409,7.221)rectangle++(0.04,0.04) (0.635,7.717)rectangle++(0.04,0.04)
  (0.846,8.13)rectangle++(0.04,0.04) (1.042,8.473)rectangle++(0.04,0.04) (1.222,8.757)rectangle++(0.04,0.04) (1.387,8.991)rectangle++(0.04,0.04) (1.537,9.183)rectangle++(0.04,0.04)
  (1.671,9.34)rectangle++(0.04,0.04) (2.117,4.143)rectangle++(0.04,0.04) (0.778,3.242)rectangle++(0.04,0.04) (0.789,4.32)rectangle++(0.04,0.04) (0.841,5.237)rectangle++(0.04,0.04)
  (0.922,6.014)rectangle++(0.04,0.04) (1.022,6.671)rectangle++(0.04,0.04) (1.134,7.226)rectangle++(0.04,0.04) (1.25,7.693)rectangle++(0.04,0.04) (1.759,3.729)rectangle++(0.04,0.04)
  (-0.433,4.014)rectangle++(0.04,0.04) (-0.21,5.025)rectangle++(0.04,0.04) (0.021,5.876)rectangle++(0.04,0.04) (0.25,6.59)rectangle++(0.04,0.04) (0.474,7.188)rectangle++(0.04,0.04)
  (0.688,7.687)rectangle++(0.04,0.04) (0.89,8.102)rectangle++(0.04,0.04) (1.079,8.448)rectangle++(0.04,0.04) (1.253,8.733)rectangle++(0.04,0.04) (1.412,8.97)rectangle++(0.04,0.04)
  (1.557,9.164)rectangle++(0.04,0.04) (-1.173,3.338)rectangle++(0.04,0.04) (-0.865,4.48)rectangle++(0.04,0.04) (-0.559,5.439)rectangle++(0.04,0.04) (0.839,2.124)rectangle++(0.04,0.04)
  (0.796,3.368)rectangle++(0.04,0.04) (0.809,4.427)rectangle++(0.04,0.04) (0.863,5.327)rectangle++(0.04,0.04) (0.944,6.089)rectangle++(0.04,0.04) (1.044,6.734)rectangle++(0.04,0.04)
  (1.154,7.279)rectangle++(0.04,0.04) (1.27,7.737)rectangle++(0.04,0.04) (1.387,8.122)rectangle++(0.04,0.04) (1.849,3.843)rectangle++(0.04,0.04) (1.723,4.789)rectangle++(0.04,0.04)
  (1.654,5.598)rectangle++(0.04,0.04) (1.627,6.288)rectangle++(0.04,0.04) (1.633,6.876)rectangle++(0.04,0.04) (1.661,7.376)rectangle++(0.04,0.04) (1.704,7.799)rectangle++(0.04,0.04)
  (1.758,8.157)rectangle++(0.04,0.04) (1.819,8.459)rectangle++(0.04,0.04) (1.175,2.287)rectangle++(0.04,0.04) (-0.231,3.619)rectangle++(0.04,0.04) (-0.054,4.682)rectangle++(0.04,0.04)
  (1.166,2.853)rectangle++(0.04,0.04) (1.103,3.974)rectangle++(0.04,0.04) (1.094,4.93)rectangle++(0.04,0.04) (1.125,5.743)rectangle++(0.04,0.04) (1.184,6.433)rectangle++(0.04,0.04)
  (1.261,7.017)rectangle++(0.04,0.04) (1.351,7.51)rectangle++(0.04,0.04) (1.446,7.926)rectangle++(0.04,0.04) (1.544,8.275)rectangle++(0.04,0.04) (1.641,8.568)rectangle++(0.04,0.04)
  (1.736,8.814)rectangle++(0.04,0.04) (1.826,9.018)rectangle++(0.04,0.04) (2.019,4.116)rectangle++(0.04,0.04) (1.879,5.014)rectangle++(0.04,0.04) (1.795,5.783)rectangle++(0.04,0.04)
  (1.755,6.44)rectangle++(0.04,0.04) (1.747,7)rectangle++(0.04,0.04) (1.763,7.476)rectangle++(0.04,0.04) (-1.227,4.754)rectangle++(0.04,0.04) (1.321,2.6)rectangle++(0.04,0.04)
  (1.225,3.754)rectangle++(0.04,0.04) (1.189,4.738)rectangle++(0.04,0.04) (1.198,5.576)rectangle++(0.04,0.04) (1.239,6.288)rectangle++(0.04,0.04) (1.302,6.891)rectangle++(0.04,0.04)
  (1.381,7.402)rectangle++(0.04,0.04) (1.467,7.832)rectangle++(0.04,0.04) (1.558,8.195)rectangle++(0.04,0.04) (1.65,8.5)rectangle++(0.04,0.04) (1.74,8.755)rectangle++(0.04,0.04)
  (1.827,8.968)rectangle++(0.04,0.04) (1.91,9.146)rectangle++(0.04,0.04) (1.987,9.294)rectangle++(0.04,0.04) (2.058,9.417)rectangle++(0.04,0.04) (2.124,9.518)rectangle++(0.04,0.04)
  (2.184,9.602)rectangle++(0.04,0.04) (2.238,9.67)rectangle++(0.04,0.04) (2.287,9.726)rectangle++(0.04,0.04) (2.331,9.772)rectangle++(0.04,0.04) (0.386,2.256)rectangle++(0.04,0.04)
  (0.514,2.485)rectangle++(0.04,0.04) (0.534,3.688)rectangle++(0.04,0.04) (0.599,4.709)rectangle++(0.04,0.04) (0.695,5.575)rectangle++(0.04,0.04) (-1.015,4.179)rectangle++(0.04,0.04)
  (-0.698,5.189)rectangle++(0.04,0.04) (-0.388,6.034)rectangle++(0.04,0.04) (-0.091,6.741)rectangle++(0.04,0.04) (0.023,2.499)rectangle++(0.04,0.04) (0.118,3.72)rectangle++(0.04,0.04)
  (0.247,4.753)rectangle++(0.04,0.04) (0.397,5.627)rectangle++(0.04,0.04) (0.561,6.363)rectangle++(0.04,0.04) (0.729,6.982)rectangle++(0.04,0.04) (0.897,7.502)rectangle++(0.04,0.04)
  (1.06,7.937)rectangle++(0.04,0.04) (1.216,8.3)rectangle++(0.04,0.04) (1.364,8.603)rectangle++(0.04,0.04) (1.501,8.854)rectangle++(0.04,0.04) (1.628,9.062)rectangle++(0.04,0.04)
  (1.744,9.234)rectangle++(0.04,0.04) (-1.182,3.353)rectangle++(0.04,0.04) (-0.872,4.493)rectangle++(0.04,0.04) (-0.564,5.45)rectangle++(0.04,0.04) (-0.264,6.252)rectangle++(0.04,0.04)
  (0.024,6.921)rectangle++(0.04,0.04) (0.295,7.478)rectangle++(0.04,0.04) (0.548,7.941)rectangle++(0.04,0.04) (0.365,2.922)rectangle++(0.04,0.04) (0.425,4.066)rectangle++(0.04,0.04)
  (0.522,5.035)rectangle++(0.04,0.04) (0.643,5.855)rectangle++(0.04,0.04) (0.778,6.547)rectangle++(0.04,0.04) (0.921,7.13)rectangle++(0.04,0.04) (1.066,7.62)rectangle++(0.04,0.04)
  (1.209,8.031)rectangle++(0.04,0.04) (1.346,8.374)rectangle++(0.04,0.04) (1.477,8.66)rectangle++(0.04,0.04) (-1.53,4.934)rectangle++(0.04,0.04) (-1.105,5.851)rectangle++(0.04,0.04)
  (-0.707,6.614)rectangle++(0.04,0.04) (-0.339,7.247)rectangle++(0.04,0.04) (0,7.769)rectangle++(0.04,0.04) (0.308,8.2)rectangle++(0.04,0.04) (0.588,8.554)rectangle++(0.04,0.04)
  (-1.666,4.746)rectangle++(0.04,0.04) (0.082,3.465)rectangle++(0.04,0.04) (0.206,4.538)rectangle++(0.04,0.04) (0.355,5.445)rectangle++(0.04,0.04) (-1.049,4.089)rectangle++(0.04,0.04)
  (-0.73,5.114)rectangle++(0.04,0.04) (-0.418,5.972)rectangle++(0.04,0.04) (-0.119,6.689)rectangle++(0.04,0.04) (0.164,7.287)rectangle++(0.04,0.04) (1.785,3.405)rectangle++(0.04,0.04)
  (1.651,4.419)rectangle++(0.04,0.04) (1.578,5.286)rectangle++(0.04,0.04) (1.551,6.026)rectangle++(0.04,0.04) (-0.958,4.428)rectangle++(0.04,0.04) (-0.639,5.398)rectangle++(0.04,0.04)
  (-0.329,6.21)rectangle++(0.04,0.04) (-0.034,6.888)rectangle++(0.04,0.04) (0.245,7.452)rectangle++(0.04,0.04) (0.504,7.921)rectangle++(0.04,0.04) (0.743,8.309)rectangle++(0.04,0.04)
  (0.962,8.629)rectangle++(0.04,0.04) (2.024,3.844)rectangle++(0.04,0.04) (1.872,4.783)rectangle++(0.04,0.04) (1.78,5.587)rectangle++(0.04,0.04) (1.735,6.274)rectangle++(0.04,0.04)
  (1.723,6.86)rectangle++(0.04,0.04) (1.737,7.358)rectangle++(0.04,0.04) (1.768,7.781)rectangle++(0.04,0.04) (1.812,8.139)rectangle++(0.04,0.04) (1.864,8.442)rectangle++(0.04,0.04)
  (1.92,8.697)rectangle++(0.04,0.04) (1.977,8.912)rectangle++(0.04,0.04) (2.035,9.092)rectangle++(0.04,0.04) (-0.225,3.985)rectangle++(0.04,0.04) (-0.034,4.992)rectangle++(0.04,0.04)
  (-1.022,3.927)rectangle++(0.04,0.04) (-0.714,4.975)rectangle++(0.04,0.04) (-0.41,5.853)rectangle++(0.04,0.04) (-0.117,6.588)rectangle++(0.04,0.04) (0.162,7.201)rectangle++(0.04,0.04)
  (0.424,7.71)rectangle++(0.04,0.04) (0.349,2.895)rectangle++(0.04,0.04) (0.41,4.043)rectangle++(0.04,0.04) (0.508,5.016)rectangle++(0.04,0.04) (0.63,5.84)rectangle++(0.04,0.04)
  (0.767,6.535)rectangle++(0.04,0.04) (0.911,7.12)rectangle++(0.04,0.04) (1.057,7.612)rectangle++(0.04,0.04) (1.201,8.024)rectangle++(0.04,0.04) (1.34,8.369)rectangle++(0.04,0.04)
  (1.471,8.656)rectangle++(0.04,0.04) (1.594,8.895)rectangle++(0.04,0.04) (1.709,9.093)rectangle++(0.04,0.04) (1.814,9.257)rectangle++(0.04,0.04) (1.91,9.392)rectangle++(0.04,0.04)
  (2.102,4.207)rectangle++(0.04,0.04) (1.953,5.088)rectangle++(0.04,0.04) (1.861,5.843)rectangle++(0.04,0.04) (1.814,6.488)rectangle++(0.04,0.04) (-0.926,3.373)rectangle++(0.04,0.04)
  (-0.655,4.5)rectangle++(0.04,0.04) (-0.379,5.448)rectangle++(0.04,0.04) (-1.181,3.952)rectangle++(0.04,0.04) (-0.848,5.003)rectangle++(0.04,0.04) (-0.523,5.882)rectangle++(0.04,0.04)
  (-0.211,6.617)rectangle++(0.04,0.04) (0.083,7.229)rectangle++(0.04,0.04) (0.358,7.738)rectangle++(0.04,0.04) (0.611,8.159)rectangle++(0.04,0.04) (0.844,8.507)rectangle++(0.04,0.04)
  (1.055,8.793)rectangle++(0.04,0.04) (1.246,9.028)rectangle++(0.04,0.04) (1.418,9.22)rectangle++(0.04,0.04) (1.572,9.377)rectangle++(0.04,0.04) (-1.665,5.093)rectangle++(0.04,0.04)
  (-1.213,5.992)rectangle++(0.04,0.04) (-0.794,6.738)rectangle++(0.04,0.04) (-0.408,7.355)rectangle++(0.04,0.04) (-0.054,7.865)rectangle++(0.04,0.04) (0.266,8.283)rectangle++(0.04,0.04)
  (-1.667,4.633)rectangle++(0.04,0.04) (-1.234,5.601)rectangle++(0.04,0.04) (-0.827,6.406)rectangle++(0.04,0.04) (-1.464,4.057)rectangle++(0.04,0.04) (-1.084,5.103)rectangle++(0.04,0.04)
  (-0.72,5.978)rectangle++(0.04,0.04) (-0.375,6.706)rectangle++(0.04,0.04) (-0.053,7.311)rectangle++(0.04,0.04) (0.245,7.813)rectangle++(0.04,0.04) (0.519,8.227)rectangle++(0.04,0.04)
  (0.768,8.569)rectangle++(0.04,0.04) (0.993,8.849)rectangle++(0.04,0.04) (1.196,9.078)rectangle++(0.04,0.04) (1.378,9.265)rectangle++(0.04,0.04) (1.539,9.416)rectangle++(0.04,0.04)
  (1.683,9.538)rectangle++(0.04,0.04) (1.81,9.636)rectangle++(0.04,0.04) (1.922,9.715)rectangle++(0.04,0.04) (2.02,9.777)rectangle++(0.04,0.04) (2.106,9.826)rectangle++(0.04,0.04)
  (2.181,9.864)rectangle++(0.04,0.04) (2.246,9.893)rectangle++(0.04,0.04) (-1.25,3.391)rectangle++(0.04,0.04) (-0.929,4.528)rectangle++(0.04,0.04) (-0.61,5.483)rectangle++(0.04,0.04)
  (-1.229,3.915)rectangle++(0.04,0.04) (-0.891,4.974)rectangle++(0.04,0.04) (-0.56,5.859)rectangle++(0.04,0.04) (-0.244,6.599)rectangle++(0.04,0.04) (0.054,7.215)rectangle++(0.04,0.04)
  (0.333,7.727)rectangle++(0.04,0.04) (-1.023,3.244)rectangle++(0.04,0.04) (-0.742,4.394)rectangle++(0.04,0.04) (-0.457,5.361)rectangle++(0.04,0.04) (1.381,2.892)rectangle++(0.04,0.04)
  (1.287,3.999)rectangle++(0.04,0.04) (0.84,3.026)rectangle++(0.04,0.04) (0.833,4.134)rectangle++(0.04,0.04) (0.871,5.077)rectangle++(0.04,0.04) (0.941,5.877)rectangle++(0.04,0.04)
  (1.033,6.554)rectangle++(0.04,0.04) (1.138,7.126)rectangle++(0.04,0.04) (1.25,7.608)rectangle++(0.04,0.04) (1.365,8.013)rectangle++(0.04,0.04) (1.478,8.352)rectangle++(0.04,0.04)
  (1.588,8.637)rectangle++(0.04,0.04) (1.693,8.874)rectangle++(0.04,0.04) (1.792,9.071)rectangle++(0.04,0.04) (1.884,9.235)rectangle++(0.04,0.04) (1.969,9.371)rectangle++(0.04,0.04)
  (2.046,9.482)rectangle++(0.04,0.04) (2.116,9.574)rectangle++(0.04,0.04) (2.179,9.65)rectangle++(0.04,0.04) (2.236,9.711)rectangle++(0.04,0.04) (2.287,9.761)rectangle++(0.04,0.04)
  (2.332,9.802)rectangle++(0.04,0.04) (2.372,9.835)rectangle++(0.04,0.04) (2.408,9.861)rectangle++(0.04,0.04) (-1.618,5.346)rectangle++(0.04,0.04) (-1.163,6.205)rectangle++(0.04,0.04)
  (-0.743,6.917)rectangle++(0.04,0.04) (-0.357,7.505)rectangle++(0.04,0.04) (1.9,3.328)rectangle++(0.04,0.04) (1.746,4.349)rectangle++(0.04,0.04) (1.656,5.223)rectangle++(0.04,0.04)
  (1.614,5.97)rectangle++(0.04,0.04) (1.609,6.606)rectangle++(0.04,0.04) (1.63,7.147)rectangle++(0.04,0.04) (1.669,7.606)rectangle++(0.04,0.04) (1.72,7.994)rectangle++(0.04,0.04)
  (1.78,8.323)rectangle++(0.04,0.04) (-1.404,4.924)rectangle++(0.04,0.04) (-0.998,5.838)rectangle++(0.04,0.04) (-1.374,3.913)rectangle++(0.04,0.04) (-1.014,4.977)rectangle++(0.04,0.04)
  (-0.665,5.867)rectangle++(0.04,0.04) (-0.333,6.61)rectangle++(0.04,0.04) (-0.021,7.228)rectangle++(0.04,0.04) (0.269,7.741)rectangle++(0.04,0.04) (0.536,8.165)rectangle++(0.04,0.04)
  (0.78,8.515)rectangle++(0.04,0.04) (1.002,8.803)rectangle++(0.04,0.04) (1.201,9.039)rectangle++(0.04,0.04) (1.381,9.231)rectangle++(0.04,0.04) (-1.668,5.029)rectangle++(0.04,0.04)
  (-1.219,5.937)rectangle++(0.04,0.04) (-0.801,6.692)rectangle++(0.04,0.04) (-0.415,7.316)rectangle++(0.04,0.04) (-0.062,7.832)rectangle++(0.04,0.04) (0.258,8.256)rectangle++(0.04,0.04)
  (0.547,8.603)rectangle++(0.04,0.04) (0.807,8.887)rectangle++(0.04,0.04) (1.039,9.118)rectangle++(0.04,0.04) (1.246,9.305)rectangle++(0.04,0.04) (-1.708,5.018)rectangle++(0.04,0.04)
  (1.443,2.542)rectangle++(0.04,0.04) (1.326,3.699)rectangle++(0.04,0.04) (1.273,4.688)rectangle++(0.04,0.04) (1.267,5.53)rectangle++(0.04,0.04) (1.296,6.246)rectangle++(0.04,0.04)
  (-1.05,4.424)rectangle++(0.04,0.04) (-0.718,5.398)rectangle++(0.04,0.04) (-0.396,6.213)rectangle++(0.04,0.04) (-0.023,3.641)rectangle++(0.04,0.04) (0.124,4.692)rectangle++(0.04,0.04)
  (0.291,5.579)rectangle++(0.04,0.04) (-1.088,4.104)rectangle++(0.04,0.04) (-1.022,3.554)rectangle++(0.04,0.04) (1,2.415)rectangle++(0.04,0.04) (0.944,3.609)rectangle++(0.04,0.04)
  (0.945,4.626)rectangle++(0.04,0.04) (0.986,5.491)rectangle++(0.04,0.04) (1.055,6.224)rectangle++(0.04,0.04) (1.144,6.844)rectangle++(0.04,0.04) (1.244,7.368)rectangle++(0.04,0.04)
  (1.35,7.809)rectangle++(0.04,0.04) (1.458,8.18)rectangle++(0.04,0.04) (1.564,8.491)rectangle++(0.04,0.04) (-1.478,4.917)rectangle++(0.04,0.04) (-1.062,5.835)rectangle++(0.04,0.04)
  (-0.671,6.598)rectangle++(0.04,0.04) (-0.309,7.232)rectangle++(0.04,0.04) (0.024,7.755)rectangle++(0.04,0.04) (0.329,8.187)rectangle++(0.04,0.04) (1.989,3.616)rectangle++(0.04,0.04)
  (-0.367,4.034)rectangle++(0.04,0.04) (-0.153,5.04)rectangle++(0.04,0.04) (0.069,5.886)rectangle++(0.04,0.04) (0.292,6.597)rectangle++(0.04,0.04) (0.51,7.192)rectangle++(0.04,0.04)
  (0.719,7.689)rectangle++(0.04,0.04) (0.916,8.103)rectangle++(0.04,0.04) (1.101,8.447)rectangle++(0.04,0.04) (1.271,8.732)rectangle++(0.04,0.04) (1.428,8.968)rectangle++(0.04,0.04)
  (1.57,9.162)rectangle++(0.04,0.04) (2.085,4.085)rectangle++(0.04,0.04) (1.933,4.985)rectangle++(0.04,0.04) (1.008,3.362)rectangle++(0.04,0.04) (0.989,4.414)rectangle++(0.04,0.04)
  (1.015,5.308)rectangle++(0.04,0.04) (1.073,6.068)rectangle++(0.04,0.04) (1.153,6.711)rectangle++(0.04,0.04) (1.246,7.254)rectangle++(0.04,0.04) (1.347,7.713)rectangle++(0.04,0.04)
  (1.451,8.098)rectangle++(0.04,0.04) (1.555,8.421)rectangle++(0.04,0.04) (0.357,2.389)rectangle++(0.04,0.04) (0.397,3.612)rectangle++(0.04,0.04) (0.48,4.651)rectangle++(0.04,0.04)
  (0.592,5.53)rectangle++(0.04,0.04) (0.722,6.273)rectangle++(0.04,0.04) (0.862,6.9)rectangle++(0.04,0.04) (-1.266,4.471)rectangle++(0.04,0.04) (-0.9,5.448)rectangle++(0.04,0.04)
  (0.08,3.25)rectangle++(0.04,0.04) (0.196,4.355)rectangle++(0.04,0.04) (0.338,5.29)rectangle++(0.04,0.04) (1.27,3.059)rectangle++(0.04,0.04) (1.2,4.146)rectangle++(0.04,0.04)
  (1.184,5.072)rectangle++(0.04,0.04) (1.207,5.86)rectangle++(0.04,0.04) (1.258,6.529)rectangle++(0.04,0.04) (1.328,7.096)rectangle++(0.04,0.04) (1.411,7.574)rectangle++(0.04,0.04)
  (1.5,7.978)rectangle++(0.04,0.04) (1.592,8.317)rectangle++(0.04,0.04) (1.683,8.602)rectangle++(0.04,0.04) (1.773,8.841)rectangle++(0.04,0.04) (1.858,9.04)rectangle++(0.04,0.04)
  (1.939,9.206)rectangle++(0.04,0.04) (2.014,9.344)rectangle++(0.04,0.04) (2.084,9.458)rectangle++(0.04,0.04) (2.147,9.552)rectangle++(0.04,0.04) (2.205,9.629)rectangle++(0.04,0.04)
  (2.257,9.693)rectangle++(0.04,0.04) (2.304,9.745)rectangle++(0.04,0.04) (2.346,9.787)rectangle++(0.04,0.04) (2.383,9.822)rectangle++(0.04,0.04) (2.417,9.849)rectangle++(0.04,0.04)
  (0.569,3.162)rectangle++(0.04,0.04) (0.608,4.261)rectangle++(0.04,0.04) (0.685,5.194)rectangle++(0.04,0.04) (1.206,3.119)rectangle++(0.04,0.04) (1.148,4.199)rectangle++(0.04,0.04)
  (1.141,5.12)rectangle++(0.04,0.04) (1.173,5.902)rectangle++(0.04,0.04) (-0.999,4.333)rectangle++(0.04,0.04) (-0.678,5.319)rectangle++(0.04,0.04) (-0.366,6.144)rectangle++(0.04,0.04)
  (-0.067,6.834)rectangle++(0.04,0.04) (0.024,2.512)rectangle++(0.04,0.04) (0.119,3.731)rectangle++(0.04,0.04) (0.248,4.762)rectangle++(0.04,0.04) (0.399,5.634)rectangle++(0.04,0.04)
  (0.562,6.369)rectangle++(0.04,0.04) (0.73,6.988)rectangle++(0.04,0.04) (0.898,7.507)rectangle++(0.04,0.04) (1.062,7.941)rectangle++(0.04,0.04) (1.218,8.303)rectangle++(0.04,0.04)
  (1.365,8.605)rectangle++(0.04,0.04) (-1.538,4.902)rectangle++(0.04,0.04) (-1.114,5.825)rectangle++(0.04,0.04) (-0.716,6.592)rectangle++(0.04,0.04) (-0.347,7.228)rectangle++(0.04,0.04)
  (-0.008,7.754)rectangle++(0.04,0.04) (0.301,8.187)rectangle++(0.04,0.04) (0.581,8.543)rectangle++(0.04,0.04) (0.834,8.835)rectangle++(0.04,0.04) (1.06,9.072)rectangle++(0.04,0.04)
  (1.262,9.265)rectangle++(0.04,0.04) (1.441,9.421)rectangle++(0.04,0.04) (1.599,9.547)rectangle++(0.04,0.04) (1.739,9.647)rectangle++(0.04,0.04) (1.862,9.726)rectangle++(0.04,0.04)
  (1.969,9.789)rectangle++(0.04,0.04) (0.531,3.148)rectangle++(0.04,0.04) (0.718,2.69)rectangle++(0.04,0.04) (0.715,3.854)rectangle++(0.04,0.04) (0.76,4.843)rectangle++(0.04,0.04)
  (0.838,5.683)rectangle++(0.04,0.04) (-0.839,3.297)rectangle++(0.04,0.04) (-0.583,4.432)rectangle++(0.04,0.04) (-0.321,5.387)rectangle++(0.04,0.04) (1.371,2.928)rectangle++(0.04,0.04)
  (1.281,4.03)rectangle++(0.04,0.04) (1.248,4.971)rectangle++(0.04,0.04) (1.257,5.771)rectangle++(0.04,0.04) (-0.956,4.323)rectangle++(0.04,0.04) (-0.642,5.309)rectangle++(0.04,0.04)
  (-0.336,6.134)rectangle++(0.04,0.04) (-1.318,4.096)rectangle++(0.04,0.04) (-0.959,5.13)rectangle++(0.04,0.04) (-0.612,5.995)rectangle++(0.04,0.04) (-0.283,6.717)rectangle++(0.04,0.04)
  (0.026,7.317)rectangle++(0.04,0.04) (-1.505,4.397)rectangle++(0.04,0.04) (-1.105,5.394)rectangle++(0.04,0.04) (-0.726,6.225)rectangle++(0.04,0.04) (-0.37,6.917)rectangle++(0.04,0.04)
  (-0.04,7.49)rectangle++(0.04,0.04) (0.263,7.964)rectangle++(0.04,0.04) (0.54,8.355)rectangle++(0.04,0.04) (0.791,8.677)rectangle++(0.04,0.04) (1.017,8.94)rectangle++(0.04,0.04)
  (1.22,9.154)rectangle++(0.04,0.04) (1.401,9.329)rectangle++(0.04,0.04) (-1.685,5.052)rectangle++(0.04,0.04) (-1.232,5.958)rectangle++(0.04,0.04) (-0.811,6.71)rectangle++(0.04,0.04)
  (-0.424,7.332)rectangle++(0.04,0.04) (-0.069,7.845)rectangle++(0.04,0.04) (0.253,8.267)rectangle++(0.04,0.04) (0.544,8.613)rectangle++(0.04,0.04) (0.804,8.896)rectangle++(0.04,0.04)
  (0.037,2.793)rectangle++(0.04,0.04) (0.141,3.968)rectangle++(0.04,0.04) (0.276,4.964)rectangle++(0.04,0.04) (0.431,5.804)rectangle++(0.04,0.04) (0.597,6.513)rectangle++(0.04,0.04)
  (0.765,7.108)rectangle++(0.04,0.04) (0.933,7.607)rectangle++(0.04,0.04) (1.775,3.641)rectangle++(0.04,0.04) (1.652,4.62)rectangle++(0.04,0.04) (1.587,5.457)rectangle++(0.04,0.04)
  (1.565,6.171)rectangle++(0.04,0.04) (1.575,6.779)rectangle++(0.04,0.04) (1.607,7.295)rectangle++(0.04,0.04) (1.656,7.733)rectangle++(0.04,0.04) (-1.301,4.785)rectangle++(0.04,0.04)
  (-0.917,5.715)rectangle++(0.04,0.04) (1.523,2.856)rectangle++(0.04,0.04) (1.406,3.963)rectangle++(0.04,0.04) (1.352,4.909)rectangle++(0.04,0.04) (1.343,5.714)rectangle++(0.04,0.04)
  (1.368,6.4)rectangle++(0.04,0.04) (1.417,6.981)rectangle++(0.04,0.04) (1.481,7.474)rectangle++(0.04,0.04) (0.028,2.599)rectangle++(0.04,0.04) (-0.503,3.466)rectangle++(0.04,0.04)
  (-0.782,3.538)rectangle++(0.04,0.04) (-0.526,4.634)rectangle++(0.04,0.04) (-1.034,3.764)rectangle++(0.04,0.04) (-0.731,4.837)rectangle++(0.04,0.04) (-0.43,5.737)rectangle++(0.04,0.04)
  (-0.138,6.49)rectangle++(0.04,0.04) (-0.045,3.674)rectangle++(0.04,0.04) (0.107,4.721)rectangle++(0.04,0.04) (0.277,5.605)rectangle++(0.04,0.04) (0.458,6.349)rectangle++(0.04,0.04)
  (0.641,6.975)rectangle++(0.04,0.04) (0.822,7.499)rectangle++(0.04,0.04) (0.996,7.938)rectangle++(0.04,0.04) (1.162,8.303)rectangle++(0.04,0.04) (1.318,8.608)rectangle++(0.04,0.04)
  (1.977,3.921)rectangle++(0.04,0.04) (1.835,4.85)rectangle++(0.04,0.04) (1.752,5.645)rectangle++(0.04,0.04) (1.713,6.325)rectangle++(0.04,0.04) (1.706,6.904)rectangle++(0.04,0.04)
  (1.724,7.396)rectangle++(0.04,0.04) (1.759,7.814)rectangle++(0.04,0.04) (0.03,2.645)rectangle++(0.04,0.04) (0.129,3.844)rectangle++(0.04,0.04) (0.261,4.858)rectangle++(0.04,0.04)
  (0.414,5.715)rectangle++(0.04,0.04) (0.579,6.437)rectangle++(0.04,0.04) (0.747,7.045)rectangle++(0.04,0.04) (0.915,7.554)rectangle++(0.04,0.04) (1.077,7.981)rectangle++(0.04,0.04)
  (1.233,8.337)rectangle++(0.04,0.04) (1.379,8.633)rectangle++(0.04,0.04) (1.515,8.879)rectangle++(0.04,0.04) (1.641,9.083)rectangle++(0.04,0.04) (-1.59,5.048)rectangle++(0.04,0.04)
  (-1.152,5.951)rectangle++(0.04,0.04) (-0.743,6.701)rectangle++(0.04,0.04) (-0.366,7.321)rectangle++(0.04,0.04) (-0.02,7.834)rectangle++(0.04,0.04) (0.294,8.256)rectangle++(0.04,0.04)
  (0.578,8.602)rectangle++(0.04,0.04) (0.833,8.885)rectangle++(0.04,0.04) (1.061,9.115)rectangle++(0.04,0.04) (1.264,9.301)rectangle++(0.04,0.04) (1.445,9.452)rectangle++(0.04,0.04)
  (1.604,9.572)rectangle++(0.04,0.04) (1.744,9.669)rectangle++(0.04,0.04) (1.867,9.745)rectangle++(0.04,0.04) (1.974,9.805)rectangle++(0.04,0.04) (2.197,4.301)rectangle++(0.04,0.04)
  (2.037,5.164)rectangle++(0.04,0.04) (1.936,5.904)rectangle++(0.04,0.04) (1.879,6.537)rectangle++(0.04,0.04) (1.857,7.078)rectangle++(0.04,0.04) (1.859,7.538)rectangle++(0.04,0.04)
  (1.88,7.929)rectangle++(0.04,0.04) (1.913,8.261)rectangle++(0.04,0.04) (1.954,8.541)rectangle++(0.04,0.04) (2,8.778)rectangle++(0.04,0.04) (2.049,8.978)rectangle++(0.04,0.04)
  (2.099,9.145)rectangle++(0.04,0.04) (2.148,9.286)rectangle++(0.04,0.04) (2.195,9.403)rectangle++(0.04,0.04) (2.239,9.501)rectangle++(0.04,0.04) (2.281,9.583)rectangle++(0.04,0.04)
  (2.32,9.65)rectangle++(0.04,0.04) (2.356,9.706)rectangle++(0.04,0.04) (2.389,9.752)rectangle++(0.04,0.04) (2.418,9.79)rectangle++(0.04,0.04) (2.141,4.399)rectangle++(0.04,0.04)
  (1.994,5.25)rectangle++(0.04,0.04) (1.902,5.979)rectangle++(0.04,0.04) (1.854,6.602)rectangle++(0.04,0.04) (1.838,7.134)rectangle++(0.04,0.04) (1.845,7.586)rectangle++(0.04,0.04)
  (1.87,7.971)rectangle++(0.04,0.04) (1.755,3.924)rectangle++(0.04,0.04) (-0.475,4.051)rectangle++(0.04,0.04) (-0.244,5.059)rectangle++(0.04,0.04) (-0.007,5.906)rectangle++(0.04,0.04)
  (0.228,6.616)rectangle++(0.04,0.04) (0.456,7.211)rectangle++(0.04,0.04) (0.674,7.707)rectangle++(0.04,0.04) (0.879,8.12)rectangle++(0.04,0.04) (-1.522,4.716)rectangle++(0.04,0.04)
  (-1.108,5.665)rectangle++(0.04,0.04) (-0.717,6.456)rectangle++(0.04,0.04) (-0.353,7.112)rectangle++(0.04,0.04) (-0.018,7.656)rectangle++(0.04,0.04) (0.289,8.104)rectangle++(0.04,0.04)
  (0.567,8.473)rectangle++(0.04,0.04) (0.819,8.776)rectangle++(0.04,0.04) (1.045,9.023)rectangle++(0.04,0.04) (1.247,9.224)rectangle++(0.04,0.04) (2.138,4.024)rectangle++(0.04,0.04)
  (1.976,4.931)rectangle++(0.04,0.04) (0.99,3.362)rectangle++(0.04,0.04) (0.974,4.414)rectangle++(0.04,0.04) (1.002,5.309)rectangle++(0.04,0.04) (1.062,6.069)rectangle++(0.04,0.04)
  (1.143,6.712)rectangle++(0.04,0.04) (1.238,7.256)rectangle++(0.04,0.04) (1.34,7.714)rectangle++(0.04,0.04) (-1.353,4.722)rectangle++(0.04,0.04) (-0.964,5.664)rectangle++(0.04,0.04)
  (-0.595,6.449)rectangle++(0.04,0.04) (1.667,3.071)rectangle++(0.04,0.04) (1.537,4.14)rectangle++(0.04,0.04) (1.47,5.053)rectangle++(0.04,0.04) (1.449,5.833)rectangle++(0.04,0.04)
  (1.275,3.416)rectangle++(0.04,0.04) (1.218,4.449)rectangle++(0.04,0.04) (1.211,5.329)rectangle++(0.04,0.04) (1.24,6.077)rectangle++(0.04,0.04) (1.295,6.712)rectangle++(0.04,0.04)
  (1.367,7.25)rectangle++(0.04,0.04) (1.45,7.704)rectangle++(0.04,0.04) (1.538,8.087)rectangle++(0.04,0.04) (1.629,8.408)rectangle++(0.04,0.04) (-1.449,4.913)rectangle++(0.04,0.04)
  (-1.037,5.83)rectangle++(0.04,0.04) (0.017,2.376)rectangle++(0.04,0.04) (0.108,3.615)rectangle++(0.04,0.04) (0.234,4.664)rectangle++(0.04,0.04) (0.383,5.552)rectangle++(0.04,0.04)
  (1.33,3.12)rectangle++(0.04,0.04) (1.253,4.195)rectangle++(0.04,0.04) (1.231,5.112)rectangle++(0.04,0.04) (1.248,5.892)rectangle++(0.04,0.04) (1.295,6.554)rectangle++(0.04,0.04)
  (1.36,7.116)rectangle++(0.04,0.04) (1.439,7.59)rectangle++(0.04,0.04) (1.524,7.99)rectangle++(0.04,0.04) (1.613,8.327)rectangle++(0.04,0.04) (1.702,8.61)rectangle++(0.04,0.04)
  (1.789,8.846)rectangle++(0.04,0.04) (1.872,9.044)rectangle++(0.04,0.04) (1.951,9.209)rectangle++(0.04,0.04) (2.024,9.346)rectangle++(0.04,0.04) (2.092,9.459)rectangle++(0.04,0.04)
  (2.155,9.553)rectangle++(0.04,0.04) (2.211,9.63)rectangle++(0.04,0.04) (2.263,9.693)rectangle++(0.04,0.04) (2.309,9.745)rectangle++(0.04,0.04) (2.35,9.787)rectangle++(0.04,0.04)
  (2.387,9.821)rectangle++(0.04,0.04) (2.152,4.398)rectangle++(0.04,0.04) (0.837,3.299)rectangle++(0.04,0.04) (0.841,4.367)rectangle++(0.04,0.04) (-0.733,3.967)rectangle++(0.04,0.04)
  (-0.466,4.998)rectangle++(0.04,0.04) (-0.199,5.863)rectangle++(0.04,0.04) (1.475,3.048)rectangle++(0.04,0.04) (1.374,4.128)rectangle++(0.04,0.04) (1.33,5.05)rectangle++(0.04,0.04)
  (1.331,5.835)rectangle++(0.04,0.04) (-0.859,3.337)rectangle++(0.04,0.04) (-0.599,4.467)rectangle++(0.04,0.04) (-0.333,5.417)rectangle++(0.04,0.04) (-0.068,6.214)rectangle++(0.04,0.04)
  (0.188,6.881)rectangle++(0.04,0.04) (0.433,7.437)rectangle++(0.04,0.04) (1.791,3.495)rectangle++(0.04,0.04) (1.66,4.495)rectangle++(0.04,0.04) (1.589,5.351)rectangle++(0.04,0.04)
  (1.562,6.081)rectangle++(0.04,0.04) (1.569,6.702)rectangle++(0.04,0.04) (1.599,7.23)rectangle++(0.04,0.04) (1.646,7.678)rectangle++(0.04,0.04) (1.704,8.057)rectangle++(0.04,0.04)
  (1.769,8.376)rectangle++(0.04,0.04) (1.836,8.645)rectangle++(0.04,0.04) (-1.462,4.999)rectangle++(0.04,0.04) (-1.045,5.903)rectangle++(0.04,0.04) (-0.654,6.656)rectangle++(0.04,0.04)
  (-0.292,7.28)rectangle++(0.04,0.04) (-0.004,1.926)rectangle++(0.04,0.04) (0.072,3.234)rectangle++(0.04,0.04) (-0.63,3.599)rectangle++(0.04,0.04) (-0.394,4.681)rectangle++(0.04,0.04)
  (-0.15,5.591)rectangle++(0.04,0.04) (0.094,6.354)rectangle++(0.04,0.04) (0.332,6.994)rectangle++(0.04,0.04) (0.56,7.527)rectangle++(0.04,0.04) (0.775,7.972)rectangle++(0.04,0.04)
  (0.975,8.342)rectangle++(0.04,0.04) (1.16,8.648)rectangle++(0.04,0.04) (1.33,8.9)rectangle++(0.04,0.04) (1.484,9.108)rectangle++(0.04,0.04) (1.624,9.279)rectangle++(0.04,0.04)
  (1.749,9.418)rectangle++(0.04,0.04) (1.861,9.532)rectangle++(0.04,0.04) (1.961,9.624)rectangle++(0.04,0.04) (2.05,9.698)rectangle++(0.04,0.04) (2.128,9.757)rectangle++(0.04,0.04)
  (2.197,9.805)rectangle++(0.04,0.04) (-1.644,5.295)rectangle++(0.04,0.04) (-1.187,6.163)rectangle++(0.04,0.04) (-0.765,6.882)rectangle++(0.04,0.04) (-0.377,7.477)rectangle++(0.04,0.04)
  (-0.024,7.966)rectangle++(0.04,0.04) (-1.04,3.213)rectangle++(0.04,0.04) (-0.758,4.369)rectangle++(0.04,0.04) (1.171,2.633)rectangle++(0.04,0.04) (1.098,3.787)rectangle++(0.04,0.04)
  (1.083,4.771)rectangle++(0.04,0.04) (1.109,5.609)rectangle++(0.04,0.04) (1.165,6.319)rectangle++(0.04,0.04) (1.241,6.921)rectangle++(0.04,0.04) (1.329,7.429)rectangle++(0.04,0.04)
  (1.425,7.858)rectangle++(0.04,0.04) (1.523,8.218)rectangle++(0.04,0.04) (1.621,8.521)rectangle++(0.04,0.04) (1.717,8.774)rectangle++(0.04,0.04) (1.808,8.986)rectangle++(0.04,0.04)
  (1.894,9.162)rectangle++(0.04,0.04) (1.974,9.308)rectangle++(0.04,0.04) (0,2.015)rectangle++(0.04,0.04) (0.079,3.309)rectangle++(0.04,0.04) (0.197,4.406)rectangle++(0.04,0.04)
  (0.342,5.333)rectangle++(0.04,0.04) (0.501,6.116)rectangle++(0.04,0.04) (0.669,6.774)rectangle++(0.04,0.04) (0.837,7.328)rectangle++(0.04,0.04) (1.717,3.565)rectangle++(0.04,0.04)
  (1.599,4.558)rectangle++(0.04,0.04) (1.54,5.406)rectangle++(0.04,0.04) (1.523,6.13)rectangle++(0.04,0.04) (1.537,6.746)rectangle++(0.04,0.04) (1.574,7.269)rectangle++(0.04,0.04)
  (1.066,2.285)rectangle++(0.04,0.04) (0.995,3.495)rectangle++(0.04,0.04) (0.984,4.527)rectangle++(0.04,0.04) (1.015,5.405)rectangle++(0.04,0.04) (1.077,6.15)rectangle++(0.04,0.04)
  (1.159,6.781)rectangle++(0.04,0.04) (1.254,7.313)rectangle++(0.04,0.04) (1.356,7.762)rectangle++(0.04,0.04) (1.461,8.14)rectangle++(0.04,0.04) (1.565,8.457)rectangle++(0.04,0.04)
  (1.667,8.722)rectangle++(0.04,0.04) (1.763,8.943)rectangle++(0.04,0.04) (1.854,9.127)rectangle++(0.04,0.04) (1.939,9.28)rectangle++(0.04,0.04) (0.199,2.246)rectangle++(0.04,0.04)
  (0.257,3.497)rectangle++(0.04,0.04) (0.356,4.558)rectangle++(0.04,0.04) (0.483,5.457)rectangle++(0.04,0.04) (0.626,6.215)rectangle++(0.04,0.04) (0.779,6.854)rectangle++(0.04,0.04)
  (0.934,7.391)rectangle++(0.04,0.04) (1.087,7.841)rectangle++(0.04,0.04) (1.236,8.218)rectangle++(0.04,0.04) (1.377,8.532)rectangle++(0.04,0.04) (1.509,8.793)rectangle++(0.04,0.04)
  (1.633,9.01)rectangle++(0.04,0.04) (1.746,9.189)rectangle++(0.04,0.04) (1.849,9.337)rectangle++(0.04,0.04) (1.943,9.459)rectangle++(0.04,0.04) (2.028,9.559)rectangle++(0.04,0.04)
  (2.104,9.64)rectangle++(0.04,0.04) (2.172,9.706)rectangle++(0.04,0.04) (2.232,9.759)rectangle++(0.04,0.04) (2.285,9.802)rectangle++(0.04,0.04) (2.332,9.837)rectangle++(0.04,0.04)
  (2.374,9.864)rectangle++(0.04,0.04) (2.41,9.886)rectangle++(0.04,0.04) (0.043,2.927)rectangle++(0.04,0.04) (0.152,4.083)rectangle++(0.04,0.04) (0.29,5.06)rectangle++(0.04,0.04)
  (0.447,5.886)rectangle++(0.04,0.04) (1.405,3.199)rectangle++(0.04,0.04) (1.32,4.259)rectangle++(0.04,0.04) (1.29,5.163)rectangle++(0.04,0.04) (-0.821,4.209)rectangle++(0.04,0.04)
  (-0.532,5.207)rectangle++(0.04,0.04) (-0.246,6.043)rectangle++(0.04,0.04) (0.03,6.743)rectangle++(0.04,0.04) (0.293,7.326)rectangle++(0.04,0.04) (0.54,7.812)rectangle++(0.04,0.04)
  (0.769,8.215)rectangle++(0.04,0.04) (0.98,8.548)rectangle++(0.04,0.04) (1.173,8.823)rectangle++(0.04,0.04) (2.047,3.93)rectangle++(0.04,0.04) (1.895,4.855)rectangle++(0.04,0.04)
  (1.803,5.647)rectangle++(0.04,0.04) (1.756,6.324)rectangle++(0.04,0.04) (1.743,6.901)rectangle++(0.04,0.04) (1.756,7.392)rectangle++(0.04,0.04) (1.786,7.81)rectangle++(0.04,0.04)
  (1.828,8.163)rectangle++(0.04,0.04) (1.878,8.462)rectangle++(0.04,0.04) (1.933,8.713)rectangle++(0.04,0.04) (-1.459,5.03)rectangle++(0.04,0.04) (-1.041,5.93)rectangle++(0.04,0.04)
  (-0.65,6.679)rectangle++(0.04,0.04) (-0.288,7.299)rectangle++(0.04,0.04) (0.045,7.812)rectangle++(0.04,0.04) (0.349,8.234)rectangle++(0.04,0.04) (1.998,3.629)rectangle++(0.04,0.04)
  (1.841,4.601)rectangle++(0.04,0.04) (1.747,5.433)rectangle++(0.04,0.04) (1.7,6.145)rectangle++(0.04,0.04) (1.689,6.751)rectangle++(0.04,0.04) (1.703,7.267)rectangle++(0.04,0.04)
  (1.736,7.705)rectangle++(0.04,0.04) (1.782,8.076)rectangle++(0.04,0.04) (1.835,8.39)rectangle++(0.04,0.04) (1.893,8.654)rectangle++(0.04,0.04) (1.953,8.876)rectangle++(0.04,0.04)
  (2.013,9.063)rectangle++(0.04,0.04) (2.072,9.219)rectangle++(0.04,0.04) (2.127,9.35)rectangle++(0.04,0.04) (2.18,9.458)rectangle++(0.04,0.04) (2.229,9.549)rectangle++(0.04,0.04)
  (2.275,9.623)rectangle++(0.04,0.04) (2.316,9.685)rectangle++(0.04,0.04) (2.354,9.736)rectangle++(0.04,0.04) (2.388,9.777)rectangle++(0.04,0.04) (2.419,9.812)rectangle++(0.04,0.04)
  (2.446,9.839)rectangle++(0.04,0.04) (2.471,9.862)rectangle++(0.04,0.04) (2.492,9.88)rectangle++(0.04,0.04) (2.155,4.433)rectangle++(0.04,0.04) (2.007,5.278)rectangle++(0.04,0.04)
  (1.072,3.435)rectangle++(0.04,0.04) (1.046,4.473)rectangle++(0.04,0.04) (1.066,5.357)rectangle++(0.04,0.04) (1.118,6.107)rectangle++(0.04,0.04) (1.193,6.742)rectangle++(0.04,0.04)
  (1.281,7.279)rectangle++(0.04,0.04) (1.378,7.732)rectangle++(0.04,0.04) (1.478,8.114)rectangle++(0.04,0.04) (1.579,8.434)rectangle++(0.04,0.04) (1.677,8.702)rectangle++(0.04,0.04)
  (1.772,8.926)rectangle++(0.04,0.04) (1.861,9.112)rectangle++(0.04,0.04) (1.944,9.267)rectangle++(0.04,0.04) (2.021,9.395)rectangle++(0.04,0.04) (2.09,4.233)rectangle++(0.04,0.04)
  (1.943,5.11)rectangle++(0.04,0.04) (1.038,3.389)rectangle++(0.04,0.04) (1.016,4.435)rectangle++(0.04,0.04) (1.038,5.325)rectangle++(0.04,0.04) (1.094,6.081)rectangle++(0.04,0.04)
  (1.171,6.722)rectangle++(0.04,0.04) (1.262,7.263)rectangle++(0.04,0.04) (1.361,7.719)rectangle++(0.04,0.04) (1.463,8.103)rectangle++(0.04,0.04) (0.032,2.685)rectangle++(0.04,0.04)
  (0.132,3.877)rectangle++(0.04,0.04) (0.265,4.886)rectangle++(0.04,0.04) (0.419,5.739)rectangle++(0.04,0.04) (1.372,3.164)rectangle++(0.04,0.04) (1.291,4.231)rectangle++(0.04,0.04)
  (0.55,2.21)rectangle++(0.04,0.04) (0.554,3.453)rectangle++(0.04,0.04) (0.79,2.754)rectangle++(0.04,0.04) (0.78,3.905)rectangle++(0.04,0.04) (0.817,4.885)rectangle++(0.04,0.04)
  (0.887,5.715)rectangle++(0.04,0.04) (0.981,6.419)rectangle++(0.04,0.04) (1.088,7.013)rectangle++(0.04,0.04) (1.203,7.514)rectangle++(0.04,0.04) (1.321,7.935)rectangle++(0.04,0.04)
  (1.438,8.288)rectangle++(0.04,0.04) (-1.458,4.853)rectangle++(0.04,0.04) (-1.047,5.78)rectangle++(0.04,0.04) (-0.661,6.551)rectangle++(0.04,0.04) (-0.302,7.191)rectangle++(0.04,0.04)
  (0.029,7.721)rectangle++(0.04,0.04) (0.331,8.158)rectangle++(0.04,0.04) (0.605,8.517)rectangle++(0.04,0.04) (0.853,8.811)rectangle++(0.04,0.04) (1.075,9.052)rectangle++(0.04,0.04)
  (1.274,9.247)rectangle++(0.04,0.04) (1.451,9.405)rectangle++(0.04,0.04) (1.607,9.533)rectangle++(0.04,0.04) (1.745,9.635)rectangle++(0.04,0.04) (1.866,9.716)rectangle++(0.04,0.04)
  (1.973,9.78)rectangle++(0.04,0.04) (2.066,9.83)rectangle++(0.04,0.04) (-1.672,5.276)rectangle++(0.04,0.04) (-1.213,6.147)rectangle++(0.04,0.04) (-0.787,6.87)rectangle++(0.04,0.04)
  (1.793,3.108)rectangle++(0.04,0.04) (0.2,2.231)rectangle++(0.04,0.04) (0.257,3.485)rectangle++(0.04,0.04) (0.356,4.548)rectangle++(0.04,0.04) (0.482,5.448)rectangle++(0.04,0.04)
  (0.626,6.207)rectangle++(0.04,0.04) (0.778,6.847)rectangle++(0.04,0.04) (0.933,7.385)rectangle++(0.04,0.04) (1.086,7.837)rectangle++(0.04,0.04) (1.234,8.214)rectangle++(0.04,0.04)
  (1.376,8.529)rectangle++(0.04,0.04) (1.508,8.79)rectangle++(0.04,0.04) (1.631,9.008)rectangle++(0.04,0.04) (1.745,9.188)rectangle++(0.04,0.04) (1.848,9.336)rectangle++(0.04,0.04)
  (-1.192,3.361)rectangle++(0.04,0.04) (-0.881,4.501)rectangle++(0.04,0.04) (-0.571,5.457)rectangle++(0.04,0.04) (-0.27,6.258)rectangle++(0.04,0.04) (0.019,6.926)rectangle++(0.04,0.04)
  (0.291,7.482)rectangle++(0.04,0.04) (0.544,7.945)rectangle++(0.04,0.04) (0.778,8.327)rectangle++(0.04,0.04) (0.993,8.643)rectangle++(0.04,0.04) (1.187,8.903)rectangle++(0.04,0.04)
  (1.363,9.117)rectangle++(0.04,0.04) (1.521,9.291)rectangle++(0.04,0.04) (1.662,9.432)rectangle++(0.04,0.04) (1.788,9.547)rectangle++(0.04,0.04) (1.9,9.64)rectangle++(0.04,0.04)
  (-1.662,5.207)rectangle++(0.04,0.04) (-1.207,6.088)rectangle++(0.04,0.04) (-0.784,6.82)rectangle++(0.04,0.04) (-0.396,7.424)rectangle++(0.04,0.04) (-0.042,7.923)rectangle++(0.04,0.04)
  (-1.645,4.504)rectangle++(0.04,0.04) (-1.221,5.49)rectangle++(0.04,0.04) (-0.82,6.312)rectangle++(0.04,0.04) (-0.447,6.994)rectangle++(0.04,0.04) (-0.102,7.559)rectangle++(0.04,0.04)
  (0.213,8.025)rectangle++(0.04,0.04) (1.962,3.558)rectangle++(0.04,0.04) (1.808,4.542)rectangle++(0.04,0.04) (1.716,5.385)rectangle++(0.04,0.04) (1.672,6.105)rectangle++(0.04,0.04)
  (1.663,6.718)rectangle++(0.04,0.04) (1.68,7.24)rectangle++(0.04,0.04) (-0.194,3.791)rectangle++(0.04,0.04) (-0.016,4.827)rectangle++(0.04,0.04) (0.177,5.7)rectangle++(0.04,0.04)
  (0.377,6.434)rectangle++(0.04,0.04) (0.575,7.05)rectangle++(0.04,0.04) (1.678,3.452)rectangle++(0.04,0.04) (1.562,4.463)rectangle++(0.04,0.04) (1.504,5.328)rectangle++(0.04,0.04)
  (1.49,6.064)rectangle++(0.04,0.04) (1.507,6.691)rectangle++(0.04,0.04) (1.546,7.224)rectangle++(0.04,0.04) (1.601,7.674)rectangle++(0.04,0.04) (1.666,8.055)rectangle++(0.04,0.04)
  (1.8,3.894)rectangle++(0.04,0.04) (1.683,4.834)rectangle++(0.04,0.04) (1.622,5.638)rectangle++(0.04,0.04) (1.602,6.323)rectangle++(0.04,0.04) (1.612,6.907)rectangle++(0.04,0.04)
  (1.645,7.403)rectangle++(0.04,0.04) (1.692,7.823)rectangle++(0.04,0.04) (1.749,8.178)rectangle++(0.04,0.04) (1.811,8.477)rectangle++(0.04,0.04) (1.877,8.729)rectangle++(0.04,0.04)
  (1.942,8.941)rectangle++(0.04,0.04) (1.987,4.127)rectangle++(0.04,0.04) (1.852,5.025)rectangle++(0.04,0.04) (1.028,3.351)rectangle++(0.04,0.04) (1.005,4.404)rectangle++(0.04,0.04)
  (1.028,5.299)rectangle++(0.04,0.04) (1.084,6.059)rectangle++(0.04,0.04) (1.162,6.703)rectangle++(0.04,0.04) (1.253,7.248)rectangle++(0.04,0.04) (1.353,7.707)rectangle++(0.04,0.04)
  (1.456,8.093)rectangle++(0.04,0.04) (1.833,3.853)rectangle++(0.04,0.04) (1.71,4.798)rectangle++(0.04,0.04) (1.644,5.606)rectangle++(0.04,0.04) (1.619,6.296)rectangle++(0.04,0.04)
  (1.626,6.883)rectangle++(0.04,0.04) (1.514,3.658)rectangle++(0.04,0.04) (0.729,3.011)rectangle++(0.04,0.04) (0.738,4.127)rectangle++(0.04,0.04) (0.79,5.074)rectangle++(0.04,0.04)
  (0.872,5.878)rectangle++(0.04,0.04) (-1.048,4.272)rectangle++(0.04,0.04) (-0.722,5.269)rectangle++(0.04,0.04) (-0.405,6.104)rectangle++(0.04,0.04) (-0.102,6.8)rectangle++(0.04,0.04)
  (0.183,7.381)rectangle++(0.04,0.04) (0.448,7.862)rectangle++(0.04,0.04) (0.693,8.261)rectangle++(0.04,0.04) (-1.585,4.709)rectangle++(0.04,0.04) (1.352,2.51)rectangle++(0.04,0.04)
  (1.247,3.676)rectangle++(0.04,0.04) (1.205,4.671)rectangle++(0.04,0.04) (1.209,5.518)rectangle++(0.04,0.04) (1.225,3.301)rectangle++(0.04,0.04) (1.171,4.353)rectangle++(0.04,0.04)
  (1.167,5.249)rectangle++(0.04,0.04) (1.2,6.011)rectangle++(0.04,0.04) (1.258,6.658)rectangle++(0.04,0.04) (1.334,7.205)rectangle++(0.04,0.04) (1.42,7.667)rectangle++(0.04,0.04);
\fill[piece2]
  (1.511,8.056)rectangle++(0.04,0.04) (-1.395,4.821)rectangle++(0.04,0.04) (-0.995,5.75)rectangle++(0.04,0.04) (-0.618,6.524)rectangle++(0.04,0.04) (-0.267,7.166)rectangle++(0.04,0.04)
  (0.058,7.698)rectangle++(0.04,0.04) (-1.58,4.478)rectangle++(0.04,0.04) (-1.166,5.466)rectangle++(0.04,0.04) (-0.775,6.289)rectangle++(0.04,0.04) (-0.409,6.973)rectangle++(0.04,0.04)
  (-0.071,7.539)rectangle++(0.04,0.04) (0.239,8.007)rectangle++(0.04,0.04) (1.955,3.56)rectangle++(0.04,0.04) (1.802,4.544)rectangle++(0.04,0.04) (1.711,5.387)rectangle++(0.04,0.04)
  (1.668,6.107)rectangle++(0.04,0.04) (1.66,6.72)rectangle++(0.04,0.04) (1.677,7.242)rectangle++(0.04,0.04) (-1.193,4.692)rectangle++(0.04,0.04) (-0.828,5.632)rectangle++(0.04,0.04)
  (-0.481,6.416)rectangle++(0.04,0.04) (-0.154,7.069)rectangle++(0.04,0.04) (0.149,7.611)rectangle++(0.04,0.04) (0.429,8.06)rectangle++(0.04,0.04) (0.685,8.43)rectangle++(0.04,0.04)
  (0.917,8.734)rectangle++(0.04,0.04) (1.127,8.984)rectangle++(0.04,0.04) (1.315,9.187)rectangle++(0.04,0.04) (2.121,4.032)rectangle++(0.04,0.04) (1.962,4.939)rectangle++(0.04,0.04)
  (1.863,5.716)rectangle++(0.04,0.04) (1.81,6.38)rectangle++(0.04,0.04) (1.791,6.947)rectangle++(0.04,0.04) (1.798,7.429)rectangle++(0.04,0.04) (1.824,7.839)rectangle++(0.04,0.04)
  (1.728,3.888)rectangle++(0.04,0.04) (1.622,4.832)rectangle++(0.04,0.04) (1.57,5.638)rectangle++(0.04,0.04) (1.557,6.326)rectangle++(0.04,0.04) (1.575,6.911)rectangle++(0.04,0.04)
  (1.613,7.408)rectangle++(0.04,0.04) (1.665,7.828)rectangle++(0.04,0.04) (1.726,8.183)rectangle++(0.04,0.04) (1.792,8.483)rectangle++(0.04,0.04) (1.861,8.735)rectangle++(0.04,0.04)
  (0.036,2.771)rectangle++(0.04,0.04) (0.139,3.95)rectangle++(0.04,0.04) (0.274,4.948)rectangle++(0.04,0.04) (0.429,5.791)rectangle++(0.04,0.04) (0.594,6.501)rectangle++(0.04,0.04)
  (0.763,7.099)rectangle++(0.04,0.04) (0.93,7.6)rectangle++(0.04,0.04) (1.092,8.019)rectangle++(0.04,0.04) (1.247,8.368)rectangle++(0.04,0.04) (1.392,8.659)rectangle++(0.04,0.04)
  (1.528,8.901)rectangle++(0.04,0.04) (1.652,9.101)rectangle++(0.04,0.04) (1.766,9.266)rectangle++(0.04,0.04) (1.87,9.402)rectangle++(0.04,0.04) (1.963,9.513)rectangle++(0.04,0.04)
  (0.041,2.876)rectangle++(0.04,0.04) (0.148,4.04)rectangle++(0.04,0.04) (0.285,5.024)rectangle++(0.04,0.04) (1.211,2.996)rectangle++(0.04,0.04) (1.147,4.094)rectangle++(0.04,0.04)
  (1.136,5.03)rectangle++(0.04,0.04) (1.165,5.827)rectangle++(0.04,0.04) (1.221,6.502)rectangle++(0.04,0.04) (1.296,7.074)rectangle++(0.04,0.04) (1.382,7.557)rectangle++(0.04,0.04)
  (1.475,7.965)rectangle++(0.04,0.04) (1.57,8.307)rectangle++(0.04,0.04) (1.665,8.595)rectangle++(0.04,0.04) (1.757,8.835)rectangle++(0.04,0.04) (1.844,9.036)rectangle++(0.04,0.04)
  (0.038,2.812)rectangle++(0.04,0.04) (0.143,3.985)rectangle++(0.04,0.04) (0.227,2.481)rectangle++(0.04,0.04) (-0.443,3.48)rectangle++(0.04,0.04) (-0.239,4.572)rectangle++(0.04,0.04)
  (-0.023,5.492)rectangle++(0.04,0.04) (0.198,6.265)rectangle++(0.04,0.04) (-1.251,4.222)rectangle++(0.04,0.04) (-0.896,5.235)rectangle++(0.04,0.04) (-0.555,6.081)rectangle++(0.04,0.04)
  (-1.347,4.044)rectangle++(0.04,0.04) (-0.985,5.088)rectangle++(0.04,0.04) (-0.636,5.96)rectangle++(0.04,0.04) (-0.304,6.688)rectangle++(0.04,0.04) (0.007,7.293)rectangle++(0.04,0.04)
  (0.295,7.795)rectangle++(0.04,0.04) (0.56,8.21)rectangle++(0.04,0.04) (0.803,8.552)rectangle++(0.04,0.04) (1.022,8.834)rectangle++(0.04,0.04) (1.22,9.064)rectangle++(0.04,0.04)
  (2.102,3.987)rectangle++(0.04,0.04) (1.944,4.901)rectangle++(0.04,0.04) (1.846,5.684)rectangle++(0.04,0.04) (1.794,6.354)rectangle++(0.04,0.04) (1.777,6.925)rectangle++(0.04,0.04)
  (1.785,7.412)rectangle++(0.04,0.04) (1.626,3.797)rectangle++(0.04,0.04) (1.532,4.758)rectangle++(0.04,0.04) (-0.692,4.174)rectangle++(0.04,0.04) (-0.423,5.172)rectangle++(0.04,0.04)
  (-0.155,6.009)rectangle++(0.04,0.04) (0.106,6.71)rectangle++(0.04,0.04) (0.357,7.296)rectangle++(0.04,0.04) (0.593,7.783)rectangle++(0.04,0.04) (0.813,8.188)rectangle++(0.04,0.04)
  (0.393,2.954)rectangle++(0.04,0.04) (0.45,4.091)rectangle++(0.04,0.04) (0.96,2.853)rectangle++(0.04,0.04) (0.928,3.983)rectangle++(0.04,0.04) (0.946,4.945)rectangle++(0.04,0.04)
  (1,5.761)rectangle++(0.04,0.04) (1.078,6.453)rectangle++(0.04,0.04) (1.172,7.038)rectangle++(0.04,0.04) (1.276,7.532)rectangle++(0.04,0.04) (1.384,7.947)rectangle++(0.04,0.04)
  (0.031,2.664)rectangle++(0.04,0.04) (0.131,3.859)rectangle++(0.04,0.04) (-0.752,3.734)rectangle++(0.04,0.04) (-0.492,4.8)rectangle++(0.04,0.04) (-0.229,5.696)rectangle++(0.04,0.04)
  (0.031,6.447)rectangle++(0.04,0.04) (0.282,7.075)rectangle++(0.04,0.04) (0.521,7.599)rectangle++(0.04,0.04) (0.744,8.034)rectangle++(0.04,0.04) (0.952,8.395)rectangle++(0.04,0.04)
  (1.143,8.694)rectangle++(0.04,0.04) (1.317,8.941)rectangle++(0.04,0.04) (1.475,9.143)rectangle++(0.04,0.04) (1.617,9.309)rectangle++(0.04,0.04) (1.745,9.444)rectangle++(0.04,0.04)
  (1.859,9.554)rectangle++(0.04,0.04) (1.96,9.643)rectangle++(0.04,0.04) (2.049,9.714)rectangle++(0.04,0.04) (2.128,9.771)rectangle++(0.04,0.04) (2.198,9.817)rectangle++(0.04,0.04)
  (2.258,9.852)rectangle++(0.04,0.04) (2.312,9.88)rectangle++(0.04,0.04) (2.358,9.902)rectangle++(0.04,0.04) (2.398,9.919)rectangle++(0.04,0.04) (2.433,9.931)rectangle++(0.04,0.04)
  (2.463,9.94)rectangle++(0.04,0.04) (2.489,9.947)rectangle++(0.04,0.04) (2.511,9.952)rectangle++(0.04,0.04) (2.53,9.955)rectangle++(0.04,0.04) (2.547,9.956)rectangle++(0.04,0.04)
  (2.561,9.957)rectangle++(0.04,0.04) (2.166,4.463)rectangle++(0.04,0.04) (2.017,5.304)rectangle++(0.04,0.04) (1.077,3.443)rectangle++(0.04,0.04) (1.051,4.48)rectangle++(0.04,0.04)
  (1.07,5.362)rectangle++(0.04,0.04) (1.122,6.111)rectangle++(0.04,0.04) (1.196,6.746)rectangle++(0.04,0.04) (1.284,7.282)rectangle++(0.04,0.04) (1.381,7.735)rectangle++(0.04,0.04)
  (1.481,8.115)rectangle++(0.04,0.04) (1.836,3.863)rectangle++(0.04,0.04) (1.713,4.806)rectangle++(0.04,0.04) (1.646,5.613)rectangle++(0.04,0.04) (1.622,6.302)rectangle++(0.04,0.04)
  (1.628,6.888)rectangle++(0.04,0.04) (1.657,7.386)rectangle++(0.04,0.04) (1.702,7.808)rectangle++(0.04,0.04) (1.757,8.165)rectangle++(0.04,0.04) (1.818,8.466)rectangle++(0.04,0.04)
  (1.881,8.719)rectangle++(0.04,0.04) (1.946,8.932)rectangle++(0.04,0.04) (1.984,4.126)rectangle++(0.04,0.04) (-0.477,4.134)rectangle++(0.04,0.04) (-0.242,5.129)rectangle++(0.04,0.04)
  (-0.003,5.966)rectangle++(0.04,0.04) (0.234,6.667)rectangle++(0.04,0.04) (0.463,7.254)rectangle++(0.04,0.04) (0.682,7.744)rectangle++(0.04,0.04) (0.887,8.151)rectangle++(0.04,0.04)
  (1.078,8.489)rectangle++(0.04,0.04) (1.976,3.844)rectangle++(0.04,0.04) (1.831,4.784)rectangle++(0.04,0.04) (1.746,5.59)rectangle++(0.04,0.04) (1.705,6.278)rectangle++(0.04,0.04)
  (1.698,6.864)rectangle++(0.04,0.04) (1.716,7.363)rectangle++(0.04,0.04) (1.751,7.786)rectangle++(0.04,0.04) (1.797,8.144)rectangle++(0.04,0.04) (1.851,8.447)rectangle++(0.04,0.04)
  (1.909,8.702)rectangle++(0.04,0.04) (-1.134,3.369)rectangle++(0.04,0.04) (-0.831,4.506)rectangle++(0.04,0.04) (-0.529,5.459)rectangle++(0.04,0.04) (-0.233,6.258)rectangle++(0.04,0.04)
  (-1.326,4.139)rectangle++(0.04,0.04) (-0.964,5.168)rectangle++(0.04,0.04) (-1.226,3.787)rectangle++(0.04,0.04) (-0.893,4.864)rectangle++(0.04,0.04) (-0.567,5.766)rectangle++(0.04,0.04)
  (-0.253,6.52)rectangle++(0.04,0.04) (1.644,3.162)rectangle++(0.04,0.04) (1.522,4.218)rectangle++(0.04,0.04) (-0.579,4.066)rectangle++(0.04,0.04) (-0.332,5.075)rectangle++(0.04,0.04)
  (1.296,2.866)rectangle++(0.04,0.04) (0.558,2.809)rectangle++(0.04,0.04) (0.584,3.962)rectangle++(0.04,0.04) (0.653,4.94)rectangle++(0.04,0.04) (0.751,5.769)rectangle++(0.04,0.04)
  (0.866,6.47)rectangle++(0.04,0.04) (0.993,7.061)rectangle++(0.04,0.04) (1.124,7.558)rectangle++(0.04,0.04) (1.256,7.976)rectangle++(0.04,0.04) (1.384,8.325)rectangle++(0.04,0.04)
  (1.507,8.617)rectangle++(0.04,0.04) (1.624,8.861)rectangle++(0.04,0.04) (1.733,9.063)rectangle++(0.04,0.04) (1.833,9.23)rectangle++(0.04,0.04) (-0.208,3.994)rectangle++(0.04,0.04)
  (0.006,2.126)rectangle++(0.04,0.04) (0.088,3.403)rectangle++(0.04,0.04) (0.208,4.485)rectangle++(0.04,0.04) (0.354,5.4)rectangle++(0.04,0.04) (0.515,6.172)rectangle++(0.04,0.04)
  (0.682,6.822)rectangle++(0.04,0.04) (0.851,7.368)rectangle++(0.04,0.04) (-1.368,4.562)rectangle++(0.04,0.04) (-0.982,5.528)rectangle++(0.04,0.04) (-0.616,6.335)rectangle++(0.04,0.04)
  (-0.272,7.005)rectangle++(0.04,0.04) (0.046,7.561)rectangle++(0.04,0.04) (0.34,8.022)rectangle++(0.04,0.04) (0.607,8.401)rectangle++(0.04,0.04) (0.85,8.713)rectangle++(0.04,0.04)
  (1.069,8.968)rectangle++(0.04,0.04) (1.265,9.176)rectangle++(0.04,0.04) (1.44,9.345)rectangle++(0.04,0.04) (1.596,9.482)rectangle++(0.04,0.04) (1.734,9.592)rectangle++(0.04,0.04)
  (1.855,9.68)rectangle++(0.04,0.04) (-1.678,5.206)rectangle++(0.04,0.04) (-1.362,3.661)rectangle++(0.04,0.04) (-0.978,3.415)rectangle++(0.04,0.04) (-0.697,4.538)rectangle++(0.04,0.04)
  (0.758,2.126)rectangle++(0.04,0.04) (0.727,3.373)rectangle++(0.04,0.04) (0.751,4.434)rectangle++(0.04,0.04) (0.813,5.335)rectangle++(0.04,0.04) (0.903,6.099)rectangle++(0.04,0.04)
  (1.009,6.744)rectangle++(0.04,0.04) (1.125,7.288)rectangle++(0.04,0.04) (1.246,7.746)rectangle++(0.04,0.04) (-1.377,4.71)rectangle++(0.04,0.04) (-0.984,5.655)rectangle++(0.04,0.04)
  (-0.613,6.442)rectangle++(0.04,0.04) (-0.265,7.097)rectangle++(0.04,0.04) (0.056,7.639)rectangle++(0.04,0.04) (0.351,8.087)rectangle++(0.04,0.04) (0.62,8.456)rectangle++(0.04,0.04)
  (-1.64,4.733)rectangle++(0.04,0.04) (-1.207,5.685)rectangle++(0.04,0.04) (-0.801,6.477)rectangle++(0.04,0.04) (-0.424,7.133)rectangle++(0.04,0.04) (-0.077,7.677)rectangle++(0.04,0.04)
  (0.239,8.124)rectangle++(0.04,0.04) (0.526,8.492)rectangle++(0.04,0.04) (0.785,8.794)rectangle++(0.04,0.04) (1.017,9.039)rectangle++(0.04,0.04) (1.224,9.239)rectangle++(0.04,0.04)
  (1.407,9.4)rectangle++(0.04,0.04) (1.57,9.53)rectangle++(0.04,0.04) (1.714,9.634)rectangle++(0.04,0.04) (0.042,2.893)rectangle++(0.04,0.04) (0.149,4.054)rectangle++(0.04,0.04)
  (0.287,5.036)rectangle++(0.04,0.04) (0.443,5.865)rectangle++(0.04,0.04) (0.609,6.564)rectangle++(0.04,0.04) (0.778,7.151)rectangle++(0.04,0.04) (0.945,7.644)rectangle++(0.04,0.04)
  (-1.023,3.295)rectangle++(0.04,0.04) (-0.74,4.438)rectangle++(0.04,0.04) (-0.454,5.398)rectangle++(0.04,0.04) (-0.172,6.203)rectangle++(0.04,0.04) (0.1,6.875)rectangle++(0.04,0.04)
  (0.3,2.798)rectangle++(0.04,0.04) (-0.497,3.556)rectangle++(0.04,0.04) (-0.282,4.639)rectangle++(0.04,0.04) (-0.056,5.551)rectangle++(0.04,0.04) (0.172,6.316)rectangle++(0.04,0.04)
  (0.397,6.958)rectangle++(0.04,0.04) (0.613,7.495)rectangle++(0.04,0.04) (0.819,7.942)rectangle++(0.04,0.04) (1.011,8.314)rectangle++(0.04,0.04) (1.19,8.623)rectangle++(0.04,0.04)
  (1.354,8.878)rectangle++(0.04,0.04) (1.504,9.088)rectangle++(0.04,0.04) (1.64,9.261)rectangle++(0.04,0.04) (1.762,9.403)rectangle++(0.04,0.04) (1.872,9.518)rectangle++(0.04,0.04)
  (1.969,9.612)rectangle++(0.04,0.04) (2.056,9.687)rectangle++(0.04,0.04) (2.133,9.748)rectangle++(0.04,0.04) (2.164,4.326)rectangle++(0.04,0.04) (2.01,5.187)rectangle++(0.04,0.04)
  (1.914,5.925)rectangle++(0.04,0.04) (1.862,6.556)rectangle++(0.04,0.04) (1.842,7.094)rectangle++(0.04,0.04) (1.848,7.552)rectangle++(0.04,0.04) (1.87,7.942)rectangle++(0.04,0.04)
  (1.905,8.272)rectangle++(0.04,0.04) (1.948,8.551)rectangle++(0.04,0.04) (1.996,8.787)rectangle++(0.04,0.04) (2.046,8.985)rectangle++(0.04,0.04) (2.096,9.152)rectangle++(0.04,0.04)
  (2.146,9.291)rectangle++(0.04,0.04) (2.193,9.408)rectangle++(0.04,0.04) (2.238,9.505)rectangle++(0.04,0.04) (0.041,2.875)rectangle++(0.04,0.04) (0.148,4.039)rectangle++(0.04,0.04)
  (0.285,5.023)rectangle++(0.04,0.04) (0.441,5.855)rectangle++(0.04,0.04) (0.607,6.555)rectangle++(0.04,0.04) (0.776,7.144)rectangle++(0.04,0.04) (0.943,7.637)rectangle++(0.04,0.04)
  (1.105,8.05)rectangle++(0.04,0.04) (1.259,8.395)rectangle++(0.04,0.04) (1.404,8.681)rectangle++(0.04,0.04) (1.538,8.919)rectangle++(0.04,0.04) (1.662,9.116)rectangle++(0.04,0.04)
  (1.775,9.278)rectangle++(0.04,0.04) (1.878,9.412)rectangle++(0.04,0.04) (1.97,9.521)rectangle++(0.04,0.04) (2.053,9.61)rectangle++(0.04,0.04) (2.128,9.683)rectangle++(0.04,0.04)
  (2.194,9.741)rectangle++(0.04,0.04) (2.252,9.789)rectangle++(0.04,0.04) (2.304,9.827)rectangle++(0.04,0.04) (2.349,9.857)rectangle++(0.04,0.04) (2.389,9.88)rectangle++(0.04,0.04)
  (2.167,4.41)rectangle++(0.04,0.04) (0.838,3.304)rectangle++(0.04,0.04) (-0.508,3.757)rectangle++(0.04,0.04) (-0.283,4.81)rectangle++(0.04,0.04) (-0.051,5.696)rectangle++(0.04,0.04)
  (1.417,3.049)rectangle++(0.04,0.04) (1.324,4.131)rectangle++(0.04,0.04) (1.288,5.055)rectangle++(0.04,0.04) (1.295,5.841)rectangle++(0.04,0.04) (1.295,3.383)rectangle++(0.04,0.04)
  (1.234,4.42)rectangle++(0.04,0.04) (1.223,5.303)rectangle++(0.04,0.04) (1.25,6.055)rectangle++(0.04,0.04) (1.302,6.693)rectangle++(0.04,0.04) (1.372,7.233)rectangle++(0.04,0.04)
  (1.63,3.669)rectangle++(0.04,0.04) (1.53,4.649)rectangle++(0.04,0.04) (1.485,5.487)rectangle++(0.04,0.04) (1.479,6.201)rectangle++(0.04,0.04) (1.503,6.808)rectangle++(0.04,0.04)
  (1.51,3.616)rectangle++(0.04,0.04) (1.426,4.61)rectangle++(0.04,0.04) (1.394,5.457)rectangle++(0.04,0.04) (1.401,6.179)rectangle++(0.04,0.04) (1.436,6.792)rectangle++(0.04,0.04)
  (1.514,3.598)rectangle++(0.04,0.04) (0.714,3)rectangle++(0.04,0.04) (0.247,2.382)rectangle++(0.04,0.04) (0.303,3.611)rectangle++(0.04,0.04) (0.4,4.654)rectangle++(0.04,0.04)
  (0.524,5.536)rectangle++(0.04,0.04) (0.665,6.281)rectangle++(0.04,0.04) (0.814,6.908)rectangle++(0.04,0.04) (0.966,7.436)rectangle++(0.04,0.04) (1.116,7.878)rectangle++(0.04,0.04)
  (1.82,3.735)rectangle++(0.04,0.04) (1.694,4.698)rectangle++(0.04,0.04) (1.626,5.522)rectangle++(0.04,0.04) (-0.834,3.363)rectangle++(0.04,0.04) (-0.576,4.488)rectangle++(0.04,0.04)
  (0.758,2.138)rectangle++(0.04,0.04) (0.728,3.383)rectangle++(0.04,0.04) (0.752,4.443)rectangle++(0.04,0.04) (0.814,5.342)rectangle++(0.04,0.04) (0.904,6.105)rectangle++(0.04,0.04)
  (1.01,6.749)rectangle++(0.04,0.04) (1.126,7.292)rectangle++(0.04,0.04) (1.247,7.75)rectangle++(0.04,0.04) (1.368,8.134)rectangle++(0.04,0.04) (1.486,8.455)rectangle++(0.04,0.04)
  (1.599,8.724)rectangle++(0.04,0.04) (1.706,8.947)rectangle++(0.04,0.04) (1.213,2.268)rectangle++(0.04,0.04) (1.12,3.476)rectangle++(0.04,0.04) (1.089,4.506)rectangle++(0.04,0.04)
  (1.103,5.383)rectangle++(0.04,0.04) (1.151,6.127)rectangle++(0.04,0.04) (1.221,6.758)rectangle++(0.04,0.04) (1.306,7.292)rectangle++(0.04,0.04) (-1.27,4.632)rectangle++(0.04,0.04)
  (-0.896,5.584)rectangle++(0.04,0.04) (-0.541,6.379)rectangle++(0.04,0.04) (-0.207,7.04)rectangle++(0.04,0.04) (0.104,7.588)rectangle++(0.04,0.04) (0.389,8.042)rectangle++(0.04,0.04)
  (0.651,8.416)rectangle++(0.04,0.04) (0.887,8.724)rectangle++(0.04,0.04) (1.101,8.976)rectangle++(0.04,0.04) (1.293,9.182)rectangle++(0.04,0.04) (1.464,9.349)rectangle++(0.04,0.04)
  (1.616,9.484)rectangle++(0.04,0.04) (1.751,9.593)rectangle++(0.04,0.04) (1.87,9.68)rectangle++(0.04,0.04) (1.974,9.75)rectangle++(0.04,0.04) (2.066,9.805)rectangle++(0.04,0.04)
  (2.146,9.847)rectangle++(0.04,0.04) (2.216,9.881)rectangle++(0.04,0.04) (2.276,9.906)rectangle++(0.04,0.04) (2.329,9.925)rectangle++(0.04,0.04) (2.375,9.94)rectangle++(0.04,0.04)
  (2.414,9.95)rectangle++(0.04,0.04) (2.181,4.429)rectangle++(0.04,0.04) (2.029,5.273)rectangle++(0.04,0.04) (1.933,5.997)rectangle++(0.04,0.04) (1.881,6.617)rectangle++(0.04,0.04)
  (1.861,7.145)rectangle++(0.04,0.04) (1.551,3.762)rectangle++(0.04,0.04) (1.467,4.732)rectangle++(0.04,0.04) (1.434,5.56)rectangle++(0.04,0.04) (1.439,6.265)rectangle++(0.04,0.04)
  (1.472,6.864)rectangle++(0.04,0.04) (1.523,7.371)rectangle++(0.04,0.04) (1.587,7.801)rectangle++(0.04,0.04) (1.659,8.164)rectangle++(0.04,0.04) (1.735,8.469)rectangle++(0.04,0.04)
  (1.811,8.725)rectangle++(0.04,0.04) (1.886,8.94)rectangle++(0.04,0.04) (1.959,9.12)rectangle++(0.04,0.04) (2.027,9.27)rectangle++(0.04,0.04) (2.092,9.394)rectangle++(0.04,0.04)
  (2.152,9.498)rectangle++(0.04,0.04) (2.207,9.583)rectangle++(0.04,0.04) (2.257,9.654)rectangle++(0.04,0.04) (2.302,9.712)rectangle++(0.04,0.04) (2.343,9.759)rectangle++(0.04,0.04)
  (2.38,9.798)rectangle++(0.04,0.04) (2.147,4.392)rectangle++(0.04,0.04) (1.999,5.243)rectangle++(0.04,0.04) (1.906,5.973)rectangle++(0.04,0.04) (1.857,6.597)rectangle++(0.04,0.04)
  (1.84,7.129)rectangle++(0.04,0.04) (1.847,7.582)rectangle++(0.04,0.04) (1.871,7.967)rectangle++(0.04,0.04) (1.907,8.294)rectangle++(0.04,0.04) (1.951,8.57)rectangle++(0.04,0.04)
  (1.999,8.802)rectangle++(0.04,0.04) (2.049,8.998)rectangle++(0.04,0.04) (-1.499,5.108)rectangle++(0.04,0.04) (-1.072,5.998)rectangle++(0.04,0.04) (1.611,2.875)rectangle++(0.04,0.04)
  (1.482,3.976)rectangle++(0.04,0.04) (1.417,4.916)rectangle++(0.04,0.04) (1.399,5.719)rectangle++(0.04,0.04) (1.415,6.401)rectangle++(0.04,0.04) (1.457,6.98)rectangle++(0.04,0.04)
  (1.515,7.471)rectangle++(0.04,0.04) (1.585,7.886)rectangle++(0.04,0.04) (-1.346,4.802)rectangle++(0.04,0.04) (-0.954,5.731)rectangle++(0.04,0.04) (-0.584,6.506)rectangle++(0.04,0.04)
  (-0.238,7.15)rectangle++(0.04,0.04) (0.081,7.683)rectangle++(0.04,0.04) (0.374,8.123)rectangle++(0.04,0.04) (0.641,8.486)rectangle++(0.04,0.04) (0.882,8.784)rectangle++(0.04,0.04)
  (1.099,9.027)rectangle++(0.04,0.04) (1.293,9.225)rectangle++(0.04,0.04) (0.46,3.075)rectangle++(0.04,0.04) (0.512,4.192)rectangle++(0.04,0.04) (0.601,5.139)rectangle++(0.04,0.04)
  (0.714,5.94)rectangle++(0.04,0.04) (0.842,6.617)rectangle++(0.04,0.04) (-1.21,4.412)rectangle++(0.04,0.04) (-0.854,5.395)rectangle++(0.04,0.04) (-0.512,6.216)rectangle++(0.04,0.04)
  (-0.189,6.9)rectangle++(0.04,0.04) (0.113,7.469)rectangle++(0.04,0.04) (0.393,7.94)rectangle++(0.04,0.04) (-1.571,4.589)rectangle++(0.04,0.04) (-1.154,5.559)rectangle++(0.04,0.04)
  (-0.761,6.368)rectangle++(0.04,0.04) (-0.394,7.039)rectangle++(0.04,0.04) (-0.056,7.595)rectangle++(0.04,0.04) (0.254,8.055)rectangle++(0.04,0.04) (0.536,8.432)rectangle++(0.04,0.04)
  (0.791,8.742)rectangle++(0.04,0.04) (1.02,8.996)rectangle++(0.04,0.04) (1.224,9.202)rectangle++(0.04,0.04) (1.407,9.369)rectangle++(0.04,0.04) (-1.693,5.061)rectangle++(0.04,0.04)
  (-1.238,5.966)rectangle++(0.04,0.04) (1.622,2.831)rectangle++(0.04,0.04) (1.49,3.938)rectangle++(0.04,0.04) (1.422,4.884)rectangle++(0.04,0.04) (1.402,5.691)rectangle++(0.04,0.04)
  (1.417,6.377)rectangle++(0.04,0.04) (1.457,6.96)rectangle++(0.04,0.04) (1.515,7.454)rectangle++(0.04,0.04) (1.584,7.872)rectangle++(0.04,0.04) (1.659,8.224)rectangle++(0.04,0.04)
  (1.737,8.52)rectangle++(0.04,0.04) (1.815,8.769)rectangle++(0.04,0.04) (1.891,8.977)rectangle++(0.04,0.04) (1.964,9.151)rectangle++(0.04,0.04) (2.033,9.296)rectangle++(0.04,0.04)
  (-1.565,5.164)rectangle++(0.04,0.04) (-1.126,6.048)rectangle++(0.04,0.04) (-0.717,6.782)rectangle++(0.04,0.04) (-0.341,7.39)rectangle++(0.04,0.04) (0.004,7.891)rectangle++(0.04,0.04)
  (0.317,8.303)rectangle++(0.04,0.04) (0.599,8.641)rectangle++(0.04,0.04) (0.853,8.918)rectangle++(0.04,0.04) (0.037,2.796)rectangle++(0.04,0.04) (0.688,2.51)rectangle++(0.04,0.04)
  (-0.012,1.747)rectangle++(0.04,0.04) (0.057,3.082)rectangle++(0.04,0.04) (0.17,4.214)rectangle++(0.04,0.04) (0.311,5.171)rectangle++(0.04,0.04) (0.469,5.979)rectangle++(0.04,0.04)
  (0.635,6.66)rectangle++(0.04,0.04) (0.804,7.232)rectangle++(0.04,0.04) (-0.122,3.77)rectangle++(0.04,0.04) (0.045,4.806)rectangle++(0.04,0.04) (0.806,2.188)rectangle++(0.04,0.04)
  (0.77,3.424)rectangle++(0.04,0.04) (0.789,4.476)rectangle++(0.04,0.04) (0.848,5.369)rectangle++(0.04,0.04) (0.933,6.126)rectangle++(0.04,0.04) (1.036,6.766)rectangle++(0.04,0.04)
  (1.149,7.306)rectangle++(0.04,0.04) (1.267,7.76)rectangle++(0.04,0.04) (1.385,8.142)rectangle++(0.04,0.04) (-1.436,4.815)rectangle++(0.04,0.04) (1.36,2.565)rectangle++(0.04,0.04)
  (1.257,3.722)rectangle++(0.04,0.04) (1.215,4.709)rectangle++(0.04,0.04) (1.219,5.551)rectangle++(0.04,0.04) (1.256,6.265)rectangle++(0.04,0.04) (1.404,3.455)rectangle++(0.04,0.04)
  (1.33,4.477)rectangle++(0.04,0.04) (1.307,5.349)rectangle++(0.04,0.04) (1.323,6.09)rectangle++(0.04,0.04) (1.366,6.72)rectangle++(0.04,0.04) (1.428,7.254)rectangle++(0.04,0.04)
  (1.501,7.705)rectangle++(0.04,0.04) (1.582,8.085)rectangle++(0.04,0.04) (0.454,2.955)rectangle++(0.04,0.04) (0.502,4.09)rectangle++(0.04,0.04) (0.588,5.052)rectangle++(0.04,0.04)
  (0.7,5.867)rectangle++(0.04,0.04) (0.827,6.555)rectangle++(0.04,0.04) (0.963,7.135)rectangle++(0.04,0.04) (1.102,7.623)rectangle++(0.04,0.04) (1.239,8.031)rectangle++(0.04,0.04)
  (1.372,8.373)rectangle++(0.04,0.04) (1.499,8.659)rectangle++(0.04,0.04) (1.619,8.896)rectangle++(0.04,0.04) (1.729,9.093)rectangle++(0.04,0.04) (1.831,9.256)rectangle++(0.04,0.04)
  (1.925,9.391)rectangle++(0.04,0.04) (2.01,9.501)rectangle++(0.04,0.04) (2.086,9.592)rectangle++(0.04,0.04) (2.155,9.666)rectangle++(0.04,0.04) (2.216,9.726)rectangle++(0.04,0.04)
  (2.27,9.775)rectangle++(0.04,0.04) (2.318,9.814)rectangle++(0.04,0.04) (2.361,9.845)rectangle++(0.04,0.04) (2.399,9.87)rectangle++(0.04,0.04) (-1.249,3.404)rectangle++(0.04,0.04)
  (-0.928,4.539)rectangle++(0.04,0.04) (-0.609,5.492)rectangle++(0.04,0.04) (-0.301,6.289)rectangle++(0.04,0.04) (-0.006,6.954)rectangle++(0.04,0.04) (0.271,7.507)rectangle++(0.04,0.04)
  (0.528,7.966)rectangle++(0.04,0.04) (0.765,8.346)rectangle++(0.04,0.04) (0.982,8.66)rectangle++(0.04,0.04) (1.179,8.918)rectangle++(0.04,0.04) (1.357,9.129)rectangle++(0.04,0.04)
  (2.102,4.031)rectangle++(0.04,0.04) (1.945,4.938)rectangle++(0.04,0.04) (1.849,5.716)rectangle++(0.04,0.04) (1.798,6.381)rectangle++(0.04,0.04) (1.781,6.948)rectangle++(0.04,0.04)
  (1.79,7.431)rectangle++(0.04,0.04) (1.816,7.841)rectangle++(0.04,0.04) (1.855,8.188)rectangle++(0.04,0.04) (1.902,8.482)rectangle++(0.04,0.04) (1.954,8.73)rectangle++(0.04,0.04)
  (2.008,8.938)rectangle++(0.04,0.04) (2.062,9.113)rectangle++(0.04,0.04) (-1.521,5.133)rectangle++(0.04,0.04) (-1.09,6.02)rectangle++(0.04,0.04) (-0.688,6.757)rectangle++(0.04,0.04)
  (-0.316,7.367)rectangle++(0.04,0.04) (0.024,7.871)rectangle++(0.04,0.04) (0.333,8.286)rectangle++(0.04,0.04) (0.612,8.626)rectangle++(0.04,0.04) (-1.677,4.765)rectangle++(0.04,0.04)
  (-1.237,5.714)rectangle++(0.04,0.04) (1.56,2.783)rectangle++(0.04,0.04) (1.435,3.899)rectangle++(0.04,0.04) (1.373,4.853)rectangle++(0.04,0.04) (1.359,5.667)rectangle++(0.04,0.04)
  (0.016,2.353)rectangle++(0.04,0.04) (0.106,3.596)rectangle++(0.04,0.04) (0.232,4.648)rectangle++(0.04,0.04) (1.124,2.911)rectangle++(0.04,0.04) (1.07,4.026)rectangle++(0.04,0.04)
  (0.872,2.981)rectangle++(0.04,0.04) (0.858,4.096)rectangle++(0.04,0.04) (0.891,5.043)rectangle++(0.04,0.04) (0.957,5.847)rectangle++(0.04,0.04) (1.045,6.528)rectangle++(0.04,0.04)
  (1.147,7.103)rectangle++(0.04,0.04) (1.257,7.588)rectangle++(0.04,0.04) (1.37,7.996)rectangle++(0.04,0.04) (1.482,8.338)rectangle++(0.04,0.04) (1.891,3.906)rectangle++(0.04,0.04)
  (1.762,4.841)rectangle++(0.04,0.04) (1.689,5.641)rectangle++(0.04,0.04) (1.659,6.323)rectangle++(0.04,0.04) (1.661,6.904)rectangle++(0.04,0.04) (1.515,3.671)rectangle++(0.04,0.04)
  (1.433,4.656)rectangle++(0.04,0.04) (1.402,5.496)rectangle++(0.04,0.04) (1.409,6.212)rectangle++(0.04,0.04) (1.444,6.82)rectangle++(0.04,0.04) (1.498,7.335)rectangle++(0.04,0.04)
  (1.565,7.771)rectangle++(0.04,0.04) (1.639,8.139)rectangle++(0.04,0.04) (1.824,3.904)rectangle++(0.04,0.04) (1.704,4.841)rectangle++(0.04,0.04) (1.64,5.643)rectangle++(0.04,0.04)
  (1.617,6.327)rectangle++(0.04,0.04) (1.626,6.91)rectangle++(0.04,0.04) (1.656,7.404)rectangle++(0.04,0.04) (1.702,7.824)rectangle++(0.04,0.04) (1.757,8.178)rectangle++(0.04,0.04)
  (1.818,8.477)rectangle++(0.04,0.04) (1.883,8.729)rectangle++(0.04,0.04) (1.947,8.941)rectangle++(0.04,0.04) (2.011,9.118)rectangle++(0.04,0.04) (2.071,9.266)rectangle++(0.04,0.04)
  (2.129,9.39)rectangle++(0.04,0.04) (2.183,9.492)rectangle++(0.04,0.04) (2.233,9.577)rectangle++(0.04,0.04) (-1.589,5.257)rectangle++(0.04,0.04) (-1.142,6.128)rectangle++(0.04,0.04)
  (-0.728,6.851)rectangle++(0.04,0.04) (-0.347,7.449)rectangle++(0.04,0.04) (0.001,7.941)rectangle++(0.04,0.04) (0.316,8.346)rectangle++(0.04,0.04) (0.6,8.678)rectangle++(0.04,0.04)
  (-1.691,4.773)rectangle++(0.04,0.04) (-1.248,5.721)rectangle++(0.04,0.04) (-0.834,6.509)rectangle++(0.04,0.04) (-1.487,4.076)rectangle++(0.04,0.04) (-1.103,5.12)rectangle++(0.04,0.04)
  (-0.735,5.993)rectangle++(0.04,0.04) (-0.388,6.719)rectangle++(0.04,0.04) (-0.063,7.323)rectangle++(0.04,0.04) (-1.522,4.382)rectangle++(0.04,0.04) (-1.121,5.381)rectangle++(0.04,0.04)
  (-0.739,6.215)rectangle++(0.04,0.04) (-0.382,6.909)rectangle++(0.04,0.04) (-1.491,4.24)rectangle++(0.04,0.04) (-1.1,5.26)rectangle++(0.04,0.04) (-0.727,6.111)rectangle++(0.04,0.04)
  (-0.376,6.82)rectangle++(0.04,0.04) (-0.049,7.408)rectangle++(0.04,0.04) (0.253,7.895)rectangle++(0.04,0.04) (0.528,8.297)rectangle++(0.04,0.04) (0.779,8.627)rectangle++(0.04,0.04)
  (1.005,8.898)rectangle++(0.04,0.04) (1.208,9.119)rectangle++(0.04,0.04) (1.389,9.299)rectangle++(0.04,0.04) (1.551,9.445)rectangle++(0.04,0.04) (0.487,3.104)rectangle++(0.04,0.04)
  (0.536,4.215)rectangle++(0.04,0.04) (0.096,2.315)rectangle++(0.04,0.04) (0.562,2.431)rectangle++(0.04,0.04) (0.573,3.64)rectangle++(0.04,0.04) (0.834,2.794)rectangle++(0.04,0.04)
  (0.819,3.938)rectangle++(0.04,0.04) (0.851,4.911)rectangle++(0.04,0.04) (-0.846,4.077)rectangle++(0.04,0.04) (-0.558,5.095)rectangle++(0.04,0.04) (-0.273,5.949)rectangle++(0.04,0.04)
  (1.505,3.047)rectangle++(0.04,0.04) (0.313,2.4)rectangle++(0.04,0.04) (0.36,3.624)rectangle++(0.04,0.04) (0.449,4.662)rectangle++(0.04,0.04) (0.566,5.541)rectangle++(0.04,0.04)
  (0.7,6.283)rectangle++(0.04,0.04) (0.844,6.909)rectangle++(0.04,0.04) (0.992,7.435)rectangle++(0.04,0.04) (1.138,7.876)rectangle++(0.04,0.04) (1.28,8.246)rectangle++(0.04,0.04)
  (1.416,8.554)rectangle++(0.04,0.04) (-1.518,4.902)rectangle++(0.04,0.04) (-1.096,5.823)rectangle++(0.04,0.04) (-0.701,6.59)rectangle++(0.04,0.04) (-0.335,7.226)rectangle++(0.04,0.04)
  (0.002,7.751)rectangle++(0.04,0.04) (-0.057,3.804)rectangle++(0.04,0.04) (0.101,4.832)rectangle++(0.04,0.04) (0.277,5.699)rectangle++(0.04,0.04) (0.283,2.672)rectangle++(0.04,0.04)
  (0.345,3.856)rectangle++(0.04,0.04) (-0.662,3.266)rectangle++(0.04,0.04) (-0.434,4.398)rectangle++(0.04,0.04) (-0.195,5.352)rectangle++(0.04,0.04) (0.046,6.153)rectangle++(0.04,0.04)
  (0.283,6.825)rectangle++(0.04,0.04) (0.511,7.386)rectangle++(0.04,0.04) (1.769,3.502)rectangle++(0.04,0.04) (1.642,4.503)rectangle++(0.04,0.04) (1.573,5.358)rectangle++(0.04,0.04)
  (1.549,6.087)rectangle++(0.04,0.04) (1.558,6.708)rectangle++(0.04,0.04) (1.591,7.236)rectangle++(0.04,0.04) (-1.207,4.674)rectangle++(0.04,0.04) (-0.841,5.618)rectangle++(0.04,0.04)
  (-0.493,6.405)rectangle++(0.04,0.04) (-0.165,7.06)rectangle++(0.04,0.04) (0.14,7.604)rectangle++(0.04,0.04) (0.421,8.054)rectangle++(0.04,0.04) (0.678,8.425)rectangle++(0.04,0.04)
  (0.911,8.73)rectangle++(0.04,0.04) (1.121,8.981)rectangle++(0.04,0.04) (1.31,9.185)rectangle++(0.04,0.04) (1.479,9.351)rectangle++(0.04,0.04) (2.142,4.101)rectangle++(0.04,0.04)
  (1.983,4.996)rectangle++(0.04,0.04) (1.883,5.764)rectangle++(0.04,0.04) (1.829,6.42)rectangle++(0.04,0.04) (1.809,6.98)rectangle++(0.04,0.04) (1.815,7.457)rectangle++(0.04,0.04)
  (1.839,7.862)rectangle++(0.04,0.04) (1.875,8.205)rectangle++(0.04,0.04) (1.92,8.496)rectangle++(0.04,0.04) (1.97,8.741)rectangle++(0.04,0.04) (2.022,8.947)rectangle++(0.04,0.04)
  (2.074,9.12)rectangle++(0.04,0.04) (2.126,9.266)rectangle++(0.04,0.04) (2.175,9.387)rectangle++(0.04,0.04) (2.222,9.488)rectangle++(0.04,0.04) (2.266,9.572)rectangle++(0.04,0.04)
  (2.307,9.642)rectangle++(0.04,0.04) (2.344,9.7)rectangle++(0.04,0.04) (2.378,9.747)rectangle++(0.04,0.04) (2.409,9.786)rectangle++(0.04,0.04) (2.437,9.818)rectangle++(0.04,0.04)
  (1.312,2.32)rectangle++(0.04,0.04) (1.205,3.515)rectangle++(0.04,0.04) (1.163,4.536)rectangle++(0.04,0.04) (1.168,5.405)rectangle++(0.04,0.04) (1.207,6.144)rectangle++(0.04,0.04)
  (1.269,6.77)rectangle++(0.04,0.04) (1.347,7.3)rectangle++(0.04,0.04) (1.435,7.748)rectangle++(0.04,0.04) (-1.344,4.746)rectangle++(0.04,0.04) (-0.954,5.684)rectangle++(0.04,0.04)
  (-0.586,6.466)rectangle++(0.04,0.04) (-0.242,7.116)rectangle++(0.04,0.04) (1.79,3.279)rectangle++(0.04,0.04) (1.651,4.312)rectangle++(0.04,0.04) (0.875,3.168)rectangle++(0.04,0.04)
  (0.868,4.254)rectangle++(0.04,0.04) (0.906,5.178)rectangle++(0.04,0.04) (0.975,5.961)rectangle++(0.04,0.04) (1.065,6.624)rectangle++(0.04,0.04) (1.168,7.184)rectangle++(0.04,0.04)
  (-0.983,3.316)rectangle++(0.04,0.04) (-0.705,4.454)rectangle++(0.04,0.04) (-0.423,5.41)rectangle++(0.04,0.04) (-0.146,6.212)rectangle++(0.04,0.04) (-1.301,4.147)rectangle++(0.04,0.04)
  (-0.942,5.173)rectangle++(0.04,0.04) (-0.596,6.031)rectangle++(0.04,0.04) (-0.268,6.746)rectangle++(0.04,0.04) (0.04,7.341)rectangle++(0.04,0.04) (0.326,7.835)rectangle++(0.04,0.04)
  (0.588,8.243)rectangle++(0.04,0.04) (0.827,8.579)rectangle++(0.04,0.04) (1.044,8.855)rectangle++(0.04,0.04) (1.24,9.081)rectangle++(0.04,0.04) (1.415,9.266)rectangle++(0.04,0.04)
  (1.571,9.416)rectangle++(0.04,0.04) (1.71,9.537)rectangle++(0.04,0.04) (1.832,9.634)rectangle++(0.04,0.04) (1.941,9.712)rectangle++(0.04,0.04) (2.036,9.774)rectangle++(0.04,0.04)
  (2.119,9.822)rectangle++(0.04,0.04) (-1.661,5.284)rectangle++(0.04,0.04) (-1.203,6.154)rectangle++(0.04,0.04) (0.595,2.407)rectangle++(0.04,0.04) (-0.361,3.534)rectangle++(0.04,0.04)
  (-0.168,4.615)rectangle++(0.04,0.04) (1.162,2.814)rectangle++(0.04,0.04) (1.098,3.942)rectangle++(0.04,0.04) (1.089,4.903)rectangle++(0.04,0.04) (1.12,5.72)rectangle++(0.04,0.04)
  (1.178,6.413)rectangle++(0.04,0.04) (1.256,7)rectangle++(0.04,0.04) (1.345,7.496)rectangle++(0.04,0.04) (1.078,2.26)rectangle++(0.04,0.04) (1.004,3.474)rectangle++(0.04,0.04)
  (0.99,4.509)rectangle++(0.04,0.04) (1.02,5.389)rectangle++(0.04,0.04) (1.08,6.136)rectangle++(0.04,0.04) (1.162,6.769)rectangle++(0.04,0.04) (1.256,7.303)rectangle++(0.04,0.04)
  (1.357,7.754)rectangle++(0.04,0.04) (1.462,8.133)rectangle++(0.04,0.04) (1.566,8.45)rectangle++(0.04,0.04) (1.667,8.716)rectangle++(0.04,0.04) (1.763,8.939)rectangle++(0.04,0.04)
  (1.854,9.123)rectangle++(0.04,0.04) (1.939,9.277)rectangle++(0.04,0.04) (2.017,9.404)rectangle++(0.04,0.04) (2.088,9.509)rectangle++(0.04,0.04) (2.153,9.595)rectangle++(0.04,0.04)
  (2.212,9.666)rectangle++(0.04,0.04) (2.264,9.724)rectangle++(0.04,0.04) (2.312,9.771)rectangle++(0.04,0.04) (2.353,9.809)rectangle++(0.04,0.04) (2.391,9.84)rectangle++(0.04,0.04)
  (0.008,2.168)rectangle++(0.04,0.04) (0.091,3.439)rectangle++(0.04,0.04) (0.213,4.515)rectangle++(0.04,0.04) (0.359,5.426)rectangle++(0.04,0.04) (0.52,6.194)rectangle++(0.04,0.04)
  (0.688,6.84)rectangle++(0.04,0.04) (1.613,3.437)rectangle++(0.04,0.04) (1.507,4.453)rectangle++(0.04,0.04) (1.457,5.321)rectangle++(0.04,0.04) (1.449,6.061)rectangle++(0.04,0.04)
  (1.472,6.69)rectangle++(0.04,0.04) (1.516,7.224)rectangle++(0.04,0.04) (1.576,7.676)rectangle++(0.04,0.04) (1.644,8.058)rectangle++(0.04,0.04) (1.718,8.379)rectangle++(0.04,0.04)
  (1.793,8.65)rectangle++(0.04,0.04) (1.868,8.877)rectangle++(0.04,0.04) (-1.506,5.05)rectangle++(0.04,0.04) (-1.08,5.949)rectangle++(0.04,0.04) (1.6,2.864)rectangle++(0.04,0.04)
  (0.522,2.878)rectangle++(0.04,0.04) (0.557,4.021)rectangle++(0.04,0.04) (0.632,4.992)rectangle++(0.04,0.04) (0.735,5.814)rectangle++(0.04,0.04) (0.855,6.509)rectangle++(0.04,0.04)
  (0.985,7.095)rectangle++(0.04,0.04) (1.119,7.587)rectangle++(0.04,0.04) (1.252,8.001)rectangle++(0.04,0.04) (1.382,8.347)rectangle++(0.04,0.04) (1.506,8.636)rectangle++(0.04,0.04)
  (1.624,8.876)rectangle++(0.04,0.04) (1.733,9.076)rectangle++(0.04,0.04) (1.834,9.241)rectangle++(0.04,0.04) (1.926,9.378)rectangle++(0.04,0.04) (0.04,2.858)rectangle++(0.04,0.04)
  (-0.556,3.519)rectangle++(0.04,0.04) (-0.334,4.61)rectangle++(0.04,0.04) (-0.102,5.528)rectangle++(0.04,0.04) (0.132,6.299)rectangle++(0.04,0.04) (0.362,6.945)rectangle++(0.04,0.04)
  (0.584,7.485)rectangle++(0.04,0.04) (0.793,7.935)rectangle++(0.04,0.04) (0.989,8.309)rectangle++(0.04,0.04) (1.171,8.62)rectangle++(0.04,0.04) (-1.576,4.869)rectangle++(0.04,0.04)
  (-1.147,5.798)rectangle++(0.04,0.04) (-0.745,6.57)rectangle++(0.04,0.04) (-0.373,7.211)rectangle++(0.04,0.04) (-0.031,7.74)rectangle++(0.04,0.04) (0.281,8.177)rectangle++(0.04,0.04)
  (1.992,3.603)rectangle++(0.04,0.04) (1.835,4.579)rectangle++(0.04,0.04) (-0.599,4.196)rectangle++(0.04,0.04) (-0.344,5.186)rectangle++(0.04,0.04) (-0.087,6.018)rectangle++(0.04,0.04)
  (0.165,6.715)rectangle++(0.04,0.04) (0.406,7.298)rectangle++(0.04,0.04) (0.635,7.783)rectangle++(0.04,0.04) (0.849,8.186)rectangle++(0.04,0.04) (1.047,8.521)rectangle++(0.04,0.04)
  (1.228,8.797)rectangle++(0.04,0.04) (1.394,9.024)rectangle++(0.04,0.04) (2.071,4.019)rectangle++(0.04,0.04) (1.919,4.929)rectangle++(0.04,0.04) (-0.659,4.281)rectangle++(0.04,0.04)
  (-0.391,5.261)rectangle++(0.04,0.04) (-1.143,3.913)rectangle++(0.04,0.04) (-0.817,4.968)rectangle++(0.04,0.04) (-0.498,5.852)rectangle++(0.04,0.04) (-0.192,6.59)rectangle++(0.04,0.04)
  (0.099,7.205)rectangle++(0.04,0.04) (0.37,7.717)rectangle++(0.04,0.04) (0.621,8.141)rectangle++(0.04,0.04) (0.851,8.491)rectangle++(0.04,0.04) (1.061,8.779)rectangle++(0.04,0.04)
  (1.251,9.016)rectangle++(0.04,0.04) (1.422,9.21)rectangle++(0.04,0.04) (1.575,9.368)rectangle++(0.04,0.04) (1.711,9.496)rectangle++(0.04,0.04) (1.832,9.599)rectangle++(0.04,0.04)
  (1.939,9.682)rectangle++(0.04,0.04) (1.283,2.277)rectangle++(0.04,0.04) (1.18,3.48)rectangle++(0.04,0.04) (1.14,4.507)rectangle++(0.04,0.04) (1.147,5.382)rectangle++(0.04,0.04)
  (1.198,3.26)rectangle++(0.04,0.04) (1.147,4.32)rectangle++(0.04,0.04) (1.145,5.222)rectangle++(0.04,0.04) (1.18,5.989)rectangle++(0.04,0.04) (1.241,6.64)rectangle++(0.04,0.04)
  (1.318,7.19)rectangle++(0.04,0.04) (1.406,7.655)rectangle++(0.04,0.04) (1.499,8.047)rectangle++(0.04,0.04) (1.594,8.376)rectangle++(0.04,0.04) (1.687,8.652)rectangle++(0.04,0.04)
  (1.778,8.883)rectangle++(0.04,0.04) (1.865,9.076)rectangle++(0.04,0.04) (1.946,9.236)rectangle++(0.04,0.04) (2.021,9.369)rectangle++(0.04,0.04) (2.091,9.479)rectangle++(0.04,0.04)
  (2.154,9.57)rectangle++(0.04,0.04) (2.211,9.644)rectangle++(0.04,0.04) (-1.607,5.266)rectangle++(0.04,0.04) (-1.157,6.137)rectangle++(0.04,0.04) (-0.741,6.859)rectangle++(0.04,0.04)
  (-0.357,7.456)rectangle++(0.04,0.04) (-0.027,3.766)rectangle++(0.04,0.04) (0.125,4.798)rectangle++(0.04,0.04) (0.296,5.67)rectangle++(0.04,0.04) (0.476,6.404)rectangle++(0.04,0.04)
  (0.659,7.02)rectangle++(0.04,0.04) (0.839,7.537)rectangle++(0.04,0.04) (1.012,7.969)rectangle++(0.04,0.04) (1.177,8.329)rectangle++(0.04,0.04) (1.331,8.629)rectangle++(0.04,0.04)
  (1.475,8.878)rectangle++(0.04,0.04) (0.466,3.04)rectangle++(0.04,0.04) (0.516,4.162)rectangle++(0.04,0.04) (0.602,5.113)rectangle++(0.04,0.04) (0.714,5.918)rectangle++(0.04,0.04)
  (0.842,6.598)rectangle++(0.04,0.04) (0.977,7.171)rectangle++(0.04,0.04) (1.115,7.652)rectangle++(0.04,0.04) (1.252,8.056)rectangle++(0.04,0.04) (1.384,8.394)rectangle++(0.04,0.04)
  (1.51,8.676)rectangle++(0.04,0.04) (1.628,8.91)rectangle++(0.04,0.04) (1.738,9.105)rectangle++(0.04,0.04) (1.84,9.266)rectangle++(0.04,0.04) (1.932,9.398)rectangle++(0.04,0.04)
  (2.016,9.508)rectangle++(0.04,0.04) (2.092,9.597)rectangle++(0.04,0.04) (2.16,9.67)rectangle++(0.04,0.04) (2.22,9.729)rectangle++(0.04,0.04) (0.55,3.146)rectangle++(0.04,0.04)
  (0.591,4.248)rectangle++(0.04,0.04) (0.67,5.183)rectangle++(0.04,0.04) (-0.008,1.834)rectangle++(0.04,0.04) (0.064,3.155)rectangle++(0.04,0.04) (0.179,4.276)rectangle++(0.04,0.04)
  (0.321,5.223)rectangle++(0.04,0.04) (0.479,6.023)rectangle++(0.04,0.04) (0.646,6.697)rectangle++(0.04,0.04) (0.815,7.263)rectangle++(0.04,0.04) (1.703,3.548)rectangle++(0.04,0.04)
  (1.587,4.544)rectangle++(0.04,0.04) (1.529,5.395)rectangle++(0.04,0.04) (-0.827,4.299)rectangle++(0.04,0.04) (-0.534,5.284)rectangle++(0.04,0.04) (-0.244,6.109)rectangle++(0.04,0.04)
  (0.034,6.798)rectangle++(0.04,0.04) (0.299,7.373)rectangle++(0.04,0.04) (0.547,7.852)rectangle++(0.04,0.04) (0.777,8.248)rectangle++(0.04,0.04) (0.988,8.576)rectangle++(0.04,0.04)
  (1.181,8.846)rectangle++(0.04,0.04) (1.355,9.068)rectangle++(0.04,0.04) (1.512,9.25)rectangle++(0.04,0.04) (-1.649,5.057)rectangle++(0.04,0.04) (-1.201,5.961)rectangle++(0.04,0.04)
  (-0.785,6.711)rectangle++(0.04,0.04) (-0.401,7.332)rectangle++(0.04,0.04) (-0.05,7.844)rectangle++(0.04,0.04) (0.269,8.266)rectangle++(0.04,0.04) (0.557,8.611)rectangle++(0.04,0.04)
  (0.816,8.894)rectangle++(0.04,0.04) (1.047,9.123)rectangle++(0.04,0.04) (1.253,9.309)rectangle++(0.04,0.04) (-1.708,5.02)rectangle++(0.04,0.04) (-1.253,5.932)rectangle++(0.04,0.04)
  (-0.83,6.688)rectangle++(0.04,0.04) (-0.44,7.314)rectangle++(0.04,0.04) (-0.084,7.831)rectangle++(0.04,0.04) (0.24,8.256)rectangle++(0.04,0.04) (0.532,8.604)rectangle++(0.04,0.04)
  (0.794,8.888)rectangle++(0.04,0.04) (1.028,9.12)rectangle++(0.04,0.04) (1.237,9.307)rectangle++(0.04,0.04) (1.421,9.457)rectangle++(0.04,0.04) (1.584,9.578)rectangle++(0.04,0.04)
  (1.727,9.674)rectangle++(0.04,0.04) (1.853,9.75)rectangle++(0.04,0.04) (1.963,9.81)rectangle++(0.04,0.04) (2.059,9.856)rectangle++(0.04,0.04) (2.142,9.891)rectangle++(0.04,0.04)
  (2.214,9.918)rectangle++(0.04,0.04) (-1.665,5.321)rectangle++(0.04,0.04) (-1.204,6.185)rectangle++(0.04,0.04) (-0.778,6.902)rectangle++(0.04,0.04) (-0.388,7.494)rectangle++(0.04,0.04)
  (1.901,3.319)rectangle++(0.04,0.04) (1.747,4.341)rectangle++(0.04,0.04) (1.656,5.217)rectangle++(0.04,0.04) (1.614,5.964)rectangle++(0.04,0.04) (1.608,6.601)rectangle++(0.04,0.04)
  (1.629,7.143)rectangle++(0.04,0.04) (1.668,7.602)rectangle++(0.04,0.04) (1.687,3.807)rectangle++(0.04,0.04) (1.584,4.764)rectangle++(0.04,0.04) (1.535,5.583)rectangle++(0.04,0.04)
  (1.526,6.28)rectangle++(0.04,0.04) (1.546,6.873)rectangle++(0.04,0.04) (1.587,7.377)rectangle++(0.04,0.04) (1.642,7.803)rectangle++(0.04,0.04) (1.705,8.163)rectangle++(0.04,0.04)
  (1.774,8.466)rectangle++(0.04,0.04) (1.844,8.722)rectangle++(0.04,0.04) (1.914,8.936)rectangle++(0.04,0.04) (1.982,9.115)rectangle++(0.04,0.04) (2.047,9.265)rectangle++(0.04,0.04)
  (2.109,9.389)rectangle++(0.04,0.04) (2.166,9.493)rectangle++(0.04,0.04) (2.218,9.578)rectangle++(0.04,0.04) (2.267,9.649)rectangle++(0.04,0.04) (2.31,9.707)rectangle++(0.04,0.04)
  (2.35,9.755)rectangle++(0.04,0.04) (2.385,9.794)rectangle++(0.04,0.04) (-1.607,5.328)rectangle++(0.04,0.04) (-1.155,6.19)rectangle++(0.04,0.04) (1.668,2.893)rectangle++(0.04,0.04)
  (1.531,3.989)rectangle++(0.04,0.04) (1.459,4.926)rectangle++(0.04,0.04) (1.435,5.725)rectangle++(0.04,0.04) (1.447,6.405)rectangle++(0.04,0.04) (1.484,6.982)rectangle++(0.04,0.04)
  (1.538,7.472)rectangle++(0.04,0.04) (1.604,7.886)rectangle++(0.04,0.04) (1.677,8.235)rectangle++(0.04,0.04) (1.752,8.529)rectangle++(0.04,0.04) (1.828,8.776)rectangle++(0.04,0.04)
  (1.903,8.982)rectangle++(0.04,0.04) (1.975,9.155)rectangle++(0.04,0.04) (2.042,9.299)rectangle++(0.04,0.04) (2.106,9.419)rectangle++(0.04,0.04) (2.165,9.518)rectangle++(0.04,0.04)
  (2.218,9.6)rectangle++(0.04,0.04) (2.267,9.667)rectangle++(0.04,0.04) (2.312,9.723)rectangle++(0.04,0.04) (-1.606,5.3)rectangle++(0.04,0.04) (-1.155,6.166)rectangle++(0.04,0.04)
  (-0.737,6.883)rectangle++(0.04,0.04) (-0.354,7.476)rectangle++(0.04,0.04) (-0.004,7.965)rectangle++(0.04,0.04) (0.313,8.367)rectangle++(0.04,0.04) (0.599,8.696)rectangle++(0.04,0.04)
  (0.855,8.963)rectangle++(0.04,0.04) (1.083,9.181)rectangle++(0.04,0.04) (1.285,9.357)rectangle++(0.04,0.04) (1.465,9.498)rectangle++(0.04,0.04) (1.623,9.611)rectangle++(0.04,0.04)
  (1.762,9.701)rectangle++(0.04,0.04) (1.883,9.771)rectangle++(0.04,0.04) (1.989,9.826)rectangle++(0.04,0.04) (2.082,9.869)rectangle++(0.04,0.04) (0.043,2.925)rectangle++(0.04,0.04)
  (0.152,4.081)rectangle++(0.04,0.04) (0.29,5.059)rectangle++(0.04,0.04) (0.447,5.884)rectangle++(0.04,0.04) (-0.852,3.263)rectangle++(0.04,0.04) (-0.596,4.404)rectangle++(0.04,0.04)
  (-0.333,5.363)rectangle++(0.04,0.04) (-0.071,6.168)rectangle++(0.04,0.04) (0.185,6.842)rectangle++(0.04,0.04) (0.428,7.405)rectangle++(0.04,0.04) (0.028,2.59)rectangle++(0.04,0.04)
  (0.125,3.796)rectangle++(0.04,0.04) (0.256,4.818)rectangle++(0.04,0.04) (0.408,5.681)rectangle++(0.04,0.04) (0.016,2.355)rectangle++(0.04,0.04) (0.106,3.598)rectangle++(0.04,0.04)
  (0.232,4.65)rectangle++(0.04,0.04) (1.124,2.912)rectangle++(0.04,0.04) (1.07,4.026)rectangle++(0.04,0.04) (1.068,4.976)rectangle++(0.04,0.04) (1.105,5.783)rectangle++(0.04,0.04)
  (1.168,6.467)rectangle++(0.04,0.04) (1.25,7.047)rectangle++(0.04,0.04) (-1.228,4.573)rectangle++(0.04,0.04) (-0.863,5.532)rectangle++(0.04,0.04) (-0.515,6.333)rectangle++(0.04,0.04)
  (-0.186,7)rectangle++(0.04,0.04) (0.119,7.554)rectangle++(0.04,0.04) (0.402,8.012)rectangle++(0.04,0.04) (0.66,8.39)rectangle++(0.04,0.04) (0.894,8.702)rectangle++(0.04,0.04)
  (1.106,8.957)rectangle++(0.04,0.04) (1.296,9.165)rectangle++(0.04,0.04) (1.466,9.335)rectangle++(0.04,0.04) (1.617,9.472)rectangle++(0.04,0.04) (2.156,4.156)rectangle++(0.04,0.04)
  (1.997,5.042)rectangle++(0.04,0.04) (1.897,5.802)rectangle++(0.04,0.04) (1.842,6.452)rectangle++(0.04,0.04) (1.822,7.007)rectangle++(0.04,0.04) (1.827,7.479)rectangle++(0.04,0.04)
  (1.638,3.819)rectangle++(0.04,0.04) (1.543,4.777)rectangle++(0.04,0.04) (1.5,5.595)rectangle++(0.04,0.04) (1.497,6.292)rectangle++(0.04,0.04) (1.522,6.884)rectangle++(0.04,0.04)
  (1.567,7.387)rectangle++(0.04,0.04) (1.625,7.812)rectangle++(0.04,0.04) (1.692,8.172)rectangle++(0.04,0.04) (1.763,8.475)rectangle++(0.04,0.04) (1.835,8.729)rectangle++(0.04,0.04)
  (1.907,8.943)rectangle++(0.04,0.04) (1.976,9.121)rectangle++(0.04,0.04) (2.042,9.27)rectangle++(0.04,0.04) (2.105,9.394)rectangle++(0.04,0.04) (2.162,9.497)rectangle++(0.04,0.04)
  (2.216,9.582)rectangle++(0.04,0.04) (2.265,9.652)rectangle++(0.04,0.04) (2.309,9.71)rectangle++(0.04,0.04) (2.349,9.757)rectangle++(0.04,0.04) (2.384,9.796)rectangle++(0.04,0.04)
  (2.416,9.828)rectangle++(0.04,0.04) (2.445,9.853)rectangle++(0.04,0.04) (2.47,9.873)rectangle++(0.04,0.04) (2.492,9.89)rectangle++(0.04,0.04) (2.512,9.903)rectangle++(0.04,0.04)
  (2.158,4.442)rectangle++(0.04,0.04) (2.01,5.285)rectangle++(0.04,0.04) (1.918,6.008)rectangle++(0.04,0.04) (1.868,6.627)rectangle++(0.04,0.04) (1.851,7.154)rectangle++(0.04,0.04)
  (1.857,7.603)rectangle++(0.04,0.04) (1.88,7.985)rectangle++(0.04,0.04) (1.916,8.308)rectangle++(0.04,0.04) (1.958,8.581)rectangle++(0.04,0.04) (2.006,8.812)rectangle++(0.04,0.04)
  (2.055,9.006)rectangle++(0.04,0.04) (-0.395,3.285)rectangle++(0.04,0.04) (-0.206,4.404)rectangle++(0.04,0.04) (-0.002,5.348)rectangle++(0.04,0.04) (0.21,6.142)rectangle++(0.04,0.04)
  (0.422,6.808)rectangle++(0.04,0.04) (0.629,7.367)rectangle++(0.04,0.04) (0.827,7.833)rectangle++(0.04,0.04) (1.014,8.221)rectangle++(0.04,0.04) (1.189,8.543)rectangle++(0.04,0.04)
  (1.35,8.81)rectangle++(0.04,0.04) (1.498,9.031)rectangle++(0.04,0.04) (1.632,9.213)rectangle++(0.04,0.04) (1.754,9.362)rectangle++(0.04,0.04) (1.863,9.484)rectangle++(0.04,0.04)
  (2.13,4.214)rectangle++(0.04,0.04) (1.977,5.093)rectangle++(0.04,0.04) (-0.683,4.324)rectangle++(0.04,0.04) (-0.41,5.299)rectangle++(0.04,0.04) (-0.139,6.116)rectangle++(0.04,0.04)
  (0.125,6.801)rectangle++(0.04,0.04) (0.376,7.372)rectangle++(0.04,0.04) (0.612,7.847)rectangle++(0.04,0.04) (0.832,8.242)rectangle++(0.04,0.04) (1.035,8.568)rectangle++(0.04,0.04)
  (1.22,8.838)rectangle++(0.04,0.04) (1.388,9.06)rectangle++(0.04,0.04) (1.54,9.241)rectangle++(0.04,0.04) (1.677,9.39)rectangle++(0.04,0.04) (1.799,9.51)rectangle++(0.04,0.04)
  (1.907,9.608)rectangle++(0.04,0.04) (2.003,9.687)rectangle++(0.04,0.04) (2.088,9.75)rectangle++(0.04,0.04) (2.162,9.8)rectangle++(0.04,0.04) (2.228,9.84)rectangle++(0.04,0.04)
  (2.285,9.871)rectangle++(0.04,0.04) (2.335,9.895)rectangle++(0.04,0.04) (2.378,9.914)rectangle++(0.04,0.04) (2.416,9.928)rectangle++(0.04,0.04) (2.448,9.938)rectangle++(0.04,0.04)
  (2.477,9.946)rectangle++(0.04,0.04) (2.501,9.951)rectangle++(0.04,0.04) (2.521,9.954)rectangle++(0.04,0.04) (2.539,9.957)rectangle++(0.04,0.04) (2.554,9.958)rectangle++(0.04,0.04)
  (2.167,4.462)rectangle++(0.04,0.04) (2.018,5.302)rectangle++(0.04,0.04) (0.001,2.021)rectangle++(0.04,0.04) (0.079,3.314)rectangle++(0.04,0.04) (0.198,4.41)rectangle++(0.04,0.04)
  (0.342,5.336)rectangle++(0.04,0.04) (0.502,6.119)rectangle++(0.04,0.04) (0.669,6.777)rectangle++(0.04,0.04) (-1.275,4.411)rectangle++(0.04,0.04) (-0.909,5.397)rectangle++(0.04,0.04)
  (-0.559,6.22)rectangle++(0.04,0.04) (-0.229,6.905)rectangle++(0.04,0.04) (0.08,7.475)rectangle++(0.04,0.04) (0.364,7.947)rectangle++(0.04,0.04) (0.625,8.336)rectangle++(0.04,0.04)
  (0.863,8.657)rectangle++(0.04,0.04) (1.077,8.92)rectangle++(0.04,0.04) (1.27,9.135)rectangle++(0.04,0.04) (1.443,9.31)rectangle++(0.04,0.04) (1.597,9.452)rectangle++(0.04,0.04)
  (2.154,4.147)rectangle++(0.04,0.04) (1.994,5.035)rectangle++(0.04,0.04) (1.894,5.796)rectangle++(0.04,0.04) (1.84,6.447)rectangle++(0.04,0.04) (1.82,7.003)rectangle++(0.04,0.04)
  (1.825,7.476)rectangle++(0.04,0.04) (1.848,7.878)rectangle++(0.04,0.04) (1.883,8.218)rectangle++(0.04,0.04) (-1.363,4.923)rectangle++(0.04,0.04) (-0.964,5.835)rectangle++(0.04,0.04)
  (-1.368,3.919)rectangle++(0.04,0.04) (-1.008,4.982)rectangle++(0.04,0.04) (-0.66,5.871)rectangle++(0.04,0.04) (-0.328,6.613)rectangle++(0.04,0.04) (-0.017,7.23)rectangle++(0.04,0.04)
  (0.273,7.743)rectangle++(0.04,0.04) (0.54,8.167)rectangle++(0.04,0.04) (1.959,3.659)rectangle++(0.04,0.04) (1.809,4.628)rectangle++(0.04,0.04) (1.721,5.458)rectangle++(0.04,0.04)
  (1.679,6.167)rectangle++(0.04,0.04) (1.672,6.771)rectangle++(0.04,0.04) (1.481,3.647)rectangle++(0.04,0.04) (0.327,2.466)rectangle++(0.04,0.04) (0.375,3.679)rectangle++(0.04,0.04)
  (0.867,2.757)rectangle++(0.04,0.04) (-0.124,3.325)rectangle++(0.04,0.04) (0.838,2.575)rectangle++(0.04,0.04) (0.813,3.752)rectangle++(0.04,0.04) (0.839,4.753)rectangle++(0.04,0.04)
  (0.901,5.602)rectangle++(0.04,0.04) (0.988,6.322)rectangle++(0.04,0.04) (1.09,6.93)rectangle++(0.04,0.04) (1.202,7.444)rectangle++(0.04,0.04) (1.317,7.875)rectangle++(0.04,0.04)
  (1.432,8.237)rectangle++(0.04,0.04) (1.545,8.541)rectangle++(0.04,0.04) (-1.492,4.923)rectangle++(0.04,0.04) (1.394,2.573)rectangle++(0.04,0.04) (1.286,3.727)rectangle++(0.04,0.04)
  (1.24,4.713)rectangle++(0.04,0.04) (1.022,3.152)rectangle++(0.04,0.04) (0.993,4.235)rectangle++(0.04,0.04) (1.011,5.156)rectangle++(0.04,0.04) (1.063,5.939)rectangle++(0.04,0.04)
  (1.139,6.601)rectangle++(0.04,0.04) (1.23,7.162)rectangle++(0.04,0.04) (-1.256,4.592)rectangle++(0.04,0.04) (-0.886,5.55)rectangle++(0.04,0.04) (-0.533,6.349)rectangle++(0.04,0.04)
  (1.635,3.065)rectangle++(0.04,0.04) (1.51,4.136)rectangle++(0.04,0.04) (1.447,5.052)rectangle++(0.04,0.04) (1.43,5.832)rectangle++(0.04,0.04) (1.446,6.496)rectangle++(0.04,0.04)
  (1.487,7.06)rectangle++(0.04,0.04) (1.544,7.538)rectangle++(0.04,0.04) (1.612,7.942)rectangle++(0.04,0.04) (1.686,8.282)rectangle++(0.04,0.04) (1.762,8.569)rectangle++(0.04,0.04)
  (1.838,8.809)rectangle++(0.04,0.04) (1.913,9.01)rectangle++(0.04,0.04) (1.984,9.179)rectangle++(0.04,0.04) (2.051,9.319)rectangle++(0.04,0.04) (2.114,9.435)rectangle++(0.04,0.04)
  (2.172,9.531)rectangle++(0.04,0.04) (-1.59,5.236)rectangle++(0.04,0.04) (-1.144,6.111)rectangle++(0.04,0.04) (-0.73,6.836)rectangle++(0.04,0.04) (-0.35,7.436)rectangle++(0.04,0.04)
  (-0.002,7.931)rectangle++(0.04,0.04) (0.313,8.338)rectangle++(0.04,0.04) (0.598,8.671)rectangle++(0.04,0.04) (0.853,8.942)rectangle++(0.04,0.04) (1.08,9.163)rectangle++(0.04,0.04)
  (1.283,9.342)rectangle++(0.04,0.04) (1.462,9.485)rectangle++(0.04,0.04) (1.62,9.6)rectangle++(0.04,0.04) (2.188,4.181)rectangle++(0.04,0.04) (-0.452,4.183)rectangle++(0.04,0.04)
  (-0.219,5.17)rectangle++(0.04,0.04) (0.018,5.999)rectangle++(0.04,0.04) (1.484,3.123)rectangle++(0.04,0.04) (1.384,4.192)rectangle++(0.04,0.04) (1.342,5.104)rectangle++(0.04,0.04)
  (-0.791,3.341)rectangle++(0.04,0.04) (-0.541,4.468)rectangle++(0.04,0.04) (-1.002,3.728)rectangle++(0.04,0.04) (-0.705,4.805)rectangle++(0.04,0.04) (-0.409,5.709)rectangle++(0.04,0.04)
  (-0.121,6.465)rectangle++(0.04,0.04) (0.153,7.096)rectangle++(0.04,0.04) (0.412,7.622)rectangle++(0.04,0.04) (0.653,8.058)rectangle++(0.04,0.04) (0.875,8.42)rectangle++(0.04,0.04)
  (1.982,3.784)rectangle++(0.04,0.04) (1.834,4.734)rectangle++(0.04,0.04) (1.746,5.546)rectangle++(0.04,0.04) (1.704,6.241)rectangle++(0.04,0.04) (1.696,6.833)rectangle++(0.04,0.04)
  (1.713,7.336)rectangle++(0.04,0.04) (1.747,7.763)rectangle++(0.04,0.04) (1.793,8.125)rectangle++(0.04,0.04) (1.847,8.431)rectangle++(0.04,0.04) (1.905,8.689)rectangle++(0.04,0.04)
  (-1.459,5.02)rectangle++(0.04,0.04) (-0.436,3.384)rectangle++(0.04,0.04) (-0.237,4.49)rectangle++(0.04,0.04) (-0.024,5.422)rectangle++(0.04,0.04) (0.194,6.206)rectangle++(0.04,0.04)
  (0.411,6.864)rectangle++(0.04,0.04) (0.622,7.414)rectangle++(0.04,0.04) (0.823,7.873)rectangle++(0.04,0.04) (1.012,8.255)rectangle++(0.04,0.04) (1.188,8.573)rectangle++(0.04,0.04)
  (1.351,8.836)rectangle++(0.04,0.04) (1.499,9.052)rectangle++(0.04,0.04) (1.634,9.231)rectangle++(0.04,0.04) (1.756,9.377)rectangle++(0.04,0.04) (1.866,9.496)rectangle++(0.04,0.04)
  (1.963,9.593)rectangle++(0.04,0.04) (2.051,9.672)rectangle++(0.04,0.04) (2.128,9.735)rectangle++(0.04,0.04) (2.196,9.786)rectangle++(0.04,0.04) (2.256,9.827)rectangle++(0.04,0.04)
  (2.308,9.859)rectangle++(0.04,0.04) (2.354,9.884)rectangle++(0.04,0.04) (2.394,9.903)rectangle++(0.04,0.04) (2.429,9.918)rectangle++(0.04,0.04) (2.459,9.929)rectangle++(0.04,0.04)
  (2.485,9.938)rectangle++(0.04,0.04) (2.508,9.944)rectangle++(0.04,0.04) (2.527,9.948)rectangle++(0.04,0.04) (2.544,9.951)rectangle++(0.04,0.04) (2.166,4.458)rectangle++(0.04,0.04)
  (-0.515,4.234)rectangle++(0.04,0.04) (-0.27,5.216)rectangle++(0.04,0.04) (-0.023,6.04)rectangle++(0.04,0.04) (0.22,6.731)rectangle++(0.04,0.04) (0.454,7.309)rectangle++(0.04,0.04)
  (0.676,7.791)rectangle++(0.04,0.04) (0.884,8.191)rectangle++(0.04,0.04) (1.077,8.523)rectangle++(0.04,0.04) (1.254,8.798)rectangle++(0.04,0.04) (1.416,9.024)rectangle++(0.04,0.04)
  (1.562,9.21)rectangle++(0.04,0.04) (1.694,9.362)rectangle++(0.04,0.04) (1.812,9.487)rectangle++(0.04,0.04) (1.918,9.587)rectangle++(0.04,0.04) (2.011,9.669)rectangle++(0.04,0.04)
  (2.159,4.283)rectangle++(0.04,0.04) (2.004,5.15)rectangle++(0.04,0.04) (1.907,5.894)rectangle++(0.04,0.04) (1.855,6.53)rectangle++(0.04,0.04) (1.835,7.072)rectangle++(0.04,0.04)
  (1.841,7.534)rectangle++(0.04,0.04) (1.864,7.927)rectangle++(0.04,0.04) (1.899,8.259)rectangle++(0.04,0.04) (1.942,8.541)rectangle++(0.04,0.04) (1.991,8.778)rectangle++(0.04,0.04)
  (2.041,8.978)rectangle++(0.04,0.04) (2.092,9.146)rectangle++(0.04,0.04) (2.142,9.286)rectangle++(0.04,0.04) (2.19,9.404)rectangle++(0.04,0.04) (2.235,9.502)rectangle++(0.04,0.04)
  (2.278,9.584)rectangle++(0.04,0.04) (2.317,9.651)rectangle++(0.04,0.04) (2.353,9.707)rectangle++(0.04,0.04) (-0.251,4.053)rectangle++(0.04,0.04) (-0.054,5.052)rectangle++(0.04,0.04)
  (0.154,5.892)rectangle++(0.04,0.04) (0.365,6.598)rectangle++(0.04,0.04) (0.572,7.19)rectangle++(0.04,0.04) (0.771,7.685)rectangle++(0.04,0.04) (1.813,3.619)rectangle++(0.04,0.04)
  (1.683,4.6)rectangle++(0.04,0.04) (1.613,5.439)rectangle++(0.04,0.04) (1.586,6.155)rectangle++(0.04,0.04) (1.592,6.764)rectangle++(0.04,0.04) (1.622,7.282)rectangle++(0.04,0.04)
  (1.668,7.721)rectangle++(0.04,0.04) (1.724,8.092)rectangle++(0.04,0.04) (1.802,3.914)rectangle++(0.04,0.04) (1.686,4.851)rectangle++(0.04,0.04) (1.625,5.652)rectangle++(0.04,0.04)
  (-0.865,4.368)rectangle++(0.04,0.04) (-0.563,5.344)rectangle++(0.04,0.04) (0.014,2.309)rectangle++(0.04,0.04) (0.102,3.559)rectangle++(0.04,0.04) (0.227,4.617)rectangle++(0.04,0.04)
  (0.376,5.512)rectangle++(0.04,0.04) (0.537,6.266)rectangle++(0.04,0.04) (0.705,6.901)rectangle++(0.04,0.04) (-1.295,4.442)rectangle++(0.04,0.04) (0.052,3.443)rectangle++(0.04,0.04)
  (0.18,4.521)rectangle++(0.04,0.04) (0.331,5.432)rectangle++(0.04,0.04) (0.497,6.2)rectangle++(0.04,0.04) (0.668,6.846)rectangle++(0.04,0.04) (0.839,7.389)rectangle++(0.04,0.04)
  (1.007,7.843)rectangle++(0.04,0.04) (1.167,8.222)rectangle++(0.04,0.04) (1.319,8.539)rectangle++(0.04,0.04) (1.46,8.801)rectangle++(0.04,0.04) (2.008,3.991)rectangle++(0.04,0.04)
  (1.865,4.908)rectangle++(0.04,0.04) (1.779,5.693)rectangle++(0.04,0.04) (1.202,3.464)rectangle++(0.04,0.04) (1.158,4.492)rectangle++(0.04,0.04) (-0.703,4.052)rectangle++(0.04,0.04)
  (-0.437,5.068)rectangle++(0.04,0.04) (-0.171,5.922)rectangle++(0.04,0.04) (0.089,6.636)rectangle++(0.04,0.04) (0.339,7.234)rectangle++(0.04,0.04) (0.575,7.731)rectangle++(0.04,0.04)
  (0.796,8.145)rectangle++(0.04,0.04) (1,8.487)rectangle++(0.04,0.04) (1.188,8.77)rectangle++(0.04,0.04) (1.358,9.004)rectangle++(0.04,0.04) (1.512,9.195)rectangle++(0.04,0.04)
  (1.651,9.351)rectangle++(0.04,0.04) (1.775,9.479)rectangle++(0.04,0.04) (1.886,9.582)rectangle++(0.04,0.04) (1.984,9.666)rectangle++(0.04,0.04) (2.071,9.733)rectangle++(0.04,0.04)
  (2.168,4.309)rectangle++(0.04,0.04) (2.013,5.172)rectangle++(0.04,0.04) (1.915,5.912)rectangle++(0.04,0.04) (1.862,6.545)rectangle++(0.04,0.04) (1.843,7.085)rectangle++(0.04,0.04)
  (1.847,7.545)rectangle++(0.04,0.04) (1.87,7.935)rectangle++(0.04,0.04) (1.905,8.266)rectangle++(0.04,0.04) (1.947,8.546)rectangle++(0.04,0.04) (1.995,8.783)rectangle++(0.04,0.04)
  (2.045,8.982)rectangle++(0.04,0.04) (2.095,9.149)rectangle++(0.04,0.04) (2.145,9.289)rectangle++(0.04,0.04) (2.192,9.406)rectangle++(0.04,0.04) (2.237,9.504)rectangle++(0.04,0.04)
  (2.28,9.585)rectangle++(0.04,0.04) (2.319,9.652)rectangle++(0.04,0.04) (2.355,9.708)rectangle++(0.04,0.04) (2.388,9.754)rectangle++(0.04,0.04) (2.418,9.791)rectangle++(0.04,0.04)
  (2.444,9.822)rectangle++(0.04,0.04) (2.146,4.411)rectangle++(0.04,0.04) (-0.508,4.221)rectangle++(0.04,0.04) (-0.265,5.204)rectangle++(0.04,0.04) (-0.02,6.03)rectangle++(0.04,0.04)
  (0.222,6.723)rectangle++(0.04,0.04) (0.456,7.302)rectangle++(0.04,0.04) (-1.424,4.475)rectangle++(0.04,0.04) (-1.034,5.457)rectangle++(0.04,0.04) (-0.663,6.276)rectangle++(0.04,0.04)
  (-0.315,6.957)rectangle++(0.04,0.04) (0.009,7.522)rectangle++(0.04,0.04) (0.306,7.99)rectangle++(0.04,0.04) (1.943,3.572)rectangle++(0.04,0.04) (1.792,4.555)rectangle++(0.04,0.04)
  (1.703,5.396)rectangle++(0.04,0.04) (1.661,6.115)rectangle++(0.04,0.04) (1.655,6.727)rectangle++(0.04,0.04) (1.673,7.248)rectangle++(0.04,0.04) (-1.195,4.692)rectangle++(0.04,0.04);
\fill[piece2]
  (-0.83,5.632)rectangle++(0.04,0.04) (-0.483,6.417)rectangle++(0.04,0.04) (-1.405,4.124)rectangle++(0.04,0.04) (-1.031,5.158)rectangle++(0.04,0.04) (-0.673,6.022)rectangle++(0.04,0.04)
  (-0.333,6.741)rectangle++(0.04,0.04) (-0.016,7.34)rectangle++(0.04,0.04) (0.278,7.836)rectangle++(0.04,0.04) (0.548,8.245)rectangle++(0.04,0.04) (0.793,8.583)rectangle++(0.04,0.04)
  (1.015,8.86)rectangle++(0.04,0.04) (1.215,9.086)rectangle++(0.04,0.04) (1.394,9.271)rectangle++(0.04,0.04) (1.554,9.421)rectangle++(0.04,0.04) (1.695,9.542)rectangle++(0.04,0.04)
  (1.82,9.639)rectangle++(0.04,0.04) (1.931,9.716)rectangle++(0.04,0.04) (2.18,4.274)rectangle++(0.04,0.04) (2.022,5.142)rectangle++(0.04,0.04) (1.922,5.886)rectangle++(0.04,0.04)
  (1.867,6.522)rectangle++(0.04,0.04) (1.846,7.066)rectangle++(0.04,0.04) (1.849,7.528)rectangle++(0.04,0.04) (1.871,7.921)rectangle++(0.04,0.04) (1.905,8.254)rectangle++(0.04,0.04)
  (1.947,8.536)rectangle++(0.04,0.04) (1.994,8.774)rectangle++(0.04,0.04) (-0.221,3.952)rectangle++(0.04,0.04) (1.005,2.674)rectangle++(0.04,0.04) (0.959,3.829)rectangle++(0.04,0.04)
  (0.836,2.918)rectangle++(0.04,0.04) (0.825,4.043)rectangle++(0.04,0.04) (0.861,5)rectangle++(0.04,0.04) (0.93,5.812)rectangle++(0.04,0.04) (1.02,6.499)rectangle++(0.04,0.04)
  (1.125,7.08)rectangle++(0.04,0.04) (1.621,3.583)rectangle++(0.04,0.04) (1.519,4.577)rectangle++(0.04,0.04) (1.472,5.426)rectangle++(0.04,0.04) (1.171,3.343)rectangle++(0.04,0.04)
  (1.127,4.391)rectangle++(0.04,0.04) (1.131,5.283)rectangle++(0.04,0.04) (1.171,6.042)rectangle++(0.04,0.04) (1.235,6.685)rectangle++(0.04,0.04) (1.315,7.229)rectangle++(0.04,0.04)
  (1.404,7.688)rectangle++(0.04,0.04) (1.499,8.075)rectangle++(0.04,0.04) (1.595,8.4)rectangle++(0.04,0.04) (1.69,8.672)rectangle++(0.04,0.04) (1.781,8.9)rectangle++(0.04,0.04)
  (1.867,9.09)rectangle++(0.04,0.04) (-1.551,5.092)rectangle++(0.04,0.04) (-1.117,5.987)rectangle++(0.04,0.04) (-0.712,6.73)rectangle++(0.04,0.04) (-1.512,4.143)rectangle++(0.04,0.04)
  (-1.122,5.178)rectangle++(0.04,0.04) (-0.749,6.042)rectangle++(0.04,0.04) (-0.397,6.762)rectangle++(0.04,0.04) (-0.069,7.36)rectangle++(0.04,0.04) (0.233,7.855)rectangle++(0.04,0.04)
  (0.51,8.263)rectangle++(0.04,0.04) (0.371,2.957)rectangle++(0.04,0.04) (0.431,4.095)rectangle++(0.04,0.04) (-0.748,3.837)rectangle++(0.04,0.04) (-0.485,4.887)rectangle++(0.04,0.04)
  (-0.219,5.77)rectangle++(0.04,0.04) (0.043,6.509)rectangle++(0.04,0.04) (0.294,7.127)rectangle++(0.04,0.04) (0.533,7.643)rectangle++(0.04,0.04) (0.029,2.622)rectangle++(0.04,0.04)
  (0.127,3.824)rectangle++(0.04,0.04) (0.259,4.841)rectangle++(0.04,0.04) (0.412,5.701)rectangle++(0.04,0.04) (0.576,6.426)rectangle++(0.04,0.04) (0.021,2.457)rectangle++(0.04,0.04)
  (0.114,3.683)rectangle++(0.04,0.04) (0.242,4.723)rectangle++(0.04,0.04) (0.393,5.601)rectangle++(0.04,0.04) (0.555,6.341)rectangle++(0.04,0.04) (0.724,6.964)rectangle++(0.04,0.04)
  (0.891,7.487)rectangle++(0.04,0.04) (1.055,7.924)rectangle++(0.04,0.04) (1.211,8.29)rectangle++(0.04,0.04) (1.359,8.594)rectangle++(0.04,0.04) (1.497,8.847)rectangle++(0.04,0.04)
  (1.624,9.056)rectangle++(0.04,0.04) (-1.587,5.04)rectangle++(0.04,0.04) (0.113,2.569)rectangle++(0.04,0.04) (0.197,3.776)rectangle++(0.04,0.04) (0.316,4.798)rectangle++(0.04,0.04)
  (0.458,5.662)rectangle++(0.04,0.04) (0.614,6.39)rectangle++(0.04,0.04) (0.775,7.003)rectangle++(0.04,0.04) (-1.304,4.475)rectangle++(0.04,0.04) (-0.931,5.453)rectangle++(0.04,0.04)
  (-0.576,6.268)rectangle++(0.04,0.04) (-0.241,6.947)rectangle++(0.04,0.04) (0.071,7.511)rectangle++(0.04,0.04) (0.358,7.978)rectangle++(0.04,0.04) (0.622,8.363)rectangle++(0.04,0.04)
  (0.861,8.68)rectangle++(0.04,0.04) (-0.128,3.916)rectangle++(0.04,0.04) (0.046,4.93)rectangle++(0.04,0.04) (0.234,5.785)rectangle++(0.04,0.04) (0.428,6.504)rectangle++(0.04,0.04)
  (0.622,7.108)rectangle++(0.04,0.04) (0.81,7.613)rectangle++(0.04,0.04) (0.991,8.035)rectangle++(0.04,0.04) (1.874,3.736)rectangle++(0.04,0.04) (1.74,4.697)rectangle++(0.04,0.04)
  (1.664,5.519)rectangle++(0.04,0.04) (1.633,6.221)rectangle++(0.04,0.04) (1.635,6.819)rectangle++(0.04,0.04) (-0.19,3.748)rectangle++(0.04,0.04) (-0.014,4.79)rectangle++(0.04,0.04)
  (0.177,5.668)rectangle++(0.04,0.04) (-1.128,4.1)rectangle++(0.04,0.04) (-0.797,5.126)rectangle++(0.04,0.04) (-0.474,5.985)rectangle++(0.04,0.04) (-0.166,6.703)rectangle++(0.04,0.04)
  (0.125,7.3)rectangle++(0.04,0.04) (0.396,7.796)rectangle++(0.04,0.04) (0.646,8.207)rectangle++(0.04,0.04) (0.875,8.547)rectangle++(0.04,0.04) (-1.613,4.799)rectangle++(0.04,0.04)
  (-1.182,5.74)rectangle++(0.04,0.04) (-0.777,6.522)rectangle++(0.04,0.04) (-0.402,7.171)rectangle++(0.04,0.04) (-0.057,7.708)rectangle++(0.04,0.04) (1.915,3.436)rectangle++(0.04,0.04)
  (1.763,4.44)rectangle++(0.04,0.04) (-0.583,4.155)rectangle++(0.04,0.04) (-0.331,5.151)rectangle++(0.04,0.04) (-0.078,5.988)rectangle++(0.04,0.04) (0.171,6.689)rectangle++(0.04,0.04)
  (0.411,7.275)rectangle++(0.04,0.04) (0.638,7.763)rectangle++(0.04,0.04) (0.851,8.17)rectangle++(0.04,0.04) (1.048,8.506)rectangle++(0.04,0.04) (1.229,8.785)rectangle++(0.04,0.04)
  (1.393,9.014)rectangle++(0.04,0.04) (1.543,9.202)rectangle++(0.04,0.04) (1.677,9.357)rectangle++(0.04,0.04) (1.798,9.482)rectangle++(0.04,0.04) (1.905,9.584)rectangle++(0.04,0.04)
  (2.001,9.667)rectangle++(0.04,0.04) (2.085,9.733)rectangle++(0.04,0.04) (2.159,9.786)rectangle++(0.04,0.04) (2.225,9.828)rectangle++(0.04,0.04) (2.282,9.861)rectangle++(0.04,0.04)
  (2.332,9.886)rectangle++(0.04,0.04) (2.375,9.906)rectangle++(0.04,0.04) (2.413,9.922)rectangle++(0.04,0.04) (2.446,9.933)rectangle++(0.04,0.04) (2.474,9.942)rectangle++(0.04,0.04)
  (-1.623,5.374)rectangle++(0.04,0.04) (-1.167,6.229)rectangle++(0.04,0.04) (-1.488,3.959)rectangle++(0.04,0.04) (-1.108,5.02)rectangle++(0.04,0.04) (1.373,2.679)rectangle++(0.04,0.04)
  (1.272,3.818)rectangle++(0.04,0.04) (-0.539,3.94)rectangle++(0.04,0.04) (-0.303,4.967)rectangle++(0.04,0.04) (-0.061,5.83)rectangle++(0.04,0.04) (0.179,6.554)rectangle++(0.04,0.04)
  (0.412,7.16)rectangle++(0.04,0.04) (0.635,7.666)rectangle++(0.04,0.04) (0.844,8.087)rectangle++(0.04,0.04) (1.039,8.436)rectangle++(0.04,0.04) (1.218,8.725)rectangle++(0.04,0.04)
  (1.382,8.964)rectangle++(0.04,0.04) (1.531,9.16)rectangle++(0.04,0.04) (1.666,9.321)rectangle++(0.04,0.04) (1.787,9.453)rectangle++(0.04,0.04) (1.894,9.559)rectangle++(0.04,0.04)
  (1.99,9.646)rectangle++(0.04,0.04) (2.076,9.716)rectangle++(0.04,0.04) (2.151,9.771)rectangle++(0.04,0.04) (2.217,9.816)rectangle++(0.04,0.04) (2.275,9.851)rectangle++(0.04,0.04)
  (2.325,9.879)rectangle++(0.04,0.04) (2.369,9.9)rectangle++(0.04,0.04) (2.408,9.916)rectangle++(0.04,0.04) (2.441,9.929)rectangle++(0.04,0.04) (0.044,2.933)rectangle++(0.04,0.04)
  (0.152,4.088)rectangle++(0.04,0.04) (0.291,5.065)rectangle++(0.04,0.04) (0.447,5.889)rectangle++(0.04,0.04) (-1.126,4.194)rectangle++(0.04,0.04) (-0.792,5.206)rectangle++(0.04,0.04)
  (-0.467,6.053)rectangle++(0.04,0.04) (-0.157,6.76)rectangle++(0.04,0.04) (0.135,7.349)rectangle++(0.04,0.04) (1.804,3.41)rectangle++(0.04,0.04) (1.668,4.422)rectangle++(0.04,0.04)
  (-0.596,4.133)rectangle++(0.04,0.04) (-0.344,5.133)rectangle++(0.04,0.04) (-0.089,5.973)rectangle++(0.04,0.04) (0.161,6.677)rectangle++(0.04,0.04) (0.402,7.265)rectangle++(0.04,0.04)
  (0.63,7.756)rectangle++(0.04,0.04) (0.843,8.163)rectangle++(0.04,0.04) (1.041,8.501)rectangle++(0.04,0.04) (1.223,8.781)rectangle++(0.04,0.04) (1.389,9.011)rectangle++(0.04,0.04)
  (1.538,9.2)rectangle++(0.04,0.04) (1.674,9.355)rectangle++(0.04,0.04) (1.794,9.481)rectangle++(0.04,0.04) (1.902,9.583)rectangle++(0.04,0.04) (1.998,9.666)rectangle++(0.04,0.04)
  (-1.65,5.23)rectangle++(0.04,0.04) (-1.195,6.108)rectangle++(0.04,0.04) (1.652,2.869)rectangle++(0.04,0.04) (1.517,3.968)rectangle++(0.04,0.04) (0.006,2.122)rectangle++(0.04,0.04)
  (0.525,2.373)rectangle++(0.04,0.04) (0.539,3.592)rectangle++(0.04,0.04) (-0.622,3.758)rectangle++(0.04,0.04) (-0.381,4.815)rectangle++(0.04,0.04) (-0.133,5.704)rectangle++(0.04,0.04)
  (-1.191,4.049)rectangle++(0.04,0.04) (-0.853,5.085)rectangle++(0.04,0.04) (-0.524,5.953)rectangle++(0.04,0.04) (-0.209,6.677)rectangle++(0.04,0.04) (0.087,7.28)rectangle++(0.04,0.04)
  (0.363,7.781)rectangle++(0.04,0.04) (0.618,8.195)rectangle++(0.04,0.04) (0.851,8.537)rectangle++(0.04,0.04) (2.014,3.802)rectangle++(0.04,0.04) (0.706,3.153)rectangle++(0.04,0.04)
  (0.724,4.248)rectangle++(0.04,0.04) (0.783,5.178)rectangle++(0.04,0.04) (0.87,5.966)rectangle++(0.04,0.04) (0.976,6.633)rectangle++(0.04,0.04) (-1.189,4.44)rectangle++(0.04,0.04)
  (1.239,2.548)rectangle++(0.04,0.04) (1.153,3.713)rectangle++(0.04,0.04) (1.126,4.706)rectangle++(0.04,0.04) (1.143,5.551)rectangle++(0.04,0.04) (1.192,6.269)rectangle++(0.04,0.04)
  (1.261,6.877)rectangle++(0.04,0.04) (1.345,7.391)rectangle++(0.04,0.04) (1.437,7.825)rectangle++(0.04,0.04) (1.532,8.19)rectangle++(0.04,0.04) (1.628,8.496)rectangle++(0.04,0.04)
  (1.721,8.753)rectangle++(0.04,0.04) (1.811,8.967)rectangle++(0.04,0.04) (1.896,9.146)rectangle++(0.04,0.04) (2.043,4.156)rectangle++(0.04,0.04) (1.9,5.047)rectangle++(0.04,0.04)
  (1.815,5.81)rectangle++(0.04,0.04) (1.773,6.462)rectangle++(0.04,0.04) (1.393,3.611)rectangle++(0.04,0.04) (1.326,4.61)rectangle++(0.04,0.04) (1.309,5.461)rectangle++(0.04,0.04)
  (1.329,6.186)rectangle++(0.04,0.04) (1.375,6.801)rectangle++(0.04,0.04) (1.439,7.322)rectangle++(0.04,0.04) (-1.253,4.663)rectangle++(0.04,0.04) (-0.88,5.61)rectangle++(0.04,0.04)
  (-0.526,6.4)rectangle++(0.04,0.04) (-0.193,7.057)rectangle++(0.04,0.04) (0.116,7.602)rectangle++(0.04,0.04) (0.029,2.617)rectangle++(0.04,0.04) (0.127,3.819)rectangle++(0.04,0.04)
  (-0.744,3.725)rectangle++(0.04,0.04) (-0.486,4.793)rectangle++(0.04,0.04) (-0.223,5.689)rectangle++(0.04,0.04) (-1.204,4.029)rectangle++(0.04,0.04) (-0.864,5.069)rectangle++(0.04,0.04)
  (-0.534,5.939)rectangle++(0.04,0.04) (-0.012,3.61)rectangle++(0.04,0.04) (0.132,4.665)rectangle++(0.04,0.04) (0.297,5.556)rectangle++(0.04,0.04) (0.472,6.307)rectangle++(0.04,0.04)
  (0.652,6.938)rectangle++(0.04,0.04) (-0.11,3.737)rectangle++(0.04,0.04) (0.054,4.777)rectangle++(0.04,0.04) (0.235,5.655)rectangle++(0.04,0.04) (1.373,3.106)rectangle++(0.04,0.04)
  (1.289,4.181)rectangle++(0.04,0.04) (1.261,5.099)rectangle++(0.04,0.04) (1.273,5.88)rectangle++(0.04,0.04) (-0.976,4.347)rectangle++(0.04,0.04) (-0.658,5.33)rectangle++(0.04,0.04)
  (-0.348,6.153)rectangle++(0.04,0.04) (-0.052,6.84)rectangle++(0.04,0.04) (0.227,7.413)rectangle++(0.04,0.04) (0.487,7.888)rectangle++(0.04,0.04) (0.728,8.281)rectangle++(0.04,0.04)
  (0.947,8.606)rectangle++(0.04,0.04) (1.147,8.874)rectangle++(0.04,0.04) (2.063,3.934)rectangle++(0.04,0.04) (1.909,4.857)rectangle++(0.04,0.04) (1.814,5.649)rectangle++(0.04,0.04)
  (1.766,6.325)rectangle++(0.04,0.04) (1.752,6.902)rectangle++(0.04,0.04) (1.763,7.393)rectangle++(0.04,0.04) (1.792,7.809)rectangle++(0.04,0.04) (-0.623,3.097)rectangle++(0.04,0.04)
  (-0.725,3.442)rectangle++(0.04,0.04) (-0.817,3.491)rectangle++(0.04,0.04) (-0.557,4.596)rectangle++(0.04,0.04) (-0.292,5.525)rectangle++(0.04,0.04) (0.232,2.644)rectangle++(0.04,0.04)
  (-0.008,1.845)rectangle++(0.04,0.04) (0.065,3.164)rectangle++(0.04,0.04) (0.18,4.283)rectangle++(0.04,0.04) (0.322,5.23)rectangle++(0.04,0.04) (0.481,6.029)rectangle++(0.04,0.04)
  (-0.095,3.643)rectangle++(0.04,0.04) (0.063,4.696)rectangle++(0.04,0.04) (0.239,5.585)rectangle++(0.04,0.04) (0.424,6.334)rectangle++(0.04,0.04) (-1.225,4.278)rectangle++(0.04,0.04)
  (-0.872,5.281)rectangle++(0.04,0.04) (-0.532,6.12)rectangle++(0.04,0.04) (-0.21,6.82)rectangle++(0.04,0.04) (0.092,7.401)rectangle++(0.04,0.04) (-1.511,4.426)rectangle++(0.04,0.04)
  (-1.109,5.419)rectangle++(0.04,0.04) (-0.728,6.246)rectangle++(0.04,0.04) (-0.371,6.935)rectangle++(0.04,0.04) (-0.04,7.506)rectangle++(0.04,0.04) (0.264,7.978)rectangle++(0.04,0.04)
  (0.541,8.367)rectangle++(0.04,0.04) (0.792,8.686)rectangle++(0.04,0.04) (1.019,8.948)rectangle++(0.04,0.04) (1.222,9.161)rectangle++(0.04,0.04) (1.403,9.334)rectangle++(0.04,0.04)
  (1.563,9.474)rectangle++(0.04,0.04) (1.706,9.587)rectangle++(0.04,0.04) (1.831,9.677)rectangle++(0.04,0.04) (1.941,9.748)rectangle++(0.04,0.04) (2.038,9.804)rectangle++(0.04,0.04)
  (2.122,9.848)rectangle++(0.04,0.04) (2.196,9.883)rectangle++(0.04,0.04) (2.259,9.909)rectangle++(0.04,0.04) (-1.654,5.327)rectangle++(0.04,0.04) (-1.384,3.689)rectangle++(0.04,0.04)
  (-1.031,4.787)rectangle++(0.04,0.04) (-1.157,3.699)rectangle++(0.04,0.04) (0.041,3.372)rectangle++(0.04,0.04) (0.168,4.461)rectangle++(0.04,0.04) (0.319,5.381)rectangle++(0.04,0.04)
  (0.484,6.157)rectangle++(0.04,0.04) (0.655,6.811)rectangle++(0.04,0.04) (0.827,7.359)rectangle++(0.04,0.04) (1.725,3.569)rectangle++(0.04,0.04) (1.607,4.561)rectangle++(0.04,0.04)
  (1.546,5.409)rectangle++(0.04,0.04) (1.529,6.132)rectangle++(0.04,0.04) (1.339,3.492)rectangle++(0.04,0.04) (1.276,4.51)rectangle++(0.04,0.04) (1.263,5.379)rectangle++(0.04,0.04)
  (1.286,6.118)rectangle++(0.04,0.04) (1.336,6.745)rectangle++(0.04,0.04) (1.403,7.276)rectangle++(0.04,0.04) (1.481,7.725)rectangle++(0.04,0.04) (1.566,8.103)rectangle++(0.04,0.04)
  (-1.395,4.841)rectangle++(0.04,0.04) (-0.995,5.767)rectangle++(0.04,0.04) (-0.152,3.123)rectangle++(0.04,0.04) (-0.647,3.536)rectangle++(0.04,0.04) (-0.41,4.627)rectangle++(0.04,0.04)
  (-0.166,5.546)rectangle++(0.04,0.04) (0.079,6.317)rectangle++(0.04,0.04) (0.317,6.962)rectangle++(0.04,0.04) (0.546,7.501)rectangle++(0.04,0.04) (0.762,7.951)rectangle++(0.04,0.04)
  (0.963,8.324)rectangle++(0.04,0.04) (1.15,8.633)rectangle++(0.04,0.04) (1.32,8.888)rectangle++(0.04,0.04) (1.476,9.098)rectangle++(0.04,0.04) (1.616,9.271)rectangle++(0.04,0.04)
  (1.742,9.412)rectangle++(0.04,0.04) (1.855,9.526)rectangle++(0.04,0.04) (1.956,9.619)rectangle++(0.04,0.04) (2.045,9.695)rectangle++(0.04,0.04) (2.124,9.755)rectangle++(0.04,0.04)
  (2.193,9.803)rectangle++(0.04,0.04) (2.254,9.841)rectangle++(0.04,0.04) (2.307,9.871)rectangle++(0.04,0.04) (2.354,9.894)rectangle++(0.04,0.04) (2.394,9.912)rectangle++(0.04,0.04)
  (2.43,9.926)rectangle++(0.04,0.04) (2.46,9.936)rectangle++(0.04,0.04) (2.486,9.943)rectangle++(0.04,0.04) (-1.621,5.377)rectangle++(0.04,0.04) (-1.165,6.231)rectangle++(0.04,0.04)
  (-0.743,6.939)rectangle++(0.04,0.04) (-0.356,7.524)rectangle++(0.04,0.04) (-1.617,4.366)rectangle++(0.04,0.04) (-1.202,5.372)rectangle++(0.04,0.04) (-1.312,3.783)rectangle++(0.04,0.04)
  (-0.966,4.864)rectangle++(0.04,0.04) (1.318,2.681)rectangle++(0.04,0.04) (1.225,3.822)rectangle++(0.04,0.04) (1.192,4.796)rectangle++(0.04,0.04) (1.203,5.625)rectangle++(0.04,0.04)
  (1.245,6.33)rectangle++(0.04,0.04) (1.309,6.926)rectangle++(0.04,0.04) (1.388,7.431)rectangle++(0.04,0.04) (1.475,7.857)rectangle++(0.04,0.04) (1.566,8.216)rectangle++(0.04,0.04)
  (1.851,3.902)rectangle++(0.04,0.04) (1.727,4.839)rectangle++(0.04,0.04) (1.66,5.64)rectangle++(0.04,0.04) (1.202,3.427)rectangle++(0.04,0.04) (1.157,4.461)rectangle++(0.04,0.04)
  (-0.696,4.045)rectangle++(0.04,0.04) (-0.432,5.062)rectangle++(0.04,0.04) (-0.167,5.917)rectangle++(0.04,0.04) (0.093,6.632)rectangle++(0.04,0.04) (0.342,7.23)rectangle++(0.04,0.04)
  (0.578,7.728)rectangle++(0.04,0.04) (0.798,8.142)rectangle++(0.04,0.04) (1.002,8.485)rectangle++(0.04,0.04) (1.189,8.768)rectangle++(0.04,0.04) (1.359,9.002)rectangle++(0.04,0.04)
  (1.513,9.193)rectangle++(0.04,0.04) (1.651,9.35)rectangle++(0.04,0.04) (1.776,9.478)rectangle++(0.04,0.04) (1.886,9.581)rectangle++(0.04,0.04) (1.984,9.665)rectangle++(0.04,0.04)
  (2.071,9.732)rectangle++(0.04,0.04) (2.147,9.785)rectangle++(0.04,0.04) (2.215,9.828)rectangle++(0.04,0.04) (2.273,9.861)rectangle++(0.04,0.04) (2.325,9.887)rectangle++(0.04,0.04)
  (2.369,9.908)rectangle++(0.04,0.04) (2.408,9.923)rectangle++(0.04,0.04) (2.441,9.934)rectangle++(0.04,0.04) (2.47,9.943)rectangle++(0.04,0.04) (2.495,9.949)rectangle++(0.04,0.04)
  (2.171,4.447)rectangle++(0.04,0.04) (2.021,5.289)rectangle++(0.04,0.04) (1.927,6.011)rectangle++(0.04,0.04) (1.877,6.629)rectangle++(0.04,0.04) (1.858,7.155)rectangle++(0.04,0.04)
  (1.863,7.604)rectangle++(0.04,0.04) (1.665,3.851)rectangle++(0.04,0.04) (0.759,3.083)rectangle++(0.04,0.04) (0.766,4.186)rectangle++(0.04,0.04) (0.816,5.124)rectangle++(0.04,0.04)
  (0.897,5.919)rectangle++(0.04,0.04) (0.997,6.591)rectangle++(0.04,0.04) (1.516,3.459)rectangle++(0.04,0.04) (1.424,4.476)rectangle++(0.04,0.04) (1.388,5.344)rectangle++(0.04,0.04)
  (1.391,6.083)rectangle++(0.04,0.04) (1.423,6.711)rectangle++(0.04,0.04) (1.476,7.244)rectangle++(0.04,0.04) (1.542,7.694)rectangle++(0.04,0.04) (-1.313,4.756)rectangle++(0.04,0.04)
  (-0.928,5.691)rectangle++(0.04,0.04) (-0.563,6.471)rectangle++(0.04,0.04) (-0.222,7.119)rectangle++(0.04,0.04) (0.094,7.656)rectangle++(0.04,0.04) (0.384,8.1)rectangle++(0.04,0.04)
  (0.648,8.466)rectangle++(0.04,0.04) (0.887,8.766)rectangle++(0.04,0.04);
\fill[piece3]
  (-2.06,4.121)rectangle++(0.04,0.04) (-2.016,3.136)rectangle++(0.04,0.04) (-1.444,3.471)rectangle++(0.04,0.04) (-0.444,2.453)rectangle++(0.04,0.04) (-2.085,4.192)rectangle++(0.04,0.04)
  (-1.829,3.725)rectangle++(0.04,0.04) (-2.014,3.297)rectangle++(0.04,0.04) (-0.739,3.067)rectangle++(0.04,0.04) (-1.43,2.675)rectangle++(0.04,0.04) (-1.474,2.503)rectangle++(0.04,0.04)
  (-1.621,2.737)rectangle++(0.04,0.04) (-1.975,3.371)rectangle++(0.04,0.04) (-1.293,1.876)rectangle++(0.04,0.04) (-1.056,2.965)rectangle++(0.04,0.04) (-2.041,3.692)rectangle++(0.04,0.04)
  (-0.336,1.934)rectangle++(0.04,0.04) (-0.291,2.634)rectangle++(0.04,0.04) (-2.122,4.257)rectangle++(0.04,0.04) (-2.097,3.337)rectangle++(0.04,0.04) (-2.067,3.405)rectangle++(0.04,0.04)
  (-1.41,3.346)rectangle++(0.04,0.04) (-0.349,2.731)rectangle++(0.04,0.04) (-0.125,1.669)rectangle++(0.04,0.04) (-0.647,2.224)rectangle++(0.04,0.04) (-0.325,2.549)rectangle++(0.04,0.04)
  (-0.715,2.354)rectangle++(0.04,0.04) (-1.607,3.618)rectangle++(0.04,0.04) (-0.229,2.636)rectangle++(0.04,0.04) (-1.272,2.677)rectangle++(0.04,0.04) (-1.046,2.932)rectangle++(0.04,0.04)
  (-1.932,3.426)rectangle++(0.04,0.04) (-1.818,2.824)rectangle++(0.04,0.04) (-1.93,3.19)rectangle++(0.04,0.04) (-0.711,2.979)rectangle++(0.04,0.04) (-2.123,4.09)rectangle++(0.04,0.04)
  (-2.021,3.126)rectangle++(0.04,0.04) (-0.699,2.433)rectangle++(0.04,0.04) (-1.165,3.312)rectangle++(0.04,0.04) (-1.352,2.372)rectangle++(0.04,0.04) (-0.752,2.493)rectangle++(0.04,0.04)
  (-0.079,1.955)rectangle++(0.04,0.04) (-0.666,2.422)rectangle++(0.04,0.04) (-0.637,2.874)rectangle++(0.04,0.04) (-2.004,3.7)rectangle++(0.04,0.04) (-2.176,3.911)rectangle++(0.04,0.04)
  (-1.809,2.606)rectangle++(0.04,0.04) (-0.647,2.493)rectangle++(0.04,0.04) (-0.799,1.989)rectangle++(0.04,0.04) (-0.128,2.191)rectangle++(0.04,0.04) (-1.663,3.421)rectangle++(0.04,0.04)
  (-1.67,3.568)rectangle++(0.04,0.04) (-1.468,3.529)rectangle++(0.04,0.04) (-1.991,3.946)rectangle++(0.04,0.04) (-1.453,3.504)rectangle++(0.04,0.04) (-0.741,2.586)rectangle++(0.04,0.04)
  (-1.446,3.224)rectangle++(0.04,0.04) (-2.045,3.576)rectangle++(0.04,0.04) (-2.178,4.15)rectangle++(0.04,0.04) (-1.196,2.997)rectangle++(0.04,0.04) (-1.492,2.594)rectangle++(0.04,0.04)
  (-0.611,2.792)rectangle++(0.04,0.04) (-1.304,3.116)rectangle++(0.04,0.04) (-2.08,4.381)rectangle++(0.04,0.04) (-0.724,2.743)rectangle++(0.04,0.04) (-1.561,3.114)rectangle++(0.04,0.04)
  (-0.654,2.911)rectangle++(0.04,0.04) (-1.698,3.108)rectangle++(0.04,0.04) (-0.791,2.828)rectangle++(0.04,0.04) (-2.08,3.885)rectangle++(0.04,0.04) (-1.639,2.313)rectangle++(0.04,0.04)
  (-2.009,3.483)rectangle++(0.04,0.04) (-1.708,2.556)rectangle++(0.04,0.04) (-1.561,3.459)rectangle++(0.04,0.04) (-1.233,1.991)rectangle++(0.04,0.04) (-2.137,4.101)rectangle++(0.04,0.04)
  (-0.367,2.324)rectangle++(0.04,0.04) (-1.711,3.511)rectangle++(0.04,0.04) (-1.617,3.504)rectangle++(0.04,0.04) (-1.889,3.095)rectangle++(0.04,0.04) (-1.877,3.013)rectangle++(0.04,0.04)
  (-1.18,1.82)rectangle++(0.04,0.04) (-2.042,3.715)rectangle++(0.04,0.04) (-1.395,1.937)rectangle++(0.04,0.04) (-1.53,2.66)rectangle++(0.04,0.04) (-2.064,3.651)rectangle++(0.04,0.04)
  (-0.67,2.531)rectangle++(0.04,0.04) (-1.097,2.311)rectangle++(0.04,0.04) (-2.014,3.639)rectangle++(0.04,0.04) (-1.86,2.847)rectangle++(0.04,0.04) (-1.235,3.117)rectangle++(0.04,0.04)
  (-2.089,4.196)rectangle++(0.04,0.04) (-2.19,3.853)rectangle++(0.04,0.04) (-0.668,2.531)rectangle++(0.04,0.04) (-1.864,3.407)rectangle++(0.04,0.04) (-2.124,3.698)rectangle++(0.04,0.04)
  (-1.923,3.932)rectangle++(0.04,0.04) (-1.977,3.124)rectangle++(0.04,0.04) (-0.999,2.867)rectangle++(0.04,0.04) (-1.343,3.2)rectangle++(0.04,0.04) (-1.433,2.952)rectangle++(0.04,0.04)
  (-0.016,1.993)rectangle++(0.04,0.04) (-1.126,2.597)rectangle++(0.04,0.04) (-1.313,3.378)rectangle++(0.04,0.04) (-2.107,3.819)rectangle++(0.04,0.04) (-1.196,3.017)rectangle++(0.04,0.04)
  (-2.145,4.176)rectangle++(0.04,0.04) (-1.305,3.07)rectangle++(0.04,0.04) (-0.029,2.01)rectangle++(0.04,0.04) (-2.115,4.287)rectangle++(0.04,0.04) (-1.753,3.72)rectangle++(0.04,0.04)
  (-1.27,3.294)rectangle++(0.04,0.04) (-0.967,3.151)rectangle++(0.04,0.04) (-1.94,3.454)rectangle++(0.04,0.04) (-0.758,2.768)rectangle++(0.04,0.04) (-1.186,2.668)rectangle++(0.04,0.04)
  (-0.038,1.965)rectangle++(0.04,0.04) (-2.115,4.286)rectangle++(0.04,0.04) (-0.36,2.398)rectangle++(0.04,0.04) (-0.717,2.034)rectangle++(0.04,0.04) (-1.911,3.501)rectangle++(0.04,0.04)
  (-1.802,3.735)rectangle++(0.04,0.04) (-1.891,3.609)rectangle++(0.04,0.04) (-1.166,2.884)rectangle++(0.04,0.04) (-1.94,3.414)rectangle++(0.04,0.04) (-1.713,3.759)rectangle++(0.04,0.04)
  (-0.656,2.636)rectangle++(0.04,0.04) (-2.016,4.073)rectangle++(0.04,0.04) (-0.688,2.164)rectangle++(0.04,0.04) (-2.129,4.222)rectangle++(0.04,0.04) (-0.441,2.893)rectangle++(0.04,0.04)
  (-1.784,3.822)rectangle++(0.04,0.04) (-1.573,2.356)rectangle++(0.04,0.04) (-1.011,2.727)rectangle++(0.04,0.04) (-0.594,2.543)rectangle++(0.04,0.04) (-0.22,1.967)rectangle++(0.04,0.04)
  (-0.654,2.403)rectangle++(0.04,0.04) (-0.234,2.301)rectangle++(0.04,0.04) (-1.804,3.409)rectangle++(0.04,0.04) (-1.279,3.319)rectangle++(0.04,0.04) (-1.863,3.851)rectangle++(0.04,0.04)
  (-2.177,3.998)rectangle++(0.04,0.04) (-0.97,3.163)rectangle++(0.04,0.04) (-2.085,4.204)rectangle++(0.04,0.04) (-1.95,2.925)rectangle++(0.04,0.04) (-2.022,3.371)rectangle++(0.04,0.04)
  (-0.823,2.96)rectangle++(0.04,0.04) (-0.137,2.275)rectangle++(0.04,0.04) (-1.708,3.673)rectangle++(0.04,0.04) (-2.124,3.725)rectangle++(0.04,0.04) (-0.712,2.513)rectangle++(0.04,0.04)
  (-0.221,1.75)rectangle++(0.04,0.04) (-1.799,3.653)rectangle++(0.04,0.04) (-2.091,4.214)rectangle++(0.04,0.04) (-1.699,3.753)rectangle++(0.04,0.04) (-2.007,3.844)rectangle++(0.04,0.04)
  (-2.131,3.606)rectangle++(0.04,0.04) (-1.466,3.37)rectangle++(0.04,0.04) (-1.743,2.872)rectangle++(0.04,0.04) (-2.09,3.66)rectangle++(0.04,0.04) (-1.593,2.268)rectangle++(0.04,0.04)
  (-0.695,3.025)rectangle++(0.04,0.04) (-1.794,3.852)rectangle++(0.04,0.04) (-2.109,4.33)rectangle++(0.04,0.04) (-1.569,2.056)rectangle++(0.04,0.04) (-2.106,4.361)rectangle++(0.04,0.04)
  (-1.553,3.537)rectangle++(0.04,0.04) (-1.913,3.782)rectangle++(0.04,0.04) (-2.165,4.228)rectangle++(0.04,0.04) (-1.813,3.755)rectangle++(0.04,0.04) (-1.943,3.111)rectangle++(0.04,0.04)
  (-0.741,3.056)rectangle++(0.04,0.04) (-1.221,2.401)rectangle++(0.04,0.04) (-2.066,3.754)rectangle++(0.04,0.04) (-1.717,3.727)rectangle++(0.04,0.04) (-1.927,3.116)rectangle++(0.04,0.04)
  (-1.97,3.216)rectangle++(0.04,0.04) (-1.901,3.022)rectangle++(0.04,0.04) (-1.61,2.492)rectangle++(0.04,0.04) (-0.34,1.851)rectangle++(0.04,0.04) (-2.118,4.183)rectangle++(0.04,0.04)
  (-1.952,2.918)rectangle++(0.04,0.04) (-1.41,2.148)rectangle++(0.04,0.04) (-1.427,3.258)rectangle++(0.04,0.04) (-2.151,4.216)rectangle++(0.04,0.04) (-2.003,4.087)rectangle++(0.04,0.04)
  (-2.113,3.489)rectangle++(0.04,0.04) (-1.748,3.452)rectangle++(0.04,0.04) (-1.66,2.59)rectangle++(0.04,0.04) (-1.502,3.491)rectangle++(0.04,0.04) (-1.873,3.105)rectangle++(0.04,0.04)
  (-1.493,3.545)rectangle++(0.04,0.04) (-1.985,3.999)rectangle++(0.04,0.04) (-0.387,2.117)rectangle++(0.04,0.04) (-1.739,3.534)rectangle++(0.04,0.04) (-2.022,4.022)rectangle++(0.04,0.04)
  (-2.143,3.621)rectangle++(0.04,0.04) (-2.13,3.579)rectangle++(0.04,0.04) (-2.096,4.351)rectangle++(0.04,0.04) (-2.171,4.241)rectangle++(0.04,0.04) (-1.558,2.023)rectangle++(0.04,0.04)
  (-1.047,2.969)rectangle++(0.04,0.04) (-0.905,3.092)rectangle++(0.04,0.04) (-1.274,3.193)rectangle++(0.04,0.04) (-1.989,3.996)rectangle++(0.04,0.04) (-1.907,2.905)rectangle++(0.04,0.04)
  (-2.167,4.128)rectangle++(0.04,0.04) (-1.205,3.108)rectangle++(0.04,0.04) (-2.059,3.698)rectangle++(0.04,0.04) (-1.395,1.929)rectangle++(0.04,0.04) (-0.626,2.423)rectangle++(0.04,0.04)
  (-1.45,3.325)rectangle++(0.04,0.04) (-1.176,1.987)rectangle++(0.04,0.04) (-1.06,2.97)rectangle++(0.04,0.04) (-0.335,2.65)rectangle++(0.04,0.04) (-0.777,2.095)rectangle++(0.04,0.04)
  (-1.023,2.206)rectangle++(0.04,0.04) (-1.514,2.746)rectangle++(0.04,0.04) (-1.5,2.502)rectangle++(0.04,0.04) (-2.113,3.85)rectangle++(0.04,0.04) (-0.245,2.697)rectangle++(0.04,0.04)
  (-1.103,2.985)rectangle++(0.04,0.04) (-2.095,4.298)rectangle++(0.04,0.04) (-1.965,2.939)rectangle++(0.04,0.04) (-0.208,2.607)rectangle++(0.04,0.04) (-1.03,2.412)rectangle++(0.04,0.04)
  (-1.961,3.521)rectangle++(0.04,0.04) (-2.086,4.376)rectangle++(0.04,0.04) (-0.658,2.847)rectangle++(0.04,0.04) (-2.044,3.8)rectangle++(0.04,0.04) (-1.91,3.985)rectangle++(0.04,0.04)
  (-0.405,1.905)rectangle++(0.04,0.04) (-1.174,2.523)rectangle++(0.04,0.04) (-1.581,2.777)rectangle++(0.04,0.04) (-2.155,4.168)rectangle++(0.04,0.04) (-1.183,3.104)rectangle++(0.04,0.04)
  (-1.065,2)rectangle++(0.04,0.04) (-2.064,4.092)rectangle++(0.04,0.04) (-1.917,3.92)rectangle++(0.04,0.04) (-0.94,3.054)rectangle++(0.04,0.04) (-1.488,3.334)rectangle++(0.04,0.04)
  (-1.853,3.085)rectangle++(0.04,0.04) (-0.389,2.129)rectangle++(0.04,0.04) (-0.325,2.598)rectangle++(0.04,0.04) (-0.893,3.019)rectangle++(0.04,0.04) (-1.263,2.995)rectangle++(0.04,0.04)
  (-1.883,3.253)rectangle++(0.04,0.04) (-1.381,2.89)rectangle++(0.04,0.04) (-0.264,2)rectangle++(0.04,0.04) (-0.597,2.768)rectangle++(0.04,0.04) (-1.836,3.709)rectangle++(0.04,0.04)
  (-1.941,3.102)rectangle++(0.04,0.04) (-0.178,2.551)rectangle++(0.04,0.04) (-0.646,2.362)rectangle++(0.04,0.04) (-1.85,3.391)rectangle++(0.04,0.04) (-2.046,3.425)rectangle++(0.04,0.04)
  (-1.106,3.232)rectangle++(0.04,0.04) (-1.915,3.954)rectangle++(0.04,0.04) (-1.432,2.018)rectangle++(0.04,0.04) (-1.267,3.202)rectangle++(0.04,0.04) (-2.028,3.999)rectangle++(0.04,0.04)
  (-1.913,2.901)rectangle++(0.04,0.04) (-1.14,2.882)rectangle++(0.04,0.04) (-0.999,1.961)rectangle++(0.04,0.04) (-1.855,3.327)rectangle++(0.04,0.04) (-0.951,2.884)rectangle++(0.04,0.04)
  (-0.284,2.407)rectangle++(0.04,0.04) (-1.803,3.743)rectangle++(0.04,0.04) (-1.056,3.098)rectangle++(0.04,0.04) (-2.08,3.806)rectangle++(0.04,0.04) (-0.725,2.701)rectangle++(0.04,0.04)
  (-1.143,2.334)rectangle++(0.04,0.04) (-0.943,2.873)rectangle++(0.04,0.04) (-2.064,3.794)rectangle++(0.04,0.04) (-2.187,3.992)rectangle++(0.04,0.04) (-2.195,3.925)rectangle++(0.04,0.04)
  (-2.187,3.832)rectangle++(0.04,0.04) (-1.166,3.256)rectangle++(0.04,0.04) (-1.006,2.476)rectangle++(0.04,0.04) (-2.083,3.859)rectangle++(0.04,0.04) (-1.441,1.959)rectangle++(0.04,0.04)
  (-1.351,2.376)rectangle++(0.04,0.04) (-1.987,3.503)rectangle++(0.04,0.04) (-0.61,2.789)rectangle++(0.04,0.04) (-0.823,2.741)rectangle++(0.04,0.04) (-1.065,2.94)rectangle++(0.04,0.04)
  (-0.684,2.925)rectangle++(0.04,0.04) (-1.852,3.892)rectangle++(0.04,0.04) (-1.362,2.868)rectangle++(0.04,0.04) (-2.111,4.26)rectangle++(0.04,0.04) (-1.749,3.66)rectangle++(0.04,0.04)
  (-1.567,3.551)rectangle++(0.04,0.04) (-2.006,4.092)rectangle++(0.04,0.04) (-2.181,4.167)rectangle++(0.04,0.04) (-2.15,4.303)rectangle++(0.04,0.04) (-2.093,4.091)rectangle++(0.04,0.04)
  (-1.788,3.849)rectangle++(0.04,0.04) (-0.43,1.927)rectangle++(0.04,0.04) (-0.509,2.337)rectangle++(0.04,0.04) (-1.345,3.256)rectangle++(0.04,0.04) (-2.118,4.34)rectangle++(0.04,0.04)
  (-0.751,2.821)rectangle++(0.04,0.04) (-1.95,4.041)rectangle++(0.04,0.04) (-2.174,3.847)rectangle++(0.04,0.04) (-0.379,2.863)rectangle++(0.04,0.04) (-0.609,2.091)rectangle++(0.04,0.04)
  (-0.983,3.209)rectangle++(0.04,0.04) (-2.047,3.714)rectangle++(0.04,0.04) (-2.135,4.325)rectangle++(0.04,0.04) (-2.118,4.356)rectangle++(0.04,0.04) (-0.626,2.96)rectangle++(0.04,0.04)
  (-1.797,3.789)rectangle++(0.04,0.04) (-1.945,3.117)rectangle++(0.04,0.04) (-1.972,3.215)rectangle++(0.04,0.04) (-1.902,3.022)rectangle++(0.04,0.04) (-1.682,3.38)rectangle++(0.04,0.04)
  (-1.352,3.21)rectangle++(0.04,0.04) (-1.901,3.878)rectangle++(0.04,0.04) (-2.123,3.614)rectangle++(0.04,0.04) (-0.876,2.904)rectangle++(0.04,0.04) (-0.607,2.247)rectangle++(0.04,0.04)
  (-1.717,3.627)rectangle++(0.04,0.04) (-0.978,2.521)rectangle++(0.04,0.04) (-0.872,1.919)rectangle++(0.04,0.04) (-1.876,3.686)rectangle++(0.04,0.04) (-2.151,4.276)rectangle++(0.04,0.04)
  (-1.735,2.377)rectangle++(0.04,0.04) (0.018,1.919)rectangle++(0.04,0.04) (-0.144,1.685)rectangle++(0.04,0.04) (-1.122,2.852)rectangle++(0.04,0.04) (-1.995,3.557)rectangle++(0.04,0.04)
  (-1.397,2.933)rectangle++(0.04,0.04) (-1.295,3.338)rectangle++(0.04,0.04) (-0.397,2.834)rectangle++(0.04,0.04) (-1.706,3.526)rectangle++(0.04,0.04) (-2.162,3.983)rectangle++(0.04,0.04)
  (-1.819,2.62)rectangle++(0.04,0.04) (-0.88,2.35)rectangle++(0.04,0.04) (-1.513,3.592)rectangle++(0.04,0.04) (-1.674,3.607)rectangle++(0.04,0.04) (-2.105,4.346)rectangle++(0.04,0.04)
  (-0.382,2.047)rectangle++(0.04,0.04) (-1.618,3.628)rectangle++(0.04,0.04) (-2.161,4.192)rectangle++(0.04,0.04) (-2.098,4.317)rectangle++(0.04,0.04) (-1.662,3.252)rectangle++(0.04,0.04)
  (-2.163,4.093)rectangle++(0.04,0.04) (-2.086,4.309)rectangle++(0.04,0.04) (-1.559,2.057)rectangle++(0.04,0.04) (-1.788,3.439)rectangle++(0.04,0.04) (-1.273,1.934)rectangle++(0.04,0.04)
  (-2.138,4.245)rectangle++(0.04,0.04) (-2.001,4.045)rectangle++(0.04,0.04) (-2.11,3.485)rectangle++(0.04,0.04) (-1.852,2.812)rectangle++(0.04,0.04) (-1.908,3.987)rectangle++(0.04,0.04)
  (-1.213,2.747)rectangle++(0.04,0.04) (-1.996,3.543)rectangle++(0.04,0.04) (-1.844,2.834)rectangle++(0.04,0.04) (-2.165,4.076)rectangle++(0.04,0.04) (-2.085,3.313)rectangle++(0.04,0.04)
  (-2.107,3.553)rectangle++(0.04,0.04) (-0.723,2.682)rectangle++(0.04,0.04) (-1.778,3.745)rectangle++(0.04,0.04) (-1.982,3.999)rectangle++(0.04,0.04) (-0.355,2.838)rectangle++(0.04,0.04)
  (-0.141,1.804)rectangle++(0.04,0.04) (-0.93,2.475)rectangle++(0.04,0.04) (-2.095,4.334)rectangle++(0.04,0.04) (-1.638,3.531)rectangle++(0.04,0.04) (-1.79,2.876)rectangle++(0.04,0.04)
  (-1.567,2.483)rectangle++(0.04,0.04) (-2.078,4.141)rectangle++(0.04,0.04) (-1.379,3.21)rectangle++(0.04,0.04) (-1.558,2.604)rectangle++(0.04,0.04) (-2.074,4.187)rectangle++(0.04,0.04)
  (-2.157,3.634)rectangle++(0.04,0.04) (-1.268,2.757)rectangle++(0.04,0.04) (-0.203,2.308)rectangle++(0.04,0.04) (-0.896,2.885)rectangle++(0.04,0.04) (-2.134,4.133)rectangle++(0.04,0.04)
  (-1.176,3.161)rectangle++(0.04,0.04) (-1.747,3.518)rectangle++(0.04,0.04) (-2.045,3.445)rectangle++(0.04,0.04) (-1.326,1.876)rectangle++(0.04,0.04) (-1.076,2.06)rectangle++(0.04,0.04)
  (-0.444,2.201)rectangle++(0.04,0.04) (-1.982,4.084)rectangle++(0.04,0.04) (-0.38,1.885)rectangle++(0.04,0.04) (-1.889,3.807)rectangle++(0.04,0.04) (-2.162,3.832)rectangle++(0.04,0.04)
  (-1.207,3.093)rectangle++(0.04,0.04) (-1.979,4.05)rectangle++(0.04,0.04) (-1.77,3.803)rectangle++(0.04,0.04) (-1.364,2.019)rectangle++(0.04,0.04) (-1.537,2.676)rectangle++(0.04,0.04)
  (-2.021,3.521)rectangle++(0.04,0.04) (-1.943,3.059)rectangle++(0.04,0.04) (-2.165,3.956)rectangle++(0.04,0.04) (-0.748,2.797)rectangle++(0.04,0.04) (-1.384,2.627)rectangle++(0.04,0.04)
  (-1.246,2.197)rectangle++(0.04,0.04) (-0.842,1.786)rectangle++(0.04,0.04) (-2.111,4.318)rectangle++(0.04,0.04) (-2.2,4.018)rectangle++(0.04,0.04) (-0.321,2.397)rectangle++(0.04,0.04)
  (-0.348,2.779)rectangle++(0.04,0.04) (-1.456,3.339)rectangle++(0.04,0.04) (-1.244,2.731)rectangle++(0.04,0.04) (-2.104,3.869)rectangle++(0.04,0.04) (-1.913,3.991)rectangle++(0.04,0.04)
  (-1.983,3.133)rectangle++(0.04,0.04) (-1.894,3.01)rectangle++(0.04,0.04) (-0.331,2.801)rectangle++(0.04,0.04) (-0.597,2.359)rectangle++(0.04,0.04) (-2.098,4.008)rectangle++(0.04,0.04)
  (-1.658,3.579)rectangle++(0.04,0.04) (-0.394,2.124)rectangle++(0.04,0.04) (-0.652,1.966)rectangle++(0.04,0.04) (-0.265,2.16)rectangle++(0.04,0.04) (-1.796,3.393)rectangle++(0.04,0.04)
  (-1.843,3.709)rectangle++(0.04,0.04) (-0.641,2.775)rectangle++(0.04,0.04) (-1.792,3.275)rectangle++(0.04,0.04) (-2.03,4.168)rectangle++(0.04,0.04) (-0.681,2.707)rectangle++(0.04,0.04)
  (-1.743,3.537)rectangle++(0.04,0.04) (-0.643,2.964)rectangle++(0.04,0.04) (-2.096,3.971)rectangle++(0.04,0.04) (-1.148,3)rectangle++(0.04,0.04) (-1.126,2.861)rectangle++(0.04,0.04)
  (-0.731,2.404)rectangle++(0.04,0.04) (-1.508,3.587)rectangle++(0.04,0.04) (-1.473,2.362)rectangle++(0.04,0.04) (-1.321,2.818)rectangle++(0.04,0.04) (-1.275,2.246)rectangle++(0.04,0.04)
  (-1.001,3.22)rectangle++(0.04,0.04) (-1.514,3.114)rectangle++(0.04,0.04) (-2.087,4.38)rectangle++(0.04,0.04) (-2.2,3.95)rectangle++(0.04,0.04) (-0.362,2.201)rectangle++(0.04,0.04)
  (-0.572,2.057)rectangle++(0.04,0.04) (-2.029,3.897)rectangle++(0.04,0.04) (-1.899,2.885)rectangle++(0.04,0.04) (-2.128,4.328)rectangle++(0.04,0.04) (-1.422,2.939)rectangle++(0.04,0.04)
  (-1.193,2.676)rectangle++(0.04,0.04) (-1.845,3.22)rectangle++(0.04,0.04) (-0.141,1.682)rectangle++(0.04,0.04) (-1.839,3.497)rectangle++(0.04,0.04) (-1.513,2.282)rectangle++(0.04,0.04)
  (-2.138,3.995)rectangle++(0.04,0.04) (-1.333,3.342)rectangle++(0.04,0.04) (-0.863,2.432)rectangle++(0.04,0.04) (-1.523,2.8)rectangle++(0.04,0.04) (-2.027,3.533)rectangle++(0.04,0.04)
  (-1.756,3.642)rectangle++(0.04,0.04) (-1.694,2.622)rectangle++(0.04,0.04) (-1.66,3.408)rectangle++(0.04,0.04) (-1.982,4.089)rectangle++(0.04,0.04) (-2.081,4.26)rectangle++(0.04,0.04)
  (-2.178,4.208)rectangle++(0.04,0.04) (-2.194,4.126)rectangle++(0.04,0.04) (-2.041,4.057)rectangle++(0.04,0.04) (-2.187,4.121)rectangle++(0.04,0.04) (-0.055,2.363)rectangle++(0.04,0.04)
  (-0.85,2.721)rectangle++(0.04,0.04) (-2.014,3.68)rectangle++(0.04,0.04) (-1.466,3.531)rectangle++(0.04,0.04) (-1.453,2.359)rectangle++(0.04,0.04) (-0.945,3.112)rectangle++(0.04,0.04)
  (-1.161,3.308)rectangle++(0.04,0.04) (-0.896,2.604)rectangle++(0.04,0.04) (-1.375,2.571)rectangle++(0.04,0.04) (-1.446,2.493)rectangle++(0.04,0.04) (-1.445,2.469)rectangle++(0.04,0.04)
  (-1.122,3.237)rectangle++(0.04,0.04) (-1.076,3.082)rectangle++(0.04,0.04) (-0.288,2.705)rectangle++(0.04,0.04) (-0.303,2.667)rectangle++(0.04,0.04) (-1.952,3.726)rectangle++(0.04,0.04)
  (-2.165,4.226)rectangle++(0.04,0.04) (-2.167,3.633)rectangle++(0.04,0.04) (-2.184,3.904)rectangle++(0.04,0.04) (-0.93,3.165)rectangle++(0.04,0.04) (-1.437,3.301)rectangle++(0.04,0.04)
  (-0.613,2.902)rectangle++(0.04,0.04) (-1.165,2.395)rectangle++(0.04,0.04) (-1.972,3.513)rectangle++(0.04,0.04) (-1.762,3.399)rectangle++(0.04,0.04) (-1.79,2.84)rectangle++(0.04,0.04)
  (-1.279,3.305)rectangle++(0.04,0.04) (-1.95,3.854)rectangle++(0.04,0.04) (-1.364,3.134)rectangle++(0.04,0.04) (-1.695,2.848)rectangle++(0.04,0.04) (-1.984,4.081)rectangle++(0.04,0.04)
  (-1.664,2.37)rectangle++(0.04,0.04) (-0.425,2.356)rectangle++(0.04,0.04) (-1.744,3.546)rectangle++(0.04,0.04) (-1.441,3.455)rectangle++(0.04,0.04) (-1.046,2.516)rectangle++(0.04,0.04)
  (-1.456,3.461)rectangle++(0.04,0.04) (-1.212,2.015)rectangle++(0.04,0.04) (-0.902,2.372)rectangle++(0.04,0.04) (-1.876,3.376)rectangle++(0.04,0.04) (-1.908,3.053)rectangle++(0.04,0.04)
  (-1.961,3.21)rectangle++(0.04,0.04) (-2.16,3.883)rectangle++(0.04,0.04) (-1.695,3.744)rectangle++(0.04,0.04) (-1.926,3.12)rectangle++(0.04,0.04) (-1.757,3.391)rectangle++(0.04,0.04)
  (-1.889,3.709)rectangle++(0.04,0.04) (-0.739,3.067)rectangle++(0.04,0.04) (-2.084,4.058)rectangle++(0.04,0.04) (-0.275,2.743)rectangle++(0.04,0.04) (-0.826,2.549)rectangle++(0.04,0.04)
  (-1.678,3.042)rectangle++(0.04,0.04) (-0.078,2.275)rectangle++(0.04,0.04) (-0.935,2.371)rectangle++(0.04,0.04) (-0.304,2.721)rectangle++(0.04,0.04) (-1.488,3.519)rectangle++(0.04,0.04)
  (-2.089,4.379)rectangle++(0.04,0.04) (-0.222,2.673)rectangle++(0.04,0.04) (-1.42,2.937)rectangle++(0.04,0.04) (-2.026,3.551)rectangle++(0.04,0.04) (-1.757,3.706)rectangle++(0.04,0.04)
  (-1.336,2)rectangle++(0.04,0.04) (-0.809,1.722)rectangle++(0.04,0.04) (-1.819,3.314)rectangle++(0.04,0.04) (-1.786,2.821)rectangle++(0.04,0.04) (-1.582,3.601)rectangle++(0.04,0.04)
  (-1.344,3.432)rectangle++(0.04,0.04) (-1.846,3.112)rectangle++(0.04,0.04) (-2.025,3.396)rectangle++(0.04,0.04) (-1.309,1.87)rectangle++(0.04,0.04) (-1.499,3.266)rectangle++(0.04,0.04)
  (-1.143,3.286)rectangle++(0.04,0.04) (-1.967,3.457)rectangle++(0.04,0.04) (-1.066,3.133)rectangle++(0.04,0.04) (-1.049,2.033)rectangle++(0.04,0.04) (-1.494,2.716)rectangle++(0.04,0.04)
  (-2.063,3.657)rectangle++(0.04,0.04) (-2.035,3.27)rectangle++(0.04,0.04) (-1.941,4.016)rectangle++(0.04,0.04) (-1.989,3.134)rectangle++(0.04,0.04) (-2.174,4.086)rectangle++(0.04,0.04)
  (-1.331,3.167)rectangle++(0.04,0.04) (-1.906,3.268)rectangle++(0.04,0.04) (-2.114,4.22)rectangle++(0.04,0.04) (-1.936,3.982)rectangle++(0.04,0.04) (-1.897,2.909)rectangle++(0.04,0.04)
  (-1.157,1.793)rectangle++(0.04,0.04) (-1.502,3.57)rectangle++(0.04,0.04) (-1.468,2.36)rectangle++(0.04,0.04) (-1.943,3.999)rectangle++(0.04,0.04) (-2.108,4.364)rectangle++(0.04,0.04)
  (-2.17,3.648)rectangle++(0.04,0.04) (-1.763,2.561)rectangle++(0.04,0.04) (-2.022,3.881)rectangle++(0.04,0.04) (-2.174,3.833)rectangle++(0.04,0.04) (-2.168,3.72)rectangle++(0.04,0.04)
  (-0.435,1.931)rectangle++(0.04,0.04) (-0.926,2.237)rectangle++(0.04,0.04) (-1.08,2.206)rectangle++(0.04,0.04) (-1.672,2.971)rectangle++(0.04,0.04) (-1.411,3.227)rectangle++(0.04,0.04)
  (-1.651,3.427)rectangle++(0.04,0.04) (-2.075,3.584)rectangle++(0.04,0.04) (-2.102,4.37)rectangle++(0.04,0.04) (-1.747,2.405)rectangle++(0.04,0.04) (-1.884,3.145)rectangle++(0.04,0.04)
  (-2.169,4.136)rectangle++(0.04,0.04) (-2.032,3.128)rectangle++(0.04,0.04) (-1.889,3.607)rectangle++(0.04,0.04) (-0.414,2.458)rectangle++(0.04,0.04) (-0.925,2.952)rectangle++(0.04,0.04)
  (-1.737,3.087)rectangle++(0.04,0.04) (-0.677,2.697)rectangle++(0.04,0.04) (-0.422,2.885)rectangle++(0.04,0.04) (-2.113,4.211)rectangle++(0.04,0.04) (-2.192,3.854)rectangle++(0.04,0.04)
  (-1.653,2.287)rectangle++(0.04,0.04) (-0.659,2.916)rectangle++(0.04,0.04) (-0.091,2.006)rectangle++(0.04,0.04) (-1.596,3.493)rectangle++(0.04,0.04) (-1.468,2.327)rectangle++(0.04,0.04)
  (-1.467,3.516)rectangle++(0.04,0.04) (-2.059,4.206)rectangle++(0.04,0.04) (-2.093,4.342)rectangle++(0.04,0.04) (-1.318,3.142)rectangle++(0.04,0.04) (-1.903,3.266)rectangle++(0.04,0.04)
  (-1.653,2.535)rectangle++(0.04,0.04) (-1.011,1.766)rectangle++(0.04,0.04) (-0.801,2.673)rectangle++(0.04,0.04) (-1.147,3.277)rectangle++(0.04,0.04) (-2.022,3.595)rectangle++(0.04,0.04)
  (-1.567,2.267)rectangle++(0.04,0.04) (-0.924,1.727)rectangle++(0.04,0.04) (-1.602,3.643)rectangle++(0.04,0.04) (-0.962,2.842)rectangle++(0.04,0.04) (-0.953,1.993)rectangle++(0.04,0.04)
  (-2.09,4.341)rectangle++(0.04,0.04) (-0.803,2.949)rectangle++(0.04,0.04) (-0.788,2.554)rectangle++(0.04,0.04) (-2.132,4.2)rectangle++(0.04,0.04) (-0.367,2.788)rectangle++(0.04,0.04)
  (-0.819,2.118)rectangle++(0.04,0.04) (-1.555,3.185)rectangle++(0.04,0.04) (-0.708,2.975)rectangle++(0.04,0.04) (-0.592,2.807)rectangle++(0.04,0.04) (-1.463,3.467)rectangle++(0.04,0.04)
  (-1.485,3.48)rectangle++(0.04,0.04) (-2.158,4.149)rectangle++(0.04,0.04) (-2.162,3.626)rectangle++(0.04,0.04) (-1.348,2.851)rectangle++(0.04,0.04) (-1.944,3.984)rectangle++(0.04,0.04)
  (-1.446,2.019)rectangle++(0.04,0.04) (-0.835,1.701)rectangle++(0.04,0.04) (-1.799,3.674)rectangle++(0.04,0.04) (-1.665,3.487)rectangle++(0.04,0.04) (-1.893,3.088)rectangle++(0.04,0.04)
  (-1.641,3.463)rectangle++(0.04,0.04) (0.001,2.048)rectangle++(0.04,0.04) (-0.873,2.342)rectangle++(0.04,0.04) (-1.943,3.522)rectangle++(0.04,0.04) (-1.933,3.066)rectangle++(0.04,0.04)
  (-1.438,2.18)rectangle++(0.04,0.04) (-1.469,3.362)rectangle++(0.04,0.04) (-1.418,2.325)rectangle++(0.04,0.04) (-0.158,2.427)rectangle++(0.04,0.04) (-1.705,3.27)rectangle++(0.04,0.04)
  (-2.147,3.892)rectangle++(0.04,0.04) (-2.085,4.382)rectangle++(0.04,0.04) (-0.417,1.916)rectangle++(0.04,0.04) (-2.109,4.103)rectangle++(0.04,0.04) (-0.932,3.168)rectangle++(0.04,0.04)
  (-1.861,3.299)rectangle++(0.04,0.04) (-1.652,2.547)rectangle++(0.04,0.04) (-1.014,1.769)rectangle++(0.04,0.04) (-1.666,3.396)rectangle++(0.04,0.04) (-2.148,3.9)rectangle++(0.04,0.04)
  (-1.553,3.526)rectangle++(0.04,0.04) (-2.159,4.189)rectangle++(0.04,0.04) (-1.235,3.225)rectangle++(0.04,0.04) (-1.706,3.596)rectangle++(0.04,0.04) (-2.164,3.987)rectangle++(0.04,0.04)
  (-1.935,3.897)rectangle++(0.04,0.04) (-2.152,3.734)rectangle++(0.04,0.04) (-1.147,3.158)rectangle++(0.04,0.04) (-1.206,2.689)rectangle++(0.04,0.04) (-1.227,2.242)rectangle++(0.04,0.04)
  (-0.348,2.477)rectangle++(0.04,0.04) (-2.048,4.177)rectangle++(0.04,0.04) (-0.251,2.709)rectangle++(0.04,0.04) (-1.741,3.297)rectangle++(0.04,0.04) (-1.883,3.828)rectangle++(0.04,0.04)
  (-2.178,4.113)rectangle++(0.04,0.04) (-1.706,2.34)rectangle++(0.04,0.04) (-1.842,3.781)rectangle++(0.04,0.04) (-2.032,4.062)rectangle++(0.04,0.04) (-1.753,3.688)rectangle++(0.04,0.04)
  (-1.829,2.886)rectangle++(0.04,0.04) (-1.571,3.561)rectangle++(0.04,0.04) (-1.261,2.011)rectangle++(0.04,0.04) (-1.142,3.084)rectangle++(0.04,0.04) (-2.141,4.241)rectangle++(0.04,0.04)
  (-1.726,2.372)rectangle++(0.04,0.04) (-0.983,1.714)rectangle++(0.04,0.04) (-0.663,1.74)rectangle++(0.04,0.04) (-0.349,2.405)rectangle++(0.04,0.04) (-1.96,3.674)rectangle++(0.04,0.04)
  (-2.025,3.28)rectangle++(0.04,0.04) (-1.984,3.93)rectangle++(0.04,0.04) (-1.985,3.118)rectangle++(0.04,0.04) (-2.143,3.782)rectangle++(0.04,0.04) (-1.204,2.866)rectangle++(0.04,0.04)
  (-2.14,4.082)rectangle++(0.04,0.04) (-1.275,3.335)rectangle++(0.04,0.04) (-1.71,2.89)rectangle++(0.04,0.04) (-0.392,2.269)rectangle++(0.04,0.04) (-1.585,3.146)rectangle++(0.04,0.04)
  (-1.156,1.917)rectangle++(0.04,0.04) (-2.116,3.983)rectangle++(0.04,0.04) (-1.48,1.978)rectangle++(0.04,0.04) (-0.831,2.906)rectangle++(0.04,0.04) (-2.073,4.235)rectangle++(0.04,0.04)
  (-0.752,3.071)rectangle++(0.04,0.04) (-1.398,2.98)rectangle++(0.04,0.04) (-1.801,3.045)rectangle++(0.04,0.04) (-1.009,2.671)rectangle++(0.04,0.04) (-0.917,1.945)rectangle++(0.04,0.04)
  (-0.149,2.304)rectangle++(0.04,0.04) (-0.863,2.863)rectangle++(0.04,0.04) (-1.719,3.081)rectangle++(0.04,0.04) (-0.715,3.009)rectangle++(0.04,0.04) (-0.574,2.755)rectangle++(0.04,0.04)
  (-1.946,4.024)rectangle++(0.04,0.04) (-1.905,2.914)rectangle++(0.04,0.04) (-1.592,2.473)rectangle++(0.04,0.04) (-0.94,2.41)rectangle++(0.04,0.04) (-0.629,2.11)rectangle++(0.04,0.04)
  (-1.617,3.584)rectangle++(0.04,0.04) (-0.724,2.198)rectangle++(0.04,0.04) (-0.737,1.906)rectangle++(0.04,0.04) (-2.044,4.105)rectangle++(0.04,0.04) (-0.991,3.212)rectangle++(0.04,0.04)
  (-1.186,3.322)rectangle++(0.04,0.04) (-2.117,4.34)rectangle++(0.04,0.04) (-1.97,3.984)rectangle++(0.04,0.04) (-1.616,3.672)rectangle++(0.04,0.04) (-1.507,3.364)rectangle++(0.04,0.04)
  (-2.144,4.27)rectangle++(0.04,0.04) (-2.034,3.99)rectangle++(0.04,0.04) (-2.189,4.065)rectangle++(0.04,0.04) (-1.509,3.581)rectangle++(0.04,0.04) (-2.101,4.328)rectangle++(0.04,0.04)
  (-0.601,2.472)rectangle++(0.04,0.04) (-1.863,3.666)rectangle++(0.04,0.04) (-1.86,3.645)rectangle++(0.04,0.04) (-1.549,2.307)rectangle++(0.04,0.04) (-1.208,2.103)rectangle++(0.04,0.04)
  (-1.958,3.485)rectangle++(0.04,0.04) (-2.052,4.059)rectangle++(0.04,0.04) (-1.487,2.01)rectangle++(0.04,0.04) (-1.882,3.767)rectangle++(0.04,0.04) (-2.043,4.191)rectangle++(0.04,0.04)
  (-0.698,2.631)rectangle++(0.04,0.04) (-1.201,2.888)rectangle++(0.04,0.04) (-1.268,2.282)rectangle++(0.04,0.04) (-1.561,2.728)rectangle++(0.04,0.04) (-2.027,4.153)rectangle++(0.04,0.04)
  (-2.076,3.332)rectangle++(0.04,0.04) (-1.303,1.845)rectangle++(0.04,0.04) (-0.358,2.632)rectangle++(0.04,0.04) (-0.777,2.086)rectangle++(0.04,0.04) (-0.159,2.315)rectangle++(0.04,0.04)
  (-2.116,4.195)rectangle++(0.04,0.04) (-0.406,2.842)rectangle++(0.04,0.04) (-0.841,2.121)rectangle++(0.04,0.04) (-0.363,2.078)rectangle++(0.04,0.04) (-2.01,4.1)rectangle++(0.04,0.04)
  (-0.756,2.808)rectangle++(0.04,0.04) (-1.387,2.628)rectangle++(0.04,0.04) (-0.982,1.848)rectangle++(0.04,0.04) (-1.184,3.316)rectangle++(0.04,0.04) (-0.618,2.347)rectangle++(0.04,0.04)
  (-1.628,3.035)rectangle++(0.04,0.04) (-1.998,3.402)rectangle++(0.04,0.04) (-1.201,3.357)rectangle++(0.04,0.04) (-1.703,2.902)rectangle++(0.04,0.04) (-1.831,3.005)rectangle++(0.04,0.04)
  (-2.075,4.258)rectangle++(0.04,0.04) (-1.718,2.388)rectangle++(0.04,0.04) (-2.154,4.059)rectangle++(0.04,0.04) (-1.902,3.63)rectangle++(0.04,0.04) (-1.345,1.95)rectangle++(0.04,0.04)
  (-1.429,3.188)rectangle++(0.04,0.04) (-1.562,2.592)rectangle++(0.04,0.04) (-1.958,3.359)rectangle++(0.04,0.04) (-0.105,2.404)rectangle++(0.04,0.04) (-0.968,2.392)rectangle++(0.04,0.04)
  (-2.01,3.952)rectangle++(0.04,0.04) (-1.789,2.638)rectangle++(0.04,0.04) (-0.737,2.209)rectangle++(0.04,0.04) (-1.473,2.781)rectangle++(0.04,0.04) (-1.788,3.011)rectangle++(0.04,0.04)
  (-1.908,3.775)rectangle++(0.04,0.04) (-1.384,1.981)rectangle++(0.04,0.04) (-1.784,3.792)rectangle++(0.04,0.04) (-1.006,2.911)rectangle++(0.04,0.04) (-1.238,2.323)rectangle++(0.04,0.04)
  (-1.103,2.769)rectangle++(0.04,0.04) (-0.962,1.944)rectangle++(0.04,0.04) (-0.255,2.618)rectangle++(0.04,0.04) (-0.592,2.451)rectangle++(0.04,0.04) (-2.122,4.293)rectangle++(0.04,0.04)
  (-2.185,3.759)rectangle++(0.04,0.04) (-1.871,3.576)rectangle++(0.04,0.04) (-1.692,3.35)rectangle++(0.04,0.04) (-1.058,3.084)rectangle++(0.04,0.04) (-1.644,2.877)rectangle++(0.04,0.04)
  (-2.007,3.894)rectangle++(0.04,0.04) (-2.084,4.294)rectangle++(0.04,0.04) (-2.164,3.643)rectangle++(0.04,0.04) (-1.342,3.209)rectangle++(0.04,0.04) (-2.035,3.579)rectangle++(0.04,0.04)
  (-0.244,2.686)rectangle++(0.04,0.04) (-1.985,4.042)rectangle++(0.04,0.04) (-1.913,2.912)rectangle++(0.04,0.04) (-2.089,4.353)rectangle++(0.04,0.04) (-2.082,4.27)rectangle++(0.04,0.04)
  (-1.166,3.199)rectangle++(0.04,0.04) (-1.328,2.35)rectangle++(0.04,0.04) (-1.034,2.83)rectangle++(0.04,0.04) (-2.01,3.913)rectangle++(0.04,0.04) (-2.159,3.733)rectangle++(0.04,0.04)
  (-0.847,2.522)rectangle++(0.04,0.04) (-0.441,2.9)rectangle++(0.04,0.04) (-0.892,2.733)rectangle++(0.04,0.04) (-0.91,1.985)rectangle++(0.04,0.04) (-1.467,2.726)rectangle++(0.04,0.04)
  (-2.117,3.864)rectangle++(0.04,0.04) (-2.06,3.291)rectangle++(0.04,0.04) (-1.818,2.788)rectangle++(0.04,0.04) (-1.11,1.784)rectangle++(0.04,0.04) (-0.331,2.772)rectangle++(0.04,0.04)
  (-0.52,2.411)rectangle++(0.04,0.04) (-1.342,2.951)rectangle++(0.04,0.04) (-1.578,3.563)rectangle++(0.04,0.04) (-0.523,2.719)rectangle++(0.04,0.04) (-1.515,2.886)rectangle++(0.04,0.04)
  (-1.922,3.928)rectangle++(0.04,0.04) (-0.671,3.008)rectangle++(0.04,0.04) (-2.01,4.125)rectangle++(0.04,0.04) (-2.058,4.032)rectangle++(0.04,0.04) (-1.667,2.346)rectangle++(0.04,0.04)
  (-1.869,3.144)rectangle++(0.04,0.04) (-0.947,3.179)rectangle++(0.04,0.04) (-2.084,4.378)rectangle++(0.04,0.04) (-2.169,3.651)rectangle++(0.04,0.04) (-0.958,3.132)rectangle++(0.04,0.04)
  (-1.636,2.897)rectangle++(0.04,0.04) (-0.141,1.683)rectangle++(0.04,0.04) (-1.663,3.397)rectangle++(0.04,0.04) (-0.65,2.745)rectangle++(0.04,0.04) (-0.931,2.479)rectangle++(0.04,0.04)
  (-2.116,4.333)rectangle++(0.04,0.04) (-1.25,3.274)rectangle++(0.04,0.04) (-0.828,2.903)rectangle++(0.04,0.04) (-0.942,2.037)rectangle++(0.04,0.04) (-1.931,3.491)rectangle++(0.04,0.04)
  (-1.696,2.569)rectangle++(0.04,0.04) (-1.029,1.764)rectangle++(0.04,0.04) (-1.493,3.567)rectangle++(0.04,0.04) (-1.764,3.501)rectangle++(0.04,0.04) (-2.164,4.188)rectangle++(0.04,0.04)
  (-1.941,4.024)rectangle++(0.04,0.04) (-2.183,4.12)rectangle++(0.04,0.04) (-0.874,2.344)rectangle++(0.04,0.04) (-1.98,4.06)rectangle++(0.04,0.04) (-1.503,3.581)rectangle++(0.04,0.04)
  (-1.947,3.935)rectangle++(0.04,0.04) (-2.096,3.476)rectangle++(0.04,0.04) (-1.947,3.045)rectangle++(0.04,0.04) (-2.138,3.778)rectangle++(0.04,0.04) (-1.431,1.928)rectangle++(0.04,0.04)
  (-0.358,2.706)rectangle++(0.04,0.04) (-1.636,3.109)rectangle++(0.04,0.04) (-1.846,3.04)rectangle++(0.04,0.04) (-0.983,2.453)rectangle++(0.04,0.04) (-2.051,4.085)rectangle++(0.04,0.04)
  (-1.819,2.655)rectangle++(0.04,0.04) (-0.904,2.689)rectangle++(0.04,0.04) (-2.126,4.297)rectangle++(0.04,0.04) (-0.319,2.463)rectangle++(0.04,0.04) (-2.117,4.308)rectangle++(0.04,0.04)
  (-2.045,3.156)rectangle++(0.04,0.04) (-0.733,3.048)rectangle++(0.04,0.04) (-1.909,3.457)rectangle++(0.04,0.04) (-1.687,2.567)rectangle++(0.04,0.04) (-1.496,2.445)rectangle++(0.04,0.04)
  (-1.752,2.963)rectangle++(0.04,0.04) (-1.576,2.504)rectangle++(0.04,0.04) (-1.636,2.725)rectangle++(0.04,0.04) (-1.312,2.169)rectangle++(0.04,0.04) (-1.481,3.351)rectangle++(0.04,0.04)
  (-0.924,3.141)rectangle++(0.04,0.04) (-1.858,3.296)rectangle++(0.04,0.04) (-0.976,2.797)rectangle++(0.04,0.04) (-1.794,3.842)rectangle++(0.04,0.04) (-1.894,3.938)rectangle++(0.04,0.04)
  (-1.424,2.02)rectangle++(0.04,0.04) (-1.811,3.124)rectangle++(0.04,0.04) (-1.613,2.524)rectangle++(0.04,0.04) (-1.265,2.134)rectangle++(0.04,0.04) (-0.062,2.063)rectangle++(0.04,0.04)
  (-2.087,4.06)rectangle++(0.04,0.04) (-0.637,2.402)rectangle++(0.04,0.04) (-1.063,2.292)rectangle++(0.04,0.04) (-0.356,2.319)rectangle++(0.04,0.04) (-1.371,3.21)rectangle++(0.04,0.04)
  (-1.368,2.313)rectangle++(0.04,0.04) (-2.036,4.136)rectangle++(0.04,0.04) (-1.124,2.892)rectangle++(0.04,0.04) (-2.097,4.109)rectangle++(0.04,0.04) (-1.69,2.354)rectangle++(0.04,0.04)
  (-2.081,4.382)rectangle++(0.04,0.04) (-2.142,3.515)rectangle++(0.04,0.04) (-1.572,2.229)rectangle++(0.04,0.04) (-0.709,2.184)rectangle++(0.04,0.04) (-2.057,3.848)rectangle++(0.04,0.04)
  (-0.681,2.215)rectangle++(0.04,0.04) (-1.903,3.928)rectangle++(0.04,0.04) (-2.151,3.738)rectangle++(0.04,0.04) (-2.177,3.82)rectangle++(0.04,0.04) (-2.042,4.05)rectangle++(0.04,0.04)
  (-1.62,3.506)rectangle++(0.04,0.04) (-1.337,3.331)rectangle++(0.04,0.04) (-2.107,4.36)rectangle++(0.04,0.04) (-2.187,3.765)rectangle++(0.04,0.04) (-0.001,2.213)rectangle++(0.04,0.04)
  (-1.534,3.077)rectangle++(0.04,0.04) (-0.862,2.807)rectangle++(0.04,0.04) (-0.366,2.684)rectangle++(0.04,0.04) (-0.215,2.45)rectangle++(0.04,0.04) (-1.999,3.785)rectangle++(0.04,0.04)
  (-1.604,2.309)rectangle++(0.04,0.04) (-1.02,2.824)rectangle++(0.04,0.04) (-1.433,2.591)rectangle++(0.04,0.04) (-2.007,3.518)rectangle++(0.04,0.04) (-0.686,2.916)rectangle++(0.04,0.04)
  (-1.871,3.861)rectangle++(0.04,0.04) (-2.122,4.346)rectangle++(0.04,0.04) (-2.028,4.115)rectangle++(0.04,0.04) (-0.715,2.502)rectangle++(0.04,0.04) (-1.328,2.58)rectangle++(0.04,0.04)
  (-1.891,3.61)rectangle++(0.04,0.04) (-1.85,3.913)rectangle++(0.04,0.04) (-1.711,3.732)rectangle++(0.04,0.04) (-1.927,3.117)rectangle++(0.04,0.04) (-0.637,2.957)rectangle++(0.04,0.04)
  (-2.12,4.089)rectangle++(0.04,0.04) (-1.69,2.346)rectangle++(0.04,0.04) (-2.069,4.205)rectangle++(0.04,0.04) (-1.843,2.673)rectangle++(0.04,0.04) (-1.919,3.17)rectangle++(0.04,0.04)
  (-1.78,2.786)rectangle++(0.04,0.04) (-1.549,2.469)rectangle++(0.04,0.04) (-1.918,3.788)rectangle++(0.04,0.04) (-2.12,3.605)rectangle++(0.04,0.04) (-1.525,3.358)rectangle++(0.04,0.04)
  (-1.772,3.574)rectangle++(0.04,0.04) (-1.816,2.866)rectangle++(0.04,0.04) (-1.564,3.242)rectangle++(0.04,0.04) (-1.594,2.581)rectangle++(0.04,0.04) (-0.894,3.086)rectangle++(0.04,0.04)
  (-1.849,3.292)rectangle++(0.04,0.04) (-1.247,1.888)rectangle++(0.04,0.04) (-2.096,4.245)rectangle++(0.04,0.04) (-2.035,3.15)rectangle++(0.04,0.04) (-0.366,2.287)rectangle++(0.04,0.04)
  (-1.545,3.188)rectangle++(0.04,0.04) (-1.41,3.049)rectangle++(0.04,0.04) (-1.34,2.275)rectangle++(0.04,0.04) (-1.389,2.447)rectangle++(0.04,0.04) (-2.079,3.752)rectangle++(0.04,0.04)
  (-1.885,2.856)rectangle++(0.04,0.04) (-2.164,4.016)rectangle++(0.04,0.04) (-0.271,2.111)rectangle++(0.04,0.04) (-2.12,4.221)rectangle++(0.04,0.04) (-1.19,3.352)rectangle++(0.04,0.04)
  (-1.977,3.461)rectangle++(0.04,0.04) (-2.106,3.572)rectangle++(0.04,0.04) (-1.371,1.891)rectangle++(0.04,0.04) (-0.291,2.478)rectangle++(0.04,0.04) (-1.98,3.953)rectangle++(0.04,0.04)
  (-2.167,4.232)rectangle++(0.04,0.04) (-2.14,3.494)rectangle++(0.04,0.04) (-1.921,3.906)rectangle++(0.04,0.04) (-1.42,3.028)rectangle++(0.04,0.04) (-0.592,2.856)rectangle++(0.04,0.04)
  (-0.357,2.542)rectangle++(0.04,0.04) (-1.549,3.186)rectangle++(0.04,0.04) (-1.582,2.574)rectangle++(0.04,0.04) (-1.528,3.323)rectangle++(0.04,0.04) (-2.085,4.381)rectangle++(0.04,0.04)
  (-2.104,3.355)rectangle++(0.04,0.04) (-0.433,2.54)rectangle++(0.04,0.04) (-2.002,4)rectangle++(0.04,0.04) (-1.462,2.008)rectangle++(0.04,0.04) (-1.975,3.464)rectangle++(0.04,0.04)
  (-1.93,3.055)rectangle++(0.04,0.04) (-1.287,3.34)rectangle++(0.04,0.04) (-1.147,2.028)rectangle++(0.04,0.04) (-0.768,2.469)rectangle++(0.04,0.04) (-1.625,3.571)rectangle++(0.04,0.04)
  (-0.877,2.476)rectangle++(0.04,0.04) (-1.349,2.552)rectangle++(0.04,0.04) (-1.224,2.188)rectangle++(0.04,0.04) (-1.253,2.938)rectangle++(0.04,0.04) (-1.036,1.947)rectangle++(0.04,0.04)
  (-2.08,4.381)rectangle++(0.04,0.04) (-1.891,3.773)rectangle++(0.04,0.04) (-1.38,1.984)rectangle++(0.04,0.04) (-2.113,3.909)rectangle++(0.04,0.04) (-0.91,3.051)rectangle++(0.04,0.04)
  (-1.849,3.797)rectangle++(0.04,0.04) (-1.759,3.712)rectangle++(0.04,0.04) (-1.834,2.889)rectangle++(0.04,0.04) (-0.694,2.326)rectangle++(0.04,0.04) (-1.852,3.384)rectangle++(0.04,0.04)
  (-1.801,2.827)rectangle++(0.04,0.04) (-1.929,3.192)rectangle++(0.04,0.04) (-0.612,2.378)rectangle++(0.04,0.04) (-0.822,2.716)rectangle++(0.04,0.04) (-0.698,2.576)rectangle++(0.04,0.04)
  (-1.52,2.843)rectangle++(0.04,0.04) (-1.801,3.015)rectangle++(0.04,0.04) (-1.741,2.777)rectangle++(0.04,0.04) (-1.542,2.474)rectangle++(0.04,0.04) (-2.115,3.848)rectangle++(0.04,0.04)
  (-2.19,3.994)rectangle++(0.04,0.04) (-0.457,2.905)rectangle++(0.04,0.04) (-1.356,2.69)rectangle++(0.04,0.04) (-1.631,2.78)rectangle++(0.04,0.04) (-2.113,4.353)rectangle++(0.04,0.04)
  (-2.034,4.138)rectangle++(0.04,0.04) (-2.15,3.631)rectangle++(0.04,0.04) (-0.927,3.007)rectangle++(0.04,0.04) (-1.858,3.926)rectangle++(0.04,0.04) (-1.878,2.91)rectangle++(0.04,0.04)
  (-2.015,4.046)rectangle++(0.04,0.04) (-1.192,3.056)rectangle++(0.04,0.04) (-0.692,3.019)rectangle++(0.04,0.04) (-2.08,3.898)rectangle++(0.04,0.04) (-1.504,3.04)rectangle++(0.04,0.04)
  (-2.039,3.556)rectangle++(0.04,0.04) (-0.342,1.853)rectangle++(0.04,0.04) (-0.637,2.894)rectangle++(0.04,0.04) (-1.168,2.388)rectangle++(0.04,0.04) (-1.041,2.835)rectangle++(0.04,0.04)
  (-2.008,4.096)rectangle++(0.04,0.04) (-2.069,3.327)rectangle++(0.04,0.04) (-0.917,2.813)rectangle++(0.04,0.04) (-1.963,3.909)rectangle++(0.04,0.04) (-0.947,2.555)rectangle++(0.04,0.04)
  (-0.671,2.674)rectangle++(0.04,0.04) (-1.127,2.339)rectangle++(0.04,0.04) (-1.366,2.49)rectangle++(0.04,0.04) (-1.213,2.173)rectangle++(0.04,0.04) (-1.537,2.718)rectangle++(0.04,0.04)
  (-0.697,2.284)rectangle++(0.04,0.04) (-1.753,3.212)rectangle++(0.04,0.04) (-0.929,2.398)rectangle++(0.04,0.04) (-1.983,4.062)rectangle++(0.04,0.04) (-1.474,2.027)rectangle++(0.04,0.04)
  (-0.633,2.676)rectangle++(0.04,0.04) (-0.844,2.032)rectangle++(0.04,0.04) (-0.154,2.311)rectangle++(0.04,0.04) (-1.204,2.632)rectangle++(0.04,0.04) (-0.891,2.361)rectangle++(0.04,0.04)
  (-1.387,2.948)rectangle++(0.04,0.04) (-1.62,3.685)rectangle++(0.04,0.04) (-1.68,2.65)rectangle++(0.04,0.04) (-0.272,2.578)rectangle++(0.04,0.04) (-0.746,2.094)rectangle++(0.04,0.04)
  (-1.717,3.626)rectangle++(0.04,0.04) (-0.59,2.877)rectangle++(0.04,0.04) (-1.373,2.665)rectangle++(0.04,0.04) (-0.989,1.86)rectangle++(0.04,0.04) (-2.132,4.144)rectangle++(0.04,0.04)
  (-1.518,3.303)rectangle++(0.04,0.04) (-1.414,2.304)rectangle++(0.04,0.04) (-0.468,2.47)rectangle++(0.04,0.04) (-1.626,3.691)rectangle++(0.04,0.04) (-1.412,2.927)rectangle++(0.04,0.04)
  (-1.611,3.555)rectangle++(0.04,0.04) (-0.68,2.609)rectangle++(0.04,0.04) (-0.836,2.007)rectangle++(0.04,0.04) (-0.697,2.978)rectangle++(0.04,0.04) (-1.711,3.113)rectangle++(0.04,0.04)
  (-1.35,3.256)rectangle++(0.04,0.04) (-1.729,3.346)rectangle++(0.04,0.04) (-1.778,2.836)rectangle++(0.04,0.04) (-1.408,3.22)rectangle++(0.04,0.04) (-0.17,2.483)rectangle++(0.04,0.04)
  (-1.818,3.702)rectangle++(0.04,0.04) (-0.182,1.717)rectangle++(0.04,0.04) (-2.045,3.913)rectangle++(0.04,0.04) (-0.803,2.93)rectangle++(0.04,0.04) (-0.223,2.539)rectangle++(0.04,0.04)
  (-1.669,3.436)rectangle++(0.04,0.04) (-0.46,2.551)rectangle++(0.04,0.04) (-1.063,2.354)rectangle++(0.04,0.04) (-2.053,3.755)rectangle++(0.04,0.04) (-1.304,3.118)rectangle++(0.04,0.04)
  (-2.095,4.254)rectangle++(0.04,0.04) (-0.118,1.663)rectangle++(0.04,0.04) (-2.08,4.09)rectangle++(0.04,0.04) (-1.861,3.872)rectangle++(0.04,0.04) (-0.952,2.551)rectangle++(0.04,0.04)
  (-1.152,2.262)rectangle++(0.04,0.04) (-1.126,2.169)rectangle++(0.04,0.04) (-0.607,2.565)rectangle++(0.04,0.04) (-1.771,3.412)rectangle++(0.04,0.04) (-1.603,3.589)rectangle++(0.04,0.04)
  (-2.163,4.152)rectangle++(0.04,0.04) (-2.092,3.322)rectangle++(0.04,0.04) (-0.798,2.902)rectangle++(0.04,0.04) (-1.199,2.36)rectangle++(0.04,0.04) (-1.817,3.182)rectangle++(0.04,0.04)
  (-1.878,3.032)rectangle++(0.04,0.04) (-1.185,1.824)rectangle++(0.04,0.04) (-1.484,2.668)rectangle++(0.04,0.04) (-1.956,3.372)rectangle++(0.04,0.04) (-1.955,3.73)rectangle++(0.04,0.04)
  (-1.984,4.09)rectangle++(0.04,0.04) (-0.657,2.771)rectangle++(0.04,0.04) (-1.288,3.23)rectangle++(0.04,0.04) (-1.119,2.003)rectangle++(0.04,0.04) (-0.196,2.272)rectangle++(0.04,0.04)
  (-0.92,2.39)rectangle++(0.04,0.04) (-1.95,3.524)rectangle++(0.04,0.04) (-0.905,2.694)rectangle++(0.04,0.04) (-1.633,3.388)rectangle++(0.04,0.04) (-1.453,2.299)rectangle++(0.04,0.04)
  (-0.91,1.761)rectangle++(0.04,0.04) (-2.132,4.181)rectangle++(0.04,0.04) (-0.67,3.012)rectangle++(0.04,0.04) (-0.375,2.234)rectangle++(0.04,0.04) (-0.677,1.994)rectangle++(0.04,0.04)
  (-0.675,1.872)rectangle++(0.04,0.04) (-0.957,2.181)rectangle++(0.04,0.04) (-0.957,2.503)rectangle++(0.04,0.04) (-1.96,3.533)rectangle++(0.04,0.04) (-1.998,4.041)rectangle++(0.04,0.04)
  (-1.471,2.019)rectangle++(0.04,0.04) (-2.145,4.101)rectangle++(0.04,0.04) (-2.196,3.932)rectangle++(0.04,0.04) (-2.015,4.023)rectangle++(0.04,0.04) (-1.999,3.127)rectangle++(0.04,0.04)
  (-1.644,2.496)rectangle++(0.04,0.04) (-1.912,3.903)rectangle++(0.04,0.04) (-0.931,3.007)rectangle++(0.04,0.04) (-2.033,3.7)rectangle++(0.04,0.04) (-1.39,1.935)rectangle++(0.04,0.04)
  (-0.139,1.768)rectangle++(0.04,0.04) (-0.838,2.26)rectangle++(0.04,0.04) (-1.501,3.51)rectangle++(0.04,0.04) (-1.469,3.444)rectangle++(0.04,0.04) (-0.367,2.777)rectangle++(0.04,0.04)
  (-0.817,2.115)rectangle++(0.04,0.04) (-1.906,3.979)rectangle++(0.04,0.04) (-2.154,3.742)rectangle++(0.04,0.04) (-1.776,2.579)rectangle++(0.04,0.04) (-1.513,2.434)rectangle++(0.04,0.04)
  (-0.408,2.228)rectangle++(0.04,0.04) (-1.23,2.579)rectangle++(0.04,0.04) (-2.055,4.123)rectangle++(0.04,0.04) (-0.109,2.46)rectangle++(0.04,0.04) (-2.088,4.362)rectangle++(0.04,0.04)
  (-1,2.47)rectangle++(0.04,0.04) (-2.11,4.289)rectangle++(0.04,0.04) (-1.096,3.225)rectangle++(0.04,0.04) (-1.578,3.544)rectangle++(0.04,0.04) (-0.798,2.852)rectangle++(0.04,0.04)
  (-0.639,2.814)rectangle++(0.04,0.04) (-2.092,3.963)rectangle++(0.04,0.04) (-1.47,1.98)rectangle++(0.04,0.04) (-1.547,2.657)rectangle++(0.04,0.04) (-1.857,3.868)rectangle++(0.04,0.04)
  (-0.381,2.362)rectangle++(0.04,0.04) (-1.612,3.639)rectangle++(0.04,0.04) (-1.332,2.976)rectangle++(0.04,0.04) (-2.073,4.022)rectangle++(0.04,0.04) (-1.481,1.997)rectangle++(0.04,0.04)
  (-1.977,3.834)rectangle++(0.04,0.04) (-2.128,3.606)rectangle++(0.04,0.04) (-0.004,2.198)rectangle++(0.04,0.04) (-1.366,2.868)rectangle++(0.04,0.04) (-1.278,3.271)rectangle++(0.04,0.04)
  (-0.393,2.791)rectangle++(0.04,0.04) (-1.326,2.679)rectangle++(0.04,0.04) (-1.345,3.433)rectangle++(0.04,0.04) (-1.183,2.035)rectangle++(0.04,0.04) (-1.513,2.7)rectangle++(0.04,0.04)
  (-2.145,4.014)rectangle++(0.04,0.04) (-2.122,3.473)rectangle++(0.04,0.04) (-1.852,2.808)rectangle++(0.04,0.04) (-2.144,4.282)rectangle++(0.04,0.04) (-1.563,2.038)rectangle++(0.04,0.04)
  (-1.569,2.655)rectangle++(0.04,0.04) (-1.964,3.365)rectangle++(0.04,0.04) (-1.289,1.878)rectangle++(0.04,0.04) (-2.103,4.312)rectangle++(0.04,0.04) (-1.829,3.829)rectangle++(0.04,0.04)
  (-1.739,2.644)rectangle++(0.04,0.04) (-2.145,3.937)rectangle++(0.04,0.04) (-2.191,3.923)rectangle++(0.04,0.04) (-1.386,3.261)rectangle++(0.04,0.04) (-0.138,2.399)rectangle++(0.04,0.04)
  (-1.414,2.883)rectangle++(0.04,0.04) (-1.306,2.241)rectangle++(0.04,0.04) (-1.916,3.817)rectangle++(0.04,0.04) (-0.563,2.842)rectangle++(0.04,0.04) (-1.643,3.305)rectangle++(0.04,0.04)
  (-1.617,2.581)rectangle++(0.04,0.04) (0.005,1.945)rectangle++(0.04,0.04) (-1.321,3.157)rectangle++(0.04,0.04) (-1.986,3.796)rectangle++(0.04,0.04) (-1.331,3.032)rectangle++(0.04,0.04)
  (-1.076,1.95)rectangle++(0.04,0.04) (-1.27,3.148)rectangle++(0.04,0.04) (-1.683,2.861)rectangle++(0.04,0.04) (-1.918,3.205)rectangle++(0.04,0.04) (-2.166,3.964)rectangle++(0.04,0.04)
  (-1.701,3.671)rectangle++(0.04,0.04) (-0.691,2.167)rectangle++(0.04,0.04) (-1.023,2.236)rectangle++(0.04,0.04) (-1.331,2.488)rectangle++(0.04,0.04) (-0.165,2.437)rectangle++(0.04,0.04)
  (-2.055,3.965)rectangle++(0.04,0.04) (-0.69,2.856)rectangle++(0.04,0.04) (-1.957,3.578)rectangle++(0.04,0.04) (-1.759,3.594)rectangle++(0.04,0.04) (-1.882,3.813)rectangle++(0.04,0.04)
  (-0.988,2.743)rectangle++(0.04,0.04) (-1.272,3.179)rectangle++(0.04,0.04) (-1.102,1.996)rectangle++(0.04,0.04) (-1.902,3.913)rectangle++(0.04,0.04) (-0.155,2.526)rectangle++(0.04,0.04)
  (-0.848,2.317)rectangle++(0.04,0.04) (-2.131,4.251)rectangle++(0.04,0.04) (-2.097,3.336)rectangle++(0.04,0.04) (-1.308,1.841)rectangle++(0.04,0.04) (-1.596,3.55)rectangle++(0.04,0.04)
  (-1.106,3.022)rectangle++(0.04,0.04) (-1.9,3.95)rectangle++(0.04,0.04) (-0.243,1.91)rectangle++(0.04,0.04) (-0.566,1.953)rectangle++(0.04,0.04) (-1.888,3.743)rectangle++(0.04,0.04)
  (-2.164,4.227)rectangle++(0.04,0.04) (-2.079,4.068)rectangle++(0.04,0.04) (-1.495,2.006)rectangle++(0.04,0.04) (-2.148,4.147)rectangle++(0.04,0.04) (-2.089,3.322)rectangle++(0.04,0.04)
  (-0.347,2.566)rectangle++(0.04,0.04) (-1.375,3.083)rectangle++(0.04,0.04) (-2.03,3.566)rectangle++(0.04,0.04) (-1.814,3.709)rectangle++(0.04,0.04) (-2.011,3.296)rectangle++(0.04,0.04)
  (-1.569,3.568)rectangle++(0.04,0.04) (-0.621,2.187)rectangle++(0.04,0.04) (-1.611,3.013)rectangle++(0.04,0.04) (-1.713,2.799)rectangle++(0.04,0.04) (-0.952,2.988)rectangle++(0.04,0.04)
  (-0.998,2.468)rectangle++(0.04,0.04) (-0.697,2.953)rectangle++(0.04,0.04) (-1.902,3.943)rectangle++(0.04,0.04) (-1.427,2.019)rectangle++(0.04,0.04) (-1.695,2.912)rectangle++(0.04,0.04)
  (-1.558,2.504)rectangle++(0.04,0.04) (-0.814,2.867)rectangle++(0.04,0.04) (-0.374,2.803)rectangle++(0.04,0.04) (-1.767,3.406)rectangle++(0.04,0.04) (-2.167,4.145)rectangle++(0.04,0.04)
  (-1.711,2.348)rectangle++(0.04,0.04) (-1.245,2.081)rectangle++(0.04,0.04) (-1.34,2.427)rectangle++(0.04,0.04) (-1.786,3.58)rectangle++(0.04,0.04) (-2.096,3.592)rectangle++(0.04,0.04)
  (-1.82,3.488)rectangle++(0.04,0.04) (-1.292,1.938)rectangle++(0.04,0.04) (-0.689,2.274)rectangle++(0.04,0.04) (-0.713,2.509)rectangle++(0.04,0.04) (-1.132,3.178)rectangle++(0.04,0.04)
  (-1.074,2.028)rectangle++(0.04,0.04) (-1.303,2.444)rectangle++(0.04,0.04) (-2.006,3.95)rectangle++(0.04,0.04) (-1.656,3.494)rectangle++(0.04,0.04) (-1.951,4.04)rectangle++(0.04,0.04)
  (-1.755,3.688)rectangle++(0.04,0.04) (-1.573,3.238)rectangle++(0.04,0.04) (-2.007,3.426)rectangle++(0.04,0.04) (-2.176,3.974)rectangle++(0.04,0.04) (-2.031,4.102)rectangle++(0.04,0.04)
  (-0.261,2.545)rectangle++(0.04,0.04) (-1.262,2.657)rectangle++(0.04,0.04) (-1.861,3.699)rectangle++(0.04,0.04) (-2.108,4.312)rectangle++(0.04,0.04) (-1.564,2.053)rectangle++(0.04,0.04)
  (-1.978,4.068)rectangle++(0.04,0.04) (-2.141,3.627)rectangle++(0.04,0.04) (-2.065,4.21)rectangle++(0.04,0.04) (-2.157,3.636)rectangle++(0.04,0.04) (-1.88,2.829)rectangle++(0.04,0.04)
  (-1.938,3.186)rectangle++(0.04,0.04) (-1.894,3.021)rectangle++(0.04,0.04) (-1.01,2.64)rectangle++(0.04,0.04) (-1.128,2.865)rectangle++(0.04,0.04) (-1.962,3.797)rectangle++(0.04,0.04)
  (-0.186,2.588)rectangle++(0.04,0.04) (-0.955,2.517)rectangle++(0.04,0.04) (-0.866,1.923)rectangle++(0.04,0.04) (-1.584,3.547)rectangle++(0.04,0.04) (-2.086,4.31)rectangle++(0.04,0.04)
  (-1.749,3.624)rectangle++(0.04,0.04) (-1.655,3.438)rectangle++(0.04,0.04) (-1.246,1.965)rectangle++(0.04,0.04) (-1.665,2.922)rectangle++(0.04,0.04) (-2.151,4.237)rectangle++(0.04,0.04)
  (-0.058,2.374)rectangle++(0.04,0.04) (-0.528,2.255)rectangle++(0.04,0.04) (-1.701,3.548)rectangle++(0.04,0.04) (-2.101,4.304)rectangle++(0.04,0.04) (-1.956,3.972)rectangle++(0.04,0.04)
  (-0.756,2.227)rectangle++(0.04,0.04) (-2.028,3.755)rectangle++(0.04,0.04) (-0.845,3.018)rectangle++(0.04,0.04) (-0.919,3.115)rectangle++(0.04,0.04) (-1.467,2.657)rectangle++(0.04,0.04)
  (-2.151,4.165)rectangle++(0.04,0.04) (-2.086,4.219)rectangle++(0.04,0.04) (-0.149,2.545)rectangle++(0.04,0.04) (-1.118,3.127)rectangle++(0.04,0.04) (-2.124,3.972)rectangle++(0.04,0.04)
  (-2.191,3.925)rectangle++(0.04,0.04) (-2.157,4.28)rectangle++(0.04,0.04) (-2.091,4.292)rectangle++(0.04,0.04) (-2.087,4.26)rectangle++(0.04,0.04) (-1.462,2.986)rectangle++(0.04,0.04)
  (-1.09,1.91)rectangle++(0.04,0.04) (-1.603,3.427)rectangle++(0.04,0.04) (-0.199,1.897)rectangle++(0.04,0.04) (-0.554,1.96)rectangle++(0.04,0.04) (-0.95,3.164)rectangle++(0.04,0.04)
  (-0.996,3.22)rectangle++(0.04,0.04) (-1.649,3.234)rectangle++(0.04,0.04) (-1.781,3.83)rectangle++(0.04,0.04) (-1.731,2.651)rectangle++(0.04,0.04) (-0.239,2.022)rectangle++(0.04,0.04)
  (-1.933,3.823)rectangle++(0.04,0.04) (-1.637,3.448)rectangle++(0.04,0.04) (-2.065,4.185)rectangle++(0.04,0.04) (-0.671,3.012)rectangle++(0.04,0.04) (-1.08,3.02)rectangle++(0.04,0.04)
  (-1.022,2.005)rectangle++(0.04,0.04) (-1.291,2.448)rectangle++(0.04,0.04) (-1.899,3.939)rectangle++(0.04,0.04) (-0.026,2.235)rectangle++(0.04,0.04) (-0.607,2.075)rectangle++(0.04,0.04)
  (-0.747,3.068)rectangle++(0.04,0.04) (-0.428,2.871)rectangle++(0.04,0.04) (-1.166,3.206)rectangle++(0.04,0.04) (-1.674,3.269)rectangle++(0.04,0.04) (-0.039,2.171)rectangle++(0.04,0.04)
  (-0.958,2.844)rectangle++(0.04,0.04) (-0.739,3.053)rectangle++(0.04,0.04) (-1.828,3.3)rectangle++(0.04,0.04) (-1.571,3.568)rectangle++(0.04,0.04) (-1.786,2.89)rectangle++(0.04,0.04)
  (-0.387,2.753)rectangle++(0.04,0.04) (-2.011,4.018)rectangle++(0.04,0.04) (-1.913,2.906)rectangle++(0.04,0.04) (-1,2.589)rectangle++(0.04,0.04) (-2.051,3.955)rectangle++(0.04,0.04)
  (-1.687,3.308)rectangle++(0.04,0.04) (-0.276,2.716)rectangle++(0.04,0.04) (-0.34,2.516)rectangle++(0.04,0.04) (-1.459,2.877)rectangle++(0.04,0.04) (-2.068,3.671)rectangle++(0.04,0.04)
  (-0.823,2.971)rectangle++(0.04,0.04) (-2.111,4.345)rectangle++(0.04,0.04) (-1.874,2.691)rectangle++(0.04,0.04) (-0.76,2.722)rectangle++(0.04,0.04) (-1.463,2.988)rectangle++(0.04,0.04)
  (-0.838,2.945)rectangle++(0.04,0.04) (-2.135,4.175)rectangle++(0.04,0.04) (-0.847,2.503)rectangle++(0.04,0.04) (-1.531,2.813)rectangle++(0.04,0.04) (-1.55,3.273)rectangle++(0.04,0.04)
  (-1.637,3.531)rectangle++(0.04,0.04) (-0.989,2.458)rectangle++(0.04,0.04) (-1.545,2.789)rectangle++(0.04,0.04) (-0.122,2.364)rectangle++(0.04,0.04) (-0.244,2.508)rectangle++(0.04,0.04)
  (-1.295,3.175)rectangle++(0.04,0.04) (-1.693,3.719)rectangle++(0.04,0.04) (-1.327,2.018)rectangle++(0.04,0.04) (-0.811,1.728)rectangle++(0.04,0.04) (-0.951,2.127)rectangle++(0.04,0.04)
  (-1.061,2.179)rectangle++(0.04,0.04) (-0.662,2.224)rectangle++(0.04,0.04) (-0.732,1.926)rectangle++(0.04,0.04) (-0.222,2.397)rectangle++(0.04,0.04) (-1.583,3.08)rectangle++(0.04,0.04)
  (-1.927,3.238)rectangle++(0.04,0.04) (-2.146,4.289)rectangle++(0.04,0.04) (-1.737,2.38)rectangle++(0.04,0.04) (-1.472,2.405)rectangle++(0.04,0.04) (-0.047,2.087)rectangle++(0.04,0.04)
  (-1.815,3.732)rectangle++(0.04,0.04) (-0.261,1.784)rectangle++(0.04,0.04) (-0.864,2.241)rectangle++(0.04,0.04) (-2.109,4.348)rectangle++(0.04,0.04) (-2.049,3.162)rectangle++(0.04,0.04)
  (-1.658,2.494)rectangle++(0.04,0.04) (-0.95,2.836)rectangle++(0.04,0.04) (-1.713,3.636)rectangle++(0.04,0.04);
\fill[piece4]
  (2.412,3.408)rectangle++(0.04,0.04) (1.147,2.217)rectangle++(0.04,0.04) (2.244,2.563)rectangle++(0.04,0.04) (1.435,1.251)rectangle++(0.04,0.04) (2.296,3.017)rectangle++(0.04,0.04)
  (1.087,2.122)rectangle++(0.04,0.04) (1.469,0.894)rectangle++(0.04,0.04) (2.406,3.455)rectangle++(0.04,0.04) (1.552,1.597)rectangle++(0.04,0.04) (1.926,1.801)rectangle++(0.04,0.04)
  (0.198,1.363)rectangle++(0.04,0.04) (2.046,2.611)rectangle++(0.04,0.04) (1.883,1.749)rectangle++(0.04,0.04) (1.94,2.071)rectangle++(0.04,0.04) (0.951,1.881)rectangle++(0.04,0.04)
  (2.166,2.927)rectangle++(0.04,0.04) (0.76,1.111)rectangle++(0.04,0.04) (1.385,1.723)rectangle++(0.04,0.04) (0.257,1.204)rectangle++(0.04,0.04) (1.715,1.988)rectangle++(0.04,0.04)
  (1.762,2.522)rectangle++(0.04,0.04) (2.21,2.842)rectangle++(0.04,0.04) (0.447,1.687)rectangle++(0.04,0.04) (1.739,2)rectangle++(0.04,0.04) (1.436,2.218)rectangle++(0.04,0.04)
  (0.745,1.592)rectangle++(0.04,0.04) (2.166,2.903)rectangle++(0.04,0.04) (2.183,2.479)rectangle++(0.04,0.04) (1.475,1.447)rectangle++(0.04,0.04) (2.415,3.414)rectangle++(0.04,0.04)
  (2.232,3.176)rectangle++(0.04,0.04) (2.243,2.562)rectangle++(0.04,0.04) (0.713,1.836)rectangle++(0.04,0.04) (2.369,3.207)rectangle++(0.04,0.04) (0.85,0.648)rectangle++(0.04,0.04)
  (0.794,0.833)rectangle++(0.04,0.04) (1.62,2.01)rectangle++(0.04,0.04) (1.772,2.518)rectangle++(0.04,0.04) (0.768,1.727)rectangle++(0.04,0.04) (0.761,0.838)rectangle++(0.04,0.04)
  (2.125,2.739)rectangle++(0.04,0.04) (0.88,1.227)rectangle++(0.04,0.04) (1.613,2.204)rectangle++(0.04,0.04) (1.018,1.85)rectangle++(0.04,0.04) (1.283,1.397)rectangle++(0.04,0.04)
  (2.182,2.821)rectangle++(0.04,0.04) (1.906,2.76)rectangle++(0.04,0.04) (0.806,1.807)rectangle++(0.04,0.04) (1.668,1.522)rectangle++(0.04,0.04) (2.363,2.974)rectangle++(0.04,0.04)
  (2.163,2.478)rectangle++(0.04,0.04) (1,1.997)rectangle++(0.04,0.04) (2.34,2.78)rectangle++(0.04,0.04) (1.182,0.893)rectangle++(0.04,0.04) (1.133,1.21)rectangle++(0.04,0.04)
  (2.21,2.318)rectangle++(0.04,0.04) (2.251,3.173)rectangle++(0.04,0.04) (1.122,2.145)rectangle++(0.04,0.04) (1.959,2.479)rectangle++(0.04,0.04) (1.254,1.611)rectangle++(0.04,0.04)
  (2.397,3.228)rectangle++(0.04,0.04) (1.735,1.312)rectangle++(0.04,0.04) (1.779,1.769)rectangle++(0.04,0.04) (0.211,1.317)rectangle++(0.04,0.04) (2.184,2.872)rectangle++(0.04,0.04)
  (1.794,1.749)rectangle++(0.04,0.04) (1.708,1.655)rectangle++(0.04,0.04) (2.434,3.209)rectangle++(0.04,0.04) (1.106,2.183)rectangle++(0.04,0.04) (1.02,1.233)rectangle++(0.04,0.04)
  (1.222,0.828)rectangle++(0.04,0.04) (1.538,2.181)rectangle++(0.04,0.04) (1.563,1.159)rectangle++(0.04,0.04) (1.946,2.686)rectangle++(0.04,0.04) (1.068,1.999)rectangle++(0.04,0.04)
  (1.583,1.379)rectangle++(0.04,0.04) (1.041,0.857)rectangle++(0.04,0.04) (1.39,0.786)rectangle++(0.04,0.04) (1.686,2.361)rectangle++(0.04,0.04) (1.911,2.687)rectangle++(0.04,0.04)
  (1.072,1.994)rectangle++(0.04,0.04) (2.428,3.356)rectangle++(0.04,0.04) (0.214,0.621)rectangle++(0.04,0.04) (1.417,1.844)rectangle++(0.04,0.04) (2.297,3.198)rectangle++(0.04,0.04)
  (1.817,2.702)rectangle++(0.04,0.04) (1.429,1.064)rectangle++(0.04,0.04) (1.477,1.53)rectangle++(0.04,0.04) (1.402,2.119)rectangle++(0.04,0.04) (1.024,0.766)rectangle++(0.04,0.04)
  (2.259,2.819)rectangle++(0.04,0.04) (1.368,0.799)rectangle++(0.04,0.04) (2.21,2.325)rectangle++(0.04,0.04) (1.969,2.815)rectangle++(0.04,0.04) (1.06,1.406)rectangle++(0.04,0.04)
  (1.038,0.827)rectangle++(0.04,0.04) (2.028,1.891)rectangle++(0.04,0.04) (2.125,2.299)rectangle++(0.04,0.04) (2.299,2.868)rectangle++(0.04,0.04) (0.776,1.913)rectangle++(0.04,0.04)
  (2.176,2.925)rectangle++(0.04,0.04) (2.263,3.199)rectangle++(0.04,0.04) (1.709,1.481)rectangle++(0.04,0.04) (2.344,3.103)rectangle++(0.04,0.04) (0.908,0.762)rectangle++(0.04,0.04)
  (1.5,0.909)rectangle++(0.04,0.04) (0.138,0.555)rectangle++(0.04,0.04) (2.16,2.82)rectangle++(0.04,0.04) (0.107,0.529)rectangle++(0.04,0.04) (1.178,1.562)rectangle++(0.04,0.04)
  (1.62,2.274)rectangle++(0.04,0.04) (1.546,1.232)rectangle++(0.04,0.04) (1.529,1.187)rectangle++(0.04,0.04) (0.833,1.642)rectangle++(0.04,0.04) (0.229,0.634)rectangle++(0.04,0.04)
  (0.893,1.314)rectangle++(0.04,0.04) (2.343,3.267)rectangle++(0.04,0.04) (1.898,1.857)rectangle++(0.04,0.04) (1.401,2.214)rectangle++(0.04,0.04) (1.769,2.048)rectangle++(0.04,0.04)
  (1.88,1.595)rectangle++(0.04,0.04) (2.378,3.444)rectangle++(0.04,0.04) (2.308,3.042)rectangle++(0.04,0.04) (0.878,0.587)rectangle++(0.04,0.04) (1.95,2.234)rectangle++(0.04,0.04)
  (0.678,1.704)rectangle++(0.04,0.04) (1.297,1.84)rectangle++(0.04,0.04) (0.784,0.886)rectangle++(0.04,0.04) (0.527,1.168)rectangle++(0.04,0.04) (1.443,1.76)rectangle++(0.04,0.04)
  (1.616,2.323)rectangle++(0.04,0.04) (0.745,1.652)rectangle++(0.04,0.04) (2.09,2.564)rectangle++(0.04,0.04) (1.78,2.615)rectangle++(0.04,0.04) (0.79,1.749)rectangle++(0.04,0.04)
  (1.291,1.863)rectangle++(0.04,0.04) (2.338,3.282)rectangle++(0.04,0.04) (0.551,1.826)rectangle++(0.04,0.04) (1.331,1.843)rectangle++(0.04,0.04) (2.181,2.441)rectangle++(0.04,0.04)
  (1.377,1.254)rectangle++(0.04,0.04) (1.857,2.559)rectangle++(0.04,0.04) (2.428,3.144)rectangle++(0.04,0.04) (2.074,2.976)rectangle++(0.04,0.04) (2.293,3.244)rectangle++(0.04,0.04)
  (1.298,1.286)rectangle++(0.04,0.04) (2.113,2.23)rectangle++(0.04,0.04) (1.472,1.578)rectangle++(0.04,0.04) (1.482,1.824)rectangle++(0.04,0.04) (0.271,1.253)rectangle++(0.04,0.04)
  (1.673,2.154)rectangle++(0.04,0.04) (2.38,3.476)rectangle++(0.04,0.04) (0.906,2.061)rectangle++(0.04,0.04) (0.782,1.446)rectangle++(0.04,0.04) (1.915,2.546)rectangle++(0.04,0.04)
  (1.589,1.162)rectangle++(0.04,0.04) (2.38,3.478)rectangle++(0.04,0.04) (0.907,2.061)rectangle++(0.04,0.04) (1.164,1.45)rectangle++(0.04,0.04) (2.417,3.402)rectangle++(0.04,0.04)
  (1.145,2.217)rectangle++(0.04,0.04) (1.08,1.773)rectangle++(0.04,0.04) (2.429,3.156)rectangle++(0.04,0.04) (2.023,2.1)rectangle++(0.04,0.04) (2.148,2.251)rectangle++(0.04,0.04)
  (1.668,2.088)rectangle++(0.04,0.04) (2.402,3.245)rectangle++(0.04,0.04) (1.952,2.301)rectangle++(0.04,0.04) (2.388,3.385)rectangle++(0.04,0.04) (1.767,1.582)rectangle++(0.04,0.04)
  (2.144,2.145)rectangle++(0.04,0.04) (2.393,3.13)rectangle++(0.04,0.04) (0.824,1.989)rectangle++(0.04,0.04) (2.331,2.665)rectangle++(0.04,0.04) (1.628,2.52)rectangle++(0.04,0.04)
  (1.421,1.069)rectangle++(0.04,0.04) (2.355,3.329)rectangle++(0.04,0.04) (2.37,3.435)rectangle++(0.04,0.04) (2.429,3.246)rectangle++(0.04,0.04) (0.846,2.022)rectangle++(0.04,0.04)
  (2.177,2.934)rectangle++(0.04,0.04) (2.141,3.027)rectangle++(0.04,0.04) (1.391,1.323)rectangle++(0.04,0.04) (2.215,2.329)rectangle++(0.04,0.04) (2.276,2.594)rectangle++(0.04,0.04)
  (0.368,1.644)rectangle++(0.04,0.04) (0.388,0.92)rectangle++(0.04,0.04) (2.376,3.082)rectangle++(0.04,0.04) (2,2.889)rectangle++(0.04,0.04) (1.133,0.688)rectangle++(0.04,0.04)
  (2.29,3.154)rectangle++(0.04,0.04) (1.67,1.575)rectangle++(0.04,0.04) (1.819,1.577)rectangle++(0.04,0.04) (0.871,1.766)rectangle++(0.04,0.04) (2.347,3.279)rectangle++(0.04,0.04)
  (0.549,1.827)rectangle++(0.04,0.04) (0.975,1.185)rectangle++(0.04,0.04) (0.904,1.547)rectangle++(0.04,0.04) (1.233,0.826)rectangle++(0.04,0.04) (2.417,3.124)rectangle++(0.04,0.04)
  (1.32,2.319)rectangle++(0.04,0.04) (1.606,1.585)rectangle++(0.04,0.04) (2.271,2.85)rectangle++(0.04,0.04) (1.699,1.856)rectangle++(0.04,0.04) (1.574,1.545)rectangle++(0.04,0.04)
  (1.688,1.239)rectangle++(0.04,0.04) (1.348,1.283)rectangle++(0.04,0.04) (0.875,1.63)rectangle++(0.04,0.04) (1.561,1.781)rectangle++(0.04,0.04) (1.556,1.632)rectangle++(0.04,0.04)
  (1.98,2.262)rectangle++(0.04,0.04) (0.18,0.592)rectangle++(0.04,0.04) (1.417,1.836)rectangle++(0.04,0.04) (1.994,2.739)rectangle++(0.04,0.04) (2.301,2.63)rectangle++(0.04,0.04)
  (1.44,1.816)rectangle++(0.04,0.04) (1.705,1.81)rectangle++(0.04,0.04) (2.154,2.376)rectangle++(0.04,0.04) (1.519,1.153)rectangle++(0.04,0.04) (1.706,2.413)rectangle++(0.04,0.04)
  (2.345,3.203)rectangle++(0.04,0.04) (1.116,2.166)rectangle++(0.04,0.04) (1.521,0.938)rectangle++(0.04,0.04) (0.436,1.335)rectangle++(0.04,0.04) (1.363,1.728)rectangle++(0.04,0.04)
  (1.991,2.727)rectangle++(0.04,0.04) (1.595,1.195)rectangle++(0.04,0.04) (2.035,2.801)rectangle++(0.04,0.04) (2.365,3.402)rectangle++(0.04,0.04) (2.337,2.787)rectangle++(0.04,0.04)
  (1.276,2.251)rectangle++(0.04,0.04) (0.422,1.302)rectangle++(0.04,0.04) (2.364,2.793)rectangle++(0.04,0.04) (2.37,3.43)rectangle++(0.04,0.04) (1.555,2.501)rectangle++(0.04,0.04)
  (0.702,0.633)rectangle++(0.04,0.04) (0.833,1.395)rectangle++(0.04,0.04) (2.262,2.502)rectangle++(0.04,0.04) (2.252,3.182)rectangle++(0.04,0.04) (1.124,2.146)rectangle++(0.04,0.04)
  (1.825,2.499)rectangle++(0.04,0.04) (1.796,1.574)rectangle++(0.04,0.04) (2.041,2.836)rectangle++(0.04,0.04) (0.979,0.919)rectangle++(0.04,0.04) (1.375,1.973)rectangle++(0.04,0.04)
  (1.403,1.082)rectangle++(0.04,0.04) (2.37,3.379)rectangle++(0.04,0.04) (0.573,1.857)rectangle++(0.04,0.04) (1.488,1.764)rectangle++(0.04,0.04) (1.885,2.616)rectangle++(0.04,0.04)
  (2.079,2.601)rectangle++(0.04,0.04) (0.962,0.923)rectangle++(0.04,0.04) (1.292,1.14)rectangle++(0.04,0.04) (0.515,1.33)rectangle++(0.04,0.04) (1.93,2.386)rectangle++(0.04,0.04)
  (2.371,3.22)rectangle++(0.04,0.04) (1.116,2.174)rectangle++(0.04,0.04) (2.077,2.202)rectangle++(0.04,0.04) (1.53,1.865)rectangle++(0.04,0.04) (0.65,1.536)rectangle++(0.04,0.04)
  (2.314,2.732)rectangle++(0.04,0.04) (2.322,2.966)rectangle++(0.04,0.04) (2.285,3.245)rectangle++(0.04,0.04) (2.087,2.979)rectangle++(0.04,0.04) (1.77,1.871)rectangle++(0.04,0.04)
  (1.348,1.231)rectangle++(0.04,0.04) (1.964,2.683)rectangle++(0.04,0.04) (0.78,1.803)rectangle++(0.04,0.04) (1.6,1.38)rectangle++(0.04,0.04) (0.918,1.265)rectangle++(0.04,0.04)
  (0.926,1.552)rectangle++(0.04,0.04) (2.245,2.591)rectangle++(0.04,0.04) (1.459,2.362)rectangle++(0.04,0.04) (1.672,1.814)rectangle++(0.04,0.04) (1.199,1.995)rectangle++(0.04,0.04)
  (1.492,2.158)rectangle++(0.04,0.04) (1.847,2.256)rectangle++(0.04,0.04) (2.312,3.258)rectangle++(0.04,0.04) (1.353,2.327)rectangle++(0.04,0.04) (1.317,2.038)rectangle++(0.04,0.04)
  (1.658,2.348)rectangle++(0.04,0.04) (2.032,2.087)rectangle++(0.04,0.04) (1.601,1.189)rectangle++(0.04,0.04) (2.033,2.364)rectangle++(0.04,0.04) (1.452,2.304)rectangle++(0.04,0.04)
  (1.813,1.511)rectangle++(0.04,0.04) (2.166,3.004)rectangle++(0.04,0.04) (2.323,3.141)rectangle++(0.04,0.04) (1.952,2.469)rectangle++(0.04,0.04) (1.413,1.582)rectangle++(0.04,0.04)
  (1.431,1.536)rectangle++(0.04,0.04) (2.271,2.848)rectangle++(0.04,0.04) (1.293,2.252)rectangle++(0.04,0.04) (1.683,2.374)rectangle++(0.04,0.04) (2.146,2.978)rectangle++(0.04,0.04)
  (2.262,2.998)rectangle++(0.04,0.04) (2.415,3.424)rectangle++(0.04,0.04) (1.998,2.222)rectangle++(0.04,0.04) (0.305,1.483)rectangle++(0.04,0.04) (1.711,1.28)rectangle++(0.04,0.04)
  (2.118,2.292)rectangle++(0.04,0.04) (2.173,2.474)rectangle++(0.04,0.04) (2.433,3.335)rectangle++(0.04,0.04) (0.936,1.135)rectangle++(0.04,0.04) (1.715,2.02)rectangle++(0.04,0.04)
  (1.618,2.381)rectangle++(0.04,0.04) (1.902,1.619)rectangle++(0.04,0.04) (2.4,3.407)rectangle++(0.04,0.04) (2.433,3.249)rectangle++(0.04,0.04) (2.397,3.414)rectangle++(0.04,0.04)
  (2.397,2.921)rectangle++(0.04,0.04) (1.808,2.1)rectangle++(0.04,0.04) (1.983,2.779)rectangle++(0.04,0.04) (2.065,2.195)rectangle++(0.04,0.04) (1.971,2.798)rectangle++(0.04,0.04)
  (0.47,1.614)rectangle++(0.04,0.04) (1.835,1.797)rectangle++(0.04,0.04) (1.4,2.198)rectangle++(0.04,0.04) (1.77,1.814)rectangle++(0.04,0.04) (1.923,1.776)rectangle++(0.04,0.04)
  (1.388,2.205)rectangle++(0.04,0.04) (1.766,2.212)rectangle++(0.04,0.04) (2.388,3.483)rectangle++(0.04,0.04) (0.907,2.064)rectangle++(0.04,0.04) (0.425,1.161)rectangle++(0.04,0.04)
  (1.035,1.166)rectangle++(0.04,0.04) (1.592,2.209)rectangle++(0.04,0.04) (1.878,1.834)rectangle++(0.04,0.04) (0.215,1.358)rectangle++(0.04,0.04) (0.651,0.664)rectangle++(0.04,0.04)
  (0.908,0.549)rectangle++(0.04,0.04) (2.092,2.687)rectangle++(0.04,0.04) (0.421,1.619)rectangle++(0.04,0.04) (0.476,0.797)rectangle++(0.04,0.04) (0.893,1.388)rectangle++(0.04,0.04)
  (1.532,1.153)rectangle++(0.04,0.04) (0.076,1.105)rectangle++(0.04,0.04) (2.183,2.577)rectangle++(0.04,0.04) (2.041,2.289)rectangle++(0.04,0.04) (1.164,1.515)rectangle++(0.04,0.04)
  (2.322,3.232)rectangle++(0.04,0.04) (1.142,1.336)rectangle++(0.04,0.04) (2.254,2.671)rectangle++(0.04,0.04) (1.89,2.749)rectangle++(0.04,0.04) (0.805,1.802)rectangle++(0.04,0.04)
  (0.366,1.072)rectangle++(0.04,0.04) (2.386,3.485)rectangle++(0.04,0.04) (1.849,1.857)rectangle++(0.04,0.04) (2.335,2.68)rectangle++(0.04,0.04) (0.383,1.68)rectangle++(0.04,0.04)
  (2.174,2.825)rectangle++(0.04,0.04) (2.137,3.019)rectangle++(0.04,0.04) (2.42,3.428)rectangle++(0.04,0.04) (1.368,2.371)rectangle++(0.04,0.04) (2.408,2.997)rectangle++(0.04,0.04)
  (1.393,1.28)rectangle++(0.04,0.04) (2.435,3.339)rectangle++(0.04,0.04) (1.69,2.617)rectangle++(0.04,0.04) (2.414,3.408)rectangle++(0.04,0.04) (1.807,1.588)rectangle++(0.04,0.04)
  (0.156,1.281)rectangle++(0.04,0.04) (1.546,1.129)rectangle++(0.04,0.04) (1.918,2.254)rectangle++(0.04,0.04) (1.744,1.564)rectangle++(0.04,0.04) (2.42,3.147)rectangle++(0.04,0.04)
  (1.9,2.306)rectangle++(0.04,0.04) (1.459,2.281)rectangle++(0.04,0.04) (0.283,0.68)rectangle++(0.04,0.04) (1.164,1.474)rectangle++(0.04,0.04) (0.221,1.086)rectangle++(0.04,0.04)
  (1.777,1.994)rectangle++(0.04,0.04) (2.31,3.108)rectangle++(0.04,0.04) (1.519,2.446)rectangle++(0.04,0.04) (1.223,1.014)rectangle++(0.04,0.04) (0.839,1.557)rectangle++(0.04,0.04)
  (1.454,0.883)rectangle++(0.04,0.04) (1.498,1.384)rectangle++(0.04,0.04) (0.146,1.152)rectangle++(0.04,0.04) (1.493,1.915)rectangle++(0.04,0.04) (2.391,2.967)rectangle++(0.04,0.04)
  (2.181,3.098)rectangle++(0.04,0.04) (1.416,1.789)rectangle++(0.04,0.04) (2.129,2.712)rectangle++(0.04,0.04) (0.764,1.844)rectangle++(0.04,0.04) (0.752,1.369)rectangle++(0.04,0.04)
  (1.64,2.211)rectangle++(0.04,0.04) (2.408,3)rectangle++(0.04,0.04) (1.791,2.687)rectangle++(0.04,0.04) (2.107,2.456)rectangle++(0.04,0.04) (1.003,1.983)rectangle++(0.04,0.04)
  (0.751,1.45)rectangle++(0.04,0.04) (0.721,0.645)rectangle++(0.04,0.04) (2.341,3.144)rectangle++(0.04,0.04) (1.332,2.312)rectangle++(0.04,0.04) (1.074,1.823)rectangle++(0.04,0.04)
  (0.42,0.797)rectangle++(0.04,0.04) (2.371,3.324)rectangle++(0.04,0.04) (2.033,2.527)rectangle++(0.04,0.04) (1.27,2.165)rectangle++(0.04,0.04) (1.619,1.177)rectangle++(0.04,0.04)
  (2.18,2.603)rectangle++(0.04,0.04) (2.345,2.79)rectangle++(0.04,0.04) (1.834,1.538)rectangle++(0.04,0.04) (0.609,0.64)rectangle++(0.04,0.04) (0.07,0.742)rectangle++(0.04,0.04)
  (1.996,1.834)rectangle++(0.04,0.04) (1.387,2.221)rectangle++(0.04,0.04) (2.289,3.016)rectangle++(0.04,0.04) (1.585,1.88)rectangle++(0.04,0.04) (2.302,2.632)rectangle++(0.04,0.04)
  (2.373,3.498)rectangle++(0.04,0.04) (2.378,3.022)rectangle++(0.04,0.04) (2.301,3.039)rectangle++(0.04,0.04) (0.816,1.95)rectangle++(0.04,0.04) (0.769,1.403)rectangle++(0.04,0.04)
  (1.836,1.56)rectangle++(0.04,0.04) (2.419,3.323)rectangle++(0.04,0.04) (0.55,1.856)rectangle++(0.04,0.04) (1.26,1.418)rectangle++(0.04,0.04) (0.912,1.64)rectangle++(0.04,0.04)
  (1.202,0.777)rectangle++(0.04,0.04) (1.298,0.841)rectangle++(0.04,0.04) (1.834,2.513)rectangle++(0.04,0.04) (2.385,3.48)rectangle++(0.04,0.04) (2,1.883)rectangle++(0.04,0.04)
  (2.379,3.469)rectangle++(0.04,0.04) (2.234,3.178)rectangle++(0.04,0.04) (1.537,2.448)rectangle++(0.04,0.04) (1.152,1.249)rectangle++(0.04,0.04) (0.893,1.59)rectangle++(0.04,0.04)
  (2.366,2.971)rectangle++(0.04,0.04) (1.409,0.838)rectangle++(0.04,0.04) (1.32,2.008)rectangle++(0.04,0.04) (2.405,3.34)rectangle++(0.04,0.04) (1.926,2.825)rectangle++(0.04,0.04)
  (1.926,2.747)rectangle++(0.04,0.04) (1.082,2.008)rectangle++(0.04,0.04) (2.164,2.939)rectangle++(0.04,0.04) (2.191,3.092)rectangle++(0.04,0.04) (0.582,0.947)rectangle++(0.04,0.04)
  (0.908,1.434)rectangle++(0.04,0.04) (2.401,3.304)rectangle++(0.04,0.04) (0.548,1.847)rectangle++(0.04,0.04) (1.801,2.258)rectangle++(0.04,0.04) (0.345,1.44)rectangle++(0.04,0.04)
  (2.391,3.443)rectangle++(0.04,0.04) (1.955,2.146)rectangle++(0.04,0.04) (1.393,1.364)rectangle++(0.04,0.04) (1.39,2.085)rectangle++(0.04,0.04) (2.374,3.498)rectangle++(0.04,0.04)
  (1.034,0.943)rectangle++(0.04,0.04) (0.505,1.233)rectangle++(0.04,0.04) (0.252,0.857)rectangle++(0.04,0.04) (1.419,0.806)rectangle++(0.04,0.04) (0.775,1.551)rectangle++(0.04,0.04)
  (0.682,1.309)rectangle++(0.04,0.04) (0.639,1.237)rectangle++(0.04,0.04) (1.228,1.755)rectangle++(0.04,0.04) (2.262,2.999)rectangle++(0.04,0.04) (0.812,1.934)rectangle++(0.04,0.04)
  (1.268,1.512)rectangle++(0.04,0.04) (1.514,0.915)rectangle++(0.04,0.04) (2.251,3.107)rectangle++(0.04,0.04) (1.11,2.132)rectangle++(0.04,0.04) (2.387,3.488)rectangle++(0.04,0.04)
  (0.601,1.888)rectangle++(0.04,0.04) (2.299,3.14)rectangle++(0.04,0.04) (1.859,2.336)rectangle++(0.04,0.04) (2.081,2.226)rectangle++(0.04,0.04) (1.429,2.289)rectangle++(0.04,0.04)
  (1.22,1.574)rectangle++(0.04,0.04) (1.838,1.55)rectangle++(0.04,0.04) (1.557,2.33)rectangle++(0.04,0.04) (1.049,1.865)rectangle++(0.04,0.04) (2.02,2.726)rectangle++(0.04,0.04)
  (1.61,1.564)rectangle++(0.04,0.04) (0.179,1.224)rectangle++(0.04,0.04) (0.298,0.77)rectangle++(0.04,0.04) (1.531,1.142)rectangle++(0.04,0.04) (1.717,1.668)rectangle++(0.04,0.04)
  (1.734,2.486)rectangle++(0.04,0.04) (0.766,1.712)rectangle++(0.04,0.04) (2.409,3.432)rectangle++(0.04,0.04) (1.425,0.896)rectangle++(0.04,0.04) (1.083,1.806)rectangle++(0.04,0.04)
  (2.375,3.39)rectangle++(0.04,0.04) (2.018,2.916)rectangle++(0.04,0.04) (1.691,1.233)rectangle++(0.04,0.04) (1.832,2.574)rectangle++(0.04,0.04) (2.295,3.227)rectangle++(0.04,0.04)
  (1.537,2.462)rectangle++(0.04,0.04) (1.317,2.086)rectangle++(0.04,0.04) (1.033,1.777)rectangle++(0.04,0.04) (1.868,1.566)rectangle++(0.04,0.04) (1.26,1.596)rectangle++(0.04,0.04)
  (2.385,2.973)rectangle++(0.04,0.04) (1.817,1.919)rectangle++(0.04,0.04) (1.199,2.033)rectangle++(0.04,0.04) (2.209,3.023)rectangle++(0.04,0.04) (2.375,3.178)rectangle++(0.04,0.04)
  (1.909,2.093)rectangle++(0.04,0.04) (1.976,2.782)rectangle++(0.04,0.04) (0.801,1.827)rectangle++(0.04,0.04) (0.743,1.374)rectangle++(0.04,0.04) (2.272,2.799)rectangle++(0.04,0.04)
  (1.896,2.086)rectangle++(0.04,0.04) (2.23,2.743)rectangle++(0.04,0.04) (1.098,1.332)rectangle++(0.04,0.04) (2.103,2.703)rectangle++(0.04,0.04) (0.775,0.991)rectangle++(0.04,0.04)
  (0.552,1.189)rectangle++(0.04,0.04) (0.958,1.476)rectangle++(0.04,0.04) (0.252,1.033)rectangle++(0.04,0.04) (0.635,1.087)rectangle++(0.04,0.04) (1.621,2.164)rectangle++(0.04,0.04)
  (1.469,2.229)rectangle++(0.04,0.04) (2.426,3.14)rectangle++(0.04,0.04) (1.525,1.599)rectangle++(0.04,0.04) (2.267,2.747)rectangle++(0.04,0.04) (2.131,2.377)rectangle++(0.04,0.04)
  (1.597,1.566)rectangle++(0.04,0.04) (0.898,1.726)rectangle++(0.04,0.04) (1.354,0.848)rectangle++(0.04,0.04) (1.694,2.366)rectangle++(0.04,0.04) (1.039,1.895)rectangle++(0.04,0.04)
  (2.199,2.552)rectangle++(0.04,0.04) (2.305,3.268)rectangle++(0.04,0.04) (1.21,0.969)rectangle++(0.04,0.04) (1.306,1.531)rectangle++(0.04,0.04) (0.928,1.67)rectangle++(0.04,0.04)
  (2.188,2.461)rectangle++(0.04,0.04) (2.127,2.554)rectangle++(0.04,0.04) (2.415,3.067)rectangle++(0.04,0.04) (1.822,1.72)rectangle++(0.04,0.04) (1.845,1.789)rectangle++(0.04,0.04)
  (0.645,1.004)rectangle++(0.04,0.04) (2.405,3.458)rectangle++(0.04,0.04) (2.4,3.391)rectangle++(0.04,0.04) (1.364,2.362)rectangle++(0.04,0.04) (0.789,1.607)rectangle++(0.04,0.04)
  (2.288,3.13)rectangle++(0.04,0.04) (2.353,2.786)rectangle++(0.04,0.04) (0.792,1.017)rectangle++(0.04,0.04) (1.597,1.447)rectangle++(0.04,0.04) (2.237,2.808)rectangle++(0.04,0.04)
  (2.216,2.856)rectangle++(0.04,0.04) (2.26,3.196)rectangle++(0.04,0.04) (1.126,2.151)rectangle++(0.04,0.04) (1.52,2.173)rectangle++(0.04,0.04) (0.363,1.346)rectangle++(0.04,0.04)
  (1.674,2.174)rectangle++(0.04,0.04) (1.465,2.237)rectangle++(0.04,0.04) (0.746,1.602)rectangle++(0.04,0.04) (1.273,1.832)rectangle++(0.04,0.04) (2.003,2.733)rectangle++(0.04,0.04)
  (2.005,2.842)rectangle++(0.04,0.04) (0.478,1.634)rectangle++(0.04,0.04) (1.053,1.548)rectangle++(0.04,0.04) (2.202,2.992)rectangle++(0.04,0.04) (2.289,2.533)rectangle++(0.04,0.04)
  (1.902,1.779)rectangle++(0.04,0.04) (2.263,2.69)rectangle++(0.04,0.04) (1.692,1.57)rectangle++(0.04,0.04) (0.169,1.247)rectangle++(0.04,0.04) (2.327,3.169)rectangle++(0.04,0.04)
  (1.553,1.594)rectangle++(0.04,0.04) (0.908,1.724)rectangle++(0.04,0.04) (1.015,1.638)rectangle++(0.04,0.04) (2.367,3.044)rectangle++(0.04,0.04) (2.117,2.246)rectangle++(0.04,0.04)
  (2.371,2.946)rectangle++(0.04,0.04) (2.407,3.462)rectangle++(0.04,0.04) (0.899,2.064)rectangle++(0.04,0.04) (1.328,1.931)rectangle++(0.04,0.04) (1.184,0.719)rectangle++(0.04,0.04)
  (1.53,2.164)rectangle++(0.04,0.04) (1.793,2.53)rectangle++(0.04,0.04) (2.212,2.333)rectangle++(0.04,0.04) (2.338,3.049)rectangle++(0.04,0.04) (2.186,3.102)rectangle++(0.04,0.04)
  (2.014,2.882)rectangle++(0.04,0.04) (0.488,1.645)rectangle++(0.04,0.04) (1.31,1.803)rectangle++(0.04,0.04) (0.98,1.722)rectangle++(0.04,0.04) (1.479,2.093)rectangle++(0.04,0.04)
  (2.349,3.321)rectangle++(0.04,0.04) (2.216,2.663)rectangle++(0.04,0.04) (1.351,1.096)rectangle++(0.04,0.04) (1.779,1.809)rectangle++(0.04,0.04) (0.921,1.804)rectangle++(0.04,0.04)
  (1.529,1.186)rectangle++(0.04,0.04) (1.542,1.817)rectangle++(0.04,0.04) (1.762,2.491)rectangle++(0.04,0.04) (1.613,1.446)rectangle++(0.04,0.04) (2.173,2.997)rectangle++(0.04,0.04)
  (1.519,2.414)rectangle++(0.04,0.04) (2.398,3.385)rectangle++(0.04,0.04) (2.224,2.649)rectangle++(0.04,0.04) (0.391,1.644)rectangle++(0.04,0.04) (2.187,2.898)rectangle++(0.04,0.04)
  (0.898,0.958)rectangle++(0.04,0.04) (1.628,2.017)rectangle++(0.04,0.04) (2.308,2.897)rectangle++(0.04,0.04) (2.229,2.632)rectangle++(0.04,0.04) (2.119,2.127)rectangle++(0.04,0.04)
  (1.512,0.928)rectangle++(0.04,0.04) (2.341,3.295)rectangle++(0.04,0.04) (1.694,2.603)rectangle++(0.04,0.04) (0.457,1.495)rectangle++(0.04,0.04) (1.29,1.774)rectangle++(0.04,0.04)
  (1.268,1.318)rectangle++(0.04,0.04) (2.318,3.023)rectangle++(0.04,0.04) (2.327,3.049)rectangle++(0.04,0.04) (2.334,2.884)rectangle++(0.04,0.04) (1.292,2.267)rectangle++(0.04,0.04)
  (2.376,3.499)rectangle++(0.04,0.04) (2.326,3.328)rectangle++(0.04,0.04) (2.14,2.144)rectangle++(0.04,0.04) (0.929,0.929)rectangle++(0.04,0.04) (1.7,1.262)rectangle++(0.04,0.04)
  (2.164,2.988)rectangle++(0.04,0.04) (2.376,3.461)rectangle++(0.04,0.04) (0.595,1.878)rectangle++(0.04,0.04) (2.27,3.076)rectangle++(0.04,0.04) (1.725,2.159)rectangle++(0.04,0.04)
  (0.328,1.397)rectangle++(0.04,0.04) (2.367,3.039)rectangle++(0.04,0.04) (2.225,3.157)rectangle++(0.04,0.04) (2.051,2.467)rectangle++(0.04,0.04) (2.114,2.608)rectangle++(0.04,0.04)
  (1.173,0.882)rectangle++(0.04,0.04) (2.046,2.763)rectangle++(0.04,0.04) (1.571,1.133)rectangle++(0.04,0.04) (2.095,2.409)rectangle++(0.04,0.04) (1.291,1.65)rectangle++(0.04,0.04)
  (1.281,0.844)rectangle++(0.04,0.04) (1.813,1.719)rectangle++(0.04,0.04) (1.168,1.999)rectangle++(0.04,0.04) (1.198,0.835)rectangle++(0.04,0.04) (0.174,0.586)rectangle++(0.04,0.04)
  (1.612,2.058)rectangle++(0.04,0.04) (1.808,1.918)rectangle++(0.04,0.04) (2.428,3.15)rectangle++(0.04,0.04) (1.849,2.108)rectangle++(0.04,0.04) (0.663,1.655)rectangle++(0.04,0.04)
  (2.089,2.072)rectangle++(0.04,0.04) (2.018,1.869)rectangle++(0.04,0.04) (1.522,1.911)rectangle++(0.04,0.04) (2.102,2.083)rectangle++(0.04,0.04) (1.409,2.27)rectangle++(0.04,0.04)
  (2.371,2.873)rectangle++(0.04,0.04) (2.228,2.928)rectangle++(0.04,0.04) (2.138,3.028)rectangle++(0.04,0.04) (1.252,1.609)rectangle++(0.04,0.04) (1.891,2.588)rectangle++(0.04,0.04)
  (1.305,1.665)rectangle++(0.04,0.04) (0.253,1.169)rectangle++(0.04,0.04) (0.992,1.421)rectangle++(0.04,0.04) (0.622,1.327)rectangle++(0.04,0.04) (2.371,3.501)rectangle++(0.04,0.04)
  (2.405,3.378)rectangle++(0.04,0.04) (2.193,3.122)rectangle++(0.04,0.04) (1.534,2.435)rectangle++(0.04,0.04) (2.148,2.721)rectangle++(0.04,0.04) (2.319,2.666)rectangle++(0.04,0.04)
  (2.287,2.734)rectangle++(0.04,0.04) (1.473,2.388)rectangle++(0.04,0.04) (1.673,1.619)rectangle++(0.04,0.04) (1.916,2.416)rectangle++(0.04,0.04) (0.725,1.735)rectangle++(0.04,0.04)
  (1.297,1.851)rectangle++(0.04,0.04) (1.837,1.581)rectangle++(0.04,0.04) (2.138,2.379)rectangle++(0.04,0.04) (1.071,0.977)rectangle++(0.04,0.04) (1.137,1.771)rectangle++(0.04,0.04)
  (1.646,2.298)rectangle++(0.04,0.04) (2.317,3.053)rectangle++(0.04,0.04) (2.228,3.157)rectangle++(0.04,0.04) (1.347,2.299)rectangle++(0.04,0.04) (2.305,2.72)rectangle++(0.04,0.04)
  (1.269,2.235)rectangle++(0.04,0.04) (1.068,1.798)rectangle++(0.04,0.04) (1.009,1.68)rectangle++(0.04,0.04) (2.073,2.055)rectangle++(0.04,0.04) (0.619,1.691)rectangle++(0.04,0.04)
  (2.267,2.495)rectangle++(0.04,0.04) (1.048,0.679)rectangle++(0.04,0.04) (0.677,0.872)rectangle++(0.04,0.04) (1.915,2.095)rectangle++(0.04,0.04) (1.219,1.26)rectangle++(0.04,0.04)
  (2.06,2.794)rectangle++(0.04,0.04) (2.381,3.499)rectangle++(0.04,0.04) (1.678,1.873)rectangle++(0.04,0.04) (0.255,1.316)rectangle++(0.04,0.04) (2.396,3.476)rectangle++(0.04,0.04)
  (0.596,1.887)rectangle++(0.04,0.04) (1.52,1.66)rectangle++(0.04,0.04) (1.673,2.055)rectangle++(0.04,0.04) (2.054,2.219)rectangle++(0.04,0.04) (2.306,3.257)rectangle++(0.04,0.04)
  (1.952,2.297)rectangle++(0.04,0.04) (2.336,3.104)rectangle++(0.04,0.04) (2.397,3.203)rectangle++(0.04,0.04) (2.068,2.536)rectangle++(0.04,0.04) (1.471,2.334)rectangle++(0.04,0.04)
  (0.416,1.373)rectangle++(0.04,0.04) (0.305,0.867)rectangle++(0.04,0.04) (2.371,3.501)rectangle++(0.04,0.04) (1.713,2.633)rectangle++(0.04,0.04) (1.928,2.716)rectangle++(0.04,0.04)
  (1.493,1.679)rectangle++(0.04,0.04) (1.847,1.587)rectangle++(0.04,0.04) (0.542,1.543)rectangle++(0.04,0.04) (0.333,0.941)rectangle++(0.04,0.04) (0.196,0.743)rectangle++(0.04,0.04)
  (1.434,1.86)rectangle++(0.04,0.04) (2.381,3.446)rectangle++(0.04,0.04) (1.879,2.097)rectangle++(0.04,0.04) (0.288,1.422)rectangle++(0.04,0.04) (2.171,2.82)rectangle++(0.04,0.04)
  (2.326,3.31)rectangle++(0.04,0.04) (2.371,2.809)rectangle++(0.04,0.04) (1.64,2.542)rectangle++(0.04,0.04) (2.081,2.104)rectangle++(0.04,0.04) (0.26,1.476)rectangle++(0.04,0.04)
  (2.369,3.307)rectangle++(0.04,0.04) (1.873,1.856)rectangle++(0.04,0.04) (1.403,2.211)rectangle++(0.04,0.04) (2.029,1.884)rectangle++(0.04,0.04) (0.586,1.645)rectangle++(0.04,0.04)
  (2.416,3.441)rectangle++(0.04,0.04) (2.361,3.41)rectangle++(0.04,0.04) (1.212,0.897)rectangle++(0.04,0.04) (2.437,3.267)rectangle++(0.04,0.04) (0.85,2.028)rectangle++(0.04,0.04)
  (1.081,1.686)rectangle++(0.04,0.04) (2.394,3.455)rectangle++(0.04,0.04) (0.591,1.882)rectangle++(0.04,0.04) (2.208,2.467)rectangle++(0.04,0.04) (2.384,3.356)rectangle++(0.04,0.04)
  (2.015,2.913)rectangle++(0.04,0.04) (0.496,1.653)rectangle++(0.04,0.04) (2.015,2.031)rectangle++(0.04,0.04) (2.182,2.841)rectangle++(0.04,0.04) (0.786,1.883)rectangle++(0.04,0.04)
  (0.768,0.526)rectangle++(0.04,0.04) (1.101,1.653)rectangle++(0.04,0.04) (2.419,3.197)rectangle++(0.04,0.04) (1.463,0.905)rectangle++(0.04,0.04) (2.021,2.429)rectangle++(0.04,0.04)
  (0.713,1.76)rectangle++(0.04,0.04) (2.354,3.012)rectangle++(0.04,0.04) (2.235,2.756)rectangle++(0.04,0.04) (1.898,2.757)rectangle++(0.04,0.04) (1.526,1.704)rectangle++(0.04,0.04)
  (1.697,1.265)rectangle++(0.04,0.04) (2.325,3.267)rectangle++(0.04,0.04) (1.989,2.36)rectangle++(0.04,0.04) (2.13,2.984)rectangle++(0.04,0.04) (0.499,1.7)rectangle++(0.04,0.04)
  (2.126,2.685)rectangle++(0.04,0.04) (1.842,1.592)rectangle++(0.04,0.04) (0.647,0.618)rectangle++(0.04,0.04) (0.484,1.082)rectangle++(0.04,0.04) (2.161,2.863)rectangle++(0.04,0.04)
  (2.31,3.278)rectangle++(0.04,0.04) (1.089,0.918)rectangle++(0.04,0.04) (1.126,1.763)rectangle++(0.04,0.04) (0.683,1.429)rectangle++(0.04,0.04) (2.394,3.483)rectangle++(0.04,0.04)
  (2.319,2.828)rectangle++(0.04,0.04) (1.442,0.844)rectangle++(0.04,0.04) (1.069,0.946)rectangle++(0.04,0.04) (1.864,1.756)rectangle++(0.04,0.04) (0.797,0.752)rectangle++(0.04,0.04)
  (2.274,2.707)rectangle++(0.04,0.04) (2.064,2.938)rectangle++(0.04,0.04) (0.825,1.879)rectangle++(0.04,0.04) (2.325,2.973)rectangle++(0.04,0.04) (2.306,3.282)rectangle++(0.04,0.04)
  (0.556,1.817)rectangle++(0.04,0.04) (1.078,1.596)rectangle++(0.04,0.04) (1.234,1.876)rectangle++(0.04,0.04) (2.131,2.749)rectangle++(0.04,0.04) (1.557,1.248)rectangle++(0.04,0.04)
  (0.666,1.024)rectangle++(0.04,0.04) (2.081,2.501)rectangle++(0.04,0.04) (1.982,2.458)rectangle++(0.04,0.04) (0.725,1.759)rectangle++(0.04,0.04) (1.045,1.614)rectangle++(0.04,0.04)
  (1.953,1.792)rectangle++(0.04,0.04) (2.018,2.054)rectangle++(0.04,0.04) (0.255,1.448)rectangle++(0.04,0.04) (1.97,2.04)rectangle++(0.04,0.04) (1.561,1.422)rectangle++(0.04,0.04)
  (1.15,1.914)rectangle++(0.04,0.04) (0.346,1.188)rectangle++(0.04,0.04) (2.37,2.89)rectangle++(0.04,0.04) (1.434,1.81)rectangle++(0.04,0.04) (1.443,2.156)rectangle++(0.04,0.04)
  (2.276,2.607)rectangle++(0.04,0.04) (1.987,2.467)rectangle++(0.04,0.04) (2.419,3.304)rectangle++(0.04,0.04) (2.224,2.353)rectangle++(0.04,0.04) (2.13,2.135)rectangle++(0.04,0.04)
  (1.963,2.795)rectangle++(0.04,0.04) (1.159,1.035)rectangle++(0.04,0.04) (2.333,3.028)rectangle++(0.04,0.04) (1.507,2.437)rectangle++(0.04,0.04) (1.818,2.559)rectangle++(0.04,0.04)
  (0.427,1.517)rectangle++(0.04,0.04) (0.637,0.997)rectangle++(0.04,0.04) (2.395,3.456)rectangle++(0.04,0.04) (0.861,1.011)rectangle++(0.04,0.04) (1.167,1.748)rectangle++(0.04,0.04)
  (2.307,2.967)rectangle++(0.04,0.04) (1.119,1.337)rectangle++(0.04,0.04) (2.328,3.154)rectangle++(0.04,0.04) (0.517,1.792)rectangle++(0.04,0.04) (0.407,0.994)rectangle++(0.04,0.04)
  (2.103,2.421)rectangle++(0.04,0.04) (1.986,2.085)rectangle++(0.04,0.04) (2.014,2.211)rectangle++(0.04,0.04) (0.968,1.921)rectangle++(0.04,0.04) (1.819,2.463)rectangle++(0.04,0.04)
  (2.002,2.811)rectangle++(0.04,0.04) (1.809,2.637)rectangle++(0.04,0.04) (2.416,3.144)rectangle++(0.04,0.04) (1.819,2.102)rectangle++(0.04,0.04) (1.883,2.668)rectangle++(0.04,0.04)
  (1.725,2.158)rectangle++(0.04,0.04) (2.331,3.291)rectangle++(0.04,0.04) (2.429,3.211)rectangle++(0.04,0.04) (1.107,2.182)rectangle++(0.04,0.04) (0.428,1.242)rectangle++(0.04,0.04)
  (2.414,3.225)rectangle++(0.04,0.04) (1.718,1.865)rectangle++(0.04,0.04) (2.157,2.254)rectangle++(0.04,0.04) (2.05,2.9)rectangle++(0.04,0.04) (0.818,1.868)rectangle++(0.04,0.04)
  (1.894,2.385)rectangle++(0.04,0.04) (2.229,3.119)rectangle++(0.04,0.04) (1.757,2.174)rectangle++(0.04,0.04) (0.69,1.65)rectangle++(0.04,0.04) (1.663,1.518)rectangle++(0.04,0.04)
  (1.726,2.465)rectangle++(0.04,0.04) (1.783,2.065)rectangle++(0.04,0.04) (2.063,2.637)rectangle++(0.04,0.04) (1.173,1.524)rectangle++(0.04,0.04) (0.233,1.101)rectangle++(0.04,0.04)
  (2.142,2.14)rectangle++(0.04,0.04) (1.914,2.078)rectangle++(0.04,0.04) (2.39,3.275)rectangle++(0.04,0.04) (1.465,1.262)rectangle++(0.04,0.04) (0.519,1.393)rectangle++(0.04,0.04)
  (2.424,3.085)rectangle++(0.04,0.04) (0.391,0.773)rectangle++(0.04,0.04) (2.233,2.808)rectangle++(0.04,0.04) (2.431,3.3)rectangle++(0.04,0.04) (1.346,2.351)rectangle++(0.04,0.04)
  (2.388,3.387)rectangle++(0.04,0.04) (1.678,1.58)rectangle++(0.04,0.04) (1.352,1.277)rectangle++(0.04,0.04) (1.731,2.415)rectangle++(0.04,0.04) (2.082,2.898)rectangle++(0.04,0.04)
  (1.135,1.131)rectangle++(0.04,0.04) (1.893,1.923)rectangle++(0.04,0.04) (0.928,1.845)rectangle++(0.04,0.04) (1.036,1.665)rectangle++(0.04,0.04) (0.988,1.649)rectangle++(0.04,0.04)
  (1.894,1.61)rectangle++(0.04,0.04) (1.558,2.343)rectangle++(0.04,0.04) (1.982,1.972)rectangle++(0.04,0.04) (2.048,2.873)rectangle++(0.04,0.04) (1.091,2.052)rectangle++(0.04,0.04)
  (1.82,2.486)rectangle++(0.04,0.04) (2.367,3.398)rectangle++(0.04,0.04) (1.601,1.644)rectangle++(0.04,0.04) (1.634,2.046)rectangle++(0.04,0.04) (0.678,1.597)rectangle++(0.04,0.04)
  (1.495,2.041)rectangle++(0.04,0.04) (2.369,3.049)rectangle++(0.04,0.04) (2.41,3.177)rectangle++(0.04,0.04) (1.103,2.173)rectangle++(0.04,0.04) (1.714,2.145)rectangle++(0.04,0.04)
  (2.331,3.291)rectangle++(0.04,0.04) (2.276,2.647)rectangle++(0.04,0.04) (1.1,1.002)rectangle++(0.04,0.04) (1.432,0.862)rectangle++(0.04,0.04) (2.374,3.396)rectangle++(0.04,0.04)
  (1.368,2.359)rectangle++(0.04,0.04) (0.787,1.607)rectangle++(0.04,0.04) (1.723,2.156)rectangle++(0.04,0.04) (1.632,2.4)rectangle++(0.04,0.04) (2.273,3.172)rectangle++(0.04,0.04)
  (1.119,2.148)rectangle++(0.04,0.04) (2.25,2.924)rectangle++(0.04,0.04) (1.308,2.262)rectangle++(0.04,0.04) (1.802,1.922)rectangle++(0.04,0.04) (1.872,2.649)rectangle++(0.04,0.04)
  (1.07,1.98)rectangle++(0.04,0.04) (0.741,1.463)rectangle++(0.04,0.04) (2.348,3.271)rectangle++(0.04,0.04) (1.351,2.334)rectangle++(0.04,0.04) (2.179,2.835)rectangle++(0.04,0.04)
  (0.45,1.677)rectangle++(0.04,0.04) (1.956,2.551)rectangle++(0.04,0.04) (1.852,1.594)rectangle++(0.04,0.04) (0.543,1.546)rectangle++(0.04,0.04) (1.028,1.542)rectangle++(0.04,0.04)
  (1.498,1.513)rectangle++(0.04,0.04) (1.866,2.592)rectangle++(0.04,0.04) (0.428,1.537)rectangle++(0.04,0.04) (1.513,2.004)rectangle++(0.04,0.04) (2.381,3.135)rectangle++(0.04,0.04)
  (0.503,1.801)rectangle++(0.04,0.04) (1.333,1.832)rectangle++(0.04,0.04) (1.8,1.536)rectangle++(0.04,0.04) (0.143,1.267)rectangle++(0.04,0.04) (2.079,2.544)rectangle++(0.04,0.04)
  (2.433,3.291)rectangle++(0.04,0.04) (1.344,2.35)rectangle++(0.04,0.04) (1.688,2.392)rectangle++(0.04,0.04) (1.66,2.427)rectangle++(0.04,0.04) (0.413,1.444)rectangle++(0.04,0.04)
  (1.025,1.5)rectangle++(0.04,0.04) (0.958,1.616)rectangle++(0.04,0.04) (1.423,0.871)rectangle++(0.04,0.04) (1.524,2.205)rectangle++(0.04,0.04) (1.468,1.222)rectangle++(0.04,0.04)
  (1.956,2.687)rectangle++(0.04,0.04) (1.607,1.157)rectangle++(0.04,0.04) (1.5,1.885)rectangle++(0.04,0.04) (1.813,2.051)rectangle++(0.04,0.04) (2.061,2.873)rectangle++(0.04,0.04)
  (2.084,2.944)rectangle++(0.04,0.04) (2.146,3.025)rectangle++(0.04,0.04) (2.012,2.874)rectangle++(0.04,0.04) (1.745,2.087)rectangle++(0.04,0.04) (1.188,0.969)rectangle++(0.04,0.04)
  (2.025,2.031)rectangle++(0.04,0.04) (2.116,2.963)rectangle++(0.04,0.04) (1.447,1.593)rectangle++(0.04,0.04) (2.346,3.31)rectangle++(0.04,0.04) (1.823,2.718)rectangle++(0.04,0.04)
  (2.038,1.9)rectangle++(0.04,0.04) (2.431,3.336)rectangle++(0.04,0.04) (0.867,2.042)rectangle++(0.04,0.04) (2.421,3.242)rectangle++(0.04,0.04) (1.356,0.888)rectangle++(0.04,0.04)
  (2.217,3.038)rectangle++(0.04,0.04) (2.36,2.938)rectangle++(0.04,0.04) (2.398,3.197)rectangle++(0.04,0.04) (1.916,2.81)rectangle++(0.04,0.04) (2.016,2.845)rectangle++(0.04,0.04)
  (1.811,2.642)rectangle++(0.04,0.04) (0.906,0.62)rectangle++(0.04,0.04) (2.329,2.745)rectangle++(0.04,0.04) (1.27,2.243)rectangle++(0.04,0.04) (2.407,3.235)rectangle++(0.04,0.04)
  (0.847,2.015)rectangle++(0.04,0.04) (2.399,3.468)rectangle++(0.04,0.04) (2.195,3.127)rectangle++(0.04,0.04) (0.529,1.751)rectangle++(0.04,0.04) (2.286,2.803)rectangle++(0.04,0.04)
  (1.734,1.871)rectangle++(0.04,0.04) (2.058,2.857)rectangle++(0.04,0.04) (1.086,2.051)rectangle++(0.04,0.04) (1.17,0.657)rectangle++(0.04,0.04) (2.249,2.922)rectangle++(0.04,0.04)
  (2.367,3.419)rectangle++(0.04,0.04) (1.705,2.622)rectangle++(0.04,0.04) (1.52,2.307)rectangle++(0.04,0.04) (0.401,1.379)rectangle++(0.04,0.04) (1.107,0.793)rectangle++(0.04,0.04)
  (1.71,2.343)rectangle++(0.04,0.04) (1.653,2.423)rectangle++(0.04,0.04) (2.087,2.896)rectangle++(0.04,0.04) (0.48,1.668)rectangle++(0.04,0.04) (1.432,0.831)rectangle++(0.04,0.04)
  (0.778,1.558)rectangle++(0.04,0.04) (2.419,3.429)rectangle++(0.04,0.04) (1.337,1.28)rectangle++(0.04,0.04) (1.154,1.859)rectangle++(0.04,0.04) (1.654,2.311)rectangle++(0.04,0.04)
  (2.144,2.723)rectangle++(0.04,0.04) (1.996,2.847)rectangle++(0.04,0.04) (2.375,3.443)rectangle++(0.04,0.04) (1.73,2.164)rectangle++(0.04,0.04) (1.248,2.062)rectangle++(0.04,0.04)
  (1.742,1.454)rectangle++(0.04,0.04) (1.135,1.944)rectangle++(0.04,0.04) (1.665,2.32)rectangle++(0.04,0.04) (2.234,2.49)rectangle++(0.04,0.04) (0.345,1.608)rectangle++(0.04,0.04)
  (2.191,2.585)rectangle++(0.04,0.04) (2.339,3.109)rectangle++(0.04,0.04) (1.675,2.579)rectangle++(0.04,0.04) (1.069,1.282)rectangle++(0.04,0.04) (1.601,2.227)rectangle++(0.04,0.04)
  (0.366,1.381)rectangle++(0.04,0.04) (2.09,2.361)rectangle++(0.04,0.04) (2.059,2.906)rectangle++(0.04,0.04) (1.096,2.061)rectangle++(0.04,0.04) (1.711,1.434)rectangle++(0.04,0.04)
  (1.325,1.687)rectangle++(0.04,0.04) (2.232,2.56)rectangle++(0.04,0.04) (0.722,0.618)rectangle++(0.04,0.04) (1.759,2.011)rectangle++(0.04,0.04) (2.129,2.955)rectangle++(0.04,0.04)
  (1.81,2.662)rectangle++(0.04,0.04) (2.295,2.763)rectangle++(0.04,0.04) (0.412,1.69)rectangle++(0.04,0.04) (0.394,0.943)rectangle++(0.04,0.04) (0.779,0.637)rectangle++(0.04,0.04)
  (0.683,0.628)rectangle++(0.04,0.04) (2.274,3.088)rectangle++(0.04,0.04) (1.735,1.312)rectangle++(0.04,0.04) (0.831,1.701)rectangle++(0.04,0.04) (1.812,2.428)rectangle++(0.04,0.04)
  (1.484,2.29)rectangle++(0.04,0.04) (1.296,2.05)rectangle++(0.04,0.04) (1.028,1.767)rectangle++(0.04,0.04) (2.432,3.252)rectangle++(0.04,0.04) (2.127,2.742)rectangle++(0.04,0.04)
  (1.693,1.848)rectangle++(0.04,0.04) (1.92,1.845)rectangle++(0.04,0.04) (1.067,0.647)rectangle++(0.04,0.04) (1.95,2.62)rectangle++(0.04,0.04) (0.424,1.566)rectangle++(0.04,0.04)
  (1.166,0.695)rectangle++(0.04,0.04) (0.785,1.486)rectangle++(0.04,0.04) (1.334,1.503)rectangle++(0.04,0.04) (1.426,1.531)rectangle++(0.04,0.04) (2.359,3.247)rectangle++(0.04,0.04)
  (1.123,2.177)rectangle++(0.04,0.04) (2.371,3.501)rectangle++(0.04,0.04) (1.903,1.623)rectangle++(0.04,0.04) (1.142,1.995)rectangle++(0.04,0.04) (0.37,1.206)rectangle++(0.04,0.04)
  (2.053,2.621)rectangle++(0.04,0.04) (2.426,3.189)rectangle++(0.04,0.04) (2.186,3.111)rectangle++(0.04,0.04) (1.345,2.285)rectangle++(0.04,0.04) (1.932,2.661)rectangle++(0.04,0.04)
  (2.401,3.426)rectangle++(0.04,0.04) (1.597,1.388)rectangle++(0.04,0.04) (2.432,3.299)rectangle++(0.04,0.04) (1.346,2.351)rectangle++(0.04,0.04) (2.393,3.226)rectangle++(0.04,0.04)
  (1.963,1.83)rectangle++(0.04,0.04) (0.582,1.62)rectangle++(0.04,0.04) (2.037,2.685)rectangle++(0.04,0.04) (1.284,0.753)rectangle++(0.04,0.04) (2.179,2.788)rectangle++(0.04,0.04)
  (1.146,1.078)rectangle++(0.04,0.04) (2.149,2.586)rectangle++(0.04,0.04) (1.489,1.088)rectangle++(0.04,0.04) (1.344,2.054)rectangle++(0.04,0.04) (1.21,1.222)rectangle++(0.04,0.04)
  (2.238,2.488)rectangle++(0.04,0.04) (1.447,2.346)rectangle++(0.04,0.04) (0.422,1.369)rectangle++(0.04,0.04) (2.123,2.563)rectangle++(0.04,0.04) (1.469,2.344)rectangle++(0.04,0.04)
  (1.998,2.187)rectangle++(0.04,0.04) (0.295,1.474)rectangle++(0.04,0.04) (2.229,2.585)rectangle++(0.04,0.04) (0.72,1.838)rectangle++(0.04,0.04) (2.176,2.919)rectangle++(0.04,0.04)
  (1.509,2.403)rectangle++(0.04,0.04) (2.194,2.863)rectangle++(0.04,0.04) (1.376,1.289)rectangle++(0.04,0.04) (1.573,2.255)rectangle++(0.04,0.04) (1.525,1.685)rectangle++(0.04,0.04)
  (1.661,1.737)rectangle++(0.04,0.04) (1.366,1.452)rectangle++(0.04,0.04) (0.184,1.134)rectangle++(0.04,0.04) (2.408,3.41)rectangle++(0.04,0.04) (1.947,2.461)rectangle++(0.04,0.04)
  (1.019,1.104)rectangle++(0.04,0.04) (1.587,2.2)rectangle++(0.04,0.04) (2.255,2.629)rectangle++(0.04,0.04) (1.687,1.25)rectangle++(0.04,0.04) (0.08,1.169)rectangle++(0.04,0.04)
  (2.045,2.624)rectangle++(0.04,0.04) (1.651,2.489)rectangle++(0.04,0.04) (1.068,1.911)rectangle++(0.04,0.04) (1.422,1.352)rectangle++(0.04,0.04) (2.39,3.47)rectangle++(0.04,0.04)
  (1.014,0.765)rectangle++(0.04,0.04) (1.716,2.333)rectangle++(0.04,0.04) (1.792,2.562)rectangle++(0.04,0.04) (2.237,3.126)rectangle++(0.04,0.04) (1.115,2.133)rectangle++(0.04,0.04)
  (1.825,2.497)rectangle++(0.04,0.04) (2.243,2.628)rectangle++(0.04,0.04) (1.262,2.211)rectangle++(0.04,0.04) (1.52,1.682)rectangle++(0.04,0.04) (2.264,3.14)rectangle++(0.04,0.04)
  (1.941,2.003)rectangle++(0.04,0.04) (1.595,2.399)rectangle++(0.04,0.04) (1.334,1.045)rectangle++(0.04,0.04) (1.923,2.052)rectangle++(0.04,0.04) (1.029,1.309)rectangle++(0.04,0.04)
  (2.371,3.293)rectangle++(0.04,0.04) (2.288,3.257)rectangle++(0.04,0.04) (0.551,1.807)rectangle++(0.04,0.04) (1.535,2.058)rectangle++(0.04,0.04) (1.254,2.018)rectangle++(0.04,0.04)
  (2.376,2.98)rectangle++(0.04,0.04) (1.302,2.289)rectangle++(0.04,0.04) (0.428,1.318)rectangle++(0.04,0.04) (2.08,2.085)rectangle++(0.04,0.04) (1.952,2.279)rectangle++(0.04,0.04)
  (1.533,1.856)rectangle++(0.04,0.04) (1.987,2.746)rectangle++(0.04,0.04) (2.193,3.078)rectangle++(0.04,0.04) (0.84,1.936)rectangle++(0.04,0.04) (1.671,1.524)rectangle++(0.04,0.04)
  (1.342,0.797)rectangle++(0.04,0.04) (0.005,0.97)rectangle++(0.04,0.04) (1.259,1.623)rectangle++(0.04,0.04) (1.334,1.205)rectangle++(0.04,0.04) (2.214,2.339)rectangle++(0.04,0.04)
  (1.753,2.173)rectangle++(0.04,0.04) (0.328,1.407)rectangle++(0.04,0.04) (2.426,3.15)rectangle++(0.04,0.04) (1.513,2.465)rectangle++(0.04,0.04) (2.023,2.798)rectangle++(0.04,0.04)
  (2.283,2.646)rectangle++(0.04,0.04) (1.631,2.513)rectangle++(0.04,0.04) (1.514,2.282)rectangle++(0.04,0.04) (0.395,1.371)rectangle++(0.04,0.04) (2.403,3.398)rectangle++(0.04,0.04)
  (0.885,2.049)rectangle++(0.04,0.04) (1.328,1.926)rectangle++(0.04,0.04) (1.109,1.234)rectangle++(0.04,0.04) (2.274,3.025)rectangle++(0.04,0.04) (0.488,1.747)rectangle++(0.04,0.04)
  (2.051,2.692)rectangle++(0.04,0.04) (2.008,1.992)rectangle++(0.04,0.04) (1.32,0.817)rectangle++(0.04,0.04) (2.383,3.458)rectangle++(0.04,0.04) (2.096,2.119)rectangle++(0.04,0.04)
  (0.261,1.483)rectangle++(0.04,0.04) (1.698,2.183)rectangle++(0.04,0.04) (1.78,2.544)rectangle++(0.04,0.04) (1.787,2.068)rectangle++(0.04,0.04) (0.711,0.62)rectangle++(0.04,0.04)
  (1.341,1.704)rectangle++(0.04,0.04) (2.047,2.576)rectangle++(0.04,0.04) (1.478,2.338)rectangle++(0.04,0.04) (2.322,3.259)rectangle++(0.04,0.04) (2.019,2.038)rectangle++(0.04,0.04)
  (2.172,3.038)rectangle++(0.04,0.04) (1.107,2.105)rectangle++(0.04,0.04) (1.063,1.745)rectangle++(0.04,0.04) (1,1.669)rectangle++(0.04,0.04) (0.873,0.841)rectangle++(0.04,0.04)
  (0.851,1.464)rectangle++(0.04,0.04) (2.43,3.154)rectangle++(0.04,0.04) (2.392,3.309)rectangle++(0.04,0.04) (0.866,2.028)rectangle++(0.04,0.04) (1.562,1.695)rectangle++(0.04,0.04)
  (0.223,1.243)rectangle++(0.04,0.04) (1.891,1.785)rectangle++(0.04,0.04) (1.73,2.511)rectangle++(0.04,0.04) (1.413,0.803)rectangle++(0.04,0.04) (1.237,1)rectangle++(0.04,0.04)
  (2.383,3.451)rectangle++(0.04,0.04) (1.783,2.235)rectangle++(0.04,0.04) (1.601,1.245)rectangle++(0.04,0.04) (1.547,2.269)rectangle++(0.04,0.04) (1.915,2.67)rectangle++(0.04,0.04)
  (0.197,0.606)rectangle++(0.04,0.04) (2.4,3.378)rectangle++(0.04,0.04) (2.087,2.993)rectangle++(0.04,0.04) (1.978,2.505)rectangle++(0.04,0.04) (0.994,1.341)rectangle++(0.04,0.04)
  (2.242,3.055)rectangle++(0.04,0.04) (2.388,3.285)rectangle++(0.04,0.04) (0.891,0.621)rectangle++(0.04,0.04) (2.236,2.943)rectangle++(0.04,0.04) (1.312,2.264)rectangle++(0.04,0.04)
  (1.995,1.872)rectangle++(0.04,0.04) (2.39,3.405)rectangle++(0.04,0.04) (0.398,0.778)rectangle++(0.04,0.04) (2.331,2.674)rectangle++(0.04,0.04) (2.011,2.292)rectangle++(0.04,0.04)
  (1.849,1.84)rectangle++(0.04,0.04) (2.347,3.106)rectangle++(0.04,0.04) (0.825,1.975)rectangle++(0.04,0.04) (1.949,2.601)rectangle++(0.04,0.04) (0.959,0.665)rectangle++(0.04,0.04)
  (2.049,2.74)rectangle++(0.04,0.04) (1.344,0.876)rectangle++(0.04,0.04) (1.302,1.54)rectangle++(0.04,0.04) (2.234,3.069)rectangle++(0.04,0.04) (1.498,1.252)rectangle++(0.04,0.04)
  (0.419,0.796)rectangle++(0.04,0.04) (0.556,1.072)rectangle++(0.04,0.04) (2.124,2.478)rectangle++(0.04,0.04) (0.358,1.577)rectangle++(0.04,0.04) (1.562,1.136)rectangle++(0.04,0.04)
  (1.998,2.222)rectangle++(0.04,0.04) (2.254,2.751)rectangle++(0.04,0.04) (2.223,3.141)rectangle++(0.04,0.04) (2.078,2.391)rectangle++(0.04,0.04) (0.34,1.544)rectangle++(0.04,0.04)
  (2.326,3.181)rectangle++(0.04,0.04) (1.552,1.598)rectangle++(0.04,0.04) (0.909,1.725)rectangle++(0.04,0.04) (2.138,2.138)rectangle++(0.04,0.04) (1.644,1.19)rectangle++(0.04,0.04)
  (1.782,1.818)rectangle++(0.04,0.04) (1.131,1.017)rectangle++(0.04,0.04) (2.226,2.48)rectangle++(0.04,0.04) (2.312,3.037)rectangle++(0.04,0.04) (1.088,2.129)rectangle++(0.04,0.04)
  (2.107,2.856)rectangle++(0.04,0.04) (0.466,1.663)rectangle++(0.04,0.04) (0.378,0.95)rectangle++(0.04,0.04) (0.948,0.629)rectangle++(0.04,0.04) (1.413,1.788)rectangle++(0.04,0.04)
  (2.378,3.422)rectangle++(0.04,0.04) (0.584,1.87)rectangle++(0.04,0.04) (1.334,1.855)rectangle++(0.04,0.04) (2.428,3.291)rectangle++(0.04,0.04) (0.54,1.851)rectangle++(0.04,0.04)
  (0.784,1.324)rectangle++(0.04,0.04) (1.539,1.639)rectangle++(0.04,0.04) (1.973,2.498)rectangle++(0.04,0.04) (1.72,1.563)rectangle++(0.04,0.04) (1.82,1.952)rectangle++(0.04,0.04)
  (2.195,2.838)rectangle++(0.04,0.04) (2.432,3.296)rectangle++(0.04,0.04) (2.361,3.408)rectangle++(0.04,0.04) (1.934,2.829)rectangle++(0.04,0.04) (2.378,2.968)rectangle++(0.04,0.04)
  (1.067,2.127)rectangle++(0.04,0.04) (1.878,1.823)rectangle++(0.04,0.04) (1.398,2.207)rectangle++(0.04,0.04) (1.48,1.862)rectangle++(0.04,0.04) (2.374,3.408)rectangle++(0.04,0.04)
  (0.891,2.045)rectangle++(0.04,0.04) (1.654,1.786)rectangle++(0.04,0.04) (2,2.328)rectangle++(0.04,0.04) (2.065,2.101)rectangle++(0.04,0.04) (1.32,1.281)rectangle++(0.04,0.04)
  (2.293,3.177)rectangle++(0.04,0.04) (1.909,2.406)rectangle++(0.04,0.04) (1.216,0.867)rectangle++(0.04,0.04) (2.29,3.161)rectangle++(0.04,0.04) (0.524,1.784)rectangle++(0.04,0.04)
  (0.404,0.995)rectangle++(0.04,0.04) (0.201,0.774)rectangle++(0.04,0.04) (1.564,1.959)rectangle++(0.04,0.04) (2.277,2.852)rectangle++(0.04,0.04) (2.236,2.497)rectangle++(0.04,0.04)
  (2.375,3.449)rectangle++(0.04,0.04) (1.159,2.218)rectangle++(0.04,0.04) (2.31,3.21)rectangle++(0.04,0.04) (1.122,2.162)rectangle++(0.04,0.04) (1.527,1.916)rectangle++(0.04,0.04)
  (0.974,1.781)rectangle++(0.04,0.04) (0.335,1.111)rectangle++(0.04,0.04) (0.615,0.752)rectangle++(0.04,0.04) (0.866,1.403)rectangle++(0.04,0.04) (2.024,1.874)rectangle++(0.04,0.04)
  (0.584,1.642)rectangle++(0.04,0.04) (0.355,0.976)rectangle++(0.04,0.04) (0.203,0.756)rectangle++(0.04,0.04) (2.198,2.412)rectangle++(0.04,0.04) (1.758,2.605)rectangle++(0.04,0.04)
  (1.87,2.272)rectangle++(0.04,0.04) (1.988,2.796)rectangle++(0.04,0.04) (2.195,3.081)rectangle++(0.04,0.04) (1.112,2.117)rectangle++(0.04,0.04) (2.03,2.751)rectangle++(0.04,0.04)
  (1.666,2.505)rectangle++(0.04,0.04) (0.78,1.702)rectangle++(0.04,0.04) (1.539,1.148)rectangle++(0.04,0.04) (1.346,2.069)rectangle++(0.04,0.04) (2.189,2.414)rectangle++(0.04,0.04)
  (2.358,2.977)rectangle++(0.04,0.04) (1.904,2.786)rectangle++(0.04,0.04) (1.347,1.71)rectangle++(0.04,0.04) (1.077,1.423)rectangle++(0.04,0.04) (0.22,1.052)rectangle++(0.04,0.04)
  (2.337,2.923)rectangle++(0.04,0.04) (1.79,2.673)rectangle++(0.04,0.04) (2.153,2.607)rectangle++(0.04,0.04) (0.389,1.615)rectangle++(0.04,0.04) (1.061,1.529)rectangle++(0.04,0.04)
  (2.386,3.486)rectangle++(0.04,0.04) (1.165,2.227)rectangle++(0.04,0.04) (2.033,2.763)rectangle++(0.04,0.04) (2.192,3.08)rectangle++(0.04,0.04) (2.198,2.748)rectangle++(0.04,0.04)
  (1.846,1.929)rectangle++(0.04,0.04) (1.453,1.562)rectangle++(0.04,0.04) (0.591,1.455)rectangle++(0.04,0.04) (2.292,3.122)rectangle++(0.04,0.04) (0.863,0.664)rectangle++(0.04,0.04)
  (1.893,2.106)rectangle++(0.04,0.04) (0.288,1.428)rectangle++(0.04,0.04) (2.392,3.324)rectangle++(0.04,0.04) (2.391,3.422)rectangle++(0.04,0.04) (0.892,2.052)rectangle++(0.04,0.04)
  (1.585,1.784)rectangle++(0.04,0.04) (1.424,2.169)rectangle++(0.04,0.04) (1.802,2.523)rectangle++(0.04,0.04) (1.663,2.455)rectangle++(0.04,0.04) (1.669,2.433)rectangle++(0.04,0.04)
  (0.763,1.687)rectangle++(0.04,0.04) (0.715,1.336)rectangle++(0.04,0.04) (0.249,0.937)rectangle++(0.04,0.04) (2.021,2.611)rectangle++(0.04,0.04) (0.862,0.678)rectangle++(0.04,0.04)
  (1.389,0.856)rectangle++(0.04,0.04) (2.407,3.446)rectangle++(0.04,0.04) (1.704,2.629)rectangle++(0.04,0.04) (2.369,3.4)rectangle++(0.04,0.04) (1.823,1.591)rectangle++(0.04,0.04)
  (2.372,2.98)rectangle++(0.04,0.04) (0.461,1.762)rectangle++(0.04,0.04) (1.839,1.54)rectangle++(0.04,0.04) (0.76,1.111)rectangle++(0.04,0.04) (0.582,1.211)rectangle++(0.04,0.04)
  (1.887,2.375)rectangle++(0.04,0.04) (0.916,0.739)rectangle++(0.04,0.04) (1.546,1.212)rectangle++(0.04,0.04) (1.549,2.26)rectangle++(0.04,0.04) (2.009,2.782)rectangle++(0.04,0.04)
  (0.46,1.62)rectangle++(0.04,0.04) (2.335,2.974)rectangle++(0.04,0.04) (1.768,1.48)rectangle++(0.04,0.04) (1.249,1.606)rectangle++(0.04,0.04) (0.949,1.674)rectangle++(0.04,0.04)
  (2.095,2.816)rectangle++(0.04,0.04) (1.667,2.519)rectangle++(0.04,0.04) (2.016,2.809)rectangle++(0.04,0.04) (2.395,3.1)rectangle++(0.04,0.04) (0.491,1.797)rectangle++(0.04,0.04)
  (0.873,0.801)rectangle++(0.04,0.04) (0.076,0.849)rectangle++(0.04,0.04) (0.616,1.011)rectangle++(0.04,0.04) (0.916,1.452)rectangle++(0.04,0.04) (1.689,1.787)rectangle++(0.04,0.04)
  (0.61,1.553)rectangle++(0.04,0.04) (0.326,0.961)rectangle++(0.04,0.04) (2.153,2.291)rectangle++(0.04,0.04) (1.537,1.87)rectangle++(0.04,0.04) (0.276,1.278)rectangle++(0.04,0.04)
  (1.974,2.336)rectangle++(0.04,0.04) (2.379,3.059)rectangle++(0.04,0.04) (1.502,0.922)rectangle++(0.04,0.04) (2.011,2.075)rectangle++(0.04,0.04) (0.942,1.894)rectangle++(0.04,0.04)
  (1.886,1.594)rectangle++(0.04,0.04) (0.656,0.984)rectangle++(0.04,0.04) (1.3,1.494)rectangle++(0.04,0.04) (1.231,0.839)rectangle++(0.04,0.04) (1.795,2.25)rectangle++(0.04,0.04)
  (1.81,1.577)rectangle++(0.04,0.04) (0.545,1.533)rectangle++(0.04,0.04) (2.23,2.466)rectangle++(0.04,0.04) (2.3,2.958)rectangle++(0.04,0.04) (1.471,0.87)rectangle++(0.04,0.04)
  (1.317,2.02)rectangle++(0.04,0.04) (2.124,2.126)rectangle++(0.04,0.04) (2.247,2.944)rectangle++(0.04,0.04) (0.47,1.721)rectangle++(0.04,0.04) (2.374,3.342)rectangle++(0.04,0.04)
  (0.562,1.849)rectangle++(0.04,0.04) (1.333,1.848)rectangle++(0.04,0.04) (1.894,2.385)rectangle++(0.04,0.04) (2.139,2.128)rectangle++(0.04,0.04) (1.861,2.688)rectangle++(0.04,0.04)
  (1.889,1.615)rectangle++(0.04,0.04) (1.25,1.606)rectangle++(0.04,0.04) (2.392,3.124)rectangle++(0.04,0.04) (1.803,2.7)rectangle++(0.04,0.04) (2.379,3.46)rectangle++(0.04,0.04)
  (2.196,3.126)rectangle++(0.04,0.04) (0.528,1.751)rectangle++(0.04,0.04) (0.394,0.988)rectangle++(0.04,0.04) (2.359,3.19)rectangle++(0.04,0.04) (2.388,3.484)rectangle++(0.04,0.04)
  (2.39,3.114)rectangle++(0.04,0.04) (2.188,2.53)rectangle++(0.04,0.04) (1.628,2.49)rectangle++(0.04,0.04) (0.782,1.69)rectangle++(0.04,0.04) (0.339,1.039)rectangle++(0.04,0.04)
  (0.957,1.41)rectangle++(0.04,0.04) (0.624,1.317)rectangle++(0.04,0.04) (0.932,1.196)rectangle++(0.04,0.04) (1.585,1.224)rectangle++(0.04,0.04) (1.116,1.883)rectangle++(0.04,0.04)
  (1.84,1.815)rectangle++(0.04,0.04) (2.428,3.155)rectangle++(0.04,0.04) (1.66,1.875)rectangle++(0.04,0.04) (2.209,2.678)rectangle++(0.04,0.04) (1.607,1.255)rectangle++(0.04,0.04)
  (2.387,3.48)rectangle++(0.04,0.04) (2.388,3.273)rectangle++(0.04,0.04) (2.348,3.372)rectangle++(0.04,0.04) (2.019,2.475)rectangle++(0.04,0.04) (1.99,2.82)rectangle++(0.04,0.04)
  (1.421,1.16)rectangle++(0.04,0.04) (2.228,2.922)rectangle++(0.04,0.04) (1.31,2.259)rectangle++(0.04,0.04) (2.397,3.304)rectangle++(0.04,0.04) (2.366,3.423)rectangle++(0.04,0.04)
  (2.174,2.427)rectangle++(0.04,0.04) (2.414,3.264)rectangle++(0.04,0.04) (2.14,3.054)rectangle++(0.04,0.04) (1.122,1.321)rectangle++(0.04,0.04) (0.911,1.598)rectangle++(0.04,0.04)
  (1.796,2.422)rectangle++(0.04,0.04) (1.281,2.115)rectangle++(0.04,0.04) (2.22,2.546)rectangle++(0.04,0.04) (0.362,1.618)rectangle++(0.04,0.04) (1.632,1.227)rectangle++(0.04,0.04)
  (0.081,1.149)rectangle++(0.04,0.04) (2.284,2.737)rectangle++(0.04,0.04) (2.084,2.616)rectangle++(0.04,0.04) (1.647,2.492)rectangle++(0.04,0.04) (2.149,2.265)rectangle++(0.04,0.04)
  (1.555,1.347)rectangle++(0.04,0.04) (2.02,2.229)rectangle++(0.04,0.04) (1.872,1.585)rectangle++(0.04,0.04) (0.576,0.941)rectangle++(0.04,0.04) (2.246,2.7)rectangle++(0.04,0.04)
  (0.745,1.866)rectangle++(0.04,0.04) (1.404,1.186)rectangle++(0.04,0.04) (2.259,3.118)rectangle++(0.04,0.04) (1.526,2.442)rectangle++(0.04,0.04) (1.726,1.717)rectangle++(0.04,0.04)
  (1.828,2.265)rectangle++(0.04,0.04) (1.531,1.848)rectangle++(0.04,0.04) (2.295,2.858)rectangle++(0.04,0.04) (1.574,1.68)rectangle++(0.04,0.04) (0.217,1.242)rectangle++(0.04,0.04)
  (1.599,2.002)rectangle++(0.04,0.04) (1.192,1.017)rectangle++(0.04,0.04) (1.919,2.411)rectangle++(0.04,0.04) (1.47,2.299)rectangle++(0.04,0.04) (1.489,1.346)rectangle++(0.04,0.04)
  (2.048,2.809)rectangle++(0.04,0.04) (1.313,2.215)rectangle++(0.04,0.04) (2.394,3.466)rectangle++(0.04,0.04) (1.832,2.739)rectangle++(0.04,0.04) (0.811,1.787)rectangle++(0.04,0.04)
  (1.679,2.271)rectangle++(0.04,0.04);

        \node[etiq, firebrick3, right] at (0.15,0.4) {$w_1(A)$};
        \node[etiq, steelblue4, right] at (2.7,6) {$w_2(A)$};
        \node[etiq, coveraccent!80!black, left] at (-2.4,2.2) {$w_3(A)$};
        \node[etiq, vertpiece, right] at (2.6,1.7) {$w_4(A)$};
      \end{scope}
    \end{tikzpicture}}
  \small Vingt-quatre coefficients ; les folioles du bas sont elles-mêmes
  de petites fougères.
\end{frame}

\begin{frame}{Remarque 3.7 --- La compression fractale d'images}
  \begin{Rem}{Remarque 3.7}
    \begin{itemize}
      \item La fougère : vingt-quatre nombres, contre des centaines de
        milliers pour une image matricielle.
      \item Fin des années 1980 : Barnsley et son étudiant Arnaud Jacquin
        en tirent la compression fractale. L'image est découpée en petits
        blocs ; pour chacun, on cherche une application affine qui
        l'obtient à partir d'un bloc plus grand de la même image (principe
        du collage).
      \item Utilisée commercialement, notamment dans l'encyclopédie
        Encarta de Microsoft, avant d'être supplantée par JPEG.
    \end{itemize}
  \end{Rem}
\end{frame}
